% BEGIN EMBEDDED BIBLIOGRAPHY
% Bibliography data are part of this source; the compiler writes its temporary input.
\begin{filecontents*}[overwrite]{\jobname.bib}
@article{Wilson1974,
  author = {Wilson, Kenneth G.},
  title = {Confinement of quarks},
  journal = {Physical Review D},
  volume = {10},
  pages = {2445--2459},
  date = {1974},
  doi = {10.1103/PhysRevD.10.2445}
}

@article{OS1973,
  author = {Osterwalder, Konrad and Schrader, Robert},
  title = {Axioms for {Euclidean} {Green}'s functions},
  journal = {Communications in Mathematical Physics},
  volume = {31},
  number = {2},
  pages = {83--112},
  date = {1973},
  doi = {10.1007/BF01645738}
}

@article{B-CMP109,
  author = {Ba{\l}aban, Tadeusz},
  title = {Renormalization group approach to lattice gauge field theories. I. Generation of effective actions in a small field approximation and a coupling constant renormalization in four dimensions},
  journal = {Communications in Mathematical Physics},
  volume = {109},
  number = {2},
  pages = {249--301},
  date = {1987},
  doi = {10.1007/BF01215223}
}

@article{B-CMP119,
  author = {Ba{\l}aban, Tadeusz},
  title = {Convergent renormalization expansions for lattice gauge theories},
  journal = {Communications in Mathematical Physics},
  volume = {119},
  number = {2},
  pages = {243--285},
  date = {1988},
  doi = {10.1007/BF01217741}
}

@article{B-CMP95,
  author = {Ba{\l}aban, Tadeusz},
  title = {Propagators and renormalization transformations for lattice gauge theories. {I}},
  journal = {Communications in Mathematical Physics},
  volume = {95},
  number = {1},
  pages = {17--40},
  date = {1984},
  doi = {10.1007/BF01215753}
}

@article{B-CMP98,
  author = {Ba{\l}aban, Tadeusz},
  title = {Averaging operations for lattice gauge theories},
  journal = {Communications in Mathematical Physics},
  volume = {98},
  number = {1},
  pages = {17--51},
  date = {1985},
  doi = {10.1007/BF01211042}
}

@article{B-CMP99a,
  author = {Ba{\l}aban, Tadeusz},
  title = {Spaces of regular gauge field configurations on a lattice and gauge fixing conditions},
  journal = {Communications in Mathematical Physics},
  volume = {99},
  number = {1},
  pages = {75--102},
  date = {1985},
  doi = {10.1007/BF01466594}
}

@article{B-CMP99b,
  author = {Ba{\l}aban, Tadeusz},
  title = {Propagators for lattice gauge theories in a background field},
  journal = {Communications in Mathematical Physics},
  volume = {99},
  number = {3},
  pages = {389--434},
  date = {1985},
  doi = {10.1007/BF01240355}
}

@online{Dyatlov-155,
  author={Dyatlov, Semyon},
  title={Lecture notes for 18.155: distributions, elliptic regularity, and applications to {PDEs}},
  date={2026-10-20},
  url={https://math.mit.edu/~dyatlov/18.155/155-notes.pdf},
  urldate={2026-10-21}
}
@online{MIT-AnalysisII-2005,
  author={{MIT OpenCourseWare}},
  title={Analysis {II}: Lecture 36},
  note={Student notes, course taught by Victor Guillemin},
  date={2005},
  url={https://ocw.mit.edu/courses/18-101-analysis-ii-fall-2005/f738081bb28c1f7434cd9241c55a8ef7_lecture36.pdf},
  urldate={2026-10-21}
}
@online{MIT-LieGroups-2024,
  author={{MIT OpenCourseWare}},
  title={Lie groups and Lie algebras {II}: Lecture 11},
  note={Course taught by Pavel Etingof},
  date={2024},
  url={https://ocw.mit.edu/courses/18-755-lie-groups-and-lie-algebras-ii-spring-2024/mit18_755_s24_lec11.pdf},
  urldate={2026-10-21}
}
@article{FernandezProcacci2007,
  author = {Fern{\'a}ndez, Roberto and Procacci, Aldo},
  title = {Cluster expansion for abstract polymer models. New bounds from an old approach},
  journal = {Communications in Mathematical Physics},
  volume = {274},
  pages = {123--140},
  date = {2007},
  doi = {10.1007/s00220-007-0279-2}
}

@misc{Tserunyan2026,
  author = {Tserunyan, Anush},
  title = {Introduction to Descriptive Set Theory},
  date = {2026-10-19},
  note = {Lecture notes, version of 19 October 2026},
  url = {https://www.math.mcgill.ca/atserunyan/Teaching_notes/dst_lectures.pdf}
}

@article{B-CMP102a,
  author = {Ba{\l}aban, Tadeusz},
  title = {Ultraviolet stability of three-dimensional lattice pure gauge field theories},
  journal = {Communications in Mathematical Physics},
  volume = {102},
  pages = {255--275},
  date = {1985},
  doi = {10.1007/BF01229380}
}
@article{King1986I,
  author = {King, Chris},
  title = {The {U(1)} {Higgs} model. {I}. The continuum limit},
  journal = {Communications in Mathematical Physics},
  volume = {102},
  pages = {649--677},
  date = {1986},
  doi = {10.1007/BF01221651}
}
@article{King1986II,
  author = {King, Chris},
  title = {The {U(1)} {Higgs} model. {II}. The infinite volume limit},
  journal = {Communications in Mathematical Physics},
  volume = {103},
  pages = {323--349},
  date = {1986},
  doi = {10.1007/BF01206942}
}
@article{OS1975,
  author = {Osterwalder, Konrad and Schrader, Robert},
  title = {Axioms for {Euclidean} {Green}'s functions. {II}},
  journal = {Communications in Mathematical Physics},
  volume = {42},
  pages = {281--305},
  date = {1975},
  doi = {10.1007/BF01608978}
}
@online{QED3FieldStrength,
  shorthand = {QEM1},
  title = {The continuum field-strength construction for {QED} in three dimensions},
  note = {Preprint, version at commit 732255e},
  url = {https://github.com/ethanwsussman/Yang--Mills-LLM-experiment/blob/732255ef81e87d982739a18ab0699610cee4b4e4/QED3/QEM1.pdf},
  urldate = {2026-10-22}
}
@online{ScalarQED-YM4Conditional,
  shorthand = {YM4-cond},
  title = {Pure {SU(2)} {Yang--Mills} theory in four dimensions: separated-point {Osterwalder--Schrader} axioms and {Euclidean} invariance under ultraviolet, infrared and complex-rotation hypotheses},
  note = {Preprint, version at commit cbedf8b},
  url = {https://github.com/marikgoldstein/scalar-qed/blob/cbedf8b/ym/su2_4d_full/ym_su2_4d_full.tex},
  urldate = {2026-10-22}
}
@online{ScalarQED-YM4Blocks,
  shorthand = {YM4-blocks},
  title = {{Osterwalder--Schrader} type properties of bounded block observables of four-dimensional {SU(2)} lattice {Yang--Mills} theory},
  note = {Preprint, version at commit cbedf8b},
  url = {https://github.com/marikgoldstein/scalar-qed/blob/cbedf8b/ym/su2_4d_blocks/su2_4d_blocks.tex},
  urldate = {2026-10-22}
}
@online{ScalarQED-Main,
  shorthand = {SQED},
  title = {{Osterwalder--Schrader} axioms for three-dimensional scalar quantum electrodynamics at fixed photon mass},
  note = {Preprint, version at commit 29c4ac8},
  url = {https://github.com/marikgoldstein/scalar-qed/blob/29c4ac8/main_paper/main_paper_latest.pdf},
  urldate = {2026-10-22}
}
@online{ScalarQED-Rotations,
  shorthand = {SQED-rot},
  title = {Rotation invariance of the three-dimensional abelian {Higgs} model: the curvature sector},
  note = {Preprint, version at commit f1c37fc},
  url = {https://github.com/marikgoldstein/scalar-qed/blob/f1c37fc/main_paper/supp_material/rotations_proofs.pdf},
  urldate = {2026-10-22}
}

\end{filecontents*}
% END EMBEDDED BIBLIOGRAPHY
\documentclass[11pt,letterpaper,reqno]{amsart}
% BEGIN INLINE manuscript-style
\makeatletter
\NeedsTeXFormat{LaTeX2e}[2020/10/01]
% Use with \documentclass[11pt,letterpaper,reqno]{amsart}.
% This file controls presentation only. It supplies no mathematical content.

\RequirePackage[T1]{fontenc}
\RequirePackage{lmodern}
\RequirePackage{microtype}
\RequirePackage[letterpaper,margin=1in,headheight=15pt]{geometry}
\RequirePackage{amsmath,amsthm,amssymb,mathtools}
\RequirePackage{mathrsfs}
\RequirePackage{enumitem}
\RequirePackage{perpage}
\RequirePackage{xcolor}
\definecolor{manuscriptred}{rgb}{0.6,0.1,0.1}
\definecolor{manuscriptblue}{rgb}{0.1,0.3,0.5}
\definecolor{manuscriptmagenta}{HTML}{60445C}
\RequirePackage[breaklinks=true,colorlinks=true,
  linkcolor=manuscriptred,citecolor=manuscriptblue,urlcolor=manuscriptmagenta]{hyperref}
\RequirePackage{cleveref}
\hypersetup{pdfauthor={}}

% Retain the AMS paragraph, title, heading and running-head conventions.
% Equations use one global Arabic sequence.
\setcounter{tocdepth}{1}
\renewcommand{\thefigure}{\Roman{figure}}
\renewcommand{\thetable}{\Roman{table}}
\MakePerPage{footnote}
\renewcommand{\thefootnote}{%
  \ifnum\value{footnote}>9\arabic{footnote}\else\fnsymbol{footnote}\fi}

% Upright statements. Main results use letters; supporting statements
% share a section-local sequence.
\theoremstyle{definition}
\newtheorem{ms@theorem}{Theorem}
\ExplSyntaxOn
\cs_new:Npn \ymtheoremnumber { \int_to_Alph:n { \value{ms@theorem} } }
\ExplSyntaxOff
\renewcommand{\thems@theorem}{\ymtheoremnumber}
\newtheorem*{ms@theoremstar}{Theorem}
\newtheorem{ms@theoremzero}{Theorem}
\renewcommand{\thems@theoremzero}{0}
\newtheorem{ms@proposition}{Proposition}[section]
\newtheorem{ms@lemma}[ms@proposition]{Lemma}
\newtheorem{ms@definition}[ms@proposition]{Definition}
\newtheorem{ms@corollary}[ms@proposition]{Corollary}
\newtheorem*{ms@corollarystar}{Corollary}

\newcommand{\ms@majorstart}{%
  \pushQED{\qed}\renewcommand{\qedsymbol}{$\blacksquare\square$}}
\newcommand{\ms@minorstart}{%
  \pushQED{\qed}\renewcommand{\qedsymbol}{$\blacksquare$}}
\newenvironment{theorem}[1][]{\ms@majorstart\ms@theorem[#1]}
  {\popQED\endms@theorem}
\newenvironment{theoremzero}[1][]{\ms@majorstart\ms@theoremzero[#1]}
  {\popQED\endms@theoremzero}
\newenvironment{theorem*}[1][]{\ms@majorstart\ms@theoremstar[#1]}
  {\popQED\endms@theoremstar}
\newenvironment{proposition}[1][]{\ms@minorstart\ms@proposition[#1]}
  {\popQED\endms@proposition}
\newenvironment{lemma}[1][]{\ms@minorstart\ms@lemma[#1]}
  {\popQED\endms@lemma}
\newenvironment{definition}[1][]{\ms@minorstart\ms@definition[#1]}
  {\popQED\endms@definition}
\newenvironment{corollary}[1][]{\ms@majorstart\ms@corollary[#1]}
  {\popQED\endms@corollary}
\newenvironment{corollary*}[1][]{\ms@majorstart\ms@corollarystar[#1]}
  {\popQED\endms@corollarystar}

\theoremstyle{remark}
\newtheorem{remark}[ms@proposition]{Remark}
\newtheorem*{remark*}{Remark}
% The ordinary amsthm proof environment retains its open square.

\crefname{ms@theoremzero}{theorem}{theorems}
\Crefname{ms@theoremzero}{Theorem}{Theorems}
\crefname{ms@theorem}{theorem}{theorems}
\Crefname{ms@theorem}{Theorem}{Theorems}
\crefname{ms@proposition}{proposition}{propositions}
\Crefname{ms@proposition}{Proposition}{Propositions}
\crefname{ms@lemma}{lemma}{lemmas}
\Crefname{ms@lemma}{Lemma}{Lemmas}
\crefname{ms@definition}{definition}{definitions}
\Crefname{ms@definition}{Definition}{Definitions}
\crefname{ms@corollary}{corollary}{corollaries}
\Crefname{ms@corollary}{Corollary}{Corollaries}
\crefname{remark}{remark}{remarks}
\Crefname{remark}{Remark}{Remarks}

% Bibliography configuration for this manuscript.
  \RequirePackage[backend=biber,style=alphabetic,backref=true,
    giveninits=true,isbn=false,date=year]{biblatex}
  \addbibresource{\jobname.bib}
  \setcounter{biburlnumpenalty}{100}
  \setcounter{biburllcpenalty}{100}
  \setcounter{biburlucpenalty}{100}
  \renewcommand*{\bibfont}{\footnotesize}
  \renewbibmacro{in:}{}
  \DeclareFieldFormat{pages}{#1}
  % Use the public alphabetic style. No private author-hash overrides;
  % repeated author names are printed in full.
% Alternative labels for a single displayed identity.
\newcommand{\equationalias}[1]{\ltx@label{{#1}}}
\makeatother
% END INLINE manuscript-style
\DeclareMathOperator{\tr}{tr}
\DeclareMathOperator{\Lip}{Lip}
\title[Scalar fields from lattice Yang--Mills theory]{Scalar fields from three-dimensional lattice Yang--Mills theory}
\date{}
\hypersetup{pdftitle={Scalar fields from three-dimensional lattice Yang-Mills theory},pdfauthor={}}
\begin{document}
\begin{abstract}
We construct continuum correlations of the centered scalar plaquette-energy
observable for the three-dimensional \(\mathrm{SU}(2)\) Wilson measure.
Its formal classical counterpart is the Yang--Mills action density
\(\operatorname{tr}(F^2)\).
For a fixed admissible scale factor \(L\), on every torus
\(T_{2L^m}\), \(m\in\mathbb Z\), connected correlations of order
at least two converge along the full cutoff sequence at every positive
physical inverse coupling, provided the test supports have empty
common intersection. For canonical first-passage data at sufficiently
small reference coupling and tending cutoffs, limits through periodic
Gibbs states exist subsequentially and admit canonical clustering
components whose Schwinger functionals satisfy the
Osterwalder--Schrader axioms on Schwartz tests flat at the coincidence
diagonals. A point-insertion bound reconstructs a Hermitian scalar
Wightman field in each component. Under a further restriction on the
reference coupling, a uniform short-distance two-point law determines
the physical inverse coupling and proves nontriviality.
The proof combines a local complex-source expansion, reflection and
chessboard estimates, comparison of consecutive cutoffs, and analytic
continuation of the separated Euclidean kernels.
\end{abstract}
\maketitle
\tableofcontents

\section{Introduction}\label{sec:introduction}

We study a scalar observable whose formal classical counterpart is
\(\operatorname{tr}(F^2)\), the Yang--Mills action density.
Its quantum correlations are defined through centered Wilson
plaquette energies.

Let \(L\ge13\) be odd and \(a_K=L^{-K}\) for \(K\in\mathbb Z\).
On \(a_K\mathbb Z^3\), a configuration assigns
\(U(b)\in\mathrm{SU}(2)\) to every oriented
bond, with \(U(\bar b)=U(b)^{-1}\). For a positive plaquette \(p\),
let \(x(p)\) be its barycentre and put
\[
  a_p(U)=1-\tfrac12\operatorname{tr}\partial U(p).
\]
Here \(\partial U(p)\) is the ordered product around \(p\), and
\(a_p\) is the gauge-invariant scalar plaquette energy.
The Wilson action is normalized as in \cite[\S III]{Wilson1974}.
On the infinite lattice $a_K\mathbb Z^3$, a periodic Gibbs
state is a probability law invariant under a finite-index
subgroup of lattice translations whose conditional law on
each finite set of bonds has density proportional to the
Wilson factor for the plaquettes meeting that set.  The
boundary bonds are held fixed.  Such states exist at every
$K$ and $\beta>0$.  For a compactly supported test $f$, put
\[
 \mathcal O_K(f)=a_K^{-1}\sum_p f(x(p))a_p,
 \qquad \mathcal O_K^\nu(f)=\mathcal O_K(f)-\nu(\mathcal O_K(f)).
\]
The sum runs over positive plaquettes in all three orientations.
We call $\mathcal O_K^\nu$ the smeared, centered scalar
plaquette-energy density. The plaquette definition fixes the observable
and its normalization;
\eqref{eq:action-field-strength-splitting} gives the exact lattice relation.
The notation \(\operatorname{tr}(F^2)\) does not denote the pointwise
square of a separately constructed continuum field strength.

For $n\ge1$, let $\mathcal S_{\ne}(\mathbb R^{3n})$ consist
of complex Schwartz functions whose derivatives of every order
vanish whenever two arguments coincide.  Let
$\mathcal S_+(\mathbb R^{3n})$ be its subspace of functions
also flat outside $\{x_{13}>0,\ldots,x_{n3}>0\}$, and set
$\mathcal S_{\ne}(\mathbb R^0)=\mathcal S_+(\mathbb R^0)=\mathbb C$.
We use the seminorms
\[
 |\varphi|_k=\sup_{x,\,|\alpha|\le k}
       (1+|x|^2)^{k/2}|\partial^\alpha\varphi(x)|.
\]
For a Euclidean isometry $h$ and a vector $v$, put
\[
 \begin{split}
 (h\cdot\varphi)(x_1,\ldots,x_n)
    &=\varphi(h^{-1}x_1,\ldots,h^{-1}x_n),\\
 (\tau_v\varphi)(x_1,\ldots,x_n)
    &=\varphi(x_1-v,\ldots,x_n-v).
 \end{split}
\]
Write $\theta(x_1,x_2,x_3)=(x_1,x_2,-x_3)$ and
$(\Theta\varphi)(x_1,\ldots,x_n)
=\varphi(\theta x_1,\ldots,\theta x_n)$.
These actions are the identity in degree zero.
The following theorem states the Euclidean axioms of
Osterwalder and Schrader \cite[\S\S 2--3]{OS1973},
with an explicit growth bound.

\begin{theoremzero}[Osterwalder--Schrader theorem for the plaquette-energy density]
\label{thm:scalar-os}
There exist an odd integer $L\ge13$, a physical inverse coupling
$b>0$, integers $K_j\to\infty$, and periodic Gibbs states
$\rho_j,\nu_j$ at spacing $a_{K_j}=L^{-K_j}$ and inverse
coupling $\beta_j=L^{K_j}b$ with the following properties.
For every tuple of smooth compactly supported functions
$f_1,\ldots,f_n$ with pairwise disjoint supports, the limit
\[
 \mathsf S_n(f_1\otimes\cdots\otimes f_n)
   =\lim_{j\to\infty}\rho_j\!\left(
          \prod_{i=1}^n\mathcal O_{K_j}^{\nu_j}(f_i)\right)
\]
exists.  These moments extend uniquely to continuous linear
functionals on $\mathcal S_{\ne}(\mathbb R^{3n})$, with
$\mathsf S_0(z)=z$ for $z\in\mathbb C$.  They have the
following properties.
\begin{enumerate}
\item[\textup{(E0$'$)}] \emph{Regularity.}  There is a finite
$C>0$ such that, for every $n\ge1$ and
$\varphi\in\mathcal S_{\ne}(\mathbb R^{3n})$,
\[
 |\mathsf S_n(\varphi)|\le2C^n(n!)^8|\varphi|_{8n}.
\]
Moreover,
$\mathsf S_n(\bar\varphi)=\overline{\mathsf S_n(\varphi)}$.
\item[\textup{(E1)}] \emph{Euclidean invariance.}  For every
Euclidean isometry $h$,
$\mathsf S_n(h\cdot\varphi)=\mathsf S_n(\varphi)$.
\item[\textup{(E2)}] \emph{Reflection positivity.}  For every
$N\ge0$ and $\varphi_n\in\mathcal S_+(\mathbb R^{3n})$,
$0\le n\le N$,
\[
 \sum_{n,n'=0}^{N}
 \mathsf S_{n+n'}(\Theta\bar\varphi_n\otimes\varphi_{n'})\ge0.
\]
\item[\textup{(E3)}] \emph{Symmetry.}  The functional
$\mathsf S_n$ is unchanged by permutations of its $n$ arguments.
\item[\textup{(E4)}] \emph{Clustering.}  For $n,n'\ge0$,
$\varphi\in\mathcal S_+(\mathbb R^{3n})$,
$\psi\in\mathcal S_+(\mathbb R^{3n'})$ and
$0\ne v\in\mathbb R^2\times\{0\}$,
\[
 \lim_{s\to\infty}
 \mathsf S_{n+n'}(\Theta\bar\varphi\otimes\tau_{sv}\psi)
   =\mathsf S_n(\Theta\bar\varphi)\mathsf S_{n'}(\psi).
\]
\end{enumerate}
On their reflection Hilbert space these functionals reconstruct
a Hermitian scalar Wightman field on $\mathbb R^{1+2}$ with
positive energy, a unique invariant vacuum, a dense common
Poincar\'e-invariant polynomial domain, covariance, local
commutativity and tempered vacuum distributions.  Its Euclidean
continuations recover the separated moments and the functionals
$\mathsf S_n$ on the stated spaces.  The Hilbert space has
dimension at least two.
\end{theoremzero}

\begin{proof}
Fix a positive source radius and choose the reference coupling
$g>0$ below the ceiling in \cref{thm:short-distance-law}.
Choose $b$ in the nonempty interval $I(g)$ of
\cref{thm:gibbs-continuum} and take a volume-first continuum
limit at this label.  Select one of its canonical clustering
components.  Its integrating states $\rho_j$ are obtained by
invariant conditioning of the original states $\nu_j$, which
remain in the centering.  The component construction and the
Euclidean properties are \cref{IV-26-41,IV-28-48,IV-78-15}.
The point-insertion bounds supplied there give the reconstruction
by \cref{IV-76-59}; nontriviality is \cref{IV-69-39}.
\end{proof}

The Euclidean assertions use the flat Schwartz spaces specified
above and leave coincident-point contact terms unspecified.
The following theorems give the lattice limits, parameter ranges
and short-distance normalization.

For \(m\in\mathbb Z\), \(K+m\ge1\), and \(\ell=2L^m\),
the torus lattice is
\[
 \Lambda_{K,\ell}=L^{-K}\mathbb Z^3/\ell\mathbb Z^3
 \subset T_\ell=(\mathbb R/\ell\mathbb Z)^3.
\]
On this finite lattice, put
\(A(c,U)=\sum_p c(x(p))a_p(U)\) and use normalized product
Haar measure \(dU\). A complex source \(w\) changes the
partition function to
\begin{equation}\label{III-16-25}
 Z_{K,\ell,\beta}[w]
   =\int \exp[-A(\beta-L^Kw,U)]\,dU .
\end{equation}

\begin{theorem}[The continuum limit on a fixed torus]
\label{thm:fixed-torus-continuum}
There is an odd integer $L\ge13$ for which the following holds.
Let $m\in\mathbb Z$, $\ell=2L^m$, $b>0$ and $n\ge2$.
For real Lipschitz functions $f_1,\ldots,f_n$ on $T_\ell$
whose supports have empty common intersection, the connected
correlation
\[
 S_n^{(\ell,b)}(f_1,\ldots,f_n)
 =\lim_{K\to\infty}
 \left.\partial_{z_1}\cdots\partial_{z_n}
 \log\frac{Z_{K,\ell,L^Kb}[\sum_i z_if_i]}
               {Z_{K,\ell,L^Kb}[0]}\right|_{z=0}
\]
exists along the full sequence of sufficiently fine lattices.
Convergence is locally uniform in $b>0$, with the summable
cutoff rate of \eqref{IV-63-26} after a fixed dilation.
The limit is smooth in $b$, and
\[
 \partial_b^\nu S_n^{(\ell,b)}(f_1,\ldots,f_n)
 =(-1)^\nu S_{n+\nu}^{(\ell,b)}
                      (\underbrace{1,\ldots,1}_{\nu},f_1,\ldots,f_n).
\]
For every $s\in\mathbb Z$, with $f^{(s)}(x)=f(L^{-s}x)$,
\[
 S_n^{(\ell,b)}(f_1,\ldots,f_n)
 =L^{sn}S_n^{(L^s\ell,L^sb)}(f_1^{(s)},\ldots,f_n^{(s)}).
\]
If $K_j\to\infty$ and $L^{-K_j}\beta_j\to b>0$, the
same limit is obtained along $(K_j,\beta_j)$.
For pairwise disjoint test supports these correlations are
the restrictions of unique distributions on the space of
distinct points of $T_\ell^n$. Different positive labels $b$
give different limiting families on the same torus.
\end{theorem}


The proof is given by Theorem~\ref{IV-63-9} and
Corollary~\ref{cor:iv-all-sides-cutoff}; the distributional
assertion follows from Section~\ref{sec:distribution-extension}
and finite dilation. Corollary~\ref{cor:positive-torus-labels}
proves the distinction of positive labels. The powers of the side in the cutoff
estimate are retained throughout that proof.


The lattice conventions and block average are specified in
\cref{sec:finite-lattice,sec:block-averaging,sec:nonabelian-mean}. All distances on \(T_\ell\) in the source norm are
physical distances. For \(M=L^{m'}\), define
\begin{equation}\label{III-16-17}
 \|w\|_{\flat;T_\ell}
 =\max\left\{\sup_{T_\ell}|w|,
       \frac{40M}{\sqrt3}\Lip_\infty(w;T_\ell)\right\},
 \qquad
 \Lip_\infty(w;T_\ell)
 =\sup_{x\ne y}\frac{|w(x)-w(y)|}{|x-y|_\infty}.
\end{equation}
Holomorphy on a ball of sources means holomorphy on every
finite-dimensional affine slice of that ball.

At a fixed admissible threshold, we use the coefficient-determined
real coupling recursion of Theorem~\ref{II-203-9}, with its canonical
auxiliary data. Write it as
\[
 x_0(\beta)=\beta,\qquad
 x_j(\beta)=L^{-1}x_{j-1}(\beta)-b_j,\qquad |b_j|\le B_\sharp.
\]
The increments in this recursion are computed at zero source.
Set \(\bar B=LB_\sharp/(L-1)\) and, for \(g>0\), set
\(X_g=L(g^{-2}+B_\sharp)\). The first-passage index at \(g\) is
the least index \(K\ge0\) for which \(x_K(\beta)\le X_g\).
We call \((K,\beta)\) a first-passage datum when this index is positive.



A separated tuple consists of real $C_c^1(\mathbb R^3)$
functions with pairwise disjoint supports.  Its limiting
moment is denoted by $\sigma(f_1,\ldots,f_n)$, and
$\sigma(\varnothing)=1$.

\begin{theorem}[Infinite-volume limits and reconstruction]
\label{thm:gibbs-continuum}
The canonical construction constants can be fixed, with
$L$ as above, so that for every $\varrho>0$ there is
$\gamma_\varrho>0$ with the following properties.  Fix
$0<g\le\gamma_\varrho$, using this one assignment of the
coupling recursion throughout.

Let $(K_j,\beta_j)$ be first-passage data at $g$, with
$K_j\to\infty$, and let $\nu_j$ be any periodic Gibbs
states at these data.  There is a common subsequence on
which
\[
 \sigma(f_1,\ldots,f_n)
   =\lim_j\nu_j\!\left(\prod_{i=1}^n
                              \mathcal O_{K_j}^{\nu_j}(f_i)\right)
\]
exists for every separated tuple, and on which the positive
labels $a_{K_j}\beta_j$ converge.  The limit has zero
one-point function, is invariant under all Euclidean
isometries, and is reflection positive at every plane.
Its reflection Hilbert space is separable and its joint
energy--momentum spectrum lies in $\{(E,q):E\ge|q|\}$.

Every such limit has a canonical component measure
$(\mathsf X,\lambda,(\sigma_\xi)_{\xi\in\mathsf X})$ as
in Theorem~\ref{IV-35-36}: $\mathsf X$ is compact
metrizable, $\lambda$ has full support,
$\xi\mapsto\sigma_\xi(X)$ is continuous for each tuple,
and $\sigma(X)=\int\sigma_\xi(X)\,d\lambda(\xi)$.
The components are limits of the same lattice sequence
after invariant conditioning, with the original laws
$\nu_j$ retained for centering.  Each component is
Euclidean invariant and reflection positive, has
$\sigma_\xi(f)=c_\xi\int f$ with a common bound on
$|c_\xi|$, and satisfies
\[
 |\sigma_\xi(X\cdot\tau_vY)-\sigma_\xi(X)\sigma_\xi(Y)|
       \le C_g(X,Y)|v|^{-1}
       \qquad(|v|\ge r(X,Y)),
\]
uniformly in $\xi$.  The measure is unique in the precise
sense of Proposition~\ref{IV-36-9}.

The component moments extend uniquely to Schwinger
functionals on $\mathcal S_{\ne}$ satisfying
\[
 |\mathsf S_n^{\sigma_\xi}(\varphi)|
          \le2(C_E^g)^n(n!)^8|\varphi|_{8n},
\]
with the common constant in \eqref{eq:gibbs-schwinger-growth}
and the Schwartz seminorm \eqref{IV-26-4}.  They obey
Euclidean invariance, symmetry, reflection positivity and
Schwartz clustering in the exact test-space formulation
of Proposition~\ref{IV-28-48}.  On each component's
reflection Hilbert space they reconstruct one Hermitian
scalar Wightman field on $\mathbb R^{1+2}$, with positive
energy, a unique invariant vacuum, a dense common
Poincar\'e-invariant polynomial domain, covariance,
local commutativity and tempered vacuum distributions.
Its Euclidean continuations recover all the separated
moments and the flat-Schwartz Schwinger functionals.

The original limit also has this reconstruction, except
that vacuum uniqueness and clustering are not asserted.
Its zero-energy space is $L^2(\mathsf X,\lambda)$, with
the vacuum vector corresponding to the constant $1$.
All growth and insertion constants are independent of
the sequence and of its components.

The interval
$I(g)=(g^{-2}+\bar B,X_g-\bar B]$ is nonempty.
For every $b\in I(g)$ there are such limits with
$a_{K_j}\beta_j=b$ at every index, obtained by taking
the volume limit of cubic torus measures first.
They are also limits along a suitable sequence of
growing tori with the same $(K_j,\beta_j)$.
\end{theorem}

The infinite-volume conclusion is subsequential.  The
fixed-torus estimate retains a power of the volume and
does not give uniqueness of these infinite-volume
limits.  Full rotation invariance for an arbitrary
growing-torus limit has the additional criterion in
Section~\ref{sec:growing-torus-rotation}.

The observable throughout is the scalar plaquette-energy
density defined above.  The reconstructed field is a
scalar field.  The construction makes no identification
with a tensor-valued field strength and asserts no mass gap.

Theorem~\ref{IV-78-15} proves the component and reconstruction
assertions.  Existence and the comparison of limiting orders
are proved in Section~\ref{sec:continuum-subsequences}.

For $0<\vartheta<1$, let $\mathfrak T_E(d;\vartheta)$ consist
of pairs of nonzero nonnegative Lipschitz functions $(f,h)$
on $E=\mathbb R^3$ or $T_\ell$ such that
\begin{gather*}
 \operatorname{supp}f\subset\overline B(x,d),\qquad
 \operatorname{supp}h\subset\overline B(y,d),\qquad
 d\le\operatorname{dist}(\operatorname{supp}f,
                         \operatorname{supp}h)\le6d,\\
 \int f\ge\vartheta\|f\|_\infty d^3,\qquad
 \int h\ge\vartheta\|h\|_\infty d^3,\qquad
 \operatorname{Lip}f\le\|f\|_\infty/(\vartheta d),\qquad
 \operatorname{Lip}h\le\|h\|_\infty/(\vartheta d).
\end{gather*}
On a torus also require $20d\le\ell$. The two supports then
lie in one ball of radius $5d$, and the kernel integral
\[
 \mathcal K_{\rm free}(f,h)
       =\iint \frac{f(x)h(y)}{|x-y|^6}\,dx\,dy
\]
uses their common Euclidean lift in the torus case.
For a moment functional $\sigma$, write
$\kappa_2^\sigma(f,h)=\sigma(f,h)-\sigma(f)\sigma(h)$.

\begin{theorem}[The short-distance two-point law]
\label{thm:short-distance-law}
There is a further ceiling $0<\gamma_W\le\gamma_\varrho$,
independent of $\vartheta$, with the following property.
Fix $0<g\le\gamma_W$ and $0<\vartheta<1$.
There are $d_0=d_0(g,\vartheta)>0$ and a nondecreasing
function $\varepsilon_{g,\vartheta}:(0,d_0]\to[0,1/3]$
with $\varepsilon_{g,\vartheta}(d)\to0$ as $d\downarrow0$
such that every limit $\sigma$ of
Theorem~\ref{thm:gibbs-continuum} and every one of its
canonical components obeys
\[
 \kappa_2^\sigma(f,h)
   =\frac9{16\pi^2}\,b^{-2}\mathcal K_{\rm free}(f,h)
      (1+e_{\sigma,f,h}),\qquad
 |e_{\sigma,f,h}|\le\varepsilon_{g,\vartheta}(d),
\]
for every $0<d\le d_0$ and every $C^1$ pair
$(f,h)\in\mathfrak T_{\mathbb R^3}(d;\vartheta)$.
Here $b$ is the limiting label of the defining lattice
sequence; the components have the same label as their
parent. The error bound is uniform in the limit,
component, test pair and its position.

The label is uniquely determined by the connected
two-point function. Along every defining sequence at $g$
the labels converge to this same value. The connected
two-point function is positive on each indicated pair,
and each reflection Hilbert space has dimension at least
two.
\end{theorem}


Theorem~\ref{IV-69-29} proves the law and label recovery;
Corollary~\ref{IV-69-39} gives their component versions and
nontriviality. The finite-cutoff estimate and its torus and
growing-torus limits are proved in
Section~\ref{sec:two-point-transfer}.

\subsection{The source estimates behind the construction}

The finite-lattice construction begins in
\cref{sec:finite-renormalisation-overview}.
Regular critical backgrounds, their Gaussian fluctuation
measures and a positive history decomposition give the exact
renormalisation theorem, \cref{II-203-9}; its proof is completed
in \cref{sec:finite-renormalisation-proof}.
The complex-source construction begins in
\cref{sec:local-source-energies} and extends these measures
to a complex plaquette coefficient. The quadratic part of a
source changes the Gaussian damping and must be retained
separately; the remaining terms are expanded about the
conditional mean. The resulting local source energies have
uniform bounds, although their complete sum need not be
bounded uniformly in the cutoff.

Separated source derivatives avoid this last difficulty.  A local
energy contributes to a mixed derivative only when its localization
domain meets all the source supports.  At sufficiently fine scales,
this forces a large domain and its exponential decay makes the
scale sum converge.  Reflection positivity and chessboard
estimates turn these local bounds into bounds on cumulants,
uniform in the cutoff and volume.  A countable determining family
then gives simultaneous subsequential limits on all separated
tuples; flatness at the coincidence diagonals permits their
Schwartz extensions.

Full convergence on a fixed torus requires a separate comparison.
We discard histories whose first large-field region occurs too
early and compare the remaining flows after their early
interactions have decayed.  Classical backgrounds on consecutive
meshes are compared in a common gauge.  The resulting real-source
estimate is continued to complex sources and differentiated after
subtracting the local contact terms.  This gives a summable
cutoff difference.  Dilation places any positive physical inverse
coupling in the interval where the comparison applies.

For infinite-volume Gibbs states, reflection at coordinate and
diagonal planes restricts the translation spectrum.  The slab
insertion bound yields analytic point fields and controls their
growth near a narrow slab.  Angular continuation then proves
full rotation invariance.  These estimates hold for every limit
in a Gibbs frame, including limits obtained by invariant
conditioning.  The canonical component measure therefore has
rotation-invariant components with a uniform clustering bound.

Finally the point-insertion estimate gives polynomial energy
bounds for smeared fields.  Their ordered products have tempered
boundary values.  Complex rotation gives Lorentz covariance,
and continuation through spacelike configurations gives locality.
The clustering estimate makes the vacuum of each component
unique.

For the short-distance law we transport the covariance through
the small-field history. The linearized curvature and its Gaussian
covariance produce the squared dipole kernel. Four fluctuation
cumulants control the departure from this Gaussian calculation;
locality and the rarity of early large-field regions control the
remaining history. A choice of intermediate scales makes all
these errors small before the two-state comparison transfers the
estimate to arbitrary admissible volumes and periodic Gibbs states.
Shrinking one fixed test pair then recovers the physical label.

\section{Finite renormalisation and critical backgrounds}\label{sec:finite-renormalisation-overview}

For the finite renormalisation construction, write
\(A(U)=A(1,U)\), so the bare measure is
\(e^{-\beta A(U)}\,dU\).
Blocking by an odd integer \(L\) gives a tower
\(\mathbb T^{(0)},\ldots,\mathbb T^{(K)}\); on
\(\mathbb T^{(j)}\) each side has \(2L^{m+K-j}\)
sites. Write \(N_j=\#\mathbb T^{(j)}\).
Thus the final side length is \(\ell=2L^m\).
We use the nonabelian block average \(\mathcal M\)
of \cref{sec:nonabelian-mean}: the symmetrized contour
mean on its small-curvature domain, with the global
real extension specified in \eqref{I-3-49}--\eqref{I-3-50}.

The theorem concerns a finite number of steps on
each such torus. Its local expansion uses regular
critical backgrounds and retains the fluctuation
coordinates outside the integration regions.

\begin{theorem}[Finite renormalisation]\label{II-203-9}
There is a choice of an odd blocking factor \(L\ge13\),
integers \(r\ge2\), \(\mathfrak p_0\ge6r+1\),
a power \(M=L^{m'}\), \(m'\in\mathbb Z_{\ge1}\),
and positive construction constants, including
\[
 A_0,\ A_1,\ B_\sharp,\ \kappa,\ E_0,\ B,\ B_{\rm fr},
 \ C_{\rm maj},\ c_\star,\
 c_{\rm low},\ \gamma,
 \qquad 0<\gamma<1,
\]
fixed independently of the lattice and the number of
steps, for which the following holds.
Let \(K\ge1\), \(m\ge1\), and \(x_0,\ldots,x_K\) satisfy
\[
 x_j\ge\gamma^{-2}\quad(0\le j\le K),\qquad
 x_j=L^{-1}x_{j-1}-b_j,\quad |b_j|\le B_\sharp
                                  \quad(1\le j\le K).
\]
Put \(\beta=x_0\), \(g_j=x_j^{-1/2}\),
\(p_0(g)=A_0(\log g^{-2})^{\mathfrak p_0}\), and
\[
 R(x)=\min\{L^a:a\in\mathbb Z_{\ge0},\
                              L^a\ge(\log x)^r\}.
\]
Assume \(MR(x_K)\le L^m\).
For every \(0\le k\le K\), the effective measure
\(\rho_k\,dV_k=(\mathcal M^k)_*(e^{-\beta A}\,dU)\)
has the following properties.

\begin{enumerate}[label=(\roman*)]
\item It is the sum of finite positive history measures,
\[
 \rho_k\,dV_k=\sum_{h\in\mathcal H_k}\tau_h,\qquad
 \sum_h\mathsf t(h)=Z_{K,\ell,\beta}
             :=\int e^{-\beta A}\,dU\in(0,1],
\]
where \(\mathsf t(h)=\tau_h(\mathcal V_k)\).
The histories and their regions are defined by the
local tests in \cref{sec:history-decomposition}.

\item On its completed free and stored coordinates,
each history measure is the marginal of the
nonnegative integrable density
\[
 \mathsf{Stay}^{\rm T}_h\,\chi_k(\Omega_k)
                    \exp(A_{k,h}+\Phi_{k,h})\,d\mathscr M_h.
\]
The exponent \(A_{k,h}\) is
\[
 -A(\beta_k(\cdot),\mathsf U)
       +\mathbf E_k(\mathsf U)
       +\mathbf E_k^{\rm bd}(\mathsf U,\mathfrak q)
       -\mathcal E_{k,\mathcal S(h)},
\]
with the finite sums in \eqref{II-99-12}--\eqref{II-99-19};
\(\Phi_{k,h}\) is the sum of the retained quadratic and
pivot exponents. On the support of the tests these two
exponents are defined outside a Borel
\(\mathscr M_h\)-null set and are independent of the
admitted adapted representative \(\mathsf U\).
The density is zero off that support.

\item The pointed regular terms are holomorphic on
their fixed analytic cores, depend only on the
background pair, are gauge invariant and conjugation
real, and satisfy
\[
 |E^{(j)}(X,P,z)|\le E_0e^{-(\kappa+1)d_j(X)}.
\]
They depend only on \(x_0,\ldots,x_{j-1}\).
The created boundary terms have norm at most
\(Be^{-\kappa d_j(X)}\) on their layered output
domains. At fixed stored data, the sum of the
oscillations of the retained exponents at \((j,X)\)
is at most
\(B_{\rm fr}\mathfrak b_{j-1}^2e^{-\kappa d_j(X)}\),
where \(\mathfrak b_i=A_1(\log x_i)^{\mathfrak p_0}\).
Earlier functions and their domains are unchanged
on extending a history or the coupling prefix.

\item For \(1\le k\le K\), each history obeys
\[
 e^{\mathcal E_{k,\mathcal S(h)}}\mathsf t(h)
 \le\exp\!\left[
 C_{\rm maj}\left(N_k+\sum_{j=1}^k\#_j Z_j(h)\right)
       -c_\star\sum_{j=1}^kp_0(g_j)^2n_j^{\rm cr}(h)\right].
\]
Here \(Z_j=\Lambda_j^c\), and \(n_j^{\rm cr}\) counts
the cells created at step \(j\).
The vacuum-normalized partition function satisfies
\[
 e^{-C_{\rm low}(x_k)N_k}
       \le e^{-E_{(k)}}Z_{K,\ell,\beta}
       \le e^{(C_{\rm maj}+2)N_k},
 \qquad
 C_{\rm low}(x)=c_{\rm low}+\tfrac92\log x.
\]
Here \(E_{(k)}\) is the accumulated vacuum constant
defined in \eqref{II-97-8}.
\end{enumerate}

There is in addition a unique coefficient-determined
real coupling sequence from every \(\beta\ge\gamma^{-2}\).
It is defined successively by
\[
 x_0=\beta,\qquad
 b_{j+1}=\mathsf{cf}(\widehat E^{(j+1)}[x_0,\ldots,x_j]),
 \qquad x_{j+1}=L^{-1}x_j-b_{j+1}.
\]
Each coefficient is defined before the next coupling,
is independent of the lattice, and satisfies
\(|b_{j+1}|\le B_\sharp\).
Here \(\widehat E^{(j+1)}\) is the regular family
of \cref{II-147-20}, and \(\mathsf{cf}\) is its
second-moment coefficient \eqref{II-56-18}.
The sequence reaches \(1<x_J<\gamma^{-2}\) in
finitely many steps. For every \(K<J\), the preceding
conclusions apply on every torus with
\(MR(x_K)\le L^m\).
\end{theorem}

The principal issue is preservation of the density
under a change of critical background. The small-field
integral gives local regular terms, while the terms
meeting a large-field boundary remain attached to
their original scale. The kept quadratic forms retain
coercivity on the stored fluctuations. A positive
change of variables identifies the resulting density
and proves its integrability at each step.
The failed local tests then provide a gain for each
created large-field cell; that gain absorbs the number
of histories. The proof of \cref{II-203-9} is completed
in \cref{sec:finite-renormalisation-proof}.
The weighted bounds \eqref{II-203-24} and
\eqref{II-203-25} retain quantitative gains for
large-field components and their earlier deposits;
their parameter ranges and thresholds are stated
with those bounds.

\subsection{Critical backgrounds}


Fix an odd \(L\ge13\) and a finite periodic cubic lattice tower
\(\mathbb T^{(0)}\supset\cdots\supset\mathbb T^{(K)}\).
Successive spacings have ratio \(L\). At each level the three
sides have the same number of sites, twice a power of \(L\).
In level-\(k\) units, \(1\le k\le K\),
\[
 \eta=L^{-k},\qquad \xi_j=L^{j-k}\quad(0\le j\le k).
\]
A configuration assigns \(U(b)\in\mathrm{SU}(2)\) to oriented
bonds, with \(U(\bar b)=U(b)^{-1}\). Write \(\partial U(p)\)
for the ordered plaquette product and
\[
 A_{\Omega_0}(U)
       =\sum_{p\subset\Omega_0}
                       \left(1-\tfrac12\operatorname{tr}\partial U(p)\right).
\]
Here inner plaquettes have all corners in \(\Omega_0\).
Use operator norms, and let \(\mathfrak g\) be the traceless
Hermitian \(2\times2\) matrices,
\(\mathfrak g^c=\mathfrak g\oplus i\mathfrak g\).

A decreasing sequence
\(\Omega_0\supset\cdots\supset\Omega_k\), with
\(\Omega_{-1}=\Omega_0\), \(\Omega_{k+1}=\emptyset\), is
\emph{separated} with parameters \(\mathfrak M,R_1\), powers
of \(L\), if each \(\Omega_j\) is a union of cubes of side
\(\mathfrak M\xi_j\) in the fixed centred partition, is empty
when no such cubes exist, and
\[
 \operatorname{dist}_\infty(\Omega_j^c,\Omega_{j+1})
                         >R_1\mathfrak M\xi_j\qquad(0\le j<k).
\]
Complements are closed; distance to the empty set is infinite.
Let \(\Gamma_j^\flat\) be the set of level-\(j\) bonds with
both endpoints in \(\Omega_0\), at least one in \(\Omega_j\),
and neither in \(\Omega_{j+1}\). The reduced determining set is
\(\mathbf B^\flat=\bigcup_j\Gamma_j^\flat\).

For the iterated nonabelian mean \(\mathcal M^j\) defined in
\cref{sec:nonabelian-mean,sec:multiscale-constraints}, put
\[
 M_{\mathbf B^\flat}(U)
       =\bigl((\mathcal M^jU)|_{\Gamma_j^\flat}\bigr)_{j=0}^k .
\]
A datum \(V=(V_j)\) assigns group elements to these rows,
with inversion under reversal. A level-\(j\) plaquette is
\emph{tested} if it has a bond in \(\Gamma_j^\flat\) and every
bond outside \(\Gamma_j^\flat\) has both endpoints in
\(\Omega_{j-1}\setminus\Omega_j\). Its holonomy uses \(V_j\)
on retained bonds and \(\mathcal M(V_{j-1})\) elsewhere.
Criticality over \(V\) means membership in
\(M_{\mathbf B^\flat}^{-1}(V)\) and vanishing derivative of
\(A_{\Omega_0}\) along every differentiable curve therein.
Residual gauges equal \(I\) at all retained-bond endpoints.
The axial condition is
\(\operatorname{Ax}^{\circ}(\mathbf B^\flat,I)\) of
\cref{I-3-95}, with its fixed coordinate combs.

For a bond or plaquette \(\mathfrak e\), let
\(\mathsf L(\mathfrak e)=\xi_j\) for the largest \(j\) such
that \(\Omega_j\) contains an endpoint or corner.
A bond is \emph{full} if all incident plaquettes are inner.
The class \(\mathfrak U'_k(r)\) consists of real configurations
on inner bare bonds satisfying
\[
 |\partial U(p)-I|<r\eta^2\mathsf L(p)^{-2},\qquad
 |(D_U^{\mathsf T}\partial U)(b)|
                            <r\eta^2\mathsf L(b)^{-3}
\]
at inner plaquettes and full bonds. Here \(D_U^{\mathsf T}\)
is the covariant plaquette divergence \eqref{I-2-5},
including scalar matrix parts. The action current is
\[
 J(U)=D_U^{\mathsf T}F_U,\qquad
 F_U(p)=\frac{\partial U(p)-\partial U(p)^{-1}}{2i\eta^2}.
\]

Relative to \(U_0\), put
\[
 F(A)=\eta^{-1}A_{\Omega_0}(e^{i\eta A}U_0),\qquad
 \|A\|_{X,\mathbf B}
     =\max\{|A|_{(-1)},|\nabla^{U_0}A|_{(-2)}\}.
\]
Here \(|f|_{(-s)}=\sup\mathsf L^s|f|\). Use only inner
differences, weighted by the larger local scale of their
two bonds, with the zero-extension convention of \cref{I-45-14}.
Use the mean-Landau projection \(\mathsf R_{\mathsf m}\) of
\eqref{eq:bg-projection-operators}, the dimensionless relative
straightened mean \(\mathcal Q\) of
\cref{sec:actual-local-nonlinear-maps}, and the inverse
datum twist \(B_\natural\) of
\eqref{eq:background-inverse-twist}.
Transposes use bare weight \(\eta^3\) and row weight
\(\xi_c^{-1}\), where \(\xi_c=\xi_j\) at level \(j\).

\begin{theorem}[Canonical critical background]\label{I-58-22}
There are constants
\[
 a_{\rm orb}>0,\qquad B_{\mathrm{rad}}\ge\max\{4,L^2\},
 \qquad C_\partial\ge4
\]
and powers \(\mathfrak M_{\mathrm G},\rho_{\mathrm G}\) of
\(L\), depending only on \(L\), with the following property.
For every power \(\rho\ge\rho_{\mathrm G}\) of \(L\), there
are positive numbers
\(\varepsilon_{\mathrm G}(\rho),r_{\mathrm G}(\rho),
K_{\mathrm G}(\rho)\), depending only on \(L,\rho\).
Suppose
\[
 \mathfrak M\ge\mathfrak M_{\mathrm G},\qquad R_1\ge L,\qquad
 R_1\mathfrak M\ge L^2\rho,
\]
and \(\mathbb T^{(k)}\) has at least \(6\rho\) sites in
each direction. Let the region sequence be separated with
these parameters. Let \(V_0\) be real reduced data whose
tested holonomies are defined and satisfy
\[
 |\partial V_0-I|<\varepsilon_1,\qquad
              0<\varepsilon_1\le\varepsilon_{\mathrm G}(\rho).
\]

For every real row field \(B\) with
\(\sup|B|<r_{\mathrm G}(\rho)\), the critical configurations
over \(e^{iB}V_0\) in \(\mathfrak U'_k(a_{\rm orb})\)
form one nonempty residual gauge orbit. This orbit is
contained in
\[
 \mathfrak U'_k\bigl(
        B_{\mathrm{rad}}(\varepsilon_1+C_\partial\sup|B|)\bigr).
\]
Here the perturbed datum is defined on positive rows by
\((e^{iB}V_0)(c)=e^{iB(c)}V_0(c)\), and on reversed rows
by inversion. Denote the unique axial representative by
\(U^\circ(B)\), and set \(U_0=U^\circ(0)\).

There is exactly one holomorphic map \(\Phi_{V_0}\) from
the complex row ball
\(\mathcal B=\{B:\sup|B|<r_{\mathrm G}(\rho)\}\) to
\(\mathrm{SL}(2,\mathbb C)\)-valued configurations on the
inner bare bonds which equals \(U^\circ(B)\) at real \(B\).
It has the form
\[
 \Phi_{V_0}(B)=(e^{i\eta A(B)}U_0)^{h(B)},\qquad
                       A(0)=0,\quad h(0)=I ,
\]
where \(A,h\) are holomorphic and are respectively real
and unitary at real \(B\). There is a holomorphic multiplier
\(\lambda(B)\), real at real \(B\), such that
\begin{gather*}
 \mathsf R_{\mathsf m}D^{\mathsf T}A=0,\qquad
 \mathcal Q(A)=B_\natural,\qquad
 F'(A)=\mathcal Q'(A)^{\mathsf T}\lambda(B),\\
 \|A(B)\|_{X,\mathbf B}
  +|D^{\mathsf T}DA(B)|_{(-3)}
  +\sup_x|h(B)(x)-I|
                    \le K_{\mathrm G}(\rho)\sup|B|,\\
 \sup|\lambda(B)|
                    \le K_{\mathrm G}(\rho)(\varepsilon_1+\sup|B|).
\end{gather*}
The covariant difference operators use the reference \(U_0\), and
\(D^{\mathsf T}D\) is the inner-plaquette curl followed
by its transpose. All actual complex averages are defined,
and
\[
               M_{\mathbf B^\flat}(\Phi_{V_0}(B))=e^{iB}V_0 .
\]
At every inner plaquette and full bond, respectively,
\begin{align*}
 |\partial\Phi_{V_0}(B)(p)-I|
  &<\bigl(B_{\mathrm{rad}}\varepsilon_1
             +K_{\mathrm G}(\rho)\sup|B|\bigr)
                                      \eta^2\mathsf L(p)^{-2},\\
 |J(\Phi_{V_0}(B))(b)|
  &<\bigl(B_{\mathrm{rad}}\varepsilon_1
             +K_{\mathrm G}(\rho)\sup|B|\bigr)
                                                \mathsf L(b)^{-3}.
\end{align*}
\end{theorem}

The uniqueness assertion concerns the specified regular
class. It does not require the critical point to minimize
the action over the entire fibre. All radii and bounds in
the theorem are independent of the number of layers, the
lattice spacing and the volume, once \(L,\rho\) are fixed.
The proof is given in \cref{sec:canonical-background}.

\subsection{Fluctuations, response and Gaussian coefficients}

The fluctuation theorem applies to the natural regular
pair class of \cref{def:natural-regular-pair}, which controls
curvature, a complex potential and its first covariant
differences, the current, and real gauges on cubes contained
in the union of two successive layers. For an additional separated region
and any subset of its coarse sites,
\cref{I-73-90} constructs both the Landau and axial
fluctuation covariances on the corresponding nonpivot
coordinates. Under its stated geometric thresholds and
smallness bounds, either covariance \(\mathcal C\) satisfies
\[
 |\mathcal C(b,b')|\le B_0e^{-\delta_{\rm cv}|b_--b'_-|_1}.
\]
It is the sum of an absolutely convergent expansion over
tree walks. Each term is holomorphic in the pair restricted
to its carrier, and the sum of the terms with \(m\) edges
is bounded by
\(2^{-m-1}B_0e^{-\delta_{\rm cv}|b_--b'_-|_1}\).
The constants are independent of the number of layers,
the volume and the chosen subset. These are covariances
of the constrained fluctuation form specified in that
theorem.

The same local expansions allow parameters to turn off
couplings across prescribed walls. The nonlinear response
of \cref{I-88-110} is jointly holomorphic in those parameters,
the local regular pair and the small coarse datum. It has
weighted Lipschitz bounds, hence exponential decay of the
response to a change in a distant datum row. When wall
parameters vanish, the solution splits into the corresponding
components; each component depends only on the pair and
datum on its stated local region. The parameter domains
and their associated datum radii are specified in that
theorem. At parameters equal to one, the response is
identified with the critical background, in its
mean-Landau coordinates, by \cref{I-88-127,I-88-131}.
The converse uses the reduced determining set and the
common small ball fixed by the criticality radius
\(\varepsilon_{\rm crit}\) and the response radius.

There is a separate uniform statement for the Gaussian
normalizers near the trivial configuration. For the
one-step normalizer \(\Phi^\natural_l\), let
\(\pi^{(l)}_{\mu\nu}(u)\) be its infinite-volume second
variation on individual unit-lattice bonds, with the
colour factor removed. \Cref{I-109-104} proves that this
kernel is real and exponentially decreasing, and that
\[
 b_l^0=-\tfrac12\sum_{u\in\mathbb Z^3}
                            u_\mu^2\pi^{(l)}_{\nu\nu}(u),
 \qquad \mu\ne\nu,
\]
converges absolutely and is independent of the distinct
directions. It depends only on \(l,L\), and
\[
 \sup_{k\ge0}\left|
       \sum_{l=1}^kL^{-(k-l)}b_l^0\right|\le c_2(L)<\infty .
\]
The nonlinear Gaussian normalizers are germs at the
trivial configuration; their neighbourhoods may depend
on the height and the finite volume. The bound displayed
above is uniform in the height.

\subsection{The proof}

The block averages and gauge conditions follow the
multiscale framework of Ba{\l}aban
\cite{B-CMP98,B-CMP99a}; his background propagators provide
the corresponding analytic setting \cite{B-CMP99b}.
We begin with the local averaging maps and gauge
representatives, keeping the actual averaging constraint
throughout the construction.

The linear estimates are proved in local gauges and
assembled over separated regions. Coercivity controls the
gauge and average constraints, while the gap between
successive regions makes the localization error small.
This gives a constrained inverse with exponential decay
and uniform control of one covariant difference.
The nonlinear equation is then solved with its initial
current retained. A correction imposes the gauge condition
and the prescribed coarse averages.

To construct the real background, we descend the given
data to the fine lattice and then solve upward through
the levels. After each step a regularity improvement
restores a fixed radius. Comparing two solutions in a
common gauge proves uniqueness of the regular orbit.
The real background supplies the local regular gauges
needed by the complex response, and the axial condition
identifies its continuation with \(\Phi_{V_0}\).

The covariance expansion adds the fluctuation coordinate
to the same constrained system. Local terms and current
multipliers can then be organized into convergent tree
walks. Their carrier dependence is what makes the
decoupled response local. Finally, Gaussian gluing makes
the accumulated coefficients telescope. A Ward identity
reduces their quadratic part to a potential response;
the remainder is controlled by a trace with three
propagators. The remaining massive coefficient is computed
in a constant abelian field and bounded uniformly as the
mass decreases.

\section{The finite-lattice model and its sources}\label{sec:finite-lattice}
\subsection{Configurations and the Wilson measure}
Let $L\ge3$ be odd, let $K,m$ be positive integers, and put
\[
a_j=L^{-j}\quad(j\ge0),\qquad \ell=2L^m,\qquad
T_\ell=\mathbb R^3/\ell\mathbb Z^3,\qquad
\Lambda_{K,\ell}=a_K\mathbb Z^3/\ell\mathbb Z^3.
\]
The group is $G=\mathrm{SU}(2)$, and $\tr$ denotes the ordinary
matrix trace, so that $\tr\mathbf1=2$.
Distances used in Lipschitz norms on $T_\ell$ are the quotient
Euclidean distances in physical units. Write $a=a_K$ and $N=\ell/a=2L^{m+K}$ when the
parameters are fixed.

A bond is an ordered pair $b=\langle x,x+\epsilon ae_\mu\rangle$,
where $x\in\Lambda_{K,\ell}$, $\epsilon\in\{-1,1\}$ and
$1\le\mu\le3$. It is positive when $\epsilon=1$, and $\bar b$
denotes its reverse. A configuration assigns $U_b\in G$ to each
bond, subject to $U_{\bar b}=U_b^{-1}$. Thus the configuration
space $\mathcal V_0$ is the finite product of copies of $G$ indexed
by the positive bonds. We equip it with the product Haar probability
measure $dU$.

For $\mu<\nu$, let $p=p_{\mu\nu}(x)$ be the elementary plaquette
based at $x$, with barycentre
$x(p)=x+\tfrac a2(e_\mu+e_\nu)$. Its boundary product is
\begin{equation}\label{I-0-3}
\begin{split}
U_p={}&U(\langle x,x+ae_\mu\rangle)
U(\langle x+ae_\mu,x+ae_\mu+ae_\nu\rangle)\\
&\quad\cdot U(\langle x+ae_\nu,x+ae_\mu+ae_\nu\rangle)^{-1}
U(\langle x,x+ae_\nu\rangle)^{-1}.
\end{split}
\end{equation}
Each plaquette is counted once in sums over $p$. There are
$\mathsf n=3N^3$ plaquettes and the same number of positive bonds.
We use Wilson's plaquette action \cite[\S III]{Wilson1974}, with
the following normalization:
\begin{equation}\label{I-0-4}
a_p(U)=1-\tfrac12\tr U_p,\qquad A_K(U)=\sum_p a_p(U).
\end{equation}
For $V\in G$, $V+V^{-1}=(\tr V)\mathbf1$ and $V^{-1}=V^*$,
so $\tr V\in[-2,2]$ and $0\le A_K\le2\mathsf n$.
For $\beta\in\mathbb R$, define
\begin{equation}\label{I-0-13}
\rho_0(U)=e^{-\beta A_K(U)},\qquad
Z_{K,\ell,\beta}=\int\rho_0\,dU,\qquad
\mu_{K,\ell,\beta}(dU)=Z_{K,\ell,\beta}^{-1}\rho_0(U)\,dU.
\end{equation}
The continuous density on the compact configuration space satisfies
$e^{-2|\beta|\mathsf n}\le\rho_0,Z_{K,\ell,\beta}
\le e^{2|\beta|\mathsf n}$; for $\beta\ge0$ the bounds improve to
$e^{-2\beta\mathsf n}\le\rho_0,Z_{K,\ell,\beta}\le1$.
Thus $\mu_{K,\ell,\beta}$ is a Borel probability measure.
We write $\langle F\rangle_{K,\ell,\beta}=\int F\,d\mu_{K,\ell,\beta}$.

A gauge transformation $u:\Lambda_{K,\ell}\to G$ acts by
\[
U^u(\langle x,y\rangle)=u(x)U(\langle x,y\rangle)u(y)^{-1}.
\]
Cancellation along a plaquette gives $U^u_p=u(x)U_pu(x)^{-1}$,
so $a_p$, $A_K$ and $\rho_0$ are invariant. Left and right Haar
invariance on each bond then gives gauge invariance of $dU$ and
$\mu_{K,\ell,\beta}$.

\subsection{The observable and its source}
For a complex weight $c$ on plaquettes, or a function on $T_\ell$
evaluated at their barycentres, set
\begin{equation}\label{I-2-12}
A(c,U)=\sum_p c(p)a_p(U),\qquad
A^\xi(c,U)=\xi^{-1}A(c,U).
\end{equation}
Here $\xi$ is the spacing of the same lattice in the chosen units.
The scalar observable used below is the weighted action in physical
units,
\begin{equation}\label{I-2-13}
\mathcal O(f)=A^{a_K}(f,U)
=a_K^{-1}\sum_p f(x(p))a_p(U).
\end{equation}
For real continuous $f$ it is a real continuous function on
$\mathcal V_0$, and
\begin{equation}\label{eq:finite-observable-bound}
|\mathcal O(f)|\le2a_K^{-1}\mathsf n\|f\|_\infty.
\end{equation}

From now on take $\beta>0$. Its physical form is
\begin{equation}\label{I-2-21}
\mathsf b=a_K\beta,
\qquad \beta A_K(U)=\mathsf b A^{a_K}(1,U).
\end{equation}
For real Lipschitz functions $f_1,\dots,f_n$ on $T_\ell$,
$n\ge1$, and $z\in\mathbb C^n$, put $w=\sum_i z_i f_i$.
The source changes the physical inverse coupling to $\mathsf b-w$:
\begin{equation}\label{I-2-22}
\rho_0^w(U)=\exp[-A^{a_K}(\mathsf b-w,U)]
=\rho_0(U)\exp\Bigl(\sum_i z_i\mathcal O(f_i)\Bigr).
\end{equation}
Write
$Z(z)=Z_{K,\ell,\beta}[w]=\int\rho_0^w\,dU$.
The positive sign of the source term in \eqref{I-2-22} fixes the
cumulant convention. In particular a constant source $w=\zeta$
replaces the bare inverse coupling by $\beta-L^K\zeta$.

\subsection{Generating functions and cumulants}
We recall the finite-measure argument that justifies differentiation
of $Z$. Let $X_1,\dots,X_n$ be bounded complex measurable
functions on a probability space $(\Omega,\nu)$, and put
\[
\Gamma(z)=\int e^{\sum_i z_iX_i}\,d\nu,\qquad c_i=\sup|X_i|.
\]
The logarithm near the origin is the branch with value zero there.
The joint cumulant is
\begin{equation}\label{I-0-14}
\langle X_1;\dots;X_n\rangle_\nu
=\left.\partial_{z_1}\cdots\partial_{z_n}\log\Gamma(z)\right|_{z=0}.
\end{equation}

\begin{lemma}\label{I-0-17}
The function $\Gamma$ is entire and satisfies
\[
|\Gamma(z)|\le e^{\sum_i|z_i|c_i},\qquad
|\Gamma(z)-1|\le e^{\sum_i|z_i|c_i}-1.
\]
On $D=\{z:\sum_i|z_i|c_i<\log2\}$ the logarithmic series
defines a holomorphic logarithm of $\Gamma$ with value zero at the
origin. In particular \eqref{I-0-14} is well defined.

Let $\mathcal P$ range over the partitions of $\{1,\dots,n\}$
into nonempty blocks, and write
$m_B=\int\prod_{i\in B}X_i\,d\nu$ and
$\kappa_B=\langle X_i;i\in B\rangle_\nu$. Then
\begin{align}
\langle X_1;\dots;X_n\rangle_\nu
&=\sum_{\mathcal P}(-1)^{|\mathcal P|-1}(|\mathcal P|-1)!
  \prod_{B\in\mathcal P}m_B,\label{eq:cumulant-partitions}\\
\int\prod_{i=1}^nX_i\,d\nu
&=\sum_{\mathcal P}\prod_{B\in\mathcal P}\kappa_B.
\label{eq:moment-partitions}
\end{align}
Thus cumulants are symmetric and complex multilinear, are real for
real variables, and are unchanged by adding constants when $n\ge2$.
For centered variables, \eqref{eq:moment-partitions} reduces to the
sum over partitions with no singleton blocks.

If $\nu=T_*\nu'$ for a measurable map $T:\Omega'\to\Omega$
and a probability measure $\nu'$ on $\Omega'$, then
\[
\langle X_1;\dots;X_n\rangle_\nu
=\langle X_1\circ T;\dots;X_n\circ T\rangle_{\nu'}.
\]
For fixed continuous $X_i$ on a compact metrizable space, the
cumulants are also continuous under weak convergence of Borel
probability measures.
\end{lemma}

\begin{proof}
Normal convergence of the exponential series proves analyticity
and the bounds; on $D$, the logarithmic series also converges
normally. The coefficient of $z_1\cdots z_n$ in $(\Gamma-1)^j$ is
$j!\sum_{|\mathcal P|=j}\prod_{B\in\mathcal P}m_B$.
Substitution in the logarithmic and exponential series gives
\eqref{eq:cumulant-partitions} and \eqref{eq:moment-partitions}.
These imply symmetry, multilinearity, reality and weak continuity,
since each $m_B$ converges under the stated hypotheses.
Adding constants adds only a linear term to $\log\Gamma$;
centering removes singleton cumulants. Pushforwards preserve $\Gamma$.
\end{proof}

Apply the lemma with $\nu=\mu_{K,\ell,\beta}$ and
$X_i=\mathcal O(f_i)$. Since $Z(z)=Z(0)\Gamma(z)$,
$Z$ is entire and
\begin{equation}\label{I-2-23}
\langle\mathcal O(f_1);\dots;\mathcal O(f_n)\rangle_{K,\ell,\beta}
=\left.\partial_{z_1}\cdots\partial_{z_n}\log Z(z)\right|_{z=0}.
\end{equation}
Here the branch is chosen so that $\log Z(0)=\log Z_{K,\ell,\beta}$.
By \eqref{eq:finite-observable-bound}, one sufficient domain for
this logarithm is
\[
2a_K^{-1}\mathsf n\sum_i|z_i|\|f_i\|_\infty<\log2.
\]
Its size depends on both the cutoff and the volume.

\subsection{Source bounds and units of length}
Fix $M>0$ independently of $K$, $\ell$ and $\beta$, and set
\begin{equation}\label{I-0-5}
\|f\|_\sharp=\max\{\|f\|_\infty,40M\Lip f\}.
\end{equation}
The Lipschitz constant is taken in physical units. For $\varrho>0$
define
\[
\mathbb D_\varrho(f_1,\dots,f_n)
=\left\{z\in\mathbb C^n:\sum_i|z_i|\|f_i\|_\sharp<\varrho\right\}.
\]
This is the source domain used in the following density estimates.

In level-$k$ units, $0\le k\le K$, a physical length $r$ is
$L^{K-k}r$, and the bare lattice has spacing $\eta_k=L^{-k}$.
Thus $A^{\eta_k}(c,U)=L^kA(c,U)$.

\begin{lemma}\label{I-2-24}
Let $\beta>0$, $\mathsf b=a_K\beta$, and let $f_1,\dots,f_n$
be real Lipschitz functions on $T_\ell$. For
$z\in\mathbb D_\varrho(f_1,\dots,f_n)$ and $w=\sum_i z_if_i$,
\begin{equation}\label{I-2-26}
\|w\|_\infty<\varrho,\qquad
\Lip w<\frac{\varrho}{40M}.
\end{equation}
For every configuration $U$,
\begin{equation}\label{I-2-27}
\rho_0(U)^{1+\varrho/\mathsf b}
\le |\rho_0^w(U)|
\le\rho_0(U)^{1-\varrho/\mathsf b}.
\end{equation}
In level-$k$ units the same density is
\begin{equation}\label{I-2-25}
\rho_0^w(U)=\exp[-A(\beta-L^Kw,U)]
=\exp[-A^{\eta_k}(L^{-k}\beta-L^{K-k}w,U)].
\end{equation}
The Lipschitz constant of $L^{K-k}w$, measured in level-$k$
units, is $\Lip w$. Its variation between any two points of a
closed cube of side $M$ in those units is less than
$\sqrt3\varrho/40$.
\end{lemma}

\begin{proof}
The triangle inequality gives \eqref{I-2-26}, and linearity
in the weight gives \eqref{I-2-25}. Since $a_p\ge0$,
\[
 |\rho_0^w|=\rho_0e^{A^{a_K}(\operatorname{Re}w,U)},\qquad
 |A^{a_K}(\operatorname{Re}w,U)|\le\varrho A^{a_K}(1,U),
\]
which proves \eqref{I-2-27}. Rescaling source and distance by
$L^{K-k}$ preserves the Lipschitz quotient; a cube of side $M$
has diameter at most $\sqrt3M$.
\end{proof}

For $\xi'=\upsilon\xi$, one has
$A^{\xi'}(c,U)=\upsilon^{-1}A^\xi(c,U)$ and
$A^{\xi'}(\upsilon c,U)=A^\xi(c,U)$.
The physical weight corresponding to a level-$k$ weight $c$
is $a_{K-k}c$; its weight in the unscaled plaquette sum is $L^kc$.

\subsection{Dilation of the finite lattice}

\begin{lemma}\label{I-2-28}
Let $\beta>0$, let $f_1,\dots,f_n$ be real Lipschitz functions
on $T_\ell$, and put $w=\sum_i z_if_i$ for $z\in\mathbb C^n$.
Let $b\in\mathbb Z$ with $0\le b<K$. For $h:T_\ell\to\mathbb C$
write $h^{(b)}(x)=h(L^{-b}x)$ on $T_{L^b\ell}$. Then
\begin{equation}\label{I-2-29}
Z_{K,\ell,\beta}[w]
=Z_{K-b,L^b\ell,\beta}[L^bw^{(b)}],
\end{equation}
and
\begin{equation}\label{I-2-30}
\begin{split}
\langle\mathcal O(f_1);\dots;\mathcal O(f_n)\rangle_{K,\ell,\beta}
={}&L^{bn}\langle\mathcal O(f_1^{(b)});\dots;
\mathcal O(f_n^{(b)})\rangle_{K-b,L^b\ell,\beta}.
\end{split}
\end{equation}
On the right, the observables use spacing $a_{K-b}$ and the
dilated torus. With the same $M$ on both tori,
$\|f_i^{(b)}\|_\sharp\le\|f_i\|_\sharp$, so that
\[
z\in\mathbb D_\varrho(f_1,\dots,f_n)
\quad\Longrightarrow\quad
L^bz\in\mathbb D_{L^b\varrho}(f_1^{(b)},\dots,f_n^{(b)}).
\]
The physical inverse coupling is multiplied by $L^b$.
\end{lemma}

\begin{proof}
Dilation preserves orientations, plaquette products and product
Haar measure. At corresponding barycentres,
$w^{(b)}(L^bx(p))=w(x(p))$ and $a_{K-b}^{-1}L^b=a_K^{-1}$,
so the sourced coefficients agree, proving \eqref{I-2-29}.
Differentiating $Z(z)=\widetilde Z(L^bz)$ in the local logarithms
gives \eqref{I-2-30}. The remaining assertions follow from
$\|f_i^{(b)}\|_\infty=\|f_i\|_\infty$,
$\Lip f_i^{(b)}=L^{-b}\Lip f_i$ and
$a_{K-b}\beta=L^b\mathsf b$.
\end{proof}

\section{Block averaging and image measures}\label{sec:block-averaging}
The block map replaces fine bond variables by coarse ones while
preserving gauge covariance. On configurations with small plaquette
curvature, contour transports lie close together and their logarithmic
mean depends analytically on the fine field. Two further rules extend
the map to all real configurations. We prove its measurability before
using it to transform measures.

\subsection{Blocks and contour comparison}
For integers $j\ge-m$, put
\[
 a_j=L^{-j},\qquad
 \Lambda_{j,\ell}=a_j\mathbb Z^3/\ell\mathbb Z^3.
\]
Thus the level-$k$ lattice is
$\mathbb T^{(k)}=\Lambda_{K-k,\ell}$, $0\le k\le K$.
Its configuration space $\mathcal V_k$ consists of the $G$-valued
bond assignments satisfying $U(\bar b)=U(b)^{-1}$. The norm
$|x|_\infty$ on $T_\ell$ is the least sup norm of a lift of $x$.
The cube centered at $y\in\Lambda_{j,\ell}$ and its block of order
$n\ge0$ are
\[
 \Delta_j(y)=\{x:|x-y|_\infty\le a_j/2\},\qquad
 B^n(y)=\Delta_j(y)\cap\Lambda_{j+n,\ell}.
\]
Write $B(y)=B^1(y)$ and $h=(L-1)/2$. The specified lattice of
centers determines the scale suppressed in this notation. In particular,
\[
 B^n(y)=\left\{y+a_{j+n}\sum_{\mu=1}^3 r_\mu e_\mu:
 r_\mu\in\mathbb Z,\quad |r_\mu|\le\frac{L^n-1}{2}\right\}.
\]
These $L^{3n}$ points are distinct. The blocks centered on
$\Lambda_{j,\ell}$ partition $\Lambda_{j+n,\ell}$, and
\[
 B^{n+n'}(y)=\bigcup_{x\in B^n(y)}B^{n'}(x),\qquad
 \Delta_j(y)=\bigcup_{x\in B^n(y)}\Delta_{j+n}(x).
\]
For $j>-m$, every fine bond either has both endpoints in one
block of order $n$, or joins two such blocks whose centers are
the endpoints of a unique coarse bond of the same orientation
and direction. In the second case we say that the fine bond
lies over the coarse bond.

These assertions follow from the unique decomposition
$N=L^nq+r$, $|r|\le(L^n-1)/2$, of an integer $N$. Applied in
each coordinate, it gives the partition into blocks. Two admissible
coordinate vectors cannot represent the same site: their differences
have absolute value at most $L^n-1$, whereas the period in fine
lattice units is $2L^{m+j+n}\ge2L^n$. Applying the decomposition
twice gives the nesting identity. The intervals of length $a_{j+n}$
centered at the consecutive fine sites fill the interval of length
$a_j$, proving the cube identity. Finally, increasing one fine
coordinate by one stays in the same block except at
$r_\mu=(L^n-1)/2$; there it enters the block centered at
$y+a_je_\mu$. Reversal handles negative bonds. Since the coarse
lattice has at least six sites along each axis when $j>-m$,
the coarse bond is determined by its endpoints.

If $X$ is a union of cubes $\Delta_j(y)$, a fine site
$x\in B^n(y)$ belongs to $X$ exactly when its center $y$ does.
Indeed, $|x-y|_\infty\le(a_j-a_{j+n})/2<a_j/2$, whereas
distinct cube centers are at sup distance at least $a_j$.
Thus $x$ lies in the interior of $\Delta_j(y)$ and in no other
coarse cube. This observation permits the same block construction
inside a region.

Let $j>-m$ and $y\in\Lambda_{j,\ell}$. For
$x=y+a_{j+1}\sum_\mu r_\mu e_\mu\in B(y)$, write
$S=\{\mu:r_\mu\ne0\}$. For each ordering $\sigma$ of $S$,
let $\Gamma_\sigma(y,x)$ traverse the $|r_{\sigma(1)}|$ bonds
in direction $\operatorname{sign}(r_{\sigma(1)})e_{\sigma(1)}$,
then those in direction $\operatorname{sign}(r_{\sigma(2)})e_{\sigma(2)}$,
and so on. Following \cite[p.~252]{B-CMP109}, the contour family is
\[
 \mathcal G(y,x)=\{\Gamma_\sigma(y,x):\sigma\text{ orders }S\}.
\]
It consists of $|S|!\le6$ paths; $\mathcal G(y,y)$ consists of
the empty path. For a path $\Gamma=(b_1,\ldots,b_r)$ put
$U(\Gamma)=U(b_1)\cdots U(b_r)$, with value $\mathbf1$ for the
empty path. If $c$ is a coarse bond and $x\in B(c_-)$, set
$x'=x+(c_+-c_-)$ and write $[x,x']$ for the straight fine path
between them. We abbreviate $U[c]=U([c_-,c_+])$.

To compare these products, take the spacing of a lattice
$\Lambda_{r,\ell}$, $r>-m$, as the unit of length. A step sequence
$s=(s_1,\ldots,s_q)$ has entries in
$\{\pm e_1,\pm e_2,\pm e_3\}$. Its path from $z$ is denoted
by $\Gamma(z;s)$. Define
\begin{equation}\label{I-3-6}
 \Sigma(z;s)=\left\{z+\sum_{i\in I}s_i:
 I\subset\{1,\ldots,q\}\right\}.
\end{equation}
All rearrangements of $s$ have the same endpoint and the same set
$\Sigma(z;s)$, which contains every site visited by their paths.
We use the operator norm for matrices.

\begin{lemma}[Rearrangement of steps]\label{I-3-7}
Let $U$ be a $G$-valued configuration, let $s$ have length $q$,
and let $s'$ be a rearrangement of $s$. If $\omega\ge0$ and
$|U_p-\mathbf1|\le\omega$ for every plaquette whose four corners
lie in $\Sigma(z;s)$, then
\begin{equation}\label{I-3-8}
 \bigl|U(\Gamma(z;s'))U(\Gamma(z;s))^{-1}-\mathbf1\bigr|
 \le(1+\omega)^{q(q-1)/2}-1.
\end{equation}
\end{lemma}

\begin{proof}
First exchange two adjacent steps. The initial path before them is
unchanged, and the common final path cancels from the quotient of
the two transports. If $P$ is the initial transport, this quotient
is $PHP^{-1}$, where $H$ traverses the four-edge path formed by
the two orders of the exchanged steps. Equal or opposite steps
give $H=\mathbf1$. Otherwise $H$ is a conjugate of $U_p$ or
$U_p^{-1}$ for the plaquette with those four corners. All corners
belong to $\Sigma(z;s)$. Unitary invariance of the operator norm
and $|V^{-1}-\mathbf1|=|V-\mathbf1|$ for $V\in G$ therefore give
\begin{equation}\label{I-3-9}
 U(\Gamma(z;s''))=(\mathbf1+E)U(\Gamma(z;s)),
 \qquad |E|\le\omega.
\end{equation}
Any permutation of $q$ entries is a product of at most
$q(q-1)/2$ adjacent exchanges: sorting a permutation removes one
inversion at each exchange. Each intermediate sequence has the
same set $\Sigma(z;s)$, so the plaquette hypothesis continues
to apply. After $T$ exchanges the quotient is
$(\mathbf1+E_T)\cdots(\mathbf1+E_1)$. Expanding this product and
using submultiplicativity gives
\[
 \left|(\mathbf1+E_T)\cdots(\mathbf1+E_1)-\mathbf1\right|
 \le\prod_{t=1}^T(1+|E_t|)-1
 \le(1+\omega)^{q(q-1)/2}-1.
\]
\end{proof}

\begin{lemma}[Contours in a double block]\label{I-3-10}
Let $j>-m$ and express lengths in units of $a_{j+1}$. If
$x=y+\sum_\mu r_\mu e_\mu\in B(y)$, all paths in
$\mathcal G(y,x)$ are rearrangements of the same
$q=\sum_\mu|r_\mu|\le3(L-1)/2$ steps. Their common subset-sum set is
\[
 \left\{y+\sum_{\mu\in S}t_\mu
 \operatorname{sign}(r_\mu)e_\mu:
 0\le t_\mu\le|r_\mu|,\quad t_\mu\in\mathbb Z\right\}
 \subset B(y).
\]
The $|S|!$ paths are distinct and shortest.

For a positive coarse bond $c=\langle c_-,c_-+Le_\nu\rangle$,
translation by $Le_\nu$ identifies
$\mathcal G(c_-,x)$ with $\mathcal G(c_+,x')$,
where $x'=x+Le_\nu$. If $\Gamma'$ is the translate of $\Gamma$,
the paths $\Gamma\cup[x,x']$ and $[c_-,c_+]\cup\Gamma'$
are rearrangements of at most $(5L-3)/2$ steps. Their common
subset-sum set is contained in $B(c_-)\cup B(c_+)$; in particular,
the straight paths have $L$ bonds and lie in this double block.
The straight path associated with the reversed bond is the
reverse of $[x,x']$.
\end{lemma}

\begin{proof}
Each ordering uses $|r_\mu|$ copies of the signed step in
direction $\mu$. Choosing a subset of these copies gives exactly
the displayed site set. Distinct orderings first differ by steps
in different coordinate directions and hence yield distinct paths.
The period of the fine lattice is
$N'=2L^{m+j+1}\ge2L^2$. A path from $y$ to $x$ has a signed
displacement congruent to $r_\mu$ modulo $N'$ in direction
$\mu$. If it used fewer than $|r_\mu|$ steps in that direction,
the difference from $r_\mu$ would be a nonzero multiple of $N'$
of absolute value less than $2|r_\mu|<N'$. This is impossible,
so the prescribed paths are shortest.

The translated endpoints have the same displacement vector,
which proves the correspondence of contour families. The two
concatenations place the $L$ steps in direction $e_\nu$ after
or before the same $q$ contour steps. Relative to $c_-$, every
coordinate of their common subset-sum set lies in $[-h,h]$ except the
$\nu$ coordinate, which lies in $[-h,L+h]$. These are exactly
the integer coordinate ranges of $B(c_-)\cup B(c_+)$. They
describe distinct sites because $N'\ge2L^2>2L$. The length is
$q+L\le(5L-3)/2$, and reversing the straight path reverses its
bonds in reverse order.
\end{proof}

\subsection{The average on a double block}
Continue to use the operator norm on $\mathbb C^{2\times2}$. Put
$G^c=\mathrm{SL}(2,\mathbb C)$, and
$\mathfrak g=\{A:A^*=A,\ \tr A=0\}$; its complexification
$\mathfrak g^c$ is the space of traceless complex matrices.
The series
\[
\log X=\sum_{r\ge1}\frac{(-1)^{r-1}}r(X-\mathbf1)^r
\qquad(|X-\mathbf1|<1)
\]
will be used near the identity. On $G\setminus\{-\mathbf1\}$
we also use the intrinsic logarithm, the unique traceless
skew-Hermitian logarithm whose eigenvalues lie in $i(-\pi,\pi)$.
These logarithms agree on their common domain. The matrix estimates
and logarithm properties needed here are proved in
\cref{sec:nonabelian-mean}.

For a finite family $\mathbf V=(V_1,\dots,V_n)$ of invertible
matrices, put
\begin{equation}\label{I-3-26}
d(\mathbf V)=\max_{i,j}|V_iV_j^{-1}-\mathbf1|.
\end{equation}
A mean is an invertible matrix $V$ satisfying
\begin{equation}\label{I-3-23}
\max_i|V_iV^{-1}-\mathbf1|<1,
\qquad \sum_i\log(V_iV^{-1})=0.
\end{equation}
For a family in $G$ or $G^c$, a mean in that group is required
to belong to the same group. We use the logarithmic mean in
Ba{\l}aban's presentation \cite[(0.10), p.~253]{B-CMP109},
where the equivalent definition is attributed to Federbush.

\begin{proposition}[The mean of a small-diameter family]
\label{I-3-28}
Let $n\ge1$, $\mathbf V\in\mathrm{GL}(2,\mathbb C)^n$ and
$d=d(\mathbf V)\le1/40$. On the ball $|Y|\le1/8$, the map
\begin{equation}\label{I-3-29}
 \Phi_{\mathbf V}(Y)
 =Y+\frac1n\sum_i\log(V_iV_1^{-1}e^{-Y})
\end{equation}
is defined, maps the ball into itself, and has Lipschitz constant
at most $2/5$. Its unique fixed point $Y_*$ satisfies $|Y_*|\le2d$.
The matrix
\begin{equation}\label{I-3-30}
 \operatorname{av}(\mathbf V)=e^{Y_*}V_1
\end{equation}
satisfies
\begin{equation}\label{I-3-31}
 \max_i|V_i\operatorname{av}(\mathbf V)^{-1}-\mathbf1|\le4d,
 \qquad
 \sum_i\log(V_i\operatorname{av}(\mathbf V)^{-1})=0.
\end{equation}
It is the only invertible solution of the mean equation for which
$\min_i|V_iV^{-1}-\mathbf1|\le1/10$ and
$\max_i|V_iV^{-1}-\mathbf1|<1$.
It is independent of the ordering of the family and equals $V$ on
the constant family $(V,\ldots,V)$.

If $V_i\in G^c$ for every $i$, then $Y_*\in\mathfrak g^c$ and
$\operatorname{av}(\mathbf V)\in G^c$.
If $V_i\in G$ for every $i$, then $Y_*\in i\mathfrak g$ and
$\operatorname{av}(\mathbf V)\in G$; in this case it is the only
mean in $G$ satisfying \eqref{I-3-23}, without the additional
$1/10$ condition. We denote this mean also by
$\operatorname{Av}(\mathbf V)$.

The maps $\mathbf V\mapsto Y_*$ and
$\mathbf V\mapsto\operatorname{av}(\mathbf V)$ are holomorphic
in the matrix entries on the open set
\[
 \mathcal O_n=\{\mathbf V\in\mathrm{GL}(2,\mathbb C)^n:
                    d(\mathbf V)<1/40\}.
\]
Their restrictions to $G^n\cap\mathcal O_n$ are real-analytic.
For $\mathbf V\in(G^c)^n\cap\mathcal O_n$ and $W,W'\in G$,
\begin{equation}\label{I-3-24}
 \operatorname{av}(W\mathbf VW')
 =W\operatorname{av}(\mathbf V)W'.
\end{equation}
\end{proposition}

The proof in \cref{sec:nonabelian-mean} constructs the mean by
a contraction with constant $2/5$, uniform in the number of matrices.
The uniqueness statement makes the average independent of the matrix
chosen as the reference point in that construction.

Fix $0\le k<K$ and a coarse bond $c$ of $\mathbb T^{(k+1)}$.
Let $\mathsf P_c$ be the bonds of $\mathbb T^{(k)}$ whose
endpoints lie in $B(c_-)\cup B(c_+)$, and let $\mathsf P_c^+$
be its positive bonds. Write $U[c]=U([c_-,c_+])$.
For $0<\omega\le1$, set
\[
\mathcal S_c(\omega)=\{U\in\mathcal V_k:
 |\partial U(p)-\mathbf1|<\omega\text{ for every plaquette
 with all corners in }B(c_-)\cup B(c_+)\}.
\]
Here $\partial U(p)=U_p$. We fix
\begin{equation}\label{I-3-34}
\omega_M=\frac1{200L^2},\qquad
\mathcal S_c=\mathcal S_c(\omega_M).
\end{equation}
Whenever the contour family has a unique mean in $G$, define
\begin{equation}\label{I-3-48}
U(y,x)=\operatorname{Av}(\{U(\Gamma):\Gamma\in\mathcal G(y,x)\}),
\qquad U(x,y)=U(y,x)^{-1}.
\end{equation}
The notation $\operatorname{Av}$ denotes that unique group mean;
on families of diameter at most $1/40$ it equals
$\operatorname{av}$. This is the averaged contour variable
of \cite[(0.11), p.~253]{B-CMP109}.

\begin{lemma}\label{I-3-37}
Let $0\le k<K$ and $0<\omega\le\omega_M$.
For $U\in\mathcal S_c(\omega)$, $y\in\{c_-,c_+\}$ and
$x\in B(y)$, the contour family has diameter at most $1/160$.
Its unique mean satisfies
\begin{equation}\label{I-3-40}
|U(\Gamma)U(y,x)^{-1}-\mathbf1|\le1/40
\quad(\Gamma\in\mathcal G(y,x)).
\end{equation}
If $c$ is positive and $x'=x+(c_+-c_-)$, then
\begin{equation}\label{I-3-41}
W_x=U(c_-,x)U([x,x'])U(c_+,x')^{-1}U[c]^{-1}
\end{equation}
belongs to $G$ and satisfies $|W_x-\mathbf1|<1/14$.
The sets $\mathcal S_c(\omega)$ are open, and the contour means
are real-analytic there.
\end{lemma}

\begin{proof}
Two contour paths from $y$ to $x$ have rearranged step sequences
of common length $r\le3(L-1)/2$. All subset sums of these steps
stay in $B(y)$. Thus \cref{I-3-7} gives
\begin{equation}\label{I-3-38}
|U(\Gamma')U(\Gamma)^{-1}-\mathbf1|
\le(1+\omega)^{9L^2/8}-1
\le e^{9/1600}-1<1/160.
\end{equation}
The mean assertion and \eqref{I-3-40} follow from
\cref{I-3-28}.

Choose $\Gamma\in\mathcal G(c_-,x)$ and let $\Gamma'$ be its
translate in $\mathcal G(c_+,x')$. The paths
$\Gamma\cup[x,x']$ and $[c_-,c_+]\cup\Gamma'$ have rearranged
step sequences of length at most $(5L-3)/2$. Their subset sums
stay in the double block. Applying \cref{I-3-7} again gives
\begin{equation}\label{I-3-42}
P_x=U(\Gamma)U([x,x'])U(\Gamma')^{-1}U[c]^{-1},
\qquad |P_x-\mathbf1|\le e^{1/64}-1<1/60.
\end{equation}
Write
$U(c_-,x)=(\mathbf1+E_1)U(\Gamma)$ and
$U(c_+,x')^{-1}=U(\Gamma')^{-1}(\mathbf1+E_2)$.
Equation \eqref{I-3-40} and unitarity give $|E_1|,|E_2|\le1/40$.
Consequently
\[
W_x=(\mathbf1+E_1)P_x(\mathbf1+U[c]E_2U[c]^{-1}),
\]
and expansion of this product gives
\begin{equation}\label{I-3-43}
|W_x-\mathbf1|
\le\left(\frac{41}{40}\right)^2\frac{61}{60}-1
=\frac{6541}{96000}<\frac1{14}.
\end{equation}

Each plaquette transport is continuous, so its finitely many
strict curvature inequalities define an open set. Path transports
extend holomorphically to invertible complex bond variables.
Their contour families depend holomorphically on these variables,
and the condition $d<1/40$ is open. Composition with the mean
of \cref{I-3-28} therefore gives a holomorphic extension of each
contour mean near every point of $\mathcal S_c(\omega)$.
On the real configurations the values lie in $G$, proving the
claimed real-analyticity.
\end{proof}

\subsection{The global real map and its analytic branch}
For a positive coarse bond $c$, define $\mathcal M(U)(c)$ by
the following three cases, and set
$\mathcal M(U)(\bar c)=\mathcal M(U)(c)^{-1}$.
On $\mathcal S_c$, put
\begin{equation}\label{I-3-49}
\mathcal M(U)(c)=
\exp\left(L^{-3}\sum_{x\in B(c_-)}\log W_x\right)U[c].
\end{equation}
The logarithms and means in this formula exist by
\cref{I-3-37}. It is the symmetrized average
of \cite[(0.12), p.~254]{B-CMP109}, in dimension three.

For $x\in B(c_-)$, $\Gamma\in\mathcal G(c_-,x)$ and
$\Gamma'\in\mathcal G(c_+,x')$, let
\[
Y_{x,\Gamma,\Gamma'}=
U(\Gamma)U([x,x'])U(\Gamma')^{-1}U[c]^{-1},
\]
and let $\mathcal N_c$ be the set where at least one of these
matrices equals $-\mathbf1$. On
$\mathcal V_k\setminus(\mathcal S_c\cup\mathcal N_c)$, put
\begin{equation}\label{I-3-50}
\begin{split}
\mathcal M(U)(c)&=\exp\left(L^{-3}\sum_{x\in B(c_-)}
                              H_x(U)\right)U[c],\\
H_x(U)&=\frac{\displaystyle
 \sum_{\Gamma\in\mathcal G(c_-,x)}
 \sum_{\Gamma'\in\mathcal G(c_+,x')}
 \log Y_{x,\Gamma,\Gamma'}}
 {|\mathcal G(c_-,x)|\,|\mathcal G(c_+,x')|}.
\end{split}
\end{equation}
Here the logarithm is the intrinsic logarithm on
$G\setminus\{-\mathbf1\}$; this formula follows
\cite[(0.4), p.~253]{B-CMP109}. On the remaining set
$\mathcal N_c\setminus\mathcal S_c$, put $\mathcal M(U)(c)=U[c]$.
Each exponent is traceless and skew-Hermitian, so all three
formulas take values in $G$.

For complex configurations, we use \eqref{I-3-49} on its analytic
domain. More precisely, let $\mathcal O_c$ be the configurations
in $(G^c)^{\mathsf P_c^+}$ whose contour families all have diameter
less than $1/40$. Define their contour means using
$\operatorname{av}$, and let
\[
\mathcal O'_c=\{\mathbf U\in\mathcal O_c:
                   |W_x(\mathbf U)-\mathbf1|<1
                   \text{ for all }x\in B(c_-)\}.
\]
On this open set, define
\begin{equation}\label{I-3-53}
\mathcal M(\mathbf U)(c)=
\exp\left(L^{-3}\sum_{x\in B(c_-)}\log W_x(\mathbf U)\right)
\mathbf U[c].
\end{equation}
Define the value on a reversed coarse bond by inversion.
The map is holomorphic: transports, inversion, the contour means and
the logarithmic series are holomorphic in the positive-bond
coordinates. The logarithms of the $G^c$-valued matrices have
trace zero, so \eqref{I-3-53} lies in $G^c$.
The restriction of every $U\in\mathcal S_c$ lies in
$\mathcal O'_c$, where \eqref{I-3-53} agrees with
\eqref{I-3-49}. For $U_0\in\mathcal S_c$, assign traceless complex
matrices $A(b)$ on $\mathsf P_c^+$ and define
$(e^{iA}U_0)(b)=e^{iA(b)}U_0(b)$ there, with the values on reversed
bonds defined by inversion. Then
$A\mapsto\mathcal M(e^{iA}U_0)(c)$ is holomorphic near zero.

\begin{proposition}\label{I-3-55}\label{I-3-58}
For $0\le k<K$, the map
$\mathcal M:\mathcal V_k\to\mathcal V_{k+1}$ is Borel measurable.
At a positive bond $c$ it is continuous on each of the three
sets defining it, and real-analytic on $\mathcal S_c$.
Its value and the applicable case depend only on $U|_{\mathsf P_c}$.
For every $G$-valued gauge transformation $u$,
\begin{equation}\label{I-3-56}
\mathcal M(U^u)=\mathcal M(U)^{u|_{k+1}},
\end{equation}
and the applicable case is unchanged. For complex configurations,
$\mathcal O'_c$ is invariant under such gauge transformations,
and the same covariance holds on that domain.
\end{proposition}

\begin{proof}
The set $\mathcal S_c$ is open. Each $Y_{x,\Gamma,\Gamma'}$
is continuous and there are finitely many of them, so
$\mathcal N_c$ is closed. The three defining sets are therefore
Borel. On the first, analyticity follows from the complex extension.
On the second, continuity follows from continuity of the intrinsic
logarithm away from $-\mathbf1$; on the third the formula is the
path product $U[c]$. The preimage of an open subset of $G$ is thus
a finite union of Borel sets. The finite product of these bond maps
is Borel as a map to $\mathcal V_{k+1}$.

All paths and plaquettes used at $c$ have their bonds in
$\mathsf P_c$, proving locality. For a path from $y$ to $x$,
cancellation of the internal gauge factors gives
\[
U^u(\Gamma)=u(y)U(\Gamma)u(x)^{-1}.
\]
Equation \eqref{I-3-24} therefore gives
$U^u(y,x)=u(y)U(y,x)u(x)^{-1}$ wherever the contour mean
is used. The quantities $W_x$ and $Y_{x,\Gamma,\Gamma'}$
transform by conjugation by $u(c_-)$, and each plaquette
transforms by conjugation at its base point. Unitary conjugation
preserves the operator norm, the equation $Y=-\mathbf1$, and
both logarithms. It follows that $\mathcal S_c$, $\mathcal N_c$
and all three cases are gauge invariant. Conjugating the exponent
and using $U^u[c]=u(c_-)U[c]u(c_+)^{-1}$ proves
\eqref{I-3-56}. The value at the reversed bond follows by inversion.
The same calculation for complex configurations preserves the
contour diameters and $|W_x-\mathbf1|$, proving invariance of
$\mathcal O'_c$ and covariance there.
\end{proof}

\subsection{Partial configurations and iteration}
If $U$ is given on a reversal-invariant set $\mathcal B$ of fine
bonds, its average is defined at precisely the coarse bonds $c$
with $\mathsf P_c\subset\mathcal B$, using any extension of $U$
to the whole lattice. Locality makes this independent of the
extension. Define $\mathcal M^n$ by composing the maps one level
at a time, with $\mathcal M^0=\mathrm{id}$.

Let $0\le j\le j+n\le K$, and let $c$ be a bond of
$\mathbb T^{(j+n)}$. Denote by $\mathsf D_n(c)$ the bonds of
$\mathbb T^{(j)}$ with both endpoints in
$B^n(c_-)\cup B^n(c_+)$. Then
$(\mathcal M^nU)(c)$ is defined exactly when
$\mathsf D_n(c)\subset\mathcal B$, and depends only on that
restriction. To prove this by induction, the required geometric
identity is
\begin{equation}\label{eq:determining-bonds-recursion}
\mathsf D_{n+1}(c)=\bigcup_{c'\in\mathsf P_c}\mathsf D_n(c').
\end{equation}
The nesting of blocks gives inclusion from right to left. Conversely,
each fine bond in the left side either joins two order-$n$ blocks
with adjacent centers in the double block of $c$, or lies within
one such block. In the first case it belongs to
$\mathsf D_n(c')$ for the bond between those centers. In the second,
choose a neighbor of that center within the double block; the bond
to this neighbor provides the required $c'$. Every site of a double
block has such a neighbor. This proves the identity and the induction,
starting with $\mathsf D_0(c)=\{c,\bar c\}$.

Iteration of \eqref{I-3-56} gives, for a bare gauge transformation,
\[
\mathcal M^n(U^u)=(\mathcal M^nU)^{u|_n}.
\]
The identity also holds for partial configurations on their domains,
by extending a configuration and using locality. For a region made
of closed coarse cubes, all fine sites in the block of a coarse site
of the region also lie in the region. Hence these partial averages
are defined on every inner coarse bond when the original field is
given on the inner fine bonds. Here an inner bond has both endpoints
in the region.

Complex iteration uses \eqref{I-3-53} at every participating bond.
For any finite family of iterates, the common domain is open and
the iterates are holomorphic there. Indeed, the domain of a single
step is open in $(G^c)^{\mathsf P_c^+}$. Inductively,
intersect the domains of the previous-level values on $\mathsf P_c$
and take the inverse image of $\mathcal O'_c$ under their joint
holomorphic map. On this open set the next iterate is its composition
with \eqref{I-3-53}. Applying the same strict inequalities and
formulas in the ambient invertible positive-bond coordinates gives
an open set on which the iterates are holomorphic. The actual
domain is its intersection with the product of copies of $G^c$.
The complex and real iterates agree whenever every
participating real configuration lies in the corresponding
$\mathcal S_c$.

\subsection{Linearization}
Fix $0\le k<K$ and a positive bond $c$ of $\mathbb T^{(k+1)}$.
Regard a bond field $A$ as an antisymmetric function on bonds,
$A(\bar b)=-A(b)$, with values in $\mathfrak g^c$, and write
$A(\Gamma)$ for its sum along a path. Put
\[
\bar A(y,x)=\frac1{|\mathcal G(y,x)|}
                   \sum_{\Gamma\in\mathcal G(y,x)}A(\Gamma).
\]
The derivative as a map on bond variables is
\begin{equation}\label{I-3-52}
(\mathbf Q^{\mathrm{lin}}A)(c)=L^{-3}\sum_{x\in B(c_-)}
 [\bar A(c_-,x)+A([x,x'])-\bar A(c_+,x')].
\end{equation}
When fine bonds have length one and coarse bonds have length $L$,
the normalized operator on one-form components is
\begin{equation}\label{I-3-51}
Q=L^{-1}\mathbf Q^{\mathrm{lin}}.
\end{equation}
It fixes constant components: the two contour sums cancel by
translation, and $A([x,x'])=LA_\nu$ at a bond in direction $\nu$.

\begin{lemma}\label{I-3-59}
Suppose $t\mapsto\mathbf U_t$ is a curve of $G^c$-valued
configurations on $\mathsf P_c$, differentiable at zero, with
$\mathbf U_0=\mathbf1$ and $\dot{\mathbf U}_0(b)=iA(b)$.
For small $|t|$ the complex average is defined, and
\begin{equation}\label{I-3-60}
\mathcal M(\mathbf1)(c)=\mathbf1,\qquad
\left.\frac d{dt}\mathcal M(\mathbf U_t)(c)\right|_{t=0}
=i(\mathbf Q^{\mathrm{lin}}A)(c).
\end{equation}
For a curve on $\mathsf D_n(c)$ the same assertion holds for
the $n$th iterate, with the composition
$\mathbf Q^{\mathrm{lin}}_{j\to j+n}$ on the right.
For real curves these averages agree, near zero, with the global
real averages.
\end{lemma}

\begin{proof}
The identity configuration lies in $\mathcal O'_c$, and all its
contour means and relative loop variables equal $\mathbf1$.
Openness gives the first assertion. A path product has derivative
$iA(\Gamma)$ at zero. Differentiating the equation defining the
contour mean, whose value is $\mathbf1$ at zero, gives
\[
0=\sum_{\Gamma\in\mathcal G(y,x)}
   [iA(\Gamma)-\dot{\mathbf U}_0(y,x)],
\qquad \dot{\mathbf U}_0(y,x)=i\bar A(y,x).
\]
Here the differential of $\log$ at $\mathbf1$ is the identity.
The derivative of $W_x$, and of $\log W_x$, is therefore
\[
i[\bar A(c_-,x)+A([x,x'])-\bar A(c_+,x')-A([c_-,c_+])].
\]
After averaging over the $L^3$ sites, the exponent has derivative
$i[\mathbf Q^{\mathrm{lin}}A(c)-A([c_-,c_+])]$.
The final factor $\mathbf U_t[c]$ contributes the opposite
reference-path term, proving \eqref{I-3-60}.
For iterates apply this calculation successively on the finitely
many determining bonds and use the chain rule. Each intermediate
configuration equals the identity at zero. For real curves,
continuity puts every intermediate plaquette inside its fixed
small-curvature bound for sufficiently small $|t|$, so the real
and complex branches agree throughout the finite iteration.
\end{proof}

There is one useful coefficient in \eqref{I-3-52}. The corridor
$B(c)$ consists of the fine bonds directed from $B(c_-)$ to $B(c_+)$.
It has $L^2$ members, all directed as $c$. In fine-lattice units
its initial sites have $\nu$-coordinate $(L-1)/2$ relative to $c_-$
and transverse coordinates between $-(L-1)/2$ and $(L-1)/2$.
Contour paths stay within a block, so only the straight paths
$[x,x']$ contribute to a corridor bond. Exactly the $L$ starting
sites on its longitudinal line contribute, each once. Thus, writing
$\mathbf Q^{\mathrm{lin}}A(c)=\sum_b\mathsf w_c(b)A(b)$
over positive fine bonds, its unique real coefficients vanish
outside $\mathsf P_c$ and satisfy
\begin{equation}\label{I-3-62}
\mathsf w_c(b)=L^{-2}\quad(b\in B(c)),\qquad
\sum_{b\in B(c)}\mathsf w_c(b)=1.
\end{equation}

\subsection{Image measures and block fields}
For $0\le k<K$ and a finite complex Borel measure $\nu$ on
$\mathcal V_k$, define
$(\mathbf T\nu)(E)=\nu(\mathcal M^{-1}E)$ for Borel
$E\subset\mathcal V_{k+1}$.

\begin{lemma}[Image measures]\label{I-4-67}
The measure $\mathbf T\nu$ is finite and complex, depends linearly
on $\nu$, preserves its total mass, and satisfies
\begin{equation}\label{I-3-64}
\int F\,d(\mathbf T\nu)=\int F(\mathcal M U)\,d\nu
\end{equation}
for every bounded Borel $F$ on $\mathcal V_{k+1}$, with
$|\mathbf T\nu|\le\mathbf T|\nu|$.
For $0\le j\le l\le K$ and $\nu$ on $\mathcal V_j$,
$\mathbf T^{l-j}\nu$ is the image under $\mathcal M^{l-j}$, so these iterates preserve mass and
compose according to their levels. In particular, for $\nu$ on
$\mathcal V_0$,
\begin{equation}\label{I-3-65}
\int F\,d(\mathbf T^k\nu)=\int F(\mathcal M^kU)\,d\nu,
\qquad 0\le k\le K.
\end{equation}
\end{lemma}

\begin{proof}
The map $\mathcal M$ is Borel. The integral formula follows
from indicators by bounded approximation; taking total variations
over Borel partitions gives $|\mathbf T\nu|\le\mathbf T|\nu|$.
Linearity, mass preservation and iteration follow from the
corresponding properties of inverse images.
\end{proof}

Positivity is preserved. If $\nu$ is gauge invariant, extend
a coarse gauge transformation $v$ by the identity to a fine
gauge $u$. Then $\mathcal M(U^u)=(\mathcal M U)^v$ and
\eqref{I-3-64} give gauge invariance of $\mathbf T\nu$.
For a probability measure $\nu$ on $\mathcal V_0$,
$\mathbf T^k\nu$ is the law
of $\mathcal M^kU$ when $U$ has law $\nu$; in particular the
Wilson measure gives a gauge-invariant probability measure
at each coarse level. A representation by a Haar density
requires a further argument.


\section{The nonabelian mean}\label{sec:nonabelian-mean}

Writing $T_i=V_iV_1^{-1}$, the linear terms in
$Y+\log(T_ie^{-Y})$ cancel when $T_i=\mathbf1$.
The following estimates quantify the error when $T_i$ is close to
$\mathbf1$. They also give the uniqueness of the mean on $G$.
At the end of the section, \cref{lem:unique-group-mean} establishes
the Borel and covariance properties of arbitrary unique group
means needed for gauge normalization in \cref{sec:gauge-lifts}.

\begin{lemma}[Exponential and logarithm estimates]\label{I-3-15}
For complex $2\times2$ matrices, the following hold.
If $|Z|,|Z'|\le r$, then
\begin{align*}
 |e^Z|&\le e^{|Z|}, & |e^Z-\mathbf1|&\le e^{|Z|}-1,\\
 |e^Z-e^{Z'}|&\le e^r|Z-Z'|,
 &|e^Z-e^{Z'}-(Z-Z')|&\le(e^r-1)|Z-Z'|.
\end{align*}
Also $|We^Z-\mathbf1|\le(1+|W-\mathbf1|)e^{|Z|}-1$.
For $|E|<1$,
\[
 |\log(\mathbf1+E)|\le-\ln(1-|E|)\le\frac{|E|}{1-|E|},
 \qquad |(\mathbf1+E)^{-1}-\mathbf1|
 \le\frac{|E|}{1-|E|},
\]
and $e^{\log(\mathbf1+E)}=\mathbf1+E$.
For $|Z|<\ln2$, one has $|e^Z-\mathbf1|<1$ and
$\log(e^Z)=Z$.

If $|A-\mathbf1|,|B-\mathbf1|\le\varepsilon<1$, then
\begin{equation}\label{I-3-16}
 \log A-\log B
 =\int_0^1(\mathbf1+t(A-\mathbf1))^{-1}(A-B)
                (\mathbf1+t(B-\mathbf1))^{-1}\,dt
\end{equation}
and
\begin{equation}\label{I-3-17}
 \begin{split}
 |\log A-\log B-(A-B)|&\le
       \frac{\varepsilon}{1-\varepsilon}|A-B|,\\
 |\log A-\log B|&\le\frac1{1-\varepsilon}|A-B|.
 \end{split}
\end{equation}
For unitary $A,B$ with $|A-\mathbf1|,|B-\mathbf1|<1$,
the sharper uniform estimate is
\[
 |\log A-\log B|\le2|A-B|.
\]
Finally, $\det e^Z=e^{\operatorname{tr}Z}$.
If $X\in G^c$ and $|X-\mathbf1|<1$, then
$\log X\in\mathfrak g^c$; if $X\in G$ as well, then
$\log X\in i\mathfrak g$.
\end{lemma}

\begin{proof}
The operator norm is submultiplicative and unchanged by unitary
multiplication. All power series used below converge absolutely in
this norm. Multiplying the exponential series gives
$e^Ze^{-Z}=\mathbf1$; conjugation commutes termwise with both the
exponential and the logarithm whenever the latter series converges
before and after conjugation. The identity
\[
 Z^q-Z'^q=\sum_{j=0}^{q-1}Z^j(Z-Z')Z'^{q-1-j}
\]
bounds the $q$th difference by $qr^{q-1}|Z-Z'|$.
Summing the exponential series, and then omitting its linear term,
gives the exponential estimates. The estimate with $W$ follows by
writing $We^Z-\mathbf1=(W-\mathbf1)e^Z+(e^Z-\mathbf1)$.
The logarithmic and inverse bounds follow from their power series
and the Neumann series.

For completeness, the inverse identities can be checked along
$0\le t\le1$. Put $g(t)=\log(\mathbf1+tE)$ and
$R(t)=(\mathbf1+tE)^{-1}$. Termwise differentiation is justified
by $|E|<1$, and gives $g'=ER$, $R'=-ER^2$.
These matrices commute, being limits of polynomials in $E$.
It follows that $(e^{g(t)}R(t))'=0$, and its value at $0$ is
$\mathbf1$. Thus $e^{\log(\mathbf1+E)}=\mathbf1+E$.
If $|Z|<\ln2$, the bound
$|e^{tZ}-\mathbf1|\le e^{|Z|}-1<1$ permits the same
differentiation of the logarithmic series; it yields
$(\log(e^{tZ}))'=Z$ and hence $\log(e^Z)=Z$.

Set $E=A-\mathbf1$, $F=B-\mathbf1$,
$R_t=(\mathbf1+tE)^{-1}$, and
$S_t=(\mathbf1+tF)^{-1}$.
Integrating the Neumann series gives
$\log A=\int_0^1 ER_t\,dt$, and similarly for $B$.
Since $E$ commutes with $R_t$,
\[
 ER_t-FS_t=R_t\bigl(E(\mathbf1+tF)
                     -(\mathbf1+tE)F\bigr)S_t
          =R_t(A-B)S_t,
\]
which proves \eqref{I-3-16}.
Moreover,
\[
 |R_t|,|S_t|\le(1-t\varepsilon)^{-1},\qquad
 |R_t-\mathbf1|,|S_t-\mathbf1|
       \le t\varepsilon(1-t\varepsilon)^{-1}.
\]
The decomposition
$R_t(A-B)S_t-(A-B)=(R_t-\mathbf1)(A-B)S_t
 +(A-B)(S_t-\mathbf1)$ therefore bounds its norm by
$[(1-t\varepsilon)^{-2}-1]|A-B|$.
Integration, using
$\int_0^1(1-t\varepsilon)^{-2}\,dt=(1-\varepsilon)^{-1}$,
gives both estimates in \eqref{I-3-17}.

When $A$ is unitary and $|A-\mathbf1|<1$, each eigenvalue
$e^{i\theta}$ has $\cos\theta>1/2$.
Unitary diagonalization gives
\[
 |(\mathbf1+t(A-\mathbf1))^{-1}|\le(1-t/2)^{-1}.
\]
The same bound applies to $B$, so \eqref{I-3-16} gives
$|\log A-\log B|\le|A-B|\int_0^1(1-t/2)^{-2}\,dt
=2|A-B|$.

Triangularizing $Z$ proves $\det e^Z=e^{\operatorname{tr}Z}$.
If $|X-\mathbf1|<1$, its two eigenvalues lie in the disc
$|\lambda-1|<1$, where the scalar logarithm has imaginary part
in $(-\pi/2,\pi/2)$.
The trace of the matrix logarithm is the sum of these scalar
logarithms, as follows by taking the trace termwise in its series.
For $\det X=1$, this sum belongs to $2\pi i\mathbb Z$ and
has imaginary part in $(-\pi,\pi)$, so it is zero.
If $X$ is also unitary, unitary diagonalization shows that its
logarithm is skew-Hermitian. Thus it belongs to $i\mathfrak g$.
\end{proof}

The exponential and the local logarithm are holomorphic in the
matrix entries by local uniform convergence of their power series;
products and inversion on $\mathrm{GL}(2,\mathbb C)$ are
holomorphic as well. Local uniform limits preserve holomorphy
by Cauchy's formula in each coordinate. These facts apply in the
charts $A\mapsto e^AX_0$ of $G^c$. A matrix-valued holomorphic
map with values in $G^c$ is therefore holomorphic into $G^c$;
its restriction to $G$ is real-analytic when it takes values in $G$.

For the definition of the averaging map away from the small-field
set, we also use the intrinsic logarithm on $G\setminus\{-\mathbf1\}$.

\begin{lemma}[The intrinsic logarithm]\label{I-3-13}
Every $U$ in this set has a unique representation $U=e^{iH}$ with
$H\in\mathfrak g$ and $|H|<\pi$. Put $\log U=iH$.
Writing $s(\theta)=\sin\theta/\theta$, with $s(0)=1$, this gives
\[
 H=\frac{U-U^{-1}}{2i\,s(\arccos(\operatorname{tr}U/2))}.
\]
In particular, the intrinsic logarithm is continuous, agrees with
the logarithmic series when $|U-\mathbf1|<1$, and satisfies
\[
 \log(WUW^{-1})=W(\log U)W^{-1},\qquad
 \log(U^{-1})=-\log U\qquad(W\in G).
\]
\end{lemma}

\begin{proof}
A trace-zero Hermitian $H$ satisfies
$H^2=|H|^2\mathbf1$, so
$e^{iH}=\cos|H|\,\mathbf1+i\,s(|H|)H$.
The characteristic equation of $U\in G$ gives
$U+U^{-1}=(\operatorname{tr}U)\mathbf1$.
Thus $U=t\mathbf1+iJ$, where
$t=\operatorname{tr}U/2\in\mathbb R$, $J\in\mathfrak g$, and
$t^2+|J|^2=1$.
For $U\ne-\mathbf1$, set $\theta=\arccos t\in[0,\pi)$.
Then $H=J/s(\theta)$ satisfies $|H|=\theta$ and $e^{iH}=U$.
Conversely the exponential formula determines $\theta$ from the
trace and then determines $H$, proving uniqueness and the displayed
formula. Continuity follows from this formula, while conjugation
and inversion follow from uniqueness.
For $|U-\mathbf1|<1$, unitary diagonalization shows that the
eigenphases lie in $(-\pi/3,\pi/3)$; the logarithmic series has
these eigenphases multiplied by $i$, and therefore gives the same
logarithm.
\end{proof}

The numerical bounds follow by substituting the indicated rational
arguments into
\[
 e^t\le1+t+\frac{t^2}{2}
          +\frac{t^3}{6(1-t/4)}\qquad(0\le t<4),
\]
and using
$-\ln(1-x)\le x+x^2/2+x^3/(3(1-x))$ for $0\le x<1$.

\begin{proof}[Proof of \cref{I-3-28}]
Write $T_i=V_iV_1^{-1}$, so that $|T_i-\mathbf1|\le d$
and $|T_i|\le1+d$. For $|Y|\le1/8$ put $M_i(Y)=T_ie^{-Y}$.
Then
\[
 |M_i(Y)-\mathbf1|\le(1+d)e^{|Y|}-1
 \le\frac{41}{40}e^{1/8}-1<\frac16.
\]
In particular, all logarithms in \eqref{I-3-29} are defined.
For $\Delta=Y-Y'$, the exponential estimate gives
\[
 e^{-Y}-e^{-Y'}=-\Delta+\sigma,
 \qquad |\sigma|\le(e^{1/8}-1)|\Delta|.
\]
By \eqref{I-3-17} with $\varepsilon=1/6$,
\[
 \log M_i(Y)-\log M_i(Y')
 =M_i(Y)-M_i(Y')+\rho_i,
 \qquad
 |\rho_i|\le\frac15(1+d)e^{1/8}|\Delta|.
\]
Consequently,
\[
 \Delta+\log M_i(Y)-\log M_i(Y')
 =-(T_i-\mathbf1)\Delta+T_i\sigma+\rho_i.
\]
Averaging yields the contraction estimate
\begin{equation}\label{I-3-32}
 \begin{split}
 |\Phi_{\mathbf V}(Y)-\Phi_{\mathbf V}(Y')|
 &\le\left[d+(1+d)(e^{1/8}-1)
                  +\frac15(1+d)e^{1/8}\right]|Y-Y'|\\
 &\le\frac25|Y-Y'|.
 \end{split}
\end{equation}
Here $e^{1/8}<1.1332$ makes the bracket less than $2/5$.
At the origin,
$|\Phi_{\mathbf V}(0)|\le d/(1-d)\le1/39$; hence
\[
 |\Phi_{\mathbf V}(Y)|\le\frac1{39}+\frac25\frac18
 =\frac1{39}+\frac1{20}<\frac18.
\]
The ball is complete, so the contraction principle gives its
unique fixed point. The same estimates give
\[
 |Y_*|\le\frac53\frac d{1-d}\le2d.
\]

Since $e^{Y_*}$ is invertible with inverse $e^{-Y_*}$,
$V_i\operatorname{av}(\mathbf V)^{-1}=M_i(Y_*)$.
The fixed-point equation is the logarithmic equation in
\eqref{I-3-31}. Its bound follows from
\[
 |M_i(Y_*)-\mathbf1|
 \le(1+d)e^{2d}-1\le e^{3d}-1
 \le3de^{3d}\le4d,
\]
where $e^{3/40}<1.078$.

Suppose an invertible $V'$ satisfies the two mean conditions and
$|V_1V'^{-1}-\mathbf1|\le1/10$.
For $X=V'V_1^{-1}$, the inverse estimate gives
$|X-\mathbf1|\le1/9$. Thus
\[
 Y=\log X,\qquad e^Y=X,\qquad
 |Y|\le\ln(9/8)<1/8.
\]
The mean equation says $\Phi_{\mathbf V}(Y)=Y$.
It follows that $Y=Y_*$ and $V'=\operatorname{av}(\mathbf V)$.
If the $1/10$ condition holds at $V_j$ instead, apply this
argument first with $V_j$ numbered first. The mean constructed
with that numbering lies within $4d\le1/10$ of every member,
so the argument with $V_1$ numbered first identifies the two
constructions. This also proves independence of the numbering.
For a constant family, $T_i=\mathbf1$ and
$\Phi_{\mathbf V}(Y)=Y+\log(e^{-Y})=0$ on the ball; hence
$Y_*=0$ and its mean is $V_1$.

If every $V_i\in G^c$, then the ball's intersection with
$\mathfrak g^c$ is preserved by $\Phi_{\mathbf V}$:
for trace-zero $Y$, each $T_ie^{-Y}$ has determinant one and
its logarithm has trace zero. If every $V_i\in G$, the same
argument preserves the intersection with $i\mathfrak g$,
since $e^{-Y}$ and $T_ie^{-Y}$ are unitary there and their
logarithms are skew-Hermitian. Starting the contraction at zero
proves the assertions about $Y_*$ and the mean in both cases.

Now suppose the family and an arbitrary admissible mean $V'$
lie in $G$. Set $U_i=V_iV'^{-1}$. Then
$|U_i-U_1|=|(T_i-\mathbf1)U_1|\le d$, and the unitary
logarithm estimate gives $|\log U_i-\log U_1|\le2d$.
Since $\sum_i\log U_i=0$,
\[
 |\log U_1|
 =\left|\frac1n\sum_i(\log U_1-\log U_i)\right|
 \le2d\le\frac1{20}.
\]
It follows that
$|U_1-\mathbf1|\le e^{1/20}-1<1/10$.
The preceding local uniqueness argument therefore applies to
every admissible mean in $G$.

The diameter is continuous on $\mathrm{GL}(2,\mathbb C)^n$,
so $\mathcal O_n$ is open. On
$\mathcal O_n\times\{|Y|<1/5\}$,
\[
 |V_iV_1^{-1}e^{-Y}-\mathbf1|
 <\frac{41}{40}e^{1/5}-1<1.
\]
Thus $\Phi_{\mathbf V}(Y)$ is jointly holomorphic there.
The iterates $Y_0=0$, $Y_{q+1}=\Phi_{\mathbf V}(Y_q)$
are holomorphic on $\mathcal O_n$ and remain in $|Y|\le1/8$.
Their convergence is uniform on that whole open set, because
\[
 |Y_q(\mathbf V)-Y_*(\mathbf V)|
 \le\frac18\left(\frac25\right)^q.
\]
Hence $Y_*$ and $e^{Y_*}V_1$ are holomorphic, with the stated
real-analytic restrictions.

Finally, unitary left and right multiplication leave the diameter
unchanged. They replace each relative matrix by
\[
 (WV_iW')(WVW')^{-1}=W(V_iV^{-1})W^{-1}.
\]
Unitary conjugation preserves its distance from $\mathbf1$ and
commutes with the logarithmic series. Thus
$W\operatorname{av}(\mathbf V)W'$ is a mean of
$W\mathbf VW'$ within $4d\le1/10$ of every member.
The local uniqueness already proved gives \eqref{I-3-24}.
\end{proof}

\begin{lemma}[Unique group means]\label{lem:unique-group-mean}
Let $n\ge1$.
The set $\mathsf E_n\subset G^n$ of families having exactly one
group mean is Borel, and $\operatorname{Av}:\mathsf E_n\to G$
is Borel. Consequently the domain of each averaged contour
variable $U(y,x)$ is Borel, and that variable is Borel on its domain.
These contour assertions hold for configurations on any reversal-invariant
bond set containing all bonds of $B(y)$.

The set $\mathsf E_n$ is invariant under left and right
multiplication of the family by fixed elements $W,W'\in G$, and
\[
 \operatorname{Av}(W\mathbf VW')=W\operatorname{Av}(\mathbf V)W'.
\]
\end{lemma}

\begin{proof}
For $q\ge2$, let
\[
 K_q=\left\{(\mathbf V,V)\in G^n\times G:
 \max_i|V_iV^{-1}-\mathbf1|\le1-q^{-1},\quad
 \sum_i\log(V_iV^{-1})=0\right\}.
\]
These are increasing compact sets: the logarithms are continuous
on the indicated closed subdomain of their open domain.
Every admissible mean belongs to one of these sets. Also set
\[
 K'_q=\{(\mathbf V,V,V'):(\mathbf V,V),(\mathbf V,V')\in K_q,
                                      \ |V-V'|\ge q^{-1}\}.
\]
These sets are compact as well. Projection onto $G^n$ maps each
of them to a compact set. Thus the families with at least one
mean form $\bigcup_q\operatorname{pr}K_q$, and those with at
least two form $\bigcup_q\operatorname{pr}K'_q$. The latter
identity follows by taking $q$ large enough for both distinct
means and their positive separation. Their difference is
$\mathsf E_n$, hence is Borel. For closed $F\subset G$,
\[
 \operatorname{Av}^{-1}(F)=\mathsf E_n\cap
 \bigcup_q\operatorname{pr}\bigl(K_q\cap(G^n\times F)\bigr)
\]
is Borel. Finally, the contour transports form a continuous map
from the configuration space to $G^n$, depending only on the
bonds of $B(y)$. Taking the inverse image of $\mathsf E_n$
and composing with $\operatorname{Av}$ proves the assertions
about the contour variables.

For covariance, the map $V\mapsto WVW'$ is a bijection between
the admissible group means of $\mathbf V$ and those of
$W\mathbf VW'$. Indeed, the relative matrices are conjugated
by $W$, which preserves their norms and commutes with the
logarithmic series in \eqref{I-3-23}. The inverse bijection
uses $W^{-1}$ and $(W')^{-1}$. It preserves the number of means
and gives the asserted formula when this number is one.
\end{proof}


\section{Gauge normalization and lifts}\label{sec:gauge-lifts}
The contour mean permits a separate integration over each
noncentral gauge variable in a block. This gives a normalized
gauge factor. An explicit lift then realizes arbitrary coarse
bond values as fine configurations whose block averages recover
those values exactly. Both constructions are used in Ba{\l}aban's
renormalization expansion \cite[(1.3), (1.5)--(1.7)]{B-CMP119}.

\subsection{Axial gauge and its normalized factor}
\begin{definition}[Block axial gauge]\label{I-3-66}
Let $0\le k<K$ and let $P''$ be a set of sites of
$\mathbb T^{(k+1)}$. Write $B(P'')=\bigcup_{y\in P''}B(y)$.
A configuration $V$ on a reversal-invariant bond set containing
all bonds of these blocks is in the block axial gauge on $P''$
if all its contour means are defined and
\begin{equation*}
 V(y,x)=\mathbf1\qquad(y\in P'',\ x\in B(y)\setminus\{y\}).
\end{equation*}
\end{definition}
We use the averaged-contour gauge factor of
Ba{\l}aban \cite[(0.14)--(0.16), pp.~254--255]{B-CMP109}.
Its normalization integrates independently over the gauge
value at every noncentral site of these blocks.

For $\mathsf x>0$ and $0<\epsilon\le1$, put
\begin{equation}\label{I-3-67}
 \mathsf z(\mathsf x,\epsilon)=
 \int_G\mathbf1_{\{|v-\mathbf1|<\epsilon\}}
 e^{-\mathsf x(1-\frac12\tr v)}\,dv
\end{equation}
and define
\begin{equation}\label{I-3-68}
 \mathfrak G(V;P'';\mathsf x,\epsilon)=
 \prod_{y\in P''}\prod_{x\in B(y)\setminus\{y\}}
 \frac{\mathbf1_{\{|V(y,x)-\mathbf1|<\epsilon\}}}
      {\mathsf z(\mathsf x,\epsilon)}
 e^{-\mathsf x(1-\frac12\tr V(y,x))}.
\end{equation}
Set this factor equal to zero when any required mean is undefined.
Here $\mathsf z>0$: the integrand is positive on a nonempty open
set, and such a set has positive Haar measure since finitely many
of its translates cover $G$. The factor is bounded and Borel by
\cref{lem:unique-group-mean}; each of its factors lies
between zero and $\mathsf z^{-1}$.

For parameters $x_k>0$ and $\varepsilon_k\in(0,1]$, write
\begin{equation}\label{I-3-69}
 \mathfrak G_k(V;P'')=
       \mathfrak G(V;P'';x_k,\varepsilon_k).
\end{equation}
On the domain where
the means exist, the associated gauge-fixing function is
\begin{equation}\label{I-3-70}
 G_k(P'';V)=\sum_{y\in P''}\sum_{x\in B(y)\setminus\{y\}}
       \left(1-\tfrac12\tr V(y,x)\right).
\end{equation}
Every summand is $\tfrac12|V(y,x)-\mathbf1|^2$, so this
function is nonnegative and vanishes precisely in axial gauge.

We use invariance of normalized Haar measure under inversion.
Indeed, its image $\kappa$ under inversion is right invariant:
for bounded Borel $\varphi$ and fixed $v$, left invariance gives
\[
 \int\varphi(uv)\,\kappa(du)
 =\int\varphi((v^{-1}u)^{-1})\,du
 =\int\varphi(u)\,\kappa(du).
\]
Integrating $\varphi(uv)$ against $\kappa(du)\,dv$ in the two
orders, right invariance of $\kappa$ and left invariance of $dv$
give $\int\varphi\,d\kappa=\int\varphi\,dv$. Thus $\kappa=dv$.

\begin{proposition}
Suppose all the means $V(y,x)$ in \eqref{I-3-68} exist. Let
$u$ equal $\mathbf1$ outside $B(P'')\setminus P''$. Then
\begin{equation}\label{I-3-71}
 \int \mathfrak G(V^u;P'';\mathsf x,\epsilon)
                  \prod_{x\in B(P'')\setminus P''}du(x)=1.
\end{equation}
\end{proposition}

\begin{proof}
By \cref{lem:unique-group-mean}, left and right multiplication
by elements of $G$ preserve the unique-mean domain.
Consequently all transformed means exist and
$V^u(y,x)=V(y,x)u(x)^{-1}$, since $u(y)=\mathbf1$.
The blocks are disjoint, and their centers are their only sites
of $\mathbb T^{(k+1)}$. Hence each integration variable occurs
in precisely one factor. Haar invariance under inversion and
translation makes the integral of this factor equal to
$\mathsf z/\mathsf z=1$. Taking the product proves
\eqref{I-3-71}.
\end{proof}

\subsection{Lifting coarse configurations}
An explicit lift preserves a coarse value on every crossing bond
and assigns the identity within each block. Its averages recover
the same assignment at every intermediate scale.

Let $0\le k\le K$, and let $V$ be a $G$- or $G^c$-valued
configuration on a reversal-invariant set $\mathcal C$ of bonds
of $\mathbb T^{(k)}$. Define its bare lift by
\begin{equation}\label{I-3-73}
 (Q^{s*}_kV)(b)=
 \begin{cases}
 \mathbf1,&b_-,b_+\in B^k(y)\text{ for some }y\in\mathbb T^{(k)},\\
 V(c),&b\text{ lies over }c\in\mathcal C.
 \end{cases}
\end{equation}
It is defined on precisely these bonds. Reversal leaves the first
case unchanged and replaces $c$ by $\bar c$ in the second, so
it is a configuration. At $k=0$ this is the identity map.

\begin{proposition}\label{I-3-72}
For $0\le j\le k$ the partial iterate
$Q^{s*}_{k\to j}V:=\mathcal M^j(Q^{s*}_kV)$ is defined on
exactly the bonds of $\mathbb T^{(j)}$ lying inside an order-$(k-j)$
block or over a bond of $\mathcal C$. There,
\begin{equation}\label{I-3-74}
 (Q^{s*}_{k\to j}V)(b)=
 \begin{cases}
 \mathbf1,&b_-,b_+\in B^{k-j}(y),\\
 V(c),&b\text{ lies over }c\in\mathcal C.
 \end{cases}
\end{equation}
In particular $\mathcal M^k(Q^{s*}_kV)=V$ on $\mathcal C$.
Every participating complex average lies in its analytic domain.
For real $V$, the small-curvature case applies at every step and
the real and complex averages agree.
\end{proposition}

\begin{proof}
Let $W_j$ denote the right side of \eqref{I-3-74} on the stated
domain. We have $W_0=Q^{s*}_kV$ and $W_k=V$. To prove the
assertion, fix $j<k$ and a positive bond $c'$ of
$\mathbb T^{(j+1)}$. Its endpoints lie in unique order-$(k-j-1)$
blocks centered at $y_-,y_+\in\mathbb T^{(k)}$. Block nesting
gives $B(c'_\pm)\subset B^{k-j}(y_\pm)$.

If $y_-=y_+$, then $W_j=\mathbf1$ on the entire double block
of $c'$, and its average is $\mathbf1=W_{j+1}(c')$.
Otherwise $c'$ lies over $c=\langle y_-,y_+\rangle$.
Every bond of its double block is internal to one of
$B(c'_-),B(c'_+)$ or is a corridor bond or its reverse.
Internal bonds lie within order-$(k-j)$ blocks; corridor bonds
lie over $c$, and their reverses over $\bar c$.
Since the corridor is nonempty, all the required fine bonds
belong to the domain of $W_j$ if and only if $c\in\mathcal C$.
This is exactly the domain condition for $W_{j+1}(c')$.

When this condition holds, internal values are $\mathbf1$ and
all corridor values are $V(c)$. A plaquette within either block
has holonomy $\mathbf1$; a plaquette crossing between them has
two crossing edges whose factors cancel, so again has holonomy
$\mathbf1$. Each contour stays in one block, hence has transport
$\mathbf1$. Each straight path $[x,x']$ crosses the corridor
once, hence has transport $V(c)$. Thus all contour families are
constant with mean $\mathbf1$, and every relative loop $W_x$
in the averaging formula is $\mathbf1$. The complex formula
is defined and gives $V(c)=W_{j+1}(c')$. For real values the
plaquette calculation places the field in $\mathcal S_{c'}$.
Negative bonds follow by inversion. This proves the induction.
\end{proof}


\section{Regions and multiscale constraints}\label{sec:multiscale-constraints}
Constraints at different scales are assigned to the successive
layers of a decreasing sequence of regions. To evaluate a
constraint on a coarse bond, both of its fine endpoint blocks
must stay in the initial region. We give this domain criterion,
then describe the removal of inward bonds and a choice of nested
regions that supplies the required separation. The determining
data and region towers follow Ba{\l}aban's construction
\cite[(2.10)--(2.13), pp.~256--257]{B-CMP119}.

\subsection{Regions and determining data}
Write $\pi[a_r]=\{\Delta_r(y):y\in\Lambda_{r,\ell}\}$,
$r\ge-m$. A region for this partition is a union of its closed
cubes, including the empty union. Its complement $X^c$ means
the union of the other cubes. As a point set this is the closure
of $T_\ell\setminus X$ and is independent of refinement.
Indeed, excluded cube interiors are disjoint from $X$ and dense
in their cubes; the union of the excluded cubes is closed and
contains the point-set complement. A lattice site at the partition scale or a finer scale lies in exactly one
of $X,X^c$, by the block membership established above.
Unions, intersections and differences of regions mean the corresponding
operations on their cube families at a specified common refinement. Two cubes touch when
their intersection is nonempty. We use
$\operatorname{dist}_\infty(\cdot,\emptyset)=+\infty$.

For two cubes in the same standard partition, of side $s$
centered at $y,y'$,
\begin{equation}\label{eq:cube-distance}
 \operatorname{dist}_\infty(\Delta(y),\Delta(y'))
                  =\max\{0,|y-y'|_\infty-s\}.
\end{equation}
For the lower bound use the triangle inequality. For equality,
choose a least lift $v$ of $y'-y$ and move from each center
toward the other by
$t_\mu=\operatorname{sign}(v_\mu)\min\{s/2,|v_\mu|/2\}$
in coordinate $\mu$. The resulting points belong to their cubes
and their difference has coordinates of modulus
$\max\{0,|v_\mu|-s\}$. Distances between nonempty regions of
this partition are consequently attained integer multiples of $s$.
The diameter of a cube in the sup metric is at most $s$.

\begin{definition}\label{I-3-77}
Let $0\le k\le K$ and let
$\Omega_0\supset\Omega_1\supset\cdots\supset\Omega_k$
be regions, with $\Omega_j$ a union of cubes of side at least
$a_{K-j}$. Empty members are allowed. Set $\Omega_{k+1}=\emptyset$
and, for $0\le j\le k$, let
\[
 \Lambda_j(\mathbf B)=(\Omega_j\setminus\Omega_{j+1})^{(j)},
 \qquad
 \Gamma_j=\{\text{inner bonds of }\Omega_0\text{ at level }j
                    \text{ with an endpoint in }\Lambda_j(\mathbf B)\}.
\]
Here $X^{(j)}=X\cap\mathbb T^{(j)}$, and the region difference
is taken at cube side $a_{K-j}$. Thus its sites at level $j$
are precisely those in $\Omega_j$ and outside $\Omega_{j+1}$.
The determining set is $\mathbf B=\bigcup_{j=0}^k\Gamma_j$;
its support is $\Omega_1$ (empty when $k=0$).
A datum is a family $V=(V_j)_{j=0}^k$ of $G$-valued configurations on
$\Gamma_j$. For a configuration $U$ on the inner bare bonds
of $\Omega_0$, put
\[
 M_{\mathbf B}(U)=\bigl((\mathcal M^jU)|_{\Gamma_j}\bigr)_{j=0}^k
\]
whenever all indicated values exist. A node is an endpoint of
one of the bonds in $\mathbf B$.
\end{definition}

\begin{proposition}[Domain of the determining map]\label{prop:determining-domain}
For real configurations, the map $M_{\mathbf B}$ is defined
on every configuration on the inner bare bonds of $\Omega_0$
if and only if
\begin{equation}\label{eq:determining-domain}
 B^j(c_-),B^j(c_+)\subset\Omega_0
             \qquad(0\le j\le k,\ c\in\Gamma_j).
\end{equation}
This is a condition on the sequence of regions. Each of the
following conditions suffices:
\begin{enumerate}[label=(\roman*)]
\item $\Omega_0=T_\ell$;
\item $\Omega_0$ is a region for $\pi[a_{K-k}]$;
\item in level-$0$ units,
$\operatorname{dist}_\infty(\Omega_{j-1}^c,\Omega_j)>L^j$
for $1\le j\le k$.
\end{enumerate}
Complex configurations use the complex averaging domains at
every participating bond, in addition to the geometric condition.

\end{proposition}

\begin{proof}
The exact determining bond set
$\mathsf D_j(c)$ consists of the bare bonds with both endpoints
in $B^j(c_-)\cup B^j(c_+)$. Every site of this double block
is an endpoint of such a bond: for $j=0$ use $c$, and for
$j\ge1$ use a neighbor within its block. Thus every determining
bond is inner to $\Omega_0$ precisely when both blocks lie there.
Conditions (i) and (ii) give this by the previously proved block
membership property. Under (iii), choose an endpoint
$y\in\Lambda_j(\mathbf B)$ of $c$ and denote the other by $y'$.
For $j=0$ both endpoints already lie in $\Omega_0$. For $j\ge1$,
the level-$j$ cube centered at $y$ belongs to $\Omega_j$ and
touches the cube centered at $y'$. Take a common point $z$.
Every point of the latter cube is within distance $L^j$ of
$z\in\Omega_j$. The strict gap excludes this entire cube from
$\Omega_{j-1}^c$, so it lies in $\Omega_{j-1}\subset\Omega_0$.
Both required blocks are therefore contained in $\Omega_0$.
\end{proof}

We also use the coordinate formula for $M_{\mathbf B}$ when
the configuration is given on a larger bond domain, evaluating
each average on that domain. In particular, a full configuration
$V\in\mathcal V_k$ induces the data
\begin{equation}\label{I-3-75}
 M_{\mathbf B}(Q^{s*}_kV)
 =\bigl((\mathcal M^j(Q^{s*}_kV))|_{\Gamma_j}\bigr)_{j=0}^k
 =\bigl((Q^{s*}_{k\to j}V)|_{\Gamma_j}\bigr)_{j=0}^k,
 \qquad\mathbf B=\bigcup_{j=0}^k\Gamma_j.
\end{equation}
All these values are defined by \cref{I-3-72}. When \eqref{eq:determining-domain}
holds, the same datum is obtained from the restriction of the lift to
the inner bare bonds of $\Omega_0$, by locality of the averages.


The constraint set consists of the configurations on the inner
bare bonds of $\Omega_0$ satisfying $M_{\mathbf B}(U)=V$:
\[
 C(\mathbf B,V)=\{U:M_{\mathbf B}(U)=V\}.
\]
Its residual gauge group is
\[
 \mathcal G^{\mathrm{res}}(\mathbf B)
 =\{u:\Omega_0^{(0)}\to G:\ u=\mathbf1\text{ at every node}\}.
\]
Set $C(\mathbf B,V)=\emptyset$ if the real map is not defined
on this domain. The restricted action is
$A_{\Omega_0}(U)=\sum_{p\subset\Omega_0}a_p(U)$, with the sum
over inner bare plaquettes. The residual group preserves the
constraint set: iterated covariance gives
\[
 (\mathcal M^j(U^u))(c)
 =u(c_-)(\mathcal M^jU)(c)u(c_+)^{-1}
 =(\mathcal M^jU)(c)\qquad(c\in\Gamma_j).
\]
Extend $u$ by the identity outside $\Omega_0$ when applying that
covariance formula. The action is invariant under every gauge
transformation, since every plaquette transforms by conjugation.

\begin{lemma}\label{I-3-76}
In level-$k$ units put $\xi=\eta_k=L^{-k}$ and
$A^{\eta_k}_{\Omega_0}=L^kA_{\Omega_0}$.
For $k<K$,
$A^{\eta_{k+1}}_{\Omega_0}=L A^{\eta_k}_{\Omega_0}$.
If $r\ge0$, $U$ is real and
$|\partial U(p)-\mathbf1|\le r\xi^2$ for every inner bare
plaquette of $\Omega_0$, then
\[
 A^{\eta_k}_{\Omega_0}(U)
 \le\tfrac12r^2\xi^3\#\{p\subset\Omega_0\}
 \le\tfrac32r^2\operatorname{vol}(\Omega_0),
\]
where the volume is in level-$k$ units.
\end{lemma}

\begin{proof}
The rescaling identity follows from $\eta_{k+1}^{-1}=L\eta_k^{-1}$.
For the bound, $a_p(U)=\tfrac12|\partial U(p)-\mathbf1|^2
\le\tfrac12r^2\xi^4$. Sum over inner plaquettes and multiply
by $\xi^{-1}$. Each plaquette has a base site in $\Omega_0$
and one of three direction pairs, so there are at most three
per site. The region contains one bare cube of volume $\xi^3$
per bare site, proving the last inequality.
\end{proof}

\subsection{Removing inward bonds}
When $k\ge1$, removing inward bonds leaves the noncentral
sites of blocks centered in $\Omega_k^{(k)}$ outside the node set, as proved in
\cref{I-3-81}(iv). The residual group can then vary the gauge
values at these sites while preserving the retained constraints.
The sequence of regions is kept fixed.

\begin{definition}[Bond sets]\label{I-3-79}\label{I-3-80}
Let $\mathbf B=\bigcup_{j=0}^k\Gamma_j$ be the determining set
of $\Omega_0\supset\cdots\supset\Omega_k$, and put
$\Omega_{k+1}=\emptyset$. If $\mathbf C\subset\mathbf B$ is
closed under reversal, its members are $\mathbf C\cap\Gamma_j$.
Data, the averaging map, nodes, the constraint set and the residual
gauge group are defined by restricting the corresponding definitions
to these members. Thus
\[
 M_{\mathbf C}(U)
 =\bigl((\mathcal M^jU)|_{\mathbf C\cap\Gamma_j}\bigr)_{j=0}^k,
 \qquad
 C(\mathbf C,V)=\{U:M_{\mathbf C}(U)=V\},
\]
where $U$ is a configuration on the inner bare bonds of $\Omega_0$.
The map is defined when every displayed average is defined. The
site sets remain those of the given sequence.

For $j<k$, a bond of $\Gamma_j$ is \emph{inward} if it has an
endpoint in $\Omega_{j+1}$. Write $\mathfrak F_{\mathrm{in}}$
for the set of all inward bonds. For a reversal-invariant subset
$\mathfrak F\subset\mathfrak F_{\mathrm{in}}$, the associated
\emph{bond set} and the \emph{reduced determining set} are
\[
 \mathbf B_{\mathfrak F}=\mathbf B\setminus\mathfrak F,
 \qquad
 \mathbf B^\flat=\mathbf B\setminus\mathfrak F_{\mathrm{in}},
 \qquad
 \Gamma_j^\flat=\Gamma_j\setminus\mathfrak F_{\mathrm{in}}.
\]
The full set is $\mathbf B_\emptyset=\mathbf B$.
\end{definition}

The residual group preserves $C(\mathbf C,V)$ by the same
covariance calculation: each retained bond has both endpoints
among the nodes fixed by this group. The restricted action
remains gauge invariant.

\begin{lemma}[Members and nodes]\label{I-3-81}
Let $\mathbf B$ be as above.
\begin{enumerate}[label=(\roman*)]
\item An ordered pair of sites is a bond of at most one level
$0\le j\le K$. Hence the members $\Gamma_j$ are pairwise
disjoint. Each member and $\mathfrak F_{\mathrm{in}}$ are closed
under reversal.
\item The set $\Gamma_j^\flat$ consists of the bonds of
$\mathbb T^{(j)}$ with both endpoints in $\Omega_0$, at least
one endpoint in $\Omega_j$, and neither endpoint in
$\Omega_{j+1}$. Every set $\mathbf B_{\mathfrak F}$ contains $\mathbf B^\flat$,
and therefore contains every level-$j$ bond with both endpoints
in $\Lambda_j(\mathbf B)$.
\item Suppose $k<K$, and append a region
$\Omega'_{k+1}\subset\Omega_k$ whose partition cubes have
side at least the spacing of $\mathbb T^{(k+1)}$. Write
$\Gamma'_j$ and $\Gamma'^{\flat}_j$ for the new full and
reduced members. For $j<k$ these agree with the old members.
The top member $\Gamma'_{k+1}=\Gamma'^{\flat}_{k+1}$ consists
of the level-$(k+1)$ bonds with both endpoints in $\Omega_0$
and an endpoint in $\Omega'_{k+1}$. Moreover,
\[
 \Gamma_k^\flat\setminus\Gamma'^{\flat}_k
 =\{b\text{ of level }k:b_-,b_+\in\Omega_0,
                 \ \{b_-,b_+\}\cap\Omega'_{k+1}\ne\emptyset\}.
\]
For the full sets, $\Gamma_k\setminus\Gamma'_k$ instead
consists of those $b\in\Gamma_k$ whose endpoints each lie
in $\Omega'_{k+1}$ or outside $\Omega_k$. If
\[
 \operatorname{dist}_\infty(\Omega_k^c,\Omega'_{k+1})>L^k
\]
in level-$0$ units, this is exactly the set of level-$k$ bonds
with both endpoints in $\Omega'_{k+1}$.
\item Every node of $\mathbf B^\flat$ belongs to
$\bigcup_{j=0}^k\Lambda_j(\mathbf B)$. In particular, if
$k\ge1$ and $Y\subset\Omega_k^{(k)}$, then
$B(Y)\setminus Y$ contains no node of $\mathbf B^\flat$.
Here $B(Y)$ is the union of the level-$(k-1)$ blocks centered
at the sites of $Y$.
\end{enumerate}
\end{lemma}

\begin{proof}
Use level-$0$ units. If the same ordered pair is a bond of levels
$q,q'$, then
$\epsilon L^qe_\mu-\epsilon'L^{q'}e_{\mu'}$ is a period
of the torus. Every coordinate has modulus at most
$2L^K<2L^{m+K}$, so this period vanishes. Thus $q=q'$.
The endpoint conditions for the full members and for inward bonds
are invariant under reversal, proving (i).

Removing the inward bonds from $\Gamma_j$ leaves exactly the
bonds described in (ii), including at $j=k$ because
$\Omega_{k+1}=\emptyset$. Since
$\mathfrak F\subset\mathfrak F_{\mathrm{in}}$, every
$\mathbf B_{\mathfrak F}$ contains $\mathbf B^\flat$.
Both endpoints of a bond in $\Lambda_j(\mathbf B)$ lie in
$\Omega_j$ and outside $\Omega_{j+1}$, giving the last claim.

The members with $j<k$ depend only on
$\Omega_0,\Omega_j,\Omega_{j+1}$, which have not changed.
At level $k$, the old reduced member equals the old full member,
and its new reduced member deletes exactly the bonds meeting
$\Omega'_{k+1}$. The new full member deletes a bond precisely
when neither endpoint belongs to $\Omega_k\setminus\Omega'_{k+1}$.
This gives the two descriptions in (iii). Under the distance
condition, a deleted full-set bond has an endpoint in
$\Omega'_{k+1}$, and its other endpoint, at distance $L^k$,
cannot lie outside $\Omega_k$. Both therefore lie in
$\Omega'_{k+1}$. The converse follows from the same descriptions.

For (iv), let $x$ be an endpoint of a bond of $\Gamma_i^\flat$.
Then $x\in\Omega_0$ and $x\notin\Omega_{i+1}$. There is
a unique $i'\le i$ with $x\in\Omega_{i'}$ and
$x\notin\Omega_{i'+1}$. The inclusion
$\mathbb T^{(i)}\subset\mathbb T^{(i')}$ gives
$x\in\Lambda_{i'}(\mathbf B)$.
If $y\in Y$ and $x\in B(y)\setminus\{y\}$, block
containment gives $x\in\Omega_k$. It cannot belong to a
site set of index less than $k$, and it is not a site of
$\mathbb T^{(k)}$, since $B(y)\cap\mathbb T^{(k)}=\{y\}$.
It is therefore not a node.
\end{proof}

\subsection{Separated towers}
Fix powers $\mathfrak M,R_1,M_1$ of $L$, with nonnegative
integer exponents, $R_1\ge L$ and $M_1\ge R_1\mathfrak M$.
For $s\in\{L^r:r\in\mathbb N_0\}$ and $0\le q\le K$,
an $s$-cube at level $q$ has physical side $s a_{K-q}$;
such cubes exist exactly when $s\le L^{m+K-q}$.
Let $\Omega$ be a union of cubes of side at least
$M_1a_{K-i}$, with $0\le i\le k\le K$. A tower to an integer
depth $n\ge0$, $i+n\le k$, is a decreasing sequence
$\Omega=\Omega_i\supset\cdots\supset\Omega_{i+n}$
built by either of the following rules, for $i\le q<i+n$.
In these rules lengths are in level-$i$ units.
\begin{enumerate}[label=(\roman*)]
\item Under $\mathfrak r_1$, take all $\mathfrak M$-cubes
at level $q+1$ contained in $\Omega_q$ whose distance from
$\Omega_q^c$ is strictly greater than
$R_1\mathfrak M L^{q-i}$.
\item Under $\mathfrak r_2$, take all $M_1$-cubes at level
$q+1$ contained in $\Omega_q$ whose distance from
$\Omega_q^c$ is at least $M_1L^{q+1-i}$.
\end{enumerate}
The union is empty when the required cubes do not exist.
Each rule selects the largest union of its specified cubes
satisfying its own condition.

More generally, a decreasing sequence satisfies the separation
condition of $\mathfrak r_1$ if each member is a union of
$\mathfrak M$-cubes at its level, is empty when no such cubes
exist, and
\begin{equation}\label{eq:tower-separation}
 \operatorname{dist}_\infty(\Omega_q^c,\Omega_{q+1})
                    >R_1\mathfrak M L^{q-i}.
\end{equation}
Both tower rules satisfy this condition. Indeed, the gap of a
finite union is the minimum of the gaps of its cubes, and
$M_1L^{q+1-i}\ge R_1\mathfrak M L^{q+1-i}
>R_1\mathfrak M L^{q-i}$. Every $M_1$-cube is a union
of $\mathfrak M$-cubes. The constant whole-torus sequence
satisfies the condition exactly when
$\mathfrak M\le L^{m+K-i-n}$; then its complementary regions
are empty and all gaps are infinite. The constant whole-torus sequence will also be allowed when this
size condition fails; in that case it is not called separated.

For $i=0$, separation gives
\[
 \operatorname{dist}_\infty(\Omega_{j-1}^c,\Omega_j)
 >R_1\mathfrak M L^{j-1}\ge L^j,
\]
so \eqref{eq:determining-domain} holds. The region
$\Omega_q\setminus\Omega_{q+1}$ is the seam of level $q$.
Write $\mathbf B_j(\Omega)$ for the determining set of the
$\mathfrak r_2$-tower from level zero to depth $j$; at $j=0$
it is the determining set of $\Omega$ alone.

For a region $X$ in a fixed partition, let $X^{\sim-1}$ be
the union of its cubes that do not touch $X^c$, and set
$X^{\sim-2}=(X^{\sim-1})^{\sim-1}$.

\begin{proposition}\label{prop:explicit-tower}
Let $1\le j\le k\le K$, suppose $M_1\le L^{m+K-j}$,
and let $\Omega$ be a union of $M_1$-cubes at level $j$.
In level-zero units set $s_q=M_1L^q$.
Its $\mathfrak r_2$-tower from level zero satisfies
\[
 \Omega_1=\bigcup\{\Delta\text{ of side }s_1:\Delta\subset\Omega,
                       \ \operatorname{dist}_\infty(\Delta,\Omega^c)\ge s_1\},
\]
and, for $2\le q\le j$,
\[
 \Omega_q=\bigcup\{\Delta\text{ of side }s_q:\Delta\subset\Omega,
                       \ \operatorname{dist}_\infty(\Delta,\Omega^c)\ge2s_q\}.
\]
With layers measured in $M_1$-cubes at level $j$,
$\Omega^{\sim-2}\subset\Omega_j$, with equality for $j\ge2$.
\end{proposition}

\begin{proof}
For a cube $\Delta$ of side $s_q$, put
$D(\Delta)=\operatorname{dist}_\infty(\Delta,\Omega^c)$.
Both $\Omega$ and $\Omega^c$ are unions of cubes at this scale.
Thus $D(\Delta)$ belongs to $s_q\mathbb Z\cup\{+\infty\}$,
and every point of $\Delta$ lies within $D(\Delta)+s_q$
of $\Omega^c$ when the distance is finite.
Set $t_1=s_1$ and $t_q=2s_q$ for $q\ge2$, and let
$\Omega_q^\star$ be the union of the cubes with
$\Delta\subset\Omega$ and $D(\Delta)\ge t_q$.
The first step is the definition of $\mathfrak r_2$.
Suppose $q\ge2$ and $\Omega_{q-1}=\Omega_{q-1}^\star$.

Every point at distance at most $s_{q-1}$ from $\Omega^c$
lies in $\Omega_{q-1}^c$. For strict inequality each containing
cube fails the threshold. In the equality case take a shortest
segment toward $\Omega^c$; its remaining points satisfy the
strict inequality, and $\Omega_{q-1}^c$ is closed.
Conversely every point of $\Omega_{q-1}^c$ is within
$2s_{q-1}$ of $\Omega^c$: a cube there either lies in
$\Omega^c$ or has $D<t_{q-1}\le2s_{q-1}$, hence
$D\le s_{q-1}$ by quantization. Its diameter supplies the
remaining $s_{q-1}$.

If $D(\Delta)\ge2s_q$, its subcubes meet the previous
threshold, so $\Delta\subset\Omega_{q-1}$. The preceding
distance bound gives
\[
 \operatorname{dist}_\infty(\Delta,\Omega_{q-1}^c)
 \ge2s_q-2s_{q-1}\ge s_q,
\]
so $\Delta$ is selected by the rule. Conversely, if selected
$\Delta$ had $D(\Delta)<2s_q$, quantization would give
$\delta=D(\Delta)\le s_q$. Choose points $x\in\Delta$
and $z\in\Omega^c$ realizing $\delta$. If
$\delta\le s_{q-1}$, then $x\in\Omega_{q-1}^c$.
Otherwise the point on a shortest segment from $z$ to $x$
at distance $s_{q-1}$ from $z$ lies in $\Omega_{q-1}^c$
and is at distance $\delta-s_{q-1}<s_q$ from $x$.
Both cases contradict the required gap $s_q$.
This proves $\Omega_q=\Omega_q^\star$ by induction.

At the final scale, a cube remains after one layer is removed
exactly when its center is at least $2s_j$ from every center
of a complementary cube. It remains after two removals exactly
when those distances are at least $3s_j$. For the latter
criterion, a complementary center within $2s_j$ gives an
intermediate center within $s_j$ of both centers: move each
coordinate toward the complementary center by the smaller of
$s_j$ and that coordinate's displacement in a least lift.
The intermediate cube is removed at the first step and touches
the original cube. Conversely, if all complementary centers
are at least $3s_j$ away, every cube touching the original
has center at least $2s_j$ from them and remains at the first
step. Formula \eqref{eq:cube-distance} identifies this
two-layer criterion with $D(\Delta)\ge2s_j$.
The claimed inclusion and equality now follow from the formulas
for $\Omega_j$.
\end{proof}
\section{Axial gauge representatives}\label{sec:axial-representatives}
\subsection{Combs and tower gauge conditions}

A gauge transformation changes a transport only at its endpoints.
We shall use this fact successively inside the blocks of a tower.
Here the paths are fixed coordinate-order paths. The block
gauge of \cref{I-3-66} instead uses averaged contour variables.

We use Ba{\l}aban's coordinate-order combs
\cite[(1.7), p.~18]{B-CMP95} and the relative axial gauge
conditions imposed at successive scales
\cite[(64)--(69), p.~29]{B-CMP98}
\cite[(1.15), (1.19), pp.~78--79]{B-CMP99a}, with the
centered blocks and bond sets defined above.

\begin{definition}\label{I-3-89}
On $\mathbb T^{(q)}$, write $\xi$ for the spacing in the
chosen units and $N_q=2L^{m+K-q}$ for the number of sites
in each direction. If
\[
 x-y=\xi\sum_{\mu=1}^3 n_\mu e_\mu,
 \qquad |n_\mu|<N_q/2,
\]
the integers $n_\mu$ are uniquely determined. The comb from
$y$ to $x$ adjusts the coordinates in the order $3,2,1$:
\begin{equation}\label{I-3-90}
 \Gamma_{y,x}
 =[y,(y_1,y_2,x_3)]
  \cup[(y_1,y_2,x_3),(y_1,x_2,x_3)]
  \cup[(y_1,x_2,x_3),x].
\end{equation}
Each segment is the straight lattice path with the prescribed
signed displacement. An empty segment contributes no bonds.
Coordinates in this formula may be taken in a lift of any
containing box whose sides are less than half the torus side.

Let $W_3$ be the group of signed coordinate permutations.
For $r\in W_3$, the comb in the frame $r$ is
\[
 \Gamma^r_{y,x}=r\Gamma_{r^{-1}y,r^{-1}x}.
\]
Every $r\in W_3$ preserves the lattices, bonds, and standard
cube partitions. Unless a frame is specified, we use the
coordinate frame in \eqref{I-3-90}.

Let $\mathbf B$ be any bond set of the decreasing sequence
$\Omega_0\supset\cdots\supset\Omega_k$, and let $U_0,U$
be $G$-valued configurations on the inner bare bonds of $\Omega_0$.
Here $\mathbf B$ may be full, reduced, or obtained by deleting
any reversal-invariant subset of the inward bonds;
$\Lambda_j(\mathbf B)$ always denotes the site set of the
underlying sequence.
We write $U\in\operatorname{Ax}(\mathbf B,U_0)$ if
\begin{equation}\label{I-3-93}
 (\mathcal M^nU)(\Gamma_{z,x})
       =(\mathcal M^nU_0)(\Gamma_{z,x})
\end{equation}
whenever
\[
 \begin{gathered}
 1\le j\le k,\qquad x_j\in\Lambda_j(\mathbf B),\qquad
 0\le n<j,\\
 z\in\mathbb T^{(n+1)}\cap B^{j-n-1}(x_j),
 \qquad x\in B(z).
 \end{gathered}
\]
These are the tower axial gauge conditions relative to $U_0$.
For $U_0=\mathbf1$ we call them the absolute tower axial
gauge conditions. Replacing every comb by its image in the
frame $r$ defines $\operatorname{Ax}^r(\mathbf B,U_0)$.
\end{definition}

The strict half-period restriction makes the comb unambiguous:
two possible displacements differ in each coordinate by an
integer multiple of $N_q$, while their difference has modulus
less than $N_q$. If $z\in\mathbb T^{(n+1)}$ and $x\in B(z)$,
the offsets have modulus at most $(L-1)/2<N_n/2$.
Thus $\Gamma_{z,x}$ is one of the contour paths already
defined in the block, and it stays in $B(z)$.

The transports in \eqref{I-3-93} are defined on these blocks,
as verified in the proof of \cref{I-3-96} below.
At $x=z$ the path is empty and the condition is an identity.

\subsection{A representative of each residual orbit}
Fix a bond set $\mathbf B$ of a decreasing sequence
$\Omega_0\supset\cdots\supset\Omega_k$, and put
$\Omega_{k+1}=\emptyset$. Its site sets are
$\Lambda_j=\Lambda_j(\mathbf B)$.
The tower gauge fixes transports within the blocks associated
with these site sets. We first impose its conditions only at
endpoints that the residual gauge group is free to change.

\begin{definition}\label{I-3-95}
For $x_j\in\Lambda_j$, call $B^j(x_j)$ its \emph{pillar},
with top $x_j$. If $x\ne x_j$ belongs to this pillar, define
its level and parent by
\[
 x\in\mathbb T^{(n(x))}\setminus\mathbb T^{(n(x)+1)},
 \qquad x\in B(z(x)),\quad
 z(x)\in\mathbb T^{(n(x)+1)}.
\]
Here $0\le n(x)<j$. The ascent from $x$ is obtained by
successively taking parents until reaching $x_j$. Its
\emph{free part} stops at the first node of $\mathbf B$ or
at the top.

For configurations $U_0,U$ on the inner bare bonds of $\Omega_0$,
write $U\in\operatorname{Ax}^{\circ}(\mathbf B,U_0)$ when
the conditions \eqref{I-3-93} hold for every tower triple
$(n,z,x)$ with $x$ not a node. Thus the full tower gauge
$\operatorname{Ax}(\mathbf B,U_0)$ imposes these conditions
also when $x$ is a node. The choice $U_0=\mathbf1$ gives the
absolute gauges.
\end{definition}

The definitions of level, parent and ascent are unambiguous;
their geometric properties are proved below. The following
condition suffices for a representative in every residual orbit:
\begin{equation*}\tag{T}\label{I-3-97}
\begin{gathered}
 \text{every site of $\Lambda_j$, $1\le j\le k$,
 is a node of $\mathbf B$,}\\
 \text{and every site of $\Lambda_0$ incident to an inner
 bare bond of $\Omega_0$ is a node.}
\end{gathered}
\end{equation*}
The second condition permits isolated bare sites, on which a
gauge transformation has no effect on the configuration.

\begin{proposition}[Tower gauge in a residual orbit]\label{I-3-96}
Suppose \eqref{I-3-97} holds, and let $U_0$ be any
configuration on the inner bare bonds of $\Omega_0$.
Every orbit of $\mathcal G^{\mathrm{res}}(\mathbf B)$
on these configurations meets
$\operatorname{Ax}^{\circ}(\mathbf B,U_0)$ in exactly
one configuration.

For a given configuration $U$, its representative is $U^u$.
The gauge $u$ is uniquely specified by $u=\mathbf1$ at
the nodes and at $\Lambda_0$, and by the recursion
\begin{equation}\label{I-3-98}
 u(x)=(\mathcal M^nU_0)(\Gamma_{z,x})^{-1}
       u(z)(\mathcal M^nU)(\Gamma_{z,x}),
 \qquad (n,z)=(n(x),z(x)),
\end{equation}
at all remaining sites. This recursion is performed in each
pillar in decreasing order of level. If a residual gauge $w$
also satisfies $U^w\in\operatorname{Ax}^{\circ}(\mathbf B,U_0)$,
then $w=u$ except possibly at isolated sites of $\Lambda_0$
that are not nodes.
\end{proposition}

\begin{proof}
We first verify the geometry used in the recursion.
The regions $\Omega_j$ and $\Omega_{j+1}$ are unions of
partition cubes at level $j$. The bare sites in the cube
centered at $x_j\in\mathbb T^{(j)}$ form $B^j(x_j)$.
Consequently, the bare sites of the layer
$\Omega_j\setminus\Omega_{j+1}$ are the disjoint union of
$B^j(x_j)$ over $x_j\in\Lambda_j$. As these layers partition
the bare sites of $\Omega_0$, so do all the pillars.

A non-top site of $B^j(x_j)$ lies at a unique level
$0\le n<j$. Nesting of blocks gives
\[
 B^j(x_j)\cap\mathbb T^{(n)}
 =B^{j-n}(x_j)
 =\bigcup_{z\in B^{j-n-1}(x_j)} B(z).
\]
The last union is disjoint, so the parent exists uniquely
and belongs to the same pillar. It is either the top or a
site of strictly greater level. Hence every ascent reaches
the top after at most $j$ parent steps.

The nontrivial tower triples $(n,z,x)$, with $x\ne z$,
are exactly $(n(x),z(x),x)$ for non-top bare sites $x$.
Indeed $B(z)\cap\mathbb T^{(n+1)}=\{z\}$ identifies
the level of $x$, and disjointness of the pillars identifies
its top. Conversely the displayed decomposition of a pillar
gives the tower triple of every non-top site. A triple with
$x=z$ imposes only the identity transport along an empty path.

For a triple belonging to $B^j(x_j)$, the comb lies in
$B(z)\subset B^{j-n}(x_j)$. Every bond $c$ of this comb
has $B^n(c_\pm)\subset B^j(x_j)\subset\Omega_0$ at
the level of bare sites. Thus $\mathcal M^nU$ and
$\mathcal M^nU_0$ are defined on the comb by the partial-domain
criterion proved with \eqref{eq:determining-bonds-recursion}.

Extend any gauge $v$ by $\mathbf1$ outside $\Omega_0$.
Gauge covariance and cancellation of the factors at internal
vertices give
\[
 (\mathcal M^n(U^v))(\Gamma_{z,x})
 =v(z)(\mathcal M^nU)(\Gamma_{z,x})v(x)^{-1}.
\]
It follows that $U^v\in\operatorname{Ax}^{\circ}(\mathbf B,U_0)$
if and only if $v$ satisfies \eqref{I-3-98}, with $v$ in
place of $u$, at every site
that is neither a top nor a node. Prescribing $v$ at the
tops and nodes determines it at every other site, since the
parent has greater level.

Under \eqref{I-3-97}, prescribe $u=\mathbf1$ at the nodes
and at $\Lambda_0$. Every top is now assigned a value.
Use \eqref{I-3-98} at each remaining site after its parent
has been assigned. The pillars are disjoint, so every site
receives one value. The resulting gauge is residual and its
transform lies in the required gauge.

If $w$ is another residual gauge with this property, it
agrees with $u$ at all positive-level tops and at all nodes.
The same recursion gives equality throughout their pillars.
The only possible exceptions are non-node sites of
$\Lambda_0$. By \eqref{I-3-97}, none is an endpoint of an
inner bare bond. Thus $w$ and $u$ agree at the endpoints of
every bond of the configuration, and $U^w=U^u$.
\end{proof}

The recursion also records the dependence on $U$ and $U_0$.
For a site $x$ that is neither a node nor a top, let
$y_0=x,\ldots,y_f$ be its free ascent and put
\[
 \tau_s(W)=(\mathcal M^{n(y_s)}W)
                    (\Gamma_{y_{s+1},y_s}),\qquad 0\le s<f,
 \qquad
 \mathcal T_W(x)=\tau_{f-1}(W)\cdots\tau_0(W).
\]
For any gauge $v$ satisfying the recursion with $v$ in place of $u$,
successive substitution gives
\begin{equation}\label{I-3-99}
 v(x)=\mathcal T_{U_0}(x)^{-1}v(y_f)\mathcal T_U(x).
\end{equation}
The ordering is from the endpoint of the free ascent down
to $x$. For the normalized gauge of \cref{I-3-96},
$u(y_f)=\mathbf1$, and hence
$u(x)=\mathcal T_{U_0}(x)^{-1}\mathcal T_U(x)$.
If $x\in B^j(x_j)$, this expression uses only averages
of orders $n<j$ on the indicated combs, whose sites lie in
$B^{j-n(y_s)}(x_j)$. At the prescribed nodes and at
$\Lambda_0$, $u$ is constant. This is the locality of the
gauge construction.
\begin{lemma}\label{lem:tower-top-nodes}
Condition \eqref{I-3-97} holds for every bond set of a
sequence with $\Omega_0=T_\ell$, and for every bond set
of a decreasing sequence satisfying
\eqref{eq:tower-separation}. For the full determining set,
the level-zero clause of \eqref{I-3-97} always holds.
\end{lemma}

\begin{proof}
Use level-zero units. For $x_j\in\Lambda_j$, choose its
block centered at $z\in\Lambda_{K-j-1,\ell}$. This auxiliary
coarser lattice exists also at $j=K$: it is
$\Lambda_{-1,\ell}$, and $m\ge1$. Write
\[
 x_j=z+L^j\sum_{\mu=1}^3 n_\mu e_\mu,
 \qquad |n_\mu|\le h=(L-1)/2.
\]
Put $\epsilon=-1$ if $n_1>0$ and $\epsilon=1$ otherwise.
Then $x'=x_j+\epsilon L^je_1$ is a different site of the
same next-level block, because $h\ge1$. Since
$x_j\notin\Omega_{j+1}$ and this region is a union of
next-level cubes, $x'\notin\Omega_{j+1}$; at $j=k$
this follows from $\Omega_{k+1}=\emptyset$.
Whenever $x'\in\Omega_0$, the bond joining $x_j$ to $x'$
lies in $\Gamma_j^\flat$ and therefore in every bond set
of the sequence. It makes $x_j$ a node.

For $\Omega_0=T_\ell$ the required inclusion is automatic.
For a separated sequence and $j\ge1$,
\[
 |x'-x_j|_\infty=L^j
 <\operatorname{dist}_\infty(\Omega_{j-1}^c,\Omega_j),
\]
so $x'\in\Omega_{j-1}\subset\Omega_0$.
For $j=0$, suppose $x_0$ has an inner-bond neighbor $x''$.
If $x''\notin\Omega_1$, that bond itself belongs to
$\Gamma_0^\flat$. If $x''\in\Omega_1$, then
\[
 |x'-x''|_\infty\le2<L\le R_1\mathfrak M
 <\operatorname{dist}_\infty(\Omega_0^c,\Omega_1),
\]
so again $x'\in\Omega_0$. This proves both clauses of
\eqref{I-3-97} for every bond set in the stated cases.
For the full determining set, any inner bare bond with an
endpoint in $\Lambda_0$ already belongs to $\Gamma_0$,
which proves the final assertion without separation.
\end{proof}
\subsection{Conditions at fixed nodes}

For a separated sequence and the full determining set, an inward
bond fixes a site inside the next block, and the block center
is also a node. The transport between them is therefore invariant
under the residual gauge group. The following proposition identifies exactly the
tower conditions that remain after gauge fixing outside
the nodes.

Under \eqref{I-3-97}, let $U^u$ be the representative in
$\operatorname{Ax}^{\circ}(\mathbf B,U_0)$.
Its orbit meets the full tower gauge
$\operatorname{Ax}(\mathbf B,U_0)$ if and only if $U^u$
satisfies the omitted conditions at nontrivial triples whose
endpoint $x$ is a node. The intersection then consists only
of $U^u$. This follows because the full gauge is a subset
of the gauge outside the nodes.

\begin{proposition}\label{I-3-94}
Suppose that the decreasing sequence satisfies the
separation condition \eqref{eq:tower-separation}, in
level-zero units, and let $\mathbf B$ be any of its bond
sets. The nontrivial conditions in \eqref{I-3-93} omitted
from $\operatorname{Ax}^\circ(\mathbf B,U_0)$ are exactly
those at triples
\[
 (n,z,x)=(j-1,x_j,b_+),\qquad 1\le j\le k,
\]
where $b\in\mathbf B$ is an inward bond of level $j-1$,
oriented with $b_-\in\Omega_{j-1}\setminus\Omega_j$ and
$b_+\in\Omega_j$, and $b_+\in B(x_j)$.
For every such bond,
\[
 B(x_j)\subset\Omega_j,\qquad
 x_j\in\Lambda_j(\mathbf B),\qquad b_+\ne x_j.
\]
Both $x_j$ and $b_+$ are nodes of $\mathbf B$. No bond of
$\Gamma_{x_j,b_+}$ belongs to the full determining set,
and the transport
\[
 (\mathcal M^{j-1}U)(\Gamma_{x_j,b_+})
\]
is constant on the residual-gauge orbit of $U$.

In particular, a residual orbit meets the absolute tower
axial gauge if and only if all these transports are
$\mathbf1$. The intersection then consists of its unique
representative in
$\operatorname{Ax}^\circ(\mathbf B,\mathbf1)$.
\end{proposition}

\begin{proof}
Write $h=(L-1)/2$. Since $\Omega_j$ is a union of cubes
of side at least $L^j$, the block containing $b_+$ lies
in $\Omega_j$. Thus $b_-\notin B(x_j)$. Every
level-$(j-1)$ neighbor of $x_j$ lies in this block, so
$b_+\ne x_j$. If $j<k$, then
\[
 |x_j-b_-|_\infty
 \le(h+1)L^{j-1}<L^j
 \le R_1\mathfrak M L^j
 <\operatorname{dist}_\infty(\Omega_j^c,\Omega_{j+1}).
\]
As $b_-\in\Omega_j^c$, this excludes
$x_j\in\Omega_{j+1}$. The same exclusion is automatic
at $j=k$, because $\Omega_{k+1}=\emptyset$.
Consequently $x_j\in\Lambda_j(\mathbf B)$ and is a
node by \cref{lem:tower-top-nodes}.
The endpoint $b_+$ is a node because $b$ is retained.

The comb stays inside $B(x_j)\subset\Omega_j$, so its
bonds have no endpoint in
$\Lambda_{j-1}(\mathbf B)$. They therefore do not belong
to the full member $\Gamma_{j-1}$. A bond of this lattice
is a bond of no other level, so none belongs to the full
determining set. Gauge covariance of the averages gives
\[
 (\mathcal M^{j-1}U^w)(\Gamma_{x_j,b_+})
 =w(x_j)(\mathcal M^{j-1}U)(\Gamma_{x_j,b_+})w(b_+)^{-1}.
\]
For a residual transformation both endpoint factors are
$\mathbf1$, proving the invariance.

It remains to identify all the omitted conditions. An
inward bond as above gives a triple belonging to the
pillar with top $x_j$, and its terminal point is a node.
Conversely, suppose that a triple $(n,z,x)$ belongs to
the pillar with top $x_j$ and that $x\ne z$ is a node.
The pillar decomposition gives $n=n(x)<j$ and
$x\in\Omega_j\setminus\Omega_{j+1}$. Some retained
bond $c$ of level $i$ has endpoint $x$. Since
$x\in\mathbb T^{(i)}$ but
$x\notin\mathbb T^{(n(x)+1)}$, we have $i\le n(x)<j$.
The other endpoint of $c$ must lie in
$\Omega_i\setminus\Omega_{i+1}$, whereas
$x\in\Omega_j\subset\Omega_{i+1}$. Thus $c$ is inward.
Applying the preceding geometry to $c$ shows that $x$
lies in a pillar with top $x_{i+1}\in\Lambda_{i+1}$,
and is a noncentral site of $B(x_{i+1})$. Distinct
pillars are disjoint; hence $i+1=j$,
$x_{i+1}=x_j$, $n=j-1$, and $z=x_j$.

By \cref{I-3-96}, the orbit has a unique representative
satisfying the conditions outside the nodes. It satisfies
the full tower conditions precisely when it satisfies
the remaining conditions just classified. Their orbit
invariance proves the final assertion.
\end{proof}
For $\mathbf B^\flat$, all nodes belong to
$\bigcup_j\Lambda_j$ by \cref{I-3-81}, whereas a nontrivial
tower triple has an endpoint that is not a pillar top.
Thus no nontrivial condition is omitted:
\[
 \operatorname{Ax}^{\circ}(\mathbf B^\flat,U_0)
 =\operatorname{Ax}(\mathbf B^\flat,U_0).
\]
Every ascent for the reduced set has its entire length free.
When \eqref{I-3-97} holds, each residual orbit therefore
has exactly one representative in the full tower gauge.

For the constant sequence $\Omega_0=\cdots=\Omega_k=T_\ell$,
the layers below $k$ are empty, its bond set consists of all
level-$k$ bonds, and its nodes are precisely
$\mathbb T^{(k)}$. A nontrivial tower triple has endpoint
in $\mathbb T^{(n)}\setminus\mathbb T^{(n+1)}$ for
$n<k$, so again no condition is omitted. Condition
\eqref{I-3-97} holds by \cref{lem:tower-top-nodes}.

Finally, if $U_0\in\operatorname{Ax}^{\circ}(\mathbf B,\mathbf1)$,
the defining transports of $U_0$ at all the retained triples
are $\mathbf1$. Hence
$\operatorname{Ax}^{\circ}(\mathbf B,U_0)
 =\operatorname{Ax}^{\circ}(\mathbf B,\mathbf1)$,
and the representatives in these two gauges agree.
\section{The current and regular critical configurations}
\label{sec:current-critical}
The current is the first variation of the Wilson action with
respect to a bond potential. We establish this identity with
the spacing explicit, and then use the current to define
regular configurations in a multiscale constraint set.
The covariant derivatives and current follow Ba{\l}aban
\cite[pp.~390--392, (3.3)--(3.12)]{B-CMP99b}; we use the
ordinary two-by-two trace and display its factor one-half
explicitly.

\subsection{Forms and the current}
Let $\Lambda$ be one of the lattices $\Lambda_{j,\ell}$
with $j>-m$, and let $\xi$ be its spacing in the units in use.
It has more than two sites along every coordinate axis,
so each ordered nearest-neighbor pair has a unique direction
and orientation.

\begin{definition}\label{I-2-2}
A one-form $B$ assigns an element of $\mathfrak g^c$ to
each positive bond and is extended by $B(\bar b)=-B(b)$.
Write $B_\mu(x)=B(\langle x,x+\xi e_\mu\rangle)$.
A two-form assigns a value $F_{\mu\nu}(x)\in\mathfrak g^c$
to each plaquette $p_{\mu\nu}(x)$, $\mu<\nu$, and is
extended by $F_{\nu\mu}=-F_{\mu\nu}$. A site function
has values in $\mathfrak g^c$. These objects are real when
their values lie in $\mathfrak g$.

For $X,Y\in\mathfrak g^c$, put
$\langle X,Y\rangle=\tfrac12\tr(XY)$. On site functions,
one-forms and two-forms, respectively, use the pairings
\[
 (\lambda,\lambda')_\xi=\xi^3\sum_x
                  \langle\lambda(x),\lambda'(x)\rangle,
 \quad
 (B,C)_\xi=\xi^3\sum_b\langle B(b),C(b)\rangle,
 \quad
 (F,H)_\xi=\xi^3\sum_p\langle F(p),H(p)\rangle.
\]
The bond sum is over positive bonds, and the plaquette sum
counts each plaquette once. These pairings are complex
bilinear. Their restrictions to real functions and forms
are real inner products.

For a real or complex configuration $U$, set
\[
 (\nabla^U_\mu f)(x)=\xi^{-1}
 \bigl[\operatorname{Ad}(U(\langle x,x+\xi e_\mu\rangle))
                     f(x+\xi e_\mu)-f(x)\bigr].
\]
The derivative acts componentwise on forms. Its exterior
derivative on one-forms is
\[
 (D_UB)_{\mu\nu}=\nabla^U_\mu B_\nu-\nabla^U_\nu B_\mu.
\]
The transpose $D_U^{\mathsf T}$ is defined by
$(D_UB,F)_\xi=(B,D_U^{\mathsf T}F)_\xi$.
\end{definition}

The pairings are nondegenerate: if $X\ne0$ then
$X^*\in\mathfrak g^c$ and
$\langle X,X^*\rangle=\tfrac12\tr(XX^*)>0$.
Thus the transpose exists and is unique on these finite
dimensional spaces. For real $U$ it is the real adjoint,
also denoted $D_U^*$ on two-forms.

For $p=p_{\mu\nu}(x)$, $\mu<\nu$, denote its positive
bonds by
\begin{equation}\label{I-2-3}
\begin{aligned}
 b_1&=\langle x,x+\xi e_\mu\rangle,&
 b_2&=\langle x+\xi e_\mu,x+\xi e_\mu+\xi e_\nu\rangle,\\
 b_3&=\langle x+\xi e_\nu,x+\xi e_\mu+\xi e_\nu\rangle,&
 b_4&=\langle x,x+\xi e_\nu\rangle.
\end{aligned}
\end{equation}
Thus $\partial p=(b_1,b_2,\bar b_3,\bar b_4)$ and
\begin{equation}\label{I-2-4}
 (D_UB)(p)=\xi^{-1}\bigl[B(b_1)
 +\operatorname{Ad}(U(b_1))B(b_2)
 -\operatorname{Ad}(U(b_4))B(b_3)-B(b_4)\bigr].
\end{equation}
The transpose is
\begin{equation}\label{I-2-5}
 (D_U^{\mathsf T}F)_\mu(x)=\xi^{-1}
 \sum_{\nu\ne\mu}\left[
 F_{\mu\nu}(x)
 -\operatorname{Ad}(U(\langle x-\xi e_\nu,x\rangle)^{-1})
                         F_{\mu\nu}(x-\xi e_\nu)\right].
\end{equation}
Indeed, write $(D_UB,F)_\xi$ as the sum over ordered
pairs $\mu\ne\nu$ of
$\xi^3\langle\nabla^U_\nu B_\mu,F_{\nu\mu}\rangle$.
In its translated term use
$\langle\operatorname{Ad}(V)X,Y\rangle
 =\langle X,\operatorname{Ad}(V^{-1})Y\rangle$
and replace $x$ by $x-\xi e_\nu$. The coefficient of
$B_\mu(x)$ is the right side of \eqref{I-2-5}, since
$F_{\nu\mu}=-F_{\mu\nu}$.

For a complex configuration $U$, define its field strength
and current by
\begin{equation}\label{I-2-7}
 F_U(p)=\frac{\partial U(p)-\partial U(p)^{-1}}{2i\xi^2},
 \qquad J(U)=D_U^{\mathsf T}F_U.
\end{equation}
For $X\in\mathrm{SL}(2,\mathbb C)$, Cayley--Hamilton gives
\begin{equation}\label{I-2-6}
 X+X^{-1}=(\tr X)\mathbf1.
\end{equation}
Hence $F_U$ is traceless, and for real $U$ it is Hermitian.
The formulas for $D_U$ and its transpose preserve real
forms, so $J(U)$ is real in this case.

Each positive bond belongs to four plaquettes: for each
transverse direction $\nu$, the two plaquettes based at
$x$ and $x-\xi e_\nu$. Formula \eqref{I-2-5} shows that
$J(U)$ at this bond uses only the variables on those four
plaquettes. If $U$ is given on a reversal-invariant partial
bond set, $F_U(p)$ is defined when all four bonds of $p$
are given, and $J(U)(b)$ is defined when all four plaquettes
containing $b$ are defined. These domain conditions will be
retained for regional configurations.

\subsection{Variation of the action}
For a complex configuration, use the same polynomial action
$a_p(U)=1-\tfrac12\tr\partial U(p)$ and
$A^\xi(c,U)=\xi^{-1}\sum_pc(p)a_p(U)$.
An inverse in $G^c$ is polynomial in the matrix entries by
\eqref{I-2-6}. The perturbation $e^{it\xi B}U$ is defined
at positive bonds by
\[
 (e^{it\xi B}U)(b)=e^{it\xi B(b)}U(b),
\]
and at reverse bonds by inversion. It is real for real
$t$, $B$ and $U$.

\begin{proposition}[The current as a first variation]\label{I-2-14}
Let $U$ be a complex configuration on $\Lambda$, let $B$
be a $\mathfrak g^c$-valued one-form, and let $c$ be a
complex plaquette weight. Then
\begin{equation}\label{I-2-16}
\begin{aligned}
 \left.\frac{d}{dt}\right|_{t=0}
 A^\xi(c,e^{it\xi B}U)
 &=(D_UB,cF_U)_\xi
  =(B,D_U^{\mathsf T}(cF_U))_\xi,\\
 \left.\frac{d}{dt}\right|_{t=0}
 A^\xi(e^{it\xi B}U)&=(B,J(U))_\xi.
\end{aligned}
\end{equation}
At the identity configuration,
\[
 \left.\frac{d^2}{dt^2}\right|_{t=0}
 A^\xi(c,e^{it\xi B})=(cD_{\mathbf1}B,D_{\mathbf1}B)_\xi.
\]

Let $\Omega_0$ be a region whose cubes have side at least
the lattice spacing, and let $U'$ be defined on its inner
bonds. If $B$ vanishes outside the bonds where $J(U')$ is
defined, then
\begin{equation}\label{I-2-17}
 \left.\frac{d}{dt}\right|_{t=0}
 A^\xi_{\Omega_0}(e^{it\xi B}U')=(B,J(U'))_\xi,
\end{equation}
where the pairing is summed over that current domain.

For a complex gauge transformation $u$,
\[
 F_{U^u}(p)=\operatorname{Ad}(u(x_{\mathrm b}(p)))F_U(p),
 \qquad
 J(U^u)(b)=\operatorname{Ad}(u(b_-))J(U)(b)
\]
at positive bonds. The weighted action is gauge invariant.
The same covariance holds on every partial domain where
these quantities are defined.
\end{proposition}

\begin{proof}
The SU(2) matrix identity underlying the calculation extends
to $G^c$. For $X\in G^c$, put
$\theta=\tfrac12\tr X$ and $H=(X-X^{-1})/(2i)$.
Equation \eqref{I-2-6} gives
$X=\theta\mathbf1+iH$ and
$H^2=(1-\theta^2)\mathbf1$. Applying this to the plaquette
variable gives
\begin{equation}\label{I-2-15}
\begin{aligned}
 \partial U(p)&=(1-a_p(U))\mathbf1+i\xi^2F_U(p),\\
 \xi^4\langle F_U(p),F_U(p)\rangle
 &=a_p(U)(2-a_p(U)).
\end{aligned}
\end{equation}
Multiplication by $c(p)/(2\xi)$ and summation yield
\begin{equation}\label{eq:action-field-strength-splitting}
 A^\xi(c,U)=\tfrac12(cF_U,F_U)_\xi
       +\frac1{2\xi}\sum_pc(p)a_p(U)^2.
\end{equation}

For the first variation, use the four positive bonds in
\eqref{I-2-3}, and write $U_i=U(b_i)$, $B_i=B(b_i)$
and $P=\partial U(p)$. The perturbed plaquette variable is
\[
 e^{it\xi B_1}U_1e^{it\xi B_2}U_2
 U_3^{-1}e^{-it\xi B_3}U_4^{-1}e^{-it\xi B_4}.
\]
Its derivative at zero equals
\[
 i\xi\bigl[B_1P+(\operatorname{Ad}(U_1)B_2)P
             -P\operatorname{Ad}(U_4)B_3-PB_4\bigr].
\]
By cyclicity of the trace and \eqref{I-2-4}, the trace
of this expression is $i\xi^2\tr(P(D_UB)(p))$.
The scalar part of $P$ in \eqref{I-2-15} has zero pairing
with the traceless matrix $(D_UB)(p)$. The trace is therefore
$-\xi^4\tr(F_U(p)(D_UB)(p))$, and
\[
 \left.\frac{d}{dt}\right|_{t=0}a_p(e^{it\xi B}U)
   =\xi^4\langle(D_UB)(p),F_U(p)\rangle.
\]
Multiplying by $\xi^{-1}c(p)$ and summing proves
\eqref{I-2-16}, first by the definition of the transpose
and then by \eqref{I-2-7}.

For the regional identity, extend $U'$ by $\mathbf1$
outside the inner bonds and take $c$ to be the indicator
of the inner plaquettes. Then
$A^\xi_{\Omega_0}(e^{it\xi B}U')
 =A^\xi(c,e^{it\xi B}U)$.
By \eqref{I-2-4}, $D_UB$ vanishes at a plaquette unless
one of its bonds carries a nonzero value of $B$.
Every such bond is in the current domain of $U'$, so all
four plaquettes containing it are inner. Thus $c=1$
wherever $D_UB$ can be nonzero, and
\[
 (D_UB,cF_U)_\xi=(D_UB,F_U)_\xi=(B,J(U))_\xi.
\]
On the support of $B$, the four plaquettes defining the
current use only $U'$, so $J(U)=J(U')$ there. This proves
\eqref{I-2-17}.

At $U=\mathbf1$, expansion of the same four factors gives
\[
 \partial(e^{it\xi B})(p)
    =\mathbf1+it\xi^2(D_{\mathbf1}B)(p)+O(t^2),
 \qquad
 F_{e^{it\xi B}}(p)=t(D_{\mathbf1}B)(p)+O(t^2).
\]
Also $a_p(e^{it\xi B})=O(t^2)$, because its value and
first derivative vanish at zero. Equation
\eqref{eq:action-field-strength-splitting} therefore gives
\[
 A^\xi(c,e^{it\xi B})
   =\tfrac12t^2(cD_{\mathbf1}B,D_{\mathbf1}B)_\xi+O(t^3),
\]
which proves the second-variation formula. These finite
products of matrix exponentials are entire in $t$.

Finally, a plaquette variable is conjugated at its base
point by a gauge transformation. This proves the action
invariance and the transformation of $F_U$. In
\eqref{I-2-5}, set $y=x-\xi e_\nu$ and use
\[
 U^u(\langle y,x\rangle)^{-1}u(y)
       =u(x)U(\langle y,x\rangle)^{-1}.
\]
Each summand defining $J(U^u)_\mu(x)$ is conjugated by
$u(x)$, which proves the current covariance. The calculation
uses only the four incident plaquettes, so it also applies
on the stated partial domains.
\end{proof}

For a general variation on the inner bonds of a region,
the weighted identity \eqref{I-2-16}, with the indicator
of its inner plaquettes, remains the regional variation
formula. The unweighted pairing with $J(U')$ in
\eqref{I-2-17} uses the stated support condition.

If the spacing in a second system of units is
$\xi'=\upsilon\xi$, the same bond perturbation is written
with $B'=\upsilon^{-1}B$. Directly from the formulas,
\[
 F'_U=\upsilon^{-2}F_U,\qquad
 J'(U)=\upsilon^{-3}J(U),\qquad
 (B',J'(U))_{\xi'}=\upsilon^{-1}(B,J(U))_\xi.
\]
This agrees with
$A^{\xi'}(c,U)=\upsilon^{-1}A^\xi(c,U)$:
the pairing contributes the cell weight $\xi'^3$, while
its two arguments contribute $\upsilon^{-1}$ and
$\upsilon^{-3}$.

\subsection{Regularity at the local block scale}
Let $\mathbf B$ be a bond set of
$\Omega_0\supset\cdots\supset\Omega_k$.

\begin{definition}\label{def:local-regularity}
For a bare plaquette or bond $\mathfrak e$ meeting
$\Omega_0$ at a corner or endpoint, respectively, let $j$
be the largest index for which such a corner or endpoint
lies in $\Omega_j$. Its local block side, in level-$k$
units, is
\[
 \mathsf L_{\mathbf B}(\mathfrak e)=L^{j-k}.
\]
Put $\eta_k=L^{-k}$. A real configuration on the inner
bare bonds of $\Omega_0$ is \emph{regular at radius} $r>0$
if
\begin{equation}\label{eq:regional-regularity}
\begin{aligned}
 |\partial U(p)-\mathbf1|
     &<r\eta_k^2\mathsf L_{\mathbf B}(p)^{-2}
       &&\text{for every inner bare plaquette }p,\\
 |J(U)(b)|&<r\mathsf L_{\mathbf B}(b)^{-3}
       &&\text{wherever the current is defined.}
\end{aligned}
\end{equation}
The current is computed with bare spacing $\eta_k$.
Denote this set by $\mathfrak U_k(\mathbf B,r)$.

\end{definition}


For the constant whole-torus sequence all local block sides
equal one. Thus $\mathfrak U_k(\mathbf B,r)$ consists of
the configurations with
\[
 |\partial U(p)-\mathbf1|<r\eta_k^2,
 \qquad |J(U)(b)|<r
\]
at every bare plaquette and positive bond.
The following comparison relates this condition to a
bound on the divergence of the plaquette variables.

\begin{lemma}\label{I-3-113}
Let $1\le k\le K$, put $\eta=\eta_k$, and let
$\mathbf B$ be the constant-sequence determining set.
Define $\mathfrak U'_k(\mathbf B,r)$ to be the set of real
configurations satisfying
\[
 |\partial U(p)-\mathbf1|<r\eta^2,
 \qquad |(D_U^{\mathsf T}\partial U)(b)|<r\eta^2
\]
at every bare plaquette and positive bond. Here
$D_U^{\mathsf T}$ is the formula \eqref{I-2-5}, extended
to matrix-valued two-forms; the plaquette variables are
extended by $X_{\nu\mu}=-X_{\mu\nu}$ after setting
$X_{\mu\nu}(x)=\partial U(p_{\mu\nu}(x))$ for $\mu<\nu$.
Then, for every $r>0$,
\begin{equation}\label{I-3-114}
 \mathfrak U'_k(\mathbf B,r)
 \subset\mathfrak U_k(\mathbf B,r)
 \subset\mathfrak U'_k(\mathbf B,r+r^2).
\end{equation}
\end{lemma}

\begin{proof}
Write $\partial U(p)=c(p)\mathbf1+i\eta^2F_U(p)$,
where $c(p)=\tfrac12\tr\partial U(p)$ is real.
Since conjugation fixes scalar matrices, the extension
of \eqref{I-2-5} gives
\[
 (D_U^{\mathsf T}\partial U)(b)
             =a(b)\mathbf1+i\eta^2J(U)(b).
\]
Here $a(b)$ is the sum of two terms
$\pm\eta^{-1}(c(p)-c(p'))$, one for each transverse
direction, with $p,p'$ the two incident plaquettes in
that coordinate plane. For real $a$ and $H\in\mathfrak g$,
$|a\mathbf1+iH|^2=a^2+|H|^2$; indeed
$H^2=|H|^2\mathbf1$ and
$(a\mathbf1-iH)(a\mathbf1+iH)=(a^2+|H|^2)\mathbf1$.
Thus
\[
 |(D_U^{\mathsf T}\partial U)(b)|^2
           =a(b)^2+\eta^4|J(U)(b)|^2,
\]
which proves the first inclusion.

If $U\in\mathfrak U_k(\mathbf B,r)$, the SU(2) identity
$1-c(p)=\tfrac12|\partial U(p)-\mathbf1|^2$
gives $0\le1-c(p)<\tfrac12r^2\eta^4$.
Consequently $|c(p)-c(p')|<\tfrac12r^2\eta^4$ and
$|a(b)|<r^2\eta^3$. Since $\eta\le1$,
\[
 |(D_U^{\mathsf T}\partial U)(b)|
 \le |a(b)|+\eta^2|J(U)(b)|
 <(r^2\eta+r)\eta^2
 \le(r+r^2)\eta^2.
\]
The plaquette bound with radius $r+r^2$ follows from
the bound with radius $r$, proving the second inclusion.
\end{proof}

\subsection{Critical orbits}
Return to a general bond set $\mathbf B$ of
$\Omega_0\supset\cdots\supset\Omega_k$ and its constraint
set $C(\mathbf B,V)$.

\begin{definition}\label{def:regular-critical}
A configuration $U\in C(\mathbf B,V)$ is \emph{critical}
if
\[
 \left.\frac{d}{dt}\right|_{t=0}A_{\Omega_0}(U_t)=0
\]
for every differentiable curve $t\mapsto U_t$ in
$C(\mathbf B,V)$ with $U_0=U$. Here differentiability means
entrywise differentiability on an open real interval
containing zero. Fix a radius $\varepsilon_{\mathrm v}>0$,
and set
\[
 \mathcal C(\mathbf B,V;r)
   =\{U\in C(\mathbf B,V):U\text{ is critical and }
                     U\in\mathfrak U_k(\mathbf B,r)\},
 \qquad
 \mathcal C(\mathbf B,V)=
                  \mathcal C(\mathbf B,V;\varepsilon_{\mathrm v}).
\]
If $\mathcal C(\mathbf B,V)$ is nonempty and consists
of one orbit of $\mathcal G^{\mathrm{res}}(\mathbf B)$,
call it the \emph{minimal orbit}. This term specifies an orbit
of regular critical points and does not assert global action
minimization. Write $U_{\mathbf B}(V)$
for any element of this orbit; an invariant function of
$U_{\mathbf B}(V)$ has its common value on the orbit.
\end{definition}

Criticality is unchanged if the action is multiplied by the
positive constant $L^k$, and hence can be tested with
$A^{\eta_k}_{\Omega_0}$.
The definition uses curves in the constraint set and does
not require it to have been identified as a smooth manifold.

\begin{proposition}\label{prop:critical-orbit-gauge}
The set $\mathcal C(\mathbf B,V;r)$ is invariant under
$\mathcal G^{\mathrm{res}}(\mathbf B)$.
Suppose the minimal orbit exists and condition
\eqref{I-3-97} holds. For every real reference configuration
$U_0$ on the inner bare bonds of $\Omega_0$, there is exactly one element of this orbit in
$\operatorname{Ax}^{\circ}(\mathbf B,U_0)$.
Denote it by $U^\circ_{\mathbf B}(V;U_0)$, and set
$U^\circ_{\mathbf B}(V)=U^\circ_{\mathbf B}(V;\mathbf1)$.

The absolute tower representative
$U^{\mathrm{ax}}_{\mathbf B}(V)$ exists precisely when
$U^\circ_{\mathbf B}(V)$ satisfies the omitted node
conditions, and then the two representatives agree.
For a reduced bond set under \eqref{I-3-97}, or for
the constant whole-torus sequence, they are both defined
and equal.
\end{proposition}

\begin{proof}
Plaquette and current covariance, with invariance of the operator
norm under unitary conjugation, preserve regularity.
A residual gauge maps $C(\mathbf B,V)$ and its differentiable
curves bijectively onto themselves and preserves the action;
it therefore preserves criticality. If the minimal orbit exists,
\cref{I-3-96} gives its unique outside-node representative.
The full tower-gauge criterion and its reduced and constant
cases are the corresponding conclusions of
\cref{sec:axial-representatives}.
\end{proof}

If $U_0\in\operatorname{Ax}^{\circ}(\mathbf B,\mathbf1)$,
then $\operatorname{Ax}^{\circ}(\mathbf B,U_0)
=\operatorname{Ax}^{\circ}(\mathbf B,\mathbf1)$, so $U^\circ_{\mathbf B}(V;U_0)=U^\circ_{\mathbf B}(V)$.
For a full or partially reduced separated set, the omitted
conditions are exactly the retained seam conditions described
in \cref{I-3-94}.

\subsection{Backgrounds induced by block fields}
The use of constrained critical orbits as background fields
follows Ba{\l}aban \cite[(0.21), p.~256]{B-CMP109};
localization by a tower determining set follows
\cite[(2.12)--(2.16), pp.~256--257]{B-CMP119}.
For $V\in\mathcal V_k$, take the constant sequence
$\Omega_0=\cdots=\Omega_k=T_\ell$. Its determining map
is $\mathcal M^k$ and its bond set consists of all
level-$k$ bonds. Whenever the minimal orbit at $V$ exists,
write
\[
 U_k(V)=U_{\mathbf B}(V).
\]
This notation specifies the orbit through any of its elements;
its invariant quantities, including the action, are functions
of $V$. The unique absolute tower representative is
$U^{\mathrm{ax}}_{\mathbf B}(V)=U^\circ_{\mathbf B}(V)$.
More generally, a sequence beginning at $T_\ell$ and block
fields $V_j\in\mathcal V_j$ supply the datum
$V=((V_j)|_{\mathbf B\cap\Gamma_j})_{j=0}^k$ on a specified
bond set. The same notation $U_k(V)$ is used when that
datum has a minimal orbit.

Let $\Omega$ be a union of $M_1$-cubes at level $k$ and
let $\mathbf B_k(\Omega)$ be the full determining set
of its $\mathfrak r_2$-tower. The localized background is
\begin{equation}\label{I-3-78}
 U_{k,\Omega}(V)=
 U_{\mathbf B_k(\Omega)}
       \bigl(M_{\mathbf B_k(\Omega)}(Q^{s*}_kV)\bigr),
\end{equation}
provided the indicated minimal orbit exists. This is a
configuration on the inner bare bonds of $\Omega$, specified
up to the residual group. Its unique outside-node tower
representative is obtained by replacing $U_{\mathbf B_k(\Omega)}$
in \eqref{I-3-78} by $U^\circ_{\mathbf B_k(\Omega)}$.

The datum in \eqref{I-3-78} is defined for every block field
$V$. Indeed the tower gaps, in level-zero units, satisfy
\[
 \operatorname{dist}_\infty(\Omega_{j-1}^c,\Omega_j)
       \ge M_1L^j>L^j,
\]
so \cref{prop:determining-domain} applies.
By locality of averaging, the restriction of $Q^{s*}_kV$
to the inner bare bonds gives the datum \eqref{I-3-75}.
At a member bond $b$ of level $j$, its value is either
$\mathbf1$ or $V(c)$ for the level-$k$ bond above $b$.
In the latter case
$b_\pm\in B^{k-j}(c_\pm)$; since $b_\pm\in\Omega$
and $\Omega$ is a union of level-$k$ partition cubes,
$c_\pm\in\Omega$. Thus the induced datum, and the
conditional background orbit and its selected representative,
depend only on the inner level-$k$ bonds of $V$ in $\Omega$.

The reduced background $U^\flat_{k,\Omega}(V)$ is defined
by replacing $\mathbf B_k(\Omega)$ with
$\mathbf B_k(\Omega)^\flat$ in \eqref{I-3-78}, whenever
the corresponding minimal orbit exists.
The domain and locality arguments are
unchanged; its absolute and outside-node tower
representatives agree when the minimal orbit exists.

\section{Gauge averages and the Landau condition}
\label{sec:landau-gauge}
To specify a Landau representative relative to a background,
we normalize its gauge transformation by transported block
averages. Their derivative gives the admissible infinitesimal
gauge transformations and hence the projection in the Landau
condition. The averages follow Ba{\l}aban
\cite[(78)--(81), p.~30]{B-CMP98}.

\subsection{Transported gauge averages}
Let $\mathbf B$ be a bond set of
$\Omega_0\supset\cdots\supset\Omega_k$, let $\xi$ be the
bare spacing, and let $U$ be a real configuration on the
inner bare bonds of $\Omega_0$. The averages
$\mathcal M^nU$ needed below are defined on the inner bonds
of $B^{j-n}(y)$ for $y\in\Lambda_j(\mathbf B)$ and $n<j$.
Indeed, for a bond $c$ with endpoints in $B^{j-n}(y)$,
nesting gives $B^n(c_\pm)\subset B^j(y)\subset\Omega_0$
at the level of bare sites. The partial-domain criterion
proved with \eqref{eq:determining-bonds-recursion} therefore
applies.

\begin{definition}\label{I-3-100}
For $0\le n<K$ and $z\in\mathbb T^{(n+1)}$, let $W$ be a configuration
on the bonds of $B(z)$ and let $v$ be a $G^c$-valued
function on its sites. Define
\begin{equation}\label{I-3-101}
 (\bar R_Wv)(z)=v(z)\exp\left\{
 L^{-3}\sum_{x\in B(z)}
 \log\left[v(z)^{-1}
       \operatorname{Ad}(W(\Gamma_{z,x}))v(x)\right]\right\}.
\end{equation}
Here and below the logarithm is its series at $\mathbf1$.
This average is defined when every argument $X$ of a
logarithm satisfies $|X-\mathbf1|<1$. Its value lies in
$G^c$, and lies in $G$ when $W$ and $v$ are real.

Put $\bar R_U^{(0)}v=v$ and
\[
 \bar R_U^{(n+1)}v
   =\bar R_{\mathcal M^nU}(\bar R_U^{(n)}v).
\]
Every stage of an iterated average must be defined.
A gauge transformation $v$ satisfies the
\emph{averaged normalization relative to $(\mathbf B,U)$}
if
\begin{equation}\label{I-3-102}
 (\bar R_U^{(j)}v)(y)=\mathbf1,
 \qquad y\in\Lambda_j(\mathbf B),\quad 0\le j\le k.
\end{equation}
In particular, $v=\mathbf1$ on $\Lambda_0(\mathbf B)$.

For a site function $\lambda$, set
\begin{equation}\label{I-3-103}
 (q_W\lambda)(z)=L^{-3}\sum_{x\in B(z)}
             \operatorname{Ad}(W(\Gamma_{z,x}))\lambda(x).
\end{equation}
Put $q_U^{(0)}=\mathrm{id}$ and
$q_U^{(n+1)}=q_{\mathcal M^nU}q_U^{(n)}$.
The admissible infinitesimal gauge transformations form
the space
\begin{equation}\label{I-3-104}
 \mathcal N(\mathbf B,U)=
 \left\{\lambda:\ (q_U^{(j)}\lambda)(y)=0
       \text{ for }y\in\Lambda_j(\mathbf B),\ 0\le j\le k
 \right\},
\end{equation}
where $\lambda$ is a site function on the bare sites of
$\Omega_0$.
\end{definition}

\begin{proposition}\label{prop:gauge-average-derivative}
With the background averages fixed,
$(\bar R_U^{(j)}v)(y)$ depends on $v$ only through its
values on the pillar $B^j(y)$. Moreover, for every site
function $\lambda$, the averages of $e^{it\lambda}$ and
their logarithms are defined for sufficiently small
$|t|$ wherever all required background averages exist, and
\[
 \left.\frac d{dt}\right|_{t=0}
       \frac1i\log\bar R_U^{(n)}(e^{it\lambda})
       =q_U^{(n)}\lambda.
\]
This is the linearization used in the Landau condition
\cite[p.~80]{B-CMP99a}.
\end{proposition}

\begin{proof}
Each logarithm in \eqref{I-3-101} is traceless, so its
exponential has determinant one. For real $W$ and $v$,
the logarithms lie in $i\mathfrak g$; their average
exponentiates to an element of $G$. This proves the
group-valued assertions. Locality follows inductively
from
\[
 \bigcup_{x\in B(z)}B^n(x)=B^{n+1}(z).
\]

For the derivative, suppose
$w_t=\bar R_U^{(n)}(e^{it\lambda})$ satisfies
$w_0=\mathbf1$ and $\dot w_0=i\mu$, where
$\mu=q_U^{(n)}\lambda$; this holds at $n=0$.
Set $W=\mathcal M^nU$. Every argument of a logarithm
in $\bar R_Ww_t$ equals $\mathbf1$ at $t=0$ and is
therefore in the logarithm domain for small $|t|$.
Its derivative is
\[
 i\left[\operatorname{Ad}(W(\Gamma_{z,x}))\mu(x)
                    -\mu(z)\right].
\]
The differential of $\log$ at $\mathbf1$ is the
identity. Since $B(z)$ has $L^3$ sites, the exponent
in \eqref{I-3-101} has derivative
$i[(q_W\mu)(z)-\mu(z)]$. The prefactor $w_t(z)$
contributes $i\mu(z)$, and these base-point terms
cancel:
\[
 \left.\frac d{dt}\right|_{t=0}(\bar R_Ww_t)(z)
  =i\mu(z)+i[(q_W\mu)(z)-\mu(z)]
  =i(q_U^{(n+1)}\lambda)(z).
\]
This proves the induction. Taking the logarithm at
$\mathbf1$ gives the asserted derivative.
\end{proof}

\subsection{Constraint mismatch}
For a one-form $A$ on the inner bare bonds of $\Omega_0$
and $c\in\mathbf B\cap\Gamma_j$, define
\begin{equation}\label{I-3-105}
 Q_j(U,\xi A)(c)=\frac1i\log\left[
  \mathcal M^j(e^{i\xi A}U)(c)\mathcal M^j(U)(c)^{-1}
                         \right].
\end{equation}
Both averages must be defined, and the logarithm argument
must lie within distance one of $\mathbf1$. For real $A$
use the global real block average; for nonreal $A$ use
the complex branch for the first average. When the
directional derivative exists, put
\[
 Q'_j(U)A=\left.\frac d{dt}\right|_{t=0}Q_j(U,t\xi A).
\]
It follows that
\[
 \mathcal M^j(e^{it\xi A}U)(c)
   =\exp\{it(Q'_j(U)A)(c)+o(t)\}\mathcal M^j(U)(c).
\]
Indeed, $Q_j(U,0)=0$, and exponentiating the defining
logarithm gives this equality from
$Q_j(U,t\xi A)=tQ'_j(U)A+o(t)$.

The derivative includes the coarse bond length
\cite[(1.37), p.~82]{B-CMP99a},
\cite[(3.110), p.~417]{B-CMP99b}.
For $\Omega_0=T_\ell$, $U=\mathbf1$, and a constant
one-form $A_\mu(x)=a_\mu$, it is
\[
 (Q'_j(\mathbf1)A)(c)=L^j\xi a_\mu
       \quad\text{if $c$ has direction $\mu$.}
\]
To see this, apply \cref{I-3-59} to
$e^{it\xi A}$. One application of
$\mathbf Q^{\mathrm{lin}}$ multiplies a constant
component by $L$, so $j$ iterations give $L^j\xi A$.
The real and complex block averages agree near the
identity for real $A$ by the same lemma.

\subsection{The Landau condition}
The Landau condition requires the divergence to be orthogonal
to the image of the admissible infinitesimal gauge transformations
under the Dirichlet Laplacian. The site operators follow
Ba{\l}aban \cite[pp.~390--393, (3.3), (3.8)]{B-CMP99b}.
Let $\Omega_0$ be a region whose cubes have side at least
the bare spacing $\xi$. A bond of $\Omega_0$ has at least
one endpoint there; an inner bond has both endpoints there.
Initially let $U$ be a real or complex configuration on
all bonds of $\Omega_0$. Site functions on $\Omega_0$ and
one-forms on its bonds are extended by zero outside their
respective domains.

\begin{definition}\label{I-2-8}
The covariant gradient sends site functions to one-forms
on the bonds of $\Omega_0$ and is given on positive bonds by
\begin{equation}\label{I-2-9}
 (D_U\lambda)(b)=\xi^{-1}
 \bigl[\operatorname{Ad}(U(b))\lambda(b_+)-\lambda(b_-)\bigr].
\end{equation}
Its transpose for the bilinear pairings is denoted $D_U^*$:
$(D_U^*A,\lambda)_\xi=(A,D_U\lambda)_\xi$.
At a site $x\in\Omega_0$ it is
\begin{equation}\label{I-2-10}
 (D_U^*A)(x)=\xi^{-1}\sum_{\mu=1}^3
 \left[\operatorname{Ad}(U(\langle x-\xi e_\mu,x\rangle)^{-1})
              A_\mu(x-\xi e_\mu)-A_\mu(x)\right].
\end{equation}
The covariant Laplacian on site functions is
$\Delta_U=D_U^*D_U$.
Here $D_U$ acts on site functions, whereas the exterior
derivative in \eqref{I-2-4} acts on one-forms; the domains
also distinguish their transposes. For real $U$, $D_U^*$
is the adjoint on the real spaces.
\end{definition}

Every bond entering \eqref{I-2-10} has an endpoint in
$\Omega_0$. The formula follows from \eqref{I-2-9}, trace
invariance and a shift of the summation variable in the
translated term. The zero extension of $\lambda$ leaves
only sites of $\Omega_0$ in the resulting sum; nondegeneracy
of the pairings gives uniqueness of the transpose.
Extend a gauge $u$ on $\Omega_0$ by the identity outside,
and let $\operatorname{Ad}(u)$ act at a site on site functions
and at the initial point of a positive bond on one-forms.
Then
\[
 D_{U^u}\operatorname{Ad}(u)=\operatorname{Ad}(u)D_U,
 \qquad
 D_{U^u}^*\operatorname{Ad}(u)=\operatorname{Ad}(u)D_U^*,
 \qquad
 \Delta_{U^u}\operatorname{Ad}(u)=\operatorname{Ad}(u)\Delta_U.
\]
Cancellation of endpoint factors proves the first identity;
trace invariance gives the second, and composition the third.
The same algebra applies to complex gauges.

Let $\mathbf B$ be a bond set of a decreasing sequence
$\Omega_0\supset\cdots\supset\Omega_k$ satisfying the
separation condition at scale $\mathfrak r_1$, or the
determining set of the constant sequence
$\Omega_0=\cdots=\Omega_k=T_\ell$. Let $U_0$ be a real
configuration on the inner bare bonds of $\Omega_0$.
No regularity assumption on $U_0$ is needed for the
following definitions.

Extend $U_0$ by the identity to bonds with exactly one
endpoint in $\Omega_0$, site functions by zero outside
$\Omega_0$, and inner one-forms by zero to the other bonds.
The intermediate gradient in $\Delta_{U_0}=D_{U_0}^*D_{U_0}$
is still computed on all incident bonds. At a site of
$\Omega_0$,
\[
\begin{split}
 (\Delta_{U_0}\lambda)(x)=\xi^{-2}\sum_{\mu=1}^3
 \bigl[2\lambda(x)
 &-\operatorname{Ad}(U_0(\langle x-\xi e_\mu,x\rangle)^{-1})
                  \lambda(x-\xi e_\mu)\\
 &-\operatorname{Ad}(U_0(\langle x,x+\xi e_\mu\rangle))
                  \lambda(x+\xi e_\mu)\bigr].
\end{split}
\]
Each background value on a non-inner bond multiplies a
vanishing exterior value of $\lambda$, so this expression
does not depend on the chosen background extension.
The same holds for $D_{U_0}^*A$, since $A$ vanishes on
non-inner bonds, and, at $x\in\Omega_0$, for
\[
 (\nabla^{U_0}_\mu A_\nu)(x)=\xi^{-1}
 \bigl[\operatorname{Ad}(U_0(\langle x,x+\xi e_\mu\rangle))
 A_\nu(x+\xi e_\mu)-A_\nu(x)\bigr]:
\]
if $x+\xi e_\mu$ is outside $\Omega_0$, the one-form
component there vanishes. We use these Dirichlet conventions
throughout the Landau construction.

\begin{definition}\label{I-3-106}
Following Ba{\l}aban \cite[p.~80, (1.27); p.~82, (1.38)]{B-CMP99a},
let $R(U_0)=R(\mathbf B,U_0)$ be the orthogonal projection
in the real space of site functions, with inner product
$(\cdot,\cdot)_\xi$, onto
\[
 \Delta_{U_0}
 \bigl(\mathcal N(\mathbf B,U_0)
       \cap\{\text{real site functions}\}\bigr).
\]
Extend $R(U_0)$ complex-linearly to complex site functions.
A one-form $A$ satisfies the Landau condition relative to
$(\mathbf B,U_0)$ if
\begin{equation}\label{I-3-107}
 R(U_0)D_{U_0}^*A=0.
\end{equation}
For real $A$, this is equivalent to
\[
 (D_{U_0}^*A,\Delta_{U_0}\lambda)_\xi=0
 \quad\text{for every real }
 \lambda\in\mathcal N(\mathbf B,U_0).
\]
\end{definition}

For real $U_0$, the averages and site operators preserve real
functions. Their displayed image is a subspace of a finite-dimensional
real inner-product space, so its orthogonal projection exists and
has the orthogonal complement as kernel, giving the equivalence.

Let $U$ be a real configuration on the inner bare bonds
of $\Omega_0$. A \emph{Landau pair of $U$ relative to
$(\mathbf B,U_0)$} is a real one-form $A$ on these bonds
and a $G$-valued gauge transformation $u$ on the bare
sites of $\Omega_0$ such that
\begin{equation}\label{I-3-110}
 \begin{gathered}
 U=(e^{i\xi A}U_0)^u,\qquad R(U_0)D^*_{U_0}A=0,\\
 (\bar R_{U_0}^{(j)}u)(y)=\mathbf1,
      \qquad y\in\Lambda_j(\mathbf B),\quad 0\le j\le k.
 \end{gathered}
\end{equation}
The last condition includes the definedness of every
averaged normalization. The first two relations are equivalent to
\begin{equation}\label{I-3-108}
 e^{i\xi A}U_0=U^{u^{-1}},
 \qquad R(U_0)D^*_{U_0}A=0.
\end{equation}
We call $e^{i\xi A}$ a \emph{Landau representative of $U$
relative to $(\mathbf B,U_0)$}.
As in \cref{sec:current-critical}, the product
$e^{i\xi A}U_0$ is defined on positive bonds and extended
to reversed bonds by inversion of the whole product.

Set $\Omega_{k+1}=\emptyset$. The small-field conditions
at radius $\alpha>0$ are the following inequalities:
for every $0\le j\le k$ and every positive inner bare bond
$b=\langle x,x+\xi e_\nu\rangle$ with
$\{b_-,b_+\}\cap(\Omega_j\setminus\Omega_{j+1})\ne\emptyset$,
\begin{equation}\label{I-3-109}
 L^j\xi\,|A(b)|<\alpha,
 \qquad
 (L^j\xi)^2\max_{1\le\mu\le3}
       |(\nabla^{U_0}_\mu A_\nu)(x)|<\alpha.
\end{equation}
The direction $\nu$ is fixed by $b$; the maximum is only
over $\mu$. Both inequalities are imposed for every layer
containing an endpoint of $b$.

The constraint mismatch of the Landau representative is the family
$Q_j(U_0,\xi A)$ on $\mathbf B\cap\Gamma_j$,
$0\le j\le k$, whenever these maps are defined.
Fix a radius $\alpha>0$. If there is exactly one Landau
pair whose first member satisfies \eqref{I-3-109} at
that radius, denote this member by
$A^{\mathrm{La}}(U;\mathbf B,U_0)$.

For a real one-form $A$ on the inner bare bonds of
$\Omega_0$, the Wilson Hessian at $U_0$ is the quadratic
form
\begin{equation}\label{I-3-112}
 \begin{aligned}
 (A,\Delta^{\mathrm W}_{U_0}A)_\xi
   &=\left.\frac{d^2}{dt^2}\right|_{t=0}
                   A^\xi_{\Omega_0}(e^{it\xi A}U_0)\\
   &=\xi^{-1}\left.\frac{d^2}{dt^2}\right|_{t=0}
                   A_{\Omega_0}(e^{it\xi A}U_0).
 \end{aligned}
\end{equation}
At the identity background, the weighted second-variation
formula of \cref{I-2-14}, with the indicator of the inner
plaquettes as weight, gives
\[
 (A,\Delta^{\mathrm W}_{\mathbf1}A)_\xi
  =\xi^3\sum_{p\text{ inner to }\Omega_0}
       \langle(D_{\mathbf1}A)(p),(D_{\mathbf1}A)(p)\rangle.
\]
Here $D_{\mathbf1}$ is the exterior derivative on
one-forms.

\section{Continuation from real background data}
\label{sec:background-continuation}
Complex tower configurations enter the domains for analytic
fluctuations, following Ba{\l}aban
\cite[pp.~261--263, (1.10)--(1.16)]{B-CMP109}.
They are obtained by holomorphic continuation of the selected
representatives of real regular critical orbits. The two
existence hypotheses below define this continuation; the
identity principle then proves its uniqueness and continues
the block constraint on the specified domains.
For the reduced determining sets and radii of
\cref{I-58-22}, these existence hypotheses are established
in \cref{sec:canonical-background}.

\subsection{Perturbations of a real datum}
Let $\mathbf B$ be a bond set of a decreasing sequence
$\Omega_0\supset\cdots\supset\Omega_k$ satisfying the
separation condition of the rule $\mathfrak r_1$, or
the determining set of the constant sequence
$\Omega_0=\cdots=\Omega_k=T_\ell$.
A complex datum on $\mathbf B$ is a family
$\mathbf V=(\mathbf V_j)_{j\le k}$ of $G^c$-valued
functions on its members, satisfying
$\mathbf V_j(\bar b)=\mathbf V_j(b)^{-1}$. It is real
when its values lie in $G$.

Fix a real datum $V_0$. A perturbation
$B=(B_j)_{j\le k}$ assigns an element of $\mathfrak g^c$
to each positive bond of each member. Define
$e^{iB}V_0$ on positive bonds by
\[
 (e^{iB}V_0)_j(b)=e^{iB_j(b)}V_{0,j}(b),
 \qquad b\in\mathbf B\cap\Gamma_j,
\]
and on reversed bonds by inversion of the whole product.
Real perturbations have values in $\mathfrak g$ and
give real data.

For $r>0$, put
\[
 \mathcal B_r=
 \left\{B:\ \sup_{j,b}|B_j(b)|<r\right\}.
\]
Since $\mathfrak g^c=\mathfrak g\oplus i\mathfrak g$ over
$\mathbb R$, a real basis identifies the finite-dimensional
perturbation space with $\mathbb C^N$ and its real slice with
$\mathbb R^N$. The set $\mathcal B_r$ is open and convex, hence
connected. Holomorphy means holomorphy of the matrix entries on
positive inner bonds; inversion gives the values on their reverses.

\begin{definition}\label{I-3-115}
The continuation hypotheses at $(V_0,r_{\mathrm B})$,
where $r_{\mathrm B}>0$, are the following:
\begin{enumerate}[label=\textup{(\roman*)}]
\item For every real $B\in\mathcal B_{r_{\mathrm B}}$,
the set $\mathcal C(\mathbf B,e^{iB}V_0)$ is nonempty
and consists of one orbit of
$\mathcal G^{\mathrm{res}}(\mathbf B)$.
\item The map
$B\mapsto U^\circ_{\mathbf B}(e^{iB}V_0)$ on those
real perturbations is the restriction of a holomorphic
map from $\mathcal B_{r_{\mathrm B}}$ to the
$G^c$-valued configurations on the inner bare bonds of
$\Omega_0$.
\end{enumerate}
Under these hypotheses, write
$\mathcal B(V_0)=\mathcal B_{r_{\mathrm B}}$ and denote
the holomorphic extension by $\Phi_{V_0}$. Its value is
the \emph{representative continued from $V_0$}:
\begin{equation}\label{I-3-118}
 U_{\mathbf B}(V_0;B)
   =\Phi_{V_0}(B)=U^\circ_{\mathbf B}(V_0;B),
       \qquad B\in\mathcal B(V_0).\qedhere
\end{equation}
\end{definition}

The real representatives used in \textup{(ii)} are defined by
\cref{prop:critical-orbit-gauge}: the stated geometry
implies \eqref{I-3-97}. They are the outside-node
representatives. At real data where the absolute tower
representative exists, it agrees with this choice; this
holds in particular for a reduced bond set and for the
constant whole-torus sequence.

\begin{proposition}\label{I-3-117}
Under the continuation hypotheses, $\Phi_{V_0}$ is
unique. If those hypotheses hold at two radii, the two
extensions agree on the smaller parameter neighborhood.

Let $\mathcal D\subset\mathcal B(V_0)$ be the set on
which all the complex block averages
$M_{\mathbf B}(\Phi_{V_0}(B))$ are defined. Then
$\mathcal D$ is open and this map is holomorphic there.
Let $\mathcal D'\subset\mathcal D$ be connected and
open. Suppose there is a nonempty relatively open set
$E\subset\mathcal D'\cap\mathbb R^N$ such that the
real and complex averages of
$U^\circ_{\mathbf B}(e^{iB}V_0)$ agree for every
$B\in E$. Then, with the averages on the left taken
in the complex branch,
\[
 M_{\mathbf B}(U_{\mathbf B}(V_0;B))=e^{iB}V_0,
                 \qquad B\in\mathcal D'.
\]
\end{proposition}

\begin{proof}
A holomorphic map on a connected open subset of $\mathbb C^N$
that vanishes on a nonempty relatively open real slice vanishes
identically. At a point of that slice its Taylor coefficients
are the real derivatives divided by their factorials, hence vanish.
The set where the map vanishes locally is open and closed:
at a limit point all derivatives vanish by continuity, so the
local power series is again zero. Connectedness proves the claim,
entrywise for a matrix-valued map.

Apply this principle to the difference of two proposed extensions
on $\mathcal B(V_0)$, or on the smaller ball when two radii are
given. This proves uniqueness and radius compatibility.

By the finite-iteration construction in \cref{sec:block-averaging},
the domain of $M_{\mathbf B}$ is
$(G^c)^{\mathsf P}\cap\mathcal W$, where $\mathsf P$ is the set
of positive inner bare bonds of $\Omega_0$ and $\mathcal W$
is open in the ambient matrix coordinates; the averages extend
holomorphically to $\mathcal W$. Thus
$\mathcal D=\Phi_{V_0}^{-1}(\mathcal W)$ is open and
$M_{\mathbf B}\circ\Phi_{V_0}$ is holomorphic there.

For real $B\in E$, the continued representative belongs to
$C(\mathbf B,e^{iB}V_0)$, so its real block averages equal
$e^{iB}V_0$. By the assumed agreement of branches, its complex
averages have the same value. The two holomorphic maps therefore
agree on the nonempty relatively open real slice $E$ of the
connected domain $\mathcal D'$. The preceding principle proves
their equality on $\mathcal D'$; inversion gives the values on
reversed bonds.
\end{proof}

\subsection{Tower configurations}
Let $\mathbf B'$ be a bond set of a decreasing sequence
with first region $X$, depth $n$, and the separation
condition at scale $\mathfrak r_1$.
Let $U$ be a real configuration and $A'$ a complex
one-form on the inner bare bonds of $X$, and put
$\mathbf U=e^{i\xi A'}U$. Thus
$\mathbf U(b)=e^{i\xi A'(b)}U(b)$ on positive bonds,
with the inverse value on reversed bonds.

We distinguish two averaged data. The datum
$M'(\mathbf U)=M_{\mathbf B'}^{\mathrm C}(\mathbf U)$ uses
the complex branch at every step, including when
$\mathbf U$ is real. The real datum
$V_0=M'(U)=M_{\mathbf B'}^{\mathrm R}(U)$ uses the global
real averages. Assume that the following conditions hold:
\begin{enumerate}[label=\textup{(\roman*)}]
\item The continuation of \cref{I-3-115} exists
at the real datum $V_0$ with radius $r_{\mathrm B}>0$.
\item Every average defining $M'(\mathbf U)$ exists, and
\[
 |M'(\mathbf U)_j(b)V_{0,j}(b)^{-1}-\mathbf1|<1
\]
for every bond $b$ of every member of $\mathbf B'$,
$0\le j\le n$. At positive bonds define
\[
 B_j(b)=\frac1i\log
       [M'(\mathbf U)_j(b)V_{0,j}(b)^{-1}].
\]
\item The resulting family satisfies $B\in\mathcal B_{r_{\mathrm B}}$.
\end{enumerate}
The tower configuration and its current are
\begin{equation}\label{I-3-119}
 U_n(M'(\mathbf U))
   :=U_{\mathbf B'}(V_0;B),
 \qquad
 J_n(M'(\mathbf U)):=J(U_n(M'(\mathbf U))).
\end{equation}
The current is taken on its domain from \eqref{I-2-7},
using the bare spacing $\eta_n=L^{-n}$ in level-$n$ units.
By \cref{I-3-117}, uniqueness on the smaller of two
admissible balls makes \eqref{I-3-119} independent of
the radius used in these conditions.

These quantities are functions of the pair $(A',U)$,
through $M'(\mathbf U)$ and the real base datum $M'(U)$.
The complex configuration alone does not determine that
pair: for a real one-form $A''$ and a real configuration
$W$,
\[
 e^{i\xi A''}(e^{-i\xi A''}W)=W=e^{i\xi0}W.
\]

\begin{proposition}\label{prop:tower-real-base}
Suppose that the conditions defining \eqref{I-3-119}
hold with $A'=0$, and write
$V_{\mathrm C}=M_{\mathbf B'}^{\mathrm C}(U)$ and
$V_0=M_{\mathbf B'}^{\mathrm R}(U)$.
Then $V_{\mathrm C}$ is real, the family $B$ is real,
and $e^{iB}V_0=V_{\mathrm C}$. The regular critical set
$\mathcal C(\mathbf B',V_{\mathrm C})$ is nonempty and
consists of one residual orbit. Moreover,
\[
 U_n(M'(\mathbf U))=U^\circ_{\mathbf B'}(V_{\mathrm C}),
\]
so the tower configuration is real and its current is
the current of this representative. Finally,
\[
 U_n(M'(\mathbf U))=U^\circ_{\mathbf B'}(V_0)
 \quad\Longleftrightarrow\quad
 V_{\mathrm C}=V_0
 \quad\Longleftrightarrow\quad B=0.
\]
These equalities hold if every configuration averaged
in computing $M_{\mathbf B'}^{\mathrm R}(U)$ belongs to
the set $\mathcal S_c$ for the bond $c$ at which it is
averaged.
\end{proposition}

\begin{proof}
For a real input $W$ in the domain $\mathcal O'_c$, the contour
means lie in $G$ by \cref{I-3-28}. The relative matrices
$W_x(W)$ therefore lie in $G$ and satisfy
$|W_x(W)-\mathbf1|<1$; their logarithms lie in $i\mathfrak g$
by \cref{I-3-15}. Formula \eqref{I-3-53}, and inversion
at reversed bonds, preserve $G$. Iteration proves that
$V_{\mathrm C}$ is real.

At each positive member bond,
$g=V_{\mathrm C,j}(b)V_{0,j}(b)^{-1}$ belongs to $G$
and satisfies $|g-\mathbf1|<1$. The logarithm therefore
gives $B_j(b)\in\mathfrak g$ and $e^{iB_j(b)}=g$.
The inverse convention at reversed bonds yields
$e^{iB}V_0=V_{\mathrm C}$. Since $B$ is a real point
of the continuation ball, its defining hypotheses
give the asserted critical orbit and
\[
 U_{\mathbf B'}(V_0;B)
   =U^\circ_{\mathbf B'}(e^{iB}V_0)
   =U^\circ_{\mathbf B'}(V_{\mathrm C}).
\]

The logarithmic definition gives $B=0$ exactly when
$V_{\mathrm C}=V_0$, in which case the representatives
coincide. Conversely, the real/global averages of
$U^\circ_{\mathbf B'}(V)$ equal $V$ for every real
datum for which this representative is defined.
Equal representatives for $V_{\mathrm C}$ and $V_0$
therefore imply equal data. If every real averaging
step lies in its set $\mathcal S_c$, \eqref{I-3-49}
and \eqref{I-3-53} agree at that step.
Induction through the levels gives
$V_{\mathrm C}=V_0$.
\end{proof}

\section{Summability in the multiscale distance}
\label{sec:multiscale-sums}

We next develop quantitative bounds for the action and averaging
map, followed by real gauge-fixed inverse estimates. We first show
that exponential decay in the multiscale distance can be summed
over all determining bonds with a bound independent of the volume
and the number of layers. Separation between successive regions
compensates for the increasing density of sites at finer scales.

Assume $L\ge13$ throughout this section. Fix $1\le n\le K$ and a decreasing sequence
$\Omega_0\supset\cdots\supset\Omega_n$ satisfying the separation
condition of \eqref{eq:tower-separation} with initial level zero.
We measure lengths in level-$n$ units and set
\[
 \xi_q=L^{q-n}\quad(0\le q\le n),\qquad
 \Omega_{n+1}=\emptyset.
\]
Thus $\Omega_q$ is a union of $\mathfrak M$-cubes of side
$\mathfrak M\xi_q$, and
\[
 \operatorname{dist}_\infty(\Omega_q^c,\Omega_{q+1})
       >R_1\mathfrak M\xi_q\quad(0\le q<n).
\]
Here $\mathfrak M,R_1$ are powers of $L$, with
$\mathfrak M\ge1$ and $R_1\ge L$, as in the separation condition.
Empty members are allowed. Denote the full determining set by
$\mathbf B$ and its reduced set by $\mathbf B^\flat$, with members
$\Gamma_q^\flat$.

Let
\[
 \mathcal L_q=\Omega_q\setminus\Omega_{q+1},\qquad
 \Lambda_q=\Lambda_q(\mathbf B)
           =\mathbb T^{(q)}\cap\mathcal L_q,
 \qquad \Lambda_*=\bigcup_{q=0}^n\Lambda_q.
\]
The region difference is cube-wise. Its point set is the closure
of the point-set difference, and the sets $\Lambda_q$ are disjoint;
these facts are proved in \cref{I-9-20} below. For each positively
oriented $c\in\Gamma_q^\flat$, choose
\[
 y_c=\begin{cases}
 c_-,&c_-\in\Lambda_q,\\
 c_+,&c_-\notin\Lambda_q,
 \end{cases}
 \qquad y_{\bar c}=y_c.
\]
By \cref{I-3-81}, one endpoint of $c$ belongs to $\Lambda_q$;
hence $y_c\in\Lambda_q$.

For torus points, $|y-y'|_1$ and $|y-y'|_\infty$ denote
the minima of the corresponding vector norms over lifts of
$y-y'$, measured in these units. Both minima are attained by
choosing each coordinate of the lift with least absolute value.
The sum of least lifts is a lift of the sum, so the vector
triangle inequalities give the two metric triangle inequalities.
The same least lift also gives $|y-y'|_\infty\le|y-y'|_1$.

For $x\in T_\ell$ and $v\in\mathbb R^3$, write
$\sigma(x,v)=\{x+sv:0\le s\le1\}$. The length of a set
$S\subset T_\ell$ is
\begin{equation}\label{I-1-10}
 \operatorname{length}(S)
 =\inf\left\{\sum_{i=1}^r|v_i|_1:
  r\ge0,\ S\subset\bigcup_{i=1}^r\sigma(x_i,v_i)\right\},
 \qquad \inf\emptyset=+\infty,
\end{equation}
where $x_i\in T_\ell$, $v_i\in\mathbb R^3$, and the vectors
are expressed in level-$n$ units. In particular, the covering is
finite and $|v_i|_1=\sum_\mu|v_{i,\mu}|$.
For a finite union $\Gamma$ of closed segments, set
\begin{equation}\label{I-9-17}
 |\Gamma|_{\mathbf B^\flat}
 =\sum_{q=1}^n\xi_q^{-1}
       \operatorname{length}(\Gamma\cap\mathcal L_q)
   +\xi_0^{-1}\operatorname{length}(\Gamma\setminus\Omega_1),
\end{equation}
with point-set intersections and difference. Define
\[
 d_{\mathbf B^\flat}(y,y')
 =\inf\bigl\{|\Gamma|_{\mathbf B^\flat}:
       \Gamma\text{ is a connected finite union of closed segments}
       \text{ containing }y,y'\bigr\}.
\]
The competitors may lie anywhere in $T_\ell$. The closed layers
may meet along faces, and \eqref{I-9-17} retains each of their
contributions there.

\begin{proposition}[Sums over determining bonds]\label{I-10-125}
For $a>0$, the series
\[
 \mathsf S(a)=54L^3\sum_{r\ge1}(2r+1)^3e^{-a(r-1)/2}
\]
converges. For the separated sequence above and $y\in\Lambda_*$,
the following estimates hold.
\begin{enumerate}[label=(\alph*)]
\item For every $y'\in\Lambda_*$,
\[
 d_{\mathbf B^\flat}(y,y')\ge|y-y'|_\infty.
\]
\item If $aR_1\mathfrak M\ge4(3\log L+\log2)$, then
\[
 \sum_{c\in(\mathbf B^\flat)^+}
      e^{-a d_{\mathbf B^\flat}(y,y_c)}\le\mathsf S(a).
\]
\item If $aR_1\mathfrak M\ge8(3\log L+\log2)$ and $D\ge0$, then
\[
 \sum_{\substack{c\in(\mathbf B^\flat)^+\\
                  d_{\mathbf B^\flat}(y,y_c)\ge D}}
      e^{-a d_{\mathbf B^\flat}(y,y_c)}
       \le e^{-aD/2}\mathsf S(a/2).
\]
\end{enumerate}
The sums count positively oriented bonds, with no scale weight.
The bound $\mathsf S(a)$ depends only on $L$ and $a$; in particular,
it is independent of the torus size, the number of layers, and
the choice of separated sequence.
\end{proposition}

The proof in \cref{subsec:multiscale-sum-proof} combines two
estimates for a connecting set: its displacement costs length
at the coarsest scale it reaches, and each intermediate layer
costs at least $R_1\mathfrak M$ in weighted length. The resulting
exponential factors control both the spatial shell count and
the sum over scales.

\subsection{Layers and weighted length}

\begin{lemma}\label{I-9-20}
For the sequence above, the point set of $\mathcal L_q$ is
\[
 \overline{\{x\in T_\ell:x\in\Omega_q,
                           \ x\notin\Omega_{q+1}\}}.
\]
It is independent of the standard partition used to take the
cube-wise difference. Moreover,
\[
 \Omega_0=\bigcup_{q=0}^n\mathcal L_q,\qquad
 \mathcal L_q\subset\Omega_q,\qquad
 \mathcal L_q\subset\Omega_{q+1}^c.
\]
Every bare site of $\Omega_0$ belongs to exactly one layer, and
the sets $\Lambda_q$ are pairwise disjoint. Every bond of
$\Gamma_q^\flat$ has an endpoint in $\Lambda_q$.
\end{lemma}

\begin{proof}
Let $X_1,X_2$ be regions for one standard partition, and let
$D$ be their cube-wise difference. If $x\in X_1$ and
$x\notin X_2$, a cube of $X_1$ containing $x$ cannot be a cube
of $X_2$, so $x\in D$. Since $D$ is a finite union of closed
cubes, it contains the closure of this point-set difference.
Conversely, let $\Delta$ be a cube of $D$. Its interior points
belong to no other cube of the partition and therefore do not
belong to $X_2$. Every point of $\Delta$ is a limit of these
interior points: writing $\Delta=\{y+w:|w|_\infty\le s/2\}$,
use $y+(1-j^{-1})w$ as $j\to\infty$. Thus $D$ is contained
in the same closure. This description does not refer to the
partition. Applying it to $X_1=\Omega_q$, $X_2=\Omega_{q+1}$
proves the first assertion.

The cubes of $\mathcal L_q$ are cubes of $\Omega_q$ and of
$\Omega_{q+1}^c$, giving the two inclusions. If $x\in\Omega_0$,
choose the largest $q$ such that $x\in\Omega_q$. Then
$x\notin\Omega_{q+1}$, so $x\in\mathcal L_q$. This proves
the covering assertion.

For $0\le i\le q$, a site of $\mathbb T^{(i)}$ belongs to
$\mathcal L_q$ exactly when it belongs to $\Omega_q$ and
not to $\Omega_{q+1}$, by the site convention in
\cref{I-3-77}. Taking $i=0$ gives a unique layer for each
bare site. If $q<q'$ and $y\in\Lambda_q\cap\Lambda_{q'}$,
then $y\notin\Omega_{q+1}$ but
$y\in\Omega_{q'}\subset\Omega_{q+1}$, a contradiction.
Finally, \cref{I-3-81} characterizes a reduced member bond
as having an endpoint in $\Omega_q$ and neither endpoint
in $\Omega_{q+1}$. That endpoint is a site of
$\mathbb T^{(q)}$ in $\mathcal L_q$, hence in $\Lambda_q$.
\end{proof}

\begin{lemma}\label{I-9-22}
For finite unions $\Gamma,\Gamma'$ of closed segments,
\[
 |\Gamma\cup\Gamma'|_{\mathbf B^\flat}
 \le|\Gamma|_{\mathbf B^\flat}+|\Gamma'|_{\mathbf B^\flat}.
\]
The distance is symmetric and satisfies
\[
 d_{\mathbf B^\flat}(y,y'')
 \le d_{\mathbf B^\flat}(y,y')+d_{\mathbf B^\flat}(y',y'')
 \qquad(y,y',y''\in T_\ell).
\]
\end{lemma}

\begin{proof}
The length in \eqref{I-1-10} is monotone under inclusion,
since a cover of a larger set covers every subset.
Combining covers of $S$ and $S'$ and taking infima gives
\[
 \operatorname{length}(S\cup S')
 \le\operatorname{length}(S)+\operatorname{length}(S').
\]
Apply this inequality to the intersections with each layer
and to the point-set differences from $\Omega_1$, multiply
by the positive weights in \eqref{I-9-17}, and add. This
proves weighted subadditivity.

The competitors defining $d_{\mathbf B^\flat}(y,y')$ are
unchanged when $y,y'$ are interchanged, proving symmetry.
If $\Gamma$ connects $y$ to $y'$ and $\Gamma'$ connects
$y'$ to $y''$, their union is connected and contains $y,y''$.
Indeed, two connected sets with a common point cannot be
separated into two nonempty relatively open subsets of their
union. Weighted subadditivity therefore gives
\[
 d_{\mathbf B^\flat}(y,y'')
 \le|\Gamma\cup\Gamma'|_{\mathbf B^\flat}
 \le|\Gamma|_{\mathbf B^\flat}+|\Gamma'|_{\mathbf B^\flat}.
\]
Taking the two infima separately proves the triangle inequality.
\end{proof}

\subsection{Proof of the summation estimate}
\label{subsec:multiscale-sum-proof}

\begin{proof}[Proof of \cref{I-10-125}]
For $a>0$, the ratio of consecutive terms of the series
defining $\mathsf S(a)$ tends to $e^{-a/2}<1$, so the series
converges. Set $w=R_1\mathfrak M$.

\emph{Lower bounds on connecting sets.}
We first record a consequence of the definition of length.
Suppose $g:T_\ell\to\mathbb R$ satisfies
$|g(z)-g(z')|\le |z-z'|_1$, and a bounded interval $I$ is
contained in $g(S)$ for some $S\subset T_\ell$. Then
\[
 \operatorname{length}(S)\ge |I|.
\]
Indeed, let $S\subset\bigcup_{j=1}^r\sigma(x_j,v_j)$ be a
finite segment cover. The image of each segment under $g$
is an interval $I_j$ of length at most $|v_j|_1$, since
\[
 |g(x_j+sv_j)-g(x_j+s'v_j)|
                  \le |s-s'|\,|v_j|_1.
\]
The intervals $I_j$ cover $I$. Integrating
$\mathbf1_I\le\sum_j\mathbf1_{I_j}$ gives
$|I|\le\sum_j|I_j|\le\sum_j|v_j|_1$. Taking the infimum
over covers proves the assertion; if there is no finite cover,
the asserted lower bound holds because the length is infinite.

Let $y\in\Lambda_q$ and $y'\in\Lambda_{q'}$, and let
$\Gamma$ be a connected finite union of segments containing
$y$ and $y'$. Let $t$ be the largest index for which
$\Gamma\cap\Omega_t\ne\emptyset$. This index exists since
$y\in\Omega_q$, and $t\ge\max\{q,q'\}$.
Every point of $\Gamma$ either lies outside $\Omega_1$ or
belongs to one of the layers $\mathcal L_i$, $1\le i\le t$:
for a point in $\Omega_1$, take the largest index of a region
containing it. Thus
\[
 \Gamma=\bigcup_{i=1}^t(\Gamma\cap\mathcal L_i)
                         \cup(\Gamma\setminus\Omega_1).
\]
The sets on the right can share layer faces. Since
$\xi_i^{-1}\ge\xi_t^{-1}$ for $i\le t$, finite
subadditivity of length gives
\[
 |\Gamma|_{\mathbf B^\flat}
                     \ge\xi_t^{-1}\operatorname{length}(\Gamma).
\]
Apply the preceding observation to $g(z)=|z-y|_1$.
Its image on the connected set $\Gamma$ contains the interval
from $0$ to $|y-y'|_1$, and hence
\begin{equation}\label{eq:multiscale-spatial-length}
 |\Gamma|_{\mathbf B^\flat}
             \ge\xi_t^{-1}|y-y'|_\infty.
\end{equation}
In particular, $\xi_t\le1$ implies the first assertion of
the proposition after taking the infimum over $\Gamma$.

For an index $q<i<t$, consider
$f(z)=\operatorname{dist}_\infty(z,\Omega_{i+1})$.
The region $\Omega_{i+1}$ is nonempty because $\Gamma$
meets $\Omega_t\subset\Omega_{i+1}$, and $f$ is
$1$-Lipschitz for the sup metric, hence also for the
$\ell^1$ metric. The site $y$ lies outside
$\Omega_{q+1}\supset\Omega_i$, so $y\in\Omega_i^c$.
Separation gives
\[
 f(y)\ge\operatorname{dist}_\infty
                  (\Omega_i^c,\Omega_{i+1})>w\xi_i.
\]
The function $f$ vanishes at a point of
$\Gamma\cap\Omega_t$. By connectedness, it takes every
value in $(0,w\xi_i)$ on $\Gamma$. Consequently the set
\[
 S_i=\{z\in\Gamma:0<f(z)<w\xi_i\}
\]
has length at least $w\xi_i$. A point of $S_i$ lies outside
$\Omega_{i+1}$ and outside $\Omega_i^c$: the latter follows
from the strict separation inequality. It therefore lies in
$\Omega_i\setminus\Omega_{i+1}$ as a point-set difference,
and in the layer $\mathcal L_i$. Monotonicity of length gives
\[
 \xi_i^{-1}\operatorname{length}(\Gamma\cap\mathcal L_i)
                                                           \ge w.
\]
There are $(t-q-1)_+$ indices strictly between $q$ and $t$,
where $s_+=\max\{s,0\}$. Applying the same argument from
$y'$ gives the two bounds
\begin{equation}\label{eq:multiscale-layer-length}
 |\Gamma|_{\mathbf B^\flat}\ge w(t-q-1)_+,
 \qquad
 |\Gamma|_{\mathbf B^\flat}\ge w(t-q'-1)_+.
\end{equation}
Take one half of the lower bound
\eqref{eq:multiscale-spatial-length} and one quarter of each
lower bound in \eqref{eq:multiscale-layer-length}. We obtain
\[
 |\Gamma|_{\mathbf B^\flat}
 \ge\frac12\xi_t^{-1}|y-y'|_\infty
       +\frac w4\bigl((t-q-1)_++(t-q'-1)_+\bigr).
\]
Exponentiating, and bounding the term for this $t$ by the sum
over all possible $t$, gives a bound independent of $\Gamma$.
Taking the supremum of $e^{-a|\Gamma|_{\mathbf B^\flat}}$
over the connecting sets yields
\begin{equation}\label{eq:multiscale-depth-bound}
 e^{-a d_{\mathbf B^\flat}(y,y')}
 \le\sum_{t=\max\{q,q'\}}^n
 e^{-\frac a2\xi_t^{-1}|y-y'|_\infty}
 e^{-\frac{aw}4((t-q-1)_++(t-q'-1)_+)}.
\end{equation}

\emph{Summing sites and levels.}
Fix $t\ge q'$ and a lift $\widetilde y$ of $y$.
Every site $y'\in\mathbb T^{(q')}$ satisfying
$|y-y'|_\infty<r\xi_t$, $r\ge1$, has a lift in the open
cube of radius $r\xi_t$ centered at $\widetilde y$.
The lifted lattice has spacing
$\xi_{q'}=L^{q'-t}\xi_t$. In each coordinate the cube
contains at most $2rL^{t-q'}+1$ lattice values. Therefore
\[
 \#\{y'\in\mathbb T^{(q')}:|y-y'|_\infty<r\xi_t\}
       \le (2rL^{t-q'}+1)^3
       \le (2r+1)^3L^{3(t-q')}.
\]
Partitioning $\Lambda_{q'}$ into the shells
$r-1\le\xi_t^{-1}|y-y'|_\infty<r$ gives
\begin{equation}\label{eq:multiscale-shell-sum}
 \sum_{y'\in\Lambda_{q'}}
       e^{-\frac a2\xi_t^{-1}|y-y'|_\infty}
 \le L^{3(t-q')}
              \sum_{r\ge1}(2r+1)^3e^{-a(r-1)/2}.
\end{equation}
Under the hypothesis of the second assertion,
$L^3e^{-aw/4}\le\frac12$. Hence
\[
 \begin{aligned}
 L^{3(t-q')}e^{-\frac{aw}4(t-q'-1)_+}
                       &\le L^3\,2^{-(t-q'-1)_+},\\
 e^{-\frac{aw}4(t-q-1)_+}&\le2^{-(t-q-1)_+}.
 \end{aligned}
\]
The layer site sets are disjoint. Sum
\eqref{eq:multiscale-depth-bound} over $y'\in\Lambda_*$,
use \eqref{eq:multiscale-shell-sum}, and interchange the
finite sums over $q'$ and $t$. The result is
\[
 \begin{aligned}
 \sum_{y'\in\Lambda_*}e^{-a d_{\mathbf B^\flat}(y,y')}
 &\le L^3\sum_{t=q}^n2^{-(t-q-1)_+}
       \sum_{q'=0}^t2^{-(t-q'-1)_+}
       \sum_{r\ge1}(2r+1)^3e^{-a(r-1)/2}\\
 &\le9L^3\sum_{r\ge1}(2r+1)^3e^{-a(r-1)/2},
 \end{aligned}
\]
since each of the two sums of powers of $2$ is at most
$\sum_{j\ge0}2^{-(j-1)_+}=3$.

If $y'=y_c$ for a positive bond $c\in\Gamma_{q'}^\flat$,
then $y'\in\Lambda_{q'}$. Disjointness of the layer site
sets implies that all bonds anchored at a given site belong
to this one level. A site of its lattice is incident to at most
six positive bonds. It follows that
\[
 \sum_{c\in(\mathbf B^\flat)^+}
                   e^{-a d_{\mathbf B^\flat}(y,y_c)}
 \le6\sum_{y'\in\Lambda_*}e^{-a d_{\mathbf B^\flat}(y,y')}
 \le\mathsf S(a).
\]

\emph{The tail estimate.}
For $d\ge D$, one has
$e^{-ad}\le e^{-aD/2}e^{-ad/2}$. The stronger lower bound
on $aw$ is precisely the hypothesis of the sum estimate
with $a/2$ in place of $a$. Therefore
\[
 \begin{aligned}
 \sum_{\substack{c\in(\mathbf B^\flat)^+\\
                   d_{\mathbf B^\flat}(y,y_c)\ge D}}
                e^{-a d_{\mathbf B^\flat}(y,y_c)}
 &\le e^{-aD/2}
       \sum_{c\in(\mathbf B^\flat)^+}
                e^{-\frac a2 d_{\mathbf B^\flat}(y,y_c)}\\
 &\le e^{-aD/2}\mathsf S(a/2).
 \end{aligned}
\]
\end{proof}

\section{Local estimates for the action}\label{sec:action-estimates}

The current formula gives the derivative of the action along
multiplicative variations. We now estimate this derivative in
exponential coordinates, separating its covariant second-order
part, the correction due to the reference curvature, and a
quadratic remainder. The estimates are local and uniform in
the lattice spacing.

Throughout this section assume $L\ge13$, fix $0\le k\le K$,
and use level-$k$ scaling, with $\eta=L^{-k}$. Let $U_0$ be
a real configuration on the bare lattice. Write
$D=D_{U_0}$ and $\nabla=\nabla^{U_0}$; the transpose
$D^{\mathsf T}$ below acts on two-forms. For a complex bare
one-form $A$, set
\[
 F(A)=A^\eta(e^{i\eta A}U_0)
     =\eta^{-1}\sum_p a_p(e^{i\eta A}U_0).
\]
At positive bonds the perturbed configuration has value
$e^{i\eta A(b)}U_0(b)$; reverse bonds are defined by inversion.
The gradient $F'(A)$ is characterized by
\[
 \left.\frac{d}{dt}\right|_{t=0}F(A+t\delta A)
       =(F'(A),\delta A)_\eta
\]
for all complex one-forms $\delta A$. The pairing is the
complex bilinear pairing of \cref{I-2-2}. Its nondegeneracy
gives a unique gradient, as verified below.

For a plaquette $p$, let $\mathsf b(p)$ be its four positive
bonds. For a positive bond $b$, let $\mathsf{pl}(b)$ be the
four plaquettes containing it and put
$\mathsf{st}(b)=\bigcup_{p\in\mathsf{pl}(b)}\mathsf b(p)$.
For a set $S$ of positive bonds, write
\[
 s_S(A)=\max_{b\in S}|A(b)|,
 \qquad
 g_S(A)=\max|\nabla_\lambda A_\kappa(x)|,
 \qquad
 \mathsf n_S(A)=\max\{s_S(A),g_S(A)\},
\]
where the second maximum is over covariant differences
between two bonds of $S$, and is zero if there are none.
Thus the two bonds in such a difference are
$\langle x,x+\eta e_\kappa\rangle$ and
$\langle x+\eta e_\lambda,x+\eta e_\lambda+\eta e_\kappa\rangle$.
Use subscripts $p$ and $b$ for $S=\mathsf b(p)$ and
$S=\mathsf{st}(b)$, respectively, and set
\begin{equation}\label{I-10-12}
 \|A\|_X=\max\{\sup|A|,\sup|\nabla A|\}.
\end{equation}

\begin{proposition}[Local first-variation estimates]\label{I-10-42}
For every real reference configuration $U_0$, the maps
$F$ and $F'$ are entire in the complex form $A$, and
$F'(0)=J(U_0)$. Put
\[
 \mathcal W(A)=F'(A)-J(U_0)-D^{\mathsf T}DA,
 \qquad
 \Delta'=\mathcal W'(0),
 \qquad
 \mathcal V(A)=\mathcal W(A)-\Delta'A.
\]
The values $F'(A)(b)$, $(\Delta'A)(b)$ and
$\mathcal V(A)(b)$ depend on $A$ only through its restriction
to $\mathsf{st}(b)$.

There is an absolute constant $C_4\ge64$ such that, with
$\rho_1=1/128$, the following hold whenever
\[
 0\le\alpha_0\le1,
 \qquad
 |\partial U_0(p)-\mathbf1|\le\alpha_0\eta^2
 \quad\text{at every bare plaquette }p.
\]
The constant is independent of the volume, the cutoffs,
the lattice spacing and $U_0$.
\begin{enumerate}
\item[(i)] The first variation decomposes as
\[
 F'(A)=J(U_0)+D^{\mathsf T}DA+\Delta'A+\mathcal V(A).
\]
The operator $\Delta'$ is linear and symmetric for
$(\cdot,\cdot)_\eta$, and
\[
 |(\Delta'A)(b)|\le C_4\alpha_0s_b(A)
\]
for every complex $A$. Its value at $b$ vanishes if
$\partial U_0(p)=\mathbf1$ for all
$p\in\mathsf{pl}(b)$. The operator
$\Delta^{\mathrm W}_{U_0}=D^{\mathsf T}D+\Delta'$
represents the Wilson Hessian:
\[
 (A,\Delta^{\mathrm W}_{U_0}A)_\eta
   =\left.\frac{d^2}{dt^2}\right|_{t=0}F(tA).
\]
The map $\mathcal V$ is entire,
$\mathcal V(0)=0$ and $\mathcal V'(0)=0$. If $A$ is complex
and $s_b(A)\le\rho_1$, then
\[
 |\mathcal V(A)(b)|\le C_4s_b(A)\mathsf n_b(A).
\]
In particular, for every complex $A$,
\[
 \sup|\Delta'A|\le C_4\alpha_0\sup|A|.
\]
For complex $A$ with $\sup|A|\le\rho_1$,
\[
 \sup|\mathcal V(A)|\le C_4\sup|A|\,\|A\|_X.
\]

\item[(ii)] Let $A$ be real and put $E=e^{i\eta A}U_0$.
With $\operatorname{ad}_XY=XY-YX$ and
\[
 \mathsf f(z)=\frac{e^z-1}{z}
             =\sum_{q\ge0}\frac{z^q}{(q+1)!},
\]
where the power series defines the value at zero and on
linear maps, one has
\[
 F'(A)(b)=\mathsf f(-i\eta\operatorname{ad}_{A(b)})J(E)(b).
\]
If $\eta\sup|A|\le1/8$, then
\begin{align*}
 |J(E)(b)|&\le(1+4\eta|A(b)|)\,|F'(A)(b)|,\\
 |F'(A)(b)|&\le(1+2\eta|A(b)|)\,|J(E)(b)|.
\end{align*}

\item[(iii)] Let $A$ be real and put $E=e^{i\eta A}U_0$.
At any plaquette $p$ with $s_p(A)\le\rho_1$,
\[
 |\partial E(p)-\mathbf1|
 \le |\partial U_0(p)-\mathbf1|
       +\bigl(2g_p(A)+C_4s_p(A)^2\bigr)\eta^2.
\]
Consequently, if $\sup|A|\le\rho_1$, the same bound holds
with $g_p(A)$ and $s_p(A)$ replaced by
$\sup|\nabla A|$ and $\sup|A|$.

\item[(iv)] If $A_1,A_2$ are complex and
$\|A_1\|_X,\|A_2\|_X\le r\le\rho_1/4$, then
\[
 \sup|\mathcal V(A_1)-\mathcal V(A_2)|
       \le4C_4r\,\|A_1-A_2\|_X.\qedhere
\]
\end{enumerate}
\end{proposition}

The nonlinear error of a single plaquette is quadratic in
$A$. To estimate its exterior transpose uniformly in $\eta$,
we compare adjacent plaquettes after transport to the same
base point. Their difference contains a factor $\eta$,
which cancels the inverse spacing in $D^{\mathsf T}$.
The calculation is given in
\cref{subsec:action-plaquette-calculation};
\cref{subsec:action-remainder-proof} then extracts the
linear curvature correction and quadratic remainder.

\subsection{The coordinate gradient}

\begin{proof}[Entire dependence and locality in \cref{I-10-42}]
The action is a finite polynomial in the entries of
$e^{\pm i\eta A(b)}$, with coefficients determined by $U_0$.
It is therefore entire. Choose a real basis
$t_1,t_2,t_3$ of $\mathfrak g$ with
$\langle t_a,t_{a'}\rangle=\delta_{aa'}$, and write
$A(b)=\sum_a A^a(b)t_a$. This is also a complex basis of
$\mathfrak g^c$. Since the pairing on one-forms has the
factor $\eta^3$, the form
\[
 F'(A)(b)=\eta^{-3}\sum_a
                \frac{\partial F}{\partial A^a(b)}(A)t_a
\]
satisfies the defining derivative identity and is entire.
If a form pairs to zero with every variation, variations
supported at one bond and equal to one basis vector show
that each of its coordinates is zero. This proves uniqueness.
Equation \eqref{I-2-16} at $U_0$ gives
$F'(0)=J(U_0)$.

Let $\delta_bY$ be supported at the positive bond $b$, with
value $Y\in\mathfrak g^c$. Only the four plaquettes
containing $b$ contribute to its variation, so
\begin{equation}\label{I-10-43}
 \eta^3\langle F'(A)(b),Y\rangle
 =\eta^{-1}\sum_{p\in\mathsf{pl}(b)}
     \left.\frac{d}{dt}\right|_{t=0}
          a_p(e^{i\eta(A+t\delta_bY)}U_0).
\end{equation}
Every summand uses $A$ only on $\mathsf b(p)$. Testing
against the basis vectors proves the asserted locality of
$F'$. The formula for $D^{\mathsf T}$ in \eqref{I-2-5}
uses only those same four plaquettes; each value of $DA$
uses their four bonds. Thus $\mathcal W(A)(b)$ depends only
on $A|_{\mathsf{st}(b)}$. Its derivative at zero and the
remainder have the same dependence. Finally,
$\mathcal W$ is entire and $\mathcal W(0)=0$; hence
$\mathcal V=\mathcal W-\mathcal W'(0)A$ is entire and
$\mathcal V(0)=\mathcal V'(0)=0$.
\end{proof}

\subsection{The plaquette calculation}
\label{subsec:action-plaquette-calculation}

We use the four positive bonds $b_1,\dots,b_4$ of
\eqref{I-2-3}, with $\xi=\eta$. Transport the values of $A$ to
the base point of $p=p_{\mu\nu}(x)$ by setting
\[
 \alpha_1=A(b_1),\qquad
 \alpha_2=\operatorname{Ad}(U_0(b_1))A(b_2),\qquad
 \alpha_3=\operatorname{Ad}(U_0(b_4))A(b_3),\qquad
 \alpha_4=A(b_4).
\]
Since $U_0$ is unitary, $|\alpha_j|=|A(b_j)|$.
Moving the reference transports past the exponentials gives
\begin{equation}\label{I-10-40}
 \partial(e^{i\eta A}U_0)(p)
 =e^{i\eta\alpha_1}e^{i\eta\alpha_2}\,
       \partial U_0(p)\,
       e^{-i\eta\alpha_3}e^{-i\eta\alpha_4}.
\end{equation}
The covariant differences satisfy
\begin{equation}\label{eq:action-alpha-differences}
\begin{aligned}
 \alpha_2-\alpha_4&=\eta(\nabla_\mu A_\nu)(x),&
 \alpha_3-\alpha_1&=\eta(\nabla_\nu A_\mu)(x),\\
 \eta(DA)(p)&=\alpha_1+\alpha_2-\alpha_3-\alpha_4.
\end{aligned}
\end{equation}
Consequently
\begin{equation}\label{eq:action-exterior-local}
 |(DA)(p)|\le2g_p(A),\qquad
 \eta|(DA)(p)|\le4s_p(A).
\end{equation}
These identities follow directly from the definition of
$\nabla$ and \eqref{I-2-4}. They also show that
$(D\,\delta_bY)(p)=0$ when $p\notin\mathsf{pl}(b)$, where
$\delta_bY$ has value $Y$ at $b$ and zero at the other positive
bonds.

For complex matrices $X_1,\dots,X_4$ sufficiently close to zero,
put
\begin{align*}
 \Lambda(X_1,\dots,X_4)
 &=\log(e^{-X_3}e^{-X_4}e^{X_1}e^{X_2})
       -(X_1+X_2-X_3-X_4),\\
 \mathsf N(\beta)
 &=\frac{\Lambda(i\eta\beta_1,\dots,i\eta\beta_4)}
                       {i\eta^2}.
\end{align*}
The order of the four factors comes from cyclically moving the
last two factors of \eqref{I-10-40} to the front inside the
trace. We write
\[
 N_p=\mathsf N(\alpha_1,\dots,\alpha_4),\qquad
 \Theta_p=(DA)(p)+N_p.
\]

\begin{lemma}\label{I-10-38}
If $s_p(A)\le s\le1/64$, then $N_p,\Theta_p\in\mathfrak g^c$,
and $i\eta^2\Theta_p$ is the series logarithm of
$e^{-i\eta\alpha_3}e^{-i\eta\alpha_4}
 e^{i\eta\alpha_1}e^{i\eta\alpha_2}$.
Moreover,
\begin{equation}\label{eq:action-plaquette-remainder}
 |N_p|\le64s^2,\qquad
 \eta|\Theta_p|\le5s,\qquad
 |\Theta_p|\le2g_p(A)+64s^2.
\end{equation}
For $|\beta_j|\le s$, the function $\mathsf N$ satisfies
\begin{equation}\label{eq:action-plaquette-derivative}
 \left|
  \frac{\partial\mathsf N}{\partial\beta_j}(\beta)Z
 \right|\le64s|Z|.
\end{equation}

Let $\lambda\in\{\mu,\nu\}$,
$p'=p_{\mu\nu}(x-\eta e_\lambda)$,
$u=U_0(\langle x-\eta e_\lambda,x\rangle)^{-1}$, and
$S=\mathsf b(p)\cup\mathsf b(p')$. Then
\[
 |\operatorname{Ad}(u)\alpha_j(p')-\alpha_j(p)|
       \le3\eta g_S(A).
\]
If $s_S(A)\le s\le1/64$, then
\begin{equation}\label{eq:action-neighboring-remainders}
 |\operatorname{Ad}(u)N_{p'}-N_p|
       \le768\eta s\,g_S(A).
\end{equation}
\end{lemma}

\begin{proof}
First consider matrices $Z_1,\dots,Z_n$, $n\le4$, and put
$\sigma=\sum_j|Z_j|\le1$. Multiplication of the exponential
series bounds all terms of degree at least two by
$e^\sigma-1-\sigma\le\sigma^2$. Hence
\[
 \left|e^{Z_1}\cdots e^{Z_n}-\mathbf1-\sum_jZ_j\right|
       \le\sigma^2,\qquad
 \left|e^{Z_1}\cdots e^{Z_n}-e^{\sum_jZ_j}\right|
       \le2\sigma^2.
\]
For $\sigma\le1/4$, both products in the second difference
are within $e^{1/4}-1<1/2$ of $\mathbf1$.
The logarithmic estimate in \cref{I-3-15}, and
$\log e^{\sum_jZ_j}=\sum_jZ_j$, therefore give
\[
 \left|\log(e^{Z_1}\cdots e^{Z_n})-\sum_jZ_j\right|
       \le4\sigma^2.
\]
Applied in the order $-X_3,-X_4,X_1,X_2$, this proves that
$\Lambda$ is holomorphic on the ball
$|X|:=\sum_j|X_j|<1/4$, with
$|\Lambda(X)|\le4|X|^2$. Conjugation by a unitary matrix
commutes with this logarithm and preserves its domain.
If the $X_j$ are traceless, the product has determinant one;
its series logarithm is traceless by \cref{I-3-15}.
Thus $\Lambda$ is traceless as well.

We also need a derivative bound that vanishes at the origin.
On $|X|<1/4$, the preceding estimate gives
$|\Lambda(X)|\le1/4$. The one-variable Cauchy estimate along
$X+t\,\delta X$ yields
\[
 |\Lambda'(X)\delta X|
 \le\frac{(1/4)|\delta X|}{1/4-|X|}
 \le2|\delta X|,\qquad |X|<1/8,
\]
where $|\delta X|=\sum_j|\delta X_j|$.
The quadratic estimate gives $\Lambda'(0)=0$.
Apply the Schwarz lemma on the ball of radius $1/8$ to
$X\mapsto\Lambda'(X)\delta X$, by restricting to a complex
line through zero and then applying a linear functional of
norm one. This gives
\begin{equation}\label{eq:action-log-derivative}
 |\Lambda'(X)\delta X|
       \le16|X|\,|\delta X|,\qquad |X|<1/8.
\end{equation}

If $|\beta_j|\le s\le1/64$, then
$\sum_j|i\eta\beta_j|\le4\eta s\le1/16$.
The quadratic bound for $\Lambda$ gives
$|\mathsf N(\beta)|\le64s^2$.
For the derivative in $\beta_j$, the variation of $i\eta\beta$
has the single nonzero entry $i\eta Z$. Therefore
\eqref{eq:action-log-derivative} gives
\[
 \left|
  \frac{\partial\mathsf N}{\partial\beta_j}(\beta)Z
 \right|
 \le\eta^{-2}\,16(4\eta s)\eta|Z|
 =64s|Z|.
\]
At $\beta=\alpha$, tracelessness and
\eqref{eq:action-alpha-differences} imply
\begin{align*}
 i\eta^2\Theta_p
 &=i\eta(\alpha_1+\alpha_2-\alpha_3-\alpha_4)
                  +\Lambda(i\eta\alpha)\\
 &=\log(e^{-i\eta\alpha_3}e^{-i\eta\alpha_4}
             e^{i\eta\alpha_1}e^{i\eta\alpha_2}).
\end{align*}
The inverse identity for the series logarithm proves the
assertion about its exponential.
Together with \eqref{eq:action-exterior-local}, the bound on
$N_p$ gives
$\eta|\Theta_p|\le4s+64\eta s^2\le5s$ and
$|\Theta_p|\le2g_p(A)+64s^2$.

It remains to compare the two plaquettes. For
$\kappa\in\{\mu,\nu\}$, the definition of the covariant
difference gives
\[
 \operatorname{Ad}(u)A_\kappa(x-\eta e_\lambda)-A_\kappa(x)
 =-\eta\operatorname{Ad}(u)
             (\nabla_\lambda A_\kappa)(x-\eta e_\lambda).
\]
The two bonds involved belong to $S$. This gives the bound
$\eta g_S(A)$ for $\alpha_1$ and $\alpha_4$.
For $\alpha_2$, use \eqref{eq:action-alpha-differences} at
both plaquettes:
\begin{align*}
 \operatorname{Ad}(u)\alpha_2(p')-\alpha_2(p)
 &=\operatorname{Ad}(u)\alpha_4(p')-\alpha_4(p)\\
 &\quad+\eta\bigl[
   \operatorname{Ad}(u)(\nabla_\mu A_\nu)(x-\eta e_\lambda)
                   -(\nabla_\mu A_\nu)(x)\bigr].
\end{align*}
Each of the last two covariant differences is between bonds
of one of the two plaquettes. Their norms are at most
$g_S(A)$, so the bound is $3\eta g_S(A)$.
The same computation with $\alpha_3-\alpha_1$ gives the
bound for $\alpha_3$.

Unitary covariance of $\mathsf N$ gives
\[
 \operatorname{Ad}(u)N_{p'}
 =\mathsf N\bigl(\operatorname{Ad}(u)\alpha_1(p'),\dots,
                \operatorname{Ad}(u)\alpha_4(p')\bigr).
\]
Every entry on the line segment joining this quadruple to
$(\alpha_j(p))_j$ has norm at most $s$.
Integrating \eqref{eq:action-plaquette-derivative} along that
segment yields
\[
 |\operatorname{Ad}(u)N_{p'}-N_p|
 \le64s\sum_{j=1}^4
         |\operatorname{Ad}(u)\alpha_j(p')-\alpha_j(p)|
 \le768\eta s\,g_S(A).
\]
\end{proof}

To differentiate the plaquette action, write
\[
 c_p=\tfrac12\tr\partial U_0(p),\qquad F_p=F_{U_0}(p).
\]
The reference configuration is real. The SU(2) decomposition
in \eqref{I-2-15} gives
\begin{equation}\label{eq:action-reference-plaquette}
\begin{gathered}
 \partial U_0(p)=c_p\mathbf1+i\eta^2F_p,\qquad -1\le c_p\le1,\\
 1-c_p=\tfrac12|\partial U_0(p)-\mathbf1|^2,\qquad
 \eta^2|F_p|\le|\partial U_0(p)-\mathbf1|.
\end{gathered}
\end{equation}
Indeed, the eigenvalues of a matrix in SU(2) are
$e^{i\theta},e^{-i\theta}$, so $c_p=\cos\theta$ and
$|\partial U_0(p)-\mathbf1|^2=2(1-c_p)$.
Also $\eta^4|F_p|^2=1-c_p^2\le2(1-c_p)$.

Introduce the entire scalar functions
\begin{gather*}
 \mathrm{co}(\zeta)=\sum_{q\ge0}\frac{(-\zeta)^q}{(2q)!},
 \qquad
 \mathrm{si}(\zeta)=\sum_{q\ge0}\frac{(-\zeta)^q}{(2q+1)!},\\
 \zeta_p=\eta^4\langle\Theta_p,\Theta_p\rangle .
\end{gather*}
For a traceless two-by-two matrix $\Theta$,
Cayley--Hamilton gives
$\Theta^2=\langle\Theta,\Theta\rangle\mathbf1$.
Hence
\[
 e^{i\eta^2\Theta}
 =\mathrm{co}(\eta^4\langle\Theta,\Theta\rangle)\mathbf1
  +i\eta^2\mathrm{si}(\eta^4\langle\Theta,\Theta\rangle)\Theta .
\]
By cyclicity of the trace in \eqref{I-10-40}, and by
\cref{I-10-38},
\[
 \tfrac12\tr\partial(e^{i\eta A}U_0)(p)
 =\tfrac12\tr\bigl[\partial U_0(p)e^{i\eta^2\Theta_p}\bigr].
\]
Using \eqref{eq:action-reference-plaquette}, we obtain
\begin{equation}\label{eq:action-plaquette-trace}
 a_p(e^{i\eta A}U_0)
 =1-c_p\mathrm{co}(\zeta_p)
       +\eta^4\mathrm{si}(\zeta_p)
                      \langle F_p,\Theta_p\rangle .
\end{equation}
The logarithm calculation used here holds on the open set
$4\eta s_p(A)<1/4$, so the identity may be differentiated at each $A$ with
$s_p(A)\le1/64$.

Set
\begin{align*}
 M_p^{\mathrm r}
 &=-c_p\bigl(2\mathrm{co}'(\zeta_p)+1\bigr)\Theta_p
       +\bigl(\mathrm{si}(\zeta_p)-1\bigr)F_p\\
 &\quad+2\eta^4\mathrm{si}'(\zeta_p)
                     \langle F_p,\Theta_p\rangle\Theta_p,
 \qquad
 M_p=F_p+c_p\Theta_p+M_p^{\mathrm r}.
\end{align*}
All these matrices are traceless. Differentiation of
\eqref{eq:action-plaquette-trace}, with
$\delta\zeta_p=2\eta^4\langle\Theta_p,\delta\Theta_p\rangle$,
gives
\begin{align*}
 \delta a_p
 &=\eta^4\left\langle
       -2c_p\mathrm{co}'(\zeta_p)\Theta_p
       +\mathrm{si}(\zeta_p)F_p,\delta\Theta_p\right\rangle\\
 &\quad+2\eta^8\mathrm{si}'(\zeta_p)
        \langle F_p,\Theta_p\rangle
        \langle\Theta_p,\delta\Theta_p\rangle\\
 &=\eta^4\langle M_p,\delta\Theta_p\rangle .
\end{align*}

For $b=b_j\in\mathsf b(p)$ and a variation $\delta_bY$, let
\begin{gather*}
 \mathsf L_{p,b}Y
 =\frac{\partial\mathsf N}{\partial\beta_j}(\alpha)\,Y_j,\\
 Y_1=Y_4=Y,\quad
 Y_2=\operatorname{Ad}(U_0(b_1))Y,\quad
 Y_3=\operatorname{Ad}(U_0(b_4))Y.
\end{gather*}
Then $\delta\Theta_p=(D\,\delta_bY)(p)+\mathsf L_{p,b}Y$.
The map $\mathsf L_{p,b}$ preserves $\mathfrak g^c$ and,
by \eqref{eq:action-plaquette-derivative},
$|\mathsf L_{p,b}Y|\le64s|Y|$ if $s_p(A)\le s\le1/64$.
Its transpose is taken in the complex bilinear form
$\langle\cdot,\cdot\rangle$. To bound its norm, we use
\[
 |X|\le2\sup_{\substack{Y\in\mathfrak g^c\\ |Y|\le1}}
                    |\langle X,Y\rangle|.
\]
For $X\ne0$, choose $Y=X^*/|X|$; then
$\langle X,Y\rangle=\tr(XX^*)/(2|X|)\ge|X|/2$.
Since $|\langle Z,W\rangle|\le|Z|\,|W|$, it follows that
\[
 |\mathsf L_{p,b}^{\mathsf T}Z|
 \le2\sup_{|Y|\le1}|\langle Z,\mathsf L_{p,b}Y\rangle|
 \le128s|Z|.
\]
We have proved the local differential formula
\begin{equation}\label{I-10-41}
\begin{aligned}
 \frac{d}{dt}\,
 a_p(e^{i\eta(A+t\delta_bY)}U_0)\Big|_{t=0}
 &=\eta^4\langle M_p,(D\,\delta_bY)(p)\rangle\\
 &\quad+\eta^4\langle\mathsf L_{p,b}^{\mathsf T}M_p,Y\rangle,\\
 |\mathsf L_{p,b}^{\mathsf T}Z|&\le128s|Z|.
\end{aligned}
\end{equation}

Finally, $|\zeta_p|\le\eta^4|\Theta_p|^2\le(5\eta s)^2\le1$.
The scalar power series show, for $|\zeta|\le1$, that
\[
 |2\mathrm{co}'(\zeta)+1|\le|\zeta|,\qquad
 |\mathrm{si}(\zeta)-1|\le|\zeta|,\qquad
 |\mathrm{si}'(\zeta)|\le1.
\]
For example, the absolute sums bounding the respective
coefficients are
$\sum_{q\ge2}2q/(2q)!$,
$\sum_{q\ge1}1/(2q+1)!$, and
$\sum_{q\ge1}q/(2q+1)!$, each less than one.
As $|c_p|\le1$, these inequalities give
\begin{equation}\label{eq:action-trace-remainder}
 |M_p^{\mathrm r}|
 \le\eta^4\bigl(|\Theta_p|^3+3|F_p|\,|\Theta_p|^2\bigr).
\end{equation}

\subsection{The linear correction and nonlinear remainder}
\label{subsec:action-remainder-proof}

We prove the expansion and local estimates in
\cref{I-10-42}. The entire-gradient and locality argument
above defines
$\mathcal W(A)=F'(A)-J(U_0)-D^{\mathsf T}DA$ and its
linear and higher-order terms $\Delta'A$ and $\mathcal V(A)$.
In the estimates below $C$ is an absolute constant whose
value may change.

Assume the reference curvature bound. By
\eqref{eq:action-reference-plaquette},
\begin{equation}\label{eq:action-reference-smallness}
 0\le1-c_p\le\tfrac12\alpha_0^2\eta^4,\qquad
 |F_p|\le\alpha_0.
\end{equation}
For $s_b(A)\le1/64$, insert \eqref{I-10-41} into the
local gradient identity \eqref{I-10-43} and transpose $D$.
The factors $\eta^{-1}\eta^4=\eta^3$ agree with
the pairing of forms, and nondegeneracy gives
\begin{equation}\label{I-10-44}
\begin{aligned}
 F'(A)(b)&=(D^{\mathsf T}M)(b)+\mathcal T(A)(b),\\
 \mathcal T(A)(b)&=
     \sum_{p\in\mathsf{pl}(b)}\mathsf L_{p,b}^{\mathsf T}M_p .
\end{aligned}
\end{equation}
Here $M$ need only be specified on $\mathsf{pl}(b)$:
extend it by zero before taking the transpose.
The value at $b$ is independent of this extension by
\eqref{I-2-5}. Also,
\[
 |\mathcal T(A)(b)|
       \le4\cdot128s_b(A)\max_{p\in\mathsf{pl}(b)}|M_p|.
\]

Fix $b=\langle x,x+\eta e_\mu\rangle$ and abbreviate
$s=s_b(A)\le1/64$, $g=g_b(A)$, and
$\mathsf n=\mathsf n_b(A)$.
On the incident plaquettes,
\[
 M_p=F_p+(DA)(p)+(c_p-1)(DA)(p)+c_pN_p+M_p^{\mathrm r}.
\]
Since $J(U_0)=D^{\mathsf T}F_{U_0}$, subtracting the
constant and exterior-Laplacian terms in \eqref{I-10-44}
gives
\begin{equation}\label{eq:action-four-remainders}
\begin{aligned}
 \mathcal W(A)(b)
 &=\bigl(D^{\mathsf T}((c-1)DA)\bigr)(b)
  +\bigl(D^{\mathsf T}(cN)\bigr)(b)\\
 &\quad+(D^{\mathsf T}M^{\mathrm r})(b)
  +\mathcal T(A)(b).
\end{aligned}
\end{equation}
Formula \eqref{I-2-5} expresses $(D^{\mathsf T}X)(b)$
as two terms of the form
$\pm\eta^{-1}[\operatorname{Ad}(u_\nu)X(p')-X(p)]$,
where $p'$ is the translate of $p$ by $-\eta e_\nu$ and
$u_\nu=U_0(\langle x-\eta e_\nu,x\rangle)^{-1}$.
Both plaquettes belong to $\mathsf{pl}(b)$.
In particular,
\[
 |(D^{\mathsf T}X)(b)|
       \le4\eta^{-1}\max_{p\in\mathsf{pl}(b)}|X(p)|.
\]

For the first term of \eqref{eq:action-four-remainders},
\eqref{eq:action-exterior-local} and
\eqref{eq:action-reference-smallness} yield
\[
 \bigl|\bigl(D^{\mathsf T}((c-1)DA)\bigr)(b)\bigr|
 \le4\eta^{-1}\cdot\tfrac12\alpha_0^2\eta^4
                     \cdot4\eta^{-1}s
 \le8\alpha_0s.
\]
For the second term, the bound on $N_p$ alone would leave
an inverse power of $\eta$. Compare the two values before
applying the transpose:
\[
 c_{p'}\operatorname{Ad}(u_\nu)N_{p'}-c_pN_p
 =c_{p'}[\operatorname{Ad}(u_\nu)N_{p'}-N_p]
       +(c_{p'}-c_p)N_p.
\]
The first bracket is at most $768\eta sg$ by
\eqref{eq:action-neighboring-remainders}.
Both $1-c_p$ and $1-c_{p'}$ belong to
$[0,\alpha_0^2\eta^4/2]$, so
$|c_{p'}-c_p|\le\alpha_0^2\eta^4/2$.
Since $|c_{p'}|\le1$ and $|N_p|\le64s^2$, the two
differences in the transpose satisfy
\[
 \bigl|\bigl(D^{\mathsf T}(cN)\bigr)(b)\bigr|
 \le2\eta^{-1}(768\eta sg+32\alpha_0^2\eta^4s^2)
 \le C s\mathsf n .
\]

For the third term, use \eqref{eq:action-trace-remainder},
$|F_p|\le\alpha_0$, and $\eta|\Theta_p|\le5s$:
\begin{align*}
 |(D^{\mathsf T}M^{\mathrm r})(b)|
 &\le4\eta^3
    \max_{p\in\mathsf{pl}(b)}
       (|\Theta_p|^3+3\alpha_0|\Theta_p|^2)\\
 &\le4(125s^3+75\alpha_0s^2)
 \le C s(\alpha_0+\mathsf n).
\end{align*}
For the last term,
$|\Theta_p|\le2g+64s^2\le3\mathsf n$.
Furthermore,
\[
 |M_p^{\mathrm r}|
 \le(\eta|\Theta_p|)^2(|\Theta_p|+3\alpha_0)
 \le|\Theta_p|+3\alpha_0s,
\]
because $\eta\le1$ and
$(\eta|\Theta_p|)^2\le25s^2\le s\le1$.
It follows that $|M_p|\le\alpha_0+C\mathsf n$, and the
bound following \eqref{I-10-44} gives
$|\mathcal T(A)(b)|\le Cs(\alpha_0+\mathsf n)$.
Thus the four terms establish
\begin{equation}\label{I-10-45}
 |\mathcal W(A)(b)|
       \le C s_b(A)(\alpha_0+\mathsf n_b(A)),
 \qquad s_b(A)\le1/64.
\end{equation}

We now separate the linear term. For fixed $A$ and $b$,
write the entire expansion
$\mathcal W(tA)(b)=\sum_{q\ge1}t^qw_q$.
If $s_b(A)=0$, locality gives
$\mathcal W(tA)(b)=\mathcal W(0)(b)=0$, and all
coefficients vanish. Otherwise, put
$T=1/(64s_b(A))$.
For $|t|\le T$, \eqref{I-10-45} gives
\[
 |\mathcal W(tA)(b)|
 \le C|t|s_b(A)
             \bigl(\alpha_0+|t|\mathsf n_b(A)\bigr).
\]
For any $0<R\le T$, the Cauchy estimate for the linear
coefficient therefore gives
\[
 |w_1|\le Cs_b(A)
             \bigl(\alpha_0+R\mathsf n_b(A)\bigr).
\]
Letting $R$ tend to zero removes the quadratic part and
yields
$|w_1|\le C'\alpha_0s_b(A)$.
For the coefficients of order at least two, use the
largest available circle, $R=T$:
\[
 |w_q|\le C T^{1-q}s_b(A)
             \bigl(\alpha_0+T\mathsf n_b(A)\bigr),
 \qquad q\ge2.
\]
If $s_b(A)\le1/128$, then $T\ge2$ and
$\sum_{q\ge2}T^{1-q}\le2T^{-1}$. Consequently
\begin{align*}
 |\mathcal V(A)(b)|
 &\le\sum_{q\ge2}|w_q|\\
 &\le2Cs_b(A)\bigl(\alpha_0T^{-1}+\mathsf n_b(A)\bigr)
 \le C''s_b(A)\mathsf n_b(A).
\end{align*}
The last inequality uses
$\alpha_0T^{-1}=64\alpha_0s_b(A)\le64\mathsf n_b(A)$.
Choose
\[
 \rho_1=1/128,\qquad C_4=\max\{64,C',C''\}.
\]
This proves the stated local bounds for $\Delta'A=w_1$
and $\mathcal V(A)$. The global bounds follow from
$s_b(A)\le\sup|A|$ and $\mathsf n_b(A)\le\|A\|_X$.
No smallness of $g_b(A)$ was used.

The derivative of $F'$ at zero is
$D^{\mathsf T}D+\Delta'$, since $\mathcal V'(0)=0$.
For complex $A,B$,
\[
 \bigl(B,(D^{\mathsf T}D+\Delta')A\bigr)_\eta
 =\frac{\partial^2}{\partial t\,\partial u}
                         F(tA+uB)\Big|_{t=u=0}.
\]
Mixed derivatives commute, so this form is symmetric.
The form $(B,D^{\mathsf T}DA)_\eta=(DB,DA)_\eta$ is
symmetric as well; hence $\Delta'$ is symmetric.
Taking $B=A$ identifies $D^{\mathsf T}D+\Delta'$ with
the Hessian.

Finally, suppose $\partial U_0(p)=\mathbf1$ for all
$p\in\mathsf{pl}(b)$. Then $c_p=1$ and $F_p=0$ there.
In \eqref{eq:action-four-remainders} the first term
vanishes. At $tA$, the bounds for $N_p$,
$M_p^{\mathrm r}$ and $\mathsf L_{p,b}^{\mathsf T}M_p$
show that each remaining term is $O(|t|^2)$ as $t\to0$:
$N_p$ is quadratic,
$|M_p^{\mathrm r}|\le\eta^4|\Theta_p|^3$, and
$|\mathsf L_{p,b}^{\mathsf T}M_p|\le128s|M_p|$
with $s=|t|s_b(A)$ and $M_p=\Theta_p+M_p^{\mathrm r}$.
The constants here may depend on the fixed $A$ and $\eta$.
Thus the coefficient $w_1$ vanishes, which is the
locally flat assertion.

\subsection{The current and the real plaquette bound}

\begin{proof}[Completion of \cref{I-10-42}]
We prove (ii) by differentiating exponential coordinates.
For complex matrices $X,Y$,
\[
 \left.\frac{d}{dt}\right|_{t=0}e^{X+tY}
       =\bigl[\mathsf f(\operatorname{ad}_X)Y\bigr]e^X.
\]
Indeed, integrating the derivative of
$s\mapsto e^{s(X+tY)}e^{(1-s)X}$ gives
\[
 e^{X+tY}-e^X
       =t\int_0^1e^{s(X+tY)}Ye^{(1-s)X}\,ds.
\]
After differentiation at $t=0$, the right side becomes
$[\int_0^1 e^{sX}Ye^{-sX}\,ds]e^X$. Left and right
multiplication by $X$ commute, so termwise integration of
the exponential series yields
\[
 \int_0^1 e^{sX}Ye^{-sX}\,ds
       =\sum_{q\ge0}\frac{\operatorname{ad}_X^qY}{(q+1)!}
       =\mathsf f(\operatorname{ad}_X)Y.
\]

Now let $A$ be real, $E=e^{i\eta A}U_0$, and let
$\delta A$ be complex. The formula at each positive bond
shows that the curves
\[
 t\longmapsto e^{i\eta(A+t\delta A)}U_0,
 \qquad
 t\longmapsto e^{it\eta B}E,
 \qquad
 B(b)=\mathsf f(i\eta\operatorname{ad}_{A(b)})\delta A(b),
\]
have the same derivative at zero. A commutator has trace
zero, so $B$ is a complex one-form. Applying the first
variation formula \eqref{I-2-16} at $E$ gives
\[
 (F'(A),\delta A)_\eta
      =\bigl(J(E),\mathsf f(i\eta\operatorname{ad}_A)
                                      \delta A\bigr)_\eta.
\]
Cyclicity of the trace gives
\[
 \langle[A(b),X],Z\rangle
      =-\langle X,[A(b),Z]\rangle.
\]
Thus the transpose of
$\mathsf f(i\eta\operatorname{ad}_{A(b)})$ for the
bilinear pairing is
$\mathsf f(-i\eta\operatorname{ad}_{A(b)})$. The
nondegeneracy of the pairing proves the identity in (ii).

For the norm bounds, use the operator norm on linear maps
of the matrix space. If $\|z\|\le1/4$, the power series
and $(q+1)!\ge2\cdot3^{q-1}$ give
\[
 \|\mathsf f(z)-I\|
 \le\frac12\|z\|\sum_{q\ge1}12^{-(q-1)}
 =\frac6{11}\|z\|\le\frac3{22}.
\]
In particular $\|\mathsf f(z)-I\|\le\|z\|$. The
Neumann series gives
\[
 \|\mathsf f(z)^{-1}-I\|
 \le\frac{(6/11)\|z\|}{1-3/22}
 =\frac{12}{19}\|z\|\le2\|z\|.
\]
If $\eta\sup|A|\le1/8$, then
$z=-i\eta\operatorname{ad}_{A(b)}$ satisfies
$\|z\|\le2\eta|A(b)|\le1/4$. Applying these two
bounds to the identity in (ii), and then to its inverse,
gives the stated factors $1+2\eta|A(b)|$ and
$1+4\eta|A(b)|$.

For (iii), fix a plaquette $p$ with $s_p(A)\le\rho_1$
and use the transported variables of \cref{I-10-38}.
They are real. Put
\[
 X=e^{i\eta\alpha_1}e^{i\eta\alpha_2},
 \qquad
 Y=e^{-i\eta\alpha_3}e^{-i\eta\alpha_4}.
\]
These matrices are unitary, and \eqref{I-10-40} gives
$\partial E(p)=X\partial U_0(p)Y$. Consequently,
\begin{align*}
 |\partial E(p)-\mathbf1|
 &\le|\partial U_0(p)-\mathbf1|+|XY-\mathbf1|\\
 &=|\partial U_0(p)-\mathbf1|+|YX-\mathbf1|.
\end{align*}
Here the last equality follows by unitary conjugation.
Since $s_p(A)\le\rho_1\le1/64$, \cref{I-10-38}
identifies the series logarithm of $YX$ as
$i\eta^2\Theta_p$. The real logarithm assertion in
\cref{I-3-15} gives $\eta^2\Theta_p\in\mathfrak g$.
For $H\in\mathfrak g$, its eigenvalues are
$\pm|H|$, and hence
\[
 |e^{iH}-\mathbf1|=2|\sin(|H|/2)|\le|H|.
\]
The bound on $\Theta_p$ in \cref{I-10-38} now gives
\[
 |YX-\mathbf1|
 \le\eta^2|\Theta_p|
 \le\bigl(2g_p(A)+64s_p(A)^2\bigr)\eta^2.
\]
This proves the local estimate because $C_4\ge64$;
the global form follows from
$s_p(A)\le\sup|A|$ and $g_p(A)\le\sup|\nabla A|$.

Finally, let $\delta=A_1-A_2$. If $\delta=0$, (iv) is
immediate. Otherwise put
$A_t=tA_1+(1-t)A_2$, $0\le t\le1$, and
$R=r/\|\delta\|_X$. For $|z|=R$,
\[
 \|A_t+z\delta\|_X\le2r\le\rho_1,
 \qquad
 \sup|\mathcal V(A_t+z\delta)|\le C_4(2r)^2.
\]
The Cauchy formula for the entire function
$z\mapsto\mathcal V(A_t+z\delta)$ therefore gives
\[
 \sup|\mathcal V'(A_t)\delta|
 \le R^{-1}C_4(2r)^2
 =4C_4r\,\|\delta\|_X.
\]
Integrating this estimate in
$\mathcal V(A_1)-\mathcal V(A_2)
 =\int_0^1\mathcal V'(A_t)\delta\,dt$
proves (iv).
\end{proof}

\section{Quantitative bounds for averaging}
\label{sec:quantitative-averaging}

We estimate the curvature of successive block averages and
their distance from straight bare transports. The centered
mean and corner comb constructions have the same leading
curvature bound: one averaging step multiplies the fine
curvature bound by $L^2$, with a quadratic error. After
rescaling, these errors are summable over the averaging steps.

Assume $L\ge13$ throughout this section. At one step,
$0\le n<K$, use level-$n$ scaling, so that fine bonds have
length one. For $y\in\mathbb T^{(n+1)}$, put
\[
 o_{\mathsf m}(y)=y,
 \qquad
 o_{\mathsf c}(y)=y-\mathsf h(e_1+e_2+e_3),
 \qquad \mathsf h=(L-1)/2.
\]
The symbols $\mathsf m$ and $\mathsf c$ denote the centered
mean and corner comb schemes. For $x\in B(y)$, the former
uses the contour mean already constructed in
\cref{sec:block-averaging}, while the latter uses the comb
of \eqref{I-3-90}:
\[
 W_{\mathsf m}(y,x)
   =\operatorname{av}\bigl((W(\Gamma))_{\Gamma\in\mathcal G(y,x)}\bigr),
 \qquad
 W_{\mathsf c}(y,x)=W(\Gamma_{o_{\mathsf c}(y),x}).
\]
The centered variable is defined when the contour family
has diameter less than $1/40$. The corner construction uses
Ba{\l}aban's logarithmic contour average with the lowest
corner as block origin \cite[(42)--(43), pp.~23--24]{B-CMP98}.

Let $c=\langle c_-,c_-+Le_\mu\rangle$ be a positive
coarse bond. Write
$\mathsf D_c=B(c_-)\cup B(c_+)$ and let $\mathsf P_c$
be all fine bonds whose endpoints belong to $\mathsf D_c$.
Write $\mathsf P_c^+$ for its positive bonds.
For $x\in B(c_-)$, put $x'=x+Le_\mu$. Denote straight
transport from $o_\star(c_-)$ to $o_\star(c_+)$ by
$W_\star[c]$, where $\star\in\{\mathsf m,\mathsf c\}$.
For a $G^c$-valued configuration on $\mathsf P_c$, set
\begin{equation}\label{I-10-58}
\begin{aligned}
 W_x^\star
   &=W_\star(c_-,x)W([x,x'])
                   W_\star(c_+,x')^{-1}W_\star[c]^{-1},\\
 \mathcal M_\star(W)(c)
   &=\exp\left(L^{-3}\sum_{x\in B(c_-)}\log W_x^\star\right)
                    W_\star[c].
\end{aligned}
\end{equation}
This formula is used when all its contour variables are
defined and $|W_x^\star-\mathbf1|<1$ for every $x$; all
logarithms in it are series logarithms. Define the value at
$\bar c$ by inversion. For the centered scheme this is
\eqref{I-3-53}. Write $\mathcal M_\star^0=\mathrm{id}$
and $\mathcal M_\star^{j+1}=\mathcal M_\star\mathcal M_\star^j$
where defined.

Let $o_\star^{(j)}$ be the composition of the base maps
at levels $j-1,\ldots,0$, with $o_\star^{(0)}=\mathrm{id}$.
In bare units,
\[
 o_{\mathsf m}^{(j)}(y)=y,
 \qquad
 o_{\mathsf c}^{(j)}(y)
       =y-\tfrac12(L^j-1)(e_1+e_2+e_3),
\]
so the latter is the lowest corner of $B^j(y)$.
For a level-$j$ bond $c$, write $U[\![c]\!]_\star$ for
the transport of the bare configuration $U$ along the
straight path from $o_\star^{(j)}(c_-)$ to
$o_\star^{(j)}(c_+)$. Finally, for real $W$, put
\[
 \varpi(W)=\max_p|\partial W(p)-\mathbf1|,
 \qquad
 \varpi_c(W)=\max_{p\text{ with corners in }\mathsf D_c}
                            |\partial W(p)-\mathbf1|.
\]

\begin{proposition}[Iterated averages]\label{I-10-101}
There are positive constants $c^{\mathrm{SO}}_2$,
$\varpi'_L$ and $\kappa_4$, depending only on $L$, with
the following property. Let $1\le k\le K$, use level-$k$
scaling with $\eta=L^{-k}$, and let
$0<\alpha_0\le c^{\mathrm{SO}}_2$. Suppose the real bare
configuration $U$ satisfies
\[
 |\partial U(p)-\mathbf1|\le\alpha_0\eta^2
 \qquad\text{at every bare plaquette }p.
\]
For each $\star\in\{\mathsf m,\mathsf c\}$ and
$0\le n\le k$, the iterate $\mathcal M_\star^nU$ is
defined and real, and
\begin{equation}\label{I-10-102}
\begin{aligned}
 \varpi(\mathcal M_\star^nU)
      &\le2\alpha_0(L^n\eta)^2\le\varpi'_L,\\
 |\mathcal M_\star^nU(c)-U[\![c]\!]_\star|
      &\le\kappa_4\alpha_0(L^n\eta)^2
      \qquad(c\text{ a bond of }\mathbb T^{(n)}).
\end{aligned}
\end{equation}
For the centered scheme,
$\mathcal M_{\mathsf m}^nU=\mathcal M^nU$, and the
small-curvature branch \eqref{I-3-49} applies at every
participating bond of every step. Moreover, $\mathcal M^k$
is real-analytic on
\[
 \left\{U\in\mathcal V_0:
       |\partial U(p)-\mathbf1|<c^{\mathrm{SO}}_2\eta^2
                  \text{ for every bare plaquette }p\right\}.
\]

Under the same global hypothesis, the estimates have the
following local form. For a plaquette or bond $\mathfrak e$
of $\mathbb T^{(n)}$, put
\[
 \mathsf Z(\mathfrak e)
     =\bigcup_{y\text{ a vertex of }\mathfrak e}B^n(y),
 \qquad
 \alpha(\mathfrak e)
     =\max_{p\text{ bare, with corners in }\mathsf Z(\mathfrak e)}
                  \eta^{-2}|\partial U(p)-\mathbf1|,
\]
with an empty maximum equal to zero. Then the curvature
bound in \eqref{I-10-102} at a plaquette $\mathfrak e$
holds with $\alpha(\mathfrak e)$ in place of $\alpha_0$,
and so does the transport bound at a bond $\mathfrak e$.
The constants are chosen in
\cref{subsec:averaging-single-step,subsec:averaging-iteration}.
\end{proposition}

The one-step proof uses a local axial gauge and cancellation
of the contour terms around a coarse plaquette, as in
Ba{\l}aban's argument \cite[(47)--(51), pp.~25--26]{B-CMP98}.
We give the calculation for both schemes. The surviving
linear term is an average of sums over $L\times L$ squares;
it has total weight $L^2$. A uniform holomorphic neighborhood
controls the quadratic error. The normalized curvature
recursion then has a bounded product, while a separate
recursion controls the straight transports.

\subsection{Domains and base sites}

\begin{lemma}\label{I-10-59}
The contours used by either scheme stay in their block,
and the straight paths in \eqref{I-10-58} stay in the
double block. The corner contours translate with the block,
and $o_{\mathsf c}^{(j)}(y)$ is its lowest bare corner.

For either scheme, the domain of \eqref{I-10-58} is open
in the positive-bond coordinates, and the contour variables
and the average extend holomorphically to an open subset
of the ambient matrix space. At real configurations in
this domain their values lie in $G$.

Let $u:\mathsf D_c\to G$. The domain is invariant under
the real gauge transformation $u$, and
\begin{equation}\label{I-10-60}
\begin{aligned}
 (W^u)_\star(y,x)
   &=u(o_\star(y))W_\star(y,x)u(x)^{-1},\\
 \mathcal M_\star(W^u)(c)
   &=u(o_\star(c_-))\mathcal M_\star(W)(c)
                                  u(o_\star(c_+))^{-1}.
\end{aligned}
\end{equation}
For the centered scheme, if $W$ is real and
$\varpi_c(W)<\omega_M$, its average is the
small-curvature branch \eqref{I-3-49}; this restriction is
real-analytic.

At a level-$j$ bond $c$, the value and definedness of
$\mathcal M_\star^jU(c)$ depend only on the bare bonds
with endpoints in $B^j(c_-)\cup B^j(c_+)$.
\end{lemma}

\begin{proof}
In level-$n$ units, $B(y)$ has integer coordinates
$y+\sum_\nu t_\nu e_\nu$, $|t_\nu|\le\mathsf h$.
The corner comb adjusts coordinates monotonically from
$y-\mathsf h\sum_\nu e_\nu$ to $x$, in the order
$3,2,1$, and thus stays in $B(y)$. Its formula translates
with the block. The centered contours have the same
containment property by their definition. The double
block has coordinate intervals
$[-\mathsf h,\mathsf h+L]$ in direction $\mu$ and
$[-\mathsf h,\mathsf h]$ in the other directions, relative
to $c_-$. Each interval contains at most $2L$ sites,
and $2L<L^{m+K-n}$, so these charts are injective.
Thus all the straight paths in
\eqref{I-10-58} lie in it. Composing the corner translations
in bare units gives
\[
 o_{\mathsf c}^{(j)}(y)
 =y-\mathsf h(1+L+\cdots+L^{j-1})\sum_\nu e_\nu
 =y-\tfrac12(L^j-1)\sum_\nu e_\nu,
\]
the lowest site of $B^j(y)$.

On $\mathrm{GL}(2,\mathbb C)^{\mathsf P_c^+}$, extend
bond variables by inversion on reverse bonds. Each corner
transport is then a product of matrix variables and their
inverses. Its domain imposes only the finitely many open
conditions $|W_x^{\mathsf c}-\mathbf1|<1$, so the average
is holomorphic there. For the centered scheme, these
properties were proved for \eqref{I-3-53}, using the
holomorphic mean in \cref{I-3-28}. At $G^c$-valued points
the logarithms are traceless, so their exponential and the
average lie in $G^c$. For real configurations the contour
variables and loops lie in $G$, their series logarithms
belong to $i\mathfrak g$ by \cref{I-3-15}, and the average
lies in $G$.

The first identity in \eqref{I-10-60} follows by cancellation
along a path for the corner scheme and by \eqref{I-3-24}
for the centered scheme. Writing $u_-=u(o_\star(c_-))$,
it gives
\[
 (W^u)_x^\star=u_-W_x^\star u_-^{-1}.
\]
Unitary conjugation preserves the norm condition defining
each logarithm and commutes with its series. The centered
contour diameters are preserved as well. Substitution in
\eqref{I-10-58} proves the second identity and one direction
of domain preservation; applying it to $u^{-1}$ gives the
reverse direction. Centered branch agreement and
real-analyticity follow from \cref{I-3-37,I-3-55}.

Finally, the one-step formula uses only the double block.
If locality holds after $j$ steps, every value needed for
the next step lies on a level-$j$ bond $c'$ whose endpoints
belong to $B(c_-)\cup B(c_+)$. Block nesting gives
\[
 B^j(c'_- )\cup B^j(c'_+)
      \subset B^{j+1}(c_-)\cup B^{j+1}(c_+).
\]
This proves the assertion, including definedness, by
induction from $j=0$. Reverse bonds follow by inversion.
\end{proof}

\subsection{Path estimates and one averaging step}
\label{subsec:averaging-path-estimates}

We use fine-lattice units in this subsection. A nonparallel
exchange of two adjacent steps changes a path transport by
one plaquette variable, up to unitary conjugation. Counting
these exchanges gives bounds proportional to the curvature.
The two uses below are a gauge with small bond variables
on a box and a comparison of the average with straight
transport between its block origins.

A bond function $a$ takes values in $\mathfrak g^c$ on
positive bonds and is extended by $a(\bar b)=-a(b)$.
It is real when its values lie in $\mathfrak g$.
We write $a(\Gamma)=\sum_{b\in\Gamma}a(b)$ and
$e^{ia}(b)=e^{ia(b)}$, with no spacing factor in the exponent.

For $V\in G$, put
$\theta(V)=\arccos(\tfrac12\tr V)\in[0,\pi]$.
Its eigenvalues are $e^{\pm i\theta(V)}$, so
\[
 |V-\mathbf1|=2\sin\frac{\theta(V)}2,\qquad
 |V-\mathbf1|\le\theta(V)\le\frac\pi2|V-\mathbf1|.
\]
By \cref{I-3-13}, the intrinsic logarithm satisfies
$|i^{-1}\log V|=\theta(V)$ when $V\ne-\mathbf1$.
If $|V-\mathbf1|<1$, this is the series logarithm and
$\theta(V)<\pi/3$.

\begin{lemma}[Path exchanges]\label{I-10-48}
Let $W$ be a $G$-valued configuration on a reversal-symmetric
set of bonds of $\mathbb T^{(n)}$, $0\le n\le K$.
Suppose that $\Gamma_0,\ldots,\Gamma_N$ arise successively
by exchanging adjacent steps, with every intermediate path
in the domain of $W$. Suppose also that every plaquette
associated with a nonparallel exchange lies in that domain.
List these plaquettes, with repetitions, as $p_1,\ldots,p_{N'}$.
Then
\begin{equation}\label{I-10-49}
 |W(\Gamma_N)-W(\Gamma_0)|
 \le\sum_{j=1}^{N'}|\partial W(p_j)-\mathbf1|.
\end{equation}
\end{lemma}

\begin{proof}
The adjacent-step calculation in the proof of \cref{I-3-7}
is local to the exchanged steps. Parallel exchanges leave
the transport unchanged; for a nonparallel exchange the
quotient of the two transports is a unitary conjugate of
$\partial W(p_j)$ or its inverse. Thus its distance from
$\mathbf1$ is $|\partial W(p_j)-\mathbf1|$.
Summing the differences of consecutive transports proves
\eqref{I-10-49}.
\end{proof}

\begin{lemma}[Comb gauge on a box]\label{I-10-51}
Let $0\le n\le K$. In level-$n$ units, let $\mathsf Y$ be the lattice sites
in a coordinate box whose sides are less than half the
torus side. Suppose that $\mathsf Y$ has at most $s+1$
sites in each direction, where $s\ge0$.
Let $W$ be $G$-valued on all bonds with endpoints in
$\mathsf Y$, and suppose, for some $\varpi\ge0$, that
$|\partial W(p)-\mathbf1|\le\varpi$ whenever all four
corners of $p$ lie in $\mathsf Y$.
For $y\in\mathsf Y$ define
$h(x)=W(\Gamma_{y,x})$, using the comb \eqref{I-3-90}.
Then $h(y)=\mathbf1$, and for
$b=\langle x,x+e_\nu\rangle$ with endpoints in $\mathsf Y$,
\begin{equation}\label{I-10-52}
 |W^h(b)-\mathbf1|
 \le\varpi\sum_{\nu'<\nu}|x_{\nu'}-y_{\nu'}|
 \le\min\{2s,|x-y|_1\}\varpi.
\end{equation}
If $2s\varpi<2$, then $W^h=e^{ia^0}$ for a real bond
function $a^0$ with $\sup|a^0|\le\pi s\varpi$.
If $\Gamma_1,\Gamma_2$ are paths in $\mathsf Y$
with the same initial and final points and with $\ell_1,\ell_2$
bonds, respectively, then
\[
 |W(\Gamma_1)-W(\Gamma_2)|
 \le2s\varpi(\ell_1+\ell_2).
\]
\end{lemma}

\begin{proof}
Use chart coordinates in the containing box.
No bond with endpoints in $\mathsf Y$ wraps in these
coordinates: its two lifts differ by the indicated unit
step, since a nonzero period would make their separation
at least half the torus side.
In particular, the combs stay in $\mathsf Y$.
The comb from $y$ to itself is empty, and
\[
 W^h(b)=W(\Gamma_{y,x}\cup b)
                  W(\Gamma_{y,x+e_\nu})^{-1}.
\]
The comb adjusts coordinates in the order $3,2,1$.
Move the final step $e_\nu$ backwards through its
$\sum_{\nu'<\nu}|x_{\nu'}-y_{\nu'}|$ predecessors in
directions smaller than $\nu$. It then prolongs the
steps in direction $\nu$, or cancels the last step
$-e_\nu$. The resulting transport is that of
$\Gamma_{y,x+e_\nu}$.
Every intermediate site and every exchanged plaquette
has coordinates between those of $y$ and those of $x$
or $x+e_\nu$; hence they remain in $\mathsf Y$.
Equation \eqref{I-10-49} proves the first bound.
At most two coordinates occur in the sum, each of size
at most $s$. Their total is also at most $|x-y|_1$:
the chart differences are the least coordinate lifts
because they are smaller than half a period.

If $2s\varpi<2$, every positive bond value of $W^h$
avoids $-\mathbf1$. Define
$a^0(b)=i^{-1}\log W^h(b)$ by the intrinsic logarithm,
and extend by oddness to reversed bonds. Then
\[
 |a^0(b)|=\theta(W^h(b))
 \le\frac\pi2|W^h(b)-\mathbf1|
 \le\pi s\varpi.
\]

For the last assertion, the common endpoints give
the same unitary factors in the gauge transforms
of both path transports. Thus
\[
 |W(\Gamma_1)-W(\Gamma_2)|
 =|W^h(\Gamma_1)-W^h(\Gamma_2)|.
\]
Every factor in either transformed path is within
$2s\varpi$ of $\mathbf1$, by \eqref{I-10-52},
also for reversed bonds. Telescoping each unitary
product gives
$|W^h(\Gamma_i)-\mathbf1|\le2s\varpi\ell_i$.
The triangle inequality proves the claim.
\end{proof}

\begin{lemma}[Spread about the contour mean]\label{I-10-53}
Let $0\le n<K$, $y\in\mathbb T^{(n+1)}$, and let $V$
be $G$-valued on the bonds inside $B(y)$. Put
\begin{equation}\label{I-10-54}
 \varpi_y=\max_{p\subset B(y)}|\partial V(p)-\mathbf1|,
 \qquad \mathsf i=\frac34(L-1)^2,
\end{equation}
where $p\subset B(y)$ means that all four corners lie
in $B(y)$. For $x\in B(y)$, the family
$(V(\Gamma))_{\Gamma\in\mathcal G(y,x)}$ has diameter
at most $\mathsf i\varpi_y$.
If this family has exactly one mean $V(y,x)$, then
for every $\Gamma\in\mathcal G(y,x)$,
\begin{equation}\label{I-10-55}
\begin{split}
 |V(\Gamma)V(y,x)^{-1}-\mathbf1|
 &\le\theta\bigl(V(\Gamma)V(y,x)^{-1}\bigr)\\
 &\le\frac{2\pi}{3\sqrt3}\,\mathsf i\varpi_y
   =\frac{\pi}{2\sqrt3}(L-1)^2\varpi_y .
\end{split}
\end{equation}
\end{lemma}

\begin{proof}
Write $x-y=\sum_\nu n_\nu e_\nu$, with
$|n_\nu|\le(L-1)/2$. A contour takes the steps in
each direction together; two contours differ only
in the order of these groups.
Changing one order to another exchanges each pair
of steps in differently ordered directions once.
The number of nonparallel exchanges is therefore at most
\[
 \sum_{\nu<\nu'}|n_\nu|\,|n_{\nu'}|
 \le3\left(\frac{L-1}{2}\right)^2=\mathsf i.
\]
All intermediate coordinates lie between those of $y$
and $x$, so the exchanges stay inside $B(y)$.
Equation \eqref{I-10-49} gives the asserted diameter.

We need the following estimate for the inverse exponential:
if $Y,Y'\in\mathfrak g$ and
$|Y|,|Y'|\le\vartheta<\pi$, with $\vartheta>0$, then
\begin{equation}\label{eq:averaging-inverse-exponential}
 |Y-Y'|\le\frac{\vartheta}{\sin\vartheta}
                         |e^{iY}-e^{iY'}|.
\end{equation}
Put $a=|Y|$, $a'=|Y'|$, and
$\langle Y,Y'\rangle=qaa'$, $-1\le q\le1$;
if $aa'=0$, take any $q$ in this interval.
The exponential formula in the proof of \cref{I-3-13}
gives
\[
\begin{split}
 f(q):=|e^{iY}-e^{iY'}|^2
   &=2-2\cos a\cos a'-2q\sin a\sin a',\\
 g(q):=|Y-Y'|^2&=a^2+a'^2-2qaa'.
\end{split}
\]
Both functions are affine in $q$, and
\[
\begin{array}{ll}
 f(1)=4\sin^2\frac{a-a'}2,&g(1)=(a-a')^2,\\
 f(-1)=4\sin^2\frac{a+a'}2,&g(-1)=(a+a')^2.
\end{array}
\]
The function $\sin u/u$ decreases on $(0,\pi)$:
the numerator $u\cos u-\sin u$ of its derivative
vanishes at zero and has derivative $-u\sin u<0$.
Since $|a-a'|/2\le(a+a')/2\le\vartheta$,
we have
$f(\pm1)\ge(\sin\vartheta/\vartheta)^2g(\pm1)$.
Affineness gives the same inequality for all
$q\in[-1,1]$, proving
\eqref{eq:averaging-inverse-exponential}.
The same monotonicity gives
$\vartheta/\sin\vartheta<2\pi/(3\sqrt3)$
when $0<\vartheta<\pi/3$.

Let $M=V(y,x)$ and index the contour family by $\sigma$.
The defining equation of its mean says that
\[
 Z_\sigma=i^{-1}\log\bigl(V(\Gamma_\sigma)M^{-1}\bigr)
 \in\mathfrak g,\qquad
 |Z_\sigma|<\pi/3,\qquad
 \sum_\sigma Z_\sigma=0.
\]
Take $\vartheta=\max_\sigma|Z_\sigma|$; the case
$\vartheta=0$ is immediate. Unitary invariance and
\eqref{eq:averaging-inverse-exponential} imply
\[
 |Z_\sigma-Z_{\sigma'}|
 \le\frac{2\pi}{3\sqrt3}
       |V(\Gamma_\sigma)-V(\Gamma_{\sigma'})|
 \le\frac{2\pi}{3\sqrt3}\,\mathsf i\varpi_y.
\]
If the family has $N$ members, then
$Z_\sigma=N^{-1}\sum_{\sigma'}(Z_\sigma-Z_{\sigma'})$.
This proves the angle bound in \eqref{I-10-55};
the norm bound follows from $|V-\mathbf1|\le\theta(V)$.
\end{proof}

The mean contributes two endpoint errors when a coarse
bond is compared with straight transport. The corner
scheme uses actual path transports at these endpoints.
The following constants cover both choices:
\[
 \kappa^{\mathrm{SO}}_1=\frac{\pi}{2\sqrt3}(L-1)^2,
 \qquad
 \kappa^{\mathrm{SO}}_2=2L(L-1)+2\kappa^{\mathrm{SO}}_1,
 \qquad
 \kappa_3=\frac\pi2\kappa^{\mathrm{SO}}_2.
\]

\begin{lemma}[One-step transport estimate]\label{I-10-61}
Let $0\le n<K$, let $c$ be a positive bond of
$\mathbb T^{(n+1)}$, and let $W$ be $G$-valued on
$\mathsf P_c$. Suppose
\[
 \varpi:=\varpi_c(W)
 <\min\{\omega_M,(\kappa^{\mathrm{SO}}_2)^{-1}\}.
\]
For either scheme $\star$, its contour variables and
$\mathcal M_\star(W)(c)$ are defined and lie in $G$.
Moreover, for $y=c_-$ or $c_+$, $x\in B(y)$, and
$\Gamma\in\mathcal G(y,x)$,
\[
 \theta\bigl(W(\Gamma)W_{\mathsf m}(y,x)^{-1}\bigr)
 \le\kappa^{\mathrm{SO}}_1\varpi.
\]
For $x\in B(c_-)$ we have
$|W^\star_x-\mathbf1|\le\kappa^{\mathrm{SO}}_2\varpi$,
and
\[
 \mathcal M_\star(W)(c)=e^{iY}W_\star[c],
 \qquad Y\in\mathfrak g,\qquad |Y|\le\kappa_3\varpi.
\]
\end{lemma}

\begin{proof}
The centered contour variables are the unique means
by \cref{I-10-59}, since $\varpi<\omega_M$.
Their angle bound follows from \cref{I-10-53} on each
endpoint block. The corner contours are transports
along combs contained in those blocks.

Let $\mu$ be the direction of $c$ and put $x'=x+Le_\mu$.
For the centered scheme choose any
$\Gamma\in\mathcal G(c_-,x)$; for the corner scheme
take $\Gamma=\Gamma_{o_{\mathsf c}(c_-),x}$.
In either case
$\Gamma'=\Gamma+Le_\mu$ is the corresponding contour
from $o_\star(c_+)$ to $x'$.
Compare
\[
 \Gamma\cup[x,x']
 \quad\hbox{and}\quad
 [o_\star(c_-),o_\star(c_+)]\cup\Gamma'.
\]
Move the $L$ final steps $e_\mu$ in the first path
to its beginning. Each passes at most $2(L-1)$
steps in directions different from $\mu$.
Thus at most $2L(L-1)$ exchanges are nonparallel.
All intermediate sites and exchanged plaquette corners
are of the form $z+je_\mu$, where $z$ lies on $\Gamma$
and $0\le j\le L$. They lie in $\mathsf D_c$, so
\eqref{I-10-49} gives
\begin{equation}\label{eq:averaging-corridor-exchange}
 \bigl|W(\Gamma)W([x,x'])
             -W_\star[c]W(\Gamma')\bigr|
 \le2L(L-1)\varpi.
\end{equation}
For the corner scheme the two contour transports here
are exactly its contour variables. Multiplication
by their unitary inverses gives
$|W^{\mathsf c}_x-\mathbf1|\le2L(L-1)\varpi$.
For the centered scheme, insert these two path
transports in
\[
 |W^{\mathsf m}_x-\mathbf1|
 =\bigl|W_{\mathsf m}(c_-,x)W([x,x'])
            -W_{\mathsf m}[c]W_{\mathsf m}(c_+,x')\bigr|.
\]
Each replacement costs at most
$\kappa^{\mathrm{SO}}_1\varpi$ by the angle bound.
Together with \eqref{eq:averaging-corridor-exchange},
this proves the bound with $\kappa^{\mathrm{SO}}_2$.

In particular, $|W^\star_x-\mathbf1|<1$ for every
$x\in B(c_-)$. All logarithms defining the average
therefore exist, including those for the corner scheme.
They belong to $i\mathfrak g$, and
\[
 Y=L^{-3}\sum_{x\in B(c_-)}i^{-1}\log W^\star_x
\]
satisfies
\[
 |Y|\le L^{-3}\sum_{x\in B(c_-)}
          \frac\pi2|W^\star_x-\mathbf1|
 \le\kappa_3\varpi.
\]
Formula \eqref{I-10-58} gives the asserted factorization
and its $G$-valuedness.
\end{proof}

\subsection{The logarithm near the trivial configuration}
\label{subsec:averaging-holomorphic}

The comb gauge in \cref{I-10-51} turns small plaquette
variables on a box into small bond variables. In that
gauge, holomorphy gives a quadratic remainder after
linearizing the average. We first identify this
linearization; its contour terms will cancel around
a coarse plaquette.

For a bond function $a$ and $x\in B(y)$, write
\[
 \bar a_{\mathsf m}(y,x)
   =\frac1{|\mathcal G(y,x)|}
      \sum_{\Gamma\in\mathcal G(y,x)}a(\Gamma),
 \qquad
 \bar a_{\mathsf c}(y,x)
   =a(\Gamma_{o_{\mathsf c}(y),x}),
\]
where $a(\Gamma)$ is the oriented sum of $a$ along
the path.

\begin{lemma}[Logarithm of a small average]\label{I-10-93}
There is $r_*\in(0,1/16]$, depending only on $L$, with
the following property. Let $0\le n<K$, let $c$ be a
positive bond of $\mathbb T^{(n+1)}$, and let
$\star\in\{\mathsf m,\mathsf c\}$. On the space of
complex bond functions $a$ on $\mathsf P_c$ with
$\sup|a|<r_*$, the series logarithm
\[
 \Phi_\star(a)(c)
   =i^{-1}\log\mathcal M_\star(e^{ia})(c)
\]
is defined and holomorphic, vanishes at $a=0$, and
has norm at most $1$. In the chart $x\mapsto x-c_-$,
with fine bonds of length one, this map depends only
on $L$, the scheme, and the direction of $c$.

The same radius is chosen for the following relative maps.
In these units write
$c=\langle c_-,c_-+Le_\nu\rangle$ and
\[
 b_0(c)=\left\langle c_-+\frac{L-1}{2}e_\nu,\,
                         c_-+\frac{L+1}{2}e_\nu\right\rangle .
\]
For $X\in\mathfrak g^c$, let $A_{c,X}$ equal $X$ on this
positive fine bond and zero on the other positive fine bonds;
configuration values on reverse bonds are obtained by inversion.
For complex $a^0$ on $\mathsf P_c$ with $\sup|a^0|<r_*$
and $|X|<r_*$, all contour variables, averages and series
logarithms in
\[
 F_{\star,c}(a^0,X)
 =i^{-1}\log\left[
   \mathcal M_\star(e^{iA_{c,X}}e^{ia^0})(c)
                  \mathcal M_\star(e^{ia^0})(c)^{-1}\right]
\]
are defined. This map is jointly holomorphic in $(a^0,X)$
and has norm at most $1$. Its expression in the same chart
depends only on $L$, the scheme and the direction of $c$.

The derivative of $\Phi_\star$ at zero is
\begin{equation}\label{I-10-100}
\begin{split}
 (\mathbf Q^{\mathrm{lin}}_\star a)(c)
 ={}&L^{-3}\sum_{x\in B(c_-)}
 \bigl[\bar a_\star(c_-,x)+a([x,x'])\\
 &\hspace{6em}-\bar a_\star(c_+,x')\bigr].
\end{split}
\end{equation}
For the centered scheme this is \eqref{I-3-52}.
For either scheme, $y\in\{c_-,c_+\}$, and $x\in B(y)$,
the contour variable has derivative
\[
 \left.\frac d{dt}(e^{ita})_\star(y,x)\right|_{t=0}
 =i\bar a_\star(y,x).
\]
\end{lemma}

\begin{proof}
In fine-lattice units the chart $x\mapsto x-c_-$
identifies $\mathsf D_c$ with the double block of
$\langle0,Le_\mu\rangle$ in $\mathbb Z^3$, where
$\mu$ is the direction of $c$. Blocks, contour
families, combs and straight paths become their
fixed counterparts in this chart.
Thus there are six families of finitely many
functions to consider: two schemes and three
directions. None depends on the level or on the
location of $c$.

By \cref{I-10-59}, the contour variables and averages
are holomorphic in the matrix entries on an open set
containing the trivial configuration. The maps
$a^0\mapsto e^{ia^0}$ and
$(a^0,X)\mapsto e^{iA_{c,X}}e^{ia^0}$ are holomorphic
and take values in $G^c$. The pivot bond lies between
the two blocks, so its perturbation changes no contour
variable within either block. Thus the relative map above
is the relative one-step logarithm with unchanged endpoint
corrections. Products and inverses are holomorphic, and
at $(a^0,X)=(0,0)$ all arguments of the series logarithms
equal $\mathbf1$. The logarithmic maps are consequently
defined and jointly holomorphic in a neighborhood of zero,
where they vanish. Choose a radius on which all these maps,
including the relative maps, have norm at most $1$.
Taking the minimum over the six finite families and $1/16$
gives one common $r_*$ depending only on $L$.

For the centered scheme, \eqref{I-3-60} and its proof give
both derivatives. For the corner scheme, differentiation of
the path product gives
$\left.\frac d{dt}(e^{ita})_{\mathsf c}(y,x)\right|_{t=0}
=i\bar a_{\mathsf c}(y,x)$.
Let $a_{\mathsf c}[c]$ be the sum of $a$ along the straight
path from $o_{\mathsf c}(c_-)$ to $o_{\mathsf c}(c_+)$.
The derivative at zero of the loop variable in \eqref{I-10-58} is
\[
 i\bigl[\bar a_{\mathsf c}(c_-,x)+a([x,x'])
       -\bar a_{\mathsf c}(c_+,x')-a_{\mathsf c}[c]\bigr].
\]
The differential of $\log$ at $\mathbf1$ is the identity.
Averaging this expression therefore gives
$i(\mathbf Q^{\mathrm{lin}}_{\mathsf c}a)(c)-ia_{\mathsf c}[c]$
as the derivative of the exponent. Its last term cancels the
derivative $ia_{\mathsf c}[c]$ of the final path factor in
\eqref{I-10-58}. Hence
$\left.\frac d{dt}\mathcal M_\star(e^{ita})(c)\right|_{t=0}
=i(\mathbf Q^{\mathrm{lin}}_\star a)(c)$ for both schemes.
Applying the logarithm proves
$\Phi_\star'(0)a=\mathbf Q^{\mathrm{lin}}_\star a$.
\end{proof}

When this lemma is applied after the real gauge
transformation $W^h=e^{ia}$, the coarse endpoint
frames are $h(o_\star(c_-))$ and $h(o_\star(c_+))$,
as in \eqref{I-10-60}. Consequently the plaquette
of the coarse average changes by unitary conjugation.
This permits the curvature estimate below to be
proved using the small potential $a$ on the four
blocks surrounding the coarse plaquette.

\subsection{The one-step curvature bound}
\label{subsec:averaging-single-step}

Let $0\le n<K$ and use level-$n$ scaling. For a coarse
plaquette $p'=p_{\mu\nu}(y_1)$ and $x\in B(y_1)$, let
$S_x(p')$ consist of the fine plaquettes
\[
 p_{\mu\nu}(x+ie_\mu+je_\nu),\qquad 0\le i,j<L.
\]
Define the nonnegative kernel
\begin{equation}\label{I-10-95}
 \mathsf W(p',p)
   =L^{-3}\#\{x\in B(y_1):p\in S_x(p')\},
 \qquad
 (\mathsf Wf)(p')=\sum_p\mathsf W(p',p)f(p).
\end{equation}

\begin{lemma}\label{I-10-96}
Every $S_x(p')$ has $L^2$ plaquettes, whose corners lie
in the union of the four blocks of the corners of $p'$.
Moreover,
\[
 (\mathsf Wf)(p')
     =L^{-3}\sum_{x\in B(y_1)}\sum_{p\in S_x(p')}f(p),
 \qquad
 \sum_p\mathsf W(p',p)=L^2,
 \qquad
 \sum_{p'}\mathsf W(p',p)=L^{-1}.
\]
Consequently, for $f\ge0$,
\[
 \sup\mathsf Wf\le L^2\sup f,
 \qquad
 \sum_{p'}(\mathsf Wf)(p')=L^{-1}\sum_pf(p),
 \qquad
 \sum_{p'}(\mathsf Wf)(p')^2\le L\sum_pf(p)^2.
\]
\end{lemma}

\begin{proof}
The $L^2$ base points $x+ie_\mu+je_\nu$, $0\le i,j<L$,
are distinct, since the fine torus has more than $L$ sites
in each direction. The corners of their plaquettes have
$0\le i,j\le L$. Relative to $y_1$, their coordinates
therefore belong to $[-\mathsf h,\mathsf h+L]$ in the
two plaquette directions and to $[-\mathsf h,\mathsf h]$
in the third. These are the coordinate intervals of the
union of the four blocks. The formula for $\mathsf Wf$
follows by exchanging the count over $x$ with the sum
over $p$. Taking $f=1$ proves the row sum.

Fix $p=p_{\mu\nu}(z)$. Only coarse plaquettes with the
same two directions contribute to its column. The condition
$p\in S_x(p_{\mu\nu}(y_1))$ holds precisely when
\[
 x=z-ie_\mu-je_\nu\in B(y_1),\qquad 0\le i,j<L.
\]
These are $L^2$ distinct sites, each in exactly one coarse
block. Summing the count in \eqref{I-10-95} over $y_1$
therefore gives $L^{-3}L^2=L^{-1}$.

The first two consequences follow by summation, since
$\mathsf W\ge0$. For the last, Cauchy--Schwarz and the
row sum give
\[
 (\mathsf Wf)(p')^2
   \le L^2\sum_p\mathsf W(p',p)f(p)^2.
\]
Summing over $p'$ and using the column sum proves the
factor $L^2L^{-1}=L$.
\end{proof}

With $r_*$ from \cref{I-10-93} and the constants of
\cref{I-10-61}, fix
\[
\begin{aligned}
 \varpi_L&=\min\left\{\tfrac12\omega_M,
       \tfrac12(\kappa^{\mathrm{SO}}_2)^{-1},
       \frac{r_*}{4\pi L}\right\},\\
 \varpi'_L&=\min\left\{\varpi_L,\frac{r_*}{8\pi L}\right\},
 \qquad
 C^{\mathrm{SO}}_0=(2\pi L)^2(35L^2+40r_*^{-2}).
\end{aligned}
\]
These are positive constants depending only on $L$, with
$C^{\mathrm{SO}}_0\ge1$.

\begin{lemma}[One-step curvature estimate]\label{I-10-97}
Let $p'=p_{\mu\nu}(y_1)$ be a plaquette of
$\mathbb T^{(n+1)}$, with corners $y_1,y_2,y_3,y_4$
in cyclic order, and put $\mathsf Y=\bigcup_{i=1}^4B(y_i)$.
Let $W$ be a real configuration on the fine bonds whose
endpoints belong to $\mathsf Y$, and suppose
\[
 \varpi=\max_{p\text{ with corners in }\mathsf Y}
                  |\partial W(p)-\mathbf1|\le\varpi'_L.
\]
For either scheme the average is defined at the four
bonds of $p'$, and
\begin{equation}\label{I-10-99}
 |\partial\mathcal M_\star(W)(p')-\mathbf1|
 \le\sum_p\mathsf W(p',p)|\partial W(p)-\mathbf1|
                    +C^{\mathrm{SO}}_0\varpi^2.
\end{equation}
Only plaquettes with corners in $\mathsf Y$ contribute
to the sum. In particular,
\begin{equation}\label{I-10-98}
 |\partial\mathcal M_\star(W)(p')-\mathbf1|
                    \le L^2\varpi+C^{\mathrm{SO}}_0\varpi^2.
\end{equation}
\end{lemma}

\begin{proof}
The union $\mathsf Y$ is a coordinate box with $2L$
sites in directions $\mu,\nu$ and $L$ in the third
direction. Its sides are less than half the fine torus
side, since $2L<L^{m+K-n}$.
Apply \cref{I-10-51} with $s=2L-1$.
Its logarithm condition holds because
\[
 2(2L-1)\varpi\le4L\varpi_L\le r_*/\pi<2.
\]
We obtain a real gauge transformation $u$ on $\mathsf Y$
such that $W^u=e^{ia}$, with $a$ real and
\[
 \sup|a|\le\pi(2L-1)\varpi\le2\pi L\varpi\le r_*/4.
\]
The double block of each boundary bond of $p'$ lies in
$\mathsf Y$. Its curvature bound is at most
$\varpi_L<\min\{\omega_M,(\kappa^{\mathrm{SO}}_2)^{-1}\}$,
so \cref{I-10-61} defines the four coarse variables.
By \eqref{I-10-60}, gauge transformation conjugates their
plaquette variable by $u(o_\star(y_1))$. It likewise
conjugates each fine plaquette variable. Unitary
conjugation preserves all norms in \eqref{I-10-99},
so we may replace $W$ by $W^u=e^{ia}$.

Write $s_a=\sup|a|$, and let
\[
 c^{(1)}=\langle y_1,y_2\rangle,\quad
 c^{(2)}=\langle y_2,y_3\rangle,\quad
 c^{(3)}=\langle y_4,y_3\rangle,\quad
 c^{(4)}=\langle y_1,y_4\rangle.
\]
Put $\phi(c)=\Phi_\star(a)(c)$ at these four positive
bonds. The average is real, and its series logarithm
exists by \cref{I-10-93}, so $\phi(c)\in\mathfrak g$.
That lemma bounds $\Phi_\star$ by one on the complex
ball of radius $r_*$ and gives $\Phi_\star(0)=0$.
The one-variable Schwarz bound on each complex line,
followed by the supremum over norm-one linear functionals,
gives $|\Phi_\star(a)(c)|\le s_a/r_*$.
The Cauchy coefficient formula on the same lines bounds
the terms of degree $q$ by $(s_a/r_*)^q$. Since the
derivative is \eqref{I-10-100}, summing from $q=2$ gives
\[
 |\phi(c)|\le s_a/r_*,
 \qquad
 |\phi(c)-(\mathbf Q^{\mathrm{lin}}_\star a)(c)|
       \le\frac{s_a^2}{r_*(r_*-s_a)}
       \le2s_a^2r_*^{-2}.
\]

We use the elementary exponential-product estimate
\[
 |e^{X_1}\cdots e^{X_q}-e^{X_1+\cdots+X_q}|
       \le2\sigma^2,
 \qquad \sigma=\sum_j|X_j|\le1,\quad q\le4.
\]
Indeed, multiplying the exponential series bounds the
remainder after its constant and linear terms by
$e^\sigma-1-\sigma\le\sigma^2$; the exponential of
the sum has the same bound. For the coarse plaquette
the four exponents are $i\phi(c^{(1)})$,
$i\phi(c^{(2)})$, $-i\phi(c^{(3)})$ and
$-i\phi(c^{(4)})$, whose norms sum to at most
$4s_a/r_*\le1$. Their sum is skew-Hermitian, and
$|e^{iH}-\mathbf1|\le|H|$ for $H\in\mathfrak g$.
It follows that
\begin{align*}
 |\partial\mathcal M_\star(e^{ia})(p')-\mathbf1|
 &\le|\phi(c^{(1)})+\phi(c^{(2)})
              -\phi(c^{(3)})-\phi(c^{(4)})|\\
 &\quad+32\max_j|\phi(c^{(j)})|^2.
\end{align*}

We next compute the same signed combination of the linear
terms. For $x\in B(y_1)$ put
\[
 x^{(1)}=x,\quad x^{(2)}=x+Le_\mu,\quad
 x^{(3)}=x+Le_\mu+Le_\nu,\quad x^{(4)}=x+Le_\nu.
\]
These translations are bijections from $B(y_1)$ onto the
four corner blocks. Thus \eqref{I-10-100} gives
\begin{align*}
 &(\mathbf Q^{\mathrm{lin}}_\star a)(c^{(1)})
 +(\mathbf Q^{\mathrm{lin}}_\star a)(c^{(2)})
 -(\mathbf Q^{\mathrm{lin}}_\star a)(c^{(3)})
 -(\mathbf Q^{\mathrm{lin}}_\star a)(c^{(4)})\\
 &=L^{-3}\sum_{x\in B(y_1)}\Bigl\{
   \bar a_\star(y_1,x)+a([x,x^{(2)}])
                         -\bar a_\star(y_2,x^{(2)})\\
 &\qquad+\bar a_\star(y_2,x^{(2)})+a([x^{(2)},x^{(3)}])
                         -\bar a_\star(y_3,x^{(3)})\\
 &\qquad-\bar a_\star(y_4,x^{(4)})-a([x^{(4)},x^{(3)}])
                         +\bar a_\star(y_3,x^{(3)})\\
 &\qquad-\bar a_\star(y_1,x)-a([x,x^{(4)}])
                         +\bar a_\star(y_4,x^{(4)})\Bigr\}.
\end{align*}
At each corner the same contour sum occurs once with
each sign, for either choice of scheme. The remaining
expression is
\[
 L^{-3}\sum_{x\in B(y_1)}a(\partial S_x),
 \qquad
 a(\partial S_x)=a([x,x^{(2)}])+a([x^{(2)},x^{(3)}])
                -a([x^{(4)},x^{(3)}])-a([x,x^{(4)}]).
\]
Every interior bond of $S_x$ occurs once with each sign
in the sum of its plaquette boundaries. Hence
\[
 a(\partial S_x)=\sum_{p\in S_x}a(\partial p),
 \qquad
 a(\partial p)=a(b_1)+a(b_2)-a(b_3)-a(b_4),
\]
with the four positive bonds ordered as in \eqref{I-2-3}.

At a fine plaquette, the exponential-product estimate gives
\[
 |e^{ia(\partial p)}-\partial e^{ia}(p)|\le32s_a^2,
\]
since the four norms sum to at most $4s_a\le1$.
For $H\in\mathfrak g$ with $|H|\le\pi$, the spectral formula and
$\sin u\ge u-u^3/6$ give
\[
 |H|\le|e^{iH}-\mathbf1|+\tfrac1{24}|H|^3.
\]
Applying this with $H=a(\partial p)$, whose norm is at
most $4s_a\le1/16$, yields
\begin{align*}
 |a(\partial p)|
 &\le|\partial W(p)-\mathbf1|+32s_a^2
                                  +\tfrac{64}{24}s_a^3\\
 &\le|\partial W(p)-\mathbf1|+35s_a^2.
\end{align*}
All plaquettes of $S_x$ lie in $\mathsf Y$ by
\cref{I-10-96}, so this estimate applies to them.
Combining the four Taylor errors, the coarse product
error, and the square identity gives
\begin{align*}
 |\partial\mathcal M_\star(W)(p')-\mathbf1|
 &\le L^{-3}\sum_{x\in B(y_1)}\sum_{p\in S_x}
                          |a(\partial p)|+40s_a^2r_*^{-2}\\
 &\le\sum_p\mathsf W(p',p)|\partial W(p)-\mathbf1|
                       +(35L^2+40r_*^{-2})s_a^2.
\end{align*}
Since $s_a\le2\pi L\varpi$, the last term is at most
$C^{\mathrm{SO}}_0\varpi^2$, proving \eqref{I-10-99}.
Every nonzero kernel entry uses a plaquette in
$\mathsf Y$, where the fine curvature is at most
$\varpi$. The row sum $L^2$ therefore gives
\eqref{I-10-98}.
\end{proof}

\subsection{Iteration and local estimates}
\label{subsec:averaging-iteration}

With the constants of \cref{I-10-61,I-10-97}, set
\[
 c^{\mathrm{SO}}_2
   =\min\{\tfrac12\varpi'_L,L^2/(4C^{\mathrm{SO}}_0)\},
 \qquad
 \kappa_4=4\kappa_3L^{-2}.
\]
We also write $\varepsilon_\flat:=c^{\mathrm{SO}}_2$.
The quadratic curvature error contributes a factor
proportional to $L^{-2(k-n)}$ at step $n$. Its sum is
bounded independently of the number of steps. The
transport error obeys a separate recursion, since
each coarse straight path contains $L$ fine bonds.

\begin{proof}[Proof of \cref{I-10-101}]
Write $W_n=\mathcal M_\star^nU$ whenever this iterate is
defined, and set
\[
 \varpi_n=\varpi(W_n)=\alpha_0(L^n\eta)^2p_n.
\]
The hypothesis gives $p_0\le1$. Suppose that $n<k$
and $p_n\le2$. Then
\[
 \varpi_n\le2\alpha_0(L^n\eta)^2
       \le2c^{\mathrm{SO}}_2\le\varpi'_L.
\]
In particular,
$\varpi_c(W_n)\le\varpi_n\le\varpi_L
 <\min\{\omega_M,(\kappa^{\mathrm{SO}}_2)^{-1}\}$
at every positive bond of level $n+1$.
Thus \cref{I-10-61} defines the next real iterate.
Apply \cref{I-10-97} to the four blocks at each
level-$(n+1)$ plaquette. It gives
\[
 \varpi_{n+1}
 \le L^2\varpi_n(1+C^{\mathrm{SO}}_0L^{-2}\varpi_n),
\]
and consequently
\[
 p_{n+1}
 \le p_n\bigl(1+2C^{\mathrm{SO}}_0L^{-2}
                    \alpha_0L^{-2(k-n)}\bigr).
\]
If $p_j\le2$ for $0\le j\le n$, then
$\sum_{\ell\ge1}L^{-2\ell}\le2L^{-2}$ yields
\begin{align*}
 p_{n+1}
 &\le\prod_{j=0}^{n}
       \bigl(1+2C^{\mathrm{SO}}_0L^{-2}
                         \alpha_0L^{-2(k-j)}\bigr)\\
 &\le\exp(4C^{\mathrm{SO}}_0L^{-4}\alpha_0)
 \le e^{1/L^2}<2.
\end{align*}
Here $\alpha_0\le L^2/(4C^{\mathrm{SO}}_0)$ bounds
the exponent by $1/L^2$.
Induction proves existence of all iterates and the
curvature bound in \eqref{I-10-102}.

For the transport bound, let
\[
 \mathsf d_n
   =\max_c|W_n(c)-U[\![c]\!]_\star|.
\]
We have $\mathsf d_0=0$. Reversing a bond replaces
both matrices by their inverses; for unitary $X,Y$,
$|X^{-1}-Y^{-1}|=|X-Y|$. It suffices, therefore,
to estimate positive bonds.
For a positive level-$(n+1)$ bond $c$,
\cref{I-10-61} gives
\[
 W_{n+1}(c)=e^{iY}(W_n)_\star[c],
 \qquad
 |e^{iY}-\mathbf1|\le |Y|\le\kappa_3\varpi_n.
\]
The straight transport $(W_n)_\star[c]$ has $L$
factors $W_n(c'')$, each within $\mathsf d_n$ of
$U[\![c'']\!]_\star$. The corresponding bare paths
join to the straight path defining $U[\![c]\!]_\star$:
the origins satisfy
$o_\star^{(n+1)}=o_\star^{(n)}\circ o_\star$,
and $o_\star^{(n)}$ is the identity in the centered
scheme and a fixed translation in the corner scheme.
Telescoping the two products of unitary matrices gives
\[
 |W_{n+1}(c)-U[\![c]\!]_\star|
 \le2\kappa_3\alpha_0(L^n\eta)^2+L\mathsf d_n.
\]
Hence, for $n\ge1$,
\begin{align*}
 \mathsf d_n
 &\le2\kappa_3\alpha_0\eta^2
                  \sum_{j<n}L^{n-1-j}L^{2j}\\
 &\le2\kappa_3\alpha_0\eta^2
                  \frac{L^{2n-1}}{L-1}\\
 &\le4\kappa_3L^{-2}\alpha_0(L^n\eta)^2
   =\kappa_4\alpha_0(L^n\eta)^2.
\end{align*}
The same bound holds at $n=0$ because its left side
vanishes.

For $\star=\mathsf m$, the curvature bound gives
$\varpi(W_n)\le\varpi'_L<\omega_M$ at every step.
The branch agreement in \cref{I-10-59}, applied
successively, therefore gives
$\mathcal M_{\mathsf m}^nU=\mathcal M^nU$ and puts
each bond in the branch \eqref{I-3-49}.
Every $U$ in the open set stated in
\cref{I-10-101} satisfies the hypotheses with
$\alpha_0=c^{\mathrm{SO}}_2$.
All intermediate configurations there remain in
the real-analytic domains of \cref{I-10-59}.
Thus $\mathcal M^k$ is real-analytic on that set
as a composition of real-analytic maps. The set
is open because it is defined by finitely many
strict inequalities between continuous functions.

We now prove the local refinements. All iterates
have already been defined by the global hypothesis.
The local coefficients in the statement satisfy
$0\le\alpha(\mathfrak e)\le\alpha_0$.
Set
\[
 \pi_0=1,\qquad
 \pi_{n+1}
   =\pi_n\bigl(1+2C^{\mathrm{SO}}_0L^{-2}
                       \alpha_0L^{-2(k-n)}\bigr).
\]
The preceding product estimate gives $\pi_n<2$
for $0\le n\le k$. We claim that
\[
 |\partial W_n(p')-\mathbf1|
       \le\alpha(p')(L^n\eta)^2\pi_n
\]
at every level-$n$ plaquette $p'$.
At $n=0$ this is the definition of $\alpha(p')$.
Suppose it holds at level $n$, and let $p'$ be
a level-$(n+1)$ plaquette with corners $y_i$.
If a level-$n$ plaquette $p$ has all its corners
in $\bigcup_iB(y_i)$, block nesting gives
\[
 \mathsf Z(p)\subset\bigcup_iB^{n+1}(y_i)
       =\mathsf Z(p'),
 \qquad \alpha(p)\le\alpha(p').
\]
The local curvature bound $\varpi$ used in
\cref{I-10-97} therefore satisfies
\[
 \varpi\le\alpha(p')(L^n\eta)^2\pi_n
       \le2\alpha_0L^{-2(k-n)}\le\varpi'_L.
\]
By \eqref{I-10-98},
\begin{align*}
 |\partial W_{n+1}(p')-\mathbf1|
 &\le L^2\varpi(1+C^{\mathrm{SO}}_0L^{-2}\varpi)\\
 &\le\alpha(p')(L^{n+1}\eta)^2\pi_n
       \bigl(1+2C^{\mathrm{SO}}_0L^{-2}
                         \alpha_0L^{-2(k-n)}\bigr)\\
 &=\alpha(p')(L^{n+1}\eta)^2\pi_{n+1}.
\end{align*}
This proves the local curvature claim, including
the case $\alpha(p')=0$.

For bonds, put
\[
 \delta_0=0,\qquad
 \delta_{n+1}=L\delta_n+2\kappa_3L^{2n}.
\]
The geometric sum used above gives
$\delta_n\le4\kappa_3L^{-2}L^{2n}$.
We claim that
\[
 |W_n(c)-U[\![c]\!]_\star|
        \le\delta_n\alpha(c)\eta^2
\]
at every positive level-$n$ bond.
Both sides vanish for $n=0$.
Let $c$ be a positive level-$(n+1)$ bond.
Every level-$n$ plaquette $p$ with corners in
$\mathsf D_c$ satisfies
$\mathsf Z(p)\subset\mathsf Z(c)$, so the
local curvature bound just proved gives
\[
 \varpi_c(W_n)\le2\alpha(c)(L^n\eta)^2.
\]
The $L$ bonds $c''$ of the straight base path
have their endpoints in $\mathsf D_c$ by
\cref{I-10-59}. Hence
$\mathsf Z(c'')\subset\mathsf Z(c)$ and
$\alpha(c'')\le\alpha(c)$.
The induction hypothesis bounds each fine-bond
transport error by $\delta_n\alpha(c)\eta^2$.
Using the local estimate
$|Y|\le\kappa_3\varpi_c(W_n)$ in
\cref{I-10-61} and telescoping the same products
as before, we obtain
\begin{align*}
 |W_{n+1}(c)-U[\![c]\!]_\star|
 &\le2\kappa_3\alpha(c)(L^n\eta)^2
                 +L\delta_n\alpha(c)\eta^2\\
 &=\delta_{n+1}\alpha(c)\eta^2.
\end{align*}
Reversal preserves the left side and
$\alpha(\bar c)=\alpha(c)$.
The bounds $\pi_n<2$ and
$\delta_n\eta^2\le\kappa_4(L^n\eta)^2$
give both local refinements.
\end{proof}

\subsection{Corner averages in a small gauge}\label{subsec:corner-average-defect}

The local inverse construction uses corner averages. We compare
their linear parts with the straight averages at the identity,
keeping the entry bound needed for the weighted transpose.
The endpoint correction and linearization follow Ba{\l}aban's
averaging construction \cite[pp.~34--39, (110), (121)--(127)]{B-CMP98}.

We use the one-step blocks and corner transports defined above.
For a complex Lie-algebra-valued function $a$ on positive fine bonds,
put $\widetilde W(b)=e^{ia(b)}W(b)$ and define its value on reversed
bonds by inversion. Set
\[
\begin{aligned}
 v_{\mathsf c}(W,a)(y)
   &=\exp\left(L^{-3}\sum_{x\in B(y)}
       \log\bigl(\widetilde W_{\mathsf c}(y,x)
                        W_{\mathsf c}(y,x)^{-1}\bigr)\right),\\
 \mathcal Q^1_{\mathsf c}(W,a)(c)
   &=i^{-1}\log\bigl[v_{\mathsf c}(W,a)(c_-)^{-1}
       \mathcal M_{\mathsf c}(e^{ia}W)(c)
       v_{\mathsf c}(W,a)(c_+)
       \mathcal M_{\mathsf c}(W)(c)^{-1}\bigr].
\end{aligned}
\]
All logarithms are series logarithms. The endpoint factors
remove the infinitesimal comb averages from the bond average.
Write $\mathbf Q_{\mathsf c}(W)=d_a\mathcal Q^1_{\mathsf c}(W,0)$.
The following argument gives a uniform domain for this derivative
at real configurations whose plaquette variables are sufficiently
close to the identity.

Fix level-$n$ units, $0\le n\le K$, so $\eta=L^{-n}$, and put
$\xi_i=L^i\eta$. For $1\le i\le K$, let $\mathsf Y$ be one
block $B^i(y)$ or the double block
$\mathsf Y_c=B^i(c_-)\cup B^i(c_+)$ of a positive level-$i$ bond.
Assume that each coordinate interval of $\mathsf Y$ has at
most half as many sites as the corresponding torus direction.
Extend $U$ by coordinatewise folding, as specified in the proof,
and denote the extension by $\widetilde U$.

Let $W_q=\mathcal M_{\mathsf c}^q\widetilde U$.
For $x\in B^i(y)$, let $x_q$ be the unique level-$q$ site
with $x\in B^q(x_q)$, so $x_i=y$ and $x_0=x$. Define
\[
 T^{\mathsf c}_U(y,x)
   =(W_{i-1})_{\mathsf c}(x_i,x_{i-1})\cdots
                           (W_0)_{\mathsf c}(x_1,x_0).
\]
On a double block, define the dimensionless bond kernel by
\[
 Q_i^{\mathsf c}[U]
   =\eta\mathbf Q_{\mathsf c}(W_{i-1})\cdots
                       \mathbf Q_{\mathsf c}(W_0),
\]
restricted to the output $c$ and to input bonds in $\mathsf Y_c$.
Equivalently, it is the derivative at zero of the iteration
$a_0=\eta A$, $a_{q+1}=\mathcal Q^1_{\mathsf c}(W_q,a_q)$.
Only a neighborhood of zero is needed to take this derivative.
All kernels use counting measure on positive bonds; row and
column norms mean the supremum of the sums of entry norms.
At level zero, set $T^{\mathsf c}_U(y,y)=\mathbf1$ and
$Q_0^{\mathsf c}[U](c,b)=\eta\delta_{cb}$.

\begin{proposition}\label{I-32-37}
There are positive constants $c_{\mathsf c},\kappa_7^{\mathsf c}$
and $\kappa_\delta$, depending only on $L$, with the following
property. Let $1\le i\le K$, $0<\alpha\le c_{\mathsf c}$ and
$\varepsilon>0$. Suppose that $U=e^{i\eta A^\circ}$ on the
positive bonds in $\mathsf Y$, with $A^\circ$ real, and
\[
 \xi_i|A^\circ|\le\varepsilon,\qquad
 |\partial U(p)-\mathbf1|\le\alpha L^{-2i}
       \quad\text{for every plaquette in }\mathsf Y.
\]
The local averages just defined exist and depend only on
the data in $\mathsf Y$. On a single block,
\[
 |T^{\mathsf c}_U(y,x)-\mathbf1|
      \le\kappa_7^{\mathsf c}\alpha+3\varepsilon
       \qquad(x\in B^i(y)).
\]
On a double block, $Q_i^{\mathsf c}[U](c,b)=0$ unless both
endpoints of $b$ belong to $\mathsf Y_c$, and
\[
\begin{aligned}
 \sum_b|Q_i^{\mathsf c}[U](c,b)-Q_i^{\mathsf c}[\mathbf1](c,b)|
       &\le\kappa_\delta(\alpha+\varepsilon)\xi_i,\\
 |Q_i^{\mathsf c}[U](c,b)-Q_i^{\mathsf c}[\mathbf1](c,b)|
       &\le\kappa_\delta(\alpha+\varepsilon)\eta^3\xi_i^{-2}.
\end{aligned}
\]
The kernels also satisfy
\[
 \sum_b|Q_i^{\mathsf c}[U](c,b)|\le2\xi_i,
 \qquad |Q_i^{\mathsf c}[U](c,b)|\le2\eta^3\xi_i^{-2}.
\]
At level zero both defects vanish.
\end{proposition}

\begin{proof}
In bare lattice
coordinates, write each coordinate interval as
$\{s,\ldots,s+N_0-1\}\subset\mathbb Z/N\mathbb Z$,
where $2N_0\le N$. The folding map in that coordinate is
\[
 s+t\longmapsto
 \begin{cases}
 s+t,&0\le t<N_0,\\
 s+2N_0-1-t,&N_0\le t<2N_0,\\
 s,&2N_0\le t<N.
 \end{cases}
\]
Apply it in each coordinate and denote the resulting map by $F$.
A collapsed bond receives the value $\mathbf1$; every other bond
receives the value of $U$ on its image, in the image orientation.
Denote this extension by $\widetilde U=F^*U$.
Its restriction to $\mathsf Y$ is $U$.

Each coordinate increment of $F$ is $0$, $1$ or $-1$.
Thus a folded plaquette collapses, giving holonomy $\mathbf1$,
or maps onto a plaquette in $\mathsf Y$ with a possible change
of orientation and base point. Its holonomy is then a conjugate
of that plaquette holonomy or its inverse. The curvature bound
therefore holds everywhere for $\widetilde U$.
In level-$i$ units the bare spacing is $\eta_i=L^{-i}$;
\cref{I-10-101}, with $\alpha\le c^{\mathrm{SO}}_2$, gives
\[
 \varpi_q:=\varpi(W_q)\le2\alpha L^{-2(i-q)},\qquad
 \sum_{q<i}\varpi_q\le4\alpha L^{-2}.
\]
Nested blocks and the locality of one averaging step imply that
all factors contributing to a value on $\mathsf Y$ use only its
original data. Indeed, a contributing level-$q$ bond has its
order-$q$ endpoint blocks inside $\mathsf Y$; the same holds
for every bond used in the next finer step. This also proves
the support assertion for the derivative kernel.

We first estimate one derivative. Put
\[
 (\mathbf S_{\mathsf c}(W)a)(c)
 =L^{-3}\sum_{x\in B(c_-)}\sum_{b\in[x,x']}
 \operatorname{Ad}\bigl(W_{\mathsf c}(c_-,x)
                              W([x,b_-])\bigr)a(b),
 \qquad x'=x+Le_\mu,
\]
where $\mu$ is the direction of $c$.
For real $W$, the transports are unitary. There are $L^4$
summands at a fixed output $c$, and $L$ pairs $(c,x)$ for a
fixed input bond. Hence
\[
 \|\mathbf S_{\mathsf c}(W)\|_{\mathrm{row}}\le L,
 \qquad
 \|\mathbf S_{\mathsf c}(W)\|_{\mathrm{col}}\le L^{-2}.
\]
At $W=\mathbf1$, differentiation of the two endpoint factors
cancels the endpoint contour means in \eqref{I-10-100}.
Thus $\mathbf Q_{\mathsf c}(\mathbf1)
=\mathbf S_{\mathsf c}(\mathbf1)=\mathbf Q^s$, where
$\mathbf Q^s a=L^{-3}\sum_{x\in B(c_-)}a([x,x'])$.

For $W=e^{ia^0}$, use the finite double-block charts of
\cref{I-10-93}, now for the straightening expressions in
both variables $(a^0,a)$. At $(0,0)$ every logarithm has
argument $\mathbf1$. Holomorphy and continuity therefore give
a common radius $r_{\mathsf c}\in(0,r_*]$, depending only
on $L$, such that on $\sup|a^0|,\sup|a|<r_{\mathsf c}$
the logarithms have norm at most one and
the adjoint transports in $\mathbf S_{\mathsf c}$ have norm
at most two. Cauchy's formula in $a(b)$ bounds the kernel of
$\mathbf Q_{\mathsf c}$ by $r_{\mathsf c}^{-1}$ there.
The kernel of $\mathbf S_{\mathsf c}$ is bounded by $2L^{-2}$.
Their difference is holomorphic in $a^0$ and vanishes at zero;
the Schwarz bound on each complex line gives
\[
 |(\mathbf Q_{\mathsf c}(e^{ia^0})
       -\mathbf S_{\mathsf c}(e^{ia^0}))(c,b)|
 \le(r_{\mathsf c}^{-1}+2)r_{\mathsf c}^{-1}\sup|a^0|.
\]
For real $W$ with curvature at most $\varpi$, the comb gauge
of \cref{I-10-51} on $\mathsf D_c$ gives
$W^h=e^{ia^0}$ and $\sup|a^0|\le2\pi L\varpi$, provided
$\varpi\le r_{\mathsf c}/(4\pi L)$.
The gauge transforms of both kernels are obtained by left
composition with $\operatorname{Ad}(h(o_{\mathsf c}(c_-)))$
and right composition with $\operatorname{Ad}(h(b_-))^{-1}$.
This follows directly from \eqref{I-10-60} in the definitions;
the maps are isometries, so their norms are unchanged.
The derivative is consequently defined at $W$ and
\[
 \|\mathbf Q_{\mathsf c}(W)-\mathbf S_{\mathsf c}(W)\|
                  _{\mathrm{row},\mathrm{col}}
       \le\kappa_9^{\mathsf c}\varpi,
 \qquad \kappa_9^{\mathsf c}=\kappa_9^{\mathsf c}(L).
\]
Here a row has at most $6L^3$ entries; a fine bond belongs
to at most six coarse double blocks. The same estimates
hold after restriction to local boxes.

Choose $c_{\mathsf c}$ small enough that
\[
 c_{\mathsf c}\le c^{\mathrm{SO}}_2,\qquad
 2c_{\mathsf c}\le r_{\mathsf c}/(4\pi L),\qquad
 8\kappa_9^{\mathsf c}c_{\mathsf c}\le\log2.
\]
All one-step derivatives at $W_q$ are now defined.
For real $a$ close to zero, the arguments of the series
logarithms in the straightening formula are $G$-valued;
their logarithms belong to $i\mathfrak g$ by \cref{I-3-15}.
Thus $\mathcal Q^1_{\mathsf c}(W_q,a)$ is $\mathfrak g$-valued.
Its trace is holomorphic and vanishes on the real bond space;
the power series in a real basis therefore vanishes also for
complex $a$ near zero. Each step consequently takes values in
$\mathfrak g^c$ and vanishes at zero. Their finite composition is
defined near zero, and the chain rule gives the stated product
formula for $Q_i^{\mathsf c}$.
Products
are ordered from $q=i-1$ to $q=0$. Expanding the product
of $\mathbf Q_{\mathsf c}(W_q)
=\mathbf S_{\mathsf c}(W_q)+E_q$ and using the preceding
row and column bounds gives
\[
\begin{aligned}
 \left\|\prod_{q<i}\mathbf Q_{\mathsf c}(W_q)
        -\prod_{q<i}\mathbf S_{\mathsf c}(W_q)\right\|_{
          \mathrm{row}}
   &\le8\kappa_9^{\mathsf c}L^{-3}\alpha L^i,\\
 \left\|\prod_{q<i}\mathbf Q_{\mathsf c}(W_q)
        -\prod_{q<i}\mathbf S_{\mathsf c}(W_q)\right\|_{
          \mathrm{col}}
   &\le8\kappa_9^{\mathsf c}\alpha L^{-2i}.
\end{aligned}
\]
For example, the row norm of the difference is bounded by
$L^i[\exp(\kappa_9^{\mathsf c}L^{-1}
\sum_q\varpi_q)-1]$; the column estimate uses $L^2$ in
place of $L^{-1}$. The chosen threshold permits
$e^t-1\le2t$. The products without subtraction are bounded
by $2L^i$ and $2L^{-2i}$, respectively.

It remains to compare the leading transports with the identity.
An entry of $\prod_{q<i}\mathbf S_{\mathsf c}(W_q)$ is a
sum of $L^{-3i}\operatorname{Ad}(\tau)$ over nested choices
of a starting site in a block and a bond on its straight path.
There are $L^{4i}$ such choices from one output. There are
at most $L^i$ choices ending at a fixed input bond: at each
upward step its starting site is $b_- -s e_\mu$,
$0\le s<L$, in the units of that step, and its parent block
and parent bond are then determined.

Each factor of $\tau$ follows a corner comb and an initial
straight segment at level $q$, of total length at most $4L$.
Replace every coarse bond of that path by the bare straight
path between its order-$q$ corners. By \eqref{I-10-102},
the change in its transport is bounded by
$4L\kappa_4\alpha(L^q\eta_i)^2$.
The resulting bare paths concatenate, stay in $\mathsf Y_c$
by block nesting, and have total length at most
$\sum_{q<i}4L^{q+1}\le5L^i$.
The sum of the transport errors is at most
$8\kappa_4L^{-1}\alpha$.
The same construction for $T^{\mathsf c}_U(y,x)$ uses only
the comb portions and stays in $B^i(y)$.
For either construction, compare the resulting path with a
bare path $\Gamma$ in the same box, with the same endpoints
and at most $16L^i$ bonds. The path estimate of
\cref{I-10-51}, with $s=2L^i-1$ and curvature
$\alpha L^{-2i}$, bounds the change by $84\alpha$.
Thus there is $\kappa_7^{\mathsf c}=\kappa_7^{\mathsf c}(L)$ such that
\[
 |\tau-U(\Gamma)|\le\kappa_7^{\mathsf c}\alpha,
 \qquad
 |T^{\mathsf c}_U(y,x)-U(\Gamma)|
                         \le\kappa_7^{\mathsf c}\alpha.
\]
For the scalar transport choose the bare comb from the lowest
corner of $B^i(y)$ to $x$, of length at most $3L^i$.
For the bond transport choose the comb from the lowest corner
of $\mathsf Y_c$ to $b_-$, of length at most $4L^i$.
Telescoping products and $|e^{iX}-\mathbf1|\le|X|$ for real
$X$ give $|U(\Gamma)-\mathbf1|\le3\varepsilon$ and
$4\varepsilon$, respectively. This proves the scalar claim.
For the leading bond kernel, the corresponding adjoint
transport differs from the identity by at most
$2\kappa_7^{\mathsf c}\alpha+8\varepsilon$.
For the difference between the leading bond kernel and its value
at the identity, the two counts give a row bound
$(2\kappa_7^{\mathsf c}\alpha+8\varepsilon)L^i$ and an
entry bound $(2\kappa_7^{\mathsf c}\alpha+8\varepsilon)L^{-2i}$.

Add the product errors and multiply by $\eta$ in the
definition of $Q_i^{\mathsf c}$. Since
$\eta L^i=\xi_i$ and $\eta L^{-2i}=\eta^3\xi_i^{-2}$,
these estimates give the asserted defect bounds for a constant
$\kappa_\delta$ depending only on $L$.
The bounds without subtraction follow from the product bounds
above. The level-zero identities give the last assertion.
\end{proof}

\subsection{Solving the averaging constraint at one bond}

Fix a positive bond $c=\langle c_-,c_-+Le_\nu\rangle$ of
$\mathbb T^{(k+1)}$, where $0\le k<K$, and measure lengths in
level-$k$ units. Put $h=(L-1)/2$. The middle bond
\[
 b_0(c)=\langle c_-+he_\nu,c_-+(h+1)e_\nu\rangle
\]
is the pivot bond. For a reversal-symmetric bond set $S$, write
$\mathcal V(S)=G^{S^+}$, with inverse values on reversed bonds.
Write $\mathsf P_c^\circ$ for the bonds of the double block other
than $b_0(c)$ and its reverse. For a configuration
$w$ on $\mathsf P_c^\circ$ and $x\in G$, let $w\sqcup x$ denote its
completion with value $x$ at $b_0(c)$ and $x^{-1}$ at the reverse.
Define the partial map $\pi_c(w,v)$ to be the unique $x\in G$ satisfying
$\mathcal M(w\sqcup x)(c)=v$, whenever exactly one such $x$ exists.
The average here is the global real map, with all three branches of
\cref{sec:block-averaging}.

The local analytic inverse alone does not exclude other solutions
on the other branches. The next proposition supplies this uniqueness
and the measurability needed when the eliminated variables are recovered.
Set
\begin{equation}\label{II-6-10}
 \mathsf c_{\mathrm{piv}}=\frac52L^3,
 \qquad \omega_{\mathrm p}=\frac{\omega_M}{16L^4}
                         =\frac1{3200L^6}.
\end{equation}

\begin{proposition}[The pivot variable]\label{II-6-12}
Let $L\ge13$ be odd. Suppose $\omega,\omega'\ge0$,
$\omega+\omega'\le\omega_{\mathrm p}$, and $G_1$ is a real
configuration on $\mathsf P_c$ with
$|\partial G_1(p)-\mathbf1|\le\omega$ on every plaquette of the
double block. For every real configuration $w$ on $\mathsf P_c^\circ$
such that
\[
 |w(b)G_1(b)^{-1}-\mathbf1|\le\omega'
 \qquad(b\in\mathsf P_c^\circ),
\]
there is exactly one $x\in G$ satisfying
\begin{equation}\label{II-6-15}
 \mathcal M(w\sqcup x)(c)=\mathcal M(G_1)(c).
\end{equation}
It obeys
\begin{equation}\label{II-6-14}
 |xG_1(b_0)^{-1}-\mathbf1|\le\frac52L^3\omega',
\end{equation}
and $\varpi_c(w\sqcup x)<\omega_M$, so that the small-curvature
branch applies to the solution. In particular, these conclusions
hold for every $W$ whose average agrees with that of $G_1$ and
whose nonpivot values satisfy the displayed bound, without any
initial assumption on $W(b_0)$ or on the plaquettes of $W$.

On the full space $\mathcal V(\mathsf P_c^\circ)\times G$, the
domain of $\pi_c$ is Borel and $\pi_c$ is a Borel map on its domain.
This last assertion requires no smallness hypothesis. All constants
in \eqref{II-6-10} are independent of the level and the torus.
\end{proposition}

\begin{proof}
We first record the estimates used below. For $g\in G$, let
$\theta(g)\in[0,\pi]$ be its angle. Then
\[
 |g-\mathbf1|\le\theta(g)\le\frac\pi2|g-\mathbf1|,
 \qquad \theta(gg')\le\theta(g)+\theta(g').
\]
The angle is invariant under inversion and conjugation, and
$\theta(e^{iA})\le|A|$ for $A\in\mathfrak g$, by its eigenvalues.
For the
second inequality, only $\theta(g)+\theta(g')<\pi$ needs proof.
Use
\[
 gg'-\mathbf1=\tfrac12(g-\mathbf1)(g'+\mathbf1)
              +\tfrac12(g+\mathbf1)(g'-\mathbf1)
\]
and $|g\mp\mathbf1|=2\sin(\theta(g)/2),
2\cos(\theta(g)/2)$, respectively. Monotonicity of sine on
$[0,\pi/2]$ gives the result. Unitary multiplication preserves
the norm, and telescoping gives a bound by the sum of the bond
errors for any path product. Also
\[
 |e^{iA}-e^{iB}|\le|A-B|\qquad(A,B\in\mathfrak g),
\]
by integrating
$\frac{d}{ds}(e^{isA}e^{i(1-s)B})
=e^{isA}i(A-B)e^{i(1-s)B}$.
Together with \eqref{eq:averaging-inverse-exponential}, this gives
the corresponding logarithm Lipschitz bound
$\vartheta/\sin\vartheta$ when both angles are at most
$\vartheta<\pi$.

We also need a sharper comparison of two means. Suppose
$A,B\in G$ are means of the families $(U_j)$ and $(V_j)$
in $G$, and
\[
 |U_jA^{-1}-\mathbf1|,\ |V_jB^{-1}-\mathbf1|\le\rho\le1/12,
 \qquad \delta=\max_j|U_jV_j^{-1}-\mathbf1|.
\]
Put $X_j=U_jA^{-1}$, $Y_j=V_jB^{-1}$, $E_j=U_jV_j^{-1}$,
$g=BA^{-1}$ and $a=|g-\mathbf1|$. Since $X_j=E_jY_jg$,
$|X_j-Y_j|\le\delta+a$. The mean equations and
\eqref{I-3-17} give
\[
 0=\sum_j(E_j-\mathbf1)Y_jg+n(g-\mathbf1)
       +\sum_j(Y_j-\mathbf1)(g-\mathbf1)+\sum_jR_j,
 \qquad |R_j|\le\frac\rho{1-\rho}(\delta+a).
\]
Consequently $a\le\delta+\rho a+\rho(\delta+a)/(1-\rho)$, and
\begin{equation}\label{II-6-21}
 |AB^{-1}-\mathbf1|
 \le\frac{\delta}{(1-\rho)^2-\rho}\le(1+4\rho)\delta.
\end{equation}
The last inequality follows from
$(1+4\rho)(1-3\rho)\ge1$.

Let $\Gamma_-$ and $\Gamma_+$ be the portions of $[c_-,c_+]$
before and after the pivot. They each have $h$ bonds, all nonpivot.
Use the straight transport $u$ as unknown:
\begin{equation}\label{II-6-24}
 \begin{gathered}
 u_0=G_1[c],\qquad v_0=\mathcal M(G_1)(c),\\
 x(u)=w(\Gamma_-)^{-1}u\,w(\Gamma_+)^{-1},
 \qquad U_u=w\sqcup x(u).
 \end{gathered}
\end{equation}
This parametrizes every completion once, and $U_u[c]=u$.
Telescoping the two half-paths gives
\begin{equation}\label{II-6-25}
 |x(u)-G_1(b_0)|\le(L-1)\omega'+|u-u_0|.
\end{equation}
The $L$ sites on the axis through $c$ have identical contour
transports along that axis, and their loops $W_x$ and
$Y_{x,\Gamma,\Gamma'}$ equal $\mathbf1$ by cancellation.
Write $B'$ for the other $L^3-L$ sites and put $q=1-L^{-2}$.
For $x\in B'$, the path $[x,x']$ avoids the pivot, as do the
contours inside either block. Thus
\[
 T_x(w)=w(c_-,x)w([x,x'])w(c_+,x')^{-1},\qquad
 T_{x,\Gamma,\Gamma'}(w)=w(\Gamma)w([x,x'])w(\Gamma')^{-1}
\]
depend only on $w$, and the corresponding loops of $U_u$ are
$T_x(w)u^{-1}$ and $T_{x,\Gamma,\Gamma'}(w)u^{-1}$.

Here all the means defining $w(y,x)$ exist. Indeed,
\begin{equation}\label{II-6-23}
 16L^4\omega+8L^3\omega'\le\omega_M,
\end{equation}
so $\omega+4\omega'\le\omega_M/L^3$. The internal plaquettes
of $w$ are bounded by $\omega+4\omega'$. By
\cref{I-10-53,I-3-28}, their contour families have unique means,
and their distances from those means are at most
\[
 \kappa_1^{\mathrm{SO}}(\omega+4\omega')
 \le\rho_L:=\frac1{200L^3}.
\]
The same holds for $G_1$. Applying \eqref{II-6-21} to the two
families and telescoping their paths yields, with $C_L=1+4\rho_L$,
\[
 |w(y,x)-G_1(y,x)|\le C_L|x-y|_1\omega',
 \qquad
 |T_x(w)-T_x(G_1)|\le(2C_L|x-c_-|_1+L)\omega'\le4L\omega'.
\]
For the individual path products the angle error is at most
$2\pi L\omega'$, since their paths have at most $4L-3$ bonds.
Moreover
\[
 \sum_{x\in B'}|x-c_-|_1=\tfrac14(L^2-1)(3L^2-1).
\]
Since $C_L(3L^2-1)<3L^2$, the averaged error satisfies
\begin{equation}\label{II-6-35}
 d_w:=L^{-3}\sum_{x\in B'}|T_x(w)-T_x(G_1)|
 \le\frac52\frac{L^2-1}{L}\omega'.
\end{equation}

To exclude distant solutions, write $\mathcal M(U)(c)=e^{i\zeta(U)}U[c]$.
On the first branch, $\zeta$ is $L^{-3}$ times the sum over $B'$
of the intrinsic logarithms divided by $i$; on the second branch
use the mean over the individual contour pairs instead; on the
third put $\zeta=0$. The omitted axial terms vanish. For $G_1$,
\cref{I-10-59,I-10-61} give the first branch and
\[
 \theta(W_x(G_1))\le\kappa_3\omega,\quad
 \theta(Y_{x,\Gamma,\Gamma'}(G_1))
       \le(2\kappa_1^{\mathrm{SO}}+\kappa_3)\omega,\quad
 |\zeta(G_1)|\le q\kappa_3\omega.
\]
The second bound follows by inserting the two endpoint contour
errors around $W_x(G_1)$. The constants defined above satisfy
$\kappa_1^{\mathrm{SO}}<0.907L^2$,
$\kappa_3<5.992L^2$, and
$2\kappa_1^{\mathrm{SO}}+2\kappa_3<14L^2$.
For every $u\in G$, each of the first two branches therefore gives
\[
 |\zeta(U_u)|\le q\bigl[\theta(uu_0^{-1})+2\pi L\omega'
                  +(2\kappa_1^{\mathrm{SO}}+\kappa_3)\omega\bigr],
\]
and the third satisfies this inequality because its exponent is zero.
Apply angle subadditivity to
$uu_0^{-1}=e^{-i\zeta(U_u)}
[\mathcal M(U_u)(c)v_0^{-1}]e^{i\zeta(G_1)}$ to obtain
\begin{equation}\label{II-6-36}
 L^{-2}\theta(uu_0^{-1})
 \le\theta(\mathcal M(U_u)(c)v_0^{-1})
              +q(2\pi L\omega'+14L^2\omega).
\end{equation}
Every solution thus lies in the closed ball
\begin{equation}\label{II-6-26}
 \mathcal U=\{u\in G:|u-u_0|\le t\},\qquad
 t=(L^2-1)(14L^2\omega+2\pi L\omega').
\end{equation}

For $u\in\mathcal U$, a plaquette has at most one pivot edge.
Using \eqref{II-6-25} on that edge gives
\[
 \varpi_c(U_u)\le\omega+(L+2)\omega'+t<\omega_M.
\]
Indeed, the coefficients of $\omega$ and $\omega'$ are respectively
less than $16L^4$ and $8L^3$; if both vanish the conclusion is immediate.
Thus only the first branch occurs throughout $\mathcal U$.
The loop angles there obey
\[
 \theta(W_x(U_u))\le2\pi L\omega'+\kappa_3\omega+\tfrac\pi2t
 \le2\omega_M=:\vartheta_L,
 \qquad \theta(W_x(G_1))\le\vartheta_L.
\]
For the middle inequality, the two coefficients are at most
$22L^4$ and $9.87L^3$, bounded by twice those in \eqref{II-6-23}.

Set $\varkappa=\vartheta_L/\sin\vartheta_L$ and
$\Psi(u)=e^{-i\zeta(U_u)}v_0$. The logarithm bound now gives
\begin{equation}\label{II-6-42}
 |\Psi(u)-\Psi(u')|\le q\varkappa|u-u'|,
 \qquad |\Psi(u_0)-u_0|\le\varkappa d_w.
\end{equation}
Here $\vartheta_L=1/(100L^2)$. The inequality
$\sin s\ge s-s^3/6$ for $0\le s\le1$ gives
$s/\sin s\le1+s^2/5$, hence
\[
 \varkappa\le1+\frac1{50000L^4},\quad
 1-q\varkappa\ge L^{-2}\left(1-\frac1{50000L^2}\right)>0,
 \quad \frac{\varkappa}{1-q\varkappa}\le L^2+10^{-4}.
\]
For the last bound, put $\epsilon_1=1/(50000L^2)$ and use
$(1+\epsilon_1)/(1-\epsilon_1)\le1+3\epsilon_1$.
Equations \eqref{II-6-26} and \eqref{II-6-35} also give
\[
 (1-q\varkappa)t
 \ge\left(1-\frac1{50000L^2}\right)
             2\pi\frac{L^2-1}{L}\omega'
 \ge\left(1+\frac1{50000L^4}\right)
             \frac52\frac{L^2-1}{L}\omega'
 \ge\varkappa d_w.
\]
Thus $\Psi$ maps the complete space $\mathcal U$ into itself
and is a contraction. Its unique fixed point $u_*$ satisfies
\begin{equation}\label{II-6-45}
 |u_*-u_0|\le\frac{\varkappa d_w}{1-q\varkappa}
 \le(L^2+10^{-4})\frac52\frac{L^2-1}{L}\omega'.
\end{equation}
On $\mathcal U$, the fixed point equation is precisely
$\mathcal M(U_u)(c)=v_0$, and \eqref{II-6-36} puts every solution
on every branch in this ball. This proves existence and global
uniqueness. Finally, \eqref{II-6-25} and \eqref{II-6-45} give
\[
 |x(u_*)-G_1(b_0)|
 \le\left[L-1+\frac52(L^2+10^{-4})\frac{L^2-1}{L}\right]\omega'
 \le\frac52L^3\omega'.
\]

For measurability, no reference configuration is needed. The space
$\mathcal V(\mathsf P_c)$ is compact. With $\mathcal N_c$ as in
\cref{sec:block-averaging}, exhaust its three branches by
\begin{equation}\label{II-6-46}
 \begin{aligned}
 K_n^{\mathrm a}&=\{U:\varpi_c(U)\le(1-1/n)\omega_M\},\\
 K_n^{\mathrm b}&=\{U:\varpi_c(U)\ge\omega_M,
            \ |Y_{x,\Gamma,\Gamma'}(U)+\mathbf1|\ge1/n
                            \text{ for every }x,\Gamma,\Gamma'\},\\
 K^{\mathrm c}&=\{U:\varpi_c(U)\ge\omega_M\}\cap\mathcal N_c.
 \end{aligned}
\end{equation}
These sets are compact, and the average is continuous on each
separately by \cref{I-3-55}. Put
$F=\mathcal V(\mathsf P_c^\circ)\times G$, and let
$\Sigma_n\subset F\times G$ consist of the solutions
$(w,v,x)$ with $w\sqcup x\in K_n^{\mathrm a}\cup K_n^{\mathrm b}
\cup K^{\mathrm c}$. Each $\Sigma_n$ is a finite union of compact
graphs, and these sets increase to the full solution set $\Sigma$.
The pairs admitting at least one solution form
$F_{\ge1}=\bigcup_n\operatorname{pr}_F\Sigma_n$, a Borel set.
The pairs admitting two distinct solutions form the Borel set
\[
 F_{\ge2}=\bigcup_n\operatorname{pr}_F
 \{(w,v,x,x'):(w,v,x),(w,v,x')\in\Sigma_n,
                                  \ |x-x'|\ge1/n\},
\]
since each set being projected is compact. The use of one index
$n$ loses no pair: the $\Sigma_n$ increase, and any positive
distance eventually exceeds $1/n$. Therefore
$\operatorname{dom}\pi_c=F_{\ge1}\setminus F_{\ge2}$ is Borel.
For any closed $E\subset G$,
\begin{equation}\label{II-6-49}
 \pi_c^{-1}(E)=\operatorname{dom}\pi_c\cap
       \bigcup_n\operatorname{pr}_F(\Sigma_n\cap(F\times E))
\end{equation}
is Borel, again by compactness. Since closed sets generate the
Borel sets, this proves the last assertion.
\end{proof}


\section{Energy under repeated block averaging}\label{sec:averaging-energy}
Throughout this section, $L\ge13$ is odd.
The small-field decomposition tests the curvature of block averages.
A failed test must therefore impose a cost on the original Wilson
action. The estimate below gives this implication without assuming
that the original configuration is regular. Its proof uses the
one-step curvature estimate only while the energy remains below a
fixed threshold.

For a plaquette $p'$ of $\mathbb T^{(t)}$ and $0\le l\le t$, let
\[
 \mathsf Y_{t,l}(p')=\bigcup_{y\text{ a corner of }p'}B^{t-l}(y).
\]
Let $\mathcal N_l(p')$ be the set of plaquettes of
$\mathbb T^{(l)}$ whose four corners lie in this set, and put
$\mathcal N_l(S)=\bigcup_{p'\in S}\mathcal N_l(p')$ for a set $S$
of plaquettes of $\mathbb T^{(t)}$. In this notation the level $t$
is determined by the argument. In particular, membership in
$\mathcal N_l(S)$ requires all four corners to lie in the cone of
one plaquette of $S$.

With the constants $C^{\mathrm{SO}}_0$ and $\varpi'_L$ of
\cref{I-10-97}, fix
\begin{equation}\label{II-6-56}
 \begin{gathered}
 C_{\mathrm{en}}=\sqrt{12}\,C^{\mathrm{SO}}_0L^{-1/2},\\
 \mathsf e_L=\min\left\{\frac L2(\varpi'_L)^2,
       \frac{(\log2)^2(\sqrt L-1)^2}{8C_{\mathrm{en}}^2}\right\},
 \qquad e_*=\frac12\mathsf e_L.
 \end{gathered}
\end{equation}
These positive constants depend only on $L$.

\begin{proposition}[Energy of block averages]\label{II-6-57}
Let $1\le t\le K$, let $S$ be a set of plaquettes of
$\mathbb T^{(t)}$, and let $U\in\mathcal V_0$. Then
\begin{equation}\label{II-6-58}
 L^t\sum_{p\in\mathcal N_0(S)}a_p(U)
 \ge\min\left\{\frac14\sum_{p'\in S}
     |\partial\mathcal M^tU(p')-\mathbf1|^2,e_*\right\}.
\end{equation}
If $S$ consists of inner plaquettes of a region $X$ for the
level-$t$ cube partition, the sum on the left is at most $L^tA_X(U)$.
Consequently, let $\mathsf R$ be a set of bare sites and let
$\sigma_{\mathsf R}(U)$ be the sum of $a_p(U)$ over bare plaquettes
whose corners lie in $\mathsf R$. If $1\le l\le K$ and
$L^l\sigma_{\mathsf R}(U)<e_*$, then every plaquette $p'$ of
$\mathbb T^{(l)}$ with $\mathsf Y_{l,0}(p')\subset\mathsf R$ satisfies
\begin{equation}\label{II-6-59}
 |\partial\mathcal M^lU(p')-\mathbf1|
 \le2\bigl(L^l\sigma_{\mathsf R}(U)\bigr)^{1/2}.
\end{equation}
\end{proposition}

\begin{proof}
We first record the geometry needed to iterate the one-step estimate.
For $l<t$,
\begin{equation}\label{eq:energy-cone-recursion}
 \mathcal N_t(S)=S,\qquad
 \mathcal N_l(S)=\bigcup_{p''\in\mathcal N_{l+1}(S)}\mathcal N_l(p'').
\end{equation}
Here the cones on the right are those of one step. To check the
second identity, fix $p'=p_{\mu\nu}(y)\in S$ and put $r=L^{t-l}$.
In units of $a_{K-l}$, its cone is a rectangular box with $2r$
sites in the $\mu$ and $\nu$ directions and $r$ sites in the third
direction. It does not wrap around the torus: $2r$ is less than
half the number of sites in each direction. At consecutive levels
the boxes are related by replacing each coarse site by its block
of $L^3$ sites. Thus a fine cone of a coarse plaquette in the
coarser box is contained in the finer box.

Conversely, use coordinates $s_\rho$ starting from a corner of the
finer box. If a fine plaquette has directions $\kappa,\kappa'$, and
the coarser box has $R_\rho$ sites in direction $\rho$, choose its
coarse base coordinates to be
\[
 q_\rho^*=\min\{\lfloor s_\rho/L\rfloor,R_\rho-2\}
       \quad(\rho=\kappa,\kappa'),\qquad
 q_\rho^*=\lfloor s_\rho/L\rfloor\quad(\rho\ne\kappa,\kappa').
\]
When $l+1<t$, all $R_\rho\ge L$, so these are coordinates of a
plaquette in the coarser box. The two fine coordinates of the
plaquette in either of its directions belong to the union of the
blocks based at $q_\rho^*$ and $q_\rho^*+1$. Its remaining coordinate
belongs to the block based at $q_\rho^*$. This proves the reverse
inclusion. When $l+1=t$, the coarser box consists of the four
corners of $p'$ and the identity is immediate. Taking unions over
$p'\in S$ proves \eqref{eq:energy-cone-recursion}.

A fine plaquette belongs to the cones of at most twelve coarse
plaquettes at any fixed higher level. Indeed, its base point belongs
to a unique block of that level, whose center must be a corner of
each such coarse plaquette. There are at most four plaquettes through
a given corner in each of the three coordinate planes. Therefore,
for every nonnegative function $g$ on level-$l$ plaquettes,
\begin{equation}\label{II-6-74}
 \sum_{p'\in S}\sum_{p\in\mathcal N_l(p')}g(p)
 \le12\sum_{p\in\mathcal N_l(S)}g(p).
\end{equation}
If the corners of $p'$ belong to a region $X$ made of level-$t$
cubes, their blocks of every finer level belong to $X$. This proves
the asserted inclusion of the fine plaquettes in $X$.

We next prove a one-step estimate. For a real configuration $W$,
write $f_W(p)=|\partial W(p)-\mathbf1|$ and use the counting
$\ell^2$ norm on sets of plaquettes. Let $S$ now consist of
plaquettes of $\mathbb T^{(n+1)}$, and suppose
\[
 \varpi=\max_{p\in\mathcal N_n(S)}f_W(p)\le\varpi'_L.
\]
For $p'\in S$, put
$\varpi_{p'}=\max_{p\in\mathcal N_n(p')}f_W(p)$.
The four double blocks at the bonds of $p'$ lie in its cone.
Hence \cref{I-10-97} applies to the restriction of $W$ there and
agrees with the regular branch of the globally defined average:
\[
 f_{\mathcal M(W)}(p')\le(\mathsf Wf_W)(p')
                             +C^{\mathrm{SO}}_0\varpi_{p'}^2.
\]
The kernel $\mathsf W$ is supported in the one-step cone. Its
row and column sums in \cref{I-10-96} give
$\|\mathsf Wf_W\|_S\le\sqrt L\|f_W\|_{\mathcal N_n(S)}$,
by applying the full operator estimate to $f_W$ restricted to that
set and extended by zero. Moreover, \eqref{II-6-74} gives
\[
 \sum_{p'\in S}\varpi_{p'}^4
 \le\varpi^2\sum_{p'\in S}\sum_{p\in\mathcal N_n(p')}f_W(p)^2
 \le12\varpi^2\|f_W\|_{\mathcal N_n(S)}^2.
\]
The triangle inequality therefore yields
\begin{equation}\label{II-6-77}
 \|f_{\mathcal M(W)}\|_S
 \le\sqrt L(1+C_{\mathrm{en}}\varpi)\|f_W\|_{\mathcal N_n(S)}.
\end{equation}

Return to the original $t,S,U$, and put
\[
 E_l=\sum_{p\in\mathcal N_l(S)}
       |\partial\mathcal M^lU(p)-\mathbf1|^2,
 \qquad
 \varpi_l=\max_{p\in\mathcal N_l(S)}
       |\partial\mathcal M^lU(p)-\mathbf1|.
\]
Empty sums and maxima are zero. We have $\varpi_l\le\sqrt{E_l}$,
$E_0=2\sum_{\mathcal N_0(S)}a_p(U)$, and
$E_t=\sum_{p'\in S}|\partial\mathcal M^tU(p')-\mathbf1|^2$.
By \eqref{eq:energy-cone-recursion} and \eqref{II-6-77}, as long as
$\varpi_j\le\varpi'_L$ for $j<l$,
\begin{equation}\label{II-6-82}
 E_l^{1/2}\le L^{l/2}
       \exp\left(C_{\mathrm{en}}\sum_{j<l}\varpi_j\right)E_0^{1/2}.
\end{equation}

Suppose $L^tE_0\le\mathsf e_L$. We prove
$E_l\le2L^lE_0$ by induction. It holds at $l=0$. If it holds
through $l<t$, then, for $0\le j\le l$,
\[
 \varpi_j\le\sqrt{2\mathsf e_L}\,L^{-(t-j)/2}\le\varpi'_L,
 \qquad
 C_{\mathrm{en}}\sum_{j\le l}\varpi_j
 \le\frac{C_{\mathrm{en}}\sqrt{2\mathsf e_L}}{\sqrt L-1}
 \le\frac12\log2.
\]
The first and last inequalities use the two bounds in
\eqref{II-6-56}. Equation~\eqref{II-6-82} now gives the claimed
bound at $l+1$. In particular $E_t\le2L^tE_0$.
Thus, if $L^tE_0\le\mathsf e_L$, then
$\frac14E_t\le\frac12L^tE_0$; otherwise
$\frac12L^tE_0>e_*$. These alternatives prove \eqref{II-6-58}.

Finally, apply \eqref{II-6-58} to the single plaquette $p'$ at
level $l$. Its fine cone lies in $\mathsf R$, so
\[
 \min\left\{\tfrac14|\partial\mathcal M^lU(p')-\mathbf1|^2,e_*\right\}
 \le L^l\sigma_{\mathsf R}(U)<e_*.
\]
The minimum must be its first entry, proving \eqref{II-6-59}.
\end{proof}


\section{Regional gauge fixing and the minimizing lift}
\label{sec:regional-linear-operators}

We construct a field with prescribed straight block averages that
minimizes the curl energy. The ambiguity consists precisely of gradients
whose scalar block averages vanish. A relative Landau condition selects
one field from this affine space. The scalar projection, gauge-fixed
inverse and minimizing-lift formula follow Ba{\l}aban's construction
\cite[pp.~24--35, especially (1.103)]{B-CMP95}.
We give the proof for separated regions and for both choices of
boundary bonds.

Fix an odd integer $L\ge13$, positive integers $m,K$, and an
integer $0\le t\le K$. Use the torus of \cref{sec:finite-lattice}
and the nested lattices of \cref{sec:block-averaging},
with physical torus side $\ell=2L^m$,
and measure lengths in level-$t$ units. Thus
\[
 \xi_n=L^{n-t},\qquad \xi=\xi_0=L^{-t}.
\]
Let $\mathfrak M=L^p$ for an integer $p\ge0$, and let $R_1$ be a power of $L$
with $R_1\ge L$. Suppose that the required partitions into
$\mathfrak M\xi_n$-cubes exist, and fix regions
$\Omega_0\supset\cdots\supset\Omega_t$ which are unions of these
cubes and satisfy
\[
 \Omega_0\ne\emptyset,\qquad
 \operatorname{dist}_\infty(\Omega_n^c,\Omega_{n+1})
       >R_1\mathfrak M\xi_n\quad(0\le n<t).
\]
Here the complement is the closed-cube complement used in
\cref{sec:multiscale-sums}; set $\Omega_{t+1}=\emptyset$.
Empty later regions, the whole torus, and disconnected regions
are allowed. Write
\[
 \mathcal L_n=\overline{\Omega_n\setminus\Omega_{n+1}},\qquad
 \Lambda_n=\mathbb T^{(n)}\cap\mathcal L_n,\qquad
 \Lambda_*=\bigcup_{n=0}^t\Lambda_n,
\]
where the difference inside the closure is a point-set difference.
The sets $\Lambda_n$ are disjoint by \cref{I-9-20}.
Retain the centered cubes and blocks
\[
 \Delta^n(y)=\{x:|x-y|_\infty\le\tfrac12\xi_n\},\qquad
 B^n(y)=\Delta^n(y)\cap\mathbb T^{(0)},\qquad
 B(y)=\Delta^n(y)\cap\mathbb T^{(n-1)}\quad(n\ge1).
\]
The level in $B(y)$ will always be specified. Also write
$B^n(Y)=\bigcup_{y\in Y}B^n(y)$.

Let $\mathsf S=\mathbb T^{(0)}$, let $\mathsf E$ be the
positive bare bonds $(x,\mu)=\langle x,x+\xi e_\mu\rangle$,
and let $\mathsf F$ be the bare plaquettes $p_{\mu\nu}(x)$
with $\mu<\nu$. The real spaces $\mathcal H^0,\mathcal H^1,
\mathcal H^2$ of functions on these sets have inner products
$\langle f,g\rangle=\xi^3\sum fg$.
Extend a bond function by $A(\bar b)=-A(b)$ and put
$A(\Gamma)=\sum_{b\in\Gamma}A(b)$. The ambient differences are
\[
 \partial_\mu f(x)=\xi^{-1}\bigl(f(x+\xi e_\mu)-f(x)\bigr),
 \qquad
 (\partial\lambda)_\mu=\partial_\mu\lambda,\qquad
 (\partial A)_{\mu\nu}=\partial_\mu A_\nu-\partial_\nu A_\mu.
\]
Adjoints refer to these inner products, and
$\Delta=\partial^*\partial$ on $\mathcal H^0$.

Put $\mathsf S_0=\Omega_0\cap\mathsf S$. Choose one of the
two sets of positive bonds
\[
 \mathsf E_0^{\mathrm{in}}
   =\{b:b_-,b_+\in\mathsf S_0\},\qquad
 \mathsf E_0^{\mathrm{st}}
   =\{b:\{b_-,b_+\}\cap\mathsf S_0\ne\emptyset\},
 \qquad
 \mathsf E_0\in\{\mathsf E_0^{\mathrm{in}},
                     \mathsf E_0^{\mathrm{st}}\}.
\]
We call these the inner and star frames. Let $\chi_0,\chi_1$
be multiplication by the indicators of $\mathsf S_0,\mathsf E_0$,
respectively, and identify
$\mathcal H^i(\Omega_0)=\chi_i\mathcal H^i$, $i=0,1$,
with functions extended by zero to the whole lattice. In particular,
every ambient derivative below is taken after this extension.
Define
\[
 \mathrm d=\chi_1\partial\big|_{\mathcal H^0(\Omega_0)},\qquad
 \mathrm d^*=\chi_0\partial^*\big|_{\mathcal H^1(\Omega_0)},
 \qquad
 \Delta_\Omega=\chi_0\Delta\chi_0
                    \big|_{\mathcal H^0(\Omega_0)}.
\]
We also use the compressed curl operator
$(\partial^*\partial)(\Omega_0)
 =\chi_1\partial^*\partial\chi_1|_{\mathcal H^1(\Omega_0)}$.
Thus
\[
 \langle\lambda,\Delta_\Omega\lambda\rangle
       =\|\partial\lambda\|^2,\qquad
 \langle A,(\partial^*\partial)(\Omega_0)A\rangle
       =\|\partial A\|^2.
\]
The second norm includes every plaquette having a bond in the frame.
In the inner frame $\Delta_\Omega$ need not equal
$\mathrm d^*\mathrm d$.

For $1\le n\le t$, take the coarse positive rows
\[
 \Lambda_n^{\mathrm b}
 =\{c\text{ positive in }\mathbb T^{(n)}:
      \{c_-,c_+\}\cap\Omega_n\ne\emptyset,\ 
      \{c_-,c_+\}\cap\Omega_{n+1}=\emptyset\},
\]
and set
$\Lambda_0^{\mathrm b}
 =\{c\in\mathsf E_0:\{c_-,c_+\}\cap\Omega_1=\emptyset\}$.
In particular, a row may point out of $\Omega_n$; rows entering
$\Omega_{n+1}$ are excluded. The level-zero rows include boundary
bonds in the star frame. Let $\mathcal K^0,\mathcal K^1$ be the
real function spaces on $\Lambda_*$ and
$\bigcup_{n=0}^t\Lambda_n^{\mathrm b}$, with inner products
\[
 \langle\omega,\omega'\rangle
   =\sum_{n=0}^t\xi_n^3\sum_{y\in\Lambda_n}\omega(y)\omega'(y),
 \qquad
 \langle B,B'\rangle
   =\sum_{n=0}^t\xi_n^3\sum_{c\in\Lambda_n^{\mathrm b}}B(c)B'(c).
\]
The row sets are disjoint because their bond lengths $\xi_n$
are distinct.

For ambient scalar and bond functions define the straight averages
\begin{equation}\label{I-11-20}
 \begin{split}
 (\mathsf Q'_n\lambda)(y)
    &=L^{-3n}\sum_{x\in B^n(y)}\lambda(x),\\
 (\mathsf Q_nA)(c)
    &=L^{-4n}\sum_{x\in B^n(c_-)}
           A([x,x+\xi_ne_\mu]),
       \qquad c=\langle c_-,c_-+\xi_ne_\mu\rangle .
 \end{split}
\end{equation}
Here the path contains $L^n$ bare bonds. Restrict these values to
$y\in\Lambda_n$ and $c\in\Lambda_n^{\mathrm b}$ to obtain
\[
 \mathsf Q':\mathcal H^0(\Omega_0)\longrightarrow\mathcal K^0,
 \qquad
 \mathsf Q:\mathcal H^1(\Omega_0)\longrightarrow\mathcal K^1.
\]
These symbols always denote \eqref{I-11-20} in this section.
At level zero both averages are the identity.

Fix $a>0$ and put
\[
 \begin{gathered}
 a_0=a,\qquad a_n=\frac{a(1-L^{-2})}{1-L^{-2n}}\quad(n\ge1),
 \\
 (\mathsf a'\omega)(y)=a_n\xi_n^{-2}\omega(y)\ (y\in\Lambda_n),
 \qquad
 (\mathsf aB)(c)=a\xi_n^{-2}B(c)\ (c\in\Lambda_n^{\mathrm b}).
 \end{gathered}
\]
Let $N=\ker\mathsf Q'$, let $\mathsf R$ be the orthogonal
projection in $\mathcal H^0(\Omega_0)$ onto $\Delta_\Omega N$,
and put $\mathsf P=1-\mathsf R$. Define
\[
 \Delta'_a=\Delta_\Omega+\mathsf Q'^*\mathsf a'\mathsf Q',
 \qquad
 \Delta_a=(\partial^*\partial)(\Omega_0)
                +\mathrm d\mathsf R\mathrm d^*
                +\mathsf Q^*\mathsf a\mathsf Q.
\]

\begin{proposition}\label{I-11-30}
For the separated sequence and either frame above, the following hold.
\begin{enumerate}[label=(\alph*)]
\item For every $\lambda\in\mathcal H^0(\Omega_0)$,
\begin{equation}\label{I-11-31}
 \begin{aligned}
 \langle\lambda,\Delta'_a\lambda\rangle
 &\ge\sum_{n=0}^t\min\{a_n,4\}\xi_n^{-2}
           \sum_{x\in B^n(\Lambda_n)}\xi^3|\lambda(x)|^2,\\
 \langle\lambda,\Delta_\Omega\lambda\rangle
 &\ge4\sum_{n=1}^t\xi_n^{-2}
           \sum_{x\in B^n(\Lambda_n)}\xi^3|\lambda(x)|^2
             \qquad(\lambda\in N).
 \end{aligned}
\end{equation}
In a term with $n\ge1$, the number $4$ may be replaced by
$4L^{2n}\sin^2(\pi/(2L^n))\in[4,\pi^2)$.
Consequently $\Delta'_a$ is positive definite, its inverse
$\mathsf G'=(\Delta'_a)^{-1}$ exists and is positive definite,
and $\Delta_\Omega$ is injective on $N$.

\item If $y\in\Lambda_n$ and $x\in B^n(y)$, then
$(\mathsf Q'^*\omega)(x)=\omega(y)$.
Moreover $\mathsf Q'\mathsf Q'^*=1$, $\mathsf Q'$ is onto,
and both $\mathsf Q'\mathsf G'^2\mathsf Q'^*$ and
$\mathsf Q'\mathsf G'\mathsf Q'^*$ are positive definite.

\item
\[
 \{A\in\mathcal H^1(\Omega_0):
                  \partial A=0,\ \mathsf QA=0\}=\mathrm dN,
\]
and $\mathrm d$ is injective on $N$.

\item $\Delta_a$ is positive definite. Its inverse
$\mathsf G=\Delta_a^{-1}$ exists and is positive definite.

\item $\mathsf Q$ is onto, and
$\mathsf Q\mathsf G\mathsf Q^*$ is positive definite on $\mathcal K^1$.
\end{enumerate}
\end{proposition}

The scalar assertions follow from the gap of the Neumann difference
operator on each block; their proof is in
\cref{sec:regional-linear-scalar}. The proof of (c) in
\cref{sec:regional-linear-kernel} starts with an ambient closed
one-form and its affine potential on the covering lattice.
The zero-average equations make its block means agree across
successive layers; subtracting those means cancels the affine
increments. This identifies the joint kernel of curl and the coarse
average, including the boundary restrictions in each frame.
The kernel identity and scalar positivity give the vector positivity
in (d). A triangular construction on block-face corridors proves
surjectivity in (e), and hence positivity of
$\mathsf Q\mathsf G\mathsf Q^*$.

\begin{proposition}\label{I-11-33}
In the same setting, the projection, gauge correction and minimizing
lift have the following properties.
\begin{enumerate}[label=(\roman*)]
\item Every $\lambda\in N$ is supported in $\Omega_1$.
Its gradient and the gradient of its ambient Laplacian are supported
on inner bonds of $\Omega_0$; $\Delta\lambda$ is supported on the
bare sites at graph distance at most one from $\Omega_1$.
Moreover
\[
 \mathrm d\lambda=\partial\lambda,\qquad
 \mathrm d^*\mathrm d\lambda
       =\Delta_\Omega\lambda=\Delta'_a\lambda=\Delta\lambda,
 \qquad
 \mathrm d\Delta_\Omega\lambda=\partial\Delta\lambda .
\]
The map $\Delta_\Omega:N\longrightarrow\operatorname{ran}\mathsf R$
is bijective.

\item The projection is given by
\[
 \mathsf P=\mathsf G'\mathsf Q'^*
      (\mathsf Q'\mathsf G'^2\mathsf Q'^*)^{-1}
          \mathsf Q'\mathsf G',\qquad \mathsf R=1-\mathsf P.
\]
Furthermore
\[
 \mathsf R\mathsf G'\mathsf Q'^*=0,\qquad
 \mathsf Q'\mathsf G'\mathsf R=0,\qquad
 \mathsf G'\mathsf R\mathcal H^0(\Omega_0)=N.
\]
For fixed geometry and frame, $\mathsf R,\mathsf P$ are independent
of $a$.

\item For $A\in\mathcal H^1(\Omega_0)$, the unique $\lambda\in N$
such that $\mathsf R\mathrm d^*(A-\mathrm d\lambda)=0$ is
\[
 \lambda=\mathsf G'\mathsf R\mathrm d^*A .
\]
It is also the unique minimizer over $N$ of the full ambient norm
$\|\partial^*(A-\partial\lambda)\|^2$.
The correction preserves $\mathsf QA$ and $\partial A$.

\item For $\lambda\in N$,
$\Delta_a\mathrm d\lambda=\mathrm d\Delta'_a\lambda$.
The inverse identities are
\begin{equation}\label{I-11-38}
 \mathsf G\mathrm d\mathsf R=\mathrm d\mathsf G'\mathsf R,
 \qquad
 \mathsf R\mathrm d^*\mathsf G=\mathsf R\mathsf G'\mathrm d^* .
\end{equation}
They imply
\[
 \mathsf R\mathrm d^*\mathsf G\mathrm d
   =\mathrm d^*\mathsf G\mathrm d\mathsf R=\mathsf R,\qquad
 \mathsf R\mathrm d^*\mathsf G\mathsf Q^*=0,\qquad
 \mathsf Q\mathsf G\mathrm d\mathsf R=0.
\]

\item Define
\[
 \mathsf H=\mathsf G\mathsf Q^*
                         (\mathsf Q\mathsf G\mathsf Q^*)^{-1}.
\]
Then $\mathsf Q\mathsf H=1$ and $\mathsf R\mathrm d^*\mathsf H=0$.
For $B\in\mathcal K^1$, $\mathsf HB$ is the unique minimizer of
$\|\partial A\|^2$ subject to
$\mathsf QA=B$ and $\mathsf R\mathrm d^*A=0$.
With only the constraint $\mathsf QA=B$, the set of minimizers is
$\mathsf HB+\mathrm dN$.
For fixed geometry and frame, $\mathsf H$ is independent of $a$, and
\begin{equation}\label{I-11-39}
 \begin{split}
 (\mathsf Q\mathsf G\mathsf Q^*)^{-1}
     &=\mathsf a+\mathsf H^*(\partial^*\partial)(\Omega_0)\mathsf H,\\
 \langle B,(\mathsf Q\mathsf G\mathsf Q^*)^{-1}B\rangle
     &=\langle B,\mathsf aB\rangle+\|\partial\mathsf HB\|^2 .
 \end{split}
\end{equation}
\end{enumerate}
\end{proposition}

The support identities are proved in
\cref{sec:regional-linear-geometry}, and the projection formula in
\cref{sec:regional-linear-scalar}. The gauge correction and
minimizing-lift assertions are proved in
\cref{sec:regional-linear-minimizers}.

\subsection{Blocks and the regional boundary}
\label{sec:regional-linear-geometry}

For a level-$n$ positive bond $c=(y,\mu)$, let
\[
 \operatorname{cor}(c)
 =\{(z,\mu)\in\mathsf E:
       z\in B^n(y),\ z+\xi e_\mu\in B^n(y+\xi_ne_\mu)\}.
\]
These are the bare bonds crossing the common face of the two blocks.
In particular, $\operatorname{cor}(c)=\{c\}$ when $n=0$.
Write $\mathsf S_0^\partial$ for the sites of $\mathsf S_0$
having a bare nearest neighbour outside $\mathsf S_0$.
For $r\ge0$, let $\Omega_1^{+r}$ be the set of bare sites
reachable from a bare site of $\Omega_1$ by at most $r$ bonds.
It is empty if $\Omega_1$ is empty.

\begin{lemma}\label{I-11-27}
The separated sequence and either frame have the following properties.
\begin{enumerate}[label=(\roman*)]
\item The blocks $B^n(y)$, $y\in\Lambda_n$, $0\le n\le t$,
partition $\mathsf S_0$. Write $n(x),y(x)$ for the unique
indices such that $x\in B^{n(x)}(y(x))$.
If two sites of $\mathsf S_0$ are bare nearest neighbours,
their indices $n(x)$ differ by at most one.

\item Let $1\le n\le t$ and let $y\in\mathbb T^{(n)}\setminus\Omega_n$
be a level-$n$ nearest neighbour of a site of $\Omega_n$.
Then
\[
 \Delta^n(y)\subset\Omega_{n-1},\qquad
 B(y)\subset\Lambda_{n-1},\qquad
 B^n(y)=\bigsqcup_{z\in B(y)}B^{n-1}(z).
\]
Every endpoint of a row in $\Lambda_n^{\mathrm b}$ either
belongs to $\Lambda_n$ or is such an outer site.
For $n=0$, each endpoint belongs to $\Lambda_0$ or lies
outside $\mathsf S_0$; the latter case occurs only in the
star frame. A positive level-$n$ bond with both endpoints
in $\Lambda_n$ belongs to $\Lambda_n^{\mathrm b}$.

\item All bare bonds on the paths defining
$(\mathsf Q_nA)(c)$ have endpoints in
$B^n(c_-)\cup B^n(c_+)$. For $c\in\Lambda_n^{\mathrm b}$
and $n\ge1$, this union lies in
$\Omega_{n-1}\setminus\Omega_{n+1}$, and all these bonds
are inner to $\Omega_0$.
The corridor contains $L^{2n}$ bonds. Each of them occurs
in exactly $L^n$ of the paths defining $(\mathsf Q_nA)(c)$
and in none of the paths for another positive level-$n$ bond.

\item
\[
 \Omega_1^{+3}\subset\mathsf S_0,\qquad
 \Omega_1^{+2}\cap\mathsf S_0^\partial=\emptyset,\qquad
 \mathsf S_0^\partial\subset\Lambda_0.
\]
Every bare bond with an endpoint in $\Omega_1^{+2}$
is inner to $\Omega_0$.
\end{enumerate}
\end{lemma}

\begin{proof}
The nonempty region $\Omega_0$ contains a defining
$\mathfrak M$-cube. Its center is a bare site, so
$\mathsf S_0\ne\emptyset$.
The coordinate description of blocks in
\cref{sec:block-averaging} shows that a site of
$\mathbb T^{(j)}$, $j\le n$, lies in a unique level-$n$
cube, and that its membership in $\Omega_n$ is the
membership of that cube's center. Indeed, $L$ is odd,
so no such site lies on a face of a level-$n$ cube.
The regions $\Omega_n$ and $\Omega_{n+1}$ are unions
of level-$n$ cubes. By the layer-site description in
\cref{I-9-20}, a bare site $x$ of $\Omega_0$ has a unique
index $n$ with $x\in\Omega_n\setminus\Omega_{n+1}$;
the center of its level-$n$ block then belongs to $\Lambda_n$.
Conversely, $B^n(y)\subset\Omega_n\setminus\Omega_{n+1}$
when $y\in\Lambda_n$. Blocks at one level are disjoint,
and the point-set differences for different levels are
disjoint. This proves the partition. For $t=0$ it is
the partition into singleton blocks.

If $x,x'\in\mathsf S_0$ are bare neighbours and
$n=n(x)>n(x')$, then $x\in\Omega_n$ and $x'\notin\Omega_n$.
Separation gives
\[
 \operatorname{dist}_\infty(x',\Omega_n)
 \le\xi\le\xi_{n-1}
 <R_1\mathfrak M\xi_{n-1}
 <\operatorname{dist}_\infty(\Omega_{n-1}^c,\Omega_n).
\]
Thus $x'\in\Omega_{n-1}\setminus\Omega_n$, and
$n(x')=n-1$.

For (ii), choose a neighbour $y'\in\Omega_n\cap\mathbb T^{(n)}$
of the outer site $y$. Translation by $y'-y$ maps
$\Delta^n(y)$ to $\Delta^n(y')\subset\Omega_n$.
Every point of $\Delta^n(y)$ is consequently within
sup distance $\xi_n$ of $\Omega_n$, whereas
\[
 \operatorname{dist}_\infty(\Omega_{n-1}^c,\Omega_n)
 >R_1\mathfrak M\xi_{n-1}\ge\xi_n.
\]
The cube therefore lies in $\Omega_{n-1}$.
Its level-$(n-1)$ sites lie outside $\Omega_n$,
since its center $y$ does. Hence $B(y)\subset\Lambda_{n-1}$,
and the decomposition of $B^n(y)$ follows from block nesting.
The assertions about row endpoints now follow from the
definition of $\Lambda_n^{\mathrm b}$.

For (iii), put $h_n=(L^n-1)/2$ and write
$x=c_-+\xi\sum_\nu k_\nu e_\nu$ with $|k_\nu|\le h_n$.
Along $[x,x+\xi_ne_\mu]$, only $k_\mu$ changes.
The sites before crossing the common face lie in $B^n(c_-)$;
after crossing it they lie in $B^n(c_+)$, since subtracting
$L^n$ from the $\mu$-coordinate brings it back into
$[-h_n,h_n]$. If an endpoint of $c$ lies in $\Lambda_n$,
its block is contained in $\Omega_n\setminus\Omega_{n+1}$.
The block at an outer endpoint lies in $\Omega_{n-1}$
by (ii), and lies outside $\Omega_n$ by the block membership
property. This proves the support assertion.

A bond $(z,\mu)$ belongs to $\operatorname{cor}(c)$
precisely when
\[
 z=c_-+\xi\left(h_ne_\mu+
                    \sum_{\nu\ne\mu}k_\nu e_\nu\right),
 \qquad |k_\nu|\le h_n.
\]
There are $L^{2n}$ choices. Such a bond lies on the path
from $x$ to $x+\xi_ne_\mu$ precisely for
$x=z-i\xi e_\mu$, $0\le i<L^n$.
These are the $L^n$ sites in $B^n(c_-)$ on the line
through $z$ in direction $\mu$. They are distinct, since
the torus period in each direction is at least $2\xi_n$.
A path of another direction cannot contain this bond.
If a path of direction $\mu$ from another level-$n$ block
contains it, its initial site is one of these same $L^n$
sites. Disjointness of the blocks forces the initial block,
and hence the coarse bond, to be the original one.

Finally, the assertions in (iv) hold when $\Omega_1$ is empty.
Otherwise $t\ge1$, and each site of $\Omega_1^{+3}$ is within
sup distance $3\xi$ of $\Omega_1$. Since
\[
 3\xi\le L\xi\le R_1\mathfrak M\xi
 <\operatorname{dist}_\infty(\Omega_0^c,\Omega_1),
\]
it lies in $\mathsf S_0$.
All neighbours of a site of $\Omega_1^{+2}$ therefore lie
in $\mathsf S_0$. Such a site cannot belong to
$\mathsf S_0^\partial$, and every bond incident to it is inner.
In particular, a boundary site lies in
$\mathsf S_0\setminus\Omega_1=\Lambda_0$.
\end{proof}

For $\lambda\in N$, the identity $\mathsf Q'_0=1$
implies that $\lambda$ vanishes on
$\Lambda_0=\mathsf S_0\setminus\Omega_1$.
Consequently $\partial\lambda$ is supported on bonds incident
to $\Omega_1$, $\Delta\lambda$ is supported in $\Omega_1^{+1}$,
and $\partial\Delta\lambda$ is supported on bonds incident
to $\Omega_1^{+1}$.
The collar in \cref{I-11-27} places all these sites and bonds
inside their respective regional spaces. Thus
\begin{equation}\label{eq:regional-null-support}
 \begin{gathered}
  \mathrm d\lambda=\partial\lambda,\qquad
  \mathrm d^*\mathrm d\lambda
       =\Delta_\Omega\lambda
       =\Delta'_a\lambda=\Delta\lambda,\\
  \mathrm d\Delta_\Omega\lambda=\partial\Delta\lambda
  \qquad(\lambda\in N).
 \end{gathered}
\end{equation}
Here $\Delta'_a\lambda=\Delta_\Omega\lambda$ also uses
$\mathsf Q'\lambda=0$.

\subsection{The scalar inverse and the Landau projection}
\label{sec:regional-linear-scalar}

The nesting of blocks expresses a scalar average at one level
as the average of the preceding level's averages:
\[
 (\mathsf Q'_n\lambda)(y)
    =L^{-3}\sum_{z\in B(y)}(\mathsf Q'_{n-1}\lambda)(z)
       \quad(n\ge1).
\]
Indeed, $B^n(y)$ is the disjoint union of $B^{n-1}(z)$ over
$z\in B(y)$, and the averaging factors multiply.
Telescoping the differences along each straight path also gives
\begin{equation}\label{I-11-23}
 \begin{split}
 (\mathsf Q_n\partial\lambda)(c)
 &=L^{-4n}\xi^{-1}\sum_{x\in B^n(c_-)}
               \bigl(\lambda(x+\xi_ne_\mu)-\lambda(x)\bigr)\\
 &=\xi_n^{-1}
       \bigl((\mathsf Q'_n\lambda)(c_+)
                   -(\mathsf Q'_n\lambda)(c_-)\bigr).
 \end{split}
\end{equation}
For the second equality, translation by $\xi_ne_\mu$ maps
$B^n(c_-)$ onto $B^n(c_+)$, and $L^{-n}\xi^{-1}=\xi_n^{-1}$.
Finally, the inner products and weights give
\[
 \begin{split}
 \langle\lambda,\mathsf Q'^*\mathsf a'\mathsf Q'\lambda\rangle
     &=\sum_{n=0}^t a_n\xi_n
                    \sum_{y\in\Lambda_n}|(\mathsf Q'_n\lambda)(y)|^2,\\
 \langle A,\mathsf Q^*\mathsf a\mathsf QA\rangle
     &=\sum_{n=0}^t a\xi_n
                    \sum_{c\in\Lambda_n^{\mathrm b}}|(\mathsf Q_nA)(c)|^2.
 \end{split}
\]
In each term the product of the coarse volume $\xi_n^3$ and
the inverse-square weight is $a_n\xi_n$ or $a\xi_n$, respectively.

\begin{proof}[Proof of \cref{I-11-30}(a),(b)]
Fix a level $n$ and put $J=L^n$. The graph Laplacian of the path
with vertices $0,\ldots,J-1$ has eigenvectors
\[
 u_j^{(k)}=\cos\left(\frac{\pi k(j+\tfrac12)}{J}\right),
 \qquad 0\le k<J,
\]
with eigenvalues $4\sin^2(\pi k/(2J))$.
Indeed, the cosine identity gives
$u_{j+1}^{(k)}+u_{j-1}^{(k)}
 =2\cos(\pi k/J)u_j^{(k)}$, while
$u_{-1}^{(k)}=u_0^{(k)}$ and
$u_J^{(k)}=u_{J-1}^{(k)}$ give the same eigenvalue equation at
the two endpoints. The eigenvalues are distinct and
$u_0^{(k)}=\cos(\pi k/(2J))>0$, so these vectors form a basis.
Their threefold tensor products form an eigenbasis for the
nearest-neighbour box $\{0,\ldots,J-1\}^3$.
The zero eigenspace consists of constants; for $J>1$ the first
positive eigenvalue is
\[
 4\sin^2\left(\frac{\pi}{2J}\right)\ge4J^{-2},
 \qquad
 4J^2\sin^2\left(\frac{\pi}{2J}\right)<\pi^2.
\]
The inequalities follow from
$\sin s\ge2s/\pi$ on $[0,\pi/2]$ and $\sin s<s$ for $s>0$.

The coordinate description of $B^n(y)$ identifies it with
the translated box
$\{-(J-1)/2,\ldots,(J-1)/2\}^3$.
Internal box edges map to distinct positive bare bonds.
Write $b\subset B^n(y)$ for these bonds and
$\bar\lambda_y=(\mathsf Q'_n\lambda)(y)$.
Subtracting this mean makes the function orthogonal to constants
and leaves its differences unchanged. The box gap therefore gives
\[
 \sum_{b\subset B^n(y)}\xi^3|(\partial\lambda)(b)|^2
 \ge4\xi_n^{-2}
       \sum_{x\in B^n(y)}\xi^3|\lambda(x)-\bar\lambda_y|^2.
\]
For $n=0$ both sides vanish, since the block is a singleton.
For $n\ge1$ we may retain the exact eigenvalue, replacing the
coefficient $4$ by $4L^{2n}\sin^2(\pi/(2L^n))$.

By \cref{I-11-27}, the blocks partition $\mathsf S_0$.
Their internal bonds are distinct, so the sum of the left sides
is at most
$\|\partial\lambda\|^2
 =\langle\lambda,\Delta_\Omega\lambda\rangle$.
Since $|B^n(y)|\xi^3=L^{3n}\xi^3=\xi_n^3$,
\[
 \langle\lambda,\mathsf Q'^*\mathsf a'\mathsf Q'\lambda\rangle
  =\sum_{n=0}^t\sum_{y\in\Lambda_n}
       a_n\xi_n^{-2}
         \sum_{x\in B^n(y)}\xi^3|\bar\lambda_y|^2.
\]
On each block the mean and the mean-zero part are orthogonal:
\[
 \sum_{x\in B^n(y)}\xi^3|\lambda(x)|^2
 =\sum_{x\in B^n(y)}\xi^3|\lambda(x)-\bar\lambda_y|^2
       +\sum_{x\in B^n(y)}\xi^3|\bar\lambda_y|^2.
\]
Adding the two bounds proves the first inequality of
\eqref{I-11-31}. For $\lambda\in N$ all block means vanish,
and the block gaps with $n\ge1$ give the second inequality.
The exact block gap gives its stated refinement and the corresponding
refinement of the first inequality.

The operator $\Delta'_a$ is symmetric, and every coefficient
$\min\{a_n,4\}$ is positive. The block partition makes the first
lower bound strictly positive for nonzero $\lambda$.
Thus $\Delta'_a$ and its inverse $\mathsf G'$ are positive definite.
If $\lambda\in N$ and $\Delta_\Omega\lambda=0$, the second bound
makes $\lambda$ vanish on every block of positive level.
It also vanishes on $\Lambda_0$ since $\mathsf Q'_0=1$.
This proves injectivity on $N$.

For $\omega\in\mathcal K^0$, the definition of the coarse inner
product gives
\[
 \begin{split}
 \langle\omega,\mathsf Q'\lambda\rangle
 &=\sum_{n=0}^t\xi_n^3
      \sum_{y\in\Lambda_n}\omega(y)L^{-3n}
                              \sum_{x\in B^n(y)}\lambda(x)\\
 &=\xi^3\sum_{x\in\mathsf S_0}\omega(y(x))\lambda(x).
 \end{split}
\]
Hence $(\mathsf Q'^*\omega)(x)=\omega(y(x))$.
Averaging this block-constant function gives
$\mathsf Q'\mathsf Q'^*\omega=\omega$, proving surjectivity.
In particular, $\mathsf Q'^*$ is injective. For $\omega\ne0$,
\[
 \langle\omega,\mathsf Q'\mathsf G'\mathsf Q'^*\omega\rangle
 =\langle\mathsf Q'^*\omega,\mathsf G'\mathsf Q'^*\omega\rangle>0,
 \qquad
 \langle\omega,\mathsf Q'\mathsf G'^2\mathsf Q'^*\omega\rangle
 =\|\mathsf G'\mathsf Q'^*\omega\|^2>0.
\]
Both operators are symmetric, which proves (b).
\end{proof}

\begin{proof}[Proof of \cref{I-11-33}(i),(ii)]
The support identities were proved in
\eqref{eq:regional-null-support}. The injectivity just established,
together with the definition of $\mathsf R$, shows that
$\Delta_\Omega$ maps $N$ bijectively onto
$\operatorname{ran}\mathsf R$.

The projection calculation applies to any positive definite symmetric
operator $D$ on $\mathcal H^0(\Omega_0)$ satisfying
$D\lambda=\Delta_\Omega\lambda$ for $\lambda\in N$.
Put $S=D^{-1}\mathsf Q'^*$. Since $\mathsf Q'^*$ is injective,
$S^*S=\mathsf Q'D^{-2}\mathsf Q'^*$ is positive definite.
The operator $S(S^*S)^{-1}S^*$ is symmetric, is the identity
on $\operatorname{ran}S$, and vanishes on its orthogonal complement.
Moreover
\[
 (\operatorname{ran}S)^\perp
   =\ker S^*
   =\{f:\mathsf Q'D^{-1}f=0\}
   =DN=\Delta_\Omega N.
\]
Thus
\begin{equation}\label{I-11-36}
 \mathsf P
   =D^{-1}\mathsf Q'^*
       (\mathsf Q'D^{-2}\mathsf Q'^*)^{-1}\mathsf Q'D^{-1},
 \qquad
 \mathsf R=1-\mathsf P .
\end{equation}
For $D=\Delta'_a$ the hypotheses hold by
\cref{I-11-30}(a) and \eqref{eq:regional-null-support},
and this gives the formula in (ii) for every $a>0$.
The space $\Delta_\Omega N$ is independent of $a$, so its
orthogonal projections $\mathsf R,\mathsf P$ are also independent
of $a$.

If $\Omega_0\ne T_\ell$, one may also take $D=\Delta_\Omega$.
Indeed, a nonempty complementary cube contains its bare center,
so $\mathsf S_0\ne\mathsf S$. If
$\langle\lambda,\Delta_\Omega\lambda\rangle
 =\|\partial\lambda\|^2=0$, then $\lambda$ is constant on
the connected bare torus and vanishes outside $\mathsf S_0$;
hence $\lambda=0$. This proves the required positive definiteness.

With $D=\Delta'_a$, the preceding calculation identifies
$\operatorname{ran}(\mathsf G'\mathsf Q'^*)$ with
$\operatorname{ran}\mathsf P$. Therefore
$\mathsf R\mathsf G'\mathsf Q'^*=0$; taking adjoints gives
$\mathsf Q'\mathsf G'\mathsf R=0$.
Finally,
$\operatorname{ran}\mathsf R=\Delta_\Omega N=\Delta'_aN$,
so applying $\mathsf G'$ gives
$\mathsf G'\mathsf R\mathcal H^0(\Omega_0)=N$.
\end{proof}

\subsection{The curl kernel and the coarse averages}
\label{sec:regional-linear-kernel}

\begin{proof}[Proof of \cref{I-11-30}(c)--(e)]
Let $\lambda\in N$. By \eqref{eq:regional-null-support},
$\mathrm d\lambda=\partial\lambda$, so
$\partial\mathrm d\lambda=0$.
For $c\in\Lambda_n^{\mathrm b}$, telescoping gives
\[
 (\mathsf Q\mathrm d\lambda)(c)
 =\xi_n^{-1}\bigl(
    (\mathsf Q'_n\lambda)(c_+)
       -(\mathsf Q'_n\lambda)(c_-)\bigr).
\]
Both endpoint averages vanish. At an endpoint in $\Lambda_n$
this is the condition $\mathsf Q'\lambda=0$.
At an outer endpoint $y$ with $n\ge1$,
\[
 (\mathsf Q'_n\lambda)(y)
 =L^{-3}\sum_{z\in B(y)}
                      (\mathsf Q'_{n-1}\lambda)(z)=0,
\]
because $B(y)\subset\Lambda_{n-1}$ by \cref{I-11-27}(ii).
At a level-zero endpoint outside $\mathsf S_0$, the average
is the zero extension of $\lambda$.
Thus $\mathsf Q\mathrm d\lambda=0$.
If $\mathrm d\lambda=\partial\lambda=0$, then $\lambda$
is constant on the connected bare torus. Its mean on any
layer block is zero, and at least one such block exists
because $\mathsf S_0\ne\emptyset$. Hence $\lambda=0$.
This proves injectivity and one inclusion in the kernel formula.

\emph{An affine potential on the covering lattice.}
Suppose now that $A\in\mathcal H^1(\Omega_0)$ satisfies
$\partial A=0$ and $\mathsf QA=0$.
We first write its zero extension as
$A=\partial\mu+h$ on the full torus, with $h$ constant
in each component.
The ambient summation-by-parts identity
\cite[(1.21), p.~21]{B-CMP95} is
\begin{equation}\label{I-11-18}
 \|\partial B\|^2+\|\partial^*B\|^2
 =\sum_{\mu,\nu}\|\partial_\mu B_\nu\|^2
 \qquad(B\in\mathcal H^1).
\end{equation}
To check it with the present normalizations, translations
preserve the $\xi^3$-weighted sums, and therefore
\[
 \partial_\mu^*f(x)
       =\xi^{-1}\bigl(f(x-\xi e_\mu)-f(x)\bigr).
\]
All forward and backward difference operators commute.
Consequently
\[
 \langle\partial_\mu B_\nu,\partial_\nu B_\mu\rangle
 =\langle\partial_\nu^*B_\nu,\partial_\mu^*B_\mu\rangle.
\]
Expanding
$\frac12\sum_{\mu,\nu}
 \|\partial_\mu B_\nu-\partial_\nu B_\mu\|^2$
and summing this identity gives
$\|\partial B\|^2
 =\sum_{\mu,\nu}\|\partial_\mu B_\nu\|^2
       -\|\sum_\nu\partial_\nu^*B_\nu\|^2$,
which is \eqref{I-11-18}.

The symmetric nonnegative operator
$\Delta^{(1)}=\partial^*\partial+\partial\partial^*$
on ambient one-forms consequently has kernel exactly
the forms with constant components.
Let $h$ be the orthogonal projection of $A$ onto this
kernel. The vector $A-h$ belongs to the range of
$\Delta^{(1)}$, so $A-h=\Delta^{(1)}B$ for some $B$.
Since $\partial(A-h)=0$ and $\partial\partial=0$,
\[
 \partial\partial^*\partial B=0,\qquad
 \|\partial^*\partial B\|^2
   =\langle\partial B,\partial\partial^*\partial B\rangle=0.
\]
Thus $A=\partial\mu+h$, where $\mu=\partial^*B$.
Write $h_\nu=c_\nu$.

Let $\ell_t=\ell L^{K-t}$ be the side length in level-$t$
units, and let $\mathsf p:\mathbb R^3\to T_\ell$ be the
corresponding quotient map, with period group
$\mathsf V=\ell_t\mathbb Z^3$.
On the covering lattice $\xi\mathbb Z^3$, define
\[
 \widetilde\mu(\widetilde x)
   =\mu(\mathsf p\widetilde x)+c\cdot\widetilde x.
\]
Then
\[
 \partial\widetilde\mu=A\circ\mathsf p,\qquad
 \widetilde\mu(\widetilde x+v)
      =\widetilde\mu(\widetilde x)+c\cdot v
      \quad(v\in\mathsf V).
\]
The potential and its block means have the same affine increments.
Subtracting the means will give a periodic potential.

For $\widetilde y\in\xi_n\mathbb Z^3$, put
\[
 \begin{split}
 \widetilde B^n(\widetilde y)
   &=\left\{\widetilde x\in\xi\mathbb Z^3:
       |\widetilde x-\widetilde y|_\infty
                          \le\tfrac12(\xi_n-\xi)\right\},\\
 \widetilde\varphi_n(\widetilde y)
   &=L^{-3n}\sum_{\widetilde x\in\widetilde B^n(\widetilde y)}
                              \widetilde\mu(\widetilde x).
 \end{split}
\]
The covering map sends this lifted block bijectively onto
$B^n(\mathsf p\widetilde y)$.
At a fixed level these lifted blocks partition $\xi\mathbb Z^3$:
each integer has a unique representation
$L^ni+i'$ with $|i'|\le(L^n-1)/2$.
Two bare neighbours belong either to the same lifted block
or to blocks with nearest-neighbour centers. Indeed, if
two such integer representations differ by one, then
$L^n|i-j|\le1+|i'-j'|\le L^n$; the other coordinates
of the two sites agree.

If $\widetilde c$ is a lift of a positive level-$n$ bond $c$
of direction $\nu$, telescoping on its parallel bare paths gives
\begin{equation}\label{eq:regional-lifted-averages}
 (\mathsf Q_nA)(c)
  =\xi_n^{-1}\bigl(
       \widetilde\varphi_n(\widetilde c_+)
          -\widetilde\varphi_n(\widetilde c_-)\bigr).
\end{equation}
In fact, each path sum is
$\xi^{-1}(\widetilde\mu(\widetilde x+\xi_ne_\nu)
-\widetilde\mu(\widetilde x))$.
The factor $L^{-4n}$ in \eqref{I-11-20}, followed by
the block sum, leaves $L^{-n}\xi^{-1}=\xi_n^{-1}$.

\emph{Matching the means across levels.}
For $\widetilde x\in\mathsf p^{-1}(\mathsf S_0)$, let
$n=n(\mathsf p\widetilde x)$ and let
$\widetilde y(\widetilde x)\in\xi_n\mathbb Z^3$
be the center of its lifted level-$n$ block. Its image
is $y(\mathsf p\widetilde x)\in\Lambda_n$.
Define
\[
 \widetilde\varphi(\widetilde x)
       =\widetilde\varphi_n(\widetilde y(\widetilde x)).
\]
We show that this function has equal values at any two
bare neighbours $\widetilde x,\widetilde x'$ whose images
lie in $\mathsf S_0$.
Write $n=n(\mathsf p\widetilde x)$,
$n'=n(\mathsf p\widetilde x')$, and
$\widetilde y=\widetilde y(\widetilde x)$,
$\widetilde y'=\widetilde y(\widetilde x')$;
exchange the sites if necessary so that $n\ge n'$.

If $n=n'$, there is nothing to prove when
$\widetilde y=\widetilde y'$.
Otherwise the centers are nearest neighbours at level $n$,
and their images lie in $\Lambda_n$.
The bond between them, with positive orientation, belongs
to $\Lambda_n^{\mathrm b}$ by \cref{I-11-27}(ii).
The condition $\mathsf QA=0$ and
\eqref{eq:regional-lifted-averages} give equality of the means.

Suppose $n'<n$. By \cref{I-11-27}(i), $n'=n-1$.
Let $\widetilde y''$ be the center of the lifted level-$n$
block containing $\widetilde x'$.
The image of $\widetilde x'$ lies outside $\Omega_n$,
so the image of $\widetilde y''$ also lies outside $\Omega_n$.
In particular $\widetilde y''\ne\widetilde y$.
These centers are nearest neighbours, since their blocks
contain the neighbours $\widetilde x'$ and $\widetilde x$.
The level-$n$ bond between their images has one endpoint
in $\Lambda_n$ and the other outside $\Omega_n$;
it belongs to $\Lambda_n^{\mathrm b}$.
Therefore
\[
 \widetilde\varphi_n(\widetilde y'')
       =\widetilde\varphi_n(\widetilde y).
\]

By \cref{I-11-27}(ii), the outer block at
$\mathsf p\widetilde y''$ is a union of level-$(n-1)$
blocks with centers in $\Lambda_{n-1}$.
The corresponding lifted decomposition is
\[
 \widetilde B^n(\widetilde y'')
  =\bigsqcup_{\substack{
        \widetilde z=\widetilde y''+
                  \xi_{n-1}\sum_\nu i_\nu e_\nu\\
        |i_\nu|\le(L-1)/2}}
             \widetilde B^{n-1}(\widetilde z).
\]
To see the bounds in this decomposition directly, write
each coordinate $k$ of $\xi^{-1}(\widetilde x-\widetilde y'')$
as $L^{n-1}i+i'$ with $|i'|\le(L^{n-1}-1)/2$.
Then
\[
 \begin{array}{ll}
 |i|\le(L-1)/2
   &\Longrightarrow
     |k|\le\frac12L^{n-1}(L-1)+\frac12(L^{n-1}-1)
          =\frac12(L^n-1),\\[2pt]
 |i|\ge(L+1)/2
   &\Longrightarrow
     |k|\ge\frac12L^{n-1}(L+1)-\frac12(L^{n-1}-1)
          =\frac12(L^n+1).
 \end{array}
\]
Thus precisely these $L^3$ lifted centers occur.
Their images are the distinct sites of
$B(\mathsf p\widetilde y'')\subset\Lambda_{n-1}$.
The rectangular grid of centers is connected by
level-$(n-1)$ nearest-neighbour bonds.
Each such bond projects, after choosing positive orientation,
to a row in $\Lambda_{n-1}^{\mathrm b}$.
Equation \eqref{eq:regional-lifted-averages} and
$\mathsf QA=0$ therefore show that
$\widetilde\varphi_{n-1}$ is constant on these centers.
Its mean over them is $\widetilde\varphi_n(\widetilde y'')$.
The center $\widetilde y'$ is among them, because
$\widetilde x'$ belongs to its level-$(n-1)$ block
and to the displayed level-$n$ block. Hence
\[
 \widetilde\varphi(\widetilde x')
  =\widetilde\varphi_{n-1}(\widetilde y')
  =\widetilde\varphi_n(\widetilde y'')
  =\widetilde\varphi_n(\widetilde y)
  =\widetilde\varphi(\widetilde x).
\]
This proves equality along every bare bond with both
endpoints above $\mathsf S_0$.

\emph{The regional potential and its periods.}
Set
\[
 \widetilde\lambda(\widetilde x)=
 \begin{cases}
  \widetilde\mu(\widetilde x)
                -\widetilde\varphi(\widetilde x),
       &\mathsf p\widetilde x\in\mathsf S_0,\\
  0,&\mathsf p\widetilde x\notin\mathsf S_0.
 \end{cases}
\]
On a bond with both endpoints above $\mathsf S_0$,
the equality of neighbouring means gives
$\partial\widetilde\lambda=A\circ\mathsf p$.
On a bond with neither endpoint there, both sides vanish:
the image bond lies outside even the star frame.
For a boundary bond, its endpoint $x$ inside $\mathsf S_0$
belongs to $\mathsf S_0^\partial\subset\Lambda_0$.
At any lift $\widetilde x$ of this endpoint, the level-zero
block is a singleton, so
$\widetilde\varphi(\widetilde x)=\widetilde\mu(\widetilde x)$.
Thus $\widetilde\lambda$ vanishes at both endpoints.
The field $A$ also vanishes on this boundary bond.
In the inner frame this follows from zero extension.
In the star frame the bond belongs to
$\Lambda_0^{\mathrm b}$, since neither endpoint lies
in $\Omega_1$; there $\mathsf Q_0A=A$ and
$\mathsf QA=0$.
We have therefore proved
$\partial\widetilde\lambda=A\circ\mathsf p$ on every bare bond.

For $v\in\mathsf V$, the ratio $\ell_t/\xi_n$ is an integer
at every level. Translating a lifted level-$n$ block by $v$
translates its center by $v$, and the affine periodicity
of $\widetilde\mu$ gives
\[
 \widetilde\varphi_n(\widetilde y+v)
       =\widetilde\varphi_n(\widetilde y)+c\cdot v.
\]
The layer assignment is unchanged by a period, so
$\widetilde\varphi(\widetilde x+v)
=\widetilde\varphi(\widetilde x)+c\cdot v$
above $\mathsf S_0$.
These increments cancel those of $\widetilde\mu$.
Outside the invariant set $\mathsf p^{-1}(\mathsf S_0)$,
$\widetilde\lambda$ is zero. Hence it is periodic everywhere
and descends to $\lambda\in\mathcal H^0(\Omega_0)$
with $\partial\lambda=A$ and $\mathrm d\lambda=A$.
On a lifted layer block it is the difference between
$\widetilde\mu$ and the mean of $\widetilde\mu$
on that block. Its mean is zero.
The block projects bijectively, so $\mathsf Q'\lambda=0$.
This completes the reverse inclusion in (c).

\emph{Surjectivity of the coarse average.}
For $c\in\Lambda_n^{\mathrm b}$ let $E_c$ be the indicator
one-form of $\operatorname{cor}(c)$.
It belongs to $\mathcal H^1(\Omega_0)$:
for $n\ge1$ the corridor bonds are inner by
\cref{I-11-27}(iii), and for $n=0$ the corridor is $\{c\}$.
The path count in that lemma gives, for
$c'\in\Lambda_n^{\mathrm b}$,
\[
 (\mathsf Q_nE_c)(c')
    =L^{-4n}L^nL^{2n}\delta_{cc'}
    =L^{-n}\delta_{cc'}.
\]
If $n'<n$, choose the endpoint $y\in\Lambda_n$ of $c$.
Every corridor bond has an endpoint in
$B^n(y)\subset\Omega_n\subset\Omega_{n'+1}$.
In contrast, every bond in a path contributing to a
level-$n'$ row average has both endpoints outside
$\Omega_{n'+1}$, by \cref{I-11-27}(iii) when $n'\ge1$.
For $n'=0$, the same assertion is the defining condition
on $\Lambda_0^{\mathrm b}$.
Consequently $(\mathsf Q_{n'}E_c)(c')=0$ for
$c'\in\Lambda_{n'}^{\mathrm b}$.

Given a coarse field $B\in\mathcal K^1$, seek
$A=\sum_{n=0}^t\sum_{c\in\Lambda_n^{\mathrm b}}\beta_cE_c$.
The equation $\mathsf QA=B$ reads
\[
 L^{-n'}\beta_{c'}
 +\sum_{n<n'}\sum_{c\in\Lambda_n^{\mathrm b}}
       \beta_c(\mathsf Q_{n'}E_c)(c')
   =B(c'),
 \qquad c'\in\Lambda_{n'}^{\mathrm b}.
\]
Solve first at $n'=0$ and then at each increasing level.
At level $n'$, only coefficients from lower levels
appear beside the nonzero diagonal coefficient $L^{-n'}$.
This determines a regional field $A$ with
$\mathsf QA=B$, proving surjectivity.

\emph{Positivity of the vector operator.}
For $A\in\mathcal H^1(\Omega_0)$, the defining quadratic form is
\[
 \langle A,\Delta_aA\rangle
   =\|\partial A\|^2+\|\mathsf R\mathrm d^*A\|^2
          +\langle\mathsf QA,\mathsf a\mathsf QA\rangle
   \ge0.
\]
If it vanishes, then $\partial A=0$ and $\mathsf QA=0$,
since the bond weight $\mathsf a$ is positive.
The kernel formula gives $A=\mathrm d\lambda$ with
$\lambda\in N$.
By \eqref{eq:regional-null-support},
$\mathrm d^*\mathrm d\lambda=\Delta_\Omega\lambda$,
which belongs to the range $\Delta_\Omega N$ of $\mathsf R$.
The vanishing gauge term therefore gives
$\Delta_\Omega\lambda=0$.
The scalar bound \eqref{I-11-31} on $N$ implies
$\lambda=0$; when $t=0$, $N=\{0\}$ already.
Thus $\Delta_a$ and its inverse $\mathsf G$ are positive definite.
Finally, surjectivity of $\mathsf Q$ makes $\mathsf Q^*$
injective, and for $B\ne0$,
\[
 \langle B,\mathsf Q\mathsf G\mathsf Q^*B\rangle
       =\langle\mathsf Q^*B,\mathsf G\mathsf Q^*B\rangle>0.
\]
This proves the remaining inverse assertion in (e).
\end{proof}

\subsection{Gauge correction and constrained minimization}
\label{sec:regional-linear-minimizers}

\begin{proof}[Proof of \cref{I-11-33}(iii)--(v)]
For $\lambda\in N$, \eqref{eq:regional-null-support} gives
$\mathrm d^*\mathrm d\lambda=\Delta_\Omega\lambda
 \in\operatorname{ran}\mathsf R$. Consequently
\[
 \mathsf R\mathrm d^*(A-\mathrm d\lambda)
      =\mathsf R\mathrm d^*A-\Delta_\Omega\lambda .
\]
Since $\Delta_\Omega:N\to\operatorname{ran}\mathsf R$ is
bijective and $\mathsf G'\mathsf R$ maps into $N$, this expression
vanishes exactly for
$\lambda_0=\mathsf G'\mathsf R\mathrm d^*A$:
\[
 \Delta_\Omega\lambda_0
   =\Delta'_a\mathsf G'\mathsf R\mathrm d^*A
   =\mathsf R\mathrm d^*A .
\]
By \cref{I-11-30}(c), $\partial\mathrm d\lambda=0$ and
$\mathsf Q\mathrm d\lambda=0$ for every $\lambda\in N$.
Thus subtracting $\mathrm d\lambda_0$ preserves both the curl
and the coarse averages.

To compare full divergences, use
$\partial\lambda=\mathrm d\lambda$ and the support of
$\Delta\lambda$ inside $\mathsf S_0$.
The parts on $\mathsf S_0$ and its complement are orthogonal, so
\[
 \|\partial^*(A-\partial\lambda)\|^2
 =\|(1-\chi_0)\partial^*A\|^2
       +\|\mathrm d^*A-\Delta_\Omega\lambda\|^2 .
\]
Put $f=\mathrm d^*A$. Since
$\mathsf Rf=\Delta_\Omega\mathsf G'\mathsf Rf$,
the orthogonal decomposition $f=\mathsf Pf+\mathsf Rf$ gives
\[
 \|f-\Delta_\Omega\lambda\|^2
  =\langle f,\mathsf Pf\rangle
       +\|\Delta_\Omega(\lambda-\mathsf G'\mathsf Rf)\|^2.
\]
The first term and the contribution outside $\mathsf S_0$
are independent of $\lambda$. Injectivity of $\Delta_\Omega$
on $N$ proves the asserted uniqueness of the minimizer of
the ambient divergence norm.

For the inverse identities, take $\lambda\in N$. The curl and
average terms in $\Delta_a\mathrm d\lambda$ vanish, while
$\mathrm d^*\mathrm d\lambda=\Delta_\Omega\lambda$ lies in
$\operatorname{ran}\mathsf R$. It follows that
\[
 \Delta_a\mathrm d\lambda
  =\mathrm d\mathsf R\Delta_\Omega\lambda
  =\mathrm d\Delta_\Omega\lambda
  =\mathrm d\Delta'_a\lambda .
\]
Apply this to $\lambda=\mathsf G'\mathsf Rf$ and then apply
$\mathsf G$. This gives the first identity in \eqref{I-11-38};
the second is its adjoint. The same identities imply
\[
 \mathsf Q\mathsf G\mathrm d\mathsf R
      =\mathsf Q\mathrm d\mathsf G'\mathsf R=0,
 \qquad
 \mathrm d^*\mathsf G\mathrm d\mathsf R
      =\mathrm d^*\mathrm d\mathsf G'\mathsf R
      =\Delta'_a\mathsf G'\mathsf R=\mathsf R.
\]
Here $\mathsf G'\mathsf Rf\in N$ allows the use of both
the zero-average gradient identity and
\eqref{eq:regional-null-support}. Taking adjoints gives
$\mathsf R\mathrm d^*\mathsf G\mathsf Q^*=0$ and
$\mathsf R\mathrm d^*\mathsf G\mathrm d=\mathsf R$.

It remains to identify the minimizing lift.
Set
\[
 X=(\partial^*\partial)(\Omega_0)+\mathrm d\mathsf R\mathrm d^*.
\]
This operator is symmetric and nonnegative, with
$\langle A,XA\rangle
 =\|\partial A\|^2+\|\mathsf R\mathrm d^*A\|^2$.
On $\ker\mathsf Q$ it is positive definite, since there
$\langle Z,XZ\rangle=\langle Z,\Delta_aZ\rangle$.
The existence and positivity of
$\mathsf Q\mathsf G\mathsf Q^*$ have already been proved in
\cref{I-11-30}. The formula for $\mathsf H$ therefore gives
\[
 \mathsf Q\mathsf H=1,\qquad
 \Delta_a\mathsf H
      =\mathsf Q^*(\mathsf Q\mathsf G\mathsf Q^*)^{-1}.
\]
For $Z\in\ker\mathsf Q$ and $B\in\mathcal K^1$,
\[
 \begin{split}
 \langle Z,X\mathsf HB\rangle
 &=\langle Z,\Delta_a\mathsf HB\rangle
                  -\langle\mathsf QZ,\mathsf aB\rangle\\
 &=\langle\mathsf QZ,
               (\mathsf Q\mathsf G\mathsf Q^*)^{-1}B\rangle
                  -\langle\mathsf QZ,\mathsf aB\rangle=0.
 \end{split}
\]
Every $A$ with $\mathsf QA=B$ is uniquely of the form
$A=\mathsf HB+Z$ with $Z\in\ker\mathsf Q$, and hence
\[
 \langle A,XA\rangle
     =\langle\mathsf HB,X\mathsf HB\rangle+\langle Z,XZ\rangle.
\]
Thus $\mathsf HB$ is the unique minimizer of
$\langle A,XA\rangle$ on that affine space.
This characterization uses only $X$ and $\mathsf Q$.
Both are independent of $a$, because $\mathsf R$ is; consequently
$\mathsf H$ is independent of $a$.

The inverse identities give
\[
 \mathsf R\mathrm d^*\mathsf H
   =\mathsf R\mathrm d^*\mathsf G\mathsf Q^*
                         (\mathsf Q\mathsf G\mathsf Q^*)^{-1}=0.
\]
If also $\mathsf R\mathrm d^*A=0$, the preceding minimization
identity reads
$\|\partial A\|^2\ge\|\partial\mathsf HB\|^2$,
with equality only at $A=\mathsf HB$.
For an arbitrary $A$ with $\mathsf QA=B$, use the gauge
correction from (iii) and put $A'=A-\mathrm d\lambda_0$.
Then $\mathsf QA'=B$, $\mathsf R\mathrm d^*A'=0$ and
$\partial A'=\partial A$. Therefore the minimum curl energy
is attained exactly when $A'=\mathsf HB$, or equivalently
$A\in\mathsf HB+\mathrm dN$.
Conversely, every field in this affine space has average $B$
and the same curl as $\mathsf HB$.

Finally,
\[
 \mathsf H^*\Delta_a\mathsf H
   =(\mathsf Q\mathsf H)^*
                   (\mathsf Q\mathsf G\mathsf Q^*)^{-1}
   =(\mathsf Q\mathsf G\mathsf Q^*)^{-1}.
\]
Substituting
$\Delta_a=X+\mathsf Q^*\mathsf a\mathsf Q$,
using $\mathsf Q\mathsf H=1$ and
$\mathsf R\mathrm d^*\mathsf H=0$, gives
\eqref{I-11-39}.
\end{proof}

\section{Coercivity and decay of the inverse}
\label{sec:one-order-inverse}

The curl, the projected divergence and the straight block averages
control the full first-order norm of a one-form, with constants
independent of the fine spacing. The whole-torus form bound at $a=1$
is Ba{\l}aban's \cite[Proposition~1.1, (1.90), p.~33]{B-CMP95}.
We prove it in position space by estimating the divergence omitted
by the gauge projection. Exponential conjugation, suggested in
\cite[p.~36]{B-CMP95}, then gives decay for the inverse of each
compression to a bond set.

In this section let $L\ge2$ and $k\ge0$ be integers, put
$N=L^k$ and $\eta=N^{-1}$, and fix $a>0$.
We work in dimension three on the fine torus and its coarse lattice
\[
 \mathbb L_\eta
   =\eta\mathbb Z^3/
        (\ell_1\mathbb Z\times\ell_2\mathbb Z\times\ell_3\mathbb Z),
 \qquad
 \mathbb L_1
   =\mathbb Z^3/
        (\ell_1\mathbb Z\times\ell_2\mathbb Z\times\ell_3\mathbb Z),
 \qquad \ell_\mu\in\mathbb N,\quad \ell_\mu\ge1.
\]
The distance $|x-x'|_1$ is the least $\ell^1$ norm of a lift of
$x-x'$, and $\operatorname{dist}(U,U')$ is the induced distance
between nonempty sets.

All fields are real. Site functions, tensors
$F_{\mu_1\cdots\mu_r}(x)$, one-forms $A_\mu(x)$ and
two-forms $F_{\mu\nu}(x)$, $\mu<\nu$, have the inner products
$\langle F,F'\rangle=\eta^3\sum FF'$.
A positive bond is indexed by $(x,\mu)$, including when
$x+\eta e_\mu=x$ on the torus; the direction remains part of its
index. For a tensor, multiplication by a site function, and the
indicator $\chi_U$ of a site set, act at the base point of each
component. A bond-set indicator $\chi_S$ acts on the indexed
one-form components.

Use the forward differences
\[
 \partial_\mu=\eta^{-1}(\tau_\mu-1),\qquad
 (\tau_\mu F)(x)=F(x+\eta e_\mu),\qquad
 \partial_\mu^*=\eta^{-1}(\tau_\mu^{-1}-1).
\]
The translations are commuting isometries, so these are the
adjoint differences for the stated inner products.
The exterior derivative is
$(\partial f)_\mu=\partial_\mu f$ on scalars and
$(\partial A)_{\mu\nu}=\partial_\mu A_\nu-\partial_\nu A_\mu$
on one-forms. In particular,
$\partial^*A=\sum_\mu\partial_\mu^*A_\mu$.
The componentwise gradient and Laplacian are
\[
 (\nabla F)_{\mu\mu_1\cdots\mu_r}
      =\partial_\mu F_{\mu_1\cdots\mu_r},
 \qquad
 \Delta=\nabla^*\nabla=\sum_{\mu=1}^3\partial_\mu^*\partial_\mu .
\]
On scalars, $\nabla=\partial$.
Coarse fields have the same conventions with spacing and volume
weight one; their differences are denoted by $\partial_1$.

Identify torus fields with periodic fields on the covering lattice.
For $y\in\mathbb Z^3$ put
\[
 B(y)=y+\eta\{0,\ldots,N-1\}^3.
\]
Its image is the corner block at the class of $y$ in $\mathbb L_1$.
Every sum over $B(y)$ or over neighboring blocks below is a sum
on the covering lattice. When summing over $y$, take one coarse
period. Thus repeated images of points in two neighboring blocks
retain their multiplicities.
Define
\begin{equation}\label{I-11-63}
 \begin{split}
 (\mathsf Q'f)(y)
   &=\sum_{x\in B(y)}\eta^3 f(x),\\
 (\mathsf QA)_\mu(y)
   &=\sum_{x\in B(y)}\eta^3
                \sum_{j=0}^{N-1}\eta\,A_\mu(x+j\eta e_\mu).
 \end{split}
\end{equation}
Their adjoints use the fine weight $\eta^3$ and coarse weight one.
Let $\mathsf R$ be the orthogonal projection of the full scalar
space onto $\Delta\ker\mathsf Q'$, and put $\mathsf P=1-\mathsf R$.
On one-forms define
\begin{equation}\label{I-11-64}
 \Delta_a=\partial^*\partial+\partial\mathsf R\partial^*
                       +a\mathsf Q^*\mathsf Q .
\end{equation}

For a set $S$ of positive bonds, let $\mathcal H_S$ be the space
of one-forms vanishing off $S$ and set
$\Delta_a(S)=\chi_S\Delta_a\chi_S|_{\mathcal H_S}$.
The projection $\mathsf R$ in this operator is formed on the
whole torus before compression. The inverse, once its existence
is established, is denoted by $G(S)$ and extended by zero on
$\mathcal H_S^\perp$. Every derivative of $G(S)J$ is the ambient
derivative of this zero extension.

The constants used in the estimates are
\begin{equation}\label{I-12-13}
 C_3=\frac{30375}{16},\qquad \gamma'=(C_3+1)^{-2},
\end{equation}
\begin{equation}\label{I-12-18}
 \begin{gathered}
 \mathsf b_{\mathrm B}=C_3+2,\qquad
 C_h=\sqrt{\mathsf b_{\mathrm B}}+2\sqrt3\,\mathsf b_{\mathrm B},\\
 C'_{\mathrm B}=\max\{1+24C_h^2,\ 2C_h^2a^{-1}\},\qquad
 \gamma_{0,\mathrm B}=(25C'_{\mathrm B}+2a^{-1})^{-1},
 \end{gathered}
\end{equation}
\begin{equation}\label{I-12-27}
 c_P=548208\,\gamma'^{-2},\qquad c_V=c_P+6+3a,
\end{equation}
and
\begin{equation}\label{I-12-30}
 \delta_0=\frac{\gamma_{0,\mathrm B}}{16c_V}.
\end{equation}

\begin{proposition}\label{I-12-17}
Every one-form $A$ on $\mathbb L_\eta$ satisfies
\begin{equation}\label{I-12-19}
 \begin{split}
 \|\mathsf P\partial^*A\|
    &\le C_h\bigl(\|\mathsf QA\|+2\sqrt3\,\|\partial A\|\bigr),\\
 \langle A,\Delta_aA\rangle
    &\ge\gamma_{0,\mathrm B}
                       \bigl(\|\nabla A\|^2+\|A\|^2\bigr).
 \end{split}
\end{equation}
The constants are positive and depend only on $a$; in particular,
$\gamma_{0,\mathrm B}\le1/25$.
\end{proposition}

\pagebreak[3]
\begin{proposition}\label{I-12-29}
For every set $S$ of positive bonds the following hold.
\begin{enumerate}[label=(\alph*)]
\item The operator $\Delta_a(S)$ is symmetric and
\[
 \langle A,\Delta_a(S)A\rangle
       \ge\gamma_{0,\mathrm B}
                    \bigl(\|\nabla A\|^2+\|A\|^2\bigr)
            \qquad(A\in\mathcal H_S).
\]
It is invertible on $\mathcal H_S$, and its inverse is symmetric
and positive definite on that space.

\item With the zero extension specified above,
\[
 \|G(S)\|,\quad \|\nabla G(S)\|,\quad
 \|G(S)\nabla^*\|,\quad\|\nabla G(S)\nabla^*\|
                 \le\gamma_{0,\mathrm B}^{-1}.
\]
The first two operators have one-form inputs; the last two have
rank-two tensor inputs. Their outputs are one-forms or rank-two
tensors according to whether the left factor is $1$ or $\nabla$.

\item For $\mathcal O_1\in\{1,\nabla\}$,
$\mathcal O_2\in\{1,\nabla^*\}$ and arbitrary nonempty sets
$U,U'$ of fine sites,
\begin{equation}\label{I-12-32}
 \bigl\|\chi_U\mathcal O_1G(S)\mathcal O_2\chi_{U'}\bigr\|
       \le4\gamma_{0,\mathrm B}^{-1}
                     e^{-\delta_0\operatorname{dist}(U,U')}.
\end{equation}
The operator domains and targets are those in (b).
\end{enumerate}
The constants are independent of $k$, $L$, $S$ and the periods
$\ell_1,\ell_2,\ell_3$.
\end{proposition}

The proof of \cref{I-12-17} is in
\cref{sec:one-order-coercivity}. The averaging estimates in
\cref{sec:one-order-averages} compare a face average with the
straight average; the scalar estimate in
\cref{sec:one-order-scalar} controls the inverse appearing in
the projection $\mathsf P$. Together they bound
$\mathsf P\partial^*A$ by the curl and the block averages.
For \cref{I-12-29}, the conjugation estimate in
\cref{sec:one-order-conjugation} retains half the form lower
bound, uniformly in the fine spacing.
\Cref{sec:one-order-decay} then applies it on $\mathcal H_S$
and chooses a weight determined by the distance to $U'$.

\subsection{Straight averages and face averages}
\label{sec:one-order-averages}

The corner blocks partition the fine torus. To see this, choose
a lift $\tilde x\in\eta\mathbb Z^3$ of a site $x$ and take
the unique integer vector $\tilde y$ with
$\tilde y_\mu\le\tilde x_\mu<\tilde y_\mu+1$.
Changing the lift by a period changes $\tilde y$ by the same
integer period. Thus the class of $\tilde y$ is well defined.
The $N^3$ points $\tilde y+\eta n$, $0\le n_\mu<N$, have
distinct images: their coordinate differences have modulus
less than one, while a nonzero period has a coordinate of
modulus at least one. Each block therefore has volume
$N^3\eta^3=1$. Taking the adjoint of its mean gives
\[
 (\mathsf Q'^*\omega)(x)=\omega(y)
                      \qquad(x\in B(y)).
\]

\begin{lemma}\label{I-12-6}
On the covering lattice put
\[
 D_\mu(y)=B(y)\cup B(y+e_\mu),\qquad
 \mathrm w(s)=\min\{s+\eta,\ 2-\eta-s\}.
\]
Then
\begin{equation}\label{I-12-7}
 (\mathsf QA)_\mu(y)
    =\sum_{x'\in D_\mu(y)}\eta^3
                    \mathrm w(x'_\mu-y_\mu)A_\mu(x').
\end{equation}
For $s\in\{0,\eta,\ldots,2-\eta\}$ the weights lie in $[0,1]$,
and
\[
 \sum_{x'\in D_\mu(y)}\eta^3\mathrm w(x'_\mu-y_\mu)=1,\qquad
 \mathrm w(s)+\mathrm w(s+1)=1
                \quad(0\le s\le1-\eta,\ s\in\eta\mathbb Z).
\]
For a coarse one-form $H$ and $x'\in B(y)$, with
$s=x'_\mu-y_\mu$ computed in the chosen block lift,
\[
 (\mathsf Q^*H)_\mu(x')
      =\mathrm w(s)H_\mu(y)+\mathrm w(s+1)H_\mu(y-e_\mu).
\]
\end{lemma}

\begin{proof}
In \eqref{I-11-63}, collect the pairs $(x,j)$ giving the
same point $x'=x+j\eta e_\mu$ on the cover.
Write $x_\mu-y_\mu=m_0\eta$ and $x'_\mu-y_\mu=m\eta$.
The pairs satisfy $m_0+j=m$, with $0\le m_0,j<N$.
Their number is
$\min\{m+1,2N-1-m\}=\eta^{-1}\mathrm w(m\eta)$,
for $0\le m\le2N-1$; at the last value it is zero.
Each pair has weight $\eta^4$, proving \eqref{I-12-7}.
There are $N^4$ pairs in total, so the weights sum to
$\eta^4N^4=1$. Their nonnegativity and upper bound follow
from the count. For $0\le s\le1-\eta$,
$\mathrm w(s)=s+\eta$ and
$\mathrm w(s+1)=1-\eta-s$, proving the two-weight identity.

To compute the adjoint, sum the product of \eqref{I-12-7}
with $H_\mu(y)$ over one coarse period.
The summand indexed by $(y,x')$ is invariant when both points
are translated by the same period.
Each orbit of pairs has exactly one representative with $y$
in a chosen coarse period, and exactly one with $x'$ in a
chosen fine period. We may therefore exchange the two sums,
retaining multiplicities. If $x'\in B(y_0)$ on the cover,
then $x'\in D_\mu(y)$ precisely for
$y=y_0$ and $y=y_0-e_\mu$.
Their weights are $\mathrm w(s)$ and $\mathrm w(s+1)$.
This proves the adjoint formula even when those two coarse
sites have the same torus image.
\end{proof}

For the proof of coercivity we also use the outgoing-face average
\[
 (\widetilde{\mathsf Q}A)_\mu(y)
   =\sum_{\substack{x\in B(y)\\x_\mu=y_\mu+1-\eta}}
                      \eta^2 A_\mu(x).
\]
The equality in the summation condition is imposed on the cover.

\begin{lemma}\label{I-12-8}
For a scalar field $f$ and a one-form $A$ the following hold.
\begin{enumerate}[label=(\roman*)]
\item $\mathsf Q'\mathsf Q'^*=1$,
$\|\mathsf Q'\|=\|\mathsf Q'^*\|=1$ and $\|\mathsf Q\|\le1$.
The operator $\Pi'=\mathsf Q'^*\mathsf Q'$ is the orthogonal
projection onto the block-constant functions.
\item
$\mathsf Q\partial=\partial_1\mathsf Q'$ and
$\mathsf Q'\partial^*=\partial_1^*\widetilde{\mathsf Q}$.
\item $\|f-\Pi'f\|^2\le\tfrac32\|\partial f\|^2$.
\item
$\|(\mathsf Q-\widetilde{\mathsf Q})A\|^2
                        \le12\|\nabla A\|^2$.
\item $\|A\|^2\le24\|\nabla A\|^2+2\|\mathsf QA\|^2$.
\end{enumerate}
\end{lemma}

\begin{proof}
The formula for $\mathsf Q'^*$ and the unit volume of a block
give $\mathsf Q'\mathsf Q'^*=1$ and
$\|\mathsf Q'^*\omega\|=\|\omega\|$.
Thus $\mathsf Q'^*$ is an isometry and its adjoint has norm one.
The symmetric operator $\Pi'$ is idempotent and has range
$\operatorname{ran}\mathsf Q'^*$, the block-constant functions.
For the bond average, Cauchy--Schwarz and \cref{I-12-6} give
\[
 |(\mathsf QA)_\mu(y)|^2
   \le\sum_{x'\in D_\mu(y)}
             \eta^3\mathrm w(x'_\mu-y_\mu)|A_\mu(x')|^2.
\]
Sum over $y$ and exchange covering pairs as in that lemma.
Each point receives the two weights
$\mathrm w(s)$ and $\mathrm w(s+1)$, whose sum is one.
Summing also over $\mu$ proves $\|\mathsf QA\|\le\|A\|$.

For (ii), the straight path telescopes:
\[
 \begin{split}
 (\mathsf Q\partial f)_\mu(y)
 &=\sum_{x\in B(y)}\eta^3
             \bigl(f(x+e_\mu)-f(x)\bigr)\\
 &=(\mathsf Q'f)(y+e_\mu)-(\mathsf Q'f)(y).
 \end{split}
\]
For the divergence, sum its backward difference over a block:
\[
 (\mathsf Q'\partial^*A)(y)
   =\sum_\mu\sum_{x\in B(y)}
                 \eta^2\bigl(A_\mu(x-\eta e_\mu)-A_\mu(x)\bigr).
\]
For fixed transverse coordinates the inner sum in direction
$\mu$ telescopes to the face at $x_\mu=y_\mu-\eta$
minus the face at $x_\mu=y_\mu+1-\eta$.
These are the outgoing faces of $B(y-e_\mu)$ and $B(y)$,
respectively. The result is
$\sum_\mu((\widetilde{\mathsf Q}A)_\mu(y-e_\mu)
              -(\widetilde{\mathsf Q}A)_\mu(y))
 =(\partial_1^*\widetilde{\mathsf Q}A)(y)$.

We prove the remaining estimates by one path calculation on
the covering lattice. Let
\[
 \mathcal B=\prod_{\nu=1}^3 I_\nu,\qquad
 I_\nu=\{\alpha_\nu,\alpha_\nu+\eta,\ldots,
                            \alpha_\nu+\lambda_\nu-\eta\},
 \qquad \lambda_\nu\in\{1,2\},
\]
and let $I_\nu^-$ omit the last point of $I_\nu$.
Let $v$ be a scalar function on $\mathcal B$.
For $x,x'\in\mathcal B$ and a permutation $\sigma$ of
$\{1,2,3\}$, replace the coordinates of $x$ by those of $x'$
in the order $\sigma(1),\sigma(2),\sigma(3)$.
On the $i$th segment, write $z^{(i)}(s)$ for the point whose
coordinates $\sigma(j)$ are $x'_{\sigma(j)}$ for $j<i$,
$s$ for $j=i$, and $x_{\sigma(j)}$ for $j>i$.
The increment along this segment is at most
$\sum_{s\in I_{\sigma(i)}^-}
          \eta|\partial_{\sigma(i)}v(z^{(i)}(s))|$.
Since $\sum_{s\in I_\nu^-}\eta\le\lambda_\nu$,
two applications of Cauchy--Schwarz give
\begin{equation}\label{I-12-9}
 |v(x)-v(x')|^2
  \le3\sum_{i=1}^3\lambda_{\sigma(i)}
          \sum_{s\in I_{\sigma(i)}^-}
                   \eta|\partial_{\sigma(i)}v(z^{(i)}(s))|^2.
\end{equation}
All the differences are internal to the covering box. Set
\[
 \|\nabla v\|_{\mathcal B}^2
   =\sum_\nu\sum_{\substack{z\in\mathcal B\\z_\nu\in I_\nu^-}}
                      \eta^3|\partial_\nu v(z)|^2.
\]

Let $p(x)=\prod_\nu p_\nu(x_\nu)$ and
$p'(x')=\prod_\nu p'_\nu(x'_\nu)$ be products of probability
weights on the intervals, and suppose
\begin{equation}\label{I-12-10}
 p'_{\sigma(j)}\le\eta\quad(j<3),\qquad
 p_{\sigma(j)}\le\eta\quad(j>1).
\end{equation}
Multiply \eqref{I-12-9} by $p(x)p'(x')$ and sum.
In its $i$th term the sums over $x_{\sigma(j)}$ for $j\le i$
and over $x'_{\sigma(j)}$ for $j\ge i$ give one.
The two weights left over are among
$p'_{\sigma(j)}$ with $j<i$ and $p_{\sigma(j)}$ with $j>i$,
so each is at most $\eta$.
Together with the segment factor $\eta$ they contribute at most
$\eta^3$. The remaining variables enumerate each internal bond
of direction $\sigma(i)$ once. Hence
\begin{equation}\label{I-12-11}
 \sum_{x,x'\in\mathcal B}p(x)p'(x')|v(x)-v(x')|^2
       \le3\max_\nu\lambda_\nu\,\|\nabla v\|_{\mathcal B}^2.
\end{equation}

For (iii), take $\mathcal B=B(y)$, $v=f$ and
$p_\nu=p'_\nu=\eta$. All interval lengths are one.
The left side of \eqref{I-12-11} is
\[
 \sum_{x,x'\in B(y)}\eta^6|f(x)-f(x')|^2
       =2\sum_{x\in B(y)}\eta^3
                       |f(x)-(\mathsf Q'f)(y)|^2.
\]
Summing the resulting bound over disjoint blocks gives (iii),
because their internal bonds form a subset of all fine bonds.

For (iv), fix $y,\mu$, take
$\mathcal B=D_\mu(y)$, $v=A_\mu$ and choose $\sigma(1)=\mu$.
Let $p_\mu$ be a point mass at $y_\mu+1-\eta$,
let $p_\nu=\eta$ for $\nu\ne\mu$, and set
\[
 p'_\mu(s)=\eta\,\mathrm w(s-y_\mu),\qquad
 p'_\nu=\eta\quad(\nu\ne\mu).
\]
The weights $p'$ have total mass one and each $p'_\nu\le\eta$,
by \cref{I-12-6}. The only possibly larger weight is $p_\mu$,
which is exempt from \eqref{I-12-10} because $\sigma(1)=\mu$.
The two weighted means of $v$ are
$(\widetilde{\mathsf Q}A)_\mu(y)$ and $(\mathsf QA)_\mu(y)$.
Their squared difference is at most the left side of
\eqref{I-12-11}, by Cauchy--Schwarz. Since the box has side
two in direction $\mu$ and side one in the others,
\[
 |((\mathsf Q-\widetilde{\mathsf Q})A)_\mu(y)|^2
                      \le6\|\nabla A_\mu\|_{D_\mu(y)}^2.
\]
Sum over $y$ using covering bonds and their base points.
An internal bond of a block $B(y_0)$ belongs to
$D_\mu(y_0)$ and $D_\mu(y_0-e_\mu)$.
A bond crossing a block face in direction $\mu$ belongs
only to $D_\mu(y_0)$, and a bond crossing a face in another
direction belongs to none of these boxes.
Thus each bond is counted at most twice, including with
multiplicity when period images coincide.
Summing over $\mu$ proves (iv).

For (v), keep $v,\sigma,p'$ as above and replace $p$ by
the uniform weight $\eta^3$ on $B(y)$, zero on the rest
of $D_\mu(y)$. It again satisfies \eqref{I-12-10}.
Cauchy--Schwarz applied to the $p'$ mean, followed by
\eqref{I-12-11}, yields
\[
 \sum_{x\in B(y)}\eta^3
          |A_\mu(x)-(\mathsf QA)_\mu(y)|^2
                       \le6\|\nabla A_\mu\|_{D_\mu(y)}^2.
\]
The block has volume one, so
\[
 \|A_\mu\|_{B(y)}^2
     \le12\|\nabla A_\mu\|_{D_\mu(y)}^2
                         +2|(\mathsf QA)_\mu(y)|^2.
\]
The same multiplicity bound gives (v) after summing over
$y$ and $\mu$.
\end{proof}

\subsection{The scalar inverse and its coarse averages}
\label{sec:one-order-scalar}

For a scalar function or a one-form $F$, write
\[
 \|F\|_1^2=\|\nabla F\|^2+\|F\|^2.
\]
The scalar operator used to represent $\mathsf P$ has
unit averaging coefficient:
\[
 \Delta'=\Delta+\mathsf Q'^*\mathsf Q'.
\]
The projection formula requires a uniform lower bound
for the coarse Gram operator defined below.

\begin{lemma}\label{I-12-12}
The scalar operator $\Delta'$ satisfies
\[
 \langle f,\Delta'f\rangle\ge\tfrac13\|f\|_1^2.
\]
Its inverse $\mathsf G'=(\Delta')^{-1}$ exists and obeys
$\|\mathsf G'g\|_1\le3\|g\|$.
On coarse scalar functions, the operator
\[
 \mathsf M=\mathsf Q'\mathsf G'^2\mathsf Q'^*
\]
satisfies $\gamma'\le\mathsf M\le9$.
Moreover,
\begin{equation}\label{I-12-14}
 \mathsf P
 =\mathsf G'\mathsf Q'^*\mathsf M^{-1}\mathsf Q'\mathsf G'.
\end{equation}
\end{lemma}

\begin{proof}
Let $\Pi'=\mathsf Q'^*\mathsf Q'$ be the orthogonal
projection onto the block-constant functions.
By \cref{I-12-8}(i),(iii),
\[
 \begin{split}
 \langle f,\Delta'f\rangle
 &=\|\partial f\|^2+\|\Pi'f\|^2\\
 &\ge\tfrac12\|\partial f\|^2
          +\tfrac13\|f-\Pi'f\|^2+\|\Pi'f\|^2
 \ge\tfrac13\|f\|_1^2.
 \end{split}
\]
Thus $\Delta'$ is positive definite.
For $u=\mathsf G'g$,
\[
 \tfrac13\|u\|_1^2
 \le\langle u,\Delta'u\rangle
 =\langle u,g\rangle\le\|u\|_1\|g\|,
\]
which proves the inverse bound. The isometry of
$\mathsf Q'^*$ then gives
$\langle\omega,\mathsf M\omega\rangle
=\|\mathsf G'\mathsf Q'^*\omega\|^2\le9\|\omega\|^2$.

\emph{A normalized block profile.}
Define
\[
 \mathsf g(t)=
 \begin{cases}
  t^2(1-t)^2,&0\le t\le1,\\
  0,&t\notin[0,1].
 \end{cases}
\]
This function is continuously differentiable,
$0\le\mathsf g\le1/16$, and its derivative is Lipschitz
with constant $2$, since
$|\mathsf g''(t)|=|2-12t+12t^2|\le2$ on $[0,1]$.
Consequently
\begin{equation}\label{I-12-15}
 \begin{gathered}
 \bigl|\mathsf g(t+\eta)+\mathsf g(t-\eta)-2\mathsf g(t)\bigr|
 \le\int_0^\eta2(2s)\,ds=2\eta^2,\\
 m_N:=\frac1N\sum_{j=0}^{N-1}
             \mathsf g\left(\frac{j+\frac12}{N}\right)
       =\frac1{30}+\frac7{240N^4}\ge\frac1{30}.
 \end{gathered}
\end{equation}
The identity for $m_N$ follows by expanding
$\mathsf g(t)=t^2-2t^3+t^4$ and summing the powers of
$j+\frac12$; it holds also for $N=1$.

For a lifted coarse site $y\in\mathbb Z^3$ and
$x\in\eta\mathbb Z^3$, put
\[
 \psi_y(x)=m_N^{-3}
              \prod_{\mu=1}^3\mathsf g(x_\mu-y_\mu+\eta/2).
\]
The half-mesh displacement makes $\psi_y$ vanish at
every fine site outside the covering block $B(y)$.
On that block its average is
\[
 \sum_{x\in B(y)}\eta^3\psi_y(x)
 =m_N^{-3}\prod_{\mu=1}^3
      \left(\eta\sum_{j=0}^{N-1}
                   \mathsf g((j+\tfrac12)\eta)\right)=1.
\]
Regard a coarse function $\omega$ as a periodic function
on $\mathbb Z^3$, and define
\[
 \psi(x)=\sum_{y\in\mathbb Z^3}\omega(y)\psi_y(x).
\]
At any covering-lattice site, at most one summand is nonzero.
For each period $v$,
$\psi_{y+v}(x+v)=\psi_y(x)$, so $\psi$ is periodic
and defines a fine-torus function with
$\mathsf Q'\psi=\omega$.

\emph{Second differences and the coarse lower bound.}
Fix a direction $\mu$ and $x\in B(y_0)$, and set
\[
 t=x_\mu-y_{0,\mu}+\eta/2,\qquad
 c(x)=m_N^{-3}\prod_{\nu\ne\mu}
                \mathsf g(x_\nu-y_{0,\nu}+\eta/2)
       \le30^3\,16^{-2}.
\]
Along this coordinate line,
$\psi=c(x)\sum_{n\in\mathbb Z}
\omega(y_0+ne_\mu)\mathsf g(t-n)$.
A second difference of $\mathsf g(t-n)$ can be nonzero
only if $-\eta<t-n<1+\eta$.
For $t\in\{\eta/2,3\eta/2,\ldots,1-\eta/2\}$,
this allows $n=0$, $n=1$ at $t=1-\eta/2$,
and $n=-1$ at $t=\eta/2$. Hence \eqref{I-12-15} gives
\[
 \begin{split}
 |(\partial_\mu^*\partial_\mu\psi)(x)|
 \le2c(x)\bigl(&|\omega(y_0)|
  +\mathbf1_{\{t=1-\eta/2\}}|\omega(y_0+e_\mu)|\\
  &+\mathbf1_{\{t=\eta/2\}}|\omega(y_0-e_\mu)|\bigr).
 \end{split}
\]
Each of the three block-constant functions on the right,
with the indicators omitted, has norm $\|\omega\|$.
This follows from the isometry of $\mathsf Q'^*$ and
invariance of the coarse sum under translation.
It remains true when neighbouring coarse labels coincide
modulo a period; the estimate uses their three contributions.
The triangle inequality therefore yields
\[
 \|\partial_\mu^*\partial_\mu\psi\|
      \le6\,30^3\,16^{-2}\|\omega\|,\qquad
 \|\Delta\psi\|
      \le18\,30^3\,16^{-2}\|\omega\|
       =C_3\|\omega\|.
\]
Since $\Delta'\psi=\Delta\psi+\mathsf Q'^*\omega$,
we obtain $\|\Delta'\psi\|\le(C_3+1)\|\omega\|$.
The symmetry of $\Delta'$ now gives
\[
 \begin{split}
 \|\omega\|^2
 &=\langle\psi,\mathsf Q'^*\omega\rangle
   =\langle\Delta'\psi,\mathsf G'\mathsf Q'^*\omega\rangle\\
 &\le(C_3+1)\|\omega\|\,
                       \|\mathsf G'\mathsf Q'^*\omega\|.
 \end{split}
\]
Thus
$\langle\omega,\mathsf M\omega\rangle
=\|\mathsf G'\mathsf Q'^*\omega\|^2
\ge(C_3+1)^{-2}\|\omega\|^2=\gamma'\|\omega\|^2$.

Finally, set $\mathcal T=\mathsf G'\mathsf Q'^*$.
The scalar operator $\Delta'$ is positive,
$\mathsf Q'$ is onto, and
$\Delta'=\Delta$ on $\ker\mathsf Q'$.
These verify the hypotheses of the range calculation
leading to \eqref{I-11-36}. In the present spaces it reads
\[
 (\operatorname{ran}\mathcal T)^\perp
   =\ker(\mathsf Q'\mathsf G')
   =\Delta'\ker\mathsf Q'
   =\Delta\ker\mathsf Q'=\operatorname{ran}\mathsf R.
\]
Since $\mathcal T^*\mathcal T=\mathsf M>0$,
the orthogonal projection onto
$\operatorname{ran}\mathcal T$ is
$\mathcal T\mathsf M^{-1}\mathcal T^*$.
It equals $1-\mathsf R=\mathsf P$, proving
\eqref{I-12-14}.
\end{proof}

\subsection{Recovering the omitted divergence}
\label{sec:one-order-coercivity}

\begin{proof}[Proof of \cref{I-12-17}]
The periodic difference operators commute, so the
summation-by-parts identity \eqref{I-11-18} holds with
the present spacing and periods:
\[
 \|\nabla A\|^2=\|\partial A\|^2+\|\partial^*A\|^2,
 \qquad \partial\partial f=0.
\]

The kernel of the scalar Laplacian consists of constants:
if $\Delta f=0$, then
$\|\partial f\|^2=\langle f,\Delta f\rangle=0$,
and a path of bonds joins any two sites of the torus.
Since $\langle\partial^*A,1\rangle=0$, there is a unique
mean-zero function $\lambda_{\mathrm c}$ satisfying
$\Delta\lambda_{\mathrm c}=\partial^*A$.
Indeed, the range of the finite-dimensional symmetric
operator $\Delta$ is the orthogonal complement of its
kernel, and its restriction to that complement is
injective. The transverse part
$A^{\mathrm T}=A-\partial\lambda_{\mathrm c}$ satisfies
\[
 \partial^*A^{\mathrm T}=0,\qquad
 \partial A^{\mathrm T}=\partial A,\qquad
 \|\nabla A^{\mathrm T}\|=\|\partial A\|.
\]

Put $\varpi=\mathsf Q'\lambda_{\mathrm c}$.
The two identities in \cref{I-12-8}(ii) yield
\[
 \partial_1\varpi
       =\mathsf QA-\mathsf QA^{\mathrm T},
 \qquad
 \partial_1^*\widetilde{\mathsf Q}A^{\mathrm T}=0.
\]
The face average therefore removes the transverse term
from the following pairing:
\[
 \begin{split}
 \|\partial_1\varpi\|^2
   &=\langle\partial_1\varpi,\mathsf QA\rangle
     -\langle\partial_1\varpi,
                   (\mathsf Q-\widetilde{\mathsf Q})A^{\mathrm T}\rangle
     -\langle\varpi,
                   \partial_1^*\widetilde{\mathsf Q}A^{\mathrm T}\rangle\\
   &\le\|\partial_1\varpi\|
             \bigl(\|\mathsf QA\|
                         +2\sqrt3\,\|\nabla A^{\mathrm T}\|\bigr).
 \end{split}
\]
Here \cref{I-12-8}(iv) bounds the difference of the two
averages. It follows that
\[
 \|\partial_1\varpi\|
      \le\|\mathsf QA\|+2\sqrt3\,\|\partial A\|.
\]

To estimate $h=\mathsf P\partial^*A$, use the scalar
projection formula \eqref{I-12-14}. It gives
\[
 h=\mathsf G'\mathsf Q'^*\omega,\qquad
 \omega=\mathsf M^{-1}\mathsf Q'\mathsf G'\partial^*A,
 \qquad
 \|h\|^2=\langle\omega,\mathsf M\omega\rangle
                       \ge\gamma'\|\omega\|^2.
\]
Moreover, $\Delta'h=\mathsf Q'^*\omega$, so with
$\omega'=\omega-\mathsf Q'h$ we have
\[
 \Delta h=\mathsf Q'^*\omega',\qquad
 \|\Delta h\|=\|\omega'\|
      \le(\gamma'^{-1/2}+1)\|h\|
      =\mathsf b_{\mathrm B}\|h\|,
 \qquad
 \|\partial h\|^2
      =\langle h,\Delta h\rangle
      \le\mathsf b_{\mathrm B}\|h\|^2.
\]
We used the isometry of $\mathsf Q'^*$ and the contraction
property of $\mathsf Q'$ from \cref{I-12-8}(i).

The same identities identify the pairing that controls $h$:
\[
 \omega'
   =\mathsf Q'\mathsf Q'^*\omega'
   =\mathsf Q'\partial^*\partial h
   =\partial_1^*\widetilde{\mathsf Q}\partial h.
\]
Since $\mathsf P$ is an orthogonal projection and
$\partial^*A^{\mathrm T}=0$, this implies
\[
 \begin{split}
 \|h\|^2
   &=\langle\partial^*A,h\rangle\\
   &=\langle A^{\mathrm T},\partial h\rangle
                    +\langle\partial\lambda_{\mathrm c},\partial h\rangle\\
   &=\langle\lambda_{\mathrm c},\mathsf Q'^*\omega'\rangle
     =\langle\varpi,\partial_1^*
                              \widetilde{\mathsf Q}\partial h\rangle
     =\langle\partial_1\varpi,\widetilde{\mathsf Q}\partial h\rangle.
 \end{split}
\]
The second differences obey
\[
 \begin{split}
 \|\nabla\partial h\|^2
   &=\sum_{\mu,\nu}\|\partial_\mu\partial_\nu h\|^2\\
   &=\sum_{\mu,\nu}
        \langle h,\partial_\nu^*\partial_\mu^*
                              \partial_\mu\partial_\nu h\rangle
     =\langle h,\Delta^2h\rangle=\|\Delta h\|^2,
 \end{split}
\]
again by commutation of the translations.
Consequently \cref{I-12-8}(i),(iv) gives
\[
 \begin{split}
 \|\widetilde{\mathsf Q}\partial h\|
   &\le\|\mathsf Q\partial h\|
                              +2\sqrt3\,\|\nabla\partial h\|\\
   &\le\|\partial h\|+2\sqrt3\,\|\Delta h\|
     \le C_h\|h\|.
 \end{split}
\]
Inserting this in the pairing yields
$\|h\|\le C_h\|\partial_1\varpi\|$.
Together with the preceding bound on $\partial_1\varpi$,
this proves the first inequality in \eqref{I-12-19}.

It remains to estimate the full norm of $A$.
The quadratic form is
\[
 \mathcal E:=\langle A,\Delta_aA\rangle
    =\|\partial A\|^2+\|\mathsf R\partial^*A\|^2
                                           +a\|\mathsf QA\|^2.
\]
Orthogonality of $\mathsf R$ and $\mathsf P=1-\mathsf R$,
the ambient identity above, and the estimate just proved give
\[
 \begin{split}
 \|\nabla A\|^2
   &=\|\partial A\|^2+\|\mathsf R\partial^*A\|^2
                                  +\|\mathsf P\partial^*A\|^2\\
   &\le(1+24C_h^2)\|\partial A\|^2
        +\|\mathsf R\partial^*A\|^2+2C_h^2\|\mathsf QA\|^2\\
   &\le C'_{\mathrm B}\mathcal E .
 \end{split}
\]
The bond-average estimate \cref{I-12-8}(v) now implies
\[
 \|A\|^2
   \le24\|\nabla A\|^2+2\|\mathsf QA\|^2
   \le(24C'_{\mathrm B}+2a^{-1})\mathcal E .
\]
Adding the last two bounds gives
$\|\nabla A\|^2+\|A\|^2
 \le\gamma_{0,\mathrm B}^{-1}\mathcal E$,
as required. The constants in \eqref{I-12-18} depend only
on $a$, and $C'_{\mathrm B}\ge1$ gives
$\gamma_{0,\mathrm B}^{-1}\ge25$.
\end{proof}

\subsection{Exponential conjugation}
\label{sec:one-order-conjugation}

Let $\theta$ be a real function on $\mathbb L_\eta$ such that
\[
 |\theta(x)-\theta(x')|\le\delta|x-x'|_1.
\]
On coarse fields use its restriction to $\mathbb L_1$.
For an operator $\mathcal O$ between any of these spaces, put
$\mathcal O_\theta=e^{-\theta}\mathcal Oe^\theta$,
with multiplication on its respective domain and target.
In particular, write $\Delta_{a,\theta}=(\Delta_a)_\theta$.
Only the Lipschitz constant of $\theta$ enters the following
estimates.

\begin{lemma}\label{I-12-26}
Suppose $0<\delta\le1/8$ and put $\epsilon_0=8\delta$.
\begin{enumerate}[label=(\roman*)]
\item For
\[
 \mathcal O\in
 \{\partial,\partial^*,\nabla,\nabla^*,
      \mathsf Q',\mathsf Q'^*,\mathsf Q,\mathsf Q^*\},
 \qquad
 \|\mathcal O_\theta-\mathcal O\|\le\epsilon_0.
\]
Here $\partial$ acts on scalars, $\partial^*$ on one-forms,
and $\nabla$ is the componentwise gradient, with its
corresponding adjoint.

\item If $5\epsilon_0\le1/6$, then
\[
 \|\mathsf G'_\theta g\|_1\le6\|g\|,\qquad
 \|\mathsf G'_\theta-\mathsf G'\|\le90\epsilon_0,\qquad
 \|\mathsf M_\theta-\mathsf M\|\le1755\epsilon_0.
\]
If, in addition, $1755\epsilon_0\le\gamma'/2$, then
\[
 \|\mathsf M_\theta^{-1}\|\le2\gamma'^{-1},\qquad
 \|\mathsf M_\theta^{-1}-\mathsf M^{-1}\|
       \le3510\gamma'^{-2}\epsilon_0.
\]

\item If $c_P\epsilon_0\le1$, then
\[
 \|\mathsf P_\theta-\mathsf P\|\le c_P\epsilon_0,
 \qquad \|\mathsf P_\theta\|\le2,
\]
and for all one-forms $A,A'$,
\[
 \bigl|\langle A',(\Delta_{a,\theta}-\Delta_a)A\rangle\bigr|
       \le c_V\epsilon_0\|A'\|_1\|A\|_1.
\]
\end{enumerate}
\end{lemma}

The scalar hypotheses in (ii) will also be useful when the
stronger condition in (iii) is unavailable.

\begin{proof}
Conjugation preserves products and inverses, and
\[
 (\mathcal O^*)_\theta=(\mathcal O_{-\theta})^*.
\]
Every indicator commutes with multiplication by $e^\theta$.
These identities hold on the finite tori without any
restriction on the magnitude of $\theta$.

\emph{Differences and averages.}
For a forward difference,
\[
 (\partial_\mu)_\theta=\partial_\mu+\beta_\mu\tau_\mu,
 \qquad
 \beta_\mu(x)=\eta^{-1}
       \bigl(e^{\theta(x+\eta e_\mu)-\theta(x)}-1\bigr).
\]
The quotient distance between the two sites is at most $\eta$,
so $|\beta_\mu(x)|\le2\delta$. Here we used
$|e^s-1|\le e^{|s|}-1\le2|s|$ for $|s|\le1$,
and $\delta\eta\le1$.
Since translations preserve the norm,
\[
 \|(\nabla_\theta-\nabla)F\|^2
       \le12\delta^2\|F\|^2.
\]
In particular, $\|\nabla_\theta-\nabla\|
\le2\sqrt3\,\delta\le\epsilon_0$; the same holds for
$\partial$ on scalars.

The kernels of $\mathsf Q'$ and $\mathsf Q$ are nonnegative.
In a scalar block average every contributing site has
distance less than $3$ from its coarse site. In the straight
one-form average the distance is less than $4$.
The covering formulas remain valid with all multiplicities
when some of their images coincide on the torus.
Consequently, component by component,
\[
 \begin{split}
 |(\mathsf Q_\theta-\mathsf Q)A|
       &\le(e^{4\delta}-1)\mathsf Q|A|,\\
 |(\mathsf Q'_\theta-\mathsf Q')f|
       &\le(e^{3\delta}-1)\mathsf Q'|f|.
 \end{split}
\]
The contraction bounds in \cref{I-12-8} give
$\|\mathsf Q_\theta-\mathsf Q\|\le8\delta$ and
$\|\mathsf Q'_\theta-\mathsf Q'\|\le6\delta$.
Applying the same estimates to $-\theta$ and taking
adjoints proves (i).

\emph{The scalar inverse.}
For scalar functions $f,g$,
\[
 \langle g,\Delta'_\theta f\rangle
   =\langle\partial_{-\theta}g,\partial_\theta f\rangle
     +\langle\mathsf Q'_{-\theta}g,\mathsf Q'_\theta f\rangle .
\]
Subtract the unweighted identity. Part (i) and
$\|\mathsf Q'\|\le1$ show that the absolute difference is at most
\[
 \begin{split}
 &\epsilon_0\bigl(\|g\|\,\|\partial f\|
                     +\|\partial g\|\,\|f\|\bigr)
       +\epsilon_0^2\|g\|\,\|f\|\\
 &\hspace{25mm}
       +(2\epsilon_0+\epsilon_0^2)\|g\|\,\|f\|
 \le5\epsilon_0\|g\|_1\|f\|_1.
 \end{split}
\]
Indeed, the sum in parentheses is at most
$\|g\|_1\|f\|_1$ by the two-dimensional Cauchy--Schwarz
inequality, and $\epsilon_0\le1$.
If $5\epsilon_0\le1/6$, \cref{I-12-12} therefore gives
\[
 \langle f,\Delta'_\theta f\rangle\ge\tfrac16\|f\|_1^2.
\]
The inverse is $\mathsf G'_\theta$ by conjugation.
Setting $f=\mathsf G'_\theta g$ in this inequality yields
$\|\mathsf G'_\theta g\|_1\le6\|g\|$.
The resolvent identity reads
\[
 \mathsf G'_\theta-\mathsf G'
       =\mathsf G'(\Delta'-\Delta'_\theta)\mathsf G'_\theta .
\]
Pairing with an arbitrary scalar $h$, using the symmetry
of $\mathsf G'$, and applying the preceding form estimate,
we obtain
\[
 \bigl|\langle h,(\mathsf G'_\theta-\mathsf G')g\rangle\bigr|
 \le5\epsilon_0\|\mathsf G'h\|_1
                         \|\mathsf G'_\theta g\|_1
 \le90\epsilon_0\|h\|\,\|g\|.
\]

\emph{The coarse inverse and the projection.}
We retain the order of the factors when estimating products.
If $\|\mathcal O_j\|\le m_j$,
$\|(\mathcal O_j)_\theta\|\le2m_j$ and
$\|(\mathcal O_j)_\theta-\mathcal O_j\|\le\rho_j$, then
\[
 \prod_{j=1}^r(\mathcal O_j)_\theta-\prod_{j=1}^r\mathcal O_j
 =\sum_{i=1}^r
   \left(\prod_{j<i}(\mathcal O_j)_\theta\right)
   \bigl((\mathcal O_i)_\theta-\mathcal O_i\bigr)
   \left(\prod_{j>i}\mathcal O_j\right).
\]
Thus
\begin{equation}\label{I-12-28}
 \left\|\prod_{j=1}^r(\mathcal O_j)_\theta
                    -\prod_{j=1}^r\mathcal O_j\right\|
       \le\sum_{i=1}^r2^{i-1}\rho_i\prod_{j\ne i}m_j.
\end{equation}
Apply \eqref{I-12-28} to
$\mathsf M=\mathsf Q'\mathsf G'^2\mathsf Q'^*$, using
$(m_i)=(1,3,3,1)$ and
$(\rho_i)=\epsilon_0(1,90,90,1)$. The preceding bounds give
$\|(\mathcal O_i)_\theta\|\le2m_i$, and hence
\[
 \|\mathsf M_\theta-\mathsf M\|
 \le1701\epsilon_0\le1755\epsilon_0.
\]
If $1755\epsilon_0\le\gamma'/2$, then
$\mathsf M\ge\gamma'$ and the preceding estimate imply
$\|\mathsf M_\theta\omega\|\ge(\gamma'/2)\|\omega\|$.
Thus $\|\mathsf M_\theta^{-1}\|\le2\gamma'^{-1}$, and
\[
 \mathsf M_\theta^{-1}-\mathsf M^{-1}
   =\mathsf M_\theta^{-1}
       (\mathsf M-\mathsf M_\theta)\mathsf M^{-1}
\]
gives the inverse-difference bound in (ii).

Now suppose $c_P\epsilon_0\le1$. Since $0<\gamma'\le1$
and $c_P=548208\gamma'^{-2}$, both scalar thresholds in
(ii) hold. Apply \eqref{I-12-28} to
$\mathsf P=\mathsf G'\mathsf Q'^*\mathsf M^{-1}
\mathsf Q'\mathsf G'$ with
\[
 (m_i)=(3,1,\gamma'^{-1},1,3),\qquad
 (\rho_i)=\epsilon_0(90,1,3510\gamma'^{-2},1,90).
\]
Bounding $2^{i-1}$ by $16$ gives
\[
 \|\mathsf P_\theta-\mathsf P\|
 \le144\gamma'^{-1}(62+3510\gamma'^{-1})\epsilon_0
 \le c_P\epsilon_0.
\]
The last inequality uses $\gamma'\le1$. Since $\mathsf P$
is an orthogonal projection,
$\|\mathsf P_\theta\|\le1+c_P\epsilon_0\le2$.

\emph{The vector form.}
By the Hodge identity used in
\cref{sec:one-order-coercivity},
\[
 \Delta_a=\nabla^*\nabla-\partial\mathsf P\partial^*
                                  +a\mathsf Q^*\mathsf Q.
\]
Conjugation and the adjoint rule therefore give
\[
 \begin{split}
 \langle A',\Delta_{a,\theta}A\rangle
   ={}&\langle\nabla_{-\theta}A',\nabla_\theta A\rangle\\
     &-\langle(\partial^*)_{-\theta}A',
                         \mathsf P_\theta(\partial^*)_\theta A\rangle
       +a\langle\mathsf Q_{-\theta}A',\mathsf Q_\theta A\rangle .
 \end{split}
\]
After subtracting the unweighted terms, the first term
contributes at most
$(\epsilon_0+\epsilon_0^2)\|A'\|_1\|A\|_1
\le2\epsilon_0\|A'\|_1\|A\|_1$.
The last contributes at most
$a(2\epsilon_0+\epsilon_0^2)\|A'\|\,\|A\|
\le3a\epsilon_0\|A'\|_1\|A\|_1$.
For the middle term set
\[
 f=\partial^*A,\quad f'=\partial^*A',\qquad
 r=((\partial^*)_\theta-\partial^*)A,\quad
 r'=((\partial^*)_{-\theta}-\partial^*)A'.
\]
The Hodge identity implies $\|f\|\le\|\nabla A\|$,
and (i) gives $\|r\|\le\epsilon_0\|A\|$;
the analogous bounds hold for $f'$ and $r'$.
Expanding the difference with these four terms bounds it by
\[
 \begin{split}
 &c_P\epsilon_0\|f'\|\,\|f\|
       +2\bigl(\|r'\|\,\|f\|+\|f'\|\,\|r\|
                                  +\|r'\|\,\|r\|\bigr)\\
 &\hspace{15mm}
 \le(c_P+2+2\epsilon_0)\epsilon_0\|A'\|_1\|A\|_1
 \le(c_P+4)\epsilon_0\|A'\|_1\|A\|_1.
 \end{split}
\]
Adding the three estimates gives
$c_V=c_P+6+3a$, as claimed.
\end{proof}

\subsection{The compressed inverse and its decay}
\label{sec:one-order-decay}

\begin{proof}[Proof of \cref{I-12-29}]
For $A\in\mathcal H_S$, compression does not change its
quadratic form:
\[
 \langle A,\Delta_a(S)A\rangle
       =\langle A,\Delta_aA\rangle
       \ge\gamma_{0,\mathrm B}\|A\|_1^2
\]
by \cref{I-12-17}. The operator is symmetric on the
finite-dimensional space $\mathcal H_S$ and is positive
definite there. It therefore has a symmetric positive
inverse on that space. This includes $S=\varnothing$,
where $\mathcal H_S$ is the zero space.
Extending the inverse by zero gives
\[
 G(S)=G(S)\chi_S=\chi_SG(S)\chi_S.
\]

For a one-form $J$, set $A=G(S)J$. Then
$\Delta_a(S)A=\chi_SJ$, and hence
\[
 \gamma_{0,\mathrm B}\|A\|_1^2
       \le\langle A,J\rangle\le\|A\|_1\|J\|.
\]
For a rank-two tensor $J$, instead set
$A=G(S)\nabla^*J$. The same argument, followed by
summation by parts on the whole torus, gives
\[
 \gamma_{0,\mathrm B}\|A\|_1^2
       \le\langle A,\nabla^*J\rangle
       =\langle\nabla A,J\rangle
       \le\|A\|_1\|J\|.
\]
Both cases yield $\|A\|_1\le\gamma_{0,\mathrm B}^{-1}\|J\|$,
which proves the four bounds in (b).
The derivative in the second case is that of the ambient
zero extension; no derivative intrinsic to $S$ occurs.

For the decay estimate, use the rate $\delta_0$ in
\eqref{I-12-30} and set
\[
 \epsilon_0=8\delta_0
       =\frac{\gamma_{0,\mathrm B}}{2c_V}.
\]
Since $\gamma_{0,\mathrm B}\le1$ and $1\le c_P\le c_V$,
we have $\delta_0\le1/8$ and $c_P\epsilon_0\le1$.
Let $\theta$ be any real function with Lipschitz constant
at most $\delta_0$. Multiplication by $e^{\pm\theta}$
preserves $\mathcal H_S$ and its orthogonal complement.
Thus
\[
 \Delta_a(S)_\theta
   =\chi_S\Delta_{a,\theta}\chi_S|_{\mathcal H_S},
 \qquad
 G(S)_\theta=e^{-\theta}G(S)e^\theta
\]
are inverse on $\mathcal H_S$, and $G(S)_\theta$ is zero
on $\mathcal H_S^\perp$.
By \cref{I-12-26}(iii), for $A\in\mathcal H_S$,
\[
 \begin{split}
 \langle A,\Delta_a(S)_\theta A\rangle
 &\ge\langle A,\Delta_a(S)A\rangle
                    -c_V\epsilon_0\|A\|_1^2\\
 &\ge\tfrac12\gamma_{0,\mathrm B}\|A\|_1^2.
 \end{split}
\]

Take $\mathcal O_2=1$ or $\nabla^*$ and set
$A=G(S)_\theta(\mathcal O_2)_\theta J$,
where $J$ is a one-form or a rank-two tensor, respectively.
Then $A\in\mathcal H_S$ and
\[
 \Delta_a(S)_\theta A
       =\chi_S(\mathcal O_2)_\theta J.
\]
When $\mathcal O_2=1$, the pairing of the right-hand side
with $A$ is at most $\|A\|\,\|J\|$ in absolute value.
When $\mathcal O_2=\nabla^*$, the adjoint conjugation rule
and \cref{I-12-26}(i) instead give
\[
 \begin{split}
 \bigl|\langle A,(\nabla^*)_\theta J\rangle\bigr|
 &=\bigl|\langle\nabla_{-\theta}A,J\rangle\bigr|\\
 &\le(\|\nabla A\|+\epsilon_0\|A\|)\|J\|
 \le\sqrt2\,\|A\|_1\|J\|.
 \end{split}
\]
Consequently, in both cases,
$\|A\|_1\le2\sqrt2\,\gamma_{0,\mathrm B}^{-1}\|J\|$.
For $\mathcal O_1=1$ or $\nabla$, part (i) also gives
\[
 \|(\mathcal O_1)_\theta A\|\le\sqrt2\,\|A\|_1.
\]
We have therefore proved all four conjugated bounds
\[
 \bigl\|(\mathcal O_1)_\theta G(S)_\theta
                         (\mathcal O_2)_\theta\bigr\|
       \le4\gamma_{0,\mathrm B}^{-1}.
\]

Finally, for the given nonempty set $U'$ choose
\[
 \theta(x)=-\delta_0\operatorname{dist}(x,U').
\]
This function has Lipschitz constant at most $\delta_0$.
It vanishes on $U'$ and satisfies
$\theta(x)\le-\delta_0\operatorname{dist}(U,U')$
for $x\in U$. Undoing the conjugation gives
\[
 \begin{split}
 \chi_U\mathcal O_1G(S)\mathcal O_2\chi_{U'}
   ={}&\chi_Ue^\theta
       (\mathcal O_1)_\theta G(S)_\theta
       (\mathcal O_2)_\theta e^{-\theta}\chi_{U'}.
 \end{split}
\]
The two exterior multiplication operators have norms at most
$e^{-\delta_0\operatorname{dist}(U,U')}$ and $1$,
respectively. The preceding four bounds now imply
\eqref{I-12-32}.
\end{proof}

\section{Interior estimates for the scalar equation}
\label{sec:scalar-interior}
\label{I-13-40}

We next obtain pointwise and H\"older bounds from a local
$L^2$ bound and the lattice Poisson equation.
Fix $\xi>0$ and work on $\xi\mathbb Z^3$ with real functions
and real fields $F=(F_\mu)_{\mu=1}^3$.
The difference operators and the positive Laplacian are
\[
 \begin{gathered}
 \partial_\mu f(x)=\xi^{-1}\bigl(f(x+\xi e_\mu)-f(x)\bigr),
 \qquad
 \partial_\mu^*f(x)=\xi^{-1}\bigl(f(x-\xi e_\mu)-f(x)\bigr),\\
 \nabla^*F=\sum_\mu\partial_\mu^*F_\mu,\quad
 \Delta=\sum_\mu\partial_\mu^*\partial_\mu .
 \end{gathered}
\]
We write $\nabla f=(\partial_\mu f)_\mu$,
$|x|_1=\sum_\mu|x_\mu|$ and
$|x|_\infty=\max_\mu|x_\mu|$.
For a set $X$ of sites put
\[
 \|f\|_X^2=\xi^3\sum_{x\in X}|f(x)|^2,\qquad
 |f|_X=\sup_{x\in X}|f(x)|.
\]
Write $|f|=|f|_{\xi\mathbb Z^3}$ and
$|F|=\max_\mu|F_\mu|$ for the global suprema.
A function on a subset is extended by zero when evaluated
outside that subset.

\begin{proposition}\label{I-13-47}
There are an absolute constant $c_I>0$ and constants
$c_I(\alpha)>0$, depending only on $0<\alpha<1$, with the
following property. Let $r\ge16\xi$, let $x_0\in\xi\mathbb Z^3$,
and set
\[
 \mathsf B_s=\{x\in\xi\mathbb Z^3:|x-x_0|_\infty\le s\}.
\]
Suppose that $u$ is a function on $\mathsf B_r$, that $g$
and the field $F$ are defined on $\xi\mathbb Z^3$, and that
\[
 \Delta u=g+\nabla^*F\qquad\text{on }\mathsf B_{r-\xi}.
\]
Put
\[
 \mathsf N=r^{-3/2}\|u\|_{\mathsf B_r}
                       +r^2|g|_{\mathsf B_r}+r|F|.
\]
For all $x,x'\in\mathsf B_{r/4}$ the following estimates hold.
\begin{enumerate}[label=(\alph*)]
\item For every $0<\alpha<1$,
\[
 |u(x)|\le c_I\mathsf N,\qquad
 |u(x)-u(x')|
      \le c_I(\alpha)
               \left(\frac{|x-x'|_1}{r}\right)^\alpha\mathsf N.
\]
\item If $F=0$, then for every direction $\nu$ and every
$0<\alpha<1$,
\[
 r|\partial_\nu u(x)|\le c_I\mathsf N,\qquad
 r|\partial_\nu u(x)-\partial_\nu u(x')|
      \le c_I(\alpha)
               \left(\frac{|x-x'|_1}{r}\right)^\alpha\mathsf N.
\]
\end{enumerate}
The constants are independent of $\xi$, $r$, $x_0$, $u$, $g$
and $F$.
\end{proposition}

The proof is given in \cref{sec:scalar-interior-proof}.
Multiplication by a cutoff produces a finitely supported
function whose Laplacian is the prescribed forcing plus
terms supported in an annulus away from $\mathsf B_{r/4}$.
The free kernel constructed below represents both contributions.
The kernel and its first differences have the local summability
bounds needed for the forcing term.
Separation from the annulus permits further differences of
the other term, with bounds in terms of
$\|u\|_{\mathsf B_r}$ and $|F|$.
When $F=0$, the potential of the local forcing also has the derivative bounds
in part~(b).

\subsection{The free scalar kernel}
\label{sec:free-scalar-kernel}

For functions $k,g$ on $\xi\mathbb Z^3$ write
\[
 (k*g)(x)=\xi^3\sum_{z\in\xi\mathbb Z^3}k(x-z)g(z)
\]
whenever the sum converges absolutely.
A mixed difference $\partial^\beta$ denotes any product
of $|\beta|$ operators chosen from the forward and adjoint
differences $\partial_\mu,\partial_\mu^*$; the empty product
is the identity.

\begin{lemma}\label{I-13-45}
There are an absolute constant $c_K>0$ and constants
$c_K(\alpha)>0$, depending only on $0<\alpha<1$, such that
for every $\xi>0$ there is a function $C_{0,\mathrm{lap}}$
on $\xi\mathbb Z^3$ with the following properties.
\begin{enumerate}[label=(\roman*)]
\item For every finitely supported function $v$,
\[
 C_{0,\mathrm{lap}}*(\Delta v)=v.
\]
\item For every mixed difference of order $|\beta|\le3$,
\[
 |\partial^\beta C_{0,\mathrm{lap}}(x)|
      \le c_K(|x|_1+\xi)^{-1-|\beta|}
                  \qquad(x\in\xi\mathbb Z^3).
\]
\item For $s'\ge s\ge\xi$ and $p\in\{1,2\}$,
\[
 \begin{split}
 \sum_{|z|_1\le s}\xi^3(|z|_1+\xi)^{-p}
                  &\le c_Ks^{3-p},\\
 \sum_{s<|z|_1\le s'}\xi^3|z|_1^{-3}
                  &\le c_K\bigl(1+\log(s'/s)\bigr).
 \end{split}
\]
\item Suppose that a kernel $k$ satisfies
\[
 |k(z)|\le(|z|_1+\xi)^{-2},\qquad
 |\partial_\nu k(z)|\le(|z|_1+\xi)^{-3}
                 \qquad(z\in\xi\mathbb Z^3,\ 1\le\nu\le3).
\]
Let $r\ge\xi$, let $x,x'\in\xi\mathbb Z^3$ satisfy
$|x-x'|_1\le3r$, and let $G$ vanish whenever
$|z-x|_1>6r$. Then
\[
 \begin{split}
 |(k*G)(x)|&\le c_Kr|G|,\\
 |(k*G)(x)-(k*G)(x')|
      &\le c_K(\alpha)\,r
                \left(\frac{|x-x'|_1}{r}\right)^\alpha|G|
                    \qquad(0<\alpha<1).
 \end{split}
\]
\end{enumerate}
\end{lemma}

\begin{proof}
The letter $c$ denotes an absolute constant that may change
between estimates. We choose $c_K$ large enough to cover
the finitely many bounds below.

\emph{The one-dimensional heat kernel.}
Put $\mathrm i=\sqrt{-1}$ and
$\varpi(k)=2\xi^{-2}(1-\cos(\xi k))$.
For $a,a'\ge0$ integers with $a+a'\le3$, define
\[
 q_T^{a,a'}(s)=\frac1{2\pi}
       \int_{-\pi/\xi}^{\pi/\xi}
       e^{\mathrm iks}
       \bigl[\xi^{-1}(e^{\mathrm i\xi k}-1)\bigr]^a
       \bigl[\xi^{-1}(e^{-\mathrm i\xi k}-1)\bigr]^{a'}
       e^{-T\varpi(k)}\,dk.
\]
Here $s\in\xi\mathbb Z$ and $T\ge0$.
Write $q_T=q_T^{0,0}$.
The forward and adjoint differences on $\xi\mathbb Z$
multiply $e^{\mathrm iks}$ by the two factors in brackets,
so $q_T^{a,a'}=\partial^a\partial^{*a'}q_T$.
Differentiation under the integral on its compact interval
gives
\[
 \frac{d}{dT}q_T=-\partial^*\partial q_T,\qquad
 q_0(s)=\xi^{-1}\mathbf 1_{\{0\}}(s).
\]

\emph{The contour estimate.}
For $T>0$ set
$\vartheta=\max\{T,\xi^2\}$ and
$\sigma=\operatorname{sign}(s)\vartheta^{-1/2}$.
The integrand defining $q_T^{a,a'}$ is entire in $k$ and
has period $2\pi/\xi$, since $s/\xi$ is an integer.
Cauchy's theorem on a rectangle, whose vertical sides
cancel, therefore permits the replacement
$k\mapsto k+\mathrm i\sigma$.
Here $|\xi\sigma|\le1$, $T\sigma^2\le1$, and
\[
 |e^{\mathrm i(k+\mathrm i\sigma)s}|
                  =e^{-|s|/\sqrt\vartheta}.
\]
Since $(\cosh b-1)/b^2$ increases with $|b|$, we have
\[
 \begin{split}
 \operatorname{Re}\varpi(k+\mathrm i\sigma)
   &=\varpi(k)\cosh(\xi\sigma)
             -2\xi^{-2}\bigl(\cosh(\xi\sigma)-1\bigr)\\
   &\ge\varpi(k)-2(\cosh1-1)\sigma^2.
 \end{split}
\]
Consequently
$|e^{-T\varpi(k+\mathrm i\sigma)}|
 \le e^2e^{-T\varpi(k)}$.
For either choice of sign,
\[
 \begin{split}
 \left|\xi^{-1}
             \bigl(e^{\pm\mathrm i\xi(k+\mathrm i\sigma)}-1\bigr)\right|
   &\le\xi^{-1}
       \bigl(e|e^{\pm\mathrm i\xi k}-1|
                              +|e^{\mp\xi\sigma}-1|\bigr)\\
   &\le e\bigl(\varpi(k)^{1/2}+\vartheta^{-1/2}\bigr).
 \end{split}
\]
For $|k|\le\pi/\xi$,
\[
 \frac2\pi|k|
       \le\varpi(k)^{1/2}
       =\frac2\xi\left|\sin\frac{\xi k}{2}\right|
       \le|k|.
\]
It follows that
\[
 \begin{split}
 |q_T^{a,a'}(s)|
   &\le e^5e^{-|s|/\sqrt\vartheta}
       \int_{-\pi/\xi}^{\pi/\xi}
       (|k|+\vartheta^{-1/2})^{a+a'}
                                  e^{-4Tk^2/\pi^2}\,dk\\
   &\le c\,\vartheta^{-(1+a+a')/2}
                                  e^{-|s|/\sqrt\vartheta}.
 \end{split}
\]
For the last inequality, when $T\ge\xi^2$ substitute
$k=T^{-1/2}k'$ and bound the resulting Gaussian integral
over the real line. When $T<\xi^2$, the integrand is at
most $((\pi+1)/\xi)^{a+a'}$ on an interval of length
$2\pi/\xi$. These bounds are uniform for $a+a'\le3$.

\emph{The three-dimensional heat kernel.}
Define
$p_T(x)=\prod_{\mu=1}^3q_T(x_\mu)$.
Then
\[
 \frac{d}{dT}p_T=-\Delta p_T,\qquad
 p_0(x)=\xi^{-3}\mathbf 1_{\{0\}}(x).
\]
A mixed difference acts on the $\mu$th factor by
$\partial^{a_\mu}\partial^{*a'_\mu}$, where
$\sum_\mu(a_\mu+a'_\mu)=|\beta|$.
Multiplying the one-dimensional estimates yields
\begin{equation}\label{I-13-46}
 \begin{gathered}
 |\partial^\beta p_T(x)|
       \le c\,\vartheta^{-(3+|\beta|)/2}
                           e^{-|x|_1/\sqrt\vartheta},\\
 \vartheta=\max\{T,\xi^2\},\quad |\beta|\le3.
 \end{gathered}
\end{equation}

\emph{Time integration and the inverse identity.}
Define
\[
 C_{0,\mathrm{lap}}(x)=\int_0^\infty p_T(x)\,dT
\]
and put $\varkappa=3+|\beta|\in\{3,4,5,6\}$.
By \eqref{I-13-46},
\[
 \int_0^{\xi^2}|\partial^\beta p_T(x)|\,dT
        \le c\,\xi^{2-\varkappa}e^{-|x|_1/\xi}.
\]
For $x\ne0$, the substitution $T=|x|_1^2u^{-2}$ gives
\[
 \begin{split}
 \int_{\xi^2}^\infty|\partial^\beta p_T(x)|\,dT
   &\le c\int_0^\infty T^{-\varkappa/2}
                         e^{-|x|_1/\sqrt T}\,dT\\
   &=2c|x|_1^{2-\varkappa}
                  \int_0^\infty u^{\varkappa-3}e^{-u}\,du
     \le c|x|_1^{2-\varkappa}.
 \end{split}
\]
For $x=0$, the same integral is bounded by
\[
 c\int_{\xi^2}^\infty T^{-\varkappa/2}\,dT
       =\frac{2c}{\varkappa-2}\xi^{2-\varkappa}.
\]
All these integrals converge since $\varkappa\ge3$.
For a nonzero lattice site, $|x|_1\ge\xi$; moreover,
$(1+b)^{\varkappa-2}e^{-b}$ is bounded for $b\ge0$,
uniformly over the four possible values of $\varkappa$.
Combining the three estimates proves part~(ii).
They also show absolute convergence of the integral
defining $C_{0,\mathrm{lap}}$ and of the integrals of its
mixed differences, and hence
\[
 \partial^\beta C_{0,\mathrm{lap}}(x)
        =\int_0^\infty\partial^\beta p_T(x)\,dT.
\]

If $v$ is finitely supported, shifting the summation index
gives $(\Delta p_T)*v=p_T*(\Delta v)$, so
\[
 \int_0^{T_0}p_T*(\Delta v)\,dT
    =-\int_0^{T_0}\frac{d}{dT}(p_T*v)\,dT
    =v-p_{T_0}*v.
\]
For $T_0\ge\xi^2$, \eqref{I-13-46} gives
\[
 |p_{T_0}*v|
       \le cT_0^{-3/2}\xi^3\sum_z|v(z)|\longrightarrow0
                         \qquad(T_0\longrightarrow\infty).
\]
Letting $T_0\to\infty$ proves part~(i).

\emph{Lattice shell sums.}
For $s\ge\xi$, the number of sites with $|z|_1\le s$
is at most
$(2s/\xi+1)^3\le(3s/\xi)^3$.
Thus a shell $s_1<|z|_1\le2s_1$, with $2s_1\ge\xi$,
contributes at most $216s_1^{3-p}$ to
$\sum_z\xi^3|z|_1^{-p}$.
Cover $0<|z|_1\le s$ by the shells with
$s_1=2^{-i-1}s$, $i\ge0$; those with $2s_1<\xi$
are empty. For $p=1,2$,
\[
 \sum_{i\ge0}(2^{-i-1}s)^{3-p}\le s^{3-p}.
\]
For nonzero $z$, replace $(|z|_1+\xi)^{-p}$ by the
larger quantity $|z|_1^{-p}$, and add the contribution
$\xi^{3-p}\le s^{3-p}$ at $z=0$.
This proves the first estimate in part~(iii).
The annulus $s<|z|_1\le s'$ is covered by at most
$1+\log_2(s'/s)$ shells with $s_1=2^is$.
Each contributes at most $216$ when $p=3$, which proves
the second estimate in part~(iii).

\emph{The convolution modulus.}
The first estimate in part~(iv) follows from part~(iii)
with $p=2$ and $s=6r$.
For the second, the case $x=x'$ is immediate.
Otherwise put $\delta=|x-x'|_1\ge\xi$.
For $|z-x|_1\le2\delta$, we also have
$|z-x'|_1\le3\delta$.
Applying part~(iii) to these two balls bounds their
contribution to the convolution difference by
$5c_K\delta|G|$.

For $|z-x|_1>2\delta$, choose a shortest lattice path
$x=x^0,x^1,\ldots,x^{\delta/\xi}=x'$.
Every point $x''$ on this path satisfies
\[
 |x''-z|_1\ge|x-z|_1-\delta\ge\tfrac12|x-z|_1.
\]
Each increment
$k(x^{i+1}-z)-k(x^i-z)$ equals
$\pm\xi\partial_\nu k$ at one of its endpoints.
Summing the increments therefore gives
\[
 |k(x-z)-k(x'-z)|\le8\delta|x-z|_1^{-3}.
\]
Since $G(z)=0$ for $|z-x|_1>6r$, the second estimate
in part~(iii), with $s=2\delta$ and $s'=6r$, bounds
this contribution by
\[
 8c_K\delta\bigl(1+\log^+(3r/\delta)\bigr)|G|,
 \qquad \log^+t=\max\{\log t,0\}.
\]
Here $2\delta\le6r$ by hypothesis.
Writing $b=\delta/r\in(0,3]$, we conclude that
\[
 \begin{split}
 |(k*G)(x)-(k*G)(x')|
   &\le c_K\delta\bigl[5+8(1+\log^+(3r/\delta))\bigr]|G|\\
   &=c_Kr b^\alpha\,
           b^{1-\alpha}\bigl[5+8(1+\log^+(3/b))\bigr]|G|.
 \end{split}
\]
For $0<\alpha<1$, the last factor depending on $b$
is bounded on $(0,3]$ by a constant depending only
on $\alpha$. This proves part~(iv).
\end{proof}

\subsection{The cutoff argument}
\label{sec:scalar-interior-proof}

\begin{proof}[Proof of \cref{I-13-47}]
Use the fixed cutoff $\eta_0$ defined by
\[
 \begin{gathered}
 \phi(s)=
 \begin{cases}
  e^{-1/s},&s>0,\\
  0,&s\le0,
 \end{cases}
 \qquad
 \vartheta_0(s)=\frac{\phi(s)}{\phi(s)+\phi(1-s)},\\
 \eta_0(s)=\vartheta_0(4-6|s|).
 \end{gathered}
\]
Every derivative of $e^{-1/s}$ is a polynomial in $1/s$
times $e^{-1/s}$ and tends to zero as $s\downarrow0$,
so $\phi$ is smooth. The denominator in $\vartheta_0$
is positive, and the resulting fixed smooth cutoff satisfies
$0\le\eta_0\le1$, equals one on $[-1/2,1/2]$, and vanishes
outside $[-2/3,2/3]$. It is constant near zero, so composition
with $|s|$ preserves smoothness. Its derivative bounds are absolute.
In the proof, $c$ denotes absolute constants and
$c(\alpha)$ constants depending only on $\alpha$;
their values may increase.

\emph{The localized equation.}
Write $(\tau_{\pm\mu}h)(x)=h(x\pm\xi e_\mu)$.
Set
\[
 \zeta(x)=\prod_{\mu=1}^3
          \eta_0\bigl((x_\mu-x_{0,\mu})/r\bigr).
\]
Extend $u$ by zero outside $\mathsf B_r$ and put $v=\zeta u$.
Then $\zeta=1$ on $\mathsf B_{r/2}$ and vanishes outside
$\mathsf B_{2r/3}$. The mean value theorem and the
second-difference formula give
\[
 |\partial_\mu\zeta|\le\|\eta_0'\|_\infty r^{-1},
 \qquad
 |\Delta\zeta|\le3\|\eta_0''\|_\infty r^{-2}.
\]
Since $\xi\le r/16$, the functions
$\partial_\mu\zeta$, $\tau_{-\mu}\partial_\mu\zeta$ and
$\Delta\zeta$ vanish outside the annulus
\[
 \mathsf Z
   =\{z\in\xi\mathbb Z^3:
          7r/16\le|z-x_0|_\infty\le3r/4\}.
\]

We use the product identities with their shifts written
explicitly. If $x'\sim x$ denotes the six sites
$x\pm\xi e_\mu$, then
$\Delta f(x)=\xi^{-2}\sum_{x'\sim x}(f(x)-f(x'))$.
For functions $h,\psi$ and a field $H$ this gives
\begin{equation}\label{I-13-42}
 \Delta(h\psi)(x)
 =h(x)\Delta\psi(x)
       -\xi^{-2}\sum_{x'\sim x}(h(x')-h(x))\psi(x'),
\end{equation}
and
\begin{equation}\label{I-13-43}
 \begin{split}
 \xi^{-2}\sum_{x'\sim x}(h(x')-h(x))\psi(x')
  ={}&-\sum_{\mu=1}^3
       \partial_\mu^*\bigl[(\partial_\mu h)
                              (\psi+\tau_\mu\psi)\bigr](x)\\
     &+\psi(x)\Delta h(x).
 \end{split}
\end{equation}
To check the second identity, fix $\mu$ and write $e=\xi e_\mu$.
The corresponding adjoint difference is
\[
 \begin{split}
 &\partial_\mu^*\bigl[(\partial_\mu h)
                            (\psi+\tau_\mu\psi)\bigr](x)\\
 &\quad=\xi^{-2}\bigl[
       (h(x)-h(x-e))(\psi(x-e)+\psi(x))\\
 &\hspace{45mm}
       -(h(x+e)-h(x))(\psi(x)+\psi(x+e))\bigr].
 \end{split}
\]
Subtracting the sum of these expressions from
$\psi(x)\Delta h(x)$ cancels all terms containing $\psi(x)$
and leaves the left side of \eqref{I-13-43}.
The divergence product identity is
\begin{equation}\label{I-13-44}
 h\,\partial_\mu^*H_\mu
       =\partial_\mu^*(hH_\mu)
          +\tau_{-\mu}\bigl((\partial_\mu h)H_\mu\bigr).
\end{equation}
Indeed, the difference of its first two terms at $x$ is
$\xi^{-1}(h(x)-h(x-e))H_\mu(x-e)$, which is its last term.

Apply \eqref{I-13-42} and \eqref{I-13-43} with
$h=\zeta$, $\psi=u$. On the support of $\zeta$ the
equation for $u$ holds, since $2r/3\le r-\xi$.
Using \eqref{I-13-44} for its divergence term gives
the equation on the whole lattice
\[
 \begin{split}
 \Delta v&=\zeta g+G_0
                  +\sum_{\mu=1}^3\partial_\mu^*(\zeta F_\mu+G_\mu),\\
 G_0&=\sum_{\mu=1}^3
              \tau_{-\mu}\bigl((\partial_\mu\zeta)F_\mu\bigr)
                  -u\Delta\zeta,\\
 G_\mu&=(\partial_\mu\zeta)(u+\tau_\mu u).
 \end{split}
\]
All terms on the right have finite support. By
\cref{I-13-45}(i), and a shift of the convolution sum,
\[
 C_{0,\mathrm{lap}}*(\partial_\mu^*G)
           =(\partial_\mu^*C_{0,\mathrm{lap}})*G
 \qquad\text{for finitely supported }G.
\]
Consequently
\[
 \begin{split}
 v&=v_1+v_2,\\
 v_1&=C_{0,\mathrm{lap}}*(\zeta g)
       +\sum_{\mu=1}^3
               (\partial_\mu^*C_{0,\mathrm{lap}})*(\zeta F_\mu),\\
 v_2&=C_{0,\mathrm{lap}}*G_0
       +\sum_{\mu=1}^3
               (\partial_\mu^*C_{0,\mathrm{lap}})*G_\mu .
 \end{split}
\]

\emph{The terms supported in the annulus.}
The functions $G_0,G_\mu$ vanish outside $\mathsf Z$.
There are at most $(3r/\xi)^3$ sites in $\mathsf B_r$,
so Cauchy--Schwarz gives
\[
 \xi^3\sum_{\mathsf B_r}|u|
        \le(3r)^{3/2}\|u\|_{\mathsf B_r}.
\]
The same bound applies to the shifted values of $u$
occurring in $G_\mu$, because the sites in question
and their shifts belong to $\mathsf B_r$.
The derivative bounds for $\zeta$ therefore imply
\[
 \begin{split}
 \xi^3\sum_z|G_0(z)|
    &\le c\bigl(r^2|F|+r^{-1/2}\|u\|_{\mathsf B_r}\bigr)
      \le c r\,\mathsf N,\\
 \xi^3\sum_z|G_\mu(z)|
    &\le c r^{1/2}\|u\|_{\mathsf B_r}
      \le c r^2\,\mathsf N .
 \end{split}
\]
If $x\in\mathsf B_{3r/8}$ and $z\in\mathsf Z$, then
$|x-z|_1\ge|x-z|_\infty\ge r/16$.
Differences commute with these convolutions. For any
product $\partial^\beta$ of at most two forward or
adjoint differences, \cref{I-13-45}(ii) gives on this set
\[
 \begin{split}
 |(\partial^\beta C_{0,\mathrm{lap}})(x-z)|
       &\le c_K(r/16)^{-1-|\beta|},\\
 |(\partial^\beta\partial_\mu^*
                            C_{0,\mathrm{lap}})(x-z)|
       &\le c_K(r/16)^{-2-|\beta|}.
 \end{split}
\]
The second bound uses kernel differences of order
$|\beta|+1\le3$. Combining them with the two sum bounds yields
\[
 r^{|\beta|}|(\partial^\beta v_2)(x)|
          \le c\,\mathsf N
 \qquad(x\in\mathsf B_{3r/8},\ |\beta|\le2).
\]
Join $x,x'\in\mathsf B_{r/4}$ by a shortest lattice path
inside that box. Summing the first differences of $v_2$,
and then the second differences for $\partial_\nu v_2$,
along the path gives
\[
 |v_2(x)-v_2(x')|
     +r|\partial_\nu v_2(x)-\partial_\nu v_2(x')|
       \le c\,\frac{|x-x'|_1}{r}\,\mathsf N.
\]
Since $|x-x'|_1/r\le3$, the last factor is at most
$3^{1-\alpha}(|x-x'|_1/r)^\alpha$.
Thus $v_2$ and its first differences have the required
interior bounds.

\emph{The local right-hand side.}
For $x\in\mathsf B_{r/4}$, both $\zeta g$ and
$\zeta F_\mu$ vanish at sites $z$ with $|z-x|_1>3r$.
The kernel bound and the shell estimate
\cref{I-13-45}(ii),(iii), with $p=1$ and $s=3r$, imply
\[
 |(C_{0,\mathrm{lap}}*(\zeta g))(x)|
          \le c r^2|g|_{\mathsf B_r}.
\]
The normalized kernels
$c_K^{-1}\partial_\nu C_{0,\mathrm{lap}}$ and
$c_K^{-1}\partial_\mu^*C_{0,\mathrm{lap}}$
satisfy the hypotheses of \cref{I-13-45}(iv), by its
mixed-difference bounds. That clause applies with the
same radius $r$: its support requirement is the weaker
condition $|z-x|_1\le6r$, and
$|x-x'|_1\le3r$ for $x,x'\in\mathsf B_{r/4}$.
It gives
\[
 \begin{split}
 |\partial_\nu(C_{0,\mathrm{lap}}*(\zeta g))(x)|
     &\le c r|g|_{\mathsf B_r},\\
 |\partial_\nu(C_{0,\mathrm{lap}}*(\zeta g))(x)
       -\partial_\nu(C_{0,\mathrm{lap}}*(\zeta g))(x')|
     &\le c(\alpha)r
           \left(\frac{|x-x'|_1}{r}\right)^\alpha
                     |g|_{\mathsf B_r}.
 \end{split}
\]
For the divergence-source terms it gives
\[
 \begin{split}
 |((\partial_\mu^*C_{0,\mathrm{lap}})
                                      *(\zeta F_\mu))(x)|
     &\le c r|F|,\\
 |((\partial_\mu^*C_{0,\mathrm{lap}})
                                      *(\zeta F_\mu))(x)
   -((\partial_\mu^*C_{0,\mathrm{lap}})
                                      *(\zeta F_\mu))(x')|
     &\le c(\alpha)r
           \left(\frac{|x-x'|_1}{r}\right)^\alpha |F|.
 \end{split}
\]
Summing the gradient bound for
$C_{0,\mathrm{lap}}*(\zeta g)$ along a shortest path in
$\mathsf B_{r/4}$ also gives its difference bound
$c r|x-x'|_1|g|_{\mathsf B_r}$.
Using $|x-x'|_1/r\le
3^{1-\alpha}(|x-x'|_1/r)^\alpha$ once more, and summing
the three divergence-source terms, we obtain
\[
 \begin{split}
 |v_1(x)|
     &\le c\bigl(r^2|g|_{\mathsf B_r}+r|F|\bigr),\\
 |v_1(x)-v_1(x')|
     &\le c(\alpha)
       \left(\frac{|x-x'|_1}{r}\right)^\alpha
                   \bigl(r^2|g|_{\mathsf B_r}+r|F|\bigr).
 \end{split}
\]
When $F=0$, only $C_{0,\mathrm{lap}}*(\zeta g)$ remains
in $v_1$. Its gradient and gradient-difference estimates
above, multiplied by $r$, are bounded by
$c r^2|g|_{\mathsf B_r}$ and
$c(\alpha)(|x-x'|_1/r)^\alpha r^2|g|_{\mathsf B_r}$,
respectively.

For $x\in\mathsf B_{r/4}$, both $x$ and $x+\xi e_\nu$
belong to $\mathsf B_{r/2}$, where $u=v$.
Thus $u=v_1+v_2$ there and
$\partial_\nu u=\partial_\nu v_1+\partial_\nu v_2$
at every observation site. The bounds for $v_1$ and $v_2$
give (a), and their first-difference bounds give (b)
when $F=0$.
\end{proof}

\section{Regularity of the one-scale propagator}
\label{sec:one-scale-regularity}

Fix integers $L\ge2$ and $j\ge0$, put $\xi=L^{-j}$,
and let $N\ge4$ be an even positive multiple of $L$.
We work on
\[
 \mathsf T=\mathbb R^3/N\mathbb Z^3,
 \qquad \mathsf T_h=h\mathbb Z^3/N\mathbb Z^3
 \quad(h=\xi,1).
\]
Distances $|x-y|_1$ and $|x-y|_\infty$ are the least
lifted $\ell^1$ and $\ell^\infty$ distances on this torus.
All fields in this section are real. Write
$\mathcal F_\xi^m$ for the $m$-component functions on
$\mathsf T_\xi$, and $\mathcal A_\xi=\mathcal F_\xi^3$
for the one-forms. We use the $\xi^3$-weighted inner
products, summing over all components, and the associated
norm $\|\cdot\|$. In pointwise estimates $|f(x)|$ denotes
the maximum of the component moduli, and
$|f|=\sup_x|f(x)|$.

The differences, scalar gradient, divergence and curl have
the conventions of the preceding sections, now with spacing
$\xi$. Thus $\partial_\mu f(x)=\xi^{-1}
(f(x+\xi e_\mu)-f(x))$, $\nabla$ differentiates every
component, and $\Delta=\nabla^*\nabla$ acts componentwise.
The adjoints of $\nabla$ and $\nabla\nabla$ on one-forms
are $\nabla^*$ and $\nabla^*\nabla^*$, with domains $\mathcal F_\xi^9$ and
$\mathcal F_\xi^{27}$, respectively; these domains are
tensor spaces.

Use the corner blocks and straight averages of
\eqref{I-11-63}, with $k=j$, $\eta=\xi$ and all three
torus periods equal to $N$. We write these averages as
$\mathsf Q'_j$ and $\mathsf Q_j$, and the blocks as
\[
 B^j(y)=\mathsf T_\xi\cap\bigl(y+[0,1)^3\bigr),
 \qquad y\in\mathsf T_1.
\]
Let $\mathsf R_j$ be the orthogonal projection in the
scalar space onto $\Delta\ker\mathsf Q'_j$, and set
$\mathsf P_j=1-\mathsf R_j$. The unit-weight operator and
its inverse are
\[
 \Delta_{1,j}
   =\partial^*\partial+\partial\mathsf R_j\partial^*
                         +\mathsf Q_j^*\mathsf Q_j,
 \qquad G_j=\Delta_{1,j}^{-1}.
\]
The projection is formed on the whole torus. The inverse
exists by \cref{I-12-17}, applied with the parameters just
specified and $a=1$.

For $u\in\mathsf T$ and $r>0$ put
$U_r(u)=\{x\in\mathsf T_\xi:|x-u|_\infty\le r\}$.
For $0\le\alpha<1$ define
\[
 \|f\|_\alpha
  =\sup_{\substack{x,x'\in\mathsf T_\xi\\
                       0<|x-x'|_1\le1}}
       |x-x'|_1^{-\alpha}|f(x)-f(x')|,
\]
with value zero if there are no such pairs. In particular,
$\|f\|_0$ is an oscillation seminorm.
Let $\gamma_{\mathrm f}$ be $\gamma_{0,\mathrm B}$ of
\eqref{I-12-18} at $a=1$, let $\delta_{\mathrm f}$ be
$\delta_0$ of \eqref{I-12-30} at $a=1$, and let
$\gamma'$ be the scalar constant in \eqref{I-12-13}.
Set
\[
 \delta_P=\frac{\gamma'}{28080},
 \qquad
 \delta_* =\tfrac12\min\{\delta_{\mathrm f},\delta_P\}.
\]

\begin{proposition}\label{I-15-163}\label{I-15-121}
There are positive constants $c$, $c_\alpha$
$(0\le\alpha<1)$, $c'_\varepsilon$ $(0<\varepsilon<1)$,
and $c''_{\alpha,\varepsilon}$
$(\alpha\ge0,\ \varepsilon>0,\ \alpha+\varepsilon<1)$,
depending only on the displayed exponents, with the
following properties. In particular, the constants are
independent of $L,j,N$ and of the centers below.
Let $u,u'\in\mathsf T$, let the real function $\zeta$
vanish outside $U_1(u)$, and let
$J\in\mathcal A_\xi$, $J'\in\mathcal F_\xi^9$ and
$J''\in\mathcal F_\xi^{27}$ vanish outside $U_1(u')$.
Write $E=e^{-\delta_*|u-u'|_1}$. Then, for every
$x\in U_1(u)$,
\begin{enumerate}[label=\textup{(I\arabic*)},leftmargin=*]
\item The value, gradient and Laplacian satisfy
\[
 |(G_jJ)(x)|,\quad |(\nabla G_jJ)(x)|,\quad
 |(\Delta G_jJ)(x)|\le cE|J|,
 \qquad |(G_j\nabla^*J')(x)|\le cE|J'|.
\]
\item For $0\le\alpha<1$,
\[
 \begin{split}
 \|\zeta\nabla G_jJ\|_\alpha
   &\le c_\alpha E(\|\zeta\|_\alpha+|\zeta|)|J|,\\
 \|\zeta G_j\nabla^*J'\|_\alpha
   &\le c_\alpha E(\|\zeta\|_\alpha+|\zeta|)|J'|.
 \end{split}
\]
\item For $0<\varepsilon<1$,
\[
 |(\nabla G_j\nabla^*J')(x)|
    \le c'_\varepsilon E(\|J'\|_\varepsilon+|J'|).
\]
\item If $\alpha\ge0$, $\varepsilon>0$ and
$\alpha+\varepsilon<1$, then
\[
 \|\zeta\nabla G_j\nabla^*J'\|_\alpha
  \le c''_{\alpha,\varepsilon}E
       (\|\zeta\|_\alpha+|\zeta|)
       (\|J'\|_{\alpha+\varepsilon}+|J'|).
\]
\item The six $L^2$ estimates are
\[
 \begin{split}
 \|\zeta G_jJ\|,\quad \|\zeta\nabla G_jJ\|,\quad
 \|\zeta\nabla\nabla G_jJ\|
       &\le cE|\zeta|\|J\|,\\
 \|\zeta G_j\nabla^*J'\|,\quad
 \|\zeta\nabla G_j\nabla^*J'\|
       &\le cE|\zeta|\|J'\|,\\
 \|\zeta G_j\nabla^*\nabla^*J''\|
       &\le cE|\zeta|\|J''\|.
 \end{split}
\]
\end{enumerate}
\end{proposition}

These localized regularity and decay estimates have the
form of Ba\l aban's estimates in
\cite[Proposition~1.2]{B-CMP95}. We obtain them from the
Hilbert-space bounds and scalar interior estimates proved
above. For $f=G_jJ$, the componentwise equation is
\[
 \Delta f
   =J+\partial\mathsf P_j\partial^*f
                        -\mathsf Q_j^*\mathsf Q_jf.
\]
The first task is to control the projection term by the
decay of $f$ and its first differences. This gives a local
bound for $\Delta f$, and hence the second-difference
$L^2$ estimate. Scalar interior estimates then give the
pointwise and H\"older bounds. For a divergence source,
one further derivative reaches the critical decay of the
scalar kernel; the H\"older difference of the source
supplies the needed cancellation.
The proof of \textup{(I5)} is in
\cref{sec:one-scale-L2}; the proof of
\textup{(I1)--(I4)}, and the choice of a common $c$, are
completed in \cref{sec:one-scale-regularity-proof}.

\subsection{Decay of the scalar projection}
\label{sec:one-scale-projection}

We first record the support of an adjoint-average column.
It will also localize the term $\mathsf Q_j^*\mathsf Q_jf$
in the componentwise equation. For a coarse bond
$b=\langle y,y+e_\mu\rangle$, let $e_b$ be its unit
coordinate vector. The straight average is
\[
 (\mathsf Q_jA)_\mu(y)
  =\sum_{x\in B^j(y)}\xi^3
       \sum_{t=0}^{L^j-1}\xi A_\mu(x+t\xi e_\mu).
\]
Collecting the coefficient of $A_\nu(x)$ gives
\begin{equation}\label{eq:one-scale-adjoint-column}
 \begin{split}
 (\mathsf Q_j^*e_b)_\nu(x)
   &=\begin{cases}
      \xi\,\#\{0\le t<L^j:x-t\xi e_\mu\in B^j(y)\},
                                      &\nu=\mu,\\
      0,                              &\nu\ne\mu,
     \end{cases}\\
 0\le(\mathsf Q_j^*e_b)_\mu(x)&\le1,
 \qquad \operatorname{supp}(\mathsf Q_j^*e_b)
                   \subset y+[0,2)^3,
 \qquad \|\mathsf Q_j^*e_b\|\le1.
 \end{split}
\end{equation}
The box in this formula is understood as its image on
the torus. The coefficient is at most $\xi L^j=1$;
its integral is
\[
 \sum_x\xi^3(\mathsf Q_j^*e_b)_\mu(x)
    =\sum_{x\in B^j(y)}\xi^3
                         \sum_{t=0}^{L^j-1}\xi=1.
\]
Its squared norm is therefore at most one. The support
assertion follows directly from the paths in the average.

We use the following constants:
\begin{equation}\label{I-15-132}
 \begin{split}
 C_P&=414\gamma'^{-1},\qquad
 \sigma(\delta)=3\coth^3(\delta/2)\quad(\delta>0),\\
 c_*&=e^{12}+4\gamma_{\mathrm f}^{-1}e^{18}
       \bigl(13+4\sqrt3+\sqrt3C_P\sigma(\delta_*)\bigr).
 \end{split}
\end{equation}
These constants and the decay rates depend on none of
$L,j,N$. The explicit formulas for the preceding scalar
gap and vector bounds give
$\gamma',\gamma_{\mathrm f},\delta_{\mathrm f},\delta_P\le1$.
Here $C_P$ is distinct from the perturbation constant
$c_P$ in \eqref{I-12-27}.

\begin{proposition}\label{I-15-131}
For nonempty sets $U,U'\subset\mathsf T_\xi$,
\begin{equation}\label{I-15-134}
 \begin{split}
 \|\chi_U\partial\mathsf P_j\chi_{U'}\|
       &\le C_Pe^{-\delta_P\operatorname{dist}(U,U')},\\
 \|\chi_U\mathsf P_j\chi_{U'}\|
       &\le C_Pe^{-\delta_P\operatorname{dist}(U,U')},\\
 \|\chi_U\Delta\mathsf P_j\chi_{U'}\|
       &\le3C_Pe^{-\delta_P\operatorname{dist}(U,U')}.
 \end{split}
\end{equation}
The norms are operator norms between the appropriate
scalar and one-form spaces. For every scalar function
$g$, the function $\Delta\mathsf P_jg$ is constant on
each block, and
\[
 |\Delta\mathsf P_jg|_{B^j(y)}
       =\|\chi_{B^j(y)}\Delta\mathsf P_jg\|.
\]
Suppose, in addition, that $0<\delta\le\delta_P/2$,
$A\ge0$, $u'\in\mathsf T$ and
\[
 \|\chi_{B^j(y)}g\|\le A e^{-2\delta|y-u'|_1}
                 \qquad(y\in\mathsf T_1).
\]
Then, for every $u\in\mathsf T$, $r>0$ and
$y\in\mathsf T_1$,
\begin{equation}\label{I-15-136}
 \begin{split}
 \|\chi_{U_r(u)}\partial\mathsf P_jg\|
   &\le C_Pe^{3(r+1)}\sigma(\delta)A
                                 e^{-\delta|u-u'|_1},\\
 \|\chi_{B^j(y)}\mathsf P_jg\|
   &\le C_Pe^6\sigma(\delta)A e^{-\delta|y-u'|_1},\\
 |\Delta\mathsf P_jg|_{B^j(y)}
   &\le3C_Pe^6\sigma(\delta)A e^{-\delta|y-u'|_1}.
 \end{split}
\end{equation}
\end{proposition}

\begin{proof}
Use \cref{I-12-8,I-12-12} and \eqref{I-12-14} with
$\mathsf Q'=\mathsf Q'_j$, and put
\[
 \Delta'=\Delta+\mathsf Q_j'^*\mathsf Q'_j,\qquad
 \mathsf G'=(\Delta')^{-1},\qquad
 \mathsf M=\mathsf Q'_j\mathsf G'^2\mathsf Q_j'^*.
\]
Then $\mathsf M\ge\gamma'$ and
$\mathsf P_j=\mathsf G'\mathsf Q_j'^*\mathsf M^{-1}
\mathsf Q'_j\mathsf G'$. For a $\delta_P$-Lipschitz function
$\theta$, use the conjugation convention of
\cref{sec:one-order-conjugation}. Here
\[
 \epsilon_0=8\delta_P=\gamma'/3510,\qquad
 \delta_P\le\tfrac18,\quad
 5\epsilon_0\le\tfrac16,\quad
 1755\epsilon_0=\tfrac12\gamma'.
\]
Thus \cref{I-12-26}(i),(ii) gives
\[
 \begin{gathered}
 \|\partial_\theta-\partial\|\le\epsilon_0,\qquad
 \|(\mathsf Q'_j)_\theta\|,
 \|(\mathsf Q_j'^*)_\theta\|\le2,\\
 \|\mathsf G'_\theta g\|_1\le6\|g\|,\qquad
 \|\mathsf M_\theta^{-1}\|\le2\gamma'^{-1}.
 \end{gathered}
\]
Since $\epsilon_0\le1$,
$\|\partial_\theta\mathsf G'_\theta g\|
\le\sqrt2\|\mathsf G'_\theta g\|_1\le6\sqrt2\|g\|$.
Conjugating the formula for $\mathsf P_j$ and multiplying
these bounds yields
\[
 \| (\partial\mathsf P_j)_\theta\|
    \le6\sqrt2\cdot2\cdot2\gamma'^{-1}\cdot2\cdot6
    =288\sqrt2\gamma'^{-1}\le C_P,
 \qquad
 \|(\mathsf P_j)_\theta\|\le288\gamma'^{-1}\le C_P.
\]
The Laplacian can be handled without a second derivative
estimate for $\mathsf G'$. Indeed,
$\Delta=\Delta'-\mathsf Q_j'^*\mathsf Q'_j$ and
$\Delta'\mathsf G'=1$ give
\begin{equation}\label{I-15-137}
 \Delta\mathsf P_j
   =\mathsf Q_j'^*
       \bigl(\mathsf M^{-1}\mathsf Q'_j\mathsf G'
                              -\mathsf Q'_j\mathsf P_j\bigr).
\end{equation}
Consequently,
\[
 \| (\Delta\mathsf P_j)_\theta\|
  \le2\bigl(2\gamma'^{-1}\cdot2\cdot6
                         +2\cdot288\gamma'^{-1}\bigr)
  =1200\gamma'^{-1}\le3C_P.
\]
Choose $\theta(x)=-\delta_P\operatorname{dist}(x,U')$.
Undoing conjugation as in \cref{sec:one-order-decay}
multiplies each of these three operator bounds by at most
$e^{-\delta_P\operatorname{dist}(U,U')}$.
This proves \eqref{I-15-134}. Formula \eqref{I-15-137}
also places $\Delta\mathsf P_jg$ in the range of
$\mathsf Q_j'^*$, the functions constant on each block.
Since $\xi^3\#B^j(y)=1$, its block supremum and block
$L^2$ norm agree.

For the estimates with block inputs we need two elementary
consequences of the torus geometry. First,
\begin{equation}\label{I-15-138}
 \operatorname{dist}(U_r(u),U_{r'}(u'))
       \ge |u-u'|_1-3(r+r')
\end{equation}
whenever both sets are nonempty. This is the triangle
inequality, since each point in $U_r(u)$ lies at
$\ell^1$ distance at most $3r$ from $u$.
Second, for every $\lambda>0$ and every continuous
torus point $u$,
\[
 \sum_{y\in\mathsf T_1}e^{-\lambda|u-y|_1}
       \le\coth^3(\lambda/2)\le\sigma(\lambda).
\]
To see this, choose least-distance integer lifts of the
coarse sites; they are distinct, so the torus sum is
bounded by the corresponding sum over $\mathbb Z^3$.
For the fractional part $t\in[0,1]$ of one coordinate,
\[
 \sum_{n\in\mathbb Z}e^{-\lambda|n-t|}
    =\frac{e^{-\lambda t}+e^{-\lambda(1-t)}}
            {1-e^{-\lambda}}
    \le\coth(\lambda/2).
\]
The numerator is convex in $t$, and hence is largest at
an endpoint. Taking the product over the three
coordinates gives the assertion. The triangle inequality
then gives the convolution bound
\[
 \begin{split}
 &\sum_{y\in\mathsf T_1}
          e^{-2\delta|u-y|_1-2\delta|y-u'|_1}\\
 &\hspace{1cm}\le
 e^{-\delta|u-u'|_1}
       \sum_{y\in\mathsf T_1}e^{-\delta|u-y|_1}
 \le\sigma(\delta)e^{-\delta|u-u'|_1}.
 \end{split}
\]

If $0<\delta''\le\delta'\le1$, \eqref{I-15-138}
implies
\[
 e^{-\delta'\operatorname{dist}(U_r(u),U_{r'}(y'))}
  \le e^{3(r+r')}e^{-\delta''|u-y'|_1}.
\]
For $|u-y'|_1\le3(r+r')$ the right side is at least
one; otherwise apply \eqref{I-15-138} and then
$\delta'\ge\delta''$. Decompose
$g=\sum_{y'}\chi_{B^j(y')}g$ and apply this inequality
with $\delta'=\delta_P$, $\delta''=2\delta$ and $r'=1$,
using $B^j(y')\subset U_1(y')$. The first bound in
\eqref{I-15-134} gives
\[
 \begin{split}
 \|\chi_{U_r(u)}\partial\mathsf P_jg\|
 &\le C_Pe^{3(r+1)}A
       \sum_{y'\in\mathsf T_1}
         e^{-2\delta|u-y'|_1-2\delta|y'-u'|_1}\\
 &\le C_Pe^{3(r+1)}\sigma(\delta)A
                                      e^{-\delta|u-u'|_1}.
 \end{split}
\]
If $U_r(u)$ is empty the assertion is immediate.
Apply the other two bounds in \eqref{I-15-134} with
$r=r'=1$, and use $B^j(y)\subset U_1(y)$. This proves
the remaining bounds in \eqref{I-15-136}, with the block
norm replaced by its equal supremum for
$\Delta\mathsf P_jg$.
\end{proof}

\subsection{Localized second differences and the six $L^2$ bounds}
\label{sec:one-scale-L2}

Four of the six estimates in \textup{(I5)} follow
directly from the first-order inverse bounds. The new
estimate is the bound for $\nabla\nabla G_j$; its
adjoint gives the sixth one. We obtain the second
differences from the componentwise Laplacian, using a
cutoff whose derivative bounds remain valid on the
smallest permitted tori.

The specialization of \cref{I-12-29}(c) to the full
bond set and $a=1$ is
\begin{equation}\label{I-15-133}
 \|\chi_U\mathcal O_1G_j\mathcal O_2\chi_{U'}\|
    \le4\gamma_{\mathrm f}^{-1}
                 e^{-\delta_{\mathrm f}\operatorname{dist}(U,U')},
 \qquad
 \mathcal O_1\in\{1,\nabla\},\quad
 \mathcal O_2\in\{1,\nabla^*\},
\end{equation}
for nonempty site sets $U,U'$. Its hypotheses hold with
$k=j$, $\eta=\xi$ and periods $N$. This also identifies
the inverse in that result with $G_j$.

We first establish a local second-difference estimate
for an arbitrary one-form $f$. If $j=0$, then
\[
 \|\chi_{U_1(u)}\nabla\nabla f\|
                       \le12\|\chi_{U_3(u)}f\|.
\]
If $j\ge1$, then
\begin{equation}\label{I-15-135}
 \|\chi_{U_1(u)}\nabla\nabla f\|
 \le\|\chi_{U_3(u)}\Delta f\|
       +12\|\chi_{U_4(u)}f\|
       +4\sqrt3\|\chi_{U_4(u)}\nabla f\|.
\end{equation}
For $j=0$, the difference
\[
 \partial_\kappa\partial_\mu f_i(x)
  =f_i(x+e_\kappa+e_\mu)-f_i(x+e_\kappa)
                             -f_i(x+e_\mu)+f_i(x)
\]
uses four values in $U_3(u)$ when $x\in U_1(u)$.
Its square is bounded by four times the sum of their
squares. Sum over $x$, the components $i$, and the nine
pairs $(\kappa,\mu)$. Translation of each summation set
into $U_3(u)$ gives
$\|\chi_{U_1(u)}\nabla\nabla f\|^2
 \le144\|\chi_{U_3(u)}f\|^2$.

For $j\ge1$ we have $2\xi\le1$. Let $\phi$ be the
even function on $\mathbb R$ given, for $s\ge0$, by
\[
 \phi(s)=
 \begin{cases}
  1,&0\le s\le2,\\
  1-2(s-2)^2,&2\le s\le5/2,\\
  2(3-s)^2,&5/2\le s\le3,\\
  0,&s\ge3.
 \end{cases}
\]
It is continuously differentiable, $|\phi'|\le2$,
and $\phi'$ is Lipschitz with constant $4$.
For $t\in\mathbb R/N\mathbb Z$ write $|t|_N$ for its
distance from zero and put $q(t)=\phi(|t|_N)$.
The same derivative bounds hold for this periodic
function. Indeed, $q=1$ if $N=4$. For $N\ge6$ there is
no difficulty at zero, where $\phi$ is constant; at the
antipodal point $q$ vanishes in a neighborhood if
$N\ge8$, while for $N=6$ it is locally
$2(N/2-|t|_N)^2$. In the latter case its derivative
vanishes at the antipode and has Lipschitz constant $4$
across it.

Define
\[
 \Theta(x)=\prod_{\mu=1}^3q(x_\mu-u_\mu).
\]
Then $0\le\Theta\le1$, $\Theta=1$ on $U_2(u)$ and
$\Theta=0$ outside $U_3(u)$. The derivative bounds give
$|\partial_\mu\Theta|\le2$. Moreover,
\[
 \begin{split}
 |q(t+\xi)+q(t-\xi)-2q(t)|
  &=\left|\int_0^\xi(q'(t+s)-q'(t-s))\,ds\right|\\
  &\le\int_0^\xi8s\,ds=4\xi^2,
 \end{split}
\]
so $|\partial_\mu^*\partial_\mu\Theta|\le4$ and
$|\Delta\Theta|\le12$. Since $\xi\le1$, both
$\partial_\mu\Theta$ and $\Delta\Theta$ vanish
outside $U_4(u)$.

The second differences at a point of $U_1(u)$ use only
values in $U_2(u)$, where $\Theta=1$. Applying
componentwise the identity proved in
\cref{sec:one-order-coercivity} gives
\[
 \|\chi_{U_1(u)}\nabla\nabla f\|
    \le\|\nabla\nabla(\Theta f)\|
    =\|\Delta(\Theta f)\|.
\]
The product rule needed here is
\begin{equation}\label{I-15-139}
 \Delta(\Theta f)
  =\Theta\Delta f+(\Delta\Theta)f
       -\sum_{\mu=1}^3
           \bigl(w_\mu+w_\mu(\,\cdot-\xi e_\mu)\bigr),
 \qquad w_\mu=(\partial_\mu\Theta)\partial_\mu f.
\end{equation}
For completeness, if $e=\xi e_\mu$, expand
\[
 \begin{split}
 \Theta(x\pm e)f(x\pm e)
   &=\Theta(x)f(x\pm e)
       +(\Theta(x\pm e)-\Theta(x))f(x)\\
   &\quad+(\Theta(x\pm e)-\Theta(x))
                              (f(x\pm e)-f(x)).
 \end{split}
\]
The final product is $\xi^2w_\mu(x)$ for the upper
sign and $\xi^2w_\mu(x-e)$ for the lower sign.
Substitution in
$\xi^2\Delta F(x)=\sum_\mu
(2F(x)-F(x+\xi e_\mu)-F(x-\xi e_\mu))$
proves \eqref{I-15-139}, including both translated
terms. Now
\[
 \|w_\mu\|\le2\|\chi_{U_4(u)}\partial_\mu f\|,
 \qquad
 \|w_\mu(\,\cdot-\xi e_\mu)\|=\|w_\mu\|,
\]
and Cauchy--Schwarz gives
$\sum_\mu\|\chi_{U_4(u)}\partial_\mu f\|
 \le\sqrt3\|\chi_{U_4(u)}\nabla f\|$.
Taking norms in \eqref{I-15-139}, and using the support
and derivative bounds for $\Theta$, proves
\eqref{I-15-135}.

\begin{proof}[Proof of \cref{I-15-163}\textup{(I5)}]
We prove all six estimates with the constant $c_*$ in
\eqref{I-15-132}. Each $U_1(u)$ is nonempty: every
coordinate of a torus point lies within $\xi/2\le1$
of a fine site. Since $\zeta$ vanishes outside $U_1(u)$,
\[
 \|\zeta F\|\le|\zeta|\|\chi_{U_1(u)}F\|.
\]
Apply \eqref{I-15-133} and \eqref{I-15-138} with
$r=r'=1$. For instance,
\[
 \|\zeta G_jJ\|
    \le4\gamma_{\mathrm f}^{-1}e^{6\delta_{\mathrm f}}
       e^{-\delta_{\mathrm f}|u-u'|_1}|\zeta|\|J\|.
\]
The same calculation with the other choices of
$\mathcal O_1,\mathcal O_2$ proves the estimates for
$\zeta\nabla G_jJ$, $\zeta G_j\nabla^*J'$ and
$\zeta\nabla G_j\nabla^*J'$. Their constants are at
most $c_*$, and $\delta_{\mathrm f}\ge\delta_*$.

It remains to estimate two output derivatives and then
take the adjoint. Put $f=G_jJ$. For $w\in\mathsf T$
and $r>0$, the same first-order bounds give
\begin{equation}\label{I-15-140}
 \|\chi_{U_r(w)}f\|,\quad
 \|\chi_{U_r(w)}\nabla f\|
  \le4\gamma_{\mathrm f}^{-1}e^{3(r+1)}
                    e^{-\delta_{\mathrm f}|w-u'|_1}\|J\|.
\end{equation}
If $U_r(w)$ is empty the left sides vanish; otherwise
use \eqref{I-15-138} and $\delta_{\mathrm f}\le1$.
When $j=0$, the local estimate just proved yields
\[
 \|\chi_{U_1(u)}\nabla\nabla f\|
  \le48\gamma_{\mathrm f}^{-1}e^{12}
                  e^{-\delta_{\mathrm f}|u-u'|_1}\|J\|
  \le c_*e^{-\delta_*|u-u'|_1}\|J\|.
\]

Suppose now that $j\ge1$. The componentwise equation
is
\[
 \Delta f
    =J+\partial\mathsf P_j\partial^*f
                           -\mathsf Q_j^*\mathsf Q_jf.
\]
We estimate its three terms on $U_3(u)$. The source
term vanishes there unless $U_3(u)$ meets $U_1(u')$,
in which case $|u-u'|_1\le12$. Thus
\[
 \|\chi_{U_3(u)}J\|
             \le e^{12}e^{-\delta_*|u-u'|_1}\|J\|.
\]
Next, \eqref{eq:one-scale-adjoint-column} shows that
$(\mathsf Q_j^*B)(x)$ uses only coarse bonds $b$ for
which $x\in b_-+[0,2)^3$. The straight average
$(\mathsf Q_jf)(b)$ uses values $f(z)$ in the same
box. Hence $|z-x|_\infty<2$, and
\[
 \chi_{U_3(u)}\mathsf Q_j^*\mathsf Q_jf
   =\chi_{U_3(u)}\mathsf Q_j^*\mathsf Q_j
                                  \chi_{U_5(u)}f.
\]
Since $\|\mathsf Q_j\|\le1$ by \cref{I-12-8}(i),
\[
 \|\chi_{U_3(u)}\mathsf Q_j^*\mathsf Q_jf\|
                         \le\|\chi_{U_5(u)}f\|.
\]

Finally, the divergence is
\[
 (\partial^*f)(x)
      =-\sum_{\mu=1}^3
                 (\partial_\mu f_\mu)(x-\xi e_\mu).
\]
The translated block $B^j(y)-\xi e_\mu$ lies in
$U_1(y)$. Cauchy--Schwarz in the three components,
followed by translation of the sums and
\eqref{I-15-140}, gives
\[
 \begin{split}
 \|\chi_{B^j(y)}\partial^*f\|
   &\le\sqrt3\|\chi_{U_1(y)}\nabla f\|\\
   &\le4\sqrt3\gamma_{\mathrm f}^{-1}e^6
                  e^{-\delta_{\mathrm f}|y-u'|_1}\|J\|.
 \end{split}
\]
Since $2\delta_*\le\delta_{\mathrm f},\delta_P$,
\eqref{I-15-136} applies to $g=\partial^*f$ with
$\delta=\delta_*$ and
$A=4\sqrt3\gamma_{\mathrm f}^{-1}e^6\|J\|$.
At $r=3$ it gives
\[
 \|\chi_{U_3(u)}\partial\mathsf P_j\partial^*f\|
   \le4\sqrt3\gamma_{\mathrm f}^{-1}C_Pe^{18}
          \sigma(\delta_*)e^{-\delta_*|u-u'|_1}\|J\|.
\]
Insert these three estimates into \eqref{I-15-135},
and use \eqref{I-15-140} at radii $4$ and $5$ for
the remaining terms. This gives
\[
 \begin{split}
 &\|\chi_{U_1(u)}\nabla\nabla G_jJ\|\\
 &\quad\le
 \Bigl[e^{12}+4\gamma_{\mathrm f}^{-1}
      \bigl(e^{18}+\sqrt3C_Pe^{18}\sigma(\delta_*)
                      +12e^{15}+4\sqrt3e^{15}\bigr)\Bigr]
             e^{-\delta_*|u-u'|_1}\|J\|\\
 &\quad\le c_*e^{-\delta_*|u-u'|_1}\|J\|,
 \end{split}
\]
where the last inequality uses $e^{15}\le e^{18}$.
Multiplication by $|\zeta|$ proves the estimate with
two output derivatives.

The positive operator $G_j$ is symmetric. Therefore
\[
 \chi_{U_1(u)}G_j\nabla^*\nabla^*\chi_{U_1(u')}
   =\bigl(\chi_{U_1(u')}\nabla\nabla G_j
                                  \chi_{U_1(u)}\bigr)^*.
\]
Exchange $u$ and $u'$ in the bound just proved and
take the operator norm. This proves the final estimate
\[
 \|\zeta G_j\nabla^*\nabla^*J''\|
     \le c_*e^{-\delta_*|u-u'|_1}|\zeta|\|J''\|,
\]
and completes \textup{(I5)}.
\end{proof}

\subsection{H\"older divergence data}
\label{sec:holder-divergence-interior}

We strengthen the derivative estimate in \cref{I-13-47} by allowing
a H\"older continuous divergence source. We use the lattice,
differences and norms of \cref{sec:scalar-interior}, with
arbitrary spacing $\xi>0$, and write
\[
 \mathsf B_s(x_0)
    =\{x\in\xi\mathbb Z^3:|x-x_0|_\infty\le s\}.
\]
For a function $H$ on a set
$X\subset\xi\mathbb Z^3$, put
\[
 [H]_{\beta,r;X}
   =\sup_{\substack{z,z'\in X\\z\ne z'}}
      \left(\frac r{|z-z'|_1}\right)^\beta
         |H(z)-H(z')|,
 \qquad 0<\beta<1,\quad r>0,
\]
with value zero if there are no such pairs. Write
$[H]_{\beta,r}=[H]_{\beta,r;\xi\mathbb Z^3}$ for a function
defined on the whole lattice. For fields these seminorms and
the local supremum $|F|_X$ are the maxima over the components.
In particular, a seminorm with subscript $X$ makes no assertion
about the zero extension across the boundary of $X$.

\begin{proposition}\label{I-15-148}
For every $0<\beta<1$ there is a constant $c_{II}(\beta)>0$
with the following property. Let $r\ge32\xi$,
$x_0\in\xi\mathbb Z^3$ and
$\mathsf B_s=\mathsf B_s(x_0)$.
Let $u$ be a function on $\mathsf B_r$, and let $g$ and
$F=(F_\mu)_{\mu=1}^3$ be defined on $\xi\mathbb Z^3$, with
\[
 \Delta u=g+\nabla^*F
       \qquad\text{on }\mathsf B_{r-\xi}.
\]
Put
\[
 \mathsf N_\beta
    =r^{-3/2}\|u\|_{\mathsf B_r}
       +r^2|g|_{\mathsf B_r}
       +r\bigl(|F|_{\mathsf B_r}
                    +[F]_{\beta,r;\mathsf B_r}\bigr).
\]
For all $x,x'\in\mathsf B_{r/8}$ and every direction $\nu$,
\begin{equation}\label{I-15-149}
 \begin{split}
 r|\partial_\nu u(x)|
      &\le c_{II}(\beta)\mathsf N_\beta,\\
 r|\partial_\nu u(x)-\partial_\nu u(x')|
      &\le c_{II}(\beta)
          \left(\frac{|x-x'|_1}{r}\right)^\beta
             \mathsf N_\beta.
 \end{split}
\end{equation}
The constant is independent of $\xi,r,x_0,u,g,F$.
\end{proposition}

The proof, given after \cref{I-15-145}, subtracts the potential
of a cutoff of $F$. A derivative of this potential contains
two differences of the free kernel and therefore has a
critical singularity of order $-3$. The following estimate
uses its cancellation: summing a difference over a box leaves
two boundary faces. Subtracting the value of the source at
the observation point then gives a bound uniform in the mesh.

\begin{lemma}\label{I-15-145}
For every $0<\beta<1$ there is a constant $c_S(\beta)>0$
with the following property. Suppose that $k_1$ is a function
on $\xi\mathbb Z^3$ satisfying, for every $z$ and all
directions $\lambda,\lambda'$,
\[
 \begin{split}
 |k_1(z)|&\le(|z|_1+\xi)^{-2},\\
 |\partial_\lambda k_1(z)|&\le(|z|_1+\xi)^{-3},\\
 |\partial_{\lambda'}\partial_\lambda k_1(z)|
       &\le(|z|_1+\xi)^{-4}.
 \end{split}
\]
Let $k=\partial_\nu k_1$ for some direction $\nu$.
If $r\ge\xi$, $x_0\in\xi\mathbb Z^3$, and $G$ is a function
on the whole lattice vanishing outside $\mathsf B_r(x_0)$, then for all
$x,x'\in\mathsf B_r(x_0)$,
\begin{equation}\label{I-15-146}
 \begin{split}
 |(k*G)(x)|
    &\le c_S(\beta)\bigl([G]_{\beta,r}+|G|\bigr),\\
 |(k*G)(x)-(k*G)(x')|
    &\le c_S(\beta)
        \left(\frac{|x-x'|_1}{r}\right)^\beta
           \bigl([G]_{\beta,r}+|G|\bigr).
 \end{split}
\end{equation}
The constant depends only on $\beta$.
\end{lemma}

\begin{proof}
In this proof $c$ denotes constants depending only on $\beta$.
All sums are over $\xi\mathbb Z^3$ unless otherwise specified.

\emph{Cancellation on boxes.}
For $w\in\xi\mathbb Z^3$ and a positive integer multiple $s$
of $\xi$, set
\[
 \mathsf S_{s,w}(x)
     =\xi^3\sum_{z\in\mathsf B_s(w)}k(x-z).
\]
We claim that, if $x,x'\in\mathsf B_{s/2}(w)$, then
\begin{equation}\label{I-15-147}
 |\mathsf S_{s,w}(x)|\le72,\qquad
 |\mathsf S_{s,w}(x)-\mathsf S_{s,w}(x')|
       \le144\,\frac{|x-x'|_1}{s}.
\end{equation}
Fix the two coordinates of $z$ other than $z_\nu$, and let
$a_x(t)$ be $k_1(x-z)$ when $z_\nu=t$.
Since
$k(x-z)=\xi^{-1}\bigl(a_x(z_\nu-\xi)-a_x(z_\nu)\bigr)$,
the sum in the $\nu$ direction telescopes:
\[
 \sum_{t=w_\nu-s,w_\nu-s+\xi,\ldots,w_\nu+s}
        \bigl(a_x(t-\xi)-a_x(t)\bigr)
       =a_x(w_\nu-s-\xi)-a_x(w_\nu+s).
\]
Consequently
\[
 \mathsf S_{s,w}(x)
    =\xi^2\sum_{(z_\mu)_{\mu\ne\nu}}
       \bigl(a_x(w_\nu-s-\xi)-a_x(w_\nu+s)\bigr).
\]
There are $(2s/\xi+1)^2\le(3s/\xi)^2$ terms in this
remaining sum. At either face the $\nu$ coordinate of
$x-z$ has modulus at least $s/2$; the lower face is at
$z_\nu=w_\nu-s-\xi$. Thus
\[
 |\mathsf S_{s,w}(x)|
    \le\xi^2(3s/\xi)^2\,2(s/2)^{-2}=72.
\]

Join $x$ to $x'$ by a shortest lattice path whose coordinates
change monotonically. The entire path lies in
$\mathsf B_{s/2}(w)$. A step changes a boundary value of
$k_1$ by $\pm\xi\partial_\lambda k_1$, evaluated at one
of its endpoints. The distance from that endpoint to either
face is still at least $s/2$. Summing over the
$|x-x'|_1/\xi$ steps therefore gives
\[
 |\mathsf S_{s,w}(x)-\mathsf S_{s,w}(x')|
     \le\xi^2(3s/\xi)^2\,2|x-x'|_1(s/2)^{-3}
     =144\,\frac{|x-x'|_1}{s}.
\]
This proves \eqref{I-15-147}.

\emph{The two shell sums.}
For every $s>0$ we shall use
\[
 \begin{split}
 \xi^3\sum_{0<|z|_1\le s}|z|_1^{\beta-3}
       &\le\frac{216}{2^\beta-1}s^\beta,\\
 \xi^3\sum_{|z|_1>s}|z|_1^{\beta-4}
       &\le\frac{216}{1-2^{\beta-1}}s^{\beta-1}.
 \end{split}
\]
Indeed, the shell $\rho<|z|_1\le2\rho$ has at most
$(6\rho/\xi)^3$ points if $2\rho\ge\xi$, and is empty
otherwise. Its contributions to the two sums are at most
$216\rho^\beta$ and $216\rho^{\beta-1}$, respectively.
For the first sum take $\rho=2^{-i-1}s$, $i\ge0$; for
the second take $\rho=2^is$. The resulting geometric series
have sums $s^\beta/(2^\beta-1)$ and
$s^{\beta-1}/(1-2^{\beta-1})$.

\emph{The pointwise bound.}
Let $\hat r$ be the least multiple of $\xi$ not smaller
than $r$. Then $r\le\hat r<r+\xi\le2r$.
Set $s=2\hat r$. For $x\in\mathsf B_r(x_0)$ the support
of $G$ is contained in $\mathsf B_s(x)$, so
\[
 (k*G)(x)
   =\xi^3\sum_{\substack{z\in\mathsf B_s(x)\\z\ne x}}
        k(x-z)\bigl(G(z)-G(x)\bigr)
       +G(x)\mathsf S_{s,x}(x).
\]
The first term is bounded by
\[
 \begin{split}
 [G]_{\beta,r}r^{-\beta}
      \xi^3\sum_{0<|z-x|_1\le12r}|z-x|_1^{\beta-3}
   &\le\frac{216\,12^\beta}{2^\beta-1}
                          [G]_{\beta,r},
 \end{split}
\]
because $|z-x|_1\le3s\le12r$ on the box.
The second is at most $72|G|$ by \eqref{I-15-147}.
This proves the first estimate in \eqref{I-15-146}, with
\[
 c_1=\max\left\{\frac{216\,12^\beta}{2^\beta-1},72\right\}.
\]

\emph{The H\"older bound.}
Put $\delta=|x-x'|_1$ and suppose $\delta>0$.
If $\delta>r/4$, the pointwise estimate at $x$ and $x'$
gives
\[
 |(k*G)(x)-(k*G)(x')|
     \le2c_1\,4^\beta(\delta/r)^\beta
                   \bigl([G]_{\beta,r}+|G|\bigr).
\]
We may therefore assume $\delta\le r/4$.
Set $s_1=2\delta$ and $s_2=4\hat r$.
Both are positive integer multiples of $\xi$, and
$x'\in\mathsf B_{s_1/2}(x)\subset\mathsf B_{s_2/2}(x)$.
Subtracting $G(x)$ on the common box gives
\[
 \begin{split}
 (k*G)(x)-(k*G)(x')
    ={}&\xi^3\sum_{z\in\mathsf B_{s_2}(x)}
       \bigl(k(x-z)-k(x'-z)\bigr)\bigl(G(z)-G(x)\bigr)\\
     &+G(x)\bigl(\mathsf S_{s_2,x}(x)
                         -\mathsf S_{s_2,x}(x')\bigr).
 \end{split}
\]
Since $s_2\ge4r$, the last term is at most
$36|G|\delta/r\le36|G|(\delta/r)^\beta$.

For $z\notin\mathsf B_{s_1}(x)$ we have
$|z-x|_1>2\delta$. Every point $x''$ on a shortest path
from $x$ to $x'$ satisfies
$|x''-z|_1\ge|x-z|_1-\delta\ge|x-z|_1/2$.
Each step changes $k(\,\cdot-z)$ by a first difference
evaluated at $x''-z$ for one of its endpoints.
The second-difference bound for $k_1$ therefore gives
\[
 |k(x-z)-k(x'-z)|
        \le16\delta\,|x-z|_1^{-4}.
\]
The far part of the sum is consequently bounded by
\[
 \begin{split}
 16\delta [G]_{\beta,r}r^{-\beta}
       \xi^3\sum_{|z-x|_1>2\delta}|z-x|_1^{\beta-4}
   &\le\frac{16\cdot216\,2^{\beta-1}}
                 {1-2^{\beta-1}}
           [G]_{\beta,r}(\delta/r)^\beta.
 \end{split}
\]

For the near part, recenter the terms with $k(x'-z)$
at $G(x')$. This part is $\mathsf n-\mathsf n'$, where
\[
 \begin{split}
 \mathsf n
    &=\xi^3\sum_{\substack{z\in\mathsf B_{s_1}(x)\\z\ne x}}
           k(x-z)\bigl(G(z)-G(x)\bigr),\\
 \mathsf n'
    &=\xi^3\sum_{\substack{z\in\mathsf B_{s_1}(x)\\z\ne x'}}
           k(x'-z)\bigl(G(z)-G(x')\bigr)
        +\bigl(G(x')-G(x)\bigr)\mathsf S_{s_1,x}(x').
 \end{split}
\]
Here $|z-x|_1\le6\delta$ and
$|z-x'|_1\le7\delta\le9\delta$.
The first shell sum gives
\[
 |\mathsf n|
     \le\frac{216\,6^\beta}{2^\beta-1}
                       [G]_{\beta,r}(\delta/r)^\beta;
\]
the sum in $\mathsf n'$ has the same bound with $9^\beta$
in place of $6^\beta$.
Finally, $x'\in\mathsf B_{s_1/2}(x)$, so
\eqref{I-15-147} bounds the remaining term of
$\mathsf n'$ by
$72[G]_{\beta,r}(\delta/r)^\beta$.
Combining the near, far and boundary terms proves the
second estimate in \eqref{I-15-146}. Enlarging the constant
to cover both estimates gives $c_S(\beta)$.
\end{proof}

\begin{proof}[Proof of \cref{I-15-148}]
The constants denoted by $c$ below depend only on $\beta$,
the constants in \cref{I-13-45,I-13-47,I-15-145}, and the
fixed cutoff $\eta_0$ of \cref{sec:scalar-interior-proof}.
In particular, they are independent of the mesh and radius.

\emph{Cutting off the divergence source.}
Set
\[
 \zeta(x)=\prod_{\mu=1}^3
          \eta_0\bigl((x_\mu-x_{0,\mu})/r\bigr).
\]
Then $0\le\zeta\le1$, $\zeta=1$ on $\mathsf B_{r/2}$,
and $\zeta$ vanishes outside $\mathsf B_{2r/3}$.
The product inequality for factors in $[0,1]$ and the
mean value theorem give
\[
 |\zeta(z)-\zeta(z')|
       \le\|\eta_0'\|_\infty\,|z-z'|_1/r.
\]
Define $G_\mu=\zeta F_\mu$ on $\mathsf B_r$ and
$G_\mu=0$ elsewhere. This is a finitely supported field,
with $|G_\mu|\le|F|_{\mathsf B_r}$ and
\begin{equation}\label{I-15-150}
 [G_\mu]_{\beta,r}
     \le[F]_{\beta,r;\mathsf B_r}
           +\max\{3,\|\eta_0'\|_\infty\}|F|_{\mathsf B_r}.
\end{equation}
To verify the global seminorm in this estimate, put
$t=|z-z'|_1/r$.
If both points lie in $\mathsf B_r$, then
\[
 |G_\mu(z)-G_\mu(z')|
   \le |F_\mu(z)-F_\mu(z')|
          +|\zeta(z)-\zeta(z')|\,|F|_{\mathsf B_r},
\]
and
\[
 |\zeta(z)-\zeta(z')|
    \le\min\{1,\|\eta_0'\|_\infty t\}
    \le\max\{1,\|\eta_0'\|_\infty\}t^\beta.
\]
If one point lies in $\mathsf B_{2r/3}$ and the other
outside $\mathsf B_r$, then $t>1/3$, so the difference
of the two values of $G_\mu$ is at most
$|F|_{\mathsf B_r}\le3^\beta t^\beta|F|_{\mathsf B_r}$.
In every remaining case not already covered, both values
vanish. Dividing by $t^\beta$ proves \eqref{I-15-150}.

\emph{The potential of the cutoff field.}
Put
\[
 u_1=\sum_{\mu=1}^3
          (\partial_\mu^*C_{0,\mathrm{lap}})*G_\mu.
\]
For finitely supported fields, forward and adjoint differences
commute with convolution by shifting the summation index, as in
\cref{sec:scalar-interior-proof}. Hence \cref{I-13-45}(i) gives
$\Delta u_1=\sum_\mu\partial_\mu^*G_\mu$.
For $x\in\mathsf B_{r/2-\xi}$, the points
$x$ and $x-\xi e_\mu$ lie in $\mathsf B_{r/2}$,
where $G_\mu=F_\mu$. Consequently $u_2=u-u_1$ satisfies
\[
 \Delta u_2=g
        \qquad\text{on }\mathsf B_{r/2-\xi}.
\]

Fix $\mu$ and take
$k_1=c_K^{-1}\partial_\mu^*C_{0,\mathrm{lap}}$.
The mixed differences of orders one, two and three in
\cref{I-13-45}(ii) show that $k_1$ satisfies all three
hypotheses of \cref{I-15-145}, and the kernel hypotheses
of \cref{I-13-45}(iv).
For $x\in\mathsf B_r$, the support of $G_\mu$ lies
within $\ell^1$ distance $6r$ of $x$.
The pointwise estimate in part~(iv) therefore gives
\[
 |u_1(x)|\le3c_K^2r\,|F|_{\mathsf B_r}
                 \qquad(x\in\mathsf B_r).
\]
There are at most $(2r/\xi)^3$ points in
$\mathsf B_{r/2}$, since $r\ge\xi$. Hence
\[
 \|u_1\|_{\mathsf B_{r/2}}
      \le(2r)^{3/2}\,3c_K^2r\,|F|_{\mathsf B_r}.
\]

Moreover,
\[
 \partial_\nu u_1
     =\sum_\mu
          (\partial_\nu\partial_\mu^*
                       C_{0,\mathrm{lap}})*G_\mu,
 \qquad
 \partial_\nu\partial_\mu^*C_{0,\mathrm{lap}}
          =c_K\partial_\nu k_1.
\]
Apply \cref{I-15-145} to this $k_1$ and each $G_\mu$,
and then use \eqref{I-15-150}. For $x,x'\in\mathsf B_r$
we obtain
\[
 \begin{split}
 r|\partial_\nu u_1(x)|
    &\le cr\bigl([F]_{\beta,r;\mathsf B_r}
                              +|F|_{\mathsf B_r}\bigr),\\
 r|\partial_\nu u_1(x)-\partial_\nu u_1(x')|
    &\le c\left(\frac{|x-x'|_1}{r}\right)^\beta
          r\bigl([F]_{\beta,r;\mathsf B_r}
                              +|F|_{\mathsf B_r}\bigr).
 \end{split}
\]

\emph{The remaining scalar equation.}
Apply \cref{I-13-47}(b) to $u_2$ on
$\mathsf B_{r/2}$, with right-hand side $g$ and zero
divergence field. Its radius condition holds because
$r/2\ge16\xi$, and its equation was verified above.
Its controlling quantity satisfies
\[
 \begin{split}
 \mathsf N
    &=(r/2)^{-3/2}\|u_2\|_{\mathsf B_{r/2}}
                       +(r/2)^2|g|_{\mathsf B_{r/2}}\\
    &\le 2^{3/2}r^{-3/2}
          \bigl(\|u\|_{\mathsf B_r}
                     +\|u_1\|_{\mathsf B_{r/2}}\bigr)
                      +r^2|g|_{\mathsf B_r}\\
    &\le c\mathsf N_\beta.
 \end{split}
\]
Thus, for $x,x'\in\mathsf B_{r/8}$,
\[
 \begin{split}
 (r/2)|\partial_\nu u_2(x)|
    &\le c_I c\mathsf N_\beta,\\
 (r/2)|\partial_\nu u_2(x)-\partial_\nu u_2(x')|
    &\le c_I(\beta)c
          \left(\frac{2|x-x'|_1}{r}\right)^\beta
                    \mathsf N_\beta.
 \end{split}
\]
At these sites the identity $u=u_1+u_2$ also holds at
the forward-shifted sites used by $\partial_\nu$.
Adding the estimates for the two potentials proves
\eqref{I-15-149}, after increasing $c_{II}(\beta)$.
\end{proof}

\subsection{Lifts and pointwise projection estimates}
\label{sec:one-scale-pointwise}

The scalar interior estimates apply on the covering lattice.
The balls of radius $16$ and $32$ used below may cover the
torus more than once. We therefore count the lifted blocks,
retaining repeated images on the torus.
Let $\mathsf p:\xi\mathbb Z^3\to\mathsf T_\xi$ be the covering
map and write $f^\uparrow=f\circ\mathsf p$. Differences commute
with this lift. Fix $w\in\mathsf T_\xi$, a lift $w_0$, and an
integer $r\ge1$. The set
\[
 \mathsf Y=
 \{\widetilde y\in\mathbb Z^3:
       w_{0,\mu}-r-1<\widetilde y_\mu\le w_{0,\mu}+r
           \text{ for }\mu=1,2,3\}
\]
has $(2r+1)^3$ elements. Its unit boxes cover
$\mathsf B_r(w_0)$, and
$|\mathsf p\widetilde y-w|_1\le3(r+1)$.
Each lifted block maps bijectively to
$B^j(\mathsf p\widetilde y)$, even when different lifted
blocks have the same image. Consequently,
\[
 \begin{split}
 \|f^\uparrow\|_{\mathsf B_r(w_0)}
   &\le\sum_{\widetilde y\in\mathsf Y}
            \|\chi_{B^j(\mathsf p\widetilde y)}f\|,\\
 |f^\uparrow|_{\mathsf B_r(w_0)}
   &\le\max_{\widetilde y\in\mathsf Y}
            |f|_{B^j(\mathsf p\widetilde y)}.
 \end{split}
\]
The first inequality follows by summing the squared block
norms and then bounding the square root by their sum.
Thus, if $0<\delta\le1$ and
$\|\chi_{B^j(y)}f\|\le A e^{-\delta|y-u'|_1}$ for every
coarse site $y$, then
\begin{equation}\label{eq:one-scale-lift-L2}
 \|f^\uparrow\|_{\mathsf B_r(w_0)}
       \le(2r+1)^3e^{3(r+1)}A e^{-\delta|w-u'|_1}.
\end{equation}
If instead $|f(x)|\le A e^{-\delta|x-u'|_1}$ at every
fine site, then
\begin{equation}\label{eq:one-scale-lift-sup}
 |f^\uparrow|_{\mathsf B_r(w_0)}
       \le e^{3r}A e^{-\delta|w-u'|_1}.
\end{equation}
Both estimates follow from the triangle inequality and
$\delta\le1$.

We will also compare nearby pairs. For
$z,z'\in\xi\mathbb Z^3$ with $|z-z'|_1\le1$, every nonzero
period translate of $z-z'$ has norm at least
$N-|z-z'|_1\ge3$. Hence
$|\mathsf pz-\mathsf pz'|_1=|z-z'|_1$.
Conversely, after fixing a lift $x_0$ of a torus site $x$,
choose a least-length representative of $x'-x$ to obtain a
lift $x'_0$ satisfying
$|x_0-x'_0|_1=|x-x'|_1$.

The Laplacian of $\mathsf P_jg$ is constant on each unit
block. The block estimates proved above therefore give
the bounded right-hand side needed in \cref{I-13-47}(b).

\begin{lemma}\label{I-15-151}
Let $A\ge0$, $u'\in\mathsf T$, and suppose that the real
scalar function $g$ satisfies
\[
 \|\chi_{B^j(y)}g\|
       \le A e^{-2\delta_*|y-u'|_1}
       \qquad(y\in\mathsf T_1).
\]
There are an absolute constant $c_{P1}$ and constants
$c_{P1}(\alpha)$ depending only on $0<\alpha<1$ such that
\begin{equation}\label{I-15-152}
 \begin{split}
 |(\partial\mathsf P_jg)(x)|
    &\le c_{P1}A e^{-\delta_*|x-u'|_1},\\
 |(\partial\mathsf P_jg)(x)-(\partial\mathsf P_jg)(x')|
    &\le c_{P1}(\alpha)|x-x'|_1^\alpha
                   A e^{-\delta_*|x-u'|_1}
 \end{split}
\end{equation}
whenever $x,x'\in\mathsf T_\xi$ and $|x-x'|_1\le1$.
\end{lemma}

\begin{proof}
Put $h=\mathsf P_jg$ and
$A_1=C_Pe^6\sigma(\delta_*)A$.
The block estimates \eqref{I-15-136} give
\[
 \|\chi_{B^j(y)}h\|\le A_1e^{-\delta_*|y-u'|_1},
 \qquad
 |\Delta h|_{B^j(y)}\le3A_1e^{-\delta_*|y-u'|_1}.
\]
Here $2\delta_*\le\delta_P$, as required in that estimate.
Fix a lift $x_0$ of $x$. The lifting bounds with $r=16$
yield
\[
 \begin{split}
 \|h^\uparrow\|_{\mathsf B_{16}(x_0)}
      &\le33^3e^{51}A_1e^{-\delta_*|x-u'|_1},\\
 |\Delta h^\uparrow|_{\mathsf B_{16}(x_0)}
      &\le3e^{51}A_1e^{-\delta_*|x-u'|_1}.
 \end{split}
\]
For the second bound, use the block supremum and
$|\mathsf p\widetilde y-x|_1\le51$ in the covering above.
Apply \cref{I-13-47}(b) to $h^\uparrow$, with right-hand
side $\Delta h^\uparrow$ and zero divergence field.
The radius is admissible since $\xi\le1$, and its data norm is
\[
 \mathsf N\le(33^3/64+768)e^{51}A_1
                          e^{-\delta_*|x-u'|_1}.
\]
Choose a lift $x'_0$ at distance $|x-x'|_1\le1$ from
$x_0$. Both points lie in $\mathsf B_4(x_0)$.
The derivative estimates of that proposition give
\eqref{I-15-152}, with $c_{P1}$ absolute and
$c_{P1}(\alpha)$ depending only on $\alpha$.
\end{proof}

\subsection{Pointwise and H\"older bounds for the inverse}
\label{sec:one-scale-regularity-proof}

\begin{proof}[Proof of the pointwise and H\"older conclusions of
\cref{I-15-163}]
Throughout the proof, $c$ denotes a positive absolute constant
which may increase, and $c(\beta)$ may also depend on the
displayed exponent. Fix the source center $u'$ and put
\[
 E(x)=e^{-\delta_*|x-u'|_1}.
\]
Since $\delta_*\le1$, we have
$E(z)\le e^{|z-x|_1}E(x)$.
Let $J$ and $J'$ vanish outside $U_1(u')$, with the ranks
specified in the proposition. The weighted volume of
$U_1(u')$ is at most $27$: each coordinate interval contains
at most $2/\xi+1\le3/\xi$ sites. Thus
\[
 \|J\|\le9|J|,\qquad \|J'\|\le9\sqrt3\,|J'|.
\]
Consider either
\[
 (f,n)=(G_jJ,|J|)
       \quad\text{or}\quad
 (f,n)=(G_j\nabla^*J',|J'|).
\]
The first-order decay estimate \cref{I-12-29}(c) and
$\operatorname{dist}(U_1(v),U_1(u'))
\ge|v-u'|_1-6$ imply
\begin{equation}\label{I-15-154}
 \|\chi_{U_1(v)}f\|,\quad
 \|\chi_{U_1(v)}\nabla f\|
       \le A_0e^{-\delta_{\mathrm f}|v-u'|_1},
 \qquad
 A_0=36\sqrt3\,\gamma_{\mathrm f}^{-1}e^6n .
\end{equation}
This estimate holds for every continuous center $v$.

\emph{The componentwise equation.}
Write $J^\natural=J$ or $\nabla^*J'$ according to the chosen
pair. The definition of the inverse gives
\begin{equation}\label{I-15-155}
 \Delta f=J^\natural+h-k,
 \qquad
 h=\partial\mathsf P_j\partial^*f,
 \qquad k=\mathsf Q_j^*\mathsf Q_jf.
\end{equation}
Both lower-order terms have pointwise bounds with the
required decay. Indeed,
$\partial^*f(x)=-\sum_\mu\partial_\mu f_\mu(x-\xi e_\mu)$
and $B^j(y)-\xi e_\mu\subset U_1(y)$, so
\[
 \|\chi_{B^j(y)}\partial^*f\|
       \le\sqrt3\,\|\chi_{U_1(y)}\nabla f\|
       \le\sqrt3 A_0e^{-2\delta_*|y-u'|_1}.
\]
We used $2\delta_*\le\delta_{\mathrm f}$ in the last
inequality. Apply \cref{I-15-151} to obtain
\begin{equation}\label{I-15-156}
 |h(x)|\le c nE(x).
\end{equation}

For the averaging term, let $x\in B^j(y)$.
The adjoint formula \eqref{eq:one-scale-adjoint-column}
shows that only the two coarse bonds of direction $\mu$
with initial points $y$ and $y-e_\mu$ contribute to
$(\mathsf Q_j^*B)_\mu(x)$, each with coefficient in $[0,1]$.
In fact, $x-t\xi e_\mu$ lies in one of those two blocks
for $0\le t<L^j$.
For either initial point $y''$, the column
$\mathsf Q_j^*e_b$ has norm at most one and is supported in
$y''+[0,2)^3\subset U_1(y''+\mathbf e)$, where
$\mathbf e=(1,1,1)$. Hence
\[
 |(\mathsf Q_jf)(b)|
   =|\langle\mathsf Q_j^*e_b,f\rangle|
   \le A_0 e^{-\delta_{\mathrm f}|y''+\mathbf e-u'|_1}.
\]
For these two choices $|x-y''-\mathbf e|_1\le3$.
Since $\delta_*\le\delta_{\mathrm f}$, summing their
contributions gives
\begin{equation}\label{I-15-157}
 |k(x)|\le2e^3A_0E(x)\le c nE(x).
\end{equation}

\emph{An ordinary source.}
Suppose $f=G_jJ$.
Where $J(x)\ne0$, the distance $|x-u'|_1$ is at most
three, and $|J(x)|\le e^3|J|E(x)$.
Together with \eqref{I-15-155}--\eqref{I-15-157}, this gives
\begin{equation}\label{I-15-158}
 |(\Delta G_jJ)(x)|\le c|J|E(x).
\end{equation}
Fix a lift $x_0$ of $x$ and a component $f_\mu$.
The block bound obtained from \eqref{I-15-154} and
\eqref{eq:one-scale-lift-L2} gives
\[
 \|f_\mu^\uparrow\|_{\mathsf B_{16}(x_0)}
       \le33^3e^{51}A_0E(x)\le c|J|E(x).
\]
Equation \eqref{I-15-158} and
\eqref{eq:one-scale-lift-sup} bound
$|\Delta f_\mu^\uparrow|_{\mathsf B_{16}(x_0)}$
by the same expression with a larger constant.
Apply both parts of \cref{I-13-47} at radius $16$,
with this right-hand side and zero divergence field.
For $|x-x'|_1\le1$, choose the nearest lift $x'_0$;
both lifted points lie in $\mathsf B_4(x_0)$.
The resulting estimates are
\begin{equation}\label{I-15-159}
 \begin{split}
 |(G_jJ)(x)|,\quad |(\nabla G_jJ)(x)|
      &\le c|J|E(x),\\
 |(\nabla G_jJ)(x)-(\nabla G_jJ)(x')|
      &\le c(\alpha)|x-x'|_1^\alpha|J|E(x),
             \qquad 0<\alpha<1.
 \end{split}
\end{equation}

\emph{A divergence source.}
Now let $f=G_j\nabla^*J'$ and fix a component $\nu$.
For a lift $x_0$ of $x$ and $r\in\{16,32\}$, define
$F_\mu=(J'_{\mu\nu})^\uparrow$ on $\mathsf B_r(x_0)$
and $F_\mu=0$ outside that ball. Put
$g=(h-k)_\nu^\uparrow$.
For $z\in\mathsf B_{r-\xi}(x_0)$, both $z$ and
$z-\xi e_\mu$ lie in $\mathsf B_r(x_0)$. Therefore
\[
 \Delta f_\nu^\uparrow=g+\nabla^*F
                  \quad\text{on }\mathsf B_{r-\xi}(x_0).
\]
If $F$ is nonzero, some $z\in\mathsf B_r(x_0)$ has
$\mathsf pz\in U_1(u')$, which implies
$|x-u'|_1\le3(r+1)$. It follows, also when $F=0$, that
\[
 |F|\le e^{3(r+1)}|J'|E(x).
\]
The bounds \eqref{I-15-154}, \eqref{I-15-156},
\eqref{I-15-157} and the lifting estimates give
\[
 \|f_\nu^\uparrow\|_{\mathsf B_r(x_0)}
            +|g|_{\mathsf B_r(x_0)}\le c|J'|E(x)
               \qquad(r=16,32).
\]
At radius $16$, \cref{I-13-47}(a) applies with the global
bound on this restricted field $F$. Comparing nearest
lifts as before, it gives
\begin{equation}\label{I-15-160}
 \begin{split}
 |(G_j\nabla^*J')(x)|&\le c|J'|E(x),\\
 |(G_j\nabla^*J')(x)-(G_j\nabla^*J')(x')|
     &\le c(\alpha)|x-x'|_1^\alpha|J'|E(x)
                \qquad(0<\alpha<1).
 \end{split}
\end{equation}

To estimate the gradient, take $r=32$ and $0<\beta<1$.
For two sites of $\mathsf B_{32}(x_0)$ at distance at most
one, their torus distance equals their lifted distance.
The corresponding difference of $F$ is therefore at most
$|z-z'|_1^\beta\|J'\|_\beta$.
At distances greater than one it is at most
$2|J'|\le2|z-z'|_1^\beta|J'|$.
Consequently,
\[
 [F]_{\beta,32;\mathsf B_{32}(x_0)}
       \le32^\beta(\|J'\|_\beta+2|J'|).
\]
This seminorm vanishes when $F=0$. In the other case the
preceding support argument supplies the decay factor, so
\[
 [F]_{\beta,32;\mathsf B_{32}(x_0)}
       \le32e^{99}(\|J'\|_\beta+2|J'|)E(x).
\]
We may now apply \cref{I-15-148}, since $32\ge32\xi$.
Its data norm is at most
$c(\beta)(\|J'\|_\beta+|J'|)E(x)$.
The nearest lifts of $x,x'$ belong to
$\mathsf B_4(x_0)=\mathsf B_{32/8}(x_0)$, and hence
\begin{equation}\label{I-15-161}
 \begin{split}
 |(\nabla G_j\nabla^*J')(x)|
     &\le c(\beta)(\|J'\|_\beta+|J'|)E(x),\\
 |(\nabla G_j\nabla^*J')(x)
                 -(\nabla G_j\nabla^*J')(x')|
     &\le c(\beta)|x-x'|_1^\beta
                  (\|J'\|_\beta+|J'|)E(x).
 \end{split}
\end{equation}

\emph{Multiplication by the observation cutoff.}
Let $\zeta$ vanish outside $U_1(u)$ and put
$E=e^{-\delta_*|u-u'|_1}$.
For $x\in U_2(u)$, we have $E(x)\le e^6E$.
The pointwise conclusions follow from
\eqref{I-15-158}--\eqref{I-15-161}, using
$\beta=\varepsilon$ for the gradient with divergence input.
For the H\"older conclusions, let $q$ denote, in turn,
$\nabla G_jJ$, $G_j\nabla^*J'$ or
$\nabla G_j\nabla^*J'$, and put
\[
 m=|J|,\quad |J'|,\quad
       \|J'\|_{\alpha+\varepsilon}+|J'|,
\]
respectively. For $0<\alpha<1$, the preceding estimates give
\[
 |q(x)|\le c_0mE,\qquad
 |q(x)-q(x')|\le c_0|x-x'|_1^\alpha mE
\]
when $x\in U_2(u)$ and $|x-x'|_1\le1$.
For the third choice use $\beta=\alpha+\varepsilon<1$
and $|x-x'|_1^{\alpha+\varepsilon}\le|x-x'|_1^\alpha$.
The constant $c_0$ depends only on the indicated exponents.

If neither point lies in $U_1(u)$, the difference of
$\zeta q$ vanishes. Otherwise both lie in $U_2(u)$, and
\[
 \begin{split}
 |\zeta(x)q(x)-\zeta(x')q(x')|
   &\le|\zeta(x)-\zeta(x')|\,|q(x)|
             +|\zeta(x')|\,|q(x)-q(x')|\\
   &\le c_0|x-x'|_1^\alpha
                    (\|\zeta\|_\alpha+|\zeta|)mE.
 \end{split}
\]
This proves the asserted seminorm bounds for $\alpha>0$,
including pairs crossing the support boundary.
For $\alpha=0$ use
$\|\zeta q\|_0\le2|\zeta|\sup_{U_1(u)}|q|$;
in the third case take $\beta=\varepsilon$ in
\eqref{I-15-161}.
Together with the six $L^2$ bounds proved in
\cref{sec:one-scale-L2}, these estimates prove
\cref{I-15-163}, after enlarging its common constant $c$
to include the $L^2$ constant.
\end{proof}

\section{Two-scale propagators}
\label{sec:two-scale-propagators}

Retain the real, three-dimensional setting of
\cref{sec:one-scale-regularity}: $L\ge2$ and $j\ge0$ are
integers, $\xi=L^{-j}$, and the torus side $N$ is a positive
even multiple of $L$. In addition to $\mathsf T_\xi$ and
$\mathsf T_1$, use $\mathsf T_L=L\mathbb Z^3/N\mathbb Z^3$.
Write $\mathcal F_h$ for the real scalar functions on
$\mathsf T_h$ and $\mathcal A_h=\mathcal F_h^3$ for its
one-forms. Their inner products have weight $h^3$.
The fine differences are $\partial,\nabla,\Delta$.
At spacing $1$ or $L$, the first differences carry a
subscript $1$ or $L$. The notation $\Delta^1$ denotes
the componentwise Laplacian at spacing $1$;
$\Delta_1$ below denotes the two-scale operator on the
fine lattice. Distances, component maximum norms,
$U_r(u)$ and the seminorms $\|\cdot\|_\alpha$ retain
their preceding meanings.

For $y'\in\mathsf T_L$ let
\[
 B(y')=\mathsf T_1\cap(y'+[0,L)^3).
\]
The averages from spacing $1$ to spacing $L$ are
\[
 \begin{split}
  (\mathsf Q'_{1,L}g)(y')
    &=L^{-3}\sum_{y\in B(y')}g(y),\\
  (\mathsf Q_{1,L}B)_\mu(y')
    &=L^{-4}\sum_{y\in B(y')}\sum_{t=0}^{L-1}
                                  B_\mu(y+te_\mu).
 \end{split}
\]
The two indices distinguish this step from
$\mathsf Q_j:\mathcal A_\xi\to\mathcal A_1$.
Both block families partition their respective fine
lattices. Telescoping the straight paths gives
\[
 \mathsf Q_j\partial=\partial_1\mathsf Q'_j,
 \qquad
 \mathsf Q_{1,L}\partial_1
                         =\partial_L\mathsf Q'_{1,L}.
\]
In each $B(y')$ use the coordinate tree $\mathcal T(y')$
whose paths from $y'$ change coordinates in the order
$3,2,1$. Thus its direction-$\mu$ bonds have their
coordinates of index less than $\mu$ equal to those of
$y'$. Every path lies inside its block.

\begin{definition}\label{I-15-117}
Let $\Lambda'\subset\mathsf T_L$ be arbitrary and put
$\Lambda=\bigcup_{y'\in\Lambda'}B(y')$.
Let $\Lambda^*$ consist of the unit bonds with at least
one endpoint in $\Lambda$, and let $\Lambda^{c*}$ consist
of those with neither endpoint there.
Let $\Lambda'^*$ consist of the spacing-$L$ bonds with
at least one endpoint in $\Lambda'$; it includes bonds
leaving $\Lambda'$.
Define
\[
 \begin{split}
  N(\mathsf Q')=\{\lambda\in\mathcal F_\xi:\ &
       \mathsf Q'_j\lambda=0
                    \text{ on }\mathsf T_1\setminus\Lambda,\\[-2pt]
       &\mathsf Q'_{1,L}\mathsf Q'_j\lambda=0
                    \text{ on }\Lambda'\}.
 \end{split}
\]
Let $\mathsf R$ be the orthogonal projection onto
$\Delta N(\mathsf Q')$, and put $\mathsf P=1-\mathsf R$.
For $B\in\mathcal A_1$ set
\[
 \mathsf q(B)
   =\sum_{b\in\Lambda^{c*}}|B(b)|^2
        +L\sum_{c\in\Lambda'^*}
                      |(\mathsf Q_{1,L}B)(c)|^2 .
\]
The symmetric operator $\Delta_1$ on $\mathcal A_\xi$
is defined by
\begin{equation}\label{I-15-118}
 \langle A,\Delta_1A\rangle
   =\|\partial A\|^2+\|\mathsf R\partial^*A\|^2
                            +\mathsf q(\mathsf Q_jA).
\end{equation}
\end{definition}

\begin{proposition}\label{I-16-14}
The operator $\Delta_1$ is positive definite.
Its inverse $G=\Delta_1^{-1}$ satisfies every estimate
\textup{(I1)--(I5)} of \cref{I-15-163}, with $G$ in
place of $G_j$ and
$E=e^{-\delta_{\mathrm{per}}|u-u'|_1}$.
Here $\delta_{\mathrm{per}}>0$ depends only on $L$;
the constants $c$, $c_\alpha$, $c'_\varepsilon$ and
$c''_{\alpha,\varepsilon}$ depend only on $L$ and
their displayed exponents. They are independent of
$j,N,\Lambda'$, the continuous centers and the fields.
The support conditions, tensor source spaces and exponent
ranges are exactly those in \cref{I-15-163}, including
the oscillation seminorm at $\alpha=0$ and all six
localized $L^2$ estimates.
\end{proposition}

The quadratic form separates into a fluctuation with zero
unit-scale average, an axial field on $\mathsf T_1$, and a
scalar gauge correction. The fluctuation is controlled by
$G_j$. The other two terms require lower bounds on their
constrained spaces; the axial lower bound follows by
combining estimates inside the blocks with estimates
across their faces.
Completing the scalar square then expresses $G-G_j$ in
terms of maps from unit bonds to fine bonds and inverses
on the unit lattice. Exponential decay of these inverses
and the one-scale regularity estimates give the conclusion.
The decomposition and representation are proved in
\cref{sec:two-scale-representation}; the decay estimates
and the proof of \cref{I-16-14} are completed in
\cref{sec:two-scale-estimates}.


\subsection{Averages and constrained extensions}
\label{sec:two-scale-averages}

The averaging maps control the energy of a prescribed coarse field
in both directions. Averaging does not increase the curl energy,
while every unit-scale field has a fine-scale extension with
uniformly bounded energy. The scalar estimate below gives the gap
on fields with zero mean in each $L$-block.

\begin{lemma}\label{I-15-165}
The averages satisfy the following estimates.
\begin{enumerate}
\item[(a)] For $A\in\mathcal A_\xi$,
 \[
  \|\partial_1\mathsf Q_jA\|\le\|\partial A\|.
 \]
\item[(b)] For $\lambda\in\mathcal F_\xi$,
 \[
  \|\Delta^1\mathsf Q'_j\lambda\|^2\le3\|\Delta\lambda\|^2.
 \]
\item[(c)] For every $B\in\mathcal A_1$ there is
 $A\in\mathcal A_\xi$ such that
 \[
  \mathsf Q_jA=B,\qquad
  \langle A,\Delta A\rangle\le\gamma_1\|B\|^2,
  \qquad \gamma_1=3\pi^2(\pi/2)^8.
 \]
 In particular, $\mathsf Q_j$ is onto.
\item[(d)] If $\omega\in\mathcal F_1$ and
 $\mathsf Q'_{1,L}\omega=0$, then
 \[
  \|\Delta^1\omega\|\ge4L^{-2}\|\omega\|.
 \]
\end{enumerate}
Moreover,
\begin{equation}\label{eq:two-scale-coarse-contraction}
 \|\mathsf Q_{1,L}B\|\le\|B\|,
 \qquad B\in\mathcal A_1,
\end{equation}
where the norm on the left has the volume weight $L^3$.
\end{lemma}

\begin{proof}
Put $n=L^j$, so that $n\xi=1$. We use the telescoping identity
\[
 f(x+e_\mu)-f(x)
   =\sum_{s=0}^{n-1}\xi\,\partial_\mu f(x+s\xi e_\mu).
\]
For $\mu<\nu$, insert the path average defining
$\mathsf Q_j$ into the coarse curl and telescope in both
directions. Since $B^j(y+e_\mu)=B^j(y)+e_\mu$, this gives
\[
 (\partial_1\mathsf Q_jA)_{\mu\nu}(y)
  =\sum_{x\in B^j(y)}\xi^3
      \sum_{s,t=0}^{n-1}\xi^2
       (\partial A)_{\mu\nu}
          (x+s\xi e_\mu+t\xi e_\nu).
\]
The positive weights sum to one. The Cauchy--Schwarz
inequality therefore bounds the square by the same average
of $|(\partial A)_{\mu\nu}|^2$. Sum over $y$.
The blocks partition $\mathsf T_\xi$, and translation
preserves its weighted sum, so
\[
 \begin{split}
 \sum_y|(\partial_1\mathsf Q_jA)_{\mu\nu}(y)|^2
 &\le\sum_{s,t=0}^{n-1}\xi^2
      \sum_{x\in\mathsf T_\xi}\xi^3
       |(\partial A)_{\mu\nu}
           (x+s\xi e_\mu+t\xi e_\nu)|^2\\
 &=\|(\partial A)_{\mu\nu}\|^2.
 \end{split}
\]
Summing over the three pairs proves (a).

For (b), commutation with the scalar average gives
\[
 \Delta^1\mathsf Q'_j\lambda
   =\partial_1^*\mathsf Q_j\partial\lambda
   =\sum_{\mu=1}^3
        \partial_{1,\mu}^*(\mathsf Q_j\partial\lambda)_\mu.
\]
The same double telescoping, now in one direction, gives
\[
 \begin{split}
 &\bigl(\partial_{1,\mu}^*
                   (\mathsf Q_j\partial\lambda)_\mu\bigr)(y)\\
 &\qquad=-\sum_{x\in B^j(y-e_\mu)}\xi^3
       \sum_{s,t=0}^{n-1}\xi^2
        (\partial_\mu\partial_\mu\lambda)
                  (x+(s+t)\xi e_\mu).
 \end{split}
\]
Thus its norm is at most
$\|\partial_\mu\partial_\mu\lambda\|
 =\|\partial_\mu^*\partial_\mu\lambda\|$; the equality
uses the translation relating forward and adjoint differences.
Consequently,
\[
 \|\Delta^1\mathsf Q'_j\lambda\|^2
 \le3\sum_\mu\|\partial_\mu^*\partial_\mu\lambda\|^2
 \le3\|\Delta\lambda\|^2.
\]
For the last inequality, commutation of the differences gives
\[
 \|\Delta\lambda\|^2
   =\sum_{\mu,\nu}
      \langle\partial_\mu^*\partial_\mu\lambda,
             \partial_\nu^*\partial_\nu\lambda\rangle
   =\sum_{\mu,\nu}\|\partial_\mu\partial_\nu\lambda\|^2.
\]

To construct the extension in (c), temporarily use complex
functions and the Hermitian inner products. Set
\[
 \mathcal P_N=\frac{2\pi}{N}\mathbb Z^3\cap[-\pi,\pi)^3,
 \qquad \epsilon_p(x)=e^{\mathrm ip\cdot x}.
\]
The restrictions to $\mathsf T_1$ are its $N^3$ characters.
They remain distinct characters on $\mathsf T_\xi$ and are
orthogonal there, with squared norm $N^3$ on both lattices.
Moreover,
\[
 \Delta\epsilon_p=\Delta(p)\epsilon_p,\qquad
 \Delta(p)=\sum_{\nu=1}^3
       4\xi^{-2}\sin^2(\xi p_\nu/2)
       \le|p|_2^2\le3\pi^2.
\]
For the one-dimensional average, write
\[
 g(\theta)=\sum_{k=0}^{n-1}\xi e^{\mathrm i\theta k\xi}.
\]
We have $g(0)=1$, and, when $0<|\theta|\le\pi$,
\[
 |g(\theta)|
   =\frac{\xi|\sin(\theta/2)|}{|\sin(\xi\theta/2)|}
   \ge\frac2\pi.
\]
Here $|\sin(\theta/2)|\ge|\theta|/\pi$ and
$|\sin(\xi\theta/2)|\le\xi|\theta|/2$.
The block average contributes three factors and the path
average one more. If $A_\mu=\epsilon_p$ and the other
components vanish, then
\[
 (\mathsf Q_jA)_\mu=w_\mu(p)\epsilon_p|_{\mathsf T_1},
 \qquad w_\mu(p)=g(p_\mu)\prod_{\nu=1}^3g(p_\nu),
 \qquad |w_\mu(p)|\ge(2/\pi)^4.
\]
Expand
$B_\mu=\sum_{p\in\mathcal P_N}\beta_\mu(p)
                   \epsilon_p|_{\mathsf T_1}$
and set
\[
 A_\mu^{\mathbb C}
   =\sum_{p\in\mathcal P_N}
        \beta_\mu(p)w_\mu(p)^{-1}\epsilon_p.
\]
Then $\mathsf Q_jA^{\mathbb C}=B$, and orthogonality gives
\[
 \begin{split}
 \langle A^{\mathbb C},\Delta A^{\mathbb C}\rangle
 &=\sum_{\mu,p}N^3\Delta(p)
       |w_\mu(p)|^{-2}|\beta_\mu(p)|^2\\
 &\le3\pi^2(\pi/2)^8
       \sum_{\mu,p}N^3|\beta_\mu(p)|^2
  =\gamma_1\|B\|^2.
 \end{split}
\]
Since $\mathsf Q_j$ and $B$ are real,
$A=\operatorname{Re}A^{\mathbb C}$ still has
$\mathsf Q_jA=B$. Finally,
$\langle A,\Delta A\rangle=\|\nabla A\|^2
 \le\|\nabla A^{\mathbb C}\|^2$, which proves (c).

If $\mathsf Q'_{1,L}\omega=0$, the mean of $\omega$
on every $B(y')$ vanishes. The path spectrum computed in
the proof of \cref{I-11-30}(a),(b), with three tensor
factors, gives the gap
$4\sin^2(\pi/(2L))\ge4L^{-2}$ for this box.
The blocks have $L$ sites along each edge and the torus
period satisfies $N\ge2L$, so their internal bonds are
distinct. Summing the box inequalities and discarding
the bonds joining different blocks gives
\[
 4L^{-2}\|\omega\|^2
 \le\|\partial_1\omega\|^2
 =\langle\omega,\Delta^1\omega\rangle
 \le\|\omega\|\,\|\Delta^1\omega\|,
\]
which proves (d).

For \eqref{eq:two-scale-coarse-contraction}, apply
Cauchy--Schwarz to the probability weights in
$\mathsf Q_{1,L}$:
\[
 |(\mathsf Q_{1,L}B)_\mu(y')|^2
 \le\sum_{y\in B(y')}L^{-3}
       \sum_{t=0}^{L-1}L^{-1}|B_\mu(y+te_\mu)|^2.
\]
Multiply by the coarse volume $L^3$ and sum over $y',\mu$.
For each $t$, the blocks followed by translation by $te_\mu$
cover $\mathsf T_1$ once. Each squared component therefore
occurs $L$ times with weight $L^{-1}$, giving the claim.
\end{proof}

\subsection{Constrained extensions and the gauge correction}
\label{sec:two-scale-minimizers}

To separate the two scales, we first extend prescribed unit-scale
averages to the fine lattice. The scalar extension minimizes the
squared Laplacian; the vector extension minimizes the one-scale
gauge-fixed energy. Their Euler equations identify the additional
part of the scalar projection and remove the mixed term in the
vector energy.

Write
\[
 N_j=\ker\mathsf Q'_j,
 \qquad
 D_j=\partial^*\partial+\partial\mathsf R_j\partial^*,
 \qquad
 \Delta_{1,j}=D_j+\mathsf Q_j^*\mathsf Q_j.
\]
Thus $G_j=\Delta_{1,j}^{-1}$ is the one-scale inverse already
constructed. Inverses of operators compressed to a subspace will
always be extended by zero on its orthogonal complement; this
convention includes the zero subspace.

\begin{lemma}\label{I-15-168}
The following constructions are defined in the setting of
\cref{sec:two-scale-propagators}.
\begin{enumerate}
\item[(a)] For every $\varphi\in\mathcal F_1$ there is a unique
minimizer $H'_j\varphi$ of $\|\Delta\lambda\|^2$ subject to
$\mathsf Q'_j\lambda=\varphi$. The map $H'_j$ is linear and
\[
 \mathsf Q'_jH'_j=I,
 \qquad \mathsf R_j\Delta H'_j=0.
\]
Set $\Delta'_j=(\Delta H'_j)^*(\Delta H'_j)$. Then
\[
 \langle\varphi,\Delta'_j\varphi\rangle
       \ge\tfrac13\|\Delta^1\varphi\|^2.
\]

\item[(b)] The scalar extension satisfies
\begin{equation}\label{I-15-169}
 \begin{aligned}
  \partial H'_j&=G_j\mathsf Q_j^*\partial_1,\\
  \partial\Delta H'_j&=\Delta G_j\mathsf Q_j^*\partial_1,\\
  \Delta H'_j&=\partial^*G_j\mathsf Q_j^*\partial_1.
 \end{aligned}
\end{equation}

\item[(c)] Let
\[
 W=\{\omega\in\mathcal F_1:
       \omega=0\text{ on }\mathsf T_1\setminus\Lambda,
       \ \mathsf Q'_{1,L}\omega=0\},
\]
and let $\iota_W:W\hookrightarrow\mathcal F_1$ be the inclusion.
On $W$,
\[
 \langle\omega,\Delta'_j\omega\rangle
       \ge\frac{16}{3L^4}\|\omega\|^2.
\]
Consequently
\[
 C'=\iota_W(\iota_W^*\Delta'_j\iota_W)^{-1}\iota_W^*
\]
is defined, and
\begin{equation}\label{I-15-170}
 \begin{aligned}
  N(\mathsf Q')&=N_j\oplus H'_jW,\\
  \mathsf R&=\mathsf R_j+\Delta H'_jC'H_j'^*\Delta.
 \end{aligned}
\end{equation}
The two terms in the second line are orthogonal projections
onto the mutually orthogonal spaces $\Delta N_j$ and
$\Delta H'_jW$, respectively.

\item[(d)] Define
\[
 \mathcal A^{\mathrm{ax}}
   =\{B\in\mathcal A_1:
       B=0\text{ on }\mathcal T(y')\text{ for every }y'\in\Lambda'\}.
\]
For every $B\in\mathcal A_1$ there is a unique $\omega\in W$
such that $B+\partial_1\omega\in\mathcal A^{\mathrm{ax}}$.

\item[(e)] For $\omega\in W$ and $A\in\mathcal A_\xi$,
\[
 \begin{aligned}
  \partial(A-\partial H'_j\omega)&=\partial A,\\
  \mathsf R_j\partial^*(A-\partial H'_j\omega)
       &=\mathsf R_j\partial^*A,\\
  \mathsf q\bigl(\mathsf Q_j(A-\partial H'_j\omega)\bigr)
       &=\mathsf q(\mathsf Q_jA).
 \end{aligned}
\]

\item[(f)] The operator $D_j$ is positive definite on
$\ker\mathsf Q_j$. Let $\iota:\ker\mathsf Q_j
\hookrightarrow\mathcal A_\xi$ be the inclusion and set
\[
 \widetilde G_j=\iota(\iota^*D_j\iota)^{-1}\iota^*,
 \qquad S=(\mathsf Q_jG_j\mathsf Q_j^*)^{-1}.
\]
For every $B\in\mathcal A_1$ there is a unique minimizer
$H_jB$ of $\langle A,D_jA\rangle$ subject to
$\mathsf Q_jA=B$. Define the symmetric operator $\Delta_j$ by
$\langle B,\Delta_jB\rangle=\langle H_jB,D_jH_jB\rangle$.
Then
\begin{equation}\label{I-15-171}
 \begin{aligned}
  H_j&=G_j\mathsf Q_j^*S,\\
  \widetilde G_j
     &=G_j-G_j\mathsf Q_j^*S\mathsf Q_jG_j,\\
  \Delta_j&=S-I.
 \end{aligned}
\end{equation}
Moreover, with $\gamma_1$ as in \cref{I-15-165},
\[
 \|\partial_1B\|^2
   \le\langle B,\Delta_jB\rangle
   \le\gamma_1\|B\|^2,
 \qquad I\le S\le(1+\gamma_1)I,
\]
and $\langle A',D_jH_jB\rangle=0$ for
$A'\in\ker\mathsf Q_j$.

\item[(g)] Put
\[
 F(B)=\mathsf q(B)+\langle B,\Delta_jB\rangle,
 \qquad \gamma''_0=(1+2916L^4)^{-1}.
\]
For $B\in\mathcal A^{\mathrm{ax}}$,
\begin{equation}\label{I-15-172}
 \mathsf q(B)+\|\partial_1B\|^2
       \ge\gamma''_0\|B\|^2,
\end{equation}
and hence
\begin{equation}\label{I-15-173}
 F(B)\ge\gamma''_0\|B\|^2.
\end{equation}
Let $T_F$ be the symmetric operator of the form $F$, and let
$\iota_{\mathrm{ax}}:\mathcal A^{\mathrm{ax}}
\hookrightarrow\mathcal A_1$ be the inclusion. The operator
\[
 \widetilde C
   =\iota_{\mathrm{ax}}
      (\iota_{\mathrm{ax}}^*T_F\iota_{\mathrm{ax}})^{-1}
       \iota_{\mathrm{ax}}^*
\]
is defined.
\end{enumerate}
\end{lemma}

\begin{proof}
\emph{The scalar extension.}
Since $\mathsf Q'_j\mathsf Q_j'^*=I$, the affine constraint
in (a) is nonempty and has direction space $N_j$. If
$\nu\in N_j$ and $\Delta\nu=0$, then
$\|\partial\nu\|^2=\langle\nu,\Delta\nu\rangle=0$,
so $\nu$ is constant. Its block average is zero, and hence
$\nu=0$. Thus the squared Laplacian is strictly positive
on the direction space. Its unique minimizer is characterized by
\[
 \mathsf Q'_j\lambda=\varphi,
 \qquad
 \langle\Delta\lambda,\Delta\nu\rangle=0
       \quad(\nu\in N_j).
\]
The second condition is $\mathsf R_j\Delta\lambda=0$.
These equations also show that the minimizer depends linearly
on $\varphi$. By \cref{I-15-165}(b),
\[
 \|\Delta^1\varphi\|^2
   =\|\Delta^1\mathsf Q'_jH'_j\varphi\|^2
   \le3\|\Delta H'_j\varphi\|^2,
\]
which proves the last assertion of (a).

On scalar functions, $\partial\partial=0$ and
$\partial^*\partial=\Delta$. Therefore
\[
 D_j\partial H'_j\varphi
   =\partial\mathsf R_j\Delta H'_j\varphi=0,
 \qquad
 \mathsf Q_j\partial H'_j\varphi
   =\partial_1\mathsf Q'_jH'_j\varphi
   =\partial_1\varphi.
\]
It follows that
$\Delta_{1,j}\partial H'_j\varphi
=\mathsf Q_j^*\partial_1\varphi$.
Applying $G_j$ gives the first identity in
\eqref{I-15-169}. Commutation of $\partial$ with $\Delta$
gives the second, and application of $\partial^*$ gives
the third.

\emph{The additional scalar projection.}
For $\omega\in W$, part (a) and \cref{I-15-165}(d) give
\[
 \langle\omega,\Delta'_j\omega\rangle
   \ge\tfrac13\|\Delta^1\omega\|^2
   \ge\frac{16}{3L^4}\|\omega\|^2.
\]
Thus the compressed inverse defining $C'$ exists.
The mixed average constraint is equivalent to
$\mathsf Q'_j\lambda\in W$: on every block outside
$\Lambda$, the unit-site averages already vanish, so their
coarse average vanishes as well. For $\lambda\in N(\mathsf Q')$,
\[
 \lambda-H'_j\mathsf Q'_j\lambda\in N_j.
\]
Both $N_j$ and $H'_jW$ lie in $N(\mathsf Q')$, and their
intersection is zero because $\mathsf Q'_jH'_j=I$.
This proves the direct sum in \eqref{I-15-170}.

The equation $\mathsf R_j\Delta H'_j=0$ makes its
two images under $\Delta$ orthogonal. The map
$Z=\Delta H'_j\iota_W$ is injective by the lower bound
just proved, and
$Z^*Z=\iota_W^*\Delta'_j\iota_W$.
Consequently the orthogonal projection onto its range is
\[
 Z(Z^*Z)^{-1}Z^*
   =\Delta H'_jC'H_j'^*\Delta.
\]
Adding the projection onto $\Delta N_j$ gives the
second identity in \eqref{I-15-170}.

\emph{Choosing the axial representative.}
Write $\Gamma_{y',y}$ for the coordinate-tree path from
$y'$ to $y$, and $B(\Gamma)$ for the signed sum of $B$
along a path. On a block $B(y')$, vanishing on the tree
is equivalent to vanishing of these sums along all
$\Gamma_{y',y}$: each tree bond extends the root path
to its initial endpoint, so its value is the difference
of the two root-path sums. Along such a path,
\[
 (\partial_1\omega)(\Gamma_{y',y})=\omega(y)-\omega(y').
\]
The required tree condition and the zero block mean therefore
force
\[
 \omega(y')=L^{-3}\sum_{y\in B(y')}B(\Gamma_{y',y}),
 \qquad
 \omega(y)=\omega(y')-B(\Gamma_{y',y}).
\]
Define $\omega$ by these formulas on the blocks with
$y'\in\Lambda'$ and by zero outside $\Lambda$.
Its mean on every block is zero, so $\omega\in W$,
and the tree condition holds. The same formulas establish
uniqueness, proving (d).

The first two identities in (e) follow from
$\partial\partial=0$ and
$\mathsf R_j\partial^*\partial H'_j
=\mathsf R_j\Delta H'_j=0$.
For the last one, the change in the unit-scale average is
$\mathsf Q_j\partial H'_j\omega=\partial_1\omega$.
It vanishes on $\Lambda^{c*}$ because both endpoints of
each such bond lie outside $\Lambda$. Its coarse average
also vanishes, since
\[
 \mathsf Q_{1,L}\partial_1\omega
   =\partial_L\mathsf Q'_{1,L}\omega=0.
\]
These are exactly the two sets of values entering
$\mathsf q$, which proves its invariance.

\emph{The vector extension.}
The form of $D_j$ is nonnegative. On $\ker\mathsf Q_j$
it agrees with the positive form of $\Delta_{1,j}$, so
the compression of $D_j$ is positive definite. By
\cref{I-15-165}(c), $\mathsf Q_j$ is onto. Thus
$\mathsf Q_j^*$ is injective and
$\mathsf Q_jG_j\mathsf Q_j^*$ is positive definite;
in particular, $S$ exists.

The constrained-square calculation in
\cref{sec:regional-linear-minimizers}, with $X=D_j$,
$\mathsf Q=\mathsf Q_j$ and $a=1$, now applies.
Indeed, $A_B=G_j\mathsf Q_j^*SB$ satisfies
$\mathsf Q_jA_B=B$ and
$\langle A',D_jA_B\rangle=\langle\mathsf Q_jA',SB\rangle=0$
for $A'\in\ker\mathsf Q_j$. Positivity on that kernel makes
$A_B=H_jB$ the unique minimizer, and
\[
 \langle H_jB,D_jH_jB\rangle
 =\langle H_jB,\Delta_{1,j}H_jB\rangle-\|B\|^2
 =\langle B,(S-I)B\rangle.
\]
This proves the first and third identities in \eqref{I-15-171}
and the mixed-term identity in (f).

For the second identity, set
$C_1=G_j-G_j\mathsf Q_j^*S\mathsf Q_jG_j$.
Then $\mathsf Q_jC_1=0$ and, for every
$A'\in\ker\mathsf Q_j$,
\[
 \langle A',D_jC_1J\rangle
   =\langle A',\Delta_{1,j}C_1J\rangle
   =\langle A',J\rangle.
\]
These equations uniquely characterize the inverse of the
compressed $D_j$, extended by zero: the difference of two
solutions lies in $\ker\mathsf Q_j$ and has zero
$D_j$ energy. Hence $C_1=\widetilde G_j$.

Curl contraction in \cref{I-15-165}(a) gives
\[
 \langle B,\Delta_jB\rangle
    \ge\|\partial H_jB\|^2
    \ge\|\partial_1\mathsf Q_jH_jB\|^2
    =\|\partial_1B\|^2.
\]
For the upper bound, choose the extension $A$ from
\cref{I-15-165}(c). Since
$\Delta-D_j=\partial\mathsf P_j\partial^*\ge0$,
minimality gives
\[
 \langle B,\Delta_jB\rangle
    \le\langle A,D_jA\rangle
    \le\langle A,\Delta A\rangle
    \le\gamma_1\|B\|^2.
\]
The bounds on $S$ now follow from $S=I+\Delta_j$.

\emph{The axial lower bound.}
Let $B\in\mathcal A^{\mathrm{ax}}$, and write
\[
 B=B_1+B_2,
 \qquad B_1=\chi_{\Lambda^*}B,
 \qquad B_2=\chi_{\Lambda^{c*}}B.
\]
Every tree bond in a block belonging to $\Lambda$ has
both endpoints in $\Lambda$. Thus $B_1$ vanishes on
the trees and outside $\Lambda^*$, and the supported
axial estimate \cref{I-15-166}, proved in
\cref{sec:two-scale-axial-gap}, applies to $B_1$.
Write $\varkappa=(108L^4)^{-1}$ and
$F_0(B)=\mathsf q(B)+\|\partial_1B\|^2$.
The weighted coarse contraction gives
\[
 L\sum_{c\in\Lambda'^*}
      |(\mathsf Q_{1,L}B_2)(c)|^2
   \le L^{-2}\|\mathsf Q_{1,L}B_2\|^2
   \le L^{-2}\|B_2\|^2.
\]
Here the coarse scalar product has weight $L^3$.
Also
\[
 \|\partial_1B_2\|^2
    \le\langle B_2,\Delta^1B_2\rangle
    \le12\|B_2\|^2.
\]
Apply the supported estimate to $B_1=B-B_2$, and bound
each square of a difference by twice the sum of the
two squares. It follows that
\[
 \begin{aligned}
 \varkappa\|B_1\|^2
 &\le2\left(
       L\sum_{c\in\Lambda'^*}|(\mathsf Q_{1,L}B)(c)|^2
       +\|\partial_1B\|^2\right)
       +(2L^{-2}+24)\|B_2\|^2\\
 &\le27F_0(B).
 \end{aligned}
\]
Indeed, $\|B_2\|^2$ is the first term in
$\mathsf q(B)$, and $2+2L^{-2}\le3$ since $L\ge2$.
The supports of $B_1$ and $B_2$ are disjoint, so
\[
 \|B\|^2=\|B_1\|^2+\|B_2\|^2
   \le(1+27\varkappa^{-1})F_0(B)
   =(1+2916L^4)F_0(B).
\]
This proves \eqref{I-15-172}. Part (f) gives
$F(B)\ge F_0(B)$, hence \eqref{I-15-173} and the
existence of $\widetilde C$.
\end{proof}

\subsection{The two-scale representation}
\label{sec:two-scale-representation}

The constrained extensions give coordinates in which the
two-scale energy is a sum of three positive forms. The
additional divergence term is a scalar square; shifting
that square produces the correction to the one-scale
inverse.

\begin{proposition}\label{I-16-5}
Let the operators be those of \cref{I-15-168}.
\begin{enumerate}
\item[(a)] The map
\[
 \ker\mathsf Q_j\times\mathcal A^{\mathrm{ax}}\times W
    \longrightarrow\mathcal A_\xi,
 \qquad
 (A',B,\omega)\longmapsto A'+H_jB-\partial H'_j\omega
\]
is a linear bijection. In these coordinates,
\begin{equation}\label{I-16-6}
 \begin{split}
 \langle A,\Delta_1A\rangle
    &=\langle A',D_jA'\rangle+F(B)
          +\langle\omega-\omega_0,
                     \Delta'_j(\omega-\omega_0)\rangle,\\
 \omega_0&=C'H_j'^*\Delta\partial^*(A'+H_jB).
 \end{split}
\end{equation}
In particular, $\Delta_1$ is positive definite.

\item[(b)] Put
$\mathsf K=\partial\Delta H'_jC'H_j'^*\partial^*$.
The inverse $G=\Delta_1^{-1}$ satisfies
\begin{equation}\label{I-16-7}
 \begin{split}
 \langle J,GJ\rangle
   ={}&\langle J,\partial H'_jC'H_j'^*\partial^*J\rangle\\
     &+\langle (I-\mathsf K)J,
          (\widetilde G_j+H_j\widetilde C H_j^*)
                         (I-\mathsf K)J\rangle.
 \end{split}
\end{equation}
Equivalently,
\[
 G=\partial H'_jC'H_j'^*\partial^*
      +(I-\mathsf K)^*
          (\widetilde G_j+H_j\widetilde C H_j^*)
                               (I-\mathsf K).
\]

\item[(c)] Define maps from $\mathcal A_1$ to
$\mathcal A_\xi$ and operators on $\mathcal A_1$ by
\[
 \begin{gathered}
 X=G_j\mathsf Q_j^*,\qquad Y=G_j\Delta X,\\
 \Theta=\partial_1C'\partial_1^*,\qquad
 V=-S+S\widetilde C S,\\
 \mathsf N_1=X^*\Delta X,\qquad
 \mathsf N_2=(\Delta X)^*G_j\Delta X.
 \end{gathered}
\]
Then
\begin{equation}\label{I-16-8}
 \begin{split}
 G&=G_j+XU_0X^*-X\Theta Y^*-Y\Theta X^*,\\
 U_0&=\Theta+V-\Theta\mathsf N_1V-V\mathsf N_1\Theta
        +\Theta\mathsf N_2\Theta
        +\Theta\mathsf N_1V\mathsf N_1\Theta.
 \end{split}
\end{equation}
If $\Lambda'=\varnothing$, then $W=\{0\}$,
$C'=0$, $\widetilde C=S^{-1}$, $V=0$ and $G=G_j$.
\end{enumerate}
\end{proposition}

\begin{proof}
\emph{The coordinates and the energy.}
Given $A\in\mathcal A_\xi$, choose the unique
$\omega\in W$ for which
\[
 B=\mathsf Q_jA+\partial_1\omega\in\mathcal A^{\mathrm{ax}};
\]
this is possible by \cref{I-15-168}(d). Set
$A'=A+\partial H'_j\omega-H_jB$.
Since $\mathsf Q_j\partial H'_j=\partial_1$ and
$\mathsf Q_jH_j=I$, we have
$\mathsf Q_jA'=\mathsf Q_jA+\partial_1\omega-B=0$.
This proves surjectivity. Conversely, any such
representation of $A$ satisfies
$\mathsf Q_jA+\partial_1\omega=B\in\mathcal A^{\mathrm{ax}}$.
The uniqueness of its axial correction determines
$\omega$, then $B$, and finally $A'$. The map is linear,
so it is the stated bijection.

Put $A_0=A'+H_jB$ and
$\mathsf R_W=\Delta H'_jC'H_j'^*\Delta$.
The projection identity \eqref{I-15-170} is an orthogonal
sum. Hence
\[
 \begin{split}
 \langle A,\Delta_1A\rangle
   ={}&\|\partial A\|^2
       +\|\mathsf R_j\partial^*A\|^2
       +\mathsf q(\mathsf Q_jA)
       +\|\mathsf R_W\partial^*A\|^2.
 \end{split}
\]
The first three terms are unchanged when
$A=A_0-\partial H'_j\omega$ is replaced by $A_0$,
by \cref{I-15-168}(e). Since $\mathsf Q_jA_0=B$,
their sum is
\[
 \begin{split}
 \langle A_0,D_jA_0\rangle+\mathsf q(B)
   ={}&\langle A',D_jA'\rangle
      +2\langle A',D_jH_jB\rangle\\
      &+\langle H_jB,D_jH_jB\rangle+\mathsf q(B)
    =\langle A',D_jA'\rangle+F(B).
 \end{split}
\]
The mixed term is zero by the minimizing equation in
\cref{I-15-168}(f).

It remains to compute the last divergence term. On
scalars $\partial^*\partial=\Delta$, so
\[
 \partial^*A=\partial^*A_0-\Delta H'_j\omega,
 \qquad
 \mathsf R_W\partial^*A_0=\Delta H'_j\omega_0.
\]
The vector $\omega_0$ lies in $W$. For $\omega\in W$,
the compressed inverse satisfies $C'\Delta'_j\omega=\omega$;
therefore
\[
 \mathsf R_W\Delta H'_j\omega
  =\Delta H'_jC'(\Delta H'_j)^*(\Delta H'_j)\omega
  =\Delta H'_j\omega.
\]
It follows that
\[
 \|\mathsf R_W\partial^*A\|^2
   =\|\Delta H'_j(\omega_0-\omega)\|^2
   =\langle\omega-\omega_0,
                  \Delta'_j(\omega-\omega_0)\rangle,
\]
which proves \eqref{I-16-6}.

All three terms in that identity are nonnegative.
Their lower bounds on the three constrained spaces
are given in \cref{I-15-168}(c,f,g). If the sum is
zero, then $A'=0$ and $B=0$, so $\omega_0=0$;
the scalar term then forces $\omega=0$.
Thus $A=0$, proving positivity and the existence of $G$.

\emph{The inverse.}
We use the following finite-dimensional square completion.
If a symmetric operator $T$ is positive definite on
a subspace $\mathcal V_0$, and $C_0$ is its compressed
inverse extended by zero, then, for $v\in\mathcal V_0$,
\[
 2\langle v,J\rangle-\langle v,Tv\rangle
   =\langle J,C_0J\rangle
       -\langle v-C_0J,T(v-C_0J)\rangle.
\]
Indeed, $C_0J\in\mathcal V_0$ and
$\langle v,TC_0J\rangle=\langle v,J\rangle$
on that subspace. The supremum over $v$ is therefore
$\langle J,C_0J\rangle$, attained at $v=C_0J$.
Applying this first to $T=\Delta_1$ on $\mathcal A_\xi$
gives
\[
 \langle J,GJ\rangle
   =\sup_{A\in\mathcal A_\xi}
       \bigl(2\langle A,J\rangle-\langle A,\Delta_1A\rangle\bigr).
\]

Use the coordinates in (a) and write
$\omega=\omega_0+\omega'$. Symmetry of $C'$ and
$\Delta$ gives
\[
 \begin{split}
 \langle\partial H'_j\omega_0,J\rangle
  &=\langle\partial H'_jC'H_j'^*\Delta\partial^*A_0,J\rangle\\
  &=\langle A_0,\partial\Delta H'_jC'H_j'^*\partial^*J\rangle
   =\langle A_0,\mathsf KJ\rangle.
 \end{split}
\]
The expression to be maximized thus becomes
\[
 \begin{split}
 &-2\langle\omega',H_j'^*\partial^*J\rangle
          -\langle\omega',\Delta'_j\omega'\rangle\\
 &\quad+2\langle A',(I-\mathsf K)J\rangle
          -\langle A',D_jA'\rangle\\
 &\quad+2\langle B,H_j^*(I-\mathsf K)J\rangle-F(B).
 \end{split}
\]
For fixed $A'$ and $B$, translation by $\omega_0\in W$
is a bijection of $W$. Consequently the three variables
range independently over $W$, $\ker\mathsf Q_j$ and
$\mathcal A^{\mathrm{ax}}$. The three suprema are
\[
 \begin{gathered}
 \langle H_j'^*\partial^*J,C'H_j'^*\partial^*J\rangle,\\
 \langle (I-\mathsf K)J,\widetilde G_j(I-\mathsf K)J\rangle,\\
 \langle H_j^*(I-\mathsf K)J,
                   \widetilde C H_j^*(I-\mathsf K)J\rangle.
 \end{gathered}
\]
Their sum is \eqref{I-16-7}. Both sides of its operator
version are symmetric and have this quadratic form,
so they agree.

\emph{Expansion about the one-scale inverse.}
By \eqref{I-15-169},
$\partial H'_j=X\partial_1$ and
$\partial\Delta H'_j=\Delta X\partial_1$.
It follows that
\[
 \partial H'_jC'H_j'^*\partial^*=X\Theta X^*,
 \qquad \mathsf K=\Delta X\Theta X^*.
\]
By \eqref{I-15-171}, $H_j=XS$ and
$\widetilde G_j=G_j-XSX^*$. Thus
\[
 \widetilde G_j+H_j\widetilde C H_j^*
    =G_j+XVX^*.
\]
Substitution into (b), using symmetry of $\Theta$,
gives
\[
 G=X\Theta X^*
     +\bigl(I-X\Theta(\Delta X)^*\bigr)
        (G_j+XVX^*)
        \bigl(I-\Delta X\Theta X^*\bigr).
\]
We keep the order of the factors when multiplying.
Since $G_j$ and $\Delta$ are symmetric,
\[
 (\Delta X)^*G_j=Y^*,\qquad
 (\Delta X)^*X=X^*\Delta X=\mathsf N_1,\qquad
 (\Delta X)^*G_j\Delta X=\mathsf N_2.
\]
The two cross products and the remaining double product are
\[
 \begin{aligned}
 X\Theta(\Delta X)^*(G_j+XVX^*)
    &=X\Theta Y^*+X\Theta\mathsf N_1VX^*,\\
 (G_j+XVX^*)\Delta X\Theta X^*
    &=Y\Theta X^*+XV\mathsf N_1\Theta X^*,\\
 X\Theta(\Delta X)^*(G_j+XVX^*)\Delta X\Theta X^*
    &=X\Theta\mathsf N_2\Theta X^*
       +X\Theta\mathsf N_1V\mathsf N_1\Theta X^*.
 \end{aligned}
\]
The first two are subtracted and the last is added.
Collecting the terms with outer factors $X$ and $X^*$
gives exactly \eqref{I-16-8}.

If $\Lambda'=\varnothing$, then
$\mathcal A^{\mathrm{ax}}=\mathcal A_1$ and
$\mathsf q(B)=\|B\|^2$. Since $\Delta_j=S-I$,
we have $F(B)=\langle B,SB\rangle$ and
$\widetilde C=S^{-1}$, so $V=0$.
Also $W=\{0\}$, hence $C'=0$, $\Theta=0$ and
$U_0=0$. Formula \eqref{I-16-8} then gives $G=G_j$.
\end{proof}

In \eqref{I-16-8}, the maps $X$ and $Y$ have unit-bond
inputs and fine-bond outputs, while the coefficients
$U_0$ and $\Theta$ act on unit bonds. The estimates
in \cref{sec:two-scale-estimates} use this distinction
to transfer the one-scale regularity bounds to the
three correction terms.

\subsection{The axial-gauge gap}
\label{sec:two-scale-axial-gap}

The tree condition controls bonds inside a block by its curl.
For bonds crossing a block face, the curl controls the
variation along the face and the coarse average controls
the mean. We keep a separate estimate on the tangential
face bonds, since these are the error terms in the
crossing-face calculation.

\begin{lemma}\label{I-15-166}\label{I-15-57}
Suppose that $B\in\mathcal A_1$ vanishes outside
$\Lambda^*$ and on $\mathcal T(y)$ for every
$y\in\Lambda'$. Then
\begin{equation}\label{I-15-59}
 L\sum_{c\in\Lambda'^*}|(\mathsf Q_{1,L}B)(c)|^2
       +\|\partial_1B\|^2
 \ge\frac1{108}L^{-4}\|B\|^2.
\end{equation}
\end{lemma}

\begin{proof}
Put $F=\partial_1B$. Coordinates in a block or in two
adjacent blocks are taken on a lift of their coarse
sites. Since $N\ge2L$, the sites in these coordinate
boxes are distinct. For $y\in\mathsf T_L$, define
\[
 \begin{split}
 \Phi^-_\mu(y)&=\{x\in B(y):x_\mu=y_\mu\},\\
 \Phi^+_\mu(y)&=\{x\in B(y):x_\mu=y_\mu+L-1\},\\
 X_y^\mu&=\sum_{(x,\mu)\subset B(y)}|B_\mu(x)|^2,
 \qquad X_y=\sum_\mu X_y^\mu,\qquad
 S_y=\sum_{p\subset B(y)}|F(p)|^2.
 \end{split}
\]
Here a bond or plaquette is inside a block when all
its vertices are there. For a face $\Phi$, let
$\Phi^\circ$ be its bonds, with both endpoints in
$\Phi$, and put
$E_\Phi=\sum_{b\in\Phi^\circ}|B(b)|^2$.

For a coarse bond $c=(y,\mu)$, write
\[
 \Phi'(c)=\Phi^+_\mu(c_-),\qquad
 \Phi''(c)=\Phi^-_\mu(c_+),\qquad
 B(c)=\{(x,\mu):x\in\Phi'(c)\}.
\]
Thus $B(c)$ consists of the unit bonds crossing
from $B(c_-)$ to $B(c_+)$. Translation by $e_\mu$
maps $\Phi'(c)$ and its tangential bonds bijectively
onto $\Phi''(c)$ and its tangential bonds.
Let $\mathcal P(c)$ consist of the plaquettes with
corners
$x,x+e_\mu,x+e_\mu+e_\nu,x+e_\nu$, where
$(x,\nu)\in\Phi'(c)^\circ$. There is exactly one
such plaquette for each bond of $\Phi'(c)^\circ$.
Set
\[
 \begin{aligned}
 f_c(x)&=B_\mu(x)\quad(x\in\Phi'(c)),&
 \bar f_c&=L^{-2}\sum_{x\in\Phi'(c)}f_c(x),\\
 Y_c&=\sum_{x\in\Phi'(c)}|f_c(x)|^2,&
 S_c&=\sum_{p\in\mathcal P(c)}|F(p)|^2,\\
 D_c&=\sum_{(x,\nu)\in\Phi'(c)^\circ}
                   |f_c(x+e_\nu)-f_c(x)|^2.
 \end{aligned}
\]

\emph{Bonds inside a block.}
Fix a block $B(y)$ on whose tree $B$ vanishes.
For a bond $(x,\iota)\subset B(y)$, put
\[
 x^{(i)}=(y_1,\ldots,y_i,x_{i+1},\ldots,x_3),
 \qquad 0\le i<\iota.
\]
All these sites and their translates by $e_\iota$
remain in the block. The bond
$(x^{(\iota-1)},\iota)$ lies in the tree, so its
value is zero. For $i<\iota$, join $x^{(i)}$
to $x^{(i-1)}$ in direction $i$. At consecutive
sites $w,w+e_i$ on this segment, both
$(w,i)$ and $(w+e_\iota,i)$ lie in the tree:
their coordinates of index less than $i$ equal
those of $y$. Hence
\[
 F(w,i,\iota)
   =B_i(w)+B_\iota(w+e_i)
        -B_i(w+e_\iota)-B_\iota(w)
   =B_\iota(w+e_i)-B_\iota(w).
\]
Telescope along these segments. If
$\varpi(x,\iota)$ denotes the resulting set of
plaquettes, then
$B_\iota(x)=\sum_{p\in\varpi(x,\iota)}\pm F(p)$
and
$\#\varpi(x,\iota)\le(\iota-1)(L-1)$. Therefore
\begin{equation}\label{I-15-47}
 |B_\iota(x)|^2
 \le(\iota-1)(L-1)
             \sum_{p\in\varpi(x,\iota)}|F(p)|^2.
\end{equation}
For $\iota=1$, the bond already lies in the tree.

For a fixed plaquette $(w,i,\iota)$ with $i<\iota$,
its occurrence requires $w_m=y_m$ for $m<i$.
The coordinates of $x$ with index greater than
$i$ must equal those of $w$; there are at most
$L-1$ choices of $x_i$, and the first $i-1$
coordinates are arbitrary. Its multiplicity is
thus at most $(L-1)L^{i-1}$, and hence at most
$(L-1)L^{\iota-2}$. Summing \eqref{I-15-47}
for $\iota=2,3$ gives
\[
 X_y^\iota
 \le(\iota-1)(L-1)^2L^{\iota-2}
       \sum_{i<\iota}\ \sum_{(w,i,\iota)\subset B(y)}
                       |F(w,i,\iota)|^2.
\]
The plaquette sets for different $\iota$ are
disjoint, since $\iota$ is their larger direction.
Consequently,
\begin{equation}\label{eq:two-scale-block-energy}
 X_y\le2(L-1)^2L\,S_y\le3L^3S_y.
\end{equation}

\emph{Tangential face bonds.}
Let $\Phi=\{x\in B(y):x_\mu=s\}$ be either face
normal to direction $\mu$. Its bonds have
direction $\iota\ne\mu$. In the preceding count
for $(w,i,\iota)$, impose $x_\mu=s$.
If $\mu<i$, one of the first $i-1$ coordinates is
fixed, giving at most $(L-1)L^{i-2}\le L-1$
choices. If $\mu=i$, there are at most
$L^{i-1}\le L$ choices. If $\mu>i$, then
$i<\iota$, $i<\mu$ and $\iota\ne\mu$ force
$i=1$, so there are at most $L-1$ choices.
Thus the multiplicity is at most $L$ in every
case. Summing \eqref{I-15-47} over the face,
and again using the disjoint plaquette sets,
gives the sharper estimate
\begin{equation}\label{eq:two-scale-face-energy}
 E_\Phi\le2(L-1)L\,S_y.
\end{equation}

\emph{Bonds crossing a face.}
This calculation does not require the tree
condition. For $(x,\nu)\in\Phi'(c)^\circ$,
the plaquette relation is
\[
 f_c(x+e_\nu)-f_c(x)
   =B_\nu(x+e_\mu)-B_\nu(x)\mp F(p),
 \qquad p\in\mathcal P(c).
\]
Apply the three-term Cauchy--Schwarz inequality,
then sum over the tangential bonds. The
translation bijection between the faces gives
\begin{equation}\label{eq:two-scale-corridor-curl}
 S_c\ge\tfrac13D_c-E_{\Phi'(c)}-E_{\Phi''(c)}.
\end{equation}

The coarse average also sees the mean of the
crossing bonds. Put $\beta=y_\mu$, on the chosen
lift. In
\[
 (\mathsf Q_{1,L}B)(c)
    =L^{-4}\sum_{x\in B(c_-)}
                      \sum_{k=0}^{L-1}B_\mu(x+ke_\mu),
\]
a component at coordinate $z_\mu=\beta+t$
occurs $t+1$ times for $0\le t\le L-1$,
and $2L-1-t$ times for $L\le t\le2L-2$.
The terms at $t=L-1$ therefore contribute
$L^{-1}\bar f_c$. The other terms are
\[
 \begin{split}
 \mathsf a_c
   &=L^{-4}\sum_{x\in B(c_-)\setminus\Phi'(c)}
                         (x_\mu-\beta+1)B_\mu(x),\\
 \mathsf b_c
   &=L^{-4}\sum_{x\in B(c_+)\setminus\Phi^+_\mu(c_+)}
                         (2L-1+\beta-x_\mu)B_\mu(x).
 \end{split}
\]
They involve only direction-$\mu$ bonds inside
the indicated blocks. Each sum has at most
$L^3$ coefficients, all of modulus at most
$L^{-3}$, so
\[
 |\mathsf a_c|^2\le L^{-3}X^\mu_{c_-},
 \qquad
 |\mathsf b_c|^2\le L^{-3}X^\mu_{c_+}.
\]
Since
$L^{-1}\bar f_c=(\mathsf Q_{1,L}B)(c)
                   -\mathsf a_c-\mathsf b_c$,
we obtain
\begin{equation}\label{I-15-49}
 |(\mathsf Q_{1,L}B)(c)|^2
 \ge\tfrac13L^{-2}|\bar f_c|^2
                  -L^{-3}(X^\mu_{c_-}+X^\mu_{c_+}).
\end{equation}

The face graph is the product of two paths with
$L$ vertices. Taking two tensor factors in the
path spectrum from the proof of
\cref{I-11-30}(a),(b) shows that its first
positive eigenvalue is
\[
 \kappa_L=4\sin^2\frac{\pi}{2L}\ge4L^{-2}.
\]
Its constant eigenspace is removed by subtracting
the face mean. Hence
\[
 \begin{split}
 D_c&\ge\kappa_L
          \sum_{x\in\Phi'(c)}|f_c(x)-\bar f_c|^2,\\
 Y_c&=L^2|\bar f_c|^2+
          \sum_{x\in\Phi'(c)}|f_c(x)-\bar f_c|^2.
 \end{split}
\]
Multiply \eqref{eq:two-scale-corridor-curl} by
$\tfrac14L^{-2}$ and add \eqref{I-15-49}.
The inequality
$\tfrac14L^{-2}\kappa_L\ge L^{-4}$
makes the coefficients of the mean and the
zero-mean part agree. We conclude that
\begin{equation}\label{I-15-53}
 \begin{split}
 \tfrac14L^{-2}S_c+|(\mathsf Q_{1,L}B)(c)|^2
 &\ge\tfrac13L^{-4}Y_c\\
 &\quad-\tfrac14L^{-2}
                 (E_{\Phi'(c)}+E_{\Phi''(c)})\\
 &\quad-L^{-3}(X^\mu_{c_-}+X^\mu_{c_+}).
 \end{split}
\end{equation}

\emph{Summation over the region.}
Every unit bond is either inside one block or
crosses from that block to its unique neighbour
in the bond direction. Thus it belongs to
exactly one of the sets of internal bonds or
$B(c)$. Since $B$ vanishes outside $\Lambda^*$,
\[
 X=\sum_{y\in\Lambda'}X_y,\qquad
 Y=\sum_{c\in\Lambda'^*}Y_c,\qquad
 \|B\|^2=X+Y.
\]
If $y\notin\Lambda'$, every internal bond of
$B(y)$ has both endpoints outside $\Lambda$,
so $X_y$ and all its $E_\Phi$ vanish.
Write
\[
 S_{\mathrm{in}}=\sum_{y\in\Lambda'}S_y,\qquad
 S_{\mathrm{co}}=\sum_{c\in\Lambda'^*}S_c,\qquad
 T=L\sum_{c\in\Lambda'^*}
                    |(\mathsf Q_{1,L}B)(c)|^2.
\]
Each face occurs as $\Phi'(c)$ or $\Phi''(c)$
for exactly one coarse bond. Each block is the
initial block of one bond in each direction and
the final block of one bond in each direction.
Therefore \eqref{eq:two-scale-block-energy} and
\eqref{eq:two-scale-face-energy} give
\[
 \begin{split}
 \sum_{c\in\Lambda'^*}
       (E_{\Phi'(c)}+E_{\Phi''(c)})
       &\le12(L-1)L\,S_{\mathrm{in}},\\
 \sum_{\substack{c\in\Lambda'^*\\c=(y,\mu)}}
       (X^\mu_{c_-}+X^\mu_{c_+})
       &\le2X\le4(L-1)^2L\,S_{\mathrm{in}}.
 \end{split}
\]
The six faces account for the first factor $12$;
the two occurrences of each directional block
sum account for the second estimate. Summing
\eqref{I-15-53} now yields
\[
 \begin{split}
 \tfrac14L^{-2}S_{\mathrm{co}}+L^{-1}T
 &\ge\tfrac13L^{-4}Y
      -\left(3\frac{L-1}{L}
             +4\frac{(L-1)^2}{L^2}\right)S_{\mathrm{in}}\\
 &\ge\tfrac13L^{-4}Y-7S_{\mathrm{in}}.
 \end{split}
\]
Multiply by $1/36$. Since
$L^{-2}/144\le1$, $L^{-1}/36\le1$ and
$7/36\le1/4$, it follows that
\[
 S_{\mathrm{co}}+T
       \ge\frac1{108}L^{-4}Y-\tfrac14S_{\mathrm{in}}.
\]
On the other hand, the block bound gives
\[
 \tfrac34S_{\mathrm{in}}
       \ge\tfrac14L^{-3}X
       \ge\frac1{108}L^{-4}X.
\]
Adding proves
\[
 T+S_{\mathrm{in}}+S_{\mathrm{co}}
      \ge\frac1{108}L^{-4}(X+Y)
      =\frac1{108}L^{-4}\|B\|^2.
\]
Internal plaquettes in different blocks are
distinct. A plaquette in $\mathcal P(c)$ has
two bonds in $B(c)$ and its other two bonds
inside the adjacent blocks; it therefore
belongs to no other $\mathcal P(c')$ and is
inside no block. Consequently
$S_{\mathrm{in}}+S_{\mathrm{co}}\le\|F\|^2$,
which proves \eqref{I-15-59}.
\end{proof}

\subsection{Exponential kernels and regular columns}
\label{sec:two-scale-kernels}

The scalar correction is constrained by a zero mean on
each selected block. Its inverse can therefore be
estimated with exponential weights that are constant on
blocks. We first prove the corresponding matrix estimate,
then explain how regularity passes through the maps from
unit bonds to fine bonds.

Let $\mathbb Y$ and $\mathbb Y'$ each be the sites or the
positive bonds of $\mathsf T_1$. A bond has the position
of its initial point. These spaces have counting-measure
inner products, so their indicator vectors $e_s$ are
orthonormal. For $\delta,c>0$, write
\[
 U(s,s')=\langle e_s,Ue_{s'}\rangle,
 \qquad
 \mathfrak U(\delta,c)
  =\{U:|U(s,s')|\le ce^{-\delta|s-s'|_1}
                               \text{ for all }s,s'\}.
\]
All maps in this subsection are real and linear. Recall
$\sigma(\delta)=3\coth^3(\delta/2)$.

\begin{lemma}\label{I-15-126}
\begin{enumerate}[label=\textup{(\roman*)},leftmargin=*]
\item For any continuous centers $u,u'$,
\[
 \sum_{s\in\mathbb Y}
       e^{-\delta|u-s|_1-\delta|s-u'|_1}
 \le \sigma(\delta/2)e^{-\delta|u-u'|_1/2}.
\]
Consequently, for composable maps,
\[
 \mathfrak U(\delta,c_1)\mathfrak U(\delta,c_2)
 \subset
 \mathfrak U\bigl(\delta/2,\sigma(\delta/2)c_1c_2\bigr).
\]
\item Let $T$ be symmetric and belong to
$\mathfrak U(\delta_1,c_1)$ on $\ell^2(\mathbb Y)$.
Partition $\mathbb Y$ into cells of diameter at most $D$.
Suppose that a subspace $\mathcal H$ is invariant under
multiplication by every real function constant on the
cells, and that
\[
 \langle v,Tv\rangle\ge\gamma\|v\|^2
       \quad(v\in\mathcal H),\qquad \gamma>0.
\]
For its inclusion $\iota$, the compressed inverse
\[
 C=\iota(\iota^*T\iota)^{-1}\iota^*
\]
exists, has norm at most $\gamma^{-1}$, and satisfies
\[
 |C(s,s')|
   \le \frac2\gamma e^{\theta D}
                              e^{-\theta|s-s'|_1}
\]
for a $\theta>0$ depending only on
$c_1,\delta_1,\gamma,D$. The inverse on a
zero-dimensional space is extended by zero.
\end{enumerate}
\end{lemma}

\begin{proof}
Each unit site is the position of at most three indices of
$\mathbb Y$. The scalar torus sum in
\cref{sec:one-scale-projection} and the triangle inequality
give \textup{(i)}: three times the scalar bound
$\coth^3(\delta/4)$ is $\sigma(\delta/2)$.
Expanding the matrix product gives the kernel conclusion.

For \textup{(ii)}, let $\Pi$ project onto $\mathcal H$.
Fix $s_0\in\mathbb Y$ and put
\[
 \rho(s)=\min_{z\text{ in the cell of }s}|z-s_0|_1 .
\]
The cell diameter and triangle inequality imply
\[
 \rho(s_0)=0,\qquad
 \rho(s)\ge|s-s_0|_1-D,\qquad
 |\rho(s)-\rho(s')|\le|s-s'|_1+D.
\]
For the last inequality, choose a minimizing $z'$ in the
cell of $s'$ and use
$\rho(s)\le|s-s_0|_1
 \le|s-s'|_1+D+\rho(s')$, then exchange $s,s'$.
Multiplication $E=e^{\theta\rho}$ and its inverse preserve
$\mathcal H$. They also preserve its orthogonal complement,
since they are symmetric, and hence commute with $\Pi$.

The lower bound for $T$ makes $\iota^*T\iota$ injective
and thus invertible in finite dimension. If
$v=(\iota^*T\iota)^{-1}g$, then
$\gamma\|v\|^2\le\langle v,g\rangle$; this proves
$\|C\|\le\gamma^{-1}$.
The entries of $ETE^{-1}-T$ are bounded in modulus by
\[
 c_1e^{-\delta_1|s-s'|_1}
                  \bigl(e^{\theta(|s-s'|_1+D)}-1\bigr).
\]
Both their row sums and their column sums are at most
\[
 m(\theta)
  =3c_1\sum_{z\in\mathbb Z^3}
       e^{-\delta_1|z|_1}
                    \bigl(e^{\theta(|z|_1+D)}-1\bigr).
\]
Indeed, least-distance lifts of distinct positions are
distinct, and there are at most three indices at each
position. If a matrix $B$ has row sums of $|B|$ at most
$m_1$ and column sums at most $m_2$, Cauchy--Schwarz in
the double sum gives
\[
 \begin{split}
 |\langle g,Bf\rangle|
 &\le
 \left(\sum_{s,s'}|B(s,s')||g(s)|^2\right)^{1/2}
 \left(\sum_{s,s'}|B(s,s')||f(s')|^2\right)^{1/2}\\
 &\le(m_1m_2)^{1/2}\|g\|\|f\|.
 \end{split}
\]
Thus $\|ETE^{-1}-T\|\le m(\theta)$.
For $0\le\theta\le\delta_1/2$ the series is dominated
by a summable exponential; it is continuous and tends
to zero as $\theta\downarrow0$.
Choose $0<\theta\le\delta_1/2$ with
$m(\theta)\le\gamma/2$.

On $\mathcal H$ the conjugated form has lower bound
$\gamma/2$, so
$T_\theta=\iota^*ETE^{-1}\iota$ is invertible with
$\|T_\theta^{-1}\|\le2/\gamma$.
Commutation with $\Pi$ shows that
$\iota T_\theta^{-1}\iota^*=ECE^{-1}$.
Its $(s,s_0)$ entry is $e^{\theta\rho(s)}C(s,s_0)$,
whose modulus is at most $2/\gamma$.
The lower bound for $\rho(s)$ proves the stated kernel
estimate.
\end{proof}

To retain the fine-lattice regularity, fix $\delta>0$
and consider a map $Y:\mathcal A_1\to\mathcal A_\xi$
and a family $\mathsf k=(k,(k_\alpha)_{0\le\alpha<1})$
of positive constants.
We write $Y\in\mathfrak X(\delta,\mathsf k)$ if each
column $Ye_b$ satisfies, for every $u\in\mathsf T$,
$x\in U_1(u)$ and real $\zeta$ supported in $U_1(u)$,
\[
 \begin{gathered}
 |(Ye_b)(x)|,\quad |(\nabla Ye_b)(x)|,\quad
 |(\Delta Ye_b)(x)|\le kE_b,\\
 \|\zeta\nabla Ye_b\|_\alpha
       \le k_\alpha E_b(\|\zeta\|_\alpha+|\zeta|),
 \qquad
 \|\zeta\nabla\nabla Ye_b\|\le kE_b|\zeta|,
 \qquad E_b=e^{-\delta|u-b_-|_1}.
 \end{gathered}
\]
The last norm uses the fine weight $\xi^3$.
Multiplication of a family by a constant below is
componentwise. Put
\begin{equation}\label{I-16-10}
 \begin{split}
  C^{\mathfrak X}_1
    &=\max\{9\sigma(\delta/2),e^{3\delta}\},\\
  C^{\mathfrak X}_2&=\sqrt3\,\sigma(\delta/2),\\
  C^{\mathfrak X}_3&=243e^{3\delta}\sigma(\delta/2).
 \end{split}
\end{equation}

\begin{lemma}\label{I-16-9}
\begin{enumerate}[label=\textup{(\roman*)},leftmargin=*]
\item If $Y\in\mathfrak X(\delta,\mathsf k)$ and
$U\in\mathfrak U(\delta,c)$ acts on $\mathcal A_1$, then
\[
 YU\in\mathfrak X\bigl(\delta/2,
                           \sigma(\delta/2)c\,\mathsf k\bigr).
\]
\item If $Y,Z\in\mathfrak X(\delta,\mathsf k)$, then
$Z^*Y$, $Z^*\Delta Y$ and
$(\partial^*Z)^*\partial^*Y$ belong to
$\mathfrak U(\delta/2,C^{\mathfrak X}_1k^2)$, and
$\mathsf Q_jY\in\mathfrak U(\delta,C^{\mathfrak X}_1k)$.
\item Suppose that $O$ satisfies
\textup{(I1)--(I5)} of \cref{I-15-163} with rate
$\delta$ and constants
$c,c_\alpha,c'_\varepsilon,c''_{\alpha,\varepsilon}$.
If $f:\mathcal A_1\to\mathcal A_\xi$ is linear and
\[
 |(fe_b)(x)|\le ke^{-\delta|u-b_-|_1}
                  \quad(x\in U_1(u))
\]
for every $u$ and $b$, then
\[
 Of\in\mathfrak X\bigl(\delta/2,
                   C^{\mathfrak X}_2 k(c,(c_\alpha))\bigr).
\]
In particular, this applies to $f=\mathsf Q_j^*$ with
$k=e^{9\delta}$.
\item If $Y,Z\in\mathfrak X(\delta,\mathsf k)$, then
$YZ^*$ satisfies all of \textup{(I1)--(I5)} with rate
$\delta/2$ and constants
\[
 c=c'_\varepsilon=C^{\mathfrak X}_3k^2,\qquad
 c_\alpha=c''_{\alpha,\varepsilon}
                  =C^{\mathfrak X}_3k(k+k_\alpha).
\]
\end{enumerate}
\end{lemma}

\begin{proof}
The blocks $B^j(y)\subset U_1(y)$ partition the fine
lattice and have weighted volume one. Also
\[
 \sum_{x\in U_1(u)}\xi^3\le(2+\xi)^3\le27.
\]
These volume bounds account for the component factors
in the proof.

For \textup{(i)}, expand
$YUe_{b'}=\sum_bU(b,b')Ye_b$.
Each defining norm or seminorm is subadditive and
homogeneous. The convolution bound of
\cref{I-15-126}\textup{(i)} therefore gives all five
column estimates, with the same factors of $\zeta$.

For \textup{(ii)}, split the inner product into blocks:
\[
 \begin{split}
 |(Z^*Y)(b,b')|
 &\le
 \sum_{y\in\mathsf T_1}
       \sum_{x\in B^j(y)}\xi^3
                      3|(Ze_b)(x)||(Ye_{b'})(x)|\\
 &\le3k^2\sum_y
           e^{-\delta|y-b_-|_1-\delta|y-b'_-|_1}.
 \end{split}
\]
The same bound holds with $\Delta Y$ in place of $Y$.
For the divergence product, the adjoint difference gives
\[
 |(\partial^*Ye_b)(x)|
   \le\sum_{\mu=1}^3
          |(\partial_\mu(Ye_b)_\mu)(x-\xi e_\mu)|
   \le3k e^{-\delta|y-b_-|_1},
       \quad x\in B^j(y),
\]
since $x-\xi e_\mu\in U_1(y)$. Its inner product
therefore has coefficient $9k^2$. Convolution proves
the three kernel bounds. Finally,
$(\mathsf Q_jYe_{b'})(b)$ is a probability average of
values in $b_-+[0,2)^3$, which is contained in
$U_1(b_-+(1,1,1))$. The change of center gives
\[
 |(\mathsf Q_jYe_{b'})(b)|
 \le e^{3\delta}k e^{-\delta|b_--b'_-|_1}.
\]

For \textup{(iii)}, decompose
$fe_b=\sum_y\chi_{B^j(y)}fe_b$.
The term indexed by $y$ has supremum at most
$ke^{-\delta|y-b_-|_1}$ and $L^2$ norm at most
$\sqrt3\,ke^{-\delta|y-b_-|_1}$.
Apply \textup{(I1)} to its value, gradient and Laplacian,
the first estimate of \textup{(I2)} to its localized
gradient seminorm, and the third estimate of
\textup{(I5)} to its localized second differences.
Summing by \cref{I-15-126}\textup{(i)} gives,
respectively, the coefficients
$\sigma(\delta/2)ck$,
$\sigma(\delta/2)c_\alpha k$ and
$\sqrt3\,\sigma(\delta/2)ck$.
They are bounded by the asserted family.
For $f=\mathsf Q_j^*$, \eqref{eq:one-scale-adjoint-column} gives a
column bounded by one and supported in $b_-+[0,2)^3$.
If it is nonzero at an $x\in U_1(u)$, then
$|u-b_-|_1\le9$, which proves the required hypothesis
with $k=e^{9\delta}$.

For \textup{(iv)}, write
\[
 YZ^*J=\sum_b Ye_b\,\langle Ze_b,J\rangle .
\]
Let the three source fields have the support prescribed
in \cref{I-15-163}, and put
$E'_b=e^{-\delta|u'-b_-|_1}$.
The volume bound and Cauchy--Schwarz give
\[
 \begin{aligned}
 |\langle Ze_b,J\rangle|
       &\le81kE'_b|J|,
 &|\langle Ze_b,J\rangle|
       &\le9kE'_b\|J\|,\\
 |\langle\nabla Ze_b,J'\rangle|
       &\le243kE'_b|J'|,
 &|\langle\nabla Ze_b,J'\rangle|
       &\le9\sqrt3\,kE'_b\|J'\|.
 \end{aligned}
\]
Adjointness identifies the second row with the
coefficients for $\nabla^*J'$.
For the last source, use the second-difference column
bound with the indicator $\chi$ of $U_1(u')$:
\[
 \begin{split}
 |\langle Ze_b,\nabla^*\nabla^*J''\rangle|
  &=|\langle\chi\nabla\nabla Ze_b,J''\rangle|\\
  &\le kE'_b\|J''\|.
 \end{split}
\]
No seminorm of this indicator is involved.

On the output side, the defining column bounds give
the three pointwise quantities and the localized
gradient seminorm, together with
\[
 \|\zeta Ye_b\|,\ \|\zeta\nabla Ye_b\|
      \le9\sqrt3\,kE_b|\zeta|,
 \qquad
 \|\zeta\nabla\nabla Ye_b\|\le kE_b|\zeta|.
\]
We also need the localized value seminorm:
\begin{equation}\label{I-16-11}
 \|\zeta Ye_b\|_\alpha
    \le e^{3\delta}kE_b(\|\zeta\|_\alpha+|\zeta|).
\end{equation}
To prove it, put $g=Ye_b$ and take
$0<|x-x'|_1\le1$. If either cutoff value is nonzero,
exchange the points so that $x'\in U_1(u)$.
A shortest coordinate path from $x'$ to $x$ lies in
the union of the boxes $U_1(u+v)$ with
$v\in\{-1,0,1\}^3$. On this path the value and gradient
of $g$ are bounded by $e^{3\delta}kE_b$.
Telescoping along the path yields
\[
 |g(x)-g(x')|
 \le |x-x'|_1 e^{3\delta}kE_b
 \le |x-x'|_1^\alpha e^{3\delta}kE_b .
\]
Apply this and the value bound to
\[
 \zeta(x)g(x)-\zeta(x')g(x')
   =(\zeta(x)-\zeta(x'))g(x)
                    +\zeta(x')(g(x)-g(x')) .
\]
If both cutoff values vanish the difference is zero.
This proves \eqref{I-16-11}, including $\alpha=0$
and pairs crossing the support of $\zeta$.

Every estimate for $YZ^*$ now pairs one of the output
bounds with one of the source coefficient bounds and
sums over $b$. For example,
\[
 \begin{split}
 (\nabla YZ^*\nabla^*J')(x)
    &=\sum_b(\nabla Ye_b)(x)
                              \langle\nabla Ze_b,J'\rangle,\\
 \|\zeta YZ^*\nabla^*J'\|_\alpha
    &\le\sum_b\|\zeta Ye_b\|_\alpha
                              |\langle\nabla Ze_b,J'\rangle|.
 \end{split}
\]
Convolution contributes $\sigma(\delta/2)$ and the
rate $\delta/2$.
For \textup{(I1)}, \textup{(I3)} and all six members
of \textup{(I5)}, the coefficient products are at most
$243k^2$. In the two seminorm estimates they are at most
$243k\max\{k_\alpha,e^{3\delta}k\}$.
For \textup{(I3)} and \textup{(I4)} we have used only
$|J'|$, which is bounded by the stated
seminorm-plus-supremum expression. The constants in
the lemma dominate these products.
\end{proof}

\subsection{The two-scale estimates}
\label{sec:two-scale-estimates}

The operator of the coarse penalty has finite range.
We record its normalization, since the scalar product
on $\mathcal A_L$ carries the weight $L^3$.

\begin{lemma}\label{I-16-12}
On $\mathcal A_1$ put
\begin{equation}\label{I-16-13}
 T_{\mathsf q}
  =\chi_{\Lambda^{c*}}
           +L^{-2}\mathsf Q_{1,L}^*
                    \chi_{\Lambda'^*}\mathsf Q_{1,L}.
\end{equation}
Then $\langle B,T_{\mathsf q}B\rangle=\mathsf q(B)$,
$0\le T_{\mathsf q}\le1+L^{-2}$, and
$T_{\mathsf q}\partial_1\omega=0$ for $\omega\in W$.
Moreover,
\[
 \begin{split}
 \Delta_1
   &=\partial^*\partial+\partial\mathsf R\partial^*
                              +\mathsf Q_j^*T_{\mathsf q}\mathsf Q_j\\
   &=\Delta-\partial\mathsf P\partial^*
                              +\mathsf Q_j^*T_{\mathsf q}\mathsf Q_j,
 \qquad T_F=T_{\mathsf q}+\Delta_j .
 \end{split}
\]
The unit-bond entries of $T_{\mathsf q}$ lie in
$[0,1+2L^{-5}]$ and vanish when
$|b_--b'_-|_1\ge4L$. In particular,
$T_{\mathsf q}\in\mathfrak U(\delta,2e^{4\delta L})$
for every $\delta>0$.
\end{lemma}

\begin{proof}
The second term of the quadratic form of
\eqref{I-16-13} is
\[
 L^{-2}\|\chi_{\Lambda'^*}\mathsf Q_{1,L}B\|^2
       =L\sum_{c\in\Lambda'^*}
                         |(\mathsf Q_{1,L}B)(c)|^2 .
\]
The contraction of $\mathsf Q_{1,L}$ gives the upper
bound. For $\omega\in W$, its gradient vanishes on
$\Lambda^{c*}$ and
$\mathsf Q_{1,L}\partial_1\omega
 =\partial_L\mathsf Q'_{1,L}\omega=0$.
This proves the annihilation assertion.
The operator identities follow from the quadratic
forms and the previously proved lattice identity
$\Delta=\partial^*\partial+\partial\partial^*$.

Write
$(\mathsf Q_{1,L}B)(c)=\sum_bq(c,b)B(b)$.
For a coarse bond of direction $\mu$,
\[
 q(c,b)=L^{-4}
  \#\{(y,t):y\in B(c_-),\ 0\le t<L,\
                         y+te_\mu=b_-\}
\]
when $b$ has direction $\mu$, and the coefficient is
zero otherwise. Thus $0\le q(c,b)\le L^{-3}$.
A nonzero coefficient places $b_-$ in $B(c_-)$ or
$B(c_+)$, so for fixed $b$ at most two coarse bonds
contribute. The entries of the second term of
\eqref{I-16-13} are consequently
\[
 L\sum_{c\in\Lambda'^*}q(c,b)q(c,b')
                         \in[0,2L^{-5}].
\]
If such an entry is nonzero, the two positions lie in
one path-averaging domain: their coordinate differences
are at most $2L-2$ in direction $\mu$ and $L-1$ in
the other directions. Their torus distance is therefore
at most $4L-4<4L$.
The indicator term is diagonal. Since
$1+2L^{-5}\le2$, the exponential bound follows.
\end{proof}

\begin{proof}[Proof of \cref{I-16-14}]
The decomposition in \cref{I-16-5} proves positivity
and gives \eqref{I-16-7} and \eqref{I-16-8}.
The scalar and axial lower bounds in
\cref{I-15-168} give
\[
 \|C'\|\le\frac{3L^4}{16},\qquad
 \|\widetilde C\|\le1+2916L^4.
\]

By \cref{I-15-163} and \cref{I-16-9}\textup{(iii)},
$X=G_j\mathsf Q_j^*$ belongs to a class $\mathfrak X$
with an absolute rate and with constants depending
only on the indicated exponent.
In particular its columns have a pointwise Laplacian
bound. Apply the same result to $f=\Delta X$, decreasing
the input rate if necessary: it gives
$Y=G_j\Delta X$ with the same kind of uniformity.
All subsequent rates may be decreased and constants
increased when finitely many estimates are combined.

Next,
$S^{-1}=\mathsf Q_jG_j\mathsf Q_j^*=\mathsf Q_jX$
has an exponential kernel by
\cref{I-16-9}\textup{(ii)}.
The bound $S^{-1}\ge(1+\gamma_1)^{-1}$ from
\cref{I-15-168} is absolute. Apply
\cref{I-15-126}\textup{(ii)} to $T=S^{-1}$ on
all unit bonds, with singleton cells and $D=0$.
It gives an exponential kernel for $S$ with absolute
constants. The same follows for $\Delta_j=S-1$.

The entries of $\partial_1$ and its adjoint are
$0$ or $\pm1$ and vanish at distance greater than one.
By \eqref{I-15-169},
\[
 \Delta'_j
  =\partial_1^*(\partial^*X)^*(\partial^*X)\partial_1 .
\]
Thus \cref{I-16-9}\textup{(ii)} and kernel
multiplication give exponential decay for $\Delta'_j$
with absolute constants.
For its compression to $W$, use the cells $B(y')$,
which have diameter $3(L-1)$. Both the support condition
and the zero block means defining $W$ are preserved by
block-constant multiplication. The lower bound is
$16/(3L^4)$. Hence \cref{I-15-126}\textup{(ii)}
gives decay for $C'$ with constants depending only on
$L$. Kernel multiplication gives the same conclusion
for $\Theta=\partial_1C'\partial_1^*$.

By \cref{I-16-12}, $T_F=T_{\mathsf q}+\Delta_j$
has an exponential kernel with constants depending only
on $L$. The axial space is a coordinate subspace of
the unit-bond space, so it is invariant under every
multiplication. Apply \cref{I-15-126}\textup{(ii)}
with singleton cells and lower bound
$(1+2916L^4)^{-1}$. This gives the required decay of
$\widetilde C$, and hence of $V=-S+S\widetilde CS$.

The two maps
$\mathsf N_1=X^*\Delta X$ and $Y^*\Delta X$ have
exponential kernels with absolute constants by
\cref{I-16-9}\textup{(ii)}.
Taking the adjoint preserves these bounds, so they
also hold for
$\mathsf N_2=(\Delta X)^*Y=(Y^*\Delta X)^*$.
The six terms defining $U_0$ in \eqref{I-16-8}
are finite products of these kernels and those of
$\Theta,V$. Repeated use of
\cref{I-15-126}\textup{(i)} gives $U_0$ an exponentially
decaying kernel with constants depending only on $L$.

Finally, \cref{I-16-9}\textup{(i)} puts
$XU_0$, $X\Theta$ and $Y\Theta$ in classes
$\mathfrak X$ with constants depending only on $L$
and the exponents. After decreasing their rates to a
common positive rate, part \textup{(iv)} gives every
estimate \textup{(I1)--(I5)} for
\[
 XU_0X^*,\qquad X\Theta Y^*,\qquad Y\Theta X^* .
\]
Equation \eqref{I-16-8} and the estimates for $G_j$
give the same family for $G$.
Signs in that identity do not affect the norm bounds.

Choose $\delta_{\mathrm{per}}>0$ as the minimum of the
rates used above and enlarge the estimate constants for $G$.
The kernel estimates and the three inverse estimates use only
the unindexed constants, so this rate is independent of the
H\"older exponents. All bounds are uniform in $j,N,\Lambda'$;
their constants depend only on $L$ and the displayed exponents,
as stated in \cref{I-16-14}.
\end{proof}


\section{The global Dirichlet inverse}\label{sec:global-dirichlet}

We construct the inverse on a separated multiscale region by joining
the interior two-scale inverses of \cref{sec:two-scale-propagators}
and massive boundary inverses.
The gauge projection is nonlocal, so its replacement by the local
projections contributes a separate error. We estimate this error
before summing the Neumann series. The resulting bounds are uniform
in the lattice spacing, the volume and the number of layers.

Fix an odd integer $L\ge13$, positive integers $K,m$, and
$0\le t\le K$. In level-$t$ coordinates put
\[
 \xi=L^{-t},\qquad \xi_n=L^n\xi,\qquad
 T=\mathbb R^3/(\mathsf S\mathbb Z)^3,
 \qquad \mathsf S=2L^{m+K-t}.
\]
Thus the bare lattice has spacing $\xi$ and $\xi_t=1$.
Let $\mathfrak M,R_1$ be positive powers of $L$ and let
$\Omega_0\supset\cdots\supset\Omega_t$, with
$\Omega_0\ne\emptyset$ and $\Omega_{t+1}=\emptyset$, satisfy
\[
 \begin{gathered}
 \Omega_n\text{ is a union of cubes of }\pi[\mathfrak M\xi_n]
                         \quad(0\le n\le t),\\
 \operatorname{dist}_\infty(\Omega_n^c,\Omega_{n+1})
                         >R_1\mathfrak M\xi_n
                         \quad(0\le n<t).
 \end{gathered}
\]
Here $\pi[r]$ is the standard partition into cubes of side $r$
centered at the sites of $r\mathbb Z^3$ modulo the periods.
A complementary region has the convention of
\cref{sec:multiscale-constraints}. Require also that $\mathsf S$
be an even multiple of $\mathfrak M\xi_{t_*}$, where $t_*$ is
the last index with $\Omega_{t_*}\ne\emptyset$. Empty final
layers and regions equal to $T$ are allowed.

Write
\[
 \mathcal L_n=\overline{\Omega_n\setminus\Omega_{n+1}},
 \qquad \Lambda_n=\xi_n\mathbb Z^3\cap\mathcal L_n,
 \qquad \Lambda_* =\bigcup_{n=0}^t\Lambda_n,
\]
with the lattice understood modulo the periods of $T$.
For $y\in\Lambda_n$ put $\xi_y=\xi_n$ and let $\Delta(y)$
be the centered cube of side $\xi_y$; put
$\widetilde\Delta(y)=\{|x-y|_\infty\le\xi_y\}$.
The bare sites of $\Delta(y)$ form the block $B^n(y)$.
These blocks partition the bare sites of $\Omega_0$.
Let $y(x)$ be the center of the block containing $x$.
Use the weighted distance
\[
 d_{\mathbf B}(x,x')=\inf_{\Gamma\ni x,x'}
 \left(\sum_{n=1}^t\xi_n^{-1}
          \operatorname{length}(\Gamma\cap\mathcal L_n)
       +\xi^{-1}\operatorname{length}(\Gamma\setminus\Omega_1)\right),
\]
where $\Gamma$ ranges over connected finite unions of segments,
as in \cref{sec:multiscale-sums}.

\subsection{The operator and the estimates}
\label{subsec:global-dirichlet-statement}

We index bonds only in the positive orientations.
Let $\mathfrak b_0$ contain every bare bond with both endpoints in
$\Omega_0$, and contain only bonds with at least one endpoint there.
The spaces $\mathcal S$, $\mathcal A$ and $\mathcal Y$ consist,
respectively, of real functions on the bare sites of $\Omega_0$,
on $\mathfrak b_0$, and on $\Lambda_*$. Extend the first two by
zero. Their inner products have weights $\xi^3$, $\xi^3$ and
$\xi_y^3$. The spaces $\mathcal A^{(r)}$, $r=1,2$, consist of
functions on all bare bonds with $r$ additional direction indices;
their inner products have weight $\xi^3$ and sum over all indices.

For a positive bond $b$, its localization point is $b_-$, except
that it is $b_+$ if only $b_+$ lies in $\Omega_0$. Support of a
bond field and the indicators $\chi_y$ of $\Delta(y)$ refer to
this point. Multiplication by a site function $\zeta$, however,
uses $\zeta(b_-)$. Differences act on the zero extensions:
$\nabla_\nu=\xi^{-1}(T_\nu-I)$ and
$\Delta=\sum_\nu\nabla_\nu^*\nabla_\nu$.
The adjoints $\nabla^*$ and $\nabla^*\nabla^*$ take their
values in $\mathcal A$ by restriction to $\mathfrak b_0$.
Let $\chi_{\mathcal S}$ and $\chi_{\mathcal A}$ be the orthogonal
projections given by restriction and zero extension onto these spaces.
Use the scalar gradient and exterior derivative $\partial$ from
\cref{sec:one-order-inverse}, and set
\[
 \Delta_{\mathcal S}=\chi_{\mathcal S}\Delta\chi_{\mathcal S},
 \qquad \partial_0=\chi_{\mathcal A}\partial\chi_{\mathcal S}.
\]
The scalar averaging map and the scalar Dirichlet operator are
\[
 (Q'\lambda)(y)=L^{-3n}\sum_{x\in B^n(y)}\lambda(x)
                         \quad(y\in\Lambda_n),
 \qquad
 \Delta'_a=\Delta_{\mathcal S}+Q'^*\mathsf a'Q',
 \qquad \mathsf G'=(\Delta'_a)^{-1},
\]
where $\mathsf a'$ multiplies by $a_n\xi_n^{-2}$ on
$\Lambda_n$, $a_0=a_1=1$ and
$a_{n+1}=a_n/(L^{-2}+a_n)$ for $n\ge1$. In particular,
$1-L^{-2}\le a_n\le1$. The scalar inverse exists by the
block energy estimate proved below. Let $\mathsf R$ be the
orthogonal projection onto $\Delta_{\mathcal S}\ker Q'$ and
put $\mathsf P=I-\mathsf R$.

For $n\ge1$, let $\Gamma_n$ be the positive level-$n$ bonds
with an endpoint in $\Omega_n$ and none in $\Omega_{n+1}$;
let $\Gamma_0$ be the bonds of $\mathfrak b_0$ with no endpoint
in $\Omega_1$. For $c=(c_-,\mu)\in\Gamma_n$, put
\begin{equation}\label{I-16-63}
 (QA)(c)=L^{-4n}\sum_{x\in B^n(c_-)}
                  \sum_{j=0}^{L^n-1}A_\mu(x+j\xi e_\mu).
\end{equation}
Define the symmetric nonnegative operator $Q^*aQ$ by its form
\[
 (A,Q^*aQA)=\sum_{n=0}^t\xi_n\sum_{c\in\Gamma_n}|(QA)(c)|^2.
\]
Here the unit coefficient $a=1$ is fixed; $Q^*aQ$ denotes the
operator defined by this form. The vector operator is
\begin{equation}\label{I-16-64}
 \Delta_a(\Omega_0)=
  \chi_{\mathcal A}\partial^*\partial\chi_{\mathcal A}
  +\partial_0\mathsf R\partial_0^*+Q^*aQ.
\end{equation}
For a site or bond field, use the componentwise supremum norm
$|f|$ and the seminorm
\[
 \|f\|_{\alpha,r}
   =\sup_{0<|x-x'|_1\le r}
         \left(\frac r{|x-x'|_1}\right)^\alpha|f(x)-f(x')|,
 \qquad \|f\|_\alpha=\|f\|_{\alpha,1}.
\]
The comparison is between the same components at translated
sites or bonds, using initial points for the bond distances.
For $\alpha=0$ this is the corresponding
oscillation seminorm. Write $\|f\|$ for the weighted $L^2$ norm.

\begin{proposition}\label{I-17-5}
There are positive constants $M_0,\gamma,\delta,C$, depending
only on $L$, with the following property. Suppose the separated
sequence above satisfies $\mathfrak M\ge M_0$.
Every positive-power choice $R_1\ge L$ is allowed.
Then $\Delta_a(\Omega_0)\ge\gamma I$, and its inverse
$\mathsf G$ is symmetric and positive definite.

Let $y,y'\in\Lambda_*$, put
$E=e^{-\delta d_{\mathbf B}(y,y')}$, and let $x$ be a bare
bond localized in $\Delta(y)$.
For $J$ supported in $\Delta(y')$,
\begin{equation}\label{eq:global-dirichlet-pointwise}
 \begin{split}
 |\mathsf GJ(x)|&\le C\xi_y^2E|J|,\\
 |\nabla\mathsf GJ(x)|,
 |\mathsf G\nabla^*J(x)|&\le C\xi_yE|J|,\\
 |\Delta\mathsf GJ(x)|&\le CE|J|.
 \end{split}
\end{equation}
Here $J\in\mathcal A$ except in the divergence-input bound,
where $J\in\mathcal A^{(1)}$.
For $0\le\alpha<1$, and a real site function $\zeta$
supported in $\widetilde\Delta(y)$,
\begin{equation}\label{eq:global-dirichlet-holder}
 \|\zeta\nabla\mathsf GJ\|_\alpha,
 \|\zeta\mathsf G\nabla^*J\|_\alpha
 \le C_\alpha\xi_y^{1-\alpha}
       (\|\zeta\|_{\alpha,\xi_y}+|\zeta|)E|J|,
\end{equation}
with the same respective input spaces and supports.

For $J\in\mathcal A^{(1)}$ supported in
$\widetilde\Delta(y')$, one has
\begin{align}
 |\nabla\mathsf G\nabla^*J(x)|
 &\le C_\varepsilon E
       (\|J\|_{\varepsilon,\xi_{y'}}+|J|),
 &&0<\varepsilon<1,
 \label{eq:global-dirichlet-mixed}\\
 \|\zeta\nabla\mathsf G\nabla^*J\|_\alpha
 &\le C_{\alpha,\varepsilon}\xi_y^{-\alpha}
      (\|\zeta\|_{\alpha,\xi_y}+|\zeta|)E
      (\|J\|_{\alpha+\varepsilon,\xi_{y'}}+|J|),
 \label{eq:global-dirichlet-mixed-holder}
\end{align}
where the second bound assumes $0\le\alpha<1$,
$\varepsilon>0$, $\alpha+\varepsilon<1$, and the same
condition on $\zeta$ as in \eqref{eq:global-dirichlet-holder}.
The constants with subscripts depend only on $L$ and those exponents.

For $J$ supported in $\Delta(y')$ and a real site function
$\zeta$ supported in $\Delta(y)$, the six bounds are
\begin{equation}\label{eq:global-dirichlet-L2}
 \begin{aligned}
 \|\zeta\mathsf GJ\|&\le C\xi_y^2|\zeta|E\|J\|,\\
 \|\zeta\nabla\mathsf GJ\|,
 \|\zeta\mathsf G\nabla^*J\|
      &\le C\xi_y|\zeta|E\|J\|,\\
 \|\zeta\nabla\mathsf G\nabla^*J\|,
 \|\zeta\nabla\nabla\mathsf GJ\|,
 \|\zeta\mathsf G\nabla^*\nabla^*J\|
      &\le C|\zeta|E\|J\|.
 \end{aligned}
\end{equation}
An input with $r$ adjoint differences belongs to
$\mathcal A^{(r)}$, with $\mathcal A^{(0)}=\mathcal A$.
The same inequalities hold with $\chi_y$ in place of $\zeta$
and $1$ in place of $|\zeta|$.

Finally, there are a finite family of cubes $\mathfrak D$,
local inverses $\mathsf G_\square$, a quadratic partition
$(h_\square)$, and remainder blocks $K_{\square',\square}$,
constructed below, such that
\begin{equation}\label{I-17-6}
 \mathsf G=\sum_{i=0}^\infty\mathsf G_0\mathcal R^i
          =\sum_{i=0}^\infty(\mathcal R^*)^i\mathsf G_0
          =\sum_{i=0}^\infty\ 
            \sum_{\omega\in\mathfrak D^{2i+1}}\mathsf G_\omega.
\end{equation}
The series converge in operator norm.
Here $\mathsf G_0=\sum_\square h_\square\mathsf G_\square
h_\square$, $\Delta_a\mathsf G_0=I-\mathcal R$, and
\[
 \mathcal R=\sum_{\square',\square}
    K_{\square',\square}\mathsf G_\square h_\square.
\]
For $\omega=(\square_0,\ldots,\square_{2i})$, set $|\omega|=i$ and define
\[
 \mathsf G_\omega=
 h_{\square_0}\mathsf G_{\square_0}h_{\square_0}
 \prod_{k=1}^i
   \bigl(K_{\square_{2k-1},\square_{2k}}
       \mathsf G_{\square_{2k}}h_{\square_{2k}}\bigr).
\]
The factors are ordered with increasing $k$ from left to right.
For every $i$, each bound in
\eqref{eq:global-dirichlet-pointwise}--\eqref{eq:global-dirichlet-L2}
also bounds the sum of the
corresponding moduli, norms or seminorms over
$\omega\in\mathfrak D^{2i+1}$, with its right side multiplied
by $2^{-i}$. Use $\mathsf G_\omega$ for undifferentiated
input and for $\nabla\mathsf G\nabla^*$; use
$\mathsf G_\omega^*$ for $\mathsf G\nabla^*$ and
$\mathsf G\nabla^*\nabla^*$.
The same bounds therefore hold for $\mathsf G_0\mathcal R^i$
and $(\mathcal R^*)^i\mathsf G_0$, respectively.

All constants are independent of $K,m,t,\mathfrak M,R_1$,
the regions, the bond frame, and the number of layer changes.
\end{proposition}

We first obtain the scalar inverse and the differentiated gauge
projection directly from energy estimates. These estimates control
the discrepancy between global and local gauge conditions.
The local two-scale bounds then give a parametrix whose remainder
has norm $O(\mathfrak M^{-1})$ in exponentially weighted block
norms. The final argument sums individual products of remainder
blocks, which gives the absolute walk bounds in
\eqref{I-17-6} as well as the operator estimates.

\subsection{Geometry and summation over cells}
\label{subsec:global-cell-geometry}

The scalar estimates below will also be used on auxiliary tori.
For these estimates allow any bare spacing $\xi>0$ and torus
side $\mathsf S$ that is an even multiple of
$\mathfrak M\xi_s$, where $s$ is the last nonempty level.
Keep the centered partitions and separation condition above.
Discarding empty final regions changes neither the blocks,
the operators nor the weighted distance. No normalization
$\xi_s=1$ is imposed in these supporting estimates.

This extension preserves the elementary grid facts used earlier.
Indeed, $\mathsf S/\xi_n$ is an even integer at every
nonempty level, the ratios of successive spacings are the odd
integer $L$, and $\mathsf S\ge2\mathfrak M\xi_n$.
Thus the coordinate and nesting argument in
\cref{sec:block-averaging} and the block-partition proof of
\cref{I-11-27} apply: fine sites lie on no cell face, and
the bare block at $y\in\Lambda_n$ consists of
\[
 y+\xi k,\qquad
 |k_\mu|\le\tfrac12(L^n-1),\quad k\in\mathbb Z^3.
\]
Its inner-bond graph is the product of three paths.
The proof of the metric properties in \cref{I-9-22} uses
only finite segment covers and the torus triangle inequality,
so applies in this setting as well.

Put $\rho(x)=\xi_{y(x)}$ on the bare sites of $\Omega_0$
and $\rho(x)=\xi$ elsewhere. In this subsection write
$d=d_{\mathbf B}$, and extend the empty regions beyond $s$.

\begin{lemma}\label{lem:global-cell-geometry}
The following estimates hold uniformly in the number of layers.
\begin{enumerate}[label=(\alph*)]
\item For neighboring bare sites $x,x'$,
\[
 \frac{\rho(x')}{\rho(x)}\in\{L^{-1},1,L\},\qquad
 d(x,x')\le\frac\xi2\bigl(\rho(x)^{-1}+\rho(x')^{-1}\bigr).
\]
In particular, the latter bound is at most
$\min\{1,L\xi/\rho(x),L\xi/\rho(x')\}$.
For $x\in B^n(y)$, one has $d(x,y)\le3/2$.

\item Suppose $\mathfrak M\ge L+10$ and
$R_1\mathfrak M\ge L(L+9)$. The cube
\[
 C(y)=\{x:|x-y|_\infty\le(L+9)\xi_y\}
\]
has a lift to a cube in $\mathbb R^3$ on which the covering
map is injective; its lifted bare sites preserve nearest-neighbor
relations. If $y\in\Lambda_n$,
then $C(y)\cap\Omega_{n+2}=\emptyset$; if $n\ge1$, also
$C(y)\subset\Omega_{n-1}$. Every bare site $x\in C(y)$ satisfies
\[
 L^{-1}\xi_y\le\rho(x)\le L\xi_y,
 \qquad d(x,y)\le L\xi_y^{-1}|x-y|_1.
\]

\item For $y\in\Lambda_j$ and $y'\in\Lambda_{j'}$,
\begin{equation}\label{eq:global-layer-gap}
 d(y,y')\ge R_1\mathfrak M(|j-j'|-1)_+.
\end{equation}
Consequently, for real $q$ and
$\varrho=\log L/(R_1\mathfrak M)$,
\begin{equation}\label{eq:global-scale-ratio}
 (\xi_y/\xi_{y'})^q
       \le L^{|q|}e^{|q|\varrho d(y,y')}.
\end{equation}

\item Put
\begin{equation}\label{I-13-121}
 c_S(\kappa)=8L^3(2+4/\kappa)^3,
 \qquad \mathsf r_S(\kappa)=\frac{2(3\log L+1)}\kappa.
\end{equation}
For $\kappa>0$ and $R_1\mathfrak M\ge\mathsf r_S(\kappa)$,
and for $y,y'\in\Lambda_*$ in the second estimate,
\begin{align}
 \sup_{y\in\Lambda_*}\sum_{y'\in\Lambda_*}e^{-\kappa d(y,y')}
       &\le c_S(\kappa),\label{I-13-122}\\
 \sum_{z\in\Lambda_*}e^{-2\kappa d(y,z)-2\kappa d(z,y')}
       &\le c_S(\kappa)e^{-\kappa d(y,y')}.
       \label{I-13-123}
\end{align}
\end{enumerate}
\end{lemma}

\begin{proof}
\emph{Neighbors and local charts.}
Give a site outside $\Omega_0$ level zero. If two neighboring
sites had levels $n$ and $n'\ge n+2$, they would belong to
$\Omega_{n+1}^c$ and $\Omega_{n+2}$, respectively. Separation
would give a distance greater than $R_1\mathfrak M\xi_{n+1}$,
contrary to their distance $\xi$. This proves the scale comparison.
Each half of their connecting bond, apart from its midpoint,
lies inside the bare cube of its endpoint, where the weighted
length density is $\rho(x)^{-1}$ or $\rho(x')^{-1}$.
The midpoint has zero length, so adding the two halves proves
the distance bound. A coordinate path from a bare site to its
cell center stays in that cell and has length at most
$3\xi_y/2$, proving the last assertion of (a).

For (b), $\mathsf S\ge2\mathfrak M\xi_y
\ge2(L+10)\xi_y$. The lift of $C(y)$ lies in an open cube
of side $2(L+10)\xi_y$, on which the covering map is injective.
Two lifted sites with neighboring images differ by a bare
coordinate vector: any nonzero period would have a coordinate
of size at least $\mathsf S-\xi>2(L+9)\xi_y$.
Separation from $\Omega_{n+2}$, and from
$\Omega_{n-1}^c$ when $n\ge1$, gives the stated containments.
The levels of bare sites in $C(y)$ therefore lie between
$n-1$ and $n+1$; at level zero, sites outside $\Omega_0$
also have $\rho=\xi_y$. This proves the bounds for $\rho$.
A monotone coordinate path in the lifted cube has
$|x-y|_1/\xi$ steps. Apply (a) to these steps to obtain
the asserted local metric bound.

\emph{Crossing layers.}
The connecting-set argument in
\cref{subsec:multiscale-sum-proof} applies with these same
partitions and separation. Specifically, if a connected finite
union of segments $\Gamma$ joins $y\in\Lambda_j$ to
$y'\in\Lambda_{j'}$, put
$n=\max\{i:\Gamma\cap\Omega_i\ne\emptyset\}$.
Then $n\ge\max\{j,j'\}$ and
\begin{equation}\label{eq:global-connecting-set}
 |\Gamma|_{\mathbf B}\ge\xi_n^{-1}|y-y'|_1,
 \qquad
 |\Gamma|_{\mathbf B}\ge
 R_1\mathfrak M(n-\min\{j,j'\}-1)_+.
\end{equation}
For clarity, the first inequality uses the $1$-Lipschitz
function $z\mapsto|z-y|_1$ in the interval-image argument
there, before the weaker sup-distance bound is taken.
For the second, each intermediate separating layer contributes
at least $R_1\mathfrak M$ to the weighted length: the
distance to the next region runs through an interval of length
$R_1\mathfrak M\xi_i$ within layer $i$.
Taking the infimum over $\Gamma$ gives
\eqref{eq:global-layer-gap}. Equation
\eqref{eq:global-scale-ratio} follows from
$\xi_y/\xi_{y'}=L^{j-j'}$ and
$|j-j'|\le1+(|j-j'|-1)_+$.

\emph{The site sum.}
Take half of each lower bound in
\eqref{eq:global-connecting-set}, exponentiate, and sum
over the possible largest visited levels. This gives
\[
 e^{-\kappa d(y,y')}
 \le\sum_{n\ge\max\{j,j'\}}
   e^{-\kappa\xi_n^{-1}|y-y'|_1/2}
   e^{-\kappa R_1\mathfrak M(n-\min\{j,j'\}-1)_+/2}.
\]
Choose distinct least-distance lifts of the sites in
$\Lambda_{j'}$ to $\xi_{j'}\mathbb Z^3$.
The elementary one-dimensional bound
\[
 \sum_{k\in\mathbb Z}e^{-\beta|u+hk|}
       \le\frac2{1-e^{-\beta h}}
       \le2+\frac2{\beta h}
\]
implies, by multiplication over the coordinates,
\[
 \sum_{y'\in\Lambda_{j'}}
        e^{-\kappa\xi_n^{-1}|y-y'|_1/2}
       \le(2+4/\kappa)^3L^{3(n-j')}.
\]
Put $p=n-\min\{j,j'\}$. The hypothesis gives
$L^3e^{-\kappa R_1\mathfrak M/2}\le e^{-1}$, so the
last lattice factor times the layer factor is at most
\[
 L^3(2+4/\kappa)^3e^{-(p-1)_+}.
\]
For each $p$, the pair $(j',n)$ is determined by $j'-j$
and $p$, with $|j'-j|\le p$; hence there are at most
$2p+1$ choices. Finally,
\[
 \sum_{p\ge0}(2p+1)e^{-(p-1)_+}<8.
\]
This proves \eqref{I-13-122}. For the convolution bound,
the triangle inequality gives
\[
 e^{-2\kappa(d(y,z)+d(z,y'))}
       \le e^{-\kappa d(y,y')}e^{-\kappa d(y,z)};
\]
sum in $z$ and apply the site estimate.
\end{proof}

\subsection{The scalar input}
\label{sec:multiscale-scalar-input}

Use the auxiliary-torus setting of
\cref{subsec:global-cell-geometry}. For $y\in\Lambda_n$ write
$n(y)=n$ and $B_y=B^n(y)$, and let $S_0$ denote the bare
sites of $\Omega_0$. The scale function $\rho$ is as above.
Write $\langle\cdot,\cdot\rangle_{\mathrm c}$ for the coarse
inner product. The averaging adjoint is
$(Q'^*h)(x)=h(y(x))$, since the fine and coarse weights are
$\xi^3$ and $\xi_y^3$. Thus $Q'Q'^*=I$.

Put $\mathsf W_{\mathrm{av}}=Q'^*\mathsf a'Q'$.
The operator $\Delta'_a=\Delta_{\mathcal S}+\mathsf W_{\mathrm{av}}$
uses the full scalar Dirichlet form, independently of the
bond frame $\mathfrak b_0$. In this subsection $\nabla$ and
$\nabla^*$ denote the ambient scalar gradient and its adjoint;
scalar data are restricted by $\chi_{\mathcal S}$ after a
divergence is taken. All differences use zero extensions.

\begin{proposition}\label{I-13-20}
There are positive constants $\gamma_{\mathrm{lap}}$ and
$\delta_{0,\mathrm{lap}}$, depending only on $L$, with
\[
 \gamma_{\mathrm{lap}}\le1-L^{-2},\qquad
 6\delta_{0,\mathrm{lap}}(L^2+1)
                        \le\tfrac14\gamma_{\mathrm{lap}},
\]
for which the following assertions hold for every nonempty
separated region $\Omega_0$ in the setting above.
\begin{enumerate}[label=(\alph*)]
\item The operator $\Delta'_a$ is symmetric and positive
definite. For every scalar $v$ supported in $S_0$,
\begin{equation}\label{I-13-21}
 \langle v,\Delta'_a v\rangle\ge\|\nabla v\|^2,
 \qquad
 \langle v,\Delta'_a v\rangle
       \ge\gamma_{\mathrm{lap}}\|\rho^{-1}v\|^2.
\end{equation}
Consequently $\mathsf G'$ exists and is symmetric and positive
definite.
\item Suppose
\begin{equation}\label{I-13-59}
 \mathfrak M\ge L+10,\qquad
 R_1\mathfrak M\ge
 \max\{L(L+9),5\delta_{0,\mathrm{lap}}^{-1}\log L\}.
\end{equation}
There is a constant $C_{\mathrm{lap}}$, depending only on $L$,
with the following property. Let $y,y'\in\Lambda_*$, let
$\lambda$ be supported in $B_{y'}$, and let $F_\nu(x)=0$
for $x\notin B_{y'}$. Regard $F_\nu(x)$ as the value of
$F$ on $[x,x+\xi e_\nu]$, and put
$\nabla^*F=\sum_\nu\partial_\nu^*F_\nu$. For every bare
site $x$ with $|x-y|_\infty\le2\xi_y$ and every $\nu$,
\begin{equation}\label{I-13-22}
\begin{aligned}
 |(\mathsf G'\lambda)(x)|
   &\le C_{\mathrm{lap}}\xi_y^2
       e^{-\delta_{0,\mathrm{lap}}d_{\mathbf B}(y,y')/2}
       |\lambda|,\\
 |\partial_\nu(\mathsf G'\lambda)(x)|
   &\le C_{\mathrm{lap}}\xi_y
       e^{-\delta_{0,\mathrm{lap}}d_{\mathbf B}(y,y')/2}
       |\lambda|,\\
 |(\mathsf G'\chi_{\mathcal S}\nabla^*F)(x)|
   &\le C_{\mathrm{lap}}\xi_y
       e^{-\delta_{0,\mathrm{lap}}d_{\mathbf B}(y,y')/2}|F|.
\end{aligned}
\end{equation}
The suprema on the right are over sites and, for $F$, sites
and directions. For $\beta\in\{1/2,3/4\}$ and a real function
$\zeta$ on the bare torus which vanishes when
$|x-y|_\infty\ge\xi_y$, one also has
\begin{equation}\label{eq:global-scalar-divergence-holder}
 \|\zeta\mathsf G'\chi_{\mathcal S}\nabla^*F\|_\beta
 \le C_{\mathrm{lap}}\xi_y^{1-\beta}
      (\|\zeta\|_{\beta,\xi_y}+|\zeta|)
      e^{-\delta_{0,\mathrm{lap}}d_{\mathbf B}(y,y')/2}|F|.
\end{equation}
The constants are independent of $\xi$, the torus side, the
number of layers, $\mathfrak M$, $R_1$, and the regions
satisfying these hypotheses. All differences and seminorms
use zero extensions, including at the boundary of $S_0$.
\end{enumerate}
\end{proposition}

An exponential weight first gives an $L^2$ estimate. The
scalar interior estimate then gives the value and gradient
near an output cell. The separation between layers absorbs
the ratio of the input and output cell sizes into half of
the exponential decay.

\begin{proof}
\emph{Coercivity and the weighted estimate.}
The path-product calculation in
\cref{sec:regional-linear-scalar} gives, for the mean
$\bar v=(Q'v)(y)$ on a cell,
\[
 \xi\sum_{b\subset B_y}|v(b_+)-v(b_-)|^2
        \ge4\xi_y^{-2}\|v-\bar v\|_{B_y}^2.
\]
For a one-site cell the right side is zero. The same
calculation applies here because the inner-bond graph of
$B_y$ is the product of three paths of $L^{n(y)}$ vertices.
Since
\[
 \langle v,\mathsf W_{\mathrm{av}}v\rangle
    =\sum_y a_{n(y)}\xi_y|\bar v|^2,
 \qquad
 \|v\|_{B_y}^2=\|v-\bar v\|_{B_y}^2+\xi_y^3|\bar v|^2,
\]
and $\gamma_{\mathrm{lap}}\le a_n\le1$, summing the cell
inequalities proves the second bound in \eqref{I-13-21}.
The omitted bond terms are nonnegative. The first bound
follows from $\mathsf W_{\mathrm{av}}\ge0$. Positivity and invertibility
follow on the finite-dimensional space of functions on $S_0$.

We next prove a weighted estimate for
$u=\mathsf G'(f+\chi_{\mathcal S}\nabla^*F)$, where $f$ is supported
in $S_0$ and $F$ vanishes on bonds with neither endpoint in
$S_0$. Let $\varphi$ be a function on the torus satisfying
$|\varphi(x)-\varphi(x')|\le d_{\mathbf B}(x,x')$ and let
$0\le\theta\le\delta_{0,\mathrm{lap}}$. On a bond meeting
$S_0$, define $\varphi(b)$ at its initial endpoint if that
endpoint belongs to $S_0$, and at its final endpoint otherwise.
We claim that
\begin{equation}\label{I-13-54}
 \|\rho^{-1}e^{\theta\varphi}u\|
 \le\frac6{\gamma_{\mathrm{lap}}}
       \bigl(\|\rho e^{\theta\varphi}f\|
                       +\|e^{\theta\varphi}F\|\bigr).
\end{equation}
The last norm uses the bond values $\varphi(b)$ just defined.

Put $v=e^{\theta\varphi}u$, extended by zero. Pairing the
equation with $e^{\theta\varphi}v$ and using the ambient
adjoint gives
\begin{equation}\label{I-13-55}
 \langle e^{\theta\varphi}v,
                  \Delta'_a e^{-\theta\varphi}v\rangle
 =\langle v,e^{\theta\varphi}f\rangle
                 +\langle\nabla(e^{\theta\varphi}v),F\rangle.
\end{equation}
We first estimate the conjugated form for any real $v$
supported in $S_0$ and any $0\le\sigma\le1$, without
the narrower restriction on $\theta$. For an inner bond,
write $v_\pm=v(b_\pm)$ and
$h=\sigma(\varphi(b_+)-\varphi(b_-))$. Its contribution
before the factor $\xi$ is
\[
 (v_+-v_-)^2-2v_+v_-(\cosh h-1).
\]
At an exit bond it is the square of the value inside $S_0$.
The neighboring-site bounds in \cref{lem:global-cell-geometry}(a) give
\[
 |\varphi(b_+)-\varphi(b_-)|
     \le\min\{1,L\xi/\rho(b_+),L\xi/\rho(b_-)\}.
\]
Thus $\cosh h-1\le h^2$ and the six bonds incident to each
site give a total loss at most
$6L^2\sigma^2\|\rho^{-1}v\|^2$ in the Laplacian form.

On $B_y$, put $\psi=\varphi-\varphi(y)$ and
$\epsilon=e^{3\sigma/2}-1$. Since $|\psi|\le3/2$ there,
\[
 \left|(Q'e^{\sigma\varphi}v)(y)
              (Q'e^{-\sigma\varphi}v)(y)
                  -(Q'v)(y)^2\right|
 \le\epsilon(2+\epsilon)(Q'|v|)(y)^2.
\]
Here $\epsilon(2+\epsilon)=e^{3\sigma}-1$,
and $\xi_y^3(Q'|v|)(y)^2\le\|v\|_{B_y}^2$.
Adding the two losses gives the estimate
\begin{equation}\label{eq:scalar-conjugation-loss}
 \begin{aligned}
 \langle e^{\sigma\varphi}v,
                   \Delta'_a e^{-\sigma\varphi}v\rangle
 &\ge\langle v,\Delta'_a v\rangle\\
 &\quad-\bigl(6L^2\sigma^2+e^{3\sigma}-1\bigr)
                                      \|\rho^{-1}v\|^2.
 \end{aligned}
\end{equation}
Now take $\sigma=\theta\le\delta_{0,\mathrm{lap}}$.
The smallness condition gives $e^{3\theta}-1\le6\theta$.
Using the mean of the two bounds in \eqref{I-13-21},
we obtain
\begin{equation}\label{I-13-56}
 \langle e^{\theta\varphi}v,
                  \Delta'_a e^{-\theta\varphi}v\rangle
 \ge\tfrac12\|\nabla v\|^2
                  +\tfrac14\gamma_{\mathrm{lap}}\|\rho^{-1}v\|^2.
\end{equation}

For the last term of \eqref{I-13-55}, set
$h=\theta(\varphi(b_+)-\varphi(b_-))$ on an inner bond.
Then
\[
 \nabla(e^{\theta\varphi}v)(b)
 =e^{\theta\varphi(b)}
       \left(\nabla v(b)+\xi^{-1}(e^h-1)v_+\right),
 \qquad
 \xi^{-1}|e^h-1|\le e\theta L\rho(b_+)^{-1}.
\]
At an exit bond there is only the weighted boundary value,
which is already included in $\|\nabla v\|$. A site is the
final endpoint of three positive bonds. Since
$e\theta L\sqrt3\le1$, Cauchy--Schwarz gives
\[
 |\langle\nabla(e^{\theta\varphi}v),F\rangle|
 \le\|e^{\theta\varphi}F\|
                 \bigl(\|\nabla v\|+\|\rho^{-1}v\|\bigr).
\]
Writing $X=\|\rho^{-1}v\|$, $Y=\|\nabla v\|$,
$a=\|\rho e^{\theta\varphi}f\|$ and
$b=\|e^{\theta\varphi}F\|$, we obtain
\[
 \tfrac12Y^2+\tfrac14\gamma_{\mathrm{lap}}X^2
 \le(a+b)X+bY
 \le\tfrac18\gamma_{\mathrm{lap}}X^2
       +\frac2{\gamma_{\mathrm{lap}}}(a+b)^2
       +\tfrac12Y^2+\tfrac12b^2.
\]
Hence $X^2\le32\gamma_{\mathrm{lap}}^{-2}(a+b)^2$,
which proves \eqref{I-13-54}.

\emph{From local energy to pointwise and H\"older bounds.}
Assume now \eqref{I-13-59}. Use the cube $C(y)$ of \cref{lem:global-cell-geometry}(b),
and put
\[
 B^*(y)=\{x\in S_0:|x-y|_\infty\le(L+8)\xi_y\},
\]
with norms over bare sites or bonds. Set
\[
 N_y=\xi_y^{-3/2}\|u\|_{B^*(y)}
                       +\xi_y^2|f|_{C(y)}+\xi_y|F|_{C(y)}.
\]
Here $|F|_{C(y)}$ takes the supremum over bonds with an
endpoint in $C(y)$. For $x,x'$ within $2\xi_y$ of $y$
in the supremum distance and $\beta\in\{1/2,3/4\}$, the local
estimates we need are
\begin{equation}\label{eq:scalar-local-pointwise}
 \begin{aligned}
 |u(x)|&\le C_E N_y,\\
 |u(x)-u(x')|&\le C_E
       (|x-x'|_1/\xi_y)^\beta N_y,\\
 \xi_y|\partial_\nu u(x)|&\le C_E N_y
                    &&\text{if }|F|_{C(y)}=0,
 \end{aligned}
\end{equation}
with $C_E$ depending only on $L$, after taking the maximum of
the constants for the two exponents.

If $n(y)=0$, every site within $3\xi$ of $y$ either lies
outside $S_0$, where $u=0$, or belongs to $B^*(y)$. Its
value is at most $\xi^{-3/2}\|u\|_{B^*(y)}$. This bounds
the value and, by the two endpoint values, $\xi$ times its
first difference. Distinct bare sites have distance at least
$\xi=\xi_y$, which gives the H\"older bound as well.

If $n(y)\ge1$, \cref{lem:global-cell-geometry}(b)
identifies the bare sites of $C(y)$ with those of a Euclidean
cube and places the cube inside $\Omega_0$.
Take $r=8\xi_y$. Then $r\ge16\xi$, and all neighbors of
the cube $\mathsf B_r$ centered at $y$ lie in $C(y)$.
There the equation reads
\[
 \Delta u=g+\nabla^*F,\qquad g=f-\mathsf W_{\mathrm{av}}u.
\]
Every cell meeting $\mathsf B_r$ lies in $B^*(y)$: its
side is at most $L\xi_y$. Its side is also at least
$\xi_y/L$, so the averaging inequality gives
\[
 |\mathsf W_{\mathrm{av}}u|_{\mathsf B_r}
       \le L^{7/2}\xi_y^{-7/2}\|u\|_{B^*(y)}.
\]
Apply \cref{I-13-47} after restricting $F$ to the bonds
with an endpoint in $\mathsf B_r$. This restriction leaves
its divergence on $\mathsf B_r$ unchanged. The quantity
$\mathsf N$ in that proposition is bounded by
$C(L)N_y$. Its value and $\beta$-H\"older estimates on
$\mathsf B_{r/4}$ give the first two inequalities of
\eqref{eq:scalar-local-pointwise}; if $|F|_{C(y)}=0$, its
gradient estimate gives the third. This proves the local
estimates, including the boundary case at level zero.

\emph{Decay between cells.}
For this step allow $f$ to be supported in the sites of
$B_{y'}$ and their neighbors, and $F$ on bonds with an
endpoint in $B_{y'}$; retain the support conditions at
$S_0$ used in \eqref{I-13-54}. Set
$\varphi(x)=d_{\mathbf B}(x,y')$ and
$\theta=\delta_{0,\mathrm{lap}}$. The metric triangle
inequality permits this weight. On the support of $f$,
$\varphi\le5/2$ and $\rho\le L\xi_{y'}$; the same
bound for $\varphi(b)$ holds on the support of $F$.
There are at most $7(\xi_{y'}/\xi)^3$ such sites and
$6(\xi_{y'}/\xi)^3$ such bonds. Therefore
\eqref{I-13-54} implies
\begin{equation}\label{I-13-62}
 \|\rho^{-1}e^{\delta_{0,\mathrm{lap}}\varphi}u\|
 \le C\bigl(\xi_{y'}^{5/2}|f|+\xi_{y'}^{3/2}|F|\bigr).
\end{equation}
Throughout this step $C$ depends only on $L$.

For $x\in B^*(y)$, the local geometric bounds give
$\rho(x)\le L\xi_y$ and
$d_{\mathbf B}(x,y)\le3L(L+8)$. Hence
\begin{align*}
 \xi_y^{-3/2}\|u\|_{B^*(y)}
 &\le C\left[
     \xi_y^2\left(\frac{\xi_{y'}}{\xi_y}\right)^{5/2}|f|
     +\xi_y\left(\frac{\xi_{y'}}{\xi_y}\right)^{3/2}|F|
          \right]e^{-\delta_{0,\mathrm{lap}}d_{\mathbf B}(y,y')}.
\end{align*}
The layer-crossing bound \eqref{eq:global-layer-gap} and
\eqref{I-13-59} imply, for $0\le p\le5/2$,
\[
 \left(\frac{\xi_{y'}}{\xi_y}\right)^p
 \le L^p L^{p(|n(y)-n(y')|-1)_+}
 \le L^p e^{\delta_{0,\mathrm{lap}}d_{\mathbf B}(y,y')/2}.
\]
It follows that the first term in $N_y$ is at most
\[
 C(\xi_y^2|f|+\xi_y|F|)
           e^{-\delta_{0,\mathrm{lap}}d_{\mathbf B}(y,y')/2}.
\]

The remaining terms in $N_y$ vanish unless the support of
$f$, or an endpoint of a bond on which $F$ is nonzero,
meets $C(y)$. In that event choose such an endpoint or
site $x_1$. It belongs to $B_{y'}$ or is a neighbor of a
site in $B_{y'}$. Thus
\[
 d_{\mathbf B}(y,y')
    \le d_{\mathbf B}(y,x_1)+d_{\mathbf B}(x_1,y')
    \le3L(L+9)+\tfrac52.
\]
The smallness condition on $\delta_{0,\mathrm{lap}}$ gives
$1\le e^3e^{-\delta_{0,\mathrm{lap}}d_{\mathbf B}(y,y')/2}$.
Consequently all of $N_y$ is bounded by the preceding
display, with a larger constant $C$. Applying
\eqref{eq:scalar-local-pointwise} proves the value and
$\beta$-H\"older estimates for $u$, and its gradient estimate
when $F=0$.

Finally take $(f,F)=(\lambda,0)$ and then $(f,F)=(0,F)$.
In the second application, the support at initial sites in
$B_{y'}$ is contained in the incident-bond support used in
the proof, and every such bond has an endpoint in $S_0$.
These two applications give all three inequalities in
\eqref{I-13-22}. Their range includes both endpoints of
every bond meeting $B_y$, since each such endpoint lies
within $\xi_y/2+\xi\le2\xi_y$ of $y$.

For the cutoff estimate put $u=\mathsf G'\chi_{\mathcal S}\nabla^*F$
and $E=e^{-\delta_{0,\mathrm{lap}}d_{\mathbf B}(y,y')/2}$.
The local bounds just proved have constants
$D_0,D_1\le C\xi_yE|F|$ for the value and $\beta$-H\"older
difference. If $\zeta(x)\ne0$ and $h=|x-x'|_1\le\xi_y$,
both points lie within $2\xi_y$ of $y$, and the product identity gives
\[
 |\zeta(x)u(x)-\zeta(x')u(x')|
 \le\bigl(|\zeta|D_1+\|\zeta\|_{\beta,\xi_y}D_0\bigr)
                           (h/\xi_y)^\beta.
\]
If both cutoff values vanish there is nothing to prove;
if $h>\xi_y$, use $2|\zeta|D_0\le2|\zeta|D_0(h/\xi_y)^\beta$.
Dividing by $h^\beta$ and taking the supremum over
$0<h\le1$ proves \eqref{eq:global-scalar-divergence-holder}.
\end{proof}

\subsection{The scalar Gram inverse and differentiated projection}
\label{sec:multiscale-gram}

Here $\partial$ maps scalar functions to all positive star
bonds of $\Omega_0$, with adjoint $\partial^*$ for the fine
inner products of weight $\xi^3$. Let $\Xi$ multiply a fine
function by $\xi_x=\rho(x)$ and a coarse function by $\xi_y$.
Put
\[
 \mathsf T=\mathsf G'Q'^*,\qquad
 \mathsf V=\mathsf T^*\mathsf T=Q'\mathsf G'^2Q'^*,\qquad
 \mathsf C=\mathsf V^{-1},\qquad S=\Xi^{-2}\mathsf V\Xi^{-2}.
\]
The block coordinates and weighted pairings are exactly those
used in the proofs of \cref{I-11-30}(b) and \eqref{I-11-36}.
Since $Q'Q'^*=I$ and $\mathsf G'$ is positive definite, the
same positivity calculation gives $\mathsf V>0$ on the
auxiliary tori. The projection calculation also applies:
$\Delta'_a=\Delta_{\mathcal S}$ on $\ker Q'$.
For a coarse operator $K$, distinguish its counting kernel
$K(y,y')$ from its matrix
\[
 K[y,y']=\langle e_y,Ke_{y'}\rangle_{\mathrm c}
       =\xi_y^{3/2}K(y,y')\xi_{y'}^{-3/2},
 \qquad e_y=\xi_y^{-3/2}\mathbf1_{\{y\}}.
\]
The normalization by $\Xi^{-2}$ removes the fourth-order scale
of $\mathsf V$. We prove a uniform gap for $S$ using functions
supported on individual blocks with prescribed averages. Exponential
conjugation then gives decay of its inverse.

Use the localization point fixed above also for all star bonds.
Call a bond attached to $y$ when this point lies in $B_y$.

\begin{proposition}\label{I-14-33}\label{I-14-48}
There are positive constants
$\gamma_{0,\mathrm{gp}},\delta_1,\delta_2,\mathfrak r_4,
\mathfrak m_5,\mathfrak r_5,B_{5,\mathrm{gp}}$, depending only
on $L$, with the following properties for every separated
sequence in the stated setting.

\textup{(i)} The Gram operator satisfies
$\mathsf V\ge\gamma_{0,\mathrm{gp}}\Xi^4$.
If $R_1\mathfrak M\ge\mathfrak r_4$, then
\begin{equation}\label{I-14-34}
 \begin{split}
 |S^{-1}[y,y']|&\le\frac2{\gamma_{0,\mathrm{gp}}}
                   e^{-3\delta_1d_{\mathbf B}(y,y')},\\
 |\mathsf C(y,y')|&\le\frac2{\gamma_{0,\mathrm{gp}}}
        \xi_y^{-7/2}\xi_{y'}^{-1/2}
                   e^{-3\delta_1d_{\mathbf B}(y,y')}.
 \end{split}
\end{equation}
In particular,
\begin{equation}\label{I-14-35}
 |\mathsf C(y,y')|\le\frac{2L^{1/2}}{\gamma_{0,\mathrm{gp}}}
             \xi_y^{-4}e^{-2\delta_1d_{\mathbf B}(y,y')}.
\end{equation}
The weighted kernel
$\xi_y^3\mathsf C(y,y')$ is symmetric in $y,y'$.

\textup{(ii)} The orthogonal projections
$\mathsf P,\mathsf R=1-\mathsf P$
satisfy
\[
 \mathsf P=\mathsf G'Q'^*\mathsf CQ'\mathsf G',
 \qquad \operatorname{ran}\mathsf R=\Delta_{\mathcal S}\ker Q'.
\]
The range of $\mathsf R$, regarded as functions on the whole bare
lattice by zero extension, is supported in $\Omega_1$ and its
bare nearest neighbours. If two separated sequences have the same
$\Omega_1,\ldots,\Omega_t$, their kernels $\mathsf R$ agree at all
bare sites of the intersection of their outer regions.

\textup{(iii)} Suppose that $\mathfrak M\ge\mathfrak m_5$ and
$R_1\mathfrak M\ge\mathfrak r_5$. For $x\in B_y$, $x'\in B_{y'}$,
and star bonds $b,b'$ attached to $y,y'$, respectively,
\begin{equation}\label{I-14-49}
 \begin{gathered}
 |\mathsf P(x,x')|,\quad |(\partial\mathsf P)(b,x')|,\quad
 |(\mathsf P\partial^*)(x,b')|,\quad
 |(\partial\mathsf P\partial^*)(b,b')|\\
 \le B_{5,\mathrm{gp}}
       \left(\frac{\xi}{\xi_{y'}}\right)^3
       [1,\xi_y^{-1},\xi_y^{-1},\xi_y^{-2}]
       e^{-\delta_2d_{\mathbf B}(y,y')}.
 \end{gathered}
\end{equation}
The four factors in brackets correspond to the four kernels.
All kernels are taken against counting measure, and
$\mathsf R(x,x')=\mathbf1[x=x']-\mathsf P(x,x')$.
\end{proposition}

The mixed derivative in \eqref{I-14-49} has output scale
$\xi_y^{-2}$ even when the input block has a different size;
the scale-ratio estimate absorbs that difference into exponential
decay.

\phantomsection\label{sec:multiscale-gram-proof}

\begin{proof}
\emph{Weighted scalar energy.}
Put $c_{0,\mathrm{gp}}=1-L^{-2}$ and choose
$0<\varepsilon_0\le1$, depending only on $L$, such that
\begin{equation}\label{eq:multiscale-scalar-smallness}
 6L^2\varepsilon_0^2+e^{3\varepsilon_0}-1
          \le\tfrac14c_{0,\mathrm{gp}}.
\end{equation}
Write $A=\Delta'_a$ and $G=\mathsf G'$ in this proof.
The block argument proving \eqref{I-11-31} uses only the
coordinates and pairings just checked, so
$\langle u,Au\rangle\ge c_{0,\mathrm{gp}}\|\Xi^{-1}u\|^2$.
Let $\theta$ be real and satisfy
$|\theta(x)-\theta(x')|\le\varepsilon_0d_{\mathbf B}(x,x')$.
Apply \eqref{eq:scalar-conjugation-loss} with
$\sigma=\varepsilon_0$ and $\varphi=\theta/\varepsilon_0$.
The loss is at most
$c_{0,\mathrm{gp}}\|\Xi^{-1}u\|^2/4$.
Use half of the original form to bound $\|\partial u\|^2$
and half to bound $c_{0,\mathrm{gp}}\|\Xi^{-1}u\|^2$.
This gives
\begin{equation}\label{eq:multiscale-scalar-conjugation}
 \langle u,e^\theta Ae^{-\theta}u\rangle
 \ge\tfrac12\|\partial u\|^2
       +\tfrac14c_{0,\mathrm{gp}}\|\Xi^{-1}u\|^2.
\end{equation}
If $v=e^\theta Ge^{-\theta}f$, its pairing with $f$ yields
\begin{equation}\label{eq:multiscale-scalar-weighted-L2}
 \|\Xi^{-1}v\|\le\frac4{c_{0,\mathrm{gp}}}\|\Xi f\|,
 \qquad
 \|\partial v\|\le\sqrt{\frac8{c_{0,\mathrm{gp}}}}\|\Xi f\|.
\end{equation}

For $f$ supported in $B_{y'}$, take
$\theta=\varepsilon_0d_{\mathbf B}(\,\cdot\,,y')$ and
$v=e^\theta Gf$. Put $r=d_{\mathbf B}(y'',y')$.
The weight is at most $3\varepsilon_0/2$ on $B_{y'}$,
at least $\varepsilon_0(r-3/2)$ on $B_{y''}$, and at least
$\varepsilon_0(r-5/2)$ at the endpoints in $\Omega_0$ of
bonds touching $B_{y''}$. For an internal bond, writing
$\theta_\pm=\theta(b_\pm)$,
\[
 \partial(e^{-\theta}v)(b)=e^{-\theta_-}
 \left(\partial v(b)
    -\xi^{-1}(1-e^{\theta_--\theta_+})v(b_+)\right),
 \qquad
 \xi^{-1}|1-e^{\theta_--\theta_+}|
       \le2\varepsilon_0L\xi_{b_+}^{-1}.
\]
For an exit bond there is only the weighted difference term.
Each site is the terminal endpoint of three positive bonds.
The triangle inequality therefore gives
\[
 \|\partial Gf\|_{\operatorname{st}(B_{y''})}
 \le e^{-\varepsilon_0(r-5/2)}
       \bigl(\|\partial v\|
               +2\varepsilon_0L\sqrt3\,\|\Xi^{-1}v\|\bigr).
\]
Together with \eqref{eq:multiscale-scalar-weighted-L2},
applied to $e^\theta f$, and
$\|\Xi e^\theta f\|\le
e^{3\varepsilon_0/2}\xi_{y'}\|f\|$, this proves
\begin{equation}\label{eq:multiscale-scalar-block-L2}
 \left(\|\Xi^{-1}Gf\|_{B_{y''}}^2
       +\|\partial Gf\|_{\operatorname{st}(B_{y''})}^2\right)^{1/2}
 \le C_{3,\mathrm{gp}}\xi_{y'}
                 e^{-\varepsilon_0r}\|f\|.
\end{equation}
Here $C_{3,\mathrm{gp}}$ depends only on $L$.
In particular,
$\|Gf\|_{B_{y''}}\le C_{3,\mathrm{gp}}\xi_{y''}\xi_{y'}
e^{-\varepsilon_0r}\|f\|$.

\emph{Upper bound for the normalized Gram matrix.}
Let $\chi_y$ be the indicator of $B_y$. Since
$Q'^*e_y=\xi_y^{-3/2}\chi_y$,
\begin{equation}\label{I-13-128}
 S[y,y']=\xi_y^{-7/2}\xi_{y'}^{-7/2}
                   \langle G\chi_y,G\chi_{y'}\rangle.
\end{equation}
The blocks partition $\Omega_0$, and
$\|\chi_y\|=\xi_y^{3/2}$. Applying the last bound on each
block gives, with $r=d_{\mathbf B}(y,y'')$ and
$r'=d_{\mathbf B}(y'',y')$,
\[
 |S[y,y']|\le C_{3,\mathrm{gp}}^2
 \sum_{y''}\frac{\xi_{y''}}{\xi_y}
           \frac{\xi_{y''}}{\xi_{y'}}e^{-\varepsilon_0(r+r')}.
\]
If $R_1\mathfrak M\ge8(3\log L+1)/\varepsilon_0$, the
scale-ratio estimate \eqref{eq:global-scale-ratio} bounds the quotients by
$L^2e^{\varepsilon_0(r+r')/4}$. The triangle inequality
and \eqref{I-13-122} at rate $\varepsilon_0/4$ then give
\begin{equation}\label{eq:normalized-gram-upper}
 |S[y,y']|\le C_Se^{-\varepsilon_0d_{\mathbf B}(y,y')/2},
 \qquad
 C_S=C_{3,\mathrm{gp}}^2L^2c_S(\varepsilon_0/4).
\end{equation}

\emph{The uniform Gram gap.}
Define $g(s)=1+\cos(2\pi s)$ for $|s|\le1/2$ and
$g(s)=0$ otherwise. Then $0\le g\le2$, $g$ is continuously
differentiable, and $g'$ is Lipschitz with constant $4\pi^2$.
For $h>0$,
\[
 |g(s+h)+g(s-h)-2g(s)|
 \le\int_0^h4\pi^2(2r)\,dr=4\pi^2h^2.
\]
For $n(y)\ge1$ put
$\varphi_y(x)=\prod_{\mu=1}^3g((x_\mu-y_\mu)/\xi_y)$,
using coordinate representatives of least absolute value on
the torus; for $n(y)=0$ put $\varphi_y=\chi_y$.
The torus side is at least $2L\xi_y$, so the choice of
representative is immaterial wherever a factor is nonzero.
The odd integer $N=\xi_y/\xi$ gives block coordinates
$-(N-1)/2,\ldots,(N-1)/2$. Points outside the block have
some coordinate of modulus at least $(N+1)\xi/2$;
thus $\varphi_y$ is supported on $B_y$. For $N>1$ these
coordinates form a complete set of residues modulo $N$, and
the sum of $\cos(2\pi m/N)$ over them is zero. Taking the
product of the three one-dimensional sums gives
\begin{equation}\label{I-13-139}
 (Q'\varphi_y)(y')=\mathbf1[y=y'].
\end{equation}
The same identity holds for the singleton blocks.

Let $B_y^+$ consist of $B_y$ and its bare nearest neighbours.
Each coordinate second difference of $\varphi_y$ consists of
two factors bounded by $2$ and a second difference of $g$
with step $\xi/\xi_y$. A difference crossing the chosen
torus representative has both relevant factors zero, so the
same estimate applies there. For a singleton block use
$|\Delta \chi_y|\le6\xi^{-2}$. Hence
\begin{equation}\label{I-13-140}
 |\Delta \varphi_y|\le C_\varphi\xi_y^{-2},\qquad
 \operatorname{supp}\Delta \varphi_y\subset B_y^+,\qquad
 \|\Delta \varphi_y\|^2\le7C_\varphi^2\xi_y^{-1}.
\end{equation}
Here $C_\varphi$ is an absolute constant, and the last
inequality uses the bound $N^3+6N^2\le7N^3$ on the
number of sites of $B_y^+$.

Given $\omega'\in\mathcal Y$, let
$\psi=\sum_y\omega'(y)\varphi_y$. Then
$Q'\psi=\omega'$ and
\[
 A\psi=\chi_{\mathcal S}\Delta \psi
       +\sum_y a_{n(y)}\xi_y^{-2}\omega'(y)\chi_y.
\]
This identity uses the full star form: every bond on which
$\psi$ has a nonzero difference is included. Each bare site
belongs to at most seven sets $B_y^+$, since its block and
the blocks of its six neighbours are the only possibilities.
Thus \eqref{I-13-140} and the disjointness of the blocks imply
\[
 \|\chi_{\mathcal S}\Delta \psi\|
       \le7C_\varphi\|\Xi^{-2}\omega'\|_{\mathrm c},
 \qquad
 \|A\psi\|\le C_{2,\mathrm{gp}}\|\Xi^{-2}\omega'\|_{\mathrm c}.
\]
Here $C_{2,\mathrm{gp}}\ge1$ is an absolute constant.
Take $\omega'=\Xi^4\omega$. Symmetry of $A$ and the
averaging adjoint identity give
\[
 \|\Xi^2\omega\|_{\mathrm c}^2
 =\langle Q'\psi,\omega\rangle_{\mathrm c}
 =\langle A\psi,GQ'^*\omega\rangle
 \le C_{2,\mathrm{gp}}\|\Xi^2\omega\|_{\mathrm c}
                         \|GQ'^*\omega\|.
\]
Dividing and squaring, with the zero vector treated separately,
proves, for $\gamma_{0,\mathrm{gp}}=C_{2,\mathrm{gp}}^{-2}$,
\begin{equation}\label{eq:normalized-gram-gap}
 \langle\omega,\mathsf V\omega\rangle_{\mathrm c}
       \ge\gamma_{0,\mathrm{gp}}\|\Xi^2\omega\|_{\mathrm c}^2,
 \qquad S\ge\gamma_{0,\mathrm{gp}}.
\end{equation}

\emph{Decay of the inverse.}
Put $\kappa=\varepsilon_0/2$ and choose $\kappa'>0$ so that
\begin{equation}\label{eq:direct-gram-smallness}
 \begin{gathered}
 \kappa'\le\kappa/4,\qquad
 \frac{4\kappa'C_Sc_S(\kappa/2)}{e\kappa}
                    \le\tfrac12\gamma_{0,\mathrm{gp}}.
 \end{gathered}
\end{equation}
Choose $0<3\delta_1\le\min\{\kappa',\varepsilon_0/16\}$
and $\mathfrak r_4$ large enough that
\[
 \mathfrak r_4\ge\frac{8(3\log L+1)}{\varepsilon_0},
 \qquad \frac{\log L}{\mathfrak r_4}\le\frac{\delta_1}{20}.
\]
These constants depend only on $L$. Under
$R_1\mathfrak M\ge\mathfrak r_4$, the upper Gram bound and
the site sum at rate $\kappa/2$ apply.
Fix $y'$ and conjugate the matrix of $S$ on counting
$\ell^2(\Lambda_*)$ by
$\theta(y)=\kappa'd_{\mathbf B}(y,y')$; thus
$S_\theta=e^\theta Se^{-\theta}$.
For $r=d_{\mathbf B}(y_1,y_2)$,
\[
 |(S_\theta-S)[y_1,y_2]|
 \le\kappa'r e^{\kappa'r}C_Se^{-\kappa r}
 \le\frac{4\kappa'C_S}{e\kappa}e^{-\kappa r/2}.
\]
Here $\kappa'\le\kappa/4$ and
$re^{-\kappa r/4}\le4/(e\kappa)$.
The row and column sums are at most
$4\kappa'C_Sc_S(\kappa/2)/(e\kappa)\le\gamma_{0,\mathrm{gp}}/2$.
The Cauchy--Schwarz row-and-column estimate bounds the
operator norm by the same number. Thus
$\langle v,S_\theta v\rangle\ge
(\gamma_{0,\mathrm{gp}}/2)\|v\|_{\ell^2}^2$, and
$\|S_\theta^{-1}\|\le2/\gamma_{0,\mathrm{gp}}$.
Undoing the conjugation gives the first bound in
\eqref{I-14-34}, since $\kappa'\ge3\delta_1$.
Finally,
\[
 \mathsf C(y,y')
 =\xi_y^{-7/2}\xi_{y'}^{-1/2}S^{-1}[y,y']
 =\xi_y^{-4}(\xi_y/\xi_{y'})^{1/2}S^{-1}[y,y'].
\]
The first identity gives the second bound in \eqref{I-14-34};
the scale-ratio estimate and $\varrho/2\le\delta_1$ give
\eqref{I-14-35}. Symmetry of $\mathsf V^{-1}$ in the weighted
coarse inner product gives the asserted weighted symmetry.

\emph{Projection and outer-region independence.}
The projection formula and range are \eqref{I-11-36}, with
$D=A$. By \eqref{eq:regional-null-support}, every member of
$N=\ker Q'$ vanishes outside $\Omega_1$, and its
ambient Laplacian is supported in $\Omega_1$ and its bare
nearest neighbours. These neighbours lie in $\Omega_0$.
Hence $\Delta_{\mathcal S} N=\Delta N$ as spaces of functions
on the full bare lattice. Moreover $N$ consists exactly of
the functions supported in $\Omega_1$ whose averages on
the positive-level blocks vanish; this description involves
only $\Omega_1,\ldots,\Omega_t$. The ambient orthogonal
projection onto $\Delta N$ is therefore independent of
the outer region. Its restriction to
$\mathcal S$ is $\mathsf R$, which proves the
assertion about its kernels.

\emph{The differentiated kernels.}
Use the constants $\delta_{0,\mathrm{lap}},C_{\mathrm{lap}}$
of \cref{I-13-20}. Set
\[
 \delta_2=\tfrac18\min\{\delta_{0,\mathrm{lap}},\delta_1\}.
\]
Choose $\mathfrak m_5,\mathfrak r_5$ so their hypotheses
imply \eqref{I-13-59} and
\[
 \mathfrak r_5\ge\max\{\mathfrak r_4,
                         10(\log L+1)/\delta_2\}.
\]
For a block centre $y_1$ put
$F_0(x;y_1)=(G\chi_{y_1})(x)$ and
$F_1(b;y_1)=(\partial G\chi_{y_1})(b)$.
The pointwise scalar estimate gives
\begin{equation}\label{I-14-50}
 |F_0(x;y_1)|\le C_{\mathrm{lap}}\xi_y^2
                e^{-\delta_{0,\mathrm{lap}}d_{\mathbf B}(y,y_1)/2},
 \qquad
 |F_1(b;y_1)|\le C_{\mathrm{lap}}\xi_y
                e^{-\delta_{0,\mathrm{lap}}d_{\mathbf B}(y,y_1)/2}.
\end{equation}
For an exit bond the endpoint $x$ in $\Omega_0$ lies outside
$\Omega_1$, by separation, and therefore $\xi_y=\xi$.
The other endpoint value of $G\chi_{y_1}$ is zero, so
$|F_1(b;y_1)|=\xi^{-1}|F_0(x;y_1)|$; this proves the
second bound for exit bonds as well.

The right factors require the volume weights. Symmetry of
$G$ and the adjoint relation give
\[
 (Q'G\mathbf1_{\{x'\}})(y_2)
       =(\xi/\xi_{y_2})^3F_0(x';y_2),\qquad
 (Q'G\partial^*\mathbf1_{\{b'\}})(y_2)
       =(\xi/\xi_{y_2})^3F_1(b';y_2).
\]
For the second identity, pair $G\chi_{y_2}$ with
$\partial^*\mathbf1_{\{b'\}}$, move $\partial$ to the
first factor, and use the fine bond weight $\xi^3$.
Since $Q'^*\omega=\sum_y\omega(y)\chi_y$, each
of the four projection kernels is
\begin{equation}\label{I-14-51}
 \sum_{y_1,y_2}F_\alpha(\,\cdot\,;y_1)\,
       \mathsf C(y_1,y_2)
       (\xi/\xi_{y_2})^3F_\beta(\,\cdot'\,;y_2),
       \qquad \alpha,\beta\in\{0,1\}.
\end{equation}
Here $\alpha=1$ denotes the left derivative, $\beta=1$ the
right adjoint derivative, and the corresponding arguments
are bonds.

Put $\delta=2\delta_2$ and
$r_1=d_{\mathbf B}(y,y_1)$,
$r_2=d_{\mathbf B}(y_1,y_2)$,
$r_3=d_{\mathbf B}(y_2,y')$.
By \eqref{I-14-34} and \eqref{I-14-50}, the summand is
bounded by $2C_{\mathrm{lap}}^2/\gamma_{0,\mathrm{gp}}$ times
\[
 \xi^3\xi_y^{2-\alpha}\xi_{y_1}^{-7/2}
       \xi_{y_2}^{-7/2}\xi_{y'}^{2-\beta}
       e^{-2\delta(r_1+r_2+r_3)}.
\]
The scale product has the exact factorization
\begin{equation}\label{I-14-52}
 \xi_y^{2-\alpha}\xi_{y_1}^{-7/2}
       \xi_{y_2}^{-7/2}\xi_{y'}^{2-\beta}
 =\xi_{y'}^{-3}\xi_y^{-\alpha-\beta}
       \left(\frac{\xi_y}{\xi_{y_1}}\right)^{2+\beta}
       \left(\frac{\xi_{y'}}{\xi_{y_1}}\right)^{3/2-\beta}
       \left(\frac{\xi_{y'}}{\xi_{y_2}}\right)^{7/2}.
\end{equation}
Equation \eqref{eq:global-scale-ratio} bounds the last three factors by
$L^7e^{5\varrho(r_1+r_2+r_3)}$. Our threshold gives
$5\varrho\le\delta$. By the triangle inequality,
\[
 e^{-2\delta(r_1+r_2+r_3)}
       e^{5\varrho(r_1+r_2+r_3)}
 \le e^{-\delta d_{\mathbf B}(y,y')/2}
       e^{-\delta r_1/2}e^{-\delta r_2/2}.
\]
Sum first over $y_2$ and then over $y_1$. Equation \eqref{I-13-122}
at rate $\delta/2=\delta_2$ applies because
$\mathsf r_S(\delta_2)\le10(\log L+1)/\delta_2
\le R_1\mathfrak M$. It contributes $c_S(\delta_2)^2$.
The remaining factors are
$(\xi/\xi_{y'})^3\xi_y^{-\alpha-\beta}$, proving
\eqref{I-14-49} with $B_{5,\mathrm{gp}}$ depending only on $L$.
\end{proof}

For the local comparisons below, put
$\mathcal P=\partial_0\mathsf P\partial_0^*$ and choose common constants
\begin{equation}\label{I-17-23}
 \begin{gathered}
 \delta_{\mathrm u}=\delta_2,\qquad
 B_{\mathrm u}=\max\{1,C_{\mathrm{lap}},
                 2L^{1/2}/\gamma_{0,\mathrm{gp}},B_{5,\mathrm{gp}}\},\\
 \mathfrak m_{\mathrm u}=\mathfrak m_5,\qquad
 \mathfrak r_{\mathrm u}=\mathfrak r_5.
 \end{gathered}
\end{equation}
The choices of $\mathfrak m_5,\mathfrak r_5$ imply the scalar
and Gram hypotheses, and $\delta_2$ is no larger than their
rates. Thus \eqref{I-13-22}, \eqref{I-14-35} and the fourth
kernel bound in \eqref{I-14-49} hold with the common rate
$\delta_{\mathrm u}$ and constant $B_{\mathrm u}$.
Since $\partial_0=\chi_{\mathcal A}\partial\chi_{\mathcal S}$,
the kernel of $\mathcal P$ is the restriction of
$\partial\mathsf P\partial^*$ to $\mathfrak b_0$ in both
variables. The localization points agree, so this restriction
preserves the fourth bound.

\subsection{The local models}
\label{subsec:global-local-models}

We first choose a quadratic partition whose members meet at most
two consecutive layers. On its interior supports the local model
is a two-scale operator on a larger torus. At the outer boundary
we use a compressed curl operator with a mass. Throughout this
subsection assume $\mathfrak M\ge L^5$; all further lower bounds
on $\mathfrak M$ will depend only on $L$.

Use $\vartheta_0,\eta_0$ from
\cref{sec:scalar-interior-proof}, and fix
\[
 h(s)=\cos\!\left(\tfrac\pi2
                 \vartheta_0(4|s|-\tfrac32)\right),\qquad
 \varpi(v)=\prod_{\mu=1}^3
                  \frac{\eta_0(v_\mu)}{\int\eta_0},\qquad
 \varpi_r(v)=r^{-3}\varpi(v/r).
\]
Thus $h$ is smooth, equals one for $|s|\le3/8$, and vanishes
for $|s|\ge5/8$. The identity
$\vartheta_0(s)+\vartheta_0(1-s)=1$ gives
$h(s)^2+h(s-1)^2=1$ for $0\le s\le1$. All other integer
translates vanish there, so
$\sum_{i\in\mathbb Z}h(s-i)^2=1$ for every $s$.
The smooth nonnegative function $\varpi$ has integral one and
support in $\{|v|_\infty\le2/3\}$. Convolutions below use
the periodic lift of an indicator function.

For $0\le n\le t_*$ put $W_n=\mathfrak M\xi_n$ and
$r_n=W_n/(6L)$. Let $\mathcal Z_n$ be the corners of the
cubes of $\pi[W_n]$. For $z\in\mathcal Z_n$ set
\[
 \square=\{x:|x-z|_\infty\le W_n\},\qquad
 \widetilde\square^{i}=\{x:|x-z|_\infty\le(1+i)W_n\},
 \qquad
 g_\square(x)=\prod_{\mu=1}^3
                         h((x_\mu-z_\mu)/W_n).
\]
Coordinate differences in this formula are their periodic
representatives. A cube is indexed by $(n,z)$, even if two
such cubes have the same point set. This convention matters
when the torus side is $2W_n$.

Put $d_n(x)=\operatorname{dist}_\infty(x,\Omega_n^c)$,
with values $+\infty$ and $0$ for a full and an empty region,
respectively, and define
\[
 \Psi_n=\mathbf1_{\{d_n\ge3r_n/2\}}*\varpi_{r_n/2},
 \qquad \sigma_n=\sin(\pi\Psi_n/2),\qquad
 \tau_n=\cos(\pi\Psi_n/2),\qquad \sigma_0=1.
\]
The definitions for $n=t_*+1$ use the empty region.
For a cube of level $n$ set
\begin{equation}\label{I-16-67}
 h_\square=g_\square\sigma_n\tau_{n+1},\qquad
 S_\square=\operatorname{supp}h_\square,\qquad
 S_\square^{(a)}=\{x:\operatorname{dist}_\infty(x,S_\square)
                                                       \le ar_n\}.
\end{equation}
Let $\mathfrak D_n$ contain the cubes for which
$h_\square(b_-)\ne0$ for some $b\in\mathfrak b_0$, and
put $\mathfrak D=\bigcup_n\mathfrak D_n$. The two cutoffs are
\[
 \zeta_\square=\mathbf1_{S_\square^{(1/4)}}*\varpi_{r_n/12},
 \qquad
 \theta_\square=\mathbf1_{S_\square^{(7/12)}}*\varpi_{r_n/12}.
\]
A cube in $\mathfrak D_0$ is a boundary cube if
$\widetilde\square^{2}\cap\Omega_1=\emptyset$.
All other cubes in $\mathfrak D$ are called interior cubes.

\begin{lemma}\label{I-17-14}
The following properties hold uniformly in the number of layers.
\begin{enumerate}[label=(\alph*)]
\item The functions take values in $[0,1]$, and
\begin{equation}\label{I-16-68}
 \sum_{n,z}h_\square(x)^2=1\quad(x\in T),\qquad
 \sum_{\square\in\mathfrak D}h_\square(b_-)^2=1
                                      \quad(b\in\mathfrak b_0).
\end{equation}
For $\phi=h_\square,\zeta_\square,\theta_\square$, its
first and second derivatives are bounded by $c_\flat r_n^{-1}$
and $c_\flat r_n^{-2}$. The same bounds hold for first and
second lattice differences; in particular
$|\Delta\phi|\le3c_\flat r_n^{-2}$ and
$|\phi(x)-\phi(x')|\le c_\flat r_n^{-1}|x-x'|_1$.
The constant $c_\flat$ is absolute. Moreover,
\[
 \zeta_\square=1\text{ on }S_\square^{(1/6)},\quad
 \operatorname{supp}\zeta_\square\subset S_\square^{(1/3)},
 \qquad
 \theta_\square=1\text{ on }S_\square^{(1/2)},\quad
 \operatorname{supp}\theta_\square\subset S_\square^{(2/3)}.
\]

\item For a cube of level $n$ and center $z$,
\[
 S_\square^{(5/6)}\subset\{|x-z|_\infty<5W_n/6\},\qquad
 d_{n+1}\le W_n/2\text{ on }S_\square^{(5/6)}.
\]
For $n\ge1$, also $d_n\ge r_n/6$ there.
If $y\in\Lambda_m$ and
$S_\square\cap\{x:|x-y|_\infty\le6\xi_y\}\ne\emptyset$,
then $n\in\{m-1,m\}$ and $y\in S_\square^{(1/6)}$.
For a fixed $y$ there are at most $16$ such cubes.

\item An interior cube of level $n\ge1$ satisfies
$\square\subset\Omega_{n-1}$. For every interior cube,
\[
 \widetilde\square^{3}\cap\Omega_{n+2}=\emptyset,
 \qquad S_\square^{(5/6)}\subset\mathcal L_n\cup\mathcal L_{n+1},
 \qquad
 \operatorname{dist}_\infty(S_\square^{(5/6)},\Omega_0^c)
                                                    >\mathfrak M\xi.
\]
If $\widetilde\square^{2}$ meets $\Omega_{n+1}$, then
$\widetilde\square^{3}\subset\Omega_n$ and the torus side
is at least $26W_n$. For a boundary cube,
$\operatorname{dist}_\infty(\square,\Omega_1)\ge2\mathfrak M\xi$.

\item Let $r>0$ and $j\ge0$. Suppose the open sup ball of radius $r$ about $x$ is
disjoint from $\Omega_{j+1}$. Then
\begin{equation}\label{eq:local-cube-escape}
 d_{\mathbf B}(x,x')\ge r/\xi_j
                     \quad\text{if }|x-x'|_\infty\ge r.
\end{equation}
Consequently, if $\square$ is interior,
$\Delta(y'')\cap S_\square\ne\emptyset$, and
$\Delta(y)\not\subset S_\square^{(1/6)}$, then
\begin{equation}\label{eq:local-cube-buffer}
 d_{\mathbf B}(y,y'')\ge\frac{\mathfrak M}{72L^2}.
\end{equation}
\end{enumerate}
\end{lemma}

\begin{proof}
At each level the period ratio $\mathsf S/W_n$ is an even
integer at least two. Periodizing the identity for $h$ and
taking its product in the three coordinates gives
$\sum_{z\in\mathcal Z_n}g_\square^2=1$, including the
period ratio two. The distance $d_n$ is $1$-Lipschitz when
finite. Since the convolution radius is at most $r_n/3$,
$\sigma_n=0$ on $\{d_n\le r_n\}$ and $\sigma_n=1$ on
$\{d_n\ge2r_n\}$, with the complementary assertions for
$\tau_n$. Separation implies $\sigma_n=1$ on the support
of $\sigma_{n+1}$. Therefore
\[
 \sum_z h_\square^2
       =\sigma_n^2(1-\sigma_{n+1}^2)
       =\sigma_n^2-\sigma_{n+1}^2.
\]
Summing over $n$ telescopes to one. Discarding cubes outside
$\mathfrak D$ does not change this sum at $b_-$.

Differentiation of a convolution of an indicator with
$\varpi_r$ gives bounds $C_k r^{-k}$ for $k=1,2$, by the
$L^1$ norms of the corresponding derivatives of $\varpi$.
The chain and product rules give the asserted bounds for
$\sigma_n,\tau_n,g_\square,h_\square$. Only fixed cutoff
norms occur; $W_n\ge r_n$ and $r_{n+1}\ge r_n$.
Each lattice difference is an average of the corresponding
continuous derivative along one or two coordinate intervals.
This also proves the bound for $\Delta\phi$. Integrating
along a shortest coordinate path proves the Lipschitz bound.
The support of $\varpi_{r_n/12}$ is contained in the sup
ball of radius $r_n/12$; adding and subtracting this radius
from $r_n/4$ and $7r_n/12$ proves the cutoff assertions.

For the support bounds, use
\[
 S_\square\subset\{|x-z|_\infty\le5W_n/8\}
            \cap\{d_{n+1}\le2r_{n+1}\},
 \qquad S_\square\subset\{d_n\ge r_n\}\quad(n\ge1).
\]
Enlargement by $5r_n/6$ gives (b), because
$2r_{n+1}=W_n/3$ and $r_n=W_n/(6L)$.
In particular it does not reach $\Omega_{n+2}$.

To prove the overlap bound, choose
$x\in S_\square$ with $|x-y|_\infty\le6\xi_m$.
Since $y\notin\Omega_{m+1}$,
$d_{m+1}(x)\le6\xi_m<r_{m+1}$, excluding $n\ge m+1$;
for $n>m+1$ separation gives the same contradiction.
If $n\le m-2$, separation at $y\in\Omega_m$ gives
$d_{m-1}(x)>(\mathfrak M-6)\xi_m>2r_{m-1}$.
This contradicts the upper bound for $d_{n+1}$ when
$n=m-2$; when $n<m-2$, it places $x$ in $\Omega_{n+2}$,
giving the same contradiction by separation. Thus
$n\in\{m-1,m\}$. Now
\[
 |z-y|_\infty\le5W_n/8+6\xi_m<W_n,
 \qquad 6\xi_m\le r_n/6.
\]
In each coordinate an open interval of length $2W_n$ contains
at most two possible corner coordinates. There are at most
eight centers at each of the two levels, and the second
inequality gives $y\in S_\square^{(1/6)}$.

Here is the geometric point needed at an interior cube. If
$n\ge1$, or if $\widetilde\square^{2}$ meets
$\Omega_{n+1}$, some $W_n$-cube with corner $z$ belongs to
$\mathcal L_n$. Indeed, take $x\in S_\square$ and a
$W_n$-cube through $x$. It lies in $\Omega_n$: for $n\ge1$
this follows from $d_n(x)>0$; in the other case every point
of $\widetilde\square^{3}$ is within $7W_n<R_1W_n$
of a point of $\Omega_{n+1}$, so belongs to $\Omega_n$.
If the chosen cube lies in $\Omega_{n+1}$, use
$d_{n+1}(x)\le W_n/3$ to choose a point of
$\Omega_{n+1}^c$ within $W_n/3$ of $x$. It is within
$23W_n/24$ of $z$. A $W_n$-cube through it has corner
$z$, lies in $\Omega_n$ by separation, and does not lie
in $\Omega_{n+1}$. It is the required cube.
The assertion about its corner follows directly from the
half-integer grid: a center within $3W_n/2$ of $z$ differs
from $z$ by $\pm W_n/2$ in each coordinate.

Thus $z\in\Omega_n$. For $n\ge1$, separation from
$\Omega_{n-1}^c$ gives $\square\subset\Omega_{n-1}$.
If $\widetilde\square^{2}$ meets $\Omega_{n+1}$,
the argument just given proves
$\widetilde\square^{3}\subset\Omega_n$; also
$\Omega_{n+1}\ne\emptyset$ forces
$\mathsf S\ge2\mathfrak M\xi_{n+1}\ge26W_n$.
Since $d_{n+1}(z)<W_n$, throughout
$\widetilde\square^{3}$ one has
$d_{n+1}<5W_n<R_1\mathfrak M\xi_{n+1}$; this excludes
$\Omega_{n+2}$. The support therefore meets only layers
$n,n+1$. For $n\ge1$ its distance from $\Omega_0^c$
is larger than $R_1\mathfrak M\xi$. At $n=0$ an interior
cube has $\widetilde\square^{3}\subset\Omega_0$.
Every point of $\Omega_0^c$ is at distance at least $4W_0$
from $z$: otherwise a $W_0$-cube of the complement through
that point has center within $9W_0/2$ of $z$; the
half-integer grid places its center within $7W_0/2$, so
the cube lies in $\widetilde\square^{3}$, a contradiction.
Subtracting $5W_0/6$ proves the boundary margin in (c).
For a boundary cube, the defining exclusion of
$\Omega_1$ from $\widetilde\square^{2}$ immediately
gives its distance at least $2\mathfrak M\xi$ from
$\square$.

For (d), fix $0<r'<r$ and use the clipped distance
$f(z)=\min\{|z-x|_\infty,r'\}$.
On the portion of a connecting set where $0<f<r'$, all
points lie outside $\Omega_{j+1}$ and the weighted length
density is at least $\xi_j^{-1}$. The interval-image
calibration in \cref{subsec:multiscale-sum-proof}, applied
to this $1$-Lipschitz function, gives weighted length at
least $r'/\xi_j$. Taking the infimum and then $r'\uparrow r$
proves \eqref{eq:local-cube-escape}.
For \eqref{eq:local-cube-buffer}, choose
$x''\in\Delta(y'')\cap S_\square$ and
$x\in\Delta(y)\setminus S_\square^{(1/6)}$.
The ball of radius $r_n/6$ about $x''$ is contained in
$S_\square^{(1/6)}$ and misses $\Omega_{n+2}$. Thus
\[
 d_{\mathbf B}(x,x'')\ge\frac{r_n}{6\xi_{n+1}}
       =\frac{\mathfrak M}{36L^2}.
\]
For any point $z\in\Delta(y)$, the segment from $y$ to $z$
stays in the cell interior except possibly at its endpoint.
Its weighted length is $\xi_y^{-1}|z-y|_1\le3/2$;
the endpoint has zero length. These cell-to-center bounds give
$d_{\mathbf B}(y,y'')\ge\mathfrak M/(36L^2)-3
\ge\mathfrak M/(72L^2)$, since $\mathfrak M\ge L^5$.
\end{proof}

We also need to compare the cell of a nearby bare site with
the original cell. If $y\in\Lambda_*$, $x\in\Omega_0$ is
a bare site, $|x-y|_\infty\le2\xi_y$, and $y_*=y(x)$,
then
\begin{equation}\label{eq:nearby-cell-comparison}
 L^{-1}\xi_y\le\xi_{y_*}\le L\xi_y,\qquad
 d_{\mathbf B}(y,y_*)\le9L.
\end{equation}
Indeed $x$ lies in the chart $C(y)$ of
\cref{lem:global-cell-geometry}(b). Its scale comparison
gives the first assertion, and its coordinate-path bound gives
$d_{\mathbf B}(x,y)\le L\xi_y^{-1}|x-y|_1\le6L$.
The cell-to-center bound is $d_{\mathbf B}(x,y_*)\le3/2$;
the triangle inequality proves the second assertion.

For an interior cube of level $n$, fix a lift $\widetilde z$
of its center, put $W=W_n$, and define
\[
 T_\square=\mathbb R^3/(8W\mathbb Z)^3,
 \qquad O=\widetilde z+(-W,W)^3,
 \qquad \iota_\square=\pi_\square\circ(\pi|_O)^{-1}.
\]
Let $\Upsilon_\square$ be the union in $T_\square$ of the
images of the $W$-cubes contained in
$\widetilde z+[-3W,3W]^3$ whose images in $T$ belong to
$\Omega_{n+1}$. Thus this intersection is taken cube-wise.
Both covering maps are injective on $O$; the second is
injective on the larger closed cube of side $6W$.
The maps identify the lattices and centered partitions of
spacings $\xi_j$, $j\le n+1$, within $O$, since both
periods are multiples of these spacings. On this set,
$x\in\Omega_{n+1}$ if and only if
$\iota_\square x\in\Upsilon_\square$.

Use on $T_\square$ the sequence with members $0,\ldots,n$
equal to $T_\square$, member $n+1$ equal to
$\Upsilon_\square$, and all later members empty. Its bare
bond frame is the whole torus. Attach the corresponding spaces
$\mathcal S_\square,\mathcal Y_\square,\mathcal A_\square$
and operators
$Q'_\square,Q_\square,\mathsf R_\square,
\mathsf P_\square,\mathsf G'_\square$ to this sequence and put
\[
 \mathcal P_\square=\partial\mathsf P_\square\partial^*,
 \qquad
 \Delta_{a,\square}
   =\partial^*\partial+\partial\mathsf R_\square\partial^*
                          +Q_\square^*aQ_\square
   =\Delta-\mathcal P_\square+Q_\square^*aQ_\square.
\]
Here and below $a=1$.
For a boundary cube, put $T_\square=T$ and
$\iota_\square=I$. Let $\mathcal A_\square$ consist of
the fields on those bonds of $\mathfrak b_0$ with both
endpoints in $\widetilde\square^{3}$, and let
$\chi_\square$ be its orthogonal projection. Set
\[
 \Delta_{a,\square}
       =\chi_\square\partial^*\partial\chi_\square+\xi^{-2}I
                 \quad\text{on }\mathcal A_\square.
\]
Let $\mathcal S_\square$ consist of the site functions in
$\mathcal S$ supported in $\widetilde\square^{3}$.
All boundary fields are extended by zero, and their higher
tensor source spaces contain all bare bonds.

For either kind write
$|o-o'|_\square=\xi_n^{-1}|o-o'|_{1,T_\square}$.
A cutoff $\phi$ supported in $S_\square^{(5/6)}$ defines
two mutually adjoint transfers, by
\[
 (\phi v)(b)=\phi(b_-)v(\iota_\square b),\qquad
 (\phi A)(\iota_\square b)=\phi(b_-)A(b),
\]
with value zero elsewhere. The same convention applies to
scalar fields and tensors. The common weight $\xi^3$ and
the one-to-one correspondence of sites and bonds prove
adjointness. A transferred smooth cutoff is extended by
zero off $\pi_\square(O)$; it remains smooth with the same
derivative bounds, since its support stays within the smaller
cube of radius $19W/24$. In particular transfers introduce
no boundary distribution or change in the cutoff estimates.

Under these identifications, for fields supported in
$S_\square^{(2/3)}$ at an interior cube,
\begin{equation}\label{I-17-15}
 \begin{gathered}
 \Delta_{\mathcal S}=\Delta,\quad
 \partial_0=\partial,\quad\partial_0^*=\partial^*,\quad
 \chi_{\mathcal A}\partial^*\partial\chi_{\mathcal A}
                                               =\partial^*\partial,\\
 Q'^*\mathsf a'Q'=Q'^*_\square\mathsf a'Q'_\square,
 \qquad Q^*aQ=Q_\square^*aQ_\square.
 \end{gathered}
\end{equation}
Their outputs are supported in $S_\square^{(5/6)}$.
To check the averages, every cell meeting this larger set
has level $n$ or $n+1$, lies within $O$, and corresponds
to the cell of the same level in the local sequence.
Indeed the set is inside $\Omega_n$ and outside
$\Omega_{n+2}$, with the margins in \cref{I-17-14}; its
distance to the boundary of $\square$ exceeds the side
$\xi_{n+1}$ of such a cell. The converse follows from
the same margins and the equality of the membership tests
for $\Omega_{n+1}$ and $\Upsilon_\square$.

A determining bond whose endpoint blocks meet
$S_\square^{(2/3)}$ also has level $n$ or $n+1$.
For a smaller level its blocks lie in $\Omega_n^c$.
For level $n+2$ their distance to $\Omega_{n+1}^c$ is
greater than $(R_1\mathfrak M-2L)\xi_{n+1}>W/2$;
for a larger level their blocks lie in $\Omega_{n+2}$.
Both alternatives contradict the support margins.
The remaining blocks lie in $S_\square^{(5/6)}$, because
$2\xi_{n+1}\le r_n/6$, and their determining bonds satisfy
the same endpoint tests in the two models. The straight-path
formula \eqref{I-16-63} is unchanged by $\iota_\square$,
so its matrix weights, and hence the quadratic averaging
forms, agree. The scalar average agrees cell by cell.
Finally the difference operators involve only neighboring
sites. Their support enlargement is at most $\xi$, that
of the scalar averaging term at most $\xi_{n+1}$, and
that of the bond averaging term at most $2\xi_{n+1}$.
This proves \eqref{I-17-15} and its support assertion.

We shall also use, for bare sites in the interior of a cube,
\begin{equation}\label{eq:local-cube-distance}
 d_{\mathbf B}(x,x')\le L|x-x'|_\square,
 \qquad |x-x'|_\square\le6\mathfrak M.
\end{equation}
For an interior cube the coordinate path in $O$ stays in
$\Omega_{n-1}$ when $n\ge1$, so every step has weighted
length at most $L\xi/\xi_n$ by
\cref{lem:global-cell-geometry}(a). At level zero use the
bound one per step. On $T_\square$ this lifted path has
length $|x-x'|_1$, since its coordinate displacements are
less than $2W$, half the local period. For a boundary cube
use a shortest coordinate path in $T$ and again the bound
one per step. This proves both estimates.

The scalar and Gram estimates apply on each interior local torus
with constants independent of the cube. Set
\(\delta_{\mathrm{pr}}=\delta_{\mathrm u}/L\) and
\(\mathfrak m_{\mathrm{pr}}=\mathfrak m_{\mathrm u}\), using
\eqref{I-17-23}, and assume
\(\mathfrak M\ge\max\{L^5,L\mathfrak m_{\mathrm{pr}}\}\).
Write \(\mathsf q(o,r)=\{x:|x-o|_\infty\le r\}\) on the
torus in question. For local cells \(\Delta(w),\Delta(w')\),
put \(E_{w,w'}=e^{-\delta_{\mathrm{pr}}|w-w'|_\square}\).
There is \(B_{\mathrm{pr}}\), depending only on \(L\), such that
\begin{equation}\label{eq:local-scalar-Gram}
\begin{aligned}
 |(\mathsf G'_\square\lambda)(x)|
   &\le B_{\mathrm{pr}}\xi_n^2E_{w,w'}|\lambda|,&
 |(\nabla\mathsf G'_\square\lambda)(x)|
   &\le B_{\mathrm{pr}}\xi_nE_{w,w'}|\lambda|,\\
 |(\mathsf G'_\square\nabla^*F)(x)|
   &\le B_{\mathrm{pr}}\xi_nE_{w,w'}|F|,&
 |\mathsf C_\square(w,w')|
   &\le B_{\mathrm{pr}}\xi_n^{-4}E_{w,w'},\\
 |\mathcal P_\square(b,b')|
   &\le B_{\mathrm{pr}}(\xi/\xi_{w'})^3
                         \xi_n^{-2}E_{w,w'} .
\end{aligned}
\end{equation}
Here \(x\in\Delta(w)\), the site function \(\lambda\) and
site--direction function \(F\) are supported in \(\Delta(w')\),
\(b,b'\) have localization points in the two respective cells, and
\[
 \mathsf C_\square
   =(Q'_\square\mathsf G'^2_\square Q'^*_\square)^{-1},
 \qquad
 \mathcal P_\square=\partial\mathsf P_\square\partial^*.
\]
The last two kernels use counting measure.

Indeed, the local sequence is separated with block number
\(\mathfrak M'=\mathfrak M/L\) and last index \(n+1\).
Its torus side is \(8\mathfrak M'\xi_{n+1}\), and
\(\Upsilon_\square\) is a union of
\(\mathfrak M'\xi_{n+1}\)-cubes. All preceding regions are
the whole torus. If \(\Upsilon_\square\) is empty, discard
that final member; the period condition still holds.
Thus its separation conditions hold for every
gap number; choose an auxiliary \(R'_1\) with
\(R'_1\mathfrak M'\ge\mathfrak r_{\mathrm u}\).
Since \(\mathfrak M'\ge\mathfrak m_{\mathrm u}\),
\cref{I-13-20,I-14-48} and \eqref{I-17-23} apply.
This auxiliary choice imposes no restriction on the original
gap number. The local layers have indices \(n,n+1\), whence
\[
 \xi_n\le\xi_w\le L\xi_n,
 \qquad d_\square(w,w')\ge
 \xi_{n+1}^{-1}|w-w'|_1=L^{-1}|w-w'|_\square.
\]
The distance bound follows by calibrating a connecting set with
the distance from one endpoint, as in
\cref{lem:global-cell-geometry}.
Substitution gives \eqref{eq:local-scalar-Gram}, for example with
\(B_{\mathrm{pr}}\ge L^2B_{\mathrm u}\).

The vector estimates below allow both adjacent cell sizes.
This will let us localize at a cell without changing the local
model or distinguishing which of its two layers contains it.

\begin{proposition}\label{I-17-38}
There are \(\delta_{\mathrm{loc}}\in(0,1]\) and
\(B_{\mathrm l}\ge1\), depending only on \(L\), and
\(B_{\mathrm l}(\alpha)\), \(B'_{\mathrm l}(\varepsilon)\),
and \(B'_{\mathrm l}(\alpha,\varepsilon)\), depending only on
\(L\) and the displayed exponents, with the following properties.
Assume \(\mathfrak M\ge\max\{L^5,L\mathfrak m_{\mathrm{pr}}\}\),
and let \(\square\in\mathfrak D_n\) be of either kind.
Then \(\Delta_{a,\square}\) is positive definite and its inverse
\(\mathsf G_\square\) is symmetric. Take
\(\lambda_1,\lambda_2\in\{\xi_n,\xi_{n+1}\}\),
\(o,o'\in T_\square\), and put
\(E=e^{-\delta_{\mathrm{loc}}|o-o'|_\square}\).
Let \(\zeta\) be a site function supported in
\(\mathsf q(o,\lambda_1)\), multiplying through the initial point.
Let \(J,J',J''\) belong respectively to
\(\mathcal A_\square,\mathcal A_\square^{(1)},
\mathcal A_\square^{(2)}\) and be supported in
\(\mathsf q(o',\lambda_2)\). Let \(0\le\alpha<1\),
\(\varepsilon>0\), and \(\alpha+\varepsilon<1\).

For boundary cubes all one-forms are extended by zero,
\(\mathcal A_\square^{(r)}=\mathcal A^{(r)}\), and
\(\nabla^*,\nabla^*\nabla^*\) are the adjoints of
\(\nabla,\nabla\nabla\) on \(\mathcal A_\square\).
Localization may use the base point relative either to
\(\Omega_0\) or to \(\Omega_0\cap\widetilde\square^{3}\).
For a bond \(x\) in \(\mathsf q(o,2\lambda_1)\) we have
\begin{align*}
 |(\mathsf G_\square J)(x)|
    &\le B_{\mathrm l}\xi_n^2E|J|,&
 |(\nabla\mathsf G_\square J)(x)|
    &\le B_{\mathrm l}\xi_nE|J|,\\
 |(\mathsf G_\square\nabla^*J')(x)|
    &\le B_{\mathrm l}\xi_nE|J'|,&
 |(\Delta\mathsf G_\square J)(x)|
    &\le B_{\mathrm l}E|J|,\tag{R1}\\
 |(\nabla\mathsf G_\square\nabla^*J')(x)|
    &\le B'_{\mathrm l}(\varepsilon)E
                    (\|J'\|_{\varepsilon,\lambda_2}+|J'|).
                                                    \tag{R3}
\end{align*}
The localized H\"older bounds are
\begin{align*}
 \xi_n^{-1}\|\zeta\mathsf G_\square J\|_{\alpha,\lambda_1},\quad
 \|\zeta\nabla\mathsf G_\square J\|_{\alpha,\lambda_1}
  &\le B_{\mathrm l}(\alpha)\xi_n
           (\|\zeta\|_{\alpha,\lambda_1}+|\zeta|)E|J|,\\
 \|\zeta\mathsf G_\square\nabla^*J'\|_{\alpha,\lambda_1}
  &\le B_{\mathrm l}(\alpha)\xi_n
           (\|\zeta\|_{\alpha,\lambda_1}+|\zeta|)E|J'|,\tag{R2}\\
 \|\zeta\nabla\mathsf G_\square\nabla^*J'\|_{\alpha,\lambda_1}
  &\le B'_{\mathrm l}(\alpha,\varepsilon)
      (\|\zeta\|_{\alpha,\lambda_1}+|\zeta|)E
      (\|J'\|_{\alpha+\varepsilon,\lambda_2}+|J'|).
                                                        \tag{R4}
\end{align*}
Finally,
\begin{align*}
 \xi_n^{-2}\|\zeta\mathsf G_\square J\|,\quad
 \xi_n^{-1}\|\zeta\nabla\mathsf G_\square J\|,\quad
 \|\zeta\nabla\nabla\mathsf G_\square J\|
      &\le B_{\mathrm l}|\zeta|E\|J\|,\\
 \xi_n^{-1}\|\zeta\mathsf G_\square\nabla^*J'\|,\quad
 \|\zeta\nabla\mathsf G_\square\nabla^*J'\|
      &\le B_{\mathrm l}|\zeta|E\|J'\|,\\
 \|\zeta\mathsf G_\square\nabla^*\nabla^*J''\|
      &\le B_{\mathrm l}|\zeta|E\|J''\|.\tag{R5}
\end{align*}
All constants are uniform in the spacing, torus, cube, gap number
and number of layers.
\end{proposition}

\begin{proof}
For an interior cube we first identify its operator with the
two-scale operator of \cref{I-16-14}. First convert that
corner-block operator to centered blocks in its normalized
coordinates, with spacings $\eta=L^{-n},1,L$.
There are two translations:
one changes the origin of each fine block, and the other changes
the origin of the coarse blocks. Put
\[
 \eta=L^{-n},\quad s=(1-\eta)/2,\quad s'=(L-1)/2,
 \quad e=e_1+e_2+e_3,\quad
 (\tau_vf)(x)=f(x-v),\quad \tau=\tau_{s'e}\tau_{se}.
\]
Since \(L\) is odd, both shifts belong to their respective
lattices. Superscript \(\mathrm c\) denotes centered blocks. Since
the centered fine and coarse blocks are the corresponding
corner blocks shifted by \(-se\) and \(-s'e\),
\[
 \mathsf Q_n^{\mathrm c}=\mathsf Q_n\tau_{se},\qquad
 \mathsf Q_{1,L}^{\mathrm c}=\mathsf Q_{1,L}\tau_{s'e},
 \qquad
 \tau_{s'e}\mathsf Q_n^{\mathrm c}=\mathsf Q_n\tau,
 \qquad
 \mathsf Q_{1,L}^{\mathrm c}\mathsf Q_n^{\mathrm c}
                    =\mathsf Q_{1,L}\mathsf Q_n\tau.
\]
The same identities hold for the scalar averages. The last two
use the commutation of the fine average with unit-lattice
translations. The coarse set \(\Lambda'\) remains arbitrary;
its union of unit blocks becomes \(\Lambda_{\mathrm c}=\Lambda-s'e\).
Thus both its membership tests and the two zero-average
conditions are translated by \(\tau\). Equivalently,
\(N^{\mathrm c}=\tau^{-1}N\). The Laplacian commutes with
\(\tau\), so the same relation holds for their Laplacian images;
orthogonality gives \(\mathsf R^{\mathrm c}=\tau^{-1}\mathsf R\tau\).
The coarse penalty satisfies
\(\mathsf q^{\mathrm c}(\mathsf Q_n^{\mathrm c}A)
=\mathsf q(\mathsf Q_n\tau A)\).
Translations commute with the differences and preserve the
inner products. Hence
\[
 \Delta_1^{\mathrm c}=\tau^{-1}\Delta_1\tau,
 \qquad G^{\mathrm c}=\tau^{-1}G\tau.
\]
The estimates of \cref{I-16-14} pass to \(G^{\mathrm c}\):
the centers translate along with the supports, while their
distance, the norms and the H\"older seminorms are unchanged.

Next rescale by \((\mathscr S_nf)(X)=f(\xi_nX)\).
The resulting torus has side \(8\mathfrak M\), an even multiple
of \(L\), and bare spacing \(\eta\). Its coarse subset
\(\Lambda'\) consists of the level-\(n+1\) sites in
\(\Upsilon_\square\). The decomposition into centered blocks
is exact: in each coordinate every integer
\(|k|\le(L^{n+1}-1)/2\) has a unique expression
\[
 k=L^ni+i',\qquad |i|\le(L-1)/2,
 \qquad |i'|\le(L^n-1)/2.
\]
It gives \(Q'_\square=\mathsf Q_n^{\prime\mathrm c}\)
outside the coarse region and
\(Q'_\square=\mathsf Q_{1,L}^{\prime\mathrm c}
\mathsf Q_n^{\prime\mathrm c}\) on its coarse sites,
after rescaling. The analogous bond identity follows by
writing each straight-path index uniquely as \(L^ni+i'\)
with \(0\le i<L\) and \(0\le i'<L^n\).
The determining bonds of level \(n\) have neither endpoint
in the coarse region; those of level \(n+1\) have an endpoint
in its coarse-site subset. Thus their tests are precisely
the tests defining \(\mathsf q\).

Every scalar product acquires the same factor \(\xi_n^3\),
and a difference acquires \(\xi_n^{-1}\). The quadratic
averaging penalty becomes
\[
 \xi_n\sum_{b:\ b_-,b_+\notin\Lambda_{\mathrm c}}
                 |(\mathsf Q_n^{\mathrm c}\mathscr S_nA)(b)|^2
 +\xi_{n+1}\sum_{c\in\Lambda'^*}
       |(\mathsf Q_{1,L}^{\mathrm c}
                    \mathsf Q_n^{\mathrm c}\mathscr S_nA)(c)|^2,
\]
where the first sum is over unit bonds.
The zero-average spaces and their Laplacian images correspond,
so the common scalar-product factor also preserves orthogonal
projection. Consequently
\begin{equation}\label{I-16-73}
 \mathsf R_\square=\mathscr S_n^{-1}\mathsf R^{\mathrm c}\mathscr S_n,
 \qquad
 \Delta_{a,\square}=\xi_n^{-2}
          \mathscr S_n^{-1}\Delta_1^{\mathrm c}\mathscr S_n,
 \qquad
 \mathsf G_\square=\xi_n^2
          \mathscr S_n^{-1}G^{\mathrm c}\mathscr S_n.
\end{equation}
In particular the local operator is positive definite.
An expression with \(r\) differences, including adjoint
differences, acquires the factor \(\xi_n^{2-r}\).
Moreover
\[
 |\mathscr S_nf|=|f|,\qquad
 \|\mathscr S_nf\|_\alpha=\|f\|_{\alpha,\xi_n},\qquad
 \|f\|=\xi_n^{3/2}\|\mathscr S_nf\|.
\]
The common factor in the last identity cancels in operator
bounds. Thus \cref{I-16-14} gives (R1)--(R5) at radius
\(\xi_n\), except for the first entry of (R2), which we
derive below. This application uses no restriction on
\(\Lambda'\) beyond the one already admitted in that theorem.

To allow the larger radius \(L\xi_n\), use a partition at
scale \(\xi_n\). For the corners
\(w\in\xi_n(\mathbb Z+\tfrac12)^3\) on \(T_\square\), set
\[
 \varphi_w(x)=\prod_{\mu=1}^3
                h((x_\mu-w_\mu)/\xi_n)^2.
\]
The identity for \(h\) gives \(\sum_w\varphi_w=1\).
Each member lies in \([0,1]\), is supported in
\(\mathsf q(w,5\xi_n/8)\), and has
\(|\nabla\varphi_w|\le C_\varphi/\xi_n\).
By \eqref{I-17-34},
\(\|\varphi_w\|_{\alpha,\xi_n}\le C_\varphi\).
Let \(W(o,\lambda)\) consist of those members meeting
\(\mathsf q(o,2\lambda)\). Since the grid spacing is \(\xi_n\)
and \(\lambda\le L\xi_n\), both the number of these members
and their centers' normalized distances from \(o\) are bounded
by constants depending only on \(L\).
Decompose \(J=\sum_{w'}\varphi_{w'}J\) and
\(\zeta=\sum_w\varphi_w\zeta\), and likewise for the tensors.
The product and change-of-scale inequalities
\eqref{I-17-34}--\eqref{I-17-35} give
\[
 \begin{gathered}
 \|\varphi_w\zeta\|_{\alpha,\xi_n}
       \le\|\zeta\|_{\alpha,\lambda_1}+C_\varphi|\zeta|,
 \qquad
 \|\varphi_{w'}J'\|_{\varepsilon,\xi_n}
       \le\|J'\|_{\varepsilon,\lambda_2}+C_\varphi|J'|,\\
 \|F\|_{\alpha,\lambda_1}
       \le L\max\{\|F\|_{\alpha,\xi_n},2|F|\}.
 \end{gathered}
\]
For any rate \(\delta\le1\), the triangle inequality bounds
the exponential between \(w,w'\) by \(C(L)\) times the
exponential between \(o,o'\). Applying the radius-\(\xi_n\)
estimates to all pairs and summing proves (R1),
(R3), (R5), (R4), and the last two entries of (R2).
The change of seminorm scale also uses their sup-norm bounds;
for (R4) apply (R3) with \(\alpha+\varepsilon\) in place
of \(\varepsilon\). Every bond in
\(\mathsf q(o,2\lambda_1)\) belongs to a radius-\(\xi_n\)
cube of this partition, which proves the larger output domain.

For the remaining entry of (R2), apply \eqref{I-17-34}
to \(X=\mathsf q(o,\lambda_1)\). Its enlarged cube is
\(X^{+\lambda_1}=\mathsf q(o,2\lambda_1)\), and the required
coordinate-path estimate applies since
\(5\lambda_1\le5L\xi_n<8\mathfrak M\xi_n\). Thus
\[
 \|\zeta\mathsf G_\square J\|_{\alpha,\lambda_1}
 \le\|\zeta\|_{\alpha,\lambda_1}
                 \sup_{X^{+\lambda_1}}|\mathsf G_\square J|
    +|\zeta|\lambda_1
                 \sup_{X^{+\lambda_1}}|\nabla\mathsf G_\square J|.
\]
Now (R1) and \(\lambda_1\le L\xi_n\) give the assertion.

For a boundary cube the mass gives a direct proof.
Here \(n=0\). On \(\mathcal A_\square\) put
\(\mathsf t=\xi^2\chi_\square\partial^*\partial\chi_\square\).
The ambient difference identity \eqref{I-11-18} and
\(\|\tau_\mu-I\|\le2\) imply
\(0\le\partial^*\partial\le\Delta\le12\xi^{-2}\).
Consequently
\[
 \Delta_{a,\square}
    =\xi^{-2}\bigl(13I-(12I-\mathsf t)\bigr)\ge\xi^{-2}I,
 \qquad
 \mathsf G_\square
    =\xi^2\sum_{i\ge0}13^{-i-1}(12I-\mathsf t)^i.
\]
For counting kernels \(|O(b,b')|\le\|O\|\), since both
coordinate vectors have squared norm \(\xi^3\).
The entry \(\mathsf t(b,b')\) vanishes unless the two bonds
belong to a common plaquette, and hence unless
\(|b_--b'_-|_1\le2\xi\). The \(i\)-th power therefore
vanishes at distances greater than \(2i\xi\), while its
norm is at most \(12^i\). Summing the remaining geometric
series gives
\begin{equation}\label{I-17-39}
 |\mathsf G_\square(b,b')|
 \le\xi^2(12/13)^{|b_--b'_-|_1/(2\xi)}
       =\xi^2e^{-2\vartheta|b_--b'_-|_1/\xi},
 \qquad \vartheta=\tfrac14\log(13/12)>0.
\end{equation}
The kernels are extended by zero outside \(\mathcal A_\square\).
Every operator in (R1)--(R5) is obtained by taking at most two
differences of this kernel. These differences are sums of at
most \(36\) shifts by at most \(2\xi\), with the appropriate
factor \(\xi^{-1}\) or \(\xi^{-2}\). Their kernels are bounded
by \(C\xi^{2-r}e^{-2\vartheta|b_--b'_-|_1/\xi}\), where
\(r\) is the total number of differences.

Either allowed base point differs from the initial point by
at most \(\xi\). There are at most \(6(4L+1)^3\) bonds
whose base points lie in \(\mathsf q(o,2L\xi)\). For output
and input bonds in the required cubes,
\[
 |b_--b'_-|_1/\xi
     \ge|o-o'|_\square-3(3L+2).
\]
Summing the kernel over the input bonds gives (R1) and (R3).
The row and column bounds and \cref{I-17-29}(e) give all six
estimates in (R5). The powers of \(\xi\) are unchanged by
these finite sums. Finally distinct sites are at least \(\xi\)
apart and \(\lambda_1\le L\xi\), so
\(\|F\|_{\alpha,\lambda_1}\le2L|F|\).
If \((\zeta F)(b)\ne0\), its initial point lies in
\(\mathsf q(o,\lambda_1)\) and its base point in
\(\mathsf q(o,2\lambda_1)\). Thus the sup-norm estimates
just proved give (R2) and (R4) as well.

Choose \(\delta_{\mathrm{loc}}\le
\min\{1,\delta_{\mathrm{per}},\vartheta\}\), and enlarge
the finitely many constants to cover both kinds of cubes.
The interior rate is that of \cref{I-16-14}; the boundary
rate is absolute. This gives precisely the stated uniformity.
\end{proof}

\subsection{Combining cell estimates}
\label{I-17-27}

For $p\in\{2,\infty\}$ define
\begin{equation}\label{I-17-28}
 \begin{split}
 N_p(O;y,y')&=\|\chi_yO\chi_{y'}\|_{p\to p},\\
 \|O\|_{p,\kappa}
   &=\sup_y\sum_{y'\in\Lambda_*}
            e^{\kappa d_{\mathbf B}(y,y')}N_p(O;y,y').
 \end{split}
\end{equation}
The domain and range may be $\mathcal S$, $\mathcal A$,
$\mathcal Y$, site functions with a direction index, or the
subspace of $\mathcal A^{(r)}$ supported in $\Omega_0$.
An operator with values in an ambient tensor space is restricted
to that subspace when its cell norm is used. On each of these
spaces the $\chi_y$ are mutually orthogonal projections with
sum $I$. Write $O\in\mathfrak K_q(\kappa,A)$ when
$N_\infty(O;y,y')\le A\xi_y^q e^{-\kappa d_{\mathbf B}(y,y')}$.
Products below are asserted only when their domains match.

\begin{lemma}\label{I-17-29}
Write $d=d_{\mathbf B}$. The following rules hold.
\begin{enumerate}[label=(\alph*)]
\item For $\kappa\ge0$,
\[
 \begin{gathered}
 \|\chi_yOf\|_p\le\sum_{y'}N_p(O;y,y')\|\chi_{y'}f\|_p,\\
 N_p(O_1O_2;y,y')\le\sum_zN_p(O_1;y,z)N_p(O_2;z,y'),\\
 \|O_1O_2\|_{p,\kappa}\le
       \|O_1\|_{p,\kappa}\|O_2\|_{p,\kappa},\qquad
 N_p(O;y,y')\le\|O\|_{p,\kappa}e^{-\kappa d(y,y')}.
 \end{gathered}
\]
Also $\|I\|_{p,\kappa}=1$. If $O$ is an endomorphism and
$\|O\|_{p,\kappa}\le1/2$, then
$(I-O)^{-1}=\sum_{i\ge0}O^i$, with
$\|O^i\|_{p,\kappa}\le2^{-i}$ and
$\|(I-O)^{-1}\|_{p,\kappa}\le2$.

\item If $0\le\kappa<\kappa'$,
$R_1\mathfrak M\ge\mathsf r_S(\kappa'-\kappa)$ and
$N_p(O;y,y')\le Ae^{-\kappa'd(y,y')}$, then
$\|O\|_{p,\kappa}\le Ac_S(\kappa'-\kappa)$.
For $\kappa>0$ and $R_1\mathfrak M\ge\mathsf r_S(\kappa/2)$,
\begin{equation}\label{I-17-30}
 \sum_z e^{-\kappa d(y,z)-\kappa d(z,y')}
       \le c_S(\kappa/2)e^{-\kappa d(y,y')/2}.
\end{equation}

\item If $q\in\mathbb R$, $\epsilon>0$ and
$\epsilon R_1\mathfrak M\ge|q|\log L$, then
$\xi_{y'}^q\le L^{|q|}\xi_y^q e^{\epsilon d(y,y')}$.

\item Suppose $\kappa>0$,
$R_1\mathfrak M\ge\mathsf r_S(\kappa/4)$ and
$\kappa R_1\mathfrak M/4\ge|q_2|\log L$.
If $O_j\in\mathfrak K_{q_j}(\kappa,A_j)$ for $j=1,2$, then
\[
 O_1O_2\in\mathfrak K_{q_1+q_2}
       \bigl(\kappa/2,L^{|q_2|}c_S(\kappa/4)A_1A_2\bigr).
\]

\item For maps between fine site or bond spaces, including
direction indices,
\[
 \begin{split}
 N_2(O;y,y')&\le
   \bigl(N_\infty(O;y,y')N_\infty(O^*;y',y)\bigr)^{1/2},\\
 N_2(O^*;y',y)&=N_2(O;y,y').
 \end{split}
\]
The same assertions hold for restrictions by arbitrary set
indicators, including maps between $\mathcal A_\square$ and
$\mathcal A$.

\item For multipliers $\phi,\phi'$ on the range and domain,
and set indicators $\chi,\chi'$ on fine site or bond spaces,
\[
 \|\chi(O\phi'-\phi O)\chi'\|_{\infty\to\infty}
 \le\|\chi O\chi'\|_{\infty\to\infty}
       \sup_{\chi(x)\chi'(x')\ne0}|\phi'(x')-\phi(x)|.
\]

\item If $N_2(O;y,y')\le k(y,y')$ with $k\ge0$ symmetric,
then $\|O\|_{2\to2}\le\sup_y\sum_{y'}k(y,y')$.
\end{enumerate}
\end{lemma}

\begin{proof}
Insert $I=\sum_z\chi_z$ and use the triangle inequality to
obtain the first two inequalities in (a). The metric triangle
inequality then gives submultiplicativity of the weighted row
sums. Each block is bounded by its row sum. Since the cell
projections partition each nonzero space, the identity has norm
one; the usual geometric series in this finite-dimensional
normed algebra proves the inverse assertion.

The first assertion of (b) is \eqref{I-13-122}, and its second
is \eqref{I-13-123} at rate $\kappa/2$. Part (c) is
\eqref{eq:global-scale-ratio}. To prove (d), apply (a) and
transfer $\xi_z^{q_2}$ to the output scale by (c). The remaining
sum is bounded by
\[
 \sum_z e^{-3\kappa d(y,z)/4-\kappa d(z,y')}
 \le e^{-\kappa d(y,y')/2}\sum_z e^{-\kappa d(y,z)/4}
 \le c_S(\kappa/4)e^{-\kappa d(y,y')/2}.
\]

For (e), the fine inner products all have weight $\xi^3$.
Thus the adjoint matrix is the transpose, and Cauchy--Schwarz
in the double sum bounds its $\ell^2$ norm by the square root
of the product of its absolute row and column sums. These sums
are the two displayed $\ell^\infty$ norms. This argument also
applies after arbitrary restrictions; adjoint equality gives
the second assertion. In (f), each matrix entry $O(x,x')$ is
multiplied by $\phi'(x')-\phi(x)$, so the absolute row sums
give the bound.

Finally, put $\bar k=\sup_y\sum_{y'}k(y,y')$. Orthogonality,
(a), Cauchy--Schwarz and symmetry give
\[
 \|Of\|^2\le\sum_y
     \left(\sum_{y'}k(y,y')\|\chi_{y'}f\|\right)^2
 \le\bar k\sum_{y,y'}k(y,y')\|\chi_{y'}f\|^2
 \le\bar k^2\|f\|^2.
\]
\end{proof}

We also need the product rules with their lattice shifts.
Write $f^{\pm\nu}=T_{\pm\nu}f$. For a set $X$, let
$X^{+\lambda}=\{x:\operatorname{dist}_\infty(x,X)\le\lambda\}$,
and let $\|\cdot\|_{\alpha,\lambda;X}$ restrict both arguments
of the seminorm to $X$; for bond fields these are the initial
points.

\begin{lemma}\label{I-17-31}
\begin{enumerate}[label=(\alph*)]
\item On a bare torus, for site or bond fields,
\begin{equation}\label{I-17-32}
 \begin{split}
 \nabla_\nu(\phi f)
   &=\phi^{+\nu}\nabla_\nu f+(\nabla_\nu\phi)f,\\
 \phi\nabla_\nu^*f
   &=\nabla_\nu^*(\phi f)+(\nabla_\nu\phi)^{-\nu}f^{-\nu}.
 \end{split}
\end{equation}
The commutator $\mathsf L(\phi)=\phi\Delta-\Delta\phi$ is
antisymmetric and satisfies
\begin{equation}\label{I-17-33}
 \begin{split}
 \mathsf L(\phi)f
 &=\sum_\nu\bigl[(\nabla_\nu\phi)\nabla_\nu f
             +(\nabla_\nu^*\phi)\nabla_\nu^*f\bigr]
                      -(\Delta\phi)f\\
 &=-\sum_\nu\nabla_\nu^*
          \bigl[(\nabla_\nu\phi)(f+f^{+\nu})\bigr]
                      +(\Delta\phi)f.
 \end{split}
\end{equation}

\item For $0\le\alpha<1$ and $0<\lambda'\le\lambda$,
\begin{equation}\label{I-17-34}
 \begin{aligned}
 \|\phi f\|_{\alpha,\lambda}
 &\le |f|_{X^{+\lambda}}\|\phi\|_{\alpha,\lambda}
            +|\phi|\|f\|_{\alpha,\lambda;X^{+\lambda}}
           &&(\operatorname{supp}\phi\subset X),\\
 \|f\|_{\alpha,\lambda;X}&\le\lambda\sup_X|\nabla f|.
 \end{aligned}
\end{equation}
The second inequality holds when $X$ is the torus, or a cube
of radius $\lambda''$ with $2\lambda''+\lambda$ smaller than
the torus side. Also
\begin{equation}\label{I-17-35}
 \begin{split}
 \|f\|_{\alpha,\lambda'}&\le\|f\|_{\alpha,\lambda},\\
 \|f\|_{\alpha,\lambda}
   &\le(\lambda/\lambda')^\alpha
                   \max\{\|f\|_{\alpha,\lambda'},2|f|\}.
 \end{split}
\end{equation}

\item For $\mathfrak M\ge L^5$, $h=h_\square$ and
$J\in\mathcal A^{(1)}$, transfer to the local model gives
\begin{equation}\label{I-17-36}
 h\nabla^*J=\nabla^*(hJ)+\mathsf m_h(J),\qquad
 \mathsf m_h(J)=\sum_\nu(\nabla_\nu h)^{-\nu}J_\nu^{-\nu}.
\end{equation}
On the right, $hJ$ and $\mathsf m_h(J)$ are moved by
$\iota_\square$ without an additional multiplier. The
divergence is the adjoint in $\mathcal A_\square$.
For a boundary cube the term $\mathsf m_h(J)$ is also
restricted by $\chi_{\mathcal A}$.
\end{enumerate}
\end{lemma}

\begin{proof}
Expanding $T_\nu-I$ gives \eqref{I-17-32}.
For the commutator, at a site or initial bond point $x$,
\[
 (\phi\Delta f-\Delta(\phi f))(x)
   =\xi^{-2}\sum_{x'\sim x}(\phi(x')-\phi(x))f(x').
\]
Substituting $f(x')=f(x')-f(x)+f(x)$ gives the first form
in \eqref{I-17-33}. Applying the second rule of
\eqref{I-17-32} to the divergence in its second form gives
the same neighbor sum. Antisymmetry follows because both
$\phi$ and $\Delta$ are symmetric.

For the product seminorm, expand
$\phi(x)f(x)-\phi(x')f(x')$ into the difference of $\phi$
times $f(x')$ and $\phi(x)$ times the difference of $f$.
If either point belongs to $X$, both belong to $X^{+\lambda}$;
otherwise the product vanishes. For the gradient bound, a
coordinate path of length $|x-x'|_1$ stays in the cube or
torus and bounds the difference by
$|x-x'|_1\sup_X|\nabla f|$. The stated period condition
ensures that least-distance lifts in the cube have this
coordinate displacement. Multiplying by
$(\lambda/|x-x'|_1)^\alpha$ proves the bound.
Splitting pairs at distance $\lambda'$ proves
\eqref{I-17-35}.

For (c), apply the second identity in \eqref{I-17-32} in
each direction. All terms have initial points within $\xi$
of $S_\square$, and their bonds have endpoints within
$2\xi$. The support margin $2\xi\le5r_n/6$ places these
bonds in the faithful local chart. For an interior cube
they lie inside $\Omega_0$, so $\iota_\square$ commutes
with the required shifts and the divergence is unrestricted.
For a boundary cube the bonds lie in $\widetilde\square^{3}$;
the global and local restrictions agree there. This gives
the compression specified in (c).
\end{proof}

\subsection{Comparison of the global and local gauge terms}
\label{sec:global-local-gauge}

Write $\mathcal P=\partial_0\mathsf P\partial_0^*$ and, on an
interior local torus, $\mathcal P_\square
=\partial\mathsf P_\square\partial^*$. These operators need
not agree on forms supported in the common region. Their
difference can nevertheless be expressed through commutators
supported where an intermediate cutoff changes. The distance
from that set to the input and output gives the following estimate.

\begin{lemma}\label{I-17-43}
There are $c_{\mathrm F}>0$, $C_{\mathrm F}$ and thresholds
depending only on $L$ such that, if $\mathfrak M$ and
$R_1\mathfrak M$ exceed these thresholds, then every interior
cube $\square\in\mathfrak D_n$ and every $J\in\mathcal A$
supported in $S_\square^{(1/3)}$ satisfy
\begin{equation}\label{I-17-45}
 \sup_{b\text{ localized in }S_\square^{(1/3)}}
       |((\mathcal P-\mathcal P_\square)J)(b)|
 \le C_{\mathrm F}\xi_n^{-2}e^{-c_{\mathrm F}\mathfrak M}|J|.
\end{equation}
Here the local expression is transferred by $\iota_\square$.
One may take $c_{\mathrm F}=\delta_{\mathrm{pr}}/(288L)$.
\end{lemma}

\begin{proof}
We use $\mathsf T=\mathsf G'Q'^*$ and
$\mathsf C=(\mathsf T^*\mathsf T)^{-1}$, and attach a
subscript $\square$ to the corresponding local operators.
For a local cell centre $w$, write $B_w$ for the bare sites
of its cell $\Delta(w)$.
Let $\theta=\theta_\square$, and let
$\Theta:\mathcal S_\square\to\mathcal S$ be transfer by
$\theta$. Its adjoint is the opposite transfer. The cells meeting
$\operatorname{supp}\theta\subset S_\square^{(2/3)}$ are common
to the two settings. Thus multiplication by $\theta(y)$ at their
centres defines $\bar\Theta:\mathcal Y_\square\to\mathcal Y$;
it is zero on all other cells. The opposite transfer is
$\bar\Theta^*$, because corresponding cells have the same
weight $\xi_y^3$.

Set
\[
 \begin{aligned}
 K_\theta&=\Theta\Delta'_{a,\square}-\Delta'_a\Theta,
 & Z&=\Theta Q'^*_\square-Q'^*\bar\Theta,
 &Y_\square&=Z-K_\theta\mathsf T_\square,\\
 K_\theta^\sharp&=\Theta^*\Delta'_a-\Delta'_{a,\square}\Theta^*,
 &Z^\sharp&=\Theta^*Q'^*-Q'^*_\square\bar\Theta^*,
 &Y&=Z^\sharp-K_\theta^\sharp\mathsf T.
 \end{aligned}
\]
The maps $Y_\square:\mathcal Y_\square\to\mathcal S$ and
$Y:\mathcal Y\to\mathcal S_\square$ enter the two resolvent
identities
\begin{equation}\label{I-17-46}
 \Theta\mathsf T_\square
   =\mathsf T\bar\Theta+\mathsf G'Y_\square,\qquad
 \Theta^*\mathsf T
   =\mathsf T_\square\bar\Theta^*+\mathsf G'_\square Y.
\end{equation}
For example, multiply
$\Delta'_a\Theta=\Theta\Delta'_{a,\square}-K_\theta$
on the left by $\mathsf G'$ and on the right by
$\mathsf T_\square$, and use
$\Delta'_{a,\square}\mathsf T_\square=Q'^*_\square$.
Writing $\mathsf V=\mathsf T^*\mathsf T=\mathsf C^{-1}$
and $\mathsf V_\square=\mathsf C_\square^{-1}$, we obtain
\[
 \begin{split}
 \mathsf V\bar\Theta-\bar\Theta\mathsf V_\square
    &=Y^*\mathsf G'_\square\mathsf T_\square
                       -\mathsf T^*\mathsf G'Y_\square,\\
 \mathsf C\bar\Theta-\bar\Theta\mathsf C_\square
    &=\mathsf C(\mathsf T^*\mathsf G'Y_\square
                -Y^*\mathsf G'_\square\mathsf T_\square)
                   \mathsf C_\square.
 \end{split}
\]
Substitute these formulas and
$\mathsf P=\mathsf T\mathsf C\mathsf T^*$ into the difference
of projections. The ordered products give
\begin{equation}\label{I-17-47}
 \begin{split}
 \mathsf P\Theta-\Theta\mathsf P_\square
  &=
  \mathsf T(\mathsf C\bar\Theta-\bar\Theta\mathsf C_\square)
          \mathsf T_\square^*
    +\mathsf T\mathsf CY^*\mathsf G'_\square
    -\mathsf G'Y_\square\mathsf C_\square\mathsf T_\square^*\\
  &=\mathsf T\mathsf CY^*\mathsf G'_\square\mathsf R_\square
       -\mathsf R\mathsf G'Y_\square
                    \mathsf C_\square\mathsf T_\square^* .
 \end{split}
\end{equation}

We next locate the commutators. At a bare site $u\in\Omega_0$, put
$j=n(y(u))$ and identify the common cells by $\iota_\square$.
The local identities \eqref{I-17-15} give
\begin{equation}\label{I-17-48}
 \begin{split}
 (K_\theta f)(u)
   &=(\mathsf L(\theta)f)(u)
    +a_j\xi_j^{-2}L^{-3j}
      \sum_{u'\in B_{y(u)}}(\theta(u)-\theta(u'))f(u'),\\
 (Z\omega)(u)&=(\theta(u)-\theta(y(u)))\omega(y(u)).
 \end{split}
\end{equation}
The same formulas hold for the sharp operators, with the
domains reversed. Let $\mathcal T$ be the set of centres
$y(u)$ for which $\theta$ is not constant on
$B_{y(u)}\cup\{u\pm\xi e_\nu\}_\nu$, and put
$T_\theta=\bigcup_{y\in\mathcal T}B_y$. All four operators
vanish outside $T_\theta$. Since $\theta=1$ on
$S_\square^{(1/2)}$ and vanishes outside $S_\square^{(2/3)}$,
the cutoff margins and $\mathfrak M\ge L^5$ imply that
these cells are contained in
$S_\square^{(5/6)}$ and disjoint from $S_\square^{(5/12)}$.
They have levels $n,n+1$. Their number, and the total number
of local cells, are at most $m=(8\mathfrak M)^3$.

The first identity in \eqref{I-17-33}, the cutoff derivative
bounds, and $a_j\le1$ imply
\begin{equation}\label{I-17-49}
 |K_\theta f(u)|\le C\mathfrak M^{-1}
       \bigl(\xi_n^{-1}\sup|\nabla f|
                         +\xi_n^{-2}\sup|f|\bigr),\qquad
 |Z\omega(u)|\le C\mathfrak M^{-1}|\omega(y(u))|.
\end{equation}
The suprema are over the cell of $u$ and its bare neighbours.
Indeed the averaging term in \eqref{I-17-48} is bounded by
$3c_\flat r_n^{-1}\xi_j^{-1}\sup|f|$, and its analogue
for $Z$ by $3c_\flat r_n^{-1}\xi_j|\omega(y(u))|$.
These estimates apply to the sharp operators as well.

Assume $|J|\le1$. For a local cell centre $w$, let $D(w)$ be
its least normalized local distance from a cell meeting the
$\xi$-neighbourhood of $S_\square^{(1/3)}$. Decompose the
input into these cells and apply the scalar and Gram estimates
of \eqref{eq:local-scalar-Gram}. Using
$\mathsf T_\square^*=Q'_\square\mathsf G'_\square$,
the sup-norm contraction of $Q'_\square$, and
$Q'^*_\square\omega=\sum_w\omega(w)\chi_w$, gives
\[
 \begin{aligned}
 |\mathsf T_\square^*\partial^*J(w)|
   &\le Cm\xi_n e^{-\delta_{\mathrm{pr}}D(w)},\\
 \omega=\mathsf C_\square\mathsf T_\square^*\partial^*J:
 \quad|\omega(w)|
   &\le Cm^2\xi_n^{-3}e^{-\delta_{\mathrm{pr}}D(w)},\\
 f=\mathsf T_\square\omega:
 \quad\sup_{B_w}(|f|+\xi_n|\nabla f|)
   &\le Cm^3\xi_n^{-1}e^{-\delta_{\mathrm{pr}}D(w)},\\
 g=\mathsf G'_\square\mathsf R_\square\partial^*J:
 \quad\sup_{B_w}|g|
   &\le Cm^4\xi_n e^{-\delta_{\mathrm{pr}}D(w)}.
 \end{aligned}
\]
The last line follows from
$g=\mathsf G'_\square\partial^*J-\mathsf G'_\square f$.
Each convolution retains the displayed exponential because
$D(w)\le|w-w'|_\square+D(w')$; bounding the number of cells
by $m$ accounts for its additional factor.

A cell with centre in $\mathcal T$, or a neighbouring cell needed in
\eqref{I-17-49}, is separated from the input cells by
at least $r_n/12-4\xi_{n+1}\ge r_n/24$ in sup distance.
The local torus distances equal those of the common lifts
there. Hence $D(w)\ge\mathfrak M/(144L)$ on these cells.
With $E=e^{-\delta_{\mathrm{pr}}\mathfrak M/(144L)}$,
\begin{equation}\label{I-17-50}
 |Y_\square\omega|\le Cm^3\xi_n^{-3}E,\qquad
 \sup_{T_\theta}|g|\le Cm^4\xi_n E .
\end{equation}

The global scalar and Gram estimates, with the common constants
of \eqref{I-17-23}, give
\[
 \mathsf G',\mathsf T,\mathsf T^*\in
       \mathfrak K_2(\delta_{\mathrm u},B_{\mathrm u}),\qquad
 \nabla\mathsf G',\nabla\mathsf T\in
       \mathfrak K_1(\delta_{\mathrm u},B_{\mathrm u}),\qquad
 \mathsf C\in\mathfrak K_{-4}(\delta_{\mathrm u},B_{\mathrm u}).
\]
Choose the additional threshold
$R_1\mathfrak M\ge\mathsf r_S(\delta_{\mathrm u}/16)$.
Then $(\delta_{\mathrm u}/16)R_1\mathfrak M
\ge2(3\log L+1)$, so \cref{I-17-29} applies successively
to $\nabla\mathsf T\cdot\mathsf C$,
$\mathsf T^*\cdot\mathsf G'$,
$\mathsf C\cdot(\mathsf T^*\mathsf G')$, and
$\nabla\mathsf T\cdot(\mathsf C(\mathsf T^*\mathsf G'))$.
The resulting orders are $-3,4,0,1$, and their rates
are respectively $\delta_{\mathrm u}/2,\delta_{\mathrm u}/2,
\delta_{\mathrm u}/4,\delta_{\mathrm u}/8$. Thus
\begin{equation}\label{I-17-52}
 \nabla\mathsf T\mathsf C\in\mathfrak K_{-3}(\delta_{\mathrm u}/2,C),
 \qquad
 \nabla\mathsf R\mathsf G'
  =\nabla\mathsf G'-\nabla\mathsf T
                      \bigl(\mathsf C(\mathsf T^*\mathsf G')\bigr)
       \in\mathfrak K_1(\delta_{\mathrm u}/8,C).
\end{equation}

For $Y[u;y_1]=(Y\mathbf1_{\{y_1\}})(u)$, the sharp versions
of \eqref{I-17-48}--\eqref{I-17-49} now show that
\[
 |Y[u;y_1]|\le C\mathfrak M^{-1}
     e^{-\delta_{\mathrm u}d_{\mathbf B}(y(u),y_1)}
       \quad(u\in T_\theta),
 \qquad Y[u;y_1]=0\quad(u\notin T_\theta).
\]
Here the neighbouring cells remain in $S_\square^{(5/6)}$,
have sides at most $L\xi_n$, and have distance at most $9L$
from $y(u)$ by \eqref{eq:nearby-cell-comparison}. The adjoint uses the two
different volume weights:
\[
 (Y^*g)(y_1)=\xi_{y_1}^{-3}\sum_u\xi^3Y[u;y_1]g(u).
\]
Counting the $(\xi_{y'}/\xi)^3$ sites in a cell and applying
\eqref{eq:global-scale-ratio} with exponent $3$ gives
\[
 |Y^*g(y_1)|\le C\mathfrak M^{-1}
      \sum_{y'\in\mathcal T}
       e^{-3\delta_{\mathrm u}d_{\mathbf B}(y',y_1)/4}
       \sup_{T_\theta}|g|.
\]

At the output bonds in \eqref{I-17-45}, and at the endpoints
of the input bonds, $\theta=1$. Consequently
$\Theta\partial^*J=\partial_0^*J$ and
$\partial\Theta f=\partial f$ at those output bonds.
Equation \eqref{I-17-47} therefore yields
\[
 (\mathcal P-\mathcal P_\square)J
   =\partial\mathsf T\mathsf CY^*g
       -\partial\mathsf R\mathsf G'Y_\square\omega
\]
there. Apply \eqref{I-17-52}, \eqref{I-17-50}, and
the site sum \eqref{I-13-122}. The second term has size at
most $Cm^4\xi_n^{-2}E$; for the first, summing first over
$y_1$ and then over the at most $m$ centres in $\mathcal T$
gives $Cm^5\xi_n^{-2}E$. With
$c_{\mathrm F}=\delta_{\mathrm{pr}}/(288L)$,
$m^5E=(8\mathfrak M)^{15}e^{-2c_{\mathrm F}\mathfrak M}
\le Ce^{-c_{\mathrm F}\mathfrak M}$, proving
\eqref{I-17-45}.
\end{proof}

\subsection{The cutoff remainder}
\label{sec:global-cutoff-remainder}

For $\square\in\mathfrak D_n$, write $h=h_\square$ and
$\zeta=\zeta_\square$, with transfers understood as in the
local construction. Define maps $\mathcal A_\square\to\mathcal A$ by
\[
 \mathsf K_\square=h\Delta_{a,\square}-\Delta_a(\Omega_0)h,
 \qquad
 \mathsf K_\square^{\mathrm{loc}}=
 \begin{cases}
   \mathsf L(h)+\mathsf L_Q(h),&\square\text{ interior},\\
   \chi_{\mathcal A}\mathsf L'(h),&\square\text{ boundary},
 \end{cases}
\]
where $\mathsf L_Q(h)=hQ_\square^*aQ_\square
-Q_\square^*aQ_\square h$ and
$\mathsf L'(h)=h\partial^*\partial-\partial^*\partial h$.
For an interior cube put
\begin{equation}\label{I-16-71}
 \begin{split}
 K_{\square,\square}
   &=\mathsf K_\square^{\mathrm{loc}}
      +\zeta(\mathcal P-\mathcal P_\square)h
      +\zeta(\mathcal P_\square h-h\mathcal P_\square),\\
 K_{\square',\square}
   &=h_{\square'}^2(1-\zeta)\mathcal Ph
                         \quad(\square'\ne\square).
 \end{split}
\end{equation}
For a boundary cube set $K_{\square,\square}
=\mathsf K_\square^{\mathrm{loc}}$ and
$K_{\square',\square}=0$ for $\square'\ne\square$.
Thus the error in patching the local inverses consists of a
local commutator, a change of gauge projection, a commutator
with the local gauge projection, and the part of the global
projection outside $\zeta$.

\begin{proposition}\label{I-17-54}\label{I-17-68}
There are $\mathfrak m_{\mathrm H},\mathfrak r_{\mathrm H}$,
$\kappa_{\mathrm H}>0$ and $C_{\mathrm H}\ge1$, depending
only on $L$, such that the following holds when
$\mathfrak M\ge\mathfrak m_{\mathrm H}$ and
$R_1\mathfrak M\ge\mathfrak r_{\mathrm H}$.
The operators
\[
 \mathsf G_0=\sum_\square h_\square\mathsf G_\square h_\square,
 \qquad
 \mathcal R=\sum_{\square',\square}
         K_{\square',\square}\mathsf G_\square h_\square
\]
satisfy $\mathsf G_0^*=\mathsf G_0$ and
$\Delta_a(\Omega_0)\mathsf G_0=I-\mathcal R$.
For $y,y'\in\Lambda_*$ put
$e_{\mathrm H}=e^{-\kappa_{\mathrm H}d_{\mathbf B}(y,y')}$.
Then:
\begin{enumerate}[label=(H\arabic*)]
\item For $p=2,\infty$,
\[
 \sum_{\square',\square}
  N_p(K_{\square',\square}\mathsf G_\square h_\square;y,y')
       \le C_{\mathrm H}\mathfrak M^{-1}e_{\mathrm H}.
\]
\item The adjoint blocks satisfy
\[
 \sum_{\square',\square}
  N_\infty(h_\square\mathsf G_\square K_{\square',\square}^*;y,y')
       \le C_{\mathrm H}\mathfrak M^{-1}e_{\mathrm H}.
\]
The same bound holds for
$\sum_\square N_\infty(h_\square\mathsf G_\square
\mathsf K_\square^*;y,y')$.
\item If $0\le\alpha<1$, $\psi$ is a function of bare sites
supported in $\widetilde\Delta(y)$, and $A\in\mathcal A$, then
\[
 \begin{split}
 &\sum_{\square',\square}
   \|\psi h_\square\mathsf G_\square
                      K_{\square',\square}^*A\|_{\alpha,\xi_y}\\
 &\qquad\le C_{\mathrm H}(\alpha)\mathfrak M^{-1}
       (\|\psi\|_{\alpha,\xi_y}+|\psi|)
       \sum_{y''}e^{-\kappa_{\mathrm H}d_{\mathbf B}(y,y'')}
                         |\chi_{y''}A|.
 \end{split}
\]
In particular this bounds $\|\psi\mathcal R^*A\|_{\alpha,\xi_y}$.
\item If $0<\varepsilon<1$ and
$J\in\mathcal A^{(1)}$ is supported in
$\widetilde\Delta(y')$, then
\[
 \begin{split}
 \sum_{\square',\square}
   |\chi_yK_{\square',\square}\mathsf G_\square
                         h_\square\nabla^*J|
 &\le C'_{\mathrm H}(\varepsilon)\mathfrak M^{-1}
      \xi_{y'}^{-1}e_{\mathrm H}
      (\|J\|_{\varepsilon,\xi_{y'}}+|J|),\\
 \sum_{\square',\square}
   \|\chi_yK_{\square',\square}\mathsf G_\square
                         h_\square\nabla^*J\|
 &\le C_{\mathrm H}\mathfrak M^{-1}
      \xi_{y'}^{-1}e_{\mathrm H}\|J\|.
 \end{split}
\]
\end{enumerate}
The constants with an exponent as argument depend only on
$L$ and that exponent. One may choose
\begin{equation}\label{I-17-69}
 \kappa_{\mathrm H}=
 \min\left\{\frac{\delta_{\mathrm{loc}}}{2L},
             \frac{\delta_{\mathrm{pr}}}{2L},
             \frac{\delta_{\mathrm u}}4,
             \frac{c_{\mathrm F}}{12L}\right\}.
\end{equation}
\end{proposition}

\begin{proof}
\emph{Local cells and transfers.}
\label{I-17-60}
For an interior cube, let $w$ range over its local coarse
centres. For a boundary cube let $w$ range over all bare
sites of the ambient torus, including those outside
$\Omega_0$, with $\xi_w=\xi_n=\xi$.
Write $B_w$ for the corresponding bare cell and $\chi_w$
for localization to it. In either case
$\xi_w\in\{\xi_n,\xi_{n+1}\}$ and
$f=\sum_w\chi_wf$ on the local spaces; localization preserves
those spaces. Cells of the global sequence meeting
$S_\square^{(5/6)}$ lie inside $\square$ and are local
cells with the same multiplications. For an interior cube
this is the common-cell assertion proved with \eqref{I-17-15};
for a boundary cube the set is disjoint from $\Omega_1$,
so only bare cells occur.

For $s\in S_\square^{(5/6)}$, the sup ball of radius
$4\xi_{n+1}$ about $s$ lies in the common lifted interior:
its distance from the boundary is positive because
$4\xi_{n+1}<\mathfrak M\xi_n/6$.
Two points in this ball differ by at most $8\xi_{n+1}$
in each coordinate, no more than half the side of either
torus. The map $\iota_\square$ therefore preserves their
$\ell^1$ and sup distances. In particular it preserves
the scaled seminorms at scales at most $\xi_{n+1}$ of
functions supported in $S_\square^{(5/6)}$, supports in
$\mathsf q(o,\lambda)$ for $o$ in that set and
$\lambda\le4\xi_{n+1}$, and the
$2\xi_{n+1}$-neighbourhoods of cells meeting that set.

We also use the following uniform local lattice sum.
The centres $w$ are distinct points of the lattice of
spacing $\xi_n$. Choose a lift $\widetilde o$ of $o$ and
a point $o_0\in\xi_n\mathbb Z^3$ within sup distance
$\xi_n/2$ of it. For each $w$ choose a closest lift
$\widetilde w$. The vectors
$v=(\widetilde w-o_0)/\xi_n$ are distinct and
$|o-w|_\square\ge|v|_1-3$. Consequently, for decreasing
$F\ge0$,
\begin{equation}\label{I-17-61}
 \sum_wF(|o-w|_\square)
 \le\sum_{v\in\mathbb Z^3}F(\max\{|v|_1-3,0\}),
 \qquad
 \sum_we^{-\mu|o-w|_\square}
       \le e^{3\mu}\coth^3(\mu/2).
\end{equation}
The second bound follows by multiplying the three
one-dimensional geometric sums.

\emph{The identity and the local commutator.}
For an interior cube, the support of $hv$ is separated from
$\Omega_0^c$. The local identities \eqref{I-17-15} and the
Hodge identity give
\[
 \Delta_a(\Omega_0)(hv)
    =\Delta(hv)-\mathcal P(hv)+Q_\square^*aQ_\square(hv),
 \qquad
 h\Delta_{a,\square}v
    =h\Delta v-h\mathcal P_\square v
                            +hQ_\square^*aQ_\square v.
\]
Indeed $\partial\partial^*(hv)
=\partial_0(\mathsf P+\mathsf R)\partial_0^*(hv)$;
the term with $\mathsf R$ cancels the gauge term of
$\Delta_a$. Since $\zeta h=h$, subtraction yields
\begin{equation}\label{I-17-55}
 \mathsf K_\square
  =\mathsf K_\square^{\mathrm{loc}}
    +(1-\zeta)\mathcal Ph
    +\zeta(\mathcal P-\mathcal P_\square)h
    +\zeta(\mathcal P_\square h-h\mathcal P_\square).
\end{equation}
The bonds occurring in $\Delta(hv)$ have both endpoints
in $S_\square^{(5/6)}\subset\Omega_0$, so the displayed
local commutators take values in $\mathcal A$.

For a boundary cube $n=0$. The support of $hv$ and the
plaquettes meeting it lie within $2\xi$ of $S_\square$,
at distance at least $2\mathfrak M\xi-2\xi$ from $\Omega_1$.
They consequently meet no averaging bond of positive level:
the level-one blocks lie within $2\xi_1$ of $\Omega_1$,
and the blocks at level $j\ge2$ lie in
$\Omega_{j-1}\subset\Omega_1$ by separation.
On the remaining bonds $Q$ is the identity, so
$Q^*aQ(hv)=\xi^{-2}hv$. The support assertion of
\cref{I-14-48}(ii) also gives
$\partial_0\mathsf R\partial_0^*(hv)=0$:
the divergence of $hv$ vanishes on $\Omega_1$ and its
bare nearest neighbours, which contain the range of
$\mathsf R$.
The mass terms therefore cancel between the global and
local operators, leaving
$\mathsf K_\square=\chi_{\mathcal A}\mathsf L'(h)$.
Here the restrictions agree on every plaquette meeting
$\operatorname{supp}(hv)$, by the boundary margin in
\cref{I-17-14}.

The quadratic partition and $\zeta=1$ on $S_\square$
now give $\mathsf K_\square=\sum_{\square'}K_{\square',\square}$.
Thus
\[
 \begin{split}
 \Delta_a(\Omega_0)\mathsf G_0
  &=\sum_\square
     (h_\square\Delta_{a,\square}-\mathsf K_\square)
                     \mathsf G_\square h_\square\\
  &=\sum_\square h_\square^2-\mathcal R=I-\mathcal R.
 \end{split}
\]
The local inverses are symmetric and opposite transfers
by $h_\square$ are adjoint, which proves
$\mathsf G_0^*=\mathsf G_0$.

We will need the local commutator in two forms. Let
$\chi_X^+$ indicate the $2\xi_{n+1}$-neighbourhood of $X$.
For $p=2,\infty$,
\begin{equation}\label{I-17-56}
 \begin{split}
 \|\chi_X\mathsf K_\square^{\mathrm{loc}}v\|_p
   &\le C\mathfrak M^{-1}
       \bigl(\xi_n^{-1}\|\chi_X^+\nabla v\|_p
                 +\xi_n^{-2}\|\chi_X^+v\|_p\bigr),\\
 (\mathsf K_\square^{\mathrm{loc}})^*A
   &=\nabla^*\mathsf w_\square(A)+\mathsf v_\square(A),\\
 \xi_n\|\chi_X\mathsf w_\square(A)\|_p
   +\xi_n^2\|\chi_X\mathsf v_\square(A)\|_p
   &\le C\mathfrak M^{-1}\|\chi_X^+A\|_p .
 \end{split}
\end{equation}
All four outputs in this display are supported in the
$2\xi_{n+1}$-neighbourhood of $S_\square$, contained in
$S_\square^{(1/6)}$. To verify the estimates, write
$(Q_\square v)(c)=\sum_b w_c(b)v(b)$. Expanding the
averaging and curl commutators gives
\begin{align}
 (\mathsf L_Q(h)v)(b)
 &=\xi^{-3}\sum_{j,c\in\Gamma_j}\xi_j w_c(b)
        \sum_{b'}w_c(b')
          (h(b_-)-h(b'_-))v(b'),\label{I-17-57}\\
 (\mathsf L'(h)v)(b)
 &=\xi^{-2}\sum_{\mathsf f\ni b}\sum_{b'\in\partial\mathsf f}
          \pm(h(b_-)-h(b'_-))v(b').\label{I-17-58}
\end{align}
Only $j=n,n+1$ occur in the first line. The initial point
of $b$ belongs to a unique block of that level; a relevant
coarse bond has this block centre as one endpoint and the
direction of $b$, so there are at most two choices.
The straight-path average \eqref{I-16-63} has nonnegative
weights, total weight one, and $w_c(b)\le L^{-3j}$.
Consequently
$\xi^{-3}\xi_j w_c(b)\le\xi_j^{-2}$ and
$\sum_{b'}w_c(b')=1$. The two fine bonds in a nonzero
summand are at sup distance at most $2\xi_{n+1}$.
Their $\ell^1$ distance is at most $4\xi_{n+1}$;
the Lipschitz bound on $h$ consequently makes the absolute
row sum at most $C\mathfrak M^{-1}\xi_n^{-2}$.
In the second line four plaquettes meet $b$, with three
nontrivial terms each; their initial points are at
$\ell^1$ distance at most $2\xi$. Since $n=0$ there,
its row sum has the same bound.
Both matrices are antisymmetric in the equal fine weights,
so their column sums obey the same estimate. The Schur
bound in \cref{I-17-29} proves their localized
$p=2,\infty$ estimates.

For $\mathsf L(h)$, the first form of \eqref{I-17-33}
gives the first line of \eqref{I-17-56}, since
$r_n^{-1}=6L\mathfrak M^{-1}\xi_n^{-1}$.
Its second form supplies the required divergence:
for interior cubes,
\[
 \mathsf w_\square(A)_\nu
       =(\nabla_\nu h)(A+A^{+\nu}),\qquad
 \mathsf v_\square(A)=-(\Delta h)A-\mathsf L_Q(h)A.
\]
For boundary cubes take
$\mathsf w_\square(A)=0$ and
$\mathsf v_\square(A)=-\chi_{\mathcal A}\mathsf L'(h)A$.
The same derivative and row-sum estimates prove the last
line of \eqref{I-17-56}. The shifts have length $\xi$;
the averaging matrices have range $2\xi_{n+1}$, and
$2\xi_{n+1}\le r_n/6$. This proves the support assertion
and also $(\mathsf K_\square^{\mathrm{loc}})^*
=-\mathsf K_\square^{\mathrm{loc}}$ under the local identification.

\emph{The three gauge terms.}
\label{I-17-62}
For interior local cells $w,w'$ and $p=2,\infty$, we claim
\begin{equation}\label{I-17-64}
 \|\chi_w(\mathcal P_\square h-h\mathcal P_\square)
                      \chi_{w'}\|_{p\to p}
 \le C\mathfrak M^{-1}\xi_n^{-2}
       (1+|w-w'|_\square)e^{-\delta_{\mathrm{pr}}|w-w'|_\square}.
\end{equation}
The counting kernel in \eqref{eq:local-scalar-Gram} has the factor
$(\xi/\xi_{w'})^3\xi_n^{-2}$. There are
$3(\xi_{w'}/\xi)^3$ fine bonds in the input cell,
so its row sum is at most
$3B_{\mathrm{pr}}\xi_n^{-2}
e^{-\delta_{\mathrm{pr}}|w-w'|_\square}$.
For points in the two cells,
$|x-x'|_1\le\xi_n(|w-w'|_\square+3L)$.
Multiplication by $h(x')-h(x)$ and
$\operatorname{Lip}(h)\le c_\flat/r_n$ proves
\eqref{I-17-64} for $p=\infty$.
The commutator is antisymmetric; its opposite block has
the same bound. Schur's estimate gives $p=2$.

By \cref{I-17-43}, the two mutually adjoint transfers
\[
 \zeta(\mathcal P-\mathcal P_\square)h,\qquad
 h(\mathcal P-\mathcal P_\square)\zeta
\]
have norm at most $C_{\mathrm F}\xi_n^{-2}
e^{-c_{\mathrm F}\mathfrak M}$ for $p=\infty$.
Their input and output are supported in $S_\square^{(1/3)}$,
where initial points are localization points.
Applying Schur's estimate to this pair proves the same
bound for $p=2$.

For the global gauge operator, the fourth kernel in
\eqref{I-14-49}, the at most
$6(\xi_{y''}/\xi)^3$ input bonds in a global cell,
and symmetry give
\[
 \begin{split}
 N_\infty(\mathcal P;y,y'')
      &\le6B_{\mathrm u}\xi_y^{-2}
                    e^{-\delta_{\mathrm u}d_{\mathbf B}(y,y'')},\\
 N_2(\mathcal P;y,y'')
      &\le6B_{\mathrm u}\xi_y^{-1}\xi_{y''}^{-1}
                    e^{-\delta_{\mathrm u}d_{\mathbf B}(y,y'')}.
 \end{split}
\]
If $(\delta_{\mathrm u}/4)R_1\mathfrak M\ge2\log L$,
the scale-ratio estimate gives, in either order,
\begin{equation}\label{I-17-65}
 N_p(\mathcal P;y,y''),\ N_p(\mathcal P;y'',y)
  \le6B_{\mathrm u}L^2\xi_{y''}^{-2}
               e^{-3\delta_{\mathrm u}d_{\mathbf B}(y,y'')/4}.
\end{equation}

We combine these estimates in a form which can be applied
to every local inverse. Suppose $0<\delta_\natural\le1$,
$y'\in\Lambda_*\cap S_\square^{(1/6)}$, and
$u\in\mathcal A_\square$ satisfies
\begin{equation}\label{I-17-66}
 \xi_n^{-2}\|\chi_{B_w}^+u\|_p+
 \xi_n^{-1}\|\chi_{B_w}^+\nabla u\|_p
     \le Ue^{-\delta_\natural|w-y'|_\square}
\end{equation}
for every local cell $w$. Set
\begin{equation}\label{I-17-63}
 \kappa_\natural=\min\left\{
   \frac{\delta_\natural}{2L},\frac{\delta_{\mathrm{pr}}}{2L},
   \frac{\delta_{\mathrm u}}4,\frac{c_{\mathrm F}}{12L}\right\}.
\end{equation}
Under the preceding thresholds and
$R_1\mathfrak M\ge\mathsf r_S(\kappa_\natural/2)$, we have
\begin{equation}\label{I-17-67}
 \sum_{\square'}\|\chi_yK_{\square',\square}u\|_p
 \le C(\delta_\natural)\mathfrak M^{-1}U
                         e^{-\kappa_\natural d_{\mathbf B}(y,y')}.
\end{equation}
Here and below the constants depend only on $L$ and the
indicated exponent or rate.

The local term vanishes unless $B_y$ meets
$S_\square^{(1/6)}$. It is then a common local cell, so
\eqref{I-17-56} and \eqref{I-17-66} give
$C\mathfrak M^{-1}Ue^{-\delta_\natural|y-y'|_\square}$.
By \eqref{eq:local-cube-distance},
$d_{\mathbf B}(y,y')\le L|y-y'|_\square$.
For the comparison term, decompose $u$ into local cells.
The local lattice sum and the preceding operator bound give
\[
 \|\chi_y\zeta(\mathcal P-\mathcal P_\square)hu\|_p
       \le Ce^{-c_{\mathrm F}\mathfrak M}U.
\]
If this is nonzero, $y,y'$ lie in the interior of $\square$,
where $d_{\mathbf B}(y,y')\le6L\mathfrak M$. Thus
\[
 e^{-c_{\mathrm F}\mathfrak M}
 \le e^{-c_{\mathrm F}\mathfrak M/2}
                 e^{-c_{\mathrm F}d_{\mathbf B}(y,y')/(12L)}.
\]
For the local gauge commutator, put
$\mu=\frac12\min\{\delta_{\mathrm{pr}},\delta_\natural\}$.
Equations \eqref{I-17-64} and \eqref{I-17-66} give
\[
 \begin{split}
 &\|\chi_y\zeta(\mathcal P_\square h-h\mathcal P_\square)u\|_p\\
 &\quad\le C\mathfrak M^{-1}U
       \sum_w(1+|y-w|_\square)
        e^{-\delta_{\mathrm{pr}}|y-w|_\square
                         -\delta_\natural|w-y'|_\square}\\
 &\quad\le C\mathfrak M^{-1}U
                         e^{-(\mu/L)d_{\mathbf B}(y,y')}.
 \end{split}
\]
Indeed the exponent is at least
$\mu|y-y'|_\square+\mu|y-w|_\square$, and
$(1+s)e^{-\mu s}\le(2/\mu)e^{-\mu s/2}$, whose local
lattice sum is finite uniformly in the torus.

There are at most $16$ cubes $\square'$ whose cutoffs
meet a bond localized in $B_y$, so the remaining sum is
at most $16\|\chi_y(1-\zeta)\mathcal Phu\|_p$.
It vanishes unless $B_y\not\subset S_\square^{(1/6)}$.
For an exiting bond this follows as well: its localization
point differs from its initial point by $\xi$, whereas
$S_\square^{(5/6)}$ is at distance greater than
$\mathfrak M\xi$ from $\Omega_0^c$.
Decompose $hu$ over the cells $y''$ meeting $S_\square$.
They are common cells with $\xi_{y''}\ge\xi_n$, and the
buffer estimate \eqref{eq:local-cube-buffer} gives
$d_{\mathbf B}(y,y'')\ge\mathfrak M/(72L^2)$.
Using \eqref{I-17-65} and \eqref{I-17-66}, their sum is at most
\[
 CU\sum_{y''}
 e^{-3\delta_{\mathrm u}d_{\mathbf B}(y,y'')/4
                  -(\delta_\natural/L)d_{\mathbf B}(y'',y')}.
\]
Reserve $\delta_{\mathrm u}/4$ of the first rate for the
buffer. The convolution estimate \eqref{I-17-30}, at
rate $2\kappa_\natural\le
\min\{\delta_{\mathrm u}/2,\delta_\natural/L\}$, then gives
\[
 C U e^{-\delta_{\mathrm u}\mathfrak M/(288L^2)}
          c_S(\kappa_\natural)
          e^{-\kappa_\natural d_{\mathbf B}(y,y')}.
\]
The exponentially small factors in the comparison and
off-support terms are bounded by $C\mathfrak M^{-1}$.
Every retained spatial rate is at least $\kappa_\natural$;
this proves \eqref{I-17-67}.

\emph{Application to local inverses.}
Take $\delta_\natural=\delta_{\mathrm{loc}}$, so that
$\kappa_\natural=\kappa_{\mathrm H}$.
Let $f=\chi_{y'}f$. There are at most $16$ cubes with
$h_\square f\ne0$: their supports meet
$\widetilde\Delta(y')$. For each,
$y'\in S_\square^{(1/6)}$ and
$\xi_{y'}\in\{\xi_n,\xi_{n+1}\}$.
The $2\xi_{n+1}$-neighbourhood of a local cell $B_w$,
including the initial points of its bonds, lies in
$\mathsf q(w,3\xi_{n+1})$. Cover this by $27$ cubes
$\mathsf q(o_i,\xi_{n+1})$, with
$o_i=w+2\xi_{n+1}k$, $k\in\{-1,0,1\}^3$.
Their centres satisfy $|o_i-w|_\square\le6L$.
Apply the value and gradient estimates of
\cref{I-17-38}(R1), or (R5) for $p=2$, to
$u=\mathsf G_\square h_\square f$ on these cubes.
Since the input is supported in
$\mathsf q(y',\xi_{y'})$, we obtain \eqref{I-17-66} with
$U=C\|f\|_p$. Equation \eqref{I-17-67} bounds the sum
over $\square'$ by
$C\mathfrak M^{-1}e_{\mathrm H}\|f\|_p$.
To pass to the sum of block norms in (H1), observe that
for each of the at most $16$ input cubes there are
at most $17$ nonzero output blocks: the diagonal block,
and at most $16$ cubes whose supports meet
$\widetilde\Delta(y)$. Each block norm is bounded by the
preceding estimate; summing this fixed number proves (H1).

For (H4), the transferred product rule \eqref{I-17-36} gives
\[
 h\nabla^*J=\nabla^*(hJ)+\mathsf m_h(J),\qquad
 \mathsf m_h(J)=\sum_\nu
                    (\nabla_\nu h)^{-\nu}J_\nu^{-\nu},
\]
with the boundary compression prescribed there.
The left side vanishes unless $S_\square$ meets
$\mathsf q(y',3\xi_{y'})$; again there are at most $16$
such cubes, with $y'\in S_\square^{(1/6)}$ and
$\xi_{y'}\in\{\xi_n,\xi_{n+1}\}$.
Require $\mathfrak M\ge6L^2c_\flat$. By
\eqref{I-17-34},
\[
 \|hJ\|_{\varepsilon,\xi_{y'}}
       \le\|J\|_{\varepsilon,\xi_{y'}}+|J|,
 \qquad
 \|\mathsf m_h(J)\|_p
       \le3c_\flat r_n^{-1}\|J\|_p
       \le\xi_n^{-1}\|J\|_p.
\]
The form $hJ$ is supported in $\mathsf q(y',\xi_{y'})$.
The shifted correction is supported in
$\mathsf q(y',3\xi_{y'})$; partition its bonds among the
$27$ cubes of radius $\xi_{y'}$ covering that cube.
Use the same $27$-cube cover on the output.
The divergence value and gradient estimates (R1), (R3)
of \cref{I-17-38}, applied to $hJ$, and its value and
gradient estimates applied to the correction pieces, give
\eqref{I-17-66} with
\[
 U=C(\varepsilon)\xi_n^{-1}
                 (\|J\|_{\varepsilon,\xi_{y'}}+|J|)
       \quad(p=\infty).
\]
For $p=2$, the four corresponding entries of (R5) give
$U=C\xi_n^{-1}\|J\|$. The displaced input and output centres
cost at most $e^{12L}$ in these estimates.
Now \eqref{I-17-67}, the finite cube count, and
$\xi_n^{-1}\le L\xi_{y'}^{-1}$ prove both assertions of (H4).

\emph{The adjoint remainder.}
For (H2) and (H3), fix an output cube $\square$ for which
$h_\square$ meets $\widetilde\Delta(y)$. There are at most
$16$ choices, and for each
$y\in S_\square^{(1/6)}$ and
$\xi_y\in\{\xi_n,\xi_{n+1}\}$.
For local forms $\mathsf w\in\mathcal A_\square^{(1)}$
and $\mathsf s\in\mathcal A_\square$, decompose into local
cells and apply the value and divergence entries of
(R1) and (R2) in \cref{I-17-38}. The product rule gives
$\|\psi h\|_{\alpha,\xi_y}\le
\|\psi\|_{\alpha,\xi_y}+|\psi|$ and $|\psi h|\le|\psi|$.
We obtain
\begin{equation}\label{I-17-70}
 \begin{split}
 |\chi_yh\mathsf G_\square(\nabla^*\mathsf w+\mathsf s)|
       &\le C\Sigma(\mathsf w,\mathsf s),\\
 \|\psi h\mathsf G_\square(\nabla^*\mathsf w+\mathsf s)\|
                       _{\alpha,\xi_y}
       &\le C(\alpha)(\|\psi\|_{\alpha,\xi_y}+|\psi|)
                                  \Sigma(\mathsf w,\mathsf s),\\
 \Sigma(\mathsf w,\mathsf s)
       &=\sum_we^{-\delta_{\mathrm{loc}}|y-w|_\square}
            \bigl(\xi_n|\chi_w\mathsf w|
                         +\xi_n^2|\chi_w\mathsf s|\bigr).
 \end{split}
\end{equation}
The support and seminorm transfers used here are those
verified at the start of the proof.

Let $A=\chi_{y'}A$. Taking adjoints in \eqref{I-16-71},
and using \eqref{I-17-56}, gives
\[
 \begin{split}
 K_{\square,\square}^*A
   &=\nabla^*\mathsf w_\square(A)+\mathsf v_\square(A)
       +h(\mathcal P-\mathcal P_\square)\zeta A
       -(\mathcal P_\square h-h\mathcal P_\square)\zeta A,\\
 K_{\square',\square}^*A
   &=h\mathcal P(1-\zeta)h_{\square'}^2A
                                    \quad(\square'\ne\square).
 \end{split}
\]
Only the first two local terms occur for a boundary cube.
We estimate the contribution of each term to $\Sigma$.
For the local terms, \eqref{I-17-56} gives
$\xi_n|\mathsf w_\square(A)|
+\xi_n^2|\mathsf v_\square(A)|
\le C\mathfrak M^{-1}|A|$.
They are supported within $2\xi_{n+1}$ of both
$B_{y'}$ and $S_\square$. Unless
$B_{y'}\cap S_\square^{(1/3)}\ne\emptyset$ they vanish.
Otherwise $B_{y'}$ is a common cell and the support meets
at most $(7L)^3$ local cells, whose centres obey
$|w-y'|_\square\le12L$. Their contribution is at most
\[
 C\mathfrak M^{-1}
  e^{-\delta_{\mathrm{loc}}|y-y'|_\square}|A|
 \le C\mathfrak M^{-1}
  e^{-(\delta_{\mathrm{loc}}/L)d_{\mathbf B}(y,y')}|A|.
\]

The comparison term also vanishes unless
$B_{y'}$ meets $S_\square^{(1/3)}$. Its norm bound,
multiplied by $\xi_n^2$ in $\Sigma$ and summed with
\eqref{I-17-61}, contributes at most
\[
 Ce^{-c_{\mathrm F}\mathfrak M}|A|
 \le Ce^{-c_{\mathrm F}\mathfrak M/2}
       e^{-c_{\mathrm F}d_{\mathbf B}(y,y')/(12L)}|A|.
\]
For the local gauge commutator, put
$\mu=\frac12\min\{\delta_{\mathrm{pr}},\delta_{\mathrm{loc}}\}$.
The factor $\xi_n^{-2}$ in \eqref{I-17-64} cancels the
$\xi_n^2$ in $\Sigma$. Its contribution is bounded by
\[
 C\mathfrak M^{-1}|A|
  \sum_we^{-\delta_{\mathrm{loc}}|y-w|_\square}
     (1+|w-y'|_\square)e^{-\delta_{\mathrm{pr}}|w-y'|_\square}
 \le C\mathfrak M^{-1}
        e^{-(\mu/L)d_{\mathbf B}(y,y')}|A|.
\]
This is the same local convolution used in
\eqref{I-17-67}, with the polynomial factor on the
input distance.

Finally put $A'=h_{\square'}^2A$.
Then $|A'|\le|A|$, $A'=\chi_{y'}A'$, and at most $16$
cubes $\square'$ give a nonzero $A'$.
The form $h\mathcal P(1-\zeta)A'$ is supported in common
cells $y''$ meeting $S_\square$ and vanishes unless
$B_{y'}\not\subset S_\square^{(1/6)}$.
For these cells \eqref{eq:local-cube-buffer} gives
$d_{\mathbf B}(y'',y')\ge\mathfrak M/(72L^2)$.
Use the reverse-order bound in \eqref{I-17-65},
with $\xi_{y''}^{-2}\le\xi_n^{-2}$. After cancellation
against $\xi_n^2$ in $\Sigma$, the contribution is at most
\[
 \begin{split}
 &C|A|\sum_{y''}
    e^{-(\delta_{\mathrm{loc}}/L)d_{\mathbf B}(y,y'')
                         -3\delta_{\mathrm u}d_{\mathbf B}(y'',y')/4}\\
 &\qquad\le
 C e^{-\delta_{\mathrm u}\mathfrak M/(288L^2)}
       c_S(\kappa_{\mathrm H})
       e^{-\kappa_{\mathrm H}d_{\mathbf B}(y,y')}|A|.
 \end{split}
\]
Here we again reserve $\delta_{\mathrm u}/4$ for the
buffer and apply \eqref{I-17-30} to the remaining two rates.

All four contributions are bounded by
$C\mathfrak M^{-1}e_{\mathrm H}|A|$.
Insert this bound into the first line of \eqref{I-17-70},
take the supremum over $A=\chi_{y'}A$ with $|A|\le1$,
and sum the at most $16\cdot17$ pairs of cubes.
This proves (H2). For (H3), write
$A=\sum_{y''}\chi_{y''}A$ and apply the second line of
\eqref{I-17-70} to each component. Sum first over those
same finitely many cube pairs and then over $y''$.
The additional assertions for $\mathsf K_\square^*$ and
$\mathcal R^*$ follow by their defining sums and the
triangle inequality.

All requirements on $\mathfrak M$ used above are the local
model and gauge-comparison thresholds together with
$\mathfrak M\ge6L^2c_\flat$.
Choose their maximum as $\mathfrak m_{\mathrm H}$.
Choose $\mathfrak r_{\mathrm H}$ to dominate the scalar,
Gram and gauge-comparison thresholds and
$\mathsf r_S(\kappa_{\mathrm H}/2)$.
The condition needed in \eqref{I-17-65} is already implied
by $\mathsf r_S(\delta_{\mathrm u}/16)$.
These are finitely many constants depending only on $L$.
\end{proof}

\subsection{Estimates for the approximate inverse}

The local bounds pass through the partition with constants
independent of the number of cubes. The point is that only
boundedly many cutoffs meet an output cell; derivatives of
these cutoffs have the scale needed by the lower-order local
estimates.

\begin{lemma}\label{I-17-71}
Suppose
$\mathfrak M\ge\max\{L^5,L\mathfrak m_{\mathrm{pr}},6L^2c_\flat\}$.
Put $\kappa_{\mathrm a}=\delta_{\mathrm{loc}}/L$, let
$\widetilde\chi_y$ indicate the bonds with localization point
in $\mathsf q(y,2\xi_y)$, and write
\[
 \Sigma_p(f)=\sum_z e^{-\kappa_{\mathrm a}d_{\mathbf B}(y,z)}
                         \|\chi_zf\|_p,\qquad
 e_{\mathrm a}=e^{-\kappa_{\mathrm a}d_{\mathbf B}(y,y')}.
\]
Let $\zeta$ be supported in $\widetilde\Delta(y)$, and let
$0\le\alpha<1$, $\varepsilon>0$, $\alpha+\varepsilon<1$.
With constants depending only on $L$ and the displayed exponents,
the following estimates hold for $f\in\mathcal A$:
\begin{enumerate}[label=(G\arabic*)]
\item
$\xi_y^{-2}|\widetilde\chi_y\mathsf G_0f|$,
$\xi_y^{-1}|\widetilde\chi_y\nabla\mathsf G_0f|$ and
$|\widetilde\chi_y\Delta\mathsf G_0f|$
are at most $C\Sigma_\infty(f)$.
\item
$\xi_y^{-1}\|\zeta\nabla\mathsf G_0f\|_{\alpha,\xi_y}
 \le C(\alpha)(\|\zeta\|_{\alpha,\xi_y}+|\zeta|)
                                  \Sigma_\infty(f)$.
\item
$\xi_y^{-2}\|\zeta\mathsf G_0f\|$,
$\xi_y^{-1}\|\zeta\nabla\mathsf G_0f\|$ and
$\|\zeta\nabla\nabla\mathsf G_0f\|$
are at most $C|\zeta|\Sigma_2(f)$.
\end{enumerate}
For $J\in\mathcal A^{(1)}$ supported in
$\widetilde\Delta(y')$, the differentiated-data estimates are
\begin{enumerate}[label=(G\arabic*),start=4]
\item
\[
 \begin{split}
 \xi_y^{-1}|\widetilde\chi_y\mathsf G_0\nabla^*J|
       &\le Ce_{\mathrm a}|J|,\\
 \xi_y^{-1}\|\zeta\mathsf G_0\nabla^*J\|_{\alpha,\xi_y}
       &\le C(\alpha)(\|\zeta\|_{\alpha,\xi_y}+|\zeta|)
                                      e_{\mathrm a}|J|.
 \end{split}
\]
\item
\[
 \begin{split}
 |\widetilde\chi_y\nabla\mathsf G_0\nabla^*J|
   &\le C(\varepsilon)e_{\mathrm a}
                         (\|J\|_{\varepsilon,\xi_{y'}}+|J|),\\
 \|\zeta\nabla\mathsf G_0\nabla^*J\|_{\alpha,\xi_y}
   &\le C(\alpha,\varepsilon)
         (\|\zeta\|_{\alpha,\xi_y}+|\zeta|)e_{\mathrm a}
                         (\|J\|_{\alpha+\varepsilon,\xi_{y'}}+|J|).
 \end{split}
\]
\item
$\xi_y^{-1}\|\zeta\mathsf G_0\nabla^*J\|$,
$\|\zeta\nabla\mathsf G_0\nabla^*J\|$ and
$\|\zeta\mathsf G_0\nabla^*\nabla^*J\|$
are at most $C|\zeta|e_{\mathrm a}\|J\|$.
In the third estimate $J\in\mathcal A^{(2)}$ with the same support.
\end{enumerate}
Each estimate remains valid if its left side is replaced by
the sum over $\square$ of the corresponding expression with
$h_\square\mathsf G_\square h_\square$ in place of $\mathsf G_0$.
\end{lemma}

\begin{proof}
\emph{Support and scale.}
For $\square\in\mathfrak D_n$, put $h=h_\square$.
A difference of order at most two moves an initial point by
at most $2\xi$; the localization point differs from the
initial point by at most $\xi$. Thus an output in
$\mathsf q(y,2\xi_y)$ can occur only if $S_\square$ meets
$\mathsf q(y,6\xi_y)$. The local support geometry gives at
most $16$ such cubes, with
\[
 \xi_n\le\xi_y\le L\xi_n,\qquad y\in S_\square^{(1/6)}.
\]
Likewise, $h\chi_zf$ requires an intersection with
$\mathsf q(z,2\xi_z)$, and a term containing
$h\nabla^*J$ or $h\nabla^*\nabla^*J$ requires one with
$\mathsf q(y',4\xi_{y'})$. The corresponding input scales
belong to $\{\xi_n,\xi_{n+1}\}$ and the centers to
$S_\square^{(1/6)}$. The local metric comparison then gives
\[
 e^{-\delta_{\mathrm{loc}}|y-z|_\square}
       \le e^{-\kappa_{\mathrm a}d_{\mathbf B}(y,z)},
\]
and the same inequality with $y'$ in place of $z$.
The faithful transfer of these supports also preserves the
suprema and seminorms used below.

\emph{One product calculation.}
The cutoff bounds show that the functions
\begin{equation}\label{I-17-72}
 h,\qquad \xi_n\nabla_\nu h,\qquad
 \xi_n^2\nabla_\nu\nabla_\rho h
\end{equation}
and their shifts have suprema at most $1$, under the stated
threshold. For the first two, their seminorms at radii at most
$\xi_{n+1}$ are also at most $1$: use
\eqref{I-17-34} and $r_n^{-1}=6L\mathfrak M^{-1}\xi_n^{-1}$.

Use \eqref{I-17-32} once or twice on the output and
\eqref{I-17-36} on the input. For the second input divergence,
the exact identity is
\begin{equation}\label{I-17-74}
 \begin{split}
 h\nabla_\nu^*\nabla_\rho^*J_{\nu\rho}
 &=\nabla_\nu^*\nabla_\rho^*(hJ_{\nu\rho})\\
 &\quad+\nabla_\nu^*
       \bigl((\nabla_\rho h)^{-\rho}J_{\nu\rho}^{-\rho}\bigr)
       +\nabla_\rho^*
       \bigl((\nabla_\nu h)^{-\nu}J_{\nu\rho}^{-\nu}\bigr)\\
 &\quad+(\nabla_\rho\nabla_\nu h)^{-\nu-\rho}
                          J_{\nu\rho}^{-\nu-\rho}.
 \end{split}
\end{equation}
This follows by three applications of the second rule in
\eqref{I-17-32}; the boundary compression is that of
\eqref{I-17-36}.

Consequently every required expression
$O h\mathsf G_\square h O'g$ is a sum of a bounded number
of terms
\[
 \xi_n^{-i}\phi\,O_1\mathsf G_\square O_1'
                    \bigl(\xi_n^{-i'}\phi'g^{\mathrm{sh}}\bigr).
\]
Here $O\in\{I,\nabla,\Delta,\nabla\nabla\}$,
$O'\in\{I,\nabla^*,\nabla^*\nabla^*\}$,
$g=\chi_zf$ or $J$, and $\phi,\phi'$ are shifts of the
functions in \eqref{I-17-72}. The differential orders of
$O_1,O_1'$ have decreased by $i,i'$, respectively.
A backward output difference can occur for $O=\Delta$;
use $\nabla_\nu^*u=-(\nabla_\nu u)^{-\nu}$.
The input shift is at most $2\xi$.

If $O'=I$, the input remains $h\chi_zf$.
Otherwise $\phi'J^{\mathrm{sh}}$ is supported in
$\mathsf q(y',4\xi_{y'})$. Split this support, bond by bond,
among the $64$ cubes of radius $\xi_{y'}$ centered at
$y'+\xi_{y'}k$, $k\in\{\pm1,\pm3\}^3$.
The one term for which an input H\"older seminorm is needed
has $i'=0$: its input is $hJ$, left unsplit, and
\[
 \|hJ\|_{\beta,\xi_{y'}}
       \le\|J\|_{\beta,\xi_{y'}}+|J|\qquad(0<\beta<1).
\]
The shifted inputs of lower order need only their suprema
or $L^2$ norms.

\emph{Application of the local estimates.}
Apply \cref{I-17-38} with output radius $\xi_y$ and input
radius $\xi_z$ or $\xi_{y'}$; these are the two admissible
scales just identified. Where a cutoff is required, replace
$\zeta$ by $\zeta\phi$. For the first two types of
coefficients in \eqref{I-17-72}, the product rule gives
\[
 \|\zeta\phi\|_{\alpha,\xi_y}+|\zeta\phi|
      \le2(\|\zeta\|_{\alpha,\xi_y}+|\zeta|).
\]
Second derivatives of $h$ occur only in (G1) and (G3),
where no seminorm of that coefficient is needed.

The value and derivative bounds (R1), (R3) give (G1)
and the supremum parts of (G4), (G5); (R2) gives (G2)
and the seminorm in (G4). For (G5), apply (R3), (R4)
to the term $i=i'=0$, and (R1), (R2) to the lower-order
terms. The $L^2$ estimates (R5) give (G3), (G6).
For suprema involving backward differences, cover
$\mathsf q(y,4\xi_y)$ by the analogous $64$ output cubes.
Each displaced center has local distance at most $9L$
from the original one, so these covers cost only a constant
depending on $L$.

The powers of $\xi_n$ from the coefficients and local
estimates combine to the displayed output power of $\xi_y$,
up to a factor at most $L^2$. Convert the exponentials by
the metric comparison above, sum over the at most $16$
contributing cubes, and for (G1)--(G3) also over the input
cells $z$. This proves the stronger assertions with the
sum of local expressions on the left. The estimates for
$\mathsf G_0$ follow by the triangle inequality.
\end{proof}

\subsection{Proof of the global inverse estimates}
\label{I-17-75}

\begin{proof}[Proof of \cref{I-17-5}]
Return to the physical setting of the proposition, in which
$\xi_y\le1$. Fix the constants in the scalar, local-model,
gauge-comparison and remainder estimates before choosing
$\mathfrak M$. Decrease the remainder rate, if necessary, so
that $0<\kappa_{\mathrm H}\le\min\{1,\kappa_{\mathrm a}/2\}$,
and put
\begin{equation}\label{I-17-76}
 \kappa=\kappa_{\mathrm H}/2,\qquad
 \delta=\kappa/8=\kappa_{\mathrm H}/16.
\end{equation}
Choose $M_0=M_0(L)$ large enough to meet every preceding
lower bound on $\mathfrak M$, and so that $LM_0$ meets
every preceding lower bound on $R_1\mathfrak M$. Require also
\begin{equation}\label{eq:global-final-thresholds}
 M_0\ge\max\{2,2C_{\mathrm H}c_S(\kappa_{\mathrm H}/2)\},
 \qquad LM_0\ge\mathsf r_S(\kappa_{\mathrm H}/8).
\end{equation}
There are finitely many such conditions and all their constants
depend only on $L$; in particular $C_{\mathrm H}$ is independent
of $\mathfrak M$. Thus the choice is not circular.
For every $R_1\ge L$ and $\mathfrak M\ge M_0$ all conditions
hold. The last one also gives
$\delta R_1\mathfrak M\ge3\log L+1\ge\log L$.
Below write $d=d_{\mathbf B}$.

\emph{Convergence and absolute walk bounds.}
Set
\[
 E_{\square',\square}
       =K_{\square',\square}\mathsf G_\square h_\square,
 \qquad
 n_p(y,y')=\sum_{\square',\square}
                        N_p(E_{\square',\square};y,y'),
\]
and define $n_p^*(y,y')$ by replacing each $E$ by $E^*$.
The transfers are adjoint and the local inverses symmetric.
Hence $n_2^*(y,y')=n_2(y',y)$, and (H1), (H2) of
\cref{I-17-68} give
\[
 n_p(y,y'),\ n_p^*(y,y')
       \le C_{\mathrm H}\mathfrak M^{-1}
                         e^{-\kappa_{\mathrm H}d(y,y')}
       \qquad(p=2,\infty).
\]
By \eqref{I-13-122} and \eqref{eq:global-final-thresholds},
their weighted row sums at rate $\kappa$ are at most $1/2$.
Write $n_p^{(i)}$ and $n_p^{*(i)}$ for their matrix powers,
with identity matrix at $i=0$. Submultiplicativity gives
\begin{equation}\label{I-17-80}
 n_p^{(i)}(y,y'),\ n_p^{*(i)}(y,y')
             \le2^{-i}e^{-\kappa d(y,y')}.
\end{equation}

Since $\mathcal R=\sum E_{\square',\square}$,
$\|\mathcal R\|_{p,\kappa}$ and
$\|\mathcal R^*\|_{p,\kappa}$ are at most $1/2$.
The block calculus proves convergence of both Neumann series.
The identity $\Delta_a(\Omega_0)\mathsf G_0=I-\mathcal R$
then gives a right inverse on the finite-dimensional space
$\mathcal A$, and therefore
\[
 \mathsf G=\Delta_a(\Omega_0)^{-1}
          =\mathsf G_0(I-\mathcal R)^{-1}.
\]
The operator $\Delta_a(\Omega_0)$ is symmetric and nonnegative
by its defining form. Invertibility makes it positive definite,
and $\mathsf G$ is symmetric. Taking adjoints gives
$\mathsf G=(I-\mathcal R^*)^{-1}\mathsf G_0$.
Expanding each finite product gives exactly the walks in
\eqref{I-17-6}, and reversing the factors gives
$\mathsf G_\omega^*$. Convergence in the finite-dimensional
block norm implies operator-norm convergence.

To carry estimates through this expansion, fix an output cell
$y$ and a seminorm $\mathsf p$ such that
\[
 \sum_\square\mathsf p
       (h_\square\mathsf G_\square h_\square f)
       \le c\sum_z e^{-\kappa_{\mathrm a}d(y,z)}
                                      \|\chi_zf\|_p.
\]
Applying the block multiplication rule to the $i$ remainder
factors, then \eqref{I-17-80} and \eqref{I-17-30}, yields
\begin{equation}\label{I-17-82}
 \begin{split}
 \sum_{|\omega|=i}\mathsf p(\mathsf G_\omega f)
 &\le2^{-i}c\sum_{z,y'}
       e^{-\kappa d(y,z)-\kappa d(z,y')}\|\chi_{y'}f\|_p\\
 &\le2^{-i}c\,c_S(\kappa/2)
       \sum_{y'}e^{-\kappa d(y,y')/2}\|\chi_{y'}f\|_p.
 \end{split}
\end{equation}
All sums over walks at fixed $i$ are finite. Applying this
inequality to the normalized seminorms in (G1)--(G3) of
\cref{I-17-71} proves
the value, gradient and Laplacian bounds for undifferentiated
input, the three corresponding $L^2$ bounds, and
\begin{equation}\label{I-17-83}
 \|\zeta\nabla\mathsf GJ\|_{\alpha,\xi_y}
 \le C(\alpha)\xi_y(\|\zeta\|_{\alpha,\xi_y}+|\zeta|)
                     e^{-\kappa d(y,y')/2}|J|.
\end{equation}
The same estimates hold for sums of walk seminorms with a
factor $2^{-i}$, before summing over $i$.
For the versions with $\chi_y$, use the indicator $\zeta_y$
of $\widetilde\Delta(y)$, acting at initial points:
$\|\chi_yF\|\le\|\zeta_yF\|$.

\emph{Differentiated data without an output derivative.}
Let $J$ be supported in $\widetilde\Delta(y')$, and take
\[
 (\mathsf J,q,p)=(\nabla^*J,1,\infty),\quad
 (\nabla^*J,1,2),\quad
 (\nabla^*\nabla^*J,0,2).
\]
The last case uses $J\in\mathcal A^{(2)}$; the others use
$\mathcal A^{(1)}$. By (G4), (G6), the sum of local
$p$-norms in cell $z$ is at most
$C\xi_z^q e^{-\kappa_{\mathrm a}d(z,y')}\|J\|_p$.
Use the adjoint expansion and \eqref{I-17-80}, then transfer
$\xi_z^q$ to $\xi_y^q$ with loss $e^{\kappa d(y,z)/4}$.
The site sum gives
\begin{equation}\label{I-17-84}
 \begin{split}
 \sum_{|\omega|=i}\|\chi_y\mathsf G_\omega^*\mathsf J\|_p
 &\le2^{-i}CL\xi_y^q
      \sum_z e^{-3\kappa d(y,z)/4-\kappa d(z,y')}\|J\|_p\\
 &\le2^{-i}CLc_S(\kappa/4)\xi_y^q
                          e^{-\kappa d(y,y')/2}\|J\|_p.
 \end{split}
\end{equation}
This proves the remaining undifferentiated-output supremum
and $L^2$ estimates. For $\zeta$ supported in $\Delta(y)$,
$\|\zeta F\|\le|\zeta|\|\chi_yF\|$: every retained bond
whose initial point is in that cell is localized there.

For the H\"older seminorm use (G4) at $i=0$.
For $i\ge1$, separate the leftmost adjoint remainder block
from $\mathsf G_\omega^*$, apply (H3), and then apply
\eqref{I-17-84} to the remaining walk of length $i-1$.
Since $\mathfrak M^{-1}\le1/2$, this bounds the sum by
\[
 \begin{split}
 &2^{-i}C(\alpha)(\|\zeta\|_{\alpha,\xi_y}+|\zeta|)
       \sum_z e^{-\kappa_{\mathrm H}d(y,z)}
                  \xi_z e^{-\kappa d(z,y')/2}|J|\\
 &\hspace{1em}\le2^{-i}C(\alpha)\xi_y
      (\|\zeta\|_{\alpha,\xi_y}+|\zeta|)
                       e^{-\kappa d(y,y')/4}|J|.
 \end{split}
\]
Here the scale transfer costs $e^{\kappa d(y,z)/4}$,
and \eqref{I-17-30} at rate $\kappa/2$ sums the remainder.
Summing in $i$ gives
\begin{equation}\label{I-17-85}
 \|\zeta\mathsf G\nabla^*J\|_{\alpha,\xi_y}
 \le C(\alpha)\xi_y(\|\zeta\|_{\alpha,\xi_y}+|\zeta|)
                         e^{-\kappa d(y,y')/4}|J|.
\end{equation}

\emph{The mixed derivative.}
For $0<\beta<1$ put
$[J]_\infty=\|J\|_{\beta,\xi_{y'}}+|J|$ and
$[J]_2=\|J\|$. Use in \eqref{I-17-82} one of the three
normalized gradient seminorms
\[
 \xi_y^{-1}|\widetilde\chi_y\nabla F|,\qquad
 \frac{\|\zeta\nabla F\|_{\alpha,\xi_y}}
          {\xi_y(\|\zeta\|_{\alpha,\xi_y}+|\zeta|)},\qquad
 \frac{\|\zeta\nabla F\|}{\xi_y|\zeta|},
\]
with $(p,c)=(\infty,C)$, $(\infty,C(\alpha))$, $(2,C)$,
respectively; the zero multiplier is immediate.
For $i\ge1$, separate the rightmost forward remainder
block acting on $\nabla^*J$. Apply \eqref{I-17-82} to
the remaining $i-1$ factors and (H4) to that last block.
Convolution at rate $\kappa/2$ gives
\[
 \sum_{|\omega|=i}\mathsf p(\mathsf G_\omega\nabla^*J)
 \le2^{-i}cC(\beta)\xi_{y'}^{-1}
                       e^{-\kappa d(y,y')/4}[J]_p
 \le2^{-i}cC(\beta)\xi_y^{-1}
                       e^{-\kappa d(y,y')/8}[J]_p.
\]
For $p=2$ the constant is independent of $\beta$.
The last scale transfer uses $\delta R_1\mathfrak M\ge\log L$.
For $i=0$ use (G5), (G6).
Take $\beta=\varepsilon$ for the supremum bound,
$\beta=\alpha+\varepsilon$ for the H\"older bound, and
sum in $i$. After restoring the normalizing factors this
gives the mixed supremum and $L^2$ bounds, together with
\begin{equation}\label{I-17-86}
 \|\zeta\nabla\mathsf G\nabla^*J\|_{\alpha,\xi_y}
 \le C(\alpha,\varepsilon)
      (\|\zeta\|_{\alpha,\xi_y}+|\zeta|)
      e^{-\delta d(y,y')}
      (\|J\|_{\alpha+\varepsilon,\xi_{y'}}+|J|).
\end{equation}

\emph{The unscaled seminorms and coercivity.}
Since $\xi_y\le1$, \eqref{I-17-35} gives
\[
 \|F\|_\alpha\le\xi_y^{-\alpha}
                       (\|F\|_{\alpha,\xi_y}+2|F|).
\]
For $F=\zeta\nabla\mathsf GJ$ or
$\zeta\nabla\mathsf G\nabla^*J$, the supremum is controlled
by the corresponding estimate on $\mathsf q(y,2\xi_y)$
already proved. For $F=\zeta\mathsf G\nabla^*J$, a bond
contributing to this product is localized in a cell centered
at some $y^\star$ satisfying \eqref{eq:nearby-cell-comparison}:
\[
 d(y,y^\star)\le9L,\qquad
 L^{-1}\xi_y\le\xi_{y^\star}\le L\xi_y.
\]
There are at most $(2L(L+2))^3$ such cells: their disjoint
interiors have side at least $\xi_y/L$ and lie in
$\mathsf q(y,(L+2)\xi_y)$. Applying \eqref{I-17-84}
to them bounds the supremum by
$C|\zeta|\xi_y e^{-\kappa d(y,y')/2}|J|$.
The same finite sum bounds the sums of walk suprema.
Combining these bounds with \eqref{I-17-83},
\eqref{I-17-85} and \eqref{I-17-86} proves the asserted
ordinary H\"older estimates, with their factors
$\xi_y^{1-\alpha}$ and $\xi_y^{-\alpha}$.

The $L^2$ estimate for $\mathsf G$ and $\xi_y\le1$ give
$N_2(\mathsf G;y,y')\le Ce^{-\kappa d(y,y')/2}$.
This symmetric bound has uniformly bounded row sums by
\eqref{I-13-122}. Thus \cref{I-17-29}(g) gives a uniform
bound on $\|\mathsf G\|_{2\to2}$; its reciprocal is a
positive lower bound $\gamma$ for $\Delta_a(\Omega_0)$.

Every term estimate above has the factor $2^{-i}$, and
every decay rate is at least $\delta$.
Choose the constants in the proposition to be twice the
largest relevant term constants. This choice bounds both
the sum over $i$ and each sum over walks of length $i$.
The forward and adjoint orientations are exactly those
used in the three preceding parts of the proof. This proves
all assertions of the proposition.
\end{proof}

\section{Coarse covariance and the minimizing lift}
\label{sec:coarse-covariance}

The inverse in \cref{sec:global-dirichlet} acts on the fine
field. Its coarse covariance determines the minimizing lift
of prescribed coarse averages. We estimate this covariance and its
inverse.
The additional ingredient is a right inverse of the averaging map
with a uniform energy bound. It gives a spectral gap after
normalizing by the local scale; exponential conjugation then gives
decay of the inverse covariance.

Keep the real, three-dimensional, level-$t$ setting of
\cref{sec:global-dirichlet}, with $L\ge13$ odd and $a=1$.
Here $\mathfrak b_0$ is either the inner frame or the star frame.
Write $\Gamma=\bigsqcup_{n=0}^t\Gamma_n$ for the rows of $Q$,
put $\xi_c=\xi_n$ for $c\in\Gamma_n$, and give row fields the
inner product
\[
 \langle B,B'\rangle=\sum_{c\in\Gamma}\xi_c^3B(c)B'(c).
\]
The adjoint $Q^*$ uses this product and the fine product of weight
$\xi^3$. Let $\mathsf m$ multiply a row field by $\xi_c$.
Thus the averaging term of \eqref{I-16-64} is
$Q^*\mathsf m^{-2}Q$. Set
\begin{equation}\label{I-18-21}
 \mathsf Z=Q\mathsf GQ^*,\qquad
 \mathsf H=\mathsf GQ^*\mathsf Z^{-1}.
\end{equation}
These are the operators of \cref{I-11-33}(v) at $a=1$:
the row sets and centered averages coincide with
\eqref{I-11-20}, and their inner products and Dirichlet
projections agree. In particular, $\mathsf HB$ is the unique
curl-energy minimizer subject to $QA=B$ and
$\mathsf R\partial_0^*A=0$. With only $QA=B$, the minimizers
are $\mathsf HB+\partial\ker Q'$.

For $c\in\Gamma_n$, let $y(c)=c_-$ if $c_-\in\Omega_n$,
and $y(c)=c_+$ otherwise. Then $y(c)\in\Lambda_n$.
Use $d_{\mathbf B}(y,c)=d_{\mathbf B}(y,y(c))$ and
$d_{\mathbf B}(c,c')=d_{\mathbf B}(y(c),y(c'))$.
For an operator $K$ on row fields, distinguish its matrix in the
orthonormal basis $e_c=\xi_c^{-3/2}\delta_c$ from its counting
kernel:
\begin{equation}\label{I-18-22}
 [K]_{cc'}=\langle e_c,Ke_{c'}\rangle,\qquad
 K(c,c')=(K\delta_{c'})(c)
       =\left(\frac{\xi_{c'}}{\xi_c}\right)^{3/2}[K]_{cc'}.
\end{equation}
\begin{proposition}\label{I-18-76}
There are positive constants $M_1,\gamma,\theta,C$, depending
only on $L$, such that the following holds for every separated
sequence in this setting with $\mathfrak M\ge M_1$ and every
positive power $R_1\ge L$ of $L$. The coarse covariance
$\mathsf Z$ is symmetric and positive definite, and
\begin{equation}\label{I-18-55}
 \gamma\|\mathsf mB\|^2
       \le\langle B,\mathsf ZB\rangle
       \le\|\mathsf mB\|^2
       \qquad(B\in L^2(\Gamma)).
\end{equation}
Its inverse satisfies
\begin{equation}\label{I-18-77}
 \begin{split}
 |[\mathsf Z^{-1}]_{cc'}|
   &\le2\gamma^{-1}\xi_c^{-1}\xi_{c'}^{-1}
                 e^{-\theta d_{\mathbf B}(c,c')},\\
 |\mathsf Z^{-1}(c,c')|
   &\le C\xi_c^{-2}e^{-\theta d_{\mathbf B}(c,c')/2}.
 \end{split}
\end{equation}
The constants are independent of the torus, $K,m,t$, the regions,
their number of layer changes, the frame, and $\mathfrak M,R_1$.
\end{proposition}

\subsection{A quantitative lift of the coarse data}
\label{sec:coarse-covariance-gap}

We first construct a form for each row. Its averages at that level
are exact; its only possible error is on the next level, at a row
crossing the boundary of that level. This error can be removed in
one step.

\begin{proposition}\label{I-18-42}\label{I-18-53}
For the separated sequence and bond frame fixed above, with
$\mathfrak M,R_1\ge L$, there are a numerical constant $C_E$ and
a linear map $\mathsf E$ from the coarse row space to $\mathcal A$
such that
\begin{equation}\label{I-18-43}
 Q\mathsf E=I,\qquad
 \|\nabla\mathsf EB\|^2
       \le C_E\|\mathsf m^{-1}B\|^2
       =C_E\sum_{n=0}^t\sum_{c\in\Gamma_n}\xi_n|B(c)|^2.
\end{equation}
Consequently the covariance bound \eqref{I-18-55} holds with
$\gamma=(C_E+1)^{-1}$.
\end{proposition}

\begin{proof}
\emph{The averaging weights.}
All one-dimensional weights and profiles below have domain $\mathbb Z$.
Use $Q_n$ for the average \eqref{I-16-63} on all positive
level-$n$ bonds, applied to the zero extension of $A$.
For $c=(y,\mu)$ put
$N_n=L^n$ and $h_n=(N_n-1)/2$. For any odd integer $N$, write
$h=(N-1)/2$ and define
\[
 \tau_N(s)=\#\{i\in\mathbb Z:0\le i<N,\ |s-i|\le h\}.
\]
Counting the occurrences of each bare bond in
\eqref{I-16-63} gives
\begin{equation}\label{I-18-15}
 \begin{split}
 (Q_nA)(c)&=\sum_x w_c(x)A_\mu(x),\\
 w_c(y+\xi k)&=N_n^{-4}\tau_{N_n}(k_\mu),
 \quad -h_n\le k_\mu\le3h_n+1,\quad
             |k_\nu|\le h_n\ (\nu\ne\mu),
 \end{split}
\end{equation}
with $w_c=0$ elsewhere. These sites are the two endpoint blocks
$B^n(c_-)\cup B^n(c_+)$ and are distinct on the torus: each
coordinate range has length less than $2\xi_n$, whereas the
period is at least $2\xi_n$. There are $N_n^4$ terms before
collecting equal sites, and $0\le\tau_{N_n}\le N_n$. Hence
\[
 0\le w_c\le N_n^{-3},\qquad \sum_xw_c(x)=1.
\]
Moreover, $\tau_N$ vanishes outside $[-h,3h]$, so every
contributing bare bond has both endpoints in those two blocks.

Let $B_{n+1}(z)$ denote the level-$n$ sites in the centered
block of side $\xi_{n+1}$ about $z$. Splitting each such block
into level-$n$ blocks and writing a path index as
$i=i_1N_n+i_0$, with $0\le i_1<L$ and $0\le i_0<N_n$, gives
\begin{equation}\label{I-18-18}
 (Q_{n+1}A)(c')
     =\sum_{c''}w^{(1)}_{c'}(c'')(Q_nA)(c''),\qquad
 0\le w^{(1)}_{c'}\le L^{-3},\qquad
 \sum_{c''}w^{(1)}_{c'}(c'')=1.
\end{equation}
Indeed, after this splitting the average is
$L^{-4}\sum_{z\in B_{n+1}(c'_-)}\sum_{i_1=0}^{L-1}
(Q_nA)_\mu(z+i_1\xi_ne_\mu)$; the same count with $N=L$
proves the assertions about the weights. A nonzero weight
requires that $c''$ be parallel to $c'$ and that both endpoints
of $c''$ lie in $B_{n+1}(c'_-)\cup B_{n+1}(c'_+)$.

\emph{Forms with one prescribed average.}
For odd $N\ge3$, define
\[
 \begin{aligned}
 u(s)&=\frac6{(N+1)(N+2)}
       \left(\frac{(N+1)^2}{4}-s^2\right)
                 &&(|s|\le h+1),\\
 u(s)&=0 &&(|s|>h+1),\\
 v^\pm(s)&=u(s)(1+\beta_\pm s),\qquad
 \beta_+=\frac{10}{N+3},\qquad
 \beta_-=-\frac{10(N+1)}{(N-1)(N+3)}.
 \end{aligned}
\]
The profiles vanish at $s=\pm(h+1)$, so their supports are
contained in $[-h,h]$. The sums of the second and fourth powers
of the integers from $-h$ to $h$ give
\[
 \sum_su(s)=N,\qquad \sum_s s u(s)=0,\qquad
 \sum_s s^2u(s)=\frac{N(N-1)(N+3)}{20}.
\]
Consequently
\[
 \sum_sv^\pm(s)=N,\qquad
 \sum_s s v^+(s)=\frac{N(N-1)}2,\qquad
 \sum_s s v^-(s)=-\frac{N(N+1)}2.
\]
For $|s|\le h$ the triangular weights are
$\tau_N(s)=s+(N+1)/2$ and
$\tau_N(s+N)=(N-1)/2-s$. Thus
\begin{equation}\label{I-18-44}
 \begin{pmatrix}
  \sum_s\tau_N(s)v^+(s)&\sum_s\tau_N(s)v^-(s)\\
  \sum_s\tau_N(s+N)v^+(s)&\sum_s\tau_N(s+N)v^-(s)
 \end{pmatrix}=N^2I_2.
\end{equation}
The formulas also give, with a numerical constant $C$,
\[
 |u|+|v^\pm|\le C,\qquad
 |u(s+1)-u(s)|+|v^\pm(s+1)-v^\pm(s)|\le\frac CN.
\]
Indeed, $|\beta_\pm|\le C/N$ and
$|\beta_\pm|(h+1)\le C$; the difference of $u$ is bounded by
$C/N$, and
$v(s+1)-v(s)=(u(s+1)-u(s))(1+\beta(s+1))+\beta u(s)$.
The zero values at $\pm(h+1)$ make these bounds valid also
across the support boundary.

At level $n\ge1$, take $N=N_n$ and let
$\varphi^\pm_{y,\mu}$ have only its $\mu$-component nonzero,
given by
\[
 (\varphi^\pm_{y,\mu})_\mu(y+\xi k)
      =v^\pm(k_\mu)\prod_{\nu\ne\mu}u(k_\nu)
          \quad(|k|_\infty\le h_n),
\]
and zero off $B^n(y)$. Only the two level-$n$ bonds
$c^+=(y,\mu)$ and $c^-=(y-\xi_ne_\mu,\mu)$ of direction
$\mu$ have averaging weights meeting this block. They are
distinct even when the period is $2\xi_n$.
By \eqref{I-18-15}, \eqref{I-18-44} and $\sum u=N_n$,
\begin{equation}\label{I-18-45}
 Q_n\varphi^+_{y,\mu}=\delta_{c^+},\qquad
 Q_n\varphi^-_{y,\mu}=\delta_{c^-}
\end{equation}
on all positive level-$n$ bonds.

For $c\in\Gamma_n$, $n\ge1$, use the chosen endpoint $y(c)$:
take $\varphi_c=\varphi^+_{c_-,\mu}$ if $y(c)=c_-$ and
$\varphi_c=\varphi^-_{c_+,\mu}$ otherwise. For
$c\in\Gamma_0$ take $\varphi_c=\delta_c$.
These forms belong to $\mathcal A$: for $n\ge1$ their nonzero
bonds start in $B^n(y(c))\subset\Omega_n\subset\Omega_1$,
and separation places their other endpoints in $\Omega_0$;
for $n=0$ this follows from $c\in\mathfrak b_0$.

There is a numerical constant $C_\varphi$ such that
\begin{equation}\label{I-18-46}
 \|\nabla\varphi_c\|^2\le C_\varphi\xi_c.
\end{equation}
To check this including the zero extension, let $n\ge1$.
For each direction $\nu$, the derivative can be nonzero at
only $(N_n+1)N_n^2\le\frac43N_n^3$ sites. The extra layer
of sites is outside $B^n(y(c))$, since
$N_n+1<2N_n$ does not wrap around the torus.
The product formula and the profile bounds give
\[
 |\nabla_\nu(\varphi_c)_\mu|\le C\xi_n^{-1}
                   \qquad(1\le\nu\le3).
\]
Multiplying their squares by the site weight $\xi^3$ and
the number of sites proves \eqref{I-18-46}, uniformly in $n$.
At level zero the energy of $\delta_c$ is $6\xi$, which is
covered by increasing the same constant.

\emph{Removing the error between adjacent levels.}
Call a row $c\in\Gamma_n$ crossing if exactly one endpoint
lies in $\Omega_n$. The block geometry of \cref{I-11-27}
shows that the averages of $\varphi_c$, for $c\in\Gamma_n$,
vanish at every other level except possibly $n+1$.
For completeness, if $1\le n'<n$, the averaging paths of a
row in $\Gamma_{n'}$ have their initial points outside
$\Omega_{n'+1}$, whereas the support of $\varphi_c$ starts
inside $\Omega_n$. For $n'=0<n$, a row has no endpoint in
$\Omega_1$. If $n'\ge n+2$, its averaging paths lie inside
$\Omega_{n'-1}\subset\Omega_{n+1}$, whereas
$\varphi_c$ has no initial point there. These assertions
also cover $n=0$ because $c$ has no endpoint in $\Omega_1$.
At its own level the average is $\delta_c$ by
\eqref{I-18-45}, or by $Q_0=I$.

For $c'\in\Gamma_{n+1}$, \eqref{I-18-18} gives
\begin{equation}\label{I-18-47}
 (Q\varphi_c)(c')=w^{(1)}_{c'}(c)=:\mathsf N(c',c).
\end{equation}
If this entry is nonzero, both endpoints of $c$ lie in the
two endpoint blocks of $c'$, and neither lies in
$\Omega_{n+1}$. Hence they lie in the block of the outer
endpoint of $c'$: $c'$ is crossing. By
\cref{I-11-27}(ii), that block consists of level-$n$ sites
in $\Omega_n\setminus\Omega_{n+1}$, so $c$ is not crossing.

Set $\mathsf N(c',c)=0$ for all other pairs and put
\[
 \mathsf E_0B=\sum_c B(c)\varphi_c.
\]
Then $Q\mathsf E_0=I+\mathsf N$. The matrix $\mathsf N$
uses only noncrossing input rows and has only crossing
output rows, so $\mathsf N^2=0$. Therefore
\[
 \mathsf E=\mathsf E_0(I-\mathsf N),\qquad Q\mathsf E=I.
\]
The row sums of $\mathsf N$
are at most one and its entries at most $L^{-3}$.
Thus, for $c'\in\Gamma_{n+1}$,
\[
 \xi_{n+1}|(\mathsf NB)(c')|^2
 \le \xi_{n+1}\sum_c\mathsf N(c',c)|B(c)|^2
 \le L^{-2}\sum_{c:\,\mathsf N(c',c)\ne0}
                                      \xi_n|B(c)|^2.
\]
Each $c$ occurs for at most two rows $c'$: its two endpoints
lie in one outer block, whose center can be an endpoint of
only two bonds parallel to $c$. Hence
\[
 \|\mathsf m^{-1}\mathsf NB\|^2
       \le2L^{-2}\|\mathsf m^{-1}B\|^2.
\]
Thus $I-\mathsf N$ is uniformly bounded in this weighted norm.

Put $\beta=(I-\mathsf N)B$. At any bare site and in any
one component, at most two of the forms $\varphi_c$ are
nonzero, since the blocks partition the bare sites of
$\Omega_0$; outside that region only a level-zero form
can contribute. A first difference therefore has at most
four nonzero summands. By Cauchy--Schwarz and
\eqref{I-18-46},
\[
 \|\nabla\mathsf EB\|^2
 \le4\sum_c|\beta(c)|^2\|\nabla\varphi_c\|^2
 \le4C_\varphi\|\mathsf m^{-1}\beta\|^2
 \le C_E\|\mathsf m^{-1}B\|^2
\]
for a numerical constant $C_E$, proving \eqref{I-18-43}.

\emph{The covariance gap.}
The quadratic form of $\Delta_a$ and the ambient identity
\eqref{I-11-18} give
\[
 \|\mathsf m^{-1}QA\|^2
 \le \langle A,\Delta_a A\rangle
 \le\|\nabla A\|^2+\|\mathsf m^{-1}QA\|^2,
                  \qquad A\in\mathcal A.
\]
Indeed, the gauge term is
$\|\mathsf R\partial_0^*A\|^2\le\|\partial^*A\|^2$,
and $\|\partial A\|^2+\|\partial^*A\|^2=\|\nabla A\|^2$.
Completing the square about $\mathsf GQ^*B$ yields
\[
 \langle B,\mathsf ZB\rangle=\max_{A\in\mathcal A}
          \{2\langle QA,B\rangle-\langle A,\Delta_a A\rangle\}.
\]
For $A=s\mathsf E\mathsf m^2B$, the expression in braces
is at least
$[2s-(C_E+1)s^2]\|\mathsf mB\|^2$.
Taking $s=\gamma=(C_E+1)^{-1}$ proves the lower bound
in \eqref{I-18-55}. For every $A$ the same expression is at most
\[
 2\langle QA,B\rangle-\|\mathsf m^{-1}QA\|^2
 =\|\mathsf mB\|^2-
                   \|\mathsf m^{-1}QA-\mathsf mB\|^2
 \le\|\mathsf mB\|^2,
\]
which proves the upper bound.
\end{proof}

\subsection{Decay of the coarse inverse}
\label{sec:coarse-covariance-decay}

\begin{proof}[Proof of the decay bounds in \cref{I-18-76}]
Write $d=d_{\mathbf B}$ and let $\delta_3>0$ be the rate in
\cref{I-17-5}. The averaging weights in \eqref{I-18-15}
localize each column of $Q^*$ in a bounded number of cells.
We first record how their positions and scales compare.

For $c\in\Gamma_n$, $n\ge1$, the bare sites of
$\Delta^n(c_-)\cup\Delta^n(c_+)$ are, by
\cref{I-11-27}\textup{(ii),(iii)},
partitioned into two level-$n$ blocks if both endpoints belong
to $\Omega_n$, and otherwise into one level-$n$ block and the
$L^3$ level-$(n-1)$ blocks of the outer endpoint.
Let $\mathcal C(c)$ be their centers. For $n=0$ put
$\mathcal C(c)=\{y(c)\}$. Thus
\[
 |\mathcal C(c)|\le1+L^3,\qquad
 \xi_z\le\xi_c,\qquad d(z,y(c))\le4L
       \quad(z\in\mathcal C(c)).
\]
For the distance bound, join $y(c)$ to the midpoint of $c$
inside its level-$n$ cell. This segment has weighted length
$1/2$. If the other endpoint also lies in $\Omega_n$, the
second half has the same length. Otherwise, a center
$z\in B(c_+)$ or $B(c_-)$ in the outer block can be joined
to that midpoint by a segment of $\ell^1$ length at most
$2\xi_n$. Apart from the midpoint, this segment lies in the
interior of the outer level-$n$ cell, which belongs only to
layer $n-1$. Its weighted length is therefore at most $2L$.
The endpoints have zero length, so the same estimates apply
with the closed-layer convention. At level zero the distance
vanishes.

Every row with position $y(c)=z$ is one of the six positive
bonds incident to $z$ at its level. Hence
\cref{lem:global-cell-geometry}\textup{(d)} gives
\begin{equation}\label{eq:coarse-row-sum}
 \sup_{y\in\Lambda_*}\sum_{c\in\Gamma}e^{-\kappa d(y,c)}
 \le6c_S(\kappa)
 \quad\text{if }R_1\mathfrak M\ge\mathsf r_S(\kappa).
\end{equation}
In particular this holds with $y=y(c')$ for any row $c'$.

For $c\in\Gamma_n$ of direction $\mu$, put
$J_c=Q^*\delta_c$ and $A^c=\mathsf GJ_c$.
The adjoint is taken for the bare weight $\xi^3$ and the
coarse weight $\xi_c^3$. Thus \eqref{I-18-15} gives
\[
 \langle A,J_c\rangle
   =\xi_c^3\sum_xw_c(x)A_\mu(x),\qquad
 (J_c)_\mu(x)=L^{3n}w_c(x)
\]
on the admitted bare bonds; all other components and all bonds
outside the frame have value zero.
Since $0\le w_c\le L^{-3n}$, we have $0\le J_c\le1$.
For $n\ge1$, every contributing bond is inner to $\Omega_0$
by \cref{I-11-27}\textup{(iii)}, so its localization point
is its initial point. Restricting $J_c$ to the cells just
described writes it as a sum of at most $1+L^3$ fields,
each of supremum at most one and supported in a cell
$\Delta(z)$ with $z\in\mathcal C(c)$.
For $n=0$, $J_c=\delta_c$ and the same assertion holds
with one cell.

Apply \eqref{eq:global-dirichlet-pointwise} and
\eqref{eq:global-dirichlet-holder} to these fields.
Since $d(y,z)\ge d(y,c)-4L$, there are constants
$C_1=C_1(L)$ and $C_{1,\alpha}=C_{1,\alpha}(L,\alpha)$
such that
\begin{equation}\label{I-18-73}
 \begin{aligned}
 |A^c(b)|&\le C_1\xi_y^2e^{-\delta_3d(y,c)},\\
 |(\partial_\nu A^c_\mu)(x)|
     &\le C_1\xi_ye^{-\delta_3d(y,c)},\\
 \|\zeta\nabla A^c\|_\alpha
     &\le C_{1,\alpha}\xi_y^{1-\alpha}
       (\|\zeta\|_{\alpha,\xi_y}+|\zeta|)
                       e^{-\delta_3d(y,c)}.
 \end{aligned}
\end{equation}
Here $b$ is localized in $\Delta(y)$, $x$ is a bare site
of $\Delta(y)$, and $\zeta$ vanishes outside
$\widetilde\Delta(y)$, as in the statement.

The weights $w_c$ are nonnegative and sum to one.
Consequently $(QA)(c)$ is bounded by the largest value of
$|A|$ on the contributing bonds. Use \eqref{I-18-73},
$\xi_z\le\xi_c$ and $d(z,c')\ge d(c,c')-4L$
in the identity $\mathsf Z(c,c')=(QA^{c'})(c)$.
This gives a constant $C_{\mathsf Z}=C_{\mathsf Z}(L)$ with
\[
 |\mathsf Z(c,c')|
      \le C_{\mathsf Z}\xi_c^2e^{-\delta_3d(c,c')}.
\]
Normalize the covariance in the orthonormal basis of
\eqref{I-18-22}:
\[
 S=[\mathsf m^{-1}\mathsf Z\mathsf m^{-1}],
 \qquad S_{cc'}=\xi_c^{-1}\xi_{c'}^{-1}[\mathsf Z]_{cc'}.
\]
The gap \eqref{I-18-55} says $\gamma I\le S\le I$.
Moreover, \eqref{I-18-22} and the last bound give
\[
 |S_{cc'}|
  \le C_{\mathsf Z}
       (\xi_c/\xi_{c'})^{5/2}e^{-\delta_3d(c,c')}.
\]
The matrix $S$ is symmetric, so the same estimate holds
with the reciprocal scale ratio. Taking the smaller of
the two bounds yields
\begin{equation}\label{I-18-75}
 \gamma I\le S\le I,\qquad
 |S_{cc'}|\le C_{\mathsf Z}e^{-\delta_3d(c,c')}.
\end{equation}

We apply the conjugation calculation following
\eqref{eq:direct-gram-smallness} to the row distance $d$.
The summation constant here is
$B_{\mathrm{row}}=6c_S(\delta_3/2)$. Choose
\begin{equation}\label{eq:coarse-conjugation-rate}
 \theta=\min\left\{\frac{\delta_3}{4},
       \frac{\gamma\delta_3}
            {4C_{\mathsf Z}B_{\mathrm{row}}}\right\}>0.
\end{equation}
Choose the threshold $M_1$ in the statement at least as
large as $M_0$ from \cref{I-17-5}, and so that,
for every $\mathfrak M\ge M_1$ and $R_1\ge L$,
\begin{equation}\label{eq:coarse-decay-thresholds}
 R_1\mathfrak M\ge\mathsf r_S(\delta_3/4),
 \qquad \theta R_1\mathfrak M\ge\log L.
\end{equation}
This requires only a larger number depending on $L$:
it suffices to impose these inequalities with $R_1\mathfrak M$
replaced by $LM_1$.

Fix $c_0\in\Gamma$, set $\psi(c)=d(c,c_0)$, and put
$S_\theta=e^{\theta\psi}Se^{-\theta\psi}$.
The triangle inequality gives
$|\psi(c)-\psi(c')|\le d(c,c')$. By the cited calculation,
\eqref{I-18-75} and the actual row sum \eqref{eq:coarse-row-sum},
\[
 \|S_\theta-S\|
 \le\frac{4C_{\mathsf Z}\theta B_{\mathrm{row}}}
          {e\delta_3}
 \le\frac{\gamma}{e}<\frac{\gamma}{2}.
\]
Thus $\|S_\theta^{-1}\|\le2/\gamma$, and undoing the
conjugation gives
$|(S^{-1})_{cc'}|\le2\gamma^{-1}e^{-\theta d(c,c')}$.

Undoing the normalization gives
$[\mathsf Z^{-1}]_{cc'}
 =\xi_c^{-1}\xi_{c'}^{-1}(S^{-1})_{cc'}$,
which is the first bound in \eqref{I-18-77}.
By \eqref{I-18-22}, the counting kernel consequently satisfies
\[
 |\mathsf Z^{-1}(c,c')|
 \le2\gamma^{-1}\xi_c^{-2}
       (\xi_{c'}/\xi_c)^{1/2}e^{-\theta d(c,c')}.
\]
The scale-ratio estimate \eqref{eq:global-scale-ratio},
with exponent $1/2$, and the second condition in
\eqref{eq:coarse-decay-thresholds} imply
\[
 (\xi_{c'}/\xi_c)^{1/2}e^{-(\theta/2)d(c,c')}
       \le L^{1/2}.
\]
This proves the second bound in \eqref{I-18-77}, with
constant $2\gamma^{-1}L^{1/2}$.
\end{proof}

\section{The constrained system at the trivial configuration}
\label{sec:trivial-constrained-system}

The scalar inverse, the fine-field inverse and the coarse covariance
solve the linear constrained system at the trivial configuration.
We prove a bound for this system which includes its gauge-fixing
variables and multipliers. The forcing may contain a divergence;
the estimate therefore uses the regularity of the inverses as well
as their pointwise decay.

Keep the three-dimensional physical torus and level-$t$ coordinates
of \cref{sec:global-dirichlet}: $L\ge13$ is odd,
$K,m\ge1$, $0\le t\le K$, $\xi=L^{-t}$ and $\xi_n=L^n\xi$.
We use the inner bond frame and take the whole torus as ambient
region. Let $\mathfrak m$ be a power of $L$ and $\mathsf g\ge L^2$.
Assume that $\Omega_0\supset\cdots\supset\Omega_t$ are unions of
the cubes of $\pi[\mathfrak m\xi_n]$ at their respective levels and
\[
 \operatorname{dist}_\infty(\Omega_n^c,\Omega_{n+1})
       >\mathsf g\xi_n\qquad(0\le n<t).
\]
Empty regions are allowed; $\Omega_n=\emptyset$ when
$\mathfrak m>L^{m+K-n}$, and $\Omega_{t+1}=\emptyset$.
The cells $\Delta(y)$, tops $\Lambda_*$, row set
$\Gamma=\bigsqcup_n\Gamma_n$ and distance $d_{\mathbf B}$ are
those of \cref{sec:global-dirichlet,sec:coarse-covariance}.
In particular, the rows are positive and have no endpoint in
the next region. Write $d=d_{\mathbf B}$.

Let $\mathcal S,\mathcal A,\mathcal Y,\mathcal R$ be the spaces of
$\mathfrak g^c$-valued functions on the bare sites of $\Omega_0$,
its positive inner bare bonds, $\Lambda_*$ and $\Gamma$,
respectively. Here $\mathfrak g^c$ is the space of traceless
complex $2\times2$ matrices, with matrix operator norm $|\cdot|$.
Their pairings are the sums of $\frac12\tr(XY)$ with weights
$\xi^3,\xi^3,\xi_y^3,\xi_c^3$, respectively. They are complex
bilinear. All transposes $T^{\mathsf T}$ below use these pairings.

Apply the real operators of the preceding sections to each Lie
algebra coordinate, retaining their names
$Q,Q',\mathsf G',\mathsf C,\mathsf P,\mathsf R,\mathsf G,
\mathsf Z,\mathsf H$ and their weights $a=1$.
Set
\[
 D=\partial_0,\qquad
 \mathcal L=\chi_{\mathcal A}\partial^*\partial\chi_{\mathcal A},
 \qquad
 \mathsf L_1=\mathcal L+DD^{\mathsf T}
                    +Q^{\mathsf T}\mathsf m^{-2}Q.
\]
Here $\mathcal L$ is the ambient curl operator compressed to
inner bonds, and $\mathsf m$ multiplies a row field by $\xi_c$.
Scalar fields and bond fields are extended by zero; $\nabla_\nu$
is $\xi^{-1}(T_\nu-I)$ on these extensions, with its output
restricted to the corresponding inner index set. Thus
$(D\lambda)(b)=\xi^{-1}(\lambda(b_+)-\lambda(b_-))$ and
$D^{\mathsf T}A=\sum_\mu\nabla_\mu^{\mathsf T}A_\mu$.
The scalar Dirichlet operator in $\Delta'_a$ includes all six
neighbors, including the zero-extension terms at the boundary.

On
$\mathcal X=\mathcal A\times\mathcal S^2\times\mathcal Y
                 \times\mathcal S^2\times\mathcal Y\times\mathcal R$,
write $\Psi=(A;v,w,\sigma;p,q,r;\omega)$ and define
\begin{equation}\label{I-63-64}
 \mathbb K_0^{[5]}\Psi=
 \begin{pmatrix}
 \mathsf L_1A-D(w+p)-Q^{\mathsf T}\omega\\
 \Delta'_av+w-D^{\mathsf T}A\\
 \Delta'_aw-Q'^{\mathsf T}\sigma\\
 Q'v\\
 \Delta'_ap+Q'^{\mathsf T}r\\
 p+\Delta'_aq\\
 -Q'q\\
 QA
 \end{pmatrix}.
\end{equation}
Let $\mathbb K_0^{[4]}$ be the coordinate compression obtained by
removing $\omega$ and the last equation. The other two compressions
are
\[
 \mathbb K_0^{[2]}(v,w,\sigma)
   =(\Delta'_av+w,\ \Delta'_aw-Q'^{\mathsf T}\sigma,\ Q'v),
 \qquad
 \mathbb K_0^{[1]}w=\Delta'_aw.
\]
Denote their spaces by $\mathcal X^{[i]}$, with
$\mathcal X^{[5]}=\mathcal X$.

A bare site's position is the top of its cell; the position of an
inner bond, or of a bond with an additional direction index, is
that of its initial point. A top is its own position, and a row
has the position $y(c)$ of \cref{sec:coarse-covariance}.
For a function on any of these index sets put
\[
 |f|_y=\max_{\mathrm{position}(i)=y}|f(i)|,\qquad
 |f|_{s;y}=\xi_y^{-s}|f|_y,\qquad \max\emptyset=0.
\]
The solution seminorm is
\begin{equation}\label{I-64-3}
 \begin{split}
 |\Psi|_y=\max\{&
 |A|_{-1;y},|\nabla A|_{-2;y},
 |v|_{0;y},|Dv|_{-1;y},|w|_{-2;y},|Dw|_{-3;y},\\
 &|\sigma|_{-4;y},|p|_{-2;y},|q|_{0;y},|Dq|_{-1;y},
 |r|_{-4;y},|\omega|_{-3;y}\}.
 \end{split}
\end{equation}
For $i=4,5$ the forcing space is
$\mathcal F^{[i]}=\mathcal X^{[i]}\times\mathcal A$.
Writing its elements as $g=(f_1,\ldots,f_8;X_5)$ when $i=5$, set
\[
 \iota g=(f_1,f_2,f_3,f_4,f_5+D^{\mathsf T}X_5,f_6,f_7,f_8)
\]
and
\begin{equation}\label{I-64-4}
 \begin{split}
 |g|_y=\max\{&
 |f_1|_{-3;y},|f_2|_{-2;y},|f_3|_{-4;y},|f_4|_{0;y},
 |f_5|_{-4;y},\\
 &|X_5|_{-3;y},|f_6|_{-2;y},|f_7|_{0;y},|f_8|_{-1;y}\}.
 \end{split}
\end{equation}
For $i=4$ omit $\omega,f_8$ from these formulas.
For $i=2$ the unknown and forcing are $(v,w,\sigma)$ and
$(f_2,f_3,f_4)$; for $i=1$ they are $w$ and $f_3$.
In these two cases $\mathcal F^{[i]}=\mathcal X^{[i]}$,
$\iota=I$, and the seminorms retain precisely the entries involving
the displayed components.

\begin{proposition}\label{I-64-24}
There are a positive power $\mathfrak M_\dagger$ of $L$ and
positive constants $\delta_{\mathrm{cs}},B_{\mathrm{cs}}$,
depending only on $L$, with the following property.
For every sequence above with
\[
 \mathfrak m\ge\mathfrak M_\dagger,\qquad
 \mathsf g\ge L\mathfrak M_\dagger,
\]
all four operators $\mathbb K_0^{[i]}$, $i\in\{1,2,4,5\}$,
are invertible. Their inverses
$\mathbb E_0^{[i]}=(\mathbb K_0^{[i]})^{-1}$ satisfy
\begin{equation}\label{eq:trivial-system-decay}
 |\mathbb E_0^{[i]}\iota g|_y
 \le B_{\mathrm{cs}}\sup_{y'\in\Lambda_*}
       e^{-\delta_{\mathrm{cs}}d(y,y')}|g|_{y'}
 \qquad(g\in\mathcal F^{[i]},\ y\in\Lambda_*).
\end{equation}
The constants are uniform in $K,m,t$, the lattice spacing,
the torus and the sequence, including its number of layers.
For $\Omega_0=\emptyset$ all spaces are zero and the assertion
uses the unique map on the zero space.
In particular, the result applies to the separated sequences of
\cref{sec:global-dirichlet} whenever
$\mathfrak M\ge\mathfrak M_\dagger$, for every power $R_1\ge L$,
by taking $\mathfrak m=\mathfrak M$ and
$\mathsf g=R_1\mathfrak M$.
\end{proposition}

The algebraic reduction in \cref{sec:trivial-system-elimination}
proves invertibility of each compression. The estimate is proved
in \cref{sec:trivial-system-estimates}, with the constants chosen
in \cref{sec:trivial-system-calculus}. Its additional regularity
step concerns $p''=\mathsf R\mathsf G'D^{\mathsf T}X_5$:
a local half-H\"older bound permits its gradient to enter the
mixed estimate for $\mathsf G$.

For the local estimates below use the exponents
$\beta\in\{1/2,3/4\}$.
Let $\Delta^{3/2}(y)=\{|x-y|_\infty\le3\xi_y/4\}$.
The base point of a site is the site itself and that of a bond,
also with an extra direction index, is its midpoint. Put
\begin{equation}\label{eq:trivial-system-holder-norm}
 \begin{split}
 [f]_{\beta,y}
 &=\sup_{\substack{i\ne i'\ \mathrm{like}\\
       \mathrm{base}(i),\mathrm{base}(i')\in\Delta^{3/2}(y)}}
       \left(\frac{\xi_y}{|i-i'|_1}\right)^{\beta}|f(i')-f(i)|,\\
 |f|^{\beta}_{s;y}
 &=\xi_y^{-s}\left(
    \sup_{\mathrm{base}(i)\in\Delta^{3/2}(y)}|f(i)|
       +[f]_{\beta,y}\right).
 \end{split}
\end{equation}
Here like bond indices have the same direction and, when present,
the same extra direction; $|i-i'|_1$ is the distance between their
initial points, equivalently between their midpoints. Empty suprema
are zero. Positions in \eqref{I-64-3}--\eqref{I-64-4} continue to use
initial points.

\subsection{Identification at the trivial configuration}
\label{sec:trivial-system-identification}

The averages and the action Hessian at the identity give the
operators in \eqref{I-63-64}. We check their normalization and
boundary conventions.

Choose a Hermitian basis $t_1,t_2,t_3$ of $\mathfrak g$ with
$\frac12\tr(t_jt_k)=\delta_{jk}$. The coordinate expression
$f=\sum_jf^jt_j$ makes the pairing nondegenerate, by testing against
$\sum_j\overline{f^j}t_j$, and gives the transpose identity
\[
 (K\otimes1)^{\mathsf T}=K^*\otimes1
\]
for every real operator $K$ used here. In particular, its real
inverse extends coordinatewise to the inverse on these spaces.
The finite-dimensional norm comparison needed below is
\[
 \left|\tfrac12\tr(t_jX)\right|\le|X|,
 \qquad
 |X|\le\sum_{j=1}^3
       (|\operatorname{Re}X^j|+|\operatorname{Im}X^j|).
\]
Indeed $|t_j|=1$, since its eigenvalues are $\pm1$.
Thus the real pointwise, difference and H\"older estimates apply
to these fields, with at most a fixed factor from six real
coordinates.

All contours and averages of the identity configuration are the
identity. Hence the scalar transported block average is the
arithmetic mean, and its iteration is the existing $Q'$.
For the bond average, the endpoint factors in the straightening
must also be differentiated. In the fine-lattice coordinates of
\cref{I-10-93}, put
\[
 \sigma_a(y)=L^{-3}\sum_{x\in B(y)}\bar a_{\mathsf m}(y,x),
 \qquad
 v_s(y)=\exp\left(L^{-3}\sum_{x\in B(y)}
             \log (e^{isa})_{\mathsf m}(y,x)\right).
\]
For $s$ near zero these logarithms are defined, and the contour
derivative in \cref{I-10-93} gives
$v_s(y)=1+is\sigma_a(y)+O(s^2)$.
The one-step straightened average at the identity is
\[
 i^{-1}\log\left[
   v_s(c_-)^{-1}\mathcal M(e^{isa})(c)v_s(c_+)\right].
\]
By \eqref{I-10-100}, its derivative at $s=0$ is
\begin{equation}\label{eq:trivial-straightened-derivative}
 \mathbf Q^{\mathrm{lin}}_{\mathsf m}a(c)
       -\sigma_a(c_-)+\sigma_a(c_+)
   =L^{-3}\sum_{x\in B(c_-)}a([x,x']).
\end{equation}
Here $x\mapsto x'$ bijects the two endpoint blocks, so both
endpoint path means cancel. Iteration gives
\[
 L^{-3n}\sum_{x\in B^n(c_-)}
                a([x,x+\xi_ne_\mu]).
\]
To see the induction, decompose $B^n(c_-)$ into the
$B^{n-1}(z)$, and each length-$\xi_n$ path into $L$ successive
length-$\xi_{n-1}$ paths. Translation carries
$B^{n-1}(z)$ onto $B^{n-1}(z+i\xi_{n-1}e_\mu)$, so the
inner path sums concatenate before the block sums are taken.
Applying this formula to $a=\xi A$ and dividing the output by
$\xi_n$ gives the factor $L^{-4n}$ in \eqref{I-16-63}.
At level zero the map is the identity. All these paths are inner
by \cref{I-11-27}. Thus the straightened linearisation is precisely
$Q$. Its positive rows are the existing $\Gamma_n$: removing
inward bonds leaves exactly the bonds with an endpoint in
$\Omega_n$ and none in $\Omega_{n+1}$. The scalar weights start
at $a_0=a_1=1$ and obey the existing recursion, while the row
weight is $\xi_c^{-2}$.

The scalar difference at the identity is
\[
 \Delta^\circ f(x)=\xi^{-2}
    \sum_{\nu=1}^3\bigl(2f(x)-f(x+\xi e_\nu)-f(x-\xi e_\nu)\bigr),
\]
with zero extension. It is therefore $\Delta_{\mathcal S}$,
and the inner gradient is $D=\partial_0$.
At the identity, the exterior derivative \eqref{I-2-4} is the
ambient curl $\partial$ applied coordinatewise. By
\cref{I-2-14}, with $c=1$ on the whole torus,
\[
 \left.\frac{d^2}{ds^2}\right|_{s=0}
       A^\xi(1,e^{is\xi A})=(\partial A,\partial A)_\xi.
\]
Polarization identifies the ambient Hessian with
$\partial^*\partial\otimes1$. Restrict its bilinear form to
inner bond fields extended by zero to obtain $\mathcal L$.
The curl energy still sums over plaquettes crossing the boundary
of $\Omega_0$.

These identities identify $\Delta'_a,\mathsf C,\mathsf P,\mathsf R$
with their coordinatewise extensions from
\cref{sec:multiscale-gram}, and give
\begin{equation}\label{I-64-13}
 \mathsf L_1-D\mathsf PD^{\mathsf T}
   =\mathcal L+D\mathsf RD^{\mathsf T}
                  +Q^{\mathsf T}\mathsf m^{-2}Q
   =\mathsf G^{-1},\qquad
 Q\mathsf GQ^{\mathsf T}=\mathsf Z.
\end{equation}
Under the bounds of \cref{I-64-24}, the threshold choice in
\cref{sec:trivial-system-calculus} makes
\cref{I-11-30,I-18-76} applicable. Their inverses exist, and the
same formulas give $\mathsf H=\mathsf GQ^{\mathsf T}\mathsf Z^{-1}$.
No positivity of the complex bilinear form is used.

\subsection{Elimination of the auxiliary variables}
\label{sec:trivial-system-elimination}

We solve the two scalar constraint blocks before solving for $A$.
Put
\[
 \mathsf N=\mathsf G'Q'^{\mathsf T}\mathsf C,\qquad
 \mathsf P=\mathsf NQ'\mathsf G',\qquad
 \mathsf R=I-\mathsf P.
\]
For the three equations
\[
 \Delta'_av+w=k_2,\qquad
 \Delta'_aw-Q'^{\mathsf T}\sigma=f_3,\qquad Q'v=f_4,
\]
substitution of
$w=\mathsf G'(f_3+Q'^{\mathsf T}\sigma)$ and
$v=\mathsf G'(k_2-w)$ into the last equation gives
\[
 \mathsf C^{-1}\sigma
       =Q'\mathsf G'k_2-Q'\mathsf G'^2f_3-f_4.
\]
Consequently the unique solution is
\begin{equation}\label{I-63-68}
 \begin{split}
 \sigma&=\mathsf C(Q'\mathsf G'k_2-Q'\mathsf G'^2f_3-f_4),\\
 w&=\mathsf Pk_2+\mathsf R\mathsf G'f_3-\mathsf Nf_4,
 \qquad v=\mathsf G'(k_2-w).
 \end{split}
\end{equation}
This proves invertibility of $\mathbb K_0^{[2]}$ by taking
$k_2=f_2$; that of $\mathbb K_0^{[1]}$ follows from
$\mathsf G'=(\Delta'_a)^{-1}$.

For the other scalar block, put $h_5=f_5+D^{\mathsf T}X_5$.
Substitute $q=\mathsf G'(f_6-p)$ and
$p=\mathsf G'(h_5-Q'^{\mathsf T}r)$ into $-Q'q=f_7$.
This gives
\begin{equation}\label{I-63-69}
 \begin{split}
 r&=\mathsf C(Q'\mathsf G'^2h_5-Q'\mathsf G'f_6-f_7),\\
 p&=\mathsf R\mathsf G'h_5+\mathsf Pf_6+\mathsf Nf_7,
 \qquad q=\mathsf G'(f_6-p).
 \end{split}
\end{equation}
In particular this block is independent of $A$.

For the full system, apply \eqref{I-63-68} with
$k_2=f_2+D^{\mathsf T}A$ and put
\begin{equation}\label{I-63-71}
 \varphi=\mathsf Pf_2+\mathsf R\mathsf G'f_3-\mathsf Nf_4,
 \qquad F_{\mathrm{eff}}=f_1+D\varphi+Dp.
\end{equation}
Then
\begin{equation}\label{I-63-72}
 w=\mathsf PD^{\mathsf T}A+\varphi,\qquad
 v=\mathsf G'(f_2+D^{\mathsf T}A-w),
\end{equation}
and the first equation, by \eqref{I-64-13}, becomes
$\mathsf G^{-1}A=F_{\mathrm{eff}}+Q^{\mathsf T}\omega$.
For $i=4$ the last term is absent, so
$A=\mathsf GF_{\mathrm{eff}}$.
For $i=5$, imposing $QA=f_8$ gives
\begin{equation}\label{I-63-73}
 \omega=\mathsf Z^{-1}(f_8-Q\mathsf GF_{\mathrm{eff}}),
 \qquad
 A=\mathsf G(F_{\mathrm{eff}}+Q^{\mathsf T}\omega).
\end{equation}
Together with \eqref{I-63-68}--\eqref{I-63-72}, these formulas
determine every component uniquely and solve its equation.
They prove invertibility of both remaining compressions.

\subsection{Local estimates for the constrained system}
\label{sec:trivial-system-calculus}

We use the scalar, vector and coarse inverse estimates already proved,
with a common decay rate. Write $\delta_G$ for the rate in
\cref{I-17-5}, and use $\delta_{0,\mathrm{lap}}$, $\delta_1$ and
$\theta$ from \cref{I-13-20}, \eqref{I-14-35} and
\eqref{I-18-77}. Set
\begin{equation}\label{eq:trivial-system-rates}
 \delta_*=\min\{1,\delta_{0,\mathrm{lap}}/2,2\delta_1,
                   \delta_G,\theta/2\},
 \qquad \delta_{\mathrm{cs}}=\delta_*/2.
\end{equation}
Choose a power $\mathfrak M_\dagger$ of $L$ so large that
\begin{equation}\label{eq:trivial-system-thresholds}
 \begin{split}
 \mathfrak M_\dagger&\ge\max\{L^2,L+10,M_0,M_1\},\\
 L\mathfrak M_\dagger&\ge\max\{L(L+9),
       5\delta_{0,\mathrm{lap}}^{-1}\log L,\mathfrak r_4,
       16\delta_*^{-1}\log L,\mathsf r_S(\delta_*/4)\}.
 \end{split}
\end{equation}
Here $M_0,M_1$ are the thresholds in \cref{I-17-5,I-18-76}.
If $\mathfrak m\ge \mathfrak M_\dagger$ and $\mathsf g\ge L\mathfrak M_\dagger$,
refining the region cubes to side $\mathfrak M_\dagger\xi_n$ puts the same
sequence in the preceding setting with $(\mathfrak M,R_1)=(\mathfrak M_\dagger,L)$.
The layers, operators and metric are unchanged. All the cited estimates
apply, and \cref{lem:global-cell-geometry} applies with this choice.
The empty-region case has only zero spaces; below assume $\Omega_0\ne\emptyset$.
Constants denoted by $C$ depend only on $L$ and may increase.
All top indices below belong to $\Lambda_*$.

For the finite component spaces and local seminorms above, write
$T\in\mathfrak O(\delta_{\mathrm{cs}},C)$ when
\begin{equation}\label{eq:trivial-system-sup-decay}
 |Tf|_y\le C\sup_z e^{-\delta_{\mathrm{cs}}d(y,z)}|f|_z.
\end{equation}
Composition multiplies these constants without changing the rate:
insert the two bounds and use $d(y,z)+d(z,z')\ge d(y,z')$.
Finite sums add their constants. A bound by the maximum over
$d(y,z)\le R$ gives \eqref{eq:trivial-system-sup-decay} with the
additional factor $e^{\delta_{\mathrm{cs}}R}$.

Only one summation is needed to pass from the preceding localized
estimates to these bounds. Suppose $f=\sum_z h_z$ and
$N_y(Th_z)\le C\xi_y^a e^{-\delta_*d(y,z)}k_z$, where $N_y$
is a target seminorm. For $|s|\le4$, the scale comparison and
\eqref{eq:trivial-system-thresholds} give
$(\xi_z/\xi_y)^s\le L^4 e^{\delta_*d(y,z)/4}$.
Consequently \eqref{I-13-122} gives
\begin{equation}\label{eq:trivial-system-weighted-sum}
 \begin{split}
 \xi_y^{-s-a}N_y(Tf)
 &\le CL^4\sum_z e^{-3\delta_*d(y,z)/4}\xi_z^{-s}k_z\\
 &\le CL^4c_S(\delta_*/4)
       \sup_z e^{-\delta_{\mathrm{cs}}d(y,z)}\xi_z^{-s}k_z.
 \end{split}
\end{equation}
The witnesses $c_S,\mathsf r_S$ are those of \eqref{I-13-121}.
All losses below are absorbed in $C$; subsequent compositions do not
repeat this summation or reduce the decay rate.

We abbreviate a bound of this kind between weights $s,s'$ by
$T:s\to s'$. From \cref{I-13-20,I-17-5},
\eqref{I-14-35} and \eqref{I-18-77},
\eqref{eq:trivial-system-weighted-sum} gives the bounds used below:
\begin{equation}\label{eq:trivial-system-basic-bounds}
 \begin{gathered}
 \mathsf G':-4\to-2,\quad -2\to0,
 \qquad \nabla\mathsf G':-4\to-3,\quad -2\to-1,\\
 \mathsf G'\nabla^{\mathsf T}:-3\to-2,\quad -1\to0,
 \qquad \mathsf C:0\to-4,\qquad
 \mathsf Z^{-1}:-1\to-3,\\
 \mathsf G:-3\to-1,\qquad
 \nabla\mathsf G:-3\to-2,\qquad
 \mathsf G\nabla^{\mathsf T}:-2\to-1.
 \end{gathered}
\end{equation}
For $\mathsf C$ a cell contains just its one top index; for
$\mathsf Z^{-1}$ at most six rows have a given position, by
\eqref{eq:coarse-row-sum}. These are counting-kernel bounds, with
output factors $\xi_y^{-4}$ and $\xi_y^{-2}$, respectively.

The scalar average and its weighted transpose act within one cell:
$Q'$ takes its arithmetic mean, and
$(Q'^{\mathsf T}\nu)(x)=\nu(y(x))$. Thus both preserve every
weight. The bounds for $Q,Q^{\mathsf T}$ follow from
\eqref{I-18-15} and the double-block geometry in
\cref{sec:coarse-covariance-decay}. The averaging weights are
nonnegative with sum one, and their weighted-adjoint coefficients
are at most one. A fine bond belongs to double blocks of at most
two consecutive levels and, at either level, to those of at most
six incident rows. Their positions are within distance $C(L)$,
and their scales differ by at most $L$. Hence $Q,Q^{\mathsf T}$
preserve each weight $|s|\le4$ in the sense of
\eqref{eq:trivial-system-sup-decay}. This also explains why no
additional volume factor enters the coarse constraint.

Two local comparisons convert the Hölder estimates to the norms
of this section. If a cell meets $\widetilde\Delta(y)$, take a
bare site of that cell within $\xi/2$ of a point of intersection.
It lies within $2\xi_y$ of $y$, so
\eqref{eq:nearby-cell-comparison} gives
\begin{equation}\label{eq:trivial-system-nearby}
 L^{-1}\xi_y\le\xi_z\le L\xi_y,
 \qquad d(y,z)\le9L
 \quad\text{if }\Delta(z)\cap\widetilde\Delta(y)\ne\emptyset.
\end{equation}
For like indices with base points in $\Delta^{3/2}(y)$, the
coordinate path between their sites or initial points keeps the
corresponding base points in that cube. Its successive differences
telescope, using the departure or arrival index according to the
sign of the step. If their total length is $h$, then
$(\xi_y/h)^\beta h\le(9/2)\xi_y$ for both selected exponents. At nonzero levels these indices are inner by
the local chart; at level zero there is no distinct site pair,
and the only possible bond pair in one direction is
$\langle y-\xi e_\mu,y\rangle$, $\langle y,y+\xi e_\mu\rangle$.
Thus
\begin{equation}\label{eq:trivial-system-comb-holder}
 [f]_{\beta,y}\le\tfrac92\xi_y
       \max_{\Delta(z)\cap\Delta^{3/2}(y)\ne\emptyset}
                 |\nabla f|_z.
\end{equation}
At levels zero and one the simpler bound
$[f]_{\beta,y}\le2L\sup_{\Delta^{3/2}(y)}|f|$ follows from
the minimum nonzero distance $\xi\ge\xi_y/L$.

Choose a fixed real smooth cutoff
$\zeta_y^*$ supported in $|x-y|_\infty<9\xi_y/10$, equal to one
for $|x-y|_\infty\le4\xi_y/5$, with
$\|\zeta_y^*\|_{\beta,\xi_y}+|\zeta_y^*|\le C$ for both exponents.
Such torus cutoffs are products of even one-dimensional cutoffs
in the least coordinate distance; their support stays away from
the half-period, so they are smooth. At a positive level,
$3\xi_y/4+\xi/2\le4\xi_y/5$, so the cutoff is one at the
initial points of all indices with base points in $\Delta^{3/2}(y)$.
Consequently
\begin{equation}\label{eq:trivial-system-output-holder}
 \sup_{\Delta^{3/2}(y)}|u|+[u]_{\beta,y}
 \le\xi_y^\beta\|\zeta_y^*u\|_\beta
                        +3\sup_{\Delta^{3/2}(y)}|u|.
\end{equation}
Pairs at distance at most one enter the cutoff seminorm;
the remaining pairs are bounded by twice the supremum, since
$\xi_y\le1$. At level zero the left side is at most three
times the supremum, by the minimum nonzero distance $\xi_y$.

For scalar divergence smoothing, let $h_z$ be the input restricted
to indices of position $z$, with zero extension to the bare lattice.
The cutoff estimate \eqref{eq:global-scalar-divergence-holder}
and its pointwise companion, applied at the nearby cells in
\eqref{eq:trivial-system-nearby}, bound the enlarged local norm
of $\mathsf G'\nabla^{\mathsf T}h_z$ by
$C\xi_y e^{-\delta_*d(y,z)}|h_z|_z$.
Use \eqref{eq:trivial-system-output-holder} at positive levels
and the supremum bound at level zero. Summing by
\eqref{eq:trivial-system-weighted-sum} gives
\begin{equation}\label{eq:trivial-system-scalar-holder}
 \mathsf G'\nabla^{\mathsf T}:
       (-3)\longrightarrow(-2,\beta),
\end{equation}
where the pair denotes the enlarged H\"older norm.
The same assertion holds for $\mathsf G'D^{\mathsf T}$:
set $X_A(x,\mu)=A(\langle x,x+\xi e_\mu\rangle)$ on inner
bonds and zero otherwise. Then
$D^{\mathsf T}A=\nabla^{\mathsf T}X_A$ and
$|X_A|_{s;y}\le|A|_{s;y}$. Similarly every gradient bound for
$\mathsf G'$ restricts to $D\mathsf G'$.

For the mixed vector estimate a partition at midpoint base points
is needed. Choose $0\le\phi_y\le1$ by translating and rescaling
a single fixed smooth product cutoff at scale $\xi_y$, so that
$\phi_y=1$ on $\Delta(y)$, $\phi_y=0$ outside the interior of
$\Delta^{3/2}(y)$, and
\[
 |\phi_y(p)-\phi_y(p')|
 \le C_\phi |p-p'|_1/\xi_y
\]
with $C_\phi$ independent of $y$ and the region sequence. Put
\[
 \zeta_y(p)=\frac{\phi_y(p)}{\sum_z\phi_z(p)}\qquad(p\in \Omega_0).
\]
The denominator is at least one. If $p\in\Delta^{3/2}(y)$,
$p'\in \Omega_0$, $|p-p'|_1\le\xi_y$, and $\phi_z$ is nonzero at
either point, then $|z-y|_\infty\le2\xi_y+\xi_z$.
Separation excludes levels differing by two or more: the gap to
the larger level exceeds $L$ times its spacing, whereas this
distance is less than three times that spacing. Counting lattice
points at the three remaining levels gives at most $C(L)$ such
centers, all with $\xi_z\ge\xi_y/L$. The quotient rule, or
$|a/h-a'/h'|\le|a-a'|+|h-h'|$ for
$0\le a,a'\le1\le h,h'$, yields
\[
 |\zeta_y(p)-\zeta_y(p')|
       \le C(L)|p-p'|_1/\xi_y.
\]

Evaluate this partition at bond midpoints. For a field $X$ with
one extra direction index, the two-point product identity gives
\begin{equation}\label{eq:trivial-system-partition-holder}
 \|\zeta_zX\|_{\beta,\xi_z}+|\zeta_zX|
       \le C\xi_z^s|X|^\beta_{s;z}
\end{equation}
on inner indices. Indeed, when both base points are in the
support cube use the Hölder difference of $X$ and the displayed
Lipschitz bound. If only one is in that cube, the product at the
other is zero and the same Lipschitz bound controls the difference;
if neither is there, both products vanish.

The initial point of a bond with midpoint in this support lies
in the closed $\widetilde\Delta(z)$. At positive levels use
$3\xi_z/4+\xi/2<\xi_z$; at level zero the initial point may
be $z$ or $z-\xi e_\mu$. In particular the incoming bond must
not be discarded. Extend $\zeta_zX$ by zero to all bare bonds.
For a pair at distance at most $\xi_z$ meeting the support,
all endpoints lie within $3\xi_z$ of $z$. At positive levels
the gap places them in $\Omega_0$, so the inner-index Hölder bound
applies. At level zero a distinct pair has distance exactly
$\xi_z$, and a pair crossing the index boundary contributes at
most the supremum. Thus zero extension costs at most a factor
two in \eqref{eq:trivial-system-partition-holder}.
Since $X=\sum_z\zeta_zX$, apply
\eqref{eq:global-dirichlet-mixed} at exponent $1/2$ to these
zero extensions and then \eqref{eq:trivial-system-weighted-sum}.
Restriction of the output to inner indices gives
\begin{equation}\label{eq:trivial-system-mixed-holder}
 \nabla\mathsf G\nabla^{\mathsf T}:
       (-2,1/2)\longrightarrow(-2).
\end{equation}

The same output cutoff gives the two further bounds needed below:
\begin{equation}\label{eq:trivial-system-gradient-holder}
 \nabla\mathsf G:(-3)\longrightarrow(-2,1/2),\qquad
 \nabla\mathsf G\nabla^{\mathsf T}:
             (-2,3/4)\longrightarrow(-2,1/2).
\end{equation}
For the first estimate apply
\eqref{eq:trivial-system-output-holder}, with its level-zero
alternative, to $\nabla\mathsf Gh_z$ and use
\eqref{eq:global-dirichlet-holder} with exponent $1/2$.
Its pointwise companion bounds the enlarged supremum at the
nearby cells in \eqref{eq:trivial-system-nearby}.
Summation by \eqref{eq:trivial-system-weighted-sum} proves
the first bound.

For the second, put $J_z=\zeta_zX$, with the zero extension just
checked, and $V_z=\nabla\mathsf G\nabla^{\mathsf T}J_z$.
The partition estimate at $\beta=3/4$ and
\eqref{eq:global-dirichlet-mixed-holder}, with
$(\alpha,\varepsilon)=(1/2,1/4)$, give
\[
 \xi_y^{1/2}\|\zeta_y^*V_z\|_{1/2}
 \le C e^{-\delta_*d(y,z)}\xi_z^{-2}|X|^{3/4}_{-2;z}.
\]
The enlarged supremum obeys the same bound by
\eqref{eq:global-dirichlet-mixed} at exponent $3/4$, applied
at nearby cells. That supremum also handles level zero.
Now \eqref{eq:trivial-system-weighted-sum}, with $s=-2$ and
$a=0$, proves the second bound. Both estimates retain
$\delta_{\mathrm{cs}}$ and the thresholds
\eqref{eq:trivial-system-thresholds}.

\subsection{Bounds for the solution}
\label{sec:trivial-system-estimates}

\begin{proof}[Quantitative bound in \cref{I-64-24}]
Invertibility and the solution formulas were proved in
\cref{sec:trivial-system-elimination}. Fix a source $g$.
We say that a component has weight $s$ if its norm $|\cdot|_{s;y}$
is bounded by
$C\sup_z e^{-\delta_{\mathrm{cs}}d(y,z)}|g|_z$;
Hölder weight $(s,1/2)$ means the same bound in the enlarged norm.
The latter dominates the ordinary norm, since the base point of
an index lies in the cell of its position.

Composing \eqref{eq:trivial-system-basic-bounds} and the averaging
bounds gives
\[
 \begin{array}{c|cc}
 T&\text{bound for }T&\text{bound for }\nabla T\\ \hline
 \mathsf N&0\to-2&0\to-3\\
 \mathsf P&-2\to-2&-2\to-3\\
 \mathsf PD^{\mathsf T}&-1\to-2&-1\to-3\\
 \mathsf R\mathsf G'&-4\to-2&-4\to-3.
 \end{array}
\]
The last line uses $\mathsf R\mathsf G'=\mathsf G'-\mathsf P\mathsf G'$.
The same gradient bounds hold after restriction to $D$.

In addition,
\begin{equation}\label{eq:trivial-system-projected-holder}
 \mathsf R\mathsf G'D^{\mathsf T}:
          (-3)\longrightarrow(-2,\beta).
\end{equation}
For $\mathsf G'D^{\mathsf T}$ this is
\eqref{eq:trivial-system-scalar-holder}. For the subtracted term
$f=\mathsf P\mathsf G'D^{\mathsf T}X$, the displayed compositions
bound $f$ in weight $-2$ and $\nabla f$ in weight $-3$.
The comb estimate and scale comparison give
\[
 |f|^{\beta}_{-2;y}
  \le C\max_{d(y,z)\le9L}
                  \{|f|_{-2;z},|\nabla f|_{-3;z}\}.
\]
Indeed the supremum term costs at most $L^2$, and
\eqref{eq:trivial-system-comb-holder} costs at most
$(9/2)L^3$ after multiplying by $\xi_y^2$.
The local comparison rule following
\eqref{eq:trivial-system-sup-decay} proves
\eqref{eq:trivial-system-projected-holder} without a scalar
mixed-derivative estimate.

For $i=4,5$, split the scalar solution as
\begin{equation}\label{eq:trivial-system-p-split}
 p=p'+p'',\qquad
 p'=\mathsf R\mathsf G'f_5+\mathsf Pf_6+\mathsf Nf_7,
 \qquad p''=\mathsf R\mathsf G'D^{\mathsf T}X_5.
\end{equation}
The source weights and these bounds show that $p$ has weight $-2$,
$Dp'$ weight $-3$, and $p''$ Hölder weight $(-2,\beta)$.
The formulas for $q$ and $r$ then give weights $0,-1,-4$ for
$q,Dq,r$. Equation \eqref{I-63-71} likewise gives weights $-2,-3$
for $\varphi,D\varphi$.

The effective vector source is
$F_{\mathrm{eff}}=f_1+D\varphi+Dp'+Dp''$.
The derivative of $p''$ is handled as a divergence, rather than
estimated in the supremum norm. For an inner bond
$b=\langle x,x+\xi e_\mu\rangle$, define
\begin{equation}\label{eq:trivial-system-shift}
 X^{(p)}_\nu(b)=-\delta_{\mu\nu}p''(b_+).
\end{equation}
If $b-\xi e_\mu$ is inner, the sole nonzero component in the
divergence gives
\[
 (\nabla^{\mathsf T}X^{(p)})(b)
   =\xi^{-1}\bigl[p''(b_+)-p''(b_-)\bigr]=(Dp'')(b).
\]
Otherwise its predecessor value is zero and
\[
 Z^{(p)}(b):=(Dp''-\nabla^{\mathsf T}X^{(p)})(b)
              =-\xi^{-1}p''(x).
\]
In this case $x\in\Lambda_0$: a site in $\Omega_1$ has all
its neighbors in $\Omega_0$ by separation. Thus the position of $b$ is
$x$, $\xi_x=\xi$, and
$|Z^{(p)}|_{-3;x}\le|p''|_{-2;x}$.
At the other positions this correction vanishes. Hence
\begin{equation}\label{eq:trivial-system-vector-source}
 F_{\mathrm{eff}}=f'_1+\nabla^{\mathsf T}X^{(p)},
 \qquad f'_1=f_1+D\varphi+Dp'+Z^{(p)},
\end{equation}
where $f'_1$ has weight $-3$.

We check that the endpoint shift in \eqref{eq:trivial-system-shift}
preserves the required Hölder weight. For a scalar function $u$,
let $X^{(u)}_\nu(b)=-\delta_{\mu\nu}u(b_+)$.
If the midpoint of $b$ lies in $\Delta^{3/2}(y)$, both endpoints
lie in the closed $\widetilde\Delta(y)$. This follows from
$3\xi_y/4+\xi/2\le\xi_y$ at positive levels; at level zero
the possible bonds are the incoming and outgoing bonds at $y$.
Consequently the supremum of $X^{(u)}$ is controlled by the
suprema of $u$ in cells meeting $\widetilde\Delta(y)$.
Their scales and distances obey \eqref{eq:trivial-system-nearby}.
At levels zero and one the lattice-spacing bound following
\eqref{eq:trivial-system-comb-holder} controls the Hölder seminorm
as well.

At a level $j\ge2$, choose $b^\sharp=b_+$ if that endpoint is
in $\Delta^{3/2}(y)$, and $b^\sharp=b_-$ otherwise; the latter
endpoint then lies in this cube. For $z=y(b_+)$ both endpoints
lie in $\Delta^{3/2}(z)$: the cell of $b_+$ contains its closed
bare cube, and
$|b_\pm-z|_\infty\le\xi_z/2+\xi/2\le3\xi_z/4$,
since $\xi_z\ge L\xi$.
For like bonds $b,b'$ put $r=|b_--b'_-|_1\ge\xi$ and split
\[
 \begin{split}
 u(b'_+)-u(b_+)
  &=[u(b'_+)-u(b'^\sharp)]
    +[u(b'^\sharp)-u(b^\sharp)]\\
  &\quad+[u(b^\sharp)-u(b_+)].
 \end{split}
\]
The middle pair lies in $\Delta^{3/2}(y)$ at distance at most
$r+2\xi\le3r$. Each outer pair lies in the enlarged cube of
the endpoint's position and has distance $\xi$.
Multiplication by $(\xi_y/r)^{\beta}$ costs at most $3^\beta$
for the middle Hölder quotient and at most $L^{\beta}$ for each
outer quotient. Finally $\xi_y^2\le L^2\xi_z^2$.
Together with the low-level case this proves
\begin{equation}\label{eq:trivial-system-shift-holder}
 |X^{(u)}|^{\beta}_{-2;y}
       \le C\max_{\Delta(z)\cap\widetilde\Delta(y)\ne\emptyset}
                         |u|^{\beta}_{-2;z}.
\end{equation}
The local comparison rule now transfers the Hölder weight of
$p''$ to $X^{(p)}$.

Apply \eqref{eq:trivial-system-basic-bounds} and
\eqref{eq:trivial-system-mixed-holder} to
\eqref{eq:trivial-system-vector-source}. They give weights $-1,-2$
for $\mathsf GF_{\mathrm{eff}}$ and
$\nabla\mathsf GF_{\mathrm{eff}}$.
For $i=4$ this is the estimate for $A$.
For $i=5$ use the coarse-constraint formulas \eqref{I-63-73}.
Since $f_8$ and $Q\mathsf GF_{\mathrm{eff}}$ have weight $-1$,
$\omega$ has weight $-3$. The bounds for $Q^{\mathsf T}$,
$\mathsf G$ and $\nabla\mathsf G$ preserve the weights $-1,-2$
of $A,\nabla A$.

It remains to recover the first scalar subsystem. Its formulas give
$w=\mathsf PD^{\mathsf T}A+\varphi$, so $w,Dw$ have weights
$-2,-3$. To bound $v$ we must separately check $D^{\mathsf T}A$.
At a site $x$, in each direction whose predecessor bond is inner,
\[
 (\nabla_\mu^{\mathsf T}A_\mu)(x)
                  =-(\nabla_\mu A_\mu)(x-\xi e_\mu).
\]
If that bond is missing, the term instead equals
$-\xi^{-1}A_\mu(x)$, and vanishes if the outgoing bond is also
missing. In the nonzero exceptional case $x\in\Lambda_0$, so
$\xi_x^2\xi^{-1}|A_\mu(x)|\le|A|_{-1;x}$.
In the first case the predecessor's position has scale within
a factor $L$ of the position of $x$, and distance at most $9L$,
by \eqref{eq:nearby-cell-comparison}. Summing the three directions
therefore gives
\[
 |D^{\mathsf T}A|_{-2;y}
    \le C\max_{d(y,z)\le9L}
                       \{|\nabla A|_{-2;z},|A|_{-1;z}\}.
\]
Thus $D^{\mathsf T}A$ has weight $-2$. The elimination formulas
and the scalar bounds now give weights $0,-1,-4$ for $v,Dv,\sigma$.
For $i=2$ the same scalar calculation uses $A=0$ and $w=\varphi$;
for $i=1$, $w=\mathsf G'f_3$. These are the solutions of the
compressed systems, not restrictions of the full inverse.

Every entry in the solution norm has now been bounded at its
prescribed weight. Taking the largest of the finitely many constants
gives $B_{\mathrm{cs}}=B_{\mathrm{cs}}(L)$ and the estimate in
\cref{I-64-24}. The rates and $\mathfrak M_\dagger$ were fixed before this
constant was chosen. \end{proof}

\section{Complex averages near the identity}
\label{sec:flat-complex-averaging}
We estimate centered straightened averages and their first derivatives.
The derivative kernels retain a factor $\eta^3$, which compensates
the volume factor in the weighted transpose. The endpoint correction
and its iteration follow Ba{\l}aban's averaging construction
\cite[pp.~34--39, (110), (121)--(127)]{B-CMP98}; we use centered
contours. The differentiated iteration follows the method of
\cite[pp.~39--43, Propositions~5--7]{B-CMP98}.

For one averaging step, use level-$n$ units, $0\le n<K$, so
that fine bonds have length one. We use the centered scheme and
write $\mathcal M=\mathcal M_{\mathsf m}$ and
$W(y,x)=W_{\mathsf m}(y,x)$. For a positive bond $c$ of
$\mathbb T^{(n+1)}$, let
$\mathsf D_c=B(c_-)\cup B(c_+)$ and let $\mathsf P_c^+$ be its
positive inner fine bonds. All field coordinates below are on
positive bonds; configurations on reverse bonds are given by
inversion. Suprema are over $\mathsf P_c^+$ when the output is
at $c$.

Let $W$ be a complex configuration and put $\widetilde W=e^{ia}W$.
Whenever the contour variables, averages and series logarithms
exist, define
\begin{equation}\label{eq:straightening-maps}
\begin{aligned}
 \zeta(W,a)(y,x)
   &=i^{-1}\log\bigl(\widetilde W(y,x)W(y,x)^{-1}\bigr),\\
 \mathsf v(W,a)(y)
   &=\exp\left(iL^{-3}\sum_{x\in B(y)}\zeta(W,a)(y,x)\right),\\
 \mathcal Q^1(W,a)(c)
   &=i^{-1}\log\bigl[\mathsf v(W,a)(c_-)^{-1}
       \mathcal M(\widetilde W)(c)\mathsf v(W,a)(c_+)
       \mathcal M(W)(c)^{-1}\bigr],
\end{aligned}
\end{equation}
where $y=c_-$ or $c_+$ and $x\in B(y)$. Thus
\begin{equation}\label{eq:straightening-factorization}
 \mathcal M(e^{ia}W)(c)
 =\mathsf v(W,a)(c_-)\,e^{i\mathcal Q^1(W,a)(c)}
       \mathcal M(W)(c)\mathsf v(W,a)(c_+)^{-1}.
\end{equation}
Write $\mathbf Q(W)=d_a\mathcal Q^1(W,0)$ and
$C^1(W,a)=\mathcal Q^1(W,a)-\mathbf Q(W)a$.
For these kernels, $\|T\|_{\mathrm{row}}=\sup_c\sum_e|T(c,e)|$
and $\|T\|_{\mathrm{col}}=\sup_e\sum_c|T(c,e)|$, where the
entry norm is the operator norm on $\mathfrak g^c$.
On a box $\mathsf Y$ of level-$n$ sites, restrict the coarse
indices to $\mathsf D_c\subset\mathsf Y$ and the fine indices
to bonds with both endpoints in $\mathsf Y$.

\subsection{Iterated averages}
Fix $1\le k\le K$ and use level-$k$ units, so that the bare
spacing is $\eta=L^{-k}$. For a positive level-$k$ bond $c$, set
\[
 Y^{(n)}=B^{k-n}(c_-)\cup B^{k-n}(c_+),\qquad 0\le n\le k.
\]
For $n<k$, the bonds in $Y^{(n)}$ have both endpoints there;
at $n=k$ the bond set consists of $c$ alone. Suprema at level
$n$ are over these positive bonds. Bare fields are extended by
zero when a whole-lattice expression is used.

For a configuration $\mathbf U$ on the bare bonds in $Y^{(0)}$,
write $W_n=\mathcal M^n\mathbf U$ whenever these complex averages
exist. Define the straightened iterates by
\begin{equation}\label{eq:flat-straightened-recursion}
 a_0=\eta A,\qquad
 a_{n+1}=\mathcal Q^1(W_n,a_n),\qquad
 \mathcal Q_n(\mathbf U;A)=a_n.
\end{equation}
Thus $a_n$ is a field on the bonds in $Y^{(n)}$; definedness is
part of the proposition below. Put
\begin{equation}\label{eq:flat-averaging-products}
 Q_k(\mathbf U)=\eta\mathbf Q(W_{k-1})\cdots\mathbf Q(W_0),
 \qquad Q'_k(\mathbf U)=\mathsf q_{W_{k-1}}\cdots\mathsf q_{W_0},
\end{equation}
where
\[
 (\mathsf q_W\lambda)(y)
   =L^{-3}\sum_{x\in B(y)}\operatorname{Ad}(W(y,x))\lambda(x).
\]
The first product takes its value at $c$ and the second at
$y=c_-$ or $c_+$. Kernels are taken against counting measure.
At the identity, these are the straight bond average and the
arithmetic site average identified in
\cref{sec:trivial-system-identification}, in the present units.

\begin{proposition}\label{prop:flat-complex-averaging}
There are constants $c_{\mathrm{av}}>0$ and
$C_{\mathrm{av}},C_2,D\ge1$, depending only on $L$, with the
following properties. Let $A',A$ be $\mathfrak g^c$-valued bare
fields on $Y^{(0)}$, with
$\|A'\|_\infty,\|A\|_\infty<c_{\mathrm{av}}$, and put
$\mathbf U=e^{i\eta A'}$.
\begin{enumerate}[label=\textup{(\roman*)}]
\item The complex averages $W_n$ are defined for $0\le n\le k$.
They are holomorphic in $A'$ and have the form $W_n=e^{iB_n}$,
where $B_n$ is holomorphic, $\mathfrak g^c$-valued, and
\[
 \|B_n\|_\infty\le D L^n\eta\|A'\|_\infty.
\]
\item The iterates in \eqref{eq:flat-straightened-recursion} are
defined and jointly holomorphic, and
\[
 \|a_n\|_\infty\le2L^n\eta\|A\|_\infty.
\]
The derivative of $A\mapsto\mathcal Q_k(\mathbf U;A)$ at zero
is $Q_k(\mathbf U)$.
\item The kernel of $Q_k(\mathbf U)$ is supported on the bare
bonds in $Y^{(0)}$. Writing $a'=\|A'\|_\infty$, one has
\begin{equation}\label{eq:flat-complex-linear-kernels}
\begin{aligned}
 \sum_b|Q_k(\mathbf U)(c,b)|
   &\le\min\{2,1+C_{\mathrm{av}}a'\},\\
 |Q_k(\mathbf U)(c,b)|
   &\le\eta^3\min\{2,1+C_{\mathrm{av}}a'\},\\
 \sum_b|(Q_k(\mathbf U)-Q_k(\mathbf1))(c,b)|
   &\le C_{\mathrm{av}}a',\\
 |(Q_k(\mathbf U)-Q_k(\mathbf1))(c,b)|
   &\le C_{\mathrm{av}}\eta^3 a'.
\end{aligned}
\end{equation}
If $x_j$ is the level-$j$ site with $x\in B^j(x_j)$, then
\begin{equation}\label{eq:flat-complex-scalar-transport}
\begin{aligned}
 T_{\mathbf U}(y,x)
   &=W_{k-1}(y,x_{k-1})\cdots W_0(x_1,x),\\
 (Q'_k(\mathbf U)\lambda)(y)
   &=L^{-3k}\sum_{x\in B^k(y)}
        \operatorname{Ad}(T_{\mathbf U}(y,x))\lambda(x),\\
 |\operatorname{Ad}(T_{\mathbf U}(y,x))-1|
   &\le C_{\mathrm{av}}a'.
\end{aligned}
\end{equation}
\item For
$K_A(c,b)=\partial\mathcal Q_k(\mathbf U;A)(c)/\partial A(b)$,
\begin{equation}\label{eq:flat-complex-differentiated-kernel}
\begin{aligned}
 |K_A(c,b)|
   &\le\eta^3\bigl(1+C_{\mathrm{av}}a'
                         +C_2\|A\|_\infty\bigr),\\
 |K_A(c,b)-Q_k(\mathbf U)(c,b)|
   &\le C_2\eta^3\|A\|_\infty.
\end{aligned}
\end{equation}
\end{enumerate}
The constants are independent of the torus, the bond and $k$.
All maps are determined by the indicated local data.
\end{proposition}

\subsection{One-step estimates}\label{sec:one-step-straightening}

\begin{lemma}\label{lem:one-step-straightening}
There is $r_{\mathrm{st}}>0$, depending only on $L$, such that
\[
 r_{\mathrm{st}}\le
 \min\left\{r_*,\frac{\omega_M}{16},
                   \frac1{16\kappa^{\mathrm{SO}}_2}\right\}.
\]
Put
\[
 \kappa_{\mathrm{st}}=r_{\mathrm{st}}^{-1},\qquad
 K_{\mathrm{st}}=24L^3r_{\mathrm{st}}^{-2},\qquad
 H_{\mathrm{st}}=11r_{\mathrm{st}}^{-2}.
\]
For $W=e^{ib}$, $\|b\|_\infty,\|a\|_\infty<r_{\mathrm{st}}$,
the maps in \eqref{eq:straightening-maps} and
$i^{-1}\log\mathsf v(W,a)$ are defined and jointly holomorphic.
The three logarithmic maps take values in $\mathfrak g^c$,
vanish when $a=0$, and have norm at most
$\min\{1,\kappa_{\mathrm{st}}\|a\|_\infty\}$.
For real $b,a$ they take values in $\mathfrak g$ and
$\mathsf v(W,a)$ takes values in $G$; the averages agree with
the real small-curvature branch. In particular,
\begin{equation}\label{eq:straightening-contour-bound}
 W(y,x)=e^{i\zeta_b(y,x)},\qquad
 \zeta_b(y,x)=\zeta(\mathbf1,b)(y,x),\qquad
 |\zeta_b(y,x)|\le\kappa_{\mathrm{st}}\|b\|_\infty.
\end{equation}

For $h:\mathsf D_c\to G$, the transformed expressions are
defined and satisfy
\begin{equation}\label{eq:straightening-covariance}
\begin{aligned}
 \zeta(W^h,\operatorname{Ad}(h)a)(y,x)
   &=\operatorname{Ad}(h(y))\zeta(W,a)(y,x),\\
 \mathsf v(W^h,\operatorname{Ad}(h)a)(y)
   &=h(y)\mathsf v(W,a)(y)h(y)^{-1},\\
 \mathcal Q^1(W^h,\operatorname{Ad}(h)a)(c)
   &=\operatorname{Ad}(h(c_-))\mathcal Q^1(W,a)(c).
\end{aligned}
\end{equation}
Here $(\operatorname{Ad}(h)a)(e)=\operatorname{Ad}(h(e_-))a(e)$.

Let $\mathbf S=\mathbf Q(\mathbf1)$. Then
\[
 (\mathbf Sa)(c)=L^{-3}\sum_{x\in B(c_-)}a([x,x']),\qquad
 \|\mathbf S\|_{\mathrm{row}}\le L,\quad
 \|\mathbf S\|_{\mathrm{col}}\le L^{-2},
\]
where $x'=x+Le_\mu$ for a bond $c$ in direction $\mu$.
For $\|b\|_\infty<r_{\mathrm{st}}/16$ and
$\|a\|_\infty\le r_{\mathrm{st}}/2$,
\begin{equation}\label{eq:straightening-one-step-bounds}
\begin{aligned}
 \|\mathbf Q(e^{ib})-\mathbf S\|_{\mathrm{row}},\quad
 \|\mathbf Q(e^{ib})-\mathbf S\|_{\mathrm{col}}
   &\le K_{\mathrm{st}}\|b\|_\infty,\\
 |C^1(e^{ib},a)(c)|&\le H_{\mathrm{st}}\|a\|_\infty^2,\\
 |\partial_{a(e)}C^1(e^{ib},a)(c)|
   &\le H_{\mathrm{st}}\|a\|_\infty.
\end{aligned}
\end{equation}
The two remainder bounds at $c$ use only the supremum on
$\mathsf P_c^+$. All these kernels vanish unless
$e\in\mathsf P_c^+$. Their bounds hold on the whole lattice
and with the box restrictions above. In particular,
\begin{equation}\label{eq:straightening-derivative-column}
 \|d_a\mathcal Q^1(e^{ib},a)\|_{\mathrm{col}}
 \le L^{-2}+K_{\mathrm{st}}\|b\|_\infty
                   +6H_{\mathrm{st}}\|a\|_\infty.
\end{equation}
\end{lemma}

\begin{proof}
Use the finite-chart argument of \cref{I-10-93}, now for the
three centered directions and the two variables $(b,a)$.
By \cref{I-10-59}, the defining contour and averaging maps are
holomorphic near the identity configuration. At $(b,a)=(0,0)$,
all arguments of logarithms in \eqref{eq:straightening-maps}
are the identity. Products, inverses and series logarithms
therefore give a common neighborhood on which the three
logarithmic maps are holomorphic and bounded by one.
Choose $r_{\mathrm{st}}$ within this neighborhood and the
displayed upper bounds. For fixed $b$, the three maps vanish
at $a=0$; Cauchy's formula on complex lines gives the bound
$r_{\mathrm{st}}^{-1}\|a\|_\infty$.

If $b,a$ are real, telescoping the four unitary factors around
a plaquette gives
\[
 |\partial(e^{ia}e^{ib})(p)-\mathbf1|
 \le4(\|a\|_\infty+\|b\|_\infty)
 <8r_{\mathrm{st}}.
\]
The chosen radius and \cref{I-10-61,I-10-59} give the real
values and branch agreement. The $i^{-1}$ multiples of these
logarithms near the identity belong to $\mathfrak g$.
Their traces are holomorphic functions of $(b,a)$ vanishing
on the real slice, so they vanish throughout the joint domain,
by the real-slice uniqueness argument in the proof of
\cref{I-3-117}.
This proves the $\mathfrak g^c$ assertion.
Equation \eqref{eq:straightening-contour-bound} follows by
exponentiating the first definition at $(\mathbf1,b)$.

The identity $(e^{ia}W)^h=e^{i\operatorname{Ad}(h)a}W^h$
and \eqref{I-10-60} conjugate the first logarithm by $h(y)$
and the last by $h(c_-)$. Unitary conjugation preserves the
domain and commutes with the series logarithm and exponential;
this proves \eqref{eq:straightening-covariance}. Differentiation
gives, in particular,
\[
 \mathbf Q(W^h)(c,e)
 =\operatorname{Ad}(h(c_-))\mathbf Q(W)(c,e)
                         \operatorname{Ad}(h(e_-))^{-1},
\]
and the corresponding identities for the Taylor coefficients.
The conjugating maps are isometries.

The identity derivative is
\eqref{eq:trivial-straightened-derivative}. Each pair consisting
of $x\in B(c_-)$ and a bond of $[x,x']$ contributes $L^{-3}$
to its kernel. There are $L^3L$ such pairs for fixed $c$.
For a fixed fine bond $e$ in direction $\mu$, its possible
starting points are $x=e_--je_\mu$, $0\le j<L$; each belongs
to exactly one block $B(c_-)$. Thus there are $L$ pairs for
fixed $e$, giving the row and column bounds for $\mathbf S$.
Also, $\mathsf D_c$ has $2L^3$ sites, so
$\#\mathsf P_c^+\le6L^3$. A fine bond lies either inside one
block, in which case it belongs to the six incident double
blocks, or between two adjacent blocks, in which case it
belongs only to their double block. Consequently it belongs
to at most six sets $\mathsf P_c^+$.

Put $r=r_{\mathrm{st}}$ temporarily. For fixed $c,e$, the map
$b\mapsto\mathbf Q(e^{ib})(c,e)$ is holomorphic for
$\|b\|_\infty<r$ and bounded by $r^{-1}$: apply Cauchy's
formula in the $a(e)$ direction at zero to the joint map.
A second application bounds its derivative in the $b$
direction by $4r^{-2}$ on $\|b\|_\infty\le3r/4$.
Integration from zero to $b$ gives
\[
 |\mathbf Q(e^{ib})(c,e)-\mathbf S(c,e)|
 \le4r^{-2}\|b\|_\infty.
\]
The row and column counts just proved give the first line of
\eqref{eq:straightening-one-step-bounds} with the stated
$K_{\mathrm{st}}$.

For fixed $b$, Cauchy's formula also gives
\[
 |C^1(e^{ib},a)(c)|
 \le\frac{\|a\|_\infty^2}{r(r-\|a\|_\infty)}.
\]
On $\|a\|_\infty<3r/4$, the derivative
$\partial_{a(e)}\mathcal Q^1$ has norm at most $4/r$, and
its value at zero has norm at most $1/r$. Thus
$\partial_{a(e)}C^1$ is bounded by $8/r$ there and vanishes
at zero. Cauchy's formula on this ball bounds it by
$32\|a\|_\infty/(3r^2)$. For $\|a\|_\infty\le r/2$
these bounds are at most the asserted ones with
$H_{\mathrm{st}}=11r^{-2}$. The definitions give the stated
support. Summing the last bound over the at most six coarse
bonds incident to a fixed fine bond proves
\eqref{eq:straightening-derivative-column}. Restricting either
index set to the stated box can only decrease these sums.
\end{proof}
\begin{proof}[Proof of \cref{prop:flat-complex-averaging}]
Write $r=r_{\mathrm{st}}$, $\kappa=\kappa_{\mathrm{st}}$ and
$H=H_{\mathrm{st}}$ for the constants
of \cref{lem:one-step-straightening}. Put $\theta_n=L^n\eta$.
The sums used below are
\begin{equation}\label{eq:flat-iteration-geometric-sum}
 \sum_{j<n}\theta_j\le2L^{-1}\theta_n,
 \qquad \sum_{n<k}\theta_n\le2L^{-1}.
\end{equation}
Choose
\[
 c_0\le\min\{r/4,L^2\log2/(8H)\},\qquad
 D=4(8\kappa+2),\qquad
 c_\sharp=\min\{c_0,r/(16D)\}.
\]
We first work on $\|A'\|_\infty<c_\sharp$. The smaller radius
$c_{\mathrm{av}}$ will be chosen below.

\emph{The background averages.}
At the reference configuration $\mathbf1$, define
$a'_0=\eta A'$ and $a'_{n+1}=\mathcal Q^1(\mathbf1,a'_n)$.
The identity row bound and the quadratic remainder give, for
$t_n=\|a'_n\|_\infty/\theta_n$,
\[
 t_{n+1}\le t_n(1+HL^{-1}\theta_n t_n).
\]
As long as $t_j\le2t_0$ for $j\le n$, multiplication and
\eqref{eq:flat-iteration-geometric-sum} give
\[
 t_{n+1}\le t_0\exp(4HL^{-2}t_0)\le2t_0.
\]
Also $2t_0\le r/2$, so the same one-step domain is available
at the next stage. Induction proves the recursion and
$\|a'_n\|_\infty\le2\theta_n\|A'\|_\infty$.

Let $s_n=\mathsf v(\mathbf1,a'_n)=e^{i\sigma_n}$ and set
$v_0=1$, $v_{n+1}=v_n s_n$ on the next coarser sites.
The one-step estimate gives
$\|\sigma_n\|_\infty\le\kappa\|a'_n\|_\infty$.
At a positive level-$n$ bond $b$, the candidate
\[
 X_n(b)=v_n(b_-)e^{ia'_n(b)}v_n(b_+)^{-1}
\]
is a product of $2n+1$ exponentials. The sum of their exponent
norms is at most
\[
 E_n=2\sum_{j<n}\|\sigma_j\|_\infty+\|a'_n\|_\infty
     \le(8\kappa L^{-1}+2)\theta_n\|A'\|_\infty.
\]
Thus $|X_n-1|\le e^{E_n}-1\le2E_n<1/2$, and its series
logarithm defines $B_n=i^{-1}\log X_n$ with
\[
 \|B_n\|_\infty\le4E_n
      \le D\theta_n\|A'\|_\infty<r/16.
\]
These fields are holomorphic. For real $A'$, all factors belong
to $G$, and the logarithms are in $i\mathfrak g$; the
real-slice uniqueness argument in the proof of \cref{I-3-117}
then gives tracelessness for complex $A'$.

We now identify the candidate with the actual average.
At $n=0$, $X_0=e^{i\eta A'}$. If $W_n=X_n=e^{iB_n}$, its
restriction to every next-step double block lies in the domain
of \cref{lem:one-step-straightening}; hence $W_{n+1}$ is
defined and holomorphic. For real $A'$, the straightening
identity and unitary gauge covariance give
\[
 \mathcal M\bigl((e^{ia'_n})^{v_n}\bigr)
   =\bigl(e^{ia'_{n+1}}\bigr)^{v_n s_n}=X_{n+1}.
\]
Both sides are holomorphic on the same complex ball, so this
identity holds there. This proves the induction and (i).

The construction can first be made on the whole lattice with
the identity reference and the zero extension of $A'$.
The block identity
$B^{n+1}(y)=\bigcup_{z\in B(y)}B^n(z)$ implies that the one-step double block of every bond in the
prescribed bond set of $Y^{(n+1)}$ lies in $Y^{(n)}$.
Descending through these blocks shows that the restriction
uses only $A'$ in $Y^{(0)}$. The same argument applies to
all subsequent recursions and their derivatives.

\emph{Linear kernels and transports.}
The row and column norms are submultiplicative, by summing
$|(ST)(i,j)|\le\sum_l|S(i,l)||T(l,j)|$.
Set $\beta_n=\|B_n\|_\infty$ and
$\beta=\sum_{n<k}\beta_n\le2DL^{-1}a'$.
The one-step estimates give
\[
 \|\mathbf Q(W_n)\|_{\mathrm{row}}\le L+K_{\mathrm{st}}\beta_n,
 \qquad
 \|\mathbf Q(W_n)\|_{\mathrm{col}}\le L^{-2}+K_{\mathrm{st}}\beta_n.
\]
Since $\eta L^k=1$ and $\eta L^{-2k}=\eta^3$, multiplication
yields
\begin{equation}\label{eq:flat-kernel-product-bound}
 \|Q_k(\mathbf U)\|_{\mathrm{row}}
      \le e^{K_{\mathrm{st}}L^{-1}\beta},\qquad
 |Q_k(\mathbf U)(c,b)|\le\eta^3e^{K_{\mathrm{st}}L^2\beta}.
\end{equation}
Choose the positive radius so that
\begin{equation}\label{eq:flat-averaging-radius}
 c_{\mathrm{av}}\le
 \min\left\{\frac{c_\sharp}{4},
          \frac{\log2}{8K_{\mathrm{st}}DL},
          \frac{\log2}{96HL}\right\}.
\end{equation}
Then the exponents in \eqref{eq:flat-kernel-product-bound} are
at most $\tfrac14\log2$. Using $e^s\le1+2s$ for
$0\le s\le\log2$ gives the first two bounds of
\eqref{eq:flat-complex-linear-kernels}, with
$C_{\mathrm{av}}\ge4K_{\mathrm{st}}DL$. Their holomorphic differences
from the value at $A'=0$ vanish there and are bounded by $4$
in the row norm and $4\eta^3$ entrywise. The Schwarz estimate
on each complex line gives the other two bounds when
$C_{\mathrm{av}}\ge4/c_{\mathrm{av}}$.

Expanding the product of the scalar means gives the formula
in \eqref{eq:flat-complex-scalar-transport}: the chains of
nested blocks correspond bijectively to $x\in B^k(y)$.
The one-step contour estimate writes each factor as
$W_n(z,x')=e^{i\zeta_n(z,x')}$, with
$|\zeta_n|\le\kappa\beta_n$. Consequently, if
$z=\kappa\beta$, then
\[
 |\operatorname{Ad}(T_{\mathbf U}(y,x))-1|
   \le e^{2z}-1\le4z\le8\kappa DL^{-1}a'.
\]
Here $z\le1/8$ follows from
$c_{\mathrm{av}}\le c_\sharp/4\le r/(64D)$ and
$\kappa r=1$. This proves (iii), after increasing
$C_{\mathrm{av}}$ if necessary.

\emph{The second perturbation.}
Now use the actual complex bases $W_n$ in
\eqref{eq:flat-straightened-recursion}. For
$t_n=\|a_n\|_\infty/\theta_n$, the one-step estimates give
\[
 t_{n+1}\le t_n
       (1+K_{\mathrm{st}}L^{-1}\beta_n+HL^{-1}\theta_n t_n).
\]
The previous multiplication argument, with the additional
background term, gives
\[
 t_{n+1}\le t_0\exp(K_{\mathrm{st}}L^{-1}\beta+4HL^{-2}t_0)
       \le2t_0
\]
whenever the preceding bounds hold. The radius choice makes
the exponent smaller than $\log2$ and keeps
$\|a_n\|_\infty\le2c_{\mathrm{av}}<r/2$.
This proves the recursion and bound in (ii). Joint holomorphy
and tracelessness follow at each step from the one-step result.
Since $a_n=0$ at $A=0$, the chain rule gives
\eqref{eq:flat-averaging-products}.

For the derivative at general $A$, the same rule gives
\[
 K_A=\eta(\mathcal Q^1)'(W_{k-1},a_{k-1})
                 \cdots(\mathcal Q^1)'(W_0,a_0).
\]
By \eqref{eq:straightening-derivative-column} and
$\sum_{n<k}\|a_n\|_\infty\le4L^{-1}\|A\|_\infty$,
\[
 |K_A(c,b)|
   \le\eta^3\exp\bigl(K_{\mathrm{st}}L^2\beta+24HL\|A\|_\infty\bigr).
\]
The exponent is at most $\tfrac12\log2$. The exponential
estimate proves the first bound of
\eqref{eq:flat-complex-differentiated-kernel} with
$C_2\ge48HL$; it also gives $|K_A|\le2\eta^3$.
For fixed $A'$, the holomorphic map $A\mapsto K_A-K_0$
vanishes at zero and is bounded by $4\eta^3$ on the ball
of radius $c_{\mathrm{av}}$. The same Schwarz estimate gives
the second bound with $C_2\ge4/c_{\mathrm{av}}$.
All the choices above depend only on $L$.
\end{proof}

\section{The constrained system near a flat configuration}
\label{sec:small-constrained-system}

We extend \cref{I-64-24} to small complex perturbations, allowing
independent current and multiplier fields. The perturbation of the
scalar equation for $p$ contains a divergence. Keeping that divergence
in the forcing space permits the estimate without a derivative of $p$.
The other new term is the curvature remainder in the action Hessian.

Keep the torus, separated sequence, inner frame, spaces, pairings,
positions and weights of \cref{sec:trivial-constrained-system}, and
assume throughout that
$\mathfrak m\ge\mathfrak M_\dagger$ and
$\mathsf g\ge L\mathfrak M_\dagger$.
Put $Y=\Omega_0$. Let the ambient region $\mathsf A$ contain all bare
sites within sup distance $\xi$ of $Y$, and let $\mathbf U$ and
$\mathbf J$ be a complex configuration and a $\mathfrak g^c$-valued
one-form on its inner bare bonds. Let $\mathsf j\in\mathcal R$.

\subsection{The operators and the inverse estimate}

For $\mathsf U_\nu(x)=\operatorname{Ad}
\mathbf U(\langle x,x+\xi e_\nu\rangle)$, covariant differences are
\[
 \nabla_\nu f(x)=\xi^{-1}
       \bigl(\mathsf U_\nu(x)f(x+\xi e_\nu)-f(x)\bigr),\qquad
 \nabla_\nu^{\mathsf T}f(x)=\xi^{-1}
       \bigl(\mathsf U_\nu(x-\xi e_\nu)^{-1}f(x-\xi e_\nu)-f(x)\bigr).
\]
Here $f$ is a site field or a component of a bond field, indexed by
its initial point. A term evaluated outside its index set is zero
and uses no transport. In particular,
$(Dv)_\mu=\nabla_\mu v$ on inner bonds and
$D^{\mathsf T}A=\sum_\mu\nabla_\mu^{\mathsf T}A_\mu$.
The Dirichlet Laplacian has the same omission convention and retains
its full diagonal:
\begin{equation}\label{I-63-28}
 \Delta^\circ f(x)=\xi^{-2}
 \left(6f(x)-\sum_{\nu=1}^3
   \bigl[\mathsf U_\nu(x)f(x+\xi e_\nu)
       +\mathsf U_\nu(x-\xi e_\nu)^{-1}f(x-\xi e_\nu)\bigr]\right).
\end{equation}

At a positive-level row $c\in\Gamma_n$, rescale its double block
to level-$n$ units and use the straightened map of
\cref{sec:flat-complex-averaging}. The physical constraint is
$\xi_c^{-1}\mathcal Q_n(\mathbf U;\xi_c A)(c)$.
Its linear and quadratic Taylor terms define $QA(c)$ and
$C^{(2)}(A)(c)$, respectively. Thus, with
$W_j=\mathcal M^j\mathbf U$,
\begin{equation}\label{eq:small-system-physical-averages}
 (QA)(c)=\frac{\xi}{\xi_c}
        \bigl(\mathbf Q(W_{n-1})\cdots\mathbf Q(W_0)A\bigr)(c).
\end{equation}
At a top $y\in\Lambda_n$, $Q'\lambda(y)$ is the ordered scalar
mean of \eqref{eq:flat-complex-scalar-transport}, on the single block $B^n(y)$
in these units. At level zero both averages are the identity and
$C^{(2)}=0$. These definitions are made wherever the local maps
exist; the proposition below establishes their existence under
its smallness hypothesis.

Define symmetric operators by their quadratic forms:
\begin{equation}\label{eq:small-system-quadratic-operators}
 \begin{split}
 (A,\Delta_{\mathsf A}(\mathbf U)A)_\xi
   &=\left.\frac{d^2}{ds^2}\right|_{s=0}
       \xi^{-1}\sum_{p\subset\mathsf A}a_p(e^{is\xi A}\mathbf U),\\
 (A,\Delta^{(2)}[\mathsf j]A)_\xi
   &=2(C^{(2)}(A),\mathsf j)_{\mathbf B}.
 \end{split}
\end{equation}
The first form is on the inner bonds of $\mathsf A$; the second
is on $\mathcal A$. Both are homogeneous quadratic polynomials.
Polarization and the nondegenerate bilinear pairings give unique
symmetric operators. Compress the first to $\mathcal A$ by extending
its arguments by zero, and denote the result by $\Delta_Y(\mathbf U)$.
Set
\[
 \mathsf L=\Delta_Y(\mathbf U)+DD^{\mathsf T},\qquad
 \mathsf L_1=\mathsf L+Q^{\mathsf T}\mathsf m^{-2}Q,
 \qquad \Delta'_a=\Delta^\circ+Q'^{\mathsf T}\mathsf a'Q'.
\]

For a bond field $k$, let
\[
 (\mathsf E_-^{\mathsf T}k)(x)=\sum_{b_-=x}k(b),\qquad
 (\mathsf E^{\mathsf T}k)(x)=\sum_{b_-=x}k(b)
                 +\sum_{b_+=x}\operatorname{Ad}(\mathbf U(b))^{-1}k(b),
\]
where the bonds are positive and inner in $Y$. Write
$v_-(b)=v(b_-)$ and $v_+(b)=\operatorname{Ad}(\mathbf U(b))v(b_+)$,
and put
\begin{equation}\label{I-63-63}
 \begin{split}
 \mathsf m_J(v)&=\tfrac i2[\mathbf J,v_-+v_+],\\
 \mathsf g_0(A,v)&=\tfrac12\mathsf E^{\mathsf T}(i[A,\mathbf J])
                   -\tfrac12\mathsf E_-^{\mathsf T}(i[Dv,\mathbf J]),\\
 \mathsf x(A,v)&=-\tfrac i2[\mathbf J,v_-]
                   -\Delta^{(2)}[\mathsf j](A-Dv).
 \end{split}
\end{equation}
Define $\mathbb K^{[5]}$ from \eqref{I-63-64} by replacing its
operators with those just defined and subtracting
$\mathsf m_J(v)+\Delta^{(2)}[\mathsf j](A-Dv)$ in the first slot
and $\mathsf g_0(A,v)+D^{\mathsf T}\mathsf x(A,v)$ in the fifth.
The same coordinate compressions define $\mathbb K^{[i]}$ for
$i=1,2,4$. In \eqref{I-64-3} and the forcing injection use these
covariant differences; write $|\Psi|_{\mathbf U,y}$ and
$\iota_{\mathbf U}$ for the resulting seminorm and injection.
The forcing seminorm \eqref{I-64-4} is unchanged.

Put $\rho=3L+6$ and
$\mathsf N(y)=\bigcup_{d(y,z)\le\rho}\Delta(z)$.
We call $(\mathbf U,\mathbf J,\mathsf j)$ \emph{$\epsilon$-small
in the trivial gauge} if $\mathbf U(b)=e^{i\xi\mathbf A(b)}$
on the positive bare bonds with both endpoints within sup distance
$\xi$ of $Y$, and, for every $y\in\Lambda_*$,
\begin{equation}\label{eq:small-system-hypothesis}
 \xi_y|\mathbf A(b)|\le\epsilon,\qquad
 \xi_y^2|\mathbf A(b+\xi e_\nu)-\mathbf A(b)|\le\epsilon\xi,
 \qquad |\mathbf J|_{-3;y},\ |\mathsf j|_{-3;y}\le\epsilon.
\end{equation}
The first two inequalities apply when $b_-\in\mathsf N(y)$,
or when $b_-\notin Y$ and
$\operatorname{dist}_\infty(b_-,\mathsf N(y))\le\xi$.
Both endpoints of $b$, and of the shifted bond in the second
inequality, must be in the specified collar.
We call the family \emph{$\epsilon$-small} if a $G$-valued gauge
on that collar makes it $\epsilon$-small in the trivial gauge.
The current transforms by conjugation at $b_-$, and the row field
by conjugation at $c_-$, also when its position is $c_+$.

\begin{proposition}\label{I-65-40}
There are constants $\epsilon_*>0$ and $B_*>0$, depending only
on $L$, such that every $\epsilon_*$-small family above defines
all the local averages and operators, and each
$\mathbb K^{[i]}$, $i\in\{1,2,4,5\}$, is invertible.
Its inverse $\mathbb E^{[i]}$ satisfies
\begin{equation}\label{eq:small-system-decay}
 |\mathbb E^{[i]}\iota_{\mathbf U}g|_{\mathbf U,y}
 \le B_*\sup_z e^{-\delta_{\mathrm{cs}}d(y,z)}|g|_z
 \qquad(g\in\mathcal F^{[i]},\ y\in\Lambda_*).
\end{equation}
In the trivial gauge, with $\mathbb V^{[i]}$ defined in
\eqref{I-65-24},
\begin{equation}\label{I-65-41}
 \begin{split}
 \mathbb E^{[i]}\iota_0
    &=\mathbb E_0^{[i]}\iota_0
          (I-\mathbb V^{[i]}\mathbb E_0^{[i]}\iota_0)^{-1},\\
 \mathbb V^{[i]}\mathbb E_0^{[i]}\iota_0
    &\in\mathfrak O(\delta_{\mathrm{cs}},1/4).
 \end{split}
\end{equation}
The constants are uniform in the torus, lattice spacing and
separated sequence, including its number of layers.
The operators and their inverses depend on $\mathbf U,\mathbf J$
only at bonds with both endpoints in the collar, and on $\mathsf j$.
For $Y=\emptyset$ the assertions use the unique maps on the zero
spaces.
\end{proposition}

We first prove the local averaging and Hessian estimates. Their
combination gives
$\mathbb K_0-\mathbb K=\iota_0\mathbb V$, with $\mathbb V$
small in the forcing norm. The proof of \cref{I-65-40} is then
a geometric series on the forcing space, followed by the change
from $\iota_0$ to $\iota_{\mathbf U}$.

In the local estimates below the family is $\epsilon$-small in
the trivial gauge, with
\[
 0<\epsilon\le\epsilon_0:=\min\{1/8,c_{\mathrm{av}}/(2L)\}.
\]
Constants denoted by $C$ depend only on $L$ and may increase.
Subscript $0$ denotes the trivial configuration; in particular,
$|\Psi|_{0,y}$ is \eqref{I-64-3} with ordinary differences.

\subsection{Local averaging estimates}
\label{sec:small-system-averaging}

We first record the local geometry used below and in the coefficient
estimates. Write $y(x)$ for the position of a bare site $x\in Y$.
The closed bare cube centered at $x$ is contained in $\Delta(y(x))$,
by the nested partitions. Thus, if $x,x'\in Y$ and
$|x-x'|_\infty\le\xi$, their cells meet. Joining their centers to
a point of intersection costs at most $3/2$ in each cell for the
weighted length. The layer gap in \cref{lem:global-cell-geometry}
then gives
\begin{equation}\label{eq:small-system-touching-cells}
 d(y(x),y(x'))\le3,\qquad
 L^{-1}\xi_{y(x)}\le\xi_{y(x')}\le L\xi_{y(x)}.
\end{equation}
In particular, $\Delta(y(x'))\subset\mathsf N(y(x))$.
For bare sites, the separation of $\Omega_1$ from $Y^c$, with
$\mathsf g>3$, also gives
\begin{equation}\label{eq:small-system-boundary-layer}
 x\in Y,\quad x'\notin Y,\quad |x-x'|_\infty\le3\xi
 \quad\Longrightarrow\quad x\in\Lambda_0,\quad \xi_{y(x)}=\xi.
\end{equation}
For $t=0$ the latter conclusion holds throughout $Y$.

The complex bilinear pairing requires a small distinction from a
Hilbert space adjoint. For $X,Y\in\mathfrak g^c$ and a linear map
$T$ on $\mathfrak g^c$,
\begin{equation}\label{eq:small-system-pairing-bounds}
 |\langle X,Y\rangle|\le |X|\,|Y|,\qquad
 \tfrac12|X|\le\sup_{|Y|\le1}|\langle X,Y\rangle|,
 \qquad |T^{\mathsf T}|\le2|T|.
\end{equation}
Indeed, in Pauli coordinates $X=\sum_a x_a\sigma_a$ one has
$|x|_2\le|X|\le\sqrt2|x|_2$, since the eigenvalues of $X^*X$
are $|x|_2^2\pm|i\bar x\mathbin\times x|$.
For $X\ne0$, take
$Y=(\sqrt2|x|_2)^{-1}\sum_a\bar x_a\sigma_a$ for the lower
bound; the upper bound follows from the trace pairing. Applying
these two bounds to $\langle T^{\mathsf T}X,Y\rangle
=\langle X,TY\rangle$ proves the last assertion.

\begin{lemma}\label{lem:small-system-averages}
All the local maps defining $Q,Q'$ and $C^{(2)}$ exist.
There is a constant $C=C(L)$ such that, for $|s|\le4$,
\begin{equation}\label{eq:small-system-average-bounds}
 \begin{split}
 |Tf|_{s-\ell;y}
   &\le C\max_{d(y,z)\le8L}|f|_{s;z},\\
 |\delta Tf|_{s-\ell;y}
   &\le C\epsilon\max_{d(y,z)\le8L}|f|_{s;z},
 \end{split}
\end{equation}
where $\delta T=T(\mathbf U)-T(\mathbf1)$,
$\ell=0$ for $T=Q,Q^{\mathsf T},Q',Q'^{\mathsf T}$, and
$\ell=2$ for
$T=Q^{\mathsf T}\mathsf m^{-2}Q,Q'^{\mathsf T}\mathsf a'Q'$.
The input $f$ belongs to the corresponding site, top, bond or row
space, and all maxima are over $z\in\Lambda_*$.
Moreover,
\begin{equation}\label{eq:small-system-current-bounds}
 \begin{split}
 |\Delta^{(2)}[\mathsf j]u|_{-3;y}
   &\le C\epsilon\max_{d(y,z)\le8L}|u|_{-1;z},\\
 |\mathsf m_J(v)|_{-3;y}
   &\le C\epsilon\max_{d(y,z)\le8L}|v|_{0;z},\\
 |\mathsf g_0(A,v)|_{-4;y}
   &\le C\epsilon\max_{d(y,z)\le8L}
                  \bigl(|A|_{-1;z}+|Dv|_{-1;z}\bigr),\\
 |\mathsf x(A,v)|_{-3;y}
   &\le C\epsilon\max_{d(y,z)\le8L}
           \bigl(|v|_{0;z}+|A|_{-1;z}+|Dv|_{-1;z}\bigr).
 \end{split}
\end{equation}
\end{lemma}

\begin{proof}
Let $c$ be a positive level-$n$ bond with an endpoint
$y_0\in\Lambda_n$, $n\ge1$, and put
$\mathsf Y_c=B^n(c_-)\cup B^n(c_+)$.
Here $c$ may be inward; this will be needed for the scalar mean.
Every site of $\mathsf Y_c$ is within sup distance $3\xi_n/2$
of $y_0$. Since $\mathsf g\ge L^2$, separation puts these sites
in $\Omega_{n-1}$. Since $y_0\notin\Omega_{n+1}$, the next
separation puts them outside $\Omega_{n+2}$, when that set is
nonempty. Consequently their levels are among $n-1,n,n+1$,
and all bonds used in this double block are inner in $Y$.
The own-cell part of \eqref{eq:small-system-hypothesis} gives
\[
 \xi_n\sup_{\mathsf Y_c}|\mathbf A|
 \le L\epsilon\le c_{\mathrm{av}}/2.
\]
After rescaling to level-$n$ units, the bare spacing is
$\eta=\xi/\xi_n$ and
$\mathbf U=e^{i\eta(\xi_n\mathbf A)}$.
Thus \cref{prop:flat-complex-averaging} applies with the flat
reference $\mathbf1$. It constructs the actual iterated averages
and the straightened constraint on this block. In
\eqref{eq:small-system-physical-averages} the physical derivative
is precisely its $Q_n$: the input factor $\xi_n$ cancels the
output factor $\xi_n^{-1}$.

For a top $y\in\Lambda_n$, use the full-frame bond
$c_y=\langle y,y+\xi_ne_1\rangle$ in this argument. It need not
be a row. Every contour used in the scalar mean at $y$ lies in
the nested single block $B^n(y)$; the restrictions of the actual
averages are determined there, successively from finer levels.
The resulting scalar mean is therefore independent of this
auxiliary double block. This proves its definedness as well as
that of the row maps. At level zero the assertions are immediate.

Write $Q(c,b)$ for the counting kernel. The estimates
\eqref{eq:flat-complex-linear-kernels} and
\eqref{eq:flat-complex-scalar-transport} give
\[
 \sum_b|Q(c,b)|\le C,\quad
 |Q(c,b)|\le C(\xi/\xi_c)^3,\quad
 \sum_b|\delta Q(c,b)|\le C\epsilon,\quad
 |\delta Q(c,b)|\le C\epsilon(\xi/\xi_c)^3.
\]
The scalar transport has norm at most $C$ and differs from the
identity by at most $C\epsilon$. The weighted pairings give
\[
 \begin{split}
 (Q^{\mathsf T}\omega)(b)
   &=\sum_c(\xi_c/\xi)^3Q(c,b)^{\mathsf T}\omega(c),\\
 (Q'^{\mathsf T}\sigma)(x)
   &=\operatorname{Ad}(T_{\mathbf U}(y,x))^{\mathsf T}
          \sigma(y),\qquad x\in B^n(y).
 \end{split}
\]
In the second line the averaging factor cancels because
$\xi_y^3L^{-3n}=\xi^3$. Transposes of kernel entries are taken
for $\langle\cdot,\cdot\rangle$ and obey
\eqref{eq:small-system-pairing-bounds}.

Only twelve positive-level rows can contribute at a fixed bare
bond. Indeed, \cref{I-11-27} partitions a reduced row's double
block into cells of levels $n,n-1$. If $b_-\in\mathsf Y_c$,
then $n\in\{n(b_-),n(b_-)+1\}$. At either level the unique
level-$n$ block containing $b_-$ forces $c$ to be one of its
six incident positive bonds. There is at most one additional
identity row at level zero.

For these reduced rows, the cells $\mathcal C(c)$ described in
\cref{sec:coarse-covariance-decay} have sides $\xi_c$ or
$\xi_c/L$, and their centers are within $4L$ of the row position.
Thus any two of their centers are within $8L$. The weighted
transpose cancels the kernel factor $(\xi/\xi_c)^3$, and a
change of position costs at most $L^4$ for $|s|\le4$.
This proves the first four bounds in
\eqref{eq:small-system-average-bounds}. The two scalar maps
in fact preserve the position.
For the bond potential, expand once as
\[
 (Q^{\mathsf T}\mathsf m^{-2}Qf)(b)
 =\sum_c(\xi_c/\xi)^3\xi_c^{-2}Q(c,b)^{\mathsf T}
                       \sum_{b'}Q(c,b')f(b').
\]
Both $b$ and $b'$ belong to the same $\mathsf Y_c$, and
$\xi_c^{-2}\le\xi_{y(b)}^{-2}$. The same count and scale
comparison prove the weight loss of two with radius $8L$.
For the scalar potential, all indices have the same position
and $0<\mathsf a'(y)\le\xi_y^{-2}$.
Finally, expand each potential difference in its two factors,
for example
\[
 \delta(Q^{\mathsf T}\mathsf m^{-2}Q)
  =\delta Q^{\mathsf T}\mathsf m^{-2}Q(\mathbf U)
    +Q(\mathbf1)^{\mathsf T}\mathsf m^{-2}\delta Q.
\]
Each summand keeps its two fine indices in the same double
block and gains the factor $\epsilon$. This proves all of
\eqref{eq:small-system-average-bounds} without enlarging the
support radius by composition.

It remains to track the physical factor in the quadratic term.
For $n\ge1$ set
$f_c(\widetilde A)=\mathcal Q_n(\mathbf U;\widetilde A)(c)$,
and let $K_{\widetilde A}(c,b)$ be its derivative in
$\widetilde A(b)$. Define
$K'_c(u)(c,b)=\left.\frac{d}{ds}\right|_{s=0}K_{su}(c,b)$.
By \eqref{eq:flat-complex-differentiated-kernel},
\[
 |K'_c(u)(c,b)|\le C_2(\xi/\xi_c)^3
                             \sup_{\mathsf Y_c}|u|.
\]
The second derivative of $f_c$ is symmetric, so its quadratic
part is $\frac12\sum_bK'_c(u)(c,b)u(b)$. Rescaling this part
and then polarizing \eqref{eq:small-system-quadratic-operators}
gives
\begin{equation}\label{eq:small-system-multiplier}
 \begin{split}
 C^{(2)}(u)(c)
   &=\tfrac12\xi_c\sum_bK'_c(u)(c,b)u(b),\\
 (\Delta^{(2)}[\mathsf j]u)(b)
   &=\xi^{-3}\sum_{\substack{c\in\Gamma_n,\ n\ge1\\
                               b_-,b_+\in\mathsf Y_c}}
          \xi_c^4 K'_c(u)(c,b)^{\mathsf T}\mathsf j(c).
 \end{split}
\end{equation}
The factor $\xi_c^4$ consists of the physical factor $\xi_c$
and the row-pairing weight $\xi_c^3$. The kernel vanishes when
either endpoint of $b$ is outside the double block; it depends
on $u$ and $\mathbf U$ only on that block. Using
\eqref{eq:small-system-pairing-bounds}, each summand in the
second line has norm at most
$2C_2\xi_c|\mathsf j(c)|\sup_{\mathsf Y_c}|u|$.
For $y=y(b)$,
\[
 \xi_c|\mathsf j(c)|\le\epsilon\xi_y^{-2},\qquad
 \sup_{\mathsf Y_c}|u|
    \le L\xi_y^{-1}\max_{d(y,z)\le8L}|u|_{-1;z}.
\]
The twelve-row count proves the first bound in
\eqref{eq:small-system-current-bounds}. There is no level-zero
contribution because $C^{(2)}$ vanishes there.

Finally, $\epsilon\le1/8$ and the own-cell amplitude bound give
$|\operatorname{Ad}(\mathbf U(b))^{\pm1}|\le e^{2\epsilon}<2$
on inner bonds. Apply $|[X,Y]|\le2|X|\,|Y|$ to
\eqref{I-63-63}. Endpoint positions satisfy
\eqref{eq:small-system-touching-cells}. For example, a current
term at an incoming bond to a site of position $y$ is bounded by
$2\epsilon L^4\xi_y^{-4}|A|_{-1;y(b)}$ after transport.
There are at most three incoming and three outgoing bonds.
The outgoing $Dv$ terms have the same weight; the two terms
of $\mathsf m_J$ have weight $-3$. These observations prove
the bounds for $\mathsf m_J$ and $\mathsf g_0$. The bound for
$\mathsf x$ follows by applying the multiplier estimate to
$u=A-Dv$ and estimating its direct current term. This proves
\eqref{eq:small-system-current-bounds}.
\end{proof}

\subsection{The complex Hessian and its curvature remainder}
\label{sec:complex-hessian}

The Hessian differs from the covariant curl squared by terms containing
the plaquette holonomy minus the identity. Adding the gradient-divergence
term gives the componentwise Dirichlet Laplacian with a local curvature
remainder. The calculation must be made in the ambient region before
compression to the inner bonds of $Y$.

\begin{lemma}\label{I-65-4}
Let $\mathbf U$ be a complex configuration on the inner bonds of $\mathsf A$.
For an inner plaquette $p=p_{\mu\nu}(x)$, $\mu<\nu$, put
$Y_p=\partial\mathbf U(p)$, $y_0(p)=\tfrac12\tr Y_p$, and
$\mathsf U_\mu(x)=\operatorname{Ad}(\mathbf U(\langle x,x+\xi e_\mu\rangle))$.
For a one-form $A$ on the inner bonds of $\mathsf A$, set
\[
\begin{aligned}
 a_1&=A_\mu(x),&
 a_2&=\mathsf U_\mu(x)A_\nu(x+\xi e_\mu),\\
 a_3&=\operatorname{Ad}(Y_p)\mathsf U_\nu(x)A_\mu(x+\xi e_\nu),&
 a_4&=\operatorname{Ad}(Y_p)A_\nu(x).
\end{aligned}
\]
With $\varsigma=(1,1,-1,-1)$ and $S_p(A)=\sum_i\varsigma_i a_i$,
the ambient Hessian satisfies
\begin{equation}\label{I-65-5}
\begin{split}
 (A,\Delta_{\mathsf A}(\mathbf U)A)_\xi
 =\xi^3\sum_p\Bigl\{
 y_0(p)\langle\xi^{-1}S_p(A),\xi^{-1}S_p(A)\rangle
 +i\sum_{i<j}\varsigma_i\varsigma_j
       \langle[a_i,a_j],F_{\mathbf U}(p)\rangle\Bigr\},
\end{split}
\end{equation}
where the sum is over the inner plaquettes of $\mathsf A$.

Suppose $\mathsf A$ contains every bare site within sup distance $\xi$
of $Y$. On every plaquette $p$ containing an inner bond of $Y$, assume
$|Y_p-\mathbf1|\le1$ and
$|\operatorname{Ad}(\mathbf U(b'))^{\pm1}|\le2$ on its four bonds.
For an inner bond $b$ of $Y$, put
\[
 \mathsf f^\circ(b)=\xi^{-2}\max_{p\ni b}|Y_p-\mathbf1|,
 \qquad
 |A|_{(b)}=\max_{\substack{p\ni b,\ b'\in p\\ b'\text{ inner in }Y}}
                      |A(b')|.
\]
Then, with an absolute constant $C$,
\begin{equation}\label{I-65-6}
 \mathsf L=\Delta_Y(\mathbf U)+DD^{\mathsf T}
     =\Delta^\circ+\mathsf Z_{\mathrm{curv}},
 \qquad
 |(\mathsf Z_{\mathrm{curv}}A)(b)|
       \le C\mathsf f^\circ(b)|A|_{(b)}.
\end{equation}
Both $\Delta_Y(\mathbf U)$ and $\mathsf Z_{\mathrm{curv}}$ at $b$ depend only on
the four plaquettes containing $b$ and their bonds.
\end{lemma}

\begin{proof}
Moving the transports to the right in the plaquette product, as in
\eqref{I-10-40}, gives the algebraic identity
\[
 \partial(e^{it\xi A}\mathbf U)(p)
 =e^{it\xi a_1}e^{it\xi a_2}e^{-it\xi a_3}e^{-it\xi a_4}Y_p.
\]
The quadratic term in an ordered product of exponentials is
\[
 \frac12\sum_iX_i^2+\sum_{i<j}X_iX_j
 =\frac12\Bigl(\sum_iX_i\Bigr)^2
    +\frac12\sum_{i<j}[X_i,X_j].
\]
Taking $X_i=it\xi\varsigma_i a_i$, the coefficient of $t^2$ in
$a_p(e^{it\xi A}\mathbf U)$ is therefore
\[
 \frac{\xi^2}{4}\tr\left(
   S_p(A)^2Y_p+\sum_{i<j}\varsigma_i\varsigma_j[a_i,a_j]Y_p
                         \right).
\]
Since $S_p$ is traceless, $S_p^2=\langle S_p,S_p\rangle\mathbf1$.
Also $Y_p=y_0\mathbf1+i\xi^2F_{\mathbf U}(p)$ by \eqref{I-2-6}, and
commutators are traceless. Thus this coefficient is $\xi^4/2$ times
the brace in \eqref{I-65-5}. The action contributes a factor $\xi^{-1}$,
and its second derivative contributes a factor two, proving that identity
for every complex $\mathbf U$.

Write $D^{(1)}$ for the ambient curl,
$(D^{(1)}A)_{\mu\nu}=\nabla_\mu A_\nu-\nabla_\nu A_\mu$.
The signed sum has the exact decomposition
\[
\begin{split}
 \xi^{-1}S_p(A)&=(D^{(1)}A)_{\mu\nu}(x)+\varrho_p(A),\\
 \varrho_p(A)&=\xi^{-1}(I-\operatorname{Ad}(Y_p))
       \bigl[\mathsf U_\nu(x)A_\mu(x+\xi e_\nu)+A_\nu(x)\bigr].
\end{split}
\]
Consequently $\Delta_{\mathsf A}(\mathbf U)
=D^{(1)\mathsf T}D^{(1)}+Z_1$, where
\[
\begin{split}
 (A,Z_1A)_\xi=\xi^3\sum_p\biggl\{
 &(y_0-1)\langle D^{(1)}A,D^{(1)}A\rangle
 +2y_0\langle D^{(1)}A,\varrho_p(A)\rangle\\
 &+y_0\langle\varrho_p(A),\varrho_p(A)\rangle
 +i\sum_{i<j}\varsigma_i\varsigma_j
       \langle[a_i,a_j],F_{\mathbf U}(p)\rangle\biggr\}.
\end{split}
\]
Here the curl in a summand is evaluated at $p$.

Put $\kappa_p=|Y_p-\mathbf1|$. For a determinant-one two-by-two matrix,
$Y_p^{-1}-\mathbf1$ is the adjugate of $Y_p-\mathbf1$ and has the same
operator norm. Hence, when $\kappa_p\le1$,
\[
 |y_0-1|\le\kappa_p,\qquad
 \xi^2|F_{\mathbf U}(p)|\le\kappa_p,\qquad
 |I-\operatorname{Ad}(Y_p^{\pm1})|\le3\kappa_p.
\]
The last inequality follows by writing
$Y_pXY_p^{-1}-X=(Y_p-\mathbf1)XY_p^{-1}
+X(Y_p^{-1}-\mathbf1)$ and using $|Y_p^{\pm1}|\le2$.
If $|A|_{(p)}$ denotes the maximum on the four bonds of $p$, the
transport bounds give
\[
 |a_i|\le C|A|_{(p)},\qquad
 |D^{(1)}A(p)|\le C\xi^{-1}|A|_{(p)},\qquad
 |\varrho_p(A)|\le C\xi^{-1}\kappa_p|A|_{(p)}.
\]
Polarize the form of $Z_1$ against a one-form supported at $b$, of value
$e$ with $|e|\le1$. Each contributing plaquette contains $b$;
there are at most four. Every term is bounded by
$C\xi^{-2}\kappa_p|A|_{(p)}$, using
$\kappa_p^2\le\kappa_p$ and $|[X,X']|\le2|X||X'|$.
The factors $\xi^3$ in the form and the one-bond pairing cancel.
The complex bilinear dual bound in \eqref{eq:small-system-pairing-bounds}
then gives
\[
 |(Z_1A)(b)|\le C\xi^{-2}\max_{p\ni b}\kappa_p
                         \max_{p\ni b}|A|_{(p)}.
\]
Only the touching plaquettes require the bounds used here.

It remains to combine curl and divergence. At a bond whose neighboring
bonds and plaquettes are available in the ambient region, direct
summation gives
\[
 \bigl((D^{(1)\mathsf T}D^{(1)}+DD^{\mathsf T})A\bigr)_\rho
 =\sum_\nu\nabla_\nu^{\mathsf T}\nabla_\nu A_\rho
   +\sum_{\nu\ne\rho}[\nabla_\rho,\nabla_\nu^{\mathsf T}]A_\nu.
\]
Indeed, the off-diagonal terms are
$-\nabla_\nu^{\mathsf T}\nabla_\rho A_\nu
+\nabla_\rho\nabla_\nu^{\mathsf T}A_\nu$; in the remaining direction,
$\nabla_\rho\nabla_\rho^{\mathsf T}
=\nabla_\rho^{\mathsf T}\nabla_\rho$.
For $\rho\ne\nu$, the mixed commutator is
\[
\begin{split}
 ([\nabla_\rho,\nabla_\nu^{\mathsf T}]f)(x)
 =\xi^{-2}\bigl[&\mathsf U_\rho(x)
            \mathsf U_\nu(x+\xi e_\rho-\xi e_\nu)^{-1}\\
       &-\mathsf U_\nu(x-\xi e_\nu)^{-1}
            \mathsf U_\rho(x-\xi e_\nu)\bigr]
                  f(x+\xi e_\rho-\xi e_\nu).
\end{split}
\]
The bracket is
$(\operatorname{Ad}(Y')-I)\mathsf U_\nu(x-\xi e_\nu)^{-1}
\mathsf U_\rho(x-\xi e_\nu)$, where $Y'$ is the holonomy based at $x$
around this plaquette. For its standard holonomy $Y_p$,
$Y'=VY_p^{\pm1}V^{-1}$ with $V$ one of its bond variables.
Thus the bracket has norm at most $C\kappa_p$. This plaquette contains
both the output bond and the bond where $A_\nu$ is evaluated.

Finally extend $A$ by zero from the inner bonds of $Y$. Every plaquette
touching an inner bond $b$ of $Y$, and every bond in the preceding
calculation, lies in $\mathsf A$ by the collar hypothesis.
The ambient divergence of $A$ is supported on the sites of $Y$, so
compressing the ambient $DD^{\mathsf T}$ gives the regional operator.
The ambient sum $\sum_\nu\nabla_\nu^{\mathsf T}\nabla_\nu A_\rho$
compresses to $\Delta^\circ A_\rho$: its diagonal remains
$6\xi^{-2}$, with missing off-diagonal values zero.
Adding the mixed-commutator bound to the bound for $Z_1$ proves
\eqref{I-65-6}. The plaquettes crossing the boundary of $Y$ remain in
this calculation. Plaquettes containing no inner bond of $Y$ contribute
zero after polarization and zero extension. Hence the compressed
operators use only the stated touching plaquettes and are unchanged
if the ambient region is replaced by another region containing the
same collar.
\end{proof}

\subsection{The perturbation in the forcing space}
\label{sec:small-system-defect}

Write $\delta T=T_{\mathbf U}-T_0$ for each operator depending on
the configuration, and put
\[
 \mathsf e_\nu(x)=\xi^{-1}(\mathsf U_\nu(x)-I),\qquad
 \mathsf e'_\nu(x)=\xi^{-1}(\mathsf U_\nu(x)^{-1}-I).
\]
For a site field $f$, or a component of a bond field, define
\[
 \begin{gathered}
 X^{[f]}_\nu(x)=
 \begin{cases}
 \mathsf e_\nu(x)f(x+\xi e_\nu)-\mathsf e'_\nu(x)f(x),
                         &x+\xi e_\nu\text{ is an index},\\
 0,                       &\text{otherwise},
 \end{cases}\\
 (\mathsf cf)(x)=-\xi^{-1}
   \left(\sum_\nu'\mathsf e_\nu(x-\xi e_\nu)
               +\sum_\nu''\mathsf e'_\nu(x)\right)f(x).
 \end{gathered}
\]
The primed sum retains the directions with $x-\xi e_\nu$ in
the index set; the double-primed sum retains those with
$x+\xi e_\nu$ in it. For $p\in\mathcal S$, $X^{[p]}$ is a
one-form on the inner bonds of $Y$.

For $\Psi=(A;v,w,\sigma;p,q,r;\omega)$, define
$\mathbb V^{[5]}\Psi=(f_1,\ldots,f_8;X_5)$ by
\begin{equation}\label{I-65-24}
 \begin{aligned}
 f_1={}&-\delta\Delta^\circ A-\mathsf Z_{\mathrm{curv}}A
          -\delta(Q^{\mathsf T}\mathsf m^{-2}Q)A
          +\delta D(w+p)\\
       &\quad+\mathsf m_J(v)+\Delta^{(2)}[\mathsf j](A-Dv)
          +\delta Q^{\mathsf T}\omega,\\
 f_2={}&-\delta\Delta'_av+\delta D^{\mathsf T}A,\qquad
 f_3=-\delta\Delta'_aw+\delta Q'^{\mathsf T}\sigma,\\
 f_4={}&-\delta Q'v,\\
 f_5={}&-\mathsf cp-\delta(Q'^{\mathsf T}\mathsf a'Q')p
          -\delta Q'^{\mathsf T}r+\mathsf g_0(A,v)
          +\delta D^{\mathsf T}\mathsf x(A,v),\\
 X_5={}&-X^{[p]}+\mathsf x(A,v),\\
 f_6={}&-\delta\Delta'_aq,\qquad
 f_7=\delta Q'q,\qquad f_8=-\delta QA.
 \end{aligned}
\end{equation}
Here $D=D_{\mathbf U}$ and
$\delta\Delta'_a=\delta\Delta^\circ+
\delta(Q'^{\mathsf T}\mathsf a'Q')$.
Delete absent components and terms for the other compressions.
Let $\mathsf n g$ have fifth component
$\delta D^{\mathsf T}X_5$ and all other components zero;
put $\mathsf n=0$ for $i=1,2$.

\begin{lemma}\label{I-65-25}
There is $C_V=C_V(L)\ge1$ such that, for each of the four systems,
\begin{equation}\label{eq:small-system-defect-bound}
 \mathbb K_0^{[i]}-\mathbb K^{[i]}=\iota_0\mathbb V^{[i]},
 \qquad
 |\mathbb V^{[i]}\Psi|_y
       \le C_V\epsilon\max_{d(y,z)\le9L}|\Psi|_{0,z}.
\end{equation}
The two solution seminorms obey
\begin{equation}\label{eq:small-system-norm-comparison}
 (1+8\epsilon)^{-1}|\Psi|_{0,y}
       \le |\Psi|_{\mathbf U,y}
       \le(1+8\epsilon)|\Psi|_{0,y},
\end{equation}
and the injection satisfies
\begin{equation}\label{eq:small-system-injection}
 \iota_{\mathbf U}=\iota_0(I+\mathsf n),\qquad
 |\mathsf ng|_y\le12L^3\epsilon\max_{d(y,z)\le3}|g|_z.
\end{equation}
\end{lemma}

\begin{proof}
If $x$ lies in the cell at $y$, each connection bond based within
sup distance $\xi$ of $x$ is tested by
\eqref{eq:small-system-hypothesis}: an initial point in $Y$ lies
in a touching cell, while an initial point outside $Y$ lies in
the specified exterior collar. The same observation applies to
the shifted pairs below. Since
$\mathsf U_\nu=e^{i\xi\operatorname{ad}\mathbf A_\nu}$ and
$|\operatorname{ad}X|\le2|X|$, the exponential series gives
\begin{equation}\label{eq:small-system-connection-bounds}
 \begin{split}
 |\mathsf e_\nu|,|\mathsf e'_\nu|&\le4|\mathbf A_\nu|,\qquad
 |\mathsf e_\nu+\mathsf e'_\nu|\le8\xi|\mathbf A_\nu|^2,\\
 |\mathsf U_\nu^{\pm1}|&\le1+4\xi|\mathbf A_\nu|,\\
 |\mathsf e_\nu(x')-\mathsf e_\nu(x)|
       &\le4|\mathbf A_\nu(x')-\mathbf A_\nu(x)|.
 \end{split}
\end{equation}
The last bound holds for $\mathsf e'_\nu$ as well.
For example, these inequalities follow from
$|e^T-I|\le2|T|$, $|e^T+e^{-T}-2I|\le2|T|^2$ and
$|e^T-e^{T'}|\le2|T-T'|$ when $|T|,|T'|\le1/4$;
here $2\xi|\mathbf A_\nu|\le2\epsilon\le1/4$.

When $x+\xi e_\nu$ is an index, the same-point identities
\[
 \nabla_\nu f=\mathsf U_\nu\nabla^0_\nu f+\mathsf e_\nu f,
 \qquad
 \nabla^0_\nu f=\mathsf U_\nu^{-1}\nabla_\nu f+\mathsf e'_\nu f
\]
bound a difference of weight $s-1$ by $1+4\epsilon$ times
the other difference, plus $4\epsilon$ times the value of
weight $s$. If that neighbor is absent, the differences agree.
Each differentiated field in \eqref{I-64-3} occurs with its
value, proving \eqref{eq:small-system-norm-comparison}.
Also
\[
 (\delta Df)(b)=\mathsf e_\mu(b_-)f(b_+),\qquad
 (\delta D^{\mathsf T}A)(x)
   =\sum_\mu'\mathsf e'_\mu(x-\xi e_\mu)A_\mu(x-\xi e_\mu),
\]
where the sum retains the inner incoming bonds. These operators
lower the weight by one with bound $C\epsilon$ on neighboring
cells. At weight $-3$, the three terms and adjacent-scale ratio
give $12L^3\epsilon$, proving \eqref{eq:small-system-injection}.

For the Laplacian, subtraction in \eqref{I-63-28} gives
\[
 \delta\Delta^\circ f(x)=-\xi^{-1}\sum_\nu
    \bigl[\mathsf e_\nu(x)f(x+\xi e_\nu)
        +\mathsf e'_\nu(x-\xi e_\nu)f(x-\xi e_\nu)\bigr],
\]
with missing terms omitted. At an index with all six neighbors,
insert $f(x+\xi e_\nu)=f(x)+\xi\nabla^0_\nu f(x)$ and its
backward counterpart to obtain
\begin{equation}\label{I-65-27}
 \begin{split}
 \delta\Delta^\circ f(x)
   ={}&-\sum_\nu\bigl[\mathsf e_\nu(x)\nabla^0_\nu f(x)
        -\mathsf e'_\nu(x-\xi e_\nu)\nabla^0_\nu f(x-\xi e_\nu)\bigr]\\
     &-\xi^{-1}\sum_\nu
           [\mathsf e_\nu(x)+\mathsf e'_\nu(x-\xi e_\nu)]f(x).
 \end{split}
\end{equation}
At every index, expansion of the two terms of
$\nabla_\nu^{0\mathsf T}X^{[f]}_\nu$ instead gives
\begin{equation}\label{I-65-28}
 \delta\Delta^\circ f=\sum_\nu
       \nabla_\nu^{0\mathsf T}X^{[f]}_\nu+\mathsf cf.
\end{equation}
The primed sums in $\mathsf c$ cancel exactly the diagonal
terms introduced by this expansion, including at missing neighbors.

The zeroth-order coefficient in \eqref{I-65-27} is small at
the physical scale: split it into
$\mathsf e_\nu(x)+\mathsf e'_\nu(x)$ and the difference of
$\mathsf e'_\nu$ at the two neighboring bonds. By
\eqref{eq:small-system-connection-bounds} and the first-difference
hypothesis, its norm after multiplication by $\xi^{-1}$ is
at most $C\epsilon\xi_y^{-2}$. The coefficient of
$\mathsf cf$ has the same bound. If a neighbor is missing,
there is a site outside $Y$ within $2\xi$ of $x$; hence
$\xi_y=\xi$ by \eqref{eq:small-system-boundary-layer}.
Each remaining coefficient is then bounded directly by
$4\epsilon\xi_y^{-2}$. Thus, for $s\in\{-2,-1,0\}$ and
with maxima over $d(y,z)\le3$,
\begin{equation}\label{eq:small-system-laplacian-bounds}
 \begin{split}
 |\delta\Delta^\circ f|_{s-2;y}
   &\le C\epsilon\max_z
        (|f|_{s;z}+|\nabla^0 f|_{s-1;z}),\\
 |X^{[f]}|_{s-1;y}+|\mathsf cf|_{s-2;y}
   &\le C\epsilon\max_z|f|_{s;z}.
 \end{split}
\end{equation}
In the first line, derivatives are needed only where both
indices belong to the field's index set. The second line
applies to $p$ without a derivative estimate.

For a plaquette touching a bond based at $x\in\Delta(y)$,
write its holonomy as $e^{X_1}e^{X_2}e^{X_3}e^{X_4}$, with the
four signed exponents $\pm i\xi\mathbf A$ in boundary order.
The parallel pairs and \eqref{eq:small-system-hypothesis} give
\[
 \sum_j|X_j|\le4\epsilon\xi/\xi_y,\qquad
 \left|\sum_jX_j\right|\le2\epsilon\xi^2/\xi_y^2.
\]
The exponential remainder bound therefore yields
\[
 |\partial\mathbf U(p)-I|
 \le(2\epsilon+16\epsilon^2)\xi^2/\xi_y^2
 \le4\epsilon\xi^2/\xi_y^2.
\]
The transported coefficients have norm at most two by
\eqref{eq:small-system-connection-bounds}. Applying
\cref{I-65-4}, and comparing the touching cells, gives
\[
 |\mathsf Z_{\mathrm{curv}}A|_{-3;y}
        \le C\epsilon\max_{d(y,z)\le3}|A|_{-1;z}.
\]
In particular $\delta\mathsf L=\delta\Delta^\circ+
\mathsf Z_{\mathrm{curv}}$, since the curvature term vanishes
at the identity.

Subtracting the system at $\mathbf U$ from the trivial system
now gives \eqref{I-65-24}. In its fifth slot, the relevant identity is
\[
 \begin{split}
 -\delta\Delta'_ap-\delta Q'^{\mathsf T}r
       +\mathsf g_0+D^{\mathsf T}\mathsf x
   ={}&-\mathsf cp-\delta(Q'^{\mathsf T}\mathsf a'Q')p
       -\delta Q'^{\mathsf T}r+\mathsf g_0
       +\delta D^{\mathsf T}\mathsf x\\
     &+D_0^{\mathsf T}(-X^{[p]}+\mathsf x).
 \end{split}
\]
This proves the operator identity in
\eqref{eq:small-system-defect-bound} for all four compressions.

The differentiated fields $A,v,w,q$ have weights $-1,0,-2,0$;
the first line of \eqref{eq:small-system-laplacian-bounds}
controls their Laplacian differences in weights $-3,-2,-4,-2$.
For $p$, the second line controls $\mathsf cp$ in weight $-4$
and $X^{[p]}$ in weight $-3$. The averaging, potential and
current estimates of \cref{lem:small-system-averages}, together
with the first-order bounds above, give precisely the other
weights in \eqref{I-64-4}.

For the composed term $\delta D^{\mathsf T}\mathsf x$ at a site
$x$, every contributing multiplier block contains both endpoints
of an incoming bond, hence $x$ itself. Its input positions and
the position of $x$ are in the same row block; their distances
are at most $8L$. Its weight-$-3$ bound is multiplied by a
coefficient of size at most $4\epsilon/\xi_y$, giving weight
$-4$. The direct current term uses neighboring cells. Thus
all inputs in the defect lie within $9L$ of the output, with
bound $C\epsilon$; here $\epsilon^2\le\epsilon$.
There are only finitely many component terms and the row
incidence is bounded independently of the layer count.
Their maximum fixes $C_V$ and proves the remaining estimate.
\end{proof}

\subsection{Proof of the inverse estimate}
\label{sec:small-system-inverse}

\begin{proof}[Proof of \cref{I-65-40}]
First work in the trivial gauge. Enlarge $B_{\mathrm{cs}}$ to
$\max\{1,B_{\mathrm{cs}}\}$ and choose $0<\epsilon_*\le\epsilon_0$
so that
\[
 C_V\epsilon_*B_{\mathrm{cs}}e^{9L\delta_{\mathrm{cs}}}\le\tfrac14,
 \qquad 12L^3\epsilon_*e^{3\delta_{\mathrm{cs}}}\le1.
\]
All local operators are defined by
\cref{lem:small-system-averages}. Suppress $i$ in the proof.
The finite-range and composition bounds following
\eqref{eq:trivial-system-sup-decay}, together with
\cref{I-64-24,I-65-25}, give
\[
 T:=\mathbb V\mathbb E_0\iota_0
       \in\mathfrak O(\delta_{\mathrm{cs}},1/4).
\]
Consequently $\sum_{k\ge0}T^k$ converges in the supremum norm.
Applying the same decay bound to each term and summing gives
$(I-T)^{-1}\in\mathfrak O(\delta_{\mathrm{cs}},4/3)$.

If $\mathbb K\Psi=0$, then
$\Psi=\mathbb E_0\iota_0\mathbb V\Psi$. Hence
$g=\mathbb V\Psi$ satisfies $g=Tg$, so $g=0$ and $\Psi=0$.
Finite dimensionality makes $\mathbb K$ invertible. Moreover,
\[
 \mathbb K\mathbb E_0\iota_0g'=
       \iota_0(I-T)g',
\]
which proves \eqref{I-65-41} by setting $g'=(I-T)^{-1}g$.
The injection correction in \eqref{eq:small-system-injection}
has decay constant at most one. Combining it with the norm
comparison gives \eqref{eq:small-system-decay}, since
$2\cdot(4/3)B_{\mathrm{cs}}\cdot2\le16B_{\mathrm{cs}}$.
Take $B_*=16B_{\mathrm{cs}}$.

For a general small family, let $u$ be its trivializing gauge
on the bare collar and continue it arbitrarily to $\mathsf A$.
Conjugation acts at sites, at the initial points of bonds, and
at the initial endpoints of rows. The unitary covariance of
the straightening maps gives covariance of their derivatives
and quadratic Taylor coefficients. In particular, the defined
local averages in the trivial gauge, transported back by $u$,
are the actual averages of $\mathbf U$: the contour products,
their logarithms and the successive averages obey the same
conjugation identities. This establishes their domains as well
as their values.

Gauge invariance of the ambient action, differentiated twice,
gives covariance of its Hessian. The invariant trace pairing
preserves transposes and polarization; the differences and the
commutators in \eqref{I-63-63} obey the same covariance.
Thus, writing $\mathcal U=\operatorname{Ad}(u)$ on the respective
component spaces,
\[
 \mathbb K[\mathbf U^u,\operatorname{Ad}(u)\mathbf J,
                   \operatorname{Ad}(u)\mathsf j]
       =\mathcal U\mathbb K[\mathbf U,\mathbf J,\mathsf j]\mathcal U^{-1},
 \qquad \iota_{\mathbf U^u}\mathcal U=\mathcal U\iota_{\mathbf U}.
\]
Unitarity preserves the local solution and forcing seminorms,
so the bound in the trivial gauge proves the stated bound.

The compressed Hessian uses only plaquettes touching inner bonds
of $Y$, all of whose bonds have both endpoints in the collar.
The averages and every other coefficient use inner bonds of
$Y$ and the row field. Hence the continuation of $u$ has no
effect, and $\mathbb K$, and therefore its inverse, have the
asserted dependence.
\end{proof}

\section{Local regular gauges}\label{sec:local-regular-gauges}

Small curvature gives a gauge in which the bond variables are close
to the identity. To control the first differences of the resulting
potential uniformly in the spacing, we also use a bound on the
divergence of the plaquette variables. The construction smooths a
comb section, proves that the smoothed section has no zero, and
normalizes it to obtain the gauge.

\begin{theorem}\label{I-28-75}
Fix odd $L\ge13$. There are absolute constants $c_*>0$ and $B_*\ge1$ with the
following property. Use level-$k$ units, where $1\le k\le K$
and $\eta=L^{-k}$. Let $Q$ be a closed cube of positive integer
side $l$, formed from cubes of the level-$k$ partition, and let
$\widetilde Q$ be the concentric cube of side $l+4$, less than
half the torus period. Suppose $\alpha>0$, $l\alpha\le c_*$ and that
$U$ is a real configuration on the bare lattice, with
\[
 |\partial U(p)-\mathbf1|<\alpha\eta^2,
 \qquad |(D_U^{\mathsf T}\partial U)(b)|<\alpha\eta^2
\]
for every plaquette with a corner in $\widetilde Q$ and every
bond with an endpoint there, respectively.
There are a real gauge $u$ on the sites of $\widetilde Q$ and
a real potential $A_0$ on the bonds with both endpoints in
$\widetilde Q$ such that $U^{u^{-1}}=e^{i\eta A_0}$ there and
\[
 |A_0|,\quad |\nabla^\eta A_0|,\quad
 |D_{\mathbf1}^{\mathsf T}D_{\mathbf1}A_0|<B_*l\alpha
 \quad\text{at every positive bond with an endpoint in }Q.
\]
Here $\nabla^\eta$ is the ordinary first difference, and
$D_{\mathbf1}^{\mathsf T}D_{\mathbf1}$ is ordinary curl
followed by its transpose.
\end{theorem}

\begin{proof}
\emph{Reduction to a box.}
Lift $\widetilde Q$ to the covering lattice. Its side is less
than half the period, so the lift identifies its sites, bonds
and plaquettes with
\[
 \Omega=o+\eta\{0,\ldots,N-1\}^3,\qquad s=N\eta=l+4.
\]
The faces lie halfway between bare sites because $L$ is odd.
We have $0<\eta\le1/13$ and $5\le s\le5l$; put $\hat\mu=\eta e_\mu$.
Write $\mathfrak B$ for the positive inner bonds and $d(x)$ for
the minimum coordinate distance to an extreme site face.
Set $\Omega^{(j)}=\{x:d(x)\ge j\eta\}$.
A bond is full if all four plaquettes containing it lie in the box.
At such a bond the divergence $D_U^{\mathsf T}\partial U$
computed in the box is the given divergence. This is the
matrix-valued expression of \eqref{I-2-5}, with the
antisymmetric extension of the plaquette variables, rather
than the field-strength current $J(U)$.

We prove that $s\alpha\le c_\square$ gives a gauge with
$|A_0|<B_\square s\alpha$ on $\mathfrak B$ and
$|\nabla^\eta A_0|,|D_{\mathbf1}^{\mathsf T}D_{\mathbf1}A_0|
<B_\square s\alpha$ at base sites of depth
at least $3/2$. This suffices with $c_*=c_\square/5$ and
$B_*=5B_\square$: sites of $Q$ have depth at least $2$, and
a positive bond with an endpoint in $Q$ has base depth at least
$2-\eta\ge3/2$. The translated bonds in its first differences
also lie in the box.

\emph{Smoothing the comb section.}
Use the real algebra
$\mathbb H=\{a\mathbf1+iX:a\in\mathbb R, X\in\mathfrak g\}$,
with scalar part $\operatorname{Re}(a\mathbf1+iX)=a$.
For $q\in\mathbb H$,
$q^*q=|q|^2\mathbf1$, its unit sphere is $G$, and
$\langle p,q\rangle_{\mathbb H}=\operatorname{Re}(p^*q)$
is a real inner product with norm $|q|$.
These identities follow from $X^2=|X|^2\mathbf1$ for
$X\in\mathfrak g$. In particular, multiplication by an element
of $G$ is an isometry.

Let $\phi_0(x)$ be transport from $x$ to $o$, lowering the
coordinates in the order $1,2,3$. Moving a bond through these
paths inserts at most $2(N-1)$ plaquettes, as in the comb
calculation of \cref{I-10-51}. Hence
\[
 |\phi_0(x)^{-1}U(x,x+\hat\mu)\phi_0(x+\hat\mu)-\mathbf1|
 \le2s\alpha\eta.
\]
On quaternion-valued sections use the left action
\[
 (\mathcal D\phi)_\mu(x)
 =\eta^{-1}\bigl(U(x,x+\hat\mu)\phi(x+\hat\mu)-\phi(x)\bigr).
\]
This is a different representation from the adjoint action in
$D_U$. Its adjoint is taken with the real inner product above
and sums over the sites and positive inner bonds of the box.
In particular,
\[
 (\mathcal D^*\mathcal D\phi)(x)
 =\eta^{-2}\sum_{x'\sim x, x'\in\Omega}
          \bigl(\phi(x)-U(x,x')\phi(x')\bigr).
\]
The positive operator $\mathcal D^*\mathcal D+1$ is invertible.
Let
\begin{equation}\label{I-28-32}
 (\mathcal D^*\mathcal D+1)\phi=\phi_0,
 \qquad \psi=\mathcal D\phi,
 \qquad \varepsilon_0=\|\mathcal D\phi_0\|_\infty\le2s\alpha.
\end{equation}
At a maximum of $|\phi|$, the inner product of
$\mathcal D^*\mathcal D\phi$ with $\phi$ is nonnegative.
Equation \eqref{I-28-32} therefore gives $\|\phi\|_\infty\le1$.

\emph{The equation for $\psi$ and its boundary terms.}
Let $n(b)$ count the missing transverse neighbours of a bond
$b=(x,\mu)$; it is zero exactly when the bond is full.
For its $\mu$ component define
\[
 \begin{split}
 (\mathcal L\psi)_\mu(x)=\eta^{-2}\Bigl[&
 \sum_{\nu\ne\mu}\sum_{\pm: x\pm\hat\nu\in\Omega}
    \bigl(\psi_\mu(x)-U(x,x\pm\hat\nu)\psi_\mu(x\pm\hat\nu)\bigr)\\
 &+2\psi_\mu(x)-U(x,x+\hat\mu)\psi_\mu(x+\hat\mu)
                 -U(x,x-\hat\mu)\psi_\mu(x-\hat\mu)\Bigr].
 \end{split}
\]
Bond values outside $\mathfrak B$, and every product containing
such a value, are zero. Thus the transverse
directions have a Neumann boundary condition and the longitudinal
direction has a Dirichlet boundary condition.

Write $Y_{\mu\nu}(x)$ for the plaquette holonomy in the order
$\mu,\nu,-\mu,-\nu$, minus $\mathbf1$, and put
$Y^-_{\mu\nu}(x)=\operatorname{Ad}(U(x,x-\hat\nu))
Y_{\mu\nu}(x-\hat\nu)$ when that plaquette is present.
Commuting the two differences in \eqref{I-28-32} gives
\begin{equation}\label{eq:gauge-smoothed-difference}
 (\mathcal L+1)\psi=\mathcal D\phi_0+M-W.
\end{equation}
The terms on the right are
\[
 \begin{split}
 M_\mu(x)=\eta^{-3}\Bigl[&
 \sum_{\nu\ne\mu: x+\hat\nu\in\Omega}
 Y_{\mu\nu}(x)
 \bigl(\phi(x)+\eta\psi_\nu(x)
             +\eta U(x,x+\hat\nu)\psi_\mu(x+\hat\nu)\bigr)\\
 &-\sum_{\nu\ne\mu: x-\hat\nu\in\Omega}
 Y^-_{\mu\nu}(x)\bigl(\phi(x)+\eta\psi_\mu(x)\bigr)\Bigr],\\
 W_\mu(x)=\eta^{-2}&\sum_{\nu\ne\mu: x-\hat\nu\in\Omega}
 \bigl((\mathbf1+Y^-_{\mu\nu}(x))^{-1}-\mathbf1\bigr)\\
 &\hspace{12mm}\cdot U(x,x-\hat\nu)U(x-\hat\nu,x-\hat\nu+\hat\mu)
                   \psi_\nu(x-\hat\nu+\hat\mu).
 \end{split}
\]
For completeness, the forward commutation uses
\[
 \begin{split}
 U(x,x+\hat\mu)\psi_\nu(x+\hat\mu)-\psi_\nu(x)
 ={}&U(x,x+\hat\nu)\psi_\mu(x+\hat\nu)-\psi_\mu(x)\\
 &+\eta^{-1}Y_{\mu\nu}(x)
       U(x,x+\hat\nu)U(x+\hat\nu,x+\hat\nu+\hat\mu)
                          \phi(x+\hat\nu+\hat\mu).
 \end{split}
\]
The backward commutation is this identity at $x-\hat\nu$,
transported to $x$; its reversed plaquette contributes
$(\mathbf1+Y^-_{\mu\nu})^{-1}-\mathbf1$. Together with
the longitudinal second difference these are
\eqref{eq:gauge-smoothed-difference}, including absent neighbours.

Put $\Sigma=\|\psi\|_\infty$. The curvature bound gives
$|Y|,|Y^-|,|(\mathbf1+Y^-)^{-1}-\mathbf1|\le\alpha\eta^2$.
The terms containing $\psi$ in $M$ and $W$ are consequently
bounded by $6\alpha\Sigma$ and $2\alpha\Sigma$.
For a full bond, the remaining coefficient satisfies
\[
 \sum_{\nu\ne\mu}(Y_{\mu\nu}-Y^-_{\mu\nu})
 =\eta(D_U^{\mathsf T}\partial U)_\mu(x)
   +2\sum_{\nu<\mu}\bigl(c_{\nu\mu}(x)
                   -c_{\nu\mu}(x-\hat\nu)\bigr)\mathbf1,
\]
where $c_{\nu\mu}=\tfrac12\tr\partial U(p_{\nu\mu})$.
This follows by replacing an oppositely oriented plaquette
$g^{-1}$ by $2c\mathbf1-g$. Since
$1-c=\tfrac12|g-\mathbf1|^2$, the norm of this coefficient
is at most $\alpha\eta^3+2\alpha^2\eta^4$.
At a nonfull bond there are at most three remaining plaquettes,
so the bound $3\alpha\eta^2n(b)$ suffices. We have proved
\begin{equation}\label{eq:gauge-commutator-bound}
 |M_\mu(x)|\le\alpha(1+2\alpha\eta)
             +3\alpha\eta^{-1}n(b)+6\alpha\Sigma,
 \qquad |W_\mu(x)|\le2\alpha\Sigma.
\end{equation}

\emph{A uniform bound at the boundary.}
Let $\mathcal L_0$ be $\mathcal L$ with all transports replaced
by $1$, acting on the sites carrying the $\mu$ bonds. Its
maximum principle and the inequality
\[
 \langle\psi_\mu,\mathcal L\psi_\mu\rangle_{\mathbb H}
       \ge |\psi_\mu|\mathcal L_0|\psi_\mu|
\]
give, for $\alpha\le1/16$,
\[
 (\mathcal L_0+1)|\psi_\mu|
 \le\varepsilon_0+2\alpha+8\alpha\Sigma
                       +3\alpha\eta^{-1}n(b).
\]
At a zero of $\psi_\mu$ the same inequality holds directly
because the off-diagonal coefficients of $\mathcal L_0$ are
nonpositive.

The boundary forcing is controlled by a one-dimensional barrier.
Choose $\kappa>0$ so that $2\sinh(\kappa\eta/2)=\eta$, and set
\[
 \beta(j)=
 \frac{\cosh(\kappa\eta(j-(N-1)/2))}{\sinh(\kappa\eta N/2)}.
\]
For $0\le j<N$, $0<\beta(j)\le2$: indeed $\kappa\ge4/5$
and $N\eta\ge5$. Its second difference plus $\beta$ vanishes
at interior sites, whereas at either endpoint the Neumann
second difference plus $\beta$ equals $\eta^{-1}$.
For $\nu\ne\mu$ put $\beta_\nu(x)=\beta((x_\nu-o_\nu)/\eta)$.
The longitudinal Dirichlet terms are nonnegative, so
\[
 (\mathcal L_0+1)\sum_{\nu\ne\mu}\beta_\nu
                  \ge\eta^{-1}n(b).
\]
Comparison with
$\varepsilon_0+2\alpha+8\alpha\Sigma
 +3\alpha\sum_{\nu\ne\mu}\beta_\nu$ gives
$\Sigma\le\varepsilon_0+14\alpha+8\alpha\Sigma$.
Consequently,
\begin{equation}\label{I-28-48}
 \Sigma\le2\varepsilon_0+28\alpha\le10\epsilon,
 \qquad \epsilon=s\alpha.
\end{equation}

\emph{The section does not vanish.}
In the gauge $\phi_0^{-1}$ write
$\chi=\phi_0^{-1}\phi-\mathbf1$,
$\widetilde\psi=\phi_0^{-1}\psi$ and
$\widetilde U=U^{\phi_0^{-1}}$. Extend $\widetilde\psi$ by zero
outside $\mathfrak B$, with the same zero-product convention. Equation \eqref{I-28-32} becomes
\[
 \chi(x)=\eta^{-1}\sum_\nu
  \bigl(\widetilde\psi_\nu(x)-\widetilde\psi_\nu(x-\hat\nu)\bigr)
 -\eta^{-1}\sum_\nu
  (\widetilde U(x,x-\hat\nu)-\mathbf1)
                     \widetilde\psi_\nu(x-\hat\nu).
\]
The last term is bounded by $3\varepsilon_0\Sigma$, and
$|\chi(x+\hat\nu)-\chi(x)|\le\eta(\Sigma+\varepsilon_0)$;
the latter follows from the definition of $\psi$ and
$|\phi|\le1$.
For any site $x_0$ choose a product $K$ of $k$ consecutive
sites in each direction, containing $x_0$, where
$k=\lceil\eta^{-1}\rceil$. It fits in $\Omega$ because
$k\eta<1+\eta\le2\le s$.
Summing the divergence over $K$ leaves two faces in each
direction, hence
\[
 \left|k^{-3}\sum_K\chi\right|
 \le\frac{6\Sigma}{k\eta}+3\varepsilon_0\Sigma
 \le6\Sigma+3\varepsilon_0\Sigma.
\]
Every site of $K$ is at path length at most $3(k-1)\eta<3$
from $x_0$. Comparing its value with the mean yields
\begin{equation}\label{I-28-50}
 \|\phi-\phi_0\|_\infty
 \le9\Sigma+3\varepsilon_0+3\varepsilon_0\Sigma\le\tfrac12
\end{equation}
when $\epsilon\le1/1000$. Thus
$\rho=|\phi|$ lies in $[1/2,1]$ and $u=\phi/\rho$ is
$G$-valued.

\emph{Ordinary first differences.}
Put
\[
 V=U^{u^{-1}},\qquad \omega_\mu=\phi^{-1}\psi_\mu,
 \qquad
 \psi'(x,x')=\eta^{-1}u(x)^{-1}
                   \bigl(U(x,x')\phi(x')-\phi(x)\bigr).
\]
Then $\omega$ is quaternion-valued,
$|\omega|\le2\Sigma\le20\epsilon$, and $|\psi'|\le\Sigma$.
The definition gives
\[
 \rho(x')V(x,x')=\rho(x)\mathbf1+\eta\psi'(x,x'),
 \qquad
 V(x,x+\hat\mu)=
       \frac{\mathbf1+\eta\omega_\mu(x)}
            {|\mathbf1+\eta\omega_\mu(x)|}.
\]
At $x\in\Omega^{(2)}$ all six neighbours occur. Substitute
$\psi_\mu=\phi\omega_\mu$ into $\mathcal L\psi$ and use
\eqref{I-28-32} to obtain
\begin{equation}\label{I-28-55}
 \begin{split}
 \rho\Delta^\eta\omega_\mu={}&u^{-1}(\mathcal L\psi)_\mu
 +(\rho\mathbf1-u^{-1}\phi_0)\omega_\mu\\
 &+\eta^{-1}\sum_{x'\sim x}
        \psi'(x,x')\bigl(\omega_\mu(x')-\omega_\mu(x)\bigr).
 \end{split}
\end{equation}
Here $\Delta^\eta$ is the positive ordinary Laplacian.
The coefficient of $\omega_\mu(x)$ arises from
$u^{-1}\mathcal D^*\mathcal D\phi=u^{-1}\phi_0-\rho\mathbf1$;
the remaining differences give the displayed first-order term.
Using \eqref{eq:gauge-smoothed-difference},
\eqref{eq:gauge-commutator-bound} with $n(b)=0$, and
\eqref{I-28-48} and \eqref{I-28-50}, we obtain
\begin{equation}\label{eq:gauge-omega-inequality}
 |\Delta^\eta\omega_\mu(x)|
 \le55\epsilon+120\epsilon
            \max_{x'\sim x}\eta^{-1}
                    |\omega_\mu(x')-\omega_\mu(x)|.
\end{equation}

We give the absorption argument, so that the estimate is uniform
as $\eta$ tends to zero. If a quaternion-valued function $f$ on
$\Omega^{(1)}$ satisfies $|f|\le w$ and
$|\Delta^\eta f|\le a+\sigma\max_{x'\sim x}
\eta^{-1}|f(x')-f(x)|$ on $\Omega^{(2)}$, let
\[
 M_f=\max_{x\in\Omega^{(2)},\nu}
         \min\{1,d(x)\}|\nabla^\eta_\nu f(x)|.
\]
Choose a maximizing site $x_1$ and put $d_1=\min\{1,d(x_1)\}$.
If $d_1<64\eta$, then $M_f\le128w$.
Otherwise the cube of radius $r=d_1/4\ge16\eta$ about $x_1$
and its neighbours have depth at least $d_1/2$.
All differences in its equation are bounded by $2M_f/d_1$.
Apply \cref{I-13-47}(b), with $F=0$, to the real scalar
function $\langle e,f\rangle_{\mathbb H}$ for each unit
quaternion $e$. The normalized $L^2$ norm is at most $6w$.
Taking the supremum over $e$ gives, with $c_I\ge1$,
\[
 M_f\le4c_I(6w+a/16)+\tfrac12c_I\sigma M_f.
\]
Thus, if $c_I\sigma\le1$,
$M_f\le\max\{128w,8c_I(6w+a/16)\}$.
Apply this to \eqref{eq:gauge-omega-inequality} with
$w=20\epsilon$, $a=55\epsilon$, and $\sigma=120\epsilon$.
For $\epsilon\le(120c_I)^{-1}$ it follows that
\[
 |\nabla^\eta_\nu\omega_\mu(x)|\le2560c_I\epsilon
 \qquad(d(x)\ge1).
\]

Finally, for $q=a\mathbf1+iX\in\mathbb H$ with $a>0$, set
\[
 \Theta(q)=a^{-1}X\int_0^1
                \frac{dt}{1+t^2|X|^2/a^2}.
\]
The elementary formula for $e^{iX}$ shows that
$e^{i\Theta(q)}=q/|q|$ and $\Theta(\mathbf1)=0$.
The map is smooth on $a>0$ and has an absolute Lipschitz
constant $c_\Theta\ge1$ on $|q-\mathbf1|\le1/2$.
Set $A_0=\eta^{-1}\Theta(\mathbf1+\eta\omega)$ and
$C_M=2560c_I$. Then $e^{i\eta A_0}=V$,
$|A_0|\le20c_\Theta\epsilon$, and
$|\nabla^\eta A_0|\le c_\Theta C_M\epsilon$ at depth
at least $1$. The arguments of $\Theta$ lie in its stated
ball because $20\eta\epsilon\le1/2$.

We also need the ordinary curl-divergence of this potential.
For $a=(a_1,a_2,a_3,a_4)\in\mathfrak g^4$, put
\[
 P(a)=e^{ia_1}e^{ia_2}e^{-ia_3}e^{-ia_4}
       -\mathbf1-i(a_1+a_2-a_3-a_4).
\]
Since $P(0)=P'(0)=0$, the maximum of $\|P''\|$ on
$\max_i|a_i|\le1$ gives an absolute $c_P\ge1$ such that
\[
 |P(a)|\le c_P\tau^2,\qquad
 |P(a)-P(a')|\le c_P\tau\max_i|a_i-a_i'|
 \quad(\max_i|a_i|,\max_i|a_i'|\le\tau\le1).
\]
Let $b=(x,\mu)$ be full, and let $a_0$ and $a_1$ be the
maximum of $|A|$ and of the ordinary first differences
between bonds in its four plaquettes. For $V=e^{i\eta A}$,
expand each plaquette as
$\partial V=\mathbf1+i\eta^2D_{\mathbf1}A+P$,
with the orientation signs in the linear term. In the
divergence, the identity terms cancel. At each of the two
directions $\nu\ne\mu$, the three errors are bounded by
\[
 4\eta^3a_0a_1,\qquad c_P\eta^3a_0a_1,\qquad
 2c_P\eta^3a_0^3.
\]
Indeed, $|D_{\mathbf1}A|\le2a_1$,
$|\operatorname{Ad}(V(x,x-\hat\nu))-1|\le2\eta a_0$,
and the arguments of $P$ at the two translated plaquettes
differ by at most $\eta^2a_1$. Dividing by $\eta^3$ yields
\begin{equation}\label{eq:gauge-ordinary-curl-divergence}
 |(D_{\mathbf1}^{\mathsf T}D_{\mathbf1}A)(b)|
 \le\eta^{-2}|(D_V^{\mathsf T}\partial V)(b)|
       +2(4+c_P)a_0a_1+4c_Pa_0^3,
 \qquad \eta a_0\le1.
\end{equation}
The full plaquette divergence transforms by conjugation
under a real gauge, so its norm is unchanged.

If $d(x)\ge3/2$, every initial point of the four-plaquette
star has depth at least $3/2-\eta\ge1$. Thus for the
constructed potential we may use
$a_0\le20c_\Theta\epsilon\le1$ and
$a_1\le c_\Theta C_M\epsilon$ in
\eqref{eq:gauge-ordinary-curl-divergence}. Since
$V=U^{u^{-1}}$, its first term is at most $\alpha\le\epsilon$.
We obtain
\[
 |D_{\mathbf1}^{\mathsf T}D_{\mathbf1}A_0|
 \le\bigl[1+2(4+c_P)c_\Theta C_M+80c_Pc_\Theta\bigr]
       \epsilon.
\]
Choose
\[
 c_\square\le
 \min\{1/1000,(120c_I)^{-1},(20c_\Theta)^{-1}\},
\]
and choose $B_\square$ larger than $1$, $20c_\Theta$,
$c_\Theta C_M$ and
$1+2(4+c_P)c_\Theta C_M+80c_Pc_\Theta$.
The box estimates follow. With
$c_*=c_\square/5$ and $B_*=5B_\square$, they give all
three asserted bounds on $Q$.
\end{proof}

\section{The inverse for regular configurations}
\label{sec:regular-system-inverse}

Ba{\l}aban studied regularity and decay of lattice propagators
in a background gauge field \cite{B-CMP99b}.
Local gauges reduce the constrained system to the small-field
inverse of \cref{I-65-40}. Near the boundary of the ambient region,
the local system instead has only the bare scale; its two massive
operators can be inverted directly. A partition of unity combines
these local inverses. The error is small because its output lies
away from the support of the input cutoff, so local exponential
decay supplies a factor exponentially small in the block size.

Fix odd $L\ge13$. Keep the physical scales $\xi=L^{-t}$, $\xi_n=L^n\xi$ and the
spaces, pairings, positions and norms of
\cref{sec:small-constrained-system}. Here $(\Omega_n)_{n\le t}$
is separated with block parameter $\mathfrak M$ and relative gap
parameter $R_1\ge L$, both powers of $L$, and the ambient region
is $\mathsf A=\Omega_0$. We assume
\[
 \ell\ge2L\mathfrak M\xi_t.
\]
The current $\mathbf J$ and the row multiplier $\mathsf j$ remain
independent variables. Empty index sets retain the conventions
already used. Throughout this section $C$ denotes a constant
depending only on $L$, whose value may increase between estimates.

\begin{definition}\label{def:natural-regular-pair}
Let $\mathsf Q_n^\natural$ consist of the embedded compact cubes
which are unions of $\mathfrak M$-cubes at level $n$, have side
at most $12\mathfrak M\xi_n$, and lie in
$\mathcal L_n\cup\mathcal L_{n+1}$. There is no requirement that
such a cube meet $\mathcal L_n$.
For $\alpha_0,\alpha_1,\gamma_0>0$, the class
$\mathfrak P_t^\natural((\Omega_n);\alpha_0,\alpha_1,\gamma_0)$
consists of pairs $(\mathbf U,\mathbf J)$ on the inner bare bonds
of $\Omega_0$ for which
$\mathbf U=e^{i\xi A'}U$, with $U$ $G$-valued and $A'$ a
$\mathfrak g^c$-valued one-form, satisfying the following conditions
for every $0\le n\le t$.
\begin{enumerate}[label=\textup{(\roman*)}]
\item If the corners of a bare plaquette $p$ lie in $\Omega_0$
and its base point lies in $\mathcal L_n$, then
\[
 |\partial U(p)-I|<\alpha_0\xi^2\xi_n^{-2},\qquad
 |\partial\mathbf U(p)-I|<\alpha_0\xi^2\xi_n^{-2}.
\]
\item For each $\mathsf q\in\mathsf Q_n^\natural$, there are a
cubic box $\mathsf q\subset\square\subset\Omega_0$, a $G$-valued
gauge $u$ on its bare sites, and a real one-form $A$ on its inner
bonds such that
\[
 U^u=e^{i\xi A},\qquad
 \xi_n|A|\le\alpha_0,\qquad
 \xi_n^2|\nabla^\xi A|\le\alpha_0.
\]
The difference in this display is the ordinary first difference.
\item For each positive inner bond $b$ with
$b_-\in\mathcal L_n$ and each defined covariant difference
based at $x\in\mathcal L_n$,
\[
 \xi_n|A'(b)|<\alpha_1,\qquad
 \xi_n^2|(\nabla^U_\mu A'_\nu)(x)|<\alpha_1.
\]
The second condition is imposed exactly when the three bonds
used by that difference belong to the inner bonds of $\Omega_0$.
\item At every positive inner bond with $b_-\in\mathcal L_n$,
$|\mathbf J(b)|<\gamma_0\xi_n^{-3}$.
\end{enumerate}
\end{definition}

We use the operators of \cref{sec:small-constrained-system} with
this ambient region. Their averaging maps will be shown to exist
in \cref{lem:regular-local-inverse}. We suppress $[i]$ for a fixed
$i\in\{1,2,4,5\}$, and write $|\Psi|_y$ and $\iota$ for the
covariant solution seminorm and forcing injection. The forcing
space is $\mathcal F=\mathcal F^{[i]}$.

\begin{theorem}\label{I-65-83}
There are a power $\mathfrak M_*$ of $L$ and positive constants
$a_{\mathrm{reg}},\delta_{\mathrm{reg}},B_{\mathrm{reg}}$,
depending only on $L$, with the following property. Suppose
$\mathfrak M\ge\mathfrak M_*$,
$(\mathbf U,\mathbf J)\in
\mathfrak P_t^\natural((\Omega_n);\alpha_0,\alpha_1,\gamma_0)$,
and
\[
 0<\alpha_0,\alpha_1,\gamma_0\le a_{\mathrm{reg}},\qquad
 |\mathsf j|_{-3;y}\le\gamma'\le a_{\mathrm{reg}}
 \quad(y\in\Lambda_*),\qquad \gamma'\ge0.
\]
Then each of the four systems is defined and invertible. For
every $g\in\mathcal F$ and $y\in\Lambda_*$,
\begin{equation}\label{eq:regular-system-decay}
 \mathbb E=\mathbb K^{-1},
 \qquad
 |\mathbb E\iota g|_y\le B_{\mathrm{reg}}
       \sup_z e^{-\delta_{\mathrm{reg}}d(y,z)}|g|_z.
\end{equation}
The constants are uniform in the volume and number of layers.

\end{theorem}

The proof occupies the rest of the section. The local inverses
decay at twice the rate used in \eqref{eq:regular-system-decay}.
One half of this decay is retained in the global estimate; the
other half is spent on the distance between the input cutoff and
the support of its commutator.

\subsection{Pieces and cutoffs}
\label{subsec:regular-pieces}

Put $s_n=\mathfrak M\xi_n$, and let $\mathcal C_n$ be the
partition into cubes of side $s_n$. For a cube $\kappa\in\mathcal C_n$,
write $\kappa^{[r]}$ for its concentric enlargement of side
$(1+2r)s_n$. The enlargements used below are embedded cubes, since
$\ell\ge2L\mathfrak M\xi_t$ and $L\ge13$.
Fix the local parameters
\[
 \mathfrak m_\flat=\mathfrak M_\dagger,
 \qquad
 \mathsf g_\flat=L\mathfrak M_\dagger.
\]
In choosing $\mathfrak M_*$ we require
\begin{equation}\label{eq:piece-margin-conditions}
 \mathfrak M\ge L\mathfrak m_\flat,\qquad
 3\mathsf g_\flat+1\le\tfrac12\mathfrak M,\qquad
 \mathfrak M>30L,\qquad
 \mathfrak M\ge16(\rho+1)L^2,
 \qquad \rho=3L+6.
\end{equation}
These conditions involve only $L$. Further increases of
$\mathfrak M_*$ will control the cutoff derivatives and the remainder.

\begin{definition}\label{def:regular-pieces}
For $1\le n\le t$, divide the cubes of $\mathcal C_n$ contained
in $\mathcal L_n$ according as their distance from $\Omega_n^c$
is less than $2s_n$ or at least $2s_n$; denote their unions by
$\mathrm{Buf}_n$ and $\mathrm{Deep}_n$, respectively. Put
$\mathrm{Deep}_0=\mathcal L_0$ and $\mathrm{Buf}_{t+1}=\emptyset$.
The pieces of size $n$ are
\[
 \mathcal P_n
 =\{\kappa\in\mathcal C_n:
       \kappa\subset\mathrm{Deep}_n\cup\mathrm{Buf}_{n+1}\},
 \qquad \mathcal P=\bigcup_{n=0}^t\mathcal P_n.
\]
A piece has a unique size $n(\kappa)$, since the cube sides are
distinct. A size-zero piece with
$\kappa^{[2]}\cap\Omega_1=\emptyset$ is a boundary piece;
all other pieces are regular.

Fix a nonincreasing $\chi\in C^\infty(\mathbb R)$ equal to $1$
on $(-\infty,0]$ and to $0$ on $[1,\infty)$. If $z$ and $s$
are the center and side of $\kappa$, set
\[
 h_{\kappa;a,b}(x)
 =\prod_{\mu=1}^3
  \chi\!\left(\frac{|x_\mu-z_\mu|-(1/2+a)s}{(b-a)s}\right),
 \qquad 0\le a<b,
\]
using least coordinate distances on the torus, and put
\begin{equation}\label{eq:piece-cutoffs}
 \widetilde h_\kappa=h_{\kappa;1/3,2/3},\qquad
 \phi_\kappa=h_{\kappa;0,1/12},\qquad
 \theta_\kappa=\frac{\phi_\kappa}{\sum_{\kappa'}\phi_{\kappa'}}.
\end{equation}
The last expression is taken on the open set where its denominator
is positive; it contains $\Omega_0$ by \cref{lem:piece-geometry}.

For a regular piece $\kappa\in\mathcal P_n$, let
$e_m=2\mathsf g_\flat\xi_m$ and define
\begin{equation}\label{eq:piece-local-sequence}
 \begin{split}
 \Omega^\kappa_{n+1}&=\Omega_{n+1}\cap\kappa^{[1]},\\
 \Omega^\kappa_n&=\kappa^{[1]}+[-e_n,e_n]^3,\\
 \Omega^\kappa_m&=\Omega^\kappa_{m+1}+[-e_m,e_m]^3
                         \quad(0\le m<n).
 \end{split}
\end{equation}
Here $\Omega_{t+1}=\emptyset$; if $n=t$ the first line is
only this empty auxiliary member. For a boundary piece put
$\Omega^\kappa_0=\kappa^{[1]}\cap\Omega_0$ and
$\Omega^\kappa_m=\emptyset$ for $m\ge1$.
Intersections in these definitions are cube-wise intersections in
$\mathcal C_n$ (in $\mathcal C_0$ for a boundary piece).
Complete each local sequence by empty members to the same top
index $t$, retaining the spacing $\xi=L^{-t}$.
Local objects carry a superscript or subscript $\kappa$.
\end{definition}

For the local system define the row multiplier purely by restriction:
\begin{equation}\label{eq:piece-row-multiplier}
 \mathsf j^\kappa(c)=
 \begin{cases}
 \mathsf j(c),&\text{the midpoint of }c\text{ lies in }\kappa^{[5/6]},\\
 0,&\text{otherwise}.
 \end{cases}
\end{equation}
The first case uses the common rows in \cref{lem:piece-geometry}.
Let $\mathbb K_\kappa$ be the system for this local sequence,
with row multiplier $\mathsf j^\kappa$ and the ambient Hessian
on $\Omega_0$. Its inverse $\mathbb E_\kappa$ is constructed in
\cref{lem:regular-local-inverse}.

\subsection{Metric and geometry of the pieces}
\label{subsec:piece-geometry}

The local regions are smaller than the original regions. To compare
their distances, we must account for the extra length charged to
segments lying on a common face of two closed layers.

\begin{lemma}\label{lem:layer-density-distance}
For a separated sequence with powers $\mathfrak m,R\ge L$ of $L$
as block and relative gap parameters, put
\[
 i(x)=\max\{m\ge1:x\in\Omega_m\},\qquad
 g(x)=\xi_{i(x)}^{-1},
\]
where $i(x)=0$ when $x\notin\Omega_1$. Then
\begin{equation}\label{eq:layer-density-distance}
 d_{\mathbf B}(x,x')
   =\inf_{\Gamma\ni x,x'}\int_\Gamma g\,d\ell^1,
\end{equation}
where the infimum is over connected finite unions of closed
segments anywhere in the torus. Consequently, for two such
sequences with the same spacings and
$\widehat\Omega_m\subset\Omega_m$ for every $m$, their distances
satisfy $\widehat d\ge d$.
\end{lemma}

\begin{proof}
We first relate the integral to the outer length in \eqref{I-1-10}.
Cut any finite segment family into pieces of sup length less than
$\ell/4$. Lifts identify the intersection of two pieces with a point
or a Euclidean segment. Subdividing at intersection endpoints and
removing repeated overlaps gives injectively parametrized segments
$p_i+s v_i$, $0\le s\le1$, which meet only in finitely many points.
The trace of a cube on each segment is a finite union of intervals:
only finitely many translates of a lifted cube meet it, and each
translate imposes linear inequalities on $s$. The same holds for
regions, layers and their complements.

For finite unions of parameter intervals $J_i\subset[0,1]$, with
endpoints included or omitted, set
$Z=\bigcup_i\{p_i+s v_i:s\in J_i\}$. We claim that
\begin{equation}\label{eq:segment-outer-length}
 \operatorname{length}(Z)=\sum_i|v_i|_1|J_i|.
\end{equation}
The segments over the closures of the intervals give the upper
bound. For the reverse bound, fix a coordinate $\mu$. For a segment
with $v_\mu\ne0$, let $n(p,v,J;s)$ count the parameters in $J$
whose $\mu$-coordinate equals $s$ modulo $\ell$. By lifting that
coordinate and dividing its image into periods,
\[
 \int_0^\ell n(p,v,J;s)\,ds=|v_\mu|\,|J|.
\]
Except at finitely many coordinate values, the number of points of
$Z$ in the slice $x_\mu=s$ is the sum of these counts for its
nonoverlapping segments. Every finite segment cover
$(p'_k,v'_k)$ of $Z$ supplies at least as many parameter hits in
that slice. Exclude also the finitely many slices containing a
covering segment constant in this coordinate. Integration gives
$\sum_i|v_{i,\mu}||J_i|\le\sum_k|v'_{k,\mu}|$.
Sum over $\mu$ and take the infimum over covers to prove
\eqref{eq:segment-outer-length}.

For the finite-valued densities used here, define the integral on
$\Gamma$ by summing $|v_i|_1\int_0^1g(p_i+s v_i)\,ds$ over
such a nonoverlapping representation. A common subdivision and
an affine change of parameter show that this value is independent
of the representation. Formula \eqref{eq:segment-outer-length}
and \eqref{I-9-17} give
\[
 |\Gamma|_{\mathbf B}=\int_\Gamma w\,d\ell^1,
 \qquad
 w=\sum_{m=1}^t\xi_m^{-1}\mathbf1_{\mathcal L_m}
           +\xi_0^{-1}\mathbf1_{T_\ell\setminus\Omega_1}.
\]
For a representation with overlaps the corresponding sum is an
upper bound. Integrals of nonnegative densities are monotone under
inclusion, by retaining the appropriate subdivided edges.

At every point, $g$ is one of the nonnegative terms defining $w$.
Moreover $w=g$ except on the boundaries of $\Omega_m$, $m\ge2$:
an additional layer contribution at a point of maximum depth
lies in $\mathcal L_{m-1}\cap\Omega_m\subset\partial\Omega_m$.
Thus $w\ge g$ gives the inequality $\ge$ in
\eqref{eq:layer-density-distance}.

For the other inequality, a connected finite segment union
containing distinct $x,x'$ contains an injective polygonal path
between them. Indeed, subdivide at endpoints and intersections;
the resulting finite edge graph is connected, and a chain with
the fewest edges has no repeated vertex. Its integral of $g$ is
no larger. Cut this path at the endpoints of all its parameter
intervals in $q\cap\Omega_m^c$, where $m\ge2$ and $q$ is a
defining $\mathfrak m$-cube of $\Omega_m$. Each resulting segment
either meets all these faces only at its endpoints or lies in one
such set $q\cap\Omega_m^c$.

For a segment $[u,u+v]$ of the latter kind, separation gives
$q\cap\Omega_{m+1}=\emptyset$: every point of $q$ is within
$\mathfrak m\xi_m<R\mathfrak m\xi_m$ of $\Omega_m^c$.
Hence $g=\xi_m^{-1}$ throughout $q$, and $w=\xi_m^{-1}$ in
its interior. Lift the segment to the Euclidean cube of side
$\mathfrak m\xi_m$ centered at $z$. This is possible in one lift
because different translates of that cube are positively separated.
Put $a=z-u$ and $a'=z-u-v$ in this lift. For $0<\sigma<1$
replace the segment by the three successive segments with vectors
$\sigma a$, $(1-\sigma)v$, and $-\sigma a'$. Except at the two
original endpoints they lie in the interior of $q$, and their
weighted cost is at most
\[
 \xi_m^{-1}\bigl(\sigma|a|_1+(1-\sigma)|v|_1
                          +\sigma|a'|_1\bigr)
 \le \xi_m^{-1}|v|_1+3\sigma\mathfrak m.
\]
The other path segments have equal $w$- and $g$-costs. If there
are $r$ face segments, the modified connected union therefore has
weighted length at most $\int_\Gamma g\,d\ell^1+3\sigma\mathfrak m r$.
Letting $\sigma\downarrow0$ proves the reverse inequality. Points
have zero length, so $x=x'$ causes no exception; if $\Omega_0$
is empty, both densities are the constant $\xi_0^{-1}$.
Finally, region inclusion decreases the maximum-depth index and
increases $g$, proving the distance comparison.
\end{proof}



\begin{lemma}\label{lem:sharp-row-neighborhood}
For an inner-frame row $c\in\Gamma_n$ let
\[
 \Delta_c=\Delta^n(c_-)\cup\Delta^n(c_+),\qquad
 \mathsf Y_c=B^n(c_-)\cup B^n(c_+).
\]
The row box $\Delta_c$ is a union of top cells of side $\xi_n$ or $\xi_n/L$.
The positions and base points of $c$, of the bare sites in
$\mathsf Y_c$, and of bonds with both endpoints there belong to
this union. Any two points of $\Delta_c$ can be joined inside it
at weighted cost at most $3L+3$. In particular,
$\Delta_c\subset\mathsf N(y)$ for every top cell of the union.

For a sequence with gap parameter $\mathsf g\ge L^2$, if
$d(y,y')\le\rho=3L+6$, then
\begin{equation}\label{eq:sharp-neighborhood-size}
 L^{-1}\xi_y\le\xi_{y'}\le L\xi_y,\qquad
 |y-y'|_\infty\le\rho L\xi_y,\qquad
 \mathsf N(y)\subset\{x:|x-y|_\infty\le(\rho+1)L\xi_y\}.
\end{equation}
If a real $C^1$ function $h$ satisfies
$\xi_y|\partial_\mu h|\le\vartheta^{-1}$ on every cell
$\Delta(y)$, then, for every connected finite segment union
$\Gamma\subset\Omega_0$ and $p,p'\in\Gamma$,
\begin{equation}\label{eq:cutoff-weighted-lipschitz}
 |h(p)-h(p')|\le\vartheta^{-1}|\Gamma|_{\mathbf B}.
\end{equation}
\end{lemma}

\begin{proof}
By \cref{I-11-27}\textup{(ii)}, at least one endpoint of $c$
belongs to $\Lambda_n$. Its level-$n$ cube is its top cell;
the other is either another such cell or a union of level-$(n-1)$
cells. There is at most one outer endpoint, and at level zero
both endpoints of an inner row belong to $\Lambda_0$.
Bare sites have their positions in these cells, and the row box
is convex, so it contains all the relevant midpoints.
For $p\in\Delta^n(e)$, $p'\in\Delta^n(e')$, with
$e,e'\in\{c_-,c_+\}$, use
$[p,e]\cup[e,e']\cup[e',p']$. Each endpoint cube receives
at most $3\xi_n$ of $\ell^1$ length. A segment issuing from a
lattice site meets cell faces only finitely often, since the faces
have half-integral bare coordinates and $L$ is odd. The weights
are therefore $\xi_n^{-1}$, or $L\xi_n^{-1}$ in the one possible
outer cube. This proves the cost bound and the assertion about
$\mathsf N(y)$.

Since $\rho<L^2\le\mathsf g$, the layer-gap estimate
\eqref{eq:global-layer-gap}, with $\mathsf g$ in place of
$R_1\mathfrak M$, gives the scale comparison in
\eqref{eq:sharp-neighborhood-size}. For $y\in\Lambda_j$, a
connecting set of weighted length less than $\mathsf g$ cannot
meet $\Omega_{j+2}$. Otherwise the distance from $\Omega_{j+2}$
takes every value between $0$ and $\mathsf g\xi_{j+1}$ on that
set. The corresponding points lie in $\mathcal L_{j+1}$, and
the interval-image argument in \cref{subsec:multiscale-sum-proof}
charges length at least $\mathsf g\xi_{j+1}$ there, a contradiction.
All its points consequently count with weight at least
$(L\xi_y)^{-1}$. Taking the infimum gives
$|y-y'|_\infty\le L\xi_y d(y,y')$. Add the radius
$\xi_{y'}/2\le L\xi_y/2$ of the endpoint cell to obtain the
last inclusion.

Finally, extract a polygonal path in $\Gamma$ as in the preceding
proof. At a point of maximum depth $m$, the derivative bound is
$\vartheta^{-1}\xi_m^{-1}$, while the closed-layer density is at
least $\xi_m^{-1}$. Integrating the derivative of $h$ on each
segment and summing proves \eqref{eq:cutoff-weighted-lipschitz}.
\end{proof}

\begin{lemma}\label{lem:piece-geometry}
There are constants $N_{\mathrm G},C_{\mathrm G}\ge1$, depending
only on $L$, with the following properties.
\begin{enumerate}[label=(\alph*)]
\item The pieces cover $\Omega_0$ and have disjoint interiors.
For a regular $\kappa\in\mathcal P_n$,
\begin{equation}\label{eq:piece-layer-margin}
 \kappa^{[2]}\subset\mathcal L_n\cup\mathcal L_{n+1},
 \qquad
 \operatorname{dist}_\infty(\kappa^{[1]},\Omega_{n+2})
                              >9s_{n+1}.
\end{equation}
For a boundary piece, $\kappa^{[2]}\cap\Omega_0\subset\mathcal L_0$,
and the same distance inequality holds. Empty final regions are
understood throughout. If $\kappa'$ meets $\kappa^{[1]}$, their
sizes differ by at most one. Each of the following numbers is
at most $N_{\mathrm G}$: pieces meeting $\kappa^{[1]}$,
enlargements $\kappa'^{[1]}$ containing a fixed point, and pieces
with $\kappa'^{[1]}\cap\kappa^{[1]}\cap\Omega_0\ne\emptyset$.

\item A regular local sequence is
$(\mathfrak m_\flat,\mathsf g_\flat)$-separated; its region
$\Omega^\kappa_0$ and its bare collar lie in $\kappa^{[3/2]}
\subset\Omega_0$. A boundary local sequence is
$(\mathfrak M,R_1\mathfrak M)$-separated. For every piece,
\begin{equation}\label{eq:piece-region-comparison}
 \Omega^\kappa_m\subset\Omega_m,\qquad
 d^\kappa\ge d.
\end{equation}
The local and global indices with basepoint in $\kappa^{[5/6]}$
coincide, with their positions, cells, kinds and weights.

\item The functions in \eqref{eq:piece-cutoffs} are smooth on
an open neighborhood of $\Omega_0$, and
\[
 \sum_\kappa\theta_\kappa=1,\qquad
 0\le\theta_\kappa\le1,\qquad
 \operatorname{supp}_{\Omega_0}\theta_\kappa
                              \subset\kappa^{[1/12]}.
\]
For either the global or the local sequence, for every top $y$
and $h\in\{\theta_\kappa,\widetilde h_\kappa\}$,
\begin{equation}\label{eq:piece-cutoff-derivatives}
 |h|\le1,\qquad
 \xi_y|\partial h|\le C_{\mathrm G}/\mathfrak M,\qquad
 \xi_y^2|\partial^2h|\le C_{\mathrm G}/\mathfrak M
                  \quad\hbox{on }\mathsf N(y).
\end{equation}
Here the maxima range over coordinate derivatives, and the
neighborhood is taken in the sequence concerned.

\item Define global sets of tops
\begin{equation}\label{eq:piece-support-sets}
 \begin{split}
 S'_\kappa&=\{y:\mathsf N(y)\cap\kappa^{[1/12]}
                                        \ne\emptyset\},\\
 S_\kappa&=\{y:\widetilde h_\kappa
                   \text{ is not constant on }\mathsf N(y)\}.
 \end{split}
\end{equation}
Both lie in $\kappa^{[5/6]}$, and
\begin{equation}\label{eq:piece-support-separation}
 d(y,y')\ge D_0:=\frac{\mathfrak M}{8L}
                 \qquad(y\in S_\kappa,\ y'\in S'_\kappa).
\end{equation}
If $\mathsf N^\kappa(y)$ meets $\kappa^{[2/3]}$, then its
tops and cells are global ones and
\begin{equation}\label{eq:piece-neighborhood-comparison}
 \mathsf N^\kappa(y)\subset\kappa^{[19/24]},\qquad
 \mathsf N^\kappa(y)\subset\mathsf N(y).
\end{equation}
All local indices at these positions have basepoints in
$\kappa^{[5/6]}$. In particular a local neighborhood meeting
$\kappa^{[1/12]}$, or on which $\widetilde h_\kappa$ is not
constant, has its top in $S'_\kappa$, respectively $S_\kappa$.
\end{enumerate}
\end{lemma}

\begin{proof}
\emph{Cover and margins.}
By the nested partitions and \cref{I-9-20}, each layer is the
union of the cubes of $\mathcal C_n$ contained in it. Splitting
these cubes into $\mathrm{Deep}_n$ and $\mathrm{Buf}_n$, and
then subdividing the latter at level $n-1$, proves the covering
assertion. Each piece lies either in $\mathrm{Deep}_n$ or in
$\mathrm{Buf}_{n+1}$, exclusively. In the second case its unique
parent in $\mathcal C_{n+1}$ lies in $\mathrm{Buf}_{n+1}$.
Pieces at different sizes have disjoint interiors: either they
lie in different layers, or the smaller is in a buffer parent
and the larger in a different deep cube at that parent's level.

For a regular piece there are three cases. If $n\ge1$ and
$\kappa\subset\mathrm{Deep}_n$, its $2s_n$ margin gives
$\kappa^{[2]}\subset\Omega_n$, and
$\kappa\subset\Omega_{n+1}^c$ separates it from $\Omega_{n+2}$
by more than $13s_{n+1}$. If
$\kappa\subset\mathrm{Buf}_{n+1}$, its distance from
$\Omega_n^c$ exceeds $13s_n$. Every point of its parent is
within $3s_{n+1}$ of $\Omega_{n+1}^c$, so its distance from
$\Omega_{n+2}$ exceeds $10s_{n+1}$.
Finally, if $n=0$ and $\kappa\subset\mathcal L_0$ is regular,
then $\kappa^{[2]}$ meets $\Omega_1$. Every point of $\kappa$
is within $3s_0$ of $\Omega_1$, and hence more than $10s_0$
from $\Omega_0^c$; its distance from $\Omega_2$ exceeds
$13s_1$. These facts prove
$\kappa^{[2]}\subset\Omega_n$, exclude $\Omega_{n+2}$ there,
and give \eqref{eq:piece-layer-margin}. Dividing its cubes according
to membership in $\Omega_{n+1}$ proves the asserted two-layer
containment. For a boundary piece the definition excludes
$\Omega_1$ from $\kappa^{[2]}$ and separation gives the same
margin from $\Omega_2$.

If a piece $\kappa'$ meets $\kappa^{[1]}$, a point of intersection
is interior to $\kappa^{[2]}$, so its layer has index $n$ or
$n+1$. A piece in layer $m$ has size $m$ or $m-1$; thus the
sizes differ by at most one. Counting cubes of side $s_n/L$
that meet $\kappa^{[1]}$ bounds the first number in (a).
At a point of $\Omega_0$, a piece containing that point shows
that all enlargements through it have at most three sizes;
at each size their number is bounded by a fixed coordinate count.
Outside $\Omega_0$ only size-zero boundary pieces can occur.
If two enlargements meet in $\Omega_0$, their sizes differ by
at most two, and all the points of $\kappa'$ lie in the image
of a cube concentric with $\kappa$ of side $(4L^2+3)s_n$.
Their disjoint interiors have volumes at least $(s_n/L^2)^3$.
The volume ratio bounds the last number by a constant depending
only on $L$, and one such constant $N_{\mathrm G}$ covers all
three counts.

\emph{Local sequences and common indices.}
Since $\mathfrak M\ge L\mathfrak m_\flat$ and both are powers
of $L$, $\kappa^{[1]}$ is a union of
$\mathfrak m_\flat$-cubes at every level at most $n+1$.
Each $e_m$ is an integer multiple of $\mathfrak m_\flat\xi_m$.
Thus \eqref{eq:piece-local-sequence} gives nested regions of the
required grids, with successive gaps at least
$e_m=2\mathsf g_\flat\xi_m>\mathsf g_\flat\xi_m$.
Completing by empty members imposes no further gap condition.
Moreover
\[
 \sum_{m=0}^n e_m\le3\mathsf g_\flat\xi_n
                         \le\tfrac12s_n-\xi.
\]
The region $\Omega^\kappa_0$ and its bare collar therefore lie
in $\kappa^{[3/2]}\subset\kappa^{[2]}\subset\Omega_n$.
This proves local separation, the collar assertion and
$\Omega^\kappa_m\subset\Omega_m$. A boundary local sequence
has only its level-zero region, a union of $\mathfrak M$-cubes,
and is separated with the stated parameters. The distance
comparison follows from \cref{lem:layer-density-distance}.

At every interior point of $\kappa^{[1]}$ the lists of memberships
in $\Omega_m$ and $\Omega^\kappa_m$ agree. For $m\le n$ this
follows from the containments just proved, for $m=n+1$ from the
cube-wise intersection, and for larger $m$ from the exclusion of
$\Omega_{n+2}$. The boundary case uses the same argument with
only $m=0$. Cube-wise intersection agrees with point-set
membership in this interior, since the grids are nested.

Any global row box meeting $\kappa^{[1]}$ has level at most
$n+1$. Indeed, a row of level $m\ge n+2$ has an endpoint in
$\Omega_m$, and its box lies within $2\xi_m$ of that endpoint,
whereas \eqref{eq:piece-layer-margin} and separation put
$\kappa^{[1]}$ more than $9s_{m-1}$ from $\Omega_m$.
This is impossible since $2\xi_m=2L\xi_{m-1}<s_{m-1}$.
The same level bound holds for a global cell meeting
$\kappa^{[1]}$, and for every local row or top by construction.
For an index with basepoint in $\kappa^{[5/6]}$, its sites,
endpoints, position and cell thus lie within
$\xi/2+3\xi_{n+1}/2<s_n/6$ of that basepoint, inside
$\kappa^{[1]}$. Their region memberships agree. This determines
the same bare-site, inner-bond, top and row indices, with the same
positions, cells, kinds and weights. In particular no point on a
partition face is being treated as a bare site.

\emph{The partition and its derivatives.}
Each $\phi_\kappa$ is $1$ on its piece and vanishes outside
$\kappa^{[1/12]}$, so the cover gives
$\Phi:=\sum\phi_\kappa\ge1$ on $\Omega_0$.
The cutoff factors are constant near the coordinate cusp and
the period cut; elsewhere they are smooth. Hence $\Phi$ is
smooth and positive on an open neighborhood of $\Omega_0$,
and the partition and support assertions follow.

For either sequence, let $x\in\mathsf N(y)$ lie in the cell
of a nearby top of level $j$. Then $\xi_j\ge\xi_y/L$ by
\cref{lem:sharp-row-neighborhood}, and the maximum global depth
at $x$ is at least $j$, also for a local cell because
$\Omega^\kappa_j\subset\Omega_j$.
If a cutoff belonging to $\kappa'\in\mathcal P_{n'}$ is not
constant near $x$, then $x\in\kappa'^{[2]}$ and
$x\notin\Omega_{n'+2}$. Consequently
\[
 \frac{\xi_y}{s_{n'}}\le\frac{L^2}{\mathfrak M}.
\]
At most $N_{\mathrm G}$ nonconstant cutoff factors occur at a
point. Their first and second derivatives have bounds
$C/s_{n'}$ and $C/s_{n'}^2$. Differentiating
$\theta_\kappa=\phi_\kappa/\Phi$, with $\Phi\ge1$,
therefore gives bounds $C(L)/\mathfrak M$ and
$C(L)/\mathfrak M^2$ after multiplication by $\xi_y$ and
$\xi_y^2$. This proves \eqref{eq:piece-cutoff-derivatives}.

\emph{Separated supports.}
Put $\varrho=(\rho+1)L^2\xi_n\le s_n/16$.
For $y\in S_\kappa\cup S'_\kappa$, the set $\mathsf N(y)$
meets $\kappa^{[2/3]}$. Such a top has level at most $n+1$.
If its level were at least $n+3$, all its nearby cells would
lie in $\Omega_{n+2}$; if it were $n+2$, its distance from
$\kappa^{[1]}$ would exceed $9s_{n+1}$, whereas
$(\rho+1)L\xi_{n+2}<9s_{n+1}$. Both contradict the intersection.
Thus \eqref{eq:sharp-neighborhood-size} places $\mathsf N(y)$
within $\varrho$ of $y$.
For $y'\in S'_\kappa$ some point of this neighborhood lies in
$\kappa^{[1/12]}$. For $y\in S_\kappa$ some point lies in
$\kappa^{[2/3]}\setminus\operatorname{int}\kappa^{[1/3]}$:
otherwise, since the neighborhood has diameter at most
$2\varrho<s_n/3$, the cutoff would be identically $0$ or $1$
there. These observations put both tops in $\kappa^{[5/6]}$.

Let $z$ be the center of $\kappa$ and
$f(x)=\min\{|x-z|_\infty,3s_n/2\}$. The cube
$\kappa^{[1]}$ contains no point of $\Omega_{n+2}$, so every
point there counts with weight at least $\xi_{n+1}^{-1}$.
Apply the interval-image argument of
\cref{subsec:multiscale-sum-proof} to $f$ on a connecting set:
between its endpoint values below $3s_n/2$, its points lie in
that cube. It gives $|f(y)-f(y')|\le\xi_{n+1}d(y,y')$.
For the two tops just considered,
\[
 f(y)-f(y')\ge\tfrac14s_n-2\varrho\ge\tfrac18s_n,
\]
which is \eqref{eq:piece-support-separation}.

Every local top has scale at most $\xi_{n+1}$.
If $\mathsf N^\kappa(y)$ meets $\kappa^{[2/3]}$, its points
are within $2\varrho\le s_n/8$ of that cube, hence in
$\kappa^{[19/24]}$. Its tops and cells are consequently among
the common ones already identified. The inequality $d\le d^\kappa$
gives $\mathsf N^\kappa(y)\subset\mathsf N(y)$.
Indices at these positions have basepoints at most
$\xi_{n+1}/2\le s_n/24$ beyond these cells or tops, so they
lie in $\kappa^{[5/6]}$. The last assertions in (d) follow
from this inclusion and the cutoff supports.
\end{proof}


\subsection{The actual averages and the local inverses}
\label{sec:regular-local-inverses}

The local gauge is used to verify the domains of the averaging maps
and to apply \cref{I-65-40}. The maps themselves remain the actual
nested averages of \(\mathbf U\). Near the physical boundary there
is only one layer; the mass terms then give a separate inverse
without a small gauge potential.

\begin{lemma}\label{lem:regular-local-inverse}
\label{I-65-66}\label{I-65-70}
The constants can be chosen so that all the averages defining the
global and local systems exist. On the rows and tops identified in
\(\kappa^{[5/6]}\), their relative constraint germs, linear averages
and quadratic coefficients agree.
In particular, the entries of the averaging kernels
\(Q,Q'\), their weighted transposes, and the quadratic coefficient
kernels agree at the corresponding indices.
For every piece \(\kappa\) and \(i\in\{1,2,4,5\}\), the local
operator \(\mathbb K_\kappa^{[i]}\), with ambient region
\(\mathsf A=\Omega_0\), is invertible. Its inverse satisfies
\begin{equation}\label{eq:regular-local-inverse}
 |\mathbb E_\kappa^{[i]}\iota_\kappa g|_{\mathbf U,y}
 \le B_{\mathrm{loc}}
       \sup_{z\in\Lambda_*^\kappa}
          e^{-2\delta_{\mathrm{reg}}d^\kappa(y,z)}|g|_z,
 \qquad y\in\Lambda_*^\kappa.
\end{equation}
Here the constants depend only on \(L\), and the empty-region
convention is unchanged.

For use in the cutoff estimates, set
\[
 \mathsf f(b)=\xi^{-2}\max_p|\partial\mathbf U(p)-\mathbf1|,
\]
where \(p\) is an inner plaquette of \(\Omega_0\) with a corner
within sup distance \(\xi\) of \(b_-\). On the global and local
systems,
\begin{equation}\label{eq:regular-local-coefficients}
 |\operatorname{Ad}(\mathbf U(b))^{\pm1}|\le2,
 \qquad \xi_y^2\mathsf f(b)\le L^2\alpha_0\le1,
 \qquad |\mathbf J|_{-3;y}\le\gamma_0,
 \qquad |\mathsf j^\kappa|_{-3;y}\le\gamma'.
\end{equation}
In the curvature inequality \(b\) is an inner bond at position
\(y\); the adjoint bound holds on all inner ambient bonds.
For the global system replace \(\mathsf j^\kappa\) by \(\mathsf j\).
\end{lemma}

\begin{proof}
\emph{A small potential on a regular piece.}
Let \(n=n(\kappa)\). By \cref{lem:piece-geometry},
\(\kappa^{[2]}\) belongs to the prescribed natural cube family,
and the bare \(\xi\)-collar of
\(Y_\kappa=\Omega_0^\kappa\subset\kappa^{[3/2]}\) lies in
\(\kappa^{[2]}\subset\Omega_0\).
The real gauge condition gives a unitary \(u\) and a real field
\(A_0\) on these bonds such that
\[
 U^u=e^{i\xi A_0},\qquad
 \xi_n|A_0|\le\alpha_0,\qquad
 \xi_n^2|\nabla^\xi A_0|\le\alpha_0.
\]
Write \(A''=\operatorname{Ad}(u)A'\). The class bounds and unitary
covariance give
\[
 \mathbf U^u=e^{i\xi A''}e^{i\xi A_0},\qquad
 \xi_n|A''|\le\alpha_1,\qquad
 \xi_n^2|\nabla_{U^u}A''|\le\alpha_1.
\]
Indeed, \(\kappa^{[2]}\) meets only layers \(n,n+1\).
The difference between the covariant and ordinary derivatives is
\(\xi^{-1}(e^{i\xi\operatorname{ad}A_{0,\nu}}-I)A''(\cdot+\xi e_\nu)\).
Thus, for \(\alpha_0\le1/8\),
\[
 \xi_n^2|\nabla^\xi A''|
 \le\alpha_1+4\alpha_0\alpha_1\le2\alpha_1.
\]

Put \(Z=e^{i\xi A''}e^{i\xi A_0}\). If
\(\alpha_0+\alpha_1\le1/8\), then
\(|Z-I|\le e^{\xi(|A''|+|A_0|)}-1\le1/7\).
The series logarithm defines the traceless field
\(\mathbf A_\kappa=(i\xi)^{-1}\log Z\), and
\[
 \xi_n|\mathbf A_\kappa|\le2(\alpha_0+\alpha_1).
\]
The power-series telescoping identity gives
\(|\log Z-\log Z'|\le2|Z-Z'|\) when both matrices lie within
\(1/2\) of \(I\), and
\[
 |e^Xe^Y-e^{X'}e^{Y'}|
 \le e^{1/4}(|X-X'|+|Y-Y'|)
 \quad\text{if } |X|,|Y|,|X'|,|Y'|\le1/8.
\]
Apply these inequalities at two translated bonds. They give
\(\xi_n^2|\nabla^\xi\mathbf A_\kappa|
 \le4(2\alpha_1+\alpha_0)\), and hence
\begin{equation}\label{eq:regular-piece-potential}
 \mathbf U^u=e^{i\xi\mathbf A_\kappa},\qquad
 \xi_{n+1}|\mathbf A_\kappa|,
 \ \xi_{n+1}^2|\nabla^\xi\mathbf A_\kappa|
 \le8L^2(\alpha_0+\alpha_1).
\end{equation}
The difference bound is asserted only when both translated bonds
lie in \(\kappa^{[2]}\). Fix this gauge for the rest of the argument.

\emph{Domains and agreement of the actual maps.}
Every global row midpoint and scalar top belongs to a piece.
By \cref{lem:piece-geometry,lem:sharp-row-neighborhood}, a row
so covered has level \(j\le n+1\), and its bare double block
lies inside \(\kappa^{[1]}\). For a scalar top \(y\in\Lambda_j\),
the auxiliary full-frame bond
\(c_y=\langle y,y+\xi_je_1\rangle\) may be inward and need not
be a row. Its double block lies within \(3\xi_j/2\) of \(y\),
as in the proof of \cref{lem:small-system-averages}; the same
piece margins put it inside \(\kappa^{[1]}\).

For a regular piece choose the radius so that
\begin{equation}\label{eq:regular-averaging-radius}
 \alpha_0+\alpha_1\le1/8,
 \qquad 8L^2(\alpha_0+\alpha_1)<c_{\mathrm{av}}/2.
\end{equation}
Rescale a level-\(j\) double block by \(\xi_j\). Its fine spacing
is \(\eta=L^{-j}\), and its potential is
\(\xi_j\mathbf A_\kappa\), of norm less than
\(c_{\mathrm{av}}/2\). Thus
\cref{prop:flat-complex-averaging} defines the actual nested
averages and an open relative-constraint germ at zero on this block.
The input and output scale factors in
\eqref{eq:small-system-physical-averages} cancel in its first
derivative. The same argument covers every local row, since its
level is at most \(n+1\) and its bare block is inner in
\(Y_\kappa\subset\kappa^{[3/2]}\).
On a boundary piece all rows and tops have level zero, so
\(Q=Q'=I\) and \(C^{(2)}=0\).

Unitary covariance \eqref{I-10-60} and the straightening covariance
\eqref{eq:straightening-covariance} return these definitions to the original gauge.
They preserve the defining domains and the series logarithms,
not only the values. Reversed bond variables and averages are
obtained by inversion. For the relative logarithm,
\(\|a_j\|_\infty\le2c_{\mathrm{av}}<\log2\) from the proof of
\cref{prop:flat-complex-averaging} keeps both exponentials in
the logarithm neighborhood, with logarithms \(ia_j\) and
\(-ia_j\).

By \cref{I-10-59}, the value and definedness of each iterate
depend only on its bare double block. Block nesting gives the same
successive inputs and domain tests for the global and local data.
Thus different covering pieces or gauges evaluate the same
recursion. The equality holds on an open relative-constraint
neighborhood, so the linear and quadratic coefficients agree too.
For \(Q'\), all scalar contours remain in \(B^j(y)\), and the
ordered transport formula
\eqref{eq:flat-complex-scalar-transport} shows independence of the
auxiliary double block. Finally, common indices have the same
levels and pairing weights by \cref{lem:piece-geometry}; their
weighted transposes therefore agree as well.

\emph{Coefficient bounds and the regular inverse.}
For an inner bond at a local top \(y\) of level \(j\), its initial
point belongs to \(\Omega_j^\kappa\subset\Omega_j\); its global
level is consequently at least \(j\). This proves the current
bound in \eqref{eq:regular-local-coefficients}. The truncated row
field has the global values at common rows, with the same weights,
and vanishes elsewhere, proving its bound.
The base point of a plaquette entering \(\mathsf f(b)\) is within
\(2\xi\) of \(b_-\). For \(j\ge1\), separation places it in
\(\Omega_{j-1}\), so the complex curvature condition gives
\[
 |\partial\mathbf U(p)-I|
 \le L^2\alpha_0\xi^2\xi_y^{-2}.
\]
For \(j=0\) use \(\alpha_0\) instead. Also
\(\xi|A'|\le\alpha_1\) on every ambient bond, so
\(|\operatorname{Ad}(\mathbf U)^{\pm1}|\le e^{2\alpha_1}\le2\).
These arguments apply to the global indices as well.

Every local top has \(\xi_y\le\xi_{n+1}\). Thus
\eqref{eq:regular-piece-potential} verifies the literal neighborhood
conditions \eqref{eq:small-system-hypothesis}, including its collar,
with
\[
 \epsilon=\max\{8L^2(\alpha_0+\alpha_1),\gamma_0,\gamma'\}.
\]
Choose \(a_{\mathrm{reg}}\) so that \(\epsilon\le\epsilon_*\).
The local separation parameters are
\(\mathfrak m_\flat=\mathfrak M_\dagger\) and
\(\mathsf g_\flat=L\mathfrak M_\dagger\), and the collar is
contained in \(\Omega_0\). Hence \cref{I-65-40} applies to every
regular piece, with rate \(\delta_{\mathrm{cs}}\) and constant
\(B_*\), uniformly in the number of empty higher layers.
The absolute-kernel estimates proving the first bound of
\eqref{eq:small-system-average-bounds} and
\eqref{eq:small-system-multiplier} are preserved by this unitary
gauge. They therefore hold for the local operators in the original
gauge, with uniform constants. Since those proofs sum absolute
kernel norms, the bounds also allow an inserted scalar factor
bounded in modulus, as will be used in the cutoff calculation.

\emph{The massive boundary inverse.}
Let now \(\kappa\) be a boundary piece and put \(Y=Y_\kappa\).
The tops are the bare sites of \(Y\), the rows are its positive
inner bare bonds, and all pairing weights are \(\xi^3\).
Since there are no positive layers, \eqref{I-9-17} gives
\(d^\kappa(y,z)=\xi^{-1}|y-z|_1\), with periodic distance on
the ambient torus. In particular,
\[
 Q=Q'=I,\qquad \Delta^{(2)}[\mathsf j^\kappa]=0,\qquad
 \Delta'_a=\Delta^\circ+\xi^{-2},\qquad
 \mathsf L_1=\mathsf L+\xi^{-2}.
\]
The multiplier therefore drops out. The natural-class hypotheses
on the whole ambient region imply
\[
 |\partial U(p)-I|\le\alpha_0,\qquad
 \xi|A'(b)|\le\alpha_1,\qquad
 \xi^3|\mathbf J(b)|\le\gamma_0\le1.
\]

We first bound the two massive operators
\[
 T_1=I+\xi^2\Delta_Y(\mathbf U),\qquad
 T_2=I+\xi^2\Delta^\circ_{\mathbf U}.
\]
For real \(U\), the transports in \eqref{I-65-5} are isometries.
Moreover, \(0\le y_0(p)\le1\) when \(\alpha_0\le1/2\),
\(|S_p(A)|^2\le4\sum_i|a_i|^2\), and
\(|F_U(p)|\le\xi^{-2}\alpha_0\). The commutator part obeys
\[
 \sum_{i<j}|\langle[a_i,a_j],F_U(p)\rangle|
 \le3\xi^{-2}\alpha_0\sum_i|a_i|^2.
\]
Every bond lies in at most four plaquettes. Thus, for a real field
\(A\) on the inner bonds of \(Y\), extended by zero,
\begin{equation}\label{eq:boundary-hessian-form}
 -12\alpha_0\|A\|_\xi^2
 \le (A,\xi^2\Delta_Y(U)A)_\xi
 \le(16+12\alpha_0)\|A\|_\xi^2.
\end{equation}
This is the ambient form compressed to \(Y\); plaquettes crossing
the boundary of \(Y\) remain in the sum. The real Hessian is
symmetric, and \eqref{I-63-28} writes
\(\xi^2\Delta_U^\circ=6I-\mathsf h\), where \(\mathsf h\)
is symmetric with at most six isometric entries per row and column.
Consequently the two real massive operators have spectra in
\([1-12\alpha_0,17+12\alpha_0]\) and \([1,13]\), respectively.

For the following norm calculation use the Hermitian pairing
obtained by complexifying the real pairing. Its norm and the
matrix norm on \(\mathfrak g^c\) are equivalent with absolute
constants, as in \eqref{eq:small-system-pairing-bounds}.
Write \(T_k=T_{k,0}+E_k\), with \(T_{k,0}\) evaluated at \(U\).
Moving the unitary transports to the right in a plaquette gives
\[
 |\partial\mathbf U(p)-\partial U(p)|\le e^{4\alpha_1}-1,
 \qquad
 |\operatorname{Ad}(\mathbf U(b))^{\pm1}
      -\operatorname{Ad}(U(b))^{\pm1}|
       \le e^{2\alpha_1}-1.
\]
The entries of \(\xi^2\Delta_Y\) and
\(\xi^2\Delta^\circ\) are fixed polynomials in these transports
and holonomies, by \eqref{I-65-5} and \eqref{I-63-28}.
All arguments stay in a fixed bounded set, so
\(|E_k(i,i')|_2\le C_E\alpha_1\), with an absolute \(C_E\ge1\).
The subscript denotes the induced Hermitian operator norm.

For bonds, let \(\varrho(i,i')\) be graph distance through bare
bonds of the ambient torus, one step joining bonds in a common
plaquette. For sites use the bare nearest-neighbor graph.
The matrices \(T_{k,0}\) and \(E_k\) have range one for these
distances, with at most thirteen nonzero entries in each row and
column. This follows for the Hessian by polarization of
\eqref{I-65-5}, and for the scalar operator from
\eqref{I-63-28}. Also
\(\varrho(i,i')\ge d^\kappa(y(i),y(i'))/2\).

Choose \(\max\{\alpha_0,\alpha_1\}\le1/24\) and
\(C_E\alpha_1\le1/52\).
Both \(T_{k,0}\) then have spectra in \([1/2,18]\).
The Schur bound gives \(\|E_k\|_2\le1/4\); with
\(\tau=4/37\) we obtain
\[
 \|I-\tau T_k\|_2\le\frac{35}{37}+\frac1{37}
       =q_\partial<1,\qquad q_\partial=\frac{36}{37}.
\]
Hence \(T_k^{-1}=\tau\sum_{m\ge0}(I-\tau T_k)^m\).
The \((i,i')\) entry vanishes in every term with
\(m<\varrho(i,i')\), and the remaining terms form a geometric
tail. With
\(\delta_\partial=-\frac14\log q_\partial>0\), this proves
\begin{equation}\label{eq:boundary-massive-decay}
 |T_k^{-1}(i,i')|\le Cq_\partial^{\varrho(i,i')}
 \le C e^{-2\delta_\partial d^\kappa(y(i),y(i'))}.
\end{equation}
At each position there are at most three bond indices and one
site index. Summing one half of the exponential in
\eqref{eq:boundary-massive-decay} over periodic lattice
representatives gives
\[
 |T_k^{-1}f|_y\le C
   \sup_z e^{-\delta_\partial d^\kappa(y,z)}|f|_z;
\]
the constant is bounded by a multiple of
\(\sum_{z\in\mathbb Z^3}e^{-\delta_\partial|z|_1}\), independently
of the torus and of \(Y\).

\emph{Elimination of the boundary slots.}
In the following formulas all operators are those of the local
system, and set
\(\mathsf g(A,v)=\mathsf g_0(A,v)+D^{\mathsf T}\mathsf x(A,v)\).
For \(i=4,5\), slots two through seven of
\(\mathbb K_\kappa\Psi=\iota_\kappa g\) give
\begin{equation}\label{eq:boundary-system-elimination}
 \begin{aligned}
 v&=f_4,&
 w&=f_2+D^{\mathsf T}A-\Delta'_af_4,&
 \sigma&=\Delta'_aw-f_3,\\
 q&=-f_7,&
 p&=f_6+\Delta'_af_7,&
 r&=f_5+D^{\mathsf T}X_5+\mathsf g(A,v)-\Delta'_ap.
 \end{aligned}
\end{equation}
For \(i=5\), the last slot gives \(A=f_8\), and the first gives
\[
 \omega=\mathsf L_1A-D(w+p)-\mathsf m_J(v)-f_1.
\]
For \(i=4\), substitution of \(w\) into the first slot gives
\[
 (\Delta_Y(\mathbf U)+\xi^{-2})A
   =f_1+D(f_2-\Delta'_af_4)+Dp+\mathsf m_J(f_4),
\]
whose inverse is \(\xi^2T_1^{-1}\).
The \(i=2\) system uses the first line of
\eqref{eq:boundary-system-elimination} with \(A=0\).
For \(i=1\), \(w=\xi^2T_2^{-1}f_3\).
These formulas solve all four square systems uniquely.

The remaining factors in these formulas have uniformly bounded
range in units of \(\xi\). Their kernels have exactly the powers
of \(\xi\) prescribed by the input and output weights in
\eqref{I-64-3}--\eqref{I-64-4}: a difference contributes
\(\xi^{-1}\), a Laplacian \(\xi^{-2}\), and the current
\(\xi^{-3}\). The coefficient bound \(\xi^3|\mathbf J|\le1\)
therefore suffices for \(\mathsf m_J,\mathsf g_0,\mathsf x\).
Each component, including the differences in the solution norm,
is a finite sum of such products with at most one factor
\(\xi^2T_k^{-1}\). The weighted-supremum composition estimate of
\cref{sec:trivial-system-calculus} and
\eqref{eq:boundary-massive-decay} give the local inverse bound
with rate \(\delta_\partial\) and a constant \(B_\partial\).

Choose
\[
 2\delta_{\mathrm{reg}}
   =\min\{\delta_{\mathrm{cs}},\delta_\partial,1\},
 \qquad B_{\mathrm{loc}}=\max\{B_*,B_\partial\}.
\]
The regular and boundary estimates now imply
\eqref{eq:regular-local-inverse}. All radius restrictions above
are finitely many inequalities on \(a_{\mathrm{reg}}\) depending
only on \(L\).
\end{proof}

\begin{lemma}\label{lem:piece-operator-agreement}
Let the global and local systems be those constructed in
\cref{lem:regular-local-inverse}, with the common ambient region
$\Omega_0$ and the multiplier \eqref{eq:piece-row-multiplier}.
A field supported at indices with basepoints in $\kappa^{[2/3]}$
is identified uniquely with a field of the other system by zero
extension. Under this identification,
\begin{equation}\label{eq:piece-operator-agreement}
 \mathbb K\Psi=\mathbb K_\kappa\Psi,
 \qquad
 \operatorname{supp}(\mathbb K\Psi)
       \subset\{\text{indices with basepoint in }\kappa^{[5/6]}\}.
\end{equation}
The two solution seminorms of $\Psi$ agree at common tops;
at every other top the respective seminorm is zero. At an output
index with basepoint $z\in\kappa^{[5/6]}$, the system reads
$\Psi$ only at basepoints within sup distance $s_{n(\kappa)}/6$
of $z$.
\end{lemma}

\begin{proof}
The domain assertion of \cref{lem:regular-local-inverse} identifies
the actual recursive averages and their relative constraint germs,
not only their estimates. Thus $Q,Q'$ and $C^{(2)}$ have identical
entries at common rows and tops; the weighted transposes agree
when the weights do. We now check the remaining support and
boundary issues.

The differences, their transposes, the Dirichlet Laplacian and the
current terms in \eqref{I-63-63} use only endpoints and one bare
step. Polarizing the ambient plaquette formula shows that the
Hessian couples only bonds of a common plaquette, using the
configuration on that plaquette. Both compressed Hessians are
entries of this one ambient Hessian, including plaquettes crossing
the local boundary. The averages connect a row to its double
block $\mathsf Y_c$, or a top to its bare block. The quadratic
multiplier connects two bonds in one $\mathsf Y_c$. All their
basepoints lie in the row box or cell of
\cref{lem:sharp-row-neighborhood}.

Inspection of \eqref{I-63-64} and \eqref{I-63-63} therefore shows
that an entry of the system is a sum of kernel products along
chains passing through at most one such box or cell and at most
two additional bare steps. The longest chains occur in
$D^{\mathsf T}\mathsf x$ and in
$\Delta^{(2)}[\mathsf j](A-Dv)$; taking $Dv$ in the latter
stays inside $\mathsf Y_c$. If either end of a chain lies in
$\kappa^{[5/6]}$, its box meets $\kappa^{[1]}$, hence has level
at most $n+1$ by the geometric proof. Its sup diameter is at
most $2\xi_{n+1}$. All the objects along the chain are thus
within $2\xi_{n+1}+2\xi\le4\xi_{n+1}<s_n/6$ of either
endpoint. This proves the stated input range and shows that a
field supported in $\kappa^{[2/3]}$ has output supported in
$\kappa^{[5/6]}$.

For such an output, add the possible $\xi/2$ from a bond midpoint
to its endpoint and $\xi_{n+1}/2$ to a position. The total is
at most $5\xi_{n+1}<s_n/6$, so all these sites and positions
lie inside $\kappa^{[1]}$. Their region memberships, indices,
weights and averaging blocks agree. The local current and
configuration are restrictions of the same ambient data, and the
Hessian entries agree by the common ambient construction.

It remains to check the discarded rows of the multiplier.
If a global row midpoint lies outside $\kappa^{[5/6]}$ and its
box meets $\kappa^{[1]}$, the row has level at most $n+1$ and
the box is within $\xi_{n+1}$ of its midpoint. Its distance
from $\kappa^{[2/3]}$ is therefore greater than
$s_n/6-\xi_{n+1}>\xi$. If the box does not meet
$\kappa^{[1]}$, it does not meet the support either.
In both cases $A$ and $v$, and hence $A-Dv$, vanish on the
bare double block of the row. Its contribution to
$\Delta^{(2)}[\mathsf j](A-Dv)$ is zero. All remaining rows
are common rows with the same multiplier value and quadratic
coefficient. This proves \eqref{eq:piece-operator-agreement}.

For the seminorm assertion, a nonzero value of a component or
of one of $\nabla A,Dv,Dw,Dq$ has basepoint within $\xi$ of
$\kappa^{[2/3]}$. Every index used by that difference is within
$2\xi$ of this cube and lies in $\kappa^{[5/6]}$.
The position of such an index is within
$3\xi/2+\xi_{n+1}/2\le2\xi_{n+1}<s_n/6$ of the support.
It is therefore a common top with the same weight. Shifted bonds
are indices in both sequences or in neither, and the transports
are those of the same configuration. Each entry of
\eqref{I-64-3} agrees, while entries outside these positions
vanish. Finally, the common-index identification in
\cref{lem:piece-geometry} gives existence and uniqueness of the
zero extensions used throughout.
\end{proof}

\subsection{The cutoff error}
\label{subsec:regular-cutoff}

Multiplication by a real function $h$ is always evaluated at
index base points, in particular at bond and row midpoints.
We write $[T,h]=Th-hT$. The scalar variable $p$ has no derivative
in \eqref{I-64-3}; accordingly its Laplacian commutator must be
kept in divergence form.

For a scalar field $f(x)$, or a bond component indexed by its
initial point $x$, write $\hat x$ for its base point. Put
\[
 k_\nu(x)=\xi^{-1}\bigl(h(\hat x+\xi e_\nu)-h(\hat x)\bigr),
 \qquad
 X^{h[f]}_\nu(x)=k_\nu(x)
             \bigl(\mathsf U_\nu(x)f(x+\xi e_\nu)+f(x)\bigr)
\]
when the shifted index exists, and put both expressions equal
to zero otherwise. With the same convention for a missing
previous index, define
\[
 (\mathsf c^hf)(x)=\xi^{-1}\sum_\nu
                     \bigl(k_\nu(x)-k_\nu(x-\xi e_\nu)\bigr)f(x).
\]
Direct expansion of \eqref{I-63-28} gives, including at the boundary,
\begin{equation}\label{eq:regular-laplacian-cutoff}
 [\Delta^\circ,h]f
       =\sum_\nu\nabla_\nu^{\mathsf T}X^{h[f]}_\nu+\mathsf c^hf.
\end{equation}
For $f=p$, the first term is $D^{\mathsf T}X^{h[p]}$.

For any of the linear maps of $(A,v)$ in \eqref{I-63-63}, use
$\mathsf t^h(A,v)=\mathsf t(hA,hv)-h\mathsf t(A,v)$.
Define $\mathbb C_h\Psi\in\mathcal F$ by taking the corresponding
components of $[\mathbb K,h]\Psi$ in all $f_j$ slots except
$f_5$, and replacing that slot by
\begin{equation}\label{eq:regular-cutoff-fifth-slot}
 \begin{split}
 f_5={}&\mathsf c^hp+[Q'^{\mathsf T}\mathsf a'Q',h]p
        +[Q'^{\mathsf T},h]r-\mathsf g_0^h(A,v)
                         -[D^{\mathsf T},h]\mathsf x(A,v),\\
 X_5={}&X^{h[p]}-\mathsf x^h(A,v).
 \end{split}
\end{equation}
The smaller systems retain just their existing slots.

\begin{lemma}\label{lem:regular-cutoff}
On a local system from \cref{lem:regular-local-inverse}, let
$\vartheta\ge1$ and let $h$ be real and $C^2$ on an open
neighborhood of its region, with
\[
 |h|\le1,\qquad \xi_y|\partial h|\le\vartheta^{-1},\qquad
 \xi_y^2|\partial^2h|\le\vartheta^{-1}
                         \quad\hbox{on }\mathsf N^\kappa(y).
\]
There is a constant $C_{\mathbb C}$ depending only on $L$ such that
\begin{equation}\label{eq:regular-cutoff-bound}
 [\mathbb K_\kappa,h]=\iota_\kappa\mathbb C_h,
 \qquad
 |\mathbb C_h\Psi|_y\le C_{\mathbb C}\vartheta^{-1}
              \max_{d^\kappa(y,z)\le\rho}|\Psi|_z.
\end{equation}
The left seminorm is zero if $h$ is constant on
$\mathsf N^\kappa(y)$. Also $|h\Psi|_y\le2|\Psi|_y$ and
$|hg|_y\le|g|_y$.
\end{lemma}

\begin{proof}
We omit the local superscript. Neighboring cells have scale
ratios between $L^{-1}$ and $L$. The line segments joining a
site to an incident bond midpoint, or joining neighboring
indices of one component, lie in $\mathsf N(y)$ at the output
position. The derivative assumptions therefore give
\[
 |k_\nu(x)|\le\vartheta^{-1}\xi_y^{-1},\qquad
 \xi^{-1}|k_\nu(x)-k_\nu(x-\xi e_\nu)|
                         \le\vartheta^{-1}\xi_y^{-2}
\]
when the two shifted indices exist. If one is absent, its
position is at level zero, so $\xi=\xi_y$ gives the same bound
with an absolute constant. In the interior, expanding the
Laplacian commutator in first differences instead gives
\[
 [\Delta^\circ,h]f(x)
 =-\sum_\nu\bigl(k_\nu(x)\nabla_\nu f(x)
      -k_\nu(x-\xi e_\nu)\nabla_\nu^{\mathsf T}f(x)\bigr)
      -\mathsf c^hf(x).
\]
Consequently, for the weights used in the system,
\begin{equation}\label{eq:regular-cutoff-laplacian-bounds}
 \begin{split}
 |[\Delta^\circ,h]f|_{s-2;y}
   &\le C\vartheta^{-1}\max_{d(y,z)\le\rho}
              (|f|_{s;z}+|\nabla f|_{s-1;z}),\\
 |X^{h[f]}|_{s-1;y}+|\mathsf c^hf|_{s-2;y}
   &\le C\vartheta^{-1}\max_{d(y,z)\le\rho}|f|_{s;z}.
 \end{split}
\end{equation}
At missing neighbors the first bound follows directly from
$\xi=\xi_y$.

The first-order calculation is equally local. For example,
\begin{equation}\label{eq:regular-divergence-cutoff}
 \begin{split}
 [D^{\mathsf T},h]X(x)
 ={}&\xi^{-1}\sum_{b_+=x}(h(\hat b)-h(x))
                         \operatorname{Ad}(\mathbf U(b))^{-1}X(b)\\
   &-\xi^{-1}\sum_{b_-=x}(h(\hat b)-h(x))X(b).
 \end{split}
\end{equation}
The analogous formula for $[D,h]$ has the two endpoint values
and the midpoint value of $h$. Since the adjoint transports
have norm at most two, these commutators lower the weight by
one with a factor $C\vartheta^{-1}$.

All averaging terms can be treated by the same inserted
difference. If $T(i,i')$ is one of their kernels, then
\[
 ([T,h]f)(i)=\sum_{i'}T(i,i')
                         (h(\hat i')-h(\hat i))f(i').
\]
The row terms have both indices in a single row box; scalar
terms have them in a single cell. By
\cref{lem:sharp-row-neighborhood}, these sets lie in the output
neighborhood and the difference of $h$ is at most
$\rho\vartheta^{-1}$. The kernel bounds in
\eqref{eq:small-system-average-bounds}, now available by
\cref{lem:regular-local-inverse}, show that the commutators with
$Q,Q^{\mathsf T},Q',Q'^{\mathsf T}$ preserve weights, while
those with $Q^{\mathsf T}\mathsf m^{-2}Q$ and
$Q'^{\mathsf T}\mathsf a'Q'$ lower them by two. Expanding each
potential within its common row or cell keeps the same support
neighborhood; one does not compose two enlarged neighborhoods.

The formula \eqref{eq:small-system-multiplier} treats
$\Delta^{(2)}$ in the same way. Together with \eqref{I-63-63}
and the current bounds, it gives weights $-3,-4,-3$ for
$\mathsf m_J^h(v),\mathsf g_0^h(A,v),\mathsf x^h(A,v)$,
respectively, with bounds $C\vartheta^{-1}$ times the neighboring
solution seminorms. In the first slot the additional terms are
\[
 -[\Delta^{(2)},h](A-Dv)+\Delta^{(2)}[D,h]v.
\]
They have weight $-3$ with the same bound. In the composed terms
here and in $[D^{\mathsf T},h]\mathsf x$, the common row box
contains the output bond or its endpoint, and contains all
input indices. Thus the sharp row neighborhood still applies.

For a regular piece its bare collar lies in the ambient region,
so \eqref{I-65-6} gives
$[\mathsf L,h]=[\Delta^\circ,h]+[\mathsf Z_{\mathrm{curv}},h]$.
The curvature kernel joins bonds of a common plaquette and has
size at most $C\mathsf f(b)$. Its commutator therefore lowers
weight $-1$ to $-3$ with the required factor, since
$\xi_y^2\mathsf f(b)\le1$. For a boundary piece every scale is
$\xi$. The Hessian formula \eqref{I-65-5} bounds its kernel by
$C\xi^{-2}$ on pairs of bonds of a common plaquette; the same
inserted difference proves the assertion. The kernel of
$DD^{\mathsf T}$ is treated by the two first-order formulas.

Finally, \eqref{eq:regular-laplacian-cutoff} and
\[
 (D^{\mathsf T}\mathsf x)^h
   =[D^{\mathsf T},h]\mathsf x+D^{\mathsf T}\mathsf x^h
\]
give the fifth component of
$[\mathbb K,h]\Psi=\iota\mathbb C_h\Psi$.
The other components are their defining commutators. Apply the
first bound of \eqref{eq:regular-cutoff-laplacian-bounds} to
$A,v,w,q$ and the second to $p$; the bounds just proved give
each forcing weight in \eqref{I-64-4}. Every term contains a
difference of values of $h$ within $\mathsf N(y)$, proving the
vanishing assertion. The product rule for a first difference,
and $\vartheta^{-1}\le1$, give $|h\Psi|_y\le2|\Psi|_y$;
the forcing norm contains no derivatives, so $|hg|_y\le|g|_y$.
\end{proof}

The input cutoff also requires a divergence correction. For
$i=4,5$, let $h\star g$ agree with $hg$ except in its fifth
scalar component, where
\begin{equation}\label{eq:regular-forcing-cutoff}
 (h\star g)_5=hf_5-[D^{\mathsf T},h]X_5.
 \qquad \iota(h\star g)=h\iota g.
\end{equation}
For $i=1,2$, set $h\star g=hg$. The identity follows by adding
$D^{\mathsf T}(hX_5)$ to the displayed component.
Formula \eqref{eq:regular-divergence-cutoff} also proves
\begin{equation}\label{eq:regular-forcing-cutoff-bound}
 |\theta_\kappa\star g|_y
 \le(1+C\vartheta^{-1})
   \max_{\substack{z\in S'_\kappa\,,\ d(y,z)\le3}}|g|_z,
 \qquad \vartheta=\mathfrak M/C_{\mathrm G}.
\end{equation}
Indeed, the midpoint-to-endpoint difference contributes
$\xi\vartheta^{-1}\xi_y^{-1}/2$, the factor $\xi^{-1}$
cancels it, and the weight of $X_5$ is $-3$. Adjacent positions
have scale ratio at most $L$ and distance at most three.
A nonzero cutoff value or difference has a site or midpoint
in $\kappa^{[1/12]}$; both relevant neighborhoods meet this
set. Thus both the output position and the input positions
belong to $S'_\kappa$. The same statements hold in the local
sequence, using \eqref{eq:piece-neighborhood-comparison}.
Taking $\mathfrak M_*$ larger makes the factor in
\eqref{eq:regular-forcing-cutoff-bound} at most two.

\subsection{The parametrix and convergence}
\label{subsec:regular-parametrix}

Fix $a_{\mathrm{reg}}$ and $\delta=\delta_{\mathrm{reg}}$ as in
\cref{lem:regular-local-inverse}. Enlarge $\mathfrak M_*$ to
satisfy the geometric margins, the cutoff requirements, and
\begin{equation}\label{eq:regular-remainder-constants}
 \begin{split}
 C_q&=2C_{\mathbb C}C_{\mathrm G}B_{\mathrm{loc}}
                                      e^{2\rho+8},\qquad
 B_T=4B_{\mathrm{loc}}e^{3\delta},\\
 q_{\mathfrak M}&=C_q\mathfrak M^{-1}e^{-\delta D_0},
 \qquad N_{\mathrm G}q_{\mathfrak M}\le\tfrac12,\qquad
 B_{\mathrm{reg}}=2N_{\mathrm G}B_T.
 \end{split}
\end{equation}
The contraction condition follows, for example, from
$\mathfrak M\ge2N_{\mathrm G}C_q$. All these choices depend
only on $L$.

Define the local remainder on the original system space by
\[
 \mathsf R_\kappa
   =-[\mathbb K_\kappa,\widetilde h_\kappa]
                      \mathbb E_\kappa\theta_\kappa.
\]
Restrictions and zero extensions use the common indices in the
enlarged pieces.

For $g\in\mathcal F$, restrict its entries to the common
indices in $\kappa^{[5/6]}$ and extend by zero to the other
local indices; denote this restriction by $g|_\kappa$.
Define, with all operators on the right taken in the local system,
\begin{equation}\label{eq:regular-lifted-remainder}
 \begin{split}
 \mathsf T_\kappa g
   &=\widetilde h_\kappa\mathbb E_\kappa\iota_\kappa
                                  (\theta_\kappa\star(g|_\kappa)),\\
 \mathsf R^\sharp_\kappa g
   &=-\mathbb C_{\widetilde h_\kappa}\mathbb E_\kappa\iota_\kappa
                                  (\theta_\kappa\star(g|_\kappa)).
 \end{split}
\end{equation}
Their supported outputs are identified with global fields and
extended by zero. The input correction uses only indices in
$\kappa^{[5/6]}$, by \eqref{eq:regular-forcing-cutoff-bound}.
For the output correction, the local neighborhoods meeting the
support of $\widetilde h_\kappa$ lie in $\kappa^{[19/24]}$.
Moving from a site to an incident bond changes the base point
by at most $\xi$, and $\xi\le s_{n(\kappa)}/24$.
Thus the divergence injections and their cutoffs still use
common indices in $\kappa^{[5/6]}$. In particular,
\begin{equation}\label{eq:regular-lift-identities}
 \mathsf R_\kappa\iota=\iota\mathsf R^\sharp_\kappa,
 \qquad
 \widetilde h_\kappa\mathbb E_\kappa\theta_\kappa\iota
                                         =\mathsf T_\kappa.
\end{equation}
These are identities of the actual forcing injections, with
the midpoint convention of \eqref{eq:regular-divergence-cutoff}.

\begin{lemma}\label{lem:regular-remainder}
For every $g\in\mathcal F$,
\begin{equation}\label{eq:regular-remainder-bound}
 \begin{split}
 |\mathsf R^\sharp_\kappa g|_y
    &\le q_{\mathfrak M}
       \sup_{z\in S'_\kappa}e^{-\delta d(y,z)}|g|_z,\\
 |\mathsf T_\kappa g|_y
    &\le B_T\sup_{z\in S'_\kappa}e^{-\delta d(y,z)}|g|_z.
 \end{split}
\end{equation}
The first seminorm vanishes outside $S_\kappa$; the second
vanishes at tops outside $\kappa^{[5/6]}$.
\end{lemma}

\begin{proof}
Let $u=\mathbb E_\kappa\iota_\kappa
(\theta_\kappa\star(g|_\kappa))$.
The local inverse bound, the cutoff bound and
$d^\kappa\ge d$ give, at every common top $v$,
\[
 |u|_v\le2B_{\mathrm{loc}}e^{6\delta}
       \sup_{z\in S'_\kappa}e^{-2\delta d(v,z)}|g|_z.
\]
Indeed the input cutoff moves a position by distance at most
three. Applying \eqref{eq:regular-cutoff-bound}, and then
$d(y,v)\le d^\kappa(y,v)\le\rho$ for its output neighborhood,
yields
\[
 |\mathsf R^\sharp_\kappa g|_y
 \le\frac{2C_{\mathbb C}C_{\mathrm G}B_{\mathrm{loc}}}
               {\mathfrak M}
        e^{2\delta\rho+6\delta}
        \sup_{z\in S'_\kappa}e^{-2\delta d(y,z)}|g|_z.
\]
This seminorm vanishes outside $S_\kappa$. On that set the
support separation gives $d(y,z)\ge D_0$ for every
$z\in S'_\kappa$. Hence
\[
 e^{-2\delta d(y,z)}
       \le e^{-\delta D_0}e^{-\delta d(y,z)}.
\]
Since $\delta\le1/2$, the first estimate follows from
\eqref{eq:regular-remainder-constants}.

For $\mathsf T_\kappa$, use the local inverse at rate $\delta$
and $|\widetilde h_\kappa u|_y\le2|u|_y$.
The input cutoff contributes at most $2e^{3\delta}$, giving
$B_T$. Its values are supported in $\kappa^{[2/3]}$.
The derivatives in its solution seminorm can reach adjacent
indices, whose positions remain in $\kappa^{[5/6]}$ by
\eqref{eq:piece-neighborhood-comparison}. This proves the
second support assertion and also justifies the comparison
of local and global solution seminorms used above.
\end{proof}

\begin{proof}[Proof of \cref{I-65-83}]
Set
\[
 \mathsf{Par}=\sum_\kappa
              \widetilde h_\kappa\mathbb E_\kappa\theta_\kappa,
 \qquad
 \mathsf R=\sum_\kappa\mathsf R_\kappa,
 \qquad
 \mathsf R^\sharp=\sum_\kappa\mathsf R^\sharp_\kappa.
\]
At a given top at most $N_{\mathrm G}$ remainder terms have
nonzero seminorm. Thus
$\mathsf R^\sharp\in\mathfrak O(\delta,N_{\mathrm G}q_{\mathfrak M})$,
and \eqref{eq:regular-remainder-constants} makes it a contraction.
The local/global operator agreement and the partition of unity
give
\begin{equation}\label{eq:regular-parametrix-identity}
 \begin{split}
 \mathbb K(\widetilde h_\kappa\mathbb E_\kappa\theta_\kappa F)
    &=\theta_\kappa F-\mathsf R_\kappa F,\\
 \mathbb K\mathsf{Par}&=I-\mathsf R,
 \qquad\mathsf R\iota=\iota\mathsf R^\sharp.
 \end{split}
\end{equation}
Here $\widetilde h_\kappa=1$ on the support of
$\theta_\kappa$. Given $F\in\mathcal X$, take $g=(F;0)$ so
that $\iota g=F$; for $i=1,2$ take $g=F$. The convergent series
$g_\infty=\sum_{m\ge0}(\mathsf R^\sharp)^mg$ satisfies
$(I-\mathsf R^\sharp)g_\infty=g$. Therefore
\[
 \mathbb K\mathsf{Par}\iota g_\infty
     =\iota(I-\mathsf R^\sharp)g_\infty=F.
\]
Surjectivity on the finite-dimensional space $\mathcal X$
proves invertibility. By \eqref{eq:regular-lift-identities},
\[
 \mathbb E\iota
   =\left(\sum_\kappa\mathsf T_\kappa\right)
           \sum_{m\ge0}(\mathsf R^\sharp)^m.
\]
At each top, at most $N_{\mathrm G}$ of the solution seminorms
$|\mathsf T_\kappa g|_y$ are nonzero, so
$\sum_\kappa\mathsf T_\kappa\in
\mathfrak O(\delta,N_{\mathrm G}B_T)$.
The composition rule following \eqref{eq:trivial-system-sup-decay}
therefore bounds the $m$th summand by
\[
 N_{\mathrm G}B_T(N_{\mathrm G}q_{\mathfrak M})^m
       \sup_z e^{-\delta d(y,z)}|g|_z.
\]
Since there are finitely many tops and $\max_y|\cdot|_y$ is a norm,
this bound also proves absolute operator convergence.
Summing the geometric series gives
$2N_{\mathrm G}B_T=B_{\mathrm{reg}}$ and proves
\eqref{eq:regular-system-decay}.
\end{proof}

\section{The inverse with corner averages}
\label{sec:corner-inverse}

The real background construction uses the corner averages of
\cref{subsec:corner-average-defect}. At the identity, the resulting
linear system is a restriction of the flat system already inverted
in \cref{I-64-24}. Local gauges and a partition of unity give an
inverse at a regular configuration. Positivity requires a further
comparison with the positive local inverses.

Work in level-$n$ units, with $1\le n\le K$,
$\eta=L^{-n}$ and $\xi_j=L^j\eta$. The torus has side
$2L^{m+K-n}$ in these units. Let
$\Omega_0\supseteq\cdots\supseteq\Omega_n$ be a sequence
of regions, with $\Omega_{n+1}=\Omega_{n+2}=\emptyset$.
Use its layers, tops, cells, local block side and weighted distance
from \cref{def:local-regularity,sec:multiscale-sums}.
Put $d=d_{\mathbf B}$. Write $E_0$ for its bare sites, $S_\alpha$ for its positive inner
bare bonds, and $S_\beta$ for the positive bare bonds with at
least one endpoint in $E_0$.

The base site of a bond is its initial endpoint when that endpoint
belongs to $E_0$, and its final endpoint otherwise. A bond
difference has the base site of its first bond. A fine index is
assigned to the top whose cell contains its base site; a top
index is assigned to itself. Thus $|f|_y$ is the supremum of the
values assigned to $y$, with value zero on an empty set. This
convention allows six bonds at one base site near the outer
boundary. All fine fields are extended by zero.

For a real configuration $U$ on all bare bonds, let
$D_S$ be its covariant gradient into either $S=S_\alpha$ or
$S=S_\beta$, and let $D_S^*$ be its adjoint for the fine
pairing $\eta^3\sum\langle\cdot,\cdot\rangle$. The scalar
Dirichlet Laplacian is $\Delta_U=D_{S_\beta}^*D_{S_\beta}$;
it retains all six diagonal terms, including those for exterior
neighbors. Covariant differences of a bond field are taken when
both translated bonds belong to $S$.

The scalar and reduced bond averages are
\[
 \begin{split}
 (\mathsf Q'_{\mathsf c}\lambda)(y)
   &=L^{-3j}\sum_{x\in B^j(y)}
        \operatorname{Ad}(T^{\mathsf c}_U(y,x))\lambda(x),
          \qquad y\in\Lambda_j,\\
 (\mathsf Q_{\mathsf c}A)(c)
   &=\sum_{b\in S_\alpha}Q_j^{\mathsf c}[U](c,b)A(b),
          \qquad c\in\mathbf B^\flat\cap\Gamma_j^+.
 \end{split}
\]
Here $\mathbf B^\flat$ consists of the positive interior and
outward determining bonds, and the transports and kernels are
those of \cref{subsec:corner-average-defect}. The scalar top
pairing has weight $\xi_y^3$ and the bond-row pairing has weight
$\xi_c^{-1}$. Every star below uses these pairings. In particular,
$\mathsf Q'_{\mathsf c}(\mathsf Q'_{\mathsf c})^*=I$:
on each cell the transport is orthogonal and there are $L^{3j}$
sites. At level zero the two averages are $I$ and $\eta I$.

Let $\varkappa_0=\varkappa_1=1$ and
$\varkappa_{j+1}=\varkappa_j/(L^{-2}+\varkappa_j)$ for $j\ge1$,
and let $\mathsf a$ multiply a level-$j$ scalar top by
$\varkappa_j\xi_j^{-2}$. Set
\[
 \mathsf A'_{\mathsf c}
   =\Delta_U+(\mathsf Q'_{\mathsf c})^*\mathsf a\mathsf Q'_{\mathsf c},
 \qquad
 \mathsf R_{\mathsf c}
   =\text{orthogonal projection onto }
                         \Delta_U\ker\mathsf Q'_{\mathsf c}.
\]
Let $\mathsf H_S$ be the Hessian of the ambient Wilson action
$\eta^{-1}\sum_p a_p(e^{i\eta A}U)$ at $A=0$, compressed to
fields supported in $S$. Plaquettes crossing the boundary remain
in this definition. Define
\[
 \mathsf M_S=(\mathsf Q_{\mathsf c})^*\mathsf Q_{\mathsf c}
                  +\eta^{-2}\chi_{S\setminus S_\alpha},
 \qquad
 \mathsf A_{\mathsf c}[S]
     =\mathsf H_S+D_S\mathsf R_{\mathsf c}D_S^*+\mathsf M_S.
\]
The last mass term acts only on the touching boundary bonds.

For a power $M$ of $L$, use the pieces of
\cref{def:regular-pieces} with $s_j=M\xi_j$. This definition
uses only the regions and cubes. A size-zero piece whose double
enlargement misses $\Omega_1$ is a boundary piece; all other
pieces are regular. The needed cover and cutoff properties are
proved below under the hypotheses of the theorem.


\begin{theorem}\label{I-32-62}
Let $L\ge13$ be odd. There are a power $M_\diamond$ of $L$
and constants $\varepsilon_\diamond,B_\diamond,\delta_\diamond>0$,
depending only on $L$, with the following property. With the objects just defined, suppose that
\begin{enumerate}[label=(\roman*)]
\item $M\ge M_\diamond$ is a power of $L$, $\rho\ge LM$,
each $\Omega_j$ is a union of $M$-cubes at level $j$, and
\[
 \operatorname{dist}_\infty(\Omega_j^c,\Omega_{j+1})
       >\rho\xi_j\quad(0\le j<n).
\]
Moreover $\Omega_1\ne\emptyset$, and the level-$n$ torus has
at least $6\rho$ sites in each coordinate direction.
\item $U$ is $G$-valued on all bare bonds, and, for some
$0\le\alpha\le L^{-2}\min\{1,c_{\mathsf c}\}$,
\[
 |\partial U(p)-\mathbf1|
       \le\alpha\eta^2\mathsf L_{\mathbf B}(p)^{-2}
       \quad\text{on every inner plaquette of }\Omega_0.
\]
\item For every regular piece $\kappa$ of size $j$, there are
a real gauge $v_\kappa$ on $\kappa^{[1]}$ and a real field
$A_\kappa$ on its inner bonds such that
\[
 U^{v_\kappa}=e^{i\eta A_\kappa},\qquad
 \xi_j|A_\kappa|\le\varepsilon,\qquad
 \xi_j^2|\nabla^\eta A_\kappa|\le\varepsilon,
 \qquad 0<\varepsilon\le\varepsilon_\diamond.
\]
The difference bound is required when both translated bonds
belong to $\kappa^{[1]}$.
\end{enumerate}
Then $\mathsf A'_{\mathsf c}$ is positive definite. Its inverse
$\mathsf G'_{\mathsf c}$ and its covariant gradient satisfy
\begin{equation}\label{eq:corner-scalar-total}
 \begin{split}
 |\mathsf G'_{\mathsf c} f|_y
     &\le B_\diamond\xi_y^2
                  \sup_z e^{-\delta_\diamond d(y,z)}|f|_z,\\
 |D\mathsf G'_{\mathsf c} f|_y
     &\le B_\diamond\xi_y
                  \sup_z e^{-\delta_\diamond d(y,z)}|f|_z.
 \end{split}
\end{equation}
Here $D$ takes values on $S_\beta$.

If the plaquette bound in \textup{(ii)} holds at every ambient
bare plaquette having a corner in $\Omega_0$, then, for
$S=S_\alpha$ and $S=S_\beta$, the operator
\[
 \mathsf A_{\mathsf c}[S]=\mathsf H_S+D\mathsf R_{\mathsf c}D^*+\mathsf M_S
\]
is positive definite. Writing $\mathsf G_{\mathsf c,S}=\mathsf A_{\mathsf c}[S]^{-1}$,
we have
\begin{equation}\label{eq:corner-vector-total}
 \begin{split}
 |\mathsf G_{\mathsf c,S}J|_y
     &\le B_\diamond\xi_y^2
                  \sup_z e^{-\delta_\diamond d(y,z)}|J|_z,\\
 |\nabla^U\mathsf G_{\mathsf c,S}J|_y
     &\le B_\diamond\xi_y
                  \sup_z e^{-\delta_\diamond d(y,z)}|J|_z.
 \end{split}
\end{equation}
The second line uses only pairs $(b,\mu)$ for which both
$b$ and $b+\eta e_\mu$ belong to $S$. All suprema range over
$\Lambda_*$. The constants are independent of the volume,
the scale, the regions and the number of layers.
\end{theorem}

The proof occupies the rest of this section. All constants denoted by $C$
depend only on $L$ and may increase between estimates.

\subsection{The local system and its flat inverse}

The averages used in the theorem are defined under its curvature
hypothesis. A level-$q$ double block lies in $\Omega_0$ and
meets no layer below $q-1$. Its inner curvature is therefore at
most $L^2\alpha\eta^2\xi_q^{-2}$. Since
$L^2\alpha\le c_{\mathsf c}$, the folding and double-block
construction in the proof of \cref{I-32-37} applies. Construct
the rows in its axial gauge and use covariance to return to $U$.
At level zero they are already the stated identity maps.

Suppress the subscript $S$ on $D_S$. For
$\Psi=(A;v,w,\sigma)$ define
\[
 \mathbb K_{\mathsf c}^{[3]}\Psi
 =\bigl((\mathsf H_S+DD^*+\mathsf M_S)A-Dw,
       \mathsf A'_{\mathsf c}v+w-D^*A,
       \mathsf A'_{\mathsf c}w-(\mathsf Q'_{\mathsf c})^*\sigma,
       \mathsf Q'_{\mathsf c}v\bigr).
\]
The scalar system $\mathbb K_{\mathsf c}^{[2]}$ consists of the
last three equations with $A=0$, and
$\mathbb K_{\mathsf c}^{[1]}w=\mathsf A'_{\mathsf c}w$.
Use the solution and forcing seminorms
\[
 \begin{split}
 |\Psi|_y=\max\{&\xi_y|A|_y,\xi_y^2|\nabla^UA|_y,
             |v|_y,\xi_y|Dv|_y,\\
             &\xi_y^2|w|_y,\xi_y^3|Dw|_y,
                                      \xi_y^4|\sigma(y)|\},\\
 |F|'_y=\max\{&\xi_y^3|f_1|_y,\xi_y^2|f_2|_y,
                             \xi_y^4|f_3|_y,|f_4(y)|\}.
 \end{split}
\]
Delete the absent entries for the scalar systems, whose gradient
norms always use $S_\beta$. For the four-component system they
use the chosen frame $S$. Near a regular piece both frames
coincide.

We record the elimination needed to recover the vector inverse.
If $\mathsf G'_{\mathsf c}=(\mathsf A'_{\mathsf c})^{-1}$ exists,
then
\[
 \mathsf N_{\mathsf c}
  =\mathsf Q'_{\mathsf c}(\mathsf G'_{\mathsf c})^2
                                     (\mathsf Q'_{\mathsf c})^*
\]
is positive definite, because $(\mathsf Q'_{\mathsf c})^*$ is
injective. The orthogonal complement of
$\Delta_U\ker\mathsf Q'_{\mathsf c}
 =\mathsf A'_{\mathsf c}\ker\mathsf Q'_{\mathsf c}$ is
$\operatorname{ran}(\mathsf G'_{\mathsf c}
(\mathsf Q'_{\mathsf c})^*)$. Hence
\[
 I-\mathsf R_{\mathsf c}
   =\mathsf G'_{\mathsf c}(\mathsf Q'_{\mathsf c})^*
       \mathsf N_{\mathsf c}^{-1}\mathsf Q'_{\mathsf c}
                                        \mathsf G'_{\mathsf c}.
\]
Write $e_1J=(J,0,0,0)$. With this forcing, the last three equations give
$w=(I-\mathsf R_{\mathsf c})D^*A$, and the first becomes
$\mathsf A_{\mathsf c}[S]A=J$. This identifies the desired
vector block once the system is inverted.

Replace the first region by the whole torus and set $U=I$;
write $\mathbb K_0$ for this reference system. Its distance
is unchanged, since weighted length outside $\Omega_1$ has
weight $\eta^{-1}$ for either first region. At the identity,
the corner transports are trivial and the dimensionless bond
average is $\mathsf mQ_0$, where $Q_0$ is the physical flat
average of \cref{sec:trivial-constrained-system} and
$(\mathsf mr)(c)=\xi_cr(c)$. In fact, differentiation of the
endpoint factors cancels the transverse comb terms, and
concatenation of the remaining straight segments gives the
physical average multiplied by $\xi_c$. The two row pairings
then give
\[
 Q_{\mathsf c,0}^*=Q_0^{\mathsf T}\mathsf m^{-3},
 \qquad Q_{\mathsf c,0}^*Q_{\mathsf c,0}
                     =Q_0^{\mathsf T}\mathsf m^{-2}Q_0.
\]
The scalar averages, scalar weights and projections agree as
well. Thus the flat four-component system is the restriction
of $\mathbb K_0^{[4]}$ in \cref{I-64-24} to $p=q=r=0$.
For forcing $(f_1,f_2,f_3,f_4,0,0,0)$ and $X_5=0$,
\eqref{I-63-69} gives precisely this restriction. The retained
seminorms are the ones displayed above. The scalar subsystems
are the corresponding scalar systems there. Consequently there
are $B_1,\delta_1>0$, depending only on $L$, such that
\[
 |\mathbb E_0F|_y^0
       \le B_1\sup_z e^{-\delta_1d(y,z)}|F|_z',
 \qquad \mathbb E_0=\mathbb K_0^{-1},
\]
provided $M$ and $\rho$ meet the thresholds of that theorem.
The first-source-to-vector block, denoted by $G_0$, is the positive flat inverse
identified in \cref{I-64-13,I-17-5}.

Call tops $z,y$ near, and write $z\sim y$, if their cells
contain bare sites at sup distance at most $\eta$, or sites
in one reduced-row double block. The cell geometry gives
$d(y,z)\le6L$ and $L^{-1}\le\xi_z/\xi_y\le L$.
These are the only comparisons needed in the local estimates.


\subsection{Cutoffs and locality}

For the pieces of \cref{def:regular-pieces}, with $s_j=M\xi_j$,
use the quadratic partition
\begin{equation}\label{eq:corner-quadratic-partition}
 h_\kappa=\frac{\phi_\kappa}
                  {(\sum_{\kappa'}\phi_{\kappa'}^2)^{1/2}},
 \qquad \phi_\kappa=h_{\kappa;0,1/12},
 \qquad \widetilde h_\kappa=h_{\kappa;1/3,2/3}.
\end{equation}
Here multiplications act at the base site of an index. In particular,
a touching bond whose initial endpoint is outside $\Omega_0$ is
multiplied at its final endpoint. The cover and two-layer argument
in \cref{lem:piece-geometry} apply to these same pieces. They give
\[
 \sum_\kappa h_\kappa^2=1,\qquad
 \operatorname{supp}h_\kappa\subset\kappa^{[1/12]},\qquad
 \widetilde h_\kappa=1\text{ on }\kappa^{[1/3]},\qquad
 \operatorname{supp}\widetilde h_\kappa\subset\kappa^{[2/3]}.
\]
At each point at most $N_G$ enlarged pieces $\kappa^{[1]}$ occur.
One may take $N_G=256$: at a point of layer $q$, their sizes are
among $q-2,q-1,q,q+1$, and there are at most $4^3$ centers at each
size. Differentiating $\phi_\kappa(\sum\phi_{\kappa'}^2)^{-1/2}$,
whose denominator is at least one, gives constants $C_G=C_G(L)$
such that, for $\theta=M/C_G\ge L^2$,
\begin{equation}\label{eq:corner-cutoff-differences}
 \begin{split}
 |h(x')-h(x)|&\le\frac{\eta}{\theta\xi(x)}
                    \qquad(|x-x'|_\infty\le\eta),\\
 |h(x+\eta e_\nu)+h(x-\eta e_\nu)-2h(x)|
       &\le\frac{\eta^2}{\theta\xi(x)^2}.
 \end{split}
\end{equation}
The sites in either formula belong to $\Omega_0$, and
$h$ is either cutoff. For the first estimate, integrate its first
derivative along $[x,x']\subset\Delta_0(x)\cup\Delta_0(x')$;
the two scales differ by at most $L$. Integrating the second
derivative against $(\eta-|t|)_+$ gives the second estimate.

For a regular piece of size $j$, its double enlargement lies in
$\mathcal L_j\cup\mathcal L_{j+1}$ and is more than $8M\eta$
from $\Omega_0^c$. For a boundary piece,
$\operatorname{dist}_\infty(\kappa^{[1]},\Omega_1)\ge M\eta$.
The first assertion follows from the three cases in the cover
argument: a deep cube of positive size, a subdivided buffer cube,
or a size-zero cube whose double enlargement meets $\Omega_1$.
For the second, $\kappa^{[2]}$ and $\Omega_1$ are disjoint unions
of $M$-cubes, and their grid forces the displayed margin.

On a regular $\kappa^{[1]}$, the original and reference indices
have the same cells, positions and base sites. A corner double
block meeting $\kappa^{[3/4]}$ has level $j$ or $j+1$ and lies
in $\kappa^{[1]}$. Indeed, lower-level endpoints lie outside
$\Omega_j$, while the gap from $\Omega_{j+1}^c$ excludes levels
at least $j+2$. Its diameter is less than
$2\xi_{j+1}\le s_j/12$ when $M\ge24L$.
Thus either system, acting on a field supported in $\kappa^{[r]}$,
$0\le r\le2/3$, produces a field supported in
$\kappa^{[r+1/12]}$ and uses $U$ only inside $\kappa^{[1]}$.
Conversely, its values in $\kappa^{[r]}$ depend only on the
field in $\kappa^{[r+1/12]}$. The scalar averages couple one
cell, the row energy couples bonds in one double block, and
the differential and Hessian terms have smaller range.

\subsection{Comparison on a regular piece}

Extend a local gauge $v_\kappa$ by the identity outside
$\kappa^{[1]}$, and let $\mathcal U_\kappa$ act on a fine
field by $\operatorname{Ad}(v_\kappa)$ at its initial site.
On a level-$q$ scalar top it acts at $o_{\mathsf c}^{(q)}(y)$,
and on a bond row at $o_{\mathsf c}^{(q)}(c_-)$. The position
used by a cutoff remains the top or the assigned base site.
Covariance of the corner transports follows by cancellation of
the intermediate gauge factors. Differentiating the straightened
average gives the same covariance for its bond kernel. The folded
extensions commute with these gauge changes. Together with gauge
invariance of the Wilson action, this proves covariance of
$\mathbb K_{\mathsf c}$ and preservation of supports and pairings.
The covariant seminorms are preserved as well.

Put $\widehat\varepsilon=L^2\varepsilon$. On a regular piece
of size $j$, all relevant cells have level $j$ or $j+1$, so
\[
 \xi_q|A_\kappa|\le\widehat\varepsilon,
 \qquad \xi_q^2|\nabla^\eta A_\kappa|
                                  \le\widehat\varepsilon.
\]
In the plaquette product the terms of first order are the two
ordinary differences; the remaining terms are bounded by
$16\eta^2\sup|A_\kappa|^2$. Thus, for $\varepsilon\le1/16$,
\[
 |\partial U^{v_\kappa}(p)-I|
   \le(2\varepsilon+16\varepsilon^2)\eta^2\xi_j^{-2}
   \le3\widehat\varepsilon\eta^2\xi_q^{-2}.
\]
We take
\[
 \varepsilon_0\le L^{-2}
       \min\{(4L^2)^{-1},\min(1,c_{\mathsf c})/3\}.
\]
Thus \cref{I-32-37} applies on every required double block.

For a field supported in $\kappa^{[2/3]}$, identify the original
and reference indices there and set
\[
 V_\kappa
    =\mathbb K_0-\mathcal U_\kappa\mathbb K_{\mathsf c}
                                      \mathcal U_\kappa^{-1}.
\]
There is $C_V=C_V(L)$ such that
\begin{equation}\label{eq:corner-local-defect}
 |V_\kappa\Psi|_y'
       \le C_V\varepsilon\max_{z\sim y}|\Psi|_z^0.
\end{equation}
Its output is supported in $\kappa^{[3/4]}$.
Here and below the corresponding scalar estimates are obtained
by deleting the absent entries; all differential stencils are
inside $\kappa^{[1]}$, where the two frames agree.

We justify the estimate by grouping the differential and
averaging terms. The identities of \cref{I-65-25}, with
$\xi=\eta$ and $A'=A_\kappa$, apply to the first differences
and their adjoints. In the Laplacian, use the cancellation
\eqref{I-65-27}. Its zeroth-order coefficient is the sum of
$e(x)+e'(x)$ and $e'(x-\eta e_\mu)-e'(x)$; the first is
bounded by $4\eta|A_\kappa(x)|^2$ and the second by twice
the ordinary increment of $A_\kappa$. Consequently the
Laplacian defect sends value/gradient weight $(s,s-1)$ to
weight $s-2$ with bound $C\widehat\varepsilon$, for
$s=-1,0,-2$. This retains the first-difference hypothesis
and incurs no negative power of $\eta$ in the constant.
The ambient plaquettes in \cref{I-65-4} are the same on this
interior stencil. That formula and the displayed curvature
bound give the same estimate for
$\mathsf H_S+DD^*$ relative to the componentwise flat
Laplacian.

For the scalar average, \cref{I-32-37} bounds the difference
of each transport from the identity by
$t\le C\widehat\varepsilon$. Since its coefficients have
total mass one, both the scalar average and its weighted
adjoint differ from their flat counterparts by at most $t$.
Thus the defect of
$(\mathsf Q'_{\mathsf c})^*\mathsf a\mathsf Q'_{\mathsf c}$
is at most $2t\xi_q^{-2}$ on one cell. The linear scalar
average and adjoint in the other two equations satisfy the
corresponding bounds without this factor.

Write $Q=\mathsf Q_{\mathsf c}$ and $Q_0=Q_{\mathsf c,0}$
in the row calculation. The weighted adjoint gives
\[
 \begin{split}
 ((Q^*Q-Q_0^*Q_0)A)(b)
   =\eta^{-3}\sum_c\xi_c^{-1}
       \bigl[&(Q(c,b)-Q_0(c,b))^{\mathsf T}(QA)(c)\\
            &+Q_0(c,b)^{\mathsf T}((Q-Q_0)A)(c)\bigr].
 \end{split}
\]
The absolute row and entry bounds from \cref{I-32-37} are
$2\xi_c$ and $2\eta^3\xi_c^{-2}$; their defects have the
same scale factors times $C\widehat\varepsilon$. A bond
in this piece belongs to at most twelve relevant rows, at
levels $j,j+1$. Their cells have scales within a factor $L$
of the output scale and tops near its top. The displayed
formula is therefore bounded by
$C\widehat\varepsilon\xi_y^{-3}
\max_{z\sim y}|\Psi|_z^0$. The added boundary mass is absent
on a regular piece. These estimates prove
\eqref{eq:corner-local-defect} in every forcing slot.

The difference between covariant and ordinary first differences
is $e$ times a neighboring value. Its contribution to a solution
seminorm is at most $2L^2\widehat\varepsilon$ times a neighboring
seminorm. Our cap on $\varepsilon_0$ consequently gives
\[
 |\mathcal U_\kappa^{-1}\Psi|_y
       \le\tfrac32\max_{z\sim y}|\Psi|_z^0,
 \qquad
 |\mathcal U_\kappa\Psi|_y^0
       \le\tfrac32\max_{z\sim y}|\Psi|_z.
\]

The flat commutator with either cutoff satisfies
\begin{equation}\label{eq:corner-flat-commutator}
 |[\mathbb K_0,h]\Psi|_y'
       \le C_C\theta^{-1}\max_{z\sim y}|\Psi|_z^0,
 \qquad \theta=M/C_G,
\end{equation}
for $C_C=C_C(L)$. To see the second-order cancellation, put
$k_\nu(x)=\eta^{-1}(h(x+\eta e_\nu)-h(x))$. Then
\[
 \begin{split}
 [\Delta,h]f(x)
   =-\sum_\nu\bigl(&k_\nu(x)\nabla_\nu^\eta f(x)
          +k_\nu(x-\eta e_\nu)
                         \nabla_\nu^\eta f(x-\eta e_\nu)\bigr)\\
      &-\eta^{-1}\sum_\nu
             (k_\nu(x)-k_\nu(x-\eta e_\nu))f(x).
 \end{split}
\]
The first and second cutoff differences bound these terms with
the required weights. The first-order commutators contain only
one cutoff difference. For the averages, the oscillation of $h$
is at most $3/\theta$ on a cell and $4L/\theta$ on a double
block, by summing its increments along coordinate paths.
Insert these oscillations in the scalar and bond kernels above.
At most eighteen flat reduced rows lie over a bond, so their
sum has a bound depending only on $L$. This proves
\eqref{eq:corner-flat-commutator}; we may enlarge $C_C$ to
be at least $300L^4$. The product rule also gives
$|h\Psi|_y^0\le2\max_{z\sim y}|\Psi|_z^0$.


\subsection{Boundary inverses and positivity}

\begin{proof}[Proof of \cref{I-32-62}]
\emph{Boundary pieces.}
For a boundary piece $\kappa$, let
\[
 E^\kappa=E_0\cap\kappa^{[1]},\qquad
 S^\kappa=\{b\in S:\{b_-,b_+\}\cap E_0\subset E^\kappa\}.
\]
Let $\chi_\kappa$ restrict the vector component to $S^\kappa$
and the three scalar components to $E^\kappa$, and set
$\mathbb K_\kappa=\chi_\kappa\mathbb K_{\mathsf c}\chi_\kappa$ on its
range. For the scalar systems delete the absent components.
Use the reduced seminorm
\[
 |\Psi|_y^{\mathrm r}
   =\max\{\eta|A|_y,|v|_y,\eta^2|w|_y,\eta^4|\sigma(y)|\}.
\]
There are $B_b,\delta_b>0$, depending only on $L$, such that,
for $M$ sufficiently large, these compressions are invertible and
\begin{equation}\label{eq:corner-boundary-inverse}
 |\mathbb E_\kappa F|_y^{\mathrm r}
       \le B_b\sup_z e^{-\delta_b d(y,z)}|F|'_z,
 \qquad \mathbb E_\kappa=\mathbb K_\kappa^{-1}\chi_\kappa.
\end{equation}
For the vector system we use the additional ambient plaquette
bound. Its first-source-to-vector block is
\[
 \mathsf G_\kappa\chi_{S^\kappa},\qquad
 \mathsf G_\kappa=
       (\chi_{S^\kappa}\mathsf H_S\chi_{S^\kappa}+\eta^{-2})^{-1},
\]
where $\mathsf G_\kappa$ is positive definite.
If $|h|\le1$ and $h=0$ outside $\kappa^{[2/3]}$, then
\[
 |h\Psi|_y\le2\max_{z\sim y}|\Psi|_z^{\mathrm r}.
\]
If $h$ also satisfies \eqref{eq:corner-cutoff-differences}, then
\[
 |[\mathbb K_\kappa,h]\Psi|'_y
       \le25\theta^{-1}\max_{z\sim y}|\Psi|_z^{\mathrm r}.
\]
For a field supported in $\kappa^{[r]}$, $0\le r\le2/3$,
the compression agrees with the full system and its output
is supported in $\kappa^{[r+1/12]}$. The reverse dependence
statement holds on the same two sets.
We prove these local assertions directly.
The margin from $\Omega_1$ makes every site of $E^\kappa$
a level-zero top. There $\mathsf Q'_{\mathsf c}=I$, $\mathsf Q_{\mathsf c}=\eta I$, and the
row energy, including the touching-bond term, is $\eta^{-2}I$.
No positive-level row meets $S^\kappa$: its double block is
either contained in $\Omega_1$ or lies within $3L\eta/2$ of
an endpoint in $\Omega_1$. The definition of $S^\kappa$ also
gives $D^*\chi_{S^\kappa}=\chi_{E^\kappa}D^*\chi_{S^\kappa}$.
Put
\[
 D_\kappa=\chi_{S^\kappa}D\chi_{E^\kappa},\qquad
 A'_\kappa=\chi_{E^\kappa}\Delta_U\chi_{E^\kappa}+\eta^{-2},
 \qquad W_\kappa=\chi_{S^\kappa}\mathsf H_S\chi_{S^\kappa}+\eta^{-2}.
\]
Scalar zero extension retains all six diagonal contributions
of $\Delta_U$. The compressed vector system is
\[
 ((W_\kappa+D_\kappa D_\kappa^*)A-D_\kappa w,
    A'_\kappa v+w-D_\kappa^*A,
    A'_\kappa w-\sigma,v).
\]
Consequently its solution is
\begin{equation}\label{eq:corner-boundary-solution}
 \begin{split}
 v&=f_4,\qquad
 A=W_\kappa^{-1}[f_1+D_\kappa(f_2-A'_\kappa f_4)],\\
 w&=f_2+D_\kappa^*A-A'_\kappa f_4,\qquad
 \sigma=A'_\kappa w-f_3.
 \end{split}
\end{equation}
For the three-slot scalar system set $A=0$; for the one-slot
system use $w=(A'_\kappa)^{-1}f_3$.

The plaquette Hessian formula in \cref{I-65-4} gives
$\mathsf H_S\ge-12\alpha\eta^{-2}$. Indeed, each of the four bond
variables occurs in three commutator pairs in one plaquette,
and each bond belongs to four plaquettes. All contributing
plaquettes have a corner in $\Omega_0$, which is precisely
where the additional hypothesis is used. Since
$12\alpha\le1/2$, we have
\[
 W_\kappa\ge\tfrac12\eta^{-2},\qquad A'_\kappa\ge\eta^{-2}.
\]
The off-diagonal row sum of $\eta^2W_\kappa$ is at most $13$,
and a nonzero entry joins base sites at distance $d\le2$.
For $\eta^2A'_\kappa$ the corresponding bounds are $6$ and $1$.
Exponential conjugation therefore gives, with $t=1/300$,
\[
 |W_\kappa^{-1}(b,b')|
       \le4\eta^2e^{-t d(\mathsf x(b),\mathsf x(b'))},\qquad
 |(A'_\kappa)^{-1}(x,x')|
       \le2\eta^2e^{-t d(x,x')}.
\]
Indeed, the conjugation argument proving \eqref{I-14-34}
applies: the error row and column sums are bounded by
$13(e^{2t}-1)\le1/4$ and $6(e^t-1)\le1/2$, respectively,
half the corresponding spectral lower bounds.

At most six bonds have a given base site, including the
incoming touching bonds. Hence summing the first kernel bound
over the bonds at one site costs at most $24$, rather than
$12$. The elementary bounds
\[
 \eta^2|A'_\kappa f|_y\le13\max_{z\sim y}|f|_z,
 \quad\eta|D_\kappa f|_y\le2\max_{z\sim y}|f|_z,
 \quad\eta|D_\kappa^*A|_y\le6\max_{z\sim y}|A|_z
\]
and \eqref{eq:corner-boundary-solution}, followed by the
weighted summation estimate used in
\eqref{eq:trivial-system-weighted-sum}, prove
\eqref{eq:corner-boundary-inverse} with $\delta_b=t/2$.
The first inverse block and its positivity also follow from
\eqref{eq:corner-boundary-solution}.

All the compressed entries couple base sites within $\eta$.
For $[W_\kappa,h]$, the off-diagonal sum is at most
$12\eta^{-2}$ even without a curvature bound: the common
plaquette entry is a product of isometries with
$Y\mapsto c_pY\pm\tfrac12\eta^2i[Y,F_p]$, whose norm is at
most one since $c_p^2+\eta^4|F_p|^2=1$.
The corresponding sums for $D_\kappa D_\kappa^*$,
$A'_\kappa$ and $D_\kappa^*$ are bounded by
$12\eta^{-2},6\eta^{-2},6\eta^{-1}$.
In $[D_\kappa,h]$ one endpoint is the base site, so only
one difference remains. Multiplication by the source weights
bounds the four slots by $25,12,6,0$ times
$\theta^{-1}\max_{z\sim y}|\Psi|_z^{\mathrm r}$.

The range bound proves the two locality assertions.
Differences of $h\Psi$ use only two neighboring values and
their unitary transports, giving the factor two in the full
seminorm. For touching bonds the base sites of translated
bonds are still at sup distance at most $\eta$: the only
possible longer pair would have an intervening exterior site
with both axial neighbors interior. This is excluded because
the complementary $M$-cubes contain at least two bare sites
in that direction. Thus the argument covers both frames.

\emph{The parametrix.}
Use the flat inverse $\mathbb E_0^{[i]}$, $i=1,2,3$, identified
above, and the local defect estimate. For a fixed system suppress
$[i]$; for the vector system fix either frame and assume the
ambient plaquette bound. Let $B_1,\delta_1>0$ be the common
flat bounds, let $B_b,\delta_b$ be the boundary bounds just proved, and put
\[
 \delta_2=\min\{\delta_1,\delta_b\},\quad
 B_2=\max\{1,B_1,B_b\},\quad k_0=e^{6L\delta_2},\quad
 B_3=6N_Gk_0^2B_2.
\]
We choose $M_\diamond$ large enough for all the preceding
geometric, flat-inverse and summation thresholds, and choose
$\varepsilon_\diamond$ small enough for the local comparison.
Increase $M_\diamond$ and decrease $\varepsilon_\diamond$ so that
\begin{equation}\label{eq:corner-total-smallness}
 C_*(C_V\varepsilon+C_CC_G/M)\le\tfrac12,
 \qquad C_*=27L^2N_Gk_0^2B_2B_3.
\end{equation}
All these choices depend only on $L$; in particular we require
$\delta_2\rho\ge8\log L$ whenever $\rho\ge LM$.

For a regular piece let $\mathcal U_\kappa$ be its gauge map,
$V_\kappa=\mathbb K_0-\mathcal U_\kappa\mathbb K_{\mathsf c}
\mathcal U_\kappa^{-1}$, and $C_h=[\mathbb K_0,h]$. Define
\begin{equation}\label{I-32-61}
 \begin{aligned}
 P_\kappa&=\mathcal U_\kappa^{-1}h_\kappa
              \mathbb E_0h_\kappa\mathcal U_\kappa,
 &R_\kappa&=\mathcal U_\kappa^{-1}
        (V_\kappa h_\kappa-C_{h_\kappa})
                         \mathbb E_0h_\kappa\mathcal U_\kappa,
 \end{aligned}
\end{equation}
and on a boundary piece set
$P_\kappa=h_\kappa\mathbb E_\kappa h_\kappa$ and
$R_\kappa=-[\mathbb K_\kappa,h_\kappa]
                         \mathbb E_\kappa h_\kappa$.
The local identifications and support margins give
\[
 \mathbb K_{\mathsf c}P_\kappa=h_\kappa^2-R_\kappa,
 \qquad \mathbb K_{\mathsf c}P=I-R,
 \qquad P=\sum_\kappa P_\kappa,\quad R=\sum_\kappa R_\kappa.
\]
For example, on a regular piece expand
$(\mathbb K_0-V_\kappa)(h_\kappa\Phi)$, where
$\Phi=\mathbb E_0h_\kappa\mathcal U_\kappa F$.
Both remainder terms are supported in $\kappa^{[1/6]}$;
the boundary identity follows from the compression and locality
proved above.

In the exponential supremum classes of
\eqref{eq:trivial-system-sup-decay}, the local estimates give
\[
 P_\kappa\in\mathfrak O(\delta_2,3k_0^2B_2),\qquad
 R_\kappa\in\mathfrak O
   (\delta_2,2k_0^2B_2(C_V\varepsilon+C_C/\theta)).
\]
The factor three combines the cutoff bound two with the
gauge norm comparison $3/2$. Two neighboring-top shifts cost
$k_0^2$, since each has distance at most $6L$. On boundary
pieces the corresponding constants are $2k_0B_b$ and
$25k_0B_b/\theta$, bounded by the displayed choices.
A nonzero output seminorm has its top in $\kappa^{[1]}$,
so at most $N_G$ pieces contribute at any top. Thus
\[
 P\in\mathfrak O(\delta_2,B_3/2),\qquad
 R\in\mathfrak O(\delta_2,1/2).
\]
The composition rule following \eqref{eq:trivial-system-sup-decay}
now sums the geometric series. Since $\mathbb K_{\mathsf c}P(I-R)^{-1}=I$
on a finite-dimensional space, the full system is invertible and
\begin{equation}\label{I-32-64}
 \mathbb E_{\mathsf c}=\mathbb K_{\mathsf c}^{-1}=P(I-R)^{-1},\qquad
 |\mathbb E_{\mathsf c}F|_y\le B_3\sup_z e^{-\delta_2d(y,z)}|F|'_z.
\end{equation}

For $i=1$, the nonnegative symmetric operator $\mathsf A'_{\mathsf c}$ is
therefore positive definite. For $i=2$, this also makes
$\mathsf N_{\mathsf c}=\mathsf Q'_{\mathsf c}(\mathsf G'_{\mathsf c})^2(\mathsf Q'_{\mathsf c})^*$ positive definite, and
\eqref{I-32-64} gives the scalar three-component estimate used below. The elimination identity gives, for $i=3$,
\[
 [\mathbb E_{\mathsf c}e_1J]_A=\mathsf G_{\mathsf c,S}J.
\]
For the scalar inverse, the solution value has factor
$\xi_y^{-2}$ and the forcing has factor $\xi_z^4$ in
\eqref{I-32-64}; its gradient has factor $\xi_y^{-3}$.
For the vector inverse these three factors are
$\xi_y^{-1}$, $\xi_z^3$ and $\xi_y^{-2}$, respectively.
Absorb the input scale
pointwise under the supremum using
\[
 (\xi_z/\xi_y)^s e^{-\delta_2d(y,z)}
       \le L^4e^{-\delta_2d(y,z)/2},\qquad |s|\le4,
\]
by the scale-distance bound and $\delta_2\rho\ge8\log L$.
Substitution into \eqref{I-32-64} proves
\eqref{eq:corner-scalar-total} and
\eqref{eq:corner-vector-total} with
$B_\diamond=L^4B_3$ and $\delta_\diamond=\delta_2/2$.

It remains to prove positivity of $\mathsf A_{\mathsf c}[S]$.
Let $P_A$ be the first-source-to-vector block of $P$. Then
\[
 P_A=\sum_{\kappa\text{ regular}}
    \mathcal U_\kappa^{-1}h_\kappa G_0h_\kappa\mathcal U_\kappa
   +\sum_{\kappa\text{ boundary}}h_\kappa G_\kappa h_\kappa.
\]
The flat $G_0$ and each $G_\kappa$ are positive definite;
unitary gauge covariance and $\sum h_\kappa^2=1$ show that
$P_A$ is symmetric positive definite.

Put
\[
 \|a\|_1=\max_y\{\xi_y|a|_y,\xi_y^2|\nabla^Ua|_y\},
 \qquad \|\Psi\|=\max_y|\Psi|_y.
\]
If $\Psi=\mathbb E_{\mathsf c}e_1J=(a;v,w,\sigma)$, elimination and
the scalar three-slot bound imply
\begin{equation}\label{eq:corner-slaved-components}
 \|\Psi\|\le3L^2B_3\|a\|_1.
\end{equation}
At an interior site, $D^*a$ is a sum of three transported
backward first differences, so that
$\xi_y^2|D^*a|_y\le3L^2\|a\|_1$.
A site with an exterior neighbor, and all its interior
neighbors, lie in the zero layer by the gap from $\Omega_1$.
There $D^*a$ has at most six incident-bond terms, and
$\eta^2|D^*a|\le6\|a\|_1\le3L^2\|a\|_1$.
Applying the three-slot inverse to $(D^*a,0,0)$ proves
\eqref{eq:corner-slaved-components} for both frames.

The operator $\mathsf A_{\mathsf c}[S]P_A$ is similar to the symmetric
matrix $P_A^{1/2}\mathsf A_{\mathsf c}[S]P_A^{1/2}$. Let
$\mathsf A_{\mathsf c}[S]P_AJ=zJ$ with $J\ne0$, and set
$a=\mathsf G_{\mathsf c,S}J$. Thus $P_AJ=za$.
On a regular piece apply $\mathbb E_0$ to
$\mathbb K_0(h_\kappa\mathcal U_\kappa
\widetilde h_\kappa\Psi)$ and then multiply by
$\mathcal U_\kappa^{-1}h_\kappa$. Locality and
$\widetilde h_\kappa=1$ near the support of $h_\kappa$ give
\[
 h_\kappa^2\Psi=P_\kappa e_1J+\Theta_\kappa,
 \quad
 \Theta_\kappa=\mathcal U_\kappa^{-1}h_\kappa\mathbb E_0
   (h_\kappa V_\kappa+C_{h_\kappa})
                    \mathcal U_\kappa\widetilde h_\kappa\Psi.
\]
For a boundary piece the same identity holds with
$\Theta_\kappa=h_\kappa\mathbb E_\kappa
[\mathbb K_\kappa,h_\kappa]\widetilde h_\kappa\Psi$.
The two cutoff bounds and the local comparison give
\[
 |\Theta_\kappa|_y
       \le9B_2(C_V\varepsilon+C_C/\theta)\|\Psi\|.
\]
For boundary pieces use the reduced norm before applying
$\mathbb E_\kappa$; this gives $50B_b/\theta$, which is
bounded by the displayed constant. At most $N_G$ terms
have nonzero seminorm at a given top. Summing and taking the
vector component yields
\begin{equation}\label{I-32-65}
 (1-z)a=\Bigl[\sum_\kappa\Theta_\kappa\Bigr]_A,
 \qquad
 |1-z|\|a\|_1
       \le C_*(C_V\varepsilon+C_C/\theta)\|a\|_1
       \le\tfrac12\|a\|_1.
\end{equation}
Since $a\ne0$, every eigenvalue $z$ is at least $1/2$.
The symmetric matrix $P_A^{1/2}\mathsf A_{\mathsf c}[S]P_A^{1/2}$
and hence $\mathsf A_{\mathsf c}[S]$ are positive definite.
\end{proof}

\section{The inverse on the original region}
\label{sec:regional-inverse}

The action on a region contains only its inner plaquettes. To
pass from the ambient inverse to this action, we first apply
\cref{I-32-62} on a collar of the positive-level regions, where
the two Hessians agree. A second inverse near the outer boundary
then gives the inverse on the whole region. Exponential
conjugation supplies its decay estimates.

Use the preceding level-$n$ units and a separated sequence with
parameters $\mathfrak M,R_1$, where $R_1\ge L$ and both are
powers of $L$. Put $X=\Omega_0\cap\mathbb T^{(0)}$.
Let $\mathfrak b$ and $\mathfrak p$ be the inner positive
bonds and inner plaquettes, and let $D$ be the covariant
gradient from $X$ to $\mathfrak b$. Thus
\[
 \Delta_U=D^*D+\eta^{-2}n_{\partial},
\]
where $n_{\partial}(x)$ counts the missing bare neighbours.
Let $H_U$ be the Hessian of $\eta^{-1}A_{\Omega_0}$; only
plaquettes in $\mathfrak p$ enter this Hessian. Write $Q',Q$
for the corner scalar and reduced bond averages of the preceding
section. Their fine, scalar-top and bond-row pairings have
weights $\eta^3,\xi_y^3,\xi_c^{-1}$. In particular, the
level-zero bond average is $\eta I$.

Let $\Xi$ multiply by the local scale and retain the scalar
multiplier $a|_{\Lambda_j}=\varkappa_j\xi_j^{-2}$ from the
preceding section. Define $R_\perp$ to be the orthogonal
projection onto $\Delta_U\ker Q'$, and set
\[
 A'=\Delta_U+Q'^*aQ',\qquad
 A^\perp=H_U+DR_\perp D^*+Q^*Q.
\]
Once the indicated inverses exist, put
\[
 \begin{aligned}
 G'&=(A')^{-1},&N&=Q'G'^2Q'^*,&C&=N^{-1},\\
 P&=G'Q'^*CQ'G',&R&=I-P,&
 A_0&=H_U+DRD^*+Q^*Q,\qquad G_0=A_0^{-1}.
 \end{aligned}
\]
We will prove that $R=R_\perp$ and $A_0=A^\perp$.

All fine fields are extended by zero. In $|v|_y$, bonds and
differences are assigned by their initial or base site, and
$|v|_{[s]}=\sup_y\xi_y^{-s}|v|_y$. For a bond field,
\[
 (\nabla^U_\nu A_\mu)(x)=
 \begin{cases}
 \eta^{-1}\bigl(\operatorname{Ad}(U_\nu(x))
 A_\mu(x+\eta e_\nu)-A_\mu(x)\bigr),&
 \langle x,x+\eta e_\nu\rangle\in\mathfrak b,\\
 -\eta^{-1}A_\mu(x),&\text{otherwise}.
 \end{cases}
\]
The zero extension also applies when the translated bond is
missing. Scalar gradients are evaluated on $\mathfrak b$.

\begin{theorem}\label{I-42-86}
Fix odd $L\ge13$. There are a power $M_{\mathrm{ps}}$ of $L$
and constants $g_{\mathrm p},B_{\mathrm{ps}},a_{\mathrm{ps}}>0$
and $0<\delta_{\mathrm{ps}}\le1$, depending only on $L$, with
the following property. Suppose
\[
 \mathfrak M\ge M_{\mathrm{ps}},\qquad
 R_1\mathfrak M\ge g_{\mathrm p},
\]
and that the level-$n$ torus has at least $6LM_{\mathrm{ps}}$
sites in each direction. Let $U$ be real on the inner bonds.
For some $0<\alpha\le a_{\mathrm{ps}}$, assume
\[
 |\partial U(p)-I|<\alpha\eta^2\mathsf L_{\mathbf B}(p)^{-2},
 \qquad |(D_U^{\mathsf T}\partial U)(b)|
       <\alpha\eta^2\mathsf L_{\mathbf B}(b)^{-3}
\]
at every inner plaquette $p$ and every full positive bond $b$.
The second expression is the matrix-valued plaquette divergence
of \eqref{I-2-5}, including its scalar part.
Then $A',N,A_0$ are positive definite for their stated pairings.
For $f$ supported over one top $y'$,
\begin{align*}
 |G'f|_y&\le B_{\mathrm{ps}}\xi_y^2e^{-\delta_{\mathrm{ps}}d(y,y')}
                  \sup|f|,&
 |DG'f|_y&\le B_{\mathrm{ps}}\xi_y e^{-\delta_{\mathrm{ps}}d(y,y')}
                  \sup|f|,\\
 |G_0f|_y&\le B_{\mathrm{ps}}\xi_y^2e^{-\delta_{\mathrm{ps}}d(y,y')}
                  \sup|f|,&
 |\nabla^UG_0f|_y&\le B_{\mathrm{ps}}\xi_y e^{-\delta_{\mathrm{ps}}d(y,y')}
                  \sup|f|.
\end{align*}
The scalar and vector lines apply on their respective spaces.
The counting kernel of $C$ satisfies
\[
 |C(y,y')|\le B_{\mathrm{ps}}e^{-\delta_{\mathrm{ps}}d(y,y')}
       \min\{\xi_y^{-4},\xi_y^{-7/2}\xi_{y'}^{-1/2}\}.
\]
All constants are uniform in the depth, lattice spacing and regions.
\end{theorem}

\begin{proof}
We use the summation estimates \eqref{eq:global-scale-ratio}
and \eqref{I-13-122} in the following form. For $a>0$, put
\[
 c_{\mathrm{sum}}(a)=L^4c_S(a/2),\qquad
 w_{\mathrm{sum}}(a)=\max\{\mathsf r_S(a/2),8a^{-1}\log L\}.
\]
If $R_1\mathfrak M\ge w_{\mathrm{sum}}(a)$, then, for $|s|\le4$,
\[
 \sum_{y'}e^{-a d(y,y')}(\xi_y/\xi_{y'})^s
       \le c_{\mathrm{sum}}(a).
\]
Indeed the scale ratio costs at most $L^4e^{a d/2}$.
All occurrences below of a summation threshold refer to
$w_{\mathrm{sum}}$; the final choice of the gap includes them.

\emph{The collar.}
Fix a power $M=M_{\mathrm{ps}}\ge\max\{M_\diamond,L^2\}$ of $L$,
and put $\rho=LM$. We will choose a power $g_0\ge L^2M$ and
require $R_1\mathfrak M\ge Lg_0$. Suppose first that
$\Omega_1\ne\emptyset$. Let $\Omega_0^{\mathrm a}$ be the union
of the level-zero $M$-cubes at sup distance at most $g_0\eta$
from $\Omega_1$. It contains this neighbourhood and lies within
$(g_0+M)\eta\le2g_0\eta$ of $\Omega_1$. Hence
\[
 \operatorname{dist}_\infty((\Omega_0^{\mathrm a})^c,\Omega_1)>g_0\eta,
 \qquad
 \operatorname{dist}_\infty(\Omega_0^{\mathrm a},\Omega_0^c)
                     >(L-2)g_0\eta.
\]
The auxiliary sequence $(\Omega_0^{\mathrm a},\Omega_1,\ldots,\Omega_n)$
satisfies the hypotheses on regions in \cref{I-32-62} with gap
$\rho=LM$: the higher regions are unions of $M$-cubes by refinement,
and their original gaps exceed $\rho\xi_j$. The torus floor is
$6\rho$. Write $X_{\mathrm a}$ for its bare sites and
$S_\alpha,S_\beta$ for its inner and touching positive bonds.

Every ambient plaquette touching the collar is an original inner
plaquette, and every bond in $S_\beta$ is full and belongs to
$\mathfrak b$. Also $\Omega_1$ and its first two neighbour layers
lie in the collar. Extend $U$ by the identity outside the original
inner bonds. On the collar and its neighbouring plaquettes this
extension changes neither tested quantity nor the local block side:
the positive members of the sequence are unchanged. In particular,
the extended field satisfies the ambient curvature hypothesis of
\cref{I-32-62} on the collar.

Here is the required local gauge check. A cube $Q\subset\Omega_q$
of side $l\xi_q$, where $q\ge1$ and $l$ is a multiple of $L$,
has a concentric enlargement of side $(l+4L)\xi_q$. Provided
this is smaller than half the torus period, its entire $2\eta$
neighbourhood lies in $\Omega_{q-1}$: its points are at distance
at most $(2L^2+2)\xi_{q-1}$ from $Q$, less than the separation.
The local block sides of the tested plaquettes and bonds are
therefore at least $\xi_q/L$. In level-$q$ units the hypotheses
of \cref{I-28-75} hold with parameter $L^3\alpha$; the divergence
acquires a factor $\xi_q$ under this change of units.
Apply that theorem on the smaller enlargement of side
$(l+4)\xi_q$ and restrict to $Q$. With its constants $B_*,c_*$,
put
\[
 B_{\mathrm p}=\max\{1,LB_*/21\},\qquad
 c_{\mathrm p}=\min\{1,21c_*/L\}.
\]
The resulting potential satisfies
\begin{equation}\label{eq:regional-local-gauge}
 \xi_q|A|,\ \xi_q^2|\nabla^\eta A|,\
 \xi_q^3|D_{\mathbf1}^{\mathsf T}D_{\mathbf1}A|
       <21L^2B_{\mathrm p}l\alpha,
 \qquad 21L^2l\alpha\le c_{\mathrm p}.
\end{equation}
For $q=0$ use the comb gauge of \cref{I-10-51} in the
enlarged box. Since it has at most $5l$ sites per side,
$|U^h-I|\le10l\alpha$. Real logarithms $U^h=e^{i\eta A}$
give $\eta|A|<16l\alpha$ and
$\eta^2|\nabla^\eta A|<32l\alpha$.
The ordinary curl--divergence is $\eta^{-2}$ times four
signed sums of four potential values. Its stencil lies in
the $2\eta$ neighbourhood of $Q$, hence in the enlarged box.
Thus
\[
 \eta^3|D_{\mathbf1}^{\mathsf T}D_{\mathbf1}A|
       <4\cdot4\cdot16l\alpha=256l\alpha
       <21L^2B_{\mathrm p}l\alpha.
\]
This proves \eqref{eq:regional-local-gauge} also for $q=0$.

Now let $Q$ have side $sM\xi_j$, $1\le s\le12$, and meet
the auxiliary layer $j$. For $j=0$ use $q=0,l=sM$ above;
the enlargement lies within $14M\eta\le g_0\eta$ of the collar.
For $j\ge1$ use $q=j-1,l=sLM$; the enlargement lies within
$14LM\xi_{j-1}$ of $\Omega_j$ and hence in $\Omega_{j-1}$.
Its side is at most $16M$ in the present units, less than half
the period. Changing the two estimates from scale $\xi_{j-1}$
to $\xi_j$ costs at most $L^2$. Thus these gauges satisfy
\[
 U^v=e^{i\eta A},\qquad
 \xi_j|A|,\ \xi_j^2|\nabla^\eta A|<\kappa_\square\alpha,
 \qquad \kappa_\square=252L^5B_{\mathrm p}M,
\]
if $252L^3M\alpha\le c_{\mathrm p}$. A regular piece's
$3M\xi_j$ enlargement either meets layer $j$, or lies wholly
in layer $j+1$. In the second case it is contained in a cube
of side $2M\xi_{j+1}$ meeting that layer; restrict its gauge.
Both bounds at scale $\xi_j$ follow. Choose the radius so that
\begin{equation}\label{eq:collar-gauge-radius}
 \alpha\le L^{-2}\min(1,c_{\mathsf c}),\qquad
 252L^3M\alpha\le c_{\mathrm p},\qquad
 \kappa_\square\alpha\le\varepsilon_\diamond.
\end{equation}
These are exactly the local gauge hypotheses of \cref{I-32-62}.

We next identify its operators. Positive-level cells and tops
are unchanged. Nested locality from \cref{I-32-37} identifies
their scalar transports and the restrictions of $Q'$; both
scalar Laplacians retain the full diagonal $6\eta^{-2}$.
The two kernels of $Q'$ coincide under zero extension: their
elements vanish in layer zero. Their Laplacian images are the
same subspace, supported in $\Omega_1$ and its bare neighbours.
Thus their orthogonal projections agree there.

The positive-level interior and outward rows also coincide.
Their double blocks lie in the collar, within $2L\eta$ of
$\Omega_1$ at level one, and inside $\Omega_{j-1}$ at higher
levels. Nested locality identifies their kernels, with the
same normalization $\eta\prod_{i<j}\mathbf Q_{\mathsf c}$
and row energy $\sum_c\xi_c^{-1}|QA(c)|^2$.
Every ambient plaquette contributing on $S_\alpha$ or $S_\beta$
is original inner, so the Hessian forms agree. On the support
of the projection all six incident bonds lie in $S_\alpha$;
the projected divergence forms agree as well. The level-zero
rows in $S_\beta\setminus S_\alpha$ contribute
$\xi_0^{-1}|\eta A(b)|^2=\eta|A(b)|^2$, precisely the
additional mass term in \cref{I-32-62}.

Consequently its scalar operator is the compression $A'_{\mathrm a}$
of $A'$ to $X_{\mathrm a}$, and its two vector operators are the
compressions $A_{\mathrm a}[S]$ of $A^\perp$ to $S$.
They are positive definite. With $B_0=\max(1,B_\diamond)$,
their inverses obey
\begin{equation}\label{eq:regional-collar-input}
 \begin{gathered}
 |G'_{\mathrm a}g|_{[2]}\le B_0\|g\|_\infty,\qquad
 |DG'_{\mathrm a}g|_{[1]}\le LB_0\|g\|_\infty,\\
 |G_{\mathrm a,S}g|_{[2]}\le L^2B_0\|g\|_\infty,\qquad
 |\nabla^UG_{\mathrm a,S}g|_{[1]}\le LB_0\|g\|_\infty.
 \end{gathered}
\end{equation}
Here scalar gradients are taken on $S_\alpha$, and vector
differences where both translated bonds lie in $S$. These follow
directly from the supremum bounds of \cref{I-32-62}. On a
touching bond whose initial endpoint is outside the collar both
endpoints have scale $\eta$, so changing its assigned endpoint
does not change the scale. In what follows write $G_{\mathrm a}$
for $G_{\mathrm a,S_\alpha}$.

\emph{Two cutoffs and the scalar inverse.}
Put $d_1(x)=\eta^{-1}\operatorname{dist}_\infty(x,\Omega_1)$
and $t=(g_0-5)/2$. Choose $\zeta_1,\zeta_2$ with
$\zeta_1^2+\zeta_2^2=1$, equal to $(1,0)$ for $d_1\le t$
and to $(0,1)$ for $d_1\ge2t$, by using the cosine and sine
of $\pi(d_1-t)/(2t)$ in between. Their difference on a bare
bond is at most $2/t$. If $\Omega_1$ is empty, use
$(\zeta_1,\zeta_2)=(0,1)$ and omit the collar terms.

Let $A'_{\mathrm l}$ be the compression to $\{d_1>t-2\}$.
There $A'=\Delta_U+\eta^{-2}$, so its diagonal is
$7\eta^{-2}$ and its off-diagonal row sum is at most
$6\eta^{-2}$. The maximum principle gives
$\|G'_{\mathrm l}g\|_\infty\le\eta^2\|g\|_\infty$.
Both cutoffs commute with $Q'^*aQ'$. Their commutators with
$\Delta_U$ are supported in $t-1<d_1<2t+1$, where all
neighbours are in layer zero and in the collar, and satisfy
\[
 |[\Delta_U,\zeta_i]v(x)|
 \le12t^{-1}\eta^{-2}\max_{x'\sim x}|v(x')|.
\]
Thus
\[
 P'=\zeta_1G'_{\mathrm a}\zeta_1
        +\zeta_2G'_{\mathrm l}\zeta_2,
 \qquad A'P'=I+E',\qquad
 \|E'\|_{\infty\to\infty}\le12t^{-1}(B_0+1).
\]
For $t\ge24(B_0+1)$ this proves invertibility. Since
$A'=D^*D+\eta^{-2}n+Q'^*aQ'$ is symmetric and nonnegative,
it is positive definite. The product rule for $D$ gives
\begin{equation}\label{eq:regional-global-sup}
 \max\{|G'g|_{[2]},|DG'g|_{[1]}\}
       \le b'\|g\|_\infty,\qquad b'=2(L+4)B_0.
\end{equation}
Also $Q'Q'^*=I$. The projection calculation in
\cref{sec:multiscale-gram} therefore applies with these
weighted adjoints: $N>0$, $P$ is orthogonal, and
$\operatorname{ran}R=\Delta_U\ker Q'$.
Every member of $\ker Q'$ vanishes in layer zero. Consequently
$R$ is supported in $\Omega_1$ and its bare neighbours, and
\[
 DRD^*+\sum_{j\ge1}Q_j^*Q_j
\]
is supported, in both variables, where $d_1\le2L$.

\emph{A smooth distance.}
Choose a smooth even function $\psi$, with $0\le\psi\le1$, equal to one on
$[-1/2,1/2]$ and zero outside $(-3/4,3/4)$, and put
$c_\psi=\max(1,\|\psi'\|_\infty,\|\psi''\|_\infty)$.
For $z\in\Lambda_*$ set
\[
 \Psi_z(x)=\prod_{\mu=1}^3\psi((x-z)_\mu/\xi_z),\qquad
 \phi_z=\Psi_z/\sum_w\Psi_w,\qquad
 f_{y'}(x)=\sum_z\phi_z(x)d(z,y').
\]
Use least coordinate lifts; the supports make this choice
immaterial. The denominator is at least one on $\Omega_0$.
At any point within distance $\eta$ of a layer-$j$ site,
only levels $j-1,j,j+1$ occur, and at most eight centres
occur at each level. Each contributing centre has weighted
distance at most $9L/2$ from that point: its supporting cube
meets only those three layers, and a coordinate path has
length at most $9\xi_z/4$ and weight at most $2L/\xi_z$.
The quotient rule bounds their first and second derivatives
by $C(L)\xi_j^{-1}$ and $C(L)\xi_j^{-2}$, respectively.
Subtract one fixed $d(z_0,y')$ when
differentiating $\sum_z\phi_z=1$. The resulting differences
are at most $9L$. Integrating the first and second derivative
bounds along bare segments proves
\begin{equation}\label{eq:regional-smoothed-distance}
 \begin{gathered}
 |f_{y'}(x)-d(x,y')|\le s_0,\\
 |f_{y'}(x')-f_{y'}(x)|\le s_1\eta/\xi(x),\\
 |f_{y'}(x+\eta e_\mu)+f_{y'}(x-\eta e_\mu)-2f_{y'}(x)|
          \le s_2\eta^2/\xi(x)^2,
 \end{gathered}
\end{equation}
where $s_0=9L/2$ and $s_1,s_2$ depend only on $L$ and
our fixed bump. The last estimate is used only
when all three sites are in $X$.
Two sites in a cell, adjacent bare sites, and two sites in
a row double block have weighted distance at most $4L$.
Thus the oscillation of $f_{y'}$ on any such set is at most
$s_3=4L+2s_0=13L$; in general it is at most
$d(y(x),y(x'))+s_3$.

For $\varphi=\theta f_{y'}$ with
$0<\theta\le(s_1+s_3)^{-1}$, Taylor's formula gives the
following conjugation bound:
\begin{equation}\label{eq:regional-scalar-conjugation}
 \|(e^\varphi A'e^{-\varphi}-A')u\|_\infty
 \le\epsilon'\theta\max\{|u|_{[2]},|Du|_{[1]}\},
\end{equation}
where
$\epsilon'=(12L^2+2)s_3+12Ls_1+3(s_2+2s_1^2)$.
Indeed, outside $\Omega_1$ we have $\xi(x)=\eta$, and
every existing neighbour has scale at most $L\eta$; the direct
row bound therefore suffices. Inside $\Omega_1$,
write each transported neighbour as $u(x)$ plus its covariant
difference. The latter terms contribute $12Ls_1$; pairing
the two opposite neighbours in each direction bounds the
coefficient of $u(x)$ by $3(s_2+2s_1^2)$.
The average term contributes $2s_3$, using its cell support.
The remaining direct neighbour bound is $12L^2s_3$.

Take $\theta_1=\min\{(s_1+s_3)^{-1},(2\epsilon'b')^{-1}\}$.
By \eqref{eq:regional-global-sup} the conjugation error
composed with $G'$ has supremum norm at most $1/2$.
Its Neumann inverse costs a factor two. On the source cell
$e^\varphi\le3$, and on the output cell
$e^{-\varphi}\le3e^{-\theta_1d(y,y')}$.
The product rule for $D(e^{-\varphi}v)$ adds at most
$2L^2|v|_{[2]}$ to the normalized derivative. Hence the
scalar assertions hold with
\[
 \delta_1=\theta_1,\qquad B_1=18(1+2L^2)b'.
\]

\emph{The Gram gap.}
Put $K=Q'G'Q'^*$. Its inverse has the variational formula
\[
 (\nu,K^{-1}\nu)_{\mathbf B}
   =\min_{Q'\lambda=\nu}(\lambda,A'\lambda)_\eta.
\]
On a positive-level cell use
$\psi_y(x)=\prod_\mu(1-3|(x-y)_\mu|/\xi_y)_+$.
It vanishes on the outermost sheet, its mean is at least
$8^{-3}$, and its first differences are at most $3/\xi_y$.
Let $h_y(x)$ be the bare comb transport from $y$ to $x$.
The curvature-only transport estimate in the proof of
\cref{I-32-37}, applied to the path from the lowest corner
$o$ through $y$ to $x$, gives
\[
 \bigl|\operatorname{Ad}(T_U^{\mathsf c}(y,x)h_y(x)^{-1})
       -\operatorname{Ad}(U(\Gamma_{o,y}))\bigr|
            \le2\kappa_7^{\mathsf c}\alpha\le\tfrac12.
\]
Consequently
\[
 S_y=|B^j(y)|^{-1}\sum_{x\in B^j(y)}\psi_y(x)
         \operatorname{Ad}(T_U^{\mathsf c}(y,x)h_y(x)^{-1})
 \quad\hbox{satisfies}\quad \|S_y^{-1}\|\le2\cdot8^3.
\]
Set $\lambda=\nu$ on layer zero and
$\lambda(x)=\psi_y(x)\operatorname{Ad}(h_y(x))^{-1}S_y^{-1}\nu(y)$
on the other cells. Then $Q'\lambda=\nu$. The comb gauge
bound \cref{I-10-51} gives
$|U^{h_y}-I|\le2\alpha\eta/\xi_y$, hence
$|D\lambda_y|\le4\xi_y^{-1}|S_y^{-1}\nu(y)|$ if
$4\alpha\le1$. Distinct positive-level test functions have
no neighbouring nonzero values, so their energy cross terms
vanish. There are at most $3(\xi_y/\eta)^3$ positive bonds
per cell. Layer zero contributes at most $12\eta^{-2}$,
and the average term at most $\Xi^{-2}$. Therefore
\[
 K^{-1}\le C_K\Xi^{-2},\qquad C_K=3\cdot2^{24}+13.
\]

We use twice the matrix inversion calculation in
\cref{sec:multiscale-gram-proof}: a positive symmetric
matrix with a uniform gap and an exponentially decaying
kernel has an exponentially decaying inverse, by conjugating
with $e^{\theta d(\cdot,y')}$ and bounding the
error by its row and column sums. Remove the coarse weights
by $\Phi\mu(y)=\xi_y^{3/2}\mu(y)$.
For $\widetilde K=\Xi^{-1}K\Xi^{-1}$ we have
$\widetilde K\ge C_K^{-1}$ and
\[
 |\widetilde K^\Phi(y,y')|
       \le B_1e^{-\delta_1d(y,y')}.
\]
Indeed, the scalar estimate first gives the extra factor
$(\xi_y/\xi_{y'})^{5/2}$; symmetry gives its reciprocal,
and the smaller is at most one. The inversion calculation gives constants $B_K,\delta_K>0$
for its inverse kernel. Absorbing the scale ratio and summing
rows and columns gives $\|\Xi^2K^{-1}\|\le n_1$ for a fixed
$n_1$, in the weighted coarse Hilbert norm.
Since $Q'^*Q'$ is an orthogonal projection,
$N\ge K^2$ as quadratic forms. Applying the preceding norm
bound to $K\Xi^{-2}\mu$ yields
\[
 \widetilde N:=\Xi^{-2}N\Xi^{-2}\ge n_1^{-2}.
\]
Two applications of the scalar estimate give
\[
 |\widetilde N^\Phi(y,y')|
 \le B_1^2L^5c_{\mathrm{sum}}(\delta_1/4)
                    e^{-\delta_1d(y,y')/2}.
\]
Here the two scale ratios have powers $3/2$ and $7/2$;
the same scale-ratio and convolution estimates used in
\cref{sec:multiscale-gram-proof} absorb them. Apply the
matrix inversion calculation once more. Undoing $\Phi$
and $\Xi^{-2}$ gives the factor
$\xi_y^{-7/2}\xi_{y'}^{-1/2}$; absorbing the remaining
half-power ratio gives the factor $\xi_y^{-4}$ at half
the rate. Denote common constants for these two bounds by
$B_2,\delta_2>0$.

Put $V=DG'Q'^*$ and $Y=DPD^*=VCV^*$.
The differentiated scalar bound gives $|Ve_z|_y\le
B_1\xi_y e^{-\delta_1d(y,z)}|e|$.
Taking its weighted transpose, and counting at most
$3\xi_{y'}^3\eta^{-3}$ positive bonds per cell, gives
\[
 |(V^*A)(z)|\le3B_1\xi_z^{-3}\sum_{y'}\xi_{y'}^4
       e^{-\delta_1d(z,y')}|A|_{y'}.
\]
In the product $VCV^*$ the scale factor is
$\xi_y\xi_z^{-4}\xi_{z'}^{-3}\xi_{y'}^4$.
It equals $\xi_y^{-2}(\xi_y/\xi_z)^3
(\xi_{z'}/\xi_z)(\xi_{y'}/\xi_{z'})^4$.
Absorbing these ratios and taking the two exponential sums
proves, for fixed $B_Y,\delta_Y>0$,
\begin{equation}\label{eq:regional-projector-decay}
 |YA|_y\le B_Y\xi_y^{-2}\sum_{y'}
       e^{-\delta_Yd(y,y')}|A|_{y'}.
\end{equation}

\emph{The regional vector inverse.}
The one-plaquette Hessian calculation gives
$H_U=D^{\mathsf T}D+H'_U$, where
\[
 \sum_{b'}|(D^{\mathsf T}D)(b,b')|\le16\eta^{-2},
 \qquad \sum_{b'}|H'_U(b,b')|
       \le20\alpha\mathsf L_{\mathbf B}(b)^{-2}.
\]
Only bonds in a common inner plaquette contribute. The
counts are four plaquettes at a bond, times four entries
for the first bound and at most five curvature factors
for the second. Also $Q_0^*Q_0=\eta^{-2}$ when $d_1\ge2$.
Compress $A_0$ to the bonds whose endpoints satisfy
$d_1>t-3$. The nonlocal part vanishes there, so
$A_{\mathrm l}=H_U+\eta^{-2}$ on zero extensions.
For $40\alpha\le1$, its spectrum lies in
$[\tfrac12,34]\eta^{-2}$. The operator
$J=I-(4/69)\eta^2A_{\mathrm l}$ has norm at most $67/69$
and range one in bare sup distance. Its geometric series
therefore gives
\[
 \|G_{\mathrm l}g\|_\infty\le c_{\mathrm l}\eta^2\|g\|_\infty,
 \qquad c_{\mathrm l}=6\sum_{z\in\mathbb Z^3}(67/69)^{|z|_\infty}.
\]

First establish positivity separately. The nonlocal part
commutes with each $\zeta_i$, since it is supported where
$d_1\le2L<t$. For
$C_i=\zeta_i^2A_0-2\zeta_iA_0\zeta_i+A_0\zeta_i^2$,
\[
 A_0=\sum_i\zeta_iA_0\zeta_i+\tfrac12\sum_iC_i.
\]
The kernel of $C_i$ is
$(\zeta_i(b_-)-\zeta_i(b'_-))^2H_U(b,b')$.
It is supported in the layer-zero transition strip and
has row sum at most $132t^{-2}\eta^{-2}$.
The collar bound in \eqref{eq:regional-collar-input},
symmetry and positivity imply
$A_{\mathrm a}\ge(L^2B_0)^{-1}\Xi^{-2}$: the symmetric
positive matrix $\Xi^{-1}G_{\mathrm a}\Xi^{-1}$ is similar
to $\Xi^{-2}G_{\mathrm a}$, whose supremum norm is at most
$L^2B_0$. Combining this with the layer-zero lower bound
and $\zeta_1^2+\zeta_2^2=1$, and taking $t\ge17LB_0$,
gives
\[
 (A,A_0A)_\eta\ge(2L^2B_0)^{-1}
                         (\Xi^{-1}A,\Xi^{-1}A)_\eta.
\]

Now form $P_0=\zeta_1G_{\mathrm a}\zeta_1+
\zeta_2G_{\mathrm l}\zeta_2$. Its error is supported in
the same transition strip. The single commutators have
row bound $66t^{-1}\eta^{-2}$ there, so
\[
 A_0P_0=I+E_0,\qquad
 \|E_0\|_{\infty\to\infty}
        \le66t^{-1}(L^2B_0+c_{\mathrm l}).
\]
For $t\ge132(L^2B_0+c_{\mathrm l})$, the same Neumann
and product-rule argument as for $G'$ yields
\[
 \max\{|G_0g|_{[2]},|\nabla^UG_0g|_{[1]}\}
       \le b''\|g\|_\infty,\qquad
 b''=4(L^2B_0+c_{\mathrm l}).
\]

Use the same smooth weight $\varphi=\theta f_{y'}$.
Since $R=I-P$,
\[
 A_0=(D^{\mathsf T}D+DD^*)+H'_U+Q^*Q-Y.
\]
For $x\in\Omega_1$ all neighbours needed in the curl--divergence
identity from the proof of \cref{I-65-4} are inner.
Its diagonal part is the scalar covariant Laplacian on each
component; its two mixed transports differ by at most
$2\alpha\eta^2\xi(x)^{-2}$. Thus the same opposite-neighbour
cancellation used in \eqref{eq:regional-scalar-conjugation}
applies to the vector components. In layer zero direct row
bounds suffice. For $Q^*Q$, retain the adjoint factor:
\[
 (Q^*Qv)(b)=\eta^{-3}\sum_j\xi_j^{-1}
       \sum_{c\in\Gamma_j^\flat}Q(c,b)^{\mathsf T}
                       \sum_{b'}Q(c,b')v(b').
\]
The corner product bounds give row sum $2\xi_j$
and level-$j$ column sum $2\eta^3\xi_j^{-2}$. On its
double block the weight oscillation is at most $2s_3\theta$
and $|u|\le L^2\xi_j^2|u|_{[2]}$. Hence each of the at
most three levels contributes at most
$8L^2s_3\theta|u|_{[2]}$.
Finally \eqref{eq:regional-projector-decay},
$|e^r-1|\le|r|e^{|r|}$, and one exponential sum give
\[
 \|(e^\varphi A_0e^{-\varphi}-A_0)u\|_\infty
 \le\epsilon''\theta
        \max\{|u|_{[2]},|\nabla^Uu|_{[1]}\},
\]
provided $\theta\le\delta_Y/4$, where one may take
\[
 \epsilon''=12Ls_1+3(s_2+2s_1^2)+115L^2s_3
       +3B_YL^2(s_3+2/\delta_Y)c_{\mathrm{sum}}(\delta_Y/4).
\]
Choose additionally
$\theta\le(2\epsilon''b'')^{-1}$. The same conjugated
Neumann inverse and the same source and output weight
bounds as in the scalar case give the vector estimates,
with $B_3=18(1+2L^2)b''$ and rate $\delta_3=\theta$.

All constants used so far depend only on $L$ and the fixed
collar constants. Choose the power $g_0$ so that
$g_0\ge L^2M$, $t=(g_0-5)/2$ exceeds the four bounds
used above, and $Lg_0$ exceeds the finitely many summation
and scale-ratio thresholds. Put $g_{\mathrm p}=Lg_0$.
Choose $a_{\mathrm{ps}}$ below the restrictions in \eqref{eq:collar-gauge-radius}, $(4\kappa_7^{\mathsf c})^{-1}$,
and $1/40$. Finally enlarge $B_{\mathrm{ps}}$ and decrease
$\delta_{\mathrm{ps}}$ to cover the three pairs of bounds.
The empty-region case is vacuous; the case $\Omega_1=\emptyset$
was included by omitting the collar terms.
\end{proof}


\subsection{The constrained inverse}

The remaining constraint is imposed on the reduced corner rows.
We construct a right inverse with fields supported in the cores
of the averaging cells. Its energy bound gives a uniform gap for
the row Gram matrix, and hence the constrained inverse and the
weighted transpose bounds used in the nonlinear equation.

Keep the regional operators of \cref{I-42-86}. Put
\[
 M_Q=QG_0Q^*,\qquad H_0=G_0Q^*M_Q^{-1},\qquad
 \widetilde G_0=G_0-H_0QG_0.
\]
The existence of $M_Q^{-1}$ is part of the following proposition.
For a bond or plaquette field, let
$|f|_{(-s)}=\sup_{\mathfrak e}\mathsf L_{\mathbf B}(\mathfrak e)^s
|f(\mathfrak e)|$; for sites use $\xi(x)$. A difference uses
the largest local scale at the endpoints of its two bonds. Set
\[
 \|A\|_X=\max\{|A|_{(-1)},|\nabla^UA|_{(-2)}\},\qquad
 \|\omega\|_1^2=\sum_c\xi_c^{-1}|\omega(c)|^2.
\]
Below $D$ on bond fields is the covariant curl into inner
plaquettes, whereas $D^*$ on bond fields is the divergence
into sites. The region and zero-extension conventions remain
those fixed above.

\begin{proposition}\label{I-44-50}\label{I-44-60}\label{I-45-14}
There are constants $0<a_Q\le a_{\mathrm{ps}}$, $\gamma_Q>0$,
$w_S,c_\delta\ge1$, and $K_1\ge4$, depending only on $L$,
with the following property. Assume the hypotheses of
\cref{I-42-86} with $\alpha\le a_Q$ and
$R_1\mathfrak M\ge w_S$. Then $M_Q\ge\gamma_Q I$ in the
row pairing. Put $K_7=K_1^2c_\delta$. We have
$QH_0=I$, $Q\widetilde G_0=0$, and
\[
 \|H_0\beta\|_X\le K_1\sup|\beta|,\qquad
 \|\widetilde G_0f\|_X\le K_7|f|_{(-3)}.
\]
For real bond fields $f$, inner-plaquette fields $h$, site
fields $g$, and row fields $\omega$,
\[
 \begin{gathered}
 \sup|H_0^*f|\le K_1|f|_{(-3)},\qquad
 \sup|(DH_0)^*h|\le3K_1|h|_{(-2)},\qquad
 \sup|(D^*H_0)^*g|\le3K_1|g|_{(-2)},\\
 |Rg|_{(-2)}\le(1+K_1^3)|g|_{(-2)},\qquad
 |Q^*\omega|_{(-3)}\le2(L^3+2)\sup|\omega|.
 \end{gathered}
\]
Adjoints use the fine weight $\eta^3$ and row weight $\xi_c^{-1}$.
All bounds are uniform in the volume, lattice spacing, regions
and number of layers.
\end{proposition}

\begin{proof}

Temporarily retain all positive determining rows, including
the inward ones, and write $Q^+$ for this corner average.
The weight of each row is still $\xi_c^{-1}$. For a level-$j$
row $c$, choose $y_c=c_-$ when $c_-\in\Lambda_j$ and
$y_c=c_+$ otherwise. There are at most six rows with a
given top. Put $\vartheta_c=L$ for an outward row and
$\vartheta_c=1$ otherwise. Every plaquette in its double
block $\mathsf D_c$ has curvature at most
$\vartheta_c^2\alpha L^{-2j}$, and $\mathsf D_c\subset\Omega_0$.
Thus the folding, locality, product and path calculations in
the proof of \cref{I-32-37} apply whenever
$L^2\alpha\le c_{\mathsf c}$.

\emph{A lift supported in the cell cores.}
For $j\ge1$ put $N=L^j$, $n=(N-1)/2$ and $w=\lfloor N/8\rfloor$.
Use the tents
\[
 h(P)=(w-|P-n|)_+,
 \qquad \varphi_\pm(P)=(w-|P-(n\pm w)|)_+,
\]
and set
\[
 g^+=(n+w)\varphi_+-(n-w)\varphi_-,\qquad
 g^-=(n+w+1)\varphi_--(n-w+1)\varphi_+.
\]
Their moments are
\[
 \sum h=w^2,\qquad
 \sum g^+(P)(P+1)=\sum g^-(P)(N-1-P)=2Nw^3,
\]
\[
 \sum g^+(P)(N-1-P)=\sum g^-(P)(P+1)=0.
\]
Indeed, a tent centred at $p$ has mass $w^2$ and first
moment $pw^2$. Also $N/18\le w\le N/8$. All three
profiles vanish outside $|P-n|\le2w-1$, and their support,
together with one neighbouring sheet, lies in
$1\le P\le N-2$ and $N/L\le P\le N-1-N/L$.

Let $z_y$ be the lowest corner of the cell of $y=y_c$, and
write $P_\nu(x)=\eta^{-1}(x-z_y)_\nu$. Let
$p_y(x)=U(\Gamma_{z_y,x})^{-1}$ for the bare comb in the
cell. If $c$ has direction $\mu$, put $\epsilon_c=+$ and
$g_c=I$ when $y=c_-$; when $y=c_+$ put $\epsilon_c=-$ and
$g_c=U([z_y-\xi_je_\mu,z_y])^{-1}$. The latter path lies
in $\mathsf D_c$. On positive bonds of direction $\mu$ set
\[
 \psi_c(b)=\frac{N^2}{2\eta w^7}
  g^{\epsilon_c}(P_\mu(b_-))\prod_{\nu\ne\mu}h(P_\nu(b_-)),
 \qquad
 (Z_ce)(b)=\psi_c(b)\operatorname{Ad}(p_y(b_-)g_c)e,
\]
and set these fields to zero elsewhere. At level zero use
$Z_ce=\eta^{-1}\mathbf1_{\{c\}}e$. Put
$Z\omega=\sum_cZ_c\omega(c)$ and $\Theta=Q^+Z$.

The profile bounds give, with $a_1=18^4/2$,
\[
 \xi_j|\psi_c|\le a_1,\qquad
 \xi_j^2\eta^{-1}|\psi_c(b+\eta e_\nu)-\psi_c(b)|
       \le18^5/2,\qquad
 \eta N^{-2}\sum_b|\psi_c(b)|\le18.
\]
For $f(x)=\operatorname{Ad}(p_y(x)g_c)e$, the comb comparison
in \cref{I-10-51} gives
\[
 |U_\nu(x)p_y(x+\eta e_\nu)p_y(x)^{-1}-I|
       \le2\alpha/N,
 \qquad |\nabla^U_\nu f(x)|\le4\alpha\xi_j^{-1}|e|.
\]
The loop uses at most two strips, each of width at most
$N-1$, in a cell whose plaquettes have curvature at most
$\alpha N^{-2}$. The product rule therefore yields
\begin{equation}\label{eq:corner-test-size}
 |Z_ce|\le a_1\xi_j^{-1}|e|,
 \qquad |\nabla^UZ_ce|\le18^5\xi_j^{-2}|e|,
\end{equation}
for $\alpha\le1$. Every nonzero term, and every site or
plaquette entering its first differences, lies in the core
$\mathsf K(y)=\{n-2w\le P_\nu\le n+2w\}$.

\emph{Correction of the row values.}
Write $\Theta^=$ for the same-level part and
$\Theta^\uparrow$ for the part from level $j$ to level $j+1$.
Locality gives
\[
 \Theta=\Theta^=+\Theta^\uparrow,
 \qquad (\Theta^\uparrow)^2=0.
\]
Here is the support argument. A row can meet a positive-level
core at its own level only through one of its endpoints.
A row at a higher level must be outward: its outer block
lies in the immediately lower layer. A row at a lower level
can enter the cell only in its outermost order-$(j-1)$ sheet,
where the core-supported profile vanishes. The level-zero
case uses $Z_ce=\eta^{-1}\mathbf1_{\{c\}}e$ and gives the
same alternatives. Two consecutive upward couplings would
place a point of $\Omega_{j+1}^c$ and a point of
$\Omega_{j+2}$ at distance at most
$(1+(L-1)/2+L)\xi_{j+1}<L^2\xi_{j+1}$, contradicting
the separation of the regions.

The absolute entry bound from \cref{I-32-37} gives
$|\Theta^\uparrow(c,c')|\le36L^{-2}$. There are at most
$6L^3$ entries in each row and six in each column. The
column sum in the weighted Schur test carries the additional
factor $\xi_{c'}/\xi_c=L^{-1}$. Hence
\[
 \|\Theta^\uparrow\|\le216/L<17.
\]
For the same-level part, a nested straight chain ending at a
bond of direction $\mu$ in the cell of $y$ occurs $P_\mu+1$
times from $\langle y,y+\xi_je_\mu\rangle$ and
$N-1-P_\mu$ times from $\langle y-\xi_je_\mu,y\rangle$.
These counts follow by choosing its longitudinal starting
site; the transverse coordinates are fixed. The displayed
moments consequently give
\[
 \eta N^{-3}\sum_b\psi_{c'}(b)n_c(b)=\delta_{cc'},
 \qquad \eta N^{-3}\sum_b|\psi_{c'}(b)|n_c(b)\le18.
\]

The curvature-only product estimate in the proof of
\cref{I-32-37} writes the row as its leading chain sum plus
a remainder of row norm at most
$\kappa_S\vartheta_c^2\alpha\xi_j$, where one may take
$\kappa_S=8\kappa_9^{\mathsf c}L^{-3}$.
Compare each chain transport with the path from the lowest
corner of $c_-$ to $z_{y_{c'}}$, followed by the comb to
$b_-$. It stays in $\mathsf D_c$ and has at most $4N$ bonds.
The same proof bounds its error by
$\kappa_7^{\mathsf c}\vartheta_c^2\alpha$.
After multiplication by $p_{y_{c'}}(b_-)g_{c'}$, the comparison
transport is constant in $b$. On the diagonal it is the
identity; off the diagonal its coefficient is the vanishing
moment. Therefore
\[
 |(\Theta^=-I)(c,c')|
 \le\kappa_\Theta L^2\alpha,
 \qquad
 \kappa_\Theta=36\kappa_7^{\mathsf c}+\kappa_Sa_1.
\]
There are at most twelve same-level entries in a row and
six in a column. Thus $\|\Theta^=-I\|\le1/36$ if
$324\kappa_\Theta L^2\alpha\le1$. At level zero the
same-level part is exactly the identity. It follows that
\[
 \Theta=(I+\Theta^\uparrow)
  [I+(I-\Theta^\uparrow)(\Theta^=-I)],
 \qquad \|\Theta^{-1}\|\le36,\quad\|\Theta\|\le19.
\]
Indeed, the second factor differs from the identity by at
most $18/36$, while $(I+\Theta^\uparrow)^{-1}=I-\Theta^\uparrow$.

\emph{Energy of the lift.}
Use the nonnegative form
\[
 \begin{split}
 \mathcal N(A)={}&\eta^3\sum_{p\ {\mathrm{inner}}}|DA(p)|^2
 +\eta\sum_{p\ {\mathrm{crossing}}}
                  \Bigl(\sum_{b\subset\partial p}|A(b)|\Bigr)^2\\
 &+\|D^*A\|_\eta^2+\|\mathsf L_{\mathbf B}^{-1}A\|_\eta^2
                              +\|Q^+A\|_1^2.
 \end{split}
\]
The crossing sum is over non-inner plaquettes containing an
inner bond; it uses the zero extension of $A$, and no value
of $U$ outside the region. A plaquette or site meeting a
positive-level profile meets profiles from just one core.
All its bonds lie in that cell. Otherwise only level-zero
profiles contribute. This follows from the one-sheet core
margin and disjointness of the cells.

For a core at $y$, put
$s_\mu=\sum_{c:y_c=y,\,\mathrm{dir}(c)=\mu}|\omega(c)|$.
There are at most two summands, so
$\sum_\mu s_\mu^2\le2\sum_{c:y_c=y}|\omega(c)|^2$.
The value and difference bounds \eqref{eq:corner-test-size},
with at most $N^3$ base sites, bound the curl, divergence
and weighted-value energies by
\[
 (8\cdot18^{10}+6\cdot18^{10}+2a_1^2)
           \sum_{c:y_c=y}\xi_c^{-1}|\omega(c)|^2.
\]
For level zero the corresponding constants are $16,12,1$:
each bond belongs to four plaquettes and touches two sites,
with at most six bonds at a site. The row energy is at most
$19^2\|\omega\|_1^2$. Consequently, with $c_Z=15\cdot18^{10}$,
\[
 \mathcal N(Z\omega)\le c_Z\|\omega\|_1^2,
 \qquad Z^\sharp=Z\Theta^{-1},\qquad
 Q^+Z^\sharp=I,\qquad
 \mathcal N(Z^\sharp\omega)\le36^2c_Z\|\omega\|_1^2.
\]

The regional Hessian calculation gives
$H_U=D^{\mathsf T}D+\Delta'$ and, with $C_H=20$, the bound
\[
 |\Delta'A(b)|\le C_H\alpha\mathsf L_{\mathbf B}(b)^{-2}
                     \max_{b'\in\operatorname{st}(b)}|A(b')|.
\]
Only inner plaquettes occur. The star consists of $b$ and
the other bonds of these plaquettes, so it is symmetric and
has at most thirteen members; their local scales differ by
at most $L$. At an incomplete star the local scale is
$\eta$, so the same bound follows directly from the
one-plaquette formula. Applying $2ab\le a^2+b^2$ to
$a(b)=\mathsf L_{\mathbf B}(b)^{-1}|A(b)|$ gives, for $\alpha L\le1$,
\[
 |(A,\Delta'A)_\eta|\le13C_H\|\mathsf L_{\mathbf B}^{-1}A\|_\eta^2,
 \qquad (A,G_0^{-1}A)_\eta\le(1+13C_H)\mathcal N(A).
\]
The second inequality also uses that $R$ is orthogonal and
the reduced row energy is a partial sum of the full row energy.

Take $A=Z^\sharp\omega$. Cauchy--Schwarz for the positive
operator $G_0$ yields
\[
 \|\omega\|_1^2=(A,Q^{+*}\omega)_\eta
 \le\gamma_Q^{-1/2}\|\omega\|_1
       (\omega,Q^+G_0Q^{+*}\omega)_1^{1/2},
 \qquad
 \gamma_Q=[36^2(1+13C_H)c_Z]^{-1}.
\]
This proves the full-row lower bound. If $J$ restricts to
the interior and outward rows, $J^*$ is zero extension and
$\|J^*\omega\|_1=\|\omega\|_1$. Hence
$M_Q=JQ^+G_0Q^{+*}J^*\ge\gamma_Q I$.
All smallness requirements hold after choosing
\[
 a_Q\le\min\{a_{\mathrm{ps}},L^{-2}\min(1,c_{\mathsf c}),
                         (324\kappa_\Theta L^2)^{-1}\}.
\]


Write $d=d_{\mathbf B}$, $B=B_{\mathrm{ps}}$ and
$\delta=\delta_{\mathrm{ps}}$ for the regional constants.
The cells meeting the double block of a reduced row $c$
form a set $\mathsf N(c)$ of at most $L^3+1$ tops, and
\begin{equation}\label{eq:corner-row-neighbourhood}
 L^{-1}\xi_c\le\xi_y\le L\xi_c,
 \qquad d(y_c,y)\le L+1\quad(y\in\mathsf N(c)).
\end{equation}
For an interior row there are just its two endpoint cells.
For an outward row there is the inner endpoint cell and
the $L^3$ cells partitioning its outer block. A coordinate
path first moving transversely in the inner cell and then
along the row has weighted length at most $L+1$.

\emph{The row Schur inverse.}
For data $e_c$ at one row, the adjoint formula is
\[
 (Q^*e_c)(b)=\eta^{-3}\xi_c^{-1}Q(c,b)^{\mathsf T}e.
\]
The corner entry bound in \cref{I-32-37} gives $|Q^*e_c|\le2\xi_c^{-3}|e|$,
supported over $\mathsf N(c)$. Splitting it into its cell
parts and applying \cref{I-42-86} therefore gives
\begin{equation}\label{eq:corner-GQ-column}
 \begin{split}
 |G_0Q^*e_c|_y&\le C_1\xi_y^2\xi_c^{-3}
                         e^{-\delta d(y,y_c)}|e|,\\
 |\nabla^UG_0Q^*e_c|_y&\le C_1\xi_y\xi_c^{-3}
                         e^{-\delta d(y,y_c)}|e|,
 \end{split}
\end{equation}
where $C_1=4(L^3+1)Be^{(L+1)\delta}$ is sufficient.
Using $|QA(c)|\le2\xi_c\max_{y\in\mathsf N(c)}|A|_y$
once more yields, with $C_M=4L^2e^{(L+1)\delta}C_1$,
\[
 |M_Q(c,c')|\le C_M(\xi_c/\xi_{c'})^3
                              e^{-\delta d(y_c,y_{c'})}.
\]
Weighted symmetry says
$|M_Q(c,c')|=(\xi_c/\xi_{c'})|M_Q(c',c)|$.
Thus the scale factor can be replaced by
$\min\{(\xi_c/\xi_{c'})^3,(\xi_{c'}/\xi_c)^2\}$,
which is at most $(\xi_c/\xi_{c'})^{1/2}$.

We use the exponential conjugation argument already proved
for positive matrices in \cref{sec:multiscale-gram-proof}.
For the row pairing, its symmetric normalization is
$\omega(c)\mapsto\xi_c^{-1/2}\omega(c)$.
Let $c_*=c_{\mathrm{sum}}(\delta/2)$ and choose
\[
 0<\delta_*\le\min\{\delta/4,
                       \gamma_Q\delta/(48C_Mc_*)\}.
\]
Provided the gap exceeds the summation threshold at
$\delta/2$, the perturbation under conjugation by
$e^{\delta_*d(y_c,y_{c'})}$ has norm at most $\gamma_Q/2$:
there are at most six rows per top and
\[
 (e^{\delta_*s}-1)e^{-\delta s}
            \le(4\delta_*/\delta)e^{-\delta s/2}.
\]
The lower bound just proved and the Neumann inverse give
\[
 |M_Q^{-1}(c,c')|\le2\gamma_Q^{-1}
       (\xi_c/\xi_{c'})^{1/2}e^{-\delta_*d(y_c,y_{c'})}.
\]

Insert this in $H_0=G_0Q^*M_Q^{-1}$ and use
\eqref{eq:corner-GQ-column}. The intermediate scale factor is
\[
 \xi_y^2\xi_c^{-3}(\xi_c/\xi_{y'})^{1/2}
   =\xi_y^{-1}(\xi_y/\xi_c)^{5/2}
                          (\xi_y/\xi_{y'})^{1/2}.
\]
The triangle inequality leaves at least $\delta/2$ on
$d(y,y_c)$. The scale-ratio sum at exponent $5/2$ is
therefore bounded by $c_*$. The derivative estimate has
$\xi_y^{-2}$ in place of $\xi_y^{-1}$. Taking
\[
 B_H^\circ=\max\{B,72\gamma_Q^{-1}C_1c_*\},\qquad
 B_H=L^{1/2}B_H^\circ,
\]
we obtain the two bounds with $B_H^\circ$, rate $\delta_*$,
and the remaining factor $(\xi_y/\xi_{y'})^{1/2}$.
Choose
\[
 \delta_S=\min\{\delta_*/2,(L+2)^{-1}\}.
\]
If $(\delta_*-\delta_S)R_1\mathfrak M\ge\tfrac12\log L$,
the scale-ratio estimate absorbs that factor, giving
\begin{equation}\label{eq:corner-H-local}
 \begin{split}
 |H_0\beta|_y&\le B_H\xi_y^{-1}\sum_{y'}
             e^{-\delta_Sd(y,y')}|\beta|_{y'},\\
 |\nabla^UH_0\beta|_y&\le B_H\xi_y^{-2}\sum_{y'}
             e^{-\delta_Sd(y,y')}|\beta|_{y'}.
 \end{split}
\end{equation}
Here $|\beta|_y$ is the maximum over rows with top $y$.
The identity $QH_0=I$ follows from the definition of $M_Q$.

\emph{Weighted norms and transposes.}
Increase $w_S$ to include the preceding gap conditions,
the weighted threshold $w_{\mathrm{sum}}(\delta_S/4)$, and
$4\delta_S^{-1}\log L$. Put
\[
 c_\delta=\max\{1,c_{\mathrm{sum}}(\delta_S/4)\},
 \qquad \overline B=B_Hc_\delta,
 \qquad K_1=4L^7\overline B.
\]
For $|s|\le4$ the sum
$\sum_{y'}e^{-\delta_Sd(y,y')/4}(\xi_y/\xi_{y'})^s$
is at most $c_\delta$. For an object over $y$, its local
scale lies between $\xi_y$ and $L\xi_y$. For $s\ge0$, therefore,
\[
 |f|_y\le\xi_y^{-s}|f|_{(-s)},\qquad
 |f|_{(-s)}\le L^s\sup_y\xi_y^s|f|_y.
\]
The second bound is an equality without $L^s$ for site
fields. These comparisons also apply to differences.
If one of the translated bonds is not inner, its nonzero
neighbour has scale $\eta$, and the normalized difference
is bounded by $|A|_{(-1)}$. Hence the zero-extension
differences give the same $X$ norm as its inner pairs.
The regional bounds and \eqref{eq:corner-H-local} now give
\[
 \|G_0f\|_X\le L^2\overline B|f|_{(-3)}
                 \le\tfrac14K_1|f|_{(-3)},\qquad
 \|H_0\beta\|_X\le L^2\overline B\sup|\beta|
                 \le K_1\sup|\beta|.
\]

For the transposes, fix $c$ and put $A=H_0e_c$.
The row pairing gives the factor $\xi_c$:
\[
 \langle e,H_0^*f(c)\rangle=\xi_c(A,f)_\eta.
\]
There are at most $3(\xi_y/\eta)^3$ positive bonds and
the same number of inner plaquettes over $y$. Using
\eqref{eq:corner-H-local} and $|DA|_y\le2|\nabla^UA|_y$ gives
\begin{equation}\label{eq:corner-H-transposes}
 \begin{split}
 |H_0^*f(c)|&\le3B_H\sum_y\frac{\xi_c}{\xi_y}
        e^{-\delta_Sd(y,y_c)}\xi_y^3|f|_y,\\
 |(DH_0)^*h(c)|&\le6B_H\sum_y\frac{\xi_c}{\xi_y}
        e^{-\delta_Sd(y,y_c)}\xi_y^2|h|_y.
 \end{split}
\end{equation}
To estimate $D^*A(x)$, consider the two bonds in each
coordinate direction at $x$. If both are inner, their
contribution is a transported first difference at
$x-\eta e_\mu$. Its cell scale differs by at most $L$
and its top distance by less than three; since
$e^{4\delta_S}<3/2$, its modulus is at most
$\tfrac32L^2B_H\xi_y^{-2}e^{-\delta_Sd(y,y_c)}|e|$.
If exactly one is inner, both endpoints of that bond lie
in layer zero. The contribution is $\eta^{-1}|A(b)|$
and obeys the same bound by the value estimate. If neither
is inner it vanishes. Summing the three directions and
then the site pairing gives
\[
 |(D^*H_0)^*g(c)|\le\tfrac92L^2B_H
   \sum_y\frac{\xi_c}{\xi_y}
        e^{-\delta_Sd(y,y_c)}\xi_y^2|g|_y.
\]
The weighted sum at exponent one bounds the right sides
by $3\overline B|f|_{(-3)}$,
$6\overline B|h|_{(-2)}$, and
$\tfrac92L^2\overline B|g|_{(-2)}$, respectively.
These are the asserted transpose bounds.

For the scalar projection, write
$Pg=G'Q'^*CQ'G'g$. The scalar regional estimates and
the same weighted sums give successively
\[
 \sup|Q'G'g|\le\overline B|g|_{(-2)},\qquad
 |Q'^*C\mu|_{(-4)}\le\overline B\sup|\mu|,
 \qquad |G'f|_{(-2)}\le\overline B|f|_{(-4)}.
\]
Here $|Q'\lambda(y)|\le|\lambda|_y$ and
$|Q'^*\mu|_y=|\mu(y)|$. Thus
$|Pg|_{(-2)}\le\overline B^3|g|_{(-2)}$ and the bound
for $R=I-P$ follows.

At each level the corner product column bound is
$\sum_c|Q(c,b)|\le2\eta^3\xi_j^{-2}$.
If $\mathsf L_{\mathbf B}(b)=\xi_n$, only levels $n-1,n,n+1$ can
contribute. Multiplying by the adjoint weight and summing
therefore gives
\[
 \eta^{-3}\mathsf L_{\mathbf B}(b)^3
       \sum_c\xi_c^{-1}|Q(c,b)|\le2(L^3+2),
\]
which proves the bound on $Q^*$. The row support similarly
gives $\sup|QA|\le2L|A|_{(-1)}$.

Finally $Q\widetilde G_0=0$ follows from $QH_0=I$.
We also need a local form of its decay. Put $\beta=QG_0f$.
A row with top $z$ has support at distance at most $L+1$
from $z$, cell scales at most $L\xi_z$, and row sum at most
$2\xi_z$. The regional inverse bound therefore gives
\[
 |\beta|_z\le12L^2B_H\xi_z^3
       \sum_{y'}e^{-\delta_Sd(z,y')}|f|_{y'}.
\]
Insert this in \eqref{eq:corner-H-local}. The triangle
inequality gives
\[
 d(y,z)+d(z,y')\ge\tfrac12d(y,y')+\tfrac12d(y,z),
 \qquad
 \xi_y^{-1}\xi_z^3=\xi_y^2(\xi_z/\xi_y)^3.
\]
The remaining sum over $z$ is at most $c_\delta$.
For the derivative the output factor is $\xi_y$ instead
of $\xi_y^2$. Since $\widetilde G_0=G_0-H_0QG_0$ and
$B_H+12L^2B_H\overline B\le K_1^2$, this proves
\begin{equation}\label{eq:corner-Gtilde-local}
 \begin{split}
 |\widetilde G_0f|_y&\le K_1^2\xi_y^2
        \sum_{y'}e^{-\delta_Sd(y,y')/2}|f|_{y'},\\
 |\nabla^U\widetilde G_0f|_y&\le K_1^2\xi_y
        \sum_{y'}e^{-\delta_Sd(y,y')/2}|f|_{y'}.
 \end{split}
\end{equation}
For the global norm, the value part of the regional bound gives
\[
 \sup|\beta|\le2L^2\overline B|f|_{(-3)}
                      \le\tfrac12K_1|f|_{(-3)}.
\]
Hence
\[
 \|\widetilde G_0f\|_X
 \le(\tfrac14K_1+\tfrac12K_1^2)|f|_{(-3)}
 \le K_1^2|f|_{(-3)}\le K_7|f|_{(-3)}.
\]
\end{proof}


\section{The nonlinear regional equation}\label{sec:regional-nonlinear}

We now construct a critical configuration from nearby reduced
averaging data. The regional inverse supplies the linear part.
The nonlinear constraint uses centered averages, which must first
be compared with the corner averages in that inverse. A gauge
correction will then impose the Landau condition while prescribing
the induced gauge at every block top.

Keep the separated regional setting of \cref{I-45-14}, write
$U_0$ for the real reference configuration, and abbreviate
$\mathsf L=\mathsf L_{\mathbf B}$. The curvature parameter is
denoted by $\alpha_0$. Constants and radii below depend only on
$L$; the radii may be decreased to satisfy the preceding estimates.
For a general sequence, define $\mathfrak U'_k(\mathbf B,r)$
by the same inner-plaquette bound as
\eqref{eq:regional-regularity}, together with
\[
 |(D_U^{\mathsf T}\partial U)(b)|
       <r\eta^2\mathsf L(b)^{-3}
       \quad\text{at every full bond }b.
\]
This agrees with the earlier definition for the constant sequence.

\subsection{The action on the original region}

Let $\mathcal P(b)$ be the inner plaquettes containing $b$,
and let $\operatorname{st}(b)$ consist of their positive bonds.
At a free scale $\ell\ge\eta$, put
\[
 s_b^\ell(A)=\ell\max_{\operatorname{st}(b)}|A|,
 \qquad
 x_b^\ell(A)=\max\{s_b^\ell(A),
       \ell^2\max_{\operatorname{st}(b)}|\nabla^{U_0}A|\}.
\]
The second maximum uses only differences determined by that
star. Define $s_p^\ell,g_p^\ell$ in the same way on a plaquette,
with $g_p^\ell$ the scaled first-difference maximum. Empty
maxima are zero. A missing superscript means
$\ell=\mathsf L(b)$ or $\mathsf L(p)$, respectively.
The adjacent-layer geometry gives
$s_b\le L|A|_{(-1)}$, $x_b\le L^2\|A\|_X$ and
$\max\{s_p,g_p\}\le L\|A\|_X$.

\begin{proposition}\label{I-40-63}
Suppose $U_0$ is real and
$|\partial U_0(p)-I|\le\alpha_0\eta^2\mathsf L(p)^{-2}$
at every inner plaquette, where $0\le\alpha_0\le1/2$.
In the fine pairing of weight $\eta^3$, set
\[
 F(A)=\eta^{-1}\sum_{p\in\mathfrak p}a_p(e^{i\eta A}U_0),
 \quad J_0=F'(0),\quad
 \Delta'=H_{U_0}-D^*D,\quad
 \mathcal V(A)=F'(A)-J_0-H_{U_0}A.
\]
Here $D$ in $D^*D$ is the curl into the original inner
plaquettes. There is an absolute $C_{\mathrm{ec}}\ge1$ such
that, with $\rho_{\mathrm{ec}}=1/256$,
\begin{equation}\label{eq:regional-action-expansion}
 F'(A)=J_0+D^*DA+\Delta'A+\mathcal V(A).
\end{equation}
The remainder is entire, depends at $b$ only on its star,
and satisfies $\mathcal V(0)=\mathcal V'(0)=0$.
The operators preserve real fields, and
\begin{align*}
 |\Delta'A|_{(-3)}&\le C_{\mathrm{ec}}L\alpha_0|A|_{(-1)},\\
 |\mathcal V(A)|_{(-3)}
    &\le C_{\mathrm{ec}}L^3|A|_{(-1)}\|A\|_X
       &&\bigl(L|A|_{(-1)}\le\rho_{\mathrm{ec}}\bigr),\\
 |\mathcal V(A_1)-\mathcal V(A_2)|_{(-3)}
    &\le4C_{\mathrm{ec}}L^4r\|A_1-A_2\|_X
       &&\bigl(\|A_i\|_X\le r\le\rho_{\mathrm{ec}}/(4L^2)\bigr).
\end{align*}
Locally, if the incident curvature is at most
$\alpha_0\eta^2\ell^{-2}$, and $\ell=\eta$ for a nonfull
bond, these bounds follow from
\begin{equation}\label{eq:regional-action-local}
 \begin{split}
 \ell^3|\Delta'A(b)|&\le C_{\mathrm{ec}}\alpha_0s_b^\ell(A),\\
 \ell^3|\mathcal V(A)(b)|
      &\le C_{\mathrm{ec}}s_b^\ell(A)x_b^\ell(A)
                       \quad(s_b^\ell(A)\le\rho_{\mathrm{ec}}),\\
 \ell^3|\mathcal V(A_1)(b)-\mathcal V(A_2)(b)|
      &\le4C_{\mathrm{ec}}r x_b^\ell(A_1-A_2)
                       \quad(x_b^\ell(A_i)\le r\le\rho_{\mathrm{ec}}/4).
 \end{split}
\end{equation}
For $E=e^{i\eta A}U_0$,
\begin{equation}\label{eq:regional-current-conversion}
 F'(A)(b)=\mathsf f(-i\eta\operatorname{ad}_{A(b)})
                       J_{\Omega_0}(E)(b).
\end{equation}
The current uses only the inner plaquettes. If
$\eta|A(b)|\le1/8$, its norm and the gradient norm differ
by factors at most $1+4\eta|A(b)|$ and $1+2\eta|A(b)|$,
respectively. At a full bond $J_{\Omega_0}(U_0)=J(U_0)$;
at a nonfull bond $\eta^3|J_0(b)|\le3\alpha_0$.
For real $A$ with $s_p^\ell(A)\le\rho_{\mathrm{ec}}$,
\begin{equation}\label{eq:regional-plaquette-change}
 |\partial E(p)-I|\le|\partial U_0(p)-I|
   +[2g_p^\ell(A)+C_{\mathrm{ec}}s_p^\ell(A)^2]
                                      \eta^2\ell^{-2}.
\end{equation}
\end{proposition}

\begin{proof}
The entire local expansion in \cref{I-10-42} applies to each
inner plaquette. Its linear coefficient is the Hessian
$H_{U_0}$, by the one-plaquette calculation in \cref{I-65-4}.
That calculation bounds the local row of $\Delta'$ by
$C_{\mathrm{ec}}\alpha_0\ell^{-2}$, with an absolute constant
also valid for complex variations. No exterior plaquette is added.

For the nonlinear term, apply the local estimate
\eqref{eq:complex-star-remainder} at spacing
$\widehat\eta=\eta/\ell$ and perturbation $\widehat A=\ell A$.
The reference transports are isometries, its curvature bound is
$\alpha_0\widehat\eta^2$, and an incomplete star has
$\widehat\eta=1$. Thus all hypotheses of that estimate hold.
The scaled action is $\ell F$, while a variation of
$\widehat A$ contributes $\ell^{-1}$; the ratio of pairing
weights therefore multiplies its gradient by $\ell^3$.
This gives the second line of \eqref{eq:regional-action-local}.
The local Cauchy-disc argument following
\eqref{eq:complex-star-remainder} gives its third line.
The layer comparisons preceding the proposition give the
three global estimates.

The exponential differential in \cref{I-10-42} is computed
bond by bond, so restricting the plaquette sum gives
\eqref{eq:regional-current-conversion} with its stated factors.
The real plaquette calculation there, at the same scaled
spacing, gives \eqref{eq:regional-plaquette-change} after
enlarging $C_{\mathrm{ec}}$. A nonfull bond has at most three
inner plaquettes and has local scale $\eta$; each contributes
at most $\alpha_0\eta^{-3}$ to $J_0$. This proves the
remaining current bound. Reality follows by differentiation
of the real action in real coordinates.
\end{proof}

We record the regularity conversion used after solving the
equation. The proof of \cref{I-3-113} is local to the four
plaquettes at a full bond. At scale $\ell=\mathsf L(b)$,
its scalar term obeys
\[
 |a(b)|\le r^2\eta^3\ell^{-4}\le r^2\eta^2\ell^{-3},
\]
since every incident plaquette has local scale at least $\ell$.
Thus the same proof gives, for a general separated sequence,
\begin{equation}\label{I-40-66}
 \mathfrak U'_k(\mathbf B,r)\subset\mathfrak U_k(\mathbf B,r)
       \subset\mathfrak U'_k(\mathbf B,r+r^2).
\end{equation}
Suppose $U_0\in\mathfrak U'_k(\mathbf B,\alpha_0)$,
$A$ is real, and
\[
 \max\{\|A\|_X,|D^*DA|_{(-3)}\}\le\alpha'_2,
 \qquad \alpha_2=L^2\alpha'_2.
\]
If
\begin{equation}\label{eq:regional-regularity-thresholds}
 \begin{gathered}
 \alpha_2\le\rho_{\mathrm{ec}},\quad C_{\mathrm{ec}}\alpha_2\le1,
 \quad4\alpha_2+C_{\mathrm{ec}}(\alpha_0+\alpha_2)\le\tfrac12,\\
 0<2\alpha_0+3\alpha_2\le\tfrac13,
 \end{gathered}
\end{equation}
then
\begin{equation}\label{I-40-65}
 e^{i\eta A}U_0\in
       \mathfrak U'_k(\mathbf B,2(2\alpha_0+3\alpha_2)).
\end{equation}
Indeed the plaquette bound has radius at most
$\alpha_0+3\alpha_2$. At a full bond the action expansion
and current conversion bound the current by
$(1+4\alpha_2)(1+C_{\mathrm{ec}}\alpha_2)
(\alpha_0+\alpha_2)\mathsf L(b)^{-3}
\le(2\alpha_0+3\alpha_2)\mathsf L(b)^{-3}$.
For $r=2\alpha_0+3\alpha_2$, the scalar estimate just proved
gives divergence radius at most $r+r^2<2r$, as required.



\subsection{The centered average on a fixed ball}

The straightened mean average must be defined on a ball whose
radius is independent of the depth. For a set $X$ of level-$j$
sites, put $X_i=B^{j-i}(X)$, $0\le i\le j$, and let
$\mathsf P(X_i)$ be its inner positive bonds. Block nesting gives
$X_i=B(X_{i+1})$, so every double block needed at one stage lies
in the preceding member. These restrictions will be called the
tower over $X$.

Put $\Omega_{-1}=\Omega_0$ and
\[
 Y_j=\{y\in\mathbb T^{(j)}:B^j(y)\subset\Omega_{j-1}\},
 \qquad E_j=\mathsf P(Y_j).
\]
For $y\in Y_j$, set $\vartheta_j(y)=1$ if $y\in\Omega_j$
and $\vartheta_j(y)=L$ otherwise; at a bond take the maximum
of the endpoint values. In particular $\vartheta_c=1$ on
interior reduced rows and $L$ on outward rows. The aligned
regions imply that every bond or plaquette in $B^j(X)$,
for $X\subset Y_j$, is inner and has scale at least
$\xi_j/\max_X\vartheta_j$. Its finer towers lie in the
corresponding regions.

\begin{proposition}\label{I-39-36}\label{I-39-43}
There are positive $a_{\mathsf m},c_{\mathsf m}$ and
$C_{\mathsf m},\kappa_{6,\mathsf m}\ge1$, depending only on
$L$, with the following property. Suppose $U_0$ is real and
\[
 |\partial U_0(p)-I|\le\alpha_0\eta^2\mathsf L(p)^{-2},
 \qquad 0<\alpha_0\le L^{-2}a_{\mathsf m},
\]
at every inner bare plaquette. The actual centered straightened
averages $a_j(A)$ and accumulated factors $v_j[A]$ are
holomorphic on
\[
 \mathcal D_{\mathsf m}=
        \{A:|A|_{(-1)}<c_{\mathsf m}/L\},
\]
and are real for real $A$. For $c\in E_j$, let
$\mathsf P_c^j$ be the inner bonds of its endpoint double block
and put $\tau_c(A)=\xi_j\sup_{\mathsf P_c^j}|A|$. Then
\[
 \begin{gathered}
 |a_j(A)(c)|\le2\tau_c(A)\le2\vartheta_c|A|_{(-1)},\\
 a_j(A)=Q_j^{\mathsf m}A+C_j^{\mathsf m}(A),\qquad
 |C_j^{\mathsf m}(A)(c)|\le C_{\mathsf m}\tau_c(A)^2,\\
 \sum_b|Q_j^{\mathsf m}(c,b)|
       \le\xi_j(1+C_{\mathsf m}\vartheta_c^2\alpha_0)
       \le2\xi_j,\qquad
 \sum_b|Q_j^{\mathsf m}(c,b)|\mathsf L(b)^{-1}\le2\vartheta_c.
 \end{gathered}
\]
Here $Q_0^{\mathsf m}=\eta I$. Write $K_A(c,b)=\partial_{A(b)}a_j(A)(c)$.
On $|A|_{(-1)}\le c_{\mathsf m}/(4L)$,
\[
 \begin{split}
 |(K_A-Q^{\mathsf m})\delta A(c)|
      &\le4C_{\mathsf m}\tau_c(A)\tau_c(\delta A),\\
 |K_A\delta A(c)|&\le\tfrac83\vartheta_c|\delta A|_{(-1)}.
 \end{split}
\]
On the whole complex ball $\mathcal D_{\mathsf m}$ the kernels
are supported in $\mathsf P_c^j$ and satisfy
\begin{equation}\label{eq:regional-mean-column}
 \sum_{c\in E_j}|K_A(c,b)|\le2\eta^3\xi_j^{-2},\qquad
 \mathsf L(b)^3\sum_j\xi_j^{-1}
       \sum_{c\in\Gamma_j^{\flat,+}}|K_A(c,b)|
           \le2(L^3+2)\eta^3.
\end{equation}
The factor $v_j[A](y)$ depends only on $A$ in $B^j(y)$,
belongs to $G$ for real $A$, and satisfies
\[
 |v_j[A](y)^{\pm1}-I|
     \le4\kappa_{6,\mathsf m}\xi_j
                        \sup_{\mathsf P(B^j(y))}|A|.
\]
On a tower over $X\subset Y_j$, set
$\vartheta=\max_X\vartheta_j$,
$\alpha=\vartheta^2\alpha_0$, and
$a=\xi_j\sup_{\mathsf P(B^j(X))}|A|$.
The intermediate reference curvatures and perturbations obey
\begin{equation}\label{eq:regional-mean-tower}
 \varpi_i\le2\alpha L^{-2(j-i)},\qquad
 \sup|a_i|\le2aL^{i-j}.
\end{equation}
All constants are independent of the number of layers.
\end{proposition}

\begin{proof}
\emph{The one-step chart.}
By \cref{lem:one-step-straightening}, the centered chart has
a fixed complex radius $r_{\mathsf m}>0$ near the identity.
For a real $W$ of sufficiently small double-block curvature,
the comb gauge of \cref{I-10-51} gives $W^h=e^{ib}$ with
$\sup|b|\le2\pi L\varpi_c(W)<r_{\mathsf m}/2$.
The covariance formula \eqref{eq:straightening-covariance}
therefore gives the same perturbation ball at $W$.
Decrease the radius so that the logarithms have norm at most
one. Their vanishing at zero and Cauchy's estimate give
constants $\kappa_{5,\mathsf m},\kappa_{6,\mathsf m}$ with
\[
 |\log\mathsf v(W,a)|,\ |\mathcal Q^1_{\mathsf m}(W,a)|
       \le\kappa_{6,\mathsf m}\sup|a|,
 \quad |C^1_{\mathsf m}(W,a)|\le\kappa_{5,\mathsf m}(\sup|a|)^2,
\]
and $|\partial_{a(b)}C^1_{\mathsf m}(W,a)|
\le\kappa_{5,\mathsf m}\sup|a|$ on the half ball.
These maps are real on real inputs.

The leading transport kernel is
\[
 \mathbf S_{\mathsf m}(W)(c,b)=L^{-3}
   \sum_{\substack{x\in B(c_-)\\ b\in[x,x']}}
       \operatorname{Ad}(W_{\mathsf m}(c_-,x)W([x,b_-])),
 \qquad x'=x+Le_\mu.
\]
Its row and column bounds are $L,L^{-2}$. At the identity
the endpoint terms cancel, so
$\mathbf Q_{\mathsf m}(I)=\mathbf S_{\mathsf m}(I)$.
The Schwarz estimate in the proof of \cref{I-32-37}, now
using the centered chart, bounds each difference entry by
$(r_{\mathsf m}^{-1}+2)r_{\mathsf m}^{-1}\sup|b|$.
There are at most $6L^3$ entries in one row and six double
blocks containing one fine bond. Gauge covariance thus gives
a common $\kappa_9=\kappa_{9,\mathsf m}$ such that
\begin{equation}\label{eq:regional-mean-step}
 \begin{split}
 \|\mathbf Q_{\mathsf m}(W)-\mathbf S_{\mathsf m}(W)\|
       _{\mathrm{row},\mathrm{col}}&\le\kappa_9\varpi,\\
 \|d_a\mathcal Q^1_{\mathsf m}(W,a)\|_{\mathrm{col}}
       &\le L^{-2}+\kappa_9\varpi+6\kappa_5\sup|a|.
 \end{split}
\end{equation}
We abbreviate $\kappa_5=\kappa_{5,\mathsf m}$ below.
Single-block factors have the same domain: fold into an
adjacent block as in \cref{I-32-37}, extend $a$ by zero,
and use their dependence only on the original block.

\emph{Iteration.}
Consider a tower with initial curvature at most
$\alpha L^{-2j}$ and initial perturbation $a_0$ with
$\tau=L^j\sup|a_0|<c_{\mathsf m}$.
The local proof of \cref{I-10-101} gives
\[
 \varpi_{i+1}\le L^2\varpi_i+C_0^{\mathrm{SO}}\varpi_i^2,
 \qquad \varpi_i\le2\alpha L^{-2(j-i)},
 \qquad \sum_{i<j}\varpi_i\le4\alpha L^{-2},
\]
after decreasing $a_{\mathsf m}$ to stay in the one-step
domain. The averages agree with the real branch by the same
lemma. Choose the radii also so that
\[
 2c_{\mathsf m}\le r_{\mathsf m}/2,\qquad
 4\kappa_9a_{\mathsf m}\le\tfrac12\log2,
 \qquad24\kappa_5Lc_{\mathsf m}\le\tfrac12\log2.
\]
For $t_i=L^{-i}\sup|a_i|$, the expansion gives
$t_{i+1}\le t_i(1+\kappa_9L^{-1}\varpi_i+\kappa_5L^{i-1}t_i)$.
If $t_i\le2t_0$ through level $n<j$, multiplication yields
\[
 t_{n+1}\le t_0
    \exp(4\kappa_9L^{-3}\alpha+4\kappa_5L^{-2}\tau)
       \le2t_0.
\]
This closes the induction and all its domain conditions.
Every composition is holomorphic on the entire initial ball.
Restricting to an endpoint tower proves locality.

For the derivative product $T_{j-1}\cdots T_0$, where
$T_i=d_a\mathcal Q^1_{\mathsf m}(W_i,a_i)$, we have
\[
 \|T_i\|_{\mathrm{col}}\le L^{-2}(1+\varepsilon_i),\qquad
 \varepsilon_i=L^2(2\kappa_9\alpha L^{-2(j-i)}
                                  +12\kappa_5\tau L^{i-j}).
\]
Since $\sum_{i<j}\varepsilon_i
\le4\kappa_9\alpha+24\kappa_5L\tau\le\log2$,
submultiplicativity gives a column bound $2L^{-2j}$.
The single-block factors are holomorphic compositions on
this same domain.

\emph{The regional estimates.}
The geometry above gives the tower hypotheses with
$\alpha=\vartheta^2\alpha_0\le a_{\mathsf m}$ and
$a_0=\eta A$, since
$\tau\le\vartheta|A|_{(-1)}<c_{\mathsf m}$.
This proves the domain, value and intermediate bounds.
At zero the chain rule gives
$Q_j^{\mathsf m}=\eta\mathbf Q_{\mathsf m}(W_{j-1})
\cdots\mathbf Q_{\mathsf m}(W_0)$.
Its row sum is at most
$\xi_j\exp(4\kappa_9L^{-3}\vartheta_c^2\alpha_0)$.
Choose $C_{\mathsf m}\ge8\kappa_9L^{-3}$ and shrink
$a_{\mathsf m}$ so that $C_{\mathsf m}a_{\mathsf m}\le1$.
The row bounds follow, including the weighted bound since
$\mathsf L(b)\ge\xi_j/\vartheta_c$ on its support.

For the remainder use the local ball $\tau_c(A)<c_{\mathsf m}$,
where $|a_j(A)(c)|\le2\tau_c(A)$. Cauchy's estimate gives
$|C_j^{\mathsf m}(A)(c)|\le4\tau_c(A)^2/c_{\mathsf m}$
on the half ball; outside it the value and linear bounds
give $|C_j^{\mathsf m}(A)(c)|\le4\tau_c(A)
\le8\tau_c(A)^2/c_{\mathsf m}$.
Enlarge $C_{\mathsf m}$ to at least $8/c_{\mathsf m}$,
and decrease $a_{\mathsf m}$ again so that $C_{\mathsf m}a_{\mathsf m}\le1$.
On the quarter ball apply Cauchy's formula to
$C_j^{\mathsf m}(A+z\delta A)(c)$ with
$|z|<\tau_c(A)/\tau_c(\delta A)$; the quadratic bound
on this disc proves the stated derivative estimate.
Zero local seminorms are handled by locality.
Direct Cauchy estimation of $a_j$ on the global ball of
radius $c_{\mathsf m}/L$ gives the factor $8/3$.

The derivative product on the tower over all $Y_j$ gives
$\sum_{c\in E_j}|K_A(c,b)|\le2\eta L^{-2j}$.
If $\mathsf L(b)=\xi_n$, only levels $n-1,n,n+1$ contribute
to reduced rows; multiplication by $\xi_n^3\xi_j^{-1}$
and summation proves \eqref{eq:regional-mean-column}.
Finally write the single-step factors as $e^{iX_i}$.
Their logarithm estimates give
\[
 \sum_{i<j}|X_i(y)|\le4\kappa_{6,\mathsf m}L^{j-1}
                       \sup_{\mathsf P(B^j(y))}|\eta A|.
\]
Choose $4\kappa_{6,\mathsf m}c_{\mathsf m}\le1$.
The product estimate $|\prod e^{\pm iX_i}-I|
\le e^{\sum|X_i|}-1$, with reversed order for the inverse,
then gives the asserted bound for $v_j^{\pm1}$.
\end{proof}



\subsection{Comparison of the two averages}

The inverse just constructed uses corner averages. The nonlinear
constraint uses the centered averages, so its linearization must
be compared with the corner operator. There is a curvature-sized
error even after their output frames have been aligned.

Keep the regional setting of \cref{I-42-86}. For $y\in Y_j$, let $o_j(y)$ be the lowest corner of its order-$j$
block and put
\[
 t_j(y)=U_0(\Gamma_{y,o_j(y)}),\qquad
 (T\mu)(y)=\operatorname{Ad}(t_j(y))\mu(y),\qquad
 (T\omega)(c)=\operatorname{Ad}(t_j(c_-))\omega(c).
\]
The path is the bare comb within the block. Thus $t_0=I$,
and $T$ preserves the sup norms and the scalar-top and row
pairings. Write $Q'_{\mathrm m}$ for the scalar average formed
from the accumulated centered transports, and $Q_{\mathrm m}$
for the derivative $Q^{\mathsf m}$ in \cref{I-39-43}. The
corner averages retain the notation $Q',Q$.
For a reduced row $c$ set $\vartheta_c=1$ if it is interior
and $\vartheta_c=L$ if it is outward.

\begin{lemma}\label{I-44-26}
There are $a_F>0$ and $\kappa_F\ge1$, depending only on $L$,
such that, under the regional hypotheses with $\alpha_0\le a_F$,
\begin{equation}\label{eq:mean-corner-comparison}
 Q'_{\mathrm m}=T(Q'+E'),\qquad Q_{\mathrm m}=T(Q+E).
\end{equation}
Both errors are real, vanish at level zero and have the same
cell or double-block support as the corresponding averages.
For $y\in\Lambda_j$ and $x\in B^j(y)$,
\[
 |E'(y,x)|\le\kappa_F\alpha_0 L^{-3j}.
\]
For a positive reduced row $c$ of level $j$,
\begin{equation}\label{eq:mean-corner-row-column}
 \sum_b|E(c,b)|\le\kappa_F\vartheta_c^2\alpha_0\xi_j,
 \qquad
 \sum_{c\in\Gamma_j^{\flat,+}}|E(c,b)|
       \le\kappa_FL^2\alpha_0\eta^3\xi_j^{-2}.
\end{equation}
Consequently
\begin{equation}\label{eq:mean-corner-weighted-column}
 \mathsf L_{\mathbf B}(b)^3
     \sum_j\xi_j^{-1}\sum_{c\in\Gamma_j^{\flat,+}}|E(c,b)|
       \le\kappa_FL^2(L^3+2)\alpha_0\eta^3.
\end{equation}
The constants are independent of the number of layers.
\end{lemma}

\begin{proof}
We give the comparison for both choices of root in each block.
The corner estimate in \cref{I-32-37} and the centered
estimate \eqref{eq:regional-mean-step} give a common
$\kappa_9=\kappa_9(L)$ with
\[
 \|\mathbf Q_\star(W)-\mathbf S_\star(W)\|
           _{\mathrm{row},\mathrm{col}}
       \le\kappa_9\varpi(W),\qquad \star\in\{\mathrm m,\mathrm c\}.
\]
Choose $a_F\le L^{-2}\min\{a_{\mathsf m},c_{\mathsf c}\}$,
and decrease it so that both one-step charts are available
throughout every contributing tower.

For a level-$j$ row let
$\mathsf D_c=B^j(c_-)\cup B^j(c_+)$ and
$a_c=\vartheta_c^2\alpha_0$. Every plaquette there has curvature
at most $a_cL^{-2j}$. The nested restrictions and folding
argument of \cref{I-32-37}, followed by \cref{I-10-101}, give
\[
 \varpi_i\le2a_cL^{-2(j-i)},\qquad
 \sum_{i<j}\varpi_i\le4a_cL^{-2}.
\]
All factors contributing to row $c$ use only $\mathsf D_c$.
For a fixed input bond, take the union of the towers of rows
whose double blocks contain it; the same bounds hold there
with $\bar a_j(b)\le L^2\alpha_0$ in place of $a_c$.

Write $Q_j^{\star,\mathrm r}$ for the difference between
$Q_j^\star$ and the product of its leading kernels, multiplied
by $\eta$. The ordered-product calculation already used in
\cref{I-32-37} gives
\[
 \begin{split}
 \sum_b|Q_j^{\star,\mathrm r}(c,b)|
 &\le\eta L^j
       [\exp(4\kappa_9L^{-3}a_c)-1],\\
 \sum_c|Q_j^{\star,\mathrm r}(c,b)|
 &\le\eta L^{-2j}
       [\exp(4\kappa_9\bar a_j(b))-1].
 \end{split}
\]
Require $4\kappa_9L^2a_F\le\log2$. Since
$e^u-1\le2u$ on this interval, these are bounded by
$\kappa_Sa_c\xi_j$ and
$\kappa_S\bar a_j(b)\eta^3\xi_j^{-2}$, respectively,
where $\kappa_S=8\kappa_9$. This retains the column estimate
before any sum over levels is taken.

The leading products in the two schemes have the same nested
chains. Compare their transports to a single bare path in
$\mathsf D_c$. In the centered scheme first replace the contour
mean by one of its defining contour transports. The mean estimate
in \cref{I-10-61} bounds this error by
$\kappa_1^{\mathrm{SO}}\varpi_i
 \le2\kappa_1^{\mathrm{SO}}a_cL^{-2(j-i)}$.
Each one-step contour and initial straight segment has length
at most $4L$. The path-replacement argument in the proof of
\cref{I-32-37}, using \eqref{I-10-102}, therefore costs at most
$(4L\kappa_4+2\kappa_1^{\mathrm{SO}})a_cL^{-2(j-i)}$
at level $i$. Summing these errors and using the comb gauge
to compare the resulting bare paths gives a common
$\kappa_7(L)$ such that
each accumulated transport differs from a same-endpoint bare
path of length at most $16L^j$ by at most $\kappa_7a_c$.
For scalar transports use $\alpha_0$ in place of $a_c$.

For a scalar output $y$, take the comb from $o_j(y)$ to $x$.
Prepending the path from $y$ to $o_j(y)$ gives the comparison
path for the centered transport. Both lie in $B^j(y)$ and
have length at most $5L^j$. After multiplying by $t_j(y)^{-1}$,
the two adjoint transports differ by at most $4\kappa_7\alpha_0$.
The scalar averaging factor $L^{-3j}$ gives the asserted
kernel estimate for $E'$.

For the bond kernels take the comb from $o_j(c_-)$ to $b_-$;
its length is at most $4L^j$. The centered comparison path,
with the frame path prepended, has length at most $6L^j$.
The difference of the aligned adjoint transports is again at
most $4\kappa_7a_c$. There are $L^{4j}$ chains from one row
and at most $L^j$ chains ending at one input bond, as counted
in the proof of \cref{I-32-37}. Multiplication by $\eta L^{-3j}$
therefore gives the respective bounds
\[
 4\kappa_7a_c\xi_j,\qquad
 4\kappa_7\bar a_j(b)\eta^3\xi_j^{-2}.
\]
Adding the two derivative remainders proves
\eqref{eq:mean-corner-row-column} with
$\kappa_F=2\kappa_S+4\kappa_7$.

If $\mathsf L_{\mathbf B}(b)=\xi_i$, only levels
$i-1,i,i+1$ have a row double block containing $b$.
The last sum is thus bounded by
$\kappa_FL^2\alpha_0\eta^3(L^3+1+L^{-3})$, which proves
\eqref{eq:mean-corner-weighted-column}. At level zero both
schemes are $\eta I$ on bonds and $I$ on sites, so their
errors vanish. Reality and the support assertions follow
from the corresponding real kernels and their nested locality.
\end{proof}


\subsection{Prescribing the induced gauge}

Put $a_{\rm in}=\min\{a_Q,a_F,L^{-2}a_{\mathsf m},1/2\}$.
Assume the preceding regional geometry and
$U_0\in\mathfrak U'_k(\mathbf B,\alpha_0)$, where
$\alpha_0\le a_{\rm in}$.
Put
\[
 \|\lambda\|_1=\max\{\sup|\lambda|,|D\lambda|_{(-1)}\},
 \qquad
 \|\lambda\|_Y=\max\{\|\lambda\|_1,|\Delta\lambda|_{(-2)}\}.
\]
For $w=e^{i\lambda}$, write
$e^{i\eta A^w}U_0=(e^{i\eta A}U_0)^w$, using the small
logarithm at each inner bond. Its induced factor at a top
$y\in\Lambda_j$ is
\[
 \mathcal R_{\mathbf B}(U_0,A;w)(y)
       =v_j[A^w](y)^{-1}w(y)v_j[A](y).
\]
Here $v_j$ is the accumulated factor of the mean chart;
at level zero it equals $1$.

\begin{lemma}\label{I-46-44}
There are $c_6>0$ and $K_3$, depending only on $L$,
such that, if $\alpha_0\le c_6$,
$A_0$ is a real bond field with $\|A_0\|_X\le c_6$, and
$r:\Lambda_*\to SU(2)$ satisfies $\sup|r-1|\le c_6$,
there is a real site function $\lambda$ for which
$w=e^{i\lambda}$ satisfies
\[
 RD^*A_0^w=0,\qquad
 \mathcal R_{\mathbf B}(U_0,A_0;w)=r
       \quad\hbox{on }\Lambda_*,
\]
and
\[
 \|\lambda\|_Y\le
 K_3\bigl(|D^*A_0|_{(-2)}+\sup|r-1|\bigr),
 \qquad
 |A_0^w-A_0|_{(-1)}\le2C_{\rm gc}\|\lambda\|_Y.
\]
In particular, $w=r$ on $\Lambda_0$.
\end{lemma}

\begin{proof}
\emph{The scalar system.}
Use the regional scalar operators
\[
 G'=(\Delta+Q'^*aQ')^{-1},\qquad
 N=Q'G'^2Q'^*,\qquad C=N^{-1},\qquad
 P=G'Q'^*C Q'G',\quad R=1-P.
\]
The bare and top pairings have weights $\eta^3$ and
$\xi_y^3$, respectively. The real comb transports are
isometries, and their pillars partition the bare sites;
hence $Q'Q'^*=1$. In particular,
$\operatorname{ran}R=\ker(Q'G')$. Define
\[
 H'=G'Q'^*a+G'^2Q'^*C(1-Q'G'Q'^*a).
\]
Then
\[
 Q'H'=1,\qquad R\Delta H'=0,\qquad
 R\Delta G'f=f\quad(f\in\operatorname{ran}R).
\]
Indeed, if $\tau=(1-Q'G'Q'^*a)\nu$ and $\sigma=C\tau$,
then $Q'H'\nu=\nu$ and
$\Delta H'\nu=G'Q'^*\sigma$, which $R$ annihilates.
For $f\in\ker(Q'G')$, the identity
$\Delta G'=1-Q'^*aQ'G'$ gives $\Delta G'f=f$.
It follows that
\[
 (f,\nu)\longmapsto G'f+H'\nu
\]
is the inverse of $\lambda\mapsto(R\Delta\lambda,Q'\lambda)$:
the displayed identities give a right inverse, and the
dimensions agree because $Q'G'$ is onto.

Increase the common scalar estimate constant $K_1\ge1$
if necessary. The kernel bounds and weighted sums give
\[
 \begin{gathered}
 \|G'f\|_Y\le(1+K_1)|f|_{(-2)},\qquad
 |Rf|_{(-2)}\le(1+K_1^3)|f|_{(-2)},\\
 \|H'\nu\|_Y\le K_2\sup|\nu|,\qquad K_2=(1+K_1)^4.
 \end{gathered}
\]
The two weight conversions used here are
$|G'g|_{(-2)}\le K_1|g|_{(-4)}$ and
$|(C\nu)(y)|\le K_1\xi_y^{-4}\sup|\nu|$.
They bound $P$ by three factors $K_1$.
For $H'$, its first term is bounded by $1+K_1$;
the second is bounded by $(1+K_1)K_1^2(1+K_1)$.
Their sum is at most $(1+K_1)^4$.
The $\Delta$ component of the first estimate follows
from $\Delta G'f=f-Q'^*aQ'G'f$.

We need these identities for the inner-bond divergence,
also when the level-zero values are prescribed.
If $n(x)$ is the number of missing neighbours, then
$\Delta=D^*D+\eta^{-2}n$.
Let $\Lambda_0^\circ$ consist of the level-zero sites
whose neighbours in $\Omega_0$ are all in $\Lambda_0$.
For $f$ supported there,
$(G')^{-1}f=\Delta f+Q'^*aQ'f$ is supported on
$\Lambda_0$ and therefore equals $Q'^*\mu$ for a top
function $\mu$. Thus $f=G'Q'^*\mu$ and $Rf=0$.
Separation puts every site with $n(x)>0$ in
$\Lambda_0^\circ$: it and its neighbours are within
$2\eta$ of $\Omega_0^c$, while $\Omega_1$ is at distance
greater than $R_1\mathfrak M\eta\ge L^2\eta$.
Consequently
\[
 RD^*D=R\Delta,\qquad
 \|\lambda\|_{Y^\circ}\le7\|\lambda\|_Y,\qquad
 \|\lambda\|_{Y^\circ}
   :=\max\{\|\lambda\|_1,|D^*D\lambda|_{(-2)}\}.
\]
The factor $7$ follows from $n\le6$ and $\xi(x)=\eta$
where $n(x)>0$.

\emph{The divergence remainder.}
For $|A|_{(-1)},\|\lambda\|_1\le\rho_2=1/64$, the
small bond logarithm is real analytic and satisfies
\begin{equation}\label{eq:gauge-correction-divergence}
 \begin{split}
 D^*A^{e^{i\lambda}}
    &=D^*A-D^*D\lambda+\mathcal N(A,\lambda),\\
 |\mathcal N(A,\lambda)|_{(-2)}
    &\le C_{\rm gc}
       (\|A\|_X+\|\lambda\|_{Y^\circ})
                   \|\lambda\|_{Y^\circ},\\
 |A^{e^{i\lambda}}-A|_{(-1)}
    &\le2C_{\rm gc}\|\lambda\|_1 ,
 \end{split}
\end{equation}
with $C_{\rm gc}=93$.
Here is the cancellation giving the middle estimate.
At a site put $l=\lambda(x)$, and at an adjacent bond
put $u=\eta A$ and $d=\eta D\lambda$, transporting
incoming values to $x$. The outgoing and incoming
logarithms are
\[
 \Psi_{\rm out}=\frac1i\log(e^{il}e^{iu}e^{-i(l+d)}),
 \qquad
 \Psi_{\rm in}=\frac1i\log(e^{i(l-d)}e^{iu}e^{-il}).
\]
Both have the expansion
\[
 \Psi=\operatorname{Ad}(e^{il})u-k(l)d+\Psi^{(2)},
 \qquad
 |\Psi^{(2)}|\le14|d|\max\{|u|,|d|\},
\]
where $k(l)=(e^{i\operatorname{ad}l}-1)/(i\operatorname{ad}l)$,
defined by its power series at zero.
The derivative in $(u,d)$ gives the first two terms;
the remainder vanishes when $d=0$, and its derivative
at $(u,d)=0$ vanishes. The small logarithm estimate
and the mixed-vanishing Cauchy bound give the displayed
estimate with its factor $|d|$.
Also $|\operatorname{Ad}(e^{il})-1|\le3|l|$ and
$|k(l)-1|\le3|l|$.
In the signed sum at $x$, the linear terms combine
into
$(\operatorname{Ad}(e^{il})-1)D^*A
 -(k(l)-1)D^*D\lambda$.
The at most six quadratic terms contribute
$84\max|D\lambda|\max\{|A|,|D\lambda|\}$.
Multiplying by $\xi(x)^2$, using
$\xi(x)\le\mathsf L(b)$ at every adjacent inner bond
and $|D^*A|_{(-2)}\le3\|A\|_X$, proves the estimate.
The same expansion proves the last bound.
Only existing inner bonds enter this calculation.

\emph{The induced factors.}
There are $c_{\rm top}>0$ and $C_{\rm top}$ depending
only on $L$ such that, for
$|A|_{(-1)},\|\lambda\|_1\le L^{-1}c_{\rm top}$,
all induced top factors are defined and
\begin{equation}\label{eq:gauge-correction-top}
 \begin{split}
 T^{-1}\frac1i
       \log\mathcal R_{\mathbf B}(U_0,A;e^{i\lambda})
       &=Q'\lambda+q(A,\lambda),\\
 \sup|q(A,\lambda)|
 &\le[\kappa_F\alpha_0+
       C_{\rm top}(|A|_{(-1)}+\|\lambda\|_1)]
          \|\lambda\|_1 .
 \end{split}
\end{equation}
The maps in this display are real analytic on a
neighbourhood of the closed admitted set.
We choose $c_{\rm top}\le c_{\mathsf m}/4$.
For $w=e^{i\lambda}$ the real angle estimate at a bond gives
\[
 |A^w|_{(-1)}\le |A|_{(-1)}+\tfrac\pi2|D\lambda|_{(-1)}
       <c_{\mathsf m}/L.
\]
Indeed the angle of $w(b_-)e^{i\eta A(b)}
U_0(b)w(b_+)^{-1}U_0(b)^{-1}$ is at most
$\eta|A(b)|+(\pi/2)\eta|D\lambda(b)|$; it is in the
small logarithm branch. Thus both $A$ and $A^w$ lie in the
mean domain of \cref{I-39-43}. We prove the top-factor
estimate uniformly in the depth.

On one mean block let $r$ be the incoming gauge factor and put
\[
 S=\sup|r-1|,\qquad
 g=\max_b\theta(r(b_-)^{-1}\operatorname{Ad}(W(b))r(b_+)).
\]
Let $\varpi$ bound the reference plaquettes.
The one-step factor satisfies
\begin{equation}\label{eq:gauge-correction-one-step}
 \begin{split}
 S'&\le(1+\kappa\varpi)S+\kappa g,\\
 g'&\le Lg+\kappa(\varpi S+g^2+|a|g),\\
 \rho'&=q_W^{\mathsf m}\rho+u,\qquad
 |u|\le\kappa(S+|a|)(g+\varpi S),
 \end{split}
\end{equation}
where $\rho=(1/i)\log r$ and $\kappa=\kappa(L)$.
To see the last estimate, write
$\widehat r(x)=\operatorname{Ad}(W(y,x))r(x)$ and
$\delta_x=(1/i)\log(r(y)^{-1}\widehat r(x))$.
Comparison with an ordered path in the block gives
$|\delta_x|\le C(L)(g+\varpi S)$: the path contributes
its length times $g$, while the change of transport
contributes $\varpi S$ by the conjugation estimate.
The mean-chart formula gives
$r'(y)=r(y)e^{i\tau(y)}$ with
$\tau(y)=L^{-3}\sum_x\delta_x+
 O_L(|a|(g+\varpi S))$.
Expanding the two logarithms then gives the first and
third estimates in \eqref{eq:gauge-correction-one-step}.
For the second, compare the two endpoint formulas
along each straight path of length $L$ joining the
blocks. Before averaging their logarithms, its
transported gauge difference has angle at most
$Lg+C(L)\varpi S$. The endpoint corrections are
the averages of $\delta_x$, up to the displayed error.
The quadratic logarithm errors are bounded by
$C(L)(g+\varpi S)^2$; with the fixed small radii
this is at most $C(L)(g^2+\varpi S)$.
This proves the second estimate with its leading
coefficient exactly $L$.
The estimates at a single block follow by reflection
into an adjacent double block: reflected bonds retain
their values up to reversal, and the joining corridor
has identity reference and zero potential. Curvature,
gauge-difference and size bounds are preserved, and
locality leaves the factor at the original block
unchanged.

Apply these estimates in the tower ending at a top
$y\in\Lambda_j$. The preceding mean-chart estimates give
\[
 \begin{gathered}
 \varpi_i\le2\alpha L^{-2(j-i)},\qquad
 |a_i|\le2aL^{i-j},\qquad
 a:=\xi_j\sup_{B^j(y)}|A|\le |A|_{(-1)},\\
 \epsilon:=\max\{\sup_{B^j(y)}|\lambda|,
                 \xi_j\sup_{B^j(y)}|D\lambda|\}
                 \le\|\lambda\|_1 .
 \end{gathered}
\]
where initially $S_0\le\epsilon$ and
$g_0\le(\pi/2)L^{-j}\epsilon$.
For the latter inequality use
$\theta(e^{-i\lambda(x)}
 e^{i(\lambda(x)+\eta D\lambda(b))})
 \le(\pi/2)\eta|D\lambda(b)|$.
Here the bond suprema range over bonds with both ends
in $B^j(y)$. At tops the tower lies entirely in the
corresponding pillar, so its scale-conversion factor
is $1$. The reference curvature hypothesis supplies
$\alpha\le L^2\alpha_0$. We reduce the reference radius
so that $\alpha\le\alpha_*\le a_{\mathsf m}$, with
$\alpha_*$ chosen below.

Put $G_i=L^{j-i}g_i$ and $S_*=2+32\kappa$.
The recurrences become
\[
 \begin{split}
 G_{i+1}&\le
 [1+\kappa L^{-1}L^{i-j}(G_i+2a)]G_i
       +2\kappa\alpha L^{i-j-1}S_i,\\
 S_{i+1}&\le(1+\kappa\varpi_i)S_i+
                         \kappa L^{i-j}G_i.
 \end{split}
\]
Choose the tower radii so that
$20\kappa L^{-2}c_{\rm top}\le1/4$,
$4\kappa S_*L^{-2}\alpha_*\le1$, and
$4\kappa L^{-2}\alpha_*\le1/4$,
also keeping $2c_{\rm top}$, $8c_{\rm top}$ and
$S_*c_{\rm top}$ within the one-step chart radii.
Using $\sum_{i<j}L^{i-j}\le2L^{-1}$, discrete
Gronwall closes the simultaneous induction
\[
 G_i\le8\epsilon,\qquad S_i\le S_*\epsilon.
\]
Indeed the multiplicative errors in each recurrence
sum to at most $1/4$, and their inhomogeneous sums
are at most $\epsilon$ and $16\kappa L^{-1}\epsilon$,
respectively; $e^{1/4}<4/3$ gives the stated bounds.
All intermediate logarithms remain in their charts.
Iterating the last relation in
\eqref{eq:gauge-correction-one-step} gives
\[
 \rho_j=q_{j-1}^{\mathsf m}\cdots q_0^{\mathsf m}\lambda
          +C'_j,\qquad
 |C'_j|\le C_{\rm top}(a+\epsilon)\epsilon.
\]
Each real scalar mean average has norm at most one.
The sum of its remainders is bounded by
\[
 \kappa(S_*\epsilon+2a)
       \sum_{i<j}
        (8L^{i-j}\epsilon+
                    2\alpha L^{-2(j-i)}S_*\epsilon).
\]
This proves the asserted bound uniformly in $j$.
To identify this factor, let $a'_i,v'_i$ be the averages and
accumulated factors of $A^w$. The factorization
\eqref{eq:straightening-factorization} and real gauge covariance
give inductively
\[
 e^{ia'_i}W_i=(e^{ia_i}W_i)^{r_i},\qquad
 v'_i r_i=w v_i[A].
\]
At the next level use
$r_{i+1}=\mathsf v(W_i,a'_i)^{-1}r_i\mathsf v(W_i,a_i)$;
multiplication proves the second identity at level $i+1$.
The small logarithm identifies the first identity with the
recursion for $a'_i$. At $i=j$ the second identity gives
$r_j=v_j[A^w]^{-1}w v_j[A]$, precisely the induced factor.
Nested locality identifies the tower values with those of
the sequence.
All means and logarithms satisfy strict domain
inequalities at the admitted real pairs; their finite
compositions are therefore real analytic on a
neighbourhood of those pairs.

The iterated scalar mean average is
$Q'_{\rm mean}=T(Q'+E')$, where the earlier
mean/corner comparison gives
$|(E'\lambda)(y)|\le\kappa_F\alpha_0|\lambda|_y$.
Thus $q=E'\lambda+T^{-1}C'_j$ on $\Lambda_j$,
which proves \eqref{eq:gauge-correction-top}.
On $\Lambda_0$ the induced factor is $e^{i\lambda}$,
$Q'\lambda=\lambda$ and $T=1$, so $q=0$.

\emph{Solving the two equations.}
Set
\[
 K'=1+K_1+K_2,\qquad
 K_3=4(1+K_1^3)K',\qquad K''=4K_3 .
\]
Choose $c_6\le\min\{a_{\rm in},L^{-2}\alpha_*\}$ small enough that
\[
 \begin{gathered}
 K''c_6\le\min\{\rho_2,L^{-1}c_{\rm top}\},\\
 7C_{\rm gc}(1+K_1^3)K'(1+7K'')c_6\le\tfrac12,\\
 [\kappa_F+C_{\rm top}(1+K'')]K'c_6\le\tfrac12 .
 \end{gathered}
\]
All these choices precede the sequence and its depth.
Put $m_0=|D^*A_0|_{(-2)}+\sup|r-1|$ and
$s=4(1+K_1^3)m_0$. Since
$|D^*A_0|_{(-2)}\le3\|A_0\|_X$, we have
$m_0\le4c_6$ and $K's\le K''c_6$.
On the real max-norm ball
\[
 f\in\operatorname{ran}R,\quad |f|_{(-2)}\le s,
 \qquad \sup|\nu|\le s,
\]
put $\lambda=G'f+H'\nu$ and define
\[
 \Psi(f,\nu)=
 \left(R[D^*A_0+\mathcal N(A_0,\lambda)],
       T^{-1}\frac1i\log r-q(A_0,\lambda)\right).
\]
The logarithm is defined and its norm is at most
$(\pi/2)\sup|r-1|$. The scalar estimates give
$\|\lambda\|_Y\le K's$ and
$\|\lambda\|_{Y^\circ}\le7K's$.
The chosen caps and the two remainder estimates imply
\[
 |\Psi_1|_{(-2)}\le\tfrac14s+\tfrac12s,
 \qquad
 \sup|\Psi_2|\le\tfrac\pi8s+\tfrac12s .
\]
This continuous map sends the ball into its strict
interior and is $C^2$ on the open ball.

For completeness, the finite-dimensional fixed-point
fact used here follows by a degree integral.
If such a map had no fixed point, put
$F_t(p)=p-t\Psi(p)$ and choose a $C^1$ bump $\omega$
of integral one supported sufficiently near zero.
The functions $\omega(F_t)\det DF_t$ are supported
in a common compact subset of the open ball.
If $C_{ji}$ are the cofactors of $DF_t$, then
$\sum_i\partial_iC_{ji}=0$, by equality of mixed
derivatives and alternation of the determinant, and
\[
 \partial_t[\omega(F_t)\det DF_t]
   =-\sum_i\partial_i
        \left[\omega(F_t)\sum_j C_{ji}\Psi_j\right].
\]
Its integral is constant in $t$, equal to $1$ at
$t=0$ and to $0$ at $t=1$, a contradiction.
The zero-radius case is immediate.

Take a fixed point and its $\lambda=G'f+H'\nu$.
The scalar identities give $R\Delta\lambda=f$ and
$Q'\lambda=\nu$. The first fixed-point equation,
\eqref{eq:gauge-correction-divergence}, and
$RD^*D=R\Delta$ give $RD^*A_0^w=0$.
The second equation and
\eqref{eq:gauge-correction-top} identify the two
small logarithms and hence the induced top factors.
Finally $\|\lambda\|_Y\le K's=K_3m_0$ and
\eqref{eq:gauge-correction-divergence} give the bounds.
\end{proof}


\subsection{The nonlinear rows and gauge directions}

Keep the real background $U_0$, the reduced corner average $Q$,
and the weighted spaces of \cref{I-45-14}. Write
$a_j(A)=\mathcal Q_j^{\mathsf m}(U_0;A)$ for the straightened
averages. On the positive interior and outward bonds of
$\Gamma_j^{\flat,+}$ put $\mathcal Q(A)(c)=a_j(A)(c)$ and
$K_A=\mathcal Q'(A)$. The frame rotation
$(T\omega)(c)=\operatorname{Ad}(t_j(c_-))\omega(c)$ preserves
both the row inner product and the supremum norm. Define
\[
 \mathcal C(A)=T^{-1}\mathcal Q(A)-QA,\qquad
 \widehat K_A=T^{-1}K_A=Q+\mathcal C'(A).
\]
The mean/corner comparison gives
\[
 \mathcal C(A)=EA+T^{-1}(C_j(A))_{j\le k},\qquad \mathcal C'(0)=E.
\]
In particular the first term is linear and is retained throughout.

\begin{lemma}\label{I-45-23}\label{I-45-28}
Put $C_3=68/c_{\mathsf m}$ and
\[
 \epsilon(\alpha,r)=\kappa_FL^2\alpha
                         +4(C_{\mathsf m}+1)Lr,\qquad
 \epsilon^*(\alpha,r)=\kappa_FL^2\alpha+C_3Lr.
\]
The map $\mathcal Q$ is real-analytic on
$|A|_{(-1)}<c_{\mathsf m}/L$. If $|A|_{(-1)},|A_i|_{(-1)}\le r\le c_{\mathsf m}/(4L)$,
then, with $\epsilon=\epsilon(\alpha_0,r)$ and
$\epsilon^*=\epsilon^*(\alpha_0,r)$,
\[
 \begin{gathered}
 \sup|\mathcal C(A)|\le L\epsilon|A|_{(-1)},\qquad
 \sup|\mathcal C'(A)\delta A|\le L\epsilon|\delta A|_{(-1)},\\
 \sup|\mathcal C(A_1)-\mathcal C(A_2)|\le L\epsilon|A_1-A_2|_{(-1)},\\
 |\mathcal C'(A)^*\omega|_{(-3)}\le3L^3\epsilon^*\sup|\omega|,\\
 |(\mathcal C'(A_1)^*-\mathcal C'(A_2)^*)\omega|_{(-3)}
       \le2L^4C_3|A_1-A_2|_{(-1)}\sup|\omega|.
 \end{gathered}
\]
\end{lemma}

\begin{proof}
\Cref{I-39-43,I-44-26}
give, for a level-$j$ row $c$,
\[
 |\mathcal C'(A)\delta A(c)|\le
 (\kappa_FL^2\alpha_0\vartheta_c
       +4C_{\mathsf m}\vartheta_c^2r)|\delta A|_{(-1)},
 \qquad \vartheta_c\in\{1,L\}.
\]
Indeed $E(c,b)$ is supported in the double block $\mathsf D_c$,
its row sum is at most
$\kappa_FL^2\alpha_0\xi_j$, and every bond there
has local scale at least $\xi_j/\vartheta_c$.
The remaining term is the derivative of the quadratic remainder
$C_j$. Since $\mathcal C(0)=0$, integration along line segments proves
the first three inequalities.

For the transpose estimates use the weighted column norm
\[
 \|B\|_{\mathrm{col}}
  =\sup_b\mathsf L(b)^3\sum_c\xi_c^{-1}|B(c,b)|.
\]
At a bond of local scale $\xi_n$, only rows of levels
$n-1,n,n+1$ occur. The straightened kernel satisfies
$\sum_{c\in\Gamma_j^{\flat,+}}|K_A(c,b)|\le2\eta^3\xi_j^{-2}$,
also on the complex ball $|A|_{(-1)}<c_{\mathsf m}/L$. Hence
\[
 \|T^{-1}K_A\|_{\mathrm{col}}
       \le2(L^3+2)\eta^3.
\]
Apply the Cauchy derivative estimate on the ball of radius
$c_{\mathsf m}/(2L)$ and integrate on the ball of radius $c_{\mathsf m}/(4L)$.
With the stated, larger constant $C_3$ this gives
\[
 \|\mathcal C'(A_1)-\mathcal C'(A_2)\|_{\mathrm{col}}
 \le C_3L(L^3+2)\eta^3|A_1-A_2|_{(-1)}.
\]
The column estimate in the mean/corner comparison gives
$\|E\|_{\mathrm{col}}\le
\kappa_FL^2\alpha_0(L^3+2)\eta^3$.
Thus
$\|\mathcal C'(A)\|_{\mathrm{col}}\le3L^3\epsilon^*\eta^3$.
Finally
$|B^*\omega|_{(-3)}\le\eta^{-3}\|B\|_{\mathrm{col}}\sup|\omega|$,
and $L(L^3+2)\le2L^4$.
\end{proof}

The straightening leaves an endpoint factor at an outward
bond. Put $W_{0,j}=\mathcal M^jU_0$ and
$V_0=M_{\mathbf B^\flat}U_0$.
For an outward level-$j$ bond $c$ with outer endpoint $z$,
let $v_z(q)=\mathsf v(W_{0,j-1},q|_{B(z)})(z)$. The restriction
uses only interior rows of level $j-1$. The one-step
straightening estimate gives
$\theta(v_z(q))\le\kappa_{6,\mathsf m}\sup_{B(z)}|q|$.
Define $\Xi(q)=q$ on interior rows and
\[
 \Xi(q)(c)=
 \begin{cases}
  i^{-1}\log[v_z(q)e^{iq(c)}],&z=c_-,\\
  i^{-1}\log[e^{iq(c)}W_{0,j}(c)v_z(q)^{-1}W_{0,j}(c)^{-1}],&z=c_+.
 \end{cases}
\]
This is real-analytic when the interior values are less than
$r_{\mathsf m}$ and $|q(c)|+\kappa_{6,\mathsf m}\sup_{\mathrm{interior}}|q|<\pi$.
Define the gauge-removed row logarithm by
$\mathcal Q^\flat=\Xi\circ\mathcal Q$ and put
$K_A^\flat=(\mathcal Q^\flat)'(A)$. The straightening recursion
gives
\[
 M_{\mathbf B^\flat}(e^{i\eta A}U_0)
       =(e^{i\mathcal Q^\flat(A)}V_0)^{v[A]},\qquad
 v[A]|_{\Lambda_j}=v_j[A].
\]
The bounds $|a_j(c)|\le2\vartheta_c|A|_{(-1)}$ place
$\mathcal Q(A)$ in the indicated open set under the following
hypotheses.


\begin{lemma}\label{I-48-90}
If $|A|_{(-1)}\le r\le c_{\mathsf m}/(4L)$ and
$2(L+\kappa_{6,\mathsf m})r\le1$, then $K_A^\flat$ and $K_A$ have the
same kernel, and their adjoints have the same range. Also
$\mathcal Q^\flat(A)=0$ if and only if $\mathcal Q(A)=0$.
\end{lemma}

\begin{proof}
If $d\Xi(q)\delta q=0$, its interior components give
$\delta q=0$ there, so all $v_z$ have zero variation.
At each outward row the remaining derivative is that of
$X\mapsto i^{-1}\log(ve^{iX})$ or
$X\mapsto i^{-1}\log(e^{iX}v')$, with fixed $v,v'\in G$.
These maps are local diffeomorphisms on the specified
logarithm branch. Thus $\delta q=0$ also on outward rows.
Finite dimension makes $d\Xi(q)$ bijective. The chain rule
now gives the kernel and transpose-range assertions.
Finally $v_z(0)=I$. If $\Xi(q)=0$, its interior components
vanish, hence $v_z(q)=I$, and its outward components then
give $q=0$. This proves the last assertion.
\end{proof}

Let $\delta_\lambda A$ be the derivative at $t=0$ of the
bare gauge change $A^{e^{it\lambda}}$, and set
\[
 \varrho_A(\lambda)=\frac1i\left.\frac d{dt}
       \mathcal R_{\mathbf B}(U_0,A;e^{it\lambda})\right|_{t=0}
\]
at the block tops. The scalar operators $G',H'$ satisfy
\[
 (R\Delta,Q')^{-1}(f,\nu)=G'f+H'\nu
 \quad(f\in\operatorname{ran}R),
\]
with $\|G'f\|_Y\le(1+K_1)|f|_{(-2)}$ and
$\|H'\nu\|_Y\le K_2\sup|\nu|$. Recall that
$R\Delta=RD^*D$ and $\|\lambda\|_{Y^\circ}\le7\|\lambda\|_Y$.

\begin{lemma}\label{I-46-46}\label{I-46-49}\label{I-46-50}
After decreasing $c_6$ in \cref{I-46-44}, when
$A$ is real and $\alpha_0,\|A\|_X\le c_6$, the map
\[
 \lambda\longmapsto
 (RD^*\delta_\lambda A,\varrho_A(\lambda)|_{\Lambda_*})
\]
is a linear bijection from the real site functions onto
$\operatorname{ran}R$ times the real functions on the tops.
In particular,
\[
 \{RD^*\delta_\lambda A:\varrho_A(\lambda)|_{\Lambda_*}=0\}
          =\operatorname{ran}R.
\]
\end{lemma}

\begin{proof}
The divergence and induced-gauge estimates, differentiated
at the identity gauge, give
\[
 \begin{aligned}
 D^*\delta_\lambda A&=-D^*D\lambda+\dot N_A\lambda,
 &|\dot N_A\lambda|_{(-2)}
       &\le C_{\rm gc}\|A\|_X\|\lambda\|_{Y^\circ},\\
 T^{-1}\varrho_A(\lambda)&=Q'\lambda+\dot q_A\lambda,
 &\sup|\dot q_A\lambda|
       &\le(\kappa_F\alpha_0
                 +C_{\rm top}|A|_{(-1)})\|\lambda\|_1.
 \end{aligned}
\]
Parametrize the bare functions by
$\lambda=-G'f+H'\nu$. Put $K'=1+K_1+K_2$ and give
$(f,\nu)$ the norm $\max\{|f|_{(-2)},\sup|\nu|\}$.
Then $\|\lambda\|_Y\le K'\|(f,\nu)\|$, and the displayed
map, after rotation of its second component, is
\[
 (f,\nu)\longmapsto
 (f+R\dot N_A\lambda,\nu+\dot q_A\lambda).
\]
Choose $c_6<\min\{L^{-1}c_{\rm top},\rho_2\}$ and
\[
 [7(1+K_1^3)C_{\rm gc}+
       \kappa_F+C_{\rm top}]K'c_6\le\tfrac12.
\]
The perturbation has norm at most $1/2$, so its Neumann
series converges. This proves bijectivity and both stated
consequences.
\end{proof}

\begin{lemma}\label{I-46-60}
Suppose the hypotheses of \cref{I-46-46,I-48-90} hold,
$3K_1L^3\epsilon^*(\alpha_0,r)<1$, and
$E_A=e^{i\eta A}U_0$ satisfies
\[
 |\partial E_A(p)-I|<L^{-2}a_{\mathsf m}
                         \eta^2\mathsf L(p)^{-2}
\]
at every inner plaquette. Then $E_A$ is critical for the
action on its actual reduced averaging fibre if and only if
\[
 F'(A)\in\operatorname{ran}K_A^*
             =\operatorname{ran}K_A^{\flat*}.
\]
Moreover $(F'(A),\delta_\lambda A)_\eta=0$ for every real
$\lambda$, and $K_A\delta_\lambda A=0$ whenever
$\varrho_A(\lambda)$ vanishes on the tops. Criticality is
preserved by a fixed real gauge transformation.
\end{lemma}

\begin{proof}
The operator
$H_0^*\widehat K_A^*=I+H_0^*\mathcal C'(A)^*$ is invertible by
\cref{I-45-14,I-45-28}. Hence $K_A^*$ is injective and
$K_A$ is onto. By \cref{I-48-90}, the physical row map
$x(A)=e^{i\mathcal Q^\flat(A)}V_0$ also has surjective
derivative and $\ker dx(A)=\ker K_A$.
By \cref{I-39-43}, the curvature condition places the actual
averages in their analytic domain. Their factorization is
\[
 n(A):=M_{\mathbf B^\flat}(E_A)=x(A)^{v[A]},
 \qquad v[A]|_{\Lambda_j}=v_j[U_0;A].
\]
For a variation $\delta A$ write
$\mu=i^{-1}(dv[A]\delta A)v[A]^{-1}$. Differentiation gives
\[
 dn(A)\delta A=(dx(A)\delta A)^{v[A]}
                         +\operatorname{rot}_{n(A)}\mu.
\]
For $z$ a coarse configuration, write
$\operatorname{rot}_z\mu=(d/dt)z^{e^{it\mu}}|_{t=0}$.
Gauge covariance gives
\[
 dn(A)\delta_\lambda A
       =\operatorname{rot}_{n(A)}(\lambda|_{\Lambda_*}),\qquad
 dx(A)\delta_\lambda A
       =\operatorname{rot}_{x(A)}\varrho_A(\lambda).
\]
The first relation shows that $dn(A)$ is onto: after using
surjectivity of $dx(A)$, remove the unwanted rotation by a
bare gauge direction equal to $\mu$ on the tops.

Let $\mathfrak G$ be the space of bare gauge directions.
If $K_A\delta A=0$, the preceding formula for $dn$ and a
bare extension of $\mu$ give
$\delta A\in\ker dn(A)+\mathfrak G$.
Conversely, if $dn(A)\delta A=0$, then
\[
 dx(A)\delta A=-\operatorname{rot}_{x(A)}
                                  (\operatorname{Ad}(v[A]^{-1})\mu).
\]
By \cref{I-46-46}, choose $\lambda$ whose induced top values
are $\operatorname{Ad}(v[A]^{-1})\mu$.
Then $\delta A+\delta_\lambda A\in\ker dx(A)=\ker K_A$.
Consequently
$\ker dn(A)+\mathfrak G=\ker K_A+\mathfrak G$.

Gauge invariance of the action implies
$F'(A)\perp\mathfrak G$. The finite-dimensional submersion
theorem realizes every vector of $\ker dn(A)$ as the velocity
of a curve in the actual fibre. Thus criticality is equivalent
to $F'(A)\perp\ker dn(A)$, and hence to
$F'(A)\perp\ker K_A$. This is the stated multiplier condition.
The second gauge-covariance identity also proves the claim
when the induced top values vanish. Finally a fixed gauge
maps each actual averaging fibre diffeomorphically onto the
corresponding transformed fibre and preserves the action;
it therefore preserves criticality.
\end{proof}


\subsection{Solving the inhomogeneous equation}

The reference need not be critical. Keep the regional action
$F$ of \cref{I-40-63} and write $J_0=F'(0)$,
$j=|J_0|_{(-3)}$. The action remainder in this proof is
$V(A)=\mathcal V(A)$. Its local and weighted estimates are
those of \cref{I-40-63}; in particular the real plaquette
bound has radius $\alpha_0+3L^2r$ if $\|A\|_X\le r$,
$Lr\le\rho_{\mathrm{ec}}$ and $C_{\mathrm{ec}}L^2r\le1$.

\begin{proposition}\label{I-48-91}
There are $c_8,r_0>0$, depending only on $L$, with the
following property. Assume the preceding separated regional
setting with real $U_0$ and $\alpha_0\le c_8$. Put
\[
 K_8=K_1(16+9K_1^3),\qquad
 \Theta=4K_8,\qquad K_5=2(K_7+K_1).
\]
For a real row field $B^\natural$, set
$\beta=\sup|B^\natural|$ and $r_1=K_5(j+\beta)$.
If $r_1\le r_0$, the system
\[
 F'(A)+DRD^*A=K_A^*\omega,\qquad
 \mathcal Q(A)=B^\natural
\]
has exactly one real solution with $\|A\|_X\le r_0$.
Its multiplier is unique without a separate smallness
assumption, and
\[
 \begin{gathered}
 \|A\|_X\le r_1,\qquad \sup|\omega|\le\Theta r_1,\qquad
 RD^*A=0,\qquad F'(A)=K_A^*\omega,\\
 |D^*DA|_{(-3)}\le3(1+L^3)\Theta r_1.
 \end{gathered}
\]
The configuration $E=e^{i\eta A}U_0$ is critical on its
actual reduced averaging fibre.

In particular, suppose $B$ is prescribed on the reduced
rows and vanishes below the top level. Take $B^\natural=B$.
For any real gauge $h$ satisfying
$h=v_j[U_0;A]^{-1}$ on every $\Lambda_j$, the configuration
$E^h$ is critical on
\[
 C(\mathbf B^\flat,e^{iB}M_{\mathbf B^\flat}U_0).
\]
The values of $h$ away from these tops may be chosen arbitrarily.
\end{proposition}

\begin{proof}
Write $\|\cdot\|=\|\cdot\|_X$ and suppress the subscript on
$H_0$. We choose $c_8\le c_6$ and $r_0\le c_6$ small enough
that the action estimates, the row estimates and
\cref{I-48-90} apply on $\|A\|\le r_0$; also choose them so
that $C_{\mathrm{ec}}L^3c_8\le1$ and
\[
 c_8+3L^2r_0<L^{-2}a_{\mathsf m}.
\]
All further smallness requirements below concern constants
depending only on $L$ and quantities that tend to zero with
$(c_8,r_0)$.

Put $\Delta_R=F''(0)+DRD^*$. The exact regional Hessian
identity gives $G_0^{-1}=\Delta_R+Q^*Q$. The definitions of
$H$ and $\widetilde G_0$, and symmetry, imply
\[
 QH=I,\quad Q\widetilde G_0=0,\quad
 \widetilde G_0Q^*=0,\quad
 \widetilde G_0\Delta_R=I-HQ,\quad
 \ker\widetilde G_0=\operatorname{ran}Q^*.
\]
For the last identity, $\widetilde G_0 f=0$ implies
$G_0f=G_0Q^*(QG_0Q^*)^{-1}QG_0f$, and $G_0$ is invertible.

\emph{Eliminating the multiplier.}
Let $\widehat B=T^{-1}B^\natural$ and put
\[
 \mathcal G(A)=F'(A)+DRD^*A,\qquad
 S_A=I+H^*\mathcal C'(A)^*,\qquad
 w(A)=S_A^{-1}H^*\mathcal G(A).
\]
Decrease $c_8,r_0$ so that, with
$\epsilon^*=\epsilon^*(\alpha_0,r_0)$,
$3K_1L^3\epsilon^*\le1/8$. Then
$\|S_A^{-1}\|_{\infty\to\infty}\le8/7$ throughout the ball.
The field equation becomes
\[
 \Psi(A)=\widetilde G_0[-J_0-V(A)+\mathcal C'(A)^*w(A)]
                           +H[\widehat B-\mathcal C(A)].
\]
More precisely, the original system is equivalent on this
ball to $A=\Psi(A)$ and $\omega=Tw(A)$. To see both
directions, first apply $\widetilde G_0$ to the field
equation and use the constraint; applying $H^*$ gives
$S_A\widehat\omega=H^*\mathcal G(A)$.
Conversely, a fixed point gives
$QA=\widehat B-\mathcal C(A)$. The displayed identities then show
that
$\mathcal G(A)-(Q^*+\mathcal C'(A)^*)w(A)$ belongs to
$\operatorname{ran}Q^*$. Its image under $H^*$ is zero
by the definition of $w(A)$; since $H^*Q^*=I$, it is zero.

The transpose estimates are needed here in addition to the
field estimates. Since $|DA|_{(-2)}\le2\|A\|$ and
$|D^*A|_{(-2)}\le3\|A\|$,
\[
 H^*\mathcal G(A)=H^*J_0+(DH)^*DA
              +H^*(\Delta'A+V(A))+(D^*H)^*RD^*A
\]
gives
\[
 \begin{aligned}
 \sup|H^*\mathcal G(A)|&\le K_8(j+\|A\|+|V(A)|_{(-3)}),\\
 \sup|H^*(\mathcal G(A_1)-\mathcal G(A_2))|
  &\le K_8(\|A_1-A_2\|+|V(A_1)-V(A_2)|_{(-3)}).
 \end{aligned}
\]

\emph{The contraction.}
Let $j\le r\le r_0$ and $\|A\|,\|A_i\|\le r$.
By further reducing $r_0$, the action estimates give
\[
 |V(A)|_{(-3)}\le r/8,\qquad
 |V(A_1)-V(A_2)|_{(-3)}\le\tfrac12\|A_1-A_2\|.
\]
It follows that
$\sup|w(A)|\le(17/7)K_8r\le3K_8r$.
Writing $w_i=w(A_i)$, subtraction gives
\[
 S_{A_1}(w_1-w_2)
  =H^*(\mathcal G(A_1)-\mathcal G(A_2))
     -H^*(\mathcal C'(A_1)^*-\mathcal C'(A_2)^*)w_2.
\]
Choose $6K_1L^4C_3r_0\le1$. The preceding estimates then give
\[
 \sup|w_1-w_2|\le\tfrac87
       [\tfrac32K_8+6K_1L^4C_3K_8r]\|A_1-A_2\|
       \le3K_8\|A_1-A_2\|.
\]
With $\epsilon=\epsilon(\alpha_0,r_0)$ and
$\epsilon^*=\epsilon^*(\alpha_0,r_0)$, it follows that
\[
 \begin{aligned}
 \|\Psi(A)\|&\le K_7j+K_1\beta
       +[K_7(C_{\mathrm{ec}}L^3r+9K_8L^3\epsilon^*)+K_1L\epsilon]r,\\
 \|\Psi(A_1)-\Psi(A_2)\|
 &\le\{K_7[4C_{\mathrm{ec}}L^4r+6K_8L^4C_3r+9K_8L^3\epsilon^*]
                         +K_1L\epsilon\}\|A_1-A_2\|.
 \end{aligned}
\]
Decrease $c_8,r_0$ once more so that the braces in the
second line, at $r=r_0$ and $\alpha_0=c_8$, are at most
$1/2$. The bracket in the first line is then also at most
$1/2$. Consequently
\[
 \|\Psi(A)\|\le\tfrac12K_5(j+\beta)+\tfrac12r,\qquad
 \|\Psi(A_1)-\Psi(A_2)\|\le\tfrac12\|A_1-A_2\|.
\]
Apply the contraction principle on the balls of radii
$r_0$ and $r_1=K_5(j+\beta)$. Their fixed points agree,
and the multiplier is determined by $w(A)$. If $r_1=0$,
then $J_0=B^\natural=0$ and $A=0$ is the fixed point
directly. In all cases
\[
 \|A\|\le r_1,\qquad
 \sup|\omega|\le\tfrac{17}{28}\Theta r_1.
\]

\emph{The gauge condition and criticality.}
The curvature estimate and the choice of $c_8,r_0$ verify
the hypotheses of \cref{I-46-60}. For every $\lambda$ with
$\varrho_A(\lambda)=0$ at the tops, pair the field equation
with $\delta_\lambda A$. Both its action term and its
multiplier term vanish. Since $R$ is an orthogonal projection,
\[
 0=(RD^*A,RD^*\delta_\lambda A)_\eta.
\]
By \cref{I-46-50} the second arguments fill
$\operatorname{ran}R$, which contains $RD^*A$. Thus
$RD^*A=0$, so $F'(A)=K_A^*\omega$ and \cref{I-46-60}
gives criticality on the actual averaging fibre.

For the second derivative estimate, use
\[
 D^*DA=(Q^*+\mathcal C'(A)^*)T^{-1}\omega-J_0-\Delta'A-V(A).
\]
Our choices give $3L^3\epsilon^*\le1/(8K_1)\le1/32$,
$|\Delta'A|_{(-3)}\le r_1$ and $|V(A)|_{(-3)}\le r_1/8$.
The reduced-row bound for $Q^*$ therefore yields
\[
 |D^*DA|_{(-3)}
 \le[2(L^3+2)+1/32]\tfrac{17}{28}\Theta r_1
                         +\tfrac{17}{8}r_1
 \le3(1+L^3)\Theta r_1.
\]

Finally suppose $B$ vanishes below the top level. On the
interior rows $a_j=B_j$. Each endpoint factor $v_z$ at an
outward level-$j$ row depends only on the interior rows of
level $j-1$, where $a_{j-1}=0$; hence $v_z=I$.
The factorization of the actual averages and
$h=v_j^{-1}$ at the level-$j$ tops then give
$M_{\mathbf B^\flat}(E^h)=e^{iB}M_{\mathbf B^\flat}U_0$.
The top sets are disjoint, so these prescriptions admit a
real gauge extension. Criticality follows from gauge
invariance, as proved in \cref{I-46-60}.
\end{proof}


\subsection{Regularity on the gauge slice}

Put
\[
 C_{\mathrm{abs}}=C_{\mathrm{ec}}(L^3+L)
          +3L(\kappa_F L^2+4C_{\mathsf m}L),
\]
and choose
\[
 0<c_7\le\min\left\{
 a_{\rm in},\frac{c_{\mathsf m}}{4L},
 \frac1{(2+C_{\mathrm{ec}})L},
 \frac1{C_{\mathrm{ec}}L^3},
 \frac2{K_1C_{\mathrm{abs}}},
 \frac{\rho_{\mathrm{ec}}}{L}\right\}.
\]
The constants are fixed before the sequence and the fields.

\begin{lemma}\label{I-47-4}
In the preceding regional setting, let
$U_0\in\mathfrak U'_k(\mathbf B,\alpha_0)$,
$0<\alpha_0\le c_7$, and let $A$ be real. Suppose
\[
 |A|_{(-1)}\le c_7,\qquad RD^*A=0,\qquad
 E=e^{i\eta A}U_0\in\mathfrak U'_k(\mathbf B,\alpha_0).
\]
If $|a_j(A)(c)|\le\beta$ at every positively oriented
interior or outward row $c\in\Gamma_j^{\flat,+}$, including level
zero, then
\[
 \|A\|_X\le2K_1(\alpha_0+4\beta),\qquad
 |D^*DA|_{(-3)}\le8K_1(\alpha_0+4\beta).
\]
\end{lemma}

\begin{proof}
Put $r=|A|_{(-1)}$. The choice $r\le c_{\mathsf m}/(4L)$ lies
strictly inside the mean chart. We use the actual corner
rows $Q$ restricted to the interior and outward bonds.
Their adjoint satisfies
\begin{equation}\label{eq:reduced-corner-adjoint}
 |Q^*\omega|_{(-3)}\le2(1+L^{-3})\sup|\omega|.
\end{equation}
To see this, extend $\omega$ by zero to inward rows. If
$\mathsf L(b)=\xi_i$, locality permits only levels
$i-1,i,i+1$ in the column at $b$. A double block of an
interior or outward row of level $i-1$ misses $\Omega_i$,
whereas an endpoint of $b$ lies in $\Omega_i$; hence that
level contributes zero. At each remaining level $j$ the
column bound and the pairings give the factor
$\eta^{-3}\xi_j^{-1}(2\eta^3\xi_j^{-2})=2\xi_j^{-3}$.
Multiplication by $\xi_i^3$ proves
\eqref{eq:reduced-corner-adjoint}, including the endpoint
levels by omitting absent indices.

Write $J'=D^*DA$. At a full bond the local action
decomposition, the bounds on the currents of $U_0$ and $E$,
and the chart-gradient estimate give
\[
 \mathsf L(b)^3|J'(b)|
 \le[2+(2+C_{\mathrm{ec}})Lr]\alpha_0
          +C_{\mathrm{ec}}L^3r\|A\|_X
 \le3\alpha_0+C_{\mathrm{ec}}L^3r\|A\|_X.
\]
Here $s_b(A)\le Lr\le\rho_{\mathrm{ec}}$ and
$x_b(A)\le L^2\|A\|_X$. At a nonfull inner bond,
the bond and every bond of its inner plaquettes are
level-zero interior rows. On these bonds
$\eta|A|=|a_0(A)|\le\beta$. There are at most three
inner plaquettes in the divergence, with four values in
each curl. Since the real transports preserve norm,
\[
 \eta^3|J'(b)|\le3\cdot4\beta=12\beta.
\]
The expression is zero if there is no inner plaquette.
Consequently
\begin{equation}\label{eq:slice-second-order-bound}
 |J'|_{(-3)}\le
 \max\{3\alpha_0+C_{\mathrm{ec}}L^3r\|A\|_X,12\beta\}.
\end{equation}

For a row $c\in\Gamma_j^{\flat,+}$ put
$\widehat a(c)=\operatorname{Ad}(t_j(c_-))^{-1}a_j(A)(c)$.
The mean/corner comparison and the mean Taylor remainder
give
\[
 QA=\widehat a-\mathcal C(A),\qquad
 \sup|\widehat a|\le\beta,\qquad
 \sup|\mathcal C(A)|
   \le L(\kappa_F L^2\alpha_0+4C_{\mathsf m}Lr)r.
\]
Indeed the row of the linear difference applied to $A$ is
at most $\kappa_F L^2\alpha_0\vartheta_c r$, and the
mean remainder is at most $C_{\mathsf m}\vartheta_c^2r^2$;
$\vartheta_c\le L$ gives the displayed common bound.
The rotations preserve norms, and both sides are real.

The slice condition and $H_{U_0}=D^*D+\Delta'$ now give
\[
 (H_{U_0}+DRD^*+Q^*Q)A=f_A,\qquad
 f_A=J'+\Delta'A+Q^*(\widehat a-\mathcal C(A)).
\]
Use \eqref{eq:slice-second-order-bound},
$|\Delta'A|_{(-3)}\le C_{\mathrm{ec}}L\alpha_0r$ and
\eqref{eq:reduced-corner-adjoint}. Since
$\alpha_0,r\le c_7$ and $r\le\|A\|_X$,
\[
 |f_A|_{(-3)}
 \le3\alpha_0+[12+2(1+L^{-3})]\beta
            +C_{\mathrm{abs}}c_7\|A\|_X.
\]
Apply $G_0$ and its weighted estimate. The last term
contributes at most $\|A\|_X/2$, so it can be absorbed:
\[
 \|A\|_X\le\frac{K_1}{2}(3\alpha_0+16\beta)
             \le2K_1(\alpha_0+4\beta).
\]
Finally $C_{\mathrm{ec}}L^3r\le1$ in
\eqref{eq:slice-second-order-bound} gives the asserted
bound for $D^*DA$, using $K_1\ge1$.
\end{proof}


\section{Background inverses and current multipliers}
\label{sec:background-inverses}

The multiplier in the background operator is obtained from the
first constrained inverse. We first estimate this multiplier,
then apply \cref{I-65-83} again with that value. The weighted
transpose introduces a quotient of local scales; exponential
decay absorbs it uniformly in the number of layers.

Keep the geometry, natural regular-pair class, spaces and
bilinear pairings of \cref{sec:regular-system-inverse}, with
all the straightened rows of the determining set. In particular,
$\mathsf A=Y=\Omega_0$, and write
$\Delta(\mathbf U)=\Delta_{\Omega_0}(\mathbf U)$ for the
Hessian in \eqref{eq:small-system-quadratic-operators}.
It sums over the inner plaquettes of $\Omega_0$.
All operators below are evaluated at $\mathbf U$.
Whenever the indicated inverses exist, set
\begin{equation}\label{eq:bg-projection-operators}
 \begin{gathered}
 \mathsf G'=(\Delta'_a)^{-1},\qquad
 \mathsf C'=(Q'\mathsf G'^2Q'^{\mathsf T})^{-1},\qquad
 \mathsf P=\mathsf G'Q'^{\mathsf T}\mathsf C'Q'\mathsf G',
 \qquad \mathsf R=I-\mathsf P,\\
 \mathsf V=\mathsf G'\mathsf RD^{\mathsf T},\qquad
 \mathsf S=I-D\mathsf V.
 \end{gathered}
\end{equation}
The current form and its symmetric representing operator are
\begin{equation}\label{eq:bg-current-form}
 \begin{split}
 \mathfrak q_{\mathbf J}(A,v)
   &=(i[v_-,A]-i[A,v_+]-i[v_-,Dv],\mathbf J)_\xi,\\
 (A,\Delta'_\pi[\mathbf J]A)_\xi
   &=\mathfrak q_{\mathbf J}(A,\mathsf VA).
 \end{split}
\end{equation}
Here $v_\pm$ are as in \eqref{I-63-63}. The second line is a
homogeneous quadratic polynomial in $A$, so polarization with
the nondegenerate pairing defines $\Delta'_\pi[\mathbf J]$.
For an independent row multiplier $\mathsf j$ put
\begin{equation}\label{eq:bg-reduced-operator}
 \mathcal L_{\mathsf j}
  =\Delta(\mathbf U)-\Delta'_\pi[\mathbf J]
     -\mathsf S^{\mathsf T}\Delta^{(2)}[\mathsf j]\mathsf S
     +D\mathsf RD^{\mathsf T}+Q^{\mathsf T}\mathsf m^{-2}Q.
\end{equation}
The row weight $\mathsf m^{-2}(c)=\xi_c^{-2}$ is unchanged.
Define successively
\begin{equation}\label{eq:bg-physical-inverses}
 \begin{gathered}
 \mathsf G=\mathcal L_0^{-1},\qquad
 \mathsf H=\mathsf GQ^{\mathsf T}
                   (Q\mathsf GQ^{\mathsf T})^{-1},
 \qquad \mathsf j_* =\mathsf H^{\mathsf T}\mathbf J,\\
 \mathsf G_1=\mathcal L_{\mathsf j_*}^{-1},\qquad
 \mathsf H_1=\mathsf G_1Q^{\mathsf T}
                   (Q\mathsf G_1Q^{\mathsf T})^{-1},
 \qquad \mathsf j_{1*}=\mathsf H_1^{\mathsf T}\mathbf J.
 \end{gathered}
\end{equation}
The proposition below establishes existence in this order.
In particular, the second inverse is taken only after
$\mathsf j_*$ has been bounded.

\begin{proposition}\label{I-84-25}
For every odd $L\ge13$ there are positive constants
$a_{\mathrm{bg}},C_{\mathrm{bg}}$ and a power
$\mathfrak M_{\mathrm{bg}}$ of $L$ with the following property.
Let $0\le t\le K$, let $\mathfrak M\ge\mathfrak M_{\mathrm{bg}}$
and $R_1\ge L$ be powers of $L$, and let $(\Omega_n)_{n\le t}$
be a separated sequence with these parameters on a torus of
side $\ell\ge2L\mathfrak M\xi_t$.
Suppose
\[
 (\mathbf U,\mathbf J)\in
 \mathfrak P_t^\natural((\Omega_n);\alpha_0,\alpha_1,\gamma_0),
 \qquad 0<\alpha_0,\alpha_1,\gamma_0\le a_{\mathrm{bg}}.
\]
Then all inverses in
\eqref{eq:bg-projection-operators}--\eqref{eq:bg-physical-inverses}
exist. With
$\mathbb E^{[i]}(\mathsf j)
  =\mathbb K^{[i]}[\mathbf U,\mathbf J,\mathsf j]^{-1}$,
one has
\begin{equation}\label{eq:bg-inverse-blocks}
 \mathsf H=[\mathbb E^{[5]}(0)]_{A,8},\qquad
 \mathsf H_1=[\mathbb E^{[5]}(\mathsf j_*)]_{A,8},
\end{equation}
and, for every top $y$,
\begin{equation}\label{eq:bg-total-multipliers}
 |\mathsf j_*|_{-3;y},\ |\mathsf j_{1*}|_{-3;y}
       \le C_{\mathrm{bg}}\gamma_0\le a_{\mathrm{reg}}.
\end{equation}
All four systems at $\mathsf j_*$ are invertible and satisfy
\eqref{eq:regular-system-decay} with the same
$\delta_{\mathrm{reg}},B_{\mathrm{reg}}$.
The quadratic correction in $\mathsf G_1$ obeys
\[
 (A,\Delta^{(2)}[\mathsf j_*]A)_\xi
      =2(\mathsf HC^{(2)}(A),\mathbf J)_\xi,
 \qquad
 \mathsf G_1^{-1}
      =\mathsf G^{-1}
          -\mathsf S^{\mathsf T}\Delta^{(2)}[\mathsf j_*]\mathsf S.
\]
If the class decomposition has $A'=0$ and $\mathbf J$ is real,
then $\mathsf H,\mathsf H_1$ and both multipliers are real.
All constants are independent of the lattice spacing, volume
and number of layers.
\end{proposition}

\subsection{Elimination of the auxiliary variables}

Write $\mathsf g(A,v)=\mathsf g_0(A,v)+D^{\mathsf T}\mathsf x(A,v)$,
with the terms in \eqref{I-63-63}.
The scalar substitutions in \cref{sec:trivial-system-elimination}
also apply to the present operators. The current terms enter when
the resulting expressions are substituted into the first equation.

\begin{lemma}\label{I-63-67}
Let $(\mathbf U,\mathbf J,\mathsf j)$ satisfy the hypotheses of
\cref{I-65-83}, with ambient region $\mathsf A=Y=\Omega_0$.
The scalar inverses $\mathsf G'$ and $\mathsf C'$ in
\eqref{eq:bg-projection-operators} exist. The operators
$\mathcal L_{\mathsf j}$ and
$Q\mathcal L_{\mathsf j}^{-1}Q^{\mathsf T}$ are invertible, and
\[
 [\mathbb E^{[4]}]_{A,1}=\mathcal L_{\mathsf j}^{-1},\qquad
 [\mathbb E^{[5]}]_{A,8}
 =\mathcal L_{\mathsf j}^{-1}Q^{\mathsf T}
       (Q\mathcal L_{\mathsf j}^{-1}Q^{\mathsf T})^{-1}.
\]
Here the inverses $\mathbb E^{[i]}$ are taken at this triple.
\end{lemma}

\begin{proof}
The operator $\Delta'_a=\mathbb K^{[1]}$ is invertible by
\cref{I-65-83}. With $\mathsf G'=(\Delta'_a)^{-1}$, the scalar elimination
preceding \eqref{I-63-68}, with $k_2=f_2$, gives
\[
 Q'\mathsf G'^2Q'^{\mathsf T}\sigma
       =Q'\mathsf G'f_2-Q'\mathsf G'^2f_3-f_4.
\]
The same substitutions recover $v,w$ uniquely from $\sigma$,
so solutions correspond bijectively to solutions of this equation.
Its right-hand side ranges over the whole scalar coarse space as
$f_4$ varies. Invertibility of $\mathbb K^{[2]}$ therefore proves
invertibility of $Q'\mathsf G'^2Q'^{\mathsf T}$ and defines
$\mathsf C'$. The substitutions \eqref{I-63-68}--\eqref{I-63-69}
now hold with these operators and $\mathsf C'$ in place of $\mathsf C$.

The identity
$\langle\operatorname{Ad}(V)X,Y\rangle
 =\langle X,\operatorname{Ad}(V)^{-1}Y\rangle$ for $V\in G^c$
and reversal of each neighboring pair in \eqref{I-63-28} show that
$\Delta^\circ$ is symmetric. The diagonal terms, including those
at the boundary, are symmetric as well. Thus $\Delta'_a$,
$\mathsf G'$, $\mathsf C'$, $\mathsf P$ and $\mathsf R$ are symmetric,
and
\[
 \mathsf V^{\mathsf T}=D\mathsf R\mathsf G'.
\]

Put $\Delta^{(2)}=\Delta^{(2)}[\mathsf j]$ for this proof.
Cyclicity of the trace rewrites \eqref{eq:bg-current-form} as
\[
 \mathfrak q_{\mathbf J}(A,v)
  =(v_-+v_+,i[A,\mathbf J])_\xi
        -(v_-,i[Dv,\mathbf J])_\xi.
\]
Its half derivative in the $A$ direction $\dot A$ is
$(\dot A,\mathsf m_J(v))_\xi$.
For its half derivative in $v$, move the endpoint factors to the site pairing.
The derivatives of the endpoint factors give
$(\dot v,\mathsf g_0(A,v))_\xi$, while differentiating $Dv$
gives $(D\dot v,-\tfrac i2[\mathbf J,v_-])_\xi$.
Adding the quadratic form of $\Delta^{(2)}$ consequently gives
\begin{equation}\label{eq:bg-current-derivatives}
\begin{split}
 &\frac12\,d_{(A,v)}
   \bigl[\mathfrak q_{\mathbf J}(A,v)
       +(A-Dv,\Delta^{(2)}(A-Dv))_\xi\bigr](\dot A,\dot v)\\
 &\qquad=(\dot A,\mathsf m_J(v)
                       +\Delta^{(2)}(A-Dv))_\xi
                    +(\dot v,\mathsf g(A,v))_\xi.
\end{split}
\end{equation}
In the last term we used
$(D\dot v,\mathsf x(A,v))_\xi
  =(\dot v,D^{\mathsf T}\mathsf x(A,v))_\xi$.

Let
$\mathsf W_{\mathsf j}
 =\Delta'_\pi[\mathbf J]+\mathsf S^{\mathsf T}\Delta^{(2)}\mathsf S$.
Its quadratic form is the form in \eqref{eq:bg-current-derivatives}
evaluated at $(A,\mathsf VA)$.
Since $\mathsf SA=A-D\mathsf VA$, the chain rule and
$\mathsf V^{\mathsf T}=D\mathsf R\mathsf G'$ give
\begin{equation}\label{eq:bg-schur-correction}
 \mathsf W_{\mathsf j}A
  =\mathsf m_J(\mathsf VA)+\Delta^{(2)}\mathsf SA
                +D\mathsf R\mathsf G'\mathsf g(A,\mathsf VA).
\end{equation}
Indeed, pairing the right-hand side with $\dot A$ gives half the
derivative of this quadratic form. Its symmetric representing
operator is unique.

Take a right-hand side with only its first and, for $i=5$, eighth
components possibly nonzero. The scalar substitutions give
\[
 v=\mathsf VA,\qquad w=\mathsf PD^{\mathsf T}A,\qquad
 p=\mathsf R\mathsf G'\mathsf g(A,\mathsf VA).
\]
They also determine
$\sigma=\mathsf C'Q'\mathsf G'D^{\mathsf T}A$,
$r=\mathsf C'Q'\mathsf G'^2\mathsf g(A,\mathsf VA)$ and
$q=-\mathsf G'p$ uniquely.
Substitute these expressions in the first equation and use
\eqref{eq:bg-schur-correction}. Since $\mathsf R=I-\mathsf P$,
that equation becomes
\[
 \mathcal L_{\mathsf j}A=f_1+Q^{\mathsf T}\omega,
\]
with the last term omitted for $i=4$.
Here the Hessian is $\Delta(\mathbf U)$: both the ambient region
and $Y$ are $\Omega_0$, so no compression changes its plaquette sum.

For $i=4$, the equivalence just proved recovers every auxiliary
component from $A$. Invertibility of $\mathbb K^{[4]}$ therefore
gives a unique solution of $\mathcal L_{\mathsf j}A=f_1$ for each
$f_1$, proving the first block identity.
For $i=5$, the remaining equation $QA=f_8$ becomes
\[
 Q\mathcal L_{\mathsf j}^{-1}Q^{\mathsf T}\omega
       =f_8-Q\mathcal L_{\mathsf j}^{-1}f_1.
\]
Again all components are recovered uniquely from a solution.
Invertibility of $\mathbb K^{[5]}$, with $f_8$ arbitrary, proves
that the operator on the left is invertible. Taking $f_1=0$
gives the second block identity.
\end{proof}

\subsection{The multiplier bound}
\label{sec:background-multiplier}

Write $\delta=\delta_{\mathrm{reg}}$ and
$S_\delta=c_S(\delta/4)$, with $c_S$ from \eqref{I-13-121}.
Choose a power $\mathfrak M_{\mathrm{bg}}\ge\mathfrak M_*$
of $L$ so large that, for every $\mathfrak M\ge\mathfrak M_{\mathrm{bg}}$
and $R_1\ge L$,
\begin{equation}\label{eq:bg-thresholds}
 R_1\mathfrak M\ge\mathsf r_S(\delta/4),\qquad
 \frac{\log L}{R_1\mathfrak M}\le\frac\delta4.
\end{equation}
It suffices to impose these inequalities at $R_1=L$.
Set
\[
 C_{\mathrm{bg}}=\max\{1,6LB_{\mathrm{reg}}S_\delta\},
 \qquad a_{\mathrm{bg}}=a_{\mathrm{reg}}/C_{\mathrm{bg}}.
\]
These choices depend only on $L$.

\begin{proof}[Multiplier bound in \cref{I-84-25}]
Let $\mathsf j$ be any independent multiplier with
$|\mathsf j|_{-3;y}\le a_{\mathrm{reg}}$ for every top, and put
$T=[\mathbb E^{[5]}(\mathsf j)]_{A,8}$.
The inverse exists by \cref{I-65-83}, since
$a_{\mathrm{bg}}\le a_{\mathrm{reg}}$.
We estimate $T^{\mathsf T}\mathbf J$ once, then use the estimate
at $\mathsf j=0$ and at the induced multiplier.

Fix a positive row $c$ and $e\in\mathfrak g^c$.
Let $e_c$ equal $e$ at $c$ and zero at the other positive rows.
The weighted transpose gives
\begin{equation}\label{eq:bg-transpose-test}
 \begin{split}
 \xi_c^3\langle e,(T^{\mathsf T}\mathbf J)(c)\rangle
  &=(e_c,T^{\mathsf T}\mathbf J)_{\mathbf B}
    =(Te_c,\mathbf J)_\xi\\
  &=\xi^3\sum_b\langle(Te_c)(b),\mathbf J(b)\rangle,
 \end{split}
\end{equation}
where the sum is over positive inner bare bonds.
The source $g=(e_8e_c;0)$ satisfies $\iota g=e_8e_c$ and,
by \eqref{I-64-4},
\[
 |g|_{y(c)}=\xi_c|e|,\qquad |g|_y=0\quad(y\ne y(c)).
\]
Here $y(c)$ is the position of $c$ and $\xi_{y(c)}=\xi_c$.
The solution seminorm dominates $|A|_{-1;y}=\xi_y|A|_y$.
Thus \eqref{eq:regular-system-decay} gives
\[
 \xi_y|(Te_c)(b)|
 \le B_{\mathrm{reg}}e^{-\delta d(y,y(c))}\xi_c|e|
 \qquad\bigl(\operatorname{position}(b)=y\bigr).
\]

At most $3(\xi_y/\xi)^3$ positive bare bonds have position $y$,
and the regular-pair condition bounds the current there by
$\gamma_0\xi_y^{-3}$.
Their cell volume cancels the current weight in
\eqref{eq:bg-transpose-test}; the input and output weights leave
one quotient $\xi_c/\xi_y$. The upper pairing bound in
\eqref{eq:small-system-pairing-bounds} therefore yields
\[
 \xi_c^3|\langle e,(T^{\mathsf T}\mathbf J)(c)\rangle|
 \le3B_{\mathrm{reg}}\gamma_0|e|
       \sum_{y\in\Lambda_*}\frac{\xi_c}{\xi_y}
                          e^{-\delta d(y,y(c))}.
\]
The scale comparison \eqref{eq:global-scale-ratio} and
\eqref{eq:bg-thresholds} imply
\[
 \frac{\xi_c}{\xi_y}
 \le L\exp\!\left(\frac{\log L}{R_1\mathfrak M}d(y,y(c))\right)
 \le Le^{\delta d(y,y(c))/4}.
\]
Consequently the sum is at most
$L\sum_y e^{-\delta d(y,y(c))/4}\le LS_\delta$
by \eqref{I-13-122}.
The lower pairing bound in
\eqref{eq:small-system-pairing-bounds} contributes a factor of two
when recovering the matrix norm. Taking the maximum over rows
of position $y$, for which $\xi_c=\xi_y$, gives
\begin{equation}\label{eq:bg-multiplier-bound}
 |T^{\mathsf T}\mathbf J|_{-3;y}
 \le6LB_{\mathrm{reg}}S_\delta\gamma_0
 \le C_{\mathrm{bg}}\gamma_0.
\end{equation}

At $\mathsf j=0$, \cref{I-63-67} gives $\mathsf G,\mathsf H$
and identifies $T=\mathsf H$. Thus the first block identity holds
and $|\mathsf j_*|_{-3;y}\le C_{\mathrm{bg}}\gamma_0
\le a_{\mathrm{reg}}$.
\end{proof}

\subsection{The second inverse}
\label{sec:background-second-inverse}

\begin{proof}[Completion of the proof of \cref{I-84-25}]
We have $|\mathsf j_*|_{-3;y}\le a_{\mathrm{reg}}$ at every top.
Thus \cref{I-65-83} applies to the triple
$(\mathbf U,\mathbf J,\mathsf j_*)$. It gives all four inverses
and their claimed decay. Apply \cref{I-63-67} at this triple to
obtain $\mathsf G_1$, $(Q\mathsf G_1Q^{\mathsf T})^{-1}$ and
the second block identity in \eqref{eq:bg-inverse-blocks}.
The defining weighted transpose gives
\[
 2(C^{(2)}(A),\mathsf H^{\mathsf T}\mathbf J)_{\mathbf B}
       =2(\mathsf HC^{(2)}(A),\mathbf J)_\xi.
\]
Together with \eqref{eq:small-system-quadratic-operators},
this identifies the quadratic correction. Subtracting
\eqref{eq:bg-reduced-operator} at $\mathsf j=0$ from its value
at $\mathsf j_*$ gives the stated formula for $\mathsf G_1^{-1}$.

At this admissible multiplier, take $T=\mathsf H_1$ in
\eqref{eq:bg-multiplier-bound}. It gives
$|\mathsf j_{1*}|_{-3;y}\le C_{\mathrm{bg}}\gamma_0
\le a_{\mathrm{reg}}$, proving the second bound in
\eqref{eq:bg-total-multipliers}.

\emph{Real data.}
Suppose the class decomposition has $A'=0$ and $\mathbf J$ is real.
On a regular piece, the fixed unitary gauge in the proof of
\cref{lem:regular-local-inverse} writes $U^u=e^{i\xi A_0}$ with
$A_0$ real. The covering argument there places every global row
double block and scalar block in a regular or boundary piece,
and identifies the actual relative-constraint germ before
returning to the original gauge.
In the rescaled block, the proof of
\cref{prop:flat-complex-averaging} gives actual bases
$W_n=e^{iB_n}$ with $B_n$ real and
$\|B_n\|_\infty<r_{\mathrm{st}}/16$.
For a sufficiently small real relative perturbation, its iterates
start real and satisfy
$\|a_n\|_\infty\le2c_{\mathrm{av}}<r_{\mathrm{st}}/2$.
The real assertion of \cref{lem:one-step-straightening} therefore
applies at each step and keeps $a_n$ real.
The first and second Taylor coefficients of the actual constraint
are consequently real. The fixed unitary covariance
\eqref{eq:straightening-covariance} preserves the real form and
returns these coefficients to the original gauge.
The agreement of germs in \cref{lem:regular-local-inverse}
identifies them with the global $Q$ and $C^{(2)}$.
The scalar transport formula \eqref{eq:flat-complex-scalar-transport}
makes $Q'$ real as well. On boundary pieces
$Q=Q'=I$ and $C^{(2)}=0$.

For unitary $U$, its adjoint transports preserve $\mathfrak g$,
and $i[X,Y]\in\mathfrak g$ for $X,Y\in\mathfrak g$.
Thus the covariant differences and the current commutators are real.
The positive real pairing weights make the transpose of a real map
real. The action is real on real perturbations of $U$, so its
Hessian is real; likewise $\Delta^{(2)}[\mathsf j]$ is real for
real $\mathsf j$, by its quadratic form and polarization.
It follows from \eqref{I-63-63} and the system equations that
$\mathbb K^{[i]}$ is real at every such real triple.
Its inverse is real by uniqueness of solutions.
At $\mathsf j=0$, the block identity
\eqref{eq:bg-inverse-blocks} makes $\mathsf H$ real, and therefore
$\mathsf j_*=\mathsf H^{\mathsf T}\mathbf J$ is real.
Apply the same argument at this real multiplier to obtain real
$\mathsf H_1$ and
$\mathsf j_{1*}=\mathsf H_1^{\mathsf T}\mathbf J$.

\end{proof}

\section{The nonlinear response}\label{sec:nonlinear-response}

The nonlinear correction is obtained in two steps. First we
solve $X=\mathsf C(A'-\mathsf HX)$ and denote its solution by
$\mathsf D(A')$. Its derivative enters the reduced field
$\mathsf W$ defined below. The quadratic bound for this field
makes $\mathbf A=-\mathsf G\mathsf W(\mathsf KB'+\mathbf A)$ a
contraction; the response is then
$\mathcal L(\mathsf KB'+\mathbf A)$, where
$\mathcal L(A')=A'-\mathsf H\mathsf D(A')$.
We prove this implication under explicit bounds on the
nonlinear maps and the four linear operators.

Use the separated layers, cells $\Delta(y)$, reduced positive
rows $c$, and distance $d=d_{\mathbf B^\flat}$ of
\cref{sec:multiscale-sums}. A bare bond lies over the unique
$y\in\Lambda_*$ whose cell contains its initial site; write
$y_b=y$. Rows are placed at $y_c$ as in that section, and
$\xi_c$ is the spacing of the row. Let $\mathcal A$ and
$\mathcal X$ be the complex Lie-algebra-valued fields on
these bare bonds and rows, respectively. In this section the
row variables are dimensionless, and the bilinear pairings are
\begin{equation}\label{I-86-6}
 (A,f)_\xi=\xi^3\sum_b\langle A(b),f(b)\rangle,
 \qquad (X,Y)_1=\sum_c\xi_c^{-1}\langle X(c),Y(c)\rangle.
\end{equation}
Every transpose below uses these pairings. In particular, the
row weight here is $\xi_c^{-1}$, whereas the physical row
pairing in \cref{sec:background-inverses} has weight $\xi_c^3$.
The letters $Q,\mathsf H,\mathsf S,\mathsf K,\mathsf G$ below
denote maps satisfying the stated assumptions.

For $A,f\in\mathcal A$ and $X\in\mathcal X$ put
\[
 |A|_y=\xi_y\max_{y_b=y}|A(b)|,\qquad
 |f|_y^\sharp=\xi_y^3\max_{y_b=y}|f(b)|,\qquad
 |X|_y=\max_{y_c=y}|X(c)|.
\]
An empty maximum is zero. Let $\mathsf g_y$ be finite
seminorms on $\mathcal A$, and set
$\|A\|_y=\max\{|A|_y,\mathsf g_y(A)\}$.
Fix $\theta_0>0$, put $\rho_0=3L+2$, and define
\[
 N(y)=\{z:d(y,z)\le2\rho_0\},\qquad
 \mathsf e=\exp(2\rho_0\theta_0).
\]
Admissible weights are positive functions with
$w(y)\le e^{\theta_0d(y,z)}w(z)$ for all $y,z$.
Write $|A|_w=\sup_y w(y)|A|_y$, and likewise
$|f|_w^\sharp$, $|X|_w$, and $\|A\|_w$.
At $w=1$ write $\sup|A|$, $\sup|f|^\sharp$,
$\sup|X|$, and $\|A\|$. These are finite-dimensional
norms; the zero-dimensional cases are included. The weight
comparison we shall use is
\begin{equation}\label{eq:response-neighborhood-weight}
 w(y)\max_{z\in N(y)}h(z)
       \le\mathsf e\sup_z w(z)h(z)\qquad(h\ge0).
\end{equation}

\subsection{Hypotheses and the response theorem}

Fix $c_\flat,C_2,C_3,C_{\mathcal V}\ge1$ and
$c_4,\rho_1\in(0,1]$. Let $Q:\mathcal A\to\mathcal X$
be linear, and let
\[
 \mathsf C:\{\sup|A|<c_4\}\longrightarrow\mathcal X,
 \qquad
 \mathcal V:\{\sup|A|<2\rho_1\}\longrightarrow\mathcal A
\]
be holomorphic. Assume the following local estimates for
every $y$ and every argument in the specified domains.
All maxima in these hypotheses are over $z\in N(y)$.
\begin{enumerate}[label=(\roman*)]
\item The linear map and its transpose satisfy
\[
 |QA|_y\le c_\flat\max |A|_z,
 \qquad |Q^{\mathsf T}X|_y^\sharp
                 \le c_\flat\max |X|_z.
\]
\item $\mathsf C(0)=\mathsf C'_0=0$, and
\[
 |\mathsf C(A)|_y\le C_2\max |A|_z^2.
\]
If $\sup|A_i|\le r\le c_4/4$, then
\[
 |\mathsf C(A_1)-\mathsf C(A_2)|_y
       \le4C_2r\max |A_1-A_2|_z.
\]
For $\sup|A|,\sup|A_i|\le c_4/2$,
\begin{equation}\label{eq:response-average-transpose}
 \begin{split}
 |\mathsf C'_A{}^{\mathsf T}X|_y^\sharp
     &\le C_3\max |A|_z\max |X|_z,\\
 |(\mathsf C'_{A_1}{}^{\mathsf T}
       -\mathsf C'_{A_2}{}^{\mathsf T})X|_y^\sharp
     &\le C_3\max |A_1-A_2|_z\max |X|_z.
 \end{split}
\end{equation}
\item $\mathcal V(0)=\mathcal V'(0)=0$. For
$\sup|A|\le\rho_1$,
\[
 |\mathcal V(A)|_y^\sharp
       \le C_{\mathcal V}\max |A|_z\max\|A\|_z.
\]
If $\|A_i\|\le r\le\rho_1/4$, then
\[
 |\mathcal V(A_1)-\mathcal V(A_2)|_y^\sharp
       \le4C_{\mathcal V}r\max\|A_1-A_2\|_z.
\]
\end{enumerate}

Let $\mathsf a\ge0$ act on rows by multiplication, and fix
$\mathbf J\in\mathcal A$ with
$\gamma_0:=\sup|\mathbf J|^\sharp\le1$. Assume that the
linear maps
\[
 \mathsf H,\mathsf K:\mathcal X\to\mathcal A,
 \quad \mathsf S:\mathcal X\to\mathcal X,
 \quad \mathsf G:\mathcal A\to\mathcal A
\]
obey, for one finite $\hat{\mathsf b}\ge\max\{1,\sup\mathsf a\}$,
every admissible $w$, every $X\in\mathcal X$, and every
$f\in\mathcal A$,
\begin{equation}\label{I-86-50}
 \begin{gathered}
 \|\mathsf HX\|_w,\ \|\mathsf KX\|_w,\ |\mathsf SX|_w
                      \le\hat{\mathsf b}|X|_w,\qquad
 |(\mathsf S-\mathsf a)X|_w\le2\hat{\mathsf b}|X|_w,\\
 \|\mathsf Gf\|_w\le\hat{\mathsf b}|f|_w^\sharp,
 \qquad |\mathsf H^{\mathsf T}f|_w
                      \le8\hat{\mathsf b}|f|_w^\sharp.
 \end{gathered}
\end{equation}
Choose the radii by
\begin{equation}\label{I-86-46}
 \begin{gathered}
 c_W=\min\left\{\frac{c_4}{8},\frac1{32C_2\mathsf e},
              \frac1{32C_3\mathsf e},\frac{\rho_1}{8}\right\},\\
 C_W=\frac{72C_3\mathsf e}{c_W}
       +32c_\flat C_2\mathsf e^2+30c_\flat C_3\mathsf e^2
       +40C_{\mathcal V}\mathsf e,\\
 \varepsilon_3=\frac{c_W}{\hat{\mathsf b}},\qquad
 \varepsilon_B=\min\left\{\frac{c_W}{4\hat{\mathsf b}^2},
                         \frac1{8C_W\hat{\mathsf b}^4}\right\}.
 \end{gathered}
\end{equation}

\begin{proposition}\label{I-86-52}
Under the hypotheses above, for every $\sup|A'|\le\varepsilon_3$
there is a unique solution $X=\mathsf D(A')$ of
\begin{equation}\label{I-86-53}
 X=\mathsf C(A'-\mathsf HX),\qquad
 \sup|X|\le\varepsilon_3/\hat{\mathsf b}.
\end{equation}
It satisfies $\sup|\mathsf D(A')|\le4C_2\sup|A'|^2$ and is
holomorphic for $\sup|A'|<\varepsilon_3$.
Put $\mathcal L(A')=A'-\mathsf H\mathsf D(A')$.
For $\sup|A_i'|\le\varepsilon\le\varepsilon_3$,
\begin{equation}\label{eq:response-first-differences}
 \begin{gathered}
 |\mathsf D(A_1')-\mathsf D(A_2')|_w
       \le16C_2\mathsf e\varepsilon|A_1'-A_2'|_w,\\
 \tfrac12|A_1'-A_2'|_w
       \le|\mathcal L(A_1')-\mathcal L(A_2')|_w
       \le2|A_1'-A_2'|_w,\\
 \|\mathcal L(A_1')-\mathcal L(A_2')\|_w
       \le2\|A_1'-A_2'\|_w.
 \end{gathered}
\end{equation}
On the open ball, write $\mathfrak D(A')=\mathsf D'(A')$.
With $T=\mathcal L(A')$,
\begin{equation}\label{I-86-59}
 \mathfrak D(A')^{\mathsf T}
   =\mathsf C'_T{}^{\mathsf T}
          (I+\mathsf H^{\mathsf T}\mathsf C'_T{}^{\mathsf T})^{-1}.
\end{equation}
\end{proposition}

For $\sup|A'|<\varepsilon_3$, abbreviate
$\mathsf D=\mathsf D(A')$, $T=\mathcal L(A')$, and
$\mathfrak D=\mathfrak D(A')$, and define
\begin{equation}\label{I-86-62}
 \begin{split}
 \mathsf W(A')={}&
  -(\mathfrak D^{\mathsf T}-\mathsf C''[A']^{\mathsf T})
                                    \mathsf H^{\mathsf T}\mathbf J
  -Q^{\mathsf T}(\mathsf S-\mathsf a)\mathsf D\\
 &-\mathfrak D^{\mathsf T}(\mathsf S-\mathsf a)QA'
  +\mathfrak D^{\mathsf T}(\mathsf S-\mathsf a)\mathsf D
  +\mathcal V(T)
  -\mathfrak D^{\mathsf T}\mathsf H^{\mathsf T}\mathcal V(T).
 \end{split}
\end{equation}
Here $\mathsf C''[A']=\frac{d}{dt}\mathsf C'_{tA'}|_{t=0}$.
The subtraction in the first term removes the linear part of
$\mathfrak D^{\mathsf T}$; its quantitative effect is proved
below.

\begin{theorem}\label{I-87-35}
Under the hypotheses of this subsection, for each
$B'\in\mathcal X$ with $\sup|B'|\le\varepsilon_B$ the system
\begin{equation}\label{I-87-37}
 \begin{gathered}
 X=\mathsf C(\mathsf KB'+\mathbf A-\mathsf HX),\qquad
 \mathbf A+\mathsf G\mathsf W(\mathsf KB'+\mathbf A)=0,\\
 \mathbf H=\mathsf KB'+\mathbf A-\mathsf HX
 \end{gathered}
\end{equation}
has exactly one solution with
$\sup|X|\le\varepsilon_3/\hat{\mathsf b}$ and
$\|\mathbf A\|\le\hat{\mathsf b}\varepsilon_B$.
Write $\mathbf A_0(B')$ and $\mathbf H(B')$ for its second
and third components. They are holomorphic for
$\sup|B'|<\varepsilon_B$, vanish at zero, and satisfy
\begin{equation}\label{I-87-38}
 \|\mathbf A_0(B')\|\le8C_W\hat{\mathsf b}^5\sup|B'|^2
      \le\hat{\mathsf b}\sup|B'|,\qquad
 \|\mathbf H(B')\|\le4\hat{\mathsf b}\sup|B'|.
\end{equation}
For every admissible $w$ and $\sup|B_i'|\le\varepsilon_B$,
\begin{equation}\label{eq:response-weighted-bounds}
 \begin{gathered}
 \|\mathbf A_0(B_1')-\mathbf A_0(B_2')\|_w
       \le\hat{\mathsf b}|B_1'-B_2'|_w,\\
 \|\mathbf H(B_1')-\mathbf H(B_2')\|_w
       \le4\hat{\mathsf b}|B_1'-B_2'|_w,\\
 \|\mathbf H(B')-\mathsf KB'\|_w
       \le10C_W\hat{\mathsf b}^5\sup|B'|\,|B'|_w.
 \end{gathered}
\end{equation}
If a nonempty $\Sigma\subset\Lambda_*$ contains the support
of $B'$ in the row positions, then
\begin{equation}\label{eq:response-support-decay}
 \begin{split}
 \|\mathbf H(B')\|_y
    &\le4\hat{\mathsf b}e^{-\theta_0d(y,\Sigma)}\sup|B'|,\\
 \|\mathbf H(B')-\mathsf KB'\|_y
    &\le10C_W\hat{\mathsf b}^5e^{-\theta_0d(y,\Sigma)}\sup|B'|^2.
 \end{split}
\end{equation}
\end{theorem}

\subsection{The change of variables and its derivative}

\begin{proof}[Proof of \cref{I-86-52}]
Fix $\sup|A'|\le\varepsilon\le\varepsilon_3$.
On the ball $\sup|X|\le\varepsilon/\hat{\mathsf b}$,
\[
 \sup|A'-\mathsf HX|\le2\varepsilon\le c_4/4,
 \qquad
 \sup|\mathsf C(A'-\mathsf HX)|\le4C_2\varepsilon^2
                       \le\varepsilon/\hat{\mathsf b}.
\]
The Lipschitz constant of this map of $X$ is at most
$8C_2\hat{\mathsf b}\varepsilon\le1/4$. The contraction
theorem gives a unique fixed point in each of these nested
balls. Uniqueness in the largest ball identifies them, and
the image estimate at $\varepsilon=\sup|A'|$ gives the
quadratic bound. Iteration from zero has increments bounded
by $4^{-n}\varepsilon_3/\hat{\mathsf b}$, uniformly in
$A'$ on the open ball. The iterates are holomorphic and
remain in the domain of $\mathsf C$, so their limit is
holomorphic.

Put $X_i=\mathsf D(A_i')$ and $T_i=A_i'-\mathsf HX_i$.
The local difference estimate and
\eqref{eq:response-neighborhood-weight} give
\[
 |X_1-X_2|_w\le8C_2\mathsf e\varepsilon
       \bigl(|A_1'-A_2'|_w+\hat{\mathsf b}|X_1-X_2|_w\bigr).
\]
Absorb the last term, whose coefficient is at most $1/4$.
This proves the first bound in
\eqref{eq:response-first-differences}. Moreover,
\[
 \|\mathsf H(X_1-X_2)\|_w
       \le16C_2\mathsf e\hat{\mathsf b}\varepsilon
                       |A_1'-A_2'|_w
       \le\tfrac12|A_1'-A_2'|_w.
\]
The remaining bounds follow by subtracting this correction
from $A_1'-A_2'$.

For $\sup|T|\le2\varepsilon_3$ set
$M_T=\mathsf C'_T{}^{\mathsf T}$. By
\eqref{eq:response-average-transpose}, for
$\sup|T|\le2\varepsilon\le2\varepsilon_3$,
\begin{equation}\label{eq:response-M-bounds}
 \begin{gathered}
 |M_T\varphi|_w^\sharp\le2C_3\mathsf e\varepsilon|\varphi|_w,\\
 |(M_{T_1}-M_{T_2})\varphi|_w^\sharp
       \le C_3\mathsf e|T_1-T_2|_w\sup|\varphi|.
 \end{gathered}
\end{equation}
The second estimate holds for
$\sup|T_i|\le2\varepsilon_3$.
Consequently $\mathsf H^{\mathsf T}M_T$ has weighted
operator norm at most $16C_3\mathsf e\hat{\mathsf b}
\varepsilon_3\le1/2$. Its Neumann series gives
\[
 \Psi_T=(I+\mathsf H^{\mathsf T}M_T)^{-1},\qquad
 |\Psi_T\varphi|_w\le2|\varphi|_w.
\]
Differentiate $\mathsf D(A')=\mathsf C(A'-\mathsf H\mathsf D(A'))$
and transpose. This gives
$\mathfrak D^{\mathsf T}(I+\mathsf H^{\mathsf T}M_T)=M_T$,
which proves \eqref{I-86-59}.
\end{proof}

We record within the argument the derivative estimates needed
for \eqref{I-86-62}. If $\sup|A_i'|\le\varepsilon<\varepsilon_3$
and $\mathfrak D_i=\mathfrak D(A_i')$, then
\begin{equation}\label{eq:response-derivative-bounds}
 \begin{gathered}
 |\mathfrak D_i^{\mathsf T}\varphi|_w^\sharp
       \le4C_3\mathsf e\varepsilon|\varphi|_w,\\
 |(\mathfrak D_1^{\mathsf T}-\mathfrak D_2^{\mathsf T})
                      \varphi|_w^\sharp
       \le8C_3\mathsf e|A_1'-A_2'|_w\sup|\varphi|.
 \end{gathered}
\end{equation}
Indeed, the first follows from $\mathfrak D_i^{\mathsf T}=M_i\Psi_i$.
For the second use
\[
 \Psi_1-\Psi_2=-\Psi_2\mathsf H^{\mathsf T}(M_1-M_2)\Psi_1,
 \quad
 \mathfrak D_1^{\mathsf T}-\mathfrak D_2^{\mathsf T}
  =(I-M_2\Psi_2\mathsf H^{\mathsf T})(M_1-M_2)\Psi_1.
\]
Here $M_2\Psi_2\mathsf H^{\mathsf T}$ has weighted source
norm at most $32C_3\mathsf e c_W\le1$, while
\eqref{eq:response-M-bounds} and
$|T_1-T_2|_w\le2|A_1'-A_2'|_w$ bound the last two factors
by $4C_3\mathsf e|A_1'-A_2'|_w\sup|\varphi|$.
At zero, $M_0=0$, $\Psi_0=I$, and $\mathcal L'(0)=I$;
thus $\mathfrak D(0)=0$ and
\begin{equation}\label{eq:response-linear-cancellation}
 \left.\frac d{dt}\mathfrak D(tA')^{\mathsf T}\right|_{t=0}
                   =\mathsf C''[A']^{\mathsf T}.
\end{equation}

\subsection{The reduced field and the second contraction}

\begin{proof}[Proof of \cref{I-87-35}]
First we establish the estimate that makes the second
contraction possible. The field $\mathsf W$ is holomorphic
on $\sup|A'|<\varepsilon_3$ and vanishes at zero. Indeed,
all its factors are holomorphic there, and
$\sup|\mathcal L(A')|<2\varepsilon_3\le\rho_1/4$
keeps the argument in the domain of $\mathcal V$.
We shall show, for $\|A_i'\|\le\varepsilon<\varepsilon_3$,
\begin{equation}\label{I-87-29}
 |\mathsf W(A_1')-\mathsf W(A_2')|_w^\sharp
       \le C_W\hat{\mathsf b}(1+\gamma_0\hat{\mathsf b})
                    \varepsilon\|A_1'-A_2'\|_w.
\end{equation}
The case $\varepsilon=0$ is immediate. Otherwise write
$\delta=\|A_1'-A_2'\|_w$, $\delta_0=|A_1'-A_2'|_w\le\delta$,
$X_i=\mathsf D(A_i')$, and $T_i=\mathcal L(A_i')$.
The bounds already proved give
\[
 \sup|X_i|\le4C_2\varepsilon^2,\quad
 |X_1-X_2|_w\le16C_2\mathsf e\varepsilon\delta_0,
 \quad \|T_i\|\le2\varepsilon,
 \quad \|T_1-T_2\|_w\le2\delta.
\]

\emph{The cancellation in the current term.}
Put $\mathsf N(A')=\mathfrak D(A')^{\mathsf T}
-\mathsf C''[A']^{\mathsf T}$ and
$\varphi=\mathsf H^{\mathsf T}\mathbf J$.
Differentiating \eqref{eq:response-average-transpose} at zero
gives
\[
 |\mathsf C''[A']^{\mathsf T}\varphi|_w^\sharp
       \le C_3\mathsf e|A'|_w\sup|\varphi|.
\]
Together with \eqref{eq:response-derivative-bounds}, this
bounds the difference of $\mathsf N$ at any two points of
the open amplitude ball by
$9C_3\mathsf e|A_1'-A_2'|_w\sup|\varphi|$.
By \eqref{eq:response-linear-cancellation},
$\mathsf N(0)=\mathsf N'(0)=0$.
Let $R=\varepsilon_3/\varepsilon>1$ and set
\[
 g(z)=(\mathsf N(zA_1')-\mathsf N(zA_2'))\varphi
                       \qquad(|z|<R).
\]
Then $g(z)/z$ extends holomorphically across zero, vanishes
there, and has weighted source norm at most
$9C_3\mathsf e\delta_0\sup|\varphi|$.
Apply scalar Schwarz after each norm-one linear functional
and take the supremum. At $z=1$ this gains the factor
$R^{-1}=\hat{\mathsf b}\varepsilon/c_W$.
Since $\sup|\varphi|\le8\hat{\mathsf b}\gamma_0$, the
first term of \eqref{I-86-62} has difference bounded by
\begin{equation}\label{eq:response-current-difference}
 \frac{72C_3\mathsf e}{c_W}\gamma_0\hat{\mathsf b}^2
                              \varepsilon\delta_0.
\end{equation}

\emph{The remaining terms.}
For terms containing $\mathfrak D^{\mathsf T}$ use the
single product estimate
\begin{equation}\label{I-87-32}
 |\mathfrak D_1^{\mathsf T}\varphi_1
       -\mathfrak D_2^{\mathsf T}\varphi_2|_w^\sharp
 \le8C_3\mathsf e\delta_0\sup|\varphi_1|
       +4C_3\mathsf e\varepsilon|\varphi_1-\varphi_2|_w.
\end{equation}
The local bounds for $Q,Q^{\mathsf T}$ cost at most
$c_\flat\mathsf e$ in weighted norms and $c_\flat$ at
weight one. Applying them, \eqref{I-86-50}, and the
displayed bounds for $X_i,T_i$ gives the following estimates
for the differences of terms two through six in
\eqref{I-86-62}:
\[
\begin{array}{c|c}
 \text{term}&\text{bound in }|\cdot|_w^\sharp\\ \hline
 -Q^{\mathsf T}(\mathsf S-\mathsf a)X
    &32c_\flat C_2\mathsf e^2\hat{\mathsf b}\varepsilon\delta_0\\
 -\mathfrak D^{\mathsf T}(\mathsf S-\mathsf a)QA'
    &24c_\flat C_3\mathsf e^2\hat{\mathsf b}\varepsilon\delta_0\\
 \mathfrak D^{\mathsf T}(\mathsf S-\mathsf a)X
    &6C_3\mathsf e\varepsilon\delta_0\\
 \mathcal V(T)&16C_{\mathcal V}\mathsf e\varepsilon\delta\\
 -\mathfrak D^{\mathsf T}\mathsf H^{\mathsf T}\mathcal V(T)
    &24C_{\mathcal V}\mathsf e\varepsilon\delta
\end{array}
\]
For the fourth term, \eqref{I-87-32} first gives
$(64+128\mathsf e)C_2C_3\mathsf e\hat{\mathsf b}
\varepsilon^2\delta_0$; use
$32C_2\mathsf e\hat{\mathsf b}\varepsilon\le1$.
For the last two terms, $\|T_i\|\le2\varepsilon\le\rho_1/4$
permits the local Lipschitz bound for $\mathcal V$ and gives
$\sup|\mathcal V(T_i)|^\sharp\le4C_{\mathcal V}\varepsilon^2$.
The last application of \eqref{I-87-32} gives
$(256+512\mathsf e)C_3C_{\mathcal V}\mathsf e\hat{\mathsf b}
\varepsilon^2\delta$; use
$32C_3\mathsf e\hat{\mathsf b}\varepsilon\le1$.
Adding the table and \eqref{eq:response-current-difference},
with $c_\flat,\mathsf e,\hat{\mathsf b}\ge1$, proves
\eqref{I-87-29}. Taking $A_2'=0$ also yields
\begin{equation}\label{I-87-30}
 |\mathsf W(A')|_w^\sharp
       \le C_W\hat{\mathsf b}(1+\gamma_0\hat{\mathsf b})
                          \|A'\|\,\|A'\|_w
             \qquad(\|A'\|<\varepsilon_3).
\end{equation}

\emph{The fixed point.}
Let $\sup|B'|\le r\le\varepsilon_B$, and consider
$\Phi_{B'}(\mathbf A)=-\mathsf G\mathsf W(\mathsf KB'+\mathbf A)$
on $\|\mathbf A\|\le\hat{\mathsf b}r$.
Its argument has strong norm at most
$2\hat{\mathsf b}r\le\varepsilon_3/2$.
Since $\gamma_0\le1\le\hat{\mathsf b}$,
\eqref{I-87-29} and \eqref{I-87-30} imply
\begin{equation}\label{I-87-36}
 \begin{gathered}
 \|\Phi_{B'}(\mathbf A)\|
   \le2C_W\hat{\mathsf b}^3
                (\hat{\mathsf b}\sup|B'|+\|\mathbf A\|)^2
   \le8C_W\hat{\mathsf b}^5r^2\le\hat{\mathsf b}r,\\
 \|\Phi_{B'}(\mathbf A_1)-\Phi_{B'}(\mathbf A_2)\|_w
   \le4C_W\hat{\mathsf b}^4r\|\mathbf A_1-\mathbf A_2\|_w
   \le\tfrac12\|\mathbf A_1-\mathbf A_2\|_w.
 \end{gathered}
\end{equation}
The contraction theorem in the finite-dimensional strong norm
gives a unique fixed point. As before, the nested balls have
the same fixed point, denoted $\mathbf A_0(B')$.
Taking $r=\sup|B'|$ proves its bounds in \eqref{I-87-38}.
The first equation in \eqref{I-87-37} then uniquely determines
$X=\mathsf D(\mathsf KB'+\mathbf A_0(B'))$, and the last gives
$\mathbf H(B')=\mathcal L(\mathsf KB'+\mathbf A_0(B'))$.
Conversely, any solution in the stated balls must arise this
way. The factor-two estimate for $\mathcal L$ proves the
remaining bound in \eqref{I-87-38}.
Holomorphic iteration from zero has successive differences
at most $2^{-n}\hat{\mathsf b}\varepsilon_B$ in the strong
norm. Thus $\mathbf A_0$, and consequently $\mathbf H$, are
holomorphic on the open datum ball.

For two data put $A_i'=\mathsf KB_i'+\mathbf A_0(B_i')$ and
$r=\max_i\sup|B_i'|$. Their fixed-point equations give
\[
 \|\mathbf A_0(B_1')-\mathbf A_0(B_2')\|_w
 \le\tfrac12\bigl(\hat{\mathsf b}|B_1'-B_2'|_w
                  +\|\mathbf A_0(B_1')-\mathbf A_0(B_2')\|_w\bigr).
\]
Absorption proves the first bound in
\eqref{eq:response-weighted-bounds}, and then
$\|A_1'-A_2'\|_w\le2\hat{\mathsf b}|B_1'-B_2'|_w$.
Applying \eqref{eq:response-first-differences} proves the
bound for $\mathbf H$.

Finally set $A'=\mathsf KB'+\mathbf A_0(B')$.
The preceding bounds at the zero datum imply
$\|A'\|\le2\hat{\mathsf b}\sup|B'|$ and
$\|A'\|_w\le2\hat{\mathsf b}|B'|_w$. Hence
\eqref{I-87-30} and the first-correction bound give
\[
 \begin{split}
 \|\mathbf A_0(B')\|_w
     &\le8C_W\hat{\mathsf b}^5\sup|B'|\,|B'|_w,\\
 \|\mathsf H\mathsf D(A')\|_w
     &\le64C_2\mathsf e\hat{\mathsf b}^3\sup|B'|\,|B'|_w.
 \end{split}
\]
Since $64C_2\mathsf e\le2C_W$ and
$\mathbf H(B')-\mathsf KB'=\mathbf A_0(B')-\mathsf H\mathsf D(A')$,
these prove the last bound in
\eqref{eq:response-weighted-bounds}. The weight
$w(y)=e^{\theta_0d(y,\Sigma)}$ is admissible by the triangle
inequality and equals one on $\Sigma$; thus
$|B'|_w=\sup|B'|$. This proves
\eqref{eq:response-support-decay}.
\end{proof}


\section{The local nonlinear maps}\label{sec:actual-local-nonlinear-maps}

We verify the local hypotheses of \cref{sec:nonlinear-response}
for the relative averages and the action of a regular pair.
The averaging estimates follow from the fixed local gauges used
in \cref{lem:regular-local-inverse}. For the action, the plaquette
calculation of \cref{sec:action-estimates} must allow a complex
reference and the incomplete plaquette stars at the boundary.

Keep the geometry and natural regular-pair class of
\cref{sec:regular-system-inverse}, with
$\mathfrak M\ge\mathfrak M_*$. Use the dimensionless row
pairing, the norms, and $N(y)$ of \cref{sec:nonlinear-response}.
For a pair $(\mathbf U,\mathbf J)$ define
\[
 \mathsf g_y(A)=\sup_{b,\nu}\xi_y^2
                |(\nabla^{\mathbf U}_\nu A)(b)|,
 \qquad \|A\|_y=\max\{|A|_y,\mathsf g_y(A)\},
\]
where $b$ lies over $y$ and the translated bond and connecting
bond used by the difference are both inner bare bonds of
$\Omega_0$. Let $\mathcal Q(A)(c)$ be the relative iterate
\eqref{eq:flat-straightened-recursion} on the double block of
$c$, starting with $a_0=\xi A$ and evaluated at level $n(c)$.
Its definedness is part of the proposition below. Set
\[
 Q=\mathcal Q'(0),\qquad \mathsf C(A)=\mathcal Q(A)-QA,
 \qquad
 F(A)=\xi^{-1}\sum_{p\subset\Omega_0}a_p(e^{i\xi A}\mathbf U),
\]
where the sum is over the inner bare plaquettes. With the
gradient taken in the bare pairing of \eqref{I-86-6}, put
\[
 \mathcal V(A)=F'(A)-F'(0)-F''(0)A.
\]
The action and its gradient are entire: each plaquette term
is a finite product of matrix exponentials and fixed reference
matrices, and the pairing is nondegenerate.

\begin{proposition}\label{I-87-84}
There are constants $a_{\mathrm{nl}}>0$, $c_4,\rho_1\in(0,1]$
and $c_\flat,C_2,C_3,C_{\mathcal V}\ge1$, depending only on $L$,
with the following property. If
\[
 (\mathbf U,\mathbf J)\in
 \mathfrak P_t^\natural((\Omega_n);\alpha_0,\alpha_1,\gamma_0),
 \qquad 0<\alpha_0,\alpha_1,\gamma_0\le a_{\mathrm{nl}},
\]
then $\mathcal Q$ is defined and holomorphic on
$\{\sup|A|<c_4\}$. The maps $Q,\mathsf C,\mathcal V$ satisfy
hypotheses \textup{(i)--(iii)} of
\cref{sec:nonlinear-response} with these constants and the
seminorms above. Also $\sup|\mathbf J|^\sharp\le\gamma_0\le1$.
The constants are uniform in the volume and number of layers.
\end{proposition}

\begin{proof}
\emph{The averaging maps.}
For a row $c$ let $\mathsf P_c$ be the positive bare bonds
whose endpoints lie in $B^{n(c)}(c_-)\cup B^{n(c)}(c_+)$.
The double-block geometry in \cref{lem:sharp-row-neighborhood}
gives
\begin{equation}\label{eq:actual-average-geometry}
 \#\mathsf P_c\le6(\xi_c/\xi)^3,
 \qquad L^{-1}\xi_c\le\xi_{y_b}\le L\xi_c
       \quad(b\in\mathsf P_c).
\end{equation}
We shall use the slightly sharper distance bound
$d(y_c,y_b)\le(5L+3)/2<\rho_0$.
For $n(c)\ge1$, join $y_c$ to $b_-$ by coordinate segments
inside the double block, of total length at most $5\xi_c/2$.
The block lies in $\Omega_{n(c)-1}$, so this path has weighted
length at most $5L/2$. The path from $b_-$ to its cell center
has weighted length at most $3/2$. The coordinate paths cross
cell walls only finitely often, as in
\cref{lem:sharp-row-neighborhood}. For $n(c)=0$ one has $b=c$
and $y_b=y_c$. At each fixed level a bare bond belongs to at
most six double blocks: if its endpoints lie in one block,
the row must be incident to that block; otherwise its two
blocks determine the row. The scale comparison in
\eqref{eq:actual-average-geometry} allows at most three levels.
Thus each bare bond lies in at most $18$ sets $\mathsf P_c$.

Denote by $c_{\mathrm{av}},C_{\mathrm{av}},C_{\mathrm{var}}$
the constants of \cref{prop:flat-complex-averaging}, writing
$C_{\mathrm{var}}$ for its constant $C_2$. Choose
\[
 a_{\mathrm{nl}}=\min\left\{a_{\mathrm{reg}},\frac1{16},
           \frac{c_{\mathrm{av}}}{64L^2},\frac1{2L^2}\right\},
 \qquad c_4=\min\left\{1,\frac{c_{\mathrm{av}}}{2L}\right\}.
\]
For a row of positive level choose a covering regular piece
and fix its unitary gauge $u$ from
\cref{lem:regular-local-inverse}. By
\eqref{eq:regular-piece-potential},
\[
 \mathbf U^u=e^{i\xi\mathbf A_\kappa},\qquad
 \xi_c\sup_{\mathsf P_c}|\mathbf A_\kappa|
       \le8L^2(\alpha_0+\alpha_1)
       \le c_{\mathrm{av}}/4.
\]
On the whole local ball
$\mathcal U_c=\{\xi_c\sup_{\mathsf P_c}|A|<c_{\mathrm{av}}\}$,
rescale by $\xi_c$, with $\eta=\xi/\xi_c$ and
$\hat A=\xi_c\operatorname{Ad}(u)A$.
\Cref{prop:flat-complex-averaging} defines every relative
iterate of $e^{i\eta\hat A}e^{i\eta\xi_c\mathbf A_\kappa}$.
Pulling back by $u$ gives the actual recursion at every point
of $\mathcal U_c$: its first input is $\xi A$, and
\eqref{I-10-60} and \eqref{eq:straightening-covariance}
identify each successive contour, endpoint correction and
series logarithm. Unitary conjugation preserves their domains.
Block nesting keeps all inputs within $\mathsf P_c$.
Different covering pieces consequently evaluate the same
recursion, since their ungauged initial data and successive
operations coincide. This also proves locality on $\mathsf P_c$.
For a global field,
\[
 \xi_c\sup_{\mathsf P_c}|A|
       \le L\max_{b\in\mathsf P_c}|A|_{y_b}
       \le L\sup|A|<Lc_4\le c_{\mathrm{av}}/2.
\]
All rows are therefore defined on the asserted ball. At level
zero, including every row of a boundary piece, the recursion
is simply $\mathcal Q(A)(c)=\xi A(c)$.

The counting-measure derivative kernel
$K_A(c,b)=\partial\mathcal Q(A)(c)/\partial A(b)$ vanishes
off $\mathsf P_c$. The input rescaling contributes $\xi_c$
to the kernel in
\eqref{eq:flat-complex-differentiated-kernel}; hence its factor
is $\xi_c\eta^3=\xi^3\xi_c^{-2}$. The fixed unitary gauges
preserve its norm. Thus, throughout $\mathcal U_c$,
\begin{equation}\label{I-87-77}
 \begin{split}
 |K_A(c,b)|&\le C_Q\xi^3\xi_c^{-2},\\
 |K_A(c,b)-K_0(c,b)|
    &\le C_Q\xi^3\xi_c^{-2}\,
                         \xi_c\sup_{\mathsf P_c}|A|,
 \end{split}
\end{equation}
where one may take
$C_Q=\max\{1,C_{\mathrm{var}},
1+C_{\mathrm{av}}c_{\mathrm{av}}/2
 +C_{\mathrm{var}}c_{\mathrm{av}}\}$.
At level zero the kernel is $\xi I$ on its single bond, and
its difference is zero. The recursion vanishes at $A=0$, so
$\mathsf C(0)=\mathsf C'_0=0$.
The physical average in
\eqref{eq:small-system-physical-averages} divides this
recursion by $\xi_c$; thus
$(QA)(c)=\xi_c(Q_{\mathrm{phys}}A)(c)$.

The row count and \eqref{I-87-77} give
$|QA|_y\le6C_QL\max_{N(y)}|A|$.
For transposes the pairing is essential:
\begin{equation}\label{eq:actual-average-transpose}
 (Q^{\mathsf T}X)(b)=\xi^{-3}
       \sum_{c:b\in\mathsf P_c}\xi_c^{-1}
                                  K_0(c,b)^{\mathsf T}X(c).
\end{equation}
There are at most $18$ terms, and the transpose of a matrix
kernel has norm at most twice its norm, by the pairing estimate
used in \eqref{eq:small-system-pairing-bounds}. Multiplying
by $\xi_{y_b}^3$ gives
$|Q^{\mathsf T}X|_y^\sharp\le36C_QL^3\max_{N(y)}|X|$.
We take $c_\flat=36C_QL^3$.

Put $A_s=A_2+s(A_1-A_2)$ and integrate the derivative of
$\mathsf C$ along this segment. The second bound in
\eqref{I-87-77} yields
\[
 |\mathsf C(A_1)(c)-\mathsf C(A_2)(c)|
 \le6C_Q\int_0^1\xi_c\sup_{\mathsf P_c}|A_s|\,ds
                       \;\xi_c\sup_{\mathsf P_c}|A_1-A_2|.
\]
For $A_2=0$ the integral supplies a factor $1/2$; for
$\sup|A_i|\le r$ it is at most $Lr$.
These are the quadratic and difference estimates with
$C_2=3C_QL^2$. Applying \eqref{eq:actual-average-transpose}
to $K_A-K_0$ gives
\[
 |\mathsf C'_A{}^{\mathsf T}X|_y^\sharp
 \le36C_QL^4\max_{N(y)}|A|\max_{N(y)}|X|.
\]
Indeed, if $b\in\mathsf P_c$, the triangle inequality puts
every input bond of $\mathsf P_c$ over $N(y_b)$.
For the difference of two such transposes, use the whole ball
$\mathcal U_c$. The holomorphic kernel $K_A-K_0$ is bounded
there by $C_Q\xi^3\xi_c^{-2}c_{\mathrm{av}}$. If
$\sup|A_i|\le c_4/2$ and $\delta=A_1-A_2$ does not vanish
on $\mathsf P_c$, the disc
\[
 A_s+z\delta,\qquad
 |z|<\frac{c_{\mathrm{av}}}
                 {2\xi_c\sup_{\mathsf P_c}|\delta|}
\]
lies in $\mathcal U_c$. Cauchy's estimate followed by
integration in $s$ gives
$|K_{A_1}-K_{A_2}|\le
2C_Q\xi^3\xi_c^{-2}\xi_c\sup_{\mathsf P_c}|\delta|$.
The same transpose sum proves the required difference bound
with $C_3=72C_QL^4$. If $\delta$ vanishes locally, the kernel
difference is zero.

\emph{The action remainder.}
We first give the local complex-reference estimate needed for
rescaling. At spacing $0<\eta\le1$, let $\mathsf S_b$ be a
subset of the plaquettes incident to a positive bond $b$, and
let $\mathsf N_b$ be their positive bonds. Suppose either that
all four incident plaquettes are present or that $\eta=1$.
Assume on these bonds and plaquettes that
\[
 |\operatorname{Ad}(\mathbf U(b'))^{\pm1}|\le2,
 \qquad |\partial\mathbf U(p)-I|\le\alpha\eta^2,
 \qquad 0\le\alpha\le\tfrac12.
\]
Let $s=\sup_{\mathsf N_b}|A|$, let $g$ be the largest norm
of a covariant difference involving only bonds of
$\mathsf N_b$, and put $\mathsf n=\max(s,g)$.
Write $r(A)$ for the gradient at $b$ of
$\eta^{-1}\sum_{p\in\mathsf S_b}a_p(e^{i\eta A}\mathbf U)$,
using the $\eta^3$ pairing, and expand
$r(tA)=\sum_{q\ge0}t^qr_q(A)$.
We claim that, for an absolute $C_{\mathrm a}\ge1$,
\begin{equation}\label{eq:complex-star-remainder}
 \left|\sum_{q\ge2}r_q(A)\right|
       \le C_{\mathrm a}s\mathsf n,
 \qquad s\le\rho_{\mathrm a}:=1/256.
\end{equation}

Use the transported variables $\alpha_j$, the logarithmic
remainder $N_p$, and $\Theta_p=DA(p)+N_p$ from
\cref{subsec:action-plaquette-calculation}, now with
$\mathbf U$ in place of $U_0$. The transport bounds give
$|\alpha_1|,|\alpha_4|\le s$ and
$|\alpha_2|,|\alpha_3|\le2s$. The general complex-matrix
calculation leading to \eqref{eq:action-log-derivative} gives,
for $s\le1/128$,
\[
 |N_p|\le144s^2,\qquad
 \eta|\Theta_p|\le9s,\qquad
 |\Theta_p|\le2g+144s^2,
 \qquad |\mathsf L_{p,b}|\le192s.
\]
Here $\mathsf L_{p,b}$ is the derivative correction in
\eqref{I-10-41}; the last bound includes the possible factor
two from transport of the varied bond.
Put $c_p=\tfrac12\tr\partial\mathbf U(p)$ and
$F_p=(i\eta^2)^{-1}(\partial\mathbf U(p)-c_pI)$.
Cayley--Hamilton gives
$c_p^2+\eta^4\langle F_p,F_p\rangle=1$.
Also $2i\eta^2F_p=\partial\mathbf U(p)-\partial\mathbf U(p)^{-1}$.
The Neumann estimate bounds
$|\partial\mathbf U(p)^{-1}-I|$ by $2\alpha\eta^2$.
Together with $|c_p-1|\le\alpha\eta^2$, these identities give
\[
 |F_p|\le\tfrac32\alpha,\qquad
 |c_p|\le\tfrac32,
 \qquad |1+c_p|\ge\tfrac32,
 \qquad |1-c_p|\le\tfrac32\alpha^2\eta^4.
\]
Thus the algebraic trace and differential identities
\eqref{eq:action-plaquette-trace} and \eqref{I-10-44} still
hold, with
\[
 |M_p^{\mathrm r}|
       \le\eta^4(2|\Theta_p|^3+3|F_p||\Theta_p|^2).
\]
Subtracting $r(0)+D^{\mathsf T}DA(b)$ gives the same four
terms as \eqref{eq:action-four-remainders}. Here $D$ maps
the bonds of $\mathsf N_b$ to $\mathsf S_b$, and its transpose
obeys $|D^{\mathsf T}X(b)|\le6\eta^{-1}\sup|X|$.
The first, third and fourth terms are bounded, respectively,
by $54\alpha s$, $Cs^2(s+\alpha)$ and $Cs(\alpha+\mathsf n)$,
using the displayed bounds and
$|\mathsf L_{p,b}^{\mathsf T}|\le2|\mathsf L_{p,b}|$.

For $D^{\mathsf T}(cN)$, suppose first that all four
plaquettes are present. Pair opposite incident plaquettes
$p,p'$ in the transpose and let $v$ be the connecting inverse
bond transport. The identities
\eqref{eq:action-alpha-differences} give
\[
 |\operatorname{Ad}(v)\alpha_j(p')-\alpha_j(p)|
          \le5\eta g\qquad(1\le j\le4).
\]
For $j=1,4$ the bound is $2\eta g$; for $j=2,3$ insert
those identities, giving the three contributions
$2\eta g,2\eta g,\eta g$.
Both the original and transported quadruples have sum of
norms at most $12\eta s<1/8$. The series logarithm commutes
with this conjugation on their common domain. Integrating
\eqref{eq:action-log-derivative} along the segment between
the quadruples therefore gives
\[
 |\operatorname{Ad}(v)N_{p'}-N_p|
             \le3840\eta sg.
\]
Together with $|c_p-c_{p'}|\le3\alpha^2\eta^4$, this
bounds $D^{\mathsf T}(cN)$ by $Cs\mathsf n$.
If the star is incomplete, $\eta=1$, and the direct bound
$6\sup|cN|\le Cs^2$ suffices. Consequently
\[
 |r(A)-r(0)-D^{\mathsf T}DA(b)|\le Cs(\alpha+\mathsf n),
 \qquad s\le1/128.
\]
For fixed $A$ with $0<s\le1/256$, apply the coefficient
estimate used after \eqref{I-10-45} on the circle
$|t|=T=(128s)^{-1}\ge2$. The sum of the coefficients of
order at least two is at most
$2Cs(\alpha/T+\mathsf n)\le130Cs\mathsf n$.
These are exactly the coefficients $r_q$, $q\ge2$.
For $s=0$ all coefficients vanish by locality. This proves
\eqref{eq:complex-star-remainder}.

Return to a bare bond $b$ over $y$, and set $n=n(y)$.
The inner incident plaquettes use only bonds over sites $z$
with $d(y,z)\le5$ and $|n(z)-n|\le1$.
For the distance, join the two initial sites along at most
two bare edges, at cost at most two, and join each to its cell
center, at cost at most $3/2$. The adjacent-scale assertion
follows from the layer separation, as in
\cref{lem:sharp-row-neighborhood}: the corners lie within
$\xi$ of $b_-$ and cannot cross two layers.
If $n\ge1$, all four plaquettes are inner, since
$b_-\in\Omega_1$ and
$\operatorname{dist}_\infty(\Omega_0^c,\Omega_1)>
R_1\mathfrak M\xi>\xi$.
At $n=0$ the rescaled spacing below is one.

Rescale by $\xi_n$, putting
$\eta=\xi/\xi_n$ and $\hat A=\xi_nA$.
The regular-pair conditions give
\[
 |\operatorname{Ad}(\mathbf U(b'))^{\pm1}|
       \le e^{2\alpha_1}\le2,
 \qquad |\partial\mathbf U(p)-I|
       \le L^2\alpha_0\eta^2,
 \qquad L^2\alpha_0\le\tfrac12.
\]
The local estimate thus applies with $\alpha=L^2\alpha_0$.
The rescaled local star action is $\xi_n$ times the original
one, whereas a variation of $\hat A$ varies $A$ by $\xi_n^{-1}$
times as much. These factors cancel. Its derivative at $b$
equals that of the full action, and the ratio of the pairing
weights is $(\xi/\eta)^3=\xi_n^3$. Consequently
\[
 \xi_n^3F'(A)(b)=r(\hat A),\qquad
 \xi_n^3\mathcal V(A)(b)=\sum_{q\ge2}r_q(\hat A).
\]
On this star the local seminorms satisfy
\[
 s(\hat A)\le L\max_{N(y)}|A|,
 \qquad \mathsf n(\hat A)\le L^2\max_{N(y)}\|A\|.
\]
Indeed, the rescaled differences are $\xi_n^2$ times the
original differences, and every difference on the star is
one admitted in the corresponding $\mathsf g_z$.
It follows from \eqref{eq:complex-star-remainder} that
\[
 |\mathcal V(A)|_y^\sharp
       \le C_{\mathrm a}L^3
                 \max_{N(y)}|A|\max_{N(y)}\|A\|
 \quad\text{if }L\sup|A|\le\rho_{\mathrm a}.
\]
Take $\rho_1=\rho_{\mathrm a}L^{-2}$ and
$C_{\mathcal V}=C_{\mathrm a}L^4$.

Finally, suppose $\|A_i\|\le r\le\rho_1/4$ and put
$\delta=A_1-A_2$. If $\mathsf n(\hat\delta)>0$, on the disc
$A_s+z\delta$ of radius $R=L^2r/\mathsf n(\hat\delta)$ one has
$\mathsf n(\hat A_s+z\hat\delta)\le2L^2r\le\rho_{\mathrm a}/2$.
The local holomorphic remainder is bounded there by
$C_{\mathrm a}(2L^2r)^2$. Cauchy's estimate and integration
in $s$ give
\[
 \xi_y^3|\mathcal V(A_1)(b)-\mathcal V(A_2)(b)|
 \le4C_{\mathrm a}L^2r\,\mathsf n(\hat\delta)
 \le4C_{\mathcal V}r\max_{N(y)}\|\delta\|.
\]
If $\mathsf n(\hat\delta)=0$, locality makes the difference zero.
The identities $\mathcal V(0)=\mathcal V'(0)=0$ follow from
its definition, and the bound on $\mathbf J$ is the
regular-pair condition at the initial site of each bond.
\end{proof}

\section{The projected propagator}\label{sec:projected-propagator}

The nonlinear response uses the part of the background inverse
left after subtracting its scalar and coarse constraints. After the
coarse subtraction the operator is already a block of the constrained inverse.
To estimate the scalar subtraction, we place two copies of that
system in a triangular system. Its coupling contains a derivative
in the first equation. A local H\"older estimate permits this
derivative to be treated as a divergence before the local inverses
are combined.

Keep the setting and operators of \cref{sec:background-inverses}.
In this section write $Q_{\mathrm{phys}}$ for its physical linear
average; the dimensionless average in
\cref{sec:actual-local-nonlinear-maps} is
$Q=\mathsf mQ_{\mathrm{phys}}$.
Define
\begin{equation}\label{eq:projected-propagator-definition}
 \mathsf E_1=\mathsf G_1D\mathsf RD^{\mathsf T}\mathsf G_1,
 \qquad
 \mathfrak G=\mathsf G_1-\mathsf E_1
                     -\mathsf H_1Q_{\mathrm{phys}}\mathsf G_1.
\end{equation}
The projection $\mathsf R$ is the scalar operator in
\eqref{eq:bg-projection-operators}.

\begin{proposition}\label{prop:projected-propagator}
There are a positive constant $a_{\mathrm{pr}}$, a power
$\mathfrak M_{\mathrm{pr}}$ of $L$, and a constant $B_{\mathrm{pr}}$,
depending only on $L$, with the following property. In the
geometry of \cref{I-84-25}, suppose
$\mathfrak M\ge\mathfrak M_{\mathrm{pr}}$ and
\[
 (\mathbf U,\mathbf J)\in
 \mathfrak P_t^\natural((\Omega_n);\alpha_0,\alpha_1,\gamma_0),
 \qquad 0<\alpha_0,\alpha_1,\gamma_0\le a_{\mathrm{pr}}.
\]
Then the operators in \eqref{eq:projected-propagator-definition}
exist and, for $T=\mathsf E_1$ or $\mathfrak G$,
\begin{equation}\label{eq:projected-propagator-bound}
 \max\{|Tf|_{-1;y},|\nabla Tf|_{-2;y}\}
 \le B_{\mathrm{pr}}\sup_z
       e^{-\delta_{\mathrm{reg}}d(y,z)}|f|_{-3;z}.
\end{equation}
Here $f$ is any bare bond field and $y\in\Lambda_*$.
The derivative is taken at $\mathbf U$. The constants are
uniform in the spacing, volume and number of layers.
\end{proposition}

\begin{proof}
\emph{The two local estimates.}
We first allow an independent multiplier $\mathsf j$ and use only
the system $[4]$, suppressing this superscript. Set $\rho=3L+6$.
For a site or bond field, including a bond field with an additional
direction, define the transported half-H\"older norm by
\begin{equation}\label{eq:enhanced-local-holder}
 \begin{split}
 [f]_{\mathbf U,1/2;y}
 &=\sup_{i\ne i'}
   \left(\frac{\xi_y}{|i-i'|_1}\right)^{1/2}
   |\operatorname{Ad}(\mathbf U(\Gamma_{x,x'}))f(i')-f(i)|,\\
 |f|^{1/2}_{\mathbf U,s;y}
 &=\xi_y^{-s}\left(\sup_{\Delta^{3/2}(y)}|f|
                                  +[f]_{\mathbf U,1/2;y}\right).
 \end{split}
\end{equation}
The pairs have the same kind and both base points in
$\Delta^{3/2}(y)$, as in \eqref{eq:trivial-system-holder-norm}.
The points $x,x'$ are their sites or initial endpoints, and
$\Gamma_{x,x'}$ is the coordinate comb \eqref{I-3-90}.
Above level zero its containing box lies in the preceding region,
so all its bonds are inner. At level zero a distinct bond pair
consists of the incoming and outgoing bonds in one direction,
and the connecting comb is one inner bond. Empty suprema are zero.
The subscript $0$ denotes the identity configuration.

Let $X_1(b,\nu)$ be a $\mathfrak g^c$-valued bond tensor and put
\begin{equation}\label{eq:enhanced-local-norms}
 \begin{split}
 g^+&=(g,X_1),\qquad
 \iota^+g^+=\iota g+e_1\nabla^{\mathsf T}X_1,\\
 |g^+|^{1/2}_y
   &=\max\{|g|_y,|X_1|^{1/2}_{\mathbf U,-2;y}\},\\
 |\Psi|^{+1/2}_y
   &=\max\{|\Psi|_y,|\nabla A|^{1/2}_{\mathbf U,-2;y}\}.
 \end{split}
\end{equation}
The transpose uses the bare bilinear pairing. The additional
divergence enters the first equation; $\iota$ retains its
fifth-slot divergence. We will prove, on every local piece,
\begin{equation}\label{eq:enhanced-local-estimates}
 \begin{split}
 |\mathbb E_\kappa\iota_\kappa^+g^+|_y
 &\le B_{\mathrm e}\sup_z e^{-2\delta_{\mathrm{reg}}d^\kappa(y,z)}
                                      |g^+|^{1/2}_z,\\
 |\mathbb E_\kappa\iota_\kappa g|^{+1/2}_y
 &\le B_{\mathrm e}\sup_z e^{-2\delta_{\mathrm{reg}}d^\kappa(y,z)}
                                      |g|_z,
 \end{split}
\end{equation}
with $B_{\mathrm e}\ge1$ depending only on $L$, after decreasing
the admissible smallness radius.

At the identity these estimates follow from the elimination in
\cref{sec:trivial-system-elimination,sec:trivial-system-estimates}.
With $p=p'+p''$ as in \eqref{eq:trivial-system-p-split}, the
effective vector source for $g^+$ is
\[
 f'_1+\nabla^{\mathsf T}(X^{(p)}+X_1),\qquad
 f'_1=f_1+D\varphi+Dp'+Z^{(p)}.
\]
The boundary correction $Z^{(p)}$ is exactly the one in
\eqref{eq:trivial-system-vector-source}; it has weight $-3$.
For the first estimate use the half-H\"older bounds for
$p''$ and $X^{(p)}$, and then
\eqref{eq:trivial-system-basic-bounds} and \eqref{eq:trivial-system-mixed-holder}.
These bound $A$ and $\nabla A$ in weights $-1,-2$.
The remaining components follow from the same scalar substitutions
and bounds used in \cref{sec:trivial-system-estimates}.
For the second estimate $X_1=0$. Now use exponent $3/4$ in
\eqref{eq:trivial-system-projected-holder} and \eqref{eq:trivial-system-shift-holder},
and apply \eqref{eq:trivial-system-gradient-holder} to $f'_1$
and $X^{(p)}$. This bounds the additional half-H\"older norm
of $\nabla A$. Both estimates have rate $\delta_{\mathrm{cs}}$.

To transfer them to a regular piece, work in its small gauge and
suppress $\kappa$. Require
$\epsilon\le\epsilon_{\mathrm e}:=\min\{\epsilon_*,2^{-11}\}$.
We record the local geometry needed for the transported norms.
Let $\mathsf M(y)$ consist of the tops of cells meeting
$\widetilde\Delta(y)$, together, at level zero, with the positions
of $y-2\xi e_\mu$ when both preceding sites belong to the region.
Then
\begin{equation}\label{eq:enhanced-local-neighborhood}
 z\in\mathsf M(y)\quad\Longrightarrow\quad
 d(y,z)\le\tfrac32(L+2)=\rho/2,\qquad
 L^{-1}\xi_y\le\xi_z\le L\xi_y.
\end{equation}
For a meeting cell, join its intersection point to $y$ and to
its own top. The first segment costs at most $3/2$ inside
$\Delta(y)$ and $3L/2$ outside it: the latter part has
$\ell^1$-length at most $3\xi_y/2$ and lies in the preceding
region. The segment in the meeting cell costs at most $3/2$.
This uses \eqref{eq:layer-density-distance}; at level zero
the outer cost only improves. The exceptional point is reached
by two neighboring cells, at cost at most six. The scale bound
is \eqref{eq:sharp-neighborhood-size}.
Adding neighbors to $\mathsf M(y)$ gives a set $\mathsf M^+(y)$
of radius at most $3L/2+6$ with the same scale comparison.

The comb in \eqref{eq:enhanced-local-holder} lies in
$\widetilde\Delta(y)\cap Y\subset\mathsf N(y)$. The sum of
the norms of its adjoint exponents is at most
$2\epsilon|x-x'|_1/\xi_y\le9\epsilon$.
Thus its transport and inverse differ from $I$ by at most
$4\epsilon|x-x'|_1/\xi_y$. Comparing the two differences gives
\begin{equation}\label{eq:enhanced-holder-comparison}
 (1+18\epsilon)^{-1}|f|^{1/2}_{0,s;y}
 \le |f|^{1/2}_{\mathbf U,s;y}
 \le(1+18\epsilon)|f|^{1/2}_{0,s;y}.
\end{equation}

The required shift and product bounds follow directly from this
norm. For an ordinary index shift by at most $\xi$ in each
coordinate, split pairs at distance $\xi_y/(8L)$.
The more distant pairs, and all pairs at levels zero and one,
are bounded by $C(L)$ times the supremum. For a closer pair
at a higher level, both shifted indices exist and their base
points lie in the enlarged cube of the first shifted index's
position: their separation is at most $\xi_z/8$ and the first
base point lies in $\Delta(z)$. Its H\"older norm applies.
The positions lie in $\mathsf M^+(y)$, and the scale factors
cost only powers of $L$. Expanding the difference of a product
also shows that a coefficient bounded by $a_0$ with differences
bounded by $a_1|x-x'|_1/\xi_y$ has multiplier norm at most
$a_0+9a_1$. Missing shifted indices occur only in the low-level
case already bounded by the supremum.

Put $\mathsf U_\nu=\operatorname{Ad}(\mathbf U_\nu)$ and
$\mathsf e_\nu=(\mathsf U_\nu-I)/\xi$.
The bounds \eqref{eq:small-system-connection-bounds}, summed
along the comb, give
\[
 |\mathsf e_\nu|\le4\epsilon\xi_y^{-1},\qquad
 |\mathsf e_\nu(x)-\mathsf e_\nu(x')|
       \le4\epsilon\xi_y^{-2}|x-x'|_1.
\]
When the shifted bond exists,
$(\nabla_{\mathbf U,\nu}-\nabla_{0,\nu})A(b)
 =\mathsf e_\nu(b_-)A(b+\xi e_\nu)$; otherwise the difference
is zero. The preceding product and shift bounds, followed by
the ordinary comb estimate, bound its half-H\"older norm by
$C(L)\epsilon\max_{d(y,z)\le2\rho}|\Psi|_{0,z}$.
Together with \eqref{eq:small-system-norm-comparison} and
\eqref{eq:enhanced-holder-comparison}, this compares the two
enhanced solution norms by a nearby maximum with constant $C(L)$,
in either direction.

The injection difference contains no derivative of its input.
With $\mathsf e'_\nu=(\mathsf U_\nu^{-1}-I)/\xi$,
\[
 ((\nabla_{\mathbf U}^{\mathsf T}-\nabla_0^{\mathsf T})X_1)(b)
   =\sum_\nu\mathsf e'_\nu(b_--\xi e_\nu)
                          X_1(b-\xi e_\nu,\nu),
\]
where missing predecessors contribute zero. Combining this with
\eqref{eq:small-system-injection} gives
\[
 \iota_{\mathbf U}^+g^+
 =\iota_0^+\bigl(g^++(\mathsf n^+g^+,0)\bigr),\qquad
 |\mathsf n^+g^+|_y\le C(L)\epsilon
                 \max_{d(y,z)\le3}|g^+|^{1/2}_{0,z}.
\]
The two nonzero entries have the required weights $-3,-4$,
since the coefficients have weight $-1$ and the tensors
$X_1,X_5$ have weights $-2,-3$.

Use the existing factorization
$\mathbb K_0-\mathbb K=\iota_0\mathbb V$ and put
$\Psi_0=\mathbb E_0\iota_0^+g^+$.
The identities
\begin{equation}\label{eq:enhanced-local-resolvents}
 \mathbb E\iota_0^+g^+
   =\Psi_0+\mathbb E\iota_0\mathbb V\Psi_0,
 \qquad
 \Psi=\mathbb E\iota_0g
   =\mathbb E_0\iota_0(g+\mathbb V\Psi)
\end{equation}
transfer the two estimates. In the first identity bound $\Psi_0$
strongly by the identity-configuration estimate, then apply
\eqref{eq:small-system-defect-bound} and the already proved
strong bound for $\mathbb E\iota_0$.
In the second, bound $\Psi$ strongly first and apply the
identity-configuration enhanced-output estimate to
$g+\mathbb V\Psi$. Composition and nearby maxima preserve
$\delta_{\mathrm{cs}}$ by
\eqref{eq:trivial-system-sup-decay}. The injection and norm
comparisons just proved return both estimates to $\mathbf U$.
Unitary gauge covariance preserves all transported norms and
intertwines the injections.

On a boundary piece $\xi_y=\xi$. The minimum nonzero distance
and the bound $|\mathsf U^{\pm1}|\le2$ imply
\[
 |f|^{1/2}_{\mathbf U,s;y}\le C(L)
       \max_{d(y,z)\le3}|f|_{s;z},\qquad
 |\nabla^{\mathsf T}X_1|_{-3;y}\le9
       \max_{d(y,z)\le3}|X_1|_{-2;z}.
\]
Hence $\iota^+g^+=\iota\widetilde g$ with
$|\widetilde g|_y\le10\max_{d(y,z)\le3}|g^+|^{1/2}_z$,
and $|\Psi|^{+1/2}_y\le C(L)\max_{d(y,z)\le6}|\Psi|_z$.
The massive boundary inverse of
\cref{lem:regular-local-inverse} proves both estimates.
Taking the larger constant and the common rate
$2\delta_{\mathrm{reg}}$ proves
\eqref{eq:enhanced-local-estimates}.

\emph{The coupling of the two systems.}
For $\Psi'=(A';v',w',\sigma';p',q',r')$ set
\[
 f=D^{\mathsf T}A'-w',\qquad
 \mathbb J\Psi'=e_1Df,\qquad
 \mathbb K^{\mathrm d}(\Psi,\Psi')
       =(\mathbb K\Psi-\mathbb J\Psi',\mathbb K\Psi').
\]
The prime here distinguishes the second copy of the system.
For $b=\langle x,x+\xi e_\mu\rangle$ put
\begin{equation}\label{eq:doubled-gradient-divergence}
 X^{(f)}_\nu(b)=-\delta_{\mu\nu}\mathsf U_\mu(x)f(b_+),
 \qquad
 c^{(f)}(b)=
 \begin{cases}
 -\xi^{-1}f(x),&b-\xi e_\mu\text{ is not inner},\\
 0,&b-\xi e_\mu\text{ is inner}.
 \end{cases}
\end{equation}
Expansion of the transpose gives
$Df=\nabla^{\mathsf T}X^{(f)}+c^{(f)}$.
The correction is supported at level zero, where $\xi_y=\xi$,
and therefore $|c^{(f)}|_{-3;y}\le|f|_{-2;y}$.
Write $\mathsf j^+\Psi'$ for the enhanced forcing with just
these two entries. Then
\begin{equation}\label{eq:doubled-coupling-bound}
 \mathbb J=\iota^+\mathsf j^+,\qquad
 |\mathsf j^+\Psi'|^{1/2}_y
       \le C_J\max_{d(y,z)\le\rho}|\Psi'|^{+1/2}_z.
\end{equation}

When both adjacent bonds exist, a summand of $D^{\mathsf T}A'$
is $-\mathsf U_\nu(x-\xi e_\nu)^{-1}
\nabla_\nu A'_\nu(x-\xi e_\nu)$.
Otherwise its site is at level zero and direct endpoint values
give the same weight. Thus
\[
 |f|_{-2;y}\le(12L^2+1)\max_{d(y,z)\le3}|\Psi'|_z,
 \qquad |X^{(f)}|_{-2;y}
       \le2L^2\max_{d(y,z)\le3}|f|_{-2;z}.
\]
At levels zero and one, the H\"older quotient costs at most
$3L$ times the enlarged supremum, so these bounds suffice.
At higher levels all the shifted indices are inner, and
\[
 X^{(f)}_\mu(b)=
 \sum_\nu\mathsf U_\mu(x)
   \mathsf U_\nu(x+\xi e_\mu-\xi e_\nu)^{-1}
   (\nabla_\nu A'_\nu)(x+\xi e_\mu-\xi e_\nu)
       +\mathsf U_\mu(x)w'(b_+).
\]
The coefficients in the sum are bounded by four, with ordinary
differences bounded by $16\epsilon_{\mathrm e}|x-x'|_1/\xi_y$.
The shift and product bounds above, followed by
\eqref{eq:enhanced-holder-comparison}, therefore use exactly
the half-H\"older norm of $\nabla A'$.
For $\bar w(b)=w'(b_+)$, direct subtraction gives
\[
 (\nabla_\nu\bar w)(b)
  =(Dw')(\langle b_+,b_++\xi e_\nu\rangle)
    +\xi^{-1}(\mathsf U_\nu(x)-\mathsf U_\nu(b_+))
                                      w'(b_++\xi e_\nu).
\]
The coefficient difference is at most
$4\epsilon_{\mathrm e}\xi/\xi_z$ at a position $z$.
Consequently the comb bound controls $\bar w$ by the existing
entries $w',Dw'$. Let
$T=\operatorname{Ad}(\mathbf U(\Gamma_{x,x'}))$.
For the transported product $\mathsf U_\mu\bar w$,
conjugating the coefficient by this transport
changes it by at most
$28\epsilon_{\mathrm e}|x-x'|_1/\xi_y$; this follows by
inserting the two factors $T-I,T^{-1}-I$ and then the ordinary
coefficient difference. Since $28\epsilon_{\mathrm e}\le1$,
the product estimate applies. The positions in the shifted sum
lie in $\mathsf M^+(y)$, and those in the $w'$ estimate lie
within six of $\mathsf M(y)$. Both radii are at most $\rho$.
This proves \eqref{eq:doubled-coupling-bound}. Boundary pieces
use only its low-level case.

The triangular inverse is
\[
 (\mathbb K^{\mathrm d})^{-1}(F,F')
       =(\mathbb EF+\mathbb E\mathbb J\mathbb EF',\mathbb EF').
\]
Use the maximum of the two ordinary forcing norms on
$\mathcal F^{[4]}\times\mathcal F^{[4]}$, and the maximum
of the two strong solution norms. Apply the second estimate
of \eqref{eq:enhanced-local-estimates} to the second copy,
then \eqref{eq:doubled-coupling-bound}, and then the first
estimate to the first copy. The local doubled inverse has
rate $2\delta_{\mathrm{reg}}$ and constant
\begin{equation}\label{eq:doubled-local-constant}
 B_{\mathrm{d,loc}}
    =B_{\mathrm e}(1+C_Je^\rho B_{\mathrm e}).
\end{equation}

\emph{The additional cutoff commutator.}
Take $h,\vartheta$ as in \cref{lem:regular-cutoff} and put
\[
 k_b^-=(h(\hat b)-h(b_-))/\xi,\qquad
 k_b^+=(h(b_+)-h(\hat b))/\xi,\qquad
 \phi_h=[D^{\mathsf T},h]A'.
\]
The midpoint convention gives
\begin{equation}\label{eq:doubled-commutator}
 \begin{split}
 [\mathbb J,h]\Psi'&=e_1([D,h]f+D\phi_h),\\
 ([D,h]f)(b)&=k_b^+\mathsf U(b)f(b_+)+k_b^-f(b_-),\\
 \phi_h(x)&=-\sum_{b_+=x}k_b^+\mathsf U(b)^{-1}A'(b)
                         -\sum_{b_-=x}k_b^-A'(b).
 \end{split}
\end{equation}
The first term has weight $-3$ bounded by
$C(L)\vartheta^{-1}\max_{d(y,z)\le6}|\Psi'|_z$.
At level zero the direct estimate for the last term is
\[
 |D\phi_h|_{-3;y}\le\tfrac{27}{2}\vartheta^{-1}
                       \max_{d(y,z)\le3}|A'|_{-1;z}.
\]
At a positive level the cancellation between incoming transports
is needed. Write $x^-=x-\xi e_\nu$ for a fixed direction $\nu$.
For $b=\langle x,x+\xi e_\mu\rangle$, expansion gives
$\xi D\phi_h(b)=-\sum_\nu(T^-_\nu+T^+_\nu)$, where
\[
 T^-_\nu=\xi k^-_\nu(x+\xi e_\mu)\nabla_\mu A'_\nu(x)
           +(k^-_\nu(x+\xi e_\mu)-k^-_\nu(x))A'_\nu(x)
\]
and
\begin{align*}
 T^+_\nu={}&\mathsf U_\nu(x^-)^{-1}
 \bigl\{\xi k^+_\nu(x^-+\xi e_\mu)\nabla_\mu A'_\nu(x^-)
       +(k^+_\nu(x^-+\xi e_\mu)-k^+_\nu(x^-))A'_\nu(x^-)\bigr\}\\
 &+k^+_\nu(x^-+\xi e_\mu)\mathsf W_\nu A'_\nu(x^-+\xi e_\mu),\\
 \mathsf W_\nu={}&\mathsf U_\mu(x)\mathsf U_\nu(x^-+\xi e_\mu)^{-1}
                      -\mathsf U_\nu(x^-)^{-1}\mathsf U_\mu(x^-).
\end{align*}
For $\nu=\mu$ the last coefficient vanishes. Otherwise it is
the difference of the transports along the two sides of a
plaquette. If its variable is $P$, the adjoint bounds give
$|\mathsf W_\nu|\le48|P-I|$.
The small-potential bounds imply
\[
 |P-I|\le8\epsilon_{\mathrm e}\xi^2\xi_y^{-2},\qquad
 |\mathsf W_\nu|\le\xi^2\xi_y^{-2}.
\]
Indeed the sum of the four exponents is bounded by
$2\epsilon_{\mathrm e}\xi^2\xi_y^{-2}$ and their quadratic
remainder by $16\epsilon_{\mathrm e}^2\xi^2\xi_y^{-2}$;
$384\epsilon_{\mathrm e}\le1$ gives the second inequality.
All these plaquettes lie in the collar where the small-potential
bounds hold. The cutoff assumptions give
$|k^\pm|\le\vartheta^{-1}\xi_y^{-1}/2$ and adjacent differences
at most $\vartheta^{-1}\xi\xi_y^{-2}/2$.
Substitution in the expansion therefore yields
\[
 |D\phi_h(b)|\le\vartheta^{-1}
     \left(\tfrac92\xi_y^{-1}\max|\nabla A'|
                         +6\xi_y^{-2}\max|A'|\right).
\]
The nearby scale comparison proves the required weight $-3$
bound without a negative power of the bare spacing in its constant.

Let $c_h\Psi'$ be the ordinary forcing whose only entry is
$[D,h]f+D\phi_h$. Increase $C_J$ to include this estimate.
With the original commutator map $\mathbb C_h$ set
\[
 \mathbb C_h^{\mathrm d}(\Psi,\Psi')
       =(\mathbb C_h\Psi-c_h\Psi',\mathbb C_h\Psi').
\]
Writing $\iota^{\mathrm d}=\iota\oplus\iota$, we have
\begin{equation}\label{eq:doubled-cutoff-bound}
 [\mathbb K^{\mathrm d},h]=\iota^{\mathrm d}\mathbb C_h^{\mathrm d},
 \qquad
 |\mathbb C_h^{\mathrm d}\vec\Psi|_y
 \le(C_{\mathbb C}+C_J)\vartheta^{-1}
                 \max_{d(y,z)\le\rho}|\vec\Psi|_z.
\end{equation}
This seminorm vanishes when $h$ is constant on $\mathsf N(y)$.
The construction is gauge covariant, so it applies in the original
local gauges. Its output is ordinary doubled forcing.

\emph{The total inverse and the physical multiplier.}
Use the same pieces and cutoffs as in
\cref{subsec:regular-parametrix}, and write
$\delta=\delta_{\mathrm{reg}}$. Choose
\begin{equation}\label{eq:doubled-global-constants}
 \begin{split}
 C_q^{\mathrm d}&=2(C_{\mathbb C}+C_J)C_{\mathrm G}
                              B_{\mathrm{d,loc}}e^{2\rho+8},\\
 B_T^{\mathrm d}&=4B_{\mathrm{d,loc}}e^{3\delta},\qquad
 q_{\mathfrak M}^{\mathrm d}
       =C_q^{\mathrm d}\mathfrak M^{-1}e^{-\delta D_0},\qquad
 B_{\mathrm d}=2N_{\mathrm G}B_T^{\mathrm d}.
 \end{split}
\end{equation}
Enlarge the power threshold to include all earlier thresholds and
$2N_{\mathrm G}C_q^{\mathrm d}$.
The ordinary cutoff correction $\theta_\kappa\star g$ acts
on each forcing copy. The new coupling uses only $w'$ at the
two endpoints and $A'$ on incident bonds. These indices lie
within two bare spacings of the support, hence in the existing
common-index margins between $\kappa^{[2/3]}$ and
$\kappa^{[5/6]}$. Thus the local/global operator agreement
and the lifted identities \eqref{eq:regular-lift-identities}
hold on the product space.

The proof of \cref{lem:regular-remainder} now applies with
$B_{\mathrm{loc}}$ replaced by $B_{\mathrm{d,loc}}$ and
$C_{\mathbb C}$ by $C_{\mathbb C}+C_J$.
The commutator support and its separation from the input are
unchanged, so its remainder constant is
$q_{\mathfrak M}^{\mathrm d}$. The overlap bound gives
$N_{\mathrm G}q_{\mathfrak M}^{\mathrm d}\le1/2$.
The same Neumann argument therefore proves
\begin{equation}\label{eq:doubled-total-bound}
 |(\mathbb K^{\mathrm d})^{-1}\iota^{\mathrm d}(g,g')|_y
 \le B_{\mathrm d}\sup_z e^{-\delta d(y,z)}
                                      \max\{|g|_z,|g'|_z\}.
\end{equation}

To apply this at the actual multiplier, take
\[
 a_{\mathrm d}=\min\{a_{\mathrm{reg}},\epsilon_{\mathrm e}/(16L^2)\},
 \qquad
 a_{\mathrm{pr}}=\min\{a_{\mathrm{bg}},
                       a_{\mathrm d}/\max(1,C_{\mathrm{bg}})\}.
\]
Include $\mathfrak M_{\mathrm{bg}}$ in
$\mathfrak M_{\mathrm{pr}}$.
By \eqref{eq:bg-total-multipliers},
$|\mathsf j_*|_{-3;y}\le C_{\mathrm{bg}}\gamma_0\le a_{\mathrm d}$.
The existing local gauge estimate consequently gives
\[
 \max\{8L^2(\alpha_0+\alpha_1),\gamma_0,
                         \sup_y|\mathsf j_*|_{-3;y}\}
       \le\epsilon_{\mathrm e}.
\]
The truncated multiplier has the same bound on every regular
piece. On boundary pieces its quadratic coefficient vanishes.
Thus all local estimates used in \eqref{eq:doubled-total-bound}
apply at this physical triple.

Finally apply the doubled inverse to $(0,e_1f)$.
The scalar elimination in \cref{I-63-67} gives
$A'=\mathsf G_1f$ and $w'=\mathsf PD^{\mathsf T}A'$.
The first-copy $A$ component is therefore
$\mathsf G_1D\mathsf RD^{\mathsf T}\mathsf G_1f=\mathsf E_1f$.
This proves its bound with constant $B_{\mathrm d}$.
The same elimination for system $[5]$, with eighth input zero,
gives
\[
 [\mathbb E^{[5]}(\mathsf j_*)]_{A,1}
       =\mathsf G_1-\mathsf H_1Q_{\mathrm{phys}}\mathsf G_1.
\]
Its strong bound is \eqref{eq:regular-system-decay}.
Subtracting $\mathsf E_1$ proves
\eqref{eq:projected-propagator-bound} with
$B_{\mathrm{pr}}=B_{\mathrm{reg}}+B_{\mathrm d}$.
\end{proof}


\subsection{Application to the nonlinear response}
\label{sec:actual-response}

We can now apply \cref{I-87-35} to a regular pair. The row
variables in that theorem are dimensionless, so the physical
inverse blocks must first be rescaled.
Keep $\mathsf G,\mathsf H,\mathsf H_1$ from
\cref{sec:background-inverses}, and write ${}^*$ for the
transpose with physical row weight $\xi_c^3$.
With $\mathsf mX(c)=\xi_cX(c)$, put
\begin{equation}\label{eq:actual-response-operators}
 \begin{gathered}
 H=\mathsf H\mathsf m^{-1},\qquad
 K=\mathsf H_1\mathsf m^{-1},\\
 S=\mathsf m^3
       (Q_{\mathrm{phys}}\mathsf GQ_{\mathrm{phys}}^*)^{-1}
       \mathsf m^{-1}.
 \end{gathered}
\end{equation}
The fourth operator is $\mathfrak G$ of
\eqref{eq:projected-propagator-definition}.
The superscript ${}^{\mathsf T}$ denotes transpose for the dimensionless pairing
\eqref{I-86-6}.

\begin{proposition}\label{I-88-53}
In the geometry of \cref{I-84-25}, suppose
$\mathfrak M\ge\mathfrak M_{\mathrm{pr}}$ and
\[
 (\mathbf U,\mathbf J)\in
 \mathfrak P_t^\natural((\Omega_n);\alpha_0,\alpha_1,\gamma_0),
 \qquad 0<\alpha_0,\alpha_1,\gamma_0
                  \le\min\{a_{\mathrm{nl}},a_{\mathrm{pr}}\}.
\]
Use the nonlinear maps and seminorms of \cref{I-87-84} and set
\[
 \theta_0=\delta_{\mathrm{reg}}/2,\qquad
 \hat{\mathsf b}=\max\{1,B_{\mathrm{pr}},C_{\mathrm{bg}}\}.
\]
Define $c_W,C_W,\varepsilon_3,\varepsilon_B$ by
\eqref{I-86-46} with these constants. For
$\sup|B'|\le\varepsilon_B$, the system \eqref{I-87-37}, with
\[
 (\mathsf H,\mathsf S,\mathsf K,\mathsf G,\mathsf a)
                       =(H,S,K,\mathfrak G,I),
\]
has a unique solution $(X,\mathbf A_0,\mathbf H)$ in the balls
\[
 \sup|X|\le c_W\hat{\mathsf b}^{-2},\qquad
 \|\mathbf A_0\|\le\hat{\mathsf b}\varepsilon_B.
\]
Its components are holomorphic for $\sup|B'|<\varepsilon_B$.
The bounds \eqref{I-87-38},
\eqref{eq:response-weighted-bounds} and
\eqref{eq:response-support-decay} hold with these constants.
Moreover,
\begin{equation}\label{I-88-66}
 \mathcal Q(\mathbf H(B'))=B'+Q\mathbf A_0(B').
\end{equation}
All constants depend only on $L$.
\end{proposition}

\begin{proof}
We verify the linear bounds \eqref{I-86-50}; the local nonlinear
hypotheses are \cref{I-87-84}.
The two row pairings satisfy
$(X,Y)_{\mathbf B}=(X,\mathsf m^4Y)_1$.
Since $Q_{\mathrm{phys}}=\mathsf m^{-1}Q$, this gives
$Q_{\mathrm{phys}}^*=Q^{\mathsf T}\mathsf m^3$. Hence
\[
 S=(Q\mathsf GQ^{\mathsf T})^{-1},\qquad
 H=\mathsf GQ^{\mathsf T}S,\qquad
 K=\mathsf G_1Q^{\mathsf T}
                   (Q\mathsf G_1Q^{\mathsf T})^{-1}.
\]
In particular, $QH=QK=I$, and the physical quadratic term
$Q_{\mathrm{phys}}^*\mathsf m^{-2}Q_{\mathrm{phys}}$
is $Q^{\mathsf T}Q$, so the dimensionless multiplier is
$\mathsf a=I$.

Apply the inverse $\mathbb E^{[5]}(0)$ to a forcing with
$f_8=\mathsf m^{-1}X$ and all other components zero.
Its norm in \eqref{I-64-4} is exactly $|X|_y$.
The elimination in \cref{I-63-67} gives
$A=HX$ and $\mathsf m^3\omega=SX$.
The corresponding solution norms in \eqref{I-64-3} dominate
$\|HX\|_y$ and $|SX|_y$: the actual seminorm uses only the
covariant differences already present there.
At the multiplier $\mathsf j_*$, the same forcing gives $A=KX$.
Thus \eqref{eq:regular-system-decay} proves
\[
 \|HX\|_y,\ \|KX\|_y,\ |SX|_y
 \le B_{\mathrm{reg}}
            \sup_z e^{-\delta_{\mathrm{reg}}d(y,z)}|X|_z.
\]
\Cref{prop:projected-propagator} gives the corresponding bound
for $\mathfrak G$ with constant $B_{\mathrm{pr}}$.
Multiplication by an admissible weight proves their weighted
bounds, since $\theta_0\le\delta_{\mathrm{reg}}$ and
$B_{\mathrm{reg}}\le B_{\mathrm{pr}}\le\hat{\mathsf b}$.
Also $|(S-I)X|_w\le2\hat{\mathsf b}|X|_w$.

For the transpose, repeat the weighted-pairing calculation
\eqref{eq:bg-transpose-test}--\eqref{eq:bg-multiplier-bound} with
$H=\mathsf H\mathsf m^{-1}$, row weight $\xi_c^{-1}$ and the
current bound $\gamma_0$ replaced at each top by $|f|_y^\sharp$.
The preceding column estimate then gives
\begin{equation}\label{eq:actual-response-transpose}
 |H^{\mathsf T}f|_{y'}
 \le6B_{\mathrm{reg}}\sum_y\frac{\xi_{y'}}{\xi_y}
           e^{-\delta_{\mathrm{reg}}d(y,y')}|f|_y^\sharp.
\end{equation}
By \eqref{eq:global-scale-ratio} and \eqref{eq:bg-thresholds},
the scale ratio costs at most
$Le^{\delta_{\mathrm{reg}}d(y,y')/4}$.
The weight costs at most $e^{\theta_0d(y,y')}$, with
$\theta_0=\delta_{\mathrm{reg}}/2$, leaving decay at rate
$\delta_{\mathrm{reg}}/4$. The sum bound
\eqref{I-13-122} therefore yields
\[
 |H^{\mathsf T}f|_w
 \le6LB_{\mathrm{reg}}c_S(\delta_{\mathrm{reg}}/4)|f|_w^\sharp
 \le C_{\mathrm{bg}}|f|_w^\sharp
 \le8\hat{\mathsf b}|f|_w^\sharp.
\]

\Cref{I-87-35} now gives the asserted solution and bounds;
\cref{I-86-52} also gives holomorphy of its first component
$X=\mathsf D(KB'+\mathbf A_0)$.
The argument of $\mathsf D$ has supremum at most
$2\hat{\mathsf b}\varepsilon_B\le\varepsilon_3/2$,
so it lies in that map's holomorphic domain.
Finally $X=\mathsf C(\mathbf H)$ and
$\mathbf H=KB'+\mathbf A_0-HX$ imply
$Q\mathbf H=B'+Q\mathbf A_0-X$.
Adding $\mathsf C(\mathbf H)=X$ proves \eqref{I-88-66}.
\end{proof}


\section{Constrained critical response}\label{sec:critical-response}

We now identify the response as a constrained critical configuration.
The fixed-point equation first gives stationarity on the gauge slice
at fixed straightened average.
Gauge invariance removes the remaining multiplier in the gauge
direction. A second use of gauge invariance converts this equation
for the straightened average into criticality for the original
block average.

Use the geometry, regular-pair hypotheses and constants of
\cref{I-88-53}. The constraint in this section uses the reduced
rows $\mathbf B^\flat$. Let $M_{\mathbf B^\flat}$ be the family
of actual iterated averages on these rows, and let
$C(\mathbf B^\flat,V)$ be its real fiber.
The pair is $(U_0,J)$, where
$U_0$ is $G$-valued and
\begin{equation}\label{eq:critical-background-hypothesis}
 J=F'(0),\qquad J=Q^{\mathsf T}\lambda_0
                   \quad\text{for some }\lambda_0\in\mathcal X.
\end{equation}
Here $F$ is the inner-plaquette action of
\cref{sec:actual-local-nonlinear-maps}; thus its gradient includes
the incomplete plaquette stars at the boundary. Membership of
$(U_0,J)$ in the stated regular-pair class, including its current
bound, remains an assumption. Keep $H,K,S,\mathfrak G$ from
\cref{sec:actual-response}, and write
$\Pi=I-D\mathsf V$ for the field projection $\mathsf S$ in
\eqref{eq:bg-projection-operators}.

With the constants of \cref{lem:one-step-straightening}, set
\begin{equation}\label{eq:critical-response-radius}
 r_{\mathrm b}=\min\{r_{\mathrm{st}}/2,
                 [16(1+\kappa_{\mathrm{st}})]^{-1}\},\qquad
 \varepsilon_{\mathrm{crit}}=\min\{\varepsilon_B,r_{\mathrm b}/2\}.
\end{equation}
These radii depend only on $L$.

\begin{theorem}\label{I-88-93}
Under these hypotheses the response $T(B')=\mathbf H(B')$
satisfies, for
$\sup|B'|<\varepsilon_{\mathrm{crit}}$,
\begin{equation}\label{I-88-94}
 \mathcal Q(T)=B',\qquad \mathsf RD^{\mathsf T}T=0,
 \qquad F'(T)=(Q+\mathsf C'_T)^{\mathsf T}\lambda(B').
\end{equation}
The multiplier is uniquely determined and holomorphic, with
\begin{equation}\label{eq:critical-response-multiplier}
 \lambda(B')=(I+H^{\mathsf T}\mathsf C'_T{}^{\mathsf T})^{-1}
                         H^{\mathsf T}F'(T).
\end{equation}
For real $B'$, both $T$ and $\lambda$ are real, and
$E=e^{i\xi T}U_0$ is a critical configuration of
$C(\mathbf B^\flat,M_{\mathbf B^\flat}(E))$.
\end{theorem}

The last assertion concerns the actual block averages. Their
relation to $B'$ involves the endpoint factors in
\eqref{I-88-80} below; these factors are needed at outward rows.

\subsection{Gauge identities for the linear operators}

We begin with the scalar identities underlying the gauge slice.
Let $\mathcal N=\ker Q'$ in the space of bare site fields.
The formula for $\mathsf P$ in \eqref{eq:bg-projection-operators}
gives $\mathsf P^2=\mathsf P$, $\mathsf R^2=\mathsf R$ and
$Q'\mathsf G'\mathsf R=0$. These operators are symmetric.
For $\chi\in\mathcal N$ one has
$\Delta'_a\chi=\Delta^\circ\chi$, and consequently
\[
 \mathsf G'\Delta^\circ\chi=\chi,\qquad
 \mathsf R\Delta^\circ\chi=\Delta^\circ\chi.
\]
Conversely, $\chi_f=\mathsf G'\mathsf Rf$ lies in $\mathcal N$
and $\Delta^\circ\chi_f=\mathsf Rf$. Thus
\begin{equation}\label{eq:critical-scalar-projection}
 \operatorname{ran}\mathsf R=\Delta^\circ\mathcal N,
 \qquad \Delta^\circ\mathsf G'\mathsf R=\mathsf R.
\end{equation}
The boundary term in the scalar Laplacian is
\[
 (\Delta^\circ-D^{\mathsf T}D)\chi(x)
                 =\xi^{-2}(6-n(x))\chi(x),
\]
where $n(x)$ counts the inner neighbors of $x$. A site with
a missing neighbor lies in $\Lambda_0$, where $Q'$ is the
identity. Hence this term vanishes for $\chi\in\mathcal N$.
It follows that $\mathsf VD\chi=\chi$ and $\Pi D\chi=0$ on
$\mathcal N$. Since $\mathsf V$ takes values in $\mathcal N$,
$\mathsf VD\mathsf V=\mathsf V$, so $\Pi^2=\Pi$.
Also $\Pi A=A$ whenever $\mathsf RD^{\mathsf T}A=0$.

The other identity we need is
\begin{equation}\label{eq:critical-row-gradient}
                         QD\chi=0\qquad(\chi\in\mathcal N).
\end{equation}
To prove it, consider one averaging step with reference $W$,
in units in which the fine spacing is one. For
$h_s=e^{is\rho}$ write $W^{h_s}=e^{ia_s}W$. The logarithm
gives $a_s=-s\mathsf d_W\rho+O(s^2)$. The first formula in
\eqref{eq:straightening-maps} and contour covariance give
\[
 \zeta(W,a_s)(y,x)
 =s\{\rho(y)-\operatorname{Ad}(W(y,x))\rho(x)\}+O(s^2).
\]
Therefore the induced gauge after one step satisfies
\begin{equation}\label{eq:critical-induced-gauge}
 \mathsf v(W,a_s)(y)^{-1}h_s(y)
                      =\exp(is\mathsf q_W\rho(y)+O(s^2)).
\end{equation}
Substitution into \eqref{eq:straightening-factorization},
followed by differentiation at $s=0$, proves
\[
 \mathbf Q(W)\mathsf d_W\rho
                =\mathsf d_{\mathcal M(W)}\mathsf q_W\rho.
\]
This holds first for real fields and then by complex linearity
for every $\rho$. All maps are defined on the local domains
used in \cref{I-87-84}.
Iterating, with $W_j=\mathcal M^jU_0$ and
$\mu_j=\mathsf q_{W_{j-1}}\cdots\mathsf q_{W_0}\chi$, gives
\[
 (Q_{\mathrm{phys}}D\chi)(c)
   =\xi_c^{-1}\{
      \operatorname{Ad}(W_{n(c)}(c))\mu_{n(c)}(c_+)
                           -\mu_{n(c)}(c_-)\}.
\]
At an interior row both scalar means vanish when $Q'\chi=0$.
At an outward level-$j$ row, the outer endpoint $z$ has
$B(z)\subset\Lambda_{j-1}$. Its level-$j$ mean is an average
of vanishing level-$(j-1)$ means, so it also vanishes.
The local means agree on overlapping towers by the nested
transport formula \eqref{eq:flat-complex-scalar-transport}.
At level zero the assertion follows from $\chi=0$ on
$\Lambda_0$. Multiplying by $\mathsf m$ proves
\eqref{eq:critical-row-gradient}.

Next put $\Delta=F''(0)$ and
$\Delta_\pi=\Delta-\Delta'_\pi[J]$.
For a site field $v$ the gauge coordinate near $(A,v)=(0,0)$ is
\[
 A^v(b)=(i\xi)^{-1}
       \log(e^{iv_-}e^{i\xi A(b)}e^{-iv_+}),
 \qquad v_+=\operatorname{Ad}(U_0(b))v(b_+).
\]
Expanding the three exponentials and then the logarithm to
second order gives
\begin{equation}\label{eq:critical-gauge-expansion}
 \begin{split}
 A^v&=A-Dv+\tfrac12N(A,v)+O((|A|+|v|)^3),\\
 N(A,v)&=i[v_-,A]-i[A,v_+]-i[v_-,Dv].
 \end{split}
\end{equation}
Indeed the quadratic part of
$\log(e^Xe^Ye^Z)$ is
$\tfrac12([X,Y]+[X,Z]+[Y,Z])$; here
$v_+-v_-=\xi Dv$.
Every inner plaquette transforms by conjugation, so
$F(A^v)=F(A)$. Comparing the first two homogeneous terms gives
\[
 (Dv,J)_\xi=0,\qquad
 (A-Dv,\Delta(A-Dv))_\xi
       =(A,\Delta A)_\xi-(N(A,v),J)_\xi.
\]
With $v=\mathsf VA$, the definition
\eqref{eq:bg-current-form} and polarization therefore give
\begin{equation}\label{eq:critical-hessian-projection}
                         \Delta_\pi=\Pi^{\mathsf T}\Delta\Pi.
\end{equation}

For $O=\mathsf G$ or $\mathsf G_1$, its inverse in
\eqref{eq:bg-reduced-operator} is
\[
 O^{-1}=\Delta_\pi-\Pi^{\mathsf T}X_0\Pi
             +D\mathsf RD^{\mathsf T}+Q^{\mathsf T}Q,
\]
where $X_0=0$ or $\Delta^{(2)}[\mathsf j_*]$.
Equations \eqref{eq:critical-row-gradient} and
\eqref{eq:critical-hessian-projection} imply
$O^{-1}D\chi=D\Delta^\circ\chi$ for $\chi\in\mathcal N$.
Take $\chi=\mathsf G'\mathsf Rf$ to obtain
$OD\mathsf R=D\mathsf G'\mathsf R$. Applying $Q$ and
$\mathsf RD^{\mathsf T}$, and transposing the first identity,
now yields
\begin{equation}\label{eq:critical-linear-identities}
 \begin{gathered}
 \mathsf RD^{\mathsf T}H=\mathsf RD^{\mathsf T}K=0,\qquad
 Q\mathsf G_1D\mathsf R=0,\qquad
 \mathsf RD^{\mathsf T}\mathsf G_1D\mathsf R=\mathsf R,\\
 Q\mathfrak G=\mathsf RD^{\mathsf T}\mathfrak G=0,
 \qquad \mathfrak GQ^{\mathsf T}=\mathfrak GD\mathsf R=0.
 \end{gathered}
\end{equation}
For the last line, substitute the first line into
$\mathfrak G=\mathsf G_1-\mathsf G_1D\mathsf RD^{\mathsf T}
\mathsf G_1-KQ\mathsf G_1$. This operator is symmetric:
the middle term is symmetric, and
$KQ\mathsf G_1=\mathsf G_1Q^{\mathsf T}
(Q\mathsf G_1Q^{\mathsf T})^{-1}Q\mathsf G_1$ is symmetric.
In particular, with
$\mathcal K=\ker Q\cap\ker(\mathsf RD^{\mathsf T})$, one has
\begin{equation}\label{eq:critical-projected-inverse}
 \mathfrak G\mathsf G_1^{-1}
       =I-\mathsf G_1D\mathsf RD^{\mathsf T}-KQ,
 \qquad \mathfrak G\mathsf G_1^{-1}\delta=\delta
                         \quad(\delta\in\mathcal K).
\end{equation}

\subsection{The action on the slice}

Call $\ker(\mathsf RD^{\mathsf T})$ the slice, and put
\[
 \Delta_1=\mathsf G_1^{-1}-D\mathsf RD^{\mathsf T}-Q^{\mathsf T}Q
             =\Delta_\pi-\Pi^{\mathsf T}
                               \Delta^{(2)}[\mathsf j_*]\Pi.
\]
For $\sup|A'|<\varepsilon_3$, write
$D_0=\mathsf D(A')$, $\mathcal L(A')=A'-HD_0$ and
$\mathcal L'=I-H\mathfrak D(A')$. The first contraction and
$QH=I$ give
\[
 \mathcal Q(\mathcal L(A'))=QA',\qquad
 (Q+\mathsf C'_{\mathcal L(A')})\mathcal L'=Q,
 \qquad \mathsf RD^{\mathsf T}\mathcal L'=\mathsf RD^{\mathsf T}.
\]
Differentiating the lower Lipschitz bound in
\eqref{eq:response-first-differences} shows that
$\sup|\mathcal L'\delta|\ge\tfrac12\sup|\delta|$.
Thus $\mathcal L'$ is invertible and maps $\mathcal K$ onto
\[
 \mathcal K_T=\ker(Q+\mathsf C'_T)\cap
                         \ker(\mathsf RD^{\mathsf T}),
 \qquad T=\mathcal L(A').
\]
It also maps the slice onto itself, and $\mathcal L$ preserves
the slice.

Let $\mathsf C^{(2)}(A)=\tfrac12\mathsf C''[A]A$ be the
homogeneous quadratic term of $\mathsf C$. In terms of the
physical coefficient $C^{(2)}$ in \cref{I-84-25},
$\mathsf C^{(2)}=\mathsf m C^{(2)}$ and
$H\mathsf C^{(2)}=\mathsf H C^{(2)}$.
Let $R_F(A)=F(A)-F(0)-(J,A)_\xi-\tfrac12(A,\Delta A)_\xi$,
so that $R_F'=\mathcal V$. From
$\mathsf G^{-1}H=Q^{\mathsf T}S$ and
\eqref{eq:critical-linear-identities} we get
$\Delta_\pi H=Q^{\mathsf T}(S-I)$ and
$H^{\mathsf T}\Delta_\pi H=S-I$.
Define
\begin{equation}\label{eq:critical-reduced-action}
 \begin{split}
 V(A')={}&-(D_0-\mathsf C^{(2)}(A'),H^{\mathsf T}J)_1
              -(QA',(S-I)D_0)_1\\
 &+\tfrac12(D_0,(S-I)D_0)_1+R_F(\mathcal L(A')).
 \end{split}
\end{equation}
Differentiating its four terms gives exactly the six terms
of $\mathsf W$ in \eqref{I-86-62}: the derivatives of
$D_0$, $\mathsf C^{(2)}(A')$ and $\mathcal L(A')$ are
$\mathfrak D$, $\mathsf C''[A']$ and $I-H\mathfrak D$,
respectively, and $S$ is symmetric. Hence $V'=\mathsf W$.

For $A'$ in the slice, both $A'$ and $T=\mathcal L(A')$
are fixed by $\Pi$. Insert $T=A'-HD_0$ into the Taylor
expansion defining $R_F$ and use
\eqref{eq:critical-hessian-projection}. The quadratic term
in $(J,HD_0)_\xi$ is
\[
 (J,H\mathsf C^{(2)}(A'))_\xi
          =\tfrac12(A',\Delta^{(2)}[\mathsf j_*]A')_\xi
\]
by \cref{I-84-25}; it changes $\Delta_\pi$ into $\Delta_1$.
The remaining terms are \eqref{eq:critical-reduced-action}, so
\begin{equation}\label{I-88-75}
 F(\mathcal L(A'))=F(0)+(J,A')_\xi
                         +\tfrac12(A',\Delta_1A')_\xi+V(A').
\end{equation}
Differentiation along the slice gives the identity needed below:
\begin{equation}\label{I-88-76}
 (F'(T),\mathcal L'\delta)_\xi
       =(J+\Delta_1A'+\mathsf W(A'),\delta)_\xi
           \qquad(\mathsf RD^{\mathsf T}\delta=0).
\end{equation}

\subsection{Endpoint factors and gauge directions}

The successive endpoint corrections in
\eqref{eq:straightening-factorization} must be compared at
the two different levels of an outward row. On each local
tower put
\[
 a_0=\xi A,\quad a_{j+1}=\mathcal Q^1(W_j,a_j),\quad
 v_0=1,\quad v_{j+1}(x)=v_j(x)\mathsf v(W_j,a_j)(x).
\]
The local definedness and restriction agreement proved in
\cref{I-87-84} apply for $\sup|A|<c_4/4$.
For $y\in\Lambda_n$ set $v[A](y)=v_n(y)$. Each factor uses
only the nested block below $y$, so this definition is
independent of the incident row. These finite products are
holomorphic and are $G$-valued at real $A$.
For real $A$, iteration of
\eqref{eq:straightening-factorization} and real gauge
covariance gives
\[
 \mathcal M^n(e^{i\xi A}U_0)(c)
   =v_n(c_-)e^{i\mathcal Q(A)(c)}V_0(c)v_n(c_+)^{-1},
 \qquad V_0=M_{\mathbf B^\flat}(U_0).
\]
At an interior row both endpoint factors are $v[A]$.
At an outward level-$n$ row with outer endpoint $z$,
$z\in\Lambda_{n-1}$ and
\[
 v_n(z)=v[A](z)\theta_z(\mathcal Q(A)),\qquad
 \theta_z(\beta)=\mathsf v(W_{n-1},\beta)(z).
\]
The last factor sees only the interior level-$(n-1)$ rows in
$B(z)$, on which $\beta$ is its bond input.

The seam factors and the following twist are defined for
$\sup|\beta|<r_{\mathrm b}$, with $r_{\mathrm b}$ as in
\eqref{eq:critical-response-radius}. The twist is unchanged at interior rows; at an
outward row it is given by the series logarithm of
\begin{equation}\label{eq:critical-seam-twist}
 (e^{i\beta^\natural}V_0)(c)=
 \begin{cases}
 e^{i\beta(c)}V_0(c)\theta_z(\beta)^{-1},&z=c_+,\\
 \theta_z(\beta)e^{i\beta(c)}V_0(c),&z=c_-.
 \end{cases}
\end{equation}
To verify the domain, use the local unitary gauges of
\cref{I-87-84}. There $U_0^u=e^{i\xi A_\kappa}$, with
$\xi A_\kappa=i^{-1}\log(U_0^u)$ bondwise. The principal
logarithm of a unitary matrix is skew-Hermitian, so $A_\kappa$ is real. The background recursion in
\cref{prop:flat-complex-averaging} has real fields and unitary
endpoint factors. Its identification with the actual averages
puts $W_{n-1}$, in the same local gauge, in the form $e^{iB_{n-1}}$
with $B_{n-1}$ real and $\|B_{n-1}\|_\infty<r_{\mathrm{st}}/16$.
Thus \cref{lem:one-step-straightening} gives
$\theta_z=e^{i\zeta_z}$ and
$|\zeta_z(\beta)|\le\kappa_{\mathrm{st}}\sup|\beta|$ for
$\sup|\beta|<r_{\mathrm{st}}$; unitary conjugation preserves
this bound.
Each matrix whose logarithm occurs in
\eqref{eq:critical-seam-twist}, after multiplication by
$V_0(c)^{-1}$, is a product of two exponentials whose
exponent norms sum to less than $1/16$. Its logarithm is
therefore defined, and $|\beta^\natural(c)|<1/8$.
The twist is holomorphic and real at real $\beta$.

Write $\mathsf T^\natural_\beta=d_\beta(\beta^\natural)$.
This derivative is invertible throughout $\sup|\beta|<r_{\mathrm b}$.
Indeed, freeze $\zeta_z$ in each outward-row logarithm.
Right or left multiplication by the opposite exponential
gives a local inverse in that row variable; all the intervening
logarithms remain near the identity. Thus the row-diagonal
part $T_{\mathrm d}$ of the derivative is invertible.
The remaining part $T_{\mathrm u}$ depends only on interior
rows and takes values only at outward rows. Hence
$(T_{\mathrm d}^{-1}T_{\mathrm u})^2=0$, and
\[
 (\mathsf T^\natural_\beta)^{-1}
       =(I-T_{\mathrm d}^{-1}T_{\mathrm u})T_{\mathrm d}^{-1}
\]
is holomorphic. For real $A$ with $\sup|A|<c_4/4$ and
$\sup|\mathcal Q(A)|<r_{\mathrm b}$, substitution of the seam
factors into the iterated average proves
\begin{equation}\label{I-88-80}
 M_{\mathbf B^\flat}(e^{i\xi A}U_0)
                =(e^{i\mathcal Q(A)^\natural}V_0)^{v[A]}.
\end{equation}
The left-hand side is the original real block average.

For real $A$ near zero and $w_\tau=e^{i\tau\chi}$, let
$A^{w_\tau}$ be the local logarithmic coordinate determined by
$e^{i\xi A^{w_\tau}}U_0=(e^{i\xi A}U_0)^{w_\tau}$, and set
\[
 X_\chi(A)=\partial_\tau A^{w_\tau}|_0,\qquad
 \varrho(A;\chi)(y)=i^{-1}\partial_\tau
 [v[A^{w_\tau}](y)^{-1}w_\tau(y)v[A](y)]\big|_0.
\]
These maps are real-linear in $\chi$ and real-analytic in $A$.
They satisfy
\begin{equation}\label{eq:critical-gauge-origin}
                X_\chi(0)=-D\chi,\qquad \varrho(0;\chi)=Q'\chi.
\end{equation}
The first identity follows from the logarithm at the identity.
For the second, the induced gauge at level $j$ is
$v_j[A^{w_\tau}]^{-1}w_\tau v_j[A]$, by the recursion for
the endpoint factors. At $A=0$, apply
\eqref{eq:critical-induced-gauge} at each successive level:
its first-order coefficient is carried from $\rho$ to
$\mathsf q_{W_j}\rho$. Iteration gives $Q'\chi$ at each
layer node by \eqref{eq:flat-averaging-products}; at level
zero the node factor is $1$ and the derivative is $\chi$.

For a row $c$ and $|q|<1/4$, define the linear map
\[
 \ell_{c,q}[\sigma_-,\sigma_+]
   =\left.\frac{d}{d\tau}\right|_0 i^{-1}\log
       [e^{i\tau\sigma_-}e^{iq}V_0(c)
                         e^{-i\tau\sigma_+}V_0(c)^{-1}].
\]
It is holomorphic in $q$. Put $K_A=Q+\mathsf C'_A$.
Gauge covariance of the actual average and
\eqref{I-88-80} show that the twisted datum transforms by
$\mathcal R=v[A^w]^{-1}wv[A]$. Differentiating gives
\begin{equation}\label{eq:critical-gauge-law}
 (\mathsf T^\natural_{\mathcal Q(A)}K_AX_\chi(A))(c)
  =\ell_{c,\mathcal Q(A)^\natural(c)}
           [\varrho(A;\chi)(c_-),\varrho(A;\chi)(c_+)].
\end{equation}
If $\varrho(A;\chi)=0$, invertibility of the twist derivative
gives $K_AX_\chi(A)=0$. Plaquette gauge invariance gives
$(F'(A),X_\chi(A))_\xi=0$.

The map
\begin{equation}\label{eq:critical-gauge-bijection}
 \chi\longmapsto
       (\varrho(A;\chi),\mathsf RD^{\mathsf T}X_\chi(A))
\end{equation}
is a bijection from the real bare-site-field space
$\mathcal S_{\mathbb R}$ onto the product of the real fields
on $\Lambda_*$ and $\mathsf R\mathcal S_{\mathbb R}$, for $A$
in a neighborhood of zero that may depend on the finite data.
At zero it is $(Q'\chi,-\mathsf RD^{\mathsf T}D\chi)$.
Its kernel is zero by \eqref{eq:critical-scalar-projection}
and $\mathsf G'\Delta^\circ\chi=\chi$ on $\mathcal N$.
The map $Q'$ is onto: its transported means are over disjoint
blocks, and one can prescribe each mean by setting a single
site value in that block. Moreover
$\operatorname{ran}\mathsf R=\Delta^\circ\mathcal N$, with
$\Delta^\circ$ injective on $\mathcal N$.
Thus the domain and target have the same real dimension.
The determinant at zero is nonzero, and remains nonzero in
a neighborhood by continuity. No uniform size of that
neighborhood is needed.

\subsection{Stationarity and the original constraint}

\begin{proof}[Proof of \cref{I-88-93}]
The response bounds give
$\|T\|\le4\hat{\mathsf b}\varepsilon_B\le\varepsilon_3\le c_4/8$.
All coefficients in the construction are real at $(U_0,J)$.
Indeed the contour maps and the action are real on real
fields, so their Taylor coefficients are real. The scalar
operators, their inverses and $\mathsf R$ are then real;
the current form, the two successive inverse operators and
their row Gram inverses are real as well. Transposes preserve
reality for the real pairings. Thus $H,K,S,\mathfrak G$ are
real. Both contractions, started at zero, have real iterates
for real data, so their unique limits $T$ and $\mathbf A_0$
are real.

Put $A'=KB'+\mathbf A_0(B')$. Since
$\mathbf A_0=-\mathfrak G\mathsf W(A')$,
\eqref{eq:critical-linear-identities} gives
$\mathbf A_0\in\mathcal K$, $QA'=B'$ and
$\mathsf RD^{\mathsf T}A'=0$.
Also $T=\mathcal L(A')$ and
$\sup|A'|\le2\hat{\mathsf b}\varepsilon_B\le\varepsilon_3/2$.
The first two equations of \eqref{I-88-94} follow from the
identities for $\mathcal L$.

For $\delta\in\mathcal K$, the two constraint terms in
$\Delta_1$ vanish when paired with $\delta$, and
$\mathsf G_1^{-1}KB'$ lies in $\operatorname{ran}Q^{\mathsf T}$.
By symmetry and \eqref{eq:critical-projected-inverse},
\[
 (\delta,\Delta_1A')_\xi
   =(\delta,\mathsf G_1^{-1}\mathbf A_0)_\xi
   =-(\mathfrak G\mathsf G_1^{-1}\delta,\mathsf W(A'))_\xi
   =-(\delta,\mathsf W(A'))_\xi.
\]
Also $(J,\delta)_\xi=0$ by
\eqref{eq:critical-background-hypothesis}.
Equation \eqref{I-88-76} and
$\mathcal L'\mathcal K=\mathcal K_T$ show that $F'(T)$
annihilates $\mathcal K_T$. The annihilator of this kernel
is the range of the transposed constraint map, so
\begin{equation}\label{I-88-95}
               F'(T)=K_T^{\mathsf T}\widetilde\lambda
                                     +D\mathsf R\widetilde\phi.
\end{equation}
Multiplication by $H^{\mathsf T}$ identifies
$\widetilde\lambda$ with \eqref{eq:critical-response-multiplier}.
The inverse there exists by the same Neumann estimate as
\eqref{I-86-59}: its perturbation has norm at most
$8\hat{\mathsf b}C_3\sup|T|\le8C_3c_W\le1/4$.
It is holomorphic and real at real data. In particular
$K_T^{\mathsf T}$ is injective and the asserted multiplier,
once stationarity is proved, is unique.

Define the holomorphic residual
$z(B')=F'(T)-K_T^{\mathsf T}\lambda(B')$.
By \eqref{I-88-95} it lies in $\operatorname{ran}D\mathsf R$.
For real $B'$ it is real, hence $z=D\mathsf R\phi$ for a
real $\phi$. For sufficiently small real $B'$, $T(B')$ lies
in the neighborhood of \eqref{eq:critical-gauge-bijection}.
Choose $\chi$ there with
\[
 \varrho(T;\chi)=0,\qquad
 \mathsf RD^{\mathsf T}X_\chi(T)=\mathsf R\phi.
\]
The gauge law and action invariance give
$(z,X_\chi(T))_\xi=0$. On the other hand,
\[
 (z,X_\chi(T))_\xi
   =(\mathsf R\phi,\mathsf RD^{\mathsf T}X_\chi(T))_\xi
   =(\mathsf R\phi,\mathsf R\phi)_\xi.
\]
The real site pairing is positive definite, so $z=0$.
A holomorphic function vanishing on the real points of a
ball vanishes on its complexification: its Taylor coefficients
are obtained by real difference quotients. Applying the
one-variable identity theorem to $\zeta\mapsto z(\zeta B')$
then gives $z(B')=0$ throughout the chosen ball. This proves
the third equation of \eqref{I-88-94}.

We next prove criticality with the original block average fixed.
For a complex
site field $\sigma$ on $\Lambda_*$ put
\begin{equation}\label{eq:critical-coarse-ward}
 Z(B';\sigma)=\left(\lambda(B'),
   (\mathsf T^\natural_{B'})^{-1}
    \bigl(\ell_{c,B'^\natural(c)}
                 [\sigma(c_-),\sigma(c_+)]\bigr)_c\right)_1.
\end{equation}
This is holomorphic in $B'$ and linear in $\sigma$.
For small real $B'$, the first component of the bijection
\eqref{eq:critical-gauge-bijection} runs over every real
$\sigma$. Equations \eqref{eq:critical-gauge-law} and
\eqref{I-88-94}, with action invariance, imply
\[
 Z(B';\varrho(T;\chi))
       =(\lambda,K_TX_\chi(T))_1
       =(F'(T),X_\chi(T))_\xi=0.
\]
Complex linearity in $\sigma$ and the preceding continuation
argument show that $Z(B';\sigma)=0$ for every $\sigma$ on
the entire chosen ball.

Now fix real $B'$ and a differentiable curve $E_\tau$ in
$C(\mathbf B^\flat,M_{\mathbf B^\flat}(E))$ with
$E_0=E=e^{i\xi T}U_0$. For small $\tau$, its local logarithm
defines a differentiable real curve $A_\tau$ with
$E_\tau=e^{i\xi A_\tau}U_0$ and $A_0=T$.
The strict domain bounds above let us apply
\eqref{I-88-80} along the curve. Constancy of the actual
averages gives
\[
 e^{i\mathcal Q(A_\tau)^\natural}V_0
     =(e^{iB'^\natural}V_0)^{h_\tau},\qquad
 h_\tau=v[A_\tau]^{-1}v[T].
\]
Set $\sigma=i^{-1}\dot h_0$. Differentiating its local
logarithm, just as in \eqref{eq:critical-gauge-law}, yields
\[
 K_T\dot A_0=(\mathsf T^\natural_{B'})^{-1}
       \bigl(\ell_{c,B'^\natural(c)}
                 [\sigma(c_-),\sigma(c_+)]\bigr)_c.
\]
Since the physical action is $\xi F$, its derivative along
this arbitrary constraint curve is
\[
 \left.\frac{d}{d\tau}\right|_0 A_{\Omega_0}(E_\tau)
      =\xi(F'(T),\dot A_0)_\xi
      =\xi Z(B';\sigma)=0.
\]
This proves the required constrained criticality.
\end{proof}

\section{Compatible data at successive depths}\label{sec:compatible-descent}

We extend a reduced datum through one interface and then repeat the
construction down to the bare lattice. Near an interface the extension
increases the curvature bound by a fixed factor. The separation of the regions
ensures that the next interface still sees the original datum, so this
factor is independent of the number of levels.

Use bare lengths, so the spacing at level $j$ is $L^j$. Let $1\le k\le K$, let $L\ge13$ be odd, and
suppose that $\mathfrak M\ge1,R_1\ge L$ are powers of $L$,
with $R_1\mathfrak M\ge L^2$. Let
$\Omega_0\supset\cdots\supset\Omega_k$ be a decreasing
sequence of regions. Each $\Omega_j$ is a union of
$\mathfrak M$-cubes at level $j$, with $\Omega_j=\emptyset$
when no such cubes exist, and
\[
 \operatorname{dist}_\infty(\Omega_j^c,\Omega_{j+1})
       >R_1\mathfrak M L^j,\qquad 0\le j<k.
\]
Assume also that $\mathbb T^{(k)}$ has at least $6L$ sites in
each coordinate direction. We allow empty regions and put
$\Omega_{-1}=\Omega_0$, $\Omega_{k+1}=\emptyset$.
The complementary region $\Omega_j^c$ is the closure of the
point-set complement; in particular, it contains $\partial\Omega_j$.
Distance to the empty set is infinite. Differences below are
point-set differences.

\subsection{The descent statement}

Write $\Gamma_j^\flat$ for the reduced members of the original
sequence. For $0\le n\le k$, the reduced determining set of
$(\Omega_0,\ldots,\Omega_n)$ is
\[
 \mathbf B^{(n),\flat}
   =\bigcup_{j<n}\Gamma_j^\flat\ \cup\ \Gamma_n^{(n)},\qquad
 \Gamma_n^{(n)}
   =\{c\text{ of level }n:\ c_\pm\in\Omega_0,\
                         \{c_-,c_+\}\cap\Omega_n\ne\emptyset\}.
\]
Indeed, the lower members depend only on
$\Omega_0,\Omega_j,\Omega_{j+1}$, while the last region of the
truncated sequence has no successor. Thus
$\mathbf B^{(k),\flat}=\mathbf B^\flat$, and
$\mathbf B^{(0),\flat}$ consists of all inner bare bonds of
$\Omega_0$. At a descent from $n$ to $n-1$, the newly assigned
bonds are exactly the level-$(n-1)$ bonds with an endpoint in
$\Omega_n$. Separation puts both endpoints of these bonds in
$\Omega_{n-1}$.

For a reduced determining set at depth $n$, denote its member
of level $j$ by $E_j$. A plaquette $p$ of level $j$ is
\emph{tested} if at least one of its bonds belongs to $E_j$,
and every bond either belongs to $E_j$ or has both endpoints
in $\Omega_{j-1}\setminus\Omega_j$. Its holonomy
$\partial V(p)$ is the ordered plaquette product using
$V_j(c)$ in the first case and
$\mathcal M(V_{j-1})(c)$ in the second.
The two cases are disjoint, since a bond of $E_j$ meets
$\Omega_j$.

\begin{definition}\label{I-50-59}
For $\varepsilon>0$ and $C\ge1$, let
$\mathcal D(\mathbf B^\flat;\varepsilon,C)$ be the real data
whose tested holonomies satisfy
$|\partial V(p)-\mathbf1|<C\varepsilon$,
and satisfy the smaller bound
$|\partial V(p)-\mathbf1|<\varepsilon$ whenever, for some
$1\le j\le k$, the tested plaquette has level $j-1$ or
$j$ and
$\operatorname{dist}_\infty(p,\partial\Omega_j)\le6L^j$.
\end{definition}
The same definition applies to every truncated sequence.
In particular, a real datum whose tested holonomies have distance
less than $\varepsilon$ from $\mathbf1$ belongs to this class.

\begin{proposition}[Compatible lower data]\label{I-50-61}
There are $C_\downarrow\ge1$ and $\varepsilon_\downarrow>0$,
depending only on $L$, with the following property.
Let $\varepsilon>0$, $V\in\mathcal D(\mathbf B^\flat;\varepsilon,C_\downarrow)$ and
$C_\downarrow\varepsilon\le\varepsilon_\downarrow$.
There are real data $V^{(n)}$ on
$\mathbf B^{(n),\flat}$, $0\le n\le k$, with
$V^{(k)}=V$, such that:
\begin{enumerate}[label=\textup{(\roman*)}]
\item $V^{(n)}_j=V_j$ on $\Gamma_j^\flat$ for $j\le n$.
\item For $1\le n\le k$ and $c\in\Gamma_n^{(n)}$,
\begin{equation}\label{eq:descent-compatible-data}
 |\mathcal M(V^{(n-1)}_{n-1})(c)
                  V^{(n)}_n(c)^{-1}-\mathbf1|
       \le C_\downarrow\varepsilon.
\end{equation}
\item Every tested holonomy of $V^{(n)}$ is defined and satisfies
$|\partial V^{(n)}(p)-\mathbf1|<C_\downarrow\varepsilon$.
Except for level-$n$ plaquettes with a corner in
$\Omega_{n+1}$, its tested plaquettes are tested for the
original datum and
$\partial V^{(n)}(p)=\partial V(p)$.
\end{enumerate}
Moreover $V^{(n)}\in
\mathcal D(\mathbf B^{(n),\flat};\varepsilon,C_\downarrow)$ for
$1\le n\le k$.
The constants are independent of $k$ and of the number
of changes in the sequence of regions.
\end{proposition}

\subsection{Extension through one interface}

The averages used to define tested holonomies exist for every real datum. A second-case
bond has $j\ge1$. Its endpoints are corners of a plaquette
meeting $\Omega_j$, hence lie within $L^j$ of $\Omega_j$.
Their blocks lie within $3L^j/2<L^{j+1}$ of $\Omega_j$,
hence in $\Omega_{j-1}$ by separation, and are disjoint
from $\Omega_j$ by block membership in a region.
Every fine bond of the double block is therefore an interior
bond of $E_{j-1}$, where $V_{j-1}$ is prescribed.
The definition of the real averaging map applies on this
double block.

For a site $z$ of level $k-1$, let $[z]$ be its block center
at level $k$. A coarse site $y\in\Omega_k$ is called inner
if every coarse site $y'$ with $|y'-y|_\infty=L^k$
also belongs to $\Omega_k$. For a bond or plaquette
$\mathfrak e$ at level $k-1$ or $k$, let
\[
 \varepsilon_V^\sharp(\mathfrak e;s)
   =\max\{|\partial V(p)-\mathbf1|:\ p\text{ tested at level }
          k-1\text{ or }k,\
          \operatorname{dist}_\infty(p,\mathfrak e)\le sL^k\},
 \qquad
 \varepsilon_V^\sharp(\mathfrak e)
      =\varepsilon_V^\sharp(\mathfrak e;3).
\]
An empty maximum is zero.

\begin{lemma}[Extension through the last interface]\label{I-50-33}
The constants in \cref{I-50-61} can be chosen so that the
following one-step construction is possible.
If every tested holonomy of a real datum $V$ on
$\mathbf B^\flat$ satisfies
$|\partial V(p)-\mathbf1|\le\varepsilon_\downarrow$, then
there is a real datum $V^\downarrow$ on
$\mathbf B^{(k-1),\flat}$ such that:
\begin{enumerate}[label=\textup{(\roman*)}]
\item $V^\downarrow_j=V_j$ on $\Gamma_j^\flat$ for $j<k$.
\item For every $c\in\Gamma_k$,
\[
 |\mathcal M(V^\downarrow_{k-1})(c)V_k(c)^{-1}-\mathbf1|
       \le C_\downarrow\varepsilon_V^\sharp(c).
\]
The discrepancy is zero if $c_-,c_+\in\Omega_k$.
\item All tested holonomies of $V^\downarrow$ are defined.
Except for level-$(k-1)$ plaquettes having a corner in
$\Omega_k$, their plaquettes are also tested for $V$, and
their holonomies equal those of $V$.
For a level-$(k-1)$ plaquette $p$ having a corner in
$\Omega_k$,
\[
 |\partial V^\downarrow(p)-\mathbf1|
       \le C_\downarrow\varepsilon_V^\sharp(p).
\]
If every block center of a corner of $p$ is inner, its
holonomy is either $\mathbf1$ or $\partial V_k(p')$ for a
coarse plaquette $p'$ whose four bonds belong to $\Gamma_k$.
\end{enumerate}
\end{lemma}

\begin{proof}
Use level-$(k-1)$ lengths for this proof and put
$\mathsf h=(L-1)/2$. Thus fine bonds have length one and
coarse bonds length $L$.

\emph{Construction.}
Let $\mathsf K$ consist of the coarse sites outside
$\Omega_k$ which share a coarse plaquette with a site of
$\Lambda_k=\Omega_k\cap\mathbb T^{(k)}$, and put
$\mathsf Z=\Lambda_k\cup\mathsf K$.
Each site of $\mathsf K$ lies within $L$ of $\Omega_k$.
Its partition cube, of side $L$, consequently lies within
$3L/2<L^2\le R_1\mathfrak M$ of $\Omega_k$ and hence in
$\Omega_{k-1}$. Block membership then gives
\[
 B(\mathsf K)\subset\Omega_{k-1}\setminus\Omega_k,
 \qquad B(\Lambda_k)=\Omega_k\cap\mathbb T^{(k-1)}.
\]
All fine bonds with endpoints in $B(\mathsf Z)$ belong to
$\Gamma_{k-1}^{(k-1)}$. Those with both endpoints in
$B(\mathsf K)$ are retained interior bonds; the others have
an endpoint in $B(\Lambda_k)$ and are new.
Conversely every new bond has this form: its endpoint in
$\Omega_k$ has center in $\Lambda_k$, and the other center
is the same site or a nearest neighbor and thus belongs
to $\mathsf Z$.

Define a coarse field on the bonds with endpoints in
$\mathsf Z$ by
\[
 \mathsf V(c)=
 \begin{cases}
 V_k(c),&\{c_-,c_+\}\cap\Lambda_k\ne\emptyset,\\
 \mathcal M(V_{k-1})(c),&c_-,c_+\in\mathsf K.
 \end{cases}
\]
Every coarse plaquette meeting $\Lambda_k$ has its corners
in $\mathsf Z$ and is tested, with holonomy
$\partial\mathsf V(p')=\partial V(p')$.
On the fine bonds with endpoints in $B(\mathsf Z)$ put
\[
 N(b)=
 \begin{cases}
 \mathbf1,&[b_-]=[b_+],\\
 \mathsf V(\langle[b_-],[b_+]\rangle),&[b_-]\ne[b_+].
 \end{cases}
\]
On the corresponding sites define
\begin{equation}\label{I-50-31}
 \sigma(z)=
 \begin{cases}
 V_{k-1}([z],z),&[z]\in\mathsf K,\\
 \mathbf1,&z\in B(\Lambda_k).
 \end{cases}
\end{equation}
The contour variables in this formula are defined when
$\varepsilon_\downarrow$ is sufficiently small, by
\cref{I-10-53}: all plaquettes inside an outer block are
tested fine plaquettes. Likewise the small-curvature
formula for $\mathcal M(V_{k-1})$ applies on double blocks
joining two sites of $\mathsf K$, by \cref{I-10-61}.
Preserve $V_j$ on every $\Gamma_j^\flat$, $j<k$, and set
\begin{equation}\label{I-50-32}
 V^\downarrow_{k-1}(b)
       =\sigma(b_-)^{-1}N(b)\sigma(b_+)
\end{equation}
on the new bonds. These prescriptions cover the lower
determining set and respect bond reversal.

For any coarse bond with endpoints in $\mathsf Z$, the
field $N$ is trivial inside its two blocks and constant,
equal to $\mathsf V(c)$, across its corridor. Its
double-block plaquettes have holonomy $\mathbf1$.
Every contour transport inside either block is $\mathbf1$,
and every parallel crossing path has transport
$\mathsf V(c)$. Therefore
\[
 \varpi_c(N)=0,\qquad \mathcal M(N)(c)=\mathsf V(c).
\]
For a fine plaquette $p=p_{\mu\nu}(z)$ meeting $\Omega_k$,
block centers have the form
\[
 [z+ae_\mu+a'e_\nu]
   =[z]+a\alpha Le_\mu+a'\alpha'Le_\nu,
 \qquad a,a'\in\{0,1\},\quad \alpha,\alpha'\in\{0,1\}.
\]
Here $\alpha=1$ precisely when $z_\mu-[z]_\mu=\mathsf h$,
and similarly for $\alpha'$. Thus $\partial N(p)=\mathbf1$
if $\alpha\alpha'=0$. Otherwise it equals
$\partial\mathsf V([p])=\partial V([p])$, where
$[p]=p_{\mu\nu}([z])$ is a tested coarse plaquette meeting
$\Lambda_k$. This also shows that every such fine plaquette
has its corners in $B(\mathsf Z)$.

\emph{Comparison with the prescribed outer field.}
Put $\widetilde N=(V^\downarrow_{k-1})^\sigma$ on these
fine bonds. It equals $N$ on the new bonds.
We use two consequences of \cref{I-10-53,I-10-61}.
For a real field $W$ on a block with
$\varpi_y(W)<\omega_M$, the contour family has diameter
at most $\tfrac34(L-1)^2\varpi_y(W)<1/40$, and
\cref{I-3-28} defines its mean. For a bond inside the block,
\[
 |W(y,b_-)W(b)W(y,b_+)^{-1}-\mathbf1|
          \le2\kappa^{\mathrm{SO}}_1\varpi_y(W).
\]
To see this, reverse $b$ if necessary so that a
coordinate-ordered contour to $b_+$ ends with $b$.
Its prefix is a contour to $b_-$.
Each contour differs from its mean by at most
$\kappa^{\mathrm{SO}}_1\varpi_y(W)$, and unitary
multiplication preserves this distance.

For a corridor bond $b$ of a coarse bond $c$, assume
$\varpi_c(W)<\min\{\omega_M,(\kappa^{\mathrm{SO}}_2)^{-1}\}$.
Then
\[
 |W(c_-,b_-)W(b)W(c_+,b_+)^{-1}-\mathcal M(W)(c)|
       \le\kappa_\downarrow\varpi_c(W),
 \qquad
 \kappa_\downarrow
       =2\kappa^{\mathrm{SO}}_1+\kappa^{\mathrm{SO}}_2+\kappa_3.
\]
Indeed, take $c$ positive, put $x=b_-$,
$x^*=b_+$ and $x^+=x+Le_\mu$. Contours from $c_+$ to
$x^*$ and $x^+$ chosen with direction $\mu$ last have
transports satisfying
$W(\Gamma^*)W([x^*,x^+])=W(\Gamma^+)$.
Replacing them by their means costs at most
$2\kappa^{\mathrm{SO}}_1\varpi_c(W)$. The resulting
transport is $W^{\mathsf m}_xW_{\mathsf m}[c]$, which differs from
$W_{\mathsf m}[c]$ by at
most $\kappa^{\mathrm{SO}}_2\varpi_c(W)$; the latter differs from
$\mathcal M(W)(c)$ by at most
$\kappa_3\varpi_c(W)$.

Apply the first estimate to retained bonds within an outer
block and the second to retained bonds between two outer
blocks. Their plaquettes are tested and lie within $2L$ of
the bond. Consequently, with $K_\downarrow=2+2\kappa_\downarrow$,
\begin{equation}\label{I-50-34}
 |\widetilde N(b)N(b)^{-1}-\mathbf1|
       \le K_\downarrow\varepsilon_V^\sharp(b;2).
\end{equation}
On bonds with both endpoints in $\Omega_k$ the two fields
are equal.

\emph{Averages and plaquettes.}
Fix $c\in\Gamma_k$ and let
$d_c=\sup_{b\in\mathsf P_c}
|\widetilde N(b)N(b)^{-1}-\mathbf1|$.
Choosing $\varepsilon_\downarrow$ sufficiently small
allows us to write
$\widetilde N=e^{ia}N$ on $\mathsf P_c$, with $a$ real and
\[
 \|a\|_\infty\le(\pi/2)d_c\le r_{\mathrm{st}}/2.
\]
To estimate its average, apply the box gauge of \cref{I-10-51}
to the flat field $N$ on this double block. Its sides are less
than half the fine torus period, so there is a real gauge $h$
with $N^h=\mathbf1$. On positive bonds
$\widetilde N^h=e^{i\operatorname{Ad}(h)a}$, and conjugation
preserves $\|a\|_\infty$.
Apply \cref{lem:one-step-straightening} at the identity and
return by \eqref{eq:straightening-covariance}. Thus
$v_\pm=\mathsf v(N,a)(c_\pm)$ lie in $G$,
$Q=\mathcal Q^1(N,a)(c)$ lies in $\mathfrak g$, and
$|i^{-1}\log v_\pm|,|Q|\le\kappa_{\mathrm{st}}\|a\|_\infty$.
Unitary conjugation preserves the logarithm domains. Also
$\varpi_c(\widetilde N)\le4\|a\|_\infty
<\min\{\omega_M,(\kappa^{\mathrm{SO}}_2)^{-1}\}$,
so the averages agree with the real small-curvature branch.
The factorization \eqref{eq:straightening-factorization} now gives
\[
 \mathcal M(\widetilde N)(c)
   =v_-e^{iQ}\mathcal M(N)(c)v_+^{-1},\qquad
 |\mathcal M(\widetilde N)(c)-\mathcal M(N)(c)|
   \le3\kappa_{\mathrm{st}}\|a\|_\infty.
\]
The last estimate follows by telescoping the three unitary
perturbations. Reversing $c$ inverts both averages and leaves
their distance unchanged. Since $\sigma=1$ at coarse sites, using
$V_{k-1}(y,y)=\mathbf1$ on $\mathsf K$, gauge covariance
identifies this discrepancy with that for $V^\downarrow$.
Every fine bond of $\mathsf P_c$ lies within $\mathsf h$
of $c$; hence \eqref{I-50-34} gives
$d_c\le K_\downarrow\varepsilon_V^\sharp(c)$.
This proves (ii) with
$C_\downarrow\ge(3\pi/2)\kappa_{\mathrm{st}}K_\downarrow$.
If $c_\pm\in\Omega_k$, both endpoint blocks lie in
$\Omega_k$ and $\widetilde N=N$ on $\mathsf P_c$, giving
the asserted equality.

If $p$ is a fine plaquette with a corner in $\Omega_k$,
its four bonds belong to the new highest member and
$\partial V^\downarrow(p)$ is conjugate to
$\partial\widetilde N(p)$. Telescoping the four unitary
factors gives
\[
 |\partial V^\downarrow(p)-\mathbf1|
 \le4K_\downarrow\varepsilon_V^\sharp(p;2)
                   +|\partial N(p)-\mathbf1|
 \le(4K_\downarrow+1)\varepsilon_V^\sharp(p).
\]
For the last inequality, a nontrivial $\partial N(p)$ is
the tested coarse holonomy of $[p]$, which lies within
$\mathsf h$ of $p$. If every corner of $p$ has inner
block center, all its corners lie in $B(\Lambda_k)$;
there $\sigma=1$ and its bonds are new, so its holonomy
equals $\partial N(p)$. In this case all four bonds
of $[p]$, when it occurs, belong to $\Gamma_k$.

At levels below $k-1$, the members and the regions entering
the definition of a tested plaquette are unchanged.
At level $k-1$, a plaquette without a corner in $\Omega_k$
contains no new bond; its first-case bonds and second-case
bonds are also unchanged. On a second-case bond its
holonomy uses an average of interior bonds of the preceding
member, all of whose values were preserved. Thus these
holonomies equal those of $V$. Every remaining tested
plaquette is a fine plaquette meeting $\Omega_k$, already
treated above. This proves (iii). All the smallness requirements
are satisfied by choosing
\[
 \begin{split}
 C_\downarrow&=\max\{1+4K_\downarrow,
                    (3\pi/2)\kappa_{\mathrm{st}}K_\downarrow\},\\
 0<\varepsilon_\downarrow&<
 \min\{\omega_M,(\kappa^{\mathrm{SO}}_2)^{-1},
              r_{\mathrm{st}}/(\pi K_\downarrow)\}.
 \end{split}
\]
These choices depend only on $L$.
\end{proof}

\subsection{Iteration without loss of radius}

\begin{proof}[Proof of \cref{I-50-61}]
We retain the equality assertion in (iii) while applying
\cref{I-50-33} successively. Separation gives
\[
 \operatorname{dist}_\infty(\partial\Omega_j,\Omega_{j+1})
       >L^{j+2}.
\]
All segments used below may be taken in local Euclidean
coordinates: the torus side is at least
$6L^{k+1}\ge78L^n$, so sup balls of radius $7L^n$
are isometric to their Euclidean lifts.

Suppose that $V^{(n)}$ has been constructed with (i) and
(iii), starting with $V^{(k)}=V$. Its tested holonomies have distance from $\mathbf1$ less than
$C_\downarrow\varepsilon\le\varepsilon_\downarrow$, so set
$V^{(n-1)}=(V^{(n)})^\downarrow$.
The truncated sequence satisfies the same geometric
hypotheses.

We first bound the local sizes in the nonexact cases of
the one-step lemma. Let $\mathfrak e$ be either a bond
of $\Gamma_n^{(n)}$ with an endpoint that is not inner,
or a level-$(n-1)$ plaquette meeting $\Omega_n$ with a
corner whose block center is not inner.
There are $z_0\in\mathfrak e\cap\Omega_n$ and a coarse
site $y'\notin\Omega_n$ with
$|z_0-y'|_\infty\le2L^n$.
For a bond this follows by taking an outside site at
distance at most $L^n$ from its noninner endpoint.
For a plaquette, take such a site within $L^n$ of the
noninner block center; the additional distance to $z_0$
is at most $\tfrac{L+1}{2}L^{n-1}<L^n$.
The segment from $z_0$ to $y'$ meets
$\partial\Omega_n$ at a point $q$ with
$|q-z_0|_\infty\le2L^n$.

A plaquette $p$ entering
$\varepsilon_{V^{(n)}}^\sharp(\mathfrak e)$ lies within
$3L^n$ of $\mathfrak e$, and every point of
$\mathfrak e$ lies within $L^n$ of $z_0$. Hence
\[
 \operatorname{dist}_\infty(p,\partial\Omega_n)
       \le3L^n+L^n+2L^n=6L^n.
\]
If $p$ has level $n$, all its corners are within
$7L^n<L^{n+2}$ of this boundary and thus outside
$\Omega_{n+1}$. If $p$ has level $n-1$, it is
already outside the exceptional family in (iii).
In both cases (iii) identifies its holonomy with that
of the original datum, and the smaller bound in
\cref{I-50-59}, with $j=n$, applies. Therefore
\begin{equation}\label{I-50-62}
 \varepsilon_{V^{(n)}}^\sharp(\mathfrak e)<\varepsilon.
\end{equation}

The one-step preservation property proves (i) at depth
$n-1$. For a coarse bond with both endpoints inner, the
discrepancy in (ii) vanishes. Every other coarse bond
is covered by \eqref{I-50-62}; the one-step estimate
therefore proves \eqref{eq:descent-compatible-data}.

Consider a level-$(n-1)$ plaquette $p$ meeting
$\Omega_n$. Separation puts all its bonds in
$\Gamma_{n-1}^{(n-1)}$, so it is tested.
If all its block centers are inner, its new holonomy
is identity or a tested holonomy at depth $n$, whose distance
from $\mathbf1$ is less than $C_\downarrow\varepsilon$ by induction.
Otherwise the one-step bound and \eqref{I-50-62}
give the same bound. All other tested holonomies
are unchanged by the one-step lemma. Their levels
are at most $n-1$, so they were not exceptional at
depth $n$ and already equal the original holonomies.
This proves (iii) and completes the induction.

It remains to check the smaller bound for a truncated
datum. A level-$n$ plaquette meeting $\Omega_{n+1}$
satisfies
\[
 \operatorname{dist}_\infty(p,\partial\Omega_n)
       >L^{n+2}-L^n>6L^n.
\]
It therefore cannot be one of the plaquettes tested
for the smaller radius at depth $n$: its level forces
$j=n$ in that condition.
Thus every plaquette tested for the smaller radius has unchanged
original holonomy, by (iii), and inherits its distance bound
$|\partial V^{(n)}(p)-\mathbf1|<\varepsilon$.
\end{proof}

The compatibility estimate applies directly to a configuration
realizing the lower datum. If
$M_{\mathbf B^{(n-1),\flat}}(U)=V^{(n-1)}$, then its
retained rows of every level $j<n$ equal those of
$V^{(n)}$. For $c\in\Gamma_n^{(n)}$, separation puts
the fine double block $\mathsf P_c$ in
$\Gamma_{n-1}^{(n-1)}$, so locality gives
\[
 \mathcal M^nU(c)
       =\mathcal M(V^{(n-1)}_{n-1})(c).
\]
Consequently its remaining datum discrepancy is bounded
by \eqref{eq:descent-compatible-data}.

\section{Existence and uniqueness of the real background}
\label{sec:real-background}

\subsection{Improving the regularity radius}

The following improvement keeps the regularity radius uniform
when a new level is added to the region sequence.
We continue to use the reduced determining set $\mathbf B^\flat$.

\begin{proposition}[Improvement of the regularity radius]
\label{I-49-38}
There are $B_{\mathrm{rad}}\ge\max\{4,L^2\}$ and a power
$\rho_{\mathrm{rad}}\ge L^2$ of $L$, depending only on $L$, with the
following property. Let $\rho\ge\rho_{\mathrm{rad}}$ be a power
of $L$, and assume the separated-region setting, with
\[
 \mathfrak M\ge M_{\rm ps},\qquad R_1\ge L,\qquad R_1\mathfrak M\ge L^2\rho,
 \qquad \#\mathbb T^{(k)}\text{ in each direction}\ge6\rho.
\]
There is $a_{\mathrm{rad}}(\rho)>0$, depending only on $L,\rho$,
such that the following holds.
Let $V$ be real reduced data with all tested holonomies
defined and satisfying
$|\partial V(p)-\mathbf1|<\varepsilon_1$, where
$\varepsilon_1>0$. If $U$ is a critical configuration in
$C(\mathbf B^\flat,V)\cap\mathfrak U'_k(\varepsilon_0)$,
where $0<\varepsilon_0\le a_{\mathrm{rad}}(\rho)$, then
\[
                 U\in\mathfrak U'_k(B_{\mathrm{rad}}\varepsilon_1).
\]
The constants are independent of the number of levels and
of the number of strict inclusions between successive regions.
\end{proposition}

\begin{proof}
We prove the improvement
\begin{equation}\label{eq:reset-one-improvement}
 U\in\mathfrak U'_k\!
 \left(\max\{4\varepsilon_1,\,
              3K^\sharp\varepsilon_1+3\theta\varepsilon_0\}\right),
 \qquad
 \theta=C_b\rho^2e^{-\delta\rho/24}+C_c\rho^4\varepsilon_0,
\end{equation}
when $0<\varepsilon_1\le\varepsilon_0$ and $\varepsilon_0$
is sufficiently small as specified below.
Here $\delta=\delta_S$ is the exponent in \eqref{eq:corner-H-local}, and $K^\sharp,C_b,C_c$ depend only on $L$.

Fix an inner plaquette or a full bond $\mathfrak e$ of
positive local index $j$, and choose a corner
$x_0\in\Omega_j\setminus\Omega_{j+1}$.
Put $n=j$ if
$\operatorname{dist}_\infty(x_0,\Omega_{j+1})>5\rho\xi_j$,
and $n=j+1$ otherwise; set
\[
 \xi=\xi_n,\quad \widehat y=[x_0]_n,\quad
 N=\{|x-\widehat y|_\infty\le4\rho\xi\},\quad
 d_\infty(x)=\xi^{-1}|x-\widehat y|_\infty,\quad
 d_1(x)=\xi^{-1}|x-\widehat y|_1 .
\]
Then
\begin{equation}\label{eq:reset-two-layer-cube}
 N\subset\Omega_{n-1},\qquad N\cap\Omega_{n+1}=\varnothing,
 \qquad \operatorname{dist}_\infty(N,\Omega_0^c)>8\rho\xi.
\end{equation}
Indeed, if $n=j$, separation leaves
$(L-\tfrac92)\rho\xi>8\rho\xi$ between $N$ and
$\Omega_{j-1}^c$; the defining distance excludes
$\Omega_{j+1}$.
If $n=j+1$, choose a point of $\Omega_{j+1}$ within
$5\rho\xi_j$ of $x_0$. The remaining separation is
$(L^2-5-\tfrac92L)\rho\xi_j\ge8L\rho\xi_j$.
Since $x_0\notin\Omega_{j+1}$, the next separation excludes
$\Omega_{j+2}$. This also proves $n\le k$.
The torus side, being twice a power of $L$ and at least
$6\rho$, is at least $2L\rho$ in level-$k$ units.
Thus $N$ has a single coordinate chart and side less than
one third of the torus side.

Every site in $N$ belongs to layer $n-1$ or $n$.
The weighted distance therefore satisfies
\begin{equation}\label{eq:reset-distance}
 \begin{aligned}
 d_{\mathbf B}(y,y')&\ge s'-s
   &&\text{if }d_\infty(y)\le s,\ d_\infty(y')\ge s',
          \quad0\le s\le s'\le4\rho,\\
 d_{\mathbf B}(y,y')&\ge\xi^{-1}|y-y'|_1
   &&\text{if }d_\infty(y),d_\infty(y')\le\rho/3.
 \end{aligned}
\end{equation}
For the first inequality, any connecting set crosses the
annulus, where both layer weights are at least $\xi^{-1}$.
For the second, a connecting set contained in $N$ has
weighted length at least its coordinate $\ell^1$ distance.
If it leaves $N$, the first inequality gives length
at least $11\rho/3$, greater than the possible
$\ell^1$ distance $2\rho$ between the two endpoints.
This includes $n=1$, where the finer part is the
zero layer.

\emph{A reference configuration trivial near $\widehat y$.}
Choose a cube $\square$ of the $\rho\xi$ partition
containing $\widehat y$, and let $\square_s$ be its concentric
enlargement of side $(\rho+2s)\xi$.
The regional estimate \eqref{eq:regional-local-gauge}, including its ordinary
curl-divergence conclusion, applies on $\square_\rho$.
Use level $n$ if this cube lies in $\Omega_n$, and level
$n-1$ otherwise. The dimensionless side is $3\rho$ or
$3L\rho$; the required enlarged cube lies in $N$.
After extension of the local gauge by the identity, write
$U^\sharp=U^{u^{-1}}$. Gauge covariance preserves
criticality and the tested-holonomy bounds. On $\square_\rho$,
\[
 U^\sharp=e^{i\eta A_0},\qquad
 \xi|A_0|,\ \xi^2|\nabla^\eta A_0|,\
 \xi^3|D_{\mathbf1}^{\mathsf T}D_{\mathbf1}A_0|
       <a,\qquad a=63L^6B_{\mathrm p}\rho\varepsilon_0 .
\]
The $L^3$ conversion in the third bound is included in
$a$. The local-gauge hypotheses follow from
$63L^3\rho\varepsilon_0\le c_{\mathrm p}$ and from the
chosen lower bound on $\rho$.

Take a product cutoff $\chi$, equal to one on
$\square_{\rho/3}$ and zero off $\square_{2\rho/3}$, with
first and second differences bounded by
$c_\varphi/(\rho\xi)$ and $c_\varphi/(\rho\xi)^2$.
Set $A^\natural=\chi A_0$, $A'=(1-\chi)A_0$ on
$\square_\rho$, extend $A^\natural$ by zero, and define
$U_0=e^{i\eta A'}$ there,
$U_0=U^\sharp$ elsewhere.
The cutoff is zero well before the boundary, and
\[
                         U^\sharp=e^{i\eta A^\natural}U_0 .
\]
The product rule for $D_{\mathbf1}^{\mathsf T}D_{\mathbf1}
(fA_0)$ has, for each transverse direction, one term
$(\nabla f)DA_0$, two terms $(\nabla f)\nabla A_0$,
and two terms $(\nabla^2f)A_0$, besides
$fD_{\mathbf1}^{\mathsf T}D_{\mathbf1}A_0$.
Their bounds are $2,1,1,1,1$ times
$c_\varphi a\xi^{-3}$. Hence, with
$\alpha_2=(1+12c_\varphi)a$,
\[
 \xi|fA_0|\le a,\qquad
 \xi^2|\nabla^\eta(fA_0)|\le(1+c_\varphi)a,\qquad
 \xi^3|D_{\mathbf1}^{\mathsf T}D_{\mathbf1}(fA_0)|
       \le\alpha_2,\qquad f=\chi,1-\chi .
\]
The local real action estimates give curvature and current
at most $3\alpha_2$ in their respective scale norms.
The full matrix divergence is at most
$(1+3\alpha_2)3\alpha_2\le4\alpha_2$.
Outside this calculation's domain a complete plaquette
star agrees with that of $U^\sharp$.
Consequently
\begin{equation}\label{eq:reset-reference}
 U_0\in\mathfrak U'_k(\alpha'_0),\qquad
 \alpha'_0=6\alpha_2,\qquad
 |A^\natural|_{(-1)}\le\alpha_2,\quad
 \|A^\natural\|_X\le2\alpha_2.
\end{equation}
For the last bound the difference between the covariant
and ordinary derivatives is at most $2a^2\xi^{-2}$.
In particular $U_0=\mathbf1$ on $d_\infty\le\rho/3$.
By locality its averages through level $n$ are the identity
on bonds whose endpoints satisfy $d_\infty\le\rho/3-1$.
On bare stars in this smaller region,
\begin{equation}\label{eq:reset-vanishing}
                J_0:=F'(0)=0,\qquad \Delta'=0.
\end{equation}

\emph{Prescribing the local datum gauge.}
Let $\rho'=\lfloor\rho/3\rfloor$,
$Y=\{y\in\mathbb T^{(n)}:d_\infty(y)\le\rho'\}$,
and $Y^\circ=Y\setminus\Omega_n$.
For $W_i=\mathcal M^iU^\sharp$, define
\[
 z(y)=W_n(\Gamma_{\widehat y,y})\quad(y\in Y),\qquad
 z(x)=z(y)W_{n-1}(y,x)\quad
                     (x\in B(y),\ y\in Y^\circ).
\]
The prescriptions agree at the centers, and prescribe
$z$ at tops of levels $n$ and $n-1$:
$Y\cap\Omega_n\subset\Lambda_n$ and
$B(Y^\circ)\subset\Lambda_{n-1}$.
Extend $z$ by the identity at the other tops.

Here are the estimates needed from this construction.
A coarse plaquette meeting $\Omega_n$ has the holonomy
of the datum. For a coarse plaquette outside $\Omega_n$,
the fine plaquettes in its four blocks are tested, so
one averaging step bounds its holonomy by
$L^2\varepsilon_1+C\varepsilon_1^2\le2L^2\varepsilon_1$.
The comb estimate therefore gives
\[
 |W_n^z(c)-\mathbf1|
       \le2L^2\max_{\partial c}d_1\,\varepsilon_1 .
\]
In each fine block $B(y)$, $y\in Y^\circ$, covariance gives
$W_{n-1}^z(y,x)=\mathbf1$ exactly.
Within that block the contour comparison bounds a bond
by $C\varepsilon_1$. Across two blocks the same comparison
gives
\[
 |W_{n-1}^z(b)-\mathbf1|
       \le |W_n^z(c)-\mathbf1|+C\varepsilon_1 ,
\]
where $c$ joins their centers.
To see the latter estimate, compare the two averaged
contour transports with paths passing through $b$.
Moving the crossing step through the transverse path
uses at most $L-1$ plaquettes. Each contributes at most
$\varepsilon_1$; the two contour errors and the
coarse-average error are also $C(L)\varepsilon_1$.
Thus all the fine and coarse bonds just described have
size at most $C(L)(1+d_1)\varepsilon_1$ in this gauge.

The chart bounds before gauging imply that the coarse
bonds of $W_n$ here have size at most
$\kappa_4L^2\varepsilon_0+a$, and the fine bonds inside
one block have size at most $\kappa_4\varepsilon_0+a$.
The comb has at most $3\rho'$ bonds and a fine contour
at most $3(L-1)/2$ bonds. The contour comparison therefore
gives
\[
 \sup|z-\mathbf1|\le
 \zeta:=\rho(\kappa_4L^2\varepsilon_0+a)
       +\tfrac32L(\kappa_4\varepsilon_0+a)
       +2\kappa_1^{\mathrm{SO}}\varepsilon_0
       \le C(L)\rho a .
\]
Apply \cref{I-46-44} to
$A^\natural$, with the values
$r(y)=z(y)v_j[A^\natural](y)$ for $y\in\Lambda_j$.
Its smallness hypotheses follow from
$\sup|r-\mathbf1|\le\zeta+C_v\alpha_2$, where
$C_v=4\kappa_{6,\mathsf m}$.
It gives a real chart $A$ and a gauge $w$ such that
\[
 E=e^{i\eta A}U_0=(U^\sharp)^w,\qquad RD^*A=0,\qquad
 |A|_{(-1)}\le
 a_f:=\alpha_2+
       2C_{\rm gc}K_3(6\alpha_2+\zeta+C_v\alpha_2)
       \le C(L)\rho a .
\]
The induced frame is $z$ at every top.
At an outward row its outer-end correction is also
trivial in the region under consideration: the finer
reference contours are the identity and the
corresponding contours of $W_{n-1}^z$ are exactly
the identity. Thus the one-step factor is $\mathbf1$,
not merely small. The reduced chart row
$\mathcal Q(A)$ consequently satisfies
\begin{equation}\label{eq:reset-row-data}
 |\mathcal Q(A)(c)|\le C(L)(1+d_1(y_c))\varepsilon_1
       \quad\text{if }d_\infty(y_c)\le\rho/3-3,
 \qquad \sup|\mathcal Q(A)|\le b_f:=2La_f .
\end{equation}
The first bound follows by the real logarithm estimate
$|\log W|\le(\pi/2)|W-\mathbf1|$.
Only interior rows at level $n-1$ and interior or outward
rows at level $n$ occur here.

Since $E$ and $U_0$ lie in
$\mathfrak U'_k(\alpha'_0)$ and $RD^*A=0$, \cref{I-47-4} gives, with $B_1=8K_1$,
\begin{equation}\label{eq:reset-global-chart}
 \|A\|_X,\ |D^*DA|_{(-3)}
      \le r_f:=B_1(\alpha'_0+4b_f)
      \le C(L)\rho^2\varepsilon_0 .
\end{equation}
The reference current has $|J_0|_{(-3)}\le3\alpha'_0$:
at nonfull bonds the inner star has at most three
plaquettes and local scale $\eta$.
Criticality and the small-chart transversality give
$F'(A)=K_A^*\omega$.

\emph{Localizing the Euler equation.}
Use the corner frame of the reduced Schur operators,
writing $\widehat q=T^{-1}\mathcal Q(A)$ and
$\widehat\omega=T^{-1}\omega$. The map $T$ is an isometry
row by row. With $\mathcal C$ denoting the nonlinear
mean/corner remainder, the equations become
\[
 \begin{split}
 J_0+D^*DA+\Delta'A+\mathcal V(A)
        &=(Q^*+\mathcal C'(A)^*)\widehat\omega,\\
 QA&=\widehat q-\mathcal C(A).
 \end{split}
\]
Apply $H_0^*$, using $QH_0=1$:
\begin{equation}\label{eq:reset-multiplier}
 \widehat\omega
   =H_0^*J_0+(DH_0)^*DA+
      H_0^*(\Delta'A+\mathcal V(A)
                      -\mathcal C'(A)^*\widehat\omega).
\end{equation}
The action and chart estimates give
\[
 \sup|\mathcal C(A)|\le C r_f^2,\quad
 |\mathcal C'(A)^*\omega|_{(-3)}
       \le C r_f\sup|\omega|,\quad
 |\Delta'A|_{(-3)}+|\mathcal V(A)|_{(-3)}
       \le C r_f^2.
\]
Here $\alpha'_0\le r_f$ has been used.
For $r_f$ sufficiently small, the last multiplier term
in \eqref{eq:reset-multiplier} is absorbed, and
$\sup|\widehat\omega|\le C r_f$.
The Schur identities now give
\begin{equation}\label{eq:reset-schur}
 A=H_0\widehat q-\widetilde G_0J_0+A_2,\qquad
 A_2=\widetilde G_0(-\mathcal V(A)
                    +\mathcal C'(A)^*\widehat\omega)
                         -H_0\mathcal C(A),\qquad
 \|A_2\|_X\le C r_f^2.
\end{equation}

We spell out how decay improves this global bound.
Write $d=d_{\mathbf B}$ and put
\[
 \operatorname{rt}(y,y')=\max\{\xi_y/\xi_{y'},\xi_{y'}/\xi_y\},\qquad
 D_\rho=\rho/6-3,\qquad
 e_\rho=e^{-\delta D_\rho/4}\le2e^{-\delta\rho/24}.
\]
The already proved summation estimate implies
\begin{equation}\label{eq:reset-sums}
 \begin{split}
 \sum_{y'}e^{-\delta d(y,y')}(1+d(y,y'))
                  \operatorname{rt}(y,y')^4&\le C,\\
 \sum_{d(y,y')\ge D_\rho}e^{-\delta d(y,y')/2}
                  \operatorname{rt}(y,y')^4&\le Ce_\rho .
 \end{split}
\end{equation}
Indeed, extract $e^{-\delta D_\rho/4}$ in the second
sum; in the first absorb $1+d$ into three quarters
of the exponential. The remaining sums are the same
uniform top sums as before.

For $d_\infty(y)\le\rho/6$, split the terms in
\eqref{eq:reset-schur} at $d_\infty(y')=\rho/3-3$.
In the near part, \eqref{eq:reset-distance} gives
$1+d_1(y')\le(1+d_1(y))(1+d(y,y'))$.
In the far part the weighted distance is at least
$D_\rho$. The current $J_0$ vanishes in the near part
by \eqref{eq:reset-vanishing}.
The bounds \eqref{eq:corner-H-local} and
\eqref{eq:corner-Gtilde-local}, followed by
\eqref{eq:reset-sums}, yield
\begin{equation}\label{eq:reset-near-A}
 \varsigma_y:=\max\{\xi_y|A|_y,\xi_y^2|\nabla A|_y\}
 \le C(1+d_1(y))\varepsilon_1
       +Ce_\rho r_f+Cr_f^2 .
\end{equation}
The scale ratio in the $\widetilde G_0$ term has power
three, which is covered by \eqref{eq:reset-sums}.

Next evaluate \eqref{eq:reset-multiplier} at a row
with $d_\infty(y_c)\le3$.
Its $H_0^*J_0$ term is a far contribution.
For $(DH_0)^*DA$, split at $d_\infty(y)=\rho/6$,
use $|DA|_y\le2|\nabla A|_y$ and
\eqref{eq:reset-near-A} in the near part, and the
global bound in the far part. The transposed kernels in \eqref{eq:corner-H-transposes}
have the additional factor $\xi_c/\xi_y$, bounded
by $\operatorname{rt}(y_c,y)$; the same sums still apply.
The remaining terms are globally $O(r_f^2)$.
It follows that
\begin{equation}\label{eq:reset-near-multiplier}
 |\widehat\omega(c)|\le C\varepsilon_1+
                         Ce_\rho r_f+Cr_f^2
                 \qquad(d_\infty(y_c)\le3).
\end{equation}

The corners of $\mathfrak e$ and of its star lie over
tops with $d_\infty\le3$, and their local scales differ
from $\mathsf L(\mathfrak e)$ by at most $L$.
If $\mathfrak e=b$ is a bond, every row containing $b$
has level $n-1$ or $n$ and its top satisfies
$d_\infty(y_c)<3$. To check the first assertion, a row
of level $i\le n-2$ would have a top inside
$N\subset\Omega_{n-1}\subset\Omega_{i+1}$.
For $i\ge n+1$, support membership first forces
$i=n+1$, and then places that top in
$N\cap\Omega_{n+1}$. Both are impossible.
By \eqref{eq:reset-vanishing}, the Euler equation at $b$
therefore gives
\[
 D^*DA(b)=(Q^*+\mathcal C'(A)^*)\widehat\omega(b)
                          -\mathcal V(A)(b),
\]
where the first term uses only the rows just located.
The transpose row bound and \eqref{eq:reset-near-multiplier}
control this term. Combining with
\eqref{eq:reset-near-A} gives
\begin{equation}\label{eq:reset-local-size}
 \varepsilon_{\rm loc}
 \le K^\sharp\varepsilon_1+
       C_b\rho^2e^{-\delta\rho/24}\varepsilon_0+
       C_c\rho^4\varepsilon_0^2,\qquad
 \varepsilon_{\rm loc}\le L^2r_f.
\end{equation}
Here $\varepsilon_{\rm loc}$ is the maximum of the
potential and first-difference sizes on the plaquette
or full star, and, in the bond case, of
$\mathsf L(b)^3|D^*DA(b)|$.
The coefficients in \eqref{eq:reset-local-size} depend
only on $L$.

\emph{Returning to gauge-invariant quantities.}
The reference is the identity on the relevant full star.
Choose the small radius so that
$\varepsilon_{\rm loc}\le\min\{(2C_{\rm ec})^{-1},2/75\}$.
The real action estimates then give
\[
 |\partial E(p)-\mathbf1|
      \le\tfrac52\varepsilon_{\rm loc}
                    \eta^2\mathsf L(\mathfrak e)^{-2},
 \qquad
 |J(E)(b)|\le\tfrac52\varepsilon_{\rm loc}
                    \mathsf L(b)^{-3}.
\]
For the second bound, use
$F'(A)=D^*DA+\mathcal V(A)$ on this star and
$|J(E)(b)|\le(1+4\eta|A(b)|)|F'(A)(b)|$.
The full-divergence conversion gives
\[
 |D_E^{\mathsf T}\partial E(b)|
 \le(1+\tfrac52\varepsilon_{\rm loc})
          \tfrac52\varepsilon_{\rm loc}\,
                      \eta^2\mathsf L(b)^{-3}
 \le\tfrac{11}4\varepsilon_{\rm loc}
                      \eta^2\mathsf L(b)^{-3}.
\]
These bounds also hold when $\varepsilon_{\rm loc}=0$,
by passage to arbitrarily small positive radii.
Gauge covariance returns them to $U$.
Since $\varepsilon_1>0$, their strict enlargement by
$3(K^\sharp\varepsilon_1+\theta\varepsilon_0)$ proves
the required bounds at every positive-index element.

At index zero, a plaquette is a tested plaquette of
$V_0$, so its error is less than $\varepsilon_1$.
A full bond of index zero has four plaquettes, of
indices zero or one. In
$S_U=\eta D_U^{\mathsf T}\partial U$ the identity matrices
cancel between each forward and backward plaquette.
Thus
\[
 |S_U(b)|<
 4\max\{\varepsilon_1,\,
       3(K^\sharp\varepsilon_1+\theta\varepsilon_0)L^{-2}\}
 \le\max\{4\varepsilon_1,\,
       3K^\sharp\varepsilon_1+3\theta\varepsilon_0\}.
\]
This proves \eqref{eq:reset-one-improvement}.

All smallness conditions used above follow from
\[
 \varepsilon_0\le
 \min\left\{L^{-2}c_{\rm curv},\
       \frac{c_{\mathrm p}}{63L^3\rho},\
       \frac{m_*}{63L^6B_{\mathrm p}\rho^2}\right\}
\]
with a sufficiently small $m_*(L)>0$.
Here $c_{\rm curv}>0$ is no larger than $a_{\mathsf m}$,
$c_{\mathsf c}$ and the real averaging curvature threshold
in \cref{I-10-101}.
More explicitly, choose $m_*$ so that $\alpha_2$ satisfies
the real-action caps; $\alpha'_0$, $2\alpha_2$ and
$\zeta+C_v\alpha_2$ satisfy the chart, gauge and slice
caps; $a_f,r_f$ lie in those same fixed chart radii;
the multiplier coefficient is at most $1/2$; and
$L^2r_f\le\min\{(2C_{\rm ec})^{-1},2/75\}$.
Every quantity in this list is bounded by a fixed
multiple of $\rho a$, so this is one positive choice
depending only on $L$.

Choose $\rho_{\mathrm{rad}}$ large enough for the preceding
geometry, the local gauge and Schur estimates, and for
$C_b\rho^2e^{-\delta\rho/24}\le1/24$ at every power
$\rho\ge\rho_{\mathrm{rad}}$ of $L$. Define
\[
 \begin{split}
 a_{\mathrm{rad}}(\rho)&=
 \min\left\{L^{-2}c_{\rm curv},\
       \frac{c_{\mathrm p}}{63L^3\rho},\
       \frac{m_*}{63L^6B_{\mathrm p}\rho^2},\
       \frac1{24C_c\rho^4}\right\},\\
 B_{\mathrm{rad}}&=\max\{4,L^2,12K^\sharp\}.
 \end{split}
\]
Then $\theta\le1/12$.
If $B_{\mathrm{rad}}\varepsilon_1\ge\varepsilon_0$ there is
nothing to prove. Otherwise \eqref{eq:reset-one-improvement}
improves the radius to
$\max\{B_{\mathrm{rad}}\varepsilon_1,\varepsilon_0/2\}$.
Repeat while the second entry is larger.
After finitely many repetitions the radius is
$B_{\mathrm{rad}}\varepsilon_1$, proving the proposition.
\end{proof}

\subsection{Ascent to a critical configuration}

Fix a power $\rho$ of $L$ large enough that
\[
 \rho\ge\max\{\rho_{\mathrm{rad}},LM_{\rm ps}\},
 \qquad L^2\rho\ge w_S.
\]
Assume $\mathfrak M\ge M_{\rm ps}$, $R_1\ge L$,
$R_1\mathfrak M\ge L^2\rho$, and at least $6\rho$ sites in each
direction of $\mathbb T^{(k)}$.
These conditions hold with the same constants for every
truncated sequence $(\Omega_0,\ldots,\Omega_n)$:
the separation is inherited, and its terminal torus
has at least as many sites.

Let $K_5,\Theta,r_0,c_8$ be the constants of the actual
critical-point solver, \cref{I-48-91}. Take $c_{\rm ref}=a_{\rm in}$ and $r_L=1/2$;
the latter is a fixed real logarithmic datum radius.
Set
\[
 \begin{split}
 r_s&=\min\{r_0,\tfrac12K_5r_L\},\\
 k_a&=5L^2\max\{1,4(1+L^3)\Theta\}K_5,\qquad
 K_9=4+6k_a,\\
 c_{\rm asc}
 &=\min\left\{\tfrac12,c_8,c_{\rm ref},
        \frac{r_s}{5K_5},\frac{\rho_{\rm ec}}{k_a},
        \frac1{C_{\rm ec}k_a},
        \frac1{2(4k_a+C_{\rm ec}(1+k_a))},
        \frac1{3(2+3k_a)}\right\}.
 \end{split}
\]
All these constants depend only on $L$.

\begin{proposition}[Existence of a regular critical configuration]
\label{I-50-67}
In the preceding geometry, put
\[
 a_{\rm ex}(\rho)=\frac1{C_\downarrow}
 \min\left\{\varepsilon_\downarrow,\,
       \frac{c_{\rm asc}}{L^3B_{\mathrm{rad}}},\,
       \frac{a_{\mathrm{rad}}(\rho)}
                         {K_9L^3B_{\mathrm{rad}}}\right\}.
\]
If $0<\varepsilon\le a_{\rm ex}(\rho)$ and
\[
       V\in\mathcal D(\mathbf B^\flat;
                           \varepsilon,C_\downarrow),
\]
then $C(\mathbf B^\flat,V)$ contains a real critical
configuration in
$\mathfrak U'_k(B_{\mathrm{rad}}C_\downarrow\varepsilon)$.
The constants are independent of the number of levels
and of changes in the region sequence.
In particular, the conclusion holds if every tested
holonomy of $V$ has distance less than $\varepsilon$
from the identity.
\end{proposition}

\begin{proof}
Write $\varepsilon'=C_\downarrow\varepsilon$.
Proposition~\ref{I-50-61} supplies data $V^{(n)}$ on
$\mathbf B^{(n),\flat}$ with $V^{(k)}=V$, agreeing with
$V$ on the retained lower rows, satisfying the terminal
mismatch estimate, and having every tested holonomy
strictly within $\varepsilon'$ of the identity.

We first prove the ascent step.
Suppose $1\le n\le k$, and a real configuration $U_0$
satisfies
\[
 \begin{gathered}
 U_0\in\mathfrak U'_{n-1}(B_{\mathrm{rad}}\varepsilon'),\\
 \mathcal M^jU_0=V^{(n)}_j
                    \quad(j<n,\ c\in\Gamma_j^\flat),\\
 |\mathcal M^nU_0(c)V^{(n)}_n(c)^{-1}-\mathbf1|
                  \le\varepsilon'
                    \quad(c\in\Gamma_n^{(n)}).
 \end{gathered}
\]
Let $j_n(\mathfrak e)$ be the local index for the
truncated sequence. Since
$j_{n-1}(\mathfrak e)=\min\{j_n(\mathfrak e),n-1\}$,
the index changes by at most one.
In scaling-independent notation the regularity bounds
are
\[
 |\partial U(p)-\mathbf1|<rL^{-2j_n(p)},\qquad
 |S_U(b)|<rL^{-3j_n(b)},\qquad
 S_U=\eta D_U^{\mathsf T}\partial U .
\]
Consequently
\[
 U_0\in\mathfrak U'_n(\alpha_0),\qquad
                 \alpha_0=L^3B_{\mathrm{rad}}\varepsilon'
                                      \le c_{\rm asc}.
\]
This includes the passage from depth zero to depth one.
Use level-$n$ lengths in what follows.

The solver requires a bound for the entire inner-star
current $F'(0)=J_{\Omega_0}(U_0)$.
At full bonds it is bounded by $\alpha_0\mathsf L^{-3}$
because $\mathfrak U'_n(\alpha_0)\subset
\mathfrak U_n(\alpha_0)$.
If an inner bond is not full, it and its incident inner
plaquettes have index zero and scale $\eta$.
Each plaquette contribution is bounded by
$\eta^{-3}|\partial U_0(p)-\mathbf1|$, and there are at
most four contributions. Thus
\[
                         \mathsf j:=|F'(0)|_{(-3)}
                                      \le4\alpha_0 .
\]
No plaquette outside the original domain is used.

Put $B_j=0$ for $j<n$, and at a positively oriented
terminal bond define
\[
 B_n(c)=\frac1i\log
       \bigl(V^{(n)}_n(c)\mathcal M^nU_0(c)^{-1}\bigr).
\]
The real logarithm exists and
\[
 \beta:=\sup|B_n|\le\tfrac\pi2\varepsilon'\le\alpha_0,
 \qquad e^{iB}M_{\mathbf B^{(n),\flat}}U_0=V^{(n)}.
\]
The nonlinear datum correction is exactly $B$.
Indeed, at an outward row its additional factor is
constructed from $B_{j-1}=0$, so that factor is the
identity. There are no edge rows in the reduced set.
Moreover
$K_5(\mathsf j+\beta)\le5K_5\alpha_0\le r_s\le r_0$
and $\beta<r_L$.
The actual solver therefore provides a real solution
with
\[
 \|A\|_X\le5K_5\alpha_0,\qquad
 |D^*DA|_{(-3)}
       \le15(1+L^3)\Theta K_5\alpha_0,\qquad
 RD^*A=0 .
\]
Choose any real gauge $h$ with
$h(y)=v_j[A](y)^{-1}$ on every $\Lambda_j$.
The top sets are disjoint, so these values extend to
all bare sites. The solver's constraint and criticality
statements give
\[
 U=(e^{i\eta A}U_0)^h
       \text{ critical in }C(\mathbf B^{(n),\flat},V^{(n)}).
\]

It remains to verify the regularity needed for the reset.
The two solver bounds have the common upper bound
\[
 \alpha'_2=
       5\max\{1,4(1+L^3)\Theta\}K_5\alpha_0 .
\]
Conversion to the local star sizes costs $L^2$:
$\alpha_2=L^2\alpha'_2=k_a\alpha_0$.
The definition of $c_{\rm asc}$ ensures
\[
 \alpha_2\le\rho_{\rm ec},\qquad
 C_{\rm ec}\alpha_2\le1,\qquad
 4\alpha_2+C_{\rm ec}(\alpha_0+\alpha_2)\le\tfrac12,
 \qquad 2\alpha_0+3\alpha_2\le\tfrac13 .
\]
The regularity transfer \eqref{I-40-65}, at
$E=e^{i\eta A}U_0$, consequently gives
\[
 |\partial E(p)-\mathbf1|
       \le(\alpha_0+3\alpha_2)\eta^2\mathsf L(p)^{-2},
 \qquad
 |J(E)(b)|\le(2\alpha_0+3\alpha_2)\mathsf L(b)^{-3}.
\]
For $r=2\alpha_0+3\alpha_2>0$, the full-divergence
conversion gives radius $r+r^2<2r$.
Thus
\[
 E\in\mathfrak U'_n(K_9\alpha_0),\qquad
 U\in\mathfrak U'_n(K_9\alpha_0).
\]
The second inclusion uses gauge invariance, and does
not require a small extension of $h$.
This is a preliminary bound; it has not yet restored
the induction radius.
Since
\[
             K_9\alpha_0
                  =K_9L^3B_{\mathrm{rad}}\varepsilon'
                         \le a_{\mathrm{rad}}(\rho),
\]
\cref{I-49-38} applies to the critical
configuration $U$ and the tested holonomies of
$V^{(n)}$. It gives
\[
                       U\in\mathfrak U'_n
                                      (B_{\mathrm{rad}}\varepsilon'),
\]
which completes the ascent step.

We now iterate.
At depth zero the reduced determining set consists
of all inner bare bonds, so put $U^{(0)}=V^{(0)}_0$.
Its plaquette errors are strictly less than
$\varepsilon'$. At a full bond, subtract the identity
from each of its four plaquettes in the two transported
differences defining $S_{U^{(0)}}$. The identities cancel,
and $|S_{U^{(0)}}|<4\varepsilon'$.
Hence
$U^{(0)}\in\mathfrak U'_0(B_{\mathrm{rad}}\varepsilon')$.

Suppose $U^{(n-1)}$ has been constructed in the required
fiber and with this radius.
Agreement of the descended data gives
$\mathcal M^jU^{(n-1)}=V_j^{(n)}$ on the retained rows
for $j<n$.
For a terminal bond $c$ of level $n$, separation puts
its fine double block inside $\Omega_{n-1}$.
Every bond needed for its average therefore belongs
to the previous terminal frame, on which
$\mathcal M^{n-1}U^{(n-1)}=V^{(n-1)}_{n-1}$.
Locality and the descent mismatch give
\[
 \mathcal M^nU^{(n-1)}(c)
       =\mathcal M(V^{(n-1)}_{n-1})(c),\qquad
 |\mathcal M^nU^{(n-1)}(c)V_n^{(n)}(c)^{-1}-\mathbf1|
       \le\varepsilon'.
\]
The ascent step supplies $U^{(n)}$.
At $n=k$ this is the asserted configuration.
Every step begins and ends with the same radius
$B_{\mathrm{rad}}\varepsilon'$, so no loss accumulates with
the depth of the sequence.
\end{proof}

\subsection{Uniqueness of the regular critical orbit}

\begin{proposition}\label{I-48-100}
Assume the separated-region setting with
\[
 \mathfrak M\ge M_{\rm ps},\qquad R_1\ge L,\qquad
 R_1\mathfrak M\ge w_S,\qquad
 \#\mathbb T^{(k)}\text{ in each direction}\ge6LM_{\rm ps}.
\]
There is $a_{\rm orb}>0$, depending only on $L$, such that
any two real critical configurations in
\[
             C(\mathbf B^\flat,V)\cap\mathfrak U'_k(\alpha),
             \qquad0<\alpha\le a_{\rm orb},
\]
belong to the same residual gauge orbit.
Here $V$ is any real reduced datum.
The constant is independent of the number of levels
and of changes in the region sequence.
\end{proposition}

\begin{proof}
We first put two regular configurations with the same
datum in a common chart. This part does not use criticality.
We will prove that, for $\alpha$ sufficiently small,
$U,U_0\in C(\mathbf B^\flat,V)\cap\mathfrak U'_k(\alpha)$
admit a real $A$ and a gauge $g$ such that
\begin{equation}\label{eq:common-chart}
 U^g=e^{i\eta A}U_0,\qquad
 g=v_j[A]\ \text{on }\Lambda_j,\qquad RD^*A=0,\qquad
 \|A\|_X\le C_{\rm sl}\alpha,\qquad
 a_j(A)=0\ \text{on }\Gamma_j^\flat .
\end{equation}
We take $C_{\rm sl}\ge1$ large enough that the previously
proved slice estimate reads
\[
 \|A\|_X\le C_{\rm sl}(\alpha+4\beta^\sharp)
\]
when both configurations have regularity radius $\alpha$,
the chart lies in its stated small radius, $RD^*A=0$,
and the raw reduced rows satisfy
$\sup|a_j(A)|\le\beta^\sharp$.

\emph{A preliminary gauge.}
Write $W_n=\mathcal M^nU$, $W_{0,n}=\mathcal M^nU_0$.
The curvature propagation estimate makes their contour
means well defined on every block used below.
Define $g_k=\mathbf1$ on $\Omega_k^{(k)}=\Lambda_k$,
and recursively, for $n=k-1,\ldots,0$, set
\[
 \begin{split}
 g_n&=\mathbf1\quad\text{on }\Lambda_n,\\
 g_n(x)&=W_{0,n}(y,x)^{-1}g_{n+1}(y)W_n(y,x),
       \quad x\in B(y),\quad y\in\Omega_{n+1}^{(n+1)} .
 \end{split}
\]
The decomposition into these blocks and $\Lambda_n$
is disjoint. At a block center both contour transports
are the identity, so $g_n$ extends $g_{n+1}$.
Thus $\widehat g=g_0$ is defined on all bare sites,
equals the identity at all original tops, and satisfies
\begin{equation}\label{eq:common-contour}
       (\mathcal M^nU^{\widehat g})(y,x)=W_{0,n}(y,x)
       \quad(y\in\Omega_{n+1}^{(n+1)},\ x\in B(y)).
\end{equation}
This follows directly from covariance and the recursion.

We need a bound on bonds of every truncated frame,
not only on the original reduced rows.
Put
\[
 \delta_n(b)=|\mathcal M^nU^{\widehat g}(b)
                                  W_{0,n}(b)^{-1}-\mathbf1|.
\]
On original interior and outward bonds $\delta_n=0$:
their endpoints are tops, and the original data agree.
At an inward bond of level $n$, the bond lies in the
corridor of an outward bond of level $n+1$.
The latter mismatch is zero. The contour defect at
the inner endpoint is zero by \eqref{eq:common-contour}.
At the outer endpoint, comparison with a contour
of at most $3(L-1)/2$ fine bonds bounds the defect
by $4\kappa_1^{\mathrm{SO}}\alpha$, since these fine bonds
are interior datum bonds and both fine plaquette
sizes are at most $2\alpha$.
The corridor comparison in the proof of \cref{I-50-33}
therefore gives
\[
                       \delta_n(b)\le C_0\alpha
\]
with $C_0=4\kappa_\downarrow+4\kappa_1^{\mathrm{SO}}$.

For a bond $b$ with both endpoints in $\Omega_n$,
let $j_b\ge n$ be the largest level containing both
endpoints. Applying the same comparison downwards gives
\[
 \delta_n(b)\le C_0\alpha+
                 4\kappa_\downarrow\alpha
                         \sum_{i=n}^{j_b-1}L^{2(i-j_b)}.
\]
Indeed, when the endpoints lie in one block their
contour defects vanish and the one-block estimate
is at most $8\kappa_1^{\mathrm{SO}}\alpha$.
When they lie in different blocks, their centers
form a bond with the same value of $j_b$.
The contour defects again vanish, and the extra
corridor error is at most
$4\kappa_\downarrow\alpha L^{2(n-j_b)}$ by curvature
propagation. At level $j_b$ one has an original
interior or inward bond, which starts the induction.
Since $\sum_{s\ge1}L^{-2s}\le2L^{-2}$, this proves
\begin{equation}\label{eq:common-preliminary-bound}
 \delta_n(b)\le C_{\rm pg}\alpha
 \quad\text{for every inner level-$n$ bond meeting }\Omega_n,
 \qquad
 C_{\rm pg}=4(1+2L^{-2})\kappa_\downarrow
                         +4\kappa_1^{\mathrm{SO}}.
\end{equation}
Bonds with just one endpoint in $\Omega_n$ are original
outward bonds and have zero mismatch.
All the comparisons use carried double blocks;
separation places their fine bonds in $\Omega_n$.
The constants thus do not depend on the depth.
Set $U'=U^{\widehat g}$. It has the same original reduced
datum as $U$ and the same regularity radius.

\emph{Normalizing one level.}
We describe the operation used in the upward induction.
At depth $j$, suppose
\[
 U'^{g_0}=e^{i\eta_jA_0}U_0,\qquad
 g_0=v_i[A_0]\ \text{on the lower tops }(i<j),\qquad
 g_0=v_{j-1}[A_0]\ \text{on the terminal tops}.
\]
The residual factor $v_i[A_0]^{-1}g_0$ is therefore
the identity on lower tops and
$\mathsf v_{j-1}[A_0]^{-1}$ on terminal tops.
Prescribe the induced gauge correction to be
\[
 r=\mathbf1\quad\text{on lower tops},\qquad
 r=\mathsf v_{j-1}[A_0]\quad\text{on terminal tops}.
\]
Every fine block of a terminal top lies in $\Omega_j$.
The mean-tower estimate at level $j-1$, followed by
the one-step top-factor bound, gives
\[
        \sup|r-\mathbf1|
                 \le(2\kappa/L)\|A_0\|_{X,\mathbf B^{(j)}}.
\]
Here $\kappa=\kappa_{6,\mathsf m}$ is the one-step factor constant
from \cref{I-39-43}, enlarged if necessary.
The prescribed-top correction supplies $w=e^{i\lambda}$,
$A=A_0^w$ and $g=wg_0$ with
\[
 \begin{gathered}
 RD^*A=0,\qquad g=v_i[A]\text{ on every depth-$j$ top},\\
 |A|_{(-1)}\le k_1\|A_0\|_X,\qquad
 k_1=1+2C_{\rm gc}K_3(3+2\kappa/L).
 \end{gathered}
\]
We enlarge $k_1$ if necessary so that
$k_1\ge\max\{1,2\kappa/L\}$.
The estimate uses $|D^*A_0|_{(-2)}\le3\|A_0\|_X$.
The normalization is exact: the new residual factor
is the prescribed $r$ times the old one.

Suppose the reduced data of $U'$ and $U_0$ at depth $j$
differ by at most $\gamma$.
Let $B_i(c)$ be the small logarithm of that discrepancy.
The normalized chart identities give $a_i(c)=B_i(c)$
at interior bonds. If $c$ is outward with outer end $z$,
they give
\begin{equation}\label{eq:common-outward}
 e^{ia_i(c)}=
 \begin{cases}
 e^{iB_i(c)}W_{0,i}(c)
       \mathsf v(W_{0,i-1},B_{i-1})(z)W_{0,i}(c)^{-1},
                                      &z=c_+,\\
 \mathsf v(W_{0,i-1},B_{i-1})(z)^{-1}e^{iB_i(c)},
                                      &z=c_-.
 \end{cases}
\end{equation}
To verify this identity, use
$\mathcal M^iU'=s_i(c_-)^{-1}e^{ia_i(c)}
W_{0,i}(c)s_i(c_+)$ with $s_i=v_i[A]^{-1}g$.
At the outer endpoint
$s_i(z)=\mathsf v_{i-1}[A](z)^{-1}s_{i-1}(z)
=\mathsf v_{i-1}[A](z)^{-1}$.
The block $B(z)$ consists of interior rows at level
$i-1$, on which $a_{i-1}=B_{i-1}$, so locality replaces
this one-step factor by the one displayed in
\eqref{eq:common-outward}.

The logarithm estimate gives $|B_i|\le\pi\gamma/2$.
The one-step factor has a logarithm of norm at most
$\kappa\sup|B_{i-1}|$ on its block. Subadditivity of
the real group angle in \eqref{eq:common-outward},
and the small branch $|a_i|<\pi$, imply
\[
             \sup_{\Gamma_i^{(j),\flat}}|a_i|
                        \le C_q\gamma,\qquad
             C_q=\tfrac\pi2(1+\kappa).
\]
Both $e^{i\eta_j A}U_0=U'^g$ and $U_0$ have regularity
radius $\alpha$ at depth $j$.
Lemma~\ref{I-47-4} therefore improves the
chart to
\begin{equation}\label{eq:common-slice-improvement}
                   \|A\|_{X,\mathbf B^{(j)}}
                         \le C_{\rm sl}(\alpha+4C_q\gamma).
\end{equation}
This bound is independent of the size of the incoming
chart, once that chart satisfies the correction caps.

\emph{Induction over the levels.}
Choose
\[
 k_2=\max\{\pi C_{\rm pg},\
                  C_{\rm sl}(1+4C_qC_{\rm pg})\}.
\]
Choose $\sigma>0$ smaller than the common weighted
chart radius, the gauge-correction radius and the
slice radius, and so small that
$2(L+\kappa)\sigma<1/4$.
Finally take $c_{\rm cmp}>0$ within all reference and
curvature thresholds already used, including that of
\cref{I-10-101}, and such that
\begin{equation}\label{eq:common-cap}
                        k_1k_2L^2c_{\rm cmp}\le\sigma .
\end{equation}
These are choices of constants depending only on $L$.

At depth zero, $\eta_0=1$ and every scale equals one.
Define
\[
 A^{(0)}(b)=\frac1i\log(U'(b)U_0(b)^{-1}),\qquad
                            g^{(0)}=\mathbf1 .
\]
For $\alpha\le c_{\rm cmp}$,
\eqref{eq:common-preliminary-bound} and the real
logarithm estimate give
$\|A^{(0)}\|_X\le\pi C_{\rm pg}\alpha\le k_2\alpha$:
a first difference contains at most two values of
$A^{(0)}$, transported by isometries.
The chart is normalized at every depth-zero site.

Suppose a normalized pair
$(A^{(j-1)},g^{(j-1)})$ has been constructed with
$\|A^{(j-1)}\|_X\le k_2\alpha$.
In level-$j$ lengths put $A_0=LA^{(j-1)}$.
The dimensionless field is unchanged:
$\eta_jA_0=\eta_{j-1}A^{(j-1)}$.
Only the new terminal local index changes, by one;
thus
\[
                     \|A_0\|_{X,\mathbf B^{(j)}}
                           \le L^2k_2\alpha .
\]
The prior terminal site set contains both the new
level-$(j-1)$ tops and the level-$j$ terminal tops.
The mean objects depend on $\eta A$, so the incoming
pair has exactly the normalization required for
the one-level operation.

The reference configurations remain in
$\mathfrak U'_j(\alpha)$: truncating the local indices
weakens the defining scale bounds. The geometry of
each truncation satisfies the same operator thresholds,
and the terminal torus has at least as many sites.
For $j<k$, \eqref{eq:common-preliminary-bound} gives
$\gamma=C_{\rm pg}\alpha$ on its reduced frame.
For $j=k$, the original reduced data agree, so
$\gamma=0$.
Condition \eqref{eq:common-cap} permits the correction,
and \eqref{eq:common-slice-improvement} gives
$\|A^{(j)}\|_X\le k_2\alpha$ at intermediate levels
and $\|A^{(k)}\|_X\le C_{\rm sl}\alpha$ at the last level.
Thus the bound is restored after every rescaling.

Set $A=A^{(k)}$ and $g=g^{(k)}\widehat g$.
The preliminary gauge is the identity on the original
tops, so the final pair remains normalized there.
At the final level all $B_i$ vanish.
On every outward block the factor in
\eqref{eq:common-outward} is consequently the identity,
and the small logarithmic branch gives $a_i(c)=0$.
Interior rows vanish as well. This proves
\eqref{eq:common-chart}.

\emph{Criticality and uniqueness.}
Now take $U_0=U_a$, $U=U_b$, where both are critical
in the same reduced fiber.
Use \eqref{eq:common-chart}, and choose $a_{\rm orb}$
so small that
\[
 \begin{gathered}
 a_{\rm orb}\le\min\{1,c_{\rm cmp},c_8,c_6,
                                      L^{-2}a_{\mathsf m}\},\\
 C_{\rm sl}a_{\rm orb}\le\min\{r_0,c_6\},\qquad
                             4K_5a_{\rm orb}\le r_0 .
 \end{gathered}
\]
The choices of $c_8,r_0$ in \cref{I-48-91} ensure
$3K_1L^3\epsilon^*\le1/8$ and the applicability of
\cref{I-48-90} throughout this ball. Thus the multiplier
equivalence of \cref{I-46-60} applies at both $A$ and zero.

The configuration $E=U_b^g=e^{i\eta A}U_a$ is critical
for its induced datum by gauge covariance.
The converse part of the criticality statement gives
$F'(A)=K_A^*\omega_b$. Together with $RD^*A=0$ and
$a_j(A)=0$, this makes $(A,\omega_b)$ a solution of
the actual solver system with datum $B^\natural=0$.
The same argument at $A=0$, using the criticality
of $U_a$, gives a solution $(0,\omega_a)$ of that
same system. Its forcing satisfies
$\mathsf j=|F'(0)|_{(-3)}\le4\alpha$ by the earlier
inner-star calculation.
The displayed caps ensure both admissibility of the
solver and $\|A\|_X\le r_0$.
Its uniqueness in that ball therefore gives $A=0$.

It follows that $U_b^g=U_a$.
The normalization in \eqref{eq:common-chart} gives
$g=v_j[0]=\mathbf1$ at all tops.
Every node of the reduced determining set is a top,
so $g$ is a residual gauge transformation.
This proves the proposition.
\end{proof}

\begin{corollary}\label{I-58-16}
In the geometry of \cref{I-50-67}, put
\[
 a_{\mathrm{can}}(\rho)=\min\left\{a_{\rm ex}(\rho),
       \frac{a_{\rm orb}}{B_{\mathrm{rad}}C_\downarrow},
       \frac{a_{\mathrm{rad}}(\rho)}{B_{\mathrm{rad}}C_\downarrow}\right\}.
\]
Let $V$ be real reduced data whose tested holonomies are defined
and have distance less than $\varepsilon$ from the identity,
where $0<\varepsilon\le a_{\mathrm{can}}(\rho)$.
The critical configurations in
$C(\mathbf B^\flat,V)\cap\mathfrak U'_k(a_{\rm orb})$
form one nonempty residual gauge orbit. This orbit lies in
$\mathfrak U'_k(B_{\mathrm{rad}}\varepsilon)$.
\end{corollary}

\begin{proof}
Proposition~\ref{I-50-67} supplies a critical configuration
with radius $B_{\mathrm{rad}}C_\downarrow\varepsilon$.
This is within both $a_{\rm orb}$ and $a_{\mathrm{rad}}(\rho)$.
Apply \cref{I-49-38} once more, now using the original tested
holonomy bound $\varepsilon$, to obtain radius
$B_{\mathrm{rad}}\varepsilon$. Proposition~\ref{I-48-100}
puts every critical configuration of radius $a_{\rm orb}$
in this orbit. Gauge invariance preserves the improved radius.
\end{proof}

\subsection{The regular pair supplied by the real construction}

Fix a power $m=m(L)$ of $L$ with
\[
 m\ge\max\{\mathfrak M_{\mathrm{pr}},L^2\}.
\]
We increase the geometric thresholds in the real existence theorem
so that $\mathfrak M\ge m$ and $\rho\ge Lm$.
The sequence can then be regarded as separated with block parameter
$m$ and relative gap $L$. Indeed, every original block is a union
of $m$-blocks, its complement is the same point set, and the
original separation exceeds $Lm\xi_q$. The layers, determining
rows, local scales and weighted distance are unchanged. In the
regular-pair notation we take $t=k$ and $\xi=\eta$. The final
lattice has at least $6Lm$ sites in each direction, which also
implies the volume hypothesis of \cref{I-88-53}.

With $B_{\mathrm p},c_{\mathrm p}$ from
\eqref{eq:regional-local-gauge}, put $C_{\mathrm P}=504L^2B_{\mathrm p}$.
Choose $\varepsilon_{\mathrm P}>0$, depending only on $L$, so that
\begin{equation}\label{eq:background-natural-radius}
 \varepsilon_{\mathrm P}\le\tfrac13,\qquad
 504L^2m\varepsilon_{\mathrm P}\le c_{\mathrm p},\qquad
 C_{\mathrm P}m\varepsilon_{\mathrm P}
       \le\min\{1,a_{\mathrm{nl}},a_{\mathrm{pr}}\}.
\end{equation}

\begin{lemma}\label{I-60-90}
Suppose $U_0$ is real and regular at radius
$0<\varepsilon\le\varepsilon_{\mathrm P}$:
\[
 |\partial U_0(p)-I|<\varepsilon\eta^2\mathsf L(p)^{-2},
 \qquad |J(U_0)(b)|<\varepsilon\mathsf L(b)^{-3},
\]
at inner plaquettes and full bonds, respectively. With
$\alpha=C_{\mathrm P}m\varepsilon$, the fixed witness $(m,L)$
gives
\[
 (U_0,F'(0))\in
 \mathfrak P_k^\natural((\Omega_q);\alpha,\alpha_1,\alpha)
 \qquad(\alpha_1>0).
\]
Here $F$ is the action summed over inner plaquettes. The assertion
includes every cube in the natural class, whether or not it
meets the lower of its two layers.
\end{lemma}

\begin{proof}
Take $A'=0$ in the definition of the class. At a plaquette based
in layer $q$, one has $\mathsf L(p)\ge\xi_q$, giving the
required curvature bound. At a full bond the action gradient is
$F'(0)(b)=J(U_0)(b)$. An inner bond which is not full lies at
the bare scale, as do all the corners of its inner plaquettes.
Each such plaquette contributes at most
$\varepsilon\eta^{-3}$ to $F'(0)(b)$: its field strength is
bounded by $\varepsilon\eta^{-2}$, and the transpose difference
contributes $\eta^{-1}$. There are at most four contributions.
Thus, for the top $y_b$ of the initial site,
\[
 \sup_b\xi_{y_b}^3|F'(0)(b)|\le4\varepsilon<\alpha.
\]
This proves the current condition, including boundary bonds.

It remains to construct the local gauges. For a full bond $b$
the quaternion decomposition of each plaquette gives
\[
 D_{U_0}^{\mathsf T}\partial U_0(b)
       =a(b)I+i\eta^2J(U_0)(b),\qquad
 |a(b)|\le\eta^{-1}
          \max_{p\ni b}|\partial U_0(p)-I|^2.
\]
Since $\mathsf L(p)\ge\mathsf L(b)\ge\eta$,
\begin{equation}\label{eq:background-matrix-divergence}
 |D_{U_0}^{\mathsf T}\partial U_0(b)|
       <2\varepsilon\eta^2\mathsf L(b)^{-3}.
\end{equation}
The scalar part in this calculation is the sum of the two
differences of $\tfrac12\operatorname{tr}\partial U_0$;
$1-\tfrac12\operatorname{tr}g=\tfrac12|g-I|^2$ for $g\in SU(2)$
gives the stated bound.

Let $q\ge1$ and let $\mathsf q\subset\Omega_q$ be a cube
of side $l\xi_q\le12m\xi_q$, aligned with the $m$-partition.
Enlarge it concentrically to side $(l+4L)\xi_q$.
The enlarged cube and its $2\eta$-neighborhood lie in
$\Omega_{q-1}$, since
\[
 2L\xi_q+2\eta\le(2L^2+2)\xi_{q-1}<Lm\xi_{q-1}.
\]
Its side is at most $16m$, whereas half the torus side is
at least $3Lm>16m$. Thus the local gauge theorem applies
after rescaling by $\xi_q$ and extending $U_0$ by $I$
outside $\Omega_0$. That extension does not change any
plaquette or difference used near the enlarged cube.
The curvature and matrix-divergence parameters in these units
are bounded by $2L^3\varepsilon$, and
\[
 2L^3l\varepsilon\le24L^3m\varepsilon
       \le(L/21)c_{\mathrm p}\le c_*.
\]
Restricting the resulting gauge to $\mathsf q$ gives
$U_0^v=e^{i\eta A}$ and
\[
 \xi_q|A|,\ \xi_q^2|\nabla^\eta A|
       <2B_*L^3l\varepsilon
       \le C_{\mathrm P}m\varepsilon=\alpha.
\]
Only ordinary first differences occur here. A natural cube is
contained in $\mathcal L_q\cup\mathcal L_{q+1}\subset\Omega_q$,
so this applies to every required cube.

At level zero no enlargement is needed. The cube has at most
$12m$ bare spacings in each direction. The comb gauge of
\cref{I-10-51}, with $s=12m$, gives $U_0^v=e^{ia^0}$ and
$|a^0|\le12\pi m\varepsilon$; its logarithm is defined since
$24m\varepsilon\le c_{\mathrm p}<2$. For $A=\eta^{-1}a^0$,
\[
 \eta|A|\le12\pi m\varepsilon<\alpha,\qquad
 \eta^2|\nabla^\eta A|\le24\pi m\varepsilon<\alpha.
\]
These estimates hold wherever the difference is defined from
bonds in the cube. This completes the natural-class conditions.
\end{proof}

\section{The canonical holomorphic background}\label{sec:canonical-background}

For a real reference $U_0$, write $W_{0,j}=\mathcal M^jU_0$
and $V_0=M_{\mathbf B^\flat}U_0$.
The response of \cref{I-88-93} prescribes the straightened
average. We first convert the actual datum into that coordinate.
Choose $\kappa_{\mathrm m}\ge\max\{\kappa_{6,\mathsf m},
\kappa_{\mathrm{st}},2L\}$. Write $\mathsf R_{\mathsf m}$ for the
mean-Landau projection $\mathsf R$ of \cref{I-88-93};
this notation distinguishes it from the corner projection used
in the real inhomogeneous solver.
For an outward row $c\in\Gamma_j$ with outer endpoint $z$, put
$\theta_z(B)=\mathsf v(W_{0,j-1},B)(z)$, where its argument
is restricted to the interior level-$(j-1)$ bonds in $B(z)$.
Define $B_\natural=B$ at interior rows and, at outward rows, by
\begin{equation}\label{eq:background-inverse-twist}
 B_\natural(c)=
 \begin{cases}
  i^{-1}\log[e^{iB(c)}V_0(c)\theta_z(B)V_0(c)^{-1}],&z=c_+,\\
  i^{-1}\log[\theta_z(B)^{-1}e^{iB(c)}],&z=c_-.
 \end{cases}
\end{equation}
On a ball $\sup|B|<r_{\mathrm{tw}}$, with
$r_{\mathrm{tw}}>0$ depending only on $L$, this is holomorphic,
real at real $B$, and
\[
 \sup|B_\natural|\le C_{\mathrm{tw}}\sup|B|,
 \qquad C_{\mathrm{tw}}\ge1.
\]
Indeed, $\theta_z(B)=e^{i\zeta_z(B)}$ with
$|\zeta_z(B)|\le\kappa_{\mathrm m}\sup|B|$.
Each logarithm is a product of two exponentials whose total
exponent norm is at most $(1+\kappa_{\mathrm m})\sup|B|$.
Choose this less than $1/16$ and use
$|\log(e^Xe^Y)|\le s+4s^2\le2s$, $s=|X|+|Y|$.
This also proves reality and the bound. Since the twist fixes
interior rows, $\theta_z(B_\natural)=\theta_z(B)$; substituting
in \eqref{eq:critical-seam-twist} shows
$(B_\natural)^\natural=B$ whenever both arguments belong to
the indicated balls.



\begin{proof}[Proof of \cref{I-58-22}]
\emph{The real reference.}
For real $B$, the value of the datum at a retained bond changes
by at most $\beta=\sup|B|$. The averaged value on the other
side of a tested plaquette is a product of three straightening
factors and the old averaged bond. Each factor changes by at
most $\kappa_{\mathrm m}\beta$; hence this averaged value
changes by at most $3\kappa_{\mathrm m}\beta$.
The four-factor product estimate for unitary matrices gives
\[
 |\partial(e^{iB}V_0)-\partial V_0|
       \le4\max\{1,3\kappa_{\mathrm m}\}\beta.
\]
The local small-field hypothesis for each such average follows
because its supporting plaquettes are tested interior
plaquettes. Choose $C_\partial$ at least the displayed
coefficient and make the radii small enough that these averages
remain defined and
\[
 \varepsilon_{\mathrm G}(\rho)+C_\partial r_{\mathrm G}(\rho)
                       \le a_{\mathrm{can}}(\rho).
\]
The real orbit assertion follows from \cref{I-58-16};
\cref{I-3-96} supplies its unique axial representative.
At $B=0$ we obtain $U_0\in
\mathfrak U'_k(B_{\mathrm{rad}}\varepsilon_1)$.
Increase the geometric floors to those of the fixed-witness
construction and choose $B_{\mathrm{rad}}\varepsilon_{\mathrm G}
\le\varepsilon_{\mathrm P}$.
Then \cref{I-60-90} puts $(U_0,F'(0))$ in the class required
by \cref{I-88-93}. Criticality and the reduced
criticality--multiplier equivalence at $A=0$ give
\[
 J_0:=F'(0)=Q^{\mathsf T}\lambda_0.
\]
Thus all hypotheses of that response theorem have now been
verified for the reference constructed here.

\emph{The response and its multiplier.}
Set
\[
 A(B)=\mathbf H(B_\natural),\qquad
 \lambda(B)=(I+H^{\mathsf T}\mathsf C'_{A(B)}{}^{\mathsf T})^{-1}
                         H^{\mathsf T}F'(A(B)).
\]
Choose $C_{\mathrm{tw}}r_{\mathrm G}<\varepsilon_{\mathrm{crit}}$.
Then \cref{I-88-93} gives the field, row and gauge equations,
holomorphy and reality. The response bound gives
$\|A(B)\|\le C\beta$. The field norm here controls
$\|A\|_{X,\mathbf B}$ with a factor depending only on $L$:
the scale of either endpoint of a bare bond or of a bond in
one of its plaquettes differs from its local scale by at most
one power of $L$, and a curl is a sum of first differences.

We record the multiplier estimate, since its initial value need
not vanish. Group the inner plaquettes by their base cells and set
\[
 E_y(A,\delta)=\eta^3\sum_{p\ {\rm over}\ y}
       \langle DA(p),D\delta(p)\rangle.
\]
There are at most $3(\xi_y/\eta)^3$ plaquettes in this sum,
and $DA=\nabla_\mu A_\nu-\nabla_\nu A_\mu$ there. Thus
\begin{equation}\label{eq:response-local-energy}
 |E_y(A,\delta)|\le12\xi_y^{-1}
                   \mathsf g_y(A)\mathsf g_y(\delta).
\end{equation}
The calculation is bilinear, including for complex fields.
The linear coefficient in the local star calculation of
\cref{I-87-84}, rescaled by $\xi_y$, gives
$|\Delta'A|^\sharp_y\le C(L)\max_{N(y)}|A|$, where
$\Delta'=F''(0)-D^{\mathsf T}D$.
Indeed, its unscaled bound is $C_{\mathrm a}\alpha\sup|\widehat A|$,
with $\alpha=L^2\alpha_0\le1/2$ and
$\sup|\widehat A|\le L\max_{N(y)}|A|$.
Taylor's formula and the exact transpose identity give
\[
 (F'(A)-J_0,\delta)_\eta
      =\sum_yE_y(A,\delta)
                   +(\Delta'A+\mathcal V(A),\delta)_\eta.
\]
Test a row $c$ against a unit vector $e$ and take $\delta=He_c$.
The row pairing contributes $\xi_c$, while the energy bound
contributes $\xi_y^{-1}$. The column estimate for $H$ from
\cref{I-88-53}, the local action bounds and the last identity
therefore bound the dual pairing by
\[
 C\sum_y\frac{\xi_c}{\xi_y}
       e^{-\delta_{\mathrm{reg}}d(y,y_c)}
                    \max_{z\in N(y)}\|A\|_z\le C\|A\|.
\]
The last step is precisely the scale-ratio summation in
\eqref{eq:actual-response-transpose}. Hence
$|H^{\mathsf T}(F'(A)-J_0)|_\infty\le C\|A\|$.
Also $|H^{\mathsf T}J_0|_\infty\le8\hat{\mathsf b}
\sup|J_0|^\sharp\le C\varepsilon_1$.
The inverse in the formula for $\lambda$ has norm at most
two on this ball. Consequently
\[
 \sup|\lambda(B)|\le C(\varepsilon_1+\beta),\qquad
 \lambda(0)=H^{\mathsf T}Q^{\mathsf T}\lambda_0=\lambda_0.
\]
The equation for the action gradient now gives
\begin{equation}\label{eq:background-second-difference}
 D^{\mathsf T}DA
   =(Q+\mathsf C'_A)^{\mathsf T}\lambda-J_0
                       -\Delta'A-\mathcal V(A).
\end{equation}
The weighted column and local action bounds control its
$(-3)$ norm by $C(\varepsilon_1+\beta)$ on the whole ball.

\emph{Complex averages and the axial gauge.}
We verify definedness of the averages before applying gauge
covariance. On a one-step double block, let
$\mathcal O_c(\theta)$ mean that each contour family has
relative diameter less than $\theta\omega_{\mathrm c}$ and
each loop entering the outer logarithm lies within $\theta$
of $I$. These are open conditions. In the three coordinate
templates of a double block, continuity at $(a^0,a)=(0,0)$
gives a radius $r_{\mathrm c}(L)>0$ such that
$e^{ia}e^{ia^0}\in\mathcal O_c(1/2)$ when
$|a^0|,|a|<r_{\mathrm c}$. A comb gauge puts any real
reference with plaquette deviations at most
$r_{\mathrm c}/(4\pi L)$ into this form with
$|a^0|\le r_{\mathrm c}/2$.

For a complex gauge $g$, relative contour products are
conjugated at their initial site. Thus a condition number
at most $\varkappa$ multiplies their diameter and the loop
distance by at most $\varkappa$. The mean transforms in the
same way: its logarithmic mean equation is conjugated, and
uniqueness in its defining neighborhood identifies the result.
It follows that the average of $W^g$ is defined and equals
$\mathcal M(W)^g$ whenever $W\in\mathcal O_c(1/2)$ and the
two endpoint condition numbers are at most two.

The mean tower bounds give, inside a pillar of level $j$,
\[
 |a_n|\le2L^{n-j}|A|_{(-1)},\qquad
 \mathsf v_i=e^{iY_i},\quad
 |Y_i|\le2\kappa_{\mathrm m}L^{i-j}|A|_{(-1)}.
\]
On a double block the same bounds hold with $j-1$ in place
of $j$. Hence the accumulated factor
$v_n=\mathsf v_0\cdots\mathsf v_{n-1}$ has condition number
at most $4/3$ if $|A|_{(-1)}<\rho_{\mathrm f}$, for a
sufficiently small $\rho_{\mathrm f}(L)>0$.
Shrinking the curvature and field radii also puts every
straightened one-step configuration in $\mathcal O_c(1/2)$.
Induction over the tower then proves, for every gauge $h$
with condition number at most $3/2$,
\begin{equation}\label{eq:background-complex-tower}
 \mathcal M^n((e^{i\eta A}U_0)^h)
                     =\mathsf x_n^{\,hv_n},\qquad
 \mathsf x_n=e^{ia_n}W_{0,n}.
\end{equation}
At the induction step the endpoint factor is $hv_n$, whose
condition number is at most two; the preceding one-step
argument therefore establishes both definedness and the identity.

Use the parent $z(x)$ and level $n(x)$ in each pillar.
Starting at its top, define
\begin{equation}\label{eq:background-pillar-gauge}
 \begin{aligned}
  k(y)&=I &&(y\in\Lambda_*),\\
  k(x)&=k(z(x))
       [\mathsf v_{n(x)}(z(x))\cdots
        \mathsf v_{n(z(x))-1}(z(x))]^{-1}
                     \mathsf x_{n(x)}(\Gamma_{z(x),x}),\\
  h(x)&=k(x)v_{n(x)}(x)^{-1}.
 \end{aligned}
\end{equation}
The second line applies only below the top. The parent has
higher level, so the recursion is well defined. All factors
are holomorphic and unitary at real $A$. Since $U_0$ is axial,
the reference transport along each comb is $I$; moving its
unitary factors to the right expresses the straightened comb
as a product of exponentials. A comb has at most
$3(L-1)/2$ bonds. Along an ascent its levels are distinct,
and each index in an accumulated factor occurs at most once.
The total norm of the exponents in $h(x)$ is therefore at most
\[
 \left[3(L-1)+2\kappa_{\mathrm m}\right]
       \sum_{i<j}L^{i-j}|A|_{(-1)}
       \le(\kappa_{\mathrm m}+5)|A|_{(-1)}.
\]
Choose this less than $1/5$. Then
$|h-I|<1/4$, $\operatorname{cn}(h)<3/2$, and $h(0)=I$.
The recursion and \eqref{eq:background-complex-tower} give
\[
 h=v_j^{-1}\text{ on }\Lambda_j,\qquad
 h(x)=h(z(x))(\mathcal M^{n(x)}(e^{i\eta A}U_0))
                              (\Gamma_{z(x),x}).
\]

Set $\Phi_{V_0}(B)=(e^{i\eta A(B)}U_0)^{h(B)}$.
For an interior row both endpoint factors $hv_j$ are $I$,
and $a_j=B_\natural=B$ there. At an outward level-$j$ row,
the outer endpoint $z$ has level $j-1$, so
$h(z)v_j(z)=\mathsf v_{j-1}(z)$.
The latter factor depends only on the interior rows in $B(z)$,
where $a_{j-1}=B_{j-1}$; it is therefore $\theta_z(B)$.
Substituting \eqref{eq:background-inverse-twist} cancels this
factor on the appropriate side of the row. This proves
$M_{\mathbf B^\flat}(\Phi_{V_0}(B))=e^{iB}V_0$ for complex
$B$, with all averages defined.

\emph{Identification at real points.}
For real $B$, the response and the gauge are real and unitary.
The criticality conclusion of \cref{I-88-93}, followed by
gauge invariance, makes $\Phi_{V_0}(B)$ critical in its actual
fiber. The bounds on $A$ and \eqref{eq:background-second-difference},
together with the regularity conversion for the local action,
put $e^{i\eta A}U_0$ in
$\mathfrak U'_k(a_{\rm orb})$ after decreasing
$\varepsilon_{\mathrm G}$ and $r_{\mathrm G}$ once more.
The class is invariant under the real gauge $h$.
The last identity for $h$ is exactly the outside-node axial
condition of \cref{I-3-96}. Thus the unique regular orbit
and its unique axial representative give
$\Phi_{V_0}(B)=U^\circ(B)$. The identity theorem, applied
to matrix entries on the complexification of the real row
space, proves uniqueness of the holomorphic continuation.

\emph{Bounds.}
All preceding smallness conditions are finite inequalities
between positive constants depending only on $L$ and $\rho$;
they can be met simultaneously. On the resulting whole ball,
$A$, $D^{\mathsf T}DA$ and $h-I$ are uniformly bounded and
vanish at zero. The Schwarz lemma on each complex line through
zero therefore bounds them by $C\beta/r_{\mathrm G}$ in
their stated norms.

For curvature, the local action estimate gives
$|\partial(e^{i\eta A}U_0)-I|
\le C\eta^2\mathsf L^{-2}$ throughout this ball.
Conjugation by $h$ costs at most $3/2$. The holomorphic map
\[
 B\longmapsto\eta^{-2}\mathsf L(p)^2
       [\partial\Phi_{V_0}(B)(p)-\partial U_0(p)]
\]
is thus uniformly bounded and vanishes at zero.
For a full bond the action-current comparison gives
\[
 |J(e^{i\eta A}U_0)(b)|
       \le\tfrac43|F'(A)(b)|,
 \qquad F'(A)=\mathcal Q'(A)^{\mathsf T}\lambda.
\]
The weighted column bound controls the latter by
$C\mathsf L(b)^{-3}$. The current transforms by
$J(E^h)(b)=\operatorname{Ad}(h(b_-))J(E)(b)$, as follows
by cancellation of the endpoint factors in its plaquette
formula. Hence
\[
 B\longmapsto\mathsf L(b)^3
       [J(\Phi_{V_0}(B))(b)-J(U_0)(b)]
\]
is also uniformly bounded and vanishes at zero. Applying
Schwarz to these two maps and adding the original regularity
bound $B_{\mathrm{rad}}\varepsilon_1$ proves the last two
estimates. Increase $K_{\mathrm G}(\rho)$ to cover all the
finite constants above.
\end{proof}

Choose a regularity radius $\varepsilon_{\mathrm v}>0$ and
smaller datum radii $\bar\varepsilon,r_B>0$ such that
\[
 \varepsilon_{\mathrm v}+\varepsilon_{\mathrm v}^2
          \le a_{\rm orb},\qquad
 \bar\varepsilon\le\varepsilon_{\mathrm G}(\rho),\qquad
 r_B\le r_{\mathrm G}(\rho),\qquad
 B_{\mathrm{rad}}(\bar\varepsilon+C_\partial r_B)
          \le\varepsilon_{\mathrm v}.
\]
For $V_0$ with tested holonomies less than $\bar\varepsilon$
and real $\sup|B|<r_B$, the orbit in \cref{I-58-22} lies in
$\mathfrak U'_k(\varepsilon_{\mathrm v})\subset
\mathfrak U_k(\varepsilon_{\mathrm v})$.
Conversely, any critical configuration regular at radius
$\varepsilon_{\mathrm v}$ lies in
\[
 \mathfrak U_k(\varepsilon_{\mathrm v})\subset
 \mathfrak U'_k(\varepsilon_{\mathrm v}
                       +\varepsilon_{\mathrm v}^2)
 \subset\mathfrak U'_k(a_{\rm orb}),
\]
and hence in that orbit. Thus, at regularity radius
$\varepsilon_{\mathrm v}$, the conditional regular-critical
set in \cref{I-3-115} is exactly the orbit constructed here. Its axial
representative is $U^\circ(B)$, and its conditional
holomorphic continuation is the restriction of $\Phi_{V_0}$,
by the uniqueness in \cref{I-3-117}.
The whole ball $\sup|B|<r_B$ is its domain of defined complex
averages, since the determining equation holds there. At real
points the real and complex averages agree.

\section{Local expansions of the background inverse}\label{sec:local-background-expansion}

Keep the separated-sequence setting, scales and torus bound of
\cref{sec:regular-system-inverse}. We use the physical fifth system
and the natural pairs of \cref{def:natural-regular-pair}.

We use the pieces and local systems from the proof of
\cref{I-65-83}. The following carriers record the data needed
by each factor. For a cube \(\mathsf q\in\mathsf Q_n^\natural\)
and an integer \(r\ge0\), let \(\mathsf q^{+r}\) be the union
of the bare \(\mathfrak M\)-cubes contained both in
\(\Omega_0\) and in the concentric cube of side
\(s(\mathsf q)+2r\mathfrak M\xi_n\).
A cell is a nonempty family of at most \(r_{\rm w}\) cubes
of \(\mathsf Q^\natural\), connected by common points; its
union is also denoted by \(X\), and
\[
 X^\natural=\bigcup_{\mathsf q\in X}\mathsf q^{+5}.
\]
The family is part of the cell. A tree walk
\(\omega=(X_0,\ldots,X_m;p)\) has
\[
 p(i)<i,\qquad X_i^\circ\cap X_{p(i)}^\circ\ne\varnothing,
 \qquad
 p(i+1)\in\{i,p(i),p^2(i),\ldots,0\}.
\]
Its length is \(m\), its carrier is
\([\omega]=\bigcup_iX_i^\natural\), and
\(d^{\rm tr}(\omega;y,z)\) is the infimum of the weighted
lengths of connected finite unions of segments containing
\(y,z\) and meeting every \(X_i\).

For the piece \(\kappa\) of size \(n\), use the following
family for \(X(\kappa)\): \(\{\kappa^{[1]}\}\) for a regular
deep piece; \(\{\kappa^{[1]},\widehat\kappa\}\) for a buffer
piece, where \(\widehat\kappa\) is its containing cube at
level \(n+1\); and the bare \(\mathfrak M\)-cubes contained
in \(\kappa^{[1]}\cap\Omega_0\) for a layer-zero piece.
Set
\[
 \mathcal N(\kappa)
 =\{\kappa':\kappa'^{[1]}\cap\kappa^{[1]}\cap\Omega_0
                    \text{ has nonempty interior}\},
 \qquad
 X^{(2)}(\kappa)=\bigcup_{\kappa'\in\mathcal N(\kappa)}
                         X(\kappa'),
\]
with the union of the families. These are cells with
at most \(27\) and \(27N\) cubes, respectively, where
\(N\) is the piece-overlap bound. Indeed, all the
layer-zero cubes share a point with the central cube,
and a buffer cube meets its containing cube. For a
neighbour \(\kappa'\), the families of \(X(\kappa')\)
and \(X(\kappa)\) meet at a common interior point of
their defining regions. This proves connectedness of
the union. The piece geometry also gives
\[
 \kappa^{[2]}\cap\Omega_0\subset X(\kappa)^\natural,
 \qquad
 X(\kappa')^\natural\subset X^{(2)}(\kappa)^\natural
          \quad(\kappa'\in\mathcal N(\kappa)).
\]
Here and below intersections of regions are taken
cube by cube; all lattice sites in the corresponding
point intersection belong to this region.

For fixed \(L,r_{\rm w}\), a cell overlaps at most
\(n_1=n_1(L,r_{\rm w})\) cells. Indeed, cubes sharing a
point have indices differing by at most two, and at each
index their side lengths and corner positions have a
bounded number of choices. A connected family of at most
\(r_{\rm w}\) cubes containing a fixed cube is therefore
chosen from a bounded set of chains of length at most
\(r_{\rm w}-1\). Put \(n_{\rm w}=4n_1\).
There are at most \(n_{\rm w}^m\) tree walks of length
\(m\) with a prescribed first cell. The depth-first parent
maps number at most \(4^m\): their successive depth drops
are nonnegative integers with sum at most \(m-1\), encoded
by fewer than \(2m\) binary symbols. Each subsequent cell
has at most \(n_1\) choices.

For a region \(W\subset\Omega_0\), define
\(\mathfrak P_t^\natural((\Omega_n);
\alpha_0,\alpha_1,\gamma_0)\restriction W\)
by \cref{def:natural-regular-pair}, with \(W\) in place
of \(\Omega_0\), \(\mathcal L_n\cap W\) in place of
\(\mathcal L_n\), and
\(\{\mathsf q\in\mathsf Q_n^\natural:\mathsf q\subset W\}\)
in place of \(\mathsf Q_n^\natural\).
Thus the pair is given on the inner bare bonds of \(W\)
and has one decomposition \(\mathbf U=e^{i\xi A'}U\)
there, with \(U\) real. The curvature bounds are required
on inner plaquettes of \(W\); the bounds on \(A'\), its
covariant differences and \(\mathbf J\) are required
only on the bonds and triples lying in \(W\).
For every indicated cube \(\mathsf q\), the real gauge
witness in that definition must have
\(\mathsf q\subset\square\subset W\).
All bounds retain their original scale \(\xi_n\) and
their original strict or non-strict inequalities.
The real factor and its gauge witnesses need not extend
beyond \(W\). We call this the \emph{local natural class}.

On a region \(W\subset\Omega_0\), call a pair
\emph{piecewise local} if its restriction to
\(X(\kappa)^\natural\) lies in the local natural
class for every piece whose carrier lies in \(W\).
This is an open set of the bond coordinates
\((\mathbf U,\mathbf J)\). In fact, at a point of a
local natural class keep its real factor \(U\) fixed
and use a local logarithm of \(\mathbf U U^{-1}\).
The finitely many strict conditions on \(A'\) and
\(\mathbf J\) persist; the gauge conditions on \(U\)
remain unchanged. Thus the local class is open,
and the piecewise class is the finite intersection
of its inverse images under restriction. Restrictions
to smaller regions preserve the piecewise condition.
A pair in the local natural class is piecewise local.

Use the constants of the regular parametrix as follows:
\[
 \begin{gathered}
 \delta=\delta_{\mathrm{reg}},\qquad
 N=N_{\mathrm G},\qquad B_4=\max\{1,B_T/4\},\\
 C_q=2C_{\mathbb C}C_{\mathrm G}B_{\mathrm{loc}}
                       e^{2\rho+8},\qquad \rho=3L+6,\\
 q=q_{\mathfrak M}
       =C_q\mathfrak M^{-1}e^{-\delta\mathfrak M/(8L)}.
 \end{gathered}
\]
Here \(B_T=4B_{\mathrm{loc}}e^{3\delta}\), as in
\eqref{eq:regular-remainder-constants}, and
\(0<\delta\le1/2\). Put
\[
 S=c_S(\delta/4)=8L^3(2+16/\delta)^3,\qquad
 q_1=S q,\qquad C_j=96LB_4NS.
\]
We increase the lower bound on \(\mathfrak M\) so that
\[
 \mathfrak M\ge32L^4,\qquad
 R_1\mathfrak M\ge\mathsf r_S(\delta/4),\qquad
 \frac{\log L}{R_1\mathfrak M}\le\frac\delta4,\qquad
 Nq_1\le\tfrac12,\qquad q,q_1\le\mathfrak M^{-9}.
\]
The summation and scale-ratio conditions are
\eqref{eq:bg-thresholds}; in particular
\(\sum_z e^{-\delta d(y,z)/4}\le S\).
Further lower bounds will be imposed in the proof.
All are bounds by constants depending only on \(L\).

\begin{proposition}\label{prop:local-multiplier-trees}
\label{I-70-63}\label{I-70-67}\label{I-70-71}
\label{I-70-74}
For sufficiently small \(a_{\rm tr}>0\) and sufficiently
large \(\mathfrak M\), the following holds for natural
pairs with radii \(\alpha_0,\alpha_1,\gamma_0\le a_{\rm tr}\).
There are local holomorphic functions \(\mathsf j_\varpi\)
and local holomorphic operators \(\mathbb E_\Xi\), defined
below, such that
\[
 \mathsf j_*=\mathsf H^{\mathsf T}\mathbf J
       =\sum_\varpi\mathsf j_\varpi,\qquad
 \mathsf j^0=\sum_{|\varpi|=0}\mathsf j_\varpi,\qquad
 |\mathsf j_*-\mathsf j^0|_{-3;y}
              \le2C_j\gamma_0Nq_1.
\]
For the physical fifth system at the multiplier
\(\mathsf j_*\),
\[
 \mathbb E_*=\mathbb K^{[5]}
             [\mathbf U,\mathbf J,\mathsf j_*]^{-1}
             =\sum_\Xi\mathbb E_\Xi .
\]
Both sums converge absolutely. Each nonzero
\(\mathbb E_\Xi\) has a tree walk \(\omega(\Xi)\) of
cells with \(r_{\rm w}=27N\). For every such tree
\(\omega\),
\begin{equation}\label{eq:physical-tree-bound}
 \sum_{\Xi:\,\omega(\Xi)=\omega}
       |\mathbb E_\Xi\iota g|_y
 \le B_{\rm tr}(C_{\rm tr}\mathfrak M^{-3})^{|\omega|}
       \sup_z e^{-\delta d^{\rm tr}(\omega;y,z)/2}|g|_z,
 \qquad B_{\rm tr}=27N\cdot8B_4 .
\end{equation}
The terms and this estimate are defined on the
piecewise local class on \([\omega]\). The map
\(\iota\) and the seminorms may be taken at any
continuation of the local configuration to
\(\Omega_0\). The operators and their values on
their stated input supports do not depend on that
continuation. Restriction from a larger local domain
gives the same terms, and restriction of a global
pair gives the terms in the displayed inverse.
\end{proposition}

\begin{proof}
\emph{The local inverse factors.}
For a piece \(\kappa\), take the system on
\(\mathfrak S^\kappa\) with ambient region
\(W_\kappa=X(\kappa)^\natural\).
The proof of \cref{lem:regular-local-inverse}
applies to these local data. For a regular piece,
\(\kappa^{[2]}\in\mathsf Q_n^\natural\) and
\(\kappa^{[2]}\subset W_\kappa\). The local
natural-class condition supplies the real gauge
on a containing cube inside \(W_\kappa\).
Restrict it to \(\kappa^{[2]}\). The logarithm
of \(e^{i\xi\operatorname{Ad}(u)A'}e^{i\xi A}\),
as in that proof, has potential and first-difference
bounds at scale \(n+1\) at most
\(8L^2(\alpha_0+\alpha_1)\). Also
\(\Omega_0^\kappa\subset\kappa^{[3/2]}\), so the
local region is interior to \(W_\kappa\).
For a layer-zero piece the one-layer inverse
applies and the multiplier does not enter.
The touching plaquettes in the Hessian have all
corners in \(\kappa^{[2]}\cap\Omega_0\).
Thus enlarging the ambient region from \(W_\kappa\)
does not change this local operator.
All the other coefficients use bonds in its
local sequence and row values with basepoints in
\(\kappa^{[5/6]}\). The same local decay and
commutator estimates therefore hold.

More explicitly, for the lifted source cutoff
\(\theta_\kappa\star\), the maps
\[
 T_\kappa=\widetilde h_\kappa\mathbb E_\kappa
                 \iota(\theta_\kappa\star\,\cdot\,),
 \qquad
 R^\sharp_\kappa
     =-\mathbb C_{\widetilde h_\kappa}
                    \mathbb E_\kappa
                 \iota(\theta_\kappa\star\,\cdot\,)
\]
satisfy the estimates used in the parametrix proof:
\(T_\kappa\) has constant \(4B_4\), and
\(R^\sharp_\kappa\) has constant \(q\), both at
rate \(\delta\). The output support of the latter
and the input support of both lie in
\(\kappa^{[5/6]}\); their separation is
\(\mathfrak M/(8L)\). This follows from
\[
 [\mathbb K_\kappa,\widetilde h_\kappa]
       =\iota\mathbb C_{\widetilde h_\kappa},
 \qquad
 |\mathbb C_{\widetilde h_\kappa}\Psi|_y
   \le C\mathfrak M^{-1}
               \max_{d^\kappa(y,z)\le\rho_0}|\Psi|_z
\]
and the local inverse bound at rate \(2\delta\),
by the same separated-support composition argument.
The identities
\(\theta_\kappa\iota=\iota(\theta_\kappa\star)\)
and \(R_\kappa\iota=\iota R^\sharp_\kappa\)
are unchanged.

These assertions hold simultaneously at any
continuation of the local configuration. Indeed,
every transport in a source-cutoff commutator,
in \(\iota R^\sharp_\kappa\), or in a difference
of \(T_\kappa g\) has both endpoints within
\(\xi_{n+1}+2\xi\le s_n/6\) of
\(\kappa^{[5/6]}\); hence it lies in
\(\kappa^{[1]}\cap\Omega_0\subset W_\kappa\).
The source seminorm contains no transport.
Consequently no comparison of unrelated
configurations is used in composing the factors.

For \(\varpi=(\kappa_0,\ldots,\kappa_m)\), set
\[
 E_\varpi=
 \widetilde h_{\kappa_0}E_{\kappa_0}\theta_{\kappa_0}
                         R_{\kappa_1}\cdots R_{\kappa_m}.
\]
The preceding identities give, with \(y_0=y\),
\begin{equation}\label{eq:local-chain-factor}
 |E_\varpi\iota g|_y
 \le8B_4q^m\sup
   e^{-\delta\sum_{i=0}^m d(y_i,y_{i+1})}|g|_{y_{m+1}},
\end{equation}
where \(y_i\in\kappa_{i-1}^{[5/6]}\cap
\kappa_i^{[5/6]}\) for \(1\le i\le m\), and
\(y_{m+1}\in\kappa_m^{[5/6]}\), all points being
tops. The operator vanishes unless these
intersections contain tops. Its output has
basepoints in \(\kappa_0^{[2/3]}\), and it uses
the input only at basepoints in \(\kappa_m^{[1/12]}\).
These support assertions follow directly from
the first and last cutoffs.

The entries of each local system are holomorphic
on the open set of triples: the derivatives and
Hessian are polynomial in bond transports, while
the averaging coefficients are Taylor coefficients
of the holomorphic averaging chart. The inverse
is holomorphic where its determinant is nonzero.
The finite intersection of these open sets for
the pieces of \(\varpi\) contains all the prescribed
piecewise local data. Thus \(E_\varpi\) is
holomorphic there. This argument uses the local
operators themselves, so restriction to a smaller
carrier or to a global pair gives the same factors.

\emph{The multiplier terms.}
In this paragraph take all factors at multiplier
zero and set \(H_\varpi=[E_\varpi^{[5]}]_{A,8}\).
Define \(\mathsf j_\varpi(c)\) by
\[
 \xi_c^3\langle e,\mathsf j_\varpi(c)\rangle
   =\xi^3\sum_b\langle H_\varpi e_c(b),\mathbf J(b)\rangle,
\]
where the sum is over bonds with basepoint in
\(\kappa_0^{[2/3]}\). Their endpoints lie in
\(X(\kappa_0)^\natural\), so every value of
\(\mathbf J\) is given. The nondegenerate bilinear
form determines a unique element. It is holomorphic,
and vanishes unless the basepoint of \(c\) lies in
\(\kappa_m^{[1/12]}\).
Put
\[
 d(\varpi;c)=
 \inf_{y_i\in\kappa_{i-1}^{[5/6]}\cap\kappa_i^{[5/6]}}
       \sum_{i=1}^m d(y_i,y_{i+1}),
 \qquad y_{m+1}=\mathsf y(c),
\]
with value zero for \(m=0\). Then
\begin{equation}\label{eq:local-multiplier-term}
 \xi_c^3|\mathsf j_\varpi(c)|
 \le N^{-1}C_j\gamma_0q_1^m
                e^{-\delta d(\varpi;c)/2}.
\end{equation}
Here is the weighted summation. The eighth-slot
source \(e_c\) has size \(\xi_c|e|\), while the
field component has weight \(\xi_y\).
There are at most \(3(\xi_y/\xi)^3\) bare bonds
at \(y\), and \(|\mathbf J(b)|\le\gamma_0\xi_y^{-3}\).
Thus \eqref{eq:local-chain-factor} gives
\[
 |\xi_c^3\langle e,\mathsf j_\varpi(c)\rangle|
 \le24B_4\gamma_0q^m|e|
    \sum_{y_0,\ldots,y_m}
       \frac{\xi_c}{\xi_{y_0}}
       e^{-\delta\sum_{i=0}^m d(y_i,y_{i+1})}.
\]
The scale-ratio estimate bounds the quotient by
\(L\exp(\delta\sum_i d(y_i,y_{i+1})/4)\).
Retain rate \(\delta/2\) for \(d(\varpi;c)\)
and sum the remaining rate \(\delta/4\), successively,
over the \(m+1\) tops. Each sum is at most \(S\).
Finally
\(|Z|\le2\sup_{|e|\le1}|\langle e,Z\rangle|\)
gives \eqref{eq:local-multiplier-term} with the
displayed \(C_j\).

For fixed \(c\), at most \(N^{m+1}\) nonzero
sequences of length \(m\) occur, by the terminal
support and the piece-neighbour bound. This proves
absolute convergence,
\(|\mathsf j_*|_{-3;y}\le2C_j\gamma_0\), and the
asserted tail bound. Summing the already proved
inverse series and taking its weighted transpose
identifies the sum with \(\mathsf H^{\mathsf T}\mathbf J\).
Locally, at a row \(c\) in \(\kappa^{[5/6]}\), define
\[
 \mathsf j^0(c)=
   \sum_{\kappa':\,\text{base}(c)\in\kappa'^{[1/12]}}
                          \mathsf j_{(\kappa')}(c).
\]
All these \(\kappa'\) are in \(\mathcal N(\kappa)\);
their carriers are in \(X^{(2)}(\kappa)^\natural\),
and at most \(N\) occur. Therefore this is a
holomorphic local function satisfying the sharper
bound \(\xi_c^3|\mathsf j^0(c)|\le C_j\gamma_0\).

\emph{The vertex and the resolvent.}
For \(c\) above level zero let \(B_c\) be the
symmetric bilinear coefficient defined by
\(C_{\mathrm{phys}}^{(2)}(A)(c)=\xi_c B_c(A,A)\).
Here the physical remainder is \(\xi_c^{-1}\mathsf C(A)(c)\),
where \(\mathsf C\) is the dimensionless remainder
of \cref{I-87-77}. Consequently
\[
 \mathsf C^{(2)}(A)(c)=\xi_c^2B_c(A,A).
\]
Polarizing the quadratic action gives
\[
 \xi^3\langle X,\Delta^{(2)}[\mathsf j']A(b)\rangle
   =2\sum_c\xi_c^4\langle B_c(X_b,A),\mathsf j'(c)\rangle.
\]
Only bonds in the double block \(\mathsf Y_c\)
enter the \(c\)-summand. With \(F=\Delta^{(2)}
[\mathsf j'](A-Dv)\), define \(\mathbb W[\mathsf j']\Psi\)
to have components \(f_1=F,X_5=-F\) and all
other components zero. Thus
\[
 \iota\mathbb W[\mathsf j']\Psi
        =(F,0,0,0,-D^{\mathsf T}F,0,0,0),\qquad
 \mathbb K[\mathsf j+\mathsf j']
        =\mathbb K[\mathsf j]-\iota\mathbb W[\mathsf j'].
\]
Let \(\mathrm{Pos}(c)\) be the positions of the
bare sites in its double block \(\mathsf Y_c\).
For an interior row these are its two endpoints;
for an outward row they are its inner endpoint
and the \(L^3\) tops in its outer block. Their
scales lie between \(\xi_c/L\) and \(\xi_c\),
and their mutual distances are at most
\(\rho_0=3L+6\). A top belongs to at most twelve
such sets: six incident rows at its own level,
and six at the next level incident to its
unique containing block. The vertex obeys
\begin{equation}\label{eq:multiplier-local-vertex}
 |\mathbb W[\mathsf j']\Psi|_y
 \le C_{\mathrm{vert}}
   \max_{c:\,y\in\mathrm{Pos}(c)}\xi_c^3|\mathsf j'(c)|
   \max_{\substack{c:\,y\in\mathrm{Pos}(c)\\
                   z\in\mathrm{Pos}(c)}}|\Psi|_z,
 \qquad d(y,z)\le\rho_0 .
\end{equation}
To verify the estimate, write
\(\beta_{bb'}(X,X')=B_c(X_b,X'_{b'})\).
Use the constant \(C_Q\) of \cref{I-87-77}. Divide its difference
estimate for \(K_{sA}-K_0\) by \(|s|\) and let \(s\to0\).
Applied to an input \(X_b\), it gives
\[
 \left|\left.\frac d{ds}\right|_{s=0}
            K_{sA}(c,b)X\right|
 \le C_Q\xi^3\xi_c^{-1}|X|
                         \sup_{\mathsf P_c}|A|.
\]
The derivative on the left is \(2\xi_c^2 B_c(X_b,A)\).
Division by \(\xi_c^2\), rather than by \(\xi_c\), gives
\[
 |2B_c(X_b,A)|\le C_Q(\xi/\xi_c)^3
                                  |X|\sup|A|.
\]
Choose the phases of the components of \(A(b')\)
separately in a fixed Pauli basis. The resulting
sum of the absolute coordinate coefficients gives
\(\sum_{b'}|\beta_{bb'}|
 \le(27/2)n_{\mathfrak g}^2C_Q(\xi/\xi_c)^3\),
where \(n_{\mathfrak g}\) bounds the fixed
finite-dimensional coordinate norms. In the
polarized identity the factor
\(\xi^{-3}\xi_c^4(\xi/\xi_c)^3=\xi_c\).
At most twelve rows contribute at a given top,
and within each \(\mathrm{Pos}(c)\),
\(\xi_c/L\le\xi_z\le\xi_c\). Applying this to
\(A-Dv\) proves \eqref{eq:multiplier-local-vertex},
for example with
\(C_{\mathrm{vert}}=1296n_{\mathfrak g}^2L^3C_Q\).
The same formula proves that the vertex is local
to its active double blocks and holomorphic.
In particular an active row based in
\(\kappa'^{[1/12]}\) has
\(\mathsf Y_c\subset\kappa'^{[5/12]}\cap\Omega_0
 \subset X(\kappa')^\natural\).
No extension enters its coefficients or its
nonzero transports.

Choose \(a_{\rm tr}\) so that \(2C_j a_{\rm tr}\)
is within the local-inverse radius. Write
\(\delta\mathsf j=\mathsf j_*-\mathsf j^0\) and
\(\mathbb E_0=\mathbb K[\mathsf j^0]^{-1}\).
By the tail estimate and \eqref{eq:multiplier-local-vertex},
\(\mathbb W[\delta\mathsf j]\mathbb E_0\iota\)
has decay rate \(\delta\) and norm at most
\[
 2C_{\mathrm{vert}} C_j\gamma_0Nq_1e^{\rho_0}B_5,\qquad B_5=16NB_4.
\]
Increase \(\mathfrak M\) to make this at most \(1/2\).
The resolvent identity then gives the convergent series
\[
 \mathbb E_*\iota
   =\mathbb E_0\iota\sum_{n\ge0}
       \bigl(\mathbb W[\delta\mathsf j]\mathbb E_0\iota\bigr)^n .
\]

\emph{Decorations and tree decay.}
For
\(\Xi=(\varpi^0,\varpi'_1,\varpi^1,\ldots,
\varpi'_n,\varpi^n)\), with
\(|\varpi'_l|\ge1\), define
\[
 \mathbb E_\Xi
   =E_{\varpi^0}\iota\mathbb W[\mathsf j_{\varpi'_1}]
      E_{\varpi^1}\cdots
       \iota\mathbb W[\mathsf j_{\varpi'_n}]E_{\varpi^n}.
\]
The trunk factors are evaluated at the local
\(\mathsf j^0\) on all their row supports, and
zero elsewhere. List first all their
\(X^{(2)}\)-cells in order. Then list the branches
\(\varpi'_n,\ldots,\varpi'_1\), each in reverse
order, using their \(X\)-cells. Consecutive trunk
cells have the preceding trunk cell as parent;
the first cell of branch \(l\) has the last
trunk cell of \(\varpi^{l-1}\) as parent, and
the other branch cells have the preceding
branch cell as parent. This defines \(\omega(\Xi)\).

Every factor is defined and holomorphic on the
piecewise local class on its carrier. The
\(X^{(2)}\)-cells contain all the extra one-piece
carriers needed for \(\mathsf j^0\). The vertex
uses only the last branch piece and its
double blocks. This proves the stated local
realization of \(\mathbb E_\Xi\).
If the product is nonzero, consecutive pieces
within every branch or trunk meet at a top.
At a vertex, a contributing active row has
its position in the interiors of the enlarged
last preceding trunk piece, first succeeding
trunk piece and last branch piece. Indeed,
its double block has diameter \(2\xi_c<s_n/3\)
at each of the relevant piece scales; the
cutoff supports are \([1/12]\) and \([2/3]\).
Thus the two trunk \(X^{(2)}\)-cells overlap,
and the first branch cell is contained in
and overlaps its prescribed parent. The parent
order is depth first by construction, so this
is a tree walk.

Write \(m_l=|\varpi^l|\), \(m'_l=|\varpi'_l|\).
Iterating \eqref{eq:local-chain-factor},
\eqref{eq:local-multiplier-term} and
\eqref{eq:multiplier-local-vertex} gives
\[
 |\mathbb E_\Xi\iota g|_y
 \le8B_4\prod_{l=0}^nq^{m_l}
       \prod_{l=1}^n(C_{\mathrm{br}}\gamma_0q_1^{m'_l})
       \sup e^{-\delta\Sigma_{\rm tr}
                -\frac\delta2\sum_l d(\varpi'_l;c_l)}
                      |g|_z,
 \qquad C_{\mathrm{br}}=8B_4 C_{\mathrm{vert}} e^{\rho_0}C_j .
\]
Here \(\Sigma_{\rm tr}\) is the sum of the
successive trunk distances, including the
jumps across vertices. Each such jump is
at most \(\rho_0\); inserting its exponential
cost explains the factor \(e^{\rho_0}\) in \(C_{\mathrm{br}}\).
Join successive tops by segment unions of
weighted length arbitrarily close to their
distance, and join each vertex top to the
position of its active row. Attach the branch
chains at that position. Their connected
union meets every cell of \(\omega(\Xi)\);
the condition \(m'_l\ge1\) ensures that the
first cell of each branch is met. Subadditivity
of weighted length gives
\[
 d^{\rm tr}(\omega(\Xi);y,z)
 \le\Sigma_{\rm tr}
       +\sum_l d(\varpi'_l;c_l)+n\rho_0 .
\]
The exponential above is consequently bounded
by \(e^{n\rho_0/2}e^{-\delta d^{\rm tr}/2}\).
Moreover
\[
 |\omega(\Xi)|=\sum_{l=0}^n m_l+\sum_{l=1}^n m'_l+2n,
 \qquad
 q_1^{m'}\le\mathfrak M^{-3(m'+2)}
                         \quad(m'\ge1).
\]
Together with \(q^{m_l}\le\mathfrak M^{-3m_l}\)
and \(\gamma_0\le1\), this proves
\[
 |\mathbb E_\Xi\iota g|_y
 \le8B_4\bigl(e^{\rho_0}\max\{1,C_{\mathrm{br}}\}
                     \mathfrak M^{-3}\bigr)^{|\omega|}
       \sup_z e^{-\delta d^{\rm tr}(\omega;y,z)/2}|g|_z .
\]

For a fixed tree with \(m\) edges, at most
\(2^m(27N)^{m+1}\) decorations give that tree.
There are \(2^m\) ways to cut its ordered list
into the trunk and reversed-branch blocks.
A prescribed \(X(\kappa)\) has at most \(27\)
preimages; a prescribed \(X^{(2)}(\kappa)\)
has at most \(27N\). For the latter claim the
family contains the distinct cube
\(\kappa^{[1]}\) for a regular piece and
\(\kappa\) for a layer-zero piece, which
determines the piece. Thus one may take
\[
 C_{\rm tr}=54Ne^{\rho_0}\max\{1,C_{\mathrm{br}}\}
\]
in \eqref{eq:physical-tree-bound}.
Choose the last lower bound on \(\mathfrak M\)
so that \(n_{\rm w}C_{\rm tr}\mathfrak M^{-3}\le1/2\),
with \(r_{\rm w}=27N\).
There are finitely many first cells, so the tree
bound gives absolute convergence in each output
seminorm. These seminorms separate the points of
the finite-dimensional system, and \(\iota\) is
onto, so the operator series converges absolutely.
Expanding \(\mathbb E_0\) and \(\delta\mathsf j\)
in the resolvent series now gives exactly these
decorated products. Absolute convergence justifies
the rearrangement and identifies their sum with
\(\mathbb E_*\).
\end{proof}

\section{Fluctuation covariances and their local expansion}\label{sec:fluctuation-local-expansion}

We measure lengths at level $k$, where $0\le k<K$, so that
$\xi=L^{-k}$ and $\xi_n=L^{n-k}$. Let $(\Omega_n)_{n\le k}$
be a separated sequence with parameters $\mathfrak M,R_1$, and let
\[
 \Omega_{k+1}=\bigcup_{y\in\Theta'}\Delta(y)\subset\Omega_k,
 \qquad
 \operatorname{dist}_\infty(\Omega_{k+1},\Omega_k^c)
       >R_1\mathfrak M,
\]
where $\Omega_{k+1}$ is a union of $\mathfrak M$-cubes at level
$k+1$. Write $\mathbf B$ and $\mathbf B^+$ for the determining
sets of the shorter and longer sequences, with all their bonds.
The pair $(\mathbf U,\mathbf J)$ belongs to the natural regular
class of the longer sequence, with radii
$\alpha_0,\alpha_1,\gamma_0$.

At level $k$ the spacing is $1$. For a coarse bond
$c=\langle c_-,c_-+\epsilon L e_\mu\rangle$,
put $h=(L-1)/2$. Its \emph{pivot bond} is
\[
 b_0(c)=\langle c_-+\epsilon h e_\mu,
                   c_-+\epsilon(h+1)e_\mu\rangle.
\]
It has the midpoint of $c$, and reversal reverses the pivot.
Write $\mathsf P_c$ for the inner bonds of the double block
$B(c_-)\cup B(c_+)$. Write $\operatorname{Ax}(B(y))$
for the positive bonds of the coordinate tree $\mathcal T(y)$
defined before \cref{I-15-117}.

\begin{lemma}\label{I-5-15}
For coarse bonds $c,c'$, one has
\[
 b_0(c')\in\mathsf P_c
       \quad\Longleftrightarrow\quad c'\in\{c,\bar c\}.
\]
The pivot of a positive coarse bond is positive. Consequently,
for a set $\Gamma$ of positive coarse bonds, the pivots $b_0(c)$,
$c\in\Gamma$, are pairwise distinct, and $b_0(c)$ is the only
one of them in $\mathsf P_c$.
\end{lemma}

\begin{proof}
The endpoints of $b_0(c')$ lie in the distinct blocks $B(c'_-)$
and $B(c'_+)$. If this pivot belongs to $\mathsf P_c$, both
endpoints lie in $B(c_-)\cup B(c_+)$. The blocks partition the
fine lattice, so $\{c'_-,c'_+\}=\{c_-,c_+\}$. Thus $c'=c$ or
$c'=\bar c$. Conversely, the pivot of either orientation has
both endpoints in the double block of $c$.

The explicit pivot formula gives it the direction and orientation
of its coarse bond. If $c,c'\in\Gamma$ and
$b_0(c')\in\mathsf P_c$, the alternative $c'=\bar c$ is excluded
by positivity. Hence $c'=c$. In particular equality of two
selected pivots implies equality of their coarse bonds.
\end{proof}

Put $\Theta=B(\Theta')$, and let $\Gamma'$ consist of the positive
level-$(k+1)$ bonds with an end point in $\Theta'$. Write
$\mathfrak X(\Theta)$ for the level-$k$ bonds with an end point
in $\Theta$, and $\check{\mathfrak X}(\Theta)$ for this set with
the pivot bonds $b_0(c)$, $c\in\Gamma'$, and their reverses removed.
For any subset $\Theta''\subset\Theta'$, set
\[
 Y=\check{\mathfrak X}(B(\Theta'')),\qquad
 Y^{\rm ax}=Y\setminus\bigcup_{y\in\Theta'}
      \bigl(\operatorname{Ax}(B(y))\cup
                    \overline{\operatorname{Ax}(B(y))}\bigr).
\]
The pivots removed in the definition of $Y$ are those associated
with the coarse bonds having an end point in $\Theta''$.
Nevertheless $Y\subset\check{\mathfrak X}(\Theta)$: a pivot
associated with any other coarse bond has neither end point in
$B(\Theta'')$. No cube condition is imposed on $\Theta''$.
All spaces in this section use positive bonds as coordinates;
their bilinear pairings and kernels use counting measure.

Let $V=\mathcal M^k\mathbf U$. The following one-step maps use
the plain average at $V$:
\begin{align*}
 Q_VB(c)&=\frac1{iL}\frac d{dt}\Big|_0
   \log\bigl(\mathcal M(e^{itB}V)(c)\mathcal M(V)(c)^{-1}\bigr),\\
 \widetilde Q_V(B)(c)&=\frac1i
   \log\bigl(\mathcal M(e^{iB}V)(c)\mathcal M(V)(c)^{-1}\bigr),\\
 C_V^{(2)}(B_1,B_2)&=\tfrac12\partial_s\partial_t\Big|_{s=t=0}
                          \widetilde Q_V(sB_1+tB_2).
\end{align*}
If $\kappa_V(c)$ is the coefficient of $B(b_0(c))$ in $Q_VB(c)$,
the elimination map $\mathsf C_V$ keeps the nonpivot coordinates
and sets
\[
 (\mathsf C_V\check B)(b_0(c))
     =-\kappa_V(c)^{-1}(Q_V\check B^0)(c).
\]
Thus $Q_V\mathsf C_V=0$. Here and below a superscript $0$ on a
field means extension by zero. Let
\[
 \mathsf b_V(B)(y,x)=\frac1i\frac d{ds}\Big|_0
       \log\bigl((e^{isB}V)(y,x)V(y,x)^{-1}\bigr),
       \qquad x\in B(y),
\]
where the contour variable is the one used in the mean chart,
and define
\[
 (B_1,\mathsf G_V^{(2)}B_2)
   =\sum_{y\in\Theta'}\sum_{x\in B(y)\setminus\{y\}}
          \langle\mathsf b_V(B_1)(y,x),
                    \mathsf b_V(B_2)(y,x)\rangle .
\]
The right inverse $h_V$ of $LQ_V=\widetilde Q_V'(0)$ is supported on the pivots and has value
$L^{-1}\kappa_V(c)^{-1}\nu(c)$ at $b_0(c)$. Put
$\mathsf D_V^{(2)}=h_VC_V^{(2)}$.

The extension $\mathsf J$ maps a field on $\mathfrak X(\Theta)$
to the row space of $\mathbf B$. These are interior top rows,
with $\xi_c=\mathsf a(c)=1$, so its transpose is restriction.
For a row field $j_1$, define the symmetric operator
$\mathsf E^{(2)}[j_1]$ by
\[
 (B_1,\mathsf E^{(2)}[j_1]B_2)
       =2(\mathsf J\mathsf D_V^{(2)}(B_1,B_2),j_1)_{\mathbf B}.
\]
The operators $Q,\mathsf G_1,\mathsf H_1$ in the next display are
the straightened operators of the shorter sequence. With
$j_* =\mathsf H^{\mathsf T}\mathbf J$ and
$j_{1*}=\mathsf H_1^{\mathsf T}\mathbf J$, set
\begin{equation}\label{eq:cv-form}
 \Delta_k=\mathsf J^{\mathsf T}
       (Q\mathsf G_1Q^{\mathsf T})^{-1}\mathsf J
                 -1-\mathsf E^{(2)}[j_{1*}].
\end{equation}
For $\mathsf Y=Y$ or $Y^{\rm ax}$ let
$\mathsf C=\mathsf C_V\chi_{\mathsf Y}^{\mathsf T}$, where
$\chi_{\mathsf Y}^{\mathsf T}$ extends by zero, and put
\[
 T_Y=\mathsf C^{\mathsf T}(\mathsf G_V^{(2)}+\Delta_k)\mathsf C,
 \qquad
 T_Y^{\rm ax}=\mathsf C^{\mathsf T}\Delta_k\mathsf C.
\]
The map $\mathsf C$ in the two formulas has domain $Y$ and
$Y^{\rm ax}$, respectively. We denote their inverses by
$C^{(k)}(Y)$ and $\widetilde C^{(k)}(Y)$.
Formula \eqref{eq:cv-form} defines the form used here, including
away from a critical pair. In particular its first term uses
the inverse of $Q\mathsf G_1Q^{\mathsf T}$.

\begin{theorem}\label{I-73-90}
There are constants $a_{\rm cv},\delta_{\rm cv},C_{\rm cv},B_0,r_{\rm cv}^{\rm sep}>0$
and a power $M_{\rm cv}$ of $L$, depending only on $L$, with the
following property. Suppose $\mathfrak M,R_1$ are powers of $L$,
\[
 \mathfrak M\ge M_{\rm cv},\quad R_1\ge L,\quad
 R_1\mathfrak M\ge r_{\rm cv}^{\rm sep},\quad
 \ell_k\ge2L\mathfrak M,\quad
 0<\alpha_0,\alpha_1,\gamma_0\le a_{\rm cv}.
\]
All the maps above are defined, every $\kappa_V(c)$ is invertible,
and $T_Y,T_Y^{\rm ax}$ are invertible. Both covariance operators
are symmetric for their bilinear pairings. Either covariance,
denoted by $\mathcal C$, has a kernel satisfying
\begin{equation}\label{eq:cv-total-decay}
 |\mathcal C(b,b')|\le B_0e^{-\delta_{\rm cv}|b_--b'_-|_1}.
\end{equation}
Its operator norms for the supremum norm, the $\ell^2$ norm of
matrix norms, and the Hermitian norm of bond fields are at most
$B=6B_0\sigma(\delta_{\rm cv})^3$.

There is an absolutely convergent expansion
\[
 \mathcal C=\sum_\omega\mathcal C_\omega
\]
over tree walks of cells of the longer sequence. Their families
have at most $r_{\rm cv}=27N^2$ cubes, where $N$ is the overlap
constant of the regular pieces. If $\omega=(X_0,\ldots,X_m)$,
and $X_t$ is the last cell of its designated trunk, then
$\mathcal C_\omega(b,b')=0$ unless $b_-\in X_0$ and
$b'_-\in X_t$. Each term is a holomorphic function of the data
on its carrier $[\omega]$, on the open piecewise local class
defined in the proof. This class contains the local natural
class of the longer sequence; restriction of a global pair
gives the corresponding global term. Uniformly on that class,
\begin{equation}\label{eq:cv-tree-decay}
 |\mathcal C_\omega(b,b')|
 \le B_0(C_{\rm cv}\mathfrak M^{-1/2})^m\mathfrak M^{-m/2}
   e^{-\delta_{\rm cv}d^{\rm tr}_{\mathbf B^+}
                    (\omega;y(b),y(b'))/2},
\end{equation}
where $y(b)=b_-$ when $b_-\in\Theta$, and $y(b)=b_+$ otherwise.
There are at most $n_{\rm cv}^m$ such tree walks with a prescribed
first cell, for a constant $n_{\rm cv}=n_{\rm cv}(L)$, and the
constants can be chosen so that
\[
 M_{\rm cv}\ge(2C_{\rm cv}n_{\rm cv})^2,
 \qquad
 \sum_{|\omega|=m}|\mathcal C_\omega(b,b')|
       \le2^{-m-1}B_0e^{-\delta_{\rm cv}|b_--b'_-|_1}.
\]
Every partial sum over a set of tree walks satisfies
\eqref{eq:cv-total-decay} and the same three norm bounds.
The constants are independent of the number of layers, their
depths, $Y$, the volume, and the pair.
\end{theorem}

\begin{proof}
We first establish a gap at the flat pair. We then append the
fluctuation coordinate to the system already inverted in
\cref{I-65-83}. Its local inverse has the same cutoff expansion.
The final expansion inserts the two current multipliers into
that system. Constants denoted by $C(L)$ below may be enlarged;
all choices precede the choice of the pair or the regions.
Take \(r_{\rm cv}^{\rm sep}\) at least the previously required
background gap thresholds and
\(\max\{\mathsf r_S(\delta_{\rm reg}/8),
8\log L/\delta_{\rm reg}\}\).

\subsection*{The one-step coefficients}
On a double block $\mathsf D_c=B(c_-)\cup B(c_+)$ write
$\mathsf P_c$ for its bonds. The one-step chart for the mean
gives, on a fixed ball $|\widetilde a|<3r_*/4$ at
$V=e^{i\widetilde a}$, the factorization
\[
 \mathcal M(e^{iB}V)(c)\mathcal M(V)(c)^{-1}
   =e^{i\sigma(c_-)}e^{iQ_{\rm str}(V,B)(c)}
                   e^{-i\operatorname{Ad}(\mathcal M V(c))\sigma(c_+)}.
\]
Each of the first two chart functions is bounded by
$|B|/r_*$ on a smaller ball, and the last adjoint has norm at
most $e^2$. For $|B|<r_1=r_*/40$ the sum of the three exponent
norms is at most $10|B|/r_*<1/4$. The logarithm estimate and
Cauchy's inequality therefore give
\[
 |LQ_V(c,b)|,\ |\mathsf b_V(y,x;b)|\le r_1^{-1},\qquad
 |C_V^{(2)}(B_1,B_2)(c)|
     \le4r_1^{-2}\sup_{\mathsf P_c}|B_1|
                         \sup_{\mathsf P_c}|B_2|.
\]
The coefficients are holomorphic in $\widetilde a$. Their
differences from the coefficients at $1$ are at most
$C(L)|\widetilde a|$. At $1$, $\kappa_1(c)=L^{-3}$ and
$\mathsf b_1(B)(y,x)=\bar B(y,x)$, the arithmetic contour mean.
After reducing a fixed radius $s_0(L)>0$,
\[
 |\kappa_V(c)^{-1}|\le2L^3,\qquad
 |\kappa_V(c)^{-1}-L^3|\le C(L)|\widetilde a|.
\]
These facts also hold after a real gauge transformation, since
its adjoints preserve the norms.

A bond belongs to at most six double blocks, each containing
at most $6L^3$ positive bonds. Consequently the kernels of
$\mathsf C_V$, $\mathsf C_V^{\mathsf T}$ and
$\mathsf G_V^{(2)}$ are bounded by a constant $C_K(L)$;
the kernel of $\mathsf E^{(2)}[j_1]$ is at most
$C_K\max|j_1(b_0(c))|$. Every such kernel vanishes unless its
two bonds belong to a common $\mathsf P_c$. A column of
$\mathsf C_V$ has at most seven entries and a row at most
$6L^3$ entries. These counts follow directly from pivot
elimination. For the gauge kernel use
\[
 \mathsf G_V^{(2)}(b,b')
    =\sum_{y,x}\mathsf b_V(y,x;b)^{\mathsf T}
                         \mathsf b_V(y,x;b');
\]
only one block and fewer than $L^3$ sites contribute. For
$\mathsf E^{(2)}$ there are at most six contributing coarse
bonds. The same formulas give $C(L)|\widetilde a|$ bounds on
the differences of $\mathsf C_V,\mathsf C_V^{\mathsf T}$ and
$\mathsf G_V^{(2)}$ from their flat values.

\subsection*{The flat gap and the additional coordinate}
At the flat pair put $M=Q\mathsf G_0Q^{\mathsf T}$ and, for
$B=\mathsf C_1\check B$ extended by zero, put
$A=\mathsf G_0Q^{\mathsf T}M^{-1}\mathsf JB$.
Then $QA=\mathsf JB$, and the existing elimination identity gives
\begin{equation}\label{eq:cv-flat-energy}
 (B,(\mathsf J^{\mathsf T}M^{-1}\mathsf J-1)B)
    =(A,\Delta_{\Omega_0}(1)A)_\xi
                     +\|\mathsf RD^{\mathsf T}A\|_\xi^2.
\end{equation}
The Hessian here is the inner-plaquette Hessian used to define
the constrained inverse. At the flat field its quadratic form
equals the sum of squared curls over the inner bare plaquettes
of \(\Omega_0\).

Let $\mathcal P'$ be the level-$k$ plaquettes with a corner at
distance at most $L$ from $\Theta$. Their corners lie in
$\Omega_k$ by the $2L$ margin of the fluctuation structure.
Iterating the flat plaquette-average identity gives
\[
 \sum_{p'\in\mathcal P'}|\bar\partial QA(p')|^2
      \le \xi^3\sum_{p\subset\Omega_0}|\partial A(p)|^2
      \le(A,\Delta_{\Omega_0}(1)A)_\xi.
\]
Indeed, set $a_0=\xi A$, $a_{n+1}=\mathbf Q^sa_n$.
Each one-step plaquette estimate contributes a factor $L$;
after $k$ steps its factor is
$L^k\xi^4=\xi^3$. Composed blocks keep all contributing bare
plaquettes inside $\Omega_0$. At $k=0$ this is the identity
$QA=A$. On the bonds of $\mathcal P'$, $QA=\mathsf JB=B$.

On one double block consider the quadratic form
\[
 n_c(B)=\sum_{p\subset\mathsf D_c}|\bar\partial B(p)|^2
       +\sum_{y=c_\pm}\sum_{x\in B(y)\setminus\{y\}}
                                  |\bar B(y,x)|^2
       +|Q_1B(c)|^2.
\]
For its axial version omit the middle sum and restrict $B$ to
vanish on the two box trees. Both forms are positive definite.
If their value is zero, plaquette telescoping on the box gives
$B=d\nu$. In the first case the contour means, and in the
second case the tree paths, make $\nu$ constant on each block.
The remaining constraint is
$Q_1B(c)=L^{-1}(\nu_+-\nu_-)$, so $B=0$.
Translation reduces the norm comparison to six finite-dimensional
forms: three directions and the ordinary and axial cases.
Thus, for some $0<C_L<\infty$,
\[
 \sum_{b\in\mathsf P_c}|B(b)|^2\le C_L n_c(B),
\]
and the same inequality holds for the axial form on its domain.

For $B=\mathsf C_1\check B$ the constraint term vanishes.
On a block whose center is outside $\Theta'$, $B$ is zero; in
the axial case it vanishes on both relevant trees. Every bond
is covered by a double block, while each plaquette and each
block is counted at most six times. Summing the last inequality
and using \eqref{eq:cv-flat-energy} yields
\begin{equation}\label{eq:cv-flat-gap}
 T_{Y,0}\ge\gamma I,\qquad T_{Y,0}^{\rm ax}\ge\gamma I,
       \qquad \gamma=(6C_L)^{-1}.
\end{equation}
We used $\|\mathsf C_1\check B\|_2\ge\|\check B\|_2$.
The axial inequality has its own null-space argument above;
no gauge term is discarded from the ordinary inequality.

Append $\check B$ to $\Psi=(A;v,w,\sigma;p,q,r;\omega)$.
The first seven equations are those of $\mathbb K^{[5]}(j)$,
and the last two are
\begin{equation}\label{eq:cv-nine-system}
 QA-\mathsf J\mathsf C\check B=f_8,\qquad
 \mathsf C^{\mathsf T}
    \bigl[(\mathsf G_V^{(2)}-1-\mathsf E^{(2)}[j_1])
                    \mathsf C\check B+\mathsf J^{\mathsf T}\omega\bigr]=f_9.
\end{equation}
The axial system uses $Y^{\rm ax}$ and omits $\mathsf G_V^{(2)}$.
We call the systems $\mathbb K^{[9]}$ and $\mathbb K^{[9a]}$.
For either one use
\[
 |(\Psi,\check B)|_y=\max(|\Psi|_y,|\check B|_{-1;y}),\quad
 |(g,f_9)|_y=\max(|g|_y,|f_9|_{-3;y}),\quad
 \iota(g,f_9)=(\iota g,f_9).
\]
A new bond has position $b_-$, scale $1$, and midpoint as base
point for cutoffs.

The previously proved elimination for $\mathbb K^{[5]}$ now
leaves only one equation. With
$M_\natural=Q\mathsf G_\natural Q^{\mathsf T}$, it is
\begin{align}
 T(j,j_1)\check B
   &=f_9-\mathsf C^{\mathsf T}\mathsf J^{\mathsf T}
       M_\natural^{-1}(f_8-Q\mathsf G_\natural\mathsf F F),
       \label{eq:cv-schur}\\
 \Psi&=\mathbb E^{[5]}(j)
                   (F,f_8+\mathsf J\mathsf C\check B),\notag
\end{align}
where $T(j,j_1)=\mathsf C^{\mathsf T}
(\mathsf G_V^{(2)}+\mathsf J^{\mathsf T}M_\natural^{-1}
\mathsf J-1-\mathsf E^{(2)}[j_1])\mathsf C$.
The axial formula omits the gauge term. Thus the enlarged
system is invertible exactly when this Schur operator is, and
its $(\check B,9)$ block is $T(j,j_1)^{-1}$.
At $(j_*,j_{1*})$ the background identities give
$\mathsf G_\natural=\mathsf G_1$, so the Schur operators are
precisely $T_Y,T_Y^{\rm ax}$.

At the flat pair the existing inverse estimate and the finite
range of $\mathsf C_1$ give
$|T_0(b,b')|\le B_Te^{-\delta_Td(b_-,b'_-)}$, uniformly in
the structure. This follows by expanding
$\mathsf C_1^{\mathsf T}(G_0^{(2)}-1+
\mathsf J^{\mathsf T}M^{-1}\mathsf J)\mathsf C_1$:
there are at most $7^2$ entries from the two elimination
columns, and their positions move a bounded distance $C(L)$.
The axial calculation deletes $G_0^{(2)}$.

Apply the conjugation argument proving \eqref{I-14-34} to
$T\ge\gamma I$. Here the error row and column sums are at
most $2aB_TS/\delta_T$ for $a\le\delta_T/2$, where
$S=\sup_b\sum_{b'}e^{-\delta_Td(b_-,b'_-)/2}<\infty$.
Taking also $a\le\gamma\delta_T/(4B_TS)$ gives
$|T^{-1}(b,b')|\le(2/\gamma)e^{-ad(b_-,b'_-)}$.
The real symmetric gap extends to the Hermitian form on
complex fields. The top-lattice summation bound makes $S$
uniform. Insert this estimate in \eqref{eq:cv-schur} and use
the existing $\mathbb E^{[5]}_0$ estimate. The two finite-range
maps $\mathsf J\mathsf C_1$ and
$\mathsf C_1^{\mathsf T}\mathsf J^{\mathsf T}$ have bounded
row sums; their rows all have scale $1$. We obtain
\[
 \mathbb E^{[9]}_0\iota,\ \mathbb E^{[9a]}_0\iota
                         \in\mathfrak O(\delta_c,B_c)
\]
with $\delta_c,B_c$ depending only on $L$.

\subsection*{Local inversion and gluing}
In a small trivial gauge the old eight equations satisfy the
perturbation estimate used in \cref{I-65-40}. The additional
terms are $\mathsf J(\mathsf C_V-\mathsf C_1)\check B$ in
slot $8$, and the differences of
\[
 \mathsf C_V^{\mathsf T}(\mathsf G_V^{(2)}-1)\mathsf C_V,
 \qquad \mathsf C_V^{\mathsf T}\mathsf E^{(2)}[j_1]\mathsf C_V,
 \qquad \mathsf C_V^{\mathsf T}\mathsf J^{\mathsf T}
\]
in slot $9$. The coefficient bounds above give
\begin{equation}\label{eq:cv-local-defect}
 \mathbb K_0-\mathbb K=\iota_0 V_c,
 \qquad
 |V_c\Phi|_y\le C_V\varepsilon
                  \max_{d(y,z)\le\rho_c}|\Phi|_z.
\end{equation}
Here the triple is $\varepsilon$-small in the earlier sense,
$|j_1|_{-3;y}\le\varepsilon$, and one may take
$\rho_c=19L+6$. In a product of three new kernels the
intermediate indices admit at most $1512L^6$ choices, and
the total displacement of their positions is at most $12L$.
Thus no count depends on the depth or the number of layers.
The already proved averaging estimate sends this small gauge
to $V=e^{i\widetilde a}$ with
$|\widetilde a|\le C_{\rm ia}L\varepsilon$ on each double
block; choose $\varepsilon$ so that this is at most $s_0$.
Reducing $\varepsilon_c(L)>0$ further makes
$C_V\varepsilon_c B_c e^{\rho_c}\le1/4$.
The same Neumann argument as before, with the unchanged forcing
injection in the ninth coordinate, gives a uniform local
inverse for both systems.

The cutoff calculation also extends without a new differential
estimate. Multiplication by $h$ uses the midpoint of the new
bond. In each new kernel the commutator inserts
$h(\hat b')-h(\hat b)$; in a product of three kernels there are
three such links. Hence, for a cutoff varying on scale
$\vartheta$,
\begin{equation}\label{eq:cv-local-commutator}
 [\mathbb K,h]=\iota\mathbb C_h,
 \qquad
 |\mathbb C_h\Phi|_y\le C_C\vartheta^{-1}
                         \max_{d(y,z)\le\rho_c}|\Phi|_z.
\end{equation}
It vanishes where $h$ is constant on the old dependency
neighborhood together with these three-link neighborhoods.
The extension $\mathsf J$ commutes with $h$, since its two
coordinates have the same midpoint.

We check that the regular-piece construction applies to the
present geometry. Truncating the longer sequence at $k$ only
merges its last two layers. Its curvature, complex perturbation
and current bounds consequently imply those of the shorter
sequence. For the real gauge condition the only possible new
cube has index $k-1$. Its side is at most
$12\mathfrak M/L<\mathfrak M$. If it is within $12\mathfrak M$
of $\Omega_k^c$, the separation excludes $\Omega_{k+1}$ and
the cube already belongs to the longer class. Otherwise it is
contained in the concentric one-cube enlargement of a
level-$k$ $\mathfrak M$-cube through one of its points. That
enlargement lies in $\Omega_k$ and belongs to the longer
class. Restrict its gauge; the weights decrease by $L^{-1}$
and $L^{-2}$. Thus the pair belongs to the shorter natural
class with the same radii.

Use the pieces $\kappa$ and cutoffs of \cref{I-65-83}, with
\[
 \Theta'^\kappa=\{y\in\Theta':\Delta(y)\subset\kappa^{[1]}\},
 \qquad \Theta''^\kappa=\Theta''\cap\Theta'^\kappa.
\]
If $\Theta'^\kappa$ is empty the local system is the previous
eight-coordinate system. Otherwise $\kappa$ has size $k$ and
lies in the deep part of $\Omega_k$. Indeed a smaller piece
lies either outside $\Omega_k$ or within $3\mathfrak M$ of its
boundary; a point of $\Omega_{k+1}$ in its one-cube enlargement
would then be within $5\mathfrak M$, contrary to the
$R_1\mathfrak M\ge13\mathfrak M$ margin.
For these pieces use the regular local sequence
\[
 \Omega_k^\kappa=\kappa^{[1]}+[-e_k,e_k]^3,
 \qquad
 \Omega_m^\kappa=\Omega_{m+1}^\kappa+[-e_m,e_m]^3,
 \qquad e_m=2R_\flat\mathfrak M_\flat\xi_m.
\]
This prescription also applies when $k=0$: a fluctuation piece
at level zero uses this enlarged regular domain, rather than
the compressed bare boundary domain. The torus and size
conditions ensure that all these cubes are embedded.

For $\mathfrak M\ge32L^2R_\flat\mathfrak M_\flat$, the
following comparisons fit inside $\kappa^{[2]}$. If
$y\in\Theta'$ and $\Delta(y)$ lies within $10L$ of
$\kappa^{[5/6]}$, then $y\in\Theta'^\kappa$. If
$c\in\Gamma'$ and $\square_c$ lies within $8L$ of
$\kappa^{[5/6]}$, then $c\in\Gamma'^\kappa$.
To see this, the preceding boundary
argument first forces the piece to have size $k$, and then
$12L<\mathfrak M/6$ puts the block in $\kappa^{[1]}$.
Membership in $\mathfrak X(\Theta),Y,Y^{\rm ax}$ therefore
agrees with local membership in this neighborhood. The pivot
condition is the same on both sides, and each tree bond is
decided by its own block. The coefficient formulas show that
the local and global new kernels agree wherever the cutoffs
use them. For $\mathsf E^{(2)}[j_1]$ each contributing pivot
also lies in $\kappa^{[5/6]}$, where the truncated multiplier
equals $j_1$.

Every $c\in\Gamma'$ occurs in one of these local structures:
choose an end point $y\in\Theta'$ and a piece containing it.
The preceding comparison makes this a top fluctuation piece,
with $y\in\Theta'^\kappa$; moreover the double cube of $c$
lies within $3L/2$ of $y$ and hence in $\kappa^{[1]}$.
Thus the local mean and pivot charts define all the global
one-step coefficients used above.

The natural-class gauge on $\kappa^{[2]}$, together with the
same conversion of ordinary to covariant first differences as
in the regular inverse proof, makes the local quadruple
$\varepsilon$-small with
\[
 \varepsilon=8L^2(\alpha_0+\alpha_1+\gamma_0+\gamma').
\]
Here both truncated multipliers are bounded by $\gamma'$;
the global scale is no smaller than the local scale, so their
weighted bounds persist. Choose $a_3(L)>0$ so that radii and
multiplier bounds at most $a_3$ give
$\varepsilon\le\varepsilon_c$; it suffices to require
$a_3\le\varepsilon_c/(32L^2)$ in addition to the earlier
regular-inverse radius restrictions. The local Hessian remains the
ambient Hessian used in the existing proof. The bare pieces
without fluctuations retain their existing compressed inverse
and every touching-frame plaquette term.

Let $\mathbb E_\kappa$ be these local inverses. The old
partition satisfies $\operatorname{supp}\theta_\kappa
\subset\kappa^{[1/12]}$, while $\widetilde h_\kappa=1$ on
$\kappa^{[1/3]}$ and is supported in $\kappa^{[2/3]}$.
For the forcing product $\star$ already used in the regular
inverse proof, set
$\theta\star(g,f_9)=(\theta\star g,\theta f_9)$. Then
\[
 \iota(\theta\star(g,f_9))=\theta\iota(g,f_9),
\]
since the correction $-[D^{\mathsf T},\theta]X_5$ in slot $5$
cancels the extra divergence, and the ninth slot is unchanged.
The local and global indices, weights, and transports agree
on $\kappa^{[5/6]}$. The new commutator is supported there.
Thus the same parametrix gives
\[
 \operatorname{Par}=\sum_\kappa\widetilde h_\kappa
                           \mathbb E_\kappa\theta_\kappa,
 \qquad \mathbb K\operatorname{Par}=1-\mathcal R,
 \qquad \mathcal R\iota=\iota\mathcal R^\sharp.
\]
The input and commutator supports are separated by
$D_0=\mathfrak M/(8L)$ in the multiscale distance. Apply
\eqref{eq:cv-local-commutator} with
$\vartheta=\mathfrak M/C_G$ and spend half the local inverse
decay across this separation. For some $\delta_4,B_8,C_q>0$,
\[
 |\mathcal R_\kappa^\sharp g|_y
   \le q'\sup_{z\in S'_\kappa}e^{-\delta_4d(y,z)}|g|_z,
 \qquad
 q'=C_q\mathfrak M^{-1}e^{-\delta_4D_0/2}.
\]
The first factor $\widetilde h_\kappa\mathbb E_\kappa\iota
(\theta_\kappa\star\cdot)$ obeys the same estimate with
$4B_8$ instead of $q'$. The old proof applies verbatim to
the remaining coordinates. For the new coordinate it uses
only the range $12L$, midpoint cutoffs and scale $1$, as
checked above.

For $Nq'\le1/2$, expand $(1-\mathcal R)^{-1}$. If
$\varpi=(\kappa_0,\ldots,\kappa_m)$, denote the resulting
product by $\mathbb E_\varpi$. It is supported in
$\kappa_0^{[2/3]}$ at its output and uses input only in
$\kappa_m^{[1/12]}$. A nonzero product has consecutive
$\kappa_{j-1}^{[5/6]}$ and $\kappa_j^{[5/6]}$ meeting at a top,
and
\begin{equation}\label{eq:cv-leg-bound}
 |\mathbb E_\varpi\iota g|_y
 \le8B_8(q')^m\sup
       e^{-\delta_4\sum_{j=0}^m d(y_j,y_{j+1})}|g|_{y_{m+1}}.
\end{equation}
Here $y_0=y$, the intermediate $y_j$ lie in these overlaps,
and $y_{m+1}\in\kappa_m^{[5/6]}$. Summation over pieces gives
$\mathbb E\iota\in\mathfrak O(\delta_4,16NB_8)$.
This proves invertibility of both enlarged systems for all
quadruples with these bounds. By
\eqref{eq:bg-total-multipliers}, reducing $a_{\rm cv}$ makes
the physical quadruple $(\mathbf U,\mathbf J,j_*,j_{1*})$
one of them. Equations \eqref{eq:cv-schur} prove existence of
the two covariances. Their symmetry follows from symmetry of
the gauge form, $\mathsf E^{(2)}[j_{1*}]$, and
$(Q\mathsf G_1Q^{\mathsf T})^{-1}$.

\subsection*{The second multiplier}
We use the local expansion from
\cref{prop:local-multiplier-trees} for the shorter sequence.
In its notation
\[
 j_*=\sum_\varpi j_\varpi=j^0+\delta j,
 \qquad
 \mathbb E^{[5]}(j_*)=\sum_\Xi\mathbb E_\Xi.
\]
The first sum has its zero-length part $j^0$; its positive-length
terms have the factor $q_1^{|\varpi|}$, where
$q_1\le\mathfrak M^{-9}$. In particular
$|\delta j|_{-3}\le C(L)\gamma_0\mathfrak M^{-9}$.
The second sum is over the first decorated sequences, with
trunk cells $X^{(2)}(\kappa)$ and branch cells $X(\kappa)$.
With \(\delta_b=\delta_{\rm reg}\), \(C_b=B_{\rm tr}\) and
\(C_{{\rm w},b}=C_{\rm tr}\), the estimate is
\begin{equation}\label{eq:cv-first-decoration-input}
 \sum_{\omega(\Xi)=\omega}|\mathbb E_\Xi\iota g|_y
 \le C_b(C_{{\rm w},b}\mathfrak M^{-3})^{|\omega|}
       \sup_z e^{-\delta_b d^{\rm tr}_{\mathbf B}
                            (\omega;y,z)/2}|g|_z.
\end{equation}
The factors and this estimate are defined on the open
piecewise local class on the carrier. Each $j_\varpi$ is
supported on rows in its last piece's $[1/12]$ cube.
The operator $\mathbb E_\Xi$ uses input only in the last
piece's $[1/12]$ cube and has output in the first piece's
$[2/3]$ cube.

Define $j_{1,\Xi}$ by its weighted transpose, not an unweighted
transpose of the kernel:
\begin{equation}\label{eq:cv-second-multiplier}
 \xi_c^3\langle e,j_{1,\Xi}(c)\rangle
       =([\mathbb E_\Xi]_{A,8}e_c,\mathbf J)_\xi.
\end{equation}
The right side is summed over the bare output bonds of the
first piece. These bonds have both end points in its carrier,
so this defines the same function on local data. It is
holomorphic there and vanishes unless the midpoint of $c$
lies in the last piece's $[1/12]$ cube.

The current bound and the bare-to-top count used for the total
multiplier give, for $g$ supported in slot $8$ at the row $c$,
\[
 \xi_c^3|\langle e,j_{1,\Xi}(c)\rangle|
       \le3\gamma_0\sum_y\xi_y^{-1}
                                      |\mathbb E_\Xi\iota g|_y,
 \qquad |g|_{y(c)}=\xi_c|e|.
\]
If $\underline d(\omega;z)$ denotes the least length of a
tree meeting all the cells and containing $z$, then
$d^{\rm tr}(\omega;y,z)$ dominates both $d(y,z)$ and
$\underline d(\omega;z)$. Split the exponential in
\eqref{eq:cv-first-decoration-input} accordingly, and use
the existing scale-ratio summation, increasing the fixed floor so that
\(R_1\mathfrak M\ge\mathsf r_S(\delta_b/8)\) and
\(\log L/(R_1\mathfrak M)\le\delta_b/8\):
\[
 \sum_y\frac{\xi_z}{\xi_y}e^{-\delta_b d(y,z)/4}
                  \le Lc_S(\delta_b/8).
\]
Taking the three Pauli coordinates in
\eqref{eq:cv-second-multiplier} yields
\begin{equation}\label{eq:cv-second-multiplier-tree}
 \sum_{\omega(\Xi)=\omega}\xi_c^3|j_{1,\Xi}(c)|
 \le C'_j\gamma_0(C_{{\rm w},b}\mathfrak M^{-3})^{|\omega|}
       e^{-\delta_b\underline d(\omega;y(c))/4}.
\end{equation}
This proves $j_{1*}=\sum_\Xi j_{1,\Xi}$ by absolute convergence
and the already identified block
$\mathsf H_1=[\mathbb E^{[5]}(j_*)]_{A,8}$.
Let $j_1^0$ be the sum of the zero-length terms and
$\delta j_1=j_{1*}-j_1^0$.

We also need the stronger small tail bound
\begin{equation}\label{eq:cv-multiplier-tails}
 |j^0|_{-3},\ |j_1^0|_{-3}\le C_j\gamma_0,
 \qquad
 |\delta j|_{-3}+|\delta j_1|_{-3}
                   \le C''_j\gamma_0\mathfrak M^{-9}.
\end{equation}
For the second tail, write
$E_* =\mathbb E^{[5]}(j_*)$, $E_0=\mathbb E^{[5]}(j^0)$,
and let $E_{00}$ be its sum of one-piece factors. The weighted
transpose of
\[
 E_*-E_{00}=(E_*-E_0)+(E_0-E_{00})
\]
gives $\delta j_1$. The resolvent identity bounds the first
difference by two total inverses and the old vertex applied
to $\delta j$; the second is the sum of positive-length
cutoff products. For a product of length $m$, replace the
overlap supremum by a sum. The scale ratio is bounded along
the entire chain by
$L\exp(\delta_b\sum_i d(y_i,y_{i+1})/4)$, with one factor
$L$, independent of $m$. Summing successively over its
$m+1$ intermediate positions gives a bound
\[
 \xi_c^3|j_{1,(\varpi)}(c)|
       \le C(L)\gamma_0q_{\mathfrak M}^{m}
                    c_S(\delta_b/4)^{m+1}
       \le C(L)\gamma_0q_1^m,
 \qquad q_1=c_S(\delta_b/4)q_{\mathfrak M}.
\]
There are at most $N^{m+1}$ contributing sequences, counted
backwards from the last piece containing the midpoint of
$c$. With $Nq_1\le1/2$ their positive-length sum is
$O(\gamma_0q_1)$. The first resolvent difference has the
same bound after the weighted-transpose estimate. This proves
the second inequality in
\eqref{eq:cv-multiplier-tails}.
For $j_1^0(c)$ at a local pair there are at most $N$ contributing
one-piece factors, all with midpoint $c\in\kappa^{[1/12]}$.
The one-piece inverse estimate and the same summation give
the first inequality on the local class as well. Choose
$a_{\rm cv}$ so that $C_j a_{\rm cv}\le a_3/2$.

\subsection*{Two vertices and the tree expansion}
Use the local inverse at
$q^0=(\mathbf U,\mathbf J,j^0,j_1^0)$ as reference. The old
vertex $\mathbb W[j']$ is extended by a zero ninth component;
the second vertex has only a ninth component,
\[
 \mathbb W_1[j'](\Psi,\check B)
       =(0,\mathsf C^{\mathsf T}\mathsf E^{(2)}[j']
                                      \mathsf C\check B).
\]
The equations give the exact identity
\begin{equation}\label{eq:cv-two-vertices}
 \mathbb K(q_*)=\mathbb K(q^0)
                 -\iota\mathbb W[\delta j]
                 -\iota\mathbb W_1[\delta j_1].
\end{equation}
The old vertex has range at most $3L+6$. For the new one the
three-kernel product above gives
\begin{equation}\label{eq:cv-second-vertex}
 |\mathbb W_1[j']\Phi|_y
 \le C_{W,1}
       \max_{d(y,y(c))\le8L}\xi_c^3|j'(c)|
       \max_{d(y,z)\le12L}|\Phi|_z.
\end{equation}
Only pivot rows occur in the first maximum, and they have
$\xi_c=1$. There are at most $(36L^3)^3$ index choices in
this product, so $C_{W,1}$ depends only on $L$.
Both vertices therefore have this form of bound, after
enlarging the constant for the first one.

Let $\mathcal N(\kappa)$ be the pieces whose $[1]$ cubes have
a common interior point with $\kappa^{[1]}$ in $\Omega_0$.
The local expansion already uses
$X^{(2)}(\kappa)=\bigcup_{\kappa'\in\mathcal N(\kappa)}
X(\kappa')$, with the union of their cube families.
For the present expansion use also
\[
 X^{(3)}(\kappa)
       =\bigcup_{\kappa'\in\mathcal N(\kappa)}X^{(2)}(\kappa').
\]
There are at most $N$ neighbors and at most $27N$ cubes in
each family of $X^{(2)}$, so $X^{(3)}$ has at most $27N^2$
cubes. The families are connected: each neighbor's
$X^{(2)}$ contains the family of $X(\kappa)$.
They are therefore cells for $r_{\rm cv}=27N^2$.
At most $27N^2$ pieces have a prescribed $X^{(3)}$ and family.
Indeed that family contains the cube $\kappa^{[1]}$ for a
regular piece, or $\kappa$ for a bare piece; these cubes
determine the piece and are distinct across the two kinds.

Every base cell and hence every one of these larger cells is
also a cell of the longer sequence, with the same carrier.
Only a cube of index $k-1$ requires checking. In a piece's
family it meets either $\mathcal L_{k-1}$ or the buffer of
$\Omega_k$, so its distance to $\Omega_k^c$ is less than
$3\mathfrak M$. The earlier cube comparison puts it in the
longer class. Its side determines its index uniquely through
$\mathfrak M\xi_n\le s<L\mathfrak M\xi_n$, so both the
index and its five-cube enlargement are unchanged. The
usual count of connected cell families and depth-first
parent maps gives $n_{\rm cv}^m$ tree walks of length $m$
with a fixed first cell.

Consider products
\begin{equation}\label{eq:cv-decoration}
 \mathbb E_\Upsilon=
 \mathbb E_{\varpi^{(0)}}\iota\mathbb W_{\beta_1}
 \mathbb E_{\varpi^{(1)}}\cdots
 \iota\mathbb W_{\beta_n}\mathbb E_{\varpi^{(n)}}.
\end{equation}
Here every $\varpi^{(l)}$ is a cutoff sequence for the system
at $q^0$. A vertex label $\beta_l$ is either a positive-length
sequence $\varpi'$ with
$\mathbb W_{\beta_l}=\mathbb W[j_{\varpi'}]$, or a first
decoration $\Xi'$ of positive tree length with
$\mathbb W_{\beta_l}=\mathbb W_1[j_{1,\Xi'}]$.
Make a trunk from the $X^{(3)}$ cells of the sequences
$\varpi^{(0)},\ldots,\varpi^{(n)}$ in their order.
Attach a branch for each vertex to the last cell preceding
that vertex. A branch of the first kind is the reversed
sequence of its $X$ cells. A branch of the second kind has
the reversed $X^{(2)}$ trunk of $\omega(\Xi')$, keeping its
own $X$ branches attached to the same cells. List the outer
branches in reverse vertex order. After each reversed inner
trunk, list its branches in increasing label order
$k'=1,\ldots,n'$, retaining their internal order and
attachment cells. This gives
the depth-first parent map $p$, with $p(i)<i$ and each next
parent an ancestor of the preceding index. Its initial
uninterrupted branch is exactly the outer trunk.

A nonzero product gives a tree walk. Consecutive cells within
a cutoff sequence overlap by \eqref{eq:cv-leg-bound}.
At an old vertex this is the overlap proved for the first
multiplier expansion. At a new vertex, nonzero entries give
bonds $b,b_1,b_2,b_3$ and a pivot $b_0(c)$, with $b$ linked
to $b_1$, $b_1,b_2\in\mathsf P_c$, and $b_2$ linked to $b_3$.
Their midpoints obey
\[
 |\widehat{b_0(c)}-\hat b|_\infty\le4L,
 \qquad
 |\widehat{b_0(c)}-\hat b_3|_\infty\le4L.
\]
They are all within $L$ of $\Omega_{k+1}$. A piece whose
$[2]$ cube contains such a point has size at least $k-1$:
otherwise that point would lie in a layer outside the
interior of $\Omega_k$. Thus its side is at least
$\mathfrak M/L\ge32L^3$. The three cutoff supports and
$4L\le(32L^3)/8$ put the pivot midpoint in the interior of
all three $[1]$ cubes. The branch's last piece is therefore
a neighbor of the preceding trunk piece, its $X^{(2)}$ cell
lies in that trunk's $X^{(3)}$ cell, and the two adjacent
trunk cells overlap. This proves all required overlaps.

If $m_l=|\varpi^{(l)}|$ and $m'_l$ is the length of the
branch label, then $m'_l\ge1$ and counting nodes gives
\begin{equation}\label{eq:cv-decoration-length}
 |\omega(\Upsilon)|
       =\sum_{l=0}^n m_l+\sum_{l=1}^n(m'_l+2).
\end{equation}
For a fixed tree walk of length $m$ there are at most
$10^m(27N^2)^{m+1}$ such products. To check the count, mark
each node as an outer-trunk start, a first-kind branch start,
a second-kind branch start, an inner branch start, or none.
There are at most $5^m$ choices. The already proved first
decoration count contributes at most $2^{m'}$ to each
second-kind branch, hence at most $2^m$ in total. Each
prescribed cell has at most $27N^2$ piece preimages. These
marks, the parent map and the selected pieces recover the
cutoff sequences and vertex labels.

We estimate the product before summing it. Write
$z_l$ for the last intermediate point of leg $l-1$ and
$y_l$ for the first point of leg $l$. The vertex bounds use
a row $c_l$ with
\[
 d(z_l,y(c_l))\le8L,\qquad d(z_l,y_l)\le12L.
\]
Apply \eqref{eq:cv-leg-bound} on each leg. A first-kind
branch contributes
$C_j\gamma_0 q_1^{m'_l}
 e^{-\delta_b\mathsf d(\varpi';c_l)/2}$, where
$\mathsf d(\varpi';c_l)$ is the least length of a chain
through its successive piece overlaps to $y(c_l)$.
A second-kind branch contributes, by
\eqref{eq:cv-second-multiplier-tree},
\[
 C'_j\gamma_0(C_{{\rm w},b}\mathfrak M^{-3})^{m'_l}
     e^{-\delta_b\underline d(\omega(\Xi');y(c_l))/4}.
\]
Join the paths used in the legs, the branch chains or trees,
and the two vertex links. Their union is connected, contains
the two endpoint positions, and meets every cell. The two
links cost at most $20L$ per vertex. Therefore, with
$\delta_5=\min(\delta_4,\delta_b/4,1)$, the product of all
exponentials is bounded by
\[
 e^{20Ln}e^{-\delta_5d^{\rm tr}_{\mathbf B}(\omega;y,y')}.
\]
This follows by choosing paths within any prescribed error
of the defining infima and letting that error tend to zero.
Every cell of a positive-length first-kind branch is met
by an overlap point; the last point is connected to $y(c_l)$.
This is the reason these branches exclude length zero.

Now $q'\le C_q\mathfrak M^{-1}$,
$q_1\le\mathfrak M^{-9}$, and
$3m'\ge m'+2$ for $m'\ge1$. Absorb the fixed vertex factors
and $e^{20L}$ into $C_\sharp(L)$. Equation
\eqref{eq:cv-decoration-length} gives
\[
 |\mathbb E_\Upsilon\iota g|_y
 \le8B_8(C_\sharp\mathfrak M^{-1})^{|\omega|}
       \sup_{y'}e^{-\delta_5d^{\rm tr}_{\mathbf B}
                                   (\omega;y,y')}|g|_{y'}.
\]
The count of products then gives
\begin{equation}\label{eq:cv-outer-tree-bound}
 \sum_{\omega(\Upsilon)=\omega}|\mathbb E_\Upsilon\iota g|_y
 \le216N^2B_8(C_{\rm cv}\mathfrak M^{-1})^{|\omega|}
       \sup_{y'}e^{-\delta_5d^{\rm tr}_{\mathbf B}
                                   (\omega;y,y')}|g|_{y'},
\end{equation}
where one may take $C_{\rm cv}=270N^2C_\sharp$.

The sum of these products is the physical inverse. Indeed
\eqref{eq:cv-multiplier-tails}, the two vertex estimates and
the total reference inverse imply
\[
 \bigl(\mathbb W[\delta j]+\mathbb W_1[\delta j_1]\bigr)
                      \mathbb E(q^0)\iota
       \in\mathfrak O(\delta_4,C(L)\gamma_0\mathfrak M^{-9}).
\]
Increase $M_{\rm cv}$ until this bound is at most $1/2$.
The resolvent expansion of \eqref{eq:cv-two-vertices} is then
\[
 \mathbb E(q_*)=\mathbb E(q^0)
       \sum_{n\ge0}\left[
        \iota\bigl(\mathbb W[\delta j]+\mathbb W_1[\delta j_1]\bigr)
                       \mathbb E(q^0)\right]^n.
\]
Insert the cutoff expansion and the two multiplier expansions.
Estimate \eqref{eq:cv-outer-tree-bound} and the tree count show
absolute convergence, so regrouping gives
$\mathbb E(q_*)=\sum_\Upsilon\mathbb E_\Upsilon$.
This equality is an identity of matrices: the forcing
injection is onto, and the spaces are finite dimensional.

\subsection*{Local terms and the covariance bounds}
We give the local meaning of the products, since choosing a
global continuation of local data would not define their
coefficients intrinsically. Use the piecewise local class
of \cref{sec:local-background-expansion} for the shorter
sequence. Its openness was proved there.

The longer local natural class lies in this piecewise class.
The pointwise inequalities only weaken on truncation. For
the cube condition, a shorter cube not in the longer class
must have index $k-1$ and distance greater than
$12\mathfrak M$ from $\Omega_k^c$. If it lies in
$X(\kappa)^\natural$, the carrier geometry forces
$\kappa$ to have size $k$: for a smaller piece every point
of its carrier is within $6\mathfrak M$ of the piece, which
is outside $\Omega_k$ or within $3\mathfrak M$ of its boundary.
For a top piece the carrier is
$\kappa^{[6]}\cap\Omega_0$. In its embedded chart, choose a
cube of side $2\mathfrak M$, a union of eight top cubes,
containing the shorter cube and contained in $\kappa^{[6]}$.
This is possible because the shorter side is less than
$\mathfrak M$. The enlarged cube stays more than
$10\mathfrak M$ from $\Omega_k^c$, so lies in the longer
natural class and in the same carrier. Its real gauge
restricts to the original cube with smaller weights. This
proves the required local inclusion without assuming a
global extension of the pair.

For a given outer decoration take $W=[\omega(\Upsilon)]$.
Define $j^0,j_1^0$ by their local formulas on the rows in the
$[5/6]$ cubes of its trunk pieces, and set them to zero on
the other rows. Their bounds in
\eqref{eq:cv-multiplier-tails} make all the local inverse
factors well defined. They are holomorphic: the coefficient
charts are holomorphic, pivot inverses have nonzero determinant,
and each finite local system remains invertible on this
open class. Only the second vertex needs an additional
restriction. If $\kappa^-$ is the last piece of the preceding
leg, use its local $\mathsf C_\kappa$ and
$\mathsf E_\kappa^{(2)}$, and truncate $j'$ to rows in
$\kappa^{-[5/6]}$. Restrict the input to the common
fluctuation bonds with midpoint in $\kappa^{-[2/3]}$, and
retain output only at common bonds in $\kappa^{-[1/12]}$.
Denote the resulting vertex by $\mathbb W_{1,\kappa^-}[j']$;
it is zero for a piece with no fluctuations.

The three links of its kernel move a midpoint by at most
$6L$. Thus, behind $\theta_{\kappa^-}$, all input, output
and coefficient indices are in precisely the common sets
checked above, and
\[
 \theta_{\kappa^-}\iota\mathbb W_{1,\kappa^-}[j']
       =\theta_{\kappa^-}\iota\mathbb W_1[j']
\]
at the restriction of a global pair. The local vertex obeys
\eqref{eq:cv-second-vertex}. Its coefficients depend only
on the given data in $\kappa^{-[2]}$. All factors of
\eqref{eq:cv-decoration}, with this replacement, are
therefore holomorphic on the piecewise local class on $W$,
independent of a continuation, and agree with the global
factors on restriction. The overlap and chain arguments
just proved use only these local formulas, so
\eqref{eq:cv-outer-tree-bound} holds on this entire class.

Define
\[
 \mathcal C_\omega
       =\sum_{\omega(\Upsilon)=\omega}
                         [\mathbb E_\Upsilon]_{\check B,9}.
\]
The sum is finite for fixed $\omega$. It gives the asserted
local holomorphic term, and the global sum is the covariance
by \eqref{eq:cv-schur}. For a ninth-coordinate input equal
to $e$ at $b'$ and zero elsewhere, $\iota g=g$,
$|g|_{b'_-}=|e|$, and all other forcing seminorms vanish.
The output satisfies
$|\check B(b)|\le|(\Psi,\check B)|_{b_-}$.
Thus \eqref{eq:cv-outer-tree-bound} becomes
\begin{equation}\label{eq:cv-tree-stronger}
 |\mathcal C_\omega(b,b')|
 \le216N^2B_8(C_{\rm cv}\mathfrak M^{-1})^m
       e^{-\delta_5d^{\rm tr}_{\mathbf B}(\omega;b_-,b'_-)}.
\end{equation}
The first and last cutoffs put $\hat b$ in the first piece's
$[2/3]$ cube and $\hat b'$ in the last piece's $[1/12]$ cube.
These pieces have side at least $32L^3$, as shown above.
Moving a midpoint by $1/2$ to its initial endpoint therefore
keeps that endpoint in the piece's $[1]$ cube and hence in
its $X$ cell. This proves the two endpoint support statements.

The weighted distance for the longer sequence is no larger
than that for the shorter sequence, including tree distances.
For the latter comparison cut a tree into segments at all
intersections and marked points. On each segment the
inverse-scale integrand decreases when the last layer is
added. Replace that segment by a path within any error of
the longer weighted distance between its endpoints; their
union remains connected and meets all marked cells.
Subadditivity and then passage to the infimum prove the
comparison. Each bond here lies in the interior of
$\Omega_k$, so $d(b_-,y(b))\le1$. Adding the two bond segments
to a tree costs at most $2$. Taking
$\delta_{\rm cv}=\delta_5/4$ and enlarging $B_0$ by $e^2$
therefore turns \eqref{eq:cv-tree-stronger} into
\eqref{eq:cv-tree-decay}.

For fixed $b$ at most $N$ first pieces have
$\hat b\in\kappa^{[2/3]}$. Moreover
$d^{\rm tr}_{\mathbf B}(\omega;b_-,b'_-)
\ge d(b_-,b'_-)\ge|b_--b'_-|_1$.
Choose $M_{\rm cv}\ge(2C_{\rm cv}n_{\rm cv})^2$ and
$B_0\ge\max(432N^3B_8,216e^2N^2B_8)$.
The count of tree walks and
\eqref{eq:cv-tree-stronger} give the length-$m$ estimate in
the statement. Summing it proves the kernel estimate for
the full covariance and for every partial sum.

Finally a site is the initial point of at most three positive
bonds, and the unit-spacing torus sum is bounded by the
$\mathbb Z^3$ sum. Thus every row sum and column sum of the
absolute kernel are at most
$3B_0\sigma(\delta_{\rm cv})^3$. This proves the supremum
bound and, by the Schur test, the $\ell^2$ bound. For a
$2\times2$ matrix $X\in\mathfrak g^c$,
\[
 \tfrac12|X|^2\le\langle X^*,X\rangle\le|X|^2;
\]
the Hermitian operator norm is consequently at most
$\sqrt2$ times that $\ell^2$ bound, and hence at most $B$.
All choices used only $L$ and the fixed constants of the
earlier inverse and multiplier estimates.
\end{proof}

\section{Local nonlinear response with decoupling parameters}\label{sec:decoupled-response}

\subsection{Local propagators with decoupling parameters}

Use the cells, their carriers, and the piecewise local class of
\cref{prop:local-multiplier-trees}. Put
\[
 \delta=\delta_{\mathrm{reg}},\qquad
 \mathring\delta_0=\delta/4,\qquad r_{\mathrm w}=27N_{\mathrm G}.
\]
The scale is the level-$t$ scale, so that $\xi_t=1$ and
$\xi=\xi_0$. All transposes below use the existing physical
pairings: weight $\xi^3$ on bare sites and bonds, and weight
$\xi_z^3$ on scalar tops and physical rows. Write
\[
 H_{\mathrm{pr}}=GQ_{\mathrm{pr}}^{\mathsf T}
        (Q_{\mathrm{pr}}GQ_{\mathrm{pr}}^{\mathsf T})^{-1},
 \quad S_{\mathrm{pr}}=(Q_{\mathrm{pr}}GQ_{\mathrm{pr}}^{\mathsf T})^{-1},
 \quad P_{\mathrm f}=D(1-\mathsf R_{\mathsf m})D^{\mathsf T}.
\]
Here $\mathsf R_{\mathsf m}$ is the projection defined by the
scalar mean average. The operators $H_{1,\mathrm{pr}}$ and
$\mathfrak G$ are those at the physical multiplier
$\mathsf j_*=H_{\mathrm{pr}}^{\mathsf T}\mathbf J$.

For a cell $X$, set
\[
 X^{\boxplus}=\bigcup_{q\in X}
 \{x:\operatorname{dist}_\infty(x,q)\le5\mathfrak M\xi_{n(q)}\}.
\]
The hull of a walk or tree $\omega$ is the union of these
sets over its cells; its carrier is $[\omega]=\Omega_0$
intersected with that hull. Let $D_\omega(y,z)$ be the
infimum of the weighted lengths of connected graphs containing
$y,z$ and meeting every cell of $\omega$.
Two cubes of a standard partition are wall adjacent when
they share a face; wall components use this adjacency.
Write $r_0=2(3L+2)$ for the response radius, using the
value $\rho_0=3L+2$ fixed in the local response estimates.

\begin{proposition}\label{prop:decoupled-local-propagators}
There are $a_{\mathrm D}>0$, $M^{\mathrm{pr}}$,
$r^{\mathrm{pr}}$, $C_{\mathrm w}$ and $B$, depending only on
$L$, for which the following holds. Assume the regular-pair
frame hypotheses, with
\[
 \alpha_0,\alpha_1,\gamma_0\le a_{\mathrm D},\quad
 \mathfrak M\ge M^{\mathrm{pr}},\quad
 R_1\mathfrak M\ge r^{\mathrm{pr}},\quad
 \ell\ge2L\mathfrak M,
\]
and the additional summation and scale-ratio requirements
\[
 \frac{\log L}{R_1\mathfrak M}\le\frac{\mathring\delta_0}{48},
 \quad R_1\mathfrak M\ge\mathsf r_S(\mathring\delta_0/16),
 \quad\mathfrak M\ge(8C_{\mathrm w}n_{1,\natural})^2.
\]
Here $n_{1,\natural}=n_1(L,r_{\mathrm w})$ is the
cell-neighbour bound from the preceding construction; $M^{\mathrm{pr}}$ also exceeds all
floors in that construction and in the doubled inverse.

Each of $H_{\mathrm{pr}},S_{\mathrm{pr}},H_{1,\mathrm{pr}},
\mathfrak G,P_{\mathrm f}$ has an absolutely convergent
expansion $O=\sum_\omega O_\omega$. Its terms are holomorphic
on the piecewise local class on $[\omega]$, depend only on
the pair and the sequence there, and agree under restriction
of larger local domains. They are gauge covariant and commute
with the real structure. A nonzero kernel has both arguments
in cells of $\omega$. The local construction includes layer
zero and all touching ambient plaquettes.

The terms obey the following bounds with
$v_m=(C_{\mathrm w}/\mathfrak M)^m$, $m=|\omega|$:
\begin{align}
 |(K_\omega e_c)(b)|
 &\le Bv_m e^{-\mathring\delta_0D_\omega(y,y_c)}|e|,
 \notag\\
 |(\nabla_{\mathbf U}K_\omega e_c)(b,\nu)|
 &\le Bv_m\xi_y^{-1}e^{-\mathring\delta_0D_\omega(y,y_c)}|e|,
 \label{eq:dec-local-row-columns}\\
 |S_{\mathrm{pr},\omega}(c,c')|
 &\le Bv_m\xi_y^{-2}e^{-\mathring\delta_0D_\omega(y,y')},
 \label{eq:dec-local-row-inverse}\\
 \max\{\xi_y^{-2}|(\mathfrak G_\omega f)(b)|,
             \xi_y^{-1}|(\nabla_{\mathbf U}\mathfrak G_\omega f)(b,\nu)|\}
 &\le Bv_m e^{-\mathring\delta_0D_\omega(y,y')}
                       \sup|f|,
 \label{eq:dec-local-projected}\\
 |P_{\mathrm f,\omega}(b,b')|
 &\le Bv_m(\xi/\xi_{y'})^3\xi_y^{-2}
                  e^{-2\mathring\delta_0D_\omega(y,y')}.
 \label{eq:dec-local-projection}
\end{align}
In the first line $K=H_{\mathrm{pr}}$ or $H_{1,\mathrm{pr}}$;
$b$ lies over $y$, $c'$ is at $y'$, and $f$ vanishes outside
the bare bonds over $y'$. Kernels use counting measure.
For every fixed initial cell,
\[
 \sum_\omega v_{|\omega|}\le2,
 \qquad \sum_\omega(|\omega|+1)v_{|\omega|}\le4.
\]

Fix a real number $\varkappa>0$ and a standard partition
$\mathsf P=\pi[\mathsf w]$ with at least three cubes in each direction and
\[
 \mathsf w\ge88r_{\mathrm w}\mathfrak M,
 \qquad\mathring\delta_0\mathsf w\ge2176\varkappa.
\]
Let $\mathsf P(\omega)$ be the cubes meeting the hull, let
$\mathsf Y\subset\mathsf P$, and put
$W=\Omega_0\cap\bigcup\mathsf Y$, with the cube-wise
intersection. For a local natural pair on $W$ and parameters
$s=0$ off $\mathsf Y$, $|s_\Delta|\le e^\varkappa$, set
\[
 O(s)=\sum_{\mathsf P(\omega)\subset\mathsf Y}
       \Bigl(\prod_{\Delta\in\mathsf P(\omega)}s_\Delta\Bigr)
       O_\omega[p|_{[\omega]}].
\]
These series converge absolutely, uniformly on the local
class and on bounded parameter sets. They are holomorphic in
the pair and in $s$, and have degree at most one in each
$s_\Delta$. The assignment of terms is fixed independently
of $s$ and of the local domain. At $s\equiv1$ on the whole
partition they equal the five global operators above.

For $\mathsf m(c)=\xi_c$, define
\[
 (H_s,S_s,K_s,\mathfrak G_s)
  =(H_{\mathrm{pr}}(s)\mathsf m^{-1},
    \mathsf m^3S_{\mathrm{pr}}(s)\mathsf m^{-1},
    H_{1,\mathrm{pr}}(s)\mathsf m^{-1},\mathfrak G(s)),
 \qquad P_s=P_{\mathrm f}(s).
\]
These are the operators of the local response structure on
$Y^1(\mathsf Y)\subset Y^2(\mathsf Y)\subset Y^3(\mathsf Y)$.
They satisfy its strong and column estimates with
\[
 b=b_0e^{8\varkappa},\quad
 \widehat b=b_{\mathrm D}e^{8\varkappa},\quad
 \mu=3\mathring\delta_0/16,\quad\theta=\mathring\delta_0/8,
 \qquad b_{\mathrm D}=b_0c_S(\mathring\delta_0/16),
\]
including $|P_sA|_w^\sharp\le\widehat b|A|_w$.
The same estimates hold for every subsum. The constants
$b_0,b_{\mathrm D}$ depend only on $L$.
Here $Y^j(\mathsf c)$ consists of the tops within sup distance
$j\mathfrak M\xi_{n(q)}$ from a cube $q$ of a cell $X$ with
$\mathsf P(X)\subset\mathsf c$.
Each operator is carried in both variables by $Y^1(\mathsf Y)$.
It is the sum of the operators belonging to the wall
components of $\{s\ne0\}$, each depending only on the pair
and parameters in that component. The sets $Y^j$ belonging
to distinct components are separated by more than the
response radius $r_0$, for $j\le4$.
\end{proposition}

\begin{proof}
\emph{The common local trees.}
The proof of \cref{prop:local-multiplier-trees} already
constructs the local fifth-system factors. The same factors
at $\mathsf j=0$ give the ordinary walks. Sum their $A,8$
and $\omega,8$ blocks to define $H_{\mathrm{pr},\omega}$
and $S_{\mathrm{pr},\omega}$; sum the $A,8$ blocks of the
decorated factors to define $H_{1,\mathrm{pr},\omega}$.
The inverse block identities identify their global sums.

For the projected operator use the doubled system. On each
piece its coupling is
\[
 \mathbb J\Psi'=e_1D(D^{\mathsf T}A'-w').
\]
It is independent of the multiplier and uses only bonds of
$\Omega_0^\kappa\subset\kappa^{[2]}\cap\Omega_0$.
Thus it is the same operator for every local ambient region
containing the carrier. The local regular gauge, the
touching-plaquette hypothesis and the layer-zero inverse
apply exactly as in the preceding local-tree proof.
The half-H\"older estimates used in the proof of
\cref{prop:projected-propagator} give the local doubled
strong inverse; the commutator
\eqref{eq:doubled-commutator} gives ordinary forcing and
vanishes where the cutoff is constant. Consequently the
lifted product proof of \eqref{eq:doubled-total-bound}
applies on these local classes, with constants $B_4^{\mathrm d}$
and
\[
 q_{\mathrm d}=C_q^{\mathrm d}\mathfrak M^{-1}
                   e^{-\delta\mathfrak M/(8L)}.
\]
Take $C_q^{\mathrm d}$ at least the larger of the global and
local commutator constants, and increase the floor until
$q_{\mathrm d}\le\mathfrak M^{-9}$.

In each decorated label $\Xi$, use the same local zero-length
multiplier $\mathsf j^0$ and the same branch multipliers
$\mathsf j_{\varpi'_k}$ as for the fifth system. The doubled
vertex is diagonal:
\[
 \mathbb W^{\mathrm d}[j](\Psi,\Psi')
      =(\mathbb W[j]\Psi,\mathbb W[j]\Psi').
\]
Indeed the coupling $\mathbb J$ cancels in the difference of
the systems at two multipliers. The support argument is
unchanged: a row used at a junction has its full averaging
set in $\kappa'^{[5/12]}$, and all three pieces at that
junction lie in the assigned carrier. The same $j^0$ and
vertex therefore have the same values under restriction or
under a continuation outside the carrier.

To check summability, write $m_k$ for the trunk lengths and
$m'_k\ge1$ for the branch lengths. The strong product estimate
is bounded by
\[
 8B_4^{\mathrm d}\prod_kq_{\mathrm d}^{m_k}
       \prod_{k\ge1}(C_b^{\mathrm d}\gamma_0q_1^{m'_k})
 \sup_z\exp\{-\delta\textstyle\sum d_{\rm trunk}
              -\frac\delta2\sum d_{\rm branch}\}\,|g|_z,
 \quad C_b^{\mathrm d}=8B_4^{\mathrm d}C_{\mathrm{vert}} e^{\rho_0}C_j.
\]
Here $q_1,C_j$ are those of the first multiplier expansion.
The vertex sees only its weak field seminorm, which is
bounded by the strong solution seminorm. The junctions have
distance at most $\rho_0$; joining them changes the exponent
by at most $\delta\rho_0$ per branch. If
$m=\sum m_k+\sum m'_k+2n$, then
$\mathfrak M^{-9m'_k}\le\mathfrak M^{-3(m'_k+2)}$.
There are at most $2^m(27N_{\mathrm G})^{m+1}$ decorated
labels for a fixed tree. Thus, after increasing a constant
$C_{\rm tr}^{\mathrm d}$,
\[
 \sum_{\Xi:\omega(\Xi)=\omega}|\mathbb E^{\mathrm d}_\Xi\iota g|_y
 \le B_{\rm tr}^{\mathrm d}
       (C_{\rm tr}^{\mathrm d}\mathfrak M^{-3})^m
       \sup_z e^{-\delta D_\omega(y,z)/2}|g|_z.
\]
Decrease $a_{\mathrm D}$ so that $|j^0|_{-3}\le a_{\mathrm d}/2$
and increase the floor so that
$4C_{\mathrm{vert}}B_5^{\mathrm d}C_jN_{\mathrm G}e^{\rho_0}q_1e\le1$.
Here $B_5^{\mathrm d}=16N_{\mathrm G}B_4^{\mathrm d}$,
and $a_{\mathrm D}$ is at most both the first-tree and
doubled-inverse radii. These choices also include the local
small-triple restriction of those inverse estimates.
The resolvent series in $\mathsf j_*-\mathsf j^0$ then
converges. Its absolute convergence permits the same
rearrangement as in the fifth-system proof, identifying the
sum with the doubled inverse at $\mathsf j_*$.

On this common set of decorated labels put
\begin{equation}\label{eq:dec-direct-projected-tree}
 \mathfrak G_\omega
  =\sum_{\Xi:\omega(\Xi)=\omega}
    \bigl([\mathbb E^{[5]}_\Xi]_{A,1}
                -[\mathbb E^{\mathrm d}_\Xi]_{A,1'}\bigr).
\end{equation}
The first slot $1'$ is in the second copy and the output $A$
is in the first. The elimination in
\cref{prop:projected-propagator} identifies the two sums as
$G_1-H_{1,\mathrm{pr}}Q_{\mathrm{pr}}G_1$ and
$G_1D\mathsf R_{\mathsf m}D^{\mathsf T}G_1$.
It follows that their difference is $\mathfrak G$.
Both terms in \eqref{eq:dec-direct-projected-tree} are
defined on the same carrier before summation; the equality
therefore also specifies a single local parameter assignment.

Here is the scale conversion in these block estimates.
A force supported over $y'$ in slot $1$ or $1'$ has forcing
seminorm $\xi_{y'}^3\sup|f|$. The output seminorm contains
$\xi_y|A|$ and $\xi_y^2|\nabla A|$. An eighth-slot column
has forcing seminorm $\xi_{c'}|e|$, while the last output
component has weight $\xi_c^3$. These give, respectively,
the factors $\xi_{y'}^3/\xi_y$, $\xi_{y'}^3/\xi_y^2$,
$\xi_{c'}/\xi_y$, $\xi_{c'}/\xi_y^2$ and
$\xi_{c'}/\xi_c^3$ before scale comparison. Absorb the
ratios of scales, of powers at most three, into half of the
rate $\delta/2$. Moving the chosen position of a row to
either end changes the distance by at most a fixed constant.
This proves \eqref{eq:dec-local-row-columns}--
\eqref{eq:dec-local-projected}, enlarging $B$ and $C_{\mathrm w}$.
The field columns in \eqref{eq:dec-direct-projected-tree}
have no eighth-slot or row-normalization factor.

\emph{The scalar mean projection.}
Let
\[
 F_\varpi=-[\mathbb E^{[2]}_\varpi]_{w,4},
 \qquad Z_{\varpi'}=\mathbb E^{[1]}_{\varpi'}Q'^{\mathsf T}.
\]
Only the values of $Z_{\varpi'}e_z$ at the final supported
tops of $F_\varpi$ are used. Take $C_T=1+C_{\mathrm{av}}c_{\mathrm{av}}/2$.
The local radius \eqref{eq:regular-averaging-radius} and
\eqref{eq:flat-complex-scalar-transport} bound the transport
adjoint by $C_T$; undoing its real gauge preserves the norm.
The extra factor $2$ below allows the bilinear transpose.
The scalar transpose column is
\[
 (Q'^{\mathsf T}e_z)(x)
   =\operatorname{Ad}(T_{\mathbf U}(z,x))^{\mathsf T}e,
 \qquad x\in\Delta(z),
\]
and is zero elsewhere, with norm at most $2C_T|e|$.
The factor $L^{-3j}$ in the average cancels the ratio
$\xi_z^3/\xi^3$ in the transpose. This column depends only
on the configuration in $\Delta(z)$; at level zero it is
the identity column.

For $\varpi=(\kappa_0,\ldots,\kappa_m)$ and
$\varpi'=(\kappa'_0,\ldots,\kappa'_{m'})$, define
\begin{equation}\label{eq:dec-projection-gluing}
 P_{\varpi,\varpi'}(b,b')
 =\sum_{z\in\Lambda_*\cap\kappa_m^{[1/12]}}
       (\xi/\xi_z)^3(DF_\varpi)(b,z)
                     (DZ_{\varpi'})(b',z)^{\mathsf T}.
\end{equation}
The last factor is transposed as a map of the Lie algebra.
The factor $(\xi/\xi_z)^3$ is required by the physical
pairings. A nonzero term forces
$\kappa'_{m'}\in\mathcal N(\kappa_m)$. Attach it to the walk
\[
 (X^{(2)}(\kappa_0),\ldots,X^{(2)}(\kappa_m),
       X^{(2)}(\kappa'_{m'-1}),\ldots,X^{(2)}(\kappa'_0)),
\]
of length $m+m'$. The base cell $X(\kappa'_{m'})$ of the
omitted junction piece is contained in $X^{(2)}(\kappa_m)$. All scalar columns
$\Delta(z)$ and all local factors in
\eqref{eq:dec-projection-gluing} lie in this walk's carrier.

The two column estimates before the sum over $z$ are
\[
 \xi_y^3|(DF_\varpi e_z)(b)|
       \le8B_4q^m e^{-\delta d(y;\varpi;z)}|e|,
 \quad
 \xi_{y'}^3|(DZ_{\varpi'}e_z)(b')|
       \le16C_TB_4q^{m'}\xi_z^4
                    e^{-\delta d(y';\varpi';z)}|e|.
\]
Write $E(z)=d(y;\varpi;z)+d(y';\varpi';z)$ and
$D=D_\omega(y,y')$. Then $E(z)\ge D,d(y,z)$. The weighted
transpose sum is bounded by
\begin{align*}
 &\sum_z(\xi/\xi_z)^3\xi_z^4\xi_y^{-3}\xi_{y'}^{-3}
                           e^{-\delta E(z)}\\
 &\qquad\le L^2c_S(\delta/4)(\xi/\xi_{y'})^3\xi_y^{-2}
                                      e^{-\delta D/2}.
\end{align*}
For completeness, the power outside the exponential on the
left is $(\xi/\xi_{y'})^3\xi_y^{-2}(\xi_z/\xi_y)$.
Scale comparison bounds the last ratio by
$L e^{\lambda d(y,z)}$, where
$\lambda=\log L/(R_1\mathfrak M)\le\delta/8$.
Use $E(z)\ge\frac12D+\frac12d(y,z)$ and sum the
remaining exponential at rate at least $\delta/4$.
This is the $a=a'=3$ case of the weighted scalar gluing
estimate; the larger displayed constant also covers both
choices of bond endpoint.

The transpose on the Lie algebra costs at most $2$.
Thus a pair contributes at most
\[
 B_Pq^{m+m'}(\xi/\xi_{y'})^3\xi_y^{-2}e^{-\delta D/2}.
\]
For a fixed walk of length $n$, there are at most
$(n+1)(27N_{\mathrm G})^{n+1}N_{\mathrm G}$ attached pairs.
Since $q\le\mathfrak M^{-9}$, $n+1\le2^n$ and the floor
may ensure $54N_{\mathrm G}\mathfrak M^{-8}\le1$, their
sum gives \eqref{eq:dec-local-projection}.
The scalar inverse identities give
\[
 \sum_\varpi F_\varpi=G'Q'^{\mathsf T}C',\qquad
 \sum_{\varpi'}Z_{\varpi'}=G'Q'^{\mathsf T},\qquad
 1-\mathsf R_{\mathsf m}
       =G'Q'^{\mathsf T}C'Q'G',
 \quad C'=(Q'G'^2Q'^{\mathsf T})^{-1}.
\]
Absolute convergence allows multiplication and the two
differences term by term. Their sum is $P_{\mathrm f}$.

\emph{Locality and summation.}
Every factor just constructed is an inverse of a finite
local matrix, a local mean derivative, a difference, or one
of the previously constructed multiplier terms. Its entries
are holomorphic on the indicated open local class. Equal
traces of two sequences on a carrier give the same pieces,
row positions and weights, cutoffs, local sequences and
neighbours. They therefore give identical matrices and
identical finite products, including the attached scalar
pairs. In particular none of these assertions requires a
regular continuation to the whole torus. Gauge changes
intertwine each matrix and preserve the physical pairings;
the real structure commutes with the local means and
inverses. These properties pass to every displayed term.

The number of walks of length $m$ with a fixed first cell
is at most $n_{1,\natural}^m$, and the number of the trees
used here is at most $(4n_{1,\natural})^m$.
Choose $C_{\mathrm w}$ to dominate all constants above.
The floor makes
$4n_{1,\natural}C_{\mathrm w}/\mathfrak M\le1/2$.
The geometric series and its first moment prove the two
mass bounds. For each fixed output top, the supports of
the first cutoffs and their first differences permit only
a bounded number of first cells. This is the same local
overlap bound used in the regular parametrix. Consequently
the series converge uniformly on the local classes, in
each finite-dimensional operator norm. All these estimates
also hold after discarding any set of terms.

\emph{The parameters.}
For a graph $\Gamma$ meeting the cells of $\omega$, the
hull lies within sup distance
\[
 \rho=(12r_{\mathrm w}+10)\mathfrak M+4
       \le23r_{\mathrm w}\mathfrak M\le23\mathsf w/88
\]
of $\Gamma$. The extra $4$ permits the first-difference
output window. A tube of radius $\rho$ around a connected
graph of ordinary length $\lambda_\Gamma$ meets at most
\[
 8\max\{1,4\lambda_\Gamma/(\mathsf w-2\rho)\}
\]
partition cubes. To see this, take a maximal $D$-separated
net in the graph, with $2D<\mathsf w-2\rho$. Its balls of
radius $D+\rho$ each meet at most eight cubes. If the net
has at least two points, its disjoint balls of radius $D/2$
each contain graph length at least $D/2$: distance to the
centre is $1$-Lipschitz and takes every value from $0$ to
$D/2$ on the connected graph. Thus the net has at most
$2\lambda_\Gamma/D$ points. Let $2D$ increase to
$\mathsf w-2\rho$. The eight-cube assertion follows by
lifting to the grid; the torus side is at least $3\mathsf w$,
so a set of sup diameter less than $\mathsf w$ meets at most
two intervals in each direction.

In the present scale ordinary length is at most weighted
length. Taking the infimum over $\Gamma$ gives
\[
 \#\mathsf P(\omega)\le8+\frac{68}{\mathsf w}D_\omega(y,y'),
 \qquad
 |s^\omega|e^{-aD_\omega(y,y')}
 \le e^{8\varkappa}
        e^{-(a-68\varkappa/\mathsf w)D_\omega(y,y')}.
\]
The parameter loss is at most $\mathring\delta_0/32$.
Each retained rate is at least $\mathring\delta_0/2$;
we may therefore use $\mathring\delta_0/4$ after inserting
the parameters and summing the masses and first-cell
overlaps. The proof used absolute values, so the same bound
holds for arbitrary subsums. Uniform convergence on bounded
parameter sets also follows directly from the finite number
of partition cubes. Each term is holomorphic and contains
each parameter at most once; grouping by the finitely many
sets $\mathsf P(\omega)$ proves the asserted polynomial
and holomorphy statements.

\emph{Normalization and the response bounds.}
The identity
$Q_{\mathrm{pr}}=\mathsf m^{-1}Q$ gives
$Q_{\mathrm{pr}}^{\mathsf T}=Q^{\mathsf T}\mathsf m^3$,
since dimensionless rows have weight $\xi_c^{-1}$.
Let $B_\sharp$ absorb the masses and local overlaps. The
parameter estimates for the normalized columns are
\begin{gather*}
 |(H_se_{c'})(b)|, |(K_se_{c'})(b)|
    \le B_\sharp e^{8\varkappa}\xi_{y'}^{-1}
                    e^{-\mathring\delta_0d(y,y')/4}|e|,\\
 |(\nabla H_se_{c'})(b,\nu)|,
 |(\nabla K_se_{c'})(b,\nu)|
    \le B_\sharp e^{8\varkappa}\xi_y^{-1}\xi_{y'}^{-1}
                    e^{-\mathring\delta_0d(y,y')/4}|e|,\\
 |S_s(c,c')|\le B_\sharp e^{8\varkappa}
      (\xi_y/\xi_{y'})e^{-\mathring\delta_0d(y,y')/4}.
\end{gather*}
Multiplying the first two lines by $\xi_y$ and $\xi_y^2$
gives the column bound with the explicit ratio
$\xi_y/\xi_{c'}$. The corresponding ratio for
$\mathfrak G_sf$, with input seminorm
$|f|_{y'}^\sharp=\xi_{y'}^3\sup|f|$, is
$(\xi_y/\xi_{y'})^3$. For $P_s$ there are at most
$3(\xi_{y'}/\xi)^3$ bare bonds over $y'$, so its kernel
bound gives
\[
 |P_sA|_y^\sharp\le3B_\sharp e^{8\varkappa}
       (\xi_y/\xi_{y'})e^{-\mathring\delta_0d(y,y')/4}|A|_{y'}
\]
when $A$ is supported over $y'$. The scale hypothesis gives
\[
 (\xi_y/\xi_{y'})^p
 \le L^3e^{\mathring\delta_0d(y,y')/16},\qquad p=1,3.
\]
There are at most six rows at a top. Summing the column
bounds and the bare blocks therefore gives the strong
response estimates with $b_0=6L^3B_\sharp$ and
$\mu=3\mathring\delta_0/16$. Multiplication by an admissible
weight of rate $\theta=\mathring\delta_0/8$ leaves
$\mu-\theta=\mathring\delta_0/16$ for the top sum, proving
the displayed value of $b_{\mathrm D}$.

The stronger column bounds before scale absorption also
give the transpose needed for the local energy. For the
dimensionless row pairing, testing $H_s^{\mathsf T}f$ at
$c$ contributes $\xi_c$; summing bare bonds over $y$
contributes at most $3\xi_y^3$. The factor $\xi_c^{-1}$ in
the first column bound cancels the row factor. Duality in
the Lie algebra costs at most $2$, and hence
\begin{equation}\label{eq:decoupled-whole-transpose}
 |H_s^{\mathsf T}f|_{y_c}
 \le6B_\sharp e^{8\varkappa}
       \sum_y e^{-\mathring\delta_0d(y,y_c)/4}|f|_y^\sharp,
 \qquad
 |H_s^{\mathsf T}f|_w\le\widehat b|f|_w^\sharp.
\end{equation}
After multiplication by the weight, the first inequality
leaves rate $\mathring\delta_0/8$, so the already imposed
summation floor at $\mathring\delta_0/16$ suffices.

Finally, if a term has $\mathsf P(\omega)\subset\mathsf Y$,
its whole carrier lies in $W$, and its endpoint tops lie in
$Y^1(\mathsf Y)$. A path from $Y^j(\mathsf c)$ to outside
$Y^{j+1}(\mathsf c)$ crosses an annulus of weighted width
at least $\mathfrak M/L^2$. Increase the floor so that this
exceeds $r_0$. For $j\le4$ the $r_0$-neighbourhood is then
inside $\operatorname{int}\bigcup\mathsf c$. This proves
the local margins; for $j\le2$ the larger distance to the
complement is at least $(5-j)\mathfrak M/L^2>3r_0$ after
the same choice. All bonds in the first covariant difference
are consequently in $W$, so its weighted norm is the
response seminorm $\mathsf g_y$ used above.

The hull is connected. The partition cubes meeting a
connected set are wall connected: all cubes at a point are
joined by changing one grid coordinate at a time, and a
separation into two wall components would disconnect the
set. Thus $s^\omega\ne0$ places its hull in one wall
component of $\{s\ne0\}$. Absolute convergence gives the
component splitting. The preceding margin shows that sites
in $Y^j$ of different components, $j\le4$, cannot be within
distance $r_0$. Restriction consistency of each local
coefficient gives the claimed component dependence.
\end{proof}

\subsection{The nonlinear response and its current}

Keep the hypotheses, parameters and local regions of
\cref{prop:decoupled-local-propagators}, and put $\eta=\xi$.
Put
\[
 (Y,Y^1,Y^2)=
 (Y^1(\mathsf Y),Y^2(\mathsf Y),Y^3(\mathsf Y)).
\]
Decrease $a_{\mathrm D}$ so that
\[
 a_{\mathrm D}\le\min\left\{a_{\mathrm{nl}},a_{\mathrm{pr}},
       \frac{\min(1,c_{\mathrm{av}})}{8(L+2\pi L^2)}\right\}.
\]
We localize the actual maps of \cref{I-87-84} as follows.
On rows over $Y^1$, use the straightened average
$\mathcal Q=Q+\mathsf C$ and put $\mathsf a=1$; on all other
rows set $Q,\mathsf C,\mathsf a$ equal to zero.
Let $F(A)=\eta^{-1}\sum_{p\subset W}a_p(e^{i\eta A}\mathbf U)$,
where the sum is over plaquettes with all corners in $W$.
On bonds over $Y$, set
\[
 \mathcal V(A)=F'(A)-F'(0)-F''(0)A,\qquad
 \Delta'A=F''(0)A-D^{\mathsf T}DA,
\]
and retain the given current there; set these three maps to
zero on the other bonds. We continue to write $\mathbf J$
for this restricted current. Use the covariant-difference
seminorm $\mathsf g_y$ on $Y^2$ and set it to zero elsewhere.
All row pairings retain weight $\xi_c^{-1}$ and all bare
pairings weight $\eta^3$.

The local margins put each retained double block and every
plaquette touching a bond over $Y^2$ inside $W$. To verify the
whole averaging ball, fix a row $c$ of positive level and use
row units $\eta_c=\xi/\xi_c$. Its double block is a coordinate
box with at most $2\eta_c^{-1}$ sites in each direction.
The real curvature there is at most $L^2\alpha_0\eta_c^2$.
The comb gauge of \cref{I-10-51}, taken entirely in this box,
gives $U^h=e^{i\eta_c B_0}$ with
$\sup|B_0|\le2\pi L^2\alpha_0$. If $A'$ is the perturbation
in $\mathbf U=e^{i\xi A'}U$, then
$\sup|\xi_c\operatorname{Ad}(h)A'|\le L\alpha_1$.
The product logarithm used in
\eqref{eq:regular-piece-potential} therefore writes
$\mathbf U^h=e^{i\eta_c B}$, where
\[
 \sup|B|\le2(L\alpha_1+2\pi L^2\alpha_0)
                    \le c_{\mathrm{av}}/4.
\]
Here the sum in parentheses is at most $1/8$, so the same
series logarithm is defined. On $\sup|A|<c_4$, the relative
field $\xi_c\operatorname{Ad}(h)A$ has norm less than
$Lc_4\le c_{\mathrm{av}}/2$.
Apply \cref{prop:flat-complex-averaging} at this complex base.
Unitary covariance of its domains and recursion identifies
its relative map with the actual straightened average.
The counting-kernel rescaling is
$\xi_c\eta_c^3=\xi^3\xi_c^{-2}$, exactly as in
\eqref{I-87-77}. At level zero the recursion is the identity.
Thus the averaging, action and transpose estimates of
\cref{I-87-84} apply with their existing constants. All gauges
and stencils remain in $W$. Fixing the comb gauge in a
neighborhood gives local holomorphy in the complex bond data;
the intrinsic recursion identifies the overlapping charts.

Choose the fixed geometric floor so that the carrier margins exceed
$3r_0$. Use $c_W,C_W$ from \eqref{I-86-46}, with the original
averaging constants. The five operators satisfy
\eqref{I-86-50} with $\widehat b=b_{\mathrm D}e^{8\varkappa}$:
the new row weight satisfies $0\le\mathsf a\le1$, hence
$|(S_s-\mathsf a)X|_w\le2\widehat b|X|_w$.
The reduced field $\mathsf W_s$ is \eqref{I-86-62} with these
localized maps and parameter operators.

For the multiplier estimate we use the quadratic energy in its
bilinear form. If the local structure is carried by
$Y\subset Y^1\subset Y^2$, put
\[
 E_y(A,\delta)=\eta^3\sum_{p\ {\rm over}\ y}
       \langle D_{\mathbf U}A(p),D_{\mathbf U}\delta(p)\rangle
       \quad(y\in Y^1),\qquad E_y=0\quad(y\notin Y^1),
 \qquad E=\sum_yE_y.
\]
All these plaquettes lie in the local region. The count and
curl estimate in \eqref{eq:response-local-energy} apply unchanged:
\[
 |E_y(A,\delta)|\le12\xi_y^{-1}\mathsf g_y(A)\mathsf g_y(\delta).
\]
This calculation is bilinear and remains valid at a complex
regular pair; it asserts no positivity there. For real
$A$ at a real pair, $E(A,A)\ge0$, and in the global case
$E(A,A)=(DA,DA)_\eta$.

The local action estimate, at level $n(y)$, gives
\[
 |\Delta'A|^\sharp_y
       \le C_E\max_{N(y)}|A|,
 \qquad
 \Delta'=F''(0)-D^{\mathsf T}D
\]
on the retained bonds. In fact its unscaled error is bounded
by $C_{\mathrm a}\alpha\,\sup|\widehat A|$,
where $\widehat A=\xi_yA$ and
$\alpha=L^2\alpha_0\le1/2$; the adjacent-level ratio
$\xi_y/\xi_{y'}\le L$ gives the displayed bound with
$C_E\ge C_{\mathrm a}L$.
By the definition of the Taylor remainder,
\begin{equation}\label{eq:response-energy-identity}
 (F'(A)-F'(0),\delta)_\eta
       =E(A,\delta)+(\Delta'A+\mathcal V(A),\delta)_\eta
\end{equation}
whenever $\delta$ vanishes off the bonds over $Y$.
Indeed, first add
$\Delta'A=F''(0)A-D^{\mathsf T}DA$ and
$\mathcal V(A)=F'(A)-F'(0)-F''(0)A$ at those bonds.
In the resulting pairing, a plaquette contributes to
$(DA,D\delta)_\eta$ only if it has a bond over $Y$.
Its base lies in $N_5(Y)\subset Y^1$, so grouping these
plaquettes by their bases gives exactly $E(A,\delta)$.

Let $H$ satisfy the column and whole-field transpose estimates
used in \cref{I-88-53}, with weighted constant $\widehat b$.
Define $H^\circledast(A)$ by
\[
 (H^\circledast(A),X)_1
       =E(A,HX)+(\Delta'A+\mathcal V(A),HX)_\eta .
\]
The nondegenerate row pairing defines it uniquely. For
$\|A_i\|\le2c_W/\widehat b$ and every admitted weight $w$,
\begin{equation}\label{eq:response-energy-dual-bound}
 |H^\circledast(A_1)-H^\circledast(A_2)|_w
 \le c_\circledast\widehat b\,\|A_1-A_2\|_w,
 \qquad c_\circledast=100(1+C_E+C_{\mathcal V})\mathsf e .
\end{equation}
Use the normalized column estimate proved in
\cref{prop:decoupled-local-propagators}, before its scale ratio
is absorbed into the exponential. For a test row $c$ over $y'$ and a bond $b$ over $y$,
\[
 |(He_c)(b)|\le B_\sharp e^{8\varkappa}\xi_c^{-1}
                    e^{-\mathring\delta_0d(y,y')/4},\qquad
 \mathsf g_y(He_c)\le B_\sharp e^{8\varkappa}
       \frac{\xi_y}{\xi_c}e^{-\mathring\delta_0d(y,y')/4}.
\]
The row pairing contributes $\xi_c$. It cancels the factor
$\xi_y^{-1}(\xi_y/\xi_c)$ in the energy term. In the lower-order
pairing the bare volume and the three bond directions leave
the factor $\xi_c\xi_c^{-1}$ as well. Thus no scale quotient
remains in either sum. The local Lipschitz estimate for
$\mathcal V$, the bound on $\Delta'$, and complex bilinear
duality bound their difference by
\[
 (24+6C_E+48C_{\mathcal V})B_\sharp e^{8\varkappa}
 \sum_y e^{-\mathring\delta_0d(y,y')/4}
                   \max_{N(y)}\|A_1-A_2\|.
\]
An admissible weight costs at most
$e^{\mathring\delta_0d(y,y')/8}$, and the neighborhood maximum
costs $\mathsf e$. The remaining sum is at most
$c_S(\mathring\delta_0/16)$. Since
$b_{\mathrm D}\ge6L^3B_\sharp c_S(\mathring\delta_0/16)$,
this proves \eqref{eq:response-energy-dual-bound} with the
displayed, larger constant. This choice does not change the
averaging constant $c_\flat$ or the contraction radii.

Write $J=F'(0)$ in the global case. Then
\eqref{eq:response-energy-identity} says
$H^{\mathsf T}(F'(A)-J)=H^\circledast(A)$.
Thus, with
\[
 \Psi_A=(I+H^{\mathsf T}\mathsf C'_A{}^{\mathsf T})^{-1},
 \quad
 \omega(A)=\Psi_A[H^{\mathsf T}J+H^\circledast(A)],
 \quad
 m_\omega=8\widehat b\,\gamma_0+2c_\circledast c_W,
\]
the previously proved inverse and transpose bounds give
\begin{gather}
 \sup|\omega(A)|\le2m_\omega,\label{eq:response-multiplier-size}\\
 |\omega(A_1)-\omega(A_2)|_w
 \le(32C_3\mathsf e\,m_\omega+2c_\circledast)
                      \widehat b\,\|A_1-A_2\|_w.
 \label{eq:response-multiplier-difference}
\end{gather}
For the second inequality subtract $\Psi_{A_i}\varphi_i$:
the inverse difference is at most
$32C_3\mathsf e\,\widehat b\,m_\omega\|A_1-A_2\|_w$,
and $\|\Psi_{A_2}\|_w\le2$ handles the remaining term by
\eqref{eq:response-energy-dual-bound}.
At a critical reference $J=Q^{\mathsf T}\lambda_0$ and
$QH=I$, one has $\omega(0)=H^{\mathsf T}J=\lambda_0$.
These estimates also retain the possibly nonzero initial
multiplier.

Define the multiplier and current using the energy already
introduced:
\begin{equation}\label{eq:decoupled-current-definition}
 \omega_s(A)=
  (I+H_s^{\mathsf T}\mathsf C'_A{}^{\mathsf T})^{-1}
       [H_s^{\mathsf T}\mathbf J+H_s^\circledast(A)],
 \qquad
 j_s(A)=f_A^{-1}(Q+\mathsf C'_A)^{\mathsf T}\omega_s(A)
                         -Q^{\mathsf T}H_s^{\mathsf T}\mathbf J .
\end{equation}
Here $f_A=f(-i\eta\operatorname{ad}A)$ acts bondwise,
$f(z)=(e^z-1)/z$, with its removable value at zero.
It satisfies $\|f_A^{-1}\|\le2$ on the present ball.
For $\|A_i\|\le2c_W/\widehat b$ one has
\begin{equation}\label{eq:decoupled-current-lipschitz}
 |j_s(A_1)-j_s(A_2)|^\sharp_w
 \le C_j\widehat b(1+\gamma_0\widehat b)
                      \|A_1-A_2\|_w,\qquad j_s(0)=0,
\end{equation}
where $C_j=C_j(L)>0$ is fixed below.
Here and below the constants are uniform in the number of
layers.

We include the estimate because differentiating the current
pointwise would lose a derivative. Put
$g_i=(Q+\mathsf C'_{A_i})^{\mathsf T}\omega_s(A_i)$ and
$m=8\widehat b\gamma_0+2c_\circledast c_W$.
The transpose estimates for $Q,\mathsf C'$ and the preceding
multiplier estimates give, for $\delta=\|A_1-A_2\|_w$,
\[
 \sup|g_i|^\sharp\le2(c_\flat+1)m,\qquad
 |g_1-g_2|^\sharp_w
 \le(c_\flat\mathsf e+1)|\omega_s(A_1)-\omega_s(A_2)|_w
                                      +2C_3\mathsf e m\delta.
\]
The series for $f$ gives $\|f(z)-I\|\le1/14$ and
$\|f(z_1)-f(z_2)\|\le(4/7)\|z_1-z_2\|$.
The inverse difference identity then gives
$|f(z_1)^{-1}-f(z_2)^{-1}|\le|z_1-z_2|$ for
$|z_i|\le1/8$. Since
$|\operatorname{ad}X|\le2|X|$ and
$\eta|A_1(b)-A_2(b)|\le|A_1-A_2|_{y_b}$,
\[
 |j_s(A_1)-j_s(A_2)|^\sharp_w
 \le2|g_1-g_2|^\sharp_w+2\delta\sup|g_2|^\sharp .
\]
Inserting the multiplier difference bound gives
\[
 [4c_\flat\mathsf e(32C_3\mathsf e\widehat b m
       +2c_\circledast\widehat b)
       +4C_3\mathsf e m+8c_\flat m]\delta.
\]
Since $c_W\le1$ and $\widehat b\ge1$, one has
$m\le(8+2c_\circledast)(1+\gamma_0\widehat b)$.
A fixed $C_j(L)$ therefore bounds the displayed expression by
the right side of \eqref{eq:decoupled-current-lipschitz}.
The inverse series and the energy formula prove joint
holomorphy. At zero $\Psi_0=I$, $H_s^\circledast(0)=0$
and $f_0=I$, which proves $j_s(0)=0$ with the subtraction
in \eqref{eq:decoupled-current-definition}.

\begin{theorem}\label{I-88-110}
In the preceding setting put
\[
 \varepsilon_{\mathrm D}
    =\min\{c_W/(4b_{\mathrm D}^2),
       1/(8C_Wb_{\mathrm D}^4),\varepsilon_{\mathrm{crit}}\}.
\]
For $\sup|B'|\le\varepsilon_{\mathrm D}e^{-32\varkappa}$
there is a unique triple $(X,\mathbf A,\mathbf H)$ in the
balls
\[
 \sup|X|\le c_W\widehat b^{-2},\qquad
 \|\mathbf A\|\le\varepsilon_{\mathrm D}b_{\mathrm D}
                                      e^{-24\varkappa}
\]
solving
\[
 A'=K_sB'+\mathbf A,\quad
 X=\mathsf C(A'-H_sX),\quad
 \mathbf A+\mathfrak G_s\mathsf W_s(A')=0,\quad
 \mathbf H=A'-H_sX.
\]
The two successive iterations converge to it, and
\[
 \|\mathbf H\|\le4\widehat b\sup|B'|,\qquad
 \sup|X|\le16C_2\widehat b^2(\sup|B'|)^2,\qquad
 \|\mathbf A\|\le8C_W\widehat b^5(\sup|B'|)^2 .
\]
Set $j(s;B')=j_s(\mathbf H(s;B'))$ and
$\mathbf H^P(s;B')=P_s\mathbf H(s;B')$.
For every weight $w\in\mathcal W_\theta$ the three maps
are Lipschitz in $B'$ with respective constants
\[
 4\widehat b,\qquad 8C_j\widehat b^3,\qquad4\widehat b^2
\]
in $\|\cdot\|_w$, $|\cdot|^\sharp_w$ and
$|\cdot|^\sharp_w$. In particular, changing only datum rows
over $y'$ gives decay $e^{-\theta d(y,y')}$ at $y$.
All three maps vanish at $B'=0$.

They are jointly holomorphic in the interior parameter
polydisc, in $B'$, and in $(A_0,\mathbf J)$ on the open
local natural class with a fixed real factor
$\mathbf U=e^{i\eta A_0}U$.
They split as sums over the wall components $\mathsf c_i$
of $\{s\ne0\}$. The response and $\mathbf H^P$ of a component
are supported over $Y^1(\mathsf c_i)$; its current is
supported over $N_{\rho_0}(Y^1(\mathsf c_i))
\subset Y^2(\mathsf c_i)$.
Each depends only on the pair in $W_{\mathsf c_i}$, the
parameters on $\mathsf c_i$, and the datum rows over
$Y^1(\mathsf c_i)$. It vanishes when those datum rows vanish.
The maps are unchanged by extending the pair to a larger
local natural class while keeping $s=0$ outside the old set.
They commute with unitary gauge transformations and with
complex conjugation; in particular real inputs give real
outputs.
\end{theorem}

\begin{proof}
The parameter estimates proved above give
\[
 \varepsilon_{\mathrm D}e^{-32\varkappa}
 \le\min\{c_W/(4\widehat b^2),
                         1/(8C_W\widehat b^4)\}.
\]
Thus \cref{I-87-35} applies with the displayed balls.
Its weighted estimate gives the response bound; the first
fixed-point equation gives
$|X|\le4C_2|A'|^2\le16C_2\widehat b^2|B'|^2$.
Composition with \eqref{eq:decoupled-current-lipschitz},
using $\gamma_0\le1\le\widehat b$, gives the current
constant. Applying $P_s$ gives the third constant.
Exponential weights centred at $y'$ give the stated decay.
Uniform convergence of both iterations gives joint
holomorphy. The domain in $(A_0,\mathbf J)$ is open:
the real-factor condition is fixed and the other conditions
are finitely many strict inequalities between continuous
functions.

We give the locality argument, including the inverse in
$\mathsf W_s$. Transposes preserve both-variable carriers
because the pairings have diagonal, nonzero weights.
The kernels of $Q,\mathsf C'_T,\mathsf C''[T]$ have range
$\rho_0$; at a row or bond near a component their transpose
sums therefore use only data within its $2\rho_0$-neighbourhood.
If two fields agree on $N_{r_0}(Y_i)$, their first
fixed-point solutions agree on $N_{\rho_0}(Y_i)$.
To see this, keep the full solution $X$ on that set and
define it elsewhere by
$\mathsf C(A''-H_iX)$ for the component system.
Since $H_i$ uses only rows in $Y_i$, the resulting vector
solves the component fixed-point equation and remains in
the same contraction ball. Uniqueness gives the claimed
agreement and agreement of $A'-H_sX$ on $N_{r_0}(Y_i)$.
Further,
\[
 (I+H_s^{\mathsf T}\mathsf C'_T{}^{\mathsf T})^{-1}
   =\sum_{k\ge0}(-H_s^{\mathsf T}
                          \mathsf C'_T{}^{\mathsf T})^k .
\]
Every term agrees near $Y_i$ with the corresponding term
for the component system: $\mathsf C'_T{}^{\mathsf T}$
uses the $\rho_0$-neighbourhood and $H_s^{\mathsf T}$ then
restricts to $Y_i$. Induction proves the assertion for
all powers. The defining six terms of $\mathsf W_s$
therefore agree over $Y_i$ with the component reduced field.

Let $\mathbf A_i$ be the component solutions and put
$\mathbf A=\sum_i\mathbf A_i$.
A $\rho_0$-neighbourhood meets at most one $Y_i$, because
different carriers are separated by more than
$r_0=2\rho_0$. Hence
$\|\mathbf A\|=\max_i\|\mathbf A_i\|$; the same observation
controls the first fixed-point arguments.
The local agreement just proved gives
\[
 -\mathfrak G_s\mathsf W_s(K_sB'+\mathbf A)
   =-\sum_i\mathfrak G_i\mathsf W_i(K_iB'+\mathbf A_i)
   =\sum_i\mathbf A_i.
\]
Uniqueness in the second contraction now gives the
decomposition of $\mathbf H$.

For the current, $H_i^\circledast$ uses only $E_y$ with
$y\in N_{\rho_0}(Y_i)$ and $\Delta',\mathcal V$ over $Y_i$.
The preceding inverse-series comparison therefore identifies
the multipliers on each component. The transpose of
$Q+\mathsf C'$ enlarges its support only to
$N_{\rho_0}(Y_i)$; these enlarged supports are disjoint.
This proves the current decomposition.
The identity $P_i\mathbf H_j=0$ for $i\ne j$ proves the
decomposition of $\mathbf H^P$.

Finally, restricting the pair to $W_{\mathsf c_i}$ changes
none of the formulas used in these comparisons.
In particular all plaquettes touching a bond over $Y_i$
and all relevant averaging blocks lie in that region.
Thus $Q,\mathsf C$ agree near $Y_i$,
$\mathcal V,\Delta',\mathbf J$ agree over $Y_i$, and
$E_y,g_y$ agree on the enlarged carrier. The same argument
proves invariance under extensions of the local pair.
Dependence on the datum only through $K_iB'$ proves the
remaining assertions, including vanishing at zero.
The walk terms, local averages, action derivatives and
energy pairings intertwine unitary gauge transformations
and conjugation. Both convergent iterations and the
inverse series therefore intertwine them as well.
\end{proof}

\subsection{Identification with the critical family}

In this subsection the determining set is the reduced set
$\mathbf B^\flat$, as in \cref{I-88-93}. All actual fibers
and gauge factors refer to these rows and their top nodes.
At $\mathsf Y=\mathsf P$ and $s=1$, all monomials
$s^\omega$ equal one. The complete expansions therefore
sum to the operators of \cref{I-88-53}:
\[
\begin{gathered}
 S=(QGQ^{\mathsf T})^{-1},\qquad H=GQ^{\mathsf T}S,\qquad
 K=G_1Q^{\mathsf T}(QG_1Q^{\mathsf T})^{-1},\\
 \mathfrak G=G_1-G_1D\mathsf R_{\mathsf m}D^{\mathsf T}G_1
                         -KQG_1,\qquad
 P_1=D(1-\mathsf R_{\mathsf m})D^{\mathsf T}.
\end{gathered}
\]
Thus the response at one is independent of the partition
and its parameter radius. It is defined on the full ball
$\sup|B'|<\varepsilon_{\mathrm D}$, using the global
estimates with $\widehat b=b_{\mathrm D}$.
In this subsection $\mathbf H(B')$ denotes that global
response, whether or not decoupling parameters are present.

\begin{proposition}\label{I-88-127}
Let $U_0$ be real, let
$(U_0,J^\circ(U_0))$ belong to the global natural class
with radii at most $a_{\mathrm D}$, and assume
$J^\circ(U_0)=Q^{\mathsf T}\lambda_0$.
Use the mean-Landau projection $\mathsf R_{\mathsf m}$ and
the row pairing with weights $\xi_c^{-1}$.
For $\sup|B'|<\varepsilon_{\mathrm D}$,
$A=\mathbf H(B')$ satisfies
\[
 \mathcal Q(A)=B',\qquad
 \mathsf R_{\mathsf m}D^{\mathsf T}A=0,\qquad
 F'(A)=(Q+\mathsf C'_A)^{\mathsf T}\omega(A).
\]
Moreover
\[
 \mathbf J+j_1(A)=J^\circ(e^{i\eta A}U_0),\qquad
 P_1A=DD^{\mathsf T}A,\qquad
 \omega(0)=\lambda_0.
\]
The multiplier obeys
\[
\begin{gathered}
 \sup|\lambda_0|\le8b_{\mathrm D}\gamma_0,\qquad
 \sup|\omega(A)|\le2m_\omega,\\
 |\omega(\mathbf H(B_1'))-\omega(\mathbf H(B_2'))|_w
 \le4(32C_3\mathsf e m_\omega+2c_\circledast)b_{\mathrm D}^2
                                      |B_1'-B_2'|_w .
\end{gathered}
\]
where $m_\omega=8b_{\mathrm D}\gamma_0+
2c_\circledast c_W$.

Conversely, if
\[
 \|A\|\le r_1:=\frac{\varepsilon_{\mathrm D}}{4c_\flat},
 \quad \mathsf R_{\mathsf m}D^{\mathsf T}A=0,\quad
 F'(A)=(Q+\mathsf C'_A)^{\mathsf T}\lambda ,
\]
then $\sup|\mathcal Q(A)|<\varepsilon_{\mathrm D}$ and
\[
 A=\mathbf H(\mathcal Q(A)),\qquad \lambda=\omega(A).
\]
In particular this applies to a real $A$ in that ball
and slice whenever $e^{i\eta A}U_0$ is a critical
configuration for its actual determining datum.
\end{proposition}

\begin{proof}
The equations in the first display are \cref{I-88-93}.
The energy identity gives
$\omega(A)=
(I+H^{\mathsf T}\mathsf C'_A{}^{\mathsf T})^{-1}
H^{\mathsf T}F'(A)$.
Since $QH=I$, the multiplier equation implies
$\omega(A)=\lambda$, and
$Q^{\mathsf T}H^{\mathsf T}J=J$ at the reference.
Consequently the current definition gives
$J+j_1(A)=f_A^{-1}F'(A)=J^\circ(e^{i\eta A}U_0)$.
The last equality is the local action derivative identity,
at every bare bond, including incomplete boundary stars.
The slice equation gives $P_1A=DD^{\mathsf T}A$.
The preceding energy bounds and the response Lipschitz
estimate give the three multiplier estimates.

For the converse set
\[
 X=\mathsf C(A),\quad A'=A+HX,\quad
 B'=\mathcal Q(A),\quad \mathbf A=A'-KB'.
\]
The local estimates and the definition of $r_1$ give
\[
 |B'|\le2c_\flat r_1=\varepsilon_{\mathrm D}/2,\quad
 |X|\le C_2r_1^2\le c_Wb_{\mathrm D}^{-2},\quad
 \|A'\|\le2r_1<c_W/b_{\mathrm D},\quad
 \|\mathbf A\|\le b_{\mathrm D}\varepsilon_{\mathrm D}.
\]
The first fixed-point uniqueness gives $X=\mathsf D(A')$.
The identities $QH=QK=I$ show $QA'=B'$ and
$Q\mathbf A=0$. The on-shell identities
$\mathsf R_{\mathsf m}D^{\mathsf T}H
=\mathsf R_{\mathsf m}D^{\mathsf T}K=0$
put both $A'$ and $\mathbf A$ in the mean-Landau slice.
Write $\mathcal L(A')=A'-H\mathsf D(A')$ and
$K_A=Q+\mathsf C'_A$.
The differentiated first fixed-point equation gives
$K_A\mathcal L'=Q$.
The reduced-action identity used in \cref{I-88-93}
therefore yields, for every $\delta$ in that slice,
\[
 (J+\Delta_1A'+\mathsf W_1(A'),\delta)_\eta
       =(F'(A),\mathcal L'\delta)_\eta
       =(Q^{\mathsf T}\lambda,\delta)_\eta .
\]
The difference lies in $\operatorname{ran}(D\mathsf R_{\mathsf m})$.
Apply $\mathfrak G$, using the already proved on-shell
relations
\[
 \mathfrak GD\mathsf R_{\mathsf m}=0,\quad
 \mathfrak GQ^{\mathsf T}=0,\quad
 \mathfrak G\Delta_1A'=A'-KQA'.
\]
Since $J=Q^{\mathsf T}\lambda_0$, this gives
$\mathbf A+\mathfrak G\mathsf W_1(A')=0$.
The triple is in the two contraction balls, so uniqueness
identifies $A=\mathbf H(B')$ and then $\lambda=\omega(A)$.

It remains to justify the multiplier for a real critical
configuration. The estimate
$\|\mathsf C'_AH\|_\infty\le4C_2b_{\mathrm D}r_1\le1/2$
shows that $K_AH=I+\mathsf C'_AH$ is invertible.
Thus $K_A$ is onto and $K_A^{\mathsf T}$ is injective.
For any real $\delta$ in $\ker K_A$, the finite-dimensional
implicit function theorem gives
\[
 A_\tau=A+\tau\delta+H\beta(\tau),\qquad
 \mathcal Q(A_\tau)=\mathcal Q(A),\qquad
 \beta(0)=\beta'(0)=0.
\]
Let $v[A]$ be the node factors in the actual-average
identity. Set $g_\tau=v[A]v[A_\tau]^{-1}$ on the reduced top set $\Lambda_*$ and $g_\tau=I$ at
other bare sites. Each top has one prescribed value, so this
defines a real gauge transformation. The formula for the actual averages,
including its outward twist, gives
\[
 M_{\mathbf B^\flat}((e^{i\eta A_\tau}U_0)^{g_\tau})
                        =M_{\mathbf B^\flat}(e^{i\eta A}U_0).
\]
Criticality and gauge invariance therefore imply
$(F'(A),\delta)_\eta=0$.
The real pairings are positive definite; hence the
annihilator of $\ker K_A$ is the range of its real
adjoint. There is a real $\lambda$ with
$F'(A)=K_A^{\mathsf T}\lambda$, unique by injectivity.
\end{proof}

\begin{corollary}\label{I-88-131}
Let $U_0=\Phi_{V_0}(0)$ be the canonical real background,
and write the mean-Landau pair of its holomorphic
continuation as
\[
 e^{i\eta\mathcal H(B)}U_0
       =\Phi_{V_0}(B)^{u_B^{-1}},\qquad
 \mathsf R_{\mathsf m}D^{\mathsf T}\mathcal H(B)=0,\qquad
 u_Bv[\mathcal H(B)]=I\quad\hbox{on }\Lambda_* .
\]
Assume its reference pair lies in the natural class with
radii at most $a_{\mathrm D}$. On the common ball
\[
 \sup|B|<\min\{r_B,\varepsilon_{\mathrm D}/C_{\mathrm{tw}}\},
\]
where $r_B$ is the radius of this pair and the inverse
datum twist, one has
\[
 \mathcal H(B)=\mathbf H(1;B_\natural),\qquad
 \mathcal Q(\mathcal H(B))=B_\natural.
\]
For real $B$ on this ball,
\[
 u_B^{-1}=v[\mathcal H(B)]\quad\hbox{on }\Lambda_*,
 \qquad
 M_{\mathbf B^\flat}(e^{i\eta\mathcal H(B)}U_0)
                 =(e^{iB}V_0)^{v[\mathcal H(B)]}.
\]
\end{corollary}

\begin{proof}
Criticality of $U_0$ gives the real initial multiplier by
the preceding argument at zero.
For real $B$ sufficiently near zero, $\mathcal H(0)=0$
puts $\mathcal H(B)$ in the radius-$r_1$ ball.
Gauge covariance preserves criticality, so the proposition
identifies it with
$\mathbf H(\mathcal Q(\mathcal H(B)))$.
Compare its actual-average identity with
$M_{\mathbf B^\flat}(e^{i\eta\mathcal H(B)}U_0)
=(e^{iB}V_0)^{v[\mathcal H(B)]}$.
Cancellation of the node factors and $V_0$ gives
$e^{i\mathcal Q(\mathcal H(B))^\natural}=e^{iB}$.
Both logarithms are in the fixed logarithm ball; hence
$\mathcal Q(\mathcal H(B))^\natural=B$ and
$\mathcal Q(\mathcal H(B))=B_\natural$.
This proves the response identity near real zero.
Both sides are holomorphic on the stated common complex
ball. Their difference vanishes on a real neighbourhood
of zero, hence on its complexification, and then along
each complex radial line on the whole ball.
The remaining equations follow from the on-shell response
identity and the normalization of $u_B$.
\end{proof}

\section{Gaussian normalizers near the trivial configuration}\label{sec:gaussian-normalizers}

\subsection{Critical germs and constrained Gaussian integrals}

Fix a height $k\ge0$ and a finite side
$\mathfrak s\in\{2L^n:n\in\mathbb N\}$. In the level-$k$
scale put
\[
 \Lambda_j=L^{j-k}\mathbb Z^3/\mathfrak s\mathbb Z^3,
 \qquad\ell_j=L^{j-k},\qquad\eta=L^{-k}.
\]
In this subsection a form $B$ is the dimensionless coordinate
in $e^{iB}U$ on the positively oriented bonds. The difference $d$ has no spacing factor, and $\mathsf r$
denotes restriction of site functions to the target level.
The inner products are unweighted sums of $\langle X,Y\rangle
=\frac12\operatorname{tr}(XY)$ on bonds or sites. Thus
$B=\ell_j A$ for the physical field $A$ at level $j$.
The covariant gauge differential is
$\mathsf D_U\lambda(b)=\operatorname{Ad}(U(b))\lambda(b_+)
-\lambda(b_-)$, without a factor $\ell_j^{-1}$.
Reversed tangent coordinates are continued by
$B(\bar b)=-\operatorname{Ad}(U(b))^{-1}B(b)$.
All assertions near the trivial configuration in this
subsection concern germs; their neighbourhoods may depend
on $k$ and $\mathfrak s$.

Let $\mathcal V_j,\mathcal A_j,\mathfrak G_j$ denote
configurations, forms and site functions on $\Lambda_j$.
Use the existing averaging map $\mathcal M$ between
successive levels, and set
\[
 A(U)=\sum_p[1-\tfrac12\operatorname{Re}\operatorname{tr}\partial U(p)],
 \qquad
 \mathrm{GF}_j(V)=\sum_{y\in\Lambda_j}
       \sum_{x\in B(y)\setminus\{y\}}
                 [1-\tfrac12\operatorname{Re}\operatorname{tr}V(y,x)].
\]
No spacing factor enters these functions. Write
$\mathsf Q=\mathcal M'(\mathbf1)$,
$\mathsf Q^{[j]}=(\mathcal M^j)'(\mathbf1)$, and
$\mathsf L_VB(x)=\frac1i\frac d{dt}|_0
\log(e^{itB}V)(y,x)$ for $x\in B(y)\setminus\{y\}$.
At $V=\mathbf1$ this is the mean contour sum
$\mathsf LB(x)=\bar B(y,x)$.

For a surjective linear map $M:E\to F$ of real Euclidean
spaces and a positive definite form $T$ on $\ker M$, use
the constrained Gaussian integral
\begin{equation}\label{eq:gaussian-coarea-normalizer}
 I(T,M)=\det(MM^{\mathsf T})^{-1/2}
                  \int_{\ker M}e^{-T(B)/2}\,dB,
 \qquad
 \mathbb S(X,M)=\begin{pmatrix}X&M^{\mathsf T}\\M&0\end{pmatrix}.
\end{equation}
The Lebesgue measure is that of the induced inner product.

\begin{theorem}\label{thm:gaussian-critical-germs}
\label{I-108-32}\label{I-108-62}
There are unique real-analytic germs $\mathcal S_j$ on
$\mathcal V_j$, $0\le j\le k$, and $p_j:\mathcal V_j\to
\mathcal V_{j-1}$, $1\le j\le k$, with
$\mathcal S_0=A$, $\mathcal S_j(\mathbf1)=0$,
$\nabla\mathcal S_j(\mathbf1)=0$, $p_j(\mathbf1)=\mathbf1$ and the following
properties. Put $\mathcal F_j=\mathcal S_{j-1}+\mathrm{GF}_j$.
For $W$ near $\mathbf1$, $p_j(W)$ is the unique critical
point near $\mathbf1$ of $\mathcal F_j$ on
$\{V:\mathcal M(V)=W\}$. Its restricted Hessian $T_{j,W}$
is positive definite; it is a strict local minimum,
is axial on $\Lambda_j$, and
\[
 \mathcal M(p_j(W))=W,\qquad
 \mathcal S_j=\mathcal S_{j-1}\circ p_j
             =\mathcal F_j\circ p_j\ge0.
\]
These germs are covariant under lattice isometries mapping
$\Lambda_j$ onto itself and,
for a coarse gauge transformation $v$ near $\mathbf1$ or
constant, under its block-constant extension $\tilde v$:
\[
 p_j(W^v)=p_j(W)^{\tilde v},\quad
 p_j(\mathsf sW)=\mathsf s p_j(W),\quad
 \mathcal S_j(W^v)=\mathcal S_j(\mathsf sW)=\mathcal S_j(W).
\]
Here each equality is read on a common neighbourhood where
its two members are defined. Their Hessians $h_j$ at
$\mathbf1$ have null space $\{B:dB=0\}$, and
\[
 h_0(B)=|dB|^2,\qquad
 h_j(B)=\min_{\mathsf QA=B}
                \{h_{j-1}(A)+|\mathsf LA|^2\}.
\]
The minimizer is $p_j'(\mathbf1)B$.

Set $c_j=p_1\circ\cdots\circ p_j$, with $c_0$ the identity.
Then
\begin{equation}\label{eq:gaussian-germ-multiplier}
 \mathcal M^jc_j=I,\quad\mathcal M^{j-1}c_j=p_j,
 \quad\mathcal S_j=A\circ c_j,
 \qquad
 \nabla A(c_j(W))=(\mathcal M^j)'(c_j(W))^{\mathsf T}
                                      \nabla\mathcal S_j(W).
\end{equation}
In particular $c_j(W)$ is critical on its iterated-average
fibre. Let $W_k=W$, $W_{j-1}=p_j(W_j)$, $U=W_0=c_k(W)$,
and $\mathcal Q_U^{[j]}=(\mathcal M^j)'(U)$. If $H_W$ is
the Hessian of $A$ on the fibre $\mathcal M^k=W$ at $U$, put
\[
 \mathbf L_W=(\mathsf L_{W_{j-1}}\mathcal Q_U^{[j-1]})_{j=1}^k,
 \quad Z_j(W)=I(T_{j,W},\mathcal M'(p_j(W))),
 \quad\Phi_j^\natural(W)=\log Z_j(W)-\log Z_j(\mathbf1).
\]
The form $H_W+|\mathbf L_W\,\cdot\,|^2$ is positive
definite on $\ker\mathcal Q_U^{[k]}$, and
\begin{equation}\label{eq:gaussian-product}
 \mathcal Z_k(W):=I(H_W+|\mathbf L_W\,\cdot\,|^2,
                                  \mathcal Q_U^{[k]})
       =\prod_{j=1}^kZ_j(W_j).
\end{equation}
Consequently the normalized logarithm
$\Psi_k=\log\mathcal Z_k-\log\mathcal Z_k(\mathbf1)$,
with $\Psi_0=0$, satisfies
\begin{equation}\label{eq:gaussian-recursion}
 \Psi_k=\Phi_k^\natural+\Psi_{k-1}\circ p_k.
\end{equation}\label{I-108-63}
All these logarithms are real analytic and invariant under
the stated gauge transformations and isometries. In
\eqref{eq:gaussian-recursion}, the function $\Psi_{k-1}$
is transported from the tower of height $k-1$ and side
$L\mathfrak s$ by the dilation $x\mapsto Lx$; forms are
relabelled without an amplitude factor.
\end{theorem}

\begin{proof}
The mean map is analytic near $\mathbf1$. Its derivative is
onto: a form supported on the pivot bonds has
$\mathsf QA(c)=L^{-2}A(b_0(c))$, so the values on the
distinct pivots prescribe every output independently.
Their contours stay inside blocks, hence
$\mathsf LA=0$ for such a form. Full rank persists near
$\mathbf1$, making all the fibres below analytic
submanifolds. The gauge-fixing function is nonnegative and
vanishes exactly on the axial slice. At an axial point,
writing $Y_x(V)=i^{-1}\log V(y,x)$ and using
$1-\cos|Y|=|Y|^2/2+O(|Y|^4)$ gives
$\nabla\mathrm{GF}_j=0$ and
$\mathrm{GF}_j''(V)(B)=|\mathsf L_VB|^2$.

Suppose $\mathcal S_{j-1}$ has been constructed and its
Hessian has the asserted null space. The form
$g(A)=h_{j-1}(A)+|\mathsf LA|^2$ is positive definite on
$\ker\mathsf Q$. Indeed, $g(A)=0$ implies $dA=0$.
A curl-free form on a torus is $d\lambda+a$ with $a$
constant: subtract its three period averages, and integrate
the remaining form along paths. Interchanging steps changes
the integral by a plaquette curl, and its periods now
vanish, so the potential is well defined. If also
$\mathsf LA=0$, then, for
$x=y+\ell_{j-1}\sum n_\mu e_\mu\in B(y)$,
\[
 \lambda(x)=\lambda(y)-\sum_\mu n_\mu a_\mu.
\]
The identities $\mathsf Qd=d\mathsf r$ and
$\mathsf Qa=La$ follow by telescoping the mean paths.
Thus $\mathsf QA=0$ implies $d\mathsf r\lambda+La=0$.
Summing around a coarse torus loop gives $a=0$, after which
$\mathsf r\lambda$, and then $\lambda$, is constant.

Use the decomposition
$\mathcal A_{j-1}=\ker\mathsf Q\oplus(\ker\mathsf Q)^\perp$.
The implicit function theorem parametrizes the fibre by
$V=\exp i(B_0+\beta(B_0,W))$, where $B_0\in\ker\mathsf Q$
and $\beta(0,\mathbf1)=0$, $\beta(B_0,\mathbf1)=O(|B_0|^2)$.
The Hessian in $B_0$ of $\mathcal F_j(V)$ at
$(0,\mathbf1)$ is $g|_{\ker\mathsf Q}$. A second application
of that theorem gives its unique analytic critical point
$B_0(W)$; positivity on a common convex coordinate ball
gives the strict minimum. Define $p_j$ by this point.
The gauge transformation equal to $p_j(W)(y,x)$ at
$x\in B(y)$ and to $1$ on $\Lambda_j$ makes it axial,
preserves the constraint and $\mathcal S_{j-1}$, and can
only decrease $\mathrm{GF}_j$. Both configurations lie in
the uniqueness neighbourhood for $W$ sufficiently close to
$\mathbf1$; hence the minimizer was already axial.

Coarse gauge transformations extended blockwise preserve the
fibre problem and its gauge-fixing function. The same holds
for the lattice isometries in the statement. Indeed, set
$(\mathsf sV)(\mathsf sb)=V(b)$. Such an isometry maps blocks,
straight paths and contour families bijectively to their
images. Path holonomies and contour means are therefore
unchanged after relabelling. For the matrices $W_x$ in the
averaging formula, with $x'=x+c_+-c_-$, reversal gives
\[
 W^{\bar c}_{x'}=V[c]^{-1}W_x^{-1}V[c],\qquad
 \log W^{\bar c}_{x'}=-V[c]^{-1}(\log W_x)V[c].
\]
Here all fields are in a sufficiently small neighbourhood
of $\mathbf1$, where both logarithms are given by their
series. Summing and exponentiating gives
$\mathcal M(V)(\bar c)=\mathcal M(V)(c)^{-1}$.
Thus $\mathcal M(\mathsf sV)=\mathsf s\mathcal M(V)$ even
when the isometry reverses a positive bond. The contour
means also show that $\mathrm{GF}_j$ is preserved, and the
plaquette trace shows that $A$ is preserved. Induction and
local uniqueness give the stated covariance and invariance.
Differentiating the constraint and criticality
at $\mathbf1$ shows that $p_j'(\mathbf1)B$ is the
$g$-orthogonal lift of $B$, giving the displayed minimum
formula for $h_j$. A null vector is therefore curl-free.
Conversely write a coarse curl-free form as $B=d\gamma+a'$.
Put $a=L^{-1}a'$ and
$\lambda(x)=\gamma(y)-\sum n_\mu a_\mu$ in each block.
Then $A=d\lambda+a$ satisfies $\mathsf QA=B$,
$\mathsf LA=0$ and $dA=0$, proving $h_j(B)=0$.
This closes the induction. The same uniqueness induction
also shows that the germs agree under truncation of the
tower and the stated dilation.

At $V=p_j(W)$, axiality removes the gradient of
$\mathrm{GF}_j$, so constrained criticality and surjectivity
give $\nabla\mathcal S_{j-1}(V)=\mathcal M'(V)^{\mathsf T}\mu$.
Differentiate $\mathcal M p_j=I$ and
$\mathcal S_j=\mathcal S_{j-1}p_j$ to obtain
$\mu=\nabla\mathcal S_j(W)$. Iterating the chain rule proves
\eqref{eq:gaussian-germ-multiplier}.

We record the two Gaussian calculations used in the product.
For $X|_{\ker M}=T$, an orthonormal decomposition
$E=\ker M\oplus(\ker M)^\perp$ gives
\[
 |\det\mathbb S(X,M)|=\det(T)\det(MM^{\mathsf T}),\qquad
 I(T,M)=(2\pi)^{(\dim E-\dim F)/2}
                             |\det\mathbb S(X,M)|^{-1/2}.
\]
These formulas show analyticity of the normalizers: the
restricted Hessian has an analytic extension to the ambient
space. Explicitly, if $q_V^1$ is the second derivative of
$i^{-1}\log(\mathcal M(e^{iB}V)W^{-1})$, that extension is
$\mathcal F_j''(V)-(\nabla\mathcal S_j(W),q_V^1)$.
To verify it, choose a second-order correction $C$ with
$\mathcal M'(V)C=-q_V^1(B,B)$ for a tangent $B$; the curve
$e^{i(tB+t^2C/2+O(t^3))}V$ lies in the fibre. Its second
derivative of $\mathcal F_j$ gives the asserted formula.
The bordered determinant is nonzero by positivity.

Next let $M:E\to F$ and $N:F\to F'$ be onto, and let $T$
be positive on $\ker NM$. For $\zeta\in\ker N$ let
$\sigma\zeta$ minimize $T$ over $M\xi=\zeta$.
Then $T(b+\sigma\zeta)=T_1(b)+T_2(\zeta)$, where
$T_1=T|_{\ker M}$. The Jacobian $J$ of this splitting
satisfies
\[
 J^2=\det((S^{-1})_{11}),\quad S=MM^{\mathsf T},\qquad
 \det(NSN^{\mathsf T})=\det(NN^{\mathsf T})\det S\,J^2.
\]
Here the block $11$ is on $\ker N$. The first identity
follows by replacing $\sigma\zeta$ by the orthogonal lift
$M^{\mathsf T}S^{-1}\zeta$, which changes the splitting by
a shear of determinant one; the second is the block inverse
formula on $F=\ker N\oplus(\ker N)^\perp$.
The Gaussian and coarea factors therefore give
\begin{equation}\label{eq:gaussian-two-stage}
 I(T,NM)=I(T_1,M)I(T_2,N).
\end{equation}

The Hessians of successive fibre minima have precisely this
Schur complement. In submersion coordinates $(u,v,w)$,
where $m(u,v,w)=(v,w)$ and $n(v,w)=w$, a critical section
satisfies $F_u=0$, so
$f_{vv}=F_{vv}-F_{vu}F_{uu}^{-1}F_{uv}$ for
$f(v,w)=F(u(v,w),v,w)$. Adding a function of $(v,w)$ adds
its $vv$ Hessian to this Schur complement. Positive fibre
Hessians at the two stages therefore give a positive total
Hessian and the minimum formula used in
\eqref{eq:gaussian-two-stage}.

Apply this calculation inductively to
$F_i=A+\sum_{j\le i}\mathrm{GF}_j\circ\mathcal M^{j-1}$,
with $m=\mathcal M^{i-1}$, $n=\mathcal M$ and the critical
section $c_{i-1}$. It gives a positive Hessian of $F_i$ on
the fibre $\mathcal M^i=W_i$ at $c_i(W_i)$, and its
normalizer is $\prod_{j\le i}Z_j(W_j)$.
Since $W_{j-1}$ is axial, each gauge-fixing summand has zero
gradient there and contributes the quadratic form
$|\mathsf L_{W_{j-1}}\mathcal Q_U^{[j-1]}B|^2$.
At $i=k$ the total Hessian is therefore
$H_W+|\mathbf L_WB|^2$, proving
\eqref{eq:gaussian-product}. Taking normalized logarithms
gives \eqref{eq:gaussian-recursion}. Gauge transformations
and isometries act orthogonally on these tangent spaces,
including the adjoint correction on reversed bonds; they
preserve both Gaussian and coarea factors. This proves the
remaining invariance assertions.
\end{proof}

\subsection{The determinant in a background gauge}

We specify the correction from the derivative of the mean
to the covariant straight average. At one level put
\begin{align*}
 (\mathsf n_VB)(y)&=L^{-3}\sum_{x\in B(y)}
  \frac1i\frac d{dt}\Big|_0
       \log[(e^{itB}V)(y,x)V(y,x)^{-1}],\\
 (\mathfrak a_V\lambda)(y)&=L^{-3}\sum_{x\in B(y)}
                \operatorname{Ad}(V(y,x))\lambda(x),
 &\mathring{\mathcal Q}_V&=\mathcal M'(V)
                              +\mathsf D_{\mathcal M(V)}\mathsf n_V.
\end{align*}
For $V_j=\mathcal M^j(U)$, use the products
$\mathfrak a_U^{[j]}=\mathfrak a_{V_{j-1}}\cdots \mathfrak a_{V_0}$ and
$\mathring{\mathcal Q}_U^{[j]}
=\mathring{\mathcal Q}_{V_{j-1}}\cdots\mathring{\mathcal Q}_{V_0}$,
and the recursion
\[
 \mathsf n_U^{[0]}=0,\qquad
 \mathsf n_U^{[j]}=\mathfrak a_{V_{j-1}}\mathsf n_U^{[j-1]}
                         +\mathsf n_{V_{j-1}}\mathcal Q_U^{[j-1]}.
\]
The restriction of site functions to $\Lambda_j$ is $r=\mathsf r$;
its transpose extends by zero. Define $q_U$ by
\[
 \frac1i\log[\mathcal M^k(e^{iB}U)\mathcal M^k(U)^{-1}]
       =\mathcal Q_U^{[k]}B+\tfrac12q_U(B,B)+O(|B|^3),
\]
and set
\[
 c_U(\lambda,B)(b)=\tfrac i2
   [\lambda(b_-)+\operatorname{Ad}(U(b))\lambda(b_+),B(b)],
 \qquad\sigma_B=r^{\mathsf T}\mathsf n_U^{[k]}B,
 \qquad\tau_UB=B+\mathsf D_U\sigma_B.
\]
At $U=\mathbf1$, abbreviate the scalar mean and correction
by $\mathfrak a^{[j]}$ and $\mathsf n^{[j]}$, and the straight
average by $\mathring{\mathsf Q}^{[j]}$. Omit $[1]$ for
one step.

For independent variables $(U,J,\mu)$ near $(\mathbf1,0,0)$,
define the symmetric operator $X_{U,J,\mu}$ by
\begin{align}
 (B,X_{U,J,\mu}B)
  ={}&A''(U)(B,B)+2(J,c_U(\sigma_B,B))
                 +(J,c_U(\sigma_B,\mathsf D_U\sigma_B))
                 -(\mu,q_U(\tau_UB,\tau_UB))\notag\\
    &+|(1-L^{3k}\mathfrak a_U^{[k]\mathsf T}\mathfrak a_U^{[k]})
                                      \mathsf D_U^{\mathsf T}B|^2.
 \label{eq:gaussian-background-X}
\end{align}
All pairings and transposes here are the unweighted ones
specified above. Put
\begin{equation}\label{eq:gaussian-background-functional}
 \mathfrak F_{\mathfrak s}(U,J,\mu)
 =-\tfrac12\log|\det\mathbb S(X_{U,J,\mu},
                                  \mathring{\mathcal Q}_U^{[k]})|
       +\log|\det\mathbb S(\mathsf D_U^{\mathsf T}\mathsf D_U,
                                                \mathfrak a_U^{[k]})|.
\end{equation}

\begin{proposition}\label{prop:gaussian-background-determinant}
\label{I-108-78}
For fixed $k\ge1$ and finite side $\mathfrak s$, both
matrices in \eqref{eq:gaussian-background-functional} are
invertible at $(\mathbf1,0,0)$. The function
$\mathfrak F_{\mathfrak s}$ is a real-analytic germ there
and is invariant under simultaneous gauge transformation
of $(U,J,\mu)$, for gauges near $\mathbf1$ or constant.
Its matrix entries are built by finite compositions and
derivatives of local analytic means and the plaquette
action, with the same local formulas on finite towers and
on the infinite tower. These formulas are covariant under
lattice isometries preserving $\Lambda_k$, with the tangent action on reversed
bonds. They admit complex analytic extensions near the
trivial triple.

For the one-step maps there is a common radius
$\rho_L\in(0,1]$, depending only on $L$: on the polydisc
$|\mathbf V(b)-1|<\rho_L$ for all bonds of a double block,
of both orientations, the complex means and contour
variables are holomorphic and within $1/2$ of $1$.
Their left derivatives, including the straightened
derivative kernels, are holomorphic on this same open
polydisc. In particular
$N(B)=i^{-1}\log\mathcal M(e^{iB})-\mathsf QB$
has $N(0)=N'(0)=0$.

For $W$ near $\mathbf1$,
\begin{equation}\label{eq:gaussian-background-representation}
 \Psi_k(W)=\mathfrak F_{\mathfrak s}
        (c_k(W),\nabla A(c_k(W)),\nabla\mathcal S_k(W))
                      -\mathfrak F_{\mathfrak s}(\mathbf1,0,0).
\end{equation}\label{I-108-80}
\end{proposition}

\begin{proof}
Differentiating gauge covariance of the mean and of the
contour variables gives the one-step identities
\[
 \mathcal M'(V)\mathsf D_V=\mathsf D_{\mathcal M(V)}r,
 \quad \mathsf n_V\mathsf D_V=\mathfrak a_V-r,
 \quad\mathring{\mathcal Q}_V\mathsf D_V
                         =\mathsf D_{\mathcal M(V)}\mathfrak a_V,
 \quad \mathfrak a_V\mathfrak a_V^{\mathsf T}=L^{-3}I.
\]
The last identity follows because blocks are disjoint and
adjoint transport preserves the bilinear Lie-algebra
pairing. Multiplication of these identities and the
recursion for $\mathsf n_U^{[j]}$ yield
\begin{equation}\label{eq:gaussian-straightening-identities}
 \mathring{\mathcal Q}_U^{[j]}
       =\mathcal Q_U^{[j]}+\mathsf D_{V_j}\mathsf n_U^{[j]},
 \quad \mathfrak a_U^{[j]}=r+\mathsf n_U^{[j]}\mathsf D_U,
 \quad \mathfrak a_U^{[j]}\mathfrak a_U^{[j]\mathsf T}=L^{-3j}I.
\end{equation}
In particular $L^{3k}\mathfrak a_U^{[k]\mathsf T}\mathfrak a_U^{[k]}$ is the
orthogonal projection onto $(\ker \mathfrak a_U^{[k]})^\perp$ for
real $U$. At $U=\mathbf1$, $\mathfrak a^{[k]}$ is the arithmetic
mean on a block of $L^{3k}$ sites and
$\mathring{\mathsf Q}^{[k]}$ is the corresponding mean
of straight path sums. Indeed the centred derivative is
$\mathsf Q=\mathring{\mathsf Q}-d\mathsf n$, and telescoping gives
$\mathring{\mathsf Q}d=d\mathfrak a$. These identities also show that
$\mathring{\mathsf Q}^{[k]}$ is onto, by the pivot lift at
each level.

The operators at the trivial triple are positive on their
constraint kernels. For the scalar operator,
$|d\lambda|^2=0$ makes $\lambda$ constant, and
$\mathfrak a^{[k]}\lambda=0$ then makes it zero. For the vector
operator write $P=L^{3k}\mathfrak a^{[k]\mathsf T}\mathfrak a^{[k]}$; its
quadratic form is
$|dB|^2+|(1-P)d^{\mathsf T}B|^2$.
If it vanishes and $\mathring{\mathsf Q}^{[k]}B=0$, write
$B=d\lambda+a_0$ with $a_0$ constant. The latter constraint
is $d\mathfrak a^{[k]}\lambda+L^ka_0=0$. Torus periods give $a_0=0$,
and $\mathfrak a^{[k]}\lambda$ is constant. Therefore
\[
 |d\lambda|^2=(\lambda,Pd^{\mathsf T}d\lambda)
 =L^{3k}(\mathfrak a^{[k]}\lambda,\mathfrak a^{[k]}d^{\mathsf T}d\lambda)=0,
\]
since the last scalar mean has zero total sum. Thus $B=0$.
The bordered determinant formula in the preceding proof
gives invertibility of both matrices.

Their entries are analytic. To check the local assertion
uniformly over translates, choose a polydisc about the
trivial configuration on one double block, small enough
that its complex means and contour variables are defined
and within $1/2$ of $1$. Translation, dilation and coordinate
permutation identify every double block with this one;
including both orientations makes the same statement valid
under reflections. The double blocks lift injectively from
the finite towers. Thus one radius depending only on $L$
serves every one-step map. Finite iteration gives the
locally defined entries above for fixed $k$, and the same
formulas on the infinite lattice. Derivatives of order at
most two still depend on only finitely many bond variables;
the resulting matrices have finite range. The determinant
and its chosen logarithm are analytic near the trivial
triple because their values there are nonzero.

Gauge transformations intertwine the means, their
derivatives, $\mathsf D_U$, $\mathsf n_U^{[k]}$, $\mathfrak a_U^{[k]}$ and
$\mathring{\mathcal Q}_U^{[k]}$. They are orthogonal for the
pairings, commute with $r^{\mathsf T}$ and intertwine
$q_U,c_U$ as well. They consequently conjugate both
bordered matrices. The same argument applies to lattice
isometries, using the tangent adjoint on any reversed bond.
It proves the asserted covariance and invariance.

We now relate the two determinants to the fibre integral.
The mixed derivative of gauge invariance, computed from
\[
 (e^{isB}U)^{e^{it\lambda}}
    =\exp i(sB-t\mathsf D_U\lambda+st\,c_U(\lambda,B)
                                  +O(t^2)+O((|s|+|t|)^3))U,
\]
gives
\begin{equation}\label{eq:gaussian-mixed-gauge-hessian}
 A''(U)(B,\mathsf D_U\lambda)
                      =(\nabla A(U),c_U(\lambda,B)).
\end{equation}
At $U=c_k(W)$, the second-order fibre calculation from the
preceding proof gives
\begin{equation}\label{eq:gaussian-fibre-hessian}
 H_W(B)=A''(U)(B,B)-(
              \nabla\mathcal S_k(W),q_U(B,B)),
                    \qquad B\in\ker\mathcal Q_U^{[k]}.
\end{equation}
Let $\mathfrak k=\{\lambda\in\mathfrak G_0:
\lambda|_{\Lambda_k}=0\}$. Gauging a curve in this fibre
by $e^{-it\kappa}$ preserves its action and changes its
tangent from $B$ to $B+\mathsf D_U\kappa$. Hence
$H_W(B+\mathsf D_U\kappa)=H_W(B)$.
Since every intermediate configuration is axial,
\[
 (\mathbf L_W\mathsf D_U\kappa)(x)=\kappa(x)-\kappa(y),
       \qquad x\in\Lambda_{j-1}\setminus\Lambda_j,
                       \quad x\in B(y).
\]
Ordering the levels in decreasing order makes this map
triangular with identity diagonal. In particular its
determinant on $\mathfrak k$ is one.

We use the following change-of-constraint calculation.
Suppose $D:\mathfrak G\to E$ and
$D_c:\mathfrak G_c\to F$ are injective, $Q:E\to F$ and
$r:\mathfrak G\to\mathfrak G_c$ are onto, and
$QD=D_cr$. Let $\mathcal E=Q^{-1}(D_c\mathfrak G_c)$,
and let $H\ge0$ on $\mathcal E$ have null space
$D\mathfrak G$. Given $n:E\to\mathfrak G_c$, set
$Q'=Q+D_cn$ and $r'=r+nD$. Suppose $r'$ is onto, and
put $\mathfrak k'=\ker r'$. Let
$L':\ker Q'\to K_g$, where $K_g$ is a Euclidean space
of dimension $\dim\mathfrak k'$. If $L'D|_{\mathfrak k'}$
is invertible, then
\begin{equation}\label{eq:gaussian-change-gauge-invariant}
 I(H+|L'\,\cdot\,|^2,Q')
       |\det(L'D|_{\mathfrak k'})|
                   \det(r'r'^{\mathsf T})^{1/2}
\end{equation}
is independent of $n,L'$.
Here is the determinant calculation. The projection
$P_\perp:F\to(D_c\mathfrak G_c)^\perp$ satisfies
$P_\perp Q'=P_\perp Q$; since the range of $Q'$ contains
$D_c\mathfrak G_c$, $Q'$ is onto and
$\mathcal E=Q'^{-1}(D_c\mathfrak G_c)$.
Put $v'=D|_{\mathfrak k'}$,
$K'=\ker Q'\ominus D\mathfrak k'$ and
$\mathcal E_\perp=\mathcal E\ominus D\mathfrak G$.
Projection $P':K'\to\mathcal E_\perp$ is an isomorphism
preserving $H$. Integrating first in the variable $\kappa$
in $B=a+v'\kappa$ removes $L'$ and contributes
\[
 (2\pi)^{\dim\mathfrak k'/2}|\det L'v'|^{-1}
       \det(v'^{\mathsf T}v')^{1/2}|\det P'|^{-1}
                \int_{\mathcal E_\perp}e^{-H/2}.
\]
To compute the remaining Jacobian put
$s'=r'^{\mathsf T}(r'r'^{\mathsf T})^{-1}$ and
$N'=Q'Q'^{\mathsf T}$. The parametrization
$\Theta(a,\kappa,\gamma)=a+D\kappa+Ds'\gamma$ of
$\mathcal E$ has, in its two orthogonal decompositions,
the determinants
\begin{align*}
 |\det\Theta|
  &=|\det P'|\det(D^{\mathsf T}D)^{1/2}
                         \det(r'r'^{\mathsf T})^{-1/2}\\
  &=\det(v'^{\mathsf T}v')^{1/2}
                         \det(D_c^{\mathsf T}N'^{-1}D_c)^{1/2}.
\end{align*}
The first follows by splitting off $D\mathfrak G$; for the
second the final orthogonal component is
$Q'^{\mathsf T}N'^{-1}D_c\gamma$.
The block inverse formula gives
\[
 \det N'\det(D_c^{\mathsf T}N'^{-1}D_c)
  =\det(D_c^{\mathsf T}D_c)
       \det(P_\perp QQ^{\mathsf T}P_\perp),
\]
where the last determinant is on the range of $P_\perp$.
Substitution in \eqref{eq:gaussian-change-gauge-invariant}
leaves only these fixed determinants, $\det(D^{\mathsf T}D)$
and the integral on $\mathcal E_\perp$; the dimension
$\dim\mathfrak k'=\dim\mathfrak G-\dim\mathfrak G_c$ is
fixed as well. This proves the claimed independence.

First assume $\mathsf D_W$ injective. Then $\mathsf D_U$
is injective: $\mathcal Q_U^{[k]}\mathsf D_U=\mathsf D_Wr$
implies $r\lambda=0$ for a null vector, and the triangular
map $\mathbf L_W\mathsf D_U|_{\mathfrak k}$ then gives
$\lambda=0$. On
$\mathcal E=\ker\mathcal Q_U^{[k]}+\mathsf D_U\mathfrak G_0$
extend $H_W$ by $H(B+\mathsf D_U\lambda)=H_W(B)$.
Its gauge invariance makes this well defined. Removing
$\mathbf L_WB$ by the invertible triangular gauge map
and using the positive total Hessian shows that $H\ge0$
with null space exactly $\mathsf D_U\mathfrak G_0$.
The change-of-constraint calculation applies.

For its first choice take $n=0,L'=\mathbf L_W$. Since
$rr^{\mathsf T}=I$ and the gauge determinant is one,
\eqref{eq:gaussian-change-gauge-invariant} equals
$\mathcal Z_k(W)$. For the second take $n=\mathsf n_U^{[k]}$,
$r'=\mathfrak a_U^{[k]}$, $Q'=\mathring{\mathcal Q}_U^{[k]}$ and
$L'=(1-L^{3k}r'^{\mathsf T}r')\mathsf D_U^{\mathsf T}$.
Then $r'r'^{\mathsf T}=L^{-3k}I$ and
$L'\mathsf D_U|_{\ker r'}=v'^{\mathsf T}v'$ is positive.
For $B\in\ker Q'$, $\tau_UB\in\ker\mathcal Q_U^{[k]}$,
so $H(B)=H_W(\tau_UB)$. Expanding this expression with
\eqref{eq:gaussian-mixed-gauge-hessian} and
\eqref{eq:gaussian-fibre-hessian} identifies
$H(B)+|L'B|^2$ with \eqref{eq:gaussian-background-X} at
$J=\nabla A(U)$, $\mu=\nabla\mathcal S_k(W)$.

Apply the bordered determinant formula to this vector
Gaussian and to the scalar pair $(\mathsf D_U^{\mathsf T}
\mathsf D_U,r')$. The second application gives
\[
 \det(v'^{\mathsf T}v')
   =|\det\mathbb S(\mathsf D_U^{\mathsf T}\mathsf D_U,r')|
                         \det(r'r'^{\mathsf T})^{-1}.
\]
Comparison of the two choices therefore yields
\begin{align*}
 \log\mathcal Z_k(W)
  -\mathfrak F_{\mathfrak s}
            (c_k(W),\nabla A(c_k(W)),\nabla\mathcal S_k(W))
  =\tfrac12(\dim\mathcal A_0-\dim\mathcal A_k)\log(2\pi)
                         +\tfrac32k\dim\mathfrak G_k\log L.
\end{align*}
Both sides extend continuously near $\mathbf1$. Configurations
with $\mathsf D_W$ injective are dense there. Indeed the
analytic determinant $\det(\mathsf D_W^{\mathsf T}\mathsf D_W)$
is not identically zero: put $W=e^{i\epsilon t_1}$ in
direction $1$, $W=e^{i\epsilon t_2}$ in direction $2$ and
$W=1$ in direction $3$, where $0<\mathfrak s\epsilon<\pi$.
A covariantly constant site function is fixed by the two
loop holonomies $e^{i\mathfrak s\epsilon t_1}$ and
$e^{i\mathfrak s\epsilon t_2}$. Their adjoints are rotations
through angle $2\mathfrak s\epsilon$ about the two distinct
Pauli axes, whose fixed subspaces intersect only at zero.
This proves injectivity for arbitrarily small $\epsilon$.
The identity consequently holds also at $W=\mathbf1$.
Subtracting that value removes its constant right-hand side
and proves \eqref{eq:gaussian-background-representation}.
\end{proof}

\section{The accumulated Gaussian coefficient}\label{sec:gaussian-coefficient}

\subsection{Second derivatives of the Gaussian determinant}

Keep the tower, means and Gaussian normalizers of
\cref{sec:gaussian-normalizers}. The infinite tower uses the
same lattices without the periodic quotient. All complex
means are their holomorphic branches near the trivial
configuration. The neighbourhoods may depend on height and
volume.

All Brillouin-zone integrals use normalized measure:
\[
 \int f(\kappa)\,d\kappa
       =(2\pi)^{-3}\int_{[-\pi,\pi]^3}f(\kappa)\,d^3\kappa,
\]
with the same convention for $p$ and $q$.
We use Bloch symbols with respect to translations by
\(\mathbb Z^3\). Each fine lattice has finitely many
indices per unit cell, carrying the pairing weights of the
operator's domain and target. The symbol is characterized by
\[
 T(e^{i\kappa\cdot x}v)=e^{i\kappa\cdot x}\widehat T(\kappa)v,
 \qquad v(x+n)=v(x).
\]
For finite-range kernels this is a finite sum of
exponentials. A holomorphic symbol on a strip has an
exponentially decreasing kernel, obtained by Fourier
inversion and a shift of the integration contour.
Products, weighted transposes and periodisation commute
with this correspondence. Indeed, the product rule is
the absolutely convergent convolution of the kernels;
periodising this convolution proves the torus rule.

We need only second derivatives of a local operator
function \(T(\zeta)\). Along
\(\zeta_\pm=e^{\pm ip\cdot x}\varepsilon_\pm\), record them
as the quadruple
\(\mathbf T=(T_{00},T_{10},T_{01},T_{11})\): the first
component is \(\widehat{T(0)}\), the next two carry
momentum \(\kappa\) to \(\kappa\pm p\), and the last
is the mixed derivative, which preserves momentum.
If \(S_r=\{\kappa:|\operatorname{Im}\kappa_\mu|\le r\}\),
give these four components the sum of their supremum
operator norms on, respectively,
\[
 S_{2\delta},\quad
 \{\kappa:\kappa,\kappa+p\in S_{2\delta}\},\quad
 \{\kappa:\kappa,\kappa-p\in S_{2\delta}\},\quad
 \mathbb R^3 ,
 \qquad p\in S_\delta.
\]
We call this quadruple a jet. Products are compositions
with the indicated momentum changes. Equivalently they
are products of
\[
 T_{00}+\epsilon_+T_{10}+\epsilon_-T_{01}
                    +\epsilon_+\epsilon_-T_{11},
 \qquad \epsilon_+^2=\epsilon_-^2=0.
\]
The resulting norm is submultiplicative, and transpose
preserves it. Three jets with zero \(00\)-component
have zero product. Put $\mathbf T_0=(T_{00},0,0,0)$.
Thus a jet is invertible whenever
\(T_{00}^{-1}\) is bounded on \(S_{2\delta}\), and
\[
 \mathbf T^{-1}
       =(1-\mathbf N+\mathbf N^2)\mathbf T_0^{-1},
 \quad \mathbf N=\mathbf T_0^{-1}\mathbf T-1,\quad
 \|\mathbf T^{-1}\|\le \iota(1+\iota\|\mathbf T\|)^2,
 \quad \iota=\sup\|T_{00}^{-1}\|.
\]
Set
\[
 \mathsf t(\mathbf T)=\int\operatorname{tr}T_{11}(\kappa)\,d\kappa,
 \qquad
 \ell(\mathbf T)=\mathsf t(\mathbf N-\tfrac12\mathbf N^2).
\]
Contour shifts in the periodic integrals give trace
cyclicity. Expanding the two nilpotent variables then yields
\[
\begin{split}
 \ell(\mathbf T)
 =\int\operatorname{tr}\big[
   T_{00}(\kappa)^{-1}T_{11}(\kappa)
   -T_{00}(\kappa)^{-1}T_{10}(\kappa-p)
    T_{00}(\kappa-p)^{-1}T_{01}(\kappa)\big]\,d\kappa,\\
 \ell(\mathbf T\mathbf S)=\ell(\mathbf T)+\ell(\mathbf S),
 \qquad
 \frac d{dt}\ell(\mathbf T_t)
                =\mathsf t(\mathbf T_t^{-1}\dot{\mathbf T}_t).
\end{split}
\]
The last formula holds for differentiable families with
uniformly bounded zero-component inverses. To see
additivity directly, conjugate \(\mathbf N_{\mathbf T}\)
by \(\mathbf S_0\), write
\(\mathbf N_{\mathbf T\mathbf S}=a+b+ab\), and use
the vanishing of cubic products and
\(\mathsf t(ab)=\mathsf t(ba)\).
The derivative follows by applying additivity to
\(\mathbf T_t^{-1}\mathbf T_{t+h}
=1+h\mathbf T_t^{-1}\dot{\mathbf T}_t+o(h)\).
Block triangular jets contribute the sum of their
diagonal blocks. Field-independent invertible diagonal
factors, including changes of the pairing weights,
contribute zero.

Suppose \(T\) is a local analytic operator function
and \(\widehat{T(0)}\) is invertible for real \(\kappa\).
Compactness of one period cube gives an invertible
strip. On a sufficiently large torus,
\[
 D^2\log\det T(0)[\zeta_+,\zeta_-]
 =\operatorname{Tr}(T_0^{-1}T^{(2)})
  -\operatorname{Tr}(T_0^{-1}T^{(1)}_+
                              T_0^{-1}T^{(1)}_-).
\]
Fourier decomposition makes the right side a sum of
the integrand defining \(\ell(\mathbf T)\). Its
normalisation by the torus volume converges to
\(\ell(\mathbf T)\) with error
\(Ce^{-c\mathfrak s}|\varepsilon_+||\varepsilon_-|\),
uniformly for real torus momenta \(p\).
Indeed, the integrand is periodic and bounded on a
strip; its Fourier coefficients decrease exponentially,
and the torus Riemann sum retains precisely the
coefficients indexed by \(\mathfrak s\mathbb Z^3\).
Different external momenta give zero trace unless their
sum is in \(2\pi\mathbb Z^3\). For real kernels the same
formula applies to \(\log|\det T|\).

Write $\Pi^F_{\mathfrak s}=F_{\mathfrak s}''(\mathbf1)$,
extended bilinearly to complex forms. For a positive bond
$b$, let $\delta_b^a$ equal $t_a$ on $b$ and zero elsewhere.
For $p\in(2\pi/\mathfrak s)\mathbb Z^3$ modulo
$2\pi\mathbb Z^3$, let $e^a_{p,\mu}(b)$ equal
$e^{ip\cdot b_-}t_a$ on bonds of direction $\mu$ and zero
on the others. An infinite-lattice bond in such a formula
is represented by its image on the torus.

We call a family \(F_{\mathfrak s}\) tempered when it
is real analytic near \(\mathbf1\), vanishes there,
is invariant under gauge transformations near
\(\mathbf1\), constant gauges, and all lattice
isometries, and its volume-normalised plane-wave
polarisation converges uniformly to
\(\delta_{aa'}\widehat\pi_{\mu\nu}(p)\), where
\(\widehat\pi\) is bounded and holomorphic on a strip.
The neighbourhood and strip may depend on the family.
Its coefficient is
\[
 \mathsf b[F]=-\partial_\mu\partial_\nu
                      \widehat\pi_{\mu\nu}(0),\qquad \mu\ne\nu.
\]
We verify that this definition agrees with the
kernel coefficient and is independent of the chosen pair.

Constant gauge invariance first gives
\(\nabla F_{\mathfrak s}(\mathbf1)=0\), since the adjoint
representation has no nonzero fixed vector.
Differentiating gauge invariance gives
\(\Pi^F[d\lambda,B]=0\). Translation invariance removes
unequal momenta, and Fourier inversion on the tori gives
\[
 \Pi^F_{\mathfrak s}[\delta^a_x,\delta^{a'}_y]
 \longrightarrow\delta_{aa'}\pi_{\mu\nu}(y_--x_-),
 \qquad
 \pi_{\mu\nu}(u)=\int\widehat\pi_{\mu\nu}(p)e^{ip\cdot u}\,dp.
\]
The error in this inversion is bounded by the uniform
plane-wave error. The strip gives exponential decay of
\(\pi\), and its reality follows from the finite-torus
polarisations. Gauge invariance passes to the limit as
\[
 \sum_\mu(e^{ip_\mu}-1)\widehat\pi_{\mu\nu}(p)=0.
\]
In the bond-midpoint convention
\(\widetilde\pi_{\mu\nu}
=e^{i(p_\mu-p_\nu)/2}\widehat\pi_{\mu\nu}\),
the lattice isometries give
\(\widetilde\pi(Rp)=R\widetilde\pi(p)R^{\mathsf T}\)
for all signed permutation matrices \(R\).
In particular \(\widetilde\pi\) is even.
Its constant term is \(cI\); the gauge relation
\(\sum_\mu\sin(p_\mu/2)\widetilde\pi_{\mu\nu}=0\)
forces \(c=0\). Reflection and permutation symmetry
leave for its quadratic term only
\[
 \alpha\delta_{\mu\nu}|p|^2+\beta p_\mu p_\nu
                           +\gamma\delta_{\mu\nu}p_\mu^2.
\]
The cubic term of the gauge relation forces
\(\beta=-\alpha\) and \(\gamma=0\). Hence
\[
 \widetilde\pi_{\mu\nu}(p)
  =\mathsf b[F](|p|^2\delta_{\mu\nu}-p_\mu p_\nu)
                                               +O(|p|^4).
\]
Termwise differentiation of its exponentially
convergent Fourier series now gives
\[
 \mathsf b[F]=\sum_u\pi_{\mu\nu}(u)u_\mu u_\nu
           =-\tfrac12\sum_u\pi_{\nu\nu}(u)u_\mu^2
                      \quad(\mu\ne\nu).
\]
The coefficient is real, additive, and independent
of the pair of distinct directions.

Write $\Phi^\natural_{l,\mathfrak s}$ for the normalizer
above with $j=k=l$ and side $\mathfrak s$.

\begin{theorem}\label{I-109-104}
Fix \(G=\mathrm{SU}(2)\), dimension three, and odd
\(L\ge13\), with the tower means and Gaussian
normalizers specified above. For each \(l\ge1\),
the one-step Gaussian normalizer
\(\Phi_l^\natural=(\Phi^\natural_{l,\mathfrak s})_{\mathfrak s}\)
is tempered. Its members are real analytic near
\(\mathbf1\), and their second differentials on
individual bonds have an infinite-volume limit
\[
 \lim_{\mathfrak s\to\infty}
 \Pi^{\Phi_l^\natural}_{\mathfrak s}
       [\delta^a_{\hat x},\delta^{a'}_{\hat y}]
 =\delta_{aa'}\pi^{(l)}_{\mu\nu}(\hat y_--\hat x_-),
\]
where \(\mu,\nu\) are the bond directions and
\(\pi^{(l)}\) is real and exponentially decreasing.
The coefficient
\[
 b_l^0=-\tfrac12\sum_{u\in\mathbb Z^3}
                                  u_\mu^2\pi^{(l)}_{\nu\nu}(u),
 \qquad \mu\ne\nu,
\]
is absolutely convergent, independent of the two
distinct directions, and depends only on \(l,L\).
There is a finite constant \(c_2(L)\) such that
\[
 \mathsf c_k:=\sum_{l=1}^kL^{-(k-l)}b_l^0
        =\mathsf b[\Psi_k],\qquad |\mathsf c_k|\le c_2(L)
                              \quad(k\ge0).
\]
One may take
\[
 c_2(L)=\sup_{k\ge0}|\mathsf b[\Psi_k]|
       \le C_{\mathrm{rem}}(L)\delta_{\mathrm{loop}}^{-2}
                                                       +C_{\mathrm{uv}}.
\]
The constants on the right are those in
\cref{I-109-38,I-109-75,I-109-99}.
\end{theorem}

For complex configurations the plaquette action means its
holomorphic extension
$A(U)=\sum_p[1-\tfrac12\operatorname{tr}U_p]$.
It agrees with the real action because an $\mathrm{SU}(2)$
matrix has real trace. On the infinite lattice its
derivatives on finitely supported forms use only the
finitely many plaquettes meeting their supports.

For $U$ whose intermediate means lie in the one-step
polydisc, set
\[
\begin{gathered}
 D_U=\eta^{-1}\mathsf D_U,\qquad
 Q_U=\eta\mathring{\mathcal Q}^{[k]}_U,\qquad
 R_U=\mathfrak a_U^{[k]},\qquad H_U=\eta^{-2}A''(U),\\
 K_U=H_U+D_UD_U^{\mathsf T},\qquad
 \Delta_U=D_U^{\mathsf T}D_U,\\
 X_U=H_U+D_U(1-R_U^{\mathsf T}R_U)D_U^{\mathsf T},
 \quad\Xi_U=X_U+Q_U^{\mathsf T}Q_U,
 \quad Y_U=\Delta_U+R_U^{\mathsf T}R_U.
\end{gathered}
\]
The transposes $\mathsf T$ in this display use the physical
weights $\eta^3$ at the fine level and $1$ at the coarse
level. Thus $(b,H_Ub)_\eta=
\eta^{-1}A''(U)(\eta b,\eta b)$.
If $t$ denotes an unweighted transpose, then
\[
 D_U^{\mathsf T}=\eta^{-1}\mathsf D_U^t,\qquad
 R_U^{\mathsf T}=\eta^{-3}\mathfrak a_U^{[k]t},\qquad
 Q_U^{\mathsf T}=\eta^{-2}\mathring{\mathcal Q}^{[k]t}_U.
\]
In particular $R_UR_U^{\mathsf T}=I$. The symmetric
idempotent $P=R_U^{\mathsf T}R_U$ gives
$(1-P)^t(1-P)=1-P$. Substitution in the corrected Hessian
at $J=\mu=0$ therefore gives
$X_{U,0,0}=\eta^2X_U$. Consequently, with unweighted
transposes on the left and physical transposes on the right,
\[
\begin{split}
 \mathbb S(X_{U,0,0},\mathring{\mathcal Q}^{[k]}_U)
 &=\begin{pmatrix}\eta^2&0\\0&\eta^{-1}\end{pmatrix}
                                  \mathbb S(X_U,Q_U),\\
 \mathbb S(\mathsf D_U^t\mathsf D_U,\mathfrak a_U^{[k]})
 &=\begin{pmatrix}\eta^2&0\\0&1\end{pmatrix}
       \mathbb S(\Delta_U,R_U)
       \begin{pmatrix}1&0\\0&\eta\end{pmatrix}.
\end{split}
\]
These diagonal factors are independent of the field.
They change logarithmic determinants only by constants
and hence leave their polarisation jets unchanged.

\subsection{A uniform strip for the free response}

Let \(\eta=L^{-k}\), \(\ell_j=L^{j-k}\), and use
pairing weight \(\ell_j^3\) at level \(j\).
In this subsection \(d\) is the exterior derivative
divided by the lattice spacing. The straight averages
on potentials and sites are denoted by \(Q\) and \(R\):
\[
 Q=\eta\,\mathring{\mathsf Q}^{[k]},\qquad
 R=\mathfrak a^{[k]},\qquad
 Qd=dR,\qquad RR^{\mathsf T}=1.
\]
Thus \(R\) is the mean over a block of side one, and
\(Q\) is the mean over the pairs consisting of a site
in the initial block and a bond on its straight path
to the translated site in the final block.
Put
\[
\begin{gathered}
 K=d^{\mathsf T}d+dd^{\mathsf T},\qquad \Delta=d^{\mathsf T}d,\\
 \Xi=K-dR^{\mathsf T}Rd^{\mathsf T}+Q^{\mathsf T}Q,\qquad
 Y=\Delta+R^{\mathsf T}R,\\
 G=(K+1)^{-1},\qquad g=(\Delta+1)^{-1}.
\end{gathered}
\]
Here \(K,G\) act on forms and \(\Delta,Y,g\) on sites.

For \(m\in\mathsf m=\{m:|m_\mu|\le(\eta^{-1}-1)/2\}\)
and \(q_m=\kappa+2\pi m\), set
\[
\begin{gathered}
 d_\mu(q)=\eta^{-1}(e^{i\eta q_\mu}-1),\qquad
 D(q)=\sum_\mu2\eta^{-2}(1-\cos\eta q_\mu),\\
 s(\tau)=\eta\sum_{|j|\le(\eta^{-1}-1)/2}e^{i\eta j\tau},
 \qquad v(q)=\prod_\nu s(q_\nu),\qquad
 u_\mu(q)=e^{i(1-\eta)q_\mu/2}s(q_\mu)v(q).
\end{gathered}
\]
The functions \(e^{2\pi im\cdot x}\) form an orthonormal
basis for a unit cell of the fine lattice. In this basis
\[
 \widehat R\lambda=\sum_mv(q_m)\widehat\lambda_m,\qquad
 (\widehat Qb)_\mu=\sum_mu_\mu(q_m)\widehat b_{\mu,m},
\]
and the transpose coefficients are \(v(-q_m)\) and
\(u_\mu(-q_m)\). The symbols of \(K,\Delta\) multiply
each coefficient by \(D(q_m)\). These formulas follow
by summing the block and straight-path geometric series;
the same sums give
\[
 u_\mu(q)d_\mu(q)=(e^{iq_\mu}-1)v(q).
\]
The iterated block formulas follow by the unique
decomposition of a centered \(L^n\)-block into centered
\(L^{n-1}\)-blocks. They also give, in absorbed lengths,
\[
 \mathsf Q^{[n]}
   =\mathring{\mathsf Q}^{[n]}-d\mathsf n^{[n]},\quad
 \|\mathring{\mathsf Q}^{[n]}\|_{\infty\to\infty}\le L^n,
 \quad \|\mathsf n^{[n]}\|_{\infty\to\infty}\le\tfrac32L^n,
 \quad \|\mathsf Q^{[n]}\|_{\infty\to\infty}\le4L^n.
\]
The correction \(\mathsf n^{[n]}\) is the sum of the
successive contour averages followed by restriction.
The identity follows by induction from
\(\mathsf Q=\mathring{\mathsf Q}-d\mathsf n\),
\(\mathsf n d=\mathfrak a-\mathsf r\).
The pivot lift in the proof of
\cref{thm:gaussian-critical-germs} is a right inverse for both
one-step means, since its contour correction vanishes.
Composing these lifts proves surjectivity of their iterates.
The unweighted transpose of restriction
is extension by zero, so
\(\mathsf r\mathsf r^{\mathsf T}=1\).

\begin{proposition}\label{I-109-36}\label{I-109-38}
There are absolute constants \(\delta_{\mathrm{loop}}>0\)
and \(C_{\mathrm{loop}}\), independent of \(k\) and \(L\),
such that the following symbols are bounded by
\(C_{\mathrm{loop}}\) on \(S_{2\delta_{\mathrm{loop}}}\):
the square-root propagators, their compositions with one
\(d\) or \(d^{\mathsf T}\), and \(Q,R\) and their transposes.
The four symbols
\[
\begin{gathered}
 G^{1/2}(t\Xi+(1-t)K+1)G^{1/2},\quad
 G^{1/2}(\Xi+\tau)G^{1/2},\\
 g^{1/2}(tY+(1-t)\Delta+1)g^{1/2},\quad
 g^{1/2}(Y+\tau)g^{1/2},
 \qquad 0\le t,\tau\le1,
\end{gathered}
\]
and the coarse symbols
\[
 Q(\Xi+\tau)^{-1}Q^{\mathsf T},\qquad
 R(Y+\tau)^{-1}R^{\mathsf T}
\]
are invertible there, with inverse norms at most
\(C_{\mathrm{loop}}\).
\end{proposition}

\begin{proof}
The essential point is coercivity at real momentum.
We give that argument before passing to a strip.
Cauchy--Schwarz in the averaging formulas gives
\(\|Q\|,\|R\|\le1\). A fine bond occurs in exactly
\(\eta^{-1}\) of the \(\eta^{-4}\) pairs for the
straight average, which cancels the physical pairing
weights. Site blocks are disjoint. With exponential
weights the same count gives
\[
 \|e^{\theta\cdot x}Qe^{-\theta\cdot x}\|
       \le e^{2|\theta|_1},\qquad
 \|e^{\theta\cdot x}Re^{-\theta\cdot x}\|
       \le e^{|\theta|_1}.
\]
In particular, for real \(\kappa\),
\(\sum_m|u_\mu(q_m)|^2\le1\) and \(\sum_mv(q_m)^2=1\).

Take \(\kappa\in[-\pi,\pi]^3\setminus\{0\}\), and write
\(D_m=D(q_m)\), \(v_m=v(q_m)\), \(u_{\mu,m}=u_\mu(q_m)\).
The sine formula for \(s\) gives
\[
\begin{gathered}
 D_m\ge4\quad(m\ne0),\qquad
 4|\kappa|^2/\pi^2\le D_0\le|\kappa|^2\le3\pi^2,\\
 v_0\ge(2/\pi)^3,\qquad |u_{\mu,0}|\ge(2/\pi)^4,\qquad
 1-v_0^2\le\pi^2D_0/48 .
\end{gathered}
\]
Consequently
\(QGQ^*\ge c_CI\) and \(RgR^*\ge c_CI\), with
\(c_C=(2/\pi)^8/(3\pi^2+1)\), by retaining the \(m=0\)
term. For site functions,
\[
 |\sum_mv_m\lambda_m|^2
   \ge\tfrac12v_0^2|\lambda_0|^2
                         -\tfrac14\sum_{m\ne0}D_m|\lambda_m|^2,
\]
so \(Y\ge c_Y(\Delta+1)\), for an absolute \(c_Y>0\).

For forms in \(\ker Q\), decompose
\[
 b_m=b_m^\perp+d(q_m)\lambda_m,\qquad
 b_m^\perp\perp d(q_m)\mathfrak g^c.
\]
Set
\[
 \chi_m=D_m\lambda_m,\quad
 \gamma=\sum_mv_m\chi_m,\quad
 \rho_m=\chi_m-v_m\gamma,\quad
 \beta=\sum_mv_m\lambda_m .
\]
Then \(\sum_mv_m\rho_m=0\), and
\[
\begin{split}
 q(b)&:=\|db\|^2+\|(1-R^{\mathsf T}R)d^{\mathsf T}b\|^2
       =\sum_mD_m|b_m^\perp|^2+\|\rho\|^2,\\
 \|b\|_1^2&:=(b,(K+1)b)
       =\sum_m(D_m+1)|b_m^\perp|^2
          +|\gamma|^2+\|\rho\|^2+\sum_mD_m^{-1}|\chi_m|^2 .
\end{split}
\]
The constraint reads
\[
 \sum_mu_{\mu,m}b_{\mu,m}^\perp+d_\mu^{\mathrm c}(\kappa)\beta=0 .
\]
Since \(d_\mu^{\mathrm c}=u_{\mu,0}d_\mu(\kappa)/v_0\),
the distance of \((u_{\mu,0}b_{\mu,0}^\perp)_\mu\) from
the coarse gradient space is at least
\((2/\pi)^4|b_0^\perp|\). The constraint bounds it by
\((\sum_{m\ne0}|b_m^\perp|^2)^{1/2}\). Thus
\[
 \|b^\perp\|^2\le C\sum_mD_m|b_m^\perp|^2,\qquad
 |\beta|\le C D_0^{-1/2}\|b^\perp\|.
\]
Moreover, with \(S=\sum_mv_m^2D_m^{-1}\ge v_0^2D_0^{-1}\),
\[
 \beta=S\gamma+
       \sum_{m\ne0}v_m\rho_m(D_m^{-1}-D_0^{-1}).
\]
The last sum is bounded by \(CD_0^{-1/2}\|\rho\|\),
using \(1-v_0^2\le CD_0\). Hence
\[
 |\gamma|\le CD_0^{1/2}(\|b^\perp\|+\|\rho\|),\qquad
 |\rho_0|\le CD_0^{1/2}\|\rho\|.
\]
These bounds control \(D_0^{-1}|\chi_0|^2\);
the remaining inverse weights are at most \(1/4\).
Substitution in the formula for \(\|b\|_1^2\) proves
\(\|b\|_1^2\le Cq(b)\) on \(\ker Q\).

For general \(b\), split off
\[
 b_0=GQ^*(QGQ^*)^{-1}Qb .
\]
Its energy is at most \(c_C^{-1}|Qb|^2\), and
\(b-b_0\in\ker Q\). The triangle inequality for
\(q^{1/2}\), together with \(q\le\|\cdot\|_1^2\),
therefore gives
\[
 c_\Xi(K+1)\le\Xi\le K+1,\qquad
 c_Y(\Delta+1)\le Y\le\Delta+1
\]
with absolute positive constants. Continuity includes
\(\kappa=0\).

For \(q=q'+i\theta\), \(|q'_\mu|\le\pi/\eta\),
\(|\theta_\mu|\le1/8\), the cosine formula gives
\[
 \operatorname{Re}(D(q)+1)\ge\tfrac9{10}(D(q')+1),\qquad
 |d_\mu(\pm q)|\le\tfrac65(|d_\mu(q')|+1/8).
\]
Thus the propagator square roots and their first
derivatives are uniformly bounded on \(S_{1/8}\).
The weighted averaging estimates give the same conclusion
for \(Q,R\). Each sandwiched operator in the statement
is a fixed polynomial in these bounded symbols, so it
has an absolute bound on \(S_{1/8}\).
Cauchy's estimate bounds its first momentum derivatives
on \(S_{1/16}\), uniformly in \(k,t,\tau\).
The real coercivity just proved and the Neumann series
therefore give its inverse on a smaller absolute strip.
The coarse Schur symbols are bounded there by their
factorisation through \(G^{1/2}\) or \(g^{1/2}\).
On the real axis they are bounded below by \(c_C/2\),
since \(\Xi+\tau\le2(K+1)\) and
\(Y+\tau\le2(\Delta+1)\).
A second Cauchy estimate and Neumann-series argument
give their inverses on a smaller strip. Fix such a
\(\delta_{\mathrm{loop}}\le1/64\) and a common
\(C_{\mathrm{loop}}\). All choices were absolute.
\end{proof}

Define
\[
 C=Q\Xi^{-1}Q^{\mathsf T},\qquad
 h=\Xi^{-1}Q^{\mathsf T}C^{-1},\qquad \nu=C^{-1}-1 .
\]
These are periodic operators with real kernels and
symbols holomorphic on a neighbourhood of the closed
strip, by compactness at each fixed \(k\).
They obey
\[
 Qh=1,\quad d^{\mathsf T}dh=Q^{\mathsf T}\nu,\quad
 (1-R^{\mathsf T}R)d^{\mathsf T}h=0,\quad d^{\mathsf T}\nu=0.
\]
Indeed, \(Xh=Q^{\mathsf T}\nu\), where
\(X=K-dR^{\mathsf T}Rd^{\mathsf T}\).
Applying \(d^{\mathsf T}\), and taking the real-momentum
Hermitian pairing with
\(\rho=(1-R^{\mathsf T}R)d^{\mathsf T}hB\), gives
\((\rho,\Delta\rho)=0\). Thus \(\rho\) is constant
at a zero momentum, and \(R\rho=0\) removes that
constant as well. Applying \(R\) then gives
\(d^{\mathsf T}\nu=0\).

For \(b=\widehat h(p)\varepsilon\),
\(n=\widehat\nu(p)\varepsilon\), and
\(\gamma=\widehat R\,\widehat{d^{\mathsf T}}b\), these
relations imply
\[
 D(q_m)b_{\mu,m}
   =u_\mu(-q_m)n_\mu+d_\mu(q_m)v(-q_m)\gamma .
\]
The sandwiched factorisation of \(h\) bounds
\(\|b\|,|n|,|\gamma|\) uniformly by \(C|\varepsilon|\).
For \(m\ne0\), the sine formula for \(s\) and the
preceding momentum bounds give
\[
 |b_m|+|f_m|\le
 C|m|_\infty^{-1}
       \prod_\mu\max\{1,2|m_\mu|-1\}^{-1}|\varepsilon|,
\]
where \(d(e^{ipx}b)=e^{ipx}f\). In the formula for
\(f_{\mu\nu,m}\), the terms containing \(\gamma\)
cancel. The displayed bound is summable independently
of \(k\), since
\(|m|_\infty\ge\prod_\mu\max(1,|m_\mu|)^{1/3}\).
It follows that the periodic potential and its curl
are uniformly bounded in supremum norm on the strip.

Finally, for fixed \(k\) and \(p\to0\),
\[
 b_m=O(|p|^2|\varepsilon|)\ (m\ne0),\qquad
 b_0=\varepsilon+O(|p||\varepsilon|),\qquad
 n=O(|p|^2|\varepsilon|).
\]
To prove this, use the trial field
\(b^{(0)}_{\mu,0}=\varepsilon_\mu/u_\mu(p)\) and
\(b^{(0)}_m=0\) otherwise. It satisfies \(Qb^{(0)}=\varepsilon\).
Its residual under \(X\) is \(O(|p|^2|\varepsilon|)\):
the zero mode has two derivatives, and every nonzero
mode contains \(v(-q_m)=O(|p|)\) times a divergence
of order \(|p|\).
The saddle matrix \(\mathbb S(X,Q)\) is invertible
at zero momentum, by coercivity on \(\ker Q\) and
surjectivity of \(Q\), and hence has bounded inverse
near zero. This proves the asserted expansions.

\subsection{The Gaussian coefficient along the critical branch}

We use absorbed bond lengths and unweighted transposes
when differentiating the critical branch. Let \(c_k\)
be its fine representative, \(p_k=\mathcal M^{k-1}c_k\)
its preceding-level representative, and \(\mathcal S_k\)
its critical action. The scaled free response \(h\)
and multiplier response \(\nu\) satisfy
\[
\begin{split}
 c_k'(\mathbf1)B-\eta hB&\in d\mathfrak G_0,\\
 (\nabla\mathcal S_k)'(\mathbf1)B&=\eta\nu B,\\
 \left.\partial_t\nabla A(c_k(e^{itB}))\right|_0
     &=\eta\mathring{\mathsf Q}^{[k]\mathsf T}\nu B,\\
 p_k'(\mathbf1)B-\mathring{\mathsf Q}^{[k-1]}(\eta hB)
     &\in d\mathfrak G_{k-1}.
\end{split}
\]
Here \(d\) in the displayed gradient spaces is unscaled.
We verify these identities, including the distinction
between the original and straightened means.

Put \(A_1=c_k'(\mathbf1)B\),
\(\mu_1=(\nabla\mathcal S_k)'(\mathbf1)B\),
\(A_2=\eta hB\), and \(\mu_2=\eta\nu B\).
Differentiate the constraint and the multiplier equation.
The first derivative of the plaquette action at
\(\mathbf1\) is zero, and its Hessian is \(d^{\mathsf T}d\).
Thus the mixed Campbell--Hausdorff term makes no
contribution, and
\[
 \mathsf Q^{[k]}A_1=B,\qquad
 d^{\mathsf T}dA_1=\mathsf Q^{[k]\mathsf T}\mu_1.
\]
The scaled response equations, after conversion to
unweighted pairings, give
\[
 \mathring{\mathsf Q}^{[k]}A_2=B,\qquad
 d^{\mathsf T}dA_2
     =\mathring{\mathsf Q}^{[k]\mathsf T}\mu_2,\qquad
 d^{\mathsf T}\mu_2=0.
\]
Use
\(\mathsf Q^{[k]}=\mathring{\mathsf Q}^{[k]}
-d\mathsf n^{[k]}\) and
\(\mathsf Q^{[k]}d=d\mathsf r\).
The correction
\(A'_2=A_2+d\mathsf r^{\mathsf T}\mathsf n^{[k]}A_2\)
satisfies the two equations for \(A_1\), with
\(\mu_2\) in place of \(\mu_1\).
Consequently \(E=A_1-A'_2\) satisfies
\[
 \mathsf Q^{[k]}E=0,\qquad
 |dE|^2=(\mathsf Q^{[k]}E,\mu_1-\mu_2)=0 .
\]
For real \(B\), \(E\) is real, so \(dE=0\).

By the curl-free decomposition used in
\cref{thm:gaussian-critical-germs}, write $E=d\lambda+a$
with $a$ constant. Its zero mean
constraint becomes \(d\mathsf r\lambda+L^ka=0\).
Summing around each coarse period gives \(a=0\).
Thus \(A_1-A_2\) is a gradient.
The mean \(\mathsf Q^{[k]}\) is onto, so its transpose
is injective and \(\mu_1=\mu_2\).
Finally differentiate \(p_k=\mathcal M^{k-1}c_k\);
both \(\mathsf Q^{[k-1]}d\) and
\(\mathsf Q^{[k-1]}-\mathring{\mathsf Q}^{[k-1]}\)
take values in gradients. This proves all four
identities, first for real \(B\), then complex linearly.

Write $T^v_{\mathfrak s},T^g_{\mathfrak s}$ for the two
saddle matrices in \eqref{eq:gaussian-background-functional}.
We apply \eqref{eq:gaussian-background-representation} to
the critical branch.
These are the vector and scalar saddle matrices of
\cref{prop:gaussian-background-determinant}. Their local
analytic formulas include the terms linear in $J,\mu$ in
\eqref{eq:gaussian-background-X}. At fixed height a common
neighbourhood makes every intermediate mean and contour
well defined. On a torus larger than four times their range,
each contour lifts uniquely to the covering lattice.
Differentiating the lifted identities, and then taking the
finite products and transposes, shows that both full saddle
matrices are the periodizations of their infinite-lattice
kernels. At real fields they have real matrices in the Pauli
basis. This proves the local and periodization hypotheses
for the logarithmic determinant calculation, including its
multiplier variables.

At zero field their Bloch matrices are
invertible for every real momentum: the vector saddle
uses the free coercivity on \(\ker Q\) and the
surjectivity of \(Q\); the scalar saddle uses the
corresponding facts for \(d^{\mathsf T}d\) and \(R\).
Compactness gives an invertible strip at each fixed
height. Hence the logarithmic determinant polarisation
formula applies to the full triple.

Denote its limiting bilinear symbol by
\[
 \mathsf p_{\mathfrak F}(p)[\varepsilon_+,\varepsilon_-]
       =-\tfrac12\ell(\mathbf T^v)+\ell(\mathbf T^g),
\]
where the jets are taken along
\(e^{\pm ipx}\varepsilon_\pm\) in the triple field.
Constant gauge invariance implies
\(D\mathfrak F_{\mathfrak s}(0)=0\): the coefficient
of each field variable would otherwise be a nonzero
fixed vector of the adjoint representation.
Differentiating the full gauge covariance, and using
this zero first derivative, gives the Ward identity
\[
 D^2\mathfrak F_{\mathfrak s}(0)
          [(d\lambda,0,0),\zeta]=0 .
\]
The first derivative of the triple in the Gaussian
formula is therefore equivalent, under its second
differential, to
\[
 \mathsf hB=
  \bigl(\eta hB,\,
        \eta\mathring{\mathsf Q}^{[k]\mathsf T}\nu B,\,
        \eta\nu B\bigr).
\]
The periodic operators defining this expression
commute with periodization. Write
\(\mathsf h(e^{ipx}\varepsilon)
=e^{ipx}\widehat{\mathsf h}(p)\varepsilon\).
The determinant trace formula then gives
\[
 \mathfrak s^{-3}
 \Pi^{\Psi_k}_{\mathfrak s}[e^a_{p,\mu},e^{a'}_{-p,\nu}]
 =\mathsf p_{\mathfrak F}(p)
       [\widehat{\mathsf h}(p)e_\mu t_a,
        \widehat{\mathsf h}(-p)e_\nu t_{a'}]
             +O(e^{-c\mathfrak s}),
\]
uniformly in the real torus momenta, at fixed \(k\).
The right side is bounded and holomorphic on a strip.
Colour invariance makes it a multiple of
\(\delta_{aa'}\) on torus momenta; density and
holomorphy give the same conclusion throughout the
strip. Together with the gauge and lattice invariances
of the Gaussian construction, this proves that
\(\Psi_k\) is tempered.

Write
\(\widehat{\mathsf h}=(\widehat{\mathsf h}^A,
\widehat{\mathsf h}^{\vartheta})\), separating the
potential from the two current components. The small
momentum estimate gives
\[
 \widehat{\mathsf h}^{\vartheta}(p)=O(|p|^2),\qquad
 \widehat{\mathsf h}^{A}(p)\varepsilon
                           =\eta\varepsilon+O(|p|).
\]
The Ward identity supplies one more power of momentum
in a mixed term. Indeed, apply it to
\(\lambda=e^{ipx}\varphi\), pass to the limiting symbol,
then set \(p=\tau w\) and \(\varphi=t_a\).
Since
\(\widehat d(\tau w)t_a=i\tau\eta(w_\mu t_a)_\mu
+O(\tau^2)\), division by \(\tau\) and passage to zero
show
\[
 \mathsf p_{\mathfrak F}(0)[(a,0),\varepsilon]=0
\]
for every constant potential \(a\); the same holds
in the second argument. Holomorphy and the expansion
of \(\widehat{\mathsf h}^A\) therefore make each
mixed potential-current term \(O(|p|^3)\), while the
current-current term is \(O(|p|^4)\). Thus
\[
 \widehat\pi^{\Psi_k}(p)=\widehat\pi^0(p)+O(|p|^3),
 \qquad
 \mathsf b[\Psi_k]=-\partial_1\partial_2
                                      \widehat\pi^0_{12}(0),
\]
where
\[
 \widehat\pi^0_{\mu\nu}(p)
 =\mathsf p_{\mathfrak F}(p)
  [(\widehat{\mathsf h}^A(p)e_\mu t_3,0),
   (\widehat{\mathsf h}^A(-p)e_\nu t_3,0)] .
\]

We next record the scale change for a tempered
functional \(F\). On a tower of side \(\mathfrak s\),
identify the preceding lattice with the unit lattice
of side \(L\mathfrak s\). Then \(F\circ p_k\), meaning
\((F_{L\mathfrak s}\circ p_k)_{\mathfrak s}\), is
tempered and
\[
                    \mathsf b[F\circ p_k]=L^{-1}\mathsf b[F].
\]
The vanishing first derivative of \(F\) and its Ward
identity reduce its composite polarisation to
\[
 \Pi^F_{L\mathfrak s}
 [\mathring{\mathsf Q}^{[k-1]}(\eta hB),
  \mathring{\mathsf Q}^{[k-1]}(\eta hB')] .
\]
For a plane wave, the straight path formula for the
average gives
\[
 \mathring{\mathsf Q}^{[k-1]}(\eta h e^a_{p,\mu})
   =L^{-1}\sum_{m'\in\mathcal M_L}\sum_\rho
       H_{m'}(p)_{\rho\mu}
                         e^a_{(p+2\pi m')/L,\rho},
\]
where \(\mathcal M_L\) is the centred set of
representatives modulo \(L\), and
\[
 H_{m'}(p)_{\rho\mu}
  =\sum_{m\equiv m'\,(\mathrm{mod}\,L)}
     u'_\rho(p+2\pi m)(\widehat h(p)e_\mu)_{\rho,m}.
\]
Here \(u'_\rho(q)\) is the average of the exponential
along the preceding-level straight paths; it is entire,
with \(u'_\rho(q)=1+O(|q|)\) near zero.
Only opposite \(m'\) pair under the translation-invariant
polarisation. The two factors \(L^{-1}\) and the volume
factor \(L^3\) give
\[
 \widehat\pi^{F\circ p_k}_{\mu\nu}(p)
  =L\sum_{m'\in\mathcal M_L}\sum_{\rho,\sigma}
     H_{m'}(p)_{\rho\mu}
     \widehat\pi^F_{\rho\sigma}((p+2\pi m')/L)
     H_{-m'}(-p)_{\sigma\nu}.
\]
The finite sum is holomorphic on a strip and inherits
the convergence of the finite-volume polarisations.
The branch covariance gives the remaining invariances,
so the composite is tempered. Moreover
\[
 H_{m'}(p)=O(|p|^2)\ (m'\ne0),\qquad
 H_0(p)=I+O(|p|),\qquad \widehat\pi^F(p)=O(|p|^2).
\]
It follows that
\(\widehat\pi^{F\circ p_k}(p)
=L\widehat\pi^F(p/L)+O(|p|^3)\), proving the asserted
coefficient identity.

The Gaussian gluing identity now yields, as families,
\[
 \Phi_l^\natural=\Psi_l-\Psi_{l-1}\circ p_l,\qquad
 \Psi_0=0 .
\]
Therefore each \(\Phi_l^\natural\) is tempered, and
\[
 \mathsf b[\Phi_l^\natural]
       =\mathsf b[\Psi_l]-L^{-1}\mathsf b[\Psi_{l-1}],
 \qquad
 \sum_{l=1}^kL^{-(k-l)}\mathsf b[\Phi_l^\natural]
       =\mathsf b[\Psi_k].
\]
The latter sum telescopes. The unique near-trivial
germs in the Gaussian construction depend only on the
tower, so \(\mathsf b[\Phi_l^\natural]\) depends only
on \(l\) and \(L\).

\subsection{A uniform determinant remainder}

We return to the physical pairings of weight $\ell_j^3$
and to differences $d$ divided by the lattice spacing.

We first bound the field derivatives which enter the
determinants. There are numbers \(\rho_*>0\) and \(C_*>0\),
depending only on \(L\), such that the scaled straight
average \(Q_U\) and scalar average \(R_U\), for
\(U=e^{i\eta a}\), are holomorphic on
\(\|a\|_\infty<\rho_*\) and have operator norms at most
\(C_*\), independently of the height.

Here is the uniform-radius argument. Write the one-step
mean in logarithmic coordinates as
\[
 A_j=\mathsf Q A_{j-1}+N(A_{j-1}),\qquad
 |N(A)|_\infty\le C_N|A|_\infty^2,\qquad A_0=\eta a .
\]
Choose $r_{\log}>0$ with $e^{r_{\log}}-1<\rho_L$.
The one-step remainder is holomorphic on that ball and
satisfies the displayed quadratic bound on its half-ball,
by its double zero and Schwarz's lemma.
The folding identity gives
\(\|\mathsf Q^{[n]}\|_{\infty\to\infty}\le4L^n\).
With \(a_j=\ell_j^{-1}|A_j|_\infty\), iteration gives
\[
 a_j\le4a_0+4C_NL^{-2}
              \sum_{i=1}^j\ell_i a_{i-1}^2,\qquad
 \sum_{i=1}^j\ell_i\le\tfrac32 .
\]
Choose \(\rho_*\) so that
\[
 384C_NL^{-2}\rho_*\le4,\quad
 16\rho_*\le r_{\log},\quad 8\rho_*\le1,\quad
 32\rho_*<\rho_L ,
\]
where \(r_{\log},\rho_L\) are the one-step logarithmic and
holomorphic radii. Induction gives
\(a_j\le8a_0\) throughout the complex ball. Thus every
intermediate configuration \(V_j=e^{iA_j}\) satisfies
\[
 |V_j-\mathbf1|_\infty\le16\ell_j|a|_\infty .
\]
These estimates also ensure that all intermediate
logarithms are the specified analytic branches.

Define the one-step straightened kernel by
\[
 (\mathring{\mathcal Q}_VB)(c)
          =\sum_{b\in\mathsf P_c^+}k_V(c,b)B(b).
\]
The one-step derivative kernels obey
\[
 |k_V-k_{\mathbf1}|+|\operatorname{Ad}V(y,x)-1|
 \le C_L|V-\mathbf1|_\infty .
\]
To see this uniformly on the smaller ball, apply Cauchy's
formula first to the left variation \(e^{itB}V'\), with
\(r'_L=\rho_L/16\). The mean derivative, its gauge
correction, and their straightened sum are bounded by
\(1/r'_L,1/r'_L,5/r'_L\), respectively.
Apply it next to \(e^{isA}\), with
\(r''_L=\log(1+3\rho_L/4)\).
Since \(|\log V|\le2|V-\mathbf1|\), subtraction at
\(s=0\) gives the kernel bound
\(20|A|/(r'_Lr''_L)\) and the transport bound
\(4|A|/r''_L\); when \(2|A|>r''_L\), the same bounds
follow from the preceding absolute bounds after increasing
the constant. Conversion to the fixed Hilbert norm costs
only \(2\sqrt2\).

A double block contains at most \(6L^3\) fine bonds,
and a fine bond belongs to at most six double blocks.
Schur's estimate therefore bounds the difference between
a scaled one-step straight average and its free value
by \(C'_L|V-\mathbf1|_\infty\). The scalar average has
disjoint blocks and obeys the same estimate. The free
factors have norm at most one. Since
\(\sum_{j=1}^k\ell_{j-1}\le3/2\), the products defining
\(Q_U,R_U\) have norms at most
\(\exp(24C'_L)\), after decreasing \(\rho_*\le1\).
This proves the assertion. Cauchy's formula on the
whole ball also gives bounds \(C_*/\rho_*\) and
\(4C_*/\rho_*^2\) for the first and mixed second
derivatives.
The products retain their block locality. Each
intermediate contour lies within its block, and the
union of the preceding-level blocks making up a
later block is precisely that block. Induction
therefore confines the row of \(Q_U\) at \(c\) and
its field dependence to bonds with both endpoints in
\(B^k(c_-)\cup B^k(c_+)\). The row of \(R_U\) at
\(y\) and its field dependence lie in \(B^k(y)\).

We apply these bounds along the plane-wave responses
\[
 \zeta_\pm=\eta e^{\pm ip\cdot x}\alpha_\pm,\qquad
 \alpha_\pm=\widehat h(\pm p)e_{\mu_\pm}t_3 .
\]
The preceding free estimate bounds both \(\alpha_\pm\)
and their covariant curls in supremum norm. We use
bold letters for jets along these two fields.
All the jets below have width \(\delta_{\mathrm{loop}}\),
with \(|p_\mu|\le\delta_{\mathrm{loop}}\).

The Hessian has a useful local form. For an oriented
plaquette with bonds \(b_1,\ldots,b_4\), put
\[
 T_1=1,\quad T_2=U(b_1),\quad
 T_3=U(b_1)U(b_2)U(b_3)^{-1},\quad T_4=U_p,
 \qquad(\varsigma_i)=(+,+,-,-).
\]
Moving all four exponentials to the left in
\((e^{itB}U)_p\) gives
\(\prod_i e^{itY_i}U_p\), where
\(Y_i=\varsigma_i\operatorname{Ad}(T_i)B(b_i)\).
Its second derivative uses
\[
 \left.\partial_t^2\prod_i e^{itY_i}\right|_0
   =-\Bigl(\sum_iY_i\Bigr)^2-\sum_{i<j}[Y_i,Y_j].
\]
The Pauli identities
\(Y^2=\langle Y,Y\rangle I\) and
\(U_p^{-1}=(\operatorname{tr}U_p)I-U_p\) now give
\[
 H_U=(d+E_U)^{\mathsf T}z_U(d+E_U)+f_U,
\]
where
\[
\begin{split}
 E_UB(p)&=\eta^{-1}\sum_i\varsigma_i
                  (\operatorname{Ad}(T_i)-1)B(b_i),\\
 z_U(p)&=\tfrac12\operatorname{tr}U_p,\qquad
 F_U(p)=\frac{U_p-U_p^{-1}}{2i\eta^2},\\
 (B,f_UB)&=\eta^3\sum_p\tfrac i2\operatorname{tr}
                         \sum_{i<j}[Y_i,Y_j]F_U(p).
\end{split}
\]
Write \(e_U=D_U-d\). At \(U=e^{i\eta a}\), its first
derivative is \(i\operatorname{ad}a\), and its second
coefficient is \(-\eta(\operatorname{ad}a)^2/2\).
The derivatives of \(E_U\) are sums of at most four
such terms. The first derivative of \(z_U-1\) vanishes;
its quadratic coefficient is
\(-\eta^4\langle da,da\rangle/2\).
The first derivative of \(f_U\) contains one \(da\);
its second derivative is bounded by
\(C(|a|^2+\eta|a||da|)|B|^2\).
Consequently the jets of
\[
 e_U,\ E_U,\ z_U,\ f_U,\ Q_U,\ R_U
\]
are bounded by a constant depending only on \(L\).
Indeed, the first four kernels couple indices in one
plaquette, with at most thirteen neighbours; their
Bloch factors on this strip are uniformly bounded.
For the last two kernels use the preceding Cauchy bounds,
the double-block overlap of six (one for scalar blocks),
and the bounded exponential factor across a block.
Their entries, and hence their jets, are holomorphic
in \(p\): each entry is a finite sum of local kernel
derivatives multiplied by exponentials in the momentum.
The \(00\) components of \(e_U,E_U,z_U-1,f_U\) vanish.

We shall also use the following trace bound:
\[
 |\mathsf t(G\mathbf B_1G\mathbf B_2G\mathbf B_3)|
 \le C_{\mathrm{tr}}\prod_{j=1}^3\|\mathbf B_j\|.
\]
The same bound holds with \(g\) in place of \(G\).
Expand the \(11\) component. Every term has three
propagators evaluated at \(\kappa\) or \(\kappa\pm p\).
Schatten H\"older followed by integral H\"older reduces
the estimate to the integral of the cube of one free
propagator. The Bloch sectors unfold to the fine
Brillouin cube, and
\[
 \int\|\widehat G(\kappa+\vartheta)\|_3^3\,d\kappa
 \le9(10/9)^3(2\pi)^{-3}
     \int_{\mathbb R^3}(1+4|q|^2/\pi^2)^{-3}\,dq .
\]
Here \(|\operatorname{Im}\vartheta_\mu|\le2\delta_{\mathrm{loop}}\)
and the previously proved lower bound for
\(1+D(q+\vartheta)\) was used. The integral is finite
and independent of \(\eta\). The finite number of jet
components is absorbed in \(C_{\mathrm{tr}}\).

Put
\[
 \mathcal L(z)=\ell(\boldsymbol\Delta+z)
                         -\tfrac12\ell(\mathbf K+z),
 \qquad
 \mathcal L^\sharp=
       \bigl(1-\partial_z+\tfrac12\partial_z^2\bigr)
                      \mathcal L(z)\big|_{z=1}.
\]
For \(z>0\) these jets are invertible at width
\(\delta_z=\min(\delta_{\mathrm{loop}},\sqrt z/5)\).
In fact, on \(S_{2\delta_z}\),
\(\operatorname{Re}(D(q)+z)\ge z/3\).
Near \(z=1\) the width is fixed, and differentiation
in the jet algebra gives
\[
 \partial_z\ell(\mathbf K+z)
       =\mathsf t((\mathbf K+z)^{-1}),\qquad
 \partial_z^2\ell(\mathbf K+z)
       =-\mathsf t((\mathbf K+z)^{-2}),
\]
and the same identities for \(\boldsymbol\Delta\).

\begin{proposition}\label{I-109-75}
Let \(\widehat\pi^0\) be the polarisation of the
Gaussian determinant evaluated on the potential
responses \(\zeta_\pm\), with the two multiplier
components set to zero. Then
\(\widehat\pi^0_{\mu_+\mu_-}-\mathcal L^\sharp\)
is holomorphic for \(|p_\mu|<\delta_{\mathrm{loop}}\)
and has absolute value at most \(C_{\mathrm{rem}}(L)\),
uniformly in \(k\).
\end{proposition}

\begin{proof}
The vector and scalar saddle matrices in the Gaussian
formula, up to field-independent changes of pairing
weights, are
\[
 \mathbb S(X_U,Q_U),\qquad \mathbb S(\Delta_U,R_U),
 \qquad X_U=\Xi_U-Q_U^{\mathsf T}Q_U.
\]
Set
\[
 \mathbf V=\mathbf Q^{\mathsf T}\mathbf Q
            -\mathbf D\mathbf R^{\mathsf T}\mathbf R\mathbf D^{\mathsf T},
 \qquad\boldsymbol\Xi=\mathbf K+\mathbf V,\qquad
 \mathbf Y=\boldsymbol\Delta+\mathbf R^{\mathsf T}\mathbf R .
\]
Multiplication by a triangular jet of determinant one
changes \(\mathbb S(\boldsymbol\Xi-\mathbf Q^{\mathsf T}
\mathbf Q,\mathbf Q)\) into
\(\mathbb S(\boldsymbol\Xi,\mathbf Q)\).
Its Schur factorisation has diagonal blocks
\(\boldsymbol\Xi\) and
\(-\mathbf Q\boldsymbol\Xi^{-1}\mathbf Q^{\mathsf T}\).
The free strip estimates make both \(00\) blocks
invertible; hence the jet factorisation is legitimate.
The scalar factorisation is identical. Additivity of
\(\ell\) therefore gives
\[
 \widehat\pi^0
 =-\tfrac12\ell(\boldsymbol\Xi)
  -\tfrac12\ell(\mathbf Q\boldsymbol\Xi^{-1}\mathbf Q^{\mathsf T})
  +\ell(\mathbf Y)
  +\ell(\mathbf R\mathbf Y^{-1}\mathbf R^{\mathsf T}).
\]

Write \(\mathbf G_U=(\mathbf K+1)^{-1}\) and
\(\mathbf P_\tau=(\boldsymbol\Xi+\tau)^{-1}\).
Integrate first in the added constraint \(\mathbf V\)
and then in the mass:
\[
 \ell(\boldsymbol\Xi)=\ell(\mathbf K+1)
 +\int_0^1\mathsf t((\mathbf K+1+t\mathbf V)^{-1}\mathbf V)\,dt
 -\int_0^1\mathsf t(\mathbf P_\tau)\,d\tau .
\]
Two uses of the resolvent identity give
\[
\begin{split}
 \mathbf P_\tau={}&\mathbf G_U+(1-\tau)\mathbf G_U^2
 +(1-\tau)^2\mathbf G_U^2\mathbf P_\tau\\
 &-\mathbf G_U\mathbf V\mathbf P_\tau
 -(1-\tau)\mathbf G_U^2\mathbf V\mathbf P_\tau .
\end{split}
\]
Thus
\[
 \ell(\boldsymbol\Xi)
 =\bigl(1-\partial_z+\tfrac12\partial_z^2\bigr)
                  \ell(\mathbf K+z)\big|_{z=1}+\rho_v ,
\]
where
\[
\begin{split}
 \rho_v={}&\int_0^1
     \mathsf t((\mathbf K+1+t\mathbf V)^{-1}\mathbf V)\,dt\\
 &-\int_0^1\bigl[
 (1-\tau)^2\mathsf t(\mathbf G_U^2\mathbf P_\tau)
 -\mathsf t(\mathbf G_U\mathbf V\mathbf P_\tau)
 -(1-\tau)\mathsf t(\mathbf G_U^2\mathbf V\mathbf P_\tau)
                         \bigr]\,d\tau .
\end{split}
\]
The same computation for
\((\boldsymbol\Delta,\mathbf R^{\mathsf T}\mathbf R,\mathbf Y)\)
defines \(\rho_g\). It remains to bound these two
remainders and the two coarse Schur determinants.

Introduce bounded jets
\[
 \boldsymbol\kappa=G^{1/2}(\mathbf K-K)G^{1/2},
 \quad x=\mathbf QG^{1/2},\quad
 y=G^{1/2}\mathbf D\mathbf R^{\mathsf T},\quad
 \omega=x^{\mathsf T}x-yy^{\mathsf T}.
\]
The local Hessian identity above proves that
\(\boldsymbol\kappa\) is bounded. Each term contains
only a bounded jet among \(e,E,z,f\), a square-root
propagator, or a propagator with one derivative;
for example its Hessian part is
\[
\begin{split}
 &(dG^{1/2})^{\mathsf T}(z-1)dG^{1/2}
 +(dG^{1/2})^{\mathsf T}zEG^{1/2}
 +(EG^{1/2})^{\mathsf T}zdG^{1/2}\\
 &\quad +(EG^{1/2})^{\mathsf T}zEG^{1/2}
                 +G^{1/2}fG^{1/2}.
\end{split}
\]
Its remaining part is the expansion of
\(G^{1/2}[(d+e)(d+e)^{\mathsf T}-dd^{\mathsf T}]G^{1/2}\).
In particular \(\kappa_{00}=0\). Set
\[
 B_t=(1+\boldsymbol\kappa+t\omega)^{-1},\qquad
 B'_\tau=(1+\boldsymbol\kappa+\omega-(1-\tau)G)^{-1}.
\]
The free strip estimates bound the inverses of the
\(00\) components uniformly in \(t,\tau\); the finite
nilpotent inverse formula therefore bounds the full
jets uniformly. We have
\(\mathbf G_U=G^{1/2}B_0G^{1/2}\) and
\(\mathbf P_\tau=G^{1/2}B'_\tau G^{1/2}\).

Cyclicity converts the four traces in \(\rho_v\) to
\[
\begin{split}
 &\mathsf t(xB_tx^{\mathsf T})
                         -\mathsf t(y^{\mathsf T}B_ty),\\
 &\mathsf t(xB'_\tau GB_0x^{\mathsf T})
                         -\mathsf t(y^{\mathsf T}B'_\tau GB_0y),\\
 &\mathsf t(xB'_\tau GB_0GB_0x^{\mathsf T})
                         -\mathsf t(y^{\mathsf T}B'_\tau GB_0GB_0y),\\
 &\mathsf t(GB_0GB_0GB'_\tau).
\end{split}
\]
The first three traces act on a coarse bond or site
fibre, of dimension at most nine, and are bounded by
the jet norm. The fourth is bounded by the
three-propagator estimate. For the scalar remainder
use
\[
 \kappa_g=g^{1/2}(\boldsymbol\Delta-\Delta)g^{1/2}
 =(dg^{1/2})^{\mathsf T}eg^{1/2}
 +(eg^{1/2})^{\mathsf T}dg^{1/2}
 +(eg^{1/2})^{\mathsf T}eg^{1/2},
 \qquad x_g=\mathbf Rg^{1/2}.
\]
There is no \(y\) term. The same proof bounds \(\rho_g\).
Finally
\(\mathbf Q\boldsymbol\Xi^{-1}\mathbf Q^{\mathsf T}
=xB'_0x^{\mathsf T}\), and the scalar counterpart
has the same form. Their jet norms and their
\(00\)-inverse norms are bounded; their fibres have
dimension at most nine. The finite formula for \(\ell\)
therefore bounds their logarithmic polarisations.

It follows that
\[
 \widehat\pi^0-\mathcal L^\sharp
 =-\tfrac12\rho_v+\rho_g
  -\tfrac12\ell(\mathbf Q\boldsymbol\Xi^{-1}\mathbf Q^{\mathsf T})
  +\ell(\mathbf R\mathbf Y^{-1}\mathbf R^{\mathsf T})
\]
is uniformly bounded. All jets are holomorphic in
\(p\), as are their inverses; the bounds justify the
trace and parameter integrals. This also proves the
asserted holomorphy.
\end{proof}

\subsection{The massive coefficient}

In this subsection the lattice spacing is one, forms have
absorbed lengths, and all pairings are unweighted. Put
\[
 \psi(q)=\sum_{\mu=1}^3\psi^{(\mu)}(q),\quad\psi^{(\mu)}(q)=2-2\cos q_\mu,\qquad
 \Upsilon_z(q)=(\psi(q)+z)^{-1}.
\]
For \(z>0\), let
\[
 \Delta_U=\mathsf D_U^{\mathsf T}\mathsf D_U,\qquad
 \mathsf K_U=A''(U)+\mathsf D_U\mathsf D_U^{\mathsf T},
\]
and define on a finite torus, near the trivial configuration,
\[
 \mathfrak L_{z,\mathfrak s}(U)
 =\left[\log\det(\Delta_{U'}+z)
       -\tfrac12\log\det(\mathsf K_{U'}+z)
   \right]_{U'=\mathbf1}^{U'=U}.
\]
At \(U=\mathbf1\) the symbols of both operators are
\(\psi(q)I\). Indeed, the scalar operator is
\(d^{\mathsf T}d\), and the form operator is
\(d^{\mathsf T}d+dd^{\mathsf T}\); the mixed terms cancel
in their Fourier symbols. Both massive matrices are
positive at \(\mathbf1\), so their real logarithms are
real analytic on a neighbourhood of \(\mathbf1\).

Their kernels are entire local functions of the absorbed
potential and commute with translations, gauge transformations
and lattice isometries. The logarithmic determinant
polarisation formula therefore makes
\(\mathfrak L_z=(\mathfrak L_{z,\mathfrak s})_{\mathfrak s}\)
a tempered functional: its zero-field inverse symbols
\((\psi+z)^{-1}I\) are holomorphic on a strip, the
finite-torus Fourier sums converge exponentially to the
corresponding integrals, and colour invariance makes the
polarisation diagonal in the Pauli basis. We write
\(\widetilde\beta(z)=\mathsf b[\mathfrak L_z]\).
This is real analytic for \(z>0\), since its polarisation
is an integral of products of \((\psi+z)^{-1}\) and
the first two derivatives of the local kernels.

\begin{proposition}\label{I-109-99}
There is an absolute constant \(C_{\mathrm{uv}}\) such that
\[
 z^{1/2}|\widetilde\beta(z)|
 +z^{3/2}|\widetilde\beta'(z)|
 +z^{5/2}|\widetilde\beta''(z)|
 \le C_{\mathrm{uv}},\qquad 0<z\le1.
\]
\end{proposition}

\begin{proof}
We first identify the coefficient by a local trace,
then compute that trace. For a real bond function
\(\theta\), put \(U_\theta=e^{i\theta t_3}\), and set
\(\vartheta^\circ_2(x)=x_1\),
\(\vartheta^\circ_1=\vartheta^\circ_3=0\).
The plaquette angles of \(f\vartheta^\circ\) are \(f\)
in the \(12\)-plane and zero in the other two planes.
Let \(\tau_x\) be the trace on indices based at \(x\).
For sufficiently small \(|f|\), define
\[
 \varrho_x(\theta)=\tau_x\log(\Delta_\theta+z)
                   -\tfrac12\tau_x\log(\mathsf K_\theta+z).
\]
These logarithms are well defined even though
\(\vartheta^\circ\) is unbounded.

To verify this and the estimates used below, write
\(t_\pm=(t_1\pm it_2)/2\) and
\(B=B^3t_3+wt_++\overline w t_-\).
The shifts \(u_\mu w(x)=e^{2i\theta_\mu(x)}w(x+e_\mu)\)
are unitary. At a plaquette \(p=p_{\mu\nu}(x)\), put
\[
 (y_1,y_2,y_3,y_4)
  =(w_\mu,u_\mu w_\nu,-e^{2if_p}u_\nu w_\mu,
                                  -e^{2if_p}w_\nu)(x),
 \qquad f_p=(d\theta)(p).
\]
Differentiating its four Wilson factors gives
\[
\begin{split}
 (B,\mathsf K_\theta B)
 ={}&\sum_p\left[
  \cos f_p\left((dB^3)(p)^2+\left|\sum_i y_i\right|^2\right)
  -2\sin f_p\sum_{i<j}\operatorname{Im}(y_i\overline y_j)
                  \right]\\
 &+\sum_x\left[(d^{\mathsf T}B^3)(x)^2
       +\left|\sum_\mu(u_\mu^*-1)w_\mu(x)\right|^2\right].
\end{split}
\]
Here \(\langle t_+,t_-\rangle=1/2\) and
\([t_+,t_-]=t_3\). Since
\(\sum_{i<j}|y_i||y_j|\le(3/2)\sum_i|y_i|^2\)
and a bond belongs to four plaquettes,
\[
 \Delta_\theta\ge0,\qquad
 \mathsf K_\theta\ge-12\sup_p|\sin f_p|
       \quad\text{if }\sup_p|f_p|\le\pi/2.
\]
In particular \(\mathsf K_{f\vartheta^\circ}+z\ge z/2\)
for small \(|f|\). Both kernels have range at most two
and uniformly bounded trigonometric-polynomial entries.

\emph{Identification by a local density.}
For either operator \(T_\theta\), the resolvent obeys
\[
 |(T_\theta+z+w)^{-1}(i,j)|
 \le (z/4+w)^{-1}e^{-c|x(i)-x(j)|},\qquad w\ge0,
\]
uniformly under the preceding lower bound. Indeed,
conjugation by \(e^{\max(-N,\min(N,\sigma\cdot x))}\)
changes its finite-range kernel by an operator of norm
at most \(C(e^{2|\sigma|}-1)\). Choose \(|\sigma|=c\)
so this is at most \(z/4\), apply the Neumann series,
and let \(N\to\infty\), choosing the direction of
\(\sigma\) along \(x(i)-x(j)\).

The identity
\[
 \log(T+z)=\int_0^\infty
       \big[(1+w)^{-1}-(T+z+w)^{-1}\big]\,dw
\]
then gives a local second-derivative kernel
\(\mathcal P(x;b,b')
 =\partial_{\theta(b)}\partial_{\theta(b')}
                         \varrho_x(0)\), with
\[
 |\mathcal P(x;b,b')|
       \le C e^{-c(|x-b_-|+|x-b'_-|)}.
\]
For completeness, if \(R_w=(T_0+z+w)^{-1}\), the
integrand of this derivative is the local trace of
\[
 R_wT^{(2)}[\delta_b,\delta_{b'}]R_w
 -R_wT^{(1)}[\delta_b]R_wT^{(1)}[\delta_{b'}]R_w
 -R_wT^{(1)}[\delta_{b'}]R_wT^{(1)}[\delta_b]R_w.
\]
The differentiated kernels have finite support near the
indicated bonds. The bound
\(C(z/4+w)^{-2}\) is integrable and justifies
differentiation under the integral.
Summing in \(x\), using trace cyclicity and
\(\partial_wR_w=-R_w^2\), yields
\(\operatorname{Tr}(R_0T^{(2)}-R_0T^{(1)}R_0T^{(1)})\).
Periodising \(R_0\) gives the finite-torus inverse; its
exponential decay shows that these finite-support traces
converge to their infinite-lattice values. Consequently
\[
 \sum_x\mathcal P(x;b,b')=\pi_{\mu\nu}(b'_--b_-).
\]

The local trace is invariant under
\(\theta\mapsto\theta+d\lambda\), by conjugation with
\(e^{-i\lambda t_3}\), and is even in \(\theta\), by
conjugation with \(it_1\).
Expand the resolvent at \(f\vartheta^\circ\) to second
order. The Taylor coefficients of the finite-range
kernel grow at most quadratically in position, and the
remainder is \(O(|f|^3(1+|x|)^3)\). The resolvent identity
with three insertions and the preceding exponential
bound make all local sums absolutely convergent, with
an integrable \(O(|f|^3(z/4+w)^{-2})\) remainder.
Thus
\[
 \varrho_0(f\vartheta^\circ)-\varrho_0(0)
 =\tfrac12f^2\sum_{b,b'}\mathcal P(0;b,b')x_1x'_1
      +O(|f|^3),
\]
where both bonds have direction two.
Gauge invariance with \(\lambda(x)=x_2\) gives
\(\sum_{b'}\mathcal P(0;b,b')=0\).
Using symmetry, translation covariance, and
\(2x_1x'_1=x_1^2+(x'_1)^2-(x_1-x'_1)^2\), the sum equals
\[
 -\tfrac12\sum_u u_1^2\pi_{22}(u)=\widetilde\beta(z).
\]
We have proved
\[
 \varrho_0(f\vartheta^\circ)-\varrho_0(0)
             =\tfrac12\widetilde\beta(z)f^2+O(|f|^3).
\]

\emph{The magnetic resolvent calculation.}
Put \(\varphi=2f\). Then \(u_1u_2=e^{i\varphi}u_2u_1\)
and \(u_3\) commutes with both. For
\[
 W_\varphi(n)=e^{-i\varphi n_1n_2/2}
                     u_1^{n_1}u_2^{n_2}u_3^{n_3},
\]
we have
\[
 W_\varphi(n)W_\varphi(n')
   =e^{i\varphi(n_1n'_2-n_2n'_1)/2}W_\varphi(n+n'),
 \qquad \tau_0 W_\varphi(n)=\delta_{n0}.
\]
For a matrix symbol \(a(q)=\sum_n a_ne^{in\cdot q}\),
write \(\operatorname{Op}_\varphi(a)=\sum_n a_nW_\varphi(n)\)
and \(\|a\|_{(j)}=\sum_n|a_n|(1+|n|)^j\).
The product symbol is
\[
 a\star b=ab-\tfrac{i\varphi}{2}\{a,b\}
                  -\tfrac{\varphi^2}{8}\{\!\{a,b\}\!\}
                  +O_{(0)}(|\varphi|^3\|a\|_{(3)}\|b\|_{(3)}),
\]
with remainder constant \(1/48\), where
\(\{a,b\}=a_1b_2-a_2b_1\) and
\(\{\!\{a,b\}\!\}=a_{11}b_{22}-2a_{12}b_{12}+a_{22}b_{11}\).
This follows term by term from the exponential remainder
\(|e^{is}-1-is+s^2/2|\le|s|^3/6\).
Also \(\|\operatorname{Op}_\varphi(a)\|\le\|a\|_{(0)}\)
and \(\tau_0\operatorname{Op}_\varphi(a)=\int\operatorname{tr}a\).

Suppose a Hermitian $N'\times N'$ trigonometric-polynomial matrix has the expansion
\(\sigma_\varphi=\psi I+\varphi\sigma_1+\varphi^2\sigma_2
 +O_{(3)}(|\varphi|^3)\), and
\(\operatorname{Op}_\varphi(\sigma_\varphi)+z\ge z/2\).
With \(U_w=(\psi+z+w)^{-1}\), a right resolvent parametrix is
\[
 r_0=U_w I,\qquad r_1=-U_w^2\sigma_1,\qquad
 r_2=U_w^3\sigma_1^2-U_w^2\sigma_2
                  +\tfrac18U_w\{\!\{U_w,\psi\}\!\}I .
\]
The coefficients of orders one and two in
\((r_0+\varphi r_1+\varphi^2r_2)\star
                 (\sigma_\varphi+z+w)-I\) vanish:
\(\{U_w,\psi\}=0\), and
\(\{-U_w^2\sigma_1,\psi\}+\{U_w,\sigma_1\}=0\).
On a strip depending on \(z\), the symbol norms of these
terms are bounded by \(C_z(z+w)^{-1}\).
The product error is therefore
\(O_{(0)}(|\varphi|^3(z+w)^{-1})\).
Multiplying by the exact inverse bounds the local-trace
error by \(C_z|\varphi|^3(z+w)^{-1}(z/2+w)^{-1}\),
which is integrable in \(w\).
Integrating the resolvent gives
\[
\begin{split}
 \tau_0\log(\operatorname{Op}_\varphi(\sigma_\varphi)+z)
 ={}&\int\operatorname{tr}\log(\psi+z)
       +\varphi\int\Upsilon_z\operatorname{tr}\sigma_1\\
 &+\varphi^2\int\left[
  -\tfrac12\Upsilon_z^2\operatorname{tr}\sigma_1^2
  +\Upsilon_z\operatorname{tr}\sigma_2
  +\tfrac{N'}{24}\Upsilon_z^2\psi_{11}\psi_{22}
                  \right]+O(|\varphi|^3).
\end{split}
\]
To check the last term, use \(\psi_{12}=0\),
\[
 \{\!\{U_w,\psi\}\!\}
  =-2U_w^2\psi_{11}\psi_{22}
       +2U_w^3(\psi_1^2\psi_{22}+\psi_2^2\psi_{11}),
\]
then integrate in \(w\). Periodic integration by parts gives
\(\int\Upsilon_z^3\psi_1^2\psi_{22}
 =\frac12\int\Upsilon_z^2\psi_{11}\psi_{22}\),
and the same formula with the indices exchanged.

\emph{The three colour sectors.}
The neutral and two charged sectors are invariant.
Complex conjugation exchanges the charged sectors,
so their local traces agree. The neutral scalar operator
is independent of \(f\); either charged scalar operator
is \(\operatorname{Op}_\varphi(\psi)\).
The neutral form operator is
\[
 d^{\mathsf T}d+dd^{\mathsf T}
                         +(\cos f-1)d_{12}^{\mathsf T}d_{12}.
\]
If \(\sigma^{\mathrm n}\) is the symbol of
\(d_{12}^{\mathsf T}d_{12}\), then
\(0\le\sigma^{\mathrm n}\le\psi I\) and
\(\operatorname{tr}\sigma^{\mathrm n}
                       =\psi^{(1)}+\psi^{(2)}\).
Its logarithmic trace changes by
\(-\frac12 f^2\int\Upsilon_z(\psi^{(1)}+\psi^{(2)})
 +O(f^4)\).

The displayed plaquette quadratic form shows that the
charged form operator is a polynomial of degree at most
two in the \(u_\mu^{\pm1}\), with coefficients analytic
in \(\varphi\). Reordering these shifts therefore writes it
as \(\operatorname{Op}_\varphi(\sigma_\varphi)\), where
\[
 \sigma_\varphi=\psi I+\varphi\sigma_1+\varphi^2\sigma_2
                          +O_{(3)}(|\varphi|^3).
\]
The matrices \(\sigma_1,\sigma_2\) are fixed Hermitian
trigonometric polynomials, independent of \(z\).
The preceding lower bound verifies the positivity
hypothesis of the parametrix calculation.
Adding the two charged scalar traces, subtracting
one half of the two charged form traces and the
neutral form trace, and using \(\varphi^2=4f^2\), gives
\[
 \widetilde\beta(z)=
 8\int\left[
   \tfrac12\Upsilon_z^2\operatorname{tr}\sigma_1^2
   -\tfrac1{24}\Upsilon_z^2\psi_{11}\psi_{22}
   -\Upsilon_z\operatorname{tr}\sigma_2\right]
 +\tfrac12\int\Upsilon_z(\psi^{(1)}+\psi^{(2)}).
\]
The linear term vanishes by evenness of the local density.

Thus \(\widetilde\beta(z)=\int(\Upsilon_z^2g_2+
\Upsilon_zg_1)\), where \(g_1,g_2\) are fixed bounded
trigonometric polynomials. Differentiating under the
integral, for \(n=0,1,2\),
\[
 |\widetilde\beta^{(n)}(z)|
 \le C\left[(n+1)!\int\Upsilon_z^{n+2}
                         +n!\int\Upsilon_z^{n+1}\right].
\]
Finally \(\psi(q)\ge4|q|^2/\pi^2\) on the period cube.
For \(j\ge2\), the substitution \(q=\sqrt z\,q'\)
gives \(\int\Upsilon_z^j\le C_jz^{3/2-j}\);
\(\int\Upsilon_z\le C_1\), because \(|q|^{-2}\)
is locally integrable in three dimensions.
These estimates prove the asserted bound.
\end{proof}

\subsection{The accumulated coefficient}

The scaled massive polarisation is related to the
unit-spacing massive coefficient by
\[
 -\partial_{p_1}\partial_{p_2}\mathcal L(z)(0)
                       =\eta\widetilde\beta(\eta^2z).
\]
To verify the normalization, use
\[
 \Delta_U=\eta^{-2}\mathsf D_U^{\mathsf T}\mathsf D_U,
 \qquad K_U=\eta^{-2}\mathsf K_U .
\]
The logarithmic determinant combination at mass \(z\)
therefore differs from the unscaled combination at
mass \(\eta^2z\) by a field-independent constant.
Under \(x\mapsto x/\eta\), the fine torus has side
\(\mathfrak s/\eta\). Expand the two response fields:
\[
 \zeta_\pm
   =\eta\sum_{m,\rho}b^\pm_{\rho,m}
                   e^{i(\pm p+2\pi m)x}e_\rho t_3 .
\]
Momentum conservation pairs only \(m\) with \(-m\).
The two field factors \(\eta\), and the torus volume
ratio \(\eta^{-3}\), give
\[
 \mathcal L(z)(p)
 =\eta^{-1}\sum_{m,\rho,\sigma}
   \widehat\pi^{(\eta^2z)}_{\rho\sigma}
                          (\eta(p+2\pi m))
                     b^+_{\rho,m}b^-_{\sigma,-m}.
\]
This identity follows first at the common torus
momenta from the finite-volume determinant formula
and then near zero by holomorphy. At fixed \(k\),
\[
 b^\pm_{\rho,m}=O(|p|^2)\quad(m\ne0),\qquad
 b^\pm_{\rho,0}=\delta_{\rho\mu_\pm}+O(|p|).
\]
Since the massive polarisation is \(O(|p|^2)\) at
zero, all nonzero modes contribute \(O(|p|^4)\), and
\[
 \mathcal L(z)(p)
   =\eta^{-1}\widehat\pi^{(\eta^2z)}_{\mu_+\mu_-}
                                             (\eta p)+O(|p|^3).
\]
Two derivatives in \(p\) prove the asserted identity.
The massive inverse symbols and the local kernel
derivatives are jointly analytic in \(z>0\) and in
\(p\) near zero. Hence mass and momentum derivatives
commute, giving
\[
 -\partial_1\partial_2\mathcal L^\sharp(0)
 =\eta\left[
   \widetilde\beta(\eta^2)-\eta^2\widetilde\beta'(\eta^2)
               +\tfrac12\eta^4\widetilde\beta''(\eta^2)
                                  \right].
\]



\begin{proof}[Proof of \cref{I-109-104}]
Temperedness and the identity for \(\mathsf c_k\)
were proved from Gaussian gluing. The kernel limit,
reality, exponential decay and the coefficient formula
follow from the tempered polarisation calculation.
In particular all the kernel sums above converge
absolutely.

For \(k\ge1\), Cauchy's estimate applied to the
uniform determinant remainder, together with the Ward
reduction of the actual Gaussian coefficient, gives
\[
 \left|\mathsf b[\Psi_k]
                   +\partial_1\partial_2\mathcal L^\sharp(0)\right|
        \le C_{\mathrm{rem}}(L)\delta_{\mathrm{loop}}^{-2}.
\]
Use the preceding scaling identity and set \(z=\eta^2\).
The massive estimate gives
\[
\begin{split}
 \eta\left|\widetilde\beta(\eta^2)
       -\eta^2\widetilde\beta'(\eta^2)
       +\tfrac12\eta^4\widetilde\beta''(\eta^2)\right|
 &\le z^{1/2}|\widetilde\beta(z)|
       +z^{3/2}|\widetilde\beta'(z)|
       +z^{5/2}|\widetilde\beta''(z)|\\
 &\le C_{\mathrm{uv}} .
\end{split}
\]
This proves the bound uniformly in \(k\); at \(k=0\)
the coefficient is zero. The supremum is therefore
finite and is the least possible bound for the
displayed accumulated sums.

\end{proof}

\section{Histories and the exact measure decomposition}
\label{sec:history-decomposition}
We decompose the effective measure into finitely many terms,
\[
 \rho_k\,dV_k=\sum_{h\in\mathcal H_k}\tau_h.
\]
A history records the failed local tests at successive scales.
The tests first select a domain on which to compare the field with
a background, then enlarge its complement until the remaining collar
passes the required bounds. Gauge factors are inserted on the selected
blocks. Their normalization, rather than a pointwise identity for the
weighted histories, gives the displayed measure identity.

Fix the odd integer $L\ge13$. Let $r,\mathfrak p_0\in\mathbb N$,
$r\ge2$, $A_0,A_1>0$, $\Theta_1\ge1$, $\varpi_T\in\mathbb N$ and
$0<\theta_c<1$. The subscript of $\theta_c$ is a mark. Take powers
of $L$ satisfying
\[
 R_1\ge L,\qquad R_1\mathfrak M\ge L^2,\qquad
 M_1\ge R_1\mathfrak M,\qquad M_2=L^2M_1,\qquad
 M\ge L^4M_1,\qquad M\ge\varpi_T.
\]
For $x_0,\ldots,x_K>1$, set $\beta=x_0$ and
\begin{equation}\label{II-1-11}
\begin{gathered}
 g_j=x_j^{-1/2},\qquad y_j=\log x_j,\qquad
 b_j=L^{-1}x_{j-1}-x_j\quad(1\le j\le K),\\
 \varepsilon_j=A_0g_jy_j^{\mathfrak p_0},\qquad
 \delta_j=A_1g_jy_j^{\mathfrak p_0},\qquad
 \delta_j^\flat=\frac{\delta_j}{4\Theta_1},\qquad
 \mathfrak b_j=\delta_j/g_j,\\
 R(x)=\min\{L^a:a\in\mathbb Z_{\ge0},\ L^a\ge(\log x)^r\},
 \qquad \check x_j=\min_{i\le j}x_i,\qquad
 \check R_j=R(\check x_j).
\end{gathered}
\end{equation}
Throughout this section assume
\begin{equation}\label{eq:history-standing}
 \varepsilon_j\le1,\qquad \delta_j\le1,\qquad
 M\check R_j\le L^{m+K-j}\qquad(0\le j\le K).
\end{equation}
For the small-coupling regime used later,
on $x\ge\gamma^{-2}$ the function $x^{-1/2}(\log x)^{\mathfrak p_0}$
decreases if $\log\gamma^{-2}>2\mathfrak p_0$; hence the first two bounds
follow from
$\max(A_0,A_1)\gamma(\log\gamma^{-2})^{\mathfrak p_0}\le1$.
The later running-coupling estimates impose a bound on $|b_j|$;
the partition below uses only \eqref{eq:history-standing}.

At step $t=k+1$, lengths are measured at level $k$.
The coarse partition $\mathcal P_t$ has side
$W_t=LM\check R_t$. A test cube $\square'$ has side
$L^2M_2\check R_t$, and its enlargement $\square'^{\sim a}$ adds
$a$ layers of cubes of side $LM_2\check R_t$.
At level $j$, a density cube has side $LM_2\check R_j$ in level-$j$
units, and its enlargement uses layers of side $M_2\check R_j$.
Thus a test cube of step $t$ is a density cube of level $t$.
The bounds above ensure that all these are standard torus partitions
and that the test partition refines $\mathcal P_t$.
For coarse regions, complements and differences are taken cube by cube;
$X^{\sim a}$ and $X^{\sim-a}$ add or remove $a$ coarse layers.
A prime or a hat denotes the $\mathcal P_t$-hull. These conventions
include the empty region and the whole torus.

We first specify the tests, including their domains. A partial map is
\emph{stably defined} at a point if it is defined and continuous on
an open neighborhood of that point. Its stable domain is open, and
its restriction there is continuous. For a density cube $\square$
at level $j\ge1$, use the reduced determining set of the tower of
$D=\square^{\sim4}$ and the datum induced by $V_j$ as in
\eqref{I-3-75}. Write
$U_{j,\square}(V_j)=U^\flat_{j,D}(V_j)$ for its reduced
minimal orbit when that orbit exists. Define
\[
 \mathsf m_{j,\square}(V_j)
   =\max_{p\text{ inner bare plaquette of }\square^{\sim}}
                   |\partial U_{j,\square}(V_j)(p)-\mathbf1|.
\]
The value is independent of the orbit representative, since plaquette
variables change by unitary conjugation. With $\eta_j=L^{-j}$, put
\begin{equation}\label{eq:history-density-test}
 \chi_j(\square)
 =\mathbf1_{\{\mathsf m_{j,\square}\text{ stably defined},\
                   \mathsf m_{j,\square}<\varepsilon_j\eta_j^2\}}.
\end{equation}
For $j=0$ use the bare plaquettes of $V_0$ in $\square^{\sim}$
and the threshold $\varepsilon_0$ directly.

For a test cube at step $t$, normalize an element of
$U_{t,\square'}(V_t)$ so that its level-$k$ contour means equal
$\mathbf1$ on
$Y_{\square'}=(\square'^{\sim3})^{(t)}$.
When the orbit exists and these means exist for one representative,
this gives a uniquely determined restriction
\[
 V^{\mathrm{bg}}_{k,\square'}
       =(\mathcal M^kU_{t,\square'}(V_t))
                    \big|_{\text{level-$k$ bonds inner to }\square'^{\sim3}}.
\]
Its conditional well-definedness is checked below.
The two cube tests are
\begin{equation}\label{eq:history-PQ-tests}
\begin{split}
 \chi_t^{\mathrm P}(\square')&=\chi_t(\square'),\\
 \chi_t^{\mathrm Q}(\square')&=
 \mathbf1_{\{V^{\mathrm{bg}}_{k,\square'}\text{ stably defined},\
 \max_{b\in(\square'^{\sim2})^{(k)*}}
 |V_k(b)V^{\mathrm{bg}}_{k,\square'}(b)^{-1}-\mathbf1|
                                      <2\delta_k\}}.
\end{split}
\end{equation}
Here $X^{(k)*}$ consists of the level-$k$ bonds with an endpoint
in $X$, with the pivot bonds of incident level-$t$ bonds removed.
Each bond tested in the second line has both endpoints in
$\square'^{\sim3}$.

An input of step $t$ is a nested family
$T_\ell\supset\Omega_1\supset\Lambda_1\supset\cdots
\supset\Omega_k\supset\Lambda_k$, where the members of index $j$
are regions for $\mathcal P_j$. The first input is empty.
For any finite collection $\mathcal Q$ of cubes and binary tests $\chi$,
write
\[
 [P;\chi]_{\mathcal Q}
   =\prod_{q\in P}(1-\chi(q))\prod_{q\in\mathcal Q\setminus P}\chi(q).
\]
Exactly one of these products is one: that of the failed cubes.

For $t\ge2$, set $Y=\Lambda_k^c$ and select the failed
$\mathrm P$-tests among test cubes in $(\widehat Y^{\sim4})^c$.
Their union is $P_t$, with indicator $\mathbf1_t^{\mathrm P}$.
Set
\begin{equation}\label{eq:history-step-domain}
 B(P_t'')=(P_t'^{\sim1}\cup\widehat Y^{\sim5})^c,
 \qquad
 \Omega_t=(Q_t'^{\sim2}\cup P_t'^{\sim4}
                              \cup\widehat Y^{\sim8})^c,
\end{equation}
where $P_t''$ is the set of level-$t$ sites of the first region,
and $Q_t$ is selected by the $\mathrm Q$-tests among test cubes in
$B(P_t'')^{\sim-1}$. Denote this second product by
$\mathbf1_t^{\mathrm Q}$.
At $t=1$, first select the failed bare density cubes $P_0$,
then the failed $\mathrm P$-tests $P_1$ in $(P_0'^{\sim})^c$;
$\mathbf1_1^{\mathrm P}$ is the product of these two indicators.
Use instead
\begin{equation}\label{eq:history-first-domain}
 B(P_1'')=(P_1'^{\sim1}\cup P_0'^{\sim2})^c,
 \qquad
 \Omega_1=(Q_1'^{\sim2}\cup P_1'^{\sim4}\cup P_0'^{\sim5})^c,
\end{equation}
with $Q_1$ selected in $B(P_1'')^{\sim-1}$ as before.
The domains are coarse regions, and $\Omega_t\subset\Lambda_k$
for $t\ge2$.

For the domains through step $t$, form the completed collared sequence
\begin{equation}\label{II-10-32}
 \Omega_0^\diamond=\Omega_1^\sim
       \supset\Omega_1\supset\cdots\supset\Omega_t,
\end{equation}
where the first enlargement uses $\mathcal P_1$. Use its reduced
determining set $\mathbf B_t^\diamond$: its level-$j$ bonds have
both endpoints in $\Omega_0^\diamond$, an endpoint in $\Omega_j$,
and no endpoint in $\Omega_{j+1}$, with $\Omega_{t+1}=\emptyset$.
The datum is the corresponding restriction of $(V_j)_{j\le t}$.
Let $U_t^\diamond$ be its minimal orbit when it exists.
Let $\mathfrak X_0$ be the level-$k$ bonds with at least one endpoint
in $\Omega_t$, and put $\mathfrak X=\Omega_t^{(k)*}$.
If the orbit exists and its level-$k$ contour means exist on the
blocks centered in $\Omega_t^{(t)}$, let
\begin{equation}\label{II-10-34}
 V_k^{\mathrm{bg}\diamond}
   =(\mathcal M^kU)\big|_{\mathfrak X_0},\qquad
 U\in U_t^\diamond,\quad
 \mathcal M^kU\text{ in block axial gauge on }\Omega_t^{(t)}.
\end{equation}
This includes the incident bonds crossing the boundary of $\Omega_t$.
The restriction is independent of the choice of $U$, as proved below.
For a positive bond $b\in\mathfrak X_0$, define
\begin{equation}\label{II-11-8}
 \mathcal A_k(b)=\frac1i\log
       [V_k(b)V_k^{\mathrm{bg}\diamond}(b)^{-1}]
\end{equation}
on the stable domain of the background, excluding the value $-\mathbf1$
inside the logarithm. The logarithm is intrinsic; extend the field by
$\mathcal A_k(\bar b)=-\mathcal A_k(b)$. Elsewhere mark it undefined.
This convention defines a partial odd field even when the background
does not exist. It is not a logarithm formula imposed again on negative
bonds.

Write $\mathfrak S=\Omega_t\setminus\Omega_t^{\sim-1}$ and let
$\mathsf K$ be the coarse cubes in $\Omega_t^{\sim-1}$.
For $c\in\mathsf K$, set
$c^\star=(c^{(k)*}\setminus\mathfrak S^{(k)*})\cap\mathfrak X$.
Define binary tests, with value zero at undefined fields, by
\[
 \chi_b=\mathbf1_{\{|\mathcal A_k(b)|<\delta_k\}},\qquad
 \chi_b^0=\mathbf1_{\{|\mathcal A_k(b)|<\delta_k^\flat\}},\qquad
 \chi_b^1=\chi_b-\chi_b^0,
 \qquad
 \operatorname{lo}(c)=\prod_{b\in c^\star}
             \mathbf1_{\{|\mathcal A_k(b)|<\theta_c\delta_k\}}.
\]
For test cubes $\square'\subset\Omega_t^{\sim-1}$, put
$\chi^{(k)}(\square')=\prod_{b\in(\square'^{\sim2})^{(k)*}}\chi_b$.
Select their failed union $R$. Its complement below is relative to
$\Omega_t^{\sim-1}$, and
\[
 \mathfrak U(R)=\bigcup_{\square'\subset R^c}
                              (\square'^{\sim2})^{(k)*}.
\]

The collar consists only of exterior bonds. For a coarse region
$\mathcal Y$, let $\mathbb B(\mathcal Y)$ be the bonds of
$\mathfrak X$ with both endpoints in $\mathcal Y$, and put
\[
 \operatorname{Col}(\mathcal Y)=
 \{b\in\mathfrak X:\mathrm e(b)\cap\mathcal Y=\emptyset,
    \ \min_{b'\in\mathbb B(\mathcal Y)}
                  \operatorname{dist}_\infty(b,b')\le\varpi_T\}.
\]
Distances between bonds are distances between their segments, and an
empty minimum is $+\infty$. Thus crossing bonds are excluded.
For a coarse set $\mathsf R$, set
$Z[\mathsf R]=(\Omega_t^c)^{\sim2}\cup\mathsf R^\sim$.
Let $\operatorname{rim}(Z)$ be the cubes of $\mathsf K$ in $Z$
which touch $Z^c$, and let $\mathsf K(I)$ consist of the cubes of
$\mathsf K$ containing an endpoint of a bond in $I$.

An answer is a pair $(I^\natural,\mathsf R^\natural)$ of subsets of
$\mathfrak X$ and $\mathsf K$. Starting with $\mathsf R_0=R'$,
explore the following sets, removing everything already explored:
\begin{equation}\label{II-4-11}
\begin{aligned}
 Z^{[n]}&=Z[\mathsf R_n],&\Lambda^{[n]}&=(Z^{[n]})^c,\\
 \mathfrak B_n&=\operatorname{Col}(\Lambda^{[n]})
                           \setminus\bigcup_{i<n}\mathfrak B_i,
 &I_n&=\mathfrak B_n\cap I^\natural,\\
 \mathfrak C_n&=\operatorname{rim}(Z^{[n]})
                           \setminus\bigcup_{i<n}\mathfrak C_i,
 &\mathsf R_n^{\mathrm{lo}}&=\mathfrak C_n\cap\mathsf R^\natural,\\
 \mathsf R_{n+1}&=\mathsf R_n\cup\mathsf K(I_n)
                                      \cup\mathsf R_n^{\mathrm{lo}}.
\end{aligned}
\end{equation}
Stop at the first $n_*$ with $I_{n_*}=\mathsf R_{n_*}^{\mathrm{lo}}
=\emptyset$. It exists, since successive nonempty answers consume
new members of finite sets. The outcome is
$\Xi=(R,I,\mathsf R^{\mathrm{lo}})$, with
$I=\bigcup_{n<n_*}I_n$ and
$\mathsf R^{\mathrm{lo}}=\bigcup_{n<n_*}\mathsf R_n^{\mathrm{lo}}$.
Its explored sets are $E^T=\bigcup_{n\le n_*}\mathfrak B_n$ and
$E^H=\bigcup_{n\le n_*}\mathfrak C_n$.
A triple is admissible if its recursion with its own answer
$(I,\mathsf R^{\mathrm{lo}})$ explores every member of $I$ and
$\mathsf R^{\mathrm{lo}}$. The output and strip are
\begin{equation}\label{II-4-27}
 Z_t=(\Omega_t^c)^{\sim2}\cup
                (R'\cup\mathsf K(I)\cup\mathsf R^{\mathrm{lo}})^\sim,
 \qquad \Lambda_t=Z_t^c,
 \qquad S_t^{\mathrm{pr}}=\mathfrak S\cup R.
\end{equation}
They satisfy $\Lambda_t\subset\Omega_t^{\sim-2}$ and
$S_t^{\mathrm{pr}}\subset\Omega_t\setminus\Lambda_t$.
Thus adjoining $\Omega_t\supset\Lambda_t$ gives the next input.

The fluctuation factor attached to an admissible triple is
\begin{equation}\label{II-11-11}
\begin{split}
 \chi^\wedge(\Xi)={}&
 \prod_{\square'\subset R}(1-\chi^{(k)}(\square'))
 \prod_{b\in\mathfrak U(R)\setminus E^T}\chi_b
 \prod_{b\in I}\chi_b^1\prod_{b\in E^T\setminus I}\chi_b^0\\
 &\times\prod_{c\in\mathsf R^{\mathrm{lo}}}(1-\operatorname{lo}(c))
       \prod_{c\in E^H\setminus\mathsf R^{\mathrm{lo}}}
                                             \operatorname{lo}(c).
\end{split}
\end{equation}
The full outcome $\mathsf o_t$ consists of $P_t,Q_t,\Xi_t$
(and $P_0$ at the first step), with indicator
$\mathbf1_t(\mathsf o_t)=\mathbf1_t^{\mathrm P}
\mathbf1_t^{\mathrm Q}\chi^\wedge(\Xi_t)$.
A history $h=(\mathsf o_1,\ldots,\mathsf o_k)\in\mathcal H_k$
successively uses the output of each outcome as its next input.
There are finitely many histories, including the empty history at
$k=0$. Write $\mathcal S(h)=(\Omega_j,\Lambda_j,S_j^{\mathrm{pr}})_{j\le k}$
for its label.

Evaluate the tests at $V_j=\mathcal M^jU$, and put
\begin{equation}\label{II-5-14}
 \mathbf1_h=\prod_{t\le k}\mathbf1_t(\mathsf o_t),\qquad
 \mathfrak G_h=\prod_{t\le k}
          \mathfrak G_{t-1}(V_{t-1};P_t''),\qquad
 \mathfrak w_h=\mathbf1_h\mathfrak G_h.
\end{equation}
For a bounded complex Borel function $\varsigma$ on $\mathcal V_0$,
define
\begin{equation}\label{II-5-15}
 \tau_h[\varsigma]=(\mathcal M^k)_*
                 (\mathfrak w_h\varsigma\,dU),
 \qquad \tau_h=\tau_h[e^{-\beta A_K}].
\end{equation}

\begin{proposition}\label{II-5-46}\label{II-5-71}\label{II-14-5}
The tests and weights above are Borel. At each bare field exactly one
history has $\mathbf1_h=1$, and
\[
 0\le\mathfrak w_h\le\mathfrak G_h
 \le\prod_{t\le k}
     \mathsf z(x_{t-1},\varepsilon_{t-1})^{-(L^3-1)\#P_t''}.
\]
Consequently every $\tau_h[\varsigma]$ is a finite complex measure,
positive when $\varsigma\ge0$.
If $\varepsilon_j\le c_2^{\mathrm{SO}}$ for $0\le j<k$ and
$\varsigma$ is gauge invariant, then
\begin{equation}\label{II-5-73}
 \sum_{h\in\mathcal H_k}\tau_h[\varsigma]
            =(\mathcal M^k)_*(\varsigma\,dU).
\end{equation}
In particular,
\[
 \rho_k\,dV_k=\sum_{h\in\mathcal H_k}\tau_h,
 \qquad \sum_{h\in\mathcal H_k}\tau_h(\mathcal V_k)
                                      =Z_{K,\ell,\beta}.
\]
Here $\rho_k\,dV_k$ denotes the effective measure of \cref{I-4-67};
no absolute-continuity assertion is implicit in this notation.
\end{proposition}

\subsection{Domains of the tests and the collar exploration}
\label{subsec:history-tests-proof}
We first justify the conditional axial normalizations used above.
Suppose an orbit for a reduced determining set is given, and let
$Y$ be a set of top sites. No site of $B(Y)\setminus Y$ is a node.
If the contour means of $\mathcal M^kU$ exist, set
\[
 u(x)=(\mathcal M^kU)(y,x)
       \quad(x\in B(y)\setminus\{y\},\ y\in Y),
 \qquad u=\mathbf1\text{ at every other bare site}.
\]
Disjointness of the blocks makes this consistent, and $u$ fixes every
node. Covariance of the contour mean gives
$ (\mathcal M^kU^u)(y,x)=\mathbf1$.
Conversely, if $U'=U^v$ and both representatives are axial on $Y$,
then $v(y)=\mathbf1$ at the top nodes and
$\mathbf1=v(y)\mathbf1v(x)^{-1}$ on each block. Thus $v=\mathbf1$
on $B(Y)$ and at all other nodes. The averaged representatives agree
on every defined bond whose endpoints lie in this union.

For the tower of $\square'^{\sim4}$, the margin
$M_2\check R_t\ge L^2M_1$ puts $\square'^{\sim3}$ in the top
member: each of its $M_1$-cubes is at distance greater than $M_1$
from the complement, and hence lies in the two-layer interior used
to form the tower. Its top sites are nodes by
\cref{lem:tower-top-nodes}. Block membership then puts all required
bare blocks inside $\square'^{\sim4}$. The preceding argument
defines $V^{\mathrm{bg}}_{k,\square'}$ uniquely. A bond with an
endpoint in $\square'^{\sim2}$ has its other endpoint in
$\square'^{\sim3}$, since its length is $1<LM_2\check R_t$.
This also verifies the bond domain in \eqref{eq:history-PQ-tests}.

For \eqref{II-10-34}, a site within $2L$ of $\Omega_t$ lies in
$\Lambda_k$ when $t\ge2$, and in $\Omega_0^\diamond$ when $t=1$.
Indeed, \eqref{eq:history-step-domain} separates $\Omega_t$ from
$\Lambda_k^c$ by at least $8W_t>2L$, while the first collar has
width $W_1>2L$. Consequently the bare blocks needed for the level-$k$
averages on $\mathfrak X_0$ lie in $\Omega_0^\diamond$.
Each site of $Y=\Omega_t^{(t)}$ has a neighbor inside the same
coarse cube, so it is the endpoint of a top datum bond and is a node.
The noncentral sites of $B(Y)$ are not nodes of the reduced set.
Finally, an endpoint $x\notin\Omega_t$ of a bond of $\mathfrak X_0$
has a level-$k$ neighbor $x'$ in its own coarse cube: in the centered
integer coordinates of that cube, change the first coordinate by one
toward its interior. Both $x,x'$ lie within $2$ of $\Omega_t$ and
outside it. The bond $\langle x,x'\rangle$ therefore belongs to the
level-$k$ member of the reduced datum, so $x$ is a node.
Thus every endpoint of $\mathfrak X_0$ lies in $B(Y)$ or is a node,
and the same normalization argument proves uniqueness on all of
$\mathfrak X_0$, including the crossing bonds. If $\Omega_t$ is empty,
the assertion concerns the empty restriction.

The stable domain convention now gives the Borel properties without
requiring backgrounds at arbitrary fields. The stable restriction
of each background or of $\mathsf m_{j,\square}$ is continuous.
For \eqref{II-11-8}, multiplication and inversion are continuous,
$G\setminus\{-\mathbf1\}$ is open, and the intrinsic logarithm
is continuous there. The passing sets of the density and background
tests, of $\chi_b$, $\chi_b^0$, $\operatorname{lo}(c)$ and of the
cube test $\chi^{(k)}(\square')$ are therefore open in the finite
product of field spaces. The indicator $\chi_b^1=\chi_b-\chi_b^0$
is Borel as a difference of Borel functions. Complements and finite
products of these indicators are Borel as well.

We will also need their covariance. Let $v$ be a gauge transformation
at level $n$ and extend it to the bare lattice by making it constant
on each order-$n$ block; denote its restrictions by $s_j$.
The covariance of the actual averages, including their domains,
gives $\mathcal M^j(U^{\tilde v})=(\mathcal M^jU)^{s_j}$.
For a test of step $t\le n$, $s_k$ is constant on every block
centered at level $t$. Gauge transformation maps the defining critical
orbit to that of the transformed datum: it preserves the constraints,
criticality and the curvature and current bounds, and conjugates the
residual groups. It preserves the axial condition on these blocks.
Uniqueness of the restrictions just proved therefore makes the
background transform as a configuration. The fluctuation logarithm
transforms by $\operatorname{Ad}s_k(b_-)$, with its domain unchanged.
The same applies to the tower backgrounds, and plaquette norms are
unchanged. The gauge transformations are homeomorphisms, so they also
preserve the stable domains. All outcome tests of step $t\le n$ are
therefore invariant. The earlier gauge factors are invariant as well,
since their contour means are conjugated by a common block value.
The plaquette part of step $n+1$ is invariant: it depends only on
$V_{n+1}$ (and on $V_0$ through conjugation when $n=0$).

We next check the only geometric point needed in the fluctuation
split. For any stage of the exploration, $R'\subset\mathsf R_n$.
If a site $z$ is within distance strictly less than $W_t$ of
$\Lambda^{[n]}$, its coarse cube touches that region. This cube
belongs to $\Omega_t^{\sim-1}$ and not to $\mathsf R_n$:
otherwise its touching cube would be in
$(\Omega_t^c)^{\sim2}\cup\mathsf R_n^\sim$.
The test cube containing $z$ is consequently outside $R$, and all
incident bonds belong to its two-layer enlargement. In particular,
they belong to $\mathfrak U(R)$. If $b$ belongs to the collar,
a nearest inner bond of $\Lambda^{[n]}$ has its whole segment in
that region. Each endpoint of $b$ is therefore within
$\varpi_T+1<W_t$ of $\Lambda^{[n]}$. This proves
\begin{equation}\label{eq:history-explored-small}
 E^T\subset\mathfrak U(R).
\end{equation}
The elementary grid facts used here remain valid on the torus:
in centered cube coordinates, touching means that each index differs
by at most one, and an inner bond stays in its containing coarse
region. The bounds on $M,M_2$ ensure that a one-bond extension lies
inside the indicated two test layers.

\begin{proof}[Proof of the partition and Borel assertions in
\cref{II-5-46}]
Fix the fields after the choices of $P_t,Q_t$. Run
\eqref{II-4-11} with $R$ the failed cube set and answer
\[
 I^\natural=\{b:\chi_b^1=1\},\qquad
 \mathsf R^\natural=\{c:\operatorname{lo}(c)=0\}.
\]
At each stage the recursion asks only for the answer on the newly
explored sets. Replacing the answer by the union of its discovered
parts therefore leaves every stage unchanged, by induction up to
the stopping time. Its outcome is admissible.
Every factor in \eqref{II-11-11} for this outcome is one.
For the factors $\chi_b^0$ this uses
\eqref{eq:history-explored-small}: the collar bonds lie in passed
cube tests, hence $\chi_b=1$, and absence from $I$ gives
$\chi_b^1=0$, so $\chi_b^0=1$.

Conversely, if an admissible triple has $\chi^\wedge=1$, the
factors on $R$ force its cubes to fail. Every bond tested by a cube
outside $R$ has $\chi_b=1$, either from its unexamined factor or
from $\chi_b^0=1$ or $\chi_b^1=1$ on $E^T$.
Thus $R$ is precisely the failed cube set. The remaining factors
give the actual answer on $E^T$ and $E^H$. Induction through the
recursion identifies this triple with the outcome just constructed.
Hence $\sum_\Xi\chi^\wedge(\Xi)=1$, with exactly one nonzero term.
The elementary failed-set identities for $P_0,P_t,Q_t$ then give
one outcome, and induction over the steps gives one history.

All index sets are finite and fixed by the chosen outcome.
The Borel assertions proved above therefore apply to each indicator
and weight. Finally, each factor of \eqref{I-3-68} is between zero
and $\mathsf z^{-1}$ and each block has $L^3-1$ noncentral sites.
This proves the stated bound. Finite pushforwards of the resulting
bounded densities give the claimed measures.
\end{proof}

\subsection{Gauge normalization on the selected histories}
\label{subsec:history-measure-proof}
A passed density test controls the plaquettes of its block field.
Precisely, if $\chi_j(\square)=1$ and
$\varepsilon_j\le c_2^{\mathrm{SO}}$, then
\begin{equation}\label{II-5-41}
 |\partial V_j(p)-\mathbf1|<2\varepsilon_j
 \quad\text{when }p\text{ has a corner in }\square.
\end{equation}
For $j=0$, the stronger bound $\varepsilon_0$ follows directly from
the test: a one-bond extension of $\square$ lies in
$\square^{\sim}$. For $j\ge1$, all corners of such a plaquette lie
in $\square^{\sim}$, and its bonds are top datum bonds of the tower
of $\square^{\sim4}$. Thus its block holonomy is that of
$\mathcal M^jU$ for an element $U$ of the tested orbit. All bare
blocks over its four corners lie in $\square^{\sim}$.

Here the curvature assumption is local, so we spell out its use in
the averaging argument. If $\mathsf Z$ is this union of four bare
blocks, call a level-$n$ bond or plaquette covered when the bare
blocks over its vertices lie in $\mathsf Z$.
Extend $U$ by $\mathbf1$ outside its original bond domain. The
actual averages on covered bonds are unchanged, by locality, and
are restrictions of a real configuration at every level.
For a covered level-$(n+1)$ plaquette, every level-$n$ plaquette
in its four blocks is covered. Thus \cref{I-10-97} applies at each
stage and gives, for the maximal covered curvature,
\[
 \gamma_{n+1}\le L^2\gamma_n+C_0^{\mathrm{SO}}\gamma_n^2.
\]
The small-curvature branch agrees with the actual average there;
reversed bonds carry the inverse value. If
$\gamma_0\le\alpha L^{-2j}$ with $0<\alpha\le c_2^{\mathrm{SO}}$,
the normalized induction of \cref{subsec:averaging-iteration} gives
\[
 \gamma_n\le\alpha L^{2(n-j)}
       \prod_{l<n}(1+\tfrac12L^{2(l-j)})
       <2\alpha L^{2(n-j)}\qquad(0\le n\le j).
\]
Its hypothesis is preserved because
$2c_2^{\mathrm{SO}}\le\varpi'_L<\omega_M$ and
$C_0^{\mathrm{SO}}c_2^{\mathrm{SO}}\le L^2/4$.
Choose $\alpha<\varepsilon_j$ above the finite maximum
$L^{2j}\mathsf m_{j,\square}(V_j)$. This proves
\eqref{II-5-41} without a global curvature assumption on $U$.

On the support of a plaquette part of step $j$, every test cube
inside $\Omega_j^{\sim4}$ passes. Indeed,
\eqref{eq:history-step-domain} and
\eqref{eq:history-first-domain} put this enlargement in the region
where the $\mathrm P$-tests were selected, and outside their failed
hull. Hence \eqref{II-5-41} holds at every plaquette with a corner
in $\Omega_j^{\sim4}$. At the first level the bare tests likewise
give the bound $\varepsilon_0$ at every plaquette with a corner
in $P_0^c$.

Fix a history $h'$ through step $t-1$ and a plaquette part of step
$t$. Let $y\in P_t''$ and consider the double block of the
level-$t$ bond $\langle y,y+Le_1\rangle$ in level-$(t-1)$ units.
Its sites lie within $L+(L-1)/2<2L<W_t$ of $y$.
For $t\ge2$, \eqref{eq:history-step-domain} puts $y$ at distance
at least $5W_t$ from the hull of $\Lambda_{t-1}^c$, so the double
block lies in $\Lambda_{t-1}\subset\Omega_{t-1}^{\sim4}$.
For $t=1$, \eqref{eq:history-first-domain} instead puts it in
$P_0^c$. On the support of
$\mathbf1_{h'}\mathbf1_t^{\mathrm P}$, its plaquette curvatures are
therefore less than $2\varepsilon_{t-1}\le\omega_M$.
The small-curvature contour estimate \cref{I-3-37} shows that every
$V_{t-1}(y,x)$ used by $\mathfrak G_{t-1}(\,\cdot\,;P_t'')$ exists.
This includes $P_t''=\emptyset$, when the gauge factor is one.

\begin{proof}[Proof of \eqref{II-5-73}]
For fixed $h'$ and plaquette part $\mathsf p$, set
$P''=P_t''(\mathsf p)$ and
$e_{\mathsf p}=\mathfrak w_{h'}\mathbf1_t^{\mathrm P}(\mathsf p)$,
with all fields evaluated at their actual bare iterates.
Summing the $\mathrm Q$ and fluctuation outcomes gives
\[
 \sum_{\mathsf o_t:\,\mathsf p(\mathsf o_t)=\mathsf p}
       \mathfrak w_{(h',\mathsf o_t)}
       =e_{\mathsf p}\mathfrak G_{t-1}(V_{t-1};P'').
\]
Let $v$ vary independently on $B(P'')\setminus P''$ at level
$t-1$ and be the identity elsewhere. Extend it constantly on the
order-$(t-1)$ bare blocks. Its restriction to level $t$ is the
identity. Consequently $V_t$ is unchanged, $V_{t-1}$ transforms
by $v$, and the covariance proved above leaves $e_{\mathsf p}$
unchanged. The function $\varsigma$ is unchanged by its assumed
gauge invariance.

For every bounded Borel function $F$ on $\mathcal V_t$, Haar
invariance on the bare bonds and Fubini's theorem now give
\begin{equation}\label{II-5-65}
\begin{split}
 &\int F(V_t)e_{\mathsf p}\,
          \mathfrak G_{t-1}(V_{t-1};P'')\varsigma\,dU\\
 &\qquad=\int F(V_t)e_{\mathsf p}\varsigma
       \left[\int\mathfrak G_{t-1}(V_{t-1}^{v};P'')\,dv\right]dU
       =\int F(V_t)e_{\mathsf p}\varsigma\,dU.
\end{split}
\end{equation}
The integrands are bounded Borel functions on finite products.
Where $e_{\mathsf p}\ne0$, the domain check just proved permits
\eqref{I-3-71}, which makes the bracket one; elsewhere the outside
factor vanishes. Summing \eqref{II-5-65} over the plaquette parts
and using their partition of one yields
\[
 \sum_{h:\,h|_{t-1}=h'}\tau_h[\varsigma]
       =\mathbf T\tau_{h'}[\varsigma].
\]
Starting with $\tau_\emptyset[\varsigma]=\varsigma\,dU$ proves
\eqref{II-5-73} by induction. Apply it to
$\varsigma=e^{-\beta A_K}$ and then take total masses, using
\cref{I-4-67}, to obtain the remaining assertions.
\end{proof}

\subsection{Regular backgrounds and the cost of a failed test}
\label{subsec:failed-test-cost}

We use the cubes and tests of \eqref{eq:history-density-test}
and \eqref{eq:history-PQ-tests}.
The point of the comparison below is that failure of a test costs either
plaquette action or gauge-fixing energy, even when the background in the
definition of the test is undefined. All estimates are local to the
four-layer enlargement of the test cube.

Choose a power $\rho_{\mathrm s}\ge\rho_{\mathrm G}$ of $L$ and impose
the following additional restrictions on the constants of the history
construction:
\begin{equation}\label{eq:test-background-constants}
\begin{gathered}
 \mathfrak M\ge\mathfrak M_{\mathrm G},\qquad
 R_1\mathfrak M\ge L^2\rho_{\mathrm s},\qquad
 B_3\ge B_{\mathrm{rad}},\\
 \varepsilon_{\mathrm v}+\varepsilon_{\mathrm v}^2\le a_{\rm orb},
 \qquad 0<\bar\varepsilon\le\varepsilon_{\mathrm G}(\rho_{\mathrm s}),
 \qquad 0<r_B\le r_{\mathrm G}(\rho_{\mathrm s}),\\
 B_3(\bar\varepsilon+C_\partial r_B)\le\varepsilon_{\mathrm v}.
\end{gathered}
\end{equation}
Here $\varepsilon_{\mathrm v}$ is the radius in the definition of the
regular critical orbit. These choices are possible: choose the geometric
parameters and $B_3$ first, then $\varepsilon_{\mathrm v}$, and finally
the two datum radii. They are independent of the lattice and the number
of steps.

\begin{proposition}[Regularity of localized backgrounds]\label{II-7-40}
Let $1\le t\le K$, $k=t-1$, and let $\square'$ be a test cube.
Put
\[
 X=\square'^{\sim4},\qquad X_3=\square'^{\sim3},\qquad
 Y=X_3^{(t)}.
\]
Suppose \eqref{eq:test-background-constants} and the standing history
conditions hold. If $0<\varepsilon_{\mathrm r}\le\bar\varepsilon$ and
\begin{equation}\label{eq:test-coarse-regularity}
 |\partial V_t(p')-\mathbf1|<\varepsilon_{\mathrm r}
 \quad\text{at every inner level-$t$ plaquette of $X$},
\end{equation}
then the orbit $U_{t,\square'}(V_t)$ exists. Every member is regular at
radius $B_3\varepsilon_{\mathrm r}$, and
\begin{equation}\label{eq:test-density-background-bound}
 \mathsf m_{t,\square'}(V_t)
       <B_3\varepsilon_{\mathrm r}\eta_t^2.
\end{equation}
The function on the left is stably defined at every such $V_t$.

If also $B_3\varepsilon_{\mathrm r}\le c^{\mathrm{SO}}_2$, then
$V^{\mathrm{bg}}_{k,\square'}$ is stably defined there, is in block
axial gauge on $Y$, and satisfies
\begin{equation}\label{eq:test-background-curvature}
 |\partial V^{\mathrm{bg}}_{k,\square'}(p)-\mathbf1|
       \le2B_3\varepsilon_{\mathrm r}L^{-2}
 \quad\text{for all level-$k$ plaquettes with corners in $X_3$}.
\end{equation}
For every inner level-$t$ bond $c$ of $X_3$, its average uses the
first branch and
\[
 \mathcal M(V^{\mathrm{bg}}_{k,\square'})(c)=V_t(c).
\]
In particular, the $\mathrm P$-test passes if
$B_3\varepsilon_{\mathrm r}\le\varepsilon_t$; under the additional
smallness condition, the $\mathrm Q$-test is the strict bond inequality
in \eqref{eq:history-PQ-tests}.
\end{proposition}

\begin{proof}
First consider induced data on any separated tower of depth $t$ whose
initial region is contained in $X$. The lift in \cref{I-3-72} has value
$\mathbf1$ on bonds within an order-$(t-j)$ block and value $V_t(c)$
on bonds crossing over $c$ at level $j$. On a level-$j$ plaquette,
crossing a block boundary in one coordinate is unaffected by a step
in the other coordinate. Thus either the four block centers are
distinct and the product is a level-$t$ plaquette variable with the
same cyclic order and base point, or the factors cancel and the
product is $\mathbf1$. All these coarse plaquettes lie in $X$.
At a tested plaquette, an edge outside the determining member uses
the average of the preceding member; by \cref{I-3-72} this is the
same lifted value. Hence every tested holonomy of the induced datum
is $\mathbf1$ or an inner plaquette variable of $V_t$ in $X$.
The datum is regular at radius $\varepsilon_{\mathrm r}$.
Its coordinates are copies of $\mathbf1$ or of coordinates of $V_t$,
so the map to induced data is continuous.

The tower of $X$ is separated. To check the remaining geometric
conditions, measure lengths at level $t$ and put
$w=M_2\check R_t$. Then $w\ge L^2M_1$ and
$M\check R_t\le L^{m+K-t}$, while $M\ge L^4M_1$.
In particular the torus has at least $6R_1\mathfrak M$ sites in
each direction. The gap between $X_3$ and $X^c$ is at least $w$;
\cref{prop:explicit-tower} therefore gives $X_3\subset\Omega_t$.
This remains true if an enlargement is the whole torus.
Block nesting gives
\[
 X_3^{(j)}=B^{t-j}(Y),\qquad 0\le j\le t.
\]
Every $y\in Y$ has a nearest neighbor in the same $w$-cube, since
$w\ge L$. Thus it is an endpoint of an inner level-$t$ bond of
$X_3$, which belongs to the top determining member. These sites are
nodes, and the sites of $B(Y)\setminus Y$ are not nodes by
\cref{I-3-81}.

Apply \cref{I-58-22} to the induced reduced datum, at
$\rho=\rho_{\mathrm s}$. There are no pinch rows in the reduced
setting. The radius comparison following its proof identifies its
orbit with the full regular critical set at radius
$\varepsilon_{\mathrm v}$: the constructed orbit lies in
$\mathfrak U'_t(B_{\mathrm{rad}}\varepsilon_{\mathrm r})
\subset\mathfrak U_t(\varepsilon_{\mathrm v})$, while
\[
 \mathfrak U_t(\varepsilon_{\mathrm v})
 \subset\mathfrak U'_t(\varepsilon_{\mathrm v}
                          +\varepsilon_{\mathrm v}^2)
 \subset\mathfrak U'_t(a_{\rm orb}).
\]
All members of this orbit have the stated linear regularity bound.
For bare plaquettes with corners in $X_3$, the local block side in
level-$t$ units is one. In particular their curvature is less than
$B_3\varepsilon_{\mathrm r}L^{-2t}$, proving
\eqref{eq:test-density-background-bound}.

The existence assertion also supplies local continuous representatives.
Indeed, near any induced datum $V_0$, write
$V'=e^{iB(V')}V_0$ on positive rows, where
$B(V')=i^{-1}\log(V'V_0^{-1})$ is the logarithm near $\mathbf1$.
After restricting the neighborhood, $\sup|B(V')|<r_B$.
The map $V'\mapsto U^\circ(B(V'))$ is continuous by
\cref{I-58-22}, and represents the regular orbit by
\eqref{eq:test-background-constants}. Composing with the continuous
map to induced data gives local selections over $V_t$. The coarse condition
\eqref{eq:test-coarse-regularity} is open. A finite maximum of
gauge-invariant plaquette norms evaluated on these selections is
continuous and independent of the selection. This proves stable
definedness of $\mathsf m_{t,\square'}$.

For the asserted level-$k$ bound, put
$\alpha=B_3\varepsilon_{\mathrm r}L^{-2}\le c^{\mathrm{SO}}_2$.
On $B^t(Y)$ the bare curvature is less than $\alpha L^{-2k}$.
The iteration proof of \cref{I-10-101} applies on these nested block
unions: every double block and every four-block plaquette cone at
the next level is contained in the preceding union. More explicitly,
if $\varpi_j$ is the maximum curvature there, \cref{I-10-97} gives
\[
 \varpi_{j+1}\le L^2\varpi_j+C^{\mathrm{SO}}_0\varpi_j^2,
 \qquad
 \varpi_j\le\alpha L^{2(j-k)}
          \prod_{a<j}(1+\tfrac12L^{2(a-k)})
          <2\alpha L^{2(j-k)}.
\]
The product is at most
$\exp(\frac12\sum_{a=0}^{k-1}L^{2(a-k)})<2$.
The definition of $c^{\mathrm{SO}}_2$ closes the induction and keeps
each participating average in the first branch. Consequently
$W=\mathcal M^kU$ has curvature at most $2\alpha$ on $B(Y)$.
For $k=0$, this statement follows directly from the bare bound.
Each $y\in Y$ is an endpoint of a double block in $B(Y)$, so
\cref{I-10-61} defines every contour mean $W(y,x)$ used below.

Normalize by the residual gauge transformation equal to $W(y,x)$
at $x\in B(y)\setminus\{y\}$, $y\in Y$, and equal to $\mathbf1$
elsewhere. On inner bonds of $X_3$ this gives
\begin{equation}\label{II-7-42}
 V^{\mathrm{bg}}_{k,\square'}(b)
 =W([b_-]_t,b_-)\,W(b)\,W([b_+]_t,b_+)^{-1}.
\end{equation}
The conditional uniqueness of the normalized background identifies
this with the field in \eqref{eq:history-PQ-tests}. Gauge covariance
preserves its curvature and its coarse average; the top constraint
sets the latter equal to $V_t$.
On each local continuous selection, all iterated averages and contour
means used in \eqref{II-7-42} belong to their analytic small-curvature
domains. They are therefore continuous. Formula \eqref{II-7-42}
and independence of the selection prove stable definedness.
The remaining assertions follow from the definitions of the tests.
\end{proof}

The comparison with a field satisfying the same coarse constraint
uses the following reconstruction estimate. It is useful to retain
the contour error explicitly.

\begin{lemma}[Reconstruction on a double block]\label{I-25-42}
Set $\kappa_{\mathsf c}=4\kappa^{\mathrm{SO}}_1+
\kappa^{\mathrm{SO}}_2+\kappa_3$. Let $c$ be a positive level-$(i+1)$
bond, $0\le i<K$, and let $W,W_0$ be real configurations on
$\mathsf P_c$ with
$\max\{\varpi_c(W),\varpi_c(W_0)\}\le\varpi\le\varpi_L$.
If
\begin{gather*}
 |W(y,x)W_0(y,x)^{-1}-\mathbf1|\le\zeta
       \quad(y=c_-,c_+,\ x\in B(y)),\\
 |\mathcal M(W)(c)\mathcal M(W_0)(c)^{-1}-\mathbf1|\le a,
\end{gather*}
then
\begin{equation}\label{I-25-45}
 |W(b)W_0(b)^{-1}-\mathbf1|
       \le a+2\kappa_{\mathsf c}\varpi+2\zeta,
       \qquad b\in\mathsf P_c.
\end{equation}
\end{lemma}

\begin{proof}
Use level-$i$ lengths. The smallness of $\varpi_L$ permits \cref{I-10-61}.
For an inner bond $b$ of $B(y)$, orient $b$ away from $y$ in
its coordinate direction. Choose a rectilinear path to $b_+$
which traverses that direction last; its last step is $b$.
Apply the contour estimate to that path and to the path with its
last step removed. Unitary multiplication preserves the norm, so
\[
 |W(b)-W(y,b_-)^{-1}W(y,b_+)|
       \le2\kappa^{\mathrm{SO}}_1\varpi.
\]
For a corridor bond $b=\langle x,x+e_\mu\rangle$, write
$x=c_-+u+\frac12(L-1)e_\mu$, with $u\perp e_\mu$, and put
$x_\perp=c_-+u$. The path from $c_-$ to $x$ that makes the
$\mu$ steps last, followed by $b$ and the reversed corresponding
path from $c_+$, is the transverse path to $x_\perp$, the straight
length-$L$ segment, and the reversed translated transverse path.
Replacing the endpoint paths by their contour means costs
$4\kappa^{\mathrm{SO}}_1\varpi$. The relative-loop estimate and
the average-versus-straight-transport estimate in \cref{I-10-61}
then give
\[
 |W(c_-,b_-)W(b)W(c_+,b_+)^{-1}-\mathcal M(W)(c)|
       \le\kappa_{\mathsf c}\varpi.
\]
Apply these two estimates also to $W_0$. Subtracting the reconstructed
products costs at most $2\zeta$ inside a block and at most
$a+2\zeta$ in the corridor. This proves respectively
$4\kappa^{\mathrm{SO}}_1\varpi+2\zeta$ and the bound
\eqref{I-25-45}. Every bond or its reverse is of one of these
types, and the latter bound dominates the former.
\end{proof}

\begin{proposition}[Cost of a failed plaquette test]\label{II-9-27}
Assume the setting of \cref{II-7-40}, and suppose
$\varepsilon_t\le B_3\bar\varepsilon$. If
$\chi_t^{\mathrm P}(\square')=0$, then some inner level-$t$
plaquette of $X=\square'^{\sim4}$ satisfies
$|\partial V_t-\mathbf1|\ge\varepsilon_t/B_3$.
For every bare configuration $U'$ on the inner bonds of $X$ with
$\mathcal M^tU'=V_t$ on its inner level-$t$ bonds,
\begin{equation}\label{II-9-28}
 L^tA_X(U')\ge\min\{c_1^{\mathrm P}\varepsilon_t^2,e_*\},
 \qquad c_1^{\mathrm P}=\tfrac14B_3^{-2}.
\end{equation}
At level zero a failed density test gives
$A_{\square^{\sim}}(V_0)\ge\frac12\varepsilon_0^2$.
\end{proposition}

\begin{proof}
If all coarse curvatures were less than $\varepsilon_t/B_3$,
\cref{II-7-40} would make the test pass. Thus the stated plaquette
exists. Extend $U'$ by $\mathbf1$ on other bare bonds.
The full bare cone under this plaquette lies in $X$, since $X$
is a union of level-$t$ cubes. Locality identifies its iterated
average with that of $U'$. Apply \cref{II-6-57} to this one
plaquette; its cone action is at most $A_X(U')$.
At level zero use $a_p(V_0)=\frac12|\partial V_0(p)-\mathbf1|^2$.
\end{proof}

For the second test define
\begin{equation}\label{II-9-10}
\begin{gathered}
 \varepsilon_{\mathrm{bg}}=
    \min\{\bar\varepsilon,c^{\mathrm{SO}}_2/B_3\},\qquad
 \widetilde\sigma_{\mathrm Q}=\frac4{5L^2}\varepsilon_{\mathrm{bg}},
 \qquad \widetilde c_2^{\mathrm Q}=\frac2{5\kappa_{\mathsf c}B_3+2},\\
 c_2^{\mathrm Q}=\tfrac14(\widetilde c_2^{\mathrm Q})^2,
 \qquad
 \sigma_{\mathrm Q}=\min\{\tfrac14\widetilde\sigma_{\mathrm Q}^2,e_*\}.
\end{gathered}
\end{equation}
These constants depend only on $L,B_3,\bar\varepsilon$.
Extend the gauge-fixing function by
\begin{equation}\label{II-9-3}
 G_k^\infty(Y;V_k)=
 \begin{cases}
 \displaystyle\frac12\sum_{y\in Y}\sum_{x\in B(y)\setminus\{y\}}
                  |V_k(y,x)-\mathbf1|^2,
       &\text{if all these contour means are defined},\\
 +\infty,&\text{otherwise}.
 \end{cases}
\end{equation}

\begin{proposition}[Cost of a failed fluctuation test]\label{II-9-33}
Assume the setting of \cref{II-7-40}, and impose the local tie
\begin{equation}\label{eq:failed-test-tie}
 V_t(c)=\mathcal M(V_k)(c)
 \quad\text{on every inner level-$t$ bond of $X$}.
\end{equation}
If $\chi_t^{\mathrm Q}(\square')=0$, set
$a_{\mathrm Q}=\min\{\widetilde c_2^{\mathrm Q}\delta_k,
\widetilde\sigma_{\mathrm Q}\}$. Then either a contour mean
$V_k(y,x)$, $y\in Y$, $x\in B(y)$, is undefined or has distance
at least $a_{\mathrm Q}$ from $\mathbf1$, or an inner level-$k$
plaquette of $X$ has curvature at least $a_{\mathrm Q}$.
For every bare configuration $U'$ on the inner bonds of $X$ such
that $\mathcal M^kU'=V_k$ on its inner level-$k$ bonds,
\begin{equation}\label{II-9-34}
 L^kA_X(U')+G_k^\infty(Y;V_k)
       \ge\min\{c_2^{\mathrm Q}\delta_k^2,\sigma_{\mathrm Q}\}.
\end{equation}
This includes $k=0$, with $\mathcal M^0$ the identity.
\end{proposition}

\begin{proof}
We first establish the comparison that makes the contrapositive
quantitative. Suppose all inner level-$k$ plaquettes of $X$ have
curvature at most $\varpi_0$, where
$0<\varpi_0\le\widetilde\sigma_{\mathrm Q}$, and suppose
$|V_k(y,x)-\mathbf1|\le\zeta$ on the indicated blocks.
The four-block cone under every inner coarse plaquette lies in $X$.
Since $C^{\mathrm{SO}}_0\widetilde\sigma_{\mathrm Q}\le1/5$,
\cref{I-10-97} and \eqref{eq:failed-test-tie} give
\[
 |\partial V_t(p')-\mathbf1|
 \le(L^2+\tfrac15)\varpi_0
 <\tfrac54L^2\varpi_0=:\varepsilon_{\mathrm r}.
\]
Here $\varepsilon_{\mathrm r}\le\varepsilon_{\mathrm{bg}}$.
By \cref{II-7-40} the background is stably defined, has the same
average as $V_k$ on inner coarse bonds of $X_3$, is in block axial
gauge on $Y$, and has curvature at most
$\frac52B_3\varpi_0$ there. Moreover
\[
 \tfrac52B_3\varpi_0
 \le2c^{\mathrm{SO}}_2L^{-2}<\varpi_L.
\]
Thus \cref{I-25-42} applies to $V_k$ and this background on each
double block of an inner coarse bond of $X_3$, with $a=0$.
Their contour defect is at most $\zeta$, so
\begin{equation}\label{II-9-32}
 |V_k(b)V^{\mathrm{bg}}_{k,\square'}(b)^{-1}-\mathbf1|
       \le5\kappa_{\mathsf c}B_3\varpi_0+2\zeta
\end{equation}
for every inner level-$k$ bond of $X_3$.
Indeed, a bond whose endpoints are in one block is covered by
any inner coarse bond incident to its center; such a bond exists
within the same $w$-cube. Otherwise its two block centers are
nearest neighbors in $Y$, and their double block covers it.
In particular \eqref{II-9-32} covers all bonds tested in
$\chi_t^{\mathrm Q}(\square')$.

If neither alternative in the statement holds, all the finite
maxima of curvature and contour error are strictly below
$a_{\mathrm Q}$. Choose $0<\varpi_0<a_{\mathrm Q}$ above the
curvature maximum and let $\zeta<a_{\mathrm Q}$ be the contour
maximum. Then
\[
 5\kappa_{\mathsf c}B_3\varpi_0+2\zeta
 <(5\kappa_{\mathsf c}B_3+2)a_{\mathrm Q}\le2\delta_k.
\]
Stable definedness and \eqref{II-9-32} make the test pass,
a contradiction. This proves the alternative.

An undefined contour gives infinite gauge-fixing energy; a contour
of size at least $a_{\mathrm Q}$ gives energy at least
$\frac12a_{\mathrm Q}^2$. In the plaquette alternative, extend
$U'$ as in the preceding proof and use \cref{II-6-57} at level
$k\ge1$, obtaining action at least
$\min\{\frac14a_{\mathrm Q}^2,e_*\}$ after multiplication by $L^k$.
For $k=0$ the individual plaquette action is at least
$\frac12a_{\mathrm Q}^2$, which implies the same lower bound.
Both energies are nonnegative. Finally,
\[
 \min\{\tfrac14a_{\mathrm Q}^2,e_*\}
 =\min\{c_2^{\mathrm Q}\delta_k^2,\sigma_{\mathrm Q}\},
\]
proving \eqref{II-9-34}.
\end{proof}

Under the ceilings $c_1^{\mathrm P}\varepsilon_t^2\le e_*$ and
$c_2^{\mathrm Q}\delta_k^2\le\sigma_{\mathrm Q}$, the minima in
\eqref{II-9-28} and \eqref{II-9-34} take their first values. If
$|b_t|\le B_\sharp$, the identities
$L^{-1}x_{t-1}=x_t+b_t$ and
$x_j\varepsilon_j^2=A_0^2y_j^{2\mathfrak p_0}$ give
\begin{align}
 L^{t-1}x_{t-1}A_X(U')
 &\ge c_1^{\mathrm P}(1-B_\sharp g_t^2)
                A_0^2y_t^{2\mathfrak p_0},\label{II-9-29}\\
 x_k\bigl(L^kA_X(U')+G_k^\infty(Y;V_k)\bigr)
 &\ge c_2^{\mathrm Q}A_1^2y_k^{2\mathfrak p_0}.\label{II-9-35}
\end{align}
Each inequality retains the corresponding failure and averaging
hypotheses of the preceding propositions.
These costs are supported in the one-layer enlargement of the
coarse cell $\mathfrak c''$ containing $\square'$. Indeed, in
level-$t$ units, every point of $X$ is within $4M_2\check R_t$
of $\square'\subset\mathfrak c''$, and
$4M_2\check R_t\le M\check R_t$.


\subsection{Profiles and the action reserved for large fields}
\label{sec:history-profiles}

The local changes of coupling are supported inside the regular
regions. Their smooth profiles have disjoint transition regions;
this preserves positivity of the resulting position-dependent
coupling. We also identify disjoint sets of plaquettes whose
action will pay for the failed tests.

Fix a nonnegative smooth function \(\psi\) on \(\mathbb R^3\),
invariant under coordinate permutations and sign changes, with
integral one and support in \([-1/4,1/4]^3\). Put
\[
 C_{\rm mo}=\sum_{\mu=1}^3\int_{\mathbb R^3}|\partial_\mu\psi(u)|\,du .
\]
For a set \(S\subset T_\ell\) and a physical length \(s>0\), define
\begin{equation}\label{eq:history-profile}
 \varphi[S;s](y)=\int_{\mathbb R^3}\psi(u)
   \mathbf1\{\operatorname{dist}_\infty(y-su,S)\ge s/2\}\,du ,
\end{equation}
where the distance from the empty set is \(+\infty\).

\begin{lemma}\label{II-21-10}
The profile is smooth, takes values in \([0,1]\), vanishes where
\(\operatorname{dist}_\infty(y,S)\le s/4\), and equals one where
\(\operatorname{dist}_\infty(y,S)\ge3s/4\). Moreover,
\begin{equation}\label{II-21-13}
 \partial^\alpha\varphi[S;s](y)
 =s^{-|\alpha|}\int(\partial^\alpha\psi)(u)
       \mathbf1\{\operatorname{dist}_\infty(y-su,S)\ge s/2\}\,du .
\end{equation}
Thus the profile has Lipschitz constant \(C_{\rm mo}/s\) for the
sup metric. If the sets \(S_i\) have mutual distance at least \(s\),
then
\[
 1-\varphi[\bigcup_iS_i;s]
       =\sum_i(1-\varphi[S_i;s]).
\]
\end{lemma}

\begin{proof}
Lift the distance function to the periodic covering space.
The set where it is at least \(s/2\) is closed. After the change
of variables \(w=\widetilde y-su\), the integral is convolution
with \(s^{-3}\psi((\widetilde y-w)/s)\). Its derivatives are
integrable and have compact support, so differentiation under
the integral gives \eqref{II-21-13}. The derivative bounds and
the Lipschitz estimate follow by integrating along a shortest
lifted segment. On the set where \(\psi\ne0\), one has
\(|u|_\infty<1/4\). The Lipschitz property of distance therefore
gives both asserted plateaux, including their endpoints.
Finally, the open \(s/2\)-neighbourhoods of the \(S_i\) are
disjoint. Their indicators add, which proves the last identity.
\end{proof}

For a history \(h\) of level \(k\), write
\[
 s_j=M\check R_j a_{K-j},\qquad
 \varphi_j=\varphi[Z_j;s_j],\qquad
 \varphi_{j,D}=\varphi[D;s_j].
\]
The last notation is used for a component \(D\) of \(Z_j\).
The cells have side \(s_j\), and \(a_{K-j}\le s_j/4\).
Hence \(\varphi_j=1\) on \(\Lambda_j^{\sim-1}\),
\(\operatorname{supp}\varphi_j\subset\Lambda_j\), and
\begin{equation}\label{II-21-19}
 1-\varphi_j=\sum_{D\subset Z_j}(1-\varphi_{j,D}).
\end{equation}
Indeed distinct components have distance at least \(s_j\).
The tents of \eqref{II-20-6} also satisfy
\begin{equation}\label{II-21-20}
 \varphi_j h_z^{(j)}=h_z^{(j)}
       \qquad(z\in\Lambda_j^{\sim-1}\cap\mathbb T^{(j)}).
\end{equation}
To see this, the cell of \(z\) has distance at least \(s_j\)
from \(Z_j\), whereas a nonzero tent is supported within
\(a_{K-j}\le s_j/4\) of \(z\).

The position-dependent bare coupling and its level-\(n\) version are
\begin{equation}\label{eq:history-weight}
 \beta_n(y)=\beta-\sum_{j=1}^nL^jb_j\varphi_j(y),
 \qquad x_n(y)=L^{-n}\beta_n(y).
\end{equation}
Here \(x_n(\cdot)\) is a function, whereas \(x_n\) without
an argument is the running coupling already fixed. In this
subsection assume comparable cell sizes,
\begin{equation}\label{II-21-25}
 \check R_j\le L\check R_{j+1}\quad(1\le j<k),
 \qquad\text{equivalently }s_1\le\cdots\le s_k .
\end{equation}
Use the endpoint conventions \(\Omega_0=\Lambda_0=T_\ell\),
\(Z_0=\emptyset\), \(\Omega_{k+1}=\Lambda_{k+1}=\emptyset\)
and \(Z_{k+1}=T_\ell\).

\begin{lemma}\label{II-21-29}
At a point where the profiles have the pattern
\(1,\ldots,1,\varphi_n,0,\ldots,0\),
\begin{equation}\label{II-21-31}
 \beta_k=(1-\varphi_n)L^{n-1}x_{n-1}+\varphi_n L^nx_n .
\end{equation}
Every point has this pattern, or all profiles vanish. Consequently
\[
 0<\min_{0\le n\le k}L^nx_n
       \le\beta_k(y)\le\max_{0\le n\le k}L^nx_n .
\]
Formula \eqref{II-21-31} holds throughout the layer
\(\mathcal L_n=\Omega_n\setminus\Omega_{n+1}\) for \(1\le n\le k\);
on \(\mathcal L_0\), \(\beta_k=\beta\).
For \(y\in Z_n\), \(\beta_k(y)=\beta_{n-1}(y)\).
On \(\Lambda_n^{\sim-1}\cap Z_{n+1}\), \(\beta_k=L^nx_n\).
For every complex bare field \(\mathbf U\),
\begin{equation}\label{II-21-34}
 A(\beta_k(\cdot),\mathbf U)
       +\sum_{j=1}^kL^jb_jA(\varphi_j,\mathbf U)
          =\beta A(\mathbf U).
\end{equation}
\end{lemma}

\begin{proof}
History separation gives
\(\operatorname{dist}_\infty(Z_i,\Omega_{i+1})\ge8s_{i+1}\).
By \eqref{II-21-25}, \(\varphi_i=1\) on \(\Omega_{i+1}\),
hence on every later \(\Lambda_j\) and every later profile's
support. Thus at most one profile can lie strictly between zero
and one. Since \(Z_i\subset Z_j\) for \(i<j\), all profiles of
index at least \(n\) vanish on \(Z_n\).
Now
\[
 L^jb_j=L^{j-1}x_{j-1}-L^jx_j,\qquad
 \beta-\sum_{j=1}^nL^jb_j=L^nx_n .
\]
Substitution gives the convex combination and positivity.
On \(\mathcal L_n\), earlier profiles equal one and later ones
vanish because \(\Omega_{n+1}^c\subset Z_{n+1}\).
The remaining constant values follow from the plateaux.
Finally the weighted action is linear in its weight; applying
\eqref{eq:history-weight} at each plaquette centre proves
\eqref{II-21-34}.
\end{proof}

We next reserve action for each component created by the tests.
A \emph{piece} is a pair \((t,D)\), where \(1\le t\le k\)
and \(D\) is a component of \(Z_t\). Write \(\operatorname{Pc}(h)\)
for the set of pieces. The inner hull \(T^\natural\) at level
\(t\) is the union of the level-\(t\) cells containing a bare
cube of \(T\). Hulls and enlargements in the following definition
are taken at level \(t\).
\begin{definition}\label{II-21-26}
Set
\[
 F_t=\begin{cases}
 P_0^\natural\cup P_1^\natural\cup Q_1^\natural,&t=1,\\
 P_t^\natural\cup Q_t^\natural,&t\ge2,
 \end{cases}
 \qquad N_{t,D}=(F_t\cap D)^\sim .
\]
The share of \((t,D)\) is the set \(\mathcal P(N_{t,D})\)
of bare plaquettes whose four corners lie in \(N_{t,D}\).
For a set of pieces \(\mathfrak u\), define
\begin{equation}\label{eq:history-energy-share}
 A_{\rm cr}(\mathfrak u;U)
   =\sum_{(t,D)\in\mathfrak u}
        L^{t-1}x_{t-1}A_{N_{t,D}}(U),\qquad
 A_{\rm cr}(U)=A_{\rm cr}(\operatorname{Pc}(h);U).
\end{equation}
\end{definition}

\begin{proposition}\label{II-21-30}
For each piece, \(N_{t,D}\subset D\cap\Omega_t^c\).
It contains the closed \(s_t\)-neighbourhood of every failed
test cube contained in \(D\), including a density cube of
\(P_0\) when \(t=1\). If \(t\ge2\), then
\[
 \operatorname{dist}_\infty(N_{t,D},\widehat{Z_{t-1}})
       \ge3s_t,\qquad N_{t,D}\subset\Lambda_{t-1}.
\]
The regions \(N_{t,D}\) are pairwise disjoint.
At every point of \(N_{t,D}\), the profiles of index less
than \(t\) equal one and those of index at least \(t\) vanish.
In particular each plaquette in its share has weight
\(L^{t-1}x_{t-1}\). Thus for every real bare field \(U\),
\begin{equation}\label{II-21-35}
 A(\beta_k(\cdot),U)\ge A_{\rm cr}(U)
       \ge A_{\rm cr}(\mathfrak u;U)\ge0 .
\end{equation}
\end{proposition}

\begin{proof}
For any region \(T\),
\[
 T\subset T^\natural\subset\widehat T
       \subset(T^\natural)^\sim .
\]
The first inclusion follows by taking the cell containing
a bare cube of \(T\); for the last, a cell meeting \(T\)
touches one of these cells. For a union of test cubes,
the cells of the inner hull are precisely those containing
a test cube, since that partition refines the cell partition.
The formulas \eqref{eq:history-step-domain} and
\eqref{eq:history-first-domain} therefore give
\(F_t^\sim\subset\Omega_t^c\subset Z_t\).
Every cell of \(N_{t,D}\) touches a cell of \(D\), so belongs
to the same component of \(Z_t\). This proves
\(N_{t,D}\subset D\cap\Omega_t^c\).
For a failed cube \(\square\subset D\), its inner hull is
contained in \(F_t\cap D\). Its one-cell enlargement contains
the closed \(s_t\)-neighbourhood of \(\square\), proving the
neighbourhood assertion also for \(P_0\).

For \(t\ge2\), the tests select \(P_t\) outside
\(\widehat{Z_{t-1}}^{\sim4}\). They select \(Q_t\) in
\(B(P_t'')^{\sim-1}\), where
\[
 B(P_t'')^c=\widehat P_t^{\,\sim1}
                   \cup\widehat{Z_{t-1}}^{\sim5}.
\]
Consequently \(F_t\) has no cell in
\(\widehat{Z_{t-1}}^{\sim4}\).
One enlargement leaves a distance of at least \(3s_t\)
from \(\widehat{Z_{t-1}}\), and hence from \(Z_{t-1}\).
This proves the asserted distance and inclusion in
\(\Lambda_{t-1}\).

At a fixed level disjointness follows from the distinct
components. If \(t<t'\), then
\[
 N_{t,D}\subset Z_t\subset Z_{t'-1}
                  \subset\widehat{Z_{t'-1}},
\]
whereas \(N_{t',D'}\) has positive distance from that hull.
This proves disjointness across levels.
At a point of \(N_{t,D}\subset Z_t\), all later profiles
vanish. For \(j<t\), the distance from \(Z_j\) is at least
\(3s_t\ge3s_j\), so \(\varphi_j=1\).
The telescoping identity in \cref{II-21-29} gives the
claimed weight.

The centre of a bare plaquette with a corner in a region
of cells lies in that region: it belongs to the bare cube
centred at that corner, which is contained in its cell.
Thus the weight statement applies at every centre in a share.
The shares are disjoint sets of plaquettes and \(a_p(U)\ge0\)
for a real field. Split the weighted action into its sums
over shares and over the remaining plaquettes. The former
is \(A_{\rm cr}(U)\); the latter is nonnegative by
\cref{II-21-29}. Removing any shares decreases that sum,
which proves \eqref{II-21-35}.
\end{proof}

\section{Convergent local cluster sums}\label{sec:cluster-sums}

A small interaction can be exponentiated after its connected
contributions have been summed absolutely. We first prove the
finite polymer criterion and then check it for domains of cubes.
The criterion is the Koteck\'y--Preiss condition, in the form
recorded by Fern\'andez and Procacci
\cite[\S2, (2.15)]{FernandezProcacci2007}.
The proof below retains a finite-order bound throughout;
in particular, the identity for the logarithm follows only
after convergence has been established.

\subsection{The pinned cluster estimate}

Let \(\mathcal P\) be a finite set with a symmetric,
reflexive incompatibility relation \(\nsim\).
An activity is a function \(h:\mathcal P\to\mathbb C\).
Its partition function is
\[
 \Xi(h)=\sum_{\substack{\mathcal Q\subset\mathcal P\\
                         \mathcal Q\text{ pairwise compatible}}}
                    \prod_{Z\in\mathcal Q}h(Z).
\]
The empty subset contributes \(1\). For \(n\ge1\), let
\(\mathcal C_n\) be the connected graphs on
\(\{1,\ldots,n\}\), and set
\begin{equation}\label{II-18-45}
 \rho^T(Z_1,\ldots,Z_n)
 =\sum_{G\in\mathcal C_n}
            \prod_{\{i,j\}\in G}(-\mathbf1_{Z_i\nsim Z_j}),
 \qquad \rho^T(Z)=1.
\end{equation}
Repetitions among the \(Z_i\) are allowed. The coefficient
vanishes unless the incompatibility graph is connected;
such a family is called a cluster. Define
\[
 S=\sum_{n\ge1}\frac1{n!}
       \sum_{(Z_1,\ldots,Z_n)\in\mathcal P^n}
          \rho^T(Z_1,\ldots,Z_n)\prod_{i=1}^n h(Z_i).
\]

\begin{theorem}[The finite polymer criterion]\label{II-18-63}
Suppose \(a:\mathcal P\to[0,\infty)\) satisfies
\begin{equation}\label{II-18-48}
 \sum_{Y\nsim Z}|h(Y)|e^{a(Y)}\le a(Z)
                    \qquad(Z\in\mathcal P).
\end{equation}
Then \(S\) converges absolutely, \(\Xi(h)=e^S\ne0\), and
\begin{equation}\label{II-18-64}
 \Theta(Z):=\sum_{n\ge1}\frac1{(n-1)!}
   \sum_{Z_2,\ldots,Z_n\in\mathcal P}
     |\rho^T(Z,Z_2,\ldots,Z_n)|\prod_{i=2}^n|h(Z_i)|
       \le e^{a(Z)}.
\end{equation}
For every nonnegative function \(f\) on \(\mathcal P\),
\begin{equation}\label{II-18-65}
 \sum_{n\ge1}\frac1{n!}\sum_{Z_1,\ldots,Z_n}
 |\rho^T(Z_1,\ldots,Z_n)|
       \sum_{l=1}^n f(Z_l)\prod_{i\ne l}|h(Z_i)|
 \le\sum_Z f(Z)e^{a(Z)}.
\end{equation}
All these assertions hold for a subsystem and for activities
whose moduli are smaller. If holomorphic activities on an open
set in \(\mathbb C^N\) have a common majorant satisfying
\eqref{II-18-48}, their cluster sum, and every sum over a fixed
subfamily of clusters, is holomorphic there.
\end{theorem}

\begin{proof}
Write \(\theta_n(Z)\) for the term of order \(n\) in
\eqref{II-18-64}, and \(\Theta_R=\sum_{n=1}^R\theta_n\).
We prove the recursion
\begin{equation}\label{II-18-58}
 \Theta_{R+1}(Z)\le
       \exp\Bigl(\sum_{Y\nsim Z}|h(Y)|\Theta_R(Y)\Bigr).
\end{equation}
Remove vertex \(1\) from a connected graph on \(n\ge2\)
vertices. Its remaining components form a partition
\(\mathcal B\) of \(\{2,\ldots,n\}\), and each block
must have a nonempty set of edges to vertex \(1\).
For a fixed block \(B\), the sum over these edge sets is
\[
 \prod_{i\in B}(1-\mathbf1_{Z_i\nsim Z_1})-1
 =-\tau_B,\qquad
 \tau_B=\mathbf1_{\{Z_i\nsim Z_1\text{ for some }i\in B\}}.
\]
Summing first within each component therefore gives the exact
identity
\begin{equation}\label{II-18-53}
 \rho^T(Z_1,\ldots,Z_n)
 =\sum_{\mathcal B}\prod_{B\in\mathcal B}
                      \bigl(-\tau_B\rho^T(Z_B)\bigr).
\end{equation}
After taking absolute values, bound
\(\tau_B\le\sum_{i\in B}\mathbf1_{Z_i\nsim Z_1}\).
Put
\[
 \beta_k(Z)=\frac1{k!}\sum_{Y_1,\ldots,Y_k}
       |\rho^T(Y_1,\ldots,Y_k)|
       \Bigl(\sum_{i=1}^k\mathbf1_{Y_i\nsim Z}\Bigr)
                                 \prod_{i=1}^k|h(Y_i)|.
\]
Symmetry in the \(Y_i\) gives
\(\beta_k(Z)=\sum_{Y\nsim Z}|h(Y)|\theta_k(Y)\).
For blocks of ordered sizes \(k_1,\ldots,k_s\), summation
over partitions contributes
\((n-1)!/(s!\prod_l k_l!)\). Consequently
\[
 \theta_n(Z)\le
  \sum_{s=1}^{n-1}\frac1{s!}
   \sum_{\substack{k_1+\cdots+k_s=n-1\\k_l\ge1}}
                        \prod_{l=1}^s\beta_{k_l}(Z).
\]
Sum this inequality for \(2\le n\le R+1\), add
\(\theta_1=1\), and enlarge the sum by allowing every
\(k_l\le R\) independently. The resulting exponential
series proves \eqref{II-18-58}. Since \(\Theta_1=1\),
\eqref{II-18-48} and induction imply
\(\Theta_R\le e^a\). Monotone convergence proves
\eqref{II-18-64}.

Let \(S_n\) denote the term of order \(n\) in \(S\),
and let \(\overline S_n\) be the same sum with every term
replaced by its modulus. Rooting at the first polymer yields
\[
 \sum_{n\ge1}\overline S_n
 =\sum_{n\ge1}\frac1n\sum_Z|h(Z)|\theta_n(Z)
 \le\sum_Z|h(Z)|e^{a(Z)}<\infty.
\]
To identify the exponential, expand
\(\prod_{i<j}(1-\mathbf1_{Z_i\nsim Z_j})\)
as a sum over graphs and partition each graph into its
connected components. If \(\Xi^{(n)}\) is the sum over
compatible families of size \(n\), this gives
\[
 \Xi^{(n)}=
 \sum_{s=1}^n\frac1{s!}
  \sum_{\substack{k_1+\cdots+k_s=n\\k_l\ge1}}
                              \prod_{l=1}^sS_{k_l}.
\]
The sum of the moduli of all products on the right, over
all orders, is at most
\(\exp(\sum_n\overline S_n)\). We may thus regroup it
by \(s\) or by \(\sum_lk_l\). The first grouping gives
\(e^S\); the second gives \(\Xi\). Reflexivity excludes
repeated polymers from a compatible family, so
\(\Xi^{(n)}=0\) for \(n>\#\mathcal P\).

For \eqref{II-18-65}, move the marked position \(l\)
to the first position. Symmetry and \(n/n!=1/(n-1)!\)
identify its left side with \(\sum_Zf(Z)\Theta(Z)\).
The same estimate bounds the absolute cluster sum over
families meeting \(\mathcal Q\subset\mathcal P\) by
\begin{equation}\label{II-18-72}
 \sum_{Z\in\mathcal Q}|h(Z)|e^{a(Z)}.
\end{equation}
Indeed, mark each occurrence of a member of \(\mathcal Q\);
every such family is counted at least once.
Restriction and domination preserve \eqref{II-18-48}.
A common majorant gives a summable majorant for every
cluster series, uniformly in the holomorphic parameters.
The Weierstrass theorem then proves the last assertion.
The empty polymer system has \(\Xi=1\) and \(S=0\).
\end{proof}

\subsection{Domains, entropy, and decay}

Let \(L\ge13\) be odd, \(m'\ge0\) an integer, and
\(j\) an integer with \(0\le j\le K\). Fix the partition \(\pi_j\) of the
torus into cubes of side \(ML^{j-K}\), where
\(M=L^{m'}\) and \(m'<m+K-j\). Two cubes are adjacent
when they share a wall. A localization domain is a nonempty
union of cubes connected by such adjacencies; their finite
set is denoted by \(\mathcal D_j\).
For \(X\in\mathcal D_j\), write \(\#X\) for its
number of cubes. Define \(d_j(X)\) as the least number
of cube edges in a connected union of edges contained
in \(X\) which meets every cube of \(X\). A single
corner is allowed and has size zero. Cycles are allowed
in this definition.

\begin{proposition}[Entropy and joining]\label{II-18-26}
For every domain \(X\),
\[
 \#X\le8(d_j(X)+1).
\]
For every cube \(\square\) and \(n\ge0\),
\begin{equation}\label{II-18-27}
 \#\{X\in\mathcal D_j:\square\subset X,\ d_j(X)=n\}
                \le2048\cdot9216^n.
\end{equation}
In particular, with \(c_{\mathrm E}=\log18432\) and
\(C_{\mathrm E}=4096\),
\begin{equation}\label{II-18-29}
 \sum_{X\in\mathcal D_j:\square\subset X}
                   e^{-c d_j(X)}\le C_{\mathrm E}
                \qquad(c\ge c_{\mathrm E}).
\end{equation}
If the graph joining indices whose domains have equal
or adjacent cubes is connected, their union is a domain and
\begin{equation}\label{II-18-31}
 d_j\Bigl(\bigcup_{i=1}^nX_i\Bigr)+4
            \le\sum_{i=1}^n(d_j(X_i)+4).
\end{equation}
\end{proposition}

\begin{proof}
A domain has a spanning connected edge set: join the
edge skeletons of its cubes along a spanning tree of the
cube adjacency graph. A connected set with \(n\) edges
has at most \(n+1\) vertices. A cube met by this set
contains one of its vertices, and a vertex belongs to
at most eight cubes. This proves the first inequality.

Choose a minimizing connected edge set for a domain
containing \(\square\), and start at a vertex in
\(\square\). Traversing every edge twice gives a walk
of \(2n\) steps visiting all its vertices. There are
at most eight initial vertices and six choices for
each step. At most \(8(n+1)\) cubes meet the vertices
of this walk. The walk together with the subset of
these cubes constituting \(X\) determines \(X\), so
the number of possibilities is bounded by
\[
 8\cdot6^{2n}\cdot2^{8(n+1)}=2048\cdot9216^n.
\]
Summing the geometric series proves \eqref{II-18-29}.

If two domains have equal or adjacent cubes, choose
vertices of their spanning edge sets in those cubes.
The two vertices can be joined inside the cubes by
at most four edges. The union of the two spanning
sets and this path spans both domains. Thus
\(d_j(X_1\cup X_2)\le d_j(X_1)+d_j(X_2)+4\).
Adding the domains in an order supplied by a spanning
tree of their incompatibility graph proves
\eqref{II-18-31} and connectivity of the union.
\end{proof}

Domains \(\gamma,\gamma'\) are incompatible if a cube
of one is equal or adjacent to a cube of the other.
This is a symmetric reflexive relation. Given
\(\mathcal P\subset\mathcal D_j\) and activities
\(z(\gamma)\), regroup the cluster series by its union:
\begin{equation}\label{II-18-78}
 \mathsf L(X)=\sum_{n\ge1}\frac1{n!}
  \sum_{\substack{\gamma_1,\ldots,\gamma_n\in\mathcal P\\
                         \gamma_1\cup\cdots\cup\gamma_n=X}}
       \rho^T(\gamma_1,\ldots,\gamma_n)
                         \prod_{i=1}^nz(\gamma_i).
\end{equation}
Let \(\overline{\mathsf L}(X)\) denote the corresponding
sum of moduli.

\begin{theorem}[Local cluster sums]\label{II-18-86}
Suppose \(\varepsilon>0\), \(r''\ge0\),
\(r\ge r''+c_{\mathrm E}+8\), and
\[
 C_{\mathrm{gas}}=C_{\mathrm E}e^{8+4r''},\qquad
 7C_{\mathrm{gas}}\varepsilon\le1,
 \qquad |z(\gamma)|\le\varepsilon e^{-rd_j(\gamma)}.
\]
Then \eqref{II-18-78} converges absolutely and vanishes
unless \(X\in\mathcal D_j\). It depends only on the
activities of domains contained in \(X\), and
\begin{equation}\label{II-18-89}
 \Xi(z)=\exp\sum_{X\in\mathcal D_j}\mathsf L(X)\ne0.
\end{equation}
Moreover, for \(X\in\mathcal D_j\),
\begin{equation}\label{II-18-90}
 |\mathsf L(X)|\le\overline{\mathsf L}(X)
       \le C_{\mathrm{gas}}\varepsilon_Xe^{-r''d_j(X)},
 \qquad
 \varepsilon_X=\max_{\substack{\gamma\in\mathcal P\\\gamma\subset X}}
                              |z(\gamma)|e^{rd_j(\gamma)}.
\end{equation}
The maximum is zero when its index set is empty.
For holomorphic activities with the displayed common
majorant, every \(\mathsf L(X)\) is holomorphic.
\end{theorem}

\begin{proof}
Multiply each activity by its required decay weight:
\[
 \widehat z(\gamma)=z(\gamma)e^{r''(d_j(\gamma)+4)},
                  \qquad a(\gamma)=\#\gamma.
\]
The first bound in \cref{II-18-26} gives
\[
 |\widehat z(\gamma)|e^{\#\gamma}
 \le\varepsilon e^{8+4r''}
                         e^{-(r-r''-8)d_j(\gamma)}.
\]
By \eqref{II-18-29}, its sum over domains containing
one cube is at most \(C_{\mathrm{gas}}\varepsilon\).
Every domain incompatible with \(\gamma\) contains
one of the at most \(7\#\gamma\) cubes equal or
adjacent to a cube of \(\gamma\). Hence
\[
 \sum_{\gamma'\nsim\gamma}
          |\widehat z(\gamma')|e^{\#\gamma'}
 \le7\#\gamma\,C_{\mathrm{gas}}\varepsilon
 \le\#\gamma.
\]
Thus \cref{II-18-63} applies to both \(\widehat z\)
and \(z\). Absolute convergence permits regrouping
by the union. A nonzero connected coefficient has
a connected incompatibility graph, so its union is
a domain by \cref{II-18-26}. This proves
\eqref{II-18-89} and the support assertion.

Fix a cube \(\square\subset X\). Every cluster with
union \(X\) has a member containing \(\square\),
and all its members are contained in \(X\).
By \eqref{II-18-31}, multiplication of its absolute
weight by \(e^{r''(d_j(X)+4)}\) is bounded by the
same cluster weight formed with \(\widehat z\).
Apply \eqref{II-18-72} to the members satisfying
\(\square\subset\gamma\subset X\). It follows that
\[
 e^{r''(d_j(X)+4)}\overline{\mathsf L}(X)
 \le\sum_{\square\subset\gamma\subset X}
                    |\widehat z(\gamma)|e^{\#\gamma}
 \le C_{\mathrm{gas}}\varepsilon_X.
\]
Dropping \(e^{-4r''}\le1\) proves
\eqref{II-18-90}. The series indexed by clusters
with union \(X\) uses only activities inside \(X\),
which proves locality. Holomorphy follows from the
common-majorant assertion of \cref{II-18-63}.
\end{proof}


\section{Gaussian conditioning and local interactions}
\label{sec:gaussian-conditioning}

The local expansion uses a Gaussian field coupled to auxiliary
real variables. Its covariance and coupling may be complex.
We control the integral by completing the square with a real
reference pair. The resulting bound charges only the variables
on which the coupling is supported. This will preserve locality
when the auxiliary field is defined on the whole torus.

\subsection{Normalization and an integrable majorant}

We use orthonormal coordinates on the real spaces of bond
fields, with three coordinates for each positive bond. Pairings
are extended bilinearly to the complexification, transposes are
bilinear transposes, and operator norms are induced by the
Hermitian Euclidean norm. Thus the formulas below apply directly
to the field spaces and pairings used earlier.

For a complex symmetric matrix \(O\) with \(\operatorname{Re}O>0\),
the integral
\[
 I(O,J)=\int_{\mathbb R^d}
              e^{-B^{\mathsf T}OB/2+J^{\mathsf T}B}\,dB
\]
converges absolutely and is holomorphic in \((O,J)\).
Indeed, locally \(\operatorname{Re}O\ge cI\) and \(|J|\le C\),
so the integrand has the integrable majorant
\(e^{-c|B|^2/2+C|B|}\). The set of such symmetric matrices
is open and convex, and its members are invertible: if
\(Ov=0\), positivity of \(\operatorname{Re}(v^*Ov)\)
implies \(v=0\). At real positive \(O\) and real \(J\),
completion of the square gives the Gaussian formula.
Holomorphic continuation on this connected domain gives
\begin{equation}\label{II-19-65}
 I(O,0)^2\det(O/2\pi)=1,
 \qquad I(O,J)=I(O,0)e^{J^{\mathsf T}O^{-1}J/2}.
\end{equation}
In particular, \(I(O,0)^{-1}\) is the unique holomorphic
square root of \(\det(O/2\pi)\) which is positive for real
positive \(O\). Uniqueness follows because the quotient of
two such roots is a continuous function with values in
\(\{1,-1\}\). We use this root in
\[
 d\mu_C(B)=\det{}^{1/2}(C^{-1}/2\pi)
                          e^{-B^{\mathsf T}C^{-1}B/2}\,dB,
 \qquad \int d\mu_C=1.
\]
The zero-dimensional integral and determinant both equal one.

Let \(C\) be a symmetric invertible matrix on \(\mathbb C^d\)
with \(\operatorname{Re}C^{-1}>0\),
and \(\Gamma:\mathbb C^D\to\mathbb C^d\). For real \(X,B\),
write
\begin{equation}\label{II-19-83}
 \Theta_{C,\Gamma}(X,B)=
 \frac{\det{}^{1/2}(C^{-1})}{(2\pi)^{(d+D)/2}}
 \exp\left[-\frac{|X|^2}{2}
 -\frac{B^{\mathsf T}C^{-1}B}{2}
 -B^{\mathsf T}\Gamma X
 -\frac{X^{\mathsf T}\Gamma^{\mathsf T}C\Gamma X}{2}\right].
\end{equation}
The root is the preceding one, with the scalar factor
\((2\pi)^{-d/2}\) placed outside it. Define
\[
 \mathfrak J(C,\Gamma)[F]
   =\int_{\mathbb R^D}\int_{\mathbb R^d}
                         \Theta_{C,\Gamma}(X,B)F(B)\,dB\,dX
\]
when the integral of the modulus is finite.

\begin{proposition}[Gaussian domination]\label{II-19-85}
Suppose \(0<\gamma\le1\), \(c\ge1\), and real matrices
\(C_{\rm r}=C_{\rm r}^{\mathsf T}\), \(\Gamma_{\rm r}\)
satisfy
\[
 C_{\rm r}^{-1}\ge\gamma I,
 \qquad
 \|C_{\rm r}^{-1}\|,\ \|\Gamma_{\rm r}\|,
                     \|C_{\rm r}\Gamma_{\rm r}\|\le c.
\]
Assume that each of
\begin{equation}\label{eq:gaussian-three-defects}
 C^{-1}-C_{\rm r}^{-1},\qquad
 \Gamma-\Gamma_{\rm r},\qquad
 \Gamma^{\mathsf T}C\Gamma
          -\Gamma_{\rm r}^{\mathsf T}C_{\rm r}\Gamma_{\rm r}
\end{equation}
has norm at most \(\widehat\varepsilon\ge0\).
Let \(P\) be a real orthogonal coordinate projection on
\(\mathbb R^D\) of rank \(q\), with
\(\Gamma=\Gamma P\) and \(\Gamma_{\rm r}=\Gamma_{\rm r}P\).
For \(\rho\ge0\), suppose
\[
 u=\widehat\varepsilon+\rho
       \le\frac1{12}\min\{\gamma,c^{-2}\}.
\]
Then \(\operatorname{Re}C^{-1}\ge(11/12)\gamma I\), and
\begin{equation}\label{eq:gaussian-local-mass}
 \int\!\int |\Theta_{C,\Gamma}(X,B)|e^{\rho|B|^2}\,dB\,dX
 \le \exp\left(\frac92\gamma^{-1}ud+6c^2uq\right).
\end{equation}
There is also the pointwise bound
\begin{equation}\label{eq:gaussian-uniform-majorant}
 |\Theta_{C,\Gamma}(X,B)|e^{\rho|B|^2}
 \le \frac{(2c)^{d/2}}{(2\pi)^{(d+D)/2}}
       \exp\left[-\frac{|X|^2}{8}
                      -\frac{\gamma}{16c^2}|B|^2\right].
\end{equation}
Consequently \(\mathfrak J(C,\Gamma)[F]\) is defined whenever
\(F\) is Borel and \(|F(B)|\le F_0e^{\rho|B|^2}\).
\end{proposition}

\begin{proof}
The first defect in \eqref{eq:gaussian-three-defects}
gives the claimed lower bound on \(\operatorname{Re}C^{-1}\).
Set
\[
 \widetilde B=B+C_{\rm r}\Gamma_{\rm r}X,
 \qquad a=2\widehat\varepsilon+4uc^2\le6c^2u\le\tfrac12.
\]
Completion of the square with the reference pair leaves
three error terms. Their absolute values sum to at most
\[
 \tfrac12\widehat\varepsilon|PX|^2
 +\tfrac12\widehat\varepsilon|B|^2
 +\widehat\varepsilon|B|\,|PX|
 \le\widehat\varepsilon(|B|^2+|PX|^2).
\]
Use
\( |B|^2\le2|\widetilde B|^2+2c^2|PX|^2\).
The real part of the exponent in
\(\Theta_{C,\Gamma}e^{\rho|B|^2}\) is therefore at most
\begin{equation}\label{eq:gaussian-reference-square}
 -\frac{|(I-P)X|^2}{2}
 -\frac{1-a}{2}|PX|^2
 -\frac12\widetilde B^{\mathsf T}
                      (C_{\rm r}^{-1}-4uI)\widetilde B.
\end{equation}
This is a real change of variables in the integral.
Integrating first in \(\widetilde B\), and then in
the two orthogonal parts of \(X\), gives the bound
\[
 \left|\frac{\det C^{-1}}{\det C_{\rm r}^{-1}}\right|^{1/2}
       \det(I-4uC_{\rm r})^{-1/2}(1-a)^{-q/2}.
\]
The three factors are bounded respectively by
\[
 e^{\widehat\varepsilon d/(2\gamma)},\qquad
 e^{4ud/\gamma},\qquad e^{aq}\le e^{6c^2uq}.
\]
For the first, factor
\(C^{-1}=C_{\rm r}^{-1}[I+C_{\rm r}(C^{-1}-C_{\rm r}^{-1})]\)
and use \(\|C_{\rm r}\|\le\gamma^{-1}\).
For the second, every eigenvalue of \(4uC_{\rm r}\)
lies in \([0,1/3]\). For the third, \(0\le a\le1/2\).
Both use \(-\log(1-s)\le2s\) on \([0,1/2]\).
This proves \eqref{eq:gaussian-local-mass}.

For a majorant independent of the reference, use \(P=I\)
in \eqref{eq:gaussian-reference-square}.
Its negative quadratic terms are bounded above by
\(-|X|^2/4-\gamma|\widetilde B|^2/4\).
Since \(|\widetilde B|^2\ge|B|^2/2-c^2|X|^2\),
retain the fraction \(1/(2c^2)\) of the second term.
The result is at most
\(-|X|^2/8-\gamma|B|^2/(16c^2)\).
Finally \(\|C^{-1}\|\le2c\), whence
\(|\det C^{-1}|^{1/2}\le(2c)^{d/2}\).
This proves \eqref{eq:gaussian-uniform-majorant}, whose
right side is integrable.
\end{proof}

The bound also proves holomorphy under the integral.
More precisely, suppose \(C,C^{-1},\Gamma\) depend
holomorphically on a parameter in an open subset of
\(\mathbb C^N\), and have references satisfying
\cref{II-19-85} with the same \(\gamma,c,\widehat\varepsilon\).
The references need not depend holomorphically on the parameter.
If \(F\) is jointly Borel, holomorphic in that parameter
for every real \(B\), and has the common bound
\(F_0e^{\rho|B|^2}\), then
\(\mathfrak J(C,\Gamma)[F]\) is holomorphic.
The determinant root and the integrand are holomorphic;
\eqref{eq:gaussian-uniform-majorant} permits Cauchy's
formula, or Morera's theorem, to pass through the integral.

\subsection{The conditioning identity}

The covariance recovered by the conditioning integral is
\(C+C\Gamma\Gamma^{\mathsf T}C\). This identity will
identify the local expansion with the whole fluctuation
integral. It requires an equality of covariances, in
addition to local estimates on the kernels.

\begin{proposition}\label{prop:gaussian-conditioning}
Under the hypotheses of \cref{II-19-85},
\begin{equation}\label{eq:conditioning-Fourier}
 \mathfrak J(C,\Gamma)[e^{if^{\mathsf T}B}]
  =e^{-f^{\mathsf T}(C+C\Gamma\Gamma^{\mathsf T}C)f/2}
                       \qquad(f\in\mathbb R^d).
\end{equation}
In particular, \(\mathfrak J(C,\Gamma)[1]=1\).
Suppose also that \(C_*\) is symmetric and invertible on
\(\mathbb C^{D_*}\), and \(\chi:\mathbb R^{D_*}\to\mathbb R^d\)
is a coordinate restriction, with
\begin{equation}\label{II-19-20}
 C+C\Gamma\Gamma^{\mathsf T}C=\chi C_*\chi^{\mathsf T}.
\end{equation}
Assume a real positive reference satisfies
\[
 \gamma I\le C_{*\rm r}^{-1}\le cI,
 \qquad \|C_*^{-1}-C_{*\rm r}^{-1}\|
                                  \le\widehat\varepsilon.
\]
Then, for every Borel \(F\) with
\(|F(B)|\le F_0e^{\rho|B|^2}\),
\begin{equation}\label{eq:conditioning-whole}
 \mathfrak J(C,\Gamma)[F]
              =\int d\mu_{C_*}(B_*)\,F(\chi B_*).
\end{equation}
Both integrals converge absolutely.
\end{proposition}

\begin{proof}
The majorant justifies integrating first in \(B\).
Apply \eqref{II-19-65} with \(O=C^{-1}\) and
\(J=if-\Gamma X\). Its quadratic \(\Gamma X\) term
cancels the corresponding term of
\eqref{II-19-83}. The remaining integral is
\[
 e^{-f^{\mathsf T}Cf/2}
          \int d\mu_I(X)e^{-if^{\mathsf T}C\Gamma X}.
\]
One more application of \eqref{II-19-65} proves
\eqref{eq:conditioning-Fourier}.

For the whole integral,
\[
 \operatorname{Re}C_*^{-1}-2\rho I
 \ge(\gamma-\widehat\varepsilon-2\rho)I
 \ge\tfrac56\gamma I.
\]
Its integrand is therefore absolutely integrable.
Integrating \(\Theta_{C,\Gamma}\) in \(X\) gives
an integrable complex density on \(\mathbb R^d\);
integrating \(\mu_{C_*}\) in the complementary
coordinates gives another. By
\eqref{eq:conditioning-Fourier}, \eqref{II-19-65}
and \eqref{II-19-20}, these finite complex measures
have identical Fourier transforms. Uniqueness of the
Fourier transform of finite measures makes their
densities equal almost everywhere. Fubini's theorem,
with the two absolute bounds above, now gives
\eqref{eq:conditioning-whole} for the stated \(F\).
\end{proof}


\section{Local Gaussian integration}\label{sec:gaussian-mayer}

We now apply the polymer criterion to a Gaussian integral with
local interactions and small-field tests. The Gaussian variables
are decoupled by parameters attached to the larger cubes. Finite
differences in these parameters produce the decay needed for the
cluster sum. We first record how the size of a connected region
changes under this passage to larger cubes.

\subsection{Enlargement and Gaussian data}

Fix \(0\le k<K\), put \(t=k+1\), and assume
\(m'<m+K-t\). Lengths in this section are measured in the
scale of level \(k\). Fine cubes have side \(M=L^{m'}\),
and big cubes have side \(LM\). For a fine region \(Y\),
let \(Y^\uparrow\) be the union of the big cubes containing
a fine cube of \(Y\), and put
\[
 Y^\sharp=(Y^\sim)^\uparrow.
\]
For a lattice bond \(b\), let \(\square(b)\) be the union
of the fine cubes containing its endpoints, and put
\(b^\sharp=\square(b)^\sharp\).

\begin{lemma}[Size after enlargement]\label{II-18-23}
For \(Y\in\mathcal D_k\), \(Y^\sharp\in\mathcal D_t\) and
\begin{equation}\label{II-18-24}
 d_t(Y^\sharp)\le\frac{d_k(Y)}{L-2}
                     \le\frac{13}{9L}d_k(Y).
\end{equation}
For every big cube \(\Delta\) and \(\lambda\ge\log18432\),
\begin{equation}\label{eq:sharp-entropy}
 \sum_{\substack{Y\in\mathcal D_k\\\Delta\subset Y^\sharp}}
                  e^{-\lambda d_k(Y)}\le (L+2)^3\,4096.
\end{equation}
The regions \(\square(b)\) and \(b^\sharp\) are connected,
and both have linear size zero. At most
\(6((L+2)M)^3\) positive bonds satisfy
\(\Delta\subset b^\sharp\).
\end{lemma}

\begin{proof}
Use integer coordinates in units of the fine cube size, with a
common origin for the two partitions. A big cube has the form
\(Q[U]=U+[0,L]^3\), \(U\in L\mathbb Z^3\), modulo the
periods of the torus. It belongs to \(Y^\sharp\) precisely
when \(Y\) has a cube of index
\(u\in U+\{-1,\ldots,L\}^3\), modulo these periods.
For a fixed fine cube, these big cubes form a rectangular
array containing its unique parent. Parent cubes of adjacent
fine cubes are equal or adjacent. Thus \(Y^\sharp\) is
wall-connected.

Let \(\Gamma\) be a connected set of fine edges and corners
in \(Y\) which meets every cube of \(Y\), with
\(|\Gamma|=d_k(Y)\). A one-corner set is allowed.
For each coordinate direction \(\mu\), choose
\(s_\mu\in\{1,\ldots,L-2\}\) so that the number
\(n_\mu(s_\mu)\) of edges of \(\Gamma\) whose initial
coordinate has residue \(s_\mu\) modulo \(L\) is minimal.
Then
\[
 \sum_{\mu=1}^3n_\mu(s_\mu)\le\frac{|\Gamma|}{L-2}.
\]
Map each fine corner by
\[
 R(v)_\mu=L\left\lceil\frac{v_\mu-s_\mu}{L}\right\rceil.
\]
This map respects the periods. A fine edge is contracted to
a point unless its initial coordinate has the chosen residue;
otherwise it maps to one coarse edge. Corners of a fine cube
map to corners of its parent, so the image is a connected
coarse edge set in \(Y^\uparrow\). It has at most
\(|\Gamma|/(L-2)\) edges.

If \(Q[U]\subset Y^\sharp\), choose a cube of \(Y\)
with index \(u\in U+\{-1,\ldots,L\}^3\). Since
\(\Gamma\) meets this cube, it contains one of its corners
\(v\in u+\{0,1\}^3\). For
\(U_\mu-1\le v_\mu\le U_\mu+L+1\) and
\(1\le s_\mu\le L-2\), the formula for \(R\) gives
\(R(v)_\mu\in\{U_\mu,U_\mu+L\}\). Thus the image
meets every cube of \(Y^\sharp\), proving the first bound
in \eqref{II-18-24}. The second follows from \(L\ge13\).

For fixed \(\Delta=Q[U]\), every \(Y\) in
\eqref{eq:sharp-entropy} contains one of the at most
\((L+2)^3\) fine cubes with indices in
\(U+\{-1,\ldots,L\}^3\). Apply the entropy estimate of
Proposition~\ref{II-18-26} to each of these cubes.
The endpoints of a nearest-neighbor bond lie in one fine cube
or in two adjacent fine cubes. These cubes have a common
corner, so \(d_k(\square(b))=0\); \eqref{II-18-24} gives
\(d_t(b^\sharp)=0\). If \(\Delta\subset b^\sharp\),
one endpoint belongs to one of the same \((L+2)^3\) cubes.
They contain at most \(((L+2)M)^3\) lattice sites, and each
site is an endpoint of six positive bonds. This proves the
last assertion.
\end{proof}

Let \(\mathfrak X\) be the nonpivot bonds of a region
\(\Omega_t\) for \(\pi_t\). Write
\(\mathfrak X\langle Z\rangle\) for the bonds with at least
one endpoint in \(Z\), and
\(\mathfrak X^\circ\langle Z\rangle\) for those with both
endpoints there. Fields are odd under bond reversal. On a
reversal-invariant set \(\mathfrak e\), use the coordinates
on \(\mathfrak e^+\), with Euclidean norm and measure
\[
 \|B\|^2=\sum_{b\in\mathfrak e^+}|B(b)|^2,
 \qquad dB=\prod_{b\in\mathfrak e^+}dB(b).
\]
There are \(3\#\mathfrak e^+\) real coordinates. In
particular, for a big region \(Z\),
\begin{equation}\label{eq:gaussian-region-dimension}
 N(Z):=3\#\mathfrak X^+\langle Z\rangle
                            \le18(LM)^3\#_{\pi_t}Z.
\end{equation}
Indeed, each bond counted has an endpoint among the
\((LM)^3\#_{\pi_t}Z\) sites of \(Z\).

\begin{definition}\label{II-19-15}
Fix a nonempty big region \(R\), and constants
\(\bar\kappa\ge1\), \(0<\gamma_{\rm r}\le1\),
\(c_{\rm r}\ge1\), and \(\hat\varepsilon\ge0\).
A decoupled Gaussian structure assigns to each fine region
\(Z_\circ\) with \(Z_\circ^\sharp\subset R\), including
the empty region, operators
\[
 C(Z_\circ,s):\mathcal H(Z_\circ)\longrightarrow
                 \mathcal H(Z_\circ),\qquad
 \Gamma(Z_\circ,s):\mathcal H(R)\longrightarrow
                 \mathcal H(Z_\circ).
\]
Here \(\mathcal H(Z)=\mathbb C^{3\#\mathfrak X^+\langle Z\rangle}\),
\(C\) is symmetric and invertible, and \(s\) ranges over
the closed polydisc \(|s(\Delta)|\le e^{\bar\kappa}\) on
\[
 \mathsf P(R,Z_\circ)=
       \{\Delta\subset R:\Delta\not\subset Z_\circ^\sharp\}.
\]
The following properties are required.
\begin{enumerate}
\item[(St1)] The kernels of \(C,C^{-1},\Gamma\) extend
holomorphically to a neighborhood of this closed polydisc.
\item[(St2)] Let \(Z_\circ^\sharp\subset Z\subset R\),
let \(\gamma_i\) be the wall-components of \(Z\), and
suppose \(s=0\) outside \(Z\). Put \(Z_i=Z_\circ\cap\gamma_i\)
and let \(s_i\) agree with \(s\) in \(\gamma_i\) and vanish
elsewhere. Then
\begin{equation}\label{II-19-17}
 C(Z_\circ,s)=\bigoplus_i C(Z_i,s_i),\qquad
 \chi_{\mathcal H(Z_i)}\Gamma(Z_\circ,s)=\Gamma(Z_i,s_i),
\end{equation}
and \(\Gamma(Z_i,s_i)\) uses only auxiliary coordinates
in \(\mathfrak X\langle\gamma_i\rangle\).
\item[(St3)] For each \(Z_\circ\) there are real
\(C_{\rm r}=C_{\rm r}^{\mathsf T}\), \(\Gamma_{\rm r}\),
independent of \(s\), such that \(\Gamma_{\rm r}\) uses
only coordinates in \(\mathfrak X\langle Z_\circ^\sharp\rangle\),
\begin{equation}\label{II-19-18}
 C_{\rm r}^{-1}\ge\gamma_{\rm r}I,
 \qquad
 \|C_{\rm r}^{-1}\|,\ \|\Gamma_{\rm r}\|,
                  \ \|C_{\rm r}\Gamma_{\rm r}\|\le c_{\rm r},
\end{equation}
and, throughout the polydisc,
\begin{equation}\label{II-19-19}
 \max\{\|C^{-1}-C_{\rm r}^{-1}\|,
        \|\Gamma-\Gamma_{\rm r}\|,
        \|\Gamma^{\mathsf T}C\Gamma-
                  \Gamma_{\rm r}^{\mathsf T}C_{\rm r}\Gamma_{\rm r}\|\}
                       \le\hat\varepsilon.
\end{equation}
\end{enumerate}
When \(R=T_\ell\), the structure is complete if it also
has a symmetric invertible whole covariance \(C_*\), with
a real reference satisfying
\[
 \gamma_{\rm r}I\le C_{*\rm r}^{-1}\le c_{\rm r}I,
 \qquad\|C_*^{-1}-C_{*\rm r}^{-1}\|\le\hat\varepsilon,
\]
and the covariance identity \eqref{II-19-20} holds for
every \(Z_\circ\) at \(s=1\). This is property (St4).
\end{definition}

All intersections here are intersections of cube sets.
The decompositions in (St2) are disjoint even for bonds
crossing the boundary of \(Z_\circ\). To see this, a fine
bond has endpoints in one big cube or two adjacent big cubes.
It cannot have endpoints in different wall-components of
\(Z\). Moreover, the touching enlargement of a fine cube
is a rectangular array of fine cubes, and its big-cube
parents form a wall-connected array. If that enlargement
lies in \(Z\), it lies in the component containing the
original fine cube. Taking unions gives
\[
 (Z_\circ\cap\gamma_i)^\sharp
           =Z_\circ^\sharp\cap\gamma_i,
 \qquad
 \mathfrak X\langle Z_\circ\rangle
           =\bigsqcup_i\mathfrak X\langle Z_\circ\cap\gamma_i\rangle.
\]
The sets \(\mathfrak X\langle\gamma_i\rangle\) are
pairwise disjoint for the same endpoint reason.

Write
\[
 \mathbb J(Z_\circ;s)[F]
      =\mathfrak J(C(Z_\circ,s),\Gamma(Z_\circ,s))[F],
\]
with the conditioning integral of
Section~\ref{sec:gaussian-conditioning}.
Under (St2), this integral factorizes over the components
whenever \(F=\prod_iF_i\), with \(F_i\) depending on
\(\mathcal H(Z_i)\), and the corresponding integrals are
absolutely convergent. Indeed, both quadratic forms and the
mixed term split into the indicated blocks. The determinant
root is multiplicative on these blocks: it is so for positive
real matrices, and then throughout the domain by holomorphy.
The unused auxiliary coordinates integrate to one. Fubini's
theorem gives the asserted product.

\subsection{The local activities and their bounds}

For the rest of this section, \(\mathfrak X\) is the set of
nonpivot bonds of the whole lattice. Let \(\mathfrak b\in(0,\infty]\)
and \(t_b\in[0,1]\). A potential family
\(\mathbb V^o=(\mathbb V^o(Y;B))\) is indexed by
\(Y\in\mathcal D_k\) with \(Y^\sharp\subset R\) and
depends only on the field on
\(\mathfrak X^\circ\langle Y\rangle\).
It is of sort \(\Sigma_{\rm v}\) if it is a Borel function
on \(|B(b)|<\mathfrak b\), and of sort \(\Sigma_{\rm q}\)
if it is Borel on all real fields. The labels \(o\) are
called origins. We allow at most four origins, and require
\(t_b=1\) for all \(b\) if any origin has sort
\(\Sigma_{\rm v}\). The abstract sort \(\Sigma_{\rm q}\)
does not require a quadratic dependence on \(B\).
Set
\begin{equation}\label{II-19-37}
 \|\mathbb V\|_{\star,r}
  =\sup_{Y,B}\frac{e^{r d_k(Y)}|\mathbb V(Y;B)|}
           {1+M^{-3}\sum_{b\in\mathfrak X^{\circ+}\langle Y\rangle}
                                                        |B(b)|^2},
\end{equation}
where \(B\) ranges over the domain of the family.

We give the constants explicitly because their smallness
conditions will be checked for the running interactions.
Let \(\gamma_2\ge0\), \(\lambda_{\rm out}>0\), and put
\begin{gather*}
 \lambda_{\rm s}=\log18432,\qquad S_{\rm s}=4096,
 \qquad w_0=(4S_{\rm s}L^3)^{-1},\\
 \mathsf C_{\rm ex}\ge
       \max\{8e^8,13/9,3+\tfrac14\log(7S_{\rm s})\}.
\end{gather*}
Define
\begin{equation}\label{II-19-46}
\begin{gathered}
 u=\hat\varepsilon+2(LM)^{-3}+\tfrac12\gamma_2,
 \qquad c_G=\mathsf C_{\rm ex}(\gamma_{\rm r}^{-1}+c_{\rm r}^2),\\
 c_{\rm v}=2+\mathsf C_{\rm ex}c_G(LM)^3u,
 \qquad c_{\rm w}=c_{\rm v}+3,\\
 \lambda_1=\lambda_{\rm out}+8c_{\rm w}
                         +\mathsf C_{\rm ex}+\lambda_{\rm s},
 \qquad\lambda_2=\lambda_1-8c_{\rm w},\\
 \lambda_V=\frac{\mathsf C_{\rm ex}}L\lambda_1+2\lambda_{\rm s},
 \qquad\lambda_0=\lambda_V-\lambda_{\rm s},\\
 w_\star=\frac{w_0e^{-4\lambda_1}}
                  {\mathsf C_{\rm ex}(L+2)^3S_{\rm s}},\qquad
 \mathsf p=e^{-\gamma_2\mathfrak b^2/2},\qquad
 \mathsf p_\star=\mathsf C_{\rm ex}((L+2)M)^3e^{4\lambda_1}\mathsf p,\\
 K_E=\mathsf C_{\rm ex}S_{\rm s}e^{8c_{\rm w}-4\lambda_1},
 \qquad K^\Pi_P=K_E,\qquad
 K_\Pi=K_E(L+2)^3S_{\rm s}e^{4\lambda_1}/w_0.
\end{gathered}
\end{equation}
When \(\mathfrak b=\infty\), set \(\mathsf p=0\).

For a reversal-invariant bond set \(\mathfrak e\), put
\[
 \mathsf{in}_{\mathfrak e}(B)
       =\prod_{b\in\mathfrak e^+}\mathbf1_{\{t_b|B(b)|<\mathfrak b\}},
 \qquad
 \mathsf{out}_{\mathfrak e}(B)
       =\prod_{b\in\mathfrak e^+}\mathbf1_{\{t_b|B(b)|\ge\mathfrak b\}}.
\]
A Mayer index \(\iota=(\mathbf D,\mathfrak e)\ne(\emptyset,\emptyset)\)
consists of a set of vertices \((Y,o)\) and a
reversal-invariant set of large-field bonds. If \(Y^{\rm v}\)
is the union of the vertex domains of sort \(\Sigma_{\rm v}\),
we require
\(\mathfrak e\cap\mathfrak X\langle Y^{\rm v}\rangle=\emptyset\)
and \(b^\sharp\subset R\) for \(b\in\mathfrak e\).
Define
\begin{equation}\label{II-19-29}
\begin{split}
 Z_\circ(\iota)&=\bigcup_{(Y,o)\in\mathbf D}Y
                     \ \cup\bigcup_{b\in\mathfrak e}\square(b),\\
 F_\iota(B)&=(-1)^{\#\mathfrak e^+}
       \mathsf{in}_{\mathfrak X\langle Y^{\rm v}\rangle}(B)
       \mathsf{out}_{\mathfrak e}(B)
            \prod_{(Y,o)\in\mathbf D}(e^{\mathbb V^o(Y;B)}-1).
\end{split}
\end{equation}
The function is defined to be zero when its small-field test
vanishes, so no potential is evaluated outside its domain.
It is Borel and depends only on
\(\mathfrak X\langle Z_\circ(\iota)\rangle\).

For \(\gamma\in\mathcal D_t\), \(\gamma\subset R\),
let \(\mathcal I(\gamma)\) consist of the indices with
\(Z_\circ(\iota)^\sharp\subset\gamma\), and put
\(\mathcal K(\gamma,\iota)=\gamma\setminus Z_\circ(\iota)^\sharp\),
as a set of big cubes. If
\(f_\iota(s)=\mathbb J(Z_\circ(\iota);s)[F_\iota]\), define
\begin{equation}\label{II-19-30}
 T(\gamma,\iota)=
 \sum_{\mathcal K'\subset\mathcal K(\gamma,\iota)}
  (-1)^{\#\mathcal K(\gamma,\iota)-\#\mathcal K'}
                              f_\iota(\mathbf1_{\mathcal K'}),
 \qquad
 \mathsf z(\gamma)=\sum_{\iota\in\mathcal I(\gamma)}T(\gamma,\iota).
\end{equation}
The index sets in these sums are finite. With the
incompatibility relation of Section~\ref{sec:cluster-sums}, set
\begin{equation}\label{II-19-31}
 E(X)=\sum_{n\ge1}\frac1{n!}
       \sum_{\substack{\gamma_1,\ldots,\gamma_n\subset R\\
                        \gamma_i\in\mathcal D_t,\ \cup_i\gamma_i=X}}
       \rho^T(\gamma_1,\ldots,\gamma_n)
                                  \prod_i\mathsf z(\gamma_i).
\end{equation}
Expanding the activities in this expression gives its refined
series, indexed also by \(\iota_i\in\mathcal I(\gamma_i)\).

The estimates keep track of specified vertices or large-field
bonds. For bounds \(A^o\ge\|\mathbb V^o\|_{\star,\lambda_V}\),
assign regions and weights to vertices, positive bonds, and
big cubes as follows:
\begin{equation}\label{II-19-49}
\begin{array}{c|ccc}
 e &(Y,o)&b&\Delta\\ \hline
 \operatorname{rg}(e)&Y^\sharp&b^\sharp&\Delta\\
 w(e)&(A^o/w_0)e^{-\lambda_0d_k(Y)}&\mathsf p&e^{-(\bar\kappa-1)}
\end{array}
\qquad
 \widehat w(e)=w(e)e^{\lambda_1(d_t(\operatorname{rg}(e))+4)}.
\end{equation}
The elements associated with \((\gamma,\iota)\) are
\(\mathcal E(\gamma,\iota)=\mathbf D\cup\mathfrak e^+\cup
\mathcal K(\gamma,\iota)\). A set of marks \(\mathfrak m\)
is any set of vertices and positive bonds; it contains no big
cubes. Put
\begin{equation}\label{II-19-50}
 \overline T(\gamma,\iota)
        =e^{c_{\rm v}\#\gamma}\prod_{e\in\mathcal E(\gamma,\iota)}w(e),
 \qquad
 \eta^{\mathfrak m}
        =\sup_{\Delta\subset R}
          \sum_{\substack{e\in\mathfrak m\\\Delta\subset\operatorname{rg}(e)}}
                                                     \widehat w(e).
\end{equation}
Here and below \(\#\gamma=\#_{\pi_t}\gamma\).
Let \(\overline z^{\mathfrak m}(\gamma)\) be the sum of
\(\overline T(\gamma,\iota)\) over indices whose element
set meets \(\mathfrak m\). Let
\(\overline E^{\mathfrak m}(X)\) be the refined absolute
cluster sum with \(T\) replaced by \(\overline T\), each
term multiplied by the number of its polymers whose element
set meets \(\mathfrak m\). Finally,
\(E^{\mathfrak m}(X)\) is the subseries of \(E(X)\)
consisting of refined terms that meet \(\mathfrak m\).

\begin{theorem}[Gaussian cluster expansion]\label{II-19-55}
Let a Gaussian structure as in Definition~\ref{II-19-15}
and at most four potential families be given. Suppose
\begin{equation}\label{II-19-57}
 u\le\mathsf C_{\rm ex}^{-1}
                    \min\{\gamma_{\rm r},c_{\rm r}^{-2}\},
 \qquad\bar\kappa\ge4\lambda_1+1,
 \qquad\mathsf p_\star\le\tfrac12,
 \qquad\|\mathbb V^o\|_{\star,\lambda_V}\le w_\star.
\end{equation}
All conditioning integrals in \eqref{II-19-30} converge
absolutely, and the refined series \eqref{II-19-31} does so
for every \(X\in\mathcal D_t\), \(X\subset R\).
If \(\|\mathbb V^o\|_{\star,\lambda_V}\le A^o\le w_\star\),
then
\begin{align}
 |T(\gamma,\iota)|&\le\overline T(\gamma,\iota),\label{eq:mayer-term-bound}\\
 \overline z^{\mathfrak m}(\gamma)
      &\le\eta^{\mathfrak m}e^{8c_{\rm w}-4\lambda_1}
                                      e^{-\lambda_2d_t(\gamma)},\label{eq:mayer-activity-bound}\\
 |E^{\mathfrak m}(X)|\le\overline E^{\mathfrak m}(X)
      &\le K_E\eta^{\mathfrak m}e^{-\lambda_{\rm out}d_t(X)}.
                                                     \label{II-19-112}
\end{align}
In particular,
\begin{equation}\label{eq:cluster-potential-bound}
 |E(X)|\le
 \left(K_\Pi\sum_o\|\mathbb V^o\|_{\star,\lambda_V}
                   +K^\Pi_P\mathsf p_\star\right)
                                      e^{-\lambda_{\rm out}d_t(X)}.
\end{equation}
When \(R=T_\ell\) and the structure is complete, the whole
integral converges absolutely and satisfies
\begin{equation}\label{II-19-39}
 \Xi:=\int d\mu_{C_*}(B)\,\mathsf{in}_{\mathfrak X}(B)
        \exp\left(\sum_o\sum_{Y\in\mathcal D_k}\mathbb V^o(Y;B)\right)
       =\exp\left(\sum_{X\in\mathcal D_t}E(X)\right)\ne0.
\end{equation}
The integrand on the left is zero where its test fails.
\end{theorem}

\begin{proof}
We first bound one Mayer integrand. Put
\[
 \alpha_Y=w_0e^{-\lambda_{\rm s}d_k(Y)},\qquad
 x_Y=1+M^{-3}\sum_{b\in\mathfrak X^{\circ+}\langle Y\rangle}|B(b)|^2.
\]
The potential bound gives
\(|\mathbb V^o(Y;B)|\le w(Y,o)\alpha_Yx_Y\).
Since \(w(Y,o)\le1\), \(\alpha_Yx_Y\ge0\), and
\(|e^z-1|\le|z|e^{|z|}\),
\[
 |e^{\mathbb V^o(Y;B)}-1|
               \le w(Y,o)e^{2\alpha_Yx_Y}.
\]
Also, for \(0\le t_b\le1\),
\[
 \mathbf1_{\{t_b|B(b)|\ge\mathfrak b\}}
              \le\mathsf p\,e^{\gamma_2|B(b)|^2/2}.
\]
For infinite threshold the left side is zero, as required.

Suppose \(Z_\circ(\iota)^\sharp\subset Z\), with \(Z\)
a big region. The entropy bound, summed over the fine cubes
of \(Z\) and at most four origins, gives
\[
 \sum_{(Y,o)\in\mathbf D}\alpha_Y\le\# Z.
\]
For a fixed bond in the interior of a vertex domain, select
the fine cube containing one of its endpoints. The same
entropy bound and \(4w_0S_{\rm s}=L^{-3}\) show
\(\sum_{(Y,o):\,b\in\mathfrak X^\circ\langle Y\rangle}
\alpha_Y\le L^{-3}\). Therefore, with
\(\rho_*=2(LM)^{-3}+\gamma_2/2\),
\begin{equation}\label{eq:mayer-integrand-bound}
 |F_\iota(B)|\le e^{2\# Z}
        \prod_{e\in\mathbf D\cup\mathfrak e^+}w(e)
                                      e^{\rho_*\|B\|^2}.
\end{equation}

Take parameters supported in \(Z\). By (St2), \(\Gamma\)
uses only auxiliary coordinates in
\(\mathfrak X\langle Z\rangle\); the same holds for
\(\Gamma_{\rm r}\). Apply Proposition~\ref{II-19-85}
with \(d\le N(Z)\), \(q\le N(Z)\), and
\(\hat\varepsilon+\rho_*=u\). Its exponential cost is at
most \(c_GuN(Z)\), since \(\mathsf C_{\rm ex}\ge8e^8\).
Using \eqref{eq:gaussian-region-dimension} and
\(\mathsf C_{\rm ex}\ge18\), we obtain
\begin{equation}\label{eq:mayer-integral-bound}
 \mathbb J^{\rm abs}(Z_\circ;s)[F_\iota]
       \le e^{c_{\rm v}\# Z}
                      \prod_{e\in\mathbf D\cup\mathfrak e^+}w(e).
\end{equation}
The same domination proves holomorphy in the decoupling
parameters. In applying the finite-dimensional bound, first
integrate the unused standard Gaussian coordinates; this
keeps its cost proportional to \(N(Z)\).

For completeness, if a holomorphic function \(g\) on the
polydisc of radius \(r>1\) has modulus at most \(M_0\),
its finite difference in a coordinate set \(\mathcal K\)
obeys
\begin{equation}\label{eq:ordinary-cauchy-difference}
 \left|\sum_{A\subset\mathcal K}
       (-1)^{\#\mathcal K-\#A}g(\mathbf1_A)\right|
                         \le M_0(r-1)^{-\#\mathcal K}.
\end{equation}
Indeed, one coordinate difference is represented by the
Cauchy kernel \(1/(z(z-1))\). Integrating these kernels
over circles of radius \(r'<r\) gives
\(M_0(r'-1)^{-\#\mathcal K}\); let \(r'\uparrow r\).
Apply this to \(f_\iota\), with all coordinates outside
\(\mathcal K(\gamma,\iota)\) set to zero, using
\eqref{eq:mayer-integral-bound} with \(Z=\gamma\).
Since \((e^{\bar\kappa}-1)^{-1}\le e^{-(\bar\kappa-1)}\),
this proves \eqref{eq:mayer-term-bound}.

We next sum the element weights through a fixed big cube.
Lemma~\ref{II-18-23} and the definition of \(\lambda_0\)
give
\begin{align*}
 \sum_{Y:\Delta\subset Y^\sharp}\widehat w(Y,o)
    &\le\frac{(L+2)^3S_{\rm s}e^{4\lambda_1}}{w_0}A^o
      \le\mathsf C_{\rm ex}^{-1},\\
 \sum_{b:\Delta\subset b^\sharp}\widehat w(b)
    &\le\mathsf p_\star,\qquad
 \widehat w(\Delta)=e^{-(\bar\kappa-1)+4\lambda_1}\le1.
\end{align*}
Let \(\mathfrak m_{\rm all}\) be all vertices and positive
bonds. With at most four origins,
\begin{equation}\label{eq:marked-weight-sum}
 \eta^{\mathfrak m_{\rm all}}
 \le H\sum_o A^o+\mathsf p_\star\le1,
 \qquad H=\frac{(L+2)^3S_{\rm s}e^{4\lambda_1}}{w_0}.
\end{equation}
The sum through \(\Delta\) over all three kinds of elements
is at most \(2\). For marks whose vertices have origins in
\(S\), the corresponding bound is
\begin{equation}\label{eq:selected-mark-weight}
 \eta^{\mathfrak m}\le H\sum_{o\in S}A^o
              +\mathsf p_\star\mathbf1_{\{\mathfrak m\text{ contains a bond}\}}.
\end{equation}

The regions of the elements of \(\mathcal E(\gamma,\iota)\)
are connected and have union \(\gamma\). The joining bound
of Proposition~\ref{II-18-26} yields
\[
 d_t(\gamma)+4
        \le\sum_{e\in\mathcal E(\gamma,\iota)}
                                   (d_t(\operatorname{rg}(e))+4).
\]
Consequently
\[
 \overline T(\gamma,\iota)
 \le e^{c_{\rm v}\#\gamma-4\lambda_1-\lambda_1d_t(\gamma)}
                      \prod_{e\in\mathcal E(\gamma,\iota)}\widehat w(e).
\]
For fixed \(\gamma\), the map \(\iota\mapsto\mathcal E(\gamma,\iota)\)
is injective: its vertices and positive bonds recover
\(\iota\). Enlarge the sum to all sets of elements whose
regions lie in \(\gamma\), choose one marked element, and
sum the others by a product of \(1+\widehat w(e)\).
Every element has a cube in its region, so
\[
 \sum_{\substack{\iota\in\mathcal I(\gamma)\\
                  \mathcal E(\gamma,\iota)\cap\mathfrak m\ne\emptyset}}
       \prod_{e\in\mathcal E(\gamma,\iota)}\widehat w(e)
       \le\eta^{\mathfrak m}\#\gamma\,e^{2\#\gamma}.
\]
Since \(\#\gamma\le e^{\#\gamma}\) and
\(\#\gamma\le8(d_t(\gamma)+1)\), this proves
\eqref{eq:mayer-activity-bound}.

It remains to apply the polymer criterion with enough decay
left for the union of a cluster. Set
\[
 \mu=\lambda_{\rm out}+1,\qquad a(\gamma)=\#\gamma,
 \qquad\overline h(\gamma)=\overline z^{\mathfrak m_{\rm all}}(\gamma),
 \qquad\widetilde h(\gamma)=\overline h(\gamma)
                                      e^{\mu(d_t(\gamma)+4)}.
\]
The activity bound and entropy estimate give, for every
\(\Delta\),
\begin{equation}\label{II-19-110}
 \sum_{\gamma:\Delta\subset\gamma}
       \overline z^{\mathfrak m}(\gamma)
              e^{\mu(d_t(\gamma)+4)+a(\gamma)}
       \le\eta^{\mathfrak m}S_{\rm s}
                    e^{8c_{\rm w}-4\lambda_1+4\mu+8}.
\end{equation}
Here \(\lambda_2-\mu-8=\mathsf C_{\rm ex}+\lambda_{\rm s}-9
\ge\lambda_{\rm s}\). The exponent on the right is
\(-24c_{\rm w}-4\lambda_{\rm s}+12-4\mathsf C_{\rm ex}\).
There are at most \(7\#\gamma\) cubes equal or adjacent
to a cube of \(\gamma\). By \eqref{eq:marked-weight-sum}
and \(7S_{\rm s}e^{12-4\mathsf C_{\rm ex}}\le1\),
\[
 \sum_{\gamma'\nsim\gamma}\widetilde h(\gamma')e^{a(\gamma')}
                         \le\#\gamma=a(\gamma).
\]
Thus Theorem~\ref{II-18-63} applies to \(\widetilde h\),
and hence also to \(\overline h\) and \(\mathsf z\).

In a refined cluster, summing the number of marked polymers
amounts to replacing, in turn, one factor \(\overline h\)
by \(\overline z^{\mathfrak m}\). For a cluster of union
\(X\), the joining estimate permits the extraction of
\(e^{-\mu(d_t(X)+4)}\). Apply the one-mark estimate
\eqref{II-18-65} to
\[
 f(\gamma)=\mathbf1_{\{\gamma\subset X\}}
       \overline z^{\mathfrak m}(\gamma)e^{\mu(d_t(\gamma)+4)}.
\]
Dropping the condition that the union equal \(X\) adds
only nonnegative terms. Then \eqref{II-19-110} gives
\[
 \overline E^{\mathfrak m}(X)
 \le e^{-\mu(d_t(X)+4)}\#X\,\eta^{\mathfrak m}S_{\rm s}
                       e^{8c_{\rm w}-4\lambda_1+4\mu+8}.
\]
Use \(\#X\le8e^{d_t(X)}\) and
\(\mathsf C_{\rm ex}\ge8e^8\) to obtain
\eqref{II-19-112}. Every Mayer index contains a vertex or
a positive large-field bond. Thus every refined term is
marked by \(\mathfrak m_{\rm all}\), which proves its
absolute convergence. Taking \(A^o=\|\mathbb V^o\|_{\star,\lambda_V}\)
in \eqref{eq:marked-weight-sum} proves
\eqref{eq:cluster-potential-bound}.

Finally let \(R=T_\ell\) and assume (St4).
Expand each \(e^{\mathbb V}\) as \(1+(e^{\mathbb V}-1)\)
and each small-field indicator outside the selected
\(\Sigma_{\rm v}\) domains as \(1-\mathsf{out}_{\{b,\bar b\}}\).
This finite expansion is exactly
\[
 \mathsf{in}_{\mathfrak X}(B)e^{\sum_{o,Y}\mathbb V^o(Y;B)}
                         =1+\sum_\iota F_\iota(B).
\]
The identity also holds where a test fails: for each fixed
set of vertices either its retained test is zero, or the
sum over its large-field subsets produces the missing zero
test. No extension of a small-field potential is used.
By \eqref{eq:mayer-integrand-bound} and
Proposition~\ref{prop:gaussian-conditioning}, each term is
absolutely integrable and its integral is
\(f_\iota(\mathbf1_{\mathsf P})\). Finite-difference inversion
gives
\[
 \Xi=1+\sum_\iota\sum_{\mathcal K\subset\mathsf P(R,Z_\circ(\iota))}
                                     f_\iota[\mathcal K\mid\emptyset].
\]
In this display the bracket denotes the alternating sum in
\eqref{II-19-30}. Let \(\gamma_i\) be the wall-components
of \(Z_\circ(\iota)^\sharp\cup\mathcal K\). Each vertex
and each large-field bond belongs to a unique component,
since its enlarged region is connected. Split \(\iota\)
accordingly. Gaussian factorization and the product of the
alternating sums show
\[
 f_\iota[\mathcal K\mid\emptyset]
       =\prod_i T(\gamma_i,\iota_i).
\]
If some \(\iota_i\) is empty, its component consists only
of cubes of \(\mathcal K\). Its factor is the nonempty
finite difference of the constant function one, hence zero.

The nonzero terms are in bijection with finite compatible
families \(\{\gamma_i\}\), each supplied with an index
\(\iota_i\in\mathcal I(\gamma_i)\). The inverse takes
the unions of the vertex sets and large-field bonds, and
sets \(\mathcal K=\bigcup_i\gamma_i\setminus Z_\circ(\iota)^\sharp\).
Compatibility makes the bond fields disjoint and preserves
the exclusion of large-field bonds from the tested vertex
domains. Thus
\[
 \Xi=\sum_{\{\gamma_i\}\text{ compatible}}\prod_i\mathsf z(\gamma_i).
\]
Theorem~\ref{II-18-63} and the absolute convergence already
proved identify its logarithm by the cluster series.
Grouping by the union gives \eqref{II-19-39}.
\end{proof}

\subsection{Comparison and parameter dependence}

We write \(E[\underline{\mathbb V};\mathbf t]\) when the
potential system and the bond parameters need to be displayed.
Let \(E^\circ[\underline{\mathbb V}]\) be the cluster sum
of the \(\Sigma_{\rm q}\) origins alone, with threshold
\(\infty\). We call this the Gaussian part; in its application
below, these origins will be quadratic forms.

\begin{corollary}\label{cor:cluster-comparison}
Assume the hypotheses of Theorem~\ref{II-19-55}.
If \(\underline{\mathbb V}\) has at most three origins
and \(\|\mathbb W\|_{\star,\lambda_V}\le w_\star\), then
\begin{equation}\label{eq:cluster-add-potential}
 |E[\underline{\mathbb V}\cup\{\mathbb W\};\mathbf t](X)
            -E[\underline{\mathbb V};\mathbf t](X)|
 \le K_\Pi\|\mathbb W\|_{\star,\lambda_V}
                                  e^{-\lambda_{\rm out}d_t(X)}.
\end{equation}
Here \(\mathbb W\) has either sort, with sort
\(\Sigma_{\rm q}\) required if \(\mathbf t\ne1\).
For any system of at most four origins,
\begin{equation}\label{eq:cluster-gaussian-part}
 |E(X)-E^\circ(X)|
 \le\left(K_\Pi\sum_{o\text{ of sort }\Sigma_{\rm v}}
                \|\mathbb V^o\|_{\star,\lambda_V}
                  +K^\Pi_P\mathsf p_\star\right)
                                  e^{-\lambda_{\rm out}d_t(X)}.
\end{equation}
For a complete structure on the torus,
\[
 \int d\mu_{C_*}(B)
       \exp\left(\sum_{o\text{ of sort }\Sigma_{\rm q}}\sum_Y
                                                \mathbb V^o(Y;B)\right)
                 =\exp\left(\sum_XE^\circ(X)\right),
\]
with an absolutely convergent integral.

Suppose in addition that \(\mathbb W\) has the sort of
an existing origin \(o_1\), and that both
\(\underline{\mathbb V}\cup\{\mathbb W\}\) and the system
obtained by replacing \(\mathbb V^{o_1}\) with
\(\mathbb V^{o_1}+\mathbb W\) satisfy the theorem's bounds.
Their activities and their cluster sums agree.
\end{corollary}

\begin{proof}
Mark the new origin. The refined terms containing no vertex
of that origin are exactly those of the old system, with
unchanged tests and Gaussian integrals. Absolute convergence
therefore identifies the difference in
\eqref{eq:cluster-add-potential} with the marked subseries.
Use \eqref{eq:selected-mark-weight} and
\eqref{II-19-112}.

For \eqref{eq:cluster-gaussian-part}, mark every
\(\Sigma_{\rm v}\) vertex and every large-field bond.
The unmarked terms are precisely the terms of the
\(\Sigma_{\rm q}\) subsystem with no large-field bonds.
For that subsystem at infinite threshold, all other terms
vanish. The same two estimates prove the bound. The infinite
threshold changes only \(\mathsf p\) and \(\mathsf p_\star\)
to zero, so Theorem~\ref{II-19-55} applies and also proves
the displayed integral identity.

For the last assertion, merge the two origin labels.
For each domain, the three nonempty choices of vertices
before merging contribute
\[
 (e^V-1)+(e^W-1)+(e^V-1)(e^W-1)=e^{V+W}-1.
\]
Their unions, sorts, and tests agree with those of the merged
vertex. Thus the Mayer functions sum along the fibers of
the merging map to the new Mayer function. The conditioning
integral and finite differences are linear, so the activities
agree. Their absolutely convergent cluster sums agree as well.
\end{proof}

\begin{proposition}\label{prop:cluster-parameters}
Fix the geometry, threshold, bond parameters, origins, and
sorts in Theorem~\ref{II-19-55}.
Let the Gaussian structure and potentials depend on
\(z\) in an open set \(\mathcal O\subset\mathbb C^N\).
Suppose the structure has the same constants for every
\(z\), the kernels \(C,C^{-1},\Gamma\) are jointly
holomorphic on an open neighborhood of the closed
decoupling polydisc times \(\mathcal O\), and each potential
is jointly Borel in \((B,z)\) and holomorphic in \(z\)
at fixed \(B\) in its domain. If
\[
 \sup_{z\in\mathcal O}\|\mathbb V^o_z\|_{\star,\lambda_V}
                                      \le A^o\le w_\star,
\]
then \(T_z(\gamma,\iota)\), \(\mathsf z_z(\gamma)\),
\(E_z(X)\), and every fixed marked part \(E^{\mathfrak m}_z(X)\)
are holomorphic in \(z\). Their refined series have the
common summable majorants of Theorem~\ref{II-19-55}.

For two systems satisfying that theorem on regions
\(R,R'\), with the same constants and the same remaining
discrete data, let \(X\subset R\cap R'\). Suppose their
potentials agree for \(Y^\sharp\subset X\), and, at every
\((Z_\circ,\mathbf1_A)\) with
\(Z_\circ^\sharp\subset X\) and
\(A\subset X\setminus Z_\circ^\sharp\), their operators
\(C\) agree and their kernels \(\Gamma\) agree on auxiliary
coordinates in \(\mathfrak X\langle X\rangle\).
Then all terms supported in \(X\), and hence \(E(X)\), agree.
\end{proposition}

\begin{proof}
The index sets are independent of \(z\). The tests are fixed,
and the zero extension in \eqref{II-19-29} is across a fixed
Borel set. Thus \(F_{\iota,z}\) is jointly Borel and is
holomorphic in \(z\) for every real \(B\).
The bound \eqref{eq:mayer-integrand-bound} is uniform in
\(z\). The common Gaussian majorant in
Proposition~\ref{II-19-85} justifies holomorphy under each
conditioning integral. The real references need not depend
holomorphically on \(z\). Finite differences and the finite
activity sums preserve holomorphy. Finally,
\(\overline E^{\mathfrak m_{\rm all}}(X)\) is a common
summable majorant of the refined series. Uniform convergence
proves the assertion for \(E_z\) and each of its fixed
marked subseries.

For locality, the Mayer indices with enlarged support in
\(X\) and their functions coincide. At parameters supported
in \(X\), (St2) makes \(\Gamma\) vanish on the auxiliary
coordinates outside \(\mathfrak X\langle X\rangle\).
Those coordinates integrate to one, leaving identical
conditioning integrals in the two systems. Their finite
differences, activities, and cluster sums therefore agree.
\end{proof}

For a system parametrized by backgrounds \(P\) that can be
restricted to big regions, the locality hypothesis takes the
following useful form. Whenever \(Z_\circ^\sharp\subset Z\subset R\)
and \(s=0\) outside \(Z\), require
\begin{equation}\label{II-19-21}
\begin{split}
 C_R(Z_\circ,s;P)
   &=C_Z(Z_\circ,s|_{\mathsf P(Z,Z_\circ)};P|_Z),\\
 \Gamma_R(Z_\circ,s;P)
   &=\Gamma_Z(Z_\circ,s|_{\mathsf P(Z,Z_\circ)};P|_Z)
                                      \chi_{\mathfrak X\langle Z\rangle}.
\end{split}
\end{equation}
Together with restriction-compatible potentials, this implies
that \(E(X;P)\) depends only on \(P|_X\).

The expansion also preserves the symmetries of its data.
Let \(\mathsf r\) be a lattice isometry preserving the
two cube partitions and the nonpivot bond set, and let
\(O_b\) be real orthogonal maps on \(\mathfrak g\), for
\(b\in\mathfrak X^+\). They define an orthogonal field map
\(\mathsf O\) by
\[
 (\mathsf OB)((\mathsf rb)^+)=O_bB(b),
\]
with odd extension to reversed bonds. Transform the tests by
\(\widetilde t_{(\mathsf rb)^+}=t_b\).
If the transformed data satisfy
\[
 \widetilde C=\mathsf O C\mathsf O^{\mathsf T},\qquad
 \widetilde\Gamma=\mathsf O\Gamma\mathsf O^{\mathsf T},\qquad
 \widetilde{\mathbb V}^{o}(\mathsf rY;\mathsf OB)
                                  =\mathbb V^o(Y;B),
\]
at the correspondingly transformed decoupling vertices,
then \(\widetilde E(\mathsf rX)=E(X)\).
In fact \(\mathsf r\) gives bijections of the Mayer indices
and compatible polymers; the field maps preserve the tests,
Lebesgue measure, and all quadratic forms in the conditioning
density. The determinant root is preserved by real orthogonal
conjugation, first for positive real matrices and then by
holomorphy. Hence the terms agree, and absolute convergence
allows their reindexing. This argument includes reflections
and orthogonal maps of determinant \(-1\).

Likewise, complex conjugation of the operators at every
decoupling vertex and of the potentials on real fields
conjugates \(T,\mathsf z,E\). The tests and finite-difference
coefficients are real, the chosen determinant root commutes
with conjugation, and the coefficients \(\rho^T\) are
integers. Thus real data give real cluster sums and real
marked parts.

\section{Local variations of the fluctuation integral}
\label{sec:cluster-variations}

The cluster expansion assigns a local term to each connected
region. To separate the changes in the interaction from the
changes in the cutoff, we shall also assign the first variation
to individual bonds. This requires a bound on the integral of
the absolute value of a finite-difference kernel.

\subsection{Insertions and changes in the interaction}

Use the setting and constants of Theorem~\ref{II-19-55}.
Unless stated otherwise, all bond parameters in this
subsection equal one. Write
\(\|\mathbb W\|=\|\mathbb W\|_{\star,\lambda_V}\)
for the potential norm. Let
\(g(Z_\circ,s;\mathbb X,B)\) denote the density
\(\Theta_{C(Z_\circ,s),\Gamma(Z_\circ,s)}(\mathbb X,B)\)
of \eqref{II-19-83}. If \(Z_\circ^\sharp\subset\gamma\), put
\begin{equation}\label{II-20-34}
 g^\gamma(Z_\circ;\mathbb X,B)
   =\sum_{A\subset\mathcal K}(-1)^{\#\mathcal K-\#A}
                         g(Z_\circ,\mathbf1_A;\mathbb X,B),
 \qquad\mathcal K=\gamma\setminus Z_\circ^\sharp.
\end{equation}

\begin{lemma}\label{II-20-33}
Suppose \(F\) is Borel and
\(\mathbb J^{\rm abs}(Z_\circ;s)[F]\le J_0<\infty\)
for all decoupling parameters in the closed polydisc that
vanish outside \(\gamma\). Then
\begin{equation}\label{II-20-35}
 \iint |g^\gamma(Z_\circ;\mathbb X,B)|\,|F(B)|
                                    \,d\mathbb X\,dB
       \le e^{-(\bar\kappa-1)\#\mathcal K}J_0.
\end{equation}
In particular the finite difference of
\(\mathbb J(Z_\circ;s)[F]\) equals
\(\iint g^\gamma F\).
\end{lemma}

\begin{proof}
The last assertion follows by linearity of the finitely many
absolutely convergent integrals. For the bound, apply the
Cauchy representation used in
\eqref{eq:ordinary-cauchy-difference} to the density before
integration. At every fixed \((\mathbb X,B)\), it is
holomorphic in the decoupling parameters. On circles of
radius \(1<r<e^{\bar\kappa}\), each coordinate difference
contributes the kernel \((re^{i\theta}-1)^{-1}\) in angular
coordinates. Taking absolute values, multiplying by
\(|F(B)|\), and applying Tonelli gives
\[
 \iint |g^\gamma|\,|F|
       \le (r-1)^{-\#\mathcal K}
                 \sup_{\theta}\mathbb J^{\rm abs}(Z_\circ;s_\theta)[F]
       \le(r-1)^{-\#\mathcal K}J_0.
\]
The angular and field integrand is measurable, since the
density is continuous in both sets of variables and \(F\)
is Borel. Let \(r\uparrow e^{\bar\kappa}\) and use
\((e^{\bar\kappa}-1)^{-1}\le e^{-(\bar\kappa-1)}\).
When \(\mathcal K\) is empty the assertion is the original
bound at \(s=0\).
\end{proof}

Let \(\underline{\mathbb V}\) have at most three origins
and satisfy \(\|\mathbb V^o\|\le w_\star\).
An insertion \(\mathbb W\) is a potential family of either
sort with finite norm, without a smallness requirement.
A Mayer index with an insertion is
\(\iota'=(\mathbf D,\mathfrak e,Y')\). Its enlarged support
includes \(Y'^\sharp\), and its tested region includes
\(Y'\) if \(\mathbb W\) has sort \(\Sigma_{\rm v}\).
The set \(\mathfrak e\) avoids the bonds incident to the
whole tested region. Unlike an ordinary Mayer index,
\((\mathbf D,\mathfrak e)\) may be empty. Put
\begin{equation}\label{II-19-33}
\begin{split}
 Z_\circ(\iota')&=
       \bigcup_{(Y,o)\in\mathbf D}Y
          \ \cup\bigcup_{b\in\mathfrak e}\square(b)\ \cup Y',\\
 F'_{\iota'}(B)&=(-1)^{\#\mathfrak e^+}
       \mathsf{in}_{\mathfrak X\langle Y^{\rm v}\rangle}(B)
       \mathsf{out}_{\mathfrak e}(B)
       \prod_{(Y,o)\in\mathbf D}(e^{\mathbb V^o(Y;B)}-1)
                                                   \mathbb W(Y';B).
\end{split}
\end{equation}
Again the function is zero where the test fails. Define
\(T'(\gamma,\iota')\) by the finite difference of the
conditioning integral of \(F'_{\iota'}\), and let
\(\mathsf z'[\underline{\mathbb V};\mathbb W](\gamma)\)
be its finite sum over \(Z_\circ(\iota')^\sharp\subset\gamma\).
The cluster sum with one insertion is
\begin{equation}\label{II-19-34}
 E'[\underline{\mathbb V};\mathbb W](X)
  =\sum_{n\ge1}\frac1{(n-1)!}
      \sum_{\substack{\gamma_1,\ldots,\gamma_n\subset R\\
               \gamma_i\in\mathcal D_t,\ \cup_i\gamma_i=X}}
          \rho^T(\gamma_1,\ldots,\gamma_n)
          \mathsf z'(\gamma_1)\prod_{i=2}^n\mathsf z(\gamma_i).
\end{equation}

\begin{theorem}[First variation]\label{II-20-54}
With these hypotheses, all terms in \eqref{II-19-34} are
defined, and its refined series converges absolutely. It is
linear in insertions of a fixed sort, and
\begin{equation}\label{II-20-56}
 |E'[\underline{\mathbb V};\mathbb W](X)|
       \le K_\Pi\|\mathbb W\|e^{-\lambda_{\rm out}d_t(X)}.
\end{equation}
It vanishes if \(\mathbb W(Y;\cdot)=0\) whenever
\(Y^\sharp\subset X\). For
\(\epsilon_\star=w_\star/\|\mathbb W\|\), with
\(\epsilon_\star=\infty\) when \(\mathbb W=0\),
the cluster sum with the additional origin
\(\epsilon\mathbb W\) is holomorphic for
\(|\epsilon|<\epsilon_\star\), and
\begin{equation}\label{II-20-55}
 \left.\frac{d}{d\epsilon}
 E[\underline{\mathbb V}\cup\{\epsilon\mathbb W\}](X)
                                                    \right|_{\epsilon=0}
                      =E'[\underline{\mathbb V};\mathbb W](X).
\end{equation}
The insertion sum has the same locality, orthogonal covariance,
and conjugation properties as its Gaussian data and potentials.
Under the parameter hypotheses of
Proposition~\ref{prop:cluster-parameters}, with a jointly Borel,
holomorphic insertion of uniformly bounded norm, it is holomorphic
in those parameters.

If an origin \(o_1\) and a family \(\mathbb W_1\) of its
sort satisfy \(\|\mathbb V^{o_1}\|,\|\mathbb W_1\|\le w_\star/2\),
then
\begin{multline}\label{II-20-59}
 |E'[\underline{\mathbb V}+_{o_1}\mathbb W_1;\mathbb W](X)
                  -E'[\underline{\mathbb V};\mathbb W](X)|\\
 \le\frac{8K_\Pi}{w_\star}
       \|\mathbb W\|\,\|\mathbb W_1\|e^{-\lambda_{\rm out}d_t(X)}.
\end{multline}
Here \(+_{o_1}\) means addition to the named existing origin.
\end{theorem}

\begin{proof}
Use the majorants in \eqref{II-19-50} with every origin
bound equal to \(w_\star\). Adjoin a new origin \(o_+\)
of the insertion's sort. The map
\[
 \jmath(\mathbf D,\mathfrak e,Y')
       =(\mathbf D\cup\{(Y',o_+)\},\mathfrak e)
\]
is a bijection from insertion indices to ordinary indices
having exactly one vertex of the new origin. It preserves
the support, tested region, and decoupling cubes. If
\(\overline T^+\) is the four-origin majorant, set
\[
 \overline T'(\gamma,\iota')
                 =w_\star^{-1}\overline T^+(\gamma,\jmath\iota').
\]

We first prove the stronger integral bound
\begin{equation}\label{II-20-49}
 T'(\gamma,\iota')=\iint g^\gamma F'_{\iota'},\qquad
 \iint|g^\gamma|\,|F'_{\iota'}|
                          \le\|\mathbb W\|\overline T'(\gamma,\iota').
\end{equation}
There is nothing to prove when \(\mathbb W=0\). Otherwise,
put \(c=w_\star/\|\mathbb W\|\) and introduce the Borel
potential
\[
 \mathbb U(Y;B)=\log(1+c|\mathbb W(Y;B)|).
\]
It has the insertion's domain and sort, and
\(\|\mathbb U\|\le w_\star\). For the system with this
fourth origin,
\[
 |F_{\jmath\iota'}(B)|=c|F'_{\iota'}(B)|.
\]
Apply \eqref{eq:mayer-integral-bound} to this admitted
system and then Lemma~\ref{II-20-33}. Dividing by \(c\)
proves \eqref{II-20-49}. The same lemma gives
\(\iint|g^\gamma|\,|F_\iota|\le\overline T(\gamma,\iota)\)
for ordinary terms.

Let \(\mathfrak m_+\) mark the new origin. In the absolute
marked cluster series for the augmented system, retain only
terms having exactly one new-origin vertex. The vertex can
occur at any of the \(n\) polymers. Permuting that polymer
to the first position preserves \(\rho^T\) and replaces
\(1/n!\) by \(1/(n-1)!\). The map \(\jmath\) and the
definition of \(\overline T'\) therefore give
\begin{equation}\label{II-20-51}
 \sum_{\text{refined insertion indices}}
 \frac{|\rho^T(\gamma_1,\ldots,\gamma_n)|}{(n-1)!}
       \overline T'(\gamma_1,\iota'_1)
                   \prod_{i=2}^n\overline T(\gamma_i,\iota_i)
 \le\frac{\overline E^{+,\mathfrak m_+}(X)}{w_\star}
 \le K_\Pi e^{-\lambda_{\rm out}d_t(X)}.
\end{equation}
The last inequality uses
\(\eta^{\mathfrak m_+}\le Hw_\star\) from
\eqref{eq:selected-mark-weight}. This is a summable majorant
independent of the insertion and of the admitted potentials.
Together with \eqref{II-20-49}, it proves absolute convergence
and \eqref{II-20-56}. Linearity follows first under each
absolutely convergent integral and then in the series.
The support assertion holds term by term, since
\(Y'^\sharp\subset\gamma_1\subset X\).

To identify the derivative, apply
Proposition~\ref{prop:cluster-parameters} to the fourth
origin \(\epsilon\mathbb W\). Its refined series is
uniformly dominated on \(|\epsilon|<\epsilon_\star\).
A term containing \(p\) new-origin vertices has the factors
\(\prod_{j=1}^p(e^{\epsilon\mathbb W(Y_j;B)}-1)\).
At \(\epsilon=0\) its derivative is zero unless \(p=1\),
in which case the single factor is replaced by
\(\mathbb W(Y';B)\). This differentiation may be performed
under the field integral: use Cauchy's formula on a fixed
circle \(|\epsilon|<\epsilon_\star\), followed by
\eqref{II-20-49}'s ordinary-term counterpart and Fubini.
The common summable majorant then permits differentiation
of the refined series by the same Cauchy formula. Exactly
one of its polymers carries the derivative; the same
permutation and factorial count as above give
\eqref{II-20-55}. If \(\mathbb W=0\), the added-origin
terms vanish and the identity holds directly.

Locality and orthogonal covariance follow by differentiating
the corresponding equalities for the augmented system.
For conjugation, differentiate along real \(\epsilon\).
For holomorphic parameters, the uniform insertion bound
provides one positive radius in \(\epsilon\); the augmented
cluster sum is jointly holomorphic, so its partial derivative
at zero is holomorphic as well.

Finally set \(a=K_\Pi\|\mathbb W\|e^{-\lambda_{\rm out}d_t(X)}\).
When \(\mathbb W_1\ne0\), let
\(r=w_\star/(2\|\mathbb W_1\|)\ge1\) and
\[
 \phi(v)=E'[\underline{\mathbb V}+_{o_1}v\mathbb W_1;\mathbb W](X).
\]
The system is admitted on \(|v|\le r\), so
\(|\phi(v)|\le a\); it is holomorphic on \(|v|<r\).
If \(r>1\), Schwarz's lemma applied to
\((\phi(rz)-\phi(0))/(2a+\delta)\), followed by
\(\delta\downarrow0\), gives
\(|\phi(1)-\phi(0)|\le2a/r\). For \(r=1\) this is
the triangle inequality. This yields the bound with
\(4K_\Pi/w_\star\), hence \eqref{II-20-59}.
The zero cases follow directly from linearity.
\end{proof}

Consider now a path \((\underline{\mathbb V}_v)_{0\le v\le1}\)
with fixed origins and sorts, at most three origins, and
\(\|\mathbb V_v^o\|\le w_\star\). Suppose its potentials
are jointly Borel in \((v,B)\), are \(C^1\) in \(v\)
at every fixed \(B\), and have velocities satisfying
\[
 V^o:=\sup_{0\le v\le1}\|\dot{\mathbb V}_v^o\|<\infty.
\]
The derivatives at the endpoints are one-sided.

\begin{proposition}\label{prop:cluster-path}
For this path,
\begin{equation}\label{II-20-58}
 E[\underline{\mathbb V}_1](X)-E[\underline{\mathbb V}_0](X)
      =\int_0^1\sum_o
                 E'[\underline{\mathbb V}_v;\dot{\mathbb V}_v^o](X)\,dv.
\end{equation}
The integrand is Borel and bounded in modulus by
\(K_\Pi\sum_o V^o\,e^{-\lambda_{\rm out}d_t(X)}\).

Suppose that the nonzero velocities have one sort and
admit a decomposition into insertions \(\dot{\mathbb V}_{v,b}\)
of that sort, jointly Borel in \((v,B)\), such that
\[
 \sum_o\dot{\mathbb V}_v^o(Y;B)
       =\sum_{b\in\mathfrak X^+}\dot{\mathbb V}_{v,b}(Y;B),
 \qquad
 \dot{\mathbb V}_{v,b}(Y;\cdot)=0
                         \quad\text{if }b\notin\mathfrak X^\circ\langle Y\rangle.
\]
If \(\overline v=\sup_{v,b}\|\dot{\mathbb V}_{v,b}\|<\infty\),
then the terms
\begin{equation}\label{II-20-22}
 E^\bullet(X,b)=\int_0^1
                  E'[\underline{\mathbb V}_v;\dot{\mathbb V}_{v,b}](X)\,dv
\end{equation}
are defined and satisfy
\begin{gather}
 \sum_{b\in\mathfrak X^+}E^\bullet(X,b)
          =E[\underline{\mathbb V}_1](X)-E[\underline{\mathbb V}_0](X),
                                                    \label{II-20-66}\\
 |E^\bullet(X,b)|\le K_\Pi\overline v\,e^{-\lambda_{\rm out}d_t(X)},
 \qquad E^\bullet(X,b)=0\quad\text{unless }b^\sharp\subset X.
                                                    \label{II-20-67}
\end{gather}
They inherit locality, covariance, and conjugation from
the path and its bond insertions. They are holomorphic in
additional parameters when the Gaussian structure, path,
velocities, and bond insertions satisfy the joint Borel
and holomorphic hypotheses above with uniform norm bounds.
\end{proposition}

\begin{proof}
The velocities are jointly Borel as limits of difference
quotients, using a forward increment except near the right
endpoint and a backward increment there. For an ordinary
Mayer index the tests are independent of \(v\), and
\[
 \partial_v(e^{V_v}-1)=\dot V_v+
                                      \dot V_v(e^{V_v}-1).
\]
Differentiate its finite product. Each summand is an
insertion integrand in one of two ways: remove the
differentiated vertex and insert its velocity, or retain
that vertex and insert its velocity. The two constructions
have the same tested region and support as the original
index. Conversely, an insertion index of a fixed origin
arises uniquely by the first construction if its vertex
is absent, and by the second if its vertex is present.

Use the pointwise fundamental theorem of calculus for
these finite products and integrate against \(g^\gamma\).
The bound \eqref{II-20-49}, with
\(\|\mathbb W\|\le V^o\), justifies Fubini over the path
and field variables. Summing the finite index sets gives
\[
 \mathsf z_{v_2}(\gamma)-\mathsf z_{v_1}(\gamma)
   =\int_{v_1}^{v_2}\sum_o
       \mathsf z'[\underline{\mathbb V}_v;\dot{\mathbb V}_v^o](\gamma)\,dv.
\]
Thus each activity is absolutely continuous with this
derivative almost everywhere. Products of the activities
are absolutely continuous. In the derivative of a cluster
product, the distinguished factor can be put first; symmetry
again converts \(n/n!\) to \(1/(n-1)!\).
Equation~\eqref{II-20-51} bounds the resulting absolute
series, uniformly in \(v\), by
\(K_\Pi\sum_o V^o e^{-\lambda_{\rm out}d_t(X)}\).
Fubini now permits summing the integrated derivatives,
which proves \eqref{II-20-58}.
The integrands are Borel: their field integrands are jointly
Borel, their integrals converge absolutely, and their
cluster series are pointwise limits of Borel partial sums.

Apply Theorem~\ref{II-20-54} to each bond insertion and
integrate its bound. Linearity in insertions of one sort
and the finite bond sum prove \eqref{II-20-66}. If
\(b^\sharp\not\subset X\), then \(b\) cannot lie in
\(\mathfrak X^\circ\langle Y\rangle\) for a domain
\(Y^\sharp\subset X\). Every insertion integrand therefore
vanishes, proving the support assertion.
The identities expressing locality, covariance, and
conjugation hold before the path integral. Under the stated
uniform parameter bounds the integrand is holomorphic
at each \(v\) and bounded by an integrable constant, so
integration also preserves holomorphy.
\end{proof}

\subsection{Changes in the sharp cutoff}

We now assume \(0<\mathfrak b<\infty\). All potentials in
this subsection have sort \(\Sigma_{\rm q}\); their fields
are therefore defined without a small-field test. There may
be up to four origins, each with norm at most \(w_\star\).
For a bond \(b\), let \(\mathbf t^b(t',t)\) take the value
\(t'\) at \(b\) and \(t\) at every other positive bond.
The term assigned to this bond will be
\begin{equation}\label{eq:cutoff-bond-definition}
 E^\chi(X,b)=\int_0^1
   \left.\partial_{t'}E[\underline{\mathbb V};
                          \mathbf t^b(t',t)](X)\right|_{t'=t}\,dt.
\end{equation}
The derivative in this formula is an almost-everywhere
derivative. Its existence requires an argument because
the tests are characteristic functions.

\begin{lemma}\label{II-20-71}
Let \(\Omega\) be a countable disjoint union of finite-dimensional
Lebesgue spaces, let \(d\in L^1(\Omega;\mathbb C)\), and let
\(\theta:\Omega\to(0,1]\), \(\theta':\Omega\to[0,1]\)
be measurable. Suppose \(\bar k\ge0\) is integrable on
\((0,1)\),
\[
 \int_\Omega|d|\mathbf1_{\{t_1<\theta\le t_2\}}
       \le\int_{t_1}^{t_2}\bar k(v)\,dv
       \quad(0\le t_1<t_2\le1),\qquad
 \int_\Omega|d|\mathbf1_{\{\theta=\theta'\}}=0.
\]
Set
\[
 H(t',t)=\int_\Omega d\mathbf1_{\{\theta\le t'\}}
                              \mathbf1_{\{\theta'\le t\}},
 \qquad
 S(t)=\int_\Omega d\mathbf1_{\{\theta'\le\theta\le t\}}.
\]
Then \(S\) is absolutely continuous, \(S(0)=0\), and
\(\partial_{t'}H(t,t)=S'(t)\) for almost every \(t\in(0,1)\).
Consequently,
\begin{equation}\label{II-20-73}
 \int_0^1\partial_{t'}H(t,t)\,dt
             =\int_\Omega d\mathbf1_{\{\theta'\le\theta\}}.
\end{equation}
\end{lemma}

\begin{proof}
Put
\[
 P_n(t)=\int_\Omega|d|\mathbf1_{\{\theta\le t\}}
                           \mathbf1_{\{|\theta-\theta'|\le1/n\}}.
\]
The marginal bound shows that \(S\) and every \(P_n\)
are absolutely continuous and vanish at zero. Moreover,
\(P_n\) and \(P_n-P_{n+1}\) are nondecreasing, so
\(P_n'\ge P_{n+1}'\ge0\) almost everywhere. Dominated
convergence and the tie hypothesis give
\[
 \int_0^1P_n'(t)\,dt=P_n(1)
       \longrightarrow\int_\Omega|d|
                              \mathbf1_{\{\theta=\theta'\}}=0.
\]
Outside one null set, \(S\) and all \(P_n\) are
differentiable and \(P_n'(t)\downarrow0\).

Fix such a \(t\). For \(0<h\le1/n\), the difference
between the increments \(H(t+h,t)-H(t,t)\) and
\(S(t+h)-S(t)\) is supported where
\(t<\theta\le t+h\) and \(t<\theta'\le\theta\).
Its modulus is at most \(P_n(t+h)-P_n(t)\).
For \(-1/n\le h<0\), the corresponding difference is
supported where \(t+h<\theta\le t\) and
\(\theta<\theta'\le t\), and its modulus is at most
\(P_n(t)-P_n(t+h)\). Thus
\[
 \limsup_{h\to0}
 \left|\frac{H(t+h,t)-H(t,t)}h-S'(t)\right|
                      \le P_n'(t)
 \quad\hbox{for every }n.
\]
Letting \(n\to\infty\) proves the derivative assertion.
The fundamental theorem for \(S\) gives \eqref{II-20-73}.
\end{proof}

\begin{theorem}[Cutoff variation]\label{II-20-87}
In the preceding setting, \eqref{eq:cutoff-bond-definition}
is defined for every \(X\in\mathcal D_t\), \(X\subset R\),
and every positive bond \(b\). The diagonal function
\(v\mapsto E[\underline{\mathbb V};(v,\ldots,v)](X)\)
is absolutely continuous, and
\begin{align}
 E(X)-E^\circ(X)&=\sum_{b\in\mathfrak X^+}E^\chi(X,b),
                                                   \label{II-20-88}\\
 |E^\chi(X,b)|&\le K^\Pi_P e^{4\lambda_1}\mathsf p
                          e^{-\lambda_{\rm out}d_t(X)}.
                                                   \label{II-20-89}
\end{align}
The term is zero unless \(b^\sharp\subset X\), and it
inherits locality, orthogonal covariance, and conjugation
from the Gaussian data and the potentials. Under the
joint measurability and holomorphy hypotheses of
Proposition~\ref{prop:cluster-parameters}, with uniform
constants and a fixed threshold, it is holomorphic in
the parameters.
\end{theorem}

\begin{proof}
Fix \(X\). A refined cluster index is
\(\alpha=(n;(\gamma_i);(\iota_i))\), where
\(\iota_i=(\mathbf D_i,\mathfrak e_i)\), and its coefficient
is \(c_\alpha=\rho^T(\gamma_1,\ldots,\gamma_n)/n!\).
On the product of the real integration spaces, write
\(\omega=(\mathbb X_i,B_i)_{i=1}^n\) and define the density
\[
 d_\alpha(\omega)=c_\alpha\prod_{i=1}^n
       g^{\gamma_i}(Z_\circ(\iota_i);\mathbb X_i,B_i)
                                             F^{\mathbf1}_{\iota_i}(B_i).
\]
Here \(\mathbf1\) denotes the family of bond parameters
equal to one. Lemma~\ref{II-20-33} and the Mayer majorants
give
\begin{equation}\label{eq:cutoff-density-majorant}
 \int|d_\alpha|\le\bar D_\alpha
       :=|c_\alpha|\prod_i\overline T(\gamma_i,\iota_i),
 \qquad\sum_\alpha\bar D_\alpha<\infty.
\end{equation}

Let \(Q_\alpha=\{(i,b):b\in\mathfrak e_i^+\}\).
Define \(\theta(y)=\min\{1,\mathfrak b/|y|\}\) for
\(y\ne0\) and \(\theta(0)=1\), and put
\(\theta_{i,b}=\theta(B_i(b))\). Since all origins have
sort \(\Sigma_{\rm q}\), the only tests in a Mayer
integrand are the out tests. Factoring them from their
values at \(\mathbf1\) yields
\begin{equation}\label{II-20-75}
 E[\underline{\mathbb V};\mathbf t](X)
       =\sum_\alpha\int d_\alpha(\omega)
             \prod_{(i,b)\in Q_\alpha}
                          \mathbf1_{\{\theta_{i,b}\le t_b\}}\,d\omega.
\end{equation}
In particular the density in this formula is independent
of the bond parameters.

We verify the marginal and tie hypotheses of
Lemma~\ref{II-20-71}. Fix a test pair \((i,b)\) and
integrate \(|d_\alpha|\) over all variables except
\(y=B_i(b)\in\mathbb R^3\). Denote the resulting
nonnegative integrable density by \(d^\perp_{\alpha,i,b}(y)\).
It is supported in \(|y|\ge\mathfrak b\). Polar
coordinates followed by \(r=\mathfrak b/v\) show that
the marginal of \(\theta_{i,b}\) has density
\[
 \bar k_{\alpha,i,b}(v)
       =\mathfrak b^3v^{-4}
          \int_{S^2}d^\perp_{\alpha,i,b}
                              ((\mathfrak b/v)u)\,d\sigma(u),
 \qquad 0<v<1.
\]
Its integral is \(\int|d_\alpha|\); the sphere
\(|y|=\mathfrak b\) is null. For distinct test pairs
the variables are distinct coordinate factors, even when
the physical bonds agree and the cluster indices differ.
The set on which their norms agree is null by Fubini,
since a sphere in \(\mathbb R^3\) is null. Thus ties
between distinct ratios have zero \(|d_\alpha|\)-mass.

Let \(\mathcal A^b(X)\) consist of the indices with at
least one test at \(b\). On each such space put
\[
 \theta^b=\max_{i:(i,b)\in Q_\alpha}\theta_{i,b},
 \qquad
 \theta'^{\,b}=\max_{(i,c)\in Q_\alpha,\ c\ne b}\theta_{i,c},
\]
where the second maximum is zero if its set is empty.
The \(\theta^b\)-marginal is dominated by
\(\bar k^b=\sum_{\alpha\in\mathcal A^b(X)}
                   \sum_{i:(i,b)\in Q_\alpha}\bar k_{\alpha,i,b}\).
It is integrable: its integral is bounded by the refined
absolute cluster sum marked at the single bond \(b\),
namely \(\overline E^{\{b\}}(X)\).
The tie \(\theta^b=\theta'^{\,b}\) is contained in a
finite union of the pairwise ties in each integration
space, and hence has zero mass.

Apply Lemma~\ref{II-20-71} on the disjoint union of
the spaces indexed by \(\mathcal A^b(X)\), with density
\(d_\alpha\), and set
\[
 S_b(v)=\sum_{\alpha\in\mathcal A^b(X)}
       \int d_\alpha\mathbf1_{\{\theta'^{\,b}\le\theta^b\le v\}}.
\]
The terms of \eqref{II-20-75} outside \(\mathcal A^b(X)\)
are independent of \(t'\). The lemma therefore proves
\eqref{eq:cutoff-bond-definition} and the useful formula
\begin{equation}\label{II-20-79}
 E^\chi(X,b)=S_b(1)
       =\sum_{\alpha\in\mathcal A^b(X)}
                    \int d_\alpha
                          \mathbf1_{\{\theta'^{\,b}\le\theta^b\}}.
\end{equation}
Each \(S_b\) is absolutely continuous and vanishes at
zero. Outside the null tie sets there is a unique largest
ratio among all test pairs. Assigning it to its physical
bond in \eqref{II-20-75} gives
\[
 E[\underline{\mathbb V};(v,\ldots,v)](X)
                         =E^\circ(X)+\sum_bS_b(v).
\]
The no-test terms are exactly \(E^\circ\); at zero every
term with a test vanishes. This proves absolute continuity
and \eqref{II-20-88}.

Every term containing a test at \(b\) has
\(b^\sharp\subset X\), which proves the support assertion.
Moreover, \eqref{II-20-79} and the marked bound give
\[
 |E^\chi(X,b)|\le\overline E^{\{b\}}(X)
       \le K_E\eta^{\{b\}}e^{-\lambda_{\rm out}d_t(X)}.
\]
The bond weight is
\(\eta^{\{b\}}\le\mathsf p e^{\lambda_1(d_t(b^\sharp)+4)}
=\mathsf p e^{4\lambda_1}\), since
\(d_t(b^\sharp)=0\). As \(K_E=K^\Pi_P\), this is
\eqref{II-20-89}. The locality, covariance, and conjugation
identities for \(E[\mathbf t]\) hold for every bond
parameter family. They therefore persist under the real
difference quotients and the integral in
\eqref{eq:cutoff-bond-definition}.

For holomorphy use \eqref{II-20-79}. The set selected
by its indicator is independent of the complex parameters.
For each fixed \(\alpha\), the integrand is jointly Borel,
pointwise holomorphic in the parameters, and its \(L^1\)
norm is bounded by \(\bar D_\alpha\) uniformly in them.
These facts imply holomorphy of its integral. Indeed,
on a polydisc of radius \(2r\) about a point \(a\),
the iterated Cauchy formula expresses its integral at
\(z\), \(|z_j-a_j|<r\), as an integral over the boundary
torus with kernel
\[
 \prod_{j=1}^N
       \frac{2re^{i\vartheta_j}}
                     {a_j+2re^{i\vartheta_j}-z_j}.
\]
This kernel has modulus at most \(2^N\).
The uniform \(L^1\) bound justifies Fubini before the
field integration. The resulting boundary integral is
holomorphic in \(z\), by the same bound on each smaller
polydisc. Finally \eqref{eq:cutoff-density-majorant}
gives uniform convergence of the series of holomorphic
integrals. The exceptional set of cutoff parameters where
the original derivative fails to exist may vary with
the complex parameter; the fixed-set formula
\eqref{II-20-79} removes it from this argument.
\end{proof}

\subsection{Interpolation in the field variables}

Fix a finite threshold \(\mathfrak b\) and
\(\mathsf C_{\mathcal A}\ge3\). For a fine domain \(Y\), put
\(J_Y=\mathfrak X^{\circ+}\langle Y\rangle\), and let
\begin{equation}\label{II-20-92}
 \mathbb D^{\rm f}(Y)=
       \{B:|B(b)|<\mathsf C_{\mathcal A}\mathfrak b
                                      \text{ for every }b\in J_Y\}.
\end{equation}
The fields are complex and the norm at each bond is the matrix
operator norm. For a family holomorphic on these polydiscs define
\begin{equation}\label{II-20-93}
 \|\mathbb V\|^{\rm f}_{\star,r}
   =\sup_Y\sup_{B\in\mathbb D^{\rm f}(Y)}
       \frac{e^{r d_k(Y)}|\mathbb V(Y;B)|}{Q_Y(B)},
 \qquad Q_Y(B)=1+M^{-3}\sum_{b\in J_Y}|B(b)|^2.
\end{equation}
Its restriction to real fields with \(|B(b)|<\mathfrak b\)
has sort \(\Sigma_{\rm v}\), with no larger norm.

Suppose \(\mathbb V^{\rm J},\mathbb V^{\rm G},\mathbb V^{\rm A},
\mathbb V'\) are such families. For bond parameters
\(\mathbf v=(v_b)\in[0,1]^{\mathfrak X^+}\), define
\((\mathbf v\cdot B)(b)=v_bB(b)\) at positive bonds, and set
\begin{equation}\label{II-20-95}
\begin{split}
 \mathbb V^{\rm int}_{\mathbf v}(Y;B)
   &=\mathbb V^{\rm J}(Y;\mathbf v\cdot B)\\
   &\quad+\int_0^1(1-\sigma)
      \left.\partial_\epsilon^2
       [\mathbb V^{\rm G}+\tau\mathbb V^{\rm A}]
            (Y;\sigma\mathbf v\cdot B+\epsilon B)
                                             \right|_{\epsilon=0}\,d\sigma,\\
 \mathbb V'_{\mathbf v}(Y;B)&=\mathbb V'(Y;\mathbf v\cdot B).
\end{split}
\end{equation}
Write \(\mathbb V_v\) for the restriction to the diagonal
\(\mathbf v=(v,\ldots,v)\), and put
\[
 \dot{\mathbb V}^{\rm int}_{v,b}
       =\left.\partial_{v_b}\mathbb V^{\rm int}_{\mathbf v}
                                             \right|_{\mathbf v=(v,\ldots,v)},
 \qquad
 \dot{\mathbb V}'_{v,b}
       =\left.\partial_{v_b}\mathbb V'_{\mathbf v}
                                             \right|_{\mathbf v=(v,\ldots,v)}.
\]
These parameters occur in the potentials; all cutoff parameters
in the cluster sums remain equal to one.

\begin{proposition}\label{II-20-103}
Assume the four norms in \eqref{II-20-93} are finite,
\(\mathbb V^{\rm J}(Y;0)=\mathbb V'(Y;0)=0\), and
\(\mathbb V^{\rm G},\mathbb V^{\rm A}\) vanish through degree
two at zero. Put
\begin{equation}\label{II-20-104}
 n_{
 \rm int}=\|\mathbb V^{\rm J}\|^{\rm f}_{\star,r}
       +\|\mathbb V^{\rm G}\|^{\rm f}_{\star,r}
       +|\tau|\|\mathbb V^{\rm A}\|^{\rm f}_{\star,r},
 \qquad n'=\|\mathbb V'\|^{\rm f}_{\star,r}.
\end{equation}
The interpolated potentials are defined and jointly continuous
in \((v,B)\), and are \(C^1\) in \(v\in[0,1]\) at every
admitted real field. They vanish at \(v=0\); for \(v>0\),
\begin{equation}\label{II-20-105}
 \mathbb V^{\rm int}_v(Y;B)=\mathbb V^{\rm J}(Y;vB)
       +v^{-2}[\mathbb V^{\rm G}+\tau\mathbb V^{\rm A}](Y;vB).
\end{equation}
For every \(v\in[0,1]\), these diagonal expressions are
holomorphic on \(\mathbb D^{\rm f}(Y)\) and satisfy
\[
 \|\mathbb V^{\rm int}_v\|^{\rm f}_{\star,r}\le n_{\rm int},
 \qquad\|\mathbb V'_v\|^{\rm f}_{\star,r}\le n'.
\]
The insertions are Borel in \((v,B)\), vanish unless
\(b\in J_Y\), and obey
\begin{align}
 \partial_v\mathbb V^{\rm int}_v
        &=\sum_b\dot{\mathbb V}^{\rm int}_{v,b},\qquad
 \partial_v\mathbb V'_v=\sum_b\dot{\mathbb V}'_{v,b},
                                                   \label{II-20-106}\\
 \|\dot{\mathbb V}^{\rm int}_{v,b}\|_{\star,r}
        &\le9n_{\rm int},\qquad
 \|\dot{\mathbb V}'_{v,b}\|_{\star,r}\le9n',
                                                   \label{eq:bond-interpolation-bound}\\
 \|\partial_v\mathbb V^{\rm int}_v\|_{\star,r}
        &\le4\|\mathbb V^{\rm J}\|^{\rm f}_{\star,r}
         +\tfrac98\bigl(\|\mathbb V^{\rm G}\|^{\rm f}_{\star,r}
                      +|\tau|\|\mathbb V^{\rm A}\|^{\rm f}_{\star,r}\bigr),
                                                   \label{II-20-107}\\
 \|\partial_v\mathbb V'_v\|_{\star,r}&\le4n'.\notag
\end{align}
The construction is linear in the four input families, with
\(\tau\mathbb V^{\rm A}\) regarded as one input. It preserves
locality, orthogonal covariance, conjugation, and holomorphic
dependence on additional parameters under uniform norm bounds.
\end{proposition}

\begin{proof}
Fix \(Y\), a real field below the threshold, and one of the
four potentials \(\mathbb V\). Write
\[
 f(\mathbf x,B)=\mathbb V(Y;\mathbf x\cdot B),\qquad
 s(\mathbf x,B)=\left.\partial_\epsilon^2
                 f(\mathbf x+\epsilon\mathbf1,B)\right|_{\epsilon=0},
 \qquad g_B(z)=\mathbb V(Y;zB).
\]
These are holomorphic wherever the fields in their definitions
belong to the polydisc. In particular \(f\) is holomorphic
on a neighborhood of \(|x_b|\le3\), since the real field
satisfies the strict bound \(|B(b)|<\mathfrak b\).
Put \(A=\|\mathbb V\|^{\rm f}_{\star,r}e^{-r d_k(Y)}Q_Y(B)\).
For \(|x_b|\le c\), \(c\ge1\), the norm bound gives
\(|f(\mathbf x,B)|\le c^2A\).

For \(\mathbf x\in[0,1]^{J_Y}\), Cauchy's estimate on
\(|\zeta|=1\) in \(f(\mathbf x+\zeta e_b,B)\) gives
\(|\partial_bf|\le4A\). Applied twice in the common
direction \(\epsilon\mathbf1\), it gives \(|s|\le8A\).
On \(|\epsilon|,|\zeta|\le1\), all coordinates of
\(\mathbf x+\zeta e_b+\epsilon\mathbf1\) have modulus
at most three. Cauchy's estimates in these two variables
therefore give \(|\partial_bs|\le18A\).
For \(x\in[0,1]\), the circles of radii one and two about
\(x\) give
\begin{equation}\label{II-20-100}
 |g'_B(x)|\le4A,\qquad |g'''_B(x)|\le\frac{3!}{2^3}9A
                                                     =\frac{27}4A.
\end{equation}

The integral in \eqref{II-20-95} is thus well defined.
On the diagonal, integration by parts gives
\[
 \int_0^1(1-\sigma)g_B''(\sigma v)\,d\sigma
          =v^{-2}\{g_B(v)-g_B(0)-vg_B'(0)\},\qquad v>0.
\]
This proves \eqref{II-20-105}. At zero the integral is
\(g_B''(0)/2=0\). For a complex field in the polydisc,
\(g_B\) is holomorphic on a disc of radius greater than
one and has modulus at most \(A\) on \(|z|\le1\).
If it vanishes to order three, successive applications
of Schwarz's lemma to \(g_B(z),g_B(z)/z,g_B(z)/z^2\)
give \(|g_B(v)|\le A v^3\). Hence
\(|v^{-2}g_B(v)|\le Av\), including its zero value at
\(v=0\). This proves the two polydisc norm bounds.
Holomorphy for \(v>0\) follows by composition, and at
\(v=0\) both functions are zero.

The integral expression is continuously differentiable in
\(v\), including at zero. Its derivative for either of
the cubic inputs is
\(\int_0^1(1-\sigma)\sigma g'''_B(\sigma v)\,d\sigma\).
Equation~\eqref{II-20-100} and
\(\int_0^1(1-\sigma)\sigma\,d\sigma=1/6\) give
the factor \(9/8\) in \eqref{II-20-107}.
Differentiating in one bond parameter instead replaces
\(g'''_B\) by \(\partial_bs\) in this formula. Its
bound gives
\[
 |\dot{\mathbb V}^{\rm int}_{v,b}(Y;B)|
   \le4A_{\rm J}+3A_{\rm G}+3|\tau|A_{\rm A},
 \qquad |\dot{\mathbb V}'_{v,b}(Y;B)|\le4A'.
\]
These bounds are stronger than
\eqref{eq:bond-interpolation-bound}. They are locally
bounded in \(B\), so differentiation and continuity
under the \(\sigma\)-integral follow by dominated
convergence. Thus all insertions are Borel. Parameters
outside \(J_Y\) do not occur, and the chain rule on the
diagonal gives \eqref{II-20-106}.

Linearity and locality hold in each defining operation.
For covariance, the complex extension of a real orthogonal
map of \(\mathfrak g\) preserves the matrix norm on
\(\mathfrak g^c\). To check this, write \(Z=A+iA'\)
with \(A,A'\) Hermitian and traceless. Cayley--Hamilton
and polarization give
\[
 |Z|^2=\langle A,A\rangle+\langle A',A'\rangle
       +2\sqrt{\langle A,A\rangle\langle A',A'\rangle
                                      -\langle A,A'\rangle^2}.
\]
Indeed \(Z^*Z=(\langle A,A\rangle+\langle A',A'\rangle)I
+i[A,A']\), and the traceless Hermitian last term has
eigenvalues equal to the positive and negative of the
displayed square-root term. This also shows \(|Z^*|=|Z|\).
Thus the transformed fields stay in their polydiscs;
bond multiplication commutes with the orthogonal maps,
after the bonds have been relabeled. Covariance and
conjugation now pass through the derivatives and real
integrals. Joint holomorphy in further parameters follows
from the same Cauchy formulas and the locally uniform
bounds. If \(J_Y\) is empty, all four inputs vanish
at its sole field, and every assertion holds with value zero.
\end{proof}

\section{Domains and local terms of the effective action}
\label{sec:effective-local-terms}

The fluctuation expansion will be iterated on domains which
control both the original field and its backgrounds on smaller
towers. We distinguish these analytic cores from the pointwise
domains used for the Gaussian construction.

\subsection{Analytic cores and restriction}

Assume \(R_1\ge L\), \(M_1\ge R_1\mathfrak M\), and
\(M\ge M_1\), with these scales powers of \(L\).
Let \(1\le j\le K\), \(m'<m+K-j+1\), and
\(X\in\mathcal D_{j-1}\). Use level-\(j\) units,
so the bare spacing is \(\xi=\eta_j=L^{-j}\).
A triple is \(P=(U,A',\mathbf J)\) on
\(\operatorname{Bd}(X)\), with \(U\) real and the two
forms complex. Its pair is
\begin{equation}\label{II-16-9}
 \mathbf p(P)=(\mathbf U,\mathbf J),\qquad
 \mathbf U=e^{i\eta_jA'}U.
\end{equation}
The forms are specified at positive bonds; the expression
for \(\mathbf U\) is extended by reversal as a configuration.
For a real gauge transformation on all endpoints of these
bonds put
\[
 P^u=(U^u,\operatorname{Ad}(u)A',\operatorname{Ad}(u)\mathbf J),
 \qquad P^{\mathfrak c}=(U,A'^*,\mathbf J^*).
\]
Adjoints are taken bondwise at positive bonds. A triple is
real when both forms are Hermitian. Its conjugate pair has
configuration \((\mathbf U^*)^{-1}\).

\begin{definition}\label{II-16-12}
The simple domain \(\mathcal Q_j(X;b_0,b_1,c_0)\) consists
of triples satisfying
\begin{enumerate}
\item \(|\partial U-1|<b_0\eta_j^2\), and every cubic
box in \(X\) of side at most \(\mathsf oLM\) admits a
real gauge \(U^w=e^{i\eta_jA}\) on its inner bonds with
\(|A|,|\nabla^{\eta_j}A|<\mathsf oLMB^{(\rm cg)}b_0\);
\item \(|A'|,|\nabla_U^{\eta_j}A'|<b_1\);
\item \(|\partial\mathbf U-1|<b_0\eta_j^2\) and
\(|\mathbf J|<c_0\).
\end{enumerate}
Except for the explicitly inner-bond gauge bounds, each
inequality is required wherever its stencil lies in
\(\operatorname{Bd}(X)\). Boxes are unions of bare cubes,
with an injective rectangular lift of equal side lengths.

The core \(\mathcal U^\circ_j(X;b_0,b_1,c_0)\) is the
subset satisfying the following additional conditions.
For every \(1\le p\le j\), every \(Z\in\mathcal D_{p-1}\)
with \(Z\subset X\), both tower rules
\(\mathfrak r_1,\mathfrak r_2\), and every truncation
\(1\le n\le p\), take the reduced determining set of
the tower of \(Z\), with all inward bonds omitted.
The averaged complex datum must exist. Relative to the
averaged real datum, its positive-bond quotients must lie
in the logarithm ball \(|V'V^{-1}-1|<1\), and their
logarithms must lie strictly inside a polydisc on which
the continuation of the real background exists.
Write \(U_n^{\mathfrak r}(Z;M'(\mathbf U))\) for this
continued background and \(J_n^{\mathfrak r}\) for its
current at spacing \(\eta_n\). With
\[
 c=b_0L^{-(j-p)},\qquad
 Z^{\wr p}=\bigcup\{\square\in\pi_{p-1}:
                          \square^{\sim2}\subset Z\},
\]
require
\begin{equation}\label{eq:analytic-core-tower-bounds}
 |\partial U_n^{\mathfrak r}-1|<cL^{-2p},\qquad
 |J_n^{\mathfrak r}|<cL^{2(n-p)}\quad\text{on }Z^{\wr p}.
\end{equation}
Require the same existence and bounds for the datum from
\(U\). The real and complex configurations in these
conditions are the ones in the specified decomposition.
When \(c_0=b_0\), omit the final radius. A pair lies in
the core if it is the pair of a triple in the core.
\end{definition}

In these tower conditions the logarithmic perturbation is
\(B_l(b)=i^{-1}\log[M'(\mathbf U)_l(b)M'(U)_l(b)^{-1}]\).
The strict continuation requirement means that some radius
\(r>\sup_{l,b}|B_l(b)|\) is admitted. The resulting background
is independent of that radius by continuation uniqueness.
The conditions retain the existence of the backgrounds;
pointwise curvature estimates alone do not assert it.

For \(j\le j'\), \(X\subset Y\), and
\(\lambda=L^{-(j'-j)}\), define
\begin{equation}\label{II-16-16}
 \rho_X^{j'\to j}(U,A',\mathbf J)
       =(U|_X,\lambda A'|_X,\lambda^3\mathbf J|_X).
\end{equation}
The pair has unchanged configuration and current scaled
by \(\lambda^3\). At the same level write \(P|_X\).

\begin{lemma}\label{II-16-46}
The simple domains and cores increase with each radius
and pass to smaller domains at the same level. More
generally, if \(1\le j\le j'\le K\),
\(m'<m+K-j'+1\), \(X\in\mathcal D_{j-1}\),
\(Y\in\mathcal D_{j'-1}\), and \(X\subset Y\), then
\begin{equation}\label{II-16-48}
 P\in\mathcal U^\circ_{j'}(Y;b_0,b_1,c_0)
 \quad\Longrightarrow\quad
 \rho_X^{j'\to j}P\in
       \mathcal U^\circ_j(X;\lambda b_0,\lambda b_1,\lambda^3c_0).
\end{equation}
The same implication holds with \(\mathcal Q\) in place
of \(\mathcal U^\circ\).
\end{lemma}

\begin{proof}
Increasing a radius weakens each defining inequality;
the indices of the tower conditions do not change.
For restriction, all stencils and boxes of the smaller
domain are stencils and boxes of the larger one, and
the smaller tower index set is a subset of the larger.

For scaling put \(\xi'=\eta_{j'}=\lambda\eta_j\).
The complex configuration is unchanged because
\(\eta_j\lambda A'=\xi'A'\). Curvature bounds acquire
the factor \(\lambda^2\le\lambda\). A box of side
at most \(\mathsf oLM\) at level \(j\) has side at most
that at level \(j'\); keep its gauge and replace its
potential \(A\) by \(\lambda A\). Potentials scale
by \(\lambda\) and their derivatives by \(\lambda^2\).
The same calculation applies to \(A'\), and the current
coordinate has the required factor \(\lambda^3\).
This proves the simple-domain assertion.

For \((p,Z)\) in a lower-level tower condition, the real
and complex configurations on \(\operatorname{Bd}^\circ(Z)\)
are unchanged. Thus the averaged data, logarithms, and
continued backgrounds are the same. Their required
constant is also unchanged:
\((\lambda b_0)L^{-(j-p)}=b_0L^{-(j'-p)}\).
Every required tower condition therefore follows from
the corresponding condition at level \(j'\).
\end{proof}

\subsection{Local terms and their norms}

\begin{definition}\label{II-16-13}
For fixed real \(U\), the slice of a core is
\(\mathcal O_U=\{(A',\mathbf J):(U,A',\mathbf J)
                         \in\mathcal U^\circ_j(X;b_0,b_1,c_0)\}\).
A function on the core is analytic if near each point
of every slice it is the restriction of a holomorphic
function on an open neighborhood in the vector space
of the two forms. It is a function of the pair if its
value is independent of the admitted decomposition.
It is gauge invariant if \(E(P^u)=E(P)\) whenever both
triples lie in the core, and conjugation real if
\(E(P^{\mathfrak c})=\overline{E(P)}\) whenever both
triples lie there.
\end{definition}

\begin{definition}\label{II-16-29}
For \(m'<m+K-j\) and \(X\in\mathcal D_j\), a local
term of level \(j\) at radii \((b_0,b_1)\) is a
bounded, analytic, gauge-invariant function of the pair
on \(\mathcal U^\circ_j(X;b_0,b_1)\). Its norm is
\(\|E\|_X=\sup_{P\in\mathcal U^\circ_j(X;b_0,b_1)}|E(P)|\),
with the supremum of an empty set equal to zero.
\end{definition}

\begin{definition}\label{II-16-31}
A pointed family consists of local terms \(E(X,\cdot,z)\)
for \(X\in\mathcal D_j\) and \(z\in X\cap\mathbb T^{(j)}\).
A family without base point consists of local terms
\(H(X,\cdot)\) for \(X\in\mathcal D_j\). Their norms are
\begin{equation}\label{II-16-32}
 \|E\|_{j,\kappa}=\sup_{X,z}e^{\kappa d_j(X)}
                                      \|E(X,\cdot,z)\|_X,
 \qquad
 \|H\|_{j,\kappa}=\sup_Xe^{\kappa d_j(X)}\|H(X,\cdot)\|_X.
\end{equation}
Unspecified members are zero. Evaluation at a real
configuration means evaluation at its pair
\((U,J^{(j)}(U))\), provided this pair lies in the core.
Evaluation at \(\mathbf1\) means evaluation at \((\mathbf1,0)\),
with the same requirement.
\end{definition}

\subsection{Assigning local terms to sites}

For \(0\le j\le K\), use level-\(j\) units and put
\(q_j=2L^{m+K-j}\). Periodize the product hat function:
\begin{equation}\label{II-20-6}
 \hat h(u)=\prod_{\nu=1}^3(1-|u_\nu|)_+,
 \qquad h_z^{(j)}(y)=\sum_{n\in\mathbb Z^3}
                         \hat h(\hat y-\hat z-q_jn).
\end{equation}
Here \(\hat y,\hat z\) are lifts, with \(\hat z\in\mathbb Z^3\).
Evaluate a tent at the midpoint of a bond and at the
barycentre of a plaquette. At step \(k+1\) write
\(h_z=h_z^{(k+1)}\).

\begin{proposition}\label{II-20-13}
The tents are well defined, take values in \([0,1]\),
and sum to one. On every bond of a level \(i\le j\),
\[
 h_z^{(j)}(b)=h_z^{(j)}(\bar b)
           =\tfrac12\{h_z^{(j)}(b_-)+h_z^{(j)}(b_+)\}.
\]
They commute with translations by level-\(j\) lattice
vectors and with signed coordinate permutations.
At step \(k+1\), in level-\(k\) units,
\[
 h_z(b)\ne0\ \Longrightarrow\
 |x_b-z|_\infty<L,\quad |b_\pm-z|_\infty\le L,
 \qquad\sum_{b\in\mathfrak X_0^+}h_z(b)=3L^3.
\]
If \(M\ge L\), then \(h_z(b)\ne0\) also implies
\(z\in b^\sharp\). Consequently, for numbers
\(E(X,b)\) vanishing unless \(b^\sharp\subset X\),
\begin{equation}\label{II-20-14}
 \sum_{z\in X\cap\mathbb T^{(k+1)}}\sum_{b\in\mathfrak X^+}
           h_z(b)E(X,b)=\sum_{b\in\mathfrak X^+}E(X,b).
\end{equation}
\end{proposition}

\begin{proof}
The period is at least \(2L\), so at most one term in
\eqref{II-20-6} is nonzero. Changing a lift reindexes
the sum. In one dimension the translates of
\(\hat h_1(u)=(1-|u|)_+\) sum to one; taking products
proves the partition identity. The hat is affine on
each unit interval. A level-\(i\) bond, expressed in
level-\(j\) units, lies in such an interval in its
varying coordinate, because \(L^{j-i}\) is an integer.
This proves the midpoint identity. Evenness and the
product formula prove the isometry assertions.

The support gives the strict midpoint distance. Its
varying coordinate is a half-integer in level-\(k\)
units, and the other coordinates are integers. Moving
by one half to either endpoint therefore gives the
stated endpoint bound. If \(M\ge L\), the centers of
the fine cubes containing \(z\) and \(b_-\) are at
distance less than \(M+L\le2M\). Their coordinate
differences are multiples of \(M\); hence these cubes
touch, and \(z\in\square(b)^\sim\subset b^\sharp\).

For any positive integer \(N\), summing the two linear
pieces of the hat gives
\[
 \sum_{n\in\mathbb Z}\hat h_1(n/N)
 =\sum_{n\in\mathbb Z}\hat h_1((n+1/2)/N)=N.
\]
At the midpoints of the bonds of a fixed direction,
one coordinate uses the shifted sum and two use the
unshifted sum, with \(N=L\). Their product is \(L^3\).
Summing the three directions proves the bond sum.
Finally every contributing tent site belongs to \(X\),
and its sum over those sites is one. Interchanging
the two finite sums proves \eqref{II-20-14}.
\end{proof}

For \(z\in\mathbb T^{(j)}\), let \(\square_z\) be its
unique cube in \(\pi_j\) and put \(X_z=\square_z^\sim\).
For \(X\in\mathcal D_j\), define
\begin{equation}\label{II-20-16}
 \mathsf A^{\mathsf t}(X,P,z)=
 \begin{cases}
  \displaystyle\eta_j^{-1}\sum_p h_z^{(j)}(p)a_p(\mathbf U),
                                  &X=X_z,\\
  0,                              &X\ne X_z.
 \end{cases}
\end{equation}
Only plaquettes with nonzero tent weight enter the sum.

\begin{lemma}\label{II-20-18}
For \(1\le j\le K\), \(m'<m+K-j\), and \(M\ge L\),
\eqref{II-20-16} is a pointed family at radii \((b_0,b_1)\),
vanishes at the trivial triple, and satisfies
\[
 \|\mathsf A^{\mathsf t}\|_{j,\kappa}
              \le\tfrac32b_0^2e^{7\kappa}.
\]
On a whole-lattice configuration at which the terms
are defined, their sum over \((X,z)\) is
\(A^{\eta_j}(\mathbf U)\).
\end{lemma}

\begin{proof}
The cubes touching \(\square_z\) are joined to it
by changing their coordinate offsets one at a time,
so \(X_z\in\mathcal D_j\). A tree of seven edges
joining the eight corners of \(\square_z\) meets
all these cubes. Hence \(d_j(X_z)\le7\).
A plaquette with nonzero tent weight has barycentre
within one of \(z\); its corners are within
\(1+\eta_j/2<L\le M\). They therefore lie in
\(X_z\), and all four bonds are available there.

For a fixed plaquette direction, the same product
calculation as above, now with two shifted coordinates
and \(N=L^j\), gives total weight \(L^{3j}\).
Thus \(\sum_ph_z^{(j)}(p)=3\eta_j^{-3}\).
For a complex matrix \(V\) of determinant one,
\(2-\operatorname{tr}V=\det(V-1)\) and
\(|\det(V-1)|\le|V-1|^2\). The core curvature bound
therefore gives
\[
 |a_p(\mathbf U)|\le\tfrac12b_0^2\eta_j^4,
 \qquad
 |\mathsf A^{\mathsf t}(X_z,P,z)|\le\tfrac32b_0^2.
\]
This and \(d_j(X_z)\le7\) prove the norm bound.
The finite sum is holomorphic in \(A'\) on every
slice, depends only on \(\mathbf U\), and is gauge
invariant by invariance of the plaquette trace.
It is therefore a local term. At the trivial triple
every plaquette action vanishes. Finally summing
the partition of unity at each plaquette proves
the action identity.
\end{proof}

\section{Local Gaussian fluctuation fields}\label{sec:local-gaussian-fields}
The fluctuation form has long range, even when the background is known only
locally. We first separate it into connected local contributions. A second
expansion, obtained from local resolvents, does the same for its inverse
square root. Inserting one parameter for each coarse cube then gives both
the covariance and the coupling to an independent Gaussian field. The
positivity estimate below is uniform in the region; the final block identity
identifies the resulting integral with the original fluctuation integral.

Fix a step $t=k+1$, $0\le k<K$, and use level-$k$ units: the spacings are
$\xi_n=L^{n-k}$, the bare spacing is $\xi=L^{-k}$, and the torus has side
$\ell_k=2L^{m+K-k}$. The sequence is the constant sequence through level $t$.
Write $\mathfrak X$ for its nonpivot level-$k$ bonds, including both
orientations. All kernels below use counting measure on positive bonds and
the pairing $\frac12\operatorname{tr}(XY)$. The transpose is bilinear;
$|\cdot|_2$ denotes the operator norm for the Hermitian norm on
$\mathfrak g^c$. Put $|b-b'|=|b_--b'_-|_\infty$ and
\[
 \|T\|_\theta^{\rm w}
 =\max\left\{\sup_b\sum_{b'}|T(b,b')|_2e^{\theta|b-b'|},
              \sup_{b'}\sum_b|T(b,b')|_2e^{\theta|b-b'|}\right\}.
\]
These norms dominate the Hilbert-space norm and are submultiplicative.
Restrictions and extensions by zero are denoted by $\chi$ and
$\chi^{\mathsf T}$.

The coarse cubes have side $LM$, where $M=L^{m'}$ and
$m'<m+K-k-1$. Choose a grain $w$ and a mesh $h$, both powers of $L$, and
a clearance $r$ such that
\begin{equation}\label{eq:gaussian-localization-scales}
 \begin{gathered}
 R_1\ge L,\quad M_1\ge R_1\mathfrak M,\quad
 M_1\le w,\quad Lw\le h,\quad Lh\le M,\\
 \rho_{\rm t}+\rho_{\rm c}+L<r,\qquad 4(r+1)\le w,
 \qquad \delta_{\rm cv}w\ge128L\log(4n_{\rm w}),\\
 \rho_{\rm t}=\frac{L(L+R_1+1)}{L-1}\mathfrak M,
 \qquad \rho_{\rm c}=22r_{\rm cv}L\mathfrak M.
 \end{gathered}
\end{equation}
The covariance floors of \cref{I-73-90} are imposed on $\mathfrak M$ and
$R_1\mathfrak M$. The constant $n_{\rm w}$ bounds the number, to the power
$N-1$, of connected sets of $N$ grain cubes containing a specified cube.
The mesh has a further lower bound $h\ge h_\sharp(L)$, fixed in the proof.
All these choices precede the lattice and the step.

Here and below, enlargement of a cube means enlargement by the indicated
number of layers of cubes of its own partition. A region is a union of
closed cubes; connectedness of a grain region allows face, edge and corner
contacts. Write $\mathfrak X\langle Z\rangle$ for the nonpivot bonds with
at least one endpoint in $Z$, and $\mathfrak X^\circ\langle Z\rangle$ for
those with both endpoints in $Z$.

\paragraph{Background domain.}\label{def:gaussian-background-domain}
For a region $W$, a pair $p=(\mathbf U,\mathbf J)$ is written
$\mathbf U=e^{i\xi A'}U$, with $U$ real. We use pairs satisfying the
following conditions at radii $(\alpha_0,\alpha_1,\alpha_2;a)$:
\begin{enumerate}
\item Both plaquette curvatures are less than $\alpha_0\xi^2L^{-2}$.
\item Every cubic box $q\subset W$ made of level-$t$ $\mathfrak M$-cubes,
of side at most $12\mathfrak ML$, lies in a box $q'\subset W$ on which
$U$ has a real gauge $U^v=e^{i\xi A}$ with
$L|A|,L^2|\nabla^\xi A|\le\alpha_0$.
\item $L|A'|,L^2|\nabla^U A'|<\alpha_1$ and
$L^3|\mathbf J|<\alpha_2$.
\item Every mesh cube $\square$ with $\square^{\sim2}\subset W$ has a
real gauge on $\square^{\sim2}$ with
$U^v=e^{i\xi A}$ and $L|A|,L^2|\nabla^\xi A|\le a$.
\end{enumerate}
All bonds and stencils in these conditions are contained in $W$.
Let $\mathcal P_h(W;\alpha_0,\alpha_1,\alpha_2;a)$ be this domain, and
put $\mathcal P_{h,*}(W)=\mathcal P_h(W;a_{\rm w},a_{\rm w},a_{\rm w};a_{\rm w})$,
where the positive $L$-dependent radius $a_{\rm w}$ is chosen below.
The conditions pass to restrictions. Their existential real-gauge
condition is retained under small complex perturbations by keeping the
real factor fixed; hence these are open domains when $\alpha_1\le\pi$.

A level-$t$ triple $P=(U,A'_P,J_P)$ is given on the bare bonds incident
to its region. Its level-$k$ reading is
\begin{equation}\label{II-38-12}
 p_{k,W}(P)=\bigl(e^{i\eta_tA'_P}U,L^{-3}J_P\bigr)|_W,
 \qquad \eta_t=L^{-t},\qquad
 (U,A')=(U,L^{-1}A'_P).
\end{equation}
The simple domain $\mathcal Q_t(X;b_0,b_1,c_0)$ of
Definition~\ref{II-16-12} is expressed in level-$t$ units.
Its bounds imply the preceding pair conditions if
\[
 12\mathfrak M\le\mathsf oLM,\qquad 5h\le\mathsf oL^2M,
 \qquad b_0\le a_{\rm G}M^{-1},\quad b_1\le\alpha_1,\quad c_0\le\alpha_2,
\]
where
\[
 a_{\rm G}=\frac{a_{\rm w}}{2\max\{1,\mathsf oLB^{(\rm cg)}\}},
\]
with $\alpha_0=a_{\rm w}/2$. Indeed the two kinds of boxes in the pair
conditions have side at most $\mathsf oLM$ after rescaling, while the
potential and gradient rescale by $L^{-1}$ and $L^{-2}$.

\paragraph{The operators.}
Let $\mathsf A$ denote the fluctuation form of \eqref{eq:cv-form},
including the gauge term. The change of units in \eqref{II-38-12} gives
the same form on the same nonpivot coordinates: the defining quadratic
form is
$ (\mathsf C_VB,(\mathsf G_V^{(2)}+\Delta_k)\mathsf C_VB)_1$.
Thus the restriction of the form and its inverse agree with those already
constructed, including their zero-extension convention. We construct
local coefficients $a_S$ and $r_\omega$ below, with
\[
 \mathsf A=\sum_{S\ \mathrm{connected}}a_S,
 \qquad \mathsf A^{-1/2}=\sum_\omega r_\omega
\]
on the whole torus. Their domains $X(S)$ and $X(\omega)$ are the coarse
cube hulls of their figures. A coefficient is multiplied once by each
parameter in its domain:
\[
 \mathsf A(s;p)=\sum_S\prod_{\Delta\subset X(S)}s(\Delta)\,a_S[p],
 \qquad r(s;p)=\sum_\omega\prod_{\Delta\subset X(\omega)}s(\Delta)\,r_\omega[p].
\]
Terms whose parameter monomial vanishes are omitted. These expressions
require the background only on the union of the surviving domains.

Let $R$ be a coarse region and let $Z_\circ$ be a level-$k$ region with
$Z_\circ^\sharp=(Z_\circ^\sim)^\uparrow\subset R$. Complete parameters
on the coarse cubes of $R\setminus Z_\circ^\sharp$ by setting
$\bar s=1$ on $Z_\circ^\sharp$ and $\bar s=0$ off $R$. Put
$\chi_\circ=\chi_{\mathfrak X\langle Z_\circ\rangle}$ and
$E_\circ=\chi_\circ^{\mathsf T}\chi_\circ$. Define
\begin{equation}\label{II-38-15}
 \mathbf C=(\chi_\circ\mathsf A(\bar s;p)\chi_\circ^{\mathsf T})^{-1},
 \qquad
 \boldsymbol\Gamma=\chi_\circ\mathsf A(\bar s;p)(1-E_\circ)
                         r(\bar s;p)\chi_{\mathfrak X\langle R\rangle}^{\mathsf T}.
\end{equation}
The reference operators $\mathbf C_{\rm r},\boldsymbol\Gamma_{\rm r}$
are these operators at $s=0$ and $p_{\rm r}=(U,0)$.

\begin{theorem}\label{II-38-61}
There are positive constants $a_{\rm w},a_{\rm Str},\theta_{\rm Str},
C_{\rm Str}$, $0<\gamma_{\rm r}\le1$, and $c_{\rm r}\ge1$, depending
only on $L$, with the following property. Assume the geometric conditions
above, $\bar\kappa\ge1$, and
\begin{equation}\label{eq:gaussian-structure-smallness}
 \theta_{\rm Str}M\ge16\bar\kappa,\qquad
 0<\alpha_0\le a_{\rm w}/2,\qquad
 \varepsilon_\circ:=e^{-\theta_{\rm Str}M}+\alpha_1+\alpha_2
       \le a_{\rm Str}.
\end{equation}
For every $p\in\mathcal P_h(R;\alpha_0,\alpha_1,\alpha_2;a_{\rm w})$,
\eqref{II-38-15} defines a local Gaussian structure with
$\widehat\varepsilon=C_{\rm Str}\varepsilon_\circ$:
\begin{enumerate}
\item The kernels of $\mathbf C,\mathbf C^{-1},\boldsymbol\Gamma$ are
jointly holomorphic in the pair and the parameters on a neighborhood of
the closed polydisc $|s(\Delta)|\le e^{\bar\kappa}$.
\item If the parameters vanish off a coarse region $Z$, the structure
is the direct sum of the structures on the wall-connected components
of $Z_\circ^\sharp\cup Z$. Each coupling has input supported on the
inner bonds of its component. Restricting the background to any region
containing all surviving domains leaves these operators unchanged.
\item The reference covariance is real and symmetric, with
$\mathbf C_{\rm r}^{-1}\ge\gamma_{\rm r}I$, and
\[
 \|\mathbf C_{\rm r}^{-1}\|,\quad
 \|\boldsymbol\Gamma_{\rm r}\|,\quad
 \|\mathbf C_{\rm r}\boldsymbol\Gamma_{\rm r}\|\le c_{\rm r}.
\]
Its coupling has input supported in
$\mathfrak X^\circ\langle Z_\circ^\sharp\rangle$. Each of
\[
 \mathbf C^{-1}-\mathbf C_{\rm r}^{-1},\qquad
 \boldsymbol\Gamma-\boldsymbol\Gamma_{\rm r},\qquad
 \boldsymbol\Gamma^{\mathsf T}\mathbf C\boldsymbol\Gamma
       -\boldsymbol\Gamma_{\rm r}^{\mathsf T}
                    \mathbf C_{\rm r}\boldsymbol\Gamma_{\rm r}
\]
has norm $\|\cdot\|^{\rm w}_{4\theta_{\rm Str}}$ at most
$\widehat\varepsilon$.
\item On the whole torus the covariance is
$\mathbf C_* =\mathsf A(p)^{-1}$, and at $s=1$
\begin{equation}\label{eq:actual-gaussian-conditioning}
 \mathbf C+\mathbf C\boldsymbol\Gamma\boldsymbol\Gamma^{\mathsf T}\mathbf C
     =\chi_\circ\mathbf C_*\chi_\circ^{\mathsf T}.
\end{equation}
The real whole covariance $\mathbf C_{*\rm r}=\mathsf A(p_{\rm r})^{-1}$
satisfies $\gamma_{\rm r}I\le\mathbf C_{*\rm r}^{-1}\le c_{\rm r}I$ and
$\|\mathbf C_*^{-1}-\mathbf C_{*\rm r}^{-1}\|\le\widehat\varepsilon$.
\end{enumerate}
Empty cores have the zero-dimensional convention. If, in addition,
$e^{-\theta_{\rm Str}M},\alpha_1+\alpha_2\le(LM)^{-3}$, then
$4C_{\rm Str}+c_{\rm r}+3$ is a defect number for the unscaled structure.
For $|\tau-1|\le\frac12(LM)^{-3}$, the scaled structure
$(\tau^{-1}\mathbf C,\tau^{1/2}\boldsymbol\Gamma)$, with the principal
root, has the same references and defect at most
\begin{equation}\label{eq:gaussian-cutoff-reserve}
 \widehat\varepsilon_* -m_0,\qquad
 \widehat\varepsilon_*=2\widehat\varepsilon+c_{\rm r}(LM)^{-3},
 \qquad m_0=\tfrac12c_{\rm r}(LM)^{-3}>0.
\end{equation}
Its whole covariance is $\tau^{-1}\mathbf C_*$.
\end{theorem}

The proof occupies the rest of this section. The cancellation which creates
the $a_S$ is what makes their parameter dependence local. The mesh
resolvents preserve that locality for the inverse square root. Their
weighted bounds then separate the exponentially small contribution from
outside the core collar from the small complex change of background.

\subsection{Connected coefficients of the form}
\label{sec:gaussian-connected-coefficients}
For a grain region $Y$, let $\Omega_0(Y)=Y$ and form its maximal
separated tower by retaining at stage $q+1$ the level-$(q+1)$
$\mathfrak M$-cubes contained in $\Omega_q(Y)$ and at distance strictly
greater than $R_1\mathfrak M\xi_q$ from its complement.
Put $S_0(b)=\{\sigma:\operatorname{dist}_\infty(b,\sigma)<r\}$ and
$Y^- =\{b\in\mathfrak X:S_0(b)\subset Y\}$. The clearance ensures that
$Y^-$ consists of fluctuation bonds of this tower.
The background is admissible on every such tower, with the same radii.
Indeed a site in layer $n$ of the meet lies in an original layer
$n'\ge n$, so its curvature, potential, derivative and current bounds
only weaken. Each natural gauge cube of the meet is contained in a
plain gauge cube of the original sequence lying in $Y$. Restrict that
gauge to the smaller cube; its containing box is then inside $Y$.
For the constant sequence the enclosing box can be seen directly.
If the natural cube is already at level $t$, use that cube. Otherwise
its side is at most $12\mathfrak M\xi_n<\mathfrak M L$.
In each coordinate enclose its interval in two adjacent intervals of
length $\mathfrak M L$: use the two intervals straddling a crossed
grid point, or, if there is no crossing, choose the neighboring interval
inside the same interval of length $\mathfrak M L^2$.
Every selected interval lies in an interval of length
$\mathfrak M L^2$ whose intersection with the original interval
contains a bare coordinate. Taking products, each
$\mathfrak M L^2$-cube used contains a bare site of the original
cube. Since
$w>\mathfrak M L$ and these sides are powers of $L$, the grain
partition is coarser than this partition; those cubes therefore lie
in $Y$. The resulting box has side $2\mathfrak M L$ and lies in $Y$.
Thus the local regularity conditions, including the containing-box
condition, hold before any operator is evaluated.
Let $g_Y$ be its
fluctuation form compressed to $Y^-$ and extended by zero, and put
$g_\emptyset=0$. Define
\begin{equation}\label{II-29-11}
 a_S=\sum_{Y\subset S}(-1)^{\#S-\#Y}g_Y.
\end{equation}
Then $\sum_{S\subset Y}a_S=g_Y$: the coefficient of $g_{Y'}$ is
$\sum_{Y'\subset S\subset Y}(-1)^{\#S-\#Y'}$, which is zero unless
$Y'=Y$. In particular the whole-torus sum is $\mathsf A$.

The tower rule has range
$e_n=(L+R_1+1)\mathfrak M\sum_{q<n}\xi_q$, so $e_{k+1}\le\rho_{\rm t}$.
Indeed membership at stage $q+1$ checks precisely the level-$q$ cubes
within $R_1\mathfrak M\xi_q$ of the containing level-$(q+1)$ cube.
Their sites are within $(L+R_1+1)\mathfrak M\xi_q$ of the original site.
Induction proves the claim, separation and maximality. Thus changing a
grain cube farther than $e_{k+1}$ from a cell hull changes neither the
tower nor its fluctuation coordinates there.

We use the form expansion obtained from the same inverse and multiplier
terms as \cref{I-73-90}. Its three kinds of terms are
\begin{equation}\label{I-73-106}
 \begin{aligned}
 A^{(1)}_\kappa&=\theta_\kappa\mathsf C^{\mathsf T}
        (\mathsf G_V^{(2)}-1-\mathsf E^{(2)}[j_1^0])\mathsf C,\\
 A^{(2)}_\Xi&=\mathsf C^{\mathsf T}\mathsf J^{\mathsf T}
                   [\mathbb E_\Xi]_{\omega,8}\mathsf J\mathsf C,\\
 A^{(3)}_{\kappa,\Xi'}&=-\theta_\kappa\mathsf C^{\mathsf T}
                       \mathsf E^{(2)}[j_{1,\Xi'}]\mathsf C.
 \end{aligned}
\end{equation}
In the third kind $|\omega(\Xi')|\ge1$: the zero-length multiplier
terms are already included in $j_1^0$.
Here $\mathsf C$ includes the extension of the coordinates under
consideration. The first term is attached to $X^{(3)}(\kappa)$, the
second to the tree of $\Xi$, and the third to $X^{(3)}(\kappa)$ followed
by the reversed trunk of $\Xi'$ and its branches. They sum to the form
by inserting the existing expansions of
$(Q\mathsf G_1Q^{\mathsf T})^{-1}$ and $j_{1*}$ into \eqref{eq:cv-form}.
Absolute convergence permits this regrouping.

The outer kernels link only bonds in a common double block, at most
$36L^3$ neighbors per bond. The first and third kinds therefore have
both endpoints in the first cell; the second has its input in the last
trunk cell. A nonzero third-kind term forces the final multiplier cell
to overlap its new first cell. Marking branch starts and reversing the
trunk reduces its decoration count to the count already used for the
multiplier expansion. Moving endpoints through the outer kernels adds
at most $8L$ to the second-kind tree length and $16L$ to the third;
the first kind has length at most $12L$. The inverse and multiplier
majorants, these finite neighbor sums and the decoration count give
\begin{equation}\label{I-73-109}
 |A_{Y,\omega}(b,b')|
 \le B_0(C_{\rm cv}\mathfrak M^{-1/2})^{|\omega|}
       \mathfrak M^{-|\omega|/2}
       e^{-\delta_{\rm cv}d^{\rm tr}(\omega;y(b),y(b'))/2}.
\end{equation}
For the third kind, if the original trunk length is $m'\ge1$, the new
length is $m'+1$; the factor $\mathfrak M^{-3m'}$ absorbs
$\mathfrak M^{-(m'+1)}$, and the at most $2^{m'}$ branch markings are
absorbed in $C_{\rm cv}$. The second kind has the stronger
factor $\mathfrak M^{-3|\omega|}$. Enlarge the covariance constants
and their fixed floor on $\mathfrak M$ together, so that the tree
count times $C_{\rm cv}\mathfrak M^{-1/2}$ remains at most
$1/2$; decrease the rate once if necessary.

The local products are defined as follows. In the first kind use the
local $\mathsf C_\kappa,\mathsf E_\kappa,\mathsf G_\kappa$, restrict
$j_1^0$ to $\kappa^{[5/6]}$, cut the input to $\kappa^{[2/3]}$ and
multiply the output by $\theta_\kappa$. In the second kind use the
local $\mathsf J\mathsf C$ at the last trunk cell, with input cut
$\kappa^{[2/3]}$ and row cut $\kappa^{[1/12]}$, followed by
$\mathbb E_\Xi$ and the local $\mathsf C^{\mathsf T}\mathsf J^{\mathsf T}$
at the first cell, with output cut $\kappa_0^{[5/6]}$.
In the third kind use the already constructed local solution
$W_{1,\kappa}[j_{1,\Xi'}]$ and multiply its ninth component by
$-\theta_\kappa$. These are exactly the local restrictions of
\eqref{I-73-106}. Their factors are holomorphic on their carriers.
Equal tower traces on the cell hulls and equal fluctuation coordinates
give equal decorations and equal local matrices, hence equal terms.

\begin{proposition}\label{II-30-49}
The coefficients $a_S$ are symmetric and holomorphic in the background
on $S^\sim$. They vanish unless $S$ is connected, and their kernels
vanish unless $S_0(b)\cup S_0(b')\subset S$. There are $L$-dependent
constants $C_a<\infty$ and $\delta_a=\delta_{\rm cv}/(4L)>0$ such that
\begin{equation}\label{eq:connected-form-bound}
 |a_S(b,b')|\le C_a\epsilon^{\#S-1}e^{-\delta_a\Lambda_S(b,b')},
 \qquad \epsilon=(2n_{\rm w})^{-1}.
\end{equation}
Here $\Lambda_S$ is the least length of a connected finite union of
segments through $b_-,b'_-$ meeting $\sigma^\sim$ for every cube
$\sigma\subset S$.
\end{proposition}

\begin{proof}
The support and symmetry follow from \eqref{II-29-11}. Expand each
$g_Y(b,b')$ in trees. Their weighted length is at least
$L^{-1}(\lambda_\omega-2)$, where $\lambda_\omega$ is the plain tree
length through $b_-,b'_-$ and the cells: all weights are at least
$L^{-1}$, and the two unit bonds connect the preferred endpoints.
Only a bounded number of first cells contain $b_-$, and the tree count
times $C_{\rm cv}\mathfrak M^{-1/2}$ is at most $1/2$.
The remaining majorants thus have a uniformly bounded sum.

Let $N(\omega)$ consist of the grain cubes within $e_{k+1}$ of the
walk hull. If $S$ contains a cube outside
$S_0(b)\cup S_0(b')\cup N(\omega)$, toggle the first such cube in a
fixed ordering. The finite-range observation and local product identity
give the same tree and kernel on the toggled tower, with opposite sign
in \eqref{II-29-11}. This is an involution; absolute convergence
justifies the cancellation.

For a survivor every cube of $S$ lies within $r+1$ of a tree through
the endpoints and the cells. For cubes in $S_0$ this is the definition;
for the others use the hull diameter $\rho_{\rm c}$ and
$e_{k+1}+\rho_{\rm c}+L<r$. Hence
\[
 \lambda_\omega\ge
       \max\{\Lambda_S(b,b'),(\#S-8)w/32\}.
\]
For the second bound, traverse a connected segment cover twice and
divide the tour into pieces of length less than $w/2$. Every cube
within $r+1\le w/4$ of the cover is within $w/2$ of one of at most
$4\lambda_\omega/w+1$ selected points. At most eight grain cubes lie
this close to a point, by counting intervals in each coordinate.
Approximation of the infimum gives the bound. The cells of a nonzero
tree form a connected set in $Y$; its intersecting grain cubes form
a connected subset of $Y$, and every surviving cube touches this
subset. Thus $S$ is connected.

There are at most $2^{\#S}$ subregions. Spend half the exponential
decay on each of the two lower bounds for $\lambda_\omega$.
The grain inequality in \eqref{eq:gaussian-localization-scales} absorbs
$2^{\#S}$ and gives \eqref{eq:connected-form-bound}; the cases
$\#S\le8$ are absorbed in $C_a$. Holomorphy follows from locally
uniform convergence of the tree sums and the finite alternating sum.
The rooted exploration-tree count of connected grain regions gives
\begin{equation}\label{eq:grain-entropy-sum}
 \sum_{S\supset S_0(b),\ S\ \mathrm{connected}}\epsilon^{\#S-1}\le2.
\end{equation}
Indeed a depth-first traversal records one of at most $n_{\rm w}^{N-1}$
connected sets of $N$ cubes through a specified cube, and
$n_{\rm w}\epsilon=1/2$.
\end{proof}

\subsection{Local resolvents and their strings}
\label{sec:gaussian-local-resolvents}
Write $\gamma>0$ for the flat gap in
\eqref{eq:cv-flat-gap}. It applies to every maximal tower above.
We first transfer it to local backgrounds, and then construct the inverse
square root without extending those backgrounds beyond their domains.

Put $c_{\rm lat}(v)=3\sum_{z\in\mathbb Z^3}e^{-v|z|_\infty}$.
Least lifts of torus sites give the same bound for sums over positive
bonds. By \eqref{eq:grain-entropy-sum}, with $\theta_a=\delta_a/2$,
every sum of coefficients with multipliers of modulus at most $\kappa$
has weighted norm at most $\kappa B_a$, where
\[
 B_a=\max\{\gamma,\,2\sqrt2 C_a c_{\rm lat}(\theta_a)\}.
\]
If all its terms have $\Lambda_S\ge D$, its unweighted norm is at most
$\kappa B_a e^{-\theta_aD}$. Here $\sqrt2$ converts the matrix norm
on $\mathfrak g^c$ to its Hermitian norm.

The complex comparison uses the disc
\begin{equation}\label{II-33-17}
 p_z=(e^{iz\xi A'}U,z\mathbf J),\qquad
 |z|<R=\frac{c_{\rm d}a}{\alpha_1+\alpha_2},\qquad c_{\rm d}=1/16.
\end{equation}
If $\alpha_0\le a/2$ and $a\le1$, this disc satisfies the first three
pair conditions at radii $(a,a,a)$; its real mesh gauges are unchanged.
Only curvature needs verification. At a plaquette based in layer $n$,
put $\vartheta=\xi/\xi_n$ and $\tau=|z|\alpha_1$.
Transport its four perturbations to the base point using $U$.
The product is $EP$, where $P=\partial U$ and
$E=e^{iX_1}e^{iX_2}e^{-iX_3}e^{-iX_4}$.
The two transverse values are expressed through the covariant
differences at the base point. Consequently
\[
 \sum_j|X_j|\le6\tau\vartheta,\qquad
 |X_1+X_2-X_3-X_4|
 \le\tau\vartheta^2(2+4\alpha_0).
\]
The latter bound uses
$|(1-\operatorname{Ad}P)X|\le2|P-1|\,|X|$.
The exponential series, with its terms of degrees zero and one removed,
has norm at most $\frac12s^2e^s$, $s=\sum|X_j|$. Thus
\begin{equation}\label{II-33-21}
 |E-1|\le\vartheta^2\tau
       (2+4\alpha_0+18\tau e^{6\tau})<\tfrac12a\vartheta^2.
\end{equation}
Together with $|P-1|<\alpha_0\vartheta^2$, this proves the claim.
It avoids changing scales at the opposite corners of the plaquette.
Schwarz's lemma, applied to scalar matrix elements and then taking the
operator norm, now bounds the difference of any holomorphic coefficient
from its value at $p_0$ by $2/R$ times its uniform majorant.

Choose $a_{\rm w}>0$ so that
$a_{\rm w}\le c_{\rm d}a_{\rm cv}\gamma/(64B_a)$ and the inherited
covariance radii hold. A real mesh gauge with bounds $a_{\rm w}$
can be viewed relative to the flat reference, at radii
$(a_{\rm w},3a_{\rm w}/2,a_{\rm w}/2)$.
Apply \eqref{II-33-17} with $a=a_{\rm cv}$.
The resulting form differs from its flat value by at most $\gamma/16$.
The same bound compares a pair in $\mathcal P_{h,*}$ with its real
reference. Unitary gauge covariance preserves both the Hermitian norm
and the form inequality.

For a mesh cube with center $y_\square$, use the periodic product cutoff
\[
 h_\square(b)=
 \prod_{\nu=1}^3
 \cos\!\left(\frac{\pi(x_b-y_\square)_\nu}{2h}\right)
 \quad\text{if }|x_b-y_\square|_\infty<h,
 \qquad h_\square(b)=0\text{ otherwise},
\]
where $x_b$ is the midpoint and least lifts are used on the support.
The one-dimensional identity $\cos^2u+\sin^2u=1$ gives
$\sum_\square h_\square^2=1$.
At most eight cutoffs are nonzero at a bond, and
$|h_\square(b)-h_\square(b')|
\le(3\pi/2h)(|b-b'|+1/2)$.
Consequently the kernel identity
\[
 O-\sum_\square h_\square O h_\square
 =\tfrac12\sum_\square(h_\square(b)-h_\square(b'))^2O(b,b')
\]
has norm at most
$18\pi^2h^{-2}(4/\theta^2+1/2)\|O\|_\theta^{\rm w}$.
This is the only localization loss in the gap estimate.

Let $\mathfrak X^\triangleright\langle W\rangle$ contain both
orientations of those positive bonds whose initial point lies in $W$.
Compress $\sum_{S^\sim\subset\square^{\sim2}}a_S$ to
$\mathfrak X^\triangleright\langle\square^\sim\rangle$, obtaining
$a^\square$. At the flat pair the coefficient sum is the tower form
for the union of all grains with $\sigma^\sim\subset\square^{\sim2}$.
The compressed bonds are deep in that union, so its gap is $\gamma$.
The two disc comparisons give a real gap at least $7\gamma/8$ and
a complex real-part gap at least $3\gamma/4$ for $a^\square$.

More generally, take any reversal-invariant bond set $Y$ and any partial
sum $A_{\mathcal W}=\sum_{S\in\mathcal W}a_S$ that is full near $Y$:
whenever $h_\square$ meets $Y$, its enlarged cube lies in the background
domain and $\mathcal W$ includes all $S$ with
$S^\sim\subset\square^{\sim2}$.
For a term excluded from this local sum and a bond starting in
$\square^\sim$, the length is at least $h-3w\ge10h/13$.
Indeed choose a point of $S^\sim$ outside $\square^{\sim2}$ and a
point of its grain enlargement on a connecting tree; their distance is
at most $3w$, while the first point is more than $h$ from
$\square^\sim$. If $h_\square(b')>0$, the same argument from the
midpoint instead gives
$\Lambda_S(b,b')\ge3h/2-3w-1/2\ge h$.
Choose $h_\sharp$ so that
\[
 B_a e^{-10\theta_a h/13}\le\gamma/32,\qquad
 18\pi^2h^{-2}(4/\theta_a^2+1/2)B_a\le\gamma/16.
\]
Localization by $\sum h_\square^2=1$ then gives real gap
$3\gamma/4$ for the compression to $Y$; the complex comparison leaves
real-part gap $\gamma/2$.

For completeness, the resolvent estimate used next follows by the same
exponential conjugation for every compression. If
$\|K\|_{\theta_0}^{\rm w}\le B$ and $\operatorname{Re}K\ge gI$, then
\begin{equation}\label{II-37-38}
 |(\mu+K)^{-1}(b,b')|_2
 \le(\mu+g/2)^{-1}
      \exp\!\left(-\frac{\theta_0g}{2B}|b-b'|\right),
 \qquad\mu\ge0.
\end{equation}
Conjugate by $e^{\theta|b-b'|}$, with $\theta=\theta_0g/(2B)$.
Convexity gives $|e^s-1|\le(\theta/\theta_0)e^{\theta_0d}$
when $|s|\le\theta d$, so the conjugated operator differs from $K$
by at most $g/2$ in norm. Its inverse has norm at most
$(\mu+g/2)^{-1}$, proving \eqref{II-37-38}.

We now keep the above $Y$ and full partial sum, writing $A_Y$ for its
compression. Put $Y^\square=Y\cap
\mathfrak X^\triangleright\langle\square^\sim\rangle$ and let
$E_\square$ be its projection. All operators in the following formulas
are extended by zero to $Y$:
\begin{align*}
 g_\square^\mu&=(\mu+E_\square a^\square E_\square)^{-1}
                     \quad\text{on }Y^\square,\\
 k_{\square,S}
 &=1_Y\bigl(a_Sh_\square-
       1_{\{S^\sim\subset\square^{\sim2}\}}h_\square a_S\bigr)E_\square,\\
 \Pi^\mu&=\sum_\square h_\square g_\square^\mu h_\square,\qquad
 e^\mu=-\sum_{\square,S}k_{\square,S}g_\square^\mu h_\square .
\end{align*}
Multiplying out, using fullness and $\sum h_\square^2=1$, gives
$(\mu+A_Y)\Pi^\mu=1_Y-e^\mu$.
The products in $\Pi^\mu(e^\mu)^n$ are indexed by strings
$\omega=(\square_0;(S_1,\square_1),\ldots,(S_n,\square_n))$:
$S_i^\sim$ must meet both neighboring mesh enlargements, or the
product vanishes. Their figures are
$F_\omega=\bigcup_i\square_i^{\sim2}\cup\bigcup_iS_i^\sim$.
Define
\begin{equation}\label{II-37-69}
 g_\omega^\mu=h_{\square_0}g_{\square_0}^\mu h_{\square_0}
       \prod_{i=1}^n(-k_{\square_i,S_i}g_{\square_i}^\mu h_{\square_i}),
 \qquad
 r_\omega=\frac1\pi\int_0^\infty\mu^{-1/2}g_\omega^\mu\,d\mu .
\end{equation}
Each product uses the background only on $F_\omega$.
Its input and output bonds have both endpoints in that figure.

Here is the convergence estimate, including the loss from the
commutators. Put
$m_S=\sqrt2 C_a\epsilon^{\#S-1}e^{-\delta_a\Lambda_S}$
on its prescribed support, and zero elsewhere. A majorant for
$k_{\square,S}$ is
\[
 \mathfrak k_{\square,S}(c,c')=
 \begin{cases}
 \dfrac{3\pi}{2h}(|c-c'|+1/2)
   1_{\{h_\square(c)+h_\square(c')>0\}}m_S(c,c'),
       &S^\sim\subset\square^{\sim2},\\
 1_{\{h_\square(c')>0\}}m_S(c,c'),&\text{otherwise}.
 \end{cases}
\]
Set $\theta_\square=3\theta_a\gamma/(8B_a)$ and
$\widehat g(c,c')=(4/\gamma)e^{-\theta_\square|c-c'|}$.
By \eqref{II-37-38},
$|g_\square^\mu|\le\gamma(4\mu+\gamma)^{-1}\widehat g$.
Choose $0<\theta_g\le\min(\delta_a/4,\theta_\square/6)$.
The preceding support and length bounds, the sixteen possible cutoffs
at a pair of bonds, and \eqref{eq:grain-entropy-sum} give
\[
 \widehat K(c,c'):=\sum_{\square,S}
       \mathfrak k_{\square,S}(c,c')e^{\theta_g\Lambda_S(c,c')}
 \le2\sqrt2 C_a e^{-\delta_a|c-c'|/2}
       \left[\frac{24\pi}{h}(|c-c'|+1/2)
                         +8e^{-\delta_a h/4}\right].
\]
Its row and column sums are bounded by a number
$\epsilon_k(h)\to0$. The kernel
$\widehat g(c,c')e^{3\theta_g|c-c'|}$ has both sums at most
$C_g/\gamma$, where $C_g=4c_{\rm lat}(\theta_\square/2)$.
Increase $h_\sharp$ so that $\epsilon_k(h)C_g/\gamma\le1/2$.

Let $\Lambda_\omega(b,b')$ be the least connecting length through
the endpoints and every piece $\square_i^{\sim2},\sigma^\sim$ of
$F_\omega$. Join the intermediate bonds in the ordered product by
coordinate paths, and join the pieces by the trees for the $S_i$.
This gives
$\Lambda_\omega\le3|b-c_0|+
\sum_i(\Lambda_{S_i}(c_{i-1},c_i')+3|c_i'-c_i|)$.
Replacing the factors of $g_\omega^\mu$ by
$h_{\square_0}\widehat g h_{\square_0}$ and
$\mathfrak k_{\square_i,S_i}\widehat g h_{\square_i}$ defines a
nonnegative majorant $m_\omega$, for which
\begin{equation}\label{II-37-96}
 \sup_b\sum_{b'}\sum_{|\omega|=n}
 m_\omega(b,b')e^{\theta_g\Lambda_\omega(b,b')}
 \le\frac{C_g}{\gamma}2^{-n}.
\end{equation}
The identical column estimate follows by summing products in the
opposite order. Moreover
$|g_\omega^\mu|\le(\gamma/(4\mu+\gamma))^{n+1}m_\omega$ and
$|r_\omega|\le\sqrt\gamma\,m_\omega/2$.
These bounds prove convergence and holomorphy of the integrals and all
string sums, locally uniformly on the background domain.

The error $e^\mu$ has row and column sums at most
$\frac12\gamma/(4\mu+\gamma)$, by the same $\widehat K\widehat g$
bound. Hence
$(\mu+A_Y)\sum_{n=0}^N\Pi^\mu(e^\mu)^n=1_Y-(e^\mu)^{N+1}$
and
$\sum_\omega g_\omega^\mu=(\mu+A_Y)^{-1}$.
Integration gives $\sum_\omega r_\omega=A_Y^{-1/2}$.
Indeed accretivity implies $\|e^{-tA_Y}\|\le e^{-\gamma t/2}$;
the Laplace resolvent formula identifies the integral with
$\pi^{-1/2}\int_0^\infty t^{-1/2}e^{-tA_Y}\,dt$.
Squaring and using $t+s=u$ gives the beta integral
$\int_0^1[v(1-v)]^{-1/2}dv=\pi$, hence its square is
$\int_0^\infty e^{-uA_Y}\,du=A_Y^{-1}$.
It is symmetric because every resolvent is. This also proves the
claimed identity on the whole torus.

\subsection{Parameters, positivity and conditioning}
\label{sec:actual-gaussian-structure}
We finish the proof of \cref{II-38-61}. The same majorants control the
form coefficients and the strings. Set
$n_S=m_S$, $n_\omega=\sqrt\gamma\,m_\omega/2$, and
$\theta_\flat=\min(\theta_g,\delta_a/2)$.
Then \eqref{eq:grain-entropy-sum} and \eqref{II-37-96} give
\begin{equation}\label{II-38-40}
 \begin{aligned}
 \sup_b\sum_{b',S}n_S(b,b')e^{\theta_\flat\Lambda_S(b,b')}
       &\le b_a:=2\sqrt2 C_a c_{\rm lat}(\delta_a/2),\\
 \sup_b\sum_{b',\omega}n_\omega(b,b')
                 e^{\theta_\flat\Lambda_\omega(b,b')}
       &\le b_r:=C_g/\sqrt\gamma .
 \end{aligned}
\end{equation}
Both estimates hold with rows and columns exchanged.
On the disc \eqref{II-33-17}, the coefficient differences therefore obey
\begin{equation}\label{eq:gaussian-term-difference}
 |a_S[p]-a_S[p_{\rm r}]|_2
       \le\tfrac13\varrho(\alpha_1+\alpha_2)n_S,\qquad
 |r_\omega[p]-r_\omega[p_{\rm r}]|_2
       \le\tfrac23\varrho(\alpha_1+\alpha_2)n_\omega,
 \quad \varrho=\frac3{c_{\rm d}a_{\rm w}}.
\end{equation}
For the first estimate use radius $c_{\rm d}a_{\rm cv}/
(\alpha_1+\alpha_2)$, with $a_{\rm cv}\ge2a_{\rm w}$; for the second
use $c_{\rm d}a_{\rm w}/(\alpha_1+\alpha_2)$ and integrate the
resolvent difference. Taking
$\alpha_1+\alpha_2\le c_{\rm d}a_{\rm w}/3$ keeps these discs larger
than one.

We need two geometric facts about the coarse domains $X(\iota)$.
Each piece of a figure has diameter at most $5h$; a tree admissible
for $\Lambda_\iota$ meets every piece. Thus the figure is contained
in the $5h$-neighborhood of that tree. The segment-cover count used
above, now at side $LM$, gives
\begin{equation}\label{eq:gaussian-parameter-count}
 \#X(\iota)\le8+\frac{32}{(L-1)M}\Lambda_\iota.
\end{equation}
Indeed a doubled tour divided at spacing $(LM-10h)/2$ gives at most
$8\max(1,4\,\mathrm{length}/(LM-10h))$ intersected cubes;
$10h\le M$ and passage to the infimum give the display.
If a figure leaves $Z_\circ^\sharp$ and the first bond meets
$Z_\circ$, its length is at least $M/2$.
To see this, a point in a coarse cube outside $Z_\circ^\sharp$
is at distance at least $M$ from $Z_\circ$: otherwise the fine cube
containing it would belong to $Z_\circ^\sim$.
The initial point of the first bond is within one of $Z_\circ$,
and a tree point lies within $5h$ of that figure point.
Hence its length is at least $M-1-5h\ge M/2$.
For a product $a_S(b,c)r_\omega(c,b')$, the same conclusion holds
for $\Lambda_S(b,c)+\Lambda_\omega(c,b')$, since their trees meet
at $c_-$.

Choose
\[
 \theta_C=\frac{\theta_\flat\gamma}{16b_a},\qquad
 b_C=\frac4\gamma c_{\rm lat}(\theta_C/2),\qquad
 0<\theta_{\rm Str}\le
             \min(\theta_\flat/18,\theta_C/8).
\]
For $|s|\le e^{2\bar\kappa}$,
\eqref{eq:gaussian-parameter-count} and
$\theta_{\rm Str}M\ge16\bar\kappa$ imply
\begin{equation}\label{eq:gaussian-parameter-cost}
 \left|\prod_{\Delta\subset X(\iota)}\bar s(\Delta)\right|
 \le\exp\{\theta_{\rm Str}M+
                       \theta_{\rm Str}\Lambda_\iota/3\}.
\end{equation}
Here $L\ge13$ absorbs $64\bar\kappa/((L-1)M)$.
An individual form term outside the core hull has length at least
$M/2$, so the factor $e^{\theta_{\rm Str}M}$ costs at most
$e^{2\theta_{\rm Str}\Lambda_S}$.
Keeping another such length factor gives
$e^{-\theta_{\rm Str}M}$; the remaining exponent is within
$\theta_\flat$ even after the weight
$e^{\theta_\flat|b-b'|/2}$ is inserted.
Terms inside the hull have parameter one and are compared by
\eqref{eq:gaussian-term-difference}. Consequently
\begin{equation}\label{II-38-53}
 \|\chi_\circ(\mathsf A(\bar s;p)-\mathsf A_{\rm r})
                       \|_{\theta_\flat/2}^{\rm w}
 \le\varrho b_a\varepsilon_\circ .
\end{equation}

The reference sum is full near $\mathfrak X\langle Z_\circ\rangle$.
In fact $h_\square(b)\ne0$ and an endpoint $x\in Z_\circ$ imply
$|y-x|_\infty<7h/2+1/2<M$ for every
$y\in\square^{\sim2}$, so every coarse cube meeting that enlarged
mesh cube lies in $Z_\circ^\sharp$.
The real partial-sum gap proved above therefore gives
$\mathbf C_{\rm r}^{-1}\ge3\gamma I/4$.
Choose
$a_{\rm Str}\le\min(c_{\rm d}a_{\rm w}/3,\gamma/(4\varrho b_a))$.
Then \eqref{II-38-53} gives
$\operatorname{Re}\mathbf C^{-1}\ge\gamma I/2$ and
$\|\mathbf C^{-1}\|_{\theta_\flat/2}^{\rm w}\le2b_a$.
The exponential resolvent bound yields
$|\mathbf C(b,b')|_2\le(4/\gamma)e^{-\theta_C|b-b'|}$.
The same estimate holds for $\mathbf C_{\rm r}$, and hence
their weighted norms at $4\theta_{\rm Str}$ are at most $b_C$.
The resolvent identity now gives
\[
 \|\mathbf C-\mathbf C_{\rm r}\|_{4\theta_{\rm Str}}^{\rm w}
 \le\varrho b_ab_C^2\varepsilon_\circ .
\]

To estimate the coupling, expand its two factors as $a_S r_\omega$,
with the projection $1-E_\circ$ between them.
If both domains lie in $Z_\circ^\sharp$, subtract the reference
product in two terms and use \eqref{eq:gaussian-term-difference}.
Otherwise their combined length is at least $M/2$.
The two parameter factors cost
$e^{2\theta_{\rm Str}M+\theta_{\rm Str}(\Lambda_S+\Lambda_\omega)/3}$;
after retaining $e^{-\theta_{\rm Str}M}$ and the desired bond weight,
the total exponent charged to length is at most
$31\theta_{\rm Str}/3\le\theta_\flat$.
Using \eqref{II-38-40} in either order proves
\[
 \|\boldsymbol\Gamma-\boldsymbol\Gamma_{\rm r}
                          \|_{4\theta_{\rm Str}}^{\rm w}
 \le\varrho b_ab_r\varepsilon_\circ,\qquad
 \|\boldsymbol\Gamma_{\rm r}\|_{\theta_\flat}^{\rm w}\le b_ab_r,\qquad
 \|\boldsymbol\Gamma\|_{4\theta_{\rm Str}}^{\rm w}\le2b_ab_r.
\]
The last inequality follows also by reducing $a_{\rm Str}$ so that
$\varrho\varepsilon_\circ\le1$.
Split
$\Gamma^{\mathsf T}C\Gamma-\Gamma_{\rm r}^{\mathsf T}C_{\rm r}\Gamma_{\rm r}$
into the three differences, one factor at a time. The preceding
bounds give
$\varrho b_a^2b_r^2b_C(3+2b_ab_C)\varepsilon_\circ$.
Taking $C_{\rm Str}$ to dominate these finitely many coefficients
proves all the stated differences. Take
$\gamma_{\rm r}=\min(1,3\gamma/4)$ and
$c_{\rm r}\ge\max(1,b_a,b_ab_r,b_ab_rb_C)$.
These choices depend only on $L$.

It remains to check the structural assertions.
For fixed background, group all terms according to the subset of
coarse parameters in their domains. There are only finitely many such
subsets, and their coefficient sums converge absolutely by
\eqref{II-38-40}. Thus $\mathsf A(s)$ and $r(s)$ are polynomials,
of degree at most one in each parameter. Their coefficients are
holomorphic by local uniform convergence; inversion is holomorphic
where the determinant is nonzero. The estimates were proved on the
larger polydisc $|s|\le e^{2\bar\kappa}$, which supplies the required
neighborhood.

Each figure is connected. Its coarse hull is wall-connected: for
grain enlargements this follows by refining the touching chain, and
successive string pieces overlap the neighboring mesh enlargements.
If a parameter vanishes off $Z$, each surviving domain therefore
lies in one wall-component of $Z_\circ^\sharp\cup Z$.
The inverse of the resulting block-diagonal form has the same blocks.
Every string input bond is inner to its figure, so the coupling
has precisely the asserted inner-bond support.
The definitions of the coefficients and strings also show restriction
consistency. Real backgrounds make all local matrices, integrals and
coefficients real. Symmetry of $\mathbf C$ follows from symmetry
of each $a_S$. These verify the locality, factorization and reference
requirements of \cref{II-19-17,II-19-18,II-19-19}.

On the whole torus put $G=\mathsf A^{-1}$ and split the coordinates
into the core and its complement. Then
$\mathbf C=A_{00}^{-1}$ and
$\boldsymbol\Gamma=A_{01}\chi_1\mathsf A^{-1/2}$.
Symmetry and $(\mathsf A^{-1/2})^2=G$ give
$\boldsymbol\Gamma\boldsymbol\Gamma^{\mathsf T}=A_{01}G_{11}A_{10}$.
From $GA=AG=1$,
$G_{10}=-G_{11}A_{10}\mathbf C$ and
$G_{00}=\mathbf C+\mathbf C A_{01}G_{11}A_{10}\mathbf C$,
which is \eqref{eq:actual-gaussian-conditioning}.
The whole-torus gap and difference estimates are the preceding ones
with the entire torus as the core. This proves the remaining Gaussian
identity in \cref{II-19-20}.

Finally replace $(\mathbf C,\boldsymbol\Gamma)$ by
$(\tau^{-1}\mathbf C,\tau^{1/2}\boldsymbol\Gamma)$.
The quadratic product $\Gamma^{\mathsf T}C\Gamma$ is unchanged;
the inverse covariance and the coupling differences are at most
$2\widehat\varepsilon+c_{\rm r}|\tau-1|$,
using $|\tau^{1/2}-1|\le|\tau-1|$ on the prescribed disc.
The whole identity scales by $\tau^{-1}$.
Since $|\tau-1|\le\frac12(LM)^{-3}$, this gives
\eqref{eq:gaussian-cutoff-reserve} with a strictly positive reserve.
When the two smallness terms are at most $(LM)^{-3}$, the unscaled
defect number is at most $4C_{\rm Str}+c_{\rm r}+3$.
For the scaled structure the bound instead reads
$(LM)^3\widehat\varepsilon_*\le4C_{\rm Str}+c_{\rm r}$, with the
strict reserve in \eqref{eq:gaussian-cutoff-reserve}.
This completes the proof.

\subsection{Local charts between lattices}\label{sec:local-lattice-charts}
At height $k$, write the infinite point lattices as
$L^{j-k}\mathbb Z^3$, $0\le j\le k+1$.
A finite lattice $\mathfrak w=(\ell,K)$ has side
$s_{\mathfrak w}=2L^{m+K-k}$ in these units; its covering
map is $\operatorname{pr}^{\mathfrak w}$.
Fix an integer $n_{\rm e}\ge1$ and require
$s_{\mathfrak w}\ge C_LM$, where $C_L\ge8n_{\rm e}L$.
The normalizer uses $n_{\rm e}=1$; a later construction may
increase this fixed halo width before choosing the lattices.

A set $\widehat X\subset\mathbb R^3$ is \emph{near} for
$\mathfrak w$ if it lies in a cuboid with all side lengths
less than $s_{\mathfrak w}/2$. Its image is called
\emph{non-wrapping}. In the infinite lattice, near means
bounded. Set
\begin{equation}\label{II-15-10}
 \widehat X^+=\{y:\operatorname{dist}_\infty(y,\widehat X)
                                  \le n_{\rm e}LM\}.
\end{equation}
The halo $X^+=\operatorname{pr}^{\mathfrak w}(\widehat X^+)$
does not depend on the near lift. Indeed, two near lifts of
a nonempty set differ by one period vector: after matching
one point, the coordinate differences between any other
matched pair have magnitude less than one period and are
periods, hence vanish. The halo lies in a cuboid with side
lengths less than $3s_{\mathfrak w}/4$, so its covering
map is injective. The empty set has empty halo.

If $\widehat X$ is near for two lattices and $v\in LM\mathbb Z^3$,
the corresponding chart is
\begin{equation}\label{II-15-12}
 \mathsf i(\operatorname{pr}^{\mathfrak w'}\hat y)
       =\operatorname{pr}^{\mathfrak w''}(\hat y+v),
       \qquad \hat y\in\widehat X^+.
\end{equation}
It acts on fields by carrying the sites and oriented bonds,
and on operators by relabelling their rows and columns.
In a halo, any displacement of norm at most $LM$ is its
actual lifted displacement: an alternative differs by a
period and would have a coordinate of magnitude at least
$s_{\mathfrak w}-LM\ge7s_{\mathfrak w}/8$.
Thus charts preserve all contained lattice stencils,
blocks, contours and pivots. Cubes contained in a halo also
lift to cubes: the difference between the inverse chart
and a lift on the cube is continuous and period-valued,
hence constant.

We compare a local construction only when its entire input
domain and the stencils it uses are in these halos.
Chart consistency means equality after \eqref{II-15-12}
whenever both inputs are admitted. It gives a unique
infinite-lattice value on every finite domain that embeds
in some member of the chosen family of tori: choose an
embedding, evaluate there, and use the zero-translation
chart to compare two choices. This defines values on every
finite domain when the family of sides is unbounded.
Bounds and slice holomorphy pass to this value, since its
evaluation on any fixed finite domain uses one such torus.

\subsection{The Gaussian normalizer}\label{sec:actual-normalizer}
We localize the normalization before expanding the interaction. The
logarithm of the determinant has a convergent series on every deep
subregion. Taking its finite difference over coarse cubes leaves only
closed chains whose support fills the prescribed region. Short chains
then have a large spatial length; long chains are controlled by the
spectral gap.

Use the constant sequence, units and localization parameters of
\cref{sec:local-gaussian-fields}, with $t=k+1$. The coarse cubes have
side $LM$. Choose a further clearance $r_N$ satisfying
\[
 \rho_{\rm t}<r_N,\qquad 4h\le r_N,\qquad
 4r_N\le LM,\qquad 2r_N+1<LM.
\]
Keep $12\mathfrak M\le\mathsf oLM$ and
$5h\le\mathsf oL^2M$. There will be an additional lower bound
$M\ge M_N(L)$. Let $\Gamma$ be all positive level-$t$ bonds,
$\mathfrak X_0$ all level-$k$ bonds, and $\mathfrak X$ their
nonpivot subset. Thus positive bonds remain the independent coordinates.
Write $P_0=(\mathbf1,0,0)$ and set
\begin{equation}\label{II-44-10}
 \alpha_N=\frac{a_N}{\max\{1,\mathsf oLMB^{(\rm cg)}\}},
 \qquad
 \mathcal Q^N(X)=\mathcal Q_t(X;\alpha_N,a_N,a_N),
\end{equation}
where $a_N>0$ is chosen below. The level-$k$ reading of $P$ is always
\eqref{II-38-12}; in particular the current is $L^{-3}J_P$.

For a level-$k$ bond $b$, let $S_N(b)$ be the union of the coarse cubes
at sup distance less than $r_N$ from its segment. For a coarse region
$Y$, put
\[
 Y_N^-=\{b\in\mathfrak X:S_N(b)\subset Y\},\qquad
 A_{Y,N}=\chi_{Y_N^-}\mathsf A_Y\chi_{Y_N^-}^{\mathsf T}.
\]
Here $\mathsf A_Y$ is the actual form for the maximal tower over $Y$,
not the restriction of the whole-torus form. If $c\in\Gamma$, let
$X(c)$ be the union of the coarse cubes containing $c_-$ and $c_+$.
For brevity write
$m_c(V)=\kappa_V(c)$ for the already defined pivot coefficient.

\begin{proposition}\label{II-44-37}
There are $L$-dependent constants
$a_N,\delta_N>0$, $M_N,E_b\ge1$, and $E_N=3L^3E_b$, such that
the following holds when $M\ge M_N$.
For every wall-connected coarse region $X$ there are local terms
$N(X,b;P)$, $b\in\mathfrak X_0^+$, on $\mathcal Q^N(X)$.
They are holomorphic on every fixed-real-background slice, depend
only on the pair represented by $P$, are gauge invariant and
conjugation real, and vanish at $P_0$. Their pointed versions are
\begin{equation}\label{II-44-21}
 N^{(t)}(X,P,z)=\sum_{b\in\mathfrak X_0^+}h_z(b)N(X,b;P).
\end{equation}
A nonzero summand has $z\in X$. Uniformly in the lattice, $k,X,P,b,z$,
\begin{equation}\label{eq:normalizer-decay}
 |N(X,b;P)|\le E_b e^{-\delta_NM d_t(X)},\qquad
 |N^{(t)}(X,P,z)|\le E_N e^{-\delta_NM d_t(X)}.
\end{equation}

On the whole torus put $V=\mathcal M^k\mathbf U$ and
$\Phi_t(P)=\sum_X\sum_{z\in X}N^{(t)}(X,P|_X,z)$. Then
\begin{equation}\label{II-44-39}
 e^{2\Phi_t(P)}
 =\frac{\det\mathsf A(P_0)}{\det\mathsf A(P)}
       \prod_{c\in\Gamma}
       \left(\frac{\det m_c(\mathbf1)}{\det m_c(V)}\right)^2.
\end{equation}
For real $P$, the form is positive and all pivot determinants are
positive. With Lebesgue measure for the real pairing on nonpivot
coordinates, define
\[
 Z_k[P]=\prod_{c\in\Gamma}|\det m_c(V)|^{-1}
       \int\exp\bigl(-\tfrac12(B,\mathsf A(P)B)_1\bigr)\,dB .
\]
Then
\begin{equation}\label{II-44-40}
 \Phi_t(P)=\log Z_k[P]-\log Z_k[P_0].
\end{equation}
The terms are unchanged by a local lattice chart carrying the background
and all the cubes of $X$. Consequently they define the corresponding
family on the infinite lattice, with the same domains and bounds.
\end{proposition}

\begin{proof}
\emph{The local series.}
The simple-domain bounds and the choice of $\alpha_N$ put the
level-$k$ pair in $\mathcal P_{h,*}(Y)$ for every nonempty coarse
$Y\subset X$, if $a_N\le a_{\rm w}$.
Indeed the two real-gauge boxes in that definition fit in the boxes
of the simple domain: their sides in level-$t$ units are at most
$12\mathfrak M$ and $5h/L$. Restriction preserves the curvature,
perturbation and current bounds. The all-meet construction in
\cref{sec:gaussian-connected-coefficients} therefore defines
$\mathsf A_Y$ from the data on $Y$ itself.

Since $r_N>r$ and $4h\le r_N$, the bonds of $Y_N^-$ are deep
enough for the localization proof in
\cref{sec:gaussian-local-resolvents}. In particular, after decreasing
$a_N$, its real-reference compression has spectrum in
$[\gamma/2,B_a]$, and its complex perturbation has norm at most
$\gamma/8$. Take
\[
 c_N=\tfrac12(B_a+\gamma/2),\qquad
 q_N=\frac{B_a-\gamma/4}{B_a+\gamma/2}<1,\qquad
 T_Y=1-c_N^{-1}A_{Y,N}.
\]
The midpoint estimate gives $\|T_Y\|\le q_N$. Thus, for
$b\in(Y_N^-)^+$, the series
\begin{equation}\label{II-44-14}
 f_b(Y;P)=3\log c_N-\sum_{n\ge1}\frac1n
                       \operatorname{tr}T_Y^n(b,b)
\end{equation}
converges normally and has modulus at most
$F_N=3|\log c_N|+3\log(1/(1-q_N))$.
The trace is the trace on the three-dimensional complex Lie algebra.
Define
\begin{equation}\label{II-44-16}
 E^N(X;P,b)=
   \sum_{S_N(b)\subset Y\subset X}(-1)^{\#X-\#Y}f_b(Y;P|_Y).
\end{equation}
All sums over regions here mean sums over their coarse cubes.

The pivot contribution is
\[
 z_c(P)=-\operatorname{Log}\det(Lm_c(\mathcal M^k\mathbf U)).
\]
Its domain and branch require a check. The cubes $X(c)$ contain a
cube of side $2L$ containing the endpoint double block and all its
bare descendants. To see this, in each transverse coordinate choose
one adjacent interval of length $L$ inside the same coarse interval;
in the direction of $c$ use the two consecutive intervals centered
at its endpoints. The coarse intervals contain $M\ge3$ such intervals.
The simple-domain gauge on this cube gives
$U^u=e^{i\eta_t A}$ with $|A|<a_N$.
The logarithm of
$e^{i\eta_t\operatorname{Ad}(u)A'_P}e^{i\eta_t A}$ is therefore
$i\xi\phi$ with $\|\phi\|_\infty<4a_N/L$.
By \cref{I-27-78}, its level-$k$ average is $e^{i\widetilde a}$,
holomorphic in $\phi$, with
$\|\widetilde a\|_\infty\le C_Q\|\phi\|_\infty$.

The pivot coefficient of the one-step derivative is holomorphic in
$\widetilde a$ and equals $L^{-3}$ at zero. The one-step derivative
bound therefore gives
$|1-L^3m_c(e^{i\widetilde a})|\le C_{\rm piv}\|\widetilde a\|_\infty$
on a fixed smaller chart, with $C_{\rm piv}$ depending only on $L$.
Here one may use \eqref{eq:straightening-one-step-bounds}: a variation
supported on the pivot changes neither endpoint contour variable,
so the straightened derivative at that coordinate is $L m_c$.
Choose $a_N$ so that $4C_Qa_N/L$ lies in this chart and
$4C_{\rm piv}C_Qa_N/L\le1/2$, as well as
$4a_N/L<c_{\rm plain}$.
Put $V=\mathcal M^k\mathbf U$ and
$T=L^3m_c(e^{i\widetilde a})$. Gauge covariance gives
$T=L^3\operatorname{Ad}(g)m_c(V)\operatorname{Ad}(g')^{-1}$
for two real group elements $g,g'$. Their adjoint matrices have
determinant one, so $\det T=\det(L^3m_c(V))$.
The bound $|1-T|\le1/2$ places the eigenvalues of $T$ in the disk
$|z-1|\le1/2$. Each has argument at most $\pi/6$, so their product
has argument at most $\pi/2$ and avoids the negative real axis.
Consequently
\begin{equation}\label{eq:normalizer-pivot-log}
 z_c(P)=6\log L-\operatorname{Log}\det(L^3m_c(V)),\qquad
 |z_c(P)-6\log L|\le3\log2.
\end{equation}
This also proves holomorphy on each slice. All averages used here
are the complex averages; at a reversed bond the value is the
inverse, and the domain comparison follows from the same fixed
real gauge.

Define $N^{\rm un}(X,b;P)$ to be $-\tfrac12E^N(X;P,b)$ when
$b$ is nonpivot and $S_N(b)\subset X$, to be $z_c(P)$ when
$b=b_0(c)$ and $X=X(c)$, and to be zero otherwise. Put
\begin{equation}\label{II-44-18}
 N(X,b;P)=N^{\rm un}(X,b;P)-N^{\rm un}(X,b;P_0).
\end{equation}
These definitions give pair dependence, holomorphy and trivial
vanishing. Gauge covariance conjugates every diagonal matrix in
\eqref{II-44-14}; traces and determinants are unchanged.
Complex conjugation preserves the convergent series and the
specified logarithm branch, giving conjugation reality.

\emph{Decay of the connected terms.}
The neighborhood of a bond has side lengths at most $2r_N+1<LM$.
It meets at most two coarse intervals in each coordinate; their
cubes have a common corner. Thus $S_N(b)$ is wall connected,
$\#S_N(b)\le8$, and $d_t(S_N(b))=0$. Likewise $X(c)$ consists
of one cube or two adjacent cubes and has size zero.

Expand each factor of $T_Y$ into the identity or
$-c_N^{-1}a_S$. A term of $T_Y^n(b,b)$ is a closed chain
$b_0=b,b_1,\ldots,b_n=b$, together with the connected grain
regions attached to its nonidentity steps. Its coarse support is
the union of the $S_N(b_i)$ and of the coarse cubes containing
those grain regions. In \eqref{II-44-16}, the sum over all $Y$
containing this support is zero unless the support is exactly $X$.
If $e_n(X,b)$ is the sum of these exact-support chains, then
\begin{equation}\label{eq:normalizer-exact-chains}
 E^N(X;P,b)=3\log c_N\,\mathbf1_{\{X=S_N(b)\}}
                         -\sum_{n\ge1}\frac{e_n(X,b)}n .
\end{equation}

Put $\theta_N=\delta_a/2=\delta_{\rm cv}/(8L)$.
The connected-form estimate and the ordered row sums used in
\cref{II-38-40} give a constant $Q_N>1$, depending only on $L$,
such that the absolute weights of all length-$n$ chains beginning
at $b$, multiplied by $e^{\theta_N\lambda}$, sum to at most
$3Q_N^n$. Here $\lambda$ is the sum of the lengths of their
connecting trees. In particular,
\begin{equation}\label{eq:normalizer-chain-majorants}
 |e_n(X,b)|\le3Q_N^n e^{-\theta_NLMd_t(X)/128},
 \qquad
 |e_n(X,b)|\le3\,2^{\#X}q_N^n.
\end{equation}
The second estimate follows directly from the finite difference
of $\operatorname{tr}T_Y^n(b,b)$.

For the length bound in the first estimate, join the step trees;
their union is connected and has length at most $\lambda$.
Every cube of $X$ lies within $r_N+1$ of this union.
Choose a maximal $LM/8$-separated set on it. If it has $N\ge2$
points, its length is at least $NLM/16$: on each covering segment,
the inverse images of the distance ranges $[0,LM/16)$ about the
chosen points are disjoint, and connectedness makes each range
have that full length. Every cube of $X$ meets a ball of radius
$r_N+1+LM/8<LM/2$ about one chosen point. Such a ball meets at
most eight coarse cubes, with a common corner. If $N=1$ then
$d_t(X)=0$; otherwise
$\#X\le8N\le128\lambda/(LM)$, hence
$LMd_t(X)\le128\lambda$.

For $d=d_t(X)\ge1$, use $\#X\le16d$ and split the sum at
$n_*=\lfloor\vartheta_NLMd\rfloor$, where
$\vartheta_N=\theta_N/(256\log Q_N)$. The first majorant in
\eqref{eq:normalizer-chain-majorants} gives
\[
 \sum_{1\le n\le n_*}\frac{|e_n|}{n}
 \le\frac{3Q_N}{Q_N-1}e^{-\theta_NLMd/256}.
\]
The second gives
\[
 \sum_{n>n_*}\frac{|e_n|}{n}
 \le\frac3{1-q_N}
       e^{16d\log2-\vartheta_NLMd\log(1/q_N)}.
\]
Increase $M_N$ so that
$\vartheta_NLM\log(1/q_N)\ge32\log2$, and take
\[
 0<\delta_N\le
 L\min\{\theta_N/256,\tfrac12\vartheta_N\log(1/q_N)\}.
\]
For $d=0$, there are at most eight cubes and
$|E^N|\le2^8F_N$. These estimates, together with
\eqref{eq:normalizer-pivot-log}, prove the bond bound after fixing
$E_b$. Subtracting the trivial value costs no extra factor:
the nonpivot coefficient has the factor $1/2$ in its definition.
The tent support lies at distance less than $L<r_N$ from $b$.
Thus a nonzero term in \eqref{II-44-21} has $z\in X$, including
the pivot case, and $\sum_bh_z(b)=3L^3$ gives the pointed bound.

\emph{Connectedness and the normalization identity.}
We must also show that finite differences on regions that are not
wall connected vanish. Let $Y=Y_1\cup Y_2$, with no wall adjacency
between the two nonempty regions. A bare site lies in one coarse
cube, and a bond or plaquette cannot join the two components.
Every smaller partition cube whose center lies in $Y_i$ has its
interior there. For $x\in Y_1$, a coordinate-monotone coarse-cube
path to a cube of $Y_2$ first leaves $Y_1$ through a cube of $Y^c$;
the distance from $x$ does not increase before that exit. Hence
$\operatorname{dist}_\infty(x,Y^c)\le
\operatorname{dist}_\infty(x,Y_2)$.
The defining distance-to-complement test for a tower cube is
therefore the same in its component as in $Y$. Induction through
the levels gives
$\Omega_q(Y)=\Omega_q(Y_1)\cup\Omega_q(Y_2)$ with disjoint site
sets and equal level assignments.

The actual form now splits. Its derivative and Hessian stencils,
averaging blocks and contour paths remain in their component.
Consequently the scalar and field averages, their weighted
transposes, the scalar inverse and projection split. Sums,
products, inverses and bilinear polarization preserve this
decomposition. Applying these operations to the displayed
construction \eqref{eq:cv-form}, including both induced multipliers
and the pivot and gauge forms, gives
\[
 \mathsf A_Y=\mathsf A_{Y_1}\oplus\mathsf A_{Y_2}.
\]
The same holds after deep compression. In \eqref{II-44-16},
$f_b(Y)$ depends only on the component containing $S_N(b)$;
summing with alternating signs over any other component gives
zero. Thus the finite difference is supported on wall-connected
$X$.

M\"obius inversion on the finite set of coarse cubes now gives
\[
 \sum_{X\supset S_N(b)}E^N(X;P|_X,b)=f_b(T_\ell;P).
\]
The tents sum to one, so the sum of all pointed terms is
\[
 \Phi_t(P)=-\frac12\sum_{b\in\mathfrak X^+}
                    [f_b(T_\ell;P)-f_b(T_\ell;P_0)]
             +\sum_{c\in\Gamma}[z_c(P)-6\log L].
\]
In finite dimension,
$\exp(\sum_bf_b(T_\ell;P))=\det\mathsf A(P)$:
the logarithm series applies to $T_{T_\ell}$, whose norm is less
than one, and its trace is the sum of the diagonal traces.
Also $\exp z_c=\det(Lm_c(V))^{-1}$ and
$m_c(\mathbf1)=L^{-3}$. This proves \eqref{II-44-39}.
For real $P$, the gap is at least $3\gamma/8$ and the pivot
determinants are real and avoid the nonpositive axis. The real
Gaussian integral is positive, and its determinant formula,
together with $\Phi_t(P)\in\mathbb R$, proves
\eqref{II-44-40} with the real logarithms.

\emph{Local charts.}
Take a finite coarse region in the infinite lattice and embed
its halo in two tori of side at least $8LM$, without wrapping.
The chart between the images preserves levels, directions,
bond midpoints and all distances at most $LM$. The simple sets
correspond: plaquette and derivative stencils are carried
bijectively, and any compact cube contained in the region lifts
to a translate of the same Euclidean cube. The latter assertion
follows because the difference between the inverse chart and
a lift is a continuous, lattice-valued function on that cube,
hence is constant.

The chart therefore carries all tower tests, paths, blocks and
pivot bonds. The actual operator construction is a renaming of
the same finite coordinates, including every inverse, so
$A_{\mathsf iY,N}(\mathsf i_*P)(\mathsf ib,\mathsf ib')
=A_{Y,N}(P)(b,b')$. It also carries $S_N(b)$ and $X(c)$.
The series \eqref{II-44-14}, the finite differences and the pivot
determinants consequently agree term by term. The tents agree
because their values are functions of $(x_b-z)/L$ on their
support. This proves chart independence. It defines the infinite
family on each finite region, preserves slice holomorphy,
reality, gauge invariance and the bounds, and is consistent
when the region is placed in any larger torus.
\end{proof}

\section{The local interaction}\label{sec:local-interaction}
The fluctuation integral has three nonlinear terms: the Jacobian of the
constraint chart, the remainder of the gauge fixing, and the remainder of
the action at its critical configuration. The first two are supported on
one or two fine cubes. We localize the third by taking a finite difference
in each cube outside the plaquette being observed. Holomorphy in these
parameters gives exponential decay in the size of the resulting region.

Use the level-$k$ units and the triple domain of
\cref{def:gaussian-background-domain}, with $0\le k<K$.
Fix $\bar\kappa\ge1$. The fine cubes have side $M=L^{m'}\ge L^2$, the coarse cubes have side
$LM$, and the torus has side at least $C_LM$, where $C_L\ge8L$.
We write $\mathcal D_k$ for the wall-connected fine-cube regions and
$d_k(Y)$ for their linear size. All forms below are specified on positive
bonds; negative bonds carry their odd extensions. Kernels use counting
measure, and
\[
 (B_1,B_2)_1=\sum_{b\in\mathfrak X^+}
       \langle B_1(b),B_2(b)\rangle,
 \qquad \langle X,Y\rangle=\tfrac12\operatorname{tr}(XY).
\]
For $P=(U,A',\mathbf J)\in\mathcal Q_{k+1}(Y;a,a,a)$, put
$\mathbf U=e^{i\eta_{k+1}A'}U$ and $V=\mathcal M^k\mathbf U$ on
the level-$k$ bonds inside $Y$. Its existence for small $a$ is proved
below. The same bare pair, when used in the level-$k$ response equations,
is $(\mathbf U,L^{-3}\mathbf J)$, with decomposition $(U,L^{-1}A')$.

\subsection{A uniform chart for the pivot variables}\label{subsec:ymm1-50-1}

We first record the plain-average consequence of the centered iteration.
\begin{proposition}\label{I-27-78}
There are positive $a_{\rm plain},c_{\rm plain}$ and $C_Q\ge8$, depending
only on $L$, such that, on a tower of height $j$ over a box,
\begin{equation}\label{I-27-82}
 \varpi(U_0)\le a_{\rm plain}L^{-2j},\quad
 \tau=L^j\sup|a_0|<c_{\rm plain}
 \quad\Longrightarrow\quad
 \left|\frac1i\log\bigl[\mathcal M^j(e^{ia_0}U_0)
                  (\mathcal M^jU_0)^{-1}\bigr]\right|\le C_Q\tau.
\end{equation}
The averages and logarithm are holomorphic in $a_0$, local to the
endpoint double-block tower, and real for real $a_0$.
Here and below a local supremum in this assertion is taken on that tower.
\end{proposition}
\begin{proof}
To see the assertion, use the centered iteration of
\cref{I-39-36,I-39-43} and its one-step factorization. Decreasing
$a_{\rm plain},c_{\rm plain}$, the reference curvatures and centered
perturbations at every intermediate stage remain in the one-step mean
chart. The accumulated factors $v_i$ are sufficiently close to the
identity that their conjugations preserve the strict logarithm domains.
On these domains, the series logarithm commutes with conjugation, so the
one-step factorization remains valid for the accumulated complex gauges.
Induction in the factorization therefore gives the complex plain average
\[
 \mathcal M^j(e^{ia_0}U_0)(c)
   =v_j(c_-)e^{ia_j(c)}W_j(c)v_j(c_+)^{-1},
 \qquad W_j=\mathcal M^jU_0.
\]
The centered estimates give $|a_j(c)|\le2\tau$ and
$|v_j(c_\pm)^{\pm1}-1|\le C_v(L)\tau$. Choose
$C_Q\ge8(1+C_v)$ and $c_{\rm plain}\le C_Q^{-1}$.
Since $W_j$ is unitary, the relative product differs from $1$ by at most
$e^{(2+2C_v)\tau}-1\le C_Q\tau/2<1/2$.
The logarithm series gives \eqref{I-27-82}. The factorization also proves
holomorphy and locality; for real arguments the complex and real means
coincide at each stage by uniqueness of the local mean.
\end{proof}

Let $r_1=r_*/40$ and choose $\varsigma_{0,\mathrm c}>0$ within the
one-step coefficient radius established in
\cref{sec:fluctuation-local-expansion}. Thus, if $V$ on a double block
$\mathsf D_c$ has a unitary gauge $V^h=e^{i\widetilde a}$ with
$\sup|\widetilde a|\le\varsigma_{0,\mathrm c}$, the functions
\[
 f_V(B)=\frac1i\log\bigl[\mathcal M(e^{iB}V)(c)
                          \mathcal M(V)(c)^{-1}\bigr],
 \qquad
 x_V(B)(y,x)=\frac1i\log\bigl[(e^{iB}V)(y,x)V(y,x)^{-1}\bigr]
\]
are holomorphic and bounded by $1$ for $\sup|B|<r_1$, and vanish at
$B=0$. Write $\ell_V=f'_V(0)=LQ_V$ and $R_V=f_V-\ell_V$.
The pivot coefficient is $\mathsf m_c(V)=Q_V(\delta_{b_0(c)}\cdot)$.
Set $H_V=(L\mathsf m_c(V))^{-1}$ and
$\mathsf h_VD=(H_VD)\delta_{b_0(c)}$. Then
\begin{equation}\label{II-41-50}
 \begin{gathered}
 |H_V|\le2L^2,\qquad \ell_V\mathsf h_V=1,\qquad
 \sup|\mathsf C_V\mathcal A|\le2L^2r_1^{-1}\sup|\mathcal A|,\\
 |R_V(B)|\le2r_1^{-2}(\sup|B|)^2,\qquad
 |R_V(B_1)-R_V(B_2)|\le8s r_1^{-2}\sup|B_1-B_2|
 \quad(\sup|B_i|\le s\le r_1/4).
 \end{gathered}
\end{equation}
The last two bounds are the Cauchy remainder and derivative bounds for
$f_V$; the coefficient bounds and the right inverse were proved with the
fluctuation form.

\begin{definition}\label{II-41-43}
Put
\begin{equation}\label{II-41-44}
 s_3=\frac{r_1^2}{32L^2},\qquad r_3=\frac{s_3}{2},\qquad
 \rho_3=\frac{s_3}{4L^2},\qquad
 C_P=\max\{C_Q,2L^2/r_1\},\qquad r_P^\circ=\frac{r_3}{2C_P}.\qedhere
\end{equation}
\end{definition}
For $\sup|B_0|<r_3$, the map
$D\mapsto R_V(B_0-\mathsf h_VD)$ maps the closed ball
$|D|\le\rho_3$ into its quarter and has Lipschitz constant at most
$1/2$: indeed $\sup|B_0-\mathsf h_VD|<s_3$ and
$2(s_3/r_1)^2=\rho_3/4$, while
$8s_3r_1^{-2}\cdot2L^2=1/2$.
Its iterates from zero converge uniformly to a holomorphic map
$\mathsf D_V(B_0)$. Consequently
\begin{equation}\label{II-41-53}
 \begin{gathered}
 \psi^{\rm lin}_V(B_0)=B_0-\mathsf h_V\mathsf D_V(B_0),\qquad
 f_V(\psi^{\rm lin}_V(B_0))=\ell_VB_0,\\
 |\mathsf D_V(B_0)|\le4r_1^{-2}(\sup|B_0|)^2,\qquad
 \sup|\mathsf h_V\mathsf D_V(B_0)|\le\tfrac18\sup|B_0|.
 \end{gathered}
\end{equation}
The first bound follows from the contraction estimate and the value of
the map at zero; the second follows from $\sup|B_0|<r_3$.
Cauchy's inequality, applied in the pivot direction, gives
\begin{equation}\label{II-41-54}
 \mathsf N_V(B_0)=H_V\partial_{B_0(b_0(c))}\mathsf D_V(B_0),
 \qquad \mathsf N_V(0)=0,\qquad
 |\mathsf N_V(B_0)|\le\tfrac14\quad(\sup|B_0|\le r_3/2).
\end{equation}
For example, on the ball of radius $r_3$ the first bound in
\eqref{II-41-53} gives the derivative bound $8r_3/r_1^2$ on its
half-ball, which after multiplication by $2L^2$ is $1/4$.

All these maps commute with unitary gauges and are real on real data.
For a holomorphic family of $V$ presented in a fixed unitary gauge, the
coefficients and $H_V$ are holomorphic, and the uniform fixed-point
iteration proves joint holomorphy. The same uniqueness proves covariance
under translations and under a coordinate chart preserving the double
block, its pivot and its ordered paths. These charts preserve the actual
holonomies defining the means, hence preserve $f_V$ before the iteration
is applied. For a set of coarse bonds, the corrections just constructed
assemble componentwise: the only pivot in $\mathsf P_c$ is $b_0(c)$.
Thus $\mathsf h_V$ and $\mathsf C_V$ restrict to the one-double-block
maps and \eqref{II-41-53} holds in every row. Any holomorphic constraint
chart with the same germ at zero agrees with this one on their common
ball, by the identity theorem.

Choose $a_P^\circ>0$, depending only on $L$, so small that
\begin{equation}\label{eq:potential-background-radius}
 a_P^\circ\le\min\{a_{\rm plain},c_{\rm plain},c_2^{\rm SO},
                  \varsigma_{0,\mathrm c}/(2C_Q+8\pi)\}.
\end{equation}
Then every $P\in\mathcal Q_{k+1}(Y;a_P^\circ,a_P^\circ,a_P^\circ)$
has the block field $V$ used above. In fact, with $V_0=\mathcal M^kU$,
\begin{equation}\label{II-41-62}
 V=e^{iq}V_0,\qquad |q|\le C_QL^{-1}\sup|A'|,
 \qquad |\partial V_0-1|\le2L^{-2}a_P^\circ.
\end{equation}
Apply \eqref{I-27-82} to the tower over each endpoint double block with
$a_0=\eta_{k+1}A'$, for which $L^k\sup|a_0|=L^{-1}\sup|A'|$;
the real curvature estimate is the tower estimate for $U$.
For $k=0$ the same conclusions are immediate from $V=\mathbf U$.
On a double block, the comb gauge of $V_0$ has a real logarithm
of norm at most $4\pi L^{-1}a_P^\circ$. In this fixed gauge,
$V=e^{iq^h}e^{ia^0}$, and the logarithm of this product has norm at most
$2(C_Q+4\pi)L^{-1}a_P^\circ\le\varsigma_{0,\mathrm c}$.
This verifies the chart hypothesis uniformly. It also proves fixed-$U$
holomorphy, since this gauge does not depend on $A'$.

\begin{definition}\label{II-41-33}
For the restriction $\mathcal A_c$ to the nonpivot bonds of
$\mathsf P_c$, define
\begin{equation}\label{II-41-34}
 \beta_c(P,\mathcal A)
   =\psi^{\rm lin}_V(\mathsf C_V\mathcal A_c)(b_0(c)),\qquad
 \mathsf{Jb}_c(P,\mathcal A)
   =\mathsf N_V(\mathsf C_V\mathcal A_c).
\end{equation}
\end{definition}
If $\alpha=\sup|\mathcal A|<r_P^\circ$, then
\begin{equation}\label{eq:potential-pivot-bounds}
 |\beta_c|\le2C_P\alpha\le1,\qquad
 |\mathsf{Jb}_c|\le\frac{\alpha}{4r_P^\circ}<\frac14.
\end{equation}
The second estimate uses \eqref{II-41-54}, its vanishing at zero, and
the Schwarz estimate on each complex radial line.

\subsection{Finite differences and the three potentials}

For disjoint subsets $\mathcal K,\mathsf N$ of a finite parameter set,
use the brackets
\begin{equation}\label{II-18-108}
 F[\mathcal K\mid\mathsf N]
  =\sum_{S\subset\mathcal K}(-1)^{\#\mathcal K-\#S}
                           F(\mathbf1_{S\cup\mathsf N}),
 \qquad
 F^*[\mathcal K\mid\mathsf N]=(F-F(0))[\mathcal K\mid\mathsf N].
\end{equation}
If $F$ is holomorphic on $|s(\Delta)|<R$, $R>1$, and $|F|\le b$, then
\begin{equation}\label{II-18-113}
 |F[\mathcal K\mid\mathsf N]|\le b(R-1)^{-\#\mathcal K},\qquad
 |F^*[\mathcal K\mid\mathsf N]|\le2b(R-1)^{-\#\mathcal K}.
\end{equation}
Indeed the power series in one coordinate gives
$|F(s^{\Delta\to1})-F(s^{\Delta\to0})|
\le b\sum_{n\ge1}R^{-n}=b/(R-1)$, uniformly in the remaining
coordinates. Apply this successively to the coordinates of $\mathcal K$.
The starred assertion applies the same estimate to $F-F(0)$.
For $R=e^{\bar\kappa}$ with $\bar\kappa\ge1$,
$(R-1)^{-1}\le e^{-(\bar\kappa-1)}$, and therefore
\begin{equation}\label{II-18-114}
 |F[\mathcal K\mid\mathsf N]|\le b e^{-(\bar\kappa-1)\#\mathcal K}.
\end{equation}
Expansion of the finite sum also gives
\begin{equation}\label{II-18-110}
 F(\mathbf1_W)-F(0)=
       \sum_{\mathcal K\subset W\setminus\mathsf N}
                  F^*[\mathcal K\mid\mathsf N]
       \qquad(\mathsf N\subset W).
\end{equation}

For $k\ge1$, suppose a shift
$\mathfrak h_Y(s;P,\mathcal A)=(\mathbf A_Y,\mathbf K_Y)$ is given
on
\begin{equation}\label{II-41-25}
 P\in\mathcal Q_{k+1}(Y;a_h,a_h,a_h),\qquad
 \sup|\mathcal A|<\varepsilon_h,\qquad |s(\Delta)|<e^{\bar\kappa}.
\end{equation}
We require
\begin{equation}\label{II-41-24}
 0<a_h,\varepsilon_h\le1,\qquad c_h\ge1,
\end{equation}
and the following
properties. For fixed $U$ the shift is jointly holomorphic, depends on
$P$ only through $(\mathbf U,\mathbf J)$, and satisfies
\begin{equation}\label{eq:shift-uniform-bound}
 |\mathbf A_Y|, |\nabla_{\mathbf U}^{\eta_k}\mathbf A_Y|,
 |\mathbf K_Y|\le c_h\sup|\mathcal A|
\end{equation}
on every available bond or stencil. It vanishes at $s=0$. If the support
of $s$ has wall-components $\gamma$, the shift is the sum of its values
at $s\mathbf1_\gamma$, and that summand vanishes on bonds not inside
$\gamma$. Extension of $s$ by zero to a larger region restricts to the
same shift. The shift is gauge covariant, real for real inputs, and
commutes with translations by $M\mathbb Z^3$ and with nonwrapping
coordinate charts whenever the transported triple belongs to its domain.
These are hypotheses at this stage; the response construction below
provides the shift with these properties.

For a level-$k$ bond $b$, let $\square(b)$ be the union of the fine
cubes containing its endpoints. For a coarse bond $c$, define
$\square(c)$ in the same way; it equals $\square(b_0(c))$ and
contains $\mathsf D_c$. For a bare plaquette $p$, let $X(p)$ be the
union of the fine cubes containing its corners. These regions belong
to $\mathcal D_k$ and have linear size zero. Indeed $\square(b)$ and
$\square(c)$ consist of one or two adjacent cubes, and $X(p)$ consists
of one, two or four cubes with a common corner. There are at most
$7\#Y$ different $X(p)\subset Y$, and at most
$12(M/\eta_k)^3$ plaquettes with any one value of $X(p)$.
The first count chooses a base cube and a subset of at most two of the
three positive directions; the second chooses one of at most four cubes,
a bare site in it, and one of three plaquette directions.

For $\mathsf t>1$, put $g=\mathsf t^{-1/2}$ and $\mathcal A=gB$.
Write
\[
 \mathcal R_3f(B)=f(B)-\sum_{n=0}^2\frac1{n!}
                  \left.\partial_z^nf(zB)\right|_{z=0},\qquad
 \sigma_1(w)=\left(\sum_{n\ge0}\frac{(-w)^n}{(2n+1)!}\right)^2.
\]
All logarithms below are their series at $1$.
For a single fine cube define
\[
 \Psi^{\rm G}(\square;P,B)=
 \sum_{y\in\square^{(k+1)}}\sum_{x\in B(y)\setminus\{y\}}
 \left(1-\tfrac12\operatorname{tr}
                 [(e^{igB}V)(y,x)V(y,x)^{-1}]\right),
\]
and set this expression to zero on other domains. The Jacobian and
gauge potentials are
\begin{align}
 \mathbb V^{\rm J}(Y;P,B)
 &=\sum_{c:\,\square(c)=Y}
       \bigl\{\operatorname{tr}\log(1-\mathsf{Jb}_c)
                    +\log\sigma_1(\langle\beta_c,\beta_c\rangle)\bigr\}
       \notag\\
 &\quad+\sum_{b\in\mathfrak X^+:\,\square(b)=Y}
                    \log\sigma_1(\langle gB(b),gB(b)\rangle),
                    \label{II-41-37}\\
 \mathbb V^{\rm G}(Y;P,B)&=-\mathsf t\mathcal R_3
                      \Psi^{\rm G}(Y;P,\cdot)(B).
                      \label{II-41-36}
\end{align}
For the action, put $a_p(W)=1-\tfrac12\operatorname{tr}\partial W(p)$
and $W_1(X_1,W)=\eta_k^{-1}\sum_{p:X(p)=X_1}a_p(W)$. If $k\ge1$,
define
\begin{equation}\label{II-41-39}
 \Psi^{\rm A}(Y;P,B)=\mathbb V^{\rm act}(Y;P,gB)
 =\sum_{X_1\subset Y}
 \bigl[s\mapsto W_1(X_1,e^{i\eta_k\mathbf A_Y(s;P,gB)}\mathbf U)
                \bigr]^*[Y\setminus X_1\mid X_1],
\end{equation}
where only regions $X_1=X(p)$ contribute. If $k=0$, put
\begin{equation}\label{II-41-40}
 \Psi^{\rm A}(Y;P,B)=\sum_{p:X(p)=Y}
              [a_p(e^{iB'}\mathbf U)-a_p(\mathbf U)],
 \quad B'|_{\mathfrak X}=gB,\quad B'(b_0(c))=\beta_c(P,gB).
\end{equation}
Finally set
\begin{align}
 \mathbb V^{\rm A}&=-\mathsf t\mathcal R_3\Psi^{\rm A},\label{II-41-41}\\
 \mathbb V^{\rm int}&=\mathbb V^{\rm J}+\mathbb V^{\rm G}
                                      +\tau\mathbb V^{\rm A}.\label{II-41-42}
\end{align}

\begin{proposition}\label{II-41-79}
For $k\ge1$ set
$a_P=\min\{a_P^\circ,a_h\}$ and
$\varepsilon_{\rm A}=\min\{r_P^\circ,\varepsilon_h,c_h^{-1}\}$;
for $k=0$ set $a_P=a_P^\circ$ and $\varepsilon_{\rm A}=r_P^\circ$.
Put
\begin{equation}\label{II-41-45}
 \varepsilon_P=\varepsilon_{\rm A}/2,\qquad
 b_{\rm A}=6(1+1200)^2M^3,\qquad
 K_P=224b_{\rm A}M^3\mathsf C_{\mathcal A}\varepsilon_{\rm A}^{-3}.
\end{equation}
Suppose $\mathsf C_{\mathcal A}\ge3$ and
$\mathsf C_{\mathcal A}\delta(\mathsf t)\le\varepsilon_P$,
where $\delta(\mathsf t)=g\mathfrak b(\mathsf t)$.
The three potentials are defined for
\[
 P\in\mathcal Q_{k+1}(Y;a_P,a_P,a_P),\qquad
 B\in\mathbb D^{\rm f}_{\mathsf t}(Y)
 :=\{\sup_{\mathfrak X^\circ\langle Y\rangle}|B|
                              <\mathsf C_{\mathcal A}\mathfrak b(\mathsf t)\}.
\]
They are jointly holomorphic on each fixed-$U$ slice, depend only on the
pair, are gauge invariant and real on real data, and have the translation
and coordinate-chart covariance specified above. The Jacobian potential
vanishes at zero; the other two vanish to third order. For
$o\in\{{\rm J},{\rm G},{\rm A}\}$,
\begin{equation}\label{II-41-80}
 |\mathbb V^o(Y;P,B)|\le K_P\delta(\mathsf t)
     e^{-(\bar\kappa-2)d_k(Y)}
       \left(1+M^{-3}\sum_{b\in\mathfrak X^{\circ+}\langle Y\rangle}
                                      |B(b)|^2\right).
\end{equation}
All constants are independent of the volume, $k$ and $\mathsf t$.
The interaction terms are asserted only on this strict cutoff domain;
their later use outside it is through the cutoff in the integral.
\end{proposition}

\begin{proof}
We estimate the functions before taking their Taylor remainders on the
larger ball $g\sup|B|<\varepsilon_{\rm A}$. Let
$\alpha=g\sup|B|$. For $|w|\le1$, the sine series gives
$|\sigma_1(w)-1|\le(2/5)|w|$, hence
$|\log\sigma_1(w)|\le|w|$. Also
$|\operatorname{tr}\log(1-N)|\le4|N|$ when $|N|\le1/4$.
Together with \eqref{eq:potential-pivot-bounds}, these bound each pivot
summand in \eqref{II-41-37} by $2\alpha/r_P^\circ$ and each nonpivot
summand by $\alpha/r_P^\circ$. At most $6(M/L)^3$ coarse bonds and
$6M^3$ nonpivot bonds have a prescribed region $Y$. Thus
\begin{equation}\label{eq:potential-prebounds}
 |\mathbb V^{\rm J}|\le7M^3(r_P^\circ)^{-1}g\sup|B|,
 \qquad |\Psi^{\rm G}|\le M^3/2,
 \qquad |\Psi^{\rm A}|\le112b_{\rm A}
                                  e^{-(\bar\kappa-2)d_k(Y)}.
\end{equation}
The first function is zero unless $d_k(Y)=0$.
For the second inequality, the relative contour matrix $X$ satisfies
$|X-1|\le C_P\alpha\le1$. Since $\det X=1$,
$|1-\operatorname{tr}X/2|=|\det(X-1)|/2\le1/2$.
There are $(M/L)^3(L^3-1)<M^3$ contour pairs in a fine cube.

For the action bound we use the following plaquette estimate, valid for
complex configurations. If $0<\eta,\rho\le1$ and the two field
components and their two cross derivatives at a plaquette are bounded by
$\rho$, then
\begin{equation}\label{II-41-47}
 |\partial(e^{i\eta A}U)-1|
 \le(1+600\eta\rho)|\partial U-1|+600\rho\eta^2.
\end{equation}
To verify it, transport the four fields to the base point and write the
new plaquette as $e^ae^{b'}Ze^{-a'}e^{-b}$, where $Z=\partial U$,
$|a|,|b|\le\eta\rho$, $|a'|,|b'|\le2\eta\rho$, and
$|a'-a|,|b'-b|\le\eta^2\rho$. The part containing $Z-1$ is bounded
by $e^{6\eta\rho}|Z-1|$. In the remaining part replace $a',b'$ by
$a,b$ and then use
$|e^ae^b-e^be^a|\le2(e^{|a|}-1)(e^{|b|}-1)$.
The two replacements and the commutator cost at most
$[e^6+e^5+2e^2(e-1)^2]\rho\eta^2<600\rho\eta^2$, while
$e^{6\eta\rho}\le1+(e^6-1)\eta\rho$.

For $k\ge1$, apply \eqref{II-41-47} with
$\rho=c_h\alpha<1$ and $\eta=\eta_k$. The triple bounds imply
$|\partial\mathbf U-1|\le\eta_k^2$, so the resulting plaquette is
within $(1+1200)\eta_k^2$ of $1$. The determinant identity above and
the plaquette count give
\[
 |W_1(X_1,e^{i\eta_k\mathbf A_Y}\mathbf U)|
 \le\eta_k^{-1}\,12(M/\eta_k)^3
                    \frac{(1+1200)^2\eta_k^4}{2}=b_{\rm A}.
\]
The starred-bracket bound is therefore
$2b_{\rm A}e^{-(\bar\kappa-1)n}$, where $n=\#(Y\setminus X_1)$.
Since $d_k(X_1)=0$, we have $d_k(Y)\le n$. Indeed enumerate the
remaining cubes so that each is adjacent to a preceding cube. An edge
tree meeting the preceding region can be extended to the new cube with
at most one edge: from a corner of the adjacent old cube, move to its
common face if necessary. Starting at a common corner of $X_1$ proves
the bound. Also $\#Y\le8(d_k(Y)+1)$, because an edge tree of $d_k(Y)$
edges has at most $d_k(Y)+1$ vertices, each in at most eight cubes.
There are at most $7\#Y\le56e^{d_k(Y)}$ roots $X_1$ to sum, which
proves the third bound in \eqref{eq:potential-prebounds}.
For $k=0$, $|B'|\le2C_P\alpha$ and
$|\nabla_{\mathbf U}B'|\le3\sup|B'|$; the small radius gives
$6C_P\alpha\le1$. The same plaquette estimate bounds the difference
in \eqref{II-41-40} by $2b_{\rm A}$. It is zero unless $d_k(Y)=0$.

All displayed functions are jointly holomorphic on the larger ball.
For a bounded holomorphic function $f$ on a ball of radius $R$,
Cauchy's coefficient estimates along the line through $B$ give
\[
 |\mathcal R_3f(B)|\le2\sup|f|\,(\sup|B|/R)^3
              \qquad(\sup|B|\le R/2).
\]
Take $R=\varepsilon_{\rm A}/g$.
The cutoff hypothesis puts $\mathbb D^{\rm f}_{\mathsf t}(Y)$ in this
half-ball. Using $\mathsf tg^2=1$ and
$g\sup|B|\le\mathsf C_{\mathcal A}\delta(\mathsf t)$ gives
\[
 \mathsf t|\mathcal R_3f(B)|
 \le2\sup|f|\,M^3\mathsf C_{\mathcal A}\varepsilon_{\rm A}^{-3}
       \delta(\mathsf t)\,M^{-3}\sum_b|B(b)|^2.
\]
This proves \eqref{II-41-80} for the gauge and action terms using
\eqref{eq:potential-prebounds}; the Jacobian bound follows directly
from its first inequality and \eqref{II-41-45}.
The vanishing orders follow from the definitions.
Pair dependence, reality and gauge invariance follow before taking the
brackets or remainders, since the local chart is covariant and traces are
invariant. Translations and coordinate charts carry every defining path,
pivot, plaquette and cube to its counterpart; restrictions and the finite
brackets commute with these bijections. This proves the remaining claims.
\end{proof}

\subsection{The quadratic gauge term and the global sums}

Write the contour derivative as
$\mathsf G^{(1)}[V](B)(y,x)=x'_V(0)B(y,x)$. It defines a symmetric
operator on the nonpivot bonds inside a fine cube by
\begin{equation}\label{II-41-31}
 (B_1,\widehat{\mathsf G}_\square(P)B_2)_1
 =\sum_{y\in\square^{(k+1)}}\sum_{x\in B(y)\setminus\{y\}}
  \langle\mathsf G^{(1)}[V](B_1)(y,x),
                         \mathsf G^{(1)}[V](B_2)(y,x)\rangle.
\end{equation}
Its kernel is holomorphic on fixed-$U$ slices and has row sums at most
\begin{equation}\label{II-41-71}
 c_{\mathsf G}=6L^6r_1^{-2}.
\end{equation}
In fact its kernel is $\sum_{y,x}\mathsf b_V(y,x;b)^{\mathsf T}
\mathsf b_V(y,x;b')$, with $|\mathsf b_V|\le r_1^{-1}$ and support
inside $B(y)$. For fixed $b$, at most one $y$, fewer than $L^3$ values
of $x$, and at most $3L^3$ bonds $b'$ contribute. The transpose costs
a factor of at most $2$ in the matrix norm, giving \eqref{II-41-71}.
Define the free quadratic potential, for arbitrary $B$, by
\begin{equation}\label{II-41-32}
 \mathbb V^{\rm q}(\square;P,B)
       =-\frac{1-\tau}{2}(B,\widehat{\mathsf G}_\square(P)B)_1,
\end{equation}
and set it to zero on other regions. Expansion of
$\tfrac12\operatorname{tr}e^{ix}=1-\tfrac12\langle x,x\rangle
+O(|x|^4)$ at $x=x_V(gB)$ gives
\begin{equation}\label{II-41-83}
 \mathbb V^{\rm G}(\square;P,B)
 =-\mathsf t\Psi^{\rm G}(\square;P,B)
             +\tfrac12(B,\widehat{\mathsf G}_\square(P)B)_1.
\end{equation}

We describe exactly when these local expressions recover the global
ones. Let $U_+$ be a real critical background over a coarse field $W$,
choose its averaged field $V=\mathcal M^kU_+$ in axial gauge, and assume
$P_W=(U_+,0,J^{(k+1)}(U_+))\in
\mathcal Q_{k+1}(T_\ell;a_P,a_P,a_P)$.
Suppose that the real fibre chart and the critical background over
$V_W(B)=e^{i\psi^{\rm lin}_V(\mathsf C_VgB)}V$ exist for
$\sup|B|<\mathfrak b(\mathsf t)$. For $k\ge1$ require, in addition,
that the full-support shift $\mathbf A=\mathbf A_{T_\ell}(1;P_W,gB)$
satisfy the background identification
\begin{equation}\label{II-41-27}
\begin{gathered}
 e^{i\eta_k\mathbf A}U_+\sim_{\rm gauge} U_k(V_W(B)),\\
 J^{(k)}(e^{i\eta_k\mathbf A}U_+)
       =L^{-3}J^{(k+1)}(U_+)+\mathbf K_{T_\ell}(1;P_W,gB).
\end{gathered}
\end{equation}
Here $U_k(V_W(B))$ denotes a representative of the critical background
over $V_W(B)$, and $\sim_{\rm gauge}$ denotes gauge equivalence.
The current equality states the normalization needed later; the action
identity below uses only the first assertion.

Under these hypotheses, the local Jacobian terms sum to the nonlinear
correction to the chart Jacobian, including the exponential-coordinate
density. The constant linear-coordinate determinant is separate.
To verify the nonlinear correction, let $\Pi$ retain the pivot coordinates and
$N=\mathsf h_V\partial\mathsf D_V(\mathsf C_VgB)$. Since
$\Pi N=N$ and each double block has exactly one pivot,
$\Pi N\Pi$ is block diagonal with blocks $\mathsf{Jb}_c$.
Thus $\operatorname{Tr}N^n=\sum_c\operatorname{tr}\mathsf{Jb}_c^n$,
and the absolutely convergent logarithm series splits as in
\eqref{II-41-37}. The remaining density factors are precisely
$\sigma_1(\langle\beta_c,\beta_c\rangle)$ at pivots and
$\sigma_1(\langle gB(b),gB(b)\rangle)$ elsewhere.
The one-double-block chart agrees with the global chart by the germ
uniqueness proved above, so no second choice of pivot variable occurs.

The gauge terms sum by sorting the pairs $(y,x)$ according to their fine
cube. On inner block bonds the pivot elimination and nonlinear correction
leave $gB$ unchanged. Since $V(y,x)=1$ in axial gauge, the sum of
$\Psi^{\rm G}$ is the gauge-fixing function $G_W(B)$, and
\begin{equation}\label{II-41-86}
 \widehat{\mathsf G}_W
   =\sum_{\square\in\pi_k}\chi_\square^{\mathsf T}
       \widehat{\mathsf G}_\square(P_W|_\square)\chi_\square.
\end{equation}
Consequently the sum of \eqref{II-41-83} is
$-\mathsf tG_W(B)+(B,\widehat{\mathsf G}_WB)_1/2$.

Finally let $\mathbb F_W(B)$ be the Wilson action of the critical
background over $V_W(B)$, in level-$k$ units. Then
\begin{equation}\label{II-41-84}
 \sum_{Y\in\mathcal D_k}\mathbb V^{\rm act}(Y;P_W|_Y,gB|_Y)
                         =\mathbb F_W(B)-\mathbb F_W(0).
\end{equation}
For $k\ge1$, fix $X_1=X(p)$ and apply \eqref{II-18-110} to
$F_{X_1}(s)=W_1(X_1,e^{i\eta_k\mathbf A_{T_\ell}(s)}U_+)$.
Consistency identifies the brackets with those in \eqref{II-41-39}.
A support component disjoint from $X_1$ does not change the fields on
its plaquettes; if $\mathcal K\cup X_1$ is disconnected, pairing
subsets in the alternating sum over any such component makes its bracket
zero. The remaining brackets are indexed exactly by connected
$Y\supset X_1$. Summing first over $Y$ and then over $X_1$ therefore
gives
$A^{\eta_k}(e^{i\eta_k\mathbf A_{T_\ell}(1)}U_+)
-A^{\eta_k}(U_+)$, since the shift at zero vanishes. Gauge invariance
and \eqref{II-41-27} identify this with the right side of
\eqref{II-41-84}; the shift at $B=0$ vanishes by its bound.
For $k=0$, the constrained background is the field itself, and every
plaquette appears exactly once in \eqref{II-41-40}.

The left side of \eqref{II-41-84} is holomorphic for
$g\sup|B|<\varepsilon_{\rm A}$. If the critical action has the
quadratic expansion
$\mathbb F_W(zB)-\mathbb F_W(0)
=z^2g^2(B,\mathsf BB)_1/2+O(z^3)$ for real $B,z$, comparison of
its first three Taylor coefficients gives
\begin{equation}\label{II-41-85}
 \sum_Y\mathbb V^{\rm A}(Y;P_W|_Y,B|_Y)
 =-\mathsf t[\mathbb F_W(B)-\mathbb F_W(0)]
                       +\tfrac12(B,\mathsf BB)_1.
\end{equation}
This is the exact action remainder used with the Gaussian fluctuation
form; the background and quadratic-expansion hypotheses must be checked
at that application.

\subsection{Construction of the shift}

We now supply the shift used above. Assume $k\ge1$ and the geometric
hypotheses of \cref{prop:decoupled-local-propagators}, with partition
width $w=M$ and parameter radius $e^{\bar\kappa}$. In particular,
\[
 M\ge88r_w\mathfrak M,\qquad
 \mathring\delta_0M\ge2176\bar\kappa.
\]
Assume also $12\mathfrak M\le\mathsf oL^2M$ and the room condition
$C_L\ge8\mathsf n_{\rm e}L$, where $\mathsf n_{\rm e}\ge1$ is the
number of coarse layers used in coordinate-chart halos.
Use the constant frame of height $k$: all its regions through level
$k$ are the torus, and its terminal region is empty. This terminal
region belongs to the response frame; the localization step still
has its entire torus available. Thus all constraint rows have level
$k$, their weights are $1$, and the bare spacing is $\eta_k$.

There are $0<c_{\mathfrak s}\le1\le C_{\mathfrak s}$, depending only
on $L$, for which one may take
\begin{equation}\label{II-51-27}
 a_{\mathfrak h}=\frac{c_{\mathfrak s}}
                      {\max\{1,\mathsf oLM B^{(\mathrm{cg})}\}},
 \qquad
 \varepsilon_{\mathfrak h}
       =c_{\mathfrak s}\min\{1,\varepsilon_{\mathrm v}\}
                         e^{-C_{\mathfrak s}\bar\kappa},
 \qquad
 \mathsf c_{\mathfrak h}=C_{\mathfrak s}
                         e^{C_{\mathfrak s}\bar\kappa}.
\end{equation}
Use $a_h=a_{\mathfrak h}$, $\varepsilon_h=\varepsilon_{\mathfrak h}$
and $c_h=\mathsf c_{\mathfrak h}$ in the preceding subsection.
These choices are uniform in $k$, the cutoff and the torus. To specify
the required smallness, set
$b_*=\max\{4b_{\mathrm D},8C_jb_{\mathrm D}^3\}$ and
$C_\times=2C_Pb_*$.
Choose $C_{\mathfrak s}\ge\max\{32,2C_Pb_*\}$ and choose
$c_{\mathfrak s}$ no larger than
\[
 \min\{a_P^\circ,r_P^\circ,La_{\mathrm D},
       \varepsilon_{\mathrm D}/(2C_P),
       (1200C_\times)^{-1},1\}.
\]
Decreasing it further preserves every assertion below.

For $P\in\mathcal Q_{k+1}(Y;a_{\mathfrak h},a_{\mathfrak h},
a_{\mathfrak h})$ and $\sup|\mathcal A|<\varepsilon_{\mathfrak h}$,
define a datum on all positive level-$k$ bonds by
\begin{equation}\label{II-51-33}
 B'_Y(P,\mathcal A)(b)=
 \begin{cases}
  \mathcal A(b),&b\text{ is a nonpivot bond inside }Y,\\
  \beta_c(P,\mathcal A),&b=b_0(c),\ \square(c)\subset Y,\\
  0,&\text{otherwise}.
 \end{cases}
\end{equation}
Here $\beta_c$ is the pivot value of
$\psi^{\rm lin}_V(\mathsf C_V\mathcal A)$ from the preceding
subsection. The first two cases exhaust the level-$k$ bonds inside
$Y$. In particular,
\begin{equation}\label{eq:shift-source-size}
 \sup|B'_Y|\le2C_P\sup|\mathcal A|
       <\varepsilon_{\mathrm D}e^{-32\bar\kappa}.
\end{equation}
Let $s$ vanish off $Y$. Form the decoupled response on this constant
frame with pair $(\mathbf U,L^{-3}\mathbf J)$ restricted to the
inner bare bonds of $Y$, and put
\begin{equation}\label{II-51-35}
 \mathfrak h_Y(s;P,\mathcal A)
   =\bigl(\mathbf H(s;B'_Y),j(s;B'_Y)\bigr)
                                    \big|_{\operatorname{Bd}(Y)}.
\end{equation}
The restriction includes the incident crossing bonds.

\begin{proposition}\label{II-51-52}
The maps \eqref{II-51-35} satisfy the shift hypotheses
\eqref{eq:shift-uniform-bound}, including their support,
restriction, gauge, reality and coordinate-chart properties, with
the constants \eqref{II-51-27}. They commute with complex conjugation.
In the real global setting described before \eqref{II-41-27}, they
satisfy that identification, provided the backgrounds there exist
as single regular critical gauge orbits.
\end{proposition}

\begin{proof}
We first verify that the response is defined. Read the triple in
level-$k$ units with decomposition $(U,L^{-1}A')$. Its natural-class
radii can be taken to be
\[
 (\alpha_0,\alpha_1,\gamma_0)
    =(L^{-1}c_{\mathfrak s},L^{-1}a_{\mathfrak h},
                                      L^{-3}a_{\mathfrak h}).
\]
They are at most $a_{\mathrm D}$. The curvature bound of the simple
domain is $a_{\mathfrak h}\eta_{k+1}^2
=L^{-2}a_{\mathfrak h}\eta_k^2$. Every cube needed by the local
natural class has side at most $12\mathfrak M$ in level-$k$ units,
hence at most $12\mathfrak M/L\le\mathsf oLM$ in level-$(k+1)$
units. The simple domain supplies its own gauge on that cube, with
logarithm and derivative bounded by
$\mathsf oLM B^{(\mathrm{cg})}a_{\mathfrak h}\le c_{\mathfrak s}$.
Changing units multiplies this logarithm by $L^{-1}$ and its
derivative by $L^{-2}$. The intermediate scale factors in the
natural-class inequalities are at most $1$. The same changes of units
apply to $A'$, while the current is multiplied by $L^{-3}$.
These observations verify all the local natural-class inequalities
on the inner bonds used by the response. Equation
\eqref{eq:shift-source-size} now permits \cref{I-88-110}.

The source is holomorphic at fixed real $U$, depends only on the pair,
and restricts consistently: each pivot value uses just its endpoint
double-block tower. The response has these properties by
\cref{I-88-110}. Its estimates, since every active scale factor is
$1$, give
\[
 \sup|\mathbf H|,\quad\sup|\nabla_{\mathbf U}^{\eta_k}\mathbf H|,
 \quad\sup|j|
 \le b_*e^{24\bar\kappa}\sup|B'_Y|
 \le\mathsf c_{\mathfrak h}\sup|\mathcal A|.
\]
The derivative estimate also applies at the incident boundary. Indeed,
the component response is supported over $Y^1(\gamma)$ and its current
over $Y^2(\gamma)$. The margin construction ensures that a bare bond
over $Y^j(\gamma)$, $j\le4$, has both endpoints in
$\operatorname{int}\gamma$. It also places the whole differentiation
stencil based over $Y^3(\gamma)$ inside that region. Thus a nonzero
stencil uses the same pair as the response estimate; outside that
margin both terms of the derivative vanish. No choice of a pair
outside $Y$ is used in this argument.

For completeness, each endpoint or plaquette corner in this assertion
is at intrinsic distance at most $\rho_0+5/2\le r_0$ from its base
site. The margin estimate in the local propagator construction puts
it in the hull of a cell $X$ with $\mathsf P(X)\subset\gamma$.
Every partition cube containing this point then belongs to
$\mathsf P(X)$; their union is a neighborhood of the point. Thus
the point lies in $\operatorname{int}\gamma$. The local $8L$ cube
used to compare the nonlinear formulas also lies in $\gamma$:
if the base site belongs to $\mathsf q^{+3}$, then
$4L\xi_y\le4L^3\xi_m\le2\mathfrak M\xi_m$, so this cube is
contained in $\mathsf q^{\boxplus}\subset\bigcup\mathsf P(X)$.
Consequently the component sum in \cref{I-88-110} gives exactly the
support and splitting assertions required here. At $s=0$ every walk
monomial vanishes, since every trail meets a partition cube; the
linear operators and both fixed-point iterates are zero, hence so
are the shift and current. Extension by zero preserves the source
rows used by each component and therefore preserves the response.
Unitary gauges preserve all the constructions. Complex conjugation
commutes with the local mean series, with the pivot iteration and
with the response iteration; it conjugates $s$ as well as the fields.
This proves the reality assertions, including conjugation covariance.

We give the coordinate-chart argument because the support at crossing
bonds matters. A translation by $M\mathbb Z^3$, or a chart of a
nonwrapping domain, is on each local lift a translation by a vector
in $\mathfrak M\mathbb Z^3$. It preserves the fine cubes, the
underlying lattices, the ordered averaging paths and the endpoint
blocks. Restrict first to one component $\gamma$ of
$\{s\ne0\}$. The map takes each cell trail inside $\gamma$ to its
translated trail with the same label and weight. Its inverse gives
the inverse correspondence of trails. Each local inverse and
projection in a walk is the inverse or projection of the translated
finite matrix, so its kernel is carried to that of the translated
walk. Products, sums, weighted transposes and the convergent inverse
series preserve this equality. The parameter monomial is unchanged
when $s$ is transported with its cubes. Thus the four operators in
the response equations are intertwined by the chart.

On the $8L$ cubes around the sites of $Y^3(\gamma)$, the local
holonomies, currents and logarithmic averages agree after transport.
Their Taylor series and derivatives therefore agree as well. Starting
the two fixed-point iterations at zero now proves equality of every
iterate, and uniform convergence proves equality of the responses
and their currents. The pivot source is carried to the corresponding
source by its already proved locality and uniqueness. This proves
transport on the bonds over $Y^2(\gamma)$. Both fields vanish on all
other incident bonds. There is no extra nonzero value on the other
side: preservation of endpoint blocks and injectivity of the chart
imply that a transported bond is over the transported margin exactly
when the original bond was over the original margin. Summing the
components proves the asserted chart covariance. As before, the
transported triple must belong to the stated domain.

It remains to identify the real global response. Put
$V=\mathcal M^kU_+$ and $B'=B'_{T_\ell}(P_W,\mathcal A)$.
The finer constraint set over $V$ is contained in the coarser one
over $W$, so $U_+$ is critical for the finer constraint. The multiplier
argument of \cref{I-88-127}, at $A=0$, gives
$J^{(k)}(U_+)=Q^{\mathsf T}\lambda_0$.
At $s=1$ the response is independent of the auxiliary partition and
its parameter radius. By \cref{I-88-127}, with
$A=\mathbf H(1;B')$ and $K_A=\mathcal Q'(A)$,
\begin{equation}\label{eq:shift-stationarity}
 \mathcal Q(A)=B',\qquad F'(A)=K_A^{\mathsf T}\omega(A),\qquad
 J^{(k)}(e^{i\eta_kA}U_+)=L^{-3}J^{(k+1)}(U_+)+j(1;B').
\end{equation}
All rows of this frame are interior. Its actual-average factorization is
\begin{equation}\label{eq:shift-actual-average}
 \mathcal M^k(e^{i\eta_kA}U_+)
                         =(e^{iB'}V)^{v[A]}.
\end{equation}

We check that stationarity in the reduced coordinates implies
criticality for this actual average. For a real node function $\sigma$
let
\[
 \ell_{B'}\sigma(c)=\left.\frac d{dt}\right|_{t=0}
 \frac1i\log\!\left[e^{it\sigma(c_-)}e^{iB'(c)}V(c)
                          e^{-it\sigma(c_+)}V(c)^{-1}\right].
\]
The function
$z(B';\sigma)=(\omega(A),\ell_{B'}\sigma)_1$ is holomorphic in
$B'$ and linear in $\sigma$. For a small real gauge variation
$w_t=e^{it\chi}$, write
$e^{i\eta_kA^{w_t}}U_+=(e^{i\eta_kA}U_+)^{w_t}$ and set
\[
 X_\chi(A)=\left.\partial_tA^{w_t}\right|_0,\qquad
 \rho(A;\chi)=\left.\frac1i\partial_t
            \bigl(v[A^{w_t}]^{-1}w_tv[A]\bigr)\right|_0.
\]
Differentiating \eqref{eq:shift-actual-average} gives
$K_AX_\chi(A)=\ell_{B'}\rho(A;\chi)$.
At $A=0$, the derivative of the accumulated mean factors gives
$\rho(0;\chi)=Q'\chi$, the scalar block average. It is onto:
the blocks $B^k(y)$ are disjoint, and one may choose one site in
each block and assign its value so that its transported average is
any prescribed node value. In finite dimensions surjectivity persists
for $A$ sufficiently close to zero. Gauge invariance of $F$ and
\eqref{eq:shift-stationarity} therefore give
\[
 z(B';\rho(A;\chi))
       =(F'(A),X_\chi(A))_{\eta_k}=0
\]
for all node values and all real $B'$ near zero. Holomorphy first
extends this equality to a complex neighborhood and then, along
complex lines, to the whole response ball $\sup|B'|<\varepsilon_{\mathrm D}$.

Now let $E_t$ be any differentiable real curve through
$E=e^{i\eta_kA}U_+$ with the actual average in
\eqref{eq:shift-actual-average} fixed. Its local logarithm gives
$E_t=e^{i\eta_kA_t}U_+$. Put
$h_t=v[A_t]^{-1}v[A]$ and
$\sigma=(1/i)\partial_th_t|_0$. The fixed-average identity yields
$K_A\dot A=\ell_{B'}\sigma$.
Consequently the derivative of its action is
$\eta_k(F'(A),\dot A)_{\eta_k}
=\eta_k z(B';\sigma)=0$.
Thus $E$ is critical for its actual average.

Choose a bare gauge whose values on the level-$k$ sites are
$v[A]^{-1}$. It makes this average $e^{iB'}V=V_W(\mathcal A)$.
Finally the estimates at $s=1$ give
\[
 \sup|A|,\quad\sup|\nabla_{U_+}^{\eta_k}A|,
 \quad\sup|j(1;B')|\le C_\times\sup|\mathcal A|=:\rho,
 \qquad 600\rho<\tfrac12\min\{1,\varepsilon_{\mathrm v}\}.
\]
The plaquette estimate \eqref{II-41-47} bounds the curvature of $E$
by $[(3/2)L^{-2}+1/2]\varepsilon_{\mathrm v}\eta_k^2
<\varepsilon_{\mathrm v}\eta_k^2$; its current is bounded by
$(L^{-3}+1/2)\varepsilon_{\mathrm v}<\varepsilon_{\mathrm v}$.
These are the regularity inequalities at level $k$. They are gauge
invariant. The assumed uniqueness of the regular critical orbit
therefore identifies $E$ with the background over $V_W(\mathcal A)$.
The current identity in \eqref{eq:shift-stationarity} proves the other
assertion of \eqref{II-41-27}.
\end{proof}

\subsection{Tower estimates and substitution in the analytic domain}
\label{sec:tower-core-substitution}

The pointwise bounds defining the simple domain do not by themselves
give the same bounds for its constrained backgrounds. We first estimate
those backgrounds, including their dependence on complex data. This
will give a quantitative margin for substituting the fluctuation into
an interaction from an earlier scale.

All determining sets in this subsection are reduced: every inward bond
is removed. We write them as $\mathbf B^\flat$, including those of
truncated towers. Fix the geometric thresholds of
\cref{I-58-22,I-88-53}, and decrease the regularity threshold and the
complex datum radius so that the canonical background belongs to the
natural regular-pair class required by \cref{I-88-53}.
Write $\bar\varepsilon$ and $r_{\mathrm B}$ for these thresholds.
They are uniform in the depth and the volume. We use $B_3\ge6$ for
the real-background regularity constant, increased if necessary.

\begin{lemma}\label{I-82-15}
For every regular real datum $V_0$ on $\mathbf B^\flat$, put
$U_0=U^\circ_{\mathbf B^\flat}(V_0)$. On
$\sup|B|<r_{\mathrm B}$ there is a holomorphic potential
$\mathcal H_{\mathbf B^\flat,V_0}(B)$, vanishing at zero, such that
\begin{equation}\label{I-82-16}
 \Phi_{\mathbf B^\flat,V_0}(B)
   =(e^{i\eta\mathcal H(B)}U_0)^{u_B},\qquad
 \mathsf R_{\mathsf m}D^{\mathsf T}\mathcal H(B)=0,
 \qquad u_B(y)v_q[U_0;\eta\mathcal H(B)](y)=I
       \quad(y\in\Lambda_q).
\end{equation}
The potential and $u_B$ are real and unitary, respectively, at real
data. For suitable $B_5\ge1$ and $0<\delta_{\mathrm{tw}}\le1$, the
derivative $h=D\mathcal H(B)e_c$ satisfies
\begin{equation}\label{I-93-47}
 \max\{\xi_y|h(b)|,\ \xi_y^2|\nabla_{U_0}h(b)|,
              \ \xi_y^3|D_{U_0}^{\mathsf T}D_{U_0}h(b)|\}
 \le B_5e^{-\delta_{\mathrm{tw}}d(y,y_c)}|e|.
\end{equation}
Here an endpoint of $b$ lies in $\Delta(y)$, $y_c$ is an endpoint
of $c$ in its layer, the first differences use only inner bonds,
and the last member is asserted where the four incident plaquettes
are inner. The constants are independent of the depth, the sequence
of regions and the volume.
\end{lemma}

\begin{proof}
Use the potential in the proof of \cref{I-58-22}:
$\mathcal H(B)=\mathbf H(B_\natural)$, with the actual datum converted
to the straightened datum by the endpoint factors. The construction
of the axial representative gives \eqref{I-82-16}. At a top, the
ascent is empty, so the axial gauge factor there is exactly the
inverse of the accumulated factor $v_q$. The identities first hold
at real data and extend to the complex row ball by holomorphy.
The fixed-point uniqueness and the gauge normalization give uniqueness
of the pair near zero, and hence throughout its connected domain.

We give the derivative argument, since the third bound in
\eqref{I-93-47} is needed below. The derivative of $B\mapsto
B_\natural$ in a direction $e_c$ is bounded by $C(L)|e|$ and
is supported in a bounded number of rows within distance $4L$
of $y_c$. Indeed, interior rows read $B(c)$ directly; an outward
row uses the one-step endpoint factor, which depends only on the
interior rows of the adjacent block. Differentiating the local
logarithm on half its domain gives the stated bound. There are no
edge rows for the reduced determining set.

Take $w(y)=e^{\theta d(y,y_c)}$ with $\theta>0$ smaller than the
rate in \eqref{eq:response-weighted-bounds}. Differentiating that
estimate bounds the weighted potential and first-difference norm
of $h$ by $C(L)|e|$. The multiplier is
\[
 \lambda(A)=(I+H^{\mathsf T}\mathsf C'_A{}^{\mathsf T})^{-1}
                         H^{\mathsf T}F'(A).
\]
Its derivative has the same weighted bound. To see this without
estimating a second difference of $h$ in advance, pair
$F''(A)h$ against a column of $H$ and use the differentiated
energy identity \eqref{eq:response-local-energy}. The curl term
is paired with the first differences of that column, and the
remaining terms have the local bounds of \cref{I-87-84}.
The column estimate and \eqref{eq:actual-response-transpose}
then bound $H^{\mathsf T}F''(A)h$ in the weighted row norm by
$C\|h\|_w$. The derivative of
$H^{\mathsf T}\mathsf C'_A{}^{\mathsf T}$ has the same bound
by the local averaging estimate. Its undifferentiated operator
has weighted norm less than $1/2$ after reducing the datum radius,
so its inverse is bounded by two. Thus
$|\dot\lambda|_w\le C\|h\|_w$.

Differentiate the stationarity equation. Since the gauge equation
also differentiates to $\mathsf R_{\mathsf m}D^{\mathsf T}h=0$,
\[
 D^{\mathsf T}Dh
   =(Q+\mathsf C'_A)^{\mathsf T}\dot\lambda
     +\bigl[D_A\mathsf C'_A{}^{\mathsf T}h\bigr]\lambda
     -\Delta'h-\mathcal V'(A)h.
\]
Every term on the right has weighted source norm at most
$C(|\dot\lambda|_w+\|h\|_w)$: this is the transpose averaging
bound and the local action-remainder bound, with their fixed
neighborhood enlargement. This proves the last estimate.
Changing from the initial endpoint to either endpoint costs at
most a factor $L^3$ in the scale and a fixed exponential factor
in the distance. Increase $B_5$ and decrease
$\delta_{\mathrm{tw}}$ to include these factors.
\end{proof}

For an inner bond $b$ over $y$, let $N_b(f)$ denote the maximum
on the left of \eqref{I-93-47}, omitting the last member when
its plaquette star is not inner. The exponential summation
bound gives finite constants $B^\sharp\le B^\flat$ such that
\begin{align}
 N_b(\mathcal H(B_1)-\mathcal H(B_2))
  &\le B^\sharp\sup|B_1-B_2|,\label{I-93-61}\\
 N_b(\mathcal H(B_1)-\mathcal H(B_2))
  &\le B^\flat e^{-\delta_{\mathrm{tw}}D/2}
                  \sup|B_1-B_2| \label{eq:tower-distant-data}
\end{align}
in the second line when the data agree at all rows within
distance $D$ of $y$. The same second estimate holds for the
linearization $\mathsf H=D\mathcal H(0)$. These inequalities
follow by integrating the derivative along the segment between
the two data and summing \eqref{I-93-47}. Cauchy's estimate on
the line $zB$ also gives
\begin{equation}\label{eq:tower-quadratic-response}
 N_b(\mathcal H(B)-\mathsf HB)
   \le 2B^\sharp r_{\mathrm B}^{-1}(\sup|B|)^2,
       \qquad \sup|B|\le r_{\mathrm B}/2.
\end{equation}

A pillar $B^q(y)$ of $\mathbf B^\flat$ is \emph{free} if
$q\ge1$ and it contains no node of $\mathbf B^\flat$ other
than its top $y$.

We need one quantitative comparison between the axial and
Landau representatives. On a free pillar of level $q$, put
$h_y=\sup_{B^q(y)}\xi_y|\mathcal H(B)|$. For a sufficiently
small $r_{\mathrm G}\le r_{\mathrm B}$,
\begin{equation}\label{I-93-66}
 \Phi(B)=W(B)^{u(B)},\qquad W(B)=e^{i\eta\mathcal H(B)}U_0,
 \qquad |u(B)(x)^{\pm1}|\le e^{c_uh_y}\le e^{2/25}
\end{equation}
on the free pillars, provided the base curvature there is in
the averaging domain. Here $c_u$ depends only on $L$.
For completeness, the gauge is
$u(B)(x)=v_q[U_0;\eta\mathcal H(B)](y)^{-1}\mathcal T_{W(B)}(x)$,
where $\mathcal T$ is the product of contour variables along
the ascent from $x$ to $y$. At level $i<q$ the logarithm of
its contour factor is bounded by
$C(L)L^{i+1-q}h_y$, by \cref{I-27-78}; the accumulated
top factor satisfies the same geometric sum bound. Summing
$\sum_{i<q}L^{i-q}\le(L-1)^{-1}$ gives the first bound in
\eqref{I-93-66}. Choose $r_{\mathrm G}$ so that
$c_uB^\sharp r_{\mathrm G}\le2/25$ and every contour logarithm
stays in its domain. At real data this is the axial gauge
normalization; both sides are holomorphic, so the identity
extends to complex data. In particular the condition number
of $u(B)(x)$ is less than $6/5$.

We record the local expansions used to pass between the potential,
curvature and current. Their constants do not deteriorate under a
change of lattice scale.

\begin{lemma}\label{I-10-120}
Let $\eta\le1/13$. On a plaquette based at $x$, number the edges
in directions $\mu,\nu,\mu,\nu$, with initial points
$x,x+\eta e_\mu,x+\eta e_\nu,x$, respectively.
Let $\phi=\max|f|$, let $\mathbf W$ be a complex configuration,
and suppose $|T||T^{-1}|\le\varkappa$ for
$T=\mathbf W_1,\mathbf W_4,\partial\mathbf W,
\partial\mathbf W\mathbf W_4$. If
$|\partial\mathbf W-I|\le1/4$ and $4\varkappa\eta\phi\le1/4$, then
$\partial(e^{i\eta f}\mathbf W)=R\partial\mathbf W$, where
\begin{equation}\label{I-10-121}
 |R-I-i\eta^2D_{\mathbf W}f|
 \le11\varkappa^2\eta^2\phi^2
       +6\varkappa\eta\phi|\partial\mathbf W-I|.
\end{equation}
On the four plaquettes incident to a bond, let $U_0$ be real,
$h=\max\{|f|,|\nabla_{U_0}f|\}\le1/2$, and
$f_0=\max\eta^{-2}|\partial U_0-I|\le1$. Then
\begin{equation}\label{I-10-122}
 |J(e^{i\eta f}U_0)-J(U_0)-D_{U_0}^{\mathsf T}D_{U_0}f|
                       \le134(h^2+f_0h).
\end{equation}
Only differences between bonds of these plaquettes enter $h$.
\end{lemma}

\begin{proof}
Number the edges of the positively oriented plaquette so that its
product is $W_1W_2W_3^{-1}W_4^{-1}$. Moving these transports to the
right expresses $R$ as four exponentials with generators
\[
 f_1,\quad\operatorname{Ad}(W_1)f_2,\quad
 -\operatorname{Ad}(\partial W\,W_4)f_3,\quad
 -\operatorname{Ad}(\partial W)f_4.
\]
Their sum is
\[
 \eta D_Wf+(1-\operatorname{Ad}(\partial W))
                         (\operatorname{Ad}(W_4)f_3+f_4).
\]
The second term is at most $6\varkappa\phi|\partial W-I|$.
For exponentials whose total exponent norm is $s\le1/4$,
the terms of degree at least two sum to at most $s^2e^s/2$.
With $s\le4\varkappa\eta\phi$ this gives \eqref{I-10-121}.

For the current, retain one more Taylor order at the real base.
Writing the four exponents as $Y_1,\ldots,Y_4$, one has
$Y_3=-Y_1-\delta_1$, $Y_4=-Y_2-\delta_2$, with
$|\delta_i|\le(15/13)\eta^2h$. The quadratic commutator is
\[
 \frac12\sum_{i<i'}[Y_i,Y_{i'}]
 =[Y_1,Y_2]+\frac12
  (-[Y_1,\delta_1]-2[Y_2,\delta_1]-[Y_2,\delta_2]
                                             +[\delta_1,\delta_2]).
\]
The second term is bounded by $5\eta^3h^2$; replacing
$\operatorname{Ad}(U_{0,1})f_2$ in the first term by $f_\nu$
costs $2\eta^3h^2$. The third-order remainder in
$(R-R^{-1})/2$ is at most $15\eta^3h^3$. Including the base
plaquette and its inverse therefore gives
\[
 F_{e^{i\eta f}U_0}
   =F_{U_0}+D_{U_0}f+i[f_\mu,f_\nu]+\rho_F,
 \qquad |\rho_F|\le24\eta(h^2+f_0h),
 \qquad |F_{e^{i\eta f}U_0}|\le3f_0+6h.
\]
Insert this into the covariant backward difference defining $J$.
For each of the two transverse directions, replacing its perturbed
transport by the real transport costs at most
$15(h^2+f_0h)$. The difference of the commutator costs at most
$4h^2$, since each transported difference of $f$ is bounded by
$\eta h$. The two remainders cost at most $48(h^2+f_0h)$.
The linear terms are exactly $J(U_0)+D_{U_0}^{\mathsf T}D_{U_0}f$.
Thus the total error is $2(15+4+48)(h^2+f_0h)$, as asserted.
\end{proof}

\begin{proposition}\label{I-93-83}
There are $a^\sharp\in(0,1/60]$ and
$K^\sharp_{\mathrm{tw}}\ge1$, uniform in the depth and volume,
with the following property. In level-$n$ units, let
$\eta=L^{-n}$ and let $\mathsf L_-(e)$ be the smallest pillar
scale among the endpoints or corners of a bond or plaquette.
Suppose $\epsilon_0+\epsilon_1\le a^\sharp$ and
\[
 |\partial U-I|,\ |\partial(e^{i\eta f}U)-I|
       <\epsilon_0\eta^2\mathsf L_-^{-2},\qquad
 |f|<\epsilon_1\mathsf L_-^{-1},
\]
where $U$ is real. The tower configuration at the averaged
datum of $e^{i\eta f}U$ exists and is holomorphic in $f$ on
an open neighborhood of these data. If the corners of a
plaquette, or all corners of the four plaquettes incident to
a bond, lie in free pillars of one level with scale $\xi$, then
\[
 |\partial U_{\mathbf B^\flat}-I|
   <K^\sharp_{\mathrm{tw}}(\epsilon_0+\epsilon_1^2)
                         \eta^2\xi^{-2},\qquad
 |J(U_{\mathbf B^\flat})|
   <K^\sharp_{\mathrm{tw}}(\epsilon_0+\epsilon_1^2)\xi^{-3}.
\]
The same assertions hold with $U$ in place of $e^{i\eta f}U$.
\end{proposition}

\begin{proof}
Let $C_Q$ be the constant in \cref{I-27-78}, and choose
\[
 0<a_{\mathrm{av}}\le
       \min\{a_{\rm plain},c_{\rm plain},c_2^{\rm SO},1\},
 \qquad (4C_Q+2)a_{\mathrm{av}}<\omega_M.
\]
This common radius is admissible for the plain average and the real
curvature estimate in \cref{I-10-101}. A real configuration with
curvature $\alpha\eta^2\mathsf L_-^{-2}$ and
$L^2\alpha\le a_{\mathrm{av}}$ induces a
$c^\natural\alpha$-regular datum, with $c^\natural=2L^2$.
At a tested level-$q$ plaquette its four factors are the
values of $\mathcal M^qU$; the averaging corridor lies in
$\Omega_{q-1}$, where the bare curvature is bounded by
$L^2\alpha L^{-2q}$. The real curvature averaging estimate
therefore gives this assertion, including outward tests.
At level zero it is the original curvature bound.

Choose
\[
 r_0=\frac1{2L}\min\{a_{\mathrm{av}},r_{\mathrm G}/C_Q\},
 \qquad
 a^\sharp\le\min\{1/60,r_0,
       \bar\varepsilon/(22c^\natural),
       a_{\mathrm{av}}/(22B_3c^\natural)\}.
\]
Equip forms with $\|w\|=\sup_b\mathsf L_-(b)|w(b)|$.
For $\|w\|<2r_0$, the plain-average logarithm
$B(w)=Q(U;\eta w)$ satisfies
\[
 M_{\mathbf B^\flat}(e^{i\eta w}U)=e^{iB(w)}V_0,
 \qquad V_0=M_{\mathbf B^\flat}U,
 \qquad \sup|B(w)|\le C_QL\|w\|<r_{\mathrm G}.
\]
Thus the tower is $\Phi_{V_0}(B(w))$ and exists on this entire
complex ball. Put $W(w)=e^{i\eta\mathcal H(B(w))}U_0$.
For the plaquette or bond in the statement, let $q\ge1$ be
the common pillar level and put $\xi=L^{q-n}$. Thus the
bare spacing in level-$q$ units is $\eta/\xi=L^{-q}\le1/13$.
The two maps
\[
 F^\partial(w)=\eta^{-2}\xi^2(\partial W(w)-I),
       \qquad F^J(w)=\xi^3J(W(w))
\]
are holomorphic there and bounded by one, after decreasing
$r_{\mathrm G}$ if necessary. Indeed,
\eqref{I-93-61} controls the potential, its first difference
and $D^{\mathsf T}D$ in the local scale. The plaquette
expansion and the current expansion then bound these maps
by
\[
 B_3c^\natural a^\sharp+3h+11h^2,
 \qquad B_3c^\natural a^\sharp+h
       +134(h^2+B_3c^\natural a^\sharp h),
\]
where $h=B^\sharp r_{\mathrm G}$ can be chosen at most $1/300$.

Write $f=f_{\mathrm R}+if_{\mathrm I}$. The projections onto
the two real parts are contractions, also after applying
$D_U$. The plaquette expansion gives
\[
 |D_Uf|\le\mathsf L_-^{-2}(2.1\epsilon_0+11\epsilon_1^2).
\]
Applying the same expansion to $f_{\mathrm R}$ and
$f_{\mathrm I}$ shows that each real configuration
$e^{i\eta f_{\mathrm R}}U$ and $e^{i\eta f_{\mathrm I}}U$
has curvature at most
$22(\epsilon_0+\epsilon_1^2)\eta^2\mathsf L_-^{-2}$.
Its induced datum is therefore regular. The real-background
bound and \eqref{I-93-66} imply, for either $F=F^\partial,F^J$,
\[
 |F(0)|\le B_3c^\natural\epsilon_0,\qquad
 |F(f_{\mathrm R})|,|F(f_{\mathrm I})|
 \le\frac65\,22B_3c^\natural(\epsilon_0+\epsilon_1^2).
\]

The Taylor terms of order at least two of a holomorphic map
bounded by one on $\|w\|<2r_0$ have sum at most
$2\|w\|^2/(2r_0)^2$ on the half ball. Its linear term is
complex linear. Subtract it at $f$, $f_{\mathrm R}$ and
$f_{\mathrm I}$ to obtain
\[
 |F(f)|\le3|F(0)|+|F(f_{\mathrm R})|+|F(f_{\mathrm I})|
                       +\frac{6\|f\|^2}{4r_0^2}
 \le(56B_3c^\natural+2r_0^{-2})(\epsilon_0+\epsilon_1^2).
\]
Finally \eqref{I-93-66} conjugates these quantities to those
of the tower, with loss less than $6/5$. Take
$K^\sharp_{\mathrm{tw}}\ge2(56B_3c^\natural+2r_0^{-2})$.
The argument at $f=0$ proves the last assertion. All choices
precede the depth, the sequence and the volume.
\end{proof}

\begin{proposition}\label{I-93-98}
The following consequences hold with the same constants.
\begin{enumerate}
\item[(f)] Let $(U,A',\mathbf J)\in
\mathcal Q_p(Z;\alpha_0,\alpha_1,c_0)$ and
$\alpha_0+\alpha_1\le a^\sharp$. For both tower rules and
every truncation $1\le n\le p$, the tower quantities exist
and, on $Z^{\wr p}$, satisfy
\begin{equation}\label{I-93-99}
 |\partial U_n^{\mathfrak r}-I|
  <K^\sharp_{\mathrm{tw}}(\alpha_0+\alpha_1^2)L^{-2p},
 \qquad
 |J_n^{\mathfrak r}|
  <K^\sharp_{\mathrm{tw}}(\alpha_0+\alpha_1^2)L^{2(n-p)}.
\end{equation}
They are holomorphic for fixed $U$ on the open simple domain.
The estimates also hold for the tower of $U$. Consequently,
if $\alpha_i\le b_i$ and
$K^\sharp_{\mathrm{tw}}(\alpha_0+\alpha_1^2)\le b_0$, then
\begin{equation}\label{I-93-100}
 \mathcal Q_j(X;\alpha_0,\alpha_1,c_0)
             \subset\mathcal U_j^\circ(X;b_0,b_1,c_0).
\end{equation}
The trivial triple belongs to every core with positive radii.

\item[(g2)] Let $1/40\le d<1$, $\vartheta=1-d$, and
\[
 0<b_0\le b_1,\qquad b_1^2\le c_\natural b_0,
 \qquad b_0+b_1\le a^\sharp,\qquad c_\natural\ge1,
 \qquad b_2=\frac{d^2b_0}{32K^\sharp_{\mathrm{tw}}(1+c_\natural)}.
\]
For $P=(U,A',\mathbf J)\in
\mathcal U_j^\circ(X;\vartheta b_0,\vartheta b_1,c_0)$,
put $\mathbf U=e^{i\eta_jA'}U$. If
$|\mathbf A|,|\nabla_{\mathbf U}\mathbf A|<b_2$ and
$|\mathbf J+\mathbf K|<b_0$, then
\[
 P\dotplus(\mathbf A,\mathbf K)
  :=\left(U,\frac1{i\eta_j}
       \log(e^{i\eta_j\mathbf A}e^{i\eta_jA'}),
                   \mathbf J+\mathbf K\right)
       \in\mathcal U_j^\circ(X;b_0,b_1).
\]
The logarithm and the substitution are locally holomorphic.
\end{enumerate}
\end{proposition}

\begin{proof}
For (f), pass to level-$n$ units and set
$\lambda=L^{-(p-n)}$, $f=\lambda A'$. The hypotheses of
\cref{I-93-83} hold with
$\epsilon_0=\lambda^2\alpha_0$ and
$\epsilon_1=\lambda\alpha_1$. On the two-layer interior the
relevant plaquettes and their stars lie in free pillars of
level $n$, so their local scale is one. Since
$\lambda^2\eta_n^2=L^{-2p}$, this is \eqref{I-93-99}.
For a subtower $(p,Z)$ of a level-$j$ domain, rescale the
triple by $(U,A',\mathbf J)\mapsto(U,\lambda_pA',
\lambda_p^3\mathbf J)$, $\lambda_p=L^{-(j-p)}$.
The bound in \eqref{I-93-99} is then at most
$\lambda_pK^\sharp_{\mathrm{tw}}(\alpha_0+\alpha_1^2)$,
which proves every clause in \eqref{I-93-100}.
At the trivial triple the datum, canonical background and
Landau potential are trivial, proving the final assertion.

For (g2), put
$R=2K^\sharp_{\mathrm{tw}}(1+c_\natural)/d$ and
$\beta=Rb_2=db_0/16$. On $|z|\le R$ write
\[
 A'_z=(i\eta_j)^{-1}
          \log(e^{iz\eta_j\mathbf A}e^{i\eta_jA'}).
\]
The logarithm is defined since
$\eta_j(\beta+\vartheta b_1)\le1/780$.
The logarithmic remainder has Lipschitz constant
$\epsilon=8\eta_j(\beta+\vartheta b_1)<0.0103\le d/2$.
Thus
\[
 |A'_z|<\vartheta b_1+(1+\epsilon)\beta<b_1.
\]
The change from $\nabla_{\mathbf U}$ to $\nabla_U$ costs
at most $\frac14\sup|\mathbf A|$. Equivariance and the same
Lipschitz bound for the logarithmic remainder give
\[
 |\nabla_UA'_z|
   <(1+\epsilon)(\vartheta b_1+5\beta/4)<b_1.
\]
For a plaquette, apply the local product expansion with
condition number at most $11/10$. Its relative increment
has norm at most $(5/2)\beta\eta_j^2$, and hence
\[
 |\partial(e^{iz\eta_j\mathbf A}\mathbf U)-I|
   <(1-d+5.01d/32)b_0\eta_j^2<b_0\eta_j^2.
\]
The real factor is unchanged and the new current is bounded
by hypothesis. The entire disc therefore lies in the simple
domain with radii $(b_0,b_1)$.

Fix a subtower $(p,Z)$ and put $\lambda_p=L^{-(j-p)}$.
Normalize either its curvature or current by the corresponding
power in \eqref{I-93-99}. By (f) the resulting holomorphic
function $g(z)$ is bounded on $|z|<R$ by
\[
 K^\sharp_{\mathrm{tw}}(\lambda_pb_0+\lambda_p^2b_1^2)
       \le K^\sharp_{\mathrm{tw}}(1+c_\natural)\lambda_pb_0.
\]
At zero the core hypothesis gives
$|g(0)|<(1-d)\lambda_pb_0$. Schwarz's lemma applied to
$g-g(0)$ yields $|g(1)-g(0)|\le d\lambda_pb_0$.
This is the required tower bound. The tower of the unchanged
real factor already satisfies the stronger bound at zero.
Local holomorphy follows from the convergent matrix-logarithm
series. Every argument is local and also applies to finite
domains in the infinite lattice.
\end{proof}

\begin{proposition}\label{II-16-69}
Use the core and tower hypotheses of \cref{I-93-98},
at level \(j\ge1\), and put \(C_{\rm os}=18\).
The half-radius form of outer substitution has
\begin{equation}\label{II-16-70}
 b_2=\frac{b_0}{128K^\sharp_{\mathrm{tw}}(1+c_\natural)}.
\end{equation}
For \(P\in\mathcal U_j^\circ(X;b_0/2,b_1/2,c_0)\),
with the radii of \cref{I-93-98}(g2),
\begin{equation}\label{II-16-71}
 |\mathbf A|,|\nabla_{\mathbf U}\mathbf A|<b_2,
 \quad |\mathbf J+\mathbf K|<b_0
 \quad\Longrightarrow\quad
 P\dotplus(\mathbf A,\mathbf K)\in
                       \mathcal U_j^\circ(X;b_0,b_1).
\end{equation}

More generally, let \(0<\tau<1\), \(c_\natural\ge1\), and
\begin{equation}\label{II-16-73}
 \begin{gathered}
 0<b_0\le b_1,\qquad b_1^2\le c_\natural b_0,\qquad
 2(b_0+b_1)\le a^\sharp,\\
 b_2=\frac{\tau b_0}{C_{\rm os}K^\sharp_{\mathrm{tw}}
                                      (1+c_\natural)}.
 \end{gathered}
\end{equation}
If \(P\in\mathcal U_j^\circ(X;(1-\tau)b_0,(1-\tau)b_1,c_0)\)
and the three bounds on the left of \eqref{II-16-71} hold,
then the same conclusion holds. In both cases the logarithmic
substitution is locally holomorphic, and its pair is
\((e^{i\eta_j\mathbf A}\mathbf U,\mathbf J+\mathbf K)\).
\end{proposition}

\begin{proof}
The first assertion is \cref{I-93-98}(g2) with \(d=1/2\).
For the second, a complex disc with twice the target radii
will give an increment proportional to \(\tau\).
We first record the pointwise estimate used on that disc.
Write \(\xi=\eta_j\le1/13\), \(\mathbf U=e^{i\xi A'}U\),
and suppose the simple input bounds are \((\beta_0,\beta_1)\).
For \(\beta_0,\beta_1,s\le1/60\) and
\(|f|,|\nabla_{\mathbf U}f|<s\), put
\[
 A^f=(i\xi)^{-1}\log(e^{i\xi f}e^{i\xi A'}).
\]
Then
\begin{equation}\label{II-16-52}
 \begin{gathered}
 |A^f-A'|<\tfrac{11}{10}s,\qquad
 |\nabla_U(A^f-A')|<\tfrac{11}{10}s,\\
 |\partial(e^{i\xi f}\mathbf U)-\partial\mathbf U|
                                      <\tfrac94s\xi^2.
 \end{gathered}
\end{equation}
All logarithms are defined since
\(\xi(s+\beta_1)\le1/390\).

Here are the increment estimates, including the cancellation
in the derivative bound. Set \(\epsilon_*=1/380\) and
\(q_*=(1+\epsilon_*)/(1-\epsilon_*)<51/50\).
At a bond, the two matrices \(e^{i\xi f}e^{i\xi A'}\)
and \(e^{i\xi A'}\) lie in the \(\epsilon_*\)-ball of \(I\).
The logarithm series gives
\(|A^f-A'|\le q_*|f|\).
At a derivative stencil, transport the neighboring values
with its real bond \(u=U_\mu(x)\), and put
\[
 a_1=i\xi f_\nu(x),\quad
 a_2=i\xi\operatorname{Ad}(u)f_\nu(x+\xi e_\mu),\quad
 c_1=i\xi A'_\nu(x),\quad
 c_2=i\xi\operatorname{Ad}(u)A'_\nu(x+\xi e_\mu).
\]
Changing from \(\mathbf U\) to \(U\) gives
\[
 |a_2-a_1|<\xi^2s(1+\omega_*),\qquad
 |c_2-c_1|<\xi^2\beta_1,\qquad
 \omega_*=2\beta_1(1+\epsilon_*)<1/29.
\]
Indeed the difference of the two transported \(f\)'s
is a conjugation by \(e^{i\xi A'_\mu(x)}\); its norm is
at most \((e^{2\xi\beta_1}-1)|f|\).
Let \(E_i=e^{a_i}e^{c_i}\), \(H_i=e^{c_i}\),
and \(D_i=E_i-H_i\). The telescoping identity for
powers in the logarithm series gives
\[
 \begin{split}
 &|(\log E_2-\log H_2)-(\log E_1-\log H_1)|\\
 &\quad\le\frac{|D_2-D_1|}{1-\epsilon_*}
 +\frac{(|E_2-E_1|+|H_2-H_1|)|D_1|}
                                     {2(1-\epsilon_*)^2}.
 \end{split}
\]
For example, write each difference of \(n\)-th powers
as a sum of \(n\) products with one increment. Differencing
these products leaves either the increment \(D_2-D_1\)
or one increment of \(E\) or \(H\); their multiplicities
are \(n\) and \(n(n-1)/2\), respectively.
The exponential difference bounds give
\[
 \begin{split}
 |\Delta|
 &\le q_*(|a_2-a_1|+\xi s|c_2-c_1|)\\
 &\quad+\tfrac12q_*^2\xi s
                    (|a_2-a_1|+2|c_2-c_1|)\\
 &<\xi^2s\left\{q_*(1+\omega_*)+q_*\xi\beta_1
       +\tfrac12q_*^2\xi[s(1+\omega_*)+2\beta_1]\right\}
 <\tfrac{11}{10}\xi^2s .
 \end{split}
\]
Here \(\Delta\) is the logarithmic difference on the left;
it equals \(i\xi^2\nabla_U(A^f-A')\).

For the plaquette bound use \cref{I-10-121} with
\(\varkappa=101/100\). The required transport condition
numbers are at most
\[
 e^{2\xi\beta_1}
      \frac{1+\beta_0\xi^2}{1-\beta_0\xi^2}<1.003
                                      <\varkappa .
\]
The two contained derivatives give \(|D_{\mathbf U}f|<2s\).
The product remainder therefore bounds the plaquette
increment by
\[
 \xi^2s[2+11\varkappa^2s+6\varkappa\xi\beta_0]
                      (1+\beta_0\xi^2)<\tfrac94\xi^2s.
\]
This proves \eqref{II-16-52}. The real-factor gauge
conditions are unchanged, so the substituted triple has
simple radii \((\beta_0+\frac94s,\beta_1+\frac{11}{10}s)\),
with its current bounded separately.

Now set
\[
 \beta_\iota=(1-\tau)b_\iota,\qquad
 \gamma_\iota=2b_\iota,\qquad
 s=\tfrac49(1+\tau)b_0,\qquad R=s/b_2>2.
\]
For \(|z|<R\), apply \eqref{II-16-52} with \(f=z\mathbf A\).
The resulting triples, with current set to zero for this
tower estimate, have simple radii \((\gamma_0,\gamma_1)\):
\[
 \beta_0+\tfrac94s=2b_0,\qquad
 \beta_1+\tfrac{11}{10}s\le2b_1,\qquad
 \gamma_0+\gamma_1\le a^\sharp.
\]
The logarithms and the tower quantities are holomorphic
on this disc by \cref{I-93-98}(f).
For a subtower \((p,Z)\), let \(\lambda=L^{-(j-p)}\).
After division by the curvature or current scale in
\eqref{I-93-99}, each tower quantity \(F(z)\) is bounded by
\(\lambda K^\sharp_{\rm tw}(\gamma_0+\gamma_1^2)\).
At zero the core hypothesis gives
\(|F(0)|<(1-\tau)\lambda b_0\).
Schwarz's lemma applied to \(F-F(0)\) gives
\[
 |F(1)-F(0)|
 \le2\lambda K^\sharp_{\rm tw}(\gamma_0+\gamma_1^2)
                        \frac{b_2}{s}
 \le\tau\lambda b_0,
\]
because \(\gamma_0+\gamma_1^2\le4(1+c_\natural)b_0\)
and \(C_{\rm os}=18\).
Thus every tower has the required radius \(b_0\).
The real-factor towers are unchanged. At \(z=1\),
\eqref{II-16-52} with \(s=b_2\) gives the simple target
radii, since \(b_2\le\tau b_0/36\); the current has
the required bound by hypothesis. This proves core
membership. The logarithm series also proves local
holomorphy and the displayed identity of the pair.
\end{proof}

\subsection{Inner substitution on an octet}
\label{sec:inner-octet-substitution}

An octet $\square$ of level $n$ is the union of the eight cubes
of $\pi_{n-1}$ meeting a vertex. Put $Z=\square^{\sim5}$.
Thus $Z$ has side $12M/L$ in level-$n$ units. We use its
$\mathfrak r_1$ tower and its truncations, with reduced determining
sets $\mathbf B^\square_q$, $q\le n$. The fields below depend only
on the configuration on the inner bonds of $Z$.

For positive radii $\alpha_0,\alpha_1$, the window domain on $Z$
consists of pairs $(U,A')$ on $\operatorname{Bd}(Z)$, with $U$ real,
whose decomposition $\mathbf U=e^{i\eta_nA'}U$ satisfies
conditions (i), (ii), and the curvature inequality in (iii) of
\cref{II-16-12}, with radii $\alpha_0,\alpha_1$. We also require
the tower conditions in that definition for $p=n$, the region $Z$,
and the rule $\mathfrak r_1$, at every truncation $1\le q\le n$,
for both the real and complex data, with $c=\alpha_0$ in
\eqref{eq:analytic-core-tower-bounds}.

There is a convenient way to put the coarse datum near the identity.
Let $x_\square$ be the centre of the octet and, in level-$n$ units,
put $y_0=x_\square+(e_1+e_2+e_3)/2$, a level-$n$ site in the octet;
let $\Gamma_{y_0,x}$ be its coordinate comb, and set
\[
 g(x)=\mathbf U(\Gamma_{y_0,x}),\qquad
 W=\mathbf U^g,\qquad
 V=M_{\mathbf B^\square_n}(W),\qquad v=g^{-1}|_{\mathbf B^\square_n}.
\]
For sufficiently small $a_\square=a_\square(L)>0$, if
$\mathbf U=e^{i\eta_nA'}U$ satisfies the simple bounds
$(\alpha_0,\alpha_1)$ on $Z$ and
$(LM)^3(\alpha_0+\alpha_1)\le a_\square$, then
\begin{equation}\label{eq:octet-comb-datum}
 M_{\mathbf B^\square_n}(\mathbf U)=V^v,\qquad
 |V-I|<K_\square M\alpha_0,\qquad
 |v(x)|\,|v(x)^{-1}|<e^{K_\square M\alpha_1},\qquad v(y_0)=I,
\end{equation}
where $K_\square=5C_Q$. These maps are holomorphic for fixed $U$.
They commute with conjugation, constant rotations, translations
preserving $\pi_{n-1}$, and the corresponding translations in local
charts.

Here are the estimates that verify the assertion, including the use
of a complex gauge. In bare units write the side of $Z$ as $N$, so
$N\eta_n=12M/L$. A comb has at most $3N$ edges. Comparing its product
with the real comb gives
$\mathsf k(g)\le e^{72(M/L)\alpha_1}\le11/10$.
In comb gauge the first-direction bonds are trivial; those of the
second direction on the initial coordinate plane are trivial; and
those of the third direction on the initial line are trivial.
Each remaining bond is consequently a product of at most $2N$
plaquette variables or their inverses. If their distance from $I$
is $\epsilon\le1/100$, the product is at distance at most
$e^{(200/99)N\epsilon}-1$ from $I$. Hence
\[
 |W-I|\le30(M/L)\alpha_0\eta_n,\qquad
 |a_2|\le33(M/L)\alpha_0\eta_n,\quad a_2=i^{-1}\log W.
\]
To justify averaging covariance, let $u$ be the real comb and put
$g'=gu^{-1}=e^{i\lambda}$. Then
\[
 |\lambda|\le54(M/L)\alpha_1<1/16,\quad
 \eta_n^{-1}|\lambda(x+\eta_ne_\mu)-\lambda(x)|
       \le235((M/L)\alpha_0+\alpha_1).
\]
The latter estimate follows by writing
$W=(\mathbf U^u)^{g'}$, taking the local logarithm, and using
$|i^{-1}\log(\mathbf U^u)|
 \le\tfrac32(\alpha_1+27(M/L)\alpha_0)\eta_n$.
Decrease $a_\square$ until all these logarithms and their averaging
corridors lie in the domain of \cref{I-27-78}. Averaging covariance
holds there at real $(a_2,\lambda)$ and therefore, by the identity
theorem on the convex logarithmic domain, at complex arguments.
Thus $M^q(W)=g_-M^q(\mathbf U)g_+^{-1}$.
The bound $|Q_q(a_2)|\le33C_Q(M/L)\alpha_0\le1/10$ gives
$|V-I|\le37C_Q(M/L)\alpha_0<K_\square M\alpha_0$.
This proves \eqref{eq:octet-comb-datum}. The construction uses finite
products, logarithms and the same local averaging branch; it gives
the asserted holomorphy and covariance. For example a real gauge
$h$ sends $V$ to $\operatorname{Ad}(h(y_0))V$ and $v(x)$ to
$h(x)v(x)h(y_0)^{-1}$.

Fix $0<\beta_{\rm s}\le1/4$. Increase $K^\sharp_{\mathrm{tw}}$ so that, with
\[
 K=K^\sharp_{\mathrm{tw}},\qquad
 c_h=2C_Q(B^\sharp)^2K_\square,
\]
one has $K\ge\max\{4K_\square,c_h,1\}$. There is a floor
$M_\flat=M_\flat(L,B_5,\delta_{\mathrm{tw}},\beta_{\rm s})$ and, after
$M\ge M_\flat$ is fixed, a positive radius $a^\flat$ independent
of the depth and volume with the properties used below. More
explicitly, choose the floor so that
\[
 \delta_{\mathrm{tw}}M\ge8,\qquad
 7C_QL^3B^\flat B^\sharp K_\square M
                  e^{-\delta_{\mathrm{tw}}M/4}\le\beta_{\rm s}.
\]
Let $r_{\mathrm I}$ be a common smaller radius for the free-pillar
gauge comparison, averaging, changes of real base datum and
restriction to a truncated tower. Choose $a^\flat\le a^\sharp$
so that
\begin{equation}\label{eq:inner-substitution-smallness}
\begin{gathered}
 (LM)^3(1+2\beta_{\rm s})a^\flat\le a_\square,\qquad
 8(C_Q+C_QLB^\sharp KM)a^\flat\le r_{\mathrm I},\qquad
 B^\sharp KMa^\flat\le10^{-4},\\
 11C_Q^2(B^\sharp)^3K_\square^2M^2a^\flat\le\beta_{\rm s} r_{\mathrm B},
 \qquad134(c_hM)^2a^\flat\le\beta_{\rm s},\\
 [2c_uc_h+4(c_u+c_w)B^\sharp K_\square+2K_\square]
                   Ma^\flat\le\beta_{\rm s}/2,\qquad c_wa^\flat\le1/6.
\end{gathered}
\end{equation}
The constant $c_w=C(L)$ bounds the sum of contour logarithms
in a truncated tower, exactly as in \eqref{I-93-66}.
These are finitely many upper bounds on a radius chosen after $M$.

\begin{proposition}\label{prop:inner-substitution}
Let $1\le j<n$, $\nu=n-j$, and let $t_\nu>0$ satisfy, with
$\Theta_s=\sup_{\nu\ge1}t_\nu L^{-s\nu}$,
\begin{equation}\label{I-93-95}
 \Theta_1\le1,\quad
 (1+10\beta_{\rm s})\Theta_2+\beta_{\rm s}/4+1/1000\le1/8,\quad
 (1+3\beta_{\rm s})\Theta_3+5\beta_{\rm s}\Theta_2+\beta_{\rm s}/8+1/500\le1/16.
\end{equation}
Suppose
\begin{equation}\label{I-93-101}
 A_0+A_1\le a^\flat,\quad a_0+a_1\le a^\flat,\quad
 2B_3^2KMa_0<a_1,\qquad A_0\le t_\nu a_0.
\end{equation}
Let $(U,A')$ belong to the window domain on $Z$ with radii
$((1+2\beta_{\rm s})A_0,(1+2\beta_{\rm s})A_1)$, and put
$\mathbf U=e^{i\eta_nA'}U$. In particular its
$\mathfrak r_1$ tower exists and has the corresponding curvature
and current bounds on $\square^{\sim3}$. Use the datum $V$ above and put
\[
 B_V=i^{-1}\log V,\qquad H_n=\mathcal H_{\mathbf B^\square_n,1}(B_V),
 \qquad Q_0(c)=Q_q(\eta_nH_n)(c)\quad(c\in\Gamma_q^j).
\]
For $\tau\in[0,1]$ and $\sup|B'|<c_3a_0$, where
$c_3=\beta_{\rm s}/(16B^\sharp)$, define
$H=\mathcal H_{\mathbf B^\square_j,1}(\tau Q_0+B')$.
All these objects exist; for fixed $U$ they are holomorphic in
$(A',B')$ on an open neighborhood of the stated domain. Moreover
$\sup|B_V|<6KMA_0/5$ and, on $\square^{\sim3}$,
\begin{equation}\label{I-93-103}
 |H|,|\nabla H|<a_1/2,\qquad
 |\partial e^{i\eta_jH}-I|<a_0\eta_j^2/8,\qquad
 |J^{(j)}(e^{i\eta_jH})|<a_0/16.
\end{equation}
For every $X\subset\square^{\sim2}$ in $\mathcal D_{j-1}$,
\begin{equation}\label{I-93-105}
 (1,H,J^{(j)}(e^{i\eta_jH}))|_X
       \in\mathcal U_j^\circ(X;a_0/4,a_1/2,a_0/4).
\end{equation}
\end{proposition}

\begin{proof}
We give the restriction argument first. Restrict a real regular
critical configuration to a subtower whose initial region lies
inside a later member of the original tower. A variation preserving
the restricted datum extends by the original field outside that
region. The crossing plaquettes are fixed because the boundary
bonds are level-zero determining bonds. Induction in the averaging
order shows that the extended variation preserves every original
datum. It is therefore critical. Near the trivial datum uniqueness
in \cref{I-58-22} identifies the two backgrounds. The same identity
extends to the common complex logarithmic ball. Axial normalization
gives its gauge as the product of contour variables, with
\begin{equation}\label{I-93-104}
 |w^{\pm1}|\le e^{c_w\sup|H|}.
\end{equation}
Each contour logarithm at level $q<j$ is bounded by
$C(L)L^{q-j}\sup|H|$; their sum proves this bound. This proves
the same statement for truncating a tower without restricting
its initial region. Changes of real base datum and small complex
node gauges follow in the same way: the identities hold at real
arguments by uniqueness and extend on their convex logarithmic
domains. In particular the tower of $\mathbf U$ is
$\Phi_{\mathbf B^\square_n,1}(B_V)^{\tilde v}$.

The logarithm estimate in \eqref{eq:octet-comb-datum} gives
$\sup|B_V|\le b_V=(33/20)K_\square MA_0<6KMA_0/5$.
The tower and its truncation have free pillars over
$\square^{\sim4}$. Hence $h_+:=\sup_{\square^{\sim4}}|H_n|
\le B^\sharp b_V$. Averaging gives
\[
 q_{\mathrm{out}}:=C_QLB^\sharp b_V\ge\sup|Q_0|,\qquad
 q_{\mathrm{in}}:=C_QB^\sharp b_VL^{-\nu}
       \ge\sup_{\mathsf P^q(c)\subset\square^{\sim4}}|Q_0(c)|.
\]
Every row in the second supremum has level $j$. From a bond over
$\square^{\sim3}$ the other rows have weighted distance at least
$D=(M/2)L^{\nu-1}$. Apply \eqref{eq:tower-distant-data} to these
rows and \eqref{I-93-61} to the near rows. The floor on $M$ gives
\[
 B^\flat e^{-\delta_{\mathrm{tw}}D/2}q_{\mathrm{out}}
          \le(\beta_{\rm s}/4)A_0L^{-2\nu},\qquad
 N_b(\mathcal H_j(\tau Q_0))\le h_1:=c_hMA_0L^{-\nu}.
\]
Splitting $Q_0$ into the near and far parts also gives, by
\eqref{eq:tower-quadratic-response},
\begin{equation}\label{eq:inner-linearization-error}
 N_b(\mathcal H_j(\tau Q_0)-\tau\mathsf H Q_0)
 \le2B^\flat e^{-\delta_{\mathrm{tw}}D/2}q_{\mathrm{out}}
       +2B^\sharp q_{\mathrm{in}}^2/r_{\mathrm B}
 \le\beta_{\rm s} A_0L^{-2\nu}.
\end{equation}

At $\tau=1$, the restriction identity and the free-pillar gauge
comparison give
\[
\begin{gathered}
 e^{i\eta_j\mathcal H_j(Q_0)}
       =(U_Z)^{G},\qquad
 G=u_j^{-1}g_ju_n^{-1}\tilde v^{-1},\\
 \mathsf k(G)
 \le\exp(2c_uh_1+2(c_u+c_w)h_++K_\square M(1+2\beta_{\rm s})A_1)\\
 \le e^{\beta_{\rm s}/2}\le\frac{1+3\beta_{\rm s}}{1+2\beta_{\rm s}}.
\end{gathered}
\]
Here $U_Z$ is the tower of $\mathbf U$; $g_j$ is the contour
gauge for the truncation. The window hypothesis therefore bounds
the curvature by $(1+3\beta_{\rm s})A_0L^{-2\nu}\eta_j^2$ and the
current by $(1+3\beta_{\rm s})A_0L^{-3\nu}$. Applying
\eqref{I-10-121} and \eqref{I-10-122} at the trivial base, and
using $134h_1^2\le\beta_{\rm s} A_0L^{-2\nu}$, yields
\[
 |\partial\mathcal H_j(Q_0)|<(1+4\beta_{\rm s})A_0L^{-2\nu},\quad
 |\partial^{\mathsf T}\partial\mathcal H_j(Q_0)|
       <(1+3\beta_{\rm s})A_0L^{-3\nu}+\beta_{\rm s} A_0L^{-2\nu}.
\]
Apply \eqref{eq:inner-linearization-error} at $1$, use linearity
of $\mathsf H$, and apply it once more at $\tau$. The two bounds
become, respectively,
$(1+8\beta_{\rm s})A_0L^{-2\nu}$ and
$(1+3\beta_{\rm s})A_0L^{-3\nu}+3\beta_{\rm s} A_0L^{-2\nu}$.
Adding $B'$ changes $N_b$ by less than $\beta_{\rm s} a_0/16$,
and hence by less than the $\beta_{\rm s} a_0/8$ allowed below. Thus
\[
 h:=\max\{|H|,|\nabla H|\}<h_1+\beta_{\rm s} a_0/8,
 \quad |\partial H|<(1+8\beta_{\rm s})A_0L^{-2\nu}+\beta_{\rm s} a_0/4,
\]
and $|\partial^{\mathsf T}\partial H|$ is bounded by
$(1+3\beta_{\rm s})A_0L^{-3\nu}+3\beta_{\rm s} A_0L^{-2\nu}+\beta_{\rm s} a_0/8$.
The first and third assumptions in \eqref{I-93-101} give
$h<a_1/2$. Also
$134h^2\le2\beta_{\rm s} A_0L^{-2\nu}+a_0/500$,
with $11h^2\le2\beta_{\rm s} A_0L^{-2\nu}+a_0/1000$ after the
same radius reduction. The local expansions and
\eqref{I-93-95} now give \eqref{I-93-103}.

Finally apply the restriction identity to every subtower
$(p,Z')$ of $X$ and both tower rules. Its gauge has condition
number at most $e^{c_wa_1}<6/5$. The curvature bound is therefore
at most $(6/5)(a_0/8)L^{-2j}$, smaller than
$(a_0/4)L^{-(j-p)}L^{-2p}$. The current bound at truncation
$n'\le p$ is at most $(6/5)(a_0/16)L^{-3(j-n')}$, smaller than
$(a_0/4)L^{-(j-p)}L^{2(n'-p)}$, since $2j-p-n'\ge0$.
The real factor is trivial, and the potential and current bounds
are those already proved. These are all the core conditions,
so \eqref{I-93-105} follows. The strict finite bounds and the
common logarithmic balls also prove the stated holomorphy.
\end{proof}

\subsection{Replacing an old interaction by an octet interaction}
\label{sec:octet-interactions}

\begin{definition}\label{II-48-4}
A system of radii of height $k$ is
$\mathfrak a=((a_{0,j},a_{1,j}))_{1\le j\le k+1}$. Put
\begin{equation}\label{II-48-5}
 \bar a_i=(1+\beta_{\rm s})a_{i,k+1},\qquad
 \varrho_j=\max_i\bar a_i/a_{i,j},\qquad
 s_{\mathfrak a}=\delta(\mathsf t)/a_{0,k+1}.
\end{equation}
For $a^\star,s^\star>0$ and $c_\natural,c_{\mathrm{sh}}\ge1$,
the system is admissible if
\begin{equation}\label{II-48-7}
\begin{gathered}
 \bar a_0+\bar a_1\le a^\star,\qquad
 c_{\mathrm{sh}}Ma_{0,j}<a_{1,j},\quad
 a_{1,j}^2\le c_\natural a_{0,j}\quad(1\le j\le k+1),\\
 a_{i,k+1}\le(1+\beta_{\rm s})\sqrt{k+2-j}\,a_{i,j},\qquad
 a_{i,j}\le\bar a_i\quad(j\le k),\qquad s_{\mathfrak a}\le s^\star.
\end{gathered}
\end{equation}
The two comparisons indexed by $j\le k$ are void at height zero.
Write $\mathcal U_j(X)=\mathcal U_j^\circ(X;a_{0,j},a_{1,j})$.
\end{definition}

Fix $0<\beta_{\rm s}\le1/5$ and $0<\theta_{\mathrm{oc}}<1$ with
$\tfrac32\theta_{\mathrm{oc}}^2<1$. The octet radii and the
system of height $k-1$ are
\begin{equation}\label{II-48-20}
 b_i^{\mathrm{oc}}=\theta_{\mathrm{oc}}a_{i,k+1},\quad
 a_{i,k}^{\mathrm{oc}}=b_i^{\mathrm{oc}}/(1+\beta_{\rm s}),\quad
 a_{i,j}^{\mathrm{oc}}=\min\{a_{i,j},b_i^{\mathrm{oc}}\}\quad(j<k).
\end{equation}
Its output radii are $b_i^{\mathrm{oc}}$, and it is admissible
with $s^\star$ replaced by $(1+\beta_{\rm s})s^\star/\theta_{\mathrm{oc}}$.
Indeed shape is preserved by a common contraction and by these
minima. For slow growth, if the minimum is $a_{i,j}$ use
$\theta_{\mathrm{oc}}^2(k+2-j)/(k+1-j)<1$; otherwise the output
bound suffices. Thus $1\le\varrho_j^{\mathrm{oc}}\le\varrho_j$ and
$\varrho_j^2L^{j-k}<1$ for $j<k$, since
$\varrho_j^2\le(6/5)^4(k+2-j)$ and $L\ge13$.

We use the canonical maps
$\mathbb H_q=\mathcal H_{\mathbf B_q^\square,1}$.
They are real at real data, commute with adjunction and constant
rotations, and commute with the translations of lattice charts
used above. Indeed transport the constraint, Landau condition
and node normalization in \eqref{I-82-16}. The averaging paths,
covariant differences and action are carried to their counterparts;
the geometric distances and norms are unchanged. The transported
pair solves the same normalized equations. Uniqueness near zero
and holomorphy identify the maps on the row ball. At real data
their exponentials are gauge transforms of the unique regular
critical orbit, by \cref{I-58-22}. Their derivative bound is
\eqref{I-93-47}. These are properties of the constructed maps.

\begin{lemma}\label{II-8-36}
Let $(\Omega_q)_{0\le q\le n}$ and $(\Omega'_q)_{0\le q\le n'}$
be separated sequences with $n'\le n$ and
$\Omega'_q\subset\Omega_q$ for $0\le q\le n'$. A real regular
critical configuration for $(\Omega_q)_{0\le q\le n}$ restricts to
a real regular critical configuration for
$(\Omega'_q)_{0\le q\le n'}$ at the same regularity radius.
This remains true after a real gauge transformation of the
restricted configuration.
\end{lemma}

\begin{proof}
Glue a configuration with the same restricted datum to the original
configuration outside $\Omega'_0$. The gluing preserves the
original datum, by induction on averaging order. At level zero
the fields agree outside $\Omega'_0$ and on its determining rows.
At level $q$, a row with an endpoint in $\Omega'_q$ and none
in $\Omega'_{q+1}$ is a determining row. If neither endpoint
lies in $\Omega'_q$, its entire averaging corridor is outside
that region, where the induction hypothesis applies. Original
determining rows have no endpoint in $\Omega_{q+1}$ and hence
none in $\Omega'_{q+1}$, so their averages are unchanged.
A plaquette crossing $\Omega'_0$ has no corner in $\Omega'_1$,
by separation. Its inner bonds are fixed level-zero determining
rows. Thus the action of the glued field minus the action inside
$\Omega'_0$ is independent of the variation. Differentiation
proves criticality. For a plaquette or complete plaquette star,
its highest layer indices satisfy $\lambda'\le\lambda$.
Restriction multiplies its bare-scale regularity quantity by
$L^{-2(\lambda-\lambda')}$ for curvature and
$L^{-3(\lambda-\lambda')}$ for current. These factors are at
most one. Real gauges preserve these norms, the fibres and the action.
\end{proof}

For an octet of level $k$, put $X'=(\square^{\sim5})^\uparrow$
and form, from $\mathbf U$ on $X'$, the fields
\[
 B_V=i^{-1}\log V,\quad A_{\square,k}=\mathbb H_k(B_V),\quad
 Q_{\square,j}=i^{-1}\log M_{\mathbf B_j^\square}
                         (e^{i\eta_kA_{\square,k}}),\quad
 A_{\square,j}=\mathbb H_j(Q_{\square,j}).
\]
For a local term $E$ on $X\subset\square^{\sim2}$ define
\begin{equation}\label{II-48-33}
 E^\square(P')=E\bigl((1,A_{\square,j},
             J^{(j)}(e^{i\eta_jA_{\square,j}}))|_X\bigr),
 \qquad \mathbf U=\mathbf U(P').
\end{equation}

\begin{proposition}\label{II-49-44}
There are $a^\star_{\mathrm{oc}}>0$, $c^{\mathrm{oc}}_{\mathrm{sh}}$
and $C_{\mathrm{oc}}$, independent of the height and volume,
such that the following holds for an admissible system with
$a^\star\le a^\star_{\mathrm{oc}}$ and
$c_{\mathrm{sh}}\ge c^{\mathrm{oc}}_{\mathrm{sh}}$.
Let $k\ge2$, $1\le j<k$, and $X\in\mathcal D_j$ satisfy
$X\subset\square^{\sim2}$. Then \eqref{II-48-33} is defined on
$\mathcal U_k^\circ(X';b_0^{\mathrm{oc}},b_1^{\mathrm{oc}})$,
and its inner triple belongs to
\begin{equation}\label{II-49-46}
 \mathcal U_j^\circ(X;a_{0,j}^{\mathrm{oc}}/4,
                    a_{1,j}^{\mathrm{oc}}/2,a_{0,j}^{\mathrm{oc}}/4).
\end{equation}
It is a local term there, depends only on $\mathbf U$ on the
inner bonds of $\square^{\sim5}$, and satisfies
\begin{equation}\label{II-49-48}
 |E^\square(P')-E(1,0)|
 \le C_{\mathrm{oc}}\|E\|_X(d_j(X)+2)^2\varrho_j^2(L\xi_j)^4,
 \qquad \xi_j=L^{j-k}.
\end{equation}
Also $\|E^\square\|_{X'}\le\|E\|_X$. The map $E\mapsto E^\square$
is linear and preserves conjugation reality and chart and translation
covariance. If $U^*$ is in an existing regular critical orbit at
level $k$ or $k+1$ and its level-$k$ pair on $X'$ belongs to the
octet core, then its level-$j$ pair on $X$ belongs to
$\mathcal U_j(X)$ and
\begin{equation}\label{II-49-49}
 E^\square(\mathbf p^{(k)}_{X'}(U^*))
                         =E(\mathbf p^{(j)}_X(U^*)).
\end{equation}
The same identity holds for the plain triples whenever they belong
to the indicated cores.
\end{proposition}

\begin{proof}
Take $c^{\mathrm{oc}}_{\mathrm{sh}}=2B_3^2K$ and decrease
$a^\star_{\mathrm{oc}}\le a^\flat$ until
\[
 \varpi:=\tfrac65KM b_0^{\mathrm{oc}}
 <\min\{c_{\rm plain}/(2LB^\sharp),r_{\mathrm B}/(2C_QLB^\sharp),
                                      1/(120MB^\sharp)\}.
\]
Also require $a^\star_{\mathrm{oc}}\le
\min\{a_{\mathrm{av}},\omega_M/2\}$. The real averaging bound
then makes its real and complex branches agree at the curvature
$b_0^{\mathrm{oc}}\eta_k^2$: at each intermediate order $q$
the curvature is at most $2b_0^{\mathrm{oc}}L^{2(q-k)}\le\omega_M$,
so the next complex averaging step agrees with the real step.

Apply \cref{prop:inner-substitution} with
$A_i=b_i^{\mathrm{oc}}/(1+2\beta_{\rm s})$,
$a_i=a_{i,j}^{\mathrm{oc}}$, $n=k$, $\tau=1$, $B'=0$.
The core includes the required window. The sequence
$t_\nu=(6/5)\sqrt{1+\nu}$ is admitted: its three suprema are
attained at $\nu=1$, and at $L=13$, $\beta_{\rm s}=1/5$ the last two
left sides of \eqref{I-93-95} are less than $0.082$ and $0.039$.
The octet radius system verifies \eqref{I-93-101}, giving
\eqref{II-49-46}. Holomorphy holds near every fixed-real-factor
slice. The maps depend only on the pair; a gauge on $X'$ induces
the constant rotation by its value at $y_0$ on the inner triple.
This proves locality and invariance. Evaluation and covariance
give the norm bound, linearity, reality and transport.

For the gain put $\nu=k-j$, $d=d_j(X)$ and
$a^\psi=a_{0,j}/(8K)$. The simple-to-core inclusion gives
\[
 \|A\|_{X,\nabla}<a^\psi,\quad |J|<a_{0,j}
 \quad\Longrightarrow\quad(1,A,J)\in\mathcal U_j(X):
\]
the plaquette bound is less than $4a^\psi\eta_j^2$, and
$K(4a^\psi+(a^\psi)^2)<a_{0,j}$ after the radius reduction.
Thus $\Psi(A,J)=E(1,A,J)$ is holomorphic on this product ball
and bounded by $\|E\|_X$. Its differential at zero vanishes:
constant $SU(2)$ invariance makes each of its linear functionals
on the traceless matrices invariant under conjugation, hence zero.
Cauchy's estimate on radial lines gives
\begin{equation}\label{II-49-57}
 |\Psi(A,J)-\Psi(0,0)|\le8\|E\|_X(N/a^\psi)^2,
 \quad \max\{\|A\|_{X,\nabla},\sup|J|\}\le N.
\end{equation}
For $N\ge a^\psi/2$ use the bound $2\|E\|_X$ instead.

Put $\sigma=B^\sharp\varpi$, $\phi_1=L^{-\nu}\sigma$,
$\phi_2=L^{-2\nu}\sigma$, and
\[
 \epsilon_\nu=C_QB^\sharp
       (B^\sharp L^{-\nu}+B^\flat L
                         e^{-\delta_{\mathrm{tw}}ML^{\nu-1}/2}).
\]
For $\sup|B|<\varpi$, let $f_B=L^{-\nu}\mathbb H_k(B)$.
The response estimate, scaling and \eqref{I-10-122} give
$|f_B|\le\phi_1$, $|\nabla f_B|\le\phi_2$ on the box of
radius $3ML^{\nu-1}+3\eta_j$ around $x_\square$, and
$|J^{(j)}(e^{i\eta_jf_B})|\le2L^{-3\nu}\sigma$ on $X$.
Its averaged datum on the truncated tower is bounded by
$C_QLB^\sharp\sup|B|$ globally and by
$C_QL^{-\nu}B^\sharp\sup|B|$ on rows within weighted distance
$ML^{\nu-1}$ of that smaller box. Splitting the response into
these near and far rows bounds its level-$j$ potential, gradient
and curl-transpose-curl by $\epsilon_\nu\sup|B|$.
Since $e^{-t}\le1/t$,
\[
 \frac{\epsilon_\nu\varpi}{a^\psi}\le C_2\varrho_j\xi_j,
 \qquad C_2=\tfrac{48}{5}K^2MC_QB^\sharp B^\flat
                         (1+2L^2/(\delta_{\mathrm{tw}}M)).
\]
If $C_2\varrho_j\xi_j<1$, both potentials lie in the ball of
$\Psi$, and $8\phi_1\le\epsilon_\nu\varpi<a^\psi$.
For real $B$ their exponentials are gauge equivalent, by
\cref{II-8-36} and uniqueness at depth $j$. Gauge invariance
identifies their values under $E$. The identity theorem extends
this equality to $\sup|B|<\varpi$. At $B=B_V$ it gives
$E^\square(P')=\Psi(f,J_f)$, where
$f=f_{B_V}$ and $J_f=J^{(j)}(e^{i\eta_jf})$.

Choose a cube centre $z\in X$ and coordinates in the box.
For $|\lambda_\mu|<2\phi_1$ put
\[
 g_\lambda(x)=\prod_{\mu=1}^3e^{i(x_\mu-z_\mu)\lambda_\mu},
 \quad T_\mu=\prod_{\kappa>\mu}e^{i(x_\kappa-z_\kappa)\lambda_\kappa},
 \quad \tilde\lambda_\mu=T_\mu^{-1}\lambda_\mu T_\mu,
\]
\[
 A_\lambda(b)=\operatorname{Ad}(g_\lambda(b_-))\mathcal L_b,
 \quad \mathcal L_b=(i\eta_j)^{-1}\log
     (e^{i\eta_jf(b)}e^{-i\eta_j\tilde\lambda_\mu(b_-)}),
 \quad J_\lambda=\operatorname{Ad}(g_\lambda)J_f.
\]
The total exponent norm in $g_\lambda$ is at most $3M\sigma<1/40$,
so its condition number is less than $1.06$. The logarithm has
Lipschitz constant $1.3$ in these variables. Directly,
$g_\lambda(x+\eta_je_\kappa)=g_\lambda(x)e^{i\eta_j\tilde\lambda_\kappa(x)}$
and $|\tilde\lambda_\mu(x+\eta_je_\kappa)-\tilde\lambda_\mu(x)|
\le2\eta_j\tilde l^2e^{2\eta_j\tilde l}$, where
$\tilde l\le2.12\phi_1$. Therefore
\[
 |\nabla_\kappa A_{\lambda,\mu}|
 \le1.06(2.02\tilde l|\mathcal L_{b'}|+1.3\phi_2+2.63\tilde l^2).
\]
It follows that $|A_\lambda|\le5\phi_1$,
$|\nabla A_\lambda|\le2\phi_2$, $|J_\lambda|\le3\phi_1$.
Thus all these pairs remain in the ball of $\Psi$. At real
$\lambda$ the gauge is unitary and $\Psi(A_\lambda,J_\lambda)=\Psi(f,J_f)$;
holomorphy extends this identity to the complex parameter ball.

Set $\lambda_\mu=f_\mu(z)$. A coordinate path gives
$|f_\mu(x)-\lambda_\mu|\le r(x)\phi_2$, where
$r(x)=\sum_\mu|x_\mu-z_\mu|$. Also
$|\tilde\lambda_\mu(x)-\lambda_\mu|
\le2.12\sigma r(x)\phi_2\le0.02r(x)\phi_2$.
The preceding difference estimate then gives
$|A_\lambda(x)|\le1.41r(x)\phi_2$ and
$|\nabla A_\lambda|\le1.5\phi_2$.
On bonds incident to $X$, $r(x)\le3.1M(d+2)$, so
\[
 \max\{\|A_\lambda\|_{X,\nabla},\sup|J_\lambda|\}
       \le N_X=5M(d+2)\phi_2,
 \qquad N_X/a^\psi\le C_3(d+2)\varrho_j\xi_j^2,
 \quad C_3=48K^2M^2B^\sharp.
\]
Equation \eqref{II-49-57} proves the gain with $8C_3^2$.
If $C_2\varrho_j\xi_j\ge1$, use $\varrho_j^2\xi_j<1$ to
obtain $1<C_2^6\varrho_j^2\xi_j^4$; the trivial bound suffices.
Taking $C_{\mathrm{oc}}=\max\{C_2^6,8C_3^2\}$ and enlarging
$\xi_j^4$ to $(L\xi_j)^4$ proves \eqref{II-49-48}.

For the physical identity, \cref{II-8-36} applies to every
truncation of the octet tower inside the global critical field.
The real and complex averages agree. Uniqueness of the regular
critical orbits identifies successively
$e^{i\eta_kA_{\square,k}}$ and $e^{i\eta_jA_{\square,j}}$
with gauge transforms of $U^*|_Z$. The currents transform by
conjugation, so the inverse gauge on $X$ carries the inner pair
to $\mathbf p_X^{(j)}(U^*)$. This proves \eqref{II-49-49}.
For plain triples restriction scales the radii by
$(\xi_j,\xi_j,\xi_j^3)$; since
$\xi_jb_i^{\mathrm{oc}}<a_{i,j}$, they belong to the required
cores and pair dependence gives the same equality.
\end{proof}

\subsection{Substitution of a local term}
\label{sec:one-term-shift}

Use the shift $(\mathbf A_Y,\mathbf K_Y)$ constructed in
\cref{II-51-52}, with constants $(a_{\mathfrak h},
\varepsilon_{\mathfrak h},c_{\mathfrak h})$. Its input at step
$k+1$ is read as $(\mathbf U,L^{-3}\mathbf J)$ in level-$k$
units. This factor will also be needed when we identify its
values on critical backgrounds.
Choose $0<c_{\mathrm{os}}\le1$ so that
\[
 128K^\sharp_{\mathrm{tw}}(1+c_\natural)c_{\mathrm{os}}\le1.
\]
Let $0<a_i^+\le a_{\mathfrak h}$ and
$0<\varepsilon\le\varepsilon_{\mathfrak h}$. For $Y\in\mathcal D_k$
write
\[
 \mathfrak D(Y)=\{(P,\mathcal A):
 P\in\mathcal U^\circ_{k+1}(Y;a_0^+,a_1^+),\quad
 \sup|\mathcal A|<\varepsilon\}.
\]
The form $\mathcal A$ is specified on the nonpivot level-$k$
bonds inside $Y$. A term with radii is a local term $f$ on
$X\in\mathcal D_j$, $1\le j\le k$, satisfying
\begin{equation}\label{II-48-14}
 0<b_0\le b_1,\quad b_1^2\le c_\natural b_0,\quad
 b_0+b_1\le a^\sharp,\qquad
 a_i^+\le2\sqrt{1+\nu}\,b_i,\quad \nu=k+1-j.
\end{equation}
Its amplitude radius is
\begin{equation}\label{II-48-15}
 \varrho_{\mathsf{tm}}=
       \frac{c_{\mathrm{os}}b_0L^{k-j}}{c_{\mathfrak h}\varepsilon}.
\end{equation}
For $X^\uparrow\subset Y$, set $\theta=L^{-\nu}$,
$\vartheta=L^{-(k-j)}$, and
\begin{equation}\label{II-50-42}
 \Phi_Y[f](s,\lambda;P,\mathcal A)
 =f\bigl(\rho_X^{k+1\to j}P\dotplus
       (\lambda\vartheta\mathbf A_Y(s),
        \lambda\vartheta^3\mathbf K_Y(s))|_X\bigr).
\end{equation}
Thus the extra potential scales by $\vartheta$, its first
difference by $\vartheta^2$, and its current by $\vartheta^3$.
Here $s$ ranges over the decoupling polydisc
$|s_\square|<e^{\bar\kappa}$.

\begin{proposition}\label{II-50-61}
For $(P,\mathcal A)\in\mathfrak D(Y)$ and
$|\lambda|<\varrho_{\mathsf{tm}}$, the substituted triple in
\eqref{II-50-42} lies in $\mathcal U_j^\circ(X;b_0,b_1)$.
The function $\Phi_Y[f]$ is bounded by $\|f\|_X$ and is
holomorphic in $(s,\lambda)$ and locally holomorphic on each
fixed-real-$U$ slice in the remaining variables. At $\lambda=0$,
$s=0$, or $\mathcal A=0$, its value is $f(\rho_XP)$.
It is local for the roots given by the cubes of $X^\uparrow$
and consistent under zero extension to larger domains.

Suppose $\varrho_{\mathsf{tm}}>1$. With the starred bracket
of \cref{II-18-108}, put
\begin{equation}\label{II-50-43}
 \mathbb V[f](Y;\mathfrak u)
   =[\Phi_Y[f](\cdot,1;\mathfrak u)]^*
      [\mathsf P_Y\setminus\mathsf P_{X^\uparrow}
                         \mid\mathsf P_{X^\uparrow}].
\end{equation}
It vanishes at $\mathcal A=0$, is linear in $f$, kills constant
terms, and satisfies, with $m=\#_{\pi_k}(Y\setminus X^\uparrow)$,
\begin{equation}\label{eq:one-term-potential-bound}
 |\mathbb V[f](Y;\mathfrak u)|
 \le2\varrho_{\mathsf{tm}}^{-1}\|f\|_X
                      (e^{\bar\kappa}-1)^{-m}
 \le2\varrho_{\mathsf{tm}}^{-1}\|f\|_Xe^{-(\bar\kappa-1)m}.
\end{equation}
Both functions inherit pair dependence, gauge invariance,
reality, conjugation, translations and local charts from $f$
and the shift. For chart transport both the original and
transported points must belong to their specified domains.

For a real block field $W$, assume its averaged background is
defined, real and admissible, with invertible pivots; assume
the fibre chart is defined at $\mathcal A$, its level-$k$
critical background exists, and
$\mathfrak u_W=(P_W,\mathcal A)\in\mathfrak D(T_\ell)$.
Then $f$ is defined on both critical orbits, is constant on
each, and
\begin{equation}\label{eq:one-term-background-sum}
 f(U_k(V_W[\mathcal A]))-f(U_{k+1}(W))
  =\sum_{Y\supset X^\uparrow}\mathbb V[f](Y;\mathfrak u_W|_Y).
\end{equation}
\end{proposition}

\begin{proof}
Since $4\sqrt{1+\nu}\le L^\nu$, restriction puts $\rho_XP$
in the half-radius core, with current radius
$c_0=\theta^3a_0^+\le b_0/338$. The shift bounds and
\eqref{II-48-15} give
\[
 |\lambda\vartheta\mathbf A|,
 |\lambda\vartheta^2\nabla_{\mathbf U}\mathbf A|
       <c_{\mathrm{os}}b_0\le b_2,
 \qquad |\lambda\vartheta^3\mathbf K|
       <c_{\mathrm{os}}b_0\le b_0/256.
\]
Thus the total current has norm less than $b_0$.
\Cref{II-16-69} proves core membership.

For completeness, local holomorphy does not require the core
itself to be open in all its variables. Fix a real $U$ and a
point in its slice. The shift has a holomorphic extension to a
neighborhood of that point. Its substituted configuration is a
product of matrix exponentials and fixed matrices. The logarithm
relative to its fixed real factor exists near the point, and
provides a holomorphic pair of local coordinates for the slice
on which $f$ is holomorphic. Their composition proves the claim.
The zero values follow from the zero values of the shift.
Every bond incident to $X$ has an endpoint in a root cube.
A component of the decoupling support disjoint from all roots
therefore cannot affect the substituted triple on $X$.
This proves root locality. Restriction and transport commute
with the bond formulas, including reversal, so consistency and
the asserted symmetries follow from those of the shift and $f$.

Subtract $\Phi_Y[f](s,0)$ before taking the starred bracket.
Amplitude Cauchy gives the bound
$2\varrho_{\mathsf{tm}}^{-1}\|f\|_X$ at $\lambda=1$;
the bracket bound then gives \eqref{eq:one-term-potential-bound}.
In particular the finite vertex formula is
\[
 \mathbb V[f](Y)=
 \sum_{X^\uparrow\subset W\subset Y}
       (-1)^{\#Y-\#W}f(T_W)
       -\mathbf1_{Y=X^\uparrow}f(T_\emptyset),
 \quad T_W=\rho_XP\dotplus\mathfrak h_Y(\mathbf1_W)|_X,
\]
where the factors $\vartheta,\vartheta^3$ in $T_W$ are those
of \eqref{II-50-42}. This also proves local holomorphy of the
potential and its inheritance assertions.

At full support, the identification of the actual shift gives
the level-$k$ background on the fibre. Its current at level $j$
is
\[
 \vartheta^3\bigl(L^{-3}J^{(k+1)}(U_+)+\mathbf K\bigr)
       =J^{(j)}(U_k(V_W[\mathcal A])).
\]
At zero support the pair is the restriction of $U_+$.
Gauge invariance gives the asserted orbit values. Finally sum
the finite differences over all subsets of cubes. Components
disjoint from the roots cancel in pairs, by root locality;
the remaining support together with the connected root set
is connected. Zero-extension consistency identifies its
coefficient with the corresponding $\mathbb V[f](Y)$.
The total is $\Phi(\mathbf1,1)-\Phi(0,1)$, proving
\eqref{eq:one-term-background-sum}.
\end{proof}

\subsection{The fourth potential}
\label{sec:fourth-potential}

Let $\mathfrak f_\circ$ be an array of pointed interactions
$G^{(j)}(X,P,z)$, counterterms $e_j$, and unpointed interactions
$H^{(j)}(X,P)$, $1\le j\le k$. Set $H^{(j)}=0$ off its slot.
Write
\[
 g_j=\sup_{X,P,z}e^{\kappa d_j(X)}|G^{(j)}(X,P,z)|,
 \quad r_j=\sup_{X,P}e^{\kappa d_j(X)}|H^{(j)}(X,P)|,
 \quad g_j^{\mathrm{ct}}=\tfrac32a_{0,j}^2|e_j|,
\]
and
\begin{equation}\label{II-52-9}
 N_{\mathfrak a,\omega}(\mathfrak f_\circ)
   =\sum_{j=1}^k\omega_{k+1-j}(g_j+g_j^{\mathrm{ct}}+r_j).
\end{equation}
All suprema are over the corresponding cores. The action is
\begin{equation}\label{II-52-7}
\begin{split}
 S[\mathfrak f_\circ](U)
 ={}&\sum_{j=1}^k\left\{
   \sum_z\sum_{X\ni z}\bigl(G^{(j)}(X,U,z)-G^{(j)}(X,1,z)\bigr)
                      -e_jA^{\eta_j}(U)\right\}\\
 &+\sum_{j=1}^k\sum_X
              \bigl(H^{(j)}(X,U)-H^{(j)}(X,1)\bigr).
\end{split}
\end{equation}
Use the tent action to put
$E_{X,z}=G^{(j)}(X,\cdot,z)-e_j\mathsf A^{\mathsf t}(X,\cdot,z)$.
It obeys
\[
 \|E_{X,z}\|_X\le g_je^{-\kappa d_j(X)}
                    +g_j^{\mathrm{ct}}\mathbf1_{X=X_z},
 \qquad \#_{\pi_j}X_z=27,\quad d_j(X_z)\le7.
\]

For $j<k$ and a level-$j$ site $z$, let $\square(z)$ be the
octet centered at the upper corner of the unique
level-$(k-1)$ partition cube containing $z$; put
$X'(z)=(\square(z)^{\sim5})^\uparrow$. Thus $z\in X'(z)$,
$\#_{\pi_k}X'(z)\le8$ and $d_k(X'(z))=0$.
Call $X$ narrow at $z$ if $j<k$ and
$X\subset\square(z)^{\sim2}$, and wide at $z$ otherwise.
A domain is narrow if it is nonwrapping and narrow at the
lexicographically first site $z_X$ in a near lift. This site
is independent of the near lift after projection, since two
such lifts differ by a period. All other domains are wide.
At level $k$ every domain is wide.

Replace the narrow terms by their octet evaluations. More
precisely, form the following four families of local terms:
\begin{equation}\label{II-52-79}
\begin{array}{c|c|c}
 \text{index and condition}&\text{level, domain}&f_{\mathfrak i}\\ \hline
 (j,z,X),\ X\text{ wide at }z&j,X&E_{X,z}\\
 (j,z),\ j<k&k,X'(z)&
  \displaystyle\sum_{z\in X\subset\square(z)^{\sim2}}
              (E_{X,z}^{\square(z)}-E_{X,z}(1))\\
 (j,X),\ X\text{ wide}&j,X&H^{(j)}(X,\cdot)\\
 (j,X),\ X\text{ narrow}&k,X'(z_X)&
        H^{(j)}(X,\cdot)^{\square(z_X)}-H^{(j)}(X,1).
\end{array}
\end{equation}
The first and third families use $(a_{0,j},a_{1,j})$;
the other two use $(b_0^{\mathrm{oc}},b_1^{\mathrm{oc}})$.
For $k=1$ the second and fourth families are empty.
Define
\begin{equation}\label{II-52-78}
 \mathbb V'[\mathfrak f_\circ](Y;\mathfrak u)
    =\sum_{\mathfrak i:X_{\mathfrak i}^\uparrow\subset Y}
                            \mathbb V[f_{\mathfrak i}](Y;\mathfrak u).
\end{equation}
At height zero this rule is zero.

Here are the quantitative choices used in the estimate.
Let $\kappa_{\mathcal D}=\log9216$ and
$C_{\mathrm E,2}=2048\sum_{n\ge0}(n+2)^2e^{-n}$.
Fix $0<\delta_{\mathrm w}<1/2$ and require
\begin{equation}\label{II-52-64}
 \delta_{\mathrm w}\kappa>\kappa_{\mathcal D}+2,
 \qquad\bar\kappa\ge\kappa+1,
 \qquad1+\beta_{\rm s}\le2\sqrt2\theta_{\mathrm{oc}}.
\end{equation}
Put
\begin{equation}\label{II-52-58}
\begin{gathered}
 C_\Phi=3C_{\mathcal A}c_{\mathfrak h}/c_{\mathrm{os}},
 \qquad C_\Sigma=9C_{\mathrm E,2},\qquad
 C_\Lambda=2C_\Sigma C',\qquad K_\Lambda=C_\Lambda C_\Phi,\\
 C'=\max\{27M^3e^{7\kappa},
       8M^3L^4(1+C_{\mathrm{oc}}\max\{C_{\mathrm E,2},81\})\}.
\end{gathered}
\end{equation}
For $k=1$ use $C'=27M^3e^{7\kappa}$. Select ceilings satisfying
\begin{equation}\label{II-52-65}
 a^\star\le\min\{a_{\mathfrak h},a^\sharp,a^\star_{\mathrm{oc}}\},
 \qquad c_{\mathrm{sh}}\ge c^{\mathrm{oc}}_{\mathrm{sh}},\qquad
             2C_\Phi s^\star\le1.
\end{equation}
At the parameter $\mathsf t$, assume $\delta(\mathsf t)\le1$
and $C_{\mathcal A}\delta(\mathsf t)\le\varepsilon_{\mathfrak h}$.
The actual localization domain is
\[
 \mathfrak D_k(Y)=\{(P,\mathcal A):
 P\in\mathcal U^\circ_{k+1}(Y;\bar a_0,\bar a_1),\quad
              \sup|\mathcal A|<C_{\mathcal A}\delta(\mathsf t)\}.
\]
In particular $\mathcal A=gB$ lies here on the open cutoff
polydisc $|B|<C_{\mathcal A}\mathfrak b(\mathsf t)$.

\begin{proposition}\label{II-52-89}
Assume the tower, shift and octet hypotheses above and let
$\mathfrak a$ satisfy \eqref{II-48-7} with the ceilings
\eqref{II-52-65}. The function \eqref{II-52-78} is defined on
$\mathfrak D_k(Y)$ and is locally holomorphic on every
fixed-real-$U$ slice. In particular it is holomorphic in $B$
throughout the cutoff polydisc. It is linear and local in the
array, vanishes at $\mathcal A=0$, depends on the input pair,
and is gauge invariant. It preserves holomorphic parameters,
reality, conjugation, translations and local charts. Constant
changes of an interaction do not change it, nor does adding a
zero slot. Restriction from a larger admissible radii system
commutes with the construction.

If positive lag weights satisfy
\begin{equation}\label{II-52-90}
 \varrho_j^2L^{-(k-j)}\le c_\omega\omega_{k+1-j}
                            \qquad(1\le j\le k),
\end{equation}
then
\begin{equation}\label{II-52-91}
 |\mathbb V'[\mathfrak f_\circ](Y;\mathfrak u)|
 \le K_\Lambda c_\omega s_{\mathfrak a}
       N_{\mathfrak a,\omega}(\mathfrak f_\circ)
                  e^{-(1-2\delta_{\mathrm w})\kappa d_k(Y)}.
\end{equation}
The constants depend only on the fixed geometric, decay and
shift constants displayed in \eqref{II-52-58}, and on
$C_{\mathrm{oc}}$ when $k\ge2$; thus they are uniform in
the lattice, the radii, the array and $\mathsf t$.

For a real $W$, assume its averaged background is defined,
real and admissible with invertible pivots, its fibre chart
is defined throughout the real cutoff, all level-$k$ backgrounds
on the real cutoff exist, and
$P_W\in\mathcal U^\circ_{k+1}(T_\ell;\bar a_0,\bar a_1)$.
Then the action is defined and constant on both critical orbits,
and for every real field below the cutoff,
\begin{equation}\label{II-52-92}
 \sum_Y\mathbb V'[\mathfrak f_\circ](Y;(P_W,gB)|_Y)
   =S[\mathfrak f_\circ](U_k(V_W(B)))
                         -S[\mathfrak f_\circ](U_{k+1}(W)).
\end{equation}
\end{proposition}

\begin{proof}
First check the amplitude radius, which is the source of the
small factor $s_{\mathfrak a}$. By admissibility,
$\bar a_i\le(1+\beta_{\rm s})^2\sqrt{k+2-j}\,a_{i,j}
\le2\sqrt{k+2-j}\,a_{i,j}$. For the octet terms, whose
level is $k$, the corresponding inequality is
$\bar a_i/b_i^{\mathrm{oc}}=(1+\beta_{\rm s})/\theta_{\mathrm{oc}}
\le2\sqrt2$. The radius and shape bounds follow from
\eqref{II-52-65} and \eqref{II-48-20}. Thus every term in
\eqref{II-52-79} satisfies \eqref{II-48-14}. Moreover
\[
 \varrho_{\mathsf{tm}}C_\Phi s_{\mathfrak a}
                  =3b_0L^{k-j}/a_{0,k+1}>1.
\]
For old radii the right side is at least $5/(2\sqrt2)$,
by slow growth and $\beta_{\rm s}\le1/5$; for octet radii it equals
$3\theta_{\mathrm{oc}}>3/(2\sqrt2)$. Hence every radius is
at least $(C_\Phi s_{\mathfrak a})^{-1}\ge2$.
The one-term proposition now proves definition, holomorphy
and the inheritance properties. All sums are finite. In
particular, parameter holomorphy is preserved by evaluation
at the fixed substituted points. The index sets, substituted
points and finite vertex differences are independent of the
chosen radii, which proves restriction compatibility.

We give the summation estimate. Write
$D_m=\max\{2L^{m-1}-2,0\}$. A domain wide at a site $z$
contains a point outside $\square(z)^{\sim2}$, whose distance
from $z$ is greater than $2M L^{k-1-j}$ in level-$j$ units.
The diameter bound $|x-z|_\infty\le M(d_j(X)+2)$ gives
$d_j(X)>2L^{k-j-1}-2$. If $X$ wraps, its diameter bound
and the period at least $8LM$ in level-$k$ units give
$d_j(X)\ge2L^{k-j+1}-2$. Consequently each term at level
$j$ has $d_j(X)\ge D_{k-j}$; octet terms are at level $k$.

Set $b_{j,X}=\sum_{\mathfrak i:j_{\mathfrak i}=j,
X_{\mathfrak i}=X}\|f_{\mathfrak i}\|_X$. We claim
\[
 b_{j,X}\le B_j(d_j(X)+1)e^{-\kappa d_j(X)},
 \qquad B_j=8M^3g_j+r_j\quad(j<k),
\]
and
\begin{align*}
 B_k={}&8M^3g_k+27M^3e^{7\kappa}g_k^{\mathrm{ct}}+r_k\\
 &+\sum_{j<k}8M^3L^4C_{\mathrm{oc}}\varrho_j^2L^{-(k-j)}
       (C_{\mathrm E,2}g_j+81g_j^{\mathrm{ct}}+C_{\mathrm E,2}r_j).
\end{align*}
For wide pointed terms count at most $8M^3(d_j(X)+1)$
sites; at $j<k$ the tent vanishes because $X_z$ is narrow
at $z$. At $j=k$ its at most $27M^3$ sites give the displayed
$e^{7\kappa}$ term. The unpointed term occurs once.

For the narrow terms use \eqref{II-49-48}. The rooted-domain
count gives, for $q>\kappa_{\mathcal D}+2$,
\begin{equation}\label{eq:fourth-rooted-tail}
 \sum_{X\ni z,\ d_j(X)\ge D}(d_j(X)+2)^2e^{-q d_j(X)}
                             \le C_{\mathrm E,2}e^{-D}.
\end{equation}
Indeed at most $2048\,9216^n$ rooted domains have $d_j(X)=n$;
multiply by $(n+2)^2e^{-qn}$ and use
$e^{-2n}\le e^{-D}e^{-n}$ for $n\ge D$.
For each pointed octet term this bounds the sum of norms by
\[
 C_{\mathrm{oc}}\varrho_j^2L^{4-4(k-j)}
                    (C_{\mathrm E,2}g_j+81g_j^{\mathrm{ct}}).
\]
The unpointed terms obey the same bound with
$C_{\mathrm E,2}r_j$. A fixed target $X'(z)$ contains at
most $8M^3L^{3(k-j)}$ possible level-$j$ sites. Multiplying
these two factors gives $L^{4-(k-j)}$, as in $B_k$.
These terms have $d_k(X'(z))=0$, completing the claim.

Fix $j$ and put $d=d_j(X)$,
$m=\#_{\pi_k}(Y\setminus X^\uparrow)$ and
$N_Y=\#_{\pi_k}Y$. Parent trees and attachment of cubes give
$d_k(Y)\le d+m$ and $N_Y\le8(d+1)+m$. Therefore
\begin{align*}
 b_{j,X}e^{-(\bar\kappa-1)m}
 &\le B_j e^{-(1-\delta_{\mathrm w})\kappa d_k(Y)}
       (d+1)e^{-\delta_{\mathrm w}\kappa d}
                         e^{-\delta_{\mathrm w}\kappa m}\\
 &\le\frac{9B_j}{N_Y}e^{-(1-\delta_{\mathrm w})\kappa d_k(Y)}
                      (d+2)^2e^{-\delta_{\mathrm w}\kappa d}.
\end{align*}
For the second inequality insert
$1\le(8(d+1)+m)/N_Y$ and use $me^{-m}\le1$.
There are $L^{3(k-j)}N_Y$ possible level-$j$ cube centres
in $Y$. Sum first over domains containing such a centre,
use \eqref{eq:fourth-rooted-tail} with $D=D_{k-j}$, and
cancel $N_Y$. It follows that
\begin{equation}\label{II-52-95}
 \sum_{j,X:X^\uparrow\subset Y}
        b_{j,X}e^{-(\bar\kappa-1)m}
 \le C_\Sigma\sum_j\varpi_{k-j}B_j
                  e^{-(1-2\delta_{\mathrm w})\kappa d_k(Y)},
 \qquad \varpi_m=L^{3m}e^{-D_m}.
\end{equation}
Here $\varpi_0=1$ and $\varpi_m\le L^{4-m}$, since
$D_m\ge2(m-1)(L-1)\ge4(m-1)\log L$ for $m\ge2$;
the case $m=1$ is equality. Also $\varrho_j\ge1$.
Using \eqref{II-52-90} in $B_j$ and $B_k$ now gives
\[
 \sum_j\varpi_{k-j}B_j
      \le c_\omega C'N_{\mathfrak a,\omega}(\mathfrak f_\circ).
\]
Multiply \eqref{II-52-95} by $2C_\Phi s_{\mathfrak a}$,
using the one-term bound and the amplitude-radius estimate.
This is \eqref{II-52-91}.

Finally apply \eqref{eq:one-term-background-sum} to every
term of \eqref{II-52-79}. For a narrow term,
\eqref{II-49-49} identifies its octet evaluation with its
value on the actual critical background, on both sides of
the step. Each old interaction occurs exactly once, as a wide
term or in an octet sum. The subtracted constants cancel.
The tents sum to $A^{\eta_j}$, so the remaining counterterms
are exactly those in \eqref{II-52-7}. This proves
\eqref{II-52-92} and the well-definedness of its orbit values.
\end{proof}

The two choices used in the iteration can now be stated without
absorbing a change of weights into the smallness hypothesis:
\begin{equation}\label{eq:fourth-weight-cases}
\begin{array}{c|c|c}
 \text{radii}&\omega_n&K\\ \hline
 \text{admissible}&(n+1)L^{-(n-1)}&(1+\beta_{\rm s})^4K_\Lambda\\
 \varrho_j=1\ (1\le j\le k)&L^{-(n-1)}&K_\Lambda.
\end{array}
\end{equation}
Indeed slow growth gives
$\varrho_j^2\le(1+\beta_{\rm s})^4(k+2-j)$, and
$\sum_{n\ge1}(n+1)L^{-(n-1)}
=(2-L^{-1})/(1-L^{-1})^2\le3$.
For any one-step rate $\lambda_V\le(1-2\delta_{\mathrm w})\kappa$,
use its normalization
\begin{equation}\label{eq:fourth-ceiling}
 \nu(\mathbb V')=(Ks_{\mathfrak a})^{-1}
    \sup_{Y,\mathfrak u}e^{\lambda_Vd_k(Y)}
                         |\mathbb V'(Y;\mathfrak u)|,
 \qquad \check s(K)=\min\{s^\star,w_*/(4KN_*)\}.
\end{equation}
Then $\nu(\mathbb V'[\mathfrak f_\circ])
\le N_{\mathfrak a,\omega}(\mathfrak f_\circ)$, and
$s_{\mathfrak a}\le\check s(K)$, $N_{\mathfrak a,\omega}\le N_*$
give the required fourth-potential reserve
$4KN_*s_{\mathfrak a}\le w_*$.


\section{One step of the local fluctuation expansion}
\label{sec:one-step-local-expansion}

We combine the Gaussian construction with the interaction
potentials and the change in the earlier local terms. The
argument in this section uses their quantitative properties
as hypotheses. The subsequent application verifies these
properties for the actual fields and identifies the resulting
integral with the renormalization step.

Fix \(0\le k<K\), \(m'<m+K-k-1\), and
\(M\ge M_1\ge R_1\mathfrak M\ge L^2\).
Let \(\mathsf t>1\), \(g=\mathsf t^{-1/2}\), and
\(\tau=x/\mathsf t\). The functions
\(\varepsilon(\mathsf t),\delta(\mathsf t)\) are the
running thresholds; write
\(\mathfrak b=\mathsf t^{1/2}\delta(\mathsf t)\).
A system of positive radii is
\(\mathfrak a=((a_{0,j},a_{1,j}))_{1\le j\le k+1}\).
For \(0<\beta_{\rm s}\le1/5\), put
\[
 \bar a_i=(1+\beta_{\rm s})a_{i,k+1},\qquad
 s_{\mathfrak a}=\frac{\delta(\mathsf t)}{a_{0,k+1}},\qquad
 \mathcal U^\circ(Y)
       =\mathcal U^\circ_{k+1}(Y;\bar a_0,\bar a_1).
\]
This core is defined for \(Y\in\mathcal D_k\), including
every \(X\in\mathcal D_{k+1}\). Require its slices
\(\mathcal O_U(Y)\) to be open. On a slice, analyticity
then means holomorphy in \((A',\mathbf J)\).

\subsection{Input bounds and constants}

\begin{definition}\label{II-53-12}
The local input consists of the following data, with
constants independent of the domain and the triple.
\begin{enumerate}
\item For every \(X\in\mathcal D_{k+1}\) and
\(P\in\mathcal U^\circ(X)\), a Gaussian structure
\((C,\Gamma)\) on \(X\) with constants
\[
 \bar\kappa\ge1,\quad0<\gamma_{\rm r}\le1,\quad
 c_{\rm r}\ge1,\quad
 \hat\varepsilon_0=\tfrac12u_{\rm G}(LM)^{-3},\quad u_{\rm G}>0.
\]
It is a function of the pair, jointly holomorphic in
the slice and decoupling parameters on a neighborhood
of the closed decoupling polydisc, and gauge covariant
on nonpivot coordinates. It is consistent with
restriction of a whole-torus triple to \(X\).
On the whole torus it is complete, with covariance \(C_*(P)\).
\item For each fine cube \(\square\), a symmetric
gauge-fixing operator \(\widehat{\mathsf G}_\square(P)\)
on its inner nonpivot bonds, with row and column sums
at most \(c_{\mathsf G}\ge1\). It depends only on the
pair, is holomorphic on each slice, and transforms by
the bondwise adjoint maps under a gauge transformation.
Its restrictions agree when the ambient domain changes.
\item Three potentials \(\mathbb V^{\rm J},\mathbb V^{\rm G},
\mathbb V^{\rm A}\) on
\(\mathcal U^\circ(Y)\times\mathbb D^{\rm f}(Y)\).
They are jointly holomorphic on every slice and field
polydisc, are functions of the pair, and satisfy
\(\mathbb V^o(Y;P^u,\operatorname{Ad}(u)B)
=\mathbb V^o(Y;P,B)\).
The J potential vanishes at zero and the G and A
potentials vanish through degree two. Uniformly in \(P\),
\begin{equation}\label{II-53-14}
 \|\mathbb V^o\|^{\rm f}_{\star,\lambda_V}
                         \le K_P\delta(\mathsf t),\qquad K_P\ge1.
\end{equation}
\item A fourth potential \(\mathbb V'\), with the same
analyticity, pair dependence and gauge invariance,
vanishing at zero. For \(K_\Lambda\ge1\), define
\begin{equation}\label{II-53-9}
 \nu(\mathbb V')=(K_\Lambda s_{\mathfrak a})^{-1}
       \sup_{Y,P,B}e^{\lambda_Vd_k(Y)}|\mathbb V'(Y;P,B)|,
\end{equation}
where \(P\in\mathcal U^\circ(Y)\) and
\(B\in\mathbb D^{\rm f}(Y)\). Require \(\nu(\mathbb V')\le N_*\).
\end{enumerate}
Families with additional complex parameters are jointly
holomorphic on these domains and satisfy the same uniform
bounds. Translation covariance and conjugation identities
are imposed when their corresponding conclusions are used.
\end{definition}

All restrictions of triples in this section are at level
\(k+1\). The actual Gaussian construction reads such a
triple at level \(k\) by \eqref{II-38-12}; in particular
its current coordinate is \(L^{-3}\mathbf J\).

Use the expansion constants \eqref{II-19-46} with
\begin{equation}\label{II-53-33}
 \hat\varepsilon_{\rm S}=(u_{\rm G}+c_{\rm r})(LM)^{-3},\qquad
 \gamma_2=2(LM)^{-3},\qquad
 \lambda_{\rm out}=(1+4\beta_{\rm s})\kappa.
\end{equation}
Then \((LM)^3u=3+c_{\rm r}+u_{\rm G}=:u^*_{\rm S}\).
In particular \(c_{\rm v},c_{\rm w}\) are independent
of \(M\). Define
\begin{equation}\label{II-53-32}
 K_{\mathsf G}=\tfrac12c_{\mathsf G}M^3,\qquad
 \sigma_* =\min\{\tfrac12(LM)^{-3},w_*/(2K_{\mathsf G})\},
\end{equation}
and
\begin{align}
 \mathfrak b_*^2
   &=\frac2{\gamma_2}\log\!\left(
            2\mathsf C_{\rm ex}((L+2)M)^3e^{4\lambda_1}\right),
                                            \label{eq:step-threshold-floor}\\
 K_{\mathbf M}
   &=3L^3\max\{K_\Pi K_{\mathsf G},\ K^\Pi_Pe^{4\lambda_1},
                          \ 32K_\Pi K_P,\ 81K_\Pi K_\Lambda\}.
                                            \label{II-53-34}
\end{align}
Take \(\mathsf C_{\mathcal A}\ge3\), \(N_*\ge1\),
\(0<a^*\le1\), and \(s^*>0\), and require
\begin{equation}\label{eq:step-numerical-conditions}
\begin{gathered}
 (LM)^3\ge\mathsf C_{\rm ex}u^*_{\rm S}
                         \max\{\gamma_{\rm r}^{-1},c_{\rm r}^2\},
 \qquad\bar\kappa\ge4\lambda_1+1,\\
 8K_Ps^*\le w_*,\qquad4K_\Lambda N_*s^*\le w_*,\qquad
 \bar a_0+\bar a_1\le a^*,\qquad s_{\mathfrak a}\le s^*,\\
 \varepsilon(\mathsf t),\delta(\mathsf t)\le1,
 \qquad\mathfrak b(\mathsf t)\ge\mathfrak b_*.
\end{gathered}
\end{equation}
The lower bound on \(\mathsf C_{\rm ex}\) remains the one
in the Gaussian cluster expansion. These conditions imply
all its smallness assumptions at the present threshold and
at threshold infinity.

\subsection{The two local contributions}

For a fine cube define
\[
 \hat{\mathsf g}_b(\square;P,B)
       =\tfrac12\langle B(b),(\widehat{\mathsf G}_\square(P)B)(b)\rangle,
 \qquad
 \mathbb V^{\rm q}(\square;P,B)
       =-(1-\tau)\sum_b\hat{\mathsf g}_b(\square;P,B).
\]
Set these families equal to zero on other fine domains,
and set the bond insertion to zero when the bond is not
inner to the cube. Use the scaled structure
\begin{equation}\label{II-53-18}
 (C_\tau,\Gamma_\tau)=(\tau^{-1}C,\tau^{1/2}\Gamma),
\end{equation}
with the principal root. Construct the interpolation of
\(\mathbb V^{\rm J},\mathbb V^{\rm G},\mathbb V^{\rm A},
\mathbb V'\) by \eqref{II-20-95}.
For each \(X,P,b\), with the Gaussian data on \(R=X\)
and the potentials evaluated at restrictions of \(P\), put
\begin{equation}\label{II-20-115}
\begin{split}
 \mathbf G_\tau(X,b;P)
   &=-(1-\tau)\int_0^1
              E'[\{\sigma\mathbb V^{\rm q}\};\hat{\mathsf g}_b](X)
                                                    \,d\sigma,\\
 E^\chi(X,b;P)&=E^\chi[\{\mathbb V^{\rm q}\}](X,b),\\
 E^\bullet[\mathbb V'](X,b;P)
   &=\int_0^1E'[\{\mathbb V^{\rm q},\mathbb V_v^{\rm int},
                         \mathbb V'_v\};
                \dot{\mathbb V}^{\rm int}_{v,b}
                             +\dot{\mathbb V}'_{v,b}](X)\,dv.
\end{split}
\end{equation}
The first integral has threshold infinity; the other
two terms have threshold \(\mathfrak b(\mathsf t)\).
The G term and the cutoff term do not depend on
\(\mathbb V'\). Distribute these terms to sites by
\begin{equation}\label{II-53-23}
 \mathbf G_\tau(X,P,z)=\sum_bh_z(b)\mathbf G_\tau(X,b;P),\qquad
 \mathbf F[\mathbb V'](X,P,z)
       =\sum_bh_z(b)(E^\chi+E^\bullet[\mathbb V'])(X,b;P).
\end{equation}

\begin{theorem}[Local one-step expansion]\label{II-53-39}
Assume the local input of Definition~\ref{II-53-12} and
\eqref{eq:step-numerical-conditions}. For
\(x\in D(\mathsf t,\sigma_*\mathsf t)\), all terms in
\eqref{II-20-115} are defined. Their bond terms vanish
unless \(b^\sharp\subset X\), and their site terms are
pointed families at output radii \((\bar a_0,\bar a_1)\).
Writing \(d=d_{k+1}(X)\), one has
\begin{align}
 |\mathbf G_\tau(X,P,z)|
     &\le K_{\mathbf M}|1-\tau|e^{-\lambda_{\rm out}d},
                                           \label{eq:step-G-bound}\\
 |\mathbf F[\mathbb V'](X,P,z)|
     &\le K_{\mathbf M}\bigl(\delta(\mathsf t)
                     +s_{\mathfrak a}\nu(\mathbb V')
                     +e^{-\gamma_2\mathfrak b^2/2}\bigr)
                                      e^{-\lambda_{\rm out}d}.
                                           \label{eq:step-F-bound}
\end{align}
The terms are jointly holomorphic in \(x\), the slice
variables, and additional holomorphic input parameters
under uniform bounds. They inherit translation covariance
and conjugation. Real input at real \(x\) gives real output.

For two admitted fourth potentials,
\begin{equation}\label{eq:step-fourth-lipschitz}
 |\mathbf F[\mathbb V']-\mathbf F[\mathbb V'']|(X,P,z)
 \le K_{\mathbf M}s_{\mathfrak a}\nu(\mathbb V'-\mathbb V'')
                                      e^{-\lambda_{\rm out}d}.
\end{equation}
On the half-radius \(x\)-disc, the difference at \(x,x'\)
is bounded by the right side of \eqref{eq:step-F-bound}
times \(2|x-x'|/(\sigma_*\mathsf t)\).

With \(\mathbb V^{\rm int}=\mathbb V^{\rm J}
+\mathbb V^{\rm G}+\tau\mathbb V^{\rm A}\), the exact
local sum rules are
\begin{align}
 \sum_z\mathbf G_\tau(X,P,z)
       &=E^\circ[\{\mathbb V^{\rm q}\}](X),
                                             \label{eq:step-G-sum}\\
 \sum_z\mathbf F[\mathbb V'](X,P,z)
       &=E[\{\mathbb V^{\rm q},\mathbb V^{\rm int},\mathbb V'\}](X)
                       -E^\circ[\{\mathbb V^{\rm q}\}](X).
                                             \label{eq:step-F-sum}
\end{align}
The site sums run over \(X\cap\mathbb T^{(k+1)}\).
For a whole-torus triple \(P_*\), these formulas give
an absolutely convergent integral and the identity
\begin{equation}\label{II-53-67}
\begin{split}
 &\int d\mu_{\tau^{-1}C_*(P_*)}(B)
       \prod_{b\in\mathfrak X^+}\mathbf1_{\{|B(b)|<\mathfrak b\}}
       \exp\!\left[\sum_{Y\in\mathcal D_k}
          (\mathbb V^{\rm q}+\mathbb V^{\rm int}+\mathbb V')
                    (Y;P_*|_Y,B|_Y)\right]\\
 &\hspace{20mm}=
       \exp\!\left[\sum_{X,z}
          (\mathbf G_\tau+\mathbf F[\mathbb V'])(X,P_*|_X,z)\right]
       \ne0.
\end{split}
\end{equation}
The integrand is zero where its test fails. Each refined
series in \eqref{II-20-115} has a summable majorant for
the integrals of the absolute values of its terms,
uniform in \(x,P\) and in bounded input families.
\end{theorem}

\begin{proof}
The disc gives \(\operatorname{Re}\tau>1/2\),
\(|\tau|<3/2\), and \(|1-\tau|K_{\mathsf G}\le w_*/2\).
For the scaled Gaussian structure, the inverse covariance
and coupling defects are at most
\(2\hat\varepsilon_0+c_{\rm r}|1-\tau|\); the product
\(\Gamma^{\mathsf T}C\Gamma\) is unchanged.
Thus all three defects are at most
\(\hat\varepsilon_{\rm S}-\tfrac12c_{\rm r}(LM)^{-3}\).
The structure remains local, and its whole covariance is
\(\tau^{-1}C_*\). This is the scaling calculation at
the end of Theorem~\ref{II-38-61}. The first numerical
condition gives \(u\le\mathsf C_{\rm ex}^{-1}
\min(\gamma_{\rm r},c_{\rm r}^{-2})\), and the threshold
floor gives \(\mathsf p_*\le1/2\).

At real fields, \(|\langle A,B\rangle|\le|A||B|\).
Using \(ab\le(a^2+b^2)/2\) and the row and column
bounds of \(\widehat{\mathsf G}_\square\) gives
\[
 \left|\tfrac12(B,\widehat{\mathsf G}_\square B)_1\right|
       \le\tfrac12c_{\mathsf G}\sum_{b\in\square}|B(b)|^2.
\]
The same upper bound holds for any single bond contribution.
Since a fine cube has \(d_k=0\), this proves
\(\|\hat{\mathsf g}_b\|_{\star,\lambda_V}\le K_{\mathsf G}\)
and \(\|\mathbb V^{\rm q}\|\le w_*/2\).
Also \(a_{0,k+1}\le1\), so
\(\delta(\mathsf t)=s_{\mathfrak a}a_{0,k+1}\le s^*\).
Proposition~\ref{II-20-103} therefore yields
\begin{equation}\label{eq:step-potential-margins}
\begin{gathered}
 \|\mathbb V_v^{\rm int}\|\le\tfrac72K_P\delta
                                       \le\tfrac7{16}w_* ,\qquad
 \|\mathbb V'_v\|\le K_\Lambda s_{\mathfrak a}\nu
                                       \le\tfrac14w_*,\\
 \|\dot{\mathbb V}^{\rm int}_{v,b}\|\le32K_P\delta,
 \qquad\|\dot{\mathbb V}'_{v,b}\|
                              \le9K_\Lambda s_{\mathfrak a}\nu,\\
 \|\dot{\mathbb V}^{\rm int}_{v,b}+\dot{\mathbb V}'_{v,b}\|
                       \le9(7/16+1/4)w_*=99w_*/16.
\end{gathered}
\end{equation}
Here all norms have rate \(\lambda_V\).

Apply the insertion theorem to the first and third
lines of \eqref{II-20-115}, and the cutoff theorem
to its second line. They prove existence, support, and
the bond bounds
\[
 |\mathbf G_\tau(X,b;P)|\le|1-\tau|K_\Pi K_{\mathsf G}e^{-\lambda_{\rm out}d},
\]
\[
 |E^\chi+E^\bullet|(X,b;P)
 \le\bigl(K^\Pi_Pe^{4\lambda_1}\mathsf p
          +32K_\Pi K_P\delta
          +9K_\Pi K_\Lambda s_{\mathfrak a}\nu\bigr)e^{-\lambda_{\rm out}d}.
\]
The tent sum over nonpivot bonds is at most \(3L^3\).
Equation~\eqref{II-53-34} now gives
\eqref{eq:step-G-bound} and \eqref{eq:step-F-bound}.

Along the first path the sum of its bond velocities
is \(\mathbb V^{\rm q}\). The path identity gives
its cluster difference from the empty system, whose
no-test cluster sum is zero. Along the third path it
gives the difference between the full system and
\(\{\mathbb V^{\rm q}\}\) at finite threshold.
The cutoff term supplies the difference between
the finite-threshold and no-test sums for this last
system. The tent identity \eqref{II-20-14} proves
\eqref{eq:step-G-sum} and \eqref{eq:step-F-sum}.
Restriction consistency identifies the local terms
with those computed from the whole-torus data.
The complete Gaussian cluster expansion then proves
\eqref{II-53-67}.

The uniform insertion majorant \eqref{II-20-51},
together with \eqref{eq:step-potential-margins},
justifies all exchanges of the path integral and
refined sums. For a cutoff term, apply
Lemma~\ref{II-20-71} separately to each refined
density in \eqref{II-20-79}. Its derivative has
integral of absolute value bounded by that density's
total variation. These masses are bounded by the
single-bond marked series, independently of the
parameters. This gives the asserted common majorant.

Holomorphy follows from the parameter version of the
insertion theorem, the fixed-density cutoff formula,
and finite tent sums. Pair dependence and gauge
invariance hold because the corresponding entries
are identical, or related by orthogonal adjoint
maps. They pass through every integral. The same
argument proves translation covariance and
conjugation, using the real covariant tents.
For real inputs, holomorphic potentials real on
real fields satisfy \(V(B^*)=\overline{V(B)}\)
throughout the connected field polydisc by the
identity theorem. Thus the output is real.
The bounds already proved show that these functions
are local terms on the output core.

It remains to prove dependence on the fourth input.
Put \(\Delta V=\mathbb V'-\mathbb V''\) and
\(\nu_\Delta=\nu(\Delta V)\le2N_*\).
Linearity of interpolation gives
\(\|\Delta V_v\|\le K_\Lambda s_{\mathfrak a}\nu_\Delta\)
and \(\|\Delta\dot V_{v,b}\|
\le9K_\Lambda s_{\mathfrak a}\nu_\Delta\).
Split the difference of insertion sums into the
change of insertion and the change of its background
system. The first costs at most
\(9K_\Pi K_\Lambda s_{\mathfrak a}\nu_\Delta e^{-\lambda_{\rm out}d}\).
For the second, the fourth-origin base norm is
at most \(w_*/4\), its change is at most \(w_*/2\),
and its insertion norm is at most \(99w_*/16\).
Equation~\eqref{II-20-59} bounds this cost by
\((99/2)K_\Pi K_\Lambda s_{\mathfrak a}\nu_\Delta
 e^{-\lambda_{\rm out}d}\).
Integrating and using \(9+99/2=117/2<81\), then
summing the tents, proves \eqref{eq:step-fourth-lipschitz}.
Finally Cauchy's derivative estimate on the half-radius
\(x\)-disc gives the stated \(x\)-dependence.
\end{proof}

For comparison with the decay rate of the input
potentials, let \(0<\delta_{\rm w}<1/2\), and define
\begin{equation}\label{II-53-35}
 L_\Pi=\frac{2\mathsf C_{\rm ex}(1+4\beta_{\rm s})}{1-2\delta_{\rm w}},
 \qquad
 \kappa_\Pi=\frac2{1-2\delta_{\rm w}}
       \left[\frac{\mathsf C_{\rm ex}}L
          (8c_{\rm w}+\mathsf C_{\rm ex}+\lambda_{\rm s})
                                      +2\lambda_{\rm s}\right].
\end{equation}
If \(L\ge L_\Pi\) and \(\kappa\ge\kappa_\Pi\),
the term proportional to \(\kappa\) in \(\lambda_V\)
and its remaining terms are each at most
\((1-2\delta_{\rm w})\kappa/2\). Hence
\(\lambda_V\le(1-2\delta_{\rm w})\kappa\).
If also \(\bar\kappa\ge4\lambda_1+1\), then
\(\lambda_1\ge\mathsf C_{\rm ex}+\lambda_{\rm s}\ge2\), and hence
\[
 \bar\kappa-2\ge4\lambda_1-1>\lambda_1
           >\lambda_{\rm out}>\kappa>\lambda_V.
\]
These choices are independent
of the lattice, step and running coupling.

\subsection{The small-field measure}\label{sec:small-field-measure}
We next identify the integral to which the Gaussian expansion applies.
All lengths and pairings in this subsection are in level-$k$ units.
Put $t=k+1$, and let $\Gamma$, $\mathfrak X_0$ and $\mathfrak X$
be the whole-lattice sets used above. If $N_j$ is the number of
level-$j$ sites, then
\[
 \#\Gamma=3N_t,\qquad
 \#\mathfrak X^+=3(N_k-N_t),\qquad
 n_{\mathcal O}=3(N_k-N_t),\qquad
 n_\perp=3\#\mathfrak X^+-n_{\mathcal O}.
\]
Here $n_{\mathcal O}$ counts the gauge variables within blocks.

Let $W$ have a real regular critical background $U_t(W)$.
Assume that its averaged contours are defined, and choose a
representative $U_+$ for which
$p=\mathcal M^kU_+$ is axial at the level-$t$ sites.
The residual-orbit and axial-representative results make $p$
independent of this choice. We work on an open set where $p$ is
defined continuously and
$\sup|\partial p-1|<\varpi_\chi$, with $\varpi_\chi>0$
sufficiently small depending only on $L$.
Use the actual constraint chart from \cref{II-41-33}:
\[
 V_W[\mathcal A]
   =e^{i\psi_p^{\rm lin}(\mathsf C_p\mathcal A)}p,\qquad
 V_W(B)=V_W[gB],\qquad g=\mathsf t^{-1/2}.
\]
There are $r_\chi>0$ and $C_\chi\ge1$, depending only on $L$,
such that this chart is defined when
$C_\chi\sup|\mathcal A|<r_\chi$, and its full logarithmic
field has norm at most $2C_\chi\sup|\mathcal A|$.
Decrease $r_\chi,\varpi_\chi$ so that all these real charts
lie in the small-curvature branch of the averages and in
the radii of \cref{II-6-12}.

Write $b=\mathfrak b(\mathsf t)$, $\delta=gb$, and
$\epsilon=\varepsilon(\mathsf t)$. Assume
\begin{equation}\label{eq:sf-chart-thresholds}
 \epsilon\le1,\qquad C_\chi\delta\le r_\chi,\qquad
 C_\chi\delta\le\epsilon .
\end{equation}
Increasing $C_\chi$ if necessary, the last inequality is
ensured by $A_0/A_1\ge C_\chi$.
Define $\chi^{\rm sf}_{\mathsf t}(V)$ to be one precisely when
$p(\mathcal M V)$ is stably defined and
\[
 \left|\frac1i\log\bigl(V(b')p(\mathcal M V)(b')^{-1}\bigr)\right|
 <\delta\qquad(b'\in\mathfrak X^+);
\]
otherwise it is zero. These are intrinsic logarithms, and no
pivot variable is tested. Let $\mathfrak G_k[\mathsf t]$ be
the block gauge-fixing factor: for each
$y\in\mathbb T^{(t)}$ and $x\in B(y)\setminus\{y\}$ its factor is
\[
 \mathsf z(\mathsf t,\epsilon)^{-1}
 \mathbf1_{\{|V(y,x)-1|<\epsilon\}}
 \exp\{-\mathsf t(1-\tfrac12\operatorname{Re}\operatorname{tr}V(y,x))\},
\]
where $\mathsf z$ is the corresponding single-group integral.
The product is zero where a contour is undefined.
For a bounded Borel function $\chi^{\rm sf}_{\mathsf t}F$,
the small-field measure is the pushforward under $\mathcal M$
of $\mathfrak G_k[\mathsf t]\chi^{\rm sf}_{\mathsf t}F\,dV$.

For the chart field
$B'=\psi_p^{\rm lin}(\mathsf C_pgB)$ put
\begin{equation}\label{eq:sf-Jacobian}
 \mathsf J_W(B)=
 \sum_{c\in\Gamma}\log\det(1-\mathsf{Jb}_c(p,gB))
 +\sum_{b'\in\mathfrak X_0^+}
       \log\sigma_1(|B'(b')|^2),\qquad
 \sigma_1(s)=\left(\frac{\sin\sqrt{s}}{\sqrt{s}}\right)^2 .
\end{equation}
The value at $s=0$ is one; all logarithms here are real.
Set
\begin{equation}\label{II-40-57}
 \mathsf{pf}_k(\mathsf t)
 =\mathsf z(\mathsf t,\epsilon)^{-(N_k-N_t)}
  (\sigma_{\rm H}\mathsf t^{-3/2})^{\#\mathfrak X^+}
  L^{-3\#\Gamma},\qquad \sigma_{\rm H}=(2\pi^2)^{-1}.
\end{equation}

At a whole-lattice real background let
$\mathsf B_W=\mathsf C_p^{\mathsf T}\Delta_k(U_+,J^\circ)\mathsf C_p$
be the action part of \eqref{eq:cv-form}, and put
\[
 \widehat{\mathsf G}_W
   =\mathsf C_p^{\mathsf T}\mathsf G^{(2)}[p]\mathsf C_p,\qquad
 \mathsf A_W=\mathsf B_W+\widehat{\mathsf G}_W,\qquad
 C_W=\mathsf A_W^{-1}.
\]
The current in this formula is
$J^\circ=L^{-3}J^{(t)}(U_+)$, as in \eqref{II-38-12}.
Take the background in the radius of \cref{II-38-61}, so that
$\mathsf A_W$ is real positive definite.
Let $S$ be gauge invariant on the minimal orbits in question,
and assume that $U_k(V_W(B))$ exists for every $\sup|B|<b$.
For the Hessian identification, assume in addition that the
constant tower at $(U_+,J^\circ)$ satisfies the hypotheses of
\cref{I-88-53,I-88-93}, and that there is an open neighbourhood
$\mathcal N$ of $p$ in the space of all real level-$k$ fields
on which $U_k(V)$ exists and its minimal orbit contains the
restored critical response at $U_+$. Thus this hypothesis
concerns arbitrary nearby data, not just the constraint fibre.
At $k=0$ it means $U_0(V)=V$ on the corresponding open
regularity neighbourhood.
Write
\[
 \mathbb F_W(B)=A^{\eta_k}(U_k(V_W(B))),\qquad
 G_W(B)=\sum_y\sum_{x\in B(y)\setminus\{y\}}
        (1-\tfrac12\operatorname{Re}\operatorname{tr}V_W(B)(y,x)).
\]
Define the three remainders directly by
\begin{align*}
 \mathbb P^J(B)&=\mathsf J_W(B),\\
 \mathbb P^G(B)&=-\mathsf tG_W(B)
                    +\tfrac12(B,\widehat{\mathsf G}_WB)_1,\\
 \mathbb P^A(B)&=-\mathsf t[\mathbb F_W(B)-\mathbb F_W(0)]
                    +\tfrac12(B,\mathsf B_WB)_1,
\end{align*}
and $s_W[S](B)=S(U_k(V_W(B)))-S(U_t(W))$.
These formulas also specify the normalization of the potentials.
Under the global identification hypotheses stated before \eqref{II-41-27},
they equal the whole totals of its three actual local potentials.
Put $P_W=(U_+,0,J^{(t)}(U_+))$ and define
\[
 Z^{(k)}(W;\mathsf A_W)
   =\prod_{c\in\Gamma}|\det m_c(p)|^{-1}
      \int_{\mathbb R^{3\#\mathfrak X^+}}
             e^{-\frac12(B,\mathsf A_WB)_1}\,dB
   =Z_k[P_W].
\]
The last equality uses the actual form at $P_W$ in
\eqref{II-44-40}.

\begin{proposition}\label{II-40-99}
Let the preceding hypotheses hold, let
$\tau=x/\mathsf t$ with $\operatorname{Re}\tau>0$, and suppose
$B\mapsto F_{x,S}(V_W(B))$ is bounded Borel on $\sup|B|<b$,
where
$F_{x,S}(V)=\exp\{-xA^{\eta_k}(U_k(V))+S(U_k(V))\}$.
The chart density is
\begin{align}
 \mathsf I[F_{x,S}](W)
 &=\mathsf{pf}_k(\mathsf t)\,
   Z^{(k)}(W;\mathsf A_W)
   e^{-xA^{\eta_k}(U_t(W))+S(U_t(W))}
   \tau^{-3\#\mathfrak X^+/2} \notag\\
 &\quad{}\times
 \int\mu_{\tau^{-1}C_W}(dB)\,
       \mathbf1_{\{\sup|B|<b\}}
 \exp\left\{\begin{aligned}
       &-\frac{1-\tau}{2}(B,\widehat{\mathsf G}_WB)_1\\
       &+\mathbb P^J(B)+\mathbb P^G(B)
        +\tau\mathbb P^A(B)+s_W[S](B)
       \end{aligned}\right\}.
 \label{II-40-100}
\end{align}
The integration is over real fields, and the integral is
absolutely convergent. If
$\chi^{\rm sf}_{\mathsf t}F_{x,S}$ is bounded Borel on the
whole configuration space, this is the Borel density of the
small-field pushforward on the open background domain.
It is continuous there if the chart integrand is continuous
in the background and locally uniformly bounded.

Moreover the uncut quadratic gauge integral satisfies
\begin{equation}\label{II-43-105}
 \int\mu_{\tau^{-1}C_W}(dB)
   e^{-(1-\tau)(B,\widehat{\mathsf G}_WB)_1/2}
       =\tau^{n_{\mathcal O}/2}.
\end{equation}
For the connected quadratic terms of the Gaussian expansion,
with the logarithm normalized at zero potential,
\begin{equation}\label{II-43-100}
 \sum_{X,z}\mathbf G_\tau(X,P_W,z)
       =\frac{n_{\mathcal O}}2\operatorname{Log}\tau .
\end{equation}
The last identity is asserted on the small $\tau$-domain
where that expansion is defined.
\end{proposition}

\begin{proof}
\emph{The fibre and its measure.}
The chart has average $W$ and unchanged nonpivot logarithms.
Conversely, a field of average $W$ passing the nonpivot test
has these logarithms as its chart coordinates. On each double
block it has the same nonpivot entries as the chart field;
\cref{II-6-12} gives equality of the remaining pivot entry,
including across the real averaging branches. Thus the chart
parametrizes exactly the tested fibre.

All within-block bonds are nonpivot. A block contour has at
most $3(L-1)/2$ bonds, so its transport changes by at most
$3L\delta/2$. The small-curvature choice makes both contour
families lie within $1/40$ of their means. If their members
are $V_i,V_i^0$, their means are $\bar V,1$, and
$e=\max_i|V_i-V_i^0|$, the logarithmic mean equations give
\[
 |\bar V-1|\left(1-\frac1{39}-\frac1{20}\right)
       \le\frac{40}{39}e,\qquad |\bar V-1|\le\frac{10}{9}e.
\]
Indeed, subtract the mean equations and use
$|\log X-\log Y-(X-Y)|\le |X-Y|/39$ in the $1/40$ ball;
then write
$V_i\bar V^{-1}-V_i^0
 =(\bar V^{-1}-1)+(V_i-1)(\bar V^{-1}-1)+(V_i-V_i^0)$.
Consequently $|V_W(B)(y,x)-1|\le C_\chi\delta$.
All gauge indicators are therefore one, and the gauge factor is
$\mathsf z(\mathsf t,\epsilon)^{-(N_k-N_t)}e^{-\mathsf tG_W(B)}$.

For clarity, the determinant in the change of variables can
be calculated in a single double block. Write
$q_U(a)=i^{-1}\log(\mathcal M(e^{ia}U)\mathcal M(U)^{-1})$
and $n_U=\partial_0q_U(0)=L m_c(U)$, where $\partial_0$
differentiates the pivot variable. The correction used in the
chart obeys
$q_U(a-A_cy_U(a))=\dot q_U(a)$.
At $\dot q_U(a)=0$, differentiation in that variable gives
\[
 \partial_0q_U(a-A_cy_U(a))(1-\partial_0y_U(a))=n_U.
\]
For $U'=e^{i(a-A_cy_U(a))}U$ and pivot value $\beta$,
differentiating the exponential gives
\[
 n_{U'}\,\mathsf{dE}(\beta)(1-\partial_0y_U(a))=n_U,\qquad
 \mathsf{dE}(Y)H=\int_0^1\operatorname{Ad}(e^{isY})H\,ds.
\]
On the line spanned by $Y$, this map is the identity; on its
orthogonal plane its determinant is $(\sin|Y|/|Y|)^2$.
Thus
\begin{equation}\label{eq:sf-pivot-Jacobian}
 (\det n_{U'})^{-1}
  =(\det n_U)^{-1}\sigma_1(|\beta|^2)
                       \det(1-\partial_0y_U(a)).
\end{equation}
The factors are positive by the chart bounds.

In exponential coordinates, normalized Haar measure is
$\sigma_{\rm H}\sigma_1(|Y|^2)\,dY$ for $|Y|<\pi$.
One can see the normalization and invariance directly:
the radial integral is
$4\pi\int_0^\pi\sin^2r\,dr=2\pi^2$, and differentiating
$e^{iT(Y)}=e^{iY}u$ gives
$\mathsf{dE}(T(Y))DT(Y)=\mathsf{dE}(Y)$.
Hence this density is invariant under right translation away
from the two omitted chart points, which have zero measure.

Fix the nonpivot entries. The small-curvature pivot domain is
open, and its averaging map is injective by \cref{II-6-12}.
A plaquette containing the pivot places this domain in one
exponential chart after the radius is decreased. The
differential calculation above makes it a local
diffeomorphism with Haar Jacobian $\det n_{U'}$.
The ordinary change-of-variables formula, followed by
\eqref{eq:sf-pivot-Jacobian}, therefore integrates out that
pivot. Distinct double blocks contain no other pivot, so this
calculation factors over $\Gamma$.
The nonpivot substitution $a=gB$ contributes
$(\sigma_{\rm H}\mathsf t^{-3/2})^{\#\mathfrak X^+}$.
It follows that the tested pushforward has density
\begin{equation}\label{eq:sf-Haar-density}
 \mathsf{pf}_k(\mathsf t)
 \prod_{c\in\Gamma}|\det m_c(p)|^{-1}
 \int_{\sup|B|<b}
      e^{-\mathsf tG_W(B)+\mathsf J_W(B)}F(V_W(B))\,dB .
\end{equation}
The factor $L^{-3\#\Gamma}$ in \eqref{II-40-57} accounts
for $\det n_U=L^3\det m_c(U)$.

This argument also gives the Borel assertion without a
choice of conditional measures. The tested set is Borel;
the intrinsic nonpivot logarithms and the continuous chart
make its inverse Borel. For a nonnegative Borel integrand,
Tonelli applies to the successive pivot and nonpivot
substitutions. Apply the same identity to the absolute value
before treating a complex integrand. On every chart fibre
the Jacobian is bounded and the field box has finite volume;
boundedness of $F(V_W(B))$ therefore gives absolute convergence.
The same bounds locally in $W$, followed by dominated
convergence, give the stated continuity.

\emph{The action Hessian and gauge directions.}
The background response identifies
$\varphi(\beta)=A^{\eta_k}(U_k(e^{i\beta}p))$
near $\beta=0$. We record its second differential in the
normalization used here. Write the response as
$e^{i\xi\widehat A(\beta)}U_+$ followed by its restoring
gauge, and let $\omega(\beta)$ be its constraint multiplier.
On this constant tower all constraint rows have spacing one,
and the row injection $\mathsf J$ has
$\mathsf J^{\mathsf T}\mathsf J=I$. The response equations are
\[
 F'(\widehat A)=K_{\widehat A}^{\mathsf T}\omega,\qquad
 \mathcal Q(\widehat A)=\mathsf J\beta,\qquad
 \mathsf RD^{\mathsf T}\widehat A=0.
\]
Differentiating at zero, with $a=\widehat A'(0)\nu$ and
$\mu=\omega'(0)\nu$, gives the saddle equations
\[
 (\Delta_1+D\mathsf RD^{\mathsf T})a-Q^{\mathsf T}\mu=0,
 \qquad Qa=\mathsf J\nu.
\]
Here $\mathsf G_1^{-1}
=\Delta_1+D\mathsf RD^{\mathsf T}+Q^{\mathsf T}Q$.
Solving these saddle equations gives
\[
 a=\mathsf H_1\mathsf J\nu,\qquad
 \mu=[(Q\mathsf G_1Q^{\mathsf T})^{-1}-I]\mathsf J\nu.
\]
The chain rule and differentiated constraint give the
envelope identity
$\varphi'(\beta)\nu=(\omega(\beta),\mathsf J\nu)_{\mathbf B}$.
Since $\omega(0)=\mathsf H_1^{\mathsf T}J^\circ$, we obtain
\[
 D\varphi(0)\nu=(\mathsf H_1\mathsf J\nu,J^\circ)_\xi,\qquad
 D^2\varphi(0)(\nu,\nu)=(\nu,\mathsf Z_0\nu)_1,\qquad
 \mathsf Z_0=\mathsf J^{\mathsf T}
                   (Q\mathsf G_1Q^{\mathsf T})^{-1}\mathsf J-I .
\]
For $k=0$, the response is $\widehat A(\beta)=\beta$;
$Q=\mathsf H_1=I$ and
$\mathsf G_1^{-1}=\Delta(U_+)+I$ give the same identities.

If $V(s)$ stays in the fibre, its background
lift has $\mathcal M^{k+1}U(s)=W$ and starts at $U_+$.
Criticality gives $D\varphi(0)B=0$ for every
$B\in\ker Q_p$. Differentiating the constraint twice gives
\[
 \ddot\beta(0)=\kappa-2\mathsf D_p^{(2)}(B,B),
       \qquad \kappa\in\ker Q_p.
\]
The chain rule and the already identified response Hessian
therefore give
\[
 \frac{d^2}{ds^2}\Big|_0\varphi(\beta(s))
  =(B,\mathsf Z_0B)_1
    -(B,\mathsf E^{(2)}[\mathsf H_1^{\mathsf T}J^\circ]B)_1
  =(B,\Delta_k(U_+,J^\circ)B)_1.
\]
The gauge-fixing term has zero first differential at the
axial field and second differential $\mathsf G^{(2)}[p]$.
Thus the free-coordinate Hessian is precisely $\mathsf A_W$,
not a separately chosen Gaussian form.

Let $\mathcal E^0$ be the real site functions vanishing at
the level-$t$ sites. Define
$j\rho=\chi_{\mathfrak X}\mathfrak d_p\rho$ and let
$(rB)(x)$ be the first variation of the block contour
at $(y,x)$ in the direction $\mathsf C_pB$, with
$(rB)(y)=0$. Covariance of the actual average gives
$Q_p\mathfrak d_p\rho=0$ and $rj\rho=\rho$:
the first contour variation is $\rho(x)-\rho(y)$
because $p(y,x)=1$. Moreover
\[
 (B_1,\widehat{\mathsf G}_WB_2)_1=(rB_1,rB_2)_0.
\]
Every real $B$ consequently has the unique decomposition
$B=B''+j\rho$ with $\rho=rB$ and $rB''=0$.
Gauge invariance along the curves
$V(s)^{e^{-is a\rho}}$, for every real $a$, shows that
the action Hessian annihilates $j\rho$. Indeed, comparing
the second derivatives gives
$2a\,q(B,j\rho)+a^2q(j\rho,j\rho)=0$ for every $a$.
The two quadratic forms therefore occupy complementary
blocks of dimensions $n_\perp,n_{\mathcal O}$.
Taking their determinant in an adapted basis, whose
change-of-basis determinant cancels in the quotient, gives
\[
 \det(\tau\mathsf B_W+\sigma\widehat{\mathsf G}_W)
   =\tau^{n_\perp}\sigma^{n_{\mathcal O}}\det\mathsf A_W .
\]
Positivity of $\mathsf A_W$ makes $\mathsf B_W$
nonnegative, and hence the form on the left has positive
real part for $\operatorname{Re}\tau,\operatorname{Re}\sigma>0$.
Continuity from $\tau=\sigma=1$ fixes the square-root branch.
The Gaussian determinant formula now proves
\eqref{II-43-105}.

\emph{The density and its logarithm.}
Compatibility of the two critical backgrounds gives
$\mathbb F_W(0)=A^{\eta_k}(U_t(W))$.
Insert the definitions of the potentials: their quadratic
terms cancel the exponent
$-\tau(B,\mathsf A_WB)_1/2$, leaving exactly the
integrand of \eqref{eq:sf-Haar-density}, times
$e^{xA^{\eta_k}(U_t(W))-S(U_t(W))}$.
Since
\[
 \mu_{\tau^{-1}C_W}(dB)
   =\tau^{3\#\mathfrak X^+/2}
     \left(\int e^{-(B,\mathsf A_WB)_1/2}\,dB\right)^{-1}
     e^{-\tau(B,\mathsf A_WB)_1/2}\,dB ,
\]
this proves \eqref{II-40-100}, with the same absolute
convergence already obtained in \eqref{eq:sf-Haar-density}.

Finally multiply the quadratic potential by
$\theta\in[0,1]$. The whole Gaussian integral is
\[
 \exp\left\{\frac{n_{\mathcal O}}2
   [\operatorname{Log}\tau-
    \operatorname{Log}(\tau+\theta(1-\tau))]\right\}.
\]
The connected expansion has the same exponential by its
whole-integral identity, is continuous in $\theta$, and
has zero total at $\theta=0$. Their logarithms thus differ
by a continuous $2\pi i\mathbb Z$-valued function vanishing
at zero. At $\theta=1$ this proves \eqref{II-43-100}.
\end{proof}

\subsection{The actual small-field step}\label{sec:actual-small-field-step}
We now use the actual covariance, potentials and normalizer in the
one-step expansion. The only loss in the fourth-potential estimate
depends on the chosen weights for the earlier terms; its smallness
ceiling must therefore be chosen with those same weights.

Keep the geometric hypotheses of \cref{II-38-61,II-44-37} and
the notation of \cref{sec:one-step-local-expansion}.
The shifts and substitutions are those of the fourth-potential
construction, with an admissible system of radii \(\mathfrak a\);
retain its octet condition when \(k\ge2\).
Choose the common radii so that
\begin{equation}\label{eq:actual-step-radii}
\begin{gathered}
 \bar a_0+\bar a_1\le a^\sharp,\qquad
 \bar a_0\le\min\{a_{\rm G}/M,\alpha_N,a_P,a_N\},\qquad
 \bar a_1\le\min\{a_P,a_N\},\\
 \theta_{\rm Str}M\ge16\bar\kappa,\qquad
 e^{-\theta_{\rm Str}M}+\bar a_1+\bar a_0
 \le\min\left\{a_{\rm Str},
           \frac{u_{\rm G}}{2C_{\rm Str}}(LM)^{-3}\right\},\\
 \mathsf C_{\mathcal A}\delta(\mathsf t)\le\varepsilon_P,
 \qquad \lambda_V\le\bar\kappa-2 .
\end{gathered}
\end{equation}
Here \(a^\sharp\) is the tower-continuation radius, and
the fourth-potential construction retains its own admissibility,
shift and rate inequalities. No existence condition on a tower is
replaced by a curvature bound.

When \(k\ge1\), the local action member is
\begin{equation}\label{II-52-11}
 \mathsf W_\zeta(X,P)
   =\zeta\eta_k^{-1}\sum_{p:X(p)=X}
          \left(1-\tfrac12\operatorname{tr}\partial\mathbf U(p)\right),
 \qquad (\mathbf U,\mathbf J)=\mathbf p(P).
\end{equation}
Here \(X(p)\) is the union of the cubes containing the corners of \(p\).
It has at most four cubes meeting at a corner, so \(d_k(X(p))=0\).
The member is gauge invariant, holomorphic, and zero at the trivial pair.
Indeed
\(\left|1-\frac12\operatorname{tr}V\right|
=\frac12|\det(V-I)|\le\frac12|V-I|^2\), and there are at most
\(12(M/\eta_k)^3\) plaquettes with a given \(X(p)\). Thus
\(\|\mathsf W_\zeta\|_{k,\kappa}\le6M^3a_{0,k}^2|\zeta|\)
on the input core. Its total is \(\zeta A^{\eta_k}(\mathbf U)\).
If the last slot is vacant, let
\(\mathfrak g=\mathfrak f_\circ\sqcup\mathsf W_\zeta\) denote the
array with this member in slot \(k\); otherwise take
\(\mathfrak g=\mathfrak f_\circ\) and \(\zeta=0\).
At \(k=0\) the input array is empty.
Select one of the weight cases
\((\omega,\mathsf K)\) of \cref{II-52-89}, so that
\begin{equation}\label{eq:actual-fourth-size}
 |\mathbb V'[\mathfrak g](Y;P,B)|
 \le \mathsf K s_{\mathfrak a}
       N_{\mathfrak a,\omega}(\mathfrak g)e^{-\lambda_Vd_k(Y)}.
\end{equation}
At \(k=0\) the array is empty and this potential is zero.
In \eqref{II-53-34} replace \(K_\Lambda\) by \(\mathsf K\),
and denote the resulting constant by \(K_{\mathbf M}(\mathsf K)\).
Use
\begin{equation}\label{eq:actual-step-ceiling}
 \check s(\mathsf K)
   =\min\{s^*,\,w_*/(4\mathsf K N_*)\},\qquad
 s_{\mathfrak a}\le\check s(\mathsf K),\qquad
 N_{\mathfrak a,\omega}(\mathfrak g)\le N_* .
\end{equation}
Assume the other conditions of
\eqref{eq:step-numerical-conditions}, with the same replacement.
The radius \(s^*\) here is small enough for the shift construction
and \(8K_Ps^*\le w_*\); it need not satisfy the fourth-potential
inequality for every case.

The actual output is
\begin{equation}\label{II-53-25}
 E^{(k+1)}(X,P,z)
  =\mathbf N^{(k+1)}(X,P,z)
       +\mathbf G_\tau(X,P,z)
       +\mathbf F[\mathbb V'[\mathfrak g]](X,P,z).
\end{equation}
The normalizer in this formula is independent of the input array
and of \(\mathsf t,x\).

\begin{theorem}\label{II-127-29}\label{II-128-31}
Under the preceding hypotheses, the actual fields satisfy the local
input conditions of \cref{II-53-12}. For
\(x\in D(\mathsf t,\sigma_*\mathsf t)\), the family
\eqref{II-53-25} is defined on every output core
\(\mathcal U^\circ(X)\), \(X\in\mathcal D_{k+1}\).
It is a pointed family, holomorphic in \(x\), the slice variables,
and holomorphic input parameters under uniform size bounds.
It has the translation and conjugation properties when the input
arrays and the selected shifts have them.

Suppose also that
\begin{align}
 4\beta_{\rm s}\kappa&\ge1,\label{II-128-3}\\
 \delta_NM&\ge\kappa+1,\label{II-128-4}\\
 E_0&\ge2E_N .\label{II-128-5}
\end{align}
Then, with \(d=d_{k+1}(X)\),
\begin{equation}\label{II-128-32}
\begin{split}
 |E^{(k+1)}(X,P,z)|
 \le\Bigl[\frac{E_0}{2}+K_{\mathbf M}(\mathsf K)
 \bigl(&|1-\tau|+\delta(\mathsf t)
       +s_{\mathfrak a}N_{\mathfrak a,\omega}(\mathfrak g)\\
       &+e^{-\gamma_2\mathfrak b(\mathsf t)^2/2}\bigr)\Bigr]
                                      e^{-(\kappa+1)d}.
\end{split}
\end{equation}
If, further,
\begin{align}
 K_{\mathbf M}(\mathsf K)\sigma_*&\le E_0/4,\label{II-128-6}\\
 K_{\mathbf M}(\mathsf K)
   \bigl(\delta(\mathsf t)+s_{\mathfrak a}N_*
                  +e^{-\gamma_2\mathfrak b(\mathsf t)^2/2}\bigr)
                                                   &\le E_0/4.\label{II-128-7}
\end{align}
the bracket in \eqref{II-128-32} is at most \(E_0\);
at \(x=\mathsf t\) it is at most \(3E_0/4\).

For a whole-torus block field \(W\), assume the background and
chart hypotheses of \cref{II-40-99}, that its triple \(P_W\)
lies in \(\mathcal U^\circ(T_\ell)\), and that the actual
potentials have the global identification stated before
\eqref{II-41-27}
and the fourth potential sums to \(s_W[S]\). In particular,
retain the open whole-data response neighbourhood required
in \cref{II-40-99}. For the density conclusion also assume that
\(\chi^{\rm sf}_{\mathsf t}F_{x,S}\) is bounded and Borel on the
whole configuration space \(\mathcal V_k\). Then the small-field
pushforward density is
\begin{equation}\label{eq:actual-step-density}
\begin{split}
 \rho^{\rm sf}_{x,S}(W)
  ={}&\mathsf{pf}_k(\mathsf t)\,Z_{\rm triv}^{(k)}
               \tau^{-3\#\mathfrak X^+/2}\\
 &{}\times\exp\left\{
  -xA^{\eta_k}(U_t(W))+S(U_t(W))
  +\sum_{X,z}E^{(k+1)}(X,P_W|_X,z)\right\}.
\end{split}
\end{equation}
All integrals in this identity are absolutely convergent.
Equivalently, remove \(\mathbf G_\tau\) from the sum and
replace the power of \(\tau\) by \(\tau^{-n_\perp/2}\).
The powers use the principal logarithm.
\end{theorem}

\begin{proof}
We first check the common domain. By
\cref{II-16-12,II-16-46},
\begin{equation}\label{II-127-27}
 \mathcal U^\circ(Y)
 \subset\mathcal Q_t(Y;\bar a_0,\bar a_1,\bar a_0)
 \subset
 \mathcal Q_t(Y;a_P,a_P,a_P)\cap\mathcal Q^N(Y).
\end{equation}
The first inclusion retains all the tower conditions of the core.
The second follows coordinate by coordinate from
\eqref{eq:actual-step-radii}. Reading the triple at level \(k\)
by \eqref{II-38-12}, its potential and derivative acquire the
factors \(L^{-1},L^{-2}\), and its current the factor \(L^{-3}\).
The gauge boxes required for the actual Gaussian structure are
contained in the boxes in the simple-domain definition, by
\(12\mathfrak M\le\mathsf oLM\) and \(5h\le\mathsf oL^2M\).
Thus \cref{II-38-61} applies with
\((\alpha_0,\alpha_1,\alpha_2)=(a_{\rm w}/2,\bar a_1,\bar a_0)\);
its defect is at most \(u_{\rm G}(LM)^{-3}/2\).

The slices used here are open. To see the point that is not contained
in the simple inequalities, fix the real factor \(U\).
Its own tower conditions are either satisfied or the slice is empty.
Under \(\bar a_0+\bar a_1\le a^\sharp\), the continued tower
configurations and currents are holomorphic in \(A'\) on the
simple-domain slice, by the tower continuation theorem.
There are only finitely many subdomains, depths and stencils.
The extra core conditions are strict inequalities on these continuous
functions; their intersection with \(|\mathbf J|<\bar a_0\) is open.

The Gaussian structure's holomorphy, restrictions and whole
covariance are \cref{II-38-61}; pair dependence is built into
its coefficients. For completeness, its covariance follows
through the construction. At each local tower, a unitary
gauge conjugates the finite operators at the initial points
of their bond coordinates. Products, bilinear transposes,
compressions and inverses preserve this rule. The connected
coefficient sum uses the same subsets and signs; the mesh
resolvent products use scalar cutoffs and the same integration
variable. Hence both coefficient families, their absolutely
convergent parameter series, and finally \(\mathbf C\) and
\(\boldsymbol\Gamma\) have the required conjugation rule.
Complex conjugation is the identical calculation with the
antilinear map \(A\mapsto A^*\) and conjugated parameters;
translations relabel the coordinates and cube indices.
These maps preserve the real pairings, so they give precisely
the covariance relations used in the Gaussian expansion.
The gauge-fixing operator and the three potentials
are those of \cref{II-41-71,II-41-79}. Their domain contains the core by
\eqref{II-127-27}, and
\(\mathsf C_{\mathcal A}\delta(\mathsf t)\le\varepsilon_P\)
puts the entire field polydisc in their domain.
Since \(\lambda_V\le\bar\kappa-2\), their estimate gives
\eqref{II-53-14}. The same positive-height constant works at
\(k=0\): replacing \(\varepsilon_{\rm A}\) by the larger
radius \(r_P^\circ\) only decreases that bound.
The fourth-potential theorem supplies its holomorphy, dependence
and vanishing, and \eqref{eq:actual-fourth-size} gives
\(\nu(\mathbb V'[\mathfrak g])\le
N_{\mathfrak a,\omega}(\mathfrak g)\) when \(\nu\) is defined
with \(\mathsf K\).
The ceiling \eqref{eq:actual-step-ceiling} is precisely what
checks \(4\mathsf K N_*s_{\mathfrak a}\le w_*\).
We can therefore apply \cref{II-53-39} with the actual data.

The normalizer bound is
\[
 |\mathbf N^{(k+1)}(X,P,z)|
       \le E_Ne^{-\delta_NMd}.
\]
Add this to \eqref{eq:step-G-bound} and
\eqref{eq:step-F-bound}. Both decay rates are at least
\(\kappa+1\), and \(E_N\le E_0/2\), proving
\eqref{II-128-32}. The two final numerical conclusions follow
by inserting \eqref{II-128-6}; at \(\tau=1\) the G term is zero.
No additional estimate is needed for this rate change.

For the whole-torus assertion, the actual potential sums are
the Jacobian, gauge and action remainders in
\cref{II-40-99}; the action identity uses the Hessian identified
in \cref{sec:small-field-measure}. The fourth sum is \(s_W[S]\)
by hypothesis. Hence \eqref{II-53-67} is exactly the integral
in \eqref{II-40-100}. By \eqref{II-44-40},
\[
 Z^{(k)}(W)=Z_{\rm triv}^{(k)}
       \exp\sum_{X,z}\mathbf N^{(k+1)}(X,P_W|_X,z).
\]
This proves \eqref{eq:actual-step-density}. Finally
\eqref{II-43-100} gives
\(\sum\mathbf G_\tau=n_{\mathcal O}\operatorname{Log}\tau/2\);
since \(3\#\mathfrak X^+=n_\perp+n_{\mathcal O}\), it gives
the equivalent expression. Absolute convergence is the one
already proved for the actual pushforward and in
\cref{II-53-39}.
\end{proof}

\subsection{Consistency on non-wrapping domains}
\label{sec:step-chart-consistency}
We require a common input array on different tori to have the
same slot and coefficient numbers, and equal local terms after
every zero-translation chart on a common bounded domain.
These conditions concern the actual input functions, not only
their norms. They imply that the fourth potentials agree under
these charts, by their local substitution construction.

The core conditions also agree under a chart. Here is the
geometric and analytic check. Cubes and their wall adjacency
are preserved, so subdomains \(Z\subset X\) correspond.
In constructing a tower, test the distance from a cube to the
complement by containment of its open \(R\)-neighborhood.
All radii used are at most the halo width: at level \(p\le t\)
they are \(M_1L^{q+1}\) or \(R_1\mathfrak ML^q\), and in
level-\(t\) units they are at most \(M\).
Thus the chart carries each stage of either tower to the
corresponding stage. The two-layer interiors are carried along
as well, since their radius is \(5ML^{p-1}/2\le ML^t\)
in bare units. The reduced determining sets, their nodes,
blocks and ordered comb contours therefore correspond.

The one-step averaging formulas then have identical matrix
arguments, including their domain tests and their logarithm
branches. Induction gives equality of every actual iterated
average. Relabelling also preserves the action, the regularity
inequalities and the residual gauge group. It consequently
identifies the real critical orbits and their tower-axial
representatives. A continuation on a polydisc is carried to a
holomorphic continuation on the identical coordinate polydisc;
the two agree at real points and hence everywhere by uniqueness.
The radius is unchanged. This proves
\begin{equation}\label{II-17-45}
 \mathsf i_*\mathcal U^{\circ,\mathfrak w'}_j(X;b_0,b_1,c_0)
  =\mathcal U^{\circ,\mathfrak w''}_j
                      (\mathsf iX;b_0,b_1,c_0),
\end{equation}
including all the strict tower-current conditions. The same
argument applies to the simple domains and gives a complex
linear bijection between corresponding slices.

For the Gaussian data it is enough to compare parameters equal
to zero or one. On each grain subregion \(Y\), the preceding
tower argument identifies the reduced rows, pivot bonds,
ordered contours and inner bond spaces. Their finite
differential and averaging matrices are relabelled by the
chart. Bilinear transposes and inverses are therefore
relabelled as well, so the local fluctuation forms agree.
The chart bijects the subsets of a grain region and preserves
their cardinalities. It consequently preserves every term
and sign of the finite connected-coefficient sum. At a pair
of bonds outside its indicated support both coefficients
vanish; applying the inverse chart checks the support in
both directions.

The same argument carries mesh strings and their figures.
Their cutoff factors, local resolvents and commutators
correspond, so the products correspond for every positive
resolvent parameter. Absolute convergence permits integration
with the inverse-square-root weight and summation over
strings. The parameter monomial is unchanged. Compression,
inversion and the product in \eqref{II-38-15} now give, for
every parameter subset
\(\mathcal K\),
\begin{equation}\label{II-54-23}
\begin{split}
 \mathbf C''(\mathsf iZ_\circ,\mathbf1_{\mathsf i\mathcal K})
   &=\mathsf i_*\mathbf C'(Z_\circ,\mathbf1_{\mathcal K}),\\
 \boldsymbol\Gamma''(\mathsf iZ_\circ,\mathbf1_{\mathsf i\mathcal K})
   &=\mathsf i_*\boldsymbol\Gamma'(Z_\circ,\mathbf1_{\mathcal K}).
\end{split}
\end{equation}
The square-root branch is the same positive-real continuation.
The actual block fields, gauge operators and potentials likewise
agree, as do the normalizer terms by \cref{II-44-37}.

We spell out why this suffices for the connected terms.
The coordinate map on the positive bond fields is a real
isometry and preserves Lebesgue measure. Under it the
conditioning integrands have the same Gaussian determinant,
quadratic pairings and cutoff tests. Their absolute integrals
are therefore equal, and so are the integrals themselves.
This includes the zero-dimensional core. The chart bijects
the Mayer indices and the subsets used in the finite
differences. It preserves incompatibility and hence the
connected graph coefficients. Corresponding terms of the
refined connected series agree, and the series are absolutely
summable.

For the interaction insertion the same assertion holds at
every interpolation parameter \(v\in(0,1)\): scaling the fields
and differentiating the indicated bond coordinate commute
with relabelling. For the cutoff insertion compare the
functions of its varied bond threshold before differentiating.
The deciding cutoff is the same on both tori; where it fails
both integrands are zero, so a potential outside its domain
is never evaluated. The common integrable majorants in
\cref{II-53-39} permit the interpolation integrals.
Finally \(h_{\mathsf iz}(\mathsf ib)=h_z(b)\), and the
nonzero tent summands lie in the halo. Thus
\begin{equation}\label{II-54-59}
 E^{(k+1),\mathfrak w''}(\mathsf iX,\mathsf i_*P,\mathsf iz)
       =E^{(k+1),\mathfrak w'}(X,P,z).
\end{equation}
The same identity holds for each of its three summands.

Let now the sides of the chosen family \(\mathfrak W\) be
unbounded, with one common input array in the preceding sense.
Every bounded domain \(\widehat X\) embeds in some member.
By \eqref{II-17-45} its core is the pullback of the core of
any such member, independent of the embedding. Define
\begin{equation}\label{II-54-79}
 \widehat E^{(k+1)}(\widehat X,\hat P,\hat z)
  =E^{(k+1),\mathfrak w}
       (\operatorname{pr}^{\mathfrak w}\widehat X,
        \operatorname{pr}^{\mathfrak w}_*\hat P,
        \operatorname{pr}^{\mathfrak w}\hat z).
\end{equation}
Equation \eqref{II-54-59} makes this well defined and unique.
It is a pointed local family: on each finite domain its
holomorphy, pair dependence and gauge invariance are those
of one finite-dimensional evaluation. It has the same
conjugation and translation properties when the input
arrays do. These statements remain valid for holomorphic
input families under uniform size bounds.

The size of a bounded domain is preserved by the chart.
Indeed, the chart and its inverse map cubes and faces of
each dimension bijectively. They therefore map the connected
edge sets meeting every cube onto each other, preserve the
number of edges, and also preserve the zero-edge case of
a single common corner. Taking the least edge count gives
equality of the two sizes. Therefore \eqref{II-128-32} passes to the infinite
model. More precisely, put
\[
 \widehat N_{\mathfrak a}(\widehat X)
   =\inf_{\mathfrak w:\,\widehat X\ {\rm near\ for}\ \mathfrak w}
                  N_{\mathfrak a,\omega}(\mathfrak g^{\mathfrak w}).
\]
This is the infimum of a nonempty set in \([0,N_*]\).
The left side of the bound is independent of the selected
embedding and the right side is affine increasing in this
size. Taking the infimum proves
\begin{equation}\label{II-128-33}
\begin{split}
 |\widehat E^{(k+1)}(\widehat X,\hat P,\hat z)|
 \le\Bigl[\frac{E_0}{2}+K_{\mathbf M}(\mathsf K)
 \bigl(&|1-\tau|+\delta(\mathsf t)
       +s_{\mathfrak a}\widehat N_{\mathfrak a}(\widehat X)\\
       &+e^{-\gamma_2\mathfrak b(\mathsf t)^2/2}\bigr)\Bigr]
                         e^{-(\kappa+1)d_{k+1}(\widehat X)} .
\end{split}
\end{equation}
The required unbounded family exists at every height:
take \(K=k+1\) and let the integer \(m\) increase until
\(2L^{m+1}\ge C_LM\), then continue increasing it.


\subsection{The quadratic coefficient of the local output}
\label{sec:local-output-coefficient}

We extract the quadratic field-strength term from the output of a step.
The estimate must be uniform in the level, whereas the neighbourhood on
which a finite-volume critical representative is defined may depend on
the volume. We first distinguish these two requirements.

Measure lengths at level \(j\ge1\), so that the coarse lattice has
spacing one and the fine lattice has spacing \(\xi=L^{-j}\).
Write \(D\) for the fine curl, \(d=\nabla^\xi\) for the site gradient,
and \(\Delta=d^*d\). The straight averages are
\[
 (Qh)(c)=\xi L^{-3j}\sum_{x\in B^j(c_-)}
                       \sum_{b\in[x,c_+-c_-+x]}h(b),\qquad
 (Q'\lambda)(y)=L^{-3j}\sum_{x\in B^j(y)}\lambda(x).
\]
Fine pairings have weight \(\xi^3\); coarse pairings have weight one.
Thus \(Qd=\partial Q'\), with \(\partial\) the coarse site gradient.
Let \(\mathsf H_N\) be the free Landau response on the torus of side
\(N=2L^a\). Its column \(h^N_{c,Y}\), together with its multiplier
\(\omega^N_{c,Y}\), solves
\begin{equation}\label{I-82-31}
 D^{\mathsf T}Dh=Q^{\mathsf T}\omega,\qquad Qh=Y\delta_c,\qquad
 (d^*h,\Delta\lambda)_\xi=0
\end{equation}
for every \(\lambda\) supported in one fine block and having zero
sum there. Choose a fixed lower bound \(N_{\mathsf H}\) for the side
so that the constant-sequence geometry of \cref{I-82-15} is admitted.
Increase it to meet the torus condition in \cref{I-58-22} and
\(N_{\mathsf H}\delta_{\mathsf H}\ge1\).
It depends only on the fixed geometric constants, not on \(j\).
The corresponding infinite-lattice response is denoted by
\(\mathsf H^{[j]}\). Below we prove its existence, periodization, and
the uniform bound
\begin{equation}\label{II-56-6}
 |\mathfrak k^c_\mu(x)|,\quad
 |\nabla^\xi_\nu\mathfrak k^c_\mu(x)|,\quad
 |(D^{\mathsf T}D\mathfrak k^c)_\mu(x)|
       \le C_{\mathsf H}e^{-\delta_{\mathsf H}|x-c_-|_\infty},
 \qquad
 (\mathsf H^{[j]}B)(b)=\sum_{c>0}\mathfrak k^c(b)B(c).
\end{equation}
The kernel is real, and \(C_{\mathsf H}\ge1\) and
\(0<\delta_{\mathsf H}\le1\) are common to every level.

For a localization domain \(X\), let
\(\mathfrak V(X)\) be the space of independent pairs
\(v=(A,\mathbf J)\) on its incident fine bonds, with norm
\[
 \|v\|_X=\max\{\sup|A|,\sup|\nabla^\xi A|,\sup|\mathbf J|\}.
\]
Only differences defined on those bonds enter the maximum.
The trivial-pair chart is \(P[v]=(I,A,\mathbf J)\). The core estimate
\cref{I-93-98} gives a uniform usable radius
\begin{equation}\label{II-17-20}
 r_0=a^{\mathrm{us}}(b_0,b_1)
   :=\min\{b_1,a^\sharp/4,b_0/(4K^\sharp_{\mathrm{tw}})\},
 \qquad
 P[\mathfrak V_{r_0}(X)]\subset
                   \mathcal U^\circ_j(X;b_0,b_1).
\end{equation}
We verify this inclusion in the proof. In particular the reading
\(\mathsf e_{X,z}(v)=E(X,P[v],z)\) of a local term is holomorphic
on this ball and is invariant whenever two points of the ball
represent gauge-equivalent pairs.
Primes on a reading denote derivatives at \(v=0\).

Fix \(M=L^{m'}\), \(m'\ge6\), and an unbounded family
\(\mathfrak W\) of tori carrying level \(j\), all with
level-\(j\) side \(N\ge N_{\mathsf H}\). At a common height
\(h\ge j-1\), require \(K\ge j\) and the room and near-domain
chart hypotheses of \cref{II-17-45}; the constant sequences must
also satisfy the hypotheses of \cref{I-82-15,I-58-22}, as fixed
in the choice of \(N_{\mathsf H}\).
Use the background constants fixed for the cores, in particular
\(0<\varepsilon_{\mathrm v}\le a_{\rm orb}\).
Let \(E^{\mathfrak w}\) be pointed local families on these tori and
\(\widehat E\) their companion on the infinite model, at the same
positive core radii \(b_0,b_1\). Suppose that
\begin{equation}\label{eq:coefficient-common-size}
 |E^{\mathfrak w}(X,P,z)|\le\mathcal A e^{-\kappa d_j(X)},
 \qquad
 |\widehat E(\widehat X,P,\hat z)|
                  \le\mathcal A e^{-\kappa d_j(\widehat X)}
\end{equation}
on their full cores. The hypotheses for the core and chart comparison
in \cref{I-93-98,II-17-45} are understood, with the same constants
at every height. Take
\(\kappa\ge\kappa_{\mathrm{am}}:=4\log9216+9\).
For a real block field \(W\), define the total, whenever its
minimal background orbit exists and all displayed pairs lie
in the term domains, by
\begin{equation}\label{II-16-38}
 \mathcal P[E^{\mathfrak w}](W)
  =\sum_z\sum_{X\ni z}E^{\mathfrak w}
       \bigl(X,(U_j(W),J^{(j)}(U_j(W)))|_{\operatorname{Bd}(X)},z\bigr).
\end{equation}
Here \(U_j(W)\) denotes the canonical axial representative,
and a local term is read as a function of its pair.
For the near-trivial totals used below, the usable-ball
argument also proves that all representatives are in the term
domains, and hence proves representative-independent definedness.
For \(B\) of polynomial growth put
\[
 \mathfrak h_X B
   =\bigl((\mathsf H^{[j]}B)|_{\operatorname{Bd}(X)},
          (D^{\mathsf T}D\mathsf H^{[j]}B)|_{\operatorname{Bd}(X)}\bigr).
\]
With \(t_1,t_2,t_3\) an orthonormal Pauli basis, define
\begin{equation}\label{II-56-15}
 \Pi(c,c')
 =\frac13\sum_{a=1}^3\sum_{\hat z\in\mathbb Z^3}
       \sum_{\widehat X\ni\hat z}
       \hat{\mathsf e}_{\widehat X,\hat z}''
          [\mathfrak h_{\widehat X}(t_a\delta_c),
           \mathfrak h_{\widehat X}(t_a\delta_{c'})].
\end{equation}
All bond sums in this subsection use positive orientations; the kernel
extends oddly in either bond.

\begin{theorem}\label{II-56-86}\label{II-129-73}
Under these hypotheses the series \eqref{II-56-15} is absolutely
convergent. Put
\(\delta_{\mathrm{am}}=\frac12\min(\delta_{\mathsf H},\kappa/M)\).
There is \(C_{\mathrm{am}}\), depending only on
\(C_{\mathsf H},\delta_{\mathsf H},\kappa,M\), such that
\begin{equation}\label{eq:coefficient-amplitude}
 |\Pi(c,c')|\le C_{\mathrm{am}}r_0^{-2}\mathcal A
                 e^{-\delta_{\mathrm{am}}|c'_--c_-|_\infty},
 \qquad
 \sum_{c'>0}|\Pi(c,c')|(1+|c'_--c_-|_\infty)^2
                  \le C_{\mathrm{am}}r_0^{-2}\mathcal A.
\end{equation}
On every finite torus \(\mathfrak w\in\mathfrak W\), the total \(\mathcal P[E^{\mathfrak w}]\)
is defined and analytic near the trivial coarse field. For fixed
infinite-lattice bonds \(c,c'\),
\begin{equation}\label{II-129-74}
 \lim_{\mathfrak w\in\mathfrak W,\ N\to\infty}
 D^2\bigl[\mathcal P[E^{\mathfrak w}](e^{iB})\bigr]_{B=0}
 [t_a\delta_{\mathrm{pr}_jc},t_{a'}\delta_{\mathrm{pr}_jc'}]
                  =\delta_{aa'}\Pi(c,c').
\end{equation}
The neighbourhood in this assertion need not be uniform in
\(\mathfrak w\).

If the finite-volume totals are invariant under lattice isometries
near the trivial field, then \(\Pi(c,c')=\Pi_{\mu\nu}(u)\), where
\(u=c'_--c_-\), and its moments of orders zero and one vanish.
Its second moments have the form
\begin{equation}\label{II-129-75}
 \sum_{u\in\mathbb Z^3}\Pi_{\mu\nu}(u)u_\rho u_\sigma
   =\mathsf{cf}(\widehat E)
       (\delta_{\mu\rho}\delta_{\nu\sigma}
        +\delta_{\mu\sigma}\delta_{\nu\rho}
        -2\delta_{\mu\nu}\delta_{\rho\sigma}),
\end{equation}
where
\begin{equation}\label{II-56-18}
 \mathsf{cf}(\widehat E)
    =\sum_u\Pi_{\mu\nu}(u)u_\mu u_\nu\quad(\mu\ne\nu),
 \qquad
 |\mathsf{cf}(\widehat E)|
               \le C_{\mathrm{am}}r_0^{-2}\mathcal A.
\end{equation}
All these moment series converge absolutely. The same coefficient
is obtained from any finite-volume functionals equal to the totals
up to constants near the trivial field. If those functionals are
real up to constants, the kernel and the coefficient are real.
\end{theorem}

The proof separates the uniform linear response from the two summations
in \eqref{II-56-15}. Gauge invariance removes the second derivative
of the background and leaves only its linear response.
Exponential decay then controls both the large-domain tail and the
periodic images.

\subsubsection*{The free response}

We give the infinite-volume argument, retaining only the three bounds
in \eqref{II-56-6}. The finite response is the derivative at zero in
\cref{I-82-15}. At the trivial field the nonlinear constraint remainder
and action remainder have zero derivative. Differentiating their
equations gives \eqref{I-82-31}; the Landau normalization gives its
last equation. The estimate \eqref{I-93-47}, and the multiplier
estimate in its proof, give uniform exponential bounds for
\(h,\nabla^\xi h,D^{\mathsf T}Dh,\omega\). Here every layer site is
at the single coarse scale, and the distance is the coarse
\(\ell^1\) distance. The base point of a fine bond is within \(1/2\)
in sup distance of its block centre. Decreasing the decay rate and
increasing the bound therefore gives \eqref{II-56-6}, with constants
independent of \(j\).

For clarity, uniqueness uses the local gauge condition, not a choice
of a limiting sequence. For a real homogeneous solution, pairing
the first equation with \(h\) gives \(\|Dh\|_\xi^2=0\).
On a torus the last equation is equivalent to \(Rd^*h=0\), since
the blockwise mean-zero functions span \(\ker Q'\). Thus \(h\)
lies in the kernel of the positive free operator
\(D^{\mathsf T}D+dRd^*+Q^{\mathsf T}\mathsf a Q\), with the
positive row weight of \cref{I-17-5}, and is zero.
Also \(Q^{\mathsf T}\omega=0\) implies \(\omega=0\):
at the fine bond crossing the central outgoing face of \(B^j(c_-)\)
in the direction of \(c\), only the row \(c\) contributes, and its
weighted transpose coefficient is one. Complex uniqueness follows
by separating real and imaginary parts.

On the infinite lattice consider an exponentially decaying homogeneous
solution. The same pairing gives \(Dh=0\). Set
\[
 \phi(x)=-\xi\sum_{n\ge0}h_1(x+n\xi e_1).
\]
Telescoping and \(Dh=0\) give \(d\phi=h\). The full line sum
\(\xi\sum_{n\in\mathbb Z}h_1(x+n\xi e_1)\) is constant in every
coordinate, again by telescoping, and is zero by transverse decay.
Thus \(\phi(x)=\xi\sum_{n<0}h_1(x+n\xi e_1)\) as well. Using
the first expression for \(x_1\ge0\) and the second for \(x_1<0\)
shows that \(\phi\) decays exponentially.
Since \(\partial Q'\phi=Qd\phi=0\), decay gives \(Q'\phi=0\).
Decompose \(\phi\) into its restrictions to the fine blocks.
Each restriction is an allowed test function in \eqref{I-82-31};
absolute summability permits their sum. Consequently
\[
 \|\Delta\phi\|_\xi^2
  =\sum_y(d^*h,\Delta(\phi\mathbf1_{B^j(y)}))_\xi=0,
 \qquad
 \|d\phi\|_\xi^2=(\phi,\Delta\phi)_\xi=0.
\]
Therefore \(h=0\), and the same crossing-face argument gives
\(\omega=0\).

Lift finite-torus columns to the infinite lattice and let \(N\to\infty\).
Their values at the countably many bonds and rows are uniformly bounded,
so a diagonal subsequence converges pointwise. Every equation and every
blockwise gauge test involves finitely many values and passes to the
limit. At a fixed row the periodic source is eventually \(Y\delta_c\).
The uniform exponential bounds pass to the limit too. The uniqueness
just proved identifies all subsequential limits. Real linearity,
translation covariance and the fact that the coefficients are scalar
in colour give a real kernel \(\mathfrak k^c\).

Absolute convergence allows periodization through the local operators
and through a blockwise gauge test. The periodized column solves the
finite-torus system, so uniqueness gives
\begin{equation}\label{eq:coefficient-response-periodization}
 h^N_{\mathrm{pr}_jc,Y}(\mathrm{pr}_jb)
      =\sum_{n\in N\mathbb Z^3}\mathfrak k^c(b+n)Y.
\end{equation}
This holds after applying \(\nabla^\xi\) or \(D^{\mathsf T}D\).
The sum over nonzero images is bounded by a convergent geometric shell
sum. In particular finite columns satisfy the counterpart of
\eqref{II-56-6} with periodic distance and a fixed larger constant,
and converge locally to the infinite columns.

We shall also use the gradient property of this response. If a coarse
site function \(\lambda\) has finite support, let \(\lambda^\flat\)
be constant with value \(\lambda(y)\) on \(B^j(y)\). Then
\(Qd\lambda^\flat=\partial\lambda\). For
\(E=\mathsf H^{[j]}\partial\lambda-d\lambda^\flat\) one has
\(QE=0\) and \(D^{\mathsf T}DE=Q^{\mathsf T}\omega\).
Pairing gives \(DE=0\). For polynomially growing \(\lambda\),
regroup the absolutely convergent kernel sum by coarse sites:
the coefficient at \(y\) is the response to \(\partial\delta_y\)
and is curl-free by the finite-support case. Applying the four-term
curl to the sum proves
\begin{equation}\label{II-129-31}
 D\mathsf H^{[j]}\partial\lambda=0,\qquad
 Q\mathsf H^{[j]}B=B.
\end{equation}
The second identity follows first for a single column and then by
absolute summation. A curl-free form on the infinite cubic lattice
has a primitive: sum along the three coordinate segments, and use
each plaquette relation to interchange adjacent coordinate steps.
Thus \(\mathsf H^{[j]}\partial\lambda=d\theta\), and its
\(D^{\mathsf T}D\) image vanishes. No decay of this primitive is
needed for a polynomial \(\lambda\).

\subsubsection*{Local derivatives and summation}

\begin{proof}[Proof of \cref{II-56-86,II-129-73}]
We first verify \eqref{II-17-20}. Write \(r=r_0\).
The fixed constants have \(K^\sharp_{\mathrm{tw}}\ge1\) and
\(a^\sharp\le1/60\), so \(r\le1/240\). If \(\|v\|_X<r\),
the flat real reference admits the identity gauge, and the potential,
its first differences and the independent current have size less
than \(r\). The only additional simple-domain condition is the
plaquette bound. With the four positive link potentials of a
plaquette written as \(Y_i=i\xi A(b_i)\), one has
\[
 |Y_i|\le s:=\xi r,\qquad
 |Y_3-Y_1|,\ |Y_4-Y_2|\le\xi^2r.
\]
Subtract \(e^{-Y_1}e^{-Y_2}\) from
\(e^{-Y_3}e^{-Y_4}\), and commute the two remaining factors.
The two difference terms are bounded by \(e^{2s}\xi^2r\);
the commutator is bounded by \(2e^{2s}\xi^2r^2\).
After multiplication by \(e^{Y_1}e^{Y_2}\), this gives
\[
 |\partial(e^{i\xi A})(p)-I|
       \le2e^{4s}(1+r)r\xi^2<3r\xi^2.
\]
Thus \(P[v]\in\mathcal Q_j(X;3r,r,b_0)\). The inequalities
\[
 3r\le b_0,\qquad r\le b_1,\qquad
 K^\sharp_{\mathrm{tw}}(3r+r^2)\le4K^\sharp_{\mathrm{tw}}r\le b_0,
 \qquad4r\le a^\sharp
\]
give \eqref{II-17-20} by \eqref{I-93-100}.
For an infinite domain use any sufficiently large near lift.
The chart comparison \cref{II-17-45} preserves these conditions
and the norm in both directions.

The same core inclusion holds for a unitary gauge transform
of such a triple. Indeed, all curvature, current and covariant
difference bounds defining \(\mathcal Q_j(X;3r,r,b_0)\)
are unchanged by unitary conjugation. A cube gauge \(w\) for
the original real factor is replaced by \(wu^{-1}\), with the
same potential and bounds, for its transform by \(u\).
Thus the transformed triple still lies in this simple domain,
and \eqref{I-93-100} puts it in the same core. Whenever
\(P[v]^u\) and \(P[w]\) represent the same pair, with
\(v,w\) in the usable ball, gauge invariance and pair dependence
of the local term therefore give \(e(w)=e(v)\).

For each reading \(e=\mathsf e_{X,z}\), Cauchy's inequality on
the bidisc \(|s_i|<r_0/(2\|v_i\|_X)\) gives
\begin{equation}\label{II-56-31}
 |e''[v_1,v_2]|
  \le4r_0^{-2}\mathcal A e^{-\kappa d_j(X)}
                  \|v_1\|_X\|v_2\|_X.
\end{equation}
Constant gauge rotations preserve this ball. Their invariance
implies \(e'=0\), and makes \(e''\) scalar in colour:
rotations by \(\pi\) about the Pauli axes kill every invariant
linear form and all off-diagonal entries of an invariant bilinear
form; rotations permuting the axes make its diagonal entries equal.
Consequently
\[
 e''[Yf,Y'f']=\langle Y,Y'\rangle e''[t_1f,t_1f']
\]
for scalar vectors \(f,f'\).

The Hessian also kills local gauge directions. For a real site
function \(\theta\) on the endpoints of the incident bonds, set
\[
 u_t=e^{it\theta},\qquad
 w(t,s)=\left(
 \frac1{i\xi}\log\bigl(u_t(b_-)e^{i\xi sA(b)}u_t(b_+)^{-1}\bigr),
 \ s\operatorname{Ad}(u_t(b_-))\mathbf J(b)\right).
\]
There are finitely many bonds, so the series logarithms and both
points \(w(t,s),sv\) lie in the ball for sufficiently small \(s,t\).
Gauge invariance gives \(e(w(t,s))=e(sv)\).
At zero, \(\partial_sw=v\) and \(\partial_tw=-(d\theta,0)\).
The mixed chain-rule term involving \(e'\) vanishes. Thus
\begin{equation}\label{eq:coefficient-local-Ward}
 e''[(d\theta,0),v]=0 .
\end{equation}
Splitting a complex \(\theta\) into its real and imaginary parts
extends the identity complex linearly.

Let \(\rho_w(X)=\inf_{x\in X}|w-x|_\infty\), with periodic
distance on a torus. A base point of an incident bond lies within
\(\xi\) of \(X\). The three response estimates therefore give
\[
 \|\mathfrak h_X(Y\delta_c)\|_X
       \le e^{1/13}C_{\mathsf H}e^{-\delta_{\mathsf H}\rho_{c_-}(X)}|Y|.
\]
Periodization gives the same estimate on a torus with a fixed
larger constant. The shell count
\(\#\{u\in\mathbb Z^3:|u|_\infty\le n\}=(2n+1)^3\)
shows that the response series for a polynomially growing input
converges in every local norm. It may be differentiated by the
finite-difference operators and inserted in any local bilinear
form term by term.

Here are the domain sums needed for both the kernel and its limit.
An edge set realizing \(d_j(X)\), together with a segment in each
endpoint cube, gives
\(\operatorname{diam}_\infty X\le M(d_j(X)+2)\).
The walk count in \cref{II-18-26} is unchanged on the infinite
lattice. A site belongs to at most eight cubes; hence
\(\kappa/4\ge\log9216+9/4\) implies
\[
 \sum_{\substack{X\ni z\\d_j(X)\ge D}}
                  e^{-\kappa d_j(X)/4}\le C_{\mathrm E,0}e^{-D}.
\]
For real \(D\ge0\), sum the geometric series starting at
\(\lceil D\rceil\). These bounds also hold for wrapping domains
on the torus. Put
\[
 T_{w,w'}(X)
   =e^{-\kappa d_j(X)}
      e^{-\delta_{\mathsf H}(\rho_w(X)+\rho_{w'}(X))},
 \qquad
 S_p(a)=\sum_{u\in\mathbb Z^3}(1+|u|_\infty)^p e^{-a|u|_\infty}.
\]
All \(S_p(a)\) used below are finite by the shell count.
For \(z\in X\), the diameter bound gives
\[
 |w-z|_\infty\le\rho_w(X)+M(d_j(X)+2),\qquad
 |w-w'|_\infty\le\rho_w(X)+\rho_{w'}(X)+M(d_j(X)+2).
\]
Since \(\delta_{\mathrm{am}}\le\delta_{\mathsf H}/2\) and
\(M\delta_{\mathrm{am}}\le\kappa/2\), distributing these two
inequalities in the exponential yields
\begin{align}
 T_{w,w'}(X)
 &\le e^\kappa
 e^{-\delta_{\mathrm{am}}(|w-z|_\infty+|w'-z|_\infty)/2}
 e^{-\kappa d_j(X)/2},\label{eq:coefficient-two-point-majorant}\\
 T_{w,w'}(X)
 &\le e^{3\kappa/2}e^{-\delta_{\mathrm{am}}|w-w'|_\infty}
       e^{-\delta_{\mathrm{am}}|z-w|_\infty/2}
       e^{-\kappa d_j(X)/4}.\label{eq:coefficient-domain-majorant}
\end{align}
The first inequality gives a product of two summable endpoint
factors after summation over \(X\ni z\). The second gives
\begin{equation}\label{eq:coefficient-domain-tail}
 \sum_z\sum_{\substack{X\ni z\\d_j(X)\ge D}}T_{w,w'}(X)
 \le C e^{-\delta_{\mathrm{am}}|w-w'|_\infty}e^{-D}.
\end{equation}
Here and below \(C\) depends only on
\(C_{\mathsf H},\delta_{\mathsf H},\kappa,M\).
For the torus, choose one nearest lift of each site; its exponential
sum is bounded by the corresponding \(S_0\).

Combining \eqref{II-56-31} with these estimates proves absolute
convergence of \eqref{II-56-15} and its exponential bound.
Summing the bound over the three directions at \(c'_-=c_-+u\),
with weight \((1+|u|_\infty)^2\), proves
\eqref{eq:coefficient-amplitude} after enlarging \(C_{\mathrm{am}}\).
The same calculation with any polynomial weight is finite.
For a fixed base point \(z\), the two endpoint factors in
\eqref{eq:coefficient-two-point-majorant} also justify all double
response sums with polynomially growing inputs.

\subsubsection*{The finite-volume Hessian and its limit}

Apply \cref{I-58-22} to the constant full-torus sequence and trivial
datum. There are no lower-layer rows or inward rows.
The auxiliary geometric parameters were fixed before \(N_{\mathsf H}\).
Choose the datum bound and then the row ball small enough that
the regular orbit lies in the selected \(\varepsilon_{\mathrm v}\)
class. Its axial representative has a holomorphic continuation
\(\Phi_I(B)\) near zero, with \(\Phi_I(0)=I\).
All bonds are inner. On a smaller finite-volume neighbourhood set
\[
 A(B)=\frac1{i\xi}\log\Phi_I(B),\qquad
 \mathbf J(B)=J(\Phi_I(B)).
\]
These maps are holomorphic and vanish at zero. There are only
finitely many domains and bonds, so their restrictions
\(w_X(B)=(A(B),\mathbf J(B))|_{\operatorname{Bd}(X)}\)
lie in \(\mathfrak V_{r_0}(X)\) on one common neighbourhood.
The background belongs to the actual minimal orbit; independence
of its representative therefore identifies the total with
\[
 \Psi_N(B)=\sum_z\sum_{X\ni z}\mathsf e^N_{X,z}(w_X(B)).
\]
This finite sum proves the asserted analyticity. No uniform lower
bound on this neighbourhood is needed.

The critical-branch calculation in \cref{sec:gaussian-coefficient},
converted from absorbed links to the present potential units, gives
\[
 A'(0)(Y\delta_c)=h^N_{c,Y}-d\theta_{c,Y},
 \qquad
 \mathbf J'(0)(Y\delta_c)=D^{\mathsf T}Dh^N_{c,Y}.
\]
That calculation already includes the endpoint correction between
the original and straightened means. The vanishing of \(e'\)
removes the second derivative of \(w_X\), and
\eqref{eq:coefficient-local-Ward} removes \(d\theta_{c,Y}\).
It follows that
\begin{equation}\label{II-129-66}
 D^2\Psi_N(0)[Y\delta_c,Y'\delta_{c'}]
  =\sum_z\sum_{X\ni z}
       (\mathsf e^N_{X,z})''
       [\mathfrak v_Xh^N_{c,Y},\mathfrak v_Xh^N_{c',Y'}],
\end{equation}
where \(\mathfrak v_Xh=(h,D^{\mathsf T}Dh)|_{\operatorname{Bd}(X)}\).
This also holds for wrapping \(X\): \(\theta_{c,Y}\) is a site
function on the entire torus.

For a near domain \(\widehat X\) projecting to \(X\), the chart
identity for the companion holds throughout the usable ball and
preserves every norm stencil. Differentiating it twice identifies
the two local Hessians. We compare their sums without dropping
the wrapping terms.

Fix \(c,c'\) on the infinite lattice, write \(u=c'_--c_-\), and take
\[
 \mathrm{FD}_N
   =\{z\in\mathbb Z^3:-N/2\le(z-c_-)_\mu<N/2,\ \mu=1,2,3\}.
\]
Projection maps this set bijectively onto the coarse torus.
A non-wrapping pointed domain has exactly one near lift with its
point in \(\mathrm{FD}_N\). A far or wrapping domain has
\(d_j(X)\ge D_N:=\max\{0,N/(2M)-2\}\), by the diameter estimate.
On each near lift, periodization gives the absolutely convergent
bilinear expansion
\[
 (\mathsf e^N_{X,z})''
 [\mathfrak v_Xh^N_{\mathrm{pr}_jc,Y},
  \mathfrak v_Xh^N_{\mathrm{pr}_jc',Y'}]
 =
 \sum_{n,n'\in N\mathbb Z^3}
 \hat{\mathsf e}_{\widehat X,\hat z}''
 [\mathfrak h_{\widehat X}(Y\delta_{c-n}),
  \mathfrak h_{\widehat X}(Y'\delta_{c'-n'})].
\]
The terms \(n=n'=0\) give the infinite sum restricted to near
domains and base points in \(\mathrm{FD}_N\). There are four errors:
the wrapping part of \eqref{II-129-66}, the nonzero image pairs,
the infinite base points outside \(\mathrm{FD}_N\), and the far
infinite domains with base point inside it.
By \eqref{eq:coefficient-domain-tail}, the first and fourth errors
are at most
\(Cr_0^{-2}\mathcal A|Y||Y'|e^{-D_N}\).

For the other two, put
\[
 t_N(c'')=\sum_{z\notin\mathrm{FD}_N}
       e^{-\delta_{\mathrm{am}}|c''_--z|_\infty/2}.
\]
Splitting the exponential into two equal factors gives
\[
 t_N(c)\le S_0(\delta_{\mathrm{am}}/4)e^{-\delta_{\mathrm{am}}N/8},
 \quad
 t_N(c')\le S_0(\delta_{\mathrm{am}}/4)
       e^{-\delta_{\mathrm{am}}(N/2-|u|_\infty)/4}.
\]
The outside-base-point error is at most
\(Cr_0^{-2}\mathcal A|Y||Y'|t_N(c)\), by
\eqref{eq:coefficient-two-point-majorant}.
For a nonzero image pair, either \(n\ne0\) or \(n'\ne0\).
The map \((z,n)\mapsto z+n\) is a bijection from
\(\mathrm{FD}_N\times(N\mathbb Z^3\setminus\{0\})\) onto
\(\mathbb Z^3\setminus\mathrm{FD}_N\).
Sum that endpoint factor first; the other is bounded by
\(S_0(\delta_{\mathrm{am}}/2)\). Thus the image-pair error is at most
\[
 Cr_0^{-2}\mathcal A|Y||Y'|
 S_0(\delta_{\mathrm{am}}/2)\bigl(t_N(c)+t_N(c')\bigr).
\]
All four errors tend to zero for fixed \(c,c'\).
The colour identity proved above identifies the limiting form with
\(\langle Y,Y'\rangle\Pi(c,c')\), proving \eqref{II-129-74}.

\subsubsection*{Gauge cancellation and the moment tensor}

Suppose the finite totals are invariant under lattice isometries.
Their Hessians then satisfy the corresponding signed bond identity.
Pass to the limit in \eqref{II-129-74}. Reversing either bond changes
its sign, so the identity holds for all oriented bonds.
Translations give \(\Pi(c,c')=\Pi_{\mu\nu}(c'_--c_-)\).

For a polynomially growing coarse site function \(\chi\),
\eqref{II-129-31} makes \(\mathsf H^{[j]}\partial\chi\) a gradient,
with zero \(D^{\mathsf T}D\) image. Every local Hessian annihilates
this vector. The summability just proved permits the response and
domain sums to be interchanged, giving the Ward identity
\begin{equation}\label{II-129-76}
 \sum_{\nu,u}
     [\chi(u+e_\nu)-\chi(u)]\Pi_{\mu\nu}(u)=0 .
\end{equation}
Taking \(\chi(u)=u_\nu\) proves the vanishing of every zeroth moment.

To handle reflections, use the difference of bond midpoints,
\(\tilde u=u+(e_\nu-e_\mu)/2\). If
\(re_\alpha=\varsigma_\alpha e_{\pi(\alpha)}\), the signed bond law is
\begin{equation}\label{II-129-77}
 \Pi_{\mu\nu}(u)
   =\varsigma_\mu\varsigma_\nu\Pi_{\pi(\mu)\pi(\nu)}(u'),
 \qquad
 u'+\tfrac12(e_{\pi(\nu)}-e_{\pi(\mu)})=r\tilde u.
\end{equation}
Indeed the positive representative of \(r\langle x,x+e_\mu\rangle\)
starts at \(r(x+e_\mu/2)-e_{\pi(\mu)}/2\).
This also shows that \(u\mapsto u'\) is an integer bijection.
For \(r=-I\) its midpoint argument is \(-\tilde u\), while the
two orientation signs multiply to one. Hence the first midpoint
moment vanishes. The zeroth moment already vanishes, so the first
moment in \(u\) vanishes too. For the same reason the second
moments in \(u\) and \(\tilde u\) agree.

Write these second moments as \(P_{\mu\nu,\rho\sigma}\).
Reflecting any one axis in \eqref{II-129-77} shows that a component
vanishes unless each index value occurs an even number of times.
Coordinate permutations leave just four possible values:
\[
 P_{\mu\mu,\mu\mu}=q_0,\quad
 P_{\mu\mu,\rho\rho}=q_1,\quad
 P_{\mu\nu,\mu\nu}=q_2,\quad
 P_{\mu\nu,\nu\mu}=q_3
 \qquad(\mu\ne\rho,\ \mu\ne\nu).
\]
Symmetry of the two moment indices gives \(q_3=q_2\).
Now take \(\chi(u)=u_\alpha u_\beta u_\gamma\) in
\eqref{II-129-76}. Its terms of degrees zero and one vanish
after summation, leaving
\begin{equation}\label{II-129-79}
 P_{\mu\alpha,\beta\gamma}
 +P_{\mu\beta,\alpha\gamma}
 +P_{\mu\gamma,\alpha\beta}=0.
\end{equation}
With all four indices equal this gives \(q_0=0\); with
\(\alpha=\mu\), \(\beta=\gamma=\rho\ne\mu\), it gives
\(q_1+2q_2=0\). Thus \(q_2\) is the coefficient in
\eqref{II-56-18}, independent of the distinct pair \(\mu,\nu\),
and \eqref{II-129-75} follows. Its bound is a part of
\eqref{eq:coefficient-amplitude}.
Finally, adding constants to the finite totals does not change
their Hessians. If they are real up to constants, those Hessians
are real; \eqref{II-129-74} and the absolutely convergent moment
sum give reality of the kernel and coefficient.
\end{proof}


\subsection{The global domain of the step}\label{sec:actual-step-global-domain}
The small-field expansion is an actual step on a uniformly small
class of block fields. We verify this before iterating it. The
verification also identifies the physical backgrounds in its local
potentials and in the quadratic subtraction.

Retain the common geometric hypotheses of
\cref{I-58-22,I-82-15,II-38-61,II-44-37,II-51-52}, at every height
under consideration. In particular, choose
\(B_3\ge B_{\rm rad}\), \(\rho_{\rm s}\ge\rho_{\rm G}\),
\(\mathfrak M\ge\mathfrak M_{\rm G}\),
\(R_1\mathfrak M\ge L^2\rho_{\rm s}\), and
\[
 \varepsilon_{\rm v}+\varepsilon_{\rm v}^2\le a_{\rm orb},\qquad
 \bar\varepsilon\le\varepsilon_{\rm G}(\rho_{\rm s}),\qquad
 r_{\rm B}\le r_{\rm G}(\rho_{\rm s}),\qquad
 B_3(\bar\varepsilon+C_\partial r_{\rm B})\le\varepsilon_{\rm v}.
\]
The level lattices have the size required by these results. We use
the same unbounded family of admitted tori as in
\cref{II-129-73}, including its floor \(N_{\rm H}\).
Increase the room constant to \(C_L\ge2\mathsf oL^2+13\).
Write \(c_{\rm rg},B_{\rm rg}\) for the constants of
\cref{I-28-75}, put \(l_M=\mathsf oLM+2\), and choose
\begin{equation}\label{eq:global-core-gauge-constant}
 B^{(\mathrm{cg})}\ge
       2B_{\rm rg}\left(1+\frac2{\mathsf oL^2}\right).
\end{equation}
All these choices precede the height and the volume.

For a real critical configuration at level \(j\), set
\[
 \mathsf m_j(U)=L^{2j}\max_p|\partial U(p)-I|.
\]
For \(1\le j\le K\) and \(0<e\le1\), let
\(\mathcal C_j(b_0,b_1;e)\) be the set of block fields whose
minimal regular orbit exists, has \(\mathsf m_j<e\), and whose
pairs on every domain lie in the cores with radii \((b_0,b_1)\).
The assertion about pairs is required for every representative.
We use \(\mathcal W_{k+1}(\mathsf t;b_0,b_1)\) for the global
domain of the actual step: its averaged axial background and the
whole cutoff chart are defined; the lower backgrounds over that
chart exist; the whole background belongs to the shift domain;
the octet restrictions needed in \cref{II-49-44} have their
minimal orbits; and every plain triple of the top background on a domain
belongs to the core with radii \((b_0,b_1)\).
These are the actual domain requirements of the preceding
construction, including the off-fibre response neighbourhood
in \cref{II-40-99}.

\begin{proposition}\label{II-42-39}
There is \(\varepsilon_{\mathcal W}>0\), independent of height
and volume, with the following properties. If \(\mathsf t>1\)
is real, \(\varepsilon(\mathsf t)\le1\), \(0<b_0,b_1\le1\),
\(\delta(\mathsf t)\le\varepsilon_{\mathcal W}\), and
\(\varpi(W)<\varepsilon_{\mathcal W}b_0\), then
\(W\in\mathcal W_{k+1}(\mathsf t;b_0,b_1)\).
Its minimal orbit is contained in
\(\mathfrak U'_{k+1}(B_3\varepsilon_{\mathcal W}b_0)\).
The averaged axial field \(V=V_k^{\rm bf}(W)\) is independent
of its construction, depends continuously on \(W\), and satisfies
\begin{equation}\label{II-42-40}
 \mathcal MV=W,\qquad
 \varpi(V)\le2B_3L^{-2}\varepsilon_{\mathcal W}b_0.
\end{equation}
For a real cutoff field \(B\), its chart is
\(V_W(B)=e^{iB'}V\), where
\begin{equation}\label{II-42-41}
 \sup|B'|<2C_{\rm db}\delta(\mathsf t),\qquad
 \mathcal M V_W(B)=W,\qquad
 \varpi(V_W(B))<K_2\varepsilon_{\mathcal W},
 \quad K_2=2B_3L^{-2}+8C_{\rm db}.
\end{equation}

If \(2b_0\le\varepsilon_{\mathcal W}\), then
\begin{equation}\label{II-42-73}
\begin{split}
 \mathcal C_j(b_0,b_1;e)=\{W:\;&U_j(W)\text{ exists, and its orbit has}\\
 &\mathsf m_j<\min(e,b_0),\quad |J^{(j)}|<b_0\}.
\end{split}
\end{equation}
On this set every plain triple belongs to its core. The set is
invariant under every isometry of the level-\(j\) lattice, and
\[
 \mathcal C_{k+1}(b_0,b_1;e)
          \subset\mathcal W_{k+1}(\mathsf t;b_0,b_1).
\]
It contains the open neighbourhood of the identity
\begin{equation}\label{II-42-72}
 \left\{W:\varpi(W)<
       \min\{\bar\varepsilon,e/B_3,\varepsilon_{\mathcal W}b_0\}
                                                        \right\}.
\end{equation}
\end{proposition}

\begin{proof}
Choose \(\varepsilon_{\mathcal W}\le\varepsilon_{\mathfrak h}\)
so small that
\begin{gather*}
 2L^2B_3K_2\varepsilon_{\mathcal W}\le\bar\varepsilon,
 \qquad L^2B_3K_2\varepsilon_{\mathcal W}\le\varepsilon_\flat,
 \qquad 2B_3L^{-2}\varepsilon_{\mathcal W}\le\varpi_{\rm db},\\
 C_{\rm db}\varepsilon_{\mathcal W}\le r_{\rm db},\qquad
 l_MB_3\varepsilon_{\mathcal W}\le c_{\rm rg},\qquad
 B_{\rm rg}l_MB_3\varepsilon_{\mathcal W}
             \le\mathsf oLM B^{(\mathrm{cg})}a_{\mathfrak h},\qquad
 B_3\varepsilon_{\mathcal W}\le a_{\mathfrak h}\le1.
\end{gather*}
Here \(a_{\mathfrak h},\varepsilon_{\mathfrak h}\) are the
fixed radii in the actual shift construction. Reducing the first
constant also imposes all preceding curvature ceilings. These
are finitely many positive bounds depending only on the fixed
geometric constants.

\Cref{I-58-22} constructs the whole minimal orbit and gives the
claimed regularity. By \cref{I-10-101}, its level-\(k\) average
has curvature at most \(2B_3L^{-2}\varepsilon_{\mathcal W}b_0\).
The contour estimates of \cref{I-10-51,I-10-53} therefore put
every contour family in its unique-mean domain. For
\(x\in B(y)\), set \(u(x)=(\mathcal M^kU)(y,x)\), and set
\(u=I\) at the other bare sites. Since \(u(y)=I\), this is a
residual gauge for the top datum. It makes all these contour
variables equal to \(I\). Two representatives so normalized
have the same averaged field: their residual gauge is determined
by these contour variables and its value \(I\) at the block
centres. This proves \eqref{II-42-40} and independence.
Restriction of criticality in \cref{II-8-36}, followed by the
uniqueness in \cref{I-58-22}, identifies the orbit over \(V\)
with this normalized subset of the top orbit.

The representative supplied by \cref{I-58-22} varies
holomorphically in logarithmic row coordinates near every
datum under consideration. The real averages and their contour
means are continuous there: a convergent sequence of means has,
by compactness of the group, a convergent subsequence, and its
limit is the unique mean of the limiting family. Applying the
preceding axial gauge proves continuity of \(V_k^{\rm bf}\)
on the indicated open curvature set.

The pivot chart estimates used in \cref{II-40-99} now apply to
\(V\). Since \(\sup|gB|<\delta(\mathsf t)\), they give the
first two conclusions of \eqref{II-42-41}. The value of the
chart does not depend on which permitted ball is used: two
holomorphic representatives of its correction have the same
germ at zero and hence agree on the smaller ball. At spacing
one, \(|\nabla B'|\le2\sup|B'|\). The real plaquette product
estimate gives
\[
 |\partial(e^{iB'}V)-I|
       \le |\partial V-I|+4\sup|B'|,
\]
which proves the last bound in \eqref{II-42-41}. Thus
\cref{I-58-22} also constructs every lower background along
the cutoff chart. At \(k=0\) the lower configuration is the
chart field itself; the current formula gives
\(|J(V_W(B))|\le4\varpi(V_W(B))\), within the regularity
radius by the same choice of \(\varepsilon_{\mathcal W}\).

We check the local conditions rather than assuming them from
whole-lattice existence. A cube of side at most \(\mathsf oLM\)
is contained in a coordinate cube of integral side
\(l<l_M\); enlarge the latter by two layers at each face.
The room inequality puts this enlargement inside a chart with
sides less than half the torus period. If the whole curvature
and current have radius \(\alpha\), \cref{I-28-75} supplies
there a real gauge with
\[
 |A|,|\nabla A|<B_{\rm rg}l_M\alpha .
\]
This is precisely the local-gauge condition of the simple set
when its right side is at most
\(\mathsf oLM B^{(\mathrm{cg})}b_0\).
The other simple-set inequalities follow from curvature and
current, and the perturbing field of a plain triple is zero.

For every localized tower used by a core or an octet, apply
\cref{II-8-36} to the restricted real critical field. At level
\(q\), its defining averages have curvature at most
\(2L^2\alpha\), by \cref{I-10-101}. The choice above puts
them inside the real-background radius, so \cref{I-58-22}
identifies their minimal orbit with the restriction. The
complex and real means agree on these real fields. At depth
\(p\), curvature and current are multiplied respectively by
\(L^{-2(k+1-p)}\) and \(L^{-3(k+1-p)}\); hence all tower
inequalities in the core hold. This proves plain-core membership for the top background
and existence of the localized minimal orbits needed along
the lower cutoff chart. A residual gauge on an
octet extends by \(I\) outside it: its boundary nodes are fixed
and its outer determining rows are unchanged. Thus the same
statements hold for every representative, not only for the
chosen axial field. The shift hypotheses now follow from
\(P_W\in\mathcal Q_{k+1}(T_\ell;
a_{\mathfrak h},a_{\mathfrak h},a_{\mathfrak h})\) and
\(\sup|gB|<\varepsilon_{\mathfrak h}\).

This also discharges the response hypothesis in
\cref{II-40-99}. The simple-set gauges just constructed verify
the natural-class conditions in the proof of \cref{II-51-52};
in level-\(k\) units the current is \(L^{-3}J^{(k+1)}\).
The actual response of \cref{I-88-53,I-88-93} therefore applies
near \(p=V\), including variations off its one-step fibre.
For \(k\ge1\) take the open set
\(\{V':\varpi(V')<\bar\varepsilon\}\) about \(p\);
\cref{I-58-22} constructs its backgrounds. The response is
real regular and critical, and uniqueness identifies it with
that constructed family. For \(k=0\) take instead the open
set of fields with both curvature and current below the
regularity ceilings; the background is the field itself.
The off-fibre derivatives and quadratic subtraction in
\cref{II-40-99} are thus derivatives of actual backgrounds.

It remains to check the assertions for \(\mathcal C_j\).
Its core condition immediately gives
\(\mathsf m_j<b_0\) and \(|J^{(j)}|<b_0\) on the whole
torus. Conversely, these two bounds imply the strengthened
regularity bound \(b_0+b_0^2\le2b_0\). The preceding local
gauge proof applies with \(\alpha=2b_0\). By
\eqref{eq:global-core-gauge-constant} and \(M\ge L^2\),
\[
 2B_{\rm rg}l_Mb_0
                \le\mathsf oLM B^{(\mathrm{cg})}b_0.
\]
The restricted critical-tower argument is unchanged. It puts
every plain triple in the core, for any \(b_1>0\), and proves
\eqref{II-42-73}. Curvature, current norms, averages and
critical orbits are carried to their counterparts by an
isometry: transport the paths, and differentiate the invariant
action. Consequently this description is isometry invariant.
Also \(\varpi(W)<2b_0\le\varepsilon_{\mathcal W}\) by
\cref{I-10-101}; applying the already proved domain assertion
with radii \((1,1)\) gives all its conditions except the
specified plain-core radii, which we have just checked.
Finally \cref{I-58-22} and the first assertion applied at
\(b_0,b_1\) prove \eqref{II-42-72}.
\end{proof}

The actual shifts of \cref{II-51-52}, the octet identifications
of \cref{II-49-44}, and the substitution identity
\eqref{eq:one-term-background-sum} now hold on this global
domain. Hence all three interaction totals and the fourth
total in \cref{II-128-31} have their physical meanings there.
The positive normalizer is
\(Z^{(k)}(W)=Z_k[P_W]\), and
\begin{equation}\label{II-147-12}
 \mathcal P[\mathbf G_\tau](W)
                         =\tfrac12 n_{\mathcal O}\operatorname{Log}\tau.
\end{equation}
Indeed the identified Hessian and gauge quadratic form are
those of \cref{II-43-100}; the latter's gauge Gaussian identity
gives this formula. With \eqref{eq:sf-chart-thresholds}
also imposed, applying \cref{II-53-39,II-40-99} gives
the nonzero chart integral and the effective-action value
\begin{equation}\label{II-147-13}
 \log(\mathsf{pf}_k(\mathsf t)Z^{(k)}(W))
  +\mathcal P[\mathbf G_\tau](W)+\mathcal P[\mathbf F](W)
  -\tfrac32\#\mathfrak X^+\operatorname{Log}\tau.
\end{equation}
The zero input array proves \eqref{II-147-12} even when a
parameter family is empty. The chart identity itself needs
no whole-space measurability assumption. To regard it as a
pushforward density, retain the bounded-Borel hypothesis on
\(\chi^{\rm sf}_{\mathsf t}F_{x,S}\) from
\cref{II-128-31}.

\subsection{Covariance of the totals}\label{sec:actual-step-total-covariance}
The individual domains of the expansion use a fixed partition.
An isometry of the coarse lattice need not preserve that partition.
To recover its covariance, first sum the terms carrying one marked
bond. Their sum is an unlocalized Gaussian expectation, to which
the isometry applies directly. The tents then recover the pointed
totals.

\begin{proposition}\label{II-55-67}
Suppose the input pointed totals are covariant under all isometries
of their level lattices on \(\mathcal C_j(a_{0,j},a_{1,j};e)\),
with a common \(e\in(0,1]\). Use the actual shifts and substitutions
and the numerical hypotheses of \cref{II-128-31}. Suppose
\(2\bar a_0\le\varepsilon_{\mathcal W}\), and let
\(W\in\mathcal C_{k+1}(\bar a_0,\bar a_1;e)\).
Then the pointed totals of the actual output are defined for
every representative of the minimal orbit, and
\begin{equation}\label{II-55-68}
 \mathcal P_{\mathsf rz}[E^{(k+1)}](\mathsf rW)
                         =\mathcal P_z[E^{(k+1)}](W)
\end{equation}
for every isometry \(\mathsf r\) of the level-\((k+1)\) lattice.
The same assertion holds separately for the normalizer and the
Gaussian term. At height zero the input condition is empty.
\end{proposition}

\begin{proof}
By \cref{II-42-39}, both \(W\) and \(\mathsf rW\) belong to
the actual global domain. Write \(V=V_k^{\rm bf}(W)\),
\(\beta b=(\mathsf rb)^+\), and define, on positive bonds,
\[
 (\mathsf O B)(\beta b)=O_bB(b),\qquad
 O_b=\begin{cases}1,&\mathsf rb>0,\\
             -\operatorname{Ad}(V(b)^{-1}),&\mathsf rb<0.
             \end{cases}
\]
Use the same formula on all fine bonds, denoting that map by
\(\mathsf r_V\). The positive pivot bonds are carried to the
positive pivot bonds, since the central pivot formula is affine
covariant. Thus \(\mathsf O\) is the restriction to the
nonpivot coordinates. Each \(O_b\) is real orthogonal for the
Pauli pairing. Consequently \(\mathsf O^{\mathsf T}=\mathsf O^{-1}\),
it preserves Lebesgue measure and every bond norm, and
\[
 \mathsf O(\mathbf t\cdot B)
                 =(\mathsf r\mathbf t)\cdot\mathsf OB,
 \qquad (\mathsf r\mathbf t)_{\beta b}=t_b.
\]
In particular it preserves the cutoff and complex polydisc
\(\mathbb D=\{\sup|B|<\mathsf C_{\mathcal A}\mathfrak b(\mathsf t)\}\).

We check its action on the actual chart. Relabelling the path
families carries means, constraints and minimal orbits to their
counterparts. Normalization of the contour means therefore gives
\(V_k^{\rm bf}(\mathsf rW)=\mathsf rV\).
The identity
\(e^{i\mathsf r_V A}\mathsf rV=\mathsf r(e^{iA}V)\)
follows also on a reversed bond from
\(V(b)^{-1}e^{-iA(b)}
 =e^{-i\operatorname{Ad}(V(b)^{-1})A(b)}V(b)^{-1}\).
Differentiating the actual constraint gives its twisted linear
identity. Since pivots are permuted, its pivot elimination and
right inverse obey
\begin{equation}\label{II-55-56}
 \mathsf C_{\mathsf rV}\mathsf O=\mathsf r_V\mathsf C_V,
 \qquad h_{\mathsf rV}\mathsf r_{\mathcal M(V)}
                                  =\mathsf r_Vh_V.
\end{equation}
The defining correction equations and uniqueness of their
holomorphic germ give the analogous identity for
\(\widetilde{\mathsf D}_V\) near zero. Both chart corrections
are defined on balls containing \(g\mathsf C_VB\), and the
twist preserves these balls. On their common smaller ball the
identity theorem extends the germ identity. Hence
\begin{equation}\label{II-55-58}
 V_{\mathsf rW}(\mathsf OB)=\mathsf rV_W(B)
\end{equation}
throughout the real cutoff chart.

The action and the sum of squared contour logarithms are
invariant under this substitution. The Jacobian is invariant
as well: differentiating the correction identity conjugates
\(h_V\partial\widetilde{\mathsf D}_V\); its trace-log series
is unchanged term by term. Each Haar factor depends only on
the squared norm of the corrected bond field, also unchanged.
The physical action Hessian and the gauge quadratic form
therefore satisfy
\[
 \mathsf B_{\mathsf rW}=\mathsf O\mathsf B_W\mathsf O^{\mathsf T},
 \qquad \widehat{\mathsf G}_{\mathsf rW}
             =\mathsf O\widehat{\mathsf G}_W\mathsf O^{\mathsf T}.
\]
Their sum, inverse and Gaussian measure transform accordingly.
The sums of the three local potentials are holomorphic on
\(\mathbb D\). They agree near zero with the physical
remainders in \cref{II-40-99}, so uniqueness of holomorphic
continuation proves their covariance on all of \(\mathbb D\).
This argument identifies the actual quadratic form; no covariance
of a chosen localization is being assumed.

Let \(\hat s_W=\sum_Y\mathbb V'[\mathfrak f](Y;P_W|_Y,\cdot)\).
We next prove
\begin{equation}\label{II-55-70}
 \hat s_{\mathsf rW}(\mathsf OB)=\hat s_W(B).
\end{equation}
At \(B=0\), descent of the top critical orbit gives
\(\mathsf m_k=L^{-2}\mathsf m_{k+1}<L^{-2}e\).
The constructed response is continuous, so this strict bound
persists with \(L^{-2}e\) on a neighbourhood of zero.
Further descent gives
\(\mathsf m_j<L^{-2(k-j)}L^{-2}e\le e\).
The actual substitution and octet identities
\cref{II-50-61,II-49-44} put all corresponding pairs in the
input cores. They do so at \(\mathsf rW\) as well. Thus, for
the fields in these two neighbourhoods, the earlier action is
the finite sum
\begin{equation}\label{II-55-65}
 S[\mathfrak f](U)=\sum_{j=1}^k
       \left(\mathcal P[E^{(j)}](\mathcal M^jU)
                         -c_j-e_j A^{\eta_j}(U)\right),
\end{equation}
where \(c_j\) is its value at the trivial pair. Each earlier
total is invariant by hypothesis, and the plaquette action is
invariant. Equation \eqref{II-55-58} and the actual fourth
sum therefore give \eqref{II-55-70} near zero. Both sides are
holomorphic on \(\mathbb D\), so they agree there. At \(k=0\)
the sum is empty and both sides are zero.

We now express the marked-bond sums without the partition.
Let \(\mu=\mu_{\tau^{-1}C_W}\),
\(q(B)=\frac12(B,\widehat{\mathsf G}_WB)_1\), and
\(q_b(B)=\frac12\langle B(b),(\widehat{\mathsf G}_WB)(b)\rangle\).
Write \(P^\natural_W\) for the total continued potential.
The interpolation of \eqref{II-20-95} has total
\begin{equation}\label{II-55-59}
\begin{split}
 \Psi_W(\mathbf t;B)={}&P^J_W(\mathbf tB)+\hat s_W(\mathbf tB)\\
 &+\int_0^1(1-r)D^2(P^G_W+\tau P^A_W)(r\mathbf tB)[B,B]\,dr.
\end{split}
\end{equation}
Set
\[
 Z_G(u)=\int e^{-u(1-\tau)q}\,d\mu,\quad
 Z_\chi(\mathbf t)=\int\chi_{\mathbf t}e^{-(1-\tau)q}\,d\mu,
 \quad
 Z_\bullet(t)=\int\chi e^{-(1-\tau)q+\Psi_W(t\mathbf1;B)}\,d\mu.
\]
All these integrals are nonzero by the Gaussian cluster theorem:
the quadratic family has size at most \(w_*/2\), and the two
interpolated families have sizes at most \(w_*/2,w_*/4\).
The insertion and cutoff identities used in \cref{II-53-39}
give, after summing over the finitely many domains,
\begin{align}
 \sum_X\mathbf G_\tau(X,b;P_W|_X)
  &=- (1-\tau)\int_0^1
       \frac{\int q_b e^{-u(1-\tau)q}\,d\mu}{Z_G(u)}\,du,
                                                    \label{eq:total-G-density}\\
 \sum_X E^\chi(X,b;P_W|_X)
  &=\int_0^1
       \frac{\partial_{t_b}Z_\chi(t\mathbf1)}{Z_\chi(t\mathbf1)}\,dt,
                                                    \label{eq:total-cut-density}\\
 \sum_X E^\bullet(X,b;P_W|_X)
  &=\int_0^1
       \frac{\int\chi e^{-(1-\tau)q+\Psi_W(t\mathbf1;B)}
                         \partial_{t_b}\Psi_W(t\mathbf1;B)\,d\mu}
            {Z_\bullet(t)}\,dt.                    \label{eq:total-int-density}
\end{align}
Derivatives in the cutoff identity exist almost everywhere;
the insertion bounds give integrability in the parameter.
The finite sum allows one common null set for all domains.
Thus these are equalities of the defined integrals, including
for complex \(\tau\), not formal differentiated expressions.

Under \(B\mapsto\mathsf OB\), the measure transforms to its
counterpart, \(q_b\) to \(q_{\beta b}\),
\(\chi_{\mathbf t}\) to \(\chi_{\mathsf r\mathbf t}\), and,
by \eqref{II-55-59} and \eqref{II-55-70}, \(\Psi_W\) to
\(\Psi_{\mathsf rW}\). The parameter coordinate \(t_b\)
becomes \(t_{\beta b}\), including its exceptional null set.
All three bond densities in the preceding displays are
therefore covariant. The tent identity
\begin{equation}\label{II-55-69}
 h_{\mathsf rz}(\beta b)=h_z(b)
\end{equation}
now proves pointed covariance for \(\mathbf G_\tau,\mathbf F\).
Indeed a marked term vanishes unless \(b^\sharp\subset X\),
whereas \(h_z(b)\ne0\) implies \(z\in b^\sharp\).
One may therefore drop \(X\ni z\) before exchanging the
domain and bond sums.

For the normalizer, the finite difference in \eqref{II-44-16}
telescopes to the whole-torus diagonal logarithm. Its nonpivot
density is \(-\tfrac12\operatorname{tr}[(\log\mathsf A_W)(b,b)]\),
and its pivot density at \(b_0(c)\) is
\(-\log|\det m_c(V)|\), in each case subtracting the value at
the trivial background. Conjugation by \(\mathsf O\) leaves
the diagonal trace unchanged and permutes the pivot
determinants. Equation \eqref{II-55-69} gives pointed
covariance here too. Adding the three contributions proves
\eqref{II-55-68}.
\end{proof}

\subsection{The real progression and its coefficients}
\label{sec:actual-regular-progression}
\label{sec:real-regular-progression}

We now iterate the actual step with prescribed coupling
increments. Its estimates use the earlier levels only
through a summable scalar norm. The covariance just proved
allows us to extract the quadratic coefficient of each
output, which will determine the intrinsic running coupling.

Choose fixed $\alpha_0,\alpha_1>0$ satisfying the radius
and rate conditions of \cref{II-128-31}, with
\begin{equation}\label{II-131-3}
 \alpha_0+\alpha_1\le a^*,\qquad
 \mathsf c_{\rm sh}M\alpha_0<\alpha_1,\qquad
 \alpha_1^2\le\mathsf c_\natural\alpha_0,\qquad
 a^*\le\min\{\alpha_N,\varepsilon_{\mathcal W}/2,a^\sharp\}.
\end{equation}
At height $h$ use the absolute radii
$a_{\iota,j}=\alpha_\iota$ for $0\le j\le h$ and
$a_{\iota,h+1}=\alpha_\iota/(1+\beta_{\rm s})$,
for $\iota=0,1$. The enlarged output radii are again
$(\alpha_0,\alpha_1)$.
Use the constant-ratio case of \cref{II-52-89},
\[
 \omega_n=L^{-(n-1)},\qquad \mathsf K=K_\Lambda,
 \qquad K_{\mathbf M}=K_{\mathbf M}(K_\Lambda),
\]
with its ceiling $\check s(K_\Lambda)$.
All geometric floors and the admitted family of tori
in \cref{sec:actual-step-global-domain} are fixed
uniformly in $h$.

Put $r_\alpha=a^{\rm us}(\alpha_0,\alpha_1)$ as in
\eqref{II-17-20}, and choose
\begin{equation}\label{eq:real-progression-constants}
 E_0\ge2E_N,\qquad
 B_\sharp\ge C_{\mathrm{am}}r_\alpha^{-2}E_0,\qquad
 N_*\ge\frac L{L-1}
           (E_0+\tfrac32\alpha_0^2B_\sharp).
\end{equation}
The one-step conditions include
$\delta_NM\ge\kappa+1$ and $4\beta_{\rm s}\kappa\ge1$.
Fix $A_1>0$ and $A_0\ge C_\chi A_1$, with $C_\chi$ as in
\eqref{eq:sf-chart-thresholds}. Write
\[
 \phi(t)=t^{-1/2}(\log t)^{\mathfrak p_0},\quad
 \varepsilon(t)=A_0\phi(t),\quad
 \delta(t)=A_1\phi(t),\quad
 \mathfrak b(t)=A_1(\log t)^{\mathfrak p_0},
\]
and
\begin{equation}\label{II-131-14}
 \Theta_c(t,N)=K_{\mathbf M}
   \left(\delta(t)+\frac{(1+\beta_{\rm s})\delta(t)}{\alpha_0}N
                     +e^{-\gamma_2\mathfrak b(t)^2/2}\right).
\end{equation}
Define the real-step majorant by
$\Xi(t,N)=\tfrac12E_0+\Theta_c(t,N)$ for $t>1$ and $N\ge0$.
The fixed comparison gives $C_\chi\delta(t)\le\varepsilon(t)$
for every $t>1$. Choose $T\ge4$ so that, for all $t\ge T$,
the remaining absolute smallness hypotheses of
\cref{II-128-31} hold and
\begin{equation}\label{eq:progression-floor}
 \begin{gathered}
 \varepsilon(t)\le1,\qquad
 \delta(t)\le\min\{1,\varepsilon_{\mathcal W},r_\chi/C_\chi\},\qquad
 \mathfrak b(t)\ge\mathfrak b_*,\\
 \frac{(1+\beta_{\rm s})\delta(t)}{\alpha_0}
       \le\check s(K_\Lambda),\qquad
 \Theta_c(t,N_*)\le E_0/2 .
 \end{gathered}
\end{equation}
Such $T$ exists: $\phi(t)\to0$, while
$\mathfrak b(t)\to\infty$. It depends only on the
fixed constants.

Let $x_0,\ldots,x_k\ge T$ be real, with
$b_i=L^{-1}x_{i-1}-x_i$ and $|b_i|\le B_\sharp$
for $1\le i\le k$. Starting with the empty array,
at height $h$ apply the actual step at $t=x=x_h$
to $(E^{(i)},b_i)_{i\le h}$, with empty last slot:
\begin{equation}\label{eq:real-progression-recursion}
 E^{(h+1)}
    =\mathbf N^{(h+1)}
        +\mathbf F[x_h,x_h;(E^{(i)},b_i)_{i\le h}].
\end{equation}
The gauge contribution is zero because $\tau=1$.

\begin{theorem}\label{II-144-26}\label{II-147-20}
Under these conditions, the real progression
\eqref{eq:real-progression-recursion} is defined through
level $k+1$. Its terms are gauge invariant, conjugation
real and holomorphic on their cores, and are covariant
under the cube-grid translations. On every admitted torus,
for $P\in\mathcal U_j^\circ(X;\alpha_0,\alpha_1)$,
\begin{equation}\label{II-144-27}
 |E^{(j)}[x_\bullet](X,P,z)|
       \le E_0e^{-(\kappa+1)d_j(X)},\qquad
 \|\widehat E^{(j)}[x_\bullet]\|_{j,\kappa+1}\le E_0,
 \quad 1\le j\le k+1 .
\end{equation}
The companions agree with the finite-volume terms on
every non-wrapping domain. Their values and domains
are unchanged upon restriction to any unbounded
subfamily of admitted tori.
The level-$j$ terms depend only on $x_0,\ldots,x_{j-1}$.

For every $e\in(0,1]$, the pointed totals are covariant
under all isometries of their level lattices on
$\mathcal C_j(\alpha_0,\alpha_1;e)$.
Their quadratic coefficients are therefore defined, and
\begin{equation}\label{eq:real-progression-coefficient}
 e_j:=\mathsf{cf}(\widehat E^{(j)})\in\mathbb R,\qquad
 |e_j|\le C_{\mathrm{am}}r_\alpha^{-2}E_0\le B_\sharp .
\end{equation}
These numbers are independent of the lattice.

On the global small-field class the whole sums satisfy
the actual chart and effective-action identities of
\cref{II-128-31} and \eqref{II-147-13}.
The pushforward-density assertion retains the
whole-space bounded-Borel hypothesis of \cref{II-40-99}.
\end{theorem}

\begin{proof}
At height zero the earlier array is empty and the fourth
potential vanishes. Suppose the assertion holds through
height $h$. The earlier local terms have norm at most
$E_0$ and the prescribed counterterms have modulus at
most $B_\sharp$. Hence the input size satisfies
\begin{equation}\label{II-144-25}
 N_{\mathfrak a^{\rm abs}_h,\omega}
             ((E^{(i)},b_i)_{i\le h})
 \le\sum_{n=1}^hL^{-(n-1)}
          (E_0+\tfrac32\alpha_0^2B_\sharp)
 \le N_* .
\end{equation}
No radius decreases at an earlier level; the enlarged
output radii are precisely the next input radii.
The actual shifts and octet substitutions thus satisfy
the hypotheses of \cref{II-128-31}, with
$s_{\mathfrak a}=(1+\beta_{\rm s})\delta(x_h)/\alpha_0$.
The Gaussian and normalizer inputs are
\cref{II-38-61,II-44-37}, the first three potentials
come from \cref{II-51-52}, and the fourth is the
substitution of \cref{II-49-44,II-52-89}.
Their physical meanings on the global class were
established above.

The normalizer contributes at most
$(E_0/2)e^{-(\kappa+1)d}$.
Since $\tau=1$, the remaining contribution is at most
$\Theta_c(x_h,N_*)e^{-(1+4\beta_{\rm s})\kappa d}$.
The rate condition and \eqref{eq:progression-floor}
prove \eqref{II-144-27} at the next level.
Each operation preserves slice holomorphy,
conjugation reality, gauge invariance and covariance
under cube-grid translations.

Local charts identify the core, shifts and response,
including their domains: the chart carries every
contained stencil and all the defining tests in both
directions. They therefore identify the constructed
term. Two sufficiently large tori give the same
answer on a fixed finite domain. This defines the
companion, and its bound passes through the chart.
Every finite domain fits in some member of any
unbounded admitted subfamily, so that the companion
does not depend on this subfamily.

To check the stronger covariance of pointed totals,
fix $e\in(0,1]$. The preceding totals have covariance
on their classes with this same radius.
The inherited curvature allowance at the next step
is $L^2e$, which is at least $e$.
Since $2\alpha_0\le\varepsilon_{\mathcal W}$,
\cref{II-55-67} gives covariance on
$\mathcal C_{h+1}(\alpha_0,\alpha_1;e)$.
The background and every core restriction used here
exist by \cref{II-42-39}. \Cref{II-129-73} therefore applies on the near-trivial
window \eqref{II-42-72}, with usable radius $r_\alpha$.
It gives \eqref{eq:real-progression-coefficient}.
The coefficient is real by conjugation of its
absolutely convergent series.
Its construction from the companion makes it
independent of the torus.

The global-domain and substitution identities apply
to these earlier terms with their prescribed
subtractions. They give the stated whole-sum and
chart identities. Finally, every operation at height
$h$ uses only the earlier terms and $x_h$.
Induction proves dependence on the stated prefix,
and completes the construction.
\end{proof}


\subsection{The running coupling determined by the local output}
\label{sec:canonical-flow}

The preceding construction allowed prescribed increments
\(b_j=L^{-1}x_{j-1}-x_j\). We now determine these increments
from the quadratic coefficients. Choose \(0<\gamma\le1/2\)
so that the hypotheses of \cref{II-147-20} hold for every
finite sequence with \(|b_j|\le B_\sharp\) and
\(x_j\ge\gamma^{-2}\), and require
\begin{equation}\label{eq:flow-stopping-ceiling}
 \gamma^{-2}>\bar B:=\frac{L}{L-1}B_\sharp,\qquad
 L^{-1}\gamma^{-2}-B_\sharp>1 .
\end{equation}
In particular the thresholds \(\varepsilon(x),\delta(x)\)
are at most one for every \(x\ge\gamma^{-2}\).

Given \(\beta\ge\gamma^{-2}\), put \(x_0=\beta\).
As long as \(x_0,\ldots,x_j\) have been constructed and
\(x_j\ge\gamma^{-2}\), form the next local output by
\eqref{eq:real-progression-recursion} and set
\begin{equation}\label{II-146-28}
 e_{j+1}=\mathsf{cf}\bigl(\widehat E^{(j+1)}
                                  [x_0,\ldots,x_j]\bigr),
 \qquad x_{j+1}=L^{-1}x_j-e_{j+1}.
\end{equation}
Stop at the first index \(J\) with \(x_J<\gamma^{-2}\).
The families and the coefficient in this definition use
only the indicated prefix, including its already determined
increments; the new coupling is not an input to them.

\begin{proposition}\label{II-132-35}\label{II-146-32}
This recursion defines a unique finite sequence of real
numbers \(x_0,\ldots,x_J\), with \(J\ge1\).
Its increments are real and \(|e_j|\le B_\sharp\).
The sequence strictly decreases and satisfies
\begin{equation}\label{II-132-36}
 \begin{split}
 x_j-L^{-(j-i)}x_i
    &=-\sum_{l=i+1}^jL^{-(j-l)}e_l,\\
 |x_j-L^{-(j-i)}x_i|
    &\le\bar B(1-L^{-(j-i)})\qquad(0\le i\le j\le J).
 \end{split}
\end{equation}
Moreover,
\begin{equation}\label{II-132-37}
 \frac{\beta}{\gamma^{-2}+\bar B}<L^J
       \le\frac{L\beta}{\gamma^{-2}-\bar B},
 \qquad
 1<L^{-1}\gamma^{-2}-B_\sharp\le x_J<\gamma^{-2}.
\end{equation}
Every prefix ending before \(J\) is an admissible sequence
for \cref{II-147-20}. Its families and coefficients agree
with those obtained from any longer such prefix.
The coupling sequence is independent of the torus.

More generally, let \((u_0,\ldots,u_n)\) be an admissible
prescribed sequence with \(u_0=\beta\) and
\(u_i\ge\gamma^{-2}\) for \(0\le i<n\).
It is a prefix of this sequence if and only if
\begin{equation}\label{eq:flow-coefficient-characterization}
 L^{-1}u_{j-1}-u_j
       =\mathsf{cf}\bigl(\widehat E^{(j)}
                                    [u_0,\ldots,u_{j-1}]\bigr)
                    \quad(1\le j\le n).
\end{equation}
\end{proposition}

\begin{proof}
At \(j=0\) the earlier interaction array is empty.
Suppose a real prefix has been constructed with all its
members at least \(\gamma^{-2}\) and all its increments
bounded by \(B_\sharp\). The hypotheses of
\cref{II-147-20} apply to this prefix. They give a
well-defined next family and a real coefficient of
modulus at most \(B_\sharp\). Thus \eqref{II-146-28}
defines the next real member. If it is still above
the stopping threshold, it meets the induction hypotheses.

Consistency follows at the same time. At height \(j\)
the construction reads precisely the old families,
the old increments and \(x_j\). Appending a coupling
does not change these data. Induction therefore identifies
each earlier finite-volume family as a function on its
entire core. The companions agree through their local
charts, so their coefficients agree as well. In
particular the next coefficient requires neither a choice
of finite torus nor a radius table at the next coupling.
This is the prefix dependence proved in
\cref{II-147-20}, applied at successive heights.

Iteration of the affine recursion gives the first line
of \eqref{II-132-36}; the geometric sum of
\(|e_l|\le B_\sharp\) gives the second. As long as
\(x_j\ge\gamma^{-2}\), it follows that
\[
 L^j(\gamma^{-2}-\bar B)\le\beta-\bar B.
\]
The positive coefficient on the left and \(L>1\)
exclude an infinite sequence above the threshold.
Thus the first passage \(J\) exists, is at least one,
and the recursion supplies \(x_J\).
For \(1\le j\le J\),
\[
 x_{j-1}-x_j
   \ge(1-L^{-1})(x_{j-1}-\bar B)>0.
\]
The last step also gives
\(x_J\ge L^{-1}\gamma^{-2}-B_\sharp>1\).

For the first-passage bounds use
\[
 L^{-J}(\beta+\bar B)\le x_J+\bar B
                       <\gamma^{-2}+\bar B
\]
and the preceding upper bound at \(j=J-1\).
These imply \eqref{II-132-37}.
Every earlier prefix has thresholds at most one,
members at least \(\gamma^{-2}\), and bounded increments,
so it is admissible.

Finally, equality of a prescribed prefix with the
constructed one identifies its next local output and
coefficient. Equation \eqref{eq:flow-coefficient-characterization}
then identifies the next coupling. Starting at
\(u_0=x_0=\beta\), induction proves sufficiency.
Necessity is \eqref{II-146-28} and prefix consistency.
\end{proof}


\section{Coordinates and tests of a history step}
\label{sec:actual-history-coordinates}

\subsection{Free coordinates and the reference measure}
\label{sec:history-coordinates}

Fix a history \(h\) of length \(k\). All bonds below are taken in
positive orientation; the values at reversed bonds are prescribed by
inversion for configurations and by a change of sign for Lie algebra
fields. Put \(\Omega_0=\mathbb T\) and
\(\Omega_{k+1}=\Lambda_{k+1}=\varnothing\). At level \(l\), write
\(\Gamma_l^{\mathrm{pr}}\) for the bonds with no endpoint in
\(\Omega_{l+1}\), and \(\Gamma_l^{\mathrm{le}}\) for those with an
endpoint in \(\Omega_l\) and none in \(\Omega_{l+1}\).
Thus
\(\Gamma_l^{\mathrm{ti}}=\Gamma_l^{\mathrm{pr}}\setminus
\Gamma_l^{\mathrm{le}}\) consists of the bonds with no endpoint in
\(\Omega_l\). The superscripts distinguish the coordinates that are
retained freely from those determined by averaging.

Let \(\mathfrak X_l^{\mathrm E}\) be the nonpivot bonds of step
\(l+1\) with no endpoint in \(\Lambda_{l+1}\). Once the strip
coordinate of that step has been defined, denote its scaled value by
\(\mathsf a'_l=g_l^{-1}\mathcal A_l^{\prime\diamond}\).
The free coordinate space is
\begin{equation}\label{II-58-20}
 \mathsf X_h^{\mathrm{fr}}
 =\prod_{l=0}^k\mathcal V(\Gamma_l^{\mathrm{le}})
   \times\prod_{l=0}^{k-1}\mathfrak g^{\mathfrak X_l^{\mathrm E+}}.
\end{equation}
We write a point as \((\mathbf v,\mathfrak q)\). Its completion
\(\mathrm{cp}_h(\mathbf v,\mathfrak q)=(\mathbf w,\mathfrak q)\)
retains the free coordinates and sets, successively,
\begin{equation}\label{II-58-5}
 w_0=v_0,\qquad
 w_l(c)=
 \begin{cases}
 v_l(c),&c\in\Gamma_l^{\mathrm{le}},\\
 \mathcal M(w_{l-1})(c),&c\in\Gamma_l^{\mathrm{ti}}.
 \end{cases}
\end{equation}
The average in the second line uses only
\(\Gamma_{l-1}^{\mathrm{pr}}\): both blocks over a bond with no
endpoint in \(\Omega_l\) lie outside that region. In particular,
completion is a Borel map defined on the whole free space. Restriction
\(\mathrm{rs}_h\) to the free coordinates satisfies
\(\mathrm{rs}_h\mathrm{cp}_h=\mathrm{id}\).
The image of completion is precisely the set on which the finitely
many tying equations in \eqref{II-58-5} hold. It is Borel, since each
equation equates Borel maps into the compact metric space \(G\).

Give the configuration coordinates Haar probability measure and the
\(\mathfrak q\)-coordinates Lebesgue measure for
\(\langle X,Y\rangle=\frac12\operatorname{tr}(XY)\). Define
\begin{equation}\label{II-58-48}
 d\mathscr M_h^{\mathrm{fr}}
 =\prod_l dv_l\prod_{l<k}d\mathfrak q_l,\qquad
 \mathscr M_h=(\mathrm{cp}_h)_*\mathscr M_h^{\mathrm{fr}}.
\end{equation}
This measure is positive and \(\sigma\)-finite. Indeed, bounding all
the finitely many Lie algebra coordinates gives Borel sets of finite
measure which exhaust the completed space up to its null complement;
completion and restriction preserve these coordinates. This last
observation matters: an arbitrary image of a \(\sigma\)-finite
measure need not be \(\sigma\)-finite.

For a bare field \(U\) for which all the required strip variables
exist, its completed coordinate is
\[
 \xi_h(U)=
 \bigl((\mathcal M^lU)|_{\Gamma_l^{\mathrm{pr}}},
             (\mathsf a'_l(U))_{l<k}\bigr).
\]
It is the completion of the corresponding free coordinate. We use
partial pushforwards on their Borel domains: for a finite measure
\(\nu\) and a Borel partial map \(F\),
\(F_*\nu(E)=\nu(\{U\in\operatorname{dom}F:F(U)\in E\})\).
In particular the history measure is
\((\xi_h)_*(\mathfrak w_h\rho_0\,dU)\).

At the step \(t=k+1\), let \(\mathfrak X_0\) consist of all
level-\(k\) bonds with an endpoint in \(\Omega_t\), including pivot
bonds. The old completed coordinates split as
\[
 \phi=(V^{\mathrm{in}},\phi^+)\in
       \mathcal V(\mathfrak X_0)\times\mathsf X^+_{h'}.
\]
Here \(\phi^+\) contains all lower proper coordinates, the
level-\(k\) coordinates with no endpoint in \(\Omega_t\), and all
old external fields. New free coordinates are
\(W\in\mathcal V(\Gamma_t^{\mathrm{le}}(h'))\) and
\(\mathsf a'\in\mathfrak g^{\mathfrak X_k^{\mathrm E+}(h')}\).
If \(v_k\) is the level-\(k\) component of \(\phi^+\), complete
the new level by
\[
 v_t(c)=W(c)\quad(c\in\Gamma_t^{\mathrm{le}}),\qquad
 v_t(c)=\mathcal M(v_k)(c)\quad(c\in\Gamma_t^{\mathrm{ti}}).
\]
These formulas define a Borel injection
\(\imath^+_{h'}(W,\mathsf a',\phi^+)\) into the new completed
space; its inverse on its image is coordinate projection.
For \(c\in\Gamma_t^{\mathrm{ti}}\) the double block lies outside
\(\Omega_t\), so this formula never calls for \(V^{\mathrm{in}}\).

\begin{proposition}\label{II-58-74}
Let \(\mathscr M^+_{h'}\) be the measure on the untouched coordinates
obtained from \(\mathscr M_h\). Then
\[
 d\mathscr M_h=dV^{\mathrm{in}}\,d\mathscr M^+_{h'},\qquad
 \mathscr M_{h'}=(\imath^+_{h'})_*
       (dW\,d\mathsf a'\,d\mathscr M^+_{h'}).
\]
Both measures in these formulas are \(\sigma\)-finite.
\end{proposition}

\begin{proof}
The old free space splits as
\(\mathcal V(\mathfrak X_0)\times\mathsf Y^+\).
Completion has the form \(\mathrm{id}\times\Pi^+\):
none of the lower tying equations involves the current
level-\(k\) inner bonds. Thus its image measure has the first
displayed product decomposition. The map \(\Pi^+\) preserves every
old Lie algebra coordinate, which proves \(\sigma\)-finiteness of
\(\mathscr M^+_{h'}\) by the same bounded-coordinate exhaustion.
Completion of the new free space is
\[
 \imath^+_{h'}\circ
       (\mathrm{id}_W,\mathrm{id}_{\mathsf a'},\Pi^+).
\]
Its image measure is the second displayed measure. These identities
hold first on measurable rectangles and then on Borel sets by
uniqueness of the product measure.
\end{proof}


\subsection{The regional fibre chart}
\label{sec:regional-history-chart}

We need the constraint chart with a fixed exterior configuration.
Let \(P''\) be a set of level-\(t\) sites, \(\Gamma\) the coarse
bonds with an endpoint in \(P''\), and \(\mathfrak X_0\) the fine
bonds with an endpoint in \(B(P'')=\bigcup_{y\in P''}B(y)\).
Let \(\mathfrak X\) be its nonpivot bonds, and let
\(\mathfrak X_{\mathrm{out}}\) be the remaining fine bonds in
the double blocks of \(\Gamma\). Fix
\(\eta^+\in\mathcal V(\mathfrak X_{\mathrm{out}})\), and write
\[
 q(V^{\mathrm{in}},\eta^+)
       =\mathcal M(V^{\mathrm{in}}\uplus\eta^+)|_\Gamma .
\]
Choose \(V_0\in\mathcal V(\mathfrak X_0)\) so that
\(V=V_0\uplus\eta^+\) has double-block plaquette deviations
\(\varpi_c(V)<\varpi_\chi\) for every \(c\in\Gamma\).
The constants \(\varpi_\chi,r_\chi,C_\chi\) are the local chart
constants used above, with \(0<r_\chi\le1\) and \(C_\chi\ge1\),
decreased or increased, as appropriate, depending only on \(L\).

\begin{proposition}\label{II-26-53}
Extend a real nonpivot field \(\mathcal A\) by zero on
\(\mathfrak X_{\mathrm{out}}\). If
\(C_\chi\sup|\mathcal A|<r_\chi\), the chart
\[
 B'=\psi_V^{\mathrm{lin}}(\mathsf C_V\mathcal A),\qquad
 V[\mathcal A]=(e^{iB'}V)|_{\mathfrak X_0}
\]
is defined, is continuous in all its data, and satisfies
\[
 B'|_{\mathfrak X}=\mathcal A,\quad
 B'|_{\mathfrak X_{\mathrm{out}}}=0,\quad
 \sup|B'|\le2C_\chi\sup|\mathcal A|,\quad
 q(V[\mathcal A],\eta^+)=q(V_0,\eta^+).
\]
Conversely, any \(V^{\mathrm{in}}\) in this fibre whose nonpivot
logarithms relative to \(V_0\) satisfy the same bound equals
\(V[\mathcal A]\), where
\(\mathcal A(b)=i^{-1}\log(V^{\mathrm{in}}(b)V_0(b)^{-1})\)
on the positive nonpivot bonds. No bound on its pivot entries is required in
this converse.

Put \(n_c(V)=L\mathsf m_c(V)\) and
\(\mathsf z_c(V)=-\log\det n_c(V)\). The determinants are positive,
so the numbers \(\mathsf z_c(V)\) are real. At \(\mathcal A=gB\) define
\begin{equation}\label{II-26-54}
 \mathsf J_V[gB]
   =\sum_{c\in\Gamma}\log\det(1-\mathsf{Jb}_c(V,gB))
       +\sum_{b\in\mathfrak X_0^+}\log\sigma_1(|B'(b)|^2).
\end{equation}
Let \(g=\mathsf t^{-1/2}>0\), \(C_\chi g r\le r_\chi\), and
let \(W\mapsto V_{0,W}\), \(W\in O\subset G^\Gamma\), be a
Borel family of such references with
\(q(V_{0,W},\eta^+)=W\). Write
\(V_W=V_{0,W}\uplus\eta^+\),
\(V_W(B)=V_W[gB]\), and
\(\mathsf J_W(B)=\mathsf J_{V_W}[gB]\).
If \(\Upsilon(V^{\mathrm{in}},W)\ge0\) is Borel and vanishes
unless \(W\in O\) and
\[
 \left|\frac1i\log(V^{\mathrm{in}}(b)V_{0,W}(b)^{-1})\right|
       <g r\quad(b\in\mathfrak X^+),
\]
then
\begin{equation}\label{II-26-55}
\begin{split}
 &\int \Upsilon(V^{\mathrm{in}},
                     q(V^{\mathrm{in}},\eta^+))\,dV^{\mathrm{in}}\\
 &\quad=(\sigma_{\mathrm H}\mathsf t^{-3/2})^{\#\mathfrak X^+}
  \int_O dW\,e^{\sum_{c\in\Gamma}\mathsf z_c(V_W)}
  \int_{\sup|B|<r}e^{\mathsf J_W(B)}
                   \Upsilon(V_W(B),W)\,dB .
\end{split}
\end{equation}
The same identity holds for complex \(\Upsilon\) if either side
with \(|\Upsilon|\) is finite.
\end{proposition}

\begin{proof}
The constraint calculation is local to a double block. On
\(\mathsf P_c\), put \(U_c=V|_{\mathsf P_c}\). The contraction
used in \cref{II-41-33} gives
\[
 \widetilde{\mathsf D}_V(\mathsf B)(c)
       =n_c(V)y_{U_c}(\mathsf B|_{\mathsf P_c}),\qquad
 \psi_V^{\mathrm{lin}}(\mathsf B)(b_0(c))
       =\mathsf B(b_0(c))-y_{U_c}(\mathsf B|_{\mathsf P_c}).
\]
The maps and bounds are unchanged when exterior nonpivot
entries are fixed at zero: each double block contains just
its own pivot. In particular
\(\|\mathsf C_V\|_{\infty\to\infty}\le C_\chi\),
the real solution has full norm at most
\(2C_\chi\sup|\mathcal A|\), and it solves the averaging
constraint at every \(c\). Its plaquette deviations are bounded
by \(\varpi_\chi+8C_\chi\sup|\mathcal A|\).
Choose the constants so that this is below the small-curvature
ceiling. The same local contraction proves continuous
dependence on \(V_0,\eta^+,\mathcal A\).

The converse uses the global uniqueness in \cref{II-6-12}.
The proposed chart and \(V^{\mathrm{in}}\) agree at every
nonpivot entry and have the same averages. For a given \(c\),
their restrictions to \(\mathsf P_c^\circ\) are within
\(r_\chi/C_\chi\) of the reference. Choose
\(\varpi_\chi+r_\chi/C_\chi\le\omega_{\mathrm p}\).
The pivot proposition then says that the same coarse value has
exactly one pivot preimage, including the other real branches
of the global average. Hence the two configurations agree also
at every pivot.

The differential and determinant computation
\eqref{eq:sf-pivot-Jacobian} applies to this same double block.
It gives
\[
 (\det n_c(V_W(B)\uplus\eta^+))^{-1}
  =e^{\mathsf z_c(V_W)}
       \sigma_1(|B'(b_0(c))|^2)
       \det(1-\mathsf{Jb}_c(V_W,gB)).
\]
All factors are positive; the local bounds give
\(\|\mathsf{Jb}_c\|\le1/2\). The pivot subspace is a direct
sum over \(c\), so its logarithmic determinant is the sum in
\eqref{II-26-54}.

To justify the integral for a merely Borel family, first fix
the nonpivot entries \(w\). The permitted pairs \((w,W)\)
form a Borel set, since the intrinsic logarithm is continuous
on \(G\setminus\{-1\}\). The chart supplies its Borel pivot
solution \(\Pi(w,W)\). On the small-curvature pivot domain,
the map from pivots to averages is a diffeomorphism onto its
image: injectivity is the pivot proposition, and its
differential is invertible. The converse just proved shows
that every nonzero value of the integrand occurs in this
domain, even though no pivot cutoff was imposed.
Changing the pivot variables to \(W\) therefore gives the
product of the inverse determinants in the preceding display.
Tonelli interchanges this integral with the nonpivot integral.
For each \(W\), substitute
\(w(b)=e^{igB(b)}V_{0,W}(b)\). The nonpivot Haar factors are
\[
 (\sigma_{\mathrm H}g^3)^{\#\mathfrak X^+}
       \prod_{b\in\mathfrak X^+}\sigma_1(g^2|B(b)|^2)\,dB.
\]
Together with the pivot factors they give precisely
\eqref{II-26-55}. All maps used in this computation are Borel,
and Tonelli requires no finiteness assumption. Applying the
positive formula to \(|\Upsilon|\), followed by Fubini and
linearity, proves the complex version.
\end{proof}

The gauge factor has the same regional interpretation. If \(V_0\)
is axial on \(P''\), then the averaged contour bound gives
\[
 |V_W(B)(y,x)-1|
 \le C_\chi g\max_{b\subset B(y)}|B(b)|
 \quad(y\in P'',\ x\in B(y)).
\]
When these bounds are below \(\epsilon\), all gauge cutoffs pass,
and the product of the \((L^3-1)\#P''\) factors is
\begin{equation}\label{eq:regional-gauge-factor}
 \mathfrak G(V_W(B);P'';\mathsf t,\epsilon)
 =\mathsf z(\mathsf t,\epsilon)^{-(L^3-1)\#P''}
          e^{-\mathsf tG_k(P'';V_W(B))}.
\end{equation}
Indeed the block contains no pivot bond; its logarithmic perturbation
is just \(gB\), so the local averaged-contour estimate applies there
directly. Multiplying the defining single-contour factors proves
the identity.

\subsection{The strip substitution}
\label{sec:history-strip-substitution}

The substitution changes the reference background of an external
bond. It must preserve the tests which determined the history.
Write \(\mathfrak S=\Omega_t\setminus\Omega_t^{\sim-1}\) and
\(S_t^{\mathrm{pr}}=\mathfrak S\cup R\), where \(R\) is the union
of failed fluctuation-test cubes. With the test-cube enlargements
already fixed, set
\[
 \mathcal N_{\mathrm s}
  =\bigcup_{\square'\subset\Omega_t^{\sim-1}\setminus R}
           (\square'^{\sim2})^{(k)*},\qquad
 \mathcal N_\ell
  =\bigcup_{\square'\subset R}(\square'^{\sim2})^{(k)*}.
\]
The two zones are
\begin{equation}\label{II-26-74}
 \mathcal Z_\theta
  =(S_t^{\mathrm{pr}(k)*}\cap\mathcal N_\ell)
      \setminus\mathcal N_{\mathrm s},\qquad
 \mathcal Z_1
  =S_t^{\mathrm{pr}(k)*}
      \setminus(\mathcal N_{\mathrm s}\cup\mathcal N_\ell).
\end{equation}
They are disjoint reversal-invariant sets of external nonpivot bonds;
put \(\mathcal Z=\mathcal Z_\theta\cup\mathcal Z_1\).
The plain zone \(\mathcal Z_1\) meets no bond used in an external
fluctuation test. On \(\mathcal Z_\theta\) we will leave all
fluctuations below the test thresholds unchanged.

Here are the geometric facts used in this assertion. Put
\(X_\flat=\Omega_t^c\cup\mathfrak S\cup R'\), where \(R'\)
is the coarse hull of \(R\). Every region \(Z\) encountered in
the exploration contains
\((\Omega_t^c)^{\sim2}\cup(R')^\sim\), hence contains
\(X_\flat^\sim\). No coarse cube of \(X_\flat\) touches a
cube of \(Z^c\). A strip bond has an endpoint in \(X_\flat\);
its length is one, whereas the coarse side is \(W_t\), so it
has no endpoint in \(Z^c\). Moreover every bond of
\(\mathfrak X\setminus\mathcal N_{\mathrm s}\) has an endpoint
in \(\mathfrak S\cup R\), by considering the test cube of
that endpoint. Consequently
\[
 \mathfrak X=\mathcal N_{\mathrm s}\sqcup\mathcal Z,\qquad
 \mathcal Z\subset\mathfrak X_{\mathrm E},\qquad
 \operatorname{dist}_\infty(b,Z^c)\ge W_t-1
       \quad(b\in\mathcal Z).
\]
The last bound also separates a zone bond from the endpoints
of integrated bonds by \(W_t-1\). Since
\(\varpi_T<W_t-1\), every explored collar lies in
\(\mathcal N_{\mathrm s}\). A rim cube touches \(Z^c\)
and cannot belong to \(R'\); its tested bonds likewise belong
to \(\mathcal N_{\mathrm s}\). This holds at each stage, so
\(E^T\subset\mathcal N_{\mathrm s}\) and
\(c^\star\subset\mathcal N_{\mathrm s}\) for \(c\in E^H\).
Finally, an interior test cube belongs either to \(R\) or to
its complement, so its tested bonds lie in
\(\mathcal N_\ell\cup\mathcal N_{\mathrm s}\), as do all the
sets \(c^\star\). By its definition, \(\mathcal Z_1\) is
disjoint from this union. This proves the claimed avoidance
of all external test bonds.

First consider one bond. Fix \(0<\delta\le1\), a smooth cutoff
\(\vartheta_{\mathrm{co}}:[0,\infty)\to[0,1]\) equal to zero
on \([0,4/3]\) and to one on \([5/3,\infty)\), and \(H\in\mathfrak g\).
In the cutoff case assume
\(8\|\vartheta'_{\mathrm{co}}\|_\infty|H|\le\delta\), and put
\[
 \theta(e^{iA'})=\vartheta_{\mathrm{co}}(|A'|/\delta)
       \quad(|A'|<\pi),\qquad \theta(-1)=1.
\]
In the plain case put \(\theta=1\), with no restriction on \(H\).
Define
\begin{equation}\label{II-26-68}
 \Psi_H(W')=W'e^{-i\theta(W')H}e^{iH},\qquad
 f_H(A')=\frac1i\log\bigl(e^{iA'}e^{-i\theta(e^{iA'})H}\bigr).
\end{equation}
The logarithm is intrinsic and is defined unless its argument is
\(-1\).

\begin{lemma}\label{II-26-71}
The map \(\Psi_H\) is a diffeomorphism of \(G\).
Its Haar Jacobian is \(1\) in the plain case and, in the cutoff
case,
\begin{equation}\label{eq:strip-bond-jacobian}
 J_H(e^{iA'})
 =1-\frac{\vartheta'_{\mathrm{co}}(|A'|/\delta)}
                  {\delta|A'|}\langle A',H\rangle
 \in[7/8,9/8],
\end{equation}
where the quotient term is zero at \(A'=0\), and \(J_H(-1)=1\).
The map \(f_H\) is continuous and injective on its domain in
\(\{|A'|<\pi\}\); its domain and its image omit at most one point
of that ball. For \(0<s\le4\delta/3\), in the cutoff case,
\[
 |f_H(A')|<s\quad\Longleftrightarrow\quad |A'|<s,
 \qquad f_H(A')=A'\ \text{there}.
\]
In both cases, wherever \(f_H\) is defined,
\[
 \bigl||f_H(A')|-|A'|\bigr|\le\theta(e^{iA'})|H|.
\]
\end{lemma}

\begin{proof}
For \(\alpha_{W'}(X)=\frac d{ds}|_{s=0}\theta(W'e^{isX})\),
differentiating
\(\frac12\operatorname{tr}(e^{iA'}e^{isX})\) gives
\[
 \alpha_{e^{iA'}}(X)=
 \frac{\vartheta'_{\mathrm{co}}(|A'|/\delta)}
      {\delta|A'|}\langle A',X\rangle .
\]
The cutoff is locally constant near \(1\) and \(-1\), so this
also proves smoothness at those points. If
\(T=e^{-i\theta(W')H}e^{iH}\), the right-trivialized differential is
\[
 X\longmapsto T^{-1}(X-\alpha_{W'}(X)H)T.
\]
Conjugation preserves volume, and the rank-one factor has determinant
\(1-\alpha_{W'}(H)\). The smallness hypothesis gives
\(|\alpha_{W'}(H)|\le1/8\), proving the Jacobian formula.
For completeness, global invertibility follows from a contraction.
Since
\(\varphi(W')=2\arcsin(|W'-1|/2)\), the cutoff has Lipschitz
constant \(2\|\vartheta'_{\mathrm{co}}\|_\infty/\delta\):
on its nonconstant interval the derivative of
\(2\arcsin(s/2)\) is less than \(2\), and elsewhere the cutoff
is constant. Moreover \(\|\vartheta'_{\mathrm{co}}\|_\infty\ge3\),
so \(|H|\le1/24\). For each \(W\in G\), the equation
\(\Psi_H(W')=W\) is equivalent to
\[
 W'=T_W(W'):=We^{-iH}e^{i\theta(W')H}.
\]
The estimate \(|e^{icH}-1|\le2|c||H|\) for \(|c|\le1\)
shows that \(T_W\) is a contraction with constant at most \(1/2\).
It has a unique fixed point on the complete metric space \(G\).
Thus \(\Psi_H\) is bijective. Its differential is invertible by
\eqref{eq:strip-bond-jacobian}, so its inverse is smooth.
The plain map is the identity.
Right translation by \(e^{-iH}\), followed by the intrinsic
logarithmic chart, gives \(f_H\). The excluded point of either
chart is unique, proving the assertions about its domain and image.
If \(|A'|\le4\delta/3\), then \(\theta(e^{iA'})=0\) and
\(f_H(A')=A'\). Conversely, injectivity gives the same equality
whenever the image has norm below that threshold.
Finally the group-angle triangle inequality applied to
\(e^{if_H(A')}=e^{iA'}e^{-i\theta H}\) and its inverse relation
gives the last estimate.
\end{proof}

For the step, fix a real shift \(\mathfrak H=(H_b)_{b\in\mathcal Z^+}\)
with
\begin{equation}\label{II-26-84}
 8\|\vartheta'_{\mathrm{co}}\|_\infty|H_b|\le\delta_k
 \quad(b\in\mathcal Z_\theta^+).
\end{equation}
Take the cutoff case at \(b\in\mathcal Z_\theta\), the plain case
at \(b\in\mathcal Z_1\), and put
\begin{equation}\label{II-26-83}
 \check\Psi(\mathsf a',\mathfrak H)(b)=
 \begin{cases}
 g_k^{-1}f_{H_b}(g_k\mathsf a'(b)),&b\in\mathcal Z^+,\\
 \mathsf a'(b),&b\in\mathfrak X^{\mathrm E+}\setminus\mathcal Z.
 \end{cases}
\end{equation}
Extend by antisymmetry. Its real domain requires
\(g_k|\mathsf a'(b)|<\pi\) on zone bonds and the definedness
of the logarithms. Set
\begin{align}
 \mathsf{Zs}_{\mathrm J}
 &=\sum_{b\in\mathcal Z_\theta^+}
           \log J_{H_b}(e^{ig_k\mathsf a'(b)}),\label{II-26-87}\\
 \mathsf{Zs}_{\mathrm L}
 &=\sum_{b\in\mathcal Z^+}
  \log\frac{\sigma_1(g_k^2|\mathsf a'(b)|^2)}
                {\sigma_1(g_k^2|\check\Psi(\mathsf a',\mathfrak H)(b)|^2)}.
 \label{eq:strip-logarithmic-density}
\end{align}

\begin{proposition}\label{II-26-92}
The strip map is injective and Borel on its domain.
Its undefined set in the zone logarithmic charts is Lebesgue-null.
For every nonnegative Borel \(F\) vanishing unless
\(g_k|\mathsf a(b)|<\pi\) at all external bonds,
\begin{equation}\label{II-26-93}
 \int F(\mathsf a)\,d\mathsf a
 =\int_{\operatorname{dom}\check\Psi}
    e^{\mathsf{Zs}_{\mathrm J}+\mathsf{Zs}_{\mathrm L}}
       F(\check\Psi(\mathsf a',\mathfrak H))\,d\mathsf a',
\end{equation}
where the domain on the right is restricted by
\(g_k|\mathsf a'(b)|<\pi\) at every external bond.
The external fluctuation-test product takes the same value at
\(g_k\check\Psi(\mathsf a',\mathfrak H)\) and \(g_k\mathsf a'\).
\end{proposition}

\begin{proof}
On one bond the Haar density in logarithmic coordinates is
\(\sigma_{\mathrm H}\sigma_1(|A|^2)\,dA\).
Apply Haar substitution by \(\Psi_H\), and then right translation
by \(e^{-iH}\). Dividing by the positive density
\(\sigma_1(|f_H(A')|^2)\), and scaling \(A'=g_k\mathsf a'\),
gives the Lebesgue factor
\[
 J_H(e^{ig_k\mathsf a'})
   \frac{\sigma_1(g_k^2|\mathsf a'|^2)}
        {\sigma_1(g_k^2|g_k^{-1}f_H(g_k\mathsf a')|^2)}.
\]
The factors \(g_k^3\) on the two sides cancel.
The product of these factors is
\(e^{\mathsf{Zs}_{\mathrm J}+\mathsf{Zs}_{\mathrm L}}\).
Apply the one-bond identity successively, using Tonelli for
nonnegative functions. Each omitted set fixes one of the finitely
many three-dimensional coordinates to a single point and is null.
Off the zones the map is the identity, so the restriction to the
logarithmic ball at every external bond loses no nonzero integrand.
This proves \eqref{II-26-93} without an integrability assumption.

The tested thresholds are
\(\delta_k,\delta_k^\flat,\theta_c\delta_k\), all below
\(4\delta_k/3\). On the cutoff zone their indicators are preserved
by Lemma~\ref{II-26-71}; off both zones the variables are unchanged.
No tested bond lies in the plain zone. Each external test is a
product or difference of these indicators, which proves the
last assertion.
\end{proof}

\subsection{Localization of the actual strip background}
\label{sec:actual-strip-localization}

The shift in the preceding substitution is the difference between two
critical backgrounds. We now construct the second background and
estimate that difference. The comparison first takes place before
averaging, up to gauge. The axial normalization then turns the estimate
at the nodes into an estimate for the two averaged configurations.

We use reduced determining sets throughout. Fix the joint averaging radius
\[
 r_{\mathrm{av}}=\min\{r_{\mathrm{st}},\kappa_{6,\mathsf m}^{-1}\},
 \qquad \kappa_6=r_{\mathrm{av}}^{-1},\qquad
 \varpi_{\mathrm{av}}=\min\{\varpi'_L,r_{\mathrm{av}}/(4\pi L)\}.
\]
Here \(r_{\mathrm{st}}\) is from \cref{lem:one-step-straightening},
and \(\kappa_{6,\mathsf m}\) is the centered contour constant in
the proof of \cref{I-39-36}. In particular \(\kappa_6\ge16\).
For a real double-block reference of curvature at most
\(\varpi_{\mathrm{av}}\), the comb gauge of \cref{I-10-51}
has logarithm of norm at most \(r_{\mathrm{av}}/2\).
The joint holomorphy, unit bound and contour estimate of
\cref{lem:one-step-straightening} therefore apply on the product
of the reference and perturbation balls of radius \(r_{\mathrm{av}}\).
Denote by \(a_{\mathrm R}>0\) and \(\varrho_{\mathrm R}^0>0\)
the regular-reference and datum thresholds for the response in
\cref{I-82-15}. Decrease the regular-reference threshold so that
\(a_{\mathrm R}\le\min\{L^{-2}a_{\mathsf m},
\varpi_{\mathrm{av}}/(2L^2)\}\), as well as its earlier ceiling.
The regional curvature estimate in \cref{I-39-36} then bounds the
intermediate reference curvatures by \(\varpi_{\mathrm{av}}\).

Set
\[
 c_{\mathcal Q}=\min\{c_{\mathsf m}/L,
                  r_{\mathrm{av}}/(64L),(8\kappa_6)^{-1}\}.
\]
The centered iteration is defined on each pillar ball
\(\xi_y\sup|a|<c_{\mathcal Q}\). The datum conversion in the proof
of \cref{I-82-15} has a uniform bound
\(\sup|B_\natural|\le C_{\mathrm{dat}}\sup|B|\), with
\(C_{\mathrm{dat}}\) depending only on \(L\).
Thus \eqref{I-87-38} gives
\(\sup_y\xi_y|\mathcal H(B)|_y\le C_H\sup|B|\), where
\(C_H=4\hat{\mathsf b}C_{\mathrm{dat}}\).
Decrease \(r_{\mathrm B}\) so that
\(C_Hr_{\mathrm B}\le c_{\mathcal Q}/4\).
The response then lies in the pillar quarter-balls used below.

The response constants in
\cref{I-82-15,I-93-47} can be chosen with a rate
\(0<\delta_{\mathrm{lc}}\le1/8\) and a constant \(B_5\) such that,
on the real datum ball of radius
\(\varrho=\min\{r_{\mathrm B},\varrho_{\mathrm R}^0\}\),
\begin{equation}\label{II-8-18}
 m_y(B)=\sum_{c>0}e^{-\delta_{\mathrm{lc}}d(y,y_c)}|B(c)|,
 \qquad
 \xi_y|\mathcal H(B)(b)|\le B_5m_y(B),
 \qquad
 |u_B(y)-1|\le4\kappa_6B_5m_y(B).
\end{equation}
The bond in the middle estimate has an endpoint in the pillar of \(y\);
the last estimate is at a top \(y\). These conclusions require
\(C_{\mathfrak P}\mathfrak M B_3\varepsilon_{\mathrm r}
\le a_{\mathrm R}\), as well as the regular-data hypotheses of
\cref{I-82-15}.
The middle estimate follows by integrating \eqref{I-93-47} along
\(sB\). The same argument, applied to the first-difference estimate
in \eqref{I-93-47}, gives
\begin{equation}\label{II-8-20}
 \xi_y^2|(\nabla^{U^\circ}\mathcal H(B))(x)|
                  \le B_5m_y(B),\qquad y(x)=y .
\end{equation}
The two bounds also hold for the difference of responses at \(B_1,B_2\),
with \(m_y(B_1-B_2)\) on the right: integrate the derivative along
the segment joining the two points in the real datum ball.
To verify the last estimate, recall that
\(u_B(y)^{-1}=v_q[\mathcal H(B)](y)\).
This factor depends only on the pillar of \(y\). On the ball
\(\xi_y\sup|a|<c_{\mathcal Q}\), its factors have logarithms \(Y_m\)
with
\[
 |Y_m|\le2\kappa_6L^{m-q}\xi_y\sup|a|,
 \qquad
 \sum_{m<q}|Y_m|\le\frac{2\kappa_6}{L-1}\xi_y\sup|a|.
\]
These are the one-step contour estimate and the centered averaging
bounds used in \cref{I-27-78}. Thus both \(v_q[a](y)^{\pm1}-1\)
are bounded by \(2\kappa_6c_{\mathcal Q}\) on this ball.
Cauchy's inequality bounds their derivatives by \(4\kappa_6\) on
the half-ball. The response takes values in the quarter-ball.
The chain rule, \eqref{I-93-47}, and integration along \(sB\)
give the last bound in \eqref{II-8-18}.

Here is the localization statement in the form needed twice below.
Let \((\Omega_i)_{i\le n}\) be a separated sequence, and let
\(\Omega\subset\Omega_0\) have complement
\(Y_1\cup\cdots\cup Y_r\), \(r=1\) or \(2\).
Associate a level \(l_s\le n\) to each part.
The \(Y_s\) are regions of \(M_1\)-cubes at level \(l_s\), and
enlargements of \(Y_s\) in this paragraph use that partition.
Suppose
\begin{equation}\label{II-8-53}
 Y_s^{\sim4}\setminus Y_s
       \subset\Omega_{l_s}\setminus\Omega_{l_s+1},
 \qquad
 \operatorname{dist}_\infty(Y_s,Y_{s'})
       \ge10M_1L^{\max(l_s,l_{s'})}\quad(s\ne s'),
\end{equation}
where these lengths are in bare-lattice units.
Set
\[
 C_s=Y_s^{\sim2}\setminus Y_s,\qquad
 C_s^{\rm w}=Y_s^{\sim3}\setminus Y_s,\qquad
 \Omega^\natural=\Omega\setminus\bigcup_s C_s.
\]
Let \(\Omega^{[i]}\) be the \(\mathfrak r_2\)-tower of \(\Omega\).
The joined sequence and datum are
\begin{equation}\label{eq:joined-localization-sequence}
 \Omega'_i=(\Omega_i\cap\Omega^\natural)
       \cup\bigcup_{l_s\ge i}(\Omega^{[i]}\cap C_s),
 \qquad
 V'_i(c)=
 \begin{cases}
 V_i(c),&c\in\Gamma'_i\cap\Gamma_i,\\
 (Q^{s*}_{l_s\to i}V_{l_s})(c),&c\in\Gamma'_i\setminus\Gamma_i .
 \end{cases}
\end{equation}
In the second line \(s\) is the part of the tower bond \(c\).
Assume that this sequence is separated, that each tower bond has
both endpoints in one wide collar, and that both data are regular
at \(\varepsilon_{\mathrm r}\).

Choose the geometric floors for the response and exponential
summation bound, in particular
\(\delta_{\mathrm{lc}}R_1\mathfrak M
\ge8(3\log L+\log2)\).
Let \(A_L=3(L-1)^2/4\) and let
\[
 C_{\rm AC}=8(L-1)A_L+
       (1+\pi/3)\bigl(L(L-1)+8A_L\bigr),\qquad
 C_{\rm al}=\frac{\pi C_{\rm AC}B_3}{L^2-1}.
\]
Choose the one-step curvature threshold \(\omega_\sharp\) below
both the mean threshold divided by \(A_L\) and the averaging
threshold. It suffices to require
\begin{equation}\label{II-8-59}
 \varepsilon_{\mathrm r}\le\epsilon_{\mathrm{lc}}
 :=\min\left\{
 \bar\varepsilon,\frac{a_{\mathrm R}}{C_{\mathfrak P}\mathfrak M B_3},
 \frac{c_2^{\mathrm{SO}}}{B_3},
 \frac{L^2\omega_\sharp}{4B_3},
 \frac{\varrho}{2C_{\rm al}}\right\}.
\end{equation}
All thresholds here are the real-background, response, mean, and
curvature-iteration thresholds already fixed; none depends on the
depth or volume.

\begin{proposition}\label{II-8-68}
Under these hypotheses, and the size condition
\(2L^{m+K-n}\ge6R_1\mathfrak M\), let \(U,U'\) be arbitrary
elements of the two minimal orbits. There are gauges \(w,w'\)
on the bare sites of \(\Omega\) and a real odd field \(a\)
on its inner bonds such that
\begin{equation}\label{II-8-69}
 (U|_\Omega)^w=e^{ia}(U')^{w'}.
\end{equation}
Write \(D(z)\) for the weighted distance in the joined sequence
from \(z\) to the endpoints of its tower bonds, and
\(D(b)=\min\{D(b_-),D(b_+)\}\). An empty set of tower bonds
has distance \(+\infty\).
For \(j(b)=\max\{j:b_-\in\Omega'_j\text{ or }b_+\in\Omega'_j\}\),
the same triple satisfies, outside the wide collars,
\begin{align}
 |a(b)|&\le C_{\mathrm{lc}}\varepsilon_{\mathrm r}
             L^{-j(b)}e^{-\delta_{\mathrm{lc}}D(b)/2},
             \label{II-8-70}\\
 w'(y)&=1,\qquad
 |w(y)-1|\le C_{\mathrm{lc}}\varepsilon_{\mathrm r}
             e^{-\delta_{\mathrm{lc}}D(y)/2}
       \quad(y\text{ a node}).\label{II-8-71}
\end{align}
One may take
\(C_{\mathrm{lc}}=4\kappa_6B_5
S(\delta_{\mathrm{lc}}/2)C_{\rm al}\),
where \(S\) is the exponential row-sum constant.
\end{proposition}

\begin{proof}
The joined sequence is subordinate to the original sequence.
The restriction statement \cref{II-8-36} shows that \(U|_\Omega\)
is regular and critical for its own datum \(V^{\rm rs}\), and
that \(V^{\rm rs}=V'\) on the common rows.
We first align the remaining data in each wide collar.

In the collar of level \(l\), regularity gives bare curvature
at most \(B_3\varepsilon_{\mathrm r}L^{-2l}\).
The curvature averaging estimate gives
\(\varpi_q\le2B_3\varepsilon_{\mathrm r}L^{2(q-l)}\) for
\(q\le l\). Recursively set \(v_l=1\) and
\[
 v_q(x)=v_{q+1}(z)(\mathcal M^qU)(z,x),\qquad x\in B(z).
\]
Gauge covariance makes all the one-step averaged contour
variables of \(\widehat W_q=\mathcal M^q(U^{v_0})\) equal to one.
We recall the resulting elementary double-block estimate.
If a configuration has curvature at most \(a\), its path
transports in a block differ by at most \(A_La\); when their
mean is one each transport is within \(4A_La\) of one.
Appending one final edge gives the bound \(8A_La\) on an
interior bond. For a crossing bond, commute the \(L\) parallel
steps past the transverse part of a path. There are at most
\(L(L-1)\) plaquette exchanges. The line transport therefore
differs from the central transport by
\((L(L-1)+8A_L)a\).
The logarithmic mean contributes at most \(\pi/3\) times this
bound, and the other \(L-1\) line bonds contribute
\(8(L-1)A_La\). Hence
\[
 |\widehat W_q(b)\widehat W_{q+1}(c)^{-1}-1|
       \le C_{\rm AC}\varpi_q
\]
when \(b\) crosses the blocks of \(c\).
Inside one block the lift is one. Descending induction consequently
gives
\[
 |\widehat W_q(b)(Q^{s*}_{l\to q}V_l)(b)^{-1}-1|
 \le C_{\rm AC}\sum_{i=q}^{l-1}\varpi_i
 \le\frac{2C_{\rm AC}B_3}{L^2-1}\varepsilon_{\mathrm r}.
\]
All blocks used are contained in the wide collar. Its top bonds
are common rows, so their averages are \(V_l\).

The wide collars are disjoint. Glue these gauges to the identity
outside them, obtaining \(g\). At the endpoints of a common row
in a wide collar, the endpoint is a level-\(l_s\) site and the
collar gauge is one. Thus the common datum is unchanged, while
the datum of \((U|_\Omega)^g\) is \(e^{iB_1}V'\), with
\[
 B_1=0\text{ on common rows},\qquad
 \sup|B_1|\le C_{\rm al}\varepsilon_{\mathrm r}\le\varrho/2.
\]
The intrinsic logarithm exists because the preceding distance
is less than two; its norm is at most \(\pi/2\) times that distance.

Let \(U^\circ\) be the normalized background at \(V'\).
Uniqueness of the regular critical orbit and \cref{I-82-15} give
residual gauges \(w_0,w_1\) with
\[
 (U|_\Omega)^g
   =\bigl((e^{i\eta_n\mathcal H(B_1)}U^\circ)^{u_{B_1}}\bigr)^{w_0},
 \qquad U'=(U^\circ)^{w_1}.
\]
Put \(w=(w_0u_{B_1})^{-1}g\), \(w'=w_1^{-1}\), and
\(a=\eta_n\mathcal H(B_1)\). This proves \eqref{II-8-69}.
Since \(B_1\) is supported on tower rows, exponential summation gives
\[
 m_y(B_1)\le C_{\rm al}\varepsilon_{\mathrm r}
       S(\delta_{\mathrm{lc}}/2)e^{-\delta_{\mathrm{lc}}D(y)/2}.
\]
For an endpoint of \(b\) in layer \(j(b)\), its top has scale
\(\xi_y=L^{j(b)-n}\) and distance at most \(3/2\) from that
endpoint. Apply \eqref{II-8-18} and absorb
\(e^{3\delta_{\mathrm{lc}}/4}\) into \(4\kappa_6\).
This gives \eqref{II-8-70}, including the factor
\(\eta_n/\xi_y=L^{-j(b)}\).
At a node outside the collars, \(w_0=w_1=g=1\), so
\(w'=1\) and \(w=u_{B_1}^{-1}\).
The last estimate of \eqref{II-8-18} gives \eqref{II-8-71}.
\end{proof}

We next verify these hypotheses for the strip.
Use level-\(k\) lengths, \(t=k+1\), and the completed collared
sequence \eqref{II-10-32}. For the analytic strip assertions below
assume the running conditions
\[
 \begin{gathered}
 |b_i|\le B_\sharp\quad(1\le i\le t),\qquad
 \bar B=\frac{LB_\sharp}{L-1},\qquad x_t>\bar B,\\
 y_t>2\mathfrak p_0,\qquad g_t\le\gamma,\qquad
 \bar B\gamma^2\le1,\qquad 2\varepsilon_t\le\bar\varepsilon .
 \end{gathered}
\]
They give \(x_0>\cdots>x_t\),
\(\varepsilon_k<\varepsilon_t\le\sqrt{2L}\varepsilon_k\),
\(\check R_t=R(x_t)\ge y_t^r\), and
\(y_t\ge-2\log\gamma\).
Define
\begin{equation}\label{II-13-5}
 \Omega^{\rm loc}=Z_t\cap\Lambda_k\cap\Omega_0^\diamond,\qquad
 (Y_1,l_1)=(\Lambda_t,t),\qquad
 (Y_2,l_2)=
 \begin{cases}(Z_k,k),&t\ge2,\\
 ((\Omega_0^\diamond)^c,0),&t=1.
 \end{cases}
\end{equation}
The first-step convention is \(\Lambda_0=\mathbb T\).
Form its collars, core, joined sequence and datum by
\eqref{eq:joined-localization-sequence}. For a component
\(\mathfrak o\) of \(\Omega^{\rm loc}\), restrict each joined
region to \(\mathfrak o\), and denote the resulting reduced
determining set and datum by \(\mathbf B^{\mathfrak o},V^{\mathfrak o}\).

\begin{lemma}\label{II-13-11}
The strip joined sequences, on the full localization region and on
each component, are separated and subordinate to the completed
collared sequence. They satisfy the geometric hypotheses of
\cref{II-8-68}. If the completed datum is regular layer by layer
at \((2\varepsilon_i)_{i\le t}\), the joined data are regular
at \(2\varepsilon_t\).
The component minimal orbits exist and their product is the minimal
orbit of the joined datum, provided
\(2\varepsilon_t\le\epsilon_{\mathrm{lc}}\).
The component datum uses only the original rows of levels \(k,t\)
whose endpoints lie in that component.
\end{lemma}

\begin{proof}
Write \(\sigma_i=M_1L^{i-k}\) and
\(\rho_s(x)=\operatorname{dist}_\infty(x,Y_s)\).
The domain and exploration rules give
\[
 \operatorname{dist}_\infty(Y_1,\Omega_t^c)\ge2W_t,\quad
 \operatorname{dist}_\infty(Y_2,\Omega_t)\ge W_t,\quad
 \operatorname{dist}_\infty(Y_1,Y_2)\ge3W_t .
\]
For \(t\ge2\), the second distance is at least \(8W_t\);
for \(t=1\), it is the one-layer enlargement
\(\Omega_0^\diamond=\Omega_1^\sim\).
Here \(\Lambda_t\) is obtained by removing the two coarse
layers about the collar and explored region; the earlier
domain excludes the eight-layer hull of \(Z_k\).
Thus \((\Omega^{\rm loc})^c=Y_1\cup Y_2\), and
\(\Omega^{\rm loc}=Z_t\setminus Z_k\) for \(t\ge2\).
The regions are unions of \(M_1\)-cubes at level \(t\).
Since \(W_t\) is larger than \(4LM_1\), a point in
\(Y_s^{\sim4}\setminus Y_s\) belongs to the old layer \(l_s\).
The displayed separation makes the two wide collars disjoint
and proves \eqref{II-8-53}.

The tower membership condition at a bare site \(x\in\Omega^{\rm loc}\)
is
\[
 x\in\Omega^{[q]}\quad\Longleftrightarrow\quad
 \rho_s(x)>a_q\sigma_q\ (s=1,2),\qquad
 a_0=0,\ a_1=1,\ a_q=2\ (q\ge2).
\]
Indeed the level-\(q\) cube at \(x\) is retained precisely when
it is at distance at least \(a_q\sigma_q\) from the complement.
On a collar the old membership condition is \(q\le l_s\);
in the core the displayed inequalities hold at every old level
containing \(x\). Therefore
\[
 \Omega'_q=\Omega_q^\diamond\cap\Omega^{[q]}.
\]
For \(q\le k\) this is \(\Omega^{[q]}\). At the last level
it is \(\Omega_t\cap\Omega^\natural\) if \(t\ge2\);
if \(t=1\) it is \(\Omega_1\setminus\Lambda_1^{\sim1}\),
with the latter enlargement in level-1 \(M_1\)-cubes.
The tower rule gives
\(\operatorname{dist}_\infty((\Omega'_q)^c,\Omega'_{q+1})
\ge\sigma_{q+1}\).
Since \(M_1\ge R_1\mathfrak M\), these are the required
separation inequalities.

Membership in a reduced determining set is decided by the endpoints.
In the core old and joined memberships agree.
If a joined row is not an old row, one endpoint is in \(C_s\);
the other is within \(L\), hence in \(C_s^{\rm w}\).
Its level is strictly below \(l_s\), for at level \(l_s\)
both endpoints meet exactly the old-row condition.
An endpoint in \(C_s^{\rm w}\setminus C_s\) would belong
to the next joined region, which a joined row excludes.
Thus both endpoints are in \(C_s\).
The lift value is one inside a level-\(l_s\) block, or the
value of an old level-\(l_s\) row whose endpoints remain in
that collar. Common rows have levels \(k,t\), since
\(\Omega^{\rm loc}\subset\Omega_k^\diamond\).
This also proves the stated datum dependence.

For completeness, consider a tested joined plaquette.
If a corner lies in \(C_s\) and its level \(q<l_s\), all
four corners lie in that collar: a corner outside it would
be in \(\Omega'_{q+1}\). Every edge value, including one
obtained by averaging the finer rows, is the corresponding
value of \(Q^{s*}_{l_s\to q}V_{l_s}\).
Its holonomy is one or an old tested holonomy at level \(l_s\),
by the lift calculation in \cref{II-7-40}.
Otherwise \(q=k\) or \(t\), and it is an old tested plaquette.
Its direct edge values are common data. For an edge represented
by finer joined rows, those rows are either all common, or
are the lift in the collar with \(l_s=q\); averaging the lift
then recovers \(V_q\). Thus all four edge values agree with
the old holonomy. Since \(\varepsilon_k\le\varepsilon_t\),
this proves regularity at \(2\varepsilon_t\).

Distinct components of a region of level-\(t\) \(M_1\)-cubes
are separated by at least \(LM_1\). A bond or plaquette of a
level at most \(t\), with all its vertices in the region,
therefore lies in one component. The determining rows,
their averaging blocks, and the plaquette action split
over components. The constraint set and residual gauge group
are products, the action is a sum, and criticality and
regularity are componentwise conditions. The regular-background
theorem applies to every separated component; their product
is the joined minimal orbit.
\end{proof}

Let \(S_t^{\rm pr}=\mathfrak S\cup R\). Every strip bond has
both endpoints in one component \(\mathfrak o\).
Let \(Y_S^{\mathfrak o}\) consist of the level-\(t\) sites
in \(\Omega_t\cap\mathfrak o\) whose blocks contain an endpoint
of such a bond. These sites are nodes of
\(\mathbf B^{\mathfrak o}\). No other site of their blocks is
a node. An endpoint \(x\) of a strip bond either belongs to
one of these blocks, or is itself a level-\(k\) node.
Define its anchor \(y(x)\) to be the block center in the first
case, and \(x\) in the second.

Here are the geometric facts which justify this prescription.
The strip is at distance at least \(W_t\) from \(Y_1,Y_2\).
If a coarse bond \(c\) has an endpoint in the strip, every
bare site \(z\) in \(B^t(c_-)\cup B^t(c_+)\) is within
\(3L/2\) of that endpoint. Therefore it is in the core,
in \(\Omega'_k\), and in no wide collar. A tower endpoint
is within \(3LM_1\) of \(Y_1\) or \(Y_2\).
Since the largest joined scale is \(L\) in the present units,
weighted distance is at least \(L^{-1}\) times sup distance.
Consequently
\begin{equation}\label{II-57-56}
 D(z)\ge M\check R_t-3M_1-2
                \ge\tfrac12M\check R_t
 \qquad(z\in B^t(c_-)\cup B^t(c_+)).
\end{equation}
The inequalities use the running geometric floors
\(M_2=L^2M_1\), \(M\ge L^2M_2\), and \(\check R_t\ge1\).
Every strip bond belongs to one such double block: if it
crosses blocks take their common coarse bond; if it stays
inside one block take a neighboring center within its test
cube. The double block belongs to a single component.

If a strip endpoint \(x\) is in \(\Omega_t\), its center is
in \(\Omega'_t\); a neighboring coarse bond in the core shows
that center is a node. Otherwise choose a nearest neighbor
of \(x\) inside its level-\(t\) block. Both are in
\(\Omega'_k\setminus\Omega'_t\), so their bond is a determining
row and \(x\) is a node. This proves the assertions about
the anchors. In particular \eqref{II-57-56} holds at both
anchors of a strip bond.

For a configuration \(W\) on the double block of a strip bond
\(b\), put
\[
 \operatorname{ax}_W(x)=
 \begin{cases}W(y(x),x),&x\in B(Y_S^{\mathfrak o}),\\
 1,&\text{otherwise},\end{cases}
 \qquad
 \operatorname{dr}_b(W)=
 \operatorname{ax}_W(b_-)W(b)\operatorname{ax}_W(b_+)^{-1}.
\]
Under a gauge \(v\), its transformation law is
\[
 \operatorname{dr}_b(W^v)
   =v(y(b_-))\operatorname{dr}_b(W)v(y(b_+))^{-1}.
\]
The anchors are fixed by residual gauges. Put
\[
 c_2^{\mathrm g}=\min\{a_{\mathsf m},c_2^{\mathrm{SO}},
                           \varpi_{\mathrm{av}}/2\},\qquad
 c_{\mathrm{strip}}=\min\{c_{\mathsf m},c_{\rm plain},
                           r_{\mathrm{av}}/4\},\qquad C_v=4\kappa_6.
\]
These are curvature, perturbation and gauge-factor constants for
the centered iteration in \cref{I-39-36,I-27-78}.
For \(2L^2B_3\varepsilon_t\le c_2^{\mathrm g}\),
curvature averaging on these double-block towers puts every
intermediate configuration in the open mean domain.
Thus the local strip background is
\begin{equation}\label{II-13-27}
 V_k^{\mathrm{bg,loc}\diamond}(b)
   =\operatorname{dr}_b(\mathcal M^kU^{\mathfrak o}),
       \qquad U^{\mathfrak o}\in
               \mathcal C(\mathbf B^{\mathfrak o},V^{\mathfrak o}).
\end{equation}
It is independent of the representative. Equivalently, choose
one for which \(\mathcal M^kU^{\mathfrak o}\) is axial on
\(Y_S^{\mathfrak o}\), and restrict its average to the strip.
The normalization is permitted because those centers are nodes
and their other block sites are not nodes.

The set of data regular layer by layer is open. Indeed the plain
tested holonomies are continuous; on their regularity domain,
the finer double blocks used by the other tested holonomies
lie in the open averaging domain. All tested holonomies are
therefore continuous there, and only finitely many strict
inequalities remain. The response \cref{I-82-15} supplies a
continuous local choice of each component representative.
Successive averaging and the contour means are continuous on
the curvature domain just described. Formula \eqref{II-13-27}
therefore proves stable definedness and continuity of the local
strip background on the regular layer-data domain.

Set
\begin{equation}\label{II-57-6}
 c_L=4+2C_v+4\kappa_6,\qquad C_{\mathrm{xp}}=\pi c_LC_{\mathrm{lc}},
 \qquad
 C_{\mathrm{cu}}=4C_{\mathrm{lc}}
                  \left(1+\frac{1+C_v+2\kappa_6}{L}\right).
\end{equation}
We shall use the following explicit smallness bounds:
\begin{equation}\label{eq:actual-strip-smallness}
 2\varepsilon_t\le\epsilon_{\mathrm{lc}},\qquad
 2L^2B_3\varepsilon_t\le c_2^{\mathrm g},\qquad
 2LC_{\mathrm{lc}}\varepsilon_t\le c_{\mathrm{strip}}/2,\qquad
 2c_LC_{\mathrm{lc}}\varepsilon_t\le1.
\end{equation}
They are uniform upper bounds on \(\varepsilon_t\).

\begin{proposition}\label{II-57-64}
Suppose the completed datum is regular layer by layer at
\((2\varepsilon_i)_{i\le t}\), the running geometry and the
localization floors above hold, and
\eqref{eq:actual-strip-smallness} holds.
At every positive strip bond the intrinsic logarithm
\begin{equation}\label{II-13-30}
 H_b=\frac1i\log\left[
 V_k^{\mathrm{bg}\diamond}(b)
             V_k^{\mathrm{bg,loc}\diamond}(b)^{-1}\right]
\end{equation}
exists and satisfies
\begin{equation}\label{II-57-65}
 |H_b|\le C_{\mathrm{xp}}\varepsilon_t
                  e^{-\delta_{\mathrm{lc}}M\check R_t/4}.
\end{equation}
If the running parameters also satisfy
\begin{equation}\label{eq:actual-strip-cutoff-smallness}
 8\|\vartheta'_{\mathrm{co}}\|_\infty C_{\mathrm{xp}}
       \sqrt{2L}\frac{A_0}{A_1}
          e^{-\delta_{\mathrm{lc}}M\check R_t/4}\le1,
\end{equation}
then \eqref{II-26-84} holds and
\(|H_b|\le\delta_k/24\).
Both backgrounds and \(H_b\) are continuous on this regular
domain, depend only on the datum of the step, and are independent
of the internal fine variables.
\end{proposition}

\begin{proof}
Use a representative \(U\) of the completed orbit whose
level-\(k\) average is axial on \(\Omega_t^{(t)}\).
For the component of \(b\), choose \(U^{\mathfrak o}\)
axial on \(Y_S^{\mathfrak o}\), and complete it by arbitrary
component representatives to an element \(U'\) of the joined
orbit. The regularity and geometry required in \cref{II-8-68}
hold by \cref{II-13-11}, with \(\varepsilon_{\mathrm r}=2\varepsilon_t\).
Put
\[
 s=2C_{\mathrm{lc}}\varepsilon_t
                       e^{-\delta_{\mathrm{lc}}M\check R_t/4}.
\]
On the bare double block of \(b\), \eqref{II-57-56} and
\eqref{II-8-70} give \(|a|\le L^{-k}s\).
At the anchors \eqref{II-8-71} gives \(w'=1\) and
\(|w-1|\le s\).
The joined field has bare curvature at most
\(2B_3\varepsilon_tL^{-2k}\) on that double block.

Apply the centered averaging factorization to
\((U|_\Omega)^w=e^{ia}(U')^{w'}\) on this tower of height \(t\),
through level \(k\). Its curvature parameter is
\(2L^2B_3\varepsilon_t\), and
\(L^t\sup|a|\le Ls<c_{\mathrm{strip}}\), by
\eqref{eq:actual-strip-smallness}.
Writing \(W_0=\mathcal M^k((U')^{w'})\) and
\(W_1=\mathcal M^k((U|_\Omega)^w)\), it gives
\[
 W_1=(e^{ia_k}W_0)^{f_k},\qquad
 \sup|a_k|\le2s,\quad \sup|a_k|\le r_{\mathrm{av}}/2,\quad
 |f_k-1|\le C_vs,\quad \varpi_c(W_0)\le\varpi_{\mathrm{av}}.
\]
These are precisely the intermediate estimates in the proof
of \cref{I-27-78}; the bounds are local to the double-block tower.
For \(k=0\) the factorization is the identity.
The one-step contour estimate gives
\[
 |\operatorname{ax}_{e^{ia_k}W_0}(x)
                  \operatorname{ax}_{W_0}(x)^{-1}-1|
       \le2\kappa_6s .
\]

Write \(y_\pm=y(b_\pm)\). Covariance of the dressed values,
the two axial normalizations, and \(w'(y_\pm)=1\) give
\[
 \begin{split}
 V_k^{\mathrm{bg}\diamond}(b)
 &=w(y_-)^{-1}f_k(y_-)
    \operatorname{ax}_{e^{ia_k}W_0}(b_-)
       (e^{ia_k}W_0)(b)
    \operatorname{ax}_{e^{ia_k}W_0}(b_+)^{-1}
       f_k(y_+)^{-1}w(y_+),\\
 V_k^{\mathrm{bg,loc}\diamond}(b)
 &=\operatorname{ax}_{W_0}(b_-)W_0(b)
                           \operatorname{ax}_{W_0}(b_+)^{-1}.
 \end{split}
\]
All factors are unitary. Comparing their seven factors bounds
the distance of the two backgrounds by
\[
 (2+2C_v+4\kappa_6+2)s=c_Ls\le1.
\]
The relative product is therefore different from \(-1\).
Its intrinsic logarithm has norm at most \(\pi c_Ls/2\),
which proves \eqref{II-57-65}.

The running relation
\(\varepsilon_t\le\sqrt{2L}\varepsilon_k
 =\sqrt{2L}(A_0/A_1)\delta_k\), together with
\eqref{eq:actual-strip-cutoff-smallness}, gives
\eqref{II-26-84}. Since
\(\|\vartheta'_{\mathrm{co}}\|_\infty\ge3\), it also gives
\(|H_b|\le\delta_k/24\).
Continuity follows from the stable backgrounds and the logarithm.
Their defining data use only completed determining rows.
No bond of \(\mathfrak X_0\) belongs to those rows: at level
\(k\) such a row excludes endpoints in \(\Omega_t\).
This proves independence of the internal variables.
\end{proof}

The cutoff-zone tests also compare the strip background with a
test-cube background. The required estimate follows from the
same localization argument:
\begin{equation}\label{II-57-33}
 |V_k^{\mathrm{bg}\diamond}(b)
                    V^{\mathrm{bg}}_{k,\square}(b)^{-1}-1|
 \le C_{\mathrm{cu}}\varepsilon_t
                    e^{-\delta_{\mathrm{lc}}M_2\check R_t/4},
 \quad b_\pm\in\square^{\sim3},\quad
                 \square^{\sim4}\subset\Omega_t .
\end{equation}
To verify its applicability, take
\(\Omega=\square^{\sim4}\), \(Y_1=\Omega^c\), \(l_1=t\)
in \cref{II-8-68}. The whole region lies in the last old
layer, so the joined sequence is exactly its
\(\mathfrak r_2\)-tower. The common rows are its level-\(t\)
rows; all lower rows are tower rows in the two-layer collar.
The joined datum is the induced lift of \(V_t\), and is
regular by \cref{II-7-40}. It thus defines exactly the
background appearing in \eqref{eq:history-PQ-tests}.
Each point of \(\square^{\sim3}\) has weighted distance at least
\(M_2\check R_t-2M_1\ge M_2\check R_t/2\) from the tower rows.
Indeed its sup distance from \(\Omega^c\) is at least
\(LM_2\check R_t\), the collar has width \(2LM_1\),
and the largest weighted scale is \(L\).

On a double block inside \(\square^{\sim3}\), set
\(s=2C_{\mathrm{lc}}\varepsilon_t
 e^{-\delta_{\mathrm{lc}}D/2}\) with this lower bound \(D\).
Now the bare field in \eqref{II-8-69} has size at most
\(L^{-t}s\), so at level \(k\) the centered factorization gives
\(\sup|a_k|\le2L^{-1}s\) and
\(|f_k-1|\le C_vL^{-1}s\).
The two backgrounds are axial on both coarse block centers.
For the gauge
\(g=(w')^{-1}f_k^{-1}w\), their relation is
\[
 g(b_-)V(b)g(b_+)^{-1}=e^{i\beta(b)}V'(b),\qquad
 |\beta|\le2L^{-1}s,\qquad
 |g(c_\pm)-1|\le(1+C_v/L)s .
\]
Their contour means are one. The contour estimate therefore
bounds \(|g(x)-1|\) by
\((1+C_v/L)s+\kappa_6\sup|\beta|\) throughout the double block.
Substitution in the bond relation bounds \(|VV'^{-1}-1|\)
by
\[
 \left(2+\frac{2(1+C_v+2\kappa_6)}L\right)s.
\]
Every bond with endpoints in \(\square^{\sim3}\) lies in such
a double block. This proves \eqref{II-57-33} with
\eqref{II-57-6}.

Suppose, in addition,
\begin{equation}\label{eq:actual-cube-cutoff-smallness}
 8C_{\mathrm{cu}}\sqrt{2L}\frac{A_0}{A_1}
              e^{-\delta_{\mathrm{lc}}M_2\check R_t/4}\le1.
\end{equation}
If a cutoff-zone bond belongs to the test set of a cube \(\square\),
then \(\square^{\sim4}\subset\Omega_t\).
In fact it is also tested by an explored cube
\(\square''\subset\Omega_t^{\sim-1}\); their distance is at
most \(4LM_2\check R_t+1\), whereas this coarse interior is
at distance \(W_t\) from \(\Omega_t^c\).
Adding the diameter of \(\square\) and its four-layer
enlargement still gives a distance less than \(W_t\), by
\(M\ge L^2M_2\).
Thus, for
\(X=V_k^{\mathrm{bg,loc}\diamond}(b)
V^{\mathrm{bg}}_{k,\square}(b)^{-1}\),
\begin{equation}\label{II-57-67}
 |X-1|\le |H_b|+\delta_k/8\le\delta_k/6.
\end{equation}
For every \(W'\in G\) this implies
\begin{equation}\label{II-57-68}
 \mathbf1_{\{|\Psi_{H_b}(W')X-1|<2\delta_k\}}
    =\mathbf1_{\{|W'X-1|<2\delta_k\}}.
\end{equation}
If \(\theta(W')=1\), the two products agree.
Otherwise \(|W'-1|<5\delta_k/3\); the second product is
within \(11\delta_k/6\) of one, and the first is within
\((5/3+1/24+1/6)\delta_k=15\delta_k/8\).
Both indicators are then one.

The running cutoff supplies
\(\check R_t\ge y_t^r\) and \(y_t\ge-2\log\gamma\).
Consequently \eqref{eq:actual-strip-cutoff-smallness} and
\eqref{eq:actual-cube-cutoff-smallness} hold uniformly in the
step if \(\gamma\le e^{-y_{\mathrm{xp}}/2}\), where
\begin{equation}\label{II-57-7}
 \begin{gathered}
 \rho=\max\{1,A_0/A_1\},\\
 \Upsilon_1=8\|\vartheta'_{\mathrm{co}}\|_\infty
                C_{\mathrm{xp}}\sqrt{2L}\rho,\qquad
 \Upsilon_2=8C_{\mathrm{cu}}\sqrt{2L}\rho,\\
 y_{\mathrm{xp}}=
 \max\left\{
 \left(\frac{4\log\Upsilon_1}{\delta_{\mathrm{lc}}M}\right)^{1/r},
 \left(\frac{4\log\Upsilon_2}{\delta_{\mathrm{lc}}M_2}\right)^{1/r}
 \right\}.
 \end{gathered}
\end{equation}
This choice is made after the geometric constants and the ratio
\(A_0/A_1\); it is independent of the history, volume and level.

\subsection{Stored tests and their readings}
\label{sec:stored-history-tests}

We now specify the factors that remain after a step. Products are
indexed by occurrences of tests, so that distinct occurrences are
retained even when they define the same function. Assume throughout
that \(\check R_l\le L\check R_{l+1}\).
At step \(t=l+1\), split each passing \(Q\)-test into its one-bond
factors, and each passing rim test into its one-bond inequalities.
A failed test remains one factor \(1-\chi\). Split the gauge factor
at \((y,x)\) into its indicator and its normalized exponential.

The stored factors are the following. Outside \(\Omega_t\), retain
the density and \(P\)-tests, the gauge factors based at
\(y\notin\Omega_t\), and the \(Q\)-tests. Also retain the density
tests of level \(l\) in \(\Omega_l\setminus\Omega_t\). Inside
\(\Omega_t\), retain the one-bond \(Q\)-tests whose tested bond
does not belong to \(\mathcal N_{\mathrm s}\), all fluctuation
inequalities at external bonds, and the failed rim and cube tests.
The fluctuation inequalities at bonds meeting \(\Lambda_t\)
will instead remain under the integral. Passing density tests
in \(\Omega_t\), gauge factors based there, and one-bond
\(Q\)-tests on \(\mathcal N_{\mathrm s}\) will be treated
separately in the support identity.

For a configuration argument \((j,b)\), its \emph{reading} is
obtained recursively from the free coordinates. A free bond is
read as its coordinate; a tied bond is read by averaging its
previous-level double block. On a zone bond of step \(t\), use
\begin{equation}\label{II-58-37}
 V_l[\mathbf v,\mathfrak q](b)
   =e^{ig_l\mathfrak q_l(b)}V_l^{\mathrm{bg,loc}\diamond}(b).
\end{equation}
Here the local background is evaluated on its determining free
coordinates in the component of
\(\Omega^{\mathrm{loc}\diamond}\) containing \(b\). At a pivot
passenger \(b_0(c)\), use the Borel pivot rule of
\cref{II-6-12}, applied to the readings on \(\mathsf P_c^\circ\)
and to the level-\(t\) coordinate \(V_t(c)\).
An equation can be used only after all its arguments have been
read and only on the domain of its right side.
The rank is zero for free coordinates and one plus the maximum
argument rank for a newly read value.

For a test \(\chi\) of configuration variables, set
\(\chi^{\mathrm{wr}}\) equal to the test of its readings when
they are all defined and equal to zero otherwise. The written
complement is \(1-\chi^{\mathrm{wr}}\), including where a reading
is undefined. A fluctuation test is written by replacing
\(|\mathcal A_l(b)|\) by \(g_l|\mathfrak q_l(b)|\).
The latter rule uses external coordinates only.
Equation \eqref{II-58-37} is a coordinate reading; in the cutoff
zone its value need not be the physical field
\(\Psi_H(e^{ig_l\mathfrak q_l(b)})
 V_l^{\mathrm{bg,loc}\diamond}(b)\).

\begin{proposition}\label{II-62-19}
Every stored factor \(f\) of step \(t=l+1\) is determined by
free configuration coordinates of levels at most \(t\) and
external coordinates \(\mathfrak q_l\), all located in one
component \(D_f\) of \(Z_t\). Its written value is Borel and
bounded. It takes values in \(\{0,1\}\), except for a normalized
gauge exponential, which takes values in
\([0,\mathsf z(x_l,\varepsilon_l)^{-1}]\).

Every configuration argument of \(f\) is of one of four kinds:
a free bond; a tied bond; a zone bond of step \(t\); or a pivot
passenger \(b_0(c)\) of that step with
\[
 \mathsf P_c^\circ
       \subset\Gamma_l^{\mathrm{pr}}\cup\mathcal Z_1 .
\]
Its reading uses only tying equations and zone and pivot
equations of step \(t\). In particular no global-background
reading and no strip equation of another step enters.
For a history of length \(k\), dependence of a stored factor on
a current coordinate \(V_k(b)\) is possible only for \(t=k\),
and then both endpoints of \(b\) lie in \(Z_k\).
\end{proposition}

\begin{proof}
We give the geometric check, since it is what prevents a stored
factor from involving the next integration variables. Use
level-\(l\) lengths and put \(w=LM_2\check R_t\). The parameter
floors imply
\[
 6w+(L-1)<W_t=LM\check R_t.
\]
Every point at distance at most \(W_t\) from \(\Omega_t^c\)
lies in \(Z_t\). A test cube contained in
\(\Omega_t^{\sim-1}\) is at distance at least \(W_t\) from
\(\Omega_t^c\). Consequently a bond whose endpoints are at
distance at most \(\rho\) from \(\Omega_t^c\), with
\(\rho+2w<W_t\), cannot belong to
\(\mathcal N_{\mathrm s}\cup\mathcal N_\ell\).
If it meets \(\Omega_t\) and is nonpivot, it therefore belongs
to the plain zone \(\mathcal Z_1\); if it does not meet
\(\Omega_t\), it is proper.
If it is a pivot \(b_0(c)\), the entire double block lies within
distance \(\rho+L-1\) of \(\Omega_t^c\).
Provided \(\rho+L-1+2w<W_t\), all its nonpivot bonds are
therefore proper or in \(\mathcal Z_1\).
The endpoints of \(c\) are outside \(\Lambda_t\), so \(c\)
is a free bond of level \(t\).

For each stored factor choose a connected set containing its
arguments as follows:
\[
\begin{array}{c|c}
 \text{factor}&\text{set containing its arguments}\\ \hline
 \text{bare density test on }\square&\square^\sim\\
 P\text{- or }Q\text{-test on }\square'&\square'^{\sim4}\\
 \text{gauge factor based at }y&\Delta_{K-t}(y)\\
 \text{carried density test on }\square&\square^{\sim4}\\
 \text{external one-bond inequality at }b&[b_-,b_+]\\
 \text{failed rim test at }c&c^\sim\\
 \text{failed cube test on }\square'&\square'^{\sim3}.
\end{array}
\]
For factors outside \(\Omega_t\) these sets lie within distance
\(4w\) of \(\Omega_t^c\). For a carried density test this uses
\(\check R_l\le L\check R_t\).
Thus the sets lie in \(Z_t\), and the preceding argument with
\(\rho=4w\) classifies all their level-\(l\) arguments.
Level-\(t\) arguments are proper, since their endpoints lie in
\(Z_t\), outside \(\Omega_{t+1}\).

For an inside one-bond \(Q\)-factor, let \(\square'\) be its
test cube. If \(\square'\) were in
\(\Omega_t^{\sim-1}\setminus R\), all its tested bonds would
belong to \(\mathfrak U(R)=\mathcal N_{\mathrm s}\).
Thus \(\square'\subset\mathfrak S\cup R\).
Its four-layer enlargement lies in \(Z_t\), because
\(4w<W_t\) and the exploration leaves a coarse collar around
\(\mathfrak S\cup R'\).
Its sole level-\(l\) argument is either proper or a zone bond;
it is nonpivot. The sets for failed rim and cube tests lie in
\(Z_t\) by the same collar, and all their tested bonds are
external. An external bond has both endpoints in \(Z_t\).
Each set in the table is connected, hence lies in a unique
component \(D_f\).

It remains to follow the equations, including their domains.
For a proper bond, each tying step lowers the level and uses
only its double block. A region made of the relevant coarse
cubes contains these double blocks whenever it contains their
endpoints. Induction on the level therefore keeps every
entering free coordinate in \(D_f\).
A zone equation has rank one: the joined background depends
only on its determining free bonds in one component
\(\mathfrak o\subset\Omega^{\mathrm{loc}\diamond}\).
Both endpoints of the zone bond lie in \(\mathfrak o\), so
\(\mathfrak o\subset D_f\).
For a pivot equation, \(V_t(c)\) has rank zero and its other
arguments are proper or plain-zone bonds. Its rank is at most
\(\max\{l,1\}+1\). The double block and all components
\(\mathfrak o\) entering it lie in the same \(D_f\).
This proves termination, the four alternatives, and the asserted
coordinate dependence.

The zone background has an open domain in its determining
coordinates and is continuous there. To see that eliminating
irrelevant coordinates preserves this assertion, fix their
values: the stable domain is a product with those coordinates,
and its section is open, with the same continuous value.
The pivot rule is Borel by \cref{II-6-12}; tying maps are
Borel. Finite composition therefore gives Borel partial
readings. The original tests are Borel functions of their
arguments: their stable domains are open, their defining
backgrounds are continuous there, and their thresholds are
strict. Zero extension and taking \(1-\chi^{\mathrm{wr}}\)
preserve measurability. The written fluctuation tests are
finite combinations of the inequalities
\(g_l|\mathfrak q_l(b)|<s\).
Their ranges, and the bound on the gauge exponential, follow
from their formulas.
Finally the coordinate dependence just proved has level at most
\(t\). Since \(t\le k\), a current coordinate can enter only
when \(t=k\), and then it is located in \(D_f\subset Z_k\).
\end{proof}

\subsection{Regularity at the layers of a history}
\label{sec:history-layer-regularity}

The global regularity radius does not suffice when the thresholds
increase with the level. We first localize the radius in terms of
the tested holonomies of the datum. Only reduced determining sets
are used in this subsection.

For a bare plaquette or a bare bond at which the current is
defined, let \(\lambda(\mathfrak e)\) be the largest index of a
corner. In bare units define
\[
 r_U(p)=L^{2\lambda(p)}|\partial U(p)-1|,\qquad
 r_U(b)=L^{3\lambda(b)}|J^{(0)}(U)(b)|.
\]
These are the scale-independent local regularity radii. At each
bare site \(x\), let \(r(x)\) be the maximum of the radii of
the elements based there, with empty maximum zero.

\begin{lemma}\label{I-56-25}
There are a power \(\sigma\) of \(L\), a constant
\(K_{\mathrm{loc}}\ge4\), and \(a_{\mathrm c}>0\), depending
only on \(L\), with the following property. Assume the separated
region setting, \(\mathfrak M\ge\mathfrak M_{\mathrm c}\),
\(R_1\mathfrak M\ge L^2\sigma\), and at least
\(6R_1\mathfrak M\) sites in each terminal-lattice direction.
Let \(U\) be critical over a reduced datum \(V\), and regular
at radius \(a_{\mathrm c}\). Put \(R_{\mathrm{loc}}=15L\sigma\)
and
\[
 r_R(x)=\max_{d_{\mathbf B}(x,z)\le R_{\mathrm{loc}}}r(z),
 \qquad
 v_R(x)=\max_{\substack{p'\text{ tested}\\
                    d_{\mathbf B}(x,p')\le R_{\mathrm{loc}}}}
                         |\partial V(p')-1|.
\]
Then
\begin{equation}\label{I-56-26}
 r(x)\le K_{\mathrm{loc}}v_R(x)+\tfrac12r_R(x).
\end{equation}
The constants are uniform in the depth and in the number of
changes of the region sequence.
\end{lemma}

\begin{proof}
We describe the localization of the proof of
\cref{I-49-38}. Its reference configuration must now have a
global bound controlled by \(r_R(x)\), rather than by the
largest radius on the original lattice.

Let \(j\) be the layer index of \(x\). If \(j=0\) and \(x\)
is farther than \(5\sigma\) bare units from \(\Omega_1\),
the plaquettes at \(x\), and those in the current at \(x\),
are tested level-zero plaquettes. Thus
\(r(x)\le4v_R(x)\).
Otherwise choose \(q=j\) when \(j\ge1\) and
\(\operatorname{dist}_\infty(x,\Omega_{j+1})>
5\sigma L^j\), and choose \(q=j+1\) in the remaining case.
Separation implies \(1\le q\le n\). In level-\(q\) units,
the cube of radius \(5\sigma\) about \(x\) lies in
\(\Omega_{q-1}\), misses both its complement and
\(\Omega_{q+1}\), and has weighted distance at most
\(15L\sigma\) from \(x\). Thus only the two layers \(q-1,q\)
occur there.

Choose fixed powers \(\mathfrak M_{\mathrm c},R_{\mathrm c}\)
large enough for the reduced inverse and slice estimates, and
put \(\rho_{\mathrm c}=R_{\mathrm c}\mathfrak M_{\mathrm c}\).
Take \(\sigma\ge L^2\rho_{\mathrm c}\). Let \(Q_0\) be the
cube of side \(\sigma\) containing \(x\), and let \(Q_s\)
be its enlargement by distance \(s\), in level-\(q\) units.
Set
\[
 s_q=\sigma-\rho_{\mathrm c},\qquad
 s_i=s_q+\sum_{i\le h<q}
            (\rho_{\mathrm c}+\mathfrak M_{\mathrm c})L^{h-q},
 \qquad
 \Theta_0=\mathbb T,\quad
 \Theta_i=Q_{s_i}\ (1\le i<q),\quad
 \Theta_q=Q_{s_q}\cap\Omega_q.
\]
The intersections are cube-wise. These are admissible
\(\mathfrak M_{\mathrm c}\)-regions, since their boundary
coordinates differ by multiples of
\(\mathfrak M_{\mathrm c}L^{i-q}\). Moreover
\[
 s_1<\sigma-\tfrac56\rho_{\mathrm c}<\sigma-2,\qquad
 s_i-s_{i+1}
       =(\rho_{\mathrm c}+\mathfrak M_{\mathrm c})L^{i-q}.
\]
Hence this local sequence satisfies the required separation.
Its index is at most the original index on \(Q_\sigma\),
and agrees with it on \(Q_{s_q}\).
The relevant cubes have a single torus chart: the side of the
torus in these units is at least \(2L^3\sigma\).

Write \(e=r_R(x)\), \(v=v_R(x)\). The cases \(e=0\) and
\(v\) above a fixed positive threshold give
\eqref{I-56-26} immediately, after increasing
\(K_{\mathrm{loc}}\). Otherwise the local curvature and current
bounds and the regional gauge estimate provide
\[
 U^{u^{-1}}=e^{i\eta A_0},\qquad
 |A_0|,\ |\nabla A_0|\le C(L)\sigma e
       \quad\text{on }Q_\sigma,\qquad \eta=L^{-q}.
\]
Indeed, on \(Q_{3\sigma}\) the two layer indices give curvature
at most \(L^2e\eta^2\) and current at most \(L^3e\);
the relation between current and the full matrix divergence
adds only its quadratic curvature term. The local gauge
hypothesis is therefore satisfied once
\(C(L)\sigma e\) is sufficiently small.
Extend \(A_0\) by zero outside the bonds inside \(Q_\sigma\),
and denote this extension by \(\bar A_0\).
For the reduced determining set of \((\Theta_i)\),
\[
 \|\bar A_0\|_X\le C(L)\sigma e,\qquad
 |J(e^{i\eta\bar A_0})|_{(-3)}\le C(L)\sigma e.
\]
Inside \(Q_{\sigma-2}\) the second bound is the original
current bound, since the new local index is no larger.
Outside that cube the local index is zero. The four factors
of each plaquette then give
\(\eta^3|J|\le16\eta\sup|\bar A_0|\), which proves the bound
also at the abrupt extension.

The extension is critical for the local reduced constraints.
To check this, a curve preserving them can vary only on bonds
meeting \(\Theta_1\subset Q_{s_1}\). Paste that curve to
\(U^{u^{-1}}\) outside \(Q_\sigma\). Inductively in the
averaging level, a bond whose endpoints have local index at
most its level has unchanged average: at an endpoint of
exactly that index this is a local constraint, and at lower
indices it follows from the previous level on the double
block. Every original reduced row has this property, since
\(\Theta_i\subset\Omega_i\). Thus the pasted curve preserves
the original datum. Every plaquette that varies lies inside
\(Q_\sigma\), and has the same action term in both curves.
The first variation of the extended configuration therefore
vanishes.

Use the coarse comb and the block contour transports, as in the
datum gauge in the proof of \cref{I-49-38}, to prescribe a
gauge \(z\) on the local tops. On the coarse box within
distance \(\sigma/4\) of \(Q_0\), the coarse plaquettes
meeting \(\Omega_q\) are tested. Outside \(\Omega_q\), their
fine plaquettes are tested, and one averaging step bounds
the coarse holonomy by \(2L^2v\).
The comb bound and the fine contour comparison consequently
give
\[
 |W^z(c)-1|\le C(L)(1+|y_c-x|_\infty)v
\]
on the local rows near \(x\). The original small chart bounds
the gauge itself by
\(\sup|z-1|\le C(L)\sigma^2e\).
The prescribed-top slice correction \cref{I-46-44}, at
reference field \(1\), now gives
\[
 E=e^{i\eta A}=(e^{i\eta\bar A_0})^w,\qquad
 R\partial^*A=0,\qquad |A|_{(-1)}\le C(L)\sigma^2e .
\]
Its induced top gauge is \(z\). At an outward row the
additional factor at the outer endpoint is exactly one,
because the fine block contour transports in the prescribed
gauge are one. Thus the reduced nonlinear row
\(B=\mathcal Q(A)\) satisfies
\begin{equation}\label{I-56-35}
 \sup|B|\le C(L)\sigma^2e,\qquad
 |B(c)|\le C(L)(1+|y_c-x|_\infty)v
 \quad\text{if }\operatorname{dist}_\infty(y_c,Q_0)
                                      \le\sigma/4-2.
\end{equation}
There are no inward or edge rows in this reduced case.
The logarithm used here is within its intrinsic branch,
by the bound on \(|A|_{(-1)}\).

For completeness, the slice equations used at this point have
the bounds
\begin{equation}\label{I-55-25}
 \|A\|_X,\ |\partial^*\partial A|_{(-3)},\
       \sup|\omega|\le r_f
       \le C(L)\sigma^2e,
\end{equation}
where \(F'(A)=(Q+\mathcal C'(A))^*\omega\).
Indeed, with \(\alpha=|J(E)|_{(-3)}\),
\(\beta=\sup|B|\), the equations
\[
 F'(A)=\partial^*\partial A+\mathcal V(A),\qquad
 QA=B-\mathcal C(A),\qquad R\partial^*A=0
\]
and the reduced inverse give
\[
 \|A\|_X\le C(L)(\alpha+\beta)
                    +C(L)|A|_{(-1)}\|A\|_X .
\]
Absorb the last term. Applying \(H_0^*\) to the multiplier
equation and using the derivative bounds on \(\mathcal C\)
then gives
\(\sup|\omega|\le C(L)\|A\|_X+
C(L)|A|_{(-1)}\sup|\omega|\), which is absorbed in the same
way. The action remainder controls
\(|\partial^*\partial A|_{(-3)}\).
These operations use the weighted row pairing
\(\sum_c\xi_c^{-1}\langle B(c),\omega(c)\rangle\).
In particular their constants do not introduce a ratio of
different layer scales.
The Schur identities, with trivial reference current, are
\begin{equation}\label{I-55-26}
 \begin{split}
 A&=H_0B+\widetilde G_0
            [\mathcal C'(A)^*\omega-\mathcal V(A)]
                   -H_0\mathcal C(A),\\
 \omega&=(\partial H_0)^*\partial A+
       H_0^*[\mathcal V(A)-\mathcal C'(A)^*\omega].
 \end{split}
\end{equation}
The remainders in these two equations are bounded by
\(C(L)r_f^2\), respectively in \(X\) and the row supremum.
This is also the specialization of
\eqref{eq:reset-multiplier} and \eqref{eq:reset-schur}
to the trivial background.

Apply the decay bounds \eqref{eq:corner-H-local} and
\eqref{eq:corner-H-transposes}. In the first equation of
\eqref{I-55-26}, split the row sum at distance
\(\sigma/4-2\) from \(Q_0\).
For \(|y-x|_\infty\le\sigma/8\), the near part uses
\eqref{I-56-35}; summability of
\((1+d)e^{-\delta d}\) bounds it by
\(C(L)(1+|y-x|_\infty)v\).
The far part is bounded by
\(C(L)e^{-\delta\sigma/16}\sigma^2e\).
Thus
\[
 [\![A]\!]_y\le C(L)(1+|y-x|_\infty)v+
 C(L)e^{-\delta\sigma/16}\sigma^2e+C(L)\sigma^4e^2 .
\]
Use this estimate in the second equation of
\eqref{I-55-26}, splitting now at
\(|y-x|_\infty=\sigma/8\). The same summation gives,
for rows within \(4L+2\) of \(x\),
\[
 |\omega(c)|\le C(L)v+
 C(L)e^{-\delta\sigma/16}\sigma^2e+C(L)\sigma^4e^2 .
\]
The transpose estimate already incorporates the row weights;
no additional power of the local spacing enters this bound.
The plaquette estimate for \(e^{i\eta A}\) and the local
adjoint bound for \(K_A^*\omega=F'(A)\) now control the
curvature and current at \(x\). The original and local
indices agree there, and the neighboring scale conversions
cost at most \(L^2\) and \(L^3\). We obtain
\[
 r(x)\le K^\sharp v+
    \bigl(C_b\sigma^2e^{-\delta\sigma/16}
                   +C_c\sigma^4e\bigr)e,
\]
where all three constants depend only on \(L\).
Choose the power \(\sigma\) first so that the exponential
coefficient is at most \(1/4\), and then choose
\(a_{\mathrm c}\) so that \(C_c\sigma^4a_{\mathrm c}\le1/4\)
and all preceding chart and slice inequalities hold.
Taking \(K_{\mathrm{loc}}\ge\max\{4,K^\sharp\}\) proves
\eqref{I-56-26}.
\end{proof}

\begin{corollary}\label{I-56-72}
Under the hypotheses of \cref{I-56-25}, suppose the tested
holonomies of level \(q\) are smaller than
\(\varepsilon^{(q)}>0\). Set
\[
 \delta^{\mathrm c}_0=\frac{\log2}{30L\sigma},\qquad
 \mathsf c^{\mathrm c}_{\mathrm{loc}}=10K_{\mathrm{loc}},
 \qquad
 \mathsf c^{\mathrm V}_{\mathrm{loc}}
                 =e\,\mathsf c^{\mathrm c}_{\mathrm{loc}}.
\]
At every corner \(x'\) of a regularity element,
\begin{equation}\label{I-56-76}
 r_U(\mathfrak e)<
 \mathsf c^{\mathrm V}_{\mathrm{loc}}
 \sup_{\mathcal L_q\ne\varnothing}
       \varepsilon^{(q)}
       e^{-\delta^{\mathrm c}_0d_{\mathbf B}(x',\mathcal L_q)}.
\end{equation}
\end{corollary}

\begin{proof}
Put
\(P(x)=\max_{p'}e^{-\delta^{\mathrm c}_0d_{\mathbf B}(x,p')}
|\partial V(p')-1|\).
The triangle inequality gives
\(P(z)\le e^{\delta^{\mathrm c}_0d_{\mathbf B}(x,z)}P(x)\).
In particular, within radius \(R_{\mathrm{loc}}\), both
\(v_R(x)\le\sqrt2P(x)\) and \(P(z)\le\sqrt2P(x)\).
For \(\tau>0\), maximize \(r(x)/(P(x)+\tau)\) on the finite
bare lattice and call the maximum \(S\). At a maximizing
site, \eqref{I-56-26} gives
\[
 S\le\sqrt2K_{\mathrm{loc}}+\frac{S}{\sqrt2},
 \qquad S\le2(\sqrt2+1)K_{\mathrm{loc}}\le5K_{\mathrm{loc}}.
\]
Let \(\tau\) decrease to zero. Moving from the base of an
element to any of its corners costs weighted distance at
most two, so
\[
 r_U(\mathfrak e)\le
       10K_{\mathrm{loc}}P(x').
\]
A tested plaquette of level \(q\) has a corner in
\(\mathcal L_q\), and its other corners are within weighted
distance \(2L\) of that layer. Hence
\[
 d_{\mathbf B}(x',\mathcal L_q)
       \le d_{\mathbf B}(x',p')+2L.
\]
Since \(2L\delta^{\mathrm c}_0\le1\), each term of the finite
profile is strictly below
\(e\varepsilon^{(q)}
e^{-\delta^{\mathrm c}_0d_{\mathbf B}(x',\mathcal L_q)}\).
This proves \eqref{I-56-76}. If there is no test, the maximum
argument gives \(r_U=0\), while the right side is positive
at every corner, which belongs to a nonempty layer.
\end{proof}

\subsection{Regularity and support of a step}
\label{sec:history-step-support}

We apply the profile estimate to the completed collared sequence.
Let
\[
 \kappa_{\rm run}=(1+\bar B\gamma^2)^{1/2}\le\sqrt2,\qquad
 C_{\rm reg}=2L^{1/2}\kappa_{\rm run}\,\mathsf c^{\mathrm V}_{\rm loc},
 \qquad
 \bar C_{\rm reg}=2(2L)^{1/2}\mathsf c^{\mathrm V}_{\rm loc}.
\]
Increase the geometric floors so that
\(16\delta^{\mathrm c}_0LM\ge\log L\), and take
\(\varepsilon_{\mathrm v}\le a_{\mathrm c}\).
The running-coupling inequalities used below are
\[
 \varepsilon_q\le\varepsilon_j\quad(q\le j),\qquad
 \varepsilon_q\le L^{(q-j)/2}\kappa_{\rm run}\varepsilon_j\quad(q>j).
\]
They hold on every initial segment of the running sequence.

\begin{proposition}\label{II-10-20}
Suppose the datum of a completed collared sequence of depth \(n\)
has tested holonomies smaller than \(2\varepsilon_q\) at level \(q\).
Under the regular-data hypotheses of \cref{I-58-22}, its regular
critical orbit exists and is unique. If \(x'\) is a corner of a
regularity element and \(x'\in\mathcal L_j\), then
\begin{equation}\label{eq:history-layer-radius}
 r_U(\mathfrak e)<C_{\rm reg}\varepsilon_j
                  \le\bar C_{\rm reg}\varepsilon_j .
\end{equation}
In particular, for \(c\in\Gamma_j^\diamond\), every bare plaquette
in \(B^j(c_-)\cup B^j(c_+)\) satisfies
\[
 |\partial U(p)-1|<
       C_{\rm reg}\varepsilon_j L^{-2(j-1)} .
\]
\end{proposition}

\begin{proof}
Existence and uniqueness follow from \cref{I-58-22}, since
\(2\varepsilon_q\le2\varepsilon_n\le\bar\varepsilon\).
Every path between layers \(j<q\) crosses the disjoint gaps
between \(\Omega_i^c\) and \(\Omega_{i+1}\), \(j<i<q\).
The gap at level \(i\), measured with the weight \(L^{-i}\),
has length at least \(8LM\check R_{i+1}\). Thus
\begin{equation}\label{II-10-12}
 d_{\mathbf B}(x',\mathcal L_q)
       \ge8LM(|q-j|-1)_+ .
\end{equation}
The same statement for \(q<j\) follows by reversing the path.
For \(q>j\) the geometric floor therefore gives
\[
 2\varepsilon_qe^{-\delta^{\mathrm c}_0
                       d_{\mathbf B}(x',\mathcal L_q)}
 \le2L^{1/2}\kappa_{\rm run}\varepsilon_j;
\]
for \(q\le j\) the left side is at most \(2\varepsilon_j\).
Apply \eqref{I-56-76}. The endpoints of a row in
\(\Gamma_j^\diamond\), and their \(j\)-blocks, lie in layers
\(j-1\) or \(j\). This follows from block membership for the
standard partitions and separation of the sequence.
The local regularity index of a plaquette in their union is
at least \(j-1\). The last assertion follows from the
definition of \(r_U\).
\end{proof}

Hereafter the running bounds include
\begin{equation}\label{eq:history-background-smallness}
 \bar C_{\rm reg}\varepsilon_t\le\varepsilon_{\mathrm v},
 \qquad L^2\bar C_{\rm reg}\varepsilon_t\le c^{\mathrm{SO}}_2,
 \qquad 2L^2\bar C_{\rm reg}\varepsilon_t\le\bar\varepsilon .
\end{equation}
We use \(\mathsf{PQ}_t(Y)\) for exact averaging on \(Y\) and
passage of the \(P,Q\) tests whose cubes are in \(Y\).
The fields in this definition need not arise from a bare field.
Put
\begin{equation}\label{II-12-11}
 c_{\mathrm Q}=\min\{(8c_{\rm piv})^{-1},(16L)^{-1}\},
 \quad
 c_{\mathrm s}=2+4(2L)^{1/2}\kappa_{\mathsf c}\bar C_{\rm reg},
 \quad
 c_{\mathrm{fl}}=\pi(1+c_{\mathrm s}A_0/A_1).
\end{equation}
Require \(A_1/A_0\le c_{\mathrm Q}\). The smallness conditions
used in the next proof are
\[
 4B_3L^{-2}\varepsilon_t+2\delta_k\le\omega_{\mathrm p},
 \quad
 12\varepsilon_t+6(L-1)\delta_k\le1/40,\quad
 (1+c_{\mathrm s})\varepsilon_t\le1 .
\]
These hold, for example, under the ceiling
\[
 \varepsilon_t\le
 \min\{\varepsilon_\flat,\bar\varepsilon/2,
       c^{\mathrm{SO}}_2/(2B_3),
       \omega_{\mathrm p}/(2+4B_3L^{-2}),1/520,
       (1+c_{\mathrm s})^{-1}\}.
\]

\begin{proposition}\label{II-12-57}
Let the completed collared datum be regular layer by layer at
\((2\varepsilon_q)_{q\le t}\), and suppose
\((V_k,V_t)\in\mathsf{PQ}_t(\Omega_t^{\sim2})\).
Then the fluctuation field \(\mathcal A_k\) is defined on every
bond incident to \(\Omega_t\), including crossing bonds.
On its nonpivot bonds,
\begin{equation}\label{eq:history-fluctuation-bound}
 |\mathcal A_k(b)|<c_{\mathrm{fl}}\delta_k .
\end{equation}
All fluctuation tests are consequently their stated inequalities
in this field; none vanishes through undefinedness. In addition,
for \(y\in\Omega_t^{(t)}\), \(x\in B(y)\),
\[
 |V_k(y,x)-1|<15L\delta_k<\varepsilon_k .
\]
\end{proposition}

\begin{proof}
Take a test cube \(\square'\subset\Omega_t^\sim\) and write
\(G_1=V^{\mathrm{bg}}_{k,\square'}\).
The neighboring \(P\) tests and \cref{II-7-40} give this
background in axial gauge, with average \(V_t\), curvature
at most \(4B_3L^{-2}\varepsilon_t\) on
\(\square'^{\sim3}\), and at most \(2L^{-2}\varepsilon_t\)
on \(\square'^\sim\). The \(Q\) test says that
\(|V_kG_1^{-1}-1|<2\delta_k\) on the nonpivot bonds incident
to \(\square'^{\sim2}\).
If a pivot bond is in \(\square'^\sim\), its whole double
block is in \(\square'^{\sim2}\). The two configurations
have the same average there. The curvature and perturbation
bounds meet \cref{II-6-12}; its global uniqueness clause,
including the other averaging cases, gives
\[
 |V_k(b)G_1(b)^{-1}-1|<2c_{\rm piv}\delta_k
                       \quad\hbox{at the pivot}.
\]
Telescoping the four factors of a plaquette yields
\begin{equation}\label{eq:history-PQ-curvature}
 |\partial V_k(p)-1|<
        8c_{\rm piv}\delta_k+2L^{-2}\varepsilon_t
                         \quad(p\subset\square'^\sim).
\end{equation}

On one block all bonds are nonpivot. Each of the paths
defining its averaged contour has length at most
\(3(L-1)/2\), so its values for \(V_k,G_1\) differ by
less than \(3(L-1)\delta_k\).
The contour family for \(G_1\) has diameter at most
\(8A_L(2L^{-2}\varepsilon_t)\); that for \(V_k\) has
diameter at most this number plus \(6(L-1)\delta_k\).
Both are at most \(1/40\).
For two such families at pointwise distance \(\zeta\),
the contraction defining the mean gives
\(|Y_*-Y'_*|\le(2/5)|Y_*-Y'_*|+(12/5)\zeta\).
Thus their logarithmic fixed points differ by at most
\(4\zeta\), and their means by at most \(5\zeta\).
Since the mean for \(G_1\) is \(1\), this proves
\(|V_k(y,x)-1|<15(L-1)\delta_k\).

Let \(G_2=p^\diamond_W\), the averaged completed background
on the incident bonds, completed by \(V_k\) outside.
For a coarse bond \(c\) incident to \(\Omega_t\),
\cref{II-10-20} and the curvature averaging estimate give
\(\varpi_c(G_2)\le2C_{\rm reg}\varepsilon_t\).
This uses the whole double block, including its part outside
\(\Omega_t\): there \(G_2=\mathcal M^kU\) by the completed
datum. Choose \(\square'\subset\Omega_t\) containing one
endpoint of \(c\); its double block lies in \(\square'^\sim\).
The configurations \(G_1,G_2\) have the same coarse average.
Their contour variables agree with \(1\) on the blocks
inside \(\Omega_t\); on an outside block the contour of
\(G_2\) is that of \(V_k\), bounded above.
The double-block comparison \cref{I-25-42} therefore gives
\[
 |G_2(b)G_1(b)^{-1}-1|
 \le4\kappa_{\mathsf c}C_{\rm reg}\varepsilon_t
                        +30(L-1)\delta_k
 \le c_{\mathrm s}\varepsilon_k .
\]
Every incident fine bond belongs to one of these double
blocks. On a nonpivot bond the triangle inequality gives
\(|V_kG_2^{-1}-1|<2\delta_k+c_{\mathrm s}\varepsilon_k\);
at a pivot replace \(2\delta_k\) by \(2c_{\rm piv}\delta_k\).
Both bounds are smaller than \(1\), by the displayed
smallness conditions. Hence the principal logarithm is
defined, and on a nonpivot bond
\[
 |\mathcal A_k(b)|\le\frac\pi2|V_k(b)G_2(b)^{-1}-1|
 <\frac\pi2(2+c_{\mathrm s}A_0/A_1)\delta_k
 \le c_{\mathrm{fl}}\delta_k .
\]
The tests use precisely these nonpivot values. This proves
their claimed interpretation and the proposition.
\end{proof}

We also need the converse regularity test for the stored coordinates.
At a tied point, the holonomy of a tested plaquette of level \(q\)
is the ordinary plaquette holonomy of the level-\(q\) field.
For a first-kind bond this is the datum definition; for a
second-kind bond it follows from the tie and the fact that
its double block lies in the preceding determining set.
Choose a density cube containing the specified corner of the
tested plaquette. The cube lies in the level-\(q\) small-field
region, or is a positive density factor stored when it left
that region. Its four-layer body is exactly one of the bodies
in \cref{II-62-19}.
Thus its written reading gives the same tested holonomy.
If \(q\ge1\), passage of this density test bounds the curvature
of its defining bare background by
\(\varepsilon_qL^{-2q}\) on the one-layer enlargement.
The \(q\)-blocks of the four corners lie there, and the
curvature averaging estimate bounds the resulting holonomy
by \(2\varepsilon_q\). At \(q=0\) the test itself gives
\(\varepsilon_0\). We have proved the implication
\begin{equation}\label{II-67-30}
 \chi_t(\Omega_t)=1,\quad \operatorname{Stay}_{h'}\ne0,
 \quad\hbox{ties satisfied}
 \quad\Longrightarrow\quad
 \mathbf V^\diamond\hbox{ is regular at }
       (\varepsilon_0,2\varepsilon_1,\ldots,2\varepsilon_t).
\end{equation}
Before the step, the same proof uses its positive plaquette
factors together with \(\operatorname{Stay}_h\ne0\).
For a bare field on the selected history, all ties and all these
positive factors hold. This proves the regularity assertion
needed to apply \cref{II-12-57} on its support.

\subsection{Removal of the redundant density tests}
\label{sec:redundant-history-density}

The density tests inside the next small-field region are
stronger than what one gets merely by comparing individual bonds.
The proof compares the two defining backgrounds on a third,
smaller tower. Their discrepancy on its lower rows is far
from the plaquette being tested.

Retain the response constants \(B_5,\delta_{\rm lc},\varrho\)
from \eqref{II-8-18}, and put
\[
 \Theta=100B_5S(\delta_{\rm lc}/2)C_{\rm AC}B_3,\qquad
 D=M_2\check R_k-2M_1,\qquad
 \lambda=\sqrt2L^{-3/2}<0.0302 .
\]
Besides the preceding hypotheses require
\begin{equation}\label{eq:density-redundancy-smallness}
 \begin{gathered}
 \frac{A_1}{A_0}\le
 \min\{c_{\mathrm Q},(36c_{\rm piv})^{-1},
             (100\pi c_{\rm piv}B_5S(\delta_{\rm lc}))^{-1}\},
 \qquad
 e^{-\delta_{\rm lc}D/2}\le\Theta^{-1},\\
 C_{\mathfrak P}\mathfrak M B_3\varepsilon_t\le a_{\mathrm R},
 \quad
 4B_3\varepsilon_t\le\omega_\sharp,\qquad
 2(1+C_{\rm AC}B_3)\varepsilon_t\le\varrho .
 \end{gathered}
\end{equation}
These bounds are uniform in the step: \(D\ge(L^2-2)M_1y_k^r\),
so the exponential condition follows by decreasing the fixed
ceiling on the coupling. We also use
\(\check R_k\le L\check R_t\).

\begin{proposition}\label{II-59-41}
If \(Y\) is a union of coarse cubes and
\((V_k,V_t)\in\mathsf{PQ}_t(Y^\sim)\), then every density
cube \(\square\subset Y\) passes its level-\(k\) density test.
For \(k\ge1\) its defining curvature satisfies
\[
 \mathsf m_{k,\square}(V_k)
                     <\tfrac1{10}\varepsilon_k\eta_k^2 .
\]
For \(k=0\) the corresponding bare curvature is
smaller than \(\varepsilon_0/3\).
\end{proposition}

\begin{proof}
Let \(\square'\) be the test cube containing \(\square\).
The ratio of enlargement scales gives
\(\square^{\sim a}\subset\square'^{\sim a}\) for \(a\le4\)
and \(\square'^{\sim4}\subset Y^\sim\).
By \eqref{eq:history-PQ-curvature},
\[
 |\partial V_k-1|<
 \left(8c_{\rm piv}\frac{A_1}{A_0}+2\lambda\right)
                  \varepsilon_k<\varepsilon_k/3
\]
on \(\square'^{\sim4}\).
This proves the assertion for \(k=0\). For \(k\ge1\),
\cref{II-7-40} gives the density background \(U^{\rm a}\)
on the tower of \(\square^{\sim4}\), with regularity radius
\(B_3\varepsilon_k\) and stable dependence on \(V_k\).
The passed \(P\) test gives the other background \(U^{\rm b}\)
on the tower of \(\square'^{\sim4}\), of depth \(t\), with
curvature smaller than
\(\varepsilon_t\eta_t^2\le\lambda\varepsilon_k\eta_k^2\)
on \(\square^\sim\).

Restrict both backgrounds to \(A=\square^{\sim2}\) and its
tower of depth \(k\). This tower is subordinate to both
original towers; \(\square^\sim\) lies in its top member.
By \cref{II-8-36} restriction preserves criticality and
regularity. Apply the axial alignment used in the proof of
\cref{II-8-68}. Its gauge is \(1\) at level \(k\).
Every original node in \(A\) is of level at least \(k\);
extending the gauge by \(1\) therefore preserves the original
constraints. The aligned data \(V^{\rm a},V^{\rm b}\) can be
compared with the datum \(V^{\rm m}\) induced by \(V_k|_A\).
At the top, \(V^{\rm a}=V^{\rm m}\) and
\(|V^{\rm b}(V^{\rm m})^{-1}-1|\le2c_{\rm piv}\delta_k\).
On the lower rows, summing the curvature bounds of the
successive averages gives the additional discrepancies
\[
 \frac{2C_{\rm AC}B_3\varepsilon_k}{L^2-1},
 \qquad
 \frac{4C_{\rm AC}B_3L^{-2}\varepsilon_t}{L^2-1},
\]
respectively. The middle datum is regular at \(\varepsilon_k\)
by \cref{II-7-40} and the curvature bound for \(V_k\).

Write \(V^\natural=e^{iB^\natural}V^{\rm m}\).
The logarithm bound \(|\log X|\le(\pi/2)|X-1|\) and
\eqref{eq:density-redundancy-smallness} put both \(B^\natural\)
inside the response ball \(\varrho\).
The lower rows have weighted distance at least \(D\) from
the tops above \(\square^\sim\). Exponential summation gives
\[
 m_y(B)\le b_{\rm n}S(\delta_{\rm lc})
       +b_{\rm f}e^{-\delta_{\rm lc}D/2}S(\delta_{\rm lc}/2),
\]
where
\[
 (b_{\rm n},b_{\rm f})=
 \begin{cases}
 (0,\pi C_{\rm AC}B_3\varepsilon_k/(L^2-1)),&\natural={\rm a},\\
 (\pi c_{\rm piv}\delta_k,
  2\pi C_{\rm AC}B_3L^{-2}\varepsilon_t/(L^2-1)),
                                      &\natural={\rm b}.
 \end{cases}
\]
Consequently \(B_5m_y(B^{\rm a})\le\varepsilon_k/100\)
and \(B_5m_y(B^{\rm b})\le\varepsilon_k/50\).

Let \(U^\circ\) represent the middle orbit.
Equations \eqref{II-8-18} and \eqref{II-8-20} bound both
the response and its first differences by \(\vartheta\)
over the plaquette, where \(\vartheta\) is the appropriate
one of these two bounds. The response represents the
perturbed orbit up to gauge. Expanding the four factors of
its plaquette, as in \cref{I-10-120}, gives
\begin{equation}\label{II-59-39}
 |\tau_\natural-\tau_{\rm m}|
 \le\eta_k^2(2\vartheta+11\vartheta^2)
                           +6\eta_k\vartheta\tau_{\rm m},
\end{equation}
where \(\tau_\natural=|\partial U^\natural(p)-1|\) and
\(\tau_{\rm m}=|\partial U^\circ(p)-1|\).
Indeed \(|D_{U^\circ}\mathcal H|\le2\vartheta\);
the quadratic remainder is at most \(11\eta_k^2\vartheta^2\),
and the curvature commutator contributes the last term.
The expansion is applicable since \(\eta_k\le1/13\),
\(\vartheta\le1/50\), and
\(\tau_{\rm m}<\eta_k^2\).
Since \(\tau_{\rm b}<0.0302\varepsilon_k\eta_k^2\),
\eqref{II-59-39} first gives
\[
 \tau_{\rm m}\le
 \frac{0.0302+0.0444}{1-6/650}\varepsilon_k\eta_k^2
 \le0.0753\varepsilon_k\eta_k^2.
\]
Applying it to \({\rm a}\) then gives
\[
 \tau_{\rm a}\le
 ((1+6/1300)0.0753+0.0211)\varepsilon_k\eta_k^2
 <0.1\varepsilon_k\eta_k^2 .
\]
There are finitely many tested plaquettes. Taking their
maximum proves the bound and, by stable definedness, passage
of the density test.
\end{proof}

\subsection{The support identity in the new coordinates}
\label{sec:history-support-identity}

Let \(\operatorname{Stay}_h\) be the product, with multiplicities,
of the written factors stored at all preceding steps. Multiply it
by the indicators of the external tiers and call the result
\(\operatorname{Stay}^{\rm T}_h\). The tiers at step \(t\) are
\(\mathfrak b_k/(4\Theta_1)\) on collar bonds,
\(\theta_c\mathfrak b_k\) on rim bonds, and
\(C_{\rm xt}\mathfrak b_k\) elsewhere.
Choose
\[
 C_{\rm xt}\ge c_{\rm fl}^{\sharp}
   :=c_{\rm fl}+c_{\rm H}A_0/A_1,\qquad
 c_{\rm H}=\pi\kappa_{\mathsf c}(2L)^{1/2}
                         \max\{2\bar C_{\rm reg},4B_3\}.
\]
In particular \(C_{\rm xt}\ge c_{\rm fl}+1/24\).
Assume all preceding smallness conditions, the actual shift
bound of \cref{II-57-64}, and
\begin{equation}\label{eq:history-support-chart-ceilings}
 2\bar C_{\rm reg}\varepsilon_t<\varpi_\chi,\qquad
 2\varepsilon_k<\omega_{\mathrm p},\qquad
 C_\chi(C_{\rm xt}+1)\delta_k\le r_\chi,\qquad
 C_{\rm xt}\delta_k<\pi .
\end{equation}
Their constants are fixed before the coupling ceiling.

The set of regular bases \(d=(W,\phi^+)\) is the open
regular-data condition on the completed collared datum,
intersected with the Borel set of tied untouched coordinates.
For such a base let \(V_0[d]\) be the averaged critical
background on the incident bonds, in axial gauge inside
\(\Omega_t\), and use the untouched field on the exterior.
Its dependence on the base is Borel, indeed continuous on
each regular-data domain. The double-block curvature estimate
in the proof of \cref{II-12-57} and
\eqref{eq:history-support-chart-ceilings} make it a reference
configuration for \cref{II-26-53}.
Its average on \(\Gamma\) is \(W\).
Set \(r=c_{\rm fl}\mathfrak b_k\).
The regional chart at \(g_kB\), \(B=(B_{\rm I},a)\),
is defined on the fixed product \(|B(b)|<r\); it is a
Borel injection with positive Jacobian
\eqref{II-26-54}. Its inverse reads the nonpivot logarithms
and the average. Its Haar substitution formula is
\eqref{II-26-55}.

Let \(f_{h'}\) be the factor selecting the outcome of the
step, including its gauge weight, and put
\[
 T=f_{h'}\operatorname{Stay}^{\rm T}_h\chi_k(\Omega_k),
 \qquad
 S=\operatorname{Stay}^{\rm T}_{h'}\chi_t(\Omega_t);
\]
at the first step use \(T=f_{h'}\).
All products are evaluated as zero when an earlier zero test
makes a subsequent partial factor unnecessary.
Let \(R\) be the product of the integration balls prescribed
by the fluctuation split, and let \(a=\check\Psi_d(a')\).

\begin{proposition}\label{II-169-33}
The actual coordinates and tests satisfy
\eqref{eq:history-coordinate-step},
\eqref{eq:history-chart-commutation}, and the chart and strip
substitution hypotheses \eqref{II-166-48} and
\eqref{II-166-50}. Every selected bare field with nonzero new
weight has a chart point. Outside the null set of untied
untouched coordinates,
\begin{equation}\label{II-169-34}
 T(\mathsf{pt}(\Theta(p,B_{\rm I})))
 =\mathbf1_R(B_{\rm I})\,S(\imath p)\,
 z(x_k,\varepsilon_k)^{-(L^3-1)\#\Omega_t^{(t)}}
                    e^{-x_kG_d(B)} .
\end{equation}
Moreover, \(S(\imath p)>0\) implies that the base is regular,
that \(a'\) is in the domain of the strip substitution,
and that \(|\check\Psi_d(a')(b)|<r\) at every external bond.
Thus the nonzero support in the integration variables is
the product required in \eqref{II-166-17}.
\end{proposition}

\begin{proof}
The forward coordinate map retains \(\phi^+\), takes the
average of \(V_k\) on the new free rows, and records the
new strip variable. This is a Borel partial map. Recursion
of the recorded strip variables gives equality of its domain
with the recursive coordinate domain and
\(\xi_{h'}=\Sigma\xi_h\). The product of the tests and gauge
factors of the selected outcome gives
\(\mathfrak w_{h'}=\mathfrak w_h(f_{h'}\circ\xi_h)\);
the dependencies of these factors are precisely the retained
proper values and the current level-\(k\) field. The chart
fixes its average, and the inverse strip substitution records
\(a'\). Hence \(\Sigma\mathsf{pt}\Theta=\imath\).
The product-measure assertion \cref{II-58-74}, followed by
the regional Haar formula and the fixed-shift Lebesgue formula,
proves the two substitution identities, with their positive
Borel Jacobians. On a selected bare field the regularity
argument following \cref{II-12-57} applies; that proposition
puts its nonpivot logarithms in the chart ball. Global pivot
uniqueness then identifies the entire field with the chart
point, without any initial bound on its pivot entries.

We prove the test identity carefully, since written readings
on a transition zone need not equal physical field values.
At a chart point all proper readings equal the physical
values. A zone reading is
\(e^{ig_ka'(b)}V^{\rm bg,loc}_k(b)\), whereas the physical
value is
\(\Psi_b^\diamond(e^{ig_ka'(b)})V^{\rm bg,loc}_k(b)\).
These agree on \(\mathcal Z_1\).
Every pivot reading used by a stored factor has all of its
nonpivot arguments in proper bonds or \(\mathcal Z_1\),
by \cref{II-62-19}. Once the physical double-block curvature
is below \(\omega_{\mathrm p}\), \cref{II-6-12} identifies
this reading with the physical pivot.
The only remaining reading in \(\mathcal Z_\theta\) occurs
in a one-bond \(Q\) test. Equality of that test at the two
readings is exactly \eqref{II-57-68}. All stored fluctuation
factors agree by the norm-test invariance of the strip map.
Thus every stored factor without a pivot argument agrees
with its written version before any extra assumption on
the pivot curvature.

Suppose the new interior density tests and these stored
factors are nonzero and \(B_{\rm I}\in R\).
The internal and external fluctuation factors give the full
selected fluctuation split. Its small bonds satisfy
\(|g_kB(b)|<\delta_k\).
For such a bond in \(\mathcal N_{\mathrm s}\), every
intersecting \(Q\)-test cube has its four-layer enlargement
inside \(\Omega_t\): the bond meets a two-layer enlargement
of a cube in \(\Omega_t^{\sim-1}\), and the coarse collar
is wider than the sum of these test layers.
The test-cube comparison proved with \cref{II-57-64}
gives
\(|V^{\rm bg}_k(b)(V^{\rm bg}_{k,\square'}(b))^{-1}-1|
 \le\delta_k\).
Since \(|e^{ig_kB(b)}-1|<\delta_k\), this \(Q\) test passes.
All other \(Q\) tests in \(\Omega_t^{\sim2}\) are stored
factors without pivot arguments. The \(P\) tests there are
either the new density tests or stored factors.
Together with the exact chart average, this proves
\((V_k,V_t)\in\mathsf{PQ}_t(\Omega_t^{\sim2})\).

Apply \cref{II-12-57}. It gives \(|B|<c_{\rm fl}\mathfrak b_k\)
even if the chart was initially constructed on the larger
ball allowed by \eqref{eq:history-support-chart-ceilings}.
Apply \cref{II-59-41} with \(Y=\Omega_t^\sim\).
Every density cube meeting a double block over \(\Omega_t\)
lies in this enlargement. Passage of its density test,
followed by curvature averaging, gives
\(\varpi_c(V_k)<2\varepsilon_k<\omega_{\mathrm p}\).
This identifies all pivot readings and proves equality of
all the stored factors.
It also removes the interior density tests; the contour
bound in \cref{II-12-57} removes the interior gauge indicators.
The small-bond argument above removes the remaining interior
\(Q\) tests.

The new external tiers hold at these points. Collar and
rim bonds are outside the zones, so the substitution fixes
them and the selected fluctuation factors impose their tiers.
Elsewhere, \(|g_ka'|\le|g_ka|+\delta_k/24\), by the
fixed-shift estimate and \cref{II-57-64}.
Thus the ordinary tier follows from
\(C_{\rm xt}\ge c_{\rm fl}+1/24\).
Earlier stored factors have exactly the same values in
\(h,h'\): their determining components lie in \(Z_s\subset Z_k\)
and all their coordinates are untouched, by \cref{II-62-19}.
The same is true of the earlier tier indicators.

Partition the factors of the old test times the step into
five disjoint lists: the newly stored factors; the internal
fluctuation indicators; the redundant interior density,
gauge and \(Q\) indicators; the interior gauge exponentials;
and the new density tests.
This partition keeps occurrences, including carried density
tests, rather than identifying equal functions.
The second product is \(\mathbf1_R\), the fifth is
\(\chi_t(\Omega_t)\), and the fourth is
\[
 \prod_{y\in\Omega_t^{(t)}}\prod_{x\in B(y)\setminus\{y\}}
 z(x_k,\varepsilon_k)^{-1}
 e^{-x_k(1-\frac12\operatorname{Re}\operatorname{tr}V_k(y,x))}
 =z(x_k,\varepsilon_k)^{-(L^3-1)\#\Omega_t^{(t)}}e^{-x_kG_d(B)}.
\]
If the second or fifth product vanishes, both sides of
\eqref{II-169-34} vanish. If a newly stored factor without a
pivot argument vanishes, it vanishes in both readings, so
both sides again vanish. In the remaining case the preceding
argument identifies every stored factor, makes the third
product and the new tier indicator equal to one, and proves
\eqref{II-169-34}. This argument does not presume that the
pivot readings are defined in the cases where a zero factor
has already settled the equality.

Finally suppose \(S(\imath p)>0\). Equation \eqref{II-67-30}
gives a regular base. Its tiers and the chart ceiling imply
\[
 |g_ka'|+|H|
 <(C_{\rm xt}+1/24)\delta_k
 <r_\chi/C_\chi\le1<\pi.
\]
The group-angle bound therefore keeps the logarithm in
\eqref{II-26-68} away from its excluded point. The strip
map is defined and gives
\(|g_ka|<(C_{\rm xt}+1/24)\delta_k\).
Set \(B_{\rm I}=0\) and use the larger chart permitted by
\eqref{eq:history-support-chart-ceilings}. Its nonpivot stored
factors agree with the written nonzero factors.
The preceding argument now improves \(|a|\) to
\(c_{\rm fl}\mathfrak b_k=r\).
The complement of the tied set has
\(\mathscr M^+\)-measure zero by completion.
Its product with Haar and Lebesgue coordinates is null,
by their \(\sigma\)-finiteness. This proves the asserted
support statement and its precise exceptional set.
\end{proof}

For a selected bare field the actual strip backgrounds exist
at every step by \cref{II-10-20,II-57-64}. Their logarithms
and the strip inverses are defined by the same bounds.
Induction therefore puts every such field in
\(\operatorname{dom}\xi_h\). This also gives the initial
support and coordinate-definedness assertion required to
iterate the measure identity; empty integration or strip
sets are one-point products and do not change the argument.


\subsection{Local terms for the strip substitution}
\label{sec:strip-transfer}

The strip change of variables produces small boundary terms even
when applied to a quadratic form which is not small. We retain
the decay of the shift before summing its finite differences.

Use the strip zones $\mathcal Z_\theta,\mathcal Z_1$,
$\mathcal Z=\mathcal Z_\theta\cup\mathcal Z_1$, and the component
$\mathfrak o(b)$ containing a zone bond. Write $\mathsf P(b)$
for the one or two adjacent coarse cubes containing its reading
blocks, $\mathsf P_{\mathfrak o}^{\rm loc}$ for the component's
localization cubes, and $\mathsf P_{\mathfrak o}^{\partial}$
for those containing an endpoint of a tower row. Set
\[
 \mathsf{Pr}(b)=\{S\in\mathcal D_t:
  \mathsf P(b)\subset S\subset\mathsf P_{\mathfrak o(b)}^{\rm loc},
  \ S\cap\mathsf P_{\mathfrak o(b)}^\partial\ne\emptyset\}.
\]
For a domain $X$ let $\mathsf{Pr}[X]$ consist of the pairs
$(b,S)$ with $b\in\mathcal Z^+\langle X\rangle$ and
$S\in\mathsf{Pr}(b)$. Here the brackets select bonds with an
endpoint in $X$. For a subset $\mathscr S$ put
$X(\mathscr S)=X\cup\bigcup_{(b,S)\in\mathscr S}S$.

The geometry of the strip gives
\begin{equation}\label{eq:strip-transfer-geometry}
 \begin{gathered}
 d_t(S)+1\ge\check R_t\quad(S\in\mathsf{Pr}(b)),\\
 d_t(X(\mathscr S))\le d_t(X)+
                  \sum_{(b,S)\in\mathscr S}(d_t(S)+4),\\
 \#X(\mathscr S)\le\#X\prod_{(b,S)\in\mathscr S}(1+\#S).
 \end{gathered}
\end{equation}
For the first inequality, a zone endpoint lies farther than
$W_t=LM\check R_t$ from any endpoint of a tower row.
For rows at $\Lambda_t$, their endpoints lie within $3LM_1$
of $\Lambda_t$. The zone excludes the two-layer enlargements
of all passing test cubes in the coarse one-layer neighborhood
of $\Lambda_t$, so its endpoints have distance greater than
$W_t+2LM_2\check R_t$ from $\Lambda_t$.
Their distance from a row endpoint is therefore greater than
$W_t$, since $M_2=L^2M_1$. For rows at $Z_k$, use
$\operatorname{dist}_\infty(\Omega_t,Z_k)\ge8W_t$ and the
row collar width $3M_1$ to obtain the lower bound
$8W_t-1-3M_1>W_t$.
A domain meeting the reading blocks and a boundary cube thus
has diameter greater than $W_t$, whereas its diameter is at
most $LM(d_t(S)+2)$. The sizes are integers, giving the first
inequality. For the others, each $S$ meets a cube of $X$.
Join their edge trees in those cubes, using at most four
extra edges for each join. The counting bound follows from
$\#(X\cup\bigcup S)\le\#X+\sum\#S$ and $\#X\ge1$.
The same argument at a fixed bond gives
\[
 d_t\bigl(\textstyle\bigcup\mathscr S\bigr)+1
       \le\sum_{S\in\mathscr S}(d_t(S)+4)-3.
\]
All these unions are wall connected. A nonempty union carrying
a zone bond contains a cube of $Z_t$ and one of $\Omega_t$.

Suppose the localized shift is given on restriction-stable
classes of pairs by $\mathbf H(b;S;p)$, vanishes outside
$\mathsf{Pr}(b)$, has local holomorphic extensions, is gauge
covariant and conjugation real, and satisfies
\begin{equation}\label{II-73-9}
 |\mathbf H(b;S;p)|\le\bar h_S
       :=a_H e^{-\mu_H(d_t(S)+1)},\qquad a_H>0.
\end{equation}
One family is used on every subregion, by restriction.
Let $c_e>0$, $N_e\ge1$ be an entropy pair:
$\sum_{S\ni\Delta}e^{-c_ed_t(S)}\le N_e$.
For $N_b\ge1$, at most $N_b$ positive zone bonds have an endpoint in any
coarse cube. We may take $c_e=c_{\rm E}$, $N_e=C_{\rm E}$
and $N_b=6(LM)^3$.

Fix $\widehat\kappa>0$ and $0<\rho_{\rm bd}\le1/8$ with
\begin{equation}\label{II-73-14}
 \mu_H\ge8(\widehat\kappa+1)+2c_e,\qquad
 128N_bN_e^2a_He^{-\check R_t}\le\rho_{\rm bd}.
\end{equation}
Put
\[
 R_S=1+\frac{\rho_{\rm bd}e^{-c_ed_t(S)}}{2N_e\bar h_S},
 \qquad r_{\rm bd}=2\rho_{\rm bd}/g_k.
\]
A local function $F_X(\mathsf a,w)$ depends on the external
field only on $\mathfrak X_{\rm E}\langle X\rangle$; $w$
denotes its other arguments. Its domain has the same locality.
Assume $\mathsf a$ is real on $\mathcal Z_\theta$, the strip
chart is defined at zero shift, and
$g_k|\mathsf a(b)|\le1/8$ on $\mathcal Z\langle X\rangle$.
Assume also that the domain contains the complex ball
of radius $r_{\rm bd}$ about $\mathsf a$ in the directions
of these zone bonds, with $F_X$ holomorphic and bounded there
by its domain norm $\|F_X\|$.

For parameters $|\lambda(b,S)|<R_S$, define
\[
 H_\lambda(b)=\sum_S\lambda(b,S)\mathbf H(b;S;p),\qquad
 f_X(\lambda)=F_X(\check\Psi(\mathsf a,H_\lambda),w).
\]
Parameters and shifts vanish off the finite admitted index set.
Define the transfer term by
\begin{equation}\label{II-73-23}
 F_X^{\rm tr}(X_1)=
 \sum_{\substack{\emptyset\ne\mathscr S\subset\mathsf{Pr}[X]\\
                        X(\mathscr S)=X_1}}
       f_X[\mathscr S\mid\emptyset],
\end{equation}
where the bracket is the finite difference at zero.
When evaluated over $X_1$, only $S\subset X_1$ occur,
and the pair is restricted to those supports.

\begin{proposition}\label{II-73-32}
Under these hypotheses the terms in \eqref{II-73-23} are
defined. If $0\le\kappa_1\le\widehat\kappa+1$ and
$\kappa'\ge\kappa_1+1$, then
\begin{equation}\label{II-73-35}
 |F_X^{\rm tr}(X_1)|\le
 2e^{-\check R_t}\|F_X\|e^{\kappa'd_t(X)}
                                      e^{-\kappa_1d_t(X_1)}.
\end{equation}
If $\|F_X\|\le b e^{-\kappa_{\rm in}d_t(X)}$ for a family
with the same room, $0\le\kappa_{\rm out}\le\widehat\kappa$
and $\kappa_{\rm in}\ge\kappa_{\rm out}+1+c_e$, then
\begin{equation}\label{II-73-37}
 \left|\sum_{X\subset X_1}F_X^{\rm tr}(X_1)\right|
 \le4N_e e^{-\check R_t}b e^{-\kappa_{\rm out}d_t(X_1)}.
\end{equation}
With $H_{\rm tot}(b)=\sum_S\mathbf H(b;S;p)$, the whole sum is
\begin{equation}\label{II-73-39}
 \sum_X F_X(\check\Psi(\mathsf a,H_{\rm tot}),w)
 =\sum_XF_X(\mathsf a,w)+
               \sum_{X_1}\sum_{X\subset X_1}F_X^{\rm tr}(X_1),
\end{equation}
provided the room hypotheses hold for every $X$ and
$g_k|\mathsf a|\le1/8$ on all zone bonds.
Every nonzero transfer is supported on a domain reaching
a component boundary and containing a cube of $Z_t$
and one of $\Omega_t$.
\end{proposition}

\begin{proof}
The entropy bound and \eqref{eq:strip-transfer-geometry} give
\[
 \sum_{S\in\mathsf{Pr}(b)}\bar h_S
 \le N_ea_He^{-\check R_t}\le\rho_{\rm bd}/128,\qquad
 \sum_{S\in\mathsf{Pr}(b)}R_S\bar h_S
                  \le65\rho_{\rm bd}/128<\rho_{\rm bd}.
\]
Thus $|H_\lambda|<\rho_{\rm bd}$.
The complex chart map
$\mathsf{cm}(A,H)=i^{-1}\log(e^{iA}e^{-iH})$ is holomorphic
for $|A|,|H|\le1/8$, and
\[
 |\mathsf{cm}(A,H)-A|
 \le\frac{e^{1/4}}{2-e^{1/4}}|H|\le2|H|.
\]
Indeed both matrices in the logarithm comparison are within
$e^{1/4}-1$ of one, their difference is at most
$e^{1/4}|H|$, and the logarithm series has Lipschitz constant
$(2-e^{1/4})^{-1}$ there. The cut-off factor in
$\check\Psi$ is fixed when $\mathsf a$ is fixed.
The shifted field therefore lies in the specified complex
ball, and $f_X$ is holomorphic with $|f_X|\le\|F_X\|$.

Cauchy's finite-difference estimate gives
$|f_X[\mathscr S\mid\emptyset]|
\le\|F_X\|\prod_{(b,S)\in\mathscr S}q(S)$, where
$q(S)=(R_S-1)^{-1}$.
For $0\le\kappa_1\le\widehat\kappa+1$ put
\[
 \omega(S)=(1+\#S)q(S)e^{\kappa_1(d_t(S)+4)}.
\]
Writing $n=d_t(S)+1$, we have
$1+\#S\le9n\le9e^{n-1}$ and $d_t(S)+4\le4n$.
The rate condition gives
\[
 \omega(S)\le
 \frac{18N_ea_H}{e\rho_{\rm bd}}\,
          e^{-3\check R_t}e^{-c_ed_t(S)}.
\]
After summation and the amplitude condition,
\[
 \sum_{S\in\mathsf{Pr}(b)}\omega(S)
                 \le(16N_b)^{-1}e^{-2\check R_t}.
\]

Let $\Theta(X,X_1)$ be the sum of
$\prod_{(b,S)\in\mathscr S}(1+\#S)q(S)$ over the nonempty
sets in \eqref{II-73-23}. By the length bound and the
product formula for subsets,
\[
 e^{\kappa_1d_t(X_1)}\Theta(X,X_1)
 \le e^{\kappa_1D}(e^\sigma-1),\qquad
 \sigma\le\tfrac12(D+1)e^{-2\check R_t},\quad D=d_t(X).
\]
Here $\#X\le8(D+1)$ and the bond count supplied the estimate
on $\sigma$. Since
$e^\sigma-1\le\sigma e^\sigma\le2e^{-2\check R_t}e^D$,
\[
 \Theta(X,X_1)\le
 2e^{-2\check R_t}e^{(\kappa_1+1)d_t(X)}
                              e^{-\kappa_1d_t(X_1)}.
\]
The cardinality bound in \eqref{eq:strip-transfer-geometry}
also gives the stronger intermediate estimate
\[
 |F_X^{\rm tr}(X_1)|
       \le\|F_X\|\frac{\#X}{\#X_1}\Theta(X,X_1)
\]
whenever the index set is nonempty.
This proves \eqref{II-73-35}. To sum over $X$, write
$\#X=\sum_{\Delta\subset X}1$. The sum at a fixed
$\Delta\subset X_1$ is bounded by the entropy estimate
because $\kappa_{\rm in}-\kappa_{\rm out}-1\ge c_e$.
Summing over $\Delta$ cancels $\#X_1$ and gives
$2N_ee^{-2\check R_t}b e^{-\kappa_{\rm out}d_t(X_1)}$,
which implies \eqref{II-73-37}.

Finally finite M\"obius summation gives
$f_X(\mathbf1)=f_X(0)+\sum_{\emptyset\ne\mathscr S}
f_X[\mathscr S\mid\emptyset]$.
Grouping by $X(\mathscr S)$ and summing the finite family
proves \eqref{II-73-39}. The geometry of each nonempty
index gives the support assertion.
\end{proof}

\subsection{The localized strip Jacobian}

For a fixed cut-off-zone bond and a real link potential $A$,
write
\[
 \mathsf{jc}_\delta(A;H)=
       \frac{\vartheta'_{\rm co}(|A|/\delta)}
                   {\delta|A|}\langle A,H\rangle,\qquad
 j_S=\mathsf{jc}_{\delta_k}(A;\mathbf H(b;S;p)),
 \qquad
 \bar j_S=\|\vartheta'_{\rm co}\|_\infty\bar h_S/\delta_k.
\]
The first expression is zero at $A=0$.
It is smooth in $A$, complex linear in the shift, and
$|j_S|\le\bar j_S$. Define
$f_b(\lambda)=\log(1-\sum_S\lambda(S)j_S)$ with its series
branch, and
\begin{equation}\label{II-73-25}
 J_b(X_1)=
 \sum_{\substack{\emptyset\ne\mathscr S\subset\mathsf{Pr}(b)\\
                         \bigcup\mathscr S=X_1}}
                     f_b[\mathscr S\mid\emptyset].
\end{equation}
In addition to \eqref{II-73-14}, suppose
\begin{equation}\label{II-73-40}
 \begin{gathered}
 3\mu_H\ge16(c_e+1),\qquad\check R_t\ge12,\qquad0<\rho_J\le1,\\
 128N_e^2\|\vartheta'_{\rm co}\|_\infty a_He^{-\check R_t}
                                            \le\rho_J\delta_k.
 \end{gathered}
\end{equation}
Then these terms are defined, and
\begin{equation}\label{II-73-41}
 |J_b(X_1)|\le\rho_J e^{-\check R_t}
                         e^{-\mu_H(d_t(X_1)+1)/2}.
\end{equation}
They vanish unless $X_1\in\mathsf{Pr}(b)$ and unless
$4\delta_k/3<|A|<5\delta_k/3$.
Their sum is the logarithm of the bond Jacobian:
\begin{equation}\label{II-73-43}
 \sum_{X_1}J_b(X_1)
   =\log\bigl(1-\mathsf{jc}_{\delta_k}(A;H_{\rm tot}(b))\bigr)
   =\log J_{H_{\rm tot}(b)}(e^{iA})
\end{equation}
when $H_{\rm tot}(b)$ is real and $|A|<\pi$.
The first equality also holds for a complex shift.

\begin{proof}
We have $\sum_S\bar j_S\le\rho_J/(128N_e)\le1/128$.
On the polydisc with radii
\[
 R_{J,S}=1+\frac{e^{-c_ed_t(S)}}{4N_e\bar j_S},
\]
the argument of the logarithm differs from one by less
than $1/128+1/4<1/2$. Thus $f_b$ is holomorphic and
$|f_b|\le\log2$. Cauchy's finite-difference estimate gives
\[
 |f_b[\mathscr S\mid\emptyset]|
 \le\log2\prod_{S\in\mathscr S}
                          4N_e\bar j_Se^{c_ed_t(S)}.
\]
Multiply by $e^{\mu_H(d_t(X_1)+1)/2}$ and use the fixed-bond
union bound in \eqref{eq:strip-transfer-geometry}.
After dropping the union constraint the result is at most
$\log2(e^\tau-1)$, where
\[
 \tau=\sum_{S\in\mathsf{Pr}(b)}
     4N_e\bar j_S
       e^{c_ed_t(S)+\mu_H(d_t(S)+4)/2}.
\]
For $n=d_t(S)+1\ge\check R_t\ge12$, the exponent left
after inserting $\bar j_S$ is
$-\mu_H n/2+3\mu_H/2+c_e(n-1)$.
It is at most $-3\mu_H n/8+c_e(n-1)$, hence at most
$-2\check R_t-c_ed_t(S)$ by \eqref{II-73-40}.
Therefore
$\tau\le\rho_Je^{-\check R_t}/32\le1/32$.
The inequality $\log2(e^\tau-1)\le2\tau$ proves
\eqref{II-73-41}, with the stronger factor $1/16$.

Finite M\"obius summation and linearity of $\mathsf{jc}$
give the first equality in \eqref{II-73-43}.
The same estimates give
$128\|\vartheta'_{\rm co}\|_\infty|H_{\rm tot}(b)|
\le\rho_J\delta_k$. The actual strip Jacobian formula
of \cref{II-26-71} then gives the second equality;
its determinant is positive, so the series is the real
logarithm. The support in $X_1$ follows from the union
geometry. The shell support and smoothness follow from
the corresponding properties of $\vartheta'_{\rm co}$,
which vanishes near zero.
\end{proof}

For the following preservation statements, assume the
corresponding properties of the original functions and
their domains. In a holomorphic family require, in addition,
that $(\zeta,\mathsf e)\mapsto
F_X(\mathsf a_\zeta+\mathsf e,w_\zeta)$ is jointly
holomorphic on the uniform room.
Both constructions are finite combinations of restrictions,
the local shift, the chart map and the original functions.
They preserve gauge invariance, conjugation reality, locality
and Borel measurability. They are holomorphic in holomorphic
families of background pairs and other arguments for which
the room is uniform and the values of the external field on
$\mathcal Z_\theta$ are fixed. On $\mathcal Z_1$ the chart
is holomorphic also in the external field. On
$\mathcal Z_\theta$ it is smooth in that real field.
Indeed the cut-off is then a fixed scalar for the first
claim, while in the second it is a smooth radial function
constant near zero. No holomorphic dependence on the
real cut-off variable is asserted.


\section{Localizing the change of background}\label{sec:localized-strip-shift}

The strip substitution uses a difference of two actual backgrounds.
We now express this difference as a sum of holomorphic local terms.
The construction first solves a reduced response equation on a set
of big cubes. Its datum measures the discrepancy of the tower rows.
Taking the finite differences in the cube parameters then produces
the local terms. At parameters equal to one, criticality and
uniqueness identify their sum with the background difference.

\subsection{The reduced response}

Use the reduced determining rows \(\mathbf B^\flat\), the strong
field seminorms of \cref{I-87-84}, and the dimensionless row pairing
with weight \(\xi_c^{-1}\). Write \(\hat Q=\mathcal Q'(0)\) in this
subsection, so that \(Q_{\rm phys}=\mathsf m^{-1}\hat Q\), where
\(\mathsf m(c)=\xi_c\). The straightened average and the action are
the ones already constructed:
\[
 \mathsf C(A)=\mathcal Q(A)-\hat QA,\qquad
 F(A)=\xi^{-1}\sum_{p\subset\Omega_0}a_p(e^{i\xi A}\mathbf U),
 \qquad \mathcal V(A)=F'(A)-F'(0)-\Delta A.
\]
Here \(\Delta=F''(0)\). In particular, the current of the pair
need not equal \(F'(0)\).

For the local estimates enlarge the neighbourhood of a top \(y\)
to \(N_L(y)\), consisting of the tops \(y'\) for which the
associated cubes have distance at most
\(L\min(\xi_y,\xi_{y'})\). It contains the earlier averaging
neighbourhood and has weighted radius
\(r_{\rm N}=3+6L^2\). Use
\[
 |A|_y=\xi_y\sup_{b\ {\rm over}\ y}|A(b)|,\qquad
 \|A\|_y=\max\{|A|_y,\xi_y^2\sup_{x\ {\rm over}\ y}
                         |\nabla_{\mathbf U}A(x)|\},\qquad
 |f|^\sharp_y=\xi_y^3\sup_{b\ {\rm over}\ y}|f(b)|.
\]
The supremum of a covariant difference uses only differences
defined in \(\Omega_0\). A subscript \(w\) denotes the supremum
after multiplication by \(w(y)\).

The proof of \cref{I-87-84}, with these enlarged neighbourhoods,
gives positive \(c_4,\rho_1\) and \(C_2,C_3,C_{\mathcal V},C_Q\ge1\),
depending only on \(L\), with the following estimates.
All maxima on the right are over \(N_L(y)\):
\begin{equation}\label{I-91-20}
 \begin{aligned}
 |\hat QA|_y&\le C_Q\max|A|_{y'},&
 |\hat Q^{\mathsf T}\varphi|^\sharp_y&\le C_Q\max|\varphi|_{y'},\\
 |\mathsf C(A)|_y&\le C_2\max|A|_{y'}^2,&
 |\mathcal V(A)|^\sharp_y&\le C_{\mathcal V}\max|A|_{y'}\max\|A\|_{y'},\\
 |\mathsf C'^{\,\mathsf T}_A\varphi|^\sharp_y
       &\le C_3\max|A|_{y'}\max|\varphi|_{y'}.
 \end{aligned}
\end{equation}
The maps vanish to the indicated orders at zero. On a ball of
strong radius \(r\le\frac14\min(c_4,\rho_1)\), their differences
satisfy the respective bounds \(4C_2r\max|A_1-A_2|\),
\(4C_{\mathcal V}r\max\|A_1-A_2\|\), and
\[
 |(\mathsf C'^{\,\mathsf T}_{A_1}
       -\mathsf C'^{\,\mathsf T}_{A_2})\varphi|^\sharp_y
 \le C_3\max|A_1-A_2|_{y'}\max|\varphi|_{y'} .
\]
They use only the pair and fields in these neighbourhoods.
Indeed the averaging kernel is supported on the double block
of a row, that block contains at most \(6(\xi_c/\xi)^3\)
positive bonds, and a bare bond belongs to at most eighteen such
blocks. In the transpose, the derivative's factor \(\xi_c\)
cancels the row weight \(\xi_c^{-1}\).
The kernel estimate \eqref{I-87-77}, its Cauchy derivative, and
the plaquette remainder estimate in \cref{I-87-84} therefore
give exactly these bounds. The latter proof also treats a bond
next to \(\Omega_0^c\): such a bond is in layer zero, where the
local spacing equals \(\xi\). No ambient plaquette outside
\(\Omega_0\) is added to \(F\).

We need two operators in addition to the normalized \(H,S\)
already constructed. Set
\begin{equation}\label{I-91-6}
 \begin{gathered}
 G_0=G[\mathbf U,0]
       =(\Delta+D\mathsf R D^{\mathsf T}
                           +\hat Q^{\mathsf T}\hat Q)^{-1},\\
 H_0=G_0\hat Q^{\mathsf T}
                  (\hat QG_0\hat Q^{\mathsf T})^{-1},
 \qquad \widetilde G=G_0-H_0\hat QG_0 .
 \end{gathered}
\end{equation}
The inverses exist by the retained regular inverse theorem at
\((\mathbf U,0)\). These are not the two operators \(K,\mathfrak G\)
in the first nonlinear response.

The local inverse expansion used in
\cref{prop:decoupled-local-propagators} gives \(G_0\) and the
physical-row lift \(H_{\rm phys}[\mathbf U,0]\).
Multiplying the latter on the right by \(\mathsf m^{-1}\)
gives \(H_0\). For \(\widetilde G\) use the terms of \(G_0\)
and the products
\[
 -H_{{\rm phys},\omega_1}[\mathbf U,0]\,
       Q_{\rm phys}(\mathbf U)\,G_{\omega_2}[\mathbf U,0].
\]
Only pairs for which the first cell of \(\omega_2\) is linked
to a cell of \(\omega_1\) contribute. The double block of each
row in the middle factor is contained in the first carrier.
Thus the product is defined on the union of the two carriers,
without a regular extension of the pair.

Here is the summation estimate for these extra products. If a
row \(c\) contributes, choose a bond \(b\) of its double block
where the middle average is estimated. A path of weighted
length at most \(4L^2\) joins \(y_c,y(b)\) and the first cell
of \(\omega_2\). If \(D\) is the length for the joined trail,
then
\[
 D\le D_{\omega_1}(y,y_c)+4L^2
                         +D_{\omega_2}(y(b),y').
\]
Splitting the two rates \(\delta/2\) into \(3\delta/8\)
and \(\delta/8\), and using
\(\delta R_1\mathfrak M\ge512\log L\), absorbs the scale
ratio and leaves a row sum at rate \(\delta/16\).
Write \(C_{\rm av,row}\) for the physical averaging row
bound in \cref{I-87-84}. There are at most six rows at
a top, so the extra factor is
bounded by
\[
 6C_{\rm av,row}L^4e^{3\delta L^2/2}
   \sum_{r\ge1}(3+2r)(2rL+1)^3e^{\delta/16-\delta r/32}.
\]
The first moment of the first walk expansion controls the
number of possible joining cells. Its bound \(4\), the mass
bound \(2\) for the second expansion, and the fixed number
of linked cells give a finite constant depending only on \(L\).
This proves the strong field and first-difference estimates
for the product. Absolute convergence identifies the sum as
\(G_0-H_0\hat QG_0\). The same argument proves local dependence,
holomorphy, gauge covariance and compatibility with conjugation.

Assign a parameter once to each big cube meeting the hull of
each such term, including a product term. Fix a real number
\(\varkappa>0\) and use a partition
\(\mathsf P=\pi[\mathsf w]\) with at least three cubes in each
direction and
\begin{equation}\label{I-91-45}
 \mathring\delta_0R_1\mathfrak M\ge512\log L,\qquad
 \mathsf w\ge88r_{\rm w}\mathfrak M,\qquad
 \mathring\delta_0\mathsf w\ge2176\varkappa.
\end{equation}
The other block and inverse floors are those of
\cref{prop:decoupled-local-propagators}. Its hull-counting
argument applies also to the joined trails. For
\(|s_\Delta|\le e^\varkappa\), the normalized operators and
arbitrary subsums satisfy
\begin{equation}\label{I-91-48}
 \begin{aligned}
 \|H(s)e_{c'}\|_y,\ \|H_0(s)e_{c'}\|_y
   &\le B_\varkappa\frac{\xi_y}{\xi_{y'}}
                      e^{-3\mathring\delta_0d(y,y')/16}|e|,\\
 |S(s)(c,c')|
   &\le B_\varkappa\frac{\xi_y}{\xi_{y'}}
                      e^{-3\mathring\delta_0d(y,y')/16},\\
 \|\widetilde G(s)f_{y'}\|_y
   &\le B_\varkappa\left(\frac{\xi_y}{\xi_{y'}}\right)^3
              e^{-3\mathring\delta_0d(y,y')/16}|f_{y'}|^\sharp ,
 \qquad B_\varkappa=B_{\flat,\mathrm s}e^{8\varkappa}.
 \end{aligned}
\end{equation}
Here \(c,c'\) are at \(y,y'\), and \(f_{y'}\) vanishes off
the bonds over \(y'\). The terms for \(H,S\) are normalized
as \(H_{\rm phys}\mathsf m^{-1}\) and
\(\mathsf m^3S_{\rm phys}\mathsf m^{-1}\).
Testing \(H(s)^{\mathsf T}f\) at \(c'\) contributes
\(\xi_{y'}\), while summing bare bonds over \(y\) contributes
at most \(3\xi_y^3\). Consequently the scale factors cancel:
\begin{equation}\label{I-91-49}
 |H(s)^{\mathsf T}f|_{y'}
 \le6B_\varkappa\sum_y
          e^{-3\mathring\delta_0d(y,y')/16}|f|^\sharp_y .
\end{equation}
For the other operators absorb
\((\xi_y/\xi_{y'})^p\le
 L^3e^{3\mathring\delta_0d(y,y')/512}\), \(0\le p\le3\).
Since at most six rows lie at a top, the four maps
\[
 \bigl(H_0(s),\widetilde G(s),H(s)^{\mathsf T},S(s)-I\bigr)
\]
have exponentially summable local norm kernels with constant
\(b_{0,\mathrm s}e^{8\varkappa}\) and rate
\(\mu=11\mathring\delta_0/64\), where
\(b_{0,\mathrm s}=6L^3B_{\flat,\mathrm s}+1\).
The identity in \(S(s)-I\) carries no decoupling parameters.

Put
\[
 \theta_0=\mathring\delta_0/8,\quad
 e_0=e^{r_{\rm N}\theta_0},\quad
 S_L(a)=9L^3\sum_{r\ge1}(2r+1)^3e^{-a(r-1)/2},\quad
 \widehat b_0=b_{0,\mathrm s}S_L(3\mathring\delta_0/64).
\]
The weighted top-counting estimate gives operator norm at
most \(\widehat b=\widehat b_0e^{8\varkappa}\) for these
four maps whenever \(w\in\mathcal W_\theta\),
\(0\le\theta\le\theta_0\).
Choose \(a_{\mathrm D,\mathrm s}\) within the common local
analytic and inverse radii, and small enough that
\(2\widehat b_0a_{\mathrm D,\mathrm s}<r_b\), with the
seam radius \(r_b\) of \eqref{eq:critical-response-radius}.
Set
\begin{equation}\label{I-91-60}
 \Gamma=C_2+C_3+C_{\mathcal V},\qquad
 \varepsilon_{\mathrm D,\mathrm s}
 =\min\left\{a_{\mathrm D,\mathrm s},
 \frac{\min(c_4,\rho_1)}{32\widehat b_0},
 \frac1{32\Gamma e_0\widehat b_0^2}\right\},\qquad
 C_{\mathrm D,\mathrm s}=2\widehat b_0 .
\end{equation}
All these constants depend only on \(L\).

\begin{proposition}\label{I-91-59}
Assume the preceding geometry and let the natural pair have
curvature and potential radii at most \(a_{\mathrm D,\mathrm s}\)
and current radius \(\gamma_0\le
\varepsilon_{\mathrm D,\mathrm s}e^{-16\varkappa}\).
For \(|B'|_\infty\le\varepsilon_{\mathrm D,\mathrm s}e^{-16\varkappa}\)
the equations
\begin{equation}\label{I-91-62}
 \begin{aligned}
 T&=H_0(s)(B'-\mathsf C(T))
      -\widetilde G(s)
              [\mathcal V(T)-\mathsf C'^{\,\mathsf T}_T\lambda],\\
 \lambda&=H(s)^{\mathsf T}\mathbf J
      +(S(s)-I)(B'-\mathsf C(T))
      +H(s)^{\mathsf T}
              [\mathcal V(T)-\mathsf C'^{\,\mathsf T}_T\lambda]
 \end{aligned}
\end{equation}
have exactly one solution in the ball
\[
 \max\{\|T\|,|\lambda-H(s)^{\mathsf T}\mathbf J|_\infty\}
 \le C_{\mathrm D,\mathrm s}\varepsilon_{\mathrm D,\mathrm s}
                                      e^{-8\varkappa}.
\]
The solution is holomorphic on the open domains, is gauge
covariant, commutes with complex conjugation, and satisfies
\[
 \max\{\|T(B'_1)-T(B'_2)\|_w,
             |\lambda(B'_1)-\lambda(B'_2)|_w\}
 \le C_{\mathrm D,\mathrm s}e^{8\varkappa}|B'_1-B'_2|_w .
\]
In particular \(T(0)=0\). Setting all parameters outside a
set of cubes to zero splits \(T\) into its wall components.
The component response depends only on the pair in that
component and on the datum rows in the neighbourhood of its
operator carrier; it vanishes if the datum on that carrier
vanishes. The same conclusions hold for local pairs, with
the response defined intrinsically by these component equations.
\end{proposition}

\begin{proof}
Set \(\lambda^0=H(s)^{\mathsf T}\mathbf J\), \(\nu=\lambda-\lambda^0\)
and \(g=|\lambda^0|_\infty\le\widehat b\gamma_0\).
For \(r=|B'|_\infty\) or any larger allowed datum radius use
the ball \(\max(\|T\|,|\nu|)\le R=2\widehat b r\).
Let
\[
 N(T,\nu)=\mathcal V(T)-\mathsf C'^{\,\mathsf T}_T(\lambda^0+\nu).
\]
The right sides of \eqref{I-91-62}, with \(\lambda^0\)
subtracted in the second equation, define a map \(\Phi_{B'}\)
of \((T,\nu)\).
For two points in the ball, \eqref{I-91-20} gives
\[
 |\mathsf C(T_1)-\mathsf C(T_2)|_w
       +|N(T_1,\nu_1)-N(T_2,\nu_2)|^\sharp_w
 \le e_0(4\Gamma R+C_3g)
          \max(\|T_1-T_2\|_w,|\nu_1-\nu_2|_w).
\]
In the transpose term split the difference into the change
of \(\mathsf C'_T\) and the change of \(\nu\); their bounds
are \(C_3e_0(g+R)\|T_1-T_2\|_w\) and
\(C_3e_0R|\nu_1-\nu_2|_w\). This explains the coefficient
in the displayed inequality.
The radius choice gives
\[
 R\le\tfrac1{16}\min(c_4,\rho_1),\qquad
 4\widehat b e_0\Gamma R\le\tfrac14,\qquad
 \widehat b e_0C_3g\le\tfrac14 .
\]
Thus \(\Phi_{B'}\) contracts by \(1/2\) in every stated
weighted norm. Its value at zero has norm at most
\(\widehat b|B'|\), so it preserves the radius-\(R\) ball.
Iteration proves existence and uniqueness. Comparing two
fixed points and absorbing the contraction term proves the
weighted estimate. The same uniform iteration proves
holomorphy; at zero both nonlinear terms vanish.

For completeness, the component argument must retain the
unparameterized identity. Each partial walk sum is carried
strictly inside its wall component, with four
\(N_L\)-neighbourhoods still inside that component.
To check this margin, let \(z\) be an endpoint in a cube
of level \(n\) occurring in a term. After \(q\le4\)
neighbourhood steps, the scale is at most \(L^2\xi_n\)
and the whole top cube is within
\[
 D_q=((q+2)L^2+qL^3)\xi_n
\]
of \(z\). Inductively, the next cube has a point within
\(L^3\xi_n\) of the preceding one. Separation from
\(\Omega_{n+3}\), since \(z\in\Omega_{n+2}^c\), prevents
its scale from exceeding \(L^2\xi_n\).
The retained block floor gives \(\mathfrak M\ge L^3\);
therefore
\(D_4+\xi\le5L^3\xi_n\le5\mathfrak M\xi_n\).
All four neighbourhoods, including the endpoints of their
bare bonds and the rows through the first three, lie in the
enlarged cube and hence in the hull.
Every partition cube containing a point of the hull is
included in its parameter set, so the hull lies in the
interior of their union.
Carriers belonging to different wall components have
disjoint such neighbourhoods.
The local nonlinear maps consequently agree with the
component maps near each carrier.
For a component \(i\), the response is supported on its
carrier. Its multiplier has the form
\(\lambda_i=-B'+\lambda_i^{\rm loc}\), where
\(\lambda_i^{\rm loc}\) is supported in the carrier's
neighbourhood. In the combined system use
\(\lambda=-B'+\sum_i\lambda_i^{\rm loc}\).
On the support of \(T_i\) this agrees with \(\lambda_i\);
the local equations therefore combine to
\eqref{I-91-62}. Uniqueness gives the splitting.
This also shows independence of unused datum rows.
All averaging blocks and plaquettes entering these
comparisons lie in the corresponding component. Thus their
formulas depend only on its pair, and equal restrictions
give the same iteration without choosing a global regular
extension. Finally the local matrices, nonlinear maps and
pairings intertwine unitary gauge transformations and
conjugation, so every iterate does as well.
\end{proof}

\subsection{Intrinsic regional response and criticality}

Here a local natural pair means the restrictions of the
inequalities in \cref{def:natural-regular-pair}, including
its containing-box gauge condition; a regular continuation
outside the region is not required. For
\(\mathfrak k\subset\mathsf P\), put
\(W_{\mathfrak k}=\Omega_0\cap\bigcup\mathfrak k\), with
cube-wise intersection. In each walk series retain exactly
the indices whose parameter sets lie in \(\mathfrak k\).
Denote the resulting operators by \(O^{\mathfrak k}(s)\).
Their terms use only the pair on \(W_{\mathfrak k}\).
Let \(Y^\circ_{\mathfrak k}\) be the union of their top
carriers and put \(Y^{[q]}_{\mathfrak k}=N_L^q(Y^\circ_{\mathfrak k})\).
The preceding margin proof gives
\(Y^{[4]}_{\mathfrak k}\subset W_{\mathfrak k}\), with
the associated bare bonds and averaging blocks understood.

Retain \(\hat Q,\mathsf C\) only at rows over
\(Y^{[2]}_{\mathfrak k}\). On bonds over this set retain
the nonlinear remainder of
\[
 F_{W_{\mathfrak k}}(A)
   =\xi^{-1}\sum_{p\subset W_{\mathfrak k}}
                           a_p(e^{i\xi A}\mathbf U).
\]
Set these three maps to zero elsewhere; call them
\((\hat Q_{\mathfrak k},\mathsf C_{\mathfrak k},
\mathcal V_{\mathfrak k})\). Replace the four operators
in \eqref{I-91-62} by their \(\mathfrak k\)-sums and the
current by its zero extension. This is the intrinsic
regional system:
\begin{equation}\label{II-70-121}
 \eqref{I-91-62}\quad\hbox{with }
 (H_0,\widetilde G,H,S;\mathsf C,\mathcal V,\mathbf J)
 \ \hbox{replaced by }\
 (H_0^{\mathfrak k},\widetilde G^{\mathfrak k},
 H^{\mathfrak k},S^{\mathfrak k};
 \mathsf C_{\mathfrak k},\mathcal V_{\mathfrak k},\mathbf J^0).
\end{equation}
As before the identity subtracted from \(S^{\mathfrak k}\)
is not truncated or parameterized.

\begin{proposition}\label{II-70-150}
The regional system has the existence, uniqueness, estimates
and holomorphic dependence of \cref{I-91-59}. Its field
component, denoted \(\mathbf H_{\mathfrak k}(s;\mathbf p;B')\),
is zero off \(Y^\circ_{\mathfrak k}\), uses \(B'\) only
on \(Y^{[1]}_{\mathfrak k}\), and vanishes if \(B'\) is
zero on \(Y^\circ_{\mathfrak k}\).
For \(\mathfrak k\subset\mathfrak k'\), zero extension of
the parameters and restriction of the pair give the same
response. It splits over wall components, intertwines
unitary gauges and conjugation, and agrees with the global
response when a regular global pair is available.
\end{proposition}

\begin{proof}
The margin places every double row block over
\(Y^{[3]}_{\mathfrak k}\), and every plaquette of \(\Omega_0\)
touching the corresponding bonds, inside \(W_{\mathfrak k}\).
Thus the local formulas used in \eqref{I-91-20} are unchanged
near the carrier. Multiplication of the row maps by an
indicator does not enlarge their transpose estimates.
The regional action remainder has the same plaquette
formula there. Consequently the truncated triple satisfies
the same local estimates.

For a field supported on the carrier, compute its regional
first difference by extending the field by zero:
\[
 \nabla_\mu^W A(b)=\xi^{-1}
 \{\operatorname{Ad}(\mathbf U(\langle b_-,b_-+\xi e_\mu\rangle))
                      A(b+\xi e_\mu)-A(b)\},
\]
where the first term is zero when the translated bond is
not in \(W\). If both bonds belong to \(W\), the joining
bond does also. The margin shows that this norm agrees
with the norm for any continuation of \(\mathbf U\).
The partial walk bounds therefore act on this intrinsic
norm. The contraction proof of \cref{I-91-59} now applies
without any regularity assumption on a continuation.
It proves uniqueness, and hence independence of any
continuation used to write the zero extensions.

For holomorphy use the same finite-dimensional local
charts as for the walk terms. In a chart \(e^a\mathbf U\)
with decomposition \(\mathbf U=e^{i\xi A'}U\), the new
potential is
\((i\xi)^{-1}\log(e^ae^{i\xi A'})\). If
\(|a|<1/5\) and \(|\xi A'|<1/6\), its argument stays
within \(1/2\) of the identity. The strict local class
inequalities persist on a smaller chart. Uniform iteration
then proves holomorphy jointly in the pair, the datum and
the cube parameters.

Support and dependence on datum rows follow directly
from the supported operators and their transposes.
For restriction, the surviving walk terms are identical;
near their carriers the two truncated triples have the
same averaging blocks and plaquettes. The component
argument at the end of \cref{I-91-59} applies, including
the term \(-B'\) in the multiplier. Uniqueness proves
restriction, splitting, and agreement with the global
system. The same iteration proves both covariance claims.
\end{proof}

The free-boundary current requires a separate convention.
A bond is \emph{full} when all four ambient plaquettes
through it belong to \(\Omega_0\). This designation refers
to \(\Omega_0\), even for a pair given on \(W_{\mathfrak k}\).
Write \(J_W(\mathbf U)=F'_W(0)\) and set
\begin{equation}\label{II-70-123}
 \mathbf p^{\rm tc}_{\mathfrak k}
   =(\mathbf U,\mathbf J^{\rm tc}_{\mathfrak k}),\qquad
 \mathbf J^{\rm tc}_{\mathfrak k}(b)=
 \begin{cases}
 \mathbf J(b),&b\text{ full},\\
 J_{W_{\mathfrak k}}(\mathbf U)(b),&b\text{ not full}.
 \end{cases}
\end{equation}
At a nonfull bond, its initial point and the base points
of its adjacent plaquettes of \(\Omega_0\) lie in layer
zero: they are within \(2\xi\) of a point outside
\(\Omega_0\), whereas deeper layers have the retained
separation. If the curvature radius there is
\(\alpha'_0\), the current obeys
\begin{equation}\label{II-70-161}
 \xi^3|J_{W_{\mathfrak k}}(\mathbf U)(b)|
             <24\alpha'_0 .
\end{equation}
Indeed, each plaquette curvature satisfies
\(\frac12\xi^{-2}|g-g^{-1}|\le\xi^{-2}|g-1|\).
The inverse identity for \(G^c\) gives
\(|g^{-1}-1|=|g-1|\); the small-potential bound makes
the adjacent adjoint transports smaller than \(2\).
The two transverse directions in the plaquette derivative
therefore give at most \(6\alpha'_0\xi^{-3}\), which
is smaller than the stated bound.

If \(24\alpha'_0\le\gamma_0\), truncation preserves the
current inequality of the natural class. Moreover it
commutes with restriction for the response:
\begin{equation}\label{II-70-163}
 \mathbf H_{\mathfrak k'}(s^0;
             \mathbf p^{\rm tc}_{\mathfrak k'};B')
 =\mathbf H_{\mathfrak k}
       (s;(\mathbf p|_{W_{\mathfrak k}})^{\rm tc}_{\mathfrak k};B').
\end{equation}
To see this, every plaquette of \(\Omega_0\) through a
bond of a surviving walk carrier already lies in
\(W_{\mathfrak k}\). Thus the two regional action currents
coincide at precisely the bonds used by those walks and
by \(H^{\mathfrak k\,\mathsf T}\mathbf J^0\).
The two systems and their uniqueness balls are identical.

At parameters one we will use the actual critical current.
The unparameterized walk sums satisfy the same bounds
with \(\varkappa=0\); hence the datum radius is
\(\varepsilon_{\mathrm D,\mathrm s}\) and
\(\|T\|\le C_{\mathrm D,\mathrm s}|B'|_\infty\).

\begin{proposition}\label{I-91-94}
Let \(U_0\) be real, let
\((U_0,J_{\Omega_0}(U_0))\) satisfy the natural pair
bounds with curvature and potential radii at most
\(a_{\mathrm D,\mathrm s}\) and current radius at most
\(\varepsilon_{\mathrm D,\mathrm s}\), and suppose
\(J_{\Omega_0}(U_0)=\hat Q^{\mathsf T}\lambda_0\).
For \(|B'|_\infty<\varepsilon_{\mathrm D,\mathrm s}\),
the reduced response at parameters one satisfies
\begin{equation}\label{I-91-95}
 \mathcal Q(T)=B',\qquad
 \mathsf R D^{\mathsf T}T=0,\qquad
 F'(T)=(\hat Q+\mathsf C'_T)^{\mathsf T}\lambda .
\end{equation}
For real \(B'\), the configuration \(e^{i\xi T}U_0\)
is critical for the actual row constraint with datum
\((e^{iB'^\natural}V_0)^{v[T]}\), where
\(V_0=M_{\mathbf B^+}U_0\), the outward-row twist is
\eqref{eq:critical-seam-twist}, and \(v[T]\) is the
accumulated node factor.
\end{proposition}

\begin{proof}
The inverse identities give
\[
 \hat QH_0=I,\quad \hat Q\widetilde G=0,\quad
 \widetilde GG_0^{-1}=I-H_0\hat Q,\quad
 \ker\widetilde G=\operatorname{ran}\hat Q^{\mathsf T},
 \qquad \hat QH=I .
\]
Applying \(\hat Q\) to the first response equation gives
\(\mathcal Q(T)=B'\). Put \(K_T=\hat Q+\mathsf C'_T\),
\(z=F'(T)-K_T^{\mathsf T}\lambda\) and
\(\phi=\mathsf RD^{\mathsf T}T\). The two equation
defects are
\begin{equation}\label{I-91-91}
 0=\widetilde G(z+D\phi),\qquad
 0=\mathsf L G'\phi-H^{\mathsf T}(z+D\phi)
\end{equation}
for a fixed linear map \(\mathsf L\).
Here is the term accounting for this map. The correction
\(\Delta'_\pi\) in the Hessian is the polarization of
the quadratic gauge expression
\eqref{eq:critical-gauge-expansion}. With
\(\Upsilon=G'\mathsf RD^{\mathsf T}\), it has the form
\[
 (A_1,\Delta'_\pi A_2)
 =b(\Upsilon A_1,A_2)+b(\Upsilon A_2,A_1)
                       +b'(\Upsilon A_1,\Upsilon A_2).
\]
The retained linear identity \(\Upsilon H=0\) gives
\(H^{\mathsf T}\Delta'_\pi T=\mathsf L G'\phi\).
Also
\[
 H^{\mathsf T}(\Delta-\Delta'_\pi+D\mathsf RD^{\mathsf T})
              =(S-I)\hat Q,\qquad H^{\mathsf T}J=\lambda_0 .
\]
Substitution of \(F'(T)=J+\Delta T+\mathcal V(T)\)
and \(\hat QT=B'-\mathsf C(T)\) proves both defects.
Their first equation says \(z+D\phi=\hat Q^{\mathsf T}\sigma\);
their second says \(\sigma=\mathsf LG'\phi\). Thus
\[
 F'(T)+D\phi=K_T^{\mathsf T}\lambda
                            +\hat Q^{\mathsf T}\mathsf LG'\phi .
\]

For real data sufficiently near zero, use the gauge
bijection \eqref{eq:critical-gauge-bijection} to select
a gauge tangent \(X_\chi(T)\) with node variation zero
and \(\mathsf RD^{\mathsf T}X_\chi(T)=\phi\).
Its norm is at most a finite constant times \(\|\phi\|\).
The gauge law gives \(K_TX_\chi(T)=0\), and action
invariance gives \((F'(T),X_\chi(T))=0\). Pairing the
last display with \(X_\chi(T)\) therefore yields
\[
 (\phi,\phi)
   =-(\mathsf LG'\phi,\mathsf C'_TX_\chi(T)),
 \qquad
 |(\mathsf LG'\phi,\mathsf C'_TX_\chi(T))|
       \le C|T|_\infty(\phi,\phi).
\]
At this step \(C\) may depend on the finite geometry.
Since \(T\to0\) with \(B'\), it follows that \(\phi=0\)
near the real origin. Holomorphy on the connected datum
ball gives \(\phi=0\) throughout it. The defects then
give \(z=0\), proving \eqref{I-91-95}.

Finally the actual averaging identity \cref{I-88-80}
gives the stated row datum. The same node Ward argument
as in \eqref{eq:critical-coarse-ward} applies: pair
\(F'(T)=K_T^{\mathsf T}\lambda\) with gauge tangents,
use the gauge bijection near zero, and continue the
resulting identity holomorphically. It annihilates the
node term in the derivative of any curve in the actual
constraint, so the derivative of the action along that
curve is zero. The radius choice above keeps the seam
twist and its invertible derivative in their retained
domain. This proves criticality, including outward rows.
\end{proof}

\subsection{The response on a joined component}

Fix a step \(t=k+1\) and a component \(\mathfrak o\) of
the localization region in \cref{II-13-11}. Its given
sequence is denoted \(\Omega^\diamond_\bullet\) and its
joined sequence \(\Omega^{\mathfrak o}_\bullet\).
For the remainder of the section use level-\(k\) lengths,
so \(\xi_n=L^{n-k}\). Let
\(\mathsf P^{\rm loc}_{\mathfrak o}\) be the big cubes
meeting \(\mathfrak o\), and let
\(\mathsf P^\partial_{\mathfrak o}\) be those containing
an endpoint of a tower row. Rows always have the inward
rows removed.

We spell out the pair classes needed for the local terms.
For \(W=\mathsf S\cap\mathfrak o\) a pair
\(\mathbf p=(e^{i\eta_kA'}U,\mathbf J)\), with \(U\)
real, belongs to
\(\mathfrak P^{\mathfrak o}(\mathsf S;\alpha'_0,\alpha'_1)\)
if, for the given sequence, it satisfies the following
inequalities. The curvature of both configurations is
smaller than \(\alpha'_0L^{-2n}\) at a plaquette based
in layer \(n\). The potential and its covariant first
differences satisfy
\(\xi_n|A'|<\alpha'_1\) and
\(\xi_n^2|\nabla^UA'|<\alpha'_1\) at their initial sites
in that layer. The current satisfies
\(\xi_n^3|\mathbf J|<\alpha_2\).
For every tested cube \(q\in\mathsf Q_n\) contained in
\(W\), \(U\) has a real gauge on a containing cube
\(q\subset q'\subset W\), in which
\(\xi_n|A|\le\alpha_0\) and
\(\xi_n^2|\nabla A|\le\alpha_0\).
Here \(\alpha_0,\alpha_1,\alpha_2\) are fixed positive
radii, and \(0<\alpha'_i\le\alpha_i\), \(i=0,1\).
Define
\begin{equation}\label{II-71-18}
 \begin{split}
 \mathfrak P^{\mathfrak o}_{\mathbf H}(\mathsf S)
 &=\bigcup_{\alpha'_0+\alpha'_1\le\varrho_{\mathbf H}}
       \mathfrak P^{\mathfrak o}(\mathsf S;\alpha'_0,\alpha'_1),\\
 \mathsf{bs}_{\mathsf S}(\mathbf p)
 &=\inf\{\alpha'_0+\alpha'_1:
       \mathbf p\in\mathfrak P^{\mathfrak o}
                          (\mathsf S;\alpha'_0,\alpha'_1)\}.
 \end{split}
\end{equation}
These classes are open: keep the real configuration and
its gauges fixed, absorb a nearby configuration into
\(A'\) with the local series logarithm, and use the
strict curvature, potential and current inequalities.
They are stable under unitary gauges, conjugation and
restriction; restriction can only decrease \(\mathsf{bs}\).

Choose the response radii in the following order:
\begin{equation}\label{II-71-59}
 \begin{gathered}
 0<\beta'_0\le\beta_0\le a_{\mathrm D,\mathrm s},\quad
 0<\beta_1\le a_{\mathrm D,\mathrm s},\quad
 0<\beta_2\le\varepsilon_{\mathrm D,\mathrm s}e^{-16\bar\kappa},
 \quad24\beta'_0\le\beta_2,\\
 0<\varrho_{\rm rw}\le\varepsilon_{\mathrm D,\mathrm s}e^{-16\bar\kappa},
 \quad \alpha_0\le\beta_0,\quad \alpha_2\le\beta_2,\quad
 \varrho_{\mathbf H}\le\min(\beta'_0,\beta_1,1/3).
 \end{gathered}
\end{equation}
Assume the inverse floors above and
\[
 M\ge88r_{\rm w}\mathfrak M,\qquad
 \mathring\delta_0M\ge2176\bar\kappa,\qquad
 m'<m+K-t .
\]
The usual strip floors \(M_1\ge R_1\mathfrak M\) and
\(M\ge L^4M_1\) remain in force.

\begin{proposition}\label{II-71-70}
For every pair in \(\mathfrak P^{\mathfrak o}_{\mathbf H}(\mathsf S)\),
the truncated pair is a natural local pair for the joined
sequence, with radii \((\beta_0,\beta_1,\beta_2)\), and
with curvature radius \(\beta'_0\).
Its response is
\begin{equation}\label{II-71-66}
 \mathbf H^{\mathfrak o}_{\mathsf S}(s;\mathbf q;B')
   =L^{-1}\mathbf H_{\mathsf S}(s;
                         (\mathbf U,L^3\mathbf J);B'),
 \qquad |B'|_\infty<\varrho_{\rm rw}.
\end{equation}
The right side uses level-\(t\) lengths and the joined
sequence. All estimates, holomorphy, restriction,
component splitting and covariance properties above apply.
In the dimensionless strong norm for the joined sequence,
\begin{equation}\label{II-71-67}
 \|\mathbf H^{\mathfrak o}_{\mathsf S}(s;\mathbf q;B')\|
       \le\mathsf A_{\rm rs}|B'|_\infty,\qquad
 \mathsf A_{\rm rs}=C_{\mathrm D,\mathrm s}e^{8\bar\kappa}.
\end{equation}
At parameters one, for the whole component and a real
critical pair \((U_0,J_{\mathfrak o}(U_0))\), this is the
critical response of \cref{I-91-94}. It also satisfies
\begin{equation}\label{eq:strip-response-one}
 \|T\|\le C_{\rm rs}|B'|_\infty,\qquad
 |F'_{\mathfrak o}(T)-J_{\mathfrak o}(U_0)|^\sharp
       \le C_{\rm rs}|B'|_\infty,\qquad
 C_{\rm rs}=(C_Q+1)e_0C_{\mathrm D,\mathrm s}.
\end{equation}
\end{proposition}

\begin{proof}
The joined layer index is no greater than the given one,
which on \(\mathfrak o\) is either \(k\) or \(t\).
Hence every curvature, potential, difference and current
inequality remains valid when its given local scale is
replaced by the joined local scale.
For the gauge condition, let a joined test cube of level
\(n\) lie in \(W\). If it is contained in \(\Omega_t\),
enlarge it, when \(n<t\), to a cube of side
\(2\mathfrak M\xi_t\) inside \(W\cap\Omega_t\).
Otherwise enlarge it, when \(n<k\), to a cube of side
\(2\mathfrak M\xi_k\) inside \(W\).
This is possible because \(W\) is a union of
level-\(t\) \(M_1\)-cubes,
\(12\mathfrak M\xi_n<\mathfrak M\xi_{n+1}\), and
the containing grid has at least twice that side.
In the second case the enlarged cube meets layer \(k\);
in the first it lies in layer \(t\). The given gauge
condition therefore supplies a containing gauge cube
inside \(W\). Its stronger scale bounds imply those
at level \(n\). This proves the natural gauge condition,
also for a joined test cube not meeting its own layer.

For the free boundary use \eqref{II-70-161}.
The current radius becomes at most
\(\max(\alpha_2,24\alpha'_0)\le\beta_2\), while all
other conditions are unchanged. The current is a
polynomial in the configuration entries and inverses,
so truncation is holomorphic and commutes with gauges
and conjugation.

On changing from level \(k\) to level \(t\), lengths
and forms are multiplied by \(L^{-1}\) and \(L\),
respectively, and currents by \(L^3\). Thus the
dimensionless inequalities and strong norms are unchanged,
and truncation commutes with this rescaling.
The partition side becomes \(M\). The strict period
condition gives at least \(2L\ge3\) cubes per direction.
Apply \cref{II-70-150} and scale the response back.
All walks for the whole component have parameter sets
in \(\mathsf P^{\rm loc}_{\mathfrak o}\), so putting its
parameters equal to one gives the unparameterized inverse.

To apply \cref{I-91-94}, a critical reference has the
required multiplier. At zero \(\hat QH=I\). The implicit
function theorem therefore makes every vector in
\(\ker\hat Q\) tangent to a curve with constant
straightened datum. The averaging identity
\cref{I-88-80} corrects this curve by its node gauge
to one with constant actual datum. Gauge invariance
and criticality make \(F'(0)\) orthogonal to
\(\ker\hat Q\), whence \(F'(0)=\hat Q^{\mathsf T}\lambda_0\).
Small curvature keeps the real and complex means in
the same unique mean chart throughout these curves.

The field estimate at parameters one was proved above.
For the current difference write
\[
 F'(T)-J=\hat Q^{\mathsf T}(\lambda-H^{\mathsf T}J)
                            +\mathsf C'^{\,\mathsf T}_T\lambda .
\]
The contraction estimate and \eqref{I-91-20}, with the
constant weight, bound the first term by
\(C_Qe_0C_{\mathrm D,\mathrm s}|B'|\).
For the second, use \(|H^{\mathsf T}J|\le\widehat b_0\beta_2\),
\(\|T\|,|\lambda-H^{\mathsf T}J|\le
C_{\mathrm D,\mathrm s}|B'|\), and the radius choice,
which makes its coefficient at most
\(e_0C_{\mathrm D,\mathrm s}\). The resulting bound is
\eqref{eq:strip-response-one}. Scaling the derivative
by \(L^{-3}\) preserves its sharp norm.
\end{proof}

\subsection{The tower source and its read-out}

We first record the averaging estimates used at the
two ends of the construction. On a double-block tower
of height \(j\), let \(U\) be real with bare curvature
at most \(\alpha L^{-2j}\), and put
\(\mathbf U=e^{ia'}U\), \(A=L^j\sup|a'|\).
With \(\kappa_6\) as above, put \(C''_{\mathrm{ia}}=4(8\kappa_6+2)\).
Choose an auxiliary radius
\[
 0<c_\sharp\le\min\{c_{\mathsf m},c_{\rm plain},1\},\qquad
 C''_{\mathrm{ia}}c_\sharp<r_{\mathrm{av}}/16.
\]
There is
\[
 0<c_{\mathrm{ia}}\le
 \min\{c_\sharp/4,a_{\mathsf m},a_{\rm plain},
                              \varpi_{\mathrm{av}}/2\},
\]
depending only on \(L\), with the following properties.
For \(\alpha<c_{\mathrm{ia}}\) and \(A<c_\sharp\), the
averages and contour variables in the first three lines below
are defined and satisfy the displayed bounds. If also
\(A,B=L^j\sup|a_0|<c_{\mathrm{ia}}\), the relative recursion
and accumulated factors below are defined:
\begin{equation}\label{II-155-25}
 \begin{gathered}
 W_{0,n}=\mathcal M^nU,\qquad
 W_n=\mathcal M^n\mathbf U=e^{ia''_n}W_{0,n},\\
 \varpi(W_{0,n})\le2\alpha L^{2(n-j)},\qquad
 |a''_n|\le C''_{\mathrm{ia}}L^{n-j}A,\\
 W_n(y,x)=e^{i\zeta_n(y,x)}W_{0,n}(y,x),\qquad
 |\zeta_n|\le\kappa_6C''_{\mathrm{ia}}L^{n-j}A,\\
 a_{n+1}=\mathcal Q^1_{\mathsf m}(W_n,a_n),\qquad
 |a_n|\le2L^{n-j}B,\qquad
 v_n=\prod_{i<n}e^{i\lambda_i},\quad
 \sum_{i<n}|\lambda_i|\le\frac{2\kappa_6}{L-1}B .
 \end{gathered}
\end{equation}
The reference curvature estimate is for \(n<j\).
All quantities are holomorphic in the perturbations;
when the input is real the two definitions of the
mean agree.

Here are the extra estimates needed beyond the real
reference iteration in \cref{I-39-36,I-27-78}.
For \(W_n\), that iteration first gives its centered
perturbation \(a'_n\) and factors \(v'_n\). The relative
product
\[
 v'_n(c_-)e^{ia'_n(c)}
       \operatorname{Ad}(W_{0,n}(c))(v'_n(c_+)^{-1})
\]
has total exponent norm at most
\((8\kappa_6L^{-1}+2)L^{n-j}A\).
Its series logarithm gives the bound for \(a''_n\).
The choices \(2c_{\mathrm{ia}}\le\varpi_{\mathrm{av}}\) and
\(C''_{\mathrm{ia}}c_\sharp<r_{\mathrm{av}}/16\) put these
references and perturbations in the complex one-step chart.
Its contour Schwarz bound gives the estimate for \(\zeta_n\).
For clarity, the one-step estimates at this complex
reference follow from the small-gauge estimates as
follows. On a double block choose a real gauge
\(W_0^h=e^{ia^0}\), \(\sup|a^0|\le r_{\mathrm{av}}/2\).
Let \(c_\ell\ge1\) be an absolute constant such that
\[
 |\log(e^Xe^Z)-X-Z|\le c_\ell|X|\,|Z|
       \qquad (|X|,|Z|\le1/16).
\]
Indeed this holomorphic remainder vanishes on either coordinate axis;
the Schwarz estimate in each variable on the \(1/12\)-ball gives
the bound. For \(\sup|a'|<r_{\mathrm c}=r_{\mathrm{av}}/16\), the logarithm
\(\widetilde a=i^{-1}\log(e^{i\operatorname{Ad}(h)a'}e^{ia^0})\)
satisfies
\[
 \sup|\widetilde a|<3r_{\mathrm{av}}/4,\qquad
 \sup|\widetilde a-a^0|\le2c_\ell\sup|a'|.
\]
By \cref{lem:one-step-straightening}, the one-step map is jointly
holomorphic and bounded by one on the product of the reference
and perturbation balls of radius \(r_{\mathrm{av}}\). Cauchy's
inequality bounds its linear kernel by \(r_{\mathrm{av}}^{-1}\),
and its derivative in the reference by \(4r_{\mathrm{av}}^{-2}\)
on the \(3r_{\mathrm{av}}/4\)-ball. Thus the change of a kernel
entry is at most \(8c_\ell r_{\mathrm{av}}^{-2}\sup|a'|\).
There are at most \(6L^3\) entries in a row and six
rows containing a bond. Together with the real row
bound this gives
\[
 \|\mathbf Q(\mathbf W)-\mathbf S(W_0)\|_{\rm row}
 \le\kappa'_9(\varpi(W_0)+\sup|a'|),\qquad
 \|\mathbf S(W_0)\|_{\rm row}\le L,
 \quad \kappa'_9=\kappa_9+48L^3c_\ell r_{\mathrm{av}}^{-2}.
\]
The Taylor remainder on the \(r_{\mathrm{av}}/2\)-ball is at most
\(2r_{\mathrm{av}}^{-2}|a|^2\). Its derivative is bounded by
\((32/3)r_{\mathrm{av}}^{-2}|a|\): apply Cauchy's inequality on
the \(3r_{\mathrm{av}}/4\)-ball and Schwarz's inequality to the
derivative, which vanishes at zero. We may therefore
take \(\kappa'_5=11r_{\mathrm{av}}^{-2}\).
Gauge conjugations here are isometries, so these
bounds hold before choosing the gauge. They give, with
\(\epsilon_n=2\alpha L^{2(n-j)}
                 +C''_{\mathrm{ia}}L^{n-j}A\),
\[
 t_{n+1}\le t_n\{1+\kappa'_9L^{-1}\epsilon_n
                    +\kappa'_5L^{-1}L^{n-j}t_n\},
 \qquad t_n=L^{j-n}\sup|a_n| .
\]
Since
\(\sum\epsilon_n\le4L^{-2}\alpha+2C''_{\mathrm{ia}}L^{-1}A\)
and \(\sum_{n<j}L^{n-j}\le2L^{-1}\), decrease
\(c_{\mathrm{ia}}\) so the resulting product is at most
two. Induction gives \(t_n\le2B\). The one-step factor
bound \(|\lambda_n|\le\kappa_6|a_n|\) then gives the
last estimate. These arguments use only the double
block: each lower double block is nested in it.
The products stay in the strict mean and logarithm
domains. Real uniqueness and holomorphic continuation
give the factorization for complex perturbations.

The remaining real comparison is with a lifted top
average. Put
\[
 A_L=\tfrac34(L-1)^2,\qquad
 C^{\mathrm{AC}}_2=8(L-1)A_L+
          (1+\pi/3)\{L(L-1)+8A_L\}.
\]
If the intermediate curvatures are \(\varpi_n\) in
the mean chart, there is a gauge, equal to one at
the top sites, which makes all averaged block contour
variables equal to one and satisfies
\begin{equation}\label{I-10-131}
 |\widehat W_q(c)(Q^{s*}_{j\to q}W_{0,j})(c)^{-1}-1|
       \le C^{\mathrm{AC}}_2\sum_{n=q}^{j-1}\varpi_n .
\end{equation}
Indeed define it downwards by
\(v_q(x)=v_{q+1}(y)W_{0,q}(y,x)\) for \(x\in B(y)\).
The values agree at the centers and gauge covariance
makes every contour mean equal to one. The path exchange
estimate bounds the diameter of a contour family by
\(A_L\varpi_q\); its individual transports are then
within \(4A_L\varpi_q\) of one. Appending a last step
shows that each internal bond is within
\(8A_L\varpi_q\) of one. For a crossing bond, moving
the straight \(L\)-step path across an axial contour
uses at most \(L(L-1)\) plaquettes. Comparing its mean
with the straight transport costs at most
\((1+\pi/3)\{L(L-1)+8A_L\}\varpi_q\); the other
\(L-1\) internal bonds cost \(8(L-1)A_L\varpi_q\).
Thus a crossing bond differs from the next average by
at most \(C^{\mathrm{AC}}_2\varpi_q\). Descending
induction, using unitarity and the composition of lifts,
proves \eqref{I-10-131}.

For a tower row \(c\) of level \(q\) and part \(s\),
let \(l=l_s>q\). With \(x^{(i)}\) the successive
ancestors, define the descent and source by
\begin{equation}\label{II-71-24}
 \begin{gathered}
 d_{q,l}(x)=
 \mathbf U_{(l-1)}(x^{(l)},x^{(l-1)})\cdots
                \mathbf U_{(q)}(x^{(q+1)},x),\\
 \mathsf{src}_{\mathsf S}[\mathbf p](c)
   =\frac1i\log X_c[\mathbf U],\qquad
 X_c=d_{q,l}(c_-)^{-1}
       (Q^{s*}_{l\to q}\mathbf U_{(l)})(c)
       d_{q,l}(c_+)\mathbf U_{(q)}(c)^{-1}.
 \end{gathered}
\end{equation}
Use this value when both endpoint blocks of order \(l\)
lie in \(\bigcup\mathsf S\), and put zero at all
other rows. A tower endpoint lies in an \(M_1\)-cube
of the wide collar, entirely in layer \(l\) of the
given sequence. Its order-\(l\) block and its \(M_1\)-cube
lie in the same big cube. Consequently the condition
just imposed is equivalent to both row endpoints
belonging to \(\bigcup\mathsf S\).

Set \(\Sigma_{\mathrm{ia}}=(2\kappa_6+1)C''_{\mathrm{ia}}\)
and reduce \(c_{\mathrm{ia}}\) so
\(\Sigma_{\mathrm{ia}}c_{\mathrm{ia}}<3/64\).
For the real field, \eqref{I-10-131} and
\(\sum_{n<l}2\alpha'_0L^{2(n-l)}\le\alpha'_0\)
give \(|X_c[U]-1|\le C^{\mathrm{AC}}_2\alpha'_0\).
Inserting the complex perturbation in both descents,
the top average and the inverse lower average gives
total exponent norm at most
\(\Sigma_{\mathrm{ia}}L^l\sup|a'|
 <\Sigma_{\mathrm{ia}}\alpha'_1\).
Thus
\[
 |X_c[\mathbf U]-1|
 \le C^{\mathrm{AC}}_2\alpha'_0
             +2\Sigma_{\mathrm{ia}}\alpha'_1<1/2.
\]
The logarithm yields, after taking the infimum of the
admissible radii,
\begin{equation}\label{II-155-35}
 |\mathsf{src}_{\mathsf S}[\mathbf p]|_\infty
   \le C_{\rm src}\mathsf{bs}_{\mathsf S}(\mathbf p),
 \qquad C_{\rm src}=2\max(C^{\mathrm{AC}}_2,2\Sigma_{\mathrm{ia}}).
\end{equation}
This also proves holomorphy and locality on the two
endpoint blocks. Gauge factors cancel through the descent,
leaving conjugation by \(u(c_-)\); conjugation of the
complex matrices commutes with every mean and logarithm.
These give covariance and reality of the source.

For a level-\(k\) bond \(b\), let \(y_\pm\) be the
centers of its two containing level-\(t\) blocks. Its
reading block is \(B^t(y_-)\cup B^t(y_+)\).
Put \(\mathsf{lv}(x)=t\), \(\mathsf y(x)=y(x)\)
for \(x\in\Omega_t\), and
\(\mathsf{lv}(x)=k\), \(\mathsf y(x)=x\) otherwise.
Let \(\operatorname{ax}_{\mathbf U}(x)\) be the
level-\(k\) averaged contour variable from \(y(x)\)
to \(x\) inside \(\Omega_t\), and one outside.
Superscript \({\rm ax}\) denotes dressing by these
factors. For \(\mathbf U'=e^{i\eta_k\mathbf A}\mathbf U\),
let \(v_\pm\) be the accumulated factors of the relative
recursion at \((\mathsf{lv}(b_\pm),\mathsf y(b_\pm))\).
Define
\begin{equation}\label{II-71-26}
 \mathsf{rdr}_b(\mathbf U;\mathbf A)
 =\operatorname{Ad}(\operatorname{ax}_{\mathbf U}(b_-))^{-1}
  \frac1i\log\left[
   (\mathcal M^k\mathbf U)^{\rm ax}(b)
   \{v_-^{-1}(\mathcal M^k\mathbf U')^{\rm ax}(b)v_+\}^{-1}
             \right].
\end{equation}
All constituents use only the reading block.
At a real configuration its real averages coincide
with the complex averages on the domains just proved.
The accumulated factors always use the complex recursion.

For \(j=t\), write \(A=L^t\sup|a'|\) and
\(B=L^t\sup|a_0|\), \(a_0=\eta_k\mathbf A\), and suppose
\(A,B<c_{\mathrm{ia}}\). The product logarithm gives
\[
 e^{ia_0}e^{ia'}U=e^{i\widetilde a}U,\qquad
 L^t\sup|\widetilde a|\le2(A+B)<4c_{\mathrm{ia}}\le c_\sharp.
\]
Thus the plain-average and contour bounds in
\eqref{II-155-25} apply also to this combined perturbation;
the relative recursion uses the smaller radius
\(c_{\mathrm{ia}}\). They show that the argument of the last
logarithm is a product of unitary factors whose product is one,
with inserted exponentials of total norm at most
\[
 (2\kappa_6+1)C''_{\mathrm{ia}}L^{-1}\{A+2(A+B)\}
                       +\frac{4\kappa_6}{L-1}B
 \le\Sigma_{\mathrm{ia}}(A+B)<3/32 .
\]
It is therefore within \(3/16\) of one.
The outer axial conjugation has norm at most \(e^{1/8}\),
so the read-out has norm less than
\(3e^{1/8}/13<1\). It is holomorphic in \(a_0\)
on \(B<c_{\mathrm{ia}}\) and zero at \(a_0=0\).
Schwarz's inequality gives
\begin{equation}\label{II-155-36}
 |\mathsf{rdr}_b(\mathbf U;\mathbf A)|
       \le L^t\sup|\eta_k\mathbf A|/c_{\mathrm{ia}}
       =C_{\rm rdr}\sup|\mathbf A|,
 \qquad C_{\rm rdr}=L/c_{\mathrm{ia}} .
\end{equation}
The covariance formula for the dressed averages and
the accumulated factors gives
\[
 \mathsf{rdr}_b(\mathbf U^u;\operatorname{Ad}(u)\mathbf A)
 =\operatorname{Ad}(u(b_-))\mathsf{rdr}_b(\mathbf U;\mathbf A).
\]
The same constituent formulas give compatibility with
complex conjugation.

\subsection{Local terms of the shift}

Set
\[
 C_{\rm stp}=\max\{C_{\rm rdr}C_{\rm src},8L^3\},
 \qquad
 c_{\rm stp}=\min\{L^{-2}c_{\mathrm{ia}},(4C^{\mathrm{AC}}_2)^{-1}\}.
\]
In addition to the response conditions require
\begin{equation}\label{II-155-21}
 \varrho_{\mathbf H}\le c_{\rm stp},\qquad
 C_{\rm stp}\varrho_{\mathbf H}\le\varrho_{\rm rw},\qquad
 C_{\rm stp}C_{\mathrm D,\mathrm s}e^{8\bar\kappa}
                                      \varrho_{\rm rw}\le1,
 \qquad \bar\kappa\ge4,\quad\check R_t\ge19 .
\end{equation}
For a zone bond \(b\), let \(\mathsf P\langle b\rangle\)
be the one or two big cubes containing its endpoints.
Their union contains its reading block. For
\(\mathsf S\supset\mathsf P\langle b\rangle\), put
\begin{equation}\label{II-155-39}
 \Phi_{b,\mathsf S}(s;\mathbf p)
 =\mathsf{rdr}_b\!\left(\mathbf U;
    \mathbf H^{\mathfrak o}_{\mathsf S}
       (s;\mathbf p^{\rm tc}_{\mathsf S};
                      \mathsf{src}_{\mathsf S}[\mathbf p])\right).
\end{equation}
Here \(\mathbf p^{\rm tc}_{\mathsf S}\) is the
free-boundary truncation in \eqref{II-70-123}.
The source estimate and \eqref{II-155-21} put its
datum in the response ball; the last inequality
puts the response in the read-out ball.
Consequently \(\Phi\) is holomorphic for
\(|s_\Delta|<e^{\bar\kappa}\), on the entire class
\(\mathfrak P^{\mathfrak o}_{\mathbf H}(\mathsf S)\), and
\begin{equation}\label{II-155-40}
 |\Phi_{b,\mathsf S}(s;\mathbf p)|
 \le C_{\rm stp}C_{\mathrm D,\mathrm s}
           e^{8\bar\kappa}\mathsf{bs}_{\mathsf S}(\mathbf p).
\end{equation}
This is a composition of holomorphic maps in local
matrix charts; in particular, no holomorphic choice
of the real reference in a decomposition is required.

Use ordinary finite differences, keeping the anchor
\(\mathsf P\langle b\rangle\) at one:
\begin{equation}\label{II-71-34}
 \mathbf H(b;\mathsf S;\mathbf p)
 =\sum_{\mathcal K\subset\mathsf S\setminus\mathsf P\langle b\rangle}
   (-1)^{\#(\mathsf S\setminus\mathsf P\langle b\rangle)-\#\mathcal K}
   \Phi_{b,\mathsf S}
        (\mathbf1_{\mathcal K\cup\mathsf P\langle b\rangle};\mathbf p).
\end{equation}
Set this expression to zero if the anchor is missing.

\begin{proposition}\label{II-155-54}
Under the response floors and radii above, the functions
\(\mathbf H(b;\mathsf S;\cdot)\) are holomorphic on
\(\mathfrak P^{\mathfrak o}_{\mathbf H}(\mathsf S)\),
commute with conjugation, and are gauge covariant at
\(b_-\). They vanish unless \(\mathsf S\) is connected,
contains \(\mathsf P\langle b\rangle\), and meets the
boundary cubes of the localization component. For
\(\mathsf S\in\mathcal D_t\),
\begin{equation}\label{II-155-55}
 |\mathbf H(b;\mathsf S;\mathbf p)|
 \le C_{\rm stp}C_{\mathrm D,\mathrm s}
        \mathsf{bs}_{\mathsf S}(\mathbf p)
             e^{-\bar\kappa(d_t(\mathsf S)+1)/4}.
\end{equation}
Thus their amplitude and decay rate may be taken to be
\begin{equation}\label{II-155-56}
 \mathsf{am}_{\mathbf H}
       =C_{\rm stp}C_{\mathrm D,\mathrm s}\varrho_{\mathbf H},
 \qquad \mu_{\mathbf H}=\bar\kappa/4 .
\end{equation}
If \(\bar\kappa/4\ge c_{\mathrm E}=\log18432\), then,
with \(C_{\mathrm E}=4096\) and the supremum of an
empty set interpreted as zero,
\begin{equation}\label{II-155-59}
 \sum_{\mathsf S\subset\mathsf P^{\rm loc}_{\mathfrak o}}
 \sup_{\substack{\mathbf p\in
             \mathfrak P^{\mathfrak o}_{\mathbf H}(\mathsf S)\\
             \mathsf{bs}_{\mathsf S}(\mathbf p)\le\varrho_{\rm sz}}}
 |\mathbf H(b;\mathsf S;\mathbf p)|
 \le C_{\rm stp}C_{\mathrm D,\mathrm s}C_{\mathrm E}
      \varrho_{\rm sz}
       e^{-(\bar\kappa/4-c_{\mathrm E})\check R_t}.
\end{equation}
\end{proposition}

\begin{proof}
Restriction of the pair, the truncated current and
the tower source commutes with restriction of the
cube parameters. The component property of the
response then says that the read-out uses only the
component containing the reading block. If
\(\mathsf S\) has another component, pairing the
subsets in \eqref{II-71-34} with and without one
cube in that component gives zero. If
\(\mathsf S\) misses the boundary cubes, every
source row is common to the two sequences; its
source is zero. The response and the read-out
both vanish at zero source. This proves the support
assertions. The same constituent properties prove
covariance and conjugation.

Every nonzero term joins a reading cube to a
boundary cube. The zone collar gives a distance
at least \(LM(\check R_t-1)\) between them. Counting
big-cube steps yields
\[
 N:=\#\mathsf S\ge d_t(\mathsf S)+1\ge\check R_t.
\]
The Cauchy integral for the ordinary finite
difference in \eqref{II-71-34} gives, since the
anchor has at most two cubes,
\[
 |\mathbf H|
 \le C_{\rm stp}C_{\mathrm D,\mathrm s}
 \mathsf{bs}_{\mathsf S}(\mathbf p)\,
       e^{8\bar\kappa}(e^{\bar\kappa}-1)^{-(N-2)}.
\]
Put \(\epsilon_0=-\log(1-e^{-4})<1/50\).
As \(N\ge19\), \(8\le(N-3)/2\), and hence
\[
 8\bar\kappa-(N-2)(\bar\kappa-\epsilon_0)
 \le-\tfrac14\bar\kappa N.
\]
This proves \eqref{II-155-55}. The connected-set
count with the prescribed anchor, followed by
the geometric sum starting at \(d_t+1=\check R_t\),
gives \eqref{II-155-59} with the displayed entropy
constants.
\end{proof}

The finite-difference identity also gives the exact
total before any estimate. If \(\mathbf p\) is given
on the whole component, locality identifies every
value in a smaller cube set with the corresponding
value on the whole component. Möbius inversion gives
\begin{equation}\label{II-155-51}
 \sum_{\mathsf S\subset\mathsf P^{\rm loc}_{\mathfrak o}}
 \mathbf H(b;\mathsf S;
                \mathbf p|_{\mathsf S\cap\mathfrak o})
 =
 \mathsf{rdr}_b\!\left(\mathbf U;
  \mathbf H^{\mathfrak o}_{\mathsf P^{\rm loc}_{\mathfrak o}}
       (\mathbf1;\mathbf p^{\rm tc};
                             \mathsf{src}[\mathbf p])\right).
\end{equation}
All sums are finite. Terms with missing anchor
are zero, and the anchored subset identity applies
to the remaining terms.

\subsection{The actual background in the local class}

We verify the local gauges needed at the physical
pair. This uses the layerwise estimate
\eqref{eq:history-layer-radius}, not a bound by the
largest threshold without its scale factors.
Let \(U\) be an element of a completed collared
background of depth \(n\), with datum regular at
\((2\varepsilon_i)_{i\le n}\). Assume the running
bounds under which \cref{II-10-20} holds, and fix
\(1\le t\le n\), \(k=t-1\). If
\(\lambda(\mathfrak e)\) is the largest layer
index among the corners of a regularity element,
then
\[
 |\partial U(p)-1|=L^{-2\lambda(p)}r_U(p),\qquad
 |J^{(u)}(U)(b)|=L^{3(u-\lambda(b))}r_U(b).
\]
The running inequalities give
\(\varepsilon_{i+1}\le L^2\varepsilon_i\);
thus \(\varepsilon_iL^{-2i}\) and
\(\varepsilon_iL^{-3i}\) decrease.
At a corner of index at least \(i\), it follows that
\begin{equation}\label{II-65-33}
 |\partial U(p)-1|<C_{\rm reg}\varepsilon_iL^{-2i},
 \qquad
 |J^{(u)}(U)(b)|<C_{\rm reg}\varepsilon_iL^{3(u-i)}.
\end{equation}
Put the current equal to zero on bonds that do
not have four plaquettes in \(\Omega^\diamond_0\);
the same bound holds there.

The matrix divergence required for the gauge
theorem also has the correct layerwise bound.
For a full bond of index \(\lambda\), each adjacent
plaquette has radius less than both
\(C_{\rm reg}\varepsilon_\lambda\) and
\(\varepsilon_{\mathrm v}<1\). Its squared curvature
is therefore smaller than
\(L^{-4\lambda}C_{\rm reg}\varepsilon_\lambda\).
The scalar/vector decomposition in
\eqref{eq:background-matrix-divergence} gives, in
level-\(n\) lengths,
\begin{equation}\label{II-65-34}
 |D_U^{\mathsf T}\partial U(b)|
 <2C_{\rm reg}\varepsilon_\lambda
                   \eta_n^2\mathsf L(b)^{-3}.
\end{equation}

For a bare gauge \(U^v=e^{ia}\), measure its size
of index \(i\) by
\[
 \max\{L^i|a|,\ L^{2i}|a_\nu(x^{+\mu})-a_\nu(x)|\}.
\]
Writing \(A=\eta_k^{-1}a\) identifies these with
\(\xi_i|A|\) and \(\xi_i^2|\nabla^\xi A|\) in
level-\(k\) lengths. Increasing the index by one
costs at most \(L^2\).

\begin{lemma}\label{II-65-43}
Put
\[
 \mathsf c_{\mathrm{lgs}}
       =8B_*L^3c_*^{-1}\bar C_{\rm reg},
 \qquad \ell_{\mathrm c}=M/L^2.
\]
Assume \(c_*\le1\), the geometric floors
\[
 R_1\mathfrak M\ge L^3,\qquad
 M\ge L^4R_1\mathfrak M,\qquad
 M\ge L^5\mathfrak M,
\]
the history period conditions, and
\(\mathsf c_{\mathrm{lgs}}M\varepsilon_t\le1\).
For \(0\le i\le t\), every compact
\(\mathfrak M\xi_i\)-aligned cube
\(q\subset\Omega^\diamond_i\) of side at most
\(12\mathfrak M\xi_i\) has a real gauge of index
\(i\) at most \(\mathsf c_{\mathrm{lgs}}M\varepsilon_t\).
Every enlarged tiny cube \(\mathfrak y^\sim\)
meeting the bonds of the step has such a gauge
of index \(t\). These gauges restrict to any
region containing the cube.
\end{lemma}

\begin{proof}
First let \(i\ge1\), and write the side of \(q\)
as \(l\xi_i\), where \(1\le l\le12\mathfrak M\).
Its concentric enlargement of side \((l+4)\xi_i\),
together with its \(2\eta_k\)-neighborhood, lies in
\(\Omega^\diamond_{i-1}\): its distance from \(q\)
is at most \((2L+2)\xi_{i-1}\), whereas the gap
is larger than \(L^3\xi_{i-1}\).
Its side is less than half the period, since
\(l+4<M\le L^{m+K-i}\).
Extend \(U\) by the identity outside
\(\Omega^\diamond_0\). No plaquette or divergence
used in this enlargement changes.

In level-\(i\) units the curvature and divergence
parameters are bounded by
\[
 \max\!\left\{2L^3C_{\rm reg}\varepsilon_{i-1},
  \max_{q'\ge i}2L^{-2(q'-i)}
                            C_{\rm reg}\varepsilon_{q'}\right\}
 =2L^3C_{\rm reg}\varepsilon_{i-1}.
\]
Indeed all relevant indices are at least \(i-1\),
and \(\varepsilon_{q'+1}\le L^2\varepsilon_{q'}\).
The smallness needed by \cref{I-28-75} follows from
\[
 2L^3C_{\rm reg}l\varepsilon_{i-1}
 \le c_*\mathsf c_{\mathrm{lgs}}M\varepsilon_t\le c_*.
\]
That theorem gives a gauge of size
\(2B_*L^3C_{\rm reg}l\varepsilon_{i-1}\),
which is at most
\(\mathsf c_{\mathrm{lgs}}M\varepsilon_t\).

For \(i=0\), write \(N\) for the side of the cube in bare units.
Its aligned faces have half-integral coordinates, so it contains
\(N\le12\mathfrak M\) bare sites in each direction. The comb gauge of
\cref{I-10-51} gives
\(|a|\le\pi(N-1)C_{\rm reg}\varepsilon_0\);
its first differences are bounded by twice this
number. Since \(N\le12\mathfrak M\) and
\((N-1)C_{\rm reg}\varepsilon_0<1\), this gives
the same required bound.

An enlarged tiny cube has side \(3\ell_{\mathrm c}\)
in level-\(k\) lengths. Each of its points lies
within \(2\ell_{\mathrm c}+1<LM\) of \(\Omega_t\).
The history gaps put it in \(\Lambda_k\subset\Omega_k\)
when \(k\ge1\), and in \(\Omega^\diamond_0\)
when \(k=0\). It is a compact cube, since
\(3\ell_{\mathrm c}+4<M\).
For \(k\ge1\), the preceding gauge argument at
index \(k\), followed by the factor \(L^2\)
for index \(t\), bounds its size by
\[
 L^2\,2B_*L^3C_{\rm reg}(3\ell_{\mathrm c})
                \varepsilon_{k-1}
 =6B_*L^3C_{\rm reg}M\varepsilon_{k-1}
 \le\mathsf c_{\mathrm{lgs}}M\varepsilon_t.
\]
At \(k=0\) the comb bound at index one is at most
\(6\pi C_{\rm reg}M\varepsilon_0\), with the same
conclusion. The stated floors and the definition
of \(\mathsf c_{\mathrm{lgs}}\) give all inequalities
used above. Restriction does not alter a gauge on
its cube.
\end{proof}

Apply this at the last step \(t=n\), and assume
\begin{equation}\label{II-65-50}
 \mathsf c_{\mathrm{lgs}}M\varepsilon_t
             \le\min\{\alpha_0,\alpha_2,\varrho_{\mathbf H}\}.
\end{equation}
For any \(\mathsf S\subset\mathsf P^{\rm loc}_{\mathfrak o}\),
use the pair restricted to
\(W=\mathsf S\cap\mathfrak o\), with decomposition
\((U|_W,0)\). A tested cube in \(W\) is contained in
\(\Omega^\diamond_i\) and has side at most
\(12\mathfrak M\xi_i\), so the containing cube
in the pair definition may be the cube itself.
The gauge bound is at most \(\alpha_0\).
Equation \eqref{II-65-33} proves the curvature
bound with \(\alpha'_0=C_{\rm reg}\varepsilon_t\)
and the current bound with \(\alpha_2\).
The perturbation bound holds for every
\(\alpha'_1>0\), because the perturbation is zero.
Since \(C_{\rm reg}<\mathsf c_{\mathrm{lgs}}M\),
one can take
\(0<\alpha'_1\le
\min(\alpha_1,\varrho_{\mathbf H}-C_{\rm reg}\varepsilon_t)\).
Letting \(\alpha'_1\) decrease to zero proves
\begin{equation}\label{II-65-51}
 \mathbf p_U^{\mathfrak o}|_{\mathsf S\cap\mathfrak o}
 \in\mathfrak P^{\mathfrak o}_{\mathbf H}(\mathsf S),
 \qquad
 \mathsf{bs}_{\mathsf S}
   (\mathbf p_U^{\mathfrak o}|_{\mathsf S\cap\mathfrak o})
                      \le C_{\rm reg}\varepsilon_t .
\end{equation}
All decompositions and gauges here are on the actual
restricted field.

\begin{corollary}\label{II-65-44}
Let the completed background have depth \(n\ge t\),
and retain the preceding running and geometric
hypotheses. If \(m'<m+K-t\) and
\begin{equation}\label{II-65-45}
 \mathsf c_{\mathrm{lgs}}M\varepsilon_t
                   \le\tfrac12\min\{\alpha_0,\alpha_2\},
\end{equation}
then, on every nonempty region \(X\) for \(\pi_t\),
the actual pair of step \(t\), with its decomposition
\((U|_X,0)\), belongs to
\begin{equation}\label{II-65-46}
 \mathscr O^+_{h|_t}
          (X;\alpha_0/2,\alpha_1,\alpha_2/2).
\end{equation}
\end{corollary}

\begin{proof}
The step-\(t\) index at a site is the minimum of
its depth-\(n\) index and \(t\).
Equation \eqref{II-65-33}, together with the
monotonicity of the thresholds, therefore gives
the curvature and current conditions with
\(\alpha_0/2,\alpha_2/2\) on
\(X\cap\Omega^\diamond_0\).
The complex part is zero. The containing-box
gauges of \cref{II-65-43} give the natural
cube condition with radius \(\alpha_0/2\);
its tiny-cube gauges give the same radius
near every step bond. These are precisely the
additional gauges in the definition of
\(\mathscr O^+\). One decomposition satisfies
all conditions.
\end{proof}

\subsection{Identification with the background difference}

Let \(\mathsf v\in\mathsf F_t^{\rm reg}\), and choose
an element \(U\) of its completed collared background
\(U^\diamond_t(\mathbf V^{\diamond+})\).
On a localization component \(\mathfrak o\), take
the pair \(\mathbf p_U^{\mathfrak o}\) consisting
of \(U\) and its given-sequence current.
Assume \eqref{II-65-50}; by \eqref{II-65-51} this
pair is in the required local class. Use the
running and smallness hypotheses of the joined
background construction, including
\(\varepsilon_t\le\epsilon_{\mathrm H}\), and require
also
\begin{equation}\label{eq:strip-criticality-ceiling}
 C_{\rm stp}C_{\rm rs}\varrho_{\rm rw}
       \le\varepsilon_{\mathrm v},\qquad
 2C_{\rm reg}\varepsilon_t\le\varepsilon_{\mathrm v}.
\end{equation}
We prove
\begin{equation}\label{II-155-61}
 \sum_{\mathsf S\subset\mathsf P^{\rm loc}_{\mathfrak o}}
 \mathbf H(b;\mathsf S;
             \mathbf p_U^{\mathfrak o}|_{\mathsf S\cap\mathfrak o})
 =\operatorname{Ad}(\mathsf{ax}_U(b_-))^{-1}
                            \mathbf H^{(k)\diamond}(b).
\end{equation}
Here the right side is the difference of the two
actual axial backgrounds used in the strip map.

Write \(U_0=U|_{\mathfrak o}\).
The restriction of the completed background is a
critical configuration for the joined sequence.
The free-boundary rule \eqref{II-70-163} makes the
truncated pair on the whole component
\((U_0,J_{\mathfrak o}U_0)\). Put
\[
 B'=\mathsf{src}[\mathbf p_U^{\mathfrak o}],\qquad
 T=\mathbf H^{\mathfrak o}_{\mathsf P^{\rm loc}_{\mathfrak o}}
        (\mathbf1;(U_0,J_{\mathfrak o}U_0);B'),\qquad
 U''=e^{i\eta_kT}U_0 .
\]
The source is real and smaller than
\(\varrho_{\rm rw}\); \cref{I-91-94} applies.
It makes \(U''\) critical for its new actual datum.
Moreover,
\begin{equation}\label{eq:strip-small-critical-change}
 \|T\|,\quad|F'_{\mathfrak o,U_0}(T)-J_{\mathfrak o}U_0|^\sharp
 \le C_{\rm rs}|B'|
 \le\varepsilon_{\mathrm v}/C_{\rm stp}.
\end{equation}

We identify the datum before invoking uniqueness.
At a joined node \(x\) of level \(q\), define
\[
 w(x)=\mathsf{dsc}_{q,n^\diamond(x)}U_0(x),
\]
where \(n^\diamond(x)\) is the index of the given
completed sequence. The tower definition gives
\(e^{iB'}V_0=w^{-1}V^{\rm loc}w\) on the joined rows.
For a common row of level \(n\), an inner endpoint
has both indices equal to \(n\), so its descent is
the identity. An outer endpoint has joined index
\(n-1\) and given index in \(\{n-1,n\}\).
In the first case the descent has no step; in the
second it starts at the center of its own block
and is again the identity.

The outward twists in the straightened constraint
are trivial. To see this at an outer endpoint \(z\)
of level \(n\), its level-\(n\) block has either no
tower row or all its lower rows in the same tower.
In the first case the source there is zero. In the
second, let the tower reach level \(l\), and put
\(g_1=\mathsf{dsc}_{n-1,l}\). On the lower block the
source gives
\[
 e^{i\beta(c)}W_0(c)=g_1(c_-)^{-1}g_1(c_+).
\]
The lower lift inside this block is the identity.
Each perturbed contour from \(z\) to \(x\) is
\(g_1(z)^{-1}g_1(x)\). But
\[
 g_1(x)=\mathsf{dsc}_{n,l}(z)W_0(z,x),\qquad
 g_1(z)=\mathsf{dsc}_{n,l}(z);
\]
the contour therefore equals the original one.
Its relative mean is zero, and its outward twist
is one. This verifies the seam condition for all
joined rows.

Let \(r\) be the node factor supplied by
\cref{I-91-94}. Set \(g=wr^{-1}\) at joined nodes
and \(g=1\) at other bare sites. Then
\(U'''=(U'')^g\) has datum \(V^{\rm loc}\) and
is critical for it. We next check its regularity,
which is not implied by criticality alone.

The completed background has normalized curvature
and current at most \(C_{\rm reg}\varepsilon_t
\le\varepsilon_{\mathrm v}/2\).
Adjacent joined indices differ by at most one.
For a plaquette of joined index \(n\), the expansion
of its four perturbed bonds gives
\[
 \begin{split}
 |\partial U''(p)-1|
 &\le|\partial U_0(p)-1|
       +\eta_k^2\bigl(2|\nabla_{U_0}T|+2|T|^2\bigr)\\
 &\le|\partial U_0(p)-1|+4L^{-2n}\|T\|.
 \end{split}
\]
Its maximum corner index is at most \(n+1\).
After normalization the added term is at most
\(4L^2\varepsilon_{\mathrm v}/C_{\rm stp}\).
Since \(C_{\rm stp}\ge8L^3\), the plaquette remains
within the \(\varepsilon_{\mathrm v}\)-bound.

For the current, differentiating the left
exponential gives the exact identity
\[
 F'_{\mathfrak o,U_0}(T)(b)
       =f(Z)J_{\mathfrak o}(U'')(b),\qquad
 f(z)=\frac{e^z-1}{z},\qquad
 Z=-i\eta_k\operatorname{ad}T(b).
\]
By \eqref{eq:strip-small-critical-change},
\(|Z|\le2/C_{\rm stp}\le1/(4L^3)\).
Thus \(|f(Z)-1|\le1/100\) and
\(|f(Z)^{-1}|\le100/99\). The conversion from a
bond's index to a corner's index costs at most
\(L^3\), so its normalized current is bounded by
\[
 \frac{100}{99}
   \left(C_{\rm reg}\varepsilon_t
          +L^3 C_{\rm rs}\varrho_{\rm rw}\right)
 \le\frac{100}{99}
        \left(\frac12+\frac18\right)\varepsilon_{\mathrm v}
 <\varepsilon_{\mathrm v}.
\]
These bounds are gauge invariant. The uniqueness
of the regular critical orbit therefore identifies
\(U'''\) with the joined background.

Finally, at the two nodes used by the read-out, the
given and joined indices agree in the core. Hence
\(w=1\), the node factors \(r\) are exactly \(v_\pm\)
in \eqref{II-71-26}, and \(g=v_\pm^{-1}\).
Substituting the identified joined background in
that equation gives the right side of
\eqref{II-155-61}. The argument of its logarithm
lies in the series-logarithm ball proved above,
so that logarithm agrees with the intrinsic real
logarithm. Equation \eqref{II-155-51} proves
\eqref{II-155-61}. In the axial representative
\(\mathsf{ax}_U=1\), its sum is precisely
\(\mathbf H^{(k)\diamond}\).

\subsection{The sum of the strip terms}

We apply the preceding identification to a chain point
\(y=(W,B,\phi^+)\) of a bare field, with regular base
\(\mathsf d\). Use the adapted representative
\(U_{\mathsf d}\) of its completed collared background.
Let \(\mathsf U=U_{\mathsf d}\vee v_0\) be the glued
configuration, let \(\mathfrak q\) be the old stored
fields, and write \(B=(B_{\rm I},\mathsf a)\).
The new stored external field is \(\mathsf a'\).
The coordinate and support results give
\begin{equation}\label{eq:strip-sum-chain-fields}
 |\mathsf a(b)|,\ |\mathsf a'(b)|
      <\mathsf{rad}_k(b)\le C_{\rm xt}\mathfrak b_k,
 \qquad g_k|\mathsf a(b)|,\ g_k|\mathsf a'(b)|<1/8
\end{equation}
under their tier ceiling \(C_{\rm xt}\delta_k\le1/8\).
Both \((\mathfrak q,\mathsf a)\) and
\((\mathfrak q,\mathsf a')\) therefore satisfy the
stored-data tiers of \(h'\). The old tiers are the
ones of \(h\), since prolongation does not change
an earlier stored set or its radius.

Let \(P_X\) be the plain triple read from \(\mathsf U\)
at level \(t\) on \(X\), and let \(\mathbf p_T\) be
the pair of \(U_{\mathsf d}\) restricted to the
localization region. Restriction of the triple,
its pair and the old stored data gives
\begin{equation}\label{II-181-22}
 P_{X_1}|_X=P_X,\qquad
 \mathbf p_X^{\rm loc}(P_X)=\mathbf p_T|_X,\qquad
 (\mathfrak q|_{X_1})|_X=\mathfrak q|_X
 \quad(X\subset X_1).
\end{equation}
For the pair identity, the glued configuration is
\(U_{\mathsf d}\) on all inner bonds of
\(\Omega^\diamond_0\). A current is computed there
only when all four plaquettes are inner, and is
zero otherwise. Thus its reading is the same
before and after gluing; the level-\(t\) current
is multiplied by \(L^{-3}\) on passage to level
\(k\). This proves the stated pair identity with
the correct current normalization.

Equation \eqref{II-65-51} applies at every component,
including components with no zone bond. It follows
that \(\mathbf p_T\), and every restriction in
\eqref{II-181-22}, belongs to the classes of
\cref{II-155-54}. Adaptation at the last layer
means that \(\mathcal M^kU_{\mathsf d}\) is axial
on \(\Omega_t^{(t)}\). Hence \eqref{II-155-61}
has no conjugating factor, and
\begin{equation}\label{II-181-19}
 H_{\rm tot}[\mathbf p_T]=\mathfrak H_{\mathsf d},
 \qquad
 \mathsf a=\check\Psi(\mathsf a',H_{\rm tot}[\mathbf p_T]).
\end{equation}
The first identity is the actual background shift
of the base. For the second, the definition of the
stored coordinate gives \(\mathsf a=\mathsf a'\)
off \(\mathcal Z\), and on \(\mathcal Z\) it gives
\[
 g_k\mathsf a(b)=
 \mathsf{cm}\bigl(g_k\mathsf a'(b),
                    \vartheta_b\mathfrak H_{\mathsf d}(b)\bigr).
\]
These are exactly the values of \eqref{II-26-83}.
The smallness proved for the actual shift is
\(8\|\vartheta'_{\rm co}\|_\infty
|\mathfrak H_{\mathsf d}(b)|\le\delta_k\), so the
real strip map and its positive Jacobians are
defined.

The following statement keeps visible the analytic
domain conditions on the local potentials. They
are the same conditions required to form the
boundary terms in \cref{sec:boundary-assembly}.

\begin{proposition}\label{II-181-24}
Use a chain point as above and the local terms
of \cref{II-155-54}, with the rate and amplitude
conditions \eqref{II-73-14} and \eqref{II-73-40}.
Assume \(P_X\) belongs to the intermediate graded
space \(\mathscr U_X\) for each \(X\), and that
the kept, rejected and cluster families satisfy
the domain and restriction hypotheses of
\cref{sec:boundary-assembly} there.
In particular the cluster family is defined on
the stored tiers and has the indicated complex
room in the new zone variables; the two quadratic
families are defined for
\(|\mathsf a_1|<2C_{\rm xt}\mathfrak b_k\).
Assume \eqref{II-151-34} and all the actual
background and chart ceilings used above.
Then every term in the following equality is
defined, and
\begin{equation}\label{II-181-30}
 \begin{split}
 \sum_{X\in\mathcal D_t}B^{\rm sr}_{h'}
          (X;\mathsf U;\mathfrak q,\mathsf a')
 ={}&\mathsf{Zs}_{\mathsf d}(\mathsf a')\\
 &+\sum_{X\in\mathcal D_t}
 \left\{[E^{\rm cl}_{h'}+B^\times_{h'}+K^{\rm kp}_{h'}]
                    (X;\mathsf U;\mathfrak q,\mathsf a)\right.\\
 &\hspace{42mm}\left.
       -[E^{\rm cl}_{h'}+B^\times_{h'}+K^{\rm kp}_{h'}]
                    (X;\mathsf U;\mathfrak q,\mathsf a')\right\}.
 \end{split}
\end{equation}
The quadratic forms in this notation do not depend
on \(\mathfrak q\). The identity holds also at
\(k=0\), when the old stored family is empty.
\end{proposition}

\begin{proof}
Both external fields are in the tiers by
\eqref{eq:strip-sum-chain-fields}; thus the two
cluster values in \eqref{II-181-30} are defined.
The quadratic domains contain both fields.
The background and stored-data restrictions in
\eqref{II-181-22} put each evaluation at the
point \((P_X;(\mathfrak q,\mathsf a')|_X)\) used
to define the strip terms.

Extend \(\mathsf a'|_{X_1}\) by zero when forming
a term on \(X_1\). Replacing this extension by
the global field \(\mathsf a'\) changes neither
the domain condition nor the value: each shifted
bond used by that term lies in \(X_1\), and each
original function over \(X\subset X_1\) reads only
its external bonds in \(X\). The same locality
holds for the cluster function's external domain.
Restriction of \(\mathbf p_T\) therefore replaces
the local pair by one global pair in all five rows.
Restriction of the old data and the triple does
the same for their further arguments.

For the four transfer rows use the families
\(F^7_X,F^8_X,F^9_X,F^{\rm E}_X\) in
\eqref{II-151-23}. They are functions over \(X\),
including their domains, and their values at
either \(\mathsf a\) or \(\mathsf a'\) are the
actual local Lebesgue, kept, rejected and cluster
terms. The room hypotheses, \eqref{II-151-34}
and \cref{II-73-32} show that every finite-difference
vertex in their transfer terms is defined.

Here is the exact finite summation, including
its locality requirement. For fixed \(X\), put
\[
 f(\mathscr S)=
 F_X\left(\check\Psi\left(\mathsf a',
       \sum_{(b,S)\in\mathscr S}\mathbf H(b;S;\mathbf p_T)\right),w\right),
 \qquad\mathscr S\subset\mathsf{Pr}[X].
\]
The same subset gives the same value in every
target \(X_1\) containing its supports. Each
nonempty subset is assigned exactly once, to
\(X(\mathscr S)\), in \eqref{II-73-23}.
Möbius summation thus gives
\[
 \sum_{X_1}F_X^{\rm tr}(X_1)
       =f(\mathsf{Pr}[X])-f(\varnothing)
       =F_X(\mathsf a,w)-F_X(\mathsf a',w).
\]
The last equality uses \eqref{II-181-19} on the
bonds read by \(F_X\), and
\(\check\Psi(\mathsf a',0)=\mathsf a'\).
If \(\mathsf{Pr}[X]\) is empty, both differences
are zero. Summing over \(X\) proves the identities
for rows \(8,9,{\rm E}\).

For row \(7\), a bond is assigned to its one
anchor \(X=\mathsf P\langle b\rangle\), so the
same calculation gives
\[
 \sum_X R^7(X)
 =\sum_{b\in\mathcal Z^+}
          \{\mathsf{ls}_b(\mathsf a)-\mathsf{ls}_b(\mathsf a')\}
 =\mathsf{Zs}_{\rm L}(\mathsf a';\mathfrak H_{\mathsf d}).
\]
The logarithm is the one in
\eqref{eq:strip-logarithmic-density};
the bounds \(g_k|\mathsf a|,g_k|\mathsf a'|<1/8\)
keep both values on its stated branch.
For row \(6\), \eqref{II-73-43} and
\eqref{II-181-19} give
\[
 \sum_XR^6(X)
 =\sum_{b\in\mathcal Z_\theta^+}
       \log J_{\mathfrak H_{\mathsf d}(b)}
                             (e^{ig_k\mathsf a'(b)})
 =\mathsf{Zs}_{\rm J}(\mathsf a';\mathfrak H_{\mathsf d}).
\]
The Jacobians are positive, so the local series
logarithms agree with these real logarithms.
Adding the five finite sums proves
\eqref{II-181-30}, since
\(\mathsf{Zs}_{\mathsf d}=\mathsf{Zs}_{\rm J}
                         +\mathsf{Zs}_{\rm L}\).
\end{proof}

If the strip region is empty, then
\(\mathcal Z=\varnothing\): the Jacobian sum and
every transfer sum are empty. Thus the strip
terms vanish on their domains. If already
\(\Omega_1=\varnothing\), the completed and
localization regions, and every stored bond set,
are empty. Their pair and stored-data spaces
are singletons; the graded conditions on the
empty background impose no restriction on a
triple. The five rows are then defined empty
sums, while the kept, rejected and cluster
terms are the zero terms of the empty-region
convention. Both sides of \eqref{II-181-30}
are zero.



\section{Gaussian fields in a history step}
\label{sec:historical-gaussian-fields}

The integration region and its exterior now share a quadratic form.
We first extend the connected coefficients to the collared background
sequence. Local inverses then separate the integrated field from its
conditional mean; the remaining quadratic form is divided into terms
kept on strip components and terms crossing between them.

\subsection{The form and its local coefficients}
\label{sec:history-form-coefficients}

Fix a history through step $t=k+1$, with $m'<m+K-t$, and use
level-$k$ units. Thus $\xi=L^{-k}$, $\xi_n=L^{n-k}$ and
$\ell_c=M/L^2$. The collared sequence is
\begin{equation}\label{II-63-5}
 \Omega^\diamond_\bullet=(\Omega_1^\sim,\Omega_1,\ldots,\Omega_t).
\end{equation}
Its fluctuation bonds $\mathfrak X$ are the nonpivot level-$k$ bonds
whose endpoints lie in $\Omega^\diamond_0$ and at least one of whose
endpoints lies in $\Omega_t$. They split into integration bonds
$\mathfrak X_{\rm I}$, having an endpoint in $\Lambda_t$, and external
bonds $\mathfrak X_{\rm E}$. Keep the grain, clearance and covariance
floors of \eqref{eq:gaussian-localization-scales}, and require
$w\le\ell_c$. All kernels and transposes use the positive-bond counting
pairing of \cref{sec:local-gaussian-fields}.
Here $a^{\rm cv}=a_{\rm cv}$ is the previously fixed covariance
radius. Every tiny cube meeting $\mathfrak X$ has its enlargement
inside $\Omega^\diamond_0$: its points are within $2\ell_c+2$ of
$\Omega_t$, whereas the history collar has width at least $LM$.

For a coarse region $R$, a pair $p=(\mathbf U,\mathbf J)$ is given on
the inner bare bonds of $R\cap\Omega^\diamond_0$, with one decomposition
$\mathbf U=e^{i\xi A'}U$, $U$ real. Write
$\mathscr O(R;\alpha_0,\alpha_1,\alpha_2)$ for the following local
class. At every contained stencil based in layer $n$,
\begin{equation}\label{eq:history-pair-bounds}
 \begin{gathered}
  |\partial U-1|, |\partial\mathbf U-1|
       <\alpha_0\xi^2\xi_n^{-2},\qquad
  \xi_n|A'|<\alpha_1,\quad
  \xi_n^2|\nabla^UA'|<\alpha_1,\quad
  \xi_n^3|\mathbf J|<\alpha_2.
 \end{gathered}
\end{equation}
A plain gauge cube of index $n$ is an embedded cube made of
level-$n$ $\mathfrak M$-cubes, of side at most
$12\mathfrak M\xi_n$, contained in the union of layers $n,n+1$
and meeting layer $n$. Each such cube $q$ contained in
$R\cap\Omega^\diamond_0$ must lie in an embedded cube
$q'\subset R\cap\Omega^\diamond_0$ with a real gauge
$U^v=e^{i\xi A}$ satisfying
$\xi_n|A|\le\alpha_0$ and
$\xi_n^2|\nabla^\xi A|\le\alpha_0$.
Finally, every tiny cube $y$ containing an integration-bond endpoint,
with $y^\sim\subset R\cap\Omega^\diamond_0$, must admit such a gauge
on $y^\sim$ at index $t$.
The class $\mathscr O^+(R;\alpha_0,\alpha_1,\alpha_2)$ imposes this
last condition for tiny cubes containing any fluctuation-bond endpoint.
These definitions use the same $U$ throughout. Bounds on derivatives
are imposed only where their complete stencil lies in the stated domain.
Restriction preserves these classes: restrict each supplied cube gauge
to the tested cube itself. Real gauge transformations and conjugation
also preserve them, as does passage to the real reference $(U,0)$.
For $\alpha_1\le\pi$ they are open. Indeed one may fix $U$ and all
its gauges, use a local inverse of the exponential at each
$i\xi A'(b)$, and preserve the finitely many strict field inequalities.

On the shorter sequence through level $k$, use the straightened scalar
and field averages and the operators $\mathsf G_1,\mathsf H_1$ already
constructed from them. Put $j_{**}=\mathsf H_1^{\mathsf T}\mathbf J$
and $V=\mathcal M^k\mathbf U$. The one-step pivot maps at $V$ use the
ordinary one-step average. For row fields supported on
$\mathfrak X_0[\Omega^\diamond_\bullet]$, define
\begin{equation}\label{II-23-21}
 \begin{split}
 \langle B_1,\Delta_k^\sharp(p)B_2\rangle
   ={}&(B_1,(Q\mathsf G_1Q^{\mathsf T})^{-1}B_2)_{\mathbf B}
          -(B_1,\mathsf aB_2)_{\mathbf B}\\
      &-2(\mathsf D_V^{(2)}(B_1,B_2),j_{**})_{\mathbf B},
 \end{split}
\end{equation}
where $\mathsf a=1$ on level-$k$ rows. Extend the inputs by zero
before applying the one-step maps, and restrict their output back to
these rows. The fluctuation form on the nonpivot coordinates is
\begin{equation}\label{II-23-22}
 \mathsf A^\sharp[\Omega^\diamond_\bullet](p)
   =\mathsf C_V^{\mathsf T}
       (\mathsf G_V^{(2)}+\Delta_k^\sharp(p))\mathsf C_V.
\end{equation}
This is the same construction as the constant-sequence form, with its
actual row spaces and zero extensions. In particular its domain requires
the nested averages, pivot inverses and the scalar and field inverses
appearing in this formula. We verify their existence before using it.

For a grain region $Y$, take the maximal tower $\Omega_\bullet(Y)$
of \cref{sec:gaussian-connected-coefficients} and form the meet
$\Omega_n^\diamond\langle Y\rangle
=\Omega_n^\diamond\cap\Omega_n(Y)$.
Compress its fluctuation form to the bonds with $S_0(b)\subset Y$,
and extend by zero to $\mathfrak X$. Denote the resulting form by
$g_Y$ and put $g_\emptyset=0$. Define $a_S$ by
\eqref{II-29-11}, with these $g_Y$.

\begin{proposition}\label{II-31-36}
Assume $\alpha_0,\alpha_1,\alpha_2\le a^{\rm cv}$ and the grain
conditions above. If $p\in\mathscr O(R;\alpha_0,\alpha_1,\alpha_2)$,
all the meet operators defining $a_S$ exist whenever $S\subset R$.
The coefficients have the symmetry, connected support and estimate
of \cref{II-30-49}, with the same constants. In particular, writing
$2\theta_a=\delta_a$, $\epsilon_\star=(2n_{\rm w})^{-1}$,
\begin{equation}\label{eq:history-form-coefficients}
 |a_S[p](b,b')|
 \le C_a\epsilon_\star^{\#S-1}
              e^{-2\theta_a\Lambda_S(b,b')}=:\beta_S(b,b').
\end{equation}
They are holomorphic and depend only on the pair on
$\Omega^\diamond_0\cap S$ and the sequence in $S^\sim$.
Their domains have the same locality. Restriction, real gauge
transformations and conjugation commute with their construction.

If $\alpha_0\le a^{\rm cv}/2$ and
$R_d=c_{\rm d}a^{\rm cv}/(\alpha_1+\alpha_2)>1$, the same coefficients
are holomorphic on $p_z=(e^{iz\xi A'}U,z\mathbf J)$ for $|z|<R_d$
and satisfy \eqref{eq:history-form-coefficients} there. For $R_d\ge2$,
\begin{equation}\label{eq:history-form-reference-difference}
 |a_S[p]-a_S[(U,0)]|\le (2/R_d)\beta_S.
\end{equation}
On the whole torus their finite sum is the actual form
$\mathsf A^\sharp[\Omega^\diamond_\bullet](p)$.
\end{proposition}

\begin{proof}
First check admissibility on a meet. A site in its layer $n$ belongs
to an original layer $j\ge n$, so every bound in
\eqref{eq:history-pair-bounds} weakens under restriction.
For the gauge condition it suffices to enclose a candidate cube
$q\subset\Omega_n^\diamond\cap Y$ in a plain cube of the original
sequence lying in $Y$. Choose the largest $j\ge n$ with
$q\subset\Omega_j^\diamond$. If $j=n$, use $q$ itself.
Otherwise the side of $q$ is less than $\mathfrak M\xi_j$.
In each coordinate enclose its interval in two adjacent intervals
of length $\mathfrak M\xi_j$: straddle a crossed grid point, or
choose the neighboring interval inside the same interval of length
$L\mathfrak M\xi_j$ if no grid point is crossed.
Every selected interval lies in an interval of length
$L\mathfrak M\xi_j$ whose intersection with the original interval
has positive length and hence contains a bare coordinate.
Thus each level-$j$ $L\mathfrak M$-cube used in the product
contains a bare site of $q$ and lies in both
$\Omega_j^\diamond$ and $Y$; the grain partition is coarser than
this partition. The enclosing cube has side $2\mathfrak M\xi_j$.
It contains a bare site outside $\Omega_{j+1}^\diamond$ and cannot
meet $\Omega_{j+2}^\diamond$, by separation and its diameter.
It is therefore a plain cube of index $j$. Its supplied gauge,
restricted to $q$, obeys the bounds of index $n$ because
$\xi_n\le\xi_j$. This proves the natural gauge condition on every
meet, including its containing-box requirement.

The regular-pair inverse and pivot results now define
\eqref{II-23-21}--\eqref{II-23-22} on every meet, with uniform radii.
The form expansion \eqref{I-73-106} applies to these actual operators.
Its local matrices and their domains depend only on the carrier,
and the tower dependence still has range $e_{k+1}\le\rho_t$.
Consequently the proof of \cref{II-30-49} applies without changing a
constant: toggle a grain cube outside the endpoint supports and the
$e_{k+1}$-neighborhood of the tree hull; the same matrix occurs with
opposite signs. Each surviving cube lies within $r+1$ of the tree,
and the two length bounds used there absorb the finite alternating
sum and the grain count. This proves connectedness and
\eqref{eq:history-form-coefficients}. No new summation estimate is
needed.

For locality, a meet with $Y\subset S$ only uses the original
sequence at bare sites of $Y$, and the pair on
$\Omega^\diamond_0\cap Y$. If two sequences agree in $S^\sim$,
all their meets with such $Y$ agree, as do the nonpivot bond sets and
the endpoint supports. Each operator in the finite alternating sum
is then the same partial map, with the same domain. The gauge and
conjugation identities follow through the nested averages, the
bilinear transposes, inverses and pivot maps. They give real
coefficients at real pairs. Finally the disc estimate
\eqref{II-33-21} verifies the regularity bounds on every meet for
$p_z$; the fixed real gauges are unchanged. Cauchy's estimate gives
\eqref{eq:history-form-reference-difference}. Finite
inclusion--exclusion gives the whole-torus identity.
\end{proof}

The same enclosing-cube argument, using the parent cubes just
identified and the coarse region $R$ in place of $Y$, applies
to every candidate cube contained in $R\cap\Omega_n^\diamond$,
without the layer-intersection condition.
It therefore verifies the gauge condition of the normalizer class in
\cref{sec:history-normalizer}. In particular
$\mathscr O(R;\alpha_0,\alpha_1,\alpha_2)
\subset\mathscr O^N(R;a_{\rm nf})$ whenever
$a_{\rm nf}\ge\max\{\alpha_0,\alpha_1,\alpha_2\}$.

Represent $a_S$ in the orthonormal Pauli coordinates of the positive
bonds. Its footprint is $S^\sim$, and its coarse hull $X(S)$ is the
union of the $LM$-cubes meeting this footprint. In a region $R$, use
only indices with $X(S)\subset R$ and set
\begin{equation}\label{II-63-8}
 \mathsf A^{+N}_R(p)
   =\sum_{\substack{X(S)\subset R\\ S^\sim\subset N}}!a_S[p].
\end{equation}
This is a finite sum. Restricting the pair and the region retains
exactly those indices whose hull remains inside the smaller region;
their values and domains are unchanged. The zero entries outside
$S_0(b)\cup S_0(b')\subset S$ are retained in these matrix
representations. On the whole torus write $\mathsf A^+$ for the
sum; it represents \eqref{II-23-22}.

\paragraph{Splitting across components.}\label{II-68-55}
If two coarse regions $Y_1,Y_2$ have no equal or face-adjacent cubes,
their meet operators, including their domains, form a direct sum:
\begin{equation}\label{II-68-56}
 \mathsf A^\sharp_{Y_1\cup Y_2}
   =\mathsf A^\sharp_{Y_1}\oplus\mathsf A^\sharp_{Y_2}.
\end{equation}
Edge and corner contacts are allowed. To see the tower splitting,
join lifts of points in different parts by a coordinate-monotone
chain of coarse cubes. Before reaching the other part the chain
passes through a cube of $(Y_1\cup Y_2)^c$; a point of that cube is
no farther from the first point than the second point is.
Thus the distance to the complement of one's own part equals the
distance to the complement of the union at every stage of the
tower recursion. Induction gives the disjoint union of the two
towers and hence of the two meets.

No inner bond or plaquette joins the parts. Each averaging block is
in the part of its center; a top row and its double block stay in
that part by the last tower margin
$R_1\mathfrak M\ge L^2>3L/2$. Therefore the actual nested averages,
their derivatives, the gradient and its bilinear transpose, the
Hessian form and the pivot maps all split, with identical branch
domains. Their sums, products and inverses split as well; an inverse
exists precisely when both blocks are invertible. This proves
\eqref{II-68-56} for the actual form, not only for its expansion.
Any compression whose retained endpoint support is wall-connected
splits in the same way. In particular the deep integration
compression and each diagonal entry of every power in its logarithm
series are those of its own component. Empty components carry the
unique operator on the zero-dimensional space.

\subsection{A local gap and the conditional inverse}
\label{sec:history-conditional-inverse}

The tiny-cube gauge condition has a specific purpose: it transfers the
flat gap to every local compression used to integrate a history field.
The comparison must be made before inverting that compression.
Write $\gamma_{\rm tr}$ for the gap in \eqref{eq:cv-flat-gap} and put
$\gamma_{\rm pos}=\min\{1,\gamma_{\rm tr}/2\}$.
The counting norm and the coefficient bound give constants
\[
 B_{\rm w}=2\sqrt2 C_a,\qquad
 \delta=\delta_{\rm cv}/(8L),\qquad
 \varepsilon_{\rm w}=
       \frac{2(\alpha_1+\alpha_2)}{c_{\rm d}a_{\rm cv}}.
\]
For each coefficient there are fixed majorants $m_S$ and lengths
$l_S=\Lambda_S$ such that
\begin{equation}\label{eq:history-leg-majorants}
 |a_S(p)|\le m_S,\qquad
 |a_S(p)-a_S(U,0)|\le\varepsilon_{\rm w}m_S,
 \qquad \sum_S m_S(b,b')e^{\delta l_S(b,b')}\le B_{\rm w}.
\end{equation}
This is \eqref{eq:history-form-coefficients} summed by
\eqref{eq:grain-entropy-sum}, with the positive-bond coordinate norm.
A connecting tree has endpoints $b_-,b'_-$; every point of $S^\sim$
is within $3w$ of it.

We henceforth take $L^2w\le\ell_c$, enlarge a fixed lower bound on
$\ell_c$, and reduce a fixed upper bound $a_{\rm pos}$ on $\alpha_0$.
All these choices depend only on $L$. The quantitative requirements
on these choices are given in the proof below.

\begin{proposition}\label{II-63-60}\label{II-64-42}
Let $p\in\mathscr O(R;\alpha_0,\alpha_1,\alpha_2)$ and choose one
of its admitted real references $(U,0)$. For every nonempty union $W$
of tiny cubes containing integration-bond endpoints, with
$W^\sim\subset R$, and every $N$ containing the closed
$2\ell_c/3$-neighborhood of $W$,
\begin{equation}\label{eq:history-local-gap}
 \left\langle f,\sum_{\substack{X(S)\subset R\\S^\sim\subset N}}
                       a_S(U,0)f\right\rangle
       \ge\gamma_{\rm pos}|f|^2
\end{equation}
for real $f$ supported on integration bonds with an endpoint in $W$.
For $p\in\mathscr O^+(R;\alpha_0,\alpha_1,\alpha_2)$ the same
assertion holds with all fluctuation bonds in place of integration bonds.

There are constants $B_{\rm f}\ge1$, $\delta_{\rm pl}>0$ depending
only on $L$ such that, when $\alpha_1+\alpha_2$ is sufficiently small,
the local inverses constructed below have a convergent connected
expansion $t_\jmath$. Its terms are holomorphic on their actual hulls,
commute with restriction, gauge transformations and conjugation, and
have fixed majorants satisfying
\begin{equation}\label{eq:history-resolvent-majorants}
 |t_\jmath(p)|\le m_\jmath,\qquad
 |t_\jmath(p)-t_\jmath(U,0)|\le\varepsilon_*m_\jmath,
 \qquad
 \sum_\jmath m_\jmath(b,b')e^{\delta_{\rm pl}l_\jmath(b,b')}
        \le B_{\rm f}.
\end{equation}
Here $\varepsilon_*=C_{\rm res}\mathsf{cn}_{\rm w}
\varepsilon_{\rm w}$, where
$\mathsf{cn}_{\rm w}=B_{\rm w}c_{\rm lat}(\delta/8)/\gamma_{\rm pos}$;
$c_{\rm res}>0$ and $C_{\rm res}$ are absolute constants.
The same bound holds for the coefficients $a_S$ with these enlarged
constants. Every term has its input and output bonds in its footprint,
which lies within $\rho_*=4\ell_c$ of a connecting tree of length
$l_\jmath$. On the whole torus,
\begin{equation}\label{II-64-44}
 \sum_\jmath t_\jmath=(A^+_{\rm II})^{-1},\qquad
 A^+=\sum_Sa_S=\operatorname{vec}A^\sharp.
\end{equation}
All assertions include the zero-dimensional case.
\end{proposition}

\begin{proof}
Fix a tiny cube $y$ whose enlargement lies in the pair domain and for
which the appropriate tiny-cube gauge is available. In that gauge
$U=e^{i\xi A}$, with $L|A|\le\alpha_0$ and
$L^2|\nabla^\xi A|\le\alpha_0$. Set
\[
 \varsigma(x)=\min\{1,\max\{0,11/3-4\operatorname{dist}_\infty(x,y)/\ell_c\}\},
 \qquad \widetilde A(b)=\varsigma(b_-)A(b)
\]
on positive inner bonds of $y^\sim$, extend by zero, and impose
antisymmetry on reversed bonds. The cutoff is one within $2\ell_c/3$
of $y$ and zero beyond $11\ell_c/12$. Its Lipschitz constant is
$4/\ell_c$. The finite-difference product rule gives
$L|\widetilde A|\le\alpha_0$ and
$L^2|\nabla^\xi\widetilde A|\le2\alpha_0$.
This remains true at a stencil crossing the boundary: both initial
points then have distance greater than $\ell_c-2\ge11\ell_c/12$
from $y$, so both cutoff values vanish.
The real plaquette estimate proved with the action bounds gives
\[
 |\partial e^{i\xi\widetilde A}-1|
       \le6\alpha_0\xi^2L^{-2}.
\]
Thus $(e^{i\xi\widetilde A},0)$ has the flat real reference and is
admitted globally at radii $(7\alpha_0,3\alpha_0,\alpha_0)$.
Choose $7\alpha_0\le a_{\rm cv}/2$ and
$8\alpha_0\le c_{\rm d}a_{\rm cv}$.
Its complete form differs from the flat form by at most
\[
 \varepsilon_\beta B_{\rm w}c_{\rm lat}(\delta),\qquad
 \varepsilon_\beta=8\alpha_0/(c_{\rm d}a_{\rm cv}).
\]
On coefficients with $S^\sim$ in the closed $2\ell_c/3$-neighborhood
of $y$, this form is the original one conjugated by the real orthogonal
matrix $\operatorname{Ad}v(b_-)$. This follows from exact coefficient
locality, including its domain, not from a bound on the total form.
The flat form is bounded below by $\gamma_{\rm tr}$:
\eqref{eq:cv-flat-gap} applies to the collared sequence, whose fluctuation
bonds have the required $2L$ margin. Its proof uses the same actual
straightened average and pivot coordinates as \eqref{II-23-22}.

Here are the two losses in passing from this comparison to a partial
sum. Put $\widehat\delta=\min(1,\delta)$ and
$\lambda=\ell_c/3-3w-1/2$. If a footprint is not contained in that
neighborhood and the output midpoint is within $\ell_c/3$ of $y$,
its connecting length is at least $\lambda$ and at least $|b-b'|$.
The Schur bound therefore gives a discarded-tail norm at most
$E_1=B_{\rm w}c_{\rm lat}(\delta/2)e^{-\delta\lambda/2}$.
For localization use
$h_y(x)=\max(0,1-3\operatorname{dist}_\infty(x,y)/\ell_c)$ and
$\chi_y=h_y/(\sum_z h_z^2)^{1/2}$ on the relevant field bonds.
The denominator is at least one. At most sixteen cutoffs occur at
a pair of bonds, and their Lipschitz bound gives
\[
 \sum_y|\chi_y(b)-\chi_y(b')|^2
       \le144\ell_c^{-2}(|b-b'|+1/2)^2.
\]
Apply the exact localization identity from
\cref{sec:gaussian-local-resolvents}. Summing the remaining
exponential bounds its error by
\[
 E_2=156B_{\rm w}c_{\rm lat}(\delta/4)
                       (\widehat\delta\ell_c)^{-2}.
\]
Localize a field on $W$ with the cutoffs indexed by tiny cubes in $W$.
Every near footprint is among the retained terms; compare its localized
form with the global gauged form and discard the two tails.
The resulting lower bound is
\[
 \gamma_{\rm tr}-\varepsilon_\beta B_{\rm w}c_{\rm lat}(\delta)
                         -2E_1-E_2.
\]
Choose $a_{\rm pos}$ and the tiny-cube floor so that the last three
terms sum to at most $\gamma_{\rm tr}/2$.
Since $3w\le\ell_c/3$ and $L^2w\le\ell_c$, these choices are
uniform in the history and volume. This proves
\eqref{eq:history-local-gap}, in its two stated domains.

We give the actual inverse expansion, retaining its localization.
For each tiny cube $y$ containing an integration-bond endpoint put
\[
 \zeta_y(x)=\min\{1,\max\{0,3-3\operatorname{dist}_\infty(x,y)/\ell_c\}\},
 \quad N_y=\{\operatorname{dist}_\infty(\cdot,y)<2\ell_c\},\quad
 \eta_y(b)=\frac{h_y(x_b)^2}{\sum_z h_z(x_b)^2}.
\]
Here $x_b$ is the midpoint. Write
$X_y=\{b\in\mathfrak X^+_{\rm I}:\operatorname{dist}_\infty(x_b,y)<\ell_c\}$,
and use a cube in a region $R$ only when every big cube meeting $N_y$
belongs to $R$. All inverses below act on $X_y$ and are then extended
by zero. With $a_S^{\rm I}=p_{\rm I}a_Sp_{\rm I}$, set
\begin{align*}
 A^y&=p_{X_y}\sum_{S^\sim\subset N_y}a_S^{\rm I}p_{X_y},
       &G_y^\mu&=(\mu+A^y)^{-1},\\
 e_{y,S}&=
 \begin{cases}[a_S^{\rm I},\zeta_y]p_{X_y},&S^\sim\subset N_y,\\
               a_S^{\rm I}\zeta_y,&S^\sim\not\subset N_y.
 \end{cases}
\end{align*}
The local gap applies to $A^y(U,0)$: take the tiny cubes inside
$y^\sim$ that meet integration bonds; their $2\ell_c/3$-neighborhood
lies in $N_y$, and they contain the initial points of $X_y$.
The complex difference has norm at most
$\varepsilon_{\rm w}B_{\rm w}c_{\rm lat}(\delta)$.
The same argument applies to the affine interpolation of all
coefficients between $p$ and $(U,0)$ on the disc
\[
 |z|\le R_{\rm w}=
 \frac{\gamma_{\rm pos}}
 {2\varepsilon_{\rm w}B_{\rm w}c_{\rm lat}(\delta)}.
\]
Take $\varepsilon_{\rm w}\mathsf{cn}_{\rm w}\le c_{\rm res}$,
with $c_{\rm res}=1/128$. Then $R_{\rm w}\ge64$, the interpolated
coefficients are bounded by $2m_S$, and every local form has
real-part gap $\gamma_{\rm pos}/2$.
If $\varepsilon_{\rm w}=0$, this interpolation is constant.
Exponential conjugation, exactly as in \eqref{II-37-38}, gives
\[
 |G_y^\mu(b,b')|
 \le\frac{2}{\mu+\gamma_{\rm pos}/2}e^{-\delta_{\rm loc}|b-b'|},
 \qquad
 \delta_{\rm loc}=\frac{\delta}{16\mathsf{cn}_{\rm w}}.
\]

A counted string is
$\jmath=(y_0;(y_1,S_1),\ldots,(y_n,S_n))$, with
$N_{y_{i-1}}\cap S_i^\sim\ne\emptyset\ne S_i^\sim\cap N_{y_i}$.
Define
\begin{equation}\label{eq:history-inverse-strings}
 t_\jmath^\mu=(-1)^n\zeta_{y_0}G_{y_0}^\mu\eta_{y_0}
                    \prod_{i=1}^n(e_{y_i,S_i}G_{y_i}^\mu\eta_{y_i}).
\end{equation}
Its footprint is the union of these $N_y$ and $S_i^\sim$; its hull
consists of the big cubes meeting that union. An uncounted string
has zero matrix, since a nonzero coefficient entry has both of its
bonds in $S_i^\sim$.

We record the convergence calculation, since it also controls the
cross-component terms below. Let $g_y(b,b')=e^{-\delta_{\rm loc}|b-b'|}$
on $X_y\times X_y$, and zero elsewhere. Bound $e_{y,S}$ by
$\widehat e_{y,S}=2m_S|\zeta_y(b')-\zeta_y(b)|1_{X_y}(b')$
in the near case, and by $2m_S\zeta_y(b')$ otherwise.
Use $\delta_{\rm pl}=\delta_{\rm loc}/8$.
For a nonzero leak followed by $\eta_y(b'')$, either a cutoff changes
outside distance $2\ell_c/3$, while $x_{b''}$ lies within
$\ell_c/3$, or a point of $S^\sim$ lies outside $N_y$.
In the first case
$l_S(b,b')+|b'-b''|>\ell_c/3-1$; in the second
$l_S(b,b')>2\ell_c/3-1/2$.
The latter follows by joining that point to a coefficient tree within
$3w\le\ell_c/3$. In both cases, after retaining the exponential
of the string length, the remaining factors are bounded by
\[
 e^{\delta l_S}\varsigma_0
   e^{-\delta|b-b'|/4}e^{-\delta_{\rm loc}|b'-b''|/4},
 \qquad
 \varsigma_0=e^{-\delta_{\rm loc}(\ell_c/3-1)/4}.
\]
Here a local connecting segment has length $|b_--b'_-|_1\le3|b-b'|$;
this accounts for the ratio $1/8$ in $\delta_{\rm pl}$.
The first resolvent contributes a nonnegative array $P$ with
$P(b,b')\le1$ and row sum at most
$c_{\rm lat}(\delta_{\rm loc}/4)$. Each further leak and resolvent
contributes an array $Q$ whose row sum is at most
\[
 q=\frac{8B_{\rm w}}{\gamma_{\rm pos}}\varsigma_0
       c_{\rm lat}(\delta/4)c_{\rm lat}(\delta_{\rm loc}/4).
\]
Choose the tiny-cube floor so that
$q\le1/12$ and $c_{\rm lat}(\delta_{\rm loc}/4)q\le1/4$.
Expanding intermediate bonds gives the weighted majorant
$2PQ^n$ for strings of length $n$. Its sum over all $n$ is less
than $3$. Dividing by $\mu+\gamma_{\rm pos}/2$ proves absolute
convergence, uniformly also on the interpolation disc.
The same Cauchy estimate as in \eqref{eq:history-form-reference-difference}
gives the difference bound in \eqref{eq:history-resolvent-majorants}.
One may take $C_{\rm res}=64$ and enlarge
$B_{\rm f}$ to $\max(B_{\rm w},C_{\rm res}/(c_{\rm res}\gamma_{\rm pos}))$.

Joining the intermediate segments to the coefficient trees gives a
connected tree of the string length. Every point of $N_{y_i}$ is
within $4\ell_c$ of its live intermediate bond. This proves the
stated footprint bound and endpoint support. All factors are defined
on the hull, so the product has exactly the claimed locality.
Gauge covariance and conjugation hold factor by factor.

Finally put $\Pi^\mu=\sum_y\zeta_yG_y^\mu\eta_y$ and
$E^\mu=\sum_{y,S}e_{y,S}G_y^\mu\eta_y$ on the whole torus.
Since $\zeta_y\eta_y=\eta_y$ and $\sum_y\eta_y=1$, multiplication
of the definitions gives $(\mu+A^+_{\rm II})\Pi^\mu=1+E^\mu$.
The preceding absolute majorant sums
$\Pi^\mu\sum_{n\ge0}(-E^\mu)^n$ and identifies its terms with
\eqref{eq:history-inverse-strings}. It is a right inverse and hence
the inverse in finite dimension. At $\mu=0$ this is
\eqref{II-64-44}. Empty integration space gives the unique matrix
on the zero-dimensional space and an empty string sum.
\end{proof}

\subsection{Kept forms and terms crossing the strips}
\label{sec:history-kept-forms}

Write $p_{\rm I},p_{\rm E}$ for the coordinate projections, and
$p_{\mathfrak k}$ for the external coordinates of a strip component.
There are two kinds of terms in the conditional quadratic form:
\begin{equation}\label{II-64-11}
 d_S=p_{\rm E}a_Sp_{\rm E},\qquad
 d_{(S_1,\jmath,S_2)}
      =-p_{\rm E}a_{S_1}p_{\rm I}t_\jmath p_{\rm I}a_{S_2}p_{\rm E}.
\end{equation}
Their footprint is the union of the footprints of their factors;
$X(\mathfrak z)$ is its big-cube hull. A term is kept on
$\mathfrak k$ when $\mathsf{kp}_t(X(\mathfrak z))=1$ and
$\mathfrak k=\mathfrak k_{X(\mathfrak z)}$; retain its compression
to $p_{\mathfrak k}$. Write $d^{\rm k}_{\mathfrak z}$ for this
compression, or zero when the term is not kept, and
$d^\times_{\mathfrak z}=d_{\mathfrak z}-d^{\rm k}_{\mathfrak z}$.
The geometry of the strip gives two facts used below: a kept hull
lies in $\mathfrak k^\sim$, and every footprint of diameter less
than $W_t=LM\check R_t$ containing an external bond is kept on its
unique strip component.
Define
\begin{align}
 D_{\mathfrak k}&=\sum_{\mathfrak z\ {\rm kept\ on}\ \mathfrak k}
                         p_{\mathfrak k}d_{\mathfrak z}p_{\mathfrak k},
 &\Delta^{\rm loc}_{\mathfrak k}
       &=\tfrac12(D_{\mathfrak k}+D_{\mathfrak k}^{\mathsf T}),
                                                        \label{II-64-13}\\
 \varphi_{\mathfrak k,X}(f;p)
       &=-\tfrac12\sum_{\substack{\mathfrak z\ {\rm kept\ on}\ \mathfrak k\\
                                  X(\mathfrak z)=X}}
                   \langle f,p_{\mathfrak k}d_{\mathfrak z}(p)p_{\mathfrak k}f\rangle,
                                                        \label{II-64-14}\\
 B^\times(X;p,f)&=-\tfrac12\sum_{X(\mathfrak z)=X}
                            \langle f,d^\times_{\mathfrak z}(p)f\rangle,
 &K^{\rm kp}(X;p,f)&=\sum_{\mathfrak k}\varphi_{\mathfrak k,X}(f;p).
                                                        \label{II-64-15}
\end{align}
The sums below establish their domains and convergence.

\begin{proposition}\label{II-64-65}\label{II-64-75}
Under the preceding local-gap and inverse hypotheses, put
$B_{\rm ch}=2c_{\rm lat}(\delta_{\rm pl}/4)^2B_{\rm f}^3$.
For any fixed decay exponent $\kappa>0$, take
$\delta_{\rm pl}LM\ge C_2(\kappa+2)$,
$c_2\delta_{\rm pl}LM\ge2$ and
$(LM)^6\le e^{c_2\delta_{\rm pl}LM}$, where $c_2>0,C_2,C_3$
are absolute constants. Set
$K_{\rm ch}=C_3(LM)^6e^{C_3(\kappa+1)}B_{\rm ch}$.
For $p\in\mathscr O(X;\alpha)$ and $\sup|f|\le b_e$, the series
in \eqref{II-64-14}--\eqref{II-64-15} converge absolutely and
\begin{align}
 |K^{\rm kp}(X;p,f)|
   &\le\tfrac12K_{\rm ch}b_e^2e^{-\kappa(d_t(X)+1)},
                                      \label{eq:history-kept-bound}\\
 |B^\times(X;p,f)|
   &\le\tfrac12C_3B_{\rm ch}b_e^2
       e^{-(\check R_t-2)}e^{-\kappa d_t(X)}.
                                      \label{eq:history-rejected-bound}
\end{align}
These functions are finite quadratic polynomials in the external
coordinates supported in $X$, with holomorphic coefficient functions
of the pair on $X$. They commute with restriction, real gauge
transformations and conjugation. A nonzero kept term requires
$\mathsf{kp}_t(X)=1$ and its unique component; a nonzero rejected term
requires $d_t(X)\ge\check R_t-2$. Both require a cube of $X$ contained
in $\Omega_t\cap Z_t$.

On the whole torus set
\begin{equation}\label{II-64-45}
 C=(A^+_{\rm II})^{-1},\quad M=-CA^+_{\rm IE},\quad
 \Delta^{\rm cond}=A^+_{\rm EE}-A^+_{\rm EI}CA^+_{\rm IE}.
\end{equation}
Then $\Delta^{\rm cond}=\sum_{\mathfrak z}d_{\mathfrak z}$ and
for $F=F_{\rm I}+f$,
\begin{equation}\label{eq:history-gaussian-completion}
 -\tfrac12\langle F,A^+F\rangle
 =-\tfrac12\langle F_{\rm I}-Mf,A^+_{\rm II}(F_{\rm I}-Mf)\rangle
  -\tfrac12\sum_{\mathfrak k}\langle f,\Delta^{\rm loc}_{\mathfrak k}f\rangle
  +\sum_XB^\times(X;p,f).
\end{equation}
Moreover
\begin{equation}\label{II-64-69}
 |M(b,b')|\le c_Me^{-\delta_M|b-b'|},\qquad
 c_M=c_{\rm lat}(\delta_{\rm pl}/2)B_{\rm f}^2,\quad
 \delta_M=\delta_{\rm pl}/2.
\end{equation}
\end{proposition}

\begin{proof}
For a triple in \eqref{II-64-11}, multiply its three nonnegative
majorants and sum the two intermediate bonds. Define its length as
the minimum sum of the three lengths over live intermediate bonds.
The coefficient trees and the inverse-string tree join at those bonds;
their union contains both endpoints and has the asserted length.
Its footprint lies within $\rho_*=4\ell_c$ of the union.
Spending one quarter of $\delta_{\rm pl}$ on each intermediate
bond sum proves
\begin{equation}\label{eq:history-chain-majorant}
 \sum_{\mathfrak z}m_{\mathfrak z}(b,b')
                e^{3\delta_{\rm pl}l_{\mathfrak z}(b,b')/4}
       \le B_{\rm ch}.
\end{equation}
The single-coefficient terms obey the same inequality. Replacing
the three factors one at a time gives the reference difference
$3\varepsilon_*m_{\mathfrak z}$; no factor is required to be
globally defined outside its own hull. In particular
\begin{equation}\label{eq:history-kept-reference-difference}
 \|\Delta^{\rm loc}_{\mathfrak k}(p)
           -\Delta^{\rm loc}_{\mathfrak k}(U,0)\|
       \le3c_{\rm lat}(3\delta_{\rm pl}/4)
                              \varepsilon_*B_{\rm ch}.
\end{equation}

Here is the conversion from connecting length to hull size.
Since $4\rho_*\le LM$, the big-cube hull of a live footprint of
length $l$ has
$n\le8+64l/(LM)$ cubes, and its diameter is at most $l+2\rho_*$.
For a fixed hull with $n$ cubes there are at most
$9(LM)^6n^2$ pairs of positive bonds. Thus its total quadratic
coefficient is bounded by
\[
 9(LM)^6n^2B_{\rm ch}
       \exp\{-3\delta_{\rm pl}\max(0,LM(n-8)/64)/4\}.
\]
Taking $\delta_{\rm pl}LM\ge C_2(\kappa+2)$ absorbs $n^2$ and
gives the first bound, using the connected-cube relation
$d_t(X)\le n-1$.
If a term is rejected, a live footprint cannot have diameter
less than $W_t$ by the keeping rule. Hence
$l\ge l_\times:=W_t-2-2\rho_*$.
Use three quarters of this lower bound and one quarter of
$LM(n-8)/64$ in the same exponential. Since
$l_\times\ge W_t/2-2$, the resulting exponent is at least
$c_2\delta_{\rm pl}W_t+(\kappa+1)n$.
This gives the intermediate bound
\[
 |B^\times(X;p,f)|
 \le\tfrac12 C_3B_{\rm ch}b_e^2(LM)^6
           e^{-c_2\delta_{\rm pl}W_t}e^{-(\kappa+1)n}.
\]
The two remaining size assumptions give
\eqref{eq:history-rejected-bound}. They also hold uniformly at
$\check R_t=1$.

All live entries have both endpoint bonds inside the footprint.
A kept hull meeting an external bond lies in its unique enlarged
strip component. A rejected hull must traverse the coarse separation;
counting its consecutive big cubes gives $d_t(X)\ge\check R_t-2$.
An external bond has an endpoint in the strip, so its containing
big cube lies in $\Omega_t\cap Z_t$ and belongs to $X$.
These are the stated support conditions. Absolute uniform convergence
of the coefficient series proves holomorphy. Only finitely many external
bonds lie in $X$, so pairing gives a finite quadratic polynomial.
Restriction and the two symmetries hold for every factor and hence
for the sums.

Insert \eqref{II-64-44} in the Schur complement; absolute convergence
allows the two coefficient sums to pass through it, giving
$\Delta^{\rm cond}=\sum d_{\mathfrak z}$. Symmetrization does not
change a quadratic form, so separating kept and rejected terms gives
\eqref{eq:history-gaussian-completion} by completing the finite square.
The two-factor sum $M=-\sum_{\jmath,S}t_\jmath p_{\rm I}a_Sp_{\rm E}$,
with half of the common exponential spent on its intermediate bond,
gives \eqref{II-64-69}.
\end{proof}

\subsection{Positivity of the kept part}
\label{sec:history-kept-positivity}

The lower bound uses $\mathscr O^+$ rather than $\mathscr O$:
the test field below contains external as well as integration bonds.
This is the only additional local-gauge requirement.

\begin{proposition}\label{II-102-10}
For sufficiently large $M$ and sufficiently small
$\alpha_1+\alpha_2$, depending only on $L$, every
$p\in\mathscr O^+(\mathfrak k^\sim;\alpha)$ satisfies, for every
complex field $f\in\operatorname{ran}p_{\mathfrak k}$,
\begin{equation}\label{II-102-11}
 \operatorname{Re}\langle\bar f,\Delta^{\rm loc}_{\mathfrak k}(p)f\rangle
                \ge\tfrac12\gamma_{\rm pos}|f|^2.
\end{equation}
For a real pair the kept matrix is real symmetric. Consequently, if
real pairs $p_{\mathfrak k}$ belong to these classes on a union $D$
of strip components, then, with $\gamma'=\gamma_{\rm pos}/4$,
every real field $f$ on their disjoint external bond sets satisfies
\begin{equation}\label{II-102-12}
 \sum_{\mathfrak k\subset D}
       \langle f,\Delta^{\rm loc}_{\mathfrak k}(p_{\mathfrak k})f\rangle
                    \ge2\gamma'\sum_b|f(b)|^2.
\end{equation}
The right side counts their disjoint external bond sets. The statement
is unchanged by extending the history beyond this step.
\end{proposition}

\begin{proof}
Take an admitted decomposition of $p$ near all bonds and work first
at its real reference. Put $R=\mathfrak k^\sim$,
$A^R=\sum_{X(S)\subset R}a_S$ and
$G^R=\sum_{X(\jmath)\subset R}t_\jmath$.
The latter is an approximate inverse with an exact interior identity.
Indeed let $\bar\eta_R$ be the sum of $\eta_y$ over the cubes usable
in $R$. Group the absolutely convergent strings by their first cube,
writing $G^R=\sum_y\zeta_yV_y$. Removing that first cube gives
\[
 V_y=G_y\eta_yY,\qquad
 Y=1-\sum_yF_yV_y,\qquad
 F_y=\sum_{X(S)\subset R}e_{y,S}
          =A^R_{\rm II}\zeta_y-\zeta_yA^y.
\]
These regroupings are legitimate even after removing $\zeta_y$:
it is strictly positive on the finite set $X_y$, and zero elsewhere.
Substitution yields
\begin{equation}\label{II-22-80}
 1-A^R_{\rm II}G^R=(1-\bar\eta_R)Y.
\end{equation}

Let $s(b)=\operatorname{dist}_\infty(b_-,\mathfrak k)$, and let
$J$ be the fluctuation bonds meeting the union of tiny cubes whose
enlargements lie in $R$. Then the external bonds of $\mathfrak k$
belong to $J$. On its complement $s\ge W_t-\ell_c$.
Also $\bar\eta_R(b)\ne1$ implies $s(b)>W_t-4\ell_c$:
a missing cutoff cube supplies a point outside $R$ within
$\ell_c/3+\ell_c+2\ell_c+1/2$ of $b_-$, whereas every cube
outside the one-coarse-layer enlargement is at distance at least
$W_t$ from $\mathfrak k$.

For real external $f$ supported on $\mathfrak k$, put
$g=A^R_{\rm IE}f$, $u=G^Rg$ and $h=g-A^R_{\rm II}u$.
Let $p_1$ project onto $J\cap\mathfrak X^+_{\rm I}$ and
$p_2=p_{\rm I}-p_1$. The test field $F=f-p_1u$ is supported on $J$.
Local positivity gives $\langle F,A^RF\rangle\ge
\gamma_{\rm pos}|F|^2\ge\gamma_{\rm pos}|f|^2$.
For the sum $\Sigma_R$ of all chains with hull in $R$, direct expansion
of this square gives
\begin{equation}\label{II-22-96}
 \langle f,\Sigma_Rf\rangle
 =\langle F,A^RF\rangle
     +\langle p_1u,h+A^R_{\rm II}p_2u\rangle-\langle g,p_2u\rangle.
\end{equation}
Here $\Sigma_R=A^R_{\rm EE}-A^R_{\rm EI}G^RA^R_{\rm IE}$.
By \eqref{II-22-80}, $h$ is supported where
$s>W_t-4\ell_c$; so is $p_2u$.

Put $B_h=B_{\rm f}c_{\rm lat}(\delta_{\rm pl}/2)$.
Both $A^R$ and $G^R$ have weighted Schur norm at most $B_h$ at rate
$\delta_{\rm pl}/2$. Conjugate by
$E(b)=e^{\delta_{\rm pl}s(b)/2}$; $s$ is 1-Lipschitz and
$s\le1$ on the support of $f$. Thus
\[
 |Eu|\le B_h^2e^{\delta_{\rm pl}/2}|f|,\qquad
 |Eh|\le(B_h+B_h^3)e^{\delta_{\rm pl}/2}|f|.
\]
The two error terms in \eqref{II-22-96} are therefore bounded by
\[
 e_P|f|^2,\qquad
 e_P=(2B_h^3+2B_h^5)
              e^{-\delta_{\rm pl}(W_t-4\ell_c-1)/2}.
\]
For external bonds $b,b'$ of $\mathfrak k$, any chain in $R$
not assigned to $\mathfrak k$ with $m_{\mathfrak z}(b,b')\ne0$
has $l_{\mathfrak z}(b,b')\ge l_\times=W_t-2-2\rho_*$.
Splitting the exponential in \eqref{eq:history-chain-majorant}
between this length and the endpoint distance bounds the difference
between $p_{\mathfrak k}\Sigma_Rp_{\mathfrak k}$ and $D_{\mathfrak k}$
by
\[
 e_\times=B_{\rm ch}c_{\rm lat}(3\delta_{\rm pl}/8)
                                      e^{-3\delta_{\rm pl}l_\times/8}.
\]
Choose $M$ so that $e_P+e_\times\le\gamma_{\rm pos}/4$ even with
$W_t=LM$, and choose $\alpha_1+\alpha_2$ so that the right side
of \eqref{eq:history-kept-reference-difference} is at most
$\gamma_{\rm pos}/4$. This proves \eqref{II-102-11}, first for real
$f$ at the reference, then for complex $f$ by splitting its real and
imaginary parts and applying the norm difference.
Conjugation gives reality at real pairs. Finally the external bond
sets of different strip components are disjoint; summing their bounds
proves \eqref{II-102-12}, including the empty union.
\end{proof}

\section{The normalizer on the integration region}
\label{sec:history-normalizer}

We localize the determinant of the Gaussian block integrated
in a history step. The nested background regions affect this block,
including its diagonal entries. Positivity follows by comparing
short connected terms with those of the constant sequence.

\subsection{The domain of the normalizer}

Work in level-$k$ units and put $t=k+1$, $\xi=L^{-k}$ and
$\xi_i=L^{i-k}$. Assume $L\ge13$ is odd,
$M=L^{m'}$, $m'\ge2$, and $m'<m+K-t$.
Let $\Omega_0\supset\cdots\supset\Omega_t$
be an admitted separated sequence, and let $\Lambda\subset\Omega_t$
be a union of coarse cubes with
$\operatorname{dist}_\infty(\Lambda,\Omega_t^c)\ge LM$.
Write $\mathfrak X_{\rm I}$ for its integration bonds and
$\Gamma^\Lambda$ for the positive level-$t$ bonds with an endpoint
in $\Lambda$. A bond is in $\mathfrak X_{\rm I}$ precisely when
it has an endpoint in $\Lambda$ and is not a pivot bond.
Indeed both endpoints then lie in $\Omega_t$, and the remaining
condition in the fluctuation coordinates is the exclusion of pivots.
In particular these are also nonpivot bonds for the constant sequence.

Keep the grain size $w$, clearance $r_{\rm gr}$ and mesh size $h$
of the connected-form construction. Impose
\begin{gather*}
 \mathfrak M\ge\max\{L,M^{\rm cv}\},\qquad
 R_1\mathfrak M\ge r^{\rm pr},\qquad
 M\ge R_1\mathfrak M\ge13\mathfrak M,\\
 Lw\le h,\qquad Lh\le\ell_c:=M/L^2,\qquad h\ge h_{\rm nf},\\
 4h\le r_N,\qquad 4(r_N+1)\le LM,\qquad M\ge M_{\rm nf}.
\end{gather*}
Here $w,h$ are powers of $L$, $M_1\le w\le M$,
$w\le L^{m+K-k-1}$, and
\[
 \rho_t+\rho_c+L<r_{\rm gr},\qquad
 4(r_{\rm gr}+1)\le w,\qquad
 \delta^{\rm cv}w\ge2^7L\log(4n_{\rm w}).
\]
The tower extent $\rho_t$ is less than $r_N$.
The constants $h_{\rm nf},M_{\rm nf}$ depend only on $L$.

For a nonempty coarse region $X$, a pair
$p=(\mathbf U,\mathbf J)$ on $X\cap\Omega_0$ belongs to
$\mathscr O^N(X;a)$ if it has one decomposition
$\mathbf U=e^{i\xi A'}U$, with $U$ real, satisfying:
\begin{enumerate}[label=(\roman*)]
\item At sites of layer $i$, wherever the expressions are
defined on $X\cap\Omega_0$,
\[
 |\partial U-1|,\ |\partial\mathbf U-1|<a\xi^2\xi_i^{-2},
 \qquad \xi_i|A'|<a,\quad \xi_i^2|\nabla^UA'|<a,\quad
 \xi_i^3|\mathbf J|<a.
\]
\item Every compact cube $q\subset X\cap\Omega_i$ which is a union
of level-$i$ $\mathfrak M$-cubes and has side at most
$12\mathfrak M\xi_i$ admits a real gauge
$U^v=e^{i\xi A}$ with $\xi_i|A|\le a$ and
$\xi_i^2|\nabla^\xi A|\le a$.
\item If a tiny cube $y$ contains an integration-bond endpoint
and $y^\sim\subset X\cap\Omega_0$, the same conclusion holds
on $y^\sim$ with $\xi_i$ replaced by $\xi_t=L$.
\end{enumerate}
The same decomposition must satisfy all three conditions.

These conditions are preserved by restriction, real gauge
transformations, conjugation and replacement by $p_{\rm r}=(U,0)$.
They define an open set when $a\le\pi$: keep $U$ and its gauges
fixed, choose the local inverse of the exponential at each
$i\xi A'(b)$, and use the finitely many strict inequalities in (i).
The inverses exist since $|\xi A'(b)|<a\le\pi$.

They imply regularity for every meet
$\Omega_\bullet\langle Y\rangle$, $Y\subset X$.
A site in layer $i$ of the meet had old layer index
$\widetilde i\ge i$, so $\xi_{\widetilde i}\ge\xi_i$ and
all scaled inequalities remain valid. Each gauge cube for the
meet lies in $X\cap\Omega_i$ and has the size in (ii);
take that cube itself as its enclosing cube. This verifies
the natural enlarged gauge-cube condition as well.

Choose $a_{\rm nf}\le\min\{a_N,a_{\rm w},a^{\rm cv}/2,1/2\}$
small enough that
\[
 R:=\frac{c_{\rm d}a^{\rm cv}}{2a_{\rm nf}}
       \ge\max\{16,16B_{\mathsf a}/\gamma^{\rm tr}\}.
\]
Here $\gamma^{\rm tr}$ is the flat gap used in
\cref{II-38-61}, and $a_{\rm w}$ is its local comparison radius.
On $p_z=(e^{iz\xi A'}U,z\mathbf J)$, $|z|<R$, the fixed real
gauges work for every meet, and \eqref{II-33-21} gives the
complex curvature bounds. Thus the connected coefficients
of the general-sequence form are defined and satisfy
\begin{equation}\label{eq:history-normalizer-walks}
 \begin{split}
 |a_S[p_z](b,b')|&\le
 C_a\epsilon_\star^{\#S-1}e^{-2\theta_a\Lambda_S(b,b')}
 =:\beta_S(b,b'),\\
 |a_S[p](b,b')-a_S[p_{\rm r}](b,b')|
   &\le(2/R)\beta_S(b,b').
 \end{split}
\end{equation}
The second bound is Cauchy's estimate. The coefficients vanish
unless $S$ is connected and contains the endpoint grain supports.
After conversion to the Hermitian norm, the weighted row and
column sums are bounded by
$B_{\rm nf}=2\sqrt2 C_a c_{\rm lat}(\theta_a)$, by
\eqref{eq:grain-entropy-sum}. Restricting to
$\Lambda_S\ge u$ gains $e^{-\theta_a u}$.
Set
\begin{equation}\label{eq:history-normalizer-constants}
 \begin{gathered}
 B_{\mathsf a}=\max\{\gamma^{\rm tr},B_{\rm nf}\},\qquad
 \gamma_{\rm nf}=\gamma^{\rm tr}/2,\qquad
 c_{\rm nf}=(B_{\mathsf a}+\gamma_{\rm nf})/2,\\
 q_N=\frac{B_{\mathsf a}-\gamma_{\rm nf}/2}
           {B_{\mathsf a}+\gamma_{\rm nf}}<1.
 \end{gathered}
\end{equation}
Increase $h_{\rm nf}$ so that the mesh localization in
\cref{II-38-61} applies and
$B_{\rm nf}e^{-\theta_a h}\le\gamma^{\rm tr}/4$.

\subsection{Positivity of the compressed operator}

Let $\mathcal B_Y$ be the integration bonds $b$ with
$S_N(b)\subset Y$. Let $A_{Y,N}(p)$ be the compression of
the fluctuation form for $\Omega_\bullet\langle Y\rangle$
to these bonds. Finite inclusion--exclusion of the defining
coefficients gives
\[
 A_{Y,N}(p_z)=
 \sum_{\substack{\emptyset\ne S\subset Y\\S\ {\rm connected}}}
                  \chi_Ya_S[p_z]\chi_Y^{\mathsf T}.
\]
Operators are extended by zero when compared on larger bond spaces.

\begin{proposition}\label{II-69-13}
For $p\in\mathscr O^N(Y;a_{\rm nf})$, $A_{Y,N}(p)$ is
complex symmetric, its real reference satisfies
\[
 \gamma_{\rm nf}I\le A_{Y,N}(p_{\rm r})\le B_{\mathsf a}I,
\]
and
\[
 \|A_{Y,N}(p)-A_{Y,N}(p_{\rm r})\|\le\gamma_{\rm nf}/4,
 \qquad \|1-c_{\rm nf}^{-1}A_{Y,N}(p)\|\le q_N.
\]
The assertions include $\mathcal B_Y=\emptyset$.
\end{proposition}

\begin{proof}
Symmetry and reality follow from the form and compression.
The row and column estimates give the upper bound and
$\|A_{Y,N}(p)-A_{Y,N}(p_{\rm r})\|\le2B_{\rm nf}/R
\le\gamma_{\rm nf}/4$. We prove the real lower bound.

Call a mesh cube $\square$ active if its cosine cutoff
$h_\square$ meets $\mathcal B_Y$, and put
$W=\bigcup_{\square\ {\rm active}}\square^{\sim2}$.
The cutoff is nonzero at $b$ exactly when its midpoint $x_b$
is within $h$ of the center of $\square$.
Every point of $W$ is therefore within $7h/2$ of such a
midpoint. Since $4h\le r_N$, $W\subset Y$.
An integration bond has an endpoint in $\Lambda$ and
$7h/2+1/2<LM$, so $W\subset\Omega_t$.
Conversely the closed $2h$ neighborhood of each $x_b$,
$b\in\mathcal B_Y$, lies in $W$.

For every mesh cube $\square'$ with
$q:=\square'^{\sim2}\subset W$, there is a tiny cube $y$ with
\begin{equation}\label{eq:history-normalizer-tiny}
 y\subset\Lambda,\qquad q\subset y^\sim\subset Y\cap\Omega_0.
\end{equation}
To verify this, lift $q$ to a product of intervals of length
$5h$. Each factor overlaps one or two coarse intervals.
Every product of overlapped coarse intervals belongs to $Y$,
since it contains a mesh cube of $q$. Choose a point in each
factor at the common coarse boundary if there are two intervals.
If there is one, choose the endpoint farther from its nearest
outer boundary, or the coarse midpoint if it lies in the factor.
The product point has distance at least $5h$ from every coarse
cube outside those products, and is within $7h/2+1/2<5h$
of an endpoint in $\Lambda$. Hence an overlapped coarse cube
belongs to $\Lambda$.

Since $5h<\ell_c$, each factor of $q$ overlaps at most two
tiny intervals. At a crossed coarse boundary choose the
adjacent tiny interval inside the selected $\Lambda$ cube;
its two neighbors cover the factor. In a single coarse
interval choose three consecutive tiny intervals containing
the factor and take the middle one. Their product is $y$.
Its enlargement lies in the overlapped coarse cubes, hence
in $Y$, and within $\ell_c<LM$ of $\Lambda$, hence in
$\Omega_0$. This proves \eqref{eq:history-normalizer-tiny}.
Every tiny cube in $\Lambda$ contains an integration-bond
endpoint: from its center choose a unit bond remaining in
the same averaging block. It is not a pivot, since pivots
join different blocks.

Condition (iii) supplies a gauge on $y^\sim$ and hence on $q$.
Conditions (i) and (ii) on $W$ are the constant-sequence
conditions at layer $t$, since $W\subset\Omega_t$.
Apply the mesh localization proof of \cref{II-38-61} to the
family $\mathcal V$ of connected $S$ with $S^\sim\subset W$.
It contains every term supported in an active
$\square^{\sim2}$, so is full near $\mathcal B_Y$.
Its real constant-sequence sum has gap $3\gamma^{\rm tr}/4$.
Both sequences agree on $S^\sim$ for $S\in\mathcal V$;
locality identifies their coefficients there.

Every remaining $S\subset Y$ reaches outside $W$.
A connecting tree for a nonzero entry comes within $3w$
of a point of $S^\sim\setminus W$ and contains $b_-$.
Because the $2h$ neighborhood of $x_b$ lies in $W$,
its length is greater than $2h-3w-1/2\ge h$.
The omitted terms have norm at most
$B_{\rm nf}e^{-\theta_a h}\le\gamma^{\rm tr}/4$.
The difference leaves gap $\gamma_{\rm nf}=\gamma^{\rm tr}/2$.

The real reference is self-adjoint, with spectrum between
$\gamma_{\rm nf}$ and $B_{\mathsf a}$. Thus
\[
 \|1-c_{\rm nf}^{-1}A_{Y,N}(p_{\rm r})\|
 \le\frac{B_{\mathsf a}-\gamma_{\rm nf}}
          {B_{\mathsf a}+\gamma_{\rm nf}}.
\]
Adding the complex perturbation divided by $c_{\rm nf}$
gives $q_N$.
\end{proof}

\subsection{Local determinant terms}

For $b\in\mathcal B_Y^+$ define
\begin{equation}\label{II-68-16}
 f_b(Y;p)=3\log c_{\rm nf}
 -\sum_{n\ge1}\frac1n
     \operatorname{tr}(1-c_{\rm nf}^{-1}A_{Y,N}(p))^n(b,b).
\end{equation}
The series converges normally, with
$|f_b|\le F_{\rm nf}:=3|\log c_{\rm nf}|-3\log(1-q_N)$.
For $S_N(b)\subset X$ put
\begin{equation}\label{II-68-17}
 E^N(X;p,b)=
 \sum_{S_N(b)\subset Y\subset X}(-1)^{\#X-\#Y}f_b(Y;p|_Y),
\end{equation}
and
\begin{equation}\label{II-68-18}
 {\rm Ld}(p)=\sum_{b\in\mathfrak X_{\rm I}^+}f_b(T_\ell;p).
\end{equation}
These expressions use the present sequence and integration
region; the constant-sequence definitions are their special case.

For $c\in\Gamma^\Lambda$ put
$z_c(p)=-\operatorname{Log}\det(Lm_c(\mathcal M^k\mathbf U))$.
Its branch and bound are those of
\eqref{eq:normalizer-pivot-log}. Indeed the endpoint blocks
fit in a cube of side $2\mathfrak M L$ inside $X(c)$.
This cube lies in $\Omega_t$, since its points are within
$LM$ of the endpoint in $\Lambda$. Condition (ii) at $t$
supplies its gauge. Restriction gives precisely the local
gauge used in that calculation, with
$L|A|\le a_{\rm nf}$ and $L|A'|<a_{\rm nf}$.
Since $a_{\rm nf}\le a_N$, the same holomorphic average,
pivot coefficient and logarithm apply.

Define
\begin{equation}\label{II-68-22}
 N^{\rm un}(X,b;p)=
 \begin{cases}
 -\tfrac12 E^N(X;p,b),
     &b\in\mathfrak X_{\rm I}^+,\ S_N(b)\subset X,\\
 z_c(p),&b=b_0(c),\ c\in\Gamma^\Lambda,\ X=X(c),\\
 0,&\text{otherwise}.
 \end{cases}
\end{equation}
The cases are disjoint. With the tents of \eqref{II-44-21}, put
\begin{align}
 N^{\rm un}(X,p,z)&=\sum_bh_z(b)N^{\rm un}(X,b;p),
                                                 \label{II-68-23}\\
 N^\partial(X;p)&=\sum_zN^{\rm un}(X,p,z).
                                                 \label{II-68-24}
\end{align}

\begin{proposition}\label{II-69-18}
After increasing $M_{\rm nf}$, constants $E_{\rm nf}$ and
$\delta_{\rm nf}>0$ depending only on $L$ give, for
$p\in\mathscr O^N(X;a_{\rm nf})$ and $d=d_t(X)$,
\begin{align*}
 |E^N(X;p,b)|&\le2^{-d}E_{\rm nf}e^{-\delta_{\rm nf}Md},\\
 |N^{\rm un}(X,b;p)|&\le E_{\rm nf}e^{-\delta_{\rm nf}Md},\\
 |N^{\rm un}(X,p,z)|&\le3L^3E_{\rm nf}e^{-\delta_{\rm nf}Md},\\
 |N^\partial(X;p)|&\le24(LM)^3E_{\rm nf}e^{-\delta_{\rm nf}Md}.
\end{align*}
The first expression vanishes when $X$ is not wall connected.
All terms are holomorphic, gauge invariant and conjugation
real. The pointed term vanishes unless $z\in X$.
\end{proposition}

\begin{proof}
The connected-chain proof of \cref{II-44-37} applies to
these matrices. We verify its inputs. For $S\subset Y$
the coefficient is the restriction of the same $a_S[p]$,
independent of $Y$. Its support $S^\sim$ lies within $3w$
of a tree realizing $\Lambda_S$ up to any positive error.
Equation \eqref{eq:history-normalizer-walks} gives the
weighted row and column bounds, and \cref{II-69-13} the
uniform contraction $q_N$. Changing a coarse cube away
from $S^\sim$ changes neither the coefficient nor its entries.
On wall-disconnected regions the tower meets, local stencils
and averaging blocks split into components; their inverses,
bilinear forms and deep compressions do likewise.

Inclusion--exclusion leaves only closed chains of coarse
support exactly $X$. Their degree-$n$ trace coefficient obeys
the two bounds \eqref{eq:normalizer-chain-majorants}:
\[
 |e_n|\le3Q_{\rm nf}^n e^{-\theta_a LMd/128},
 \qquad |e_n|\le3\,2^{\#X}q_N^n,
 \qquad Q_{\rm nf}=1+B_{\rm nf}/c_{\rm nf}.
\]
The geometry is unchanged: the support cubes lie within
$r_N+1$ of the joined trees, since $3w<r_N$, and the
total tree length is at least $LMd/128$.
For $d\ge1$ split $\sum |e_n|/n$ at
$\lfloor\vartheta LMd\rfloor$,
$\vartheta=\theta_a/(256\log Q_{\rm nf})$.
The two parts are at most
\[
 \frac{3Q_{\rm nf}}{Q_{\rm nf}-1}e^{-\theta_a LMd/256},
 \qquad
 \frac3{1-q_N}
 e^{16d\log2-\vartheta LMd\log(1/q_N)}.
\]
Choose $\delta_{\rm nf}$ below both positive rates and
increase $M_{\rm nf}$ so the remaining factor is at most
$2^{-d}$. For $d=0$, at most eight cubes occur, and
$2^8F_{\rm nf}$ suffices. Enlarge $E_{\rm nf}$ to contain
this bound, both geometric-series constants and
$6\log L+3\log2$.

Component splitting makes $f_b(Y)$ depend only on the
component containing $S_N(b)$; the alternating sum over
any other component vanishes. Pivot terms have
$d_t(X(c))=0$, giving the bond bound.
The tents satisfy $\sum_bh_z(b)=3L^3$ and $\sum_zh_z(b)=1$.
Their support is within $L$ of the bond, hence in $X$
for a nonzero term. This gives the pointed bound and support.
At most $24(d+1)(LM)^3$ bonds contribute to $N^\partial$.
For $d\ge1$ none is a pivot, and
$\tfrac12\,2^{-d}(d+1)\le1$ gives the last bound.
For $d=0$ use the bond bound directly.

Normal convergence gives holomorphy. Gauge covariance
conjugates diagonal blocks and preserves the pivot
determinant. Conjugation preserves the series and its
specified branch. Finite sums preserve these properties.
\end{proof}

\subsection{The normalization identity}

\begin{proposition}\label{II-69-19}
For $p\in\mathscr O^N(T_\ell;a_{\rm nf})$, let $A_{\rm II}(p)$
be the full block on the integration bonds. Then
\begin{align}
 \sum_XN^\partial(X;p)
 &=-\tfrac12{\rm Ld}(p)+\sum_{c\in\Gamma^\Lambda}z_c(p),
                                          \label{eq:history-normalizer-sum}\\
 \det A_{\rm II}(p)&=e^{{\rm Ld}(p)},\qquad
 \operatorname{Re}A_{\rm II}(p)\ge\tfrac34\gamma_{\rm nf}I,
                                          \label{eq:history-normalizer-det}\\
 \int e^{-(B,A_{\rm II}(p)B)_1/2}\,dB
 &=\exp\bigl(\nu_{\mathfrak g}\#\mathfrak X_{\rm I}^+
                  -\tfrac12{\rm Ld}(p)\bigr),
 \quad \nu_{\mathfrak g}=\tfrac32\log(2\pi).
                                     \label{eq:history-normalizer-integral}
\end{align}
The sum is over wall-connected coarse domains.
Empty integration sets have determinant and integral one.
\end{proposition}

\begin{proof}
Finite M\"obius summation gives
$\sum_{X\supset S_N(b)}E^N(X;p,b)=f_b(T_\ell;p)$.
Disconnected terms vanish. Summing tents and pivots proves
\eqref{eq:history-normalizer-sum}.
On the whole torus the meet is the original sequence, so
$A_{T_\ell,N}=A_{\rm II}$. The trace of its convergent
matrix logarithm is $\mathrm{Ld}$, proving the determinant
identity. The real gap and the perturbation bound give
the stated real-part bound.
The complex Gaussian formula applies, with square-root
branch $e^{-{\rm Ld}/2}$: at the real reference both are
positive, and on the disc to $p$ the Neumann logarithm
and Gaussian integral are continuous and nonzero.
This proves the integral with the stated normalization.
\end{proof}


\section{Fluctuation integrals with a conditional mean}
\label{sec:conditional-mean}

In a history step only the bonds touching $\Lambda_t$ are
integrated. The remaining field produces a conditional mean.
Its real part stays below the cutoff, but its imaginary part
must be removed from the Gaussian exponent without moving
the boundary of the small-field integration region.

\subsection{The Gaussian data}

Let $\Lambda_t\subset\Omega_t$ be regions for $\pi_t$,
$t=k+1$, and let $\mathfrak X$ be the non-pivot bonds of
$\Omega_t$. Set
\begin{equation}\label{II-114-6}
 \mathfrak X_{\rm I}=\{b\in\mathfrak X:b\text{ touches }\Lambda_t\},
 \qquad\mathfrak X_{\rm E}=\mathfrak X\setminus\mathfrak X_{\rm I}.
\end{equation}
For a bond set $\mathfrak Y$, write $\mathfrak Y\langle Z\rangle$
for its bonds touching $Z$, and $\mathfrak Y^\circ\langle Z\rangle$
for those with both endpoints in $Z$. All bond sets are
closed under reversal; positive bonds are independent
coordinates. Give the integration bonds reversal-invariant
thresholds $1\le\mathfrak b_-\le\mathfrak b_b\le\mathfrak b_+<\infty$.
For an actual step these are
\begin{equation}\label{II-114-7}
 \mathfrak b_b=\mathfrak b_k\quad(b\in\mathbb B(\Lambda_t)),
 \qquad
 \mathfrak b_b=\theta_c\mathfrak b_k
                  \quad(b\in\mathfrak x(\Lambda_t)).
\end{equation}

\begin{definition}\label{II-114-9}
A Gaussian structure with conditional mean on a region
$\mathfrak R$ consists of families $(\mathbf C,\boldsymbol\Gamma,
\mathbf m)(Z_\circ,s)$, indexed by fine regions
$Z_\circ^\sharp\subset\mathfrak R$ and the same decoupling
polydiscs as before. Put
$\mathfrak f=\mathfrak X_{\rm I}\langle Z_\circ\rangle$ and
$\mathfrak r=\mathfrak X_{\rm I}\langle\mathfrak R\rangle$.
The covariance $\mathbf C$ is symmetric and invertible on
$\mathfrak f$, the operator $\boldsymbol\Gamma$ maps fields
on $\mathfrak r$ to fields on $\mathfrak f$, and $\mathbf m$
is a field on $\mathfrak f$. Their entries, including those
of $\mathbf C^{-1}$, are holomorphic near the closed
polydisc $|s(\Delta)|\le e^{\bar\kappa}$.

When $s$ vanishes outside a coarse region $Z\supset Z_\circ^\sharp$,
the covariance is the direct sum over its wall components,
and $\boldsymbol\Gamma$ and $\mathbf m$ restrict to the
corresponding component families. The connecting operator
on a component uses auxiliary bonds touching that component
only. In particular $\boldsymbol\Gamma$ vanishes on auxiliary
bonds outside $\mathfrak X_{\rm I}\langle Z\rangle$.

For each $Z_\circ$ there are real references $C_{\rm r},
\Gamma_{\rm r},m_{\rm r}$, independent of $s$, such that,
with $S_{\rm r}=C_{\rm r}^{-1}$,
\begin{equation}\label{II-114-11}
 S_{\rm r}\ge\gamma_{\rm r}>0,\qquad
 \|S_{\rm r}\|,\ \|\Gamma_{\rm r}\|,
       \ \|C_{\rm r}\Gamma_{\rm r}\|\le c_{\rm r},
 \qquad 0<\gamma_{\rm r}\le1\le c_{\rm r}.
\end{equation}
Here $C_{\rm r}$ is symmetric, $\bar\kappa\ge1$ and
$\hat\varepsilon,\theta_{\mathbf m},\varepsilon_{\mathbf m}\ge0$.
The reference $\Gamma_{\rm r}$ uses only auxiliary bonds
touching $Z_\circ^\sharp$. Throughout the closed polydisc,
\begin{equation}\label{II-114-12}
 \|\mathbf C^{-1}-S_{\rm r}\|,\quad
 \|\boldsymbol\Gamma-\Gamma_{\rm r}\|,\quad
 \|\boldsymbol\Gamma^{\mathsf T}\mathbf C\boldsymbol\Gamma
             -\Gamma_{\rm r}^{\mathsf T}C_{\rm r}\Gamma_{\rm r}\|
       \le\hat\varepsilon,
\end{equation}
and, on every positive integration bond,
\begin{equation}\label{eq:conditional-mean-reference-bound}
 |m_{\rm r}(b)|\le\theta_{\mathbf m}\mathfrak b_b,
 \qquad |\mathbf m(b)-m_{\rm r}(b)|
              \le\varepsilon_{\mathbf m}\mathfrak b_b,
 \qquad\theta_{\mathbf m}+\varepsilon_{\mathbf m}<1.
\end{equation}
The operators and mean on an empty integration set have
their unique zero-dimensional values.
\end{definition}

Write $S=\mathbf C^{-1}$. The conditioning integral is
\begin{equation}\label{II-114-16}
\begin{split}
 \mathbb J^{\mathbf m}(Z_\circ;s)[F]
 ={}&\int d\mu_0(\mathbb X)\det{}^{1/2}(S/2\pi)
       \int dB\,F(B)\\
 &\quad\times\exp\left\{-\tfrac12
 (B-\mathbf m+\mathbf C\boldsymbol\Gamma\mathbb X,
 S(B-\mathbf m+\mathbf C\boldsymbol\Gamma\mathbb X))_1\right\}.
\end{split}
\end{equation}
The determinant has the branch fixed on matrices with
positive real part. The integral is used only after its
absolute convergence and $\operatorname{Re}S>0$ have been
proved. If $\mathfrak f=\emptyset$, it means $F(0)$.

Fix $\vartheta_*<1$ and constants satisfying
\begin{equation}\label{II-117-7}
 \theta_{\mathbf m}+\varepsilon_{\mathbf m}\le\vartheta_*,
 \qquad \mathsf C_{\rm mn}\ge18(1-\vartheta_*)^{-2}.
\end{equation}
For $\gamma_2\ge0$ put
\begin{equation}\label{II-117-5}
\begin{aligned}
 \lambda_{\rm r}&=\min\{\gamma_{\rm r}/2,(1+2c_{\rm r}^2)^{-1}\},
 &c_{\rm mn}&=\mathsf C_{\rm mn}\lambda_{\rm r}^{-1},\\
 u_0&=\mathsf C_{\rm mn}
       [\hat\varepsilon+(1+c_{\rm r})\varepsilon_{\mathbf m}^2+\gamma_2],
 &c_{\rm P}&=\mathsf C_{\rm mn}(LM)^3c_{\rm mn}u_0,\\
 \gamma_P&=\mathsf C_{\rm mn}^{-1}\gamma_2\mathfrak b_-^2.
\end{aligned}
\end{equation}
The smallness condition for this section is
\begin{equation}\label{II-117-9}
 \mathsf C_{\rm mn}u_0\le\lambda_{\rm r}.
\end{equation}
In particular $\hat\varepsilon\le\lambda_{\rm r}/16$ and,
writing $d_*=1-\vartheta_*$,
$\varepsilon_{\mathbf m}\le d_*/4$.
Indeed $\mathsf C_{\rm mn}^2
[\hat\varepsilon+(1+c_{\rm r})\varepsilon_{\mathbf m}^2+\gamma_2]
\le\lambda_{\rm r}\le1/3$ and
$\mathsf C_{\rm mn}\ge18d_*^{-2}$ give both inequalities.

\subsection{The contour estimate}\label{subsec:ymm1-56-2}

Let $\mathfrak S,\mathfrak e\subset\mathfrak f$ be disjoint.
The first set carries strict small-field tests and the second
carries failed tests. For $\varepsilon'>\varepsilon_{\mathbf m}$,
consider the tube where $B(b)$ is real and
$|B(b)|\ge\mathfrak b_b$ on $\mathfrak e$, where
$|\operatorname{Im}B(b)|<\varepsilon'\mathfrak b_b$ elsewhere,
and where additionally
$|\operatorname{Re}B(b)|+|\operatorname{Im}B(b)|<\mathfrak b_b$
on $\mathfrak S$. A function on this tube is admitted if
it is Borel and, at fixed real values on $\mathfrak e$,
holomorphic in all remaining variables. Let $D$ be the
real part of the tube and extend its tested integrand by
zero off $D$.

\begin{proposition}\label{II-117-69}
Suppose \eqref{II-117-7} and \eqref{II-117-9}. Let $G$ be
admitted on the tube and satisfy $|G|\le F_0$ there. If
$s$ vanishes outside a coarse region
$Z_\circ^\sharp\subset Z\subset\mathfrak R$,
then
\begin{equation}\label{II-117-71}
 |\mathbb J^{\mathbf m}(Z_\circ;s)[\mathbf1_DG]|
 \le F_0\exp\left\{-\mathsf C_{\rm mn}^{-1}\gamma_2
                \sum_{b\in\mathfrak e^+}\mathfrak b_b^2
               +c_{\rm mn}u_0N_{\rm I}(Z)\right\},
\end{equation}
where $N_{\rm I}(Z)=3\#\mathfrak X_{\rm I}^+\langle Z\rangle$.
The defining integrals converge absolutely, with an
integrable majorant independent of $s$ on its closed
polydisc. For a fixed bounded Borel integrand, the integral
is holomorphic in $s$ on a neighbourhood of that polydisc.
\end{proposition}

\begin{proof}
We retain the joint Gaussian form, since estimating the
conditional density alone would lose the auxiliary-field
normalization. Put
\begin{equation}\label{II-117-41}
 Q=\begin{pmatrix}
  1+\boldsymbol\Gamma^{\mathsf T}\mathbf C\boldsymbol\Gamma
      &\boldsymbol\Gamma^{\mathsf T}\\
  \boldsymbol\Gamma&S
 \end{pmatrix},\qquad
 Q_{\rm r}=\begin{pmatrix}
  1+\Gamma_{\rm r}^{\mathsf T}C_{\rm r}\Gamma_{\rm r}
      &\Gamma_{\rm r}^{\mathsf T}\\
  \Gamma_{\rm r}&S_{\rm r}
 \end{pmatrix}.
\end{equation}
Completion of the square gives $\det Q=\det S$ and
\[
 (z,Qz)=(\mathbb X,\mathbb X)_1+
 (\beta+\mathbf C\boldsymbol\Gamma\mathbb X,
             S(\beta+\mathbf C\boldsymbol\Gamma\mathbb X))_1,
 \qquad z=(\mathbb X,\beta).
\]
For real vectors,
$\|\beta\|^2\le2\|\beta+C_{\rm r}\Gamma_{\rm r}\mathbb X\|^2
+2c_{\rm r}^2\|\mathbb X\|^2$ proves
\begin{equation}\label{II-117-73}
 Q_{\rm r}\ge\lambda_{\rm r},\qquad
 \|Q-Q_{\rm r}\|\le2\hat\varepsilon,
 \qquad\det Q=\det S.
\end{equation}
The middle inequality follows by bounding the four blocks
of the difference by $\hat\varepsilon$.

For $x=(\mathbb X,\beta')$ real and $z=(\mathbb X,\beta'-i\eta)$,
the real reference and the perturbation bound give
\begin{equation}\label{II-117-44}
 \operatorname{Re}[-(z,Qz)/2]
 \le-(x,Q_{\rm r}x)/2+\hat\varepsilon\|x\|^2
                   +(c_{\rm r}/2+\hat\varepsilon)\|\eta\|^2.
\end{equation}
If $B=\beta+i\sigma$ and each $\sigma(b)$ lies on the
segment from $0$ to $\operatorname{Im}\mathbf m(b)$, this
implies
\begin{equation}\label{II-117-47}
 |e^{-(z,Qz)/2}|
 \le e^{c_{\rm r}\|\mathbf m\|^2}
       e^{-3\lambda_{\rm r}\|\mathbb X\|^2/8
          -\lambda_{\rm r}\|\beta\|^2/8},
 \qquad z=(\mathbb X,B-\mathbf m).
\end{equation}
To see this, use $\beta'=\beta-\operatorname{Re}\mathbf m$,
$\|\beta'\|^2\ge\|\beta\|^2/2-\|\operatorname{Re}\mathbf m\|^2$
and $\|\eta\|\le\|\operatorname{Im}\mathbf m\|$ in
\eqref{II-117-44}. The coefficient of each mean square is
at most $c_{\rm r}$, since $\lambda_{\rm r}\le1/3$ and
$\hat\varepsilon\le\lambda_{\rm r}/8$.
Also $\operatorname{Re}S\ge\gamma_{\rm r}-\hat\varepsilon>0$,
$\|S\|\le c_{\rm r}+1$, and
$\|\mathbf m\|^2\le\sum_b\mathfrak b_b^2$.
Thus \eqref{II-117-47}, with its mean factor replaced by
$e^{c_{\rm r}\sum_{b\in\mathfrak f^+}\mathfrak b_b^2}$ and multiplied by
$F_0((c_{\rm r}+1)/(2\pi))^{3\#\mathfrak f^+/2}$,
is a common integrable majorant. This proves absolute
convergence before any contour change.

The auxiliary field outside $Z$ integrates to one:
decoupling makes the corresponding columns of
$\boldsymbol\Gamma$ zero. Restrict both connecting
operators to $\mathfrak z=\mathfrak X_{\rm I}\langle Z\rangle$.
The compressed perturbation bounds remain valid, so
\eqref{II-117-73} holds for the resulting $Q^Z,Q^Z_{\rm r}$.
Their dimension $n$ is at most $2N_{\rm I}(Z)$.

Set $r_1=1-3d_*/4$, $r_2=1-d_*/4$ and
\[
 h(v)=\min\{1,\max\{0,2d_*^{-1}(r_2-v)\}\},\qquad
 \phi^{\mathfrak b}(\beta)=h(|\beta|/\mathfrak b).
\]
For $\nu=\operatorname{Im}\mathbf m$, move each bond by
\begin{equation}\label{II-117-62}
 \Theta(\beta)(b)=\beta(b)+i\phi_b(\beta(b))\nu(b),\qquad
 \phi_b=\begin{cases}
  0,&b\in\mathfrak e,\\
  \phi^{\mathfrak b_b},&b\in\mathfrak S,\\
  1,&\text{otherwise}.
 \end{cases}
\end{equation}
This stays in the tube. On the support of the cutoff,
$|\beta|<r_2\mathfrak b_b$ and $|\nu|\le d_*\mathfrak b_b/4$;
outside it the contour is unchanged. Its Jacobian factor
on a small-field bond is
\[
 J_b=1+i\langle\nu(b),\nabla\phi^{\mathfrak b_b}(\beta(b))\rangle,
 \qquad |J_b|\le e^{2d_*^{-2}\varepsilon_{\mathbf m}^2}.
\]
The latter estimate follows from
$|\nabla\phi^{\mathfrak b}|\le2/(d_*\mathfrak b)$ and
$\sqrt{1+a^2}\le e^{a^2/2}$.

Here the contour identity is elementary but the fixed
boundary matters. For one bond, choose the first real
coordinate parallel to $\nu(b)$. At fixed transverse
coordinates the graph
$x_1\mapsto x_1+i|\nu(b)|\phi^{\mathfrak b_b}(\beta)$
is piecewise smooth, joins the same endpoints as the real
interval, and lies in the convex complex slice of the tube.
Cauchy's theorem equates the two line integrals; the graph
derivative is exactly $J_b$. The two contours coincide
where $|\beta|\ge r_2\mathfrak b_b$, so no boundary value of
$G$ is needed. On an untested bond use a rectangular
contour and let its real endpoints tend to infinity;
\eqref{II-117-47} makes the vertical sides tend to zero.
Iterate over the finitely many bonds. Fubini applies to
every intermediate contour by the same Gaussian majorant
and the bounded Jacobian factors. The failed-test bonds
are left fixed, so a bounded Borel factor in those variables
is allowed throughout.

Write $\beta'=\beta-\operatorname{Re}\mathbf m$ and
$\eta=(1-\phi)\nu$. On a failed-test bond,
$|\beta'(b)|\ge d_*\mathfrak b_b$. Wherever $\eta(b)\ne0$,
either that bond is failed, or
$|\beta(b)|>r_1\mathfrak b_b$ and
$|\beta'(b)|\ge d_*\mathfrak b_b/4$. Consequently
\[
 \|\eta\|^2\le16d_*^{-2}\varepsilon_{\mathbf m}^2\|\beta'\|^2.
\]
Put $\Sigma_0=\hat\varepsilon+(1+c_{\rm r})
\varepsilon_{\mathbf m}^2+\gamma_2$ and
\[
 a=\hat\varepsilon+8d_*^{-2}(1+c_{\rm r})
                      \varepsilon_{\mathbf m}^2+\gamma_2
       \le8d_*^{-2}\Sigma_0.
\]
Insert \eqref{II-117-44}, the Jacobian bound and
\[
 1\le\exp\left\{\gamma_2\|\beta'\|^2
            -\gamma_2d_*^2\sum_{b\in\mathfrak e^+}\mathfrak b_b^2\right\}.
\]
The modulus of the shifted integrand is at most
\begin{equation}\label{II-117-79}
 F_0e^{(2/3)d_*^{-2}\varepsilon_{\mathbf m}^2N_{\rm I}(Z)}
 e^{-\gamma_2d_*^2\sum_{b\in\mathfrak e^+}\mathfrak b_b^2}
 e^{-(x,(Q^Z_{\rm r}-2a)x)/2},
 \qquad x=(\mathbb X,\beta').
\end{equation}
Since $\Sigma_0\le\lambda_{\rm r}/\mathsf C_{\rm mn}^2$,
$2a<\lambda_{\rm r}/2$. Thus
$A=Q^Z_{\rm r}-2a\ge\lambda_{\rm r}/2$. Translate by the
real mean and enlarge the real integration domain. Gaussian
integration cancels the powers of $2\pi$ and leaves
$|\det Q^Z|^{1/2}(\det A)^{-1/2}$.
Factor out $A$; the determinant inequality
$|\det(1+D)|\le(1+\|D\|)^n\le e^{n\|D\|}$ gives
\[
 |\det Q^Z|^{1/2}(\det A)^{-1/2}
 \le e^{2n(a+\hat\varepsilon)/\lambda_{\rm r}}
 \le e^{4(a+\hat\varepsilon)N_{\rm I}(Z)/\lambda_{\rm r}}.
\]
The remaining volume rate is bounded by
\[
 4(a+\hat\varepsilon)\lambda_{\rm r}^{-1}
       +\tfrac23d_*^{-2}\varepsilon_{\mathbf m}^2
 \le37d_*^{-2}\lambda_{\rm r}^{-1}\Sigma_0
 \le c_{\rm mn}u_0.
\]
Finally $d_*^2\ge\mathsf C_{\rm mn}^{-1}$, proving
\eqref{II-117-71}. Empty integration spaces give the same
bound directly.

For holomorphy, extend the matrix entries and the mean to
a common neighbourhood of the closed polydisc, symmetrizing
the continued $S$ block if necessary. On a smaller open
neighbourhood their joint form is within
$2\hat\varepsilon+\lambda_{\rm r}/8$ of the fixed real
reference and the mean norm is below
$(\sum_b\mathfrak b_b^2)^{1/2}$. This neighbourhood exists
by continuity and the strict bound
$\|\mathbf m\|\le\vartheta_*(\sum_b\mathfrak b_b^2)^{1/2}$.
The proof of \eqref{II-117-47} applies with
$\hat\varepsilon+\lambda_{\rm r}/16\le\lambda_{\rm r}/8$
in place of $\hat\varepsilon$. It supplies a common
integrable majorant there. The fixed-domain integral is
therefore holomorphic by dominated integration. In the
zero-dimensional case it is constant.
\end{proof}

\subsection{The terms of the expansion}\label{subsec:ymm1-56-3}

The localization uses only the part of a fine region which
contains integration bonds. For a fine region $Z_\circ$,
let $Z_\circ^{\rm I}$ be the union of its fine cubes containing
an endpoint of an integration bond, and let
$\operatorname{Gd}(Z_\circ)$ be the union of the coarse cubes
of $Z_\circ^\sharp$ which meet $\Lambda_t$. Then
\begin{equation}\label{II-116-22}
 \mathfrak X_{\rm I}\langle Z_\circ^{\rm I}\rangle
       =\mathfrak X_{\rm I}\langle Z_\circ\rangle,\qquad
 (Z_\circ^{\rm I})^\sharp\subset\operatorname{Gd}(Z_\circ)
                           \subset Z_\circ^\sharp.
\end{equation}
The first identity follows by taking the fine cube containing
the endpoint. To check the second, take such a cube $\delta$,
an endpoint $x\in\delta$, and the endpoint $x'\in\Lambda_t$
of the same bond. Every coarse cube $\Delta$ of
$\delta^\sharp$ meets $\delta$. For $y\in\Delta\cap\delta$,
$|y-x'|_\infty\le M+1<LM$. The centres of $\Delta$ and
the coarse cube of $\Lambda_t$ containing $x'$ are therefore
at distance less than $2LM$. Their distance is an integer
multiple of $LM$, so the cubes touch. Thus $\Delta$ meets
$\Lambda_t$, proving the inclusion.

The free parameters are now the coarse cubes outside this
guard. Extend $s$ by $1$ on
$\operatorname{Gd}(Z_\circ)\setminus(Z_\circ^{\rm I})^\sharp$,
obtaining the family $s^{\rm gd}$ on
$\mathsf P(\mathfrak R,Z_\circ^{\rm I})$.
This is an affine map between the corresponding closed
polydiscs.

Let $\underline{\mathbb V}=(\mathbb V^o)_o$ be a finite
system of potentials with
\begin{equation}\label{II-117-21}
 |\mathbb V^o(Y;B)|\le\mathsf u^o(Y)\le1,\qquad
 \mathsf w^o(Y)\ge e\,\mathsf u^o(Y)
\end{equation}
on the complex ranges
$|\operatorname{Re}B(b)|+|\operatorname{Im}B(b)|<\mathfrak b_b$
for $b\in\mathfrak X_{\rm I}^{\circ+}\langle Y\rangle$.
They are Borel and holomorphic there. A vertex $(Y,o)$
has $Y^\sharp\subset\mathfrak R$ and contains a source cube.
For a finite set $\mathbf D$ of vertices put
$Y^{\rm v}=\bigcup_{(Y,o)\in\mathbf D}Y$. Choose a
reversal-invariant failed set
$\mathfrak e\subset\mathfrak X_{\rm I}
 \setminus\mathfrak X_{\rm I}\langle Y^{\rm v}\rangle$
with $b^\sharp\subset\mathfrak R$. For
$\iota=(\mathbf D,\mathfrak e)\ne(\emptyset,\emptyset)$ set
\begin{equation}\label{II-115-95}
\begin{aligned}
 Z_\circ(\iota)&=\bigcup_{(Y,o)\in\mathbf D}Y
                    \ \cup\bigcup_{b\in\mathfrak e}\square(b),\\
 F_\iota(B)&=(-1)^{\#\mathfrak e^+}
     \mathsf{in}_{\mathfrak X_{\rm I}\langle Y^{\rm v}\rangle}(B)
     \mathsf{out}_{\mathfrak e}(B)
       \prod_{(Y,o)\in\mathbf D}(e^{\mathbb V^o(Y;B)}-1).
\end{aligned}
\end{equation}
Here $\mathsf{in}$ uses the strict tests
$|B(b)|<\mathfrak b_b$, and $\mathsf{out}$ their complements.
The product is assigned zero when an inner test fails,
before evaluating any potential. Define
\begin{equation}\label{II-115-97}
 f_\iota(s)=
   \mathbb J^{\mathbf m}(Z_\circ(\iota)^{\rm I};s^{\rm gd})[F_\iota].
\end{equation}
For a set $\mathcal K$ of free coarse cubes, its term and
domain are
\begin{equation}\label{II-115-98}
 T(\iota,\mathcal K)=f_\iota[\mathcal K\mid\emptyset],
 \qquad X(\iota,\mathcal K)=Z_\circ(\iota)^\sharp
                                    \cup\bigcup\mathcal K.
\end{equation}
Thus the bracket is the alternating sum at the indicators
of the subsets of $\mathcal K$.

\begin{proposition}\label{II-117-102}
Under the hypotheses of \cref{II-117-69}, every term is
defined and satisfies
\begin{equation}\label{II-117-23}
 |T(\iota,\mathcal K)|
 \le e^{c_{\rm P}N_\Lambda(X)}
       \prod_{(Y,o)\in\mathbf D}\mathsf w^o(Y)\,
       e^{-\gamma_P\#\mathfrak e^+}
       e^{-(\bar\kappa-1)\#\mathcal K}.
\end{equation}
Here $N_\Lambda(X)$ counts the coarse cubes of $X$ meeting
$\Lambda_t$. An insertion $\mathbb W(Y_1;B)$ bounded by
$\mathsf w'(Y_1)$ on its complex range gives the same
bound with the additional factor $\mathsf w'(Y_1)$ and
with $Y_1$ included in the domain and the inner tests.
The estimate also holds after multiplying a plain term's
integrand by any Borel function of the failed variables
of modulus at most one.
\end{proposition}

\begin{proof}
On the tube with
$\mathfrak S=\mathfrak X_{\rm I}\langle Y^{\rm v}\rangle$
and failed set $\mathfrak e$, the untested part of
$F_\iota$ is admitted. Indeed the variables of each
potential lie in its complex range, and
$|e^z-1|\le e|z|$ for $|z|\le1$ bounds it by
$F_0=\prod_{(Y,o)\in\mathbf D}\mathsf w^o(Y)$.
The restrictions use the same integration variables by
\eqref{II-116-22}.

If $s$ vanishes off $\mathcal K$, then $s^{\rm gd}$
vanishes outside
$Z=\operatorname{Gd}(Z_\circ)\cup\bigcup\mathcal K$.
Apply \cref{II-117-69} there. Each failed threshold is
at least $\mathfrak b_-$, and
\begin{equation}\label{II-117-94}
 N_{\rm I}(Z)\le18(LM)^3N_\Lambda(Z).
\end{equation}
For this count, every integration bond touching a coarse
cube of $Z$ also has an endpoint in a cube of $\Lambda_t$.
The same centre-distance argument as above shows that
the first cube meets $\Lambda_t$. Such a cube contains
$(LM)^3$ level-$k$ sites, each an endpoint of at most six
positive bonds. Multiplication by the three Lie-algebra
coordinates proves \eqref{II-117-94}. Since $Z\subset X$
and $\mathsf C_{\rm mn}\ge18$, \eqref{II-117-71} yields
\[
 |f_\iota(s)|\le
 F_0 e^{c_{\rm P}N_\Lambda(X)}e^{-\gamma_P\#\mathfrak e^+}.
\]
The function is holomorphic near its closed polydisc by
\cref{II-117-69} and the affine guard extension. Apply
the finite-difference Cauchy estimate on the free
coordinates in $\mathcal K$, setting all other coordinates
to zero. Its factor is
$(e^{\bar\kappa}-1)^{-\#\mathcal K}
\le e^{-(\bar\kappa-1)\#\mathcal K}$ for $\bar\kappa\ge1$.
This proves \eqref{II-117-23}.

An insertion simply adds its variables to the inner
tests and its bound to $F_0$; the ordinary vertices are
unchanged. A bounded Borel function of the failed variables
preserves admittedness because those variables are fixed
in the holomorphy assertion and in the contour deformation.
These observations prove the other two claims.
\end{proof}


\subsection{The local conditions on a graded field}
\label{sec:history-graded-gaussian-domain}

Only two local conditions on a graded field enter the Gaussian
construction. We state them explicitly, so that the additional response
conditions on a graded domain play no role in the following argument.
Let $P=(U,A',\mathbf J)$ be a triple at level $t$ on $Y$,
$\mathbf U=e^{i\eta_tA'}U$, and let $\sigma=(t,{\rm mid})$.
For a site $x$ put
$n(x)=\max\{0\le n\le t:x\in\Omega_n^\diamond\}$.
The graded radii have the form
\[
 R_\sigma^{p,i}(x)
   =c_\sigma(n(x))r_i(\sigma;n(x))L^{-pn(x)},\qquad
 0<c_\sigma(n)\le1+\beta_{\rm s}.
\]
Here $0<\beta_{\rm s}\le1/4$; more explicitly,
$c_\sigma(t)=1+\beta_{\rm s}$ and, for $n<t$,
\[
 c_\sigma(n)=1-\beta_{\rm s}(1-2^{-(t-n)})
               +\theta_{\rm s}\beta_{\rm s}2^{-(t-n)},
 \qquad 0\le\theta_{\rm s}<1.
\]
The two required conditions are the following.
\begin{enumerate}
\item[(L1)] On each contained stencil, with its indicated base point,
\[
 |\partial U-1|,\ |\partial\mathbf U-1|<R_\sigma^{2,0},
 \quad \eta_t^3|\mathbf J|<R_\sigma^{3,0},
 \quad \eta_t|A'|<R_\sigma^{1,1},
 \quad \eta_t^2|\nabla^UA'|<R_\sigma^{2,1}.
\]
The bond quantities are based at the initial endpoint. Differences
are required only when their entire stencil is in $Y$.
\item[(L2)] For each $n\le t$ and each embedded cube
$q\subset Y\cap\Omega_n^\diamond$ of side at most
$\mathsf oLM$ in level-$n$ units, $U$ has a real gauge
$U^v=e^{i\eta_n A}$ on $q$ with
\[
 |A|,\ |\nabla^{\eta_n}A|
       <\mathsf oLM B^{({\rm cg})}r_0(\sigma;n).
\]
Also, every cube $q\subset Y$ made of level-$n$
$\mathfrak M$-cubes, contained in $\Omega_n^\diamond$ and of side
at most $12\mathfrak M$ in level-$n$ units, has a containing
embedded cube $q'\subset Y$ on which $U$ has a real gauge of
index $n$ at most $a_\flat$.
\end{enumerate}
Thus (L2) includes a containing-box condition, not just a gauge on
the tested set. In its first clause we have retained only the
current-structure instance of the graded condition.
These are the local requirements of the graded domains;
we do not use their tower-response clauses here.

\begin{proposition}\label{II-87-15}\label{II-154-38}
Assume the two history gaps and $\mathsf oL^3\ge3$. Suppose $P$
satisfies (L1)--(L2), and, for $0\le n\le t$,
\begin{equation}\label{eq:history-graded-radius-test}
 \begin{gathered}
 (1+\beta_{\rm s})r_0(\sigma;n)
      \le\min\{\alpha'_0,\alpha'_2\},\qquad
 (1+\beta_{\rm s})r_1(\sigma;n)\le\alpha'_1,\\
 a_\flat\le\alpha'_0,\qquad
 \mathsf oL^3MB^{({\rm cg})}r_0(\sigma;k)\le\alpha'_0.
 \end{gathered}
\end{equation}
For every coarse region $R\subset Y$, the level-$k$ reading
\[
 p_k(P)=(\mathbf U,L^{-3}\mathbf J),\qquad
 \mathbf U=e^{i\xi(L^{-1}A')}U
\]
on $R\cap\Omega_0^\diamond$ belongs to
$\mathscr O^+(R;\alpha'_0,\alpha'_1,\alpha'_2)$ and hence to
$\mathscr O(R;\alpha'_0,\alpha'_1,\alpha'_2)$.
It belongs to the normalizer class
$\mathscr O^N(R;a)$ whenever
$a\ge\max_i\alpha'_i$.
\end{proposition}

\begin{proof}
A stencil in layer $n$ has base point of depth $l\ge n$.
The bound $c_\sigma(l)\le1+\beta_{\rm s}$ and
$L^{-pl}\le L^{-pn}$ give the four bounds
\eqref{eq:history-pair-bounds}. The change of units is exact:
$\xi L^{-1}A'=\eta_tA'$,
$\xi^3L^{-3}\mathbf J=\eta_t^3\mathbf J$, and the corresponding
first-difference expression is unchanged. The second part of (L2)
gives the plain and natural cube gauges, including their containing
boxes.

It remains to check the tiny cubes near all fluctuation bonds.
Let $y$ contain an endpoint $z$ of such a bond, and let its other
endpoint $z'$ lie in $\Omega_t$. Every point of $y^\sim$ is within
$2\ell_c+1<8LM\check R_t$ of $z'$. The second history gap therefore
puts $y^\sim$ in $\Omega_k$ when $k\ge1$; for $k=0$ its required
containment in $\Omega_0^\diamond$ is already part of the test.
The cube is embedded, since its side $3\ell_c$ is less than the
torus side. Its side in level-$k$ units is $3M/L^2\le\mathsf oLM$.
Apply the first part of (L2) with $n=k$.
Passing from index $k$ to index $t$ costs at most $L^2$ in the two
gauge bounds. The last inequality in
\eqref{eq:history-graded-radius-test} gives the required index-$t$
gauge. The same real factor $U$ is used throughout.
The normalizer assertion follows from the containing-cube argument
following \cref{II-31-36}; restriction introduces no new stencil.
\end{proof}

\subsection{The square root on the integration bonds}
\label{sec:history-gaussian-root}

For the conditional Gaussian we need the square root of the inverse
on $\mathfrak X_{\rm I}$, rather than the inverse strings used in
the kept form. We use the mesh, cutoffs $h_\square$ and string
construction of \cref{sec:gaussian-local-resolvents}, with the actual
integration bonds as its coordinate set. Choose
\[
 13w\le h,\qquad L^3\le h,\qquad Lh\le\ell_c,
\]
and enlarge the fixed lower bound for $h$ as specified below.
Use the mesh-gauge radius $a_{\rm w}\le a_{\rm cv}/64$ fixed in
\cref{sec:gaussian-local-resolvents}, and require
$\alpha_0\le\min\{a_{\rm pos},a_{\rm w}/2\}$ and
$2(\alpha_1+\alpha_2)\le c_{\rm d}a_{\rm w}$.
The following domain check is what permits that construction.

If the mesh neighborhood $\widetilde\square^{\,2}$ is contained
in $R\cap\Omega_0^\diamond$ and has a nonempty local coordinate set
$Y_\square$, there is a tiny cube $y$ meeting an integration bond
such that
\begin{equation}\label{eq:history-mesh-tiny-cover}
 \widetilde\square^{\,2}\subset y^\sim
                         \subset R\cap\Omega_0^\diamond.
\end{equation}
To see this, choose an integration endpoint $x\in\Lambda_t$
belonging to a bond of $Y_\square$ and work in the local chart.
Each coordinate interval $I$ of $\widetilde\square^{\,2}$ has
length $5h<\ell_c$ and contains the corresponding coordinate of
$x$ in its interior. Let $J_0$ be the big interval containing that
coordinate. If $I$ meets only $J_0$, choose the tiny interval
containing the coordinate and move its index, if necessary, to
the range $2,\ldots,L^3-1$. Its three-interval enlargement covers
$I$ and stays in $J_0$. If $I$ crosses a big-grid face, choose the
tiny interval in $J_0$ adjacent to that face. Its enlargement
covers $I$ and lies in the two crossed big intervals. Every big
cube so used meets the interior of $\widetilde\square^{\,2}$ and
therefore belongs to $R$. The product tiny cube is in the big
cube of $x$, hence meets the integration set. The history margin
puts its enlargement in $\Omega_0^\diamond$.

In fact the same margin gives
$\widetilde\square^{\,2}\subset\Omega_t$.
On this neighborhood the collared sequence is the constant sequence.
The pair restrictions, their domains and every coefficient $a_S$
with $S^\sim\subset\widetilde\square^{\,2}$ coincide with the
constant-sequence ones by \cref{II-31-36}.
The gauge supplied by \eqref{eq:history-mesh-tiny-cover} is precisely
the mesh gauge used there. The local form is thus a coordinate
compression of that constant-sequence form, and its real part has
gap $3\gamma_{\rm tr}/4$ on the complex disc used in
\cref{II-33-17}. This assertion concerns the inner-plaquette form
and its stated coordinate space.

Here are useful common constants for the form and root terms:
\[
 B_{\rm w}=2\sqrt2C_a,\qquad \delta=\theta_a,\qquad
 \gamma_{\rm flat}=3\gamma_{\rm tr}/8,\qquad
 \theta_\square=\frac{3\theta_a\gamma_{\rm tr}}{8B_a}.
\]
Exponential conjugation, as in \eqref{II-37-38}, gives
\[
 |g_\square^\mu(b,b')|_2
 \le(\mu+\gamma_{\rm flat})^{-1}
                          e^{-\theta_\square|b-b'|}.
\]
Put $c_1=c_{\rm lat}(\theta_\square/2)$,
$c_2=c_{\rm lat}(\delta/2)$ and
\[
 K_0=B_{\rm w}c_2
       \left[24\pi\left(\frac{2}{e\delta}+\frac12\right)
                    +\frac{16}{\delta}\right].
\]
Take $h\ge2K_0c_1/\gamma_{\rm flat}$. The commutator and remote-leg
bounds of that proof apply on the present coordinate set.
For the latter bound, if $h_\square(b)\ne0$ but
$S^\sim\not\subset\widetilde\square^{\,2}$, the connecting length
is at least
$2.5h-3w-h-\tfrac12\ge h$.
Consequently the error of the local resolvent sum has row norm
at most $K_0c_1/[h(\mu+\gamma_{\rm flat})]\le1/2$.
Its geometric series is the inverse of the actual compressed form,
not a separately defined continuation.

Use the ordered strings and the $\mu^{-1/2}/\pi$ integral
\eqref{II-37-69} to define $r_\omega$ on these local domains.
The same summation proof gives constants
\begin{equation}\label{eq:history-root-majorants}
 \begin{gathered}
 \delta_{\rm G}=\min\{\theta_\square/6,\delta\},\qquad
 B_{\rm G}=\max\{B_{\rm w},1,2c_1\gamma_{\rm flat}^{-1/2}\},\\
 \sum_\iota m_\iota(b,b')e^{\delta_{\rm G}l_\iota(b,b')}
       \le B_{\rm G},\qquad
 |T_\iota(p)-T_\iota(U,0)|
       \le\varepsilon_{\rm G}m_\iota,\\
 \varepsilon_{\rm G}=\frac{2(\alpha_1+\alpha_2)}
                              {c_{\rm d}a_{\rm w}} .
 \end{gathered}
\end{equation}
Here $T_\iota$ is a form term $a_S$, a root term $r_\omega$, or
its transpose. Their footprints lie within $\rho_{\rm G}=5h$
of a connecting tree of length $l_\iota$.
The difference estimate follows by Cauchy's estimate on the same
disc $|z|<c_{\rm d}a_{\rm w}/(\alpha_1+\alpha_2)$, assumed at
least $2$. This proves holomorphy and locality of the strings as
well, since the resolvent majorants are integrable in $\mu$.
On the whole torus the sum $R_{\rm I}=\sum_\omega r_\omega$ obeys
\begin{equation}\label{II-115-50}
 R_{\rm I}R_{\rm I}^{\mathsf T}
  =\chi_{\rm I}^{\mathsf T}\mathsf A_{\rm II}^{-1}\chi_{\rm I}.
\end{equation}
Indeed the resolvent sum is the actual inverse for every $\mu$;
the finite-dimensional square-root identity used in
\cref{II-37-96} then applies unchanged.

\subsection{The actual conditional Gaussian}
\label{sec:history-actual-conditional-gaussian}

Fix $Z_\circ^\sharp\subset R$ and put
$e=\mathfrak X_{\rm I}\langle Z_\circ\rangle$ and
$r=\mathfrak X_{\rm I}\langle R\rangle$.
Complete the parameter family by $\bar s=1$ on $Z_\circ^\sharp$,
by $s$ on the other big cubes of $R$, and by zero outside $R$.
Multiply a coefficient or string by the product of $\bar s$ over
its big-cube hull, omitting terms with zero product. Write
\begin{equation}\label{II-114-30}
 A(\bar s;p)=\sum_S\bar s(S)a_S(p),\qquad
 R_{\rm I}(\bar s;p)=\sum_\omega\bar s(\omega)r_\omega(p).
\end{equation}
Here and below the matrices use positive-bond coordinates; $\chi$
denotes restriction and $p=\chi^{\mathsf T}\chi$ denotes the
corresponding projection. Define
\begin{equation}\label{II-114-31}
 \begin{aligned}
 C&=(\chi_eA\chi_e^{\mathsf T})^{-1},&
 \Gamma&=\chi_eA(p_{\rm I}-p_e)R_{\rm I}\chi_r^{\mathsf T},\\
 M&=-C\,[\chi_e-\Gamma\chi_rR_{\rm I}^{\mathsf T}]
                                      A\chi_{\rm E}^{\mathsf T},&
 m&=M\mathsf a .
 \end{aligned}
\end{equation}
The external field $\mathsf a$ is given on external bonds whose
two endpoints lie in $R$, and is extended by zero.
For a triple $P$ use its reading $p_k(P)$; the reference objects
$C_{\rm r},\Gamma_{\rm r},M_{\rm r}$ are obtained by setting
$s=0$ and replacing that reading by $(U,0)$.
The reference mean is $m_{\rm r}=M_{\rm r}\mathsf a^{\rm re}$,
where $\mathsf a^{\rm re}=(\mathsf a+\mathsf a^*)/2$.
Empty coordinate sets have their unique zero-dimensional operators.

The small parameter below accounts separately for terms that leave
$Z_\circ^\sharp$ and for the complex displacement of the background.
Set
\begin{equation}\label{II-115-54}
 \begin{gathered}
 \gamma_0=\gamma_{\rm pos},\quad
 B_{\rm nm}=B_{\rm G}c_{\rm lat}(\delta_{\rm G}/8),\\
 \varepsilon_{\rm dc}
   =e^{64\bar\kappa}e^{-\delta_{\rm G}(M-1)/128}
                         +64\varepsilon_{\rm G},
 \qquad \bar\kappa\ge1,
 \end{gathered}
\end{equation}
and require
\begin{equation}\label{eq:history-gaussian-smallness}
 \rho_{\rm G}\le M/2,\qquad
 \delta_{\rm G}M\ge64\bar\kappa,\qquad
 B_{\rm nm}\varepsilon_{\rm dc}\le\gamma_0/2 .
\end{equation}
All preceding grain, clearance, mesh and positivity floors remain
in force. Put
\begin{equation}\label{II-115-55}
 \begin{gathered}
 \gamma_{\rm r}=\gamma_0,\qquad
 c_{\rm r}=B_{\rm nm}^2/\gamma_0,\qquad
 \hat\varepsilon=64B_{\rm nm}^5\gamma_0^{-2}\varepsilon_{\rm dc},\\
 \delta_M=\frac{\delta_{\rm G}\gamma_0}{128B_{\rm nm}},\qquad
 c_M=72B_{\rm nm}^5\gamma_0^{-2}
              c_{\rm lat}(\delta_{\rm G}/16)c_{\rm lat}(\delta_M).
 \end{gathered}
\end{equation}

\begin{theorem}\label{II-115-62}
Let the above hypotheses hold and let
$p_k(P)\in\mathscr O(R;\alpha_0,\alpha_1,\alpha_2)$.
The matrices \eqref{II-114-31}, including $C^{-1}$, are holomorphic
for $|s(\Delta)|<e^{\bar\kappa+1}$; $C$ is symmetric and invertible.
They have the local restriction and component-decoupling properties
of \cref{II-114-9}, and their real references satisfy
\eqref{II-114-11}--\eqref{II-114-12} with
\eqref{II-115-55}. Moreover,
\begin{equation}\label{II-115-64}
 |M_{\rm r}(b,b')|_2\le c_Me^{-\delta_M|b-b'|},\qquad
 |M(s;p)(b,b')-M_{\rm r}(b,b')|_2
        \le\varepsilon_{\rm dc}c_Me^{-\delta_M|b-b'|}.
\end{equation}
The second bound holds on $|s(\Delta)|\le e^{\bar\kappa}$.

Suppose, for every $b\in\mathfrak X_{\rm I}^+\langle R\rangle$,
\begin{equation}\label{II-115-21}
 \begin{split}
 c_M\sum_{b'}e^{-\delta_M|b-b'|}|\mathsf a^{\rm re}(b')|
                 &\le\theta\,\mathfrak b_b,\\
 \sqrt2c_M\sum_{b'}e^{-\delta_M|b-b'|}
              |\mathsf a(b')-\mathsf a^{\rm re}(b')|
                 &\le\varepsilon'\,\mathfrak b_b ,
 \end{split}
\end{equation}
where both sums are over positive external bonds with both endpoints
in $R$. Then
\[
 |m_{\rm r}(b)|\le\theta\mathfrak b_b,\qquad
 |m(b)-m_{\rm r}(b)|
 \le\bigl(\sqrt2\varepsilon_{\rm dc}\theta
               +(1+\varepsilon_{\rm dc})\varepsilon'\bigr)\mathfrak b_b.
\]
If $\theta+\sqrt2\varepsilon_{\rm dc}\theta+
(1+\varepsilon_{\rm dc})\varepsilon'<1$, these are a local system
of Gaussian structures with conditional mean in the sense of
\cref{II-114-9}. On the whole torus the system is complete, with
\begin{equation}\label{II-114-33}
 C_*=\mathsf A_{\rm II}^{-1},\qquad
 m_*=-\mathsf A_{\rm II}^{-1}\mathsf A_{\rm IE}\mathsf a .
\end{equation}
If also $(LM)^3\varepsilon_{\rm dc}\le1$, a defect number is
\[
 u_{\rm gs}=128B_{\rm nm}^5\gamma_0^{-2}+c_{\rm r}+3.
\]
\end{theorem}

\begin{proof}
We give the common product estimate first. It is needed for the mean,
as well as for the covariance. A factor may be $A$, $R_{\rm I}$
or $R_{\rm I}^{\mathsf T}$, with arbitrary coordinate projections
between factors. For $1\le q\le4$, define the nonnegative kernel
\[
 K_q(b,b')=B_{\rm G}^q
   \sum_{b_1,\ldots,b_{q-1}}
      \prod_{i=1}^q
          e^{-\delta_{\rm G}|b_{i-1}-b_i|/4},
 \qquad b_0=b,\quad b_q=b'.
\]
Its row and column sums are at most $B_{\rm nm}^q$, and
$K_q(b,b')\le B_{\rm nm}^q
e^{-\delta_{\rm G}|b-b'|/8}$. These follow by retaining half the
distance exponent for the endpoints and summing the remaining
factors by $c_{\rm lat}(\delta_{\rm G}/8)$.

For an individual factor of length $l_\iota$, the hull count
in \eqref{eq:gaussian-parameter-count} gives, on the open
polydisc under consideration,
$|\bar s(\iota)|\le e^{16\bar\kappa+\delta_{\rm G}l_\iota/8}$.
The count applies because its footprint is within $\rho_{\rm G}$
of a connecting tree, $\rho_{\rm G}\le M/2$, and the second
inequality of \eqref{eq:history-gaussian-smallness} absorbs the
length contribution. If the product starts on $e$ and one of
its footprints leaves $Z_\circ^\sharp$, the connected chain
must travel at least
\[
 D_*=(M-1)-\rho_{\rm G}\ge M/2-1.
\]
Indeed an endpoint in a fine cube of $Z_\circ$ has distance at
least $M-1$ from a big cube outside its sharp enlargement; the
footprint can shorten this distance by at most $\rho_{\rm G}$.
The same argument applies with the order reversed.

Split each product into chains entirely inside $Z_\circ^\sharp$
and chains leaving it. On the former, all parameters are $1$,
and telescoping the $q$ factors bounds the difference from the
real reference by $q\varepsilon_{\rm G}K_q$. On the latter,
retain the chain-length gain before summing the intermediate bonds.
The bound is
$e^{16q\bar\kappa}e^{-5\delta_{\rm G}D_*/8}K_q$.
For $q\le4$ their sum is bounded by
$\varepsilon_{\rm dc}K_q$. Thus, for every such product $\Pi_q$
with initial projection on $e$,
\begin{equation}\label{eq:history-parameter-product}
 |\Pi_{q,{\rm r}}(b,b')|\le K_q(b,b'),\qquad
 |\Pi_q(s;p)(b,b')-\Pi_{q,{\rm r}}(b,b')|
                         \le\varepsilon_{\rm dc}K_q(b,b').
\end{equation}
The corresponding row and column estimates justify all series,
locally uniformly on the open polydisc.

The reference compression $C_{\rm r}^{-1}$ has gap $\gamma_0$.
For precision, take the union $W$ of tiny cubes of $Z_\circ$
meeting integration bonds. Then $e\subset\mathfrak X_{\rm I}
\langle W\rangle$, $W^\sim\subset R$, and the set-theoretic
interior determined by $Z_\circ^\sharp$ contains the
$2\ell_c/3$-neighborhood of $W$. A footprint belongs to that set
exactly when its big-cube hull belongs to $Z_\circ^\sharp$.
Hence its partial coefficient sum is $A(\bar0;U,0)$, and
\eqref{eq:history-local-gap} gives the stated gap.
This application uses the chosen reference in $R$; it does not
require choosing it anew on $Z_\circ^\sharp$.

Apply \eqref{eq:history-parameter-product} with $q=1,2$ to
$C^{-1}$ and $\Gamma$. Their differences have norms at most
$B_{\rm nm}\varepsilon_{\rm dc}$ and
$B_{\rm nm}^2\varepsilon_{\rm dc}$. Thus $C^{-1}$ has real-part
gap $\gamma_0/2$. The resolvent identity and the three-term
expansion of
$\Gamma^{\mathsf T}C\Gamma-\Gamma_{\rm r}^{\mathsf T}
C_{\rm r}\Gamma_{\rm r}$ give respectively the remaining
bounds \eqref{II-114-11}--\eqref{II-114-12}; the latter difference
is at most $9B_{\rm nm}^5\gamma_0^{-2}\varepsilon_{\rm dc}$.
This is bounded by the chosen $\hat\varepsilon$.
Symmetry follows from that of the coefficients.

For the mean, put
$N=[\chi_e-\Gamma\chi_rR_{\rm I}^{\mathsf T}]
A\chi_{\rm E}^{\mathsf T}$, so $M=-CN$.
This is the difference of a one-factor product and a four-factor
product. Therefore
\[
 |N_{\rm r}(b,b')|_2\le2B_{\rm nm}^4e^{-\delta_{\rm G}|b-b'|/8},
 \qquad
 |N-N_{\rm r}|_2\le
     2B_{\rm nm}^4\varepsilon_{\rm dc}e^{-\delta_{\rm G}|b-b'|/8}.
\]
Exponential conjugation of $C^{-1}$ and $C_{\rm r}^{-1}$, followed
by their resolvent identity, gives, with $\delta'=2\delta_M$,
\[
 \begin{split}
 |C(b,b')|_2,\ |C_{\rm r}(b,b')|_2
       &\le4\gamma_0^{-1}e^{-\delta'|b-b'|},\\
 |C-C_{\rm r}|_2
       &\le16B_{\rm nm}\gamma_0^{-2}c_{\rm lat}(\delta'/2)
          \varepsilon_{\rm dc}e^{-\delta'|b-b'|/2}.
 \end{split}
\]
In the product $CN$, retain the rate $\delta'/2$ at the endpoints
and sum the remaining factor with $c_{\rm lat}(\delta_{\rm G}/16)$.
The identity
$CN-C_{\rm r}N_{\rm r}=(C-C_{\rm r})N_{\rm r}
+C(N-N_{\rm r})$ proves \eqref{II-115-64} with \eqref{II-115-55}.
For the fields, use
$m-m_{\rm r}=(M-M_{\rm r})\mathsf a^{\rm re}
+M(\mathsf a-\mathsf a^{\rm re})$. The real reference has equal
coordinate and operator norms; the complex comparison costs
the factor $\sqrt2$. This gives precisely the two asserted mean
bounds.

Every surviving coefficient or string has its whole connected
hull in a single wall component when the parameters vanish
outside a region $Z$. Its input and output bonds have both
endpoints in that component. Consequently $A$ and $R_{\rm I}$
split into their coordinate blocks, and multiplication and
inversion in \eqref{II-114-31} give the required decoupling of
$C,\Gamma,M$. In particular an external field is restricted
by both endpoints, rather than merely by incidence.
The identical local matrices occur under restriction, by
\cref{II-31-36} and the locality of the root strings.
This proves the local-system assertions. Restriction preserves
\eqref{II-115-21}, since its nonnegative sums only decrease.

On the torus, \eqref{II-115-50} and block multiplication give
\[
 C+C\Gamma\Gamma^{\mathsf T}C
       =\chi_e\mathsf A_{\rm II}^{-1}\chi_e^{\mathsf T},
 \qquad
 C[\chi_e-\Gamma\chi_{\rm I}R_{\rm I}^{\mathsf T}]
       =\chi_e\mathsf A_{\rm II}^{-1}\chi_{\rm I}.
\]
The second identity, multiplied by
$A\chi_{\rm E}^{\mathsf T}\mathsf a$, yields
$m=\chi_e m_*$. These are exactly the completeness identities
\eqref{II-114-13}. Finally,
$2(LM)^3\hat\varepsilon+c_{\rm r}+3\le u_{\rm gs}$ under the
stated additional hypothesis, proving the defect assertion.
\end{proof}

\subsection{The external tiers and the history step}
\label{sec:history-mean-tiers}

We now verify the mean condition for the external fields used in
the history step. The collar bonds are
$\mathfrak X_{\rm col}=\operatorname{Col}(\Lambda_t)$.
The rim bonds are external bonds incident to a coarse cube of
$\operatorname{rim}(Z_t)$; equivalently they belong to one of
the corresponding sets $c^\star$.
Put
\begin{equation}\label{II-4-46}
 {\rm rad}_k(b)=\mathfrak b_k
 \min\left(\{C_{\rm xt}\}\cup
      \left\{\frac1{4\Theta_1}:b\in\mathfrak X_{\rm col}\right\}
      \cup\{\theta_c:b\in\mathfrak X_{\rm rim}\}\right).
\end{equation}
Take $\mathsf a=\mathsf a'+\mathsf a''$, where $\mathsf a'$ is
real and $|\mathsf a'(b)|<{\rm rad}_k(b)$, while
$\mathsf a''=0$ on the rim and
$|\mathsf a''|\le\varepsilon_a\mathfrak b_k$ elsewhere.
For a field on $R$ these conditions are imposed after zero
extension from the external bonds with both endpoints in $R$.

Choose $C_e\ge C_{\rm xt}$ and define
\begin{equation}\label{II-115-82}
 \begin{aligned}
 \vartheta_{\rm col}
   &=\frac{c_Mc_{\rm lat}(\delta_M)}{4\Theta_1\theta_c},&
 \vartheta_{\rm rim}
   &=c_Mc_{\rm lat}(\delta_M/2)e^{-\delta_M\varpi_T/2},\\
 \vartheta_{\rm far}
   &=c_Mc_{\rm lat}(\delta_M/2)
                  \frac{C_e}{\theta_c}e^{-\delta_M(W_t-2)/2},&
 \vartheta_\mu^{\rm re}
   &=\vartheta_{\rm col}+\vartheta_{\rm rim}
       +\vartheta_{\rm far}(1+\varepsilon_a/C_e).
 \end{aligned}
\end{equation}
The numbers are uniform in the history and volume, apart from
the displayed width $W_t$. Define also
\begin{equation}\label{II-115-83}
 \varepsilon_\mu
   =\sqrt2\varepsilon_{\rm dc}\vartheta_\mu^{\rm re}
       +\sqrt2(1+\varepsilon_{\rm dc})
                         \vartheta_{\rm far}\varepsilon_a/C_e .
\end{equation}

\begin{proposition}\label{II-115-84}\label{II-119-53}
The above external fields satisfy \eqref{II-115-21} with
\[
 \theta=\vartheta_\mu^{\rm re},\qquad
 \varepsilon'=\sqrt2\vartheta_{\rm far}\varepsilon_a/C_e,
\]
provided $\mathfrak b_b\ge\theta_c\mathfrak b_k$.
Consequently, if \eqref{eq:history-graded-radius-test} holds at
the radii of \cref{II-115-62} and
$\vartheta_\mu^{\rm re}+\varepsilon_\mu<1$, the actual history
step gives a local Gaussian system with conditional mean
$(\theta_{\mathbf m},\varepsilon_{\mathbf m})
=(\vartheta_\mu^{\rm re},\varepsilon_\mu)$.
It has the normalizer and all-bond positivity applicability
of \cref{II-87-15}. Its local restrictions obey
\begin{equation}\label{II-119-54}
 \begin{aligned}
 C_R(Z_\circ,s;P)&=C_X(Z_\circ,s_X;P|_X),\\
 \Gamma_R(Z_\circ,s;P)
    &=\Gamma_X(Z_\circ,s_X;P|_X)\chi_{\mathfrak X_{\rm I}\langle X\rangle},
 \qquad
 m_R(Z_\circ,s;P,\mathsf a)
    =m_X(Z_\circ,s_X;P|_X,\mathsf a|_X)
 \end{aligned}
\end{equation}
when $Z_\circ^\sharp\subset X\subset R$ and $s$ vanishes outside
$X$. On the whole torus it has the covariance and mean
\eqref{II-114-33}.
\end{proposition}

\begin{proof}
If the integration set is empty there is nothing to prove.
For an external bond put
$d(b')=\min_{b\in\mathfrak X_{\rm I}^+}|b_--b'^+_-|_\infty$.
The two geometric facts needed are
\[
 d(b')\le\varpi_T-1\ \Longrightarrow\
      b'\in\mathfrak X_{\rm col},\qquad
 d(b')\le W_t-2\ \Longrightarrow\
      b'\in\mathfrak X_{\rm rim}.
\]
For the first, the minimizing integration bond has an endpoint
$z\in\Lambda_t$ at distance at most $\varpi_T$ from an endpoint
of $b'$. There is an inner nonpivot bond of $\Lambda_t$ at $z$,
by the pivot-incidence property. Thus the external bond is in
the collar. For the second, each endpoint of $b'$ is at distance
at most $1+(W_t-2)+1=W_t$ from that $z$. Its coarse cube is
outside $\Lambda_t$ and touches its coarse enlargement; hence
it is a rim cube. These arguments also show why the thresholds
are $\varpi_T-1$ and $W_t-2$.

The integer-valued distance divides the external bonds into
three disjoint sets:
$d\le\varpi_T-1$,
$\varpi_T\le d\le W_t-2$,
and $d\ge\varpi_T$ with $d>W_t-2$.
The bounds on $|\mathsf a'|$ there are respectively
$\mathfrak b_k/(4\Theta_1)$,
$\theta_c\mathfrak b_k$, and $C_e\mathfrak b_k$.
For a fixed integration bond the least-lift lattice sum is
at most $c_{\rm lat}(\delta_M)$. On the second and third
sets retain half the exponential for summation and use the
other half at the respective distances $\varpi_T$ and $W_t-2$.
This gives
\[
 c_M\sum_{b'}e^{-\delta_M|b-b'|}|\mathsf a'(b')|
 \le(\vartheta_{\rm col}+\vartheta_{\rm rim}
                       +\vartheta_{\rm far})\theta_c\mathfrak b_k .
\]
The field $\mathsf a''$ vanishes unless $d>W_t-2$; its sum is
at most
$\vartheta_{\rm far}(\varepsilon_a/C_e)\theta_c\mathfrak b_k$.
Finally
$|\mathsf a^{\rm re}|\le|\mathsf a'|+|\mathsf a''|$ and
$|\mathsf a-\mathsf a^{\rm re}|\le|\mathsf a''|$.
These prove the claimed two weighted sums, including the
factor $\sqrt2$ in the second one.

The graded-domain assertion follows from \cref{II-87-15};
the mesh condition follows from \eqref{eq:history-mesh-tiny-cover}.
Apply \cref{II-115-62}. Every nonzero term in
\eqref{II-114-30} has its hull in $X$ when $s$ vanishes outside
$X$, and its local matrix is unchanged on restriction.
This gives \eqref{II-119-54} directly from
\eqref{II-114-31}. The torus conclusion is the completeness
assertion of \cref{II-115-62}.
\end{proof}

\section{The potentials in a history step}
\label{sec:history-actual-potentials}

The potentials are evaluated at a background that is defined on a
collared region and is glued to an arbitrary field outside it.
We first show that the local arguments of the potentials lie strictly
inside their analyticity domains. The current in these arguments is
the current of the background, at the scale of the term being
evaluated. This scale dependence must be retained when a term from
an earlier step is inserted.

\subsection{The glued background lies in the analytic cores}
\label{sec:history-actual-core-membership}

Retain the running and separation hypotheses of
\cref{II-10-20}, with a collared sequence of depth \(n\).
Thus \(\bar B\gamma^2\le1\),
\(\kappa_{\rm run}=(1+\bar B\gamma^2)^{1/2}\le\sqrt2\), and
\[
 \varepsilon_i\le\varepsilon_j\quad(i\le j),\qquad
 \varepsilon_i\le\kappa_{\rm run}L^{(i-j)/2}\varepsilon_j
       \quad(i\ge j).
\]
Use the regular-gauge constants \(B_*,c_{*,\rm rg}\) of
\cref{I-28-75}, and the averaging radius \(c^{\rm SO}_2\)
of \cref{I-10-101}. Define
\begin{equation}\label{II-90-7}
 C_{\rm mb}
 =4L^3B_*\max\left\{12\mathfrak M+2,\,
        \frac{\mathsf oLM+2}{\mathsf oLM B^{(\rm cg)}}\right\},
\end{equation}
\begin{equation}\label{II-90-8}
 \epsilon_{\rm mb}=\min\left\{
  \frac{c_{*,\rm rg}}{4L^3(\mathsf oLM+2)},\
  \frac{c_{*,\rm rg}}{4L^3(12\mathfrak M+2)},\
  \frac{c^{\rm SO}_2}{\sqrt2L^2},\
  \frac{\bar\varepsilon}{2\sqrt2L^2},\
  \frac{\varepsilon_{\rm v}}{\sqrt2}\right\}.
\end{equation}
These constants are fixed before the couplings and the history.
In particular, \(C_{\rm mb}\) depends only on
\(L,\mathsf o,\mathfrak M,M,B^{(\rm cg)}\); the additional
dependence of \(\epsilon_{\rm mb}\) is on the two regularity
radii. We use the three conditions
\begin{align}
 &\mathsf oLM+6<L^{m+K-n},
       \label{II-90-11}\\
 &C\varepsilon_n\le\epsilon_{\rm mb},
       \label{II-90-12}\\
 &C_{\rm mb}C\varepsilon_n\le\alpha_0/2 .
       \label{II-90-13}
\end{align}
The first ensures that the cubes used for gauges embed in the
torus. The second controls the gauges and the tower backgrounds.
The third places the resulting triples inside half the fixed core.

\begin{proposition}\label{II-90-39}
Let \(U\) be a real configuration on
\(\operatorname{Bd}^{\circ}(\Omega_0^\diamond)\), critical over
its own determining averages, with
\[
 r_U(\mathfrak e)<C\varepsilon_{\lambda(\mathfrak e)}
\]
at each regularity element, where \(C>0\). Assume
\eqref{II-90-11} and \eqref{II-90-12}. Let \(\mathsf U\) be
any extension of \(U\) to the bare torus. For
\(1\le j\le n\) and \(X\in\mathcal D_{j-1}\) with
\(X\subset\Omega_j\), its plain triple is
\begin{equation}\label{II-90-40}
 \mathbf P_X^{(j)}(\mathsf U)
  =\bigl(U|_{\operatorname{Bd}(X)},\,0,\,
              J^{(j)}(U)|_{\operatorname{Bd}(X)}\bigr).
\end{equation}
It is independent of the extension, and
\begin{equation}\label{II-90-48}
 \mathbf P_X^{(j)}(\mathsf U)
      \in\mathcal U_j^\circ(X;b_0,b_1,c_0)
 \quad\text{if}\quad
 b_0,c_0\ge C_{\rm mb}C\varepsilon_j,\quad b_1>0 .
\end{equation}
The region \(X\) need not be a box or be contractible.
If \(1\le i\le j\) and \(X'\in\mathcal D_{i-1}\) is
contained in \(X\), then
\begin{equation}\label{II-90-49}
 \rho_{X'}^{j\to i}\mathbf P_X^{(j)}(\mathsf U)
                   =\mathbf P_{X'}^{(i)}(\mathsf U).
\end{equation}
Under \eqref{II-90-13}, all these triples belong to
\begin{equation}\label{II-90-50}
 \mathcal U_j^\circ(X;\alpha_0/2,\alpha_1/2).
\end{equation}
In particular the conclusions hold with \(C=C_{\rm reg}\)
for every element of the collared background of
\cref{II-10-20}, and for its gluing to any outer field.
\end{proposition}

\begin{proof}
We give the gauge and tower arguments, since the pointwise
curvature bound by itself would not prove core membership.
If a regularity element has a corner in
\(\Omega_i^\diamond\), its highest layer index \(\lambda\)
satisfies \(\lambda\ge i\). The running inequalities give
\begin{equation}\label{II-90-41}
 L^{-2(\lambda-i)}\varepsilon_\lambda
 \le\kappa_{\rm run}L^{-3(\lambda-i)/2}\varepsilon_i.
\end{equation}
By the definition of the local radius, therefore,
\begin{align}
 |\partial U(p)-1|
   &<\kappa_{\rm run}C\varepsilon_iL^{-2i},
       \label{II-90-42}\\
 |J^{(u)}(U)(b)|
   &<\kappa_{\rm run}C\varepsilon_iL^{-3(i-u)} .
       \label{II-90-43}
\end{align}
The stronger factors \(L^{-3(\lambda-i)/2}\) and
\(L^{-5(\lambda-i)/2}\), respectively, may be retained.
Separation puts every bond incident to
\(X\subset\Omega_1\), with its full plaquette star, inside
\(\Omega_0^\diamond\). Thus the current of any extension
agrees there with the regional current. This proves
\eqref{II-90-40}.

For a cube \(q\subset\Omega_i^\diamond\) of side \(s\)
in level-\(i\) units, \(i\ge1\), choose a cube \(q^+\)
of integer side \(\ell=\lceil s\rceil+1<s+2\) containing
\(q\). Its concentric enlargement of side \(\ell+4\)
lies within distance \(4\) of \(q\) and embeds whenever
\(s+6\le L^{m+K-i}\). The history gap places the full
plaquette stars in this enlargement in layer at least
\(i-1\). Put
\[
 w=\kappa_{\rm run}C\varepsilon_{i-1},
 \qquad \alpha'=2L^3w .
\]
There the curvature is smaller than
\(wL^2\eta_i^2\), and
\[
 |D_U^{\mathsf T}\partial U|
 \le \eta_i^{-1}(wL^2\eta_i^2)^2
           +\eta_i^2|J^{(i)}(U)|
 <2L^3w\eta_i^2 .
\]
Here \(L\eta_i\le1\) and \(w<1\). This is the
matrix-divergence calculation already used in
\eqref{eq:background-matrix-divergence}.
If \(4L^3(s+2)C\varepsilon_{i-1}\le c_{*,\rm rg}\),
then \(\ell\alpha'<c_{*,\rm rg}\). Applying the regular
cube gauge of \cref{I-28-75} at scale \(i\), and
restricting its gauge to \(q\), yields
\[
 U^v=e^{i\eta_iA},\qquad
 |A|,|\nabla^{\eta_i}A|
       <4L^3B_*(s+2)C\varepsilon_{i-1}.
\]
The finite differences here are required wherever their
stencils lie in the inner bonds of \(q\).
For \(i=0\) the comb gauge \cref{I-10-51} instead gives
\[
 |A|,|\nabla A|<2\sqrt2\pi sC\varepsilon_0
 \quad\text{if}\quad \sqrt2sC\varepsilon_0\le1.
\]
The two gauge conclusions, with \eqref{II-90-7} and
\eqref{II-90-8}, imply
\begin{align}
 |A|,|\nabla^{\eta_i}A|
  &<\mathsf oLMB^{(\rm cg)}C_{\rm mb}C\varepsilon_i
       &&(s\le\mathsf oLM), \label{II-90-44}\\
 |A|,|\nabla^{\eta_i}A|
  &<C_{\rm mb}C\varepsilon_i
       &&(s\le12\mathfrak M). \label{II-90-45}
\end{align}
For the second bound the cube is a union of the
level-\(i\) \(\mathfrak M\)-cubes. Its enlargement
embeds since the standing size condition gives
\(L^{m+K-i}\ge L^6\mathfrak M>12\mathfrak M+6\).
For the first use \eqref{II-90-11}. The two entries
of the minimum in \eqref{II-90-8} give the respective
smallness conditions. At level zero the bounds follow
from \(2\sqrt2\pi<B^{(\rm cg)}C_{\rm mb}\) and
\(24\sqrt2\pi\mathfrak M<C_{\rm mb}\).

Consider now a separated subordinate sequence of depth
\(n'\le n\), with reduced determining set \(\mathbf B'\).
By \cref{II-8-36}, the restriction \(U'\) is critical
over \(V'=M_{\mathbf B'}U'\). If \(\lambda'\) is its
local layer index, restriction and \eqref{II-90-41}
give
\[
 r_{U'}(\mathfrak e)
 <L^{-2(\lambda-\lambda')}C\varepsilon_\lambda
 \le\kappa_{\rm run}C\varepsilon_{\lambda'}
 \le e,\qquad
 e=\kappa_{\rm run}C\varepsilon_{n'}.
\]
The current has the still smaller restriction factor
\(L^{-3(\lambda-\lambda')}\).
By \eqref{II-90-8},
\[
 e\le\varepsilon_{\rm v},\qquad
 L^2e\le c^{\rm SO}_2,\qquad
 2L^2e\le\bar\varepsilon.
\]
For a tested plaquette of order \(q\), its four endpoint
blocks lie in \(\Omega'_{q-1}\). Their bare plaquettes
have curvature smaller than \((L^2e)L^{-2q}\).
The local four-corner induction in the proof of
\cref{II-195-31}, with terminal level \(q\) and
\(\alpha=L^2e\), bounds its averaged holonomy by
\(2L^2e\). Thus all tested holonomies of \(V'\)
satisfy this bound.
Hence \cref{I-58-22} identifies the restricted
configuration with the regular critical orbit of
\(V'\). The continuation at the real datum exists
on its admitted positive radius; at zero it is a real
gauge transform of \(U'\).

The analytic averages used in the tower definition are
also defined at \(U'\) and agree with the real averages.
Indeed, for a determining row \(c\) of level \(l\), put
\(D_q=B^{l-q}(c_-)\cup B^{l-q}(c_+)\), \(0\le q\le l\).
Separation puts \(D_0\) in the layer-\((l-1)\) region,
and \(B(D_{q+1})=D_q\). Every bare plaquette with corners
in \(D_0\) has curvature smaller than \((L^2e)L^{-2l}\).
Apply the local induction in the proof of \cref{II-195-31}
with terminal level \(l\) and \(\alpha=L^2e\): the one-step
estimate uses only the four subordinate corner blocks, all
contained in \(D_q\), and its products are bounded by
\(\exp(1/[2(L^2-1)])<2\). Thus, at intermediate level
\(q<l\), the averaged curvature is bounded by
\[
 2(L^2e)L^{2(q-l)}
       \le2L^{-2}c^{\rm SO}_2<\omega_M .
\]
The one-step analytic mean is therefore defined and
equals the real mean on that double block. Induction
on \(q\) proves the assertion up to row \(c\), including
the reversed row. The logarithmic perturbation of the
real datum is zero, strictly inside the continuation
radius.

Apply this argument to either tower of
\(Z\in\mathcal D_{p-1}\), \(Z\subset X\subset\Omega_j\),
truncated at \(n'\le p\le j\). The tower is separated
and subordinate to the collared sequence. Its continued
configuration is a real gauge transform of \(U|_Z\).
Every plaquette star in \(Z^{\wr p}\) lies inside \(Z\),
since the two surrounding coarse layers contain all
its corners. Real gauge invariance and
\eqref{II-90-42}, \eqref{II-90-43} give
\[
 |\partial U_{n'}^{\mathfrak r}-1|
       <\kappa_{\rm run}C\varepsilon_jL^{-2j},
 \qquad
 |J_{n'}^{\mathfrak r}|
       <\kappa_{\rm run}C\varepsilon_jL^{-3(j-n')}.
\]
These are the core bounds
\eqref{eq:analytic-core-tower-bounds} with
\(c=b_0L^{-(j-p)}\): use
\(b_0\ge C_{\rm mb}C\varepsilon_j
 >\kappa_{\rm run}C\varepsilon_j\) and
\[
 -3(j-n')=-2(j-p)+2(n'-p)-(j-n').
\]
The simple-domain conditions follow from the first
gauge bound, the same pointwise estimates, and \(A'=0\).
This proves \eqref{II-90-48}.

The restriction formula \eqref{II-90-49} follows from
\(J^{(i)}=L^{-3(j-i)}J^{(j)}\).
The running monotonicity
\(\varepsilon_j\le\varepsilon_n\) and
\eqref{II-90-13} give \eqref{II-90-50}.
Finally \cref{II-10-20} supplies criticality and the
layer bounds with \(C=C_{\rm reg}\) for every
collared background. Gluing outside
\(\Omega_0^\diamond\) does not change any of the
quantities just used.
\end{proof}


\subsection{The domains before, during and after a step}
\label{sec:history-graded-domains}

The radii at the intermediate stage leave room for both the
fluctuation and the terms inherited from earlier levels.
We describe the domains, including their tower conditions.
The two pointwise and gauge clauses are those used in
\cref{sec:history-graded-gaussian-domain}.

A structure is a stage
\(\sigma=(s,\phi)\), where \(\phi={\rm mid}\) or
\({\rm out}\); at level zero only \((0,{\rm out})\)
occurs. Order the stages by level and put
\((s,{\rm mid})\preccurlyeq(s,{\rm out})\).
For \(0<\beta_{\rm s}\le1/4\) and
\(0\le\theta_{\rm s}<1\), set
\begin{equation}\label{II-2-9}
 \begin{gathered}
 f_{s,l}=1-\beta_{\rm s}(1-2^{-(s-l)}),\qquad
 c_{(s,{\rm out})}(l)=f_{s,l},\\
 c_{(s,{\rm mid})}(s)=1+\beta_{\rm s},\qquad
 c_{(s,{\rm mid})}(l)
      =f_{s,l}+\theta_{\rm s}\beta_{\rm s}2^{-(s-l)}
           \quad(l<s).
 \end{gathered}
\end{equation}
In particular \(1-\beta_{\rm s}<c_\sigma(l)\le
1+\beta_{\rm s}\). A radius table assigns positive
numbers \(r_\iota(\sigma;l)\), \(\iota=0,1\),
to all these stages and layers. The running table is
\(r_\iota(\sigma;l)=\alpha_{\iota,l}\).
For \(x\in\Omega_0^\diamond\) put
\[
 m_\sigma(x)=\max\{l\le s:x\in\Omega_l^\diamond\},
 \qquad
 R_\sigma^{p,\iota}(x)
   =c_\sigma(l)r_\iota(\sigma;l)L^{-pl},
 \quad l=m_\sigma(x).
\]
At a common top layer use the ambient radius:
\begin{equation}\label{II-86-8}
 r_\iota(\sigma''|\sigma;l)=
 \begin{cases}
 r_\iota(\sigma;s),&l=s''=s,\\
 r_\iota(\sigma'';l),&\text{otherwise}.
 \end{cases}
\end{equation}
For a region \(Y\), write
\begin{equation}\label{II-86-10}
 Y_{\mathcal S}=Y\cap\Omega_0^\diamond,\qquad
 \operatorname{Bd}_{\mathcal S}(Y)
  =\operatorname{Bd}(Y)\cap
          \operatorname{Bd}^{\circ}(\Omega_0^\diamond).
\end{equation}
Outside this bond set the domain imposes no local
curvature or potential bound.

For \(\sigma'\preccurlyeq\sigma\), a tower address in
\(Y\in\mathcal D_{s'-1}\) consists of
\(\sigma''\preccurlyeq\sigma'\),
\(1\le i\le s''\), and
\(Z\in\mathcal D_{i-1}\) with \(Z\subset Y\) and
\(Z_{\mathcal S}\ne\emptyset\). At depth \(q\le i\)
use the sequence
\begin{equation}\label{II-86-16}
 \bigl(\Omega_l^\diamond\cap
                    \Omega_l^{\mathfrak r}(Z)\bigr)_{0\le l\le q},
 \qquad \mathfrak r\in\{\mathfrak r_1,\mathfrak r_2\}.
\end{equation}
This sequence is separated and subordinate: the
complement of an intersection is the union of the
two complements, so its gap is at least the smaller
of the two original gaps. The first member is
\(Z_{\mathcal S}\). Its tower configuration is the
analytic continuation from the real averages, with
the strict-radius convention in \cref{II-16-12}.
A core address replaces \(Z\) by
\(X'\in\mathcal D_{i-1}\) with
\(X'\subset Y\cap\Lambda_i\).

\begin{definition}\label{II-86-20}
Let \(s'\ge1\), \(Y\in\mathcal D_{s'-1}\), and
\(P=(U,A',\mathbf J)\) be a triple at level \(s'\)
on \(Y\), with \(\mathbf U=e^{i\eta_{s'}A'}U\).
The domain \(\tilde U_{\sigma'|\sigma}(Y)\) is defined
by the following four conditions, all with this
same decomposition.
\begin{enumerate}
\item[(L1)] On \(\operatorname{Bd}_{\mathcal S}(Y)\),
where the indicated stencils are contained,
\[
 \begin{gathered}
 |\partial U-1|,\ |\partial\mathbf U-1|
       <R_\sigma^{2,0},\qquad
 \eta_{s'}^3|\mathbf J|<R_\sigma^{3,0},\\
 \eta_{s'}|A'|<R_\sigma^{1,1},\qquad
 \eta_{s'}^2|\nabla_U^{\eta_{s'}}A'|
       <R_\sigma^{2,1}.
 \end{gathered}
\]
Plaquettes and differences use their base point;
bond bounds use the initial endpoint.
\item[(L2)] For every \(\sigma''\preccurlyeq\sigma'\),
\(0\le l\le s''\), and cube
\(q\subset Y\cap\Omega_l^\diamond\) of side at most
\(\mathsf oLM\) in level-\(l\) units, \(U\) has a
real gauge \(U^v=e^{i\eta_lA}\) on its inner bonds
with
\[
 |A|,|\nabla^{\eta_l}A|
      <\mathsf oLMB^{(\rm cg)}r_0(\sigma'';l).
\]
For every \(l\le s'\) and cube
\(q\subset Y_{\mathcal S}\cap\Omega_l^\diamond\)
made of level-\(l\) \(\mathfrak M\)-cubes and of
side at most \(12\mathfrak M\), there is a containing
cube \(q\subset q'\subset Y_{\mathcal S}\) with a
real gauge of index \(l\) at most \(a_\flat\).
\item[(L3)] At every tower address
\((\sigma'',i,Z)\), both rules, and every
\(1\le q\le p\le i\), the tower configurations
\(U_q^{\mathfrak r}\) and their currents exist.
On \(Z^{\wr p}\), write
\(\ell=\min\{m_{\sigma''}(x),i\}\), where \(x\)
is the base point or initial endpoint. Require
\begin{align}
 |\partial U_q^{\mathfrak r}-1|
 &<L^{i-p}c_\sigma(\ell)
          r_0(\sigma''|\sigma;\ell)L^{-2\ell},
        \label{II-86-26}\\
 \eta_q^3|J_q^{\mathfrak r}|
 &<L^{i-p}c_\sigma(\ell)
          r_0(\sigma''|\sigma;\ell)
                    L^{-2\ell-\min(q,\ell)} .
        \label{II-86-27}
\end{align}
The first inequality is imposed only when the four
plaquette bonds lie in
\(\operatorname{Bd}^{\circ}(Z_{\mathcal S})\);
the second only when the entire four-plaquette star
does. Require the same existence and bounds with
the real factor \(U\) in place of \(\mathbf U\).
\item[(L4)] At every core address \((\sigma'',i,X')\),
\begin{equation}\label{II-86-29}
 \rho_{X'}^{s'\to i}P\in
 \mathcal U_i^\circ\bigl(X';
   c_\sigma(i)r_0(\sigma''|\sigma;i),
   c_\sigma(i)r_1(\sigma''|\sigma;i)\bigr).
\end{equation}
\end{enumerate}
Write \(\tilde U_\sigma=\tilde U_{\sigma|\sigma}\).
The superscript \(1/2\) means that all entries of
the table and \(a_\flat\) are halved, with the
coefficients \(c_\sigma\) unchanged.
\end{definition}

Thus the output and intermediate domains of step \(t\)
are \(\tilde U_{(t,{\rm out})}\) and
\(\tilde U_{(t,{\rm mid})}\), respectively. A plain
triple on a label uses the regional current only
where it is defined:
\begin{equation}\label{II-86-39}
 \mathbf P^{(j)}_{Y|\mathcal S}(\mathsf U)
  =(\mathsf U|_{\operatorname{Bd}(Y)},0,\mathbf J),
 \qquad
 \mathbf J(b)=
 \begin{cases}
 J^{(j)}(\mathsf U)(b),
   &\text{the full star of \(b\) is in }
           \operatorname{Bd}^{\circ}(\Omega_0^\diamond),\\
 0,&\text{otherwise}.
 \end{cases}
\end{equation}
In particular the current is zero at the incomplete
stars along the collar. This convention commutes
with restriction and the scaling of currents.

\begin{proposition}\label{II-90-62}
Under the hypotheses of \cref{II-90-39}, suppose
the radius table up to depth \(n\) satisfies
\begin{equation}\label{eq:history-radius-floor}
 \begin{gathered}
 C_{\rm mb}C\varepsilon_l
     \le\tfrac12(1-\beta_{\rm s})r_0(\tau;l)
       \quad(0\le l\le\operatorname{level}(\tau)\le n),\\
 C_{\rm mb}C\varepsilon_n\le a_\flat/2 .
 \end{gathered}
\end{equation}
Then, for all \(\sigma'\preccurlyeq\sigma\) of
levels at most \(n\), \(s'\ge1\), and
\(Y\in\mathcal D_{s'-1}\),
\begin{equation}\label{II-90-71}
 \mathbf P^{(s')}_{Y|\mathcal S}(\mathsf U)
       \in\tilde U^{1/2}_{\sigma'|\sigma}(Y).
\end{equation}
This includes domains meeting the complement of
\(\Omega_0^\diamond\), domains that wrap around the
torus, and both output and intermediate stages.
\end{proposition}

\begin{proof}
For every applicable layer and every choice of
the two structures,
\begin{equation}\label{II-90-59}
 \kappa_{\rm run}C\varepsilon_l
 <C_{\rm mb}C\varepsilon_l
 \le\tfrac12(1-\beta_{\rm s})r_0(\tau;l)
 <\tfrac12c_\sigma(l)r_0(\tau;l).
\end{equation}
The same inequality holds with the radius selected
by \eqref{II-86-8}, since it is an entry of the same
table. At a plaquette base point of depth
\(l=m_\sigma(x)\), \eqref{II-90-42} and
\eqref{II-90-59} give the half-radius curvature
bound. For a full star, \eqref{II-90-43} gives
\[
 \eta_{s'}^3|J^{(s')}(U)(b)|
   <\kappa_{\rm run}C\varepsilon_lL^{-3l}
   <\tfrac12R_\sigma^{3,0}(b_-).
\]
For an incomplete star the current is zero.
The two potential bounds hold since \(A'=0\).
This proves (L1).

The gauges \eqref{II-90-44} prove (L2) with half
the table. For its containing-cube clause use
\eqref{II-90-45} on \(q\) itself:
\(C_{\rm mb}C\varepsilon_l\le a_\flat/2\).
At a tower address the meet \eqref{II-86-16} is
separated and subordinate. The subordinate-sequence
argument in the proof of \cref{II-90-39} therefore
shows that its tower exists and is a real gauge
transform of \(U|_{Z_{\mathcal S}}\).
For \(\ell=\min\{m_{\sigma''}(x),i\}\), the
pointwise estimates give
\[
 |\partial U_q^{\mathfrak r}-1|
       <\kappa_{\rm run}C\varepsilon_\ell L^{-2\ell},
 \qquad
 \eta_q^3|J_q^{\mathfrak r}|
       <\kappa_{\rm run}C\varepsilon_\ell L^{-3\ell}.
\]
Use \eqref{II-90-59}, \(L^{i-p}\ge1\), and
\(L^{-3\ell}\le L^{-2\ell-\min(q,\ell)}\)
to obtain (L3) at half the radii.

For a core address \(X'\subset\Lambda_i\subset\Omega_i\),
all current stars are complete, and
\[
 \rho_{X'}^{s'\to i}
       \mathbf P^{(s')}_{Y|\mathcal S}(\mathsf U)
           =\mathbf P^{(i)}_{X'}(\mathsf U).
\]
Apply \cref{II-90-39} and \eqref{II-90-59} at
layer \(i\); the radius of kind \(1\) is positive
and \(A'=0\). This proves (L4). Every condition
uses the same real factor. Empty index sets impose
no condition, which also covers \(Y_{\mathcal S}
=\emptyset\).
\end{proof}

At step \(t=k+1\) define the fluctuation box
\(\mathsf{Bx}_t(Y;C_{\rm xt};\vartheta)\), for
\(0<\vartheta\le1\), by
\begin{equation}\label{II-90-53}
 \begin{aligned}
 g_k^{-1}|\mathcal A(b)|&<3\vartheta\mathfrak b_k
       &&(b\in\mathfrak X_{\rm I}\cap
                        \mathfrak X^\circ\langle Y\rangle),\\
 g_k^{-1}|\mathcal A(b)|&<2\vartheta C_{\rm xt}\mathfrak b_k
       &&(b\in\mathfrak X_{\rm E}\cap
                        \mathfrak X^\circ\langle Y\rangle).
 \end{aligned}
\end{equation}
Here \(\mathfrak X^\circ\langle Y\rangle\) consists
of the nonpivot fluctuation bonds with both endpoints
in \(Y\). The form is specified at positive bonds
and extended by antisymmetry.

\begin{corollary}\label{II-99-39}
Let \(h'=(h,\mathsf o_t)\), \(t=k+1\), and impose
the support and smallness hypotheses of
\cref{sec:history-support-identity}, together with
\eqref{II-90-11}, \eqref{II-90-12} and
\eqref{eq:history-radius-floor} at depth \(t\),
with \(C=C_{\rm reg}\).
Let \(V\) be a selected bare field in
\(\mathsf{Sup}(h')\cap\operatorname{dom}\mathsf v_h\),
and let \(U_{\rm ad}\) be any adapted element of its
completed collared background. Set
\(\mathsf U=U_{\rm ad}\vee V\).
Every plain triple of \(\mathsf U\) belongs to the
half-radius spaces \eqref{II-90-71}, including the
output and intermediate spaces and the spaces of
earlier levels. Under \eqref{II-90-13}, its triples
on \(X\subset\Omega_j\) also belong to
\(\mathcal U_j^\circ(X;\alpha_0/2,\alpha_1/2)\).

For every \(Y\in\mathcal D_k\), the real point
\begin{equation}\label{II-99-44}
 \mathfrak u_Y=
 \left(\mathbf P^{(t)}_{Y|\mathcal S^+}(\mathsf U),\
       \mathcal A_k|_{\mathfrak X^\circ\langle Y\rangle},\
       \mathfrak q[V]|_Y\right)
 \in
 \tilde U^{1/2}_{(t,{\rm mid})}(Y)
 \times\mathsf{Bx}_t(Y;C_{\rm xt};1/2)
 \times\mathfrak Q_h^{\rm rad}(Y)
\end{equation}
is defined. If \(C_{\mathcal A}\ge2C_{\rm xt}\),
then \(g_k^{-1}\mathcal A_k|_Y\) lies in half the
field polydisc \(\mathbb D^{\rm f}_{x_k}(Y)\).
\end{corollary}

\begin{proof}
On the selected support, the passed density tests
give a completed datum with tested holonomies
smaller than \(2\varepsilon_j\); this is the
regularity argument in
\cref{sec:history-step-support}.
Thus \cref{II-10-20} gives the background and its
layer bounds, and \cref{II-90-39,II-90-62} apply
to its gluing. Cutting a history after level \(j\)
does not change any region, address or radius used
at stages of levels at most \(j\), so it does not
change these earlier spaces. The same assertion
holds with half the table and half \(a_\flat\).

The fluctuation field exists by \cref{II-12-57}
and satisfies
\[
 g_k^{-1}|\mathcal A_k|
       <c_{\rm fl}\mathfrak b_k
       \le C_{\rm xt}\mathfrak b_k .
\]
Passage of the fluctuation tests improves this to
\(g_k^{-1}|\mathcal A_k|<\mathfrak b_k\) on
the integration bonds. These two estimates imply
\eqref{II-90-53} with \(\vartheta=1/2\).
They also give
\(g_k^{-1}|\mathcal A_k|<
\tfrac12C_{\mathcal A}\mathfrak b_k\) under the
stated polydisc inequality.
The old stored variables are defined by
\(V\in\operatorname{dom}\mathsf v_h\).
The support argument in
\cref{sec:history-support-identity}, applied to
each earlier step, puts them in their external
tiers. Restriction only deletes coordinates and
preserves these inequalities. When \(k=0\) this
family is empty. This proves \eqref{II-99-44}.
\end{proof}

\subsection{The regional shift}\label{sec:history-regional-shift}

The fluctuation changes only the level-\(k\) rows touching
\(\Omega_{k+1}\). We apply the decoupled response to the shorter
sequence \(\Omega_0,\ldots,\Omega_k\), whose determining set is
denoted by \(\mathbf B\). Throughout this subsection \(k\ge1\),
lengths are in level-\(k\) units, and
\(\xi_n=L^{n-k}\). Put
\begin{equation}\label{II-92-6}
 W_Y=\Omega_0\cap Y,\qquad
 \mathfrak P^\natural(Y\mid\Omega_\bullet;a)
 =\mathfrak P^\natural_k((\Omega_n)_{n\le k};a,a,a)
                                      \restriction W_Y ,
\end{equation}
where the intersection is cube-wise. The response uses the same
assignment of walk data at every lattice and sequence of a given
height.

In addition to the standing size and room conditions, require
\begin{align}
 R_1&\ge L,\quad \mathfrak M\ge L^4,\quad
 R_1\mathfrak M\ge512\log L/\mathring\delta_0,
                                      \label{II-92-2}\\
 \mathfrak M&\ge M^{\rm pr},\quad
 R_1\mathfrak M\ge\mathfrak r^{\rm pr},\quad
 \mathfrak M\ge(8C_{\rm w}\mathsf n_{1,\natural})^2,
                                      \label{II-92-3}\\
 M&\ge88\mathsf r_{\rm w}\mathfrak M,\qquad
 \mathring\delta_0M\ge2176\bar\kappa .
                                      \label{II-92-4}
\end{align}
These are exactly the frame and walk-expansion requirements of
\cref{I-88-110}. Indeed, truncating a separated sequence preserves
separation. The partition has side \(M\), and its period is at least
\(2L^2M\); hence it has at least three cubes in each direction and
its period is at least \(2L\mathfrak M\).

We fix the constants once, before the heights and couplings. With
the response constants of \cref{I-88-110} and the pivot constants
of \eqref{II-41-44}, set
\begin{equation}\label{II-92-23}
 \begin{gathered}
 C_{\rm b}=\max\{24,5b_{\rm D},10C_jb_{\rm D}^3\},
 \qquad \rho_*=5b_{\rm D}\varepsilon_{\rm D},\\
 C_{\rm e}=\max\{24,L^3,10L^2(1+\rho_*)b_{\rm D},
                            10L^3C_jb_{\rm D}^3\},\\
 C_{\min}=\max\{32,3/a_{\rm D},1/a_P^\circ,1/r_P^\circ,
      2C_P/\varepsilon_{\rm D},2C_Pe^{L/4}C_{\rm b},
                                      (3+2C_P)C_{\rm e}\}.
 \end{gathered}
\end{equation}
Choose \(C\ge C_{\min}\), enlarging it also to preserve the
earlier single-lattice restrictions in \eqref{II-51-27}, and put
\begin{equation}\label{II-92-25}
 \begin{gathered}
 \mu_{\mathfrak h}=\mathring\delta_0/8,\qquad
 \mathsf c_{\mathfrak h}=Ce^{C\bar\kappa},\qquad
 a_{\rm m}=C^{-1}\min\{1,\varepsilon_{\rm v}\},\\
 \varepsilon_{\mathfrak h}=a_{\rm m}e^{-C\bar\kappa},
 \qquad
 a_{\mathfrak h}
   =a_{\rm m}/\max\{1,\mathsf oLM B^{(\mathrm{cg})}\}.
 \end{gathered}
\end{equation}
Thus the previously constructed single-lattice maps are restricted
to smaller domains; their defining formulas are unchanged.
The inequalities we will use are
\begin{equation}\label{eq:history-shift-reserves}
 \begin{gathered}
 3a_{\rm m}\le a_{\rm D},\qquad
 a_{\mathfrak h}\le a_P^\circ,\qquad
 \varepsilon_{\mathfrak h}\le r_P^\circ,\qquad
 2C_P\varepsilon_{\mathfrak h}
                      \le\varepsilon_{\rm D}e^{-32\bar\kappa},\\
 2C_Pe^{L/4}C_{\rm b}e^{C_{\rm b}\bar\kappa}
                  \le\mathsf c_{\mathfrak h},\qquad
 C_{\rm e}\bigl(3a_{\rm m}
           +2C_Pe^{C_{\rm e}\bar\kappa}\varepsilon_{\mathfrak h}\bigr)
                  \le\varepsilon_{\rm v}.
 \end{gathered}
\end{equation}
They follow term by term from \eqref{II-92-23}. In the last
inequality use \(C\ge(3+2C_P)C_{\rm e}\) and
\(e^{(C_{\rm e}-C)\bar\kappa}\le1\).

Let \(\mathfrak X_0\) be the level-\(k\) bonds inside
\(\Omega_0\) with an endpoint in \(B(P'')\), and let
\(\mathfrak X\) be their nonpivot bonds. Every such bond has
both endpoints in \(\Omega_k\), by separation. They are interior
rows of \(\Gamma_k\); in particular their row weight is \(1\).
Moreover
\[
 \mathfrak X_0=\mathfrak X\sqcup
      \{b_0(c),\overline{b_0(c)}:c\in\Gamma\}.
\]
The two endpoints of \(b_0(c)\) lie in \(Y\) precisely when
\(\square(c)\subset Y\). All bonds needed for its pivot formula
then lie in \(Y\). No bond of \(\mathfrak X_0\) is an outward
row of \(\mathbf B\) or joins two sites in an outer endpoint block.

A source cube meets \(\Omega_{k+1}\). Define
\(\operatorname{dom}_{\mathfrak h}(Y\mid\Omega_\bullet)\)
to consist of pairs \((P,\mathcal A)\), with
\(P=(U,A',\mathbf J)\) a level-\((k+1)\) triple on \(Y\),
which satisfy the following conditions:
\begin{enumerate}
\item \(\sup_{\mathfrak X^{\circ+}\langle Y\rangle}
                 |\mathcal A|<\varepsilon_{\mathfrak h}\).
\item The pair and its specified decomposition
\begin{equation}\label{II-92-21}
 \mathbf p^{[k]}(P)
  =(e^{i\eta_{k+1}A'}U,L^{-3}\mathbf J)
                           \big|_{\operatorname{Bd}^\circ(W_Y)},
 \qquad (U,L^{-1}A'),
\end{equation}
serve in the natural local class at radius \(a_{\rm m}\).
\item For every \(Y_1\subset Y\) all of whose cubes are source
cubes,
\(P|_{Y_1}\in\mathcal Q_{k+1}
           (Y_1;a_{\mathfrak h},a_{\mathfrak h},a_{\mathfrak h})\).
\end{enumerate}
Complete \(\mathcal A\) at the pivots by the formula
\eqref{II-51-33}, and extend the resulting field by zero to
all other rows of \(\mathbf B\). Denote it by \(B'_Y\).
The regional shift is
\begin{equation}\label{II-92-29}
 \mathfrak h_{Y\mid\Omega_\bullet}(s;P,\mathcal A)
   =(\mathbf A,\mathbf K)
   =(\mathbf H,j)\bigl(s;\mathbf p^{[k]}(P);B'_Y\bigr)
                                          \big|_{\operatorname{Bd}(Y)} ,
\end{equation}
where both forms are extended by zero off the inner bonds of
\(\Omega_0\).

\begin{proposition}\label{II-93-26}
The map \eqref{II-92-29} is defined on
\(\mathbb D(Y)\times
       \operatorname{dom}_{\mathfrak h}(Y\mid\Omega_\bullet)\).
It depends on \(P\) only through its pair, is holomorphic on
each open fibre with \(U\) fixed, and is continuous on its
Borel domain. The domain is hereditary under restriction and
invariant under unitary gauges. The shift has the corresponding
restriction, gauge and conjugation covariance.

Let \(x\) be a bare site of \(\Omega_0\), and put
\(n=\max\{j\le k:x\in\Omega_j\}\). At a positive bond based at
\(x\), or a difference stencil based there and contained in
\(\operatorname{Bd}(Y)\cap\operatorname{Bd}^{\circ}(\Omega_0)\),
one has
\begin{equation}\label{II-93-35}
 \begin{split}
 \xi_n|\mathbf A|,\quad
 \xi_n^2|\nabla_{\mathbf U}^{\eta_k}\mathbf A|,\quad
 \xi_n^3|\mathbf K|
 \ \le\ \mathsf c_{\mathfrak h}
   \sup_{b\in\mathfrak X^{\circ+}\langle Y\rangle}
        e^{-\mu_{\mathfrak h}d_{\mathbf B}(x,b_-)}
                                     |\mathcal A(b)|.
 \end{split}
\end{equation}
For a union of wall components of the parameter support, the
corresponding component response depends only on the pair and
source within that union and is supported on its inner bonds.
It vanishes if the union contains no source cube. In particular,
the shift vanishes at \(s=0\) or \(\mathcal A=0\).

If two sequences have the same traces in \(Y\), including the
last trace, their shift domains and maps on \(Y\) agree. The
same statement holds after a fine translation or a chart
which identifies these traces. If
\begin{equation}\label{II-93-44}
                  12\mathfrak M\le\mathsf oL^2M ,
\end{equation}
then on a constant sequence the domain is exactly
\eqref{II-41-25} with the constants \eqref{II-92-25}, and the
map is the retained single-lattice shift.
\end{proposition}

\begin{proof}
For \(c\in\Gamma\), both endpoint fine cubes of \(\square(c)\)
are source cubes. To see this when only one endpoint is in
\(\Omega_{k+1}\), the midpoint of \(c\) lies in the level-\((k+1)\)
cube of that endpoint, hence in \(\Omega_{k+1}\), and also in
the fine cube of the other endpoint. The last domain condition
therefore provides the pivot data on \(\square(c)\). The pivot
estimate gives
\begin{equation}\label{II-93-28}
 |\beta_c|\le2C_P
       \max_{b\in\mathfrak X^+\cap\mathsf P_c}|\mathcal A(b)|,
 \qquad
          \sup|B'_Y|\le2C_P\sup|\mathcal A|.
\end{equation}
The set in the maximum is nonempty: a bond inside the endpoint
block belonging to \(P''\) is a nonpivot bond of \(\mathsf P_c\).
The first reserve in \eqref{eq:history-shift-reserves} now
admits the completed source into \cref{I-88-110}; the second
domain condition admits its pair. This proves existence.

The pivot depends only on the pair in its double-block tower
and on the nonpivot values in \(\mathsf P_c\). This tower is
inside \(W_{\square(c)}\), by the collar separation.
Consequently it restricts consistently. The two contractions
in \cref{I-88-110} are also independent of the decomposition
of the pair. The hereditary assertion for the natural class
includes its gauge witnesses: on restriction to \(W_{Y_1}\),
replace a containing witness cube for a natural cube
\(\mathsf q\subset W_{Y_1}\) by \(\mathsf q\) itself.
The pointwise estimates restrict, and every source subdomain
of \(Y_1\) was already tested in \(Y\). This proves heredity.

At a fixed \(U\), the conditions on its gauge witnesses are
fixed and the other domain conditions are finitely many
strict continuous inequalities. Its fibre is open, and the
holomorphic pivot formula followed by \cref{I-88-110} proves
holomorphy. To check the Borel assertion as \(U\) varies,
fix a witness cube and replace a strict logarithm bound
\(r\) by \(r(1-1/q)\), \(q\in\mathbb N\).
The allowable real forms and unitary gauges then form a
compact set. Its continuous image under
\((a,u)\mapsto(e^{i\eta a})^{u^{-1}}\) is compact.
There are finitely many cubes; finite intersections and
unions, followed by this countable exhaustion for strict
bounds, prove that the witness conditions are Borel.

For continuity at \(P_0=(U_0,A'_0,\mathbf J_0)\), freeze
the real factor \(U_0\). On the finitely many bonds used by
the response and the pivot towers put, near \(P_0\),
\[
 \Lambda(P)=(i\eta_{k+1})^{-1}
                 \log(\mathbf U U_0^{-1}).
\]
Here \(\eta_{k+1}|A'_0|<1/32<\log2\), so the logarithm
is defined near the point and equals \(A'_0\) there.
The fixed-factor fibres are open. Thus the continuous
map \(P\mapsto(U_0,\Lambda(P),\mathbf J)\) stays in them
near \(P_0\); it preserves the pair and hence the pivot
and response. Their fixed-factor holomorphy proves the
claimed continuity.

For the quantitative estimate, let \(y(x)\) be the top
of the pillar containing \(x\). Then
\(d_{\mathbf B}(x,y(x))\le3/2\).
The weighted bounds in \cref{I-88-110}, first at \(y(x)\)
and then at \(x\), are bounded by
\(5b_{\rm D}e^{8\bar\kappa}\) for the field and its first
difference, and by \(10C_jb_{\rm D}^3e^{24\bar\kappa}\)
for the current, times
\[
 S_x(B'_Y)=
       \sup_{c}e^{-\mu_{\mathfrak h}d_{\mathbf B}(x,c_-)}
                                      |B'_Y(c)|.
\]
We used \(4e^{3/16}<5\) and \(8e^{3/16}<10\).
For a pivot row \(c'=b_0(c)\), select a bond \(b\) attaining
the maximum in \eqref{II-93-28}. The two initial points
are in \(\Omega_k\), at weighted distance at most \(2L-1\).
Hence
\[
 S_x(B'_Y)\le2C_Pe^{L/4}
   \sup_{b\in\mathfrak X^{\circ+}\langle Y\rangle}
          e^{-\mu_{\mathfrak h}d_{\mathbf B}(x,b_-)}
                                       |\mathcal A(b)|.
\]
The bound for nonpivot rows is immediate. The reserve
\eqref{eq:history-shift-reserves} proves \eqref{II-93-35}.
At an incident differentiation stencil the same estimate
is valid: if its base lies in the third component margin,
the entire stencil lies in the pair domain. Otherwise
both field values vanish by the response support property.
In this second case the connecting transport can be arbitrary.

The component support and restriction proof of
\cref{II-51-52} applies to this frame with its layer
weights retained. The local \(8L\xi_y\) cube lies in
the cell margin because
\(4L\xi_y\le4L^3\xi_m\le2\mathfrak M\xi_m\).
Thus the component's pivot data and pair lie inside that
component. If it has no source cube, it has no row of
\(\mathfrak X_0\), and the component datum is zero.
The two response iterations then vanish there.
This proves the support and zero assertions.
The pivot formula and both response iterations commute
with unitary gauges and complex conjugation.

For two sequences with matching traces in \(Y\), the
layer regions, natural cubes contained in \(Y\), source
cubes, pivot indices and nonpivot bonds inside \(Y\)
coincide. Thus the three domain conditions coincide.
Their local pair and completed data also coincide.
The local comparison in \cref{I-88-110}, applied to
each wall component, identifies the responses. The
same comparison applies under a chart: its local
translations carry the fine cubes, levels, walks and
double-block formulas to their counterparts. All
lengths, inequalities and gauge witnesses are preserved,
so this is equality of the transported domains as well
as the maps. The halo contains the neighborhoods used
in these comparisons: the natural-cube enlargement
has width \(5\mathfrak M\xi_n<LM\), and the layer cube
meeting \(Y\) has side \(\mathfrak M\xi_n\le M\).

For a constant sequence all fine cubes are source cubes.
The last domain condition at \(Y_1=Y\) is therefore
necessary. Conversely, the simple domain on \(Y\)
implies all its restrictions. It also implies the
natural-class condition: curvature acquires a factor
\(L^{-2}\), the potential and its derivative acquire
factors \(L^{-1}\) and \(L^{-2}\), and the current
acquires \(L^{-3}\). The natural cubes have side
at most \(12\mathfrak M/L\le\mathsf oLM\) in
level-\((k+1)\) units, so their gauges are supplied by
the simple domain with bound
\(\mathsf oLM B^{(\mathrm{cg})}a_{\mathfrak h}
\le a_{\rm m}\). This proves exact domain equality.
The completed sources and the response formulas are
then those of \eqref{II-51-35}.
\end{proof}

\subsection{Gauges on source-cube domains}
\label{sec:source-cube-domain}

The local gauges in (L2) are given on cubes at their own
scale. A source domain at step \(t=k+1\) can contain larger
cubes which cross the boundary of \(\Omega_t\). We first
glue the level-\(k\) gauges on such a cube. Alignment with
one tree gauge makes the transition functions small, and
three coordinate interpolations avoid a loss proportional
to the number of small boxes.
We call the first and second gauge requirements in (L2)
of \cref{II-86-20} its parts (a) and (b), respectively.

Put
\[
 \varsigma_{\mathsf o}=\mathsf oLM,\qquad
 \varkappa[a]=\varsigma_{\mathsf o}B^{(\mathrm{cg})}a,
 \qquad
 q_*=\frac{13L^2}{24B^{(\mathrm{cg})}}+\frac72L .
\]
Choose the constant in the cube-gauge ceiling to be
\begin{equation}\label{II-94-65}
 \mathsf K_{\mathrm{cg}}=\mathsf K_{\rm I}
 =\max\left\{16L^2,\,
      180q_*\max\left(1,\frac{13L}{\varsigma_{\mathsf o}}\right)
       \right\}.
\end{equation}
It depends only on the already fixed geometric constants,
not on the step, volume or couplings. Thus it is chosen
after \(L,\mathsf o,M,B^{(\mathrm{cg})}\), before the
radius ceilings and the choice of \(\gamma\).
For every \(0<\varkappa\le1\), this choice satisfies
\begin{equation}\label{II-94-63}
 \mathsf K_{\rm I}\ge7q_*,
 \qquad
 \frac{5.4L^2}{\mathsf K_{\rm I}}
 +9679\varkappa\frac{q_*^2}{\mathsf K_{\rm I}^2}
 +660.5\frac{Lq_*}{\varsigma_{\mathsf o}\mathsf K_{\rm I}}
 \le0.99,
\end{equation}
and
\begin{equation}\label{II-94-64}
 \varsigma_{\mathsf o}<24/L
 \quad\Longrightarrow\quad
 \mathsf K_{\rm I}\ge
       \frac{16\pi L^3}{\varsigma_{\mathsf o}B^{(\mathrm{cg})}}.
\end{equation}

\begin{proposition}\label{II-94-90}
Let \(L\ge13\), \(2\le t\le\bar k\), \(k=t-1\), and
\(\sigma=(t,\mathrm{mid})\) or \((t,\mathrm{out})\).
Use the collared sequence and separated label of
\cref{II-86-20}, with \(m'<m+K-\bar k\),
\((\ell,K)\in\mathfrak W_k\), and the two gaps
\[
 \operatorname{dist}_\infty(\Lambda_j,\Omega_j^c)
       \ge2M\check R_j,\qquad
 \operatorname{dist}_\infty(Z_j,\Omega_{j+1})
       \ge8LM\check R_{j+1}
\]
in level-\(j\) units. Let \(Y\in\mathcal D_k\) and
\(P=(U,A',\mathbf J)\in\tilde U_\sigma(Y)\). Suppose
\begin{equation}\label{eq:source-cube-ceilings}
 \begin{gathered}
 (1+\beta_{\rm s})r_\iota(\sigma;n)\le a_{\rm m}
       \quad(0\le n\le t,\ \iota=0,1),\qquad
 a_\flat\le a_{\rm m},\\
 \max\{\mathsf K_{\mathrm{cg}},L^3\}
       (1+\beta_{\rm s})r_0(\sigma;n)\le a_{\mathfrak h},
 \qquad
 L^3(1+\beta_{\rm s})r_1(\sigma;n)\le a_{\mathfrak h}
       \quad(n=k,t).
 \end{gathered}
\end{equation}
Here \(a_{\rm m},a_{\mathfrak h}\) are the radii in
\cref{II-92-25}; in particular
\(\varkappa[a_{\mathfrak h}]\le1\).
Then the pair \eqref{II-92-21}, with its specified
decomposition, belongs to the natural local class of radius
\(a_{\rm m}\) on \(Y_{\mathcal S}\). Moreover,
\begin{equation}\label{eq:source-cube-simple-domain}
 P|_{Y_1}\in
       \mathcal Q_t(Y_1;a_{\mathfrak h},a_{\mathfrak h},
                                  a_{\mathfrak h})
\end{equation}
for every \(Y_1\in\mathcal D_k\), \(Y_1\subset Y\), all of
whose fine cubes meet \(\Omega_t\).
Consequently \((P,\mathcal A)\) lies in
\(\operatorname{dom}_{\mathfrak h}(Y\mid\Omega_\bullet)\)
whenever \(\sup|\mathcal A|<\varepsilon_{\mathfrak h}\).
At \(\sigma=(t,\mathrm{mid})\), the inequalities
\eqref{eq:source-cube-ceilings} are precisely the selected
cube and natural-class ceilings of \((\mathrm C6)_t\).
\end{proposition}

Here is the gluing estimate used for cubes crossing the
last interface. All its differences are ordinary differences
on integer boxes.

\begin{lemma}\label{II-94-80}
Let \(U\) be an \(\mathrm{SU}(2)\) configuration on the
inner bonds of \(\square_N=\{0,\ldots,N-1\}^3\), and put
\(\varpi=\max_p|\partial U(p)-1|\). Let \(n_{\rm p}\ge24\)
be an integer. Suppose every integer box in \(\square_N\)
with side lengths at most \(n_{\rm p}\) has a gauge \(W\)
such that
\[
 \max_b|W(b)-1|\le\gamma_1,\qquad
 \max_{x,\mu,\nu}|W_\nu(x+e_\mu)-W_\nu(x)|\le\gamma_2.
\]
The second maximum is over the stencils with
\(x,x+e_\mu,x+e_\nu,x+e_\mu+e_\nu\) in the box, including
\(\mu=\nu\); empty maxima are zero. If
\[
 z_*=\frac{13}{24}N\varpi+\frac72\gamma_1,\qquad
 n_{\rm p}z_*\le\frac1{14},
\]
then \(U\) has on \(\square_N\) a gauge of sizes
\begin{equation}\label{eq:source-cube-gluing-bound}
 \left(6.55z_*,\
       5.4\gamma_2+433z_*^2+\frac{634}{n_{\rm p}}z_*\right).
\end{equation}
\end{lemma}

\begin{proof}
Write \(\gamma_1^\mu(W)\) for the first size restricted to
direction \(\mu\). We give the interpolation calculation
before constructing the gauges.

For hermitian \(\theta,\theta^\flat\) of norm at most
\(\vartheta<1\), and \(0\le s,s'\le1\), set
\[
 D=\theta-\theta^\flat,\quad
 Z_s=e^{is\theta}e^{-is\theta^\flat},\quad
 E_s=Z_s-1-s(Z_1-1).
\]
The identity
\[
 Z_s-Z_{s'}=i\int_{s'}^s
       e^{i\tau\theta}D e^{-i\tau\theta^\flat}\,d\tau
\]
and unitarity give
\[
 |Z_s-Z_{s'}|\le|s-s'||D|,\quad
 |D|\le\frac{|Z_1-1|}{1-\vartheta},\quad
 |E_s|\le\frac34\vartheta|D|,\quad
 |E_s-E_{s'}|\le3\vartheta|D|\,|s-s'|.
\]
For two such pairs, indexed by \(1,2\), put
\[
 N_{12}=2\vartheta|D_1-D_2|
       +(2|\theta_1-\theta_2|+|D_1-D_2|)|D_2|.
\]
Subtracting the integrands, and subtracting their values
at \(\tau=0\), gives respectively
\begin{align*}
 |Z_s^{(1)}-Z_s^{(2)}|
 &\le s|D_1-D_2|
       +\frac{s^2}{2}
          (2|\theta_1-\theta_2|+|D_1-D_2|)|D_2|,\\
 (1-\vartheta-\tfrac12|D_2|)|D_1-D_2|
 &\le |Z_1^{(1)}-Z_1^{(2)}|
                     +|\theta_1-\theta_2||D_2|,\\
 |E_s^{(1)}-E_s^{(2)}|&\le\tfrac38N_{12}.
\end{align*}
Indeed the remainder in the integrand is bounded by
\(2\tau\vartheta|D|\), and its difference by \(\tau N_{12}\).
In \(E_s\) the constant part cancels; the remaining
integrals have total coefficient
\(\frac12s(1-s)(1+2s)\le3/8\). These observations also
give the four preceding bounds.

Now let \(W^1=(W^0)^h\), where \(|h-1|\le\rho\le1/10\),
and put \(\theta=\log(h)/i\),
\(\vartheta=\pi\rho/3\). The logarithm is the real
principal branch, with \(|\theta|\le\vartheta\).
For a cutoff \(\chi\in[0,1]\), let
\(W^\chi=(W^0)^{e^{i\chi\theta}}\). At a bond in direction
\(\mu\), set
\[
 \theta^\flat=W^0_\mu(x)\theta(x+e_\mu)W^0_\mu(x)^{-1},
 \quad s=\chi(x),\quad
 \delta=\chi(x+e_\mu)-\chi(x),\quad
 \Xi=e^{-i\delta\theta^\flat}.
\]
Then
\begin{equation}\label{II-94-76}
 \begin{split}
 W^\chi_\mu
  &=(1-s)W^0_\mu+sW^1_\mu
             +E_sW^0_\mu+Z_s(\Xi-1)W^0_\mu,\\
 Z_1W^0_\mu&=W^1_\mu .
 \end{split}
\end{equation}
Suppose both \(W^j\) have directional sizes \(g_1^\mu\)
and second size \(g_2\), and write \(g_1=\max_\mu g_1^\mu\).
If \(k_1,k_2\) bound the first and second differences of
\(\chi\), put
\[
 d=\frac{2g_1}{1-\vartheta},\quad
 t_1=2g_1\vartheta+d,\quad
 d_2=\frac{2g_2+t_1d}{1-\vartheta-d/2},\quad
 e=2\vartheta d_2+(2t_1+d_2)d .
\]
For a contained difference stencil the preceding estimates
give
\[
 |D|\le d,\quad |\Delta_\nu\theta|\le t_1,\quad
 |\Delta_\nu D|\le d_2,\quad
 |\Delta_\nu\theta^\flat|\le t_1+d_2 .
\]
Here the first bound uses \(Z_1=W^1_\mu(W^0_\mu)^{-1}\);
the second uses unitary conjugation and
\(|VTV^{-1}-T|\le2|V-1||T|\); the third uses the
two-pair estimate and \(|\Delta_\nu Z_1|\le2g_2\).
Thus the three terms in \eqref{II-94-76}, grouped as its
affine part, \(E_sW^0_\mu\), and \(Z_s(\Xi-1)W^0_\mu\),
have difference bounds
\begin{align*}
 &g_2+2k_1g_1,\\
 &3k_1\vartheta d+\tfrac38e+\tfrac34\vartheta d\,g_2,\\
 &k_1\vartheta\{k_1d+d_2+\tfrac12(2t_1+d_2)d\}
       +k_2\vartheta+k_1(t_1+d_2)+k_1\vartheta g_2.
\end{align*}
The last line follows by differencing its three factors;
\(|\Delta_\nu\Xi|\le k_2\vartheta+k_1(t_1+d_2)\).
Its undifferenced counterpart is
\[
 |W^\chi_\mu-1|
 \le\left(1+\frac{3\vartheta}{2(1-\vartheta)}\right)
                          g_1^\mu+|\Delta_\mu\chi|\vartheta .
\]

To glue a row, put \(\omega=\lfloor n_{\rm p}/2\rfloor\)
and \(s_{\rm p}=n_{\rm p}-\omega\). Cover
\(\{0,\ldots,N-1\}\) by the intervals
\[
 I_j=\{js_{\rm p},\ldots,\min(js_{\rm p}+n_{\rm p},N)-1\},
\]
stopping at the first interval reaching \(N-1\).
Successive overlaps \(O_j\) contain \(\omega\) integers
and are disjoint; non-successive intervals are disjoint.
If \(o_j=(j+1)s_{\rm p}\), use on \(O_j\)
\[
 \chi_j(c)=\operatorname{ct}
                \left(\frac{c-o_j-1}{\omega-3}\right),
 \qquad
 \operatorname{ct}(s)=
 \begin{cases}
 1,&s\le0,\\
 1-3s^2+2s^3,&0\le s\le1,\\
 0,&s\ge1.
 \end{cases}
\]
For gauges \(v_j\) on these boxes, interpolate by
\(e^{i\chi_j\log(v_jv_{j+1}^{-1})/i}v_{j+1}\).
The cutoff is constant on the first and last two layers
of each overlap. The constant-gauge and interpolation
regions therefore overlap in two layers: every bond
and every difference stencil, whose coordinate range
is at most two, is contained in one such region.
No boundary stencil is omitted.

The derivative bounds of this cubic cutoff give
\[
 k_1\le\frac{25}{6n_{\rm p}},\qquad
 k_2\le\frac{46.3}{n_{\rm p}^2}.
\]
For \(g_1\le0.042\) and \(\rho\le1/10\), we have
\(\vartheta\le0.10472\), \(d\le2.234g_1\),
\(t_1\le2.4435g_1\), and
\(d_2\le2.3575g_2+6.4348g_1^2\); the denominator
is at least \(0.84836\).
Substitution in the displayed three difference bounds
gives, for gluing in direction \(1\),
\begin{equation}\label{eq:source-cube-row-gluing}
 \begin{aligned}
 \gamma_1^1&\le1.18g_1^1+4.364\rho/n_{\rm p},&
 \gamma_1^\mu&\le1.18g_1^\mu\quad(\mu=2,3),\\
 \gamma_2&\le1.75g_2+6.2g_1^2
              +21.7g_1/n_{\rm p}+48.49\rho/n_{\rm p}^2 .
 \end{aligned}
\end{equation}
The other directions obey the same bounds.

It remains to align the gauges so that this interpolation
can be performed three times. The tree paths from
\((\lfloor N/2\rfloor,\lfloor N/2\rfloor,\lfloor N/2\rfloor)\),
taken in coordinate order, give a gauge \(W^{\rm c}\)
with \(\gamma_1(W^{\rm c})\le N\varpi\). Indeed its third
component is \(1\), its second component is \(1\) on the
central third-coordinate plane, and its first component
is \(1\) on the central coordinate line. The \(23\), \(13\)
and then \(12\) plaquettes telescope to bounds
\(|x_3-c_0|\varpi\) and
\((|x_2-c_0|+|x_3-c_0|)\varpi\).

Take the product boxes \(Q_i=I_{i_1}\times I_{i_2}\times I_{i_3}\).
Multiply each given gauge \(u_i\) by a constant so that the
new gauge \(v_i\) agrees with the tree gauge at the central
site \(x_i\) of \(Q_i\). The relative gauge
\(m_i=v_i(u^{\rm c})^{-1}\) has increments bounded by
\(\gamma_1+N\varpi\). For neighboring boxes choose \(z\)
at the first overlap coordinate and the two shared central
coordinates. Its distances from the two centers in the
varying direction are at most \(1,\omega\). The transition
\(h=v_iv_{i'}^{-1}\) has increments at most \(2\gamma_1\).
Every overlap point is at distance at most
\(\omega-1+n_{\rm p}\) from \(z\), whence
\[
 |h-1|
 \le(\omega+1)N\varpi+
              (3\omega-1+2n_{\rm p})\gamma_1
 \le n_{\rm p}z_*=: \rho_1 .
\]

Glue in directions \(1,2,3\), in that order. The transitions
between neighboring glued rows remain small. To see this,
write \(h_1=e^{i\theta}h_0e^{-i\theta'}\). With
\(\theta^\flat=h_0\theta'h_0^{-1}\), the earlier identity
for \(E_s\) gives
\[
 |e^{is\theta}h_0e^{-is\theta'}-1|
 \le\max(|h_0-1|,|h_1-1|)
       +\frac{3\vartheta}{4(1-\vartheta)}|h_1-h_0|.
\]
The cutoffs in neighboring rows are identical. Thus their
new transition bounds are
\(\rho_2=1.1755\rho_1\) and
\(\rho_3=1.1755\rho_2<1/10\), since \(\rho_1\le1/14\).
This proves at each stage that the next logarithms exist.

Since \(\gamma_1\le(2/7)z_*\) and \(z_*\le1/336\),
the directional first sizes after the successive stages,
in units of \(z_*\), are bounded by
\[
 (4.702,0.338,0.338),\quad
 (5.549,5.529,0.399),\quad
 (6.548,6.525,6.502).
\]
In particular each input size is less than \(0.042\).
Using these three input bounds and
\(\rho_j/n_{\rm p}\le(1,1.1755,1.3819)_jz_*\)
in \eqref{eq:source-cube-row-gluing}, the successive second
sizes are at most
\[
 \begin{aligned}
 &1.75\gamma_2+0.507z_*^2+54.7z_*/n_{\rm p},\\
 &3.0625\gamma_2+137.97z_*^2+254.77z_*/n_{\rm p},\\
 &5.36\gamma_2+432.4z_*^2+633.3z_*/n_{\rm p}.
 \end{aligned}
\]
These bounds imply \eqref{eq:source-cube-gluing-bound}.
Only three interpolations were used, irrespective of the
number of product boxes.
\end{proof}

\begin{proof}[Proof of \cref{II-94-90}]
First consider the pair on \(W=Y_{\mathcal S}\), in
level-\(k\) units. Its potential and current are
\(L^{-1}A'\) and \(L^{-3}\mathbf J\). At a stencil of
the shorter sequence's layer \(n\le k\), its base point
has \(n'=m_\sigma(x)\ge n\). With
\(\xi=\eta_k\), \(\xi_n=L^{n-k}\), condition (L1) gives
\[
 \begin{gathered}
 |\partial U-1|,\ |\partial\mathbf U-1|
               <a_{\rm m}\xi^2\xi_n^{-2},\\
 \xi_n|L^{-1}A'|<a_{\rm m},\qquad
 \xi_n^2|\nabla_U^{\eta_k}(L^{-1}A')|<a_{\rm m},\qquad
 |L^{-3}\mathbf J|<a_{\rm m}\xi_n^{-3}.
 \end{gathered}
\]
The extra factors \(L^{-(n'-n)}\), with powers \(1,2,3\),
are at most one. For each natural cube, (L2)(b) supplies
its containing gauge cube inside \(W\). Restrict that gauge
to the original cube and change the potential from spacing
\(\eta_n\) to \(\eta_k\): the factors are \(\xi_n^{-1}\)
and \(\xi_n^{-2}\) for the potential and its difference.
The resulting index-\(n\) bounds are at most
\(a_\flat\le a_{\rm m}\). These are all the natural local
class conditions, with the specified decomposition.

Now let \(Y_1\) be as in the statement. An endpoint \(x\)
of a bond in \(\operatorname{Bd}(Y_1)\) is within
\(M+\eta_k\) of a point of \(\Omega_t\) in level-\(k\)
units. The second gap therefore excludes \(x\in Z_k\).
It follows that
\[
 x\in\Lambda_k,\quad m_\sigma(x)\in\{k,t\},\qquad
 Y_1\subset Y\cap\Omega_k^\diamond,\quad
 \operatorname{Bd}(Y_1)\subset\operatorname{Bd}_{\mathcal S}(Y).
\]
For \(a=a_{\mathfrak h}\), the \(L^3\) ceilings and
\(c_\sigma(n)\le1+\beta_{\rm s}\) give
\[
 R_\sigma^{p,\iota}(x)\le a\eta_t^p
      \quad(p=1,2,3)
\]
at the relevant stencils. Hence, on \(\operatorname{Bd}(Y_1)\),
\[
 |\partial U-1|,\ |\partial\mathbf U-1|<a\eta_t^2,
 \qquad |A'|,\ |\nabla_U^{\eta_t}A'|,\ |\mathbf J|<a .
\]
The only remaining simple-domain condition is a real gauge
on every cube \(\square\subset Y_1\) of side at most
\(\varsigma_{\mathsf o}\) in level-\(t\) units.

A closed bare cube has \(N^3\) sites, coordinatized by
\(\square_N\), and side \(N\eta_j\) in level-\(j\) units.
The chart is injective and its side is at most one bare
spacing less than the torus side. If two chart sites are
neighbors modulo the torus period, their coordinate
difference minus the prescribed unit step has modulus
strictly less than that period in every coordinate; it
must therefore be zero. Thus inner bonds, plaquettes and
all difference stencils correspond exactly to those of
\(\square_N\), including the repeated-direction stencils.
Put \(N'=\lfloor\varsigma_{\mathsf o}L^k\rfloor\).
Then
\[
 N\le\varsigma_{\mathsf o}L^t<L(N'+1).
\]
For \(f=\eta_t A\), the desired bounds are exactly
\begin{equation}\label{II-94-83}
 \sup|f|<\eta_t\varkappa[a],\qquad
 |\Delta_\mu f_\nu|<\eta_t^2\varkappa[a].
\end{equation}
This follows from
\(\Delta_\mu f_\nu=\eta_t^2\nabla_\mu^{\eta_t}A_\nu\).

Condition (L2)(a) at level \(k\) supplies gauges on every
integer box of side lengths at most \(N'\): enlarge it to
a cube of side \(\min(N,N')\) inside \(\square_N\), apply
(L2), and restrict. Their sizes are bounded by
\[
 \gamma_1=\eta_k\varkappa[r_0(\sigma;k)],\qquad
 \gamma_2=\eta_k^2\varkappa[r_0(\sigma;k)] .
\]
The curvature bound in (L1), and the
\(\mathsf K_{\mathrm{cg}}\) ceiling, give
\[
 \gamma_1\le\eta_t\frac{L\varkappa[a]}{\mathsf K_{\rm I}},
 \quad
 \gamma_2\le\eta_t^2\frac{L^2\varkappa[a]}{\mathsf K_{\rm I}},
 \quad
 \varpi<\frac{aL^2}{\mathsf K_{\rm I}}\eta_t^2 .
\]

If \(N\le N'\), the original level-\(k\) gauge suffices:
rescaling multiplies its potential by \(L\) and its first
difference by \(L^2\), while \(\mathsf K_{\rm I}\ge L^2\).
If \(N>N'\) and \(N'\le23\), then \(N<24L\).
In this case \(\varsigma_{\mathsf o}<24/L\) and
\eqref{II-94-65} implies
\(\mathsf K_{\rm I}\ge16\pi L^3/
                  (\varsigma_{\mathsf o}B^{(\mathrm{cg})})\).
For \(N=1\) there are no bonds. Otherwise the tree gauge
used above has
\[
 \gamma_1(U^u)\le N\varpi
       <\frac3{2\pi}\varkappa[a]\eta_t^2 .
\]
Its principal logarithm has size
\(<\frac12\varkappa[a]\eta_t^2\); the difference of two
values has size \(<\varkappa[a]\eta_t^2\). This proves
\eqref{II-94-83} in the small-box case.

It remains to consider \(N>N'\ge24\). The gluing parameter
satisfies
\begin{equation}\label{II-94-88}
 0<z_*\le\eta_t\varkappa[a]\frac{q_*}{\mathsf K_{\rm I}}
       \le\frac{\eta_t}{7}\le\frac1{1183}.
\end{equation}
Take
\[
 n_{\rm p}=\min\left\{N',
                       \left\lfloor(14z_*)^{-1}\right\rfloor\right\}.
\]
Then \(n_{\rm p}\ge24\), \(n_{\rm p}z_*\le1/14\), and
\[
 n_{\rm p}^{-1}
       \le\frac{25}{24}
                 \left(14z_*+\frac{L\eta_t}{\varsigma_{\mathsf o}}\right).
\]
Applying \cref{II-94-80}, the final first size is
\(<0.0056\), and the second size is at most
\[
 \eta_t^2\varkappa[a]\left[
   \frac{5.4L^2}{\mathsf K_{\rm I}}
  +9679\,\varkappa[a]\frac{q_*^2}{\mathsf K_{\rm I}^2}
  +660.5\,\frac{Lq_*}{\varsigma_{\mathsf o}\mathsf K_{\rm I}}
                    \right]
 \le0.99\,\eta_t^2\varkappa[a].
\]
Indeed \(\varkappa[a]\le1\), and the three terms in
brackets are at most \(0.3375,0.2988,0.2823\).
The same choice gives \(\mathsf K_{\rm I}\ge7q_*\)
and the small-box inequality used above.
The unitary logarithm bounds
\[
 |\log V|\le\frac\pi3|V-1|,\qquad
 |\log V-\log V'|\le
         \frac{|V-V'|}{1-\max(|V-1|,|V'-1|)}
\]
now yield \eqref{II-94-83}: its first constant is at most
\(6.86/7<1\), and its second is at most
\(0.99/(1-0.0056)<1\).
The logarithm is real and its reversal is its negative,
so it defines the required real one-form on the inner
bonds of \(\square\).

All the pointwise and cube-gauge conditions of the simple
set have now been proved. Since \(Y_1\) was arbitrary,
\eqref{eq:source-cube-simple-domain} follows. Together with
the natural-class conclusion this gives the two background
conditions of the regional shift domain; its remaining
condition is the asserted strict bound on \(\mathcal A\).
\end{proof}

\subsection{Identification of the regional shift}
\label{sec:history-regional-shift-identification}

At full coupling, the response must use the reduced
current of the regional action, whereas the stored
triple sets the current to zero at incomplete stencils.
The following argument justifies that substitution.

\begin{lemma}\label{II-93-23}
At \(Y=T_\ell\) and \(s=\mathbf1\), let two admissible
pairs have the same first component and currents differing
by \(\delta\), with \(\delta=0\) on every bond touching
\(\Omega_1\). Their responses have the same first component
\(\mathbf H\), and their reduced-current increments differ by
\begin{equation}\label{II-93-24}
                        f_{\mathbf H}^{-1}\delta-\delta .
\end{equation}
If the first pair is \((U_0,J^\circ(U_0))\), with \(U_0\)
real and critical, and the datum is supported on
\(\mathfrak X_0\), then both increments are equal.
\end{lemma}

\begin{proof}
Let \(\Gamma_0^+\) be the bare determining rows, namely
the inner bonds with no endpoint in \(\Omega_1\).
On these rows
\(\mathcal Q(A)=\xi A\); hence the nonlinear averaging
remainder \(\mathsf C\) and all its derivatives vanish.
The correction \(\Delta'_\pi[\delta]\) also vanishes.
Indeed, the site field \(v=\mathsf VA\) in
\eqref{eq:bg-current-form} has zero scalar averages and
therefore vanishes at the bare sites outside \(\Omega_1\).
On a bond supporting \(\delta\), both \(v_-\) and \(v_+\)
are zero, so every term in that quadratic form is zero.
Polarization proves the assertion.

Consequently \(G,S,H\) are unchanged when the current
changes by \(\delta\). Since \(QH=1\), on a bare row
\((HX)(c)=\xi^{-1}X(c)\). With
\(\zeta=H^{\mathsf T}\delta\), the weighted row pairing
therefore gives
\[
 \zeta(c)=
 \begin{cases}\xi^3\delta(c),&c\in\Gamma_0^+,\\
                0,&\text{otherwise},\end{cases}
 \qquad Q^{\mathsf T}\zeta=\delta .
\]
It also gives
\(\mathsf C'_T{}^{\mathsf T}\zeta
=\mathsf C''[A']^{\mathsf T}\zeta=0\).
In particular
\(\Psi_T\zeta=(1+H^{\mathsf T}\mathsf C'_T{}^{\mathsf T})^{-1}
\zeta=\zeta\).
The quadratic form defining the additional Hessian
correction changes by
\(2(\mathsf C^{(2)}(A),\zeta)_1=0\).
Thus \(G_1,H_1,\mathfrak G\) are unchanged as well.

The first contraction, defining \(\mathsf D\) and
\(\mathcal L=A'-H\mathsf D(A')\), uses only \(\mathsf C,H\).
In the six terms defining \(\mathsf W(A')\), the only
explicit current term is
\[
 -(\mathfrak D^{\mathsf T}
       -\mathsf C''[A']^{\mathsf T})H^{\mathsf T}\mathbf J .
\]
Its change is
\(-\mathsf C'_{\mathcal L}{}^{\mathsf T}\Psi_{\mathcal L}\zeta
+\mathsf C''[A']^{\mathsf T}\zeta=0\).
The second contraction consequently has the same
iterates and limit. This proves equality of \(\mathbf H\).
The current formula
\[
 \omega(A)=\Psi_A[H^{\mathsf T}\mathbf J+H^\circledast(A)],
 \qquad
 j(A)=f_A^{-1}(Q+\mathsf C'_A)^{\mathsf T}\omega(A)
                         -Q^{\mathsf T}H^{\mathsf T}\mathbf J
\]
then gives \eqref{II-93-24}.

For the last assertion, criticality supplies the multiplier
at \(A=0\), as in the proof of \cref{I-88-127}.
That proposition gives \(\mathcal Q(\mathbf H)=B'\).
On \(\Gamma_0^+\) the right side is zero, so
\(\mathbf H=0\) there. Since \(f_0=1\) and \(\delta\)
is supported on these rows,
\(f_{\mathbf H}^{-1}\delta=\delta\).
\end{proof}

\begin{proposition}\label{II-93-38}
Let \(U_+\) be in the regular critical orbit of its
datum for \(\Omega_0,\ldots,\Omega_{k+1}\), and suppose
\(V=\mathcal M^kU_+\) is axial on \(P''\).
Let \(P_+\) be its plain triple, at level \(k+1\), from
any gluing outside \(\Omega_0\). If
\((P_+,\mathcal A)\in
   \operatorname{dom}_{\mathfrak h}(T_\ell\mid\Omega_\bullet)\)
is real, define
\[
 (\mathbf A,\mathbf K)
       =\mathfrak h_{T_\ell\mid\Omega_\bullet}
                        (\mathbf1;P_+,\mathcal A),
 \qquad E=e^{i\eta_k\mathbf A}U_+ .
\]
Let \(\mathsf d\) be the shorter determining datum equal
to \(e^{iB'}V\) on \(\mathfrak X_0\), and to
\(M_{\mathbf B}(U_+)\) elsewhere. If the regular critical
orbit for \((\mathbf B,\mathsf d)\) exists, \(E\) is a
gauge transform of an element of it. At every complete
current stencil in \(\Omega_0\),
\begin{equation}\label{II-93-39}
 J^{(k)}(E)=L^{-3}J^{(k+1)}(U_+)+\mathbf K .
\end{equation}
Moreover \(\mathbf A=0\) on bonds not touching
\(\Omega_1\), and \(\mathbf K=0\) at incomplete
current stencils.
\end{proposition}

\begin{proof}
Fixing the shorter determining datum imposes all the
constraints of the longer one. Thus \(U_+\) is critical
for the shorter datum. Its reduced current equals its
full current at complete stencils. At an incomplete
stencil no corner of its adjacent plaquettes is in
\(\Omega_1\), by separation. There are at most three
remaining plaquettes, each with curvature less than
\(a_{\rm m}\), so the action derivative gives
\[
                 |J^\circ(U_+)(b)|<3a_{\rm m}\xi_0^{-3}.
\]
Consequently \((U_+,J^\circ(U_+))\) is in the natural
class at radius \(3a_{\rm m}\le a_{\rm D}\).
Its current agrees with \(L^{-3}\mathbf J_+\) on
every bond touching \(\Omega_1\); this is precisely
the scaling \(J^{(k+1)}=L^3J^{(k)}\).
The preceding lemma therefore identifies its response
with \((\mathbf A,\mathbf K)\).

The row source is zero except on interior level-\(k\)
rows. The outward-row correction to a datum uses only
that outward row and the bonds inside its outer endpoint
block; all these source values vanish. Thus the twisted
datum in the actual response theorem is \(B'\) itself.
By \cref{I-88-127}, \(E\) is critical with datum
\((e^{iB'}M_{\mathbf B}(U_+))^v\), where \(v\) consists
of the site factors. Choose \(g=v^{-1}\) at the determining
nodes and \(g=1\) at the other bare sites. Gauge covariance
then makes \(E^g\) critical with datum \(\mathsf d\).
The current identity in that proposition proves
\eqref{II-93-39}. It also shows \(\mathbf A=0\) on
\(\Gamma_0^+\). At an incomplete current stencil all
adjacent bonds are of this type, so \(E=U_+\) on its
plaquettes and its current increment is zero.

We check regularity, rather than infer it from criticality.
Put \(\beta_*=\sup|B'|\) and
\(\rho=5b_{\rm D}e^{8\bar\kappa}\beta_*\le\rho_*\).
The response estimates at the layer of the base point give
\[
 \xi_n|\mathbf A|,\ \xi_n^2|\nabla^{U_+}\mathbf A|\le\rho,
 \qquad
 \xi_n^3|\mathbf K|
                \le10C_jb_{\rm D}^3e^{24\bar\kappa}\beta_* .
\]
If \(j\) is the largest level touched by a bond or a
plaquette and \(n\) is its base-point level, separation
gives \(n\ge j-1\). Hence
\(\xi_n^{-1}\le L\mathsf L_{\mathbf B}^{-1}\).
The elementary plaquette variation estimate, applied to
\(\xi_n\mathbf A\), bounds the new curvature by
\[
 \bigl(3a_{\rm m}+2(1+\rho_*)\rho\bigr)
                                      \eta_k^2\xi_n^{-2}.
\]
The new current is bounded by
\((3a_{\rm m}+10C_jb_{\rm D}^3e^{24\bar\kappa}\beta_*)
\xi_n^{-3}\).
Multiplication by \(L^2\) and \(L^3\), respectively, is
absorbed by \(C_{\rm e}\); the last inequality of
\eqref{eq:history-shift-reserves}, with
\(\beta_*\le2C_P\varepsilon_{\mathfrak h}\), puts both
strictly within the regularity radius \(\varepsilon_{\rm v}\).
Regularity is gauge invariant. Thus \(E^g\) is in the
regular critical orbit whenever that orbit exists.
\end{proof}

\subsection{Pinned levelling of the response}\label{sec:pinned-levelling}

The accumulated endpoint factors of a response alter both its coarse
datum and its stored transports. We remove them by a gauge supported
inside the pillars. To control first differences, the prescribed values
are extended by smooth turns on smaller cells. The only cells of order
zero occur at level-one pillars with a two-step neighborhood containing
no higher layer.

Work at step \(t=k+1\), \(1\le k<K\), with
\(\xi=\eta_k=L^{-k}\), \(\eta_t=\xi/L\),
\(\xi_q=L^q\xi\), a separated reduced sequence
\(\Omega_\bullet=(\Omega_0,\ldots,\Omega_k)\), and a localization
domain \(Y\in\mathcal D_k\). The fine partition has \(M=L^{m'}\)
and \(m'<m+K-t\). Put
\[
 \mathfrak e_Y=\operatorname{Bd}(Y)\cap\operatorname{Bd}^{\circ}(\Omega_0),
 \qquad S_Y=\{\hbox{endpoints of bonds in }\operatorname{Bd}(Y)\}.
\]
For \(b\in\mathfrak e_Y^+\), let \(y(b)\) and \(n(b)\) be the
top and level of the pillar containing \(b_-\). A top \(y\) of
level \(q\ge1\) is a top in \(Y\) if \(B^q(y)\subset Y\).
Call it \emph{pinnable} if \(q\ge2\), or if \(q=1\) and every
bare site within \(2\xi\) of \(B^1(y)\) lies in a pillar of
level zero or one. Denote the tops in \(Y\) by \(\mathsf{Tp}(Y)\).
For a pinnable top, use the turning cells
\(B^{q-1}(a)\), \(a\in B(y)\); for a nonpinnable top use the
single cell \(B^1(y)\). Write \(\mathsf{Ct}(Y)\) for their centers,
\(\mathsf C(a)\) for the cell, \(o(a)\) for its order, and \(y(a)\)
for its top. These cells are disjoint and partition the pillars
of \(\mathsf{Tp}(Y)\).

For a triple \(P=(U,A',\mathbf J)\) at level \(t\), set
\(\mathbf U=e^{i\eta_tA'}U\). We shall require the following
conditions, denoted \(P\in\mathfrak T_k(Y;c)\):
\begin{equation}\label{eq:levelling-regular-triple}
 \begin{gathered}
 \xi_{y(b)}L^{-1}|A'(b)|<c,\qquad
 \xi_{y(b)}^3L^{-3}|\mathbf J(b)|<c,\qquad
 \xi_{y(b)}^2L^{-1}|(\nabla^U_\mu A'_\nu)(b_-)|<c,\\
 L^{2q}\varpi_{B^q(y)}(U)\le c\quad(y\in\mathsf{Tp}(Y)),\\
 U^w=e^{i\xi A^{\rm g}}\ \hbox{on }\mathsf P(\mathsf C(a)),\qquad
 \xi_{o(a)}|A^{\rm g}|\le c,\quad
 \xi_{o(a)}^2|\nabla A^{\rm g}|\le c\quad(o(a)\ge1).
 \end{gathered}
\end{equation}
The first line concerns \(b\in\mathfrak e_Y^+\) and contained
difference stencils. In the last line a real gauge \(w\) and a
real potential \(A^{\rm g}\) are required separately on each cell.

An admitted shift \(\mathfrak h=(\mathbf A,\mathbf K)\) is a pair
of complex one-forms on \(\operatorname{Bd}(Y)\), zero off
\(\mathfrak e_Y\). Its local size is
\begin{equation}\label{II-79-5}
 |\mathfrak h|_{\mathbf U,Y}(b)
 =\max\left\{\xi_{y(b)}|\mathbf A(b)|,\
 \xi_{y(b)}^2\max_\mu|(\nabla^{\mathbf U}_\mu\mathbf A_\nu)(b_-)|,\
 \xi_{y(b)}^3|\mathbf K(b)|\right\}.
\end{equation}
Here \(b=\langle x,x+\xi e_\nu\rangle\), and the inner maximum
uses only stencils contained in \(\mathfrak e_Y\), including
\(\mu=\nu\); an empty maximum is zero. Let
\(\|\mathfrak h\|_y\) be the maximum over bonds with \(y(b)=y\).
The neighbors of a top \(y\) are \(y\) itself and the pinnable
level-one tops \(y'\) in \(Y\) whose pillars are within \(2\xi\)
of its pillar. Put
\(\|\mathfrak h\|_y^+=\max_{y'\in\operatorname{nb}_Y(y)}
\|\mathfrak h\|_{y'}\).

For \(y\in\mathsf{Tp}(Y)\) of level \(q\), let
\(\mathsf{Rd}_Y(y)=\mathsf P(B^q(y))\), except at pinnable
level-one tops, where it consists of the bonds of \(\mathfrak e_Y\)
whose endpoints are within \(2\xi\) of \(B^1(y)\). Define
\[
 \varepsilon_y(\mathbf A)=\xi_q\sup_{\mathsf P(B^q(y))}|\mathbf A|,
 \qquad
 \widehat\varepsilon_y(\mathbf A)=
                  \xi_q\sup_{\mathsf{Rd}_Y(y)}|\mathbf A|.
\]
A top is \emph{levelled} when
\(\widehat\varepsilon_y(\mathbf A)<\epsilon\). Notice that
\(\mathsf P(B^q(y))\subset\mathsf{Rd}_Y(y)\subset\mathfrak e_Y\)
and \(\varepsilon_y(\mathbf A)\le\|\mathfrak h\|_y\).

Choose \(c_{\rm gf}>0\) and \(C_{\rm gf}\ge1\), depending only
on \(L\), as the restricted tower constants in the proof below.
Fix an even \(C^2\) function \(\phi\) which vanishes for
\(|s|\ge1/3\), with \(0\le\phi\le1\), \(\phi(0)=1\), and
\(|\phi'|,|\phi''|\le b_\phi\), where \(b_\phi\ge1\).
For example, \(\phi(s)=(1-9s^2)^3\) for \(|s|<1/3\) and zero
elsewhere has these properties with \(b_\phi=54\). Set
\begin{equation}\label{II-85-3}
 c_{\rm m}=10C_{\rm gf}^2e^{4C_{\rm gf}},\qquad
 c_{\rm rt}=\min(c_{\rm gf},1/20),\qquad
 \varepsilon_{\rm m}=\frac{c_{\rm gf}}{100c_{\rm m}},\qquad
 C_{\rm m}=200b_\phi^2L^2c_{\rm m}.
\end{equation}
Decrease \(c_{\rm gf}\), if necessary, so \(c_{\rm gf}\le1/1024\),
and fix \(0<\epsilon\le\varepsilon_{\rm m}\).
Write \(v_q^\sharp[\mathbf U;\mathbf A](y)\) for the accumulated
endpoint factor of the centered relative recursion over \(B^q(y)\).

\begin{proposition}\label{II-85-34}
For \(P\in\mathfrak T_k(Y;c_{\rm rt})\) and every admitted shift
\(\mathfrak h\), the construction below defines
\(u=u^\delta_Y[P;\mathbf A]\) on all bare sites, equal to one
off \(S_Y\), and an admitted shift
\(\mathfrak h^\delta[P]=(\mathbf A^\delta,\mathbf K^\delta)\).
At a levelled top \(y\) of level \(q\),
\[
 u(y)v_q^\sharp[\mathbf U;\mathbf A](y)=1,\qquad
 |u(x)^{\pm1}-1|\le C_{\rm m}\varepsilon_y(\mathbf A)
                    \quad(x\in B^q(y)).
\]
At a pinnable levelled top it also preserves every stored transport:
\[
 \operatorname{ax}_{(e^{i\xi\mathbf A}\mathbf U)^u}(a)
       =\operatorname{ax}_{\mathbf U}(a),\qquad a\in B(y).
\]
Moreover \(\operatorname{cn}(u)\le3/2\), and on positive bonds
\begin{equation}\label{eq:pinned-shift-identities}
 \begin{split}
 e^{i\xi\mathbf A^\delta(b)}\mathbf U(b)
   &=u(b_-)e^{i\xi\mathbf A(b)}\mathbf U(b)u(b_+)^{-1},\\
 L^{-3}\mathbf J(b)+\mathbf K^\delta(b)
   &=\operatorname{Ad}(u(b_-))
                     (L^{-3}\mathbf J(b)+\mathbf K(b)).
 \end{split}
\end{equation}
The shift is unchanged when \(u(b_-)=u(b_+)=1\), and
\begin{equation}\label{eq:pinned-local-bound}
 |\mathfrak h^\delta[P]|_{\mathbf U,Y}(b)
             \le C_{\rm m}\|\mathfrak h\|^+_{y(b)}.
\end{equation}
The construction is local on the indicated pillars and reading sets,
commutes with real gauge transformations and complex conjugation,
and is Borel in all its data. For fixed \(U\), its values on
specified pillars are holomorphic on the open set where their tops
are levelled. Its value at a bond is holomorphic when this condition
holds for every top in \(Y\) whose pillar contains either endpoint.

If \(\mathsf s\ge0\) on the bare sites of \(\Omega_0\) satisfies
\[
 \mathsf s(x)\le e^{\delta d_{\mathbf B}(x,x')}\mathsf s(x'),
 \qquad |\mathfrak h|_{\mathbf U,Y}(b)\le\mathsf s(b_-),
\]
then
\begin{equation}\label{eq:pinned-profile-bound}
 \|\mathfrak h\|^+_{y(b)}\le e^{C_{\rm m}\delta}\mathsf s(b_-),
 \qquad
 \widehat\varepsilon_{y'}(\mathbf A)
      \le L e^{C_{\rm m}\delta}\mathsf s(b_-)
 \quad(y'\in\operatorname{nb}_Y(y(b))\cap\mathsf{Tp}(Y)).
\end{equation}
All bounds persist after restriction to
\(\operatorname{Bd}(X)\), \(X^\uparrow\subset Y\), using only
the difference stencils contained there.
\end{proposition}

\begin{proof}
\emph{The prescribed values.}
On a pillar \(B^q(y)\), use the centered relative recursion
of \eqref{II-155-25} with initial perturbation \(\xi\mathbf A\).
Write \(v_i\) for its accumulated factors and
\(\mathbf U_{(i)}=\mathcal M^i\mathbf U\). For \(a\in B(y)\) put
\[
 \operatorname{ax}_{\mathbf W}(a)
           =(\mathcal M^{q-1}\mathbf W)(y,a),\qquad
 \mathbf E=e^{i\xi\mathbf A}\mathbf U,\qquad
 p(a)=\operatorname{ax}_{\mathbf U}(a)^{-1}
                v_q(y)^{-1}\operatorname{ax}_{\mathbf E}(a).
\]
The single-pillar statement is the restriction of the double-block
one. Extend the real configuration by the coordinate folding used
in \cref{I-32-37}, and extend the two perturbations by zero to an
adjacent double block. Folding preserves the maximal plaquette
defect. Every contour and lower double block entering the recursion
on the original pillar stays in it, so its values and domains are
independent of the extension.

Choose \(c_{\rm gf}\) below the radius in \eqref{II-155-25},
with enough reserve that the contour diameters are less than
\(\omega_{\rm c}/2\) and all normalized averaging logarithms
are within half their defining balls. This follows from the
same contour estimates there, applied also to
\(i^{-1}\log(e^{i\xi\mathbf A}e^{ia'})\); its norm is at most
\(2(\xi|\mathbf A|+|a'|)\). The required data are
\[
 L^{2q}\varpi_{B^q(y)}(U)\le c_{\rm gf},\qquad
 L^q\sup|a'|<c_{\rm gf},\qquad
 \varepsilon_y(\mathbf A)<c_{\rm gf}.
\]
They hold at each levelled top with \(a'=\eta_tA'\), by
\eqref{eq:levelling-regular-triple}, since
\(L^q|\eta_tA'|=\xi_qL^{-1}|A'|<c_{\rm rt}\).
Thus the actual complex
averages and transports are defined. They are holomorphic in
\((a',\mathbf A)\), equal their real versions for real data,
and satisfy
\[
 |v_i^{\pm1}-1|\le C_{\rm gf}L^{i-q}\varepsilon_y,\qquad
 S:=\operatorname{ax}_{\mathbf U}(a)
       =e^{i\zeta'(a)}\operatorname{ax}_U(a),\quad
 |\zeta'(a)|\le C_{\rm gf}L^q\sup|a'|.
\]
At the last step let
\(\zeta(a)=\xi_{\mathsf m}(\mathbf U_{(q-1)},a_{q-1})(y,a)\)
and \(\bar\zeta=L^{-3}\sum_{a\in B(y)}\zeta(a)\).
The one-step factorization gives
\begin{equation}\label{II-85-19}
 \operatorname{ax}_{\mathbf E}(a)
   =v_{q-1}(y)e^{i\zeta(a)}S v_{q-1}(a)^{-1},\qquad
 p(a)=S^{-1}e^{-i\bar\zeta}e^{i\zeta(a)}S\,v_{q-1}(a)^{-1}.
\end{equation}
Here \(|\zeta(a)|\le C_{\rm gf}L^{-1}\varepsilon_y\).
Consequently
\[
 |p(a)^{\pm1}-1|
 \le\{2C_{\rm gf}e^{4C_{\rm gf}}(1+C_{\rm gf})+C_{\rm gf}\}
                   \varepsilon_y
 \le\tfrac12c_{\rm m}\varepsilon_y<1/200.
\]
In particular, \(p(y)=v_q(y)^{-1}\), and the pin logarithm
\(\zeta^{\rm pin}(a)=i^{-1}\log p(a)^{-1}\) has norm at most
\(c_{\rm m}\varepsilon_y\). Put this logarithm at every center
over a levelled top and zero at all other centers.

The covariance of the transports also holds under a complex
gauge \(h\) with \(\operatorname{cn}(h)\le3/2\). Indeed,
path factors cancel, and left conjugation multiplies each
contour diameter and normalized logarithm bound by at most
\(3/2\). The preceding half-radius reserve keeps the transformed
objects in the same branches. Induction over the actual averages
therefore gives
\(\operatorname{ax}_{\mathbf E^h}(a)
=h(y)\operatorname{ax}_{\mathbf E}(a)h(a)^{-1}\).
The choice \(h(y)=v_q(y)^{-1}\), \(h(a)=p(a)\) makes this
equal to \(\operatorname{ax}_{\mathbf U}(a)\).

\emph{Extending the prescribed values.}
On a cell \(\mathsf C(a)=B^o(a)\), \(o\ge1\), consider real
endomorphism fields \(T\), normalized by \(T(a)=1\), satisfying
\[
 \xi_o|\nabla^U_\varsigma T|\le1,\qquad
 \xi_o^2|\nabla^U_{\varsigma'}\nabla^U_\varsigma T|\le1
\]
for all signed directions and all contained stencils. The norm on
endomorphisms here is the Euclidean operator norm. Among these fields
choose the unique minimizer of the sum of their squared Hilbert--Schmidt
norms.

This set is nonempty, with a strict feasible point. Indeed, the cell
gauge in \eqref{eq:levelling-regular-triple} gives
\(T_0(x)=\operatorname{Ad}(w(x)^{-1}w(a))\). Put
\(D_\varsigma(x)=\operatorname{Ad}(e^{i\xi A^{\rm g}_\varsigma(x)})-1\).
After conjugation by the orthogonal maps \(\operatorname{Ad}w\), its
first difference is \(D_\varsigma/\xi\), and its second difference is
\[
 \xi^{-2}\bigl[(1+D_{\varsigma'}(x))
 D_\varsigma(x+\xi\varsigma')-D_\varsigma(x)\bigr].
\]
The potential and first-difference bounds give, after multiplication
by \(\xi_o\) and \(\xi_o^2\), respectively,
\(2e^c c<1/2\) and \(2e^c c+4e^{2c}c^2<1/2\), where
\(c=c_{\rm rt}\). Negative directions follow by translating the
same identities. The feasible set is finite-dimensional, closed
and convex, and the objective is strictly convex. Its minimizer
therefore exists and is unique. A coordinate path of length at
most \(3(L^o-1)/2\) from \(a\) gives \(|T(x)|\le5/2\).
The complex-linear extension has norm at most \(5\) in the matrix
norm used for \(\mathfrak g_{\mathbb C}\), and its first and second
differences have bounds \(2\xi_o^{-1}\) and \(2\xi_o^{-2}\).

This choice is continuous in \(U\). To see this, mix a minimizer
with a positive amount of the strict feasible field \(T_0\).
That mixture remains feasible for every sufficiently nearby
configuration. It gives the upper bound for the limiting minimum;
boundedness and uniqueness then identify every subsequential
limit. Real gauges act on the feasible set by
\[
 T(x)\longmapsto
 \operatorname{Ad}(g(x))T(x)\operatorname{Ad}(g(a))^{-1}.
\]
They preserve the objective, so the minimizer is gauge covariant.
Its definition also commutes with translations of the cell.

Set
\[
 b_a(x)=\prod_{\mu=1}^3
       \phi\bigl((x_\mu-a_\mu)/\xi_o\bigr),\qquad
 \mu(x)=b_a(x)T(x)\zeta^{\rm pin}(a)
       \quad(x\in\mathsf C(a),\ o\ge1).
\]
On a singleton cell set \(\mu(a)=\zeta^{\rm pin}(a)\), and set
\(\mu=0\) elsewhere. Finally put \(u=e^{-i\mu}\).
For \(z=|\zeta^{\rm pin}(a)|\), the product rule gives
\begin{equation}\label{eq:levelling-turn-differences}
 |\mu|\le5z,\qquad
 |\xi\nabla^U_\varsigma\mu|\le7b_\phi L^{-o}z,\qquad
 |\xi^2\nabla^U_{\varsigma'}\nabla^U_\varsigma\mu|
                       \le11b_\phi^2L^{-2o}z.
\end{equation}
For completeness, if \(\tau=T\zeta^{\rm pin}(a)\), all values
transported to the same base point, the mixed product difference is
\[
 b_{12}(\tau_{12}-\tau_1-\tau_2+\tau_0)
 +(b_{12}-b_1)(\tau_1-\tau_0)
 +(b_{12}-b_2)(\tau_2-\tau_0)
 +(b_{12}-b_1-b_2+b_0)\tau_0.
\]
The four bounds are \(2\), \(2b_\phi\), \(2b_\phi\), and
\(5b_\phi^2\), times \(L^{-2o}z\). They also hold after
extension by zero: a nonzero bump lies in the inner third,
whose two-step neighborhood is inside the cell, since \(L^o\ge13\).
At a singleton, pinnability instead ensures that its two-step
neighborhood contains only levels zero and one.

The prescribed logarithms lie in the larger polydisc
\[
 |\zeta(a)|<2c_{\rm m}\varepsilon_{y(a)}
 \quad\hbox{at levelled tops with }\varepsilon_{y(a)}>0,
 \qquad \zeta(a)=0\quad\hbox{at all other centers}.
\]
Throughout this polydisc \(z<1/50\), \(|\mu|\le1/10\),
\(|u^{\pm1}-1|\le6z\), and \(\operatorname{cn}(u)\le5/4\).
If \(u\) is nontrivial at an endpoint of a bond, that bond belongs
to the reading set of a levelled top; hence
\(\xi|\mathbf A(b)|<\epsilon\). Define, on positive bonds,
\[
 \mathbf A^\delta(b)=\frac1{i\xi}\log
 \bigl(u(b_-)e^{i\xi\mathbf A(b)}
       \mathbf U(b)u(b_+)^{-1}\mathbf U(b)^{-1}\bigr),
 \qquad
 \mathbf K^\delta(b)=
 \operatorname{Ad}(u(b_-))(L^{-3}\mathbf J+\mathbf K)(b)
                         -L^{-3}\mathbf J(b).
\]
When both endpoint values of \(u\) equal one, define
\(\mathbf A^\delta(b)=\mathbf A(b)\) instead. This proviso
is essential: the original shift need not be small on an
untouched bond. Extend both forms oddly to reversed bonds.
The logarithm on a touched bond is the series logarithm; the
estimate below places its argument within \(2/5\) of one.
The inner-third support and the singleton neighborhood show
that bonds outside \(\mathfrak e_Y\) are untouched, proving
admissibility. At the centers \(T(a)=1\) and \(b_a(a)=1\),
so \(u(a)=p(a)\). A top or a stored site over a pinnable top
is either such a center or outside every turning cell.
Thus the prescribed values and the stored-transport identity
are independent of the frame away from the centers.

\emph{The first-difference estimate.}
We give the mixed estimate that controls the logarithm. For
\(|\theta_0|,|\theta_1|,|\theta_2|,|\theta_{12}|\le1/8\) and
\(|\alpha_0|,|\alpha_1|\le\rho\le1/12\), put
\[
 X_0=e^{-i\theta_0}e^{i\alpha_0}e^{i\theta_2},
 \qquad X_1=e^{-i\theta_1}e^{i\alpha_1}e^{i\theta_{12}},
 \quad
 d_i=|\theta_i-\theta_0|,\quad
 s=|\theta_{12}-\theta_1-\theta_2+\theta_0|.
\]
Then \(|X_j-1|<2/5\), and
\begin{equation}\label{eq:levelling-log-differences}
 \begin{split}
 |\log X_0|&\le3|\alpha_0|+3d_2,\\
 |\log X_1-\log X_0|
 &\le3|\alpha_1-\alpha_0|+5s+7d_1d_2+6\rho d_1.
 \end{split}
\end{equation}
Here is the product calculation, including the term which would
be lost by estimating four exponentials separately. Write
\(u_j=e^{-i\theta_j}\) and \(a_j=e^{i\alpha_j}\).
The difference \(X_1-X_0\) is the sum
\[
 \begin{split}
 &[(u_1-u_{12})-(u_0-u_2)]u_{12}^{-1}
 +(u_0-u_2)(u_{12}^{-1}-u_2^{-1})\\
 &\quad +(u_1-u_0)(a_1-1)u_{12}^{-1}
 +u_0(a_1-a_0)u_{12}^{-1}
 +u_0(a_0-1)(u_{12}^{-1}-u_2^{-1}).
 \end{split}
\]
In the exponential power series, every word in the mixed difference
either contains the mixed increment or contains both first increments.
If \(r=\max|\theta_j|\), this gives
\[
 |u_{12}-u_1-u_2+u_0|
       \le e^{3r}s+e^{5r}d_1d_2.
\]
The ordinary exponential difference estimate bounds the other four
terms. The logarithm is Lipschitz with constant \(2\) in the
\(2/5\)-ball. Substitution of \(r\le1/8\), \(\rho\le1/12\)
gives \eqref{eq:levelling-log-differences}; the same estimates
give \(|\log X_0|<.52<\log2\).

The differences in \eqref{eq:levelling-turn-differences} use \(U\);
the logarithm uses \(\mathbf U=e^{i\eta_t A'}U\).
This change preserves their scale. Indeed, put
\(E_\mu=\operatorname{Ad}(e^{i\eta_t A'_\mu})\) and let
\(\widehat E_\nu\) be its \(\nu\)-analogue transported with \(U_\mu\).
On a cell in a level-\(q\) pillar, the first line of
\eqref{eq:levelling-regular-triple} gives
\[
 |E_\mu|\le1.008,\qquad
 |E_\mu-1|\le\tfrac18L^{-q},\qquad
 |\widehat E_\nu-E_\nu|\le\tfrac18L^{-2q}.
\]
If hats denote the four \(U\)-transported values, the mixed
\(\mathbf U\)-transported difference is
\[
 E_\mu\widehat E_\nu\widehat s
 +E_\mu(\widehat E_\nu-1)\widehat d_1
 +(E_\mu\widehat E_\nu-E_\nu)\widehat d_2
 +[(E_\mu-1)(\widehat E_\nu-1)
                  +(\widehat E_\nu-E_\nu)]\theta_0.
\]
Since \(o\le q\), this yields
\[
 d_1,d_2\le9b_\phi L^{-o}z,\qquad
 s\le15b_\phi^2L^{-2o}z.
\]

Fix a contained difference stencil at \(b\), and put
\(H=\|\mathfrak h\|_{y(b)}\), \(H_+=\|\mathfrak h\|^+_{y(b)}\).
There are three cases. If \(\mu\) vanishes on its four sites,
the shift is unchanged. If a nonzero value lies in a cell of
positive order, the entire stencil lies in that cell.
Its level \(q\) satisfies \(o\le q\le o+1\), and
\[
 z\le\min(1/50,2c_{\rm m}H),\qquad
 \rho\le1.008L^{-q}H.
\]
Apply \eqref{eq:levelling-log-differences} with
\(\alpha_0=\xi\mathbf A_\nu(b_-)\) and
\(\alpha_1=\operatorname{Ad}(\mathbf U_\mu)
                   \xi\mathbf A_\nu(b_-+\xi e_\mu)\).
After multiplying by the output scales, it gives
\[
 \xi_{y(b)}|\mathbf A^\delta(b)|
                    \le(3+54b_\phi Lc_{\rm m})H,\qquad
 \xi_{y(b)}^2|\nabla^{\mathbf U}_\mu\mathbf A^\delta_\nu(b_-)|
                    \le175b_\phi^2L^2c_{\rm m}H.
\]
In the remaining case the nonzero turns are singletons.
All four sites have levels at most one; the tops of those with
nonzero turns belong to \(\operatorname{nb}_Y(y(b))\). With
\(z_\bullet\le\min(1/50,2c_{\rm m}H_+)\), use
\[
 |\theta_j|\le5z_\bullet/4,\quad
 d_1,d_2\le5z_\bullet/2,\quad s\le5z_\bullet
\]
in the same calculation. The corresponding bounds are
\((3+15Lc_{\rm m})H_+\) and \(59L^2c_{\rm m}H_+\).
The current formula in \eqref{eq:pinned-shift-identities},
the bound on \(u-1\), and
\(\xi_{y(b)}^3L^{-3}|\mathbf J(b)|<c_{\rm rt}\)
give \(\xi_{y(b)}^3|\mathbf K^\delta(b)|\le3c_{\rm m}H_+\).
All three bounds are absorbed by \(C_{\rm m}\), proving
\eqref{eq:pinned-local-bound}.

\emph{Dependence and profile bounds.}
The pin reads only the actual tower over its pillar, the levelled
test reads only \(\mathsf{Rd}_Y(y)\), and the frame reads only
the real configuration on its cell. These observations give
locality and real gauge covariance. For fixed \(U\), on the
open set where the relevant tops remain levelled, the pins and
the series logarithms are holomorphic. Where both turns become
one the logarithmic formula agrees with the inert definition,
since the relevant input is small there. Thus there is no
additional branch at that locus. Joint continuity follows on
each prescribed set of levelled tops from continuity of the
frames and of the actual averages. The finitely many levelled
sets form Borel pieces, giving the global Borel assertion.
For real data all factors are unitary. Holomorphic reflection
therefore gives \(u(\overline P,\overline{\mathfrak h})=u(P,\mathfrak h)^c\);
the real frames commute with the involution, and
\(|\log X|<\log2\) identifies the same logarithm after
conjugation. This proves the conjugation statement.

A coordinate path within a level-\(q\) pillar has weighted
length less than \(3\). Two bare sites at distance at most
\(2\xi\) in the supremum norm are joined by at most six bare
steps, of weighted length at most \(6\). Thus a top neighboring
\(y(b)\), together with a site in its pillar, is at weighted
distance at most \(3+6+3=12\) from \(b_-\).
The exceptional enlarged reading set costs at most six
further steps. Its levels are zero or one, so changing
\(\xi_{y'}\) to the input scale costs at most \(L\).
Since \(C_{\rm m}\ge18\), these facts prove
\eqref{eq:pinned-profile-bound}. Restriction to a smaller
bond domain retains precisely a subset of these stencils.
\end{proof}

The locality needed for old terms includes two exact support
statements. They apply for all complex values of the parameters;
holomorphy still requires the forced-levelled condition above.

\begin{proposition}\label{II-85-45}
Let \(\gamma_1,\ldots,\gamma_N\) be the wall-components of a
set of fine cubes in \(Y\). Suppose that the admitted shifts
\(\mathfrak h_i\) vanish unless both endpoints belong to
\(\gamma_i\), and satisfy the margin condition
\[
 \mathfrak h_i(b)=0
 \quad\hbox{if an endpoint \(p\) of \(b\) is within
 \((L+2)\xi\) of a bare site outside \(\gamma_i\)}.
\]
For \(\mathfrak h=\sum_i\lambda_i\mathfrak h_i\),
\(\lambda_i\in\mathbb C\), one has
\[
 u_Y^\delta[P;\mathbf A]
       =\prod_i u_Y^\delta[P;\lambda_i\mathbf A_i],
 \qquad
 \mathfrak h^\delta[P]
       =\sum_i(\lambda_i\mathfrak h_i)^\delta[P].
\]
At each site at most one factor differs from one. Each pinned
summand still vanishes unless both endpoints lie in its component,
and a zero input gives a zero summand.

If \(Y_1\subset Y\), \(Y_1\in\mathcal D_k\), and an admitted
\(\mathfrak h\) vanishes unless both endpoints belong to \(Y_1\),
then \(P|_{Y_1}\in\mathfrak T_k(Y_1;c_{\rm rt})\) and
\[
 u_Y^\delta[P;\mathbf A]\big|_{S_{Y_1}}
    =u_{Y_1}^\delta[P|_{Y_1};\mathbf A|_{\operatorname{Bd}(Y_1)}],
 \qquad u_Y^\delta=1\quad\hbox{on }S_Y\setminus Y_1.
\]
The corresponding shifts agree on \(\operatorname{Bd}(Y_1)\)
and are zero on its complement.
\end{proposition}

\begin{proof}
Every pillar lies in a single fine cube, and hence either in
one component or outside their union. Two bare nearest neighbors
cannot lie in different wall-components: their fine cubes are
equal or adjacent. On a pillar in \(\gamma_i\), the summed shift
equals \(\lambda_i\mathfrak h_i\); every other shift vanishes.
This also identifies its pin and frame. The only enlarged
levelled test is at a pinnable level-one top. Every bond of
that reading set lies within \(2\xi\) of its pillar; the
margin therefore annihilates all other component shifts
there. Thus the levelled test agrees as well. On all other
pillars the turn is one. These facts prove the gauge product,
and the bond formula \eqref{eq:pinned-shift-identities} proves
the sum, including the inert proviso.

To check the asserted support after turning, consider a bond
crossing out of \(\gamma_i\). A nonzero turn at its inside
endpoint cannot lie in the inner third of a positive-order
cell, since its nearest neighbor would then lie in that cell
and in \(\gamma_i\). If it is a singleton, every point of its
level-one pillar is within \(L\xi\) of the outside endpoint.
The margin forces the original shift to vanish on that whole
pillar. Its pin, and hence its turn, is one. The crossing
bond is therefore untouched.

For restriction to \(Y_1\), the tops and cells are exactly
those of \(Y\) lying in \(Y_1\); their frames and pillar
data coincide. The extra bonds of
\(\mathsf{Rd}_Y(y)\setminus\mathsf{Rd}_{Y_1}(y)\) lie outside
\(\operatorname{Bd}(Y_1)\), and the shift vanishes on them.
Thus the levelled tests coincide. On every pillar outside
\(Y_1\) the initial perturbation is zero, so its turn is one.
The bond formula gives the remaining assertions. The conditions
on \(P|_{Y_1}\) involve only a subset of the original stencils,
pillars and cells.
\end{proof}

We record the endpoint geometry used when these pins are compared
between successive histories. A stored site of a level-\(q\)
top \(y\) means a site \(a\in B(y)\); its stored transport is
\((\mathcal M^{q-1}W)(y,a)\).

\begin{lemma}\label{II-84-18}
Let \(h\) be a history through level \(k\), \(0\le l<k\), and
let \(x\) be an endpoint of an external bond of level \(l\).
Then \(x\notin\Lambda_{l+1}\). It belongs to \(\Omega_l\)
if \(l\ge1\), and to \(\Omega_1^\sim\) if \(l=0\).
Write \(y=\mathsf c_{l+1}(x)\).
For the separated truncations of this history and the one-step
extensions arising from an admissible next outcome, with the
same historical regions through level \(l+1\), the following
alternatives hold.

If \(x\notin\Omega_{l+1}\), it is a top of level \(l\).
For \(l\ge1\) it is its own stored site; for \(l=0\) it has
no stored transport. If \(x\in\Omega_{l+1}\), it is a stored
site of the pinnable top \(y\) of level \(l+1\).
In either nonempty case, the stored transport is the same
actual contour average for every such sequence:
\begin{equation}\label{II-84-21}
 \operatorname{ax}_W(x)
       =(\mathcal M^{q-1}W)(y_*,x),
 \qquad
 (q,y_*)=
 \begin{cases}
 (l,x),&x\notin\Omega_{l+1},\ l\ge1,\\
 (l+1,y),&x\in\Omega_{l+1}.
 \end{cases}
\end{equation}
Definedness in either the real or complex sense is the same.

For a subordinate separated sequence containing \(x\), its
layer \(q'\) is at most \(l+1\). If \(q'\le l\), then \(x\)
is its own top and \(B^{q'}(x)\subset B^{l+1}(y)\).
If \(q'=l+1\), it has the same pinnable top \(y\) and the
same transport in \eqref{II-84-21}.
\end{lemma}

\begin{proof}
An external bond is incident on \(\Omega_{l+1}\) and has
neither endpoint in \(\Lambda_{l+1}\). Its endpoints are
within one level-\(l\) step of \(\Omega_{l+1}\). The
historical collars have widths at least \(8LM\check R_{l+1}\)
inside \(\Omega_l\), or \(LM\check R_1\) inside
\(\Omega_1^\sim\) when \(l=0\); these widths exceed one.
This proves the first assertions.

The level-\((l+1)\) cube containing \(x\) is a cube of
\(Z_{l+1}=\Lambda_{l+1}^c\). Its distance from the next
region is at least \(8LM\check R_{l+2}\), measured at level
\(l+1\); the same bound holds for a possible next step
after the history. Thus no higher region can contain \(x\).
The layer is exactly \(l\) or \(l+1\) according as
\(x\notin\Omega_{l+1}\) or \(x\in\Omega_{l+1}\).
The block decomposition now gives the stated top and stored
site. For \(l\ge1\) the latter top is pinnable by its level.
For \(l=0\), its containing cube is at distance more than
two bare steps from every level-two region. The separated
level-one collar also contains that two-step neighborhood.
Hence every site in the neighborhood is in a layer of
level zero or one, which is exactly pinnability.

A subordinate sequence has smaller regions and inherits
the same exclusion of level \(l+2\). Its layer \(q'\le l\)
puts the level-\(l\) site \(x\) on its own top; nesting of
blocks gives \(B^{q'}(x)\subset B^{l+1}(y)\).
If \(q'=l+1\), the same containing block and two-step
argument apply. Finally every contour and lower block
used in \eqref{II-84-21} is inside the indicated pillar.
The restriction argument in the proof of \cref{II-85-34}
therefore identifies both its domain and its value.
\end{proof}

We now apply the construction to the real canonical response.
Let \(\mathsf U_0\) be a real configuration on the ambient torus,
let \(U_0\) be its restriction to the inner bonds of \(\Omega_0\),
and put \(V_0=M_{\mathbf B}(U_0)\).
Its plain triple at level \(t\) is
\[
 P_0=(\mathsf U_0,0,\mathbf J_+),\qquad
 \mathbf J_+(b)=
 \begin{cases}
 J^{(t)}(U_0)(b),&\text{the full current stencil is in }\Omega_0,\\
 0,&\text{otherwise}.
 \end{cases}
\]
Use the scale-\(k\) pairs
\(\mathbf p_1=(U_0,L^{-3}\mathbf J_+)\) and
\(\mathbf p_0=(U_0,J^\circ(U_0))\).
The condition \((\mathrm{m2})[\alpha]\) means that
\(\mathbf p_1\), with decomposition \((U_0,0)\), satisfies
the four local natural-class conditions at radius \(\alpha\).
All response constants below are those of the regional response,
with their common bound denoted by \(\mathsf C\).

\begin{proposition}\label{II-94-44}\label{II-94-49}
Suppose that \(U_0\) is critical for its reduced determining
datum, that \(P_0\) satisfies \((\mathrm{m2})[\alpha]\), and that
\(B'\) is a real admissible response datum supported on the
interior rows of level \(k\), with
\[
 \mathsf b=\sup|B'|\le\varepsilon_{\rm r}
      :=\varepsilon_{\rm D}e^{-32\bar\kappa}.
\]
Let \((\mathbf H,\mathsf j)\) be the canonical response at
parameters one and pair \(\mathbf p_1\). Assume
\begin{equation}\label{II-94-50}
 \alpha\le\min(a_{\rm D}/4,L^{-2}c_{\rm gf},c_{\rm rt}),
 \qquad
 L\mathsf C e^{\mathsf C\bar\kappa}\mathsf b
                   <\epsilon\le\varepsilon_{\rm m}.
\end{equation}
Then \(P_0\in\mathfrak T_k(T_\ell;c_{\rm rt})\), every top
is levelled, and the pinned shift is real. The configuration
\[
 U^\delta=(e^{i\xi\mathbf H}U_0)^{u^\delta}
                          =e^{i\xi\mathbf H^\delta}U_0
\]
has actual reduced averages \(e^{iB'}V_0\) and is critical
for that datum. At every complete inner current stencil,
\begin{equation}\label{II-94-51}
 J^{(k)}(U^\delta)=L^{-3}J^{(t)}(U_0)+\mathsf j^\delta .
\end{equation}
At every stored site of a pinnable top, its actual transport
is the transport of \(U_0\).
If in addition
\(\mathsf C(4\alpha+e^{\mathsf C\bar\kappa}\mathsf b)
\le\varepsilon_{\rm v}\), then \(U^\delta\) is regular at
radius \(\varepsilon_{\rm v}\); it belongs to the minimal
orbit for \(e^{iB'}V_0\) whenever that orbit exists.
\end{proposition}

\begin{proof}
The two currents agree on every bond touching \(\Omega_1\).
At an incomplete star the bond is in the bare layer, and
there are at most three remaining inner plaquettes. The
natural-class curvature bound gives
\(|J^\circ(U_0)(b)|<3\alpha\xi^{-3}\) there.
Thus \(\mathbf p_0\) belongs to the natural class at radius
\(4\alpha\), whereas \(\mathbf p_1\) belongs at radius
\(\alpha\). Both responses are defined. Their equality is
\cref{II-93-23}: the current difference is supported on bare
determining rows, on which the datum and the response vanish.
Criticality at \(A=0\) gives
\(J^\circ(U_0)=Q^{\mathsf T}\lambda_0\) by the
constraint-multiplier argument in \cref{I-88-127}.
That proposition therefore gives
\[
 \mathcal Q(\mathbf H)=B',\qquad
 E:=e^{i\xi\mathbf H}U_0\ \text{critical for }M_{\mathbf B}(E),
 \qquad
 J^{(k)}(E)=L^{-3}J^{(t)}(U_0)+\mathsf j
\]
at complete stencils. These are the same response and
current conventions as in \cref{II-93-38}.

The natural gauges also give the regular-triple conditions.
The first line of \eqref{eq:levelling-regular-triple} is
immediate for \(P_0\). Every plaquette in \(B^q(y)\) has
base-point layer \(q\), so its defect is bounded by
\(\alpha L^{-2q}\). For a turning cell in that pillar,
take the \(\mathfrak M\xi_q\)-cube containing \(y\).
It is contained in the same layer: if any of its finer
cubes were in \(\Omega_{q+1}\), nesting of the partitions
would put \(y\) there too. It is therefore a natural cube
inside \(\Omega_0\). Its containing-box gauge restricts
to the turning cell. Since \(o\le q\), the bounds at scale
\(\xi_q\) imply the required bounds at \(\xi_o\).
Consequently \(P_0\in\mathfrak T_k(T_\ell;\alpha)\).

For a row \(c\in\Gamma_q\), its double block
\(\mathsf Z=B^q(c_-)\cup B^q(c_+)\) lies in
\(\Omega_{q-1}\), with \(\Omega_{-1}=\Omega_0\).
Every base-point layer there satisfies
\(n\ge\max(q-1,0)\), and hence \(\xi_q\le L\xi_n\).
The natural curvature bound and the regional response estimate
give
\begin{equation}\label{II-94-45}
 L^{2q}\varpi_{\mathsf Z}(U_0)\le L^2\alpha,\qquad
 \xi_q\sup_{\mathsf P(\mathsf Z)}|\mathbf H|
                 \le L\mathsf C e^{\mathsf C\bar\kappa}\mathsf b.
\end{equation}
The curvature maximum is zero at \(q=0\). These bounds
place the actual double-block recursion in the domain used
in the proof of \cref{II-85-34}. The nonlinear response
average \(\mathcal Q_q(U_0,\xi\mathbf H)\) is its
straightened terminal value: at order zero the inputs
coincide, and at every later order both definitions apply
the same one-step factorization to each contained lower
double block. Reversal is imposed after the positive-bond
step in both definitions. Induction therefore identifies
the domains and the values, not only their derivatives.
Its terminal value is \(B'(c)\).

On an ordinary pillar reading set \(n=q\); on the enlarged
reading set of a pinnable level-one top \(n\ge0\).
The same estimate consequently gives
\[
 \widehat\varepsilon_y(\mathbf H)
       \le L\mathsf C e^{\mathsf C\bar\kappa}\mathsf b<\epsilon.
\]
All tops are levelled. By \cref{II-85-34},
\(u^\delta\) is real and cancels the accumulated factor
\(v_{\mathbf B}(x)\) at every top \(x\); at a bare top
both factors are one.

The cancellation at an outward row deserves attention.
The terminal double-block factorization is
\begin{equation}\label{II-85-42}
 M_{\mathbf B}(E)(c)
       =v_q(c_-)\,e^{iB'(c)}V_0(c)\,v_q(c_+)^{-1}.
\end{equation}
At an endpoint of level \(q\), \(v_q=v_{\mathbf B}\).
At the outer endpoint \(z\) of an outward row, the endpoint
level is \(q-1\), and
\[
 v_q(z)=v_{\mathbf B}(z)\,
        \mathsf v_{\mathsf m}(\mathcal M^{q-1}U_0,a_{q-1})(z).
\]
Every bond inside \(B(z)\) is an interior row of
\(\Gamma_{q-1}\). Restriction of the same recursion
identifies \(a_{q-1}\) there with \(B'\), which is zero
because \(q-1<k\). The last factor is therefore exactly
one. Thus at either endpoint of every row the factors in
\eqref{II-85-42} are cancelled by \(u^\delta\).
Real domain covariance gives
\(M_{\mathbf B}(U^\delta)=e^{iB'}V_0\).
Gauge invariance preserves criticality. Current covariance
and the second identity of \eqref{eq:pinned-shift-identities}
give \eqref{II-94-51}; the pin identity gives the stored
transports.

Finally, the regularity calculation in the proof of
\cref{II-93-38} applies with the initial natural-class
radius \(4\alpha\). The potential and its first difference
are bounded at layer \(n\) by
\(\mathsf C e^{\mathsf C\bar\kappa}\mathsf b\) after
multiplication by \(\xi_n,\xi_n^2\); the current has the
corresponding \(\xi_n^3\) bound. A plaquette or bond
meeting a higher layer loses at most one power of \(L\)
in that scale comparison. The plaquette variation estimate
and the current identity thus bound the regularity radius
by \(\mathsf C(4\alpha+e^{\mathsf C\bar\kappa}\mathsf b)\).
The constants of the regional response already absorb
these \(L^2,L^3\) losses. Regularity is unchanged by the
real gauge \(u^\delta\), which proves the final assertion.
\end{proof}

\subsection{Tower bounds in the ambient layer scale}
\label{sec:history-ambient-tower}

The tower used by a local term can end before the ambient
history does. Its local block size is then smaller than the
size against which the historical domain measures curvature.
We retain the scale factor in the next estimate. This is the
input needed to preserve the tower conditions under a shift.
All determining sets in this subsection are reduced.

Fix an ambient separated sequence
\(\Omega^\flat_0\supset\cdots\supset\Omega^\flat_i\)
and a subordinate separated sequence
\(\Omega'_0\supset\cdots\supset\Omega'_p\), where
\(1\le p\le i\) and \(\Omega'_q\subset\Omega^\flat_q\).
Write \(Z=\Omega'_0\), \(\mathsf Z=\operatorname{Bd}^\circ(Z)\),
\(\xi_q=L^q\eta\), and
\[
\begin{gathered}
 \lambda^\flat(x)=\max\{q\le i:x\in\Omega^\flat_q\},\qquad
 \lambda'(x)=\max\{q\le p:x\in\Omega'_q\},\\
 l(x)=\xi_{\lambda^\flat(x)},\quad
 \xi'(x)=\xi_{\lambda'(x)},\quad
 \varsigma(x)=l(x)/\xi'(x).
\end{gathered}
\]
The two weighted distances are denoted by \(d^\flat,d'\).
An inner plaquette or full bond \(e\) of \(Z\) is called
settled when
\(\lambda'(x_e)=\min\{\lambda^\flat(x_e),p\}\).
The base of a positively oriented bond is its initial point.
All differences below are taken only on stencils in
\(\mathsf Z\).

For a real configuration \(U\), a complex form \(w\), and
\(a>0\), let \(\mathcal Q^{\rm T}(a)\) be the domain
\begin{equation}\label{I-95-27}
 \epsilon_U(e):=\eta^{-2}l(x_e)^2|\partial U(e)-I|<a,
 \qquad
 \mathsf n_b(w):=\max\{l(x_b)|w(b)|,
                 l(x_b)^2|\nabla_Uw|_{\mathsf Z}(b)\}<a.
\end{equation}
Put
\[
 E_U^\delta(x)=\max_e e^{-\delta d^\flat(x,x_e)}\epsilon_U(e),
 \qquad
 E_w^\delta(x)=\max_b e^{-\delta d^\flat(x,x_b)}\mathsf n_b(w).
\]
The tower configuration is
\[
 U^{\rm tw}(U,w)=\Phi_{V_0}(B(w)),\qquad
 V_0=M_{\mathbf B'}U,\qquad
 B(w)(c)=Q_q(U;\eta w)(c),\quad c\in\Gamma'_q.
\]
Thus the datum is the actual average of \(e^{i\eta w}U\),
expressed relative to \(V_0\); there is no replacement of
this average by its linearization.

We specify constants used later in the radii selection.
Take the real-background and continuation thresholds of
\cref{I-58-22}, the local regularity threshold and profile
constants from the proof of \cref{I-56-72}, and the response
constants of \cref{I-82-15}. Assume their geometric lower
bounds on \(\mathfrak M\) and \(\rho=R_1\mathfrak M\),
and that the top torus has at least \(6\rho\) sites in each
direction. Take
\[
\begin{gathered}
 \delta=\min\{\delta_{\rm c},\delta_{\rm tw}\},\qquad
 \delta\rho\ge64\log L,\\
 \mathsf S(a)=54L^3\sum_{r\ge1}(2r+1)^3e^{-a(r-1)/2},
 \qquad B^\sharp_{\rm cx}=\max\{1,B_5\mathsf S(\delta)\}.
\end{gathered}
\]
Here \(\delta_{\rm c}\) is the local regularity rate and
\(B_5,\delta_{\rm tw}\) are the constants in
\eqref{I-93-47}. Let
\[
 c^\natural=2L^2,\qquad
 a^{(\rm ia)}_{\rm tw}=\min\{a_{\rm plain},c_{\rm plain}\},\qquad
 \rho_{\rm f}=\min\{c_{\rm plain},[5(\kappa_6+3)]^{-1}\}.
\]
Here \(a_{\rm plain},c_{\rm plain}\) are the constants
of \cref{I-27-78}; decrease \(a_{\rm plain}\), if needed,
so that \(a_{\rm plain}\le c_2^{\rm SO}\).
Choose \(\varepsilon_{\rm T}\) and \(r_{\rm T}\) by
\begin{align}
 \varepsilon_{\rm T}
 &=\min\left\{\frac{\bar\varepsilon}{2},
       \varepsilon^{\rm s}_{\rm G}(\rho_{\rm s}),
       \frac{a_{\rm R}}{C_{\mathfrak P}\mathfrak M B_3},
       \frac{a_{{\rm c},{\rm LI}}}{B_3},
       \frac{a_{\rm plain}}{B_3},\frac1{22L^3B_3}\right\},
       \label{I-95-39}\\
 r_{\rm T}
 &=\min\left\{\frac{r^{\rm s}_{\rm B}(\rho_{\rm s})}{2},
       \frac{\varrho^0_{\rm R}}2,
       \frac{\bar\varepsilon}{2C_\partial},
       \frac{\rho_{\rm f}}{2LB^\sharp_{\rm cx}},
       \frac1{600L^3B^\sharp_{\rm cx}},\frac1{20}\right\}.
       \label{I-95-40}
\end{align}
The entries \(a_{\rm R},\varrho^0_{\rm R}\) are the
regularity and datum thresholds for the response in
\cref{I-82-15}; \(C_{\mathfrak P}\mathfrak M\) is its
regularity conversion factor. The entries
\(a_{{\rm c},{\rm LI}},\varepsilon^{\rm s}_{\rm G},
r^{\rm s}_{\rm B},C_\partial\) are the local regularity,
continuation, and tested-holonomy constants just selected.
Thus these choices place the same real critical family
simultaneously in the local regularity, continuation, and
response domains.

Set
\begin{equation}\label{I-95-41}
 \begin{split}
 C_{\rm R}&=\mathsf c_{\rm loc}c^\natural L^4e^{7L\delta/2},
 \qquad C_{\rm X}=900L^3B_5e^{7\delta/2},\\
 C_{\rm h}&=\frac{15B_5}{4\rho_{\rm f}}e^{3\delta/2},
 \qquad C_\chi=1+L^2,\\
 r_{\rm w}&=\min\{r_{\rm T}/(\mathsf c_QL),
                         a^{(\rm ia)}_{\rm tw}/(2L)\},
 \qquad\delta_{\rm T}=\delta/4,
 \end{split}
\end{equation}
where \(\mathsf c_Q=C_Q\) is the averaging constant in
\cref{I-27-78}. Finally, let
\begin{equation}\label{I-95-57}
 \begin{split}
 C_{\rm far}&=\frac{8\mathsf c_QL^3}{\delta}
          e^{3(5L+1)\delta/8}\mathsf S(\delta/4),\\
 K_{\rm a}&=7C_{\rm R}
      +(6C_{\rm R}+5L^2/r_{\rm w})C_\chi e^{9L\delta/4}
      +(C_{\rm X}+C_{\rm h})C_{\rm far},\\
 K_{\rm T}&=(C_{\rm X}+C_{\rm h})C_{\rm far}
                  +4K_{\rm a}C_\chi e^{9L\delta/4},\\
 \bar a&=\min\{\varepsilon_{\rm T}/(5c^\natural),
          a^{(\rm ia)}_{\rm tw}/(5L^2),r_{\rm w}/2,
          (11L^2)^{-1},(2K_{\rm a})^{-1}\},
 \qquad a_{\rm T}=\frac{\bar a}{5C_\chi^2}.
 \end{split}
\end{equation}
All constants are independent of the depth, the volume, and
the number of layers. The rate \(\delta_{\rm T}\) depends
only on \(L\); \(K_{\rm a},K_{\rm T}\) depend only on
\(L,\rho_{\rm s},\bar\varepsilon\); and \(a_{\rm T}\)
may also depend on \(\mathfrak M,B_3\). In particular,
\(0<\delta_{\rm T}\le1/4\),
\(\delta_{\rm T}\rho\ge16\log L\),
\(K_{\rm T}\ge K_{\rm a}\ge1\), and
\(a_{\rm T}\le(64L^2)^{-1}\).

\begin{proposition}\label{I-95-61}
On \(\mathcal Q^{\rm T}(a_{\rm T})\), the tower and its
current are defined and, for fixed \(U\), holomorphic in
\(w\). At a settled element, put \(n=\lambda^\flat(x_e)\),
\(n'=\lambda'(x_e)=\min\{n,p\}\), and
\[
 G(U,w;e)=
 \begin{cases}
 L^{2n}(\partial U^{\rm tw}(e)-I),&e\text{ a plaquette},\\
 L^{2n+n'}\eta^3J(U^{\rm tw})(e),&e\text{ a full bond}.
 \end{cases}
\]
Then
\begin{align}
 |G(U,w;e)|
   &\le K_{\rm a}\{E_U^{\delta_{\rm T}}(x_e)
                         +E_w^{\delta_{\rm T}}(x_e)\}\le1,
                         \label{eq:ambient-tower-size}\\
 |G(U,w_1;e)-G(U,w_0;e)|
   &\le K_{\rm T} E_{w_1-w_0}^{\delta_{\rm T}}(x_e).
                         \label{I-95-62}
\end{align}
At real \(w\), the tower is the real axial representative
of its actual averaged datum. It respects gauge changes,
with the gauge at each pillar extended from its top, and
complex conjugation.
\end{proposition}

\begin{proof}
We first record the geometry that permits the comparison of
the two scales. The definitions of the weighted lengths and
separation give
\begin{equation}\label{eq:ambient-tower-geometry}
 \begin{gathered}
 d'(x,z)\ge d^\flat(x,z),\qquad
 d'(x,z)\ge\xi_p^{-1}|x-z|_\infty,\\
 |\lambda^\flat(x)-\lambda^\flat(z)|
       \le1+d^\flat(x,z)/\rho,\qquad
 |\lambda'(x)-\lambda'(z)|\le1+d'(x,z)/\rho,\\
 \varsigma(x)/\varsigma(z)\le L^{2+2d'(x,z)/\rho}.
 \end{gathered}
\end{equation}
For the first inequality, subdivide a polygonal path at
the faces of both layer partitions. Away from shared faces,
its tower density is at least the ambient density. On a
shared face, move the segment an arbitrarily small distance
into a bare cube of the deeper ambient layer, and join its
endpoints by two short segments. The added weighted length
tends to zero. Taking the path infimum proves the assertion,
including paths running along faces. The lower bound by
\(\xi_p^{-1}|x-z|_\infty\) follows from the smallest tower
density. To cross from layer \(m\) to layer \(m'>m+1\), a
connected path crosses, for each \(m<q<m'\), a strip of
width \(\rho\xi_q\) in layer \(q\). Each strip contributes
at least \(\rho\) to its weighted length. This proves the
two index bounds and then the ratio bound.

If \(|x-z|_\infty\le3l(x)\), separation also gives
\(|\lambda^\flat(x)-\lambda^\flat(z)|\le1\) and
\(d^\flat(x,z)\le9L\). A coordinate path stays in the
preceding layer and proves the latter inequality. A site
and its tower top are within distance \(3/2\). If
\(b\in\mathsf P^q(c)\), the two endpoint blocks of \(c\)
lie in \(\Omega'_{q-1}\), and their coordinate paths give
\begin{equation}\label{eq:ambient-row-geometry}
 d'(x_b,y_c)\le5L/2,\qquad \xi_q\le L\xi'(x_b).
\end{equation}

Here is the real estimate. If
\(\epsilon_W<\alpha\),
\(L^2\alpha\le a^{(\rm ia)}_{\rm tw}\), and
\(c^\natural\alpha\le\varepsilon_{\rm T}\), curvature
averaging makes \(M_{\mathbf B'}W\)
\(c^\natural\alpha\)-regular. More precisely, a tested
plaquette of order \(q\) averages bare plaquettes in
\(\Omega'_{q-1}\). On its four-corner block union their
curvatures obey
\(L^{2q}|\partial W-1|
 \le L^2\max\varsigma^{-2}\epsilon_W
 <a^{(\rm ia)}_{\rm tw}\le c_2^{\rm SO}\).
The local four-corner induction used in
\cref{II-195-31}, at terminal level \(q\), therefore
bounds its holonomy by
\(c^\natural\max\varsigma^{-2}\epsilon_W\) there.
These plaquettes are within tower distance \(7L/2\) of a
corner of the test. The profile estimate in the proof of
\cref{I-56-72}, followed by
\eqref{eq:ambient-tower-geometry}, therefore gives
\begin{equation}\label{I-95-45}
 \varsigma(x_e)^2\mathsf e_{U_0}(e)
 \le C_{\rm R}\max_f
              e^{-\delta d'(x_e,x_f)/2}\epsilon_W(f).
\end{equation}
Here \(U_0\) is the real critical representative and
\(\mathsf e_{U_0}\) measures curvature and current in the
tower block scale. Indeed, the square of the scale ratio
costs at most \(L^4e^{4(\log L)d'/\rho}\); the condition
\(\delta\rho\ge64\log L\) leaves the displayed decay.

Apply this to \(W=U\). By \eqref{I-95-39},
\(V_0=M_{\mathbf B'}U\) lies in the continuation and
response domains. Integrating \eqref{I-93-47} gives, with
\(\Sigma(\beta;x)=\sum_c e^{-\delta d'(x,y_c)}|\beta(c)|\),
\begin{equation}\label{eq:ambient-Landau-difference}
 N_b(\mathcal H(B_1)-\mathcal H(B_0))
 \le\min\{B^\sharp_{\rm cx}\sup|B_1-B_0|,
       B_5e^{\delta d'(x,y_b)}\Sigma(B_1-B_0;x)\}.
\end{equation}
Here \(y_b\) is the pillar top of the initial point of
\(b\).
The norm \(N_b\) includes \(\xi'|\mathcal H|\),
\((\xi')^2|\nabla\mathcal H|\), and, at full bonds,
\((\xi')^3|D^{\mathsf T}D\mathcal H|\).

The gauge factor in
\(\Phi(B)=(e^{i\eta\mathcal H(B)}U_0)^{h(B)}\)
depends, at a site, only on the potential in its pillar.
For a pillar of level \(q\), its product of ascent factors
has total exponent norm at most
\[
 3(L-1)\tau\sum_{m<q}L^{m-q}
       +2\kappa_6\tau\sum_{m<q}L^{m-q}
 \le(\kappa_6+3)\tau,\qquad
 \tau=\xi_q\sup_{\mathsf P(B^q(y))}|\mathcal H|.
\]
The first sum counts the comb edges; the second counts each
averaging correction once. Thus, on
\(|\mathcal H|_{(-1)}<\rho_{\rm f}\),
\(|h-I|<1/4\), \(|h^{-1}|<5/4\), and the condition
number of \(h\) is less than \(3/2\). Applying Cauchy's
estimate to this local pillar map, with radius
\(\rho_{\rm f}/2\), yields
\[
 |h(B_1)(x)h(B_0)(x)^{-1}-I|
 \le\min\{5/8,\ C_{\rm h}\Sigma(B_1-B_0;x)/3\}.
\]
The site-to-top displacement costs at most \(e^{3\delta/2}\),
already included in \(C_{\rm h}\).

Let \(\mathcal F(B)\) be curvature or current in the tower
scale at \(e\), so that \(G=\varsigma(x_e)^2\mathcal F\).
The base curvature and current in each local plaquette star
are at most \(1/22\), by \eqref{I-95-39}. Express the
potential in that local scale and put
\(N_{\rm loc}(f)=\max\{|f|,|\nabla f|,|D^{\mathsf T}Df|\}\),
omitting the last member for curvature. The expansions
\eqref{I-10-121} and \eqref{I-10-122} show that the
unconjugated curvature or current has norm less than one
on \(N_{\rm loc}(f)<1/300\). At level zero the same bound
follows directly from the four exponential factors, giving
curvature at most \(1/22+5/300\) and current at most
\(4e^{2/300}(1/22+5/300)<1\).
Cauchy's estimate on the half ball gives Lipschitz constant
\(600\). The passage from \(N_b\) to the local star costs
at most \(L^3\), and its diameter costs \(e^{7\delta/2}\).
Combining this with the preceding gauge estimate gives
\begin{equation}\label{I-95-55}
 |\mathcal F(B)|<3/2,\qquad
 |\mathcal F(B_1)-\mathcal F(B_0)|
 \le[C_{\rm X}+C_{\rm h}|\mathcal F(B_0)|]
                      \Sigma(B_1-B_0;x_e).
\end{equation}

We now pass from the datum to the original potential.
Write \(\|w\|'=\max_b\mathsf L_-(b)|w(b)|\), where
\(\mathsf L_-\) is the smaller endpoint block side in the
tower. The averaging theorem \cref{I-27-78} and
\eqref{eq:ambient-row-geometry} show that \(B(w)\) is
holomorphic on \(\|w\|'<r_{\rm w}\), with
\(\sup|B(w)|<r_{\rm T}\). Cauchy's estimate on the
half averaging ball gives
\begin{equation}\label{I-95-63}
 |B(w_1)(c)-B(w_0)(c)|
 \le2\mathsf c_QL\max_{b\in\mathsf P^q(c)}
       \varsigma(x_b)^{-1}l(x_b)|w_1(b)-w_0(b)|.
\end{equation}
This also proves existence and holomorphy of the actual
tower on the domain in the statement.

Fix a settled \(e\), and write \(x=x_e\), \(l=l(x)\),
\(\varsigma=\varsigma(x)\), and \(F(w)=G(U,w;e)\).
Let \(\chi(b)\) be one for \(|x_b-x|_\infty\le l\),
zero for \(|x_b-x|_\infty\ge2l\), and linear between.
Call bonds within distance \(3l\) near. The local geometry
above and \(|\chi(b^+)-\chi(b)|\le\eta/l\) give
\begin{equation}\label{I-95-65}
 \mathsf n(\chi v)\le C_\chi\mathsf n^{\rm nr}(v),
 \qquad \|\chi v\|'\le L\varsigma^{-1}\mathsf n(\chi v),
 \qquad
 \mathsf n^{\rm nr}(v)\le e^{9L\delta/4}E_v^{\delta_{\rm T}}(x).
\end{equation}
For real \(g\) with \(\mathsf n(g)\le\bar a\),
\eqref{I-10-121} gives
\(\epsilon_{e^{i\eta g}U}\le2\epsilon_U+3\mathsf n_p(g)\).
Consequently \eqref{I-95-45} gives
\[
 |F(g)|\le C_{\rm R}
          \{2E_U^{\delta_{\rm T}}(x)+3\mathsf n(g)\},
 \qquad |F(0)|\le C_{\rm R}E_U^{\delta_{\rm T}}(x).
\]
For a holomorphic map bounded by \(M_0\) on a radius-\(r\)
ball, Taylor's formula and Cauchy's estimate imply
\begin{equation}\label{I-95-60}
 |F(z)|\le3|F(0)|+|F(\Re z)|+|F(\Im z)|
                         +6M_0\|z\|^2/r^2
 \quad(\|z\|\le r/2).
\end{equation}
Indeed, each remainder after the linear term is bounded by
\(2M_0\|z\|^2/r^2\); substitute the real and imaginary
parts into the complex-linear first derivative. Apply this
with \(M_0=3\varsigma^2/2\), \(r=r_{\rm w}\), and
\(z=\chi w\). The factor \(\varsigma^2\) cancels against
\(\|\chi w\|'^2\). Since
\(\mathsf n(\chi w)\le\bar a\le r_{\rm w}/2\), we obtain
\begin{equation}\label{I-95-67}
 |F(\chi w)|\le7C_{\rm R}E_U^{\delta_{\rm T}}(x)
       +(6C_{\rm R}+5L^2/r_{\rm w})C_\chi
                                  \mathsf n^{\rm nr}(w).
\end{equation}

For a change \(u=w-\widetilde w\) supported outside the
distance-\(l\) neighborhood, put \(D=d'(x,x_b)\).
Settlement gives \(\varsigma\le1+D\): if
\(\varsigma>1\), the tower has ended at \(p\), and
\(D>l/\xi_p=\varsigma\). Moreover,
\[
 \frac{\varsigma}{\varsigma(x_b)}\le L^2e^{\delta D/32},
 \qquad d'(x,y_c)\ge D-5L/2.
\]
Use these inequalities in \eqref{I-95-63} and
\eqref{I-95-55}. Split off one quarter of the exponential
for summation and use
\((1+D)e^{-\delta D/4}\le4/\delta\). The remaining
exponential is at least the decay \(\delta/4\) required
for the ambient profile. The row sum in \cref{I-10-125},
with the site-to-top cost \(e^{3\delta/8}\), yields
\begin{equation}\label{I-95-68}
 |F(w)-F(\widetilde w)|
 \le[C_{\rm X}+C_{\rm h}|F(\widetilde w)|]C_{\rm far}
       \max_b e^{-\delta_{\rm T}d'(x,x_b)}l(x_b)|u(b)|.
\end{equation}
For \(C_\chi\mathsf n(w)\le\bar a\),
\eqref{I-95-67} is at most one. Apply
\eqref{I-95-68} with \(\widetilde w=\chi w\), and then
\eqref{I-95-65}. The definition of \(K_{\rm a}\) gives
\eqref{eq:ambient-tower-size}.

Finally write \(v=w_1-w_0\) and split the change of \(F\)
into the change from \(w_0+\chi v\) to \(w_1\), and the
change from \(w_0\) to \(w_0+\chi v\). The first is
bounded by \((C_{\rm X}+C_{\rm h})C_{\rm far}E_v^{\delta_{\rm T}}(x)\).
For the second let
\(n_v=C_\chi\mathsf n^{\rm nr}(v)\) and
\(E=E_U^{\delta_{\rm T}}(x)+E_{w_0}^{\delta_{\rm T}}(x)\).
If \(n_v=0\) it vanishes. Otherwise the disc
\(|z|\le1+E/n_v\) satisfies
\(C_\chi\mathsf n(w_0+z\chi v)\le5C_\chi^2a_{\rm T}=\bar a\).
The bound already proved on this disc, followed by the
one-variable Schwarz estimate for
\(F(w_0+z\chi v)-F(w_0)\), gives a bound
\(4K_{\rm a}n_v\). If \(E=0\), the same bound follows
directly at the endpoints. Equation~\eqref{I-95-65} and
the definition of \(K_{\rm T}\) prove \eqref{I-95-62}.
The reality and covariance statements follow first for
real data from uniqueness of the axial representative,
and then throughout the connected complex domain by
holomorphy.
\end{proof}

\begin{corollary}\label{I-95-81}
Let \(\Omega^\flat_\bullet\) have depth \(t\), let
\(1\le p\le i\le t\), and let \(X\) be a region for
\(\pi_0\). Assume \(M_1\ge\rho\),
\(M=L^{m'}\ge L^3M_1\), and \(m'\le m+K-p+1\).
Use either tower rule and localize by
\(\Omega'_q=\Omega^\flat_q\sqcap X_{[q]}\), \(q\le p\).
Put \(n_i(x)=\min\{\lambda_{\Omega^\flat_\bullet}(x),i\}\)
and \(l_i(x)=\eta L^{n_i(x)}\). On the domain
\eqref{I-95-27} with \(l=l_i\), the localized real and
complex towers exist. If an inner plaquette or full bond
\(e\) has its base within distance \(2\xi_p\) of
\(X^{\wr p}\), define
\[
 G_{p,X}(U,A';e)=
 \begin{cases}
 L^{2n_i(x_e)}(\partial U_{p,X}(e)-I),&e\text{ a plaquette},\\
 L^{2n_i(x_e)+\min\{p,n_i(x_e)\}-3p}J_{p,X}(e),
                                     &e\text{ a full bond},
 \end{cases}
\]
where \(J_{p,X}\) is in level-\(p\) units. Then
\begin{equation}\label{I-95-84}
 \begin{split}
 &|G_{p,X}(U,A'_1;e)-G_{p,X}(U,A'_0;e)|\\
 &\quad\le e^{3\delta_{\rm T}}K_{\rm T}
       \max_b e^{-\delta_{\rm T}d_{\mathbf B^\flat}(y_e,y_b)}
       \max\{l_i(x_b)|A'_1-A'_0|(b),
                   l_i(x_b)^2|\nabla_U(A'_1-A'_0)|_{\mathsf Z}(b)\}.
 \end{split}
\end{equation}
The tops \(y_e,y_b\) and the distance in this formula belong
to the full ambient sequence of depth \(t\), while
\(\mathsf Z=\operatorname{Bd}^\circ(\Omega^\flat_0\sqcap X)\).
In particular the estimate applies when a corner of \(e\)
lies in \(X^{\wr p}\).
\end{corollary}

\begin{proof}
The localized sequence is separated because its complement
at level \(q\) is the union of the two complementary
regions. Under the first tower rule the accumulated inward
erosion to level \(q\) is at most
\(2\rho\sum_{r<q}\xi_r\le M_1\xi_q\); under the second
it is at most
\(2M_1\sum_{1\le r\le q}\xi_r\le3M_1\xi_q\).
On the other hand,
\(\operatorname{dist}_\infty(X^{\wr p},X^c)\ge2M\xi_{p-1}\).
Thus a site within \(2\xi_p\) of this interior remains in
every \(X_{[q]}\), \(q\le p\), since
\((2M/L-2)\xi_p>3M_1\xi_p\).
At such sites the localized and ambient memberships agree,
so the element is settled. The scale-invariant current is
\(\eta^3J\), which gives the displayed power of \(L\).
Apply \eqref{I-95-62}. Passing to the deeper ambient
sequence can only decrease the weighted distance. Moving
each endpoint to its ambient pillar top costs at most
\(3/2\), and therefore costs \(e^{3\delta_{\rm T}}\).
A corner and the base of an element are within distance
\(\eta\le2\xi_p\), proving the last assertion.
\end{proof}

\subsection{The margin for substitution into an earlier domain}
\label{sec:history-graded-substitution}

Write \(K^\sharp=K^\sharp_{\rm tw}\), and use \(C_{\rm os}=18\) from \cref{II-16-69}.
Let \(t=k+1\ge2\). In this subsection the radii table is
independent of the structure: its entries are
\(\mathsf r_\iota(n)\), \(\iota=0,1\). Assume
\begin{equation}\label{eq:history-substitution-radii}
 \begin{gathered}
 \mathsf r_0(n)\le\mathsf r_1(n),\qquad
 \mathsf r_1(n)^2\le\mathsf c_\natural\mathsf r_0(n),
 \qquad \mathsf c_\natural\ge1,\\
 \mathsf r_\iota(n)\le\mathsf r_\iota(n+1),\qquad
 L^{-\nu}2^\nu\mathsf r_\iota(n+\nu)\le\mathsf r_\iota(n),
 \qquad (1+\beta_{\rm s})\mathsf r_\iota(n)\le\hat a.
 \end{gathered}
\end{equation}
The indices range over the entire table, with \(\nu\ge1\)
and \(n+\nu\) in its range. These are the conditions
(R0), (R1), (Q5.2), and (Q6). Let
\(n_i(x)=\max\{q\le i:x\in\Omega^\diamond_q\}\),
where \(\Omega^\diamond_0=\Omega_1^\sim\), and put
\begin{equation}\label{II-88-12}
 \mathsf c_0=\tfrac12(1-\theta_{\rm s})\beta_{\rm s},\qquad
 \mathsf w(n)=\mathsf c_0\,2^{-(k+1-n)}\mathsf r_0(n),
 \qquad \mathsf v_k(b)=\mathsf w(n_k(b_-)).
\end{equation}
For a level-\(j\) pair \(\mathfrak h=(\mathbf A,\mathbf K)\),
\(\mathbf U=e^{i\eta_jA'}U\), and a positive profile
\(\mathsf v\), set
\begin{equation}\label{II-88-15}
 \|\mathfrak h\|^{\mathcal S}_{\mathsf v,\mathbf U}
 =\max_{b\in\operatorname{Bd}_{\mathcal S}(X)}
 \frac{\max\{L^{n_k(b)}\eta_j|\mathbf A(b)|,
 L^{2n_k(b)}\eta_j^2|\nabla_{\mathbf U}\mathbf A|(b),
 L^{3n_k(b)}\eta_j^3|\mathbf K(b)|\}}{\mathsf v(b)}.
\end{equation}
Only positively oriented bonds enter the maximum, and only
differences defined in \(\operatorname{Bd}_{\mathcal S}(X)\)
enter its second member. In particular this norm imposes
no condition outside that bond set.

Write \(\sigma^{\rm m}=(t,{\rm mid})\). For
\(\sigma'\preccurlyeq\sigma^*=(t^*,\phi^*)
\preccurlyeq(k,{\rm out})\), the coefficients in
\eqref{II-2-9} satisfy
\begin{equation}\label{II-88-21}
 [\mathsf c_{\sigma^*}(n)-\mathsf c_{\sigma^{\rm m}}(n)]
                  \mathsf r_\iota(n)\ge2\mathsf w(n)
                         \qquad(n\le t^*).
\end{equation}
Indeed their difference is at least
\(\beta_{\rm s}2^{-(t-n)}(2^{t-t^*}-1-\theta_{\rm s})
\ge(1-\theta_{\rm s})\beta_{\rm s}2^{-(t-n)}\).
The same gap holds for the pointwise right members:
\begin{equation}\label{II-88-22}
 \mathsf R^{p,\iota}_{\sigma^*}(x)
      -\mathsf R^{p,\iota}_{\sigma^{\rm m}}(x)
 \ge2\mathsf w(n_{t^*}(x))L^{-pn_{t^*}(x)},
                       \qquad p=1,2,3.
\end{equation}
If \(n_t(x)\le t^*\), this is \eqref{II-88-21}. Otherwise
put \(n_t(x)=t^*+\nu\). The smaller-scale right member is
at most \((5/4)2^{-\nu}L^{-(p-1)\nu}\mathsf r_\iota(t^*)
L^{-pt^*}\), whereas the other is at least
\(\mathsf r_\iota(t^*)L^{-pt^*}\). Their difference is
at least \(3\mathsf r_\iota(t^*)L^{-pt^*}/8\), and
\(2\mathsf w(t^*)\le\mathsf r_0(t^*)/8\).

The profile is increasing, satisfies
\(\mathsf w(n-1)\le\mathsf w(n)/2\), and obeys
\[
 \mathsf w(n+\nu)L^{-p(n+\nu)}
     \le\mathsf w(n)L^{-pn},\qquad p=1,2,3.
\]
Consequently the norm bound with index \(n_k\) also holds
with every truncated index \(n_i\), \(i\le k\).
Furthermore, with \(\omega_{\rm T}=\delta_{\rm T}\rho\),
\begin{equation}\label{eq:history-profile-decay}
 e^{-\omega_{\rm T}(\nu-1)}\mathsf w(n+\nu)
          \le C^*\mathsf w(n),\qquad C^*=L.
\end{equation}
This follows from (Q5.2) and \(\omega_{\rm T}\ge16\log L\).
If, for some \(C_\uparrow\ge1\),
\[
 2^\nu e^{-\omega_{\rm T}(\nu-1)}\mathsf r_0(n+\nu)
       \le C_\uparrow\mathsf r_0(n)
 \qquad(n\ge0,\ \nu\ge1,\ n+\nu\le k),
\]
where \(n,\nu\) are integers, one may use
\(C^*=\min\{L,C_\uparrow\}\).

\begin{proposition}\label{II-88-43}
Assume the geometric hypotheses of \cref{I-95-81} and of
the outer-core substitution \cref{II-16-69},
\(M\ge L^4M_1\), and
\eqref{eq:history-substitution-radii}.
Let \(\sigma'=(j,\phi')\preccurlyeq\sigma^*
\preccurlyeq(k,{\rm out})\), \(j\ge1\),
\(X\in\mathcal D_{j-1}\), and
\(P\in\widetilde U_{\sigma'|\sigma^{\rm m}}(X)\).
For a shift \(\mathfrak h\) put
\(r=\|\mathfrak h\|^{\mathcal S}_{\mathsf v_k,\mathbf U}\).
Suppose
\begin{equation}\label{II-88-33}
 \hat a\le a_{\rm T},\quad 4\hat a\le a^\sharp,
 \qquad r\le\tfrac12,\qquad4C^*K_{\rm T}r\le1,
 \qquad L^2C_{\rm os}K^\sharp(1+\mathsf c_\natural)r\le3,
 \qquad L^3r\le4.
\end{equation}
Then substitution is defined on
\(\operatorname{Bd}_{\mathcal S}(X)\). If it is defined
on all of \(\operatorname{Bd}(X)\), then
\[
 P\dotplus\mathfrak h\in
                    \widetilde U_{\sigma'|\sigma^*}(X).
\]
\end{proposition}

\begin{proof}
Put \(\eta=\eta_j\),
\(\mathfrak e=\operatorname{Bd}_{\mathcal S}(X)\), and
\[
 A''=(i\eta)^{-1}\log(e^{i\eta\mathbf A}e^{i\eta A'}).
\]
We first control this logarithm and its first differences.
For \(\ell(X,Z)=\log(e^Xe^Z)-X-Z\), the exponential and
logarithm series give
\(|\ell(X,Z)|\le2(|X|+|Z|)^2\) when
\(|X|+|Z|\le1/4\). Cauchy's estimate on the ball of radius
\(|X|+|Z|\) about \((X,Z)\) gives a derivative bound
\(8(|X|+|Z|)\) on the half ball. Since \(\ell(0,Z)=0\),
the one-variable Schwarz estimate also gives
\[
 |\ell(X,Z)-\ell(X,Z')|\le16|X|\,|Z-Z'|
                  \quad(|X|,|Z|,|Z'|\le1/16).
\]
For the latter assertion apply Schwarz on \(|z|<1/16\)
to \(\ell(zX/|X|,Z)-\ell(zX/|X|,Z')\), bounded there
by \(|Z-Z'|\).

These estimates imply
\begin{equation}\label{II-88-29}
 \eta|A''-A'|(b)
 \le[1+8\eta(|\mathbf A(b)|+|A'(b)|)]\eta|\mathbf A(b)|.
\end{equation}
For a difference based at \(x\), write
\(b=\langle x,x+\eta e_\nu\rangle\),
\(e=\langle x,x+\eta e_\mu\rangle\), and
\(b^+=b+\eta e_\mu\). With
\(s=\eta|A'(b^+)|+\max\{\eta|\mathbf A(b)|,
\eta|\mathbf A(b^+)|\}\), transport all terms by \(U(e)\).
Splitting the change in \(\ell\) first in its first
variable and then in its second gives
\begin{equation}\label{II-88-30}
 \begin{split}
 \eta^2|\nabla_U(A''-A')|(b)
 &\le(1+8s)\{\eta^2|\nabla_{\mathbf U}\mathbf A|(b)
       +(e^{2\eta|A'(e)|}-1)\eta|\mathbf A(b^+)|\}\\
 &\quad+16\eta|\mathbf A(b)|\eta^2|\nabla_UA'|(b).
 \end{split}
\end{equation}
The exponential term changes transport from \(\mathbf U(e)\)
to \(U(e)\). This derivation includes \(\mu=\nu\).

Now \(\hat a\le(64L^2)^{-1}\), and
\(\eta|\mathbf A|\le r\hat a/16\) on \(\mathfrak e\).
Thus the logarithm is defined there. Neighboring sites
have layer indices differing by at most one, by separation.
Using the profile relations above in
\eqref{II-88-29} and \eqref{II-88-30}, for every \(i\le k\)
and \(n=n_i(x)\), gives
\begin{equation}\label{eq:history-substitution-increments}
 \begin{gathered}
 L^n\eta|A''-A'|\le\tfrac32r\mathsf w(n),\qquad
 L^{2n}\eta^2|\nabla_U(A''-A')|
                                  \le\tfrac32r\mathsf w(n),\\
 |\partial(e^{i\eta\mathbf A}\mathbf U)-I|
 \le|\partial\mathbf U-I|+2.03r\mathsf w(n)L^{-2n}.
 \end{gathered}
\end{equation}
For clarity, the coefficient for the derivative is bounded
by
\((1+10^{-3})(1+1.05L\hat a)+16\hat a<3/2\):
the neighboring shift is at most
\((L/2)r\mathsf w(n)L^{-n}\), and
\(e^{2\eta|A'(e)|}-1\le2.1\hat a L^{-n}\).
For the last line use \eqref{I-10-121} with condition
number \(2\). The linear term is at most
\(2r\mathsf w(n)L^{-2n}\); its two remainders are at most
\((11/1024+6/832)r\mathsf w(n)L^{-2n}\).
Multiplication by \(|\partial\mathbf U|\le1+\hat a\)
gives \(2.03\).

Assume now that substitution is defined on all bonds of
\(X\). The resulting triple is
\(P''=(U,A'',\mathbf J+\mathbf K)\). The first three
increments in (L1) are bounded, in their local units, by
\(r\mathsf w\), \(3r\mathsf w/2\), and
\(3r\mathsf w/2\); the curvature increment is bounded by
\(2.03r\mathsf w\). Since \(r\le1/2\),
\eqref{II-88-22} absorbs all of them. Real curvature is
unchanged. This proves (L1), including
\(\eta|A''|<\hat aL^{-n_{t^*}}\) and
\(\eta^2|\nabla_UA''|<\hat aL^{-2n_{t^*}}\).
The real gauges in (L2) are unchanged.

For (L3), fix an address of level \(i\), its region \(Z\),
a rule, and a depth \(q\le i\). The meet determining set
is exactly the localization in \cref{I-95-81}: its members
are \(\Omega^\diamond_r\sqcap Z_{[r]}\), \(r\le q\).
Its inner bonds lie in \(\mathfrak e\). The three pairs
\((U,A')\), \((U,A'')\), and \((U,0)\) belong to
\(\mathcal Q^{\rm T}(a_{\rm T})\), by (L1), the bounds
just proved, and \(n_i\le n_t,n_{t^*}\). The tower
constructed there is the tower in (L3): its averaged datum
is the same in the real and complex senses wherever the
field is real, by \cref{I-27-78}; changing the normalization
keeps \(\eta A'\) and \(\eta^3J\) fixed.

At a place \(f\) of (L3), with \(\bar n=n_i(x_f)\),
the inequality for depth \(q\) and interior index
\(q\le p\le i\) is exactly
\begin{equation}\label{II-88-42}
 |G_{q,Z}(U,A';f)|<L^{i-p}
                    \mathsf c_{\sigma^{\rm m}}(\bar n)
                    \mathsf r_0(\bar n).
\end{equation}
For currents this uses the power
\(2\bar n+\min\{q,\bar n\}-3q\) in
\cref{I-95-81}. The place has a corner in \(Z^{\wr p}\),
which is contained in \(Z^{\wr q}\), so that the localized
estimate applies even at its incident boundary stencils.

If \(n_i(b)>\bar n\), crossing the intervening layers
costs at least \(\rho(n_i(b)-\bar n-1)\) in weighted
distance; moving both sites to their tops costs at most
three. If \(n_i(b)\le\bar n\), monotonicity of the
profile suffices. Thus \eqref{eq:history-profile-decay}
gives
\[
 e^{-\delta_{\rm T}d(y_f,y_b)}\mathsf w(n_i(b))
       \le e^{3\delta_{\rm T}}C^*\mathsf w(\bar n).
\]
Equations~\eqref{I-95-84} and
\eqref{eq:history-substitution-increments} bound the tower
increment by
\[
 \tfrac32e^{6\delta_{\rm T}}K_{\rm T}C^*r
                          \mathsf w(\bar n)
 \le\tfrac38e^{3/2}\mathsf w(\bar n)
 <2\mathsf w(\bar n).
\]
The coefficient gap \eqref{II-88-21}, with \(L^{i-p}\ge1\),
absorbs it in \eqref{II-88-42}. The real tower is unchanged.
This proves all of (L3), including existence.

Finally fix a core address \((\sigma'',i,X')\). Set
\[
 b_\iota=\mathsf c_{\sigma^*}(i)\mathsf r_\iota(i),
 \qquad \tau=1-\frac{\mathsf c_{\sigma^{\rm m}}(i)}
                           {\mathsf c_{\sigma^*}(i)}.
\]
Then \(\tau b_0\ge2\mathsf w(i)\),
\(b_1^2\le(5\mathsf c_\natural/4)b_0\), and
\(2(b_0+b_1)\le4\hat a\le a^\sharp\).
An incident bond of \(X'\subset\Lambda_i\) can start
in layer \(i-1\). This boundary case gives the bounds
\[
 |\mathbf A_{(i)}|\le\tfrac L2r\mathsf w(i),\quad
 |\nabla_{\mathbf U}\mathbf A_{(i)}|
                         \le\tfrac{L^2}2r\mathsf w(i),\quad
 |\mathbf K_{(i)}|\le\tfrac{L^3}2r\mathsf w(i)
\]
after restriction to level \(i\); here
\(\mathsf w(i-1)\le\mathsf w(i)/2\).
The allowable potential radius in \cref{II-16-69} is at
least
\[
 \frac{8\mathsf w(i)}{5C_{\rm os}K^\sharp(1+\mathsf c_\natural)}.
\]
The last two inequalities in \eqref{II-88-33} put the
potential and derivative strictly inside this radius and
put the current increment within \(\tau b_0\).
Apply \cref{II-16-69} to the core of the restricted input.
Restriction commutes with substitution because
\(\eta_iL^{-(j-i)}=\eta_j\), and the current scales by
the cube of this factor. The resulting core is precisely
the one required in (L4) for \(\sigma^*\).
\end{proof}

For later use take
\begin{equation}\label{II-88-51}
 \mathsf c_{\rm rm}=(C_{\rm os}L^3)^{-1},\qquad
 \mathsf r_\natural=\mathsf c_{\rm rm}
       \min\{1,K_{\rm T}^{-1},
                         [K^\sharp(1+\mathsf c_\natural)]^{-1}\}.
\end{equation}

\begin{corollary}\label{II-88-50}
If \(\hat a\le\mathsf c_{\rm rm}\min\{a_{\rm T},a^\sharp\}\),
then \(\|\mathfrak h\|^{\mathcal S}_{\mathsf v_k,\mathbf U}
<\mathsf r_\natural\) implies the conclusions of
\cref{II-88-43} with \(\sigma^*=\sigma'\).
In particular, if \(P^+\in\widetilde U_{(t,{\rm mid})}(Y)\),
\(X\subset Y\), and \(\sigma'=(j,\phi')\preccurlyeq(k,{\rm out})\),
then
\[
 \rho_X^{t\to j}P^+
       \dotplus\rho_X^{k\to j}\mathfrak h^+
                       \in\widetilde U_{\sigma'}(X)
\]
whenever the restricted shift has this norm bound and
the substitution is defined on \(\operatorname{Bd}(X)\).

For a core on \(X\) whose incident bonds have both endpoints
in \(\Omega_j\), put
\(\mathsf v_j^\sharp(b)=b_0L^{n_k(b)-j}\).
If \(P\in\mathcal U_j^\circ(X;(1-\tau)b_0,(1-\tau)b_1)\)
and the radii satisfy the hypotheses of \cref{II-16-69},
then
\begin{equation}\label{II-88-52}
 \|\mathfrak h\|^{\mathcal S}_{\mathsf v_j^\sharp,\mathbf U}
             <\frac{\tau}{C_{\rm os}K^\sharp(1+\mathsf c_\natural)}
\end{equation}
implies that substitution is defined and belongs to
\(\mathcal U_j^\circ(X;b_0,b_1)\).
\end{corollary}

\begin{proof}
Use \(C^*=L\) in \cref{II-88-43}.
The definition \eqref{II-88-51} gives
\(r<L^{-3}\), \(4LK_{\rm T}r<4L^{-2}<1\), and
\(L^2C_{\rm os}K^\sharp(1+\mathsf c_\natural)r<L^{-1}\),
which verify \eqref{II-88-33}.
Restriction of a graded triple preserves (L1)--(L4):
\(\eta A'\), \(\eta^2\nabla A'\), and \(\eta^3\mathbf J\)
are unchanged; every new tower and core address was already
an address of the larger region; and a containing gauge
cube in (L2) restricts to the requested smaller cube.
Thus the restricted input lies in
\(\widetilde U_{\sigma'|(t,{\rm mid})}(X)\), as required.

For the core assertion write \(\nu_b=n_k(b)-j\ge0\).
The norm bound gives
\(|\mathbf A|<b_2\),
\(|\nabla_{\mathbf U}\mathbf A|<L^{-\nu_b}b_2\le b_2\),
and \(|\mathbf K|<L^{-2\nu_b}b_2\le b_2\), where
\(b_2=\tau b_0/[C_{\rm os}K^\sharp(1+\mathsf c_\natural)]\).
In particular \(b_2\le\tau b_0/2\), so the current stays
below \(b_0\). Apply \cref{II-16-69}.
\end{proof}


\subsection{Re-expanding an earlier term}
\label{sec:history-term-reexpansion}

We write \(C_{\mathcal A}=\mathsf C_{\mathcal A}\)
for the fixed field-cutoff constant.
We next apply the substitution margin to a term at an
earlier level. There are two domains. A regular term on
\(X\subset\Lambda_j\) is a function of the pair on the
fixed core \(\mathcal U^\circ_j(X;b_0,b_1)\), where
\((b_0,b_1)=(\alpha_0,\alpha_1)\) or
\[
 b_\iota=b_\iota^{\rm oc}
       =\frac{\theta_{\rm oc}}{1+\beta_{\rm s}}\alpha_\iota .
\]
A boundary term on a layered output space at \((j,X)\), or a formed exponent
with \(j=l+1\), is a function on
\(\widetilde U_{(j,{\rm out})}(X)\), together with its
stored coordinates of levels \(<j\), respectively level
\(l\). We use the plain substitution for a regular term
and the pinned substitution for a boundary term on a layered
space or a formed exponent. All stored coordinates remain fixed.

Write \(X^\uparrow\) for the union of level-\(k\) fine
cubes meeting \(X\). For \(Y\supset X^\uparrow\), put
\[
 \Phi_Y[f](s,\lambda;P,\mathcal A,\mathfrak q)
 =f\bigl(\rho_X^{t\to j}P
       \dotplus_{\mathcal S}
       \lambda\rho_X^{k\to j}\mathfrak h_Y(s;P,\mathcal A);
       \mathfrak q_X\bigr).
\]
In the pinned case replace \(\lambda\mathfrak h_Y\) by
\((\lambda\mathfrak h_Y)^\delta[P]\) and denote the
function by \(\Phi_Y^\delta[f]\).
Here \(\dotplus_{\mathcal S}\) performs the logarithmic
substitution on \(\operatorname{Bd}_{\mathcal S}(X)\)
and leaves the other bond components unchanged.
This convention is needed because the graded spaces
place no bound on those other components.
The restrictions of stored coordinates are
\(\mathfrak q_X=(\mathfrak q|_X)|^{<j}\), or
\((\mathfrak q|_X)|^{(l)}\), in the two cases with stored coordinates.

The localized increment is
\begin{equation}\label{II-94-25}
 \mathbb V^\star[f](Y;\mathfrak u)
   =[\,Y\setminus X^\uparrow\mid X^\uparrow\,]^*
        \Phi_Y^\star[f](\,\cdot\,,1;\mathfrak u),
\end{equation}
where \(\star\) indicates the plain or pinned choice.
It is zero when \(X^\uparrow\not\subset Y\), on the
domain of the regional shift. The starred bracket is
the one of \cref{II-18-108}; in particular it subtracts
the value at zero parameters.

Let \(\mu=\mu_{\mathfrak h}\) and
\(\delta_*=\mu M/4\). In this subsection the profile
radii \(r_\iota(n)=\mathsf r_\iota((k,\mathrm{out});n)\)
are those of the previous output structure. Bounds on
the input triple use the radii of \((t,\mathrm{mid})\).
Retain the comparison and margin hypotheses of
\cref{II-88-50} between these two structures and every
earlier output structure. In addition to the substitution
conditions above, assume the radius comparison
\begin{equation}\label{eq:history-radius-decay}
 2^\nu e^{-\delta_*L(\nu-1)}r_0(n+\nu)
       \le C_\uparrow r_0(n)
 \quad(0\le n<n+\nu\le k).
\end{equation}
For regular terms assume
\((1+\beta_{\rm s})\mathsf r_\iota((t,\mathrm{mid});n)
\le(1-\tau)b_\iota\),
\(0<b_0\le b_1\), \(b_1^2\le\mathsf c_\natural b_0\),
and \(2(b_0+b_1)\le a^\sharp\).
The source-cube ceilings of \cref{II-94-90} are imposed
with
\(\mathsf c_{\rm KL}=\max\{\mathsf K_{\rm I},L^3\}^{-1}\):
\[
 (1+\beta_{\rm s})\mathsf r_\iota((t,\mathrm{mid});n)
       \le a_{\rm m}\quad(0\le n\le t),
 \qquad
 (1+\beta_{\rm s})\mathsf r_\iota((t,\mathrm{mid});n)\le
             \mathsf c_{\rm KL}a_{\mathfrak h}\quad(n=k,t).
\]
Also assume
\[
 (1+\beta_{\rm s})\mathsf r_\iota((t,\mathrm{mid});n)
       \le c_{\rm rt},\qquad
 (1+\beta_{\rm s})\mathsf r_1(\sigma;n)\le\log2,\qquad
 a_\flat\le\min\{a_{\rm m},c_{\rm rt}\},\qquad
 \mathsf C_{\mathcal A}\delta_k\le\varepsilon_{\mathfrak h}.
\]
The first inequalities give the regularity required by
the levelling construction; the logarithm bound is imposed
at both the input and earlier output structures \(\sigma\).
For example, at a top of
level \(q\), the scaled bounds are
\(\xi_qL^{-1}|A'|\),
\(\xi_q^2L^{-1}|\nabla^\xi A'|\), and
\(\xi_q^3L^{-3}|\mathbf J|\le c_{\rm rt}\), by (L1).
For each turning cell of order \(q\), the containing
cube in (L2) supplies a gauge with the same bounds;
rescale its potential by \(\xi_q^{-1}\).
Thus \cref{II-85-34} applies to these graded triples.

The distance from a term to the new active region is
\[
 \mathsf d_k(X)=M^{-1}d_{\mathbf B}(X,\Omega_t).
\]
Put \(\mathsf a=\delta_k/\alpha_0\) for a regular term
and \(\mathsf a=\delta_k/r_0(k)\) for a term on a layered output space.
With \(\mathsf c_0=(1-\theta_{\rm s})\beta_{\rm s}/2\),
define
\begin{align}
 \mathsf m_*&=\max\left\{
  \frac{C_{\rm m}C_\uparrow}
       {\mathsf r_\natural\mathsf c_0},
  \frac L\epsilon,
  \frac{C_{\rm m}(1+\beta_{\rm s})
        C_{\rm os}K^\sharp(1+\mathsf c_\natural)}
       {\tau\theta_{\rm oc}}\right\},
       \label{eq:history-amplitude-constant}\\
 \mathsf C_{\rm sb}
 &=2\mathsf C C_\Sigma
     e^{(2C_{\rm m}+4)\mu}\mathsf m_*,
 \qquad
 \mathsf H_{\rm sb}
    =\mathsf C_{\rm sb}\mathsf C_{\mathcal A}
                         e^{\mathsf C_{\rm sb}\bar\kappa}.
       \label{II-103-11}
\end{align}
Here \(C_\Sigma=\mathsf S(\mu/2)\) is the finite
weighted row-sum constant, and
\(\mathsf c_{\mathfrak h}=\mathsf C e^{\mathsf C\bar\kappa}\).
For \(0<\delta_{\rm sb}\le\delta_*\), set
\begin{equation}\label{II-103-15}
 \varrho_X=
  [\,\mathsf H_{\rm sb}\mathsf a
                    e^{-\delta_{\rm sb}\mathsf d_k(X)}\,]^{-1}.
\end{equation}
We interpret \(e^{-\infty}=0\), so that
\(\varrho_X=\infty\) if no active source is reachable.
The constant \(\mathsf C_{\rm sb}\) is independent of
the height, volume, history, couplings, \(M\), and
\(\bar\kappa\). Its dependence is on the displayed
fixed analytic and levelling constants.

\begin{proposition}\label{II-104-17}
Under these conditions, the substituted triple belongs
to the domain of \(f\), for every
\(P\in\widetilde U_{(t,{\rm mid})}(Y)\),
\(\sup|\mathcal A|<\mathsf C_{\mathcal A}\delta_k\),
\(s\in\mathbb D(Y)\), and \(|\lambda|<\varrho_X\).
In the pinned case the substitution is locally
holomorphic on each fixed-real-factor slice at its
restriction to \(\operatorname{Bd}(X)\).

Suppose \(\varrho_X>1\), and that \(f\) is analytic,
bounded by \(b_f\), gauge invariant with its stored
coordinates, and has the tie for real triples.
Then \(\mathbb V^\star[f]\) is defined on every
admitted stored-coordinate tier, is analytic there,
and
\begin{equation}\label{II-104-18}
 |\mathbb V^\star[f](Y;\mathfrak u)|
 \le2b_f\varrho_X^{-1}
                e^{-(\bar\kappa-1)\#(Y\setminus X^\uparrow)}.
\end{equation}
It vanishes at \(\mathcal A=0\) and unless \(Y\)
contains a source cube. It is gauge invariant and
has the tie; it is Borel, conjugation real, or a
function of the pair if \(f\) has the corresponding
property. The construction is linear in \(f\).
For a formed exponent, the same assertions hold with
\(b_f=\operatorname{osc}f\), assuming invariance up to
its stored coordinates.

If all restrictions of a point \(\mathfrak u\) on
the torus belong to the indicated tiers, then
\begin{equation}\label{II-104-20}
 \sum_{Y\supset X^\uparrow}
       \mathbb V^\star[f](Y;\mathfrak u|_Y)
 =\Phi^\star_{T_\ell}[f](\mathbf1,1;\mathfrak u)
       -f(\rho_X^{t\to j}P;\mathfrak q_X).
\end{equation}
\end{proposition}

\begin{proof}
First consider triples with
\(\eta_t|A'(b)|<\log2\) at every bond of
\(\operatorname{Bd}(X)\setminus
 \operatorname{Bd}_{\mathcal S}(X)\).
At such a bond the shift is zero and
\(\log(e^{i\eta_tA'(b)})=i\eta_tA'(b)\).
Thus the ordinary substitution is defined there and
agrees with \(\dotplus_{\mathcal S}\).
On the bonds of the label, \cref{II-88-50} gives
substitution room
\[
 r_{\rm room}=\mathsf r_\natural
 \quad\hbox{or}\quad
 r_{\rm room}=
       \frac{\tau}{C_{\rm os}K^\sharp(1+\mathsf c_\natural)}
\]
for the graded and fixed-core domains, respectively.
Restriction from the mid-level space to an earlier
output space preserves (L1)--(L4), as in the proof
of \cref{II-88-50}.

We verify the decay needed to use this room.
Let \(\mathsf T_X\) be the positive-level tops whose pillars
contain an endpoint of a bond in \(\operatorname{Bd}(X)\).
For every top \(y\), put
\[
 \mathsf D_y=\min_{c\in\mathfrak X^+}d_{\mathbf B}(y,c_-),
 \qquad \mathsf D(b)=\mathsf D_{y(b)}.
\]
Call a level-\(k\) determining bond live if it has an endpoint
in \(\Omega_t\). For each such bond \(c\), let
\(y_c^{\rm lv}\) be its endpoint in \(\Omega_t\); if both
endpoints lie there, take the initial endpoint in the positive
orientation. Define
\[
 \mathsf D^{\rm lv}(b)=\min_{c\ {\rm live}}
                  d_{\mathbf B}(y(b),y_c^{\rm lv}).
\]
An empty minimum is \(+\infty\).
For every live bond, \(d_{\mathbf B}(c_-,y_c^{\rm lv})\le1\).
For a bond incident to \(X\), with \(n=n_k(b_-)\),
the triangle inequality and the pillar length bound
give
\[
 M\mathsf d_k(X)\le\tfrac52+
                 d_{\mathbf B}(y(b),y_c^{\rm lv}).
\]
Crossing the \(k-n\) separated layers gives the second
bound
\[
 8LM(k-n)\le\tfrac32+
                 d_{\mathbf B}(y(b),y_c^{\rm lv}).
\]
Indeed every layer crossing traverses a strip of width
\(8LM\) in its own units; adding these disjoint pieces
is exactly the weighted length defining
\(d_{\mathbf B}\). Taking the mean and then the
minimum over live bonds yields
\begin{equation}\label{II-88-56}
 e^{-\mu\mathsf D^{\rm lv}(b)/2}
 \le e^\mu e^{-\delta_*\mathsf d_k(X)}
                      e^{-8\delta_*L(k-n)}.
\end{equation}
Here \(\mathsf D(b)\ge\mathsf D^{\rm lv}(b)-1\).
Divide by
\(\mathsf v_k(b)=\mathsf c_0\,2^{-(k+1-n)}r_0(n)\)
and apply \eqref{eq:history-radius-decay}.
For a core use
\(\mathsf v_j^\sharp(b)=b_0L^{n-j}\ge b_0\).
For a top whose pillar meets an endpoint of a bond
incident to \(X\), choose that bond: its top is at
weighted distance less than
\(3/2+13/12+3/2<4.1\).
Consequently
\begin{align}
 \max_{b\in\operatorname{Bd}_{\mathcal S}(X)^+}
   \frac{e^{-\mu\mathsf D(b)/2}}{\mathsf v_k(b)}
 &\le\frac{C_\uparrow e^{4\mu}}
              {\mathsf c_0r_0(k)}
                e^{-\delta_*\mathsf d_k(X)},\label{II-103-30}\\
 \max_{b\in\operatorname{Bd}_{\mathcal S}(X)^+}
   \frac{e^{-\mu\mathsf D(b)/2}}{\mathsf v_j^\sharp(b)}
 &\le b_0^{-1}e^{4\mu}
                e^{-\delta_*\mathsf d_k(X)},\label{II-103-31}\\
 \max_{y\in\mathsf T_X}e^{-\mu\mathsf D_y/2}
 &\le e^{4\mu}e^{-\delta_*\mathsf d_k(X)}.\label{II-103-32}
\end{align}
All these maxima are zero if their sets are empty.

The regional response estimate, summed over source
bonds, gives
\[
 |\mathfrak h_Y|_{\mathbf U,Y}(b)
    \le c_*e^{-\mu\mathsf D(b)/2},\qquad
 c_*=2C_\Sigma\mathsf c_{\mathfrak h}
                         \mathsf C_{\mathcal A}\delta_k.
\]
For pinning, use the slowly varying profile
\(c_*e^{3\mu/4}e^{-\mu d_{\mathbf B}(x,\mathfrak X)/2}\)
in \cref{II-85-34}. Moving from a site to its top
costs at most \(3/2\) in the weighted metric.
It follows that
\begin{align}
 |(\lambda\mathfrak h_Y)^\delta|_{\mathbf U,Y}(b)
  &\le|\lambda|C_{\rm m}e^{2C_{\rm m}\mu}
          c_*e^{-\mu\mathsf D(b)/2},\label{II-96-32}\\
 \widehat\varepsilon_y(\lambda\mathbf A_Y)
  &\le|\lambda|L e^{2C_{\rm m}\mu}
          c_*e^{-\mu\mathsf D_y/2}.\label{eq:history-pin-top-bound}
\end{align}
The top estimate follows also by applying the profile
bound at the positive bare bond beginning at the top
inside its pillar.

The three entries of \(\mathsf m_*\) now give, for
\(|\lambda|<\varrho_X\),
\[
 \|(\lambda\mathfrak h_Y)^\delta|_X\|_{\mathsf v,\mathbf U}
       <r_{\rm room},\qquad
 \max_{y\in\mathsf T_X}
       \widehat\varepsilon_y(\lambda\mathbf A_Y)
        <\epsilon.
\]
For graded terms use
\(\delta_k/\mathsf a=r_0(k)\le1\);
for regular terms use
\(\delta_k/\mathsf a=\alpha_0\le1\) and
\(\alpha_0/b_0\le(1+\beta_{\rm s})/\theta_{\rm oc}\).
These are exactly the two factors allowed by
\eqref{eq:history-amplitude-constant}.
The first inequality proves membership, and the
second forces every endpoint pillar needed at
\(\operatorname{Bd}(X)\) to be levelled.
The local holomorphy clause of \cref{II-85-34}
therefore applies on those bonds. No holomorphy of
the threshold map on the entire region is asserted.
One may take local extensions of the response into
the linear space of shifts vanishing off
\(\mathfrak e_Y\); projection onto that space preserves
the actual response and keeps these extensions
admissible for the pinning map.

For fixed stored coordinates, the resulting function
\(\Phi^\star_Y[f]\) is bounded by \(b_f\) on the
\(\lambda\)-disc of radius \(\varrho_X\).
Subtract its value at \(\lambda=0\).
Schwarz's lemma and the starred Cauchy estimate
\eqref{II-18-113} give
\[
 |\mathbb V^\star[f](Y)|
 \le\frac{2b_f}{\varrho_X}
       (e^{\bar\kappa}-1)^{-\#(Y\setminus X^\uparrow)}.
\]
This implies \eqref{II-104-18} for
\(\bar\kappa>\log2\).
If only the oscillation is bounded, subtract
\(f(P_0;\mathfrak q_X)\) for any one admitted \(P_0\).
The starred bracket removes this constant. Its
choice need not be measurable, since it is used
only to prove the bound.

We justify the tie for pinned substitution before passing to
the full space. The estimates in \eqref{II-96-32} hold
throughout the larger center polydisc
\(|\zeta(a)|<2c_{\rm m}\varepsilon_{y(a)}\) of
\cref{II-85-34}; thus the same margin admits every
substitution with these independent center values.
By \cref{II-84-18}, initial sites of stored bonds
are tops or stored sites of pinnable tops. Turns
agree there for inputs with one pair.

For real input triples with one pair and a real
shift, let \(u,\widetilde u\) be their turns and
\(Q,\widetilde Q\) their substituted triples.
Then \(w=\widetilde uu^{-1}\) is unitary, equals one
at every stored initial site, and
\(\mathbf p(\widetilde Q)=\mathbf p(Q^w)\).
Gauge invariance admits \(Q^w\). The tie identifies
its value with that at \(\widetilde Q\); invariance
and the identity of \(w\) on the stored sites identify
it with the value at \(Q\).
For fixed real input triples these readings are
holomorphic on the connected product of the
\(s,\lambda,\mathcal A\) discs. Equality at real
points extends throughout this product.

For a function of the pair, allow complex input
triples and use their common independent center
variables \(\zeta\). The same margin admits both
substituted triples for every \(\zeta\); their
readings are holomorphic in \(\zeta\).
At real \(\zeta\), the unitary \(w\) argument uses
pair dependence in place of the tie. Analytic
continuation gives equality at the common actual pin:
\begin{equation}\label{II-85-37}
 \Phi^\delta[f](\widetilde P,\mathfrak h;\mathfrak q)
 =\Phi^\delta[f](P,\mathfrak h;\mathfrak q)
 \quad\text{if }\mathbf p(\widetilde P)=\mathbf p(P).
\end{equation}
For a plain local term, with no stored data, the
same continuation compares the actual pin with
the zero center. Thus plain and pinned readings agree.

For completeness, we remove the restriction on the
outer bond components. Let \(\nu_R\) set \(A'\) and
\(\mathbf J\) equal to zero on
\(\mathsf{Out}(R)=\operatorname{Bd}(R)^+
 \setminus\operatorname{Bd}_{\mathcal S}(R)\), preserving
the real factor. Membership in a graded space is
unchanged. The regional response and its domain are
unchanged on the label: the natural-class data,
all determining averages, source cubes and pivot
rows use inner bonds only. The same is true of
pinning, since every positive-level pillar and its
reading set consist of inner bonds. Thus the
substituted triples \(T_X\) satisfy
\begin{equation}\label{II-104-15}
 T_X(s,\lambda;\nu_Y P,\mathcal A)
       =\nu_XT_{X|h'}(s,\lambda;P,\mathcal A),
\end{equation}
with simultaneous definedness.

A bounded analytic \(f\) on a layered output space does not
depend on its outer \(A'\) and \(\mathbf J\) coordinates.
Each such coordinate ranges over a whole complex
line in that space. The restriction of \(f\), or of
\(f-f(P_0;\mathfrak q_X)\) in the oscillation case,
is bounded and entire on that line, and hence
constant. Applying this successively to the finitely
many outer coordinates proves
\(f(P;\mathfrak q_X)=f(\nu_XP;\mathfrak q_X)\).
Consequently
\begin{equation}\label{II-104-21}
 \mathbb V^\star[f](Y;\mathfrak u)
       =\mathbb V^\star[f](Y;\nu_Y\mathfrak u).
\end{equation}
The right side is evaluated at a triple satisfying
the earlier \(\log2\) restriction. This proves
membership and the bound on the full tiers.
Composition with this coordinate projection gives
local holomorphic extensions and Borel measurability;
the projection commutes with restriction, gauge
transformations and conjugation.

There is one additional point for the tie:
\(\nu_Y\) need not preserve equality of pairs.
For two real input triples with the same pair,
choose at each outer bond \(Z_b\in\mathfrak g\),
\(|Z_b|\le\pi\), with
\(e^{iZ_b}=U_1(b)U_2(b)^{-1}\).
Join their normalized real factors in five steps.
The intermediate triples with outer components
\[
 (e^{inZ_b/5}U_2(b),0,0),\qquad
 (e^{i(n-1)Z_b/5}U_2(b),\eta_t^{-1}Z_b/5,0)
\]
have the same pair for \(1\le n\le5\).
They belong to the same graded space, and
\(\pi/5<\log2\) keeps the second triples in the
already treated set. Equality of the two potential
values follows from the tie there; normalization
then joins successive steps. This proves the tie
on the full space. The same argument with arbitrary
input triples proves the pair property when \(f\)
is a function of the pair.

Finally, decoupling the parameters splits the
response into its wall components. Pinning preserves
that splitting and the support of each component,
by the margin and restriction clauses of
\cref{II-85-45}. Components meeting neither the
base nor a source can therefore be removed before
taking the bracket. This proves the stated
vanishing and locality. Summing the finite
Möbius differences over connected domains containing
the base leaves exactly the full-parameter value
minus the zero-parameter value. Domain consistency
identifies these values on the torus, and the shift
at zero parameters is zero. This proves
\eqref{II-104-20}. Gauge covariance, conjugation,
the tie on the restricted set and Borel measurability
follow by composing the corresponding response and
pinning properties with those of \(f\); finite
differences preserve them. Linearity is immediate
from the bracket.
\end{proof}

\subsection{Counting earlier domains}

A domain of level \(j<k\) is admitted if it meets
\(\Omega_{j+1}^c\); at level \(k\) it is admitted if it
meets \(Z_k^\sim\). Let \(n_X\) be the largest \(n\le k\)
for which \(X\) and \(\Omega_n\) share a bare cube, and
set \(n=\max(j,n_X)\), \(p=k-n\), \(q=n-j\).
A domain is placed with integer \(\nu\ge0\) if
\(d_j(X)\ge2L^{n-j-\nu}-2\) whenever \(j<n\le k\) and
\(X\) shares a bare cube with \(\mathcal L_n\).
For \(a'\ge0\), put \(D_0=D_1=0\),
\(D_q=8L^q-2\) for \(q\ge2\), and
\begin{equation}\label{II-91-8}
 C_\square(L,a')=C_{\mathrm E}
 \left(\sum_{p\ge0}(1+p)^{a'}L^{-5p}\right)
 \left(\sum_{q\ge0}(1+q)^{a'}L^{3q}e^{-D_q}\right).
\end{equation}
For \(0<\delta<1/2\), put
\(\kappa'=(1-2\delta)\kappa\), \(r_V=(1-2\delta)(\kappa-2)\).
Assume
\[
 a=2\delta\kappa\ge2(c_{\mathrm E}+1),\qquad
 \bar\kappa-1\ge\kappa,\qquad \delta_{\rm pc}L\ge10\log L.
\]
Use the smaller distance rate needed in each application
for \(\delta_{\rm pc}\). Set
\begin{equation}\label{II-91-22}
 \begin{split}
 E_q&=\max(2L^{q-\nu}-2,0),\quad
 u_p=2^pe^{-\delta_{\rm pc}Lp/2},\quad v_q=e^{-aE_q/2},\\
 \varpi_i&=L^{3i}\sum_{p+q=i}u_pv_q,\qquad
 C_{\rm h}=8C_{\mathrm E}\max\{1,[2(1-2\delta)]^{-1}\},
 \quad C'_{\rm h}=3(\nu+1)^2L^{4\nu}.
 \end{split}
\end{equation}

\begin{lemma}\label{II-91-52}
Write \(m(X,Y)=\#_{\pi_k}(Y\setminus X^\uparrow)\);
all sums below have \(X^\uparrow\subset Y\).
For admitted domains,
\begin{equation}\label{II-91-56}
 \sum_{j,X}(1+k-j)^{a'}
 e^{-\kappa d_j(X)-\delta_{\rm pc}\mathsf d_k(X)
                    -(\bar\kappa-1)m(X,Y)}
 \le8C_\square(L,a')e^{-(\kappa'-1)d_k(Y)} .
\end{equation}
This conclusion needs only \(a\ge c_{\mathrm E}+1\).
For placed domains and nonnegative numbers \(B_j\),
\begin{equation}\label{II-91-58}
\begin{gathered}
 \sum_{j,X}B_j(d_j(X)+1)
 e^{-\kappa d_j(X)-\delta_{\rm pc}\mathsf d_k(X)
                    -(\bar\kappa-1)m(X,Y)}
 \le C_{\rm h}\sum_{j=0}^k\varpi_{k-j}B_j\,e^{-r_Vd_k(Y)},\\
 \sum_{i\ge0}(i+1)L^i\varpi_i\le C'_{\rm h}.
\end{gathered}
\end{equation}
Under \((\mathrm{CR})_t\), an admitted domain has a point
at sup distance at least \(7LM\check R_t\) from
\(\Omega_t\). Thus a domain \(Y\) containing its base
and a source cube satisfies
\begin{equation}\label{eq:historical-source-distance}
 d_k(Y)\ge7L\check R_t-2.
\end{equation}
\end{lemma}

\begin{proof}
Truncate the distance to \(\Omega_r^c\) at its collar
width. Along a connected path from \(X\) to \(\Omega_t\),
this 1-Lipschitz function traverses its whole range;
the portion strictly between its endpoint values lies
in \(\mathcal L_r\). Its length is at least the collar
width. These layer portions have disjoint interiors,
so their weighted lengths add, giving
\(\mathsf d_k(X)\ge8L(k-n_X)\ge8Lp\).
If \(X\) is admitted and \(q\ge2\), its points cross
the collar between \(\Omega_{n-1}^c\) and \(\Omega_n\).
The diameter bound gives \(d_j(X)\ge8L^q-2=D_q\).
For a placed domain it gives \(d_j(X)\ge E_q\):
maximality of \(n_X\) ensures that \(X\) shares a
bare cube with \(\mathcal L_{n_X}\).

Through a fixed level-\(j\) cube, \cref{II-18-26}
bounds the admitted sum by
\(C_{\mathrm E}\sum_{p+q=k-j}L^{-8p}e^{-D_q}\).
A fine cube contains \(L^{3(k-j)}\) such cubes.
Use \((1+p+q)^{a'}\le(1+p)^{a'}(1+q)^{a'}\) and
sum to obtain \(C_\square(L,a')\) per fine cube.
The attachment inequality
\(d_k(Y)\le d_j(X)+m(X,Y)\) gives
\[
 \kappa d_j(X)+(\bar\kappa-1)m(X,Y)
 \ge\kappa'd_k(Y)+a d_j(X).
\]
Pin \(X\) by a small cube in \(Y\), use
\(\#_{\pi_k}Y\le8(d_k(Y)+1)\), and then \(d+1\le e^d\).
This proves \eqref{II-91-56}.

For placed domains absorb \(d_j(X)+1\) into \(e^{d_j(X)}\),
split \(a-1=(a/2-1)+a/2\), and use \(a/2-1\ge c_{\mathrm E}\).
The rooted sum is bounded by
\(C_{\mathrm E}\sum_{p+q=k-j}u_pv_q\).
The same pinning argument with
\((d+1)e^{-2(1-2\delta)d}
\le\max\{1,[2(1-2\delta)]^{-1}\}\) proves the first
inequality in \eqref{II-91-58}.
For the second, \(L^{4p}u_p\le(2/L)^p\).
The terms \(q\le\nu\) in
\(\sum(q+1)L^{4q}v_q\) total at most
\((\nu+1)^2L^{4\nu}\); the rest is at most
\((\nu+1)L^{4\nu}\sum_{s\ge1}(s+1)e^{-26s}\).
Together with \(\sum(p+1)(2/13)^p<1.4\), this gives \(C'_{\rm h}\).

For \(j<k\) an admitted domain meets the old excluded
region; for \(j=k\) it meets \(Z_k^\sim\).
The next collar, using \((\mathrm{CR})_t\) in the latter
case, leaves distance at least \(7LM\check R_t\).
Join that point to a point of a source cube in \(Y\)
and use \(|x-y|_\infty\le M(d_k(Y)+2)\).
\end{proof}

\subsection{The four historical potential families}
\label{sec:historical-potential-families}

Use the actual regional response of \cref{II-93-26},
its pinned version of \cref{II-85-34}, and the full
graded domains of \cref{II-104-17}. All potentials
are restricted to \(B\in\mathbb D^{\rm f}_{x_k}(Y)\),
\(\mathcal A=g_kB\), and the prescribed stored tiers,
including when a defining sum is empty.
Throughout these families the counterterm is the actual
running increment \(b_j=L^{-1}x_{j-1}-x_j\),
with \(|b_j|\le B_\sharp\).

Write
\[
 \widetilde E^{(j)}(X,P,z)=E^{(j)}(X,P,z)
 -b_j\mathbf1_{\{X=X_z\}}A^{\eta_j}(h_z^{(j)},\mathbf U).
\]
For \(z\in(\Lambda_j^{\sim-1})^{(j)}\), let
\(n=\mathsf{ol}_j(z)\) be the largest index among \(j\)
and the \(j<n\le k\) with \(X'_n(z)\subset\Omega_{n+1}\).
Call \(X\ni z\) narrow if \(n>j\) and
\(X\subset\square_n(z)^{\sim2}\), and wide otherwise.
The regular list consists of the wide members
\(\widetilde E^{(j)}(X,\cdot,z)\), the octet members
\begin{equation}\label{II-94-13}
 f_{j,z}=\sum_{X:\ z\in X\subset\square_n(z)^{\sim2}}
 \left[(\widetilde E^{(j)}_{X,z})^{\square_n(z)}
                 -\widetilde E^{(j)}_{X,z}(\mathbf1,0)\right],
\end{equation}
and the edge tents
\(-b_jA^{\eta_j}(\varphi_jh_y^{(j)},\cdot)\),
\(y\notin\Lambda_j^{\sim-1}\).
The octet sum is formed before taking its starred
bracket; its domain is the octet core at level \(n\)
on \(X'_n(z)\).

The boundary list has the plain members
\(E_{\mathcal S}^{(j),\rm ed}(X,\cdot)
=\sum_{z\notin\Lambda_j^{\sim-1}}E^{(j)}(X,\cdot,z)\)
on fixed cores, and the members
\(\mathscr B_h(j,X,\mathrm{mk})\) on layered spaces
with stored levels \(<j\).
The formed list has the old factor exponents
\(\varphi_f\) at \(j_f=l_f+1\) on \(X_f\), with stored
level \(l_f\) only. Restrict each old partial map
to its stated space before forming a member.

For an inner bare plaquette \(p\) of \(\Omega_0^\diamond\),
that is, one whose four corners lie in \(\Omega_0^\diamond\),
let \(X(p;m)\) be the union of level-\(m\) cubes
containing its corners. For \(0\le m\le k\), put
\begin{equation}\label{II-94-17}
 W_w^{(m)}(X,\mathbf U)=\eta_k^{-1}
 \sum_{\substack{p\ {\rm inner\ in}\ \Omega_0^\diamond\\
                  x_{\rm b}(p)\in\mathcal L_m,\ X(p;m)=X}}
 w(p)a_p(\mathbf U),\qquad
 p_w=\sum_{j=1}^kL^{j-k}b_j(1-\varphi_j).
\end{equation}
Use plain brackets for Wilson, regular and edge
members, and pinned brackets for boundary members on layered
spaces and for formed members. With only bases contained in \(Y\), define
\begin{equation}\label{II-94-30}
 \begin{aligned}
 \mathbb V^{\rm E}&=\sum_{\rm regular}\mathbb V[f],
 &\mathbb V^{\rm c}&=-\sum_{m,X}\mathbb V[W_{p_w}^{(m)}(X)],\\
 \mathbb V^{\rm B}&=\sum_{\rm ed}\mathbb V[f]
                      +\sum_{\rm mk}\mathbb V^\delta[f],
 &\mathbb V^\Phi&=\sum_f\mathbb V^\delta[\varphi_f].
 \end{aligned}
\end{equation}
Let \(\mathbb V^{\rm act}=\sum_{m,X}\mathbb V[W_1^{(m)}(X)]\)
and \(\mathbb V^{\rm A}=-x_k\mathcal R_3\mathbb V^{\rm act}\).
At \(k=0\) use the direct plaquette difference of
the completed pivot field. The Jacobian and gauge
potentials are the formulas of \cref{II-41-37,II-41-36}, with
the regional nonpivot set and pivot completion,
continued by zero on unused bonds. The four origins are
\begin{equation}\label{II-94-31}
 \mathbb V^{\rm int}=\mathbb V^{\rm J}
 +\mathbb V^{\rm G}+\mathbb V^{\rm A},\qquad
 \mathbb V'=\mathbb V^{\rm E}+\mathbb V^{\rm c},
 \qquad\mathbb V^{\rm B},\qquad\mathbb V^\Phi.
\end{equation}


\subsection{The action and coupling potentials}

Set
\[
 \epsilon_A=\min\{r_P^\circ,\epsilon_{\mathfrak h},
                      (3\mathsf c_{\mathfrak h})^{-1}\},
 \quad \epsilon_P=\epsilon_A/2,\quad
 \delta_{\rm W} =M\min\{\mu_{\mathfrak h},1\},
\]
with \(\epsilon_A=r_P^\circ\) and
\(\mathsf c_{\mathfrak h}=1\) at \(k=0\).
Take
\begin{equation}\label{II-106-10}
 \begin{split}
 K_\flat&=\max\{18/r_P^\circ+12C_P+6,\ 4C_P/(r_P^\circ)^2\},\\
 K_A&=\max\{672e(1+2\mathsf c_{\rm pl})^2(L^4+3),4512C_P\},
 \qquad C_{\rm seg}=8C_P .
 \end{split}
\end{equation}
Here \(\mathsf c_{\rm pl}=600\) is the plaquette product
constant. In addition to the preceding domain and rate
conditions, require
\begin{equation}\label{II-106-56}
 \begin{gathered}
 \mathsf oL\ge1,\qquad
 L^3(1+\beta_{\rm s})\mathsf r_\iota((t,\mathrm{mid});n)
       \le a_P^\circ\quad(n=k,t,\ \iota=0,1),\\
 (1+\beta_{\rm s})\mathsf r_0((t,\mathrm{mid});n)\le1
       \quad(0\le n\le k),\qquad
 C_{\mathcal A}\delta_k\le\epsilon_P .
 \end{gathered}
\end{equation}

\begin{proposition}\label{II-106-39}
On the actual graded domain,
\begin{align}
 |\mathbb V^{\rm act}(Y)|
 &\le K_A\mathsf c_{\mathfrak h}M^3
                      \sup|\mathcal A|e^{-r_Vd_k(Y)},\label{II-106-40}\\
 |\mathbb V^{\rm A}(Y)|
 &\le2K_A\mathsf c_{\mathfrak h}\epsilon_A^{-2}
             M^3x_k\sup|\mathcal A|^3e^{-r_Vd_k(Y)},\label{II-106-41}\\
 |\mathbb V^{\rm int}(Y)|
 &\le K_*\delta_ke^{-r_Vd_k(Y)}
       \left(1+M^{-3}\sum_{b\in\mathfrak X^{\circ+}\langle Y\rangle}
                                     |B(b)|^2\right),\label{II-106-44}
\end{align}
where
\(K_*=(2K_\flat+2K_A\mathsf c_{\mathfrak h}\epsilon_A^{-2})
M^6C_{\mathcal A}\).
The action and gauge remainders vanish to order three,
the Jacobian potential to order one. They are analytic,
Borel, gauge invariant, conjugation real and functions
of the pair, and vanish without a source cube.
\end{proposition}

\begin{proof}
First verify the domain at which the estimates are used.
Let \(p\) be an inner bare plaquette of
\(\Omega_0^\diamond\) appearing in \eqref{II-94-17},
with base in \(\mathcal L_m\). Its graded layer index
is \(m\), except that at \(m=k\) it may be \(t\).
All four bonds are controlled inner bonds.
Conditions (L1) and \eqref{II-106-56} therefore give
\(|\partial\mathbf U(p)-1|<L^{-2m}\).
On any wall-connected union of at most four fine cubes
containing a source cube all bond endpoints have indices \(k\) or
\(t\). Multiplication of each graded bound of order
\(q\le3\) by \(L^{qt}\) costs at most \(L^3\);
thus the bounds of \(\mathcal Q_t(\cdot;a_P^\circ)\)
hold. For its cube gauges, apply (L2)(a) at level \(k\)
and set \(A=LA_k\); its first difference is
\(L^2\nabla^{\eta_k}A_k\). The same ceiling supplies the
required level-\(t\) gauge. Finally the actual regional
response is defined for every field of norm
less than \(\epsilon_{\mathfrak h}\), by
\cref{II-93-26,II-94-90}. This proves the auxiliary
domain assertions for all points of the graded domain,
without imposing conditions on their outer coordinates.

An occupied Wilson member contains at most \(12M^3L^{3m}\)
plaquettes; its base has at most four fine cubes and
linear size zero. For the interpolated action set
\[
 \rho=e\,\mathsf c_{\mathfrak h}\sup|\mathcal A|
                         e^{-\delta_{\rm W}\mathsf d_k(X)}<1 .
\]
The layerwise response bound and the four-factor
plaquette estimate give, for \(|\lambda|\rho\le1\),
\[
 |\partial(e^{i\lambda\eta_k\mathbf A}\mathbf U)(p)-1|
       \le(1+2\mathsf c_{\rm pl})L^{-2m}.
\]
The identity
\(|1-\tfrac12\operatorname{tr}Z|\le\tfrac12|Z-1|^2\)
for \(Z\in\mathrm{SL}(2,\mathbb C)\) gives the bound
\(6(1+2\mathsf c_{\rm pl})^2M^3L^{k-m}\)
on the Wilson member. The starred Cauchy estimate on
the disc \(|\lambda|<\rho^{-1}\) therefore yields
\begin{equation}\label{II-106-37}
 |\mathbb V[W_1^{(m)}(X)](Y)|
 \le12e(1+2\mathsf c_{\rm pl})^2\mathsf c_{\mathfrak h}
 M^3\sup|\mathcal A|L^{k-m}
 e^{-\delta_{\rm W}\mathsf d_k(X)-(\bar\kappa-1)m(X,Y)} .
\end{equation}
When \(\rho=0\), take arbitrarily large discs instead.
There are at most \(7\#_{\pi_k}Y\,L^{3(k-m)}\)
occupied members at level \(m\).
Moreover \(m(X,Y)\ge d_k(Y)\) and
\(\mathsf d_k(X)\ge8L\max(k-m-1,0)\).
The geometric sum
\[
 \sum_{i\ge0}L^{4i-8\max(i-1,0)}\le L^4+3
\]
and \(\#Y\le8(d_k(Y)+1)\) prove \eqref{II-106-40}.
At \(k=0\), the pivot completion satisfies
\(\sup|B'|\le2C_P\sup|\mathcal A|\); telescoping the
four plaquette factors bounds their action difference
by \(376C_P\sup|\mathcal A|\). Multiplication by
\(12M^3\) gives the other term in \(K_A\).

For \(\mathcal A\ne0\), the function
\(\phi(z)=\mathbb V^{\rm act}(Y;P,z\mathcal A)\)
is holomorphic on \(|z|<T=\epsilon_A/\sup|\mathcal A|\ge2\)
and bounded by
\(K_A\mathsf c_{\mathfrak h}M^3\epsilon_Ae^{-r_Vd_k(Y)}\).
Cauchy's estimate for its tail from degree three gives
\(2T^{-3}\) times this bound, proving \eqref{II-106-41}.
The regional pivot calculation is the same calculation
on each pivot block with the unused nonpivot fields
zero; source-cube admission proved above supplies its
domain. Hence
\[
 |\mathbb V^{\rm J}|\le K_\flat M^3\sup|\mathcal A|,
 \qquad
 |\mathbb V^{\rm G}|\le K_\flat M^3x_k\sup|\mathcal A|^3 .
\]
These terms have \(d_k(Y)=0\) on their support.
Since \(x_kg_k^2=1\),
\(x_k\sup|\mathcal A|^3\le
 C_{\mathcal A}\delta_k\sum|B|^2\), proving
\eqref{II-106-44}.
Finite brackets preserve analyticity and symmetries.
Taylor coefficients are limits of finite differences
at real scalar multiples of \(\mathcal A\), so they
are Borel as well. The zero-field and no-source
claims follow from the response vanishing there.
\end{proof}

For the coupling potential use the enlarged layer radii
\[
 a_m=(1+\beta_{\rm s})\mathsf r_0((t,\mathrm{mid});m)
 \quad(m<k),\qquad
 a_k=(1+\beta_{\rm s})\max\{\mathsf r_0((t,\mathrm{mid});k),
                     L^{-2}\mathsf r_0((t,\mathrm{mid});t)\},
\]
and put
\begin{equation}\label{II-3-11}
 \vartheta^a(\epsilon)=3\mathsf c_{\mathfrak h}\epsilon
 \max_{0\le m\le k}a_m^{-1}
       e^{-\delta_{\rm W}L\max(k-m-1,0)/2}.
\end{equation}
Assume \(a_m\le1\) and
\(\vartheta^a(C_{\mathcal A}\delta_k)\le1\).
Strict field cutoff then makes the value at
\(\sup|\mathcal A|\) strictly less than one.

\begin{proposition}\label{II-107-44}
Suppose \(|b_j|\le B_\sharp\), and write
\(\bar B=B_\sharp L/(L-1)\).
With \(\varpi^c\) from \eqref{II-91-22} at
\(\nu=1\), \(\delta_{\rm pc}=\delta_{\rm W}/2\),
\begin{equation}\label{II-107-45}
 |\mathbb V^{\rm c}(Y)|
 \le K_c\vartheta^a(\sup|\mathcal A|)M^3\bar B
       \sum_{m=0}^k\varpi^c_{k-m}L^{k-m}a_m^2
       e^{-r_Vd_k(Y)},
 \quad K_c=12(1+2\mathsf c_{\rm pl})^2C_{\rm h}.
\end{equation}
If \(c_-\alpha_{0,m}\le a_m\le c_+\alpha_{0,m}\),
\(\alpha_{0,m}=C_0g_my_m^q\),
\(x_K\ge2\bar B\), \(y_K\ge\log(2L)\), and
\(4\sqrt L\,\mathsf c_{\mathfrak h}c_-^{-1}
 C_{\mathcal A}\vartheta_{0,k}\le1\), then
\begin{equation}\label{II-107-48}
 |\mathbb V^{\rm c}(Y)|
 \le K'_c\mathsf c_{\mathfrak h}\frac{c_+^2}{c_-}
 C_{\mathcal A}\vartheta_{0,k}M^3\bar B C_0^2
 g_k^2y_k^{2q}e^{-r_Vd_k(Y)},
\end{equation}
where
\(K'_c=8\sqrt L\,K_c C'_{\rm h}
 \sup_{i\ge0}(1+i)^{2q-1}L^{-i}\).
\end{proposition}

\begin{proof}
The proof of \eqref{II-106-37} with curvature radius
\(a_m\) bounds the interpolation by
\[
 6(1+2\mathsf c_{\rm pl})^2M^3\|p_w\|_\infty L^{k-m}a_m^2.
\]
Its scalar radius can be taken as
\[
 a_m/[e\mathsf c_{\mathfrak h}\sup|\mathcal A|
                e^{-\delta_{\rm W}\mathsf d_k(X)}].
\]
Splitting the distance exponential in half and using
\eqref{II-3-11} bounds its reciprocal by
\[
 \vartheta^a(\sup|\mathcal A|)e^{-\delta_{\rm W}\mathsf d_k(X)/2}.
\]
Occupied domains have \(d_m(X)=0\), are placed with
\(\nu=1\), and \(\|p_w\|_\infty\le\bar B\).
Apply \eqref{II-91-58} to obtain \eqref{II-107-45}.

The running-coupling bounds give
\(L^{k-m}g_m^2\le2g_k^2\),
\(y_m\le(1+k-m)y_k\), and
\(\alpha_{0,k}/\alpha_{0,m}\le
\sqrt{3/2}\,L^{(k-m)/2}\).
The distance floor absorbs the latter growth in
\eqref{II-3-11}, leaving
\(\vartheta^a(\sup|\mathcal A|)
\le4\sqrt L\,\mathsf c_{\mathfrak h}c_-^{-1}
 C_{\mathcal A}\vartheta_{0,k}\).
Use \(\sum(i+1)L^i\varpi_i^c\le C'_{\rm h}\) and
\((1+i)^{2q}\le
[\sup(1+i)^{2q-1}L^{-i}](i+1)L^i\).
This proves \eqref{II-107-48}.
\end{proof}


\subsection{Bounds for the inherited terms}

Put
\[
 \mathsf s_k^{\rm fx}=\delta_k/\alpha_0,\qquad
 \vartheta_k^{\mathsf{Tab}}=\delta_k/
             \mathsf r_0((k,\mathrm{out});k),
\]
and require
\(\mathsf H_{\rm sb}
 \max\{\mathsf s_k^{\rm fx},\vartheta_k^{\mathsf{Tab}}\}\le1/2\).
The radius comparisons in \cref{II-104-17} are imposed
for every earlier output structure used below.

Assume the old pointed terms are analytic local terms
on the fixed cores, with
\(\|E^{(j)}\|_{j,\kappa}\le\mathsf g_j\), and
\(|b_j|\le B_\sharp\). Set
\begin{equation}\label{II-108-22}
 \begin{split}
 \mathsf g'_j&=\mathsf g_j+\alpha_0^2B_\sharp,\qquad
 C_{\rm ert}^{\flat}=L^4C_{\rm oc}\max\{C_{\mathrm E,2},1176\},\\
 C_{\rm ert}&=8M^3e^{26\kappa}C_{\rm ert}^{\flat},\qquad
 B_m^{\rm E}=C_{\rm ert}
       \left(\mathsf g'_m+\sum_{j<m}L^{-(m-j)}\mathsf g'_j\right),\\
 \mathsf{size}^{\rm E}&=\sum_{m=1}^k\varpi_{k-m}B_m^{\rm E}.
 \end{split}
\end{equation}
Here and for the regular list use \(\nu=3\),
\(\delta_{\rm pc}=\delta_{\rm sb}\) in the placed count.

\begin{proposition}\label{II-108-24}
Assume \(\kappa>\kappa_{\mathcal D}+2\) and the preceding
rate and domain conditions. The regular potential is
defined on the full graded domain. Its wide and octet
parts satisfy
\begin{equation}\label{II-108-30}
 |\mathbb V^{\rm E,wd}(Y)+\mathbb V^{\rm E,oc}(Y)|
 \le C_V C_{\mathcal A}\mathsf s_k^{\rm fx}
        \mathsf{size}^{\rm E}e^{-r_Vd_k(Y)},\qquad
 C_V=2C_{\rm h}\mathsf C_{\rm sb}
                         e^{\mathsf C_{\rm sb}\bar\kappa},
\end{equation}
where
\(\mathsf{size}^{\rm E}\le
[L/(L-1)]C'_{\rm h}C_{\rm ert}\sup_j\mathsf g'_j\).
If \((\mathrm{CR})_t\) and
\(\mathsf r_0((j,\mathrm{out});j)^2
 \le r_\star^2(1+k-j)^{a'}\), its edge-tent part satisfies
\begin{equation}\label{II-108-32}
 |\mathbb V^{\rm E,et}(Y)|
 \le C'_V C_{\mathcal A}\vartheta_k^{\mathsf{Tab}}
            M^3B_\sharp r_\star^2 e^{-r_Vd_k(Y)},
\end{equation}
with
\(C'_V=3e^{26\kappa}c(a',L)C_{\rm h}C'_{\rm h}
\mathsf C_{\rm sb}e^{\mathsf C_{\rm sb}\bar\kappa}\) and
\(c(a',L)=\sup_{i\ge0}(1+i)^{a'-1}L^{-i}\).
The edge-tent part vanishes unless
\eqref{eq:historical-source-distance} holds.
\end{proposition}

\begin{proof}
The octet provision \cref{II-49-44} gives, at
\(n=\mathsf{ol}_j(z)>j\),
\[
 \|f_{j,z}\|\le C_{\rm ert}^{\flat}\mathsf g'_jL^{-4(n-j)}.
\]
Indeed, apply its bound to each term before summing;
the pointed entropy sum controls
\[
 \sum_{X\ni z}e^{-\kappa d_j(X)}(d_j(X)+2)^2,
\]
and the subtracted tent has \(d_j(X_z)\le26\),
so its contribution is at most
\[
 (3/2)28^2\alpha_0^2B_\sharp=1176\alpha_0^2B_\sharp.
\]
At fixed \(n,X'\) there are at most
\(8M^3L^{3(n-j)}\) choices of \(z\), since \(z\in X'\)
and \(X'\) has at most eight level-\(n\) cubes.
Together with the wide-term bound this yields
\begin{equation}\label{II-108-26}
 \sum_{\substack{\text{wide or octet }f\\
                         j_f=m,\ X_f=X}}\|f\|
 \le B_m^{\rm E}(d_m(X)+1)e^{-\kappa d_m(X)} .
\end{equation}

These domains are placed with \(\nu=3\).
For clarity, if \(n_0=\mathsf{ol}_j(z)<k\), maximality
provides a point of \(X'_{n_0+1}(z)\) outside
\(\Omega_{n_0+2}\), at distance at most \(2M\)
from \(z\) in level-\((n_0+1)\) units.
Every point of \(\Omega_n\), \(n\ge n_0+3\), is
therefore at distance at least \(7ML^n\) from \(z\)
in bare units. The diameter bound proves placement
for a wide member meeting these layers.
If \(n\le n_0+2\), either \(n_0=j\), when the required
lower bound is nonpositive, or wideness puts a point
outside \(\square_{n_0}(z)^{\sim2}\), at distance
greater than \(2ML^{n_0-1-j}\) from \(z\) in level-\(j\)
units. This proves the remaining bound.
An octet member cannot meet \(\Omega_n\) for
\(n\ge n_0+4\), since all its points lie within
\(2ML^{n_0}\) of \(z\); its other placement bounds
are nonpositive.

Apply \eqref{II-104-18} to each member and then
\eqref{II-91-58} to \eqref{II-108-26}.
This gives \eqref{II-108-30}.
The size bound follows by summing \(L^{-i}\) and
using the last inequality of \eqref{II-91-58}.

An edge tent at \(y\notin\Lambda_j^{\sim-1}\)
does not meet \(\Omega_{j+1}\); it is both admitted
and placed. Its base need not lie in \(\Lambda_j\),
so use the layered space, not the fixed-core case.
Every plaquette with nonzero weight
\(\varphi_jh_y^{(j)}\) has its corners in
\(X_y\cap\Lambda_j\), hence all its bonds belong to
\(\operatorname{Bd}_{\mathcal S}(X_y)\).
The function \(\widehat f(P;\mathfrak q)=f(\mathbf U_P)\)
is therefore an entire function of the complex potential
on this layered space, independent of the stored data.
The layer-\(j\) curvature bound and
\(\sum_ph_y^{(j)}(p)=3\eta_j^{-3}\) give
\( |\widehat f|\le(3/2)|b_j|
\mathsf r_0((j,\mathrm{out});j)^2\).
The pinned substituted configuration is a complex gauge
transform of the plain shifted configuration on these
bonds. The turn factors cancel around each plaquette,
leaving conjugation at its base; its trace is unchanged.
Thus the pinned reading of \(\widehat f\) is the plain
configuration reading of the edge tent. Apply
\cref{II-104-17} to this layered reading, with
\(\mathsf a=\vartheta_k^{\mathsf{Tab}}\).
At a fixed domain \(X_y\) there are
at most \(M^3\) such sites, and the weighted
plaquette estimate bounds the sum of their norms by
\((3/2)M^3B_\sharp\mathsf r_0((j,\mathrm{out});j)^2\).
Since \(d_j(X_y)\le26\), multiply by
\(e^{26\kappa}(d_j(X_y)+1)e^{-\kappa d_j(X_y)}\)
and apply the same placed count, now with
\(\mathsf a=\vartheta_k^{\mathsf{Tab}}\).
The inequality
\((1+i)^{a'}\le c(a',L)(i+1)L^i\) proves
\eqref{II-108-32}. The admitted-domain separation
gives its support assertion.
\end{proof}

For the boundary terms suppose the old terms on layered
spaces are analytic, invariant with their stored
coordinates and tied, and the edge sums are analytic
local terms on their fixed cores. Suppose each has
norm at most \(\mathsf B'e^{-\kappa d_j(X)}\)
on that whole space. For the formed family suppose
its members are in their prescribed places and
\begin{equation}\label{eq:historical-factor-input}
 \sum_{f:\,j_f=l+1,\ X_f=X}\operatorname{osc}\varphi_f
 \le B_{\rm fr}\mathfrak b_l^2e^{-\kappa d_{l+1}(X)} .
\end{equation}
Each exponent is defined and analytic on its entire
layered space and stored tier, invariant with those
stored coordinates, and tied. Its oscillation is
the supremum, at fixed stored datum, of the difference
between its values at two triples. Thus additive
functions of the stored data are immaterial.

\begin{proposition}\label{II-109-28}
Assume \((\mathrm{CR})_t\) and \(x_K\ge2B_\sharp\).
The boundary and formed potentials are defined on
all stored tiers. They vanish without a source cube
and unless \eqref{eq:historical-source-distance} holds.
For \(a'\ge0\),
\begin{equation}\label{II-109-32}
 \begin{aligned}
 y_k^{a'}|\mathbb V^{\rm B,mk}(Y)|
 &\le C_B(a')\mathsf B'C_{\mathcal A}
 \vartheta_k^{\mathsf{Tab}}e^{\mathsf C_{\rm sb}\bar\kappa}
                         e^{-r_Vd_k(Y)},\\
 y_k^{a'}|\mathbb V^{\rm B,ed}(Y)|
 &\le C_B(a')\mathsf B'C_{\mathcal A}
 \mathsf s_k^{\rm fx}e^{\mathsf C_{\rm sb}\bar\kappa}
                         e^{-r_Vd_k(Y)},
 \end{aligned}
\end{equation}
where, with \(c=1-2\delta\),
\[
 C_B(a')=\max\left\{1,\frac{16}{c}
 C_\square(L,0)\mathsf C_{\rm sb}\,
 e^{2c}(1+\log(2L))^{a'}
                 \left(\frac{a'}{7ecL}\right)^{a'}\right\},
 \qquad 0^0=1.
\]
The formed potential satisfies
\begin{equation}\label{II-111-30}
 |\mathbb V^\Phi(Y)|
 \le C_\Phi A_1^2B_{\rm fr}C_{\mathcal A}
 \vartheta_k^{\mathsf{Tab}}e^{\mathsf C_{\rm sb}\bar\kappa}
 e^{-L\check R_t}e^{-r_Vd_k(Y)},
\end{equation}
provided either \(\delta\le1/8\), or
\(\delta>1/8\) and \(2\delta\kappa\ge c_{\mathrm E}+3\).
Here
\begin{equation}\label{II-111-7}
 C_\Phi=16e^{4/7}(2\mathfrak p_0)!
 (1+2\log(2L))^{2\mathfrak p_0}
 C_\square(L,2\mathfrak p_0)\mathsf C_{\rm sb}.
\end{equation}
All these potentials inherit analyticity, the tie,
gauge invariance with the stored fields, and the
zero-field property. They inherit Borel measurability,
conjugation reality and dependence on the pair
when the input terms have those properties.
Their values are linear in the input terms.
\end{proposition}

\begin{proof}
The full-domain assertion is \cref{II-104-17};
the old domain and tier are precisely the corresponding
layered or fixed-core domain and tier in that proposition. In particular no choice of
representative outside those spaces is used.
The layered boundary and edge domains are admitted under
\((\mathrm{CR})_t\). A formed factor is in its place
exactly when its domain is admitted. Thus every
nonzero summand has the separation
\eqref{eq:historical-source-distance}.

For either boundary part, write \(\mathsf a\) for
\(\vartheta_k^{\mathsf{Tab}}\) or \(\mathsf s_k^{\rm fx}\).
The single-member bound and the per-fine-cube proof
of \eqref{II-91-56}, before absorbing \(d+1\), give
\[
 \sum|\mathbb V^\star[f](Y)|
 \le16\mathsf C_{\rm sb}C_\square(L,0)
 C_{\mathcal A}\mathsf a e^{\mathsf C_{\rm sb}\bar\kappa}
 \mathsf B'(d+1)e^{-(1-2\delta)\kappa d},
 \qquad d=d_k(Y).
\]
Split the exponent as \(r_Vd+cd+cd\).
Use \((d+1)e^{-cd}\le1/c\) and
\(d\ge7L\check R_t-2\).
The running bounds imply
\(y_k\le(1+\log(2L))\check R_t\);
maximizing \(R^{a'}e^{-7cLR}\) proves
\eqref{II-109-32}.

For formed terms use the oscillation version of
\eqref{II-104-18} and group them by \(j=l+1,X\).
The bound \eqref{eq:historical-factor-input} gives
the same sum with coefficient
\(2\mathsf H_{\rm sb}\vartheta_k^{\mathsf{Tab}}
 B_{\rm fr}\mathfrak b_{j-1}^2\).
Put \(\Upsilon_k=y_k+\log(2L)\).
Then
\(\mathfrak b_{j-1}^2\le
 A_1^2\Upsilon_k^{2\mathfrak p_0}
 (1+k-j)^{2\mathfrak p_0}\).
Apply the admitted count with rate split
\(\widehat a=c_{\mathrm E}+1\). Its output rate is
\(\kappa-c_{\mathrm E}-2\).
The difference between that rate and \(r_V\) is
\(2\delta\kappa-c_{\mathrm E}-4\delta\), at least
\(2/7\) in either stated case.
Since
\(\Upsilon_k\le(1+2\log(2L))\check R_t\),
\[
 \Upsilon_k^{2\mathfrak p_0}
 e^{-\frac27(7L\check R_t-2)}
 \le e^{4/7}(2\mathfrak p_0)!
 (1+2\log(2L))^{2\mathfrak p_0}e^{-L\check R_t}.
\]
Here \(R^ne^{-R}\le n!\) follows from the exponential
series. This proves \eqref{II-111-30}.
Subtracting any one value at fixed stored datum
reduces the oscillation estimate to the bounded case;
the starred bracket kills the subtracted value.
No measurable choice is needed for this estimate.
The remaining properties follow termwise from
\cref{II-104-17} and finite summation.
\end{proof}


\subsection{The exact finite sums}

Let \(U_+\) be a real representative of the minimal
orbit for the longer collared sequence, let
\(\mathsf U\) extend it to all bare bonds, and let
\(P_+=(\mathsf U,0,\mathbf J_+)\) be its plain triple.
Assume its restrictions lie in the graded domains
above, that \(V=\mathcal M^kU_+\) is axial on \(P''\),
and that \(\mathcal A\) is real with
\(\sup|\mathcal A|<\epsilon_P\).
At the shorter reduced determining set \(\mathbf B\),
let \(\mathsf d_z\) equal \(e^{iB'_z}V\) on the changed
top rows and \(M_{\mathbf B}(U_+)\) elsewhere; \(B'_z\)
is the actual pivot completion of \(z\mathcal A\).
Averages on these rows are the real averages.
Their agreement with the analytic response averages
follows from the local double-pillar induction in
\cref{II-195-31}, using its curvature hypothesis on
each double pillar.

\begin{proposition}\label{II-106-49}
Suppose first that \(k\ge1\), and retain the
separated-sequence and regularity ceilings
of \cref{I-58-22}, with reduced determining set
\(\mathbf B\) and
\(\bar\varepsilon\le\varepsilon_{\mathrm G}(\rho_{\rm s})\).
If the old tested datum is \(\epsilon'\)-regular and
\begin{equation}\label{eq:historical-real-segment}
 \epsilon'+8C_P\sup|\mathcal A|\le\bar\varepsilon,
\end{equation}
then the minimal orbit of \((\mathbf B,\mathsf d_z)\)
exists for every real \(z\in[-1,1]\).
Define
\(F(z)=A^{\eta_k}_{\Omega^\diamond_0}
             (U_{\mathbf B}(\mathsf d_z))\)
on this interval. It extends holomorphically to a
neighborhood of \(|z|\le1\). With the regional pivot
Jacobian \(\mathsf J^\diamond_V\), gauge functional
\(\mathbb G\), and gauge Hessian \(\widehat{\mathsf G}_V\),
the exact sums are
\begin{align}
 \sum_Y\mathbb V^{\rm J}(Y;P_+|_Y,B|_Y)
   &=\mathsf J^\diamond_V[\mathcal A],\label{II-106-50}\\
 \sum_Y\mathbb V^{\rm G}(Y;P_+|_Y,B|_Y)
   &=-x_k\mathbb G[\mathcal A]
                    +\tfrac12(B,\widehat{\mathsf G}_VB)_1,\label{II-106-51}\\
 \sum_Y\mathbb V^{\rm act}(Y;P_+|_Y,z\mathcal A|_Y)
   &=F(z)-F(0),\label{II-106-52}\\
 \sum_Y\mathbb V^{\rm A}(Y;P_+|_Y,B|_Y)
   &=-x_k\{F(1)-F(0)-F'(0)-\tfrac12F''(0)\}.\label{II-106-53}
\end{align}
For the actual counterterms \(b_j=L^{-1}x_{j-1}-x_j\),
the following sum is over the inner bare plaquettes of
\(\Omega_0^\diamond\):
\begin{equation}\label{II-107-57}
 \sum_Y\mathbb V^{\rm c}(Y)
 =-\sum_{p\ {\rm inner}}
   [\beta_k(x(p))-L^kx_k]\,
       [a_p(U_{\mathbf B}(\mathsf d_1))-a_p(U_+)].
\end{equation}
All sums are finite and evaluated at
\((P_+|_Y,\mathcal A|_Y)\); \(B=g_k^{-1}\mathcal A\).
At \(k=0\), the determining map is the identity and
\(U_{\mathbf B}(\mathsf d_z)=\mathsf d_z\). With this
convention the displayed sum identities hold by the
direct first-step calculation.
\end{proposition}

\begin{proof}
For \(k\ge1\), the completed field is real and
\(\sup|B'_z|\le2C_P|z|\sup|\mathcal A|\).
Only first-kind rows at the last level change.
A tested plaquette at a lower level is unchanged;
at the last level it is a product of four unitary
factors, each changing by at most
\(2C_P\sup|\mathcal A|\).
Telescoping the product, including reversed bonds,
gives a change at most \(8C_P\sup|\mathcal A|\).
The strict original test and
\eqref{eq:historical-real-segment} therefore give
\(\bar\varepsilon\)-regularity. The reduced determining
set, shorter separated sequence and lattice-size
hypotheses are those of \cref{I-58-22}.
Apply that theorem at datum \(V_0=\mathsf d_z\)
and perturbation \(B=0\), separately for each real
\(z\). It gives the entire real path of regular orbits.

For the Jacobian decompose the bond space into
nonpivot and pivot coordinates. Its derivative
maps into the pivot coordinates; the pivot-to-pivot
block is the direct sum of the individual matrices
\(\mathsf{Jb}_c\), since another pivot never lies in
the block of \(c\).
Thus the trace of every power is the sum of the
individual traces. Their norms are less than \(1/2\),
so the logarithmic series converges absolutely.
Summing the potentials counts each pivot at
\(\square(c)\) and each nonpivot Haar factor at
\(\square(b)\), exactly once. This proves
\eqref{II-106-50}.

Inside \(B(y)\), \(y\in P''\), there is no pivot and
the completion equals \(\mathcal A\).
The contour logarithm
\(\xi(a)=i^{-1}\log[(e^{ia}V)(y,x)V(y,x)^{-1}]\)
has first term \(\mathsf G^{(1)}[V]a\).
Since it is traceless,
\(1-\frac12\operatorname{tr}e^{i\xi(za)}
=\frac12z^2\langle\mathsf G^{(1)}a,\mathsf G^{(1)}a\rangle
+O(z^3)\).
The linear completion fixes these nonpivot bonds.
Summing over the disjoint blocks and using
\(x_kg_k^2=1\) proves \eqref{II-106-51}.

For a Wilson member the finite starred sum is its
full-parameter value minus its zero-parameter value:
all disconnected additions vanish by locality, and
restriction agrees by zero extension of decoupling
parameters. The inner bare plaquettes of \(\Omega_0^\diamond\)
are partitioned by \((m,X)\); hence summing those
differences gives the action difference of the actual
full response.
By \cref{II-93-38} it is a gauge transform of the
minimal orbit of \(\mathsf d_z\). Each plaquette
action is gauge invariant, proving
\eqref{II-106-52}. For \(k=0\), each plaquette is
instead counted directly at its \(X(p)\), and the
same identity holds with the completed field.
Each summand is holomorphic for
\(|z|\sup|\mathcal A|<\epsilon_A\), a disc of radius
greater than one. Their finite sum extends \(F-F(0)\).
Taylor subtraction commutes with the finite sum,
proving \eqref{II-106-53}.

The same finite calculation applies with any fixed
plaquette weight. For \(p_w\), the weighted action
difference enters with a minus sign.
The telescoping running-coupling identity gives
\(L^kp_w(p)=\beta_k(x(p))-L^kx_k\).
This proves \eqref{II-107-57}; no identification of
\(b_j\) with an intrinsic local coefficient is used.
\end{proof}

For the remaining sums let
\(\mathsf U^\sharp=e^{i\eta_k\mathbf A^\delta}\mathsf U\)
be the pinned full-response field. Suppose it is
real and every old plain triple at \(\mathsf U^\sharp\)
belongs to its layered output space, and to the
fixed core when \(X\subset\Lambda_j\). The corresponding
memberships at \(\mathsf U\) are required as well.
For actual glued fields these are the output and
old-orbit assertions of \cref{II-99-39,II-113-38}.
They are stronger than membership of a substituted
triple with real factor \(\mathsf U\).

\begin{proposition}\label{II-111-28}
Under these physical domain conditions, all old
terms at \(\mathsf U\) and \(\mathsf U^\sharp\) are
defined on their prescribed stored tiers, and
\begin{align}
 \mathbf E_k(\mathsf U^\sharp)-\mathbf E_k(\mathsf U)
     &=\sum_Y\mathbb V^{\rm E}(Y;\mathfrak u_+|_Y),
                                          \label{II-108-27}\\
 \mathbf E_k^{\rm bd}(\mathsf U^\sharp,\mathfrak q)
       -\mathbf E_k^{\rm bd}(\mathsf U,\mathfrak q)
     &=\sum_Y\mathbb V^{\rm B}(Y;\mathfrak u_+|_Y),
                                          \label{II-109-31}\\
 \boldsymbol\Phi_h(\mathsf U^\sharp;\mathfrak q)
       -\boldsymbol\Phi_h(\mathsf U;\mathfrak q)
     &=\sum_Y\mathbb V^\Phi(Y;\mathfrak u_+|_Y).
                                          \label{II-111-29}
\end{align}
The last identity is a separate identity of the old
formed exponents; it does not absorb them into the
regular or boundary terms.
\end{proposition}

\begin{proof}
At full parameters the current identity in
\cref{II-94-49} shows that the pinned substituted
triple and the old plain triple of
\(\mathsf U^\sharp\) have the same pair. Both are
real and admitted; the tie identifies their values
at the same stored datum. At zero parameters,
restriction of \(P_+\) is exactly the old plain
triple of \(\mathsf U\). Thus \eqref{II-104-20}
is the physical difference for every layered boundary or
formed member. For a plain local member, the
plain and pinned readings agree by the proof
of \eqref{II-85-37}.

For the regular list, every old pointed term is
either wide or occurs in exactly one octet sum.
The octet identity of \cref{II-49-44} applies
because \(X'_n(z)\subset\Omega_{n+1}\) and its
five-layer cube lies in the domain of the minimal
configuration. It identifies each pulled-back
narrow term with its original value.
The tents partition the weighted action:
deep tents have \(\varphi_jh_z=h_z\), and the
remaining tents are exactly the edge tents.
Consequently the sum of the regular members is
\(\mathbf E_k\) plus a field-independent constant;
that constant cancels in the difference.
This proves \eqref{II-108-27}.

For the boundary list apply the single-member
difference separately to layered boundary terms, with stored
levels \(<j\), and edge terms, with no stored
argument. Reordering the two finite sums proves
\eqref{II-109-31}.
For a formed factor of level \(l+1\), restrict the
stored argument to \((\mathfrak q|_{X_f})|^{(l)}\).
The physical domain assumptions put both plain
triples in its entire defined space. Thus both
old exponents are defined; their difference is the
sum of that factor's pinned potentials.
Summing first over factors and then over domains
gives \eqref{II-111-29}. Additive functions of the
stored datum cancel on both sides. This also explains
why the oscillation bound, without an absolute bound
on each exponent, suffices.
\end{proof}


\subsection{The actual families in the history expansion}

For a family \(V\), let \(\|V\|_{r_V}\) be the
supremum of \(e^{r_Vd_k(Y)}|V(Y)|\) over the full
graded field domain. A superscript \({\rm rad}\)
restricts the stored coordinates to their prescribed
tiers. Thus the two norms differ only in their stored
arguments, not in their field cutoff or graded triple.
Put
\[
 \mathbb V^{\rm rg}=\mathbb V^{\rm int}+\mathbb V',
 \qquad \mathbb V^{\rm bd}=\mathbb V^{\rm B}+\mathbb V^\Phi,
 \qquad \vartheta_{0,k}=\frac{\delta_k}{\alpha_{0,k}},
 \qquad C_{\rm sz}=\frac L{L-1}C'_{\rm h}.
\]
Suppose the old pointed norm at rate \(\kappa+1\)
is at most \(E_0\); each term on a layered space and each edge sum
has bound \(\mathsf B e^{-\kappa d_j(X)}\); and the
formed family satisfies \eqref{eq:historical-factor-input}.
The terms are Borel and conjugation real, and have
the invariance and tie required above. These are the
input bounds on the whole stated spaces, not just
at the actual real backgrounds.

For the radius specialization assume
\[
 \begin{gathered}
 \alpha_{\iota,n}=C_\iota g_ny_n^q,
 \qquad q>\mathfrak p_0,\qquad
 \alpha_{0,k}\le\alpha_0,\qquad
 c_-\alpha_{0,k}\le\mathsf r_0((k,\mathrm{out});k),\\
 \mathsf r_0((j,\mathrm{out});j)\le c_+\alpha_{0,j},
 \qquad
 c_-\alpha_{0,m}\le a_m\le(1+\beta_{\rm s})c_+\alpha_{0,m}.
 \end{gathered}
\]
Here the last inequality concerns the enlarged radii
\(a_m\) used for the coupling potential. Retain all
domain, pinning, rate and smallness inequalities of
\cref{II-104-17,II-106-39,II-107-44,II-108-24,II-109-28};
for the edge-tent estimate use \(a'=2q\).
These are the corresponding radius and cutoff requirements
in \cref{sec:common-iteration-constants}.

Define
\begin{equation}\label{II-112-26}
\begin{aligned}
 C^{\rm rg}={}&K_\flat M^3C_{\mathcal A}\alpha_0
 +K_\flat C_{\mathcal A}^3
 +2K_A\mathsf c_{\mathfrak h}\epsilon_A^{-2}C_{\mathcal A}^3
 +C_VC_{\rm sz}C_{\rm ert}C_{\mathcal A}\\
 &+K'_c\mathsf c_{\mathfrak h}(1+\beta_{\rm s})^2
       c_+^2c_-^{-1}C_{\mathcal A}M^3\bar B C_0^2
 +C'_VC_{\mathcal A}c_-^{-1}c_+^2M^3B_\sharp C_0^2,\\
 C^{\rm bd}={}&C_{\mathcal A}e^{\mathsf C_{\rm sb}\bar\kappa}
       [C_B(0)(1+c_-^{-1})+C_\Phi A_1^2c_-^{-1}],\\
 C_{\rm eff}={}&\max\{1,C^{\rm rg},C^{\rm bd}\}.
\end{aligned}
\end{equation}
These constants depend on the fixed geometric, analytic,
radius and input bounds, but not on the history,
height, volume or individual running couplings.

\begin{proposition}\label{II-119-60}
The four actual historical families are defined,
analytic and Borel on every prescribed tier, are
gauge invariant and conjugation real, and have the
tie. The regular families are independent of the
stored coordinates and depend only on the pair.
All four vanish at zero field and on a domain
without a source cube. Moreover
\begin{align}
 \|\mathbb V^{\rm int}\|_{r_V}+\|\mathbb V'\|_{r_V}
 &\le C_{\rm eff}\left[
 \frac{\delta_k}{\alpha_0}(1+E_0+\alpha_0^2B_\sharp)
 +M^3g_k\mathfrak b_k^3
 +\vartheta_{0,k}g_k^2y_k^{2q}\right],\label{II-112-53}\\
 \|\mathbb V^{\rm B}\|^{\rm rad}_{r_V}+\|\mathbb V^\Phi\|^{\rm rad}_{r_V}
 &\le C_{\rm eff}(\mathsf B+B_{\rm fr})\vartheta_{0,k}.
                                                        \label{II-112-54}
\end{align}
The norms of the combined regular and boundary
families are bounded by these respective left sides.
The sums of norms also give the per-origin weights
required in the history cluster expansion. They hold
on its conditional-mean field polydiscs whenever
those polydiscs lie in the original field cutoff.
For \(k=0\) the inherited families are empty and the
regular bound is the direct Jacobian, gauge and
action bound of \cref{II-106-39}, with
\(\epsilon_A=r_P^\circ\) and \(\mathsf c_{\mathfrak h}=1\).
\end{proposition}

\begin{proof}
The source-cube argument in \cref{II-106-39} admits
all Jacobian, gauge and action members on the full
graded domain. For the earlier terms,
\cref{II-104-17} admits every substituted triple
and every stored argument on its correct space.
The edge-tent argument in the proof of
\cref{II-108-24} supplies the same conclusion for
its plain configuration reading. The bounds on the
input terms therefore apply to every summand;
taking a finite sum leaves no additional domain
condition. This remains true when the sum is empty,
since the family was restricted to that domain in
its definition.

The invariance, tie and pair dependence follow from
the re-expansion proof, including the independent-pin
argument \eqref{II-85-37}, and from the action proof.
Borel measurability is preserved by composition and
finite summation. At fixed real factor and stored
argument each summand has a holomorphic extension
to a neighborhood of the point; intersect these
finitely many neighborhoods and add the extensions.
Restricting the stored variables to their tiers does
not change this slice. It also preserves Borel
measurability, since the preimage of a Borel set is
intersected with the restricted domain. This proves
all stated regularity properties on the tiers.
The zero-field and no-source assertions were proved
for the individual members.

For the norms, the running inequalities give
\(g_j\le g_k\) and \(y_j\le(1+k-j)y_k\).
Consequently the edge-tent hypothesis holds with
\(r_\star=c_+C_0g_ky_k^q\), and
\[
 \vartheta_k^{\mathsf{Tab}}\le c_-^{-1}\vartheta_{0,k},
 \qquad \mathsf s_k^{\rm fx}\le\vartheta_{0,k}.
\]
Use \(\mathsf g'_j\le E_0+\alpha_0^2B_\sharp\)
in \eqref{II-108-22}. The wide and octet bound is
then \(C_VC_{\rm sz}C_{\rm ert}C_{\mathcal A}\)
times \(\mathsf s_k^{\rm fx}(E_0+\alpha_0^2B_\sharp)\),
with the common factor \(e^{-r_Vd_k(Y)}\).
The Jacobian bound is
\(K_\flat M^3C_{\mathcal A}\alpha_0\mathsf s_k^{\rm fx}\).
The gauge and action bounds are respectively
\(K_\flat C_{\mathcal A}^3\) and
\(2K_A\mathsf c_{\mathfrak h}\epsilon_A^{-2}C_{\mathcal A}^3\)
times \(M^3g_k\mathfrak b_k^3\).
The last two coefficients in \(C^{\rm rg}\) multiply
\(\vartheta_{0,k}g_k^2y_k^{2q}\), by
\eqref{II-107-48} and \eqref{II-108-32}.
Taking each supremum separately and adding gives
\eqref{II-112-53}; no cancellation between origins
is used.

For the boundary sum, take \(a'=0\) in
\eqref{II-109-32}, and use the two ratio inequalities
above. For the formed sum use \eqref{II-111-30}
and \(e^{-L\check R_t}\le1\). Both coefficients are
bounded by \(C^{\rm bd}\). Taking the two suprema
separately gives \eqref{II-112-54}.
The domain restriction of the norm can only decrease
these suprema. In the conditional Gaussian application,
\cref{II-119-53} supplies the conditional mean and the
cutoff inclusion; evaluating an already defined family
on that smaller field polydisc preserves all the bounds.
For the actual families write \(\mathsf B=B\) and
\(\mathsf B_{\rm fr}^{\flat}=B_{\rm fr}^{\flat}\).
Then \cref{II-157-15} supplies
\eqref{eq:historical-factor-input} with the smaller
coefficient \(\mathsf B_{\rm fr}^{\flat}\) in place of
\(B_{\rm fr}\).
Use this coefficient and the actual boundary bound
\(\mathsf B\) in \eqref{II-112-54}.
Then the two displayed majorants are precisely
\(S_{\rm rg}(x_k)\) and \(S_{\rm bd}(x_k)\) of
\cref{sec:common-iteration-constants}, since
\(g_k=e^{-y_k/2}\), \(\mathfrak b_k=A_1y_k^{\mathfrak p_0}\)
and \(\vartheta_{0,k}=(A_1/C_0)y_k^{-(q-\mathfrak p_0)}\).
Their smallness is obtained there by the single fixed
choice of the coupling ceiling.
\end{proof}

\subsection{The actual shifted field}
\label{sec:history-actual-arrival}

We now verify that the local potentials give the change
in the old exponent on the support of a history.
Let \(h'=(h,\mathsf o_t)\), \(t=k+1\), with \(k\ge1\).
Take \(U\in\mathsf{Sup}(h')\) where the stored variables
are defined, and an adapted representative \(U_+\)
of its collared critical orbit at level \(t\).
Glue \(U_+\) to \(U\) outside
\(\Omega^\diamond_0=\Omega_1^\sim\), obtaining
\(\mathsf U\). Let
\[
 P_+=(\mathsf U,0,\mathbf J_+),\qquad
 V=\mathcal M^kU_+,\qquad
 p=V_k^{{\rm bg}\diamond}\uplus
                  V_k|_{\mathfrak X_{\rm out}},
 \qquad B=g_k^{-1}\mathcal A .
\]
Here \(\mathcal A\) is the actual free fluctuation and
\(\mathfrak q\) the old stored field. The pair read at
level \(k\) is
\((U_+,L^{-3}\mathbf J_+)\) on the inner bare bonds.
Write \(\mathbf B\) for the shorter determining set,
through level \(k\).

Use the domain, radius and smallness hypotheses of
the actual potential construction and pinned response.
The additional chart and segment inequalities are
\begin{equation}\label{eq:history-arrival-ceilings}
 \begin{gathered}
 2L^2C_{\rm reg}\varepsilon_t\le\varpi_\chi,\qquad
 C_\chi\mathsf C_{\mathcal A}\delta_k<r_\chi,\\
 2L^2C_{\rm reg}\varepsilon_t
       +C_{\rm seg}\mathsf C_{\mathcal A}\delta_k
                                      \le\bar\varepsilon .
 \end{gathered}
\end{equation}
The pair is required to lie within the radii of the
regional response and fluctuation form, and the
level-\(t\) torus has at least \(6R_1\mathfrak M\)
sites in each direction. These are among the
simultaneous requirements in
\cref{sec:common-iteration-constants}.

Let \(B'_z\) be the completed pivot field at
\((P_+,z\mathcal A)\). Denote by \(\mathsf d_z\)
the shorter datum obtained from \(M_{\mathbf B}(U_+)\)
by replacing its rows in \(\mathfrak X_0\) by
\(e^{iB'_z}V\). Set
\[
 (\mathbf A,\mathbf K)=
          \mathfrak h_{T_\ell|h'}(\mathbf1;P_+,\mathcal A),
 \qquad
 \mathsf U^{\rm pl}=e^{i\eta_k\mathbf A}\mathsf U,\qquad
 \mathsf U^\sharp=e^{i\eta_k\mathbf A^\delta}\mathsf U .
\]
The last field uses the pinned levelling.

\begin{proposition}\label{II-113-38}
For every real \(z\in[-1,1]\), the completed field is
defined and agrees with the regional chart:
\begin{equation}\label{eq:history-completed-chart}
 e^{iB'_z}V
       =V_k^{{\rm bg}\diamond}[z\mathcal A]
                               \quad\hbox{on }\mathfrak X_0 .
\end{equation}
In particular,
\[
 V_k|_{\mathfrak X_0}
       =V_k^{{\rm bg}\diamond}[\mathcal A],
 \qquad
                    \mathsf d_1=\xi^\diamond_h(U).
\]
The regular critical orbit of \((\mathbf B,\mathsf d_z)\)
exists throughout this segment.
The pinned field \(\mathsf U^\sharp\) belongs to the
old critical orbit on \(\Omega^\diamond_0\) and agrees
with \(U\) outside it. At every earlier index \((j,X)\)
used by the old sums, its plain triple lies in the
old layered output domain and, when \(X\subset\Lambda_j\),
in the required fixed analytic core.

Let
\[
 \mathbb F(z)=L^k A_{\Omega^\diamond_0}
                        (U_{\mathbf B}(\mathsf d_z)).
\]
The action is independent of the representative in
this orbit. It has a holomorphic continuation near
the segment, and at zero
\begin{equation}\label{II-113-39}
 \mathbb F'(0)=0,\qquad
 x_k\mathbb F''(0)+(B,\widehat{\mathsf G}_VB)_1
          =(B,\mathsf A^\sharp_{h'}(P_+)B)_1 .
\end{equation}
The gauge quadratic form is evaluated on the same
regional free coordinates as the fluctuation form.
\end{proposition}

\begin{proof}
The support and graded-domain results give real
\(\mathcal A\), with
\(\sup|\mathcal A|<\mathsf C_{\mathcal A}\delta_k\),
and put every restricted point
\((P_+,\mathcal A,\mathfrak q)|_Y\) in its admitted tier.
The shorter datum of \(U_+\) is regular at
\(2L^2C_{\rm reg}\varepsilon_t\).
The segment estimate for the completed datum in
\cref{II-106-49} bounds
its regularity radius by
\[
       2L^2C_{\rm reg}\varepsilon_t
                         +C_{\rm seg}|z|\sup|\mathcal A|.
\]
The last inequality of \eqref{eq:history-arrival-ceilings}
therefore gives the regular orbit at every \(z\).
The local response identifies these orbits and gives
the asserted holomorphic action near the segment.

\emph{The chart.}
On each double block of a coarse bond in \(\Gamma\),
\(|\partial p-1|<2L^2C_{\rm reg}\varepsilon_t\).
The first two inequalities of
\eqref{eq:history-arrival-ceilings} thus put
\(z\mathcal A\) in the regional fibre chart.
Its free coordinates equal \(z\mathcal A\),
and its outer coordinates equal those of \(p\).
The block field used by the local potentials is
also \(p\) on these double blocks. Indeed,
\(\mathsf U=U_+\) on their bare pillars; the
regularity estimate there is
\[
      |\partial U_+-1|
            <C_{\rm reg}\varepsilon_tL^{-2k}.
\]
All lower blocks required by the averaging induction
remain in the pillar. The local real and complex
averages consequently agree there, by the same
small-curvature induction used in
\cref{II-10-41}. No bound on the exterior gluing
is used.

For each pivot \(b_0(c)\), the single-bond completion
and the \(\Gamma\)-chart use the same restricted
configuration on \(\mathsf P_c\) and the same free
field, extended by zero outside \(\mathfrak X\).
The local pivot uniqueness of the fibre chart gives
\[
       B'_z(b_0(c))
          =\psi_p^{\rm lin}(\mathsf C_pz\mathcal A)(b_0(c)).
\]
Equality on the free and pivot bonds proves
\eqref{eq:history-completed-chart}.
The actual field \(V_k\) and the reference \(p\)
have the same next average \(V_t\), and the
intrinsic logarithm of their quotient on the free
bonds is \(\mathcal A\). The inverse-chart
statement therefore gives
\(V_k=V_k^{{\rm bg}\diamond}[\mathcal A]\) on
\(\mathfrak X_0\).
All other shorter rows were unchanged. Their
values are the restrictions of \(\mathcal M^jU\),
so \(\mathsf d_1=\xi_h^\diamond(U)\).

\emph{Arrival.}
The pinned-response result gives a real gauge
\(u\), with
\(\mathsf U^\sharp=(\mathsf U^{\rm pl})^u\),
restores the shorter datum \(\mathsf d_1\), and
leaves every noninner bond of \(\Omega^\diamond_0\)
unchanged. The prefix support inclusion places
\(U\) in \(\mathsf{Sup}(h)\), so \cref{II-10-20}
also supplies the old orbit with datum
\(\xi_h^\diamond(U)\).
By \cref{II-93-38}, the plain response is a gauge
transform of an element of this orbit. Pinning
preserves criticality and regularity and restores
the exact constraints. Its restriction is therefore
itself an element of the old regular critical orbit.
The membership statement \cref{II-99-39}, applied
to the shorter history, puts its plain triples
in the half-radius output spaces and in the
half-radius fixed cores, hence in the required
full spaces.

\emph{The second variation.}
We give the regional calculation because its row
pairing differs from the one on a constant tower.
Let \(J=J^\circ(U_+)\) be the actual regional action
current, and let \(\mathsf J\) inject the level-\(k\)
rows into the shorter row space. Write
\(\mathsf m\) for the row-spacing multiplication.
On the injected rows the spacing and the multiplier
\(\mathsf a\) are one:
\[
       \mathsf m^{-1}\mathsf J=\mathsf J,\qquad
                    \mathsf J^{\mathsf T}\mathsf a\mathsf J=I.
\]
The regional response at small arbitrary real row
data \(\beta\), followed by its restoring gauge,
has action \(\varphi(\beta)\). Its equations are
the full-parameter response equations
\cref{I-91-95}, with
\(\mathcal Q(\widehat A(\beta))=\mathsf J\beta\).
The strict regularity bounds allow a smaller
neighbourhood on which this restored response
is regular. Whenever the regular orbit is defined
it contains this response, by the criticality
argument of \cref{II-93-38}.

Differentiate these equations at zero. With
\(a=\widehat A'(0)\nu\) and
\(\mu=\mathsf m^{-3}\widehat\omega'(0)\nu\), the
saddle equations are
\[
 (\Delta_1+D\mathsf RD^{\mathsf T})a-Q^{\mathsf T}\mu=0,
 \qquad Qa=\mathsf J\nu .
\]
The inverse identities already used to construct
the response give
\[
 a=\mathsf H_1\mathsf J\nu,\qquad
 \mu=[(Q\mathsf G_1Q^{\mathsf T})^{-1}-\mathsf a]
                                                   \mathsf J\nu .
\]
The chain rule and the differentiated constraint
give the envelope identity
\[
 D\varphi(\beta)\nu
     =(\mathsf m^{-3}\widehat\omega(\beta),
                                      \mathsf J\nu)_{\mathbf B}.
\]
At zero,
\(\mathsf m^{-3}\widehat\omega(0)=\mathsf H_1^{\mathsf T}J\).
Consequently
\begin{align*}
 D\varphi(0)\nu&=(\mathsf H_1\mathsf J\nu,J)_{\eta_k},\\
 D^2\varphi(0)(\nu,\nu)
   &=(\nu,[\mathsf J^{\mathsf T}
                (Q\mathsf G_1Q^{\mathsf T})^{-1}\mathsf J-I]\nu)_1.
\end{align*}

A curve in the coarse fibre lifts to the fixed
longer constraint. Criticality of \(U_+\) thus
annihilates \(D\varphi(0)\) on \(\ker Q_p\).
For a fibre curve \(e^{i\beta(s)}p\) with velocity
\(v\), differentiating the averaging constraint
twice gives
\[
 \ddot\beta(0)=\kappa-2\mathsf D_p^{(2)}(v,v),
                                \qquad \kappa\in\ker Q_p .
\]
Inserting this into the second derivative of
\(\varphi(\beta(s))\) removes the \(\kappa\) term.
The remaining correction is
\(-2(\mathsf J\mathsf D_p^{(2)}(v,v),
                         \mathsf H_1^{\mathsf T}J)_{\mathbf B}\).
Together with the preceding quadratic form this
is exactly \(\langle v,\Delta_k^\sharp(U_+,J)v\rangle\)
from \eqref{II-23-21}. The completed chart has
velocity \(v=g_k\mathsf C_pB\). Therefore
\[
 \mathbb F'(0)=0,\qquad
 x_k\mathbb F''(0)
      =\langle\mathsf C_pB,
                      \Delta_k^\sharp(U_+,J)\mathsf C_pB\rangle .
\]

It remains to check the current in the actual
triple. The field \(J'=L^{-3}\mathbf J_+\) agrees
with \(J\) at every bond touching \(\Omega_1\).
Their difference is supported on bare determining
rows. The proof of \cref{II-93-23} shows that
\(\mathsf G_1,\mathsf H_1\) are unchanged by this
difference, and that its transpose image has
only bare-row components. The nonlinear averaging
remainder vanishes on those rows. Thus the last
term of \eqref{II-23-21} is unchanged as well.
Both forms have the same defined operators;
the graded Gaussian-domain result supplies their
admission. Finally the block contours in
\(\widehat{\mathsf G}_V\) use only the double-block
field where \(V=p\). Equation~\eqref{II-23-22}
therefore gives \eqref{II-113-39}.
\end{proof}

\subsection{The identity for the old exponent}
\label{sec:history-onshell-identity}

Let \(A_{k,h}\) and \(\Phi_h\) be the old exponent
and retained sum defined in
\eqref{II-99-19} and \cref{sec:formed-factors}.
Write \(\mathsf J_V^\diamond[\mathcal A]\) for
the logarithm of the regional pivot and Haar
Jacobian, and \(\mathbb G[\mathcal A]\) for the
block gauge-fixing action. These are the local
totals of the J and G constructions.

\begin{proposition}\label{II-112-46}
At every support point of \cref{II-113-38}, all
terms in the following identity are defined:
\begin{equation}\label{II-112-48}
 \begin{split}
 &A_{k,h}(\mathsf U^\sharp,\mathfrak q)
       +\Phi_h(\mathsf U^\sharp;\mathfrak q)
       -x_k\mathbb G[\mathcal A]
       +\mathsf J_V^\diamond[\mathcal A]\\
 &\quad=A_{k,h}(\mathsf U,\mathfrak q)
       +\Phi_h(\mathsf U;\mathfrak q)
       -\tfrac12(B,\mathsf A^\sharp_{h'}(P_+)B)_1\\
 &\qquad\quad+\sum_{Y\in\mathcal D_k}
       [\mathbb V^{\rm int}+\mathbb V'
                      +\mathbb V^{\rm B}+\mathbb V^\Phi]
                   (Y|h';(P_+,\mathcal A,\mathfrak q)|_Y).
 \end{split}
\end{equation}
At the first step the same formula holds with no
old regular, boundary or retained terms; the
level-zero action and completed fibre field give
its direct interpretation.
\end{proposition}

\begin{proof}
For an inner plaquette of \(\Omega^\diamond_0\),
the pinned and plain fields have the same trace,
and the plain field is gauge equivalent to the
critical field with datum \(\mathsf d_1\).
For a plaquette not inner to the collar all its
bond values are unchanged. To see the latter
assertion also for its inner bonds, take a corner
outside \(\Omega_1^\sim\). Every other corner
is within one bare spacing of it. Such a corner
cannot lie in \(\Omega_1\), since the collar has
width at least a cell side. Any inner bond of
the plaquette is therefore a level-zero
determining row, where \(\mathsf d_1\) is unchanged.
For every plaquette weight \(w\), it follows that
\begin{equation}\label{eq:history-weighted-action-change}
 \begin{split}
 A(w,\mathsf U^\sharp)-A(w,\mathsf U)
   =\sum_{p\ {\rm inner\ in}\ \Omega^\diamond_0}
        w(x(p))[a_p(U_{\mathbf B}(\mathsf d_1))-a_p(U_+)] .
 \end{split}
\end{equation}
In particular,
\(L^k[A(\mathsf U^\sharp)-A(\mathsf U)]
                             =\mathbb F(1)-\mathbb F(0)\).

The exact totals of J, G and A in
\cref{II-106-49}, with their Taylor subtractions,
and \eqref{II-113-39} yield
\[
 \begin{split}
 \sum_Y\mathbb V^{\rm int}(Y)
   ={}&\mathsf J_V^\diamond[\mathcal A]
       -x_k\mathbb G[\mathcal A]
       -x_k[\mathbb F(1)-\mathbb F(0)]\\
      &+\tfrac12(B,\mathsf A^\sharp_{h'}(P_+)B)_1 .
 \end{split}
\]
The counterterm identity \cref{II-107-57}, applied with
\(w_k=\beta_k(\cdot)-L^kx_k\), gives
\[
 \sum_Y\mathbb V^{\rm c}(Y)
                =-[A(w_k,\mathsf U^\sharp)-A(w_k,\mathsf U)].
\]
Here the historical counterterms are \(b_j\), so their
dimensionless weight is \(p_w=L^{-k}w_k\).
The regular-member identity \cref{II-108-27} gives
\[
 \sum_Y\mathbb V^{\rm E}(Y)
                  =\mathbf E_k(\mathsf U^\sharp)-\mathbf E_k(\mathsf U).
\]
Here one may first use the plain response; gauge
invariance of the regular members identifies it
with the pinned response. Adding the first two
action changes and using
\(\beta_k=w_k+L^kx_k\) supplies the full weighted
Wilson-action difference.

For the other origins, the pinned substitution and
\cref{II-109-31}, \eqref{II-111-29} give
\[
 \begin{split}
 \sum_Y\mathbb V^{\rm B}(Y)
    &=\mathbf E_k^{\rm bd}(\mathsf U^\sharp,\mathfrak q)
                            -\mathbf E_k^{\rm bd}(\mathsf U,\mathfrak q),\\
 \sum_Y\mathbb V^\Phi(Y)
    &=\Phi_h(\mathsf U^\sharp;\mathfrak q)
                            -\Phi_h(\mathsf U;\mathfrak q).
 \end{split}
\]
The second equality is the old-retained identity
\eqref{II-111-29}; its fixed stored data are the
same on both sides. The membership and potential
bounds prove definedness before these sums are
taken. Each sum is finite.
Substitution of \eqref{II-99-19} now gives
\eqref{II-112-48}; the vacuum constant cancels.

At \(k=0\) the shifted field is the completed
one-step chart field on \(\mathfrak X_0\) and the
original field elsewhere. The direct plaquette
Taylor expansion and gauge-fixing quadratic
form give the same J, G and A total. There are
no E, c, B or \(\Phi\) members, so the same
calculation proves the first-step assertion.
\end{proof}


\section{Cluster sums for a history step}
\label{sec:history-cluster}

We use the Gaussian data and term bounds of
Section~\ref{sec:conditional-mean}. The integration set may
be a proper subset of the non-pivot bonds. Accordingly,
the guard in the finite differences remains
$\operatorname{Gd}(Z_\circ)$; the domain assigned to a
term remains $Z_\circ^\sharp\cup\bigcup\mathcal K$.

\subsection{Factorization and the whole integral}

Call $(\iota,\mathcal K)$ rooted if every wall component
of $\operatorname{Gd}(Z_\circ(\iota))\cup\mathcal K$
contains a cube of the guard. Write
$\mathcal I^{\rm gd}(Z)$ for the rooted pairs with domain
$Z$. An empty guard permits only $\mathcal K=\emptyset$.
For a coarse wall-connected region $\gamma\subset\mathfrak R$,
set
\[
 \mathsf z(\gamma)=
   \sum_{(\iota,\mathcal K)\in\mathcal I^{\rm gd}(\gamma)}
                            T(\iota,\mathcal K).
\]
These sums are finite. Define $E(X)$ by the local cluster
series \eqref{II-18-78} with these activities. Its refined
series expands every activity into its individual terms.
For a set $\mathsf{Mk}$ of vertices and integration bonds,
$E^{\mathsf{Mk}}(X)$ is the subseries containing at least
one vertex or failed bond in $\mathsf{Mk}$.

\begin{proposition}\label{II-116-68}
Terms that are not rooted vanish. If $Z$ has wall
components $\gamma_1,\ldots,\gamma_n$, restriction to
these components is a bijection
\[
 \mathcal I^{\rm gd}(Z)
       \longrightarrow\prod_{i=1}^n\mathcal I^{\rm gd}(\gamma_i),
\]
and
\begin{equation}\label{II-116-69}
 T(\iota,\mathcal K)
   =\prod_iT(\iota|_{\gamma_i},\mathcal K\cap\gamma_i).
\end{equation}
Consequently the sum of terms with domain $Z$ is
$\prod_i\mathsf z(\gamma_i)$.

For fixed $\iota$, put $Z_\circ=Z_\circ(\iota)$.
If $Z_\circ^\sharp\subset Z\subset\mathfrak R$, let
$\mathcal K_0$ be the cubes of $Z\setminus Z_\circ^\sharp$
and $\mathsf N'$ the cubes of
$Z_\circ^\sharp\setminus\operatorname{Gd}(Z_\circ)$.
Then
\begin{equation}\label{II-116-71}
 \sum_{\mathcal K:\,X(\iota,\mathcal K)=Z}T(\iota,\mathcal K)
                    =f_\iota[\mathcal K_0\mid\mathsf N'].
\end{equation}
In particular, enlarging the guard within $Z_\circ^\sharp$
without removing $(Z_\circ^{\rm I})^\sharp$ does not change
the activities or the cluster sums, whenever the enlarged
guards respect restriction to wall components.
\end{proposition}

\begin{proof}
Put $A=Z_\circ^{\rm I}$ and $G=\operatorname{Gd}(Z_\circ)$.
For $W\subset\mathcal K$, the Gaussian parameter in
$f_\iota(\mathbf1_W)$ is the indicator of
$(W\cup G)\setminus A^\sharp$.
Suppose a wall component $\mathsf c$ of $\mathcal K\cup G$
contains a cube of $\mathcal K$ and no cube of $A^\sharp$.
By the component property of the Gaussian structure,
deleting $\mathsf c$ leaves the covariance, connecting
operator and mean on $A$ unchanged. Thus
$f_\iota(\mathbf1_W)=f_\iota(\mathbf1_{W\setminus\mathsf c})$.
Pairing the subsets in the alternating difference along
any cube of $\mathsf c\cap\mathcal K$ gives zero.
This applies in particular to every pair that is not rooted.

Every enlarged vertex domain $Y^\sharp$ and bond domain
$b^\sharp$ is wall-connected. If it lies in $Z$, it lies
in exactly one $\gamma_i$. The potential and test factors
therefore split between the components. Their integration
bonds, and the auxiliary bonds used by their connecting
operators, are disjoint. The Gaussian covariance is a
direct sum, its determinant branch is the product of the
component branches, and its mean splits by restriction.
Absolute convergence from the preceding section permits
Fubini. The resulting integral is a product of component
integrals; taking the finite difference separately in the
disjoint sets $\mathcal K\cap\gamma_i$ proves
\eqref{II-116-69}. Restriction preserves rootedness.
Conversely, gluing rooted component indices preserves
the exclusion of failed bonds from tested vertex domains
and recovers the original index. Every component is
nonvoid: a component with only parameter cubes would
contradict rootedness. This proves the bijection.

The sets with $X(\iota,\mathcal K)=Z$ are exactly
$\mathcal K=\mathcal K_0\sqcup\mathcal K'$ with
$\mathcal K'\subset\mathsf N'$. The elementary finite
difference identity
\[
 \sum_{\mathcal K'\subset\mathsf N'}
 f[\mathcal K_0\cup\mathcal K'\mid\emptyset]
                 =f[\mathcal K_0\mid\mathsf N']
\]
proves \eqref{II-116-71}. Expanding its right side, the
parameter at $W\subset\mathcal K_0$ is the indicator of
$(W\cup Z_\circ^\sharp)\setminus A^\sharp$.
This expression does not use the guard. Nonrooted terms
may be added since they vanish, so the activity is
unchanged. The assertion for cluster sums then follows
term by term.
\end{proof}

For $\mathfrak R=T_\ell$, call the Gaussian structure
complete if there are a symmetric invertible whole
covariance $\mathbf C_*$ and a whole mean $\mathbf m_*$,
and a real positive reference $\mathbf C_{*{\rm r}}$, with
\[
 \gamma_{\rm r}\le\mathbf C_{*{\rm r}}^{-1}\le c_{\rm r},
 \qquad
 \|\mathbf C_*^{-1}-\mathbf C_{*{\rm r}}^{-1}\|
                  \le\hat\varepsilon,
\]
such that, at $s\equiv1$ and with $\chi$ the restriction
to the integration bonds touching $Z_\circ$,
\begin{equation}\label{II-114-13}
 \mathbf C+\mathbf C\boldsymbol\Gamma
        \boldsymbol\Gamma^{\mathsf T}\mathbf C
       =\chi\mathbf C_*\chi^{\mathsf T},
 \qquad \mathbf m=\chi\mathbf m_*.
\end{equation}
For complete data define
\[
 d\nu(B)=\det{}^{1/2}(\mathbf C_*^{-1}/2\pi)
 e^{-(B-\mathbf m_*,\,\mathbf C_*^{-1}(B-\mathbf m_*))_1/2}\,dB.
\]
This is a finite complex measure of mass one.
Indeed $\hat\varepsilon\le\gamma_{\rm r}/2$ implies
$\operatorname{Re}\mathbf C_*^{-1}\ge\gamma_{\rm r}/2$.

\begin{proposition}\label{prop:history-whole-integral}
For complete data and bounded Borel $F$,
\begin{equation}\label{II-116-74}
 \mathbb J^{\mathbf m}(Z_\circ;\mathbf1)[F]
                       =\int F(\chi B)\,d\nu(B).
\end{equation}
Let $\mathsf{Vert}$ be the potential indices $(Y,o)$
for which $Y$ contains a fine cube meeting $\Omega_t$.
The whole tested integral, with its potential sum over
these vertices, satisfies
\begin{equation}\label{II-116-75}
 \Xi[\underline{\mathbb V}]
 :=\int d\nu(B)\,\mathsf{in}_{\mathfrak X_{\rm I}}(B)
              \exp\sum_{(Y,o)\in\mathsf{Vert}}\mathbb V^o(Y;B)
 =\sum_{\mathcal P'}\prod_{\gamma\in\mathcal P'}\mathsf z(\gamma),
\end{equation}
where $\mathcal P'$ ranges over pairwise compatible
coarse domains, including the empty family.
\end{proposition}

\begin{proof}
The convergent Gaussian integral with the linear factor
$e^{i(f,B)_1}$ gives
\begin{equation}\label{II-116-73}
 \mathbb J^{\mathbf m}[e^{i(f,B)_1}]
 =\exp\left\{i(f,\mathbf m)_1
   -\tfrac12(f,(\mathbf C+\mathbf C\boldsymbol\Gamma
                      \boldsymbol\Gamma^{\mathsf T}\mathbf C)f)_1\right\}.
\end{equation}
One first integrates $B$ and then the real auxiliary
Gaussian; Fubini is justified by the joint Gaussian
majorant. The whole measure has Fourier transform
$\exp\{i(f,\mathbf m_*)_1-(f,\mathbf C_*f)_1/2\}$.
By \eqref{II-114-13}, these transforms agree after
restriction. Uniqueness of the Fourier transform for
finite complex Borel measures gives \eqref{II-116-74},
including the case of an empty integration set.

Expand each exponential into $1+(e^{\mathbb V^o}-1)$.
For every selected vertex retain the small-field tests
on its domain; on the remaining bonds expand
$\mathsf{in}=1-\mathsf{out}$. The resulting nonconstant
terms are precisely the $F_\iota$ of \eqref{II-115-95}.
They are bounded Borel functions, with zero continuation
where a potential is undefined. The expansion is finite.
For each index, \eqref{II-116-74} and the sum of all finite
differences give
\[
 \int F_\iota(\chi B)\,d\nu(B)
      =f_\iota(\mathbf1)
      =\sum_{\mathcal K}T(\iota,\mathcal K).
\]
Only rooted pairs contribute. Group them by their domain,
use \eqref{II-116-69}, and identify the wall components
of a coarse region with a compatible polymer family.
This proves \eqref{II-116-75}.
\end{proof}

\subsection{The majorant and its decay}

Let $w^o(Y)\ge e\,u^o(Y)$, where
$|\mathbb V^o(Y)|\le u^o(Y)\le1$ on its complex range,
and assume
\begin{equation}\label{II-118-18}
 \sum_o w^o(Y)\le\varepsilon_{\mathsf w}e^{-r_{\mathsf w}d_k(Y)}.
\end{equation}
All parameters $r_H,r_{\mathsf w},\kappa''$ are nonnegative.
Choose
\[
 \mathsf C_{\rm cs}\ge
 \max\{C'_{\rm col},24C_{\rm E}^2e^{24}\}.
\]
Here $C'_{\rm col}=13/9$ is the coarsening constant in
\eqref{II-18-24} of \cref{II-18-23}, and $c_{\rm E},C_{\rm E}$
are the entropy constants used in \cref{II-18-86}. Set
\begin{equation}\label{II-118-5}
\begin{aligned}
 \alpha_{\mathsf w}
   &=L^3\varepsilon_{\mathsf w}+(LM)^3e^{-\gamma_P},\\
 \varepsilon_H
   &=\mathsf C_{\rm cs}e^{\mathsf C_{\rm cs}(c_{\rm P}+r_H)}
                                               \alpha_{\mathsf w},\\
 r_{\rm pol}&=r_H-\mathsf C_{\rm cs}(c_{\rm P}+1),
 &C_{\rm pol}&=\mathsf C_{\rm cs}e^{\mathsf C_{\rm cs}\kappa''},\\
 K_{\rm mk}&=\mathsf C_{\rm cs}
                 e^{\mathsf C_{\rm cs}(c_{\rm P}+\kappa'')},\\
 K_{\rm or}&=\mathsf C_{\rm cs}(LM)^3
                 e^{\mathsf C_{\rm cs}(c_{\rm P}+r_H+\kappa'')}.
\end{aligned}
\end{equation}
The hypotheses for the cluster bounds are
\begin{equation}\label{eq:history-cluster-rates}
\begin{gathered}
 (\mathsf C_{\rm cs}/L)r_H+c_{\rm E}\le r_{\mathsf w},
 \qquad e^{-(\bar\kappa-1)+\mathsf C_{\rm cs}r_H}
                                      \le\mathsf C_{\rm cs}^{-1},\\
 r_{\rm pol}\ge\kappa''+c_{\rm E}+\mathsf C_{\rm cs},
 \qquad C_{\rm pol}\varepsilon_H\le\tfrac12 .
\end{gathered}
\end{equation}
Let $\bar T$ denote the majorant in \eqref{II-117-23}.
Write $\bar{\mathsf z}$, $\bar E$ and
$\bar E^{\mathsf{Mk}}$ for the corresponding nonnegative
activity and refined cluster sums.

\begin{theorem}\label{II-118-20}
Under \eqref{eq:history-cluster-rates}, the refined
cluster series converge absolutely and
\begin{align}
 |\mathsf z(\gamma)|&\le\bar{\mathsf z}(\gamma)
                    \le\varepsilon_He^{-r_{\rm pol}d_t(\gamma)},
                                                    \label{II-118-22}\\
 |E(X)|&\le\bar E(X)
                    \le C_{\rm pol}\varepsilon_He^{-\kappa''d_t(X)}.
                                                    \label{II-118-23}
\end{align}
Both sums vanish if no coarse cube of their domain meets
$\Omega_t$. For complete whole-torus data,
\begin{equation}\label{II-118-24}
 \Xi[\underline{\mathbb V}]
       =\exp\sum_{X\in\mathcal D_t}E(X)\ne0.
\end{equation}

Suppose $\mathsf{Mk}$ is closed under reversal on bonds.
Put
\begin{equation}\label{II-118-14}
\begin{split}
 \eta^{\mathsf{Mk}}=\sup_\Delta\bigg[
 &\sum_{\substack{(Y,o)\in\mathsf{Mk}\\\Delta\subset Y^\sharp}}
             w^o(Y)e^{r_H(d_t(Y^\sharp)+4)}\\
 &+\#\{b\in\mathsf{Mk}^+:\Delta\subset b^\sharp\}
                                     e^{-\gamma_P+4r_H}\bigg].
\end{split}
\end{equation}
Then
\begin{equation}\label{II-118-25}
 |E^{\mathsf{Mk}}(X)|\le\bar E^{\mathsf{Mk}}(X)
                \le K_{\rm mk}\eta^{\mathsf{Mk}}e^{-\kappa''d_t(X)}.
\end{equation}
For an admitted insertion with
$|\mathbb W(Y)|\le\varepsilon'e^{-r_{\mathsf w}d_k(Y)}$,
where $\varepsilon'\ge0$ need not be small,
\begin{equation}\label{II-118-29}
 |E'[\underline{\mathbb V};\mathbb W](X)|
                \le\bar E'(X)\le K_{\rm or}\varepsilon'
                                                    e^{-\kappa''d_t(X)}.
\end{equation}
The insertion sum is linear in $\mathbb W$.
\end{theorem}

\begin{proof}
We give the counting estimates, since the guard changes
the index sets from those of Section~\ref{sec:cluster-sums}.
The pieces of a domain are $Y^\sharp$ for its vertices,
$b^\sharp$ for its positive failed bonds, and the cubes
of $\mathcal K$. Coarsening commutes with unions.
Every piece is wall-connected, and
\[
 d_t(Y^\sharp)\le(C'_{\rm col}/L)d_k(Y),
 \qquad d_t(b^\sharp)=d_t(\Delta)=0.
\]
The connected-union estimate \eqref{II-18-31} therefore gives
\begin{equation}\label{II-118-47}
 d_t(\gamma)\le
 \sum_{(Y,o)\in\mathbf D}(d_t(Y^\sharp)+4)
             +4\#\mathfrak e^++4\#\mathcal K-4.
\end{equation}
For an insertion, add $d_t(Y_1^\sharp)+4$ on the right.
In particular
\begin{equation}\label{II-118-21}
 d_t(\gamma)\le(\mathsf C_{\rm cs}/L)
                 \sum_{(Y,o)\in\mathbf D}d_k(Y)
       +\mathsf C_{\rm cs}(\#\mathbf D+\#\mathfrak e^++\#\mathcal K).
\end{equation}

Introduce the corrected weights and constants
\[
\begin{gathered}
 \widehat w^o(Y)=w^o(Y)e^{r_H(d_t(Y^\sharp)+4)},\qquad
 \widehat w_{\rm bd}=e^{-\gamma_P+4r_H},\qquad
 q_\flat=e^{-(\bar\kappa-1)+4r_H},\\
 \beta_\flat=12C_{\rm E}e^{4r_H}\alpha_{\mathsf w},\qquad
 \varepsilon_{\rm ro}=12C_{\rm E}e^{8c_{\rm P}+16}\alpha_{\mathsf w},
 \qquad r_\flat=r_H-8c_{\rm P}-16,\\
 \mathsf C_{\rm gas}=C_{\rm E}e^{8+4\kappa''}.
\end{gathered}
\]
The four rate conditions imply
\begin{equation}\label{eq:history-counting-reserve}
\begin{gathered}
 (C'_{\rm col}/L)r_H+c_{\rm E}\le r_{\mathsf w},\quad
 q_\flat\le\tfrac12,\quad\beta_\flat\le\tfrac12,\\
 r_\flat\ge\kappa''+c_{\rm E}+8,\qquad
 14\mathsf C_{\rm gas}\varepsilon_{\rm ro}\le1 .
\end{gathered}
\end{equation}
These inequalities remain valid if the amplitude is
enlarged so that the new $\alpha_{\mathsf w}$ is at most
twice the original one. To check the numerical reserve,
use $\mathsf C_{\rm cs}\ge24C_{\rm E}^2e^{24}$ and
\[
 \mathsf C_{\rm cs}^2
 e^{\mathsf C_{\rm cs}(\kappa''+c_{\rm P}+r_H)}
                \alpha_{\mathsf w}\le\tfrac12.
\]
This bounds $\beta_\flat$ and
$14\mathsf C_{\rm gas}\varepsilon_{\rm ro}$, including
that factor two. The second rate condition gives
$q_\flat\le\mathsf C_{\rm cs}^{-1}$; and
$r_\flat-r_{\rm pol}
=(\mathsf C_{\rm cs}-8)c_{\rm P}+\mathsf C_{\rm cs}-16\ge0$.

For any $v(Y)\le\varepsilon^\natural
e^{-r_{\mathsf w}d_k(Y)}$, coarsening and the entropy
bound give
\begin{equation}\label{eq:history-pinned-weight}
 \sum_{\Delta\subset Y^\sharp}v(Y)e^{r_Hd_t(Y^\sharp)}
                   \le2C_{\rm E}L^3\varepsilon^\natural .
\end{equation}
Indeed the exponent is at most $-c_{\rm E}d_k(Y)$;
at most $(L+2)^3$ fine cubes can locate the enlarged
domain at $\Delta$, and $(L+2)^3\le2L^3$.
The number of positive integration bonds with
$\Delta\subset b^\sharp$ is at most $6((L+2)M)^3$.
Thus the corrected vertex and failed-bond weights pinned
at $\Delta$ sum to at most $\beta_\flat$. Those contained
in $\gamma$ sum to at most $\beta_\flat\#\gamma$.

Multiply the product in $\bar T$ by
$e^{r_Hd_t(\gamma)}$ and use \eqref{II-118-47}.
The result is at most
$e^{c_{\rm P}\#\gamma-4r_H}$ times the product of corrected
weights. We may now sum over all subsets of the vertices,
failed bonds and cubes contained in $\gamma$.
The first two subset sums are bounded by
$e^{\beta_\flat\#\gamma}$ and the last by
$e^{q_\flat\#\gamma}$. The vertex/failed-bond pair is
nonempty, so the first exponential may be replaced by
$e^{\beta_\flat\#\gamma}-1$. Using
$e^u-1\le ue^u$, $\#\gamma\le e^{\#\gamma}$ and
$\#\gamma\le8(d_t(\gamma)+1)$ yields
\begin{equation}\label{eq:history-activity-sharp}
 \bar{\mathsf z}(\gamma)
                \le\varepsilon_{\rm ro}e^{-r_\flat d_t(\gamma)}.
\end{equation}
The same free subset sums, with one vertex or bond
required to be marked, give
\[
 \bar{\mathsf z}^{\mathsf{Mk}}(\gamma)
       \le e^{8c_{\rm P}+16}\eta^{\mathsf{Mk}}
                                      e^{-r_\flat d_t(\gamma)}.
\]
Here the sum of marked weights in $\gamma$ is at most
$\#\gamma\,\eta^{\mathsf{Mk}}$. For an insertion, the
vertex/failed-bond pair may be empty. Its extra domain
weight is summed by \eqref{eq:history-pinned-weight};
hence
\[
 \bar{\mathsf z}'(\gamma)\le
 2C_{\rm E}L^3e^{8c_{\rm P}+16}\varepsilon'
                                      e^{-r_\flat d_t(\gamma)}.
\]

Theorem~\ref{II-18-86}, with
$r=r_\flat$, $r''=\kappa''$ and
$\varepsilon=\varepsilon_{\rm ro}$, now bounds the
refined cluster sum by
$\mathsf C_{\rm gas}\varepsilon_{\rm ro}e^{-\kappa''d_t(X)}$.
Finite expansion of its activities and Tonelli for the
majorants justify the refined form of this statement.
For one marked factor use the following consequence of
the same theorem. If
$h_0\le\varepsilon_{\rm ro}e^{-r_\flat d_t}$ and
$h_1\le\varepsilon^\natural e^{-r_\flat d_t}$, the
cluster sum with exactly one $h_1$ and all other factors
$h_0$ is at most
$2\mathsf C_{\rm gas}\varepsilon^\natural e^{-\kappa''d_t(X)}$.
For $\varepsilon^\natural>0$, apply the theorem to
$h_0+\lambda h_1$ with
$\lambda=(14\mathsf C_{\rm gas}\varepsilon^\natural)^{-1}$
and retain its nonnegative terms of degree one in
$\lambda$. Their sum times $\lambda$ is at most
$e^{-\kappa''d_t(X)}/7$. The zero case is immediate.
The indicator that a cluster has a mark is bounded by
its number of marked factors. The insertion series has
$1/(n-1)!$ in place of $1/n!$; symmetrizing its singled
out factor gives the same one-factor estimate.

Finally, the definitions and the choice of
$\mathsf C_{\rm cs}$ give
\[
\begin{gathered}
 \varepsilon_{\rm ro}\le\varepsilon_H,\quad
 r_\flat\ge r_{\rm pol},\quad
 \mathsf C_{\rm gas}\le C_{\rm pol},\\
 2\mathsf C_{\rm gas}e^{8c_{\rm P}+16}\le K_{\rm mk},\\
 12C_{\rm E}\mathsf C_{\rm gas}e^{8c_{\rm P}+16}(LM)^3
                                                    \le K_{\rm or},\\
 4C_{\rm E}\mathsf C_{\rm gas}e^{8c_{\rm P}+16+4r_H}L^3
                                                    \le K_{\rm or}.
\end{gathered}
\]
These prove all the stated bounds. Each nonvoid index
has a vertex containing a source cube or a failed
integration bond touching $\Lambda_t\subset\Omega_t$.
Its domain therefore has a coarse cube meeting
$\Omega_t$, which proves the vanishing assertions.
For complete data, \eqref{II-116-75} and
\cref{II-18-86} prove \eqref{II-118-24}.
Linearity of the insertion follows first for each
integrable term and then for its absolutely convergent
series.
\end{proof}

\subsection{Regular and boundary contributions}

Split the origins into disjoint sets
$\mathsf{Or}_{\mathsf v}$ and $\mathsf{Or}_\partial$,
with respective weight bounds
$\varepsilon_{\mathsf v}e^{-r_{\mathsf w}d_k(Y)}$
and $\varepsilon_\partial e^{-r_{\mathsf w}d_k(Y)}$.
Let $E^{\rm b}$ be the subseries with a boundary origin
and put $E^{\rm r}=E-E^{\rm b}$.

\begin{corollary}\label{cor:history-regular-boundary}
\begin{equation}\label{II-118-27}
\begin{aligned}
 E^{\rm r}(X)&=E[\underline{\mathbb V}|_{\mathsf{Or}_{\mathsf v}}](X),\\
 |E^{\rm r}(X)|&\le K_{\rm or}
             (\varepsilon_{\mathsf v}+e^{-\gamma_P})e^{-\kappa''d_t(X)},\\
 |E^{\rm b}(X)|&\le K_{\rm or}\varepsilon_\partial e^{-\kappa''d_t(X)}.
\end{aligned}
\end{equation}
\end{corollary}

\begin{proof}
Removing the boundary origins removes exactly the
refined indices containing one of them; all remaining
terms, guards and domains are unchanged. Absolute
convergence proves the identity. Apply
\eqref{eq:history-activity-sharp} to the restricted
system with amplitude
$\min(\varepsilon_{\mathsf v},\varepsilon_{\mathsf w})$.
Its cluster bound is at most
\[
 12C_{\rm E}\mathsf C_{\rm gas}e^{8c_{\rm P}+16}(LM)^3
       (\varepsilon_{\mathsf v}+e^{-\gamma_P})e^{-\kappa''d_t(X)}.
\]
For the boundary part,
\eqref{eq:history-pinned-weight} gives
$\eta^\partial\le2C_{\rm E}L^3e^{4r_H}\varepsilon_\partial$.
Use the sharper marked bound in the preceding proof,
$2\mathsf C_{\rm gas}e^{8c_{\rm P}+16}\eta^\partial$,
and the last two numerical comparisons there.
\end{proof}

All these series have the same majorant when their
Gaussian entries and potentials vary with fixed bounds.
For jointly holomorphic data, the fixed-domain Gaussian
integral and its common integrable majorant prove
holomorphy of each finite-difference term; the convergent
majorant series then proves holomorphy of every sum and
every marked subseries. For jointly Borel data the
parameter integrals are Borel by measurable integration,
and the cluster sums are pointwise limits of Borel
partial sums. Orthogonal changes of the real integration
and auxiliary variables preserve the integrals. Complex
conjugation commutes with them and with the fixed
determinant branch; real data therefore give real sums.

Locality requires agreement of the thresholds on the
integration bonds internal to $X$, of every potential
whose enlarged domain lies in $X$, and of the Gaussian
entries at the $0$--$1$ decoupling parameters supported
in $X$, including the connecting-operator columns
touching $X$. Under these hypotheses, all refined
indices with union $X$, their guards and their terms
agree: auxiliary columns outside $X$ vanish and their
Gaussian integral is one. Thus the sums and the marked
subseries agree. The same statement holds for the
majorants when the local weights agree.

\subsection{Insertions and the cutoff terms}

An insertion supported on domains containing a source
cube can be recovered by adjoining it as an extra origin.
Choose $\epsilon_0>0$ such that
\[
 \epsilon_0\max_Y\sup|\mathbb W(Y)|\le1,\qquad
 e\epsilon_0\varepsilon'\le M^3e^{-\gamma_P}.
\]
The weights of $\underline{\mathbb V}\cup\{\epsilon\mathbb W\}$
then have $\alpha_\flat\le2\alpha_{\mathsf w}$ for
$|\epsilon|<\epsilon_0$. The reserve
\eqref{eq:history-counting-reserve} gives one convergent
majorant on this disc. In a term with $p$ vertices of the
new origin, the product of the new factors is $O(\epsilon^p)$.
For $p=1$ its derivative is the corresponding insertion
term, with the same domain, guard and tests. This follows
also in the integral norm, since
\[
 \left|\epsilon^{-1}(e^{\epsilon w}-1)-w\right|
                                      \le e|\epsilon|\,|w|^2
                    \quad (|\epsilon w|\le1).
\]
Cauchy's formula on a smaller circle permits termwise
differentiation of the uniformly majorized series.
Exactly one new vertex survives, so symmetrizing its
polymer position proves
\begin{equation}\label{II-118-30}
 \left.\frac{d}{d\epsilon}
 E[\underline{\mathbb V}\cup\{\epsilon\mathbb W\}](X)
 \right|_{\epsilon=0}
             =E'[\underline{\mathbb V};\mathbb W](X).
\end{equation}
The support condition ensures that every nonzero
insertion index corresponds to a permissible new vertex.

For complete data, an arbitrary admitted insertion
satisfies
\begin{equation}\label{II-118-31}
\begin{split}
 \sum_XE'[\underline{\mathbb V};\mathbb W](X)
 ={}&\Xi[\underline{\mathbb V}]^{-1}\int d\nu(B)\,
       \mathsf{in}_{\mathfrak X_{\rm I}}(B)\\
 &\hspace{8mm}\times
        e^{\sum_{(Y,o)\in\mathsf{Vert}}\mathbb V^o(Y;B)}
                         \sum_{Y_1}\mathbb W(Y_1;B).
\end{split}
\end{equation}
To see this without a support assumption, perform the
finite expansion of \cref{prop:history-whole-integral}
with a distinguished domain $Y_1$. Exactly one wall
component carries the insertion. Thus its integral is
$\sum_\gamma\mathsf z'(\gamma)\Xi(\mathcal P_0(\gamma))$,
where $\mathcal P_0(\gamma)$ consists of polymers
compatible with $\gamma$. To identify this ratio, put
$h=\mathsf z$ and
$\bar h(\gamma)=(7C_{\rm E}e^8)^{-1}e^{-r_\flat d_t(\gamma)}$.
The preceding estimates give
$r_\flat\ge c_{\rm E}+8$ and
$\varepsilon_{\rm ro}\le(14C_{\rm gas})^{-1}
 \le(14C_{\rm E}e^8)^{-1}$.
Thus $\bar h$ satisfies the polymer convergence criterion,
and $h$ lies strictly inside its open polydisc.
On this polydisc the cluster sum $S$ is holomorphic.
Differentiating its uniformly convergent series and using
symmetry in the polymer positions gives the series below
for $\partial S/\partial h(\gamma)$.
On the other hand, splitting compatible families according
to whether they contain $\gamma$ gives
$\Xi=\Xi(\mathcal P\setminus\{\gamma\})
       +h(\gamma)\Xi(\mathcal P_0(\gamma))$.
Since $\gamma$ is incompatible with itself, both partition
functions on the right are independent of $h(\gamma)$.
Differentiating this equality and $\Xi=e^S\ne0$ proves
\[
 \frac{\Xi(\mathcal P_0(\gamma))}{\Xi}
 =\sum_{n\ge1}\frac1{(n-1)!}
   \sum_{\gamma_2,\ldots,\gamma_n}
      \rho^T(\gamma,\gamma_2,\ldots,\gamma_n)
                          \prod_{i\ge2}\mathsf z(\gamma_i).
\]
Its absolute sum is at most $e^{\#\gamma}$.
There are finitely many polymers, so multiplication by
$\mathsf z'(\gamma)$, summation over $\gamma$, and grouping
by the union are justified. This proves
\eqref{II-118-31}.

For the empty potential system, allow bond parameters
$0\le t_b\le1$ in the tests:
$\mathsf{out}^{\mathbf t}_b(B)
=\mathbf1\{t_b|B(b)|\ge\mathfrak b_b\}$.
For $t_b=0$ this is zero. These tests are the original
failed tests multiplied by a bounded Borel function of
the failed variables. Proposition~\ref{II-117-102}
therefore gives the same term majorants for every
$\mathbf t$.

\begin{proposition}\label{prop:history-cutoff-members}
The empty-system cluster sum admits the bond
decomposition
\begin{equation}\label{II-118-33}
 E[\emptyset](X)=\sum_{b\in\mathfrak X_{\rm I}^+}
                  E^\chi[\emptyset](X,b),\qquad
 E^\chi[\emptyset](X,b)
 =\int_0^1
 \left.\partial_{t_b}E[\emptyset;\mathbf t](X)
                  \right|_{\mathbf t=(u,\ldots,u)}\,du .
\end{equation}
The derivative exists for every $0<u<1$, is integrable,
and
\begin{equation}\label{II-118-32}
 |E^\chi[\emptyset](X,b)|\le
               \bar E^{\{b,\bar b\}}(X).
\end{equation}
This member vanishes unless $b^\sharp\subset X$.
\end{proposition}

\begin{proof}
For $\iota=(\emptyset,\mathfrak e)$, integrate the
Gaussian density first in all variables except those
on $\mathfrak e$. The alternating sum over parameter
cubes, including the sign $(-1)^{\#\mathfrak e^+}$,
gives a continuous density $\sigma(x)$ such that
\[
 T(\iota,\mathcal K;\mathbf t)
    =\int\sigma(x)\prod_{b\in\mathfrak e^+}
                    \mathbf1\{t_b|x(b)|\ge\mathfrak b_b\}\,dx .
\]
It satisfies $|\sigma(x)|\le Ce^{-c\|x\|^2}$ for some
$c>0$. Indeed the real part of each joint Gaussian
exponent is a positive quadratic form plus a linear
term; absorbing that term leaves a Gaussian bound.
Its marginal has the same property, and the alternating
sum is finite. Moreover
\[
 \int|\sigma(x)|\prod_{b\in\mathfrak e^+}
          \mathbf1\{|x(b)|\ge\mathfrak b_b\}\,dx
                                  \le\bar T(\iota,\mathcal K).
\]
Apply the bounded-failed-factor conclusion of
\cref{II-117-102} with
$\psi=\overline\sigma/|\sigma|$ where $\sigma\ne0$,
and zero otherwise.

Use spherical coordinates separately on each
three-dimensional bond space, and set
$\zeta_b=\mathfrak b_b/|x(b)|$. The radial Jacobian is
$\mathfrak b_b^3\zeta_b^{-4}$. Consequently
\[
 T(\iota,\mathcal K;\mathbf t)
       =\int_{\prod_b[0,t_b]}\upsilon(\zeta)\,d\zeta,\qquad
 \int_{[0,1]^{\mathfrak e^+}}|\upsilon|\,d\zeta
                                          \le\bar T(\iota,\mathcal K).
\]
The Gaussian decay makes $\upsilon$ continuous on the
closed cube after setting it to zero on its zero faces:
each factor $\zeta_b^{-4}
e^{-c\mathfrak b_b^2/\zeta_b^2}$ tends to zero there.

In a product of $n$ terms take a separate copy
$\zeta_{i,b}$ for each factor. Partition this cube,
up to null sets, according to the coordinate that is
largest. On the part where coordinate $(i,b)$ is
largest, integration in that coordinate is precisely
the integral from zero to one of the corresponding
diagonal $t_b$ derivative. It follows that the sum
over bonds of the integrals of the absolute derivatives
is at most $\prod_i\bar T_i$, and the sum without
absolute values equals $\prod_iT_i(\mathbf1)$.
Only bonds occurring in some $\mathfrak e_i$ contribute.

For completeness, differentiation of the infinite
cluster series is justified before this integration.
There are finitely many Mayer indices with cubes.
For each one $\upsilon$ is bounded and $\bar T>0$, so
$\mathsf M=\max\|\upsilon\|_\infty/\bar T<\infty$.
If that index set is empty, set $\mathsf M=0$;
all the cluster sums in question are then empty.
The derivative of a product is bounded by
$\mathsf M n\prod_i\bar T_i$.
The one-mark argument in \cref{II-118-20} with
$h_1=h_0=\bar{\mathsf z}$ proves summability of the
refined majorant multiplied by $n$. Thus the derivative
series converges uniformly for $0< t_b<1$ with other
parameters fixed, and on the diagonal. Termwise
integration is also justified by the sharper
largest-coordinate estimate. Summing the identities
over refined indices proves \eqref{II-118-33};
restricting to indices containing $b$ proves
\eqref{II-118-32} and the support assertion.
\end{proof}


\subsection{The cluster terms in the interior}
\label{sec:history-interior-cluster}

The expansion in a domain contained in \(\Lambda_t\)
must reproduce the regular expansion. This requires
agreement of the local Gaussian operators and potentials,
including their domains, before the cluster series is
summed.

Fix a step \(t=k+1\), a history \(h'\) of depth \(t\),
and \(X\in\mathcal D_t\) with \(X\subset\Lambda_t(h')\).
Compare it with the empty history \(h_\emptyset\) at
the same level, with the same coupling prefix, regular
families, constants, and choices in the local constructions.
Use the geometric and analytic data of the actual history
step. In particular, \(m'<m+K-t\).
When \(k\ge1\), assume
\[
                 \check R_j\le L\check R_{j+1}
                         \qquad(1\le j<t).
\]
This is \eqref{II-21-25} through level \(t\).

Write \(\mathfrak X^{\rm 1l}\) for all nonpivot bonds at
level \(k\), and set
\[
 \mathfrak c=\mathfrak X^{{\rm 1l}\circ}\langle X\rangle,
 \qquad
 \mathfrak d=\mathfrak X^{\rm 1l}\langle X\rangle.
\]
Thus \(\mathfrak c\) has both endpoints in \(X\), whereas
\(\mathfrak d\) has at least one. For either history let
\(F_h(X,P,z)\) denote the pointed term formed from
\(E^\chi_h+E^\bullet_h\), with the tents in
\eqref{II-53-23}, at \(\tau=1\).
The regular potential system is
\(\underline{\mathbb V}^{\rm rg}
       =\{\mathbb V^{\rm int},\mathbb V'\}\).

For the sum identity below we require the following
analytic bounds at \((X,P)\). The empty-history operators
form the Gaussian structure of \cref{II-19-15} on \(X\),
and its constants satisfy the first three inequalities
of \eqref{II-19-57}. For every \(Y^\sharp\subset X\), the
four potentials J, G, A, and \(\prime\), evaluated at
\((P|_Y,g_kB)\), are holomorphic on
\(\mathbb D^{\rm f}(Y)\), with \(\mathsf C_{\mathcal A}\ge3\).
The J and \(\prime\) potentials vanish at zero, the G
and A potentials vanish through degree two, and
\begin{equation}\label{II-135-8}
 \|\mathbb V^{\rm J}\|^{\rm f}_{\star,\lambda_V}
 +\|\mathbb V^{\rm G}\|^{\rm f}_{\star,\lambda_V}
 +\|\mathbb V^{\rm A}\|^{\rm f}_{\star,\lambda_V}
       \le w_\star,\qquad
 \|\mathbb V'\|^{\rm f}_{\star,\lambda_V}\le w_\star .
\end{equation}
These are the hypotheses of the interpolation and cluster
estimates already proved. They hold on the regular analytic
cores used by the actual step, by
\cref{sec:actual-small-field-step,II-53-39} and the
common constants of \cref{sec:common-iteration-constants}.

\begin{proposition}\label{II-135-20}
On \(X\), the historical covariance and connecting
operator are the empty-history operators, as partial
maps of \(P\), and the conditional mean is zero wherever
those operators are defined. The pointed terms satisfy
\begin{equation}\label{eq:history-interior-pointed}
                  F_{h'}(X,P,z)=F_{\emptyset}(X,P,z)
\end{equation}
as partial maps for \(z\in X^{(t)}\).

If the analytic bounds above hold, all terms in
\begin{equation}\label{eq:history-interior-sum}
 \begin{split}
 E[\underline{\mathbb V}^{\rm rg}]_{h'}(X;P)
   &=\sum_{b\in\mathfrak X_{\rm I}^+}
                         (E^\chi_{h'}+E^\bullet_{h'})(X,b;P)\\
   &=\sum_{z\in X^{(t)}}F_\emptyset(X,P,z)
 \end{split}
\end{equation}
are defined and the equalities hold. If, in addition,
the old retained factors have their prescribed indices,
then the full history cluster sum has this same value.
At \(k=0\) there are no old retained factors and the
last conclusion holds without that condition.
\end{proposition}

\begin{proof}
We first compare the geometry. A nonpivot bond with an
endpoint in \(X\) is an integration bond, and no external
bond meets \(X\). For every region \(Y\subset X\),
\[
 \mathfrak X_{\rm I}\langle Y\rangle
      =\mathfrak X\langle Y\rangle
      =\mathfrak X^{\rm 1l}\langle Y\rangle,
 \qquad
 \mathfrak X_{\rm E}^\circ\langle X\rangle=\emptyset .
\]
The same equalities hold for bonds with both endpoints
in \(Y\). Every fine cube of \(X\) is a source cube.
Each contains a whole block \(B(y)\), including the
nonpivot bond \(\langle y,y+e_1\rangle\).
All bonds internal to \(X\) have threshold
\(\mathfrak b_k\).

If \(Z_\circ^\sharp\subset X\), every bond touching
\(Z_\circ\) lies inside \(Z_\circ^\sharp\). Consequently
\[
 Z_\circ^{\rm I}=Z_\circ,\qquad
 \operatorname{Gd}(Z_\circ)=Z_\circ^\sharp,\qquad
 \mathsf P^{\rm gd}(X,Z_\circ)=\mathsf P(X,Z_\circ).
\]
The guarded completion of the decoupling parameters
therefore makes no further change. For a Mayer index on a polymer
\(\gamma\subset X\), the condition
\(Z_\circ^\sharp\cup\bigcup\mathcal K=\gamma\)
forces \(\mathcal K\) to be exactly the cubes of
\(\gamma\) outside \(Z_\circ^\sharp\).
The resulting index is rooted: \(\gamma\) is connected
and the nonempty guard belongs to it. This applies
also to indices with an insertion, since their region
contains the insertion domain.

\emph{The operators.}
For a grain region \(Y\) with \(X(Y)\subset X\), every
member of its tower lies in \(Y\subset\Lambda_t\).
The meet with either collared sequence is therefore
the tower itself. The two deep-bond sets \(Y^-\)
also coincide and lie in \(\mathfrak c\). In the
alternating formula for a walk term, the compressed
operators are consequently the same:
\[
 \chi_{Y^-}\mathsf A^\sharp[\Omega_\bullet(Y)]
                      \chi_{Y^-}^{\mathsf T}.
\]
Their domains agree, since they use the same tower
and restricted pair. Summing the finite alternating
formula gives equal kernels on
\(\mathfrak c^+\times\mathfrak c^+\), zero off this
set, with simultaneous definedness.

This comparison extends to the local square root.
For a string whose completed parameter weight is
nonzero, its entire figure lies in \(X\). The sets of
bonds starting in its enlarged mesh cubes coincide
for the two histories and are subsets of \(\mathfrak c\).
The local compressed operators, their resolvents
\((\mu+\mathsf a^\square)^{-1}\), the cutoff
multiplications, and the commutator factors thus
coincide. Invertibility is one condition on one
compressed matrix. The string products have identical
kernels on \(\mathfrak c\) and vanish off it.
The integral
\[
       \frac1\pi\int_0^\infty
                     \mu^{-1/2}\mathbf g^\mu_{\boldsymbol\omega}
                                                       \,d\mu
\]
and the subsequent string series have the same
conditions of entrywise absolute convergence.
It follows that both \(\mathsf A^\sharp\) and
\(\mathbf r_{\rm I}\), with parameters supported
in \(X\), agree as partial maps after restriction
to \(\mathfrak c\), and vanish off that set.

For the Gaussian entries let
\(\mathfrak e=\mathfrak X^{\rm 1l}\langle Z_\circ\rangle
 \subset\mathfrak c\).
The completed parameter is supported in \(X\).
Both constructions read \(P\) at level \(k\), with
current \(L^{-3}\mathbf J\), and use the same matrices
in
\[
 \mathbf C=(\chi_{\mathfrak e}\mathsf A^\sharp
                                  \chi_{\mathfrak e}^{\mathsf T})^{-1},
 \qquad
 \boldsymbol\Gamma=
 \chi_{\mathfrak e}\mathsf A^\sharp
      (\mathbf1_{\mathfrak X_{\rm I}}-\mathbf1_{\mathfrak e})
                  \mathbf r_{\rm I}\chi_{\mathfrak d}^{\mathsf T}.
\]
On \(\mathfrak c\), both integration indicators are
one. This proves simultaneous definedness and equality
of \(\mathbf C,\boldsymbol\Gamma\). Moreover
\(\mathsf A^\sharp\chi_{\rm E}^{\mathsf T}=0\).
All remaining factors of the conditional-mean
operator are defined when \(\mathbf C,\boldsymbol\Gamma\)
are, so \(\mathbf M=\mathbf m=0\).

The conditioning integrals now coincide as well.
Indeed, with \(v=\boldsymbol\Gamma\mathbb X\),
symmetry of \(\mathbf C\) gives
\[
 (B+\mathbf Cv,\mathbf C^{-1}(B+\mathbf Cv))_1
  =(B,\mathbf C^{-1}B)_1+2(B,v)_1+(v,\mathbf Cv)_1.
\]
Thus the historical and single-lattice integrands
are identical, with the same determinant branch
and the same real integration variables. Their
Borel, positivity, and absolute-integrability
conditions coincide, by Tonelli's theorem.
Fubini's theorem then identifies their values.
The preceding bijection of guarded indices identifies
each finite-difference term, activity, and cluster
series, also with an insertion. Absolute convergence
of the corresponding series is the same condition.

\emph{The potentials.}
No stored datum at an earlier level has an endpoint
in \(\Lambda_t\), because it is external to
\(\Lambda_{j+1}\supset\Lambda_t\). The spaces of stored
data on \(Y\subset X\) are therefore singletons.
For \(k\ge1\), the shift domains and maps on \(Y\) are those of the
constant sequence. To check the trace condition in
\cref{sec:history-regional-shift}, use
\(\operatorname{dist}_\infty(\Lambda_t,\Omega_t^c)
\ge2LM\) in level-\(k\) scale. The blocks meeting
\(Y\) and the fivefold enlargements required there
lie within distance \(5\mathfrak M<2LM\) of
\(\Lambda_t\); all belong to every member of the
collared sequence. The last trace also agrees,
since \(Y\subset\Omega_t\).

The margin from every earlier large-field region is
\[
 \operatorname{dist}_\infty(\Lambda_t,Z_j)
    \ge8LM\check R_{j+1}\ge8M\check R_j
                            \qquad(1\le j<t),
\]
in level-\(j\) scale. Write \(s_j\) for the cell side
in a common length scale. Every level-\(j\) cell
meeting \(Y\) is inside \(\Lambda_j^{\sim-1}\).
Also \(\varphi_j=1\) at distance at most \(s_j\)
from \(Y\): translating such a point by a mollifier
argument of norm at most \(s_j/4\) leaves distance
at least \((8-1-\tfrac14)s_j\) from \(Z_j\).
In particular the excess Wilson weight vanishes on
plaquettes meeting \(Y\).

All sites of \(Y\) are in the top layer \(k\).
The regular indices with base in \(Y\) are precisely
the common wide indices and the octets whose level
is \(k\). Their functions, domains, and subtraction
at the trivial pair agree. No edge-tent or previously
created boundary index has base in \(Y\), since
such an index must meet the complement of an
earlier \(\Lambda_j^{\sim-1}\).
Nor can a retained-factor base lie in \(Y\) when
the retained factors have their prescribed indices:
an older factor must meet some \(\Omega_{l+2}^c\),
and a factor of the last old step must meet
\(Z_k^\sim\). The first possibility contradicts
the nesting, and the second the distance
\(8s_k>s_k\) from \(Z_k\).

These facts identify the four regular potentials
as partial maps. For J, the pivot paths lie within
the double blocks of the coarse bonds, hence within
their local domains; the free coordinates and pivot
solutions agree. For G, the cube, block field, and
gauge-fixing coordinates agree. For the action
potential, the common shift is substituted into
the same top-layer Wilson members; all lower-layer
members vanish. The Taylor subtraction giving A
therefore agrees. Finally the regular and
counterterm members have the same plain substituted
triples, indices, and starred brackets. This
identifies \(\mathbb V'\), with the same restriction
to the shift domain and field polydisc. At \(k=0\),
\(\mathbb V'=0\) and the completed pivot field gives
the same action comparison directly.
None of these four potentials uses the pinned
substitution for boundary members.

The B and \(\Phi\) potentials are zero because their
index sets are empty. For \(k\ge1\), their domain restrictions
remain: they are defined on the common shift
domain with \(g_k^{-1}\mathcal A\in\mathbb D^{\rm f}(Y)\).
At \(k=0\) they are zero functions restricted only by
\(g_0^{-1}\mathcal A\in\mathbb D^{\rm f}(Y)\).
Under the stated retained-index condition this
holds for both potentials. A Mayer index carrying
one of these origins consequently has a factor
\(e^0-1\). Wherever the full cluster sum is
defined, deleting such zero terms gives the
regular cluster sum.

\emph{Pointing.}
The four identical potential families give the same
interpolation and bond insertions. Together with the
Gaussian comparison, this identifies
\(E^\chi(X,b;P)\) and \(E^\bullet(X,b;P)\) for every
common bond, as partial maps.
There is a small domain issue for the other bonds.
If \(b^\sharp\not\subset X\), the empty-system cluster
sum is independent of its cutoff parameter \(t_b\).
Its derivative is zero wherever defined; if a
cutoff term for any other bond is defined, the
diagonal cluster sum is defined for almost every
parameter, which also defines this zero term.
For the potential interpolation the insertion at
such a bond is identically zero on every local
domain. Its domain contains the domain of any
other bond insertion. A defined insertion term
for another bond supplies all the Gaussian
operators and unmarked terms; replacing its
insertion by zero gives absolutely convergent
zero integrals and a zero series. Thus here too
definedness at one bond supplies definedness and
vanishing at every bond with \(b^\sharp\not\subset X\).

For \(z\in X^{(t)}\), the bond
\(b_z=\langle z,z+e_1\rangle\) is in \(\mathfrak c^+\)
and \(h_z(b_z)\ge1-(2L)^{-1}>0\).
A pointed term therefore includes at least one
common bond. The preceding observation removes
all additional bonds without changing its domain.
On the remaining bonds the weights, values,
and domains agree. This proves
\eqref{eq:history-interior-pointed}.

\emph{Summation.}
Apply \cref{II-20-103} with \(\tau=1\) to the four
families in \eqref{II-135-8}. Its path has two
origins, each of norm at most \(w_\star\), and
endpoints
\[
 \underline{\mathbb V}_1
     =\{\mathbb V^{\rm J}+\mathbb V^{\rm G}
                        +\mathbb V^{\rm A},\mathbb V'\},
 \qquad
 \underline{\mathbb V}_0=\{0,0\}.
\]
The velocities have norms at most \(4w_\star\);
the sum of the two bond insertions has norm at
most \(18w_\star\) and is Borel.
Thus \cref{prop:cluster-path} applies and gives
\[
 \sum_bE^\bullet(X,b)
       =E[\underline{\mathbb V}_1](X)
                         -E[\underline{\mathbb V}_0](X).
\]
Origins with zero potential make no contribution,
so \(E[\underline{\mathbb V}_0]=E[\emptyset]\).
The cutoff identity of \cref{II-20-87} gives
\[
 \sum_bE^\chi[\emptyset](X,b)
                   =E[\emptyset](X)-E^\circ[\emptyset](X).
\]
The last term is zero: at infinite threshold and
with no potential vertices every Mayer function
contains a vanishing outside test. The Gaussian
cluster estimates remain valid at infinite
threshold since the cutoff error is then zero.
Combining these two identities proves the first
equality in \eqref{eq:history-interior-sum}.
All bond terms vanish unless \(b^\sharp\subset X\).
The tent partition of unity therefore gives the
second equality, with no sites outside \(X\).

Finally, the four analytic families are defined
on the whole field polydisc. For \(k\ge1\), definedness of the
\(\prime\) family includes membership of the
common shift domain. The B and \(\Phi\) families
are consequently defined and zero at every real
field below the threshold. For \(k=0\), these two
families are zero on the field polydisc by definition,
so the same conclusion holds without a shift-domain
condition. The full system has
four origins with norms at most \(w_\star\), so
\cref{II-19-55} proves its absolute convergence.
Deleting the zero boundary vertices identifies it
with the regular system. This proves the last
assertion on the stated analytic domain.
\end{proof}


\section{Evaluation on critical orbits}
\label{sec:history-evaluation}

The density is a function of the determining variables and
the stored fields. Its local terms are written using a
critical configuration. We show that their values do not
depend on this choice, and that the configuration produced
by the shift gives the same old density.

\subsection{Adapted representatives}

\begin{definition}\label{II-10-37}
Let \(1\le n\le K\), and let
\(\Omega_0\supset\cdots\supset\Omega_n\) be a
separated collared sequence, with reduced determining
set \(\mathbf B\). A top \(y\in\Lambda_q(\mathbf B)\),
\(q\ge1\), has stored sites \(a\in B(y)\). The transport
of a configuration \(U'\) at such a site is
\[
 \mathsf{ax}_{U'}(a)=(\mathcal M^{q-1}U')(y,a).
\]
We call \(U'\) adapted if all these averaged contour
variables are defined and equal to \(1\). Adaptedness
at a set of sites requires this only at its stored
sites. In particular, adaptedness at layer \(q\)
means that \(\mathcal M^{q-1}U'\) is in block axial
gauge on \(\Lambda_q(\mathbf B)\).
\end{definition}

\begin{proposition}\label{II-10-41}
Under the hypotheses of \cref{II-10-20} and
\eqref{eq:history-background-smallness}, every regular
critical orbit of a completed collared sequence has
an adapted representative. Such a representative can
be chosen continuously, and hence Borel measurably,
on the set of data that pass the strict regularity
tests. Every adapted representative gives the same
averaged background on the bonds of the last step.
\end{proposition}

\begin{proof}
The pillars, parent \(z(x)\), and level \(n(x)\) of a
non-top site are as in \cref{I-3-95}. Write \(q(x)\)
for its pillar's level. The parent belongs to the
same pillar and has a strictly higher level. The
nodes of the reduced determining set are its tops:
every incident row endpoint is a top, and every top
has an incident row in its layer. For the latter
assertion, the layer is a union of the separated
partition cubes; a neighboring site within the same
cube gives the row. The argument for the bare layer
also covers every endpoint of an inner bare bond.

Put \(e=\bar C_{\rm reg}\varepsilon_n\). The layer
estimate gives
\[
 |\partial U'(p)-1|<eL^{-2q}
\]
on bare plaquettes in a pillar of level \(q\).
Fix such a pillar \(B^q(y)\), and set
\(Y_{q-1}=B(y)\) and
\(Y_l=B^{q-1-l}(Y_{q-1})\). Thus \(Y_0=B^q(y)\)
and \(B(Y_{l+1})=Y_l\). Finiteness and the strict
curvature inequalities permit \(0<e'<e\) such
that all bare plaquette bounds in this pillar
hold with \(e'\) in place of \(e\).
Apply the local induction in the proof of
\cref{II-195-31}, with terminal index \(q-1\)
and \(\alpha=e'L^{-2}\). Every double block
needed for a coarse plaquette in \(Y_{l+1}\)
lies in \(Y_l\), so no curvature bound outside
the pillar is used. The induction gives
\[
 |\partial(\mathcal M^lU')(p')-1|
     \le2e'L^{2(l-q)}<2eL^{2(l-q)},\qquad 0\le l<q.
\]
All block contours used below are therefore in
the unique analytic mean chart. Here the bound
\(e\le c^{\mathrm{SO}}_2\) follows from
\eqref{eq:history-background-smallness}.

For every non-top bare site define
\[
 \tau_{U'}(x)
       =(\mathcal M^{n(x)}U')(z(x),x).
\]
Unlike the comb transports of \cref{I-3-96}, these
are averaged contour variables. They are real
analytic in a neighborhood of \(U'\), and satisfy
\[
 \tau_{U'^w}(x)
       =w(z(x))\tau_{U'}(x)w(x)^{-1}.
\]
Indeed, every contour acquires these same endpoint
factors, and uniqueness of the group mean preserves
them. Define \(u_{U'}=1\) at the tops and recursively
set
\[
 u_{U'}(x)=u_{U'}(z(x))\tau_{U'}(x).
\]
The recursion terminates at a top. It fixes all
nodes, so \(U'^{\,u_{U'}}\) remains in the same
critical orbit, and all its \(\tau\)'s equal \(1\).
The stored sites at layer \(q\), other than the
top itself, are precisely the sites with
\(q(x)=q\) and \(n(x)=q-1\). Thus this
representative is adapted.

If two representatives in the same residual orbit
have all their \(\tau\)'s equal to \(1\), the
covariance formula gives \(w(x)=w(z(x))\).
Since \(w=1\) at the tops, induction gives \(w=1\)
on every pillar. The construction consequently
depends only on the datum. Locally, choose the
real analytic critical configuration supplied by
\cref{I-58-22}, and apply the finite recursion.
This proves continuity of the chosen representative.
The set of strictly regular data is open: the
pure tested plaquettes are products of datum
coordinates; on the set where these tests pass,
the other tested bonds are analytic one-step
averages of those coordinates. All remaining
tests are finitely many strict inequalities of
continuous functions. Hence the selection is Borel.

Finally let \(U'\) and \(U''=U'^w\) be adapted
representatives. Every endpoint of a bond in the
last step is either a node or a stored site of
a last-layer top. The transport covariance gives
\(w=1\) at these endpoints. Covariance of
\(\mathcal M^{n-1}\) then gives equality of the
averaged backgrounds, including at bonds crossing
the last region's boundary.
\end{proof}

\subsection{Rigidity at the stored sites}

\begin{lemma}\label{II-167-6}
Let \(U''=U'^w\) lie in the same residual orbit
for a separated sequence, so that \(w=1\) at its
nodes. If the transports of \(U'\) and \(U''\)
are defined and equal at a stored site \(a\),
then \(w(a)=1\).

Consequently, let \(f(P;\mathfrak q)\) be a jointly
gauge-invariant local term whose stored fields
transform only at a set \(\mathsf N_f\) of nodes
and stored sites. If the two configurations have
equal transports at the stored sites in
\(\mathsf N_f\), and their plain triples belong
to the domain of \(f\), then, at fixed stored data,
\begin{equation}\label{II-167-7}
 f(P(U'');\mathfrak q)=f(P(U');\mathfrak q).
\end{equation}
\end{lemma}

\begin{proof}
If \(y\) is the top of \(a\), covariance gives
\begin{equation}\label{II-167-5}
 \mathsf{ax}_{U'^w}(a)
       =w(y)\mathsf{ax}_{U'}(a)w(a)^{-1}.
\end{equation}
The top is a node, so \(w(y)=1\). Cancellation
of the common transport gives \(w(a)=1\).
At \(a=y\) the transport itself is always \(1\).

The first assertion gives \(w=1\) on
\(\mathsf N_f\), so \(w\cdot\mathfrak q=\mathfrak q\).
Plain triples are covariant: their configuration
is restricted, their second component is zero,
and their current is covariant at complete stencils
and zero at the other stencils. Joint invariance
therefore proves \eqref{II-167-7}.
\end{proof}

For a history \(h\) of depth \(n\), let
\[
 \mathsf{Sd}_h
  =\{b_-:b\in\mathfrak X_l^{{\rm E}+}(h),
                                      \ 0\le l<n\}.
\]
These are the only sites where a gauge transformation
acts on its stored data:
\((w\cdot\mathfrak q)_l(b)=
        \operatorname{Ad}(w(b_-))\mathfrak q_l(b)\).

\begin{lemma}\label{II-167-13}
Every site of \(\mathsf{Sd}_h\) is a node or a
stored site of a top that can be pinned. If
\(h'=(h,\mathsf o_t)\), \(t=k+1\), then
\(\mathsf{Sd}_h\cap\Omega_t=\emptyset\), and
the old stored sites have the same tops, levels
and transports for the shorter and longer
collared sequences.
\end{lemma}

\begin{proof}
Use the external-endpoint geometry of
\cref{II-84-18}. For an external bond of level
\(l\), its initial point \(x\) lies outside
\(\Lambda_{l+1}\). If \(x\in\Omega_{l+1}\),
it is stored at the pinnable top
\(\mathsf c_{l+1}(x)\) of level \(l+1\).
Otherwise it is a top of level \(l\), and a node;
at \(l=0\) it is the endpoint of an inner bare
bond and hence a bare node.

For \(l<k\), the inclusion
\(\Omega_t\subset\Lambda_k\subset\Lambda_{l+1}\)
puts \(x\) outside \(\Omega_t\).
All earlier layers of the two collared sequences
coincide. At level \(k\), their top sets differ
only inside \(\Omega_t\). Membership in this
region is constant on each block of order at
most \(t\), so the top of an old stored site
cannot enter the changed part. The transport
in either sequence is the same variable
\((\mathcal M^{q-1}U')(y,x)\).
\end{proof}

\subsection{The evaluations before and after a step}

Fix a bare field \(U\) in the support of a history
\(h\). The evaluation configurations are
\[
 \mathsf{Ev}_h[U]
   =\{U'\vee U:
      U'\in U^\diamond_n(\xi^\diamond_h(U))
                     \text{ is adapted}\}.
\]
At depth zero this set is \(\{U\}\).
The gluing uses \(U'\) on inner bonds of the
first collared region and \(U\) on the outer
bonds.

\begin{proposition}\label{II-168-27}
Assume the actual domain, response and levelling
hypotheses of \cref{sec:history-actual-potentials}.
Let \(U\) be in the support of \(h'=(h,\mathsf o_t)\),
and use its actual stored fields \(\mathfrak q^+\).
Then the following assertions hold.

Every local term in the exponent and retained-factor
sum of \(h'\) is defined at every
\(\mathsf U_+\in\mathsf{Ev}_{h'}[U]\).
Each of these two sums has one value on that set.

For \(k\ge1\), the shifted configuration
\(\mathsf U^\sharp\) associated with any such
\(\mathsf U_+\) belongs to the old critical orbit
on the inner bonds and equals \(U\) on the outer
bonds. It is adapted at \(\mathsf{Sd}_h\) for
the shorter sequence. Both old sums are defined
there, and for every \(\mathsf U^\flat\in
\mathsf{Ev}_h[U]\),
\begin{equation}\label{II-183-29}
\begin{split}
 A_{k,h}(\mathsf U^\sharp,\mathfrak q)
     &=A_{k,h}(\mathsf U^\flat,\mathfrak q),\\
 \Phi_h(\mathsf U^\sharp;\mathfrak q)
     &=\Phi_h(\mathsf U^\flat;\mathfrak q).
\end{split}
\end{equation}
Here \(\mathfrak q\) is the restriction of
\(\mathfrak q^+\) to levels less than \(k\).
The statement requires no adaptedness of
\(\mathsf U^\sharp\) at the other old sites.
\end{proposition}

\begin{proof}
The actual membership result \cref{II-99-39}
puts the plain triples of every glued critical
representative in the required core and layered
output spaces, and its stored fields in their
tiers. The same result at each earlier prefix
gives the domains of the old entries. Thus every
summand is defined, by the domain statements
for the regular, boundary and retained families.

Two critical representatives differ by a residual
gauge. Extend it by \(1\) outside the first
collared region. An outer bond has no endpoint
in \(\Omega_1\); an endpoint inside the collar
is therefore a bare top and a node. The extension
fixes both endpoints of every outer bond, so
the glued configurations differ by this same
gauge. By \cref{II-167-13}, all sites acting
on the stored data are nodes or stored sites.
Both representatives are adapted. Apply
\cref{II-167-6} to each boundary entry and each
canonical retained factor, which are jointly
invariant. The regular entries and the weighted
Wilson action are invariant without stored
data. The finite sums consequently have one value.

For the old evaluation, use the actual arrival
verification in \cref{II-113-38}.
Its datum is exactly \(\xi^\diamond_h(U)\):
the completed chart recovers the actual
level-\(k\) field on the changed rows and
leaves all other rows unchanged. Since
\(U\) also belongs to the support of the
prefix \(h\), \cref{II-10-20} supplies this
old orbit. The response identification
\cref{II-93-38} gives a regular critical
configuration up to gauge, and the pinned
levelling restores the exact constraints.
Thus \(\mathsf U^\sharp\) is a gluing from
that orbit. Its plain triples satisfy the
old domain bounds by \cref{II-99-39}.

The pinned shift preserves the transports at
the pinnable old tops. By \cref{II-167-13},
the old stored sites have the same transports
before the shift as in the longer sequence.
Those transports are \(1\), and hence remain
\(1\) after the shift. At old top sites the
transport is \(1\) identically. This proves
adaptedness at \(\mathsf{Sd}_h\).
An adapted representative of the old orbit
exists by \cref{II-10-41}. Compare it with
the shifted one by \cref{II-167-6}, term by
term as above. This proves both equalities
in \eqref{II-183-29} separately.

For the first step there is no old critical
choice: the determining map at level zero
is the identity, the old field is \(U\), and
there are no old retained factors. For an
empty first region the same statements reduce
to this unchanged-field evaluation.
\end{proof}

\subsection{The canonical retained factors}

\begin{proposition}\label{II-182-41}
At every evaluation configuration of \(h'\),
the retained-factor sum obtained from the
one-step integration is the canonical sum
of \cref{II-157-13}. This includes the
first step. The equality is an equality
of the exponents themselves.
\end{proposition}

\begin{proof}
An old factor has the same domain, tier and
formula after the history is extended. The
first collared region is unchanged, so its
plain triple is unchanged as well; and the
old components of \(\mathfrak q^+\) are
\(\mathfrak q\). This identifies every old
term separately.

For the new terms, let \(\mathsf a'=\mathfrak q^+_k\).
Actual output membership and the tier bounds
give all the domains in
\eqref{eq:formed-factor-definitions}. The
kept-form expansion and the definitions yield
\begin{equation}\label{II-182-42}
\begin{split}
 &\sum_{X\in\mathcal D_t}
       \mathsf K^{\rm kp}_{h'}(X;\mathsf U_+;\mathsf a')
       +\sum_{c\in\Gamma^\partial_t}z_c(\mathsf U_+)\\
 &\hspace{2em}
       =\sum_{f\text{ formed at }t}
                      \varphi_f(\mathsf U_+;\mathfrak q^+).
\end{split}
\end{equation}
Indeed, the level-\(k\) pair in the kept form
is exactly the pair read from the level-\(t\)
plain triple. The zero-extended external
field is the same on its domain. For an
elimination term, both expressions use the
same configuration on \(\mathsf P_c^t\);
the pivot logarithm depends only on these
bonds. Grouping these finitely many terms
by \(X(c)\) gives the second family in
\eqref{eq:formed-factor-definitions}.

The canonical family is the disjoint union
of the old family and the factors just formed.
Adding the preceding identities proves the
claim. At the first step the old family is
empty, so the same identity proves it directly.
If the first region is empty, all these sums
are empty.
\end{proof}

\section{The measure identity for one history step}
\label{sec:history-integrability}

The local expansion must represent the finite history measure.
We first give the measure argument, separating the two changes
of variables from the identities for their integrands. This also
fixes the order in which integrability is proved.

All coordinate spaces below are finite products of $G$ and
$\mathfrak g$, with positive bonds as independent coordinates.
They are Polish spaces. The reference measures in free
coordinates are products of normalized Haar measures and
Lebesgue measures. The completed-coordinate measures are
their images under the specified Borel injections. All are
$\sigma$-finite. A partial map is Borel when its domain and
its inverse images of Borel sets are Borel. A density given by
a partial map means a Borel function off a Borel null set with
integrable absolute value; its values on that null set are
immaterial. A product with a zero test is defined to be zero,
without evaluating its remaining factors.

We use two facts about Borel images. A Borel image of a Borel
subset of a Polish space is analytic and is measurable for
the completion of every finite Borel measure. It therefore
lies between two Borel sets whose difference has measure
zero. If the map is injective, its image is Borel and its
inverse on that image is Borel. These are the analytic-image
and Luzin--Souslin theorems; see
\cite[Proposition~13.2, Theorem~14.3 and Corollary~15.10]{Tserunyan2026}.

\subsection{Coordinates and the hypotheses of a step}

Fix a history $h$ at level $k$ and its extension $h'$ at
$t=k+1$. Write $\nu_h,\nu_{h'}$ for their measures with
stored coordinates. At this stage the coordinate maps and
the analytic data are inputs. Their required properties are
as follows.

The bare space carries the finite positive measure
$\rho_0\,dU$. Borel partial coordinate maps $\xi_h,\xi_{h'}$
and bounded nonnegative history weights give
\[
 \nu_h=(\xi_h)_*(\mathfrak w_h\rho_0\,dU),\qquad
 \nu_{h'}=(\xi_{h'})_*(\mathfrak w_{h'}\rho_0\,dU).
\]
There are a bounded nonnegative Borel function $f$ on the
old coordinate space $\mathsf X_h$ and a Borel partial map
$\Sigma:\mathsf X_h\rightharpoonup\mathsf X_{h'}$ such that
\begin{equation}\label{eq:history-coordinate-step}
 \mathfrak w_{h'}=\mathfrak w_h(f\circ\xi_h),\qquad
 \xi_{h'}=\Sigma\circ\xi_h,
\end{equation}
including equality of the domains in the second identity.
The first identity is required on $\operatorname{dom}\xi_h$.
Suppose that the preceding density is already known:
\begin{equation}\label{eq:history-old-density}
 \nu_h=\varrho_h\,\mathscr M_h,\qquad
 \varrho_h=s_h\exp(A_h+\Phi_h),\qquad
 \int|\varrho_h|\,d\mathscr M_h<\infty.
\end{equation}
Here $s_h\ge0$ is Borel, and $A_h,\Phi_h$ are defined and
Borel off a null set wherever $s_h\ne0$. Put $T=fs_h$.

The old coordinates split as $(V,\phi^+)$, with reference
measure $dV\,d\mathscr M^+(\phi^+)$. The new preliminary
coordinates are
\[
 p=(W,a',\phi^+),\qquad
 d\mathscr M^{\rm pre}(p)=dW\,da'\,d\mathscr M^+(\phi^+).
\]
There is an injective Borel map
$\imath:\mathsf X^{\rm pre}\to\mathsf X_{h'}$ with Borel
inverse on its image, and
$\mathscr M_{h'}=\imath_*\mathscr M^{\rm pre}$.

For a Borel set of bases $d=(W,\phi^+)\in\mathsf O$,
a chart $V_d(B)$ is defined on a fixed bounded open product
of balls $\mathsf R_r$ in the real field
$B=(B_{\rm I},a)$. Let
\[
 \mathsf Y=\{(W,B,\phi^+):d\in\mathsf O,\ B\in\mathsf R_r\},
 \qquad \mathsf{pt}(W,B,\phi^+)=(V_d(B),\phi^+).
\]
Require $\mathsf{pt}$ to be Borel and injective. Its Jacobian
$j_c$ is a finite positive Borel function, and for every
nonnegative Borel $g$ one has
\begin{equation}\label{II-166-48}
 \int_{\mathsf{pt}(\mathsf Y)}g\,d\mathscr M_h
      =\int_{\mathsf Y}j_c(y)g(\mathsf{pt}(y))\,d\lambda(y),
 \qquad d\lambda=dW\,dB\,d\mathscr M^+.
\end{equation}
Every bare field in $\operatorname{dom}\xi_h$ with
$\mathfrak w_{h'}\ne0$ must have a chain point:
$\xi_h(U)=\mathsf{pt}(y)$ for some $y\in\mathsf Y$.

The strip substitution is an injective Borel partial map
$a=\check\Psi_d(a')$, with finite positive Jacobian
$j_s(d,a')$. Set
\[
 \mathsf Q=\{(p,B_{\rm I}):d\in\mathsf O,\
       a'\in\operatorname{dom}\check\Psi_d,\
       (B_{\rm I},\check\Psi_d(a'))\in\mathsf R_r\},
\]
and
\[
 \Theta(p,B_{\rm I})=(W,(B_{\rm I},\check\Psi_d(a')),\phi^+),
 \qquad d\kappa=d\mathscr M^{\rm pre}\,dB_{\rm I}.
\]
The domains and maps are Borel. Require the exact substitution
formula and coordinate identity
\begin{align}
 \int_{\mathsf Y}g\,d\lambda
   &=\int_{\mathsf Q}j_s(g\circ\Theta)\,d\kappa,
                                  \label{II-166-50}\\
 \Sigma(\mathsf{pt}(\Theta(p,B_{\rm I})))&=\imath(p).
                                  \label{eq:history-chart-commutation}
\end{align}
The first holds for every nonnegative Borel $g$; the second
includes definedness at every point of $\mathsf Q$.

Let $S\ge0$ be the new Borel test on $\mathsf X_{h'}$, and
let $\mathsf R$ be the product of the strict integration
balls $|B_{\rm I}(b)|<\mathfrak b_b$, with
$0<\mathfrak b_b\le r$. Define
\[
 \mathsf{Pe}=\{p:S(\imath p)>0,\ d\in\mathsf O,\
 a'\in\operatorname{dom}\check\Psi_d,\ |\check\Psi_d(a')|<r\}.
\]
An empty set of integration bonds gives a one-point space of
volume one. Require
$\mathscr M^{\rm pre}(\{S\circ\imath>0\}\setminus\mathsf{Pe})=0$
and, outside a Borel $\kappa$-null set in $\mathsf Q$,
\begin{equation}\label{II-166-17}
 T(\mathsf{pt}(\Theta(p,B_{\rm I})))
 =\mathbf1_{\mathsf R}(B_{\rm I})S(\imath p)
                   c_z e^{-x_kG_d(B)},\qquad a=\check\Psi_d(a').
\end{equation}
Here $c_z>0$ and $G_d(B)$ is real Borel and finite.
In particular the set where the first two factors on the
right are nonzero is exactly $\mathsf{Pe}\times\mathsf R$.

The remaining hypotheses concern the analytic identities.
They are required whenever $y=\Theta(p,B_{\rm I})$ is the
chain point of a bare field with nonzero $\mathfrak w_{h'}$
and defined new coordinates. The attached representative
$\mathsf U=\mathsf U[d]$ and stored fields $\mathfrak q$
are Borel functions of the base. Write $\phi=\mathsf{pt}(y)$.
Evaluation of the old exponents is independent of the
representative, and
\begin{equation}\label{II-166-58}
\begin{aligned}
 A_h(\phi)-x_kG_d(B)+J_d(B)
    &=A_h(\mathsf U,\mathfrak q)
                   -\tfrac12(B,\mathsf A_dB)_1+V_d^{\rm ex}(B),\\
 \Phi_h(\phi)&=\Phi_h(\mathsf U,\mathfrak q)+V_d^\Phi(B).
\end{aligned}
\end{equation}
The old exponents at $\phi$ are used here wherever the old
density defines them. Equivalently the identities can first
be stated at an admitted representative of $\phi$.
The chart Jacobian has the form
\[
 j_c(W,B,\phi^+)=c_H\,e^{S_\Gamma(d)}e^{J_d(B)},\qquad c_H>0,
\]
where $S_\Gamma,J_d$ are real Borel functions. The symmetric
real matrix $\mathsf A_d$ has $\mathsf A_{{\rm II}}>0$,
and $V_d=V_d^{\rm ex}+V_d^\Phi$ is Borel and bounded on
$\mathsf R$ for fixed $p$.

We state the finite sum identities needed for collecting
the new action. They are the identities of the local
normalizer, conditional form, cluster expansion and strip
substitution, respectively. Put
\[
 \Delta^{\rm cond}=\mathsf A_{{\rm EE}}
       -\mathsf A_{{\rm EI}}\mathsf A_{{\rm II}}^{-1}\mathsf A_{{\rm IE}},
 \qquad \nu_{\mathfrak g}=\tfrac32\log(2\pi),
 \qquad n_{\rm I}=\#\mathfrak X_{\rm I}^+.
\]
Let $\mu_{d,a}$ be the real Gaussian probability measure with
covariance $\mathsf A_{{\rm II}}^{-1}$ and mean
$-\mathsf A_{{\rm II}}^{-1}\mathsf A_{{\rm IE}}a$.
The local totals $N,K,B^\times,E,B^{\rm sr}$ are finite Borel
values at their indicated arguments and satisfy
\begin{align}
 S_\Gamma&=S_\Lambda+S_\partial,\qquad
 N=-\tfrac12\log\det\mathsf A_{{\rm II}}+S_\Lambda,
                                           \label{eq:history-normalizer-total}\\
 K(a)+B^\times(a)&=-\tfrac12(a,\Delta^{\rm cond}a)_1,
                                           \label{eq:history-kept-total}\\
 \int\mathbf1_{\mathsf R}e^{V_d(B_{\rm I},a)}\,d\mu_{d,a}
       &=e^{E(a)},                          \label{II-166-22}\\
 B^{\rm sr}(a')&=\log j_s+
       [E+B^\times+K](a)-[E+B^\times+K](a').
                                           \label{II-166-23}
\end{align}
Require these identities and the following definitions to
have all their entries defined at such a realized point:
\begin{equation}\label{eq:history-new-exponents}
\begin{aligned}
 n_t&=\log c_H+\log c_z+\nu_{\mathfrak g}n_{\rm I},\\
 A^+(p)&=A_h(\mathsf U,\mathfrak q)+n_t
                   +N+E(a')+B^\times(a')+B^{\rm sr}(a'),\\
 \Phi^+(p)&=\Phi_h(\mathsf U,\mathfrak q)+K(a')+S_\partial.
\end{aligned}
\end{equation}
The right sides depend only on $p$, are Borel partial maps,
and give the same values for every admitted representative
of $\imath p$. In particular one admitted representative
exists whenever these conditions apply. Collecting the
finite local sums into the next regular and boundary
families must preserve these identities as complex numbers.

\subsection{Positive fibre densities and integrability}

\begin{theorem}[The measure identity]\label{II-166-30}
Under the preceding hypotheses, the new exponents are
defined and single-valued at $\mathscr M_{h'}$-almost every
point where $S$ is nonzero. The density
\begin{equation}\label{II-166-39}
 \varrho_{h'}=S\exp(A^++\Phi^+)
\end{equation}
is Borel, nonnegative and integrable after assigning zero
on a Borel null set, and
$\nu_{h'}=\varrho_{h'}\mathscr M_{h'}$.
It is strictly positive almost everywhere on $\{S>0\}$.
For every bounded complex Borel $F$,
\begin{equation}\label{II-166-32}
 \int F\,d\nu_{h'}
  =\int F(\imath p)S(\imath p)
       e^{A_h(\mathsf U,\mathfrak q)+\Phi_h(\mathsf U,\mathfrak q)}
                  \mathsf I(d,a')\,d\mathscr M^{\rm pre}(p),
\end{equation}
where
\begin{equation}\label{II-166-14}
 \mathsf I(d,a')=c_Hc_z e^{S_\Gamma(d)}j_s(d,a')
       \int_{\mathsf R}
          e^{-(B,\mathsf A_dB)_1/2+V_d(B)}\,dB_{\rm I}.
\end{equation}
The integrals are absolute. Undefined values on the
exceptional null set are assigned zero.
\end{theorem}

\begin{proof}
First make the old density nonnegative everywhere. Its
integral over every Borel set is real and nonnegative,
because $\nu_h$ is positive. Applying this to the sets
where its imaginary part has either sign and where its
real part is negative shows that it is nonnegative almost
everywhere. Enlarge the exceptional Borel null set to $N_h$
and set the density to zero there. Denote the result by
$\bar\varrho_h$. It vanishes wherever $s_h=0$ and is
strictly positive outside $N_h$ wherever $s_h>0$.
Consequently $f\bar\varrho_h$ vanishes where $T=0$, and
equals $T e^{A_h+\Phi_h}>0$ outside $N_h$ where $T>0$.

The bare identities \eqref{eq:history-coordinate-step}
give $\nu_{h'}=\Sigma_*(f\nu_h)$. Every selected bare
field has a chain point, so $f\nu_h$ is concentrated on
the Borel set $\mathsf{pt}(\mathsf Y)$.
Both changes of variables transport null sets backwards:
substitute the indicator of a null set in
\eqref{II-166-48} or \eqref{II-166-50}; the integral of a
strictly positive Jacobian over its inverse image is zero,
so that inverse image has measure zero.
Equation~\eqref{eq:history-chart-commutation} and
\eqref{II-166-50} also show that $\mathsf{pt}(y)$ lies in
$\operatorname{dom}\Sigma$ for almost every $y\in\mathsf Y$.

We next justify applying the analytic identities, which
were stated at realized bare fields. Let $\mathsf B$ be
the Borel set of bare fields with $\mathfrak w_{h'}>0$ and
defined $\xi_{h'}$. The finite measure
$\nu=(f\nu_h)|_{\operatorname{dom}\Sigma}$ is the
$\xi_h$-image of the corresponding bare measure.
The analytic-image theorem supplies Borel sets
$D'\subset\xi_h(\mathsf B)\subset D''$ with
$\nu(D''\setminus D')=0$. The bare formula gives
$\nu(\mathsf X_h\setminus D'')=0$, so $\nu$ is concentrated
on $D'$. Since its density is strictly positive on
$\{T>0\}\cap\operatorname{dom}\Sigma$ outside $N_h$,
that set outside $D'\cup N_h$ is $\mathscr M_h$-null.
Pulling the null sets back through the chart gives a
Borel set $\mathsf Y^*\subset\{T\circ\mathsf{pt}>0\}$
of full measure in that set, each of whose points is
the chain point of a field in $\mathsf B$ and avoids $N_h$.

Remove also the exceptional set of \eqref{II-166-17} and
pull $\mathsf Y^*$ back through $\Theta$.
Its intersection $\mathsf Q^*$ with
$\mathsf{Pe}\times\mathsf R$ has full product measure
there. Indeed the test identity puts every other point
of that product in $\Theta^{-1}\{T\circ\mathsf{pt}>0\}$,
and the inverse image of its omitted null set is null.
By Tonelli, the set
\[
 \mathsf{Pe}^*=
 \{p\in\mathsf{Pe}:
   \operatorname{vol}\{B_{\rm I}\in\mathsf R:
                  (p,B_{\rm I})\notin\mathsf Q^*\}=0\}
\]
is Borel and has full measure in $\mathsf{Pe}$.
For each such $p$ the good fibre has the positive volume
of $\mathsf R$. Choose one point in that fibre.
Every hypothesis that is independent of $B_{\rm I}$,
including definedness of \eqref{eq:history-new-exponents},
therefore holds at $p$. The identities involving the
integration field hold almost everywhere in its fibre.
No measurable choice of this one point is needed.

Define the nonnegative Borel function on $\mathsf Q$
\begin{equation}\label{eq:history-positive-integrand}
 D(p,B_{\rm I})=j_s(d,a')j_c(\Theta(p,B_{\rm I}))
       (f\bar\varrho_h)(\mathsf{pt}(\Theta(p,B_{\rm I}))).
\end{equation}
Using the two substitutions for nonnegative functions
and then the commuting identity gives
\begin{equation}\label{II-166-68}
 \int F\,d\nu_{h'}
      =\int_{\mathsf Q}F(\imath p)D(p,B_{\rm I})\,d\kappa
                  \qquad(F\ge0\text{ Borel}).
\end{equation}
All these equalities hold with values in $[0,\infty]$.
By the test identity, $D=0$ almost everywhere off
$\mathsf{Pe}\times\mathsf R$.

At a good fibre point, insert \eqref{II-166-17} and
\eqref{II-166-58} into \eqref{eq:history-positive-integrand}.
The result is
\begin{equation}\label{II-166-70}
\begin{split}
 D(p,B_{\rm I})={}&S(\imath p)
      e^{A_h(\mathsf U,\mathfrak q)+\Phi_h(\mathsf U,\mathfrak q)}\\
 &\quad\times c_Hc_z e^{S_\Gamma(d)}j_s(d,a')
                 e^{-(B,\mathsf A_dB)_1/2+V_d(B)}.
\end{split}
\end{equation}
For fixed $p\in\mathsf{Pe}^*$ the last expression is
bounded on the bounded field box: $V_d$ is bounded and
the quadratic form is continuous. It is therefore
integrable. Its integral is the integrand of
\eqref{II-166-32}, without $F$, and is strictly positive,
since $D>0$ on the full-measure good fibre.

Complete the square in $\mathsf A_d$:
\[
 (B,\mathsf A_dB)_1
  =(B_{\rm I}+\mathsf A_{{\rm II}}^{-1}\mathsf A_{{\rm IE}}a,
    \mathsf A_{{\rm II}}
    (B_{\rm I}+\mathsf A_{{\rm II}}^{-1}\mathsf A_{{\rm IE}}a))_1
       +(a,\Delta^{\rm cond}a)_1.
\]
Gaussian normalization and \eqref{II-166-22} give
\[
 \int_{\mathsf R}e^{-(B,\mathsf A_dB)_1/2+V_d(B)}\,dB_{\rm I}
 =\exp\{\nu_{\mathfrak g}n_{\rm I}
       -\tfrac12\log\det\mathsf A_{{\rm II}}
          +K(a)+B^\times(a)+E(a)\}.
\]
Use \eqref{II-166-23} to replace $a$ by $a'$ and
\eqref{eq:history-normalizer-total} to combine the
determinant and pivot factors. The exponent in the
fibre integral of \eqref{II-166-70} becomes exactly
$A^+(p)+\Phi^+(p)$ from
\eqref{eq:history-new-exponents}.
These are identities between complex numbers, before
taking exponentials.

Set $\Phi_{\rm step}(p)=S(\imath p)e^{A^+(p)+\Phi^+(p)}$
on $\mathsf{Pe}^*$ and zero elsewhere. It is Borel,
nonnegative, and equals the fibre integral of $D$
almost everywhere. Tonelli in \eqref{II-166-68} yields
\[
 \int F\,d\nu_{h'}
   =\int F(\imath p)\Phi_{\rm step}(p)\,d\mathscr M^{\rm pre}
                         \qquad(F\ge0\text{ Borel}).
\]
Taking $F=1$ proves
$\int\Phi_{\rm step}\,d\mathscr M^{\rm pre}
=\nu_{h'}(\mathsf X_{h'})<\infty$.
Decomposing a bounded complex $F$ into four nonnegative
parts now proves \eqref{II-166-32} with absolute convergence.

Finally $\imath(\mathsf{Pe}^*)$ is Borel. The set
$\{S>0\}\setminus\imath(\mathsf{Pe}^*)$ is
$\mathscr M_{h'}$-null, by
$\mathscr M_{h'}=\imath_*\mathscr M^{\rm pre}$ and the
assumed full measure of $\mathsf{Pe}$. On its complement,
the representative-independent exponents descend through
the Borel inverse of $\imath$. Their density pulls back
to $\Phi_{\rm step}$. The pushforward identity above
therefore gives \eqref{II-166-39}, its positivity and
integrability. If collecting the finite sums identifies
$A^+$ with the exponent $A_{t,h'}$ of the next format,
the same equality proves that format without a further
integrability assumption.
\end{proof}

\section{Boundary terms and their propagation}
\label{sec:boundary-assembly}

At each step we retain the regular expansion in the interior.
The remaining terms are attached to their level and domain.
We first estimate the terms created by one integration and
then show that the same bound suffices at every level.

\subsection{The terms created at the boundary}

Fix a history step $h'=(h,\mathsf o_t)$, $t=k+1$.
Write $\mathscr U_X$ for its graded space at the intermediate
stage, and $\mathfrak Q^{\rm rad}_{h'}(X)$ for the stored data
with the prescribed tier bounds. A point of this product is
written $(P;\mathfrak q,\mathsf a)$: $\mathfrak q$ contains the
old data and $\mathsf a$ the new external field, extended by
zero off its bonds in $X$. In particular
$|\mathsf a|<C_{\rm xt}\mathfrak b_k$.
The background pair read from $P$ in level-$k$ units is
$p_X(P)$, and its further restriction to the joined strip
region is $p_X^{\rm loc}(P)$.

For this subsection assume that these readings belong to
the normalizer class, the local form class and the classes
of the localized strip shift, respectively. Assume also
that restriction to $X_0\subset X$ preserves these memberships
and the stored-data tiers. On each fixed real-$U$ slice,
$\mathscr U_X$ is open, its readings are holomorphic, and
they are continuous on the full space of triples.
The necessary estimates for the three analytic ingredients
are
\begin{align}
 |N^\partial_h(X;p)|&\le B_Ne^{-\delta_NM d_t(X)},
                         \label{eq:boundary-normalizer-input}\\
 |K^{\rm kp}_h(X;p,\mathsf a_1)|
       &\le c_{\rm kp}b_{\rm e}^2e^{-\kappa''d_t(X)},\notag\\
 |B^\times_h(X;p,\mathsf a_1)|
       &\le c_{\rm rj}b_{\rm e}^2
                    e^{-(\check R_t-2)}e^{-\kappa''d_t(X)}
       \quad (|\mathsf a_1|\le b_{\rm e}),\label{eq:boundary-form-input}\\
 |E_h^{\rm cl}(X;\mathsf a_1;P;\mathfrak q)|
       &\le b_{\rm cs}e^{-\kappa''d_t(X)}.
                         \label{eq:boundary-cluster-input}
\end{align}
The two forms are quadratic polynomials in the external
field, with coefficients analytic in $p$, and involve
only external bonds inside $X$. The cluster function is
local in the same bonds and is holomorphic jointly in
the complex coordinates of $P$ and in the perturbations
$\mathsf e$ of a tier field, supported on zone bonds of $X$,
with
\[
 |\mathsf e|<r_{\rm bd}:=2\rho_{\rm bd}/g_k.
\]
All three ingredients are Borel measurable, gauge invariant,
conjugation real and functions of the background pair.
These are assumptions on their full indicated domains,
including every such external perturbation.

We use the following support and interior identities.
The normalizer vanishes unless $X$ contains a cube of
$\Lambda_t$. The cluster function vanishes unless $X$
and $\Omega_t$ have a bare cube in common.
The rejected form vanishes unless $X$ contains a cube of
$\Omega_t\cap Z_t$ and $d_t(X)\ge\check R_t-2$.
For $X\subset\Lambda_t$ the local normalizer and pointed
cluster terms are those of the empty history:
\begin{equation}\label{eq:boundary-interior-input}
 N^{\rm un}_h(X,p,z)=N^{\rm un}_{\emptyset}(X,p,z),
 \qquad
 E_h^{\rm cl}(X;\mathsf a_1;P;\mathfrak q)
           =\sum_zF_{\emptyset}(X,P,z).
\end{equation}
The normalizer equality includes equality of its natural
domains. The cluster equality holds throughout the admitted
analytic domain specified above; its right side is independent
of the stored data and external field.

Set
\begin{equation}\label{II-151-32}
 \kappa_T=(1+4\beta_{\rm s})\kappa,\qquad
 \kappa''=\kappa_T+1+c_e.
\end{equation}
Use the localized shift of \cref{sec:strip-transfer} with
$\widehat\kappa=\kappa_T$ and impose
\begin{equation}\label{II-151-34}
 C_{\rm xt}\delta_k\le\tfrac18,\qquad
 0<\rho_{\rm bd}\le\tfrac12C_{\rm xt}\delta_k,
\end{equation}
together with \eqref{II-73-14} and \eqref{II-73-40}.
Here $\rho_{\rm bd}$ is a budget for the current step; the common
choice is given in \eqref{eq:common-boundary-budget}.
The normalizer bound is supplied by \cref{II-69-18} with
$B_N=24(LM)^3E_{\rm nf}$ and $\delta_N=\delta_{\rm nf}$.
The other inputs will be verified for the actual history
construction before this estimate is applied.

For a target domain $X_1$, restrict the background and old
data to each $X_0\subset X_1$. Define four families over $X_0$:
\begin{align*}
 F^{\rm E}_{X_0}&=E_h^{\rm cl}(X_0;\,\cdot\,;P|_{X_0};
                                      \mathfrak q|_{X_0}),\\
 F^8_{X_0}&=K^{\rm kp}_h(X_0;p_{X_0}(P),\,\cdot\,),\\
 F^9_{X_0}&=B^\times_h(X_0;p_{X_0}(P),\,\cdot\,),\\
 F^7_{X_0}(\mathsf a_1)&=
      \sum_{\substack{b\in\mathcal Z^+\\\mathsf P(b)=X_0}}
                              \mathsf{ls}_b(\mathsf a_1).
\end{align*}
The first family is taken on its external perturbation domain.
For the second and third, restrict to
$|\mathsf a_1(b)|<2C_{\rm xt}\mathfrak b_k$ on the external
bonds meeting $X_0$. Families with $X_0\not\subset X_1$
are zero, except for $F^7$, which is independent of the
target domain. Here
\[
 \mathsf{ls}_b(\mathsf a_1)
  =-\log\sigma_1\!\left(g_k^2
             \langle\mathsf a_1(b),\mathsf a_1(b)\rangle\right),
 \qquad
 \sigma_1(w)=\left(\sum_{n\ge0}\frac{(-w)^n}{(2n+1)!}\right)^2,
\]
with its series logarithm and domain $g_k|\mathsf a_1(b)|<1$.
The reading block $\mathsf P(b)$ consists of the one or two
coarse cubes at the endpoints of $b$, so $d_t(\mathsf P(b))=0$.

The five strip terms are
\begin{equation}\label{II-151-23}
 R^6(X_1)=\sum_{b\in\mathcal Z_\theta^+}J_b(X_1),\qquad
 R^\nu(X_1)=\sum_{X_0\subset X_1}
                     (F^\nu_{X_0})^{\rm tr}(X_1),
       \quad \nu=7,8,9,{\rm E}.
\end{equation}
In $R^6$, $J_b$ is evaluated at the link argument
$A=g_k\mathsf a(b)$.
The local pair used here is $p_{X_1}^{\rm loc}(P)$.
We suppress it, the tier field and the old data from the notation.
Define
\begin{equation}\label{II-151-30}
 \begin{aligned}
 B^{\rm nm}(X)&=N^\partial_h(X;p_X(P))
       -\mathbf1_{\{X\subset\Lambda_t\}}
                    \sum_zN^{\rm un}_{\emptyset}(X,p_X(P),z),\\
 B^{\rm cs}(X)&=E_h^{\rm cl}(X;\mathsf a;P;\mathfrak q)
       -\mathbf1_{\{X\subset\Lambda_t\}}
                                      \sum_zF_{\emptyset}(X,P,z),\\
 B^{\rm rj}(X)&=B^\times_h(X;p_X(P),\mathsf a),\\
 B^{\rm sr}(X)&=\sum_{\nu=6,7,8,9,{\rm E}}R^\nu(X).
 \end{aligned}
\end{equation}
A term multiplied by the indicator is omitted when the
indicator is zero. This convention imposes no domain
condition on an unused expression.
The boundary term created at step $t$ is
\begin{equation}\label{II-151-39}
 B^{\rm mk}_t(X)=B^{\rm nm}(X)+B^{\rm cs}(X)
                             +B^{\rm rj}(X)+B^{\rm sr}(X).
\end{equation}

\begin{proposition}\label{II-151-38}
Under the assumptions above, all the terms in
\eqref{II-151-39} are defined on
$\mathscr U_X\times\mathfrak Q^{\rm rad}_{h'}(X)$.
They are holomorphic in the complex coordinates of $P$
on each fixed real-$U$ slice, Borel measurable, gauge
invariant, conjugation real and functions of the pair.
Put
\begin{equation}\label{II-151-33}
 \begin{aligned}
 b_{\rm rj}&=c_{\rm rj}C_{\rm xt}^2\mathfrak b_k^2
                                      e^{-(\check R_t-2)},\\
 b_{\rm sr}&=4e^{-\check R_t}
  \left[N_b\rho_J+N_eN_b+N_eb_{\rm cs}
   +4N_eC_{\rm xt}^2\mathfrak b_k^2
                   (c_{\rm kp}+c_{\rm rj}e^{-(\check R_t-2)})\right].
 \end{aligned}
\end{equation}
Then, with $d=d_t(X)$,
\[
 |B^{\rm mk}_t(X)|
 \le B_Ne^{-\delta_NMd}
       +(b_{\rm cs}+b_{\rm rj})e^{-\kappa''d}
       +b_{\rm sr}e^{-\kappa_Td}.
\]
If
\begin{equation}\label{II-151-37}
 \delta_NM\ge\kappa_T,\qquad
 b_{\rm cs}+b_{\rm rj}+b_{\rm sr}\le1,
\end{equation}
this gives
\begin{equation}\label{II-154-45}
 |B^{\rm mk}_t(X)|\le B_Te^{-\kappa_Td_t(X)},
 \qquad B_T=24(LM)^3E_{\rm nf}+1.
\end{equation}
The term vanishes when $X\subset\Lambda_t$, and unless
$X$ and $\Omega_t$ share a bare cube.
Every strip term vanishes when the strip is empty.
\end{proposition}

\begin{proof}
We check the room needed by each transfer.
For $F^{\rm E}$ it is part of the cluster hypothesis.
For the quadratic terms,
\[
 r_{\rm bd}\le C_{\rm xt}\delta_k/g_k
             =C_{\rm xt}\mathfrak b_k.
\]
Thus every permitted perturbation of a tier field remains
in the radius $2C_{\rm xt}\mathfrak b_k$ ball.
Polynomial dependence on that field gives holomorphy.
Locality permits extension by zero outside the bonds
meeting $X_0$ when using \eqref{eq:boundary-form-input}.
Consequently their norms are bounded by
\[
 4c_{\rm kp}C_{\rm xt}^2\mathfrak b_k^2e^{-\kappa''d_t(X_0)},
 \qquad
 4c_{\rm rj}C_{\rm xt}^2\mathfrak b_k^2
             e^{-(\check R_t-2)}e^{-\kappa''d_t(X_0)}.
\]

For the Haar term, let $s(w)=\sum_{n\ge0}(-w)^n/(2n+1)!$.
For $|w|<1$, the successive absolute terms after the
first have ratio at most $1/20$, so
$|s(w)-1|\le10/57<0.18$. Hence
$|\sigma_1(w)-1|<2/5$ and
$|\log\sigma_1(w)|\le(2/5)/(1-2/5)=2/3$.
The series is normally convergent there.
On the perturbation ball,
$g_k|\mathsf a+\mathsf e|<1/8+2\rho_{\rm bd}\le3/8$,
so every one-bond function has the required room.
For a fixed nonempty reading block $\mathsf P(b)$
there are at most $N_b$ such bonds, and its size is zero.
It follows that
$\|F^7_{X_0}\|\le N_be^{-\kappa''d_t(X_0)}$.

Equation \eqref{II-73-37} now gives
\begin{equation}\label{II-151-40}
 \begin{aligned}
 |R^7(X)|&\le4N_eN_be^{-\check R_t}e^{-\kappa_Td},\\
 |R^{\rm E}(X)|&\le4N_eb_{\rm cs}e^{-\check R_t}e^{-\kappa_Td},\\
 |R^8(X)|&\le16N_ec_{\rm kp}C_{\rm xt}^2\mathfrak b_k^2
                                  e^{-\check R_t}e^{-\kappa_Td},\\
 |R^9(X)|&\le16N_ec_{\rm rj}C_{\rm xt}^2\mathfrak b_k^2
             e^{-(\check R_t-2)}e^{-\check R_t}e^{-\kappa_Td}.
 \end{aligned}
\end{equation}
For $R^6$, at most $N_b\#X\le8N_b(d+1)$ bonds contribute.
Since $\mu_H/2\ge\kappa_T+1$, \eqref{II-73-41} and
$(d+1)e^{-(d+1)}\le e^{-1}$ give
\[
 |R^6(X)|\le4N_b\rho_Je^{-\check R_t}e^{-\kappa_Td}.
\]
Their sum is bounded by $b_{\rm sr}e^{-\kappa_Td}$.

The first two differences in \eqref{II-151-30} vanish
inside $\Lambda_t$ by \eqref{eq:boundary-interior-input}.
Outside it, the normalizer and cluster estimates apply
directly. The tier bound in \eqref{eq:boundary-form-input}
gives $|B^{\rm rj}|\le b_{\rm rj}e^{-\kappa''d}$.
This proves the asserted sum bound, and
\eqref{II-151-37} gives \eqref{II-154-45}.

The support statements for the first three kinds follow
from the assumed support and interior identities.
For a nonzero strip term there is a zone bond with
$\mathsf P(b)\subset X$ and a localization set reaching
the component boundary. The strip geometry puts a cube
of $\Omega_t\cap Z_t$ in $X$. Thus a strip term cannot
occur inside $\Lambda_t$ or away from $\Omega_t$.
If the strip is empty there are no Jacobian terms and
no nonempty difference sets in the transfer sums.

Restriction of pairs, triples and stored fields is
continuous and linear in their complex coordinates.
The normalizer and form coefficients are analytic
functions of those pairs. For the cluster term use
its joint holomorphy with the external perturbation
set to zero. The transfer construction preserves these
properties: the values of the real tier field on the
cut-off zone are fixed on each holomorphic slice.
The finite sums also preserve Borel measurability and
the three symmetry and pair-dependence identities.
This proves the stated regularity.
\end{proof}

The total local exponent before the interior subtraction is
\[
 N^\partial_h(X;p_X(P))
       +E_h^{\rm cl}(X;\mathsf a;P;\mathfrak q)
       +B^{\rm rj}(X)+B^{\rm sr}(X).
\]
If the regular family is
$E^{(t)}(X,P,z)=N^{\rm un}_{\emptyset}(X,p_X(P),z)
                    +F_{\emptyset}(X,P,z)-n_{\rm tr}(X,z)$,
subtracting
$\mathbf1_{\{X\subset\Lambda_t\}}\sum_z(E^{(t)}+n_{\rm tr})$
gives exactly \eqref{II-151-39}.
All terms in this identity are defined by the interior
identities above; the constants $n_{\rm tr}$ cancel.
Moreover the finite sum identities
\eqref{II-73-39} and \eqref{II-73-43} show that the five strip terms
sum to the actual changes of the Haar density, the two
quadratic forms and the cluster exponent, together with
the logarithm of the strip Jacobian.

\subsection{Persistence of the boundary bound}

A boundary entry has an index $(j,X,o)$, where
$o={\rm mk}$ denotes a term created by an integration
and $o={\rm ed}$ an edge part of a regular family.
At $(j,X,{\rm mk})$ its domain is the layered output
space at level $j$, with the stored data up to $j$.
At $(j,X,{\rm ed})$ its domain is the level-$j$
analyticity core, with the same stored-data convention.
These domains depend only on the history truncated
at $j$ and the radii at levels at most $j$.

The permitted indices are
\[
 \begin{split}
 \mathsf{Ix}^{\rm mk}_j
   &=\{X\in\mathcal D_j:\ X\text{ and }\Omega_j
          \text{ share a bare cube},\
                           X\not\subset\Lambda_j^{\sim-1}\},\\
 \mathsf{Ix}^{\rm ed}_j
   &=\{X\in\mathcal D_j:\ X\subset\Lambda_j,\
             X^{(j)}\setminus\Lambda_j^{\sim-1}\ne\emptyset\}.
 \end{split}
\]
For each permitted index we keep at most one entry.
An entry is holomorphic on each fixed-real slice,
Borel measurable, gauge invariant, conjugation real
and dependent only on the pair and its stated stored data.
Its norm is the supremum on this domain.

For a pointed regular family with
$\|E^{(j)}\|_{j,\kappa+1}\le E_0$, define its edge part by
\begin{equation}\label{II-153-12}
 E^{(j),{\rm ed}}(X,P)=
   \sum_{z\in X^{(j)}\setminus\Lambda_j^{\sim-1}}
                                      E^{(j)}(X,P,z).
\end{equation}
It is independent of the stored data and has the same
regularity and symmetries as its summands. The number
of base points is at most
$M^3\#X\le8M^3(d_j(X)+1)\le8M^3e^{d_j(X)}$.
Therefore
\begin{equation}\label{II-153-13}
 \|E^{(j),{\rm ed}}(X)\|\le8M^3E_0e^{-\kappa d_j(X)}.
\end{equation}
For $X\subset\Lambda_j$ outside $\mathsf{Ix}^{\rm ed}_j$
the sum is empty.

\begin{proposition}\label{II-153-25}
Suppose the one-step hypotheses above hold at every
step up to $n$, with the same constants apart from the transfer
budget $\rho_{\rm bd}$, which may depend on the step, and the regular
families satisfy $\|E^{(j)}\|_{j,\kappa+1}\le E_0$.
Suppose also that each layered output space is contained
in the intermediate graded space used for
\cref{II-151-38}.
Start with no boundary entries. At step $t$ retain the
earlier entries, add $B^{\rm mk}_t(X)$ restricted to
the output space for $X\in\mathsf{Ix}^{\rm mk}_t$,
and add \eqref{II-153-12} for
$X\in\mathsf{Ix}^{\rm ed}_t$.
Then every entry at level $j\le n$ satisfies
\[
 \|B(j,X,o)\|\le B e^{-\kappa d_j(X)},\qquad
 B=\max\{1,B_T,8M^3E_0\}.
\]
Its domain and values are unchanged at subsequent levels.
\end{proposition}

\begin{proof}
At step $t$, the new entries have level $t$ and distinct
origins. Their indices are disjoint from those of all
earlier entries, so each index still has at most one term.
The old domains, tiers and index sets are determined by
the history truncated at their level; they do not change.
An old entry therefore retains its norm and its properties.
Restricting a newly created term to the layered output
space preserves its properties and can only decrease
its supremum. Its bound is $B_Te^{-\kappa_Td_t(X)}
\le Be^{-\kappa d_t(X)}$.
For an edge entry use \eqref{II-153-13}.
These three cases prove the assertion by induction.
The constant is not accumulated over the levels,
because distinct entries remain separately indexed.
\end{proof}


\subsection{The retained factors}
\label{sec:formed-factors}

The kept quadratic forms need not be small on the stored
field. We retain their exponentials, together with the
elimination factors at external pivot bonds.
At step $t=l+1$, the factors with domain $X$ have exponents
\begin{equation}\label{eq:formed-factor-definitions}
 \varphi^{\rm kp}_{t,X}(P;\mathfrak q)
       =\varphi_{\mathfrak k_X,X}
                    (\mathsf a[\mathfrak q];p_X(P)),
 \qquad
 \varphi^{\rm el}_{t,X}(P;\mathfrak q)
       =\sum_{\substack{c\in\Gamma^\partial_t\\X(c)=X}}
                                      z_c(p_X(P)).
\end{equation}
 
The first is included only when the keeping predicate
$\mathsf{kp}_t(X)$ is one, and the second only when its
index set is nonempty. The field $\mathsf a[\mathfrak q]$
equals the stored field of level $l$ on external bonds
inside $X$ and is zero elsewhere. The pair is read in
level-$l$ units:
\[
 p_X(P)=
   (e^{i\eta_tA'}U,L^{-3}\mathbf J)
          \big|_{\operatorname{Bd}^\circ(X\cap\Omega^\diamond_0)}.
\]
We restrict each exponent to the layered output space at
level $t$ and its level-$l$ stored-data tiers.
Its name is $(\nu,t,X)$, with $\nu={\rm kp}$ or $\nu={\rm el}$.
There is at most one factor of each kind for a given $(t,X)$.

Assume the hypotheses of the kept-form expansion on
$\mathscr O(X;\alpha_0,\alpha_1,\alpha_2)$, including
its rate condition
$\delta_\star^{\rm pl}LM\ge C_2(\kappa+1)$, and suppose
that the layered-space reading belongs to this class.
In particular the kept exponent is a finite quadratic
polynomial in the external field, with holomorphic
coefficients, and its estimate is
\begin{equation}\label{II-156-34}
 |\varphi_{\mathfrak k_X,X}(\mathsf a;p)|
 \le\tfrac12K_{\rm ch}C_{\rm xt}^2\mathfrak b_l^2
                        e^{-\kappa(d_t(X)+1)}
 \quad\text{if }|\mathsf a|\le C_{\rm xt}\mathfrak b_l.
\end{equation}
Here $K_{\rm ch}=C_3(LM)^6e^{C_3\kappa}B_{\rm ch}$.
We require also
\begin{equation}\label{eq:formed-pivot-radius}
 L\max\{\alpha_0,\alpha_1\}\le a_{\rm nf}.
\end{equation}

The pivot estimate applies also to coarse bonds crossing
the boundary of $\Omega_t$. We verify this point.
Let $c\in\Gamma_t$ and choose an endpoint $y_1\in\Omega_t$.
Separation gives
$\operatorname{dist}_\infty(\Omega_t,
             (\Omega^\diamond_{t-1})^c)>R_1\mathfrak M$.
If both endpoints lie in $\Omega_t$, the two averaging
blocks fit in a cube of side $2L\mathfrak M$ inside
$X(c)\cap\Omega_t$. The layer-$t$ local gauge gives
$L|A|\le\alpha_0$ and $L|A'|<\alpha_1$ there.
If the second endpoint is outside $\Omega_t$, the blocks
fit in a cube of side $2\mathfrak M$ inside $X(c)$.
Every point of this cube lies within
$2\mathfrak M<R_1\mathfrak M$ of $y_1$, so the cube
lies in the last two layers of $\Omega^\diamond_\bullet$.
The layer-$(t-1)$ gauge and the bounds in either layer give
$L|A|\le L\alpha_0$ and $L|A'|<L\alpha_1$.
In both cases the enclosing cube and its gauge are
contained in $X(c)\cap\Omega^\diamond_0$.

Put $a=L\max\{\alpha_0,\alpha_1\}$.
In this gauge the logarithmic potential of
$e^{i\xi A'}e^{i\xi A}$ has norm less than $4a/L$.
The averaging and pivot estimates used in
\eqref{eq:normalizer-pivot-log} give a matrix
$E=L^3m_c-1$ with
$|E|\le4C_{\rm piv}C_Qa/L\le1/2$.
The trace logarithm satisfies
$|\operatorname{tr}\log(1+E)|\le6|E|$.
Its exponential is the determinant and its imaginary
part has modulus at most $\pi/2$, so it is the specified
principal determinant logarithm. Thus, with a constant
$C_z$ depending only on $L$,
\begin{equation}\label{II-68-66}
 |z_c(p)-6\log L|\le C_za.
\end{equation}
For example $C_z=24C_{\rm piv}C_Q/L$ suffices.
The local fixed-gauge formula also proves holomorphy;
real gauge transformations preserve the determinant
and conjugation preserves its branch.

If an elimination factor is present, $X$ is one coarse
cube or two adjacent cubes, so $d_t(X)=0$.
Writing $N_X=\#\{c\in\Gamma^\partial_t:X(c)=X\}$,
we have $1\le N_X\le6M^3$: there are at most $2M^3$
coarse sites in $X$, with three positive bond directions.
Consequently
\begin{equation}\label{II-156-33}
 |\varphi^{\rm el}_{t,X}(P;\mathfrak q)-6N_X\log L|
                                      \le6M^3C_za_{\rm nf}.
\end{equation}

\begin{proposition}\label{II-157-13}
Under the class membership and form hypotheses above,
the factors \eqref{eq:formed-factor-definitions}, over
all steps of a history, form a finite family with
distinct names. Each exponent is defined throughout
its stated layered space and stored-data tiers,
is Borel measurable, holomorphic in the complex
coordinates of a fixed-real slice, gauge invariant,
conjugation real and dependent only on the pair.
The family at level $k+1$ retains every factor of
the history truncated at $k$, with the same domain
and exponent.

Define the oscillation at fixed stored data by
\[
 \operatorname{osc}\varphi
   =\sup_{\mathfrak q}\sup_{P,P'}
        |\varphi(P;\mathfrak q)-\varphi(P';\mathfrak q)|.
\]
Then
\begin{equation}\label{II-156-36}
 \sum_{f\text{ at }(t,X)}\operatorname{osc}\varphi_f
 \le B_{\rm fr}^{\flat}\max\{1,\mathfrak b_l^2\}
                                  e^{-\kappa d_t(X)},\qquad
 B_{\rm fr}^{\flat}
     =K_{\rm ch}C_{\rm xt}^2+12M^3C_za_{\rm nf}.
\end{equation}
If $\mathfrak b_l\ge1$, the maximum can be replaced
by $\mathfrak b_l^2$.
Each factor contains a cube of $Z_t$; at level $k>t$
its domain therefore meets $\Omega_{t+1}^c$.
\end{proposition}

\begin{proof}
The class membership gives definedness of the kept
quadratic coefficients and the pivot logarithms at
every pair under consideration. Their compositions
with the pair and stored-field readings give the
stated regularity. Those readings commute with gauge
transformations and conjugation, so the same is true
of the exponents. Equality of pairs gives equality
of the values, including equality of their domains.

A nonzero kept coefficient is assigned to a component
$\mathfrak k_X$ of $\Omega_t\cap Z_t$ and its domain
contains a cube of this component. For an elimination
factor the endpoints of its external coarse bond
lie outside $\Lambda_t$, so $X(c)\subset Z_t$.
If $t=k$ this places the factor at the boundary of
the current region. If $t<k$, separation of
$\Omega_{t+1}$ from $Z_t$ by at least $8LM$ puts
the same cube in $\Omega_{t+1}^c$.

The construction at step $t$ uses only the history
truncated at $t$. Its layered space uses the radii
at levels at most $t$, and its stored-data tiers
use only the external bonds of step $t$.
None changes on extending the history or extending
the table of radii. Thus old names, domains and
exponents agree exactly. At every step each domain
has at most the two distinct names in
\eqref{eq:formed-factor-definitions}; finiteness follows
from finiteness of the lattice and the number of steps.

For a fixed stored datum, its zero extension obeys
$|\mathsf a[\mathfrak q]|\le C_{\rm xt}\mathfrak b_l$.
Twice \eqref{II-156-34} bounds the oscillation of a
kept exponent by
$K_{\rm ch}C_{\rm xt}^2\mathfrak b_l^2e^{-\kappa d_t(X)}$.
For an elimination exponent, subtract the same
constant $6N_X\log L$ at both backgrounds and use
\eqref{II-156-33}. Its oscillation is at most
$12M^3C_za_{\rm nf}$ and its domain has size zero.
Adding the at most two bounds proves
\eqref{II-156-36}.
\end{proof}

\subsection{Radii independent of the number of steps}

The radii table associated with a coupling sequence is
$r_\iota(\sigma;n)=\alpha_{\iota,n}
 =C_\iota x_n^{-1/2}(\log x_n)^q$ for either stage
$\sigma$ containing level $n$.
Suppose $x_n\ge\gamma^{-2}$ and $\log\gamma^{-2}>2q$.
The function $F_q(x)=x^{-1/2}(\log x)^q$ is decreasing
on this interval, so
$r_\iota(\sigma;n)\le C_\iota F_q(\gamma^{-2})$.
The following conditions therefore ensure the required
ceiling bounds at every level:
\begin{equation}\label{II-157-4}
 \begin{gathered}
 \mathsf oL^3\ge3,\qquad a_\flat\le\alpha_0,\\
 (1+\beta_{\rm s})C_0F_q(\gamma^{-2})
                            \le\min\{\alpha_0,\alpha_2\},\\
 (1+\beta_{\rm s})C_1F_q(\gamma^{-2})\le\alpha_1,\qquad
 \mathsf oL^3MB^{(\rm cg)}C_0F_q(\gamma^{-2})\le\alpha_0.
 \end{gathered}
\end{equation}
Explicitly, these bound
$(1+\beta_{\rm s})r_0((j,{\rm out});n)$ by
$\min\{\alpha_0,\alpha_2\}$,
$(1+\beta_{\rm s})r_1((j,{\rm out});n)$ by $\alpha_1$,
and $\mathsf oL^3MB^{(\rm cg)}r_0((j,{\rm mid});j-1)$
by $\alpha_0$. Restriction of the sequence restricts
this same table; hence the factor families in
\cref{II-157-13} agree when consecutive levels are
described with their own truncated tables.

Likewise the condition
\begin{equation}\label{II-157-5}
 A_1(\log\gamma^{-2})^{\mathfrak p_0}\ge1
\end{equation}
gives $\mathfrak b_l\ge1$ at every level and turns
\eqref{II-156-36} into
\begin{equation}\label{II-157-15}
 \sum_{f\text{ at }(l+1,X)}\operatorname{osc}\varphi_f
       \le B_{\rm fr}^{\flat}\mathfrak b_l^2
                                      e^{-\kappa d_{l+1}(X)}.
\end{equation}
The constant depends on $L,M,\kappa,C_{\rm xt}$ alone.
It may be enlarged uniformly to
\[
 B_{\rm fr}=C_3(LM)^6e^{C_3(\kappa+1)}
 B_{\rm ch}C_{\rm xt}^2+12M^3C_za_{\rm nf}.
\]
This choice also accommodates the one-unit rate reserve
in the later locality estimates.


\section{Vacuum constants and their large-field deficits}
\label{sec:vacuum-bounds}

The field-independent part of each integration step has a
volume contribution and a correction supported near the
large-field region. We compute both from the same local
normalizer. In this section $N_j=\#\mathbb T^{(j)}$,
$y_j=\log x_j$, and $\#_j(Y)=\#(Y\cap\mathbb T^{(j)})$.
Complements and the operations $Y^\sim,Y^{\sim-1}$ are
taken in the cell partition at the indicated level.

\subsection{The constants of one step}\label{subsec:ymm1-62-1}

We first record the unnormalized value of the Gaussian
terms in \cref{sec:actual-normalizer}. At step $j=k+1$,
with its notation, put
\begin{equation}\label{II-45-8}
 \mathsf n^{(j)}(X,b)=
 \begin{cases}
 -\frac12 E^N(X;P_0,b),
     &b\in\mathfrak X^+,\ S_N(b)\subset X,\\
 6\log L,&b=b_0(c),\ X=X(c),\\
 0,&\text{otherwise},
 \end{cases}
 \qquad
 \mathsf n^{(j)}(X,z)=\sum_bh_z(b)\mathsf n^{(j)}(X,b).
\end{equation}
All the parameters and clearances of
\cref{sec:actual-normalizer} are retained. We may enlarge
its $L$-dependent constant $E_b$ so that $E_b\ge6\log L$,
and keep $E_N=3L^3E_b$. The proof of the normalizer gives
\begin{equation}\label{II-45-65}
 \mathsf n^{(j)}(X,z)\in\mathbb R,\qquad
 |\mathsf n^{(j)}(X,z)|
       \le E_N e^{-\delta_NM d_j(X)}.
\end{equation}
Indeed, $P_0$ belongs to every normalizer domain.
The nonpivot term is bounded by
$E_b e^{-\delta_NMd_j(X)}/2$ by the chain estimates in
that proof. In the pivot case the one or two cubes
of $X(c)$ have a common corner, so $d_j(X(c))=0$.
The tents are nonnegative and sum over bonds to $3L^3$.
These facts give the bound in \eqref{II-45-65}.
The real logarithm series and the real pivot value give
reality. The same change of coordinates used for the
normalizer transports these constants between local
lattice charts. In particular, their definition does
not introduce a choice of a larger torus.

Let the real local families $E^{(j)}$ of the progression
be defined on their analytic cores, including the
trivial triple, and suppose
\begin{equation}\label{eq:vacuum-local-input}
 |E^{(j)}(X,P,z)|\le E_0e^{-(\kappa+1)d_j(X)},\qquad
 \delta_NM\ge\kappa,\quad E_0\ge2E_N,\quad
 \kappa>\log9216+2.
\end{equation}
Here and below reality means conjugation reality, so
$E^{(j)}(X,\mathbf1,z)$ is real. Define the two local
vacuum families
\[
 \mathsf{vl}^{(j)}(X,z)=E^{(j)}(X,\mathbf1,z),
 \qquad
 \mathsf{vn}^{(j)}(X,z)=\mathsf n^{(j)}(X,z).
\]
Both are bounded by $E_0e^{-\kappa d_j(X)}$.

The three extensive coefficients are
\begin{equation}\label{II-192-10}
 \mathsf{co}^{\rm b}_j=\log\sigma_{\rm H}-\tfrac32y_j,
 \quad
 \mathsf{co}^{\rm s}_j=-(L^3-1)\log\mathsf z(x_j,\varepsilon_j),
 \quad
 \mathsf{co}^{\rm i}_j=\tfrac32\log(2\pi),
 \qquad \sigma_{\rm H}=(2\pi^2)^{-1}.
\end{equation}
Assume $x_j>1$, $0<\varepsilon_j\le1$ and
$x_j\varepsilon_j^2\ge1$ at all levels used here.
The last condition holds for
$\varepsilon_j=A_0x_j^{-1/2}y_j^{\mathfrak p_0}$
as soon as $A_0(\log\gamma^{-2})^{\mathfrak p_0}\ge1$
and $x_j\ge\gamma^{-2}$.
The radial Haar formula, with $w=|u-\mathbf1|$, is
\[
 \mathsf z(x,\varepsilon)=\frac2\pi\int_0^\varepsilon
      e^{-xw^2/2}w^2\sqrt{1-w^2/4}\,dw .
\]
It follows by writing a unit quaternion as
$(\cos\theta,\sin\theta\,\omega)$ and substituting
$w=2\sin(\theta/2)$. Extending the upper integral to
$[0,\infty)$, and restricting the lower integral to
$[0,x^{-1/2}]$, gives
\begin{equation}\label{eq:vacuum-haar-bound}
 \frac1{20}x^{-3/2}\le\mathsf z(x,\varepsilon)
      \le8x^{-3/2}.
\end{equation}
For the lower bound use
$e^{-xw^2/2}\ge e^{-1/2}$ and
$\sqrt{1-w^2/4}\ge\sqrt3/2$; the resulting coefficient
is $1/(\pi\sqrt{3e})>1/20$. The upper Gaussian integral
has coefficient $\sqrt{2/\pi}<8$.
Consequently, with
\begin{equation}\label{eq:vacuum-coefficient-constant}
 C_{\rm co}=\max\{\log(2\pi^2),(L^3-1)\log20,
                                  \tfrac32\log(2\pi)\},
\end{equation}
the coefficients in \eqref{II-192-10} satisfy
\[
 |\mathsf{co}^{\rm b}_j+\tfrac32y_j|\le C_{\rm co},
 \quad
 |\mathsf{co}^{\rm s}_j-\tfrac32(L^3-1)y_j|\le C_{\rm co},
 \quad |\mathsf{co}^{\rm i}_j|\le C_{\rm co}.
\]

For $\Omega\supset\Lambda$, unions of cells of level $j+1$,
let $B_j(Y)$ be the number of positive nonpivot
level-$j$ bonds of $Y$. Write
\begin{align}
 \mathsf N^{(j)}(\Omega,\Lambda)
  &=\mathsf{co}^{\rm b}_j B_j(\Omega)
       +\mathsf{co}^{\rm s}_j\#_{j+1}(\Omega)
       +\mathsf{co}^{\rm i}_j B_j(\Lambda),
                                                    \label{II-97-4}\\
 \mathsf E^{(j)}(\Omega,\Lambda)
  &=\mathsf N^{(j)}(\Omega,\Lambda)
    +\sum_{\substack{X\in\mathcal D_{j+1}\\X\subset\Lambda}}
       \left\{\sum_{z\in X\cap\Lambda^{\sim-1}}
                      \mathsf{vl}^{(j+1)}(X,z)
                   +\sum_{z\in X}\mathsf{vn}^{(j+1)}(X,z)\right\}.
                                                    \label{II-97-6}
\end{align}
The sites in these sums are level-$j+1$ sites.
The distinction between the two ranges of $z$ is essential:
the regular terms use the interior base points, whereas the
Gaussian constants use every site of $X$.
For a label $\mathcal S=(\Omega_n,\Lambda_n,S_n)_{n\le k}$
put
\begin{equation}\label{II-97-8}
 \begin{split}
 \mathcal E_{k,\mathcal S}
      &=-\sum_{j<k}\mathsf E^{(j)}(\Omega_{j+1},\Lambda_{j+1}),\\
 E_{(k)}&=\sum_{j<k}\mathsf E^{(j)}(T_\ell,T_\ell),\qquad
 \mathfrak E^{\ge}_k=\sum_{j=k}^{K-1}\mathsf E^{(j)}(T_\ell,T_\ell),\\
 \mathfrak E'_n(\mathcal S)
      &=\mathsf E^{(n-1)}(T_\ell,T_\ell)
                     -\mathsf E^{(n-1)}(\Omega_n,\Lambda_n).
 \end{split}
\end{equation}
The last-level sum $\mathfrak E^{\ge}_k$ presupposes that
the local families have been constructed through level $K$.

\subsection{Volume and deficit bounds}

Set
\[
 C_{{\rm E},p}=2048\sum_{q\ge0}(q+2)^pe^{-q},\qquad
 E_0'=C_{{\rm E},0}E_0,\qquad E_0''=8C_{{\rm E},1}E_0,
\]
and define
\begin{equation}\label{II-97-23}
 \begin{split}
 C_{\mathfrak E}&=\frac{(6L^3-5)C_{\rm co}+2E_0'}{L^3-1},\\
 C_{\log}&=3(L^3-1),\qquad
 C_{\rm loc}=(6L^3-5)C_{\rm co}+E_0'+2E_0'',\qquad
 C_\flat=\max\{C_{\log},C_{\rm loc}\}.
 \end{split}
\end{equation}

\begin{proposition}\label{II-192-21}
Assume the preceding normalizer and local-family hypotheses
at every level in question, and $m'<m+K-j$ whenever
$\mathcal D_j$ is used. Suppose the couplings satisfy
$x_j=L^{-1}x_{j-1}-b_j$, $|b_j|\le B_\sharp$,
and $x_K\ge\max\{4,2B_\sharp\}$.
All constants in \eqref{II-97-8} are real, and
\begin{equation}\label{II-192-25}
 E_{(k)}+\mathcal E_{k,\mathcal S}
          =\sum_{n=1}^k\mathfrak E'_n(\mathcal S).
\end{equation}
Writing $Z_n=\Lambda_n^c$, one has
\begin{equation}\label{II-192-27}
 |\mathfrak E'_n(\mathcal S)|
       \le C_\flat(1+y_{n-1})\,\#_n(Z_n^\sim).
\end{equation}
For the label with $\Omega_n=\Lambda_n=T_\ell$ every
deficit is zero. If the families are given through $K$,
then
\begin{equation}\label{II-192-29}
 E_{(K)}=E_{(k)}+\mathfrak E^{\ge}_k,\qquad
 -(3y_k+C_{\mathfrak E})N_k
       \le\mathfrak E^{\ge}_k\le C_{\mathfrak E}N_k .
\end{equation}
If in addition $\check R_n\ge y_n^r$, $r\ge2$, and
$\mathsf r_n=M^3\check R_n^4$, the cell count $|Z|_n$
satisfies
\begin{equation}\label{II-192-24}
 (1+y_{n-1})\#_n(Z^\sim)
    \le27(2+\log(2L))\,\mathsf r_n|Z|_n .
\end{equation}
All these constants are independent of the number of
steps, the torus and the couplings.
\end{proposition}

\begin{proof}
The entropy bound \eqref{II-18-27} gives, for every site $z$,
\[
 \sum_{\substack{X\ni z\\d_j(X)\ge D}}
       e^{-\kappa d_j(X)}(d_j(X)+2)^p
        \le C_{{\rm E},p}e^{-D}.
\]
Indeed, at size $q$ the product of entropy and decay is
at most $2048e^{-2q}$, and for $q\ge D$ this is at most
$2048e^{-D}e^{-q}$. Thus each of the local vacuum
families $v=\mathsf{vl}^{(j)},\mathsf{vn}^{(j)}$ obeys
\begin{equation}\label{eq:vacuum-site-cube}
 \sum_{X\ni z}|v(X,z)|\le E_0',\qquad
 \sum_{X\supset\square}\sum_{z\in X}|v(X,z)|
                                        \le M^3E_0''.
\end{equation}
For the second inequality use
$\#_j(X)\le8M^3(d_j(X)+1)$ and choose a site of the
fixed cube $\square$.

There are $3(L^3-1)$ positive nonpivot bonds per
coarse block: of its $3L^3$ outgoing bonds, precisely
three are pivots. Hence
$B_j(T_\ell)=3(L^3-1)N_{j+1}$.
A missing bond of a region is an inner bond of its
complement, by the bond convention. Its initial point
lies in a coarse block of that complement. Therefore
\begin{equation}\label{eq:vacuum-missing-bonds}
 0\le B_j(T_\ell)-B_j(Y)
                  \le3(L^3-1)\#_{j+1}(Y^c).
\end{equation}
The pivot assertion used in this count follows from
the corridor incidence: a pivot is a bond of $Y$
exactly when its associated coarse bond is a bond of $Y$.

On the whole torus the logarithmic coefficients combine to
$-3(L^3-1)y_jN_{j+1}$. The remaining coefficients and
the two sums in \eqref{eq:vacuum-site-cube} give
\[
 |\mathsf E^{(j)}(T_\ell,T_\ell)
                   +3(L^3-1)y_jN_{j+1}|
       \le[(6L^3-5)C_{\rm co}+2E_0']N_{j+1}.
\]
Now $\sum_{j=k}^{K-1}N_{j+1}\le N_k/(L^3-1)$.
Also $x_{j-1}=L(x_j+b_j)\ge Lx_j/2>x_j$:
backward induction starts from $x_K\ge2B_\sharp$.
Thus $0<y_j\le y_k$ for $j\ge k$. Summing the last
display proves \eqref{II-192-29}.

To estimate a deficit, put $\Omega=\Omega_n$,
$\Lambda=\Lambda_n$, and $Z=\Lambda^c$. Its logarithmic
part is
\[
 -\tfrac32y_{n-1}
  [B_{n-1}(T_\ell)-B_{n-1}(\Omega)
                       -(L^3-1)\#_n(\Omega^c)].
\]
By \eqref{eq:vacuum-missing-bonds}, its modulus is at most
$3(L^3-1)y_{n-1}\#_n(\Omega^c)$.
The three remaining coefficient contributions total at most
$(6L^3-5)C_{\rm co}\#_n(Z)$.

There are two reasons for a local term to be missing from
the restricted sum. Either $X\not\subset\Lambda$, or
$X\subset\Lambda$ and its base point lies outside
$\Lambda^{\sim-1}$. In the first case $X$ contains
a cube of $Z$. Summing the cube estimate in
\eqref{eq:vacuum-site-cube} over all such cubes bounds
each vacuum family by $E_0''\#_n(Z)$.
The second case occurs only for $\mathsf{vl}$ and has
$z\in Z^\sim\setminus Z$; the site estimate bounds it
by $E_0'[\#_n(Z^\sim)-\#_n(Z)]$.
These two cases are disjoint. Since
$\Omega^c\subset Z\subset Z^\sim$, their sum proves
\eqref{II-192-27} with \eqref{II-97-23}.
Identity \eqref{II-192-25} is the telescoping sum of
the defining deficits; reality follows from that of
each coefficient and local value.

Finally $x_{n-1}\le(3L/2)x_n\le2Lx_n$, so
$1+y_{n-1}\le(2+\log(2L))\check R_n$ when
$y_n>1$ and $\check R_n\ge y_n^r$.
A cell has at most $27$ touching cells, including
itself. Hence $\#_n(Z^\sim)\le27\#_n(Z)$, while
$\#_n(Z)=M^3\check R_n^3|Z|_n$. This proves
\eqref{II-192-24}.
\end{proof}


\section{The sum of the regular and boundary interactions}
\label{sec:scale-sum}

A uniform bound on each local interaction would give a
contribution at every finer lattice site. The regularity
of a field in a later layer is stronger. Gauge invariance
makes its local interaction quadratic at the trivial field;
the resulting fourth power of the scale ratio compensates
for the number of finer sites. We establish this estimate
before summing over the layers.

Use the constants of \cref{I-93-98}, and assume
\[
 M\alpha_0\le\alpha_1,\qquad
 \alpha_0+\alpha_1\le a^\sharp,\qquad
 \kappa\ge\log9216+2.
\]
In this section \(1\le j\le k\le K\),
\(m'<m+K-k\), and all local families have the fixed
radii \((\alpha_0,\alpha_1)\). Write
\[
 \rho_{\rm sc}=\frac{\alpha_0}{16K^\sharp_{\rm tw}},
 \qquad
 C_\natural^{\rm sc}
 =8\max\{16K^\sharp_{\rm tw}B_*,c_*^{-1}\}^{\,2},
\]
where \(B_*,c_*\) are the regional-gauge constants in
\cref{I-28-75}, with \(c_*\le1\).
The constants below are independent of the depth and
volume and of the fixed radii.

\subsection{The quadratic gain near a regular field}

\begin{lemma}\label{II-101-27}
Use level-\(j\) lengths. For \(z\in X^{(j)}\), put
\[
 d=d_j(X),\qquad l=2M(d+2)+1,\qquad
 Q=N_j(\{z\};l/2),\quad
 \widetilde Q=N_j(\{z\};l/2+2).
\]
Suppose \(l+4<L^{m+K-j}\), and let a real field \(U\)
have curvature less than \(\theta\eta_j^2\) and current
less than \(\theta\) on \(\widetilde Q\), including every
plaquette with a corner there and every bond with an
endpoint there. If its pair on \(X\) belongs to the
analyticity core, then each local term \(E\) satisfies
\begin{equation}\label{II-101-29}
 |E(U)-E(\mathbf1)|
 \le C_\natural^{\rm sc}\|E\|_X
                       (\theta/\alpha_0)^2l^2 .
\end{equation}
\end{lemma}

\begin{proof}
The spanning-edge geometry gives
\(|x-z|_\infty\le M(d+2)\) for \(x\in X\).
Thus \(Q\) contains both endpoints of every bond of
\(\operatorname{Bd}(X)\). It is a cube of odd integer
side formed from level-\(j\) unit cubes, and
\(\widetilde Q\) lies in one torus chart.

First suppose
\(2l\theta\le c_*\) and
\(a=2B_*l\theta\le\rho_{\rm sc}\).
The current-to-divergence identity used in
\cref{I-28-75} gives, with \(\eta=\eta_j\le1/13\),
\[
 |D_U^{\mathsf T}\partial U|
 \le\eta^{-1}(\theta\eta^2)^2+\eta^2|J^{(j)}(U)|
 <2\theta\eta^2.
\]
Here \(\theta\le1/2\). Apply the regional gauge theorem
with radius \(2\theta\). On the bonds meeting \(Q\)
it gives \(U^u=e^{i\eta A}\), with
\(|A|,|\nabla^\eta A|<a\).
Restrict to \(\operatorname{Bd}(X)\), put
\(J_1=\operatorname{Ad}(u)J^{(j)}(U)\), and consider
\[
 \gamma(\zeta)=(e^{i\zeta\eta A},\zeta J_1),
       \qquad |\zeta|\le R:=\rho_{\rm sc}/a .
\]
We have \(|J_1|<\theta\le a\) and \(R\ge1\).

This entire disc belongs to the fixed core. To see the
margin explicitly, take the real reference field to be
one. For \(\phi=|\zeta|a\le\rho_{\rm sc}\le1/16\),
the four exponential plaquette factors give
\[
 |\partial(e^{i\zeta\eta A})-1|
 \le(2\phi+11\phi^2)\eta^2<4\rho_{\rm sc}\eta^2.
\]
The linear part is the curl of \(A\); the remaining
terms are bounded by the exponential-product remainder.
The perturbation and its first differences are at most
\(\rho_{\rm sc}<2\rho_{\rm sc}\), and the current is
smaller than \(4\rho_{\rm sc}\). The real reference
has trivial cube gauges. These are the simple-domain
conditions at radii \((4\rho_{\rm sc},2\rho_{\rm sc},
4\rho_{\rm sc})\). Moreover
\[
 6\rho_{\rm sc}\le a^\sharp,\quad
 4\rho_{\rm sc}\le\alpha_0,\quad
 2\rho_{\rm sc}\le\alpha_1,\quad
 K^\sharp_{\rm tw}(4\rho_{\rm sc}+4\rho_{\rm sc}^2)
                                  \le\alpha_0/2 .
\]
Thus \cref{I-93-98} and radius monotonicity put the
disc in the core at \((\alpha_0,\alpha_1)\).

The function \(f(\zeta)=E(\gamma(\zeta))\) is
holomorphic near the closed disc and bounded by
\(\|E\|_X\) there. Its derivative at zero vanishes.
Indeed conjugation by the four constant quaternions
\(1,i\sigma_1,i\sigma_2,i\sigma_3\) preserves \(E\).
Their adjoint actions sum to zero on
\(\mathfrak{su}(2)^\mathbb C\). Apply the differential
of \(E(e^{i\eta A},J)\) at \((0,0)\) to the sum of
these four conjugates of \((A,J_1)\); it equals both
zero and \(4f'(0)\).
The maximum principle for
\((f(\zeta)-f(0))/\zeta^2\) gives
\[
 |E(U)-E(\mathbf1)|
   \le2\|E\|_X(a/\rho_{\rm sc})^2 .
\]
To apply gauge invariance here, both representatives must
belong to the core. The triple representing \(\gamma(1)\)
lies in the simple domain at radii
\((4\rho_{\rm sc},2\rho_{\rm sc},4\rho_{\rm sc})\).
Its transform by \(u^{-1}\) lies in the same domain:
unitary conjugation preserves the curvature, perturbation,
and current bounds, and a cube gauge \(w\) is replaced
by \(wu\), with the same gauge-fixed field.
The preceding application of \cref{I-93-98} therefore
puts both triples in the fixed core. Gauge invariance
and pair dependence now identify \(E(U)\) with \(f(1)\).

If either smallness inequality fails, use instead
\(|E(U)-E(\mathbf1)|\le2\|E\|_X\).
In the case \(2l\theta>c_*\), multiply this bound by
\((2l\theta/c_*)^2>1\); in the case \(a>\rho_{\rm sc}\),
multiply it by \((a/\rho_{\rm sc})^2>1\).
All three cases are bounded by
\[
 2\|E\|_X(2l\theta)^2
        \max\{B_*/\rho_{\rm sc},c_*^{-1}\}^2.
\]
Since \(\alpha_0\le1\), this proves \eqref{II-101-29}.
\end{proof}

\subsection{Summation at a site and over the layers}

Fix a history of length \(k\), with
\(\Omega_0=\mathbb T\), \(\Omega_{k+1}=\varnothing\),
and layers \(\mathcal L_n=\Omega_n\setminus\Omega_{n+1}\)
in the level-\(n\) unit partition.
Let \(\mathsf U\) be a real bare field with
\begin{equation}\label{II-101-13}
 \operatorname{rad}_{\mathsf U}(\mathfrak e)
      <C_{\rm ly}\varepsilon_{\lambda(\mathfrak e)}
 \quad\text{on the collared region }\Omega_0^\diamond,
 \qquad C_{\rm ly}>0.
\end{equation}
Use the curvature and current radii of
\cref{sec:history-layer-regularity}. Assume its pair belongs
to the fixed core for every \(X\subset\Lambda_j\),
including the domains used by the edge terms.
The coupling and geometric conditions used here are
\begin{equation}\label{eq:scale-sum-running}
 \begin{gathered}
 |b_j|\le B_\sharp,\qquad
 \varepsilon_q\le\varepsilon_n\quad(q\le n),\qquad
 \varepsilon_q\le\sqrt2L^{(q-n)/2}\varepsilon_n
                                     \quad(q\ge n),\\
 \check R_n\ge L,\qquad
 M\check R_n\le L^{m+K-n},\qquad
 C_{\rm ly}\varepsilon_k\le\alpha_0 .
 \end{gathered}
\end{equation}
These are the running bounds already used for the history
backgrounds, together with the displayed fixed-radius ceiling.

Let \(\|E^{(j)}\|_{j,\kappa+1}\le E_0\).
For each permitted created boundary index, take a number
\(\mathsf b_j(X)\) satisfying
\(|\mathsf b_j(X)|\le\mathsf B e^{-\kappa d_j(X)}\).
Such numbers are the evaluations of the boundary entries
when their fields lie in the stated domains.
Use the edge sums \eqref{II-153-12}, and define
\begin{equation}\label{II-101-17}
 \begin{split}
 \mathbf E_k(\mathsf U)
 &=\sum_{j=1}^k
   \left\{\sum_{\substack{X\subset\Lambda_j\\
                       z\in X^{(j)}\cap\Lambda_j^{\sim-1}}}
     [E^{(j)}(X,\mathsf U,z)-E^{(j)}(X,\mathbf1,z)]
                    -b_jL^jA(\varphi_j,\mathsf U)\right\},\\
 \mathbf E_k^{\rm bd}(\mathsf U)
 &=\sum_{j=1}^k
   \left\{\sum_{X\in\mathsf{Ix}^{\rm mk}_j}\mathsf b_j(X)
       +\sum_{X\in\mathsf{Ix}^{\rm ed}_j}
                                  E^{(j),{\rm ed}}(X,\mathsf U)\right\}.
 \end{split}
\end{equation}
Here \(\varphi_j\) is the smoothed interior profile:
\(0\le\varphi_j\le1\), and it vanishes at distance at
most one quarter of a level-\(j\) cell side from \(Z_j\).
Only these two properties of the profile enter this estimate.

For the entropy constants \(C_{{\rm E},p}\) from
\cref{sec:vacuum-bounds}, set
\begin{equation}\label{II-101-8}
 \begin{gathered}
 C'_{\rm f}=110592,\qquad
 C_{\rm sh}^{\rm sc}
       =18C_\natural^{\rm sc}C_{{\rm E},2}M^2L^4,\qquad
 C_{\rm a}^{\rm sc}=C_{\rm sh}^{\rm sc}+2C_{{\rm E},0},\\
 C_{\rm m}^{\rm sc}
       =\tfrac{13}{12}C_{\rm a}^{\rm sc}
                            +432C'_{\rm f}+\tfrac{13}4,\qquad
 C_5=C_{\rm m}^{\rm sc}(E_0+B_\sharp+M^{-3}\mathsf B).
 \end{gathered}
\end{equation}

\begin{proposition}\label{II-101-34}
Under these assumptions, both sums in \eqref{II-101-17}
are defined, and the sum of the moduli of all their
individual summands is bounded by
\begin{equation}\label{II-101-44}
 C_5\sum_{n=0}^k\#_n(\mathcal L_n)
 \le C_5\left(N_k+L^3\sum_{j=1}^k\#_j(Z_j)\right).
\end{equation}
In particular this bounds
\(|\mathbf E_k(\mathsf U)|+|\mathbf E_k^{\rm bd}(\mathsf U)|\).
\end{proposition}

\begin{proof}
The trivial pair is in every core by \cref{I-93-98}.
The membership assumption therefore makes all local and
edge terms defined.

Fix \(z\in\Lambda_j\cap\mathcal L_n\), \(j\le n\), and put
\[
 n'=\max\{n-1,j\},\quad w=M\check R_nL^{n-j},\quad
 \theta=\sqrt2C_{\rm ly}\varepsilon_{n'}L^{-2(n'-j)}.
\]
We use level-\(j\) lengths. The cell-size condition gives
\(w\le L^{m+K-j}\), half the torus side.
The ball of radius \(2w\) about \(z\) lies in
\(\Omega_{n'}\): for \(n=j\) use the \(2M\check R_j\)
gap from \(\Lambda_j\) to \(\Omega_j^c\); for \(n>j\)
use the \(8M\check R_nL^{n-j}\) gap from \(\Omega_n\)
to \(\Omega_{n-1}^c\).
On a neighborhood of radius at most \(2w-1\), all
plaquette and current stencils thus lie in
\(\Omega_0^\diamond\) and have layer index
\(\lambda\ge n'\).
Equation \eqref{II-101-13} and the coupling ratio imply
\[
 |\partial\mathsf U-1|
 <\sqrt2C_{\rm ly}\varepsilon_{n'}L^{-2n'}
       =\theta\eta_j^2,\qquad
 |J^{(j)}(\mathsf U)|
 <\sqrt2C_{\rm ly}\varepsilon_{n'}L^{-3(n'-j)}
       \le\theta .
\]
The extra factors \(L^{-3(\lambda-n')/2}\) and
\(L^{-5(\lambda-n')/2}\) have been bounded by one.

Call \(X\ni z\) short if \(2M(d_j(X)+2)+5<w\).
The cubes in \cref{II-101-27} then fit into the preceding
regularity neighborhood and into a torus chart.
Since \(l\le3M(d+2)\), that lemma bounds a short term by
\[
 18C_\natural^{\rm sc}M^2L^4 E_0
   (C_{\rm ly}\varepsilon_{n'}/\alpha_0)^2
        L^{-4(n-j)}(d+2)^2e^{-(\kappa+1)d}.
\]
The entropy estimate in \eqref{eq:vacuum-site-cube}
and its preceding weighted bound sum the last two
factors to at most \(C_{{\rm E},2}\).
For a long domain,
\(d\ge\frac12\check R_nL^{n-j}-3>0\).
Use the bound \(2E_0e^{-(\kappa+1)d}\) and the same
entropy tail to obtain at most
\(2e^3C_{{\rm E},0}E_0e^{-\check R_nL^{n-j}/2}\).
Writing \(s=L^{n-j}\ge1\), the inequalities
\(\check R_n\ge13\) and
\(e^{-13s/4}\le24(4/13)^4s^{-4}<s^{-4}\) give
\[
 e^{-\check R_ns/2}
       \le e^{-\check R_n/4}s^{-4}.
\]
The fixed-radius ceiling and \(e^{3-\check R_n/4}\le1\)
therefore yield the site estimate
\begin{equation}\label{II-101-35}
 \sum_{X:\ z\in X\subset\Lambda_j}
 |E^{(j)}(X,\mathsf U,z)-E^{(j)}(X,\mathbf1,z)|
       \le C_{\rm a}^{\rm sc}E_0L^{-4(n-j)}.
\end{equation}
There are \(L^{3(n-j)}\#_n(\mathcal L_n)\)
level-\(j\) sites in \(\mathcal L_n\). Summing gives
\(C_{\rm a}^{\rm sc}E_0
\sum_{n=j}^kL^{-(n-j)}\#_n(\mathcal L_n)\).

For the action term, a contributing plaquette has all
corners in \(\Lambda_j\). Indeed its center is farther
than a quarter-cell side from \(Z_j\), whereas its
corners are only a half bare spacing from its center.
If its base lies in \(\mathcal L_n\), then
\(\lambda(p)\ge n\), and the layer bound gives
\[
 a_p(\mathsf U)=\tfrac12|\partial\mathsf U(p)-1|^2
       <(C_{\rm ly}\varepsilon_n)^2L^{-4n}.
\]
There are three orientations and
\(L^{3n}\#_n(\mathcal L_n)\) bare base points.
Consequently
\begin{equation}\label{II-101-40}
 |b_j|L^jA(\varphi_j,\mathsf U)
 \le3B_\sharp\sum_{n=j}^k
 L^{-(n-j)}(C_{\rm ly}\varepsilon_n)^2
                                      \#_n(\mathcal L_n).
\end{equation}

Every created-boundary domain has a cube in \(Z_j^\sim\),
since it is not contained in
\(\Lambda_j^{\sim-1}=(Z_j^\sim)^c\).
The edge indices are among these indices. For a fixed
cube, \eqref{II-18-29} bounds the sum of
\(e^{-\kappa d}\) by \(4096\). Even allowing the
27 cubes touching each cube of \(Z_j^\sim\), summation
therefore gives
\[
 \sum_{X\in\mathsf{Ix}^{\rm mk}_j}e^{-\kappa d_j(X)}
       \le C'_{\rm f}M^{-3}\#_j(Z_j^\sim).
\]
The edge bound \eqref{II-153-13} and
\(\#_j(Z_j^\sim)\le27\#_j(Z_j)\) bound the boundary
contribution by
\[
 27C'_{\rm f}(M^{-3}\mathsf B+8E_0)\#_j(Z_j).
\]

To sum these bounds, put
\(\mathsf N=\sum_{n=0}^k\#_n(\mathcal L_n)\).
The regular and action sums use
\(\sum_{a\ge0}L^{-a}\le13/12\) and
\(C_{\rm ly}\varepsilon_n\le\alpha_0\le1\).
For the boundary sum note that \(x\in Z_j\) implies
that its layer index is at most \(j\), since
\(\Omega_{j+1}\subset\Lambda_j\).
Counting bare sites and rescaling gives
\[
 \#_j(Z_j)\le
     \sum_{n=0}^jL^{-3(j-n)}\#_n(\mathcal L_n),
 \qquad
 \sum_{j=1}^k\#_j(Z_j)\le2\mathsf N .
\]
Thus the sum of moduli is at most
\[
 \left[
  \left(\tfrac{13}{12}C_{\rm a}^{\rm sc}+432C'_{\rm f}\right)E_0
  +\tfrac{13}4B_\sharp+54C'_{\rm f}M^{-3}\mathsf B
 \right]\mathsf N\le C_5\mathsf N.
\]
Finally \(\mathcal L_n\subset\Omega_{n+1}^c\subset Z_{n+1}\)
at bare sites for \(n<k\), so
\(\#_n(\mathcal L_n)\le L^3\#_{n+1}(Z_{n+1})\);
the top count is at most \(N_k\). These inequalities
prove the second part of \eqref{II-101-44}.
\end{proof}

\section{Integrable bounds for the large-field components}
\label{sec:large-field-majorant}

A failed test must force a definite contribution either to the
Wilson action, to the gauge-fixing exponent, or to a retained
quadratic form. The shares of \cref{II-21-30} prevent the same
action from being used at different steps. Within one step,
a finite covering multiplicity allows us to sum these lower
bounds over the newly large cells.


Use the failed-test estimates of \cref{II-9-27,II-9-33},
with their local averaging ties and curvature ceilings.
Write \(c_{\rm P}=c_1^{\mathrm P}=\frac14B_3^{-2}\)
and \(c_{\rm Q}=c_2^{\mathrm Q}\), where the latter
and \(\sigma_{\rm Q}\) are defined in \eqref{II-9-10}.
The extended gauge energy is \eqref{II-9-3}.
The weighted costs are \eqref{II-9-29} and \eqref{II-9-35};
the failed density test at level zero costs
\(\frac12\varepsilon_0^2\) in the Wilson action.

\subsection{Averages on the reserved regions}

\begin{proposition}\label{II-77-10}
Let the collared background of a history \(h\) exist at
the data of a bare field \(U\), and let
\(\mathsf U=U_{\rm el}\vee U\) be any glued representative.
For every piece \((t,D)\), the averages of orders \(t-1\)
and \(t\) of \(\mathsf U\) equal those of \(U\) on the
inner bonds of \(N_{t,D}\) at their respective levels.
On any region \(N\) of level-\(t\) cubes contained in \(N_{t,D}\),
these averages on the inner bonds of \(N\) can be computed
from the restriction of \(\mathsf U\) to the inner bare bonds of \(N\).
In particular these assertions hold on
\(\square'^{\sim4}\) for every failed \(P_t\)- or
\(Q_t\)-cube \(\square'\subset D\).
If \(\check R_0\le L\check R_1\), the glued field also
equals \(U\) on \(\square^\sim\) for each density cube
\(\square\subset P_0\cap D\).
\end{proposition}

\begin{proof}
For \(t\ge2\), \cref{II-21-30} places \(N_{t,D}\) inside
\(\Omega_{t-1}\) and outside \(\Omega_t\).
Its level-\((t-1)\) bonds are therefore among the collared
constraints \(\Gamma^\diamond_{t-1}\). The constrained
average of \(U_{\rm el}\) equals that of \(U\).
All bare bonds needed for this average lie inside the
collar; gluing changes none of them.
For \(t=1\), an inner bare bond of \(N_{1,D}\) is either
outside the collar, where gluing leaves \(U\) unchanged,
or belongs to \(\Gamma^\diamond_0\), where the constraint
gives the same conclusion.

The averaging blocks centered at sites of a region of
level-\(t\) cubes remain in that region. The double block
of every inner level-\(t\) bond therefore lies in
\(N_{t,D}\), so one more averaging step proves the
order-\(t\) assertion. The same block property proves
the assertion for a restriction.
The four-layer enlargement of a test cube has width
\(4M_2\check R_t a_{K-t}\le s_t\), so lies in its share
region by \cref{II-21-30}. For a density cube the width
is \(M_2\check R_0a_K\le M\check R_1a_{K-1}=s_1\).
The order-zero assertion then identifies the bare fields
and their plaquette actions there.
\end{proof}

\subsection{The failed tests near a created cell}

The created cells at step \(t\) are the cells of
\[
 \mathcal W_t^\wedge
   =P_t'\cup Q_t'\cup R_t'\cup\mathsf K(I_t)
                    \cup\mathsf R_t^{\rm lo},
\]
with \(P_0'\) added when \(t=1\).
Write \(\operatorname{Cr}(t,D)\) for the pairs
\((t,\mathfrak c)\) with \(\mathfrak c\subset D\) a created
cell, and \(\operatorname{Cr}(\mathfrak u)\) for their union
over a set of pieces \(\mathfrak u\).
By \eqref{II-4-27}, \(\mathfrak c^\sim\) lies in \(Z_t\)
for every created cell, hence in its unique component \(D\).

For a piece \((t,D)\), with \(j=t-1\), set
\[
 \mathsf Y(t,D)=P_t''\cap D\cap\Omega_t^c,\qquad
 n^{\rm g}_{t,D}=(L^3-1)\#\mathsf Y(t,D).
\]
Let \(\mathfrak B(t,D)\) be the level-\(j\) bonds whose
endpoints lie in one block \(B(y)\), \(y\in\mathsf Y(t,D)\),
and let \(\mathfrak X^{\rm st}(t,D)\) be the external
nonpivot bonds of step \(t\) with both endpoints in \(D\).
The variables of this piece range over
\begin{equation}\label{II-74-8}
 \mathsf X^{\rm fr}_{t,D}
 =\mathcal V(\mathfrak B(t,D))
               \times\mathfrak g^{\mathfrak X^{\rm st+}(t,D)},
\end{equation}
with product Haar and Lebesgue measure \(m_{t,D}\).
At a bare field \(U\), the entries are the restriction of
\(V_j=\mathcal M^jU\) and the stored field \(\mathsf a'_j\).

\begin{lemma}\label{II-76-31}
Suppose the actual history indicator is one, the actual
backgrounds and strip substitutions exist, and the
hypotheses of \cref{II-26-71} hold for these substitutions.
If \(t=1\), assume also \(\check R_0\le L\check R_1\).
For a created cell \(\mathfrak c\subset D\), one of the
following tests fails in \(\mathfrak c^{\sim2}\).
\begin{enumerate}
\item If \(\mathfrak c\subset P_t'\) or \(Q_t'\), there is
a failed test cube \(\square'\) in the corresponding set,
contained in a cell \(\mathfrak c''\) equal to or touching
\(\mathfrak c\), with
\[
 \square'^{\sim4}\subset(\mathfrak c'')^\sim
               \subset N_{t,D}\cap\mathfrak c^{\sim2}.
\]
In the \(Q_t'\) case, \(Y_{\square'}\subset\mathsf Y(t,D)\).
\item If \(t=1\) and \(\mathfrak c\subset P_0'\), a failed density cube has
\(\square^\sim\subset N_{1,D}\cap\mathfrak c^{\sim2}\).
\item Otherwise there is a positive stored bond \(b\)
with endpoints in \(\mathfrak c^{\sim2}\) such that its
new variable is undefined or
\begin{equation}\label{II-76-33}
 |\mathsf a'_j(b)|\ge
       \varkappa_{\rm w}(A_1/A_0)p_0(g_j),\qquad
 \varkappa_{\rm w}=\min\{(4\Theta_1)^{-1},\theta_c\}.
\end{equation}
\end{enumerate}
\end{lemma}

\begin{proof}
A cell of the hull of \(P_t\) or \(Q_t\) meets one of
its test cubes. The cell \(\mathfrak c''\) containing
that cube touches \(\mathfrak c\); its one-cell enlargement
is in the same component and in \(N_{t,D}\).
The test-cube enlargement has width at most \(s_t\),
so lies there and in \(\mathfrak c^{\sim2}\).
For \(Q_t\), the cube is selected in
\(B(P_t'')^{\sim-1}\); hence
\((\mathfrak c'')^\sim\subset B(P_t'')\cap\Omega_t^c\cap D\).
This contains \(Y_{\square'}\).
The argument for \(P_0\) is the same, using
\(M_2\check R_0a_K\le s_1\).
The factors of the history indicator say precisely that
these tests fail.

For a cell in \(R_t'\), choose a failed cube
\(\square'\subset R_t\) meeting it and test the bonds of
\((\square'^{\sim2})^{(j)*}\). A cell of \(\mathsf K(I_t)\)
contains an endpoint of a bond in \(I_t\); use such a bond.
For \(\mathsf R_t^{\rm lo}\), use the bonds of its failed
rim test. Each tested bond has both endpoints in
\(\mathfrak c^\sim\subset D\), and is a nonpivot bond
incident to \(\Omega_t\). For the cube test this follows
from
\[
 (L+2)M_2\check R_ta_{K-t}+a_{K-j}\le s_t;
\]
for the other two tests it follows from
\(a_{K-j}\le s_t\).
Thus all these bonds are stored. They belong to
\(\mathcal N_{\rm s}\cup\mathcal N_\ell\), so none is
in the plain zone \(\mathcal Z_1\); a bond in \(I_t\)
belongs to \(\mathcal N_{\rm s}\) and is not in
\(\mathcal Z_\theta\).

Failure gives an undefined old logarithm or one of norm
at least \(\delta_j\), \(\delta_j/(4\Theta_1)\), or
\(\theta_c\delta_j\), respectively.
Outside the strip zone the variable is unchanged.
In the cutoff zone the inverse change of variables
preserves each threshold at most \(4\delta_j/3\), by
\cref{II-26-71}. If the new logarithm is defined, the
old logarithm cannot be undefined or above such a
threshold while the new one is below it: below the
threshold the change is the identity. The only excluded
case is the point \(-1\) of the new logarithmic chart.
Divide by \(g_j\), choose the positive orientation,
and use \(\delta_j/g_j=(A_1/A_0)p_0(g_j)\).
\end{proof}

\subsection{The elementary integrals}

For \(x\varepsilon^2\ge1\), the normalizer estimates in
\cref{sec:vacuum-bounds} give
\[
 \mathsf z_\infty(x/2)
    \le8\,2^{3/2}x^{-3/2}
    \le500\,\mathsf z(x,\varepsilon).
\]
Here \(\mathsf z_\infty\) is the same single-group integral
with its indicator removed.
For any set \(Y\) of level-\((j+1)\) sites, define
\[
 \mathsf g_Y(V)=
   \mathsf z(x,\varepsilon)^{-(L^3-1)\#Y}e^{-xG_j^\infty(Y;V)} .
\]
Let \(\mathsf g_Y^{1/2}\) denote the same expression with
half the exponent and the full normalization.

\begin{lemma}\label{II-7-13}
The functions just defined are Borel and
\begin{equation}\label{II-7-14}
 \int\mathsf g_Y^{1/2}(V)\,dV
                 \le500^{(L^3-1)\#Y}.
\end{equation}
There is also a number
\(\varepsilon_{\mathsf z}=r_*^2/(20\pi L^3)>0\),
depending only on \(L\), such that curvature at most
\(\varepsilon_{\mathsf z}\) on a double block implies
\begin{equation}\label{II-7-18}
 \mathsf m_c(V)\text{ is invertible},\qquad
 |\det\mathsf m_c(V)|^{-1}\le2L^9 .
\end{equation}
Here \(r_*\) is the common joint analytic radius in
\cref{I-10-93}.
\end{lemma}

\begin{proof}
The contour means and their domains are Borel. Average
over gauge transformations equal to one at the centers
\(y\in Y\), with independent Haar values at every other
site in the blocks. The domain where all contour means
exist is gauge invariant, and each contour variable
becomes \(V(y,x)u(x)^{-1}\). Haar invariance and Tonelli
therefore make its integral the measure of that domain,
at most one, times
\[
 \left[\frac{\mathsf z_\infty(x/2)}
                 {\mathsf z(x,\varepsilon)}\right]^{(L^3-1)\#Y}.
\]
This proves \eqref{II-7-14}, without assuming existence
of the means on a set of full measure.

For the pivot bound, a perturbation supported on the pivot
bond changes no contour variable inside either block.
Thus the joint-radius assertion of \cref{I-10-93} applies to
\[
 F_c(a^0,X)=\frac1i\log\left[
  \mathcal M(e^{iA_{c,X}}e^{ia^0})(c)
                \mathcal M(e^{ia^0})(c)^{-1}\right],
 \qquad |F_c|\le1
\]
for \(\sup|a^0|<r_*, |X|<r_*\), and
\(D_XF_c(a^0,0)=L\mathsf m_c(e^{ia^0})\).
Apply Schwarz in \(a^0\) and Cauchy in \(X\) to
\(F_c(a^0,X)-F_c(0,X)\). Since
\(\mathsf m_c(1)=L^{-3}\), this gives
\[
 |\mathsf m_c(e^{ia^0})-L^{-3}|
             \le\frac{2\sup|a^0|}{Lr_*^2}.
\]
The comb gauge on a double block of side at most \(2L\)
gives \(V^u=e^{ia^0}\) with
\[
 \sup|a^0|<2\pi L\varepsilon_{\mathsf z}
              =r_*^2/(10L^2)<r_*.
\]
It follows that \(|L^3\mathsf m_c(V^u)-1|<1/5\), so
\(|\det\mathsf m_c(V^u)|>(4/5)^3L^{-9}\).
Differentiating gauge covariance of the average shows
that \(\mathsf m_c(V^u)\) is obtained from
\(\mathsf m_c(V)\) by orthogonal adjoint maps on its two
ends. Their determinant moduli are one.
Thus the same lower bound holds for
\(|\det\mathsf m_c(V)|\), proving \eqref{II-7-18}.
\end{proof}

\subsection{The product bound}

Fix \(0<\gamma_{\rm kp}\le1\). At each piece let
\(\mathscr Q_{t,D}\) be a real quadratic form satisfying
\begin{equation}\label{II-77-25}
 \mathscr Q_{t,D}(a')\ge
                2\gamma_{\rm kp}\sum_{b\in\mathfrak X^{\rm st+}(t,D)}
                                                |a'(b)|^2 .
\end{equation}
Put \(N_D=\#_t(D)\), and use \(x=x_j\), \(\varepsilon=\varepsilon_j\)
and \(Y=\mathsf Y(t,D)\) in the preceding gauge weights.
Define
\begin{equation}\label{II-77-30}
 \overline{\mathsf u}_{t,D}(V,a')
  =(2L^9)^{6N_D}\mathsf g_Y^{1/2}(V)
       \exp\left(-\frac{\gamma_{\rm kp}}4\sum_b|a'(b)|^2\right).
\end{equation}
The constants used below are
\begin{equation}\label{II-77-16}
 \begin{split}
 c_\star={}&5^{-3}\min\left\{
  \tfrac12,\tfrac12c_{\rm P},
  \tfrac12c_{\rm Q}(A_1/A_0)^2,
  \tfrac14\gamma_{\rm kp}\varkappa_{\rm w}^2(A_1/A_0)^2\right\},\\
 C_\star'={}&6\log(2L^9)+L^3
       \left(\log500+\tfrac92\log(4\pi/\gamma_{\rm kp})\right).
 \end{split}
\end{equation}

\begin{proposition}\label{II-77-63}
Assume the hypotheses of \cref{II-76-31}, existence of the
collared background, and defined stored variables. Suppose
\[
 \begin{gathered}
 |b_i|\le B_\sharp\ (1\le i\le k),\qquad
 x_k\ge2B_\sharp,\qquad p_0(g_k)\ge1,\\
 \check R_l\le L\check R_{l+1}\ (0\le l<k),\\
 \varepsilon_t\le B_3\bar\varepsilon,\quad
 c_{\rm P}\varepsilon_t^2\le e_*\ (1\le t\le k),\quad
 c_{\rm Q}\delta_j^2\le\sigma_{\rm Q}\ (0\le j<k).
 \end{gathered}
\]
Let \(\mathsf U\) be any glued representative and suppose
its level-\(j\) curvature is at most
\(\varepsilon_{\mathsf z}\) on the double blocks of a set
\(\Gamma'(t,D)\) of positive coarse bonds incident to \(D\).
For any set of pieces \(\mathfrak u\), write
\[
 z_{t,D}=\prod_{c\in\Gamma'(t,D)}
            |\det\mathsf m_c(\mathcal M^j\mathsf U)|^{-1}.
\]
Then
\begin{equation}\label{eq:history-product-majorant}
 \begin{split}
 &e^{-A_{\rm cr}(\mathfrak u;\mathsf U)}
      \prod_{(t,D)\in\mathfrak u}
       \left[\mathsf g_Y(V_j)z_{t,D}
                     e^{-\mathscr Q_{t,D}(\mathsf a'_j)/2}\right]\\
 &\hspace{12mm}\le
   \exp\left[-c_\star\sum_{(t,\mathfrak c)\in
                       \operatorname{Cr}(\mathfrak u)}p_0(g_t)^2\right]
          \prod_{(t,D)\in\mathfrak u}
                      \overline{\mathsf u}_{t,D}(V_j,\mathsf a'_j).
 \end{split}
\end{equation}
The same bound holds with the displayed pivot determinants
\(\prod_c\det(L\mathsf m_c)^{-1}\) in modulus.
For each piece the majorant is Borel and
\begin{equation}\label{II-77-77}
 \int\overline{\mathsf u}_{t,D}\,dm_{t,D}
                              \le e^{C_\star'\#_t(D)} .
\end{equation}
\end{proposition}

\begin{proof}
On a piece abbreviate the three exponents by
\[
 P_{\rm W}=L^jx_jA_{N_{t,D}}(\mathsf U),\qquad
 P_{\rm G}=x_jG_j^\infty(\mathsf Y(t,D);V_j),\qquad
 P_{\rm F}=\mathscr Q_{t,D}(\mathsf a'_j)/2 .
\]
Assign the amount \(L^jx_ja_p(\mathsf U)\) to each
plaquette in its share, half the gauge exponent at each
site of \(\mathsf Y(t,D)\), and
\(\gamma_{\rm kp}|\mathsf a'_j(b)|^2/4\) to each stored bond.
The sum of these amounts is
\[
 E_{t,D}=P_{\rm W}+\tfrac12P_{\rm G}
                    +\tfrac14\gamma_{\rm kp}\sum_b|\mathsf a'_j(b)|^2.
\]
For each created cell sum only the items whose supports
lie in \(\mathfrak c^{\sim2}\).
For \(i=k,\ldots,1\), the recurrence and increment bound give
\[
 x_{i-1}=L(x_i+b_i)\ge Lx_i/2>x_i
 \qquad\text{whenever }x_i\ge2B_\sharp .
\]
Starting from \(x_k\ge2B_\sharp\), downward induction gives
\(x_j\ge x_t\ge x_k\) for \(j\le t\le k\). Hence
\(p_0(g_j)\ge p_0(g_t)\) and
\(1-B_\sharp g_t^2\ge1/2\).
The failed-test energy estimates of \cref{II-9-27,II-9-33},
applied using \cref{II-77-10,II-76-31}, give a lower bound
\(5^3c_\star p_0(g_t)^2\) for this
local sum, in each of the density, \(P\), \(Q\), and
fluctuation cases.
Every item contains a lattice site, which lies in one
level-\(t\) cell. At most \(5^3\) two-cell enlargements
contain that site. Thus each item occurs at most \(5^3\)
times, and
\begin{equation}\label{II-77-67}
 E_{t,D}\ge c_\star\#\operatorname{Cr}(t,D)\,p_0(g_t)^2 .
\end{equation}
If a contour mean is undefined, its gauge weight is zero
and the product bound is immediate. Otherwise all these
exponents are finite.

There are at most \(6N_D\) positive coarse bonds incident
to \(D\). By \eqref{II-7-18},
\(z_{t,D}\le(2L^9)^{6N_D}\).
After paying \(E_{t,D}\), the remaining fluctuation
exponent is
\[
 P_{\rm F}-\tfrac14\gamma_{\rm kp}\sum_b|\mathsf a'_j(b)|^2
       \ge\tfrac34\gamma_{\rm kp}\sum_b|\mathsf a'_j(b)|^2
       \ge\tfrac14\gamma_{\rm kp}\sum_b|\mathsf a'_j(b)|^2.
\]
The remaining gauge exponent is \(P_{\rm G}/2\).
This proves the bound for one piece using
\eqref{II-77-67}. The shares are disjoint by
\cref{II-21-30}, so their product proves
\eqref{eq:history-product-majorant}.
The displayed determinant has modulus
\(L^{-3\#\Gamma'}z_{t,D}\le z_{t,D}\), proving its version.

For the integral, \(\mathsf Y(t,D)\subset D\) gives
\(n^{\rm g}_{t,D}\le L^3N_D\).
Every positive stored bond is determined by its initial
level-\(j\) site in \(D\) and one of three directions;
hence there are at most \(3L^3N_D\) such bonds.
Use \eqref{II-7-14}, \(x_j\varepsilon_j^2=p_0(g_j)^2\ge1\),
and the three-dimensional Gaussian integral:
\[
 \int\overline{\mathsf u}_{t,D}\,dm_{t,D}
 \le(2L^9)^{6N_D}500^{L^3N_D}
                  (4\pi/\gamma_{\rm kp})^{(9/2)L^3N_D}.
\]
This is \eqref{II-77-77}. Borel measurability follows from
that of the gauge factor and continuity of the Gaussian.
\end{proof}

\section{The mass of a history}
\label{sec:history-mass}

We now apply the component majorant to the history measure.
There are two identifications to make. The retained elimination
factors must be the pivot determinants of a glued configuration,
and the retained quadratic forms must act on exactly the stored
coordinates of the components. Once these are established, the
positive density can be bounded on a Borel set of full measure.

Fix a history \(h\) of depth \(1\le k\le K\). Throughout this
section assume the running bounds, support and strip hypotheses
of \cref{sec:history-support-identity,II-76-31,II-77-63}.
The collared background is the critical orbit of \cref{II-10-20}.
For \(U\in\mathsf{Sup}(h)\), write
\(\mathsf U=U'\vee U\) for a glued representative, with
\(U'\in U_k^\diamond(\xi_h^\diamond(U))\).
The free coordinate map is denoted by \(\mathsf v_h\); thus
\(\xi_h=\mathrm{cp}_h\mathsf v_h\).
The preceding support results give its Borel definedness
on \(\mathsf{Sup}(h)\).

\subsection{The elimination factors at the background}

Let \(\Gamma_t^\partial\) consist of the positive level-\(t\)
bonds with an endpoint in \(\Omega_t\) and no endpoint in
\(\Lambda_t\). Both endpoints belong to one component \(D\)
of \(Z_t\); denote the corresponding subset by
\(\Gamma^\partial(t,D)\). Assume \(\alpha_0,\alpha_1\le1\)
and \(m'<m+K-k\), so that \(m'<m+K-t\) for
\(1\le t\le k\). Also assume, for these \(t\),
\begin{equation}\label{eq:history-pivot-ceilings}
 L^2C_{\rm reg}\varepsilon_t\le c_2^{\rm SO},
 \qquad
 2C_{\rm reg}\varepsilon_t
        <\min\{\varepsilon_{\mathsf z},\varpi_\chi\}.
\end{equation}
Here \(\varpi_\chi\) is the reference radius of
\cref{II-26-53}; the strict ceiling keeps its logarithm
on the specified branch.

\begin{lemma}\label{II-195-31}
For \(c\in\Gamma_t^\partial\), the averages needed to
evaluate \(z_c\) on the restriction to \(X(c)\) are defined
on the analytic branch, agree with the full block field, and
\begin{equation}\label{II-195-33}
 \varpi_c(\mathcal M^{t-1}\mathsf U)
                         \le2C_{\rm reg}\varepsilon_t.
\end{equation}
Moreover
\begin{equation}\label{II-195-37}
 e^{z_c(\mathsf U)}
       =\det\bigl(L\mathsf m_c(\mathcal M^{t-1}\mathsf U)\bigr)^{-1}
       >0 .
\end{equation}
The determinant is independent of the glued representative.
\end{lemma}

\begin{proof}
Put \(Y_n=B^{t-n}(c_-)\cup B^{t-n}(c_+)\), \(0\le n\le t\).
Then \(B(Y_{n+1})=Y_n\). In level-\((t-1)\) units each
point of \(Y_{t-1}\) is within \(2L\) of \(\Omega_t\).
The separation of the regions puts these points in
\(\Lambda_{t-1}\) when \(t\ge2\), and in
\(\Omega_0^\diamond\) when \(t=1\). The blocks below them
remain in that region. Thus \(Y_0\subset\Omega_0^\diamond\)
and all its inner bare bonds belong to \(X(c)\).
The endpoint conditions also give \(c\in\Gamma_t^\diamond\).
By \cref{II-10-20}, every bare plaquette with corners in
\(Y_0\) has curvature at most
\(\alpha L^{-2t}\), where
\(\alpha=L^2C_{\rm reg}\varepsilon_t\le c_2^{\rm SO}\).

For completeness, the averaging estimate can be iterated
on these blocks alone. Let \(\omega_n\) be the maximum
plaquette deviation with corners in \(Y_n\), and set
\[
 P_n=\prod_{l=0}^{n-1}(1+\tfrac12L^{2(l-t)}).
\]
The one-step proof of \cref{I-10-97} uses only the blocks
below the four corners of the coarse plaquette. Their
union is contained in \(Y_n\). Hence induction gives
\[
 \omega_n\le\alpha L^{2(n-t)}P_n,\qquad
 P_n\le\exp\!\left(\frac1{2(L^2-1)}\right)<2 .
\]
Indeed, for \(n<t\), the right side is at most
\(2\alpha L^{-2}\), below the one-step curvature radius,
and
\[
 \omega_{n+1}\le L^2\omega_n+C_0^{\rm SO}\omega_n^2
 \le L^2\omega_n(1+\tfrac12L^{2(n-t)}).
\]
Here \(C_0^{\rm SO}\alpha\le L^2/4\).
All double blocks required at the next step lie in \(Y_n\);
their curvature is below \(\omega_M\), so the real and
analytic definitions of the average agree. This also
proves by induction that restriction to the inner bonds
of \(X(c)\cap\Omega_0^\diamond\) changes no required value.
At \(n=t-1\) the bound is \eqref{II-195-33}.
The real logarithm and positivity in \cref{II-26-53}
now give \eqref{II-195-37}.

If \(U'_2=(U'_1)^u\) are two elements of the critical
orbit, extend \(u\) by one outside the collar. The glued
fields are gauge related on its inner bonds, which contain
every bond needed by the preceding averages. Locality and
covariance therefore relate their block fields on the
double block of \(c\). Differentiating covariance gives
\[
 \mathsf m_c(V^u)
  =\operatorname{Ad}(u(c_-))\mathsf m_c(V)
                  \operatorname{Ad}(u(b_0(c)_-))^{-1}.
\]
The two adjoint maps have determinant one on
\(\mathfrak g\). This proves independence of the determinant;
invertibility also follows from \eqref{II-7-18}.
\end{proof}

\subsection{The stored forms}\label{subsec:ymm1-65-2}

Assume the hypotheses and the fixed-radius ceiling of
\cref{II-65-44} for depth \(k\) and each step \(t\le k\),
and those of \cref{II-102-10}. The radius for the local
gauge neighborhoods can be taken to be
\(r_{\mathsf N}=4R_1\mathfrak M\). It satisfies
\[
 \frac{L(L+R_1+1)}{L-1}\mathfrak M
       <r_{\mathsf N},\qquad 2r_{\mathsf N}+1<LM,
\]
since \(L\ge13\), \(R_1\ge L\), and
\(M\ge L^4M_1\ge L^4R_1\mathfrak M\).
The actual pair read from \(U'\) at step \(t-1\)
therefore lies in \(\mathscr O^+(\mathfrak k^\sim)\)
at radii \((\alpha_0/2,\alpha_1,\alpha_2/2)\), for every
strip component \(\mathfrak k\) of that step. Set
\(\gamma_{\rm kp}=\gamma_{\rm pos}/4\).

Each strip component belongs to a unique component \(D\)
of \(Z_t\), and its positive external bond set
\(E_{\mathfrak k}^+\) is contained in
\(\mathfrak X^{\rm st+}(t,D)\). These sets partition
\(\mathfrak X^{\rm st+}(t,D)\). In fact an external bond
has no endpoint in \(\Lambda_t\), so both endpoints
belong to a single component of \(Z_t\). It also has an
endpoint in its strip component. The component of \(Z_t\)
containing this bond must therefore be the one containing
that strip. The converse follows by the definition of the
external bonds.

Consequently the canonical form on a piece is
\begin{equation}\label{eq:history-piece-kept-form}
 Q_{t,D}(a;U')
   =\sum_{\mathfrak k\subset D}
      \langle a,\Delta^{\rm loc}_{\mathfrak k}
                          (p_{\mathfrak k}(U'))a\rangle .
\end{equation}
Extend \(a\) by zero off the stored bonds of the piece.
The matrices are real symmetric, and \cref{II-102-10}
gives
\begin{equation}\label{II-195-49}
 Q_{t,D}(a;U')\ge
      2\gamma_{\rm kp}
          \sum_{b\in\mathfrak X^{\rm st+}(t,D)}|a(b)|^2 .
\end{equation}
The depth-\(k\) pair is used here at every earlier step;
the coefficients and keeping predicates at step \(t\)
are those of the truncated history \(h|_t\).

Suppose now that the formed factors are precisely those
of \cref{sec:formed-factors}, retained from steps \(1,\ldots,k\).
Their total exponent \(\Phi_h(\phi)\) is defined when its
evaluation on every admitted glued representative is
defined and has the same value. The grouping by
\((t,X)\) in \eqref{eq:formed-factor-definitions} then gives,
at \(\phi=\xi_h(U)\),
\[
 \Phi_h(\phi)
  =\sum_{t=1}^k\sum_{c\in\Gamma_t^\partial}z_c(U')
        -\frac12\sum_{(t,D)}Q_{t,D}(\mathsf a'_{t,D}(U);U').
\]
All sums here are finite except for the normally
convergent coefficient series already used to define
the kept forms. Thus their regrouping preserves their
values. By \cref{II-195-31},
\begin{equation}\label{II-195-47}
 e^{\Phi_h(\xi_h(U))}
   =\prod_{t=1}^k\prod_{c\in\Gamma_t^\partial}
       \det\bigl(L\mathsf m_c(\mathcal M^{t-1}\mathsf U)\bigr)^{-1}
                                      e^{-s(U)/2},
\end{equation}
where
\[
 s(U)=\sum_{(t,D)}Q_{t,D}(\mathsf a'_{t,D}(U);U'),
 \qquad
 n(U)=\sum_{(t,D)}\sum_{b\in\mathfrak X^{\rm st+}(t,D)}
                                  |\mathsf a'_{t-1}(b)(U)|^2.
\]
The determinant product is positive and independent of
\(\mathsf U\). Since \(s(U)\) is real, injectivity of the
real exponential in \eqref{II-195-47} shows that it too
is independent of the chosen admitted representative.
It satisfies \(s(U)\ge2\gamma_{\rm kp}n(U)\), and
\(n(U)=0\) implies \(s(U)=0\).

The following form of this observation avoids choosing
a background as a function of \(U\). Where \(\Phi_h\)
is defined and \(n(U)>0\), set \(\lambda(U)=s(U)/n(U)\);
otherwise set \(\lambda(U)=2\gamma_{\rm kp}\). The real
quadratic forms
\begin{equation}\label{II-196-8}
 Q^{\rm sc}_{t,D}[U](a)
       =\lambda(U)\sum_{b\in\mathfrak X^{\rm st+}(t,D)}|a(b)|^2
\end{equation}
have the lower bound \eqref{II-195-49}. At the actual
stored fields their sum equals \(s(U)\), including the
case \(n(U)=0\). Hence \eqref{II-195-47} holds with
these forms and with every glued representative in its
determinant product.

\subsection{The density on the free coordinates}

Let
\[
 \mathcal V_0[h]=\{U:\mathfrak w_h(U)>0\},\qquad
 \check\tau_h=(\mathsf v_h)_*(\mathfrak w_h\rho_0\,dU).
\]
These are a Borel set and a finite positive Borel measure,
respectively. The history mass is
\(\mathsf t(h)=\int\mathfrak w_h\rho_0\,dU\).
All of \(\mathcal V_0[h]\) lies in the domain of
\(\mathsf v_h\), so this is also the total mass of
\(\check\tau_h\).

Assume the history has an integrable density in the
completed coordinates of the form
\begin{equation}\label{eq:history-format-density}
 \mathsf f_h(\phi)=
   \operatorname{Stay}^{\rm T}_h(\phi)\chi_k(\Omega_k)
                          e^{A_h(\phi)+\Phi_h(\phi)}
\end{equation}
outside a Borel \(\mathscr M_h\)-null set. The exponents
are partial Borel functions, and where this expression
is nonzero their admitted evaluation set is nonempty.
Assume that for every admitted glued evaluation,
\begin{equation}\label{II-196-5}
 \operatorname{Re}A_h(\phi)
  =-A(\beta_k(\cdot),\mathsf U)
    +\operatorname{Re}\{\mathbf E_k(\mathsf U)
                  +\mathbf E_k^{\rm bd}(\mathsf U,\mathfrak q)\}
    -\mathcal E_{k,\mathcal S(h)} ,
\end{equation}
with real \(\mathcal E_{k,\mathcal S(h)}\).
These are the exponents obtained by collecting the
terms in \eqref{eq:history-new-exponents}.
Finally assume the core membership, layer bounds and
family estimates of \cref{II-101-34}.

\begin{proposition}\label{II-196-9}
Under these hypotheses, \(\check\tau_h\) has a nonnegative
Borel density \(f_h\) with respect to \(\mathscr M_h^{\rm fr}\).
There is a Borel set
\(S'\subset\mathsf v_h(\mathcal V_0[h])\) of full
\(\check\tau_h\)-measure such that, for every \(v\in S'\),
one can choose \(U\in\mathcal V_0[h]\), \(\mathsf v_h(U)=v\),
and a glued representative \(\mathsf U\) for which
\begin{equation}\label{eq:history-density-majorant}
 \begin{split}
 f_h(v)\le{}&
 e^{-\mathcal E_{k,\mathcal S(h)}
       +C_5\sum_{n=0}^k\#_n(\mathcal L_n)
       -A(\beta_k(\cdot),\mathsf U)}\\
 &{}\times\prod_{(t,D)}
   \left[\mathsf g_{\mathsf Y(t,D)}(V_{t-1})
          z^{\rm dsp}_{t,D}(\mathcal M^{t-1}\mathsf U)
          e^{-Q^{\rm sc}_{t,D}[U](\mathsf a'_{t,D}(U))/2}\right],
 \end{split}
\end{equation}
where \(z^{\rm dsp}_{t,D}\) is the positive determinant
product over \(\Gamma^\partial(t,D)\) in
\eqref{II-195-47}.
\end{proposition}

\begin{proof}
If \(N\) is the exceptional set in the completed
coordinates, \(N'=\mathrm{cp}_h^{-1}(N)\) is
\(\mathscr M_h^{\rm fr}\)-null by \eqref{II-58-48}.
Pullback gives the integrable density
\(d=\mathsf f_h\circ\mathrm{cp}_h\) off \(N'\).
Indeed completion sends \(\check\tau_h\) to the
completed history measure, and
\(\mathrm{rs}_h\mathrm{cp}_h=\mathrm{id}\); apply these
identities to the indicator of a Borel set in the free
coordinates. Since \(\check\tau_h\) is positive,
integrating over the sets where \(\operatorname{Im}d\)
has either sign, and then where \(\operatorname{Re}d<0\),
shows that \(d\) is nonnegative almost everywhere.
Its modulus, extended by zero on \(N'\), is therefore
a nonnegative Borel density \(f_h\).

The image of the Borel set \(\mathcal V_0[h]\) under
the Borel map \(\mathsf v_h\) is analytic. The
analytic-image argument used in \cref{II-166-30}
supplies a Borel subset \(S_0\) of this image of full
\(\check\tau_h\)-measure. The density vanishes outside
the external tiers, because \(\operatorname{Stay}^{\rm T}_h\)
does. Intersect \(S_0\) with the tiers and remove \(N'\)
to obtain \(S'\); points with zero density may also
be removed. At every remaining \(v\) choose a bare
preimage \(U\). Put \(\phi=\mathrm{cp}_h(v)=\xi_h(U)\).

At this point \(\chi_k(\Omega_k)=1\). Each factor in
that test is a passed \(P_k\)-test of the selected
history. Furthermore,
\begin{equation}\label{II-172-20}
 \operatorname{Stay}_h(v)
      =\prod_{(t,D)}\mathsf g_{\mathsf Y(t,D)}(V_{t-1})>0.
\end{equation}
To check the identity, regard stored factors as entries
in their defining lists, so that multiplicities are
preserved. Every stored indicator, including a carried
density test, equals one at \(U\). The written tests
have the same values: the test transport preserves
their required averages and logarithms, while the
contour tests use only internal bonds of \(B(y)\),
which remain free. The only remaining factors are
\(\mathfrak G_{t-1}^{\rm e}(V_{t-1};y,x)\), with
\(y\in\mathsf Y(t,D)\) and \(x\in B(y)\setminus\{y\}\).
Each such pair occurs once. Their product is precisely
the right side of \eqref{II-172-20}; their contour
variables are defined since \(\mathfrak w_h(U)>0\).

The tiers now give
\(\operatorname{Stay}^{\rm T}_h(\phi)=\operatorname{Stay}_h(v)\).
Thus the nonzero density formula has defined exponents
and an admitted glued evaluation. Choose one.
Equations \eqref{II-195-47} and \eqref{II-196-8}
identify its formed factors. Take the modulus in
\eqref{eq:history-format-density}, use
\eqref{II-196-5}, and apply \eqref{II-101-44} to the
regular and boundary sums. This proves
\eqref{eq:history-density-majorant}.
All choices were made at a single point; the final
bound below will be a function of \(v\) alone.
\end{proof}

\subsection{Integration over the components}

For a piece \((t,D)\), let \(\operatorname{pr}_{t,D}\)
select \(V_{t-1}|_{\mathfrak B(t,D)}\) and the
stored variables on \(\mathfrak X^{\rm st+}(t,D)\)
from the free coordinates. It is a continuous
coordinate projection. The internal bonds of a block
\(B(y)\), \(y\in\mathsf Y(t,D)\), are free:
the cell containing \(y\) lies outside \(\Omega_t\);
for \(t\ge2\) it avoids the enlarged old large-field
region, so \(B(y)\subset\Lambda_{t-1}\).
For \(t=1\) use \(\Omega_0=T_\ell\).
Thus both endpoints belong to \(\Omega_{t-1}\) and
neither to \(\Omega_t\), which is the required free
bond condition.

Different pieces use disjoint coordinate factors.
Different steps carry different level indices. At
one step the blocks \(B(y)\) assigned to distinct
components cannot share a fine lattice site, because
its containing coarse block is unique.
Every stored bond belongs to a unique component of
\(Z_t\), as shown above. In particular every Lebesgue
coordinate occurs in exactly one piece; all remaining
coordinates carry Haar probability measure.
Regrouping the finite product \eqref{II-58-48}
and using Tonelli therefore gives
\begin{equation}\label{II-195-18}
 \int\prod_{(t,D)}
       \overline{\mathsf u}_{t,D}(\operatorname{pr}_{t,D}v)
                         \,d\mathscr M_h^{\rm fr}(v)
  =\prod_{(t,D)}\int\overline{\mathsf u}_{t,D}\,dm_{t,D}
  \le e^{C_\star'\sum_{t=1}^k\#_t(Z_t)} .
\end{equation}
The product measures are \(\sigma\)-finite, and the
integrands are nonnegative Borel functions. The equality
also holds for no pieces, when only a Haar probability
space remains.

\begin{proposition}\label{II-195-19}
Under the hypotheses of this section, set
\[
 C_{\rm maj}=\max\{C_5,L^3C_5+C_\star'\}.
\]
Then
\begin{equation}\label{II-195-20}
 e^{\mathcal E_{k,\mathcal S(h)}}\mathsf t(h)
 \le
 \exp\!\left[C_{\rm maj}
        \left(N_k+\sum_{t=1}^k\#_t(Z_t(h))\right)\right]
 \exp\!\left[-c_\star
       \sum_{t=1}^kp_0(g_t)^2n_t^{\rm cr}(h)\right].
\end{equation}
The constants are independent of the volume, the depth,
the running couplings and the history.
\end{proposition}

\begin{proof}
Apply \cref{II-77-63} at a point \(v\in S'\) with
\(\Gamma'(t,D)=\Gamma^\partial(t,D)\) and
the forms \eqref{II-196-8}.
The background and strip hypotheses hold on the support.
Every elimination bond has its endpoints in \(D\),
and \eqref{II-195-33} verifies its curvature bound.
The stored fields are defined at the chosen bare
preimage, and \eqref{II-195-49} is the required lower
bound for the forms. The action reserved for all
created cells is at most \(A(\beta_k(\cdot),\mathsf U)\)
by \cref{II-21-30}. The pieces partition the created
cells. Hence \eqref{eq:history-density-majorant} becomes
\[
 f_h(v)\le
 e^{-\mathcal E_{k,\mathcal S(h)}
      +C_5\sum_n\#_n(\mathcal L_n)
      -c_\star\sum_t p_0(g_t)^2n_t^{\rm cr}(h)}
       \prod_{(t,D)}
          \overline{\mathsf u}_{t,D}(\operatorname{pr}_{t,D}v).
\]
The right side depends only on \(v\) and the fixed
history. Thus no measurable selection of its auxiliary
backgrounds or forms is needed.
Integrate on \(S'\), which has full history measure,
and enlarge the nonnegative integral to the whole
free space. Equation \eqref{II-195-18} and the count
\(\sum_n\#_n(\mathcal L_n)
       \le N_k+L^3\sum_t\#_t(Z_t)\)
from \eqref{II-101-44} prove \eqref{II-195-20}.
\end{proof}


\section{Summing the large-field histories}
\label{sec:history-summation}

A failed test creates a cell whose suppression must cover
the volume it produces at every later scale. We first
locate these later cells near their creations. We then
count the choices of tests inside the created cells and
sum over arbitrary subsets of the cell partitions.

Throughout this section \(L\ge13\), \(r\ge2\), and
\(\mathfrak p_0\) is a positive integer. The geometric
constants and partitions are those of
\eqref{eq:history-standing}; in particular
\(M_2=L^2M_1\ge L^4\), \(M\ge L^6\), and
\(M/(LM_2)\) is a power of \(L\) at least \(L\).
Suppose
\[
 x_{j-1}=L(x_j+b_j),\qquad |b_j|\le B_\sharp,\qquad
 y_j=\log x_j,\qquad p_0(g_j)=A_0y_j^{\mathfrak p_0}.
\]
Write \(N_j=\#\mathbb T^{(j)}\) and
\(s_j=M\check R_j a_{K-j}\) for the physical side
of a cell of \(\mathcal P_j\). Define
\[
 Y_1=\log\max\{e,2B_\sharp\},\qquad
 Y_0=\max\left\{Y_1,\frac{2\log(2L)}{L^{1/r}-1}\right\}.
\]

\begin{lemma}\label{II-75-37}
If \(y_K\ge Y_1\), then \(x_j\) and \(y_j\) decrease
strictly with \(j\), \(\check R_j=R(x_j)\), and
\[
 1\le y_j,\qquad y_j^r\le\check R_j<Ly_j^r,\qquad
 \log(L/2)\le y_{j-1}-y_j\le\log(2L),\qquad
 K-j+1\le y_j .
\]
If \(y_K\ge Y_0\), then
\[
 \begin{aligned}
 \check R_j&\le\check R_{j'}\le L\check R_j
       &&(1\le j'\le j,\ j-j'\le2),\\
 s_{j'}&\le L^{-\lfloor(j-j')/2\rfloor}s_j
       &&(1\le j'\le j).
 \end{aligned}
\]
\end{lemma}

\begin{proof}
Put \(\bar B=LB_\sharp/(L-1)\). Since
\(x_K\ge2B_\sharp>\bar B\) when \(B_\sharp>0\),
and \(x_K>0=\bar B\) otherwise, the inequality
\[
 x_{j-1}-x_j\ge(L-1)(x_j-\bar B)
\]
gives strict decrease by descending induction.
The minimum defining \(\check x_j\) is therefore \(x_j\).
Minimality of the power of \(L\) defining \(R(x_j)\)
gives \(y_j^r\le R(x_j)<Ly_j^r\), since \(y_j\ge1\).
Also \(|b_j|/x_j\le1/2\), so taking logarithms in
\(x_{j-1}/x_j=L(1+b_j/x_j)\) gives the two
increment bounds. Summing the lower bound, and using
\(\log(L/2)>1\), gives \(K-j+1\le y_j\).

For \(j-j'\le2\), the upper increment bound gives
\[
 y_{j'}\le y_j+2\log(2L)\le L^{1/r}y_j .
\]
Thus \(y_{j'}^r\le L\check R_j\), and the minimality
of \(R(x_{j'})\) gives the claimed comparison.
Since \(s_{j'}/s_j=L^{j'-j}\check R_{j'}/\check R_j\),
one backward step costs at most one and two backward
steps cost at most \(L^{-1}\). Iteration proves the
last inequality.
\end{proof}

\subsection{Later volume assigned to created cells}

For a history \(h\), put
\[
 V_t=R_t'\cup\mathsf K(I_t)\cup\mathsf R_t^{\mathrm{lo}},
 \qquad W_t=[\mathcal W_t^\wedge(h)]_t .
\]
Here \([A]_t\) denotes the cells of a level-\(t\)
region \(A\). Thus \(W_t\) is the union of the cell
sets of \(P_t',Q_t',V_t\), and also of \(P_0'\)
when \(t=1\). Write \(n_t^{\mathrm{cr}}(h)=\#W_t\).
For a final component \(X\) of \(Z_k(h)\), let
\(n_t^{\mathrm{cr}}(h;X)\) count those cells of
\(W_t\) contained in \(X\).

The definitions of \(\Omega_t\) and \eqref{II-4-27} give
\[
 \begin{split}
 Z_1&=P_0'^{\sim7}\cup P_1'^{\sim6}
                   \cup Q_1'^{\sim4}\cup V_1^\sim,\\
 Z_t&=\widehat Z_{t-1}^{\sim10}\cup P_t'^{\sim6}
                   \cup Q_t'^{\sim4}\cup V_t^\sim
                   \quad(t\ge2).
 \end{split}
\]
In particular \(Z_t\subset Z_{t+1}\), and a component
at one level lies in a unique component at the next.

\begin{proposition}\label{II-75-45}
Suppose \(y_K\ge Y_0\). For every component \(X\)
of \(Z_k(h)\) and every choice \(r_n\ge0\),
\begin{equation}\label{II-75-47}
 \sum_{n=1}^k r_n|Z_n(h)\cap X|_n
 \le48^3\sum_{t=1}^k n_t^{\mathrm{cr}}(h;X)
                         \sum_{n=t}^k r_n .
\end{equation}
Every such component contains a created cell.
The same inequality holds for the whole regions
\(Z_n(h)\), with \(n_t^{\mathrm{cr}}(h)\) on the right.
\end{proposition}

\begin{proof}
A chain through \(a\) layers of cells of side \(s\)
joins any point of \(A^{\sim a}\) to a point of \(A\)
at distance at most \(as\); the chain stays in
\(A^{\sim a}\). Apply this observation to the displayed
formulas for \(Z_n\). A point newly covered at level
\(n\) lies within \(7s_n\) of a created cell in its
component. A point covered from \(\widehat Z_{n-1}\)
lies within \(10s_n\) of that hull, and within a
further \(s_n\) of \(Z_{n-1}\). Induction therefore
places every \(x\in Z_n\) within
\[
 \rho_{t,n}=7s_t+11\sum_{j=t+1}^n s_j
\]
of a created cell of some level \(t\le n\) in the
same component. The chains and the nesting of
components preserve this last assertion at each step.
It also proves that every nonempty final component
contains a creation.

Assign each cell of \(Z_n\cap X\), by its centre,
to one such creation. The centre of a cell assigned
to a level-\(t\) cell lies within
\(\rho_{t,n}+s_t/2\) of that cell's centre.
An interval of length \(2a\) contains at most
\(\lfloor2a/s_n\rfloor+1\) points of a lattice of
spacing \(s_n\). Applying this in the three
coordinates, with lifts on the torus, bounds the
number assigned to a fixed creation by
\[
 \left(\left\lfloor
       \frac{2\rho_{t,n}+s_t}{s_n}\right\rfloor+1\right)^3.
\]
Now \(2\rho_{t,n}+s_t\le22\sum_{j=t}^ns_j\), and
\cref{II-75-37} gives
\[
 \sum_{j=t}^ns_j\le
       s_n\sum_{u\ge0}L^{-\lfloor u/2\rfloor}
       =\frac{2L}{L-1}s_n .
\]
Since \(44L/(L-1)<48\), the displayed count is at
most \(48^3\). Consequently
\[
 |Z_n\cap X|_n\le48^3\sum_{t\le n}
                         n_t^{\mathrm{cr}}(h;X).
\]
Multiply by \(r_n\)
and interchange the finite sums to obtain
\eqref{II-75-47}. Distinct final components partition
both the earlier regions and the creations, so
summing over them gives the whole-region inequality.
\end{proof}

\subsection{Counting the outcomes in fixed created cells}

Set
\[
 N_\tau=\left(\frac{M}{LM_2}\right)^3,\qquad
 N_\delta=\left(\frac{M}{M_2}\right)^3,\qquad
 n_t^{\mathrm{bit}}=3N_\tau+N_\delta+1
                          +6L^3(M\check R_t)^3,\qquad
 \lambda_L=(6L^3+1)\log2 .
\]

\begin{lemma}\label{II-75-82}
For arbitrary sets of cells \(W_t\subset\mathcal P_t\),
\begin{equation}\label{eq:history-entropy}
 \#\{h:[\mathcal W_t^\wedge(h)]_t\subset W_t
                         \text{ for all }t\}
 \le\prod_{t=1}^k2^{n_t^{\mathrm{bit}}\#W_t},
 \qquad
 n_t^{\mathrm{bit}}\le(6L^3+1)(M\check R_t)^3 .
\end{equation}
\end{lemma}

\begin{proof}
Each of \(P_t,Q_t,R_t\) is a union of test cubes.
Its hull lies in \(W_t\), and each cell contains
exactly \(N_\tau\) test cubes. These three choices
therefore use at most \(3N_\tau\#W_t\) binary
choices. The set \(\mathsf R_t^{\mathrm{lo}}\) is
a subset of \(W_t\), using \(\#W_t\) more.

For the selected collar bonds, we must locate the
initial site in a cell of \(W_t\). At the stage at
which a bond \(b\in I_t\) is selected, it belongs to
the collar of a region \(\mathcal Y\) containing no cell
of \((\Omega_t^c)^{\sim2}\). In level-\((t-1)\) units
every endpoint of \(b\) lies within
\(\varpi_T+3\le M+3\) of a site in \(\mathcal Y\).
The cells have side \(LM\check R_t\ge13M\).
Their centres lie on the cell lattice, so the cell
containing that endpoint touches a cell of
\(\mathcal Y\). It belongs to \(\Omega_t^{\sim-1}\):
otherwise the touched cell would lie in
\((\Omega_t^c)^{\sim2}\). Thus the cell of \(b_-\)
belongs to \(\mathsf K(I_t)\subset W_t\).
There are \((LM\check R_t)^3\) sites per cell and
at most six directed bonds per initial site. Hence
\(I_t\) uses at most
\(6L^3(M\check R_t)^3\#W_t\) binary choices.

At the first step, the density cubes defining \(P_0\)
need not refine \(\mathcal P_1\). If they do,
each cell contains at most
\[
 \left(\frac{M\check R_1}{M_2\check R_0}\right)^3
 \le N_\delta
\]
of them. If the cell partition is finer, every chosen
density cube contains a cell of \(W_1\), and distinct
density cubes contain disjoint sets of cells.
There are then at most \(\#W_1\le N_\delta\#W_1\)
eligible density cubes. Both cases cost at most
\(N_\delta\#W_1\) choices.

These sets determine the history, including its
recursively explored sets. Allowing all the subsets
just counted can only increase their number, which
proves the first bound. Finally
\[
 3N_\tau+N_\delta+1
 \le M^3(3L^{-15}+L^{-12}+L^{-18})
 <M^3\le(M\check R_t)^3
\]
proves the second.
\end{proof}

\subsection{Absorbing volume and entropy}

We use two weights:
\[
 \omega_j=y_j\quad\text{in case II},\qquad
 \omega_j=\check R_j\quad\text{in case I}.
\]
Case II requires \(2\mathfrak p_0>3r+2\);
case I requires \(2\mathfrak p_0>4r+1\).
For \(a,C>0\), put
\begin{equation}\label{II-75-97}
 \begin{split}
 Y^{\mathrm{II}}(a,C)&=
 \max\left\{Y_0,\left(\frac4{aA_0}\right)^{1/\mathfrak p_0},
 \left(\frac{2CL^3M^3}{aA_0^2}\right)^{1/(2\mathfrak p_0-3r-2)}
 \right\},\\
 Y^{\mathrm I}(a,C)&=
 \max\left\{Y_0,\left(\frac4{aA_0}\right)^{1/\mathfrak p_0},
 \left(\frac{2CL^4M^3}{aA_0^2}\right)^{1/(2\mathfrak p_0-4r-1)}
 \right\}.
 \end{split}
\end{equation}

\begin{lemma}\label{II-75-99}
In either case \(\varrho\in\{\mathrm I,\mathrm{II}\}\),
if \(y_K\ge Y^\varrho(a,C)\), then
\[
 ap_0(g_t)^2\ge2p_0(g_t)
                  +C\sum_{n=t}^K\omega_n(M\check R_n)^3
                  \quad(0\le t\le K).
\]
\end{lemma}

\begin{proof}
Write \(y=y_t\). The second entry of the maximum
ensures \(ap_0(g_t)^2/2\ge2p_0(g_t)\).
By \cref{II-75-37}, the sum has at most \(y\)
terms, \(y_n\le y\), and \(\check R_n<Ly^r\).
It is at most \(L^3M^3y^{3r+2}\) in case II
and \(L^4M^3y^{4r+1}\) in case I.
The last entry of the corresponding maximum
bounds \(C\) times this expression by the other
half of \(ap_0(g_t)^2\).
\end{proof}

For \(C\ge0\), \(a>0\), define
\begin{equation}\label{II-75-53}
 \Pi_k(C,a)=e^{CN_k}
       \prod_{t=1}^k(1+e^{-ap_0(g_t)^2})^{\#\mathcal P_t}.
\end{equation}

\begin{proposition}\label{II-75-106}
Let \(c>0\), \(D\ge0\), \(0<\eth<1\), and suppose
\[
 y_K\ge Y^\varrho(\eth c,48^3(D+\lambda_L)).
\]
If a nonnegative function \(F\) on \(\mathcal H_k\)
satisfies
\[
 F(h)\le
 \exp\left[D\sum_{j=1}^k\omega_j\#_j(Z_j(h))
             -c\sum_{t=1}^kp_0(g_t)^2n_t^{\mathrm{cr}}(h)\right],
\]
then
\begin{equation}\label{II-75-109}
 \sum_{h\in\mathcal H_k}F(h)
           \le\Pi_k(0,(1-\eth)c).
\end{equation}
\end{proposition}

\begin{proof}
Apply \eqref{II-75-47} with
\(r_j=(D+\lambda_L)\omega_j(M\check R_j)^3\).
Since \(\#_j(Z_j)=(M\check R_j)^3|Z_j|_j\),
\cref{II-75-99} gives
\[
 (D+\lambda_L)\sum_j\omega_j\#_j(Z_j)
 \le\sum_t[\eth c\,p_0(g_t)^2-2p_0(g_t)]
                                      n_t^{\mathrm{cr}}(h).
\]
Also \(\omega_j\ge1\), \(n_j^{\mathrm{cr}}\le|Z_j|_j\),
and \cref{II-75-82} imply
\[
 \lambda_L\sum_j\omega_j\#_j(Z_j)
       \ge\sum_t n_t^{\mathrm{bit}}\log2\,
                                      n_t^{\mathrm{cr}}(h).
\]
Subtract this last inequality. With
\(\zeta_t=e^{-(1-\eth)cp_0(g_t)^2-2p_0(g_t)}\),
the assumed bound becomes
\[
 F(h)\le\prod_t
       (2^{-n_t^{\mathrm{bit}}}\zeta_t)^{n_t^{\mathrm{cr}}(h)}.
\]
Sort histories according to their exact sets \(W_t\)
of created cells. Equation \eqref{eq:history-entropy}
bounds the multiplicity by
\(\prod_t2^{n_t^{\mathrm{bit}}\#W_t}\), cancelling
the binary factors. Summing over all
\(W_t\subset\mathcal P_t\) therefore gives
\[
 \sum_hF(h)\le
 \prod_t\sum_{W_t\subset\mathcal P_t}\zeta_t^{\#W_t}
 =\prod_t(1+\zeta_t)^{\#\mathcal P_t}
 \le\Pi_k(0,(1-\eth)c).
\]
\end{proof}

For later use, set
\[
 Y^{\mathrm S}(a)=
 \left(\frac{3\log L+\log2}
 {2a\mathfrak p_0A_0^2\log(L/2)}\right)^{1/(2\mathfrak p_0-1)}.
\]
If \(y_K\ge\max\{Y_1,Y^{\mathrm S}(a)\}\), convexity of
\(y\mapsto y^{2\mathfrak p_0}\) and
\cref{II-75-37} give
\[
 a[p_0(g_t)^2-p_0(g_k)^2]
 \ge2a\mathfrak p_0A_0^2y_K^{2\mathfrak p_0-1}
                      (k-t)\log(L/2)
 \ge(3\log L+\log2)(k-t).
\]
Consequently
\[
 \sum_{t=1}^kL^{3(k-t)}e^{-ap_0(g_t)^2}
                  \le2e^{-ap_0(g_k)^2}.
\]
Since \(\#\mathcal P_t=N_t/(M\check R_t)^3
\le M^{-3}L^{3(k-t)}N_k\), the inequality
\(\log(1+u)\le u\) now yields
\begin{equation}\label{II-75-72}
 \Pi_k(C,a)
 \le\exp[(C+2M^{-3}e^{-ap_0(g_k)^2})N_k]
 \le e^{(C+2)N_k}.
\end{equation}

\subsection{Vacuum normalisation and weighted sums}

Let \(f:\mathcal H_k\to[0,\infty)\) and let \(e\)
be real on admissible labels. Assume
\begin{equation}\label{eq:history-summation-input}
 \begin{split}
 f(h)&\le e^{C_{\rm maj}(N_k+\sum_j\#_j(Z_j(h)))}
                 \prod_te^{-cp_0(g_t)^2n_t^{\mathrm{cr}}(h)},\\
 |e(\mathcal S)|&\le C_\flat\sum_j(1+y_{j-1})
                                      \#_j(Z_j(\mathcal S)^\sim).
 \end{split}
\end{equation}
Here \(c>0\) and \(C_{\rm maj},C_\flat\ge0\).
Define
\[
 D_0=C_{\rm maj}+27(2+\log(2L))C_\flat .
\]
The constants and free parameters below are fixed
independently of the lattice and the coupling sequence.

\begin{proposition}\label{II-75-111}
Assume \(2\mathfrak p_0>3r+2\) and \(0<\eth<1\).
If
\[
 y_K\ge
 \max\{Y^{\mathrm{II}}(\eth c,48^3(D_0+\lambda_L)),
                              Y^{\mathrm S}((1-\eth)c)\},
\]
then
\begin{equation}\label{II-75-114}
 \sum_hf(h)e^{e(\mathcal S(h))}
       \le\Pi_k(C_{\rm maj},(1-\eth)c)
       \le e^{(C_{\rm maj}+2)N_k}.
\end{equation}
More generally, for \(0\le\theta<c\), \(\varkappa\ge0\),
if
\[
 y_K\ge
 Y^{\mathrm{II}}(\eth(c-\theta),
                          48^3(D_0+\varkappa+\lambda_L)),
\]
then, with \(\mathsf N(\mathcal S)\) the number of
components of \(Z_k(\mathcal S)\),
\begin{equation}\label{II-75-115}
 \sum_h f(h)e^{e(\mathcal S(h))}
       e^{\theta p_0(g_k)^2\mathsf N(\mathcal S(h))
                           +\varkappa\#_k(Z_k(h))}
 \le\Pi_k(C_{\rm maj},(1-\eth)(c-\theta)).
\end{equation}
\end{proposition}

\begin{proof}
Enlarging a cell region by one layer multiplies its
cell count by at most \(27\). Also
\[
 1+y_{j-1}\le(2+\log(2L))y_j
\]
by \cref{II-75-37}. Thus
\(F_0(h)=e^{-C_{\rm maj}N_k}f(h)e^{e(\mathcal S(h))}\)
satisfies
\[
 F_0(h)\le
 \exp\left[D_0\sum_jy_j\#_j(Z_j)
                  -c\sum_tp_0(g_t)^2n_t^{\mathrm{cr}}\right].
\]
Apply \cref{II-75-106} in case II, then
\eqref{II-75-72}, to obtain \eqref{II-75-114}.

Each final component contains a creation, and the
creations of distinct components are disjoint.
Since \(p_0(g_t)\ge p_0(g_k)\) for \(t\le k\),
\[
 p_0(g_k)^2\mathsf N(\mathcal S(h))
             \le\sum_tp_0(g_t)^2n_t^{\mathrm{cr}}(h).
\]
Also \(\#_k(Z_k)\le\sum_jy_j\#_j(Z_j)\).
The extra exponent therefore replaces \(c\) by
\(c-\theta\) and \(D_0\) by \(D_0+\varkappa\)
in the bound for \(F_0\). A second application of
\cref{II-75-106} gives \eqref{II-75-115}.
\end{proof}

The same argument controls the additional weights
attached to the history of a component. For a final
component \(X\), define
\[
 \begin{split}
 \operatorname{Dep}_k(X;h)
   &=c_{\rm Dep}\sum_t p_0(g_t)^2n_t^{\mathrm{cr}}(h;X),\\
 \operatorname{Vol}_k(X;h)
   &=C_{\rm vol}\sum_j\check R_j\#_j(Z_j(h)\cap X),
 \end{split}
\]
and let
\[
 \Theta^{(k)}(h)=\sum_X
   [w\operatorname{Dep}_k(X;h)+c_a\operatorname{Vol}_k(X;h)],
 \quad 0\le w<1,\quad c_a,C_{\rm vol}\ge0,\quad
 0<c_{\rm Dep}\le c .
\]
The earlier components partition \(Z_j\cap X\);
thus this volume is also the sum of the weights
\(C_{\rm vol}M^3\check R_j^4|D|_j\) over the
pieces \((j,D)\) contained in \(X\).

Put \(c''=(1-\eth)(1-w)c\). If \(c_a=0\), then
\begin{equation}\label{II-75-120}
 \sum_hf(h)e^{e(\mathcal S(h))+\Theta^{(k)}(h)}
       \le\Pi_k(C_{\rm maj},c'')
       \le e^{(C_{\rm maj}+2)N_k}
\end{equation}
provided
\[
 y_K\ge\max\{Y^{\mathrm{II}}(\eth(1-w)c,
                          48^3(D_0+\lambda_L)),Y^{\mathrm S}(c'')\}.
\]
If \(c_a>0\), assume \(2\mathfrak p_0>4r+1\).
Then
\begin{equation}\label{II-75-121}
 \sum_hf(h)e^{e(\mathcal S(h))+\Theta^{(k)}(h)}
       \le\Pi_k(C_{\rm maj}+\lambda_L,c'')
       \le e^{(C_{\rm maj}+\lambda_L+2)N_k}
\end{equation}
provided
\[
 y_K\ge\max\{Y^{\mathrm I}(\eth(1-w)c,
             48^3(D_0+c_aC_{\rm vol}+\lambda_L)),Y^{\mathrm S}(c'')\}.
\]
Indeed the component partitions give
\[
 \Theta^{(k)}(h)\le
 wc\sum_tp_0(g_t)^2n_t^{\mathrm{cr}}(h)
       +c_aC_{\rm vol}\sum_j\check R_j\#_j(Z_j).
\]
For \(c_a=0\), apply \cref{II-75-106} in case II
with \((1-w)c,D_0\). For \(c_a>0\), use
\(y_j\le\check R_j\) and apply it in case I
with \((1-w)c,D_0+c_aC_{\rm vol}\).
This gives the bound with \(\Pi_k(C_{\rm maj},c'')\)
also in the second case; increasing its first
argument and applying \eqref{II-75-72} gives
\eqref{II-75-121}.

\section{A lower bound from the empty history}
\label{sec:partition-lower}

The history without large-field regions contributes a positive
measure. Its density gives the lower bound once the regular
interaction has been bounded. We first estimate the small-field
integral; the passage from a bound on bare history fields to an
almost-everywhere bound on the coarse density is treated separately.

Let \(1\le k\le K\), \(N_k=\#\mathbb T^{(k)}\), and
\(g_k=x_k^{-1/2}\). Retain the geometric and regularity assumptions
of \cref{II-7-40}, and suppose
\begin{equation}\label{eq:lower-smallness}
 \varepsilon_k\le\varepsilon_\flat,\qquad
 2\varepsilon_k\le\bar\varepsilon,\qquad
 p_0(g_k)\ge4B_3,\qquad g_k\le\bar\varepsilon/4.
\end{equation}
Write \(\mathcal V_k^\chi=\{V:\chi_k(\mathbb T)(V)=1\}\).
On this open set let \(U_k(V)\) denote the whole-torus
regular critical orbit over \(V\), and put
\begin{equation}\label{II-201-4}
 \mathsf{sf}_k(V)=
 \begin{cases}
 \exp[-L^kx_kA_K(U_k(V))],&V\in\mathcal V_k^\chi,\\
 0,&V\notin\mathcal V_k^\chi.
 \end{cases}
\end{equation}
The action is constant on the orbit, so this definition does
not require a choice of representative.

\begin{proposition}\label{II-9-53}
The orbit in \eqref{II-201-4} exists on
\(\mathcal V_k^\chi\). The function \(\mathsf{sf}_k\) is Borel,
strictly positive there, and bounded by one.
For
\[
 0<s\le\tfrac14\min\{\bar\varepsilon,B_3^{-1}\varepsilon_k\},
 \qquad
 \mathsf B_k(s)=\{V:|V(b)-1|<s\text{ for every positive bond}\},
\]
all density tests pass on \(\mathsf B_k(s)\), and
\begin{equation}\label{II-9-54}
 L^kA_K(U_k(V))<24B_3^2s^2N_k
                  \quad(V\in\mathsf B_k(s)).
\end{equation}
Consequently
\begin{equation}\label{II-9-56}
 \int\mathsf{sf}_k\,dV
 \ge
 \exp\left[-\left(3\log\frac7{s^3}
                         +24B_3^2s^2x_k\right)N_k\right].
\end{equation}
In particular \(s=g_k\) is allowed under
\eqref{eq:lower-smallness}, and the coefficient in parentheses
then equals \(\frac92\log x_k+3\log7+24B_3^2\).
\end{proposition}

\begin{proof}
Passage of the density test on the cube containing a corner
of a coarse plaquette gives
\(|\partial V-1|<2\varepsilon_k\): average the defining
background through \(k\) levels, as in the support proof
for \cref{II-5-73}. The ceiling \(\varepsilon_\flat\) is
the curvature-iteration ceiling used there.
For a whole-torus datum with coarse curvature less than
\(\varepsilon_{\rm r}\le\bar\varepsilon\), use the constant
tower \(\Omega_j=\mathbb T\). Its induced lift has tested
holonomies equal to one or to coarse plaquette holonomies,
by the lift calculation in the proof of \cref{II-7-40}.
The regular-background theorem \cref{I-58-22} therefore
gives the orbit, with bare curvature less than
\(B_3\varepsilon_{\rm r}L^{-2k}\). Its local continuous
representatives also show that its action is continuous.
Apply this with \(\varepsilon_{\rm r}=2\varepsilon_k\).
The set \(\mathcal V_k^\chi\) is open because there are
finitely many stably defined strict density tests.
Extending the positive continuous function on that set
by zero proves the Borel assertion.

If \(V\in\mathsf B_k(s)\), telescoping the four unitary
plaquette factors gives \(|\partial V-1|<4s\).
The same whole-torus argument now gives bare curvature
less than \(4B_3sL^{-2k}\).
There are \(3N_0=3L^{3k}N_k\) positive bare plaquettes,
and their action contributions are
\(\frac12|\partial U-1|^2\). Multiplication by \(L^k\)
gives \eqref{II-9-54}. For each density cube, apply
\cref{II-7-40} with \(\varepsilon_{\rm r}=4s\).
Its defining curvature is less than
\(4B_3sL^{-2k}\le\varepsilon_kL^{-2k}\), so its test passes.

The one-bond Haar ball has measure
\[
 m(s)=\frac2\pi\int_0^s w^2\sqrt{1-w^2/4}\,dw
       \ge\frac{s^3}{7}\qquad(0<s\le1).
\]
Indeed the integrand's square-root factor is at least
\(\sqrt3/2\), and \(\sqrt3/(3\pi)>1/7\).
There are \(3N_k\) positive coarse bonds. Integrating over
their product ball and using \eqref{II-9-54} proves
\eqref{II-9-56}. Finally
\(g_kp_0(g_k)=\varepsilon_k\) and
\(p_0(g_k)\ge4B_3\) put \(g_k\) in the stated range.
\end{proof}

\subsection{The density and the bound on its exponent}

Let \(h_\emptyset\) be the empty history at level \(k\).
Its regions \(\Omega_j,\Lambda_j\), \(j\le k\), are the
whole torus. Its layers are empty below \(k\), and its
last layer is \(\mathbb T\); its stored external fields
and its boundary factors are empty. Thus its reference
measure is \(dV\), and its regular exponent is the total
of the regular local families, with their vacuum values
subtracted and the running-coupling action terms added.
Denote this exponent, evaluated on a bare background \(U'\),
by \(\mathbf E_k(U')\).

The following statement isolates the exact inputs needed
from the history format and the estimate on the regular sum.
The vacuum constant \(E_{(k)}\) is that of
\cref{sec:vacuum-bounds}.

\begin{proposition}\label{II-201-44}
Suppose the empty history has a finite positive density
\(f_k\,dV=\tau_{h_\emptyset}\).
Suppose there is a Borel Haar-null set \(\mathcal N\) such
that \(f_k=0\) off \(\mathcal V_k^\chi\), and, for
\(V\in\mathcal V_k^\chi\setminus\mathcal N\), its admissible
background set is nonempty and contained in \(U_k(V)\).
Assume the value \(\mathbf e_k(V)=\mathbf E_k(U')\) is
independent of the admissible background and
\begin{equation}\label{eq:empty-history-format}
 f_k(V)=\mathsf{sf}_k(V)\exp[\mathbf e_k(V)+E_{(k)}]
       \qquad(V\in\mathcal V_k^\chi\setminus\mathcal N).
\end{equation}
Finally suppose that, whenever a bare field has nonzero
empty-history weight and \(U'\) belongs to its coarse
regular critical orbit,
\begin{equation}\label{eq:empty-history-support-bound}
 |\mathbf E_k(U')|\le C_5N_k.
\end{equation}
Then there is a real Borel function \(\mathsf{rg}_k\) on
\(\mathcal V_k^\chi\), bounded by \(C_5N_k\) almost everywhere,
such that
\begin{equation}\label{II-201-45}
 d\tau_{h_\emptyset}
   =\mathsf{sf}_k\exp[\mathsf{rg}_k+E_{(k)}]\,dV.
\end{equation}
It follows that
\begin{equation}\label{II-201-21}
 e^{-E_{(k)}}\rho_k\,dV
 \ \ge\ e^{-E_{(k)}}d\tau_{h_\emptyset}
 \ \ge\ e^{-C_5N_k}\mathsf{sf}_k\,dV .
\end{equation}
\end{proposition}

\begin{proof}
Take a nonnegative Borel representative of the density.
On \(\mathcal V_k^\chi\setminus\mathcal N\),
\eqref{eq:empty-history-format} makes it strictly positive.
Since both it and \(\mathsf{sf}_k\) are positive real numbers
and \(E_{(k)}\) is real, taking absolute values gives
\[
 \operatorname{Re}\mathbf e_k
       =\log f_k-\log\mathsf{sf}_k-E_{(k)}.
\]
Use this Borel expression to define \(\mathsf{rg}_k\)
off \(\mathcal N\), and set it equal to zero on
\(\mathcal N\cap\mathcal V_k^\chi\). Null sets do not
affect the density, so \eqref{II-201-45} holds.

It remains to transfer the support bound to Haar measure.
Set
\[
 B=\{V\in\mathcal V_k^\chi\setminus\mathcal N:
                         |\mathsf{rg}_k(V)|>C_5N_k\}.
\]
For a bare field \(U\) with nonzero empty-history weight,
if \(\mathcal M^kU\) lies outside \(\mathcal N\) then an
admissible \(U'\in U_k(\mathcal M^kU)\) exists.
Equation \eqref{eq:empty-history-support-bound} shows that
\(\mathcal M^kU\notin B\). The positive pushforward formula
therefore gives \(\tau_{h_\emptyset}(B)=0\).
The density in \eqref{II-201-45} is strictly positive on
\(B\); hence \(dV(B)=0\). For example, exhaust \(B\) by
the sets on which this density is at least \(1/n\).
On \(\mathcal N\) the exponent was set to zero, proving
the almost-everywhere bound on all of \(\mathcal V_k^\chi\).
The first inequality in \eqref{II-201-21} follows from the
positive history decomposition \cref{II-5-73}; the second
follows by integrating \(\mathsf{rg}_k\ge-C_5N_k\).
\end{proof}

Write \(\mathsf t(h)=\tau_h(\mathcal V_k)\).
Under the hypotheses just proved,
\begin{equation}\label{II-201-23}
 e^{-E_{(k)}}Z_{K,\ell,x_0}
   =\sum_{h\in\mathcal H_k}e^{-E_{(k)}}\mathsf t(h)
   \ge e^{-C_{\rm low}(x_k)N_k},\qquad
 C_{\rm low}(x)=C_5+\tfrac92\log x+3\log7+24B_3^2.
\end{equation}
Indeed retain the single nonnegative term \(h_\emptyset\)
and use \cref{II-9-53,II-201-44}.
If the vacuum constants are supplied through level \(K\),
then \(E_{(K)}=E_{(k)}+\mathfrak E^{\ge}_k\).
The upper vacuum estimate in \eqref{II-192-29} consequently
gives the same lower bound after normalization by
\(e^{-E_{(K)}}\), with coefficient
\(C_{\rm low}(x_k)+C_{\mathfrak E}\).


\section{The density at every level}
\label{sec:history-iteration}

We now collect the local terms and apply the measure identity
to the objects constructed above. The induction starts with
the finite Wilson measure. At each later level, integrability
follows from the positive change of variables in
\cref{II-166-30}; it is not an additional condition on the
new local expansion.

\subsection{A common choice of constants}
\label{sec:common-iteration-constants}

The parameters can be fixed before the number of steps, the couplings and
the torus. The order matters: the geometry and the fixed analytic radii
are chosen first, the bounds for the regular terms and their quadratic
coefficients next, and the upper bound on the coupling last.
The constants furnished by an earlier estimate always retain the values
given by that estimate.

\begin{proposition}\label{II-202-29}\label{II-202-36}
There are choices of the blocking factor, the geometric parameters,
the analytic radii, the decay rates, and numbers $E_0,B_\sharp>0$ and
$0<\gamma\le1/2$ for which the numerical hypotheses of the actual
small-field step, the strip and boundary constructions, the historical
Gaussian and normalizer estimates, and the history-mass estimates hold
simultaneously.
These choices are independent of the height, the length of the coupling
sequence and the torus. All the numerical hypotheses involving a
coupling hold uniformly for $x\ge\gamma^{-2}$; after the other constants
are fixed, $\gamma$ can be decreased.

In particular, the regular progression has the bound $E_0$, its
intrinsic increments have absolute value at most $B_\sharp$, and the
three contributions to the new boundary term satisfy, at every step,
\begin{equation}\label{II-202-32}
 \mathsf b_{\rm cs}^{\rm th}\le\tfrac13,\qquad
 \mathsf b_{\rm rj}^{\rm th}+\mathsf b_{\rm sr}^{\rm th}\le\tfrac13.
\end{equation}
The assertion concerns numerical hypotheses. The room required of a
finite torus is specified separately below.
\end{proposition}

\begin{proof}
We give the order of the choices and the inequalities which couple
different parts of the construction.

First choose integer exponents with
\[
 r\ge2,\qquad \mathfrak p_0\ge6r+1,\qquad
 q\ge41\mathfrak p_0+1,
\]
and the fixed reserve fractions. We may take
$0<\beta_{\rm s}\le1/5$ and then choose
\[
 \frac{1+\beta_{\rm s}}{2\sqrt2}\le\theta_{\rm oc}<\sqrt{2/3}.
\]
This interval is nonempty. Choose the Gaussian reserve
$1/16\le\vartheta_*<1$, and fix
$\mathsf C_{\rm mn}=18(1-\vartheta_*)^{-2}$ and the absolute
cluster constant $\mathsf C_{\rm cs}$ used in \eqref{II-118-5}.
The other reserve fractions are chosen strictly inside their
prescribed intervals. Choose the halo width and then an odd
$L\ge13$ above the finitely many combinatorial and analytic floors,
including
\[
 L\ge\frac{2\mathsf C_{\rm cs}(1+4\beta_{\rm s})}
                 {1-2\delta_{\rm sp}}.
\]
All constants depending only on $L$, including the covariance
and normalizer decay rates and gaps, are now fixed. With the constants
of \eqref{II-115-55} and \eqref{II-117-5}, put
\[
 N_{\rm gs}=64B_{\rm nm}^5\gamma_0^{-2}+c_{\rm r}+2,
 \qquad
 \bar c_P=\mathsf C_{\rm mn}^3\lambda_{\rm r}^{-1}N_{\rm gs}.
\]
These numbers depend only on $L$ and the Gaussian reserve.

For the rim estimates, let $c_M,\delta_M$ have the values in
\eqref{II-115-55}, and let $B_0,\delta_{\rm cv}$ be the constants
of \cref{I-73-90}. Fix
\[
 \begin{aligned}
 C_{\mathsf T}
   &=\max\left\{
       c_Mc_{\rm lat}(\delta_M/2),\,
       B_0^2c_{\rm lat}(\delta_{\rm cv}/2)
                  c_{\rm lat}(\delta_{\rm cv}/4)\right\},\\
 \mu_{\rm ch}
   &=\tfrac12\min\{\delta_M,\delta_{\rm cv}/2\},
       \qquad \epsilon_{\rm ch}=96^{-1/2},\\
 C_{\rm pre}&=c_{\rm coe}=\lambda_-=C_\flat^{\rm ch}=1.
 \end{aligned}
\]
These are positive finite constants fixed before
$\Theta_1,\varpi_T$ and $M$. In particular,
$C_{\mathsf T}\ge c_Mc_{\rm lat}(\delta_M/2)$ and
$\mu_{\rm ch}\le\delta_M/2$.

Next choose the response parameters and the powers
$\mathfrak M,R_1$ so that the background and local-response constructions
apply, with $\rho_{\rm gap}=L^{-2}R_1\mathfrak M$. Choose their small
radii in the order used in those constructions. In particular, choose
$\varepsilon_{\rm v}$ in the nonempty background window and then
$\bar\varepsilon,r_{\rm B}$ small enough that
\[
 \varepsilon_{\rm v}+\varepsilon_{\rm v}^2\le a_{\rm orb},
 \qquad
 B_3(\bar\varepsilon+C_\partial r_{\rm B})\le\varepsilon_{\rm v}.
\]
Here the response radii may impose further positive upper bounds on
$\bar\varepsilon,r_{\rm B}$. They involve only parameters already fixed.
For bounds depending on the regularity constant, use its
$\gamma$-independent upper bound $\bar C_{\rm reg}$.
The eventual restriction on $\bar B\gamma^2$ gives
$C_{\rm reg}\le\bar C_{\rm reg}$, so this choice does not depend on the
coupling chosen at the end.

Choose $\kappa$ above the entropy and cluster floors, including
$4\beta_{\rm s}\kappa\ge1$ and
\[
 \tfrac12(1-2\delta_{\rm sp})\kappa
 \ge2(1-2\delta_{\rm sp})+c_{\rm E}
   +\frac{\mathsf C_{\rm cs}}L
       \bigl[1+2c_{\rm E}+\mathsf C_{\rm cs}(\bar c_P+2)\bigr].
\]
Set
\[
 r_{\mathsf w}=(1-2\delta_{\rm sp})(\kappa-2),\qquad
 \kappa''=(1+4\beta_{\rm s})\kappa+1+c_{\rm E},\qquad
 \bar r_H=\kappa''+c_{\rm E}
                    +\mathsf C_{\rm cs}(\bar c_P+2).
\]
Choose $\bar\kappa$ above the required coarsened rates and require
\[
 \begin{gathered}
 \bar\kappa\ge1+\mathsf C_{\rm cs}\bar r_H+\log\mathsf C_{\rm cs},\\
 \bar\kappa/4\ge8(\kappa_T+1)+2c_{\rm E},\qquad
 3\bar\kappa/4\ge16(c_{\rm E}+1),
 \quad \kappa_T=(1+4\beta_{\rm s})\kappa.
 \end{gathered}
\]
Then choose $M_1$, the clearance, the grain and the mesh, followed
by the local response radii, the collar factor $\Theta_1$ and the
rim width $\varpi_T$. Include the requirements
\[
 \Theta_1>3C_{\mathsf T}/\epsilon_{\rm ch},\qquad
 C_{\mathsf T}e^{-\mu_{\rm ch}\varpi_T}\le\epsilon_{\rm ch}/12.
\]
Use this same clearance, grain and mesh in the regular and historical
normalizers, and use the same $r_{\mathsf N}$ below.
Their coefficient bound and real gap are also identified as in the
interior normalizer identity; thus $c_{\rm nf}=c_{\mathsf N}$.
The local inverse estimates are applied before these lengths are
enlarged; their positive rates are thus fixed. Choose the power
$M=L^{m'}$ last among these lengths.
All its requirements have the form
\[
 M\ge C,\qquad cM\ge C+p\log M,\qquad
 C M^p e^{-cM}\le\epsilon ,
 \qquad c,\epsilon>0,
\]
with fixed coefficients. They hold for all sufficiently large powers
of $L$. This includes the local-to-global normalizer error, the
resolvent error, and the polynomial losses in the historical Gaussian
construction. Also require
$(LM)^3\ge\mathsf C_{\rm mn}^{-1}\bar c_P$.
The normalizer radius and the boundary rate can then be chosen in
\[
 4\,\mathsf{ms}\le r_{\mathsf N}\le LM/4-1,\qquad
 L^{-1}C_{\rm lev}\log L\le\delta_\star^{\rm sb}
                         \le\mu_{\mathfrak h}M/4 .
\]
The lower bounds on $M$ make both intervals nonempty.

For clarity, two of the requirements included here are
$\mathsf{ms}\ge L^3$, $L\mathsf{ms}\le M/L^2$.
The small radii are subsequently reduced so that the complex
perturbation of the historical Gaussian form is bounded by its
available gap. In the notation of its construction this means
\[
 B_{\rm nm}\varepsilon_{\rm dc}\le\gamma_0/2,
 \qquad (LM)^3\varepsilon_{\rm dc}\le1.
\]
The exponentially small part of $\varepsilon_{\rm dc}$ is controlled
by the choice of $M$; its remaining part is controlled by the later
choices of $\alpha_1,\alpha_2$. Use the actual mean data
$(\theta_{\mathbf m},\varepsilon_{\mathbf m})
=(\vartheta_\mu^{\rm re},\varepsilon_\mu)$ of \cref{II-119-53},
with $C_e=\varepsilon_a=C_{\rm xt}$ in \eqref{II-115-82}.
The choice of $\theta_c$ below and the rim-width condition give
$\vartheta_{\rm col},\vartheta_{\rm rim}\le\epsilon_{\rm ch}/12$.
The far term is made at most $\epsilon_{\rm ch}/12$ by the final
coupling choice. Thus
$\theta_{\mathbf m}\le\theta_{\rm xt}:=\epsilon_{\rm ch}/3<1/16$.
Set $\varepsilon_{\rm xt}=\tfrac14(LM)^{-3/2}$.
In the same final coupling choice require
\[
 \sqrt2\,\vartheta_{\rm far}\varepsilon_a/C_e
                    \le\varepsilon_{\rm xt}.
\]
Together with $\vartheta_\mu^{\rm re}\le\theta_{\rm xt}$, this
is the external-mean admission of \cref{II-115-84} at
$(\theta_{\rm xt},\varepsilon_{\rm xt})$.
Impose also
\begin{equation}\label{eq:common-mean-admission}
 (1+c_{\rm r})\varepsilon_{\mathbf m}^{\,2}\le(LM)^{-3},
 \qquad \varepsilon_{\mathbf m}\le\vartheta_*-\theta_{\rm xt}.
\end{equation}
In \eqref{II-115-83} the collar and rim terms are multiplied by
$\varepsilon_{\rm dc}$. The remaining far term contains the factor
\[
 \exp[-\delta_M(LMR(x)-2)/2].
\]
The preceding choices of $M$ and the small radii control the former
terms, and the final choice of $\gamma$ controls the latter. These requirements give
$\theta_{\mathbf m}+\varepsilon_{\mathbf m}\le\vartheta_*$.
Increasing the torus or the number of steps makes no new choice
necessary.

Now choose $A_1,A_0,C_0,C_1$ in this order, retaining
$A_0\ge C_\chi A_1$ with the chart constant in
\eqref{eq:sf-chart-thresholds}, the other threshold comparisons,
and $C_1\ge C_0$ required by the tests and the radii.
Choose $\alpha_1$ sufficiently small and then $\alpha_0$ sufficiently
small, retaining
\[
 \alpha_0+\alpha_1\le a^*,\qquad
 \mathsf c_{\rm sh}M\alpha_0<\alpha_1.
\]
The shape constant is determined afterwards:
\[
 \mathsf c_\natural=\max\{1,\alpha_1^2/\alpha_0\}.
\]
Consequently $\alpha_1^2\le\mathsf c_\natural\alpha_0$ imposes no
lower bound on $\alpha_0$. Reduce the remaining radii
$\alpha_2,\alpha_3,\alpha_4,a_\flat$ in their order.
Their upper bounds are positive. Where such an upper bound contains
a difference, the preceding radius was chosen with a strict reserve;
for example $2\alpha_1<a^\flat$ leaves room for
$2(\alpha_0+\alpha_1)\le a^\flat$.

The rim fraction $\theta_c$ has two lower bounds,
\[
 u_a=\frac{3C_{\mathsf T}}{\epsilon_{\rm ch}\Theta_1},\qquad
 u_b=\frac{C_{\rm pre}(LM)^{3/2}}
              {\min\{\lambda_-,c_{\rm coe}\}\epsilon_{\rm ch}}\,
           e^{C_\flat^{\rm ch}\bar\kappa-\mu_{\rm ch}LM/4}.
\]
The earlier choices of $\Theta_1$ and $M$ make $u_a,u_b<1$.
Choose $\max\{u_a,u_b\}\le\theta_c<1$.
These are the only lower bounds to be reconciled with its upper
bound. The normalizer constants are now evaluated by their
defining formulas.

Choose $E_0$ large enough for the normalizer and regular-step bounds,
and then choose
\begin{equation}\label{eq:common-flow-bound}
 E_0\ge2E_N,\qquad
 B_\sharp\ge \mathsf G_\sharp
     :=C_{\rm am}r_\alpha^{-2}E_0,\qquad
 \bar B=\frac{L}{L-1}B_\sharp .
\end{equation}
Here $r_\alpha$ is the usable radius of the coefficient theorem.
The norm of the earlier array is bounded by
$L(E_0+\frac32\alpha_0^2B_\sharp)/(L-1)$.
Thus its bound is fixed before the ceiling on the coupling.
The actual-potential constants, including \(C_{\rm eff}\) in
\eqref{II-112-26}, are now evaluated with these choices of
\(E_0,B_\sharp,\bar B\). They too are fixed before \(\gamma\).
Choose the history-weight parameters in their permitted ranges;
in particular $0<c_{\rm Dep}\le c_\star$ and $0\le\theta<c_\star$.
Their mass and entropy thresholds are finite numbers at this stage.

It remains to choose $\gamma$. Put $y=\log x$ and
$F_p(x)=x^{-1/2}(\log x)^p$. The coupling-dependent upper bounds in
the construction are finite sums and products of
\begin{equation}\label{eq:common-vanishing-functions}
 F_p(x),\qquad (\log x)^{-(q-\mathfrak p_0)},\qquad
 x^a(\log x)^p e^{-c(\log x)^r},\qquad
 e^{-c(\log x)^{2\mathfrak p_0}},
 \quad c>0,
\end{equation}
with fixed exponents. They tend to zero as $x\to\infty$;
the third assertion uses $r\ge2$. Lower bounds are finite
thresholds on $y$, on positive powers of $y$, or on $R(x)$.
Take $\gamma$ sufficiently small for these inequalities to hold
for every $x\ge\gamma^{-2}$, including
\begin{gather*}
 \gamma^{-2}\ge2B_\sharp,\qquad
 \gamma^{-2}>\bar B,\qquad
 \gamma^{-2}>L(1+B_\sharp),\\
 \log\gamma^{-2}>2\mathfrak p_0,\qquad
 (L^{1/r}-1)\log\gamma^{-2}\ge\log(2L),\\
 C_0(\log\gamma^{-2})^{q-\mathfrak p_0}
                    \ge2\bar C_{\rm reg}A_0 .
\end{gather*}
This also imposes the small chart and logarithm radii, the
graded-field and strip inequalities, and the lower bounds on
$\mathfrak b(x)=A_1(\log x)^{\mathfrak p_0}$, including
$\theta_c\mathfrak b(x)\ge1$.
For example the common local-gauge and chart requirements are
\begin{gather*}
 \mathsf c_{\rm lgs}M\varepsilon(x)
 \le\min\{\alpha_0/2,\alpha_1,\alpha_2/2,a_{\rm nf}/2,
                          a^{\rm pc},a^{\rm cv},\varrho_{\mathbf H},1\},\\
 2L^2\bar C_{\rm reg}\varepsilon(x)
                   +C_{\rm seg}\mathsf C_{\mathcal A}\delta(x)
                    \le\bar\varepsilon,\\
 \mathsf C_{\mathcal A}\delta(x)
           \le\min\{\varepsilon_P^\diamond,\varepsilon_{\mathfrak h},
                                                    r_P^\circ/2\},\\
 C_\chi\max\{\mathsf c_{\rm fl},\mathsf C_{\rm xt}+1,
                      \mathsf C_{\mathcal A}\}\delta(x)<r_\chi .
\end{gather*}
The graded radii additionally require lower bounds on
$(\log x)^{q-\mathfrak p_0}$ and upper bounds on $F_q(x)$,
including $C_1^2C_0^{-1}F_q(x)\le\mathsf c_\natural$.
The strip errors have the form
$C(\log x)^{q-\mathfrak p_0}e^{-(\log x)^r}$.
These are among \eqref{eq:common-vanishing-functions} and the
corresponding positive-power lower bounds.
At step $t=k+1$, use the amplitude and rate of \eqref{II-155-56},
\[
 a_H=C_{\rm stp}C_{\mathrm D,\mathrm s}\varrho_{\mathbf H},
 \qquad \mu_H=\bar\kappa/4,
\]
and choose the transfer budgets by the rule
\begin{equation}\label{eq:common-boundary-budget}
 \rho_{{\rm bd},k}=\tfrac12 C_{\rm xt}\delta_k,\qquad
 r_{{\rm bd},k}=2\rho_{{\rm bd},k}/g_k
                =C_{\rm xt}\mathfrak b_k,\qquad \rho_J=1.
\end{equation}
The amplitude $a_H$ is fixed before $\gamma$. To check the lower
budget bounds uniformly, put $y=\log x_t$. The flow bound and
$x_t>\bar B$ give $x_0>\cdots>x_t$, so that
$\check R_t=R(x_t)\ge y^r$. Since $x_t\ge2B_\sharp$,
the same flow bound gives $y<\log x_k\le y+\log(2L)$.
For $y\ge1$, therefore,
\[
 \delta_k\ge A_1(2L)^{-1/2}e^{-y/2},\qquad
 \frac{e^{-\check R_t}}{\delta_k}
 \le\frac{\sqrt{2L}}{A_1}e^{-y^r+y/2}\longrightarrow0.
\]
With $N_e=C_{\rm E}$ and $N_b=6(LM)^3$, reducing $\gamma$ thus gives
\[
 \begin{aligned}
 128N_bN_e^2a_He^{-\check R_t}&\le\rho_{{\rm bd},k},\\
 128N_e^2\|\vartheta'_{\rm co}\|_\infty a_He^{-\check R_t}
                                      &\le\delta_k.
 \end{aligned}
\]
The chart bound $C_{\rm xt}\delta_k\le1/8$ also gives
$0<\rho_{{\rm bd},k}\le1/16$. Together with the rate choices above
and $\check R_t\ge12$, these verify \eqref{II-73-14},
\eqref{II-73-40} and \eqref{II-151-34} at every step.
Both fourth-potential ceilings retain their actual values:
\[
 \frac{(1+\beta_{\rm s})\delta(x)}{\alpha_0}
 \le\min\{\check s(K_\Lambda),\check s((1+\beta_{\rm s})^4K_\Lambda)\}.
\]
The regular progression uses the first of these; the nonconstant
radius profiles use the second.
The bound $\Theta_c(x,N_*)\le E_0/2$ in
\eqref{eq:progression-floor} is included.

We spell out the bounds for the history step, since they close the
iteration. Its regular and boundary size majorants have the forms
\begin{align*}
 S_{\rm rg}(x)
 &=C_{\rm eff}\left[
 \frac{A_1}{\alpha_0}(1+E_0+\alpha_0^2B_\sharp)F_{\mathfrak p_0}(x)
 +M^3A_1^3F_{3\mathfrak p_0}(x)
 +\frac{A_1}{C_0}F_q(x)F_{\mathfrak p_0}(x)\right],\\
 S_{\rm bd}(x)
 &=C_{\rm eff}(\mathsf B+\mathsf B_{\rm fr}^{\flat})
       \frac{A_1}{C_0}(\log x)^{-(q-\mathfrak p_0)},\\
 S_0(x)
 &=M^3K_\flat\mathsf C_{\mathcal A}A_1F_{\mathfrak p_0}(x)
  +M^3\bigl(K_\flat+2K_{\rm A}(r_P^\circ)^{-2}\bigr)
                 \mathsf C_{\mathcal A}^3A_1^3F_{3\mathfrak p_0}(x).
\end{align*}
Here $\mathsf B_{\rm fr}^{\flat}\le\mathsf B_{\rm fr}$: their
formulas differ by the factors $e^{C_3\kappa}$ and
$e^{C_3(\kappa+1)}$. The zero-height majorant is $S_0$.
All three tend to zero.
For the historical Gaussian use \eqref{II-117-5} with
\[
 \gamma_2=(LM)^{-3},\qquad
 (\mathfrak b_-,\mathfrak b_+)=(\theta_c\mathfrak b_k,\mathfrak b_k).
\]
Equations \eqref{II-115-55} and \eqref{II-117-5}, together with
the defect and mean bounds above, give
\[
 u_0\le\mathsf C_{\rm mn}(LM)^{-3}N_{\rm gs},\qquad
 c_P=\mathsf C_{\rm mn}^2\lambda_{\rm r}^{-1}(LM)^3u_0
       \le\bar c_P.
\]
The lower bound on $(LM)^3$ therefore implies
$\mathsf C_{\rm mn}u_0\le\lambda_{\rm r}$, as required in
\eqref{II-117-9}. Choose
$r_H=\kappa''+c_{\rm E}+\mathsf C_{\rm cs}(c_P+2)$.
Then $r_H\le\bar r_H$ and
$r_{\rm pol}=\kappa''+c_{\rm E}+\mathsf C_{\rm cs}$.
The preceding floors on $L,\kappa,\bar\kappa$ give
\[
 (\mathsf C_{\rm cs}/L)\bar r_H+c_{\rm E}\le r_{\mathsf w},
 \qquad
 e^{-(\bar\kappa-1)+\mathsf C_{\rm cs}\bar r_H}
                   \le\mathsf C_{\rm cs}^{-1}.
\]
Thus the three rate conditions in \eqref{eq:history-cluster-rates}
hold. For its remaining size condition, fix
\[
 K_{\rm ap}=\mathsf C_{\rm cs}^2
       e^{\mathsf C_{\rm cs}(\kappa''+\bar c_P+\bar r_H)},
 \qquad
 \Gamma_P(x)=\mathsf C_{\rm mn}^{-1}(LM)^{-3}
                    \theta_c^2A_1^2(\log x)^{2\mathfrak p_0}.
\]
Decreasing $\gamma$ imposes, uniformly for $x\ge\gamma^{-2}$,
\begin{equation}\label{II-202-31}
\begin{gathered}
 S_{\rm rg}(x),S_{\rm bd}(x)\le(12eL^3K_{\rm ap})^{-1},\\
 S_0(x)\le(6eL^3K_{\rm ap})^{-1},\qquad
 (LM)^3e^{-\Gamma_P(x)}\le(6K_{\rm ap})^{-1}.
\end{gathered}
\end{equation}
At a step $t=k+1$, the actual cutoff penalty is
$\gamma_P=\Gamma_P(x_k)$. Therefore
\[
 \mathsf b_{\rm cs}^{\rm th}
 \le K_{\rm ap}
       \left(e(S_{\rm rg}(x_k)+S_{\rm bd}(x_k))L^3
                         +(LM)^3e^{-\Gamma_P(x_k)}\right)
 \le\tfrac13 .
\]
For $k=0$ replace the sum of the two sizes by $S_0$; the
same conclusion holds.
Apply \eqref{eq:history-kept-bound} and
\eqref{eq:history-rejected-bound} at the target rate $\kappa''$, and put
\[
 c_{\rm kp}=c_{\rm kp}^{\rm th}
   =\tfrac12 C_3(LM)^6e^{C_3(\kappa''+1)}B_{\rm ch},
 \qquad
 c_{\rm rj}=c_{\rm rj}^{\rm th}=\tfrac12 C_3B_{\rm ch}.
\]
In \eqref{II-151-33} take $N_e=C_{\rm E}$,
$N_b=6(LM)^3$ and $\rho_J=1$. Define
\[
 \begin{aligned}
 c_3^\flat&=4\bigl[6(LM)^3(1+C_{\rm E})+C_{\rm E}/2\bigr],\\
 c_4^\flat&=C_{\rm xt}^2
   \bigl[e^2c_{\rm rj}+16C_{\rm E}(c_{\rm kp}+e^2c_{\rm rj})\bigr].
 \end{aligned}
\]
These constants are fixed before $\gamma$. Using
$\mathsf b_{\rm cs}^{\rm th}\le1/3\le1/2$ and
$e^{-(\check R_t-2)}\le e^2$ in \eqref{II-151-33} gives
\[
 \mathsf b_{\rm rj}^{\rm th}+\mathsf b_{\rm sr}^{\rm th}
 \le e^{-\check R_t}
          \bigl(c_3^\flat+c_4^\flat\mathfrak b_k^2\bigr).
\]
Writing $y=\log x_t$, the flow bound gives
$x_k\le(3L/2)x_t$; after $y\ge\log(2L)$,
$\mathfrak b_k^2\le A_1^2(2y)^{2\mathfrak p_0}$.
Since $\check R_t\ge y^r\ge y$ and
$(2y)^{2\mathfrak p_0}e^{-y/2}
 \le(8\mathfrak p_0/e)^{2\mathfrak p_0}$, their sum is at most
\[
 \gamma\left[c_3^\flat+
       c_4^\flat A_1^2(8\mathfrak p_0/e)^{2\mathfrak p_0}\right].
\]
Make this at most $1/3$. This proves \eqref{II-202-32};
the total is at most $2/3$, leaving the required boundary reserve.

Finally impose the finitely many thresholds from the history-mass
and history-counting estimates. They are numbers already determined
by the preceding constants. In particular the threshold for the kept
form is absorbed in $\log\gamma^{-2}$. All remaining strict
inequalities are imposed with a positive reserve.
All fixed constants and the rule \eqref{eq:common-boundary-budget}
are independent of the height, the coupling sequence and the torus.
Since every condition on $\gamma$ was imposed on the entire interval
$x\ge\gamma^{-2}$, decreasing $\gamma$ preserves them.
\end{proof}

\paragraph{Finite tori and the intrinsic coupling.}
Let $x_0,\ldots,x_k$ be an admissible sequence, prolonged to
$u_0,\ldots,u_K$, where $K\ge\max\{k,1\}$.
The room required for the history step is precisely
\begin{equation}\label{eq:iteration-torus-room}
 M R(u_K)\le L^m .
\end{equation}
The adjacent-scale inequality
$\check R_j\le L^{K-j}\check R_K$ then gives
$M\check R_j\le L^{m+K-j}$ at every level.
Our thresholds also give
$\check R_j\ge\max\{L^2,\mathsf oL+1\}$ and $M>6$, hence
\[
 L^2M\le L^{m+K-j},\qquad
 \mathsf oLM+6<L^{m+K-j},\qquad m'<m+K-j.
\]
Thus the mesh, grain, containing boxes and response charts fit on every
admissible torus. No further room condition enters the numerical
iteration.

The larger local-chart assertion requires
\begin{equation}\label{II-202-42}
 LM R(u_K)\le L^m.
\end{equation}
It then holds at every earlier level as well. It is not a consequence
of \eqref{eq:iteration-torus-room}: equality in that condition is
possible and violates \eqref{II-202-42}. Let $C_{\rm s}\ge2$ be a
fixed geometric constant with $2LC_{\rm s}\ge C_L$, where $C_L$ is
the room constant of \cref{sec:actual-step-global-domain}.
It is chosen before the couplings and the torus. A sufficient choice is
\[
 m\ge1+m'+
 \max\{\log_LR(u_K+\bar B),\lceil\log_L C_{\rm s}\rceil\}.
\]
This floor depends on the terminal coupling, but not on $K$.

With the step constructed, \eqref{II-202-32} and the uniform numerical
bounds give the same induction hypotheses at each new level.
This is the constants-only implication in
\cref{II-191-32,II-202-43}.
For the intrinsic flow, \eqref{eq:common-flow-bound} and
\cref{II-147-20} give, in order, a regular next family, its
lattice-independent real coefficient $b_{k+1}$, and
\[
 x_{k+1}=L^{-1}x_k-b_{k+1},\qquad
 |b_{k+1}|\le\mathsf G_\sharp\le B_\sharp.
\]
The flow is used while $x_k\ge\gamma^{-2}$.
The bound $\gamma^{-2}>L(1+B_\sharp)$ ensures that the first
term below this threshold is still greater than $1$.
Moreover
$|x_j-L^{-j}x_0|\le\bar B(1-L^{-j})$ and
$\gamma^{-2}>\bar B$ show that this first passage occurs in finitely
many steps. Thus the uniform construction applies to every preceding
truncation.

\subsection{The exponent and its two expressions}

Use the regular progression of \cref{II-147-20}, the boundary
terms of \cref{sec:boundary-assembly}, and the retained factors
of \cref{sec:formed-factors}. Their constants and the geometric
data are required to satisfy the quantitative hypotheses of
those constructions and of the actual history potentials and
Gaussian systems. The simultaneous choice is made in
\cref{sec:common-iteration-constants}. In this section it suffices
to keep those inequalities as hypotheses. In particular,
\(m'<m+K-j\) at every level used.

Write \(\mathcal S=\mathcal S(h)\) for the label of a history
of depth \(k\). A regular index at level \(j\) is a pair
\((z,X)\) with \(X\in\mathcal D_j\), \(X\subset\Lambda_j\),
and \(z\in X^{(j)}\cap\Lambda_j^{\sim-1}\). Denote this set
by \(\mathsf{Ix}^{\rm rs}_j(\mathcal S)\).
The edge and newly created boundary indices are those of
\cref{II-153-25}. At a glued field \(\mathsf U\), with stored
fields \(\mathfrak q\), define
\begin{equation}\label{II-99-12}
\begin{split}
 \mathbf E_k(\mathsf U)
  =\sum_{j=1}^k\biggl\{
    \sum_{(z,X)\in\mathsf{Ix}^{\rm rs}_j}
       [E^{(j)}(X,\mathsf U,z)-E^{(j)}(X,\mathbf1,z)]
       -b_jL^jA(\varphi_j,\mathsf U)\biggr\}.
\end{split}
\end{equation}
Here the regular term is evaluated at the level-\(j\) pair
of \(\mathsf U\). Put
\begin{equation}\label{II-99-14}
\begin{split}
 \mathbf E_k^{\rm bd}(\mathsf U,\mathfrak q)
  =\sum_{j=1}^k\biggl\{
   \sum_{X\in\mathsf{Ix}^{\rm mk}_j}
       B_j^{\rm mk}(X;\mathbf P^{(j)}_{X|\mathcal S}(\mathsf U),
                              (\mathfrak q|_X)|^{<j})
    +\sum_{X\in\mathsf{Ix}^{\rm ed}_j}
                              E^{(j),{\rm ed}}(X,\mathsf U)
                         \biggr\}.
\end{split}
\end{equation}
The first terms are restricted to their layered output
spaces. The boundary sum has no subtraction at the trivial
field and no profile term. The exponent is
\begin{equation}\label{II-99-19}
 A_{k,h}(\mathsf U,\mathfrak q)
 =-A(\beta_k(\cdot),\mathsf U)+\mathbf E_k(\mathsf U)
       +\mathbf E_k^{\rm bd}(\mathsf U,\mathfrak q)
       -\mathcal E_{k,\mathcal S}.
\end{equation}
Each of these finite sums is defined precisely when every
term used in it is defined. Terms with empty index sets
are omitted, without requiring an evaluation on an unused
domain. At level zero, \(A_{0,\emptyset}=-\beta A\).

For later use there is a second expression, before the
profile terms and the trivial values are separated:
\begin{equation}\label{II-99-21}
\begin{split}
 A_{k,h}(\mathsf U,\mathfrak q)
  ={}&-\beta A(\mathsf U)
       +\sum_{j=1}^k\sum_{X\subset\Lambda_j}
                      \sum_{z\in X^{(j)}}E^{(j)}(X,\mathsf U,z)\\
     &+\sum_{j=1}^k\sum_{X\in\mathsf{Ix}^{\rm mk}_j}
            B_j^{\rm mk}(X;\mathbf P^{(j)}_{X|\mathcal S}(\mathsf U),
                                       (\mathfrak q|_X)|^{<j})\\
     &+\sum_{j=1}^k
       [\mathsf N^{(j-1)}(\Omega_j,\Lambda_j)
                      +\mathsf{vn}^{(j)}\langle\Lambda_j\rangle].
\end{split}
\end{equation}
The sum \(X\subset\Lambda_j\) is over \(\mathcal D_j\).
Indeed, the regular and edge indices partition all the
pointed terms on these domains. The profile identity gives
\[
 -A(\beta_k(\cdot),\mathsf U)
          -\sum_{j=1}^kb_jL^jA(\varphi_j,\mathsf U)
       =-\beta A(\mathsf U).
\]
In \eqref{II-97-6}, the \(\mathsf{vl}\) sum is exactly the
sum of the trivial values subtracted in \eqref{II-99-12}.
These two cancellations prove \eqref{II-99-21} as an
identity of partial functions, with the same domain.
The trivial values are defined: the zero potential and
current give the trivial triple in every analytic core,
and the regular progression is defined on that core.

Let \(\Phi_{k,h}(\mathsf U;\mathfrak q)=\Phi_h(\mathsf U;\mathfrak q)\) be the sum of all
the retained exponents \eqref{eq:formed-factor-definitions}
from levels \(1,\ldots,k\). At a completed coordinate
\(\phi\), evaluate \(A_{k,h}\) and \(\Phi_{k,h}\) on the
glued adapted representatives. A value at \(\phi\) is
defined when the evaluation set is nonempty, the stored
fields lie in their tiers, and the corresponding finite
sum is defined and has the same value at every such
representative. The two sums are evaluated separately.
Gauge invariance of their constituents therefore
identifies the exponents themselves, rather than just
their exponentials. Set the density equal to zero wherever
\(\mathsf{Stay}^{\rm T}_h\chi_k(\Omega_k)=0\).

\subsection{Collecting one step}

Let \(h'=(h,\mathsf o_t)\), \(t=k+1\).
The normalizer has the following exact interior comparison.
It concerns its local functions and their domains, not only
the sum of its values on a real background.

\begin{lemma}\label{II-68-69}\label{II-68-70}
Use the same grain, clearance, mesh and normalizer clearance
in the historical and empty-history constructions.
If \(X\subset\Lambda_t\), their normalizer classes on \(X\)
coincide, and
\begin{equation}\label{II-68-71}
 N^{\rm un}_{h'}(X,p,z)=N^{\rm un}_{\emptyset}(X,p,z),
 \qquad
 N^\partial_{h'}(X;p)=N^\partial_{\emptyset}(X;p)
\end{equation}
as partial functions of the pair. On the analytic core,
the latter empty-history functions are the
constant-sequence functions of \cref{II-44-37}, before
the subtraction at the trivial pair.
\end{lemma}

\begin{proof}
Every layer region of the collared sequence contains
\(\Lambda_t\). Thus its cube-wise trace on \(X\) is \(X\),
as is the trace of the constant sequence.
The real factor, its perturbation and current are read
on the same bond set. The layer tests and containing-box
gauges defining \(\mathscr O^N(X;a)\) are identical.
The tiny-cube tests are identical too: if \(y^\sim\subset X\)
and \(y\) meets an integration bond, its other endpoint
lies in \(y^\sim\), so both endpoints lie in \(X\).
Whether such a nonpivot bond is an integration bond
depends only on whether an endpoint belongs to
\(\Lambda_t\), which is the same test in both settings.
For \(S_N(b)\subset X\) both endpoints belong to \(X\);
the same argument identifies the compression coordinates.
For \(X(c)\subset X\) it identifies the contributing
pivot bonds.

For every \(Y\subset X\), its determining tower
\(\Omega_\bullet(Y)\) lies in \(Y\). Hence
\[
 \Omega_n^\diamond\cap\Omega_n(Y)
           =\Omega_n(Y)=T_\ell\cap\Omega_n(Y).
\]
The two meets, their natural operator domains, and
their compressed matrices \(A_{Y,N}\) are therefore
identical. This uses the actual operator locality
of \cref{II-31-36}, including its domain assertion.
The logarithm series \eqref{II-68-16}, their finite
M\"obius differences, and the three cases in
\eqref{II-68-22} consequently agree. A pivot logarithm
uses only the descendants of its endpoint double
block, all contained in \(X(c)\); it is the same
function on the same branch. The tents also agree.
This proves \eqref{II-68-71}.

For the comparison with \cref{II-44-37}, the reading
of the triple is \((\mathbf U,L^{-3}\mathbf J)\),
and its restriction to \(Y\) is the reading used to
construct the constant-sequence operator there.
The clearances and coordinate sets are the same.
The flat gap in \cref{sec:gaussian-local-resolvents} is
\(\gamma^{\rm tr}\). Both constructions therefore use
\[
 B_{\mathsf a}=B_a
   =\max\{\gamma^{\rm tr},2\sqrt2C_ac_{\rm lat}(\theta_a)\},
\]
and their logarithm expansion points satisfy \(c_{\rm nf}=c_N\).
Thus the two logarithm series and finite sums are
literally the same expressions. At the trivial
triple their common pointed value is
\(\mathsf n^{(t)}(X,z)\).
The analytic-core bounds put the reading in the
stated normalizer class, as checked in the local
normalizer construction.
\end{proof}

The regular family at level \(t\) is the one already
constructed from the empty-history expansion:
\[
 E^{(t)}(X,P,z)
    =N^{\rm un}_{\emptyset}(X,P,z)
                      +F_{\emptyset}(X,P,z)-\mathsf n^{(t)}(X,z).
\]
It depends on the coupling prefix, \(t,X,z\), and the
pair, and is independent of \(h'\).
For \(X\subset\Lambda_t\), the same equality holds with
the local normalizer and cluster terms of \(h'\).
For the normalizer this is \cref{II-68-70};
for the pointed cluster terms it is
\cref{II-135-20}. These pointed identities include
equality of their natural domains. The same proposition
identifies the full cluster sum with the sum of its
pointed terms on the admitted analytic domain.

The normalizer and cluster sums in the new exponent
therefore split as
\begin{equation}\label{eq:iteration-interior-split}
\begin{split}
 N+E(a')
   ={}&\sum_{X\subset\Lambda_t}\sum_{z\in X^{(t)}}
                                     E^{(t)}(X,\mathsf U,z)
            +\mathsf{vn}^{(t)}\langle\Lambda_t\rangle\\
      &+\sum_X[B^{\rm nm}(X)+B^{\rm cs}(X)] .
\end{split}
\end{equation}
This is a finite-sum identity using \eqref{II-151-30}.
The bounds and support there restrict the last sum to
the admitted newly created boundary indices.
Adding \(B^\times(a')\) and \(B^{\rm sr}(a')\) supplies
the other two terms of \(B_t^{\rm mk}\).

The constant \(n_t\) in \eqref{eq:history-new-exponents}
is \(\mathsf N^{(k)}(\Omega_t,\Lambda_t)\).
To check the normalization, the regional fibre chart
contributes
\[
 \log c_H=(\log\sigma_{\rm H}+3\log g_k)B_k(\Omega_t)
          =(\log\sigma_{\rm H}-\tfrac32\log x_k)B_k(\Omega_t).
\]
The gauge factor \eqref{eq:regional-gauge-factor} contributes
\[
 \log c_z=-(L^3-1)\log\mathsf z(x_k,\varepsilon_k)
                                          \#_{k+1}(\Omega_t),
\]
and Gaussian integration contributes
\(\nu_{\mathfrak g}B_k(\Lambda_t)\).
These are the three terms of \eqref{II-97-4}.
Thus the first sum in \eqref{eq:history-new-exponents}
is precisely
\begin{equation}\label{II-184-14}
\begin{split}
 A^+_{h'}(\mathsf U,\mathfrak q^+)
   ={}&A_{k,h}(\mathsf U,\mathfrak q)
       +\mathsf N^{(k)}(\Omega_t,\Lambda_t)
       +\mathsf{vn}^{(t)}\langle\Lambda_t\rangle\\
      &+\sum_{X\subset\Lambda_t}\sum_{z\in X^{(t)}}
                                    E^{(t)}(X,\mathsf U,z)
       +\sum_{X\in\mathsf{Ix}^{\rm mk}_t}
          B_t^{\rm mk}(X;\mathbf P^{(t)}_{X|\mathcal S(h')}(\mathsf U),
                                            (\mathfrak q^+|_X)|^{<t}).
\end{split}
\end{equation}
Comparison with \eqref{II-99-21} gives
\begin{equation}\label{II-184-25}
                    A^+_{h'}=A_{t,h'} .
\end{equation}
This comparison includes equality of domains. Old
terms are the same functions on the same domains; only
their evaluation at the new glued field has changed.

For the other exponent, the definitions of the kept
factors and external pivot factors give
\begin{equation}\label{II-184-19}
 \Phi_{t,h'}(\mathsf U;\mathfrak q^+)
    =\Phi_{k,h}(\mathsf U;\mathfrak q)
         +\sum_XK^{\rm kp}_{h'}(X;\mathsf U,\mathsf a')
         +\sum_{c\in\Gamma_t^\partial}z_c(\mathsf U).
\end{equation}
Each retained coefficient is assigned once to its
strip component. Each external pivot is assigned once
to \(X(c)\). Thus this is also an equality of finite
sums, and it is the second identity in
\eqref{eq:history-new-exponents}.
The actual historical output membership gives every
factor its required domain at every admitted glued
representative.

\begin{proposition}\label{II-184-12}
Suppose the quantitative hypotheses of this section hold
through level \(t\), and suppose the history \(h\) has
the integrable density
\[
 \mathsf{Stay}^{\rm T}_h\chi_k(\Omega_k)
                       e^{A_{k,h}+\Phi_{k,h}}
\]
with respect to \(\mathscr M_h\).
Then the exponents of \(h'\) are defined at every point
of \(\{\mathsf{Stay}^{\rm T}_{h'}\chi_t(\Omega_t)>0\}\)
outside a Borel \(\mathscr M_{h'}\)-null set, and
\begin{equation}\label{II-184-13}
 \hat\tau^{\mathrm{co}\prime}_{h'}
    =\bigl(\mathsf{Stay}^{\rm T}_{h'}\chi_t(\Omega_t)
                         e^{A_{t,h'}+\Phi_{t,h'}}\bigr)
                                      \mathscr M_{h'} .
\end{equation}
The density is Borel, nonnegative and integrable,
and is strictly positive almost everywhere on that set.
Its local families obey the same bounds \(E_0,B,B_{\rm fr}\)
as at the preceding levels.
\end{proposition}

\begin{proof}
First suppose \(\Omega_1(h')\ne\emptyset\).
We verify the identities used by \cref{II-166-30}.
The free and completed coordinates of
\cref{II-58-74}, the regional fibre chart, and the strip
substitution give the two changes of variables and
their coordinate commutation. The support identity of
\cref{sec:history-support-identity} gives the factorization
of the tests, including the strict integration balls
and the new tiers. The fibre chart gives
\(j_c=c_He^{S_\Gamma}e^J\), while
\eqref{eq:regional-gauge-factor} gives the gauge term.
All these identities hold on the actual Borel domains.

The on-shell potential identity \cref{II-112-46} gives the two
equalities \eqref{II-166-58}. In particular its separate
identity for the old retained factors is needed here;
an identity only for the total exponent would not
identify the two new sums. The shifted field lies in the
old critical orbit and has the required transports at
the old stored sites. By \cref{II-168-27}, each of its
two old sums equals the sum evaluated at any adapted
representative. This gives the old coordinate evaluation
required in \eqref{II-166-58}.

The actual Gaussian system is supplied by
\cref{II-119-53}. On the whole torus its covariance
and mean are \eqref{II-114-33}. The potential bounds
give the hypotheses of \cref{II-118-20}, so
\eqref{II-166-22} is the convergent cluster identity.
The normalizer identity is \cref{II-69-19}.
At a real background its logarithm is the real
determinant logarithm, hence it gives exactly
\eqref{eq:history-normalizer-total}.
The kept and rejected forms sum to the conditional
Schur form, as in \cref{sec:history-form-coefficients}.
Finally \cref{II-181-24} gives the full strip sum
\eqref{II-166-23}, with the actual strip substitution.
Each statement includes definedness on the
historical domains just verified.

Equations \eqref{eq:iteration-interior-split}--\eqref{II-184-19}
identify the collected expressions with the canonical
exponents. The actual core and output membership
give their constituent domains for every glued
representative. The comparison in \cref{II-168-27} makes
their two sums representative independent, and
\cref{II-182-41} identifies the retained family at every
such representative, including at the first step.
All partial functions are Borel: their restrictions
are continuous, local series converge with common
majorants, and the adapted representatives can be
chosen Borel.
These are precisely the remaining evaluation
hypotheses of \cref{II-166-30}.
It proves \eqref{II-184-13}, its null-set assertion,
and integrability.

There is no extra assumption that almost every
coordinate point is realized by a bare field.
The proof of \cref{II-166-30} supplies this where
needed. More explicitly, the image of the Borel
domain of the bare coordinate map contains a Borel
set of full history measure. Off that set the
integrable density has integral zero on every
Borel subset. Its nonvanishing on the new support
therefore makes the complement of the realized
image a \(\mathscr M_{h'}\)-null set there.
Every representative at a remaining point has the
same exponents by the preceding orbit comparison.

If \(\Omega_1(h')=\emptyset\), all \(\Omega_j,\Lambda_j\)
are empty. There are no integration, external or
pivot coordinates, no retained factors, and no
local or boundary entries used in the exponent.
The profiles and extensive constants vanish.
The evaluation set is \(\{v_0\}\), and
\(A_{t,h'}=-\beta A(v_0)\), \(\Phi_{t,h'}=0\).
The two changes of variables reduce to their
zero-dimensional versions. Thus the same theorem
gives the density
\(\mathsf{Stay}^{\rm T}_{h'}\chi_t(\Omega_t)e^{-\beta A(v_0)}\).

The regular norm bound is \eqref{II-144-27}.
The boundary bound follows from \cref{II-153-25};
the newly created entries have the extra rate reserve,
and the old entries keep their values and domains.
For retained factors use \eqref{II-157-15}.
None of these constants is summed over the levels.
\end{proof}

\subsection{Induction and compatibility}

\begin{theorem}\label{II-184-28}\label{II-191-31}
Under the same quantitative hypotheses through level \(K\),
the canonical construction gives, for every
\(0\le k\le K\) and every \(h\in\mathcal H_k\),
\begin{equation}\label{II-184-29}
 \hat\tau^{\mathrm{co}\prime}_h
   =\bigl(\mathsf{Stay}^{\rm T}_h\chi_k(\Omega_k)
                      e^{A_{k,h}+\Phi_{k,h}}\bigr)\mathscr M_h .
\end{equation}
The right side is a nonnegative integrable Borel density
after assigning zero on a Borel null set.
The regular, boundary and retained-factor bounds are
uniform in \(k,h\) and the volume.
Truncation of a history or of the coupling sequence
preserves every earlier entry, its domain and its value.
\end{theorem}

\begin{proof}
At level zero the reference measure is product Haar
measure, the tests and the retained-factor exponential
are one, and \(A_{0,\emptyset}=-\beta A\).
This is the finite Wilson measure.
Suppose the assertion is known through level \(k\).
For each \(h'\in\mathcal H_{k+1}\), the initial history
\(h=h'|_k\) has the asserted density.
Its functions are exactly the functions used as old
data in the next step: the regular progression uses
only the coupling prefix; a boundary entry of level
\(j\) uses only the history through \(j\); and the
retained factor of level \(j\) has its fixed layered
domain and its stored fields from that step.
Restricting the table of radii gives this same domain,
by \cref{II-157-13}. The completed coordinate space,
reference measure, tests and evaluation set likewise
depend only on \(h\), not on its later extension.
Thus the induction hypothesis is exactly the old
density hypothesis of \cref{II-184-12}.

Apply that proposition to every extension, including
the empty-region case. This proves the next level,
with the same bounds. Finiteness of the level range
completes the induction. Each local construction is
specified by its canonical series, fixed chart and
fixed branch. Equality on a common prefix follows
inductively from these same formulas; no coupling
or family at a later level is used to construct an
earlier one.
\end{proof}

\begin{corollary}\label{II-191-32}\label{II-202-43}
Fix the constants of \cref{II-202-29}.
Let \(x_0,\ldots,x_K\) be admissible and suppose
\(MR(x_K)\le L^m\).
Then all conclusions of \cref{II-184-28} hold at
every level \(0\le k\le K\), with these same constants.
\end{corollary}

\begin{proof}
The constant choice imposes each coupling inequality
on the entire interval \(x\ge\gamma^{-2}\).
The adjacent-scale comparisons therefore apply at
every step. The room argument following
\cref{II-202-29} gives the required containing boxes,
meshes and strict inequality \(m'<m+K-k\).
The actual potential, Gaussian and strip estimates
then have the numerical hypotheses used in
\cref{II-184-12}. In particular,
\eqref{II-202-32} bounds the sum of the cluster,
rejected and strip boundary contributions by \(2/3\).
Together with \(\delta_NM\ge\kappa_T\), this gives
\(B_T=24(LM)^3E_{\rm nf}+1\) in
\eqref{II-154-45}. Thus the same regular,
boundary and retained-factor bounds apply at each
step, and \cref{II-184-28} applies.
\end{proof}

\section{Bounds for the constructed history measures}
\label{sec:actual-history-bounds}

We apply the history estimates to the density constructed in
\cref{II-184-28}. Throughout this section the constants satisfy
\cref{sec:common-iteration-constants}, the couplings are admissible
through \(K\), and the torus satisfies
\[
                         MR(x_K)\le L^m.
\]
Write \(y_j=\log x_j\), \(g_j=x_j^{-1/2}\),
\(N_j=\#\mathbb T^{(j)}\), and
\(\mathsf t(h)=\tau_h(\mathcal V_k)\).
All vacuum constants below are those of the same canonical
regular progression.

\subsection{The mass and the sum over histories}

\begin{proposition}\label{II-203-23}
For every \(1\le k\le K\) and every history \(h\in\mathcal H_k\),
\[
 e^{\mathcal E_{k,\mathcal S(h)}}\mathsf t(h)
 \le
 \exp\!\left[C_{\rm maj}
        \left(N_k+\sum_{j=1}^k\#_j(Z_j(h))\right)
             -c_\star\sum_{t=1}^kp_0(g_t)^2
                                      n_t^{\rm cr}(h)\right].
\]
The constants \(C_{\rm maj},c_\star>0\) are independent of the
history, the lattice, the coupling sequence and \(k\).
\end{proposition}

\begin{proof}
The density hypothesis of \cref{II-196-9} is
\eqref{II-184-29}. Equation \eqref{II-99-19} gives
\eqref{II-196-5} at every admitted glued evaluation.
For such a field the actual core membership and
layer estimates of \cref{sec:history-actual-potentials}
give the field hypotheses of \cref{II-101-34}.
The regular families satisfy \eqref{II-144-27},
and the boundary families satisfy \cref{II-153-25},
on precisely those domains.
Thus the regular and boundary sums have the bound
used in \cref{II-196-9}, with the same \(C_5\).

The retained factors are exactly those of
\cref{II-157-13}; they are not replaced by absolute-value
estimates for arbitrary quadratic polynomials.
The actual all-bond gauge class is supplied by
\cref{II-65-44}, with its fixed radius margins.
The pivot ceilings and the kept-form positivity
hypotheses of \cref{sec:history-mass} follow from
the chosen constants. Therefore \cref{II-195-31}
identifies the pivot determinants, and the stored-form
argument identifies the coercive component forms.
The history support and failed-test hypotheses are
those of \cref{sec:history-support-identity}.
All the hypotheses of the integrated mass estimate
\cref{II-195-19} are now satisfied, proving the
displayed bound.

There is no change of vacuum data on truncating
at \(k\). The regular families through \(k\), their
trivial values, and the numbers \(\mathsf n^{(j)}\)
are prefix consistent. The three coefficients of
\eqref{II-192-10} likewise use only their own level.
Consequently every term in \eqref{II-97-6} with
index \(j<k\) is the same whether the construction
ends at \(k\) or at \(K\). This proves the claimed
identification of \(\mathcal E_{k,\mathcal S(h)}\).
\end{proof}

The vacuum deficit estimate now converts the preceding
bound to the fixed normalization \(e^{-E_{(k)}}\).
Set
\[
 f(h)=e^{\mathcal E_{k,\mathcal S(h)}}\mathsf t(h),
 \qquad
 e(\mathcal S)=-E_{(k)}-\mathcal E_{k,\mathcal S}.
\]
Then
\begin{equation}\label{eq:actual-vacuum-summation}
 f(h)e^{e(\mathcal S(h))}=e^{-E_{(k)}}\mathsf t(h),
 \qquad
 e(\mathcal S)=-\sum_{j=1}^k\mathfrak E'_j(\mathcal S).
\end{equation}
By \cref{II-192-21}, this is a pair satisfying
\eqref{eq:history-summation-input} with \(c=c_\star\).
Thus \cref{II-75-111} gives
\begin{equation}\label{II-203-24}
 \sum_{h\in\mathcal H_k}e^{-E_{(k)}}\mathsf t(h)
      e^{\theta p_0(g_k)^2\mathsf N(\mathcal S(h))
                           +\varkappa\#_k(Z_k(h))}
 \le\Pi_k(C_{\rm maj},(1-\eth)(c_\star-\theta))
\end{equation}
for \(0\le\theta<c_\star\), \(\varkappa\ge0\),
\(0<\eth<1\), under the threshold
\[
 y_K\ge
 Y^{\rm II}\!\left(\eth(c_\star-\theta),
            48^3(D_0+\varkappa+\lambda_L)\right),
 \qquad
 D_0=C_{\rm maj}+27(2+\log(2L))C_\flat .
\]
The stronger finite-product bound \eqref{II-75-72}
holds when its additional \(Y^{\rm S}\) threshold
is imposed.
These free parameters are fixed before decreasing
\(\gamma\); the inequality is not asserted uniformly
over all their values at a single threshold.

For the mark \(\Theta^{(k)}\) defined before
\eqref{II-75-120}, use the same substitution
\eqref{eq:actual-vacuum-summation}.
It gives, with \(c''=(1-\eth)(1-w)c_\star\),
\begin{equation}\label{II-203-25}
\begin{aligned}
 \sum_h e^{-E_{(k)}}\mathsf t(h)e^{\Theta^{(k)}(h)}
       &\le\Pi_k(C_{\rm maj},c'')
                    \le e^{(C_{\rm maj}+2)N_k},
                         &&c_a=0,\\
 \sum_h e^{-E_{(k)}}\mathsf t(h)e^{\Theta^{(k)}(h)}
       &\le\Pi_k(C_{\rm maj}+\lambda_L,c'')
                    \le e^{(C_{\rm maj}+\lambda_L+2)N_k},
                         &&c_a>0.
\end{aligned}
\end{equation}
The parameter ranges and the respective thresholds are
exactly those of \eqref{II-75-120} and \eqref{II-75-121},
with \(c=c_\star\). In particular the second line
uses \(2\mathfrak p_0>4r+1\).

\subsection{The partition function}

\begin{proposition}\label{II-203-26}
Put
\[
 C_{\rm low}(x)=C_5+\tfrac92\log x+3\log7+24B_3^2.
\]
For the constructed measures,
\begin{equation}\label{eq:actual-partition-bounds}
 e^{-C_{\rm low}(x_k)N_k}
 \le e^{-E_{(k)}}\mathsf t(h_\emptyset)
 \le e^{-E_{(k)}}Z_{K,\ell,\beta}
 \le e^{(C_{\rm maj}+2)N_k}.
\end{equation}
As measures on the coarse fields,
\begin{equation}\label{II-203-27}
 e^{-E_{(k)}}\rho_k\,dV
 \ge e^{-E_{(k)}}\tau_{h_\emptyset}
 \ge e^{-C_5N_k}\mathsf{sf}_k\,dV .
\end{equation}
With the full vacuum constant one has
\begin{equation}\label{II-203-28}
 e^{-[C_{\rm low}(x_k)+C_{\mathfrak E}]N_k}
 \le e^{-E_{(K)}}Z_{K,\ell,\beta}
 \le e^{[3y_k+C_{\mathfrak E}+C_{\rm maj}+2]N_k}.
\end{equation}
\end{proposition}

\begin{proof}
For the empty history all regular regions are the
torus. Its external fields, boundary terms and
retained factors are empty, and its completed
reference measure is coarse product Haar measure.
The evaluation set is the regular critical orbit
of the coarse field, with its adapted representatives.
It is nonempty off the exceptional set furnished by
\cref{II-184-28}. Formula \eqref{II-99-19} reduces to
\eqref{eq:empty-history-format}, since
\(\beta_k=L^kx_k\) and
\(\mathcal E_{k,\mathcal S(h_\emptyset)}=-E_{(k)}\).
The sum \(\mathbf E_k\) is representative independent
by the same evaluation rule.

On the support of an empty history,
\cref{II-101-34} has only its final layer:
the earlier layers and all boundary contributions
are empty. It gives
\(|\mathbf E_k|\le C_5N_k\).
The actual core membership verifies its hypotheses
for every admitted critical background.
Hence \cref{II-201-44} applies and gives
\eqref{II-203-27}; integrating and using
\cref{II-9-53} gives the first inequality of
\eqref{eq:actual-partition-bounds}.
The positive history decomposition \cref{II-5-73}
gives the middle inequality and
\[
 Z_{K,\ell,\beta}=\sum_{h\in\mathcal H_k}\mathsf t(h).
 \]
Apply \eqref{II-75-114} to
\eqref{eq:actual-vacuum-summation}. Its threshold
is part of the chosen small-coupling ceiling.
This gives the upper bound.

Finally \(E_{(K)}=E_{(k)}+\mathfrak E_k^{\ge}\).
Multiplication by \(e^{-\mathfrak E_k^\ge}\)
and the two bounds in \eqref{II-192-29} prove
\eqref{II-203-28}.
\end{proof}

\section{Proof of the finite renormalisation theorem}
\label{sec:finite-renormalisation-proof}

\begin{proof}[Proof of \cref{II-203-9}]
Choose the constants as in
\cref{sec:common-iteration-constants}.
They are fixed before \(K,m\) and the coupling sequence.
The running-coupling comparison gives
\(\check R_j\le L^{K-j}R(x_K)\).
Consequently the single torus hypothesis implies
\[
 M\check R_j\le L^{m+K-j}\qquad(0\le j\le K).
\]
The lower radius bounds in the constant choice also
give \(m'<m+K-j\). Thus every finite-volume construction
and domain used above is available at these levels.
This argument uses no stronger chart-size hypothesis.

The bare Wilson density is continuous, strictly positive
and at most one on a compact product of groups, so its
mass belongs to \((0,1]\). Its gauge invariance and
the chosen curvature ceilings allow \cref{II-5-73};
this proves the positive history decomposition.
The exact density and its evaluation assertions are
\cref{II-184-28}. Its marginal on \(v_k\) is
\(\tau_h\), by the definition of the completed
coordinates. The regular, boundary and retained-factor
bounds are respectively \cref{II-147-20,II-153-25,II-157-13},
with the actual domain verifications used in that induction.

\Cref{II-203-23} proves the individual history bound.
\Cref{II-203-26} proves the partition and measure
bounds, with
\(c_{\rm low}=C_5+3\log7+24B_3^2\).
All constants are independent of the number of steps,
the history, the volume and the coupling prefix.
The weighted estimates follow separately from the
same mass bound under their stated thresholds.

Finally the coefficient-determined recursion and its
finite first passage are proved in
\cref{sec:canonical-flow}. Every prefix before first
passage has \(x_j\ge\gamma^{-2}\) and
\(|b_j|\le B_\sharp\). It therefore satisfies the
hypotheses already used, and each coefficient depends
only on the preceding prefix.
\end{proof}

\section{Local source energies and stability}\label{sec:local-source-energies}

At level \(j\), let \(\mathcal D_j\) be the nonempty wall-connected
unions of cubes of side \(ML^{j-K}\). For such a union \(X\),
\(d_j(X)\) is the least number of cube edges in a connected union
of edges contained in \(X\) and meeting every cube of \(X\);
a single corner is allowed and has size zero.
Write \(X^{\sim1}\) for its enlargement by one layer of cubes.
These are the localization domains and their size in
Proposition~\ref{II-18-26}.

The following statement fixes the source ball before the cutoff.
The effective-density construction that proves it also applies to
arbitrary admissible finite real coupling sequences.

\begin{theorem}[Local source energies and stability]\label{III-111-27}
There are choices of the averaging map and the constants preceding
the coupling threshold, including an odd integer \(L\ge13\), integers \(r\ge2\) and \(m'\ge0\), and constants
\[
 M=L^{m'},\quad B_\sharp\ge0,\quad C_{\mathrm s}\ge2,\quad
 E_0>0,\quad \kappa_e>\log9216+2,\quad C_{\mathrm S}>0
\]
with the following property. For every \(\varrho>0\), there is
\(\gamma_\varrho\in(0,\tfrac12]\) such that, after fixing the
canonical auxiliary data at this threshold, the assertions below
hold for \(0<g\le\gamma_\varrho\). The recursion and all constructed
objects use this same threshold assignment.

For every \(\beta>X_g\), the coefficient-determined recursion has
a finite first-passage index \(K\ge1\), and
\(g^{-2}<x_K(\beta)\le X_g\). Define, for \(x>1\),
\[
 R(x)=\min\{L^u:u\in\mathbb Z_{\ge0},\ L^u\ge(\log x)^r\}.
\]
Choose any integer \(m\) satisfying
\begin{equation}\label{III-1-14}
 m\ge 1+m'+
 \max\left\{\log_L R\bigl(L(g^{-2}+\bar B)\bigr),
                  \lceil\log_L C_{\mathrm s}\rceil\right\},
 \qquad \ell=2L^m .
\end{equation}
There are a real vacuum constant \(E_{K,\ell,\beta}\), local
source energies \(e_j(X;w)\), \(1\le j\le K\),
\(X\in\mathcal D_j\), and a function \(\widetilde Z(w)\), defined
for \(\|w\|_{\flat;T_\ell}<2\varrho\), such that
\begin{equation}\label{III-111-28}
 Z_{K,\ell,\beta}[w]
 =\exp\left(E_{K,\ell,\beta}
       +\sum_{j=1}^K\sum_{X\in\mathcal D_j}e_j(X;w)\right)
       \widetilde Z(w).
\end{equation}
They have the following properties.
\begin{enumerate}[label=(\roman*)]
\item The functions \(e_j(X;\cdot)\) and \(\widetilde Z\) are
holomorphic on the source ball. Each \(e_j(X;w)\) depends on
\(w\) only on \(X^{\sim1}\), and \(e_j(X;0)=0\). They satisfy
\[
 e_j(X;\bar w)=\overline{e_j(X;w)},\qquad
 \widetilde Z(\bar w)=\overline{\widetilde Z(w)}.
\]
\item The source energies and the remaining partition function obey
\begin{align}
 |e_j(X;w)|
 &\le 2M^3E_0\frac{\|w\|_{\flat;T_\ell}}{2\varrho}
          e^{-\kappa_e d_j(X)}
 \le C_{\mathrm S}e^{-\kappa_e d_j(X)},\label{III-111-21}\\
 |\widetilde Z(w)|&\le e^{C_{\mathrm S}\ell^3},\label{III-111-24}\\
 \widetilde Z(0)&\ge
       e^{-C_{\mathrm S}(1+\log g^{-2})\ell^3}.\label{III-111-29}
\end{align}
\item For two side lengths \(\ell,\ell'\) satisfying
\eqref{III-1-14}, with the same \(K,\beta\) and coupling sequence,
\[
 \left|\frac{E_{K,\ell,\beta}}{\ell^3}
       -\frac{E_{K,\ell',\beta}}{(\ell')^3}\right|
       \le C_{\mathrm S}.
\]
\end{enumerate}
All constants displayed before \(\varrho\) can be fixed
independently of \(\varrho,g,K,\beta,\ell,w\), of the level and
of the localization domain. The ceiling \(\gamma_\varrho\)
may decrease when the source radius increases.
\end{theorem}

The theorem bounds each local source energy separately.
It does not bound their sum uniformly in \(K\).
The volume bound applies to \(\widetilde Z\), after the vacuum
energies have been extracted. The positive lower bound is at
zero source. No limit as \(K\to\infty\) or zero-free region for
complex sources is asserted.

The preceding construction supplies the critical backgrounds,
the unsourced Gaussian expansion and the positive history
measures of \cref{II-203-9}. A varying complex source changes
the quadratic form as well as the interaction. Its quadratic
Taylor term must therefore be retained in the local factors,
with a relative form bound that preserves damping. The linear
and higher-order terms enter a cluster expansion about the
complex conditional mean.

The proof first controls this source decomposition for one step,
including the separate first step. It then constructs the
source-dependent regular terms, boundary terms and retained factors
through all levels. An identity of complex measures identifies
the resulting density with the sourced bare measure after block
averaging. The large-field estimates give the bound on
\(\widetilde Z\); the lower bound at zero source uses the
positive construction of \cref{II-203-9}. Finally, the first-passage
estimate gives the stated range of bare couplings and torus sizes.

\section{Sources and finite history measures}
\label{sec:sources-history-measures}

The source changes the plaquette coefficient in the bare density. The
coupling sequence, tests, history labels, coordinate maps and reference
measures remain those of the real construction. We first compare the
source with the coupling at every scale. We then record the finite
measure identities which the sourced expansion must represent.

\subsection{The source class}

All distances in a source norm are physical distances on
$T_\ell=(\mathbb R/\ell\mathbb Z)^3$, where $\ell=2L^m$.
Put $a_i=L^{-i}$. A unit of length at level $k$ has physical length
$a_{K-k}$. The number $M=L^{m'}$ is fixed independently of $K$ and
$\ell$. For a nonempty subset $S\subset T_\ell$, let
\[
 |w|_S=\sup_{x\in S}|w(x)|,\qquad
 \operatorname{Lip}_\infty(w;S)
 =\sup_{x\ne y\in S}\frac{|w(x)-w(y)|}{d_\infty(x,y)}.
\]
The second supremum is zero when $S$ is a singleton. Write
$\mathsf{Lp}(S)$ for the complex functions with finite norm
\begin{equation}\label{eq:source-subset-norm}
 \|w\|_{\flat;S}
 =\max\left\{|w|_S,\frac{40M}{\sqrt3}
                    \operatorname{Lip}_\infty(w;S)\right\}.
\end{equation}
For $\varrho_0,\varrho_1>0$, set
\begin{equation}\label{III-16-18}
 \mathsf{Src}_S(\varrho_0,\varrho_1)
 =\left\{w\in\mathsf{Lp}(S):
       |w|_S<\varrho_0,\quad
       \frac{40M}{\sqrt3}\operatorname{Lip}_\infty(w;S)
                      <\varrho_1\right\}.
\end{equation}
It is the open unit ball of the norm
\begin{equation}\label{III-16-19}
 \|w\|_{\mathsf{Src}}
 =\max\left\{\frac{|w|_S}{\varrho_0},
      \frac{40M}{\sqrt3\varrho_1}
                       \operatorname{Lip}_\infty(w;S)\right\}.
\end{equation}
In particular, equal radii give $\|w\|_{\flat;S}<\varrho$.

\begin{definition}\label{III-16-20}
A function $\Phi$ on $\mathsf{Src}_S(\varrho_0,\varrho_1)$ is
\emph{holomorphic} if, for every $w_0,f_1,\ldots,f_n\in\mathsf{Lp}(S)$,
\[
 z\longmapsto\Phi\left(w_0+\sum_{i=1}^n z_if_i\right)
\]
is holomorphic on the open set where its argument belongs to the class.
Let $\mathsf{Bc}$ be the set of bare plaquette barycentres. If
$\mathsf{Bc}\subset S$, a function depends on the source \emph{only on
$Y$} when its value is unchanged by changing $w$ away from the
barycentres in $Y$.
\end{definition}

Restriction to a subset is contractive. The source class is convex,
balanced and invariant under conjugation; taking a real or imaginary
part does not increase its norm. These facts follow by applying the
triangle inequality to values and differences. They also show that
every finite dimensional affine slice in the definition is open.
For finite $S$, the definition is ordinary holomorphy on an open
subset of $\mathbb C^S$. Thus a holomorphic function of the finitely
many plaquette values pulls back to a holomorphic function of sources
on the torus.

The Euclidean norm used for presented sources is compatible with
\eqref{III-16-17}. Namely, with
\[
 \|f\|_{\sharp;S}
 =\max\{|f|_S,40M\operatorname{Lip}_2(f;S)\},
\]
the inequalities $d_\infty\le d_2\le\sqrt3d_\infty$ give
\[
 \|f\|_{\flat;S}\le\|f\|_{\sharp;S}
                   \le\sqrt3\|f\|_{\flat;S}.
\]
Consequently a presented source $w=c+\sum_i z_if_i$ satisfies
\[
 \|w\|_{\flat;S}\le |c|+\sum_i|z_i|\|f_i\|_{\sharp;S}.
\]

\begin{lemma}\label{III-16-27}
Suppose that $\Phi$ is holomorphic on a source class and
$|\Phi|\le C$. If $\Phi(0)=0$, then
\[
 |\Phi(w)|\le C\|w\|_{\mathsf{Src}}.
\]
Without the condition at zero,
$|\Phi(w)-\Phi(0)|\le2C\|w\|_{\mathsf{Src}}$.
\end{lemma}

\begin{proof}
For $w\ne0$, restrict to $\lambda\mapsto\Phi(\lambda w)$ on
$|\lambda|<\|w\|_{\mathsf{Src}}^{-1}$ and apply the scalar Schwarz
lemma. The second assertion applies the first to $\Phi-\Phi(0)$.
\end{proof}

\subsection{Size of the source at each scale}

Fix real couplings $x_0,\ldots,x_K$, with $x_0=\beta$ and
\[
 x_j=L^{-1}x_{j-1}-b_j,\qquad |b_j|\le B_\sharp,
 \qquad \bar B=\frac{LB_\sharp}{L-1},\qquad x_K>\bar B.
\]
These are the couplings of the real construction. For a source on
the torus or on $\mathsf{Bc}$, define
\begin{equation}\label{III-16-30}
 x_j^w=x_j-L^{K-j}w,
 \qquad v_j^w=-\frac{L^{K-j}w}{x_j}.
\end{equation}
For a plaquette weight $v$, write
\[
 \operatorname{sw}_j(v)
 =\max_{d_\infty(x(p),x(p'))\le Ma_{K-j}}
                  |v(p)-v(p')|.
\]
This comparison uses the physical size of a level-$j$ block.

For a history $h$ of level $k$, its real profiles $\varphi_j$ satisfy
$0\le\varphi_j\le1$. The associated bare plaquette coefficients are
\begin{equation}\label{III-24-6}
 \beta_k=\beta-\sum_{j=1}^kL^jb_j\varphi_j,
 \qquad \beta_k^w=\beta_k-L^Kw.
\end{equation}
The profiles and increments are independent of the source.

\begin{proposition}\label{III-16-31}\label{III-24-8}
For $0\le j\le K$ and $0\le k\le K$, the preceding definitions give
\begin{equation}\label{III-16-33}
 x_j\ge L^{K-j}(x_K-\bar B),\qquad
 |v_j^w|\le\frac{|w|_S}{x_K-\bar B},\qquad
 \operatorname{sw}_j(v_j^w)
       \le\frac{M\operatorname{Lip}_\infty(w;S)}{x_j}.
\end{equation}
The last two inequalities concern the restriction to $\mathsf{Bc}$.
In particular, for $w\in\mathsf{Src}_S(\varrho_0,\varrho_1)$ their
right sides are strictly smaller than
$\varrho_0/(x_K-\bar B)$ and $\sqrt3\varrho_1/(40x_j)$.
Moreover,
\begin{equation}\label{III-24-9}
 \beta_k(y)\ge L^k(x_k-\bar B)
                \ge L^K(x_K-\bar B)>0.
\end{equation}
Writing $\mathsf a(w)=|w|_S/(x_K-\bar B)$, we have
\begin{equation}\label{III-24-10}
 \begin{split}
  |\beta_k^w-\beta_k|&\le\mathsf a(w)\beta_k,\\
  \operatorname{Re}\beta_k^w&\ge(1-\mathsf a(w))\beta_k,
  \qquad |\operatorname{Im}\beta_k^w|\le\mathsf a(w)\beta_k.
 \end{split}
\end{equation}
For real gauge fields $U$,
\begin{equation}\label{III-24-13}
 \left|e^{-A(\beta_k^w,U)}\right|
       \le e^{-(1-\mathsf a(w))A(\beta_k,U)}.
\end{equation}
The profile cancellation is the exact identity
\begin{equation}\label{III-24-12}
 -A(\beta_k^w,\mathbf U)
       -\sum_{j=1}^k b_jL^jA(\varphi_j,\mathbf U)
       =-A(\beta-L^Kw,\mathbf U),
\end{equation}
also for complex gauge fields wherever the plaquette actions are read.
\end{proposition}

\begin{proof}
Iterating the recursion from $j$ to $K$ gives
\[
 |x_K-L^{-(K-j)}x_j|
       \le\bar B(1-L^{-(K-j)}).
\]
Hence $x_j-\bar B\ge L^{K-j}(x_K-\bar B)$, which proves the
coupling bound. If two barycentres have physical distance at most
$Ma_{K-j}$, then
$L^{K-j}|w(p)-w(p')|\le M\operatorname{Lip}_\infty(w;S)$.
This proves the two source estimates.

The same recursion gives
$\beta=L^kx_k+\sum_{j=1}^kL^jb_j$. Thus
\[
 |\beta_k-L^kx_k|
 \le B_\sharp\sum_{j=1}^kL^j
 =\bar B(L^k-1),
\]
and \eqref{III-24-9} follows. Since
$|\beta_k^w-\beta_k|=L^K|w|$, the floor proves
\eqref{III-24-10}. Each real plaquette action lies in $[0,2]$;
multiply its value by the real-part bound and sum to obtain
\eqref{III-24-13}. Finally, substituting \eqref{III-24-6} and using
linearity of the weighted action proves \eqref{III-24-12}.
\end{proof}

On a region where all the profiles $\varphi_1,\ldots,\varphi_k$
equal one, the coefficient in \eqref{III-24-6} is simply
$\beta_k^w=L^kx_k^w$. In the real history geometry this includes
the prescribed one-layer interior of $\Lambda_k$ under its collar
condition. No change in that condition is made by adding $w$.

\subsection{Mirror extensions}

The main theorem also applies to sources specified on one cube.
Let $Q=\Delta_i(y_*)$ have side $a_i$, with $i\ge-m$, in the
periodic cube grid. There are $\ell/a_i=2L^{m+i}$ cubes in each
coordinate direction. For $s\in[n-\tfrac12,n+\tfrac12]$, put
$\operatorname{tri}(s)=(-1)^n(s-n)$. This continuous, two-periodic
function defines the fold
\[
 \operatorname{fo}(x)=y_*+a_i
    \left(\operatorname{tri}
           \left(\frac{\widetilde x_\mu-\widetilde y_{*,\mu}}{a_i}\right)
    \right)_{\mu=1}^3\in Q.
\]
The even number of cubes makes this independent of the lift. On
each grid cube it is a reflection followed by a translation onto
$Q$. A cube with centre $y_*+a_in$, $n\in\mathbb Z^3$, has parity
$(-1)^{n_1+n_2+n_3}$.

If $v,v'\in\mathsf{Lp}(Q)$ agree on $\partial Q$, define
$[v\mid v']$ to be $v\circ\operatorname{fo}$ on even cubes and
$v'\circ\operatorname{fo}$ on odd cubes. These prescriptions agree
on cube faces. In particular,
\[
 v^{\rm mp}=[v\mid v],\qquad
 v^{\rm cmp}=[v\mid\bar v]
       \quad\hbox{when }\operatorname{Im}v|_{\partial Q}=0.
\]

\begin{proposition}\label{III-17-18}
For the preceding extensions,
\begin{equation}\label{III-17-19}
 \|[v\mid v']\|_{\flat;T_\ell}
 =\max\{\|v\|_{\flat;Q},\|v'\|_{\flat;Q}\}.
\end{equation}
Thus mirror extension and, when defined, conjugate mirror extension
preserve the source norm. The map $(v,v')\mapsto[v\mid v']$ is
complex linear on pairs with matching boundary values.
If real Lipschitz functions $f_1,\ldots,f_n$ on $Q$ vanish on
$\partial Q$, then
\begin{equation}\label{III-17-22}
 [c+\textstyle\sum_qz_qf_q\mid c+\sum_qz'_qf_q]
 =c+\sum_qz_q[f_q\mid0]+\sum_qz'_q[0\mid f_q].
\end{equation}
For a holomorphic function $\Phi$ on a source class on the torus,
its value at \eqref{III-17-22} is holomorphic in $(c,z,z')$ on
the open set where both cube sources belong to the corresponding
cube class.
\end{proposition}

\begin{proof}
On each cube the fold is an isometry. Consequently the supremum
norm and each local Lipschitz constant are at most the corresponding
maximum for $v,v'$. Join two torus points by a straight segment
between least lifts and cut it at the finitely many grid walls.
The values agree at the cut points, and the lengths of the pieces
sum to the original length. The triangle inequality therefore gives
the same Lipschitz bound globally. Both an even and an odd cube
map onto $Q$, so restriction to those cubes gives the reverse
inequalities for the supremum and Lipschitz seminorms. This proves
\eqref{III-17-19}. Linearity gives \eqref{III-17-22}; the asserted
holomorphy is then the definition on a finite dimensional slice.
\end{proof}

For real $c$, the specialization $z'=\bar z$ gives the conjugate
mirror extension. Holomorphy in the preceding proposition concerns
the independent variables $z,z'$.

\subsection{Bare density and history measures}

Let $\mathcal V_0$ be the finite product of $\mathrm{SU}(2)$ over
positive bare bonds, with product Haar probability measure $dU$.
For a positive bare plaquette $p$, write
\[
 a_p(U)=1-\tfrac12\operatorname{tr}\partial U(p),\qquad
 A(w,U)=\sum_pw(x(p))a_p(U).
\]
Thus $0\le a_p(U)\le2$. For $\mathsf b=L^{-K}\beta$, define
\begin{equation}\label{eq:bare-source-density}
 \rho_0^w(U)=e^{-A(\beta-L^Kw,U)}
            =\rho_0(U)e^{L^KA(w,U)},\qquad
 Z_{K,\ell,\beta}[w]=\int\rho_0^w\,dU,
 \qquad \rho_0=e^{-\beta A(1,U)}.
\end{equation}
These definitions require no smallness assumption on the finitely
many source values.

Fix an admissible real history construction as in
Theorem~\ref{II-203-9}. Its finite set of histories at level $k$
is $\mathcal H_k$, its bounded nonnegative Borel weights are
$\mathfrak w_h$, and its coarse map is $\mathcal M^k$. In completed
coordinates write $\mathsf X_h^{\mathrm{co}\prime}$ for the history
space and
$\xi_h:\mathcal V_0\rightharpoonup\mathsf X_h^{\mathrm{co}\prime}$
for its Borel partial coordinate map. It is independent of $w$,
as are $\mathfrak w_h$ and all tests and stored variables. For a
bounded complex Borel density $\zeta$, put
\begin{equation}\label{eq:source-history-measures}
 \tau_h[\zeta]=(\mathcal M^k)_*(\mathfrak w_h\zeta\,dU),
 \qquad
 \widehat\tau_h[\zeta]
       =(\xi_h)_*(\mathfrak w_h\zeta\,dU).
\end{equation}
The second image is taken on $\operatorname{dom}\xi_h$; this is
the completed history measure denoted
$\hat\tau_h^{\mathrm{co}\prime}$ in \cref{II-184-28}. Set
$\tau_h^w=\tau_h[\rho_0^w]$ and
$\widehat\tau_h^w=\widehat\tau_h[\rho_0^w]$.

The coordinate construction has
\begin{equation}\label{III-24-23}\equationalias{eq:source-history-domain}
 \{U:\mathfrak w_h(U)\ne0\}\subset\operatorname{dom}\xi_h,
 \qquad p_k\xi_h(U)=\mathcal M^kU,
\end{equation}
where $p_k$ projects onto the last coarse field. The domain
inclusion follows from the actual chart construction in
Proposition~\ref{II-169-33} and the paragraph following its proof.
In statements about other
coordinate data below it is an explicit hypothesis.

\begin{proposition}\label{III-16-26}\label{III-24-22}
The measures in \eqref{eq:source-history-measures} are finite complex measures.
For a bounded Borel function $F$ on the completed space,
\begin{equation}\label{III-24-24}
 \int F\,d\widehat\tau_h^w
 =\int_{\operatorname{dom}\xi_h}
       F(\xi_h(U))\mathfrak w_h(U)\rho_0^w(U)\,dU.
\end{equation}
If \eqref{eq:source-history-domain} holds, then
\begin{equation}\label{III-24-25}
 \tau_h^w=(p_k)_*\widehat\tau_h^w,\qquad
 \widehat\tau_h^w(\mathsf X_h^{\mathrm{co}\prime})
       =\tau_h^w(\mathcal V_k)=\int\mathfrak w_h\rho_0^w\,dU.
\end{equation}
For the real construction, with its thresholds satisfying
$\varepsilon_j\le c_2^{\mathrm{SO}}$ for $j<k$,
\begin{equation}\label{III-24-26}
 \sum_{h\in\mathcal H_k}\tau_h^w
      =(\mathcal M^k)_*(\rho_0^w\,dU),\qquad
 \sum_{h\in\mathcal H_k}\tau_h^w(\mathcal V_k)
      =Z_{K,\ell,\beta}[w].
\end{equation}
The expectation in \eqref{III-24-24} is entire in the plaquette
values $w(p)$. For every multi-index $\alpha$,
\begin{equation}\label{eq:source-history-derivatives}
 \partial_w^\alpha\int F\,d\widehat\tau_h^w
 =\int_{\operatorname{dom}\xi_h}F(\xi_h(U))\mathfrak w_h(U)
       \prod_p(L^Ka_p(U))^{\alpha_p}\rho_0^w(U)\,dU.
\end{equation}
It is therefore holomorphic on every source class, including
classes of sources on the torus. The same assertions hold for
coarse history expectations and for $Z[w]$.

If $s\ge|\operatorname{Re}w|_{\mathsf{Bc}}$, then
\begin{equation}\label{III-24-27}
 \begin{split}
  |\widehat\tau_h^w|
    &\le\widehat\tau_h[|\rho_0^w|]
      \le\widehat\tau_h[\rho_0^{1-s/\mathsf b}],\\
  |\tau_h^w|&\le\tau_h[|\rho_0^w|]
      \le\tau_h[\rho_0^{1-s/\mathsf b}].
 \end{split}
\end{equation}
Here $\rho_0^{1-s/\mathsf b}$ is a bounded positive density
even if $1-s/\mathsf b$ is negative. Finally,
\begin{equation}\label{eq:source-history-conjugation}
 \rho_0^{\bar w}=\overline{\rho_0^w},\qquad
 Z[\bar w]=\overline{Z[w]},\qquad
 \widehat\tau_h^{\bar w}=\overline{\widehat\tau_h^w},\qquad
 \tau_h^{\bar w}=\overline{\tau_h^w}.
\end{equation}
\end{proposition}

\begin{proof}
The density $\rho_0^w$ is continuous on a compact finite product.
The weights are bounded and the coordinate maps are Borel.
Consequently the pushforwards are finite complex measures.
Extend $F\circ\xi_h$ by zero outside $\operatorname{dom}\xi_h$;
the extension is bounded and Borel. The pushforward formula gives
\eqref{III-24-24}. The two properties in
\eqref{eq:source-history-domain} then give \eqref{III-24-25}.
The sourced density is gauge invariant, so the finite history
partition identity \eqref{II-5-73} of \cref{II-5-46}
applies to $\rho_0^w$ and proves \eqref{III-24-26}.

For analyticity, put $X_p=L^Ka_p$ and
$\Psi=(F\circ\xi_h)^\natural\mathfrak w_h\rho_0$, where the
superscript denotes the zero extension. Then
\[
 \int F\,d\widehat\tau_h^w
      =\int\Psi\exp\left(\sum_pw(p)X_p\right)dU,
 \qquad |X_p|\le2L^K.
\]
On every compact set of source values the exponential power series
and each of its differentiated series are uniformly absolutely
bounded by a constant times
$\exp(2L^K\sum_p|w(p)|)$. Integrating these series proves entire
dependence and \eqref{eq:source-history-derivatives}. Restriction
to barycentres is complex linear, proving the assertion for torus
sources.

For any Borel partition $(E_i)$ of a Borel set $E$, nonnegativity
of $\mathfrak w_h$ gives
\[
 \sum_i\left|\int_{\xi_h^{-1}(E_i)}
                \mathfrak w_h\zeta\,dU\right|
 \le\int_{\xi_h^{-1}(E)}\mathfrak w_h|\zeta|\,dU.
\]
Taking the supremum over partitions bounds total variation by
the image of $|\zeta|$. The same proof applies to the coarse
map. Also
\[
 |\rho_0^w|=\rho_0e^{L^KA(\operatorname{Re}w,U)}
       \le e^{-(\beta-L^Ks)A(1,U)}
       =\rho_0^{1-s/\mathsf b}.
\]
This proves \eqref{III-24-27}. Conjugating the defining integrals,
whose weights and reference measure are real, proves
\eqref{eq:source-history-conjugation}.
\end{proof}

\subsection{Evaluation of the sourced density}

We use the completed-coordinate convention of
Section~\ref{sec:history-integrability}: a partial density defines a measure
when it is defined and Borel off a Borel reference-measure null
set, and its absolute value is integrable there. Values on the
exceptional set are immaterial. A zero test makes the density
zero without requiring an evaluation of the other factors.

Fix data of the history format at level $k$. They consist of the
same kinds of regular terms, boundary terms and retained factors
as in the real construction, but are not assumed to have been
constructed by a sourced induction. Write $A_{k,h}$ for their
exponent and $\Phi_h$ for the sum of their retained exponents.
The bare-field exponent has the form
\[
 A_{k,h}(\mathsf U,\mathfrak q)
 =-A(\beta_k,\mathsf U)
     +\mathbf E_k(\mathsf U)
     +\mathbf E_k^{\rm bd}(\mathsf U,\mathfrak q)
     -\mathcal E_{k,\mathcal S};
\]
see equation~\eqref{II-99-19}. Define
\begin{equation}\label{III-24-15}
 A_{k,h}^w(\mathsf U,\mathfrak q)
       =A_{k,h}(\mathsf U,\mathfrak q)+L^KA(w,\mathsf U).
\end{equation}
Thus its Wilson term is $-A(\beta_k^w,\mathsf U)$.

At $\phi=((v_n)_{n\le k},\mathfrak q)$, the evaluation set
$\mathsf{Ev}_h(\phi)$ consists of the adapted representatives
of its collared critical orbit, glued to the stored exterior
field. The value $A_{k,h}^w(\phi)$ is defined only if this set is
nonempty, $\phi$ lies in its tiers, and the finite sum in
\eqref{III-24-15} is defined and has the same value at every
representative. The retained sum $\Phi_h(\phi)$ is evaluated
separately by the same convention.

Let $\mathsf{Stay}_h$ be the product of the stored factors,
$\mathbf1_h^{\rm T}$ the tier indicator, and
$\chi_k(\Omega_k)$ the density test. Set
\begin{equation}\label{III-24-16}
 g_{k,h}^w(\phi)=
 \begin{cases}
  0,&\mathsf{Stay}_h\mathbf1_h^{\rm T}\chi_k(\Omega_k)=0,\\
  \mathsf{Stay}_h(\phi)e^{\Phi_h(\phi)+A_{k,h}^w(\phi)},
       &\text{otherwise, when both sums are defined}.
 \end{cases}
\end{equation}
On the remaining points this is a partial map. The sourced
format measure is $\mathsf{Fm}_k^w(h)=g_{k,h}^w\mathscr M_h$
when it satisfies the preceding definition of an integrable
partial density. The assertion that a history measure has this
format means both that this measure is defined and that
\begin{equation}\label{III-24-17}
 \widehat\tau_h^w=\mathsf{Fm}_k^w(h).
\end{equation}
The subsequent sourced induction proves this assertion for its
constructed data.

\begin{lemma}\label{III-24-19}
For fixed format data, $A_{k,h}^w(\phi)$ is defined exactly where
$A_{k,h}(\phi)$ is defined. On this domain,
\begin{equation}\label{III-24-21}
 A_{k,h}^w(\phi)=A_{k,h}(\phi)+L^KA(w,\mathsf U),
       \qquad \mathsf U\in\mathsf{Ev}_h(\phi),
\end{equation}
where the second term is independent of the chosen evaluation
representative and has absolute value at most
$2L^K\sum_p|w(p)|$. The partial densities for $(\mathsf{df},w)$
and $(\mathsf{df},0)$ have the same domain. At level zero,
\begin{equation}\label{III-24-18}
 \widehat\tau_\emptyset^w
       =\rho_0^w\,dU=\mathsf{Fm}_0^w(\emptyset).
\end{equation}
\end{lemma}

\begin{proof}
Two adapted representatives of the same collared orbit differ
by a residual gauge transformation. Extend that transformation
by the identity outside the collar. The fixed endpoints of the
exterior bonds are nodes, so this extension takes one glued
representative to the other; this is the gauge comparison used
in Proposition~\ref{II-168-27}. Every plaquette action, and
hence every weighted plaquette action, is therefore constant
on the evaluation set. Its stated bound follows from
$0\le a_p\le2$.

Adding this everywhere defined, constant-on-representatives
quantity changes neither the domain of a summand nor the
condition that the exponent sum have a single value. The
nonempty-set and tier conditions are unchanged. This proves
\eqref{III-24-21} and the domain assertion for the exponent.
In \eqref{III-24-16}, the test and $\Phi_h$ are unchanged, so
the density domains agree too. At level zero there are no
stored factors or tests, $A_{0,\emptyset}=-\beta A(1,U)$,
$\xi_\emptyset$ is the identity and $\mathscr M_\emptyset=dU$.
Equation \eqref{III-16-25} proves \eqref{III-24-18}.
\end{proof}

This lemma compares two sources for the same format data.
In particular it makes no assertion about pointwise evaluation
of the newly constructed retained or quadratic exponent sums.
Those require their own domain and invariance arguments.

\begin{proposition}\label{III-24-31}
Let $k\ge1$, fix $h\in\mathcal H_k$, and let $w^{\rm off}(p)$
equal $w(p)$ when no corner of $p$ lies in $\Omega_1(h)$ and
zero otherwise. Put $w^{\rm on}=w-w^{\rm off}$. Then
\begin{equation}\label{III-24-33}
 \widehat\tau_h^w
  =e^{L^KA(w^{\rm off},v_0)}\widehat\tau_h^{w^{\rm on}}.
\end{equation}
For fixed format data the same multiplier relates
$\mathsf{Fm}_k^w(h)$ and $\mathsf{Fm}_k^{w^{\rm on}}(h)$:
one is defined exactly when the other is, and their measures
then satisfy \eqref{III-24-33}. Consequently the format identity
\eqref{III-24-17} holds for $w$ exactly when it holds for
$w^{\rm on}$. If $\Omega_1(h)=\emptyset$ and the zero-source
format identity holds, it follows for every source $w$.
\end{proposition}

\begin{proof}
Every bond of a plaquette without a corner in $\Omega_1(h)$
is a proper level-zero bond. The coordinate $v_0$ records
$U$ on these bonds. Thus
\[
 A(w^{\rm off},U)=A(w^{\rm off},v_0)
           \quad\text{at }\xi_h(U).
\]
The right side is a continuous function of a compact group
coordinate. Its exponential and inverse are bounded. Inserting
$\rho_0^w=e^{L^KA(w^{\rm off},U)}\rho_0^{w^{\rm on}}$
into \eqref{III-24-24} proves the measure identity.

The same proper bonds are fixed in every evaluation
representative. Indeed, an inner such bond is part of the
fixed collared datum, and every other such bond is part of
the exterior field used in gluing. Hence
$A(w^{\rm off},\mathsf U)=A(w^{\rm off},v_0)$ throughout
$\mathsf{Ev}_h(\phi)$. The two sourced exponents and therefore
their partial densities differ by the same multiplier. The
multiplier and its inverse are bounded Borel functions, so
almost everywhere definition, Borel measurability and absolute
integrability hold for either density precisely when they
hold for the other. Multiplication of measures proves the
format identity and its converse. If $\Omega_1(h)$ is empty,
$w^{\rm on}=0$, giving the last assertion.
\end{proof}

The truncated plaquette weights $w^{\rm on}$ and $w^{\rm off}$
are used here only in finite measure identities; no preservation
of a uniform Lipschitz source class is needed for this reduction.

\section{First passage and the vacuum normalization}
\label{sec:first-passage}

The finite renormalisation theorem applies to prescribed sequences of
couplings. The coefficient recursion selects one such sequence from a
bare inverse coupling. We first stop that sequence at a prescribed
reference coupling and choose torus sizes that work for every resulting
cutoff.

Fix the constants preceding the threshold in Theorem~\ref{II-203-9} and
Proposition~\ref{II-202-29}. In particular, $L\ge13$ is odd, $r\ge2$,
$M=L^{m'}$, and $B_\sharp>0$. Put
\[
 \bar B=\frac{L}{L-1}B_\sharp,
 \qquad
 R(x)=\min\{L^a:a\in\mathbb Z_{\ge0},\ L^a\ge(\log x)^r\}
 \quad(x>1).
\]
We use the fixed geometric constant $C_{\mathrm s}\ge2$ in the torus
floor of Section~\ref{sec:common-iteration-constants}, after \eqref{II-202-42}. The choices are made
before the couplings and the lattice. Decrease $\gamma$, if necessary,
so that
\begin{equation}\label{eq:first-passage-constants}
 0<\gamma\le\tfrac12,\qquad
 \gamma^{-2}\ge2B_\sharp,\qquad
 \gamma^{-2}>\bar B,\qquad
 L^{-1}\gamma^{-2}-B_\sharp>1,\qquad
 (L^{1/r}-1)\log\gamma^{-2}\ge\log(2L).
\end{equation}
These requirements are compatible with the preceding choices by
Proposition~\ref{II-202-29}. At this threshold fix the canonical
auxiliary data. Every construction in this section uses that one
assignment, including when the reference coupling or stopping index
changes.

An admissible finite sequence is a real sequence $u_0,\ldots,u_K$ with
\begin{equation}\label{eq:admissible-prescribed-sequence}
 u_j\ge\gamma^{-2}\quad(0\le j\le K),\qquad
 u_j=L^{-1}u_{j-1}-b_j,\quad |b_j|\le B_\sharp
                                      \quad(1\le j\le K).
\end{equation}
For $K,m\ge1$ and $\ell=2L^m$, it gives an admissible lattice when
$MR(u_K)\le L^m$. These are the sequence and torus hypotheses of
Theorem~\ref{II-203-9}; prescribed increments need not be the
coefficients of the constructed regular terms.

For $\beta\ge\gamma^{-2}$, let $x_j=x_j(\beta)$ instead denote the
sequence defined by
\begin{equation}\label{III-5-15}
 x_0=\beta,\qquad
 b_{j+1}=\mathsf{cf}\bigl(\widehat E^{(j+1)}[x_0,\ldots,x_j]\bigr),
 \qquad x_{j+1}=L^{-1}x_j-b_{j+1}.
\end{equation}
Here $\widehat E^{(j+1)}$ is the regular family of
Theorem~\ref{II-144-26}, and $\mathsf{cf}$ is its second-moment
coefficient. The family and its coefficient are defined from the
displayed prefix before $x_{j+1}$ is chosen. By
Proposition~\ref{II-132-35}, this recursion gives a unique finite,
strictly decreasing real sequence $x_0,\ldots,x_J$, with
$|b_j|\le B_\sharp$,
\[
 1<x_J<\gamma^{-2},\qquad
 x_j\ge\gamma^{-2}\quad(0\le j<J).
\]
The sequence and its coefficients are independent of the torus. We call
its prefixes the canonical coupling sequences. They agree at common
indices when the reference coupling or stopping index changes at this
fixed threshold assignment.

For $0<g\le\gamma$, set
\begin{equation}\label{III-1-5}
 X_g=L(g^{-2}+B_\sharp),\qquad
 X_g+\bar B=L(g^{-2}+\bar B).
\end{equation}
A pair $(K,\beta)$ is a first-passage datum at $g$ if the least
index $j\ge0$ at which $x_j(\beta)\le X_g$ is $K\ge1$. Its label is
$\mathsf b=L^{-K}\beta$. Define
\begin{equation}\label{III-1-8}
 I(g)=(g^{-2}+\bar B,\ X_g-\bar B],
\end{equation}
and choose sides from
\begin{equation}\label{III-1-15}
 \begin{split}
 m_{\mathfrak L}(g)
   &=1+m'+\max\left\{
       \log_L R\bigl(L(g^{-2}+\bar B)\bigr),
       \lceil\log_L C_{\mathrm s}\rceil\right\},\\
 \mathfrak L(g)
   &=\{2L^m:m\in\mathbb Z_{\ge1},\ m\ge m_{\mathfrak L}(g)\}.
 \end{split}
\end{equation}
The floor is an integer, so $\mathfrak L(g)$ is nonempty. It depends
only on $g$ and the fixed construction constants.

\begin{proposition}[First passage]\label{III-5-11}
Let $0<g\le\gamma$ and $\beta>X_g$. There is a unique first-passage
datum $(K,\beta)$ at $g$, with $1\le K<J$. Its canonical sequence is
admissible and satisfies
\begin{equation}\label{III-5-12}
 g^{-2}<x_K\le X_g<x_{K-1}<\cdots<x_0=\beta.
\end{equation}
For $0\le i\le j\le K$,
\begin{equation}\label{III-5-13}
 |x_j-L^{-(j-i)}x_i|
       \le\bar B\bigl(1-L^{-(j-i)}\bigr).
\end{equation}
The label and the terminal coupling obey
\begin{equation}\label{III-5-14}
 g^{-2}+B_\sharp\frac{L-2}{L-1}
    <\mathsf b
    \le L(g^{-2}+\bar B)-\bar B L^{-K},
\end{equation}
\begin{equation}\label{III-5-17}
 \log g^{-2}<\log x_K
       \le\log g^{-2}+\log\frac{3L}{2}.
\end{equation}
Every $\ell=2L^m\in\mathfrak L(g)$ gives an admissible lattice for
this sequence. More precisely,
\begin{equation}\label{eq:first-passage-torus-room}
 LM R(x_j)\le L^{m+K-j}\qquad(0\le j\le K).
\end{equation}
Thus the finite renormalisation theorem and its larger local-chart
condition apply at every level. One set $\mathfrak L(g)$ serves all
first-passage data at $g$.
\end{proposition}

\begin{proof}
The terminal value $x_J<\gamma^{-2}\le g^{-2}<X_g$ ensures that a
first crossing exists. Since $x_0=\beta>X_g$, its index $K$ is positive.
At that crossing,
\[
 x_K=L^{-1}x_{K-1}-b_K
       >L^{-1}X_g-B_\sharp=g^{-2}.
\]
Consequently $K<J$. Strict decrease gives \eqref{III-5-12}, and the
prefix satisfies \eqref{eq:admissible-prescribed-sequence}. The affine
estimate \eqref{III-5-13} is \eqref{II-132-36}. Its applications at
$(0,K)$ and $(0,K-1)$ give
\[
 \begin{split}
 \mathsf b&\le x_K+\bar B(1-L^{-K})
          \le X_g+\bar B-\bar B L^{-K},\\
 L\mathsf b&\ge x_{K-1}-\bar B(1-L^{-(K-1)})
          >X_g-\bar B.
 \end{split}
\]
Since $L^{-1}(X_g-\bar B)=g^{-2}+B_\sharp(L-2)/(L-1)$, these are
\eqref{III-5-14}. The inequality $B_\sharp\le g^{-2}/2$ gives
$x_K\le X_g\le(3L/2)g^{-2}$ and proves \eqref{III-5-17}.

For the side estimate we use the following consequence of the last
condition in \eqref{eq:first-passage-constants}. If $X\ge\gamma^{-2}$,
$s\in\mathbb Z_{\ge0}$ and $1<X'\le L^sX$, then
\begin{equation}\label{eq:radius-scaling}
 R(X')\le L^sR(X).
\end{equation}
Indeed, with $Y=\log X$ one has
$\log L\le(L^{1/r}-1)Y$, whence induction gives
$Y+s\log L\le L^{s/r}Y$. Thus
$(\log X')^r\le L^s(\log X)^r$, and $L^sR(X)$ is an allowed power
of $L$ in the definition of $R(X')$.

Put $X=X_g+\bar B=L(g^{-2}+\bar B)$. Rearranging
\eqref{III-5-13} at $(j,K)$ yields
\[
 x_j+\bar B\le L^{K-j}(x_K+\bar B)\le L^{K-j}X.
\]
Apply \eqref{eq:radius-scaling} to obtain
$R(x_j)\le L^{K-j}R(X)$. The floor in
\eqref{III-1-14} gives $LMR(X)\le L^m$, which proves
\eqref{eq:first-passage-torus-room}. In particular, $MR(x_K)\le L^m$ is the room
condition of Theorem~\ref{II-203-9}, and the stronger inequality
$LMR(x_K)\le L^m$ is its local-chart condition
\eqref{II-202-42}.
\end{proof}

\begin{corollary}[A common interval of labels]\label{III-5-19}
For fixed $0<g\le\gamma$, the sets
\[
 \mathcal W_K(g)=\{\beta:(K,\beta)\text{ is a first-passage datum at }g\},
 \qquad K\ge1,
\]
are pairwise disjoint and have union $(X_g,\infty)$. Moreover,
\begin{equation}\label{III-5-20}
 L^K I(g)\subset\mathcal W_K(g)
 \subset L^K\left(g^{-2}+B_\sharp\frac{L-2}{L-1},
                         L(g^{-2}+\bar B)\right].
\end{equation}
The interval $I(g)$ has positive length. Thus every
$\mathsf b\in I(g)$ gives a first-passage datum $(K,L^K\mathsf b)$
for every $K\ge1$.
\end{corollary}

\begin{proof}
The first assertion follows from the existence and uniqueness of the
first crossing. The outer inclusion is \eqref{III-5-14}.
The length of $I(g)$ is
\[
 (L-1)g^{-2}+B_\sharp\frac{L(L-3)}{L-1}>0.
\]
Take $\mathsf b\in I(g)$ and put $\beta=L^K\mathsf b$. Then
$\beta\ge L\mathsf b>L(g^{-2}+\bar B)=X_g+\bar B$, so the first
crossing $K'$ exists. Its label
$\mathsf b'=L^{-K'}\beta=L^{K-K'}\mathsf b$ satisfies
\[
 L\mathsf b'>X_g-\bar B,\qquad
 \mathsf b'\le X_g+\bar B
\]
by \eqref{III-5-14}. If $K'\ge K+1$, then
$L\mathsf b'\le\mathsf b\le X_g-\bar B$, a contradiction.
If $K'\le K-1$, then
$\mathsf b'\ge L\mathsf b>X_g+\bar B$, also a contradiction.
Hence $K'=K$, proving the inner inclusion.
\end{proof}

\begin{corollary}[Stopped data]\label{III-5-21}
Let $(K,\beta)$ be a first-passage datum at $0<g\le\gamma$ and let
$\ell=2L^m\in\mathfrak L(g)$. Put $g_{(0)}=g$, and, for
$1\le i\le K-1$, define
\begin{equation}\label{III-1-7}
 g_{(i)}^{-2}=L^{-1}x_{K-i}(\beta)-B_\sharp.
\end{equation}
For $0\le i\le K-1$, set $k=K-i$. Then $(k,\beta)$ is a
first-passage datum at $g_{(i)}$, with canonical sequence
$(x_0,\ldots,x_k)$ and label $L^i\mathsf b$. Moreover,
\[
 L^i\ell\in\mathfrak L(g_{(i)}),
\]
and the lattice with cutoff $k$ and side $L^i\ell$ is admissible
and satisfies the larger local-chart condition. For $i\ge1$,
$g_{(i)}<g$ and $X_{g_{(i)}}=x_k$.
\end{corollary}

\begin{proof}
There is nothing to prove when $i=0$. For $i\ge1$,
$x_k>X_g$ gives $g_{(i)}^{-2}>g^{-2}>0$. The definition implies
$X_{g_{(i)}}=x_k<x_{k-1}$, so the first crossing at $g_{(i)}$ is
exactly $k$. Prefix consistency of \eqref{III-5-15} identifies the
canonical sequence, and $L^{-k}\beta=L^i\mathsf b$ gives its label.

The affine estimate gives
\[
 X_{g_{(i)}}+\bar B=x_k+\bar B
          \le L^i(x_K+\bar B)\le L^i(X_g+\bar B).
\]
Consequently \eqref{eq:radius-scaling} yields
$m_{\mathfrak L}(g_{(i)})\le i+m_{\mathfrak L}(g)$, including
the unchanged term $\lceil\log_L C_{\mathrm s}\rceil$ in the
maximum. Hence $L^i\ell\in\mathfrak L(g_{(i)})$. Apply
\cref{III-5-11} to this datum. Its finest side exponent is
$(m+i)+k=m+K$, as required for the stopped lattice.
\end{proof}

\subsection{The zero-source partition function}

For an admissible lattice $\bar{\mathfrak w}=((\ell,K),u_\bullet)$,
write $E_{(k)}[\bar{\mathfrak w}]$ for the accumulated vacuum
constant of its first $k$ steps, as defined in
\eqref{II-97-8}. These are real numbers attached to the constructed
regular terms. Let
\[
 Z_{K,\ell,\beta}=\int e^{-\beta A_K(U)}\,dU,
 \qquad \beta=u_0,
\]
with the Wilson action and normalized product Haar measure used in
Theorem~\ref{II-203-9}. The effective measures at all levels have
this same positive total mass.

For a canonical sequence, define
\begin{equation}\label{III-7-7}\equationalias{III-7-8}
 E_{K,\ell,\beta}=E_{(K)}[\bar{\mathfrak w}],
 \qquad
 \widetilde Z(0)=e^{-E_{K,\ell,\beta}}Z_{K,\ell,\beta}.
\end{equation}
Once the construction constants and auxiliary choices are fixed,
$E_{K,\ell,\beta}$ depends on $K,\ell,\beta$; a reference coupling
does not enter its definition.

\begin{proposition}[Vacuum stability]\label{III-7-9}
Let $u_0,\ldots,u_K$ satisfy
\eqref{eq:admissible-prescribed-sequence}, and assume
$MR(u_K)\le L^m$, where $K,m\ge1$ and $\ell=2L^m$.
For every $1\le k\le K$,
\begin{equation}\label{III-7-11}
 e^{-C_{\rm low}(u_k)(L^{K-k}\ell)^3}
 \le e^{-E_{(k)}[\bar{\mathfrak w}]}Z_{K,\ell,u_0}
 \le e^{(C_{\rm maj}+2)(L^{K-k}\ell)^3},
 \qquad C_{\rm low}(x)=c_{\rm low}+\tfrac92\log x.
\end{equation}
The constants $c_{\rm low},C_{\rm maj}$ are fixed before
$\gamma$, the couplings, the cutoff and the torus.

In particular, there is a constant $C_{\rm fp}$ with the same
independence such that every first-passage datum $(K,\beta)$ at
$0<g\le\gamma$ and every $\ell\in\mathfrak L(g)$ satisfy
\begin{equation}\label{eq:first-passage-vacuum-stability}
 e^{-C_{\rm fp}(1+\log g^{-2})\ell^3}
       \le\widetilde Z(0)
       \le e^{(C_{\rm maj}+2)\ell^3}.
\end{equation}
\end{proposition}

\begin{proof}
The sequence and room hypotheses are precisely those of
Theorem~\ref{II-203-9}. Its level-$k$ vacuum constant is
$E_{(k)}[\bar{\mathfrak w}]$, and the level-$k$ torus has
$(L^{K-k}\ell)^3$ sites. Thus \eqref{III-7-11} is the partition
bound of Proposition~\ref{II-203-26}, \eqref{eq:actual-partition-bounds}.

For canonical first-passage data, \cref{III-5-11} verifies the
hypotheses. By \eqref{III-5-17},
\[
 C_{\rm low}(x_K)
 \le c_{\rm low}+\tfrac92\log\tfrac{3L}{2}
                       +\tfrac92\log g^{-2}.
\]
Take
\begin{equation}\label{eq:first-passage-stability-constant}
 C_{\rm fp}=\max\left\{\tfrac92,
                  c_{\rm low}+\tfrac92\log\tfrac{3L}{2}\right\}.
\end{equation}
Since $\log g^{-2}>0$, the right side is at most
$C_{\rm fp}(1+\log g^{-2})$. The case $k=K$ of
\eqref{III-7-11} proves \eqref{eq:first-passage-vacuum-stability}.
\end{proof}

\subsection{Comparison of vacuum constants on two tori}

The vacuum constant depends on the torus through the local terms.
Terms whose domains lift into both tori have the same value. The
remaining domains have diameter comparable to the smaller torus,
so their exponential decay controls the difference per unit volume.

Keep one threshold assignment, one admissible sequence
$u_0,\ldots,u_K$, and its real regular progression. Let $\kappa$
be the decay rate used for the vacuum estimates in
Section~\ref{sec:vacuum-bounds}, and require
\[
 \kappa\ge2\kappa_{\mathcal D}+5,\qquad
 \kappa_{\mathcal D}=\log9216.
\]
The constants can be chosen with this rate. Let $C_L\ge8L$ be
the room constant for the local lattice charts in that construction.

\begin{proposition}\label{III-7-12}
Suppose that $\bar{\mathfrak w}_\iota=((\ell_\iota,K),u_\bullet)$,
$\ell_\iota=2L^{m_\iota}$, $\iota=1,2$, are admissible lattices
for the same sequence, and that
\[
 LM R(u_K)\le L^{m_\iota},\qquad \ell_\iota\ge C_LM.
\]
Put $\ell_*=\min(\ell_1,\ell_2)$ and
\[
 C_{\mathrm E,0}=2048\sum_{n\ge0}e^{-n},\qquad
 C_{\rm tt}=2C_{\mathrm E,0},\qquad C'_{\rm tt}=2.
\]
Then
\begin{equation}\label{III-7-13}
 \left|\frac{E_{(K)}[\bar{\mathfrak w}_1]}{\ell_1^3}
       -\frac{E_{(K)}[\bar{\mathfrak w}_2]}{\ell_2^3}\right|
 \le2C_{\rm tt}C'_{\rm tt}E_0
              e^{-\kappa\ell_*/(5M)}.
\end{equation}
In particular the right side is at most $2C_{\rm tt}C'_{\rm tt}E_0$.
The proposition permits $\ell_1=\ell_2$.
\end{proposition}

\begin{proof}
At level $j$, the cube side in physical coordinates is
$W_j=ML^{j-K}$. Call an infinite-grid domain a near lift for a
torus of side $\ell$ if it is contained in a box with all side
lengths strictly less than $\ell/2$. A torus domain is common
if it is the projection of a lift near for both tori.
The other domains will be called noncommon.

We first record the geometry used below. Projection of a near
lift identifies its cubes, their incidences and their edges, and
preserves $d_j$. Two such lifts of the same domain differ by a
period translation. For a specified lift of a site $z$ in the
domain there is exactly one near lift containing that site.
These assertions follow by starting with its cube and continuing
across shared faces. No two points separated by a nonzero period
can occur in the box, even after comparing its two possible lifts,
since each of its side lengths is less than half a period.

A connected union of $d$ cube edges lies in a box of width at
most $dW_j$ in each direction. Every cube it meets lies in the
box enlarged by $W_j$ at either end. Hence a lifted domain has
width at most $(d_j+2)W_j$. The same assertion gives a near
lift of a torus domain whenever $(d_j+2)W_j<\ell/2$:
lift its edges in a connected order, choosing each new edge
to share the lifted endpoint of an earlier edge, and lift the
cubes meeting them. The lifted union is contained in the
enlarged box and projects onto the domain. This argument also
applies when the edge union is a single corner.
Consequently every noncommon domain on either torus satisfies
\begin{equation}\label{eq:noncommon-domain-size}
 d_j(X)\ge\frac{\ell_*}{2W_j}-2.
\end{equation}
Indeed, a domain with a smaller size has a lift near for its
own torus and of width less than $\ell_*/2$, hence near for both.

The two vacuum families at level $j$ are
\[
 \mathsf{vl}^{(j)}_\iota(X,z)
       =E^{(j),\mathfrak w_\iota}(X,\mathbf1,z),\qquad
 \mathsf{vn}^{(j)}_\iota(X,z)=\mathsf n^{(j),\mathfrak w_\iota}(X,z).
\]
They are real and each is bounded by $E_0e^{-\kappa d_j(X)}$;
see Section~\ref{subsec:ymm1-62-1} and Proposition~\ref{II-192-21}.
They have the same values on a common lift. For the regular
family, dilate that lift by $\lambda_j=L^{K-j+1}$ to read it in
the level-$(j-1)$ coordinates of the infinite companion in
Theorem~\ref{II-144-26}. The trivial triple $(\mathbf1,0,0)$
belongs to every analytic core and projects to the trivial triple
on both tori. The companion identity therefore gives one value
for $\mathsf{vl}^{(j)}$ on the common lift. For
$\mathsf{vn}^{(j)}$, use its local logarithm series and pivot
term, \eqref{II-45-8} and the local-chart identity following \eqref{II-45-65}.
Their definitions are transported by the same chart. All cubes,
nonpivot bonds and tents in these definitions are carried by
the common lift, so the Gaussian vacuum values agree as well.

Write the full-torus contribution of step $j-1$ as
$\mathsf N^{(j-1)}_\iota+\mathsf S^{(j)}_\iota$, where
$\mathsf N$ is its counting term and
\[
 \mathsf S^{(j)}_\iota
 =\sum_{z\in\mathbb T^{(j)}_\iota}\sum_{X\ni z}
       \bigl(\mathsf{vl}^{(j)}_\iota(X,z)
                     +\mathsf{vn}^{(j)}_\iota(X,z)\bigr).
\]
There are $3(L^3-1)$ positive nonpivot bonds per level-$j$ site.
Thus
\[
 \ell_\iota^{-3}\mathsf N^{(j-1)}_\iota
 =L^{3(K-j)}\bigl[
   3(L^3-1)(\mathsf{co}^{\rm b}_{j-1}
                    +\mathsf{co}^{\rm i}_{j-1})
                    +\mathsf{co}^{\rm s}_{j-1}\bigr].
\]
The three coefficients in \eqref{II-192-10} depend on the
common coupling sequence and fixed constants alone. Their
contributions per volume therefore cancel.

The common-domain contributions to $\mathsf S^{(j)}$ also
cancel per volume. To see this without assuming an extra
translation law, lift a site to the infinite grid and sum over
the near lifts containing it. The preceding uniqueness and
common-value statements give one function of the lifted site
which factors through both torus coverings. Let
$\ell^*=\max(\ell_1,\ell_2)$; it is an integer multiple
$q_\iota\ell_\iota$ of each side. In a half-open cube of side
$\ell^*$ each site of the $\iota$th torus has exactly
$q_\iota^3$ lifts. Summing the common function on that cube
therefore proves equality of the two common contributions
divided by their volumes. If $\mathsf S^{(j),\rm nc}_\iota$
denotes the remaining sum, we have
\begin{equation}\label{eq:vacuum-noncommon-reduction}
 \frac{E_{(K)}[\bar{\mathfrak w}_1]}{\ell_1^3}
 -\frac{E_{(K)}[\bar{\mathfrak w}_2]}{\ell_2^3}
 =\sum_{j=1}^K
  \left(\frac{\mathsf S^{(j),\rm nc}_1}{\ell_1^3}
             -\frac{\mathsf S^{(j),\rm nc}_2}{\ell_2^3}\right).
\end{equation}

It remains to bound these sums. Every level-$j$ site lies in
the interior of one cube: $M$ is odd and the cube walls have
half-integer coordinates in level-$j$ units. The entropy bound
of Proposition~\ref{II-18-26} therefore gives, for $D\ge0$,
\[
 \sum_{\substack{X\ni z\,;\ d_j(X)\ge D}}e^{-\kappa d_j(X)}
 \le2048\sum_{n\ge D}e^{-(\kappa-\kappa_{\mathcal D})n}
 \le C_{\mathrm E,0}e^{-\kappa D/2}.
\]
The last inequality uses $\kappa/2-\kappa_{\mathcal D}>1$.
Put $t_j=\ell_*L^{K-j}/M$. The side condition and the integral
power form of the sides imply $m_\iota\ge m'+2$, whence
$t_j\ge2L^2\ge338$. In particular,
\[
 \frac12\left(\frac{t_j}{2}-2\right)\ge\frac{t_j}{5}.
\]
Use \eqref{eq:noncommon-domain-size} in the entropy tail,
multiply by the two-family bound and by the number of sites
$\ell_\iota^3L^{3(K-j)}$, and obtain
\[
 \ell_\iota^{-3}|\mathsf S^{(j),\rm nc}_\iota|
 \le C_{\rm tt}E_0L^{3(K-j)}e^{-\kappa t_j/5}.
\]
Finally put $\alpha=\kappa\ell_*/(5M)>4$. The ratio of successive
terms in $\sum_{s\ge0}L^{3s}e^{-\alpha L^s}$ is at most
\[
 \exp\bigl(3\log L-\alpha L^s(L-1)\bigr)\le\tfrac12,
\]
so the sum is at most $2e^{-\alpha}$. Summing the preceding
estimate over $j$ and the two tori in
\eqref{eq:vacuum-noncommon-reduction} proves \eqref{III-7-13}.
\end{proof}

For first-passage data, the stronger room follows from
\cref{III-5-11}. Choosing $C_{\mathrm s}\ge C_L$ in the
torus floor also gives $\ell\ge C_LM$. Hence the proposition
applies to every pair of sides in $\mathfrak L(g)$ for the
same first-passage datum and the same threshold assignment.


\section{Expansion about the conditional mean}
\label{sec:centered-expansion}

The source produces a potential with quadratic growth in the
fluctuation. Its natural argument is the fluctuation minus the
conditional mean. This distinction matters in the Gaussian estimate:
bounding that potential in terms of the unshifted field would
introduce the cutoffs on all integration bonds into the volume bound.
We first prove the estimate with growth about the mean, then use it
for the local expansion.

\subsection{The Gaussian estimate with quadratic growth}

Use the local Gaussian structure of Definition~\ref{II-114-9}
at a step $t=k+1$, where $0\le k<K$,
$M=L^{m'}$ and $m'<m+K-t$. The ambient region $\mathfrak R$
is nonempty and is a union of coarse cubes. The geometric
constants satisfy $R_1\ge L$, $M_1\ge R_1\mathfrak M$
and $M\ge M_1$, and $\bar\kappa\ge1$. Thus
$\Lambda_t\subset\Omega_t$, the
integration bonds $\mathfrak X_{\rm I}$ are the nonpivot bonds
touching $\Lambda_t$, and the other nonpivot bonds are external.
For a fine region $Z_\circ$ with $Z_\circ^\sharp\subset\mathfrak R$,
write $\mathfrak f=\mathfrak X_{\rm I}\langle Z_\circ\rangle$
and $\mathfrak r=\mathfrak X_{\rm I}\langle\mathfrak R\rangle$.
Brackets mean that a bond touches the region; a superscript
$\circ$ requires both endpoints in it. Positive bonds are the
independent coordinates. The norm $|\cdot|$ on
$\mathfrak{su}(2)^{\mathbb C}$ is the matrix operator norm, whereas
\[
 \|z\|^2=\sum_{b\in\mathfrak f^+}
       \bigl(|\operatorname{Re}z(b)|^2+|\operatorname{Im}z(b)|^2\bigr)
\]
is the Euclidean coordinate norm.

The covariance, connecting operator and mean are denoted
$(C,\Gamma,m)(Z_\circ,s)$. They are holomorphic near the closed
decoupling polydisc $|s(\Delta)|\le e^{\bar\kappa}$ and have
the restriction and component properties in the cited definition.
Their real references satisfy
\[
 S_{\rm r}=C_{\rm r}^{-1}\ge\gamma_{\rm r}>0,\qquad
 \|S_{\rm r}\|,\|\Gamma_{\rm r}\|,
                 \|C_{\rm r}\Gamma_{\rm r}\|\le c_{\rm r},
 \qquad 0<\gamma_{\rm r}\le1\le c_{\rm r},
\]
and each of $C^{-1}-S_{\rm r}$, $\Gamma-\Gamma_{\rm r}$ and
$\Gamma^{\mathsf T}C\Gamma-\Gamma_{\rm r}^{\mathsf T}
C_{\rm r}\Gamma_{\rm r}$ has norm at most $\hat\varepsilon$.
The reference connecting operator uses auxiliary bonds touching
$Z_\circ^\sharp$ only. The reversal-invariant thresholds satisfy
$1\le\mathfrak b_-\le\mathfrak b_b\le\mathfrak b_+<\infty$, and
\[
 |m_{\rm r}(b)|\le\theta_m\mathfrak b_b,
 \qquad |m(b)-m_{\rm r}(b)|\le\varepsilon_m\mathfrak b_b,
 \qquad\theta_m+\varepsilon_m\le\vartheta_*<1.
\]
In this section the number $\rho\ge0$ is a quadratic-growth
coefficient, distinct from the source radius $\varrho$. Put
\begin{equation}\label{eq:centered-gaussian-constants}
\begin{gathered}
 d_*=1-\vartheta_*,\qquad C_{\rm mn}\ge18d_*^{-2},\\
 \lambda_{\rm r}=\min\{\gamma_{\rm r}/2,(1+2c_{\rm r}^2)^{-1}\},
 \qquad c_{\rm mn}=C_{\rm mn}\lambda_{\rm r}^{-1},\\
 \Sigma_\rho=\hat\varepsilon+(1+c_{\rm r})\varepsilon_m^2
                       +\gamma_2+\rho,\qquad
 u_\rho=C_{\rm mn}\Sigma_\rho,\qquad \gamma_2\ge0.
\end{gathered}
\end{equation}
We require
\begin{equation}\label{eq:centered-gaussian-smallness}
 C_{\rm mn}u_\rho\le\lambda_{\rm r}.
\end{equation}

For disjoint reversal-invariant sets $\mathfrak S,\mathfrak e
\subset\mathfrak f$, the first carrying small-field tests and the
second failed tests, use the tube of Section~\ref{subsec:ymm1-56-2}.
On $\mathfrak e$ the field is real with $|B(b)|\ge\mathfrak b_b$;
off $\mathfrak e$ its imaginary part satisfies
$|\operatorname{Im}B(b)|<\varepsilon'\mathfrak b_b$, where
$\varepsilon'>\varepsilon_m$; on $\mathfrak S$ also
$|\operatorname{Re}B(b)|+|\operatorname{Im}B(b)|<\mathfrak b_b$.
A function is admitted on this tube if it is Borel and, for fixed
real failed variables, holomorphic in the other variables. Let $D$
be its real part. Tested integrands are extended by zero outside $D$.
The conditioning integral is the normalized joint integral
\begin{equation}\label{eq:centered-conditioning-integral}
 \mathbb J^m(Z_\circ;s)[F]
 =\int d\mu_0(\mathbb X)\det{}^{1/2}(C^{-1}/2\pi)
   \int dB\,F(B)
 e^{-\frac12(B-m+C\Gamma\mathbb X,C^{-1}(B-m+C\Gamma\mathbb X))_1}.
\end{equation}
The determinant branch is the one on matrices with positive real
part. Empty integration sets have the value $F(0)$.

\begin{proposition}\label{III-29-9}
Assume \eqref{eq:centered-gaussian-smallness}. Suppose $s$ vanishes
outside a coarse region $Z$ with
$Z_\circ^\sharp\subset Z\subset\mathfrak R$. If $G$ is admitted
on the preceding tube and, at this value of $s$,
\begin{equation}\label{III-29-10}
 |G(B)|\le F_0 e^{\rho\|B-m(Z_\circ,s)\|^2},\qquad F_0\ge0,
\end{equation}
then the conditioning integral converges absolutely and
\begin{equation}\label{III-29-11}
 |\mathbb J^m(Z_\circ;s)[\mathbf1_DG]|
 \le F_0\exp\left[-C_{\rm mn}^{-1}\gamma_2
                     \sum_{b\in\mathfrak e^+}\mathfrak b_b^2
                      +c_{\rm mn}u_\rho N_{\rm I}(Z)\right],
 \qquad N_{\rm I}(Z)=3\#\mathfrak X_{\rm I}^+\langle Z\rangle.
\end{equation}
There is also an integrable majorant for the original joint
integrands which is uniform in $s$ and in all $G$ satisfying
\eqref{III-29-10} with the same $F_0,\rho$. The majorant may
depend on the fixed finite thresholds.
\end{proposition}

\begin{proof}
We use the contour of Proposition~\ref{II-117-69}, keeping the
additional growth factor in its estimate. Write $S=C^{-1}$ and
\[
 Q=\begin{pmatrix}1+\Gamma^{\mathsf T}C\Gamma&\Gamma^{\mathsf T}\\
                   \Gamma&S\end{pmatrix}.
\]
Completion of the square and the reference bounds give
$Q_{\rm r}\ge\lambda_{\rm r}$,
$\|Q-Q_{\rm r}\|\le2\hat\varepsilon$ and $\det Q=\det S$,
as in \eqref{II-117-73}. The smallness assumption implies
\[
 \Sigma_\rho\le d_*^4\lambda_{\rm r}/324,\qquad
 \hat\varepsilon+3\rho\le\lambda_{\rm r}/108,
 \qquad16d_*^{-2}\varepsilon_m^2\le\tfrac12.
\]
In particular $\operatorname{Re}S>0$.

Let $\mathfrak b_\Sigma^2=\sum_{b\in\mathfrak f^+}\mathfrak b_b^2$.
The mean bound gives $\|m\|\le\vartheta_*\mathfrak b_\Sigma$.
For real $B$, \eqref{III-29-10} is at most
$F_0e^{2\rho\mathfrak b_\Sigma^2}e^{2\rho\|B\|^2}$.
The real-part calculation leading to \eqref{II-117-47}, now
including $2\rho\|B\|^2$, gives the common majorant
\begin{equation}\label{III-29-12}
 F_0e^{(c_{\rm r}+2\rho)\mathfrak b_\Sigma^2}
 \left(\frac{c_{\rm r}+1}{2\pi}\right)^{3\#\mathfrak f^+/2}
 e^{-3\lambda_{\rm r}\|\mathbb X\|^2/8
    -\lambda_{\rm r}\|B\|^2/8}.
\end{equation}
The inequality $\hat\varepsilon+3\rho\le\lambda_{\rm r}/108$
leaves the displayed Gaussian decay. This proves absolute
convergence before any change of contour.

The auxiliary variables outside $Z$ integrate to one, by the
component property of $\Gamma$. After their removal the joint
form, denoted $Q^Z$, has the same reference and perturbation
bounds, and its dimension $n$ is at most $2N_{\rm I}(Z)$.
Put $\nu=\operatorname{Im}m$. The contour is
$\Theta(\beta)(b)=\beta(b)+i\phi_b(\beta(b))\nu(b)$.
On failed bonds $\phi_b=0$, on untested bonds $\phi_b=1$, and
on small-field bonds
\[
 \phi_b(\beta)=\min\{1,\max\{0,2d_*^{-1}
            (1-d_*/4-|\beta|/\mathfrak b_b)\}\}.
\]
It agrees with the real contour near the boundary of each
small-field ball and stays in the tube. Its Jacobian factors
satisfy $|J_b|\le e^{2d_*^{-2}\varepsilon_m^2}$.
The contour identity in the cited proof applies to the present
integrand: along every intermediate contour,
$\|B-m\|^2\le2\|B\|^2+2\|m\|^2$, and the same Gaussian
calculation used in \eqref{III-29-12} dominates the growth.
Thus the vertical sides on untested variables vanish at infinity,
and Fubini remains valid. Failed variables stay real.

Set $\beta'=\beta-\operatorname{Re}m$ and
$\eta=(1-\phi)\nu$. The cutoff geometry gives
\[
 |\beta'(b)|\ge d_*\mathfrak b_b\quad(b\in\mathfrak e^+),
 \qquad\|\eta\|^2\le16d_*^{-2}\varepsilon_m^2\|\beta'\|^2.
\]
Since $\Theta(\beta)-m=\beta'-i\eta$, the growth factor on
the moved contour is at most
\begin{equation}\label{III-29-19}
 F_0e^{\rho(\|\beta'\|^2+\|\eta\|^2)}
                         \le F_0e^{3\rho\|\beta'\|^2/2}.
\end{equation}
This is where centering removes the dependence on the mean size
from the useful estimate. Put
\begin{equation}\label{III-29-21}
 a=\hat\varepsilon+8d_*^{-2}(1+c_{\rm r})\varepsilon_m^2
                         +\gamma_2+\tfrac32\rho
                  \le8d_*^{-2}\Sigma_\rho.
\end{equation}
Multiplying the Gaussian bound, \eqref{III-29-19}, the Jacobian
bound and
\[
 1\le\exp\left(\gamma_2\|\beta'\|^2
               -\gamma_2d_*^2\sum_{b\in\mathfrak e^+}\mathfrak b_b^2\right)
\]
bounds the moved integrand by
\[
 F_0e^{(2/3)d_*^{-2}\varepsilon_m^2N_{\rm I}(Z)}
 e^{-\gamma_2d_*^2\sum_{b\in\mathfrak e^+}\mathfrak b_b^2}
 e^{-\frac12(x,(Q^Z_{\rm r}-2a)x)},\qquad x=(\mathbb X,\beta').
\]
Here $2a<\lambda_{\rm r}/2$, so $A=Q^Z_{\rm r}-2a$
is positive. Translate by the real mean and enlarge the real
integration domain. Gaussian integration leaves the determinant
ratio, for which
\[
 |\det Q^Z|^{1/2}(\det A)^{-1/2}
 \le e^{2n(a+\hat\varepsilon)/\lambda_{\rm r}}
 \le e^{4(a+\hat\varepsilon)N_{\rm I}(Z)/\lambda_{\rm r}}.
\]
This follows by factoring out $A$ and using
$|\det(1+D)|\le e^{n\|D\|}$.
The remaining volume coefficient is bounded by
\[
 4(a+\hat\varepsilon)\lambda_{\rm r}^{-1}
       +\tfrac23d_*^{-2}\varepsilon_m^2
 \le37d_*^{-2}\lambda_{\rm r}^{-1}\Sigma_\rho
 \le c_{\rm mn}u_\rho.
\]
Also $d_*^2\ge C_{\rm mn}^{-1}$. These inequalities prove
\eqref{III-29-11}. For an empty integration set the assertion
reduces to $|G(0)|\le F_0$.
\end{proof}

\begin{lemma}\label{III-30-5}
Let $\vartheta'>\vartheta_*$ and suppose that $G(B;n)$ is
admitted in $B$ on the fixed tube for every
$|n(b)|<\vartheta'\mathfrak b_b$, is holomorphic in $n$ for
every real $B$ in the tube, and satisfies
\[
 |G(B;n)|\le F_0e^{\rho\|B-n\|^2}.
\]
Under \eqref{eq:centered-gaussian-smallness},
\[
 s\longmapsto\mathbb J^m(Z_\circ;s)
                   [\mathbf1_DG(\cdot;m(Z_\circ,s))]
\]
is holomorphic on a neighbourhood of the closed decoupling
polydisc. It satisfies \eqref{III-29-11}. The same statement
holds after the affine extension of the parameters over a guard.
\end{lemma}

\begin{proof}
The mean lies strictly inside the indicated set of centers.
The integral and its estimate follow from \cref{III-29-9} at
each $s$. To prove holomorphy, continue the finitely many entries
of $C,C^{-1},\Gamma,m$ to a common open neighbourhood and
symmetrize the continued $S$ block in $Q$. Shrink the neighbourhood
so that $\|Q-Q_{\rm r}\|<2\hat\varepsilon+\lambda_{\rm r}/8$,
$\|m\|<\mathfrak b_\Sigma$ and
$|m(b)|<\vartheta'\mathfrak b_b$. It contains the closed
polydisc by the strict mean bound and continuity.
The same joint Gaussian calculation now uses
$\hat\varepsilon+\lambda_{\rm r}/16$ in place of
$\hat\varepsilon$. Its growth condition remains valid since
\[
 \hat\varepsilon+3\rho+\lambda_{\rm r}/16
 \le(1/108+1/16)\lambda_{\rm r}<\lambda_{\rm r}/8.
\]
Thus there is a common integrable majorant on this neighbourhood.
For every real $B$ the integrand is holomorphic in $s$; outside
$D$ it is identically zero. Dominated integration proves the
claim. Guard extension is affine and hence preserves holomorphy.
The zero-dimensional case is the constant $G(0;0)$.
\end{proof}

\subsection{Local potentials and their Mayer terms}

Keep the fine and coarse cube partitions, enlargement $Y^\sharp$,
guard and source-cube convention of Sections~\ref{subsec:ymm1-56-3} and~\ref{sec:history-cluster}.
Besides ordinary potentials $V^o(Y;B)$, bounded by
$0\le u^o(Y)\le1$ on their complex small-field range, allow
potentials $V^o(Y;z,\mu)$ holomorphic on
\begin{equation}\label{III-28-12}
 |z(b)|<\varsigma_{\rm ce}\mathfrak b_b,\qquad
 |\mu(b)|<\theta_{\rm sl}\mathfrak b_b,
 \qquad b\in\mathfrak X_{\rm I}^{\circ+}\langle Y\rangle.
\end{equation}
There are finitely many origins $o$, partitioned into ordinary
and centered origins. Require
\begin{equation}\label{III-28-4}
 \vartheta_{\rm ad}=\min\{\varsigma_{\rm ce}-1,\theta_{\rm sl}\}
                         >\vartheta_*.
\end{equation}
The name centered refers to their later evaluation at
$(z,\mu)=(B-m,m)$; it imposes no vanishing assumption.
For each centered potential assume
\begin{equation}\label{III-28-26}
 q_Y(z)=M^{-3}\sum_{b\in\mathfrak X_{\rm I}^{\circ+}\langle Y\rangle}
                      |z(b)|^2,\qquad
 |V^o(Y;z,\mu)|\le a^o(Y)(1+q_Y(z)),\quad a^o(Y)\ge0.
\end{equation}
Choose tilts $0<\tau^o(Y)\le1$ with
\begin{equation}\label{III-28-28}
 \sum_{o\ {\rm centered}}\sum_{Y\supset\square}
              (\tau^o(Y)+a^o(Y))\le\tfrac12L^{-3}
                  \quad\text{for each fine cube }\square.
\end{equation}
All domains in these sums satisfy $Y^\sharp\subset\mathfrak R$.
The weights must satisfy
\begin{equation}\label{III-28-29}
 w^o(Y)\ge e\,u^o(Y)\quad(o\ {\rm ordinary}),\qquad
 w^o(Y)\ge e\,a^o(Y)/\tau^o(Y)\quad(o\ {\rm centered}).
\end{equation}

Use the Mayer indices of the ordinary system with the same
origins and zero potentials. Thus $\iota=(\mathbf D,\mathfrak e)$
is a finite set of distinct potential vertices and failed bonds,
$Y^{\rm v}=\bigcup_{(Y,o)\in\mathbf D}Y$, and no failed bond
touches $Y^{\rm v}$. The region $Z_\circ(\iota)$ is the union
of the vertex domains and the fine bond cubes. Its integration
part $Z_\circ^{\rm I}$ contains exactly the fine cubes with an
endpoint of an integration bond. The guard $G_\circ$ consists of
the coarse cubes in $Z_\circ^\sharp$ meeting $\Lambda_t$.
One has $(Z_\circ^{\rm I})^\sharp\subset G_\circ\subset
Z_\circ^\sharp$ and the same integration bonds for $Z_\circ$
and $Z_\circ^{\rm I}$, by Section~\ref{subsec:ymm1-56-3}.

For $|n(b)|<\vartheta_{\rm ad}\mathfrak b_b$, define
\begin{equation}\label{III-28-16}
\begin{split}
 F_\iota^{\rm ce}(B;n)={}&(-1)^{\#\mathfrak e^+}
  \mathsf{in}_{\mathfrak X_{\rm I}\langle Y^{\rm v}\rangle}(B)
  \mathsf{out}_{\mathfrak e}(B)\\
 &\times\prod_{\substack{(Y,o)\in\mathbf D\\o\ {\rm ordinary}}}
                (e^{V^o(Y;B)}-1)
 \prod_{\substack{(Y,o)\in\mathbf D\\o\ {\rm centered}}}
                (e^{V^o(Y;B-n,n)}-1).
\end{split}
\end{equation}
A failed inner test makes this function zero before any
potential is evaluated. Extend the free parameters by one on
$G_\circ\setminus(Z_\circ^{\rm I})^\sharp$, writing $s^{\rm gd}$,
and put
\begin{equation}\label{III-28-17}
 f_\iota^{\rm ce}(s)=\mathbb J^m(Z_\circ^{\rm I};s^{\rm gd})
       [F_\iota^{\rm ce}(\cdot;m(Z_\circ^{\rm I},s^{\rm gd}))].
\end{equation}
For a finite set $\mathcal K$ of free cubes, set
\begin{equation}\label{III-28-18}
 T^{\rm ce}(\iota,\mathcal K)
 =\sum_{W\subset\mathcal K}(-1)^{\#\mathcal K-\#W}
                            f_\iota^{\rm ce}(\mathbf1_W),\qquad
 X(\iota,\mathcal K)=Z_\circ^\sharp\cup\bigcup\mathcal K.
\end{equation}
Each summand uses its own conditional mean.

Put
\begin{equation}\label{III-28-6}
 \rho_+=(LM)^{-3},\qquad u_0^+=u_{\rho_+},\qquad
 c_{\rm P}^+=C_{\rm mn}(LM)^3c_{\rm mn}u_0^+,\qquad
 \gamma_P=C_{\rm mn}^{-1}\gamma_2\mathfrak b_-^2.
\end{equation}
In the rest of this section require
\begin{equation}\label{III-28-7}
 C_{\rm mn}u_0^+\le\lambda_{\rm r}.
\end{equation}

\begin{proposition}\label{III-30-22}
Every term in \eqref{III-28-18} is defined, and
$f_\iota^{\rm ce}$ is holomorphic near its closed polydisc.
If $N_\Lambda(X)$ counts the coarse cubes of $X$ meeting
$\Lambda_t$, then
\begin{equation}\label{III-30-23}
 |T^{\rm ce}(\iota,\mathcal K)|
 \le\bar T^+(\iota,\mathcal K)
 :=e^{c_{\rm P}^+N_\Lambda(X)}
       \prod_{(Y,o)\in\mathbf D}w^o(Y)
       e^{-\gamma_P\#\mathfrak e^+}
       e^{-(\bar\kappa-1)\#\mathcal K}.
\end{equation}
The same bound without its last factor holds for
$|f_\iota^{\rm ce}(s)|$ when $s$ vanishes off $\mathcal K$.
\end{proposition}

\begin{proof}
Remove the tests from \eqref{III-28-16}, retaining their sign,
and call the resulting product $G_\iota(B;n)$. On the tube of
\cref{III-30-5}, every vertex variable obeys
$|B(b)|\le|\operatorname{Re}B(b)|+|\operatorname{Im}B(b)|
<\mathfrak b_b$. Thus $B-n$ and $n$ lie in
\eqref{III-28-12}. The product is admitted in $B$ and
holomorphic in $n$.

For an ordinary factor, $|e^V-1|\le e|V|\le w^o(Y)$.
For a centered factor put $q=q_Y(B-n)$, $a=a^o(Y)$ and
$\tau=\tau^o(Y)$. Its domain contains a source cube, so
\eqref{III-28-28} gives $a\le1$. Since
$1+q\le\tau^{-1}e^{\tau q}$,
\begin{equation}\label{III-30-20}
 |e^V-1|\le a(1+q)e^{a(1+q)}
       \le(ea/\tau)e^{(a+\tau)q}
       \le w^o(Y)e^{(a+\tau)q}.
\end{equation}
In the sum of the last exponents, a positive bond $b$ can occur
only for domains containing the unique fine cube of its initial
site. Its coefficient is therefore at most $L^{-3}/2$ by
\eqref{III-28-28}. For a complex Lie algebra element,
$|z(b)|^2\le2(|\operatorname{Re}z(b)|^2+
|\operatorname{Im}z(b)|^2)$. Consequently
\begin{equation}\label{III-30-19}
 |G_\iota(B;n)|\le
      \prod_{(Y,o)\in\mathbf D}w^o(Y)\,
                         e^{\rho_+\|B-n\|^2}.
\end{equation}
The factor $1/2$ in \eqref{III-28-28} compensates for this
comparison of the operator and coordinate norms.

Apply \cref{III-30-5} with $\rho=\rho_+$, the inner tests of
$Y^{\rm v}$, the failed set $\mathfrak e$, and
$\vartheta'=\vartheta_{\rm ad}$. For $s$ supported on
$\mathcal K$, its guarded extension is supported on
$G_\circ\cup\bigcup\mathcal K\subset X$.
The bond count \eqref{II-117-94} is
$N_{\rm I}(Z)\le18(LM)^3N_\Lambda(Z)$; since $C_{\rm mn}\ge18$,
\eqref{III-29-11} gives the asserted bound on $f_\iota^{\rm ce}$.
The finite-difference Cauchy estimate on the free coordinates
has factor $(e^{\bar\kappa}-1)^{-\#\mathcal K}
\le e^{-(\bar\kappa-1)\#\mathcal K}$. This proves
\eqref{III-30-23}.
\end{proof}

\subsection{Factorization and convergent cluster sums}

A pair $(\iota,\mathcal K)$ is rooted if every wall component
of $G_\circ\cup\bigcup\mathcal K$ contains a cube of $G_\circ$.
For a coarse connected domain $\gamma$, define
\begin{equation}\label{III-28-19}
 \mathsf z^{\rm ce}(\gamma)
 =\sum_{\substack{(\iota,\mathcal K)\ {\rm rooted}\\
                        X(\iota,\mathcal K)=\gamma}}
                         T^{\rm ce}(\iota,\mathcal K).
\end{equation}
The sum is finite. Compatibility means that the coarse domains
do not touch. With the corresponding truncated function $\rho^T$,
put
\begin{equation}\label{III-28-20}
 E^{\rm ce}(X)=\sum_{n\ge1}\frac1{n!}
   \sum_{\gamma_1\cup\cdots\cup\gamma_n=X}
       \rho^T(\gamma_1,\ldots,\gamma_n)
             \prod_{i=1}^n\mathsf z^{\rm ce}(\gamma_i).
\end{equation}
Expanding each activity gives the refined series. Its part
$E^{{\rm ce},\mathsf{Mk}}(X)$ through a specified set of
vertices and failed bonds consists of the terms with at least
one such vertex or bond. A marked bond set is closed under reversal.

\begin{proposition}\label{III-29-61}
Under the preceding Gaussian and potential hypotheses,
nonrooted terms vanish, and a term whose domain has wall
components $\gamma_i$ factors into the corresponding terms
on those components. For a complete Gaussian structure on the
torus, with whole covariance $C_*$ and mean $m_*$ as in
\eqref{II-114-13}, the whole tested integral is
\begin{equation}\label{III-29-66}
 \begin{split}
 \Xi^{\rm ce}&=\int d\nu(B)\,\mathsf{in}_{\mathfrak X_{\rm I}}(B)
 \exp\left(\sum_{o\ {\rm ordinary},Y}V^o(Y;B)
       +\sum_{o\ {\rm centered},Y}V^o(Y;B-m_*,m_*)\right)
 \\&=\sum_{\mathcal P'}\prod_{\gamma\in\mathcal P'}
                                      \mathsf z^{\rm ce}(\gamma).
 \end{split}
\end{equation}
Here $\nu$ is the normalized complex Gaussian with covariance
$C_*$ and mean $m_*$, the vertex sums run over domains containing
a source cube, and $\mathcal P'$ runs over compatible polymer
families, including the empty family. A zero test is read before
the potential sum.
\end{proposition}

\begin{proof}
The proof of Proposition~\ref{II-116-68} applies once the
center is included in the restriction argument. If a component
of the parameter cubes is disjoint from the integration guard,
deleting it leaves $C,\Gamma,m$ on the integration bonds
unchanged. The potential product uses the same restricted mean.
The alternating sum in any cube of that component therefore
vanishes. This proves the nonrooted assertion.

When the domain splits into wall components, every enlarged
vertex domain and failed-bond domain belongs to a unique
component. The Gaussian covariance is a direct sum, the
connecting operators use disjoint auxiliary coordinates, and
the mean restricts to the component mean. Thus the centered
arguments restrict as well. The integrands, determinant branches
and finite differences factor. Absolute convergence from
\cref{III-29-9} permits Fubini. Restriction of indices to the
components is a bijection: conversely, their union preserves
the failed-bond exclusions and rootedness. This proves the
factorization of terms and of their sums.

For the whole integral, first freeze the center at $m_*$.
On the real small-field domain each ordinary potential is
bounded by its prescribed size. The centered potentials are
bounded by their quadratic-growth estimates, since the thresholds
are finite and the center is fixed. Their tested exponentials
are therefore bounded Borel functions. The identity between the conditioning
integral and the marginal of $\nu$ in
Proposition~\ref{prop:history-whole-integral} applies to each such function.
Completeness gives $m(Z_\circ,\mathbf1)=m_*|_{\mathfrak f}$.
Expand the finitely many factors into $1+(e^V-1)$ and the
remaining tests into $1-\mathsf{out}$. The resulting functions
are exactly $F_\iota^{\rm ce}(B;m_*|_{\mathfrak f})$.
Their conditioning integrals at all parameters equal to one
are $f_\iota^{\rm ce}(\mathbf1)$. Summing all finite differences
recovers this value, even though each difference has its own
center. Remove the nonrooted terms and group by wall components.
This gives \eqref{III-29-66}.
\end{proof}

Suppose, for rates $r_{\mathsf w},r_H,\kappa''\ge0$,
\begin{equation}\label{III-28-34}
 \sum_o w^o(Y)\le\varepsilon_{\mathsf w}e^{-r_{\mathsf w}d_k(Y)}.
\end{equation}
Let $c_{\mathrm E}=\log18432$, $C_{\mathrm E}=4096$ and
$C_{\rm cs}\ge\max\{C'_{\rm col},24C_{\mathrm E}^2e^{24}\}$,
where $C'_{\rm col}=13/9$ is the coarsening constant. Define
\begin{equation}\label{III-28-8}
\begin{aligned}
 r_{\rm pol}^+&=r_H-C_{\rm cs}(c_{\rm P}^++1),&
 C_{\rm pol}&=C_{\rm cs}e^{C_{\rm cs}\kappa''},\\
 \varepsilon_H^+&=C_{\rm cs}e^{C_{\rm cs}(c_{\rm P}^++r_H)}
       \bigl(L^3\varepsilon_{\mathsf w}+(LM)^3e^{-\gamma_P}\bigr),\\
 K_{\rm mk}^+&=C_{\rm cs}e^{C_{\rm cs}(c_{\rm P}^++\kappa'')},&
 K_{\rm or}^+&=C_{\rm cs}(LM)^3
                   e^{C_{\rm cs}(c_{\rm P}^++r_H+\kappa'')}.
\end{aligned}
\end{equation}
Require the four inequalities
\begin{equation}\label{eq:centered-cluster-rates}
\begin{gathered}
 (C_{\rm cs}/L)r_H+c_{\mathrm E}\le r_{\mathsf w},\qquad
 e^{-(\bar\kappa-1)+C_{\rm cs}r_H}\le C_{\rm cs}^{-1},\\
 r_{\rm pol}^+\ge\kappa''+c_{\mathrm E}+C_{\rm cs},\qquad
 C_{\rm pol}\varepsilon_H^+\le\tfrac12.
\end{gathered}
\end{equation}
Define the nonnegative activity and refined cluster majorants
by replacing each term $T^{\rm ce}$ with $\bar T^+$ and
$\rho^T$ with $|\rho^T|$. Write them as
$\bar{\mathsf z}^+$, $\bar E^+$ and $\bar E^{\mathsf{Mk},+}$.
Set
\begin{equation}\label{eq:centered-marked-size}
\begin{split}
 \eta^{\mathsf{Mk},+}=\sup_\Delta\biggl[
 &\sum_{\substack{(Y,o)\in\mathsf{Mk}\\\Delta\subset Y^\sharp}}
               w^o(Y)e^{r_H(d_t(Y^\sharp)+4)}\\
 &+\#\{b\in\mathsf{Mk}^+:\Delta\subset b^\sharp\}
                                e^{-\gamma_P+4r_H}\biggr].
\end{split}
\end{equation}

\begin{theorem}\label{III-30-36}
Under \eqref{III-28-7}, \eqref{III-28-28},
\eqref{III-28-29}, \eqref{III-28-34} and
\eqref{eq:centered-cluster-rates}, the refined series converges
absolutely and
\begin{align}
 |\mathsf z^{\rm ce}(\gamma)|&\le\bar{\mathsf z}^+(\gamma)
                  \le\varepsilon_H^+e^{-r_{\rm pol}^+d_t(\gamma)},
                                                   \label{eq:centered-activity-bound}\\
 |E^{\rm ce}(X)|&\le\bar E^+(X)
                  \le C_{\rm pol}\varepsilon_H^+e^{-\kappa''d_t(X)},
                                                   \label{eq:centered-cluster-bound}\\
 |E^{{\rm ce},\mathsf{Mk}}(X)|&\le\bar E^{\mathsf{Mk},+}(X)
                  \le K_{\rm mk}^+\eta^{\mathsf{Mk},+}e^{-\kappa''d_t(X)}.
                                                   \label{eq:centered-marked-bound}
\end{align}
The sums vanish when no coarse cube of their domain meets
$\Omega_t$. For complete whole-torus data,
\begin{equation}\label{III-30-40}
 \Xi^{\rm ce}=\exp\sum_XE^{\rm ce}(X)\ne0.
\end{equation}
These are assertions about this fluctuation integral.
\end{theorem}

\begin{proof}
Use an auxiliary ordinary system with the same origins and all
potentials zero, the weights $w^o(Y)$ of both kinds, and the
larger Gaussian error $\hat\varepsilon+\rho_+$.
The Gaussian structure still satisfies all its hypotheses:
only upper bounds are enlarged. Its value of $u_0$ is exactly
$u_0^+$, and its Gaussian volume constant is $c_{\rm P}^+$.
Its index sets, domains and majorants are precisely those just
defined. The weight decay and the four numerical hypotheses
are those of Theorem~\ref{II-118-20} at these constants.
The counting proof of that theorem therefore gives the
three majorant bounds above. By \cref{III-30-22} the actual
refined terms are dominated term by term by these majorants.
Absolute convergence permits all regroupings, proving the
activity, cluster and marked bounds. This use of the ordinary
theorem concerns its majorant series, not its zero activities.
The geometric vanishing statement follows from the same
index sets. Finally \cref{III-29-61} identifies the whole
integral with the compatible polymer sum, and the convergent
cluster identity gives \eqref{III-30-40}.
\end{proof}

\begin{proposition}\label{III-30-53}
The activities and cluster sums restrict to a smaller region
when the Gaussian structure, thresholds and potential members
restrict there. They depend only on the potential members with
$Y^\sharp\subset X$ and the Gaussian data used by the corresponding
indices. They are holomorphic in finite dimensional holomorphic
families of these data satisfying the same bounds, and Borel in
Borel families. They inherit gauge invariance, conjugation, and
dependence only on the represented pair whenever these properties
hold for the Gaussian structure and the potential members.
\end{proposition}

\begin{proof}
Each term is a finite difference of conditional integrals.
Restriction identifies their integration bonds, the guard,
the local matrices and the restricted means. It consequently
identifies the centered arguments, integrals and terms.
Holomorphic parameter families give holomorphic terms by
\cref{III-30-5}; locally one may choose its open neighbourhood
with a common integrable Gaussian majorant. The cluster
majorant is independent of the external parameters, so the
refined sums converge locally uniformly and remain holomorphic.
For Borel families, the joint integrands are Borel and the same
majorants give Borel parameter integrals and limits of partial
sums. A gauge change acts by real orthogonal adjoint maps on
the integration variables and carries the mean by the same maps.
Their Jacobians, the tests and the Gaussian branch are unchanged.
Complex conjugation conjugates the matrices, mean and potentials;
the branch based at real positive matrices conjugates with them.
These changes preserve the finite differences and the absolutely
convergent sums. Finally, if the Gaussian data and potential
members agree for two representatives of one pair, their
integrals agree term by term. This last assertion presupposes
the stated agreement; it does not supply it for arbitrary
retained quadratic factors.
\end{proof}


\section{The centered expansion on a history}
\label{sec:centered-history}

We apply the preceding expansion at a step $t=k+1$ of a fixed
history $h'$. The Gaussian system is the canonical one of
Theorem~\ref{II-115-62} and Proposition~\ref{II-115-84}. Its graded triple
domain is $\tilde U_{\sigma_t^{\rm mid}}(Y)$, and
$\mathfrak Q_h^{\rm rad}(Y)$ is the space of stored fields in
their prescribed tiers. Define
\begin{equation}\label{eq:centered-tier-domain}
 \mathfrak D_k^{\rm rad}(Y)=
 \tilde U_{\sigma_t^{\rm mid}}(Y)
 \times\{\mathcal A:|g_k^{-1}\mathcal A(b)|
                   <C_{\mathcal A}\mathfrak b_k\}
 \times\mathfrak Q_h^{\rm rad}(Y).
\end{equation}
The field inequality is imposed on all positively oriented
bonds of $\mathfrak X^\circ\langle Y\rangle$. A centered
potential has, in addition, an argument
\begin{equation}\label{III-31-10}
 \mu\in\mathbb S_k^{\rm rad}(Y)
 :=\{\mu:|\mu(b)|<\theta_{\rm sl}\mathfrak b_k,
                    \ b\in\mathfrak X_{\rm I}^{\circ+}\langle Y\rangle\}.
\end{equation}
An empty set of bonds contributes a one-point space.
These domains use the running cutoff $\mathfrak b_k$, rather
than the individual integration thresholds.

Analyticity on these domains means local holomorphic extension
in $(A',\mathbf J,\mathcal A,\mu)$ at fixed real $U$ and fixed
stored argument. A parameter family is analytic when these
extensions are jointly holomorphic in its complex parameters.
Borel measurability is with respect to the trace of the Borel
sigma algebra of the finite dimensional ambient product.
The ordinary norm is the supremum norm with weight
$e^{r d_k(Y)}$. The centered norm is
\begin{equation}\label{III-31-12}
 \|\mathbb W\|_{{\rm I}\star,r}^{\rm rad}
 =\sup_{Y,\mathfrak u,\mu}
 \frac{e^{r d_k(Y)}|\mathbb W(Y;\mathfrak u,\mu)|}
 {1+M^{-3}\sum_{b\in\mathfrak X_{\rm I}^{\circ+}\langle Y\rangle}
                         g_k^{-2}|\mathcal A(b)|^2}.
\end{equation}
Here $\mathfrak u=(P,\mathcal A,\mathfrak q)$ ranges over
\eqref{eq:centered-tier-domain}, and $\mu$ over
\eqref{III-31-10}. In particular the centered norm is a
supremum norm in the center, the external field and the stored
argument. Its permitted growth is confined to the integration field.

For $X_1\in\mathcal D_t$, write
$\mathfrak z=(\mathsf a;P;\mathfrak q)$ for a point with
$P\in\tilde U_{\sigma_t^{\rm mid}}(X_1)$,
$\mathfrak q\in\mathfrak Q_h^{\rm rad}(X_1)$ and
$|\mathsf a(b)|<C_{\mathcal A}\mathfrak b_k$ on its external
bonds. Require also the external mean bounds of
Theorem~\ref{II-115-62}, at parameters
$(\theta_{\rm xt},\varepsilon_{\rm xt})$. We denote this
point set by $\mathsf{Pt}_{h'}\langle X_1\rangle$.
In the canonical application these bounds are supplied by
Proposition~\ref{II-115-84}. For $Y^\sharp\subset X_1$, put
\begin{equation}\label{III-31-14}
 \mathfrak u_Y(z;\mathfrak z)
 =\bigl(P|_Y,g_k(z\oplus\mathsf a|_Y),\mathfrak q|_Y\bigr),
 \qquad
 \mathbb W\langle\mathfrak z\rangle^{\rm ce}(Y;z,\mu)
 =\mathbb W(Y;\mathfrak u_Y(z;\mathfrak z),\mu).
\end{equation}
The same substitution without $\mu$ evaluates an ordinary potential.

\subsection{Constants and admissible families}

Use $C_{\rm mn}$ and $C_{\rm cs}$ from the preceding section.
The constants denoted by $\mathsf C_{\rm mn}$ and
$\mathsf C_{\rm cs}$ in formulas for the history are these same
constants. Set
\begin{equation}\label{III-31-4}
 \rho_+=(LM)^{-3},\qquad
 u_0^+=u_0+C_{\rm mn}\rho_+,\qquad
 c_{\rm P}^+=c_{\rm P}+C_{\rm mn}^3\lambda_{\rm r}^{-1}.
\end{equation}
For the canonical Gaussian constants define
\begin{equation}\label{III-31-5}
 N_{\rm gs}=C_{\rm res}B_{\rm nm}^5\gamma_{\rm pos}^{-2}
                  +c_{\rm r}+2,\qquad
 \bar c_{\rm P}^+=C_{\rm mn}^3\lambda_{\rm r}^{-1}(N_{\rm gs}+1).
\end{equation}
The bound $B_{\rm nm}$ and the positive gap $\gamma_{\rm pos}$
are those of the real reference matrices in \cref{II-115-62}. Here $C_{\rm res}=64$, the coefficient in its
Gaussian error bound, and we specialize
\[
 \gamma_2=(LM)^{-3},\qquad
 (\mathfrak b_-,\mathfrak b_+)=(\theta_c\mathfrak b_k,\mathfrak b_k),
 \qquad \theta_m=\theta_{\rm xt},\qquad
 \varepsilon_m=\sqrt2\varepsilon_{\rm dc}\theta_{\rm xt}
                     +(1+\varepsilon_{\rm dc})\varepsilon_{\rm xt}.
\]
For $\kappa>2$, $0<\delta_{\rm sp}<1/2$ and
$\beta_{\rm s}>0$, put
\begin{equation}\label{eq:centered-history-rates}
 r_V=(1-2\delta_{\rm sp})(\kappa-2),\qquad
 \kappa''=(1+4\beta_{\rm s})\kappa+1+c_{\mathrm E}.
\end{equation}
Define
\begin{equation}\label{III-31-6}
\begin{gathered}
 r_{\mathsf w}^+=r_V-c_{\mathrm E},\qquad
 r_H^+=\kappa''+c_{\mathrm E}+C_{\rm cs}(c_{\rm P}^++2),\\
 \bar r_H^+=\kappa''+c_{\mathrm E}+C_{\rm cs}(\bar c_{\rm P}^++2),
 \qquad
 K_{\rm ap}^+=C_{\rm cs}^2
       e^{C_{\rm cs}(\kappa''+\bar c_{\rm P}^++\bar r_H^+)}.
\end{gathered}
\end{equation}
The constants $r_{\rm pol}^+$, $K_{\rm mk}^+$ and
$K_{\rm or}^+$ are those of \eqref{III-28-8} at $r_H=r_H^+$.
The useful ceilings are
\begin{equation}\label{III-31-8}
 w_0=(4C_{\mathrm E}L^3)^{-1},\qquad
 \varepsilon^{\rm O}=(16eL^3K_{\rm ap}^+)^{-1},\qquad
 c_M^{\rm O}=w_0\varepsilon^{\rm O}.
\end{equation}

We use the following strengthened history hypotheses. The
geometric, graded-radius and covariance hypotheses of
Theorem~\ref{II-115-62} hold on the stated point sets; in
particular the local Gaussian structure, its restrictions and
its parameter dependence are defined there. In addition,
\begin{equation}\label{eq:centered-history-smallness}
\begin{gathered}
 1\le\theta_c\mathfrak b_k\le\mathfrak b_b\le\mathfrak b_k,
 \qquad \varepsilon_{\rm dc}\le(LM)^{-3},\\
 \theta_m+\varepsilon_m\le\vartheta_*<1,
 \qquad (LM)^3\varepsilon_m^2\le1,\qquad
 1+\vartheta_*<C_{\mathcal A},\quad
 \vartheta_*<\theta_{\rm sl},\\
 (LM)^3\ge C_{\rm mn}^{-1}\bar c_{\rm P}^+,
 \qquad L\ge\frac{2C_{\rm cs}(1+4\beta_{\rm s})}{1-2\delta_{\rm sp}},\\
 \tfrac12(1-2\delta_{\rm sp})\kappa
 \ge2(1-2\delta_{\rm sp})+2c_{\mathrm E}
 +\frac{C_{\rm cs}}L[1+2c_{\mathrm E}+C_{\rm cs}(\bar c_{\rm P}^++2)],\\
 \bar\kappa\ge1+C_{\rm cs}\bar r_H^++\log C_{\rm cs},
 \qquad (LM)^3e^{-\gamma_P}\le(4K_{\rm ap}^+)^{-1}.
\end{gathered}
\end{equation}
The first inequality is required even if there are no
integration bonds. These are the strengthened Gaussian, range
and rate conditions used below.

\begin{definition}\label{III-31-23}
A frame with a center consists of finitely many analytic
ordinary families $\mathbb W^o$, $o\in\mathsf{Or}^{\rm f}$,
an onto map $p:\mathsf{Or}^{\rm f}\longrightarrow\mathsf{Or}^{\rm c}$,
and finitely many analytic centered families indexed by a
disjoint set $\mathsf{Or}^{\rm ce}$. Put
\[
 \mathbb W^{{\rm c},o'}=\sum_{p(o)=o'}\mathbb W^o,
 \qquad \mathbb W^{\rm ce}=\sum_{o\in\mathsf{Or}^{\rm ce}}\mathbb W^o.
\]
It includes sizes $s^{o'}\in[0,1]$, $s^{\rm ce}\in[0,w_0]$ with
\[
 \|\mathbb W^{{\rm c},o'}\|_{r_V}^{\rm rad}\le s^{o'},
 \qquad \|\mathbb W^{\rm ce}\|_{{\rm I}\star,r_V}^{\rm rad}
                \le s^{\rm ce},\qquad
 \varepsilon_{\mathsf w}^{\mathfrak F+}
   :=e\sum_{o'}s^{o'}+e w_0^{-1}s^{\rm ce}
   \le(4K_{\rm ap}^+L^3)^{-1}.
\]
If there is no centered origin, require $s^{\rm ce}=0$.
The fine system consists of all the members. The coarse system
consists of the ordinary sums $\mathbb W^{{\rm c},o'}$ and,
when present, the single centered sum $\mathbb W^{\rm ce}$.
The individual members need not satisfy the bounds on these sums.
\end{definition}

Define the coarse-system weights by
\begin{equation}\label{III-31-24}
\begin{aligned}
 u^{o'}(Y)&=s^{o'}e^{-r_Vd_k(Y)},&
 w^{o'}(Y)&=e u^{o'}(Y),\\
 a^{\rm ce}(Y)&=s^{\rm ce}e^{-r_Vd_k(Y)},&
 \tau^{\rm ce}(Y)&=w_0e^{-c_{\mathrm E}d_k(Y)},\\
 w^{\rm ce}(Y)&=e w_0^{-1}s^{\rm ce}e^{-r_{\mathsf w}^+d_k(Y)},&
 b_{\rm cs}^{\mathfrak F+}&=C_{\rm pol}\varepsilon_H^{\mathfrak F+},
\end{aligned}
\end{equation}
where $\varepsilon_H^{\mathfrak F+}$ is \eqref{III-28-8} at
$r_H=r_H^+$ and amplitude $\varepsilon_{\mathsf w}^{\mathfrak F+}$.

\begin{proposition}\label{III-31-40}
Under \eqref{eq:centered-history-smallness}, the coarse system
of every frame with a center, evaluated at any
$\mathfrak z\in\mathsf{Pt}_{h'}\langle X_1\rangle$, satisfies
the hypotheses of \cref{III-30-36}. The fine system has the
same Gaussian data and defined analytic conditional integrals.
Moreover
\begin{equation}\label{III-31-41}
 b_{\rm cs}^{\mathfrak F+}
 \le K_{\rm ap}^+
 [L^3\varepsilon_{\mathsf w}^{\mathfrak F+}+(LM)^3e^{-\gamma_P}]
 \le\tfrac12.
\end{equation}
The bounds and majorants are independent of the point and of
the individual thresholds within their stated ranges.
\end{proposition}

\begin{proof}
For an integration argument $|z(b)|<C_{\mathcal A}\mathfrak b_b$,
the inequality $\mathfrak b_b\le\mathfrak b_k$ and the external
field bound place $\mathfrak u_Y(z;\mathfrak z)$ in
\eqref{eq:centered-tier-domain}. Similarly
$|\mu(b)|<\theta_{\rm sl}\mathfrak b_b$ puts the center in
\eqref{III-31-10}. The substitution is linear in the fields.
To check joint analyticity, extend $\mu$ by zero to the full
set of integration bonds, treat that extension as an extra
parameter, and compose the local holomorphic extension with
restriction. This gives a jointly holomorphic function on the
required product domain. The norm \eqref{III-31-12} now gives
\[
 |\mathbb W^{\rm ce}\langle\mathfrak z\rangle^{\rm ce}(Y;z,\mu)|
       \le a^{\rm ce}(Y)(1+q_Y(z)),
\]
since $\mathcal A=g_kz$ at the integration bonds.

The entropy estimate of Proposition~\ref{II-18-26} gives
\[
 \sum_{Y\supset\square}
       [\tau^{\rm ce}(Y)+a^{\rm ce}(Y)]
 \le C_{\mathrm E}(w_0+s^{\rm ce})\le\tfrac12L^{-3}.
\]
Here $r_V\ge c_{\mathrm E}$ follows from
\eqref{eq:centered-history-smallness}. Thus the tilts fit,
and $w^{\rm ce}=e a^{\rm ce}/\tau^{\rm ce}$. The ordinary
bounds are immediate from their norms. Both types of weights
are bounded in sum by
$\varepsilon_{\mathsf w}^{\mathfrak F+}e^{-r_{\mathsf w}^+d_k(Y)}$.

The Gaussian error formula and the mean bound give
\[
 u_0^+\le C_{\rm mn}(LM)^{-3}(N_{\rm gs}+1),\qquad
 c_{\rm P}^+\le\bar c_{\rm P}^+.
\]
Consequently
$C_{\rm mn}u_0^+\le\lambda_{\rm r}$ by the lower bound on
$(LM)^3$. The inequalities for $L$ and $\kappa$ give
\[
 \frac{C_{\rm cs}}L r_H^++c_{\mathrm E}
 \le\frac12(1-2\delta_{\rm sp})\kappa+c_{\mathrm E}
   +\frac{C_{\rm cs}}L[1+2c_{\mathrm E}+C_{\rm cs}(\bar c_{\rm P}^++2)]
 \le r_V-c_{\mathrm E}.
\]
The bound on $\bar\kappa$ gives the second rate inequality in
\eqref{eq:centered-cluster-rates}, and
$r_{\rm pol}^+=\kappa''+c_{\mathrm E}+C_{\rm cs}$ gives the third.
Finally $c_{\rm P}^+\le\bar c_{\rm P}^+$ and
$r_H^+\le\bar r_H^+$ bound the exponential constant by
$K_{\rm ap}^+$, proving \eqref{III-31-41} and the last rate
inequality. The center is admitted because
$\vartheta_*<\min(C_{\mathcal A}-1,\theta_{\rm sl})$.
The canonical Gaussian hypotheses supply integrability,
restriction and completeness. These checks establish every
input to the abstract expansion.

For the fine system, each individual centered member is
holomorphic on the same domain, and each conditional integral
uses finitely many members on a compact subset of that domain.
The Gaussian integrability argument therefore applies. No
separate uniform size for an individual member is claimed.
All majorant formulas use the coarse sizes and fixed Gaussian
constants; none uses the point or an individual threshold.
\end{proof}

\subsection{Expansion and perturbation estimates}

\begin{theorem}\label{III-32-12}
Under the hypotheses of \cref{III-31-40}, the fine and coarse
systems have the same polymer activities and cluster sums.
Write the latter as $E[\mathfrak F^+]_{h'}(X|X_1;\mathfrak z)$.
They satisfy
\begin{equation}\label{III-32-15}
 |E[\mathfrak F^+]_{h'}(X|X_1;\mathfrak z)|
 \le\bar E^+(X)\le b_{\rm cs}^{\mathfrak F+}e^{-\kappa''d_t(X)}.
\end{equation}
They vanish if $X$ contains no coarse cube meeting $\Omega_t$.
Restriction to $X$ leaves the cluster sum at $X$ unchanged.
Its value depends only on the potential members with
$Y^\sharp\subset X$ and the restricted Gaussian data and
thresholds. On the whole torus,
\begin{equation}\label{III-32-17}
\begin{split}
 &\int d\nu(B)\,\mathsf{in}_{\mathfrak X_{\rm I}}(B)
 \exp\biggl[\sum_{Y,o\in\mathsf{Or}^{\rm f}}
                   \mathbb W^o\langle\mathfrak z\rangle(Y;B)
 +\sum_{Y,o\in\mathsf{Or}^{\rm ce}}
       \mathbb W^o\langle\mathfrak z\rangle^{\rm ce}
                           (Y;B-m_*,m_*)\biggr]\\
 &\hspace{25mm}=\exp\sum_XE[\mathfrak F^+]_{h'}(X;\mathfrak z|_X)\ne0.
\end{split}
\end{equation}
The sums over $Y$ contain the potential vertices of the
expansion, and the integrand is zero before evaluating its
exponent wherever a small-field test fails.

With fixed sizes, the activities and cluster sums are
holomorphic in holomorphic parameter families and Borel in
Borel families. Gauge invariance and conjugation are inherited
from the members. At fixed external field $\mathsf a$ and
stored argument $\mathfrak q$, equal represented pairs give
equal values if the members depend only on the pair. It suffices instead
that the two triples are real and the members have equal
values at twins. The marked estimate
\eqref{eq:centered-marked-bound} holds for the coarse system.
\end{theorem}

\begin{proof}
Apply \cref{III-30-36,III-30-53} to the coarse data admitted
by \cref{III-31-40}. To identify fine and coarse activities,
group the ordinary origins over each fiber of $p$ and all
centered origins into one. At any fixed domain and fixed
conditional mean, the finite algebraic identity
\[
 e^{\sum_{j=1}^nV_j}-1
       =\sum_{\emptyset\ne I\subset\{1,\ldots,n\}}
                    \prod_{j\in I}(e^{V_j}-1)
\]
expands each merged vertex. Replacing it by the chosen fine
vertices changes neither the union of their domains, the
guard, the failed-bond exclusions nor the conditional mean.
Thus the conditional integrals and finite differences agree
after this finite regrouping, at every decoupling parameter.
Rootedness and connected domains are preserved. Summing gives
equal activities, hence equal convergent cluster sums. This
argument requires bounds only on the merged system.

The restriction identities of the canonical Gaussian system
are given in \cref{II-115-84}, equation~\eqref{II-119-54}. Restriction also
identifies $\mathfrak u_Y$ and the center arguments. Therefore
the abstract restriction result applies term by term, including
to the majorants. Completeness identifies the whole mean with
each local mean at full coupling, and \cref{III-30-36} gives
\eqref{III-32-17}. The finite grouping leaves its integrand
unchanged.

For parameter dependence the substitutions in
\eqref{III-31-14} are jointly holomorphic at fixed real $U$
and fixed stored argument. The Gaussian structure has this
same property by \cref{II-115-62}. The common cluster
majorant therefore proves holomorphy of the sum. The
substitutions are continuous in ambient coordinates, so they
also preserve trace Borel measurability. Gauge changes rotate
$z$ and $\mu$ by the same real orthogonal adjoint map;
conjugation acts on both arguments and on the Gaussian mean.
The abstract covariance statements apply. Finally, the
canonical Gaussian matrices and mean depend only on the pair.
For two triples with that pair, the points in
\eqref{III-31-14} differ only in the triple. The asserted
equality of the potential values therefore identifies every
integral. In the real case the triples are twins, so the tie
suffices. No such conclusion for arbitrary complex twins is
used here.
\end{proof}

\begin{proposition}\label{III-32-22}
Split the ordinary coarse origins into old and new ones, and
regard every centered origin as new. Let $\mathfrak F_0$ be
the old ordinary subsystem and set
\begin{equation}\label{III-32-23}
 n_+=e\sum_{o'\ {\rm new}}s^{o'}+e w_0^{-1}s^{\rm ce},
 \qquad K_{\rm Lip}^+=2C_{\mathrm E}L^3e^{4r_H^+}K_{\rm mk}^+.
\end{equation}
Then $E[\mathfrak F^+]-E[\mathfrak F_0]$ is exactly the part
of the refined cluster series meeting a new coarse vertex, and
\begin{equation}\label{III-32-24}
 |E[\mathfrak F^+]_{h'}(X)-E[\mathfrak F_0]_{h'}(X)|
 \le \min(K_{\rm or}^+,K_{\rm Lip}^+)n_+e^{-\kappa''d_t(X)}.
\end{equation}
If the new members vanish for every $Y^\sharp\subset X$,
the difference at $X$ vanishes.

For two frames with the same old members and the same named
sizes, suppose each new coarse sum has at most half its named
size. Let $d^{o'}$ and $d^{\rm ce}$ be the corresponding norms
of the differences of those sums. Put
\begin{equation}\label{III-32-26}
 n_{12}=e\sum_{o'\ {\rm new}}d^{o'}+e w_0^{-1}d^{\rm ce},
 \qquad n=\max\left\{\max_{o'\ {\rm new}}\frac{d^{o'}}{s^{o'}},
                              \frac{d^{\rm ce}}{s^{\rm ce}}\right\}.
\end{equation}
Here $0/0=0$ and an empty maximum is zero. Then $0\le n\le1$,
and
\begin{align}
 |E[\mathfrak F_1^+]_{h'}(X)-E[\mathfrak F_2^+]_{h'}(X)|
 &\le4\min(K_{\rm or}^+,K_{\rm Lip}^+)n_{12}e^{-\kappa''d_t(X)},
                                                  \label{III-32-27}\\
 |E[\mathfrak F_1^+]_{h'}(X)-E[\mathfrak F_2^+]_{h'}(X)|
 &\le4b_{\rm cs}^{\mathfrak F+}n e^{-\kappa''d_t(X)}
       \le2n e^{-\kappa''d_t(X)}.                 \label{III-32-29}
\end{align}
\end{proposition}

\begin{proof}
In the refined series a term without a new vertex is a term
of the old subsystem. Removing all such terms leaves exactly
the marked part. The origin-class bound in the counting proof
of Theorem~\ref{II-118-20}, with the same majorants as in
\cref{III-30-36}, gives the constant $K_{\rm or}^+$.
For the second bound, coarsening and the first inequality in
\eqref{eq:centered-cluster-rates} give
\[
 e^{-r_{\mathsf w}^+d_k(Y)}e^{r_H^+d_t(Y^\sharp)}
       \le e^{-c_{\mathrm E}d_k(Y)}.
\]
For a coarse cube $\Delta$, at most $(L+2)^3$ fine cubes
can anchor a $Y$ with $\Delta\subset Y^\sharp$. The entropy
bound gives
\[
 \sum_{Y:\Delta\subset Y^\sharp}e^{-c_{\mathrm E}d_k(Y)}
       \le(L+2)^3C_{\mathrm E}\le2L^3C_{\mathrm E}.
\]
There are no marked failed bonds here. Thus
$\eta^{{\rm new},+}\le2C_{\mathrm E}L^3e^{4r_H^+}n_+$.
The marked estimate proves \eqref{III-32-24}. If the new
members vanish locally, every marked Mayer factor is zero.

Write $D=\min(K_{\rm or}^+,K_{\rm Lip}^+)e^{-\kappa''d_t(X)}$.
If $n_{12}=0$, the coarse systems coincide. Otherwise set
$R=n_+/(2n_{12})$. If $R\le1$, the triangle inequality
through the common old subsystem gives
$|E_1-E_2|\le2Dn_+\le4Dn_{12}$.
For $R>1$, interpolate the new members by
$\mathbb W_1+\lambda(\mathbb W_2-\mathbb W_1)$, $|\lambda|<R$.
Admit this family with sizes $s^{o'}/2+Rd^{o'}$ and
$s^{\rm ce}/2+Rd^{\rm ce}$. Their total amplitude is $n_+$;
each ordinary size is at most $n_+/e\le1$, and the centered
size is at most $w_0n_+/e\le w_0$. The frame bound is
therefore unchanged. The cluster sum is holomorphic in
$\lambda$. Its difference from its value at zero has modulus
at most $2Dn_+$ throughout the disk. Schwarz's lemma at
$\lambda=1$ gives $2Dn_+/R=4Dn_{12}$.

For the relative bound, if $n\ge1/2$, the triangle inequality
and \eqref{III-32-15} give
\[
 |E_1-E_2|\le2b_{\rm cs}^{\mathfrak F+}e^{-\kappa''d_t(X)}
          \le4b_{\rm cs}^{\mathfrak F+}n e^{-\kappa''d_t(X)}.
\]
If $0<n<1/2$, the same interpolation is admitted at the
original named sizes for $|\lambda|<(2n)^{-1}$. Apply Schwarz's
lemma with the uniform bound $2b_{\rm cs}^{\mathfrak F+}
e^{-\kappa''d_t(X)}$. This proves \eqref{III-32-29}; the case
$n=0$ was already covered.
\end{proof}


\section{The source perturbation on one lattice}
\label{sec:one-lattice-source}

Inside the small-field region the history expansion agrees
with the expansion on one lattice. We need this agreement
with an extra potential, together with a bound on the change
of the pointed interaction. The Gaussian data are unchanged.

Use the one-lattice Gaussian structure of
Definition~\ref{II-19-15}, with finite threshold
$\mathfrak b$, at $t=k+1$, $0\le k<K$. Assume $L\ge13$ is
odd, $M=L^{m'}$, $1\le m'<m+K-k-1$, and let $R$ be a
nonempty union of coarse cubes. Write $\mathfrak X$ for the
nonpivot bonds. For $Y^\sharp\subset R$ put
\[
 \mathbb D^{\rm f}(Y)=\{B:|B(b)|<C_{\mathcal A}\mathfrak b\},
 \qquad Q_Y(B)=1+M^{-3}\sum_{b\in\mathfrak X^{\circ+}\langle Y\rangle}|B(b)|^2,
 \qquad C_{\mathcal A}\ge3,
\]
and use the norm
$\|V\|^{\rm f}_{\star,r}=\sup_{Y,B}e^{r d_k(Y)}|V(Y;B)|/Q_Y(B)$.
The Gaussian constants and the quantities
$u,\lambda_1,\lambda_V,w_\star,K_\Pi,K_P^\Pi,\mathsf p_\star$
are those of \eqref{II-19-46}. Require
\begin{equation}\label{eq:one-lattice-gaussian-smallness}
 u\le C_{\rm ex}^{-1}\min(\gamma_{\rm r},c_{\rm r}^{-2}),
 \qquad\bar\kappa\ge4\lambda_1+1,\qquad
 \mathsf p_\star\le\tfrac12.
\end{equation}
Thus the Gaussian and numerical hypotheses of
Theorems~\ref{II-19-55}--\ref{II-20-87} hold. All constants in this section
refer to this one set of Gaussian data.

Let $V^{\rm J},V^{\rm G},V^{\rm A},V^{\rm iv}$ be holomorphic
on the field polydiscs. Require $V^{\rm J}(0)=V^{\rm iv}(0)=0$,
and require $V^{\rm G},V^{\rm A}$ to vanish through degree two.
Set
\[
 n_{\rm int}=\|V^{\rm J}\|^{\rm f}_{\star,\lambda_V}
          +\|V^{\rm G}\|^{\rm f}_{\star,\lambda_V}
          +\|V^{\rm A}\|^{\rm f}_{\star,\lambda_V},\qquad
 n_{\rm iv}=\|V^{\rm iv}\|^{\rm f}_{\star,\lambda_V},
\]
and suppose both are at most $w_\star$. The interpolation and
bond insertions are those of Proposition~\ref{II-20-103} at
$\tau=1$. Denote its path by
$\underline V_v=\{V_v^{\rm int},V_v^{\rm iv}\}$ and its
bond insertion by $\dot V_{v,b}$. Define
\begin{equation}\label{eq:source-bond-localization}
 E^\bullet(X,b)=\int_0^1 E'[\underline V_v;\dot V_{v,b}](X)\,dv,
 \qquad \mathcal E(X,b)=E^\chi[\emptyset](X,b)+E^\bullet(X,b).
\end{equation}
Here $E^\chi[\emptyset]$ is the cutoff contribution of the
empty potential system. It depends only on the Gaussian data.

\begin{lemma}\label{III-34-16}
The integrals in \eqref{eq:source-bond-localization} are
defined, and both terms vanish unless $b^\sharp\subset X$.
Moreover
\begin{align}
 E[\{V^{\rm J}+V^{\rm G}+V^{\rm A},V^{\rm iv}\}](X)
       &=\sum_{b\in\mathfrak X^+}\mathcal E(X,b),
                                           \label{eq:source-bond-sum}\\
 |E^\bullet(X,b)|&\le9K_\Pi(n_{\rm int}+n_{\rm iv})
                            e^{-\lambda_{\rm out}d_t(X)},\label{eq:source-bond-size}\\
 |E^\chi[\emptyset](X,b)|&\le K_P^\Pi e^{4\lambda_1}
              e^{-\gamma_2\mathfrak b^2/2}e^{-\lambda_{\rm out}d_t(X)}.
                                           \label{eq:source-bond-cutoff}
\end{align}
These objects are holomorphic in holomorphic families of the
Gaussian data and potentials with the same bounds. They
preserve orthogonal covariance and conjugation.
\end{lemma}

\begin{proof}
\Cref{II-20-103} gives an admitted path with two
origins, bounded velocities, and bond insertions of norm at
most $9(n_{\rm int}+n_{\rm iv})$. The insertion theorem and
pointed-path identity in Section~\ref{sec:cluster-variations} therefore
give a bounded Borel integrand, the estimate
\eqref{eq:source-bond-size}, and
\[
 \sum_b E^\bullet(X,b)=E[\underline V_1](X)-E[\underline V_0](X).
\]
The two zero potentials in $\underline V_0$ can be removed:
the added-potential bound in Theorem~\ref{II-19-55} has right side zero.
The cutoff identity of Theorem~\ref{II-20-87} gives
$E[\emptyset](X)=\sum_bE^\chi[\emptyset](X,b)$, since at
infinite threshold its empty-system cluster sum is bounded
by zero. This proves \eqref{eq:source-bond-sum}.
Theorem~\ref{II-20-87} also gives \eqref{eq:source-bond-cutoff}.
The pointed insertion and cutoff terms vanish unless the
marked bond cube lies in $X$. The same theorems give
holomorphy of the integrands with common majorants, so dominated
integration gives holomorphy of $E^\bullet$.
The interpolation is linear in the four input families and
commutes with rotations and conjugation. The Gaussian change
of variables proves the remaining assertions.
\end{proof}

Let $h_z$ be the tents of Proposition~\ref{II-20-13} and set
\[
 F(X,z)=\sum_{b\in\mathfrak X^+}h_z(b)\mathcal E(X,b).
\]
The tents are nonnegative. Their total over all fine bonds is
$3L^3$; each of the six coarse bonds incident to $z$ has a
different pivot bond and contributes $1/2$ there. Thus
\begin{equation}\label{III-33-20}
 \sum_{b\in\mathfrak X^+}h_z(b)=3L^3-3=:N_L.
\end{equation}
The distinctness follows because a fine bond belongs to at
most one coarse corridor. The support property of the tents
and $m'\ge1$ show that $F(X,z)=0$ for $z\notin X$ and
\begin{equation}\label{III-34-23}
 \sum_{z\in X\cap\mathbb T^{(t)}}F(X,z)
        =E[\{V^{\rm J}+V^{\rm G}+V^{\rm A},V^{\rm iv}\}](X).
\end{equation}

\begin{proposition}\label{III-34-19}
Let $W$ be holomorphic on the same field polydiscs, with
$W(0)=0$ and $n_W=\|W\|^{\rm f}_{\star,\lambda_V}$.
Write $F^+$ for the pointed term with fourth family
$V^{\rm iv}+W$. If $n_{\rm iv}<w_\star$ and
$n_{\rm iv}+n_W\le w_\star$, then
\begin{equation}\label{III-34-20}
 |F^+(X,z)-F(X,z)|
 \le N_L\frac{18K_\Pi(n_{\rm int}+w_\star)}{w_\star-n_{\rm iv}}
                 n_W e^{-\lambda_{\rm out}d_t(X)}.
\end{equation}
If $n_{\rm iv},n_W\le w_\star/2$, one also has
\begin{equation}\label{III-34-24}
 |F^+(X,z)-F(X,z)|
 \le9N_LK_\Pi\left(1+\frac{8(n_{\rm int}+n_{\rm iv})}{w_\star}\right)
                 n_W e^{-\lambda_{\rm out}d_t(X)}.
\end{equation}
\end{proposition}

\begin{proof}
The fourth family $V^{\rm iv}+\lambda W$ is admitted for
$|\lambda|\le R=(w_\star-n_{\rm iv})/n_W$ when $n_W>0$.
For a fixed bond, its $E^\bullet$ term is holomorphic on
$|\lambda|<R$ and bounded by
$A=9K_\Pi(n_{\rm int}+w_\star)e^{-\lambda_{\rm out}d_t(X)}$.
If $R=1$, the triangle inequality bounds its difference at
$1$ and $0$ by $2A$. If $R>1$, Schwarz's lemma applied to
that difference gives $2A/R$. The cutoff term is unchanged.
Multiplication by the tent mass $N_L$ proves
\eqref{III-34-20}. If $n_W=0$, the difference is zero.

For \eqref{III-34-24}, linearity of the interpolation gives
$\underline V_v^+=\underline V_v+_{\rm iv}W_v$ and
$\dot V_{v,b}^+=\dot V_{v,b}+\dot W_{v,b}$, with
$\|W_v\|\le n_W$ and $\|\dot W_{v,b}\|\le9n_W$.
Subtract the two insertion terms. The change in the insertion
is bounded by $9K_\Pi n_W e^{-\lambda_{\rm out}d_t(X)}$.
The change in the system, with its insertion fixed, is bounded
by the insertion Lipschitz estimate of
Theorem~\ref{II-20-54} by
\[
 \frac{8K_\Pi}{w_\star}\,9(n_{\rm int}+n_{\rm iv})n_W
                               e^{-\lambda_{\rm out}d_t(X)}.
\]
The half-size assumptions are exactly those needed for this
estimate. Integrate over $v\in(0,1)$ and sum with the tents.
\end{proof}

\subsection{Identification with the history terms}

An escort on one lattice is a family
$W^{\rm 1l}(Y;P,B)$, holomorphic on the full analytic triple
core and on $\mathbb D^{\rm f}(Y)$, vanishing at $B=0$.
Its norm $\|W^{\rm 1l}\|^{\rm co}_{\star,\lambda_V}$ is the
supremum of the preceding weighted quadratic norm also over
that core. Write this norm as $\theta$.

\begin{proposition}\label{III-34-30}
Suppose $k\ge1$ and the unsourced one-lattice step is defined
on the full analytic core with the Gaussian hypotheses
\eqref{eq:one-lattice-gaussian-smallness}. Suppose its three
internal families and its regular fourth family satisfy
\[
 n_{\rm int}\le3K_P\delta_k\le\tfrac38w_\star,
 \qquad n_{\rm iv}\le K_\Lambda s_{\mathfrak a}N[\mathscr L]
                         \le\tfrac14w_\star.
\]
Here $s_{\mathfrak a}=(1+\beta_{\rm s})\delta_k/\alpha_0$,
and $N[\mathscr L]$ is the weighted size of the old regular list.
If
\begin{equation}\label{III-34-31}
 K_\Lambda s_{\mathfrak a}N[\mathscr L]+\theta\le w_\star,
\end{equation}
then the step with fourth family $V'+W^{\rm 1l}$ is defined
on the same core and
\begin{equation}\label{III-34-32}
 |F^+(X,P,z)-F(X,P,z)|
     \le33(3L^3-3)K_\Pi\theta e^{-\lambda_{\rm out}d_t(X)}.
\end{equation}
It obeys the sum rule \eqref{III-34-23}, is holomorphic in
the triple and source parameters, and inherits gauge
invariance, pair dependence and conjugation from the escort.
For a history region $X\subset\Lambda_t$, the same pointed
term is obtained whenever the unperturbed history potentials
agree there with their one-lattice counterparts and the
history escort restricts to $W^{\rm 1l}$.
\end{proposition}

\begin{proof}
At fixed $P$, restriction of the escort to $(P|_Y,B)$ is
holomorphic on the field polydiscs and has norm at most
$\theta$. The fourth-family room \eqref{III-34-31} therefore
admits the data of \cref{III-34-16}. In
\eqref{III-34-20},
\[
 18\frac{n_{\rm int}+w_\star}{w_\star-n_{\rm iv}}
        \le18\frac{1+3/8}{1-1/4}=33.
\]
This proves the bound. The sum rule and the parameter
properties follow from \cref{III-34-16}. Pair dependence
identifies the Gaussian matrices and all four potential
families, hence every term. For conjugation one uses the
conjugate escort and conjugate triple together.

On $X\subset\Lambda_t$, the interior comparison in
Proposition~\ref{II-135-20} identifies the integration
bonds, matrices, thresholds and ordinary potentials with
those on one lattice. Under the stated escort agreement the
fourth family is also identical. The conditional integrals,
bond insertions and tents therefore give the same pointed
term. This comparison uses the full equality of the old
regular potentials, including their domains; it is not
inferred from an equality of their values at one field.
\end{proof}

\begin{proposition}\label{III-47-46}
At the first step $k=0$, suppose the same Gaussian and
analyticity hypotheses hold with $V'=0$ and
$n_{\rm int}\le3w_\star/8$. Let $W^{\rm 1l}$ be an escort
with $\theta\le w_\star/2$. Then
\begin{equation}\label{III-47-48}
 |F^+(X,P,z)-F(X,P,z)|
 \le36(3L^3-3)K_\Pi\theta e^{-\lambda_{\rm out}d_1(X)}
 =108(L^3-1)K_\Pi\theta e^{-\lambda_{\rm out}d_1(X)}.
\end{equation}
Suppose also that a centered history frame is admitted on
$X\subset\Lambda_1$, its local thresholds are $\mathfrak b_0$,
and its ordinary and centered additional members, evaluated
at zero center, sum to this escort on all the field polydiscs.
Then
\begin{equation}\label{III-47-47}
 E[\mathfrak F^+]_{h'}(X;\emptyset;P;\mathfrak q)
             =\sum_{z\in X\cap\mathbb T^{(1)}}F^+(X,P,z).
\end{equation}
The centered frame is assumed to contain the canonical
first-step system, whose inherited regular and boundary
families are empty. The same identities hold in holomorphic
source families with these bounds.
\end{proposition}

\begin{proof}
Use \eqref{III-34-24} with $n_{\rm iv}=0$:
$9(1+8n_{\rm int}/w_\star)\le36$.
Inside $\Lambda_1$ there is no external field on $X$.
The conditional mean is therefore zero for every decoupling
parameter. Substituting this value in every centered Mayer
factor replaces $V(Y;B-m,m)$ by $V(Y;B,0)$. The resulting
ordinary system has exactly the same terms and activities.
The interior comparison of Proposition~\ref{II-135-20}
identifies its Gaussian system and thresholds with those
on one lattice. Merge the additional members into the
escort, using the finite exponential identity in the proof
of \cref{III-32-12}. The remaining families are J, G and A,
since the inherited ones are empty at $k=0$.
Equation~\eqref{III-34-23} proves \eqref{III-47-47}.
All identifications hold at every point of the specified
domains, so they also hold in the stated parameter families.
\end{proof}


\section{Weighted potentials and their centered expansion}
\label{sec:weighted-source-calculus}

A plaquette source contributes a quadratic form at the conditional
mean and a potential for the centered integration field. The latter
must retain the oscillation of the weight: a bound by its supremum
alone would lose the small factor in the linear term. We obtain both
parts from the localized Wilson action.

\subsection{The localized weight and its Taylor coefficients}

Fix $1\le k<K$, $t=k+1$, a history $h'$, and the canonical shift
system of the step. All its localization domains, restrictions and
radii are those of Section~\ref{sec:history-actual-potentials}. For a weight $w$ on the
bare plaquettes, replace the Wilson member of layer $m$ by
\[
 W_w^{(m)}(X,\mathbf U)=L^k
       \sum_{p\in\mathsf{Pl}(m,X)}w(p)a_p(\mathbf U).
\]
Here $\mathsf{Pl}(m,X)$ consists of the inner bare plaquettes
of $\Omega_0^\diamond$ with base in layer $m$ and smallest
level-$m$ domain $X(p;m)=X$.
Let $\mathsf{Pl}\langle Y\rangle$ be their union over $m\le k$
and $X^\uparrow\subset Y$, and write
$|w|_Y=\sup_{\mathsf{Pl}\langle Y\rangle}|w|$.

The localized potential is the finite difference
\begin{equation}\label{eq:weighted-localization}
 \begin{split}
 \mathbb V^w(Y;P,\mathcal A)
 =\sum_{m=0}^k\sum_{X^\uparrow\subset Y}
   \sum_{D\subset Y\setminus X^\uparrow}
  (-1)^{m(X,Y)-\#D}
  \bigl[W_w^{(m)}(X,\mathbf U_{D\cup X^\uparrow})
                        -W_w^{(m)}(X,\mathbf U)\bigr],
 \end{split}
\end{equation}
where $D$ is a set of level-$k$ cubes,
$m(X,Y)=\#_{\pi_k}(Y\setminus X^\uparrow)$,
$\mathbf U=e^{i\eta_tA'}U_{\rm r}$, and
$\mathbf U_D=e^{i\eta_k\mathbf A_Y(1_D;P,\mathcal A)}\mathbf U$.
Every term is read only on $X$. The domain of
\eqref{eq:weighted-localization} is the domain of the localized
shift $(P,\mathcal A)$, including when the weight or the sum is
zero. Stored fields do not enter it.

For positive radii $a_m\le1$, let $\mathbb U_k^a(Y|h')$ consist of
triples for which the shift is defined for every
$\sup|\mathcal A|<\varepsilon_{\mathfrak h}$ and
\[
 |\partial\mathbf U(p)-1|<a_mL^{-2m}
\]
at every inner bare plaquette of layer $m$ whose bonds lie in $Y$.
This specifies the domains used below. Set $c_E=\log18432$ and put
\[
 \delta_* =M\min(\mu_{\mathfrak h},1),\qquad
 r=(1-2\delta)(\kappa_*-2),\qquad
 r_0=\min\{\varepsilon_{\mathfrak h},(3c_{\mathfrak h})^{-1}\},
\]
and assume
\begin{equation}\label{eq:weighted-localization-rates}
 0<\delta<\tfrac12,\quad \kappa_*>0,\quad
 2\delta\kappa_*\ge2(c_E+1),\quad
 \bar\kappa-1\ge\kappa_*,\quad \delta_*L\ge20\log L.
\end{equation}
Let $\varpi_i^c$ and $\varpi_i^w$ be the coefficients of
equation \eqref{II-91-22}, with $\nu=1$, $a=2\delta\kappa_*$,
and distance rates $\delta_*/2$ and $\delta_*$, respectively.
Equivalently, they are the coefficients $\varpi_i$ in the
placed-domain estimate of Lemma~\ref{II-91-52} at those rates.
Set
\[\begin{gathered}
 K_c=12(1+2c_{\rm pl})^2C_h(L,\delta,1),\quad c_{\rm pl}=600,
 \qquad K_0=K_c C'_h(L,\delta,1)M^3,\\
 \qquad a^\Sigma=\sum_{m=0}^k\varpi^w_{k-m}L^{k-m}a_m.
\end{gathered}\]
The constants $C_h,C'_h$ are those of the same counting lemma;
in particular $C'_h(L,\delta,1)=12L^4$.

\begin{proposition}\label{III-39-18}\label{III-39-24}
For $P\in\mathbb U_k^{\mathbf1}(Y|h')$, the function
$\mathcal A\mapsto\mathbb V^w(Y;P,\mathcal A)$ is holomorphic
on $\sup|\mathcal A|<r_0$, vanishes at zero, and obeys
\begin{equation}\label{III-39-19}
 |\mathbb V^w(Y;P,\mathcal A)|\le K_0|w|_Y e^{-rd_k(Y)}.
\end{equation}
Define its homogeneous Taylor parts and symmetric quadratic form by
\begin{align}
 T_n[w](Y;P;\mathcal A)
   &=\frac1{n!}\left.\partial_z^n
                \mathbb V^w(Y;P,z\mathcal A)\right|_{z=0},
                                                 \label{III-39-21}\\
 T_2[w](\mathcal A_1,\mathcal A_2)
   &=\tfrac12\{T_2[w](\mathcal A_1+\mathcal A_2)
                       -T_2[w](\mathcal A_1)-T_2[w](\mathcal A_2)\}.
                                                 \label{III-39-22}
\end{align}
Then $T_0=0$, each $T_n$ is a homogeneous polynomial defined for
all fields, and
\begin{align}
 |T_n[w](\mathcal A)|
   &\le K_0|w|_Y r_0^{-n}\sup|\mathcal A|^n e^{-rd_k(Y)},
                                                   \label{III-39-25}\\
 |T_{\ge N}[w](\mathcal A)|
   &\le2K_0|w|_Y r_0^{-N}\sup|\mathcal A|^N e^{-rd_k(Y)}
                  \quad(\sup|\mathcal A|\le r_0/2),\label{III-39-26}\\
 |2T_2[w](\mathcal A_1,\mathcal A_2)|
   &\le4K_0|w|_Yr_0^{-2}
             \sup|\mathcal A_1|\sup|\mathcal A_2|e^{-rd_k(Y)}.
                                                   \label{III-39-27}
\end{align}
Here $T_{\ge N}=\mathbb V^w-\sum_{n<N}T_n[w]$ on the radius-$r_0$
ball. If $P\in\mathbb U_k^a(Y|h')$, there is also the sharper
linear bound
\begin{equation}\label{III-39-28}
 |T_1[w](\mathcal A)|\le eK_cc_{\mathfrak h}M^3
                   |w|_Y a^\Sigma\sup|\mathcal A|e^{-rd_k(Y)}.
\end{equation}
All these functions are linear and local in the weight. They
vanish at domains containing no source cube, where a source cube
means a fine cube meeting $\Omega_t$. They depend on the pair
of the triple and inherit its gauge, conjugation and Borel
properties. At every fixed real factor their coefficients admit
local holomorphic extensions in the remaining triple variables.
\end{proposition}

\begin{proof}
The finite sum \eqref{eq:weighted-localization} proves linearity,
locality and independence of the stored datum. A plaquette used
in a term lies in $X\subset X^\uparrow\subset Y$.
The shift vanishes at $\mathcal A=0$ and at domains without a
source cube, so each difference then vanishes. For a fixed real
factor, take local holomorphic extensions of the finitely many
shift values at $1_D$ and intersect their neighborhoods.
Substituting them in \eqref{eq:weighted-localization} gives a
holomorphic extension of the potential. Pair dependence, gauge
covariance, conjugation and measurability pass through this same
finite formula.

For completeness, the weighted estimate uses the following two
bounds for a single Wilson member at points where
$\vartheta^a(\sup|\mathcal A|)<1$:
\[
 \begin{split}
 |\mathbb V[W_w^{(m)}(X)](Y)|
 &\le12(1+2c_{\rm pl})^2\vartheta^a(\sup|\mathcal A|)
 M^3s_m(w;Y)L^{k-m}a_m^2
 e^{-\delta_*\mathsf d_k(X)/2-(\bar\kappa-1)m(X,Y)},\\
 |\mathbb V[W_w^{(m)}(X)](Y)|
 &\le12e(1+2c_{\rm pl})^2c_{\mathfrak h}\sup|\mathcal A|
 M^3s_m(w;Y)L^{k-m}a_m
 e^{-\delta_*\mathsf d_k(X)-(\bar\kappa-1)m(X,Y)}.
 \end{split}
\]
Here $s_m(w;Y)$ is the supremum of the weight on the plaquettes
of layer $m$ used in $Y$, and $\vartheta^a$ is the radius ratio
of equation \eqref{II-3-11}. To keep this supremum local,
replace $w$ in each Wilson member by its restriction to that
member's plaquettes; the localized member is unchanged.
Here is the weighted calculation. Put
$\rho=e c_{\mathfrak h}\sup|\mathcal A|
 e^{-\delta_*\mathsf d_k(X)}$. The plaquette product estimate of
equation \eqref{II-41-47}, applied with the layer radius,
bounds the interpolated Wilson member for $|\lambda|\rho\le a_m$
by
\[
 b=6(1+2c_{\rm pl})^2M^3s_m(w;Y)L^{k-m}a_m^2.
\]
Its constant value at $\lambda=0$ cancels in the finite difference.
The rooted Cauchy estimate therefore bounds that difference by
$2bR^{-1}e^{-(\bar\kappa-1)m(X,Y)}$ for every
$1<R\le a_m/\rho$. If $\rho=0$, let $R\to\infty$.
Otherwise take $R=a_m/\rho$. The distance from an occupied
layer-$m$ base to the source gives
\[
 \rho/a_m\le(e/3)\vartheta^a(\sup|\mathcal A|)
                         e^{-\delta_*\mathsf d_k(X)/2}<1.
\]
Using this estimate yields the first displayed bound; substituting
$\rho$ without weakening its exponential yields the second.
The argument is the weighted version of the plaquette and rooted
Cauchy calculations in the proofs of Propositions~\ref{II-106-39} and~\ref{II-107-44}.
The placed-domain estimate of Lemma~\ref{II-91-52} applies:
occupied Wilson bases have $d_m(X)=0$, are placed with $\nu=1$,
and the smaller distance rate $\delta_*/2$ satisfies its
hypotheses by \eqref{eq:weighted-localization-rates}. Summing gives
\[
 \begin{split}
 |\mathbb V^w(Y;P,\mathcal A)|
 &\le K_c\vartheta^a(\sup|\mathcal A|)M^3|w|_Y
       \sum_{m=0}^k\varpi^c_{k-m}L^{k-m}a_m^2e^{-rd_k(Y)},\\
 |\mathbb V^w(Y;P,\mathcal A)|
 &\le eK_cc_{\mathfrak h}\sup|\mathcal A|M^3|w|_Y
                                  a^\Sigma e^{-rd_k(Y)}.
 \end{split}
\]
For $a_m=1$, the radius ratio is $3c_{\mathfrak h}\sup|\mathcal A|$.
It is less than one on the radius-$r_0$ ball. The sum in the first
bound is at most $C'_h(L,\delta,1)$, proving \eqref{III-39-19}.

Apply the one-variable Cauchy formula to
$z\mapsto\mathbb V^w(Y;P,z\mathcal A)$ on the disc
$|z|\sup|\mathcal A|<r_0$. This proves \eqref{III-39-25}; summing
the geometric tail proves \eqref{III-39-26}. Polarization at
normalized fields, using $4T_2(B,B')=T_2(B+B')-T_2(B-B')$, proves
\eqref{III-39-27}. For a fixed family $a$, the condition
$\vartheta^a(|z|\sup|\mathcal A|)<1$ holds for sufficiently small
$|z|$. Divide the second potential bound at $z\mathcal A$ by
$|z|$ and let $z\to0$ to prove \eqref{III-39-28}.
The Cauchy formula on a common small circle gives local
holomorphic extensions of the Taylor coefficients. Homogeneity
extends those coefficients to arbitrary fields. Taylor limits at
real scalar multiples and polarization preserve measurability
and the two symmetries.
\end{proof}

\subsection{Three source terms and the oscillation bound}

We now apply the preceding proposition at a stronger rate
$r^\sharp=(1-2\delta)(\kappa^\sharp-2)$. Assume
\eqref{eq:weighted-localization-rates} for these parameters, and
that every middle-domain triple lies in $\mathbb U_k^a(Y|h')$
with the enlarged layer radii $a_m\le1$. Let
\[
 g_k^{-1}\mathcal A=z\oplus a,\qquad \phi=\mu\oplus a,
 \qquad |z|,|a|<C_{\mathcal A}\mathfrak b_k,\qquad
 |\mu|<\theta_{\rm sl}\mathfrak b_k,
\]
where inequalities are bondwise and the center $\mu$ is on the
integration bonds. Require
\begin{equation}\label{eq:weighted-field-room}
 (C_{\mathcal A}+\theta_{\rm sl})\delta_k\le r_0/2,
 \qquad \delta_k=g_k\mathfrak b_k.
\end{equation}
Then both $g_k(z+\phi)$ and $g_k\phi$ lie in the Taylor ball.
For the canonical construction take $p_b=p_{12}(z_\square)$, where $z_\square$
is the centre of the cube of $\pi_k$ containing $b_-$.
Its four corners lie in that cube: their distance from $z_\square$ is
at most the bare spacing $\eta_k\le1\le M/2$.
Set $v_b=v-v(p_b)\mathbf1$. The bounds below also hold for any
choice of $p_b$ whose four corners lie in the same cube.
The canonical choice obeys
\begin{equation}\label{eq:source-reference-translation}
 p_{b+\mathsf v}=p_b+\mathsf v\qquad(\mathsf v\in M\mathbb Z^3),
\end{equation}
because translation takes the containing cube with centre $z_\square$
to the containing cube with centre $z_\square+\mathsf v$.
With all Taylor parts taken at $(Y;P)$, define
\begin{align}
 \mathbb W[v]
  &=-\{\mathbb V^{x_kv}(g_k(z+\phi))
                      -\mathbb V^{x_kv}(g_k\phi)\}
    +\sum_{b\in\mathfrak X^{\circ+}_{\rm I}\langle Y\rangle}
       x_kv(p_b)\{T_1[\mathbf1](g_kz_b)
                     +2T_2[\mathbf1](g_kz_b,g_k\phi)\},
                                                 \label{III-39-34}\\
 \mathbb W^1[v]&=-\sum_{b\in\mathfrak X^{\circ+}\langle Y\rangle}
                                      T_1[x_kv_b](g_k\phi_b),\notag\\
 \mathbb W^3[v]&=-T_{\ge3}[x_kv](g_k\phi).
                                                 \label{III-39-35}
\end{align}
The canonical centered member is their sum
$\mathbb W^{\rm S}[v]=\mathbb W[v]+\mathbb W^1[v]+\mathbb W^3[v]$.
Its domain consists of the pairs with
\[
 P\in\mathbb U_k^{\mathbf1}(Y|h'),\qquad
 \sup|g_k(z+\phi)|<r_0,\qquad \sup|g_k\phi|<r_0,
\]
where $g_k^{-1}\mathcal A=z\oplus a$ and $\phi=\mu\oplus a$.
The norm estimates below use the stated radial-field and seat
restrictions within this domain. The three terms are analytic
on the stated products of domains,
linear in $v$, local in the plaquettes of $Y$, and independent of
the stored fields. They inherit the pair, gauge, conjugation and
Borel properties. All vanish without a source cube. The first
vanishes at $z=0$; the other two are independent of $z$ and vanish
at $\phi=0$.

Put $q_Y(z)=M^{-3}\sum_{b\in\mathfrak X^{\circ+}_{\rm I}\langle Y\rangle}
|z(b)|^2$. For such families define
\[
 \|F\|_{{\rm I}\star,r}^{\rm rad}
 =\sup_{Y,(P,\mathcal A,\mathfrak q),\mu}
       \frac{e^{rd_k(Y)}|F(Y;P,\mathcal A,\mathfrak q;\mu)|}
                         {1+q_Y(z)},
\]
where the points and centers range over the stated tiers.
Let $D$ be the least $d_k(Y)$ over $Y\not\subset\Omega_k$
containing an integration bond, and let $D'$ be the corresponding
minimum for any bond of $\mathfrak X$; empty minima are $+\infty$.
For either empty minimum we set $e^{-c(+\infty)}=0$ for every
$c\ge0$, including $c=0$.
Let $\Theta=\max(C_{\mathcal A},\theta_{\rm sl})$, and set
\[
 C^{\rm src}=64\max\{3eK_cc_{\mathfrak h}M^6,
   12K_0r_0^{-2}\Theta M^3,6K_0r_0^{-2}M^6,
   4K_0r_0^{-3}(C_{\mathcal A}+\theta_{\rm sl})^3\}.
\]

\begin{proposition}\label{III-39-43}\label{III-40-9}\label{III-40-22}
Suppose $|v|_\infty\le\sigma$ and
$x_k|v(p)-v(p')|\le\rho'$ for inner plaquettes of $\Omega_k$
whose barycentres are at distance at most $M$, in level-$k$ units.
For $0\le r\le r^\sharp-2$, the three norms are bounded by
\begin{align}
 N_W(r)&=C^{\rm src}\{\sigma(1+g_k\mathfrak b_k^3)
       +\rho'(g_ka^\Sigma+g_k\delta_k)
       +\sigma(g_k^{-1}a^\Sigma+\mathfrak b_k)
                               e^{-(r^\sharp-r-2)D}\},
                                                    \label{III-39-48}\\
 N_1(r)&=384eK_cc_{\mathfrak h}M^6\Theta
       \{\rho'a^\Sigma\delta_k
          +\sigma g_k^{-1}a^\Sigma\mathfrak b_k
                               e^{-(r^\sharp-r-2)D'}\},
                                                    \label{III-40-11}\\
 N_3(r)&=2K_0r_0^{-3}\Theta^3\sigma g_k\mathfrak b_k^3.
                                                    \label{III-40-12}
\end{align}
Consequently their sum has norm at most
$\mathsf{sz}^{\rm src}=N_W(r_V)+N_1(r_V)+N_3(r_V)$ at any
$0\le r_V\le r^\sharp-2$. If an unsourced frame has coarse sizes
$s_{o'}$ and
\begin{equation}\label{III-40-23}
 \mathsf{sz}^{\rm src}\le w_0,\qquad
 eL^3\left(\sum_{o'}s_{o'}+w_0^{-1}\mathsf{sz}^{\rm src}\right)
                                      \le(4K^+_{\rm ap})^{-1},
\end{equation}
adding this centered member gives a frame to which
\cref{III-32-12} applies. The statement is uniform over affine
holomorphic families of weights obeying the same two bounds.
\end{proposition}

\begin{proof}
Taylor expansion and $x_kg_k^2=1$ give
\begin{equation}\label{III-39-44}
 \begin{split}
 \mathbb W[v]={}&-\sum_bT_1[x_kv_b](g_kz_b)
       -2\sum_bT_2[v_b](z_b,\phi)-T_2[v](z)\\
       &-\{T_{\ge3}[x_kv](g_k(z+\phi))
                              -T_{\ge3}[x_kv](g_k\phi)\}.
 \end{split}
\end{equation}
The first two sums are over integration bonds. Let $n_Y=\#_{\pi_k}Y$.
The entropy geometry and Cauchy--Schwarz give
\[
 n_Y\le8e^{d_k(Y)},\qquad
 \sum_b|z(b)|\le\tfrac32M^3n_Y(1+q_Y(z)),\qquad
 (\sum_b|z(b)|)^2\le3M^6n_Y(1+q_Y(z)).
\]
Let $\omega_Y=\sup_{p,b}|v(p)-v(p_b)|$, with $p$ among the
plaquettes used by the potential and $b$ among the bonds in $Y$.
Using \eqref{III-39-25}--\eqref{III-39-28}, the four terms in
\eqref{III-39-44}, after division by
$(1+q_Y(z))e^{-r^\sharp d_k(Y)}$, are bounded respectively by
\[\begin{gathered}
 \tfrac32eK_cc_{\mathfrak h}M^6n_Ya^\Sigma x_kg_k\omega_Y,
 \quad6K_0r_0^{-2}\Theta M^3n_Y\mathfrak b_k\omega_Y,\\
 \quad6K_0r_0^{-2}M^6n_Y|v|_Y,
 \quad4K_0r_0^{-3}(C_{\mathcal A}+\theta_{\rm sl})^3
                                  g_k\mathfrak b_k^3|v|_Y.
\end{gathered}\]
The two additional terms satisfy
\[
 |\mathbb W^1[v]|\le3eK_cc_{\mathfrak h}M^6\Theta
             n_Yx_k\omega_Ya^\Sigma\delta_ke^{-r^\sharp d_k(Y)},
 \quad
 |\mathbb W^3[v]|\le2K_0r_0^{-3}\Theta^3
                    |v|_Yg_k\mathfrak b_k^3e^{-r^\sharp d_k(Y)}.
\]
For $Y\subset\Omega_k$, choose $p,b$ in the definition of
$\omega_Y$, and let $Q,Q'$ be fine cubes containing a corner
of $p$ and the reference plaquette $p_b$, respectively.
Join $Q$ to $Q'$ by a nonrepeating adjacent-cube path in $Y$,
with at most $n_Y-1$ edges, and use the plaquette at each cube
centre. Consecutive barycentres, including the two endpoint
links to $p$ and $p_b$, are at distance at most $M$. Hence
$x_k\omega_Y\le(n_Y+1)\rho'\le2n_Y\rho'$; if either index set
is empty, $\omega_Y=0$.
For $Y\not\subset\Omega_k$, use $\omega_Y\le2\sigma$ and the
lower bounds $d_k(Y)\ge D$ or $D'$. Since
$n_Y^2\le64e^{2d_k(Y)}$, two units of decay absorb these counts.
Substitution gives \eqref{III-39-48}--\eqref{III-40-12}.
Subadditivity gives the bound for the sum. The two inequalities
in \eqref{III-40-23} are precisely the size and amplitude
conditions of the centered expansion. Finally each local formula
is a finite linear combination of the coordinates of the weight
with locally holomorphic coefficients. This proves joint
holomorphy and keeps the same norm bounds over the parameter set.
\end{proof}

\subsection{The regional Hessian}

The identification of the quadratic source requires the Hessian
of the constrained regional action. We use the reduced determining
set of the shorter sequence, including its outward rows. The
row pairing is
\[
 \langle\omega,\nu\rangle_{\mathfrak R}
       =\sum_c\xi_c^{-1}\langle\omega(c),\nu(c)\rangle,
 \qquad
 (\omega,\nu)_{\mathbf B}
       =\sum_c\xi_c^3\langle\omega(c),\nu(c)\rangle.
\]
Write $\mathsf m\nu(c)=\xi_c\nu(c)$ and let $\mathsf J$
extend a level-$k$ field on $\mathfrak X(\Theta)$ by zero to
the shorter rows. Thus $\mathsf m\mathsf J=\mathsf J$.
All response operators below use this shorter sequence and
these pairings. The following argument is also the regional
second-variation calculation in
Proposition~\ref{II-113-38}.

\begin{proposition}\label{II-27-28}
Consider a separated regional step $k+1$, with $R_1\ge L$,
$R_1\mathfrak M\ge L^2$, and, when $k\ge1$, at least
$6R_1\mathfrak M$ sites in each direction of the level-$k$
torus. Let $p$ be a real point of the coarse fibre, admissible
on every averaging double block, and axial on $\Theta'$.
Suppose the regular critical orbit for its shorter datum
$\mathsf d(p)$ exists, and take an element $U$ of this orbit.
Assume that the pair $(U,J^\circ(U))$ is in the natural class
of the regional fluctuation construction, that all its
response and fluctuation operators are defined, and that
its real row map and node gauge maps are analytic near zero.
The last condition holds at the regularity radius used in
Propositions~\ref{I-88-127} and~\ref{II-113-38}.
Assume also that, along every differentiable curve $V(s)$
in the coarse fibre through $p$, the shorter regular critical
orbit exists for small $|s|$, and
\begin{equation}\label{eq:regional-stationarity}
 \left.\frac d{ds}\right|_{s=0}
    A_{\Omega_0}(U_{\mathbf B^\flat}(\mathsf d(V(s))))=0.
\end{equation}

Then every real $B\in\ker Q_p$ is the velocity of a real
analytic fibre curve through $p$. Along every twice
differentiable such curve with velocity $B$, the first
variation of the regional action vanishes and
\begin{equation}\label{II-27-30}
 \left.\frac {d^2}{ds^2}\right|_{s=0}
 L^k A_{\Omega_0}(U_{\mathbf B^\flat}(\mathsf d(V(s))))
       =\langle B,\Delta_k(U,J^\circ(U))B\rangle.
\end{equation}
The form on the right is real and symmetric and is independent
of the choice of $U$ in that orbit among the admitted pairs.
In free coordinates the full fluctuation operator is
\begin{equation}\label{II-27-32}
 \mathsf A(U,J^\circ(U))
    =\mathsf C_p^{\mathsf T}\Delta_k(U,J^\circ(U))\mathsf C_p
                         +\widehat{\mathsf G}_p .
\end{equation}
Here $\widehat{\mathsf G}_p$ is the quadratic block-contour
form of the same regional step.
\end{proposition}

\begin{proof}
Put $J=J^\circ(U)$ and
$F(A)=L^kA_{\Omega_0}(e^{i\xi A}U)$, with $\xi=L^{-k}$.
Let $\mathcal Q(A)$ be the straightened shorter row map and
$K_A=\mathcal Q'(A)$. At zero its first two derivatives are
\[
 K_0=\mathsf m Q,
 \qquad \mathcal Q''(0)[A_1,A_2]
                =2\mathsf m C^{(2)}(A_1,A_2).
\]
The adjoint for the dimensionless row pairing is
$K_A^\dagger=K_A^{\mathsf T}\mathsf m^{-4}$.
In particular $K_0^\dagger=Q^{\mathsf T}\mathsf m^{-3}$.

Criticality of $U$ and surjectivity of $Q$ give a unique real
multiplier $\nu_0$ with $J=Q^{\mathsf T}\nu_0$.
Indeed a tangent vector in the kernel of $K_0$ can be integrated
to a local curve with fixed straightened rows. Restoring the
node gauges gives a curve with fixed actual datum. The action
is unchanged by this restoration, so its derivative vanishes
on $\ker K_0$. The real annihilator of this kernel is the range
of $K_0^\dagger$. Since $Q\mathsf H=Q\mathsf H_1=1$,
\[
 \nu_0=\mathsf H^{\mathsf T}J=\mathsf H_1^{\mathsf T}J,
 \qquad \omega_0=\mathsf m^3\nu_0.
\]
These are the multiplier identities of
Proposition~\ref{I-88-127} with the row-spacing factors
written out.

We recall why the response linearization uses the constrained
Hessian. Set $\Delta_1=\Delta(U)-\Delta^{(2)}$, where
\[
 (A_1,\Delta^{(2)}A_2)_\xi
                =2(C^{(2)}(A_1,A_2),\nu_0)_{\mathbf B}.
\]
Differentiating gauge invariance and row covariance at $U$
gives, for $\lambda\in\ker Q'$,
\[
 (A,\Delta(U)D\lambda)_\xi
       =\langle\omega_0,\mathcal Q''(0)[A,D\lambda]
                                                   \rangle_{\mathfrak R}
       =(A,\Delta^{(2)}D\lambda)_\xi .
\]
Thus $\Delta_1$ annihilates these gauge directions.
The projection $\mathsf S=1-D\mathsf G'\mathsf R D^{\mathsf T}$
differs from the identity by one of these directions, so
$\Delta_1\mathsf S=\Delta_1=\mathsf S^{\mathsf T}\Delta_1$.
The definitions of the corrected inverse therefore give
\[
 \mathsf G_1^{-1}
       =\Delta_1+D\mathsf RD^{\mathsf T}+Q^{\mathsf T}\mathsf aQ.
\]
Consequently the saddle system
\[
 (\Delta_1+D\mathsf RD^{\mathsf T})x-Q^{\mathsf T}\mu=f,
 \qquad Qx=g
\]
has exactly one solution. Explicitly,
\[
 \mu'=(Q\mathsf G_1Q^{\mathsf T})^{-1}
                           (g-Q\mathsf G_1f),\quad
 x=\mathsf G_1f+\mathsf G_1Q^{\mathsf T}\mu',\quad
 \mu=\mu'-\mathsf ag.
\]
The inverses here are among the assumed defined operators.

For $k\ge1$, apply the real analytic implicit function theorem
to
\[
 \Psi(A,\omega;\beta)=
 \bigl(F'(A)+D\mathsf RD^{\mathsf T}A-K_A^\dagger\omega,
                         \mathcal Q(A)-\mathsf J\beta\bigr).
\]
At $(0,\omega_0;0)$ it vanishes, and its derivative in
$(A,\omega)$ is
\[
 (x,\varpi)\longmapsto
 \bigl((\Delta_1+D\mathsf RD^{\mathsf T})x
               -Q^{\mathsf T}\mathsf m^{-3}\varpi,
                                   \mathsf mQx\bigr).
\]
The saddle calculation proves invertibility. We obtain real
analytic $\widehat A(\beta),\widehat\omega(\beta)$ near zero.

We next remove the auxiliary gauge term and restore the
actual rows; outward rows require this step. Let
$\Sigma_A\lambda$ be the derivative of the gauge action
$A\mapsto A^{e^{it\lambda}}$, and let $\Pi_A\lambda$ be the
corresponding derivatives of the iterated node gauges at the
tops. At zero,
\[
 \Sigma_0\lambda=-D\lambda,\qquad \Pi_0\lambda=Q'\lambda.
\]
The map
$\lambda\mapsto(\Pi_A\lambda,\mathsf RD^{\mathsf T}
\Sigma_A\lambda)$ is bijective near zero. At zero its kernel
is trivial by the scalar inverse on $\ker Q'$. To obtain an
arbitrary target $(\tau,f)$, first choose
$Q'\lambda_1=\tau$, then put
\[
 g=-f-\mathsf RD^{\mathsf T}D\lambda_1,
 \qquad\lambda_2=\mathsf G'\mathsf Rg.
\]
The scalar inverse identities give
$Q'\lambda_2=0$ and
$\mathsf RD^{\mathsf T}D\lambda_2=g$.
Thus $\lambda_1+\lambda_2$ has the required image.
Continuity of the determinant proves the nearby assertion.

The derivative of a row is determined by the derivatives
of its endpoint rotations. At an outward endpoint $z$ of a
level-$q$ row, the extra node factor has the form
\[
 \mathsf v_{q-1}[A](z)
     =\mathsf f_z\bigl(\mathcal Q(A)|_{B(z)}\bigr),
                         \qquad \mathsf f_z(0)=1,
\]
where every bond of $B(z)$ is a level-$(q-1)$ row.
The map $\mathsf f_z$ is holomorphic: it is the exponential
of the average of the logarithms of the block contours.
For two constrained curves with equal rotation derivatives
at the tops, equality of row derivatives at level zero is
immediate. Equality at level $q-1$ gives equal derivatives
of $\mathsf f_z$, hence equal rotations at the outward
endpoints of level $q$. Differentiating the row logarithm
then gives equal row derivatives at level $q$. Induction
proves that $\Pi_A\lambda=0$ implies
$K_A\Sigma_A\lambda=0$.

For $A=\widehat A(\beta)$, put $f=\mathsf RD^{\mathsf T}A$
and choose $\lambda$ whose image under the preceding
bijection is $(0,f)$. Pairing the first equation
$\Psi=0$ with $\Sigma_A\lambda$ gives $(f,f)_\xi=0$:
the action term vanishes by gauge invariance and the
multiplier term by $K_A\Sigma_A\lambda=0$.
Since $f$ is real, $f=0$. Hence
\begin{equation}\label{eq:regional-response-system}
 \mathcal Q(\widehat A)=\mathsf J\beta,\qquad
 F'(\widehat A)=K_{\widehat A}^\dagger\widehat\omega,
 \qquad \mathsf RD^{\mathsf T}\widehat A=0.
\end{equation}

Define $\widehat h$ to be $v_q[\widehat A]^{-1}$ at the
tops of level $q$ and the identity at other bare sites.
At an outward endpoint the factor $\mathsf f_z$ equals one,
since the lower rows of $\mathsf J\beta$ are zero.
Thus $\widehat h\,v_q=1$ at both endpoints of every row,
and
\[
 U_\beta=(e^{i\xi\widehat A(\beta)}U)^{\widehat h(\beta)}
\]
has exactly the datum $\mathsf d(e^{i\beta}p)$.
It is regular after shrinking the neighbourhood, by the
strict regularity inequalities at $U$.
For a curve through $U_\beta$ with this datum fixed, undo
$\widehat h$ and take its logarithm, giving a curve $A_s$.
Its endpoint rotations have the same top derivatives as
some gauge curve, by surjectivity of $\Pi_{A_0}$.
The row induction above then gives
$K_{A_0}\dot A=K_{A_0}\Sigma_{A_0}\lambda$.
Equation~\eqref{eq:regional-response-system} and gauge
invariance show
\[
 (F'(A_0),\dot A)_\xi
   =\langle\widehat\omega,K_{A_0}\dot A\rangle_{\mathfrak R}
   =(F'(A_0),\Sigma_{A_0}\lambda)_\xi=0.
\]
Therefore $U_\beta$ is critical, and belongs to the regular
critical orbit whenever that orbit exists.

Put $\varphi(\beta)=F(\widehat A(\beta))$.
Differentiating~\eqref{eq:regional-response-system} at zero
and using the saddle inverse gives
\[
 \widehat A'(0)=\mathsf H_1\mathsf J,
 \qquad \mathsf m^{-3}\widehat\omega'(0)
     =[(Q\mathsf G_1Q^{\mathsf T})^{-1}-\mathsf a]\mathsf J.
\]
The envelope identity is
$\varphi'(\beta)\nu
 =\langle\widehat\omega(\beta),\mathsf J\nu\rangle_{\mathfrak R}$.
Since $\mathsf m\mathsf J=\mathsf J$ and
$\mathsf J^{\mathsf T}\mathsf a\mathsf J=1$, it follows that
\begin{equation}\label{eq:regional-action-Taylor}
 \begin{split}
 \varphi'(0)\nu&=(\mathsf H_1\mathsf J\nu,J)_\xi,\\
 \varphi''(0)[\nu,\nu]
   &=(\nu,[\mathsf J^{\mathsf T}
             (Q\mathsf G_1Q^{\mathsf T})^{-1}\mathsf J-1]\nu).
 \end{split}
\end{equation}
For $k=0$ the same identities follow directly from
$\widehat A(\beta)=\mathsf J\beta$, $\widehat h=1$,
$Q=\mathsf H_1=1$, $\mathsf R=0$, and
$\mathsf G_1^{-1}=\Delta(U)+1$.

It remains to restrict this unconstrained Taylor formula
to the curved coarse fibre. Its analytic chart is
\[
 \beta_B(s)=sB-\mathsf h_p\widetilde{\mathsf D}_p(sB),
 \qquad B\in\ker Q_p,
\]
with $\widetilde{\mathsf D}_p(0)=
\widetilde{\mathsf D}'_p(0)=0$.
Thus every such $B$ is a fibre velocity. By the assumed
nearby existence and~\eqref{eq:regional-stationarity},
$\varphi'(0)$ annihilates $\ker Q_p$.
For any twice differentiable fibre curve write
$V(s)=e^{i\beta(s)}p$ using the local real logarithm.
The constraint has derivatives $LQ_p$ and $2C_p^{(2)}$,
so differentiating twice gives
\[
 \ddot\beta(0)=\kappa-2\mathsf D_p^{(2)}(B,B),
 \qquad \kappa\in\ker Q_p.
\]
The chain rule and~\eqref{eq:regional-action-Taylor} now give
\[
 \begin{split}
 \frac {d^2}{ds^2}\varphi(\beta(s))\big|_{s=0}
   &=(B,[\mathsf J^{\mathsf T}
       (Q\mathsf G_1Q^{\mathsf T})^{-1}\mathsf J-1]B)\\
   &\quad-2(\mathsf J\mathsf D_p^{(2)}(B,B),
                                      \mathsf H_1^{\mathsf T}J)_{\mathbf B}\\
   &=\langle B,\Delta_k(U,J)B\rangle.
 \end{split}
\]
This is precisely the regional fluctuation definition,
including its curved-constraint correction.
The function $\varphi(\beta(s))$ is the orbit action by
nearby existence and the preceding criticality argument,
which proves~\eqref{II-27-30}.

Finally the axial condition gives
\[
 G_k(\Theta';e^{i\beta}p)
       =\tfrac12\mathsf G_p^{(2)}(\beta,\beta)+O(|\beta|^3).
\]
The first derivative of this function is zero.
Adding its Hessian and pulling back by pivot elimination
proves~\eqref{II-27-32}. Reality and symmetry follow by
polarization of the real second derivative. The action and
its constrained Hessian depend on the orbit, so their
representing operators do not depend on the chosen admitted
representative.
\end{proof}

\subsection{Why the Wilson mixed term vanishes}

The conditional mean is defined by the full fluctuation form,
which includes gauge fixing. To use the source split we must
separate its two parts. This requires the gauge degeneracy of
the Wilson Hessian, as well as the conditional-mean equation.

Let $p$ be a real admissible configuration of the regional fibre,
with invertible pivot coefficients. Let $\Theta'$ denote the
blocks on which the gauge term is imposed and
$\Theta''\subset\Theta'$ the integration blocks. The integration
nonpivot bonds are precisely those meeting $B(\Theta'')$.
Write $C_p$ for pivot elimination, $D_p\rho$ for the gauge
direction, and $\check D_p\rho$ for its nonpivot restriction.
Site functions $\rho$ vanish off $B(\Theta'')\setminus\Theta''$.
Set
\[
 j\rho=\chi_{\rm I}\check D_p\rho,\qquad
 (bB)(x)=\operatorname{Ad}(p(y,x))^{-1}
                G_p^{(1)}(B^0)(y,x),\quad x\in B(y)\setminus\{y\},
 \quad y\in\Theta''.
\]
Here $B^0$ is the zero extension, and $bB$ is zero at the block
centers and outside these blocks. Pairings are complex bilinear;
transposes below refer to them.

\begin{lemma}\label{III-18-47}\label{III-19-8}
Let $f$ be the real background action on the regional fibre.
At a point satisfying Proposition~\ref{II-27-28}, its criticality
and Hessian have the form below. More generally, suppose $f$ is invariant under the block gauge transformations,
is critical at $p$, and has Hessian
\[
 (B_1,B_2)\longmapsto
       \langle C_pB_1,\Delta_k(U_+,J^\circ(U_+))C_pB_2\rangle
       =(B_1,H B_2)_1.
\]
Here the Hessian assertion means that every real tangent field is
the velocity of a twice differentiable fibre curve, and the
second derivative along each such curve is the displayed
quadratic form. Suppose the actual fluctuation operator is
$A=H+\widehat G_p$, with $H$ real symmetric and
$\widehat G_p$ the block-contour gauge form. If $A_{\rm II}$ is
invertible, put $m=-A_{\rm II}^{-1}A_{\rm IE}a$ and
$\bar m=m\oplus a$. Then, for every complex external field $a$,
\begin{equation}\label{III-19-12}
 b m=0,\qquad \chi_{\rm I}\widehat G_p\bar m=0,
                  \qquad \chi_{\rm I}H\bar m=0.
\end{equation}
Consequently, when the localized Wilson first and second
variations are respectively zero and $H$, the two cancellation
hypotheses in \eqref{eq:source-cancellation-hypotheses} hold.
At $k=0$ the same conclusion holds for $F_{Y,1}^{\mathbf1}$
and its polarized quadratic parts.
\end{lemma}

\begin{proof}
We first record the local geometric identities. Differentiating
the covariance of the block average and of the contour gives
\[
 LQ_pD_p\rho(c)
   =\operatorname{Ad}(\mathcal M(p)(c))\rho(c_+)-\rho(c_-)=0,
 \qquad
 G_p^{(1)}(D_p\rho)(y,x)
                  =\operatorname{Ad}(p(y,x))\rho(x).
\]
The endpoints of the coarse bonds are block centers, where
$\rho=0$. Pivot elimination is the unique solution of
$Q_p B=0$ with its given nonpivot values, so
$C_p\check D_p\rho=D_p\rho$ and $bj=1$.
The gauge directions supported in $B(\Theta'')$ have no
external nonpivot component.

Every contour path in a block uses only the nonpivot bonds of
that block; changing pivot values does not change its first
variation. Therefore the gauge form has no mixed integration--external
block, and
\[
 \widehat G_{{\rm I}{\rm E}}=0,
                    \qquad \widehat G_{{\rm I}{\rm I}}=b^{\mathsf T}b.
\]
Indeed its bilinear form is the sum of the pairings of the
first contour variations. Applying the inverse adjoint
transports in the definition of $b$ leaves these pairings
unchanged. This proves the second identity.

We next use gauge invariance of $f$. For a real tangent vector
$B$, choose a twice differentiable curve $V(s)$ in the fibre
with this velocity. For a real site function $\rho$ supported
in the block interiors and a real number $c$, the curve
$V(s)^{\exp(-isc\rho)}$ remains in the fibre and has velocity
$B+cD_p\rho$. Gauge invariance gives the same second derivative
along the two curves. Thus the Hessian $q$ on the constrained
bond fields satisfies
\[
 2c\,q(B,D_p\rho)+c^2q(D_p\rho,D_p\rho)=0
                                              \quad(c\in\mathbb R).
\]
Both coefficients vanish. Surjectivity of pivot elimination
onto the tangent space and symmetry imply
$H\check D_p\rho=0$. Splitting complex arguments into real
and imaginary parts gives the same identity over $\mathbb C$.
In particular,
\[
 H_{{\rm I}{\rm I}}j=H_{{\rm E}{\rm I}}j=0,
 \qquad j^{\mathsf T}H_{{\rm I}{\rm I}}
                    =j^{\mathsf T}H_{{\rm I}{\rm E}}=0.
\]

The conditional-mean equation is
$b^{\mathsf T}bm+H_{{\rm I}{\rm I}}m+H_{{\rm I}{\rm E}}a=0$.
Apply $j^{\mathsf T}$ and use $bj=1$ to obtain $bm=0$.
The gauge part of the equation is therefore zero by itself,
and so is the Wilson part. This proves \eqref{III-19-12}.
Finally, $2\sum_YT_2[\mathbf1](Z|_Y,\bar m|_Y)=(Z,H\bar m)_1=0$
for integration-supported $Z$; at the first step the factor two
is already included in the polarized form $F_{Y,2}$.
\end{proof}

For the regional action, the invariance used here follows by
extending a block gauge to the bare sites, taking it to be the
identity at all other sites. Covariance of the averages identifies
the two shorter constraint sets, preserves criticality and
regularity, and maps their residual gauge orbits onto one another.
The action is gauge invariant, so its value on those orbits is
unchanged. Proposition~\ref{II-27-28} supplies its Hessian at
the same point as the Gaussian. At an actual history the agreement
of the regional pair with the pair used to define that Gaussian
is still part of the applicability check.

\subsection{The exact source split}

For an arbitrary reference function $\mathsf{rf}$ on positive
bonds, let $\mathcal A_b=\mathcal A1_b$. For a field $\zeta$
vanishing on the external bonds, define
\begin{equation}\label{III-103-28}
 \begin{split}
 \mathcal W[w,\mathsf{rf}](Y;P;\mathcal A,\zeta)
 ={}&\sum_{b\in\mathfrak X^{\circ+}\langle Y\rangle}
              T_1[w-\mathsf{rf}(b)](\mathcal A_b)
              +T_{\ge3}[w](\mathcal A)+T_2[w](\zeta)\\
 &+2\sum_{b\in\mathfrak X^{\circ+}_{\rm I}\langle Y\rangle}
            T_2[w-\mathsf{rf}(b)](\zeta_b,\mathcal A-\zeta).
 \end{split}
\end{equation}
It is defined for $P\in\mathbb U_k^{\mathbf1}(Y|h')$ and
$\sup|\mathcal A|<r_0$; no bound on $\zeta$ is needed for the
polynomial terms. The notation distinguishes this positive-source
remainder from the centered member $\mathbb W^{\rm S}[v]$.

\begin{lemma}\label{III-40-14}\label{III-103-32}
Let $P$ be a triple on the torus whose restrictions lie in the
Taylor domains. Let $a$ be an external field, $m=\mathsf M a$ its
unsourced conditional mean, $\bar m=m\oplus a$, and suppose
\begin{equation}\label{eq:source-cancellation-hypotheses}
 \sum_YT_1[\mathbf1](B|_Y)=0,\qquad
 \sum_YT_2[\mathbf1](Z|_Y,\bar m|_Y)=0
\end{equation}
for all fields $B$ and all $Z$ vanishing on the external bonds.
For an integration field $B_{\rm I}$ put
$\mathcal A_Y=g_k(B_{\rm I}\oplus a)|_Y$ and
$\zeta_Y=g_k((B_{\rm I}-m)\oplus0)|_Y$.
Whenever $\sup|\mathcal A_Y|<r_0$ for every $Y$,
\begin{equation}\label{III-103-39}
 \sum_Y\mathbb V^w(Y;P|_Y,\mathcal A_Y)
 =\sum_Y\mathcal W[w,\mathsf{rf}](Y;P|_Y;\mathcal A_Y,\zeta_Y)
                         +\sum_YT_2[w](g_k\bar m|_Y).
\end{equation}
For the identification with the canonical centered member,
assume in addition, for every $Y$, that
\[
 P|_Y\in\mathbb U_k^{\mathbf1}(Y|h'),\qquad
 \sup|g_k\bar m|_Y<r_0 .
\]
Together with the assumed bound on $\mathcal A_Y$, these are
its common domain conditions. At $v=-x_k^{-1}w$ and
$\mathsf{rf}(b)=w(p_b)$, the canonical centered member evaluated
at $z=B_{\rm I}-m$, $\mu=m$ is then defined and equals the
corresponding term $\mathcal W[w,\mathsf{rf}]$.
Thus the second term in \eqref{III-103-39} is the quadratic source
at the conditional mean, localized in \cref{III-112-28}.
\end{lemma}

\begin{proof}
Expand the potential as $T_1+T_2+T_{\ge3}$. Decompose its linear
term bond by bond using $w=(w-\mathsf{rf}(b))+\mathsf{rf}(b)\mathbf1$.
After summing over $Y$, every reference contribution vanishes by
the first identity in \eqref{eq:source-cancellation-hypotheses},
applied to the field supported on that bond.
Next write $\mathcal A_Y=\zeta_Y+g_k\bar m|_Y$ and polarize the
quadratic term. Split its mixed term in the weight in the same
way. Its reference terms vanish by the second identity, again
bond by bond. What remains is \eqref{III-103-28} and the quadratic
part at $g_k\bar m$. All sums are finite. For the last
assertion, the additional domain conditions ensure that the
three centered terms are defined before they are added.
Add \eqref{III-39-44} and \eqref{III-39-35}, substitute
$x_kv=-w$, and use linearity in both the weight and the field.
The first-order sums join over all bonds, the two cubic remainders
at $g_k\phi$ cancel, and the remaining terms are exactly
\eqref{III-103-28}.
\end{proof}

The cancellation hypotheses in this lemma are checked at the
actual chain points in the measure argument. In particular the
mixed cancellation concerns the Wilson Hessian, not merely the
full Hessian with its gauge term.


\subsection{The escort on the one-lattice cores}

The one-lattice family is needed to identify the terms whose
supports remain inside the next small-field region. Write
$h_\emptyset$ for the history without large-field regions and
$\mathcal Q_{\mathfrak h}(Y)=\mathbb U_k^{\mathbf1}(Y|h_\emptyset)$.
On this history every bond is an integration bond. Define
\begin{equation}\label{III-41-18}
 \mathbb W^{\rm 1l}[v](Y;P,B)
 =-\mathbb V^{x_kv}(Y|h_\emptyset;P,g_kB)
       +\sum_bx_kv(p_b)T_1[\mathbf1](Y|h_\emptyset;P;g_kB_b).
\end{equation}
Its domain is precisely $P\in\mathcal Q_{\mathfrak h}(Y)$,
$g_k\sup|B|<r_0$. The sum is over the positive one-lattice
nonpivot bonds in $Y$.

\begin{proposition}\label{III-41-24}\label{III-41-27}
Assume the strengthened localization rates used above, the core
inclusion $\mathcal U^\circ(Y)\subset\mathcal Q_{\mathfrak h}(Y)$,
and $C_{\mathcal A}\delta_k\le r_0/2$.
On $\mathcal U^\circ(Y)\times\mathbb D^{\rm f}(Y)$,
\eqref{III-41-18} is holomorphic in $B$, vanishes at $B=0$,
depends only on the pair of $P$, and is gauge invariant.
For fixed real $U$, it is jointly holomorphic in the weight,
$(A',\mathbf J)$ and $B$ on every open core slice.
Its conjugate is $\mathbb W^{\rm 1l}[\bar v]$.
For $0\le r\le r^\sharp-2$,
\begin{equation}\label{III-41-25}
 \|\mathbb W^{\rm 1l}[v]\|^{\rm co}_{\star,r}
    \le N^\circ[v],
\end{equation}
where
\begin{equation}\label{III-41-20}
 \begin{split}
 N^\circ[v]=64\bigl[&3eK_cc_{\mathfrak h}M^6L^{-2}\alpha_0
                         g_kx_k\mathrm{sw}_k(v)\\
   &+(6K_0r_0^{-2}M^6
        +2K_0r_0^{-3}C_{\mathcal A}^3g_k\mathfrak b_k^3)
                                                     |v|_\infty\bigr].
 \end{split}
\end{equation}
The norm is taken after restricting $P\in\mathcal U^\circ(X)$
to $Y$ with $Y^\sharp\subset X$, and dividing by
$1+q_Y(B)$. Thus this is one escort on all the one-lattice cores.

Suppose also that the canonical shifts restrict to the
one-lattice shift inside $\Lambda_t(h')$. For $Y\subset\Lambda_t(h')$
and every $(P,\mathcal A)$ in the one-lattice domain,
\begin{equation}\label{III-41-30}
 \mathbb W[v](Y|h';P,\mathcal A;0)
       =\mathbb W^{\rm 1l}[v](Y;P,g_k^{-1}\mathcal A),
 \qquad \mathbb W^1[v](Y|h';P,\mathcal A;0)
       =\mathbb W^3[v](Y|h';P,\mathcal A;0)=0.
\end{equation}
This identity holds on the actual middle tiers whenever their
restrictions lie in the one-lattice domain. It gives one escort
for all such histories, with the same choice of reference plaquettes.
\end{proposition}

\begin{proof}
The core inclusion and the field room put $g_kB$ in the
half-radius Taylor ball. On the constant history there are no
plaquettes of layers $m<k$. Apply the linear Taylor estimate
with $a_k=L^{-2}\alpha_0$ and $a_m=\epsilon$ for $m<k$, then
let $\epsilon\downarrow0$. Since $\varpi^w_0=1$, this gives
\[
 |T_1[w](\mathcal A)|\le eK_cc_{\mathfrak h}M^3
        L^{-2}\alpha_0|w|_Y\sup|\mathcal A|e^{-r^\sharp d_k(Y)}.
\]
The Taylor expansion of minus \eqref{III-41-18} is
\[
 \sum_bT_1[x_k(v-v(p_b))](g_kB_b)
             +T_2[v](B)+T_{\ge3}[x_kv](g_kB).
\]
Put $d=d_k(Y)$. The cube path defining the local oscillation
has at most $d+2$ links, so
$|v-v(p_b)|_Y\le(d+2)\mathrm{sw}_k(v)$.
Also $\#_{\pi_k}Y\le8(d+1)$ and
$\sum_b|B(b)|\le12M^3(d+1)(1+q_Y(B))$.
Using $d+1\le e^d$ and $d+2\le2e^d$, the first term is bounded by
\[
 24eK_cc_{\mathfrak h}M^6L^{-2}\alpha_0g_kx_k\mathrm{sw}_k(v)
                   (1+q_Y(B))e^{-(r^\sharp-2)d}.
\]
The other terms are bounded respectively by
$K_0r_0^{-2}M^3|v|_\infty q_Y(B)e^{-r^\sharp d}$ and
$2K_0r_0^{-3}C_{\mathcal A}^3g_k\mathfrak b_k^3|v|_\infty
 e^{-r^\sharp d}$. Their sum is at most the stated bound.

The local formulas and their Taylor coefficients have
holomorphic extensions at fixed $U$. On an open core slice,
these local extensions agree with the given function on open
intersections, so the function is holomorphic there. Linearity
in the finitely many weight coordinates makes holomorphy joint.
Pair dependence and both symmetries follow term by term.
Restriction to $Y^\sharp\subset X$ preserves the core, by
Lemma~\ref{II-16-46}; it therefore preserves the bound and
all the stated properties.

For $Y\subset\Lambda_t$, there are no external bonds and the
stored datum on $Y$ is empty. The canonical shift and each
localized Wilson member agree with their one-lattice counterparts.
At center zero, $\phi=0$. Hence the two terms in
\eqref{III-39-35} vanish, the mixed Taylor term in
\eqref{III-39-34} vanishes, and the remaining expression is
\eqref{III-41-18}. This proves \eqref{III-41-30}.
\end{proof}

\subsection{The first step}

At $k=0$ the shift is the regional pivot completion itself.
We give its weighted calculation separately. Let
$\mathsf{Pl}_0(Y)=\{p\in\mathcal P(\Omega_0^\diamond):X(p)=Y\}$,
and set $|v|_{0,Y}=\sup_{\mathsf{Pl}_0(Y)}|v|$, with value zero
for an empty set. Require the following local admission condition:
\begin{equation}\label{eq:first-source-admission}
 P|_{Y_1}\in\mathcal Q_1(Y_1;a_P^\circ)
 \quad\text{whenever }Y_1\subset Y,\quad
 \#_{\pi_0}Y_1\le4,\quad Y_1\text{ contains a source cube}.
\end{equation}
The constants $a_P^\circ,r_P^\circ,C_P$ are those of
Definition~\ref{II-41-43}. On $\sup|\mathcal A|<r_P^\circ$,
complete the nonpivot field by setting the omitted nonpivot
coordinates to zero on each pivot block and applying its
analytic pivot chart. Write this completion as $B'(P,\mathcal A)$,
extended by zero outside its regional bonds. Define
\begin{equation}\label{III-48-6}
 F_Y^v[P](\mathcal A)=\sum_{p\in\mathsf{Pl}_0(Y)}v(p)
       [a_p(e^{iB'(P,\mathcal A)}\mathbf U)-a_p(\mathbf U)].
\end{equation}
This remains a partial map with the completion's domain when
the sum is empty. Put $K^0=4512C_PM^3r_P^\circ$.
Let $F^v_{Y,n}$ denote its homogeneous Taylor parts, and use the
following polarization convention at this step:
\[
 F^v_{Y,2}(A_1,A_2)
       =F^v_{Y,2}(A_1+A_2)-F^v_{Y,2}(A_1)-F^v_{Y,2}(A_2).
\]
In particular, $F^v_{Y,2}(A,A)=2F^v_{Y,2}(A)$.

\begin{proposition}\label{III-48-15}
Under \eqref{eq:first-source-admission}, $F_Y^v$ is holomorphic
on the radius-$r_P^\circ$ ball, vanishes at zero, and has modulus
at most $K^0|v|_{0,Y}$. It vanishes unless $Y$ contains a source
cube and $\mathsf{Pl}_0(Y)\ne\emptyset$; in the latter case
$d_0(Y)=0$. Its Taylor parts satisfy
\begin{align}
 |F^v_{Y,n}(A)|&\le K^0|v|_{0,Y}(\sup|A|/r_P^\circ)^n,
                                                   \label{III-48-16}\\
 |\mathcal R_NF_Y^v(A)|&\le2K^0|v|_{0,Y}(\sup|A|/r_P^\circ)^N
       \quad(\sup|A|\le r_P^\circ/2),                \label{III-48-17}\\
 |F^v_{Y,2}(A_1,A_2)|&\le4K^0|v|_{0,Y}(r_P^\circ)^{-2}
                             \sup|A_1|\sup|A_2|.   \label{III-48-18}
\end{align}
If $|\partial\mathbf U(p)-1|\le\varepsilon_\partial$ on
$\mathsf{Pl}_0(Y)$, then
\begin{equation}\label{III-48-19}
 |F^v_{Y,1}(A)|\le192C_PM^3|v|_{0,Y}
                              \varepsilon_\partial\sup|A|.
\end{equation}
The potential and its parts are linear and local in the weight,
depend on $P$ only through $\mathbf U$, and have the gauge and
conjugation laws of the Wilson action. They are Borel on the
admitted domains. At fixed real $U$, they have holomorphic
extensions on the open set where all the required pivot blocks
are admitted, jointly with the field and the weight.
\end{proposition}

\begin{proof}
Each pivot block contains at most two fine cubes and contains a
source cube. Thus \eqref{eq:first-source-admission} admits its
triple to the pivot chart. The construction of
Section~\ref{subsec:ymm1-50-1}, with the unused coordinates set to zero,
gives the same local chart and the bound
$\sup|B'|\le2C_P\sup|\mathcal A|\le1$.
Every pivot uses only its block; restriction, gauge covariance,
conjugation and dependence on the configuration pass to the
regional completion. At fixed $U$, intersection of the finitely
many open pivot domains is open. On it the completion is
holomorphic jointly with $\mathcal A$ and continuous on its
admitted domain.

A nonempty $\mathsf{Pl}_0(Y)$ requires $Y=X(p)$ for a plaquette.
Such a region consists of one, two or four fine cubes with a
common corner; it has $d_0(Y)=0$ and at most $12M^3$ plaquettes.
If it contains a source cube, admission at $Y$ gives
$|\mathbf U(b)^{\pm1}|\le e^{1/L}$. Put
$\epsilon=2C_P\sup|\mathcal A|\le1$ and telescope the four
plaquette factors. The product difference is at most
$4(e^\epsilon-1)e^{3\epsilon+4/L}\le188\epsilon$.
Since $|\operatorname{tr}Z|\le2|Z|$, summing the weighted action
differences gives
$|F_Y^v|\le4512C_PM^3|v|_{0,Y}\sup|\mathcal A|$.
This proves the uniform bound and holomorphy. If there is no
source cube, the regional field is empty and its extension is
zero. Vanishing at zero is immediate from the completion.
The one-variable Cauchy formula proves the coefficient and tail
bounds. On the bidisc
$|s_i|<r_P^\circ/(2\sup|A_i|)$, the coefficient of $s_1s_2$
in $F_Y^v(s_1A_1+s_2A_2)$ is the polarized quadratic part;
the bidisc Cauchy bound proves \eqref{III-48-18}.

For the sharper derivative, write $W=\partial\mathbf U(p)$.
Differentiation in a field $C$ gives
\[
 D_pa_p(C)=-\tfrac i2\operatorname{tr}(\Xi W)
         =-\tfrac i2\operatorname{tr}(\Xi(W-1)),
\]
where $\Xi$ is the sum of the four transported signed bond
fields. Its trace is zero and
$|\Xi|\le(1+e^{2/L}+e^{6/L}+e^{8/L})\sup|C|
 \le8\sup|C|$ for $L\ge13$.
The derivative of $B'(P,sA)$ at zero has supremum at most
$2C_P\sup|A|$. The chain rule and the plaquette count give
$12M^3\cdot8\cdot2C_P|v|_{0,Y}\varepsilon_\partial\sup|A|$,
which is \eqref{III-48-19}.
All structural assertions follow from the same finite formula
and its Cauchy coefficients. Borel measurability can also be
read from the limits of finite differences at real scalars.
\end{proof}

Let $z,a,\mu,\phi=\mu\oplus a$ be as above at $k=0$, and
assume
\begin{equation}\label{eq:first-source-room}
 (C_{\mathcal A}+\theta_{\rm sl})\delta_0\le r_P^\circ/2.
\end{equation}
Define the three centered members by
\begin{align}
 \mathbb W_0[v]={}&-[F_Y^{x_0v}(g_0(z+\phi))-F_Y^{x_0v}(g_0\phi)]
  +\sum_{b\in\mathfrak X^{\circ+}_{\rm I}\langle Y\rangle}
    x_0v(p_b)[F^{\mathbf1}_{Y,1}(g_0z_b)
                 +F^{\mathbf1}_{Y,2}(g_0z_b,g_0\phi)],
                                                   \label{III-48-23}\\
 \mathbb W^{\mathrm{src},1}_0[v]
   ={}&-\sum_{b\in\mathfrak X^{\circ+}\langle Y\rangle}
                   F^{x_0(v-v(p_b))}_{Y,1}(g_0\phi_b),\notag\\
 \mathbb W^{\mathrm{src},3}_0[v]
   ={}&-\mathcal R_3F_Y^{x_0v}(g_0\phi).                \label{III-49-4}
\end{align}
The mixed term in \eqref{III-48-23} has no factor two, by the
first-step polarization convention.

\begin{proposition}\label{III-48-27}\label{III-49-9}
Suppose all triples under consideration satisfy
\eqref{eq:first-source-admission}, the plaquette bound
$|\partial\mathbf U-1|\le a_0^{\mathsf{Tab}}$, and
\eqref{eq:first-source-room}. If $|v|_\infty\le\sigma$ and
$x_0\mathrm{sw}_0(v)\le\rho'$, the three members are defined on
the entire stated tier times the center ball. They are linear
and local in the weight, analytic and Borel, depend on the pair,
and obey the gauge and conjugation laws. They vanish unless
$d_0(Y)=0$. The last two are independent of $z$ and vanish at
$\phi=0$. At every rate $r\ge0$ their sum has norm at most
\begin{equation}\label{III-49-6}
 \begin{split}
 s_0^{\rm ce}={}&C_0^{\rm src}
 [\sigma(1+g_0\mathfrak b_0^3)
                +\rho'(g_0a_0^{\mathsf{Tab}}+g_0\delta_0)]\\
 &+9216C_PM^6\Theta\rho'a_0^{\mathsf{Tab}}\delta_0
       +2K^0(r_P^\circ)^{-3}\Theta^3\sigma g_0\mathfrak b_0^3,
 \end{split}
\end{equation}
where $\Theta=\max(C_{\mathcal A},\theta_{\rm sl})$ and
\begin{equation}\label{III-48-30}
 C_0^{\rm src}=\max\{4608C_PM^6,96K^0(r_P^\circ)^{-2}\Theta M^3,
 24K^0(r_P^\circ)^{-2}M^6,
 4K^0(r_P^\circ)^{-3}(C_{\mathcal A}+\theta_{\rm sl})^3\}.
\end{equation}
The norm has denominator $1+M^{-3}\sum_{b\in\mathfrak X_{\rm I}^{\circ+}}
|z(b)|^2$. Joint holomorphy holds for affine holomorphic weight
families with the same bounds. If the first unsourced frame has
size $s_0^{\rm rg}\le1$ and
\[
 e s_0^{\rm rg}+e w_0^{-1}s_0^{\rm ce}
                              \le(4L^3K^+_{\rm ap})^{-1},
\]
these members satisfy the size conditions for its centered expansion.
\end{proposition}

\begin{proof}
The two fields $g_0(z+\phi)$ and $g_0\phi$ lie in the half-radius
ball. On a supporting domain there are at most $12M^3$ positive
bonds. If $N=1+q_Y(z)$, then
$\sum_b|z(b)|\le6M^3N$ and
$(\sum_b|z(b)|)^2\le12M^6N$.
A plaquette in $\mathsf{Pl}_0(Y)$ can be joined to any reference
plaquette $p_b$ by at most four steps of length $M$: pass through
the centers of the at most four cubes forming $X(p)$.
Thus $x_0\sup_{p,b}|v(p)-v(p_b)|\le4\rho'$.
Taylor expansion decomposes \eqref{III-48-23} into
\[
 -\sum_bF^{x_0(v-v(p_b))}_{Y,1}(g_0z_b)
 -\sum_bF^{v-v(p_b)}_{Y,2}(z_b,\phi)-F^v_{Y,2}(z)
 -[\mathcal R_3F_Y^{x_0v}(g_0(z+\phi))
                         -\mathcal R_3F_Y^{x_0v}(g_0\phi)].
\]
Using \eqref{III-48-16}--\eqref{III-48-19}, the four terms,
divided by $N$, are bounded by the four summands
\[
 4608C_PM^6\rho'g_0a_0^{\mathsf{Tab}},\quad
 96K^0(r_P^\circ)^{-2}\Theta M^3\rho'g_0\delta_0,\quad
 24K^0(r_P^\circ)^{-2}M^6\sigma,\quad
 4K^0(r_P^\circ)^{-3}(C_{\mathcal A}+\theta_{\rm sl})^3
                                        \sigma g_0\mathfrak b_0^3.
\]
The same derivative bound, summed over all bonds, gives the
bound $9216C_PM^6\Theta\rho'a_0^{\mathsf{Tab}}\delta_0$ for
the first member of \eqref{III-49-4}. Its last member is bounded
by $2K^0(r_P^\circ)^{-3}\Theta^3\sigma g_0\mathfrak b_0^3$.
Since $d_0(Y)=0$ on every supporting domain, no decay is lost
at any rate. These inequalities prove \eqref{III-49-6}.
The structural assertions follow from the preceding proposition;
in particular their holomorphic extensions are defined on the
open set where the pivot blocks and both shifted fields are
admitted. The displayed amplitude condition implies
$s_0^{\rm ce}<w_0$ and is the required centered-frame condition.
\end{proof}

The exact first-step split is as follows. Define
\begin{equation}\label{III-49-3}
 \tfrac12(B,\Upsilon_0(P;v)B)_1
       =\sum_YF^v_{Y,2}[P|_Y](B|_Y).
\end{equation}
At a triple on the torus, let $m$ be the unsourced conditional
mean and $\bar m=m\oplus a$. Suppose its first and mixed
quadratic Wilson derivatives cancel:
\[
 \sum_YF^{\mathbf1}_{Y,1}(B|_Y)=0,\qquad
 \sum_YF^{\mathbf1}_{Y,2}(Z|_Y,\bar m|_Y)=0
 \quad(Z|_{\mathfrak X_{\rm E}}=0).
\]
Whenever all the terms are admitted and $z=B'_{\rm I}-m$ is in
the centered field polydisc, finite Taylor expansion gives
\begin{equation}\label{III-49-11}
 \begin{split}
 -\sum_Y F_Y^{x_0v}[P|_Y](g_0(B'_{\rm I}\oplus a)|_Y)
 ={}&-\tfrac12(\bar m,\Upsilon_0(P;v)\bar m)_1\\
   &+\sum_Y[\mathbb W_0+\mathbb W^{\mathrm{src},1}_0
          +\mathbb W^{\mathrm{src},3}_0][v](Y;P|_Y,z|_Y;m|_Y).
 \end{split}
\end{equation}
Indeed the reference first derivatives cancel bond by bond;
the reference mixed derivatives cancel in the same manner.
The remaining quadratic mean term is \eqref{III-49-3}, since
$x_0g_0^2=1$. The two cubic terms at $g_0\bar m$ cancel between
the first and third centered members. This is the first-step
version of \cref{III-103-32}; verification of its cancellation
hypotheses at the actual chain points is part of the measure argument.

\begin{proposition}\label{III-48-54}\label{III-49-13}
Suppose $\alpha_0,\alpha_1\le a_P^\circ$ and
$C_{\mathcal A}\delta_0\le r_P^\circ/2$. On the history without
large-field regions define
\begin{equation}\label{III-48-52}
 \mathbb W^{\rm 1l}_0[v](Y;P,B)
     =-F_Y^{x_0v}[P](g_0B)
             +\sum_bx_0v(p_b)F^{\mathbf1}_{Y,1}[P](g_0B_b).
\end{equation}
This is an escort on every core
$\mathcal U_1^\circ(Y;\alpha_0,\alpha_1)\times\mathbb D^{\rm f}(Y)$.
It is linear and local in the weight, depends on the pair, is
gauge invariant, vanishes at $B=0$, and is jointly holomorphic
in the weight, field and open fixed-real core slices. Its
conjugate is $\mathbb W^{\rm 1l}_0[\bar v]$. For every $r\ge0$,
\begin{equation}\label{III-48-56}
 \begin{split}
 \|\mathbb W^{\rm 1l}_0[v]\|^{\rm co}_{\star,r}\le N_0^\circ
 ={}&4608C_PM^6\rho'g_0L^{-2}\alpha_0
            +24K^0(r_P^\circ)^{-2}M^6\sigma\\
   &+2K^0(r_P^\circ)^{-3}C_{\mathcal A}^3\sigma g_0\mathfrak b_0^3
 \end{split}
\end{equation}
if $|v|_\infty\le\sigma$, $x_0\mathrm{sw}_0(v)\le\rho'$.
For any history and any $Y\subset\Lambda_1$ on which the triple
satisfies \eqref{eq:first-source-admission}, the sum of its three
centered source members at center zero equals
$\mathbb W^{\rm 1l}_0[v]$ on their common field domain.
\end{proposition}

\begin{proof}
The core restricts into the simple triple domain at radii
$(\alpha_0,\alpha_1,\alpha_0)$ by Lemma~\ref{II-16-46}.
Radius monotonicity gives \eqref{eq:first-source-admission}, and
its plaquette bound is $L^{-2}\alpha_0$. The field room admits
$g_0B$ to the half-radius ball. Expanding \eqref{III-48-52}
gives the first, unmixed quadratic and cubic terms in the preceding
proof, with $\phi=0$, $\varepsilon_\partial=L^{-2}\alpha_0$
and field bound $C_{\mathcal A}\mathfrak b_0$. This gives exactly
\eqref{III-48-56}. All supporting domains have $d_0(Y)=0$.
The finite formula and the open pivot extensions prove the
analyticity and symmetry claims; restriction to $Y^\sharp\subset X$
preserves the core and the bound.
Inside $\Lambda_1$ all bonds are integration bonds, the stored
datum is empty, and the regional pivot completion is the
one-lattice completion. At center zero $\phi=0$, so the two
members of \eqref{III-49-4} vanish and \eqref{III-48-23} reduces
to \eqref{III-48-52}.
\end{proof}

\section{The quadratic part of the source}
\label{sec:quadratic-source}

The quadratic part of the plaquette source is evaluated at the
conditional mean of the unsourced Gaussian. It therefore changes the
retained quadratic factors while leaving that Gaussian and its
conditional mean fixed. The resulting correction is linear in the
source. We first construct its local terms, with bounds on their
domains of definition; the measure argument will later justify their
evaluation on the histories.

\subsection{Local quadratic forms and their matrices}

Fix a step $t=k+1$ and its integration and external bonds
$\mathfrak X=\mathfrak X_{\rm I}\sqcup\mathfrak X_{\rm E}$.
Coordinate fields are indexed by positively oriented bonds and take
values in $\mathbb C^3$. All quadratic pairings in this section are
complex bilinear. We use the canonical coefficient matrices
$a_S$, the local inverse terms $t_\jmath$, and their footprints from
Proposition~\ref{II-63-60}. On the whole torus their sums are
\[
 A=\sum_Sa_S,\qquad T=\sum_\jmath t_\jmath=A_{\rm II}^{-1},
 \qquad \mathsf M=-TA_{\rm IE}.
\]
The last identity defines the conditional-mean operator. These
identities are used only at points in the admitted pair class of
\cref{II-63-60}. Every individual coefficient is evaluated on its own hull.

Let $\mathbb U(Y)$ be a set of triples $P=(U_{\rm r},A',\mathbf J)$
at level $t$ on $Y\in\mathcal D_k$. Assume that symmetric complex
bilinear forms $q_Y(P;B,B')$ satisfy
\begin{equation}\label{eq:local-quadratic-hypothesis}
 |q_Y(P;B,B)|\le k_\upsilon e^{-r_\upsilon d_k(Y)}
       \sup_{b\in\mathfrak X^\circ\langle Y\rangle}|B(b)|^2,
 \qquad k_\upsilon>0.
\end{equation}
The supremum over an empty set is zero. On each fixed real-factor
slice, the coefficients of $q_Y$ have local holomorphic extensions
at every point of $\mathbb U(Y)$. This is the precise analyticity
requirement; the sets $\mathbb U(Y)$ themselves need not be open.

To obtain open domains for the matrices, fix a real configuration
$U$ on the bare torus. For every region $R$ of the level-$t$
partition, let $\mathscr P_U(R)$ be the interior of the set of pairs
$(A',\mathbf J)$ such that
\begin{equation}\label{III-112-36}
 \begin{split}
 \Psi_R(A',\mathbf J)
   &:=(e^{i\eta_t A'}U,L^{-3}\mathbf J)\big|_{{\rm Bd}^\circ
                                  (R\cap\Omega_0^\diamond)}
                       \in\mathscr O(R;\alpha),\\
 (U,A',\mathbf J)|_Y&\in\mathbb U(Y)
       \quad\text{whenever the level-$t$ hull of $Y$ lies in $R$}.
 \end{split}
\end{equation}
Restriction maps these interiors into one another: the coordinate
restriction is a surjective linear map and hence open, and the two
conditions in \eqref{III-112-36} are preserved by restriction.
The map $\Psi_R$ is holomorphic and commutes with restriction.

For $Y\in\mathcal D_k$, let $X(Y)$ be its level-$t$ hull. Define a
symmetric matrix $\upsilon_Y$ on $\mathscr P_U(X(Y))$ by
\begin{equation}\label{III-112-37}
 \langle\operatorname{vec}B,\upsilon_Y(P)\operatorname{vec}B'\rangle
                 =-2q_Y(P|_Y;B,B').
\end{equation}
Here $P=(U,A',\mathbf J)$; the notation leaves the fixed real factor
implicit. Let
\begin{equation}\label{III-112-38}\equationalias{III-112-39}
 c_E=\log18432,\qquad C_E=4096,\qquad
 \beta_\upsilon=8C_Ee^{6\delta_\upsilon M}k_\upsilon,
 \qquad r_\upsilon-\delta_\upsilon M\ge c_E,
\end{equation}
where $0<\delta_\upsilon\le\delta_{\rm pl}$, the common rate of
the canonical coefficients and inverse terms. Lengths are measured
at level $k$.

\begin{proposition}\label{III-112-42}
The matrices $\upsilon_Y$ are well defined and holomorphic. They
are supported on $\mathfrak X^{\circ+}\langle Y\rangle$ in both
indices and admit the majorants and connecting lengths
\begin{equation}\label{eq:source-vertex-majorants}
 m_Y(b,b')=8k_\upsilon e^{-r_\upsilon d_k(Y)}
       1_{\{b,b'\in\mathfrak X^{\circ+}\langle Y\rangle\}},
 \qquad l_Y=M(d_k(Y)+6).
\end{equation}
They satisfy
\begin{equation}\label{III-112-43}
 \sum_Ym_Y(b,b')e^{\delta_\upsilon l_Y}\le\beta_\upsilon.
\end{equation}
For every live pair of bonds there is a connected graph of length
at most $l_Y$, containing their initial points and lying within
distance $M$ of every point of the footprint $Y$.
\end{proposition}

\begin{proof}
Write $c_Y=k_\upsilon e^{-r_\upsilon d_k(Y)}$ and
$N_Y(B)=\sup_{\mathfrak X^\circ\langle Y\rangle}|B|$.
Polarization and rescaling give
\[
 |q_Y(P;B,B')|\le2c_YN_Y(B)N_Y(B').
\]
If one seminorm vanishes, apply polarization to $(sB,B')$ and let
$s\to\infty$; the corresponding mixed value vanishes as well.
The matrix in \eqref{III-112-37} is therefore supported on $Y$.
For a coordinate field with value $z$ at one bond,
$N_Y(\operatorname{vec}^{-1}z)\le\sqrt2|z|$. Applying the last
inequality at two single-bond fields gives the block bound
$2\cdot2c_Y\cdot\sqrt2\cdot\sqrt2=8c_Y$.
Its coefficients are the values of $q_Y$ on coordinate basis
fields. On the interior in \eqref{III-112-36}, their local
holomorphic extensions agree with those values on an open set;
hence they define holomorphic coefficients there.

A minimal edge tree meeting every $M$-cube of $Y$ has length
$Md_k(Y)$. Join the initial points of the two bonds to that tree
inside their cubes, at a cost at most $3M$ each. Every point of
$Y$ is within $M$ of the original tree. This proves the geometric
assertion. For a fixed bond $b$, every domain contributing to
\eqref{III-112-43} contains the $M$-cube of $b_-$. The connected
domain estimate of Proposition~\ref{II-18-26}
\[
 \sum_{Y\supset\square}e^{-c d_k(Y)}\le C_E
                       \qquad(c\ge c_E)
\]
then gives
\[
 \sum_Ym_Y(b,b')e^{\delta_\upsilon l_Y}
 \le8k_\upsilon e^{6\delta_\upsilon M}
       \sum_{Y\supset\square}
         e^{-(r_\upsilon-\delta_\upsilon M)d_k(Y)}
 \le\beta_\upsilon.
\]
\end{proof}

\begin{lemma}[Quadratic parts of the plaquette weight]\label{III-113-21}
Let $s_0>0$ and $|v|_\infty\le x_k s_0$. For $k\ge1$, use the weighted
potential and its Taylor parts of \cref{III-39-24}, and set
\[
 \mathbb U(Y)=\mathbb U_k^{\mathbf1}(Y|h'),\qquad
 q^v_Y(P;B,B')=T_2[v](Y|h';P;g_kB|_Y,g_kB'|_Y).
\]
For $k=0$, use the first-step potential $F_Y^v[P]$ of
\cref{III-48-15} on $\mathbb U_0^\diamond(Y)$, and set
\[
 q^v_Y(P;B,B')=\tfrac12F^v_{Y,2}[P](g_0B|_Y,g_0B'|_Y).
\]
The polarization convention for $F_{Y,2}$ has no factor $1/2$;
its value on $(B,B)$ is twice its homogeneous quadratic part.
Then \eqref{eq:local-quadratic-hypothesis} and its local
holomorphic-extension requirement hold with
\begin{equation}\label{III-113-20}
 \begin{array}{c|c|c}
  &k_\upsilon&r_\upsilon\\ \hline
  k\ge1&K_0r_0^{-2}s_0&r_V\\
  k=0&K^0(r_P^\circ)^{-2}s_0&\text{any real number}.
 \end{array}
\end{equation}
Here $r_V$ is the rate used for the weighted-potential bound, and
\[
 r_0=\min\{\varepsilon_{\mathfrak h},(3c_{\mathfrak h})^{-1}\},
 \quad K_0=K_c C'_h(L,\delta_V,1)M^3,
 \quad K^0=4512C_PM^3r_P^\circ.
\]
The forms are linear in $v$, vanish at $v=0$, and depend only on
plaquettes with all corners in $Y$. They have the pair and gauge
properties of the weighted potential; Borel measurability and
conjugation follow under the corresponding hypotheses on the shift
system. At the first step those two properties require no shift
hypothesis.
\end{lemma}

\begin{proof}
For $k\ge1$, insert $g_kB|_Y$ in \eqref{III-39-25} with $n=2$.
The factor $x_kg_k^2=1$ cancels the bound $x_ks_0$ for the weight,
giving \eqref{eq:local-quadratic-hypothesis}. The coefficient
functions have local holomorphic extensions by the Taylor
construction; applying polarization to the extensions at
$g_k(B+B')|_Y$, $g_kB|_Y$ and $g_kB'|_Y$ gives the required
extension for $q^v_Y(P;B,B')$. At the first step, the same
calculation uses \eqref{III-48-16} and $x_0g_0^2=1$.
The potential there vanishes unless $d_0(Y)=0$, so the bound holds
with every $r_\upsilon$.
Linearity, pair dependence and covariance pass from the potential
to its Taylor coefficients. Every plaquette used by the localized
potential on $Y$ lies in a lower-scale domain whose level-$k$ hull
is contained in $Y$; thus all its corners are in $Y$. This gives
locality in the weight. The first-step sum has the same property
directly. The Borel and conjugation assertions pass through the
Taylor limits and polarization.
\end{proof}

In particular the matrix bound for these forms is
\begin{equation}\label{III-113-22}
 \beta_\upsilon=8C_Ee^{6\delta_\upsilon M}s_0
 \begin{cases}
  K_0r_0^{-2},&k\ge1,\\
  K^0(r_P^\circ)^{-2},&k=0.
 \end{cases}
\end{equation}
At the first step choose $r_\upsilon=c_E+\delta_\upsilon M$.
At later steps the inequality
$r_V-\delta_\upsilon M\ge c_E$ is a substantive rate requirement.

\subsection{The correction at the conditional mean}

We describe the localization with its numerical losses. A local
matrix family with constants $(B,\delta,\rho)$ has holomorphic
members $O_\iota$ on their hulls, nonnegative block majorants
$m_\iota$ and connecting lengths $l_\iota$ such that
\[
 |O_\iota(p)(b,b')|\le m_\iota(b,b'),\qquad
 \sum_\iota m_\iota(b,b')e^{\delta l_\iota(b,b')}\le B.
\]
Every live pair has a connected graph of length at most
$l_\iota(b,b')$ containing both initial points; the footprint is
within distance $\rho$ of this graph and contains the two bonds.
All restrictions of the family use the same indices, majorants and
lengths. These are the properties of the canonical families proved
in Proposition~\ref{II-63-60}. They also hold for the source
vertices by \cref{III-112-42}, with $\rho=M$. The canonical slack
is $4M/L^2\le M$, and $4M\le LM$.

In the following construction the two canonical families have
constants $(B_{\rm f},\delta_\upsilon,\rho_\upsilon)$, the source
family has constants $(\beta_\upsilon,\delta_\upsilon,\rho_\upsilon)$,
and $B_{\rm f}\ge1$, $4\rho_\upsilon\le LM$. Write
$p_{\rm I},p_{\rm E}$ for the coordinate projections. There are
four kinds of matrix products with exactly one source vertex:
\begin{equation}\label{III-112-18}
 \begin{split}
 d_Y&=p_{\rm E}\upsilon_Yp_{\rm E},\\
 d_{Y\jmath S}&=-p_{\rm E}\upsilon_Yp_{\rm I}t_\jmath
                                      p_{\rm I}a_Sp_{\rm E},\\
 d_{S\jmath Y}&=-p_{\rm E}a_Sp_{\rm I}t_\jmath
                                      p_{\rm I}\upsilon_Yp_{\rm E},\\
 d_{S_1\jmath_1Y\jmath_2S_2}
   &=p_{\rm E}a_{S_1}p_{\rm I}t_{\jmath_1}p_{\rm I}\upsilon_Y
                         p_{\rm I}t_{\jmath_2}p_{\rm I}a_{S_2}p_{\rm E}.
 \end{split}
\end{equation}
Each index includes its kind. The footprint of a product is the
union of the footprints of its factors, and its hull is denoted
by $X(\mathfrak z)$. Its majorant is the corresponding product
summed over intermediate bonds; its length is the minimum sum of
the constituent lengths over live intermediate bonds. Products
whose majorant vanishes are omitted.

Use the keeping rule of Section~\ref{sec:history-conditional-inverse}: a product is
kept on $\mathfrak k$ when
$\mathsf{kp}_t(X(\mathfrak z))=1$ and
$\mathfrak k=\mathfrak k_{X(\mathfrak z)}$. Let
$d^{\rm k}_{\mathfrak z}$ be its compression to that component,
or zero if it is rejected, and put
$d^\times_{\mathfrak z}=d_{\mathfrak z}-d^{\rm k}_{\mathfrak z}$.
Define
\begin{align}
 \widehat D_{\mathfrak k}
    &=\sum_{\mathfrak z\ {\rm kept\ on}\ \mathfrak k}
                        p_{\mathfrak k}d_{\mathfrak z}p_{\mathfrak k},
 &\widehat\Upsilon_{\mathfrak k}
    &=\tfrac12(\widehat D_{\mathfrak k}
                         +\widehat D_{\mathfrak k}^{\mathsf T}),
                                                    \label{III-112-20}\\
 \widehat\varphi_{\mathfrak k,X}(f;P)
    &=-\tfrac12\sum_{\substack{\mathfrak z\ {\rm kept\ on}\ \mathfrak k\\
                                  X(\mathfrak z)=X}}
              \langle f,p_{\mathfrak k}d_{\mathfrak z}(P)p_{\mathfrak k}f\rangle,
                                                    \label{III-112-21}\\
 \widehat B^\times(X;P,f)
    &=-\tfrac12\sum_{X(\mathfrak z)=X}
                           \langle f,d^\times_{\mathfrak z}(P)f\rangle.
                                                    \label{eq:source-rejected-definition}
\end{align}
Set $c_{\rm lat}(s)=3\sum_{z\in\mathbb Z^3}e^{-s|z|_\infty}$ and
\begin{equation}\label{III-112-22}
 \begin{aligned}
 c_\upsilon&=c_{\rm lat}(\delta_\upsilon/4),
 &B^\sharp_{\rm ch}&=16c_{\rm lat}(\delta_\upsilon/8)^2B_{\rm f}^3,\\
 C_\Upsilon&=c_\upsilon^2B_{\rm f}
                    c_{\rm lat}(3\delta_\upsilon/8)B^\sharp_{\rm ch},
 &K^\sharp_{\rm ch}&=C_3(LM)^6e^{C_3\kappa'}B^\sharp_{\rm ch},
 \end{aligned}
\end{equation}
where $C_2=342$, $C_3=9e^8$ and $\kappa'>0$.

\begin{theorem}[Localization of a quadratic source]\label{III-112-28}
Assume the standing geometric bounds $W_t\ge LM\ge13^7$.
The preceding series converge absolutely and locally uniformly on
their domains. Their coefficients are holomorphic, the matrices
$\widehat\Upsilon_{\mathfrak k}$ are symmetric, and
\begin{equation}\label{III-112-29}\equationalias{III-112-30}
 \sum_{\mathfrak z}m_{\mathfrak z}(b,b')
 \le c_\upsilon^2B_{\rm f}\beta_\upsilon B^\sharp_{\rm ch}
                        e^{-3\delta_\upsilon|b-b'|/8},
 \qquad
 \|\widehat\Upsilon_{\mathfrak k}\|\le C_\Upsilon\beta_\upsilon.
\end{equation}
If
\begin{equation}\label{III-112-23}
 \tfrac12\delta_\upsilon LM\ge C_2(\kappa'+1),
\end{equation}
then, for $\sup|f|\le b_e$ and $n=\#_{\pi_t}X$,
\begin{equation}\label{III-112-32}
 \begin{split}
 |\widehat\varphi_{\mathfrak k,X}(f;P)|
    &\le\tfrac12c_\upsilon^2B_{\rm f}\beta_\upsilon
               K^\sharp_{\rm ch}b_e^2e^{-\kappa'n},\\
 |\widehat B^\times(X;P,f)|
    &\le\tfrac12c_\upsilon^2B_{\rm f}\beta_\upsilon
          C_3B^\sharp_{\rm ch}b_e^2(LM)^6
                    e^{-\delta_\upsilon W_t/8}e^{-\kappa'n}.
 \end{split}
\end{equation}
These are quadratic polynomials in the external coordinates in $X$.
A kept term is zero unless $\mathsf{kp}_t(X)=1$ and
$\mathfrak k=\mathfrak k_X$. A rejected term is zero unless
$d_t(X)\ge\check R_t-2$. Either kind requires a big cube of $X$
in $\Omega_t\cap Z_t$.

On the whole torus put $\Upsilon=\sum_Y\upsilon_Y$ and
$F=(\mathsf M f,f)$. Then
\begin{equation}\label{III-112-33}
 -\tfrac12\langle F,\Upsilon F\rangle
   =\sum_{\mathfrak k,X}\widehat\varphi_{\mathfrak k,X}(f;P|_X)
        +\sum_X\widehat B^\times(X;P|_X,f),
\end{equation}
and
\begin{equation}\label{III-112-31}
 \sum_X\widehat\varphi_{\mathfrak k,X}(f;P|_X)
       =-\tfrac12\langle f,\widehat\Upsilon_{\mathfrak k}(P)f\rangle.
\end{equation}
The corrections are linear in the source vertices and depend only
on vertices with hull contained in their displayed domain. They
commute with compatible gauge transformations and conjugation.
\end{theorem}

\begin{proof}
We give the majorant argument, including the loss in the decay rate.
For two local families, multiply their majorants and sum an
intermediate bond. Spending $\delta-\delta'$ of the decay on that
bond gives a product family with bound
$c_{\rm lat}(\delta-\delta')BB'$ at rate $\delta'$.
The live coefficient trees meet at the intermediate bond, so their
union is connected, has the sum of their lengths, and retains the
larger slack. In particular the composite family
\[
 -t_{\jmath_1}p_{\rm I}\upsilon_Yp_{\rm I}t_{\jmath_2}
\]
has bound $c_\upsilon^2B_{\rm f}^2\beta_\upsilon$ at rate
$\delta_\upsilon/2$.

Set $\epsilon_0=(c_\upsilon^2B_{\rm f}\beta_\upsilon)^{-1}$.
For $R>1$, adjoin $\epsilon\upsilon_Y$ to the legs
$\lambda a_S$, and adjoin
$-\epsilon t_{\jmath_1}p_{\rm I}\upsilon_Yp_{\rm I}t_{\jmath_2}$
to the inverse terms $\lambda t_\jmath$, with
$|\lambda|<R$, $|\epsilon|<R\epsilon_0$. Each enlarged family
has bound $2RB_{\rm f}$ at rate $\delta_\upsilon/2$.
The three-factor estimate for its chains is therefore
$R^3B^\sharp_{\rm ch}$ at rate $3\delta_\upsilon/8$.
The chains containing exactly one adjoined member are precisely
the four kinds in \eqref{III-112-18}. Their majorants are
$R^{1+r}\epsilon_0m_{\mathfrak z}$, with $r=0$ or $2$.
Retaining only those nonnegative terms and letting $R\downarrow1$
proves the first inequality in \eqref{III-112-29}.

The conversion of connecting length to hull size uses the geometry
in the proof of Proposition~\ref{II-64-65}. We give the numerical
calculation needed at the present rate. A live footprint of length $l$ and slack
$\rho_\upsilon$, with $4\rho_\upsilon\le LM$, has a hull with
$n\le8+64l/(LM)$ cubes. A fixed hull has at most
$9(LM)^6n^2$ ordered pairs of positive bonds. A rejected entry has
$l\ge W_t-2-2\rho_\upsilon$. Applying these bounds to the enlarged
families, at rate $\delta_\upsilon/2$, gives under
\eqref{III-112-23}
\[
 \begin{split}
 \epsilon_0\sum_{b,b'}\sum_{X(\mathfrak z)=X}m_{\mathfrak z}(b,b')
       &\le K^\sharp_{\rm ch}e^{-\kappa'n},\\
 \epsilon_0\sum_{b,b'}\sum_{\substack{X(\mathfrak z)=X\\
                 l_{\mathfrak z}(b,b')\ge W_t-2-2\rho_\upsilon}}
           m_{\mathfrak z}(b,b')
       &\le C_3B^\sharp_{\rm ch}(LM)^6
                         e^{-\delta_\upsilon W_t/8}e^{-\kappa'n}.
 \end{split}
\]
Indeed, put $\delta=\delta_\upsilon/2$,
$l_1=\max\{0,LM(n-8)/64\}$ and $l_\times=W_t-2-2\rho_\upsilon$.
Before summing the bond pairs, the two bounds are respectively
$B^\sharp_{\rm ch}e^{-3\delta l_1/4}$ and
$B^\sharp_{\rm ch}e^{-3\delta\max(l_1,l_\times)/4}$.
The hypothesis $\delta LM\ge342(\kappa'+1)$ gives
\[
 \tfrac34\delta l_1\ge(\kappa'+1)(n-8).
\]
Since $n^2e^{-n}<1$, the first total is at most
$9e^8(LM)^6e^{8\kappa'}B^\sharp_{\rm ch}e^{-\kappa'n}$.
For the second, use $\max(l_1,l_\times)\ge l_1/4+3l_\times/4$,
$2\rho_\upsilon\le LM/2\le W_t/2$, and $W_t\ge LM\ge13^7$:
\[
 \tfrac34\delta\max(l_1,l_\times)
 \ge\delta\left(\tfrac9{16}l_\times-\tfrac3{128}LM\right)
                         +\tfrac3{1024}\delta LMn
 \ge\tfrac14\delta W_t+(\kappa'+1)n.
\]
The second total is therefore at most
$9(LM)^6B^\sharp_{\rm ch}e^{-\delta W_t/4}e^{-\kappa'n}$.
These are the displayed bounds, since $C_3=9e^8\ge8$.
Pairing with $f$ gives
\eqref{III-112-32}. The first estimate in \eqref{III-112-29},
summed over one bond index, gives the norm bound by the Schur
estimate. Compression and symmetrization do not increase it.

Every coefficient in a product is holomorphic on its own hull.
The uniform majorants prove local uniform convergence. Both live
endpoint bonds lie in the footprint, so the sums are quadratic
polynomials supported in $X$. The keeping rule gives the stated
support properties and the unique component of a kept term.

Absolute convergence permits multiplication and regrouping of the
three families. Summing the four kinds gives
\[
 \Upsilon_{\rm EE}-\Upsilon_{\rm EI}TA_{\rm IE}
       -A_{\rm EI}T\Upsilon_{\rm IE}
       +A_{\rm EI}T\Upsilon_{\rm II}TA_{\rm IE}.
\]
Since $T=A_{\rm II}^{-1}$ and $\mathsf M=-TA_{\rm IE}$, this is the
matrix of $\langle(\mathsf M f,f),\Upsilon(\mathsf M f,f)\rangle$. Splitting its
terms into kept and rejected parts proves \eqref{III-112-33};
symmetrization does not change a quadratic form and gives
\eqref{III-112-31}. Each product has one source vertex and uses
only restrictions to its hull. This proves linearity and locality.
Compatible gauge transformations commute with the coordinate
projections, and conjugation commutes with the products and the
absolutely convergent sums, proving the two symmetry assertions.
\end{proof}

\subsection{Triples, pairs, and retained factors}

We now use the graded and layered triple domains of a history
step. Write $\mathscr U^{\rm m}_X$ for the graded domain and
$\mathscr U^{\rm o}_X$ for the layered output domain. The numerical
domain hypotheses needed below are the following explicit
inclusions and chart properties:
\begin{enumerate}[label=(\roman*)]
\item $\mathscr U^{\rm o}_X\subset\mathscr U^{\rm m}_X$;
      every fixed-real-factor slice of $\mathscr U^{\rm m}_X$ is
      open, and restriction to a smaller domain preserves it.
\item If $P\in\mathscr U^{\rm m}_X$, its inner pair
      $p_X(P)=(e^{i\eta_tA'}U_{\rm r},L^{-3}\mathbf J)
       |_{{\rm Bd}^\circ(X\cap\Omega_0^\diamond)}$
      lies in $\mathscr O^+(X;\alpha)$, and
      $P|_Y\in\mathbb U(Y)$ for every $Y$ whose level-$t$ hull
      lies in $X$.
\item In these slices, $\eta_t|A'(b)|<\log2$ on the inner bonds.
\end{enumerate}
They are imposed at each step, with a single table of radii and
its truncations. These hypotheses distinguish membership in the
analytic domains from the later question of evaluating a history
measure. No independence of an evaluation representative is
assumed here.

\begin{proposition}\label{III-113-12}
Under these domain hypotheses, the source corrections are defined
on $\mathscr U^{\rm m}_X$ and their restrictions to
$\mathscr U^{\rm o}_X$ are analytic functions with stored data.
For $\mathsf{kp}_t(X)=1$, put
\begin{equation}\label{III-113-11}
 \widehat\varphi_X(P;\mathfrak q)
  =\widehat\varphi_{\mathfrak k_X,X}
                       (\operatorname{vec}a[\mathfrak q];P),
\end{equation}
where $a[\mathfrak q]$ is the stored external field, extended by
zero away from $X$. If the stored tiers obey
$\sup|\operatorname{vec}a[\mathfrak q]|<C_{\rm xt}\mathfrak b_k$,
then under \eqref{III-112-23}
\begin{equation}\label{III-113-15}
 \operatorname{osc}\widehat\varphi_X
 \le2\sup|\widehat\varphi_X|
 \le b^{\rm S}_{\rm fr}\mathfrak b_k^2e^{-\kappa'd_t(X)},
 \quad
 b^{\rm S}_{\rm fr}
   =c_\upsilon^2B_{\rm f}\beta_\upsilon K^\sharp_{\rm ch}
                           C_{\rm xt}^2e^{-\kappa'}.
\end{equation}
The corrections inherit pair dependence, gauge covariance, Borel
measurability and conjugation whenever the local forms have the
corresponding properties. They are linear and local in these
forms.
\end{proposition}

\begin{proof}
Extend the real factor of a triple on $X$ to a real configuration
$U$ on the torus. The open slice containing the triple lies in
$\mathscr P_U(X)$: this follows directly from (i), (ii) and the
definition as an interior. Thus \cref{III-112-42,III-112-28} apply.
Every chain uses only the restricted triple on its hull, so its
value is independent of the extension $U$. The uniform majorants
prove holomorphy on the open slice. A stored field enters through
finitely many coordinates; hence the restriction is analytic with
stored data. Substitute $b_e=C_{\rm xt}\mathfrak b_k$ into
\eqref{III-112-32} and use
$\#_{\pi_t}X\ge d_t(X)+1$ to obtain \eqref{III-113-15}.
The same argument works on the factor domain in every later
history, since the domains, radii and stored coordinates at this
step are unchanged. Each coefficient is a finite product of
coefficients with the asserted dependence or symmetry. The
convergent sums preserve those properties; for measurability use
their Borel partial sums.
\end{proof}

If the forms depend only on the inner pair, so do the corrections.
This assertion alone does not establish analyticity as a function
of that pair. The following local chart supplies it.

\begin{lemma}\label{III-113-35}
Let $c^\sharp$ be holomorphic on every fixed-real-factor slice of
$\mathscr U^{\rm m}_X$ and constant on triples having the same
inner pair. It descends to a function on
\begin{equation}\label{III-113-36}
 \mathscr O^{\rm S}(X)=\{p_X(P):P\in\mathscr U^{\rm m}_X\},
\end{equation}
which admits a holomorphic extension near each point of this set.
\end{lemma}

\begin{proof}
Choose $P_0=(U_0,A'_0,\mathbf J_0)$ over a pair $p_0$.
On the inner bonds define, for pairs $(\mathbf U,\mathbf j)$ near
$p_0$,
\[
 A'(b)=(i\eta_t)^{-1}\log(\mathbf U(b)U_0(b)^{-1}),
 \qquad \mathbf J(b)=L^3\mathbf j(b),
\]
and keep the other coordinates equal to those of $P_0$.
By (iii) the series logarithm is defined near $p_0$ and recovers
$A'_0$ there. This gives a holomorphic local section $\Theta$ of
$p_X$. Shrink its neighborhood so that $\Theta$ remains in the
open slice of $P_0$. Then $c^\sharp\circ\Theta$ is the required
extension; constancy on the fibers proves agreement with the
descended function. The empty inner-bond set gives a single pair
and a constant function.
\end{proof}

\begin{corollary}\label{III-113-39}
Suppose the local forms depend only on the inner pair. Then the
functions
\begin{equation}\label{III-113-38}
 \widehat K(X;p,a)=\sum_{\mathfrak k}
       \widehat\varphi_{\mathfrak k,X}(\operatorname{vec}a;P),
 \qquad
 \widehat B^\times(X;p,a)
       =\widehat B^\times(X;P,\operatorname{vec}a),
 \quad p=p_X(P),
\end{equation}
are well defined on $\mathscr O^{\rm S}(X)$ and are finite
quadratic polynomials in the external coordinates in $X$, with
coefficients analytic on that set. They have the locality,
support and symmetry properties of \cref{III-112-28}. Under
\eqref{III-112-23}, for $\sup|a|\le b_e$,
\begin{equation}\label{III-113-40}
 \begin{split}
 |\widehat K(X;p,a)|
    &\le c^{\rm S}_{\rm kp}b_e^2e^{-\kappa'd_t(X)},\\
 |\widehat B^\times(X;p,a)|
    &\le c^{\rm S}_{\rm rj}b_e^2
                       e^{-(\check R_t-2)}e^{-\kappa'd_t(X)},\\
 c^{\rm S}_{\rm kp}
    &=\tfrac12c_\upsilon^2B_{\rm f}\beta_\upsilon
                                      K^\sharp_{\rm ch}e^{-\kappa'},\\
 c^{\rm S}_{\rm rj}
    &=\tfrac12c_\upsilon^2B_{\rm f}\beta_\upsilon
                          C_3B^\sharp_{\rm ch}(LM)^6e^{-\kappa'}.
 \end{split}
\end{equation}
\end{corollary}

\begin{proof}
Apply \cref{III-113-35} to the finitely many quadratic coefficients
given by \cref{III-113-12}. For the bound use
$|\operatorname{vec}a|\le|a|$,
$\#_{\pi_t}X\ge d_t(X)+1$, and
$\delta_\upsilon LM/8\ge1$, which follows from
\eqref{III-112-23}. Since $W_t=LM\check R_t$,
$e^{-\delta_\upsilon W_t/8}\le e^{-(\check R_t-2)}$.
Only the unique component $\mathfrak k_X$ contributes to the
first sum.
\end{proof}

For a sequence of histories, add \eqref{III-113-11} to every
retained quadratic exponent and leave the elimination exponents
unchanged. Thus
\begin{equation}\label{III-113-24}
 \varphi_f^\upsilon=
 \begin{cases}
 \varphi_f+\widehat\varphi_{X_f},&f\text{ is a retained quadratic factor},\\
 \varphi_f,&f\text{ is an elimination factor}.
 \end{cases}
\end{equation}

\begin{corollary}\label{III-113-27}
Assume that $h$ has positive level and that the preceding domain
hypotheses hold at all its steps, with
one table of radii and its truncations. Suppose the unsourced
factor family has oscillation bound
$B^{\rm z}\mathfrak b_{t-1}^2e^{-\kappa d_t(X)}$ at step $t$,
where $0<\kappa\le\kappa'$. Then the family in
\eqref{III-113-24} has the same names, domains and factor spaces,
and oscillation bound
\[
 (B^{\rm z}+\bar b^{\rm S}_{\rm fr})
                   \mathfrak b_{t-1}^2e^{-\kappa d_t(X)},
 \qquad
 \bar b^{\rm S}_{\rm fr}=\max_{t\le\operatorname{level}(h)}
                                      b^{{\rm S},(t)}_{\rm fr}.
\]
It preserves prefix consistency and the analytic, Borel, pair and
gauge properties verified above. For quadratic parts linear in a
plaquette weight, the added exponent is linear in that weight and
depends only on plaquettes with all corners in $X_f$.
\end{corollary}

\begin{proof}
At a fixed step the layered space and stored tiers are unchanged
under extension of the history. The definitions therefore give
the same partial functions before and after that extension.
\Cref{III-113-12} supplies their values and analytic domains.
The triangle inequality for oscillation, \eqref{III-113-15} and
$\kappa\le\kappa'$ give the bound. Every remaining assertion
holds factor by factor; the elimination factors were unchanged.
\end{proof}

\subsection{Damping on defined history values}

The last estimate of this section concerns actual evaluations.
Let $h$ have level $n\ge1$, and let
$U\in\mathcal V_0[h]\cap\operatorname{dom}\xi_h^{\rm co\prime}$.
Let $f_t=\operatorname{vec}\mathfrak q[U]_{t-1}$ be its stored real
field, extended by zero away from the external bonds at step $t$.
Assume the fixed-radius and kept-positivity hypotheses of
Section~\ref{subsec:ymm1-65-2}, and assume
that an admitted evaluation representative $\mathsf U^\flat$
places every factor triple in its layered domain. The stored-form calculation in
equations \eqref{eq:history-piece-kept-form}--\eqref{II-195-47} gives real symmetric kept matrices and
\begin{equation}\label{eq:unsourced-piece-lower-for-source}
 \sum_{t=1}^n\sum_{\mathfrak k\in\mathfrak K_t(h)}
      \langle f_t,\Delta^{\rm loc}_{\mathfrak k}(p_t)f_t\rangle
       \ge2\gamma_{\rm kp}\mathsf n(U),
 \qquad
 \mathsf n(U)=\sum_{t=1}^n\|f_t\|^2.
\end{equation}
The external bond sets $E^+_{\mathfrak k}$ partition the stored
bonds piece by piece, so this is also the sum of the squared norms
of $f_t|_{E^+_{\mathfrak k}}$ over $t$ and $\mathfrak k$.
The determinant factors of the elimination steps have the same
value on every glued configuration of the fields of $U$ by
Lemma~\ref{II-195-31}.

\begin{proposition}\label{III-113-44}
Assume the hypotheses just stated and the quadratic construction
at each step, and put
\[
 \bar\upsilon=C_\Upsilon\max_{1\le t\le n}\beta_\upsilon^{(t)}.
\]
At every admitted evaluation representative, the product of the
modified factors is bounded by
\begin{equation}\label{III-113-48}
 \left|e^{\Phi_h^\upsilon(\mathsf U^\flat;\mathfrak q[U])}\right|
 \le
 \prod_{t=1}^n\prod_{c\in\Gamma_t^\partial(h)}
       |\det(Lm_c(\mathcal M^{t-1}\mathsf U))|^{-1}
 \exp\left[-\left(\gamma_{\rm kp}-\tfrac12\bar\upsilon\right)
                                                    \mathsf n(U)\right],
\end{equation}
where $\mathsf U$ is any glued configuration of the fields of $U$.
The same estimate holds for $\Phi_h^\upsilon(\phi)$ whenever that
value is defined, meaning that the exponent sum is defined and
constant on the evaluation set of $\phi$.
\end{proposition}

\begin{proof}
For a component $\mathfrak k$ at step $t$, sum the source chains
at the triple supplied by the evaluation representative on each
individual hull:
\begin{equation}\label{III-113-45}
 \widehat D^{\rm ev}_{\mathfrak k}(\mathsf U^\flat)
    =\sum_{\mathfrak z\ {\rm kept\ on}\ \mathfrak k}
       p_{\mathfrak k}d_{\mathfrak z}\bigl(P_{X(\mathfrak z)}
                                      (\mathsf U^\flat)\bigr)p_{\mathfrak k}.
\end{equation}
These are the compressions in \eqref{III-112-21};
$p_{\mathfrak k}$ projects onto $E^+_{\mathfrak k}$.
Thus the matrix is supported on
$E^+_{\mathfrak k}\times E^+_{\mathfrak k}$.
Every displayed value is defined by the domain hypothesis. The
majorants in \eqref{III-112-29} are independent of the triple, so
the series converges absolutely and its row and column sums give
\[
 \|\widehat D^{\rm ev}_{\mathfrak k}(\mathsf U^\flat)\|
                  \le C_\Upsilon\beta_\upsilon^{(t)}\le\bar\upsilon.
\]
This argument does not evaluate all chains at one unverified
larger-domain triple. The live endpoint bonds of each chain lie
inside its hull, so the stored field restricted to that hull has
the same values as $f_t$ in its pairing. Absolute convergence
therefore gives the exact product
\[
 e^{\Phi_h^\upsilon}
 =\prod_{t,c}\det(Lm_c(\mathcal M^{t-1}\mathsf U))^{-1}
   \exp\left[-\tfrac12\sum_{t,\mathfrak k}
       \langle f_t,(\Delta^{\rm loc}_{\mathfrak k}(p_t)
                    +\widehat D^{\rm ev}_{\mathfrak k})f_t\rangle\right].
\]
The component support and the norm bound give
\[
 \bigl|\langle f_t,\widehat D^{\rm ev}_{\mathfrak k}f_t\rangle\bigr|
 \le\bar\upsilon\|f_t|_{E^+_{\mathfrak k}}\|^2.
\]
Summing over the partition of the stored bonds shows that
the real part of the added sum is at least
$-\bar\upsilon\mathsf n(U)$. Combining this with
\eqref{eq:unsourced-piece-lower-for-source} and taking the modulus
proves \eqref{III-113-48}. If the value at $\phi$ is defined,
choose any admitted representative and apply that bound to its
common value. No constancy on an evaluation set follows from
the norm estimate itself.
\end{proof}


\subsection{The Hessian at an actual history}

\begin{lemma}\label{III-19-18}\label{III-19-31}
At an admitted real chain point of step $t=k+1$, let $U_+$ be the
collared minimal background in the required axial gauge and let
$J_+=J^\circ(U_+)$. Put $p=\mathcal M^kU_+$.
The quadratic operator in the actual conditional
Gaussian is
$A=\widehat G+\widetilde A$, where
$\widetilde A=C_p^{\mathsf T}\Delta_k(U_+,J_+)C_p$ is the Hessian of
the regional constrained action. If
$\overline m=\mathbf m\oplus\mathsf a$ is its conditional mean, then
\[
 \chi_{\rm I}\widehat G\,\overline m=0,
 \qquad \chi_{\rm I}\widetilde A\,\overline m=0.
\]
The unit-weight local potentials satisfy
\[
 \sum_Y T_1V_Y^{\rm act}(B)=0,\qquad
 \sum_Y T_{1,1}V_Y^{\rm act}(B_1,B_2)
             =(B_1,\widetilde A B_2)_1
\]
for complex directions as well as real ones. These assertions include
the first step.
\end{lemma}

\begin{proof}
The shorter determining set at this point is the reduced determining
set of the collared sequence before the step. Adding
$\Omega_t^{(t)}$ gives the longer sequence. The actual background is
the minimal orbit for this longer sequence; choosing its axial
representative does not change the local action values.
The real background and shift construction gives its strict
$2B_3\varepsilon_t$ regularity and its natural-class representation
with zero perturbation. We verify the nearby-orbit and stationarity
requirements of \cref{II-27-28} separately.
The shorter datum $d(p)$ is regular at
$2L^2C_{\rm reg}\varepsilon_t\le\bar\varepsilon$.
The separated shorter sequence satisfies the geometric hypotheses
of Theorem~\ref{I-58-22}, with the radii fixed by the real choices.
For a differentiable fibre curve $V(s)$ through $p$, set
\[
 B(s)_k(b)=-i\log(V(s)(b)p(b)^{-1}),\qquad
 B(s)_j=0\quad(j<k).
\]
On a sufficiently short interval the logarithm is defined,
$B(s)$ lies in the selector's row ball, and
$e^{iB(s)}d(p)=d(V(s))$. The real analytic selector of that theorem
therefore gives a differentiable representative $U_s$ of the shorter
regular orbit. Its constraint set is contained in the fixed longer
constraint set. At $s=0$ this orbit is contained in the longer
critical orbit, so differentiation of the action along $U_s$
gives zero. The action has the same value on every representative
of the shorter orbit. This proves the required stationarity.
When $k=0$ the shorter orbit is a singleton; its existence near
$p$ follows from the finitely many strict regularity inequalities
and continuity. The same longer-constraint argument gives
stationarity in that case.
The regional criticality and regularity conclusions of
Proposition~\ref{II-93-38} now apply.
The current is the regional action current; its extension at the
boundary is harmless by Lemma~\ref{II-93-23}.
Thus the regional calculation in Proposition~\ref{II-113-38},
restated in \cref{II-27-28}, applies at this point.

The integration bonds are exactly the nonpivot bonds in
$B(\Lambda_t^{(t)})$; the complementary nonpivot bonds are external.
These are also the two blocks used by the conditional Gaussian.
Its internal matrix is invertible by the real conditional Gaussian
gap. The gauge-degeneracy identity for the regional Hessian therefore
permits \cref{III-19-8} to be applied to this very block decomposition.
It gives the two separate mean identities displayed above.

By Proposition~\ref{II-106-49}, the finite sum of local action
potentials along a sufficiently short real fibre ray is the change
of the regional constrained action. The first derivative vanishes
by criticality, and the second derivative is the Hessian just
identified. Each local potential is holomorphic near zero and there
are finitely many of them. Comparing coefficients along real rays
therefore gives the two identities for small real directions.
Homogeneity extends them to every real direction, and equality of
the resulting complex polynomials extends them to complex directions.
Polarization yields the displayed bilinear identity. The same
argument at $k=0$ uses the first-step local action potential and its
radius $r_P^\circ$; its unit-weight potential is the actual action
potential. Thus the mean cancellation used at the first step has
the same Hessian and the same internal block.
\end{proof}


\section{The source-dependent interactions}
\label{sec:sourced-recursion}

Fix the real construction of Theorem \ref{II-203-9}, including its
couplings, histories, collars, tables of radii, test functions and
coordinate maps. These data are independent of the source. Throughout
this section the source budgets of \cref{III-105-3} are imposed.
We use them through the bounds displayed below; their simultaneous
verification is the subject of \cref{III-105-7}.
Write $t=k+1$ and $h'=(h,\mathsf o_t)$ for a history step.
Unless stated otherwise, $\Omega _1(h')\ne\emptyset$.

The regular terms are functions on the fixed cores
$\mathcal U_t^\circ(X;\alpha _0,\alpha _1)$.
The boundary terms are defined on
\[
 \mathfrak G_X(h')=\mathscr U_X(h')\times
                       \mathfrak Q^{\rm rad}_{h'}(X),
\]
where $\mathscr U_X$ is the graded space of triples of the real
construction and the second factor imposes the stored-field tiers.
A retained exponent formed at level $j=l+1$ has domain
\begin{equation}\label{eq:sourced-retained-domain}
 \mathscr P_{X_f}(h'|_j)=
 \mathsf{Sp}(\mathcal S(h'|_j);(j,X_f,{\rm mk}))
                  \times\mathfrak Q^{(l),\rm rad}_{h'|_j}(X_f).
\end{equation}
Its name, support and domain are those of the corresponding real
factor. Restriction to these domains is part of the construction.
We make no assertion about maximal analytic continuation beyond them.

\subsection{One integration and the stored terms}

Let $\mathscr L_k^{w\flat}$ contain the regular, boundary and retained
entries already constructed through level $k$. Its counterterms are
the real counterterms $b_1,\ldots,b_k$. At step $t$ use the centered
source family and its one-lattice companion from
\cref{III-39-43,III-41-24}; at $t=1$ use their first-step versions
\cref{III-48-54}. The quadratic terms of
\cref{III-113-39} are added to the real conditional forms:
\begin{equation}\label{eq:sourced-two-forms}
 K^{\rm kp,w}=K^{\rm kp}+K^{\rm kp,+},\qquad
 B^{\times,w}=B^\times+B^{\times,+}.
\end{equation}
Here the first correction vanishes unless the domain is retained.
The conditional Gaussian, its covariance and its mean are those of
the unsourced construction. Their domains and positivity are unchanged.

We denote the centered cluster function by
$E_h^{{\rm cl},w}(X;\mathsf a;P;\mathfrak q)$, the pointed term
of the empty history by $F_\emptyset^w(X,P,z)$, and the localized
normalizer by $N_\emptyset^{\rm un}(X,p_X(P),z)$.
The new regular term is
\begin{equation}\label{III-107-7}
 E^{(t),w}(X,P,z)=N_\emptyset^{\rm un}(X,p_X(P),z)
              +F_\emptyset^w(X,P,z)-n_{\rm tr}^{(t)}(X,z).
\end{equation}
Equivalently, $E^{(t),w}=\mathbf N^{(t)}+F_\emptyset^w$;
$\mathbf N^{(t)}$ includes the displayed subtraction and is
independent of $w$. Only the old regular entries, the counterterms
and the one-lattice companion enter $F_\emptyset^w$. In particular
\eqref{III-107-7} does not depend on boundary entries or on a choice
of strip representation.

For the boundary contribution use the five localized strip terms
$R^{6,w},R^{7,w},R^{8,w},R^{9,w},R^{{\rm E},w}$ of
Section \ref{sec:boundary-assembly}, inserting $K^{\rm kp,w}$,
$B^{\times,w}$ and $E_h^{{\rm cl},w}$ in the last three places.
The first two, which come from the strip Jacobian and Haar density,
are unchanged. Set
\begin{equation}\label{eq:sourced-created-boundary}
 \begin{split}
 B^{\rm nm}(X)&=N_h^\partial(X;p_X(P))
  -\mathbf1_{\{X\subset\Lambda_t\}}
                           \sum_zN_\emptyset^{\rm un}(X,p_X(P),z),\\
 B^{{\rm cs},w}(X)&=E_h^{{\rm cl},w}(X;\mathsf a;P;\mathfrak q)
  -\mathbf1_{\{X\subset\Lambda_t\}}\sum_zF_\emptyset^w(X,P,z),\\
 B^{{\rm rj},w}(X)&=B^{\times,w}(X;p_X(P),\mathsf a),\\
 B^{{\rm sr},w}(X)&=R^{6,w}(X)+R^{7,w}(X)+R^{8,w}(X)
                                  +R^{9,w}(X)+R^{{\rm E},w}(X),\\
 B_t^{{\rm mk},w}(X)&=B^{\rm nm}(X)+B^{{\rm cs},w}(X)
                                  +B^{{\rm rj},w}(X)+B^{{\rm sr},w}(X).
 \end{split}
\end{equation}
As usual an expression multiplied by a zero indicator is omitted.
Thus its domain need not be imposed when it is unused.

The new retained exponents are
\begin{equation}\label{III-101-23}
 \varphi_f^w=\begin{cases}
  \varphi_f+\widehat\varphi_{X_f}^{\widehat w_k},&f\text{ of kind }{\rm kp},\\
  \varphi_f,&f\text{ of kind }{\rm el},
 \end{cases}
 \qquad \widehat w_k(p)=L^{K-k}w(x(p)).
\end{equation}
All old entries are prolonged by ignoring the newly stored
coordinates. We restrict the three newly formed families to their
core, graded-tier and retained domains, respectively. Starting with
$\mathscr L_0^{w\flat}=(x_0)$, this defines
\begin{equation}\label{III-102-8}
 \mathscr L_t^{w\flat}
 =\left.\operatorname{Step}_t^+
          [\mathscr L_k^{w\flat};\mathsf{Ss}^w]\right|_{
           \text{cores, graded tiers and retained domains}} .
\end{equation}
This notation refers to the explicit operations
\eqref{eq:sourced-two-forms}--\eqref{III-101-23} and the centered
expansion; it does not presume that the output has the required
domains or bounds.

\begin{proposition}\label{III-102-21}
The recursion \eqref{III-102-8} is defined at every level. Its
regular entries are defined on the whole fixed core, its created
boundary entries on every $\mathfrak G_X$, and its retained
exponents on \eqref{eq:sourced-retained-domain}. On these domains
the entries are local, Borel measurable, analytic on each complex
slice, and gauge invariant. Their dependence on the source is
holomorphic on finite-dimensional source slices and is local to
the indicated domain. Conjugation takes the entry at $w$ to the
entry at $\bar w$ and conjugates its arguments.

With the notation of the source budgets,
\begin{equation}\label{eq:sourced-recursion-bounds}
 \|E^{(j),w}\|_{j,\kappa+1}\le E_0^{\rm S},\qquad
 |B_j^{{\rm mk},w}(X)|\le B_T^{\rm th}e^{-\kappa_Td_j(X)}
                         \le B^{\rm S}e^{-\kappa d_j(X)},
\end{equation}
where $\kappa_T=(1+4\beta_{\rm s})\kappa$.
The boundary term is zero on $X\subset\Lambda_j$ and off the
boundary index set. The retained exponents have the same names
and domains as at zero source, and their oscillation bound is
at most $B_{\rm fr}^{\flat}+\bar b_{\rm fr}^{\rm S}\le B_{\rm fr}^{\rm S}$.
At $w=0$ all these functions are restrictions of the canonical
real functions. The construction commutes with cutting a history
and the data at an earlier level.
\end{proposition}

\begin{proof}
We first explain why complex entries are admissible in the
potential construction of Sections \ref{sec:history-actual-potentials}--\ref{sec:history-cluster}.
Keep all geometric data and couplings fixed and compare the list
with the null list, whose interaction entries vanish.
The dependence of the inherited regular potential on the regular
entries is complex linear: its operations are restriction,
subtraction of the constant and quadratic parts, finite
localization, and sums. The inherited boundary and retained
potentials have the same property in their respective entries.
The action, gauge and Jacobian potentials are independent of
these entries. Thus the exact finite action identity for the
null list, plus the three linear differences, is the identity
for the complex list. The estimates of
Proposition \ref{II-108-24} use absolute values and Cauchy
bounds on the same domains. They apply to these differences
without a reality assumption. The conjugation identity relates
the conjugate lists, rather than declaring their values real.

Suppose the entries through level $k$ have been constructed.
Their bounds put the old list in the required box. In particular
the edge contribution is bounded by $8M^3E_0^{\rm S}$, while
the boundary and retained boxes are large enough because
\[
 B^{\rm S}\ge\max(8M^3E_0^{\rm S},B_T^{\rm th}),\qquad
 B_{\rm fr}^{\flat}+\bar b_{\rm fr}^{\rm S}\le B_{\rm fr}^{\rm S}.
\]
The ordinary potential bounds therefore hold on the full graded
domains and stored-field tiers. The weighted potentials have
these domains by \cref{III-39-43,III-111-7}; the budgets in
\cref{III-105-3} put the sum of ordinary and centered potentials
within the size allowed by \cref{III-32-12}.
The latter theorem supplies the cluster function on its full
external-field perturbation domain, including the tier values.
It also supplies its analytic and locality properties and the
interior identity with the empty history. Consequently the
right side of \eqref{III-107-7} is defined on the whole core.
The regular bound is verified below independently of the
boundary entries.

The quadratic corrections are polynomials in the external field,
with analytic coefficients on the pair image $\mathscr O^{\rm S}(X)$, by
\cref{III-113-39}. They involve only external bonds in $X$.
We use the following complex version of the boundary argument
of Proposition \ref{II-151-38}. Keep the geometric data and
domains fixed. For an ingredient $F$ define
$F^{\rm cj}(P;\mathfrak q)=\overline{F(P^{\mathfrak c};\mathfrak q)}$,
and form $(F+F^{\rm cj})/2$ and $(F-F^{\rm cj})/(2i)$.
The resulting real and imaginary systems satisfy the real
conjugation identities and the same modulus bounds on the same
domains. Also take the system with the normalizer, pointed,
cluster and quadratic ingredients zero and with the same
geometric shift. The proposition applies to these three systems.
The rows $8,9,{\rm E}$ are linear in their respective ingredients;
rows $6,7$ are fixed by the geometry. Consequently, for the full
boundary term and for the strip term,
\begin{equation}\label{III-35-18}
 B=B^{\rm re}+i(B^{\rm im}-B^0),\qquad B^0=R^6+R^7.
\end{equation}
For the other kinds and for rows $8,9,{\rm E}$ the analogous
identity has no $B^0$ term. Thus definedness on the full
graded-tier product, holomorphy, measurability, support and
the joint gauge and pair identities pass from the three real
systems to the complex system. This argument also handles
the fixed affine contribution of the strip density.

To obtain the modulus bounds without enlarging their constants,
apply the finite transfer estimate to the complex ingredients
themselves. Their domains are those of the real and imaginary
systems, and their bounds are the original sourced bounds.
To check the analytic room explicitly, use the step budget
\eqref{eq:common-boundary-budget}. The allowed external perturbation has
radius $r_{{\rm bd},k}=C_{\rm xt}\mathfrak b_k$.
Adding it to a tier field keeps the field in the ball of radius
$2C_{\rm xt}\mathfrak b_k$, on which the two quadratic estimates
hold. The cluster function is defined on that perturbation
domain by the preceding paragraph. The Haar term has
$g_k|\mathsf a+\mathsf e|<3/8$; hence its fixed series logarithm
is analytic there. The localized transfer is a finite sum of
values on precisely these domains. Its modulus estimates give
\[
 |B_t^{{\rm mk},w}(X)|\le
 B_Ne^{-\delta_NM d_t(X)}
 +(b_{\rm cs}+b_{\rm rj})e^{-\kappa''d_t(X)}
 +b_{\rm sr}e^{-\kappa_Td_t(X)}.
\]
Here $b_{\rm rj}$ and $b_{\rm sr}$ are the expressions in
Proposition \ref{II-151-38}, with the enlarged sourced form
constants. The source budgets give
$\delta_NM\ge\kappa_T$, $\kappa''\ge\kappa_T$ and
$b_{\rm cs}+b_{\rm rj}+b_{\rm sr}\le1$; hence the boundary
bound in \eqref{eq:sourced-recursion-bounds} follows.
The same proof gives the support assertions. Inside $\Lambda_t$
the normalizer and cluster subtractions cancel, while no rejected
or strip term occurs. Off the boundary index set each of the four
terms vanishes separately.

The decomposition \eqref{III-35-18} proves the domain and
structural assertions, whereas direct complex transfer proves
the bounds. Replacing the source and all arguments by their
conjugates in each formula proves the required conjugation identity.
Restriction and transfer preserve finite-slice holomorphy and
Borel measurability. For the cluster terms these properties
also follow from the locally uniform majorants in
\cref{III-32-12}. Each source plaquette read in a term over $X$
lies in $X$, so locality in the source follows through the same
operations.

For the new retained factors use \cref{III-113-27} on the full
product \eqref{eq:sourced-retained-domain}. The quadratic
correction is defined there and has oscillation at most
$\bar b_{\rm fr}^{\rm S}$ in the normalized factor bound.
Adding the real oscillation bound gives the stated box. Old
factors are prolonged by the same prefix rule as in the real
construction, so none acquires a new condition on later stored
fields. This proves all assertions at level $t$ once the regular
bound below is inserted.

At level zero the list has no interaction entries. The first
source family and escort of \cref{III-48-54} have the same
properties on the first-step domains; the preceding construction
therefore starts without an induction assumption. At source
zero all added potentials, companions and quadratic forms vanish.
Every remaining operation and every restriction is the real one,
which proves the zero-source identification. Finally the step
at level $i$ uses only data through $i$. All domains are unchanged
when a longer history is cut at $i$, including the filing spaces
and stored tiers. Induction on $i$ gives equality as partial maps
under every cut. This also proves persistence of the bounds for
older boundary entries.
\end{proof}

\subsection{Control of the regular terms}

The estimate needed for the regular recursion depends only on
the difference of two old regular arrays and on the companion.
Let $\mathfrak f,\widetilde{\mathfrak f}$ be these arrays, with the
same counterterms, and put
\begin{equation}\label{III-107-12}
 \theta=K_\Lambda\mathsf s_{\mathfrak a}
                N_{\mathfrak a,\omega}
                   (\mathfrak f-\widetilde{\mathfrak f})
          +\|\mathbb W^{\rm 1l}\|^{\rm co}_{\star,\lambda_V},
 \qquad \mathsf s_{\mathfrak a}=(1+\beta_{\rm s})\delta_k/\alpha_0.
\end{equation}
The weights of this array norm are $L^{-(k-j)}$ for the old
level-$j$ entries.

\begin{lemma}\label{III-107-11}
Let $1\le k<K$. Suppose the two old lists satisfy the regular box bound,
$K_\Lambda\mathsf s_{\mathfrak a}N_\star\le w_\star/4$,
and $\theta\le3w_\star/4$. Then their pointed terms, with
the companion attached to the first list, are defined on the
whole next core and satisfy
\[
 |F_\emptyset[\mathscr L;\mathbb W^{\rm 1l}](X,P,z)
       -F_\emptyset[\widetilde{\mathscr L};0](X,P,z)|
 \le33(3L^3-3)K_\Pi\theta e^{-\lambda_{\rm out}d_t(X)}.
\]
\end{lemma}

\begin{proof}
The regular localization rule is linear in its array. Hence
\[
 \widetilde{\mathbb W}=
       \mathbb V'[\mathfrak f-\widetilde{\mathfrak f}]
                          +\mathbb W^{\rm 1l}
\]
is an escort of size at most $\theta$. The two constructions
$F_\emptyset[\mathscr L;\mathbb W^{\rm 1l}]$ and
$F_\emptyset[\widetilde{\mathscr L};\widetilde{\mathbb W}]$
have identical local potentials, domains, Gaussian, cutoffs and
pointing. They are equal as partial maps. The regular potential
uses at most $w_\star/4$ of the allowed size; adding the escort
therefore leaves both constructions within the allowed ball.
The one-lattice difference estimate \cref{III-34-30} gives the
asserted bound. This argument also identifies their domains;
equality of formal expansion coefficients alone would not suffice.
\end{proof}

\begin{proposition}\label{III-107-14}
Put $D_0^{\rm S}=E_0^{\rm S}/4$. Let $n_k$ bound the
one-lattice companion at step $k+1$, with zero companion at
source zero. Suppose
\begin{equation}\label{eq:regular-source-budget}
 \begin{split}
 108(L^3-1)K_\Pi n_0&\le D_0^{\rm S},\\
 \theta_k^{\rm S}:=
 K_\Lambda(1+\beta_{\rm s})\frac{\delta_k}{\alpha_0}
                         \frac L{L-1}D_0^{\rm S}+n_k
 &\le\min\left(\frac12w_\star,
           \frac{D_0^{\rm S}}{33(3L^3-3)K_\Pi}\right)
                    \quad(1\le k<K).
 \end{split}
\end{equation}
Then, for every $1\le j\le K$,
\begin{equation}\label{eq:regular-source-difference}
 \|E^{(j),w}-E^{(j),0}\|_{j,\kappa+1}
                \le D_0^{\rm S}\|w\|_{\mathsf{Src}}.
\end{equation}
The families have the usual translation covariance, with the
source translated in physical coordinates along with the domain.
\end{proposition}

\begin{proof}
First omit the factor $\|w\|_{\mathsf{Src}}$ from the claim.
The first-step escort estimate \cref{III-47-46} and the first
inequality of \eqref{eq:regular-source-budget} give the base case.
Suppose the bound holds through $k$. The counterterms cancel
between the two arrays, and therefore
\[
 N_{\mathfrak a,\omega}(\mathfrak f^w-\mathfrak f^0)
 =\sum_{j=1}^kL^{-(k-j)}
                 \|E^{(j),w}-E^{(j),0}\|_{j,\kappa}
 \le\frac L{L-1}D_0^{\rm S}.
\]
In \eqref{III-107-7} the normalizer also cancels.
\Cref{III-107-11} now bounds the new difference by
$33(3L^3-3)K_\Pi\theta_k^{\rm S}$ in rate
$\lambda_{\rm out}\ge\kappa+1$. This is at most
$D_0^{\rm S}$, closing the induction.
For the zero-source terms the actual one-step bound is at most
$\Xi(x_k,N_\star)\le\Xi(\gamma^{-2},N_\star)
\le3E_0^{\rm S}/4$, the stronger numerical budget imposed
in \cref{III-105-3}. This is the estimate used in the proof
of Theorem \ref{II-144-26}, with this stronger budget;
the statement of that theorem alone only gives $E_0$.
Hence the sourced regular bound in
\eqref{eq:sourced-recursion-bounds} follows at the same time.
This is the regular estimate used in the preceding induction,
and involves no boundary estimate.

For fixed $j,X,P,z$ the difference is holomorphic in source
slices and vanishes at zero. Applying the source Schwarz
estimate \cref{III-16-27} gives
\eqref{eq:regular-source-difference}. Taking suprema preserves
the estimate.

For translation covariance use the translated old entries
and companion in the pointed expansion. The canonical reference
plaquette obeys \eqref{eq:source-reference-translation}.
Thus the weight differences $v-v(p_b)$ translate along with the
localized potentials in \cref{III-41-24,III-48-54}; their finite
formulas give the translated companion. Restriction, localization,
Gaussian integration and pointing commute with the allowed lattice
translations. Induction makes the pointed terms
equal. The possible difference of the source-independent
normalizers is zero at source zero by the real covariance, so
\eqref{III-107-7} gives covariance for every source. The same
physical translation has level-$i$ coordinates
$L^{j-i}\mathsf v$ when its level-$j$ coordinates are
$\mathsf v$. These belong to the required $LM\mathbb Z^3$
at each earlier step.
\end{proof}

\section{The exact sourced history step}
\label{sec:exact-sourced-step}

The analytic families above must be evaluated at the real
history coordinates. This requires both domain information
and a comparison of representatives. We first work at a
chain point, where the coordinate identities are pointwise,
and then pass to the reference measure.

Let $U\in\mathcal V_0[h']\cap\operatorname{dom}\xi_{h'}$
be a bare field with nonzero history weight, and
$y=(W,B,\phi^+)$ a chain point of $U$. Write
$B=(B_{\rm I},\mathsf a)$, let $\mathsf d$ be its base datum,
and let $\mathsf U$ be an adapted representative. Let
$\mathsf a=\check\Psi_{\mathsf d}(\mathsf a')$ be the real
strip change of variables and $\mathfrak q^+=(\mathfrak q,\mathsf a')$
the stored data after the step. Put $P=P_+[\mathsf d]$ and
$P_X=P|_X$.

\subsection{Domains and exact identities at chain points}

\begin{lemma}\label{III-103-37}
For every such chain point and every source in the class, all
the following values are defined. The cluster function and
the two conditional forms can be evaluated at both
$\mathsf a$ and $\mathsf a'$. For every
$B_{\rm I}$ in the integration box $\mathsf{Rg}$ the ordinary
and weighted potentials can be evaluated at
\[
 \mathfrak u_Y=(P|_Y,g_k(B_{\rm I}\oplus\mathsf a)|_Y,
                                      \mathfrak q|_Y).
\]
Their finite sum $\mathsf V_{\mathsf d}^w(B_{\rm I},\mathsf a)$
is bounded on that box. If $\mathsf m$ is the real conditional
mean, then the centered source splitting gives
\begin{equation}\label{eq:actual-centered-source-split}
 \sum_Y\mathbb V^{\widehat w_k}(Y;\mathfrak u_Y)
 =\sum_Y\mathcal W[\widehat w_k,\mathsf{rf}^w]
       (Y;P|_Y;\mathcal A_Y,\zeta_Y)
       +\mathsf{Kd}_{\mathsf d}^w(\mathsf a).
\end{equation}
Here $\mathsf{rf}^w(b)=\widehat w_k(p_b)$,
$\mathcal A_Y=g_k(B_{\rm I}\oplus\mathsf a)|_Y$ and
$\zeta_Y=g_k((B_{\rm I}-\mathsf m)\oplus0)|_Y$.
By \cref{III-103-32}, each remainder is the canonical centered
member $\mathbb W^{\rm S}[v_k^w]$ evaluated at
$z=B_{\rm I}-\mathsf m$, $\mu=\mathsf m$, with external field
$\mathsf a$. The term $\mathsf{Kd}_{\mathsf d}^w$ is the sum of the quadratic
Taylor terms at $g_k(\mathsf m\oplus\mathsf a)$; in particular
it is independent of $B_{\rm I}$.
\end{lemma}

\begin{proof}
The domain assertions are the actual geometric assertions of
Propositions \ref{II-87-15}, \ref{II-115-84} and \ref{II-182-41}.
At the chain point $P_X\in\mathscr U_X$ and
$(P_X;\mathfrak q^+|_X)\in\mathfrak G_X$.
The strip budget is $C_{\rm xt}\delta_k/2$; both the old and
new external fields lie in the allowed external domain.
Restriction to $X$ maps these points into the parameter domain
of the centered cluster function. Thus the functions just
constructed are defined at both external fields.

For $B_{\rm I}\in\mathsf{Rg}$, the real box bounds and
Proposition~\ref{II-115-84} give
$\mathfrak u_Y\in\mathfrak D_k^{\rm rad}(Y)$.
They also give $P|_Y\in\mathbb U_k^{\mathbf1}(Y)$ and
$\sup|g_kB|\le C_{\mathcal A}\delta_k<r_0/2$.
The conditional mean belongs to the prescribed mean domain,
and $B_{\rm I}-\mathsf m$ belongs to its centered-field domain.
The source radius conditions leave the room required by the
weighted Taylor expansion. Consequently all ordinary, weighted
and centered members are defined. Their estimates, summed over
the finite set of $Y$, give boundedness on $\mathsf{Rg}$.

Apply the weighted Taylor splitting at each $Y$ with the
reference coefficient $v_k^w(p_b)$ at bond $b$.
The linear terms at constant weight cancel in the sum by the
first variation of the Wilson action at the minimizing field.
For the mixed quadratic terms, \cref{III-19-8} separates the
gauge and Wilson parts of the conditional-mean equation:
$\chi_{\rm I}\widehat G_{\mathsf d}\bar m=0$ and
$\chi_{\rm I}H_{\mathsf d}\bar m=0$, where
$H_{\mathsf d}=\mathsf A_{\mathsf d}-\widehat G_{\mathsf d}$.
The regional Hessian identity has the normalization
$2\sum_YT_2[\mathbf1](Z|_Y,\bar m|_Y)
 =(Z,H_{\mathsf d}\bar m)_1$. Its right side vanishes for
every integration-supported $Z$, including fields supported
on one bond. This cancels the reference-weight mixed terms.
The remaining pure mean quadratic
terms give $\mathsf{Kd}_{\mathsf d}^w$; the remaining terms are
exactly the three centered source members. This is the sum-rule
identity of \cref{III-103-32}, and gives
\eqref{eq:actual-centered-source-split}. No integration or choice of logarithm is
used in this identity. At $k=0$ the first-step Taylor sum rules
and the source family of \cref{III-48-54} give the same formula.
\end{proof}

\begin{proposition}\label{III-103-40}
Let $\mu_{\mathsf d,\mathsf a}$ be the conditional Gaussian
with the real covariance and mean of the history step. Then
\begin{equation}\label{III-103-41}
 \int\chi^{\mathfrak b}(B_{\rm I})
   e^{\mathsf V_{\mathsf d}^w(B_{\rm I},\mathsf a)
                         -\mathsf{Kd}_{\mathsf d}^w(\mathsf a)}
                         \,d\mu_{\mathsf d,\mathsf a}
 =\exp\sum_{X\in\mathcal D_t}
        E_{h'}^{{\rm cl},w}(X;\mathsf a;P_X;\mathfrak q|_X).
\end{equation}
The integrand is zero off $\mathsf{Rg}$. Moreover
\begin{align}
 -\tfrac12(\mathsf a,\Delta_{\mathsf d}^{\rm cond}\mathsf a)_1
             +\mathsf{Kd}_{\mathsf d}^w(\mathsf a)
 &=\sum_X[K^{\rm kp,w}+B^{\times,w}]
                                (X;p_X(P_X),\mathsf a),
                                      \label{III-103-36}\\
 \sum_X B^{{\rm sr},w}(X;P_X;\mathfrak q,\mathsf a')
 &=\mathsf{Zs}_{\mathsf d}(\mathsf a')
   +\sum_X[E^{{\rm cl},w}+B^{\times,w}+K^{\rm kp,w}]
                                       (X;\mathsf a)\notag\\
 &\quad-\sum_X[E^{{\rm cl},w}+B^{\times,w}+K^{\rm kp,w}]
                                       (X;\mathsf a').
                                      \label{III-103-20}
\end{align}
All values in these identities are defined at the chain point.
\end{proposition}

\begin{proof}
The complete local Gaussian structure of Theorem \ref{II-115-62}
has the covariance and mean of $\mu_{\mathsf d,\mathsf a}$
and the same integration cutoff. Apply \cref{III-32-12} at
$(\mathsf a;P;\mathfrak q)$, whose membership in its parameter
domain was checked above. The ordinary vertices are the four
old potentials; the only additional origin is the centered
source origin. By \cref{III-103-37} their sum is the exponent
on the left of \eqref{III-103-41}. A domain containing no
integration or external source cube has empty relevant bond
set; its ordinary terms and its centered source term at $(0,0)$
vanish. Thus the sum in the centered expansion agrees with the
full sum over $Y$. This proves \eqref{III-103-41} with its
stated domains and cutoff convention.

The real conditional form satisfies the identity without
$\mathsf{Kd}^w$, by Section \ref{sec:history-integrability}. The exact quadratic
localization of \cref{III-113-39} identifies
\[
 \mathsf{Kd}_{\mathsf d}^w(\mathsf a)
       =\sum_X[K^{\rm kp,+}+B^{\times,+}]
                                   (X;p_X(P_X),\mathsf a).
\]
Adding it gives \eqref{III-103-36}. Each summand is defined
on the pair image $\mathscr O^{\rm S}(X)$ and for every external field;
no regularity of the datum outside that image is needed.

For \eqref{III-103-20} sum the five finite strip transfer
identities, equations \eqref{II-73-39}, \eqref{II-73-43}, for the
sourced forms and cluster function. Their domains were checked
in the proof of \cref{III-102-21}. The Jacobian and Haar
contributions combine into the logarithm
$\mathsf{Zs}_{\mathsf d}=\log j_{\rm s}$ of the
Jacobian of the strip change in the external Lie-algebra
coordinates. The other three contributions are precisely
their differences at $\mathsf a$ and $\mathsf a'$. This proves
the identity without exchanging an infinite sum or an integral.

The normalizer and cluster interior identities together with
\eqref{eq:sourced-created-boundary} also show that their total
local exponent equals the new vacuum constant plus the regular,
edge and created-boundary terms. The constants
$n_{\rm tr}^{(t)}$ cancel by \eqref{III-107-7}. This is an
identity of defined values, including the terms whose boundary
indicator is zero. Empty integration regions cause no exception:
their Gaussian is the point mass at zero and the same formulas
are finite evaluations.
\end{proof}

\subsection{Borel domains and evaluation representatives}

\begin{proposition}\label{III-104-21}
The evaluations of the regular, boundary and retained entries
are Borel partial maps with domains independent of $w$.
More explicitly, regular evaluation at $(j,X,z)$ has domain
\[
 \{\mathsf U:\mathbf P_X^{(j)}(\mathsf U)
                     \in\mathcal U_j^\circ(X;\alpha_0,\alpha_1)\};
\]
boundary evaluation has the inverse image of
\[
 \mathsf{Sp}(\mathcal S;(j,X,{\rm mk}))\times
 \mathfrak Q^{\rm rad}_{h'|_j}(X)
\]
under $(\mathsf U,\mathfrak q)\mapsto
(\mathbf P_{X|\mathcal S}^{(j)}(\mathsf U),\mathfrak q|_X)$;
and retained evaluation has the analogous inverse image of
\eqref{eq:sourced-retained-domain} with its stored prefix.

Put $D_X=\{\mathsf U:P_X(\mathsf U)\in\mathscr U_X\}$.
The conditional forms have domain $D_X\times\mathsf E_{h'}$,
and the cluster and strip terms have domain
\begin{equation}\label{eq:sourced-cluster-evaluation-domain}
 \{(\mathsf U,\mathfrak q,\mathsf a_1):
   \mathsf U\in D_X,\ (\mathfrak q,\mathsf a_1)|_X
                                  \in\mathfrak Q^{\rm rad}_{h'}(X)\}.
\end{equation}
Set $\mathsf r=c_{\rm fl}\mathfrak b_k$ and write
\[
 \mathsf{Rg}(\mathsf r)
 =\{B=(B_{\rm I},\mathsf a):B\text{ is a real }1\text{-form on }
       \mathfrak X,\ |B(b)|<\mathsf r
       \text{ for every }b\in\mathfrak X^+\}.
\]
The finite fluctuation potential has the Borel domain
\[
 \{(\mathsf d,B)\in\mathsf{Ba}^{\rm reg}_{h'}\times\mathsf{Rg}(\mathsf r):
       \mathfrak u_Y(\mathsf d,B)\in\mathfrak D_k^{\rm rad}(Y)
                                     \text{ for every }Y\}.
\]
\end{proposition}

\begin{proof}
The plain-triple maps are continuous and the cores and graded
spaces are Borel; the tiers are open, being defined by finitely
many strict inequalities. These facts belong to the coordinate
construction in Sections \ref{sec:history-evaluation}--\ref{sec:history-integrability} and do not use an
interaction list. By \cref{III-102-21}, each entry is defined
on the entire indicated core or product and is trace-measurable
there. Thus there is no extra source-dependent domain condition.
Composition with the continuous reading maps gives exactly the
Borel domains displayed above.

For the cluster and strip terms the continuous reading
$(P;\mathfrak q,\mathsf a_1)\mapsto
(\mathsf a_1^\uparrow;P;\mathfrak q)$ maps the graded-tier
product into the cluster domain, with strip budget
$C_{\rm xt}\delta_k/2$. The full-domain result of
\cref{III-102-21} therefore gives
\eqref{eq:sourced-cluster-evaluation-domain}. A conditional
form is a finite polynomial in $\mathsf a_1$ with analytic
pair coefficients. Pair reading is continuous on
$\mathscr U_X$, so its evaluation is Borel on
$D_X\times\mathsf E_{h'}$.

Finally $\mathsf d\mapsto P_+[\mathsf d]$ is Borel, field
restriction is linear, and the old stored coordinates are
coordinate projections. Hence all $\mathfrak u_Y$ are Borel
maps. All five potentials are defined and trace-measurable
on $\mathfrak D_k^{\rm rad}(Y)$, so their finite sum is Borel
on the displayed intersection of inverse images.
The auxiliary quantity $\mathsf{Kd}^w$ is used only where its
conditional objects and weighted Taylor parts are defined.
At chain points this has already been proved; no assertion
that its maximal domain is Borel is needed below.
\end{proof}

\begin{lemma}\label{II-168-26}
Let $U$ be in the actual support of $h'$, and let
$X\in\mathsf{Ix}^{\rm mk}_t(\mathcal S(h'))$.
Let $f_X$ be a complex function on the whole graded-tier
product $\mathfrak G_X$ satisfying the joint gauge identity
\[
 f_X(P^u;u\cdot\mathfrak q^+)=f_X(P;\mathfrak q^+)
\]
for every bare gauge transformation. At the actual stored
field $\mathfrak q^+[U]$, $f_X$ is defined and has one value
on all adapted evaluation representatives of $U$.
\end{lemma}

\begin{proof}
Actual support membership puts each representative's plain
triple in the graded space and its stored field in the tiers;
these memberships are the ones in the proof of
Corollary~\ref{II-99-39}. Two representatives differ
by a residual gauge on the first collared region. Extend
it by the identity outside that region. At each outer
bond the endpoints in the collar are bare nodes, so this
extension relates the two glued configurations as well.
It fixes nodes. It also fixes every stored site: the
transport of an adapted representative from its stored site
$a$ to the corresponding node $y$ obeys
\[
 \mathsf{ax}_{U'^u}(a)=u(y)\mathsf{ax}_{U'}(a)u(a)^{-1},
\]
and both transports and $u(y)$ are the identity.
All sites acting on the stored fields are nodes or these
stored sites, by Lemmas \ref{II-167-6} and \ref{II-167-13}.
Thus $u\cdot\mathfrak q^+[U]=\mathfrak q^+[U]$.
The plain-triple maps are gauge covariant. Apply the
joint gauge identity to the two triples and the fixed
stored field. The claimed equality follows. This proof
uses no reality condition on $f_X$.
\end{proof}

\begin{proposition}\label{III-104-17}\label{III-104-18}
At a realized chain point the new action exponent and the new
retained exponent sum are separately defined and constant on
the evaluation set $\mathsf{Ev}_{h'}(\xi_{h'}U)$.
The retained sum made directly by the integration step equals
the sum read from $\mathscr L_t^{w\flat}$ on this set.
\end{proposition}

\begin{proof}
All representatives in this evaluation set satisfy the graded
and tier conditions just used. The action exponent consists
of regular and boundary values with the actual gauge
invariance and pair dependence, together with the weighted
Wilson action. The adapted-gauge comparison of
Proposition \ref{II-168-27} makes these values independent
of the representative; the weighted action has the same
property by its plaquette gauge invariance.

Consider the additional kept form at $X$. If
$\mathsf{kp}_t(X)=0$, it vanishes. Otherwise $X$ contains a
big cube of side $LM$ in level-$k$ units, contained in
$\Omega_t\cap Z_t$, and belongs to the boundary index set.
On the graded-tier product define
\[
 f_X(P;\mathfrak q^+)=
 K^{\rm kp,+}(X;p_X(P),(\mathfrak q_k^+)^\uparrow).
\]
The quadratic locality theorem shows that this depends only
on the external bonds in $X$; zero extension elsewhere is
immaterial. Its pair covariance and the simultaneous action
on the stored field give the boundary gauge identity.
Apply \cref{II-168-26}, whose hypotheses were just checked.
Consequently $f_X$ has the same value at every adapted
representative, with the fixed stored field. This proves
representative independence of each added kept term at this
realized point.

The retained sum produced by the step is the old retained sum,
the real newly retained conditional terms, the elementary
new factors, and $\sum_XK^{\rm kp,+}(X;p_X,\mathsf a')$.
The first three parts have the actual representative identity
of the real construction, whose proof uses their gauge
identities and full domains, not their real-valuedness.
The last part has the identity just proved. Thus the sum
is separately well defined.

To compare it with the list, old factors read the identical
plain triples and the identical stored prefixes. The new
elementary factors are unchanged. For each new kept factor,
\cref{III-113-39} identifies its correction in
\eqref{III-101-23} with $K^{\rm kp,+}(X_f;p_{X_f},\mathsf a')$;
the field may be extended by zero because of external-bond
locality. There is one such factor for each domain with
$\mathsf{kp}_t(X)=1$, and the added form vanishes on the
other domains. Summing therefore gives exactly the step's
retained sum. This proves equality at every representative
of the realized point. The argument does not assert
representative independence at arbitrary points of the
coordinate space.
\end{proof}

\subsection{The measure identity for complex densities}

We give the measure argument because the positivity of the
real source-free density cannot be imposed on a complex
source. Its geometric ingredients are those of
Theorem \ref{II-166-30} and its proof.

\begin{theorem}\label{III-38-20}
Suppose the old sourced format has an integrable partial
density $\rho_h^w$ and equals
$\hat\tau_h[\rho_0^w]$. Suppose the sourced analytic families
have the domains, chain-point identities and representative
identities proved above. Then the density made by the next
integration step is an integrable partial density and
\begin{equation}\label{eq:sourced-made-measure}
 \hat\tau_{h'}[\rho_0^w]
 =\mathsf{Stay}_{h'}^{\rm T}\chi_t(\Omega_t)
               e^{A_{t,h'}^w+\Phi_{h'}^{\rm made,w}}\,\mathscr M_{h'}.
\end{equation}
Each exponent on the right is separately defined almost
everywhere where the test is nonzero. There is the variation
bound
\begin{equation}\label{III-38-26}
 \int|\rho_{h'}^{\rm made,w}|\,d\mathscr M_{h'}
 \le\|f\|_\infty\int|\rho_h^w|\,d\mathscr M_h,
\end{equation}
where $f$ is the bounded nonnegative step weight.
\end{theorem}

\begin{proof}
Write $\nu_h=\hat\tau_h[\rho_0^w]$ and suppress $w$ from
the notation temporarily. The coordinate commutation and
weight identity imply, for bounded Borel test functions,
\[
 \nu_{h'}=\Sigma_* (f\nu_h).
\]
This is a finite complex-measure identity, since the bare
source density is bounded and the real history weights are
nonnegative and bounded. Extend $\rho_h$ by zero on a Borel
null set containing its undefined points. Where the old
test is nonzero and the exponents are defined, its value
is nonzero: it is a positive test times an exponential.
This observation does not assert positivity of $\rho_h$.

We first justify using chain-point identities almost
everywhere in the coordinate change. Let $\nu_+$ be the
positive measure obtained from the bare real history weight
for $h'$ by forgetting the source density and pushing to
the old coordinates. The analytic-image argument in the
proof of Theorem~\ref{II-166-30} gives Borel sets $D'\subset D''$ such
that the actual image lies between them,
$\nu_+(D''\setminus D')=0$, and $\nu_+$ is concentrated
on $D'$. The measure $f\nu_h$ is dominated in variation
by $\|\rho_0^w\|_\infty\nu_+$, so it too vanishes on
every Borel subset outside $D'$. On the nonzero old test
outside the old null set its density $f\rho_h$ is nonzero
where $f>0$. If that set outside $D'$ had positive
$\mathscr M_h$ measure, the complex function $f\rho_h$
could not have integral zero on every Borel subset of it.
Thus it is a reference null set. The same argument removes
the part outside the domain of $\Sigma$. This is the
required realization statement; cancellation of the total
mass of a set has not been used as a positivity argument.

Use the injective chart and its positive Jacobian $j_{\rm c}$,
then the strip map and its positive Jacobian $j_{\rm s}$.
Their identities are equations \eqref{II-166-48}, \eqref{II-166-50}.
Write $p=(W,\mathsf a',\phi^+)$ for the remaining variables
and $\Theta(p,B_{\rm I})$ for the old chart point. The
pullback integrand is the Borel function
\[
 D(p,B_{\rm I})=j_{\rm s}(p)j_{\rm c}(\Theta(p,B_{\rm I}))
                          (f\rho_h)(\operatorname{pt}\Theta(p,B_{\rm I})).
\]
It is set equal to zero outside the chart and its nonzero
test. Apply the two change-of-variables formulas first to
the nonnegative function $|f\rho_h|$. This proves
\[
 \int|D|\le\|f\|_\infty\int|\rho_h|\,d\mathscr M_h<\infty.
\]
Complex Fubini is now available. The chart realization
proved above and positivity of both Jacobians imply that
almost every nonzero-test chart point is a chain point.
The inherited test identity puts
$\mathsf{Pe}\times\mathsf{Rg}$ in this nonzero-test chart
set outside a product null set, as in
equation \eqref{II-166-17}. Pulling back the omitted chart
null set through the two changes of variables therefore
leaves a set of full product measure in this product.
Tonelli gives a Borel set of full measure of $p\in\mathsf{Pe}$
whose good fibre has the full volume of $\mathsf{Rg}$.
This box has positive finite volume; in dimension zero its
volume is one. Thus each such good fibre is nonempty.

Fix such a $p$. Choose one admissible $B_{\rm I}$ to verify
the domain and representative assertions for the new
exponents. They depend on $p$, not on that choice of
$B_{\rm I}$. Their values obtained through a fixed Borel
adapted selection are Borel by \cref{III-104-21}; the
representative identity makes them their evaluation values.
No measurable choice of this auxiliary $B_{\rm I}$ is required.

The old action identity, the chart Jacobian identity and
the test identity give, at almost every point of this fibre,
the same finite algebra as equations \eqref{II-166-17}, \eqref{II-166-58}:
the $B$-dependent part is
\[
 \exp\{-\tfrac12(B,\mathsf A_{\mathsf d}B)_1
                        +\mathsf V_{\mathsf d}^w(B_{\rm I},\mathsf a)\}.
\]
The weighted action enters this identity by its exact
weighted finite sum; the old interaction potentials enter
by the linear extension established in
\cref{III-102-21}. Gaussian conditioning uses the positive
real block $(\mathsf A_{\mathsf d})_{\rm II}$.
Equation \eqref{III-103-41} evaluates its fluctuation integral
as the exponential of the cluster sum plus
$\mathsf{Kd}_{\mathsf d}^w(\mathsf a)$.
Equation \eqref{III-103-36} combines this last term with
the conditional quadratic form. Equation \eqref{III-103-20}
then replaces $\mathsf a$ by $\mathsf a'$ and absorbs
$j_{\rm s}$. The interior subtraction
\eqref{eq:sourced-created-boundary} identifies the remaining
local terms with $A_{t,h'}^w$ and the new retained sum.
Thus the fibre integral of $D$ is exactly the pullback of
the right side of \eqref{eq:sourced-made-measure}.

The last coordinate embedding is Borel and injective and
pushes its reference measure to $\mathscr M_{h'}$ on the
support of the test. Descending the preceding identity
proves \eqref{eq:sourced-made-measure}, including almost
everywhere definedness of both exponents and integrability.
The bound on $\int|D|$ and the triangle inequality for its
fibre integrals prove \eqref{III-38-26}. Empty integration
boxes have zero-dimensional Gaussian measure and are
covered by the same argument.
\end{proof}

\subsection{The stored family and induction through the levels}

\begin{theorem}\label{III-104-28}
At every real history step there is a Borel set
$N_{h'}\subset\mathsf X_{h'}^{\mathrm{co}\prime}$,
$\mathscr M_{h'}(N_{h'})=0$, independent of the source, such
that on the nonzero-test set outside $N_{h'}$ the made
retained sum and the retained sum read from
$\mathscr L_t^{w\flat}$ are separately defined and equal
for every source in the class. If the old sourced format
identity holds, then
\begin{equation}\label{eq:sourced-list-measure}
 \hat\tau_{h'}[\rho_0^w]
 =\mathsf{Stay}_{h'}^{\rm T}\chi_t(\Omega_t)
        e^{A_{t,h'}^w+\Phi_{h'}^{\mathscr L,w}}\,\mathscr M_{h'}.
\end{equation}
The right side is an integrable partial density and is
independent of the Borel choice of adapted representatives.
\end{theorem}

\begin{proof}
Apply the real format theorem and the realization argument
in Theorem \ref{II-166-30} at source zero. Its positive
density is nonzero on the test after deletion of a Borel
reference null set. Hence outside a Borel set $N_{h'}$
every nonzero-test coordinate point is $\xi_{h'}U$ for a
bare field with a chain point. This construction involves
only the real histories, coordinates and tests, so the same
$N_{h'}$ works for every source.

At each such point \cref{III-104-17} identifies the two
retained sums for every $w$. The old sourced identity and
\cref{III-38-20} make the made density integrable and exact.
The list density has the same other factors and agrees
with it off $N_{h'}$ wherever the test is nonzero. On the
zero-test set both densities are defined to be zero without
evaluating either exponent. Their integrable partial-density
classes are therefore equal, which proves
\eqref{eq:sourced-list-measure}. Each exponent is an
evaluation value, already independent of the representative;
the asserted independence of the measure follows.
\end{proof}

\begin{proposition}\label{III-104-30}\label{III-106-41}
The output density in \eqref{eq:sourced-list-measure} is
exactly the old density used at the next step, as a partial
map, with the same reference measure. At the first step
the old identity holds without any earlier construction:
\[
 \hat\tau_\emptyset[\rho_0^w]=\rho_0^w\,dU.
\]
Consequently the sourced format identity holds at every
level and for every prefix of the construction.
\end{proposition}

\begin{proof}
We first check the base step. The level-zero test is one,
there are no retained or boundary entries, and its action
is the bare real action plus $L^KA(w,U)$. Thus its density
is exactly $\rho_0^w$, which is bounded on the compact bare
configuration space.
The first-step Gaussian and coordinate identities are the
real ones. The first-step weighted Taylor splitting,
quadratic decomposition and companion are those of
\cref{III-48-54}; their actual arguments satisfy the
first-step small-field domains. The proof of
\cref{III-103-40} therefore applies with $k=0$.
There are no inherited interaction potentials or old
retained exponents to check. The first regular and created
boundary families have the domains proved in
\cref{III-102-21}, so their trivial values exist and the
interior normalizer subtraction is legitimate.
The complex measure argument and the source-independent
real null set now give \eqref{eq:sourced-list-measure}
at level one.

For compatibility, cut all data of step $t+1$ at level $t$.
The recursion \eqref{III-102-8} through that level is
unchanged: it uses the same couplings, companions, strip
rows and domain restrictions. The coordinate space,
reference measure, test and evaluation sets are also
unchanged. An old retained factor reads its formation-level
plain triple and the same stored prefix. A boundary term
reads the same filing space and stored tier. The regular
values, counterterms, vacuum values and source Wilson term
are likewise identical. Thus the two densities have exactly
the same factors and the same domains, including the
convention on zero tests. This is equality of partial maps,
not just equality almost everywhere. The level-one output
is therefore the hypothesis needed at level two, and the
same argument applies thereafter.

If $\Omega_1=\emptyset$, the real history format is multiplied
by the bounded, nonvanishing source factor read from the
stored outer field, as proved in \cref{III-24-31}. There
are no new source-dependent families in this case. This
settles the remaining histories and completes the induction.
\end{proof}

The retained-factor damping estimate is now applied to
defined evaluation values. Specifically \cref{III-113-44}
with $\bar\upsilon\le\gamma'$ and
$\gamma^\dagger=\gamma'/2$ gives the sourced damping bound
off the same real realization null set. It is not being
used to establish that the retained sum has an evaluation
value. That domain assertion was proved before the estimate.


\section{A common choice of constants}
\label{sec:sourced-common-constants}

The source estimates require extra room in the decay rates before the
coupling is made small. We first fix that room, then choose a relative
source radius, and finally choose the upper bound on the coupling.
The choices are independent of the number of steps and of the torus.
Throughout this section the constants of the real construction have
their common values in Proposition~\ref{II-202-29}. In particular,
$\lambda_{\rm s}=c_{\rm E}=\log18432$, $S_{\rm s}=C_{\rm E}=4096$,
$C_2=342$, and $C_3=9e^8$. The quantities furnished by the Gaussian,
normalizer and background estimates retain those values when the
regular and boundary size bounds are enlarged.

\subsection{Rates and fixed radii}

Set $d=1-2\delta_{\rm sp}>0$ and use the shifted pair
\begin{equation}\label{eq:source-shifted-rate}
 \kappa^\sharp=\kappa+3/d,\qquad
 r^\sharp=d(\kappa^\sharp-2)=r_V+3,
 \qquad r_V=d(\kappa-2).
\end{equation}
The quantities $\kappa''=(1+4\beta_{\rm s})\kappa+1+c_{\rm E}$,
$\bar c_P$, $\bar r_H$ and $K_{\rm ap}$ refer to the real
Gaussian construction. Define the uniform centered majorants by
\begin{gather}
 \bar c_P^+=C_{\rm mn}^3\lambda_{\rm r}^{-1}
                       (N_{\rm gs}+1),\qquad
 \bar r_H^+=\kappa''+c_{\rm E}+C_{\rm cs}(\bar c_P^++2),
 \label{eq:source-centered-majorants}\\
 K_{\rm ap}^+=C_{\rm cs}^2
            e^{C_{\rm cs}(\kappa''+\bar c_P^++\bar r_H^+)},\qquad
 \theta_{\rm sl}=(1+\vartheta_*)/2.
 \notag
\end{gather}
Here $C_{\rm mn}=18(1-\vartheta_*)^{-2}$ and
$C_{\rm cs}=C_{\rm cs,0}$ are the constants of the centered expansion.
The number $\bar c_P^+$ depends on $L$ and the reserve fractions, but
not on $\kappa$, $\bar\kappa$ or $M$.

In the order of the choices in Proposition~\ref{II-202-29}, impose
the following additional floors:
\begin{align}
 L&\ge\frac{2C_{\rm ex}(1+4\beta_{\rm s})}{d},
                                                  \label{eq:source-L-floor}\\
 \tfrac12d\kappa&\ge2d+2c_{\rm E}
       +\frac{C_{\rm cs}}L[1+2c_{\rm E}+C_{\rm cs}(\bar c_P^++2)],
                                                  \label{eq:source-centered-kappa-floor}\\
 \tfrac12d\kappa&\ge2d+3c_{\rm E}
       +\frac{C_{\rm ex}}L(8c_{\rm w}+C_{\rm ex}+c_{\rm E}),
                                                  \label{eq:source-comparison-kappa-floor}\\
 \tfrac12d\kappa&\ge2d+2c_{\rm E}+2
       +\frac{C_{\rm cs}}L[1+2c_{\rm E}+C_{\rm cs}(\bar c_P+2)],
                                                  \label{eq:source-chain-kappa-floor}\\
 \bar\kappa&\ge1+C_{\rm cs}\bar r_H^++\log C_{\rm cs},
                                                  \label{eq:source-barkappa-floor}\\
 (LM)^3&\ge C_{\rm mn}^{-1}\bar c_P^+,\qquad
 \delta_{\rm res}LM\ge2C_2(\kappa''+1).
                                                  \label{eq:source-M-floors}
\end{align}
The real construction already requires
$L\ge2C_{\rm cs}(1+4\beta_{\rm s})/d$ and
$4\beta_{\rm s}\kappa\ge1$; both requirements are retained.
Each new right side involves only constants chosen earlier than the
constant on its left. In particular the $M$ floors are imposed after
$\kappa''$ and $\delta_{\rm res}$ have been fixed; $M$ is still allowed
to run through arbitrarily large powers of $L$.

\begin{lemma}\label{III-110-17}
These floors are compatible with the real choices. Put
\begin{equation}\label{eq:source-chain-rate-choice}
 \kappa'=\kappa'',\qquad
 \delta_\upsilon=\delta^{\rm S}
              =\frac{2C_2(\kappa''+1)}{LM}.
\end{equation}
Then
\begin{gather}
 0<\delta^{\rm S}<\delta_{\mathbf M}
           \le\min(\delta_{\rm res},\delta_*^{\rm pl}),\qquad
 \tfrac12\delta^{\rm S}LM=C_2(\kappa'+1),
                                                \label{III-110-18}\\
 \lambda_V\le r_V-c_{\rm E}\le r_V-2,
 \qquad r^\sharp-\delta^{\rm S}M\ge2c_{\rm E}+5,
 \qquad \lambda_V\le r^\sharp-2.                 \notag
\end{gather}
The three weighted-potential conditions hold at
$(\delta_{\rm sp},\kappa^\sharp)$:
\begin{equation}\label{eq:source-Vc-sharp}
 2\delta_{\rm sp}\kappa^\sharp\ge a_{\rm pl},\quad
 \bar\kappa-1\ge C_\kappa\kappa^\sharp,\quad
 M\min(\mu_{\mathfrak h},1)L\ge2C_{\rm lev}\log L.
\end{equation}
The centered expansion has its fixed geometric and rate hypotheses,
and $\alpha_0,\alpha_1\le\frac12\min(a_{\mathfrak h},a_P^\circ)$.
Its Gaussian growth condition holds at $\rho_1=\rho_1^{\rm S}=(LM)^{-3}$.
\end{lemma}

\begin{proof}
The new conditions are finite lower bounds at the four upward choices
$L,\kappa,\bar\kappa,M$. Taking the maxima of the old and new floors
therefore preserves the construction in the cited proposition.
No bound on a constant already chosen is imposed at a later stage.

Write $\bar r_H=(1+4\beta_{\rm s})\kappa+1+2c_{\rm E}
                  +C_{\rm cs}(\bar c_P+2)$.
The real $L$ floor and \eqref{eq:source-chain-kappa-floor} give
\[
 \frac{C_{\rm cs}}L\bar r_H\le r_V-2c_{\rm E}-2.
\]
Since $C_{\rm cs}\ge2C_2$ and $\bar r_H\ge\kappa''+1$,
$\delta^{\rm S}M\le C_{\rm cs}\bar r_H/L$. This proves
$r_V-\delta^{\rm S}M\ge2c_{\rm E}+2$ and hence the asserted
inequality for $r^\sharp$. For the other rate, insert
$\lambda_1=(1+4\beta_{\rm s})\kappa+8c_{\rm w}+C_{\rm ex}+c_{\rm E}$
into $\lambda_V=C_{\rm ex}\lambda_1/L+2c_{\rm E}$.
Equations \eqref{eq:source-L-floor} and
\eqref{eq:source-comparison-kappa-floor} then give
$\lambda_V\le d(\kappa-2)-c_{\rm E}$.

For the range of $\delta^{\rm S}$, the real resolvent construction gives
\[
 \delta_{\mathbf M}=
 \frac{c_{\rm res}\delta_*^{\rm pl}\gamma_{\rm pos}}{B_{\rm nm}}
 \le\min(\delta_{\rm res},\delta_*^{\rm pl}),
 \qquad \delta_{\mathbf M}LM>8\bar\kappa.
\]
The last inequality follows from its chain condition:
$\mu_{\rm ch}LM/4>\bar\kappa$ and
$\mu_{\rm ch}\le\delta_{\mathbf M}/2$.
Also $\bar\kappa\ge C_{\rm cs}\bar r_H
\ge2C_2(\kappa''+1)$. Thus
$\delta^{\rm S}LM<\delta_{\mathbf M}LM$.
All divisions are by positive numbers.

The first and third inequalities in \eqref{eq:source-Vc-sharp}
are the corresponding real weighted-potential inequalities, with the
rate increased. For the second, the canonical value is $C_\kappa=1$.
The new $\bar\kappa$ floor gives
$\bar\kappa-1\ge C_{\rm cs}\bar r_H^+\ge2\kappa''\ge2\kappa$,
whereas the $\kappa$ floor gives $\kappa\ge4/d$.
Consequently $\bar\kappa-1\ge\kappa+3/d=\kappa^\sharp$.
The fixed radius bounds
are among the real analytic-radius conditions. Finally,
$\vartheta_*<\theta_{\rm sl}<1$ and
$1+\vartheta_*<2<C_{\mathcal A}$, so the center ball fits strictly
inside the field domain. Equations
\eqref{eq:source-centered-kappa-floor},
\eqref{eq:source-barkappa-floor} and \eqref{eq:source-M-floors}
are its remaining fixed rate and Gaussian-growth conditions.

For the Gaussian condition at $\rho_1^{\rm S}$, distinguish the
actual mean error $\varepsilon_{\mathbf m}$ in
\eqref{eq:common-mean-admission} from the envelope $\varepsilon_m$
defined after \eqref{III-31-5}. The external-mean admission
\cref{II-115-84} gives $\varepsilon_{\mathbf m}\le\varepsilon_m$.
Put $v=(LM)^{-3}\le1/8$. Since
$\varepsilon_{\rm dc}\le v$, $\theta_{\rm xt}=1/(3\sqrt{96})<1/24$ and
$\varepsilon_{\rm xt}=\sqrt v/4$, the definition of the envelope gives
\[
 \frac{\varepsilon_m}{\sqrt v}
 \le\sqrt{2v}\,\theta_{\rm xt}+\frac{1+v}{4}
 \le\frac5{16}<\frac12.
\]
Thus $(1+c_{\rm r})\varepsilon_m^2\le c_{\rm r}v$, since
$c_{\rm r}\ge1$. These bounds and $LM\ge13$ also give
$\theta_{\rm xt}+\varepsilon_m<1/16\le\vartheta_*$.
The defect bound and \eqref{II-115-55} give
$\widehat\varepsilon\le64B_{\rm nm}^5\gamma_0^{-2}v$.
Since $\gamma_{2,\rm gs}=\rho_1^{\rm S}=v$, we obtain
\[
 (LM)^3[\widehat\varepsilon+(1+c_{\rm r})
       \varepsilon_m^2+\gamma_{2,\rm gs}+\rho_1^{\rm S}]
 \le64B_{\rm nm}^5\gamma_0^{-2}+c_{\rm r}+2=N_{\rm gs}.
\]
Consequently the actual centered coefficient satisfies
$c_P^+\le C_{\rm mn}^3\lambda_{\rm r}^{-1}N_{\rm gs}
=\bar c_P<\bar c_P^+$, and the growth inequality follows from
$(LM)^3\ge C_{\rm mn}^{-1}\bar c_P^+$.
\end{proof}

\subsection{Source sizes and the final coupling choice}

The constants in the Taylor estimates are denoted by
$(\mathsf K_0,r_0)$ at later steps and by
$(\mathsf K^0,r_P^\circ)$ at the first step. Thus
\[
 r_0=\min\{\varepsilon_{\mathfrak h},(3c_{\mathfrak h})^{-1}\},
 \quad \mathsf K_0=12L^4M^3K_{\rm c},\quad
 \mathsf K^0=4512C_PM^3r_P^\circ.
\]
Use the coefficients $C^{\rm src}$ and $C^{\rm src}_0$ in the
general and first-step centered bounds, and set
\begin{gather}
 \mathsf k_\odot=49L^5C_{\mathcal A}
                       \max(C^{\rm src},C^{\rm src}_0),\qquad
 \varepsilon^{\rm O}=(16eL^3K_{\rm ap}^+)^{-1},
 \quad c_M^{\rm O}=w_0\varepsilon^{\rm O},
 \quad\epsilon_\times^+=\varepsilon^{\rm O}/3,\notag\\
 \mathsf m_\odot=
   \min\{c_M^{\rm O},w_\star/2,E_0/(324L^3K_\Pi)\}.
                                               \label{eq:source-fixed-smallness}
\end{gather}
For the one-source chains let $c_\upsilon,B_{\rm ch}^\sharp,
K_{\rm ch}^\sharp,C_\Upsilon$ be the quantities of
\cref{III-112-22} at \eqref{eq:source-chain-rate-choice} and at the
real resolvent bound $B_{\rm f}^{\rm res}$. Put
\begin{gather}
 \mathsf K_*^{\rm S}=\max\{\mathsf K_0r_0^{-2},
                         \mathsf K^0(r_P^\circ)^{-2}\},\quad
 b_\upsilon^{\rm S}=64S_{\rm s}e^{6\delta^{\rm S}M}\mathsf K_*^{\rm S},
 \quad m^{\rm S}=c_\upsilon^2B_{\rm f}^{\rm res}b_\upsilon^{\rm S},
 \notag\\
 \rho_1^{\rm S}=(LM)^{-3},\qquad
 \lambda_0^{\rm S}=\rho_1^{\rm S}/(16S_{\rm s}),\notag\\
 \eta^{\rm S}=\min\left\{
 \frac{\lambda_0^{\rm S}}{12eL^3K_{\rm ap}},
 \frac{w_\star}{4M^3},\frac{E_0}{2808L^3M^3K_\Pi}\right\},\qquad
 \bar n_{\rm co}^{\rm S}=M^3\eta^{\rm S},
                                               \label{eq:source-final-escort-ceiling}\\
 \mathsf m_\odot^{\rm S}=\min\{\mathsf m_\odot,E_0/(81L^3K_\Pi)\},
 \qquad \epsilon_{\rm ap}^{\rm S}
 =\min\{[24e(1+\sqrt2)L^3K_{\rm ap}]^{-1},\epsilon_\times^+\}.
 \notag
\end{gather}
All these numbers are positive and fixed before $\gamma$.
The fixed size bounds are
\begin{equation}\label{eq:source-final-box}
 E_0^{\rm S}=E_0,\qquad B^{\rm S}=B,
 \qquad B_{\rm fr}^{\rm S}=B_{\rm fr}^{\flat}+1,
 \qquad \bar b_{\rm fr}^{\rm S}=1,\qquad
 \gamma^\dagger=\gamma_{\rm kp}/2.
\end{equation}
In particular $B\ge\max\{1,8M^3E_0,24(LM)^3E_{\rm nf}+1\}$.

Here is one explicit choice of the two small relative radii.
Use the positive stability radius $\bar\sigma$ defined below and set
\begin{equation}\label{III-109-5}
\begin{split}
 \sigma^{\rm S}=\min\biggl\{&\frac12,
 \frac{\eta^{\rm S}}{8\mathsf K_*^{\rm S}},
 \frac{\gamma_{\rm kp}}{C_\Upsilon b_\upsilon^{\rm S}},
 \frac{e^{\kappa''}}{m^{\rm S}K_{\rm ch}^\sharp C_{\rm xt}^2},
 \frac{c_{\rm kp}^{\rm th}e^{\kappa''}}{m^{\rm S}K_{\rm ch}^\sharp},\\
 &\frac{c_{\rm rj}^{\rm th}e^{\kappa''}}
             {m^{\rm S}C_3B_{\rm ch}^\sharp(LM)^6},
 \frac{\sigma_*}{16},\bar\sigma,\frac{\sqrt3}{40}\bar\sigma,
 \frac{\mathsf m_\odot^{\rm S}}{4\mathsf k_\odot},
 \frac{M^3\eta^{\rm S}}{4\mathsf k_\odot},
 \frac{\sigma_{\rm src}}2\biggr\},
\end{split}
\end{equation}
and
\begin{equation}\label{III-109-6}
 \sigma_{\rm L}^{\rm S}=\min\left\{
 \frac{\eta^{\rm S}}{2K_{\rm lp}^{\rm S}},
 \frac{\eta^{\rm S}}{2K_{\rm lp}^{{\rm S},0}},
 \frac{10\mathsf m_\odot^{\rm S}}
                 {\sqrt3(1+A_1)\mathsf k_\odot},
 \frac{\min(\mathsf m_\odot,M^3\eta^{\rm S})}
                 {4(1+A_1)\mathsf k_\odot},
 \frac{20\mu_*}{\sqrt3K_{\rm L}^\vee}\right\}.
\end{equation}
The constants $K_{\rm lp}^{\rm S},K_{\rm lp}^{{\rm S},0}$ are
\begin{align*}
 K_{\rm lp}^{\rm S}&=\mathsf K_0\left[
 \frac{24\sqrt3M^3C_{\mathcal A}A_1}{5r_0}
 +\frac{48M^{3/2}(C_{\mathcal A}A_1/4+\sqrt2/2)}{5r_0^2}\right],\\
 K_{\rm lp}^{{\rm S},0}&=\mathsf K^0\left[
 \frac{3\sqrt3M^3C_{\mathcal A}A_1}{5r_P^\circ}
 +\frac{6\sqrt2M^{3/2}(C_{\mathcal A}A_1/4+\sqrt2/2)}
                       {5(r_P^\circ)^2}\right].
\end{align*}
The progression constants in these minima are
\begin{gather*}
 \varepsilon_{\rm src}^{\rm P}
       =\min\{\varepsilon_{\mathfrak h},c_{\rm os}\alpha_0/(2c_{\mathfrak h})\},
 \qquad b_{\rm A}=6\cdot1201^2M^3,\\
 K_{\rm Q}^\vee=\max\{672M^6(\alpha_0/\varepsilon_{\rm src}^{\rm P})^2,
                   2b_{\rm A}M^3/(r_P^\circ)^2\},\\
 K_{\rm L}^\vee=\max\{8064M^6\alpha_0/\varepsilon_{\rm src}^{\rm P},
                   72b_{\rm A}M^3/r_P^\circ\},\\
 K_{\mathbf M}^{\rm f}=3L^3\max\{K_Ee^{4\lambda_1},
                      32K_\Pi K_P,(117/2)K_\Pi\},\\
 \mu_* =\min\{w_\star/32,E_0/(16K_{\mathbf M}^{\rm f})\},\qquad
 \sigma_{\rm src}=\mu_*/K_{\rm Q}^\vee.
\end{gather*}

To specify $\bar\sigma$, put
$\bar N_{\mathbf M}=2c_{\mathbf M}c_{\rm lat}(\delta_{\mathbf M})$ and
\[
 s_{\rm st}(K,r)=\frac{r^2}{4KS_{\rm s}}
 \min\left\{\gamma^{\rm pc},
 \frac{\gamma_{\rm pos}}{4(1+\bar N_{\mathbf M}^2)},
 \frac{\gamma_{\rm pos}}{32(L+2)^3S_{\rm s}
                              \max(1,2\bar N_{\mathbf M}^2)}\right\}.
\]
For $(K,r)=(\mathsf K_0,r_0),(\mathsf K^0,r_P^\circ)$ let
\[
 C_{\rm kp}^+[K,r]=4c_\upsilon^2B_{\rm f}^{\rm res}C_{\rm E}
 e^{6\delta^{\rm S}M}Kr^{-2}K_{\rm ch}^\sharp e^{-\kappa''},
 \quad
 C_{\rm rj}^+[K,r]=C_{\rm kp}^+[K,r]
              \frac{C_3B_{\rm ch}^\sharp(LM)^6}{K_{\rm ch}^\sharp}.
\]
Write $\bar C_{\rm kp}^+,\bar C_{\rm rj}^+$ for their maxima over
the two choices, and take
\begin{equation}\label{III-108-43}
 \bar\sigma=\min\left\{
 \frac{\mathsf m_\odot}{4\mathsf k_\odot},
 s_{\rm st}(\mathsf K_0,r_0),s_{\rm st}(\mathsf K^0,r_P^\circ),
 \frac{c_{\rm kp}^{\rm th}}{2\bar C_{\rm kp}^+},
 \frac{c_{\rm rj}^{\rm th}}{2\bar C_{\rm rj}^+},
 \frac1{2C_{\rm xt}^2\bar C_{\rm kp}^+}\right\}.
\end{equation}
These are minima of finitely many positive numbers. The stability
ceilings therefore do not force the source radius to vanish.

\begin{proposition}\label{III-109-10}
Fix the preceding constants, and let $R>0$. There is $\gamma>0$
such that all the real numerical conditions hold and, uniformly for
$x\ge\gamma^{-2}$,
\begin{equation}\label{eq:source-radius-conditions}
 R\le\sigma^{\rm S}(x-\bar B),\qquad
 RF_{\mathfrak p_0}(x)\le\sigma_{\rm L}^{\rm S},\qquad
 F_a(x)=x^{-1/2}(\log x)^a.
\end{equation}
The following bounds can be imposed at the same choice:
\begin{gather}
 \delta(x),\ x^{-1/2}\mathfrak b(x)^3\le\epsilon_\times^+,
 \quad(C_{\mathcal A}+\theta_{\rm sl})\delta(x)
       \le\tfrac12\min(r_0,r_P^\circ),\notag\\
 S_0(x)\le\epsilon_\times^+,
 \quad (LM)^3e^{-\Gamma_P(x)}\le(6K_{\rm ap}^+)^{-1},
                                               \label{eq:source-uniform-coupling-bounds}\\
 \mathsf{sz}_{\rm rg}^{\rm S}(x),\mathsf{sz}_{\rm bd}^{\rm S}(x)
       \le\epsilon_{\rm ap}^{\rm S}.
 \notag
\end{gather}
We also impose the stronger regular-output reserve
\begin{equation}\label{eq:source-regular-output-reserve}
 K_{\mathbf M}\left(\delta(x)
 +\frac{(1+\beta_{\rm s})\delta(x)}{\alpha_0}N_\star
 +e^{-\gamma_2\mathfrak b(x)^2/2}\right)\le E_0/4
 \qquad(x\ge\gamma^{-2}).
\end{equation}
The same effective-potential bounds hold at the enlarged triple
$(\sqrt2E_0^{\rm S},\sqrt2B^{\rm S},\sqrt2B_{\rm fr}^{\rm S})$,
with the corresponding majorants used in the complex boundary
construction. Any prescribed fixed lower threshold on $\log\gamma^{-2}$
from the history-mass estimates can be included. Decreasing $\gamma$
preserves all these assertions.
\end{proposition}

\begin{proof}
Only the last, downward choice of the real construction remains.
The first condition in \eqref{eq:source-radius-conditions} holds when
$\gamma\le(\bar B+R/\sigma^{\rm S})^{-1/2}$. The second holds for
all sufficiently small $\gamma$, since
$F_{\mathfrak p_0}(x)\to0$. This is a condition on the entire interval
$[\gamma^{-2},\infty)$, so no monotonicity near its finite endpoint
is being presumed.

Every size in \eqref{eq:source-uniform-coupling-bounds} is bounded by
a finite sum of the functions $F_a$, negative powers of $\log x$,
and positive multiples of $e^{-c(\log x)^{2\mathfrak p_0}}$ with $c>0$.
The negative logarithmic power in the boundary size is
$(\log x)^{-(q-\mathfrak p_0)}$, where $q>\mathfrak p_0$.
The left side of \eqref{eq:source-regular-output-reserve} also
tends to zero. Thus the stronger ceiling $E_0/4$, used below,
can be imposed in this final choice.
The source sizes use the same expressions as the real sizes, with
$E_0,B,B_{\rm fr}$ replaced by the fixed triple in
\eqref{eq:source-final-box}. For the complex decomposition use the
triple multiplied by $\sqrt2$. The ten admission inequalities for
the effective-potential estimate contain none of these three sizes.
They therefore still hold. Its majorants at the enlarged triple tend
to zero, and making them at most $\epsilon_{\rm ap}^{\rm S}$ gives
the required estimates for both the enlarged and original triples.
Likewise $\Gamma_P(x)=C_{\rm mn}^{-1}(LM)^{-3}\theta_c^2
A_1^2(\log x)^{2\mathfrak p_0}\to\infty$.
There are finitely many requirements, all valid on a sufficiently
late interval. Their intersection with the real conditions and the
fixed history thresholds contains such an interval. This proves the
claim, including independence of $K,m$ and of any source function.
\end{proof}

\subsection{Uniform bounds at every step}

Let $u_\bullet=(x_0,\ldots,x_K)$ be an admissible sequence and set
\[
 s_0=R/(x_K-\bar B),\qquad R'=\sqrt3R/40.
\]
Use the source class of radius $R$ in physical lengths. The relative
weights obey $|v_k^w|_\infty<s_0$ and
$x_k\mathrm{sw}_k(v_k^w)<R'$ by \cref{III-16-31}.

\begin{proposition}\label{III-110-27}
At the preceding choices all numerical hypotheses of the general and
first sourced steps hold on their full analytic domains. More
precisely, for $t=k+1$, $0\le k<K$, both the centered member and its
one-lattice escort have size at most
\begin{equation}\label{III-108-52}
 U_k=\mathsf k_\odot
   [s_0(1+g_k\mathfrak b_k^3+x_t^{-1/2})+R'(g_k+\delta_k)]
 \le\tfrac34\min(\mathsf m_\odot^{\rm S},M^3\eta^{\rm S}).
\end{equation}
In particular the escort size is at most $\bar n_{\rm co}^{\rm S}$.
Writing $n_k$ for that size, the first-step reserve is explicitly
$n_0\le\bar n_{\rm co}^{\rm S}\le w_\star/4\le w_\star/2$.
The centered expansion and the boundary construction satisfy
\begin{gather}
 \varepsilon_{\mathsf w}^{\rm S}L^3\le(6K_{\rm ap}^+)^{-1},
 \quad b_{\rm cs}^{\rm S}\le\tfrac13,
 \quad b_{\rm rj}^{\rm S}+b_{\rm sr}^{\rm S}\le\tfrac12,
                                               \label{III-110-29}\\
 C_\Upsilon\beta_\upsilon^{(t)}\le\gamma_{\rm kp}/8,
 \quad b_{\rm fr}^{{\rm S},(t)}\le\tfrac18,
 \quad c_{\rm kp}^{\rm S}\le c_{\rm kp}^{\rm th}/16,
 \quad c_{\rm rj}^{\rm S}\le c_{\rm rj}^{\rm th}/16.
                                               \label{III-110-19}
\end{gather}
The regular one-lattice comparison has all its unperturbed reserves:
\begin{equation}\label{eq:source-unperturbed-reserves}
 n_{\rm int}\le3K_P\delta_k\le3w_\star/8,\qquad
 n_{\rm iv}\le K_\Lambda s_{\mathfrak a}(k)N_\star\le w_\star/4.
\end{equation}
The extra escort formed by absorbing the old regular difference has
norm at most
\begin{equation}\label{III-110-21}
 K_\Lambda s_{\mathfrak a}(k)\frac L{L-1}\frac{E_0}{4}
       +\bar n_{\rm co}^{\rm S}
 \le\min\{w_\star/2,E_0/(1404L^3K_\Pi)\}.
\end{equation}
Thus the one-lattice differences with coefficients
$33(3L^3-3)K_\Pi$ and $108(L^3-1)K_\Pi$ apply.
\end{proposition}

\begin{proof}
The table cut at $t$ is the same whether computed from $u_\bullet$
or from $u_\bullet|_t$. Its enlarged radii satisfy
$a_m^{\mathsf{Tab}}=(1+\beta_{\rm s})\alpha_{0,m}\le1$.
The weighted layer bound gives $a^\Sigma\le12L^4$.
Insert this bound in the general centered and escort estimates,
and $a_0^{\mathsf{Tab}}\le1$ in the first-step estimates.
The defining maxima of $C^{\rm src}$ and $C^{\rm src}_0$ bound
each coefficient. For the far terms use
$\mathfrak b_k\le g_k^{-1}$ and
$x_k\le(3L/2)x_t$; the threshold for $\check R_t$ gives
$g_k^{-1}e^{-8L\check R_t}\le\sqrt{3L/2}\,x_t^{-1/2}$ and
$g_k^{-2}e^{-8L\check R_t}\le(3L/2)x_t^{-1/2}$.
The two far contributions have coefficients at most
$13L^5C^{\rm src}$ and $36L^5C_{\mathcal A}C^{\rm src}$.
Their sum and the other coefficients are at most
$49L^5C_{\mathcal A}\max(C^{\rm src},C^{\rm src}_0)$.
This proves the bound by $U_k$ for the centered members, their
history restrictions, and the two one-lattice escorts.

For either $m=\mathsf m_\odot^{\rm S}$ or $m=M^3\eta^{\rm S}$,
the four terms of $U_k$ are bounded respectively by
$m/4$, $\epsilon_\times^+m/4$, $m/8$, and $m/4$.
Indeed $s_0\le\sigma^{\rm S}$,
$g_k\mathfrak b_k^3\le\epsilon_\times^+<1/2$ and
$x_t^{-1/2}\le1/2$, while
$R'(g_k+\delta_k)\le\sqrt3(1+A_1)\sigma_{\rm L}^{\rm S}/40$.
Equations \eqref{III-109-5} and \eqref{III-109-6} give the four
claims. Their sum is at most $3m/4$.

The ordinary source sizes are at most $\epsilon_\times^+$.
Thus their contribution to the amplitude, multiplied by $L^3$,
is at most $(24K_{\rm ap}^+)^{-1}$. The centered contribution is
at most $(16K_{\rm ap}^+)^{-1}$, since
$\mathsf m_\odot\le w_0\varepsilon^{\rm O}$.
Their sum is at most $(6K_{\rm ap}^+)^{-1}$.
Also $(LM)^3e^{-\gamma_P}\le(6K_{\rm ap}^+)^{-1}$, so the
centered cluster estimate yields $b_{\rm cs}^{\rm S}\le1/3$.

For the quadratic correction,
\[
 \beta_\upsilon^{(t)}
 =8C_{\rm E}e^{6\delta^{\rm S}M}
       \mathsf K^{(t)}(\mathsf R^{(t)})^{-2}s_0
 \le\tfrac18b_\upsilon^{\rm S}\sigma^{\rm S},
\]
where $(\mathsf K^{(t)},\mathsf R^{(t)})$ is either of the two
pairs used above. Substituting the corresponding four members of
\eqref{III-109-5} into the one-source-chain bounds proves
\eqref{III-110-19}. In particular each old kept or rejected
coefficient grows by at most a factor $3/2$. In the strip formula
the cluster coefficient can be any number in $[0,1/2]$; since
$N_{\rm b}^{\rm th}\ge1$, its source-independent part is also at
most $3/2$ times the old strip bound. The real bound
\eqref{II-202-32},
$b_{\rm rj}^{\rm th}+b_{\rm sr}^{\rm th}\le1/3$, therefore gives
$b_{\rm rj}^{\rm S}+b_{\rm sr}^{\rm S}\le1/2$.
The total new boundary budget is at most $5/6$.

We verify the one-lattice room separately. The real choices and
\eqref{eq:source-regular-output-reserve} give
\begin{gather*}
 N_\star\ge\frac L{L-1}(E_0+\tfrac32\alpha_0^2B_\sharp),\qquad
 s_{\mathfrak a}(k)\le s^\star
 \le\min\{w_\star/(8K_P),w_\star/(4K_\Lambda N_\star)\},\\
 K_{\mathbf M}\bigl(\delta_k+s_{\mathfrak a}(k)N_\star
              +e^{-\gamma_2\mathfrak b_k^2/2}\bigr)\le E_0/4,
 \qquad K_{\mathbf M}\ge243L^3K_\Pi K_\Lambda.
\end{gather*}
Here $\delta_k\le s_{\mathfrak a}(k)$ since
$\alpha_0\le1+\beta_{\rm s}$. These inequalities prove
\eqref{eq:source-unperturbed-reserves} and
\[
 K_\Lambda s_{\mathfrak a}(k)\frac L{L-1}\frac{E_0}{4}
 \le\min\{w_\star/16,E_0/(3888L^3K_\Pi)\}.
\]
Adding \eqref{eq:source-final-escort-ceiling} proves
\eqref{III-110-21}; the numerical inequality used is
$3888^{-1}+2808^{-1}<1404^{-1}$.
The unperturbed fourth family and the extra escort therefore have
sum of norms at most $3w_\star/4$. At the first step there is no old
regular fourth family. Its escort satisfies the still smaller bound
$\bar n_{\rm co}^{\rm S}\le w_\star/4$. The Gaussian hypotheses
and the full analytic core are the real ones, so the exact hypotheses
of \cref{III-34-30,III-47-46} now hold.
\end{proof}

\subsection{Application to the actual histories}



\begin{lemma}\label{III-105-3}\label{III-111-5}
Fix $R=2\varrho$ in the preceding construction and let the lattice
and coupling sequence be admissible for the real finite construction.
At every step, for every history, the numerical hypotheses used in
the sourced recursion hold with the data
\[
 (E_0^{\rm S},B^{\rm S},B_{\rm fr}^{\rm S},
       \bar b_{\rm fr}^{\rm S},\bar n_{\rm co}^{\rm S},
       \delta_\upsilon,\kappa',\theta_{\rm sl})
 =(E_0,B,B_{\rm fr}^{\flat}+1,1,M^3\eta^{\rm S},
       \delta^{\rm S},\kappa'',(1+\vartheta_*)/2).
\]
The conclusion is uniform in $K,m$ and in
$w\in\mathsf{Src}_{T_\ell}(2\varrho,2\varrho)$.
\end{lemma}

\begin{proof}
The real fixed-domain and finite-step hypotheses follow from
Proposition~\ref{II-202-29} and Theorem~\ref{II-203-9}. The sourced sizes and
rate conditions were established above, with the same constants at
every occurrence. In particular $\kappa'\ge\kappa>2$, and the
box conditions follow from \eqref{eq:source-final-box}.
The real regular majorant satisfies
\[
 \Xi(\gamma^{-2},N_\star)\le3E_0/4:
\]
its normalizer contribution is at most $E_0/2$, and the remaining
three terms are bounded by $E_0/4$ in the preceding proof.
This is the reserve required by the difference induction; the
weaker bound $E_0$ alone would not suffice.

The escort bound may also be read directly at its final ceiling:
\[
 U_k\le\left(\frac58+\frac{\sqrt3}{160}\right)M^3\eta^{\rm S}
 <M^3\eta^{\rm S}
 \le\min\{w_\star/4,E_0/(2808L^3K_\Pi)\}.
\]
It follows that
$108(L^3-1)K_\Pi\bar n_{\rm co}^{\rm S}\le E_0/26<E_0/4$.
For the general step, \eqref{III-110-21} gives the corresponding
bound with coefficient $33(3L^3-3)K_\Pi$.
The regular output room is bounded by
\[
 \Xi(x_k,N_\star)+27(L^3-1)K_\Pi\bar n_{\rm co}^{\rm S}
 \le E_0.
\]
The chain, formed-factor and damping requirements follow from
\eqref{III-110-18} and \eqref{III-110-19}, with
$\gamma'=\gamma_{\rm kp}$ and $\gamma^\dagger=\gamma'/2$.
All remaining real coercivity, strip and history-counting
inequalities retain their original values.

Finally \eqref{III-109-5} and \eqref{III-109-6} imply
\[
 2R\le\sigma_{\rm src}(x_K-\bar B),\qquad
 2K_{\rm L}^\vee R\gamma\le40\mu_*/\sqrt3.
\]
These are stronger than the two progression requirements at radius
$R$. Thus the source progression and the history recursion use one
common choice of constants.
\end{proof}

\begin{lemma}\label{III-105-7}\label{III-111-7}
For these constants the first-step and general-step constructions
apply on the actual source-independent domains and at every prefix
of an admissible history. Their tables, Gaussian and background
data, and analytic radii agree on overlapping prefixes. The source
class is a complex domain, stable under conjugation and under the
admitted translations, and
\[
 \mathsf a(w)=\frac{|w|_{T_\ell}}{x_K-\bar B}<\sigma^{\rm S}\le\frac12.
\]
The original density is bounded and continuous, and its exact partial
history measures and the evaluation maps have the domains used in
the sourced recursion.
\end{lemma}

\begin{proof}
The real construction supplies the same canonical table on each
prefix, the canonical shifts with their analytic and Borel
properties, and the Gaussian structures and strip maps on the full
domains. The preceding lemma verifies all additional numerical
conditions. The weighted Taylor estimates use
$(\delta_{\rm sp},\kappa^\sharp)$; the quadratic corrections use
$(\delta^{\rm S},\kappa'')$. We check their domain admission.
For a middle-tier triple, the real table's cube and natural-class
ceilings are those of Proposition~\ref{II-94-90}.
That proposition places the pair in the natural class and admits
the restriction to every union of source cubes to the pivot domain.
The canonical shift is consequently defined on the required middle
tier, and is Borel, holomorphic on its open slices, and hereditary
under restriction by Proposition~\ref{II-93-26}.
Use the enlarged radii
\[
 a_m=(1+\beta_{\rm s})\mathsf r_0((t,\mathrm{mid});m)\quad(m<k),
 \qquad
 a_k=(1+\beta_{\rm s})\max\{\mathsf r_0((t,\mathrm{mid});k),
                    L^{-2}\mathsf r_0((t,\mathrm{mid});t)\}.
\]
The layer bounds of the same table give membership in
$\mathbb U_k^a(Y|h')$, with $a_m\le1$ and $a^\Sigma\le12L^4$.
The field room imposed above admits both $g_k(z+\phi)$ and
$g_k\phi$ to the half-radius Taylor ball. Thus
\cref{III-39-43} applies on the full tiers; its quadratic
coefficient also has the domain stated in \cref{III-113-21}.

For the far-domain estimate, a bond in
$\mathfrak X^\circ\langle Y\rangle$ has an endpoint in
$\Omega_t$. Its containing fine cube lies in both $Y$ and
$\Omega_t$. If $Y\not\subset\Omega_k$, another cube of $Y$
lies in the closure of $\Omega_k^c$. The history separation
puts these cubes at distance at least $8LM\check R_t$.
An edge tree with $d_k(Y)$ edges meeting the cubes connects
two suitable points at distance at most $Md_k(Y)$.
Hence $d_k(Y)\ge8L\check R_t$, and therefore
$D\ge D'\ge8L\check R_t$. This supplies the far exponential
used in \eqref{III-108-52}.
At the first step the table's source-cube ceilings directly give
\eqref{eq:first-source-admission} for the pivot completion.
The permitted auxiliary
rate is $r_\upsilon=c_{\rm E}+\delta^{\rm S}M$, giving the required
rate inequality with equality; at later steps it is $r^\sharp$.
The centered family has one origin and no additional ordinary
origin. Its companion agrees on the small-field core by
\cref{III-41-24,III-48-54}. The centered expansion and boundary
construction consequently apply with the uniform bounds just proved.
The unit-weight subtraction and mean cancellation at an actual chain
point are justified by \cref{III-19-18}; in particular they use the
regional Hessian, not the Hessian on a constant tower.
The complex inherited and boundary families are the explicit
conjugate-pair extensions used in the recursion, so no real-valued
input theorem is being applied directly to a complex source.

For the relative size at every scale, the real flow equation gives
\[
 x_{j-1}-\bar B\ge L(x_j-\bar B),\qquad
 L^j(x_j-\bar B)\ge L^K(x_K-\bar B)>0.
\]
Thus $|v_j^w|_\infty\le\mathsf a(w)<\sigma^{\rm S}$, and the
source profile has real part at least half of its positive real
profile. The domain and symmetry assertions about sources, and the
bounded continuity of the density, follow from the source-class
construction. The partial history domains are contained in the
domains of the real coarse evaluation maps; they involve only the
history tests and charts. The exact complex measure construction
proved in the recursion section therefore applies on those domains.
Its common null set is supplied by the zero-source realization, and
is independent of $w$.

We record the conditions used by the history estimates. In the last
coupling choice put $y=\log\gamma^{-2}$ and
$\varepsilon_\gamma=A_0\gamma y^{\mathfrak p_0}$, and impose
\begin{gather*}
 y>2\mathfrak p_0,\quad A_0y^{\mathfrak p_0}\ge1,\quad
 (L^{1/r}-1)y\ge\log(2L),\\
 \varepsilon_\gamma\le
 \min\{B_3\bar\varepsilon,\bar\varepsilon/2,
                    \epsilon_{\rm xp},\varepsilon_\flat\},\quad
 c_1^{\rm P}\varepsilon_\gamma^2\le e_*,\quad
 c_2^{\rm Q}(A_1/A_0)^2\varepsilon_\gamma^2\le\sigma_{\rm Q},\\
 C_{\rm reg}\varepsilon_\gamma\le
 \min\{\varepsilon_{\rm v},L^{-2}c_2^{\rm SO},
               L^{-2}\bar\varepsilon/2,\varepsilon_{\mathsf z}/2\},\\
 2C_{\rm reg}\varepsilon_\gamma\le\varepsilon_{\mathsf z}^+,
 \quad c_{\rm lgs}M\varepsilon_\gamma
                     \le\tfrac12\min(\alpha_0,\alpha_2),\quad
 C_{\rm ly}\varepsilon_\gamma\le\alpha_0.
\end{gather*}
These requirements fit the final choice: $\varepsilon_\gamma\to0$,
and $C_{\rm reg}$ has the fixed upper bound used in the real choices.
Differentiation shows that
$\varepsilon(x)=A_0x^{-1/2}(\log x)^{\mathfrak p_0}$ is decreasing
for $\log x>2\mathfrak p_0$, so
$\varepsilon_t\le\varepsilon_\gamma$.
The flow inequality $x_{t-1}<2Lx_t$ and the third condition give
$\check R_{t-1}\le L\check R_t$.
Thus $p_0(g_t)\ge1$, $\varepsilon_t\le B_3\bar\varepsilon$,
$c_1^{\rm P}\varepsilon_t^2\le e_*$ and
$c_2^{\rm Q}\delta_t^2\le\sigma_{\rm Q}$.
These are the remaining majorant conditions
$(\mathrm R_{\rm maj})$. Together with the unchanged layer comparison
and torus room they give the fixed-radius hypotheses of
Proposition~\ref{II-101-34}. Its bound applies to the sourced sums
by the explicit complex extension in the history section.

The stronger positivity used in that section holds on actual
histories. The real local-gauge construction gives the depth-$k$
pair at every earlier strip component at radii
$(\alpha_0/2,\alpha_1,\alpha_2/2)$ by
Corollary~\ref{II-65-44}; the preceding smallness conditions are
the ones needed for this use. The local lower bound of
Proposition~\ref{II-102-10} therefore applies at each pair.
Every stored external bond belongs to exactly one strip component
in its large-field piece. Extending a piece field by zero and
summing the local bounds proves
\[
 Q_{t,D}(a;U')\ge2\gamma_{\rm kp}
            \sum_{b\in\mathfrak X^{\rm st+}(t,D)}|a(b)|^2.
\]
This is \eqref{II-195-49} at the realized pair. The depth-$k$ pair
is used at every earlier step, while the keeping predicates are
those of $h|_t$. The source-chain estimate is combined with this
bound only after the evaluation has been defined.

Use the numbers $C_5^{\rm S},c_*^{\rm S},C_{\rm maj}^{\rm S},
C_\flat^{\rm S}$ of \cref{III-25-3,III-27-5}, with the coefficient
bound $C_{\rm co}$ fixed by \eqref{eq:vacuum-coefficient-constant}. They satisfy
$C_5^{\rm S}=C_5$ and $c_*^{\rm S}=c_*/2$. In the entropy threshold put
$(c,C_{\rm maj},C_\flat)=(c_*^{\rm S},C_{\rm maj}^{\rm S},
C_\flat^{\rm S})$, $D_0=C_{\rm maj}+27(2+\log(2L))C_\flat$,
and $\lambda_L=(6L^3+1)\log2$. Then
\begin{align*}
 Y_0&=\max\{\log\max(e,2B_\sharp),2\log(2L)/(L^{1/r}-1)\},\\
 Y_{\rm b}^{\rm S}&=\max\biggl\{Y_0,
 \left(\frac4{\eth cA_0}\right)^{1/\mathfrak p_0},
 \left(\frac{2\cdot48^3(D_0+\lambda_L)L^3M^3}{\eth cA_0^2}
                         \right)^{1/(2\mathfrak p_0-3r-2)},\\
 &\hspace{4em}
 \left(\frac{3\log L+\log2}
       {2(1-\eth)c\mathfrak p_0A_0^2\log(L/2)}
                         \right)^{1/(2\mathfrak p_0-1)}\biggr\}.
\end{align*}
All denominators are positive, since $\eth\in(0,1)$, $c>0$,
$L\ge13$ and $2\mathfrak p_0-3r-2>0$. This number is finite and
fixed before $\gamma$. We include
$\log\gamma^{-2}\ge Y_{\rm b}^{\rm S}$ in the final choice;
then every $y_t$ meets the required entropy threshold.
The relative action loss has already been bounded by $1/2$.
\end{proof}


\section{History bounds and the source energies}
\label{sec:sourced-history-bounds}

We now use the constructed densities to estimate the partition
function. The complex source leaves half of the real action and
half of the retained Gaussian decay available. These two reserves
pay for the creation of a large-field cell. All estimates in this
section concern the fixed assignment and auxiliary data of
\cref{III-111-7}.

Write $N_k=(L^{K-k}\ell)^3$ and $y_j=\log x_j$. Set
\begin{equation}\label{III-25-3}
 C_5^{\rm S}=C_{\rm m}^{\rm sc}
             (E_0+B_\sharp+M^{-3}B^{\rm S}),
 \qquad
 \gamma^\dagger=\tfrac12\gamma_{\rm kp},\quad
 c_\star^{\rm S}=\tfrac12c_\star,
\end{equation}
where $C_{\rm m}^{\rm sc}$ is the constant in
equation \eqref{II-101-8}, and $c_\star,C_\star'$
are the constants in equation \eqref{II-77-16}. Define
\begin{equation}\label{III-27-5}
 \begin{split}
 C_\star^{\prime\rm S}&=C_\star'+\tfrac92L^3\log2,\\
 \mathsf C_{\rm maj}^{\rm S}
   &=\max\{C_5^{\rm S},L^3C_5^{\rm S}+C_\star^{\prime\rm S}\},\\
 C_\flat^{\rm S}
   &=\max\{3(L^3-1),(6L^3-5)C_{\rm co}
                         +(C_{{\rm E},0}+16C_{{\rm E},1})E_0\}.
 \end{split}
\end{equation}
Here $C_{\rm co}$ is the coefficient bound in
equation \eqref{eq:vacuum-coefficient-constant}, and
$C_{{\rm E},p}=2048\sum_{n\ge0}(n+2)^pe^{-n}$.
These constants are fixed before the source radius and the
coupling threshold.

\subsection{The effective density}

For a history $h\in\mathcal H_k$, let
$\mathcal S=\mathcal S(h)$ and let $\mathfrak q$ denote its
stored fields. The regular sum, boundary sum and label vacuum
constant are formed by the finite sums of
Sections \ref{sec:vacuum-bounds}--\ref{sec:scale-sum}, with $E^{(j),w}$ and the
sourced boundary families in place of the corresponding real
families. Put
\begin{equation}\label{III-111-13}
 A_{k,h}^w(\mathsf U,\mathfrak q)
 =-A(\beta_k^w,\mathsf U)
   +\mathbf E_k^w(\mathcal S;\mathsf U)
   +\mathbf E_k^{{\rm bd},w}(h;\mathsf U,\mathfrak q)
   -\mathcal E_{k,\mathcal S}^w,
 \qquad \beta_k^w=\beta_k-L^Kw.
\end{equation}
Evaluation at a history point requires that all terms be defined
on every adapted representative and have the same value there.
The retained sum $\boldsymbol\Phi_h^w$ has its own evaluation
requirement. These two requirements are imposed separately.

\begin{theorem}[The sourced effective densities]\label{III-111-9}
For every admissible finite coupling sequence, every torus with
the room required in \cref{III-111-7}, every $0\le k\le K$
and every $\|w\|_{\flat;T_\ell}<2\varrho$, the sourced
recursion gives finite complex history measures. Their sum is
the pushforward of $\rho_0^w\,dU$ by the $k$-fold block average.
For $k\ge1$ the density on the history variables is
\begin{equation}\label{III-111-12}
 \mathsf{Stay}_h^{\rm T}\chi_k(\Omega_k)
       \exp(\boldsymbol\Phi_h^w)\exp(A_{k,h}^w)
\end{equation}
with the zero-test convention. There is one reference null set
outside which, whenever $\mathsf{Stay}_h^{\rm T}\chi_k(\Omega_k)$
is nonzero, the two exponents are defined for all sources.
The local regular terms satisfy
\begin{equation}\label{III-111-14}
 \begin{split}
 |E^{(j),w}(X,P,z)|&\le E_0e^{-(\kappa+1)d_j(X)},\\
 |E^{(j),w}(X,P,z)-E^{(j),0}(X,P,z)|
 &\le\tfrac14E_0\frac{\|w\|_{\flat;T_\ell}}{2\varrho}
                         e^{-(\kappa+1)d_j(X)}.
 \end{split}
\end{equation}
The created boundary terms obey
$|B_j^{{\rm mk},w}(X)|\le B^{\rm S}e^{-\kappa d_j(X)}$;
the sum of the oscillations of retained exponents created at
$(i+1,X)$ is at most
$B_{\rm fr}^{\rm S}\mathfrak b_i^2e^{-\kappa d_{i+1}(X)}$.
The functions are Borel measurable, analytic on their stated
field domains, and holomorphic in the source. They satisfy their
respective gauge identities, with the stored fields transformed
jointly with the active fields for retained exponents.
Regular and created boundary terms depend on the source
only on $X^{\sim1}$, and a retained exponent only on its domain
$X_f$. They satisfy the conjugation laws of \cref{III-102-21}.
The regular terms satisfy the translation law of \cref{III-107-14}.
At zero source they agree with the real construction using the
same assignment. Restricting the construction to an earlier
level gives the construction for that prefix.

On every defined realized evaluation, the retained factors obey
\begin{equation}\label{III-26-23}
 |e^{\boldsymbol\Phi_h^w(\xi_h(U))}|
 \le\prod_{t=1}^k\prod_{c\in\Gamma_t^\partial(h)}
  |\det(L\mathsf m_c(\mathcal M^{t-1}\mathsf U))|^{-1}
 \exp\left(-\gamma^\dagger\sum_{(t,D)}
       \sum_{b\in\mathfrak X^{{\rm st}+}(t,D)}
                          |\mathsf a'_{t-1}(b)(U)|^2\right).
\end{equation}
Here $U\in\mathcal V_0[h]$, and $\mathsf U$ is any glued
representative. For every realized history point the regular
and boundary sums are defined and satisfy
\begin{equation}\label{III-111-20}
 |\mathbf E_k^w|+|\mathbf E_k^{{\rm bd},w}|
 \le C_5^{\rm S}\sum_{n=0}^k\#_n(\mathcal L_n)
 \le C_5^{\rm S}\left(N_k+L^3\sum_{j=1}^k\#_j(Z_j)\right).
\end{equation}
\end{theorem}

\begin{proof}
The simultaneous choice in \cref{III-111-7} gives the domains,
rate inequalities and size reserves needed at each step.
Apply \cref{III-102-21,III-107-14} and the exact measure
identities \cref{III-104-28,III-104-30}. At level zero the
density is $\rho_0^w$; the output of each step is exactly the
partial density required at the next. The inherited real
tests partition the finite bare measure, so
\cref{III-24-22} identifies the sum with the block pushforward.
The realization argument in \cref{III-104-28} supplies the
single source-independent null set. The formula for the action
is \eqref{III-111-13}; its counterterms and profiles use the
real coupling sequence.

The regular size and the uniform difference bound $E_0/4$
are \cref{III-107-14} at the reserves of \cref{III-111-7}.
For fixed $j,X,P,z$, apply \cref{III-16-27} to their
difference, which is holomorphic and vanishes at zero source.
This gives the second line of \eqref{III-111-14}.
The boundary, retained and covariance statements follow from
\cref{III-102-21}. Dependence only on the source in $X$
implies the stated dependence on the source in $X^{\sim1}$.
The retained quadratic estimate \cref{III-113-44}, applied
after the evaluation values have been established, gives
\eqref{III-26-23}. The determinant product is independent of
the glued representative: the pivot bonds partition into the
boundaries of the pieces, and their elimination factors are
the fixed coordinate Jacobians of the real history.

It remains to check the sum over scales. Use
Proposition \ref{II-101-34} with regular size $E_0$,
boundary size $B^{\rm S}$ and regular rate $\kappa+1$.
Its proof uses moduli of analytic local terms and their
quadratic vanishing at the trivial pair, not reality of their
values. The sourced families have the same gauge invariance
and analytic cores, so the same argument applies. Its
domain hypotheses hold on every realized representative:
the cores and boundary spaces at half radii are properties
of the real tests. For clarity, the summation has three
parts. The regular and action terms at a site in layer $n$
have bounds proportional to $L^{-4(n-j)}$; multiplication
by the $L^{3(n-j)}$ sites gives $L^{-(n-j)}$. Created
boundary domains meet $Z_j^{\sim}$, and their entropy sum
is at most $C'_{\rm f}M^{-3}\#_j(Z_j^{\sim})$.
Finally
\[
 \sum_{j=1}^k\#_j(Z_j)
 \le2\sum_{n=0}^k\#_n(\mathcal L_n),\qquad
 \#_n(\mathcal L_n)\le L^3\#_{n+1}(Z_{n+1})\quad(n<k).
\]
These are the geometric counts in that proposition and
give \eqref{III-111-20} with precisely \eqref{III-25-3}.
\end{proof}

\subsection{The modulus of a history contribution}

We give the loss calculation explicitly, since positivity of
the sourced measure is not available. Put $\theta=1/2$.
For a piece $(t,D)$, write $S=\sum_b|\mathsf a'(b)|^2$,
$\mathsf g=\nu e^{-\Pi}$, where $\nu>0$ and $\Pi\ge0$,
and $\zeta=|z_{t,D}^{\rm dsp}|$. The determinant estimate
of Proposition \ref{II-77-63} gives
\[
 0<\zeta\le A_D:=(2L^9)^{6\#_t(D)}.
\]
If $Q$ is a real form with $Q\ge2\theta\gamma_{\rm kp}S$,
then $Q'=\theta^{-1}Q$ satisfies the original lower bound.
The product bound of that proposition, applied to $Q'$, is
\[
 e^{-A_{\rm cr}}\prod_{(t,D)}
    \nu e^{-\Pi}\zeta e^{-Q'/2}
 \le e^{-c_\star\Sigma_{\rm cr}}
       \prod_{(t,D)} A_D\nu e^{-\Pi/2}e^{-\gamma_{\rm kp}S/4},
 \quad
 \Sigma_{\rm cr}=\sum_t p_0(g_t)^2\mathsf n_t^{\rm cr}(h).
\]
Raise it to the power $\theta$ and multiply by
$\prod(\nu e^{-\Pi}\zeta)^{1-\theta}$. Since
$\zeta\le A_D$ and $\Pi\ge0$, this gives
\begin{equation}\label{III-25-23}
 e^{-\theta A_{\rm cr}}\prod_{(t,D)}
                  \nu e^{-\Pi}\zeta e^{-Q/2}
 \le e^{-c_\star^{\rm S}\Sigma_{\rm cr}}
                  \prod_{(t,D)}\overline{\mathsf u}_{t,D}^{\rm S},
 \quad
 \overline{\mathsf u}_{t,D}^{\rm S}
       =A_D\nu e^{-\Pi/2}e^{-\gamma^\dagger S/4}.
\end{equation}
If a gauge penalty is infinite, the left side vanishes;
the preceding calculation applies to all other cases.
The piece integral uses at most $9L^3\#_t(D)$ real Gaussian
coordinates. Replacing $\gamma_{\rm kp}$ by
$\gamma^\dagger=\gamma_{\rm kp}/2$ multiplies its bound by
at most $2^{(9/2)L^3\#_t(D)}$. Thus
\begin{equation}\label{III-25-24}
 \int\overline{\mathsf u}_{t,D}^{\rm S}\,dm_{t,D}
          \le e^{C_\star^{\prime\rm S}\#_t(D)}.
\end{equation}
All curvature and failed-test hypotheses here are those of
the real history. They hold by \cref{III-111-7}; only the
quadratic lower bound has been replaced.

\begin{proposition}\label{III-26-27}
For every $h\in\mathcal H_k$, $1\le k\le K$, and every
source in the ball,
\begin{equation}\label{III-111-23}
 e^{\operatorname{Re}\mathcal E_{k,\mathcal S(h)}^w}
 |\hat\tau_h^{\mathrm{co}\prime}[\rho_0^w]|
                      (\mathsf X_h^{\mathrm{co}\prime})
 \le
 e^{\mathsf C_{\rm maj}^{\rm S}(N_k+\sum_{j=1}^k\#_j(Z_j))}
       e^{-c_\star^{\rm S}\Sigma_{\rm cr}}.
\end{equation}
\end{proposition}

\begin{proof}
Pass to the free history coordinates, using the mutually
inverse coordinate maps on their supports from
Section \ref{sec:history-mass}. Pulling back the density
\eqref{III-111-12} gives an integrable complex Borel density
$f$ relative to the free product measure. Extend it by zero
on the reference null set. The two coordinate pushforwards
are inverse on the supports, so total variation has the
same mass in the two coordinate systems.

The real history measure is concentrated on a Borel set
of realized free coordinates. The sourced measure is
dominated in total variation by
$\sup|\rho_0^w|/\min\rho_0$ times that measure, by the
bare integral formula of \cref{III-24-22}; the denominator
is positive by compactness. It therefore vanishes on
every Borel subset of the complement. Since $f$ is its
density, $|f|$ has zero integral on that complement.
Intersect with the tiers and remove the reference null
set. At every remaining point the nonzero test implies
that both exponents are defined. Choose a bare preimage
and any adapted representative at that point.

By \cref{III-24-8} and \cref{III-111-7},
$\operatorname{Re}\beta_k^w\ge\theta\beta_k>0$.
The representative is real and $a_p(\mathsf U)\ge0$;
hence
$\operatorname{Re}A(\beta_k^w,\mathsf U)
 \ge\theta A(\beta_k,\mathsf U)\ge\theta A_{\rm cr}$.
Combine \eqref{III-111-13}, \eqref{III-111-20} and
\eqref{III-26-23}, using on each piece the real form
$Q=2\gamma^\dagger S$. The stay factors provide
$\nu e^{-\Pi}$, and the determinants provide $\zeta$.
Equation \eqref{III-25-23} then gives
\[
 |f(v)|\le
 e^{-\operatorname{Re}\mathcal E_{k,\mathcal S}^w
       +C_5^{\rm S}\sum_n\#_n(\mathcal L_n)
       -c_\star^{\rm S}\Sigma_{\rm cr}}
 \prod_{(t,D)}\overline{\mathsf u}_{t,D}^{\rm S}
                                      (\operatorname{pr}_{t,D}v).
\]
The right side depends only on the free coordinates;
no measurable choice of the bare preimage is required.
Different pieces use disjoint coordinate factors, every
Lebesgue coordinate belongs to one piece, and the remaining
coordinates carry Haar probability measure. Tonelli and
\eqref{III-25-24}, equivalently the product calculation in
equation \eqref{II-195-18}, bound the integral of this
product by $e^{C_\star^{\prime\rm S}\sum_j\#_j(Z_j)}$.
Finally use \eqref{III-111-20} and \eqref{III-27-5}.
\end{proof}

\subsection{Vacuum deficits and the sum over histories}

Define the sourced vacuum data by
$\mathsf{vl}^{(j),w}(X,z)=E^{(j),w}(X,\mathbf1,z)$,
keeping the Gaussian constants $\mathsf{vn}^{(j)}$ and
the extensive coefficients of Section \ref{sec:vacuum-bounds}.
Use the finite formulas \eqref{II-97-4}--\eqref{II-97-8} to define
$E_{(k)}^w$, $\mathcal E_{k,\mathcal S}^w$ and
$\mathfrak E_n^{\prime w}(\mathcal S)$. The two ranges of
base points in those formulas are unchanged: regular
terms use the interior base points, while Gaussian
constants use every base point of the domain.

\begin{proposition}\label{III-27-18}
All these vacuum numbers are defined. For every label
$\mathcal S$ of level $k$,
\begin{equation}\label{III-111-22}
 E_{(k)}^w+\mathcal E_{k,\mathcal S}^w
       =\sum_{n=1}^k\mathfrak E_n^{\prime w}(\mathcal S),
 \qquad
 |\mathfrak E_n^{\prime w}(\mathcal S)|
   \le C_\flat^{\rm S}(1+y_{n-1})\#_n(Z_n^{\sim}).
\end{equation}
Put
\begin{equation}\label{eq:sourced-vacuum-difference}
 e_j(X;w)=\sum_{z\in X\cap\mathbb T^{(j)}}
       [E^{(j),w}(X,\mathbf1,z)-E^{(j),0}(X,\mathbf1,z)].
\end{equation}
Then
\begin{equation}\label{eq:sourced-vacuum-splitting}
 E_{(k)}^w=E_{(k)}+\sum_{j=1}^k\sum_{X\in\mathcal D_j}e_j(X;w),
 \qquad
 \widetilde Z_k(w):=e^{-E_{(k)}^w}Z_{K,\ell,\beta}[w]
\end{equation}
defines a holomorphic function with
\begin{equation}\label{eq:sourced-stability-all-levels}
 |\widetilde Z_k(w)|
 \le\sum_{h\in\mathcal H_k}e^{-\operatorname{Re}E_{(k)}^w}
       |\hat\tau_h^{\mathrm{co}\prime}[\rho_0^w]|
                     (\mathsf X_h^{\mathrm{co}\prime})
 \le e^{(\mathsf C_{\rm maj}^{\rm S}+2)N_k}.
\end{equation}
The functions $e_j$ satisfy \eqref{III-111-21} with
$\kappa_e=\kappa$, vanish at zero source and depend on
the source only on $X^{\sim1}$. They and $\widetilde Z_k$
satisfy conjugation reality.
\end{proposition}

\begin{proof}
The trivial pair belongs to every fixed analytic core.
The local vacuum values are thus defined, and their
moduli have the same bound as in the real construction.
We repeat the part of Proposition \ref{II-192-21} that
does not require their reality. The domain entropy gives
\[
 \sum_{X\ni z}|\mathsf{vl}^{(j),w}(X,z)|\le C_{{\rm E},0}E_0,
 \qquad
 \sum_{X\supset\square}\sum_{z\in X}
       |\mathsf{vl}^{(j),w}(X,z)|
           \le8M^3C_{{\rm E},1}E_0.
\]
The same estimates hold for $\mathsf{vn}^{(j)}$.
In a restricted vacuum sum, either a missing domain
contains a cube of $Z_n$, or its domain lies in
$\Lambda_n$ and its regular base point belongs to
$Z_n^{\sim}\setminus Z_n$. The cube estimate bounds
each of the two vacuum families in the first case;
the site estimate bounds the second case. The
extensive logarithmic part is bounded by
$3(L^3-1)y_{n-1}\#_n(\Omega_n^c)$, and its remaining
coefficients by $(6L^3-5)C_{\rm co}\#_n(Z_n)$.
Adding these bounds proves the inequality in
\eqref{III-111-22}. Its identity is the telescoping
sum of the defining deficits, valid for complex values.

At zero source the data are the canonical real data.
The source changes only $\mathsf{vl}$, so subtraction
of the finite whole-torus sums gives
\eqref{eq:sourced-vacuum-splitting}. By
\eqref{III-111-14} and the geometric bound
$\#_j(X)\le8M^3(d_j(X)+1)$,
\[
 |e_j(X;w)|\le
 2M^3E_0\frac{\|w\|_{\flat;T_\ell}}{2\varrho}
        (d_j(X)+1)e^{-(\kappa+1)d_j(X)}.
\]
Use $(d+1)e^{-d}\le1$ to obtain the asserted decay.
Holomorphy, locality and conjugation follow by finite
summation. The same properties for $\widetilde Z_k$
follow from \cref{III-16-26} and its defining exponential.

To sum the histories set
\[
 f(h)=e^{\operatorname{Re}\mathcal E_{k,\mathcal S(h)}^w}
       |\hat\tau_h^{\mathrm{co}\prime}[\rho_0^w]|
                             (\mathsf X_h^{\mathrm{co}\prime}),
 \qquad
 e(\mathcal S)=-\operatorname{Re}\sum_n
                          \mathfrak E_n^{\prime w}(\mathcal S).
\]
Proposition~\ref{III-26-27} gives the creation bound for
$f$, and \eqref{III-111-22} gives the deficit bound for
$e$. The product $f(h)\exp(e(\mathcal S(h)))$ is exactly
$e^{-\operatorname{Re}E_{(k)}^w}
 |\hat\tau_h^{\mathrm{co}\prime}[\rho_0^w]|
 (\mathsf X_h^{\mathrm{co}\prime})$.
Apply Proposition \ref{II-75-111} to this nonnegative
function and this real deficit, with constants
$(c_\star^{\rm S},\mathsf C_{\rm maj}^{\rm S},C_\flat^{\rm S},\eth)$.
Its hypotheses $2\mathfrak p_0>3r+2$ and the corresponding
threshold $y_K\ge\mathsf Y_{\rm b}$ are included in
\cref{III-111-7}. It bounds the sum by
\[
 e^{\mathsf C_{\rm maj}^{\rm S}N_k}
 \prod_{t=1}^k\prod_{\mathfrak c\in\mathcal P_t}
       (1+e^{-(1-\eth)c_\star^{\rm S}p_0(g_t)^2})
 \le e^{(\mathsf C_{\rm maj}^{\rm S}+2)N_k}.
\]
The triangle inequality and the finite history
decomposition give the first inequality of
\eqref{eq:sourced-stability-all-levels}.
\end{proof}

\section{Proof of the source factorisation theorem}
\label{sec:proof-source-theorem}

\begin{proof}[Proof of \cref{III-111-27}]
Choose the constants as in \cref{III-111-7}, and put
$\kappa_e=\kappa$. They precede the choice of $\varrho$.
After fixing that radius, choose $\gamma_\varrho$ and
the canonical auxiliary data at that threshold.
For $0<g\le\gamma_\varrho$ and $\beta>X_g$,
\cref{III-5-11} gives a unique finite positive first-passage
index and $g^{-2}<x_K\le X_g$. The side condition
\eqref{III-1-14} supplies all the lattice room required
in \cref{III-111-7}; arbitrarily large integral $m$ satisfy it.

Apply \cref{III-111-9,III-27-18} at $k=K$.
The real constant $E_{(K)}$ is $E_{K,\ell,\beta}$ by
\eqref{III-7-7}. The definition
$\widetilde Z=\widetilde Z_K$ and the vacuum splitting
give \eqref{III-111-28}. The local source estimates,
holomorphy and conjugation laws are those of
\cref{III-27-18}. At $k=K$ we have $N_K=\ell^3$,
so its stability estimate bounds $|\widetilde Z|$ by
$e^{(\mathsf C_{\rm maj}^{\rm S}+2)\ell^3}$.

At zero source the local source energies vanish, and
\cref{III-7-9} gives
\[
 \widetilde Z(0)=e^{-E_{K,\ell,\beta}}Z_{K,\ell,\beta}
 \ge e^{-[c_{\rm low}+(9/2)\log x_K]\ell^3}.
\]
The threshold choice includes
$B_\sharp\le\tfrac12\gamma_\varrho^{-2}$; therefore
$X_g\le(3L/2)g^{-2}$. Choose $C_{\mathrm S}$ at least
\[
 \max\{2M^3E_0,\mathsf C_{\rm maj}^{\rm S}+2,
        |c_{\rm low}|+\tfrac92\log(3L/2),\tfrac92,
        2C_{\rm tt}C'_{\rm tt}E_0\}.
\]
Since $g\le1/2$, this yields the lower bound
\eqref{III-111-29}. For two allowed tori, the sequence
and assignment are the same; \cref{III-7-12} gives
the asserted bound on their normalized vacuum constants.
Every member of the displayed maximum precedes the
source radius and threshold. This proves the stated
uniformity and completes the proof.
\end{proof}

For a stopped datum from \cref{III-5-21}, the same
argument applies to $((L^i\ell,K-i),u_\bullet|_{K-i})$
with reference coupling $g_{(i)}<g$ and the same threshold
assignment. The constants and source radius are unchanged.
This is an application on the stopped torus; it requires
no identification of sources under dilation. The affine
and mirror-periodic source families of \cref{III-17-18}
lie in the same ball, so the theorem applies to them as well.


\section{Plaquette observables and reflection positivity}
\label{sec:plaquette-reflection}

The source derivatives give correlations of a scalar observable.
To estimate them uniformly in the volume, we first prove the reflection
inequality for the Wilson measure. Rectangular tori are needed here:
the comparison changes one side at a time.

Let $a>0$, let $N_\mu\ge3$ be integers, and put
\[
 T=\prod_{\mu=1}^3\mathbb R/(aN_\mu\mathbb Z),\qquad
 \Lambda=\prod_{\mu=1}^3 a\mathbb Z/(aN_\mu\mathbb Z).
\]
Positive bonds and plaquettes have the usual coordinate orientations.
As before, $U(\bar b)=U(b)^{-1}$ and
$a_p(U)=1-\tfrac12\operatorname{tr}\partial U(p)$.
For $\beta\ge0$ define
\begin{equation}\label{IV-1-9}
 A(U)=\sum_p a_p(U),\qquad
 \mathcal O_a(f)=a^{-1}\sum_p f(x(p))a_p(U),\qquad
 d\mu_{a,T,\beta}=Z_{a,T,\beta}^{-1}e^{-\beta A(U)}\,dU.
\end{equation}
Here $dU$ is product Haar probability measure and $Z_{a,T,\beta}$
is the integral of $e^{-\beta A}$. Since $0\le a_p\le2$,
this normalization is strictly positive, and $\mathcal O_a(f)$ is
bounded for every function $f$ on the finite set of plaquette centres.
For a bounded Borel function $F$ of the configuration, write
\[
 \mathcal Z_{a,T,\beta}(F)=\int F(U)e^{-\beta A(U)}\,dU,
 \qquad \langle F\rangle_{a,T,\beta}
 =Z_{a,T,\beta}^{-1}\mathcal Z_{a,T,\beta}(F).
\]
Thus, on the cubic lattice with $a=L^{-K}$,
$Z_{K,\ell,\beta}[w]=\mathcal Z_{a,T_\ell,\beta}
(e^{\mathcal O_a(w)})$.

For $n\ge1$ the joint cumulant is
\begin{equation}\label{IV-1-23}
 S^T_{n;a,T,\beta}(f_1,\ldots,f_n)
 =\left.\partial_{z_1}\cdots\partial_{z_n}
 \log\left\langle e^{\sum_i z_i\mathcal O_a(f_i)}
 \right\rangle_{a,T,\beta}\right|_{z=0}.
\end{equation}
The expectation is entire and equals one at zero, so the logarithm
in this formula exists near zero. If $\mathcal P_n$ denotes the set
of partitions of $\{1,\ldots,n\}$, differentiating the logarithm gives
\begin{equation}\label{IV-1-25}
 S^T_n(f_1,\ldots,f_n)=
 \sum_{\pi\in\mathcal P_n}(-1)^{|\pi|-1}(|\pi|-1)!
 \prod_{B\in\pi}\left\langle\prod_{i\in B}\mathcal O_a(f_i)
 \right\rangle.
\end{equation}
Conversely, the joint moment is the sum over $\pi\in\mathcal P_n$
of the products of the cumulants indexed by its blocks. Both identities
follow by comparing the coefficient of $z_1\cdots z_n$ in the
exponential identity relating the expectation to its logarithm.
In particular, cumulants are symmetric and multilinear. For $n\ge2$
the cumulant of the random variables $\mathcal O_a(f_i)$ is unchanged
if a scalar constant is added to any of these random variables:
such a change adds a linear function of $z$ to the logarithm.

A lattice translation or coordinate reflection acts on configurations
by transporting oriented bonds. It permutes the independent Haar
variables, possibly replacing some by their inverses. Product Haar
measure is unchanged. The plaquette holonomy is carried to a conjugate
of the original holonomy or its inverse, both of which have the same
trace in $\mathrm{SU}(2)$. Hence the Wilson measure is invariant,
and $\mathcal O_a(f)$ transforms by the corresponding pullback of $f$.

\subsection{The operator across one layer}

Fix a coordinate direction, write $N$ for the number of sites in
that direction, and let $\Sigma$ be the two-dimensional cross-section.
Let $\mathsf B(\Sigma)$ be its positive bonds and set
$\mathcal X=\mathrm{SU}(2)^{\mathsf B(\Sigma)}$.
A configuration consists of the transverse fields $V^t\in\mathcal X$
and the bonds $S^t\in\mathrm{SU}(2)^\Sigma$ joining layers $t$ and
$t+1$, with $t$ taken modulo $N$. Write $A_\Sigma(V)$ for the sum
of the plaquette actions within a transverse layer and put
\begin{equation}\label{IV-3-6}
 A_\times(V,S,W)=\sum_{b\in\mathsf B(\Sigma)}
 \left(1-\tfrac12\operatorname{tr}
       V_bS_{b_+}W_b^{-1}S_{b_-}^{-1}\right).
\end{equation}
The two types of plaquettes partition all plaquettes, giving
\[
 A(U)=\sum_{t\bmod N}
       \bigl[A_\Sigma(V^t)+A_\times(V^t,S^t,V^{t+1})\bigr].
\]
On $\mathcal H=L^2(\mathcal X,dV)$ define the integral operator
\[
 \mathsf K(W,V)=\int e^{-\beta A_\times(V,S,W)}\,dS.
\]
The kernel is continuous, symmetric and bounded by one. Symmetry follows
by replacing $S$ by $S^{-1}$ and using invariance of the trace under
cyclic permutations and inversion. Thus $\mathsf K$ is self-adjoint
and Hilbert--Schmidt. Its positivity needs a separate argument.

\begin{proposition}\label{IV-3-16}
The crossing operator satisfies $\mathsf K\ge0$.
\end{proposition}
\begin{proof}
Put $\mathsf k(u)=e^{-\beta(1-\operatorname{tr}u/2)}$ and
$\mathsf K_0(V,W)=\prod_b\mathsf k(V_bW_b^{-1})$.
For $R\in\mathrm{SU}(2)^\Sigma$ let
$(V^R)_b=R_{b_-}V_bR_{b_+}^{-1}$.
The crossing weight is $\mathsf K_0(V,W^S)$, and
$\mathsf K_0(V^R,W^R)=\mathsf K_0(V,W)$. Haar invariance therefore gives
\begin{equation}\label{IV-3-18}
 \mathsf K(W,V)=\int dR\,dS\,\mathsf K_0(V^R,W^S).
\end{equation}
For $v,w\in\mathrm{SU}(2)$,
$\operatorname{tr}(vw^{-1})=\sum_{i,j}v_{ij}\overline{w_{ij}}$.
Expanding the exponential at every bond consequently yields
\[
 \mathsf K_0(V,W)=\sum_\alpha c_\alpha
      \phi_\alpha(V)\overline{\phi_\alpha(W)},\qquad c_\alpha\ge0.
\]
Here $\alpha$ specifies a nonnegative integer $n_b$ and $n_b$ pairs
of matrix indices at each bond,
\[
 c_\alpha=e^{-\beta|\mathsf B(\Sigma)|}
       \prod_b\frac{(\beta/2)^{n_b}}{n_b!},\qquad
 \phi_\alpha(V)=\prod_b\prod_{s=1}^{n_b}(V_b)_{i_{b,s}j_{b,s}}.
\]
Every matrix entry has modulus at most one. There are $4^{n_b}$
choices of indices at bond $b$, so
\[
 \sum_\alpha c_\alpha
 =e^{-\beta|\mathsf B(\Sigma)|}
       \prod_b\sum_{n\ge0}\frac{(2\beta)^n}{n!}
 =e^{\beta|\mathsf B(\Sigma)|}<\infty.
\]
This proves absolute uniform convergence. Insert the series into
\eqref{IV-3-18} and set
$\widehat\phi_\alpha(V)=\int\phi_\alpha(V^R)\,dR$. Then
\[
 \mathsf K(W,V)=\sum_\alpha c_\alpha
       \widehat\phi_\alpha(V)\overline{\widehat\phi_\alpha(W)}.
\]
Since $dV$ is a probability measure, every $\psi\in\mathcal H$
is integrable. Absolute convergence permits termwise integration, giving
\[
 \langle\psi,\mathsf K\psi\rangle
   =\sum_\alpha c_\alpha
       \left|\int\widehat\phi_\alpha(V)\psi(V)\,dV\right|^2\ge0.
\]
\end{proof}

\subsection{Slabs and the trace formula}

For an integer $q\ge2$, a slab of $q$ layers has variables
$(V^1,\ldots,V^q;S^1,\ldots,S^{q-1})$ and weight
\[
 w_q=\exp\left[-\beta\sum_{t=1}^q A_\Sigma(V^t)
       -\beta\sum_{t=1}^{q-1}A_\times(V^t,S^t,V^{t+1})\right].
\]
For a bounded Borel function $\Phi$ of these variables define
\begin{equation}\label{IV-3-10}
 \mathfrak S[\Phi](W,V)=
 \int\prod_{1<t<q}dV^t\prod_{1\le t<q}dS^t\,
       (\Phi w_q)\big|_{V^1=V,V^q=W},\qquad
 \mathfrak Y[\Phi]=\mathsf K^{1/2}\mathfrak S[\Phi]\mathsf K^{1/2}.
\end{equation}
These are Hilbert--Schmidt operators, since the first kernel is bounded
by $\|\Phi\|_\infty$ and $\mathsf K^{1/2}$ is bounded.
Reverse the order of the layers, invert the joining bonds, and conjugate
the value of $\Phi$; denote the resulting function by $\Phi^\Theta$.
This change of variables preserves Haar measure and $w_q$, and exchanges
the two endpoint fields. Consequently
\begin{equation}\label{IV-3-25}
 \mathfrak S[\Phi^\Theta]=\mathfrak S[\Phi]^*,\qquad
 \mathfrak Y[\Phi^\Theta]=\mathfrak Y[\Phi]^*.
\end{equation}

\begin{proposition}[Trace formula]\label{IV-3-26}
Partition the $N$ layers cyclically into $r\ge2$ consecutive slabs
of lengths $q_s\ge2$. Let $\Phi_s$ be a bounded Borel function
of the inner bonds of slab $s$. Then
\begin{equation}\label{IV-3-27}
 \mathcal Z_{a,T,\beta}\left(\prod_{s=1}^r\Phi_s\right)
 =\operatorname{Tr}\bigl(\mathfrak S[\Phi_r]\mathsf K\cdots
                  \mathfrak S[\Phi_1]\mathsf K\bigr)
 =\operatorname{Tr}\bigl(\mathfrak Y[\Phi_r]\cdots
                  \mathfrak Y[\Phi_1]\bigr).
\end{equation}
All the operators attached to a slab depend on its length, the
cross-section and $\beta$, but not on the total length $N$.
\end{proposition}
\begin{proof}
Keep the first and last transverse fields $V_s,W_s$ of each slab.
Every plaquette either lies inside one slab or crosses from its last
layer to the first layer of the next slab. Integrating all other
variables gives
\[
 \int\prod_{s=1}^r dV_s\,dW_s\,
       \prod_{s=1}^r\mathfrak S[\Phi_s](W_s,V_s)
                         \mathsf K(V_{s+1},W_s),\qquad V_{r+1}=V_1.
\]
The integrands are bounded and all measures are finite, so Fubini
applies. The kernel formula for a product, followed by the
Hilbert--Schmidt identity
$\operatorname{Tr}(CD)=\iint C(y,x)D(x,y)\,dx\,dy$,
gives the first trace. There are at least two Hilbert--Schmidt factors,
so the product is trace class. Splitting each $\mathsf K$ into two
square roots and moving the last square root around the trace gives
the second identity. The remaining factors form a trace-class operator,
which justifies this cyclic permutation. The definitions contain no
total layer number $N$, proving the last assertion.
\end{proof}

\begin{corollary}[Reflection at bond midpoints]\label{IV-4-8}
Suppose that $N$ is even. Let $\theta$ be reflection in a coordinate
plane at a half-integer multiple of $a$, and let $T_+$ be either open
half-torus bounded by this plane and its opposite parallel plane.
For bounded Borel functions $F,H$ depending only on bonds whose two
endpoints lie in $T_+$, the form
$\langle\overline{\theta F}H\rangle_{a,T,\beta}$ is positive
semidefinite. In particular,
\begin{equation}\label{IV-4-10}
 |\langle\overline{\theta F}H\rangle|^2
 \le\langle\overline{\theta F}F\rangle
       \langle\overline{\theta H}H\rangle.
\end{equation}
\end{corollary}
\begin{proof}
The two halves contain $q=N/2\ge2$ layers each. Their inner bond
variables identify $F,H$ with slab functions $\Phi,\Psi$.
Reflection sends the first slab to the second and reverses its layer
order, so $\overline{\theta F}$ is the placed function $\Phi^\Theta$.
By \eqref{IV-3-25} and \cref{IV-3-26},
\[
 \mathcal Z_{a,T,\beta}(\overline{\theta F}H)
 =\operatorname{Tr}\bigl(\mathfrak Y[\Phi]^*\mathfrak Y[\Psi]\bigr).
\]
This is the Hilbert--Schmidt inner product. Its positivity and
Cauchy--Schwarz inequality survive division by $Z_{a,T,\beta}>0$.
\end{proof}

\section{The chessboard estimate and comparison of volumes}
\label{sec:chessboard}

The reflection estimate bounds a local insertion by its reflected copies.
We need the resulting bound on any even number of cells, and we need
to compare it on two tori. The trace formula makes both statements
consequences of inequalities for powers of one positive operator.

\subsection{Cells and reflected copies}

Fix $a>0$, an odd integer $\lambda\ge3$, even integers
$\mathsf n_\mu\ge2$, and an offset $o\in a\mathbb Z^3$.
We call these data a board. Its torus has sides
$\ell_\mu=\mathsf n_\mu\mathsf l$, where $\mathsf l=\lambda a$.
Write
\begin{equation}\label{IV-2-3}
 \mathsf Z=\prod_{\mu=1}^3\mathbb Z/\mathsf n_\mu\mathbb Z,
 \qquad \mathsf N=\prod_\mu\mathsf n_\mu,
 \qquad h_\lambda=(\lambda-1)/2.
\end{equation}
The cell $\square_\kappa$ is the image in the torus of
$o+\mathsf l\tilde\kappa+[-\mathsf l/2,\mathsf l/2]^3$, with
$\tilde\kappa$ any integer representative of $\kappa\in\mathsf Z$.
Its sites are
$o+\mathsf l\tilde\kappa+an$, $n\in\mathsf C_\lambda
:=\{-h_\lambda,\ldots,h_\lambda\}^3$.
The cell faces contain no sites: their varying coordinate is a
half-integer multiple of $a$. Thus the cells partition the sites,
their interiors are disjoint, and a plaquette whose centre lies in
a cell interior has all its corners in that cell. A plaquette that
is not of this kind has its centre on a cell face.

A function of a cell is a bounded complex Borel function of its
inner bond variables, where an inner bond has both endpoints in
the cell. Transporting the positive bonds of $\mathsf C_\lambda$
to $\square_\kappa$ identifies such a function with a unique model
function $F_0$ on the finite product
$\mathrm{SU}(2)^{\mathsf B^+(\mathsf C_\lambda)}$; its placement
in $\square_\kappa$ is denoted by $F_0^{(\kappa)}$.
Uniqueness follows because arbitrary values of the inner bond
variables extend to a configuration on the torus. This also shows
that placing a model function preserves its supremum norm.

For $\kappa,\kappa'\in\mathsf Z$ set
$\epsilon_\mu=(-1)^{\tilde\kappa'_\mu-\tilde\kappa_\mu}$ and define
the diagonal isometry
\[
 \mathsf r_{\kappa'\kappa}(x)
 =o+\mathsf l\tilde\kappa'
    +\epsilon(x-o-\mathsf l\tilde\kappa).
\]
It carries $\square_\kappa$ to $\square_{\kappa'}$.
The evenness of each $\mathsf n_\mu$ makes the signs independent
of the representatives; changing a representative changes the
formula only by a torus period. Transport configurations and
observables by the action used in \cref{sec:plaquette-reflection},
including inversion when a bond orientation is reversed.
For a function $F$ of $\square_\kappa$, define
\begin{equation}\label{IV-2-26}
 F^{[\kappa']}=
 \begin{cases}
  \mathsf r_{\kappa'\kappa}F,&\prod_\mu\epsilon_\mu=1,\\
  \overline{\mathsf r_{\kappa'\kappa}F},&\prod_\mu\epsilon_\mu=-1,
 \end{cases}
 \qquad
 \operatorname{Fold}F=\prod_{\kappa'\in\mathsf Z}F^{[\kappa']}.
\end{equation}
These are functions of the indicated cells. The composition rule
$\mathsf r_{\kappa''\kappa'}\mathsf r_{\kappa'\kappa}
=\mathsf r_{\kappa''\kappa}$, together with multiplication of the
parity signs, gives
\begin{equation}\label{IV-2-30}
 (F^{[\kappa']})^{[\kappa'']}=F^{[\kappa'']},
 \qquad \operatorname{Fold}(F^{[\kappa']})
 =\operatorname{Fold}F.
\end{equation}
Copies preserve products and sums and commute with conjugation.
Exactly half the cells have each parity, since changing one cell
index by one interchanges the parities. Consequently
\begin{equation}\label{IV-2-31}
 \operatorname{Fold}(cF)=|c|^{\mathsf N}\operatorname{Fold}F,
 \qquad
 \operatorname{Fold}\bar F=\overline{\operatorname{Fold}F},
 \qquad
 |\operatorname{Fold}F|\le\|F\|_\infty^{\mathsf N}.
\end{equation}

\subsection{Multiple reflections}

We first give the trace inequality needed for an arbitrary even
number of slabs. For a Hilbert--Schmidt operator $A$ put
$\|A\|_{2k}=(\operatorname{Tr}(A^*A)^k)^{1/(2k)}$.
This is finite; cyclicity of the trace gives
$\|A^*\|_{2k}=\|A\|_{2k}$, and the spectral theorem for $A^*A$
shows that this number vanishes only when $A=0$.

\begin{lemma}[H\"older inequality for traces]\label{IV-3-35}
For $k\ge1$ and Hilbert--Schmidt operators $A_1,\ldots,A_{2k}$,
the product is trace class and
\begin{equation}\label{IV-3-36}
 |\operatorname{Tr}(A_1\cdots A_{2k})|
 \le\prod_{s=1}^{2k}\|A_s\|_{2k}.
\end{equation}
\end{lemma}
\begin{proof}
The trace-class property follows from the presence of two
Hilbert--Schmidt factors. Unless one factor is zero, normalize each
$A_s$ to have $\|A_s\|_{2k}=1$. Let $\mathcal A$ be the finite
set consisting of these normalized operators and their adjoints,
and maximize $|\operatorname{Tr}(B_1\cdots B_{2k})|$ over
$\mathcal A^{2k}$; denote the maximum by $m$.
An alternating tuple $(C,C^*,\ldots,C,C^*)$ has trace one,
so $m\ge1$.

For $X=B_1\cdots B_k$ and $Y=B_{k+1}\cdots B_{2k}$,
Hilbert--Schmidt Cauchy--Schwarz gives
\[
 |\operatorname{Tr}(XY)|^2
 \le\operatorname{Tr}(XX^*)\operatorname{Tr}(Y^*Y).
\]
Both traces on the right are nonnegative and correspond to
tuples in $\mathcal A^{2k}$. If the original tuple maximizes
the modulus, both reflected tuples also maximize it.
Among all maximizing tuples choose one with a longest string
of consecutive entries alternating between an operator and its
adjoint, counting positions cyclically. Let its length be $j$.
If $j\le k$, cyclically shift it to end at position $k$;
if $k<j<2k$, shift it to start at position one. Reflect the first
$k$ entries to obtain
$(B_1,\ldots,B_k,B_k^*,\ldots,B_1^*)$.
This is again a maximizer and has an alternating string of length
$2\min\{j,k\}>j$, a contradiction. Thus $j=2k$ and the
maximizing trace is one. Rescaling proves the inequality.
\end{proof}

Let $\mathbf F=(F_\kappa)_{\kappa\in\mathsf Z}$ be a family
with $F_\kappa$ a function of $\square_\kappa$, and abbreviate
\[
 Z[\mathbf F]=\mathcal Z_{a,T,\beta}\left(\prod_\kappa F_\kappa\right).
\]
For a direction $\mu$ and an index $t\bmod\mathsf n_\mu$,
write $\kappa|t$ for $\kappa$ with its $\mu$th component
replaced by $t$. Define the reflected family by
\begin{equation}\label{IV-4-3}
 (R_{\mu,t}\mathbf F)_\kappa=(F_{\kappa|t})^{[\kappa]}.
\end{equation}
The cells with $\kappa_\mu=t$ lie in a slab of $\lambda$
layers, starting just after layer
$\tau_0(t)=o_\mu/a+\lambda t-h_\lambda-1$.
Their product is a bounded Borel function $\Phi^{(t)}$ of the
inner slab bonds. Here $2\le\lambda\le\lambda\mathsf n_\mu-1$,
so these are slabs to which \cref{IV-3-26} applies.

\begin{proposition}\label{IV-4-5}
For every board, every $\beta\ge0$, and every such family,
\begin{align}
 Z[R_{\mu,t}\mathbf F]
 &=\operatorname{Tr}
   \bigl(\mathfrak Y[\Phi^{(t)}]^*
              \mathfrak Y[\Phi^{(t)}]\bigr)^{\mathsf n_\mu/2}
 =\|\mathfrak Y[\Phi^{(t)}]\|_{\mathsf n_\mu}^{\mathsf n_\mu}
 \ge0,\label{IV-4-6}\\
 |Z[\mathbf F]|&\le
 \prod_{t\bmod\mathsf n_\mu}
             Z[R_{\mu,t}\mathbf F]^{1/\mathsf n_\mu}.
 \label{IV-4-7}
\end{align}
\end{proposition}
\begin{proof}
The trace formula expresses $Z[\mathbf F]$ as the trace of
the product of the $\mathsf n_\mu$ slab operators, in reverse
cyclic order. In $R_{\mu,t}\mathbf F$, an even displacement
from slab $t$ is a translation of its slab function; an odd
displacement is a translation of the reflected, conjugated function
$(\Phi^{(t)})^\Theta$. Indeed, in the varying coordinate the
map between two slabs is respectively a translation by
$(t'-t)\mathsf l$ or reflection followed by translation, with
reflection formula $x_\mu\mapsto2o_\mu+(t+t')\mathsf l-x_\mu$.
The other coordinates are unchanged. Thus the operators in this
trace alternate between $Y=\mathfrak Y[\Phi^{(t)}]$ and $Y^*$,
by \eqref{IV-3-25}. Since $\mathsf n_\mu$ is even, their trace
is $\operatorname{Tr}(Y^*Y)^{\mathsf n_\mu/2}$, independently
of the initial parity. This proves \eqref{IV-4-6}.
Apply \cref{IV-3-35} to the original slab product and use
\eqref{IV-4-6} for each factor to obtain \eqref{IV-4-7}.
\end{proof}

Composition of copies in \eqref{IV-2-30} implies that
$R_{3,t_3}R_{2,t_2}R_{1,t_1}\mathbf F$ is the family of all
copies of $F_{(t_1,t_2,t_3)}$. In particular its product is
$\operatorname{Fold}F_{(t_1,t_2,t_3)}$.
For the family of all copies of one $F$, applying $R_{1,t}$
does not change the family. Equation \eqref{IV-4-6} therefore
shows that $\mathcal Z_{a,T,\beta}(\operatorname{Fold}F)\ge0$.
We may consequently define
\begin{equation}\label{IV-4-14}
 \|F\|_{\mathrm{ch}}
   =\langle\operatorname{Fold}F\rangle_{a,T,\beta}^{1/\mathsf N}.
\end{equation}

\begin{proposition}[Chessboard estimate]\label{IV-4-15}
For every board, every $\beta\ge0$, and every family of
bounded Borel functions $F_\kappa$ of their respective cells,
\begin{equation}\label{IV-4-16}
 \left|\left\langle\prod_{\kappa\in\mathsf Z}F_\kappa
                  \right\rangle_{a,T,\beta}\right|
 \le\prod_{\kappa\in\mathsf Z}\|F_\kappa\|_{\mathrm{ch}}.
\end{equation}
\end{proposition}
\begin{proof}
Apply \eqref{IV-4-7} in directions $1$, $2$, and $3$, in
that order. The first reflected integrals are nonnegative, so the
inequality can be applied successively under their positive roots.
The composition of copies just proved gives
\[
 |Z[\mathbf F]|
 \le\prod_{t_1,t_2,t_3}
 \mathcal Z_{a,T,\beta}
       (\operatorname{Fold}F_{(t_1,t_2,t_3)})^{1/\mathsf N}.
\]
Divide by $Z_{a,T,\beta}>0$, distributing its $\mathsf N$th
root among the $\mathsf N$ factors.
\end{proof}

\begin{proposition}\label{IV-4-17}
On functions of one cell, $\|\cdot\|_{\mathrm{ch}}$ is a
seminorm, is unchanged by conjugation or taking a mirror copy,
and satisfies
\[
 \|1\|_{\mathrm{ch}}=1,\qquad
 |\langle F\rangle_{a,T,\beta}|\le\|F\|_{\mathrm{ch}}
 \le\|F\|_\infty.
\]
For a fixed model function it is independent of the chosen cell
and of the offset of the board.
\end{proposition}
\begin{proof}
The homogeneity, conjugation, and supremum assertions follow
from \eqref{IV-2-30}--\eqref{IV-2-31} and the nonnegativity
of the folded integral. Set the other cell functions equal to
one in \cref{IV-4-15} to bound $|\langle F\rangle|$.
For the triangle inequality, expand
\[
 \operatorname{Fold}(F+H)
  =\sum_{J\subset\mathsf Z}
        \prod_{\kappa\in J}F^{[\kappa]}
        \prod_{\kappa\notin J}H^{[\kappa]}.
\]
The chessboard estimate and invariance under copies bound the
modulus of the mean of the $J$th term by
$\|F\|_{\mathrm{ch}}^{|J|}
 \|H\|_{\mathrm{ch}}^{\mathsf N-|J|}$.
Summing and taking the $\mathsf N$th root gives the triangle
inequality. Finally, changing $(o,\kappa)$ to $(o',\kappa')$
for the same model function translates its entire fold by
$o'-o+\mathsf l(\tilde\kappa'-\tilde\kappa)\in a\mathbb Z^3$.
The Wilson integral is translation invariant.
\end{proof}

\subsection{Sources and change of volume}

Let $w$ vanish outside the interior of $\square_\kappa$.
Define its copies by \eqref{IV-2-26}, acting now on functions
of position, and put
\begin{equation}\label{IV-2-43}
 w^{\mathrm{ch}}=\sum_{\kappa'}w^{[\kappa']}
       =w^++(\bar w)^-,
\end{equation}
where $w^+$ and $w^-$ are the sums of the un-conjugated
transports over respectively even and odd cells. The copies have
disjoint interiors and vanish on every cell face. In
$\mathcal O_a(w)$, only inner plaquettes of $\square_\kappa$
therefore occur, so $e^{\mathcal O_a(w)}$ is a function of that
cell. Covariance of the plaquette action and its real values give
\[
 (e^{\mathcal O_a(w)})^{[\kappa']}
   =e^{\mathcal O_a(w^{[\kappa']})},\qquad
 \operatorname{Fold}e^{\mathcal O_a(w)}
   =e^{\mathcal O_a(w^{\mathrm{ch}})}.
\]
In particular,
\begin{equation}\label{IV-4-19}
 Z_{a,T,\beta}[w^{\mathrm{ch}}]\ge0,
 \qquad
 \|e^{\mathcal O_a(w)}\|_{\mathrm{ch}}^{\mathsf N}
  =\frac{Z_{a,T,\beta}[w^{\mathrm{ch}}]}{Z_{a,T,\beta}}.
\end{equation}
Here $Z_{a,T,\beta}[w]=\mathcal Z_{a,T,\beta}(e^{\mathcal O_a(w)})$.
If $w=\sum_qz_qv_q$ for real sources $v_q$ of this cell, then
$w^{\mathrm{ch}}=\sum_q(z_qv_q^++\bar z_qv_q^-)$.
The dependence on $z$ and $\bar z$ is kept explicit.

\begin{proposition}[Comparison of volumes]\label{IV-4-20}
Fix $a>0$, an odd $\lambda\ge3$, and $\beta\ge0$.
Let $\mathsf n_\mu,\mathsf n'_\mu$ be even integers with
$2\le\mathsf n'_\mu\le\mathsf n_\mu$ in each direction.
For the same bounded Borel model function $F_0$, placed in any
cell and at any offset on the respective boards,
\begin{equation}\label{IV-4-21}
 0\le\mathcal Z_{\underline{\mathsf n}}
       (\operatorname{Fold}F_0)
 \le\mathcal Z_{\underline{\mathsf n}'}
       (\operatorname{Fold}F_0)^{
                 \prod_\mu\mathsf n_\mu/\mathsf n'_\mu}.
\end{equation}
The subscript on $\mathcal Z$ specifies the torus with sides
$a\lambda\mathsf n_\mu$; $a$, $\lambda$, and $\beta$ are unchanged.
\end{proposition}
\begin{proof}
By \cref{IV-4-17}, take offset and cell index zero. First
change only $\mathsf n_\mu$. The cross-section, the slab space
of $\lambda$ layers, and the operators $\mathsf K$ and
$\mathfrak Y$ are the same on the two tori, by \cref{IV-3-26}.
The product of the folded cell functions in slab zero is also
the same function $\Phi$ of those slab variables. To see the
last point, index its cells by the two transverse coordinates
$\varkappa$. Their reflected model functions have signs
$\epsilon_\mu=1$, $\epsilon_\nu=(-1)^{\varkappa_\nu}$
for $\nu\ne\mu$. A model coordinate $n_\mu$ is placed in
slab layer $n_\mu+h_\lambda+1$. These signs and coordinate
placements use only the transverse cell counts, $\lambda$,
and $F_0$, and hence do not involve the varying length.

Consequently \eqref{IV-4-6}, applied to the fold, gives the
two integrals as $\operatorname{Tr}B^k$ and
$\operatorname{Tr}B^{k'}$ for the same positive trace-class
operator $B=\mathfrak Y[\Phi]^*\mathfrak Y[\Phi]$, where
$k=\mathsf n_\mu/2\ge k'=\mathsf n'_\mu/2\ge1$.
If $b_i\ge0$ are the eigenvalues of $B$ and
$z=\sum_i b_i^{k'}>0$, then
$\rho_i=b_i z^{-1/k'}$ satisfies $\sum_i\rho_i^{k'}=1$.
In particular $0\le\rho_i\le1$, whence
$\sum_i\rho_i^k\le1$. Thus
\[
 \operatorname{Tr}B^k
 \le(\operatorname{Tr}B^{k'})^{k/k'}.
\]
If $z=0$, both sides vanish. This proves the comparison in
one direction. Change the three directions successively; their
exponents multiply to the exponent in \eqref{IV-4-21}.
\end{proof}

For cubic tori with spacing $a_K=L^{-K}$, write
\[
 p_\ell=\ell^{-3}\log Z_{K,\ell,\beta}.
\]
For a board of cell side $\mathsf l=a_i$, let
$\|F_0\|_{\mathrm{ch};\ell}$ be the seminorm of the model
function $F_0$ placed in any cell at any allowed offset.

\begin{corollary}[Two tori]\label{IV-4-23}
Let $\ell=2L^m$ and $\ell'=2L^{m'}$, where $m\ge m'$ are
integers, let $K+m'\ge1$, and let $\beta\ge0$. Then
\begin{equation}\label{IV-4-25}
 p_\ell\le p_{\ell'}\le0.
\end{equation}
If moreover $-m'\le i\le K-1$ and $\mathsf l=a_i$, then
every bounded Borel model function with $\lambda=L^{K-i}$
satisfies
\begin{equation}\label{IV-4-26}
 \|F_0\|_{\mathrm{ch};\ell}
 \le\|F_0\|_{\mathrm{ch};\ell'}
         e^{\mathsf l^3(p_{\ell'}-p_\ell)}.
\end{equation}
Let $\tilde w$ be a bounded complex function
on $\mathbb R^3$ vanishing outside the interior of one cell
of side $\mathsf l$ of the board with offset
$o\in a_K\mathbb Z^3$. Periodize it on the two tori to obtain
$w_\ell,w_{\ell'}$. Their mirror continuations lift to one
function on $\mathbb R^3$, and
\begin{equation}\label{IV-4-24}
 0\le Z_{K,\ell,\beta}[w_\ell^{\mathrm{ch}}]
 \le Z_{K,\ell',\beta}[w_{\ell'}^{\mathrm{ch}}]^{(\ell/\ell')^3}.
\end{equation}
\end{corollary}
\begin{proof}
For an allowed $i$, $\lambda=L^{K-i}\ge L\ge13$ is odd
and the cell counts $\ell/\mathsf l$, $\ell'/\mathsf l$
are even integers at least two. Thus \cref{IV-4-20} applies.
For the pressure assertion choose $i=-m'$, which is allowed
because $K+m'\ge1$. Apply the proposition to the constant
model function one and use $0<Z_{K,\ell,\beta}\le1$.
For general $F_0$, substitute
$\mathcal Z(\operatorname{Fold}F_0)
=Z_{K,\ell,\beta}\|F_0\|_{\mathrm{ch};\ell}^{(\ell/\mathsf l)^3}$
in \eqref{IV-4-21} and take roots to obtain \eqref{IV-4-26}.

For the source assertion, $2\mathsf l\le\ell'$, so the cell
interior embeds in both tori, and both periodizations give the
same model function for $e^{\mathcal O_a(w)}$.
Reflecting the lifted source through all cells of $\mathbb R^3$
gives a single function $\tilde w^{\mathrm{ch}}$. It has period
$2\mathsf l$ in each coordinate: shifting a cell index by an
even integer preserves its reflection sign. Thus it descends
to both tori, where it equals $w_\ell^{\mathrm{ch}}$ and
$w_{\ell'}^{\mathrm{ch}}$. Equations \eqref{IV-4-19} and
\eqref{IV-4-21} give \eqref{IV-4-24}.
\end{proof}

\section{Analytic bounds for cumulants}
\label{sec:analytic-cumulants}

The local vacuum energies may diverge as the cutoff is removed.
Their mixed derivatives at separated supports will be estimated
directly. The remaining part of a cumulant is a derivative of the
logarithm of a bounded holomorphic function. The following estimates
keep these two contributions separate and require a logarithm only
near zero.

For $c_i\ge0$ set $|z|_c=\sum_i c_i|z_i|$ and
$D_c=\{z\in\mathbb C^n:|z|_c<1\}$. We use the polydisc Cauchy bound
\begin{equation}\label{IV-5-6}
 |\partial_1\cdots\partial_n F(0)|
 \le B\prod_i r_i^{-1}
 \quad\text{if }|F|\le B\text{ on }\{|z_i|<r_i\}.
\end{equation}
Apply the one-variable Cauchy formula successively on circles of
radii smaller than $r_i$ and then increase the radii to $r_i$.
This proves the displayed bound, including for unequal radii.

\begin{lemma}\label{IV-5-7}
Let $\Gamma,\mathcal E$ be holomorphic on $D_c$, with
$\Gamma(0)=1$ and $\mathcal E(0)=0$. Suppose that
$h=e^{-\mathcal E}\Gamma$ satisfies $|h|\le Q$ there.
Then $Q\ge1$ and
\begin{equation}\label{IV-5-8}
 |h(z)-1|\le(Q+1)|z|_c.
\end{equation}
In particular $h$ has a holomorphic logarithm on
$|z|_c<(Q+1)^{-1}$, and
\begin{equation}\label{IV-5-10}
 |\partial_1\cdots\partial_n(\log\Gamma-\mathcal E)(0)|
 \le(\log2)[2(Q+1)n]^n\prod_i c_i.
\end{equation}
The logarithm of $\Gamma$ in this formula is its germ with value
zero at the origin. If
$\Gamma(z)=\mathbb E e^{\sum_i z_iX_i}$ for bounded random variables,
the mixed derivative of $\log\Gamma$ at zero is their joint cumulant.
\end{lemma}
\begin{proof}
First suppose that all $c_i$ are positive. For $z\ne0$ put
$s=|z|_c$. The function $\zeta\mapsto h(\zeta z/s)-1$ vanishes
at zero and is bounded by $Q+1$ on the unit disc. Schwarz's lemma
at $\zeta=s$ proves \eqref{IV-5-8}. On
$|z|_c\le[2(Q+1)]^{-1}$ it gives $|h-1|\le1/2$, whence the
power series for $\log(1+(h-1))$ gives $|\log h|\le\log2$.
The identity $\Gamma=e^{\mathcal E}h$ and the values at zero imply
$\log\Gamma=\mathcal E+\log h$ in a neighbourhood of zero.
The polydisc with radii
$r_i=[2(Q+1)nc_i]^{-1}$ lies in this smaller domain. Applying
\eqref{IV-5-6} proves \eqref{IV-5-10}.

For weights that may vanish, replace $c_i$ by
$c_i^\epsilon=\max(c_i,\epsilon)$. The corresponding domains are
contained in $D_c$, the bound $Q$ is unchanged, and the derivatives
at zero do not depend on $\epsilon$. Letting $\epsilon\downarrow0$
proves both assertions. In particular the mixed derivative vanishes
if any $c_i=0$. The assertion about cumulants follows from their
definition.
\end{proof}

For estimates in the number of insertions, it is useful to control
the generating function separately on every coordinate subspace.

\begin{lemma}[Bounds by the number of active variables]\label{IV-5-13}
Let $h$ be holomorphic on $D_c$, with $h(0)=1$. Suppose that for every
nonempty $B\subset\{1,\ldots,n\}$,
\[
 |h(z)|\le Q^{|B|}\quad\text{when }z_i=0\ (i\notin B),
 \qquad Q\ge1.
\]
Then
\begin{equation}\label{IV-5-15}
 |\partial_1\cdots\partial_n\log h(0)|
 \le(n-1)!(2^n-1)(eQ)^n\prod_i c_i.
\end{equation}
If $h=e^{-\mathcal E}\Gamma$ with the normalizations of
\cref{IV-5-7}, the derivative on the left is that of
$\log\Gamma-\mathcal E$.

For bounded random variables $|X_i|\le x_i$, one also has
\begin{equation}\label{IV-5-16}
 |\langle X_1;\ldots;X_n\rangle_\nu|
 \le(n-1)!(2^n-1)\prod_i x_i.
\end{equation}
If two probability laws $\nu,\nu'$ on the same space satisfy
\[
 |\nu(Y)-\nu'(Y)|\le\epsilon\|Y\|_\infty,
 \qquad Y=\prod_{i\in J}X_i,\quad\emptyset\ne J\subset\{1,\ldots,n\},
\]
then
\begin{equation}\label{IV-5-17}
 |\langle X_1;\ldots;X_n\rangle_\nu
       -\langle X_1;\ldots;X_n\rangle_{\nu'}|
 \le\epsilon\left(\sum_{\pi\in\mathcal P_n}|\pi|!\right)\prod_i x_i.
\end{equation}
\end{lemma}
\begin{proof}
For $b=|B|$ and positive $c_i$, apply \eqref{IV-5-6} on the
$B$-coordinate polydisc with radii $(bc_i)^{-1}$. This gives
\[
 |\partial_Bh(0)|\le Q^b b^b\prod_{i\in B}c_i
 \le(eQ)^b b!\prod_{i\in B}c_i.
\]
The last inequality follows from the $b$th term in the exponential
series for $e^b$. The partition formula for the logarithm, the same
coefficient calculation as in \eqref{IV-1-25}, reduces the desired
bound to
\begin{equation}\label{IV-5-14}
 \sum_{\pi\in\mathcal P_n}(|\pi|-1)!
       \prod_{B\in\pi}|B|!=(n-1)!(2^n-1).
\end{equation}
To verify this identity, first fix $k=|\pi|$. Ordering the $k$ blocks
and then their elements gives
\[
 k!\sum_{|\pi|=k}\prod_{B\in\pi}|B|!
   =n!\binom{n-1}{k-1}.
\]
Multiply this identity by $(k-1)!/k!=1/k$ and sum over $k$;
the summands on the right become
$(n-1)!\binom nk$. This proves \eqref{IV-5-14} and hence
\eqref{IV-5-15}. Vanishing weights follow by the approximation used
in \cref{IV-5-7}.

Every moment indexed by $B$ has modulus at most $\prod_{i\in B}x_i$.
The partition formula and \eqref{IV-5-14}, using $|B|!\ge1$, give
\eqref{IV-5-16}. To compare two laws, write the difference of the
products of moments for a partition with $k$ blocks as a telescoping
sum of $k$ terms. Each term is bounded by
$\epsilon\prod_i x_i$. Multiplying by the coefficient $(k-1)!$
in the partition formula and summing proves \eqref{IV-5-17}.
\end{proof}

\section{Separated derivatives of the source energies}
\label{sec:separated-source-derivatives}

The sum of the source energies need not be bounded uniformly in the
cutoff. Its mixed derivatives behave differently when two source
supports are separated. A contributing domain must meet both supports,
so at fine scales its size is large. The exponential decay in the domain
size then makes the sum over scales convergent. Combining this estimate
with the normalized partition function gives bounds for the truncated
correlations on a fixed physical torus.

Fix the constants and the source radius of \cref{III-111-27}, one
reference coupling $0<g\le\gamma_\varrho$, and a first-passage datum
$(K,\beta)$. The side $\ell$ satisfies \eqref{III-1-14}. We use the
assignment in the proof of \cref{III-111-7}, in particular
$\kappa_e=\kappa>\kappa_{\mathcal D}+2$, where
$\kappa_{\mathcal D}=\log9216$. Put $e_0(X;w)=0$. For a real Lipschitz
function on $T_\ell$ set
\begin{equation}\label{IV-1-16}
 \|f\|_\sharp=\max\{\|f\|_\infty,40M\Lip f\},
 \qquad
 \mathbb D_{\rho_*}(f_1,\ldots,f_n)
 =\left\{z:\sum_{q=1}^n|z_q|\|f_q\|_\sharp<\rho_*\right\},
 \qquad \rho_*=2\varrho.
\end{equation}
Here $\Lip$ uses Euclidean distance. Since
$|x-y|\le\sqrt3|x-y|_\infty$, the norm in \eqref{III-16-17}
satisfies $\|f\|_{\flat;T_\ell}\le\|f\|_\sharp$. Thus every source
$w(z)=\sum_qz_qf_q$ with $z\in\mathbb D_{\rho_*}$ belongs to the
source ball. The source energies are holomorphic there, vanish at zero,
depend only on the source in $X^{\sim1}$, and satisfy
\begin{equation}\label{eq:iv-source-energy-hypotheses}
 |e_j(X;w)|\le C_e e^{-\kappa_e d_j(X)},
 \qquad C_e=C_{\mathrm S}.
\end{equation}
These are precisely the properties of \cref{III-111-27} needed for
the derivative estimate.

\subsection{The size and number of contributing domains}

At depth $s=K-j$, the cubes of $\pi_j$ have physical side
$W_s=ML^{-s}$. For $Y\subset T_\ell$ let
\[
 N_s(Y)=\#\{\Delta\in\pi_{K-s}:\Delta\cap Y\ne\varnothing\},
 \qquad
 S(q)=\frac{2^{12}}{1-e^{\kappa_{\mathcal D}-q}}\quad
 (q>\kappa_{\mathcal D}).
\]

\begin{lemma}\label{IV-6-13}
If $X\in\mathcal D_{K-s}$ and $X^{\sim1}$ meets both $Y$ and
$Y'$, then
\begin{equation}\label{IV-6-14}
 d_{K-s}(X)\ge \frac{L^s\operatorname{dist}_\infty(Y,Y')}{M}-4.
\end{equation}
For $q>\kappa_{\mathcal D}$ and $D\in\mathbb R$,
\begin{equation}\label{IV-6-15}
 \sum_{\substack{X\in\mathcal D_{K-s}:\ X^{\sim1}\cap Y\ne\varnothing,\\
                         d_{K-s}(X)\ge D}}
 e^{-q d_{K-s}(X)}
 \le27N_s(Y)S(q)e^{-(q-\kappa_{\mathcal D})[D]_+}.
\end{equation}
Moreover,
\begin{equation}\label{IV-6-16}
 N_s(Y)\le(\ell L^s/M)^3,
 \qquad
 N_s(Y)\le(2RL^s/M+2)^3
 \quad\text{if }Y\subset\{x:|x-x_0|_\infty\le R\}.
\end{equation}
If $r$ is a lattice isometry induced by a signed coordinate permutation
and a translation, then $N_s(rY)\le27N_s(Y)$; if the second hypothesis
of \eqref{IV-6-16} holds, its second bound also holds for $rY$.
\end{lemma}

\begin{proof}
A connected set of $d_j(X)$ cube edges meeting all cubes of $X$
has diameter at most $d_j(X)W_s$. Every point of $X$ is within
$W_s$ of a vertex of this set. Consequently
$\operatorname{diam}_\infty X\le(d_j(X)+2)W_s$. Every point of
$X^{\sim1}$ is within $W_s$ of $X$: its cube touches a cube of
$X$, and a point of their intersection has this property. Two points
in $X^{\sim1}$ are therefore at distance at most
$(d_j(X)+4)W_s$, which proves \eqref{IV-6-14}.

Each cube meeting $Y$ touches at most $27$ cubes. If
$X^{\sim1}$ meets $Y$, the domain $X$ contains one of these at
most $27N_s(Y)$ cubes. For a fixed cube $\square$, the entropy
bound of Proposition~\ref{II-18-26} gives
\[
 \#\{X\supset\square:d_j(X)=n\}\le2048\,9216^n.
\]
Summing its geometric tail, with
$n\ge\lceil[D]_+\rceil$, gives
\[
 \sum_{\substack{X\supset\square\\d_j(X)\ge D}}
       e^{-q d_j(X)}
 \le S(q)e^{-(q-\kappa_{\mathcal D})[D]_+}.
\]
Summing over the anchor cubes proves \eqref{IV-6-15}.

There are $(\ell/W_s)^3$ cubes in the torus. To prove the local
count, lift the sup-ball into a box with side $2R$. A cube meeting
the projection of a box with side lengths $r_1,r_2,r_3$ has a
lift whose lower corner in each coordinate lies in an interval of
length $r_\mu+W_s$. This interval contains at most
$r_\mu/W_s+2$ points of the cube grid. Choose one such lift per
cube; the chosen corner determines the cube. Hence the count is at
most $\prod_\mu(r_\mu/W_s+2)$, giving \eqref{IV-6-16}.
A signed permutation and translation sends a cube to the projection
of a box with all side lengths $W_s$, so its image meets at most
$27$ partition cubes. Apply this to the cubes meeting $Y$.
The image of a sup-ball is a sup-ball of the same radius, which
proves the last assertion.
\end{proof}

We will use two elementary scale sums. For $\delta>0$ and $c\ge0$ put
\begin{equation}\label{IV-5-41}
 C'_\Sigma(\delta,c)=\frac{27}{26}(c+1)^3+
       \left(\frac6{e\delta}\right)^3
       \frac{e^{\delta c/2}}{1-e^{-\delta}},
 \qquad C''_\Sigma(\delta)=1+\frac1{1-e^{-2\delta}}.
\end{equation}
For $D>0$, $b\ge0$, and $L\ge3$,
\begin{align}
 \sum_{s\ge0}(DL^s)^3e^{-\delta[DL^s-c]_+}
     &\le C'_\Sigma(\delta,c),\label{IV-5-44}\\
 \sum_{s\ge0}(bL^s+2)^3e^{-\delta[DL^s-c]_+}
     &\le32\left(C''_\Sigma(\delta)+
                         \log_L^+\frac{c+1}{D}\right)
          +4(b/D)^3C'_\Sigma(\delta,c).
          \label{IV-5-46}
\end{align}
To verify these bounds, split at the first $s_0$ for which
$t_s=DL^s>c+1$. The preceding indices number at most
$1+\log_L^+((c+1)/D)$, and their cubic sum is at most
$\frac{27}{26}(c+1)^3$. For $s=s_0+i$,
$t_s>(c+1)L^i$ and $L^i\ge1+2i$. Thus
\[
 \sum_{s\ge s_0}e^{-\delta(t_s-c)}
     \le\frac1{1-e^{-2\delta}}.
\]
Using $t^3e^{-\delta t/2}\le(6/(e\delta))^3$ gives
\[
 \sum_{s\ge s_0}t_s^3e^{-\delta(t_s-c)}
 \le\left(\frac6{e\delta}\right)^3
             \frac{e^{\delta c/2}}{1-e^{-\delta}},
\]
which proves \eqref{IV-5-44}. Finally,
$(bL^s+2)^3\le4(b/D)^3t_s^3+32$ gives
\eqref{IV-5-46} from the preceding estimates.

\subsection{Mixed source derivatives}

For each $j$ choose any subset $\mathcal F_j\subset\mathcal D_j$.
Let $\mathcal E_{\mathcal F}(z)$ be the sum over these subsets of
terms of either of the forms
\[
 e_j(X;\textstyle\sum_qz_qf_q),\qquad
 \overline{e_j(X;\textstyle\sum_q\bar z_qf_q)}.
\]
The choice can vary with $(j,X)$. Both forms are holomorphic and
satisfy \eqref{eq:iv-source-energy-hypotheses}; this follows for
the second by conjugating the local power series. In the canonical
family the two forms agree by the conjugation law, but the estimates
below do not need that equality. Every such sum is finite,
holomorphic on $\mathbb D_{\rho_*}$, and zero at the origin.

\begin{proposition}\label{IV-6-18}
Let $n\ge2$ and suppose the real Lipschitz functions $f_q$ are
nonzero, with
$d_\infty=\operatorname{dist}_\infty(\operatorname{supp}f_1,
\operatorname{supp}f_2)>0$. Put
$\delta_e=\kappa_e-\kappa_{\mathcal D}>0$. Then
\begin{equation}\label{IV-6-19}
 \left|\partial_1\cdots\partial_n\mathcal E_{\mathcal F}(0)\right|
 \le\left(\frac n{\rho_*}\right)^n
       \prod_q\|f_q\|_\sharp\,
       27C_eS(\kappa_e)
       \sum_{s=0}^KN_s(\operatorname{supp}f_1)
             e^{-\delta_e[d_\infty L^s/M-4]_+}.
\end{equation}
The sum can be restricted to depths with
$\mathcal F_{K-s}\ne\varnothing$. It satisfies
\begin{equation}\label{IV-6-20}
 \sum_{s=0}^KN_s(\operatorname{supp}f_1)
             e^{-\delta_e[d_\infty L^s/M-4]_+}
 \le C'_\Sigma(\delta_e,4)(\ell/d_\infty)^3.
\end{equation}
If $\operatorname{supp}f_1$ lies in a sup-ball of radius $R$, it
also satisfies the volume-independent bound
\begin{equation}\label{IV-6-21}
 32\left(C''_\Sigma(\delta_e)+\log_L^+\frac{5M}{d_\infty}\right)
 +4\left(\frac{2R}{d_\infty}\right)^3C'_\Sigma(\delta_e,4).
\end{equation}
The bounds remain valid for a finite sum
$\sum_\iota c_\iota\mathcal E_{\mathcal F^\iota}$ formed at
the tuples $(r_\iota f_1,\ldots,r_\iota f_n)$, where
$c_\iota\ge0$, $\sum_\iota c_\iota\le1$, and each $r_\iota$
is a lattice isometry, with $N_s$ in \eqref{IV-6-19} replaced by
$\max_\iota N_s(r_\iota\operatorname{supp}f_1)$.
\end{proposition}

\begin{proof}
If $X^{\sim1}$ misses $\operatorname{supp}f_q$, locality makes
the corresponding energy independent of $z_q$ on
$\mathbb D_{\rho_*}$. Its full mixed derivative therefore
vanishes. In particular only domains seeing both first supports
contribute. By \cref{IV-6-13}, these satisfy
$d_j(X)\ge d_\infty L^{K-j}/M-4$.

The polydisc with radii
$r_q=\rho_*/(n\|f_q\|_\sharp)$ lies in
$\mathbb D_{\rho_*}$. Iterated one-variable Cauchy inequalities,
first at smaller radii and then in the limit, give for either term
\[
 |\partial_1\cdots\partial_n\widetilde e_j(X;0)|
 \le C_e e^{-\kappa_e d_j(X)}
          \left(\frac n{\rho_*}\right)^n\prod_q\|f_q\|_\sharp.
\]
Apply \eqref{IV-6-15} to the contributing domains at each depth,
then sum over the depths. This proves \eqref{IV-6-19}, including
its restricted-sum version. For \eqref{IV-6-20}, use the first
bound of \eqref{IV-6-16} and \eqref{IV-5-44} with
$D=d_\infty/M$ and $c=4$. For \eqref{IV-6-21}, use the second
bound of \eqref{IV-6-16} and \eqref{IV-5-46} with
$b=2R/M$. Lattice isometries preserve both source norms and the
distance between the first two supports. Apply the established
bound to each transformed tuple and take the weighted sum.
The last assertions of \cref{IV-6-13} give the stated bounds
for the maximum, so \eqref{IV-6-20} and \eqref{IV-6-21} remain
unchanged.
\end{proof}

\subsection{Truncated correlations on a fixed torus}

For the Wilson probability measure at $(K,\ell,\beta)$ define
\[
 \mathcal O_K(f)=L^K\sum_p f(x(p))a_p(U),
 \qquad
 \Gamma(z)=\left\langle e^{\sum_qz_q\mathcal O_K(f_q)}\right\rangle
         =\frac{Z_{K,\ell,\beta}[\sum_qz_qf_q]}{Z_{K,\ell,\beta}[0]}.
\]
The variables are bounded at finite cutoff, so $\Gamma$ is entire
and $\Gamma(0)=1$. Its logarithm is defined near zero. The joint
cumulant, or truncated correlation, is
\[
 S^T_{n;K,\ell,\beta}(f_1,\ldots,f_n)
       =\partial_1\cdots\partial_n\log\Gamma(0).
\]

\begin{proposition}\label{IV-6-23}
Set
\[
 Q_\ell=\exp\bigl(C_{\mathrm S}(2+\log g^{-2})\ell^3\bigr).
\]
For $n\ge2$ and real Lipschitz functions whose first two supports
have Euclidean distance at least $d>0$,
\begin{equation}\label{IV-6-24}
 \begin{split}
 |S^T_{n;K,\ell,\beta}(f_1,\ldots,f_n)|
 &\le\left(\frac n{2\varrho}\right)^n
 \left[\log2\,(2Q_\ell+2)^n
       +27C_{\mathrm S}S(\kappa_e)C'_\Sigma(\delta_e,4)
                       \left(\frac{\sqrt3\ell}{d}\right)^3\right]
       \prod_q\|f_q\|_\sharp .
 \end{split}
\end{equation}
All constants on the right are independent of $K$ and $\beta$.
\end{proposition}

\begin{proof}
If one function is zero, the mixed derivative defining the
cumulant is zero. Assume all are nonzero, and put
$\mathcal E(z)=\sum_{j,X}e_j(X;\sum_qz_qf_q)$ and
$c_q=\|f_q\|_\sharp/(2\varrho)$. The factorization and stability
bounds of \cref{III-111-27} give
\[
 h(z):=e^{-\mathcal E(z)}\Gamma(z)
       =\frac{\widetilde Z(\sum_qz_qf_q)}{\widetilde Z(0)},
 \qquad h(0)=1,\qquad |h(z)|\le Q_\ell
 \quad\left(\sum_qc_q|z_q|<1\right).
\]
Apply \cref{IV-5-7} with these weights and the bound $Q_\ell$.
It gives
\[
 |\partial_1\cdots\partial_n(\log\Gamma-\mathcal E)(0)|
 \le\log2\,(2Q_\ell+2)^n
              \left(\frac n{2\varrho}\right)^n
              \prod_q\|f_q\|_\sharp.
\]
The source-energy contribution is bounded by
\eqref{IV-6-19}--\eqref{IV-6-20}; the first two supports have
sup-distance at least $d/\sqrt3$. The triangle inequality now
gives \eqref{IV-6-24}.
\end{proof}

The estimate just proved is uniform in the ultraviolet cutoff on
a fixed torus. Its dependence on $\ell$ is still exponential.
To pass to growing tori or infinite-volume states, the normalized
partition function must instead be estimated on cells. The mixed
derivative bound \eqref{IV-6-21} is already independent of the
volume and will also apply to the cell normalisers.

\section{Cell normalization and volume-independent bounds}
\label{sec:cell-cumulants}

The bound on a fixed torus contains its volume because it uses the
stability estimate for the entire normalized partition function.
Reflection positivity permits a different normalization. We reflect a
source supported in one cell through a board, remove the source-energy
terms that see only one reflected copy, and estimate the remaining
terms by the distance between copies. The resulting bound depends on
the cell size. It can therefore be used on every larger torus.

We retain the constants and first-passage data of
\cref{sec:separated-source-derivatives}. Write
\[
 E_+=C_{\mathrm S},\qquad
 C_{\mathrm{low}}=C_{\mathrm S}(1+\log g^{-2}),\qquad
 C'_\Sigma=C'_\Sigma(\delta_e,4),\qquad
 C''_\Sigma=C''_\Sigma(\delta_e).
\]
The following argument uses the source-energy properties
\eqref{eq:iv-source-energy-hypotheses}, together with
\[
 |\widetilde Z_{\ell'}(w)|\le e^{E_+\ell'^3},\qquad
 Z_{K,\ell,\beta}\ge e^{E_\ell-C_{\mathrm{low}}\ell^3}.
\]
It thus applies with any constants satisfying these inequalities.
In this section $\mathcal O=\mathcal O_{a_K}$.
Here $E_\ell$ is the vacuum energy in \cref{III-111-27} on the
indicated torus. When comparing $\ell\ge\ell'$, put
\[
 \Delta_E=\frac{E_{\ell'}}{\ell'^3}-\frac{E_\ell}{\ell^3},
 \qquad C_{\mathrm{gap}}\ge\Delta_E.
\]
For the systems constructed in \cref{III-111-27} we may take
$C_{\mathrm{gap}}=C_{\mathrm S}$; when $\ell=\ell'$, we take
$C_{\mathrm{gap}}=0$. We do not identify the source energies on
different tori.

\subsection{The reflected source and its normalizer}

Fix $\ell=2L^m\ge\ell'=2L^{m''}$ with $m''>m'$, a cell side
$h=a_i=L^{-i}$, where $-m''\le i\le K-1$, and an offset
$o\in a_K\mathbb Z^3$. Both sides satisfy
\eqref{III-1-14}, so the source systems used above exist on them.
The cells in $\mathbb R^3$ are
\[
 Q_k=o+hk+[-h/2,h/2]^3,\qquad k\in\mathbb Z^3.
\]
Thus $h/a_K$ is an odd integer at least three and $\ell'/h$ is
an even integer. A closed carrier $A$ lies in one cell $Q_{k_0}$
at sup-distance at least $D>0$ from its boundary. Let
$\tilde v_1,\ldots,\tilde v_q$ be real Lipschitz functions vanishing
outside $A$. Write $v_i$ and $v'_i$ for their periodizations on
$T_\ell$ and $T_{\ell'}$.

These periodizations preserve the sharp norm in \eqref{IV-1-16},
where the same definition is used on $\mathbb R^3$. Indeed, the
cell has side at most half either torus side, so distances between
points in it agree with their torus distances. For a pair of points
not in the same copy of the support, a shortest joining segment meets
the boundaries of the copies containing its endpoints, where the
function vanishes. The triangle inequality, applied to the two end
segments, proves the Lipschitz bound in the other direction. The
same argument permits zero extension from a cell whenever the support
has positive distance from its boundary.

Reflect $A$ and each $\tilde v_i$ into every cell. On the reference
torus denote the resulting carriers by $A_c$ and the real functions
by $v_i^{[c]}$, where $c$ runs over the $N'=(\ell'/h)^3$ cells.
For $w=\sum_i z_iv_i$, its chessboard continuation has value
$\sum_i z_iv_i^{[c]}$ in an even cell and
$\sum_i\bar z_iv_i^{[c]}$ in an odd cell. Parity is the parity
of the number of coordinate reflections from $k_0$ to $c$.
The continuation is Lipschitz with the same supremum bound and
Lipschitz bound as $\sum_i z_i\tilde v_i$: on each cell the
reflection is an isometry, and a segment crossing cell boundaries
can be split at those boundaries, where the function is zero.
Consequently it belongs to the source ball when
$\sum_i|z_i|\|v_i\|_\sharp<\rho_*$.
Since the source vanishes on the cell boundary, every contributing
plaquette has its center in the cell interior and all its corners in
that cell. Thus $e^{\mathcal O(w)}$ is a function of the inner
cell bonds, as required for the chessboard estimate.

For a domain $X$ on $T_{\ell'}$, set
\[
 \mathcal C(X)=\{c:X^{\sim1}\cap A_c\ne\varnothing\}.
\]
If $\mathcal C(X)=\{c\}$, define the holomorphic function
\[
 \tilde e_j(X;z)=
 \begin{cases}
 e_j(X;\sum_i z_iv_i^{[c]}),&c\text{ even},\\
 \overline{e_j(X;\sum_i\bar z_iv_i^{[c]})},&c\text{ odd}.
 \end{cases}
\]
Our normalizer is
\begin{equation}\label{IV-7-16}
 \mathcal E^\circ(z)=\frac1{N'}
   \sum_{j=0}^K\sum_{X:\,|\mathcal C(X)|=1}\tilde e_j(X;z).
\end{equation}
It is holomorphic on $\mathbb D_{\rho_*}(v_1,\ldots,v_q)$ and
vanishes at zero. It is determined by the reference system, the
board, the carrier, and the function $\sum_i z_i\tilde v_i$.
In particular it is unchanged if that function is represented by
a different tuple on the same carrier. Neither the normalizer nor
its domain depends on the larger side $\ell$.

For a subset $B\subset\{1,\ldots,q\}$ put
\[
 S_c(B)=\bigcup_{i\in B}\operatorname{supp}v_i^{[c]},\qquad
 \bar N_s(B)=\max_c N_s(S_c(B)).
\]
Two distinct carriers $A_c,A_{c'}$ have sup-distance at least
$2D$, including across the periodic boundary. To see this, in a
coordinate in which their cell indices differ, their lifts lie in
intervals with centers differing by a nonzero multiple of $h$ and
half-length $h/2-D$. Every possible period shift retains a nonzero
cell-index difference, and the distance is at least
$h-2(h/2-D)=2D$. Distances between two reflected supports in a
single cell equal the original distances in $\mathbb R^3$.

The cube counts can also be read in $\mathbb R^3$. A cube of side
$W_s=ML^{-s}$ meeting one cell projects to a cube meeting its
image, and every such torus cube has a lift meeting the cell.
Two different meeting cubes cannot have the same image: their centers
differ in each coordinate by at most $h+W_s$, whereas equal images
would make these differences multiples of $\ell'$. Since
$h\le\ell'/2$ and $W_s\le M\le\ell'/(2L)$, they must all
vanish. Thus $\bar N_s(B)$ is the maximum of the corresponding
Euclidean cube counts over all reflected copies: shifting a copy by
a period of $T_{\ell'}$ preserves its reflection parity and the
cube grid, since $\ell'/h$ is even and $\ell'/W_s$ is an
integer. This maximum does not depend on $\ell'$ or $\ell$.

Define
\begin{equation}\label{IV-7-18}
 \Phi_B=C_eS(\kappa_e)\sum_{s=0}^K
  \min\left\{\left(\frac{hL^s}{M}\right)^3,
  27\bar N_s(B)e^{-\delta_e[2DL^s/M-4]_+}\right\}.
\end{equation}
The estimates below will be expressed using
\begin{equation}\label{IV-7-4}
 \mathsf U(R,d)=27C_eS(\kappa_e)
 \left[32\left(C''_\Sigma+
      \log_L^+\frac{5\sqrt3 M}{d}\right)
       +4C'_\Sigma\left(\frac{2\sqrt3 R}{d}\right)^3\right]
\end{equation}
and
\begin{equation}\label{IV-7-5}
 \Phi_1(R,D;h)=C_eS(\kappa_e)
 \left\{\frac{27}{26}+27\left[
 32\left(C''_\Sigma+\log_L^+\frac{5\min\{h,M\}}{2D}\right)
       +4C'_\Sigma\left(\frac RD\right)^3\right]\right\}.
\end{equation}
Both are nonnegative. The first increases with $R$ and decreases
with $d$; the second increases with $R,h$, decreases with $D$,
and is at most $\Phi_1(R,D;M)$.

\begin{proposition}[The cell bound]\label{IV-7-27}
For $z\in\mathbb D_{\rho_*}(v_1,\ldots,v_q)$ with $z_i=0$
outside $B$,
\begin{equation}\label{IV-7-29}
 \left\|e^{\mathcal O(w)}\right\|_{\mathrm{ch};\ell}
 \le \exp\left(\operatorname{Re}\mathcal E^\circ(z)
      +(E_++C_{\mathrm{low}}+\Delta_E)h^3+\Phi_B\right),
 \qquad E_++C_{\mathrm{low}}+\Delta_E\ge0.
\end{equation}
If each $\operatorname{supp}\tilde v_i$ is contained in a
sup-ball of radius $R$, then
$\Phi_B\le |B|\Phi_1(R,D;h)$. If one such ball contains the
whole carrier, the factor $|B|$ may be omitted. Under these radius hypotheses, suppose also that
$q\ge2$ and the first two supports have Euclidean distance at
least $d>0$. Then
\begin{equation}\label{IV-7-31}
 \begin{split}
 |\partial_1\cdots\partial_q\mathcal E^\circ(0)|
 &\le \left(\frac q{\rho_*}\right)^q\prod_i\|v_i\|_\sharp
    \,27C_eS(\kappa_e)\sum_{s=0}^K\bar N_s(\{1\})
       e^{-\delta_e[d_\infty L^s/M-4]_+}\\
 &\le \left(\frac q{\rho_*}\right)^q
       \mathsf U(R,d)\prod_i\|v_i\|_\sharp,
 \end{split}
\end{equation}
where $d_\infty$ is the sup-distance between the first two supports.
\end{proposition}

\begin{proof}
Let $N=(\ell/h)^3$. The chessboard identity and the comparison
of tori in \cref{IV-4-23} give
\[
 \left\|e^{\mathcal O(w)}\right\|_{\mathrm{ch};\ell}^{N}
 =\frac{Z_{K,\ell,\beta}[w^{\mathrm{ch}}]}{Z_{K,\ell,\beta}}
 \le \frac{Z_{K,\ell',\beta}[(w')^{\mathrm{ch}}]^{(\ell/\ell')^3}}
             {Z_{K,\ell,\beta}}.
\]
The numerator before taking the power is nonnegative. The reference
factorization and upper bound, followed by the floor on $T_\ell$,
therefore yield
\[
 \left\|e^{\mathcal O(w)}\right\|_{\mathrm{ch};\ell}
 \le \exp\left\{\frac1{N'}\operatorname{Re}
          \sum_{j,X}e_j(X;(w')^{\mathrm{ch}})
           +(E_++C_{\mathrm{low}}+\Delta_E)h^3\right\}.
\]
This argument does not require the sourced partition function to be
strictly positive. At zero source the comparison of pressures gives
$p_\ell\le p_{\ell'}$; the lower and upper vacuum bounds imply
$0\le p_{\ell'}-p_\ell\le E_++C_{\mathrm{low}}+\Delta_E$.

A domain seeing no carrier contributes zero by locality. A domain
seeing only $A_c$ has the same real part as $\tilde e_j(X;z)$,
also in an odd cell. It remains to bound, at depth $s=K-j$,
\[
 T_j=\frac1{N'}\sum_{X:\,|\mathcal C(X)|\ge2}
              |e_j(X;(w')^{\mathrm{ch}})|.
\]
Every domain contains a cube of $\pi_j$. Sum its bound first
over domains containing that cube, using the rooted count in
\cref{IV-6-13}, and then over all cubes. This bounds $T_j$ by
$C_eS(\kappa_e)(hL^s/M)^3$.
For a second bound, a nonzero summand must meet some $S_c(B)$.
It also meets a different carrier, whose distance from $A_c$ is
at least $2D$. Hence \eqref{IV-6-14} gives
$d_j(X)\ge 2DL^s/M-4$. The anchored tail estimate, first for
each $c$ and then summed over the $N'$ cells, gives
\[
 T_j\le27C_eS(\kappa_e)\bar N_s(B)
                    e^{-\delta_e[2DL^s/M-4]_+}.
\]
The factor $N'$ cancels the denominator defining $T_j$.
Taking the smaller bound at each scale proves \eqref{IV-7-29}.

We record explicitly the scale estimate needed here. For
$u,D_0>0$, $b,c\ge0$ and $\delta>0$,
\begin{equation}\label{IV-5-47}
 \begin{split}
 &\sum_{s\ge0}\min\{(uL^s)^3,
          c(bL^s+2)^3e^{-\delta[D_0L^s-4]_+}\}\\
 &\quad\le\frac{27}{26}
   +c\left[32\left(C''_\Sigma(\delta)+
       \log_L^+\frac{5\min\{u,1\}}{D_0}\right)
          +4\left(\frac b{D_0}\right)^3C'_\Sigma(\delta,4)\right].
 \end{split}
\end{equation}
For $uL^s<1$, sum the first term as a geometric progression;
its sum is at most $27/26$. For the other scales use
$(bL^s+2)^3\le4(b/D_0)^3(D_0L^s)^3+32$.
The cubic terms are controlled by \eqref{IV-5-44}. Among these
scales, those with $D_0L^s\le5$ lie in the interval
\[
 \max\{0,-\log_Lu\}\le s\le\log_L(5/D_0),
\]
so there are at most
$1+\log_L^+(5\min\{u,1\}/D_0)$ of them. The remaining
exponentials have sum at most $(1-e^{-2\delta})^{-1}$, as in
the proof of \eqref{IV-5-46}. This proves \eqref{IV-5-47}.

By the ball count in \cref{IV-6-13},
$\bar N_s(B)\le |B|(2RL^s/M+2)^3$. For nonempty $B$, apply
\eqref{IV-5-47} with $u=h/M$, $b=2R/M$, $D_0=2D/M$,
and $c=27|B|$. Its right side is at most
$|B|\Phi_1(R,D;h)/(C_eS(\kappa_e))$, with the assertion
interpreted directly from \eqref{IV-7-18} if $C_e=0$.
If $B$ is empty, $\Phi_B=0$. If the whole carrier is in one
ball, the same argument uses $c=27$.

If $v_i=0$ for some $i$, the normalizer is independent of $z_i$,
so both sides of \eqref{IV-7-31} vanish. Otherwise,
\eqref{IV-7-16} is an average of $N'$ sub-sums of
source energies at reflected tuples, with the holomorphic conjugate
used in odd cells. Apply \cref{IV-6-18} to these sub-sums.
The sharp norms and within-cell distances are preserved, and the
first support meets at most $(2RL^s/M+2)^3$ cubes at depth $s$.
Its sup-distance from the second support is at least $d/\sqrt3$.
Equation \eqref{IV-5-46}, with $b=2R/M$ and
$D_0=d/(\sqrt3 M)$, gives exactly \eqref{IV-7-31}.
\end{proof}

\subsection{Several occupied cells}

Suppose that each nonzero member of a real Lipschitz tuple
$\tilde f_1,\ldots,\tilde f_n$ on $\mathbb R^3$ is supported
in a sup-ball of radius $R$, in one cell of the board, at distance
at least $D>0$ from its boundary. Let $I_k$ be the indices in
cell $Q_k$, omitting zero functions. For each occupied cell take
the carrier $\bigcup_{i\in I_k}\operatorname{supp}\tilde f_i$
and construct \eqref{IV-7-16} on the same reference torus.
Set
\begin{equation}\label{IV-7-42}
 \mathcal E_{\underline f}(z)=\sum_{k\ \mathrm{occupied}}
                  \mathcal E_k^\circ((z_i)_{i\in I_k}).
\end{equation}
Let $f_i$ denote periodization on $T_\ell$ and put
$w_k=\sum_{i\in I_k}z_if_i$.

\begin{proposition}\label{IV-7-43}
The function \eqref{IV-7-42} is holomorphic on
$\mathbb D_{\rho_*}(f_1,\ldots,f_n)$ and vanishes at zero.
For $z_i=0$ outside $B\subset\{1,\ldots,n\}$,
\begin{equation}\label{IV-7-45}
 |e^{-\mathcal E_{\underline f}(z)}|
 \prod_{k\ \mathrm{occupied}}
       \left\|e^{\mathcal O(w_k)}\right\|_{\mathrm{ch};\ell}
 \le \exp\bigl((E_++C_{\mathrm{low}}+C_{\mathrm{gap}})h^3
                    +\Phi_1(R,D;h)\bigr)^{|B|}.
\end{equation}
If different occupied cells have different images on $T_\ell$,
the left side bounds
\begin{equation}\label{IV-7-46}
 \left|e^{-\mathcal E_{\underline f}(z)}
       \left\langle e^{\sum_i z_i\mathcal O(f_i)}\right\rangle
          _{K,\ell,\beta}\right|.
\end{equation}
The same normalizer works for all larger tori satisfying the stated
vacuum floor and gap bound. For $n\ge2$, its full mixed derivative
vanishes if two cells are occupied or one $f_i$ is zero. If all
functions occupy one cell and the first two supports have Euclidean
distance at least $d>0$, then
\begin{equation}\label{IV-7-47}
 |\partial_1\cdots\partial_n\mathcal E_{\underline f}(0)|
 \le(n/\rho_*)^n\mathsf U(R,d)\prod_i\|f_i\|_\sharp.
\end{equation}
\end{proposition}

\begin{proof}
Apply \cref{IV-7-27} in each occupied cell, with the active
indices $B\cap I_k$. A cell with no active index contributes
one. Every other cell contributes at most
\[
 \exp\bigl((E_++C_{\mathrm{low}}+C_{\mathrm{gap}})h^3
                       +|B\cap I_k|\Phi_1(R,D;h)\bigr).
\]
There are at most $|B|$ such cells, and
$E_++C_{\mathrm{low}}+C_{\mathrm{gap}}\ge0$ by
\eqref{IV-7-29}; multiplying proves \eqref{IV-7-45}.
For distinct torus cells, the chessboard inequality
\cref{IV-4-15} bounds the expectation of the product by the
product of its seminorms. Linearity of $\mathcal O$ identifies
that product with the exponential in \eqref{IV-7-46}.
The normalizers use only the fixed reference system and the
Euclidean carriers, so the claim about larger sides follows.
If two cells are occupied, every summand of
\eqref{IV-7-42} omits a variable. The same is true if some
$f_i=0$. Otherwise \eqref{IV-7-47} is \eqref{IV-7-31}.
\end{proof}

\subsection{Placement and the bound for separated supports}

\begin{lemma}[Placement by a lattice translation]\label{IV-8-17}
Let $h=a_i$ with $-m\le i\le K-1$, and suppose that the
support of each $f_q\in\operatorname{Lip}(T_\ell)$ lies in
a sup-ball of radius $R$. If $D>0$ and
\begin{equation}\label{IV-8-18}
 n(2R+2D+a_K)\le h,
\end{equation}
one lattice translation places the whole tuple in cells with margin
$D$. The translated functions have compactly supported lifts
$\tilde g_q$ in representative cells with pairwise distinct torus
images, one representative for each occupied cell. Periodization
preserves their sharp norms. The translation preserves all truncated
correlations. When two functions lie in one cell, the distance
between the supports of their lifts equals the torus distance.
\end{lemma}

\begin{proof}
Put $a=a_K$ and $\lambda=h/a$. Choose lifts of the centers
$x_1,\ldots,x_n$. For a fixed coordinate consider the $\lambda$
possible offsets $0,a,\ldots,(\lambda-1)a$. An offset is bad
for $x_q$ if that coordinate has distance less than $R+D$ from
one of the corresponding mirrors. Representatives of the bad
offsets inject into points of $a\mathbb Z$ in an open interval
of length $2(R+D)$. There are fewer than $2(R+D)/a+1$ of
them. By \eqref{IV-8-18}, fewer than $\lambda$ offsets are
bad for some $q$. Choose a good offset in each coordinate.
This gives a board offset $o'\in a\mathbb Z^3$ for which
all the balls have margin $D$. Translation by $o-o'$ places
them on the prescribed board.

Lift the translated balls to cells with indices in
$\{0,\ldots,\ell/h-1\}^3$, and extend each function by zero
outside its cell. The positive margin makes these extensions
Lipschitz, with the norm preservation proved above. Different
representative indices have different torus images. Since
$h\le\ell/2$, distances within a representative cell agree with
torus distances. Finally translation by a vector of $a_K\mathbb Z^3$
permutes bonds and plaquettes and preserves Haar measure and the
Wilson action. It therefore preserves the generating functional
and all its cumulants.
\end{proof}

For $n\ge2$ and $R>0$ set
\begin{align}
 h(n,R)&=\min\{L^b:b\in\mathbb N, L^b\ge8n(R+1)\},
                      \label{IV-8-7}\\
 Q(n,R)&=\exp\bigl((E_++C_{\mathrm{low}})h(n,R)^3
                              +\Phi_1(R,R+1;M)\bigr),
                      \label{IV-8-8}\\
 C_n^\infty(d;R)&=(n/\rho_*)^n
       \left[\log2\,(2Q(n,R)^n+2)^n+\mathsf U(R,d)\right].
                      \label{IV-8-9}
\end{align}
The minimum exists and $8n(R+1)\le h(n,R)<8Ln(R+1)$.

\begin{proposition}\label{IV-8-21}
Let $n\ge2$, $R,d>0$, and $\ell\ge16n(R+1)$. Suppose each
real Lipschitz $f_q$ on $T_\ell$ has support in a sup-ball of
radius $R$, and the first two supports have Euclidean distance
at least $d$. Then
\begin{equation}\label{IV-8-23}
 |S^T_{n;K,\ell,\beta}(f_1,\ldots,f_n)|
     \le C_n^\infty(d;R)\prod_q\|f_q\|_\sharp.
\end{equation}
The constant is independent of $K$, $\beta$, and $\ell$.
\end{proposition}

\begin{proof}
Zero functions give zero cumulant. Otherwise use the board of side
$h=h(n,R)$. Since $\ell/2=L^m\ge8n(R+1)$, the minimal
power $h$ is at most $\ell/2$, so it is an allowed cell side.
Take $D=R+1$. As $a_K<1$,
$n(2R+2D+a_K)<8n(R+1)\le h$.
Apply \cref{IV-8-17}, and use the torus itself as the reference
torus in \cref{IV-7-43}, so $C_{\mathrm{gap}}=0$.
For the translated tuple let $\Gamma$ be its generating function
and $\mathcal E$ its cell normalizer. Then
\[
 h_0=e^{-\mathcal E}\Gamma,\qquad h_0(0)=1,
 \qquad |h_0|\le Q(n,R)^n
       \quad\text{on }\mathbb D_{\rho_*}.
\]
Here \eqref{IV-7-45} and
$\Phi_1(R,D;h)\le\Phi_1(R,D;M)$ give the upper bound.
Lemma~\ref{IV-5-7} bounds the full derivative of
$\log\Gamma-\mathcal E$ by the first term of
\eqref{IV-8-9}, times the product of sharp norms.
If the tuple occupies two cells, the full derivative of
$\mathcal E$ is zero. Otherwise its bound is
\eqref{IV-7-47}; the required separation is preserved by the
placement. Adding the bounds and undoing the translation proves
\eqref{IV-8-23}.
\end{proof}

\subsection{Factorial growth in the number of insertions}

A finer use of the coordinate-slice bound avoids the rapid growth
of $Q(n,R)$ in \eqref{IV-8-9}. Define
\begin{align}
 p&=\frac{864C_eS(\kappa_e)}{\log L},\label{IV-8-4}\\
 \varkappa_\star&=5+p,\label{IV-8-10}\\
 C_Q^0&=\exp\left[E_++C_{\mathrm{low}}+C_eS(\kappa_e)
 \left\{\frac{27}{26}+27[32(C''_\Sigma+\log_L5)+4C'_\Sigma]
       \right\}\right],\label{IV-8-5}\\
 A(d)&=1+\mathsf U(1/10,d),\label{IV-8-11}\\
 B(R)&=\max\left\{
 \frac{27(4R+1)^3\,417M\,e^{5+p}}{\rho_*}
                         (2e\,3^pC_Q^0+1),
              \ 288e^4(48R+1)^3\right\}.\label{IV-8-13}
\end{align}
In particular $p\ge0$, $C_Q^0\ge1$, and $B(R)\ge1$.
For later use put
\begin{equation}\label{IV-8-12}
 A_{\mathrm{PL}}=1+27C_eS(\kappa_e)
 \left[32\left(C''_\Sigma+\log_L(5\sqrt3 M)+\frac1{\log L}\right)
                  +4C'_\Sigma(\sqrt3/5)^3\right].
\end{equation}
Then
\begin{equation}\label{IV-8-29}
 A(d)\le A_{\mathrm{PL}}\max\{1,d^{-4}\}.
\end{equation}
Indeed, for $t=\max\{1,d^{-1}\}$,
$\log_L^+(5\sqrt3 M/d)\le\log_L(5\sqrt3 M)+t/\log L$
and $d^{-3}\le t^3\le t^4$.

We first specialize the cell bound to a unit board. If a tuple is
placed with radius and margin $\eta\le1/4$, then
\[
 e^{E_++C_{\mathrm{low}}+\Phi_1(\eta,\eta;1)}
       =C_Q^0(2\eta)^{-p}=:Q_\eta.
\]
This follows from $M\ge1$ and
$\log_L(5/(2\eta))=\log_L5+\log(1/(2\eta))/\log L$.
Thus \cref{IV-7-43}, with the same torus as reference, gives
$|e^{-\mathcal E}\Gamma|\le Q_\eta^{|B|}$ on every coordinate
slice indexed by $B$. Applying the graded analytic estimate
\cref{IV-5-13} and the bound for $\partial_1\cdots\partial_n
\mathcal E(0)$ yields
\begin{equation}\label{eq:graded-cell-cumulant}
 |S^T_n(v_1,\ldots,v_n)|
 \le\left[(n-1)!2^n\left(\frac{eQ_\eta}{\rho_*}\right)^n
                  +\left(\frac n{\rho_*}\right)^n\mathsf U(\eta,d)
       \right]\prod_q\|v_q\|_\sharp,
\end{equation}
provided different occupied cells have different torus images and
one pair of supports has distance at least $d$.

\begin{proposition}\label{IV-8-27}
Let $n\ge2$ and $R,d>0$. On any allowed torus, let
$f_1,\ldots,f_n$ be real Lipschitz functions, each supported in
a sup-ball of radius $R$, with one pair of supports at Euclidean
distance at least $d$. Then
\begin{equation}\label{IV-8-28}
 |S^T_{n;K,\ell,\beta}(f_1,\ldots,f_n)|
       \le A(d)B(R)^n(n!)^{\varkappa_\star}
                           \prod_q\|f_q\|_\sharp.
\end{equation}
All constants are independent of $K$, $\beta$, and $\ell$.
\end{proposition}

\begin{proof}
By symmetry of cumulants, label a separated pair by $1,2$.
Assume all functions are nonzero. We treat coarse and fine
lattices separately.

If $a_K>1/(24n)$, the number of positive plaquette barycenters
in a sup-ball of radius $R$ is at most $3(2R/a_K+1)^3$.
For each of the three orientations, choose a lift of every
barycenter in the coordinate box about a lift of the ball center;
these choices inject into the corresponding shifted lattice.
The interval count in each coordinate gives the assertion, even
if the box is larger than the torus. Since $0\le a_p\le2$,
\[
 |\mathcal O(f)|\le6a_K^{-1}(2R/a_K+1)^3\|f\|_\infty
       \le144n^4(48R+1)^3\|f\|_\sharp.
\]
The bounded-variable part of \cref{IV-5-13} now gives
\[
 |S^T_n(f_1,\ldots,f_n)|
 \le(n-1)!\,[288(48R+1)^3]^n n^{4n}
                        \prod_q\|f_q\|_\sharp.
\]
Using $n^{4n}\le e^{4n}(n!)^4$, this is bounded by the
second term in \eqref{IV-8-13} to the power $n$, times
$(n!)^5\prod_q\|f_q\|_\sharp$.

Suppose now that $a_K\le1/(24n)$, and put $\varepsilon=1/(5n)$.
The ratio $\ell/\varepsilon=5n\ell$ is an integer. On the
periodic grid $\varepsilon\mathbb Z^3/\ell\mathbb Z^3$, define
\[
 \chi_y(x)=\prod_{\mu=1}^3
          \max\{0,1-|x_\mu-y_\mu|_\ell/\varepsilon\}.
\]
In each coordinate the adjacent linear tents sum to one, so
$\sum_y\chi_y=1$. Their products have values in $[0,1]$,
support in the sup-ball of radius $\varepsilon$ about $y$,
and Lipschitz constant at most $\sqrt3/\varepsilon$:
telescoping the product bounds its change by
$\varepsilon^{-1}\sum_\mu|x_\mu-x'_\mu|_\ell$.
Consequently
\begin{equation}\label{IV-5-34}
 \|\chi_yf\|_\sharp
       \le(1+40\sqrt3 M/\varepsilon)\|f\|_\sharp.
\end{equation}
The support of a piece is contained in the support of $f$.
If $f$ is supported in a sup-ball of radius $R$, a nonzero
piece has its node within $R+\varepsilon$ of the center.
Choosing lifts and counting each coordinate shows that there
are at most $(2R/\varepsilon+3)^3$ such pieces.

Expand every $f_q=\sum_y\chi_yf_q$ and use multilinearity of
the cumulant. For a nonzero tuple of pieces $v_1,\ldots,v_n$,
the first two supports remain separated by $d$. Moreover
\[
 n(4\varepsilon+a_K)\le\frac45+\frac1{24}<1.
\]
Thus \cref{IV-8-17} places the tuple on the unit board with
radius and margin $\varepsilon$, with distinct representatives
for distinct occupied cells. Apply \eqref{eq:graded-cell-cumulant} to that
tuple. Since $1/(6n)<\varepsilon\le1/10$,
\[
 \begin{split}
 (2R/\varepsilon+3)^3&\le[3n(4R+1)]^3,\\
 1+40\sqrt3 M/\varepsilon&\le417Mn,\\
 Q_\varepsilon&\le C_Q^0\,3^pn^p,\qquad
 \mathsf U(\varepsilon,d)\le\mathsf U(1/10,d).
 \end{split}
\]
The second inequality follows from $M\ge1$ and
$240\sqrt3<416$; the bounds intentionally use the weaker
$\varepsilon>1/(6n)$. Put $c_R=27(4R+1)^3\,417M$.
Summing the estimates for all tuples of pieces yields
\[
 \begin{split}
 |S^T_n(f_1,\ldots,f_n)|
 \le c_R^n n^{4n}\left[
 n!\left(\frac{2e\,3^pC_Q^0}{\rho_*}\right)^n n^{pn}
       +\left(\frac n{\rho_*}\right)^n\mathsf U(1/10,d)
                    \right]\prod_q\|f_q\|_\sharp.
 \end{split}
\]
For $t\ge0$, the elementary inequality
$n!\ge(n/e)^n$ gives $n^{tn}\le e^{tn}(n!)^t$.
Use it with $t=4+p$ in the first term and $t=5$ in the
second. Both terms are bounded by the first expression in
\eqref{IV-8-13} to the power $n$, times
$(n!)^{5+p}\prod_q\|f_q\|_\sharp$, with the respective
remaining factors $1$ and $\mathsf U(1/10,d)$.
Their sum is \eqref{IV-8-28}.
\end{proof}

\section{Compactness on separated test functions}
\label{sec:separated-compactness}

Uniform estimates at separated supports give one subsequence for all
orders and all test tuples. The support restriction matters in this
argument: approximating the entries must preserve their separation.
We give the compactness statement on both a fixed torus and
$\mathbb R^3$, since the latter will also be used after the volume
limit.

Let $\mathbb E=T_\ell$ or $\mathbb R^3$. For an open set
$O\subset\mathbb E$, let $\mathfrak S_n(O)$ be the tuples
$(f_1,\ldots,f_n)$ of real functions in $C_c^1(O)$ with pairwise
disjoint supports. Zero entries are allowed. Put
$\mathfrak S_0(O)=\{\varnothing\}$ and
$\mathfrak S(O)=\bigcup_{n\ge0}\mathfrak S_n(O)$.
A tuple has type $(d;R)$ if distinct supports are at distance at
least $d>0$ and each support lies in a ball of radius $R>0$;
the centres of these balls may differ. Every separated tuple has
such a type, since its supports are compact.

A map on $\mathfrak S_n(O)$ is called multilinear if its restriction
to $\prod_i C_c^1(O_i)$ is multilinear whenever the open sets
$O_i\subset O$ have pairwise disjoint closures. We use
$\|f\|_\sharp=\max(\|f\|_\infty,40M\Lip f)$ also on
$\mathbb R^3$. With
$\|f\|_{C^1}=\max(\|f\|_\infty,\|\nabla f\|_\infty)$,
integration along a shortest line segment, lifted on the torus,
gives $\|f\|_\sharp\le40M\|f\|_{C^1}$.

Fix $n_0\ge1$ and finite constants $C_n(d;R)\ge0$ for
$n\ge n_0$ and $d,R>0$. A family $\Phi=(\Phi_n)_{n\ge n_0}$
obeys these bounds if
\begin{equation}\label{IV-10-9}
 |\Phi_n(f_1,\ldots,f_n)|
 \le C_n(d;R)\prod_i\|f_i\|_\sharp
\end{equation}
on every tuple of type $(d;R)$. We impose no bound as $d$ tends
to zero.

\begin{proposition}[Simultaneous subsequence extraction]\label{IV-10-25}
Let $O_j\subset\mathbb E$ be open sets such that every compact
subset of $\mathbb E$ is contained in $O_j$ for all sufficiently
large $j$. Suppose that the real multilinear families
$\Phi^{(j)}$ on $\mathfrak S(O_j)$ satisfy
\eqref{IV-10-9} for $n\ge n_0$, with the same constants,
once $j\ge j_0(n,R)$. Then there is a subsequence along which
$\Phi^{(j)}_n(X)$ converges for every $n\ge n_0$ and every
$X\in\mathfrak S_n(\mathbb E)$. The subsequence can also be
chosen so that any prescribed countable collection of bounded
real sequences converges along it. The limiting family is
multilinear and satisfies the same bounds. Symmetry is preserved.
\end{proposition}

\begin{proof}
\emph{A countable family with controlled supports.}
Let $\mathfrak O$ consist of finite unions of rational open boxes,
using their projections on the torus. These sets have compact
closures and form a countable basis under finite unions.
For each $W\in\mathfrak O$ choose a countable $C^1$-dense subset
$\mathfrak f(W)\subset C_c^1(W)$.

Here is the separability argument with the support fixed.
For $f\in C_c^1(W)$ consider
$g_f=(f,\partial_1f,\partial_2f,\partial_3f)$ on $\overline W$.
The $C^1$ distance of two functions is the uniform distance of
these maps for the norm
$|(v_0,v')|_* =\max(|v_0|,|v'|)$.
For each integer $p\ge1$ choose a finite $1/p$-net in
$\overline W$. Classify the functions by a common modulus bound
$|g_f(x)-g_f(y)|_*\le1/q$ for $|x-y|\le1/p$, and by rational
approximations within $1/q$ to their values at the points of the
net. There are countably many classes. Each has diameter at most
$4/q$, by comparison through a nearest net point. Uniform continuity
of $g_f$ on $\overline W$ places every $f$ in such a class for each
$q$, with $p$ sufficiently large. Choosing one member of each
nonempty class therefore gives the required dense subset. For
$W=\varnothing$ take $\mathfrak f(W)=\{0\}$.

Take all tuples of all lengths with entries in
$\mathfrak f(W_i)$, where the $W_i$ have pairwise disjoint closures.
This is a countable family; enumerate it by $(X^k)_{k\ge1}$.
Given a tuple $X=(f_1,\ldots,f_n)$ of type $(d;R)$, compactness
of each support permits a choice of such sets with
\begin{equation}\label{IV-10-16}
 \operatorname{supp}f_i\subset W_i
 \subset\{x:\operatorname{dist}(x,\operatorname{supp}f_i)<d/3\}.
\end{equation}
Their mutual distances are at least $d/3$, and each lies in a ball
of radius $R+d/3$. The union of their closures is compact.
For a zero entry choose $W_i=\varnothing$.

\emph{Uniform approximation of evaluations.}
If $f_i,h_i$ have supports in $W_i$, multilinearity gives
\[
 \Phi_n(f_1,\ldots,f_n)-\Phi_n(h_1,\ldots,h_n)
 =\sum_{i=1}^n\Phi_n(h_1,\ldots,h_{i-1},f_i-h_i,
                                  f_{i+1},\ldots,f_n).
\]
Every tuple on the right retains separation $d/3$.
Choose $h_i\in\mathfrak f(W_i)$ with
$\|f_i-h_i\|_{C^1}\le\epsilon\le1$. For every family obeying
\eqref{IV-10-9} at $(n,R+d/3)$, the preceding identity gives
\begin{equation}\label{IV-10-18}
 |\Phi_n(X)-\Phi_n(h_1,\ldots,h_n)|
 \le C_n(d/3;R+d/3)\,n(40M)^n\epsilon
            (1+\max_i\|f_i\|_{C^1})^{n-1}.
\end{equation}
The constant does not depend on the approximating family or its
domain, provided the latter contains all $W_i$.

\emph{Extraction and extension.}
For each $X^k$ of length at least $n_0$, its evaluation by $\Phi^{(j)}$ is defined and
bounded for all sufficiently large $j$. Assign arbitrary values
to the finitely many undefined entries. Apply the diagonal
subsequence argument to these countably many bounded real sequences
and to the additional sequences in the statement. Each selected
evaluation then converges. For a fixed arbitrary $X$, use
\eqref{IV-10-18} to approximate it by one $X^k$, with an error
uniform in all sufficiently large $j$. The triangle inequality
shows that the evaluations at $X$ are Cauchy. This proves convergence
on every separated tuple along the same subsequence.

Each multilinearity or symmetry identity involves only finitely
many evaluations, all eventually defined, so it passes to the limit.
The same is true of \eqref{IV-10-9} for every fixed tuple and type.
\end{proof}

\begin{proposition}\label{IV-10-28}
The set of symmetric multilinear families on $\mathfrak S(\mathbb E)$
satisfying \eqref{IV-10-9} is compact and metrizable for pointwise
convergence on separated tuples.
\end{proposition}
\begin{proof}
Use the determining family $X^k$ in the preceding proof and let $n_k$
be the length of $X^k$. A metric is
\begin{equation}\label{IV-10-27}
 d(\Phi,\Psi)=\sum_{k:n_k\ge n_0}2^{-k}
        \min\{1,|\Phi_{n_k}(X^k)-\Psi_{n_k}(X^k)|\}.
\end{equation}
The triangle inequality follows term by term. If this expression
vanishes, \eqref{IV-10-18} makes every evaluation of the two families
equal, so it is a metric. That same estimate shows that every
evaluation is continuous for the metric: approximate its tuple by
one $X^k$, uniformly over all the families. Conversely, the tail
of the series is uniformly bounded by $\sum_{k>N}2^{-k}$, and its
first $N$ terms are continuous for pointwise convergence. The two
topologies thus agree. By \cref{IV-10-25}, every sequence has a
convergent subsequence whose limit still belongs to the set. Sequential
compactness in this metric proves compactness.
\end{proof}

\subsection{The fixed-torus application}

\begin{proposition}\label{IV-11-30}
Fix the source assignment, $0<g\le\gamma_\varrho$, and a side
$\ell$ satisfying \eqref{III-1-14}. Every sequence of first-passage
data $(K_j,\beta_j)$ with $K_j\to\infty$ has a subsequence along
which
\[
 S_n(X)=\lim_j S^T_{n;K_j,\ell,\beta_j}(X)
\]
exists simultaneously for all $n\ge2$ and all
$X\in\mathfrak S_n(T_\ell)$. The same subsequence may be chosen
so that $L^{-K_j}\beta_j$ converges. The limit is real,
symmetric and multilinear, and inherits \eqref{IV-6-24}.
Such sequences, and hence such limits, exist.
\end{proposition}
\begin{proof}
The restrictions of the lattice cumulants to separated $C^1$ tuples
are real, symmetric and multilinear. The coefficient multiplying
$\prod_i\|f_i\|_\sharp$ in \eqref{IV-6-24} is finite and independent
of the cutoff and bare coupling. It supplies the constants in
\eqref{IV-10-9}, independent of $R$, on the constant domains
$O_j=T_\ell$. Moreover \cref{III-5-11} bounds the physical
labels $L^{-K_j}\beta_j$ in a fixed compact interval. Apply
\cref{IV-10-25} with this one additional bounded sequence.
Every fixed-tuple inequality passes to the limit, giving the stated
bound. Finally \cref{III-5-19} supplies the nonempty interval
$I(g)=(g^{-2}+\bar B,X_g-\bar B]$ and shows that
$(K,L^K\mathsf b)$ is a first-passage datum for every $K\ge1$
and every $\mathsf b\in I(g)$. This proves existence.
\end{proof}

For a family $S=(S_n)_{n\ge2}$ define its centred moments by
\begin{equation}\label{IV-9-14}
 \sigma_S(f_1,\ldots,f_n)
 =\sum_{\substack{\pi\in\mathcal P_n\\|B|\ge2\ (B\in\pi)}}
       \prod_{B\in\pi}S_{|B|}((f_i)_{i\in B}),
 \qquad \sigma_S(\varnothing)=1.
\end{equation}
There is no term for $n=1$, so the one-point function vanishes.
The moment--cumulant identities show that the connected functions
of $\sigma_S$ are precisely $S_n$, $n\ge2$. At finite cutoff this
is the expectation of the product of
$\mathcal O_a(f_i)-\langle\mathcal O_a(f_i)\rangle$.
Since each partition sum is finite, convergence of all the cumulants
on separated tuples also gives convergence of these centred moments.

\section{Infinite-volume lattice states}\label{sec:infinite-volume-states}

We first take the volume limit at a fixed spacing. The configuration
space is compact, but the finite periodic measures do not themselves
satisfy the infinite-lattice Gibbs equations. We therefore check those
equations on each fixed set of bonds before passing to a limit.

Let $a>0$, let $\mathbb L_a=a\mathbb Z^3$, and let $\mathsf{Bd}$
be its positively oriented nearest-neighbour bonds. A configuration
is an element of
\[
 \mathcal V(\mathbb L_a)=\mathrm{SU}(2)^{\mathsf{Bd}},
 \qquad U(\bar b)=U(b)^{-1}.
\]
We use the product topology. This is a compact metrizable space.
For $B\subset\mathsf{Bd}$, write $\mathcal F(B)$ for the bounded
Borel functions depending only on $U_B$, and $\mathcal C_{\rm loc}$
for the continuous functions depending on finitely many bonds.
For a spatial set $Q$, $\mathsf{Bd}(Q)$ consists of bonds whose
closed segments lie in $Q$. The notation
$\mathcal C_{\rm loc}(Q)$ requires a finite bond support in
$\mathsf{Bd}(Q)$.

Enumerating the bonds gives a product metric with summable coordinate
weights. Replacing all but the first $N$ coordinates by the identity
approximates the identity map uniformly in this metric. Uniform
continuity therefore shows that $\mathcal C_{\rm loc}$ is uniformly
dense in $C(\mathcal V(\mathbb L_a))$. In particular, its integrals
determine a Borel probability measure. We call such a measure a law.
The laws form a compact metrizable space in the weak topology.

For a plaquette $p$, retain the notation
$a_p(U)=1-\frac12\operatorname{tr}\partial U(p)\in[0,2]$.
For bounded compactly supported $f$ on $\mathbb R^3$, put
\begin{equation}\label{IV-13-6}
 \mathcal O_a(f)=a^{-1}\sum_{p\in\mathsf{Pl}}f(x(p))a_p.
\end{equation}
The sum is finite, so $\mathcal O_a(f)\in\mathcal C_{\rm loc}$
even when $f$ is merely bounded. For a law $\nu$ define
\begin{equation}\label{IV-13-23}
 S_n^T[\nu](f_1,\ldots,f_n)
 =\langle\mathcal O_a(f_1);\ldots;\mathcal O_a(f_n)\rangle_\nu.
\end{equation}
If $X=(f_1,\ldots,f_n)$ and $\rho,\nu$ are two laws, write
\begin{equation}\label{IV-13-25}
 X_a^\nu=\prod_{q=1}^n
       \bigl(\mathcal O_a(f_q)-\nu(\mathcal O_a(f_q))\bigr),
 \qquad \sigma_{\rho,\nu}(X)=\rho(X_a^\nu),
 \qquad \sigma_\nu=\sigma_{\nu,\nu}.
\end{equation}
The empty product is one. The first index specifies the integrating
law, the second the centring law. These expressions are symmetric
and multilinear in the entries. Lattice isometries
$\mathsf r\in a\mathbb Z^3\rtimes W_3$ act by pullback on
configurations and functions, with inversion when a bond orientation
is reversed. Plaquette traces are unchanged by cyclic permutation
and inversion, hence
\[
 \mathsf r\mathcal O_a(f)=\mathcal O_a(\mathsf rf),\qquad
 \sigma_{\mathsf r_*\rho,\mathsf r_*\nu}(\mathsf rX)
 =\sigma_{\rho,\nu}(X).
\]
The same covariance holds for the cumulants.

\subsection{Gibbs equations and periodic limits}

A period lattice is a finite-index subgroup $\Gamma\subset a\mathbb Z^3$
with $\operatorname{sys}(\Gamma):=\min_{0\ne v\in\Gamma}|v|_\infty
\ge3a$. Choose one representative of every positive bond class
modulo $\Gamma$. Independent Haar variables at these representatives,
extended periodically, define $dU_\Gamma$ on the closed space
$\mathcal V_\Gamma$ of periodic configurations. Changing representatives
only permutes the variables. If $D$ contains one site of each class,
then the plaquettes based at sites of $D$ represent the plaquette
classes once. Thus
\begin{equation}\label{IV-13-30}
 S_\Gamma=\sum_{p\in\mathsf{Pl}/\Gamma}a_p,
 \quad Z_\Gamma[\Phi]=\int e^{-\beta S_\Gamma+\Phi}\,dU_\Gamma,
 \quad d\mu_\Gamma^\Phi=Z_\Gamma[\Phi]^{-1}
                        e^{-\beta S_\Gamma+\Phi}\,dU_\Gamma
\end{equation}
are independent of representatives; here $\beta\ge0$ and
$\Phi\in C(\mathcal V_\Gamma)$ is real. Write $\mu_{\Gamma,\beta}$
when $\Phi=0$. The action and Haar measure show that
$\mathsf r_*\mu_\Gamma^\Phi=\mu_\Gamma^{\mathsf r\Phi}$ whenever
$\mathsf r^\circ\Gamma=\Gamma$.

For finite $B\subset\mathsf{Bd}$, let $\mathsf{Pl}[B]$ be the
plaquettes containing a bond of $B$, and set
\begin{equation}\label{IV-14-6}
 H_B=-\beta\sum_{p\in\mathsf{Pl}[B]}a_p,
 \qquad
 (\gamma_BF)(U)=
 \frac{\int F(V\vee U_{B^c})e^{H_B(V\vee U_{B^c})}\,dV_B}
      {\int e^{H_B(V\vee U_{B^c})}\,dV_B}.
\end{equation}
Every bond is in four plaquettes, so $\#\mathsf{Pl}[B]\le4\#B$.
The denominator lies in $[e^{-2\beta\#\mathsf{Pl}[B]},1]$.
The kernel is positive, preserves one and contracts the supremum
norm. Integration over the finite set $B$ preserves continuity;
if $F$ is local, the remaining support lies in the union of its
support and the bonds of $\mathsf{Pl}[B]$, with $B$ removed.
Thus $\gamma_B\mathcal C_{\rm loc}\subset\mathcal C_{\rm loc}$.

A Gibbs state is a law satisfying
\begin{equation}\label{IV-14-9}
 \nu(\gamma_BF)=\nu(F)
 \quad(B\text{ finite},\ F\in\mathcal C_{\rm loc}).
\end{equation}
This equality then holds for every bounded Borel $F$: the law
$E\mapsto\nu(\gamma_B1_E)$ has the same integrals as $\nu$ on
$\mathcal C_{\rm loc}$, so the two laws agree.
Write $\mathcal G_{a,\beta}$ for the Gibbs states and
$\mathcal G^{\rm per}_{a,\beta}$ for those invariant under a
finite-index subgroup of $a\mathbb Z^3$. Equivalently, a periodic
state has a period number $P\in\mathbb N$, meaning invariance under
$Pa\mathbb Z^3$; a finite-index subgroup contains such a subgroup,
as the orders of the three coordinate classes in the finite quotient
have a common multiple.

We also need locally finite perturbations. Let
$\mathfrak A=(A_i)_{i\in I}$ be real local continuous functions,
each with a nonempty finite support $B_i$, such that each bond
lies in only finitely many $B_i$, counting indices. Replace $H_B$
in \eqref{IV-14-6} by
\[
 H_B^{\mathfrak A}=H_B+\sum_{i:B_i\cap B\ne\emptyset}A_i
\]
to define $\gamma_B^{\mathfrak A}$ and the corresponding Gibbs
equations.

\begin{proposition}\label{IV-14-14}
The Gibbs states form a convex weakly compact set. The subset invariant
under any specified subgroup of lattice translations is also convex
and weakly compact. Lattice isometries carry Gibbs states to Gibbs
states and preserve period numbers. If $t_j\to t$ and $\nu_j\to\nu$
weakly, with $\nu_j$ Gibbs for $t_j\mathfrak A$, then $\nu$ is Gibbs
for $t\mathfrak A$.
\end{proposition}
\begin{proof}
Each Gibbs equation is a closed affine condition on laws, since
$\gamma_BF-F$ is continuous. Invariance is likewise expressed by
$\nu(F\circ\mathsf r-F)=0$. Intersect these conditions in the compact
space of laws. Relabelling bonds, and inverting the reversed ones,
preserves product Haar measure and gives
$\gamma_{\mathsf rB}(\mathsf rF)=\mathsf r(\gamma_BF)$.

For the last assertion put
$A_B=\sum_{i:B_i\cap B\ne\emptyset}A_i$ and
$\delta_j=|t_j-t|\|A_B\|_\infty$.
The ratio of the two normalized kernel densities lies between
$e^{-2\delta_j}$ and $e^{2\delta_j}$. Consequently
\[
 \|\gamma_B^{t_j\mathfrak A}F-
       \gamma_B^{t\mathfrak A}F\|_\infty
 \le(e^{2\delta_j}-1)\|F\|_\infty\longrightarrow0.
\]
The fixed kernel applied to $F\in\mathcal C_{\rm loc}$ is continuous.
Pass to the limit in the equations for $\nu_j$.
\end{proof}

\begin{proposition}[Periodic approximations]\label{IV-14-16}
Let $\Gamma_j$ be period lattices with
$\operatorname{sys}(\Gamma_j)\to\infty$. Suppose that
$\Gamma_j$ acts on the indices of $\mathfrak A$, with
$A_{v\cdot i}=\tau_vA_i$ and $B_{v\cdot i}=B_i+v$.
Put $\Phi_j=\sum_{[i]\in I/\Gamma_j}A_i$ on $\mathcal V_{\Gamma_j}$.
Every weak limit of $\mu_{\Gamma_j}^{\Phi_j}$ is Gibbs for
$\mathfrak A$. In particular, limits of $\mu_{\Gamma_j,\beta}$
are translation-invariant Gibbs states.

At $a=a_K$, let $\widetilde\mu_{K,\ell,\beta}$ be the periodic
lift of the cubic torus measure and define
\[
 \mathcal G^{\rm tor}_{K,\beta}
 =\bigcap_{m_0\ge1}
   \overline{\{\widetilde\mu_{K,2L^m,\beta}:m\ge m_0\}}.
\]
This is a nonempty compact subset of $\mathcal G^{\rm per}_{K,\beta}$.
Every member is invariant under $a_K\mathbb Z^3\rtimes W_3$ and
has period number one.
\end{proposition}
\begin{proof}
The index action is free: $B_i+v=B_i$ with $B_i$ finite nonempty
forces $v=0$. Every orbit has an index whose support meets a fixed
bond transversal, so local finiteness makes $I/\Gamma_j$ finite.
Periodicity makes $\Phi_j$ independent of the representatives.

Fix a finite $B$ and a local continuous $F$. Choose a closed cube
$Q$ containing its bond support, all plaquettes of $\mathsf{Pl}[B]$,
and all $B_i$ meeting $B$. For large $j$, its side is smaller than
$\operatorname{sys}(\Gamma_j)$. Distinct bonds and plaquettes in
$Q$ are then incongruent modulo $\Gamma_j$. Choose a bond
transversal containing $\mathsf{Bd}(Q)$ and write its variables as
$(V_B,W)$. In these variables,
\[
 -\beta S_{\Gamma_j}+\Phi_j
       =H_B^{\mathfrak A}+H',
\]
where $H'$ is independent of $V_B$. Indeed, a remaining plaquette
class, or perturbation orbit, depending on a bond of $B$ would
have a representative already included in $H_B^{\mathfrak A}$.
For functions supported in $Q$, replacing $V_B$ in a periodic
configuration agrees on $Q$ with the nonperiodic replacement in
the definition of $\gamma_B^{\mathfrak A}$. Integrating first over
$V_B$ therefore gives
\[
 \mu_{\Gamma_j}^{\Phi_j}(\gamma_B^{\mathfrak A}F)
       =\mu_{\Gamma_j}^{\Phi_j}(F).
\]
Both functions are continuous and independent of $j$. The equality
passes to every weak limit.

For the empty perturbation family, periodic Haar measure and the
action are invariant under every lattice translation. Invariance
passes to weak limits. For the cubic sequence, the same is true
of every signed coordinate permutation. Its nested tail closures
are nonempty compact sets in the space of laws, so their intersection
is nonempty and compact. Metrizability identifies this intersection
with the subsequential limits. These limits are Gibbs by the first
part, proving all assertions.
\end{proof}

\begin{lemma}[Gauge invariance]\label{IV-14-13}
Every Gibbs state is invariant under every gauge transformation.
Two Gibbs states agreeing on gauge-invariant local continuous
functions are equal.
\end{lemma}
\begin{proof}
First let a gauge transformation differ from the identity only at
one site $x$. In the kernel for the six incident bonds, it is a
left or right translation of each Haar variable and preserves
every plaquette trace. Thus $\gamma_{B_x}(F\circ u)=\gamma_{B_x}F$,
and the Gibbs equation gives invariance. Compose these changes
for finitely many sites. On a local function, an arbitrary gauge
transformation uses only its values at the endpoints of a finite
bond support; hence invariance holds there and therefore for the law.

Averaging a local continuous $F$ over the gauge variables at those
endpoints gives a gauge-invariant local continuous $F^{\rm inv}$
with the same integral in every Gibbs state. Agreement on such
averages implies agreement on $\mathcal C_{\rm loc}$, hence equality
of the laws.
\end{proof}

\subsection{Reflection positivity for adapted periods}

The diagonal reflections require period lattices that split in the
normal direction. Let
\[
 \mathfrak N=\{e_\mu:1\le\mu\le3\}
   \cup\{e_\mu\pm e_\nu:\mu<\nu\},
 \qquad \Sigma(n,c)=\{x:n\cdot x=c\}.
\]
For coordinate normals take $c\in\frac a2\mathbb Z$, and for
diagonal normals take $c\in a\mathbb Z$. These are exactly the
positions in these families for which
$\theta_{n,c}x=x-2(n\cdot x-c)n/|n|^2$ preserves the lattice.
The coordinate mirrors are site mirrors when $c\in a\mathbb Z$
and bond mirrors otherwise.

A period lattice is adapted to $n$ if, for some $h>0$,
\begin{equation}\label{IV-15-23}
 \Gamma=\Gamma_\perp\oplus\mathbb Zv_\Gamma,
 \qquad \Gamma_\perp=\Gamma\cap n^\perp,
 \qquad v_\Gamma=h n/|n|^2.
\end{equation}
Then $n\cdot\Gamma=h\mathbb Z$, and the linear reflection preserves
$\Gamma$. If $n$ is coordinate, $h\in a\mathbb N$ and $h\ge3a$;
if $n$ is diagonal, $h\in2a\mathbb N$ and $h/2\ge3a$.
Thus $c+h/2$ is also admissible. The reflections in $c$ and $c+h/2$
differ by translation through $v_\Gamma$ and act identically on
$\mathcal V_\Gamma$.

\begin{proposition}[Two-mirror positivity]\label{IV-16-26}
Let $a>0$, $\beta\ge0$, $n\in\mathfrak N$, and let $\Gamma$
be a period lattice adapted to $n$, with height period $h$.
For admissible $c$, put $\theta=\theta_{n,c}$ and
$\mathbb S^+=\{c\le n\cdot x\le c+h/2\}$. Then
\begin{equation}\label{IV-16-27}
 \mu_{\Gamma,\beta}(\overline{\theta F}F)\ge0
 \qquad\bigl(F\in\mathcal F(\mathsf{Bd}(\mathbb S^+))\bigr).
\end{equation}
The law is $\theta$-invariant. Both assertions hold also for
$\mu_\Gamma^\Phi$ if
$\Phi=(\Phi_++\theta\Phi_+)|_{\mathcal V_\Gamma}$ with
$\Phi_+$ real continuous and dependent only on the bonds of
$\mathbb S^+$. No gauge-invariance assumption on $F$ is imposed;
bonds lying in either mirror are included.
\end{proposition}
\begin{proof}
We retain the boundary variables instead of integrating them into
either half. Write $t(x)=n\cdot x-c$. Modulo $\Gamma_\perp$, choose
representatives $R^0$ of the bonds lying in $t=0$ or $t=h/2$,
$R^+$ of the other bonds contained in $0\le t\le h/2$, and
$R^\times$ of the bonds crossing either boundary. The remaining
representatives are the reflected bonds $(\theta b)^+$,
$b\in R^+$. For a bond or plaquette, use its midpoint or
barycentre to measure height. Translation by $jv_\Gamma$ changes
this height by $jh$. Every object has a representative in the
two half-slabs, their boundaries, or the crossing sets, with height
in $(-h/2,h/2]$. Two such representatives congruent modulo $\Gamma$
differ by an element of $\Gamma_\perp$, since their height difference
has modulus less than $h$. This proves that the chosen representatives
are disjoint and exhaustive. There are finitely many of them.

Consequently the coordinates
\[
 U_0=U_{R^0},\quad U_+=U_{R^+},\quad
 U'_+=(\theta U)_{R^+},\quad U_\times=U_{R^\times}
\]
are independent Haar variables. The necessary inversions in
$U'_+$ preserve Haar measure. On the boundary bonds $\theta$ fixes
the coordinates modulo $\Gamma$, and it interchanges $U_+$ and
$U'_+$. Thus $F(U)=\widehat F(U_0,U_+)$ and
$(\theta F)(U)=\widehat F(U_0,U'_+)$.

Choose plaquette representatives in the same way. Those lying in
the boundaries contribute a weight $W_0(U_0)>0$. Those in the
positive half but not in its boundaries contribute $W_+(U_0,U_+)>0$,
and their reflected copies contribute $W_+(U_0,U'_+)$.
The remaining plaquettes meet both sides of a mirror. Their form
follows by listing the corner heights
\[
 t(x),\quad t(x)+an_\mu,\quad
 t(x)+a(n_\mu+n_\nu),\quad t(x)+an_\nu.
\]
A site mirror has none. At a diagonal mirror the only possibility
has two opposite corners $z,z'$ in the plane and two reflected
corners off it, at heights $\pm a$ relative to the plane.
With $y$ the corner in $\mathbb S^+$, set
$v_p(U)=U(\langle z,y\rangle)U(\langle y,z'\rangle)$.
These two bonds belong to the positive half and not to its boundary;
the inequality $a<h/2$ ensures that they do not reach the other mirror.
Tracing the plaquette boundary gives
\[
 e^{-\beta a_p(U)}=k_\beta(v_p(U_+)v_p(U'_+)^{-1}),
 \qquad k_\beta(g)=e^{-\beta}e^{(\beta/2)\operatorname{tr}g}.
\]

At a bond mirror, let $x,y$ be neighbouring sites in the adjacent
layer of $\mathbb S^+$. Put
$u_b=U(\langle x,y\rangle)$,
$u'_b=(\theta U)(\langle x,y\rangle)$ and
$w_x=U(\langle x,\theta^{(k)}x\rangle)$, where $\theta^{(k)}$
is reflection in the boundary under consideration. The four-bond
loop gives
\[
 e^{-\beta a_p(U)}
   =k_\beta(u_bw_y(u'_b)^{-1}w_x^{-1}).
\]
The crossing variables at the two mirrors are disjoint: their
midpoint heights differ by $h/2$, not by a multiple of $h$.
Their Haar integral at each boundary is therefore
\[
 K(u,u')=\int\prod_{b=(x,y)}
             k_\beta\bigl(u_b(w\cdot u')_b^{-1}\bigr)\,dw,
 \qquad (w\cdot u')_b=w_xu'_bw_y^{-1}.
\]

Each resulting kernel has a uniformly absolutely convergent expansion
$\sum_\alpha c_\alpha\Phi_\alpha(V)
\overline{\Phi_\alpha(V')}$ with
$c_\alpha\ge0$ and
$\sum_\alpha c_\alpha\|\Phi_\alpha\|_\infty^2<\infty$.
Indeed, the matrix-entry expansion used in \cref{IV-3-16} gives
\[
 k_\beta(gg'^{-1})
 =e^{-\beta}\sum_{N\ge0}\frac{(\beta/2)^N}{N!}
   \sum_{I,J\in\{1,2\}^N}
        g_{IJ}\overline{g'_{IJ}},
 \qquad \sum_{N,I,J}e^{-\beta}\frac{(\beta/2)^N}{N!}=e^\beta.
\]
Products and pullbacks preserve this property. For a bond boundary,
simultaneous gauge invariance and Haar invariance replace the single
average in $K$ by independent averages on its two entries. The
expansion then has features
$\Phi_\alpha^{\rm av}(u)=\int\Phi_\alpha(w\cdot u)\,dw$;
their supremum norms do not increase. This proves the assertion
also for the product $\mathcal K$ of the two boundary kernels.

After integrating $U_\times$, the full Wilson weight is
\[
 W_0(U_0)W_+(U_0,U_+)W_+(U_0,U'_+)\mathcal K(U_+,U'_+).
\]
Termwise integration is justified by the summable supremum bounds
and gives the unnormalized quadratic form
\[
 \int dU_0\,W_0(U_0)
   \sum_\alpha c_\alpha
    \left|\int \widehat F(U_0,U_+)W_+(U_0,U_+)
                    \Phi_\alpha(U_+)\,dU_+\right|^2\ge0.
\]
Division by $Z_\Gamma[0]>0$ proves \eqref{IV-16-27}.
For the perturbed measure apply this calculation to $Fe^{\Phi_+}$.
Finally $\theta$ preserves Haar measure and the action, and the
stated perturbation is $\theta$-invariant, proving invariance.
\end{proof}

\subsection{The torus-limit states}

Periodization $f^{\langle\ell\rangle}(x)=
\sum_{v\in\ell\mathbb Z^3}f(\widetilde x+v)$ is well defined
for compactly supported $f$. Grouping the plaquettes by period class
gives the exact identity
\begin{equation}\label{IV-17-14}
 \mathcal O_K(f)(\operatorname{per}_\ell U)
       =\mathcal O_{a_K}(f^{\langle\ell\rangle})(U).
\end{equation}
If the supports lie in $(-\ell/4,\ell/4)^3$, periodization preserves
their stated separation and per-entry radius bounds, and their
$\|\cdot\|_\sharp$ norms. To check the distance assertion, each
coordinate difference of two points in this cube has modulus
less than $\ell/2$, so it is already the shortest representative
of its class modulo $\ell$. The same argument and local lifting
give the supremum and Lipschitz norm identities. Periodization
commutes with the lattice isometries whenever the source and its
image are considered on the same quotient.

\begin{proposition}\label{IV-17-45}
Let $K\ge1$, $\beta\ge0$ and $\nu\in\mathcal G^{\rm tor}_{K,\beta}$.
There are sides $\ell_j\to\infty$ with
$\widetilde\mu_{K,\ell_j,\beta}\to\nu$ weakly.
Along every such sequence, for bounded compactly supported test
functions and $w=\sum_qz_qf_q$,
\begin{gather}
 \nu(e^{\mathcal O_K(w)})
   =\lim_j\frac{Z_{K,\ell_j,\beta}[w^{\langle\ell_j\rangle}]}
                         {Z_{K,\ell_j,\beta}[0]},\label{IV-17-46}\\
 S_n^T[\nu](X)
   =\lim_j S^T_{n;K,\ell_j,\beta}(X^{\langle\ell_j\rangle}),
 \qquad
 \sigma_\nu(X)
   =\lim_j\sigma_{K,\ell_j,\beta}(X^{\langle\ell_j\rangle}).
 \label{IV-17-47}
\end{gather}
Any bound on the right uniform for tuples in
$(-\ell_j/4,\ell_j/4)^3$ with prescribed separation and per-entry
radius passes to the left for tuples anywhere in $\mathbb R^3$.
The state $\nu$ is reflection positive at every coordinate site
and bond mirror, for bounded local functions of the closed positive
half-space.
\end{proposition}
\begin{proof}
The sequence exists by the tail-closure description in
\cref{IV-14-16}. Every product of the observables and their
exponentials is local continuous. Weak convergence and
\eqref{IV-17-14} therefore give convergence of all moments and
the first formula. Cumulants and centred moments are fixed finite
polynomials in these moments, proving \eqref{IV-17-47}.
A fixed compact union of supports eventually lies in the central
cube, so each uniform bound passes to the limit.

The cubic period lattice is adapted to every coordinate normal,
with height period $\ell_j$. A fixed finite bond support in
$\{x_\mu\ge c\}$ eventually lies in
$\{c\le x_\mu\le c+\ell_j/2\}$. Apply \cref{IV-16-26}
and pass to the weak limit for continuous local $F$.
To extend the resulting inequality to bounded Borel $F$ with the
same finite support, approximate its function of the finite bond
coordinates in $L^2$ of the corresponding marginal of $\nu$ by
continuous functions. Reflection invariance gives the same
$L^2$ convergence for their reflected functions. Cauchy--Schwarz
then passes the quadratic form to the limit. Reversing the side
follows by reflection invariance.
\end{proof}

The last assertion does not supply diagonal positivity. For
$n=e_\mu\pm e_\nu$, the heights of $\ell\mathbb Z^3$ are
$\ell\mathbb Z$, whereas its shortest nonzero pure-normal vector
is $\ell n$, of height $2\ell$. Thus the cubic period lattice
does not satisfy \eqref{IV-15-23} for a diagonal normal.


\subsection{Periodic arrays and their reflection}

The comparison of Gibbs states below uses two finite-volume
inequalities. The first compares a periodic array with its expectation
in a Gibbs state; the second increases its partition function by
reflecting the array across successive slabs.

For a site transversal $D$ of $\Gamma$, write
$\mathsf{Bd}_D$ for the positive bonds starting in $D$, and let
\[
 \partial_rD=\{x\in D:x+w\notin D\text{ for some }w\in a\mathbb Z^3,
                                      |w|_\infty\le ra\}.
\]
If $\Gamma\subset\Gamma'\subset a\mathbb Z^3$ have finite index,
the array of a local function $A$ is
$\mathbb A=\sum_{v\in\Gamma'/\Gamma}\tau_vA$ on
$\mathcal V_\Gamma$, with $|\mathbb A|=[\Gamma':\Gamma]$.

\begin{lemma}[Lower tangent]\label{IV-14-20}
Let $\Gamma$ be a period lattice, $\Gamma'\supset\Gamma$ a
finite-index subgroup, $D$ a site transversal of $\Gamma$, and
$\nu\in\mathcal G_{a,\beta}$ invariant under $\Gamma'$.
Let $A\in\mathcal C_{\rm loc}$ be real and $t\in\mathbb R$.
If $N_{\rm in}$ is the number of classes in $\Gamma'/\Gamma$
having a representative $v$ for which $\tau_vA$ has a support
in $\mathsf{Bd}_D$, then
\begin{equation}\label{IV-14-21}
 \log\frac{Z_\Gamma[t\mathbb A]}{Z_\Gamma[0]}
 \ge t|\mathbb A|\nu(A)-24\beta\#\partial_1D
       -2|t|\|A\|_\infty(|\mathbb A|-N_{\rm in}).
\end{equation}
If $A$ has a bond support in $x_0+[0,ra]^3$, with
$x_0\in a\mathbb Z^3$ and integer $r\ge1$, then
\[
 |\mathbb A|-N_{\rm in}\le\#\partial_rD
          \le(2r-1)^3\#\partial_1D.
\]
\end{lemma}
\begin{proof}
Put $B=\mathsf{Bd}_D$. Plaquettes whose four bonds lie in $B$
give identical action terms for $V\vee U_{B^c}$ and the periodic
extension $V^{\rm per}$. They represent distinct period classes.
Let $n_1$ count the remaining plaquettes meeting $B$, and $n_2$
the period classes not represented by an interior plaquette.
Then
\[
 \sum_{p\in\mathsf{Pl}[B]}a_p(V\vee U_{B^c})
        -S_\Gamma(V^{\rm per})=\Sigma_1-\Sigma_2,
 \qquad 0\le\Sigma_i\le2n_i.
\]
If a positive bond starts at $x\notin\partial_1D$, all plaquettes
containing it have their bonds in $B$. There are at most nine
plaquettes meeting the three positive bonds starting at a specified
$x$, and three plaquettes based at $x$. Hence
$n_1\le9\#\partial_1D$ and $n_2\le3\#\partial_1D$.

Choose one good representative for each class counted by $N_{\rm in}$,
and set $\Phi_B=t\sum_{\rm good}\tau_vA$ and
$e=|t|\|A\|_\infty(|\mathbb A|-N_{\rm in})$.
The preceding action comparison, applied to numerator and denominator,
gives for every exterior configuration
\[
 \frac{Z_\Gamma[\Phi_B]}{Z_\Gamma[0]}
       \ge e^{-2\beta(n_1+n_2)}\gamma_B(e^{\Phi_B}).
\]
Jensen's inequality for the normalized kernel, followed by the Gibbs
equation and $\Gamma'$-invariance, yields
\[
 \log\frac{Z_\Gamma[\Phi_B]}{Z_\Gamma[0]}
       \ge tN_{\rm in}\nu(A)-2\beta(n_1+n_2).
\]
On periodic configurations $|t\mathbb A-\Phi_B|\le e$;
also $|t(|\mathbb A|-N_{\rm in})\nu(A)|\le e$.
The two errors give \eqref{IV-14-21}.

For the last estimate, assign each array class the unique point of
$x_0+v+\Gamma$ in $D$. If that point is outside $\partial_rD$,
the translated cube and all its bonds lie in $\mathsf{Bd}_D$.
Different bad classes therefore give different points of
$\partial_rD$. Finally, from $x\in\partial_rD$ follow a lattice path
to a point outside $D$, moving at most $a$ in each coordinate at each
step and using at most $r$ steps. The point immediately before its first exit from $D$
lies in $\partial_1D$ and differs from $x$ by a vector of sup norm
at most $(r-1)a$. Thus $\partial_rD$ is contained in the union of
$(2r-1)^3$ translates of $\partial_1D$.
\end{proof}

Fix admissible $c_0$, $h>0$ with $v_h=hn/|n|^2\in a\mathbb Z^3$,
and put
\[
 \Sigma_j=\Sigma(n,c_0+jh),\quad\theta_j=\theta_{n,c_0+jh},
 \quad\mathbb S_j=\{(j-1)h\le n\cdot x-c_0\le jh\}.
\]
The maps
$\mathsf w_{2k+1}=\tau_{2kv_h}$ and $\mathsf w_{2k}=\theta_k$
carry $\mathbb S_1$ to $\mathbb S_j$. Define
$f^{[j]}=\mathsf w_jf$ for odd $j$ and
$f^{[j]}=\overline{\mathsf w_jf}$ for even $j$.
Writing $\vartheta=\theta_{n,c_0+h/2}$, the reflection formulas give
\begin{equation}\label{IV-15-31}
 \theta_m\mathsf w_j=\mathsf w_{2m+1-j},\quad
 \mathsf w_{j+2N}=\tau_{2Nv_h}\mathsf w_j,\quad
 \tau_{v_h}\mathsf w_j=\mathsf w_{j+1}\vartheta.
\end{equation}
All these maps preserve the lattice. The group generated by the
$\theta_j$ consists of them and the translations $\tau_{2kv_h}$.

\begin{proposition}[Reflection of arrays]\label{IV-16-31}
Suppose $\Gamma$ is adapted to $n$ with height period $2Nh$,
$N\ge1$. If a law $\mu$ on $\mathcal V_\Gamma$ is reflection
positive at every pair $\Sigma_j,\Sigma_{j+N}$, then
\begin{equation}\label{IV-16-33}
 \left|\mu\left(\prod_{j=1}^{2N}f_j^{[j]}\right)\right|
 \le\prod_{j=1}^{2N}
       \mu\left(\prod_{k=1}^{2N}f_j^{[k]}\right)^{1/(2N)}
\end{equation}
for bounded Borel functions of the bonds in $\mathbb S_1$;
each quantity under a root is nonnegative.

Suppose additionally that $\Gamma'\supset\Gamma$ has finite index
and the splitting \eqref{IV-15-23} with height period $h$.
No lower bound on its systole is required. For real
$A\in\mathcal C_{\rm loc}(\mathbb S_1)$,
let $T$ represent $(\Gamma'\cap n^\perp)/(\Gamma\cap n^\perp)$,
and set
\[
 A_T=\sum_{v\in T}\tau_vA,\qquad
 \mathbb A=\sum_{j=1}^{2N}\tau_{(j-1)v_h}A_T,\qquad
 \mathbb A^{\rm f}=\sum_{j=1}^{2N}\mathsf w_jA_T.
\]
These define functions on $\mathcal V_\Gamma$ independent of
representatives, and $|\mathbb A|=2N\#T$. For every real $t$,
\begin{equation}\label{IV-16-34}
 Z_\Gamma[t\mathbb A]\le Z_\Gamma[t\mathbb A^{\rm f}].
\end{equation}
The law $\mu_\Gamma^{t\mathbb A^{\rm f}}$ is reflection positive
at each pair $\Sigma_j,\Sigma_{j+N}$, invariant under the wall
group and under $\Gamma'\cap n^\perp$, and satisfies
\begin{equation}\label{IV-16-35}
 \mu_\Gamma^{t\mathbb A^{\rm f}}(\mathbb A^{\rm f})
       =|\mathbb A|\mu_\Gamma^{t\mathbb A^{\rm f}}(A).
\end{equation}
\end{proposition}
\begin{proof}
Let $\mathrm{cy}(f_1,\ldots,f_{2N})$ denote the integral on the
left of \eqref{IV-16-33}. The copies are cyclic modulo $2N$ on
$\mathcal V_\Gamma$. Reflection in $\Sigma_m$ changes index $j$
to $2m+1-j$ and conjugates the copy. Cauchy--Schwarz for the
positive form at that pair therefore bounds $|\mathrm{cy}|^2$
by the product of the two integrals obtained by reflecting either
half of the cyclic array into the other half. Those two integrals
are nonnegative. In particular constant arrays have nonnegative
integrals.

Assume first that the constant-array roots $\|f_j\|$ are positive.
Among the finitely many arrays with entries chosen from
$f_1,\ldots,f_{2N}$ maximize
$|\mathrm{cy}(\underline f)|/\prod_j\|f_j\|$.
The maximum is at least one and is unchanged by either half-reflection:
the product bound and maximality force both reflected arrays to
attain it. Choose a maximizing array with a longest constant cyclic
string, of length $k$. Reflect across its terminal boundary.
If $k\le N$, this doubles the string; if $N<k<2N$, it makes the
whole array constant. Maximality of $k$ therefore forces $k=2N$,
so the maximum is one. If a root vanishes, replace every $f_j$ by
$f_j+\varepsilon$ with real $\varepsilon$. Each constant-array
integral is a nonnegative polynomial of degree $2N$ with leading
coefficient one. Choose $\varepsilon\to0$ avoiding their finitely
many zeros and pass to the limit. This proves \eqref{IV-16-33}.

For the array comparison take $f=e^{tA_T}$. The splitting of both
period groups gives the displayed expression for $\mathbb A$.
For odd $j$, $\tau_{(j-1)v_h}=\mathsf w_j$; for even $j$ it is
$\mathsf w_j\vartheta$. Thus
\[
 Z_\Gamma[t\mathbb A]/Z_\Gamma[0]
       =\mathrm{cy}(f,\vartheta f,\ldots,f,\vartheta f),
 \quad
 Z_\Gamma[t\mathbb A^{\rm f}]/Z_\Gamma[0]
       =\mathrm{cy}(f,\ldots,f).
\]
The unperturbed measure has the required positivity by
\cref{IV-16-26}. Translation by $v_h$ takes the fold of $f$ to the
fold of $\vartheta f$, by \eqref{IV-15-31}; translation invariance
therefore gives $\|\vartheta f\|=\|f\|$. Inequality
\eqref{IV-16-33} proves \eqref{IV-16-34}.

For any $m$, the folded perturbation is $\Phi_++\theta_m\Phi_+$,
where $\Phi_+=t\sum_{j=m+1}^{m+N}\mathsf w_jA_T$.
Apply the perturbed assertion of \cref{IV-16-26}. This gives
positivity and invariance under every wall reflection.
Translations in $\Gamma'\cap n^\perp$ permute the transverse
classes and commute with the $\mathsf w_j$, so they preserve the
perturbation and the law. All $2N\#T$ summands in
$\mathbb A^{\rm f}$ consequently have the same expectation,
which proves \eqref{IV-16-35}.
\end{proof}

The corresponding infinite folded family is
$(\mathsf w_j\tau_vA)_{j\in\mathbb Z,\,v\in\Gamma'\cap n^\perp}$.
If $B_A$ is a finite nonempty support in the first slab, give its
members the supports $\mathsf w_j(B_A+v)$. A fixed bond meets only
the finitely many slabs containing it; within each such slab only
finitely many translates of the finite set $B_A$ contain that bond.
Thus this is a locally finite perturbation family, to which the
periodic-approximation proposition applies when the index action
and period lattice are specified.

\subsection{Domination of periodic Gibbs states}

Fix $K\ge1$, $a=a_K$ and $\beta>0$. A board on $a\mathbb Z^3$
has cells
\[
 \square_k=o+l k+[-l/2,l/2]^3,\qquad
 l=a_i,\quad i\le K-1,\quad o\in a\mathbb Z^3.
\]
The number $l/a=L^{K-i}$ is odd. Its walls are consequently
bond-midpoint mirrors. For $\ell\in\mathfrak L$ with $\ell\ge2l$,
the quotient board has an even number $\ell/l$ of cells in every
direction. A function of the inner bonds of $\square_k$ defines a
function of the corresponding torus cell; write
$\|F\|_{{\rm ch},\ell}$ for its chessboard seminorm. This
identification commutes with products, mirror copies and complex
conjugation. In particular all the conclusions of
\cref{IV-4-14,IV-4-15,IV-4-17} apply. Translation by a vector in
$2l\mathbb Z^3$ is a mirror copy without conjugation.

The following surface estimate removes the boundary condition.
\begin{proposition}\label{IV-18-28}
Let $\ell\in\mathfrak L$, $y\in a(\mathbb Z+1/2)^3$,
$D=(y+[0,\ell]^3)\cap a\mathbb Z^3$, and let $W$ be the set of
positive bonds whose initial point lies in $D$. For every
nonnegative bounded Borel function $\Psi$ of $W$ and every exterior
configuration $\xi$,
\begin{equation}\label{IV-18-29}
 e^{-36\beta(\ell/a)^2}\widetilde\mu_{K,\ell,\beta}(\Psi)
 \le\gamma_W\Psi(\xi)
 \le e^{36\beta(\ell/a)^2}\widetilde\mu_{K,\ell,\beta}(\Psi).
\end{equation}
The same inequalities hold with $\gamma_W\Psi(\xi)$ replaced by
$\nu(\Psi)$ for any Gibbs state $\nu$. For every real bounded
function $H$ of $W$,
\begin{equation}\label{IV-18-30}
 \nu(H)\le\log\widetilde\mu_{K,\ell,\beta}(e^H)
                          +36\beta(\ell/a)^2.
\end{equation}
\end{proposition}
\begin{proof}
Use the bonds of $W$ as Haar coordinates for both the Gibbs kernel
and the torus measure, and put $n=\ell/a$. In each of the three
plaquette orientations, $n(n-1)^2$ plaquettes have all their bonds
in $W$ and occur with the same action in both measures. There are
$n^3$ torus plaquettes and $n^3+2n^2$ plaquettes touching $W$ in
that orientation. Thus, after cancellation, the Gibbs action has
$N_1=12n^2-3n$ terms and the torus action has $N_2=6n^2-3n$ terms.
Since $0\le a_p\le2$, their difference is between $-2N_2$ and
$2N_1$. Comparing both numerators and both normalizing denominators
gives a density ratio between
$e^{-2\beta(N_1+N_2)}$ and $e^{2\beta(N_1+N_2)}$;
$2(N_1+N_2)=36n^2-12n\le36n^2$ proves \eqref{IV-18-29}.
Integrate the kernel inequality against $\nu$ and use its Gibbs
equation. Finally Jensen's inequality in the kernel gives
$\gamma_WH\le\log\gamma_We^H$; apply \eqref{IV-18-29} and integrate
once more to obtain \eqref{IV-18-30}.
\end{proof}

\begin{proposition}\label{IV-18-38}
Let $\nu$ be a periodic Gibbs state, let $\mathcal K_0$ be a finite
nonempty set of cell indices, and let $F_k$ be bounded Borel
functions of the inner bonds of $\square_k$. For every sequence
$\ell_j\in\mathfrak L$ tending to infinity,
\begin{equation}\label{IV-18-40}
 \left|\nu\left(\prod_{k\in\mathcal K_0}F_k\right)\right|
 \le\liminf_{j\to\infty}\prod_{k\in\mathcal K_0}
                  \|F_k\|_{{\rm ch},\ell_j}.
\end{equation}
More precisely, choose a period number $P$ of $\nu$, an even
positive multiple $q$ of $P$, and $k_0$ so that
$\mathcal K_0\subset k_0+\{0,\ldots,q-1\}^3$. Set
$\mathsf P=ql$, $m=\prod_k\|F_k\|_\infty$ and
\[
 N_\ell=\lfloor\ell/\mathsf P\rfloor^3,\qquad
 \epsilon_\ell=36\beta(\ell/a)^2/N_\ell
       \le288\beta\mathsf P^3a^{-2}\ell^{-1}
       \quad(\ell\ge\mathsf P).
\]
For $\ell\ge\mathsf P$, if $t>0$ and $2tm\le1$,
\begin{equation}\label{IV-18-39}
 \left|\nu\left(\prod_kF_k\right)\right|
 \le\prod_k\|F_k\|_{{\rm ch},\ell}+tm^2+\epsilon_\ell/t.
\end{equation}
For a real bounded function $H$ of one cell, the same packing, with
$\ell\ge\mathsf P$, gives
\begin{equation}\label{IV-18-41}
 \nu(H)\le\log\|e^H\|_{{\rm ch},\ell}+\epsilon_\ell.
\end{equation}
If $\nu$ is invariant under every reflection in the walls of the
board, then for every $\ell\ge2l$,
\begin{equation}\label{IV-18-42}
 \nu(H)\le\log\|e^H\|_{{\rm ch},\ell}
                              +36\beta l^3a^{-2}\ell^{-1}.
\end{equation}
\end{proposition}
\begin{proof}
The vectors
\[
 V_\ell=\{v\in\mathsf P\mathbb Z^3:
                    0\le v_\mu\le\ell-\mathsf P\}
\]
are periods of $\nu$, since $l/a$ is an integer and $P\mid q$.
They belong to $2l\mathbb Z^3$, so $\tau_vF_k$ is the mirror copy
in cell $k+v/l$. All these cell indices are distinct and fit in
one $\ell$-period box: division with remainder by $q$ separates
the translated blocks. Their inner bonds lie in the set $W$ of
\cref{IV-18-28}, with corner
$y=o+l(k_0-(1/2,1/2,1/2))$. Applying \eqref{IV-18-30} to the sum
of the $N_\ell$ translates of a real local function $\Phi$ gives
\begin{equation}\label{IV-18-37}
 N_\ell\nu(\Phi)\le36\beta(\ell/a)^2+
 \log\widetilde\mu_{K,\ell,\beta}
                         \left(\prod_{v\in V_\ell}e^{\tau_v\Phi}\right).
\end{equation}
Here $\#V_\ell=N_\ell$, and
$\lfloor\ell/\mathsf P\rfloor\ge\ell/(2\mathsf P)$ proves the
displayed bound on $\epsilon_\ell$.

Write $F=\prod_kF_k$, and choose $|\zeta|=1$ so that
$\zeta\nu(F)=|\nu(F)|$. If $m=0$ there is nothing to prove.
Otherwise put $G=\operatorname{Re}(\zeta F)$ and
$\Phi=\log(1+tG)$. Since $|tG|\le1/2$ and
$\log(1+x)\ge x-x^2$ on $[-1/2,1/2]$,
\[
 \nu(\Phi)\ge t|\nu(F)|-t^2m^2.
\]
Expand $\prod_{v\in V_\ell}(1+t\tau_vG)$ over subsets
$Y\subset V_\ell$. Each product $\prod_{v\in Y}\tau_vG$ is
$2^{-|Y|}$ times a sum of $2^{|Y|}$ products of mirror copies of
the $F_k$ or $\overline{F_k}$, with coefficients of modulus one.
The finite chessboard estimate and invariance of its seminorm
under copying and conjugation bound each such mean by
$\Pi_\ell^{|Y|}$, where $\Pi_\ell=\prod_k\|F_k\|_{{\rm ch},\ell}$.
Consequently
\[
 \widetilde\mu_{K,\ell,\beta}
      \left(\prod_{v\in V_\ell}e^{\tau_v\Phi}\right)
 \le(1+t\Pi_\ell)^{N_\ell}\le e^{N_\ell t\Pi_\ell}.
\]
Insert this in \eqref{IV-18-37} and divide by $N_\ell t$ to
obtain \eqref{IV-18-39}. First let $\ell=\ell_j\to\infty$ at
fixed $t$, and then let $t\downarrow0$, proving
\eqref{IV-18-40}. Applying \eqref{IV-18-37} directly to $H$ and
using the chessboard bound $\|e^H\|_{{\rm ch},\ell}^{N_\ell}$
proves \eqref{IV-18-41}. The seminorm is positive because its
defining integrand is bounded below by a positive number.

For \eqref{IV-18-42}, sum all the mirror copies of $H$ in an
$\ell$-period box. The mean of this sum is $(\ell/l)^3\nu(H)$,
while the torus exponential mean is exactly
$\|e^H\|_{{\rm ch},\ell}^{(\ell/l)^3}$. Use
\eqref{IV-18-30} and divide by $(\ell/l)^3$.
\end{proof}

\subsection{A volume threshold for local comparison}

The next threshold depends on the cutoff and coupling. No
uniform estimate for its size is asserted.
Let $W_K$ be the positive bonds with both endpoints in
$[-K,K]^3$, and let $\partial W_K$ be the exterior bonds of
plaquettes touching $W_K$. Denote the continuous density of
$\gamma_{W_K}$ by $k_K(V,\xi_{\partial W_K})$ and put
\[
 q_\nu(V)=\int k_K(V,\xi_{\partial W_K})\,d\nu(\xi),\qquad
 \varepsilon_{K,\beta}(\ell)=
 \inf_{\nu\in\mathcal G^{\rm tor}_{K,\beta}}
 \left\|\frac{q_\nu}{q_{\widetilde\mu_{K,\ell,\beta}}}-1
                                                   \right\|_\infty.
\]
\begin{proposition}\label{IV-18-50}
One has $\varepsilon_{K,\beta}(\ell)\to0$ as
$\ell\to\infty$ in $\mathfrak L$. Thus there is a least
$\ell_\flat(K,\beta)\in\mathfrak L$, greater than $4(K+1)$,
such that $\varepsilon_{K,\beta}(\ell)<e^{-L^K}$ for all
$\ell\ge\ell_\flat(K,\beta)$. For each such $\ell$ one can choose
a single $\nu\in\mathcal G^{\rm tor}_{K,\beta}$ satisfying
\[
 |\nu(Y)-\widetilde\mu_{K,\ell,\beta}(Y)|
 \le e^{-L^K}\widetilde\mu_{K,\ell,\beta}(|Y|)
\]
for every bounded Borel function $Y$ of $W_K$. If $0<r\le K-1$
and $f_1,\ldots,f_n$ are bounded functions supported in
$[-r,r]^3$, this same state satisfies
\begin{equation}\label{IV-18-51}
 \left|S^T_{n;K,\ell,\beta}
          (f_1^{\langle\ell\rangle},\ldots,f_n^{\langle\ell\rangle})
                  -S_n^T[\nu](f_1,\ldots,f_n)\right|
 \le c_n^{\rm law}e^{-L^K}
       \bigl(6L^K(2rL^K+1)^3\bigr)^n\prod_q\|f_q\|_\infty,
\end{equation}
where $c_n^{\rm law}=\sum_{\pi\in\operatorname{Part}(n)}(\#\pi)!$.
\end{proposition}
\begin{proof}
If $M$ is the number of plaquettes touching $W_K$, then
$e^{-2\beta M}\le k_K\le e^{2\beta M}$. The kernel is uniformly
continuous on the compact space of its finite interior and
exterior coordinates. Thus the functions $q_\nu$ have these same
positive bounds and are equicontinuous. If $\nu_j\to\nu$ weakly,
then $q_{\nu_j}(V)\to q_\nu(V)$ for every $V$; a finite net in
the interior coordinate space and equicontinuity upgrade this
to uniform convergence.

The Gibbs equation shows that $q_\nu$ is the $W_K$-marginal
density for every torus-limit state. It is also the marginal
density for $\widetilde\mu_{K,\ell,\beta}$ when $\ell>4(K+1)$:
the plaquettes touching $W_K$ fit in $[-K-a,K+a]^3$, whose side
is less than the period systole minus $a$, so the local
periodic Gibbs equation applies.

If $\varepsilon_{K,\beta}(\ell)$ did not tend to zero, take a
diverging sequence of sides on which it stays bounded below.
A subsequence of the corresponding laws converges to a member
$\nu_*$ of $\mathcal G^{\rm tor}_{K,\beta}$. Uniform convergence
of the marginal densities and their common positive lower bound
then give $\|q_{\nu_*}/q_{\widetilde\mu_{K,\ell,\beta}}-1\|_\infty
\to0$, a contradiction. The asserted state exists by the strict
inequality in the definition of the threshold, and the first
estimate follows by integrating the density ratio.

For the cumulant estimate, $X_q=\mathcal O_K(f_q)$ depends only
on $W_K$, since every plaquette centered in $[-r,r]^3$ lies in
$[-K,K]^3$. Counting at most $3(2rL^K+1)^3$ plaquette centers
and using $a_p\le2$ gives
$\|X_q\|_\infty\le6L^K(2rL^K+1)^3\|f_q\|_\infty$.
The density-ratio bound applies to every product of the $X_q$.
Expand the two cumulants over partitions. In a partition with
$k$ blocks, telescope the product of its $k$ moments; the error
is at most $k e^{-L^K}\prod_q\|X_q\|_\infty$ and the cumulant
coefficient has modulus $(k-1)!$. Summing over partitions and
using \eqref{IV-17-14} proves \eqref{IV-18-51}.
\end{proof}

\subsection{Transfer of cell normalizers}

Let $f_1,\ldots,f_n\in\operatorname{Lip}_c(\mathbb R^3;\mathbb R)$
be nonzero and suppose that each support lies in the interior
of one board cell, indexed by $k(q)$. Supports in the same cell
may overlap. Put
\[
 J_k=\{q:k(q)=k\},\quad w_k(z)=\sum_{q\in J_k}z_qf_q,\quad
 w(z)=\sum_qz_qf_q,\quad B(z)=\{q:z_q\ne0\},
\]
\begin{equation}\label{IV-19-3}
 \Pi_\ell(z)=\prod_{k:J_k\ne\varnothing}
            \|e^{\mathcal O_K(w_k(z))}\|_{{\rm ch},\ell},
 \qquad
 \mathbb D_\rho(\underline f)=
       \left\{z:\sum_q|z_q|\|f_q\|_\sharp<\rho\right\}.
\end{equation}
A holomorphic $E$ on $\mathbb D_\rho(\underline f)$ with $E(0)=0$
is a \emph{normalizer with bound $Q$} on a set of sides
$\mathfrak l\subset\{\ell\in\mathfrak L:\ell\ge2l\}$ if
$\Pi_\ell(z)\le|e^{E(z)}|Q$ there. It is a normalizer
\emph{for periodic states} if there is one sequence
$\ell_j\to\infty$ such that
\begin{equation}\label{IV-19-6}
 \liminf_{j\to\infty}\Pi_{\ell_j}(z)\le|e^{E(z)}|Q
       \quad\hbox{for every }z\in\mathbb D_\rho(\underline f).
\end{equation}
A \emph{graded} normalizer replaces $Q$ by $Q^{\#B(z)}$.
In particular a normalizer on an unbounded set of sides is a
normalizer for periodic states. The sequence in
\eqref{IV-19-6} is chosen before $z$; no uniform convergence of
$\Pi_{\ell_j}$ is required.

\begin{proposition}\label{IV-19-13}
A normalizer for periodic states gives, for every periodic Gibbs
state $\nu$,
\begin{equation}\label{IV-19-14}
 |e^{-E(z)}\nu(e^{\mathcal O_K(w(z))})|\le Q.
\end{equation}
A normalizer on $\mathfrak l$ gives the same bound with $\nu$
replaced by $\widetilde\mu_{K,\ell,\beta}$ if $\ell\in\mathfrak l$
and the occupied cell indices are distinct modulo
$(\ell/l)\mathbb Z^3$. Both assertions hold with
$Q^{\#B(z)}$ for a graded normalizer.

For any law $\mathfrak s$ satisfying \eqref{IV-19-14},
\begin{equation}\label{IV-19-17}
 \left|S_n^T[\mathfrak s](\underline f)
                   -\partial_1\cdots\partial_nE(0)\right|
 \le\log2\,(2Q+2)^n(n/\rho)^n\prod_q\|f_q\|_\sharp.
\end{equation}
Under the graded bound this improves to
\begin{equation}\label{IV-19-18}
 \left|S_n^T[\mathfrak s](\underline f)
                   -\partial_1\cdots\partial_nE(0)\right|
 \le(n-1)!(2^n-1)(eQ/\rho)^n\prod_q\|f_q\|_\sharp.
\end{equation}
If two laws satisfy the same ungraded normalizer bound, their
cumulants differ by at most twice the right side of
\eqref{IV-19-17}, including for $n=1$ and for overlapping supports.
\end{proposition}
\begin{proof}
Every $e^{\mathcal O_K(w_k)}$ depends on the inner bonds of its
cell, since a plaquette with center inside a cell has all its
bonds inside that cell. Apply \eqref{IV-18-40} to their product,
then \eqref{IV-19-6}, obtaining \eqref{IV-19-14}.

For a torus, move each occupied cell by an $\ell\mathbb Z^3$
translation into one period box. The images are distinct by
hypothesis. Such translations belong to $2l\mathbb Z^3$, hence
give mirror copies with unchanged chessboard seminorms; they
also leave the periodic configurations unchanged. The finite
chessboard estimate therefore bounds the exponential mean by
$\Pi_\ell(z)$. This proves the torus assertion, with the identical
argument for graded bounds. Finally the exponential moment is
entire, equals one at zero, and its logarithmic mixed derivative
is the cumulant. Apply \cref{IV-5-7,IV-5-13} to
$e^{-E(z)}\mathfrak s(e^{\mathcal O_K(w(z))})$ with weights
$\|f_q\|_\sharp/\rho$. This gives the two derivative estimates.
For two laws the common derivative of $E$ cancels.
\end{proof}

\subsection{Uniform normalizers and cumulant bounds}

Fix the source radius $\varrho$ and $0<g\le\gamma_\varrho$ as
in \cref{III-111-27}. Write
\begin{equation}\label{IV-22-3}
 \mathfrak d(g)=\{(K,\beta):K\ge1,\ \beta\in\mathcal W_K(g)\}.
\end{equation}
The admissible sides form the unbounded set $\mathfrak L(g)$ of
\eqref{III-1-14}. Throughout this subsection take
\[
 C_e=C_{\mathrm S},\quad\kappa_e=\kappa,\quad
 E_+=C_{\mathrm S},\quad
 C_{\mathrm{low}}=C_{\mathrm S}(1+\log g^{-2}),\quad
 C_{\mathrm{gap}}=C_{\mathrm S}.
\]
In particular the constants are independent of $(K,\beta)$ and
of the side. Put $A_0=E_++C_{\mathrm{low}}+C_{\mathrm{gap}}$ and
\begin{equation}\label{IV-19-22}
 \mathsf Q_1=\exp\{A_0L^3+\Phi_1(2,1;L)\},\qquad
 Q^\natural(n,R)=
       \exp\{A_0h(n,R)^3+\Phi_1(R,R+1;M)\}.
\end{equation}
Here $h(n,R)$ is \eqref{IV-8-7}. With $p$ and $C_Q^0$ as in
\eqref{IV-8-4} and \eqref{IV-8-5}, set
$C_Q=e^{C_{\mathrm{gap}}}C_Q^0$ and
$Q_s(\eta)=C_Q(2\eta)^{-p}$.

\begin{proposition}\label{IV-19-24}\label{IV-22-13}
For every $(K,\beta)\in\mathfrak d(g)$ the following normalizers
have radius $\varrho$ and work on all indicated sides, hence
also for every periodic Gibbs state.
\begin{enumerate}
\item If all supports lie in one ball $\mathbb B(x_0,2)$, on a
board of side $L$ with $x_0$ at distance at least $3$ from its
walls, there is a normalizer with bound $\mathsf Q_1$ on
$\mathfrak L(g)$, for any number of entries and any overlap.
\item If $n\ge2$, the supports are disjoint, and
$\operatorname{supp}f_q\subset\mathbb B(x_q,R)$ on a board of
side $h(n,R)$ whose walls have distance at least $2R+1$ from
every $x_q$, there is a graded normalizer with base
$Q^\natural(n,R)$ on the sides $\ell\in\mathfrak L(g)$ with
$\ell\ge2h(n,R)$.
\item If $n\ge2$, $0<\eta\le1/4$, the supports are disjoint,
and $\operatorname{supp}f_q\subset\mathbb B(y_q,\eta)$ on a
board of side one whose walls have distance at least $2\eta$
from every $y_q$, there is a graded normalizer with base
$Q_s(\eta)$ on $\mathfrak L(g)$.
\end{enumerate}
The normalizers in the last two cases satisfy
\begin{equation}\label{IV-19-9}
 |\partial_1\cdots\partial_nE(0)|
 \le(n/\varrho)^n\mathsf U(R,d)\prod_q\|f_q\|_\sharp,
\end{equation}
with $R$ replaced by $\eta$ in the third case and
$d=\min_{q\ne q'}\operatorname{dist}(\operatorname{supp}f_q,
\operatorname{supp}f_{q'})>0$.
\end{proposition}
\begin{proof}
For a board of side $l$, choose the reference side $\ell'$ to be
the least member of $\mathfrak L(g)$ with $\ell'\ge2l$. The
source system on $T_{\ell'}$ is supplied by
\cref{III-111-27} at radius $2\varrho$. The same theorem gives
$Z_{K,\ell,\beta}\ge e^{E_\ell-C_{\mathrm{low}}\ell^3}$ and
$E_{\ell'}\ell'^{-3}-E_\ell\ell^{-3}\le C_{\mathrm{gap}}$
for every larger admissible side. These are precisely the
floor and gap hypotheses of \cref{IV-7-43}. Periodization
preserves the sharp norms of functions in a cell interior,
so the parameter polydiscs agree on all these tori.

For each occupied cell construct $\mathcal E_k^\circ$ on the
reference torus, as in \eqref{IV-7-42}, and sum over the cells.
This is one holomorphic function, independent of the larger
side. If the supports lie in balls of radius $R$ with wall
margin $D$, then \eqref{IV-7-45} gives
\[
 |e^{-E(z)}|\Pi_\ell(z)
 \le\exp\{A_0l^3+\Phi_1(R,D;l)\}^{\#B(z)}.
\]
Each summand of $E$ involves only the variables in its cell.
Thus its full mixed derivative vanishes when at least two
cells are occupied. When exactly one cell is occupied,
\eqref{IV-7-47} gives \eqref{IV-19-9}, initially with
$2\varrho$ in place of $\varrho$.

For the second case take $D=R+1$ and use
$\Phi_1(R,R+1;l)\le\Phi_1(R,R+1;M)$. For the third take
$l=1$ and $R=D=\eta$. Directly from the formula for $\Phi_1$,
\[
 \Phi_1(\eta,\eta;1)
 =\Phi_1(1/2,1/2;1)+p\log(1/(2\eta)),
\]
so the base is $Q_s(\eta)$. In the first case take the single
carrier $x_0+[-2,2]^3$, with wall margin one, and use the
single-carrier estimate of \cref{IV-7-27}. Its error is
$\Phi_1(2,1;L)$ once for the carrier; it has no factor equal
to the number of entries. This gives $\mathsf Q_1$.
All admissible sides are at least $2L$, so the first and
third cases cover the whole of $\mathfrak L(g)$.

Finally restrict these functions from
$\mathbb D_{2\varrho}(\underline f)$ to
$\mathbb D_\varrho(\underline f)$. All inequalities remain
valid, and the derivative estimate only becomes weaker.
Each indicated set of sides is unbounded, which proves the
assertion about periodic states.
\end{proof}

Define
\begin{align}
 C_1^g(R)&=2\log2\,(2\mathsf Q_1+2)\varrho^{-1}
              (1+40\sqrt3M)(2R+3)^3,\label{IV-20-3}\\
 C_n^{\rm per}(d;R)&=(n/\varrho)^n
  \left[\log2\,(2Q^\natural(n,R)^n+2)^n+\mathsf U(R,d)\right],
                                                    \label{IV-20-4}\\
 C_B&=\max\{28^3\cdot961\cdot2e^5(8e)^pC_QM/\varrho,
                         \ 384\cdot64^3e^4\},\label{IV-20-5}\\
 A^{\rm per}(d)&=1+\mathsf U(1,d),\qquad
 B^{\rm per}(R)=C_B(R+1)^3,\label{IV-20-6}\\
 C_n^g(d;R)&=\min\{C_n^{\rm per}(d;R),
           A^{\rm per}(d)B^{\rm per}(R)^n(n!)^{5+p}\}.
                                                    \label{IV-22-7}
\end{align}

\begin{proposition}\label{IV-20-12}\label{IV-20-16}
\label{IV-20-24}\label{IV-22-14}
For every $(K,\beta)\in\mathfrak d(g)$ and every two periodic
Gibbs states $\nu,\nu'$, a real Lipschitz function supported in
a ball of radius $R$ satisfies
\begin{equation}\label{IV-22-15}
 |\nu(\mathcal O_K(f))-\nu'(\mathcal O_K(f))|
                         \le C_1^g(R)\|f\|_\sharp.
\end{equation}
For $n\ge2$ and a tuple of type $(d;R)$,
\begin{align}
 |S_n^T[\nu](\underline f)|
 &\le C_n^{\rm per}(d;R)\prod_q\|f_q\|_\sharp,
                                                   \label{IV-22-16}\\
 |S_n^T[\nu](\underline f)|
 &\le A^{\rm per}(d)B^{\rm per}(R)^n(n!)^{5+p}
                              \prod_q\|f_q\|_\sharp.\label{IV-22-17}
\end{align}
The one-point comparison remains true when either or both laws
are periodically extended torus measures with side in
$\mathfrak L(g)$. The two cumulant bounds hold for the torus
measure if all supports lie in $(-\ell/4,\ell/4)^3$; the first
also requires $\ell\ge2h(n,R)$, equivalently
$\ell\ge16n(R+1)$, whereas the second has no additional side
restriction. All constants are independent of cutoff, state,
side and position of the supports.
\end{proposition}
\begin{proof}
We first record the placement step. Points $x_1,\ldots,x_n$
can be put at distance at least $r+D$ from the walls of a board
of side $l$ by an offset in $a_K\mathbb Z^3$ whenever
\[
 2n(r+D)+na_K\le l.
\]
This is the offset lemma used in the proof of
\cref{IV-8-21}: in each coordinate the forbidden offsets are
$n$ intervals of total length at most $2n(r+D)$; the extra
$na_K$ allows their endpoints to be rounded to the lattice.
For one ball of radius two use $l=L\ge13$, $r+D=3$.
For $n$ balls of radius $R$ use $l=h(n,R)\ge8n(R+1)$ and
$r+D=2R+1$. For $n$ balls of radius $\eta$ use $l=1$ and
$r+D=2\eta$, provided $4n\eta+na_K\le1$.

For functions in one ball of radius two, the first normalizer
in \cref{IV-19-24} and the two-law part of
\cref{IV-19-13} give
\[
 |S_n^T[\nu](\underline f)-S_n^T[\nu'](\underline f)|
 \le2\log2\,(2\mathsf Q_1+2)^n(n/\varrho)^n
                                  \prod_q\|f_q\|_\sharp.
\]
Only one cell is occupied, so this holds as well for any two
of the stated torus or periodic Gibbs laws. For a function
supported in a ball of radius $R$, decompose it using the
unit-mesh tensor-product tents. There are at most $(2R+3)^3$
nonzero pieces, each supported in a ball of radius $\sqrt3<2$
and with sharp norm at most $(1+40\sqrt3M)\|f\|_\sharp$.
Use the preceding estimate at $n=1$ and sum. This proves
\eqref{IV-22-15}.

For a separated tuple of type $(d;R)$, place the balls on the
board of side $h(n,R)$. Apply the second normalizer with its
ungraded bound $Q^\natural(n,R)^n$. Adding
\eqref{IV-19-9} to \eqref{IV-19-17} proves
\eqref{IV-22-16}.

For the growth estimate first suppose $a_K>1/(32n)$. Counting
plaquettes in a radius-$R$ ball gives
\[
 \|\mathcal O_K(f_q)\|_\infty
 \le6a_K^{-1}(2R/a_K+1)^3\|f_q\|_\infty
 \le192\cdot64^3 n^4(R+1)^3\|f_q\|_\sharp.
\]
The bounded-variable estimate of \cref{IV-5-13} and
$n^{4n}\le e^{4n}(n!)^4$ give
\[
 |S_n^T[\mathfrak s](\underline f)|
 \le[384\cdot64^3e^4(R+1)^3]^n(n!)^5
                                  \prod_q\|f_q\|_\sharp
\]
for every law $\mathfrak s$. This is bounded by
\eqref{IV-22-17}.

Now suppose $a_K\le1/(32n)$. Put
$\eta=1/(8n)$ and $\varepsilon=\eta/\sqrt3$. Decompose each
$f_q$ into tents of mesh $\varepsilon$. Each piece has support
in a ball of radius $\eta$, contained in the original support,
and norm at most $(1+960Mn)\|f_q\|_\sharp$. There are at most
$(16\sqrt3nR+3)^{3n}$ tuples of pieces. For each such tuple,
$4a_K\le\eta$ and $4n\eta+na_K\le1$, so the side-one board
and the third normalizer are available. Its graded transfer
bound gives
\[
 |S_n^T[\nu](v_1,\ldots,v_n)|
 \le\varrho^{-n}
 \left[(n-1)!(2^n-1)(eQ_s(\eta))^n
                         +n^n\mathsf U(\eta,d)\right]
                         \prod_q\|v_q\|_\sharp.
\]
Use $Q_s(\eta)\le C_Q(8n)^p$,
$\mathsf U(\eta,d)\le\mathsf U(1,d)$ and sum the pieces.
To see the resulting exponent explicitly, bound the number
of pieces by $[28n(R+1)]^{3n}$ and their norm factor by
$(961Mn)^n$. These account for four powers of $n!$, using
$n^n\le e^nn!$. The first term in the square brackets is
bounded by $[2eC_Q(8e)^p]^n(n!)^{1+p}$; the second is bounded
by $e^n(n!)^{1+p}\mathsf U(1,d)$. The first term in $C_B$
dominates both resulting constants. This proves
\eqref{IV-22-17}.

For the torus assertions use \cref{IV-19-13}. Distinct cells
meeting $(-\ell/4,\ell/4)^3$ are incongruent modulo the torus
period whenever $\ell\ge2l$: if their indices differed by a
nonzero period, two points in these cells would have some
coordinate difference at least $\ell-l\ge\ell/2$. Tent pieces
retain this central-cube support. Thus the preceding proofs
apply to the torus at the asserted side restrictions; the
coarse-lattice bound already applies to every law. Finally,
for $\ell=2L^m$, the definition of $h(n,R)$ makes
$\ell\ge2h(n,R)$ equivalent to $\ell\ge16n(R+1)$.
\end{proof}

\section{Continuum limits on separated tuples}
\label{sec:continuum-subsequences}

The bounds just proved concern compact supports of arbitrary
position. They therefore give a common subsequence for all
separated test tuples. They do not select a unique continuum
limit on $\mathbb R^3$.

For $(K,\beta)\in\mathfrak d(g)$ the physical label
$a_K\beta$ lies in the compact interval
\begin{equation}\label{IV-22-4}
 I^{\rm lab}(g)=
 [g^{-2}+B_\sharp-L^{-1}\bar B,\ L(g^{-2}+\bar B)]
                                           \subset(0,\infty).
\end{equation}
This follows from \cref{III-5-11}. Moreover
\cref{III-5-19} shows that for every
$b\in I(g)=(g^{-2}+\bar B,X_g-\bar B]$, the pair
$(K,L^Kb)$ lies in $\mathfrak d(g)$ for every $K\ge1$.

A \emph{limit along growing tori} is a family $S=(S_n)_{n\ge2}$
obtained from $(K_j,\beta_j)\in\mathfrak d(g)$,
$K_j\to\infty$, and $\ell_j\in\mathfrak L(g)$,
$\ell_j\to\infty$, by
\[
 S_n(X)=\lim_j S^T_{n;K_j,\ell_j,\beta_j}
                         (X^{\langle\ell_j\rangle})
 \qquad(X\in\mathfrak S_n(\mathbb R^3)).
\]
A \emph{limit through periodic states} replaces these cumulants
by $S_n^T[\nu_j](X)$ with
$\nu_j\in\mathcal G^{\rm per}_{K_j,\beta_j}$. If each
$\nu_j\in\mathcal G^{\rm tor}_{K_j,\beta_j}$, call it a
\emph{volume-first limit}. A defining sequence has label $b$
when $a_{K_j}\beta_j\to b$. Labels describe defining sequences;
no uniqueness of a label from the limit family is assumed.

\begin{proposition}\label{IV-21-25}\label{IV-22-21}
Every growing-torus sequence as above, and every sequence
through periodic states, has a subsequence on which all
cumulants of order at least two converge simultaneously on
$\mathfrak S(\mathbb R^3)$. One can also require the labels to
converge. The limit is real, symmetric and multilinear and
satisfies
\begin{equation}\label{IV-22-22}
 |S_n(X)|\le C_n^g(d;R)\prod_q\|f_q\|_\sharp
                 \quad\hbox{for every tuple of type }(d;R).
\end{equation}
Limits along growing tori and volume-first limits exist with
every label $b\in I(g)$. All these limit families lie in the
common compact metrizable set of symmetric multilinear
families satisfying \eqref{IV-22-22}, with the topology of
pointwise convergence on separated tuples.

More generally, for arbitrary periodic states
$\rho_j,\nu_j\in\mathcal G^{\rm per}_{K_j,\beta_j}$ one can
choose the same subsequence so that
\[
 \delta_j(f)=\rho_j(\mathcal O_{K_j}(f))
                       -\nu_j(\mathcal O_{K_j}(f))
       \longrightarrow\delta(f),
 \qquad
 \sigma_{\rho_j,\nu_j}(X)\longrightarrow\sigma(X)
\]
for all compactly supported $C^1$ functions $f$ and all separated
tuples $X$. Here $\delta$ is real linear,
$|\delta(f)|\le C_1^g(R)\|f\|_\sharp$ on radius-$R$ balls,
and $S_n^\sigma=S_n$ for $n\ge2$, while $S_1^\sigma=\delta$.
\end{proposition}
\begin{proof}
For growing tori take the domains
$O_j=(-\ell_j/4,\ell_j/4)^3$. Every compact set is eventually
contained in $O_j$. For fixed $n,R$, the side condition
$\ell_j\ge2h(n,R)$ also holds eventually, independently of
the position of a tuple in $O_j$. Thus
\cref{IV-22-14} gives the bound table $C_n^g(d;R)$ required by
\cref{IV-10-25}. For periodic states use the constant domains
$\mathbb R^3$ and the same table at every index. Include the
bounded sequence $a_{K_j}\beta_j$ in the extraction.
This proves simultaneous convergence and inheritance of
\eqref{IV-22-22}.

For the last assertion include $\delta_j$ as the form of order
one, with the table $C_1(d;R)=C_1^g(R)$ supplied by
\eqref{IV-22-15}. The moment--cumulant identity is
\begin{equation}\label{IV-21-10}
 \sigma_{\rho_j,\nu_j}(X)=
 \sum_{\pi\in\operatorname{Part}(n)}
    \prod_{\{q\}\in\pi}\delta_j(f_q)
    \prod_{B\in\pi,\,|B|\ge2}S_{|B|}^T[\rho_j](X_B).
\end{equation}
This finite polynomial passes to the limit on each tuple.
The inverse partition identity gives its asserted cumulants.
If $\rho_j=\nu_j$, this limit is exactly the centered moment
functional $\sigma_S$ of \eqref{IV-9-14}.

For existence at a prescribed $b\in I(g)$, choose
$K_j=j$, $\beta_j=L^jb$, and either any admissible sides
tending to infinity or any
$\nu_j\in\mathcal G^{\rm tor}_{j,L^jb}$. The latter sets are
nonempty by \cref{IV-14-16}. Extraction preserves the label $b$.

The compactness assertion is \cref{IV-10-28} applied to the
table $C_n^g(d;R)$.
\end{proof}

\subsection{Comparison of the orders of limits}

\begin{proposition}\label{IV-21-28}\label{IV-21-33}
Every volume-first limit is a limit along growing tori with
the same pairs $(K_j,\beta_j)$ and the same label. The sides
can be chosen above any prescribed sequence $\lambda_j$.
Conversely, a convergent growing-torus sequence with
$\ell_j\ge\ell_\flat(K_j,\beta_j)$ has the same limit and
label as a sequence through torus-limit states at the same
pairs, without passing to a further subsequence.
\end{proposition}
\begin{proof}
Fix $K,\beta$ and $\nu\in\mathcal G^{\rm tor}_{K,\beta}$,
and take a sequence $\ell_i\to\infty$ whose periodic laws
converge to $\nu$. For a fixed ball $\mathbb B(0,r)$ there are
only finitely many plaquette centers in it. Weak convergence
gives convergence of every joint cumulant of these finitely
many plaquette variables, with repetitions allowed. For
$n\le N$, expand
\[
 S_n^T[\mathfrak s](\underline f)
 =a_K^{-n}\sum_{p_1,\ldots,p_n}
          \prod_qf_q(x(p_q))
          \langle a_{p_1};\ldots;a_{p_n}\rangle_{\mathfrak s}.
\]
The finite sums show that for all bounded functions supported
in this ball the torus--state difference is at most
$\epsilon\prod_q\|f_q\|_\infty$ once $i$ is large enough,
uniformly for $n\le N$.

At the $j$th pair choose $N=r=j$ and $\epsilon=1/j$,
and choose $\ell_j$ from a sequence defining $\nu_j$ with
$\ell_j\ge\max\{\lambda_j,j,2L^{m_{\mathfrak L}(g)}\}$.
The resulting torus--state difference tends to zero for every
fixed tuple, along the full sequence of indices. This proves
the first assertion and retains the label because the pairs
are unchanged.

For the converse choose the state $\nu_j$ of
\cref{IV-18-50} for each $j$. For a fixed tuple supported in
$[-r,r]^3$, eventually $r\le K_j-1$, so
\eqref{IV-18-51} bounds the difference by a constant times
$e^{-L^{K_j}}L^{4nK_j}$. This tends to zero. The one chosen
state works for all tuples, and the pairs again remain
unchanged.
\end{proof}

\subsection{Translation invariance and the one-point constant}

Uniform bounds permit the test functions to move with the
cutoff. More precisely, suppose all entries $f_q^{(j)}$ and
$f_q$ lie in fixed pairwise disjoint compact sets and
$\|f_q^{(j)}-f_q\|_\sharp\to0$. Telescoping one entry at a
time and using the common bound table gives
\begin{equation}\label{IV-24-33}
 \Phi_n^{(j)}(f_1^{(j)},\ldots,f_n^{(j)})
                              \longrightarrow\Phi_n(f_1,\ldots,f_n).
\end{equation}
The norms of the moving entries remain bounded. The same
conclusion holds for centered moments, by their finite
partition formula, and for the two-state moments in
\eqref{IV-21-10} using the order-one bound. Enlarging the
fixed compact supports slightly gives a useful instance:
if $v_j\to v$, then $\tau_{v_j}f\to\tau_vf$ in the sharp norm
for $f\in C_c^1$, since $f$ and its first derivatives are
uniformly continuous.

\begin{proposition}\label{IV-24-11}\label{IV-24-15}
Every growing-torus limit and every volume-first limit has
$\mathbb R^3\rtimes W_3$-invariant cumulants and centered
moments. The same holds for a limit through states whenever
both the integrating and centering laws are invariant under
$a_{K_j}\mathbb Z^3\rtimes W_3$.

In the two-state case the one-point function is
$\sigma(f)=c\int f$, for a constant $c\in\mathbb R$. Its
higher cumulants are $S_n$ and its centered functional is
$\sigma_S$. If the states obey \eqref{IV-22-15}, then
\begin{equation}\label{IV-24-19}
 |c|\le\frac{120M}{\pi}C_1^g(1).
\end{equation}
\end{proposition}
\begin{proof}
For $v\in\mathbb R^3$ choose $v_j\in a_{K_j}\mathbb Z^3$
with $v_j\to v$. The finite-cutoff cumulants and, when
applicable, two-state moments are invariant under $\tau_{v_j}$.
Use \eqref{IV-24-33} to pass to $\tau_v$. In the growing-torus
case all supports and their small translates lie in the
central cube eventually. The signed coordinate permutations
in $W_3$ preserve every lattice, so their invariance passes
directly to the limit. Torus-limit states have these lattice
invariances by \cref{IV-14-16}.

For the one-point function, lattice isometries act transitively
on plaquettes. Thus
\[
 c_j=3a_{K_j}^{-4}(\rho_j(a_p)-\nu_j(a_p))
\]
does not depend on $p$, and exactly
\begin{equation}\label{IV-24-16}
 \delta_j(f)=c_jr_j(f),\qquad
 r_j(f)=\frac{a_{K_j}^3}{3}\sum_p f(x(p)).
\end{equation}
For each orientation its centers form a translate of
$a_{K_j}\mathbb Z^3$. Partition space into half-open cubes of
side $a_{K_j}$ centered there. On each cube the pointwise
error from its center is at most
$(\sqrt3/2)a_{K_j}\operatorname{Lip}f$. Only cubes in the
closed $\sqrt3a_{K_j}$-neighborhood $N_{a_{K_j}}(f)$ of the
support contribute. Averaging over the three orientations
gives
\begin{equation}\label{IV-24-17}
 |r_j(f)-\int f|
 \le(\sqrt3/2)a_{K_j}\operatorname{Lip}f\,
                                      |N_{a_{K_j}}(f)|\to0.
\end{equation}
Choose one $f_*$ of integral one. Convergence of $\delta_j(f_*)$
and $r_j(f_*)\to1$ gives $c_j\to c=\delta(f_*)$; then
\eqref{IV-24-16} proves $\delta(f)=c\int f$ for every $f$.
The moment formula \eqref{IV-21-10} now identifies $\sigma$
as the shift of $\sigma_S$ by this constant.

For completeness, the best sharp norm of a unit-integral
$C^1$ function supported in the open unit ball is $120M/\pi$.
Indeed if $\lambda=\operatorname{Lip}f$, then
$|f(x)|\le\lambda(1-|x|)_+$ and
$\int(1-|x|)_+\,dx=\pi/3$, so
$\|f\|_\sharp\ge40M\lambda\ge120M/\pi$.
Conversely, mollify $(r-|x|)_+$ inside a ball of radius less
than one and normalize by its integral $\pi r^4/3$.
The resulting functions have sharp norm at most
$120M/(\pi r^4)$; let $r\uparrow1$. Applying
\eqref{IV-22-15} to these functions and passing to the limit
proves \eqref{IV-24-19}.
\end{proof}

\subsection{Reflection positivity in the limit}

For a reflection $\theta$ in a plane and its open positive
half-space, reflection positivity of a functional means
\begin{equation}\label{eq:iv-continuum-rp}
 \sum_{i,k}\overline{z_i}z_k\,
                    \sigma(\theta X_i\cdot X_k)\ge0
\end{equation}
for every finite collection of separated tuples $X_i$ supported
in that half-space. The concatenated tuples here are separated:
the reflected and unreflected supports lie on opposite sides.

\begin{proposition}\label{IV-24-23}\label{IV-24-25}
The centered functional of every growing-torus limit or
volume-first limit is reflection positive at every coordinate
plane, from either side. The same holds for a two-state limit
as in \cref{IV-24-15} if both its integrating and centering
laws are invariant under $a_{K_j}\mathbb Z^3\rtimes W_3$
and its integrating laws are eventually reflection positive
at one coordinate site mirror.

If, in addition, the integrating laws are eventually reflection
positive at the diagonal lattice mirror $x_1=x_2$, then the
limit functional is reflection positive at every plane with
normal $(\pm e_\mu\pm e_\nu)/\sqrt2$, $\mu<\nu$, from either
side.
\end{proposition}
\begin{proof}
At a lattice mirror let $X_i$ be supported strictly on its
positive side. A plaquette centered in this side is contained
in its closure: for a coordinate mirror this follows from the
half-lattice spacing of the centers. For a diagonal mirror,
plaquettes spanning both directions in its normal have center
height in $a\mathbb Z$ and height variation $a$ from center to
corner; those spanning only one have center height in
$a(\mathbb Z+1/2)$ and height variation $a/2$. In either case
a strictly positive center has nonnegative corner heights.
Therefore the
functions $(X_i)_a^\nu$ depend only on positive-half bonds.
Reflection invariance of the centering law gives
$\theta(X_i)_a^\nu=(\theta X_i)_a^\nu$. Lattice reflection
positivity applied to $F=\sum_i z_i(X_i)_a^\nu$ yields
\eqref{eq:iv-continuum-rp} at finite cutoff.

For a growing torus use the coordinate plane
$x_\mu=a_1/2$. This is a bond-midpoint mirror at every cutoff:
$a_1/a_K=L^{K-1}$ is odd. Every fixed positive-side support
fits in its positive torus half-slab for sufficiently large
side, and the exact periodization identity identifies the
finite moments with the lattice expressions. Their limits
therefore satisfy \eqref{eq:iv-continuum-rp} at this plane.
For a volume-first sequence use the site mirror $x_\mu=0$,
where the states are reflection positive by \cref{IV-17-45}.
The same argument applies to the assumed coordinate mirror
in the general two-state case. Translation invariance and
the action of $W_3$ then give all coordinate planes and both
sides.

For the diagonal assertion use
$\theta(x_1,x_2,x_3)=(x_2,x_1,x_3)$, which is a lattice
isometry at every cutoff. The centering laws are invariant
under it. Apply the finite-cutoff argument using the assumed
diagonal positivity of the integrating law and pass to the
limit in the finite sum. Translations and signed coordinate
permutations give all stated diagonal planes.
\end{proof}

The diagonal hypothesis in this proposition is additional.
The cubic-period construction of $\mathcal G^{\rm tor}_{K,\beta}$
has so far supplied only coordinate reflection positivity.
The adapted-period comparison argument below will supply the
diagonal assertion for the continuum functionals to which it
applies.

The same limit argument admits moving mirrors. If
$c_j\in(a_{K_j}/2)\mathbb Z$, $c_j\to c$, and the integrating
laws are reflection positive at $x_\mu=c_j$, use
$\theta_j=\tau_{2(c_j-c)e_\mu}\theta$. Fixed positive-side
supports eventually stay positive for $\theta_j$; their
reflected supports fit in disjoint compact neighborhoods of
their limiting supports. Equation~\eqref{IV-24-33} passes the
finite positive matrix to the limit. More generally, any
inequality involving finitely many such moment evaluations
passes to the limit whenever its admissible values form a
closed set.


\section{Distributional extensions and growth}
\label{sec:distribution-extension}

The bounds on separated tests define distributions away from
coincident points. Their power dependence on the separation also
controls Schwartz functions flat on the diagonals. We prove both
facts by one local Fourier construction. The dyadic partition used
below keeps the supports of the individual Fourier factors separated;
flatness then makes its fine-scale sum converge.

Write $E^n_{\ne}=\{(x_1,\ldots,x_n):x_i\ne x_j\ (i\ne j)\}$,
where $E=\mathbb R^3$ or $T_\ell$, and set
$m(x)=\min_{i\ne j}|x_i-x_j|_\infty$, with $m=\infty$ for $n=1$.
We complexify multilinear forms by complex linearity in each entry.
Thus a bound by $c\prod_i\|f_i\|_\sharp$ on real entries gives
the bound $c\prod_i(\|\operatorname{Re}f_i\|_\sharp+
\|\operatorname{Im}f_i\|_\sharp)$ on complex entries.

\subsection{A partition adapted to the diagonals}

Put $v(t)=e^{-1/t}$ for $t>0$ and $v(t)=0$ for $t\le0$,
$\vartheta(t)=v(1-16t^2)$, and $c=\int\vartheta$.
The two cutoffs
\begin{equation}\label{IV-25-9}
 \chi(t)=c^{-1}\int_{t-1/2}^{t+1/2}\vartheta(s)\,ds,
 \qquad
 \widetilde\chi(t)=c^{-1}\int_{t-1}^{t+1}\vartheta(s)\,ds
\end{equation}
have values in $[0,1]$. The support of $\chi$ is $[-3/4,3/4]$,
$\sum_{z\in\mathbb Z}\chi(t-z)=1$, and at most two summands
are nonzero. The second cutoff equals one on $[-3/4,3/4]$,
vanishes outside $[-5/4,5/4]$, and satisfies
$|\widetilde\chi'|\le16$. We shall use
\begin{equation}\label{IV-25-5}
 |v^{(k)}|\le8^k(k!)^2,\qquad
 |\vartheta^{(k)}|\le200^k(k!)^2,
 \qquad |\chi^{(k)}|\le200^k(k!)^2.
\end{equation}
Here is the derivative estimate. On $t>0$,
$v^{(k)}(t)=P_k(1/t)e^{-1/t}$, where
$P_{k+1}(s)=s^2(P_k(s)-P_k'(s))$ and $P_0=1$.
The weighted coefficient sum $A_k=\sum_j|[s^j]P_k|j!$
satisfies $A_{k+1}\le2(2k+1)^2A_k\le8(k+1)^2A_k$;
also $s^je^{-s}\le j!$. All derivatives tend to zero at
$t=0$, so this proves the first bound globally. For the second,
the chain rule gives
\[
 \vartheta^{(k)}(t)=
 \sum_{0\le2j\le k}\frac{k!}{j!(k-2j)!}
       (-32t)^{k-2j}(-16)^jv^{(k-j)}(1-16t^2).
\]
On $|t|<1/4$, use $|32t|\le8$,
$k!/(j!(k-2j)!)\le3^k j!$, and
$j!((k-j)!)^2\le(k!)^2$. The resulting sum is bounded by
$(32/31)192^k(k!)^2\le200^k(k!)^2$ for $k\ge1$.
Outside this interval it vanishes. Finally
$c\ge\tfrac14e^{-4/3}>1/16$, and differentiating
\eqref{IV-25-9} proves the last bound and the stated properties
of the cutoffs.

Let $h_q=2^{-q}/16$, $q\ge0$. On $\mathbb R^3$ use the grid
$h_q\mathbb Z^3$; on $T_\ell$ use its quotient and periodize
the cutoffs. Here $\ell=2L^m$, $m\ge0$, so $\ell/h_q$ is
an integer and the cutoff supports embed in the torus. Define
\[
 \chi_{q,z}(y)=\prod_{\mu=1}^3\chi(y_\mu/h_q-z_\mu),
 \qquad
 \widetilde\chi_{q,z}(y)=
             \prod_{\mu=1}^3\widetilde\chi(y_\mu/h_q-z_\mu),
\]
with the periodized interpretation on the torus. Let $Q_{q,z}$
and $\widetilde Q_{q,z}$ be the closed cubes of sup radii
$3h_q/4$ and $5h_q/4$ about $h_qz$.
An index tuple $\mathbf z=(z_1,\ldots,z_n)$ is good when
$D(\mathbf z):=\min_{i\ne j}|z_i-z_j|_\infty\ge6$.
Write $\chi_{q,\mathbf z}=\bigotimes_i\chi_{q,z_i}$ and
\begin{equation}\label{IV-25-18}
 G_q=\sum_{\mathbf z\text{ good}}\chi_{q,\mathbf z},
 \qquad G_{-1}=0,
 \qquad \Upsilon_q=G_q(1-G_{q-1}).
\end{equation}
The enlarged cubes in a good tuple have mutual distance at least
$(D(\mathbf z)-5/2)h_q\ge7h_q/2$ and each lies in a ball of
radius less than one. The grid identities give
\begin{gather}
 G_q=0\quad(m\le9h_q/2),\qquad
 G_q=1\quad(m\ge15h_q/2),\label{IV-25-22}\\
 G_qG_{q-1}=G_{q-1},\qquad
 \Upsilon_q=G_q-G_{q-1}\ge0,
 \qquad \sum_{q\ge0}\Upsilon_q=1\quad\text{on }E^n_{\ne}.
 \label{IV-25-23}
\end{gather}
Indeed, a nonzero grid factor has
$|x_i-h_qz_i|_\infty<3h_q/4$, so the distance of two
centres differs from that of their points by less than $3h_q/2$.
This proves the two plateaux. If $G_{q-1}\ne0$, then
$m>9h_q$, which proves the product identity. In particular,
for $q\ge1$,
\[
 \operatorname{supp}\Upsilon_q\subset
       \{9h_q/2\le m\le15h_q\}.
\]
The partition is locally finite on $E^n_{\ne}$.

For a pure derivative in any coordinate, the product rule and
\eqref{IV-25-5} yield
\begin{equation}\label{IV-25-25}
 |\partial_{i\mu}^k
       (\chi_{q,\mathbf z}(1-G_{q-1})F)|
 \le2\cdot600^k(k!)^2
       \max_{0\le j\le k}h_q^{j-k}|\partial_{i\mu}^jF|.
\end{equation}
For the derivatives of $G_q$, only two terms in the
differentiated coordinate can contribute; the sums in every
other coordinate are one. Thus
$|\partial_{i\mu}^kG_q|\le2h_q^{-k}200^k(k!)^2$.
For mixed derivatives of a sum of any subfamily of the grid
tensors, the corresponding bound is
$2^{3n}h_q^{-|\alpha|}200^{|\alpha|}(|\alpha|!)^2$.

One count will supply summability over positions. For $q\ge1$
and each unit box $\mathbf w+[0,1)^{3n}$, there are at most
\begin{equation}\label{IV-25-27}
 \binom n2\,33^3h_q^{-(3n-3)}
\end{equation}
good $\mathbf z$ with $h_q\mathbf z$ in that box and
$Q_{q,\mathbf z}$ meeting $\operatorname{supp}\Upsilon_q$.
At an intersection point some pair has distance at most $15h_q$;
the two centres therefore differ by at most $16h_q$ in each
coordinate. Choose the pair, the other $n-1$ centres, and one
of $33^3$ differences. At $q=0$ the unrestricted count is
$16^{3n}$.

\subsection{The local Fourier construction}

\begin{proposition}[Kernel on separated points]\label{IV-25-55}
Let $T$ be a real multilinear form on separated smooth compactly
supported $n$-tuples over $E$. Suppose that for every $d>0$
there is $c_T(d)<\infty$ such that
\begin{equation}\label{IV-25-31}
 |T(f_1,\ldots,f_n)|\le c_T(d)\prod_i\|f_i\|_\sharp
\end{equation}
whenever the supports have mutual distance at least $d$ and
each lies in a ball of radius one. There is a unique distribution
$\widetilde T$ on $E^n_{\ne}$ such that
$\widetilde T(\bigotimes_i f_i)=T(f_1,\ldots,f_n)$.
It has order at most $4n+1$ on compact sets, with a bound
depending only on that compact set, $M$, and the function $c_T$.
Symmetries of $T$ under permutations or isometries are inherited
by $\widetilde T$.
\end{proposition}
\begin{proof}
Fix a good box at scale $q$ and put $h=h_q$, $\lambda=5h/2$.
Lift its enlarged cubes to cubes $C_i\subset\mathbb R^3$ of
side $\lambda$. For $\omega\in\mathbb Z^3$ let
$e_\omega(t)=e^{2\pi i\omega\cdot t/\lambda}$ and let
$E_{i,\omega}$ be $\widetilde\chi_{q,z_i}e_\omega$, interpreted
through the lift and periodized if necessary. For a smooth $F$
supported in $Q_{q,\mathbf z}$, extend its lift by zero and set
\[
 \widehat F(\omega)=\lambda^{-3n}
       \int_{\prod_iC_i}F(t)\prod_i\overline{e_{\omega_i}(t_i)}\,dt,
 \qquad
 [F]_{q,k}=\max\left\{\|F\|_\infty,
       (\lambda/2\pi)^k\max_{i,\mu}
                          \|\partial_{i\mu}^kF\|_\infty\right\}.
\]
Integration by parts in a coordinate attaining
$|\omega|_\infty$ gives
\[
 |\widehat F(\omega)|\le
       2^k(1+|\omega|_\infty)^{-k}[F]_{q,k}.
\]
The real and imaginary parts of $E_{i,\omega}$ have sharp norm
at most $1109Mh^{-1}(1+|\omega|_\infty)$: the cutoff gradient
is at most $16\sqrt3/h$, and differentiating the exponential
adds at most $16\sqrt3|\omega|_\infty/h$.
Also, grouping by $|\omega|_\infty$ gives
$\sum_{\omega\in\mathbb Z^{3n}}(1+|\omega|_\infty)^{-3n-1}
\le3n8^n$. With $k=4n+1$ these bounds prove absolute
convergence of
\begin{equation}\label{IV-25-41}
 T_{q,\mathbf z}(F)=
 \sum_{\omega\in(\mathbb Z^3)^n}
       \widehat F(\omega)T(E_{1,\omega_1},\ldots,E_{n,\omega_n})
\end{equation}
and the estimate
\begin{equation}\label{IV-25-49}
 |T_{q,\mathbf z}(F)|\le
 (2\cdot10^6)^nM^nc_\square h^{-n}[F]_{q,4n+1}.
\end{equation}
Here $c_\square$ can be either $c_T(7h/2)$ or
$c_T((D(\mathbf z)-5/2)h)$, whenever the argument is finite;
for $n=1$ any valid uniform one-entry bound serves.
The numerical factor before $M^nc_\square h^{-n}[F]$ is
$6n\,283904^n\le(2\cdot10^6)^n$.
Changing a lift multiplies the entries by phases and the Fourier
coefficient by their inverse product; hence \eqref{IV-25-41}
is independent of the lifts.

The Fourier partial sums, multiplied by
$\bigotimes_i\widetilde\chi_{q,z_i}$, converge to $F$ with
all derivatives and have support in $\widetilde Q_{q,\mathbf z}$.
For completeness, integration by parts makes the Fourier series
of every derivative absolutely summable. Its sum is that
derivative: the product Fej\'er kernels are nonnegative approximate
identities, so their convolution converges uniformly to a continuous
periodic function and also to the sum of an absolutely convergent
Fourier series. The cutoff equals one on the support of $F$,
which proves the claimed convergence after multiplication.
If $F=\bigotimes_i f_i$ with
$\operatorname{supp}f_i\subset Q_{q,z_i}$, its coefficients
factor. The partial sums therefore give $T$ of the individual
Fourier approximations; their convergence in the sharp norm and
the telescoping estimate prove
$T_{q,\mathbf z}(F)=T(f_1,\ldots,f_n)$.

Define
\begin{equation}\label{IV-25-56}
 \widetilde T(F)=\sum_{q\ge0}\sum_{\mathbf z\text{ good}}
 T_{q,\mathbf z}\bigl(\chi_{q,\mathbf z}(1-G_{q-1})F\bigr).
\end{equation}
For $F$ in a fixed compact subset of $E^n_{\ne}$, only
finitely many scales and positions occur. Combining
\eqref{IV-25-25} with \eqref{IV-25-49} gives the asserted
distribution bound of order $4n+1$.
For a tensor $F=\bigotimes_i f_i$, each piece in
\eqref{IV-25-56} is the finite sum
\[
 \bigotimes_i(\chi_{q,z_i}f_i)
 -\sum_{\mathbf z'\text{ good at }q-1}
       \bigotimes_i(\chi_{q,z_i}\chi_{q-1,z_i'}f_i),
\]
where the second sum is absent at $q=0$. All entries stay
inside the original separated supports. Multilinearity and
\eqref{IV-25-23} therefore give
$\widetilde T(F)=T(f_1,\ldots,f_n)$. In making this computation,
the linear form induced on finite sums of tensors is well defined:
on the finite-dimensional spans of their factors, point evaluations
span the duals, so a tensor representing the zero function has all
its coefficients zero.

Finally, the same partition and the local Fourier approximation
approximate every test in $C_c^\infty(E^n_{\ne})$ by finite
sums of separated tensors, in all derivatives with supports in
one compact subset of $E^n_{\ne}$. This proves uniqueness.
Applying uniqueness to the pullback by a symmetry proves the
last assertion.
\end{proof}

If the form is originally defined on separated $C_c^1$ tuples
with the same bounds, these values are determined by its kernel.
Indeed, cut each entry by the $q=0$ partition, mollify the
finitely many pieces, and let the mollifier radius tend to zero.
Their supports retain positive separation, their sharp norms
converge, and the telescoping bound recovers the original value.

\subsection{Flat tests and the global bound}

For $\varphi\in\mathcal S(\mathbb R^{3n})$ put
\begin{equation}\label{IV-26-4}
 |\varphi|_k=\sup_{x,\,|\alpha|\le k}
                  (1+|x|^2)^{k/2}|\partial^\alpha\varphi(x)|.
\end{equation}
Let $\mathcal S_{\ne}(\mathbb R^{3n})$ consist of Schwartz
functions all of whose derivatives vanish whenever two entries
coincide. Let $\mathcal S_+$ be the subspace additionally flat
outside $\{x_{i3}>0\text{ for all }i\}$.

\begin{lemma}\label{IV-26-17}
Finite sums of tensors of separated smooth compactly supported
tests are dense in $\mathcal S_{\ne}$ for every norm
$|\cdot|_k$. In $\mathcal S_+$ the approximating entries may
all be supported in the open positive half-space.
For $k\ge4$ and $A\subset\{1,\ldots,n\}$, partial integration
$\psi_A(y)=\int\varphi(y,x_A)\,dx_A$ maps
$\mathcal S_{\ne}(\mathbb R^{3n})$ into
$\mathcal S_{\ne}(\mathbb R^{3(n-|A|)})$ and satisfies
\begin{equation}\label{IV-26-19}
 |\psi_A|_{k(n-|A|)}\le C_{\int}(k)^{|A|}|\varphi|_{kn},
 \qquad C_{\int}(k)=\int_{\mathbb R^3}(1+|x|^2)^{-k/2}\,dx.
\end{equation}
In particular $C_{\int}(4)=\pi^2$ and
$C_{\int}(10)=5\pi^2/64$. Also
$|\varphi\otimes\psi|_k\le|\varphi|_k|\psi|_k$.
If both factors belong to $\mathcal S_{\ne}$ and are flat outside
the negative and positive half-spaces respectively, their tensor
is flat on every diagonal.
\end{lemma}
\begin{proof}
Flatness gives the following Taylor estimate. If $\varphi$ is
flat at $x^*$, $|x-x^*|\le r$ with $r\ge1$, and
$N\ge|\alpha|+s$, then
\begin{equation}\label{IV-26-15}
 |\partial^\alpha\varphi(x)|\le
 \frac{|x-x^*|_1^s}{s!}(1+2r^2)^{N/2}
       |\varphi|_N(1+|x|^2)^{-N/2}.
\end{equation}
Apply Taylor's formula with integral remainder to
$t\mapsto\partial^\alpha\varphi(x^*+t(x-x^*))$; along this
segment $1+|x|^2\le(1+2r^2)(1+|x^*+t(x-x^*)|^2)$.
When $m(x)$ is small, replace a closest pair by its midpoint.
This gives a diagonal point $x^*$ with
$|x-x^*|_1\le3m(x)$ and $|x-x^*|\le\sqrt{3/2}m(x)$.
In the half-space case one may instead replace a smallest third
coordinate by its minimum with zero.

For the half-space partition restrict the good indices further
by $z_{i3}\ge6$, and denote the resulting sum by $G_q^+$.
The plateau and derivative estimates above hold with
$m^+(x)=\min\{m(x),x_{13},\ldots,x_{n3}\}$.
Taylor's estimate with $s=k+1$, followed by the mixed derivative
bound for the cutoffs, gives
\[
 |(1-G_Q)\varphi|_k\le C(n,k)h_Q|\varphi|_{2k+1},
\]
and the same bound with $G_Q^+$ for $\varphi\in\mathcal S_+$.
Indeed, the derivatives of the cutoff cost at most $h_Q^{-k}$,
whereas Taylor's remainder contributes $h_Q^{k+1}$.
Fix $Q$ to make this error small. Among the remaining finitely
many scales, discard grid centres outside a large box. Their
derivatives are bounded with $Q$ fixed, while one extra Schwartz
weight makes their $|\cdot|_k$ contribution tend to zero.
Only finitely many boxes remain. Approximate each by its local
Fourier tensors as in the preceding proof. Their enlarged
supports are separated, and in the half-space case have third
coordinate at least $19h_q/4$. This proves both density assertions.

For partial integrals, use
\[
 (1+|x|^2)^{kn/2}\ge
 (1+|y|^2)^{k(n-|A|)/2}
                    \prod_{i\in A}(1+|x_i|^2)^{k/2}.
\]
This permits differentiation under the integral, proves
\eqref{IV-26-19}, and preserves flatness on the remaining
diagonals. Polar coordinates give the two stated constants.
The tensor bound follows from
$1+|x|^2+|y|^2\le(1+|x|^2)(1+|y|^2)$.
At a diagonal joining the negative and positive groups, one
factor and all its derivatives vanish; diagonals inside a group
are already covered by its flatness.
\end{proof}

\begin{proposition}\label{IV-26-24}
Suppose that $T$ satisfies \eqref{IV-25-31} on $\mathbb R^3$
with $c_T(d)=c_T\max\{1,d^{-a}\}$, where $a\in\mathbb N_0$.
Its kernel has a unique extension to $\mathcal S_{\ne}$
bounded for a Schwartz norm. With $C_{\ne}=10^{53}$,
\begin{equation}\label{IV-26-26}
 |\widetilde T(\varphi)|\le
 C_{\ne}^{n+a}M^n(n!)^8c_T|\varphi|_{8n+a-1}.
\end{equation}
\end{proposition}
\begin{proof}
Use the series \eqref{IV-25-56}, now summed over all scales
and positions. Set $k=4n+1$, $s=a+4n-2$, and $N=k+s$.
For $F_{q,\mathbf z}=\chi_{q,\mathbf z}(1-G_{q-1})\varphi$,
\eqref{IV-25-25} and $\lambda/(2\pi)<2h_q/5$ give
\begin{equation}\label{IV-25-46}
 [F_{q,\mathbf z}]_{q,k}\le
 2\cdot240^k(k!)^2
 \sup_{\operatorname{supp}F_{q,\mathbf z}}
       \max_{i,\mu,\,0\le j\le k}h_q^j|\partial_{i\mu}^j\varphi|.
\end{equation}
For $q\ge1$, this support has $m\le15h_q$. The midpoint
choice in \eqref{IV-26-15} has $|x-x^*|_1\le45h_q$ and
$|x-x^*|<1$. Since $N\ge4n$, it follows that
\[
 |\partial_{i\mu}^j\varphi(x)|\le
 \frac{(45h_q)^s}{s!}3^{N/2}|\varphi|_N
                  \prod_{l=1}^n(1+|x_l|^2)^{-2},\qquad j\le k.
\]
To sum the positions, assign $h_q\mathbf z$ to its unit box
$\mathbf w+[0,1)^{3n}$. Its support lies in
$\mathbf w+[-1/20,21/20]^{3n}$, and
\[
 \sum_{\mathbf w}\sup_{x\in\mathbf w+[-1/20,21/20]^{3n}}
          \prod_l(1+|x_l|^2)^{-2}\le(81\pi^2)^n.
\]
For one factor, compare a point in the enlarged unit box with
a point in its unit cube through the common centre: the squared
weights differ by at most $81$. Integrate over that cube and
sum; then take the product in the $n$ entries.
Combining this estimate, \eqref{IV-25-27},
\eqref{IV-25-49}, and \eqref{IV-25-46} bounds the sum at
scale $q\ge1$ by
\begin{align}
 \Sigma_q\le{}&(2\cdot10^6)^nM^nc_T(2/7)^a
 2\cdot240^k(k!)^2\frac{45^s}{s!}3^{N/2}
 \binom n2\,33^3(81\pi^2)^n
 h_q^{s-a-n-(3n-3)}|\varphi|_N .\label{IV-26-31}
\end{align}
The exponent of $h_q$ is one. Thus the fine-scale sum converges.

At $q=0$, split the centres into $D(\mathbf z)\ge19$ and
$6\le D(\mathbf z)\le18$. In the first group the enlarged
supports are separated by more than one, so $c_\square=c_T$;
use the Schwartz bound without Taylor's formula and the count
$16^{3n}$. In the second group a pair in the box has distance
at most $39/32$ in the sup norm. Its midpoint has
$|x-x^*|_1\le117/32$ and $|x-x^*|\le3/2$; hence
\eqref{IV-26-15} supplies the factor
$(117/32)^s(11/2)^{N/2}/s!$, and
$c_\square=c_T(32/7)^a$.
These give, respectively,
\begin{align*}
 \Sigma_0^{\mathrm A}&\le
 (2\cdot10^6)^nM^nc_T16^{4n}
 2\cdot240^k(k!)^2(81\pi^2)^n|\varphi|_N,\\
 \Sigma_0^{\mathrm B}&\le
 (2\cdot10^6)^nM^nc_T16^{4n}
 2\cdot240^k(k!)^2(81\pi^2)^n
 (32/7)^a\frac{(117/32)^s}{s!}(11/2)^{N/2}|\varphi|_N.
\end{align*}
For $n=1$ the second group and every $q\ge1$ contribution
are empty. In general use
$(k!)^2\le25^n4^{8n}(n!)^8$,
$(2/7)^a3^{N/2}\le81^n$, and $45^s/s!\le3^{45}$
in the fine-scale bound, with $\sum_{q\ge1}h_q=1/16$.
For the second coarse group use
$(32/7)^a(11/2)^{a/2}\le11^a\le11^s$ and
$(11\cdot117/32)^s/s!\le3^{41}$.
The three sums, divided by $M^n(n!)^8c_T|\varphi|_N$,
are at most $10^{28}(10^{28})^n$,
$10^3(10^{30})^n$, and $10^{23}(10^{33})^n$.
For $n\ge2$ their sum is at most $3(10^{45})^n$;
for $n=1$ the sole coarse sum is at most $10^{33}$.
Both are bounded by $C_{\ne}^n$.

The series is therefore absolutely convergent, linear, and
satisfies the stronger bound with $C_{\ne}^n$ in place of
$C_{\ne}^{n+a}$. On compactly supported tests away from the
diagonals it is the previous finite sum. Uniqueness follows from
the density in \cref{IV-26-17}.
\end{proof}

\begin{proposition}[Growth of the Schwinger functions]\label{IV-26-37}
Suppose a symmetric family $S=(S_n)_{n\ge2}$ on separated
$C_c^1$ tuples satisfies
\[
 |S_n(f_1,\ldots,f_n)|\le
 A\max\{1,d^{-a}\}B^n(n!)^\varkappa\prod_i\|f_i\|_\sharp
\]
when the separation is at least $d$ and each support lies in a
ball of radius one, where $a\in\mathbb N$, $A,B\ge1$, and
$\varkappa\ge0$. Its centred moments $\sigma_S$ extend uniquely
to functionals $\mathsf S_n^{\sigma_S}$ on $\mathcal S_{\ne}$,
with $\mathsf S_0^{\sigma_S}=1$, $\mathsf S_1^{\sigma_S}=0$,
and
\begin{equation}\label{IV-26-39}
 |\mathsf S_n^{\sigma_S}(\varphi)|\le
 C_c^n(n!)^{\varkappa+9}|\varphi|_{kn},\qquad
 k=8+\lceil a/2\rceil,\qquad
 C_c=C_{\ne}^{1+a/2}MB\sqrt A.
\end{equation}
If a functional $\sigma$ has these connected functions and
one-point function $c\int f$, with $|c|\le\bar c$, the same
conclusion holds with $C_c$ replaced by
$C=C_c+\bar c\,C_{\int}(k)$.
\end{proposition}
\begin{proof}
In the moment partition sum, there are at most $\lfloor n/2\rfloor$
blocks and no singleton blocks. Also
$\prod_{J\in\pi}|J|!\le n!$, and the number of partitions
is at most $n!$: each partition is the cycle partition of a
permutation, choosing increasing cycles. Thus
\[
 |\sigma_S(f_1,\ldots,f_n)|\le
 A^{\lfloor n/2\rfloor}
 \max\{1,d^{-a\lfloor n/2\rfloor}\}
 B^n(n!)^{\varkappa+1}\prod_i\|f_i\|_\sharp.
\]
Apply \cref{IV-26-24} with singularity order
$a\lfloor n/2\rfloor$. The derivative order is
$8n+a\lfloor n/2\rfloor-1\le kn$, and its bound is
\eqref{IV-26-39}. For $n=1$ the functional is zero.
For an uncentred functional use the finite expansion
\begin{equation}\label{IV-26-40}
 \mathsf S_n^\sigma(\varphi)=
 \sum_{A_0\subset\{1,\ldots,n\}}c^{|A_0|}
             \mathsf S_{n-|A_0|}^{\sigma_S}(\psi_{A_0}).
\end{equation}
It agrees with the original moments on separated tensors.
By \eqref{IV-26-19}, each integrated entry costs at most
$C_{\int}(k)$, so the binomial formula gives $C^n$.
Uniqueness again follows from density.
\end{proof}

For the actual limits of \cref{IV-21-25}, put
\[
 A_g=1+\mathsf U(1,1)+p,\qquad
 C_{E,\mathrm c}^g=C_{\ne}^3M B^{\rm per}(1)\sqrt{A_g}.
\]
Then \eqref{IV-26-39} gives the common bound
\begin{equation}\label{IV-78-10}
 |\mathsf S_n^{\sigma_S}(\varphi)|
 \le(C_{E,\mathrm c}^g)^n(n!)^{14+p}|\varphi|_{10n}.
\end{equation}
Indeed, \eqref{IV-7-4} is nonincreasing in $d$, so
$1+\mathsf U(1,d)\le A_g$ when $d\ge1$.
For $d<1$, put $t=d^{-1}$.  The logarithmic term in
\eqref{IV-7-4} increases by $p\log t$, and the remaining
terms are bounded by $t^3\mathsf U(1,1)$.  Since
$\log t\le t$ and $1,t,t^3\le t^4$,
\[
 A^{\rm per}(d)\le A_g\max(1,d^{-4}).
\]
Apply \cref{IV-26-37} with $a=4$, $B=B^{\rm per}(1)$
and $\varkappa=5+p$, using \eqref{IV-22-17} and its
inherited bound \eqref{IV-22-22}.  This applies to every
growing-torus limit as well as every periodic-state limit;
no additional rotation hypothesis is used.  A constant
one-point function $c\int f$, $|c|\le\bar c$, replaces
$C_{E,\mathrm c}^g$ by
$C_{E,\mathrm c}^g+\bar c C_{\int}(10)$.

\begin{proposition}\label{IV-26-41}
Suppose the moments of a real symmetric functional $\sigma$
extend to continuous linear functionals $\mathsf S_n$ on
$\mathcal S_{\ne}(\mathbb R^{3n})$, agreeing with $\sigma$
on separated smooth tensors and satisfying
\[
 |\mathsf S_n(\varphi)|\le c_E C_E^n(n!)^\kappa
                           |\varphi|_{k_E n}
\]
for fixed $k_E\in\mathbb N$ and $c_E,C_E,\kappa\ge0$.
These extensions preserve reality and symmetry. Every
Euclidean isometry under which $\sigma$ is invariant also
leaves them invariant.
If $\sigma$ is reflection positive, then for
$\varphi_n\in\mathcal S_+(\mathbb R^{3n})$, $0\le n\le n_0$,
with $\mathcal S_+(\mathbb R^0)=\mathbb C$,
\begin{equation}\label{IV-26-42}
 \sum_{n,n'=0}^{n_0}
 \mathsf S_{n+n'}(\Theta\bar\varphi_n\otimes\varphi_{n'})
 \ge0.
\end{equation}
Here $\Theta$ reflects every time coordinate.
\end{proposition}
\begin{proof}
Conjugation and permutation preserve the Schwartz seminorms
and flatness on diagonals. The respective identities hold
on separated smooth tensors and extend by
\cref{IV-26-17}. For an isometry $rx=\rho x+v$, its
pullback preserves the flatness condition and satisfies
\begin{equation}\label{IV-26-43}
 |\varphi\circ(r^{-1})^{\times n}|_k
 \le[6(1+n|v|^2)]^{k/2}|\varphi|_k.
\end{equation}
Indeed each row of $\rho$ has sum of absolute values at
most $\sqrt3$, while
$1+|x|^2\le2(1+n|v|^2)(1+|(r^{-1})^{\times n}x|^2)$.
The chain rule gives the estimate and shows that the
pullback is Schwartz and flat on the diagonals. Density
now transfers the isometry identity.

For positivity, approximate each $\varphi_n$ by finite
sums $\varphi_n^{(j)}$ of separated positive-time smooth
tensors in the seminorm of order $k=2k_{\rm E}n_0$;
leave $\varphi_0$ unchanged. Reflection and conjugation
preserve this seminorm. The tensor estimate and the
flatness assertion in \cref{IV-26-17} give
\[
 |\Theta\bar\varphi_n\otimes\varphi_{n'}
  -\Theta\bar\varphi_n^{(j)}\otimes\varphi_{n'}^{(j)}|_k
 \le |\varphi_n-\varphi_n^{(j)}|_k|\varphi_{n'}|_k
   +|\varphi_n^{(j)}|_k|\varphi_{n'}-\varphi_{n'}^{(j)}|_k
 \longrightarrow0.
\]
Since $k_{\rm E}(n+n')\le k$, continuity permits this
limit in every term of \eqref{IV-26-42}. For each $j$,
expansion of the finite tensor sums expresses the result
as one reflection-positive quadratic form for $\sigma$,
including the empty tuple. Its limit is nonnegative.
\end{proof}

\section{Reflection positivity and the translation spectrum}
\label{sec:reflection-hilbert}

The bounds on separated correlations suffice to construct the
Hilbert space of a reflection. The additional positivity at a
diagonal plane then restricts the joint energy--momentum spectrum.
This construction uses neither coincident-point moments nor a
probability measure on continuum fields.

Let $\sigma$ be a real, symmetric, normalized functional on
$\mathfrak S(\mathbb R^3)$, multilinear on products of test spaces
with disjoint closures. Assume that its one-point function is
$\sigma(f)=c\int f$ and that its connected functions satisfy a
bound table $C_n(d;R)$ as in \cref{sec:separated-compactness}.
Thus, on tuples of separation at least $d$ whose individual
supports lie in balls of radius $R$,
\begin{equation}\label{IV-27-21}
 \sigma(X)=\sum_{A\subset\{1,\ldots,n\}}
       \prod_{i\in A}\left(c\int f_i\right)
       \sum_{\pi\in\operatorname{Part}_{\ge2}(A^c)}
                         \prod_{J\in\pi}S^\sigma_{|J|}(X_J).
\end{equation}
In particular the absolute value is bounded by the same finite
sum with $|c|\|f_i\|_1$ in place of $c\int f_i$ and
$C_{|J|}(d;R)\prod_{i\in J}\|f_i\|_\sharp$ in place of each
connected function. Multilinear telescoping shows that $\sigma$
is continuous when the entries converge in the sharp norm and
their supports remain in fixed separated compact sets.

For a unit vector $u$, put $E_+(u)=\{x:u\cdot x>0\}$,
let $\theta_u x=x-2(u\cdot x)u$, and let
$\mathcal D_u=\mathbb C\mathfrak S(E_+(u))$ be the free
complex vector space on the separated tuples in this half-space.
Concatenation is used only when the resulting tuple is separated.
The product $\theta_u\bar F\cdot G$ is always defined for
$F,G\in\mathcal D_u$. Reflection positivity means
\[
 \sigma(\theta_u\bar F\cdot F)\ge0\qquad(F\in\mathcal D_u).
\]
It makes
\begin{equation}\label{IV-27-19}
 \langle F,G\rangle_{\sigma,u}
       =\sigma(\theta_u\bar F\cdot G)
\end{equation}
a positive semidefinite Hermitian form, antilinear in its first
argument. Quotient by its null space and complete to obtain
$\mathcal H(\sigma,u)$. Write $[F]$ for a class and
$\Omega=[\varnothing]$, so that $\|\Omega\|=1$.

\begin{proposition}[The translation operators]\label{IV-27-22}
Suppose that $\sigma$ is translation invariant and reflection
positive at $u\cdot x=0$. Then $\mathcal H(\sigma,u)$ is
separable, and the maps
\[
 T_t[F]=[\tau_{tu}F]\quad(t\ge0),\qquad
 U(v)[F]=[\tau_vF]\quad(v\in u^\perp)
\]
extend to a strongly continuous semigroup of positive
self-adjoint contractions and a strongly continuous unitary
group. They commute and fix $\Omega$. There are strongly
commuting self-adjoint operators $H\ge0,P_1,P_2$ such that
\begin{equation}\label{IV-27-13}
 T_t=e^{-tH},\qquad U(v)=e^{i(v_1P_1+v_2P_2)},\qquad
 T_tU(v)=\int e^{-tE+iv\cdot q}\,d\mathcal P(E,q),
\end{equation}
where coordinates are taken in any fixed orthonormal basis of
$u^\perp$, and $\mathcal P$ is carried by
$[0,\infty)\times\mathbb R^2$. Moreover $\ker T_t=0$,
$H\Omega=P_1\Omega=P_2\Omega=0$, and
$T_t\to P_{\ker H}$ strongly as $t\to\infty$.

If an orthogonal transformation $r$ fixes $u$ and leaves
$\sigma$ invariant, $[F]\mapsto[rF]$ extends to a unitary
$U(r)$ fixing $\Omega$. It commutes with $H$ and satisfies
$U(r)U(v)U(r)^*=U(rv)$. The joint spectrum is invariant under
$(E,q)\mapsto(E,rq)$.
\end{proposition}

\begin{proof}
On $\mathcal D_u$ write $T_t^0=\tau_{tu}$ and
$U^0(v)=\tau_v$. The identities
$\theta_u\tau_{tu}=\tau_{-tu}\theta_u$ and
$\theta_u\tau_v=\tau_v\theta_u$, together with translation
invariance, give
\[
 \langle F,T_t^0G\rangle=\langle T_t^0F,G\rangle,
 \qquad
 \langle U^0(v)F,U^0(v)G\rangle=\langle F,G\rangle.
\]
For each $F=\sum_i z_iX_i$ the number
\[
 c_F=\sup_{t\ge0}|\langle F,T_t^0F\rangle|
\]
is finite. Indeed, choose common $d,R>0$ for the finitely many
tuples $\theta_uX_i\cdot X_j$. Moving their positive entries
by $tu$ does not decrease a distance across the mirror:
\[
 |y+tu-x|^2=|y-x|^2+2t\,u\cdot(y-x)+t^2\ge|y-x|^2
 \quad(u\cdot x<0<u\cdot y).
\]
The sharp and $L^1$ norms and the radii of individual supports
are unchanged. The bound following \eqref{IV-27-21}, summed
with coefficients $|z_i z_j|$, therefore bounds $c_F$.

Cauchy--Schwarz for the semidefinite form and the semigroup
law give the decisive contraction estimate
\begin{equation}\label{IV-27-12}
 \|T_t^0F\|^2=\langle F,T_{2t}^0F\rangle
       \le\|F\|\,\|T_{2t}^0F\|.
\end{equation}
After $j$ iterations,
\[
 \|T_t^0F\|\le
   \|F\|^{1-2^{-j}}\|T_{2^jt}^0F\|^{2^{-j}}
   \le\|F\|^{1-2^{-j}}c_F^{2^{-j-1}}.
\]
If $\|F\|=0$, the first inequality already gives zero.
Otherwise letting $j\to\infty$ proves
$\|T_t^0F\|\le\|F\|$. Thus these maps preserve the null
space and extend as contractions. The maps $U^0(v)$ are
invertible isometries and extend as unitaries. The algebraic
identities extend by density; symmetry gives self-adjointness,
and $T_t=T_{t/2}^2$ gives positivity.

For small translations, enlarge the finitely many supports by
less than $d/4$. They remain separated by $d/2$. Translations
of a compactly supported $C^1$ function converge in the sharp
norm, by uniform continuity of the function and its first
derivatives. Hence the continuity following
\eqref{IV-27-21} gives
$\langle F,T_t^0F\rangle\to\|F\|^2$ and
$\langle F,U^0(v)F\rangle\to\|F\|^2$ at the origin.
The identities
\[
 \|(T_t-1)[F]\|^2
 =\langle F,T_{2t}^0F\rangle
       -2\operatorname{Re}\langle F,T_t^0F\rangle+\|F\|^2,
 \qquad
 \|(U(v)-1)[F]\|^2
 =2\|F\|^2-2\operatorname{Re}\langle F,U^0(v)F\rangle
\]
prove strong continuity on the dense set of classes; the
contraction bounds extend it to the completion. If $T_t\psi=0$,
then $\|T_{t/2}\psi\|^2=\langle\psi,T_t\psi\rangle=0$.
Iteration and continuity at zero give $\psi=0$.

The spectral measure of $T_1$ is carried by $(0,1]$.
Define $H=-\log T_1$ by spectral calculus. The uniqueness of
positive square roots gives $T_{2^{-j}}=T_1^{2^{-j}}$;
the semigroup law and strong continuity then give $T_t=e^{-tH}$
for all $t\ge0$. This also proves uniqueness of $H$.
Every $U(v)$ commutes with $T_1$ and its spectral projections,
so $e^{isH}U(v)$ is a strongly continuous unitary representation
of $\mathbb R\times u^\perp$. Its joint spectral measure is
$\mathcal P$ in \eqref{IV-27-13}; its coordinate generators
are $H,P_1,P_2$. The invariance of $\Omega$ implies that its
spectral measure is concentrated at $(0,0)$. Dominated convergence
of $e^{-tE}$ gives the asserted limit of $T_t$.

For separability use the countable family of separated test
tuples constructed in \cref{sec:separated-compactness}.
Given $X\in\mathfrak S(E_+(u))$, choose $d>0$ smaller than
all separations in $\theta_uX\cdot X$. Rational neighborhoods
of its positive supports of width less than $d/3$ still lie
in $E_+(u)$, since $u\cdot x\ge d/2$ on those supports.
Approximate each entry in $C^1$ inside these neighborhoods.
The four terms in
\[
 \|[X]-[Y]\|^2
 =\sigma(\theta_uX\cdot X)-\sigma(\theta_uX\cdot Y)
  -\sigma(\theta_uY\cdot X)+\sigma(\theta_uY\cdot Y)
\]
then converge to the same value. Rational complex linear
combinations of the resulting countable set of classes are
dense. Finally $r$ commutes with $\theta_u$ and preserves
the form. Its unitary extension has the asserted covariance
on the dense set of classes, hence everywhere. Spectral
uniqueness gives the covariance of $\mathcal P$.
\end{proof}

\subsection{Changing the one-point constant}

For real $b$ define, on a separated tuple $X=(f_1,\ldots,f_n)$,
\begin{equation}\label{IV-28-13}
 \mathsf{sh}_bX=
   \sum_{A\subset\{1,\ldots,n\}}\prod_{i\in A}
                         \left(b\int f_i\right)X_{A^c}.
\end{equation}
The map extends linearly, preserves every open-set tuple space,
commutes with isometries and conjugation, and respects every
defined concatenation. Expanding the subsets gives
$\mathsf{sh}_b\mathsf{sh}_{b'}=\mathsf{sh}_{b+b'}$.
Thus it is invertible with inverse $\mathsf{sh}_{-b}$.
The moment--cumulant partition formula shows that
$\sigma' =\sigma\circ\mathsf{sh}_b$ has one-point constant
$c+b$ and the same connected functions of order at least two.

\begin{proposition}[Recentering and reflection]\label{IV-28-12}
The functionals $\sigma$ and $\sigma'$ have the same isometry
invariances and the same reflection positivities. They cluster
simultaneously. In the setting of \cref{IV-27-22}, the map
\[
 [F]_{\sigma',u}\longmapsto
                   [\mathsf{sh}_bF]_{\sigma,u}
\]
extends to a unitary which fixes the vacuum and intertwines
the translation operators and their joint spectral measures.
\end{proposition}
\begin{proof}
The algebraic identities just established imply
\begin{equation}\label{IV-28-15}
 \sigma'(\theta_u\bar F\cdot G)
 =\sigma(\theta_u\overline{\mathsf{sh}_bF}
                                  \cdot\mathsf{sh}_bG).
\end{equation}
They prove reflection positivity in both directions, using
$\mathsf{sh}_{-b}$ for the converse, and give a surjective
isometry on the quotients. The same commutation identities
intertwine translations and isometries. Uniqueness of the
spectral measure proves the last assertion. Invariance follows
also directly from the subset expansion. Each shifted clustering
difference is a finite sum of unshifted clustering differences
of subtuples, with constant integral factors; the converse uses
the inverse shift.
\end{proof}

\subsection{Positivity at a diagonal plane}

The following measure estimate is the analytic step in comparing
the Hilbert spaces of two reflections.

\begin{lemma}\label{IV-28-17}
Let $\mu$ be a finite positive measure on
$[0,\infty)\times\mathbb R$, and set $w=(E-iq)/\sqrt2$.
Suppose that, on a Hilbert space $\mathcal K$, there are
vectors $a,b$ and self-adjoint operators $H_+\ge0,Q$ with
\begin{equation}\label{IV-28-18}
 \int e^{-sw-s'\bar w}\,d\mu(E,q)
       =\langle a,e^{-sH_+}e^{-is'Q}b\rangle
       \qquad(s,s'\ge0).
\end{equation}
Then $\mu(\{|q|>E\})=0$. No commutation of $H_+$ and $Q$
is required.
\end{lemma}
\begin{proof}
For $c\in L^1(0,\infty)$ put
$p(z)=\int_0^\infty c(s)e^{-sz}\,ds$ on $\operatorname{Re}z\ge0$.
Multiply \eqref{IV-28-18} by $c(s)\overline{c(s')}$ and
integrate. Fubini is justified by the bound
$|e^{-sw-s'\bar w}|\le1$, the finiteness of $\mu$, and the
contraction and unitary bounds on the right. Spectral calculus
in the indicated order gives
\begin{equation}\label{IV-28-19}
 \int|p(w)|^2\,d\mu
 =\langle a,p(H_+)p^\sharp(Q)b\rangle
 \le\|a\|\|b\|\sup_{x\ge0}|p(x)|
                         \sup_{y\in\mathbb R}|p(iy)|,
 \qquad p^\sharp(y)=\overline{p(-iy)}.
\end{equation}
Take $z_0=\varepsilon+i\kappa$, where $\varepsilon>0$ and
$\kappa\in\mathbb R$, and
$c(s)=\Gamma(z_0)^{-1}s^{z_0-1}e^{-\eta s}$ with $\eta>0$.
The Gamma integral gives $p(z)=(z+\eta)^{-z_0}$, with the
principal logarithm. Consequently
\[
 \int |w+\eta|^{-2\varepsilon}
         e^{2\kappa\arg(w+\eta)}\,d\mu
 \le\|a\|\|b\|\eta^{-2\varepsilon}e^{|\kappa|\pi/2}.
\]
Fatou's lemma first as $\varepsilon\downarrow0$ and then
as $\eta\downarrow0$ yields
\begin{equation}\label{IV-28-20}
 \int_{w\ne0}e^{2\kappa\arg w}\,d\mu
                      \le\|a\|\|b\|e^{|\kappa|\pi/2}.
\end{equation}
For every $\delta>0$, the measure of
$\{\arg w\ge\pi/4+\delta\}$ is at most
$\|a\|\|b\|e^{-2\kappa\delta}$ for all $\kappa>0$,
and hence is zero. Taking $\kappa<0$ proves the corresponding
assertion for $\arg w\le-\pi/4-\delta$. Their countable
union is exactly $\{|q|>E\}$.
\end{proof}

\begin{proposition}[The spectral bound from two reflections]
\label{IV-28-21}
Let $\sigma$ satisfy the bounds and translation invariance of
\cref{IV-27-22}. Fix $k\in\{1,2\}$ and $\epsilon\in\{1,-1\}$.
If $\sigma$ is reflection positive for the normals $e_3$ and
$e_+=(e_3+\epsilon e_k)/\sqrt2$, then its spectral measure
in $\mathcal H(\sigma,e_3)$ satisfies
\[
 \mathcal P\{(E,q):E<|q_k|\}=0.
\]
We denote this spectral-support assertion by $H\ge|P_k|$.
\end{proposition}
\begin{proof}
Put $e_-=(e_3-\epsilon e_k)/\sqrt2$ and
$W=\{x:x_3>|x_k|\}$. For $F\in\mathbb C\mathfrak S(W)$,
both $F$ and $\theta_{e_+}\theta_{e_3}F$ lie in the positive
half-space for $e_+$. Let $\mu$ be the image of the positive
spectral measure $\langle[F],\mathcal P(\cdot)[F]\rangle$
under $(E,q)\mapsto(E,\epsilon q_k)$. The same correlation
has two representations:
\begin{equation}\label{IV-28-23}
 \begin{split}
 \int e^{-sw-s'\bar w}\,d\mu
 &=\sigma\bigl(\theta_{e_3}\bar F\cdot
                         \tau_{se_++s'e_-}F\bigr)\\
 &=\bigl\langle[\theta_{e_+}\theta_{e_3}F]_+,
                  T_s^+U^+(s'e_-)[F]_+\bigr\rangle_+.
 \end{split}
\end{equation}
Indeed $se_++s'e_-=(s+s')e_3/\sqrt2+
\epsilon(s-s')e_k/\sqrt2$, which gives the first equality
by \eqref{IV-27-13}; the second is the defining reflection
form in $\mathcal H(\sigma,e_+)$. There
$T_s^+=e^{-sH_+}$ and $U^+(s'e_-)=e^{-is'Q}$ for the
negative of the momentum generator along $e_-$. Hence
\cref{IV-28-17} shows that the spectral projection of
$\{E<|q_k|\}$ annihilates $[F]$.

For arbitrary $F\in\mathcal D_{e_3}$, a sufficiently large
positive time translation puts all its finitely many supports
in $W$. The same projection therefore annihilates $T_r[F]$.
It commutes with $T_r$, whose kernel is zero by
\cref{IV-27-22}, and so annihilates $[F]$. Density proves
that the projection is zero.
\end{proof}

\section{Slab operators and analytic translations}
\label{sec:slab-operators}

We now isolate the hypotheses used to recover rotations from
the cubic symmetries. The insertion estimate constructs bounded
operators across a slab. Together with the spectral bound, these
operators will give holomorphic correlation kernels. A separate
two-point growth estimate will exclude the nonconstant angular
harmonics.

Throughout this section, complex test functions are expanded
multilinearly into real and imaginary parts. On them we use
$\|h\|_\sharp=\max\{\|h\|_\infty,40M\operatorname{Lip}h\}$;
each of the real and imaginary parts has norm at most
$\|h\|_\sharp$. Such multilinear identities hold also for
Hilbert-space classes: pairing their difference with every
$[G]$ reduces the assertion to the multilinearity of $\sigma$,
and density makes the difference zero.

Write $y=(x_3,x_1)$, leaving $x_2$ as the third coordinate,
and let $R_t$ rotate this plane by
$(y^0,y^1)\mapsto(y^0\cos t-y^1\sin t,y^0\sin t+y^1\cos t)$.
Put $n_\upsilon=(\cos\upsilon,\sin\upsilon)$, so
$n_0=(1,0)$ corresponds to $e_3$.
A plate is $f(x)=\alpha(y)\varphi(x_2)$, where
$\varphi\in C_c^1(\mathbb R)$ is real and
$\alpha\in C_c^1(\mathbb R^2;\mathbb C)$. Its label is
$\varphi$; denote the length of the smallest interval containing
$\operatorname{supp}\varphi$ by $\lambda_\varphi$, with
$\lambda_0=0$. Set
\begin{equation}\label{IV-32-20}
 \operatorname{Prof}_t=
 \{\alpha\in C_c^1(\mathbb R^2):
   \operatorname{supp}\alpha\subset[t,2t]\times[-t,t],\
   |\alpha|\le3,\ \operatorname{Lip}\alpha\le12/t\}.
\end{equation}

\begin{definition}[Hypotheses for the rotation argument]
\label{IV-32-25}
For a real normalized functional $\sigma$ on separated tuples,
the hypotheses are as follows.
\begin{enumerate}[label=(R\arabic*),start=0]
\item $\sigma(f)=0$ for every $f\in C_c^1(\mathbb R^3)$.
\item $\sigma$ is symmetric and invariant under translations and
the cubic group $W_3$ of signed coordinate permutations.
\item For each $n\ge1$ and $d,R>0$ there is a finite
$C_n^\sigma(d;R)$, non-increasing in $d$, such that
\[
 |\sigma(f_1,\ldots,f_n)|\le C_n^\sigma(d;R)
                                      \prod_i\|f_i\|_\sharp
\]
for every tuple with separation at least $d$ and each support
in some ball of radius $R$; the centers are arbitrary.
\item $\sigma$ is reflection positive at $x_3=0$.
\item In the corresponding Hilbert space, $H\ge|P_1|$ in the
sense of spectral support.
\item There are $0<c_{\rm I}\le1/16$, $C_{\rm I}>0$ and
$p_{\rm I}\ge0$ such that, for $0<\eta\le1$, every real $h$
supported in $B(z,c_{\rm I}\eta)$ with $|z_3|\le\eta/8$, and
$F,G\in\mathbb C\mathfrak S(\{x_3>\eta/4\})$,
\begin{equation}\label{IV-32-27}
 |\sigma(\theta\bar F\cdot(h)\cdot G)|
 \le C_{\rm I}\eta^{-p_{\rm I}}\|h\|_\sharp
       \|[\tau_{-(\eta/4)e_3}F]\|
       \|[\tau_{-(\eta/4)e_3}G]\|.
\end{equation}
\item For the diagonal reflection with normal
$(e_1+e_3)/\sqrt2$, the two-point form is positive:
\[
 \sum_{a,b}\bar z_a z_b\sigma(\theta_\Sigma f_a,f_b)\ge0
 \quad\bigl(f_a\in C_c^1(\{x_1+x_3>0\})\bigr).
\]
\item There is an exponent $a_\Gamma<4$ such that for every
label $\varphi$ there is $C_\varphi<\infty$ with
\begin{equation}\label{IV-32-29}
 \sigma(\theta f,f)\le C_\varphi t^{-a_\Gamma}
 \quad(0<t\le1,\ f=\alpha\otimes\varphi,
                         \ \alpha\in\operatorname{Prof}_t).
\end{equation}
\end{enumerate}
\end{definition}

The Hilbert construction of \cref{IV-27-22} applies under
(R1)--(R3). Its proof uses only the uniform moment bound in
(R2); alternatively the partition formula bounds each connected
function by a finite sum of products of those bounds. Cubic
symmetry transports the construction to each coordinate normal.
The map $[F]^{(\mu)}\mapsto[wF]^{(\nu)}$ is unitary whenever
$we_\mu=e_\nu$, since it preserves the reflection form.
It intertwines the time and tangential translations. Thus (R4)
gives, for every ordered pair $\mu\ne\nu$,
\begin{equation}\label{IV-32-33}
 H^{(\mu)}\ge|P_\nu^{(\mu)}|.
\end{equation}

\subsection{The holomorphic translation semigroup}

For $z\in\mathbb C^3$ and $\xi\in\mathbb C^2$ put
\begin{equation}\label{IV-32-24}
 d_3(z)=\operatorname{Re}z_3-|\operatorname{Im}z_1|
                              -|\operatorname{Im}z_2|,
 \qquad \mathsf{Tb}_3=\{d_3>0\},
 \qquad z^*=(-\bar z_1,-\bar z_2,\bar z_3),
\end{equation}
and
\begin{equation}\label{IV-32-22}
 d_{\mathsf{Tb}}(\xi)=\operatorname{Re}\xi^0
                                -|\operatorname{Im}\xi^1|,
 \quad\mathsf{Tb}=\{d_{\mathsf{Tb}}>0\},
 \quad\xi^*=(\bar\xi^0,-\bar\xi^1),
 \quad\xi^\pm=\xi^0\pm i\xi^1.
\end{equation}
Complex rotations are defined by
$(R_\zeta\xi)^\pm=e^{\pm i\zeta}\xi^\pm$.
Both tubes are convex cones, stable under their respective
involutions. The map $\xi\mapsto(\xi^1,0,\xi^0)$ embeds
$\mathsf{Tb}$ in $\mathsf{Tb}_3$.

\begin{lemma}\label{IV-32-34}
Under (R1)--(R4),
\begin{equation}\label{IV-32-35}
 \mathsf T_3(z)=\int e^{-z_3E+i(z_1q_1+z_2q_2)}\,d\mathcal P(E,q)
 \quad(z\in\overline{\mathsf{Tb}_3})
\end{equation}
is a contraction semigroup, strongly continuous on the closed
tube and holomorphic in operator norm on its interior.
It satisfies $\mathsf T_3(z)^*=\mathsf T_3(z^*)$,
$\mathsf T_3(z)\Omega=\Omega$, and
$\mathsf T_3(x)[F]=[\tau_xF]$ for real $x$ with $x_3\ge0$.
Its restriction is denoted
$\mathsf T(\xi)=\mathsf T_3(\xi^1,0,\xi^0)$.
\end{lemma}
\begin{proof}
On the support of $\mathcal P$, \eqref{IV-32-33} gives
$|q_1|,|q_2|\le E$, so the integrand has modulus at most
$e^{-d_3(z)E}\le1$. Spectral calculus gives the semigroup,
adjoint, vacuum and real-translation identities. Dominated
convergence in each scalar spectral measure proves strong
continuity. On a compact subset of the open tube there is
$\delta>0$ with $d_3\ge\delta$. Every derivative of degree
$m$ in the integrand is bounded by $E^m e^{-\delta E}$.
Taylor's formula on a smaller neighborhood, with this uniform
bound, permits differentiation in operator norm. This proves
holomorphy.
\end{proof}

\subsection{Insertion across a slab}

\begin{proposition}[Slab operators]\label{IV-32-38}
Under (R1)--(R3) and (R5), for every $l>0$ and
$h\in C_c^1(\{0<x_3<l\};\mathbb C)$ there is a unique bounded
operator $B_l(h)$ such that
\begin{equation}\label{IV-32-39}
 B_l(h)[F]=[(h)\cdot\tau_{le_3}F].
\end{equation}
It is complex linear in $h$. For every compact subset $K$ of
the slab there is $C(K,l)<\infty$ with
$\|B_l(h)\|\le C(K,l)\|h\|_\sharp$ when
$\operatorname{supp}h\subset K$. It has the identities
\begin{align}
 B_l(h)&=T_cB_{c'-c}(\tau_{-ce_3}h)T_{l-c'},
 &&\operatorname{supp}h\subset\{c<x_3<c'\},\label{IV-32-41}\\
 B_l(h)^*&=B_l(\overline{\theta_{l/2}h}),
 & U(v)B_l(h)U(-v)&=B_l(\tau_vh),\label{IV-32-42}
\end{align}
where $0\le c<c'\le l$ and $v\perp e_3$.
For real $h$ supported in $B(z,c_{\rm I}\eta)$ with
$\eta/8\le z_3\le3\eta/8$, $0<\eta\le1$,
\begin{equation}\label{IV-32-44}
 \|B_{\eta/2}(h)\|\le C_{\rm I}\eta^{-p_{\rm I}}
                                             \|h\|_\sharp.
\end{equation}
If $0=c_0<\cdots<c_n$ and
$\operatorname{supp}h_j\subset\{c_{j-1}<x_3<c_j\}$, then
\begin{equation}\label{IV-32-45}
 [(h_1,\ldots,h_n)]
 =\prod_{j=1}^n B_{c_j-c_{j-1}}
                   (\tau_{-c_{j-1}e_3}h_j)\Omega,
\end{equation}
with increasing $j$ from left to right.
\end{proposition}
\begin{proof}
For a small piece as in \eqref{IV-32-44}, put
$h_0=\tau_{-(\eta/4)e_3}h$. Apply (R5) to $h_0$ and
$\tau_{(\eta/4)e_3}F,\tau_{(\eta/4)e_3}G$. Translation
invariance changes the resulting form into
\[
 \sigma(\theta\bar F\cdot(h)\cdot\tau_{(\eta/2)e_3}G),
\]
whose modulus is at most
$C_{\rm I}\eta^{-p_{\rm I}}\|h\|_\sharp\|[F]\|\|[G]\|$.
It descends to the quotient and extends continuously to the
Hilbert space. The bounded-form representation theorem gives
an operator with this bound. Pairing with every class identifies
its value on $[G]$ with $[(h)\cdot\tau_{(\eta/2)e_3}G]$.

For general compact $K$ in the slab, set
$\delta=\min_{x\in K}\min(x_3,l-x_3)>0$,
$\eta=\min(1,4\delta)$ and $b=c_{\rm I}\eta/2$.
Use the $C^1$ partition
\[
 \chi_{b,k}(x)=\prod_{\mu=1}^3\chi(x_\mu/b-k_\mu),
 \quad
 \chi(s)=\begin{cases}\cos^2(\pi s/2),&|s|\le1,\\0,&|s|>1.
 \end{cases}
\]
The identity $\chi(s)+\chi(s-1)=1$ for $0\le s\le1$
gives $\sum_k\chi_{b,k}=1$. Each support is a cube of
radius $b$, and $\operatorname{Lip}\chi_{b,k}\le\sqrt3\pi/(2b)$.
Only finitely many of these cubes meet $K$. Write
$h_k=\chi_{b,k}h$ and let $c_k$ be the projection of
$bk_3-\eta/4$ onto $[0,l-\eta/2]$. Because
$b\le\eta/32$ and $\eta/4\le\delta$, the center of
$\tau_{-c_ke_3}h_k$ has third coordinate between
$7\eta/32$ and $9\eta/32$. Its support lies in a ball
of radius $\sqrt3 b<c_{\rm I}\eta$. The small-piece
construction therefore applies, and
\[
 B_l(h)=\sum_k T_{c_k}B_{\eta/2}(\tau_{-c_ke_3}h_k)
                                      T_{l-c_k-\eta/2}
\]
has the required value on classes. The product estimate
\[
 \|h_k\|_\sharp\le
 \left(1+\frac{40\sqrt3\pi M}{c_{\rm I}\eta}\right)
                                               \|h\|_\sharp
\]
gives the compact-support bound, with an extra factor two for
the real and imaginary parts. The construction is independent
of the partition because its values on the dense set of classes
are prescribed.

Translation of the tuple in \eqref{IV-32-39} proves
\eqref{IV-32-41} and tangential covariance. Reflection in
$x_3=l/2$ interchanges its left and right entries; symmetry
and conjugation give the adjoint formula. Finally, iterating
\eqref{IV-32-39} through the ordered slabs gives
\eqref{IV-32-45}.
\end{proof}

\subsection{Removing coincident planar projections}

\begin{proposition}\label{IV-32-48}
Assume (R1)--(R4). Fix $n\ge2$ and a real angle $t$.
If $\sigma(R_tX)=\sigma(X)$ for every $n$-tuple whose
supports have pairwise disjoint projections onto the
$(x_3,x_1)$ plane, the identity holds for all separated
$n$-tuples.
\end{proposition}
\begin{proof}
First suppose that two groups of entries lie respectively
in $x_2<c-\delta$ and $x_2>c+\delta$, where $\delta>0$.
In the Hilbert space with normal $e_2$, their correlation
after a relative translation $su$, $u\perp e_2$, $|u|=1$,
is a matrix element of $T_{2\delta}^{(2)}U^{(2)}(su)$.
For its finite complex spectral measure $m$ this is
\[
 \Phi(s)=\int e^{-2\delta E+is(u_1q_1+u_3q_3)}\,dm(E,q).
\]
Since $|q_1|,|q_3|\le E$ and $|u_1|+|u_3|\le\sqrt2$,
this is holomorphic for $|\operatorname{Im}s|<\sqrt2\delta$.
Thus the relative planar translation is analytic whenever
there is such an $x_2$ gap.

Let $A_j$ be the planar projections of the supports and $I_j$
their $x_2$ interval hulls. Suppose first that every pair of
intersecting $A_j$ has a level separating its two intervals
and lying outside all the $I_j$. Induct on the number of
intersecting pairs. Choose a separating level for one pair,
and translate all entries above it by $se_1$. For all sufficiently
large $s$, the planar projections on opposite sides no longer
meet, and the number of intersecting pairs strictly decreases.
The interval hulls and their separating levels are unchanged.
The induction hypothesis gives rotation invariance for these
large $s$. Both correlations are holomorphic on a common strip:
after rotation, the relative translation is $sR_te_1$, still
perpendicular to $e_2$. The identity theorem gives equality
also at $s=0$.

For a general tuple of separation $d>0$, choose
$0<\varepsilon<d/n$ and partition each entry in the $x_2$
coordinate with $\chi(2x_2/\varepsilon-k)$.
Multilinearity gives a finite sum of tuples whose interval
hulls have length at most $\varepsilon$. If two planar
projections of a nonzero summand meet, points over their common
projection are at $x_2$ distance at least $d$. The gap between
their two interval hulls has length at least
$d-2\varepsilon>(n-2)\varepsilon$. The other $n-2$ interval
hulls cannot cover this gap, so it contains a level outside all
hulls. The preceding induction applies to each summand.
The partition commutes with $R_t$, and summing proves the
assertion. Tuples with a zero entry vanish by multilinearity.
\end{proof}

\section{The two-point rotation argument}
\label{sec:two-point-rotation}

For a plate label $\varphi$, reflection positivity gives a
positive spectral measure on the energy--momentum half-plane.
Its transform represents the two-point function. The plate
growth bound leaves only the angular frequencies $0,4,-4$.
Positivity at the coordinate and diagonal planes gives opposite
signs for the coefficient at frequency four.

\subsection{The spectral measure of a label}

Put $Y=[0,\infty)\times\mathbb R$, $V=\{(E,q):E\ge|q|\}$,
and, for a profile supported in $y^0>0$, set
\begin{equation}\label{IV-33-9}
 \widetilde\alpha(E,q)=\int\alpha(y)e^{-y^0E+iy^1q}\,dy.
\end{equation}
This transform is continuous and bounded by
$\|\alpha\|_1e^{-\varepsilon E}$ if the support has
$y^0\ge\varepsilon>0$. Convolution of profiles multiplies
their transforms.

We first work with any symmetric bilinear two-point form $s$
on separated supports which is translation invariant, satisfies
the order-two bound in (R2), and is reflection positive at
$y^0=0$. The form $s(\theta\bar f,h)$ on
$C_c^1(\{x_3>0\};\mathbb C)$ defines a one-field Hilbert
space. The proof of \cref{IV-27-22}, with single test functions
instead of tuples and with no distinguished vacuum, gives
positive time translations, tangential unitaries and a joint
spectral measure $\mathcal P^s$ on $Y$.

\begin{proposition}\label{IV-33-21}
For every real label $\varphi$ there is a unique positive
Borel measure $\rho[s,\varphi]$ on $Y$, finite on compact
sets, such that all profile transforms are in $L^2(\rho)$ and
\begin{equation}\label{IV-33-22}
 s\bigl(\theta(\bar\alpha\otimes\varphi),
                       \alpha'\otimes\varphi\bigr)
       =\int\overline{\widetilde\alpha}\,
                         \widetilde\alpha'\,d\rho.
\end{equation}
More precisely, for every Borel set $B\subset Y$,
\begin{equation}\label{IV-33-23}
 \langle[\alpha\otimes\varphi],\mathcal P^s(B)
                              [\alpha'\otimes\varphi]\rangle
 =\int_B\overline{\widetilde\alpha}\,
                         \widetilde\alpha'\,d\rho.
\end{equation}
\end{proposition}
\begin{proof}
Write $\Psi_\alpha=[\alpha\otimes\varphi]$ and
$\mu_{\alpha,\alpha'}(B)=
\langle\Psi_\alpha,\mathcal P^s(B)\Psi_{\alpha'}\rangle$.
Riemann sums for convolution converge in $C^1$, with supports
in a fixed compact subset of the positive half-plane. The
two-point bound therefore gives
\begin{equation}\label{IV-33-24}
 \Psi_{\alpha'*\alpha}
       =\widetilde\alpha'(\mathcal P^s)\Psi_\alpha.
\end{equation}
On the right, the same Riemann sums converge weakly by the
spectral representation of translations and dominated
convergence; this identifies their norm limit on the left.
Commutativity of convolution implies
\begin{equation}\label{IV-33-25}
 |\widetilde\alpha'|^2\mu_{\alpha,\alpha}
       =|\widetilde\alpha|^2\mu_{\alpha',\alpha'}.
\end{equation}

On $O_\alpha=\{\widetilde\alpha\ne0\}$ define
$d\rho_\alpha=|\widetilde\alpha|^{-2}
d\mu_{\alpha,\alpha}$. Equation \eqref{IV-33-25} makes
these local measures agree on overlaps. The sets $O_\alpha$
cover $Y$: at $(E_0,q_0)$ take
$\alpha(y)=\chi(y^0-2)\chi(y^1)e^{-iy^1q_0}$, whose
transform there is strictly positive. Compact exhaustion of
$Y$ gives a countable subcover. Partitioning it into disjoint
Borel pieces defines a positive measure $\rho$ with the
prescribed restrictions. It is locally finite, since near a
point one transform is bounded away from zero.

Equation \eqref{IV-33-25} also shows that
$\mu_{\alpha,\alpha}$ vanishes on the zeros of
$\widetilde\alpha$: intersect that zero set with each member
of the countable cover and use the nonzero transform there.
Hence $d\mu_{\alpha,\alpha}=|\widetilde\alpha|^2d\rho$
on all of $Y$. Polarization proves \eqref{IV-33-23}, and
taking $B=Y$ gives \eqref{IV-33-22}.

For uniqueness, a finite complex measure $m$ on $Y$ is
determined by its Laplace--Fourier transform. Indeed, vanishing
of $\int e^{-tE+ivq}dm$ at all integers $t\ge0$ and all
$v\in\mathbb R$ first gives vanishing against
$p(e^{-E})e^{ivq}$ for polynomials $p$. Uniform polynomial
approximation on $[0,1]$, followed by Fourier uniqueness in
$q$, gives vanishing against $g(e^{-E})\mathbf1_B(q)$ for
continuous $g$ and Borel $B$. Approximation of indicators of
closed subsets of $[0,\infty)$ and the monotone class theorem
then give $m=0$. Thus vanishing against every
$\widetilde\alpha$ also implies $m=0$, by testing the
continuous Laplace--Fourier transform against profiles.
If two proposed measures satisfy \eqref{IV-33-22}, replace
$\alpha'$ by $\alpha''*\alpha$. The preceding uniqueness
applies to the finite measure
$|\widetilde\alpha|^2(\rho'-\rho'')$ and makes it zero.
Division on $O_\alpha$ and the countable cover give
$\rho'=\rho''$.
\end{proof}

Apply this to $s=\sigma|_{\mathfrak S_2}$ under (R1)--(R4)
and write $\rho_\varphi=\rho[s,\varphi]$. The one-field
space embeds isometrically into $\mathcal H(\sigma)$ through
$[f]\mapsto[(f)]$. It intertwines translations, so
$|\widetilde\alpha|^2\rho_\varphi$ is the image of the
full spectral measure of $[(\alpha\otimes\varphi)]$ under
$(E,q_1,q_2)\mapsto(E,q_1)$. Consequently $\rho_\varphi$
is carried by $V$. Reflection of $x_1$ and uniqueness in
\cref{IV-33-21} make it invariant under $q\mapsto-q$.

\subsection{From plate bounds to a kernel bound}

Let $a_0(s)=1+\cos(2\pi(s-3/2))$ on $[1,2]$, zero
elsewhere, and let
$b_1=(3/2)\mathbf1_{[-1/3,1/3]}$, $b=b_1*b_1*b_1$.
For $1/2\le\lambda\le1$ define
$\alpha_{t,\lambda}(y)=a_0(y^0/t)b(y^1/(\lambda t))$.
These profiles belong to $\operatorname{Prof}_t$:
$\int a_0=\int b=1$, $0\le a_0\le2$, $|a_0'|\le2\pi$,
$0\le b\le9/8$, and $|b'|\le9/4$. The last bounds follow
by writing $b=(3/2)\int_{u-1/3}^{u+1/3}T(v)dv$ with
$T(v)=(9/4)(2/3-|v|)_+$.

There is an absolute $C_b<\infty$ such that every positive
locally finite measure $\rho$ carried by $V$ satisfies
\begin{equation}\label{IV-33-40}
 \Gamma[\rho](5t):=\int e^{-5tE}\,d\rho
 \le C_b t^{-4}\sup_{1/2\le\lambda\le1}
                    \int|\widetilde\alpha_{t,\lambda}|^2d\rho.
\end{equation}
Here the inequality is valid also for infinite integrals.
To see it, the transform of $b$ is
$\widehat b(q)=(\sin(q/3)/(q/3))^3$, and
$\widehat a_0(s)\ge e^{-2s}$ for $s\ge0$. Thus
\[
 |\widetilde\alpha_{t,\lambda}(E,q)|^2
      \ge(t^4/4)e^{-4tE}\widehat b(\lambda tq)^2.
\]
Let $c_6>0$ be the infimum of the means of $\sin^6$ over
intervals of length at least $1/2$. It is positive by
periodicity, compactness for lengths between $1/2$ and $\pi$,
and counting full periods for larger lengths. Set
$c_b=\min(\sin^61,3^6c_6)$. For $|s|\le3$, concavity of
sine gives $\widehat b(\lambda s)^2\ge\sin^61$; for
$|s|\ge3$, integration over $\lambda$ gives the mean of
$\sin^6$ on an interval of length $|s|/6$. Therefore
\[
 \int_{1/2}^1 2\widehat b(\lambda s)^2d\lambda
          \ge c_b(1+|s|)^{-6}.
\]
On $V$ this is at least $c_b6^{-6}e^5e^{-tE}$ when
$s=tq$, since $(1+x)^{-6}\ge6^{-6}e^5e^{-x}$.
Tonelli proves \eqref{IV-33-40} with
$C_b=4\cdot6^6e^{-5}/c_b$.

For $\rho_\varphi$, the supremum in \eqref{IV-33-40} is
finite at each fixed $t$ by (R2), so $\Gamma_\varphi(t)$
is finite for all $t>0$. Under (R7), it also gives
\begin{equation}\label{IV-33-54}
 \Gamma_\varphi(t)\le
 C'_\varphi\max(1,t^{-1})^{4+a_\Gamma},
 \qquad C'_\varphi=5^{4+a_\Gamma}C_bC_\varphi.
\end{equation}
Define
\begin{equation}\label{IV-33-29}
 k_\varphi(\xi)=\int_V e^{-\xi^0E+i\xi^1q}\,d\rho_\varphi.
\end{equation}
It is holomorphic on $\mathsf{Tb}$ and bounded by
$\Gamma_\varphi(d_{\mathsf{Tb}}(\xi))$. The integrable
majorant follows from $|q|\le E$; on smaller tubes it also
bounds all derivatives. It is even in $\xi^1$ and obeys
$\overline{k_\varphi(\xi)}=k_\varphi(\xi^*)$.

\begin{proposition}[The two-point kernels]\label{IV-33-47}
Under (R1)--(R4), $k_\varphi$ extends by quarter turns to a
real analytic function $\widetilde k_\varphi$ on
$\mathbb R^2\setminus\{0\}$. For disjointly supported profiles,
\begin{equation}\label{IV-33-50}
 \sigma(\alpha\otimes\varphi,\alpha'\otimes\varphi)
   =\iint\alpha(y)\alpha'(y')
                         \widetilde k_\varphi(y'-y)\,dy\,dy'.
\end{equation}
For two labels set
$k_{\varphi,\psi}=(k_{\varphi+\psi}-k_{\varphi-\psi})/4$,
and likewise for $\widetilde k$. Then
\begin{equation}\label{IV-33-51}
 \sigma(\alpha\otimes\varphi,\alpha'\otimes\psi)
 =\sigma(\alpha\otimes\psi,\alpha'\otimes\varphi)
 =\iint\alpha(y)\alpha'(y')
                    \widetilde k_{\varphi,\psi}(y'-y)\,dy\,dy'.
\qedhere
\end{equation}
\end{proposition}
\begin{proof}
If the profiles lie on opposite sides of $y^0=c$, translate
to $c=0$ and apply \eqref{IV-33-22}. Fubini gives
\eqref{IV-33-50} with $k_\varphi$ on the positive half-plane.
For $v^0,v^1>0$, take profiles in small neighborhoods of
$0,v$ and rotate both through $-\pi/2$. Cubic invariance
equates their correlations, and the integral representation
gives $k_\varphi(v)=k_\varphi(R_{-\pi/2}v)$ by shrinking
the neighborhoods. This equality makes the quarter-turned
definitions agree on overlaps. They cover the punctured plane
and give $\widetilde k_\varphi$. Reality follows from the
evenness in $q$.

For arbitrary disjoint profile supports of distance $d>0$,
partition each into squares of diameter less than $d/4$ using
the functions $\chi$ above. Any pair of nonzero pieces is
ordered along one of the signed coordinate axes, since one
coordinate of a vector between them has modulus at least
$d/\sqrt2$. The ordered formula applies after a quarter turn,
and its sum is \eqref{IV-33-50}.

Polarization in the labels gives the sum of the two mixed
correlations in \eqref{IV-33-51}. To see that they are equal,
use an even approximate identity $j_\delta$ on the plane.
Translation invariance and Riemann sums give
\[
 \sigma((\alpha*j_\delta)\otimes\varphi,
                      (\alpha'*j_\delta)\otimes\psi)
 =\iint\alpha(u)\alpha'(u')h_\delta(u'-u)\,du\,du',
\]
where $h_\delta(v)=\sigma(j_\delta\otimes\varphi,
(\tau_vj_\delta)\otimes\psi)$ and $\delta$ is smaller
than a quarter of the separation. Apply the same identity
to $\alpha'(-u),\alpha(-u)$ and change both integration
variables to their negatives. Its right side is unchanged.
Rotation by $\pi$ and permutation of the two entries turn
its left side into the expression with the labels exchanged.
As $\delta\downarrow0$, convolution converges in $C^1$ on
fixed separated neighborhoods. The order-two bound permits
the limit and proves the equality of the mixed correlations.
\end{proof}

\subsection{Eliminating the fourth harmonic}

\begin{proposition}\label{IV-33-53}
Under (R1)--(R4) and (R7), for each $r>0$ the function
$k_{\varphi,r}(\zeta)=k_\varphi(R_\zeta(r,0))$, initially
defined for $|\operatorname{Re}\zeta|<\pi/2$, extends to
an entire function with period $\pi/2$. There are real
$c_0(r),c_1(r)$ such that
\begin{equation}\label{IV-33-55}
 k_{\varphi,r}(\zeta)=c_0(r)+2c_1(r)\cos4\zeta,
 \qquad
 \widetilde k_\varphi(r\cos t,r\sin t)
                      =c_0(r)+2c_1(r)\cos4t.
\end{equation}
\end{proposition}
\begin{proof}
For $\zeta=t+is$ in the initial strip,
$d_{\mathsf{Tb}}(R_\zeta(r,0))=re^{-|s|}\cos t$.
The quarter-turn identity of \cref{IV-33-47} agrees on
overlapping strips by holomorphic uniqueness, so translates
by $j\pi/2$ give one entire periodic function. Move $t$
modulo $\pi/2$ into $[-\pi/4,\pi/4]$. Then
\[
 |k_{\varphi,r}(t+is)|\le
       \Gamma_\varphi(re^{-|s|}/\sqrt2)
       \le C_r e^{A|s|},\qquad A=\max(0,4+a_\Gamma)<8.
\]
For an entire $\pi/2$-periodic function with this bound,
shifting the contour of its Fourier coefficient from the
real axis to height $s$ shows that the coefficient of
$e^{4im\zeta}$ vanishes if $4|m|>A$. Only $m=0,1,-1$
remain. Evenness in $\xi^1$ equates the last two coefficients,
and reality on real points makes both remaining coefficients
real. This proves the formulas.
\end{proof}

\begin{proposition}[Rotation invariance at two plates]\label{IV-33-57}
Under (R1)--(R4), (R6) and (R7), $c_1(r)=0$ for all
labels and all $r>0$. Thus $\widetilde k_{\varphi,\psi}(v)$
depends only on $|v|$, and correlations of two plates with
disjoint planar supports are invariant under every $R_t$.
Moreover, for $v^0>0$ and real $s$,
\begin{equation}\label{IV-33-59}
 k_{\varphi,\psi}(R_{is}v)=k_{\varphi,\psi}(v).
\end{equation}
\end{proposition}
\begin{proof}
For imaginary angles the integral in \eqref{IV-33-29} is
\[
 k_{\varphi,r}(is)
 =\int_V e^{-r(E\cosh s+q\sinh s)}\,d\rho_\varphi\ge0.
\]
Dividing $c_0(r)+2c_1(r)\cosh4s$ by $\cosh4s$ and
letting $s\to\infty$ gives $c_1(r)\ge0$.

For the opposite sign, turn the two-point form through
$\pi/4$: $s^{\rm d}(f,h)=\sigma(R_{\pi/4}f,R_{\pi/4}h)$.
It satisfies translation invariance and the same local bound;
(R6) gives its coordinate reflection positivity. Hence
\cref{IV-33-21} constructs a positive measure $\rho_{\rm d}$
on $Y$. At this point no cone support for $\rho_{\rm d}$
is assumed.

Let $b_s(E,q)=(E\cosh s+q\sinh s,q\cosh s+E\sinh s)$.
It preserves $V$, has inverse $b_{-s}$, and its first
coordinate is at least $e^{-|s|}E$ on $V$. Thus
$(b_s)_*\rho_\varphi$ has finite Laplace transforms, and
its kernel at real $v^0>0$ is $k_\varphi(R_{is}v)$.
The three harmonics in \eqref{IV-33-55} give
\[
 k_\varphi(R_iv)+k_\varphi(R_{-i}v)
 = (\cosh4+1)k_\varphi(v)
       -(\cosh4-1)\widetilde k_\varphi(R_{\pi/4}v).
\]
Multiply by $\overline{\alpha(u)}\alpha'(u')$ at
$v=u'-\theta u$ and integrate. The last term represents
the turned form by \eqref{IV-33-50}; the other terms are
the spectral representations for the boosted measures.
Uniqueness in \cref{IV-33-21} therefore gives the identity
of positive measures
\begin{equation}\label{IV-33-63}
 (\cosh4-1)\rho_{\rm d}
       +(b_1)_*\rho_\varphi+(b_{-1})_*\rho_\varphi
                   =(\cosh4+1)\rho_\varphi.
\end{equation}
In particular
$0\le\rho_{\rm d}\le(\cosh4+1)(\cosh4-1)^{-1}\rho_\varphi$.
This proves both its cone support and finiteness of its
Laplace transforms. Its kernel equals
$\widetilde k_\varphi(R_{\pi/4}v)$ on real $v^0>0$,
by the integral identity against profiles. Analytic
continuation in the angle thus gives
\[
 k[\rho_{\rm d}](R_{is}(r,0))
                   =c_0(r)-2c_1(r)\cosh4s\ge0.
\]
The same limit proves $c_1(r)\le0$. Hence it vanishes.
Polarization in $\varphi\pm\psi$ gives the radiality of
the mixed kernel. Formula \eqref{IV-33-51} proves invariance
of the two-plate correlations. Holomorphic uniqueness in the
initial strip gives \eqref{IV-33-59}.
\end{proof}

\section{Point fields and plate kernels}
\label{sec:point-kernels}

The insertion estimate first gives operators smeared inside a slab.
Holomorphy of the translation semigroup improves this to a pointwise
operator bound. The essential step is a norm Cauchy estimate for
mollified insertions; weak compactness alone would not give the
products needed below. Throughout this section we assume (R1)--(R5).

Write $B_d^{\rm c}(h)=B_{2d}(\tau_{de_3}h)$ and put
\begin{equation}\label{IV-34-6}
 \wp=p_{\rm I}+4,\qquad
 C_{\rm plt}=480\cdot8^3m_2c_{\rm I}^{-4}C_{\rm I}M,
 \qquad C_{\rm pt}=320\cdot8^4m_3c_{\rm I}^{-4}C_{\rm I}M.
\end{equation}
Here the absolute constants are obtained from the radial weight
\[
 m(s)=\begin{cases}
 c_m,&0\le s\le1/3,\\
 c_m\cos^2(3\pi(s-1/3)/4),&1/3\le s\le1,\\
 0,&s\ge1,
 \end{cases}\qquad c_m=\frac{9\pi}{5\pi^2-8},
\]
by setting
$m_0=c_m$, $m_1=3\pi c_m/4$,
$m_2=\sqrt2m_0m_1+2m_0^2$ and
$m_3=m_0^3+\sqrt3m_0^2m_1$.
The function $m$ is $C^1$, nonnegative and constant near zero;
direct integration gives $2\pi\int_0^1m(s)s\,ds=1$.
Let $j_\delta^{(N)}$, $N=2,3$, be $m(2|x|/\delta)$ divided
by its integral. These are nonnegative radial $C^1$ mollifiers
supported in $|x|\le\delta/2$.

For $r>0$, the product weight
\[
 \alpha_{r,t}^{(N)}(x)=\prod_{i=1}^N
 r^{-2}m\bigl(r^{-1}\sqrt{x_i^2+t_i^2}\bigr)
\]
has supremum at most $m_0^Nr^{-2N}$ and Lipschitz constant at
most $\sqrt N m_0^{N-1}m_1r^{-2N-1}$. Iterating the one-variable
mean-value formula in polar coordinates gives
\begin{equation}\label{IV-34-16}
 u(0)=\int_{(-r,r)^N}\!dt\int\!dx\,
             \alpha_{r,t}^{(N)}(x)u(x+it)
\end{equation}
whenever $u$ is holomorphic on a neighborhood of the closed
polydisc of radius $r$. All operator integrals below converge in
norm; their integrands are continuous with compact support.

\subsection{Removing the smearing}

\begin{proposition}[Point plates]\label{IV-34-17}
For every real label $\varphi\in C_c^1(\mathbb R)$ there is a
norm-holomorphic family $B_\varphi(\xi,\xi')$, linear in
$\varphi$, on $\mathsf{Tb}\times\mathsf{Tb}$ such that
\begin{equation}\label{IV-34-18}
 \|B_\varphi(\xi,\xi')\|\le
 C_{\rm plt}(\lambda_\varphi+1)\|\varphi\|_{C^1}
 \min\{1,d_{\mathsf{Tb}}(\xi),d_{\mathsf{Tb}}(\xi')\}^{-\wp}.
\end{equation}
For $\kappa,\kappa'\in\overline{\mathsf{Tb}}$ it satisfies
\begin{align}
 \mathsf T(\kappa)B_\varphi(\xi,\xi')\mathsf T(\kappa')
 &=B_\varphi(\kappa+\xi,\xi'+\kappa'),\label{IV-34-19}\\
 B_\varphi(\xi,\xi')^*&=B_\varphi(\xi'^*,\xi^*).
 \label{IV-34-20}
\end{align}
If $\operatorname{supp}\alpha\subset\{|y^0|<d\}$, then
\begin{equation}\label{IV-34-21}
 B_d^{\rm c}(\alpha\otimes\varphi)
 =\int\alpha(y)B_\varphi(dn_0+y,dn_0-y)\,dy.
\end{equation}
These properties determine the family already among weakly
continuous families satisfying the shift law and the integral
identity for real profiles.
\end{proposition}
\begin{proof}
Put $j_\delta=j_\delta^{(2)}$ and
\[
 A_\delta=B_\delta^{\rm c}(j_\delta\otimes\varphi),\qquad
 B^\delta(\xi,\xi')=
 \mathsf T(\xi-\delta n_0)A_\delta
 \mathsf T(\xi'-\delta n_0).
\]
This is defined when both tube distances are at least $\delta$,
and holomorphic when they are larger. Radiality of the mollifier
and the slab adjoint formula give $A_\delta^*=A_\delta$.
The shift and adjoint laws follow at once. Slab localization and
tangential covariance give
\[
 B^\delta(dn_0+y,dn_0-y)
       =B_d^{\rm c}((\tau_yj_\delta)\otimes\varphi).
\]
Consequently, for continuous $\alpha$ supported in
$|y^0|<d-\delta$,
\begin{equation}\label{IV-34-25}
 \int\alpha(y)B^\delta(dn_0+y,dn_0-y)\,dy
        =B_d^{\rm c}((\alpha*j_\delta)\otimes\varphi).
\end{equation}
Indeed, convolution Riemann sums converge in $C^1$ with one
compact support, and the slab bound permits passage to the limit.

We first bound $B^\delta$ independently of $\delta$. Set
$\eta=\min(1,d_{\mathsf{Tb}}(\xi),d_{\mathsf{Tb}}(\xi'))$,
$r=c_{\rm I}\eta/8$ and $0<\delta\le r$. For unit vectors
$v,v'$, apply \eqref{IV-34-16} to
$u(z)=\langle v,B^\delta(\xi+z,\xi'-z)v'\rangle$.
The polydisc of radius $2r$ lies in its domain because
$\eta\ge128r$. Move the imaginary shift into the vectors by
the two contractions with arguments
$\xi-\eta n_0/4+it$ and $\xi'-\eta n_0/4-it$.
The remaining inner integral is a matrix element of
\[
 B_{\eta/4}^{\rm c}(\gamma_t\otimes\varphi),
 \qquad \gamma_t=\alpha_{r,t}^{(2)}*j_\delta.
\]
Its profile is supported in $[-2r,2r]^2$, with supremum
$m_0^2r^{-4}$ and Lipschitz constant
$\sqrt2m_0m_1r^{-5}$.

Partition the label as
$\varphi_k(s)=\chi(s/(4r)-k)\varphi(s)$, using the partition
in the proof of \cref{IV-32-38}. At most
$3(\lambda_\varphi+1)/(c_{\rm I}\eta)$ terms are nonzero.
Each $\gamma_t\otimes\varphi_k$ is supported in
$B(4rke_2,c_{\rm I}\eta)$ and has sharp norm at most
$40Mm_2r^{-5}\|\varphi\|_{C^1}$.
The centered version of \eqref{IV-32-44}, followed by integration
over $(-r,r)^2$, therefore gives
\[
 \|B^\delta(\xi,\xi')\|\le
 (2r)^2\frac{3(\lambda_\varphi+1)}{c_{\rm I}\eta}
 C_{\rm I}\eta^{-p_{\rm I}}40Mm_2r^{-5}\|\varphi\|_{C^1}.
\]
This is \eqref{IV-34-18} with $B^\delta$ in place of
$B_\varphi$.

To obtain norm convergence, fix $\epsilon>0$ and set
\[
 \delta_\epsilon=\frac{c_{\rm I}}{16}\min(1,\epsilon),\qquad
 b_\epsilon=C_{\rm plt}(\lambda_\varphi+1)\|\varphi\|_{C^1}
                  (\tfrac12\min(1,\epsilon))^{-\wp}.
\]
For $\delta\le\delta_\epsilon$, the holomorphic map
$z\mapsto B^\delta(\epsilon n_0+z,\epsilon n_0-z)$ is
bounded by $b_\epsilon$ on the polydisc of radius
$\epsilon/4$. Cauchy's derivative estimate gives the bound
$8b_\epsilon|y|/\epsilon$ for its change from $z=0$ to a
real $y$ with $|y|\le\epsilon/8$. Integrate against
$j_{\delta'}$ and use \eqref{IV-34-25}; interchange
$\delta,\delta'$ and use commutativity of convolution. The result is
\[
 \|B^\delta(\epsilon n_0,\epsilon n_0)
       -B^{\delta'}(\epsilon n_0,\epsilon n_0)\|
 \le\frac{4b_\epsilon}{\epsilon}(\delta+\delta').
\]
Thus $B_{[\epsilon]}=\lim_{\delta\downarrow0}
B^\delta(\epsilon n_0,\epsilon n_0)$ exists in norm.
For $0<\epsilon\le\min(d_{\mathsf{Tb}}(\xi),d_{\mathsf{Tb}}(\xi'))$,
define
\begin{equation}\label{IV-34-27}
 B_\varphi(\xi,\xi')=
 \mathsf T(\xi-\epsilon n_0)B_{[\epsilon]}
 \mathsf T(\xi'-\epsilon n_0).
\end{equation}
This is the limit of $B^\delta(\xi,\xi')$, uniformly when
both tube distances exceed $\epsilon$, so it does not depend
on the chosen $\epsilon$. Choosing a smaller $\epsilon$ locally
proves norm holomorphy directly from this formula.
The bound, shifts, adjoints and linearity pass to the limit.
In \eqref{IV-34-25}, uniform norm convergence on the support
and $\alpha*j_\delta\to\alpha$ in $C^1$ prove
\eqref{IV-34-21}.

For uniqueness, the integral identity against all real profiles
and weak continuity identify the matrix elements at
$(dn_0+y,dn_0-y)$, in particular at $(dn_0,dn_0)$.
The shift law then identifies every pair of tube points.
\end{proof}

\begin{proposition}[Point fields]\label{IV-34-30}
There is a unique weakly continuous family
$B(z,z')$, $(z,z')\in\mathsf{Tb}_3^2$, satisfying the
three-dimensional shift law and the following slab integral
formula. It is norm holomorphic and obeys
\begin{align}
 \|B(z,z')\|&\le C_{\rm pt}
 \min\{1,d_3(z),d_3(z')\}^{-\wp},
 \label{IV-34-31}\\
 \mathsf T_3(\kappa)B(z,z')\mathsf T_3(\kappa')
 &=B(\kappa+z,z'+\kappa'),\label{IV-34-32}\\
 B(z,z')^*&=B(z'^*,z^*),\label{IV-34-33}\\
 B_d^{\rm c}(h)&=\int h(x)B(de_3+x,de_3-x)\,dx.
 \label{IV-34-34}
\end{align}
Here $\kappa,\kappa'\in\overline{\mathsf{Tb}_3}$, and in
the last identity $h\in C_c^1(\{|x_3|<d\};\mathbb C)$.
\end{proposition}
\begin{proof}
Use $j_\delta^{(3)}$ in the preceding construction, with
$A_\delta=B_\delta^{\rm c}(j_\delta^{(3)})$ and
$\mathsf T_3$ in place of $\mathsf T$. The convolution
identity and the shift and adjoint laws have the same proofs.
For completeness, the change in the uniform estimate is as follows.
Put $\eta=\min(1,d_3(z),d_3(z'))$
and $r=c_{\rm I}\eta/8$. The polydisc of radius $2r$
loses at most $6r$ in tube distance. The three-variable
mean-value formula, after moving the imaginary shifts into
contractions, gives the centered insertion of
$\gamma_t=\alpha_{r,t}^{(3)}*j_\delta^{(3)}$.
Its support lies in $[-2r,2r]^3\subset B(0,c_{\rm I}\eta)$,
and $\|\gamma_t\|_\sharp\le40Mm_3r^{-7}$.
No label partition is needed. Hence the bound is
\[
 (2r)^3C_{\rm I}\eta^{-p_{\rm I}}40Mm_3r^{-7}
       =C_{\rm pt}\eta^{-\wp}.
\]
At $(\epsilon e_3,\epsilon e_3)$ use a polydisc of radius
$\epsilon/6$. With
$b_\epsilon=C_{\rm pt}(\tfrac12\min(1,\epsilon))^{-\wp}$,
Cauchy's estimate and the symmetric convolution argument give
\[
 \|B^\delta(\epsilon e_3,\epsilon e_3)
        -B^{\delta'}(\epsilon e_3,\epsilon e_3)\|
 \le6b_\epsilon(\delta+\delta')/\epsilon.
\]
Norm convergence, extension to arbitrary tube pairs, holomorphy,
the integral formula and uniqueness now follow exactly by the
limit and shift arguments above.
\end{proof}

\begin{corollary}\label{IV-34-42}
For $0\le d'<d$ and $\operatorname{supp}h\subset\{|x_3|\le d'\}$,
\begin{equation}\label{IV-34-43}
 \|B_d^{\rm c}(h)\|\le
 C_{\rm pt}\min(1,d-d')^{-\wp}\|h\|_{L^1}.
\end{equation}
If $F,F'\in\mathbb C\mathfrak S(\{x_3>d\})$, the function
\begin{equation}\label{IV-34-44}
 v_d(x)=\langle[\tau_{-de_3}F],
 B(de_3+x,de_3-x)[\tau_{-de_3}F']\rangle
\end{equation}
is real analytic for $|x_3|<d$ and satisfies
\begin{align}
 \sigma(\theta\bar F\cdot(h)\cdot F')
       &=\int h(x)v_d(x)\,dx,\label{IV-34-45}\\
 |v_d(x)|&\le C_{\rm pt}\min(1,d-|x_3|)^{-\wp}
       \|[\tau_{-de_3}F]\|\|[\tau_{-de_3}F']\|.
       \label{IV-34-46}
\end{align}
Moreover, identifying $\xi=(\xi^0,\xi^1)$ with
$(\xi^1,0,\xi^0)$,
\begin{equation}\label{IV-34-47}
 B_\varphi(\xi,\xi')
       =\int\varphi(s)B(\xi+se_2,\xi'-se_2)\,ds.
\end{equation}
\end{corollary}
\begin{proof}
Integrate \eqref{IV-34-31} to obtain \eqref{IV-34-43}.
The slab form gives \eqref{IV-34-45}; its density is the
restriction of a holomorphic matrix element on
$|\operatorname{Re}z_3|+|\operatorname{Im}z_1|
+|\operatorname{Im}z_2|<d$, and the same point bound gives
\eqref{IV-34-46}. The right side of \eqref{IV-34-47} is
weakly continuous, satisfies the plate shift law, and, by
Fubini and \eqref{IV-34-34}, represents the slab operator of
$\alpha\otimes\varphi$. Uniqueness in \cref{IV-34-17}
proves the identity.
\end{proof}

\subsection{Products and their analytic continuation}

At an interior junction of two point plates, only the sum of
the two adjoining arguments matters. More precisely, for
$\kappa_1,\kappa_2\in\mathsf{Tb}$,
$\kappa_0\in\overline{\mathsf{Tb}}$ and
$0<\epsilon\le\min(d_{\mathsf{Tb}}(\kappa_1),
d_{\mathsf{Tb}}(\kappa_2))$,
\begin{equation}\label{IV-34-50}
 \begin{split}
 B_\varphi(\xi,\kappa_1)\mathsf T(\kappa_0)
 B_\psi(\kappa_2,\xi')
 =B_\varphi(\xi,\epsilon n_0)
 \mathsf T(\kappa_1+\kappa_0+\kappa_2-2\epsilon n_0)
 B_\psi(\epsilon n_0,\xi').
 \end{split}
\end{equation}
This follows by extracting the two semigroup factors with
\eqref{IV-34-19}. Thus one may replace the adjoining arguments
by half their sum. An outer argument next to $\Omega$ may be
changed arbitrarily, since $\mathsf T(\kappa)\Omega=\Omega$.
The same statements hold for point fields and $\mathsf T_3$.

For $\xi_1,\ldots,\xi_{n-1}\in\mathsf{Tb}$, define
\begin{equation}\label{IV-34-56}
 K_n(\xi;\underline\varphi)=
 \left\langle\Omega,
 B_{\varphi_1}(\kappa,\xi_1/2)
 B_{\varphi_2}(\xi_1/2,\xi_2/2)\cdots
 B_{\varphi_n}(\xi_{n-1}/2,\kappa)\Omega\right\rangle.
\end{equation}
Here $\kappa\in\mathsf{Tb}$ is arbitrary; the preceding outer
argument identity makes the answer independent of it. Set
$K_0=1$ and $K_1(\varphi)=
\langle\Omega,B_\varphi(\kappa,\kappa)\Omega\rangle$.
For
\[
 \mathsf{Tb}^{(n)}=\{x:x_1\in\mathsf{Tb},\
                   x_{j+1}-x_j\in\mathsf{Tb}\},
\]
let $\Psi(x;\underline\varphi)$ be the operator product in
\eqref{IV-34-56}, applied to $\Omega$, with successive
differences $\xi_j=x_{j+1}-x_j$ and first outer argument
$x_1$ in place of $\kappa$; set $\Psi(\varnothing)=\Omega$.

\begin{proposition}\label{IV-34-59}
The kernels $K_n$ are holomorphic, and $\Psi$ is norm holomorphic
on $\mathsf{Tb}^{(n)}$. They are multilinear in the labels, and for $n\ge2$,
\begin{equation}\label{IV-34-60}
 |K_n(\xi;\underline\varphi)|\le
 \prod_{a=1}^n\bigl(2^\wp C_{\rm plt}
       (\lambda_{\varphi_a}+1)\|\varphi_a\|_{C^1}\bigr)
 \min\{1,\min_a d_{\mathsf{Tb}}(\xi_a)\}^{-\wp n}.
\end{equation}
For profiles whose supports are strictly ordered in time,
\begin{equation}\label{IV-34-61}
 \sigma(\alpha_1\otimes\varphi_1,\ldots,\alpha_n\otimes\varphi_n)
 =\int\prod_a\alpha_a(y_a)
 K_n(y_2-y_1,\ldots,y_n-y_{n-1};\underline\varphi)\,dY.
\end{equation}
If the first support also lies in $y^0>0$, the class of this
tuple is
\begin{equation}\label{IV-34-62}
 [(\alpha_1\otimes\varphi_1,\ldots,\alpha_n\otimes\varphi_n)]
       =\int\prod_a\alpha_a(y_a)\Psi(Y;\underline\varphi)\,dY.
\end{equation}
The vectors satisfy
\begin{equation}\label{IV-34-64}
 \begin{split}
 \langle\Psi(x;\underline\varphi),\Psi(z;\underline\psi)\rangle
 =K_{n+n'}\bigl(&(x_n-x_{n-1})^*,\ldots,(x_2-x_1)^*,\\
 &x_1^*+z_1,z_2-z_1,\ldots,z_{n'}-z_{n'-1};
 \varphi_n,\ldots,\varphi_1,\psi_1,\ldots,\psi_{n'}\bigr).
 \end{split}
\end{equation}
When one side is empty, its formula is obtained by removing
the first outer argument of the other vector; when both sides
are empty, the value is one. In addition,
\begin{equation}\label{IV-34-66}
 \overline{K_n(\xi;\underline\varphi)}
       =K_n(\bar\xi;\underline\varphi),\qquad
 K_2(\xi;\varphi,\psi)=k_{\varphi,\psi}(\xi).
\end{equation}
Under (R0), $K_1=0$.
\end{proposition}
\begin{proof}
Holomorphy follows by taking products and matrix elements of
the norm-holomorphic point plates. Choose $\kappa=n_0$ and
apply \eqref{IV-34-18} to each factor to get
\eqref{IV-34-60}. Taking adjoints reverses the first product;
at its junction with the second, \eqref{IV-34-50} replaces
$x_1^*,z_1$ by half their sum. This gives the Gram formula,
including its empty-vector variants.

For the integral formulas choose separating times
$0=c_0<c_1<\cdots<c_n$ for positive-time profiles.
The slab product \eqref{IV-32-45} and
\eqref{IV-34-21} express the class as the integral of
\[
 \prod_a\alpha_a(y_a)
 \prod_a B_{\varphi_a}(y_a-c_{a-1}n_0,c_an_0-y_a)\Omega.
\]
At each junction the arguments sum to $y_{a+1}-y_a$.
Replacing them by half this sum gives \eqref{IV-34-62}.
Taking the inner product with $\Omega$ gives
\eqref{IV-34-61}. A common translation reduces arbitrary
time-ordered supports to this case.

Reality of $\sigma$ on real tests shows that the imaginary
part of the continuous kernel vanishes on real time-ordered
configurations: test it against arbitrary products of real
profiles in small disjoint neighborhoods. The holomorphic
identity theorem gives the first identity in
\eqref{IV-34-66}. Applying the same argument to the two
integral representations of the two-point function identifies
$K_2$ with $k_{\varphi,\psi}$ first on $\mathbb R^2_+$ and
then on $\mathsf{Tb}$. Finally a profile of integral one
identifies $K_1(\varphi)$ with the one-point value, which
vanishes under (R0).
\end{proof}

For a permutation $\mathfrak s$ and a frame
$\upsilon\in\frac\pi2\mathbb Z$, set
\begin{equation}\label{IV-34-69}
 \begin{split}
 \mathcal Z(\mathfrak s,\upsilon)
 &=\{Z:z_{\mathfrak s(a+1)}-z_{\mathfrak s(a)}
                         \in R_\upsilon\mathsf{Tb},\ a<n\},\\
 K_\upsilon^{\mathfrak s}(Z;\underline\varphi)
 &=K_n\bigl(R_{-\upsilon}\partial(\mathfrak sZ);
                              \mathfrak s\underline\varphi\bigr).
 \end{split}
\end{equation}
Here $\partial Z=(z_{a+1}-z_a)_{a<n}$, and permutations act
on points and labels in the same order. Let $\mathcal Z_n$
be the union of these tubes and
$\mathcal O_n=\mathcal Z_n\cap(\mathbb R^2)^n$.

\begin{lemma}\label{IV-34-72}
The functions $K_\upsilon^{\mathfrak s}$ agree on overlaps
and define one holomorphic function
$\mathcal K_n(Z;\underline\varphi)$ on the connected open set
$\mathcal Z_n$. This set is stable under common complex
translations, permutations, quarter turns, conjugation, and
$R_{is}$ for real $s$. The kernel is invariant under the first
three operations, with labels permuted together with points,
and obeys
$\overline{\mathcal K_n(Z;\underline\varphi)}
=\mathcal K_n(\bar Z;\underline\varphi)$.
Its real domain is
\[
 \mathcal O_n=\{Y:\text{all }y_a^0\text{ are distinct, or all }
                         y_a^1\text{ are distinct}\}.
\]
For profiles whose product support lies in one tube, integration
of $K_\upsilon^{\mathfrak s}$ against their product gives
$\sigma(\alpha_1\otimes\varphi_1,\ldots,\alpha_n\otimes\varphi_n)$.
\end{lemma}
\begin{proof}
Each tube is a nonempty open convex cone, the preimage of
$(R_\upsilon\mathsf{Tb})^{n-1}$ under a surjective linear
difference map. A real configuration belongs to it precisely
when $n_\upsilon\cdot y_{\mathfrak s(a)}$ is strictly
increasing. Permuting the entries and making the quarter turn
therefore reduces its integral formula to
\eqref{IV-34-61}, by cubic invariance of $\sigma$.

Every nonempty intersection of two tubes is convex and
contains the real part of each of its points. On its real
open subset, the two kernels have equal integrals against
all products of small supported profiles. Their difference
therefore vanishes there. The holomorphic identity theorem
on the connected intersection proves agreement everywhere.
This proves that the kernels glue without a choice of
continuation path.

The stated stability follows from the difference inequalities
defining the tubes; in light-cone coordinates $R_{is}$ multiplies
the two coordinates by positive real numbers. The kernel
symmetries follow from its definition and
\eqref{IV-34-66}. There are at most $4n!$ tubes.
The complement of $\mathcal O_n$ is the finite union
\[
 \{y_a^0=y_b^0,\ y_k^1=y_l^1\},\qquad a\ne b,\ k\ne l.
\]
Each subspace has codimension two, so its finite complement
is connected. Every tube is connected and meets this real
set. Their union $\mathcal Z_n$ is therefore connected.
\end{proof}

\section{Rotation invariance of all correlations}
\label{sec:all-rotations}

We now assume (R0)--(R5). The polynomial growth of the point
plates implies that their kernels have only finitely many angular
frequencies. Positivity imposes an arithmetic restriction on the
highest frequency; this reduces the rotation problem to the
two-point function.

\subsection{Angular frequencies and their phases}

\begin{lemma}\label{IV-37-2}
Put $\mathcal M_n=\{m\in\mathbb Z:4|m|\le\wp n\}$.
For each tuple of $n\ge2$ labels there are unique holomorphic
functions $c_m$ on $\mathcal Z_n$, zero for $m\notin\mathcal M_n$,
such that
\begin{equation}\label{IV-37-4}
 \mathcal K_n(R_{is}Z;\underline\varphi)
       =\sum_{m\in\mathcal M_n}c_m(Z;\underline\varphi)e^{-4ms}.
\end{equation}
They inherit permutation and translation invariance and satisfy
\[
 c_m(R_{is}Z)=e^{-4ms}c_m(Z),\qquad
 \overline{c_m(X)}=c_{-m}(X)\quad(X\in\mathcal O_n).
\]
They extend uniquely to real-analytic functions $\widetilde c_m$
on configurations of distinct planar points, with
\begin{equation}\label{IV-37-14}
 \widetilde c_m(v+R_t(X-v);\underline\varphi)
              =e^{4imt}\widetilde c_m(X;\underline\varphi).
\end{equation}
Here a common vector $v$ is subtracted from every point.
\end{lemma}
\begin{proof}
First let the consecutive differences of a real configuration
$X$ make angles of modulus less than $\pi/8$ with $n_0$.
For every complex $\zeta=t+is$, choose a quarter-turn frame
$\upsilon$ with $|t-\upsilon|\le\pi/4$.
All the differences of $R_\zeta X$ lie in
$R_\upsilon\mathsf{Tb}$, with tube distances at least
\[
 \min_a|x_{a+1}-x_a|\cos(3\pi/8)e^{-|s|}.
\]
Thus $h_X(\zeta)=\mathcal K_n(R_\zeta X)$ is entire,
$\pi/2$-periodic, and bounded by $C_Xe^{\wp n|s|}$.
Shifting the Fourier contour vertically shows that its
coefficient at $e^{4im\zeta}$ vanishes if $4|m|>\wp n$.
Fourier uniqueness on the real circle then gives a finite
expansion, valid throughout the complex plane.

To define the coefficients at every $Z$, let
$\mu=\#\mathcal M_n$ and invert the Vandermonde matrix
$V_{lm}=e^{-4ml}$, $0\le l<\mu$. If $W=V^{-1}$, put
\[
 c_m(Z)=\sum_{l=0}^{\mu-1}W_{ml}\mathcal K_n(R_{il}Z).
\]
These functions are holomorphic. On each tube, permuted
quarter-turned configurations of the preceding narrow angular
type form a nonempty real open set. On that set the formula
recovers the Fourier coefficients. The holomorphic identity
theorem on the tube proves \eqref{IV-37-4} everywhere.
The same Vandermonde comparison proves uniqueness and all
the stated symmetry identities.

We give the phase argument explicitly because $\mathcal O_n$
need not contain an entire real rotation orbit. If $X$ is a
real point of a fixed tube, write its successive differences
in that frame as $r_an_{\vartheta_a}$,
$|\vartheta_a|<\pi/2$. For
\[
 t\in I_X=\bigcap_a(-\pi/2-\vartheta_a,\pi/2-\vartheta_a),
\]
the full vertical strip $R_{t+is}X$ remains in this tube.
Analytically continuing \eqref{IV-37-4} in the parameter and
comparing coefficients gives
$c_m(R_tX)=e^{4imt}c_m(X)$ on $I_X$.

Consider next configurations with no coincident points and no
perpendicular pair of distinct difference vectors. They form
a rotation-invariant dense open set contained in $\mathcal O_n$.
For density, the product of all squared differences and all
such scalar products is a polynomial nonzero at
$((a,a^2))_{a=1}^n$; on the line segment from any configuration
to this one it has only finitely many zeros.
For these configurations the preceding local phase identity
holds along the entire orbit, so
$t\mapsto e^{-4imt}c_m(R_tX)$ is locally constant on $\mathbb R$,
hence constant. Continuity extends the identity to every
$X,R_tX\in\mathcal O_n$.

For any configuration of distinct points, all but finitely
many angles modulo $\pi$ give distinct time coordinates.
Choose such a $t$ and define
\[
 \widetilde c_m(X)=e^{-4imt}c_m(R_tX).
\]
The preceding identity makes this independent of $t$.
Near each configuration one may keep $t$ fixed, giving a
holomorphic local extension and hence real analyticity.
Translation invariance and the same definition give
\eqref{IV-37-14}, including rotations about arbitrary centers.
\end{proof}

Set $\bar m(0)=\bar m(1)=0$. For $n\ge2$, let $\bar m(n)$
be the largest $|m|$ for which $c_m$ is nonzero for some labels,
or zero if all kernels vanish. Then
$0\le\bar m(n)\le\wp n/4$. A nonzero harmonic is nonzero at
some strictly time-ordered real configuration, after permuting
labels: otherwise the identity theorem on a tube would make
it zero. Conjugation shows that, if $\bar m(n)>0$, the
coefficient $c_{-\bar m(n)}$ is nonzero at such a configuration.
Write $\widetilde K_n=\sum_m\widetilde c_m$.

\subsection{Positivity of the highest harmonic}

\begin{lemma}\label{IV-37-18}
For all $a,b\ge0$,
\begin{equation}\label{IV-37-19}
 2\bar m(a+b)\le\bar m(2a)+\bar m(2b).
\end{equation}
For $a\ge1$, let
$t_a=\widetilde c^{(2a)}_{-\bar m(2a)}$ on labelled
configurations, where a labelled configuration is the unordered
set of its distinct points and their labels. On configurations
of $a$ points in $\mathbb R^2_+$,
\[
 G_a(S,T)=t_a(\theta S\cup T)
\]
is a positive Hermitian kernel. If $\mathsf L$ is any line,
and $S,T$ lie strictly on one side, then
\begin{equation}\label{IV-37-25}
 |t_a(\theta_{\mathsf L}S\cup T)|^2
 \le |t_a(\theta_{\mathsf L}S\cup S)|
       |t_a(\theta_{\mathsf L}T\cup T)|.
\end{equation}
\end{lemma}
\begin{proof}
For strictly time-ordered positive configurations $x,z$,
the Gram formula and $(R_{is}x)^*=R_{is}x^*$ give
\[
 \langle\Psi(R_{is}x),\Psi(R_{is}z)\rangle
 =\mathcal K_{a+b}(R_{is}(\theta x_a,\ldots,\theta x_1,
                                    z_1,\ldots,z_b)).
\]
An empty configuration gives $\Omega$. If $\bar m(a+b)=0$,
\eqref{IV-37-19} follows from nonnegativity. Otherwise, by symmetry,
take $b\ge1$. Choose a time-ordered configuration where
$c_{-\bar m(a+b)}\ne0$ and translate a separating time between
its first $a$ and last $b$ points to zero; when $a=0$, choose
this time below the first point. Cauchy--Schwarz for the two resulting vectors bounds
the squared expansion with leading exponent
$8\bar m(a+b)s$ by a product whose maximal exponent is
$4(\bar m(2a)+\bar m(2b))s$. Letting $s\to+\infty$ proves
\eqref{IV-37-19}.

For a finite family of $a$-point positive configurations,
multiply its Gram matrix at $R_{is}$ by
$e^{-4\bar m(2a)s}$ and let $s\to+\infty$.
Its limit is the matrix of $G_a$, hence positive Hermitian.
If some times agree, perturb the $k$th time by $k\epsilon$,
avoiding the finitely many exceptional $\epsilon$, and use
continuity of $\widetilde c_m$.

If the positive normal to $\mathsf L$ has angle $\upsilon$
and $v\in\mathsf L$, the map $\iota y=R_{-\upsilon}(y-v)$
takes its selected side to $\mathbb R^2_+$. The phase law gives
\[
 t_a(\theta_{\mathsf L}S\cup T)
       =e^{-4i\bar m(2a)\upsilon}G_a(\iota S,\iota T).
\]
Cauchy--Schwarz for this positive kernel proves
\eqref{IV-37-25}; the common phase has modulus one.
\end{proof}

\begin{proposition}\label{IV-37-27}
One has
\begin{equation}\label{IV-37-29}
 \bar m(2a)=a\bar m(2),\qquad
 \bar m(n)\le\frac n2\bar m(2).
\end{equation}
\end{proposition}
\begin{proof}
We first prove $a\mid2\bar m(2a)$. Let $\mathcal N_a(S)$
mean that $t_a$ vanishes on every $2a$-point labelled
configuration containing $S$.

We need a halving observation. Suppose $S'$ has $j'\le a$
points on one side of a line $\mathsf L$, and
$S\subset\theta_{\mathsf L}S'\cup S'$ satisfies
$\mathcal N_a(S)$. Fix the points and labels of
$\theta_{\mathsf L}S'$ and arbitrary labels on another
$2a-j'$ points. The resulting value of $t_a$ is real
analytic in these free points on the complement of their
coincidence sets. This domain is connected, since the excluded
affine subspaces have codimension two. On the nonempty open
subset where $a-j'$ free points lie on the negative side and
$a$ on the positive side, write the configuration as
$\theta_{\mathsf L}T\cup T'$ with $S'\subset T$.
The diagonal value $t_a(\theta_{\mathsf L}T\cup T)$ vanishes
by $\mathcal N_a(S)$; \eqref{IV-37-25} makes the mixed value
zero. The real-analytic identity theorem therefore gives
$\mathcal N_a(\theta_{\mathsf L}S')$ on all configurations.
Exchanging the two sides proves $\mathcal N_a(S')$ as well.

Suppose now $a\nmid2\bar m(2a)$, and fix any label $\varphi$.
Give all vertices of a regular $2a$-gon this label. Rotation
through $\pi/a$ preserves the labelled configuration, whereas
the phase law multiplies its top harmonic by
$e^{-4i\bar m(2a)\pi/a}\ne1$. Thus $\mathcal N_a$ holds for
every full polygon. We descend from arcs of length $j\ge2$
to arcs of length $j'=\lceil j/2\rceil$. Write the vertices
as $P_k=p+\rho n_{\omega+k\pi/a}$. Reflection in the line
through $p$ of direction
$\omega+(k_0+1/2)\pi/a$ sends $P_{k_0+k}$ to
$P_{k_0+1-k}$. The arc $P_{k_0+1},\ldots,P_{k_0+j'}$
lies strictly on one side, and its reflected union is an
arc of length $2j'$ containing an arc of length $j$.
The halving observation proves $\mathcal N_a$ for the shorter
arc. Iteration reaches a singleton at any prescribed point.
Since the label was arbitrary, $t_a$ then vanishes on every
configuration. This contradicts the definition of the
nonzero top harmonic. Hence the divisibility holds.

Put $\nu(a)=\bar m(2a)$. By \eqref{IV-37-19}, $\nu$ is a
convex integer sequence with $\nu(0)=0$ and
$\nu(a)\le\wp a/2$. Its increments are increasing bounded
integers, hence eventually equal to an integer $d$.
Thus $\nu(a)=da-c$ eventually, where
$c=\sum_k(d-(\nu(k)-\nu(k-1)))\ge0$ is a finite integer.
Since $a\mid2\nu(a)$, every sufficiently large $a$ divides
$2c$, forcing $c=0$. All increments equal $d=\nu(1)$.
This proves the even identity. Apply \eqref{IV-37-19} with
$a=\lfloor n/2\rfloor$ and $b=n-a$ to obtain the second one.
\end{proof}

\subsection{From kernels to separated tests}

\begin{proposition}\label{IV-37-30}
For arbitrary profiles with pairwise disjoint supports,
\begin{equation}\label{IV-37-31}
 \sigma(\alpha_1\otimes\varphi_1,\ldots,\alpha_n\otimes\varphi_n)
       =\int\prod_a\alpha_a(y_a)
                 \widetilde K_n(Y;\underline\varphi)\,dY.
\end{equation}
If $\bar m(n)=0$, the functional is invariant under $R_t$
on all $n$-tuples with pairwise disjoint planar projections.
\end{proposition}
\begin{proof}
Subtract the right side of \eqref{IV-37-31} from the left
and call the difference $D(\underline\alpha)$.
It is multilinear when each profile is restricted to its
own separated compact support. We first use real profiles;
complexification then gives the general statement.

Let $A_a^0,A_a^1$ be the time and space interval hulls of
the supports. Suppose each pair whose $A_a^1$ intersect
has a separating time level outside every $A_e^0$.
Induct on the number of such pairs. When there are none,
all supports lie in one permuted quarter-turned tube, so
\cref{IV-34-72} gives $D=0$.
For the induction step choose a level $c$ and split the
profiles below and above it. Translate the upper group
by $sn_{\pi/2}$. Both terms of $D$ are real analytic in
$s$. For the correlation this follows from the matrix
element of $\mathsf T(2\delta,s)$, where $2\delta>0$ is
a time gap about $c$; it is holomorphic on
$|\operatorname{Im}s|<2\delta$. For the kernel integral,
the compact product support remains in the distinct-point
domain for every real $s$. Local holomorphic extensions
on a finite cover of this compact set justify integration
and give real analyticity. For large $s$ all cross-group
spatial intervals are disjoint, so the number of
intersecting pairs decreases. The induction hypothesis
and the identity theorem imply $D=0$ also at $s=0$.

For general profiles with separation $\delta$, partition
the plane into the $C^1$ squares obtained from
$\chi(2y^0/\epsilon-j^0)\chi(2y^1/\epsilon-j^1)$, with
$0<\epsilon<\delta/(n+1)$. Each nonzero tuple of pieces
has interval lengths at most $\epsilon$. If two spatial
intervals overlap, their time intervals have a gap at
least $\sqrt{\delta^2-4\epsilon^2}>(n-2)\epsilon$.
The other $n-2$ time intervals cannot cover this gap,
so there is a level outside all intervals. The preceding
argument applies to every tuple of pieces; multilinearity
proves \eqref{IV-37-31}.

If $\bar m(n)=0$, the phase law makes
$\widetilde K_n=\widetilde c_0$ invariant under all real
rotations. Changing variables in \eqref{IV-37-31} proves
invariance for plates. For a general tuple with disjoint
planar projections, convolve each entry with the product
mollifier
\[
 J_\epsilon(x)=\prod_{\mu=1}^3
             \frac2\epsilon\chi(2x_\mu/\epsilon).
\]
For fixed sufficiently small $\epsilon$, the convolution
Riemann sums are finite sums of translates of $J_\epsilon$,
each of which is a plate. Their planar supports stay in
fixed mutually disjoint neighborhoods. Multilinearity
proves rotation invariance for these sums. They converge
in $C^1$ to the convolutions, which converge in $C^1$ to
the original entries. The telescoping consequence of (R2),
\[
 |\sigma(g_1,\ldots,g_n)-\sigma(g'_1,\ldots,g'_n)|
 \le C_n^\sigma(d;R)\sum_a\|g_a-g'_a\|_\sharp
             \prod_{b\ne a}\max(\|g_b\|_\sharp,\|g'_b\|_\sharp),
\]
on fixed separated compact supports permits both limits.
\end{proof}

\begin{theorem}[Rotation invariance]\label{IV-37-34}
Every functional satisfying (R0)--(R7) is invariant under
$\mathrm O(3)$.
\end{theorem}
\begin{proof}
The two-point result \cref{IV-33-57} and the identification
$K_2=k$ show that $K_2(R_{is}\xi)$ is independent of $s$
for real $\xi\in\mathbb R^2_+$. Comparing the coefficients
in \eqref{IV-37-4}, and then applying the identity theorem,
gives $\bar m(2)=0$. By \eqref{IV-37-29}, all $\bar m(n)$
vanish. Proposition~\ref{IV-37-30} gives invariance under
$R_t$ on tuples with disjoint planar projections, and
\cref{IV-32-48} gives it on every separated tuple.
The empty and one-point values are invariant by normalization
and (R0).

Cubic conjugation gives rotations about the other coordinate
axes. These generate $\mathrm{SO}(3)$: rotate an arbitrary
unit vector to $e_3$ by two coordinate-plane rotations,
then the remaining transformation fixes $e_3$ and is a
rotation in its perpendicular plane. A coordinate reflection
already belongs to the cubic group, so adjoining it gives
all of $\mathrm O(3)$.
\end{proof}

\section{Uniform clustering on a lattice}
\label{sec:lattice-clustering}

The estimate needed for decomposition into continuum components must
be uniform as the spacing tends to zero.  A decay estimate involving
the supremum of a lattice polynomial would not suffice.  We instead
bound correlations by the reflection norms of their local factors.
The argument uses coordinate transfer operators, the spectral
restriction imposed by a diagonal reflection, and a Fourier cutoff
which removes zero energy.

Throughout this section, $\nu$ is a translation-invariant probability
law on the configurations over $a\mathbb Z^3$, $a>0$.  We assume RP
at the six coordinate mirrors
$\Sigma(e_\mu,0)$ and $\Sigma(e_\mu,a/2)$, and at one diagonal
mirror $\Sigma(e_\mu+\epsilon e_\lambda,c)$ in each coordinate
plane, where $\epsilon\in\{1,-1\}$ and $c\in a\mathbb Z$.
Translations move these diagonal mirrors to height zero and move
the coordinate mirrors to every height in $(a/2)\mathbb Z$.
In particular these hypotheses hold for every periodic Gibbs state
by Proposition~\ref{IV-54-17}.

\subsection{Transfer operators and the lattice spectral restriction}

For a mirror $\Sigma(n,c)$ and a positive normal lattice vector $v$,
suppose RP also holds at $\Sigma(n,c+n\cdot v/2)$.  Complete the
bounded local functions supported in the upper half-space in the form
$\langle[F],[G]\rangle=\nu(\overline{\theta_{n,c}F}G)$, after
quotienting its null space.  This is the reflection Hilbert space at
the mirror.  Coordinate pictures use $n=e_\mu$ and $v=ae_\mu$;
diagonal pictures use $n=e_\mu+\epsilon e_\lambda$ and $v=an$.

\begin{lemma}\label{IV-30-5}
In these pictures the normal translation and the tangential
translations induce commuting operators
\[
 T[F]=[\tau_vF],\qquad V(w)[F]=[\tau_wF]\quad(n\cdot w=0),
 \qquad 0\le T\le1,
\]
where the $V(w)$ are unitary.  Moreover
\begin{equation}\label{IV-30-9}
 \langle[F],T^kV(w)[G]\rangle
 =\nu(\overline{\theta_{n,c}F}\,\tau_{kv+w}G),
 \qquad k\ge0,
\end{equation}
and $\|[F]\|^2\le\nu(|F|^2)\le\|F\|_\infty^2$.
In a coordinate picture there is a joint spectral measure
$\mathfrak E$ on $[0,1]\times(S^1)^2$ such that
\begin{equation}\label{IV-30-17}
 T^kV(a y)=\int t^k\zeta^y\,d\mathfrak E(t,\zeta),
 \qquad y_\mu=0.
\end{equation}
\end{lemma}

\begin{proof}
Polarization of RP, followed by a translation of an arbitrary local
function into a half-space, gives invariance under the mirror.
Ordinary Cauchy--Schwarz therefore bounds the reflection norm by
the $L^2(\nu)$ norm.  Normal translation is symmetric in the
reflection form.  Reflection Cauchy--Schwarz gives
\[
 \|[\tau_vF]\|^2
 \le\|[F]\|\,\|[\tau_{2v}F]\|.
\]
Iteration and the uniform bound by $\|F\|_\infty$ show
$\|[\tau_vF]\|\le\|[F]\|$.  Thus translation descends through
the null space to a self-adjoint contraction.  The identity
$\theta_{n,c}=\theta_{n,c+n\cdot v/2}\tau_v$ and RP at the
second mirror show $\langle[F],T[F]\rangle\ge0$.
Tangential translations preserve the form and have inverse
translations, hence induce unitaries.  Their commutation and
\eqref{IV-30-9} follow on local classes and extend by continuity.
The joint spectral theorem for the commuting normal operators
$T,V(ae_\lambda),V(ae_{\lambda'})$ gives
\eqref{IV-30-17}.
\end{proof}

The following elementary complex-variable estimate isolates the
spectral consequence of comparing a coordinate picture with a
diagonal picture.

\begin{lemma}\label{IV-30-21}
Suppose a finite positive measure $\eta$ on the closed unit disc
satisfies, for every polynomial $q$,
\begin{equation}\label{IV-30-22}
 \int|q|^2\,d\eta\le C\sup_{[0,1]}|q|\sup_{S^1}|q|.
\end{equation}
Then $\eta\{|1+z|^2<4|z|\}=0$.
\end{lemma}

\begin{proof}
Uniform Taylor approximation extends \eqref{IV-30-22} to functions
holomorphic on a disc of radius greater than one.  On the unit disc
put
\[
 \Lambda(\xi)=\log\frac{1+\xi}{1-\xi},\qquad
 g_s(z)=\cos\bigl(s\Lambda(\sqrt z)\bigr),\qquad s>0.
\]
The principal logarithm is analytic because its argument is in the
right half-plane.  Since $\Lambda$ is odd and cosine is even,
$g_s$ is a single-valued holomorphic function of $z$.
Write $A(z)=|\operatorname{Im}\Lambda(\sqrt z)|$.  Then
\[
 \sinh^2(sA(z))\le|g_s(z)|^2\le\cosh^2(sA(z)),
 \qquad 0\le A(z)<\pi/2,
\]
and $A=0$ on $[0,1)$.  Applying \eqref{IV-30-22} to
$g_s(rz)$, $0<r<1$, gives
\[
 \int\sinh^2(sA(rz))\,d\eta(z)\le C\cosh(s\pi/2).
\]
Extend $A$ to the unit circle by $A=\pi/2$ away from $1$ and
$A(1)=0$; these are its radial limits.  Bounded convergence allows
$r\uparrow1$.  For every $\delta>0$ the last inequality implies
\[
 \eta\{A\ge\pi/4+\delta\}\le16C e^{-2s\delta}
\]
for all sufficiently large $s$, so this set is null.
Inside the disc,
\[
 \tan A(z)=\frac{\sqrt{2(|z|-\operatorname{Re}z)}}{1-|z|}.
\]
Consequently $A(z)>\pi/4$ is equivalent to
$|1+z|^2<4|z|$, with the same equivalence at the boundary by the
radial values.  A countable union over $\delta>0$ proves the claim.
\end{proof}

\begin{proposition}\label{IV-30-32}
In every coordinate picture and for each tangential direction
$\lambda\ne\mu$,
\begin{equation}\label{IV-30-33}
 (1-T)^2\ge T\bigl(2-V(ae_\lambda)-V(ae_\lambda)^*\bigr).
\end{equation}
Equivalently, its joint spectral measure is supported where
$(1-t)^2\ge t|1-\zeta_\lambda|^2$.  In particular, with
\begin{equation}\label{IV-30-42}
 \mathsf k=\frac{2}{\pi}\log(1+\sqrt2)>0,
\end{equation}
the spectral measure is supported where
\begin{equation}\label{IV-30-45}
 t\le\exp\{-\mathsf k|\arg\zeta_\lambda|\}.
\end{equation}
\end{proposition}

\begin{proof}
Use the available diagonal normal
$n=e_\mu+\epsilon e_\lambda$ and write
$n^\perp=e_\mu-\epsilon e_\lambda$.
For a local upper-half-space function $F$, let $\varsigma$ be the
image of its spectral measure
$\langle[F],\mathfrak E(\cdot)[F]\rangle$ under
$(t,\zeta)\mapsto z=t\zeta_\lambda^\epsilon$.
Expansion of a polynomial $q(z)=\sum q_kz^k$ gives
\begin{equation}\label{IV-30-35}
 \int|q(z)|^2|z|^{2N}\,d\varsigma
 =\sum_{k,l}q_k\overline{q_l}\,
 \nu\bigl(\overline{\theta F}\,
       \tau_{kan+lan^\perp}\tau_{2Nae_\mu}F\bigr).
\end{equation}
Choose a diagonal mirror $\Sigma(n,c)$ beyond every endpoint of
the bonds supporting $\theta F$, then choose $N$ so that every
endpoint supporting $\tau_{2Nae_\mu}F$ lies above that mirror.
Both choices depend only on $F$.  Set
$F_1=\theta_{n,c}\theta F$ and $F_2=\tau_{2Nae_\mu}F$.
Both functions lie in its upper half-space.  In the diagonal
picture, the right side of \eqref{IV-30-35} equals
\[
 \langle[F_1],q(T_\parallel)\overline q(V_\perp)[F_2]\rangle,
\]
where $\overline q$ denotes coefficient conjugation.
Its modulus is at most
$\|F\|_\infty^2\sup_{[0,1]}|q|\sup_{S^1}|q|$.
Lemma~\ref{IV-30-21} applied to $|z|^{2N}d\varsigma$ shows
that $\varsigma$ gives no mass to $|1+z|^2<4|z|$: the weight
is strictly positive there, since $z=0$ is outside this set.
The identity
\[
 |1+t\zeta_\lambda^\epsilon|^2-4t
 =(1-t)^2-t|1-\zeta_\lambda|^2
\]
proves the spectral restriction for every local vector, hence for
the whole spectral measure by density.  Functional calculus gives
\eqref{IV-30-33}.

For $t=e^{-x}>0$ and $\zeta_\lambda=e^{iy}$, the restriction is
$\sinh(x/2)\ge|\sin(y/2)|$.  The function
$2\operatorname{arcsinh}(\sin(y/2))$ is concave on $[0,\pi]$:
its derivative is
$\cos(y/2)(1+\sin^2(y/2))^{-1/2}$, a product of nonnegative
decreasing functions.  It therefore lies above the chord
$\mathsf k y$.  This proves \eqref{IV-30-45}; the case $t=0$
is immediate.
\end{proof}

\subsection{The invariant part of a correlation}

Let $\mathcal I$ be the sigma-algebra of translation-invariant
Borel events, completed modulo $\nu$-null sets.  In a coordinate
picture let $\Pi_0$ be the orthogonal projection onto
$\ker(1-T)$.  It describes the nondecaying part of a correlation.

\begin{proposition}\label{IV-31-12}
One has $T^n\to\Pi_0$ strongly and $\Pi_0V(w)=\Pi_0$ for
tangential translations.  Moreover
\begin{equation}\label{IV-31-14}
 \langle[F],\Pi_0[G]\rangle
 =\nu\bigl(\overline{E_\nu[F\mid\mathcal I]}
                   E_\nu[G\mid\mathcal I]\bigr).
\end{equation}
Conditional expectation onto $\mathcal I$ is unchanged by
translations or coordinate reflections.  In particular,
$\nu$ is ergodic precisely when $\ker(1-T)=\mathbb C[1]$.
\end{proposition}

\begin{proof}
By the spectral theorem,
$\Pi_0=\mathfrak E(\{t=1\})$ and $T^n\to\Pi_0$ strongly.
Proposition~\ref{IV-30-32} in both tangential directions says that
$t=1$ forces $\zeta=(1,1)$, proving $\Pi_0V(w)=\Pi_0$.
Let $P_\mu$ be the $L^2(\nu)$ projection onto the functions fixed
by $\tau_{ae_\mu}$.  The mean ergodic theorem and
\eqref{IV-30-9} give, for upper-half-space functions,
\begin{equation}\label{IV-31-16}
 \langle[F],\Pi_0[G]\rangle
 =\langle\theta F,P_\mu G\rangle_{L^2(\nu)}.
\end{equation}
Since $\Pi_0V(w)=\Pi_0$, the right side is unchanged on replacing
$G$ by $\tau_wG$.  For arbitrary local $F_1,F_2$, choose
$u=kae_\mu$ large enough that
$F=\theta\tau_{-u}F_1$ and $G=\tau_uF_2$ are in the upper
half-space.  Substitution in this identity, using
$P_\mu\tau_{\pm u}=P_\mu$, gives
\[
 \langle F_1,P_\mu\tau_wF_2\rangle_{L^2(\nu)}
 =\langle F_1,P_\mu F_2\rangle_{L^2(\nu)}.
\]
Local functions are dense in $L^2(\nu)$, so
$P_\mu\tau_w=P_\mu$ and, by adjoints,
$\tau_wP_\mu=P_\mu$.  Thus a function fixed by the one normal
translation is fixed by every lattice translation.  The invariant
subspace is $L^2(\mathcal I,\nu)$, so
$P_\mu=E_\nu[\,\cdot\mid\mathcal I]$.

Reflection preserves that subspace and its orthogonal complement,
hence commutes with $P_\mu$.  Positivity of the left side of
\eqref{IV-31-16} at $G=F$ implies
$\langle\theta P_\mu F,P_\mu F\rangle_{L^2(\nu)}\ge0$.
The vectors $P_\mu F$ with $F$ in the upper half-space are dense
in the invariant subspace: translate each local function there
without changing its projection.  Reflection is consequently a
positive involution on that subspace and must be the identity;
its negative eigenspace is zero.  Formula~\eqref{IV-31-16} now
becomes \eqref{IV-31-14}.
If $\nu$ is ergodic, the right side is
$\overline{\nu(F)}\nu(G)$, so $\Pi_0$ has range $\mathbb C[1]$.
Conversely that range identity implies
$\|E_\nu[F\mid\mathcal I]\|_2=|\nu(F)|$, first for local
$F$ and then for all $L^2$ functions.  For $F=1_E$, $E\in\mathcal I$,
it says $\nu(E)^{1/2}=\nu(E)$, proving ergodicity.
\end{proof}

\subsection{A bounded density after removing zero energy}

Write $Q_r=[-ra,ra]^3$, where $r\ge1$ is an integer, and
$\theta_\mu=\theta_{e_\mu,0}$.  For a bounded Borel function $A$
of the bonds in $Q_r$, use the site picture to set
\[
 \psi_A^\mu=[\tau_{rae_\mu}A]_\mu.
\]
If $\mathsf S_8$ is the group of coordinate sign changes, define
\begin{align}
 \mathcal N_r(A)&=\max_{\mu,\,s\in\mathsf S_8}
                      \|\psi_{sA}^\mu\|^2,
 \label{IV-31-32}\\
 \mathcal N'_{r,\mu}(A)&=\max\left\{
   \|\psi_A^\mu\|^2,
   \max_{\lambda\ne\mu}
      \|\psi_{\theta_\lambda A}^\lambda\|
      \|\psi_A^\lambda\|\right\}.
 \label{IV-31-33}
\end{align}
Then $\mathcal N'_{r,\mu}(A)\le\mathcal N_r(A)
\le\nu(|A|^2)\le\|A\|_\infty^2$.
The square root of $\mathcal N_r$ is a seminorm, invariant under
$\mathsf S_8$.  Consequently the parity parts
$A_\pm=(A\pm\theta_\mu A)/2$ satisfy
$\mathcal N_r(A_\pm)\le\mathcal N_r(A)$.
Translation in the reflection form also gives
\begin{equation}\label{IV-31-38}
 \nu(\overline A\tau_xB)
 =\langle\psi_{\theta_\mu A}^\mu,
       T_\mu^{x_\mu/a-2r}V_\mu(x_\perp)\psi_B^\mu\rangle,
 \qquad x_\mu\ge2ra.
\end{equation}
The correlation at $-x$ is the conjugate of
$\nu(\overline B\tau_xA)$; coordinate reflection changes
$x_\mu$ to $-x_\mu$ and reflects both functions.  These identities
allow the same correlation to be estimated in any coordinate picture.

We require a scalar weight which vanishes at $t=1$ but controls
$1-t$ from below.  Let $z_m$ be the even, period-four sequence
$(z_0,z_1,z_2,z_3)=(0,1,-2,1)$, and let
\[
 B(u)=8(1-4|u|/\pi)_+\quad(-\pi\le u\le\pi)
\]
be extended with period $2\pi$.  Its Fourier coefficients are
\[
 b_0=1,\qquad b_m=\left(\frac{8\sin(m\pi/8)}{m\pi}\right)^2
 \quad(m\ne0).
\]
Integration by parts gives this formula; in particular the Fourier
series converges absolutely.  Put
\begin{equation}\label{IV-31-22}
 B_z(u)=\sum_{m\in\mathbb Z}z_mb_me^{imu},\qquad
 \Psi(t)=2\sum_{m\ge1}z_mb_mt^{m-1},\quad 0\le t\le1.
\end{equation}
The convention $0^0=1$ is used here.

\begin{lemma}\label{IV-31-23}
The above series converge uniformly.  One has
\[
 \sum_m|z_m|b_m=8,\qquad |B_z|\le8,\qquad
 B_z(u)=0\quad\text{if }\operatorname{dist}(u,2\pi\mathbb Z)\le\pi/4.
\]
Moreover $0\le\Psi\le8$, $\Psi(1)=0$, $\Psi(0)=2b_1>0$, and
\begin{equation}\label{IV-31-26}
 \Psi(e^{-x})\ge c_\Psi\tanh x\quad(x>0),\qquad
 c_\Psi=(1+\sqrt2)^{-1}.
\end{equation}
\end{lemma}

\begin{proof}
The identities $|z_m|=1-\cos(m\pi/2)$ and
$z_m=\cos(m\pi/2)-\cos(m\pi)$ give
\[
 \sum_m|z_m|b_m=B(0)-B(\pi/2)=8,
 \qquad B_z(u)=\tfrac12B(u+\pi/2)+\tfrac12B(u-\pi/2)-B(u+\pi).
\]
The three translated triangles have disjoint interiors of their
supports and vanish for $\operatorname{dist}(u,2\pi\mathbb Z)
\le\pi/4$.  This proves the stated bounds and, at $u=0$,
$\Psi(1)=0$.  Absolute summability gives the uniform convergence.
For $0<t<1$, the Poisson kernel
$P_t(u)=(1-t^2)/(1-2t\cos u+t^2)$ yields
\[
 t\Psi(t)=\frac1{2\pi}\int_{-\pi/4}^{\pi/4}B(v)
 \left[\tfrac12P_t(v-\pi/2)+\tfrac12P_t(v+\pi/2)-P_t(v-\pi)\right]dv.
\]
Writing $s=1+t^2$, the bracket is at least
\[
 (1-t^2)\left[\frac1s-\frac1{s+\sqrt2t}\right]
 =\frac{(1-t^2)\sqrt2t}{s(s+\sqrt2t)}:
\]
the average of $(s\mp2t\sin v)^{-1}$ is at least $1/s$, whereas
$\cos v\ge1/\sqrt2$.  Since $(2\pi)^{-1}\int B=1$, division
by $t$ gives
\[
 \Psi(t)\ge\frac{1-t^2}{1+t^2}
                  \frac{\sqrt2}{1+t^2+\sqrt2t}
 \ge c_\Psi\frac{1-t^2}{1+t^2}.
\]
This is \eqref{IV-31-26}.  Together with $\Psi(0)>0$ and
$\Psi(1)=0$ it proves positivity on $[0,1)$.
\end{proof}

Set
\begin{equation}\label{IV-31-4}
 C_{\mathrm{dn}}=\frac1{4\pi^2}
 \left[128+64\left(\frac{16}{\pi\mathsf k}+2+
                              \frac{64}{(\pi\mathsf k)^2}\right)\right].
\end{equation}

\begin{lemma}\label{IV-31-41}
If $\Phi$ is supported in $Q_r$ and $\theta_\mu\Phi=\epsilon\Phi$,
$\epsilon\in\{1,-1\}$, let
$\varrho_\psi=\langle\psi,\mathfrak E_\mu(\cdot)\psi\rangle$,
where $\psi=\psi_\Phi^\mu$.  The image of
$\Psi(t^{2r})d\varrho_\psi(t,\zeta)$ on $(S^1)^2$ has a
continuous density $h$ relative to unnormalized angular Lebesgue
measure $d\alpha$ on $(-\pi,\pi]^2$, and
\begin{equation}\label{IV-31-44}
 0\le h\le C_{\mathrm{dn}}r^2\mathcal N'_{r,\mu}(\Phi).
\end{equation}
\end{lemma}

\begin{proof}
Write $u(v)=\nu(\overline{\theta_\mu\Phi}\tau_v\Phi)
=\epsilon\nu(\overline\Phi\tau_v\Phi)$ and
\[
 K(y)=\sum_{m\in\mathbb Z}z_mb_m\,
         u\bigl(a(2rme_\mu+y)\bigr),\qquad y_\mu=0.
\]
The series is absolutely convergent.  Reflection parity makes
$u(a(-ke_\mu+y))=u(a(ke_\mu+y))$.  By
\eqref{IV-31-38}, for $k\ge2r$,
\[
 u(a(ke_\mu+y))=\langle\psi,T_\mu^{k-2r}V_\mu(ay)\psi\rangle.
\]
Pairing $m$ with $-m$, and observing $z_0=0$, therefore gives
\begin{equation}\label{IV-31-48}
 K(y)=\int\Psi(t^{2r})\zeta^y\,d\varrho_\psi(t,\zeta),
 \qquad |K(y)|\le8\|\psi\|^2.
\end{equation}
Thus $K$ is the sequence of Fourier coefficients of the weighted
momentum measure.

To bound large $y$, choose a tangential coordinate $\lambda$ with
$|y_\lambda|=|y|_\infty$.  Since $K(-y)=\overline{K(y)}$, assume
$y_\lambda\ge2r$.  Applying \eqref{IV-31-38} in the
$\lambda$ picture and summing its absolutely convergent series gives
\[
 K(y)=\epsilon\langle\psi_{\theta_\lambda\Phi}^\lambda,
  X\psi_\Phi^\lambda\rangle,
 \quad
 X=\int t^{y_\lambda-2r}
        B_z(2r\arg\zeta_\mu)\zeta_{\lambda'}^{y_{\lambda'}}
                  \,d\mathfrak E_\lambda.
\]
Where $B_z(2r\arg\zeta_\mu)$ is nonzero,
$|\arg\zeta_\mu|>\pi/(8r)$.  The lattice spectral restriction
\eqref{IV-30-45} then gives $t\le e^{-\pi\mathsf k/(8r)}$.
Consequently
\begin{equation}\label{IV-31-49}
 |K(y)|\le8e^{-\pi\mathsf k(|y|_\infty-2r)/(8r)}
                         \mathcal N'_{r,\mu}(\Phi),
 \qquad |y|_\infty\ge2r.
\end{equation}
There are at most $16r^2$ points with $|y|_\infty<2r$ and
$8j$ points with $|y|_\infty=j$.  For $c=\pi\mathsf k/(8r)$,
\[
 \sum_{j\ge2r}64j e^{-c(j-2r)}
 =64\left[\frac{2r}{1-e^{-c}}+
                 \frac{e^{-c}}{(1-e^{-c})^2}\right]
 \le64\left[2r+\frac{2r}{c}+\frac1{c^2}\right].
\]
Here $(1-e^{-c})^{-1}\le1+c^{-1}$ and
$e^{-c}(1-e^{-c})^{-2}\le c^{-2}$.
Together with \eqref{IV-31-48}, this proves
\[
 \sum_y|K(y)|\le4\pi^2C_{\mathrm{dn}}r^2\mathcal N'_{r,\mu}(\Phi).
\]
The Fourier series
$h(\zeta)=(2\pi)^{-2}\sum_yK(y)\zeta^{-y}$ is therefore
uniformly convergent and bounded by the right side of
\eqref{IV-31-44}.  Its angular measure has the same Fourier
coefficients as the weighted momentum measure by
\eqref{IV-31-48}.  Density of trigonometric polynomials gives
equality of the measures.  Positivity of the weighted measure and
continuity of $h$ imply $h\ge0$ everywhere.
\end{proof}

\subsection{The clustering estimate}

\begin{proposition}\label{IV-31-51}
Under the nine-mirror hypotheses of this section, put
\begin{equation}\label{IV-31-5}
 C_{\mathrm{dc}}=16c_\Psi^{-1}\mathsf k^{-2}C_{\mathrm{dn}}.
\end{equation}
For every $A,B$ supported in $Q_r$ and every integer $n\ge1$,
\begin{equation}\label{IV-31-52}
 0\le\langle\psi_A^\mu,T_\mu^n(1-\Pi_0)\psi_A^\mu\rangle
 \le C_{\mathrm{dc}}\mathcal N_r(A)
                  \left(\frac rn+2\frac{r^2}{n^2}\right).
\end{equation}
If $x\in a\mathbb Z^3$ and $n=|x|_\infty/a-2r\ge1$, then
\begin{multline}\label{IV-31-53}
 \left|\nu(\overline A\tau_xB)
  -\nu\bigl(\overline{E_\nu[A\mid\mathcal I]}
                        E_\nu[B\mid\mathcal I]\bigr)\right|\\
 \le C_{\mathrm{dc}}\mathcal N_r(A)^{1/2}\mathcal N_r(B)^{1/2}
                  \left(\frac rn+2\frac{r^2}{n^2}\right).
\end{multline}
In particular, for $|x|\ge4\sqrt3\,ra$,
\begin{multline}\label{IV-31-55}
 \left|\nu(\overline A\tau_xB)
  -\nu\bigl(\overline{E_\nu[A\mid\mathcal I]}
                        E_\nu[B\mid\mathcal I]\bigr)\right|\\
 \le4\sqrt3 C_{\mathrm{dc}}\,
       \mathcal N_r(A)^{1/2}\mathcal N_r(B)^{1/2}\frac{ra}{|x|}.
\end{multline}
The constants are independent of the spacing, the law, and the
observables.  For an ergodic law the subtracted term is
$\overline{\nu(A)}\nu(B)$.
\end{proposition}

\begin{proof}
Set $D_n=T_\mu^n(1-\Pi_0)=\int t^n1_{t<1}\,d\mathfrak E_\mu$.
It is positive and commutes with tangential translations.
First suppose $\theta_\mu\Phi=\pm\Phi$ and let
$\psi=\psi_\Phi^\mu$.  By Lemma~\ref{IV-31-41}, the weighted
momentum measure has a bounded density.  It gives no mass to
$\zeta=(1,1)$; since $\Psi(t^{2r})>0$ for $t<1$, the same is
true of $\varrho_\psi$ on $\{t<1,\zeta=(1,1)\}$.
For $0<t<1$ write $s=-\log t$.  The two tangential versions of
\eqref{IV-30-45} imply $s\ge\mathsf k|\alpha|_\infty$,
where $\zeta=e^{i\alpha}$.  Since
$s\mapsto e^{-ns}/\tanh(2rs)$ is decreasing and
$\Psi(t^{2r})\ge c_\Psi\tanh(2rs)$,
\[
 t^n\le c_\Psi^{-1}
 \frac{e^{-n\mathsf k|\alpha|_\infty}}
      {\tanh(2r\mathsf k|\alpha|_\infty)}\,
 \Psi(t^{2r})
\]
away from the null point $\alpha=0$.  At $t=0$, the integrand
$t^n$ is zero.  Integrating against the bounded density and using
$1/\tanh u\le1+1/u$ gives
\begin{align*}
 \langle\psi,D_n\psi\rangle
 &\le c_\Psi^{-1}C_{\mathrm{dn}}r^2\mathcal N'_{r,\mu}(\Phi)
    \int_{(-\pi,\pi]^2}
       \frac{e^{-n\mathsf k|\alpha|_\infty}}
            {\tanh(2r\mathsf k|\alpha|_\infty)}\,d\alpha\\
 &\le c_\Psi^{-1}C_{\mathrm{dn}}\mathsf k^{-2}
       \mathcal N'_{r,\mu}(\Phi)
          \left(\frac{8r^2}{n^2}+\frac{4r}{n}\right).
\end{align*}
Indeed the square integral equals
$8\int_0^\pi u e^{-n\mathsf k u}/\tanh(2r\mathsf k u)\,du$;
extend the upper limit to infinity to obtain the displayed bound.
For arbitrary $A=A_++A_-$, positivity yields
$\langle\psi_A,D_n\psi_A\rangle\le
2\langle\psi_{A_+},D_n\psi_{A_+}\rangle+
2\langle\psi_{A_-},D_n\psi_{A_-}\rangle$.
Use $\mathcal N'_{r,\mu}(A_\pm)\le\mathcal N_r(A)$ to obtain
\eqref{IV-31-52} with \eqref{IV-31-5}.

For \eqref{IV-31-53}, choose $\mu$ with
$|x_\mu|=|x|_\infty$ and first suppose $x_\mu>0$.
Insert $T_\mu^n=\Pi_0+D_n$ into \eqref{IV-31-38}.
Proposition~\ref{IV-31-12} identifies the $\Pi_0$ term with the
subtracted conditional expectation.  Cauchy--Schwarz for $D_n$,
its commutation with $V_\mu(x_\perp)$, and
\eqref{IV-31-52} bound the remainder as claimed.
For $x_\mu<0$, conjugate and interchange $A,B$.
If $|x|_\infty\ge4ra$, then
$n\ge|x|_\infty/(2a)\ge2r$, and
$r/n+2r^2/n^2\le4ra/|x|_\infty$.
Finally $|x|_\infty\ge|x|/\sqrt3$ proves
\eqref{IV-31-55}.  The ergodic assertion follows from
Proposition~\ref{IV-31-12}.
\end{proof}

\section{Decomposition by invariant conditional expectations}
\label{sec:conditional-decomposition}

A continuum limit of invariant states need not cluster.  We construct
its components from conditional expectations on the lattice.  The
essential estimate is uniform in the cutoff: two local products
approach the product of their invariant conditional expectations at
rate inverse distance.  This identifies all moments of the component
measure before a component is selected.

At cutoff $K$, let $\mathcal I$ be the Borel sigma-field invariant under
all lattice translations.  Consider a nonempty class $\mathcal G$ of
laws which are invariant under lattice isometries and reflection
positive at every lattice mirror.  Require also that
\begin{equation}\label{IV-35-2}
 \nu\in\mathcal G,\quad h\ge0\ \text{bounded and }\mathcal I
 \text{-measurable},\quad \nu(h)>0
 \quad\Longrightarrow\quad
 \nu_h:=\frac{h\nu}{\nu(h)}\in\mathcal G .
\end{equation}
For a sequence $K_j\to\infty$ and such classes $\mathcal G_j$, assume
common finite positive bounds
\begin{align}
 |S_n^T[\rho](X)|
 &\le C_n^\infty(d;R)\|X\|_\sharp
       &&(n\ge2,\ X\text{ of type }(d;R)),\label{IV-35-3}\\
 |\rho(\mathcal O_K(f))-\nu(\mathcal O_K(f))|
 &\le C_1(R)\|f\|_\sharp ,
       &&(\rho,\nu\in\mathcal G).\label{IV-35-4}
\end{align}
In the second line $\operatorname{supp}f$ lies in a ball of radius
$R$.  We take $C_1$ nondecreasing, and $C_n^\infty$ nonincreasing in
$d\in(0,\infty]$ and nondecreasing in $R$.

A \emph{frame} is a choice $\nu_j\in\mathcal G_j$.  It converges if
\begin{equation}\label{IV-35-6}
 \sigma(X)=\lim_j\nu_j(A_j(X)),\qquad
 A_j(X):=X^{\nu_j}_{K_j},
\end{equation}
exists for every separated tuple.  A limit in this frame is, more
generally, $\lim_j\rho_j(A_j(X))$ with $\rho_j\in\mathcal G_j$.
The centring law remains $\nu_j$ when the integrating law changes.
Every such limit is normalized, real, symmetric and multilinear on
separated supports.

Write $d(X)$ for the separation and $R(X)$ for the largest of the
minimal radii of balls containing individual supports.  Set
$d(X)=\infty$ for a tuple of length at most one.  Define
\begin{equation}\label{IV-35-9}
 c_n(d;R)=\sum_{\pi\in\operatorname{Part}(n)}
                 \prod_{J\in\pi}b_{|J|},\qquad
 b_1=C_1(R),\quad b_k=C_k^\infty(d;R)\ (k\ge2),\quad c_0=1 .
\end{equation}
The moment bound and the bound used in reflections are
\begin{align}
 C_{\rm mom}(X)&=c_{|X|}(d(X);R(X))\|X\|_\sharp,\label{IV-35-10}\\
 \delta(X;\bar R)&=
 \operatorname{dist}_\infty(\operatorname{supp}X,
                  \mathbb R^3\setminus(-\bar R,\bar R)^3),\label{IV-35-11}\\
 N(X;\bar R)&=
 c_{2|X|}(\min(d(X),2\delta(X;\bar R));R(X))
                         \|X\|_\sharp^2.\label{IV-35-12}
\end{align}
Here the support is contained in the indicated open cube.  Both
bounds are zero if a tuple has a zero entry, and both are one on
the empty tuple.  These conventions avoid evaluating $c_n$ at
$R=0$.  Extend $C_{\rm mom}$ to finite formal sums by the triangle
inequality.  Put
\begin{equation}\label{IV-35-13}
 k_1=4\sqrt3,\qquad k_0=4\sqrt3\,C_{\rm dc},
\end{equation}
where $C_{\rm dc}$ is the absolute constant in the lattice
clustering estimate of \cref{IV-31-51}.

\begin{proposition}[Uniform conditional bounds and clustering]
\label{IV-35-17}\label{IV-35-21}\label{IV-35-22}
For $\rho,\nu\in\mathcal G$ one has
$|\rho(X^\nu_K)|\le C_{\rm mom}(X)$.  If
$\operatorname{supp}f_i\cup\operatorname{supp}f'_i\subset Q_i$,
the sets $Q_i$ lie in balls of radius $R$ and are mutually at
distance at least $d$, then
\begin{equation}\label{IV-35-19}
 |\rho(X^\nu_K-X'^\nu_K)|
 \le c_n(d;R)\sum_i\|f_i-f'_i\|_\sharp
                    \prod_{l\ne i}\max(\|f_l\|_\sharp,\|f'_l\|_\sharp).
\end{equation}
The same bounds hold almost surely for conditional expectations
given $\mathcal I$.

Suppose $\bar R\in a_K\mathbb N$, both tuples are supported in
$(-\bar R,\bar R)^3$, and
$a_K\le2\min(\delta(X;\bar R),\delta(Y;\bar R))$.  For
$v\in a_K\mathbb Z^3$, $|v|\ge k_1\bar R$,
\begin{equation}\label{IV-35-24}
 \left|\rho((X\cdot\tau_vY)^\nu_K)
 -\rho\!\left(E_\rho[X^\nu_K\mid\mathcal I]\,
             E_\rho[Y^\nu_K\mid\mathcal I]\right)\right|
 \le k_0\sqrt{N(X;\bar R)N(Y;\bar R)}\,\frac{\bar R}{|v|}.
\end{equation}
All constants are independent of the cutoff and of the two laws.
\end{proposition}

\begin{proof}
The moment--cumulant identity for two centring laws is
\begin{equation}\label{IV-35-18}
 \rho(X^\nu_K)=\sum_{\pi\in\operatorname{Part}(n)}
                       \prod_{J\in\pi}\kappa_J,\qquad
 \kappa_{\{i\}}=\rho(\mathcal O_K(f_i))-\nu(\mathcal O_K(f_i)),
 \quad \kappa_J=S_{|J|}^T[\rho](X_J)\ (|J|\ge2).
\end{equation}
Applying \eqref{IV-35-3}--\eqref{IV-35-4} block by block gives
$c_n(d;R)\|X\|_\sharp$ on any tuple of type $(d;R)$.
Telescope the product defining $X^\nu_K-X'^\nu_K$ one factor
at a time.  Each resulting tuple remains supported in the same
sets $Q_i$; the preceding bound gives \eqref{IV-35-19}.

For a bounded Borel function $A$ and an invariant density $h$,
every version $e=E_\rho[A\mid\mathcal I]$ remains a version under
$\rho_h$, since
\[
 \rho_h(kA)=\rho(hkA)/\rho(h)=\rho(hke)/\rho(h)
\]
for every bounded invariant $k$.  If $\rho'(A)\le c$ for every
$\rho'\in\mathcal G$, condition on
$\{E_\rho[A\mid\mathcal I]>c+\varepsilon\}$.  A positive
probability for this event would contradict that bound, by
\eqref{IV-35-2}.  Applying this observation to both signs of
$X^\nu_K$ and of its difference from $X'^\nu_K$ proves the
conditional estimates.  A version can be changed to zero on
its exceptional null set to make the absolute bound hold
everywhere.

Put $r=\bar R/a_K$.  A plaquette lies within sup-norm distance
$a_K/2$ of its barycentre.  The margin assumption therefore
makes $A=X^\nu_K$ and $B=Y^\nu_K$ functions of the bonds in
$[-ra_K,ra_K]^3$.  Lattice covariance of $\mathcal O_K$ and
invariance of $\nu$ give
\[
 (\mathsf rX)^\nu_K=\mathsf r(X^\nu_K),\qquad
 (X\cdot\tau_vY)^\nu_K=A\,\tau_vB .
\]
The concatenation is separated because
$|v|_\infty\ge4\bar R$.

The reflection size of the real function $A$ is the maximum,
over coordinate directions $\mu$ and coordinate sign changes
$\mathsf g$, of
\[
 \rho\bigl((\theta_\mu\mathsf gA)
                   (\tau_{2\bar Re_\mu}\mathsf gA)\bigr).
\]
Within either reflected tuple the separation is $d(X)$.
Between the two tuples it is at least $2\delta(X;\bar R)$,
since their $\mu$-coordinates lie respectively below
$\bar R-\delta$ and above $\bar R+\delta$.
The moment bound thus makes this reflection size at most
$N(X;\bar R)$, and the same argument applies to $Y$.
The coordinate and diagonal mirror hypotheses of
\cref{IV-31-51} hold for $\rho$.  Its estimate, with
$r a_K=\bar R$, is precisely \eqref{IV-35-24}.
\end{proof}

The estimates also give continuity under small translations of
individual entries, uniformly in $K,\rho,\nu$.  Indeed,
$\|\tau_u f-f\|_\sharp\to0$ as $u\to0$, because $\nabla f$
is uniformly continuous.  The analogous fact holds for
$\|\mathsf Rf-f\|_\sharp$ as $\mathsf R\to1$ in $\mathrm O(3)$:
on a fixed ball the chain rule gives
$\nabla(\mathsf Rf)(x)=\mathsf R\nabla f(\mathsf R^{-1}x)$.
Enlarge each support by one quarter of
$\min(1,d(X))$ and apply \eqref{IV-35-19}; the enlarged
supports have at least half the original separation.
This supplies a common local modulus of continuity for all
the translations used below.

\subsection{Limits at spatial infinity}

For a converging frame put
\[
 e_j(X)=E_{\nu_j}[A_j(X)\mid\mathcal I],\qquad
 m_j(X^1,\ldots,X^q)=\nu_j\!\left(\prod_{i=1}^q e_j(X^i)\right).
\]
Choose real versions with $|e_j(X)|\le C_{\rm mom}(X)$.
Versions do not affect $m_j$.  Limits below are iterated: the
shift $v_2$ tends to infinity first, with all later shifts fixed.

\begin{proposition}[Iterated limits]\label{IV-35-29}
For any separated tuples $X^1,\ldots,X^q$, not necessarily
separated from one another, $m_j(X^1,\ldots,X^q)$ converges and
\begin{equation}\label{IV-35-35}
 \lim_j m_j(X^1,\ldots,X^q)
 =\lim_{|v_q|\to\infty}\cdots\lim_{|v_2|\to\infty}
       \sigma(X^1\cdot\tau_{v_2}X^2\cdots\tau_{v_q}X^q).
\end{equation}
Every intermediate limit is continuous in the remaining shifts
on the open set where the remaining concatenation is separated.

More precisely, put
$X'_p=X^1\cdot\tau_{v_{p+1}}X^{p+1}\cdots\tau_{v_q}X^q$
and denote the intermediate limits by $G^{(p)}$, so that
$G^{(1)}$ is the unmodified correlation.  If
$\bar R'>0$ belongs to $\bigcup_Ka_K\mathbb Z$ and
$\operatorname{supp}X'_p\cup\operatorname{supp}X^p
 \subset(-\bar R',\bar R')^3$, then
\begin{equation}\label{IV-35-30}
 |G^{(p-1)}(v_p,\ldots,v_q)-G^{(p)}(v_{p+1},\ldots,v_q)|
 \le k_0\prod_{i=2}^{p-1}C_{\rm mom}(X^i)
       \sqrt{N(X'_p;\bar R')N(X^p;\bar R')}
                         \frac{\bar R'}{|v_p|}
\end{equation}
for $|v_p|\ge k_1\bar R'$.
\end{proposition}

\begin{proof}
Define
\[
 G_j^{(p)}=\nu_j\!\left(\prod_{i=2}^p e_j(X^i)\,A_j(X'_p)\right).
\]
For $p=1$ these converge by the definition of the frame.
For $p=q$ conditional expectation gives
$G_j^{(q)}=m_j(X^1,\ldots,X^q)$.

A real invariant weight $w$ with $|w|\le\bar w$ gives
\begin{equation}\label{IV-35-32}
 \nu_j(wF)=c_+\rho_+(F)-c_-\rho_-(F),\qquad
 c_\pm=\nu_j(w_\pm),\quad c_++c_-\le\bar w .
\end{equation}
Terms with $c_\pm=0$ are omitted; otherwise
$\rho_\pm=(\nu_j)_{w_\pm}\in\mathcal G_j$.
The conditional expectations $e_j(X)$ remain versions for
$\rho_\pm$.  Apply \eqref{IV-35-24} to these two laws, centred
at $\nu_j$, with $w=\prod_{i=2}^{p-1}e_j(X^i)$.
It bounds $G_j^{(p-1)}-G_j^{(p)}$ by the right side of
\eqref{IV-35-30}, whenever the shifts and the cube side
belong to the lattice and $j$ is sufficiently large.
Using \eqref{IV-35-19} in the same signed decomposition
also gives a local modulus of continuity of $G_j^{(p)}$
in the remaining shifts, independent of $j$.

Assume $G_j^{(p-1)}$ converges.  Fix the later shifts and
$\bar R'$.  A vector $v_p$ in
$(\bigcup_Ka_K\mathbb Z)^3$ belongs to every sufficiently
fine lattice.  Passing to upper and lower limits in the
preceding estimate bounds their difference for
$G_j^{(p)}$ by twice its right side.  Letting
$v_p=ze_1$, $z\in\mathbb N$, tend to infinity makes this
difference zero.  Thus $G_j^{(p)}$ converges.  The common
modulus proves continuity of its limit and extends the
estimate from the dense set of lattice-compatible shifts
to all real shifts.  The right side tends to zero as
$|v_p|\to\infty$, identifying the resulting limit as
$G^{(p)}$.  Induction proves every assertion, including
\eqref{IV-35-35} at $p=q$.
\end{proof}

The constants in \eqref{IV-35-30} depend on the later shifts
through $X'_p$ and $\bar R'$.  No joint limit with all shifts
tending to infinity is asserted.

\subsection{Construction of the component measure}

A decomposition of the following kind is determined by the
values at spatial infinity.  Its space $\mathsf X$ is compact
and metrizable, its probability measure $\lambda$ has full
support, and each component $\sigma_\xi$ is a normalized,
real, symmetric, multilinear functional on separated tuples.
Require that each evaluation $\xi\mapsto\sigma_\xi(X)$ be
continuous, that distinct points give distinct functionals,
and that
\begin{equation}\label{eq:component-plateau-moments}
 \int_{\mathsf X}\prod_{i=1}^q\sigma_\xi(X^i)\,\lambda(d\xi)
 =\lim_{|v_q|\to\infty}\cdots\lim_{|v_2|\to\infty}
       \sigma(X^1\cdot\tau_{v_2}X^2\cdots\tau_{v_q}X^q).
\end{equation}
For $q=1$ this is the decomposition of $\sigma$ itself.

\begin{theorem}[Existence of the component measure]\label{IV-35-36}
Every converging frame in classes satisfying
\eqref{IV-35-2}--\eqref{IV-35-4} has such a decomposition.
Each component is a limit in the same frame.  More precisely,
for every $\xi\in\mathsf X$ there are invariant events $E_j$
with $\nu_j(E_j)>0$ such that $\rho_j=(\nu_j)_{1_{E_j}}$ obeys
\begin{equation}\label{IV-35-37}
 \|E_{\rho_j}[A_j(X)\mid\mathcal I]-\sigma_\xi(X)\|_{L^\infty(\rho_j)}
       \longrightarrow0
 \quad\text{for every separated }X.
\end{equation}
Every component satisfies $|\sigma_\xi(X)|\le C_{\rm mom}(X)$
and the difference estimate \eqref{IV-35-19}.
Moreover
\[
 \#\mathsf X\le
 \liminf_j\dim L^2(\mathcal V(\mathbb L_{K_j}),\mathcal I,\nu_j),
\]
with infinity allowed on either side.
\end{theorem}

\begin{proof}
Let $\mathfrak F=\{X_1,X_2,\ldots\}$ be the countable family
with controlled supports constructed in
\cref{sec:separated-compactness}.  For every nonempty tuple
$X=(f_1,\ldots,f_n)$ and $0<\varepsilon\le1$, choose
$X^{[\varepsilon]}\in\mathfrak F$ whose entries differ from
those of $X$ by at most $\varepsilon$ in $\|\cdot\|_\sharp$.
The construction permits both supports of each entry to lie
in a set contained in a ball of radius $R(X)+1/4$, with
different sets at distance at least
$\frac12\min(1,d(X))$.  Indeed, take the open neighbourhoods
within distance $\frac14\min(1,d(X))$ of the original
supports and use their countable dense test sets.
The estimate \eqref{IV-35-19}, also for conditional
expectations, then gives the error
\begin{equation}\label{IV-35-38}
 \varepsilon_X(\varepsilon)=
 c_n\!\left(\tfrac12\min(1,d(X));R(X)+\tfrac14\right)
       n\varepsilon\prod_i(\|f_i\|_\sharp+1)
       \longrightarrow0 .
\end{equation}
Choose the approximation using the fixed enumeration, so
that it is independent of the law and of $\xi$.

Consider the compact metrizable product
\[
 \Pi=\prod_{X\in\mathfrak F}[-C_{\rm mom}(X),C_{\rm mom}(X)].
\]
The map $y_j=(e_j(X))_{X\in\mathfrak F}$ is
$\mathcal I$-measurable, and
$\lambda_j=(y_j)_*\nu_j$ is a probability measure on $\Pi$.
By \cref{IV-35-29}, the integral against $\lambda_j$ of
every monomial in the coordinates converges.  Coordinate
polynomials contain the constants and separate points;
the real Stone--Weierstrass theorem makes them uniformly
dense in $C(\Pi)$.  Consequently $\lambda_j(\phi)$ converges
for every $\phi\in C(\Pi)$: approximate $\phi$ by a polynomial
and use $|\lambda_j(\phi)|\le\|\phi\|_\infty$.
The limiting positive functional, equal to one on constants,
is represented by a probability measure $\lambda$.
Let $\mathsf X=\operatorname{supp}\lambda$.
The countable base of $\Pi$ shows $\lambda(\mathsf X)=1$,
so this is a compact metrizable space with full support.

For every open $V\subset\Pi$,
$\liminf_j\lambda_j(V)\ge\lambda(V)$.  One may verify this
directly by increasing the continuous functions
$\min(1,s\operatorname{dist}(\,\cdot\,,\Pi\setminus V))$
to $1_V$.  Fix $\xi\in\mathsf X$ and set
\[
 V_s=\{\zeta:|\zeta_{X_k}-\xi_{X_k}|<1/s,\quad 1\le k\le s\}.
\]
Then $\lambda(V_s)>0$, and hence $\lambda_j(V_s)>0$
for all sufficiently large $j$.  Choose increasing
$j_s$ with this property for every $j\ge j_s$.
For $j_s\le j<j_{s+1}$ put $E_j=y_j^{-1}(V_s)$;
for the finitely many earlier indices use the whole space.
These invariant events have positive measure, and
$\rho_j=(\nu_j)_{1_{E_j}}\in\mathcal G_j$.
The unchanged versions of conditional expectation give
\[
 \|E_{\rho_j}[A_j(X_k)\mid\mathcal I]-\xi_{X_k}\|_\infty
 \le1/s\qquad(k\le s,\ j_s\le j<j_{s+1}).
\]

For an arbitrary nonempty $X$, comparison with
$X^{[\varepsilon]}$ now yields
\[
 \limsup_j\|E_{\rho_j}[A_j(X)\mid\mathcal I]
                    -\xi_{X^{[\varepsilon]}}\|_\infty
 \le\varepsilon_X(\varepsilon).
\]
In particular the numbers $\xi_{X^{[\varepsilon]}}$
are Cauchy as $\varepsilon\downarrow0$, with the difference
between two of them at most the sum of the two errors.
Define their limit to be $\sigma_\xi(X)$ and set
$\sigma_\xi(\varnothing)=1$.  Then
\[
 |\sigma_\xi(X)-\xi_{X^{[\varepsilon]}}|
 \le\varepsilon_X(\varepsilon).
\]
This bound and the preceding limsup prove \eqref{IV-35-37}.
Integration shows $\rho_j(A_j(X))\to\sigma_\xi(X)$.
It follows that $\sigma_\xi$ is a functional of the asserted
kind and inherits the moment and difference estimates.
For $X\in\mathfrak F$ it equals the coordinate $\xi_X$.

The last approximation is uniform in $\xi$.
Thus every evaluation $\sigma_\xi(X)$ is continuous, as a
uniform limit of coordinates.  Equality of all evaluations
implies equality of all coordinates, so the component map
is injective.

For tuples in $\mathfrak F$, convergence of coordinate
monomials proves \eqref{eq:component-plateau-moments}.
For general tuples replace each $X^i$ by
$(X^i)^{[\varepsilon]}$.  Telescoping the product bounds
both the component integral error and the error in $m_j$
by
\[
 \sum_i\varepsilon_{X^i}(\varepsilon)
       \prod_{l\ne i}
        (C_{\rm mom}(X^l)+\varepsilon_{X^l}(\varepsilon)).
\]
This tends to zero independently of $j$.
Empty tuples contribute a factor one.  Passing first in
$j$ and then in $\varepsilon$ proves the required identity.

Finally, if the displayed liminf dimension is a finite
integer $N$, it equals $N$ on infinitely many indices.
If $\mathsf X$ had $N+1$ distinct points, choose disjoint
open neighbourhoods of them in $\Pi$.  Each has positive
$\lambda$-measure and therefore positive $\lambda_j$-measure
for all sufficiently large $j$.  Their inverse images under
$y_j$ give $N+1$ nonzero orthogonal invariant indicator
functions, contradicting dimension $N$ on those indices.
\end{proof}

Every frame has a converging subsequence: the moment bound
and \cref{IV-10-25}, with $n_0=1$, give simultaneous
convergence on all nonempty tuples, and the empty tuple
has value one.  The component construction itself uses the
full sequence of a converging frame.

\section{Clustering components and the vacuum space}
\label{sec:component-properties}

We first show that the preceding decomposition is intrinsic to the
limit functional.  Reflection positivity will then force the
symmetries of each component, and the uniform estimate
\eqref{IV-35-24} will give clustering for every component.

\begin{proposition}[Uniqueness of the decomposition]\label{IV-36-9}
Two decompositions satisfying
\eqref{eq:component-plateau-moments} and the stated compactness,
continuity, injectivity and full-support conditions are related
by a unique measure-preserving homeomorphism which identifies
their component functionals.

If an isometry $\mathsf r$ preserves $\sigma$, it induces a unique
measure-preserving homeomorphism $h_{\mathsf r}$ of $\mathsf X$ with
\begin{equation}\label{IV-36-10}
 \sigma_{h_{\mathsf r}(\xi)}=\sigma_\xi\circ\mathsf r^{-1},
 \qquad h_{\mathsf r\mathsf r'}=h_{\mathsf r}h_{\mathsf r'}.
\end{equation}
For $s>0$ define
\begin{equation}\label{IV-36-4}
 (\operatorname{dil}_s\sigma)(f_1,\ldots,f_n)
   =s^{-n}\sigma(f_1(s\,\cdot),\ldots,f_n(s\,\cdot)).
\end{equation}
The dilated components on the same probability space decompose
$\operatorname{dil}_s\sigma$.
\end{proposition}

\begin{proof}
For any finite collection of tuples, the two component measures
have the same moments in their evaluation coordinates, by
\eqref{eq:component-plateau-moments}, allowing repeated tuples.
Both coordinate images lie in a compact box.  Polynomial
approximation and uniqueness of the representing measure show
that their finite-dimensional distributions agree.

Embed the two component spaces in
$\mathbb R^{\mathfrak S(\mathbb R^3)}$ by all evaluations.
They are compact subsets of this Hausdorff product, and their
embedding maps are homeomorphisms onto their images.  Every
basic neighbourhood of a point in the first image is determined
by finitely many evaluations.  Its inverse image is a nonempty
open set and hence has positive measure.  Equality of the
finite-dimensional distributions makes it meet the second
image.  Since that image is closed, it contains the point.
Interchanging the two spaces proves equality of the images.
Their identification is the required homeomorphism; it is
unique by injectivity of the evaluations.  Polynomials in
these evaluations are uniformly dense on either component
space, so equality of their integrals also proves preservation
of the measures.

For an isometry, the functions
$\sigma_\xi(\mathsf r^{-1}X)$ satisfy the decomposition identity:
conjugating translations changes $v$ by the orthogonal linear
part of $\mathsf r$, and therefore preserves the iterated
limits $|v|\to\infty$.  Apply uniqueness to obtain
$h_{\mathsf r}$.  Composing its defining identities proves the
group law and $h_{\mathsf r^{-1}}=h_{\mathsf r}^{-1}$.
For dilation, $f(s\,\cdot)$ is translated by $v/s$ when $f$
is translated by $v$.  Each iterated limit is thus rescaled
by $s^{-\sum_i|X^i|}$.  This is exactly the product of the
factors in \eqref{IV-36-4}.  Continuity and injectivity are
preserved because dilation of test tuples is invertible.
\end{proof}

\begin{proposition}[Reflection invariance of the components]
\label{IV-36-29}\label{IV-36-31}
Suppose a real symmetric functional $\sigma$ is translation
invariant, is reflection positive at a mirror $\Sigma=(u,c)$,
and has a decomposition as above.  Then $\sigma$ and every
component are invariant under $\theta_\Sigma$.  Each component
is also invariant under
$\mathsf g\theta_\Sigma\mathsf g^{-1}$ whenever the isometry
$\mathsf g$ preserves $\sigma$.

Consequently a Euclidean-invariant $\sigma$ which is reflection
positive at one mirror has Euclidean-invariant components.
\end{proposition}

\begin{proof}
On the positive half-space, positivity for the two formal
vectors $\varnothing,X$ implies
$\sigma(\theta_\Sigma X)=\sigma(X)$: use coefficients $1,i$
and equate imaginary parts.  Move any tuple into this
half-space by a translation parallel to $u$ and use
$\theta_\Sigma\tau_{tu}=\tau_{-tu}\theta_\Sigma$.
Translation invariance proves the same identity for all tuples.

Let $F$ be a finite formal sum of tuples and put
$F_t=\tau_{tu/2}F$.  For large $t$, $F_t$ lies in the positive
half-space.  Translation invariance and the case $q=2$ of
the decomposition identity give
\begin{equation}\label{IV-36-30}
 0\le\lim_{t\to\infty}\sigma(\theta_\Sigma\overline{F_t}\cdot F_t)
   =\int_{\mathsf X}
       \overline{\sigma_\xi(\theta_\Sigma F)}
                         \sigma_\xi(F)\,\lambda(d\xi).
\end{equation}
Indeed, term by term the expression before the limit is
$\sigma(\theta_\Sigma\bar F\cdot\tau_{tu}F)$, whose large
translation limit is prescribed by
\eqref{eq:component-plateau-moments}.
Replace $F$ by $F-\theta_\Sigma F$.  The last integral is
then
\[
 -\int_{\mathsf X}
       |\sigma_\xi(F-\theta_\Sigma F)|^2\,\lambda(d\xi).
\]
It must vanish.  The integrand is continuous, and $\lambda$
has full support; hence every component is fixed by
$\theta_\Sigma$.

Equivalently $h_{\theta_\Sigma}$ in \eqref{IV-36-10} is the
identity.  Its conjugate by $h_{\mathsf g}$ is also the
identity, proving the assertion for the conjugated mirror.
If $\sigma$ is Euclidean invariant, every mirror is such a
conjugate.  Products of two parallel reflections give
translations.  Every orthogonal map is a product of at most
three reflections: reflect its image of a basis vector back
to that vector, preserving the previously fixed vectors, and
repeat.  Thus every component is Euclidean invariant.
\end{proof}

\subsection{Properties inherited from the frame}

The following conclusions apply to every component, including
points of measure zero in the component space.  Their lattice
approximants are the conditioned laws in \eqref{IV-35-37}.

\begin{proposition}\label{IV-36-12}\label{IV-36-13}
Every limit $\sigma'$ in a frame satisfying
\eqref{IV-35-2}--\eqref{IV-35-4} is invariant under translations
and cubic isometries, and is reflection positive at every
coordinate and diagonal mirror.  Its cumulants of order at
least two are the limits of the cumulants of its integrating
laws and obey \eqref{IV-35-3}.  Its one-point function is
\begin{equation}\label{IV-36-14}
 \sigma'(f)=c_{\sigma'}\int f,\qquad
 c_{\sigma'}=\lim_j3a_{K_j}^{-4}
                      \bigl(\rho_j(a_p)-\nu_j(a_p)\bigr),
 \qquad |c_{\sigma'}|\le\bar c,
\end{equation}
where
\[
 \bar c=C_1(1)
   \inf\{\|f\|_\sharp:f\in C_c^1(B(0,1)),\ \int f=1\}.
\]
Its off-diagonal distributions have order at most $4n+1$
on compact sets.  The coordinate reflection Hilbert spaces
and their translation semigroups are therefore defined.

In the decomposition of a converging frame,
$\xi\mapsto c_\xi$ is continuous,
$|c_\xi|\le\bar c$, and
\begin{equation}\label{IV-36-15}
 \sigma(f)=0,\qquad\int_{\mathsf X}c_\xi\,\lambda(d\xi)=0.
\end{equation}
If $\kappa=\int c_\xi^2\,d\lambda$, then for disjoint $f,f'$,
\begin{equation}\label{IV-36-16}
 S_2^\sigma(f,f')
  =\int_{\mathsf X}S_2^{\sigma_\xi}(f,f')\,\lambda(d\xi)
                   +\kappa\int f\int f' .
\end{equation}
\end{proposition}

\begin{proof}
All moment and difference estimates pass to the limit by
\cref{IV-35-17}.  The moment--cumulant polynomial identifies
the limiting cumulants with those of $\sigma'$.
Translations on the lattice-compatible dense subgroup pass
directly to the limit; the common continuity estimate extends
them to all translations.  Cubic symmetries pass directly.
The coordinate and diagonal reflection arguments in
\cref{sec:continuum-subsequences} apply to the integrating
laws $\rho_j$, with the invariant centring laws $\nu_j$.
They give positivity at all the stated mirrors.

The same lattice symmetry makes
$\rho_j(a_p)-\nu_j(a_p)$ independent of the plaquette.
The barycentre Riemann sum used in
\cref{sec:continuum-subsequences} expresses the one-point
function as this difference times
$3a_{K_j}^{-4}r_j(f)$, where $r_j(f)\to\int f$.
Taking a unit-integral test first proves convergence of
the coefficient; every other test then gives
\eqref{IV-36-14}.  Apply \eqref{IV-35-4} to tests in the unit
ball and take the infimum to obtain $\bar c$.
The distribution construction of
\cref{sec:distribution-extension} uses precisely these
cumulant and one-point bounds, so its order estimate
applies.  Translation invariance and coordinate positivity
supply the hypotheses of
\cref{sec:reflection-hilbert}.

For a unit-integral $f_0$, $c_\xi=\sigma_\xi(f_0)$ is a
continuous evaluation.  Every one-point moment of the
frame is zero by its own centring, and the decomposition
identity gives \eqref{IV-36-15}.  Finally,
\[
 S_2^{\sigma_\xi}(f,f')
   =\sigma_\xi(f,f')-c_\xi^2\int f\int f'.
\]
Integrating this equality gives \eqref{IV-36-16}.
\end{proof}

\begin{proposition}[Uniform clustering of the components]\label{IV-36-17}
For the decomposition of a converging frame, let
$\bar R>0$ belong to $\bigcup_Ka_K\mathbb Z$ and contain
$\operatorname{supp}X\cup\operatorname{supp}Y$ in its open
cube.  Then, for every $\xi\in\mathsf X$ and
$|v|\ge k_1\bar R$,
\begin{equation}\label{IV-36-18}
 |\sigma_\xi(X\cdot\tau_vY)-\sigma_\xi(X)\sigma_\xi(Y)|
 \le k_0\sqrt{N(X;\bar R)N(Y;\bar R)}\,\frac{\bar R}{|v|}.
\end{equation}
The bound is uniform over all components.
\end{proposition}

\begin{proof}
Use the conditioned approximants $\rho_j$ for the given
component.  For a lattice-compatible $v$, all the hypotheses
of \eqref{IV-35-24} hold eventually.  Its first term converges
to $\sigma_\xi(X\cdot\tau_vY)$.
By \eqref{IV-35-37}, the conditional expectations of
$A_j(X)$ and $A_j(Y)$ converge in their respective
$L^\infty(\rho_j)$ norms to the constants
$\sigma_\xi(X)$ and $\sigma_\xi(Y)$.  If their errors are
$\varepsilon_j(X)$ and $\varepsilon_j(Y)$, the error in
the integral of their product is at most
\[
 \varepsilon_j(X)(|\sigma_\xi(Y)|+\varepsilon_j(Y))
       +|\sigma_\xi(X)|\varepsilon_j(Y)\longrightarrow0 .
\]
Thus \eqref{IV-36-18} follows on the dense subgroup.
Approximate any remaining $v$ by such vectors with
$|v_j|\ge k_1\bar R$; the continuity under translations
already proved gives the assertion.  The right side
contains no component-dependent quantity.
\end{proof}

\subsection{Zero energy and the component space}

For a translation-invariant, reflection-positive functional
$\sigma$ in the preceding distribution class, let
$\mathcal H(\sigma,u)$ be its reflection Hilbert space,
$\Omega=[\varnothing]$, $T_t=e^{-tH}$, and
$P_0=1_{\{0\}}(H)$.
The spectral theorem and $H\ge0$ give $T_t\to P_0$ strongly.
Consequently
\begin{equation}\label{IV-36-21}
 \lim_{t\to\infty}\sigma(\theta\bar F\cdot\tau_{tu}F')
       =\langle[F],P_0[F']\rangle .
\end{equation}
Here $F,F'$ are finite sums of tuples in the positive
half-space.  Also
$\langle\Omega,[F]\rangle=\sigma(F)$ and
$\sigma(\theta\bar F)=\overline{\sigma(F)}$.
Since these classes are dense,
\begin{equation}\label{IV-36-22}
 \ker H=\mathbb C\Omega
 \quad\Longleftrightarrow\quad
 \lim_{t\to\infty}\sigma(\theta\bar F\cdot\tau_{tu}F')
            =\overline{\sigma(F)}\sigma(F')
 \quad\text{for all }F,F'.
\end{equation}
Indeed, the right side identifies the matrix elements of
$P_0$ with those of the orthogonal projection onto
$\mathbb C\Omega$.

Applied to each component of a frame, \eqref{IV-36-18}
and its reflection invariance give the right side of
\eqref{IV-36-22}.  Thus each coordinate Hamiltonian has
a one-dimensional kernel.

\begin{theorem}[The vacuum space of the mixture]\label{IV-36-25}
Suppose $\sigma$ has a decomposition satisfying
\eqref{eq:component-plateau-moments}.  Assume its components
are translation invariant and invariant under the reflection
$\theta$ defining $\mathcal H(\sigma,u)$, and that
\[
 \sup_\xi|\sigma_\xi(X\cdot\tau_{tu}Y)
                     -\sigma_\xi(X)\sigma_\xi(Y)|\longrightarrow0
 \quad(t\to\infty)
\]
for positive-half-space tuples $X,Y$.  Then
\begin{equation}\label{IV-36-26}
 \langle P_0[F],P_0[F']\rangle
    =\int_{\mathsf X}\overline{\sigma_\xi(F)}
                             \sigma_\xi(F')\,\lambda(d\xi).
\end{equation}
The assignment
$P_0[F]\mapsto(\xi\mapsto\sigma_\xi(F))$ extends to a unitary
map $\ker H\to L^2(\mathsf X,\lambda)$, carrying $\Omega$
to $1$.  The number of components equals $\dim\ker H$
when either is finite; if $\mathsf X$ is infinite, the
kernel is a separable infinite-dimensional Hilbert space.

These conclusions apply to every coordinate reflection
Hilbert space of the limit of a converging frame.
\end{theorem}

\begin{proof}
For tuples $X,Y$ in the positive half-space, the two-factor
decomposition identity and reflection invariance give
\[
 \lim_{t\to\infty}\sigma(\theta X\cdot\tau_{tu}Y)
      =\int_{\mathsf X}\sigma_\xi(X)\sigma_\xi(Y)\,d\lambda .
\]
Combine this with \eqref{IV-36-21} and extend sesquilinearly
to obtain \eqref{IV-36-26}.
It makes the proposed map well defined and isometric.
The vectors $P_0[F]$ are dense in $\ker H$, so it extends
to an isometry with closed range, and sends $\Omega$ to $1$.

To prove surjectivity, put
$\psi_F(\xi)=\sigma_\xi(F)$ and take the uniform closure of
their span in $C(\mathsf X)$, initially for positive-half-space
$F$.  It contains constants and is closed under conjugation.
The uniform clustering assumption gives
\[
 \psi_{X\cdot\tau_{tu}Y}\longrightarrow\psi_X\psi_Y
          \quad\text{uniformly on }\mathsf X .
\]
Thus the closure is an algebra.  It separates points:
any tuple distinguishing two components may be translated
into the positive half-space without changing its values.
Stone--Weierstrass makes the closure $C(\mathsf X)$.
Since continuous functions are dense in $L^2(\lambda)$,
the isometry has dense as well as closed range.

Full support implies that disjoint nonempty open sets give
nonzero orthogonal indicator functions in $L^2(\lambda)$.
Any $n$ distinct points have $n$ such neighbourhoods.
Hence a finite component space has exactly that dimension,
and an infinite component space has infinite dimension.
Separability follows from that of the reflection Hilbert
space.  For frame limits, \cref{IV-36-12,IV-36-17} supply
both invariance and uniform clustering required above.
\end{proof}

\section{Two-layer exchange and comparison states}
\label{sec:gibbs-two-layer}

The comparison argument for diagonal reflections requires more than
positivity of a finite-volume measure.  We first prove an exchange
identity for two neighbouring layers of every periodic Gibbs state.
The proof inserts a unitary operator in the transfer matrix and compares
the resulting partition function with a Gibbs expectation.  Strict
positivity of the Fourier coefficients of the plaquette weight then
turns this operator identity into an identity for arbitrary continuous
functions of the two layers.

\subsection{The plaquette kernel and its square root}

Throughout this subsection $G=\mathrm{SU}(2)$ has normalized Haar
measure and $\beta>0$.  Set
\[
 k_\beta(g)=\exp\{-\beta(1-\tfrac12\operatorname{tr}g)\},\qquad
 k_S(u,u')=\prod_{b\in S}k_\beta(u(b)u'(b)^{-1})
\]
for a finite set $S$.  We write $K_S$ for the integral operator with
kernel $k_S$ on $L^2(G^S)$.  Inner products are antilinear in the first
argument.

For $l\in\frac12\mathbb N_0$, let $D^l$ be the representation on the
homogeneous polynomials of degree $2l$ in two variables, acting by
$\phi(z)\mapsto\phi(zg)$, with an invariant inner product obtained by
averaging any positive inner product over $G$.  Choose an orthonormal
basis and write $D^l_{mn}$ for its matrix coefficients and
$\chi_l=\operatorname{tr}D^l$ for its character.  Put
$\chi_{-1/2}=0$.  The normalized products
\[
 e_\alpha(u)=\prod_{b\in S}(2l_b+1)^{1/2}D^{l_b}_{m_bn_b}(u(b))
\]
are indexed by $\alpha=(\underline l,\underline m,\underline n)$.
Their finite linear combinations form $\mathcal T(G^S)$.

\begin{lemma}\label{IV-29-4}\label{IV-29-6}
The coefficients satisfy
\begin{equation}\label{IV-29-5}
 \int_G D^l_{mn}(g)\overline{D^{l'}_{m'n'}(g)}\,dg
 =\frac{\delta_{ll'}\delta_{mm'}\delta_{nn'}}{2l+1}.
\end{equation}
The functions $e_\alpha$ are an orthonormal basis of $L^2(G^S)$.
The algebra $\mathcal T(G^S)$ is invariant under conjugation and
uniformly dense in $C(G^S)$.  Each fixed-spin coefficient space is
invariant under conjugation.  Moreover, $\chi_l$ is real,
$\chi_l(g^{-1})=\chi_l(g)$, $|\chi_l|\le2l+1$, and
\begin{equation}\label{IV-29-7}
 (\operatorname{tr}g)^N=\sum_l m_{N,l}\chi_l(g),\quad
 m_{N,l}\ge0,\quad m_{2l,l}=1,\quad
 \sum_l(2l+1)m_{N,l}=2^N,
\end{equation}
where $m_{N,l}=0$ for $2l>N$.
\end{lemma}

\begin{proof}
In the monomial basis $z_1^jz_2^{2l-j}$, a diagonal group element with
entries $e^{i\omega},e^{-i\omega}$ acts with weights $2j-2l$.
Taking $\omega=\pi/(2l+1)$ makes these eigenvalues distinct.  An
invariant subspace therefore contains every nonzero monomial component
of each of its elements, by polynomial spectral projections.  The
matrix
$2^{-1/2}\left(\begin{smallmatrix}1&1\\-1&1\end{smallmatrix}\right)$
sends any monomial to a polynomial with a nonzero $z_1^{2l}$
coefficient, and sends $z_1^{2l}$ to a polynomial with every monomial
coefficient nonzero.  Thus the representation is irreducible.
An intertwiner has invariant kernel and image; it is zero between
different dimensions, and is scalar in one irreducible representation
by subtracting an eigenvalue.  Averaging
$D^l(g)BD^{l'}(g)^{-1}$ and taking $B$ to be a matrix unit proves
\eqref{IV-29-5}.

Let $\mathcal P_N$ be the restrictions to $G$ of polynomials of degree
at most $N$ in its four matrix entries.  The coefficients with
$2l\le N$ span a subspace of dimension
$\sum_{j=0}^N(j+1)^2=(N+1)(N+2)(2N+3)/6$.
The identity $g_{11}g_{22}-g_{12}g_{21}=1$ expresses every monomial of
degree at most $N$ as a sum of monomials of degrees $N$ and $N-1$.
The dimensions of these two homogeneous polynomial spaces add to the
same number.  Consequently $\mathcal P_N$ is exactly the orthogonal
sum of these coefficient spaces.  Conjugating an entry of an
$\mathrm{SU}(2)$ matrix gives another entry, up to sign.  Hence
$\mathcal P_N$ is conjugation invariant, as is its orthogonal
complement to $\mathcal P_{N-1}$, the spin-$N/2$ space.
Products over $S$ now give a self-adjoint algebra which contains the
constants and separates points.  Stone--Weierstrass and
\eqref{IV-29-5} prove the asserted density and completeness.

Every $g\in G$ is conjugate in $G$ to
$\operatorname{diag}(e^{i\omega},e^{-i\omega})$, so
$\chi_l(g)=\sum_{j=0}^{2l}e^{i(2j-2l)\omega}$.
This proves the reality and bounds and gives
$\chi_{1/2}\chi_l=\chi_{l+1/2}+\chi_{l-1/2}$.
Starting with $m_{0,l}=\delta_{0l}$, the recursion
$m_{N+1,l}=m_{N,l-1/2}+m_{N,l+1/2}$ proves
\eqref{IV-29-7}; evaluation at the identity gives its last equality.
\end{proof}

\begin{lemma}\label{IV-29-11}\label{IV-29-15}\label{IV-29-16}
Define
\[
 \lambda_l=\frac1{2l+1}\int_Gk_\beta(g)\chi_l(g)\,dg,
 \qquad \lambda_\alpha=\prod_{b\in S}\lambda_{l_b}.
\]
Then
\begin{equation}\label{IV-29-13}
 \frac{e^{-\beta}(\beta/2)^{2l}}{(2l)!(2l+1)}
 \le\lambda_l\le\frac{\beta^{2l}}{(2l+1)^2(2l)!},
 \qquad \sum_l(2l+1)^2\lambda_l=1.
\end{equation}
The operator $K_S$ is positive, injective and trace class, with
$K_Se_\alpha=\lambda_\alpha e_\alpha$ and
$\operatorname{Tr}K_S=1$.  Its positive square root has the continuous
real symmetric kernel
\[
 q_S(u,u')=\prod_{b\in S}q_\beta(u(b)u'(b)^{-1}),\qquad
 q_\beta=\sum_l(2l+1)\sqrt{\lambda_l}\chi_l.
\]
Writing $q_u=q_S(u,\cdot)$, we have
\begin{equation}\label{IV-29-18}
 \langle q_u,q_{u'}\rangle=k_S(u,u'),\qquad \|q_u\|=1.
\end{equation}
Both $K_S$ and its square root map each fixed-spin space to itself
by a strictly positive scalar.  In particular, $K_S$ maps
$\mathcal T(G^S)$ onto itself.  If a finite complex measure $\sigma$
on $G^S\times Y$, with $Y$ compact metrizable, annihilates
$(K_S\varphi)(u)\psi(y)$ for every $\varphi\in\mathcal T(G^S)$
and a uniformly dense collection of $\psi\in C(Y)$, then $\sigma=0$.
\end{lemma}

\begin{proof}
Expanding the exponential and using \eqref{IV-29-7} gives
\[
 (2l+1)\lambda_l
 =e^{-\beta}\sum_{N\ge2l}\frac{(\beta/2)^N}{N!}m_{N,l}.
\]
The term $N=2l$ gives the lower bound.  For the upper bound use
$m_{N,l}\le2^N/(2l+1)$ and
$\sum_{N\ge j}\beta^N/N!\le\beta^j e^\beta/j!$.
The character expansion
$k_\beta=\sum_l(2l+1)\lambda_l\chi_l$ converges absolutely
uniformly, since the exponential expansion has a summable majorant.
At the identity it gives the sum rule in \eqref{IV-29-13}.
Schur orthogonality gives the stated eigenvalues of $K_S$; their sum
with multiplicities is $1$.  Positivity, injectivity and the trace
assertions follow.

The upper bound in \eqref{IV-29-13} gives
$\sum_l(2l+1)^2\sqrt{\lambda_l}<\infty$, by the ratio test.
Thus the series defining $q_\beta$ converges uniformly.  Applying
orthogonality once more shows that its product kernel has eigenvalues
$\sqrt{\lambda_\alpha}$, and convolution of two copies gives $k_S$.
This proves \eqref{IV-29-18}.  Finally, division by the nonzero
eigenvalues is legitimate on every finite coefficient sum.  The
assertion about measures follows from this surjectivity, density of
products in $C(G^S\times Y)$, and uniqueness of representing measures.
\end{proof}

For a finite bond scheme with vertex set $X$ and bond set $S$, put
$(w\cdot u)(b)=w(b_-)u(b)w(b_+)^{-1}$ and
\[
 (R_w\varphi)(u)=\varphi(w^{-1}\cdot u),\qquad
 P\varphi=\int R_w\varphi\,dw.
\]
Haar invariance shows that $R_w$ is unitary and $P$ is the orthogonal
projection onto gauge invariant functions.  Conjugation invariance
of $q_\beta$ gives $q_S(w\cdot u,w\cdot u')=q_S(u,u')$.
Consequently $K_S$ and $K_S^{1/2}$ commute with $P$ and every $R_w$.
If $C$ has continuous kernel $c$, then $CP$ has kernel
$\int c(u,w\cdot u')\,dw$.

\subsection{Operator insertions and the finite transfer formula}

Use the third direction as time.  A configuration is specified by
horizontal layers $u_t$ on the positive bonds of $a\mathbb Z^2$ and
vertical variables $w_t(y)$ between heights $ta$ and $(t+1)a$.
The vertical plaquette over a horizontal bond $b$ has weight
\begin{equation}\label{IV-29-23}
 k_\beta\bigl(u_t(b)[(w_t\cdot u_{t+1})(b)]^{-1}\bigr).
\end{equation}
Indeed, its holonomy is
$u_t(b)w_t(b_+)u_{t+1}(b)^{-1}w_t(b_-)^{-1}$.
Horizontal plaquette actions are unchanged by a planar gauge
transformation.

For a bounded operator $X$ on $L^2(G^S)$ define
\begin{equation}\label{IV-29-25}
 \kappa_X(u,u')=\langle q_u,Xq_{u'}\rangle,
 \qquad A_X(u,u')=\frac{\kappa_X(u,u')}{k_S(u,u')},
 \qquad
 A_X^{(t)}(U)=A_X(u_t|_S,(w_t\cdot u_{t+1})|_S).
\end{equation}
The last expression uses the horizontal copies of $S$ at heights $t$
and $t+1$ and the vertical bonds at their endpoints.  In particular,
it is a continuous local function supported in that closed two-layer
slab.

\begin{proposition}\label{IV-29-28}
The map $X\mapsto A_X$ is linear, $A_1=1$, and
\[
 \|A_X\|_\infty\le e^{2\beta\#S}\|X\|.
\]
It preserves reality, tensor products on disjoint bond sets, and
relabeling of bonds.  The function $\kappa_X$ is a continuous kernel
of $K_S^{1/2}XK_S^{1/2}$.

Let $n_1,n_2,s\ge3$ be integers and
$\Gamma=a(n_1\mathbb Z\times n_2\mathbb Z\times s\mathbb Z)$.
Let $S_0$ and $X_0$ be the horizontal bonds and vertices modulo the
first two periods.  On $L^2(G^{S_0})$, let $E$ multiply by
$e^{-\beta V}$, where $V$ is the sum of the horizontal plaquette
actions in one layer, and set
\[
 Q=K_{S_0}^{1/2},\qquad M=QPEQ.
\]
Then $M\ge0$, $M$ is trace class and $Z_\Gamma[0]=\operatorname{Tr}M^s$.
At each height $t$, insert operators $X_{t,j}$ on finitely many
horizontal bond sets $S_{t,j}$. Require that the map
$(j,b)\mapsto[b]$ from the disjoint union of the $S_{t,j}$ to $S_0$
is injective. Relabel each operator by the resulting bijection onto
its image in $S_0$, and let $X_t$ be their tensor product with the
identity on the remaining bonds. Then
\begin{equation}\label{IV-29-30}
 \mu_{\Gamma,\beta}\left(\prod_{t=1}^s\prod_j A_{X_{t,j}}^{(t)}\right)
 =\frac{\operatorname{Tr}(X_1M\cdots X_sM)}{\operatorname{Tr}M^s},
 \qquad
 \left|\mu_{\Gamma,\beta}\left(\prod_{t,j}A_{X_{t,j}}^{(t)}\right)\right|
 \le\prod_{t,j}\|X_{t,j}\|.
\end{equation}
Heights in this formula are understood modulo $s$.
\end{proposition}

\begin{proof}
The map $u\mapsto q_u$ is continuous in $L^2$, by uniform continuity
of the kernel.  Hence $\kappa_X$ is continuous and
$|\kappa_X|\le\|X\|$.  Since $k_S\ge e^{-2\beta\#S}$, the norm
bound follows.  Applying the root kernel on both sides of $X$ gives
the kernel assertion, first against continuous test functions and
then in $L^2$.  Formula \eqref{IV-29-18} gives $A_1=1$.
Product Haar measure and the product formula for $q_S$ prove the tensor
and relabeling assertions.  They also prove preservation of reality.

The torus Haar measure is the product of the measures $du_t\,dw_t$,
and \eqref{IV-29-23} gives
\[
 e^{-\beta S_\Gamma(U)}
 =\prod_{t=1}^s e^{-\beta V(u_t)}k_{S_0}(u_t,w_t\cdot u_{t+1}),
 \qquad u_{s+1}=u_1.
\]
Multiplication by the insertions replaces the $t$th kernel by
$\kappa_{X_t}$.  Integrating $w_t$ gives the continuous kernel of
$F_t=EQX_tQP$.  Integrating the remaining cyclic product of kernels
therefore gives $\operatorname{Tr}(F_1\cdots F_s)$.
Each factorization needed for cyclicity contains two Hilbert--Schmidt
factors $Q$; moving $EQ$ from the beginning to the end yields
\[
 Z_\Gamma[0]\mu_{\Gamma,\beta}\left(\prod_{t,j}A_{X_{t,j}}^{(t)}\right)
 =\operatorname{Tr}(X_1M\cdots X_sM).
\]
Since $V$ is gauge invariant, $P$ commutes with $E^{1/2}$ and
$M=(E^{1/2}PQ)^*(E^{1/2}PQ)$ is positive trace class.  Empty
insertions give $Z_\Gamma[0]=\operatorname{Tr}M^s>0$.

For completeness, if $L\ge0$ is trace class and $B_1,\ldots,B_s$
are bounded, the trace H\"older inequality gives
\begin{equation}\label{IV-29-20}
 |\operatorname{Tr}(B_1L\cdots B_sL)|
 \le\prod_t\|B_t\|\operatorname{Tr}L^s.
\end{equation}
Indeed, apply Lemma~\ref{IV-3-35} to the $2s$ factors
$B_tL^{1/2},L^{1/2}$.  For positive trace class $A\le B$,
$\operatorname{Tr}A^s\le\operatorname{Tr}B^s$: in an eigenbasis
of $A$, bound each eigenvalue by $\langle e,Be\rangle$, and then
use scalar Jensen in an eigenbasis of $B$.
Applying this to
$L^{1/2}B_t^*B_tL^{1/2}\le\|B_t\|^2L$ bounds the Schatten
norm of $B_tL^{1/2}$ by
$\|B_t\|(\operatorname{Tr}L^s)^{1/(2s)}$.
This proves \eqref{IV-29-20}, and hence the second part of
\eqref{IV-29-30}.
\end{proof}

\subsection{Exchange of two neighbouring layers}

\begin{proposition}\label{IV-29-36}
Let $a>0$, $\beta>0$, and let $\nu$ be a periodic Gibbs state.
For a finite horizontal bond set $S$, an integer $t$, and every
$H\in C(G^S\times G^S)$,
\[
 \nu H(u_t|_S,(w_t\cdot u_{t+1})|_S)
 =\nu H((w_t\cdot u_{t+1})|_S,u_t|_S).
\]
More precisely, if $D$ is a bounded real operator on $L^2(G^S)$
with transpose $D^{\mathsf T}=-D$, then $\nu(A_D^{(t)})=0$.
\end{proposition}

\begin{proof}
For a real operator, $D^{\mathsf T}=D^*$, so
$X_\varepsilon=e^{\varepsilon D}$ is real unitary for real
$\varepsilon$.  Its remainder after $1+\varepsilon D$ has norm at
most $\frac12\varepsilon^2\|D\|^2e^{|\varepsilon|\|D\|}$.
By Proposition~\ref{IV-29-28}, for sufficiently small
$|\varepsilon|$ the continuous local insertion
$A_{X_\varepsilon}^{(t)}$ lies in $[1/2,3/2]$.  Set
$A=\log A_{X_\varepsilon}^{(t)}$.  Then
\begin{equation}\label{IV-29-37}
 \|A\|_\infty\le\log2,
 \qquad \|A-\varepsilon A_D^{(t)}\|_\infty\le C_D\varepsilon^2.
\end{equation}
Choose a period number $N_0$ of $\nu$ and a multiple $N'\ge3$ of
$N_0$ large enough that distinct translates of $S$ by $N'a\mathbb Z^2$
are disjoint.  For each $i\ge1$, take
$\Gamma' =N'a\mathbb Z^3$, $\Gamma=iN'a\mathbb Z^3$ and the
array $\mathbb A$ of $A$ over $\Gamma'/\Gamma$.
On the torus, $e^{\mathbb A}$ is a product of unitary insertions:
at heights $t+N'k_3$ use the $i^2$ translates by
$N'a(k_1,k_2)$, $0\le k_1,k_2<i$.  Their bond sets are pairwise
incongruent modulo the horizontal periods: a congruence would give
a nonzero $N'a\mathbb Z^2$ translation intersecting $S$, or equal
representatives.  Thus \eqref{IV-29-30} gives
\[
 \log\frac{Z_\Gamma[\mathbb A]}{Z_\Gamma[0]}\le0.
\]
Enclose the support of $A$ in a lattice cube of side $ra$.  The lower
tangent inequality~\ref{IV-14-20}, applied on the coordinate domain
of side $iN'a$, gives
\[
 0\ge i^3\nu(A)
 -\bigl(24\beta+2\log2(2r-1)^3\bigr)\#\partial_1D_i.
\]
Since $\#\partial_1D_i\le6i^2N'^2$, division by $i^3$ and
$i\to\infty$ yield $\nu(A)\le0$.
Using \eqref{IV-29-37} for both signs of $\varepsilon$ and then
letting $\varepsilon\to0$ proves $\nu(A_D^{(t)})=0$.

It remains to recover arbitrary functions of the layers.  Let $\Phi$
be a real antisymmetric trigonometric polynomial on $G^S\times G^S$,
and write its finite expansion as
\[
 \Phi(u,u')=\sum_{\alpha,\alpha'}c_{\alpha\alpha'}
 e_\alpha(u)\overline{e_{\alpha'}(u')}.
\]
Divide its coefficients by
$\sqrt{\lambda_\alpha\lambda_{\alpha'}}$ to obtain $\Phi'$.
This operation preserves reality and antisymmetry: each fixed pair
of spin spaces is conjugation invariant, and swapping the two spaces
does not change the positive divisor.  The finite-rank operator $D$
with kernel $\Phi'$ is therefore real and skew-adjoint.  Applying
$K_S^{1/2}$ in both variables gives $\kappa_D=\Phi$, and hence
\[
 \nu\left[\frac{\Phi}{k_S}
 (u_t|_S,(w_t\cdot u_{t+1})|_S)\right]=0.
\]
For real continuous $H$, the function
$k_S(u,u')[H(u,u')-H(u',u)]$ is real and antisymmetric.
Approximate it uniformly by trigonometric polynomials, then take real
parts and antisymmetrize.  Division by $k_S\ge e^{-2\beta\#S}$
preserves uniform convergence, proving the desired exchange identity
for real $H$.  Complex linearity finishes the proof.
\end{proof}

\subsection{Comparison with a state positive at a prescribed mirror}

\begin{proposition}\label{IV-29-40}
Let $a>0$, $\beta\ge0$, $n\in\mathfrak N$ and
$c\in\operatorname{Adm}(n)$.  For every periodic Gibbs state $\nu$
and every real $A\in\mathcal C_{\mathrm{loc}}(\mathbb H^+(n,c))$
there exists a periodic Gibbs state $\nu_0$, reflection positive at
$\Sigma(n,c)$, such that $\nu(A)\le\nu_0(A)$.
\end{proposition}

\begin{proof}
Choose a period number $N_0$ of $\nu$, a nonempty finite bond support
$B_A\subset\mathsf{Bd}(\mathbb H^+(n,c))$, and a lattice cube
$Q=x_0+[0,ra]^3$ containing its bonds.  Choose
$h\in2N_0a\mathbb N$ so that $B_A$ lies in
$\mathbb S_1=\{c\le n\cdot x\le c+h\}$, and put
$v_h=hn/|n|^2$.  If $n=e_\mu$, let
$\mathfrak T=N_0a(\mathbb Ze_{\mu'}\oplus\mathbb Ze_{\mu''})$.
If $n=e_\mu+s e_\nu$, $s=\pm1$, let
$\mathfrak T=N_0a\mathbb Z(e_\mu-s e_\nu)\oplus N_0a\mathbb Ze_\kappa$,
where $\kappa$ is the remaining direction.  Set
\[
 \Gamma'=\mathfrak T\oplus\mathbb Zv_h,\qquad
 \iota=[a\mathbb Z^3:\Gamma'],\qquad
 \Gamma_i=i\mathfrak T\oplus\mathbb Z(2iv_h),\quad i\ge3.
\]
These groups are adapted to $n$, with respective height periods $h$
and $2ih$.  Moreover $\Gamma'\subset N_0a\mathbb Z^3$, so $\nu$
is $\Gamma'$ invariant, and $\operatorname{sys}(\Gamma_i)\ge ia$.

In the coordinate case choose the domain
\[
 D_i=\{x\in a\mathbb Z^3:0\le x_\mu<2ih,
             \ 0\le x_{\mu'},x_{\mu''}<iN_0a\}.
\]
In the diagonal case choose
\[
 D_i=\{x\in a\mathbb Z^3:
 0\le x_\mu-sx_\nu<2iN_0a,\quad
 0\le x_\kappa<iN_0a,\quad
 0\le x_\mu+sx_\nu<2ih\}.
\]
The two diagonal coordinates divided by $a$ have the same parity.
Reducing them modulo $2iN_0$ and $2ih/a$ preserves this parity and
gives one representative per class. Thus these sets are domains of
$\Gamma_i$. In the definition of $\partial_1D_i$, the displacements
$w$ satisfy $|w|_\infty\le a$. Each ordinary coordinate changes by
at most $a$, while $x_\mu-sx_\nu$ and $x_\mu+sx_\nu$ change by
at most $2a$. Count one lattice layer at each end of an ordinary
coordinate interval and two at each end of a diagonal coordinate
interval. The parity count above gives the same number of sites
in each layer. The union bound therefore gives
\[
 \frac{\#\partial_1D_i}{\#D_i}
 \le\frac2i\left(\frac2{N_0}+\frac ah\right).
\]

Let $\mathbb A$ be the array of $A$ over $\Gamma'/\Gamma_i$, and
$\mathbb A^{\mathrm f}$ its folded array from
Section~\ref{sec:infinite-volume-states}.
Its number of terms is $|\mathbb A|=\#D_i/\iota$.
For $t>0$, the lower tangent inequality, folding inequality, and
Jensen's inequality in the folded measure give, in that order,
\[
\begin{split}
 t|\mathbb A|\nu(A)
 -\bigl(24\beta+2t\|A\|_\infty(2r-1)^3\bigr)\#\partial_1D_i
 &\le \log\frac{Z_{\Gamma_i}[t\mathbb A]}{Z_{\Gamma_i}[0]}\\
 &\le \log\frac{Z_{\Gamma_i}[t\mathbb A^{\mathrm f}]}{Z_{\Gamma_i}[0]}\\
 &\le t\mu_{\Gamma_i}^{t\mathbb A^{\mathrm f}}(\mathbb A^{\mathrm f})
 =t|\mathbb A|\mu_{\Gamma_i}^{t\mathbb A^{\mathrm f}}(A).
\end{split}
\]
The last equality is the wall symmetry of the folded measure.  Hence
\begin{equation}\label{IV-29-41}
 \nu(A)\le\mu_{\Gamma_i}^{t\mathbb A^{\mathrm f}}(A)
 +\iota\left(\frac{24\beta}{t}
       +2\|A\|_\infty(2r-1)^3\right)
       \frac{\#\partial_1D_i}{\#D_i}.
\end{equation}

We keep $t$ fixed while taking the volume limit.  To identify its
Gibbs specification, write the infinite folded family with its indices:
\begin{equation}\label{IV-29-38}
 I=\mathbb Z\times\mathfrak T,\qquad
 A_{(j,v)}=t w_j\tau_vA,\qquad
 B_{(j,v)}=w_j(B_A+v).
\end{equation}
Here $w_j$ are the alternating translations and reflections which
send the first slab to the $j$th slab.  Every bond belongs to at most
$2\#B_A$ of these supports.  For
$g=g_\perp+2ikv_h\in\Gamma_i$, define
$g\cdot(j,v)=(j+2ik,v+g_\perp)$.
The relations $w_{j+2ik}=\tau_{2ikv_h}w_j$ and
$\tau_{g_\perp}w_j=w_j\tau_{g_\perp}$ show that this free index
action translates both the functions and their supports.  Representatives
of $I/\Gamma_i$ are $\{1,\ldots,2i\}\times R_i$, where $R_i$
represents $\mathfrak T/i\mathfrak T$, and their sum on the torus is
$t\mathbb A^{\mathrm f}$.  Thus the finite-volume DLR limit applies
to one fixed locally finite perturbation family.

By compactness, a subsequence of the folded torus measures converges
to a Gibbs state $\nu_t$ for this family.  The finite measures, and
hence their limit, are invariant under
$\mathfrak T\oplus\mathbb Z(2v_h)$.  A fixed continuous local
$F$ in $\mathbb H^+(n,c)$ lies in
$\{c\le n\cdot x\le c+ih\}$ for all sufficiently large $i$.
The two-mirror positivity of the folded torus measure then gives
$\mu_{\Gamma_i}^{t\mathbb A^{\mathrm f}}(\overline{\theta F}F)\ge0$,
and this inequality passes to $\nu_t$.
The error in \eqref{IV-29-41} tends to zero, so $\nu(A)\le\nu_t(A)$.

Finally take a weak limit of a subsequence of $\nu_{1/k}$.
Continuity of the local DLR kernels as the strength of the fixed
family tends to zero makes its limit $\nu_0$ an unperturbed Gibbs
state.  The common period group, reflection positivity, and the
expectation inequality all pass to this limit.  The period group has
finite index, so $\nu_0$ is periodic, as required.
\end{proof}

\section{Symmetries of periodic Gibbs states}
\label{sec:gibbs-symmetries}

We first use the two-layer identity to obtain a common translation
group for all periodic Gibbs states.  This makes their class compact
and convex.  The comparison states can then test membership in its
reflection-positive subclass by separation.

\subsection{Temporal gauge and determination by two layers}

Fix an integer $t_0$.  Given the vertical variables $w_t$, define
\[
 \hat w_{t_0}=1,\qquad \hat w_{t+1}=\hat w_tw_t,
 \qquad \tilde u_t=\hat w_t\cdot u_t.
\]
This defines a continuous map
$\mathsf{tg}:\mathcal V(\mathbb L_a)\to(\mathcal V^{(2)})^{\mathbb Z}$:
each of its coordinates uses finitely many vertical variables and one
horizontal bond.  The gauge transformation $\hat w$ makes all vertical
bonds equal to $1$.  Write $U[\tilde u]$ for the resulting configuration,
with horizontal layers $\tilde u$ and vertical variables $1$.
For a gauge invariant local continuous function $F$,
\[
 F(U)=\tilde F(\mathsf{tg}(U)),\qquad
 \tilde F(\tilde u)=F(U[\tilde u]).
\]

For finite horizontal $S$, let $V_S(u)$ be the sum of the planar
plaquette actions which meet $S$.  It contains at most $2\#S$
terms, so $0\le V_S\le4\#S$.  Define the conditional layer kernel
\begin{equation}\label{IV-54-6}
 \tilde k_{S,t}(v;\tilde u)
 =k_S(\tilde u_{t-1}|_S,v)
   e^{-\beta V_S(v\vee\tilde u_t|_{S^c})}
   k_S(v,\tilde u_{t+1}|_S).
\end{equation}
The kernel $\tilde\gamma_{S,t}$ replaces $\tilde u_t|_S$ by $v$
and integrates with the normalized density
$\tilde k_{S,t}(v;\tilde u)\,dv$.  This density and its normalizer
are bounded below by $e^{-8\beta\#S}$ before normalization.
They depend continuously on finitely many coordinates.  A layer law
obeys the chain equations if it is fixed by every $\tilde\gamma_{S,t}$
on continuous local functions.

\begin{lemma}\label{IV-54-9}
The temporal-gauge image of every Gibbs state obeys the chain
equations.  The equations are invariant under each index reflection
$(\tilde\theta_m\tilde u)_t=\tilde u_{m-t}$.
If $\mathsf{tg}_*\nu$ is invariant under $\tilde\theta_{2t_0+1}$,
then $\nu$ is invariant under the bond reflection
$\theta_{e_3,(t_0+1/2)a}$.
\end{lemma}

\begin{proof}
The vertical plaquette formula becomes
$1-\frac12\operatorname{tr}(\tilde u_t(b)\tilde u_{t+1}(b)^{-1})$,
and the horizontal plaquette action is unchanged by $\hat w_t$.
If only the horizontal bonds $B=\{b^{(t)}:b\in S\}$ are integrated
in the DLR kernel, the variables $\hat w_s$ do not change.
The substitution $v=\hat w_t\cdot U_B$ preserves product Haar measure.
The plaquettes meeting $B$ consist of the planar plaquettes meeting
$S$, and one vertical plaquette above and one below each bond of $S$.
Their weight is precisely \eqref{IV-54-6}.  Thus
\[
 \gamma_B(\psi\circ\mathsf{tg})
 =(\tilde\gamma_{S,t}\psi)\circ\mathsf{tg},
\]
which proves the chain equations.

Symmetry of $k_S$ exchanges the two neighbouring layers in
\eqref{IV-54-6}; hence
$(\tilde\gamma_{S,m-t}\psi)\circ\tilde\theta_m
=\tilde\gamma_{S,t}(\psi\circ\tilde\theta_m)$.
This proves covariance of the equations.
With $\theta=\theta_{e_3,(t_0+1/2)a}$ one has
$\theta U[\tilde u]=U[\tilde\theta_{2t_0+1}\tilde u]$.
Consequently invariance of the layer law implies that $\nu$ and
$\theta_*\nu$ agree on all gauge invariant local continuous
functions.  Both measures are Gibbs states, so the gauge averaging
lemma~\ref{IV-14-13} makes them equal.
\end{proof}

\begin{lemma}\label{IV-54-13}
For $\beta>0$, two laws obeying the chain equations and agreeing on
two or more consecutive entire layers agree on all layers.
\end{lemma}

\begin{proof}
Suppose that the two laws agree on layers $m,\ldots,n$, with $m<n$.
We extend agreement to $n+1$.  Fix a nonempty finite horizontal set
$S$, and use the coordinates
\[
 u'=\tilde u_{n+1}|_S,\qquad v=\tilde u_n|_S,\qquad
 c=(\tilde u_n|_{S^c},\tilde u_{n-1},\ldots,\tilde u_m)\in Y.
\]
Here $Y$ is a compact countable product of copies of $G$.  Put
\[
 \Phi(v;c)=k_S(\tilde u_{n-1}|_S,v)
 e^{-\beta V_S(v\vee\tilde u_n|_{S^c})},\qquad
 N(u',c)=\int k_S(v,u')\Phi(v;c)\,dv.
\]
The assumption $m<n$ ensures that $\tilde u_{n-1}$ is among the
coordinates of $c$.  Both functions are continuous and local, with
$\Phi\ge e^{-6\beta\#S}$ and $N\ge e^{-8\beta\#S}$.
Let $\rho_i'$ be the image of the $i$th law on $(u',c)$ and set
$\sigma_i=N^{-1}\rho_i'$.  The chain equation at $(S,n)$ reads
\begin{equation}\label{IV-54-14}
 \rho_i(\psi)=\int\sigma_i(du',dc)\int
 k_S(v,u')\Phi(v;c)\psi(u',v,c)\,dv
\end{equation}
for every continuous local $\psi$.
Take $\psi=\varphi(v)\chi(c)/\Phi(v;c)$, with
$\varphi\in\mathcal T(G^S)$ and $\chi$ continuous local.
This is a function of the already equal layers.  Subtracting the two
identities gives
\[
 \int(K_S\varphi)(u')\chi(c)\,d(\sigma_1-\sigma_2)=0.
\]
The measure injectivity in Lemma~\ref{IV-29-15} yields
$\sigma_1=\sigma_2$.  Equation~\eqref{IV-54-14} then gives agreement
for all continuous local functions of $(u',v,c)$.  Enlarging $S$ to
contain the finitely many coordinates of any test on layer $n+1$
proves agreement on layers $m,\ldots,n+1$.

Reflecting the layer indices about $m+n$ and using the covariance in
Lemma~\ref{IV-54-9} proves extension to layer $m-1$.  Iterating both
extensions gives agreement on every finite block of layers.
Continuous local functions determine laws on this compact product,
so the two laws are equal.
\end{proof}

\begin{proposition}\label{IV-54-16}
Every periodic Gibbs state is invariant under every coordinate bond
reflection, and consequently under translations by $2a\mathbb Z^3$.
The periodic Gibbs states are exactly the Gibbs states invariant under
$2a\mathbb Z^3$; their set is nonempty, compact and convex.
\end{proposition}

\begin{proof}
For $\beta>0$, take the temporal gauge based at $t_0$.
Its layers $t_0,t_0+1$ are $u_{t_0}$ and
$w_{t_0}\cdot u_{t_0+1}$.  Proposition~\ref{IV-29-36} says that their
joint law is invariant under exchange.  The layer law and its
$\tilde\theta_{2t_0+1}$ reflection both obey the chain equations and
agree on these two layers.  Lemma~\ref{IV-54-13} makes them equal,
and Lemma~\ref{IV-54-9} gives the bond reflection invariance of $\nu$.
The same reasoning in each coordinate direction follows by applying
it to the coordinate permutation of $\nu$, which is again a periodic
Gibbs state.

At $\beta=0$, the DLR equation on the support of a local function
sets its expectation equal to its product Haar integral.  This
determines a unique Gibbs state, which is invariant under all lattice
isometries.

The product of the reflections at $x_\mu=a/2$ and $x_\mu=3a/2$
is translation by $2ae_\mu$.  Thus every periodic state is invariant
under $2a\mathbb Z^3$.  Conversely, this group has finite index,
so a Gibbs state invariant under it is periodic.  Compactness and
convexity follow from the DLR equations and the fixed common period
group; nonemptiness follows from the periodic volume construction.
\end{proof}

\subsection{All lattice reflections}

\begin{proposition}\label{IV-54-17}
For every $a>0$ and $\beta\ge0$, every periodic Gibbs state is
reflection positive at all coordinate site, coordinate bond, and
diagonal lattice mirrors, and is invariant under
$a\mathbb Z^3\rtimes W_3$.
In particular, the periodic Gibbs states are exactly the
translation invariant Gibbs states.  Their class is nonempty,
compact and convex, and contains every cubic torus-limit state.
\end{proposition}

\begin{proof}
Fix a mirror $\Sigma(n,c)$, $n\in\mathfrak N$, and write
$\theta=\theta_{n,c}$.  Let $\mathcal R$ be the periodic Gibbs
states reflection positive at this mirror.  It is nonempty by
Proposition~\ref{IV-29-40} with $A=0$.  It is compact and convex:
Proposition~\ref{IV-54-16} gives those properties for the ambient class,
and each condition
$\rho(\overline{\theta F}F)\in[0,\infty)$, with $F$ continuous
local in the positive half-space, is closed and affine.

Suppose $\nu\notin\mathcal R$.  There is a real continuous $f$ and
$\delta>0$ such that
$\nu(f)\ge\rho(f)+3\delta$ for every $\rho\in\mathcal R$.
Here separation can be obtained in a finite-dimensional space.
Indeed, if every finite family of continuous test integrals of $\nu$
were attained by an element of $\mathcal R$, compactness and the
finite intersection property would give an element with all the same
integrals, hence $\nu\in\mathcal R$.  Otherwise the finite vector
of integrals of $\nu$ lies outside the corresponding compact convex
set.  Its nearest point separates it by the vector joining that point
to the given vector, which gives $f$ and a positive gap.

Choose real local $A_0$ with $\|f-A_0\|_\infty\le\delta$.
Then $\nu(A_0)\ge\rho(A_0)+\delta$ for every $\rho\in\mathcal R$.
Translate $A_0$ by $2kan$, with $k$ so large that its support lies in
$\mathbb H^+(n,c)$, and call the result $A$.
All the states in question are invariant under this translation by
Proposition~\ref{IV-54-16}, so the same strict inequality holds for $A$.
Proposition~\ref{IV-29-40} now supplies $\nu_0\in\mathcal R$ with
$\nu(A)\le\nu_0(A)$, a contradiction.  This proves positivity for
continuous local functions.

The same positivity holds for bounded Borel local functions.  On
their finite compact coordinate space, approximate such a function
in $L^2$ for the sum of the marginal laws of $\nu$ and $\theta_*\nu$
by continuous functions.  Cauchy--Schwarz bounds the change in
$\nu(\overline{\theta F}F)$, so positivity passes to the limit.
Replacing $F$ by the conjugate of its reflection also gives positivity
on the negative half-space and for the opposite normal.

To obtain reflection invariance, positivity makes the form
$\nu(\overline{\theta F}F')$ Hermitian.  Taking one entry equal to
$1$ gives $\nu(F)=\nu(\theta F)$ for local $F$ in the positive
half-space.  Any local $F_0$ can be translated into that half-space
by $v=2kan$.  Since both $v$ and $-v$ are periods of $\nu$ and
$\theta\tau_v=\tau_{-v}\theta$, the same equality holds for $F_0$.
Thus $\theta_*\nu=\nu$ on all continuous local tests and hence on
all Borel sets.  The coordinate and diagonal lattice reflections
generate $a\mathbb Z^3\rtimes W_3$: adjacent site and bond mirrors
give unit translations, coordinate mirrors give sign changes, and
$x_\mu=x_\nu$ gives coordinate transpositions.
This proves the symmetry assertion.  The description and compactness
of the state class then follow from the translation-invariant DLR
equations, and the torus-limit inclusion is already known.
\end{proof}

\begin{corollary}\label{IV-54-31}
Every volume-first continuum limit of the centered moments of cubic
torus-limit Gibbs states is reflection positive at diagonal as well
as coordinate mirrors.
\end{corollary}

\begin{proof}
Proposition~\ref{IV-54-17} supplies diagonal positivity at each lattice
spacing.  The separated-test limit argument of
Section~\ref{sec:continuum-subsequences}, using the moving-entry
lemma and the plaquette support geometry, passes these inequalities
to the continuum limit.
\end{proof}

\section{Insertion estimates for periodic Gibbs states}
\label{sec:gibbs-insertions}

Fix $K\ge1$ and $\beta>0$.  For a board $(i,o)$ of side $l=a_i$,
a cell function $D$, and a sequence $\ell_j\to\infty$ of admissible
torus sides, write
\[
 \|D\|_{\mathrm{ch},\ell_\bullet}
 =\liminf_{j\to\infty}\|D\|_{\mathrm{ch},\ell_j}.
\]
The domination inequality \eqref{IV-18-40} holds for every periodic
Gibbs state and every such sequence, on every board.  We use it
together with Proposition~\ref{IV-54-17}.  All functions in the following
argument may be bounded Borel local functions; the application to
source observables uses continuous ones.

\subsection{Multiple reflections with an insertion}

We need the following consequence of the finite chessboard estimate.
Call $(i',o')$ coarser than $(i,o)$ if $i'\le i$ and
$o'-o\in l\mathbb Z^3$.  Then $d=a_{i'}/a_i=L^{i-i'}$ is odd,
and each coarse cell is the union of $d^3$ fine cells.

\begin{lemma}\label{IV-18-24}
For a coarse cell $C'$ and functions $D_k$ of its fine cells,
\begin{equation}\label{IV-18-25}
 \left\|\prod_{C_k\subset C'}D_k\right\|_{\mathrm{ch},\ell}^{\rm coarse}
 \le\prod_{C_k\subset C'}\|D_k\|_{\mathrm{ch},\ell}^{\rm fine},
 \qquad \ell\ge2a_{i'}.
\end{equation}
\end{lemma}

\begin{proof}
Writing $z=(o'-o)/l$, the fine cells of coarse cell $j$ have indices
\[
 z+dj+\{-(d-1)/2,\ldots,(d-1)/2\}^3.
\]
Every coarse wall is a fine wall.  Thus a coarse mirror copy sends
each fine cell to a fine cell by the corresponding fine mirror map;
its parity is unchanged, since parity is the determinant of that map.
The disseminated coarse product is consequently a product over all
fine cells in a torus period box.  Apply the fine chessboard inequality
to this product.  Every original factor has one copy for each coarse
cell, and its fine seminorm is unchanged by copying.  Taking the root
whose order is the number of coarse cells proves \eqref{IV-18-25}.
\end{proof}

For a coordinate bond mirror of height $c$, use the form
$\langle F',F\rangle_{\rho,c}=\rho(\overline{\theta_cF'}F)$
on the local positive-half-space algebra, and its seminorm
$\|F\|_{\rho,c}$.  Reflection positivity gives the Hermitian
symmetry and Cauchy--Schwarz inequality for this form.

\begin{proposition}\label{IV-55-9}
Let $\rho$ be a periodic Gibbs state.  Suppose that a cell $C$ of
side $l$ lies between two consecutive board walls $x_3=c$ and
$x_3=c+l$.  For a function $D$ of this cell and local functions
$F,F'$ in $\{x_3\ge c\}$,
\begin{equation}\label{IV-55-10}
 \left|\rho(\overline{\theta_cF'}D\tau_{le_3}F)\right|
 \le\|D\|_{\mathrm{ch},\ell_\bullet}
       \|F'\|_{\rho,c}\|F\|_{\rho,c}.
\end{equation}
\end{proposition}

\begin{proof}
Put $\widehat DF=D\tau_{le_3}F$ and
$D^\sharp=\overline{\theta_{c+l/2}D}$.
The middle reflection preserves the cell.  Translation invariance
and $\theta_c\theta_{c+l/2}=\tau_{-le_3}$ give
\[
 \langle F',\widehat DF\rangle_{\rho,c}
 =\langle\widehat{D^\sharp}F',F\rangle_{\rho,c}.
\]
Let $F_n=(\widehat{D^\sharp}\widehat D)^nF$ and
$s_n=\langle F,F_n\rangle_{\rho,c}$.  Moving adjoints between the
two entries gives
$s_{2p}=\|F_p\|^2$, $s_{2p+1}=\|\widehat DF_p\|^2$, and
$s_n^2\le s_{n-1}s_{n+1}$.  Hence
\begin{equation}\label{IV-55-13}
 s_1^n\le s_0^{n-1}s_n.
\end{equation}
For positive $s_1$ this follows from the increasing ratios
$s_{n+1}/s_n$; the case $s_1=0$ is immediate.

Expanding the iterates yields
\[
 s_n=\rho\left(\overline{\theta_cF}
       \prod_{r=0}^{2n-1}D_r\,\tau_{2nle_3}F\right),\qquad
 D_r=\begin{cases}\tau_{rle_3}D^\sharp,&r\text{ even},\\
                     \tau_{rle_3}D,&r\text{ odd}.
       \end{cases}
\]
These are successive mirror copies, so all their seminorms equal
$\|D\|_{\mathrm{ch},\ell}$.
Choose an odd coarsening factor $d=L^{r_0}$ large enough that $F$
is supported in a coarse cell directly above the wall $c$.
Such a board has side $l'=dl$ and center obtained from that of $C$
by adding $(d-1)le_3/2$; its bottom wall remains $c$.
For $n=n'd$, group the $2n$ functions $D_r$ into $2n'$ consecutive
coarse cells, with $F$ and its reflected copy in the two end cells.
Chessboard domination and \eqref{IV-18-25} give
\[
 s_n\le\|F\|_\infty^2
       \|D\|_{\mathrm{ch},\ell_\bullet}^{2n}.
\]
To justify the lower limit here, choose infinitely many sides with
$\|D\|_{\mathrm{ch},\ell_j}$ within any given positive error of
its lower limit; all other factors are bounded on that same
subsequence by the displayed product.
Combining this estimate with \eqref{IV-55-13}, taking $n$th roots,
and letting $n'\to\infty$ gives
$\|\widehat DF\|\le\|D\|_{\mathrm{ch},\ell_\bullet}\|F\|$.
If $\|F\|=0$, use any $n\ge2$ first to obtain $s_1=0$.
Cauchy--Schwarz now proves \eqref{IV-55-10}.
\end{proof}

\begin{proposition}\label{IV-55-19}
Let $\rho,\nu$ be periodic Gibbs states, and let $f$ be a real
Lipschitz function compactly supported in the interior of the cell
$C$ above.  Suppose $(f)$ has a normalizer $E$ of radius $\varrho$
and bound $Q\ge1$ along $\ell_\bullet$.  Then
\begin{equation}\label{IV-55-20}
 \left|\rho\left(\overline{\theta_cF'}
 [\mathcal O_K(f)-\nu\mathcal O_K(f)]\tau_{le_3}F\right)\right|
 \le\frac{2Q}{\varrho}\|f\|_\sharp
       \|F'\|_{\rho,c}\|F\|_{\rho,c}.
\end{equation}
\end{proposition}

\begin{proof}
There is nothing to prove for $f=0$.  Otherwise put
$X=\mathcal O_K(f)$ and $D_z=e^{zX-E(z)}$, for
$|z|<\varrho/\|f\|_\sharp$.
The normalizer gives
$\|D_z\|_{\mathrm{ch},\ell_\bullet}\le Q$.
The expectation with $D_z$ inserted in \eqref{IV-55-10} is
holomorphic: the exponential series in the bounded local variable
$X$ converges uniformly on every bounded $z$-disc and may be
integrated term by term.  Cauchy's estimate at zero therefore bounds
the insertion of $X-E'(0)$ by
$Q\varrho^{-1}\|f\|_\sharp\|F'\|\|F\|$.
By \eqref{IV-19-14}, $|\nu(D_z)|\le Q$, so a second Cauchy
estimate gives $|\nu X-E'(0)|\le Q\varrho^{-1}\|f\|_\sharp$.
Finally \eqref{IV-55-10} with $D=1$ bounds the remaining expectation
by $\|F'\|\|F\|$.  Adding these two contributions proves
\eqref{IV-55-20}.
\end{proof}

\subsection{Uniform constants and thin plates}

Suppose the pair $(K,\beta)$ has the source systems used in the cell
bound, on an unbounded set $\mathfrak l$ of torus sides, with the
common constants
\[
 (2\varrho;C_e,\kappa_e;E_+,C_{\mathrm{low}};C_{\mathrm{gap}}),
 \qquad \varrho>0,\quad \kappa_e>\kappa_{\mathcal D},
\]
where the other constants are nonnegative.  Here the last constant
means that their vacuum terms satisfy
$E_{\ell'}\ell'^{-3}-E_\ell\ell^{-3}\le C_{\mathrm{gap}}$
whenever $\ell\ge\ell'$ in $\mathfrak l$.
Define
\begin{equation}\label{IV-55-25}
 Q_{\mathrm{mid}}(l)=\exp\left\{
 (E_++C_{\mathrm{low}}+C_{\mathrm{gap}})l^3
 +\Phi_1(l/4,l/4;l)\right\}.
\end{equation}
The explicit formula for $\Phi_1$ gives
\begin{equation}\label{IV-56-3}
 \sup_{0<l\le1}Q_{\mathrm{mid}}(l)\le\bar Q_{\mathrm{mid}}
 :=\exp\left\{E_++C_{\mathrm{low}}+C_{\mathrm{gap}}
 +C_eS(\kappa_e)\left[\frac{27}{26}
 +27\{32(C''_\Sigma+\log_L10)+4C'_\Sigma\}\right]\right\}.
\end{equation}

\begin{lemma}\label{IV-55-32}
If all entries of a tuple are supported in the middle cube of one
cell, namely the cube of sup-radius $l/4$ about its center, then
there is one normalizer $E^\circ$, holomorphic on the source disc
of radius $2\varrho$, such that
\begin{equation}\label{IV-55-33}
 \|e^{\mathcal O_K(\sum z_qf_q)}\|_{\mathrm{ch},\ell}
 \le |e^{E^\circ(z)}|Q_{\mathrm{mid}}(l),
 \qquad \ell\in\mathfrak l,\quad\ell\ge2l,
 \quad z\in\mathbb D_\varrho(\underline f).
\end{equation}
For $n\ge2$, if the first two supports have positive sup-distance
$d_\infty$, its mixed derivative obeys
\begin{equation}\label{IV-55-34}
 |\partial_1\cdots\partial_nE^\circ(0)|
 \le\left(\frac n{2\varrho}\right)^n\prod_q\|f_q\|_\sharp
 27C_eS(\kappa_e)\sum_{s=0}^K\bar N_s(\{1\})
 e^{-\delta_e[d_\infty L^s/M-4]_+},
 \quad\delta_e=\kappa_e-\kappa_{\mathcal D}.
\end{equation}
Here $\bar N_s(\{1\})$ is the maximum count of cubes of side
$ML^{-s}$ meeting a mirror copy of the first support.
Consequently \eqref{IV-55-20} holds with $Q=Q_{\mathrm{mid}}(l)$
for every periodic pair $\rho,\nu$ and every $f$ in that middle cube.
\end{lemma}

\begin{proof}
Choose the least reference side $\ell'\in\mathfrak l$ with
$\ell'\ge2l$, and construct $E^\circ$ there, with carrier the
whole middle cube and margin $l/4$.
The vacuum floor and comparison of vacuum densities give
$\Delta_E\le C_{\mathrm{gap}}$ on every larger side in
$\mathfrak l$.  The cell bound~\ref{IV-7-27}, using the single
carrier rather than counting its entries separately, gives
\eqref{IV-55-33}.  Its normalizer depends on the reference system,
the fixed board and functions, and is independent of the larger
side.  Periodization preserves the sharp norms and identifies the
lattice and torus cell seminorms.  Thus this is one normalizer on
all the stated sides, not a choice depending on $\ell$.
The first inequality of \eqref{IV-7-31}, with source radius
$2\varrho$, gives \eqref{IV-55-34}.
An exhausting sequence in $\mathfrak l$, after discarding its
finitely many sides below $2l$, now satisfies the hypotheses of
Proposition~\ref{IV-55-19}.
\end{proof}

For $0<t\le1$ and $\lambda>0$ set
\[
 B_\pm(t,\lambda)=
 \{|x_1|\le t,\ |x_2|\le\lambda/2,\ t\le\pm x_3\le2t\},
 \qquad
 l_\lambda=\min\{L^r:r\in\mathbb N,\ L^r\ge8+2\lambda\}.
\]
Both boxes lie in the middle of the central cell of side $l_\lambda$,
and their sup-distance is at least $2t$.
For $\delta>0$ define the finite number
\begin{equation}\label{IV-55-24}
 C_{\mathrm{dep}}(\delta,c_0)
 =e^{\delta c_0}\sup_{x\ge0}(x+2)(2x+2)e^{-\delta x}
       (4+4/\delta+4/\delta^2).
\end{equation}

\begin{proposition}\label{IV-55-37}
Under the source-system assumptions above, for every periodic Gibbs
state $\rho$ and real $f$ supported in $B_+(t,\lambda)$,
\begin{equation}\label{IV-55-38}
 |S_2^T[\rho](\theta f,f)|
 \le\frac{\|f\|_\sharp^2}{\varrho^2}
 \left[4\log2(2Q_{\mathrm{mid}}(l_\lambda)+2)^2
 +27C_eS(\kappa_e)C_{\mathrm{dep}}(\delta_e,4)
                   \frac{\lambda+1+m'}t\right].
\end{equation}
If $f=\alpha\otimes\varphi$ is a plate of
Section~\ref{sec:slab-operators}, with $\alpha\in\operatorname{Prof}_t$ and
$\operatorname{supp}\varphi\subset[-\lambda/2,\lambda/2]$, then
\begin{equation}\label{IV-55-39}
 |S_2^T[\rho](\theta f,f)|\le
 \left(\frac{600M\|\varphi\|_{C^1}}{\varrho}\right)^2
 \left[4\log2(2Q_{\mathrm{mid}}(l_\lambda)+2)^2
 +27C_eS(\kappa_e)C_{\mathrm{dep}}(\delta_e,4)(\lambda+1+m')\right]t^{-3}.
\end{equation}
\end{proposition}

\begin{proof}
Apply Lemma~\ref{IV-55-32} to $(\theta f,f)$ in the central
cell of side $l_\lambda$.  The cumulant transfer estimate
\eqref{IV-19-17} bounds its difference from
$\partial_1\partial_2E^\circ(0)$ by
$4\log2(2Q_{\mathrm{mid}}(l_\lambda)+2)^2\|f\|_\sharp^2/\varrho^2$.
At depth $s$, put $x_s=tL^s/M$.  A cube grid of side $W$ has at
most $\prod_\mu(s_\mu/W+2)$ cubes meeting any translate of a
box of edge lengths $s_\mu$: in each coordinate the admissible
integer centers lie in an interval of length $s_\mu/W+1$.
Every mirror copy of the first support lies in a box with edge
lengths $(2t,\lambda,t)$.  Thus
\[
 \bar N_s(\{1\})\le(2x_s+2)(\lambda x_s/t+2)(x_s+2).
\]
In \eqref{IV-55-34}, $d_\infty\ge2t$.  Its depth sum is bounded
by
\[
 \sum_{s=0}^K(x_s+2)(2x_s+2)(\lambda x_s/t+2)
                    e^{-\delta_e[2x_s-4]_+}
 \le C_{\mathrm{dep}}(\delta_e,4)\frac{\lambda+1+m'}t.
\]
Here is the scale summation.  After extracting
$e^{4\delta_e}\sup_{x\ge0}(x+2)(2x+2)e^{-\delta_ex}$,
it remains to bound
$(\lambda/t)\sum_sx_se^{-\delta_ex_s}+2\sum_se^{-\delta_ex_s}$.
Since successive $x_s$ grow by $L\ge2$,
$\sum_{x_s\le1}x_s\le2$ and $\sum_{x_s>1}x_s^{-1}\le2$.
The inequalities $u^2e^{-u}\le2$ and $ue^{-u}\le1$ give
\[
 \sum_sx_se^{-\delta_ex_s}\le2+4/\delta_e^2,
 \qquad
 \sum_se^{-\delta_ex_s}\le1+m'+1/t+2/\delta_e.
\]
The number of scales with $x_s\le1$ is at most
$1+m'+\log_L(1/t)$, and $\log_L(1/t)\le1/t$ was used in the
second inequality.  These bounds give the stated depth estimate
and prove \eqref{IV-55-38}.

The profile bounds $\|\alpha\|_\infty\le3$ and
$\operatorname{Lip}\alpha\le12/t$ yield
$\|f\|_\infty\le3\|\varphi\|_\infty$ and
$\operatorname{Lip}f\le15\|\varphi\|_{C^1}/t$.
Consequently $\|f\|_\sharp\le600M\|\varphi\|_{C^1}/t$.
Substitute this in \eqref{IV-55-38} and use $t\le1$ in its
constant term to obtain \eqref{IV-55-39}.
\end{proof}

\section{The rotation hypotheses for volume-first limits}
\label{sec:gibbs-continuum-rotation}

Let $K_j\to\infty$, $\beta_j>0$, and let $\rho_j,\nu_j$ be
periodic Gibbs states at $(K_j,\beta_j)$.  Suppose that these pairs
have source systems on unbounded sets of sides with one set of
constants $(2\varrho;C_e,\kappa_e;E_+,C_{\mathrm{low}};C_{\mathrm{gap}})$
as above.  Suppose also that their separated moments have uniform
bound tables and that
\[
 \sigma(X)=\lim_{j\to\infty}\rho_j(X_{K_j}^{\nu_j})
\]
exists for every separated tuple.  The canonical sequences over
$\mathfrak d(g)$ satisfy the uniform assumptions; their concrete
constants are specified below.

\begin{proposition}\label{IV-56-17}
The limit $\sigma$ satisfies the insertion bound (R5) with
\begin{equation}\label{IV-56-18}
 c_{\mathrm I}=\frac1{40L},\qquad p_{\mathrm I}=0,\qquad
 C_{\mathrm I}=\frac{2\bar Q_{\mathrm{mid}}}{\varrho}.
\end{equation}
It has the inherited properties (R1)--(R3).
\end{proposition}

\begin{proof}
The moving-entry argument of Section~\ref{sec:continuum-subsequences}
gives translation and cubic invariance, the uniform moment bounds,
and coordinate reflection positivity.  The resulting Hilbert space
therefore has the contraction semigroup of
Proposition~\ref{IV-27-22}.

Fix the data in (R5): $0<\eta\le1$, a real $h$ supported in
$B(y,\eta/(40L))$ with $|y_3|\le\eta/8$, and
$F,G\in\mathbb C\mathfrak S(\{x_3>\eta/4\})$.
Choose a scale and a point, independently of $j$, with
\begin{equation}\label{IV-56-20}
 \eta/(5L)<l=a_i\le\eta/5,\qquad
 y^*\in a_{i+2}\mathbb Z^3,\quad
 |y-y^*|_\infty\le l/(2L^2),\qquad c=y_3^*-l/2.
\end{equation}
Since $L\ge13$, these choices give
\[
 \operatorname{supp}h\subset
 \{|x-y^*|_\infty<l/4\},\qquad
 -\eta/4<c<c+l<\eta/4.
\]
Indeed, the first radius is less than
$l/8+l/(2L^2)<l/4$, while
$|y_3^*|+l/2\le\eta(1/8+1/1690+1/10)<\eta/4$.
Put $F_1=\tau_{2ce_3}F$ and $F_2=\tau_{-le_3}G$.
Their finitely many compact supports lie a positive distance above
$c$; furthermore $\theta_cF_1=\theta F$ and $\tau_{le_3}F_2=G$.

For all sufficiently large $j$, $(i,y^*)$ is a lattice board and
$c$ is one of its bond walls.  In fact $K_j\ge i+2$ ensures
$y^*\in a_{K_j}\mathbb Z^3$, and the odd integer $l/a_{K_j}$
makes $c/a_{K_j}$ a half-integer.
The centered lattice polynomials associated with $F_1,F_2$ are
supported in the closed half-space above $c$: a plaquette extends
by at most $a_{K_j}/2$ in its third coordinate from its center.
Apply \eqref{IV-55-20}, with its uniform middle-cell normalizer,
to these two polynomials and $h$.  Covariance of the centering law
$\nu_j$ gives
\[
 \left|\sigma_{\rho_j,\nu_j}(\theta\bar F\cdot(h)\cdot G)\right|
 \le\frac{2Q_{\mathrm{mid}}(l)}{\varrho}\|h\|_\sharp
 \sigma_{\rho_j,\nu_j}(\theta_c\bar F_1\cdot F_1)^{1/2}
 \sigma_{\rho_j,\nu_j}(\theta_c\bar F_2\cdot F_2)^{1/2}.
\]
All the tuples on this line are separated and fixed independently
of $j$.  Taking the limit gives the same inequality for $\sigma$.

Translation invariance identifies the two limiting square roots
with $\|[\tau_{-ce_3}F_1]\|$ and
$\|[\tau_{-ce_3}F_2]\|$.  They are respectively
\[
 \|e^{-(c+\eta/4)H}[\tau_{-(\eta/4)e_3}F]\|,
 \qquad
 \|e^{-(\eta/4-c-l)H}[\tau_{-(\eta/4)e_3}G]\|.
\]
Both semigroup times are positive, so contractivity bounds these by
the norms in (R5).  Finally $l\le1$ and
$Q_{\mathrm{mid}}(l)\le\bar Q_{\mathrm{mid}}$.
This proves (R5) with \eqref{IV-56-18}.
\end{proof}

\begin{proposition}\label{IV-56-22}
The limit $\sigma$ satisfies (R7) with exponent $a_\Gamma=3$.
If $\sigma(f)=c_\sigma\int f$, one may take $C_0=0$ and, for
a nonzero plate label $\varphi$ of support length $\lambda_\varphi$,
\begin{equation}\label{IV-56-24}
\begin{split}
 C_\varphi={}&
 \left(\frac{600M\|\varphi\|_{C^1}}{\varrho}\right)^2
 \left[4\log2(2Q_{\mathrm{mid}}(l_{\lambda_\varphi})+2)^2
 +27C_eS(\kappa_e)C_{\mathrm{dep}}(\delta_e,4)
                          (\lambda_\varphi+1+m')\right]\\
 &+(6c_\sigma\lambda_\varphi\|\varphi\|_\infty)^2.
\end{split}
\end{equation}
In the self-centered case $\rho_j=\nu_j$, the last term is zero.
\end{proposition}

\begin{proof}
Translate the label in the $e_2$ direction so that its support
interval is centered at zero.  This leaves the limiting two-point
reflection form unchanged.  For $f=\alpha\otimes\varphi$,
the moment--cumulant identity is
\[
 \sigma_{\rho_j,\nu_j}(\theta f,f)
 =S_2^T[\rho_j](\theta f,f)+\delta_j(\theta f)\delta_j(f),
 \quad \delta_j(f)=\rho_j\mathcal O_{K_j}(f)-\nu_j\mathcal O_{K_j}(f).
\]
The first term is bounded by \eqref{IV-55-39}, uniformly in $j$.
The one-point limit from Section~\ref{sec:continuum-subsequences}
gives $\delta_j(f)\to c_\sigma\int f$.
The box containing $f$ has volume $2\lambda_\varphi t^2$ and
$\|f\|_\infty\le3\|\varphi\|_\infty$, so
$|c_\sigma\int f|^2\le
(6c_\sigma\lambda_\varphi\|\varphi\|_\infty)^2t^4$.
For $t\le1$ this is bounded by the same constant times $t^{-3}$.
Passing to the limit proves the upper bound, and reflection
positivity gives its nonnegativity.  For equal laws all $\delta_j$
vanish, giving the last assertion.
\end{proof}

\begin{theorem}\label{IV-56-29}\label{IV-56-39}\label{IV-56-41}
Fix $0<g\le\gamma_\varrho$.  Let $(K_j,\beta_j)\in\mathfrak d(g)$,
$K_j\to\infty$, and let $\rho_j,\nu_j$ be periodic Gibbs states.
Suppose that their moments centered at $\nu_j$ converge on every
separated tuple to $\sigma$.  Then their self-centered moments
converge to $\sigma^{\mathrm c}$, and
$\sigma^{\mathrm c}$ satisfies (R0)--(R7).  The uncentered limit
$\sigma$ satisfies (R1)--(R7) and
\begin{equation}\label{IV-56-16}
 \sigma(f_1,\ldots,f_n)=
 \sum_{A\subset\{1,\ldots,n\}}
 \prod_{i\in A}\left(c_\sigma\int f_i\right)
 \sigma^{\mathrm c}((f_i)_{i\notin A}).
\end{equation}
Both functionals are invariant under all Euclidean motions and
reflection positive at every Euclidean mirror.  Their joint
energy--momentum spectrum is contained in the forward cone
$\{(E,q):E\ge|q|\}$.

In particular these conclusions apply to every volume-first
continuum subsequence of cubic torus-limit states, with
$\rho_j=\nu_j$ and $c_\sigma=0$.
\end{theorem}

\begin{proof}
Theorem~\ref{III-111-27} provides source systems on
$\mathfrak L(g)$ at radius $2\varrho$, and
Proposition~\ref{IV-19-24} constructs their uniform normalizers.
The source-system constants are
\[
 C_e=C_{\mathrm S},\quad\kappa_e=\kappa,\quad E_+=C_{\mathrm S},
 \quad C_{\mathrm{low}}=C_{\mathrm S}(1+\log g^{-2}),
 \quad C_{\mathrm{gap}}=C_{\mathrm S}.
\]
These are the constants already verified in
Section~\ref{sec:infinite-volume-states}; none depends on $j$ or
on the side.  Propositions~\ref{IV-54-17} and \ref{IV-18-38} supply
all lattice symmetries, all mirror positivities, and chessboard
domination along any exhausting sequence in $\mathfrak L(g)$.

Let $C_n^\uparrow(d;R)$ be the non-increasing-in-$d$ envelope of
the uniform cumulant table from Proposition~\ref{IV-22-13} and
the following cumulant bounds.  The moment--cumulant formula gives
the finite moment tables
\begin{equation}\label{IV-56-38}
 C_n^{\mathrm m}(d;R)=\sum_{\pi\in\operatorname{Part}(n)}
                         \prod_{J\in\pi}C_{|J|}^\uparrow(d;R),
 \qquad
 C_n^{\mathrm c}(d;R)=\sum_{\pi\in\operatorname{Part}_{\ge2}(n)}
                         \prod_{J\in\pi}C_{|J|}^\uparrow(d;R).
\end{equation}
For $C_1^\uparrow$, use the two-state one-point bound; in the
second sum singleton blocks are absent.  Each separated subtuple
has the same lower separation and radius bounds, so these tables
apply uniformly to all the indicated states.

Cumulants are finite polynomials in moments.  Thus convergence of
$\sigma_{\rho_j,\nu_j}$ gives convergence of every cumulant of
order at least two of $\rho_j$; centering changes none of them.
Their partition sums without singleton blocks give convergence of
the self-centered moments to $\sigma^{\mathrm c}$ and (R0).
The one-point limit is $c_\sigma\int f$; the partition expansion
gives \eqref{IV-56-16}.
The moving-entry and positivity arguments give (R1)--(R3), with
the moment tables \eqref{IV-56-38}.
Proposition~\ref{IV-56-17} gives (R5), and
Proposition~\ref{IV-56-22} gives (R7), for both sequences.

For (R6), if the entries of $F$ have compact support in
$\{x_1+x_3>0\}$, their plaquette polynomials have all bonds in
that half-space once $a_{K_j}$ is smaller than the minimum positive
height of these supports.  The maximal variation of $x_1+x_3$
over a plaquette from its center is $a_{K_j}$.
Lattice diagonal positivity then gives
$\sigma_{\rho_j,\nu_j}(\theta_\Sigma\bar F\cdot F)\ge0$;
covariance of $\nu_j$ ensures the exact centered identity.
Taking the limit gives full diagonal reflection positivity and
hence (R6).  Together with the moment bounds and coordinate RP,
Proposition~\ref{IV-28-21} yields $H\ge|P_1|$, which is (R4).
There is only one Hilbert space and one joint spectral measure here:
both descriptions complete the same reflection form and have matrix
elements $\sigma(\theta\bar F\cdot\tau_{te_3+v}G)$.
Uniqueness of their Laplace--Fourier representation identifies the
spectral measures if the constructions are made separately.

The rotation theorem~\ref{IV-37-34} now applies to
$\sigma^{\mathrm c}$.  Formula~\eqref{IV-56-16} transfers its
rotational invariance to $\sigma$, because $\int f$ is Euclidean
invariant.  Transporting coordinate RP by rotations and translations
gives every mirror.  The same transport turns the component cone
bounds into the full forward cone: apply the spectral inequality
in a countable dense set of spatial directions and intersect the
corresponding closed half-cones.
\end{proof}

\section{Conditioning periodic Gibbs states}
\label{sec:gibbs-conditioning}

The decomposition argument changes a lattice law by a bounded
translation-invariant density.  Such a change preserves the Gibbs
equations, including their fixed coupling; this is the closure
condition \eqref{IV-35-2} required of the classes.

\begin{proposition}\label{IV-14-17}\label{IV-54-18}
Let $\nu\in\mathcal G^{\rm per}_{a,\beta}$ and let $h\ge0$ be
bounded and $\mathcal I$-measurable, with $\nu(h)=1$.
Then $h\nu\in\mathcal G^{\rm per}_{a,\beta}$.
The extreme points of $\mathcal G^{\rm per}_{a,\beta}$ are
exactly its ergodic members.
\end{proposition}

\begin{proof}
For a finite bond set $B$ and a continuous local function $F$,
choose a continuous local $h'$ with
$\nu(|h-h'|)<\varepsilon$.  Translate the finite support of $h'$
away from $B$, obtaining $\tau_vh'$ supported in $B^c$.
Translation invariance of $h$ and $\nu$ gives
$\nu(|h-\tau_vh'|)<\varepsilon$.  The Gibbs equation and the
fact that the kernel $\gamma_B$ leaves outside functions fixed give
\[
 \nu(\tau_vh'\,\gamma_BF)
 =\nu\bigl(\gamma_B(\tau_vh'F)\bigr)
 =\nu(\tau_vh'F).
\]
Since $\|\gamma_BF\|_\infty\le\|F\|_\infty$, replacing
$\tau_vh'$ by $h$ changes each side by at most
$\varepsilon\|F\|_\infty$.  Letting $\varepsilon\downarrow0$
proves $(h\nu)(\gamma_BF)=(h\nu)(F)$.  The density and the
original law are translation invariant, so the new Gibbs law is
translation invariant and hence periodic.

If an ergodic $\nu$ is a nontrivial convex combination of periodic
Gibbs laws, each summand has a bounded Radon--Nikodym density with
respect to $\nu$.  Invariance of the two measures makes that density
translation invariant almost everywhere.  It has an invariant
version and is constant by ergodicity; its integral is one, so the
summand equals $\nu$.  Conversely, a nontrivial invariant event
$E$ decomposes $\nu$ into its normalized restrictions to $E$ and
$E^c$.  These are distinct periodic Gibbs states by the first part.
\end{proof}

\begin{proposition}\label{IV-54-20}
There is a regular conditional distribution $(\nu^U)_U$ of
$\nu\in\mathcal G^{\rm per}_{a,\beta}$ given $\mathcal I$.
For any such distribution,
\begin{equation}\label{IV-54-21}
 \nu(F)=\int\nu^U(F)\,\nu(dU)
\end{equation}
for bounded Borel $F$, and $\nu^U$ is an ergodic periodic Gibbs
state for $\nu$-almost every $U$.
\end{proposition}

\begin{proof}
The configuration space is compact metrizable.  Choose a countable,
uniformly dense rational vector space of real continuous functions
containing $1$, and choose conditional expectations for its members.
Outside one invariant null set, these conditional expectations
satisfy all rational linear relations, preserve $1$, and have norm
at most the supremum norm.  At each remaining $U$ they extend to a
norm-one unital real functional on the continuous functions.  It is
positive: for $0\le F\le1$,
$|1-\Lambda_U(F)|=|\Lambda_U(1-F)|\le1$.
Riesz representation gives a probability measure $\nu^U$.
On the exceptional set use $\nu$ itself.  Uniform approximation
and the functional monotone-class theorem show that, for every
bounded Borel $F$, the function $U\mapsto\nu^U(F)$ is a version of
$E_\nu[F\mid\mathcal I]$.  This constructs the distribution and
also gives \eqref{IV-54-21}.

Fix a countable real family $\mathcal C_0$ of continuous local
functions, dense on each finite set of bonds and hence determining
finite Borel measures on the configuration space.  There are
countably many translations and finite bond sets.  The identities
\[
 \nu(1_E\tau_vF)=\nu(1_EF),\qquad
 \nu(1_E\gamma_BF)=\nu(1_EF),\qquad E\in\mathcal I,
\]
hold for every $F\in\mathcal C_0$: the first uses invariance, and
the second uses Proposition~\ref{IV-54-18} for the normalized
restriction to $E$ when $\nu(E)>0$.  Conditional expectation and
one countable union of null sets therefore give
\[
 \nu^U(\tau_vF)=\nu^U(F),\qquad
 \nu^U(\gamma_BF)=\nu^U(F)
\]
simultaneously for all these tests, translations and finite sets.
Uniform approximation, and the supremum-norm contraction of
$\gamma_B$, extend the identities to all continuous local tests.
Thus $\nu^U$ is a translation-invariant Gibbs state almost surely.

For ergodicity let $A_N$ average translations over
$\{0,a,\ldots,(N-1)a\}^3$.  For any invariant probability law,
these averages converge strongly in $L^2$ to conditional expectation
onto $\mathcal I$.  Indeed $A_N$ is the product of the three
commuting one-dimensional Cesaro averages, whose limiting
projections commute; their product projects onto the common
invariant subspace.  Write $g_F(U)=\nu^U(F)$ for $F\in\mathcal C_0$.
Disintegration and this convergence give
\[
 \int\nu^U(|A_NF-g_F|^2)\,\nu(dU)
 =\nu(|A_NF-g_F|^2)\longrightarrow0.
\]
A common subsequence $N_i$ makes the integrand tend to zero for
all $F\in\mathcal C_0$ outside a single null set.  Applying the
conditional-distribution identity to the countably many invariant
functions $g_F$ and $g_F^2$ gives
$\nu^U(|g_F-g_F(U)|^2)=0$ almost surely.  Thus, in $L^2(\nu^U)$,
$A_{N_i}F\to\nu^U(F)$.  The mean ergodic theorem for the already
invariant law $\nu^U$ identifies this limit with
$E_{\nu^U}[F\mid\mathcal I]$.
For any invariant event $E$, the finite measures $1_E\nu^U$ and
$\nu^U(E)\nu^U$ consequently agree on $\mathcal C_0$ and are
equal.  Evaluating at $E$ gives $\nu^U(E)=\nu^U(E)^2$.
This holds for every $E$ on the same full-measure set, because
$\mathcal C_0$ determines measures.  Hence $\nu^U$ is ergodic.
\end{proof}

\subsection{The classes used in the continuum decomposition}

For the fixed source parameters and $0<g\le\gamma$, retain the
functions $C_1^g,C_n^g$ from \eqref{IV-20-3} and
\eqref{IV-22-7}.  For finite $d,R>0$ write
\[
 C_n^{g\uparrow}(d;R)
 =\inf_{0<d'\le d,\ R'\ge R}C_n^g(d';R'),
 \qquad C_1^g(d;R)=C_1^g(R),
\]
and define the strictly positive bound system
\begin{equation}\label{IV-78-5}
 \widehat C_1^g(R)=1+C_1^g(2R),\qquad
 \widehat C_n^g(d;R)=1+C_n^{g\uparrow}(d;R),\quad
 \widehat C_n^g(\infty;R)=1+\inf_{d>0}C_n^{g\uparrow}(d;R).
\end{equation}

\begin{proposition}\label{prop:gibbs-frame-classes}
For every $(K,\beta)\in\mathfrak d(g)$ the set
$\mathcal G^{\rm per}_{K,\beta}$ satisfies the closure
condition \eqref{IV-35-2} and obeys the bound system
$\widehat{\mathsf C}^g=(\widehat C_1^g,
(\widehat C_n^g)_{n\ge2})$.
Every sequence $K_j\to\infty$, $(K_j,\beta_j)\in\mathfrak d(g)$,
and $\nu_j\in\mathcal G^{\rm per}_{K_j,\beta_j}$ therefore
defines a frame in a class system with these same bounds.
\end{proposition}

\begin{proof}
The class is nonempty by periodic-volume compactness.  Its full
lattice invariance and RP are Proposition~\ref{IV-54-17}; its closure
under the normalized invariant densities is
Proposition~\ref{IV-54-18}.  These are precisely the two class
conditions.  Proposition~\ref{IV-22-14} gives the two-state one-point
bound and both cumulant bounds uniformly over the class.
A tuple of type $(d;R)$ is also of type $(d';R')$ whenever
$d'\le d$ and $R'\ge R$, so its cumulant bound may be replaced
by the infimum defining $C_n^{g\uparrow}$.  The infimum is finite
and nonnegative; adding one ensures strict positivity and preserves
monotonicity.  Its extension at $d=\infty$ is nonincreasing in $d$
and nondecreasing in $R$ for the same reason.  For the one-point
bound, a function supported in a closed radius-$R$ ball is admissible
with radius $2R$; \eqref{IV-22-15} therefore gives the bound
$\widehat C_1^g(R)\|f\|_\sharp$.  All constants were fixed before
choosing the cutoff, the states, or the frame.
\end{proof}

\section{Point insertions and energy bounds for reconstruction}
\label{sec:reconstruction-energy}

The reconstruction uses more than reflection positivity and bounds
on separated moments.  Its additional input is an $L^1$ estimate
for one insertion between two positive-time vectors.  We first
construct point forms with polynomial energy bounds, then use
them to define smeared fields on a common domain.

Let $\sigma$ be a real, normalized, symmetric multilinear
functional on separated tuples.  Throughout this section assume:
\begin{enumerate}[label=(Rec\arabic*)]
\item $\sigma$ is invariant under all Euclidean isometries.
\item It is reflection positive at $x_3=0$.  Its reflection
Hilbert space $\mathcal H$ is separable, and its translations
are represented by $e^{-tH}$ and $\mathcal U(v)$, with $H\ge0$,
the stated strong continuity, and joint spectral measure
$\mathcal P$:
\begin{equation}\label{IV-75-12}
 e^{-tH}[F]=[\tau_{te_3}F],\qquad
 \mathcal U(v)[F]=[\tau_{(v,0)}F],\qquad
 e^{i\tau H}\mathcal U(v)
   =\int e^{i(\tau E+v\cdot q)}\,d\mathcal P(E,q).
\end{equation}
\item $\mathcal P$ is carried by $\{(E,q):E\ge|q|\}$.
\item For some $C_\wp,\wp\ge0$ and $0<c_0\le1/2$,
put $\mathrm{Sl}_\eta=\{|x_3|<c_0\eta\}$.  For
$0<\eta\le1$, $h\in C_c^1(\mathrm{Sl}_\eta)$ and
$F,F'\in\mathbb C\mathfrak S(\{x_3>\eta\})$, assume
\begin{equation}\label{IV-75-14}
 |\sigma(\theta\bar F\cdot(h)\cdot F')|
 \le C_\wp\eta^{-\wp}\|h\|_1
       \|[\tau_{-\eta e_3}F]\|
       \|[\tau_{-\eta e_3}F']\|.
\end{equation}
\end{enumerate}
We refer to these four conditions as (Rec).
Vacuum uniqueness, $\ker H=\mathbb C\Omega$, will be a separate
assumption when needed.  The numerical triple in (Rec4) is
not part of the definition of the resulting field.

Write $\vec z=(z_1,z_2)$ and define
\begin{equation}\label{IV-75-5}
 d_\circ(z)=\operatorname{Re}z_3-|\operatorname{Im}\vec z|,
 \qquad \mathrm{Tb}_\circ=\{d_\circ>0\},\qquad
 \mathcal T(z)=\int e^{-z_3E+i\vec z\cdot q}\,d\mathcal P .
\end{equation}
The last integral is defined on the closed tube.
The cone bound gives
$|e^{-z_3E+i\vec z\cdot q}|\le e^{-d_\circ(z)E}$.
Consequently $\mathcal T$ consists of contractions, is strongly
continuous on the closed tube and norm holomorphic in its
interior, and satisfies
\[
 \mathcal T(z)\mathcal T(w)=\mathcal T(z+w),\qquad
 \mathcal T(z)^*=\mathcal T(z^*),\qquad
 z^*=(-\bar z_1,-\bar z_2,\bar z_3).
\]
For norm holomorphy, expand the exponential around a point
with positive margin; the factor $e^{-d_\circ(z)E}$ bounds
every coefficient and the quadratic remainder uniformly
on a smaller ball.  Dominated convergence against scalar
spectral measures proves strong continuity at the boundary.
In the interior, $E^ke^{-d_\circ(z)E}$ is bounded for every
$k$, so the range of $\mathcal T(z)$ lies in
\[
 \mathcal D_\infty=\bigcap_{k\ge1}\operatorname{dom}H^k .
\]
The vacuum is fixed by $\mathcal T(z)$ and its adjoint.
For Minkowski coordinates $x=(t,\vec x)$ put
\begin{equation}\label{IV-75-8}
 \iota(t,\vec x)=(\vec x,-it).
\end{equation}
This identifies the forward tube
$\mathbb R^{1+2}+i\{y:y^0>|\vec y|\}$ with
$\mathrm{Tb}_\circ$, and
$\mathcal T(\iota x)=e^{itH}\mathcal U(\vec x)$ is unitary.

\subsection{Densities of bounded insertions}

\begin{lemma}\label{IV-75-22}
For $0<\eta\le1$ there is a unique complex linear map
$h\mapsto B_\eta^{\rm in}(h)$ from $L^1(\mathrm{Sl}_\eta)$
to bounded operators such that
\begin{align}
 \|B_\eta^{\rm in}(h)\|&\le
                   C_\wp\eta^{-\wp}\|h\|_1,\label{IV-75-23}\\
 B_\eta^{\rm in}(h)[F]&=
       [(\tau_{\eta e_3}h)\cdot\tau_{2\eta e_3}F]
       \quad(h\in C_c^1(\mathrm{Sl}_\eta)).\label{IV-75-24}
\end{align}
For real $h$, its adjoint is $B_\eta^{\rm in}(\theta h)$.
Each scalar matrix element has an $L^\infty$ density of
norm at most $C_\wp\eta^{-\wp}\|\psi\|\|\psi'\|$.
\end{lemma}

\begin{proof}
Apply \eqref{IV-75-14} after translating both formal vectors
by $\eta e_3$.  The bounded sesquilinear form
\begin{equation}\label{IV-75-25}
 \langle[F],B_\eta^{\rm in}(h)[F']\rangle
   =\sigma(\theta\overline{\tau_{\eta e_3}F}
                         \cdot(h)\cdot\tau_{\eta e_3}F')
\end{equation}
vanishes on null classes and therefore defines a bounded
operator by Riesz representation.  Translate the whole tuple
by $\eta e_3$ to obtain \eqref{IV-75-24}.
The support inequalities $c_0\le1/2$ make that tuple
separated.  Real linearity and the norm bound extend this
map from real $C_c^1$ functions to real $L^1$ functions.

For fixed vectors the real and imaginary parts of the
matrix element are bounded real functionals on $L^1$, hence
have a complex density $b$.  The same constant bounds
$|b|$: apply the norm estimate to normalized indicators
of small balls and use Lebesgue differentiation.  Thus
$\|b\|_\infty\le C_\wp\eta^{-\wp}\|\psi\|\|\psi'\|$.
Complex linear extension now preserves this constant,
since $|\int hb|\le\|h\|_1\|b\|_\infty$.
Reflection invariance, symmetry and reality in
\eqref{IV-75-25} give the adjoint relation on $C_c^1$,
then on $L^1$ by approximation.  The class identity,
density of classes, and $L^1$ continuity prove uniqueness.
\end{proof}

\begin{proposition}[Point insertions]\label{IV-75-27}
There is a unique family of bounded operators $B(s,s')$,
$s,s'>0$, with the shift law
\begin{equation}\label{IV-75-28}
 T_tB(s,s')T_{t'}=B(s+t,s'+t')\qquad(t,t'\ge0)
\end{equation}
and the density identity
\begin{equation}\label{IV-75-29}
 \langle\psi,B_\eta^{\rm in}(h)\psi'\rangle
 =\int h(x)\langle\psi,\mathcal U(\vec x)
        B(\eta+x_3,\eta-x_3)\mathcal U(-\vec x)\psi'\rangle\,dx .
\end{equation}
It satisfies
\begin{equation}\label{IV-75-30}
 \|B(s,s')\|\le C_\wp\min(s,s',1)^{-\wp},
 \qquad B(s,s')^*=B(s',s).
\end{equation}
It is norm continuous and commutes with spatial orthogonal
transformations acting on $\mathcal H$.  The operator density
\[
 b_\eta(x)=\mathcal U(\vec x)
                  B(\eta+x_3,\eta-x_3)\mathcal U(-\vec x)
\]
is norm continuous on $\mathrm{Sl}_\eta$,
has norm at most $C_\wp\eta^{-\wp}$ there, and has
real-analytic scalar matrix elements.
\end{proposition}

\begin{proof}
For $0<\eta'<\eta\le1$, $h\in L^1(\mathrm{Sl}_{\eta'})$
and real $w$ with $|w_3|<c_0(\eta-\eta')$, translation
covariance of \eqref{IV-75-25} gives
\begin{equation}\label{IV-75-32}
 B_\eta^{\rm in}(\tau_wh)
 =\mathcal T((\eta-\eta')e_3+w)\,
      B_{\eta'}^{\rm in}(h)\,
        \mathcal T((\eta-\eta')e_3-w).
\end{equation}
The two semigroup arguments have positive third coordinate,
and the translated support stays in $\mathrm{Sl}_\eta$.
It suffices to prove the identity on $C_c^1$ tests and
dense classes; $L^1$ continuity then proves the assertion
as written.

Fix $\psi,\psi'$ and let $b$ be their scalar density at
margin $\eta$.  Take nonnegative continuous mollifiers
$j_\varepsilon$ of integral one, supported in
$B(0,\varepsilon)$, with $\varepsilon<c_0\eta'$.
The functions
\[
 b_\varepsilon(w)=
 \langle\psi,\mathcal T((\eta-\eta')e_3+w)
 B_{\eta'}^{\rm in}(j_\varepsilon)
 \mathcal T((\eta-\eta')e_3-w)\psi'\rangle
\]
are holomorphic and bounded by
$C_\wp\eta'^{-\wp}\|\psi\|\|\psi'\|$ on
\[
 \{w\in\mathbb C^3:
       |\operatorname{Re}w_3|+|\operatorname{Im}\vec w|
                                      <\eta-\eta'\}.
\]
On the smaller real slab
$|w_3|<c_0(\eta-\eta')$, \eqref{IV-75-32}
identifies them with $\int j_\varepsilon(x-w)b(x)\,dx$.
Montel's theorem gives a locally uniformly convergent
subsequence with holomorphic limit.  Mollification also
converges to $b$ in $L^1$ on compact real subsets, so this
limit agrees with $b$ almost everywhere.
As $\eta'\downarrow0$ the smaller real slabs exhaust
$\mathrm{Sl}_\eta$.  Their continuous representatives
agree on overlaps and give one real-analytic representative
of $b$ on the whole slab.

Linearity in $\psi'$ and antilinearity in $\psi$ hold
everywhere for these representatives: they hold almost
everywhere for the original densities, and both sides are
continuous.  The original $L^\infty$ bound likewise holds
at every point.  Riesz representation thus defines bounded
operators $b_\eta(x)$ with
$\|b_\eta(x)\|\le C_\wp\eta^{-\wp}$.
This constructs the operator density from scalar densities.

Integrate \eqref{IV-75-32} against arbitrary tests.
Uniqueness of the continuous scalar densities gives
\begin{equation}\label{IV-75-33}
 b_\eta(x+w)=\mathcal T((\eta-\eta')e_3+w)\,
                 b_{\eta'}(x)\,
                     \mathcal T((\eta-\eta')e_3-w).
\end{equation}
For $0<\eta\le\min(s,s',1)$ define
\begin{equation}\label{IV-75-34}
 B(s,s')=T_{s-\eta}\,b_\eta(0)\,T_{s'-\eta}.
\end{equation}
The case $x=w=0$ of \eqref{IV-75-33} proves independence
of $\eta$.  It also proves the shift law; choosing
$\eta=\min(s,s',1)$ gives the norm bound.
For $x\in\mathrm{Sl}_\eta$, choose $\eta'>0$ with
$|x_3|<c_0(\eta-\eta')$ and apply \eqref{IV-75-33} at
$(0,x)$.  Factoring spatial translations from the semigroup
gives the stated formula for $b_\eta(x)$ and hence
\eqref{IV-75-29}.

In \eqref{IV-75-33} choose a fixed smaller margin locally.
Its semigroup factors have positive margins and are norm
holomorphic, proving norm continuity of $b_\eta$.
The same reasoning in \eqref{IV-75-34} proves norm
continuity of $B$.  The insertion adjoint identity makes
$b_\eta(x)^*=b_\eta(\theta x)$, so $b_\eta(0)$ is
self-adjoint.  This proves the adjoint relation for $B$.
Spatial rotations in \eqref{IV-75-25} give covariance of
$b_\eta$ and invariance of $b_\eta(0)$; they commute with
$T_t$, proving the remaining assertion.

For uniqueness, any bounded-operator family with the
shift law has continuous scalar densities in
\eqref{IV-75-29}: factor it through $B(\eta/2,\eta/2)$.
The density identity then determines its value at $x=0$,
hence $B(\eta,\eta)$.  The shift law determines all other
values.  The same argument shows that changing the
admissible triple $C_\wp,\wp,c_0$ does not change the
operators, by first restricting insertions to the common
smaller slab.
\end{proof}

\subsection{Resolvent regularization of the point form}

Put $R_H=(1+H)^{-1}$.  For real $l>\wp$ define
\begin{equation}\label{IV-75-40}
 A_l=\Gamma(l)^{-2}\int_0^\infty\!\!\int_0^\infty
            (ss')^{l-1}e^{-s-s'}B(s,s')\,ds\,ds'.
\end{equation}

\begin{lemma}\label{IV-75-41}
The integral is norm integrable and defines a bounded
self-adjoint operator with
\begin{equation}\label{IV-75-42}
 \|A_l\|\le C_\wp\left(1+\frac{2\Gamma(l-\wp)}{\Gamma(l)}\right).
\end{equation}
For $a>0$ and $s,s'>0$,
\begin{equation}\label{IV-75-43}
 R_H^aA_lR_H^a=A_{l+a},\qquad
 T_sA_lT_{s'}=R_H^lB(s,s')R_H^l.
\end{equation}
Consequently
\begin{equation}\label{IV-75-49}
 q_0(\psi,\psi')=
       \langle(1+H)^l\psi,A_l(1+H)^l\psi'\rangle
\end{equation}
defines a Hermitian form on
$\bigcup_{l>\wp}\operatorname{dom}H^l$, independent of the
admissible exponent used for a given pair of vectors.
It satisfies the energy bound
\[
 |q_0(\psi,\psi')|\le
       \|A_l\|\,\|(1+H)^l\psi\|\,\|(1+H)^l\psi'\|
\]
on $\operatorname{dom}H^l$, and
\begin{equation}\label{IV-75-45}
 q_0(T_s\psi,T_{s'}\psi')=\langle\psi,B(s,s')\psi'\rangle.
\end{equation}
\end{lemma}

\begin{proof}
By \eqref{IV-75-30},
$\|B(s,s')\|\le C_\wp(1+s^{-\wp}+s'^{-\wp})$.
The two Gamma integrals therefore converge for $l>\wp$
and give \eqref{IV-75-42}.  Norm continuity makes the
operator integrand measurable.  Self-adjointness follows
by interchanging $s,s'$ and using $B(s,s')^*=B(s',s)$.

Functional calculus and the scalar Gamma integral give
\[
 R_H^a=\Gamma(a)^{-1}\int_0^\infty t^{a-1}e^{-t}T_t\,dt
\]
weakly for $a>0$.  Insert this on both sides of $A_l$
and use the shift law.  Fubini is justified by the
integrable bound just proved and contractivity.
Set $s=r+t$ and $s'=r'+t'$.  The two inner integrals are
\[
 \int_0^s t^{a-1}(s-t)^{l-1}\,dt
      =s^{a+l-1}\frac{\Gamma(a)\Gamma(l)}{\Gamma(a+l)}.
\]
This proves the first identity in \eqref{IV-75-43}.
For the second, shift $B(r,r')$ by $(s,s')$, then write
$B(r+s,r'+s')=T_rB(s,s')T_{r'}$ and use the same
resolvent integral.

If $l'>l>\wp$ and both vectors lie in
$\operatorname{dom}H^{l'}$, the first identity and
$R_H^{l'-l}(1+H)^{l'}=(1+H)^l$ prove equality of the
two definitions of $q_0$.  The domains are nested, so a
common admissible exponent exists for every pair in the
stated union.  Hermiticity and the bound follow from those
of $A_l$.  Finally $(1+H)^lT_s$ is bounded for $s>0$.
Apply the second identity at $(s/2,s'/2)$ and cancel
the resolvent factors against these bounded regularized
powers.  The shift law gives \eqref{IV-75-45}.
\end{proof}

\subsection{Recovering an insertion bound after centering}

The hypotheses (Rec) are also available for a functional with
nonzero one-point function.

\begin{lemma}\label{IV-75-50}
Suppose $\sigma$ satisfies (Rec1)--(Rec2), its spectrum obeys
$E\ge|q_1|$, and $\sigma(f)=c\int f$.  If its centered
functional satisfies
\begin{equation}\label{IV-75-51}
 |\sigma^{\rm c}(\theta\bar F\cdot(h)\cdot F')|
 \le C_\wp\eta^{-\wp}\|h\|_1
 \|[\tau_{-\eta e_3}F]\|_{\sigma^{\rm c}}
 \|[\tau_{-\eta e_3}F']\|_{\sigma^{\rm c}},
\end{equation}
with the supports in (Rec4), then $\sigma$ satisfies (Rec)
with constants $(C_\wp+|c|,\wp,c_0)$.
\end{lemma}

\begin{proof}
Spatial rotations conjugate the joint spectral measure, so
$E\ge|n\cdot q|$ for every unit spatial vector $n$.
A countable dense set of such vectors gives $E\ge|q|$.
The shift identities of Section~\ref{sec:reflection-hilbert}
give
$\sigma=\sigma^{\rm c}\circ\mathsf{sh}_c$ and an isometry
$[F]_\sigma\mapsto[\mathsf{sh}_cF]_{\sigma^{\rm c}}$.
The shift commutes with translations, conjugation and
reflection, preserves supports, and is multiplicative.
Writing $F_c=\mathsf{sh}_cF$, therefore,
\begin{equation}\label{IV-75-52}
 \sigma(\theta\bar F\cdot(h)\cdot F')
 =\sigma^{\rm c}(\theta\bar F_c\cdot(h)\cdot F'_c)
       +c\int h\,\sigma^{\rm c}(\theta\bar F_c\cdot F'_c).
\end{equation}
The first term has the bound \eqref{IV-75-51} with the
$\sigma$ norms.  Cauchy--Schwarz and contractivity of
$T_\eta$ bound the second by
$|c|\|h\|_1\|[\tau_{-\eta e_3}F]\|
\|[\tau_{-\eta e_3}F']\|$.
Since $\eta\le1$ and $\wp\ge0$, their sum has the required
bound.  Centering preserves the other stated symmetries
and reflection positivity by the same shift identities.
\end{proof}

\section{Analytic continuation of the Euclidean kernels}
\label{sec:reconstruction-kernels}

The point-insertion bound gives a product estimate without
a factorial.  Rotation of the ordering direction then covers
all distinct Euclidean points, including configurations
with equal third coordinates.

Call $(h_1,\ldots,h_n)$ \emph{$\eta$-ordered along $u$}
at heights $t_1,\ldots,t_n$ if $u\in S^2$, $0<\eta\le1$,
$t_{k+1}-t_k\ge2\eta$, and
\[
 \operatorname{supp}h_k
       \subset\{|u\cdot x-t_k|<c_0\eta\}.
\]
Such tuples are separated by at least $\eta$.
For $u=e_3$, set $g_k=\tau_{-t_ke_3}h_k$ and
$s_k=t_{k+1}-t_k-2\eta\ge0$.

\begin{proposition}\label{IV-75-61}
For an $\eta$-ordered tuple along $e_3$,
\begin{equation}\label{IV-75-62}
 \sigma(h_1,\ldots,h_n)=
 \langle\Omega,B_\eta^{\rm in}(g_1)T_{s_1}
       B_\eta^{\rm in}(g_2)\cdots
       T_{s_{n-1}}B_\eta^{\rm in}(g_n)\Omega\rangle.
\end{equation}
If $t_1\ge\eta$, then
\begin{equation}\label{IV-75-63}
 [(h_1,\ldots,h_n)]=T_{t_1-\eta}
 B_\eta^{\rm in}(g_1)T_{s_1}\cdots
 B_\eta^{\rm in}(g_n)\Omega,\qquad
 \|[(h_1,\ldots,h_n)]\|
       \le\prod_k C_\wp\eta^{-\wp}\|h_k\|_1.
\end{equation}
\end{proposition}

\begin{proof}
Starting at the right, \eqref{IV-75-24} proves by downward
induction that the product beginning with
$B_\eta^{\rm in}(g_k)$ represents
$[\tau_{(\eta-t_k)e_3}(h_k,\ldots,h_n)]$.
Indeed the intervening translation has height
$s_k+\eta-t_{k+1}=-t_k-\eta$, and the insertion
adds $2\eta$ to the remaining tuple.
All these classes have supports in
$\{x_3>(1-c_0)\eta\}$.
Taking the vacuum component and using translation
invariance proves \eqref{IV-75-62}.
If $t_1\ge\eta$, translate the class forward by
$t_1-\eta$ to obtain \eqref{IV-75-63};
contractivity and \eqref{IV-75-23} give its norm bound.
\end{proof}

For later use, the scalar density formula can be iterated:
for $g_k\in L^1(\mathrm{Sl}_\eta)$ and bounded $A_0,\ldots,A_n$,
\begin{equation}\label{IV-75-60}
 \begin{split}
 &\langle\psi,A_0B_\eta^{\rm in}(g_1)A_1\cdots
                 B_\eta^{\rm in}(g_n)A_n\psi'\rangle\\
 &\qquad=\int\prod_k g_k(y_k)
       \langle\psi,A_0b_\eta(y_1)A_1\cdots
                          b_\eta(y_n)A_n\psi'\rangle\,dy .
 \end{split}
\end{equation}
The second factor is continuous and bounded by
$(C_\wp\eta^{-\wp})^n\|\psi\|\|\psi'\|\prod_j\|A_j\|$.
Thus the integrand is absolutely integrable.  Apply
\eqref{IV-75-29} to the first insertion, then to the
remaining ones; Fubini proves the formula.

\subsection{Ordered complex configurations}

For $\zeta_1,\ldots,\zeta_{n-1}\in\mathrm{Tb}_\circ$ define
\begin{equation}\label{IV-75-66}
 \mathbf W_n(\zeta_1,\ldots,\zeta_{n-1})
 =\langle\Omega,B(\eta,\eta)\mathcal T(\zeta_1-2\eta e_3)
          B(\eta,\eta)\cdots
          \mathcal T(\zeta_{n-1}-2\eta e_3)
          B(\eta,\eta)\Omega\rangle,
\end{equation}
where $0<\eta\le\min(1,\frac12\min_k d_\circ(\zeta_k))$.
For $n=1$ this is the real number
$\mathbf W_1=\langle\Omega,B(1,1)\Omega\rangle$.
On
\[
 \mathrm{Tb}^{(n)}_\circ
 =\{z:z_1\in\mathrm{Tb}_\circ,\
               z_{k+1}-z_k\in\mathrm{Tb}_\circ\}
\]
put
\begin{align}
 \eta(z)&=\min\left(1,d_\circ(z_1),
                     \tfrac12\min_{k<n}d_\circ(z_{k+1}-z_k)\right),
                                                        \label{IV-75-67}\\
 \Psi_n(z)&=\mathcal T(z_1-\eta e_3)B(\eta,\eta)
       \mathcal T(z_2-z_1-2\eta e_3)B(\eta,\eta)\cdots
       \mathcal T(z_n-z_{n-1}-2\eta e_3)B(\eta,\eta)\Omega,
                                                        \label{IV-75-68}
\end{align}
using any $0<\eta\le\eta(z)$; the inner minimum is omitted
when $n=1$.

These definitions are independent of $\eta$.
To compare $\eta'<\eta$, substitute
\[
 B(\eta,\eta)=T_{\eta-\eta'}B(\eta',\eta')T_{\eta-\eta'}
\]
at every insertion.  The semigroup law absorbs the adjacent
factors into the differences, and $T_t\Omega=\Omega$
removes the final factor.  For the scalar product it removes
the initial factor as well.

\begin{proposition}\label{IV-75-71}
The function $\mathbf W_n$ is holomorphic and $\Psi_n$ is
holomorphic in norm.  They satisfy
\begin{align}
 |\mathbf W_n(\zeta)|&\le(C_\wp\eta^{-\wp})^n
       &&\text{if }d_\circ(\zeta_k)\ge2\eta,\ 0<\eta\le1,
                                                        \label{IV-75-72}\\
 \|\Psi_n(z)\|&\le(C_\wp\eta(z)^{-\wp})^n,                 \label{IV-75-73}\\
 \mathcal T(w)\Psi_n(z)&=\Psi_n(z_1+w,\ldots,z_n+w),\notag\\
 \langle\Omega,\Psi_n(z)\rangle
       &=\mathbf W_n(z_2-z_1,\ldots,z_n-z_{n-1})           \label{IV-75-74}
\end{align}
for $w\in\overline{\mathrm{Tb}_\circ}$.
For ordered tuples,
\begin{equation}\label{IV-75-76}
 \sigma(h_1,\ldots,h_n)
 =\int\mathbf W_n(x_2-x_1,\ldots,x_n-x_{n-1})
                              \prod_k h_k(x_k)\,dx.
\end{equation}
If also $t_1\ge\eta$, then
\begin{equation}\label{IV-75-77}
 \langle\psi,[(h_1,\ldots,h_n)]\rangle
       =\int\langle\psi,\Psi_n(x)\rangle\prod_k h_k(x_k)\,dx.
\end{equation}
\end{proposition}

\begin{proof}
Contractivity and \eqref{IV-75-30} give the bounds.
Locally choose a fixed $\eta$ strictly smaller than every
allowed margin.  All semigroup factors are then norm
holomorphic; the independence of $\eta$ proves holomorphy
of the defined functions.  The semigroup law and the
vacuum identities give \eqref{IV-75-74}.

The density can be written with half the margin:
\begin{equation}\label{IV-75-75}
 b_\eta(y)=\mathcal T(y+\tfrac\eta2e_3)
        B(\tfrac\eta2,\tfrac\eta2)
                  \mathcal T(-y+\tfrac\eta2e_3).
\end{equation}
Both arguments have positive margin, since
$|y_3|<c_0\eta\le\eta/2$.  This follows directly from
the shift law and the formula for $b_\eta$.
Insert \eqref{IV-75-75} in \eqref{IV-75-60}, applied to
the products \eqref{IV-75-62}--\eqref{IV-75-63}.
With $x_k=y_k+t_ke_3$, consecutive semigroup factors
combine into
$\mathcal T(x_{k+1}-x_k-\eta e_3)$.
At the right the factor fixes $\Omega$.
At the left it either disappears in the vacuum component
or becomes $\mathcal T(x_1-\eta e_3/2)$.
These are precisely \eqref{IV-75-66} and
\eqref{IV-75-68} with margin $\eta/2$, proving the
two integral formulas.  The supports are compact and
stay inside the respective open tubes, so their
continuous integrands are integrable.
\end{proof}

\subsection{Changing the ordering direction}

For $u\in S^2$ put
\[
 d_\circ^u(z)=\operatorname{Re}(u\cdot z)
       -|\operatorname{Im}(z-(u\cdot z)u)|,\qquad
 \mathrm{Tb}^{u,\pi}_n
 =\{z:d_\circ^u(z_{\pi(k+1)}-z_{\pi(k)})>0,\ k<n\}.
\]
If $R\in\mathrm{SO}(3)$ satisfies $Re_3=u$, then
$d_\circ^u(z)=d_\circ(R^{-1}z)$.
Each tube is open and convex; it contains $\operatorname{Re}z$
whenever it contains $z$.  For $n=1$ it is all of $\mathbb C^3$.
Define on it
\begin{equation}\label{IV-75-69}
 \mathbf S_n^{R,\pi}(z)
 =\mathbf W_n\bigl(R^{-1}(z_{\pi(2)}-z_{\pi(1)}),\ldots,
                   R^{-1}(z_{\pi(n)}-z_{\pi(n-1)})\bigr).
\end{equation}

Any $n$ distinct real points admit a separating direction.
More quantitatively, if their mutual distances are at least
$\delta>0$, there is $u\in S^2$ with
\begin{equation}\label{IV-75-53}
 |u\cdot(x_i-x_j)|\ge\delta/n^2\quad(i\ne j).
\end{equation}
Indeed the spherical zone
$|u\cdot (x_i-x_j)/|x_i-x_j||<n^{-2}$ has area $4\pi/n^2$.
The union over the $\binom n2$ pairs has area less than
$2\pi$, so cannot cover $S^2$.
Ordering the projections shows that the real parts of the
tubes cover $(\mathbb R^3)^n_{\ne}$.

\begin{proposition}\label{IV-75-78}\label{IV-75-83}
There is a unique continuous kernel $\mathsf s_n$ on
$(\mathbb R^3)^n_{\ne}$ such that
\begin{equation}\label{IV-75-79}
 \sigma(f_1,\ldots,f_n)=\int\mathsf s_n(x)\prod_i f_i(x_i)\,dx
\end{equation}
for all separated tuples.  It is real analytic, symmetric
and invariant under Euclidean isometries.
It extends uniquely to a holomorphic function $\mathbf S_n$
on
$\mathrm{Tb}^{\rm all}_n=\bigcup_{u,\pi}\mathrm{Tb}^{u,\pi}_n$,
whose restriction to each tube is \eqref{IV-75-69},
independently of the choice of $R$.  This extension is
symmetric and obeys
\begin{equation}\label{IV-75-85}
 |\mathbf S_n(z)|\le(C_\wp\eta^{-\wp})^n
 \quad\text{if }d_\circ^u(z_{\pi(k+1)}-z_{\pi(k)})
                  \ge2\eta,\quad0<\eta\le1.
\end{equation}
In particular,
\begin{equation}\label{IV-75-80}
 \mathsf s_n(x)=\mathbf W_n(x_2-x_1,\ldots,x_n-x_{n-1})
                  \quad(x_{13}<\cdots<x_{n3}),
 \qquad \mathsf s_1=\mathbf W_1 .
\end{equation}
\end{proposition}

\begin{proof}
For a tuple ordered along $u$, rotate it by $R^{-1}$,
permute its entries, and apply \eqref{IV-75-76}.
Euclidean invariance and symmetry give the kernel formula
with $\mathbf S_n^{R,\pi}$.
If a real configuration belongs to two ordering regions,
choose sufficiently small balls about its entries so that
every product of tests in those balls is ordered both ways,
with one common positive margin.  The two kernel formulas
have the same value on every such product.  Continuous
functions are determined by these product tests, for example
by taking products of approximate identities.
Thus the kernels agree on each real overlap.

The real ordering regions cover all distinct points by
\eqref{IV-75-53}.  They therefore define one real-analytic
function $\mathsf s_n$.  To obtain \eqref{IV-75-79} for
arbitrary separated supports, let their least distance be
$d>0$ and choose
$\eta=\min(d/(4n^2),1)$.
Partition each test into finitely many pieces supported
in balls of radius $c_0\eta/2$.  Choose a point of the
original support in every nonzero piece.
Equation \eqref{IV-75-53} orders these chosen points with
gaps at least $4\eta$; each piece is supported in a
$c_0\eta$ slab about its chosen point.
Each resulting product tuple is therefore ordered.
Apply the formula to each product and sum by multilinearity.
For $n=1$ the same partition uses $\eta=1$.
This proves the formula on every separated tuple.
Products of arbitrarily small supported tests also prove
uniqueness.  Applying uniqueness to permuted or isometrically
transported tests gives the claimed real symmetries.

A nonempty complex overlap of two tubes is convex, and
contains the real part of each of its points.
It consequently has a nonempty open real part.
The two holomorphic pieces agree there by the preceding
argument.  The identity theorem on the connected overlap
makes them equal throughout it.  They therefore glue to
one holomorphic function on the union.
The same identity theorem on each tube proves uniqueness
and independence of $R$; it also extends the permutation
symmetry.  The bound is \eqref{IV-75-72} after rotation.
All constructions depend on $\sigma$, not on the chosen
admissible numerical triple.
\end{proof}

\subsection{Growth on tests flat at the diagonals}

The extension to the full space of tests used for Schwinger
functions requires a bound near every diagonal, not only
for products with fixed separated supports.

\begin{proposition}\label{IV-75-87}\label{IV-75-90}
Put $k_{\rm E}=\lceil\wp\rceil+4$ and
$C_{\rm E}=C_\wp2^\wp e^{2\wp}3^{k_{\rm E}/2}\pi^2$.
Then $|\mathsf s_1|\le C_\wp$, and for $n\ge2$,
\begin{equation}\label{IV-75-88}
 |\mathsf s_n(x)|\le
 C_\wp^n(2n^2)^{\wp n}
            \min(1,\min_{i\ne j}|x_i-x_j|)^{-\wp n}.
\end{equation}
For every $\varphi\in\mathcal S_{\ne}(\mathbb R^{3n})$,
\begin{equation}\label{IV-75-89}
 \int|\mathsf s_n(x)\varphi(x)|\,dx
       \le2C_{\rm E}^n(n!)^{2\wp}|\varphi|_{k_{\rm E}n}.
\end{equation}
Consequently the Schwinger functions on this space are
\begin{equation}\label{IV-75-92}
 \mathsf S_n^\sigma(\varphi)=\int\mathsf s_n(x)\varphi(x)\,dx.
\end{equation}
\end{proposition}

\begin{proof}
For a configuration with least Euclidean distance $\rho$,
choose the direction in \eqref{IV-75-53} and use
\eqref{IV-75-85} with $\eta=\min(\rho/(2n^2),1)$.
This proves \eqref{IV-75-88}; the one-point bound follows
from $B(1,1)$.

Let $k=k_{\rm E}n$, $N=\lceil\wp n\rceil$.
Then $N\le k$ and $k\ge4n$.
If $\rho\le1$ and $N\ge1$, move a closest pair to its
midpoint, obtaining a diagonal point $x^*$ with
$|x-x^*|\le1$ and $|x-x^*|_1\le\sqrt3\rho$.
Flatness at $x^*$ and Taylor's remainder, with the
weighted seminorm of Section~\ref{sec:distribution-extension},
give
\[
 |\varphi(x)|\le
 \frac{(\sqrt3\rho)^N}{N!}\,3^{k/2}
       |\varphi|_k(1+|x|^2)^{-k/2}
 \le2\,3^{k/2}\rho^N
       |\varphi|_k(1+|x|^2)^{-k/2}.
\]
The factor $3^{k/2}$ accounts for moving at most unit
distance along the Taylor segment.
Since $N\ge\wp n$, this absorbs the singular factor.
When $\rho>1$ or $\wp=0$, the unmodified seminorm bound
suffices.  Finally
\[
 n^n\le e^nn!,\qquad
 \int_{\mathbb R^{3n}}(1+|x|^2)^{-k/2}\,dx
 \le\prod_{i=1}^n\int_{\mathbb R^3}(1+|x_i|^2)^{-2}\,dx_i
 =\pi^{2n}.
\]
Combining these estimates proves \eqref{IV-75-89}.
The integral functional agrees with $\sigma$ on separated
smooth tensors by \eqref{IV-75-79}.
The tensor-density and uniqueness result of
Section~\ref{sec:distribution-extension} identifies it
with the Schwinger function.
\end{proof}

\section{Smeared fields and their boundary distributions}
\label{sec:reconstruction-fields}

We continue under (Rec).  Minkowski space has coordinates
$x=(t,\vec x)$, product $x\cdot y=t s-\vec x\cdot\vec y$,
and forward cone $V_+=\{t>|\vec x|\}$.
Write $U(x)=\mathcal T(\iota x)$ and fix an integer $l>\wp$.
For $f\in\mathcal S(\mathbb R^{1+2})$ put
\begin{equation}\label{IV-76-5}
 c_N(f)=\|A_l\|\sum_{k=0}^N\binom Nk\|\partial_t^kf\|_1.
\end{equation}
The Schwartz seminorms used above give
$\|\partial^\alpha F\|_1\le\pi^{2n}|F|_{|\alpha|+4n}$
on $\mathbb R^{3n}$, by comparison with
$\prod_i(1+|x_i|^2)^{-2}$.  In particular
\begin{equation}\label{IV-76-15}
 c_N(f)\le\|A_l\|2^N\pi^2|f|_{N+4}.
\end{equation}

We use the graph seminorms $\|(1+H)^j\psi\|$ on
$\mathcal D_\infty$.
The spectral projections $P_N=\mathcal P(\{E\le N\})$
have ranges in $\mathcal D_\infty$ and converge strongly
to the identity.  Bounded joint spectral multipliers
preserve every $\operatorname{dom}H^j$ and commute there
with $(1+H)^j$.
We will also use the following domain criterion:
\begin{equation}\label{eq:iv-domain-criterion}
 |\langle(1+H)^j\psi,\chi\rangle|\le C\|\psi\|
 \quad(\psi\in\mathcal D_\infty)
 \ \Longrightarrow\
 \chi\in\operatorname{dom}H^j,\quad
 \|(1+H)^j\chi\|\le C.
\end{equation}
Indeed, insert $\psi=(1+H)^jP_N\chi$ to bound
$\int_{E\le N}(1+E)^{2j}\,dm_\chi$ by $C^2$,
then let $N\to\infty$.

\subsection{Construction on a common domain}

\begin{proposition}\label{IV-76-25}
There is a unique complex linear map
$f\mapsto\phi(f)$ from $\mathcal S(\mathbb R^{1+2})$
to linear operators $\mathcal D_\infty\to\mathcal D_\infty$
with
\begin{equation}\label{IV-76-26}
 \langle\psi,\phi(f)\psi'\rangle
   =\int f(x)\,q_0(U(-x)\psi,U(-x)\psi')\,dx.
\end{equation}
It is independent of $l$ and of the numerical triple in
(Rec).  For $j\ge0$,
\begin{equation}\label{IV-76-28}
 \|(1+H)^j\phi(f)\psi\|
     \le c_{l+j}(f)\|(1+H)^{2l+j}\psi\|.
\end{equation}
On $\mathcal D_\infty$ it satisfies
\begin{align}
 \langle\psi,\phi(f)\psi'\rangle
    &=\langle\phi(\bar f)\psi,\psi'\rangle,\nonumber\\
 U(a)\phi(f)U(-a)&=\phi(f(\cdot-a)),\qquad
 [H,\phi(f)]=i\phi(\partial_tf).                         \label{IV-76-27}
\end{align}
\end{proposition}

\begin{proof}
The integrand in \eqref{IV-76-26} equals
\[
 f(x)\langle U(-x)(1+H)^l\psi,
                    A_lU(-x)(1+H)^l\psi'\rangle.
\]
Its second factor is bounded and continuous.
Denote the resulting sesquilinear form by $q_f$.
It obeys
\begin{equation}\label{IV-76-31}
 |q_f(\psi,\psi')|\le
 \|f\|_1\|A_l\|\|(1+H)^l\psi\|\|(1+H)^l\psi'\|.
\end{equation}
To turn this form into an operator, move energy powers
from the first vector to the second.
The map $t\mapsto e^{-itH}\xi$ is norm differentiable
for $\xi\in\operatorname{dom}H$, with derivative
$-ie^{-itH}H\xi$, by dominated convergence in the spectral
calculus.  Differentiating both vectors in the preceding
matrix element, then integrating by parts in $t$, gives
\[
 q_f(H\psi,\psi')-q_f(\psi,H\psi')
                   =i q_{\partial_tf}(\psi,\psi').
\]
The boundary terms vanish because $f$ is Schwartz and
the matrix element and its derivative are bounded.
Iteration yields
\begin{equation}\label{IV-76-33}
 q_f((1+H)^N\chi,\psi')
 =\sum_{k=0}^N\binom Nk i^k
                 q_{\partial_t^kf}(\chi,(1+H)^{N-k}\psi').
\end{equation}
Choose $\chi=R_H^l\psi$ and $N=l+j$.
Using \eqref{IV-76-31} and monotonicity of graph norms,
\begin{equation}\label{IV-76-34}
 |q_f((1+H)^j\psi,\psi')|
 \le c_{l+j}(f)\|\psi\|\|(1+H)^{2l+j}\psi'\|.
\end{equation}
At $j=0$, Riesz representation gives a unique vector
$\phi(f)\psi'$.  For every $j$, the domain criterion
\eqref{eq:iv-domain-criterion} gives its membership in
$\operatorname{dom}H^j$ and the bound \eqref{IV-76-28}.
Thus it belongs to $\mathcal D_\infty$.

Linearity, uniqueness and independence of the auxiliary
choices follow from the defining form.
Hermiticity follows from that of $q_0$.
Changing variables in \eqref{IV-76-26} proves translation
covariance.  The identity obtained by integration by parts,
tested against dense vectors in $\mathcal D_\infty$,
proves the commutator formula.
\end{proof}

\subsection{Regularized products and their continuity}

\begin{lemma}\label{IV-76-35}
\enlargethispage{2\baselineskip}
Let $w,w'\in\mathrm{Tb}_\circ$ and
$0<s\le\min(1,d_\circ(w),d_\circ(w'))$.
Then
\begin{equation}\label{IV-76-36}
 \mathcal T(w)\phi(f)\mathcal T(w')
 =\int f(x)\mathcal T(w-se_3+\iota x)
             B(s,s)\mathcal T(w'-se_3-\iota x)\,dx
\end{equation}
in matrix elements on $\mathcal H$.
The integrand after removal of $f$ is strongly continuous
and bounded by $C_\wp s^{-\wp}$.
For $w_1,\ldots,w_n\in\mathrm{Tb}_\circ$, set
$z_k(x)=\sum_{r\le k}w_r+\iota x_k$.  Then
\begin{equation}\label{IV-76-37}
 \begin{split}
 &\langle\psi,\mathcal T(w_1)\phi(f_1)\cdots
                    \mathcal T(w_n)\phi(f_n)\Omega\rangle\\
 &\qquad=\int\prod_k f_k(x_k)\langle\psi,\Psi_n(z(x))\rangle\,dx .
 \end{split}
\end{equation}
\end{lemma}

\begin{proof}
The positive margins put
$\mathcal T(w)^*\psi,\mathcal T(w')\psi'$ in
$\mathcal D_\infty$.
Apply \eqref{IV-76-26} and split $T_s$ from each
semigroup.  Equation \eqref{IV-75-45} changes the form
between these two factors into $B(s,s)$; commuting the
remaining spectral multipliers gives \eqref{IV-76-36}.
The margin is unchanged by addition of $\iota x$, so all
remaining factors are contractions.  Strong continuity
and the point bound give the stated properties.

For the product choose
$s=\min(1,d_\circ(w_1),\frac12\min_{k\ge2}d_\circ(w_k))$.
The inner minimum is omitted when $n=1$.
Split each intermediate semigroup into equal halves, and
append $T_s$ at the right, where it fixes $\Omega$.
This writes the product as a product of bounded
regularized fields, each covered by \eqref{IV-76-36}.
Successively insert those formulas.
The uniform point bounds and $\prod_k|f_k|\in L^1$
justify Fubini.  Consecutive semigroup arguments combine
to $z_{k+1}(x)-z_k(x)-2se_3$; the first is
$z_1(x)-se_3$, and the final spatial-time translation
fixes $\Omega$.  The resulting product is
$\Psi_n(z(x))$ with margin $s$, proving
\eqref{IV-76-37}.
\end{proof}

\begin{lemma}\label{IV-76-38}
Suppose $c_N(f_k^\epsilon-f_k^0)\to0$ for every $N$,
and $C_k^\epsilon$ are contractive joint spectral
multipliers converging strongly to $C_k^0$.
Then
\[
 C_1^\epsilon\phi(f_1^\epsilon)\cdots
 C_n^\epsilon\phi(f_n^\epsilon)\Omega
 \longrightarrow
 C_1^0\phi(f_1^0)\cdots C_n^0\phi(f_n^0)\Omega
\]
in every graph norm of $H$.
\end{lemma}

\begin{proof}
Work from right to left.
If $\chi^\epsilon\to\chi^0$ in all graph norms, then
\[
 \begin{split}
 &\|(1+H)^j(\phi(f^\epsilon)\chi^\epsilon
                                      -\phi(f^0)\chi^0)\|\\
 &\quad\le c_{l+j}(f^\epsilon-f^0)
                       \|(1+H)^{2l+j}\chi^\epsilon\|
       +c_{l+j}(f^0)
                       \|(1+H)^{2l+j}(\chi^\epsilon-\chi^0)\|
       \longrightarrow0.
 \end{split}
\]
The spectral multipliers commute with the graph powers.
Their strong convergence and uniform norm bound preserve
the just-proved convergence.  Induction starts at
$\chi^\epsilon=\Omega$.
\end{proof}

\subsection{Tempered boundary values}

Here we give the boundary argument explicitly; pointwise
polynomial growth near the tube boundary alone does not
justify passing a limit through a test integral.

\begin{lemma}\label{IV-76-23}
Let $\mathbf F$ be holomorphic on an open set containing
\[
 z(x;\eta)=(\iota x_j+\eta h_je_3)_{j=1}^n,
 \qquad x\in K,\quad0<\eta\le\eta_0\le1,
\]
where $K$ is compact, and suppose that for real $b,p\ge0$,
$|\mathbf F(z(x;\eta))|\le b\eta^{-p}$.
For $g\in C_c^\infty$ supported in $K$, put
$I_g(\eta)=\int g(x)\mathbf F(z(x;\eta))\,dx$ and
$D_h=\sum_jh_j\partial_{t_j}$.
If the integer $N\ge p+1/2$, then
\begin{equation}\label{eq:iv-boundary-estimate}
 \sup_{0<\eta\le\eta_0}|I_g(\eta)|
 \le 2^N\eta_0^{-(N-1/2)}
                          b\sum_{r=0}^N\|D_h^rg\|_1.
\end{equation}
The limit at $\eta=0$ exists.
\end{lemma}

\begin{proof}
On each compact interval of positive heights,
holomorphy permits differentiation under the integral.
The Cauchy--Riemann equations and integration by parts give
$I_g'=I_{-iD_hg}$.
Set $J_r=I_{(-iD_h)^rg}$.
For $0\le r\le N$, these functions obey
$J_r'=J_{r+1}$ when $r<N$, and
$|J_r(\eta)|\le b\|D_h^rg\|_1\eta^{-(N-1/2)}$.
Integrate successively from $\eta_0$ toward $\eta$.
Every integration of $\eta^{-a}$ with $a\ge3/2$
lowers the singular exponent by one at a cost at most
$2$; after $N-1$ integrations the exponent is $1/2$.
Its final integral is bounded by $2$.
The values at $\eta_0$ contribute the same bound with
the corresponding $\|D_h^rg\|_1$, giving
\eqref{eq:iv-boundary-estimate}.
Apply that estimate once more to $-iD_hg$.
Then $I_g'$ is bounded uniformly down to zero,
so $I_g$ is Lipschitz there and has a limit.
\end{proof}

\begin{proposition}\label{prop:iv-wightman-distributions}
There is a unique tempered distribution
$\mathcal W_n^\sigma$ on $\mathbb R^{3n}$ satisfying
\begin{equation}\label{IV-76-29}
 \mathcal W_n^\sigma(f_1\otimes\cdots\otimes f_n)
 =\langle\Omega,\phi(f_1)\cdots\phi(f_n)\Omega\rangle .
\end{equation}
For arbitrary $\alpha_1,\ldots,\alpha_{n-1}\in V_+$,
\begin{equation}\label{IV-76-30}
 \begin{split}
 &\mathcal W_n^\sigma(f_1\otimes\cdots\otimes f_n)\\
 &\quad=\lim_{\epsilon\downarrow0}\int\prod_k f_k(x_k)
 \mathbf W_n\bigl(\iota(x_2-x_1+i\epsilon\alpha_1),\ldots,
             \iota(x_n-x_{n-1}+i\epsilon\alpha_{n-1})\bigr)\,dx .
 \end{split}
\end{equation}
For $n=1$, $\mathcal W_1^\sigma(f)=\mathbf W_1\int f$.
\end{proposition}

\begin{proof}
First take the vacuum component of \eqref{IV-76-37},
with $w_1=e_3$ and $w_{k+1}=\iota(i\epsilon\alpha_k)$.
It identifies the integral on the right of
\eqref{IV-76-30} with
\[
 \langle\Omega,\phi(f_1)\mathcal T(w_2)
                  \phi(f_2)\cdots\mathcal T(w_n)\phi(f_n)\Omega\rangle.
\]
Its integrand is bounded by
$(C_\wp\eta^{-\wp})^n\prod_k|f_k(x_k)|$, with
$2\eta=\min(2,\epsilon\min_k(\alpha_k^0-|\vec\alpha_k|))$.
The spectral contractions tend strongly to the identity.
Lemma~\ref{IV-76-38} therefore proves the boundary formula
on tensors of Schwartz functions.

To obtain a distribution on general tests, apply
Lemma~\ref{IV-76-23} to
$\mathbf F(z)=\mathbf W_n(z_2-z_1,\ldots,z_n-z_{n-1})$
with $h_j=2j$, $\eta_0=1$, $b=C_\wp^n$,
$p=\wp n$, and $N=\lceil\wp n\rceil+1$.
Every consecutive margin is $2\eta$, so the hypothesis
follows from \eqref{IV-75-72}.
The resulting limit functional on $C_c^\infty$
has bound
\[
 |\mathcal L(F)|\le
 2^N(N+1)C_\wp^n(n(n+1))^N\pi^{2n}|F|_{N+4n}.
\]
Here $\sum_jh_j=n(n+1)$ bounds the coefficients in
$D_h^r$, and the $L^1$ derivative estimate gives the
Schwartz norm.

For a smooth cutoff $\chi$ equal to one near the origin,
$\chi(\cdot/R)F\to F$ in every Schwartz seminorm.
Thus the bound extends $\mathcal L$ continuously to
$\mathcal S(\mathbb R^{3n})$.
It agrees with \eqref{IV-76-29} first on compactly
supported tensors by the boundary identity, then on
Schwartz tensors by Lemma~\ref{IV-76-38}.
For uniqueness, finite sums of compactly supported smooth
tensors are dense in each fixed Schwartz seminorm:
cut off the test, convolve with a product mollifier,
and approximate the resulting compact integral by
Riemann sums of translates of that product mollifier.
All derivatives up to the fixed order converge uniformly
on a common compact set.  A tempered functional vanishing
on those tensors must therefore vanish.
\end{proof}

\subsection{Cyclicity}

Let $\mathcal D_0$ be the span of $\Omega$ and all
$\phi(f_1)\cdots\phi(f_n)\Omega$.
It lies in $\mathcal D_\infty$, is invariant under the
fields and translations, and is dense.

To prove the last assertion, let $\psi\perp\mathcal D_0$.
Translation covariance first gives
\[
 \langle\psi,U(a_1)\phi(f_1)\cdots
                       U(a_n)\phi(f_n)\Omega\rangle=0 .
\]
Replace the translations successively by
$\mathcal T(w_j)$, $w_j\in\mathrm{Tb}_\circ$.
For one replacement, move all preceding regularized
factors to the first vector using Hermiticity.
The first positive-margin semigroup puts that vector
in $\mathcal D_\infty$, so every move is legitimate.
The resulting expression is $\langle\chi,U(a)\chi'\rangle$.
Its vanishing for all $a$ says that the Fourier transform
of the finite joint spectral measure $m_{\chi,\chi'}$
vanishes.  Uniqueness of Fourier transforms of finite
measures gives $m_{\chi,\chi'}=0$, and hence
$\langle\chi,\mathcal T(w)\chi'\rangle=0$.
Induction proves vanishing for every fully regularized
product.

By \eqref{IV-76-37} and product approximate identities,
$\langle\psi,\Psi_n(z)\rangle=0$ throughout
$\mathrm{Tb}^{(n)}_\circ$.
Equation \eqref{IV-75-77} then makes $\psi$ orthogonal to
classes of time-ordered positive-time tuples.
We must still approximate arbitrary positive-time tuples;
their supports need not be time ordered.

Fix such a tuple $X=(f_1,\ldots,f_n)$.
Partition its entries with nonnegative functions
$\chi_{h,k}$ supported in cubes of side $2h$ about $hk$,
summing to one.  Call a product of pieces good when all
their center heights differ by at least $4h/c_0$.
For small $h$, every good product, after permutation,
is ordered with $\eta=2h/c_0$ and first height at least
$\eta$.  Its class is orthogonal to $\psi$.

Let $\xi_h$ be the sum of the remaining classes.
On the compact set consisting of the original positive
supports and their reflected negative supports,
the continuous kernel $\mathsf s_{2n}$ is bounded by
some $M_X$.  The kernel formula gives
\begin{equation}\label{IV-76-43}
 \|\xi_h\|^2\le M_X
 \left(\int\prod_i|f_i(x_i)|
       \sum_{\mathbf k\ {\rm bad}}\prod_i\chi_{h,k_i}(x_i)\,dx\right)^2 .
\end{equation}
The sum is at most one and vanishes unless some two
heights differ by less than $(4/c_0+2)h$.
Its integral is $O_X(h)$: integrate one height first
over an interval of that length, bound that test by
its supremum norm on its fixed spatial support,
then integrate the remaining tests in $L^1$.
Thus $\|\xi_h\|\to0$ and
$\langle\psi,[X]\rangle=\langle\psi,\xi_h\rangle=0$.
All tuple classes span a dense subspace of the reflection
Hilbert space, so $\psi=0$ and $\mathcal D_0$ is dense.

The translation representation already has joint spectrum
$E\ge|q|$.  In the convention
$U(x)=\int e^{ip\cdot x}\,d\mathsf E(p)$ its momentum
measure is the image of $\mathcal P$ under
$(E,q)\mapsto(E,-q)$, and is carried by
$\overline{V_+}$.  If $\ker H=\mathbb C\Omega$, the only
translation-invariant vectors are multiples of $\Omega$:
invariance under $e^{itH}$ forces spectral support at
$E=0$.  No vacuum-uniqueness assertion is used in the
preceding construction.

\section{Lorentz covariance and locality}
\label{sec:reconstruction-covariance}

Let $\mathcal L=\mathrm{SO}^+(1,2)$.
Every element is a product of spatial rotations and the boosts
\[
 B(s)(t,x_1,x_2)
   =(t\cosh s+x_1\sinh s,\ x_1\cosh s+t\sinh s,\ x_2).
\]
For completeness, $\Lambda e_0$ is a future unit timelike
vector.  A spatial rotation and a boost send $e_0$ to it.
Their inverse followed by $\Lambda$ fixes $e_0$ and hence
acts by an element of $\mathrm{SO}(2)$ on its orthogonal
spatial plane.  This proves the factorization.
These matrices have determinant one and preserve $V_+$.

\begin{proposition}\label{IV-76-44}
There is a unique strongly continuous unitary representation
$U(\Lambda)$ of $\mathcal L$ fixing $\Omega$ and satisfying
\begin{equation}\label{IV-76-47}
 U(\Lambda)\mathcal D_0=\mathcal D_0,\qquad
 U(\Lambda)\phi(f)U(\Lambda)^{-1}\psi
             =\phi(f\circ\Lambda^{-1})\psi
                   \quad(\psi\in\mathcal D_0).
\end{equation}
Together with translations it defines the strongly continuous
representation
\begin{equation}\label{IV-76-48}
 U(a,\Lambda)=U(a)U(\Lambda),\qquad
 U(a,\Lambda)U(a',\Lambda')
                  =U(a+\Lambda a',\Lambda\Lambda').
\end{equation}
Its spatial-rotation subgroup agrees with the representation
on the reflection Hilbert space.
\end{proposition}

\begin{proof}
We first prove invariance of the analytic functions:
\begin{equation}\label{IV-76-45}
 \mathbf W_n(\iota\Lambda\zeta_1,\ldots,\iota\Lambda\zeta_{n-1})
          =\mathbf W_n(\iota\zeta_1,\ldots,\iota\zeta_{n-1})
 \quad(\zeta_k\in\mathbb R^{1+2}+iV_+).
\end{equation}
The complex linear map $\widehat\Lambda=\iota\Lambda\iota^{-1}$
preserves $\mathrm{Tb}_\circ$.  Both sides, viewed as
functions of $\iota\zeta_k$, are holomorphic on its
connected product.  It suffices to compare them at real
differences $\xi_k$ of positive third coordinate.
Choose $x_1=0$, $x_{k+1}=x_k+\xi_k$.
Then $\mathbf W_n(\xi)=\mathsf s_n(x)$.
Spatial-rotation invariance follows from that of the
Euclidean kernel.

For a boost, use the complex rotation
\begin{equation}\label{IV-76-49}
 r_z(y_1,y_2,y_3)
       =(y_1\cos z+y_3\sin z,\ y_2,\ y_3\cos z-y_1\sin z).
\end{equation}
It satisfies $r_{is}=\widehat{B(s)}$.
Fix real positive-time differences $\xi_k$.
The set of $z$ for which every $r_z\xi_k$ lies in
$\mathrm{Tb}_\circ$ is open and contains the imaginary
axis, since real boosts preserve the forward tube.
A sufficiently small neighborhood of the compact segment
from $0$ to $is$ lies in this set and is connected.
For small real $z$, $r_z$ is an ordinary Euclidean
rotation preserving the ordering of these points.
Thus $\mathbf W_n(r_z\xi_1,\ldots,r_z\xi_{n-1})$
is constant there by Euclidean invariance.
The identity theorem continues this equality to $z=is$.
The identity theorem in all difference variables, then
the rotation--boost factorization, proves
\eqref{IV-76-45}.

In \eqref{IV-76-30} change variables $x_j=\Lambda y_j$
and use \eqref{IV-76-45}.  The approach vectors become
$\Lambda^{-1}\alpha_k\in V_+$, which are allowed in
that boundary formula.  Hence
\begin{equation}\label{IV-76-46}
 \mathcal W_n^\sigma(F\circ(\Lambda^{-1})^{\times n})
                         =\mathcal W_n^\sigma(F).
\end{equation}
This first holds on tensors; continuity under invertible
linear changes of Schwartz variables and tensor density
extend it to all tests.

On a polynomial vector define
\[
 U(\Lambda)\phi(f_1)\cdots\phi(f_n)\Omega
       =\phi(f_1\circ\Lambda^{-1})\cdots
                       \phi(f_n\circ\Lambda^{-1})\Omega .
\]
Inner products of two such vectors are Wightman values:
move the first string across using Hermiticity, reversing
its order and conjugating its tests.
Equation \eqref{IV-76-46} preserves every such inner
product.  The rule therefore respects all linear
relations, is isometric, and has inverse given by
$\Lambda^{-1}$.  Density of $\mathcal D_0$ gives a unitary
extension; the group law and covariance hold on its
spanning vectors and hence as stated.
This also proves uniqueness.

As $\Lambda\to1$, the time derivatives of
$f\circ\Lambda^{-1}-f$ tend to zero in $L^1$.
The chain rule expresses them as finite sums of derivatives
of $f$ after the linear change; dominated convergence,
using Schwartz decay uniformly for $\Lambda$ near $1$,
proves the assertion.
Lemma~\ref{IV-76-38} gives strong continuity first on
polynomial vectors, then on all of $\mathcal H$ by
unitarity and density.
On polynomial vectors the conjugation
$U(\Lambda)U(a)U(\Lambda)^{-1}$ translates every test
by $\Lambda a$.  This proves \eqref{IV-76-48}
and continuity of the combined representation.

Finally, the original spatial rotations commute with
$H$ and $B(s,s')$, hence with $A_l$, and preserve $q_0$.
They conjugate $U(x)$ to $U((1\oplus R)x)$.
Changing variables in \eqref{IV-76-26} therefore gives
the same field transformation as $U(1\oplus R)$.
Both fix the vacuum, so they agree on $\mathcal D_0$
and on its closure.
\end{proof}

\subsection{Interchanging spacelike-separated insertions}

\begin{proposition}\label{IV-76-52}
If $f,h\in\mathcal S(\mathbb R^{1+2})$ have spacelike
separated supports, then
\begin{equation}\label{IV-76-53}
 \phi(f)\phi(h)\psi=\phi(h)\phi(f)\psi
                         \qquad(\psi\in\mathcal D_0).
\end{equation}
\end{proposition}

\begin{proof}
First suppose all tests have compact support.
Testing against polynomial vectors and using Hermiticity
reduces the assertion to invariance of
$\mathcal W_n^\sigma(f_1\otimes\cdots\otimes f_n)$
under interchange of adjacent entries $k,k+1$ whose
supports are spacelike separated.
Set $K=\prod_j\operatorname{supp}f_j$; zero entries are
trivial.  On $K$, write
$\xi=\vec x_{k+1}-\vec x_k$, $\tau=t_{k+1}-t_k$, and put
\[
 \begin{gathered}
 \delta=\min_K(|\xi|-|\tau|)>0,\quad A=4/\delta,\quad
 \eta_0=\min((2A)^{-1},1/2),\\
 h_*=2A\max_{x\in K,\ j<n}
           (|\vec x_{j+1}-\vec x_j|+|t_{j+1}-t_j|)+6.
 \end{gathered}
\]
We deform the two adjacent heights past each other:
\begin{equation}\label{IV-76-54}
 \begin{gathered}
 h_k^\lambda=-\lambda/2,\quad h_{k+1}^\lambda=\lambda/2,\\
 h_j^\lambda=\tfrac12+(j-k-1)h_*\quad(j>k+1),\qquad
 h_j^\lambda=-\tfrac12-(k-j)h_*\quad(j<k),\\
 z_j(x;\lambda,\eta)=\iota x_j+\eta h_j^\lambda e_3,
 \qquad -1\le\lambda\le1,\quad0<\eta\le\eta_0 .
 \end{gathered}
\end{equation}
Although the vertical ordering fails during this
deformation, the points remain in the union of tilted
tubes.  Indeed choose $\sin\theta=A\eta$ and
\[
 u_\theta=\cos\theta\,e_3+\sin\theta\,(\xi/|\xi|,0).
\]
Then $\cos\theta\ge\sqrt3/2$, and direct calculation gives
\begin{equation}\label{IV-76-55}
 \begin{split}
 d_\circ^{u_\theta}(z_{j+1}-z_j)
 &=\eta(h_{j+1}^\lambda-h_j^\lambda)\cos\theta\\
 &\quad+A\eta\,\frac{\xi}{|\xi|}
                   \cdot(\vec x_{j+1}-\vec x_j)
             -A\eta|t_{j+1}-t_j|.
 \end{split}
\end{equation}
For $j=k$ this is at least $-\eta+A\eta\delta=3\eta$.
For $j\ne k$ the height gap is at least $h_*$;
the definition of $h_*$ makes the same lower bound hold.
Consequently
\begin{equation}\label{IV-76-56}
 z(x;\lambda,\eta)\in\mathrm{Tb}^{\rm all}_n,\qquad
 |\mathbf S_n(z(x;\lambda,\eta))|
                      \le(C_\wp\eta^{-\wp})^n
\end{equation}
uniformly in $x,\lambda,\eta$ in these ranges.
The witnessing direction may depend on $x$ and $\eta$;
the holomorphic function $\mathbf S_n$ is the same on
all of these tubes.

For a smooth $g$ supported in $K$, let
$I_g(\lambda,\eta)=\int g(x)\mathbf S_n(z(x;\lambda,\eta))\,dx$.
The Cauchy--Riemann equations and integration by parts give
\begin{equation}\label{IV-76-57}
 \partial_\eta I_g=I_{-iD_\lambda g},\qquad
 \partial_\lambda I_g=\eta I_{-iD'g},\qquad
 D_\lambda=\sum_jh_j^\lambda\partial_{t_j},\quad
 D'=\tfrac12(\partial_{t_{k+1}}-\partial_{t_k}).
\end{equation}
Compactness at every positive height justifies differentiation.
Lemma~\ref{IV-76-23}, with $N=\lceil\wp n\rceil+1$,
gives the uniform bound
\begin{equation}\label{IV-76-58}
 |I_g(\lambda,\eta)|\le Q(g):=
 2^N\eta_0^{-(N-1/2)}C_\wp^n
          \sum_{r=0}^N\sup_{|\lambda|\le1}\|D_\lambda^rg\|_1 .
\end{equation}
The coefficients of $D_\lambda$ are uniformly bounded,
so $Q(g)<\infty$.

Take $g=\bigotimes_jf_j$.
At $\lambda=1$ the heights increase in the original
order, and \eqref{IV-76-30} identifies the limit
$I_g(1,\eta)$ with the original Wightman value.
At $\lambda=-1$ they increase in the order with $k,k+1$
interchanged.  The symmetry of the single function
$\mathbf S_n$ and the same boundary formula identify
$I_g(-1,\eta)$ with the interchanged value.
But \eqref{IV-76-57}--\eqref{IV-76-58} imply
\[
 |I_g(1,\eta)-I_g(-1,\eta)|\le2\eta Q(D'g)\longrightarrow0.
\]
This proves the adjacent interchange for compact tests.

For polynomial vectors with other Schwartz entries,
cut off those entries and use
Lemma~\ref{IV-76-38} to pass to the limit in both
products.  Density of polynomial test vectors gives
\eqref{IV-76-53} for compactly supported $f,h$.
Finally cut off $f,h$ themselves inside their original
supports.  Their spacelike separation is preserved,
and graph-norm continuity passes the equality to the
original Schwartz functions.
\end{proof}

\begin{theorem}[Scalar reconstruction]\label{IV-76-59}
Under (Rec), the reflection Hilbert space carries a
Hermitian scalar field $\phi$ with the following properties:
\begin{enumerate}
\item The proper orthochronous Poincar\'e group acts by
a strongly continuous unitary representation $U$ fixing
the unit vector $\Omega$.  Its translation spectrum lies
in $\overline{V_+}$.
\item The dense polynomial domain $\mathcal D_0$ contains
$\Omega$ and is invariant under $U$ and every $\phi(f)$.
For $\psi,\psi'\in\mathcal D_0$,
$f\mapsto\langle\psi,\phi(f)\psi'\rangle$ is tempered,
and the Hermiticity relation in \eqref{IV-76-27} holds.
\item The field is covariant:
\[
 U(a,\Lambda)\phi(f)U(a,\Lambda)^{-1}
      =\phi\bigl(f(\Lambda^{-1}(\,\cdot-a))\bigr)
             \quad\hbox{on }\mathcal D_0,
\]
and fields with spacelike separated test supports commute
there.
\item The vacuum is cyclic.  It is the unique
translation-invariant vector up to scalar multiples
if $\ker H=\mathbb C\Omega$.
\end{enumerate}
The vacuum distributions are $\mathcal W_n^\sigma$.
Their analytic functions are $\mathbf W_n$, with
boundary values \eqref{IV-76-30}, and their Euclidean
continuations recover $\sigma$ on every separated
$C_c^1$ tuple by \eqref{IV-75-79}.
On $\mathcal S_{\ne}(\mathbb R^{3n})$ the Schwinger
functions are the absolutely convergent integrals
\eqref{IV-75-92}.
Thus, with the additional vacuum-uniqueness hypothesis,
these objects form a scalar Wightman theory.
\end{theorem}

\begin{proof}
The domain, spectrum, cyclicity and Hermiticity were proved
in Section~\ref{sec:reconstruction-fields}.
Equation \eqref{IV-76-28}, followed by
\eqref{IV-76-15}, gives tempered matrix elements.
Propositions~\ref{IV-76-44} and \ref{IV-76-52}
give covariance and locality on $\mathcal D_0$.
The boundary and Euclidean identities follow from
\eqref{IV-76-30}, \eqref{IV-75-79} and \eqref{IV-75-80}.
On differences with third coordinate at least
$\epsilon>0$ and real spatial coordinates,
\eqref{IV-75-72} bounds $\mathbf W_n$ by
$(C_\wp\min(1,\epsilon/2)^{-\wp})^n$;
in particular the Euclidean continuation has the required
polynomial growth.  Finally
\eqref{IV-75-89} proves the asserted absolute convergence.
\end{proof}

\section{Reconstruction of the Gibbs limits}
\label{sec:gibbs-reconstruction}

We now verify the reconstruction hypotheses for the lattice limits.
The estimates must hold for every limit obtained by conditioning a
frame: the canonical components need not be centered, even when
the original limit is centered.  This is why the two-state estimates
and closure of the Gibbs class under invariant conditioning were
retained in the preceding arguments.

Fix the sourced construction at radius $2\varrho$, and fix
$0<g\le\gamma_\varrho$.  In the notation of
\eqref{IV-34-6} and \eqref{IV-56-18}, put
\begin{equation}\label{IV-78-7}
 C^g_{\rm pt}=320\cdot8^4m_3(40L)^4
                  \frac{2\bar Q_{\rm mid}}{\varrho}M,
 \qquad C^g_\wp=16C^g_{\rm pt}.
\end{equation}
All these constants are fixed before choosing a sequence of states.
Let $\bar c_g$ denote the number $\bar c$ in
\eqref{IV-36-14}, evaluated at the bound system
$\widehat{\mathsf C}^g$ of Proposition~\ref{prop:gibbs-frame-classes}.

\begin{proposition}\label{IV-77-12}\label{IV-78-12}
Let $K_j\to\infty$, $(K_j,\beta_j)\in\mathfrak d(g)$ and
$\nu_j\in\mathcal G^{\rm per}_{K_j,\beta_j}$.
Every limit $\sigma'$ in the frame $(K_j,\nu_j)_j$ satisfies
\textup{(Rec)} with
\[
 (C_\wp,\wp,c_0)=(C^g_\wp+|c_{\sigma'}|,4,1/2),
 \qquad |c_{\sigma'}|\le\bar c_g.
\]
In particular the same choice $(C^g_\wp+\bar c_g,4,1/2)$
works for every such limit.  If the frame itself converges to
$\sigma$, one can use $(C^g_\wp,4,1/2)$ for $\sigma$.
\end{proposition}

\begin{proof}
The Gibbs classes form a class system with the stated uniform
bounds by Proposition~\ref{prop:gibbs-frame-classes}.  A limit in
the frame has approximants $\rho_j$ in these same classes; the
centering law remains $\nu_j$.  Theorem~\ref{IV-56-29} therefore
applies to $\sigma'$ and its centered functional.  It gives full
Euclidean invariance, reflection positivity, the translation
semigroup and its joint spectral measure, and the cone
$E\ge|q|$.  These are (Rec1)--(Rec3).  Separability is part
of the Hilbert-space construction in
Section~\ref{sec:reflection-hilbert}.  The one-point function and
its bound are \eqref{IV-36-14}.

For (Rec4), take $0<\eta\le1$,
$F,F'\in\mathbb C\mathfrak S(\{x_3>\eta\})$ and
$h\in C_c^1(\{|x_3|<\eta/2\})$.
The centered functional has (R1)--(R5) with the constants
\eqref{IV-56-18}.  Corollary~\ref{IV-34-42}, with $d=\eta$,
gives
\[
 \sigma'^{\rm c}(\theta\bar F\cdot(h)\cdot F')
       =\int h(x)v_\eta(x)\,dx,
\]
where
\[
 |v_\eta(x)|\le C^g_{\rm pt}
       \min(1,\eta-|x_3|)^{-4}
       \|[\tau_{-\eta e_3}F]\|_{\sigma'^{\rm c}}
       \|[\tau_{-\eta e_3}F']\|_{\sigma'^{\rm c}}.
\]
On the support of $h$ the first minimum exceeds $\eta/2$.
Integration thus proves
\begin{equation}\label{IV-78-14}
 \begin{split}
 |\sigma'^{\rm c}(\theta\bar F\cdot(h)\cdot F')|
 &\le C^g_\wp\eta^{-4}\|h\|_1
       \|[\tau_{-\eta e_3}F]\|_{\sigma'^{\rm c}}\,
       \|[\tau_{-\eta e_3}F']\|_{\sigma'^{\rm c}}.
 \end{split}
\end{equation}
These are the norms in the coordinate reflection Hilbert space,
as required in \eqref{IV-75-51}.  Lemma~\ref{IV-75-50}
transfers the bound to $\sigma'$ with the additional constant
$|c_{\sigma'}|$.  For the limit of the frame itself the integrating
and centering laws agree, so its one-point function is zero.
\end{proof}

The passage from clustering of separated tuples to clustering of
Schwartz tests uses the reflection Hilbert space.  Directly
translating a Schwartz approximation would give bounds depending
on the translation; the Hilbert-space argument keeps them uniform.

\begin{proposition}\label{IV-28-48}
Suppose a symmetric, normalized separated-tuple functional
$\sigma$ has the bounds of
Section~\ref{sec:reflection-hilbert}, is Euclidean invariant and
reflection positive, and has Schwinger functions satisfying
\[
 |\mathsf S_n(\varphi)|\le c_E C_E^n(n!)^\kappa
                  |\varphi|_{kn},\qquad
 \varphi\in\mathcal S_{\ne}(\mathbb R^{3n}),
\]
where $k\in\mathbb N$ and $c_E,C_E,\kappa\ge0$.
Assume also that
\[
 \sigma(X\cdot\tau_vY)\longrightarrow\sigma(X)\sigma(Y)
 \quad(|v|\to\infty)
\]
for all separated tuples $X,Y$, with the concatenation defined
for sufficiently large $|v|$.
Then the Schwinger functions satisfy Euclidean invariance,
permutation symmetry, reflection positivity
\eqref{IV-26-42}, the displayed growth bound, and
\begin{equation}\label{IV-28-40}
 \lim_{s\to\infty}
 \mathsf S_{n+n'}(\Theta\bar\varphi\otimes\tau_{sv}\varphi')
 =\mathsf S_n(\Theta\bar\varphi)\mathsf S_{n'}(\varphi')
\end{equation}
for $\varphi\in\mathcal S_+(\mathbb R^{3n})$,
$\varphi'\in\mathcal S_+(\mathbb R^{3n'})$ and
$0\ne v\in\mathbb R^2\times\{0\}$.  We refer to these five
properties as \textup{(E0$'$)--(E4)}, respectively growth,
Euclidean invariance, reflection positivity, symmetry and
clustering.  The test spaces are the flat Schwartz spaces defined
in Section~\ref{sec:distribution-extension}.
\end{proposition}

\begin{proof}
Proposition~\ref{IV-26-41} gives the first four properties.
For the last, form the algebraic direct sum
$\underline{\mathcal S}_+=\bigoplus_{n\ge0}
\mathcal S_+(\mathbb R^{3n})$, with the positive form
\[
 \langle\underline\varphi,\underline\psi\rangle_{\mathcal S}
 =\sum_{n,n'}\mathsf S_{n+n'}
                       (\Theta\bar\varphi_n\otimes\psi_{n'}).
\]
Let $\mathcal H^{\mathcal S}$ be its completion modulo the null
space.  If both sequences vanish above degree $n_0$, their
form is bounded in absolute value by
\begin{equation}\label{IV-28-49}
 \max(1,c_E)\max(1,C_E)^{2n_0}((2n_0)!)^\kappa
 \left(\sum_{n\le n_0}|\varphi_n|_{2kn_0}\right)
 \left(\sum_{n\le n_0}|\psi_n|_{2kn_0}\right).
\end{equation}
This follows term by term from the tensor-product seminorm
bound; the degree-zero term is ordinary multiplication.

Send a smooth separated positive-time tuple $X$ to the sequence
$\iota X$ with $\bigotimes X$ in its degree and zero elsewhere.
For finite sums of such tuples,
\begin{equation}\label{IV-28-50}
 \langle\iota F,\iota F'\rangle_{\mathcal S}
        =\sigma(\theta\bar F\cdot F').
\end{equation}
Lemma~\ref{IV-26-17} and \eqref{IV-28-49} show that their
classes are dense in $\mathcal H^{\mathcal S}$.
They are also dense in $\mathcal H(\sigma)$: mollify each
entry of a $C_c^1$ tuple with a common radius tending to zero,
smaller than its distance to the reflection plane and half
its least support separation.  The resulting tuples remain
separated in the positive half-space.  The moving-entry
continuity estimate gives convergence of the three terms in
\[
 \|[X]-[X_j]\|^2
 =\sigma(\theta X\cdot X)
  -2\operatorname{Re}\sigma(\theta X\cdot X_j)
  +\sigma(\theta X_j\cdot X_j).
\]
Consequently \eqref{IV-28-50} extends to a unitary
$V:\mathcal H(\sigma)\to\mathcal H^{\mathcal S}$ carrying
$\Omega$ to the constant sequence.  For a fixed tangential
translation, the seminorm estimate \eqref{IV-26-43} shows
that the same approximations converge after translation.
Thus
\[
 V\mathcal U(v')V^{-1}[\underline\varphi]
       =[\tau_{(v',0)}\underline\varphi].
\]

Clustering of tuples gives, first for tuple classes,
\[
 \langle\chi,\mathcal U(sv')\chi'\rangle
       \longrightarrow
       \langle\chi,\Omega\rangle\langle\Omega,\chi'\rangle,
       \qquad v'\ne0.
\]
The difference between the two sides has modulus at most
$2\|\chi\|\|\chi'\|$, uniformly in $s$, since
$\mathcal U(sv')$ is unitary.  Density extends this limit to
all vectors.  Apply it to the inverse images under $V$ of
the classes with $\varphi$ and $\varphi'$ in degrees $n,n'$.
The form identity gives exactly \eqref{IV-28-40}.
\end{proof}

\begin{theorem}[Gibbs limits and scalar fields]\label{IV-77-14}
\label{IV-78-15}
Fix the construction and $g$ as above.  Every sequence
$K_j\to\infty$, $(K_j,\beta_j)\in\mathfrak d(g)$ and
$\nu_j\in\mathcal G^{\rm per}_{K_j,\beta_j}$ has a subsequence
whose self-centered moments converge on all separated tuples.
Let $\sigma$ be any such limit.  It has the canonical decomposition
of Theorem~\ref{IV-35-36}, unique in the sense of
Proposition~\ref{IV-36-9}.  Every component $\sigma_\xi$ is
Euclidean invariant, reflection positive at every mirror, and
has the uniform clustering estimate \eqref{IV-36-18}.
Its one-point constant may be nonzero, with
$|c_\xi|\le\bar c_g$ and $\int c_\xi\,d\lambda=0$.

Each component satisfies \textup{(Rec)} with
$(C^g_\wp+\bar c_g,4,1/2)$ and
$\ker H_{\sigma_\xi}=\mathbb C\Omega_{\sigma_\xi}$.
Its Schwinger functions satisfy \textup{(E0$'$)--(E4)} with
the common growth estimate
\begin{equation}\label{eq:gibbs-schwinger-growth}
 |\mathsf S_n^{\sigma_\xi}(\varphi)|
 \le2(C^g_E)^n(n!)^8|\varphi|_{8n},\qquad
 C^g_E=16e^8 3^4\pi^2(C^g_\wp+\bar c_g).
\end{equation}
On its reflection Hilbert space there is a scalar Wightman
theory with unique vacuum, common invariant polynomial domain,
and the Schwinger functions of $\sigma_\xi$, as specified in
Theorem~\ref{IV-76-59}.

The original limit $\sigma$ has the same reconstruction and
growth conclusions, using $C^g_\wp$ in place of
$C^g_\wp+\bar c_g$, except that neither clustering nor vacuum
uniqueness is asserted.  Its zero-energy space is canonically
$L^2(\mathsf X,\lambda)$, with $\Omega$ corresponding to $1$.
All stated constants are independent of the sequence, its
subsequence and its components.
\end{theorem}

\begin{proof}
The cumulant compactness result, Proposition~\ref{IV-21-25}, and the
moment--cumulant formula give a common convergent subsequence.
Proposition~\ref{prop:gibbs-frame-classes} makes it a converging
frame with the fixed bound system $\widehat{\mathsf C}^g$.
Theorem~\ref{IV-35-36} and Proposition~\ref{IV-36-9} give its canonical
decomposition and uniqueness.  By Theorem~\ref{IV-35-36},
each component is a limit in this frame, obtained along the
full chosen sequence by conditioning $\nu_j$ on invariant
events of positive probability.  Proposition~\ref{IV-54-18}
keeps the conditioned laws in the periodic Gibbs classes.
The centering law is still $\nu_j$, so Proposition~\ref{IV-77-12}
applies with the possibly nonzero constant $c_\xi$.
Full Euclidean invariance and reflection positivity follow
from Theorem~\ref{IV-56-29} for these same approximants.

The uniform component clustering estimate is
\eqref{IV-36-18}; the supports of $X$ and $\tau_vY$ are
disjoint for the large translations in that estimate.
Equation~\eqref{IV-36-22} therefore gives the simple
zero-energy space of each component.  Substituting $\wp=4$
in Proposition~\ref{IV-75-87} gives
\eqref{eq:gibbs-schwinger-growth}; its integral extension is
the Schwinger extension by uniqueness on dense separated
tensors.  Proposition~\ref{IV-28-48} now supplies all five
Euclidean properties.  Theorem~\ref{IV-76-59} supplies the
Wightman theory.

For $\sigma$, the one-point constant is zero and the same
arguments apply without the clustering step.  The last
assertion on its vacuum space is Theorem~\ref{IV-36-25}.
\end{proof}


\section{A rotation criterion along growing tori}
\label{sec:growing-torus-rotation}

The lattice rotation identity has a bulk term and a term supported where
a spatial cutoff changes. Their sum has a continuum limit. Separate
control of the two terms gives a sufficient condition for rotation
invariance; no separate limit will be assumed in deriving the identity.

\subsection{The finite-lattice identity}

Fix $K,m\ge1$, $a=a_K$, $\ell=2L^m$ and $\beta>0$.
We use the central lifts in $[-\ell/2,\ell/2)^3$ when multiplying a
position by $\omega\in\mathfrak{so}(3)$. A function has margin $r$ if
its support is contained in $(-\ell/2+r,\ell/2-r)^3$. Local
displacements across a plaquette always use its local lift. Let $x_p$
and $x_b$ be the plaquette barycentre and bond midpoint. In sums
$b\subset\partial p$, each of the four underlying bonds has its
positive coordinate orientation.

For $b=\langle x,x+ae_\mu\rangle$ and a signed direction $\rho\ne\pm\mu$,
let $W_x(\mu,\rho)$ be the holonomy along the four sides of the square
starting at $x$ in directions $\mu,\rho,-\mu,-\rho$. Its underlying
plaquette is denoted by $p(b,\rho)$. If $p$ contains $b$, $W_{b,p}$ is
the holonomy around $p$ starting with $b$; the other signed direction
is determined uniquely. With $\langle A,B\rangle=\frac12\operatorname{tr}(AB)$,
put
\begin{equation}\label{IV-40-7}
 \mathbf w(V)=\frac{V-V^{-1}}{2i},\qquad
 \mathbf u_{b,\rho}=-\frac12\{
 \mathbf w(W_x(\mu,\rho))-\mathbf w(W_x(\mu,-\rho))\},\qquad
 \mathbf u_{b,\mu}=0.
\end{equation}
The indices $\rho,\lambda$ in sums below range over $1,2,3$. Define
\begin{equation}\label{IV-40-19}
\begin{aligned}
 d_{b,p,\rho}&=\langle\mathbf u_{b,\rho},\mathbf w(W_{b,p})\rangle,
 &t_\rho(p)&=\sum_{b\subset\partial p}d_{b,p,\rho},\\
 s_{\lambda\rho}(p)&=\sum_{b\subset\partial p}
       a^{-1}(x_b-x_p)_\lambda d_{b,p,\rho},
 &r_\rho(p)&=t_\rho(p)-\tfrac12(a_{p+ae_\rho}-a_{p-ae_\rho}).
\end{aligned}
\end{equation}
For functions on the torus, set
\[
 \nabla^{\rm c}_\rho h(y)=a^{-1}
       \{h(y+\tfrac a2e_\rho)-h(y-\tfrac a2e_\rho)\},\qquad
 \nabla^{\rm s}_\omega h(y)=\sum_\rho(\omega y)_\rho
       \frac{h(y+ae_\rho)-h(y-ae_\rho)}{2a}.
\]
The latter formula applies to plaquette weights as well, with
$p\pm ae_\rho$ retaining the directions of $p$.

Let $\chi$ have margin $2a$. For a plaquette weight $v$ define
\begin{equation}\label{IV-40-30}
 \mathfrak S_{\omega,\chi}[v]
   =\sum_p v(p)\chi(x_p)\sum_{\rho,\lambda}\omega_{\rho\lambda}s_{\lambda\rho}(p),
 \qquad
 \mathfrak K_{\omega,\chi}[v]
   =\sum_p a^{-1}v(p)\chi(x_p)\sum_\rho(\omega x_p)_\rho r_\rho(p).
\end{equation}
Write $\mathfrak S_\omega[v]$ and $\mathfrak K_\omega[v]$ for the same
expressions without $\chi$ when $v$ has margin $2a$. The shell
insertion is
\begin{equation}\label{IV-40-31}
\begin{split}
 h_{\omega,\chi}(p)&=a^{-1}\sum_{b\subset\partial p,\rho}
       \{\chi(x_b)-\chi(x_p)\}(\omega x_b)_\rho d_{b,p,\rho},\\
 \mathfrak B_{\omega,\chi}
   &=\beta\sum_p a_p\nabla^{\rm s}_\omega\chi(x_p)
       -\beta\sum_p h_{\omega,\chi}(p)\\
   &\quad-\tfrac32\sum_{p=p_{\mu\nu}(x),\,\mu<\nu}a_p
       \{(\omega x_p)_\mu\nabla^{\rm c}_\mu\chi(x_p)
          +(\omega x_p)_\nu\nabla^{\rm c}_\nu\chi(x_p)\}.
\end{split}
\end{equation}
These are bounded polynomial functions of the bond variables.
We say that $\chi$ covers $f_1,\ldots,f_n$ if it equals one at $x_p$
and all four $x_b$ whenever some $f_i(x_p)\ne0$.
All cumulants in the next display are taken in the Wilson measure at
$(K,\ell,\beta)$, and $\mathcal O(f)=a^{-1}\sum_p f(x_p)a_p$.
Define the bulk term by
\begin{equation}\label{IV-40-35}
\begin{split}
 \mathfrak U_{\omega,\chi}(\underline f)
  &=-\beta\langle(\mathfrak S_{\omega,\chi}
       +\mathfrak K_{\omega,\chi})[1];
       \mathcal O(f_1);\ldots;\mathcal O(f_n)\rangle\\
  &\quad+a^{-1}\sum_{i=1}^n
       \langle(\mathfrak S_\omega+\mathfrak K_\omega)[f_i];
                  (\mathcal O(f_l))_{l\ne i}\rangle.
\end{split}
\end{equation}

\begin{proposition}[Finite-lattice rotation identity]\label{IV-40-39}
If $\chi$ covers the real sources $f_1,\ldots,f_n$, then
\begin{equation}\label{IV-40-45}
 \sum_{i=1}^n S^T_{n;K,\ell,\beta}
       (f_1,\ldots,\nabla^{\rm s}_\omega f_i,\ldots,f_n)
 =\mathfrak U_{\omega,\chi}(\underline f)
       +\langle\mathfrak B_{\omega,\chi};
                   \mathcal O(f_1);\ldots;\mathcal O(f_n)\rangle.
\end{equation}
The identity also holds for complex sources.
\end{proposition}
\begin{proof}
Let $\sigma_A$, $1\le A\le3$, be the Pauli basis and let
$\partial_b^A F$ denote differentiation after replacing $U(b)$ by
$e^{it\sigma_A}U(b)$ at $t=0$. For a vector field $\eta$ set
\[
 X_b=a^{-1}\sum_\rho\eta_\rho(x_b)\mathbf u_{b,\rho},\quad
 \mathfrak v_\eta F=\sum_{b,A}\langle X_b,\sigma_A\rangle\partial_b^AF,
 \quad
 \operatorname{div}\mathfrak v_\eta
       =\sum_{b,A}\partial_b^A\langle X_b,\sigma_A\rangle.
\]
Haar invariance in the $b$ factor gives $\int\partial_b^AF\,dU=0$.
Differentiation under the integral is valid because the configuration
space is a finite product of compact groups and the functions used
here have continuous derivatives. The product rule therefore gives
\begin{equation}\label{IV-40-25}
 \int\mathfrak v_\eta F\,dU
   =-\int F\operatorname{div}\mathfrak v_\eta\,dU.
\end{equation}

For $b\subset\partial p$, its based holonomy changes to
$e^{it\sigma_A}W_{b,p}$. Since $\sigma_A^2=1$ and
$a_p=1-\frac12\operatorname{tr}W_{b,p}$, direct differentiation yields
\[
 \partial_b^Aa_p=\langle\sigma_A,\mathbf w(W_{b,p})\rangle,
 \qquad
 \partial_b^A\langle\sigma_A,\mathbf w(W_{b,p})\rangle=1-a_p.
\]
There is no sum over $A$ in the second identity. If $b$ does not
belong to $p$, the first derivative is zero. Thus, for $\rho\ne\mu_b$,
\[
 \sum_A\partial_b^A\langle\sigma_A,\mathbf u_{b,\rho}\rangle
       =\tfrac32\{a_{p(b,\rho)}-a_{p(b,-\rho)}\}.
\]
Summing the first derivative over the Pauli basis, and regrouping
the last identity by plaquettes, gives
\begin{equation}\label{IV-40-23}
\begin{split}
 \mathfrak v_\eta a_p
    &=a^{-1}\sum_{b\subset\partial p,\rho}\eta_\rho(x_b)d_{b,p,\rho},\\
 \operatorname{div}\mathfrak v_\eta
    &=-\tfrac32\sum_{p=p_{\mu\nu}(x),\,\mu<\nu}a_p
           (\nabla^{\rm c}_\mu\eta_\mu+
                         \nabla^{\rm c}_\nu\eta_\nu)(x_p).
\end{split}
\end{equation}
Indeed, the two bonds of $p$ perpendicular to direction $\rho$
have midpoints $x_p\mp\frac a2e_\rho$, and contribute the difference
$\eta_\rho(x_p-\frac a2e_\rho)-\eta_\rho(x_p+\frac a2e_\rho)$.

For a plaquette weight $v$, write $A(v)=\sum_pv(p)a_p$.
Take $\eta(y)=\chi(y)\omega y$. Split
$\chi(x_b)=\chi(x_p)+[\chi(x_b)-\chi(x_p)]$, and use
$\omega x_b=\omega x_p+\omega(x_b-x_p)$ in the first part.
The margin of $\chi$ makes this a valid equality of lifts wherever
that part is nonzero. The resulting three terms are $h_{\omega,\chi}(p)$,
$a^{-1}\chi(x_p)\sum_\rho(\omega x_p)_\rho t_\rho(p)$, and
$\chi(x_p)\sum_{\rho,\lambda}\omega_{\rho\lambda}s_{\lambda\rho}(p)$.
Insert the definition of $r_\rho$ and sum the translated plaquette
actions by parts. Since $\omega_{\rho\rho}=0$,
$(\omega(x_p\pm ae_\rho))_\rho=(\omega x_p)_\rho$ wherever
the translated cutoff is nonzero. We obtain
\begin{equation}\label{IV-40-41}
 \mathfrak v_{\chi\omega y}A(v)
 = (\mathfrak S_{\omega,\chi}+\mathfrak K_{\omega,\chi})[v]
      +\sum_p v(p)h_{\omega,\chi}(p)
      -A(\nabla^{\rm s}_\omega(v\chi)).
\end{equation}
The same zero diagonal of $\omega$ reduces the divergence in
\eqref{IV-40-23} to the last sum of \eqref{IV-40-31}.

Put $w_z=\sum_i z_if_i$. The covering condition gives
$w_z\chi=w_z$ and $w_z(p)h_{\omega,\chi}(p)=0$. Consequently
\begin{equation}\label{IV-40-43}
\begin{split}
 \mathfrak v_{\chi\omega y}A(\beta-a^{-1}w_z)
       -\operatorname{div}\mathfrak v_{\chi\omega y}
   &=(\mathfrak S_{\omega,\chi}+\mathfrak K_{\omega,\chi})
                         [\beta-a^{-1}w_z]\\
   &\quad+\mathcal O(\nabla^{\rm s}_\omega w_z)
                         -\mathfrak B_{\omega,\chi}.
\end{split}
\end{equation}
Its integral against $e^{-A(\beta-a^{-1}w_z)}dU$ is zero by
\eqref{IV-40-25}. The partition function is entire in $z$ and
positive at zero, so we may normalize in a neighborhood of zero.
For any affine insertion $Y_z=Y_0+\sum_i z_iY_i$, differentiation
of its normalized mean gives
\[
 \partial_1\cdots\partial_n\langle Y_z\rangle_{w_z}\big|_0
 =\langle Y_0;\mathcal O(f_1);\ldots;\mathcal O(f_n)\rangle
       +\sum_i\langle Y_i;(\mathcal O(f_l))_{l\ne i}\rangle.
\]
To see this, introduce an auxiliary source $u$ for a fixed insertion,
differentiate $\partial_u\log\int e^{uY+\sum_i z_i\mathcal O(f_i)}d\mu$,
and then apply the product rule to $z_i\langle Y_i\rangle_{w_z}$.
Apply this identity to the right side of \eqref{IV-40-43}, whose
normalized mean vanishes. The constant insertion is
$\beta(\mathfrak S_{\omega,\chi}+\mathfrak K_{\omega,\chi})[1]
-\mathfrak B_{\omega,\chi}$ and its $i$th linear coefficient is
$\mathcal O(\nabla^{\rm s}_\omega f_i)
-a^{-1}(\mathfrak S_\omega+\mathfrak K_\omega)[f_i]$.
Rearrangement gives \eqref{IV-40-45}.
\end{proof}

\subsection{The continuum rotation defect}

\begin{lemma}[Discretization error]\label{IV-42-24}
Let $0<a=a_K\le1$, $\omega\in\mathfrak{so}(3)$ and
$\operatorname{supp}f\subset\{|y|<\mathcal R/2\}$, with
$f\in C_c^3(\mathbb R^3)$ real. Write
$|\omega|=\max_{\rho,\lambda}|\omega_{\rho\lambda}|$ and
$\|f\|_{C^3}=\max_{|\alpha|\le3}\|\partial^\alpha f\|_\infty$.
Then
\[
 D_f=\nabla^{\rm s}_\omega f-(\omega y)\cdot\nabla f
\]
has support in the closed $a$-neighborhood of $\operatorname{supp}f$, and
\begin{equation}\label{IV-42-25}
\begin{split}
 \|D_f\|_\infty&\le|\omega|(1+\mathcal R)a^2\|f\|_{C^3},\\
 \operatorname{Lip}D_f&\le4\sqrt3|\omega|(1+\mathcal R)a\|f\|_{C^3},\\
 \|D_f\|_\sharp&\le160\sqrt3 M|\omega|(1+\mathcal R)a\|f\|_{C^3}.
\end{split}
\end{equation}
If $\mathcal R/2+2a<\ell/2$, the same statements hold after
periodization to $T_\ell$, with unchanged norms; the periodized
$D_f$ is the difference of the corresponding torus expressions.
\end{lemma}
\begin{proof}
Put $r_\rho(y)=(2a)^{-1}[f(y+ae_\rho)-f(y-ae_\rho)]-
\partial_\rho f(y)$. Taylor's formula with integral remainder gives
\[
 r_\rho(y)=\frac{a^2}{4}\int_0^1(1-s)^2
       [\partial_\rho^3f(y+sae_\rho)+\partial_\rho^3f(y-sae_\rho)]\,ds.
\]
Thus $|r_\rho|\le a^2\|f\|_{C^3}/6$. Applying the first-order
Taylor formula to $\partial_\sigma f$ gives, without using a fourth
derivative,
\[
 \partial_\sigma r_\rho(y)=\frac a2\int_0^1(1-s)
       [\partial_\rho^2\partial_\sigma f(y+sae_\rho)
          -\partial_\rho^2\partial_\sigma f(y-sae_\rho)]\,ds,
 \qquad |\partial_\sigma r_\rho|\le\tfrac a2\|f\|_{C^3}.
\]
Outside the stated neighborhood every term in $D_f=\sum_\rho
(\omega y)_\rho r_\rho$ vanishes. Inside it, $|y|\le\mathcal R/2+a$,
so $|(\omega y)_\rho|\le2|\omega|(\mathcal R/2+a)$.
The product rule therefore gives
\[
 |D_f|\le|\omega|(\mathcal R/2+a)a^2\|f\|_{C^3},\qquad
 |\partial_\sigma D_f|
 \le|\omega|a\|f\|_{C^3}(a/3+3\mathcal R/2+3a).
\]
Since $a\le1$, these imply the first two bounds. The third follows from
\[
 \|h\|_\sharp=\max(\|h\|_\infty,40M\operatorname{Lip}h).
\]
For the torus statement, the support of $D_f$ has margin $a$ and
that of $f$ has margin $2a$. The periodic lift is locally a single
translate, or zero, because these translates of the closed supports
are disjoint and locally finite. Its supremum and gradient supremum
are unchanged. The Lipschitz constant is the gradient supremum both
on $\mathbb R^3$ and on the flat torus. Finally, shifts by $\pm ae_\rho$
cannot introduce a wraparound term from $f$, which has margin $2a$.
This proves the identity for the periodized difference.
\end{proof}

Consider a sequence $K_j\to\infty$, $\ell_j\to\infty$, $\beta_j>0$,
with $K_j,m_j\ge1$ and $\ell_j=2L^{m_j}$. For $n\ge2$ and separated tuples
$X=(f_1,\ldots,f_n)$ of real $C_c^1$ functions in the central cube
$(-\ell_j/4,\ell_j/4)^3$, put
\[
 \Phi_n^{(j)}(X)=S^T_{n;K_j,\ell_j,\beta_j}(X^{\langle\ell_j\rangle}).
\]
Assume that these forms converge pointwise to a family $(S_n)_{n\ge2}$
on every separated tuple and obey a bound table: for every $n\ge2$ and $R>0$
there is $j_0(n,R)$ such that, for all $j\ge j_0(n,R)$ and all
$d>0$,
\begin{equation}\label{IV-43-3}
 |\Phi_n^{(j)}(X)|\le C_n(d;R)\prod_i\|f_i\|_\sharp
\end{equation}
whenever the supports have mutual distance at least $d$ and each
lies in a ball of radius $R$. In particular, $S$ obeys the same
table. By \cref{IV-24-11}, $S$ is translation and $W_3$ invariant.

Fix one smooth function $\chi_1:[0,\infty)\to[0,1]$ equal to one
on $[0,1]$ and zero on $[2,\infty)$, and put
$\chi_{\mathcal R}(y)=\chi_1(|y|/\mathcal R)$.
For a separated smooth tuple $\underline f$ supported in
$\{|y|<\mathcal R/2\}$, define
\begin{equation}\label{IV-43-6}
\begin{split}
 \mathfrak U^{(j)}_{\omega,\mathcal R}(\underline f)
     &=\mathfrak U_{\omega,\chi_{\mathcal R}}
                         (\underline f^{\langle\ell_j\rangle}),\\
 \mathsf{Fl}^{(j)}_{\omega,\mathcal R}(\underline f)
     &=\langle\mathfrak B_{\omega,\chi_{\mathcal R}};
           \mathcal O(f_1^{\langle\ell_j\rangle});\ldots;
           \mathcal O(f_n^{\langle\ell_j\rangle})\rangle_{K_j,\ell_j,\beta_j}.
\end{split}
\end{equation}
They are defined for all sufficiently large $j$. Indeed, writing
$d_1=\min(1,d(\underline f))$, we may require
\begin{equation}\label{IV-43-5}
 2\mathcal R+2a_{K_j}<\ell_j/2,\qquad
 2\sqrt3a_{K_j}\le\mathcal R/2,\qquad
 a_{K_j}\le d_1/4.
\end{equation}
The first condition gives the cutoff margin and the second ensures
that $\chi_{\mathcal R}$ covers the sources. The third keeps the
supports of the difference insertions separated.

\begin{proposition}[Limit of the bulk and shell sum]\label{IV-43-18}
Under these assumptions, $t\mapsto S_n(e^{t\omega}\underline f)$
is differentiable, where $(e^{t\omega}f)(y)=f(e^{-t\omega}y)$, and
\begin{equation}\label{IV-43-20}
 \dot S_n(\omega;\underline f)
 :=\frac d{dt}S_n(e^{t\omega}\underline f)\big|_{t=0}
 =-\sum_iS_n(f_1,\ldots,(\omega y)\cdot\nabla f_i,\ldots,f_n).
\end{equation}
For every such tuple and radius,
\begin{equation}\label{IV-43-21}
 \lim_{j\to\infty}
 [\mathfrak U^{(j)}_{\omega,\mathcal R}(\underline f)
       +\mathsf{Fl}^{(j)}_{\omega,\mathcal R}(\underline f)]
       =-\dot S_n(\omega;\underline f).
\end{equation}
In particular, this limit is independent of $\mathcal R$.
\end{proposition}
\begin{proof}
For each $i$, take the closed $d_1/4$-neighborhood $A_i$ of
$\operatorname{supp}f_i$. The $A_i$ have mutual distance at least
$d_*=d_1/2$ and each lies in a ball of radius $R_*=R(\underline f)+1/4$.
For large $j$ they lie in the central cube. Every $f_i$,
$(\omega y)\cdot\nabla f_i$, $D_{f_i}$, and
$\nabla^{\rm s}_\omega f_i$ is supported in $A_i$ by
\cref{IV-42-24}. Applying \eqref{IV-43-3} gives
\[
 |\Phi_n^{(j)}(f_1,\ldots,D_{f_i},\ldots,f_n)|
 \le C_n(d_*;R_*)160\sqrt3M|\omega|(1+\mathcal R)
       a_{K_j}\|f_i\|_{C^3}\prod_{l\ne i}\|f_l\|_\sharp.
\]
This tends to zero. Pointwise convergence on the fixed tuple with
entry $(\omega y)\cdot\nabla f_i$ now proves
\[
 \sum_i\Phi_n^{(j)}(f_1,\ldots,\nabla^{\rm s}_\omega f_i,\ldots,f_n)
 \longrightarrow
 \sum_i S_n(f_1,\ldots,(\omega y)\cdot\nabla f_i,\ldots,f_n).
\]
The finite identity \eqref{IV-40-45} proves the first equality
claimed in \eqref{IV-43-21}.

To identify it with an orbit derivative, let $g_i(t,y)=f_i(e^{-t\omega}y)$.
On any compact interval of $t$, the derivatives
$\partial_t^k\partial_y^\alpha g_i$ for $k\le2$, $|\alpha|\le1$
are bounded and supported in one fixed ball. Taylor's formula thus
shows that
\[
 \frac{g_i(t+h)-g_i(t)}h
       \longrightarrow-(\omega y)\cdot\nabla g_i(t)
       \quad\hbox{in }\|\cdot\|_\sharp.
\]
For $h$ sufficiently small all these supports lie in fixed,
pairwise disjoint neighborhoods of the supports at time $t$:
rotation moves a point of radius $r$ by at most $3|\omega|r|h|$.
Expand the difference of $S_n$ by a multilinear telescope and use
its bound table on these neighborhoods. Each quotient converges
to the corresponding derivative entry, and the remaining entries
converge in the sharp norm. This proves the product rule
\eqref{IV-43-20} at every $t$, with the tuple rotated by $e^{t\omega}$.
It differentiates $S_n$ itself; no derivative is interchanged with
the cutoff limit. The expression does not involve $\mathcal R$.
\end{proof}

\begin{lemma}[One rotation group and cubic symmetry]\label{IV-28-26}
A subgroup of $\mathrm O(3)$ containing $W_3$ and
$\{e^{t\omega_0}:t\in\mathbb R\}$ for some nonzero
$\omega_0\in\mathfrak{so}(3)$ is $\mathrm O(3)$.
\end{lemma}
\begin{proof}
The given one-parameter group consists of all rotations about
some unit axis $n_0$. A signed coordinate permutation takes
$n_0$ to an axis $n_1\ne\pm n_0$: change one nonzero coordinate
if there are two, and otherwise permute the nonzero coordinate.
Conjugation therefore gives every rotation about $n_1$ as well.
Let $H$ be the subgroup generated by the two axis groups,
$O=Hn_0\subset S^2$, and
$A=\{\angle(n_0,x):x\in O\}$.
Since $H$ contains all rotations about $n_0$, a point $x$ belongs
to $O$ precisely when $\angle(n_0,x)\in A$.
By conjugating this assertion, for each $x\in O$ a point $y$
belongs to $O$ precisely when $\angle(x,y)\in A$.

Rotating $n_0$ about $n_1$ from angle zero to $\pi$ traces a
continuous path whose final point differs from $n_0$.
Thus $[0,\vartheta_0]\subset A$ for some $\vartheta_0>0$.
If $\vartheta,\vartheta'\in A$ and $\vartheta+\vartheta'\le\pi$,
choose points on one great circle at angular distances
$\vartheta$ and $\vartheta+\vartheta'$ from $n_0$.
The preceding characterization first places the former in $O$
and then the latter, so $\vartheta+\vartheta'\in A$.
Every angle in $[0,\pi]$ is a finite sum of angles in
$[0,\vartheta_0]$; hence $O=S^2$.
Given $R\in\mathrm{SO}(3)$, choose $h\in H$ with $hn_0=Rn_0$.
Then $h^{-1}R$ fixes $n_0$ and is a rotation about that axis,
so $R\in H$. Finally $-1\in W_3$, giving all of $\mathrm O(3)$.
\end{proof}

For a fixed $\omega$, impose the following two conditions on the
defining sequence, with the one function $\chi_1$ fixed above:
\begin{align}
 (\mathrm{UV})_\omega:\quad&
 \mathfrak U^{(j)}_{\omega,\mathcal R}(\underline f)\longrightarrow0
 \quad\hbox{for every separated smooth tuple and every admitting radius},
 \label{IV-43-9}\\
 (\mathrm{IR\text{-}fl})_\omega:\quad&
 \liminf_{\mathcal R\to\infty}\limsup_{j\to\infty}
       |\mathsf{Fl}^{(j)}_{\omega,\mathcal R}(\underline f)|=0
 \quad\hbox{for every separated smooth tuple}.
 \label{IV-43-10}
\end{align}
An admitting radius means $\operatorname{supp}f_i\subset
\{|y|<\mathcal R/2\}$ for every $i$. These conditions concern the
sequence itself; neither the limiting family nor the bound table
occurs in their definitions.

\begin{theorem}[Rotation criterion]\label{IV-43-24}
For a convergent sequence obeying \eqref{IV-43-3}, the following
assertions hold.
\begin{enumerate}
\item If $(\mathrm{UV})_\omega$ holds, then
\begin{equation}\label{IV-43-25}
 \lim_j\mathsf{Fl}^{(j)}_{\omega,\mathcal R}(\underline f)
       =-\dot S_n(\omega;\underline f).
\end{equation}
In this case $(\mathrm{IR\text{-}fl})_\omega$ holds if and only if
all these orbit derivatives vanish.
\item If both conditions hold for one nonzero $\omega_0$, then
$S$ is invariant under $\mathrm O(3)$ on separated $C_c^1$ tuples.
\item Without either condition, $S$ is $\mathrm O(3)$ invariant
if and only if the combined limit in \eqref{IV-43-21} vanishes
for every tuple and admitting radius for one nonzero $\omega_0$.
Equivalently, it vanishes for every $\omega$.
\item If $S$ is $\mathrm O(3)$ invariant and $(\mathrm{UV})_\omega$
holds, then the shell term tends to zero at each fixed admitting
radius. In particular $(\mathrm{IR\text{-}fl})_\omega$ holds with
ordinary limits in place of both upper and lower limits.
\end{enumerate}
\end{theorem}
\begin{proof}
Under $(\mathrm{UV})_\omega$, subtract the vanishing bulk term
from \eqref{IV-43-21}. The limit in \eqref{IV-43-25} is independent
of $\mathcal R$, so the double limit condition is equivalent to
$|\dot S_n(\omega;\underline f)|=0$.
If all derivatives vanish for a fixed $\omega_0$, apply
\eqref{IV-43-20} at every rotated tuple. The differentiable
function $t\mapsto S_n(e^{t\omega_0}\underline f)$ is constant.
This proves invariance on smooth tuples. For a separated $C_c^1$
tuple, convolution with a smooth approximate identity converges
in the sharp norm and has supports in fixed, slightly enlarged,
pairwise disjoint neighborhoods. The bound table and a
multilinear telescope pass the invariance to that tuple.
The invariance group of a family is a subgroup; it already
contains $W_3$ by \cref{IV-24-11}. Lemma~\ref{IV-28-26} gives
$\mathrm O(3)$. Conversely, invariance makes every orbit
derivative zero. Combining these observations with
\eqref{IV-43-21} proves the third assertion, and with
\eqref{IV-43-25} proves the fourth.
\end{proof}

The bound table and convergence used here hold for the constructed
growing-torus limits. They establish \eqref{IV-43-21}; they do not
establish either \eqref{IV-43-9} or \eqref{IV-43-10} separately.
For a volume-first limit, rotation invariance proved by the
independent volume-first argument makes the combined limit zero
along every torus sequence defining that limit.

\section{Classical comparison between consecutive cutoffs}
\label{sec:iv-classical-comparison}

The comparison uses two maps between the lattices. Block averaging gives
the lower action bound. For the reverse bound we refine a coarse
configuration while preserving its straight bond products. The refinement
must control the action to first order in the largest plaquette angle;
the estimate below provides this control without a gauge choice.

For $W\in G$, put
\[
 \varphi(W)=\arccos\bigl(\tfrac12\operatorname{tr}W\bigr),
 \qquad a(W)=1-\tfrac12\operatorname{tr}W
             =\tfrac12|W-\mathbf1|^2.
\]
Thus $|W-\mathbf1|=2\sin(\varphi(W)/2)$, and
\begin{equation}\label{IV-57-22}
 \tfrac12\varphi(W)^2\bigl(1-\varphi(W)^2/12\bigr)
 \le a(W)\le\tfrac12\varphi(W)^2.
\end{equation}
The inequalities follow from
$u^2/2-u^4/24\le1-\cos u\le u^2/2$ for $0\le u\le\pi$.
The lower bound follows by integrating
$\sin u\ge u-u^3/6$, itself obtained by integrating
$\cos u\ge1-u^2/2$.

We will also use the following logarithmic estimate. If
$X_1,\ldots,X_r\in\mathfrak g$ and $s=\sum_j|X_j|\le1/4$, then
$W=\prod_je^{iX_j}\ne-\mathbf1$ and
\begin{equation}\label{IV-57-23}
 Y=\tfrac1i\log W\in\mathfrak g,\qquad
 |Y|\le s,\qquad
 \left|Y-\sum_jX_j\right|\le s^2e^s.
\end{equation}
Here and below the logarithm is intrinsic. To verify the first bound,
note that $\varphi(W_1W_2)\le\varphi(W_1)+\varphi(W_2)$. When the sum
on the right is less than $\pi$, the identity
\[
 W_1W_2-\mathbf1
 =\tfrac12(W_1-\mathbf1)(W_2+\mathbf1)
  +\tfrac12(W_1+\mathbf1)(W_2-\mathbf1)
\]
and $|W+\mathbf1|=2\cos(\varphi(W)/2)$ bound the left side by
$2\sin((\varphi(W_1)+\varphi(W_2))/2)$. Monotonicity of sine on
$[0,\pi/2]$ proves the assertion; when the sum is at least $\pi$,
the assertion is immediate. Since
$\varphi(e^{iX})\le|X|$ and $|Y|=\varphi(W)$, iteration gives
$|Y|\le s$. Expanding the product of exponentials in its absolutely
convergent series bounds the terms of total degree at least two by
$e^s-1-s\le s^2e^s/2$. Apply the same bound to $e^{iY}$ and subtract
the two first-order expansions to obtain \eqref{IV-57-23}.

Let $\mathbb Z_{\underline n}=\prod_{\mu=1}^3\mathbb Z/n_\mu\mathbb Z$,
where $n_\mu\ge2$. A configuration is a family $U_\mu(x)\in G$.
Write
\[
 P_{\alpha\beta}(x)
 =U_\alpha(x)U_\beta(x+e_\alpha)
   U_\alpha(x+e_\beta)^{-1}U_\beta(x)^{-1},
 \qquad
 E_{\alpha\beta}(U)=\sum_x a(P_{\alpha\beta}(x)),
\]
and set $E(U)=E_{12}(U)+E_{13}(U)+E_{23}(U)$.
The reversal of the two directions inverts $P_{\alpha\beta}$ and
does not change its angle or action. Let $\vartheta(U)$ be the largest
plaquette angle. For a site $x$, let $c(x)$ be the largest angle of
the six faces of its coordinate cube, namely the plaquettes
$P_{\alpha\beta}(x)$ and $P_{\alpha\beta}(x+e_\gamma)$ with
$\alpha<\beta$ and $\gamma\notin\{\alpha,\beta\}$.

Fix an integer $L\ge2$ and a direction $\lambda$. Enlarge the
$\lambda$ side by $L$, and let $\iota_\lambda$ multiply that coordinate
by $L$. Every site of the enlarged box has a unique expression
$x^{(j)}=\iota_\lambda x+je_\lambda$, $0\le j<L$.
If $\vartheta(U)<\pi$, define
$\Theta_{\lambda\kappa}(x)=\frac1i\log P_{\lambda\kappa}(x)$ and set
\begin{equation}\label{IV-57-50}
 \begin{aligned}
 U'_\lambda(x^{(j)})&=\mathbf1 &&(j<L-1),&
 U'_\lambda(x^{(L-1)})&=U_\lambda(x),\\
 U'_\kappa(x^{(j)})&=e^{i(j/L)\Theta_{\lambda\kappa}(x)}U_\kappa(x)
 &&(\kappa\ne\lambda).
 \end{aligned}
\end{equation}
We denote this refinement by $U'=\mathsf{Rfn}_\lambda U$.

\begin{lemma}[Refinement in one direction]\label{IV-57-51}
Suppose $\vartheta_0=\vartheta(U)\le10^{-4}$, and let $\mu,\nu$ be
the directions other than $\lambda$. The refinement preserves the
straight coarse bond products and the transverse bonds at embedded sites:
\[
 \prod_{j=0}^{L-1}U'_\lambda(x^{(j)})=U_\lambda(x),\qquad
 U'_\kappa(\iota_\lambda x)=U_\kappa(x)\quad(\kappa\ne\lambda).
\]
Its plaquettes containing the refined direction satisfy
$P'_{\lambda\kappa}(x^{(j)})=e^{i\Theta_{\lambda\kappa}(x)/L}$.
For the transverse plaquettes, put
$Y_0=\frac1i\log P_{\mu\nu}(x)$ and
$Y_1=\operatorname{Ad}(U_\lambda(x))\frac1i\log P_{\mu\nu}(x+e_\lambda)$.
Then
\begin{equation}\label{IV-57-52}
 \varphi(P'_{\mu\nu}(x^{(j)}))
 \le |(1-j/L)Y_0+(j/L)Y_1|+54c(x)^2.
\end{equation}
Consequently,
\begin{equation}\label{IV-57-53}
 \begin{aligned}
 E_{\lambda\kappa}(U')&\le L^{-1}(1+\vartheta_0^2)
                         E_{\lambda\kappa}(U),\\
 E_{\mu\nu}(U')&\le L(1+\vartheta_0^2)E_{\mu\nu}(U)
                         +221L\vartheta_0E(U),\\
 \vartheta(U')&\le\vartheta_0+54\vartheta_0^2\le2\vartheta_0.
 \end{aligned}
\end{equation}
\end{lemma}

\begin{proof}
The bond identities follow directly from \eqref{IV-57-50}.
For $j<L-1$, the two $\lambda$ bonds of a plaquette equal the identity,
so its two transverse factors give $e^{i\Theta_{\lambda\kappa}(x)/L}$.
At $j=L-1$, the same product is
$P_{\lambda\kappa}(x)e^{-i(L-1)\Theta_{\lambda\kappa}(x)/L}$,
which has the same value.

For a transverse plaquette write
\[
 \begin{gathered}
 A_1=\Theta_{\lambda\mu}(x),\quad
 A_2=\Theta_{\lambda\nu}(x+e_\mu),\quad
 A_3=\Theta_{\lambda\mu}(x+e_\nu),\quad
 A_4=\Theta_{\lambda\nu}(x),\\
 u_\mu=U_\mu(x),\quad u_\nu=U_\nu(x),\quad
 u'_\nu=U_\nu(x+e_\mu),\quad u'_\mu=U_\mu(x+e_\nu).
 \end{gathered}
\]
For $0\le t\le1$ consider
\[
 W(t)=e^{itA_1}u_\mu e^{itA_2}u'_\nu
          u_\mu'^{-1}e^{-itA_3}u_\nu^{-1}e^{-itA_4}.
\]
Then $W(j/L)=P'_{\mu\nu}(x^{(j)})$. Moving the unitary factors past
the exponentials gives
\[
 W(t)=e^{itA_1}e^{it\operatorname{Ad}(u_\mu)A_2}e^{iY_0}
       e^{-it\operatorname{Ad}(u_\nu)A_3}e^{-itA_4}.
\]
The five exponents and $Y_1$ have norms at most $c=c(x)$.
By \eqref{IV-57-23},
\[
 \tfrac1i\log W(t)
 =Y_0+t\bigl(A_1+\operatorname{Ad}(u_\mu)A_2
                 -\operatorname{Ad}(u_\nu)A_3-A_4\bigr)+e(t),
 \qquad |e(t)|\le25e^{5c}c^2\le27c^2.
\]
At $t=1$, substitution of the four defining plaquette products
cancels the three successive pairs of $\lambda$ bonds and gives
$W(1)=U_\lambda(x)P_{\mu\nu}(x+e_\lambda)U_\lambda(x)^{-1}$.
Thus $\frac1i\log W(1)=Y_1$. Subtracting this identity from the
preceding expression gives
$\frac1i\log W(t)=(1-t)Y_0+tY_1+e(t)-te(1)$.
This proves \eqref{IV-57-52} and the angle bound in
\eqref{IV-57-53}.

For an old plaquette angle $\varphi\le\vartheta_0$,
\eqref{IV-57-22} implies
\begin{equation}\label{IV-57-56}
 \tfrac12\varphi^2\le(1+\vartheta_0^2)a.
\end{equation}
Indeed, $(1+y)(1-y/12)\ge1$ for $0\le y\le1$.
Each plaquette containing $\lambda$ becomes $L$ plaquettes of angle
$\varphi/L$, which proves its energy estimate. For a transverse
plaquette put $p=|(1-t)Y_0+tY_1|$. Since $p\le c\le10^{-4}$,
\[
 a(W(t))\le\tfrac12(p+54c^2)^2
 \le\tfrac12\bigl((1-t)|Y_0|^2+t|Y_1|^2\bigr)+55c^3.
\]
The last bound uses $54c^2p+1458c^4\le55c^3$.
Summing over $j=0,\ldots,L-1$ and then $x$, the quadratic term is
$\frac L2\sum_x|Y_0(x)|^2$: the sums of $|Y_0|^2$ and $|Y_1|^2$
coincide by translation and unitary invariance. It is at most
$L(1+\vartheta_0^2)E_{\mu\nu}(U)$ by \eqref{IV-57-56}.
Each plaquette belongs to two of the coordinate cubes, and hence
\[
 \sum_xc(x)^3
 \le\vartheta_0\sum_x\sum_{\text{faces of }x}\varphi(P)^2
 \le4(1+\vartheta_0^2)\vartheta_0E(U).
\]
The remaining contribution is at most
$220(1+10^{-8})L\vartheta_0E(U)\le221L\vartheta_0E(U)$.
\end{proof}

\begin{proposition}[Prolongation]\label{IV-57-59}
Let $\vartheta_0=\vartheta(U)\le\frac14\,10^{-4}$ and define
\[
\widetilde U=\mathsf{Prl}_LU
=\mathsf{Rfn}_3\mathsf{Rfn}_2\mathsf{Rfn}_1U
\]
on $\mathbb Z_{L\underline n}$. Then
\begin{equation}\label{IV-57-60}
 \prod_{j=0}^{L-1}\widetilde U_\mu(Lx+je_\mu)=U_\mu(x),
 \qquad L E(\widetilde U)\le(1+4000L^4\vartheta_0)E(U),
 \qquad \vartheta(\widetilde U)\le8\vartheta_0.
\end{equation}
The bond product is ordered by increasing $j$.
\end{proposition}

\begin{proof}
After $i$ refinements the largest angle is at most $2^i\vartheta_0$.
Thus the input angle of each of the three steps is at most $10^{-4}$,
and \cref{IV-57-51} applies. At the step refining direction $\mu$,
the straight product becomes the old $\mu$ bond. At the other two
steps, the transverse bonds at embedded sites are preserved. Applying
these identities in their order proves the first assertion.

Write $E^{(i)}_{\alpha\beta}$ and $E^{(i)}$ for the energies after
$i$ steps, and put $\epsilon=16\vartheta_0^2\le10^{-8}$. The energy
inequalities in \eqref{IV-57-53} give, at step $i$,
\[
 E^{(i)}_{i\kappa}\le\frac{1+\epsilon}{L}E^{(i-1)}_{i\kappa},
 \qquad
 E^{(i)}_{\mu\nu}\le L(1+\epsilon)E^{(i-1)}_{\mu\nu}
             +221\cdot2^{i-1}L\vartheta_0E^{(i-1)},
\]
where $\mu,\nu\ne i$. In particular $E^{(i)}\le2LE^{(i-1)}$.
For the pairs $12$, $13$, and $23$, the step with the error term is,
respectively, the third, second, and first. The displayed recurrence
therefore yields
\[
 \begin{aligned}
 E^{(3)}_{12}&\le (1+\epsilon)^3L^{-1}E^{(0)}_{12}
                         +3536L^3\vartheta_0E^{(0)},\\
 E^{(3)}_{13}&\le (1+\epsilon)^3L^{-1}E^{(0)}_{13}
                         +884(1+\epsilon)L\vartheta_0E^{(0)},\\
 E^{(3)}_{23}&\le (1+\epsilon)^3L^{-1}E^{(0)}_{23}
                         +221(1+\epsilon)^2L^{-1}\vartheta_0E^{(0)}.
 \end{aligned}
\]
After summing and multiplying by $L$, use
$(1+\epsilon)^3-1\le4\epsilon\le\vartheta_0$,
$884(1+\epsilon)\le885$, $221(1+\epsilon)^2\le222$, and
$885L^2+223\le464L^4$. These give the coefficient $4000L^4$ in
\eqref{IV-57-60}.
\end{proof}

\subsection{Weighted averaging}

For the comparison of smeared actions we need averaging on arbitrary
sets of plaquettes. Suppose that a fine lattice of spacing $h/L$
is averaged to a lattice of spacing $h$. For a coarse plaquette $p'$,
let $\mathcal N(p')$ be the fine plaquettes all of whose corners lie
in the union of the four fine blocks centered at the corners of $p'$;
put $\mathcal N(S)=\bigcup_{p'\in S}\mathcal N(p')$.
The one-step bound \eqref{II-6-77} is
\begin{equation}\label{eq:iv-arbitrary-set-averaging}
 \sum_{p'\in S}a_{p'}(\mathcal MW)
 \le L(1+C_{\mathrm{en}}\omega)^2
             \sum_{p\in\mathcal N(S)}a_p(W),
 \qquad
 \omega=\max_{p\in\mathcal N(S)}|\partial W(p)-\mathbf1|
       \le\varpi'_L.
\end{equation}
This holds for every set $S$, not only all plaquettes. For clarity,
the reason is that the positive linear curvature kernel of
Propositions~\ref{I-10-96}--\ref{I-10-97} has counting
$\ell^2$ norm at most $\sqrt L$ after restriction to
$\mathcal N(S)$, while its nonlinear remainder at $p'$ is at most
$C^{\mathrm{SO}}_0\omega_{p'}^2$. A fine plaquette belongs to at
most twelve coarse cones: its base lies in one coarse block, whose
center must be a corner of each of the coarse plaquettes. Consequently
\[
 \sum_{p'\in S}\omega_{p'}^4
 \le12\omega^2\sum_{p\in\mathcal N(S)}|\partial W(p)-\mathbf1|^2.
\]
The triangle inequality gives the curvature bound with factor
$\sqrt L(1+C_{\mathrm{en}}\omega)$, where
$C_{\mathrm{en}}=\sqrt{12}\,C^{\mathrm{SO}}_0L^{-1/2}$;
squaring and using $2a_p=|\partial W(p)-\mathbf1|^2$ proves
\eqref{eq:iv-arbitrary-set-averaging}.

For a plaquette $p$, let $x_p$ denote its barycenter. If
$p\in\mathcal N(p')$, choose a corner $z$ of $p$ in the fine block
of a corner $y$ of $p'$. In the torus distance,
\begin{equation}\label{eq:iv-cone-barycenter-distance}
 d_\infty(x_p,x_{p'})
 \le d_\infty(x_p,z)+d_\infty(z,y)+d_\infty(y,x_{p'})
 <\frac{h}{2L}+h\le\frac32h.
\end{equation}
This estimate uses only the defining union of blocks of the cone.

\begin{lemma}[Weighted one-step action]\label{lem:iv-weighted-averaging}
Let $W$ have curvature at most $\omega\le\varpi'_L$ on the fine
lattice. Write $\eta$ and $\eta/L$ for the two dimensionless spacings,
and set
\[
 A^\eta(c,V)=\eta^{-1}\sum_p c(x_p)a_p(V),\qquad
 q=(1+C_{\mathrm{en}}\omega)^2.
\]
For every nonnegative real Lipschitz function $c$ on the torus,
\begin{equation}\label{eq:iv-weighted-averaging}
 A^\eta(c,\mathcal MW)
 \le q\left[A^{\eta/L}(c,W)
      +\tfrac32h\operatorname{Lip}_\infty(c)A^{\eta/L}(1,W)\right].
\end{equation}
\end{lemma}

\begin{proof}
Put $v=\frac32h\operatorname{Lip}_\infty(c)$ and
$S_\lambda=\{p':c(x_{p'})>\lambda\}$. By
\eqref{eq:iv-cone-barycenter-distance},
$\mathcal N(S_\lambda)\subset\{p:c(x_p)+v>\lambda\}$.
Apply \eqref{eq:iv-arbitrary-set-averaging} to $S_\lambda$ and
integrate over $\lambda\ge0$. The elementary identity
$c(x)=\int_0^\infty\mathbf1_{\{c(x)>\lambda\}}\,d\lambda$
gives
\[
 \sum_{p'}c(x_{p'})a_{p'}(\mathcal MW)
 \le Lq\sum_p(c(x_p)+v)a_p(W).
\]
Multiplication by $\eta^{-1}$ proves the assertion.
\end{proof}

\subsection{The action on a fibre and the excess of a competitor}

We specify the uniform global hypotheses used in this comparison.
Fix the geometric constants, the actual shift systems, and the octet
data of the sourced construction. We impose the hypotheses on these
data in Theorem~\ref{II-127-29} and Proposition~\ref{II-42-39}, including the
global real-background hypotheses and $\alpha_0\le a_{\rm H}$.
They are required at the radii $(\alpha_0,\alpha_1)$, for every
calibration $\tau\ge\gamma^{-2}$. The localization floors for the
real response equations~\eqref{II-8-18}--\eqref{II-8-20} are also imposed.
In particular the covariance comparison has error at most
$\gamma_{\rm r}/2$, the two interaction potentials have decay
$\lambda_V>\log9216+2$, and
$\mathsf C_{\mathcal A}\mathfrak b(\tau)\ge3$, where
\[
 \varepsilon(\tau)=A_0\tau^{-1/2}(\log\tau)^{\mathfrak p_0},\qquad
 \delta(\tau)=A_1\tau^{-1/2}(\log\tau)^{\mathfrak p_0},\qquad
 \mathfrak b(\tau)=\sqrt\tau\,\delta(\tau).
\]
We use the range on which $\varepsilon,\delta$ are continuous,
decreasing, and tend to zero. Assume
\begin{equation}\label{eq:iv-classical-room}
 K\ge1,\qquad \ell=2L^m,\qquad 2L^{m+1}\ge C_LM.
\end{equation}
Then each step below is in the global setting of the cited results.
Indeed the side in level-$n$ units is $2L^{m+K-n}\ge8M$;
the constant sequence of depth $n\le K$ is separated and its
level-$n$ lattice has at least $8R_1\mathfrak M$ sites in each
direction. These are the geometric hypotheses of the background and
response results. The actual whole-data domain in Proposition~\ref{II-42-39}
includes the open response neighborhood and identifies the actual
shifts with the constructed critical backgrounds. Thus its global
sumrules, rather than an abstract provision alone, apply here.

Write $\varpi(V)=\max_p|\partial V(p)-\mathbf1|$.
For the constant sequence of depth $n$, denote its real critical
background by $U_n(V)$, with its residual gauge freedom understood,
and put
\[
 \mathsf{ab}_n(c,V)=A^{\eta_n}(c,U_n(V)),\qquad
 \mathsf{ab}_n(V)=\mathsf{ab}_n(1,V),\qquad\eta_n=L^{-n}.
\]
These actions are independent of the representative. They are also
invariant under a gauge of the level-$n$ datum: extend that gauge
arbitrarily to the bare sites, use covariance of the constraint and
critical orbit, and then gauge invariance of each plaquette action.
At level zero $\mathsf{ab}_0(c,V)=\sum_pc(x_p)a_p(V)$ whenever the
identity background is regular. The fields $V_k^{\rm bf}(W)$ and
$V_W[\mathcal A]$ are the axial averaged background and its fibre
chart from Propositions~\ref{II-40-99}--\ref{II-42-39}; the variable
$\mathcal A$ is prescribed on the nonpivot bonds.

Let $C_{\rm db}=8L^3$, let $\varepsilon_{\mathcal W}$ be the uniform
radius of Proposition~\ref{II-42-39}, and set
\[
 \begin{gathered}
 A_L=\tfrac34(L-1)^2,\qquad
 C_{\rm AC}=8(L-1)A_L+(1+\pi/3)\bigl(L(L-1)+8A_L\bigr),\\
 C_{\rm ax}=\max\{8A_L,C_{\rm AC}\},\qquad
 C_{\rm pr}=8A_L(L-1)+C_{\rm AC},\qquad
 C'_{\rm pr}=1+16\pi C_{\rm pr},\\
 C_{\rm cp}=B_3\max\{2C'_{\rm pr},4\pi L^2\},\qquad
 C_{\rm cub}=50K_PC_{\mathrm E,1}/\mathsf C_{\mathcal A}.
 \end{gathered}
\]
Here $K_P$ is the actual potential bound, and
$C_{\mathrm E,1}=2048\sum_{j\ge0}(j+2)e^{-j}$ is the rooted
domain-sum bound of Proposition~\ref{II-18-26}.
We use the response constants $\delta_{\rm lc}$ and $B_5\ge8$ of
equations~\eqref{II-8-18}--\eqref{II-8-20}. Let $S_{\rm e}$ bound both
row and column sums of its kernel
$e^{-\delta_{\rm lc}d(y,c_-)}$, and put
\begin{equation}\label{IV-59-3}
 \begin{gathered}
 \rho_{\rm no}=\min\{r_{\mathrm B},\varrho_{\mathrm R}^0\},\qquad
 C_{\rm nrg}=108B_5^2S_{\rm e}^2(1+24C_{\rm db}^2),\\
 \varepsilon_{\rm cp}=\min\left\{
 \varepsilon_{\mathcal W}\alpha_0,
 \frac{a_{\mathrm R}}{C_{\mathfrak P}\mathfrak M B_3},
 B_3^{-1}\right\}.
 \end{gathered}
\end{equation}
All these constants are independent of the cutoff, level, and volume.

Put $K^{\rm s}_{\mathrm B}=K^{\rm s}_{\mathrm B}(\rho_{\mathrm s})$,
$C_H=4C_{\rm db}\max\{2,K^{\rm s}_{\mathrm B}\}$, and
\[
 c_\rho=\max\left\{2B_3,
       \frac{C_HA_1}{(1-L^{-3/2})A_0}\right\}.
\]
Let $\tau_\natural\ge\gamma^{-2}$ be the least calibration at which
\begin{equation}\label{eq:iv-classical-fixed-calibration}
 \begin{gathered}
 2c_\rho\varepsilon(\tau_\natural)
 \le\min\{\varepsilon_{\mathcal W}\alpha_0,c^{\rm SO}_2,
                                      \varepsilon_{\mathrm v}/4\},\qquad
 \delta(\tau_\natural)\le\varepsilon_{\mathcal W},\quad
 4C_{\rm db}\delta(\tau_\natural)
 \le\min\{1,r_{\mathrm B},r^{\rm s}_{\mathrm B}/2\}.
 \end{gathered}
\end{equation}
This minimum exists by continuity, monotonicity, and decay of the
thresholds. The same conditions hold at larger calibrations, where
also $\varepsilon(\tau)\le1$ by the global step hypotheses. Here
$r^{\rm s}_{\mathrm B}=r^{\rm s}_{\mathrm B}(\rho_{\mathrm s})$.
Define the competitor radius
\begin{equation}\label{IV-59-5}
\begin{split}
 e_{\rm cp}=\min\biggl\{&\frac{\omega_M}{2},
 \frac{\varepsilon_{\rm cp}}{2B_3},
 \frac{\gamma_{\rm r}L^2}{4\pi C_{\rm cub}C_{\rm ax}(1+2B_3)},\\
 &\frac{L^2\min\{\rho_{\rm no},(66B_5S_{\rm e})^{-1}\}}
           {2\pi C_{\rm db}C_{\rm ax}(1+2B_3)},
 \frac{L^2\delta(\tau_\natural)}{\pi C_{\rm ax}(1+2B_3)}
 \biggr\},\qquad \omega_M=(200L^2)^{-1}.
\end{split}
\end{equation}
For $0<e\le e_{\rm cp}$, set
$h=\frac\pi2C_{\rm ax}(1+2B_3)L^{-2}e$. The definition gives
\begin{equation}\label{eq:iv-competitor-smallness}
 \begin{gathered}
 2B_3e\le\varepsilon_{\rm cp},\qquad
 C_{\rm cub}h\le\gamma_{\rm r}/8,\qquad
 2C_{\rm db}h\le\rho_{\rm no}/2,\\
 B_5S_{\rm e}\,2C_{\rm db}h\le1/132,\qquad
 h\le\delta(\tau_\natural)/2.
 \end{gathered}
\end{equation}
The global radius also satisfies
$\varepsilon_{\mathcal W}\le\bar\varepsilon$ and
$4K_2\varepsilon_{\mathcal W}\le\varepsilon_{\mathrm v}$,
where $K_2=2B_3L^{-2}+8C_{\rm db}\ge8$.
Consequently $4e<\varepsilon_{\mathrm v}$.

We first record two consequences of the actual chart. For two axial
configurations $V^1,V^2$ with the same block average and respective
curvature bounds $a_1,a_2<\omega_M$,
\begin{equation}\label{IV-59-8}
 \varphi\bigl(V^1(b)V^2(b)^{-1}\bigr)
             \le\tfrac\pi2 C_{\rm ax}(a_1+a_2).
\end{equation}
To see this, use the double-block estimates in the proof of
Proposition~\ref{II-8-68}, and equation~\eqref{I-10-131}.
An internal bond differs from the identity by at most $8A_La_i$.
A crossing bond differs from its coarse average by at most
$C_{\rm AC}a_i$. Subtract the two bounds, using equality of the
coarse averages in the second case, and use
$\varphi(W)\le\frac\pi2|W-\mathbf1|$.

If $\varpi(W)<\varepsilon_{\mathcal W}\alpha_0$ and
$\sup|\mathcal A|<\delta(\tau_\natural)$, write
$V_W[\mathcal A]=e^{iB'}V_k^{\rm bf}(W)$. Then
\begin{equation}\label{IV-59-10}
 B'|_{\mathfrak X}=\mathcal A,\qquad
 \sup|B'|\le2C_{\rm db}\sup|\mathcal A|,\qquad
 \|B'\|_2^2\le(1+24C_{\rm db}^2)\|\mathcal A\|_2^2.
\end{equation}
Here the counting norm on either bond set uses positively oriented
bonds. For the last two bounds, fix a coarse bond $c$. Its pivot value
depends only on the nonpivot field in the double block $\mathsf P_c$.
Replace $\mathcal A$ by its restriction to that double block and zero
outside. At every calibration $\tau\ge\tau_\natural$ for which
the supremum of this restriction is less than $\delta(\tau)$,
the chart bound \eqref{II-42-41} gives pivot size less than
$2C_{\rm db}\delta(\tau)$. The chart is independent of calibration.
By continuity of $\delta$ and its limit zero, taking the infimum over
these calibrations bounds that pivot by
$2C_{\rm db}\max_{b\in\mathfrak X^+\cap\mathsf P_c}|\mathcal A(b)|$.
A nonpivot bond belongs to at most six double blocks, since its initial
point belongs to one block and that block is incident to six coarse
bonds. Squaring and summing the pivot bound proves
\eqref{IV-59-10}. Equality on the nonpivot bonds is the defining
property of the fibre chart; the pivot variables solve its constraint.

\begin{proposition}[Coercivity on the fibre]\label{IV-59-11}
Let $0\le k<K$ and $\tau\ge\gamma^{-2}$, and suppose that $W$ belongs
to the actual global domain
$\mathcal W_{k+1}(\tau;\alpha_0,\alpha_1)$ of
Proposition~\ref{II-42-39}. Put $g=\tau^{-1/2}$ and take a real
nonpivot field $B$ with $\sup|B|<\mathfrak b(\tau)$.
For
\[
 F_W(B)=\mathsf{ab}_k(V_W[gB]),\qquad
 G_W(B)=\sum_y\sum_{x\in B(y)\setminus\{y\}}
       \left[1-\tfrac12\operatorname{Re}\operatorname{tr}
                                  V_W[gB](y,x)\right],
\]
one has
\begin{equation}\label{IV-59-12}
 \tau\bigl(F_W(B)-F_W(0)+G_W(B)\bigr)
 \ge\bigl(\gamma_{\rm r}/4-C_{\rm cub}g\sup|B|\bigr)\|B\|_2^2.
\end{equation}
Moreover $F_W(0)=L^{-1}\mathsf{ab}_{k+1}(W)$, $G_W(B)\ge0$,
and $G_W(B)=0$ when the chart field is axial.
\end{proposition}

\begin{proof}
The global action and gauge sumrules, with their actual background
identification, give
\begin{equation}\label{IV-59-14}
 \tau\bigl(F_W(B)-F_W(0)+G_W(B)\bigr)
 =\tfrac12(B,\mathsf A_WB)_1
   -\sum_{\natural\in\{\mathrm A,\mathrm G\}}\sum_{Y\in\mathcal D_k}
                  \mathbb V^\natural(Y;P_W|_Y,B|_Y).
\end{equation}
Here $P_W$ is the plain triple of the axial top background and
$\mathsf A_W=\mathbf C_*(P_W)^{-1}$; see
equations~\eqref{II-41-27}, \eqref{II-41-84}--\eqref{II-41-85}, Theorem~\ref{II-127-29} and
Proposition~\ref{II-42-39}. All the triples on the right belong to their
cores by that proposition. The reference covariance satisfies
$(B,\mathbf C_{*\rm r}^{-1}B)_1\ge\gamma_{\rm r}\|B\|_2^2$,
and $\|\mathsf A_W-\mathbf C_{*\rm r}^{-1}\|\le\gamma_{\rm r}/2$.
The real part of the quadratic term is therefore at least
$\gamma_{\rm r}\|B\|_2^2/4$.

It remains to use the third-order vanishing of the two potentials,
which is stronger than their uniform smallness. For a domain $Y$,
put $d=d_k(Y)$, $b_Y=\max_{b\in\mathfrak X^{\circ+}\langle Y\rangle}|B(b)|$,
and $r=\mathsf C_{\mathcal A}\mathfrak b(\tau)\ge3$.
There are at most $24M^3(d+1)$ bonds in this domain.
If $b_Y=0$, its potential is zero. Otherwise
$z\mapsto\mathbb V^\natural(Y;P_W|_Y,zB|_Y)$ is holomorphic for
$|z|<r/b_Y$, vanishes to order three, and, by the actual potential
bound \eqref{II-53-14}, has modulus at most
\[
 K_P\delta(\tau)e^{-\lambda_Vd}
            \bigl(1+24(d+1)r^2\bigr)
 \le25K_P\delta(\tau)(d+1)r^2e^{-\lambda_Vd}.
\]
Apply the maximum principle to the function divided by $z^3$ on
smaller concentric discs and let their radii tend to $r/b_Y$. This gives
\[
 |\mathbb V^\natural(Y;P_W|_Y,B|_Y)|
 \le\frac{25K_P}{\mathsf C_{\mathcal A}}g\sup|B|
       (d+1)e^{-\lambda_Vd}
            \sum_{b\in\mathfrak X^{\circ+}\langle Y\rangle}|B(b)|^2.
\]
For a fixed bond, the domains in this sum contain its initial point.
Their weighted sum is at most $C_{\mathrm E,1}$ by
Proposition~\ref{II-18-26}, since $\lambda_V>\log9216+2$.
Summing over both potentials gives
$C_{\rm cub}g\sup|B|\|B\|_2^2$, proving \eqref{IV-59-12}.

The background over $V_k^{\rm bf}(W)$ is a representative of the
background over $W$, by restriction of criticality and uniqueness
as proved in Proposition~\ref{II-42-39}.
Changing the action spacing from $\eta_k$ to $\eta_{k+1}$ contributes
the factor $L$, so $F_W(0)=L^{-1}\mathsf{ab}_{k+1}(W)$.
Every term defining $G_W$ is nonnegative and vanishes in axial gauge.
\end{proof}

Call a bare configuration $U'$ a competitor of radius $e$ at level $t$
if $1\le t\le K$ and
\begin{equation}\label{IV-59-16}
 \varpi(\mathcal M^kU')\le eL^{-2(t-k)}\qquad(0\le k\le t).
\end{equation}
Its excess is
$\mathsf{exc}_t(U')=A^{\eta_t}(U')-\mathsf{ab}_t(\mathcal M^tU')$.

\begin{proposition}[Small competitors]\label{IV-59-17}
If $0<e\le e_{\rm cp}$ and $U'$ satisfies \eqref{IV-59-16}, all
backgrounds $U_k(V'_k)$, $V'_k=\mathcal M^kU'$, exist.
There are real nonpivot fields $\mathcal A_k$, $0\le k<t$, with
\begin{equation}\label{eq:iv-competitor-field}
 \sup|\mathcal A_k|\le h_k
 :=\tfrac\pi2C_{\rm ax}(1+2B_3)eL^{-2(t-k)}\le h,
\end{equation}
such that the axial representative of $V'_k$ equals
$V_{V'_{k+1}}[\mathcal A_k]$. One has
\begin{equation}\label{IV-59-19}
 \mathsf{ab}_k(V'_k)-L^{-1}\mathsf{ab}_{k+1}(V'_{k+1})
                 \ge\tfrac18\gamma_{\rm r}\|\mathcal A_k\|_2^2
\end{equation}
and
\begin{equation}\label{IV-59-20}
 \mathsf{exc}_t(U')\ge\tfrac18\gamma_{\rm r}
           \sum_{k<t}L^{t-k}\|\mathcal A_k\|_2^2\ge0.
\end{equation}
In particular $A^{\eta_t}(U_k(V'_k))$ is nonincreasing in $k$.
\end{proposition}

\begin{proof}
Since $e\le\varepsilon_{\mathcal W}\alpha_0/2<\bar\varepsilon$,
the real-background theorem constructs $U_k(V'_k)$ for $k\ge1$.
At $k=0$ the identity background is regular: its curvature is at
most $e$ and its current at most $4e<\varepsilon_{\mathrm v}$.
Fix $k<t$, put $W=V'_{k+1}$, and normalize $V'_k$ to axial gauge
by its contour means. This gauge equals the identity on the coarse
sites and preserves the average $W$.
Proposition~\ref{II-42-39} puts $W$ in the actual global domain at
$\tau_\natural$ and constructs $V_0=V_k^{\rm bf}(W)$.
Its curvature estimate, applied with any radius strictly above
$\varpi(W)$ and then taking the infimum, gives
\[
 \varpi(V_0)\le2B_3L^{-2}\varpi(W)
       \le2B_3eL^{-2(t-k)}\le L^{-2}\varepsilon_{\rm cp}.
\]
The two axial fields have the same average. Apply
\eqref{IV-59-8}; their relative intrinsic logarithm is defined
and has norm at most $h_k\le\delta(\tau_\natural)/2$.
Its restriction to the nonpivot bonds is $\mathcal A_k$.
The inverse property of the fibre chart identifies the axial
configuration with $V_W[\mathcal A_k]$.

Apply \cref{IV-59-11} with $B=\sqrt{\tau_\natural}\mathcal A_k$.
The gauge term vanishes, and
$C_{\rm cub}\sup|\mathcal A_k|\le\gamma_{\rm r}/8$ by
\eqref{eq:iv-competitor-smallness}. Divide by $\tau_\natural$
and use gauge invariance of the background action to obtain
\eqref{IV-59-19}. Finally
$A^{\eta_t}(U_k(V'_k))=L^{t-k}\mathsf{ab}_k(V'_k)$.
Multiply \eqref{IV-59-19} by $L^{t-k}$ and sum in $k$;
the left sides telescope to the excess, proving \eqref{IV-59-20}.
\end{proof}

\begin{lemma}[Plaquette differences between successive backgrounds]
\label{IV-59-21}
Under the hypotheses of \cref{IV-59-17}, for each $k<t$ one can
choose representatives $\widehat U\in\mathcal C(\mathbf B_k,V'_k)$
and $\check U\in\mathcal C(\mathbf B_{k+1},V'_{k+1})$ and elements
$u_p,R_p\in G$ such that
\begin{equation}\label{IV-59-22}
 \partial\widehat U(p)=u_pR_p\partial\check U(p)u_p^{-1},\qquad
 \eta_k^{-1}\sum_p|R_p-\mathbf1|^2
                      \le C_{\rm nrg}\|\mathcal A_k\|_2^2.
\end{equation}
\end{lemma}

\begin{proof}
Put $V_0=V_k^{\rm bf}(V'_{k+1})$ and use
\eqref{IV-59-10} to write the axial version of $V'_k$ as
$e^{iB'}V_0$. Then
\[
 \sup|B'|\le2C_{\rm db}h,\qquad
 \|B'\|_2^2\le(1+24C_{\rm db}^2)\|\mathcal A_k\|_2^2.
\]
For $k\ge1$, the real response applies at datum radius
$\varepsilon_{\rm cp}$: the base has curvature at most
$L^{-2}\varepsilon_{\rm cp}<\varepsilon_{\rm cp}$,
$C_{\mathfrak P}\mathfrak M B_3\varepsilon_{\rm cp}\le a_{\mathrm R}$,
and $\sup|B'|\le\rho_{\rm no}/2$.
The constant sequence has all pillars of level $k$ and size one
in level-$k$ units. Thus equations~\eqref{II-8-18}--\eqref{II-8-20} give
representatives related, up to gauge, by $e^{i\eta_k H}$, where
\[
 |H(b)|\le B_5m_{y(b)},\qquad
 |\nabla_{\check U} H(x)|\le B_5m_{y(x)},\qquad
 m_y=\sum_{c>0}e^{-\delta_{\rm lc}d(y,c_-)}|B'(c)|.
\]
In particular $\sup|H|\le B_5S_{\rm e}\,2C_{\rm db}h\le1/132$.
Let $\phi_p$ be the maximum of $|H|$ on the four bonds of $p$.
The real plaquette formula of Lemma~\ref{I-10-120}, with
$\eta_k\le1/13$ and base curvature at most $\eta_k^2$, gives
\[
 |R_p-\mathbf1|
 \le\eta_k^2|D_{\check U} H(p)|+11\eta_k^2\phi_p^2
                   +6\eta_k\phi_p|\partial\check U(p)-\mathbf1|
 \le3\eta_k^2B_5\max_{x\text{ corner of }p}m_{y(x)}.
\]
Here the derivative and curl are covariant with respect to the
unperturbed background $\check U$. Indeed
$|D_{\check U}H(p)|\le2B_5\max m_y$, $\phi_p\le B_5\max m_y$,
$\phi_p\le1/132$, and $2+11/132+6/13<3$.
The product-expansion hypotheses hold since
$4\eta_k\phi_p<1/4$ and $\eta_k^2<1/4$.
Every bare site is a corner of twelve plaquettes and each pillar
contains $\eta_k^{-3}$ sites. The row and column sum bounds on
the response kernel therefore imply
\[
 \sum_p\max_{x\text{ corner of }p}m_{y(x)}^2
 \le12\eta_k^{-3}\sum_ym_y^2
 \le12\eta_k^{-3}S_{\rm e}^2\|B'\|_2^2.
\]
After multiplying the squared plaquette bound by $\eta_k^{-1}$,
this is \eqref{IV-59-22}. The unperturbed representative is also
in the top background orbit, by the restriction identity used in
\cref{IV-59-11}. Undoing the level-$k$ axial gauge and the response
gauge gives the stated representatives and the conjugations $u_p$.

At $k=0$, take $\widehat U=U'$ and $\check U=V_0$.
The four-factor telescoping identity gives
$|R_p-\mathbf1|\le\sum_{b\in\partial p}|B'(b)|$.
Each bond occurs in four plaquettes, so
$\sum_p|R_p-\mathbf1|^2\le16\|B'\|_2^2$.
Since $C_{\rm nrg}\ge16(1+24C_{\rm db}^2)$, the same estimate holds.
\end{proof}

\begin{proposition}[Weighted excess]\label{IV-59-23}
For a competitor as in \cref{IV-59-17} and any bounded complex
plaquette weight $c$,
\begin{equation}\label{IV-59-24}
 \left|A^{\eta_t}(c,U')-\mathsf{ab}_t(c,\mathcal M^tU')\right|
 \le4\sup|c|\left[
   C_{\rm nrg}\gamma_{\rm r}^{-1}t\,
          \mathsf{exc}_t(U')A^{\eta_t}(U')\right]^{1/2}.
\end{equation}
\end{proposition}

\begin{proof}
For two successive backgrounds use \cref{IV-59-21} and set
$s_k=\eta_k^{-1}\sum_p|R_p-\mathbf1|^2$.
Unitary invariance and the difference of two squares give
\[
 |a_p(\widehat U)-a_p(\check U)|
 \le\tfrac12|R_p-\mathbf1|
       \bigl(|\partial\widehat U(p)-\mathbf1|
                      +|\partial\check U(p)-\mathbf1|\bigr).
\]
By Cauchy--Schwarz and the monotonicity in \cref{IV-59-17},
the difference of their weighted actions, both in level-$t$ units,
is at most
$\sup|c|\,[2L^{t-k}s_k A^{\eta_t}(U')]^{1/2}$.
Telescope these weighted actions in $k$ and apply Cauchy--Schwarz
once more. Equations \eqref{IV-59-22} and \eqref{IV-59-20} give
\[
 \sum_{k<t}L^{t-k}s_k
 \le C_{\rm nrg}\sum_{k<t}L^{t-k}\|\mathcal A_k\|_2^2
 \le8C_{\rm nrg}\gamma_{\rm r}^{-1}\mathsf{exc}_t(U').
\]
The resulting coefficient is $\sqrt{2\cdot8}=4$.
\end{proof}

\subsection{Comparison of the two background actions}

We use the holomorphic representative constants
$r^{\rm s}_{\mathrm B}$ and $K^{\rm s}_{\mathrm B}$ at the fixed
separation $\rho_{\mathrm s}$. In the notation of
Theorem~\ref{I-58-22} and its proof, these are
$r_{\mathrm G}(\rho_{\mathrm s})$ and a corresponding bound
$K_{\mathrm G}(\rho_{\mathrm s})$ for the holomorphic curvature
difference. In particular, when $\varpi(V)<\epsilon\le\bar\varepsilon$
and $B$ is real with
$\sup|B|<\min\{r_{\mathrm B},r^{\rm s}_{\mathrm B}/2\}$, the two
backgrounds exist, and their canonical representatives satisfy
\[
 |\partial U_t(e^{iB}V)(p)-\partial U_t(V)(p)|
                      \le2K^{\rm s}_{\mathrm B}\eta_t^2\sup|B|.
\]
Together with $|\partial U_t(V)(p)-\mathbf1|<B_3\epsilon\eta_t^2$,
the difference of squares and the count $3\eta_t^{-3}N_t$ of bare
plaquettes give
\begin{equation}\label{IV-59-26}
 \left|\mathsf{ab}_t(c,e^{iB}V)-\mathsf{ab}_t(c,V)\right|
 \le6K^{\rm s}_{\mathrm B}N_t\sup|c|\sup|B|
                  \bigl(B_3\epsilon+K^{\rm s}_{\mathrm B}\sup|B|\bigr).
\end{equation}
Indeed the per-plaquette difference is at most
$2K^{\rm s}_{\mathrm B}\eta_t^4\sup|B|
(B_3\epsilon+K^{\rm s}_{\mathrm B}\sup|B|)$, and the action
has the prefactor $\eta_t^{-1}$.

Consider the two cutoffs $\mathrm A=(\ell,K)$ and
$\mathrm B=(\ell,K+1)$ on the same torus, with $K\ge2$ and
\eqref{eq:iv-classical-room}. The constants and auxiliary data are
the same in both constructions. Their level-$j$ and level-$(j+1)$
point lattices coincide:
\[
 \mathbb T^{(j),\mathrm A}=\mathbb T^{(j+1),\mathrm B}
                         =\Lambda_{K-j,\ell}.
\]
At these levels the blocks, contour means, averaging map, pivot sets,
Haar measure, and gauge action agree. Their backgrounds and fibre
charts need not agree, since the bare lattices differ; we retain
the superscripts on these objects.

Fix $b\ge\gamma^{-2}$ and put
$x_j^{\mathrm A}=L^{K-j}b$ and
$x_j^{\mathrm B}=L^{K+1-j}b$. Thus
$x_{j+1}^{\mathrm B}=x_j^{\mathrm A}$.
Choose $0<c'_2\le\min\{a_{\rm plain},c^{\rm SO}_2\}$ and
$0<c_4\le c_{\rm plain}$, using the uniform curvature and field
radii in Proposition~\ref{I-27-78} and its proof.
Its proof supplies the centered estimate with coefficient two,
as well as the plain-average estimate stated there.
Define $\Theta_\star\ge\tau_\natural$ as the least calibration
for which $\epsilon=2c_\rho\varepsilon(\tau)$ obeys all of
\begin{equation}\label{IV-59-29}
 \begin{gathered}
 C_{\rm cp}\epsilon\le e_{\rm cp},\qquad
 C'_{\rm pr}B_3\epsilon\le c'_2,\qquad
 B_3\epsilon\le\varpi'_L,\\
 \tfrac\pi2 B_3\epsilon\le\tfrac14\,10^{-4},\qquad
 8\pi B_3\epsilon\le\omega_M,\\
 8\pi^2C_{\rm pr}B_3\epsilon
       \le\min\{c_4,r_{\mathrm B},r^{\rm s}_{\mathrm B}/2\},\\
 C_{\mathrm{en}}B_3\epsilon\le2/5,\qquad
 \epsilon\le\min\{1,\bar\varepsilon\}.
 \end{gathered}
\end{equation}
Each right side is positive. Since $\varepsilon$ is continuous,
decreasing, and tends to zero, such calibrations form a closed
half-line, proving that $\Theta_\star$ exists. It depends only on
the fixed constants.

Choose $1\le t\le K$ with $x_t^{\mathrm A}\ge\Theta_\star$,
put $\epsilon=2c_\rho\varepsilon(x_t^{\mathrm A})$, and take
$V\in\mathcal V_t^{\mathrm A}=\mathcal V_{t+1}^{\mathrm B}$
with $\varpi(V)<\epsilon$. Let $U^{\mathrm A}=U_t^{\mathrm A}(V)$,
$U^{\mathrm B}=U_{t+1}^{\mathrm B}(V)$, and abbreviate
\[
 \begin{gathered}
 a^{\mathrm A}=A^{\eta_t}(U^{\mathrm A}),\quad
 a^{\mathrm B}=A^{\eta_{t+1}}(U^{\mathrm B}),\quad
 N_t=(2L^{m+K-t})^3,\\
 j_t=C_{\mathrm{en}}B_3\epsilon L^{-2(t+1)},\quad
 l_t=4\pi^2C_{\rm pr}B_3\epsilon\eta_t.
 \end{gathered}
\]
Both backgrounds exist by the real-background theorem. They have
curvature less than $B_3\epsilon\eta_t^2$ and
$B_3\epsilon\eta_{t+1}^2$, respectively, so
\begin{equation}\label{eq:iv-classical-size}
 0\le a^{\mathrm A},a^{\mathrm B}
                       \le\tfrac32B_3^2\epsilon^2N_t.
\end{equation}

\begin{lemma}[The two one-sided comparisons]\label{IV-59-36}
With these data, $U'=\mathcal M U^{\mathrm B}$ is a competitor
at level $t$ in $\mathrm A$, of admitted radius $2B_3\epsilon$,
and $\mathcal M^tU'=V$. Moreover
\begin{equation}\label{IV-59-37}
 a^{\mathrm A}\le A^{\eta_t}(U')\le(1+j_t)^2a^{\mathrm B}.
\end{equation}
The reverse comparison is
\begin{equation}\label{IV-59-39}
 a^{\mathrm B}\le
   (1+4000L^4\vartheta_0)a^{\mathrm A}
                    +K_{\rm up}\epsilon^2\eta_tN_t,
 \qquad \vartheta_0=\vartheta(U^{\mathrm A}),
\end{equation}
where
\[
 K_{\rm up}=24\pi^2K^{\rm s}_{\mathrm B}C_{\rm pr}B_3^2
                         (1+4\pi^2K^{\rm s}_{\mathrm B}C_{\rm pr}).
\]
\end{lemma}

\begin{proof}
The curvature iteration of Proposition~\ref{I-10-101} applies
because $B_3\epsilon\le c^{\rm SO}_2$ and gives
\[
\varpi(\mathcal M^{k+1}U^{\mathrm B})
\le2B_3\epsilon L^{-2(t-k)}\qquad(k<t).
\]
At $k=t$ use
$\mathcal M^{t+1}U^{\mathrm B}=V$.
Since $2B_3\epsilon\le C_{\rm cp}\epsilon\le e_{\rm cp}$,
\cref{IV-59-17} proves the first inequality in \eqref{IV-59-37}.
The arbitrary-set averaging bound with all plaquettes proves the
second one, since the input curvature is at most
$B_3\epsilon L^{-2(t+1)}$.

For the other direction, the angle of $U^{\mathrm A}$ is at most
$\vartheta_0\le\frac\pi2B_3\epsilon\eta_t^2\le10^{-4}/4$.
Set $\widetilde U=\mathsf{Prl}_LU^{\mathrm A}$. By
\cref{IV-57-59},
\begin{equation}\label{eq:iv-prolonged-action}
 A^{\eta_{t+1}}(\widetilde U)
          \le(1+4000L^4\vartheta_0)a^{\mathrm A},\qquad
 \varpi(\widetilde U)\le4\pi\varpi(U^{\mathrm A}).
\end{equation}
Put $U_1=\mathcal M\widetilde U$ on the bare lattice of $\mathrm A$.
Normalize $\widetilde U$ to axial gauge by a gauge equal to one at
the coarse sites. Its straight products remain $U^{\mathrm A}(c)$.
Each of the $L-1$ internal bonds in such a product differs from
one by at most $8A_L\varpi(\widetilde U)$, and its crossing bond
differs from $U_1(c)$ by at most
$C_{\rm AC}\varpi(\widetilde U)$. The product telescoping bound
therefore gives
\begin{equation}\label{eq:iv-average-prolongation}
 |U_1(c)-U^{\mathrm A}(c)|
       \le C_{\rm pr}\varpi(\widetilde U),\qquad
 \varpi(U_1)\le C'_{\rm pr}\varpi(U^{\mathrm A}).
\end{equation}
The axial estimates apply because
$\varpi(\widetilde U)\le4\pi B_3\epsilon\eta_t^2<\omega_M$.
The second bound uses four bond differences in a plaquette and
\eqref{eq:iv-prolonged-action}.

At level zero in $\mathrm B$, \eqref{eq:iv-prolonged-action} gives
the competitor radius $4\pi L^2B_3\epsilon$. At levels $k\ge1$,
apply the curvature iteration to $U_1$ using
$C'_{\rm pr}B_3\epsilon\le c'_2\le c^{\rm SO}_2$; this gives
radius $2C'_{\rm pr}B_3\epsilon$. Thus $\widetilde U$ is a
competitor of radius $C_{\rm cp}\epsilon\le e_{\rm cp}$ at
level $t+1$ of $\mathrm B$.

The field
\[
 A(c)=\eta_t^{-1}\tfrac1i\log
                     \bigl(U_1(c)U^{\mathrm A}(c)^{-1}\bigr)
\]
is real and has supremum at most
$2\pi^2C_{\rm pr}B_3\epsilon\eta_t=l_t/2<c_4$.
The real centered averaging iteration gives
\[
 \mathcal M^tU_1=(e^{ia_t}V)^u,\qquad
                    \sup|a_t|\le2\sup|A|\le l_t,
\]
where $u$ is a real gauge on the top lattice.
This is the factorization in the proof of
Proposition~\ref{I-27-78}, also recorded in
equation~\eqref{II-155-25}; the reference curvature and field
radii are exactly those in \eqref{IV-59-29}.
Minimality for the competitor $\widetilde U$ and gauge invariance
now imply
\[
 \mathsf{ab}_{t+1}^{\mathrm B}(e^{ia_t}V)
                       \le A^{\eta_{t+1}}(\widetilde U).
\]
By \eqref{IV-59-29}, $l_t<\min\{r_{\mathrm B},r^{\rm s}_{\mathrm B}/2\}$.
The datum estimate \eqref{IV-59-26} in $\mathrm B$, with weight one,
bounds the difference between this action and $a^{\mathrm B}$ by
\[
 6K^{\rm s}_{\mathrm B}N_tl_t
           (B_3\epsilon+K^{\rm s}_{\mathrm B}l_t)
                         \le K_{\rm up}\epsilon^2\eta_tN_t.
\]
Combining this with \eqref{eq:iv-prolonged-action} proves
\eqref{IV-59-39}.
\end{proof}

Set
\begin{equation}\label{IV-59-34}
 \begin{gathered}
 K_u=K_{\rm up}+10^4L^4B_3^3+5C_{\mathrm{en}}B_3^3,\\
 K_w=6\left[2K_u+5B_3^2+
             10(C_{\rm nrg}\gamma_{\rm r}^{-1}K_uB_3^2)^{1/2}\right],
 \qquad
 \|c\|_{(t)}=\sup|c|+a_{K-t}\operatorname{Lip}_\infty(c).
 \end{gathered}
\end{equation}

\begin{proposition}[Classical action comparison]\label{IV-59-40}
For the data above, the two background actions satisfy
\begin{align}
 |a^{\mathrm B}-a^{\mathrm A}|
          &\le K_u\epsilon^2N_tL^{-t},\label{IV-59-41}\\
 |A^{\eta_{t+1}}(c,U^{\mathrm B})-A^{\eta_t}(c,U^{\mathrm A})|
          &\le K_w\|c\|_{(t)}\epsilon^2N_t\sqrt{tL^{-t}}
              \label{IV-59-42}
\end{align}
for every complex Lipschitz function $c$ on $T_\ell$.
The constants $K_u,K_w$ depend only on the fixed construction
constants, not on $K,m,b,t,V$, or $c$.
\end{proposition}

\begin{proof}
Write $q=(1+j_t)^2$. Since $j_t\le2/5$, one has
$1\le q\le2$, $q-1\le3j_t$, and
\begin{equation}\label{eq:iv-classical-q-error}
 (q-1)a^{\mathrm B}
 \le5C_{\mathrm{en}}B_3^3\epsilon^2\eta_tN_t.
\end{equation}
Here we used \eqref{eq:iv-classical-size}, $\epsilon\le1$, and
$L^{-2(t+1)}\le\eta_t$.
The lower comparison \eqref{IV-59-37} bounds
$a^{\mathrm A}-a^{\mathrm B}$ by this quantity.
In the upper comparison \eqref{IV-59-39}, the prolongation error
is at most
\[
 4000L^4\cdot\tfrac\pi2B_3\epsilon\eta_t^2
                 \cdot\tfrac32B_3^2\epsilon^2N_t
              \le10^4L^4B_3^3\epsilon^2\eta_tN_t.
\]
This proves \eqref{IV-59-41}.

For the weighted bound first suppose $0\le c\le c_{\max}$ and put
$L_c=a_{K-t}\operatorname{Lip}_\infty(c)$,
$v=\frac32\eta_tL_c$. Apply \cref{lem:iv-weighted-averaging}
to $U^{\mathrm B}$ with physical coarse spacing $a_K$.
Its curvature satisfies the required bound, and its factor is at
most $q$. Thus, for $X=A^{\eta_{t+1}}(c,U^{\mathrm B})$ and
$Y=A^{\eta_t}(c,U')$,
\[
 Y\le q(X+va^{\mathrm B}).
\]
Apply the same bound to $c_{\max}-c$, and use
$A^{\eta_t}(U')\ge a^{\mathrm A}$. Since
$0\le X\le c_{\max}a^{\mathrm B}$, the two inequalities give
\[
 -(q-1)c_{\max}a^{\mathrm B}-2va^{\mathrm B}
 \le X-Y\le c_{\max}(a^{\mathrm B}-a^{\mathrm A})
                     +(q-1)c_{\max}a^{\mathrm B}+2va^{\mathrm B}.
\]
Equations \eqref{IV-59-41}, \eqref{eq:iv-classical-q-error}, and
\eqref{eq:iv-classical-size} yield
\begin{equation}\label{eq:iv-weighted-first-difference}
 |X-Y|\le(2K_uc_{\max}+\tfrac92B_3^2L_c)
                                     \epsilon^2\eta_tN_t.
\end{equation}
For the competitor $U'$, the same inequalities give
\[
 \mathsf{exc}_t(U')\le2K_u\epsilon^2\eta_tN_t,\qquad
 A^{\eta_t}(U')\le3B_3^2\epsilon^2N_t.
\]
Apply \cref{IV-59-23} to obtain
\[
 |Y-A^{\eta_t}(c,U^{\mathrm A})|
 \le10c_{\max}(C_{\rm nrg}\gamma_{\rm r}^{-1}K_uB_3^2)^{1/2}
                             \epsilon^2N_t\sqrt{t\eta_t}.
\]
Since $\eta_t\le\sqrt{t\eta_t}$, summing this with
\eqref{eq:iv-weighted-first-difference} bounds the desired difference
by $K_w(c_{\max}+L_c)\epsilon^2N_t\sqrt{t\eta_t}/6$.

For a real signed weight, apply this result to $c+\sup|c|$ and
subtract the constant weight using \eqref{IV-59-41}.
Since $K_u\le K_w/6$, the resulting bound is at most
$K_w\|c\|_{(t)}\epsilon^2N_t\sqrt{t\eta_t}/2$.
Apply this separately to the real and imaginary parts of a complex
weight and add their bounds. This proves \eqref{IV-59-42}.
\end{proof}

\subsection{Field restrictions and comparison of the fibre origins}

Let $1\le t_*\le K-1$, with $x_{t_*}\ge\tau_\natural$, and put
\[
 e_k^{\rm c}=c_\rho\varepsilon(x_k),\qquad
 \mathsf{In}_k=\{V:\varpi(V)<2\varepsilon(x_k)\}.
\]
At level zero set $\mathsf{Cut}_0=\{V:\varpi(V)<e_0^{\rm c}\}$.
For $1\le k\le t_*$ let $\mathsf{Cut}_k$ consist of the fields
whose background exists and satisfies
$\mathsf m_k(U_k(V))<e_k^{\rm c}$, where
$\mathsf m_k(U)=\eta_k^{-2}\max_p|\partial U(p)-\mathbf1|$.
This number is independent of the representative of the critical orbit.

\begin{lemma}\label{IV-58-15}
These sets are open and gauge invariant. For $0\le k\le t_*$,
\begin{equation}\label{IV-58-16}
 \mathsf{In}_k\subset\mathsf{Cut}_k
 \subset\{\varpi<2e_k^{\rm c}\}
 \subset\{\varpi<\varepsilon_{\mathcal W}\alpha_0\}.
\end{equation}
For $1\le k\le t_*$, every field in $\mathsf{Cut}_k$ is in the
actual global step domain at every real parameter $\tau\ge\tau_\natural$.
If $k<t_*$, $W\in\mathsf{Cut}_{k+1}$, and
$|B(b)|<\mathfrak b(x_k)$ on the nonpivot bonds, then
\begin{equation}\label{IV-58-17}
 V_W(B)\in\mathsf{Cut}_k,\qquad
 \varpi(V_k^{\rm bf}(W))<2L^{-2}e_{k+1}^{\rm c}.
\end{equation}
For $1\le k\le t_*$, a background over $\mathsf{Cut}_k$ is
regular at level $k$ with radius $\alpha_0$, and
\begin{equation}\label{eq:iv-cut-action}
 |A^{\eta_k}(c,U_k(V))|
       \le\tfrac32(e_k^{\rm c})^2N_k\sup|c|.
\end{equation}
\end{lemma}

\begin{proof}
The choice of $\tau_\natural$ gives
$2e_k^{\rm c}\le\min(\varepsilon_{\mathcal W}\alpha_0,
c_2^{\rm SO},\varepsilon_v/4)$.
The real background theorem, Theorem~\ref{I-58-22}, gives
$\mathsf m_k(U_k(V))\le B_3\varpi(V)$, so
$c_\rho\ge2B_3$ proves the first inclusion.
The small-curvature averaging estimate in the proof of
Proposition~\ref{I-27-78} gives
$\varpi(\mathcal M^kU)<2\mathsf m_k(U)$ in this range.
This proves the other inclusions, including the level-zero case.
Proposition~\ref{II-42-39} now supplies the actual global domain.
It also gives continuity of the background in local logarithmic
coordinates. Thus the strict curvature condition is open inside
$\{\varpi<\bar\varepsilon\}$, which contains it by the preceding
inclusions. Gauge covariance of the background and unitary invariance
of the plaquette norm prove gauge invariance.

The restriction identity in that proposition identifies the orbit
over $V_0=V_k^{\rm bf}(W)$ with the axial representatives of the
orbit over $W$. Its curvature therefore has
$\mathsf m_k<L^{-2}e_{k+1}^{\rm c}$, and averaging proves the
second assertion of \eqref{IV-58-17}. On the real cutoff chart
$V_W(B)=e^{iB'}V_0$ one has
$\sup|B'|<2C_{\rm db}\delta(x_k)$.
The representative curvature-difference bound preceding
\eqref{IV-59-26}, with the last condition in
\eqref{eq:iv-classical-fixed-calibration}, gives
\[
 \mathsf m_k(U_k(V_W(B)))
 \le L^{-2}\mathsf m_{k+1}(U_{k+1}(W))
                        +4C_{\rm db}K^{\rm s}_{\mathrm B}\delta(x_k)
 \qquad(k\ge1).
\]
For $k=0$ the elementary plaquette product bound gives the same
inequality with $8C_{\rm db}$ in place of
$4C_{\rm db}K^{\rm s}_{\mathrm B}$. Both are bounded by $C_H$.
Finally
\[
 L^{-2}e_{k+1}^{\rm c}+C_H\delta(x_k)
 \le\left(c_\rho L^{-3/2}+C_H A_1/A_0\right)
                                  \varepsilon(x_k)
 \le e_k^{\rm c}.
\]
The first term before this estimate is strict, proving downward
stability of the cut. The regularity assertion follows from
the plain-core membership in Proposition~\ref{II-42-39}:
both the scaled curvature and current are less than $\alpha_0$.
There are $3\eta_k^{-3}N_k$ bare plaquettes, and each has
$a_p(U)<\frac12(e_k^{\rm c})^2\eta_k^4$.
Multiplication by $\eta_k^{-1}\sup|c|$ gives
\eqref{eq:iv-cut-action}.
\end{proof}

\begin{proposition}\label{IV-59-45}
Use the two partner lattices above, with top levels $t_*$ and
$t_*+1$, and assume $x^{\mathrm A}_{t_*}\ge\Theta_\star$.
Let $1\le k<t_*$ and
$W\in\mathsf{Cut}^{\mathrm A}_{k+1}
       \cap\mathsf{Cut}^{\mathrm B}_{k+2}$.
Then, on every common nonpivot bond,
\begin{equation}\label{IV-59-46}
 \varphi\left(V_{k+1}^{\mathrm{bf},\mathrm B}(W)(b)
                    V_k^{\mathrm{bf},\mathrm A}(W)(b)^{-1}\right)
 \le8c_\rho\varepsilon(x_k)
          \sqrt{K_u\gamma_{\rm r}^{-1}L^{-k}N_k}.
\end{equation}
\end{proposition}

\begin{proof}
Write $V^{\mathrm A},V^{\mathrm B}$ for the two axial backgrounds
and $\epsilon_j=2c_\rho\varepsilon(x_j)$.
Both average to $W$ and, by \cref{IV-58-15}, have curvature
less than $L^{-2}\epsilon_{k+1}\le L^{-3/2}\epsilon_k$.
The axial comparison \eqref{IV-59-8} therefore bounds their
relative intrinsic logarithm by
$h'=\pi C_{\rm ax}L^{-3/2}\epsilon_k$.
The first inequality in \eqref{IV-59-29}, together with
$C_{\rm cp}\ge4\pi L^2$,
puts $h'$ below the chart and cubic-remainder thresholds in
\cref{IV-59-11}. Thus $V^{\mathrm B}$ lies in the chart over
$W$ based at $V^{\mathrm A}$. If $A$ is its nonpivot logarithm,
that proposition yields
\[
 \tfrac18\gamma_{\rm r}\|A\|_2^2
 \le\mathsf{ab}^{\mathrm A}_k(V^{\mathrm B})
                       -L^{-1}\mathsf{ab}^{\mathrm A}_{k+1}(W).
\]
The restriction identity on lattice $\mathrm B$ gives
$\mathsf{ab}^{\mathrm B}_{k+1}(V^{\mathrm B})
=L^{-1}\mathsf{ab}^{\mathrm B}_{k+2}(W)$.
Apply \eqref{IV-59-41} at level $k$ to $V^{\mathrm B}$
and at level $k+1$ to $W$. These applications are valid since
their curvatures are below $\epsilon_k,\epsilon_{k+1}$.
The right side is at most
\[
 K_u\left(\epsilon_k^2L^{-k}N_k
       +L^{-1}\epsilon_{k+1}^2L^{-(k+1)}N_{k+1}\right)
 \le2K_u\epsilon_k^2L^{-k}N_k,
\]
where $N_{k+1}=L^{-3}N_k$ and
$\epsilon_{k+1}^2\le L\epsilon_k^2$.
Consequently $\|A\|_2\le
4\epsilon_k\sqrt{K_u\gamma_{\rm r}^{-1}L^{-k}N_k}$.
Each bond norm is bounded by this square-sum norm, and the
intrinsic logarithm has norm equal to the angle. This is
\eqref{IV-59-46}; reversing a bond leaves that angle unchanged.
\end{proof}

\subsection{Older interactions at a regular background}

Write $K_{\rm A}=K^\sharp_{\rm tw}$ and $a_{\rm A}=a^\sharp$
for the simple-domain constants of Proposition~\ref{I-93-98},
and $B_*,c_*$ for the regional-gauge constants of
Theorem~\ref{I-28-75}, with $c_*\le1$.
Assume $M\alpha_0\le\alpha_1$ and
$\alpha_0+\alpha_1\le a_{\rm A}$. Set
\begin{equation}\label{eq:iv-old-term-constants}
 \begin{gathered}
 c_{\rm A}=\min\{c_*/4,a_{\rm A}/(64K_{\rm A}B_*)\},\qquad
 C_{\rm A}^{\rm old}=2(64K_{\rm A}B_*/a_{\rm A})^2,\\
 C'_{\rm A}=6e^{\kappa_{\mathcal D}}M^2C_{\rm A}^{\rm old},\qquad
 \rho_{\rm A}=a_{\rm A}\alpha_0/(16K_{\rm A}),\\
 l_0=\lceil\log_L(M/c_{\rm A})\rceil,\qquad
 C_{\rm lg}=\sup_{l\in\mathbb N_0}L^{4l}e^{-c_{\rm A}L^l},\\
 C_{\rm old}=C'_{\rm A}
                  +2C_{\mathrm E,0}(C_{\rm lg}+L^{4l_0}).
 \end{gathered}
\end{equation}
Here $C_{\mathrm E,0}=2048\sum_{n\ge0}e^{-n}$ is a permissible
rooted-domain entropy constant. In particular $l_0\ge1$ and
$C_{\rm lg}\le(4/(ec_{\rm A}))^4$.

\begin{lemma}\label{lem:iv-old-term-decay}
Let $1\le i\le t\le K$, $16M\le L^{m+K-t}$, and let $U$
be regular at level $t$ with radius $\alpha_0$ on the torus.
Suppose its pairs belong to the fixed cores on all domains
used below. If $D^{(i)}$ is a pointed family at radii
$(\alpha_0,\alpha_1)$ and $l=t-i$, then
\[
 \sum_{z\in\mathbb T^{(i)}}\sum_{X\in\mathcal D_i:X\ni z}
 |D^{(i)}(X,U,z)-D^{(i)}(X,\mathbf1,z)|
 \le C_{\rm old}N_tL^{-l}\|D^{(i)}\|_{i,\kappa+1}.
\]
\end{lemma}

\begin{proof}
We first prove the local quadratic gain that supplies the factor
$L^{-4l}$. Put $\vartheta=L^{-l}$. For a domain of size
$d=d_i(X)\le c_{\rm A}\vartheta^{-1}$, assume
$M\alpha_0\vartheta\le c_{\rm A}\alpha_0$.
The same bare configuration is regular at level $i$ with radius
$\alpha=\alpha_0\vartheta^2$: curvature rescales by
$\vartheta^2$, and current by $\vartheta^3\le\vartheta^2$.
One has $(d+1)\vartheta\le1$. Indeed this is immediate for
$l=0$, when $d=0$, while for $l\ge1$ its left side is at
most $c_{\rm A}+L^{-1}<1$.

A connected edge set of size $d$ meeting all cubes of $X$
puts $X$ in the cube of radius $(d+1)M$ about one of its
corners. Enlarge this cube by two unit layers. Its inner side
is $q=2(d+1)M$, and
\[
 q+4\le6M\vartheta^{-1}<16M\vartheta^{-1}
                  \le L^{m+K-i}.
\]
Thus it lies in a torus chart. Also
$2q\alpha\le4c_{\rm A}\alpha_0\le c_*$.
The regional gauge theorem, applied with radius $2\alpha$,
gives on the inner cube a gauge in which
$U=e^{i\eta_i A}$ and
$|A|,|\nabla^{\eta_i}A|<\beta=4B_*(d+1)M\alpha$.
Its divergence hypothesis follows directly from
\[
 |D_U^{\mathsf T}\partial U|
 \le\eta_i^{-1}(\alpha\eta_i^2)^2
                        +\eta_i^2|J^{(i)}(U)|
 <2\alpha\eta_i^2.
\]
The gauge-transformed current $J_1$ also has norm less than
$\beta$, and $\beta\le\rho_{\rm A}$.

Consider the complex disc of pairs
$\Gamma(\zeta)=(e^{i\zeta\eta_i A},\zeta J_1)$ for
$|\zeta|\le R=\rho_{\rm A}/\beta$.
With $\phi=|\zeta|\beta\le\rho_{\rm A}\le1/16$, the
exponential-product estimate gives
\[
 |\partial(e^{i\zeta\eta_i A})-\mathbf1|
 \le(2\phi+11\phi^2)\eta_i^2<4\rho_{\rm A}\eta_i^2.
\]
The other simple-domain bounds have radii
$(4\rho_{\rm A},2\rho_{\rm A},4\rho_{\rm A})$.
Moreover
\[
 6\rho_{\rm A}\le a_{\rm A},\quad
 4\rho_{\rm A}\le\alpha_0,\quad
 2\rho_{\rm A}\le\alpha_1,\quad
 K_{\rm A}(4\rho_{\rm A}+4\rho_{\rm A}^2)\le\alpha_0.
\]
The core inclusion \eqref{I-93-100} therefore puts the
whole disc in the fixed core, with strict room in the defining
simple-domain bounds. The same holds after undoing the real
gauge. Pair dependence and gauge invariance apply there.

For a local term $D$, the function $f(\zeta)=D(\Gamma(\zeta))$
is holomorphic near the closed disc and bounded by $\|D\|_X$.
Its derivative at zero vanishes: the four constant conjugations
by $1,i\sigma_1,i\sigma_2,i\sigma_3$ preserve $D$, and their
adjoint actions sum to zero on each copy of the complexified
Lie algebra. Applying the linear differential of $D$ at the
trivial pair gives $4f'(0)=0$. The maximum principle applied
to $(f(\zeta)-f(0))/\zeta^2$ now gives
\[
 |D(U)-D(\mathbf1)|\le2\|D\|_XR^{-2}
 =C_{\rm A}^{\rm old}\|D\|_X
                   M^2(d+1)^2\vartheta^4.
\]
This is the fixed-radius argument also used in the proof of
Lemma~\ref{II-101-27}; here the displayed radius choices
give the constants in \eqref{eq:iv-old-term-constants}.

Let $H_i=\|D^{(i)}\|_{i,\kappa+1}$. At each site the norm
bound is $\|D^{(i)}(X,\cdot,z)\|_X\le H_i e^{-(\kappa+1)d}$.
There are at most $e^{\kappa_{\mathcal D}(n+1)}$ domains of
size $n$ containing the site. Since
$\kappa+1-\kappa_{\mathcal D}>1$,
\[
 \sum_{X\ni z}(d_i(X)+1)^2e^{-(\kappa+1)d_i(X)}
 \le e^{\kappa_{\mathcal D}}\sum_{n\ge0}(n+1)^2e^{-n}
 <6e^{\kappa_{\mathcal D}}.
\]
If $l\ge l_0$, the preceding quadratic bound applies to
every domain with $d_i(X)\le c_{\rm A}L^l$.
Summing it over the $L^{3l}$ sites of each level-$t$ block
and then its $N_t$ blocks gives a contribution at most
$C'_{\rm A}N_tL^{-l}H_i$.

For the remaining domains use simply
$|D^{(i)}(X,U,z)-D^{(i)}(X,\mathbf1,z)|
\le2H_i e^{-(\kappa+1)d_i(X)}$.
The rooted entropy estimate, with $\kappa+1>\kappa_{\mathcal D}+2$,
gives, for any $D\ge0$,
\[
 \sum_z\sum_{X\ni z:\ d_i(X)\ge D}
 |D^{(i)}(X,U,z)-D^{(i)}(X,\mathbf1,z)|
 \le2C_{\mathrm E,0}e^{-D}L^{3l}N_tH_i.
\]
For $l\ge l_0$ choose $D=c_{\rm A}L^l$ and use
$L^{3l}e^{-c_{\rm A}L^l}\le C_{\rm lg}L^{-l}$.
For $l<l_0$, choose $D=0$ and use
$L^{3l}\le L^{4l_0}L^{-l}$. Adding the bounds gives the
asserted constant $C_{\rm old}$.
\end{proof}

\begin{proposition}\label{IV-60-22}
For $1\le t\le t_*$ and $V\in\mathsf{Cut}_t$, let
$U=U_t(V)$ and let
$\mathfrak f=((D^{(i)},0)_{1\le i\le t},\emptyset)$ be
a bare array at the fixed radii. Then its action is defined
on the whole critical orbit, is constant there, and
\begin{equation}\label{IV-60-23}
 |S[\mathfrak f](U)|
 \le C_{\rm old}N_t\sum_{i=1}^t
                   L^{-(t-i)}\|D^{(i)}\|_{i,\kappa+1}.
\end{equation}
\end{proposition}

\begin{proof}
Plain-core membership of the level-$t$ background, its tower
restrictions, and radius monotonicity put every level-$i$ pair
in the fixed core. The defining sum for the action is therefore
\[
 S[\mathfrak f](U)=\sum_{i=1}^t\sum_z\sum_{X\ni z}
       [D^{(i)}(X,U,z)-D^{(i)}(X,\mathbf1,z)].
\]
Each summand is gauge invariant. The regularity statement of
\cref{IV-58-15} and the room assumption permit
\cref{lem:iv-old-term-decay} at every $i$. Sum its estimates.
\end{proof}

\section{Intermediate integration on one lattice}
\label{sec:iv-intermediate-integration}

We use the same real backgrounds and Gaussian structures for all
sources. Only the fourth potential is changed. This permits a
comparison of two choices of the earlier interactions before the
last integrations are performed.

At height $k\ge1$ put
\[
 r_k^{\rm src}=\min\{\varepsilon_{\mathfrak h},
                         c_{\rm os}\alpha_0/(2c_{\mathfrak h})\},
\]
and set $r_0^{\rm src}=r_P^\circ$ at height zero.
Require $C_{\mathcal A}\delta(x_k)\le r_k^{\rm src}/2$.
For $k\ge1$ define the localized weighted action by
\[
 \mathbb V_v^{\rm act}(Y;P,A)
   =\sum_{X\subset Y}\mathbb V[W_v(X,\cdot)](Y;P,A),\qquad
 W_v(X,U)=\eta_k^{-1}\sum_{p:X(p)=X}v(p)a_p(U),
\]
using the actual shift and vertex brackets of
Proposition~\ref{II-50-61}. At $k=0$ use instead
\[
 \mathbb V_v^{\rm act}(Y;P,A)
   =\sum_{p:X(p)=Y}v(p)
             [a_p(e^{iB'(P,A)}\mathbf U)-a_p(\mathbf U)],
\]
where $B'$ is the pivot completion. These are the same weighted
Wilson members and shifts used in the source calculation above.
All sums here are taken on their specified full one-lattice cores.

For $g=x_k^{-1/2}$, subtract the linear reference value at each
nonpivot bond and set
\begin{equation}\label{eq:iv-direct-source-family}
 \mathbb W^{\rm src}[v](Y;P,B)
 =x_k\mathbb V_v^{\rm act}(Y;P,gB)
       -x_k\sum_bv(p_b)D_A\mathbb V_1^{\rm act}(Y;P,0)[gB_b].
\end{equation}
Here $B_b=B(b)\delta_b$, and $p_b$ is the fixed reference
plaquette at the centre of the fine cube containing $b_-$.
Its corners lie in that cube. The reference subtraction is made
on every core, before specializing to a critical field.

Let $K_{\rm Q}^\vee,K_{\rm L}^\vee$ be at least the maxima of
\begin{equation}\label{eq:iv-direct-source-coefficients}
 \begin{array}{c|cc}
 &K_{\rm Q}&K_{\rm L}\\ \hline
 k\ge1&672M^6(\alpha_0/r_k^{\rm src})^2
       &8064M^6\alpha_0/r_k^{\rm src}\\
 k=0&2b_{\rm A}M^3(r_P^\circ)^{-2}
       &72b_{\rm A}M^3(r_P^\circ)^{-1}
 \end{array}
\end{equation}
with $b_{\rm A}=6\cdot1201^2M^3$.
For a region $R$ put
\[
 \mathsf{sz}_{k,R}(v)
 =K_{\rm Q}^\vee
     (1+2C_{\mathcal A}\delta(x_k)/r_k^{\rm src})|v|_R
                    +K_{\rm L}^\vee x_k^{1/2}\mathrm{sw}_k(v).
\]
The oscillation $\mathrm{sw}_k$ is measured between plaquettes
whose barycentres are at distance at most $M$ in level-$k$ units.

\begin{lemma}\label{lem:iv-direct-source-bound}
The source family \eqref{eq:iv-direct-source-family} is linear
in $v$, holomorphic in the source, the open triple slices, and
the field polydisc. It vanishes at $B=0$, is gauge invariant,
depends on the background pair, and is local in the plaquettes
of $Y$. Its quadratic-growth norm satisfies
\begin{equation}\label{eq:iv-direct-source-bound}
 \mathsf n_X^{\rm f}(\mathbb W^{\rm src}[v])
 \le\mathsf{sz}_{k,X}(v),\qquad
 \mathsf n^{\rm f}(\mathbb W^{\rm src}[v])
 \le\mathsf{sz}_{k,T_\ell}(v)
 \le2K_{\rm Q}^\vee|v|_\infty
                  +K_{\rm L}^\vee x_k^{1/2}\mathrm{sw}_k(v).
\end{equation}
The norm uses the rate $\lambda_V\le\bar\kappa-2$ and
$1+q_Y(B)$, with $q_Y(B)=M^{-3}\sum_b|B(b)|^2$.
On a real fibre over $W\in\mathsf{Cut}_{k+1}$,
\begin{equation}\label{eq:iv-direct-source-sum}
 \sum_Y\mathbb W^{\rm src}[v](Y;P_W|_Y,B|_Y)
 =x_k\{A^{\eta_k}(v,U_k(V_W(B)))
                         -A^{\eta_k}(v,U_{k+1}(W))\}.
\end{equation}
\end{lemma}

\begin{proof}
For $k\ge1$, each nonzero root $X=X(p)$ has at most four
fine cubes and size zero. The determinant identity for a
complex unimodular plaquette matrix and the core curvature
bound give $\|W_v(X,\cdot)\|_X\le6M^3\alpha_0^2|v|_Y$.
Its amplitude radius in Proposition~\ref{II-50-61} is
$c_{\rm os}\alpha_0/(c_{\mathfrak h}r_k^{\rm src})\ge2$.
If $N=\#_{\pi_k}Y$, there are at most $7N$ roots, and the
number of nonroot cubes is at least $\max(d_k(Y),N-4)$.
The sharper first bound in that proposition, together with
$e^{\bar\kappa}-1\ge e^{\bar\kappa-1/2}$, therefore gives
\[
 e^{r'd_k(Y)}N\sup_{|A|<r_k^{\rm src}}
       |\mathbb V_v^{\rm act}(Y;P,A)|
 \le42N^2e^{-(N-4)_+/2}M^3\alpha_0^2|v|_Y
 \le672M^3\alpha_0^2|v|_Y
\]
for $0\le r'\le\bar\kappa-1$.
The last inequality uses $N^2e^{-(N-4)_+/2}\le16$.
At $k=0$, the pivot plaquette calculation in
Proposition~\ref{II-41-79} and its proof gives instead
$\sup|\mathbb V_v^{\rm act}|\le2b_{\rm A}|v|_Y$;
the term is zero unless $d_0(Y)=0$ and $N\le4$.
The same calculation is valid with a weight inside the finite
plaquette sum, since it bounds the sum of the plaquette moduli.

Write $F_Y(v)$ for these bounds before multiplying by
$e^{r'd}N$. Cauchy's formula gives homogeneous coefficients
$|c_j[v](B)|\le F_Y(v)(\sup|B|/r_k^{\rm src})^j$.
Taylor expansion of \eqref{eq:iv-direct-source-family} is
\[
 g^{-1}\sum_bc_1[v-v(p_b)](B_b)
           +c_2[v](B)+x_k\sum_{j\ge3}c_j[v](gB).
\]
For the quadratic term use $\sup|B|^2\le M^3q_Y(B)$.
For the remainder use the half-radius bound
\[
 x_k\left|\sum_{j\ge3}c_j[v](gB)\right|
 \le2F_Y(v)(r_k^{\rm src})^{-2}\sup|B|^2
                      C_{\mathcal A}\delta(x_k)/r_k^{\rm src}.
\]
For the linear term use
$\sum_b|B(b)|\le\frac32M^3N(1+q_Y(B))$.
A cube path from a plaquette to $p_b$ has at most $d+2$
links, so $|v-v(p_b)|_Y\le(d+2)\mathrm{sw}_k(v)$;
$e^{-d}(d+2)\le2$ absorbs one unit of decay.
For $k\ge1$ the resulting coefficients are
$672M^6(\alpha_0/r_k^{\rm src})^2$ for the quadratic term
and at most $2016M^6\alpha_0/r_k^{\rm src}$ for the linear
term. For $k=0$ they are respectively
$2b_{\rm A}M^3(r_P^\circ)^{-2}$ and at most
$24b_{\rm A}M^3(r_P^\circ)^{-1}$.
Thus \eqref{eq:iv-direct-source-coefficients} bounds all three
terms and proves \eqref{eq:iv-direct-source-bound}.

The local shift and pivot formulas, their finite vertex brackets,
and Cauchy's formula give the stated parameter properties.
The reference subtraction uses only $v(p_b)$ with $p_b\subset Y$,
so it preserves locality. The actual-shift identity of
Proposition~\ref{II-50-61}, summed over the plaquette roots,
gives the weighted action difference. The subtracted linear
terms have sum zero bond by bond at $P_W$: they are the
derivative of the unweighted action on its averaging fibre,
which vanishes at the critical background. This is precisely
the criticality calculation of \cref{III-19-8}, including the
first-step pivot case. Hence \eqref{eq:iv-direct-source-sum}
holds for every weight, not just constant weights.
\end{proof}

\subsection{The recursion and its density}

In addition to the global hypotheses above, assume
\begin{gather}
 E^{\rm s}_0\ge2E_{\mathsf N},\qquad
 N^{\rm s}_\star\ge\frac L{L-1}
                (E^{\rm s}_0+\tfrac32\alpha_0^2B_\sharp),
                                                        \label{eq:iv-flow-size}\\
 K_{\mathbf M}\left(\delta(\gamma^{-2})+
          \mathsf s_\gamma N^{\rm s}_\star+
          e^{-\gamma_2\mathfrak b(\gamma^{-2})^2/2}\right)
       \le E^{\rm s}_0/4,\qquad
 8K_\Lambda\mathsf s_\gamma N^{\rm s}_\star\le w_\star,
 \qquad8K_P\mathsf s_\gamma\le w_\star .
                                                        \label{eq:iv-flow-smallness}
\end{gather}
Here $\mathsf s_\gamma=(1+\beta_{\rm s})
\delta(\gamma^{-2})/\alpha_0$. Require also
$4\beta_{\rm s}\kappa\ge1$, and put
\[
 K^{\rm d}=\tfrac{351}{2}L^3K_\Pi,\qquad
 K_{\mathbf M}^{\rm f}=3L^3\max\{K_E e^{4\lambda_1},
                         32K_\Pi K_P,\tfrac{117}{2}K_\Pi\},
 \quad
 \mu_\star=\min\{w_\star/32,E^{\rm s}_0/(16K_{\mathbf M}^{\rm f})\}.
\]
Thus $K^{\rm d}\le K_{\mathbf M}^{\rm f}$.
Set $\sigma_{\rm src}=\mu_\star/K_{\rm Q}^\vee$ and choose
$\varrho>0$ with
\begin{equation}\label{eq:iv-flow-source-room}
 2\varrho\le\sigma_{\rm src}(\mathsf b'-\bar B),\qquad
 2K_{\rm L}^\vee\varrho\gamma\le40\mu_\star/\sqrt3.
\end{equation}
The sources belong to
$\mathsf{Src}=\mathsf{Src}_{\mathsf{Bc}}(2\varrho,2\varrho)$.
For the line $x_j=L^{K-j}\mathsf b'$ write
\[
 x_j^w=x_j-L^{K-j}w,\qquad v_j^w=-w/\mathsf b',\qquad
 \mathsf s_j=(1+\beta_{\rm s})\delta(x_j)/\alpha_0.
\]
The source bounds give
$|v_j^w|_\infty\le\sigma_{\rm src}$ and
$x_j^{1/2}\mathrm{sw}_j(v_j^w)\le\mu_\star/K_{\rm L}^\vee$.
In particular, \cref{lem:iv-direct-source-bound} gives
$\mathsf n^{\rm f}(\mathbb W^{\rm src}[v_j^w])\le3\mu_\star$.

An early system at $t_0$ consists of pointed families
$E^{(i)}[w]$, $1\le i\le t_0$, on the fixed pair cores,
with $\|E^{(i)}[w]\|_{i,\kappa+1}\le E^{\rm s}_0$.
They are jointly holomorphic in $w$ and every open pair slice,
and their member at $X$ depends on $w$ only through the
plaquettes in $X$. For $j\ge t_0$ define
\begin{equation}\label{eq:iv-flow-recursion}
 \mathfrak f_j[w]=((E^{(i)}[w],0)_{1\le i\le j},\emptyset),
 \qquad
 E^{(j+1)}[w]=\mathbf N^{(j+1)}+
       \mathbf F[\mathbb V'[\mathfrak f_j[w]]
                            -\mathbb W^{\rm src}[v_j^w]].
\end{equation}
All objects in a step use its real Gaussian structure at
$x=\mathsf t=x_j$, so that $\tau=1$. The chain starts at
$t_0=0$ with the empty array; a classical start at $t_0$
sets all its early families equal to zero. A hat denotes
the latter flow.

\begin{proposition}\label{IV-60-12}
Every early system continues through level $K$ by
\eqref{eq:iv-flow-recursion}. Its new families have the same
norm, holomorphy and source-locality properties as the early
families. At each step the fourth family has norm at most
$7w_\star/32<w_\star/4$.

For $t_0\le k\le t_*$ and $V\in\mathsf{Cut}_k$ put
\[
 \Phi_k^w(V)=-A^{\eta_k}(x_k^w,U_k(V))
                              +S[\mathfrak f_k[w]](U_k(V)),
 \qquad F_k^w=\mathbf1_{\mathsf{Cut}_k}e^{\Phi_k^w}.
\]
These are bounded Borel functions. For $t_0\le k<t_*$ set
\begin{equation}\label{eq:iv-contact-increment}
 c_{k+1}(w)=\log\{\mathsf{pf}_k(x_k)
                              \mathsf Z^{(k)}[\mathbf1]\}
       +\sum_{X,z}E^{(k+1)}[w](X,\mathbf1,z),
\end{equation}
where the logarithm is real. For every
$W\in\mathsf{Cut}_{k+1}$,
\begin{equation}\label{IV-60-13}
 \mathsf I[F_k^w](W)=e^{c_{k+1}(w)+\Phi_{k+1}^w(W)}\ne0.
\end{equation}
Consequently, the finite complex measures
\[
 \nu_{t_0}^w=F_{t_0}^w\,dV,\qquad
 \nu_{k+1}^w=\mathbf1_{\mathsf{Cut}_{k+1}}
       \mathcal M_*[\mathfrak G_k[x_k]\chi_k^{\rm sf}\nu_k^w]
\]
satisfy
\begin{equation}\label{IV-60-14}
 \nu_k^w=\exp\left(\sum_{t_0<i\le k}c_i(w)\right)
                    \mathbf1_{\mathsf{Cut}_k}e^{\Phi_k^w}\,dV.
\end{equation}
The chain restarted with its own first $t_0$ families is
the same chain.
\end{proposition}

\begin{proof}
Suppose the families through $j$ have been constructed.
Their weighted array size is at most
$L E^{\rm s}_0/(L-1)\le N^{\rm s}_\star$.
The actual old-potential rule of Theorem~\ref{II-127-29}
therefore has norm at most
$K_\Lambda\mathsf s_jN^{\rm s}_\star\le w_\star/8$.
The source contributes at most $3\mu_\star\le3w_\star/32$.
The three internal potentials have total norm at most
$3K_P\delta(x_j)\le3w_\star/8$.
Thus \cref{III-34-16} applies on the full core.
Its unperturbed bound, the tent mass, and the normalization
bound of Theorem~\ref{II-127-29} give respectively
\[
 \|\mathbf F[\mathbb V'[\mathfrak f_j]]\|_{j+1,\kappa+1}
 \le K_{\mathbf M}\left(\delta(x_j)+\mathsf s_jN^{\rm s}_\star
                  +e^{-\gamma_2\mathfrak b(x_j)^2/2}\right)
 \le E^{\rm s}_0/4,
 \qquad
 \|\mathbf N^{(j+1)}\|_{j+1,\kappa+1}\le E^{\rm s}_0/2.
\]
Here the output rates exceed $\kappa+1$ by
$4\beta_{\rm s}\kappa\ge1$ and the normalizer decay
assumption. In \eqref{III-34-24}, the coefficient is at most
$162L^3K_\Pi<K^{\rm d}$ whenever
$n_{\rm int}\le3w_\star/8$, $n_{\rm iv}\le w_\star/4$
and $n_W\le w_\star/2$. The source perturbation therefore
costs at most $3K^{\rm d}\mu_\star\le3E^{\rm s}_0/16$.
The new norm is at most $15E^{\rm s}_0/16$.
The parameter statement of \cref{III-34-16}, the linear
old-potential rule and \cref{lem:iv-direct-source-bound}
give holomorphy. If two sources agree on $X$, every earlier
member used inside $X$ agrees, as does the source family;
the same local construction then gives the same new member.
This proves the induction, including the first step, whose
old array is empty.

For the density identity fix $W\in\mathsf{Cut}_{k+1}$.
The cut theorem puts $W$ in the actual global sumrule class.
The old-potential sumrule and \eqref{eq:iv-direct-source-sum}
identify the fourth-family sum with the change along the
fibre of
\[
 \mathcal S_k(U)=S[\mathfrak f_k[w]](U)
                              -x_kA^{\eta_k}(v_k^w,U).
\]
The three internal sumrules and the identity
$\mathbf C_*(P_W)=\mathsf A_W^{-1}$ identify the chart
integral with its truncated Gaussian integral. Its cluster
sum is the exponential of the sum of the pointed terms
by \eqref{III-34-23}. The Gaussian normalization, localized
by $\mathbf N^{(k+1)}$, contributes the first logarithm in
\eqref{eq:iv-contact-increment}; this is the actual
normalization identity of Theorem~\ref{II-127-29}, with
the hypotheses verified in Proposition~\ref{II-42-39}.
Hence the chart value is
\[
 \exp\left\{\log(\mathsf{pf}_k\mathsf Z^{(k)}[\mathbf1])
       +\sum_{X,z}E^{(k+1)}[w](X,U_+,z)
       -A^{\eta_k}(x_k^w,U_+)+S[\mathfrak f_k[w]](U_+)\right\}.
\]
Subtract and add the trivial-pair values in this finite sum.
The old members and the new centered member are precisely
the action of $\mathfrak f_{k+1}[w]$; the weighted action
rescales because $x_k^w/L=x_{k+1}^w$. This proves
\eqref{IV-60-13}, including its nonvanishing.

All these actions are defined on the critical orbit by the
cut theorem, and are constant there by gauge invariance.
Local analytic representatives of that orbit show that
$\Phi_k^w$ is Borel on the cut. Its absolute value is bounded
there by the cut action estimate and \cref{IV-60-22}.
Extending its exponential by zero gives the asserted bounded
Borel function. The physical image-density formula applies
to it. Every chart over $\mathsf{Cut}_{k+1}$ lies in
$\mathsf{Cut}_k$, so the inserted cut is identically one
in the chart. Induction with \eqref{IV-60-13} gives
\eqref{IV-60-14}. Equality upon restarting follows directly
from the deterministic recursion.
\end{proof}

\subsection{Decay of the early interactions}

Put $c_{\log}=\log L/\log\gamma^{-2}$ and
\[
 \Sigma_\phi=\sum_{a\ge1}L^{-a/2}(1+ac_{\log})^{\mathfrak p_0},
 \qquad C_{\rm fg}=\exp\{\tfrac{13}{72}(L-1)\Sigma_\phi\}.
\]
The series converges. The explicit logarithmic cutoff gives
\begin{equation}\label{eq:iv-flow-scale-decay}
 \mathsf s_j\le\mathsf s_\gamma
 L^{-(K-j)/2}(1+(K-j)c_{\log})^{\mathfrak p_0}.
\end{equation}
Indeed $x_j=L^{K-j}\mathsf b'$, $\mathsf b'\ge\gamma^{-2}$,
and $x^{-1/2}(\log x)^{\mathfrak p_0}$ decreases on this
range. Also $K_{\mathbf M}\ge243L^3K_\Pi K_\Lambda$ and
\eqref{eq:iv-flow-size}--\eqref{eq:iv-flow-smallness} imply
\begin{equation}\label{eq:iv-flow-contraction-constant}
 K^{\rm d}K_\Lambda\mathsf s_\gamma\frac L{L-1}\le13/72.
\end{equation}

\begin{proposition}\label{IV-60-16}
For two early systems at $t_0$ and the same source, let
\[
 \Delta_i=\|E^{(i)}[w]-\widetilde E^{(i)}[w]\|_{i,\kappa+1},
 \qquad n_j=\sum_{i\le j}L^{-(j-i)}\Delta_i.
\]
For $j\ge t_0$,
\begin{equation}\label{IV-60-18}
 n_j\le C_{\rm fg}L^{-(j-t_0)}n_{t_0}
 \le C_{\rm fg}L^{-(j-t_0)}\frac L{L-1}\max_{i\le t_0}\Delta_i.
\end{equation}
In particular, the chain and classical start obey this
bound with $\max\Delta_i\le E^{\rm s}_0$.
\end{proposition}

\begin{proof}
The normalizers and source families cancel in the difference
of the two recursions. Linearity of the old-potential rule
and its weighted-lag estimate give
$\mathsf n^{\rm f}(\mathbb V'[\mathfrak f_j-
\widetilde{\mathfrak f}_j])\le K_\Lambda\mathsf s_j n_j$.
Both fourth families have norm below $w_\star/4$, so their
difference has norm below $w_\star/2$.
Equation \eqref{III-34-24} therefore gives
$\Delta_{j+1}\le K^{\rm d}K_\Lambda\mathsf s_jn_j$.
Consequently
\[
 n_{j+1}\le(L^{-1}+K^{\rm d}K_\Lambda\mathsf s_j)n_j,
 \qquad
 n_j\le L^{-(j-t_0)}n_{t_0}
          \exp\left(LK^{\rm d}K_\Lambda
                          \sum_{i=t_0}^{j-1}\mathsf s_i\right).
\]
Equations \eqref{eq:iv-flow-scale-decay} and
\eqref{eq:iv-flow-contraction-constant} bound the last
exponential by $C_{\rm fg}$. The geometric sum bounds
$n_{t_0}$, proving the result. If $t_0=0$ the flows coincide.
\end{proof}

Define
\begin{equation}\label{eq:iv-flow-exponent-errors}
 \epsilon_{\rm fg}(t,t_0)=C_{\rm old}C_{\rm fg}
       \frac L{L-1}E^{\rm s}_0N_tL^{-(t-t_0)},\qquad
 \mathsf m_t^S=C_{\rm old}\frac L{L-1}E^{\rm s}_0N_t.
\end{equation}
Combining \cref{IV-60-16,IV-60-22}, for $t_0\le t\le t_*$
and $V\in\mathsf{Cut}_t$ we have
\begin{equation}\label{IV-60-26}
 |\Phi_t^w(V)-\widehat\Phi_t^w(V)|\le\epsilon_{\rm fg}(t,t_0),
 \qquad |S[\mathfrak f_t[w]](U_t(V))|\le\mathsf m_t^S.
\end{equation}
The same action bound holds for every early system. The
weighted Wilson actions cancel in the first difference;
this estimate concerns the field-dependent exponents only.

\subsection{A uniform fluctuation tail}

Set
\begin{equation}\label{IV-60-28}
 \kappa_{\rm b}=\frac{w_\star}{32M^3}e^{-\lambda_V},\qquad
 C_{\rm sh}=\tfrac14K^{\rm d}w_\star M^3C_{\mathrm E,1},
 \quad C_{\mathrm E,1}=2048\sum_{n\ge0}(n+2)e^{-n}.
\end{equation}
These constants are independent of $\gamma$ and of the
cutoff, torus, level and source.

\begin{proposition}\label{IV-60-29}
Let $t_0\le k<t_*$, let $b$ be a nonpivot bond at level $k$,
and let $g=x_k^{-1/2}$. Let $\Sigma\in\mathbb C$ be a scalar
such that $e^\Sigma F_k^w$ is real and nonnegative on
$\mathsf{Cut}_k$. Then, for
$W\in\mathsf{Cut}_{k+1}$ and $a\ge0$,
\begin{equation}\label{IV-60-30}
 0\le\mathsf I[e^\Sigma F_k^w
       \mathbf1\{|\mathcal A_k^\diamond(b)|\ge ga\}](W)
 \le e^{-\kappa_{\rm b}a^2+C_{\rm sh}}
                                      \mathsf I[e^\Sigma F_k^w](W),
 \qquad\mathsf I[e^\Sigma F_k^w](W)>0.
\end{equation}
The indicator is set to zero where the intrinsic fluctuation
is undefined; with this convention it is Borel.
\end{proposition}

\begin{proof}
On the stable background domain the map
$(V,W)\mapsto V(b)V_k^{\rm bf}(W)(b)^{-1}$ is continuous.
The intrinsic logarithm is continuous away from $-\mathbf1$.
Taking its norm superlevel set and composing with the Borel
map $V\mapsto(V,\mathcal M V)$ proves the Borel assertion.

Let $\mathbb Q_b(Y;P,B)=\tfrac12\operatorname{tr}(B(b)^2)$
if $Y=\square(b)$, and zero otherwise. This is holomorphic,
gauge invariant and independent of $P$. Since
$d_k(\square(b))=0$ and $|B(b)|^2\le M^3q_{\square(b)}(B)$,
\[
 \mathsf n_X^{\rm f}(\kappa_{\rm b}\mathbb Q_b)\le w_\star/32,
 \qquad
 \mathsf n_X^{\rm f}(\kappa_{\rm b}\mathbb Q_b)=0
                 \quad\hbox{unless }\square(b)^\sharp\subset X.
\]
Adding $\lambda\mathbb Q_b$, for $\lambda=0,\kappa_{\rm b}$,
leaves the fourth-family norm at most $w_\star/4$.
The Gaussian sumrule used in \cref{IV-60-12} now gives
\begin{equation}\label{eq:iv-bond-tilt-identity}
 \mathsf I[F_k^w e^{\lambda x_k|\mathcal A_k^\diamond(b)|^2}](W)
       =p_W\Xi_\lambda,\qquad
 \Xi_\lambda=\exp\sum_{X,z}(\mathbf G_1+\mathbf F_\lambda)(X,P_W|_X,z).
\end{equation}
The value assigned to the extra exponential off its domain
is immaterial. Indeed every chart lies inside the cut, where
$\mathcal A_k^\diamond(b)=gB(b)$, so the insertion is
$e^{\lambda|B(b)|^2}$. In the Gaussian calculation
\[
 p_W=\mathsf{pf}_k(x_k)Z^{(k)}(W;\mathsf A_W)
        \exp[-x_k\mathbb F_W(0)+\mathcal S_k(U_{k+1}(W))]\ne0
\]
is independent of $\lambda$. The three internal sums are
$\mathbb P^J,\mathbb P^G,\mathbb P^A[\mathsf B_W]$,
the fourth sum is $\mathsf s_W[\mathcal S_k]+\lambda|B(b)|^2$,
and the covariance is $\mathsf A_W^{-1}$. At $\tau=1$ the
quadratic correction potential vanishes. These identities
give \eqref{eq:iv-bond-tilt-identity} directly from the
physical Gaussian change of variables.

The local form of \eqref{III-34-24} gives
\[
 |(\mathbf F_{\kappa_{\rm b}}-\mathbf F_0)(X,P,z)|
 \le K^{\rm d}w_\star e^{-(\kappa+1)d_{k+1}(X)}/32,
\]
and the difference vanishes unless $\square(b)^\sharp\subset X$.
All contributing domains contain one fixed level-$(k+1)$
site $z_b$: choose it in the partition cube containing
$b_-$. Each such domain has at most
$8M^3(d_{k+1}(X)+1)$ base points. The rooted entropy bound
therefore gives
\[
 \left|\log(\Xi_{\kappa_{\rm b}}/\Xi_0)\right|
 \le\frac{K^{\rm d}w_\star}{32}
       \sum_{X\ni z_b}8M^3(d_{k+1}(X)+1)e^{-(\kappa+1)d_{k+1}(X)}
 \le C_{\rm sh},
\]
where the logarithm denotes the displayed cluster-sum
difference. Thus $|\Xi_{\kappa_{\rm b}}/\Xi_0|\le e^{C_{\rm sh}}$.

The chart measure is positive, and its fields lie in
$\mathsf{Cut}_k$. Apply
$\mathbf1\{|B(b)|\ge a\}\le
 e^{-\kappa_{\rm b}a^2}e^{\kappa_{\rm b}|B(b)|^2}$
under that measure with the nonnegative factor
$e^\Sigma F_k^w$. The denominator is nonzero by
\eqref{eq:iv-bond-tilt-identity} and nonnegative by the
same chart formula, hence positive. The quotient of the
tilted and untilted integrals is real and nonnegative,
and its modulus is at most $e^{C_{\rm sh}}$.
This proves \eqref{IV-60-30}.
\end{proof}


\subsection{Source radii and field-independent terms}

Let $f_1,\ldots,f_n$ be nonzero real Lipschitz functions
on $T_\ell$ and $w_z=\sum_i z_if_i$. Write
\[
 \begin{gathered}
 D^\flat_R=\{z:\sum_i|z_i|\|f_i\|_{\flat;T_\ell}<R\},\qquad
 r_i(t)=(8nM_t^{\rm rd}\|f_i\|_{\flat;T_\ell})^{-1},\\
 M_t^{\rm rd}=\max\{\varrho^{-1},
       \tfrac32(e_t^{\rm c})^2L^{K-t}N_t,
       \varrho^{-1}\mathsf m_t^S\}.
 \end{gathered}
\]
The polydisc $D_t^{\rm rd}=\{|z_i|<r_i(t)\}$ lies in
$D^\flat_{\varrho/8}$. Recall that
$0<\|f_i\|_\flat\le\|f_i\|_\sharp$ and
$\|w_z\|_\flat\le\sum_i|z_i|\|f_i\|_\flat$.

\begin{lemma}\label{IV-60-40}
For any flow beginning at $t_0<t$, $1\le t\le t_*$,
and $V\in\mathsf{Cut}_t$, the function
$z\mapsto\Phi_t^{w_z}(V)$ is holomorphic on
$D^\flat_{2\varrho}$, and
\begin{equation}\label{eq:iv-readout-quarter}
 |\Phi_t^{w_z}(V)-\Phi_t^0(V)|\le1/4
                    \qquad(z\in D_t^{\rm rd}).
\end{equation}
\end{lemma}

\begin{proof}
Fix a critical representative $U=U_t(V)$, independent of $z$.
The finite sum $\Psi(w)=S[\mathfrak f_t[w]](U)$ is holomorphic
and bounded by $\mathsf m_t^S$ on $\mathsf{Src}$.
The exponent difference is
$L^{K-t}A^{\eta_t}(w_z,U)+\Psi(w_z)-\Psi(0)$.
The cut action bound makes the first term at most
$M_t^{\rm rd}\|w_z\|_\flat<1/8$.
Schwarz's lemma on the complex line through $w_z$, whose
source radius is $2\varrho/\|w_z\|_\flat$, bounds the second
term by $\mathsf m_t^S\|w_z\|_\flat/\varrho<1/8$.
Their sum proves the assertion.
\end{proof}

For $n\ge2$ define the physical support spread
\[
 d_\infty(\underline f)=\min_{x_i\in\operatorname{supp}f_i}
                                      \max_{i,i'}|x_i-x_{i'}|_\infty.
\]
Compactness shows that the minimum is attained and that it
is positive exactly when the supports have no common point.
For families with the norm and locality properties of the
flows, put
\[
 c_*(z)=\sum_{j=j_0}^{j_1}\sum_{X,z'}E_*^{(j)}[w_z](X,\mathbf1,z'),
 \qquad1\le j_0\le j_1\le K.
\]

\begin{lemma}\label{IV-60-46}
The function $c_*$ is holomorphic on $D^\flat_{2\varrho}$
and $|c_*|\le2E^{\rm s}_0C_{\mathrm E,0}N_{j_0}$ there.
If $n\ge2$ and $d_\infty L^{K-j_1}\ge6M$, then
\begin{equation}\label{IV-60-48}
 |\partial_1\cdots\partial_nc_*(0)|\le\epsilon_{\rm mx}(j_1),
 \qquad
 \epsilon_{\rm mx}(j_1)=2e^3E^{\rm s}_0C_{\mathrm E,0}N_{j_1}
       e^{-d_\infty L^{K-j_1}/M}
              \prod_i\frac{n\|f_i\|_\flat}{2\varrho}.
\end{equation}
\end{lemma}

\begin{proof}
Each member is bounded by $E^{\rm s}_0e^{-(\kappa+1)d_j(X)}$.
Sum first over domains containing a fixed site and then over
levels, using $N_j=L^{-3(j-j_0)}N_{j_0}$, to obtain the
first bound. All sums are finite and holomorphic.

Apply Cauchy's formula on the polydisc with radii
$2\varrho/(n\|f_i\|_\flat)$. A differentiated member
vanishes unless, for every $i$, $X$ contains a plaquette at
whose barycentre $f_i$ is nonzero: otherwise locality makes it independent
of that coefficient. Choose one such plaquette for each $i$.
Their corners are in $X$, are within $\eta_j/2$ of their
barycentres, and have mutual sup distance at most
$M(d_j(X)+2)$ in level-$j$ units. Hence a contributing
domain satisfies
\[
 d_j(X)\ge d_\infty L^{K-j}/M-3=:D_j\ge3.
\]
The rooted tail and Cauchy's formula bound its level sum by
\[
 T_j=E^{\rm s}_0C_{\mathrm E,0}N_j
                   e^{3-d_\infty L^{K-j}/M}
                       \prod_i\frac{n\|f_i\|_\flat}{2\varrho}.
\]
Since $T_{j-1}/T_j\le L^3e^{-6(L-1)}\le1/2$,
the sum is at most $2T_{j_1}$, proving \eqref{IV-60-48}.
\end{proof}

\section{Discarding the large-field histories}
\label{sec:iv-small-field-discarding}

We now use the calibrated constants and canonical auxiliary
data of Theorem~\ref{II-203-9} together with the source
conditions above. Fix $0<\eth<1$ before decreasing $\gamma$,
and set
\[
 c'=(1-\eth)c_\star,\quad
 C_{\rm maj}^\Sigma=C_{\rm maj}+(6L^3+1)\log2,\quad
 E^{\rm s}_0=E_0,\quad
 N^{\rm s}_\star=\frac L{L-1}(E_0+\tfrac32\alpha_0^2B_\sharp),
\]
\[
 \varrho=\min\left\{\frac{20\mu_\star}{\sqrt3K_{\rm L}^\vee\gamma},
              \frac{\sigma_{\rm src}}2(\gamma^{-2}-\bar B)\right\}>0.
\]
Require $\log\gamma^{-2}$ to exceed the case-II threshold
of Proposition~\ref{II-75-111} at $c=c_\star,\eth$ and
the summability threshold $\max\{Y_1,Y^{\rm S}(c')\}$
preceding \eqref{II-75-72}. All these choices precede
$K,m,\mathsf b'$. The data are
\[
 K\ge2,\qquad \mathsf b'\ge\gamma^{-2},\qquad
 MR(\mathsf b')\le L^m,\qquad \ell=2L^m,
 \qquad1\le t_*\le K-1,\quad x_{t_*}\ge\tau_\natural.
\]
The calibration also includes the global radius, localization
and smallness conditions used in the preceding sections.

\begin{lemma}\label{lem:iv-early-history-tail}
For the real constructed history masses and $1\le t\le K$,
\begin{equation}\label{eq:iv-early-history-tail}
 \sum_{h\in\mathcal H_K:\,Z_t(h)\ne\emptyset}
               e^{-E_{(K)}}\mathsf t(h)
 \le2M^{-3}N_t e^{-c'p_0(g_t)^2}
                         \Pi_K(C_{\rm maj}^\Sigma,c').
\end{equation}
\end{lemma}

\begin{proof}
The actual masses satisfy the summation input
\eqref{eq:history-summation-input}, \eqref{eq:actual-vacuum-summation}. In the proof of
Proposition~\ref{II-75-106}, they are bounded, after
multiplication by $e^{-C_{\rm maj}N_K}$, by
\[
 \prod_{j=1}^K(2^{-n_j^{\rm bit}}\zeta_j)^{\#W_j},
 \qquad \zeta_j=e^{-c'p_0(g_j)^2-2p_0(g_j)},
\]
where $W_j$ is the exact set of created cells. The entropy
bound \eqref{eq:history-entropy} cancels the binary factors even
when the sum is restricted to prescribed creation sets.
The region $Z_t$ is nonempty exactly when some $W_j$ with
$j\le t$ is nonempty: creations persist in later regions,
and tracing any nonempty region backwards through the hull
recursion reaches a creation.
Use the union bound over these possible marked cells in
the finite subset product. The restricted sum is at most
\[
 \Pi_K(C_{\rm maj},c')
                    \sum_{j\le t}\#\mathcal P_j e^{-c'p_0(g_j)^2}.
\]
The proof of \eqref{II-75-72}, stopped at $t$, gives
\[
 \sum_{j\le t}\#\mathcal P_j e^{-c'p_0(g_j)^2}
 \le M^{-3}N_t\sum_{j\le t}L^{3(t-j)}e^{-c'p_0(g_j)^2}
 \le2M^{-3}N_te^{-c'p_0(g_t)^2}.
\]
Increase the first argument of $\Pi_K$ to
$C_{\rm maj}^\Sigma$ to obtain \eqref{eq:iv-early-history-tail}.
\end{proof}

Write $V_j=\mathcal M^jU$ and define the bare-field weight
\[
 \mathfrak w_t^{\rm sf}(U)=\mathbf1_{\mathsf{In}_t}(V_t)
     \prod_{j=0}^{t-1}
       [\mathbf1_{\mathsf{Cut}_j}\mathfrak G_j[x_j]
                                        \chi_j^{\rm sf}](V_j).
\]
It is bounded, nonnegative and Borel. Put
\begin{equation}\label{eq:iv-discard-error}
 R_t=2M^{-3}N_t\exp\left\{-c'A_0^2(\log x_t)^{2\mathfrak p_0}
               +(C_{\rm maj}^\Sigma+2+C_{\rm low}(\mathsf b'))\ell^3\right\}.
\end{equation}
Here $C_{\rm low}(x)=c_{\rm low}+(9/2)\log x$ is the
same lower-bound constant as in \cref{III-7-9}.

\begin{proposition}\label{IV-61-19}
For $1\le t\le t_*$ and every admitted source,
\begin{equation}\label{IV-61-20}
 (\mathcal M^t)_*(\mathfrak w_t^{\rm sf}\rho_0^w\,dU)
                         =\mathbf1_{\mathsf{In}_t}\nu_t^w.
\end{equation}
For every bounded, nonnegative, gauge-invariant $F$,
\begin{equation}\label{IV-61-21}
 \int F\mathfrak w_{h_\emptyset^{(t)}}\rho_0\,dU
 \le\int F\mathfrak w_t^{\rm sf}\rho_0\,dU
 \le\int F\rho_0\,dU.
\end{equation}
Consequently, with $Z=\int\rho_0\,dU$,
\begin{equation}\label{IV-61-22}
 0\le Z-\int\mathfrak w_t^{\rm sf}\rho_0\,dU\le R_tZ.
\end{equation}
Let $\mathsf{op}_n=\sum_\pi(\#\pi)!$, the sum over
partitions of $\{1,\ldots,n\}$. If $R_t\le1/2$ and
$T(z)=\int\mathfrak w_t^{\rm sf}\rho_0^{w_z}\,dU$, then
$T$ is entire, $T(0)\ge Z/2$, and
\begin{equation}\label{IV-61-23}
 \left|S^T_{n;K,\ell,L^K\mathsf b'}(\underline f)
                     -\partial_1\cdots\partial_n\log T(0)\right|
 \le4\mathsf{op}_nR_t\prod_i(6\ell^3L^{4K}\|f_i\|_\infty).
\end{equation}
\end{proposition}

\begin{proof}
The image identity follows by induction from the recursion
for $\nu_j^w$ and the rule for pushing forward a measure
multiplied by a function of its image. At level zero the
cut density is $\mathbf1_{\mathsf{Cut}_0}\rho_0^w$;
at the final level use $\mathsf{In}_t\subset\mathsf{Cut}_t$.

The empty-history weight is the displayed product for
$\mathfrak w_t^{\rm sf}$ with three additional zero-one
density and compatibility tests at each step. If that
weight is nonzero, the density tests imply
$\varpi(V_j)<2\varepsilon(x_j)$ for $j\ge1$ and
$\varpi(V_0)<\varepsilon(x_0)$. Hence all its fields satisfy
the cuts and its last field satisfies $\mathsf{In}_t$.
The two weights then coincide. This proves the first
inequality of \eqref{IV-61-21} pointwise.

For the second inequality first remove all cuts and
fluctuation tests. Let
\[
 J_j=\int F\rho_0\prod_{k<j}\mathfrak G_k[x_k](\mathcal M^kU)\,dU.
\]
Gauge transformations on the noncentral sites of the
level-$(j-1)$ blocks lift to transformations constant on
the corresponding bare blocks. They preserve $F\rho_0$
and every preceding gauge factor. Average the last gauge
factor over their product Haar measure. Its defining
normalization makes the average one on its averaging
domain and zero off that domain. Thus $J_j\le J_{j-1}$.
Fubini applies to bounded nonnegative Borel functions on
finite compact products. Iteration gives $J_t\le J_0$,
as asserted; no pointwise bound $\mathfrak w_t^{\rm sf}\le1$
is needed.

The history masses add to $Z$, and the descendants of the
empty history through level $t$ add to its mass. Those are
exactly the histories with $Z_t=\emptyset$. Apply
\cref{lem:iv-early-history-tail},
$\Pi_K(C_{\rm maj}^\Sigma,c')\le
 e^{(C_{\rm maj}^\Sigma+2)\ell^3}$, and the lower bound
$Z\ge e^{E_{(K)}-C_{\rm low}(\mathsf b')\ell^3}$.
The vacuum constants cancel, proving \eqref{IV-61-22}.

On the algebra of bounded gauge-invariant Borel functions,
the two functionals defined by $\rho_0$ and
$\mathfrak w_t^{\rm sf}\rho_0$ are positive, and their
difference is positive by \eqref{IV-61-21}. Therefore their
difference on $F$ has modulus at most $R_tZ\|F\|_\infty$.
Their normalized moments differ by at most
$4R_t\|F\|_\infty$ when $R_t\le1/2$.
For the bounded observables
$Y_i=\mathcal O(f_i)$ one has
$|Y_i|\le6\ell^3L^{4K}\|f_i\|_\infty$.
In each partition term of the cumulant formula, telescope
the product of normalized moments; there are $\#\pi$
differences, and every remaining factor is bounded by
the product of its observable bounds. Multiplication by
$(\#\pi-1)!$ and summation proves \eqref{IV-61-23}.
Finally $\rho_0^{w_z}=\rho_0e^{\sum_i z_iY_i}$ with bounded
$Y_i$, so dominated differentiation on compact coefficient
sets proves entireness and identifies the logarithmic
derivative with this cumulant. Equation \eqref{IV-61-20}
gives $T(z)=\nu_t^{w_z}(\mathsf{In}_t)$ on the source class.
\end{proof}


\section{Comparison of adjacent cutoffs}
\label{sec:iv-adjacent-cutoffs}

Fix $\ell$ and a label $\mathsf b$ in the calibrated range.
The lattices $\mathrm A=(\ell,K)$ and $\mathrm B=(\ell,K+1)$
use the lines $x_j^{\mathrm A}=L^{K-j}\mathsf b$ and
$x_j^{\mathrm B}=L^{K+1-j}\mathsf b$. At equal depth $d$
their configuration lattices, Haar measures, averaging maps,
gauge factors and fluctuation thresholds coincide; their
critical backgrounds need not coincide. Write
\[
 x_{[d]}=L^d\mathsf b,\qquad N_{[d]}=\ell^3L^{3d},\qquad
 y_{[d]}=(\log x_{[d]})^{2\mathfrak p_0}.
\]
Choose integers $1\le u<s<K$, put $t_0=K-s$, $t_1=K-u$,
and require $x_{[u]}\ge\Theta_\star$.
The partner of a level $k$ of $\mathrm A$ is $k+1$ of
$\mathrm B$; a subscript $(k)$ refers to this common depth.

On each lattice run the chain and a classical start at depth
$s$ to depth $u$. On the common final configuration space
$\mathcal X=\mathcal V(\Lambda_{u,\ell})$ set
\[
 I^\bullet(z)=\int_{\mathsf{In}_{[u]}}e^{\Phi^{\bullet,w_z}}\,dV,
 \quad \widehat I^\bullet(z)=\int_{\mathsf{In}_{[u]}}
                                      e^{\widehat\Phi^{\bullet,w_z}}\,dV,
 \quad T^\bullet=e^{c^\bullet}I^\bullet,\quad
 \widehat T^\bullet=e^{\widehat c^\bullet}\widehat I^\bullet.
\]
Here $c^\bullet$ sums all the chain increments through the
final level, while $\widehat c^\bullet$ sums the increments
strictly after the classical start. These identities are
\eqref{IV-60-14}. Define the contact function
\[
 \mathfrak c^\bullet(z)=
  \sum_{j=1}^{t_1^\bullet}\sum_{X,z'}E^{(j),\bullet}[w_z](X,\mathbf1,z')
  -\sum_{j=t_0^\bullet+1}^{t_1^\bullet}\sum_{X,z'}
                         \widehat E^{(j),\bullet}[w_z](X,\mathbf1,z').
\]
Then $c^\bullet-\widehat c^\bullet-\mathfrak c^\bullet$
is independent of $z$: it is the sum of the real Gaussian
normalization logarithms up to the starting level.

Use the common radii $r_i=r_i(t_1)$ of \cref{IV-60-40}, and put
\begin{align}
 \epsilon_{\rm fg}&=C_{\rm old}C_{\rm fg}\tfrac L{L-1}E_0
                       N_{[u]}L^{-(s-u)},\label{eq:iv-adj-forgetting}\\
 \epsilon_{\rm cl}&=4K_w c_\rho^2A_0^2
       (1+2\varrho/\mathsf b)y_{[K]}N_{[s]}\sqrt K
                                      L^{-(K-s)/2},\label{eq:iv-adj-classical}\\
 \epsilon_{\rm sh}&=3(s-u)N_{[s]}
         \exp\{C_{\rm sh}-\kappa_{\rm b}A_1^2y_{[u]}/4\},\label{eq:iv-adj-shell}\\
 \epsilon_3&=4\epsilon_{\rm cl}+6\epsilon_{\rm sh},\qquad
 M_c=4E_0C_{\mathrm E,0}N_{[s]}.\label{eq:iv-adj-contact-size}
\end{align}
In this section impose
\begin{equation}\label{eq:iv-adj-smallness}
 \epsilon_{\rm fg}\le1/8,\quad
 \epsilon_{\rm cl},\epsilon_{\rm sh}\le1/20,\quad
 256K_uc_\rho^2A_0^2N_{[s]}L^{-(K-s)}\le\gamma_{\rm r}A_1^2.
\end{equation}

\begin{proposition}\label{IV-62-42}
At every real point $x$ of the polydisc $|z_i|<r_i$,
$\widehat T^{\mathrm A}(x),\widehat T^{\mathrm B}(x)>0$ and
\begin{equation}\label{IV-62-44}
 \left|\widehat T^{\mathrm B}(x)/\widehat T^{\mathrm A}(x)-1\right|
                               \le\epsilon_3\le1/2.
\end{equation}
\end{proposition}

\begin{proof}
On $\mathcal Y=\mathcal V(\Lambda_{s,\ell})$ write
$W_k=\mathcal M^{k-t_0}W$, and let
\[
 k^\bullet(W)=\mathbf1_{\mathsf{In}_{[u]}}(W_{t_1})
     \prod_{k=t_0}^{t_1-1}
      [\mathbf1_{\mathsf{Cut}^\bullet_{(k)}}\mathfrak G_{(k)}
                                                \chi^\bullet_{(k)}](W_k),
\]
\[
 \widehat F_z^\bullet(W)=\mathbf1_{\mathsf{Cut}^\bullet_{(t_0)}}(W)
        \exp[-A^{\eta_{t_0^\bullet}}
              (L^s(\mathsf b-w_z),U_{t_0^\bullet}^\bullet(W))].
\]
Iterating the image formula gives
$\widehat T^\bullet(z)=\int k^\bullet\widehat F_z^\bullet\,dW$.
For real $x$ the start is strictly positive on its cut.
The positive chart kernel propagates nonnegativity, while
the nonzero exponential density formula makes the final
density positive on its cut. Since the common inner set is
open and contains $\mathbf1$, it has positive Haar measure.
Thus both final integrals are strictly positive.

We first bound the mass where the two weights differ.
Suppose $k^\bullet(W)>0$ and $k^\circ(W)=0$. Starting from
$W_{t_1}\in\mathsf{In}_{[u]}$, cut stability shows that
there is a last index $k$ with $\chi^\circ_{(k)}(W_k)=0$,
and its next field lies in both next-level cuts. Otherwise
all the other cuts would follow downwards from their
fluctuation tests, and $k^\circ$ would be positive.
By \cref{IV-59-45} and \eqref{eq:iv-adj-smallness}, at
each nonpivot bond the two background origins differ by
an angle at most
\[
 8c_\rho\varepsilon(x_{(k)})
       (K_u\gamma_{\rm r}^{-1}L^{-k}N_k)^{1/2}
                                      \le\delta(x_{(k)})/2.
\]
If all fluctuations of lattice $\bullet$ were smaller than
$\delta(x_{(k)})/2$, the angular triangle inequality would
make every fluctuation of lattice $\circ$ defined and
smaller than $\delta(x_{(k)})<\pi$. Its test would equal
one, a contradiction. Thus some bond has a fluctuation at
least half the threshold.

Insert the indicator of this event at its step in the
positive classical-start measure. By \cref{IV-60-29} with
$a=\mathfrak b(x_{(k)})/2$, the next measure is bounded by
$\exp(C_{\rm sh}-\kappa_{\rm b}\mathfrak b(x_{(k)})^2/4)$
times the ordinary next measure. All later positive
cut-and-image operations preserve this inequality.
There are $s-u$ steps, at most $3N_{[s]}$ nonpivot bonds
at each step, and
$\mathfrak b(x_{(k)})^2\ge A_1^2y_{[u]}$.
The union bound consequently gives both estimates
\begin{equation}\label{eq:iv-mismatch-mass}
 \int k^\bullet\mathbf1_{\{k^\circ=0\}}
                         \widehat F_x^\bullet\,dW
              \le\epsilon_{\rm sh}\widehat T^\bullet(x).
\end{equation}
This argument uses the image identity also for the marked
measure, obtained by the same induction as the unmarked
one; hence it bounds the final mass, without an assumption
that the later gauge factors are at most one.

Where both weights are positive they equal the product of
the common gauge factors. The two starts there satisfy
$\widehat F_x^{\mathrm B}=\widehat F_x^{\mathrm A}e^{q_x}$.
Apply \cref{IV-59-40} at $t=t_0$ with
$\varepsilon_{\rm r}=2c_\rho\varepsilon(x_{[s]})$ and
$c_x=L^s(\mathsf b-w_x)$. Since
$\|w_x\|_\flat<\varrho/8$,
$\|c_x\|_{(t_0)}\le L^s(\mathsf b+2\varrho)$.
The comparison estimate gives
\[
 |q_x|\le K_w\|c_x\|_{(t_0)}\varepsilon_{\rm r}^2
             N_{[s]}\sqrt{t_0}L^{-t_0/2}\le\epsilon_{\rm cl}.
\]
Split the two integrals into their common positive set and
their two exceptional sets. The latter have masses bounded
by \eqref{eq:iv-mismatch-mass}; on the common set use
$|e^{q_x}-1|\le2\epsilon_{\rm cl}$. First,
\[
 (1-\epsilon_{\rm sh})\widehat T^{\mathrm B}(x)
             \le e^{\epsilon_{\rm cl}}\widehat T^{\mathrm A}(x),
 \qquad \widehat T^{\mathrm B}(x)\le2\widehat T^{\mathrm A}(x).
\]
Then the total variation of the two weighted densities is
at most $(2\epsilon_{\rm cl}+3\epsilon_{\rm sh})
\widehat T^{\mathrm A}(x)$. This implies the weaker bound
\eqref{IV-62-44}, with the stated constants.
\end{proof}

The real estimate does not by itself control derivatives
in complex sources. The following elementary two-constants
argument supplies the required passage.

\begin{lemma}\label{IV-57-34}
If $f$ is holomorphic on the polydisc $|z_i|<r_i$, bounded
there by $B$, and bounded at its real points by
$0<\epsilon\le B$, then, for $0<\sigma\le1$,
\begin{equation}\label{IV-57-38}
 |\partial_1\cdots\partial_nf(0)|
 \le\epsilon^{1-\sigma}B^\sigma
                   (4n/(\pi\sigma))^n\prod_i r_i^{-1}.
\end{equation}
\end{lemma}

\begin{proof}
First take a one-variable unit disc. In its upper half put
$h(\zeta)=\frac2{\pi i}\operatorname{Log}\frac{1+\zeta}{1-\zeta}$.
Its real part lies between zero and one, is zero on the
diameter and tends to one on the circular arc. For $\nu>0$
apply the maximum principle to
\[
 f(\zeta)\epsilon^{-1}
       e^{-\log(B/\epsilon)h(\zeta)}
                          ((1-\zeta^2)/2)^\nu.
\]
The boundary upper limits of its modulus are at most one;
the last factor makes those at $\pm1$ zero. Let $\nu\downarrow0$.
For $|\zeta|\le\tau<1$,
\[
 \Re h(\zeta)=\frac2\pi\arctan
                   \frac{2\Im\zeta}{1-|\zeta|^2}
             \le\delta_\tau:=\frac4\pi\arctan\tau.
\]
Reflection gives the same lower-half bound, so
$|f|\le\epsilon^{1-\delta_\tau}B^{\delta_\tau}$ on the
smaller disc. Iterating in the $n$ variables gives
$|f|\le\epsilon^{(1-\delta_\tau)^n}
B^{1-(1-\delta_\tau)^n}$ on $|z_i|\le\tau r_i$.
Choose $\tau=\pi\sigma/(4n)$.
Then $\delta_\tau\le\sigma/n$ and
$(1-\delta_\tau)^n\ge1-\sigma$.
Cauchy's formula on the smaller polydisc proves the result.
The maximum principle with boundary upper limits follows
by taking a maximizing sequence in the bounded domain:
a limit on the boundary obeys the assumed bound, whereas
an interior maximum forces constancy.
\end{proof}

\begin{proposition}\label{IV-57-41}
On the common source polydisc the two chain functions are
zero-free. Moreover, for $0<\sigma\le1$,
\begin{equation}\label{eq:iv-readout-comparison}
 \begin{split}
 &\left|\partial^n\log T^{\mathrm B}(0)-\partial^n\log T^{\mathrm A}(0)
          -\partial^n(\mathfrak c^{\mathrm B}-\mathfrak c^{\mathrm A})(0)\right|\\
 &\quad\le\left[8\epsilon_{\rm fg}
       +2^{n+1}(4n/(\pi\sigma))^n
                   \epsilon_3^{1-\sigma}(1+2M_c)^\sigma\right]
                                                \prod_i r_i^{-1}.
 \end{split}
\end{equation}
Here $\partial^n=\partial_1\cdots\partial_n$.
\end{proposition}

\begin{proof}
The integrals are holomorphic by dominated differentiation:
the quarter-disc bound supplies the integrable majorant
$e^{1/4}|e^{\Phi^0}|$ for each of them.
The positive classical-start final density at zero shows
that $e^{\widehat\Phi^0}$ is a fixed nonzero complex scalar
times a positive function. Thus
$\widehat I(z)/\widehat I(0)$ is an average of
$e^{\widehat\Phi^{w_z}-\widehat\Phi^0}$ with respect to
a probability measure. The quarter-disc bound puts it in
the convex region
\[
 \Re\zeta\ge e^{-1/4}\cos(1/4),\qquad
 |\zeta|\le e^{1/4},\qquad
 |\Im\zeta|\le\tan(1/4)\Re\zeta.
\]
In this region $|\operatorname{Log}\zeta|\le1/2$.
Also \eqref{IV-60-26} implies
\[
 |I/\widehat I-1|
 \le\frac{e^{1/2}}{\cos(1/4)}(e^{\epsilon_{\rm fg}}-1)
 \le2\epsilon_{\rm fg}.
\]
Consequently the holomorphic functions
$h^\bullet=\operatorname{Log}(I^\bullet/\widehat I^\bullet)$
and $l^\bullet=\operatorname{Log}(\widehat I^\bullet/
\widehat I^\bullet(0))$ satisfy
$|h^\bullet|\le4\epsilon_{\rm fg}$,
$|l^\bullet|\le1/2$, and
$\partial^n\log T^\bullet=\partial^n(c^\bullet+h^\bullet+l^\bullet)$.

The contact bound \cref{IV-60-46} on the classical-start
families gives
$|\widehat c^\bullet(z)-\widehat c^\bullet(0)|\le M_c$.
Define
\[
 H=(\widehat c^{\mathrm B}-\widehat c^{\mathrm B}(0))
       -(\widehat c^{\mathrm A}-\widehat c^{\mathrm A}(0))
                              +l^{\mathrm B}-l^{\mathrm A}.
\]
Then $H(0)=0$, $|H|\le1+2M_c$, and on the real polydisc
\[
 H(x)=\log(\widehat T^{\mathrm B}(x)/\widehat T^{\mathrm A}(x))
       -\log(\widehat T^{\mathrm B}(0)/\widehat T^{\mathrm A}(0)).
\]
Indeed their exponentials agree, and the continuous
difference takes values in $2\pi i\mathbb Z$ and is zero
at the origin. Equation \eqref{IV-62-44} bounds this real
value by $4\epsilon_3$. Apply \cref{IV-57-34} with
$\epsilon=4\epsilon_3$ and $B=\max(2,1+2M_c)$.
The resulting constant satisfies
$4^{1-\sigma}2^\sigma\le2^{n+1}$.
Finally Cauchy's formula bounds
$\partial^n(h^{\mathrm B}-h^{\mathrm A})$ by
$8\epsilon_{\rm fg}\prod_i r_i^{-1}$.
The source-independent difference between
$c^\bullet-\widehat c^\bullet$ and $\mathfrak c^\bullet$
has zero mixed derivative. Adding the estimates proves
\eqref{eq:iv-readout-comparison}.
\end{proof}

\section{Summation of the cutoff differences}
\label{sec:iv-cutoff-summation}

Write $S^T_{n;K}(\underline f)=S^T_{n;K,\ell,L^K\mathsf b}
(\underline f)$ and retain the calibrated setting of
\cref{sec:iv-small-field-discarding}. For the comparison at
cutoffs $K,K+1$, choose
\[
 u=\left\lceil\frac K{180(3n+4)}\right\rceil,
 \qquad s=\lceil K/9\rceil,
 \qquad \mathsf p_n=2\mathfrak p_0(n+1)+\tfrac12.
\]
At this point the label is in a compact interval
$J=[\mathsf b_-,\mathsf b_+]$ satisfying
\begin{equation}\label{eq:iv-admissible-label-interval}
 \mathsf b_-\ge\gamma^{-2},\qquad MR(\mathsf b_+)\le L^m,
 \qquad \ell=2L^m,\quad m\ge1.
\end{equation}
All constants chosen before the cutoff, including the source
radius $\varrho$, are fixed throughout $J$.

\begin{theorem}\label{IV-62-56}
For an arbitrary nonzero Lipschitz tuple of order $n\ge1$,
assume the comparison hypotheses of \cref{IV-62-42} and
$R_{t_1}\le1/2$, where $t_1=K-u$. Then
\begin{equation}\label{IV-62-57}
\begin{split}
 &\left|S^T_{n;K+1}(\underline f)-S^T_{n;K}(\underline f)
       -\partial^n(\mathfrak c^{\mathrm B}
                         -\mathfrak c^{\mathrm A})(0)\right|\\
 &\quad\le
 \left[8\epsilon_{\rm fg}
   +q_n\epsilon_3^{79/80}(1+2M_c)^{1/80}\right]
                         \prod_i r_i(t_1)^{-1}\\
 &\qquad +8\mathsf{op}_nR_{t_1}
              \prod_i(6\ell^3L^{4(K+1)}\|f_i\|_\infty).
\end{split}
\end{equation}
Here $q_n=2^{n+1}(320n/\pi)^n$.
If also $n\ge2$ and $d_\infty(\underline f)L^u\ge6M$,
then $|\partial^n\mathfrak c^\bullet(0)|\le
2\epsilon_{\rm mx}(t_1)$ for each $\bullet$.

There are constants $C_n(J)>0$, $K_0(\ell,n,J)$, and,
for every $d>0$, $K_1(\ell,n,J,d)\ge K_0$, such that
all the preceding smallness hypotheses hold for $K\ge K_0$,
uniformly in $\mathsf b\in J$. For $n\ge2$, $K\ge K_1$,
and all nonzero tuples with $d_\infty(\underline f)\ge d$,
\begin{equation}\label{IV-62-60}
 |S^T_{n;K+1}(\underline f)-S^T_{n;K}(\underline f)|
 \le C_n(J)\ell^{3(n+1)}K^{\mathsf p_n}L^{-K/10}
                                      \prod_i\|f_i\|_\sharp.
\end{equation}
Here $C_n(J)$ depends only on the fixed constants, $n$,
and the endpoints of $J$. The thresholds have precisely
the additional dependencies displayed, and $K_1$ may be
chosen nonincreasing in $d$.
\end{theorem}

\begin{proof}
Apply \cref{IV-61-19} at the common final depth $u$ to
both lattices. Their discarding factors are the same,
since the final coupling, site count, and physical volume
are the same. The two errors add to the last term of
\eqref{IV-62-57}. Apply the holomorphic derivative estimate
\eqref{eq:iv-readout-comparison} with $\sigma=1/80$.
Its contact term is exactly the one displayed here, since
the omitted Gaussian normalization logarithms are source
independent. This proves \eqref{IV-62-57}. Each
$\mathfrak c^\bullet$ is a difference of two contact sums
ending at depth $u$; \cref{IV-60-46} bounds its derivative
by $2\epsilon_{\rm mx}(t_1)$.

We give the uniform estimates, including the choice of
thresholds. Put $y_{[j]}=(\log(L^j\mathsf b))^{2\mathfrak p_0}$,
$N_{[j]}=\ell^3L^{3j}$, and
\[
 \overline C_M=\max\left\{\varrho^{-1},
 \frac{3c_\rho^2A_0^2}{2\mathsf b_-},
 \varrho^{-1}C_{\rm old}\frac L{L-1}E_0\right\},\qquad
 \overline y_K=(\log(L^K\mathsf b_+))^{2\mathfrak p_0}.
\]
The source radii of \cref{IV-60-40} satisfy
\begin{equation}\label{IV-62-58}
 r_i(t_1)^{-1}\le
 8n\overline C_M\overline y_K N_{[u]}\|f_i\|_\sharp.
\end{equation}
Indeed the three entries in $M^{\rm rd}_{t_1}$ are bounded
by $\overline C_M y_{[K]}N_{[u]}$: for its classical
entry use
$(e^{\rm c}_{t_1})^2L^u=c_\rho^2A_0^2y_{[u]}/\mathsf b$,
and for the old-action entry use
$m^S_{t_1}=C_{\rm old}L(L-1)^{-1}E_0N_{[u]}$.
Also $1\le y_{[u]}\le y_{[K]}\le\overline y_K$ and
$\|f_i\|_\flat\le\|f_i\|_\sharp$.

Replace $y_{[K]}$ by $\overline y_K$ and
$1+2\varrho/\mathsf b$ by $1+2\varrho/\mathsf b_-$
in $\epsilon_{\rm cl}$ to define $\overline\epsilon_{\rm cl}$.
Define the other two majorants by
\[
 \overline\epsilon_{\rm sh}=3(s-u)N_{[s]}
  \exp\!\left[C_{\rm sh}-\tfrac14\kappa_bA_1^2
                      (\log(L^u\mathsf b_-))^{2\mathfrak p_0}\right],
\]
\[
 \overline R=2M^{-3}N_{[u]}
  \exp\!\left[-c'A_0^2(\log(L^u\mathsf b_-))^{2\mathfrak p_0}
       +(C_{\rm maj}^\Sigma+2+C_{\rm low}(\mathsf b_+))\ell^3\right].
\]
They dominate the corresponding errors throughout $J$;
the last assertion uses that $C_{\rm low}(b)$ is increasing.
Since $\log(L^u\mathsf b_-)\ge K\log L/[180(3n+4)]$
and $2\mathfrak p_0\ge2$, both majorants tend to zero
after multiplication by any $K^ae^{a'K}$.

The ceiling bounds give
\begin{equation}\label{IV-62-55}
\begin{gathered}
 u(3n+4)-s\le-19K/180+3n+4,\\
 3s-\tfrac{79}{160}(K-s)+3un\le-K/10+7/2+3n,\\
 4u-s\le-K/10+4,\qquad
 (7s-K)/2\le-K/9+7/2,\qquad
 4s-K\le-5K/9+4.
\end{gathered}
\end{equation}
For example, in the second inequality the coefficient
of $s$ is $559/160$; replacing $s$ by $K/9+1$
gives $-19K/180+559/160$, while
$3un<K/180+3n$. The others follow directly from
$K/[180(3n+4)]\le u<K/[180(3n+4)]+1$ and
$K/9\le s<K/9+1$.
Thus $1\le u<s<K$, $L^u\mathsf b_-\ge\Theta_\star$,
$\epsilon_{\rm fg}\le1/8$, and the axial comparison
condition hold eventually. The latter's left side is a
constant times $L^{4s-K}$. Also
$\overline\epsilon_{\rm cl}$ is bounded by a constant
times $K^{2\mathfrak p_0+1/2}L^{-K/9}$, so that
$\overline\epsilon_{\rm cl},\overline\epsilon_{\rm sh}\le1/20$
and $\overline R\le1/2$ eventually.

Here is an explicit constant for the asserted rate. With
$h=\log(L\mathsf b_+)>1$, take
\begin{equation}\label{eq:iv-cutoff-rate-constant}
\begin{split}
 C_n(J)={}&4(8n\overline C_M)^n h^{2\mathfrak p_0(n+1)}L^{3n+4}\\
 &\quad\times\left[8C_{\rm old}C_{\rm fg}\frac L{L-1}E_0
  +q_n\max\{1,16K_wc_\rho^2A_0^2(1+2\varrho/\mathsf b_-)\}
                          (1+8E_0C_{\mathrm E,0})\right].
\end{split}
\end{equation}
Let $P=(8n\overline C_M\overline y_KN_{[u]})^n$.
Split the right side of \eqref{IV-62-57}, after removing
$\prod_i\|f_i\|_\sharp$, into
\[
 8\epsilon_{\rm fg}P,\quad
 q_n(4\overline\epsilon_{\rm cl})^{79/80}(1+2M_c)^{1/80}P,
 \quad q_n(6\overline\epsilon_{\rm sh})^{79/80}
                 (1+2M_c)^{1/80}P,\quad
 8\mathsf{op}_n\overline R(6\ell^3L^{4(K+1)})^n.
\]
This is valid because $(a+b)^{79/80}\le a^{79/80}+b^{79/80}$.
The first two terms together are at most
$\tfrac14 C_n(J)\ell^{3(n+1)}K^{\mathsf p_n}L^{-K/10}$.
For the first, use the first inequality of \eqref{IV-62-55}
and $\overline y_K\le h^{2\mathfrak p_0}K^{2\mathfrak p_0}$.
For the second, use
\[
 (1+2M_c)^{1/80}
 \le(1+8E_0C_{\mathrm E,0})\ell^{3/80}L^{3s/80}.
\]
This combines with the $79/80$ power of
$\overline\epsilon_{\rm cl}$ to give exactly
$\ell^3L^{3s-79(K-s)/160}$; the second inequality of
\eqref{IV-62-55} gives the required exponential, and
$\overline y_K^{n+1}K^{1/2}\le
h^{2\mathfrak p_0(n+1)}K^{\mathsf p_n}$ gives its polynomial.

Choose $K_0$ so that all preceding smallness conditions
hold for every $K\ge K_0$, and the last two terms of
the four-term display are at most a further quarter of
the rate. Such a $K_0$ exists by the superexponential
decay just proved. This bounds \eqref{IV-62-57} by half
the rate, uniformly on $J$ and in the tuple.
For $n\ge2$ choose $K_1\ge K_0$ so that, for every
$K\ge K_1$,
\[
 dL^u\ge6M,\qquad
 8e^3E_0C_{\mathrm E,0}\left(\frac n{2\varrho}\right)^n
 N_{[u]}e^{-dL^u/M}
 \le\tfrac12C_n(J)\ell^{3(n+1)}K^{\mathsf p_n}L^{-K/10}.
\]
The left side decays doubly exponentially in $K$ apart
from $N_{[u]}$, so this choice exists. It bounds
$4\epsilon_{\rm mx}(t_1)$ after the same norm product
is removed. Taking the least eventual threshold makes
$K_1$ nonincreasing in $d$. This proves the theorem.
\end{proof}

\begin{corollary}\label{IV-62-74}
Under \eqref{eq:iv-admissible-label-interval}, fix $n\ge2$
and $d>0$. Put
\[
 C'_n(J)=C_n(J)\sum_{j\ge0}(1+j)^{\mathsf p_n}L^{-j/10}.
\]
For $K'\ge K\ge K_1(\ell,n,J,d)$ and every nonzero
tuple with $d_\infty(\underline f)\ge d$,
\begin{equation}\label{IV-62-75}
 |S^T_{n;K'}(\underline f)-S^T_{n;K}(\underline f)|
 \le C'_n(J)\ell^{3(n+1)}K^{\mathsf p_n}L^{-K/10}
                                      \prod_i\|f_i\|_\sharp.
\end{equation}
Thus the whole cutoff sequence has a limit, uniformly
in the label on $J$ and on every set of tuples satisfying
$d_\infty(\underline f)\ge d$ and a fixed upper bound on
$\prod_i\|f_i\|_\sharp$. The limit is continuous in the label.
The same limit is obtained whenever $K_j\to\infty$ and
$a_{K_j}\beta_j\in J$ tends to the specified label.
\end{corollary}

\begin{proof}
For $k=K+j$ one has $k\le K(1+j)$, whence
\[
 \sum_{k\ge K}k^{\mathsf p_n}L^{-k/10}
 \le K^{\mathsf p_n}L^{-K/10}
                    \sum_{j\ge0}(1+j)^{\mathsf p_n}L^{-j/10}.
\]
The last series converges by geometric domination.
Telescoping \eqref{IV-62-60} proves \eqref{IV-62-75},
and the uniform Cauchy criterion gives the limit.
At fixed cutoff, the partition function and its source
derivatives are entire functions of the bare coupling:
the configuration space is compact and all plaquette
actions are bounded. The real partition function is
strictly positive. Hence every finite-cutoff cumulant
is continuous in the positive real label, and its uniform
limit on $J$ is continuous there.
Finally write $b_j=a_{K_j}\beta_j$. The difference from
the asserted limit is bounded by the uniform cutoff
error at $b_j$, plus the difference of the two limiting
values at $b_j$ and its limit. Both tend to zero.
\end{proof}

\section{The limit at every positive physical coupling}
\label{sec:iv-all-label-cutoff-limit}

\begin{lemma}\label{IV-61-6}
The constants can be chosen so that the calibrated setting
used in \cref{IV-62-56,IV-62-74} holds on every label
interval satisfying \eqref{eq:iv-admissible-label-interval}.
The choices precede the torus, the cutoff, and the labels.
\end{lemma}

\begin{proof}
Use the actual canonical auxiliary data and the ordered
choice of Proposition~\ref{II-202-29}. Its real
background, response, Gaussian, and normalizer hypotheses
are the global hypotheses imposed above. The fixed radii
are chosen with $\alpha_0\le a_{\rm H}$, and the decay
and normalizer floors include $4\beta_{\rm s}\kappa\ge1$
and the rates used in the rooted sums. These choices
precede $\gamma$. In particular
\eqref{eq:common-flow-bound} gives $E_0\ge2E_{\mathsf N}$.
Take $E_0^{\rm s}=E_0$ and
$N_\star^{\rm s}=L(E_0+\tfrac32\alpha_0^2B_\sharp)/(L-1)$.

Every coefficient in \eqref{eq:iv-flow-smallness} is now
fixed. Its three small quantities are constant multiples
of
\[
 \gamma(\log\gamma^{-2})^{\mathfrak p_0},\qquad
 \exp[-\gamma_2A_1^2(\log\gamma^{-2})^{2\mathfrak p_0}/2],
\]
and tend to zero as $\gamma\downarrow0$. Thus the three
inequalities hold after decreasing $\gamma$. Impose the
direct-source half-radius requirement at the same stage;
its left side is $C_{\mathcal A}\delta(\gamma^{-2})$.
The response and classical calibration requirements are
positive fixed-radius ceilings on $\varepsilon$ and
$\delta$, so they hold uniformly above a finite threshold.
The derived numbers $\tau_\natural,\Theta_\star$ are then
finite; they are thresholds on the final intermediate
coupling, which tends to infinity in the cutoff comparison.

Fix $\eth$ and the history constants before this last
choice. The case-II and summability thresholds used in
\cref{lem:iv-early-history-tail} are finite and independent
of $\gamma$. They can therefore be imposed on
$\log\gamma^{-2}$ simultaneously. Also require
$\gamma^{-2}>\bar B$ and
$R(\gamma^{-2})\ge C_L/(2L)$. If $MR(\mathsf b')\le L^m$
and $\mathsf b'\ge\gamma^{-2}$, monotonicity of $R$ gives
$2L^{m+1}\ge C_LM$, as required in
\eqref{eq:iv-classical-room}. The line
$x_j=L^{K-j}\mathsf b'$ has zero increments, is admissible,
and has precisely the torus room required in
Theorem~\ref{II-203-9}. The actual mass bounds and vacuum
lower bound used in the discarding proof consequently
apply to this same line and the same vacuum constant.

Finally $K_{\rm Q}^\vee,K_{\rm L}^\vee,w_\star,E_0$
and $K_{\mathbf M}^{\rm f}$ are positive and finite.
Thus $\mu_\star,\sigma_{\rm src}>0$, and the stated
minimum defining $\varrho$ is positive. Its two defining
inequalities give \eqref{eq:iv-flow-source-room} for every
$\mathsf b'\ge\gamma^{-2}$. At zero source the added
family vanishes identically; the sourced recursion and
the canonical real progression therefore coincide step
by step. This verifies the interface used in the history
comparison. The constants and this source radius involve
no finite-lattice datum.
\end{proof}

For $s\in\mathbb Z$ and $m+s\ge1$, let
$f^{(s)}(x')=f(L^{-s}x')$ on $T_{L^s\ell}$.
Dilation multiplies both the Euclidean and sup distances
by $L^s$. Consequently
\begin{equation}\label{IV-63-5}
 \|f^{(s)}\|_\infty=\|f\|_\infty,\qquad
 \operatorname{Lip}f^{(s)}=L^{-s}\operatorname{Lip}f,
 \qquad d_\infty(\underline f^{(s)})=L^sd_\infty(\underline f).
\end{equation}
Supports are carried onto their dilates. In particular
the absence of a common support point is preserved, and
$\|f^{(s)}\|_\sharp\le\|f\|_\sharp$ when $s\ge0$.

\begin{lemma}\label{III-9-16}\label{IV-63-6}
For $K,m\ge1$, $-m<s<K$, $\beta>0$, and every
Lipschitz tuple, one has
\begin{equation}\label{IV-63-10}
 S^T_{n;K,\ell,\beta}(\underline f)
  =L^{sn}S^T_{n;K-s,L^s\ell,\beta}(\underline f^{(s)}),
 \qquad a_{K-s}\beta=L^sa_K\beta.
\end{equation}
For every compact interval $J=[b_-,b_+]\subset(0,\infty)$,
there is an integer $s\ge0$ for which
\begin{equation}\label{IV-63-7}
 L^sb_-\ge\gamma^{-2},\qquad MR(L^sb_+)\le L^{m+s}.
\end{equation}
One may take any sufficiently large integer satisfying
\begin{equation}\label{IV-63-8}
 L^s\ge\max\{\gamma^{-2}/b_-,b_+\},\qquad
 L^{m+s/2}\ge LM(4r/e)^r.
\end{equation}
\end{lemma}

\begin{proof}
For $s\ge0$, \eqref{IV-63-10} is
Lemma~\ref{I-2-28}. Explicitly, dilation is a bijection
of the bonds and plaquettes and preserves their orientations,
plaquette products, and product Haar measure. Corresponding
source coefficients agree because
$f^{(s)}(L^sx_p)=f(x_p)$ and
$a_{K-s}^{-1}L^s=a_K^{-1}$. Thus
\[
 Z_{K,\ell,\beta}\left[\sum_i z_if_i\right]
 =Z_{K-s,L^s\ell,\beta}\left[\sum_iL^sz_if_i^{(s)}\right].
\]
Differentiating local logarithms at zero gives the factor
$L^{sn}$. If $s<0$, apply the same nonnegative-dilation
identity with exponent $-s$ to the data
$(K-s,m+s,\beta,\underline f^{(s)})$; its dilated
functions are the original $f_i$. This proves the
identity in both directions. Equation \eqref{IV-63-5}
follows by taking infima over lifts for the distances
and then the minimum over the product of supports.

For \eqref{IV-63-7}, put $x=L^sb_+$ and suppose
\eqref{IV-63-8}. Then $x\ge\gamma^{-2}>e$.
The least-power definition of $R$ gives
$R(x)<L(\log x)^r$. Also
$\log x\le2s\log L=4y$, where $y=s\log L/2>0$.
Since $y^r\le(r/e)^re^y$, we obtain
\[
 MR(x)<LM(4r/e)^r L^{s/2}\le L^{m+s}.
\]
The other inequality in \eqref{IV-63-7} is already in
\eqref{IV-63-8}. Both conditions in \eqref{IV-63-8}
hold for all sufficiently large $s$, since their right
sides are fixed.
\end{proof}

\begin{theorem}\label{IV-63-9}
Fix $\ell=2L^m$, $m\ge1$. For every $\mathsf b>0$,
$n\ge2$, and real Lipschitz functions $f_1,\ldots,f_n$
with $\bigcap_i\operatorname{supp}f_i=\emptyset$, the limit
\begin{equation}\label{IV-63-3}
 S_n^{(\ell,\mathsf b)}(\underline f)
   =\lim_{K\to\infty}S^T_{n;K,\ell,L^K\mathsf b}(\underline f)
\end{equation}
exists along all positive integer cutoffs. Convergence is
locally uniform in $\mathsf b\in(0,\infty)$.

More precisely, fix $J=[b_-,b_+]\subset(0,\infty)$,
$d>0$, and an integer $s\ge0$ satisfying
\eqref{IV-63-7}. Put
\[
 C'_s=C'_n([L^sb_-,L^sb_+]),\qquad
 k_s=K_1(L^s\ell,n,[L^sb_-,L^sb_+],L^sd).
\]
For every nonzero tuple with $d_\infty(\underline f)\ge d$,
every $\mathsf b\in J$, and $K\ge k_s+s$,
\begin{equation}\label{IV-63-26}
 \left|S^T_{n;K,\ell,L^K\mathsf b}(\underline f)
                  -S_n^{(\ell,\mathsf b)}(\underline f)\right|
 \le L^{sn}C'_s(L^s\ell)^{3(n+1)}
           (K-s)^{\mathsf p_n}L^{-(K-s)/10}
                                      \prod_i\|f_i\|_\sharp.
\end{equation}
For all $K\ge1$, the finite-cutoff cumulant and its limit
are bounded in modulus by
\begin{equation}\label{IV-63-12}
 \mathsf{op}_n\prod_i(6\ell^3L^{4(k_s+s)}\|f_i\|_\infty)
 +L^{sn}C'_s(L^s\ell)^{3(n+1)}
                  k_s^{\mathsf p_n}L^{-k_s/10}
                                      \prod_i\|f_i\|_\sharp.
\end{equation}
The constants are independent of the tuple except for
its lower spread bound $d$.

For every integer $s>-m$, the scaling law is
\begin{equation}\label{IV-63-13}
 S_n^{(\ell,\mathsf b)}(\underline f)
  =L^{sn}S_n^{(L^s\ell,L^s\mathsf b)}(\underline f^{(s)}).
\end{equation}
If $K_j\to\infty$, $\beta_j>0$, and
$a_{K_j}\beta_j\to\mathsf b>0$, then
\begin{equation}\label{eq:iv-moving-label-limit}
 S^T_{n;K_j,\ell,\beta_j}(\underline f)
                  \longrightarrow S_n^{(\ell,\mathsf b)}(\underline f).
\end{equation}
No monotonicity of the sequence $K_j$ is required.
\end{theorem}

\begin{proof}
If an entry is zero, multilinearity makes every cumulant
zero. Otherwise compactness of the supports gives
$d_\infty(\underline f)>0$. For the dilated tuple,
\eqref{IV-63-5} gives spread at least $L^sd$ and norms
no larger than the original ones. The interval
$[L^sb_-,L^sb_+]$ satisfies
\eqref{eq:iv-admissible-label-interval} on $T_{L^s\ell}$.
Apply \cref{IV-62-74} there with cutoff $k=K-s$ and
multiply its estimate by $L^{sn}$ using
\eqref{IV-63-10}. This is the Cauchy estimate with the
right side of \eqref{IV-63-26}. It tends to zero uniformly
in the label, so the whole sequence converges, and passing
the second cutoff to infinity proves \eqref{IV-63-26}.

For $K\ge k_s+s$, compare to cutoff $k_s+s$ by the
same Cauchy estimate. The finite-cutoff moment bound
used in \cref{IV-61-19} bounds that term by the first
term of \eqref{IV-63-12}. For $K<k_s+s$ the same
finite bound is even smaller. This proves the uniform
bound for all cutoffs and, by passage to the limit, for
the limit itself. Taking $J=\{\mathsf b\}$ proves
existence at every positive label. Taking any compact
interval proves local uniformity and hence continuity,
since each finite-cutoff cumulant is continuous.

For any fixed integer $s>-m$, \eqref{IV-63-10} holds
for all sufficiently large $K$. The index $K-s$ also
runs through all sufficiently large integers, so passage
to the limit gives \eqref{IV-63-13}, including negative
$s$. For the last assertion put $b_j=a_{K_j}\beta_j$.
Eventually $b_j\in[\mathsf b/2,2\mathsf b]$; subtract
the limit at $b_j$ and use the uniform cutoff estimate
there. Continuity then replaces $b_j$ by $\mathsf b$.
\end{proof}

\begin{proposition}\label{IV-63-19}\label{IV-63-21}
The limits in \cref{IV-63-9} are real, symmetric, and
multilinear on separated tuples. They are invariant under
every cubic lattice isometry at any fixed scale. They
are smooth in the positive physical inverse coupling,
with
\begin{equation}\label{IV-63-23}
 \partial_{\mathsf b}^{\nu}S_n^{(\ell,\mathsf b)}(\underline f)
 =(-1)^\nu S_{n+\nu}^{(\ell,\mathsf b)}
                   (\underbrace{1,\ldots,1}_{\nu},\underline f).
\end{equation}
All derivatives of finite-cutoff cumulants converge to
these derivatives locally uniformly in $\mathsf b$,
with \eqref{IV-63-26} at order $n+\nu$ and the same
spread bound $d$.
\end{proposition}

\begin{proof}
Reality, symmetry, and multilinearity hold at each finite
cutoff and pass to the limit. The domain is preserved
under the asserted linear combinations: if only entry
$i$ changes, the common intersection of the other
supports is disjoint from each of its two possible
supports and hence from their union. A lattice isometry
at a fixed scale permutes bonds and plaquettes at every
sufficiently fine cutoff. Haar measure and the Wilson
action are invariant; reindexing the source sum gives
finite-cutoff invariance, and then invariance of the limit.

Adding constant entries $1$ leaves the support intersection
empty and does not change the norm product or the spread.
For the last fact, the enlarged diameter is at least the
old minimum; conversely choose a minimizing old point
tuple and put all new points at its first point.
At finite cutoff, differentiation in $\mathsf b$ inserts
minus the constant-source observable because
\[
 \exp[-A^{a_K}(\mathsf b-w_z,U)]
 \quad\hbox{depends on $\mathsf b$ and a constant source $t$
 through $\mathsf b-t$.}
\]
The partition function and its derivatives are entire,
and its real value at zero source is positive. On a
local zero-free neighborhood its logarithm therefore gives
\[
 \partial_{\mathsf b}^{\nu}S^T_{n;K,\ell,L^K\mathsf b}
       (\underline f)
 =(-1)^\nu S^T_{n+\nu;K,\ell,L^K\mathsf b}
                   (1,\ldots,1,\underline f).
\]
Both this sequence and its next derivative converge
uniformly on each compact positive interval by
\cref{IV-63-9}. Integrating the derivative identity on
such an interval and passing to the limit proves that
the first limit has the second as its derivative.
Iterating proves \eqref{IV-63-23} and the stated uniform
convergence with its rate.
\end{proof}

\begin{corollary}\label{IV-63-18}
Let $0<g\le\gamma$ and let $\ell=2L^m$, $m\ge1$.
The cutoff limits for separated tuples converge uniformly
for $\mathsf b\in[g^{-2}-\bar B,X_g+\bar B]$.
In particular this holds on every side in $\mathfrak L(g)$.
For $\mathsf b\in I(g)$ the pairs $(K,L^K\mathsf b)$
are first-passage data at $g$ for every $K$. More generally,
a sequence $(K_j,\beta_j)$ of first-passage data with
$K_j\to\infty$ and convergent labels has the limit
\eqref{eq:iv-moving-label-limit}.
\end{corollary}

\begin{proof}
The condition $\gamma^{-2}\ge2B_\sharp$ gives
$\bar B\le L[2(L-1)]^{-1}g^{-2}<g^{-2}$.
The displayed label interval is consequently a compact
subset of $(0,\infty)$, so \cref{IV-63-9} applies.
The assertion about $I(g)$ is \cref{III-5-19}.
The same result bounds every first-passage label between
$g^{-2}+B_\sharp-L^{-1}\bar B$ and $X_g+\bar B$;
the lower endpoint is positive. Thus a convergent sequence
of those labels has positive limit, as required in
\eqref{eq:iv-moving-label-limit}.
\end{proof}

These conclusions concern fixed tori with $m\ge1$ and
tuples with no common support point. The comparison
\eqref{IV-62-57} retains a contact subtraction for general
tuples; it does not assert convergence of the unrenormalized
one-point function or of coincident-support cumulants.


\section{Dilation and the range of torus sizes}
\label{sec:iv-general-dilation}

For this section a side is $\ell=2L^m$ with $m\in\mathbb Z$,
and a cutoff for that side is an integer $K$ with $K+m\ge1$.
The finite Wilson model uses spacing $a_K=L^{-K}$ and
$N=2L^{K+m}$ sites in each coordinate direction. In particular
$N\ge2L\ge26$, so its bonds and plaquettes have the same
incidence properties as before. For a fixed side every sufficiently
large integer $K$ is an allowed cutoff. The value $K=-m$, which
would give only two sites in each direction, is excluded.

For $z\in\mathbb Z$, write
\[
 \delta_z:T_\ell\longrightarrow T_{L^z\ell},\qquad
 \delta_zx=L^zx,\qquad f^{[z]}=f\circ\delta_z^{-1}.
\]
On $\mathbb R^3$ the same notation means $f^{[z]}(x)=f(L^{-z}x)$.
This agrees with the dilation notation used in
\cref{III-9-16,IV-63-5}.

\begin{proposition}[Finite dilation for all allowed cutoffs]\label{IV-12-2}
Let $K,m,z\in\mathbb Z$, $K+m\ge1$, $\ell=2L^m$ and $\beta>0$.
Then $K-z$ is a cutoff for $L^z\ell$, and
\begin{equation}\label{IV-12-4}
 Z_{K,\ell,\beta}[w]
       =Z_{K-z,L^z\ell,\beta}[L^zw^{[z]}],
\end{equation}
for every complex function $w$ on the torus. For $n\ge1$,
\begin{equation}\label{IV-12-5}
 S^T_{n;K,\ell,\beta}(\underline f)
       =L^{zn}S^T_{n;K-z,L^z\ell,\beta}(\underline f^{[z]}).
\end{equation}
The label changes by $a_{K-z}\beta=L^za_K\beta$.
Distances, support radii and separations are multiplied by $L^z$,
and, with the same $M$ on both tori,
\begin{equation}\label{IV-12-3}
 \|f^{[z]}\|_\infty=\|f\|_\infty,\qquad
 \operatorname{Lip}f^{[z]}=L^{-z}\operatorname{Lip}f,\qquad
 \|f^{[z]}\|_\sharp\le\max(1,L^{-z})\|f\|_\sharp.
\end{equation}
\end{proposition}
\begin{proof}
Multiplication by $L^z$ takes the coset
$x+\ell\mathbb Z^3$ onto $L^zx+L^z\ell\mathbb Z^3$.
It therefore defines a homeomorphism with inverse $\delta_{-z}$.
Taking the minimum over lifts proves the distance identity;
taking suprema of difference quotients proves \eqref{IV-12-3}.
Supports are carried to their images, since this is a
homeomorphism. The same facts hold on $\mathbb R^3$.

The new spacing is $a_{K-z}=L^za_K$, and the number of sites is
$2L^{(K-z)+(m+z)}=N$. Dilation thus bijects positively oriented
bonds and plaquettes, preserves their incidence and cyclic order,
and takes $x_p$ to $x_{\delta_zp}$. Set
$U^\delta(\delta_zb)=U(b)$. This is a permutation of the factors
of the configuration space and preserves product Haar measure.
The plaquette holonomy, and hence $a_p$, is unchanged. Therefore
the action is unchanged, while
\[
 \mathcal O_{K-z}(L^zw^{[z]})(U^\delta)
       =a_{K-z}^{-1}L^z\sum_p w(x_p)a_p(U)
       =\mathcal O_K(w)(U).
\]
Changing variables in the integral proves \eqref{IV-12-4}.
Taking mixed derivatives of its normalized logarithm at zero
proves \eqref{IV-12-5}. Equivalently, each observable has acquired
the factor $L^z$ and joint cumulants are multilinear.
\end{proof}

\begin{corollary}[The full cutoff limit on every torus side]
\label{cor:iv-all-sides-cutoff}
Let $m\in\mathbb Z$, $\ell=2L^m$, $\mathsf b>0$, $n\ge2$, and let
$f_1,\ldots,f_n$ be real Lipschitz functions on $T_\ell$ such that
$\bigcap_i\operatorname{supp}f_i=\emptyset$. Then
\[
 S_n^{(\ell,\mathsf b)}(\underline f)
       =\lim_{\substack{K\to\infty\\ K+m\ge1}}
                S^T_{n;K,\ell,L^K\mathsf b}(\underline f)
\]
exists along the whole sequence, locally uniformly for
$\mathsf b\in(0,\infty)$. For every integer $z$,
\begin{equation}\label{IV-150-43}
 S_n^{(\ell,\mathsf b)}(\underline f)
       =L^{zn}S_n^{(L^z\ell,L^z\mathsf b)}(\underline f^{[z]}).
\end{equation}
For any $K_j\to\infty$ and $\beta_j>0$ with
$a_{K_j}\beta_j\to\mathsf b>0$, the same limit is obtained from
$S^T_{n;K_j,\ell,\beta_j}(\underline f)$.
\end{corollary}
\begin{proof}
For $m\ge1$ this is \cref{IV-63-9}. For $m\le0$ set $z=1-m$.
Then $L^z\ell=2L$, and the dilated functions remain real and
Lipschitz, with empty common support intersection. At every
allowed $K$, \eqref{IV-12-5} identifies the desired cumulant with
$L^{zn}$ times the cumulant on $T_{2L}$ at cutoff $K-z$ and label
$L^z\mathsf b$. This cutoff runs through all sufficiently large
integers. Theorem~\ref{IV-63-9} gives convergence and local uniformity
on the dilated label interval, and hence on the original interval.
Its moving-label conclusion gives the last assertion in the same
way. Now both sides of \eqref{IV-12-5} have limits for every
integer $z$, so that identity yields \eqref{IV-150-43}.
\end{proof}

\subsection{Dilation of limits on \texorpdfstring{$\mathbb R^3$}{R3}}

On the infinite lattice define $U^{(z)}(L^zb)=U(b)$ and let
$\nu^{(z)}$ be the pushforward of a law $\nu$ under this map.
It is a homeomorphism of the product configuration spaces,
with inverse given by $-z$.

\begin{proposition}[Torus states under dilation]\label{IV-23-13}
For any law $\nu$ on the lattice of spacing $a_K$ and compactly
supported bounded functions $f_i$,
\begin{equation}\label{IV-23-17}
 S^T_n[\nu](\underline f)
       =L^{zn}S^T_n[\nu^{(z)}](\underline f^{[z]}).
\end{equation}
For every allowed finite torus,
\begin{equation}\label{IV-23-20}
 (\widetilde\mu_{K,\ell,\beta})^{(z)}
       =\widetilde\mu_{K-z,L^z\ell,\beta}.
\end{equation}
If $K,K-z\ge1$, dilation maps
$\mathcal G^{\rm tor}_{K,\beta}$ bijectively onto
$\mathcal G^{\rm tor}_{K-z,\beta}$.
\end{proposition}
\begin{proof}
Only finitely many plaquettes meet the support of each $f_i$.
Relabelling them gives
$\mathcal O_K(f)(U)=L^z\mathcal O_{K-z}(f^{[z]})(U^{(z)})$,
exactly as in the finite proof. Pushforward preserves the
joint law of these corresponding observables, so multilinearity
of cumulants proves \eqref{IV-23-17}.

The finite configuration map in \cref{IV-12-2} commutes with
periodic lifting: both ways assign to $L^zb$ the variable
at the class of $b$ modulo $\ell\mathbb Z^3$. Since that map
pushes forward the finite Wilson measure, it gives
\eqref{IV-23-20}. If
$\widetilde\mu_{K,2L^{m_i},\beta}\to\nu$ weakly with $m_i\to\infty$,
the configuration homeomorphism implies weak convergence of
the pushforwards. By \eqref{IV-23-20}, these are
$\widetilde\mu_{K-z,2L^{m_i+z},\beta}$; eventually $m_i+z\ge1$.
Thus $\nu^{(z)}\in\mathcal G^{\rm tor}_{K-z,\beta}$ by
\cref{IV-14-16}. Applying the same argument with $-z$ proves
surjectivity, and the configuration inverse gives injectivity.
\end{proof}

To state the limiting conclusion without imposing a reference
coupling, consider the following two classes. A limit along tori
is a family $S=(S_n)_{n\ge2}$ obtained on all separated $C_c^1$
tuples from a sequence $(K_j,\ell_j,\beta_j)$ with
$K_j\to\infty$, $\ell_j=2L^{m_j}\to\infty$, $m_j\in\mathbb Z$,
$K_j+m_j\ge1$ and $\beta_j>0$. A limit through torus states is
obtained on the same tuples from
$S^T_n[\nu_j]$, where $K_j\to\infty$, $K_j\ge1$, $\beta_j>0$ and
$\nu_j\in\mathcal G^{\rm tor}_{K_j,\beta_j}$.
In either case a label is a limit point of $a_{K_j}\beta_j$
along the defining sequence. For any family set
\begin{equation}\label{IV-23-25}
 S_n^{[z]}(X)=L^{-zn}S_n(X^{[-z]}).
\end{equation}

\begin{proposition}[Dilation of continuum limits]\label{IV-23-24}
For either of these two classes and every $z\in\mathbb Z$,
$S\mapsto S^{[z]}$ is a bijection from limits with label
$\mathsf b$ to limits with label $L^z\mathsf b$.
For limits along tori the defining sequence becomes
$(K_j-z,L^z\ell_j,\beta_j)$. For limits through torus states it
becomes $(K_j-z,\beta_j,\nu_j^{(z)})$, after discarding the finite
initial segment with $K_j-z<1$. The inverse is dilation by $-z$.
If $S$ obeys the table $C_n(d;R)$, then $S^{[z]}$ obeys
\begin{equation}\label{IV-23-8}
 C_n^{[z]}(d;R)=\max(1,L^{-z})^n
                         C_n(L^{-z}d;L^{-z}R).
\end{equation}
\end{proposition}
\begin{proof}
For a compactly supported function, changing the periodization
index from $v\in\ell\mathbb Z^3$ to $L^zv$ gives
\[
 (f^{\langle\ell\rangle})^{[z]}
       =(f^{[z]})^{\langle L^z\ell\rangle}.
\]
Use \eqref{IV-12-5} with the tuple $X^{[-z]}$ and pass to the
limit. The transformed indices satisfy
$(K_j-z)+(m_j+z)=K_j+m_j\ge1$, and both the cutoff and the side
still tend to infinity. This proves the assertion along tori.
For the other class, \cref{IV-23-13} gives the required membership
of $\nu_j^{(z)}$ and its cumulant identity gives the same limit.
Multiplication by $L^z$ carries the label sequence to
$a_{K_j-z}\beta_j$, preserving subsequential limits in both
directions. Finally $(X^{[-z]})^{[z]}=X$ and the two scalar
factors in \eqref{IV-23-25} cancel, proving the asserted inverse.

For the bound, a tuple of type $(d;R)$ is taken by $[-z]$ to
one of type $(L^{-z}d;L^{-z}R)$, and
$\|f^{[-z]}\|_\sharp\le\max(1,L^z)\|f\|_\sharp$ by
\eqref{IV-12-3}. Thus the scalar multiplying the original table
is $L^{-zn}\max(1,L^z)^n=\max(1,L^{-z})^n$.
Symmetry and multilinearity are preserved as well.
\end{proof}

The constructed growing-torus and volume-first limits belong
to these classes. The proposition describes their dilates and
labels. A shifted sequence need not be a first-passage sequence
at the original reference coupling. Properties stated with
constants at a fixed reference coupling therefore require their
own hypotheses at that coupling after dilation.




\section{Curvature projections and the Gaussian comparison}
\label{sec:curvature-projections}

The short-distance calculation uses the curvature of the Gaussian
fluctuation in an actual renormalization step.  Its covariance is an
orthogonal projection, and the projections at successive scales are
nested.  This gives an exact telescoping identity before any estimate
is made.  We establish that identity first, including its relation to
the gauge-fixed covariance, and then estimate its kernels.

\subsection{The nested projections}

Fix a lattice $(\ell,K)$.  In this subsection cochains are real, their
pairings are unweighted sums, and $d$ is the coboundary given by signed
sums along edges and plaquette boundaries.  Write $\mathsf{Ch}^i_n$
for the degree-$i$ cochains on the level-$n$ lattice and put
$\eta_n=L^{-n}$.  The level-$n$ lattice has spacing $L^n$ in bare
units.  The straight averages from the bare lattice are
\begin{align*}
 (Q^{\rm si}_n\lambda)(y)
 &=L^{-3n}\sum_{x\in B^n(y)}\lambda(x),\\
 (Q^s_na)_\mu(y)
 &=L^{-3n}\sum_{x\in B^n(y)}\sum_{i=0}^{L^n-1}
                                      a_\mu(x+ie_\mu),\\
 (Q^{\rm pl}_nF)_{\mu\nu}(y)
 &=L^{-3n}\sum_{x\in B^n(y)}\sum_{i,i'=0}^{L^n-1}
                  F_{\mu\nu}(x+ie_\mu+i'e_\nu).
\end{align*}
For an oriented plaquette the last formula has the corresponding
orientation sign.  At $n=0$ all three averages are the identity.
The same formulas define the one-step averages, denoted by boldface
$\mathbf Q^{\rm si}$, $\mathbf Q^s$, and $\mathbf Q^{\rm pl}$.

Let $\mathsf E=d\mathsf{Ch}^1_0$ and define
\begin{equation}\label{IV-44-7}
 \mathsf{Fl}_n=\{F\in\mathsf E:Q^{\rm pl}_nF=0\},\qquad
 \mathsf{Bg}_n=\mathsf E\ominus\mathsf{Fl}_n,\qquad
 \Pi_n=\operatorname{proj}_{\mathsf{Bg}_n}.
\end{equation}
In particular, $\mathsf{Bg}_0=\mathsf E$.  We use
\begin{equation}\label{IV-44-10}
 \mathcal C_{(j)}=\Pi_{j-1}-\Pi_j,\qquad
 C^{\rm free}_K=\Pi_0,\qquad
 C^{\rm bg}_{K,s}=\Pi_{K-s},\qquad
 C^{\rm fl}_{K,s}=\Pi_0-\Pi_{K-s}
                         =\sum_{j=1}^{K-s}\mathcal C_{(j)}.
\end{equation}
Kernels here and below are matrices in the oriented positive
plaquette basis; they include no lattice-volume weights.

\begin{lemma}\label{IV-44-9}
The averages satisfy the chain identities
\[
 Q^s_nd=dQ^{\rm si}_n,\qquad dQ^s_n=Q^{\rm pl}_nd,
 \qquad Q^{\rm pl}_{n+1}=\mathbf Q^{\rm pl}Q^{\rm pl}_n,
\]
and the analogous composition identities for the other two averages.
Moreover,
\begin{equation}\label{eq:curvature-average-norm}
 \|Q^{\rm pl}_nF\|^2\le L^n\|F\|^2,
 \qquad \mathsf{Fl}_n=d\ker Q^s_n.
\end{equation}
The spaces $\mathsf{Fl}_n$ increase with $n$; hence
$\Pi_n\Pi_m=\Pi_{\max\{m,n\}}$, and $\mathcal C_{(j)}$ is
the orthogonal projection onto
$\mathsf{Bg}_{j-1}\cap\mathsf{Fl}_j$.
\end{lemma}

\begin{proof}
Telescoping the edge and face sums proves the chain identities.
The entries of $Q^{\rm pl}_n$ are nonnegative after fixing an
orientation, with row sum $L^{2n}$ and column sum $L^{-n}$.
Cauchy--Schwarz in each row and summation give the first inequality
in \eqref{eq:curvature-average-norm}.

If $Q^{\rm pl}_nda=0$, the coarse cochain $Q^s_na$ is a cocycle.
On a periodic cubic lattice every real cocycle is $d\lambda+c$,
where $c$ is constant: subtract its three periods and integrate the
remaining cochain along paths.  Set $\theta(x)=\lambda(y)$ for
$x\in B^n(y)$ and let $c''=L^{-n}c$.  Then
$d(d\theta+c'')=0$ and $Q^s_n(d\theta+c'')=d\lambda+c$.
Subtracting this cocycle from $a$ proves the reverse inclusion in
$\mathsf{Fl}_n=d\ker Q^s_n$.  The forward inclusion and the
remaining assertions follow from the chain identities.
\end{proof}

For $B\in\mathsf{Ch}^1_n$, let $H_n^{\rm cv}B$ be the unique
$F\in\mathsf{Bg}_n$ satisfying $Q^{\rm pl}_nF=dB$.  Existence
uses surjectivity of $Q^s_n$: the pivot lift for the one-step
straight average is a right inverse, and composition gives a right
inverse for $Q^s_n$; see the proof of the Gaussian critical-point
construction in Theorem~\ref{thm:gaussian-critical-germs}.  Equivalently,
\begin{equation}\label{IV-44-13}
 H_n^{\rm cv}B=\Pi_n da\quad(Q^s_na=B),\qquad
 \|H_n^{\rm cv}B\|^2
   =\min_{Q^s_na=B}\|da\|^2
   \ge\eta_n\|dB\|^2.
\end{equation}
Indeed, subtracting two lifts changes their curvatures by
$\mathsf{Fl}_n$; \cref{IV-44-9} then gives the formula and its
uniqueness.  The inequality follows from
\eqref{eq:curvature-average-norm}.  This response depends only on
$dB$, has range $\mathsf{Bg}_n$, and $H_0^{\rm cv}=d$.

The response in \eqref{IV-44-13} is the curvature of the free
Landau response of equations~\eqref{I-82-31}--\eqref{II-56-6}.  More precisely,
in level-$n$ units its fine mesh is $\eta_n$, its fine curl is $D$,
and its straight potential average is $Q$.  If
\[
 D^*Dh=Q^*\omega,\qquad Qh=B,
\]
then
\begin{equation}\label{eq:cv-landau-identification}
 \eta_n^2Dh=H_n^{\rm cv}B.
\end{equation}
To check the factors, $d(\eta_nh)=\eta_n^2Dh$,
$Q^s_n(\eta_nh)=Qh$, and the fine weighted pairing is
$\eta_n^3$ times the unweighted one.  Testing the equation against
$a\in\ker Q^s_n$ makes $\eta_n^2Dh$ orthogonal to
$d\ker Q^s_n=\mathsf{Fl}_n$.  Formula
\eqref{IV-44-13} now proves \eqref{eq:cv-landau-identification}.

\subsection{The covariance of one step}

At level $k<K$, let $Q^{\rm lin}=\mathbf Q^s-d\kappa$ be
the linearized gauge-covariant average at the trivial field.  Here
$\kappa B$ is the mean of the block-path integrals used in the
one-step gauge; in particular,
\[
 Q^{\rm lin}d\theta=d(\theta|_{\mathbb T^{(k+1)}}).
\]
Let $V=\ker Q^{\rm lin}$.  The pivot parametrization
$C_{\mathbf1,\mathbb R}$ identifies the independent bond
coordinates with $V$: it keeps every nonpivot coordinate and sets
the pivot at a coarse bond $c$ equal to
$-L^2\sum_{b\ne b_0(c)}w_c(b)B(b)$.  The gauge-fixing form is
\[
 G(B,B')=\sum_y\sum_{x\in B(y)\setminus\{y\}}
                         \overline B(y,x)\overline {B'}(y,x),
\]
where the bar denotes the average of the coordinate-path integrals,
not complex conjugation.  Put
\begin{equation}\label{IV-44-16}
 B_k^{\rm cv}=\eta_k^{-1}(H_k^{\rm cv}C_{\mathbf1,\mathbb R})^*
                              H_k^{\rm cv}C_{\mathbf1,\mathbb R},
 \quad \widehat G=C_{\mathbf1,\mathbb R}^*GC_{\mathbf1,\mathbb R},
 \quad A_k^{\rm cv}=B_k^{\rm cv}+\widehat G.
\end{equation}

\begin{lemma}\label{IV-44-20}
One has
\[
 H_k^{\rm cv}V=\mathsf{Bg}_k\cap\mathsf{Fl}_{k+1},\qquad
 \ker(H_k^{\rm cv}|_V)
   =\{d\theta:\theta|_{\mathbb T^{(k+1)}}=0\}.
\]
On $V$, the kernels of $G$ and $H_k^{\rm cv}$ are complementary.
\end{lemma}

\begin{proof}
The chain identity gives
$Q^{\rm pl}_{k+1}H_k^{\rm cv}B=d\mathbf Q^sB=dQ^{\rm lin}B$,
so the first range is contained in the second.  Conversely, write
$F=H_k^{\rm cv}B$ for $F\in\mathsf{Bg}_k\cap\mathsf{Fl}_{k+1}$.
Then $Q^{\rm lin}B=d\lambda+c'$.  Choose the constant
$c''=L^{-1}c'$ and a site function with coarse values
$\theta=\lambda+\kappa c''$.  The cocycle $c''+d\theta$
has the same $Q^{\rm lin}$-average as $B$, so subtracting it
puts $B$ in $V$ without changing $H_k^{\rm cv}B$.

If $H_k^{\rm cv}B=0$, \eqref{IV-44-13} gives $dB=0$.
Thus $B=d\theta+c$; the periods of $Q^{\rm lin}B=0$
force $c=0$, and its gradient part makes $\theta$ constant on
the coarse sites.  Subtract that constant.  Finally, for $B\in V$
put $\theta(x)=\overline B(y,x)$ in $B(y)$ and $\theta(y)=0$.
The cochain $B''=B-d\theta$ has all its averaged block-path
integrals zero.  Thus $G(B'',B'')=0$, whereas $d\theta$ lies in
$\ker H_k^{\rm cv}$.  If both properties hold, the same path
formula forces $\theta=0$.  This proves the direct-sum assertion.
\end{proof}

The following elementary calculation makes the disappearance of
the gauge-fixing form from curvature covariances explicit.  Suppose
$T:V\to E$ is linear, $q\ge0$, and
$V=\ker q\oplus\ker T$.  In any coordinates $i$ on $V$, let
$M_\tau$ represent
$\tau\eta^{-1}\langle Tv,Tw\rangle+q(v,w)$.
For $\tau\ne0$ this form is nondegenerate and
\begin{equation}\label{IV-44-19}
 \eta^{-1}TiM_\tau^{-1}(Ti)^*
       =\tau^{-1}\operatorname{proj}_{\operatorname{ran}T}.
\end{equation}
Indeed, the two summands of $V$ are orthogonal for this form.
Solving the equation with right side $\langle F,Tw\rangle$
first against $w\in\ker T$ removes the $\ker T$ component.
The remaining equation says that $Tv$ is
$\eta\tau^{-1}\operatorname{proj}_{\operatorname{ran}T}F$.

\begin{proposition}\label{IV-44-21}
The flat gap $\gamma_{\rm tr}>0$ of the one-step construction
satisfies $A_k^{\rm cv}\ge\gamma_{\rm tr}I$.  For $\tau\ne0$,
\begin{equation}\label{IV-44-22}
 \eta_k^{-1}H_k^{\rm cv}C_{\mathbf1,\mathbb R}
       (\tau B_k^{\rm cv}+\widehat G)^{-1}
       C_{\mathbf1,\mathbb R}^*(H_k^{\rm cv})^*
       =\tau^{-1}\mathcal C_{(k+1)}.
\end{equation}
At the trivial pair, the Hessian and covariance of a canonical
renormalization step are the matrices in \eqref{IV-44-16}, with
the identity on the $N^2-1$ colour components, where $N=2$
denotes the matrix size of the gauge group. In particular,
the curvature covariance of that step is
$\mathcal C_{(k+1)}\otimes I_{N^2-1}$.
\end{proposition}

\begin{proof}
The flat coercivity estimate from the one-step construction in
equation~\eqref{eq:cv-flat-gap} is
$\|dB\|^2+G(B,B)\ge\gamma_{\rm tr}\|\check B\|^2$
for $B=C_{\mathbf1,\mathbb R}\check B$.
Combining this with \eqref{IV-44-13} proves the gap.
Apply \eqref{IV-44-19} and \cref{IV-44-20} to get
\eqref{IV-44-22}.

For the operator identification, differentiate the constrained
critical-point equations at the trivial pair.  The differentiated
background is the Landau response in
\eqref{eq:cv-landau-identification}.  The action has zero first
variation there; its second variation is consequently the squared
norm of this linearized curvature, namely
$\eta_k^{-1}(H_k^{\rm cv})^*H_k^{\rm cv}$ in level-$k$
coordinates.  The linearized fluctuation constraint is exactly
$Q^{\rm lin}B=0$, and the second variation of the gauge term is
$G$.  Restriction by the pivot parametrization therefore gives
$A_k^{\rm cv}\otimes I_{N^2-1}$, whose inverse is the whole
flat covariance.  These are the constrained Hessian and
fluctuation-coordinate identities proved in
Proposition~\ref{II-40-99} and equation~\eqref{II-43-105}; no
continuation away from the trivial pair is used here.
\end{proof}

This identification is used only where the canonical step exists
with its prescribed constant-sequence background and octet
properties.  These hypotheses are part of the renormalization
construction; the finite-dimensional projection identity itself
does not establish them.

\subsection{Uniform kernel bounds}

Fix the constants $N_{\mathsf H}$, $C_{\mathsf H}$, and
$\delta_{\mathsf H}$ in the uniform free-response estimate of
equations~\eqref{I-82-31}--\eqref{II-56-6}.  Increase the first to $L^{m_{\rm tp}}$, with
$m_{\rm tp}\ge2$, and decrease $\delta_{\mathsf H}$ if needed.
Their choice is independent of the level and the volume.  At an
admitted side, \eqref{eq:cv-landau-identification} and that estimate
give, in level-$n$ units,
\begin{equation}\label{IV-44-32}
 |H_n^{\rm cv}(p,c)|
 \le C^{\rm k}_{\rm tp}\eta_n^2
              e^{-\delta_{\mathsf H}|z(p)-c_-|_\infty},
 \qquad
 |H_n^{\rm cv}B(p)|\le C^{\rm cv}_{\rm tp}\eta_n^2\|B\|_\infty.
\end{equation}
The first bound also holds at $n=0$ on every side because $H_0=d$.
The second follows by summing the first: if
$s(\delta)=e^\delta\sum_{u\in\mathbb Z^3}e^{-\delta|u|_\infty}$,
then $\sum_y e^{-\delta|x-y|_\infty}\le s(\delta)$ on every
torus, and we can take
$C^{\rm cv}_{\rm tp}=3C^{\rm k}_{\rm tp}s(\delta_{\mathsf H})$.

\begin{lemma}\label{IV-44-33}
If $d\mathbf Q^sB=0$, there is $B_2$ with
\[
 dB_2=dB,\qquad \mathbf Q^sB_2=0,\qquad
 \|B_2\|_\infty^2\le384L^5\|dB\|^2.
\]
\end{lemma}

\begin{proof}
First subtract a gradient and a constant as in \cref{IV-44-9}
to arrange $\mathbf Q^sB=0$.  In each block integrate $B$
from its centre along the ordered coordinate path, subtract the
mean of this site function in the block, and subtract its gradient
from $B$.  This preserves its straight average.  Write
$\sigma(c)$ for the sum of $|dB|$ over the plaquettes in the
double block belonging to a coarse bond $c$.
Exchanging two successive segments of a coordinate path changes
its integral by a sum of curvatures over the strip between them.
For an internal bond the resulting strips are disjoint and lie
in its block.  Its new value is therefore bounded by $\sigma(c)$.

For two neighbouring parallel bonds crossing the corridor of
$c$, the intervening plaquette and the two internal side bonds
give a difference at most $3\sigma(c)$.  The corridor has
diameter $2(L-1)$ in these moves, so any two such values differ
by at most $6(L-1)\sigma(c)$.  Each of the $L^3$ straight
paths averaged by $\mathbf Q^s$ crosses the corridor once and
otherwise has $L-1$ internal bonds.  Every corridor bond is used
$L$ times.  Since the average is zero, the corridor mean has
modulus at most $(L-1)\sigma(c)$.  Every corridor value is thus
bounded by $7(L-1)\sigma(c)$.  These and the internal bonds
exhaust the lattice bonds.  There are at most $6L^3$ plaquettes
in a double block, so
$\sigma(c)^2\le6L^3\|dB\|^2$; finally $49\cdot6<384$.
\end{proof}

For an orthogonal projection $P$,
$P(p,p)=\sup_{\|F\|=1,F\in\operatorname{ran}P}|F(p)|^2$
and $|P(p,q)|^2\le P(p,p)P(q,q)$.  Applying
\cref{IV-44-33} to $F=H_k^{\rm cv}B$ in
$\mathsf{Bg}_k\cap\mathsf{Fl}_{k+1}$, then using
\eqref{IV-44-13} and \eqref{IV-44-32}, gives
\begin{equation}\label{IV-44-34}
 \mathcal C_{(k+1)}(p,p)
       \le384L^5(C^{\rm cv}_{\rm tp})^2\eta_k^3
\end{equation}
when $k\ge1$ and its side is admitted.  At $k=0$ the bound by
$1$ suffices.  On an admitted side $N$ at level $n$ one also has
\[
 \Pi_n(p,p)\le\tfrac1{16}N^2(C^{\rm cv}_{\rm tp})^2\eta_n^3.
\]
To see this, choose the least-norm solution $dB=Q^{\rm pl}_nF$.
It is orthogonal to the cocycles.  In Fourier variables its
nonzero modes satisfy
$\sum_{\mu<\nu}|k_\mu\widehat B_\nu-k_\nu\widehat B_\mu|^2
=|k|^2|\widehat B|^2$, where $k_\mu=e^{i\theta_\mu}-1$.
The least nonzero $|k|$ is at least $4/N$; hence
$\|B\|\le(N/4)\|dB\|$.  Now use
\eqref{eq:curvature-average-norm} and \eqref{IV-44-32}.

\begin{proposition}\label{IV-44-35}
There is $C^{\rm bg}_{\rm tp}$ depending only on the fixed
free-response constants and $L$ such that, on every side,
\begin{equation}\label{IV-44-36}
 |C^{\rm bg}_{K,s}(p,q)|
          \le C^{\rm bg}_{\rm tp}(a_K/a_s)^3.
\end{equation}
\end{proposition}

\begin{proof}
Put $n=K-s$ and write $\ell=2L^m$.  Relabel the same bare
periodic graph as $(2L^{m_{\rm tp}},\widetilde K)$, where
$\widetilde K=m+K-m_{\rm tp}$.  Straight averages and the
projections through any common level are unchanged under this
relabeling.  If $\widetilde K\le0$, then $n\le m_{\rm tp}$,
and $\Pi_n(p,p)\le1$ proves the result after increasing the
constant.  Otherwise set $\widetilde n=\min(n,\widetilde K)$.
All levels through $\widetilde K$ on the relabeled lattice
have admitted sides, and
\[
 \Pi_n(p,p)\le\widetilde\Pi_{\widetilde n}(p,p)
 =\widetilde\Pi_{\widetilde K}(p,p)
       +\sum_{j=\widetilde n+1}^{\widetilde K}
                                       \widetilde{\mathcal C}_{(j)}(p,p).
\]
The preceding diagonal bounds and a geometric sum bound this by
$(C^{\rm cv}_{\rm tp})^2(768L^5+L^{2m_{\rm tp}}/4)
L^{-3\widetilde n}$; the case $\widetilde n=0$ is covered
by the bound $1$.  Since $n-\widetilde n\le m_{\rm tp}$,
we can take this constant times $L^{3m_{\rm tp}}$ for
$C^{\rm bg}_{\rm tp}$.  Projection Cauchy--Schwarz gives
the off-diagonal bound.
\end{proof}

\begin{proposition}\label{IV-44-38}
There are $C_{\rm tp},C^{\rm row}_{\rm tp}<\infty$ and
$\delta_{\rm tp}>0$, depending only on the fixed free-response
constants and $L$, such that for the physical sup distance
$r=|x(p)-x(q)|_\infty$,
\begin{equation}\label{IV-44-39}
 |\mathcal C_{(j)}(p,q)|
 \le C_{\rm tp}(a_K/a_{K-j+1})^3
                       e^{-\delta_{\rm tp}r/a_{K-j+1}}.
\end{equation}
If $\delta_{\rm tp}L(L-1)r\ge(3\log L+\log2)a_s$, then
\begin{equation}\label{IV-44-40}
 |C^{\rm fl}_{K,s}(p,q)|
 \le2L^3C_{\rm tp}(a_K/a_s)^3e^{-\delta_{\rm tp}Lr/a_s}.
\end{equation}
For each fixed $t_0>0$, the same bound holds when $r\ge t_0a_s$
with $2L^3C_{\rm tp}$ replaced by
\[
 C^{\rm fl}_{\rm tp}(t_0)
   =C_{\rm tp}\sum_{i\ge1}L^{3i}
                  e^{-\delta_{\rm tp}t_0(L^i-L)}<\infty.
\]
At $k=0$, or at an admitted side of level $k$,
\begin{equation}\label{eq:cv-inverse-row}
 |(A_k^{\rm cv})^{-1}(b,b')|
   \le\frac2{\gamma_{\rm tr}}e^{-\delta_{\rm tp}|b_--b'_-|_\infty},
 \qquad
 \sum_b|(H_k^{\rm cv}C_{\mathbf1,\mathbb R})(p,b)|
          \le C^{\rm row}_{\rm tp}\eta_k^2.
\end{equation}
The distances in \eqref{eq:cv-inverse-row} are in level-$k$ units.
\end{proposition}

\begin{proof}
We give the inverse estimate and keep track of which bounds are
uniform in the side.  By \eqref{IV-44-32} and the exponential
lattice sum,
\[
 |\eta_k^{-1}(H_k^{\rm cv})^*H_k^{\rm cv}(c,c')|
 \le3(C^{\rm k}_{\rm tp})^2s(\delta_{\mathsf H}/2)
                          e^{-\delta_{\mathsf H}|c_--c'_-|_\infty/2}.
\]
The gauge matrix has entries at most $L^3$ and vanishes unless
the two bonds lie in one block.  A column of
$C_{\mathbf1,\mathbb R}$ has at most seven nonzero entries,
of modulus at most $2L^2$, within distance $2L$ of its bond.
Thus $|A_k^{\rm cv}(b,b')|\le B_{\rm tp}e^{-\delta\rho(b,b')}$
for fixed $B_{\rm tp}$ and $\delta=\delta_{\mathsf H}/2$.
Also $\sup_b\sum_{b'}e^{-\delta\rho(b,b')/2}\le
S_{\rm tp}=3s(\delta_{\mathsf H}/4)$.

For completeness, conjugate $A=A_k^{\rm cv}$ by the diagonal
weight $e^{\epsilon\rho(b,b_0)}$.  For $0\le\epsilon\le\delta/2$,
convexity gives
\[
 e^{\epsilon\rho}-1
 \le\frac{2\epsilon}{\delta}(e^{\delta\rho/2}-1).
\]
The triangle inequality therefore bounds the change in each
matrix entry by
$(2\epsilon/\delta)B_{\rm tp}e^{-\delta\rho(b,b')/2}$.
The row and column sum bound gives a change in operator norm
of at most $2\epsilon B_{\rm tp}S_{\rm tp}/\delta$.  Choose
\[
 \delta_{\rm tp}
   =\min\{\delta/2,\gamma_{\rm tr}\delta/(4B_{\rm tp}S_{\rm tp})\}.
\]
The gap in \cref{IV-44-21} and the Neumann series give norm at
most $2/\gamma_{\rm tr}$ for the conjugated inverse.  Its
$(b,b_0)$ entry proves the first bound in
\eqref{eq:cv-inverse-row}.  The pivot rows have sum at most
$1+12L^5$; summing \eqref{IV-44-32} gives its second bound.

Insert the inverse estimate in \eqref{IV-44-22}.  The pivot
rows move bonds by at most $2L$ and have bounded sums, whereas
the two response kernels contribute $\eta_k^4$.  The prefactor
$\eta_k^{-1}$ leaves $\eta_k^3$.  The triangle inequality
extracts $e^{-\delta_{\rm tp}|z(p)-z(q)|_\infty}$, and the
remaining two bond sums converge uniformly by
$\delta_{\rm tp}\le\delta_{\mathsf H}/4$.
The change from base points to barycentres costs a fixed factor.
This proves \eqref{IV-44-39} at admitted sides and at $k=0$.
At the other sides the level-$k$ distance is at most
$N_{\mathsf H}/2$; \cref{IV-44-35} and
$\mathcal C_{(k+1)}(p,p)\le\Pi_k(p,p)$ prove the same
bound after a fixed increase of $C_{\rm tp}$.

Finally write $t=r/a_s$ in the sum in \eqref{IV-44-10}.
Its terms are bounded by
$C_{\rm tp}(a_K/a_s)^3L^{3i}e^{-\delta_{\rm tp}tL^i}$,
$i\ge1$.  Under the displayed condition their successive
ratios are at most $1/2$, which proves \eqref{IV-44-40}.
For $t\ge t_0$ extract $e^{-\delta_{\rm tp}Lt}$ instead;
the remaining convergent series is the stated constant.
\end{proof}

\section{The dipole kernel and its lattice approximation}
\label{sec:dipole-comparison}

For $u\ne0$ put
\begin{equation}\label{IV-45-2}
 \mathbb K_3(u)=\frac{3\widehat u\widehat u^{\mathsf T}-I}
                         {4\pi|u|^3}
               =\nabla^2\frac1{4\pi|u|},\qquad
 \mathcal K_{\rm free}(f,h)
       =\iint f(x)h(y)|x-y|^{-6}\,dx\,dy.
\end{equation}
On a torus the distance in the integral is its least-lift distance.
We will only use pairs supported together in a ball of radius
$\ell/4$, where the integral agrees with the integral of their
compactly supported lifts.  The eigenvalues of
$3\widehat u\widehat u^{\mathsf T}-I$ are $2,-1,-1$, so
\begin{equation}\label{eq:dipole-square}
 \sum_{\lambda,\lambda'}\mathbb K_3(u)_{\lambda\lambda'}^2
                          =\frac3{8\pi^2}|u|^{-6}.
\end{equation}

\subsection{The free curvature kernel}

Represent a plaquette by $(z,\lambda)$, with $\lambda$ its dual
coordinate direction, choosing the cyclic orientation.  Its
barycentre is $x(z,\lambda)=z+\frac12(e_1+e_2+e_3-e_\lambda)$
in bare units.  Relative to the original positive plaquette
orientation let the sign be $\epsilon_p\in\{\pm1\}$.
For $q_\mu(k)=e^{ik_\mu}-1$ define
\[
 F_{\lambda\lambda'}(k)
      =\frac{\overline{q_\lambda(k)}q_{\lambda'}(k)}{|q(k)|^2},
 \qquad P(k)=I-F(k)\quad(k\ne0),\qquad P(0)=0.
\]
Let $\Pi_\infty$ and $\Pi_N$ be the curl projections on
$\mathbb Z^3$ and $(\mathbb Z/N\mathbb Z)^3$.  Their Fourier
formulas are
\begin{equation}\label{IV-45-17}
 \begin{split}
 \Pi_\infty((z,\lambda),(z',\lambda'))
  &=\delta_{pq}-\int_{[-\pi,\pi]^3}
         e^{ik\cdot(z-z')}F_{\lambda\lambda'}(k)\frac{dk}{(2\pi)^3},\\
 \Pi_N((z,\lambda),(z',\lambda'))
  &=\delta_{pq}-N^{-3}\delta_{\lambda\lambda'}
       -N^{-3}\sum_{k\in\frac{2\pi}N\mathbb Z^3\cap(-\pi,\pi]^3
                         \setminus\{0\}}
                         e^{ik\cdot(z-z')}F_{\lambda\lambda'}(k).
 \end{split}
\end{equation}
Indeed, curl has symbol $q\times$, whose range is the Hermitian
orthogonal complement of $\bar q$.  Its orthogonal projection is
$P(k)$.  On the infinite lattice, cutting out $|k|\le\varepsilon$
in the right inverse
$-|q|^{-2}\bar q\times P(k)$ shows that the range of this
multiplier is the closure of the curl range.  The cutoff tends
to the identity on that range in $L^2$.  Fourier inversion proves
\eqref{IV-45-17}.  The finite formula has $P(0)=0$ because
an exact cochain has no constant mode.  In particular,
$\epsilon_p\epsilon_q C^{\rm free}_K(p,q)=\Pi_N(p,q)$
after rescaling the bare lattice to spacing one, with $N=\ell/a_K$.

\begin{lemma}\label{IV-45-22}\label{IV-45-27}
There are absolute constants $C'_{\rm lat},C''_{\rm lat}$ such
that, for distinct plaquettes on $\mathbb Z^3$,
\begin{equation}\label{eq:infinite-dipole-error}
 |\Pi_\infty(p,q)-\mathbb K_3(x(p)-x(q))_{\lambda(p)\lambda(q)}|
                 \le C'_{\rm lat}|x(p)-x(q)|^{-5}.
\end{equation}
If $p,q$ are distinct plaquettes on the $N$-torus and
$\widetilde p,\widetilde q$ are lifts whose barycentre difference
has sup norm at most $N/2$, then
\begin{equation}\label{eq:periodic-dipole-error}
 |\Pi_N(p,q)-\Pi_\infty(\widetilde p,\widetilde q)|
                                      \le C''_{\rm lat}N^{-3}.
\end{equation}
\end{lemma}

\begin{proof}
Fix the two dual directions, put $y=z-z'$ and
$s=(e_{\lambda'}-e_\lambda)/2$, so that $u=y+s$ is the
barycentre difference.  Multiplication by the barycentre phase
removes the linear lattice error:
\[
 m(k)=e^{-ik\cdot s}F_{\lambda\lambda'}(k)
       =\frac{4\sin(k_\lambda/2)\sin(k_{\lambda'}/2)}
                    {4\sum_\mu\sin^2(k_\mu/2)},\qquad
 m^0(k)=\frac{k_\lambda k_{\lambda'}}{|k|^2}.
\]
Taylor's formula for sine and its derivatives, and
$\sum_\mu4\sin^2(k_\mu/2)\ge4|k|^2/\pi^2$ on
$[-\pi,\pi]^3$, give, for $|\alpha|\le6$,
\[
 |\partial^\alpha F(k)|\le C_\alpha|k|^{-|\alpha|},\qquad
 |\partial^\alpha(m-m^0)(k)|\le C_\alpha|k|^{2-|\alpha|}.
\]
These derivative bounds follow by differentiating the quotient:
the numerator of the second difference vanishes to order six
before division by $|k|^2\sum_\mu4\sin^2(k_\mu/2)$,
and differentiating an inverse preserves its corresponding
homogeneous bound.

Choose a radial smooth cutoff $\chi$ equal to one on $|k|\le1/2$
and zero on $|k|\ge1$.  Split the Fourier coefficient of $F$
into $(1-\chi)F$, $\chi(m-m^0)$ with phase $e^{ik\cdot u}$,
and $\chi m^0$ with that phase.  The first is a smooth periodic
symbol; five integrations by parts give $O(|u|^{-5})$.
For the second, use annuli of radius $2^{-j}$.  Their Fourier
integrals are bounded both by $C2^{-5j}$ and, after six
integrations by parts, by $C|u|^{-6}2^j$.  Splitting the sum
at $2^j=|u|$ gives $O(|u|^{-5})$.

For the last term put
$E(u)=\int e^{ik\cdot u}\chi(k)|k|^{-2}\,dk/(2\pi)^3$.
The heat-kernel identity for $|k|^{-2}$, with Fubini justified
by the integrability at zero and the rapid decay of $\check\chi$,
gives $E=(4\pi|\cdot|)^{-1}*\check\chi$.  Write
$\check\chi(v)=\varphi(|v|)$.  This function is radial,
$4\pi\int_0^\infty\varphi(t)t^2dt=1$, and
$|\varphi(t)|\le C(1+t)^{-7}$.  The mean of
$|u-v|^{-1}$ over $|v|=t$ is $1/\max(|u|,t)$, hence
\[
 E(u)=\frac1{4\pi|u|}-D(|u|),\qquad
 D(r)=\int_r^\infty\varphi(t)t^2(r^{-1}-t^{-1})\,dt.
\]
Direct differentiation gives $D'(r)=O(r^{-6})$ and
$D''(r)=O(r^{-7})$.  The Fourier transform of $\chi m^0$
is $-\partial_\lambda\partial_{\lambda'}E$; it therefore
equals $-\mathbb K_3(u)_{\lambda\lambda'}+O(|u|^{-7})$.
The minus sign in \eqref{IV-45-17} proves
\eqref{eq:infinite-dipole-error} for $|u|>4$.
For smaller nonzero barycentre differences, $|u|\ge2^{-1/2}$,
and the projection bound by one covers the remaining cases.

For the periodic comparison take $N\ge4$ and remove the ball
$|k|\le N^{-1}$ from $F$, using a smooth cutoff supported
in $|k|\le2N^{-1}$ and its $2\pi$-periodic translates.
Call the resulting smooth periodic symbol $F_N$.  It equals
$F$ at every nonzero Fourier grid point and is zero at zero.
The derivative bounds above give
$\int|\partial_\mu^4F_N|\le CN$; thus
$|\check F_N(y')|\le CN|y'|_\infty^{-4}$ for $y'\ne0$.
Its Fourier series is absolutely summable, and summing it over
the grid gives the exact identity
\[
 N^{-3}\sum_{k\ne0}e^{ik\cdot y}F(k)
                         =\sum_{n\in\mathbb Z^3}\check F_N(y+Nn).
\]
The term $n=0$ differs from $\check F(y)$ by at most $CN^{-3}$,
the volume of the removed ball.  Since
$|y|_\infty\le N/2+1/2$, for $n\ne0$ one has
$|y+Nn|_\infty\ge N|n|_\infty/4$.  The other terms are
bounded in sum by
$CN^{-3}\sum_{n\ne0}|n|_\infty^{-4}<\infty$.
The zero-mode correction in \eqref{IV-45-17} is also $O(N^{-3})$.
This proves \eqref{eq:periodic-dipole-error}; $N\le3$ follows
by increasing the constant and using the projection bound.
\end{proof}

Rescaling the last lemma gives, for a least lift $u$ of the
physical barycentre difference of distinct plaquettes,
\begin{equation}\label{IV-45-31}\equationalias{IV-45-30}
 \big|\epsilon_p\epsilon_q C^{\rm free}_K(p,q)
           -a_K^3\mathbb K_3(u)_{\lambda(p)\lambda(q)}\big|
   \le C_{\rm lat}a_K^3(a_K^2|u|^{-5}+\ell^{-3}),
\end{equation}
and consequently
\begin{equation}\label{IV-45-32}
 |C^{\rm free}_K(p,q)|\le(1+3C_{\rm lat})a_K^3|u|^{-3}.
\end{equation}
Here $C_{\rm lat}$ is absolute and
$2^{-1/2}a_K\le|u|\le\sqrt3\ell/2$.

\subsection{Separated test pairs and squared kernels}

For the test pairs $\mathfrak T_E(d;\vartheta)$ defined in the
introduction, the free kernel has the following useful bounds.
The two supports lie together in a ball of radius $5d$;
all their cross distances belong to $[d,10d]$.  Thus
\begin{equation}\label{IV-45-9}
 10^{-6}\vartheta^2\|f\|_\infty\|h\|_\infty
 \le\mathcal K_{\rm free}(f,h)
 \le(4\pi/3)^2\|f\|_\infty\|h\|_\infty.
\end{equation}
For a plaquette kernel $k$ define the squared-kernel sum
\begin{equation}\label{IV-45-34}
 \mathcal K[k](f,h)
   =\sum_{p,q}f(x(p))h(x(q))k(p,q)^2,\qquad
 \mathcal K_{K,s}=\mathcal K[C^{\rm bg}_{K,s}].
\end{equation}
There are no additional mesh factors in this sum.

\begin{proposition}\label{IV-45-38}
Suppose $8a_K\le d$, $20d\le\ell$, and
$(f,h)\in\mathfrak T_{T_\ell}(d;\vartheta)$.  Then
\begin{equation}\label{IV-45-39}
 \left|\mathcal K[k](f,h)-\frac3{8\pi^2}
                                  \mathcal K_{\rm free}(f,h)\right|
 \le C_{\rm tm}\left[\frac{a_K}{\vartheta d}
                      +(d/\ell)^3+e\right]
                                  \|f\|_\infty\|h\|_\infty
\end{equation}
for $k=C^{\rm free}_K$ and $e=0$.  The same estimate holds
for $k=C^{\rm bg}_{K,s}$ with
\[
 t=d/a_s,\qquad e=t^3e^{-\delta_{\rm tp}Lt/\sqrt3},
\]
provided $t\ge1$, or provided
$\delta_{\rm tp}L(L-1)t/\sqrt3\ge3\log L+\log2$.
The constant is independent of the mesh, side, pair, $s$, and
$\vartheta$; the free case admits an absolute constant, and
the background case depends only on the constants of
\cref{IV-44-38}.
\end{proposition}

\begin{proof}
Lift both functions to their common ball of radius $5d$.
Their cross distances are unchanged.  For a fixed direction
the barycentres form a translate of $a_K\mathbb Z^3$.
Disjoint cubes of side $a_K$ about the contributing points
lie in a ball of radius $d+\sqrt3a_K/2$, so there are at most
$18(d/a_K)^3$ contributing plaquettes for each function.
If on their product
$|k_1-k_2|\le A_1a_K^3d^{-3}$ and
$|k_2|\le A_2a_K^3d^{-3}$, it follows that
\begin{equation}\label{IV-45-42}
 |\mathcal K[k_1]-\mathcal K[k_2]|
 \le324A_1(2A_2+A_1)\|f\|_\infty\|h\|_\infty.
\end{equation}
Use \eqref{IV-45-31} to replace $C^{\rm free}_K$ by
$a_K^3\epsilon_p\epsilon_q\mathbb K_3(x(p)-x(q))_{
\lambda(p)\lambda(q)}$.  Equations
\eqref{IV-45-31}--\eqref{IV-45-42} bound the error by a
constant times $[(a_K/d)^2+(d/\ell)^3]\|f\|_\infty\|h\|_\infty$.

The resulting sum is a Riemann sum over nine products of
barycentre grids.  On every product of cubes that contributes,
all cross distances are at least $d-2\sqrt3a_K\ge d/2$.
The squared dipole entry is bounded by $Cd^{-6}$ and its
gradient by $Cd^{-7}$.  The Lipschitz bounds on $f,h$ thus
bound the difference between the integrand and its grid value
by $Ca_K\vartheta^{-1}d^{-7}\|f\|_\infty\|h\|_\infty$.
The union of contributing cubes for either function has volume
at most $8d^3$.  Integration and summation over the nine
directions give error at most
$8\cdot10^4a_K(\vartheta d)^{-1}\|f\|_\infty\|h\|_\infty$.
Equation \eqref{eq:dipole-square} identifies the integral.

For the background kernel,
$C^{\rm bg}_{K,s}-C^{\rm free}_K=-C^{\rm fl}_{K,s}$.
All contributing sup distances are at least $d/\sqrt3$.
The appropriate bound in \cref{IV-44-38} gives
$|C^{\rm fl}_{K,s}(p,q)|\le Ce\,a_K^3d^{-3}$.
Since $t^3e^{-\delta_{\rm tp}Lt/\sqrt3}$ is bounded for
$t>0$, \eqref{IV-45-42} and \eqref{IV-45-32} add an error
at most $Ce\|f\|_\infty\|h\|_\infty$.  This proves the result.
\end{proof}

\section{Differentiating the small-field iteration}
\label{sec:weighted-small-field}

We next introduce the two plaquette weights into the actual
renormalization maps.  The counterterms and the maps themselves
are kept fixed.  Differentiation then separates two effects:
the covariance created in a single step and the transport of
derivatives created at earlier steps.

Use the canonical one-lattice datum of \cref{II-203-9}, including its
absolute radii, bracket localization, and sum rule.  Throughout
this section it is assumed to satisfy its full small-coupling
hypotheses.  Thus every level under consideration is an admitted
place; at heights at least two this includes the octet property
on all the lattices in its defining class.  Write $\alpha_0,\alpha_1$
for the common output radii, $\kappa$ for the polymer decay rate,
and $\omega_n=(n+1)L^{-(n-1)}$ for the array weights.
Use the admissible-radii row of equation~\eqref{eq:fourth-weight-cases}:
our constant $K_\Lambda$ is $(1+\beta_{\rm s})^4$ times the
constant denoted by $K_\Lambda$ there.  The normalized
fourth-potential size is
\[
 \nu(\mathbb V')=(K_\Lambda s_j)^{-1}
       \sup_{Y,\mathfrak u}e^{\lambda_Vd_j(Y)}|\mathbb V'(Y;\mathfrak u)|.
\]
Choose a uniform ordinary-array ceiling
$N_\star\ge3(E_0+\tfrac32\alpha_0^2B_\sharp)$ and impose
\[
 s_j\le s^\star_\omega
 :=\min\{s^\star,w_\star/(4K_\Lambda N_\star)\},
\]
as in equation~\eqref{eq:fourth-ceiling}.
Indeed $\sum_{n\ge1}\omega_n\le3$, so the ordinary progression
has this weighted array bound even if its original iteration used
the smaller weights $L^{-(n-1)}$.  The smaller ceiling is imposed
in the final coupling choice.  The constant $K_{\mathbf M}$ is
evaluated with this weighted $K_\Lambda$.

Fix a lattice and an integer $k<K$ through which these places
exist.  For now the calibrations $t_j\ge\gamma^{-2}$ may be
independent, and the coefficients obey
$|x_j/t_j-1|<\sigma_\star$.  The ordinary progression is
$(E^{(j)},e_j)$, with $\|E^{(j)}\|_{j,\kappa}\le E_0$
and $|e_j|\le B_\sharp$.  Write $\|\cdot\|_{j,+}$ for
the same pointed-family norm at rate
$\kappa^+=(1+4\beta_s)\kappa$, and $N_j$ for the bare-array
norm at height $j$.  The localized step map is $\mathsf{St}_j$.
The fourth-potential bound and \eqref{III-34-24} give
\begin{equation}\label{IV-64-20}
 \|\mathsf{St}_j(f)-\mathsf{St}_j(f')\|_{j+1,+}
        \le\theta_jN_j(f-f'),\qquad
 s_j=\frac{(1+\beta_s)\delta(t_j)}{\alpha_0},\quad
 \theta_j=K_{\mathbf M}s_j,\quad
 N_{\star,j}=\frac{w_\star}{4K_\Lambda s_j},
\end{equation}
for arrays of size at most $N_{\star,j}$.  The map is jointly
holomorphic with the field variables, real on real data, and
its member on $X$ depends only on the restriction of the input
array to $X$.  Bracket localization is required for this last
assertion.  These statements apply with the same canonical
provision and counterterms at every parameter value.
To check the coefficient, each fourth-potential norm is at most
$w_\star/4$, their difference is at most $w_\star/2$, and the
sum of the intrinsic norms is at most $3w_\star/8$.
The cluster difference bound therefore costs at most
$162L^3K_\Pi K_\Lambda s_jN_j(f-f')$ at rate
$\lambda_{\rm out}=\kappa^+$.  This is bounded by the right
side of \eqref{IV-64-20}, since
$K_{\mathbf M}\ge243L^3K_\Pi K_\Lambda$.
The existence, joint holomorphy and actual-input identification
are those of Theorems~\ref{II-53-39}--\ref{II-127-29}; locality follows
through the bracket localization before applying the cluster map.

At $N_j(f)\le N_{\star,j}/2$, Cauchy's estimate and
\eqref{IV-64-20} give
\begin{equation}\label{IV-64-21}
 \|D\mathsf{St}_j(f)g\|_{j+1,+}\le\theta_jN_j(g),\qquad
 \|D^2\mathsf{St}_j(f)[g_1,g_2]\|_{j+1,+}
       \le\frac{2\theta_j}{N_{\star,j}}N_j(g_1)N_j(g_2).
\end{equation}
For the second inequality apply the first to $g_1$ at
$f+zg_2$, and use the circle of radius
$N_{\star,j}/(2N_j(g_2))$ in $z$.  A limit from smaller
circles covers the boundary of the stated half-ball.
For holomorphic arrays the ordinary chain rule is valid in
these norms: it follows by the Cauchy formula for each local
member, with the uniform exponentially weighted bounds.

\subsection{Existence, locality, and the first derivatives}

Let $w_1,w_2$ be real nonzero arrays on the bare plaquettes,
$w_\lambda=\lambda_1w_1+\lambda_2w_2$, and
$w^{[j]}_\lambda=L^{K-j}w_\lambda$.
Use the weight member $\mathsf W_{w^{[j]}_\lambda}$ in the
slot $\{j\}$ of the bare array.  At height zero use the
bracketed plaquette action, with the same weight, in the fourth
place of the first step.  Define
\begin{equation}\label{IV-64-28}
 f_{\lambda,j}
    =((E^{(i)}_\lambda,e_i)_{1\le i\le j},\emptyset)
                             \sqcup\mathsf W_{w^{[j]}_\lambda},
 \qquad
 E^{(j+1)}_\lambda
    =E^{(j+1)}+\mathsf{St}_j(f_{\lambda,j})
                         -\mathsf{St}_j(f_{0,j}).
\end{equation}
The first output has the identical difference formula with the
fourth-place family.  No source-dependent counterterm is
subtracted in this definition.

Here are explicit parameter domains for the iteration.  Put
\begin{gather*}
 m_{j,a}=12M^3\alpha_0^2L^{K-j}|w_a|_\infty,\qquad
 c^\flat=81C_P^2C_{\mathcal A}^2/K_\Lambda,\qquad
 d_{1,a}=c^\flat\theta_0m_{0,a},\\
 n_{j,a}=m_{j,a}+\sum_{i=1}^j\omega_{j+1-i}d_{i,a},\qquad
 d_{j+1,a}=\theta_jn_{j,a},\\
 N^{\rm rm}_j=N_{\star,j}
       -\Big(\sum_{n=1}^j\omega_n\Big)
                            (E_0+\tfrac32\alpha_0^2B_\sharp).
\end{gather*}
The recursion for $n_{j,a}$ and $d_{j+1,a}$, and the definition
of $N^{\rm rm}_j$, use $1\le j\le k$.
The canonical smallness hypotheses make $N^{\rm rm}_j>0$.
Set
\begin{equation}\label{IV-64-12}
 \Lambda_0=\{\lambda:c^\flat\sum_a|\lambda_a|m_{0,a}<N_{\star,0}\},
 \qquad
 \Lambda_j=\{\lambda\in\Lambda_{j-1}:
                       \sum_a|\lambda_a|n_{j,a}<N^{\rm rm}_j\}.
\end{equation}

\begin{proposition}\label{IV-64-34}
The weighted progression exists through $j+1$ on $\Lambda_j$.
Each $\Lambda_j$ is an open convex neighbourhood of zero,
invariant under conjugation and coordinatewise contraction.
The local members are jointly holomorphic in the parameters
and the field variables, and
\begin{equation}\label{IV-64-35}
 \|E^{(j)}_\lambda-E^{(j)}_{\lambda'}\|_{j,+}
     \le\sum_a|\lambda_a-\lambda'_a|d_{j,a}.
\end{equation}
They depend on $w_\lambda$ only through its values on the
plaquettes in their domain $X$.  Thus
$\dot E^{(j)}_a=\partial_aE^{(j)}_\lambda|_0$ vanishes unless
$X$ meets $\operatorname{supp}w_a$, and
$\ddot E^{(j)}=\partial_1\partial_2E^{(j)}_\lambda|_0$
vanishes unless it meets both supports.  Their norms and
gauge-invariance properties are those of pointed families;
in particular $\|\dot E^{(j)}_a\|_{j,+}\le d_{j,a}$.
If $x_0,\ldots,x_{j-1}$ are real, all source derivatives of
$E^{(j)}_\lambda$ at zero are real on real configurations.
The added action on a bare configuration is exactly
\begin{equation}\label{IV-64-33}
 A^{\eta_j}(w^{[j]}_\lambda,U)
                     =a_K^{-1}\sum_pw_\lambda(p)a_p(U),
\end{equation}
independent of $j$.  The canonical exact sum rule applies
to \eqref{IV-64-28} at every admitted block field.
\end{proposition}

\begin{proof}
The first-step bracket is linear in the weight and supported
on the unique elementary domain assigned to a plaquette.
In its chart the plaquette displacement is bounded by
$\rho\le2C_PC_{\mathcal A}s_0\alpha_0$.
The bracketed plaquette action is bounded by
$19\rho^2+6\rho\alpha_0L^{-2}$, which is at most
$81C_P^2C_{\mathcal A}^2s_0\alpha_0^2$ because
$s_0\le1$ and $L\ge13$.
There are at most $12M^3$ plaquettes assigned to one
elementary domain.  Divide this raw member bound by
$K_\Lambda s_0$, as required by the normalized fourth-potential
size.  The resulting size is at most
$c^\flat\sum_a|\lambda_a|m_{0,a}$.
The first-step theorem gives existence on $\Lambda_0$,
its Lipschitz estimate with $d_{1,a}$, and locality.
The chart restrictions used here are
$\alpha_\iota\le a_P^\circ$ and
$C_{\mathcal A}\delta(t_0)\le r_P^\circ/2$;
these are supplied by the canonical place conditions.

Assume the assertion through level $j$.  The norm of a weight
member is at most $6M^3\alpha_0^2L^{K-j}|w_a|_\infty$.
Since its array weight is $\omega_1=2$, the induction gives
\[
 N_j(f_{\lambda,j}-f_{\lambda',j})
       \le\sum_a|\lambda_a-\lambda'_a|n_{j,a}.
\]
The ordinary progression has array norm at most the subtracted
term in $N^{\rm rm}_j$.  Thus $\lambda\in\Lambda_j$ puts
$f_{\lambda,j}$ strictly inside the ball of radius
$N_{\star,j}$, and \eqref{IV-64-20} gives existence and
\eqref{IV-64-35} at the next level.  Holomorphy follows
from that of the step.  Equality of two weights on $X$
gives equality of every earlier member used inside $X$;
locality of the step completes the locality induction.

The parameter domains are finite intersections of strict
inequalities $\sum_ac_a|\lambda_a|<c$, with $c>0$.
Cauchy's formula, uniform in the local core and weighted
by $e^{\kappa^+d_j(X)}$, proves the derivative assertions.
Gauge invariance passes through the derivatives. When
$x_0,\ldots,x_{j-1}$ are real, the conjugation identity for
$E^{(j)}_\lambda$ also passes through source derivatives at zero.
The identity \eqref{IV-64-33} follows from
$\eta_j^{-1}L^{K-j}=a_K^{-1}$ and the weight-member sum
rule.  Adding it to the existing action gives the input to
the canonical exact step.  Its bounded Borel hypothesis holds:
at height zero the added factor is the exponential of a finite
sum of continuous plaquette functions on a compact product
of groups, and at later heights it is the action of an admitted
bare array.  The canonical rule and interaction sum rules
therefore give the exact sourced step.
\end{proof}

For the rest of the section take real running coefficients,
$t_j=x_j$, set $y_j=\log x_j$, and assume
\begin{equation}\label{IV-64-14}
 |x_j-L^{-(j-i)}x_i|\le\bar B\quad(i\le j\le k),\qquad
 x_k\ge4\bar B,\qquad y_k\ge\max(2\mathfrak p_0,1+\log L).
\end{equation}
Put $C_0^\flat=\max(2,c^\flat)$,
$C_1^\flat=\max(2,(c^\flat)^2/8)$ and
\[
 S_A=(4/3)^{1/2}\sum_{n\ge1}(n+1)^{\mathfrak p_0+1}L^{1-n/2},
 \qquad
 S_B=(4/3)^{1/2}\sum_{n\ge0}(n+1)^{\mathfrak p_0}L^{-n/2}.
\]
If $C_0^\flat S_A\theta_k\le1$, the recursion above implies
\begin{equation}\label{IV-64-37}
 \begin{gathered}
 n_{j,a}\le2m_{j,a},\qquad
 d_{j+1,a}\le2\theta_jm_{j,a}\quad(1\le j\le k),\\
 d_{j,a}\le C_0^\flat\theta_{j-1}m_{j-1,a}
                       \quad(1\le j\le k+1).
 \end{gathered}
\end{equation}
To check this induction, \eqref{IV-64-14} gives
$\frac34L^{j-i}x_j\le x_i\le\frac54L^{j-i}x_j$ and
\[
 \delta(x_i)\le(4/3)^{1/2}L^{-(j-i)/2}
                    (j-i+1)^{\mathfrak p_0}\delta(x_j).
\]
Insert this and $m_{i,a}=L^{j-i}m_{j,a}$ in the convolution
with $\omega_{j+1-i}$.  Its sum is bounded by
$C_0^\flat S_A\theta_jm_{j,a}\le m_{j,a}$.
The same comparison gives
$\sum_{i\le j}\delta(x_i)\le S_B\delta(x_j)$.

\subsection{Separation of the mixed derivative}

Let $D_j$ be the least sup distance, in level-$j$ units,
between the corners of plaquettes in the two supports, and put
\[
 \nu_j=\frac{32K_{\mathbf M}K_\Lambda}{w_\star}
                          s_j^2m_{j,1}m_{j,2},\qquad
 T(u)=\sum_{n\ge1}(n+1)^{2\mathfrak p_0+1}
              e^{-4\beta_s\kappa[L^{n-1}u-2]_+}\quad(u>0).
\]
The series converges for $u>0$; set $T(0)=+\infty$.

\begin{proposition}\label{IV-64-43}
Assume \eqref{IV-64-14}, $C_0^\flat S_A\theta_k\le1$,
and $s_j\le s^\star_\omega/2$ for $1\le j\le k$.
With $\ddot d_j=\|\ddot E^{(j)}\|_{j,+}$, one has
\begin{gather}
 \|\ddot E^{(j)}\|_{j,\kappa}
       \le e^{-4\beta_s\kappa[D_j/M-2]_+}\ddot d_j,
 \label{eq:weighted-mixed-locality}\\
 \ddot d_{j+1}\le\theta_j\sum_{i=1}^j\omega_{j+1-i}
       e^{-4\beta_s\kappa[D_i/M-2]_+}\ddot d_i+\nu_j.
 \label{IV-64-44}
\end{gather}
If also $\frac43C_1^\flat L\theta_jT(D_j/M)\le1$
for $1\le j\le k$, then
$\ddot d_j\le C_1^\flat\nu_{j-1}$, and for
$1\le\bar k\le k+1$,
\begin{equation}\label{IV-64-45}
 \sum_{j\le\bar k}\sum_{X,z}\|\ddot E^{(j)}(X,\cdot,z)\|_X
 \le8C_1^\flat M^3C_{E,1}N_{\bar k}(w_1)
   \sum_{j\le\bar k}L^{3(\bar k-j)}\nu_{j-1}
                e^{-[L^{\bar k-j}D_{\bar k}/M-2]_+}.
\end{equation}
Here $N_{\bar k}(w_1)$ counts the level-$\bar k$ $M$-cubes
containing a corner of its support.
\end{proposition}

\begin{proof}
Locality forces a nonzero mixed member on $X$ to meet both
supports.  The diameter bound for a polymer gives
$D_j\le M(d_j(X)+2)$, proving
\eqref{eq:weighted-mixed-locality}.  The second derivative
of \eqref{IV-64-28} is
\[
 \ddot E^{(j+1)}
  =D^2\mathsf{St}_j(f_{0,j})[\dot f_{1,j},\dot f_{2,j}]
                         +D\mathsf{St}_j(f_{0,j})\ddot f_j.
\]
The weighted ceiling gives
\[
 N_j(f_{0,j})\le N_\star
 \le\frac{w_\star}{4K_\Lambda s^\star_\omega}
 \le\frac{w_\star}{8K_\Lambda s_j}
 =\tfrac12N_{\star,j}.
\]
Thus the ordinary array lies in the half-ball, and
\eqref{IV-64-37} gives
$N_j(\dot f_{a,j})\le2m_{j,a}$.
Equations \eqref{IV-64-21} and
\eqref{eq:weighted-mixed-locality} prove
\eqref{IV-64-44}.  The same Cauchy estimate at the first
step gives $\ddot d_1\le(c^\flat)^2\nu_0/8$.

For $n=j+1-i$, the running comparison gives
$\nu_{i-1}\le\frac43L^n(n+1)^{2\mathfrak p_0}\nu_j$
and $D_i=L^{n-1}D_j$.  If
$\ddot d_i\le C_1^\flat\nu_{i-1}$, the convolution in
\eqref{IV-64-44} is therefore at most
$\frac43C_1^\flat L\theta_jT(D_j/M)\nu_j\le\nu_j$.
This proves the induction.  Finally every contributing
domain meets one of at most
$L^{3(\bar k-j)}N_{\bar k}(w_1)$ level-$j$ cubes.
The anchored polymer sum, including its base points, is
bounded by $8M^3C_{E,1}e^{-D}$ when $d_j(X)\ge D$.
Use $D=[L^{\bar k-j}D_{\bar k}/M-2]_+$ and sum over $j$.
\end{proof}

\subsection{Values at a regular background}

The local-term estimate needed below is quadratic in the
curvature of the background.  Suppose $U$ is regular of
level $T$ at radius $\varepsilon$, all its lower-level pairs
belong to their cores, and $M\alpha_0\le\alpha_1$,
$\alpha_0+\alpha_1\le a^\sharp$.  Set
\begin{equation}\label{IV-64-16}
 d_*(j;T,\varepsilon)
 =\min\left\{\frac{c_A\alpha_0L^{2(T-j)}}{M\varepsilon},
                 \frac{L^{m+K-j}-4}{2M}\right\}-1.
\end{equation}
For a local term $E$ of level $j\le T$, the regular-gauge
estimate of \cref{I-28-75} and the quadratic Cauchy argument
in the proof of \cref{II-101-27} give
\begin{equation}\label{IV-64-39}
 |E(U)-E(\mathbf1)|
 \le\|E\|_X\left[
 C_A(\varepsilon/\alpha_0)^2L^{-4(T-j)}M^2(1+d_j(X))^2
       +2\mathbf1_{\{d_j(X)\ge d_*(j;T,\varepsilon)\}}\right].
\end{equation}
Indeed, lowering the level changes the regularity radius to
$\varepsilon L^{-2(T-j)}$.  The regular-gauge estimate
applies if $(1+d_j(X))M$ times that radius is at most
$c_A\alpha_0$ and $2(1+d_j(X))M+4<L^{m+K-j}$.
Failure of either inequality forces the indicated lower bound
on $d_j(X)$; the bounded local-term estimate then gives the
second summand.

If $d_*=d_*(T;T,\varepsilon)\ge3$, summing
\eqref{IV-64-39} for the first derivative families gives
\begin{equation}\label{IV-64-42}\equationalias{IV-64-41}
 \sum_{j\le T}\sum_{X,z}
 |\dot E^{(j)}_a(X,U,z)-\dot E^{(j)}_a(X,\mathbf1,z)|
 \le C_fN_T(w_a)[\varepsilon^2/\alpha_0^2+e^{-d_*}]
           \alpha_0\delta(x_{T-1})\frac{|w_a|_\infty}{a_{K-T}}.
\end{equation}
Here $C_f$ depends only on the fixed construction constants.
For clarity, the level-$j$ anchored polymer sum in the first
term costs $L^{3(T-j)}$ and its regularity factor gains
$L^{-4(T-j)}$, leaving $L^{-(T-j)}d_{j,a}$.
For the second term,
$d_*(j;T,\varepsilon)\ge L^{T-j}(d_*+1)-1$ and $d_*\ge3$
give
$L^{3(T-j)}e^{-d_*(j;T,\varepsilon)}
\le L^{-(T-j)}e^{-d_*}$.
Now use \eqref{IV-64-37} and
$\sum_{j\le T}\delta(x_{j-1})\le S_B\delta(x_{T-1})$.
This proves \eqref{IV-64-42} with a constant uniform in the
level and the volume.

\section{Covariance transport along the small-field history}
\label{sec:covariance-transport}

Use the canonical real progression and history decomposition,
with $x_j=L^{-1}x_{j-1}-e_j$.  A level is used only while
the canonical density tests and background sum rule are
available.  Denote by $h_t$ the history with no large-field
region through level $t$, by $\mathcal M$ the one-step block
map, and by $g_t$ its nonnegative step factor.  For bounded
Borel functions $F$ of the original field put
\[
 m_t[F]=\tau_{h_t}[F\rho_0],\qquad m_t=m_t[1],\qquad
 \Phi_t[F]=\frac{dm_t[F]}{dm_t}.
\]
The positive history measure gives
$|m_t[F]|\le\|F\|_\infty m_t$, so the derivatives exist
with $|\Phi_t[F]|\le\|F\|_\infty$.  All quotient identities
below are identities modulo $m_t$-null sets.

\subsection{The sourced density and the conditional derivatives}

On configurations whose pairs belong to all the cores through
level $t$, set
\begin{equation}\label{IV-65-8}
 D_t^\lambda(U)=\mathcal O(w_\lambda)(U)
        +\sum_{i=1}^t\sum_{X,z}
                    (E^{(i)}_\lambda-E^{(i)})(X,U,z),\qquad
 \mathcal O(w)=a_K^{-1}\sum_pw(p)a_p.
\end{equation}
Let $\mathcal V_t^\chi$ be the open set on which the full
small-field density test at level $t$ passes.  The canonical
test bounds place it inside the class of the exact sum rule;
its minimal background $U_t(W)$ exists, belongs to all the
lower cores, and is regular of level $t$ at radius
$2B_3\varepsilon_t$.  Local terms and the plaquette action
are constant on the gauge orbit defining this background.

\begin{proposition}\label{IV-65-25}
For $\lambda\in\Lambda_{t-1}$,
\begin{equation}\label{IV-65-27}
 m_t[e^{\mathcal O(w_\lambda)}]
       =\exp\{D_t^\lambda(U_t(W))\}\,m_t(dW).
\end{equation}
Both sides are carried by $\mathcal V_t^\chi$; the expression
on the right is set to zero off that set.  The measure $m_t$
has a strictly positive bounded Borel density there.
Consequently, for $F_a=\mathcal O(w_a)$,
\begin{equation}\label{IV-65-32}
 \Phi_t[F_a]=c_{t,a}+X_{t,a}+I_{t,a},\qquad
 \Phi_t[F_1F_2]-\Phi_t[F_1]\Phi_t[F_2]
                  =\sum_{i=1}^t\sum_{X,z}\ddot E^{(i)}(X,U_t,z),
\end{equation}
where
\[
 c_{t,a}=\sum_{i\le t}\sum_{X,z}\dot E^{(i)}_a(X,\mathbf1,z),
 \quad X_{t,a}=\mathcal O(w_a)(U_t),\quad
 I_{t,a}=\sum_{i\le t}\sum_{X,z}
       [\dot E^{(i)}_a(X,U_t,z)-\dot E^{(i)}_a(X,\mathbf1,z)].
\]
\end{proposition}

\begin{proof}
Let $\mathsf{cs}_j$ be the logarithm of the positive flat
normalizing factor of step $j+1$.  The exponent of the
sourced small-field density at height $t$ is
\[
 A_t^\lambda(U)=-x_0A^{\eta_0}(U)+\mathcal O(w_\lambda)(U)
           +\sum_{i\le t}\sum_{X,z}E^{(i)}_\lambda(X,U,z)
           +\sum_{j<t}\mathsf{cs}_j.
\]
Indeed, $A^{\eta_i}=L^iA^{\eta_0}$ and the coupling recursion
give $L^tx_t+\sum_{i\le t}L^ie_i=x_0$.
In the action of the bare array, the flat values
$E^{(i)}_\lambda(X,\mathbf1,z)$ cancel those in its vacuum
factor.  The remaining terms are exactly $A_t^\lambda$.
The first sourced step and then the exact sum rule in
\cref{IV-64-34}, applied inductively, show that
$m_t[e^{\mathcal O(w_\lambda)}]$ has density
$e^{A_t^\lambda(U_t(W))}$ on $\mathcal V_t^\chi$ and zero
elsewhere.  The sum rule also proves Borel measurability
through its chart integral; it does not require a global
continuous choice of $U_t$.  The local cores and the finite
number of local terms give boundedness.  At $\lambda=0$
every term of the exponent is real, and the density is
strictly positive.  Dividing it out gives
\eqref{IV-65-27}.

For every fixed finite lattice,
\[
 |D_t^\lambda(U_t(W))|
 \le\sum_a|\lambda_a|\left(2a_K^{-1}\sum_p|w_a(p)|
                          +\sum_{i\le t}\#\mathcal Q_i\,d_{i,a}\right),
\]
uniformly in $W$.  Cauchy's formula or dominated differentiation
therefore differentiates the measure identity at zero.
The first derivative of the exponential is the first derivative
of $D_t$, and its mixed derivative is
$\partial_{12}D_t+\partial_1D_t\partial_2D_t$.
Subtracting the product gives \eqref{IV-65-32}.
Difference quotients of the Borel density show that these
versions of the quotients are Borel.  The reality of the
derivatives follows from real source parameters.
\end{proof}

In particular, on the passing set,
\begin{equation}\label{eq:cov-transport-sizes}
 |X_{t,a}|\le X^*_{t,a}
       :=2B_3^2\varepsilon_t^2L^{-4t}a_K^{-1}\sum_p|w_a(p)|,
 \qquad |I_{t,a}|\le I^*_{t,a},
\end{equation}
where $I^*_{t,a}$ is the right side of \eqref{IV-64-42}
with $T=t$ and $\varepsilon=2B_3\varepsilon_t$.
The first bound is the plaquette estimate at a regular
background.  The second requires the hypotheses of
\eqref{IV-64-42}, including $d_*(t;t,2B_3\varepsilon_t)\ge3$.

\subsection{Splitting at the first large-field region}

Let $g_t^{\rm w}$ be the sum of the nonnegative step factors
over the outcomes different from the all-small-field outcome,
with input $h_{t-1}$, and put
\[
 w_t=Z^{-1}\int g_t^{\rm w}\,dm_{t-1},\qquad
 g_{j,T}=\prod_{j<i\le T}g_i\circ\mathcal M^{i-1-j},\qquad
 m_{j|T}=Z^{-1}g_{j,T}m_j.
\]
The factors depend only on the current field because the
earlier exposed-bond sets of $h_{t-1}$ are empty.
The history product formula therefore gives
\begin{equation}\label{IV-47-34}
 m_t[F]=\mathcal M_*(g_tm_{t-1}[F]).
\end{equation}
Provided $\varepsilon_j\le\varepsilon_\flat$ through the
levels used, the test partition and its gauge-invariant
extension identity give, for bounded gauge-invariant $F$,
\begin{equation}\label{IV-47-36}
 \langle F\rangle
  =Z^{-1}\int\Phi_T[F]\,dm_T
     +Z^{-1}\sum_{t=1}^T\int\Phi_{t-1}[F]g_t^{\rm w}\,dm_{t-1}.
\end{equation}
To obtain this identity, partition the histories according
to their first outcome different from the all-small-field
outcome.  Sum all later extensions of each such prefix;
the normalized test partition leaves precisely its prefix
mass.  The sole remaining prefix is $h_T$.
Thus $w_t\ge0$, all the measures $m_{j|T}$ have the same
mass $1-\sum_{t\le T}w_t$, and
\begin{equation}\label{eq:first-large-field-mass}
 \sum_{t\le T}w_t
     =1-\mathsf t(h_T)/Z
     =\widetilde\nu_K\{Z_T\ne\emptyset\}.
\end{equation}
The last equality uses consistency of the history measures;
the region $Z_T$ is empty exactly for prefixes equal to $h_T$.

Define a step quotient by
\[
 \mathsf{Qt}_j[F]=\frac{d\mathcal M_*(Fg_jm_{j-1})}{dm_j},
 \qquad
 \mathsf{Qt}_j[F;G]=\mathsf{Qt}_j[FG]
                              -\mathsf{Qt}_j[F]\mathsf{Qt}_j[G].
\]
Positivity gives $|\mathsf{Qt}_j[F]|\le\|F\|_\infty$
and $|\mathsf{Qt}_j[F;G]|\le2\|F\|_\infty\|G\|_\infty$.
The defining pushforward gives
\begin{equation}\label{IV-65-38}
 \int\mathsf{Qt}_j[F]\,dm_{j|T}=\int F\,dm_{j-1|T}.
\end{equation}
Put
$\Phi_{j,12}=\Phi_j[F_1F_2]-\Phi_j[F_1]\Phi_j[F_2]$.
The same pushforward gives the tower identities
\begin{equation}\label{IV-65-40}
 \Phi_j[F]=\mathsf{Qt}_j[\Phi_{j-1}[F]],\qquad
 \Phi_{j,12}=\mathsf{Qt}_j[\Phi_{j-1,12}]
              +\mathsf{Qt}_j[\Phi_{j-1}[F_1];\Phi_{j-1}[F_2]].
\end{equation}
These are identities of Radon--Nikodym derivatives, obtained
by testing against bounded functions on level $j$.
They require no choice of conditional probability kernels.

\begin{proposition}\label{IV-65-49}
Fix $2\le\mathsf k\le T$ in the admitted range.  Assume the
running and regularity hypotheses for
\eqref{eq:cov-transport-sizes} through $T$, and the mixed
derivative hypotheses of \cref{IV-64-43} through $\mathsf k$.
Then
\begin{equation}\label{IV-65-50}
 \langle\mathcal O(w_1);\mathcal O(w_2)\rangle
  =\sum_{\mathsf k<j\le T}
          \int\mathsf{Qt}_j[X_{j-1,1};X_{j-1,2}]\,dm_{j|T}
                                                +\mathrm{Rest}_T.
\end{equation}
Here the error satisfies
\begin{align}\label{IV-65-51}
 |\mathrm{Rest}_T|\le{}&M_{\mathsf k}
 +2\sum_{\mathsf k\le j<T}
       (X^*_{j,1}I^*_{j,2}+I^*_{j,1}X^*_{j,2}+I^*_{j,1}I^*_{j,2})
       +2Y_{T,1}Y_{T,2}\notag\\
 &+\rho_{12}+\rho_1Y_{T,2}+\rho_2Y_{T,1}+\rho_1\rho_2.
\end{align}
The quantities in this estimate are defined entirely by
bounds already obtained:
\begin{gather*}
 Y_{j,a}=X^*_{j,a}+I^*_{j,a},\qquad
 B^T_{j,a}=Y_{j,a}+\sum_{j<i\le T}\sum_{X,z}
                              |\dot E^{(i)}_a(X,\mathbf1,z)|,\\
 M'_j=M_j\ (j\le\mathsf k),\qquad
 M'_j=M_{\mathsf k}+2\sum_{\mathsf k\le i<j}Y_{i,1}Y_{i,2}
                                                  \ (j>\mathsf k),\\
 \rho_a=\sum_{t\le T}B^T_{t-1,a}w_t,\qquad
 \rho_{12}=\sum_{t\le T}(M'_{t-1}+B^T_{t-1,1}B^T_{t-1,2})w_t.
\end{gather*}
$M_j$ is the right side of \eqref{IV-64-45} at $\bar k=j$;
at level zero set $M_0=I^*_{0,a}=0$ and
$X^*_{0,a}=2a_K^{-1}\sum_p|w_a(p)|$.
\end{proposition}

\begin{proof}
Equation \eqref{IV-65-32} gives
$|\Phi_j[F_a]-c_{j,a}|\le Y_{j,a}$ and, through
$\mathsf k$, $|\Phi_{j,12}|\le M_j$.
The tower identity and the step covariance bound propagate
the latter to $|\Phi_{j,12}|\le M'_j$ at every later level.
Center the two observables by $c_{T,a}$ and apply
\eqref{IV-47-36} to each centered observable and to their
product.  The omitted first-large-field contributions to
the means are bounded by $\rho_a$, and the one to the
product by $\rho_{12}$.  At the top the remaining centered
means are bounded by $Y_{T,a}$.  Its measure has mass at
most one, so its product-minus-product-of-means term is at
most $2Y_{T,1}Y_{T,2}$.

Integrate the second tower identity against $m_{j|T}$ and
use \eqref{IV-65-38}.  Summation from $\mathsf k+1$ to $T$
leaves the initial mixed derivative bounded by $M_{\mathsf k}$.
In each later covariance substitute
$\Phi_{j-1}[F_a]=c_{j-1,a}+X_{j-1,a}+I_{j-1,a}$.
Constants cancel; the three terms involving an $I$ are
bounded by twice the corresponding products in
\eqref{IV-65-51}.  Their measures have mass at most one.
Collecting these estimates proves the formula and its bound.
\end{proof}

\subsection{Carrying a single-step covariance to a common level}

For a pointed family $f$ of level $i'\le i$, let
$f_i[f]$ be the bare array with that member and with all
other members and counterterms zero.  Then
\[
 N_i(f_i[f])=\omega_{i+1-i'}\|f\|_{i',\kappa},\qquad
 S[f_i[f]](U)=\mathsf{Sum}[f](U)-\mathsf{Sum}[f](\mathbf1),
\]
where $\mathsf{Sum}$ denotes the finite sum over domains and
base points.  Put
\begin{equation}\label{IV-67-12}
 \dot{\mathsf{St}}_i[f]=D\mathsf{St}_i(f_{0,i})f_i[f].
\end{equation}
This is linear, preserves vanishing on domains disjoint
from a fixed set of sites, and
\begin{equation}\label{IV-67-16}
 \|\dot{\mathsf{St}}_i[f]\|_{i+1,+}
       \le\theta_i\omega_{i+1-i'}\|f\|_{i',\kappa}.
\end{equation}
It follows from \eqref{IV-64-21} and locality; the insertion
parameter is restricted to a disc whose product with
$N_i(f_i[f])$ is smaller than the remaining array room.

Let $K^{(j)}$, $j=i+1\ge2$, be the second derivative of
the step with only the two plaquette weight members inserted,
leaving its input families equal to their ordinary values:
\begin{equation}\label{IV-65-42}
 K^{(i+1)}=D^2\mathsf{St}_i(f_{0,i})
       [f_i[\mathsf W_{L^{K-i}w_1}],f_i[\mathsf W_{L^{K-i}w_2}]].
\end{equation}
It vanishes unless its domain meets both supports, and
$\|K^{(j)}\|_{j,+}\le\nu_{j-1}/4$ when
$s_{j-1}\le s^\star_\omega/2$.
Differentiating the exact one-step identity gives
\begin{equation}\label{IV-65-45}
 \mathsf{Qt}_j[X_{j-1,1};X_{j-1,2}](W)
                             =\mathsf{Sum}[K^{(j)}](U_j(W)).
\end{equation}
Indeed, the sourced quotient is the exponential of
$\sum_a\lambda_aX_{j,a}$ plus the change in the sum of the
output family.  Its logarithm's mixed derivative is exactly
the displayed sum.

Starting with $K^{(j)\to j}=K^{(j)}$, define
\begin{equation}\label{IV-67-21}
 K^{(j)\to(i+1)}=\sum_{i'=j}^i
                 \dot{\mathsf{St}}_i[K^{(j)\to i'}],\qquad
 S_n^{(j)}=\sum_{i'=j}^n\mathsf{Sum}[K^{(j)\to i'}].
\end{equation}
Set $k_{i'}^{(j)}=\|K^{(j)\to i'}\|_{i',+}$.
Equation \eqref{IV-67-16} gives
\begin{equation}\label{IV-67-24}
 k_{i+1}^{(j)}\le\theta_i\sum_{i'=j}^i
                              \omega_{i+1-i'}k_{i'}^{(j)}.
\end{equation}
If $S_\omega\theta_i\le1$, where $S_\omega=\sum_n\omega_n$,
then $k_{i'}^{(j)}\le k_j^{(j)}$ and
$k_{i+1}^{(j)}\le S_\omega\theta_i k_j^{(j)}$.
Each family still vanishes on domains disjoint from either
set of source corners.

\begin{proposition}\label{IV-67-23}
For $j\le T$ in the admitted range,
\begin{equation}\label{IV-67-26}
 \int\mathsf{Qt}_j[X_{j-1,1};X_{j-1,2}]\,dm_{j|T}
                           =\int S_T^{(j)}(U_T(W))\,dm_{T|T}(W).
\end{equation}
If also $\theta_i\le\theta_{T-1}$ and
$S_\omega\theta_i\le1$ for $j\le i<T$, then
\begin{multline}\label{IV-67-31}
 |S_T^{(j)}(U_T)-S_T^{(j)}(\mathbf1)|
 \\
 \le C_{\rm ft}k_j^{(j)}\bigg\{N_j(w_1)
       \left[\frac{4B_3^2\varepsilon_T^2}{\alpha_0^2}L^{-4(T-j)}
                        +e^{-[d_*(j;T,2B_3\varepsilon_T)]_+}\right]
       \\
 +S_\omega\theta_{T-1}N_T(w_1)
   \sum_{i'=j+1}^TL^{3(T-i')}
       \left[\frac{4B_3^2\varepsilon_T^2}{\alpha_0^2}L^{-4(T-i')}
                       +e^{-[d_*(i';T,2B_3\varepsilon_T)]_+}\right]
                     \bigg\}.
\end{multline}
This holds at every passing top-level field.  The same
right side bounds the error in replacing the right side of
\eqref{IV-67-26} by
$S_T^{(j)}(\mathbf1)\,m_{T|T}(\mathcal V_T)$.
\end{proposition}

\begin{proof}
Insert $\mu f$ into the bare array at height $i$ and
differentiate its exact sourced step.  Its input action is
$\mu(\mathsf{Sum}[f](U)-\mathsf{Sum}[f](\mathbf1))$.
The derivative of the output logarithm is the same
subtracted action at $U_{i+1}$ plus
$\mathsf{Sum}[\dot{\mathsf{St}}_i[f]](U_{i+1})$.
Adding back the constant gives the step quotient identity
\[
 \mathsf{Qt}_{i+1}[\mathsf{Sum}[f](U_i)]
  =\mathsf{Sum}[f](U_{i+1})
                   +\mathsf{Sum}[\dot{\mathsf{St}}_i[f]](U_{i+1}).
\]
These identities have bounded Borel versions, by difference
quotients of the exact chart density as in \cref{IV-65-25}.
Summing them over $i'=j,\ldots,i$ in \eqref{IV-67-21}
shows that $S_{i+1}^{(j)}(U_{i+1})$ is the step quotient of
$S_i^{(j)}(U_i)$.  Equations \eqref{IV-65-38} and
\eqref{IV-65-45} telescope to \eqref{IV-67-26}.

Apply \eqref{IV-64-39} to each carried family at the regular
top background.  The anchored polymer sum at level $i'$
costs at most $8M^5\max(C_AC_{E,3},2C_{E,1})N_{i'}(w_1)$
times its norm and the two brackets displayed above.
Keep the bound $k_j^{(j)}$ for the first level.  At later
levels use $k_{i'}^{(j)}\le S_\omega\theta_{T-1}k_j^{(j)}$
and $N_{i'}(w_1)\le L^{3(T-i')}N_T(w_1)$.
This proves \eqref{IV-67-31} with
$C_{\rm ft}=8M^5\max(C_AC_{E,3},2C_{E,1})$.
The final assertion follows because $m_{T|T}$ has mass at
most one.
\end{proof}

\section{The fluctuation measure near the trivial field}
\label{sec:fluctuation-chain}

The evaluation at the trivial top still contains the nonlinear
background maps at all intermediate levels.  We follow these maps
downwards, keeping both the potential and its current in their
analyticity cores.  Their linearized curvatures give the covariance
slices of \cref{IV-44-21}.

\subsection{Critical curves and the actual shift}

We first record the finite-dimensional consequence of criticality
needed below.  All pairings and difference quotients in this paragraph
are in level-$n$ units, with fine mesh $\xi=\eta_n$.

\begin{lemma}\label{lem:critical-curve-curvature}
Suppose $U_t$ is a curve of configurations on a finite torus,
$U_0=\mathbf1$, differentiable at zero bond by bond, and that each
$U_t$, $t\ne0$, is critical on its fibre
$\mathcal M^nU=\mathcal M^nU_t$.  Write
$U_t=\mathbf1+it\xi\dot A+o(t)$ and
$\mathcal M^nU_t=\mathbf1+itB+o(t)$.  At an admitted side,
\[
 \dot A=H_NB-\nabla^\xi\theta,
 \qquad \xi^2D\dot A=H_n^{\rm cv}B
\]
for a real site function $\theta$.
\end{lemma}

\begin{proof}
Use logarithmic coordinates $U=e^{ia}$ near the identity, and write
$\Upsilon(a)=-i\log\mathcal M^n(e^{ia})$.  The derivative
$X=D\Upsilon(0)=Q^{\rm lin}_{0\to n}$ is onto.  Indeed, the
one-step right inverse sets a field to zero on bonds within each
block and to the prescribed coarse value on bonds crossing the
corresponding face.  Each averaging path crosses that face once,
and all its internal contour terms vanish.  Composing these lifts
gives a common real right inverse $E$ of the straight and linearized
averages.  Thus
$Z(a)=E[D\Upsilon(a)E]^{-1}$ is defined near zero.

For $v\in\ker X$, put
$v_t=v-Z(a_t)D\Upsilon(a_t)v$, where $a_t=-i\log U_t$.
Then $v_t$ is tangent to the fibre and $v_t\to v$.
If $S(a)$ is the action, criticality gives $DS(a_t)v_t=0$.
Since $DS(0)=0$ and $a_t/t\to\xi\dot A$, division by $t$
gives $D^2S(0)(\xi\dot A,v)=0$ for every $v\in\ker X$.
Factoring this linear functional through $X$ gives a unique
multiplier $\dot\omega$ such that
\[
 (D\dot A,DA')_\xi
   =(\dot\omega,Q^{\rm lin}_{0\to n}(\xi A'))_1.
\]
Testing gradients shows that $\dot\omega$ annihilates coarse
gradients.  The identity
$Q^{\rm lin}_{0\to n}(\xi A')=QA'-\partial\varkappa_n(\xi A')$
therefore yields $D^*D\dot A=Q^*\dot\omega$.

For completeness the gauge normalization can be imposed without
altering this equation.  Let $\theta_1$ be constant on each fine
block, with block mean $-\varkappa_n(\xi\dot A)$.  Write
$\Theta_0=\ker Q'$ and decompose
$f=d^*(\dot A+\nabla^\xi\theta_1)$ orthogonally as
$f=f'+f''$, with $f'\in\Delta\Theta_0$ and
$f''\perp\Delta\Theta_0$.  Choose $\vartheta\in\Theta_0$
with $\Delta\vartheta=-f'$ and put $\theta=\theta_1+\vartheta$.
Then $Q'\theta=-\varkappa_n(\xi\dot A)$ and
$d^*(\dot A+\nabla^\xi\theta)=f''$ annihilates every local
Landau test.  Consequently
$h=\dot A+\nabla^\xi\theta$ satisfies $Qh=B$ and the Landau
system.  Uniqueness in equations~\eqref{I-82-31}--\eqref{II-56-6}
gives $h=H_NB$.  Taking the curl and using
\eqref{eq:cv-landau-identification} proves the assertion.
\end{proof}

At height $i\ge1$, let $\mathfrak h$ be the canonical shift
system, with constants $(a_{\mathfrak h},\varepsilon_{\mathfrak h},
c_{\mathfrak h})$.  Its full-parameter value at a block-field
background is the actual lower background.  More precisely, for
$W$ in the class of the whole-torus sum rule, $U_+$ a representative
of its minimal orbit, and $B\in\mathsf{Rg}_{x_i}$, put
\[
 (\mathbf A,\mathbf K)
  =\mathfrak h_{T_\ell}(\mathbf1;P_W,g_iB).
\]
These fields are real and
\begin{equation}\label{eq:shift-background-identity}
 \begin{gathered}
 |\mathbf A|,\ |\nabla^{\eta_i}_{U_+}\mathbf A|,
 |\mathbf K|\le c_{\mathfrak h}\sup|g_iB|,\\
 e^{i\eta_i\mathbf A}U_+=U'^u,
 \qquad U'\in\mathcal C(\mathbf B_i,V_W(B)),\\
 J^{(i)}(e^{i\eta_i\mathbf A}U_+)
       =L^{-3}J^{(i+1)}(U_+)+\mathbf K.
 \end{gathered}
\end{equation}
Here the constants defining the class of $W$ are those of this
same shift system.  To check the statement, the class provides
the real background, invertible pivot matrices and the entire
fluctuation chart, including the existence of $U_i(V_W(B))$.
Its triple and $g_iB$ lie in the shift domain.  The parameter
$\mathbf1$ lies strictly inside the decoupling polydisc, since
$e^{\bar\kappa}>1$.  The reality and field bounds of the shift
therefore apply.  At this parameter its background identity
applies with precisely $V_W(B)$; its current identity is the
last equality in \eqref{eq:shift-background-identity}.
These are the global shift properties in the construction of
Propositions~\ref{II-40-99}--\ref{II-42-39}.

\subsection{Descent inside the analytic cores}

Put
\begin{equation}\label{IV-67-35}
 \begin{aligned}
 \rho_{\rm ch}&=\min\{a_{\mathfrak h}/3,r_\alpha,1/60\},
 &\rho_{\rm fl}&=\min\{\varepsilon_{\mathfrak h},
                                      (60c_{\mathfrak h})^{-1}\},\\
 M_{\rm pl}&=\tfrac94c_{\mathfrak h}\rho_{\rm fl}
                      +3L^{-2}\rho_{\rm ch},
 &M_{\rm gr}&=L^{-1}\rho_{\rm ch}
                      +\tfrac{11}{10}c_{\mathfrak h}\rho_{\rm fl}.
 \end{aligned}
\end{equation}
For a chart point $v=(A,\mathbf J)$ at level $i+1$ and an
unscaled fluctuation $\mathcal A$, with
$\|v\|<\rho_{\rm ch}$ and $\sup|\mathcal A|<\rho_{\rm fl}$,
let $(\mathbf A^{\mathfrak h},\mathbf K^{\mathfrak h})
=\mathfrak h_{T_\ell}(\mathbf1;P[v],\mathcal A)$ and define
\begin{equation}\label{IV-67-37}
 \begin{split}
 A^{\rm lw}&=(i\eta_i)^{-1}
       \log(e^{i\eta_i\mathbf A^{\mathfrak h}}
                         e^{i\eta_{i+1}A}),\\
 v^{\rm lw}&=(A^{\rm lw},L^{-3}\mathbf J+\mathbf K^{\mathfrak h}).
 \end{split}
\end{equation}
The logarithm is taken bond by bond on the identity branch.
The shift and substitution estimates give holomorphy, reality
at real arguments, and
\begin{equation}\label{IV-67-42}
 \begin{gathered}
 |A^{\rm lw}|,\ |\nabla^{\eta_i}A^{\rm lw}|
       \le L^{-1}\|v\|+\tfrac{11}{10}c_{\mathfrak h}\sup|\mathcal A|,\\
 |\partial(e^{i\eta_iA^{\rm lw}})-\mathbf1|
       \le\eta_i^2(3L^{-2}\|v\|
                       +\tfrac94c_{\mathfrak h}\sup|\mathcal A|).
 \end{gathered}
\end{equation}
Indeed $P[v]$ has field size at most $3\rho_{\rm ch}$ and
current and potential size at most $\rho_{\rm ch}$, so it lies
both in the shift domain and in the usable core.  In the
substitution estimate its three parameters are
$L^{-1}\|v\|$, $3L^{-2}\|v\|$ and
$c_{\mathfrak h}\sup|\mathcal A|$, each less than $1/60$.
This proves the domain assertion as well as the displayed bounds.

Let $\mathsf{Cc}_n(r)$ consist of real potentials $A$ such that
\[
 \|(A,J^{(n)}(e^{i\eta_nA}))\|\le r,
 \qquad e^{i\eta_nA}=U^u,\quad
 U\in\mathcal C(\mathbf B_n,W),\quad
 \varpi(W)<\varepsilon_{\mathcal W}\alpha_0.
\]
The second condition ensures that this small chart point is an
actual background.  Write
$v_A^{[n]}=(A,J^{(n)}(e^{i\eta_nA}))$.

\begin{proposition}\label{IV-67-46}
Suppose
\[
 r<\rho_{\rm ch},\qquad \delta_i<\rho_{\rm fl},\qquad
 6L^{-2}r+8C_{\rm db}\delta_i
                              <\varepsilon_{\mathcal W}\alpha_0.
\]
For $A\in\mathsf{Cc}_{i+1}(r)$ and $B\in\mathsf{Rg}_{x_i}$,
the lowering map at $(v_A^{[i+1]},g_iB)$ satisfies
\begin{equation}\label{IV-67-47}
 A^{\rm lw}\in\mathsf{Cc}_i
       (L^{-1}r+\tfrac{11}{10}c_{\mathfrak h}\delta_i),
 \qquad v^{\rm lw}=v_{A^{\rm lw}}^{[i]}.
\end{equation}
With $\widetilde U=e^{i\eta_{i+1}A}$ and
$\widetilde U^{\rm lw}=e^{i\eta_iA^{\rm lw}}$, every admitted
bare array $\mathfrak f$ has the exact identity
\begin{equation}\label{IV-67-48}
 (\mathbb V'[\mathfrak f])^{\rm tot}(P[v_A];B)
      =S[\mathfrak f](\widetilde U^{\rm lw})
                         -S[\mathfrak f](\widetilde U).
\end{equation}
In particular this is the difference of the two pointed totals
for an inserted family, and for a plaquette weight it is
\begin{equation}\label{IV-67-49}
 a_K^{-1}\sum_pw(p)
      [a_p(\widetilde U^{\rm lw})-a_p(\widetilde U)].
\end{equation}
Also $0\in\mathsf{Cc}_n(0)$.
\end{proposition}

\begin{proof}
Choose the axial representative $U_+$ of the background over
$W$, and write $\widetilde U=U_+^u$.  Gauge covariance carries
$B$ to $\check B=\operatorname{Ad}(u)^{-1}B$, preserving its
allowed radius.  The pair of $P[v_A]$ is the pair of $P_W^u$;
the shift depends only on that pair.  Thus
\eqref{eq:shift-background-identity} at $(P_W,g_i\check B)$,
conjugated by $u$, gives both the lower minimal orbit and the
current identity in \eqref{IV-67-47}.  Equation
\eqref{IV-67-42} bounds the potential and its derivative, and
the current is bounded by $L^{-3}r+c_{\mathfrak h}\delta_i$.

The upper background has regularity number less than $3r'$
for any $r'>r$ sufficiently close to $r$.  The block-background
estimate bounds its averaged plaquette size by $6L^{-2}r'$.
The fluctuation chart has the form
$V_W(\check B)=e^{iB'}V_i^{\rm bf}(W)$ with
$\sup|B'|<2C_{\rm db}\delta_i$.  Telescoping the four unitary
plaquette factors gives
$\varpi(V_W(\check B))<6L^{-2}r'+8C_{\rm db}\delta_i$,
which is less than $\varepsilon_{\mathcal W}\alpha_0$.
This verifies the lower class condition, not only its norm bound.

The whole-torus sum rule for the fourth potential at $P_W$
is the difference of the array actions at the actual two
backgrounds.  Gauge invariance of the actions and covariance
and pair dependence of the local potentials carry it to $P[v_A]$.
This proves \eqref{IV-67-48}; the definition of the plaquette
array gives \eqref{IV-67-49}.  The identity field is flat,
critical and has zero current, proving the last assertion.
\end{proof}

If $r_T=0$ and
$r_i=L^{-1}r_{i+1}+\frac{11}{10}c_{\mathfrak h}\delta_i$,
then under the running-coupling bounds of
\eqref{IV-64-14}, extended through $T$,
\begin{equation}\label{IV-67-51}
 r_i\le\tfrac85c_{\mathfrak h}\delta_i.
\end{equation}
Indeed $\delta_{i+n}\le(5L/4)^{n/2}\delta_i$, and unwinding
the recursion gives the geometric series
\[
 \frac{11}{10}\sum_{n\ge0}(5/(4L))^{n/2}<8/5.
\]
Thus the preceding proposition can be iterated whenever its
three inequalities hold with $r=r_{i+1}$.

\subsection{The quadratic curvature and its covariance}

Define
\[
 \operatorname{ang}^{[n]}_p(A)=\eta_n(dA)(p),\qquad
 T_i^{\mathfrak h}\mathcal A
    =\eta_i d\bigl(D_{\mathcal A}\mathbf A^{\mathfrak h}(0,0)
                                                     \mathcal A\bigr),
 \qquad T_i^{\rm cv}=H_i^{\rm cv}C_{\mathbf1,\mathbb R}.
\]

\begin{lemma}\label{IV-67-45}
Suppose $\ell\ge N_{\mathsf H}$. There is a sign $\sigma_i\in\{1,-1\}$ such that
$T_i^{\mathfrak h}=\sigma_i(T_i^{\rm cv}\otimes I)$.
\end{lemma}

\begin{proof}
Constant gauge covariance and reality imply that each matrix
entry of the linearized shift is a real scalar on the colour
space.  In fact it commutes with the adjoint action of $SU(2)$:
the three half-turns about the Pauli axes make it diagonal,
and an adjoint rotation cyclically permuting those axes makes
the three diagonal entries equal.

For a real fluctuation $\mathcal A$, apply
\eqref{eq:shift-background-identity} over $W=\mathbf1$ to
$t\mathcal A$.  It gives a differentiable curve $U_t$ through
$\mathbf1$ of critical configurations, with coarse field
$V_t^{w_t}$, where
$V_t=\exp(itC_{\mathbf1}\mathcal A+o(t))$.
Absorb the part of the gauge transformation equal to the identity
on coarse sites into the minimal orbit.  The remaining gauge
is its extension from $w_t$ on the coarse sites.
By \cref{lem:critical-curve-curvature},
$T_i^{\mathfrak h}\mathcal A=(H_i^{\rm cv}\otimes I)B$,
where $V_t^{w_t}=\mathbf1+itB+o(t)$.

Fix a coarse site $y_0$, put $g_t=w_t(y_0)$ and
$\check w_t=g_t^{-1}w_t$.  The bond identity for $V_t^{w_t}$
implies $\check w_t(c_-)-\check w_t(c_+)=O(t)$.
The finite coarse graph is connected, so
$\check w_t=\mathbf1+O(t)$ everywhere.  The real site functions
$\lambda_t=(it)^{-1}\log\check w_t$ are bounded.  Along a
subsequence $g_t\to g_*$ and $\lambda_t\to\lambda_*$, and
expanding the bond identity gives
\[
 B=\operatorname{Ad}(g_*)
           (C_{\mathbf1}\mathcal A-d\lambda_*).
\]
The curvature response kills gradients, whence
$T_i^{\mathfrak h}\mathcal A
=\operatorname{Ad}(g_*)(T_i^{\rm cv}\otimes I)\mathcal A$.

Apply this to $\mathcal A=t_1a$ for a real scalar field $a$.
The two scalar response matrices, say $S,T$, satisfy
$Sa=\sigma(a)Ta$ with $\sigma(a)\in\{1,-1\}$ when $Ta\ne0$,
and $Sa=0$ otherwise.  If $Ta,Ta'$ are independent, linearity
applied to $a+a'$ forces their signs to agree.  If they are
dependent, subtract a multiple and use $\ker T\subset\ker S$.
Thus one sign works for every $a$ and, by complex linearity,
for every colour-valued field.  Its value is immaterial below.
\end{proof}

\begin{lemma}\label{IV-67-43}
Let $v,\mathcal A$ be real, $\|v\|\le r$, $\sup|\mathcal A|\le a$,
and $R=\min(\rho_{\rm ch}/r,\rho_{\rm fl}/a)\ge2$.
A quotient with zero denominator is interpreted as $\infty$,
and $\infty^{-n}=0$ for $n>0$.
With
$\Phi_p=\operatorname{ang}^{[i+1]}_p(A)+T_i^{\mathfrak h}\mathcal A(p)$,
\begin{equation}\label{IV-67-44}
 \begin{gathered}
 |a_p(e^{i\eta_iA^{\rm lw}})-\tfrac12|\Phi_p|^2|
                       \le M_{\rm pl}^2\eta_i^4R^{-3},\\
 |a_p(e^{i\eta_{i+1}A})-\tfrac12|
                 \operatorname{ang}^{[i+1]}_p(A)|^2|
                       \le M_{\rm pl}^2\eta_i^4R^{-3},\\
 |\operatorname{ang}^{[i]}_p(A^{\rm lw})-\Phi_p|
                       \le4M_{\rm gr}\eta_i^2R^{-2}.
 \end{gathered}
\end{equation}
Also $|\operatorname{ang}^{[i+1]}_p(A)|\le2L^{-2}\eta_i^2r$
and $|T_i^{\mathfrak h}\mathcal A(p)|\le2c_{\mathfrak h}\eta_i^2a$.
\end{lemma}

\begin{proof}
Scale both arguments by $\zeta$, $|\zeta|<R$.  The lowered
potential is holomorphic and has first derivative
$L^{-1}A+D_{\mathcal A}\mathbf A^{\mathfrak h}(0,0)\mathcal A$;
the derivative in $v$ vanishes since the shift is zero when
$\mathcal A=0$.  If $E(\zeta)$ is its plaquette variable minus
$I$, \eqref{IV-67-42} gives
$|E(\zeta)|\le\eta_i^2M_{\rm pl}$.
For $2\times2$ unimodular matrices,
$1-\frac12\operatorname{tr}(I+E)=\frac12\det E$.
Thus this holomorphic scalar function has modulus at most
$\frac12\eta_i^4M_{\rm pl}^2$ and its quadratic coefficient is
$\frac12\det(i\Phi_p)=\frac12|\Phi_p|^2$.
Cauchy's estimate and $R\ge2$ bound the terms of order at least
three by $M_{\rm pl}^2\eta_i^4R^{-3}$.
Set $\mathcal A=0$ for the second assertion.  Apply the same
argument to the linear curl of the lowered potential, whose
modulus is at most $2\eta_i^2M_{\rm gr}$ and first coefficient
is $\Phi_p$, to get the quadratic remainder.
The last two estimates follow by differentiating the shift's
gradient estimate at zero and by taking the curl of $A$.
\end{proof}

\subsection{Four fluctuation cumulants}

At $P[v]$ let $C_v$ be the whole covariance of the canonical
Gaussian structure.  Write $V_i^{\rm int}=V_i^J+V_i^G+V_i^A$
and $V_i^0=V'_i[f_{0,i}]$.  The probability measure used in
one real step is
\begin{equation}\label{IV-67-14}
 d\nu_{i,v}(B)=\Xi_i(v)^{-1}
 \mathbf1_{\{|B(b)|<\mathfrak b(x_i)\ \forall b\}}
 \exp\{(V_i^{\rm int}+V_i^0)^{\rm tot}(P[v];B)\}\,d\mu_{C_v}(B).
\end{equation}
The cluster identity \eqref{II-19-39} of \cref{II-19-55}
makes $\Xi_i(v)$ nonzero throughout
the complex core.  The intrinsic and ordinary-array potentials
are both of the bounded-potential sort and satisfy
\[
 \|V_i^{\rm int}\|_{\star,\lambda_V}\le3K_P\delta_i,
 \qquad \|V_i^0\|_{\star,\lambda_V}\le K_\Lambda s_iN_\star.
\]
Their sum is bounded by
\[
 n_i^V=\left(\tfrac72K_P+
       K_\Lambda(1+\beta_s)N_\star/\alpha_0\right)\delta_i.
\]
The canonical place conditions ensure that these two norms
are at most $w_\star/2$ and $w_\star/4$, respectively.
Define
\begin{equation}\label{IV-67-53}
 \begin{gathered}
 n_f=\tfrac12M^{3/2},\qquad c_C=2/\gamma_r,\\
 p_{\star,i}=C_{\rm ex}((L+2)M)^3e^{4\lambda_1}
                         e^{-\gamma_2\mathfrak b(x_i)^2/2},\\
 \epsilon_i^{\rm cu}
   =C_{E,0}(K_\Pi n_i^V+K_P^\Pi p_{\star,i})
                                  (4n_f/w_\star)^4.
 \end{gathered}
\end{equation}

\begin{proposition}\label{IV-67-56}
For real $v$ with $\|v\|<r_\alpha$, $\nu_{i,v}$ is a
probability measure with positive density on the fluctuation
box.  For components $B_a=B^{\alpha_a}(b_a)$, repetitions
allowed, the cumulants of orders at most four obey
\begin{equation}\label{IV-67-57}
 \begin{split}
 |\langle B_1;\cdots;B_n\rangle_{\nu_{i,v}}|
       &\le\epsilon_i^{\rm cu} \quad(n=1,3,4),\\
 |\langle B_1;B_2\rangle_{\nu_{i,v}}-C_v(1,2)|
       &\le\epsilon_i^{\rm cu}.
 \end{split}
\end{equation}
The covariance is holomorphic in the complex chart ball and
\begin{equation}\label{IV-67-59}
 |C_v(a,b)|\le c_C,\qquad
 |C_v(a,b)-C_0(a,b)|\le2c_C\|v\|/r_\alpha.
\end{equation}
Its value at zero satisfies
\begin{equation}\label{IV-67-60}
 C_0=(A_i^{\rm cv})^{-1}\otimes I,
 \qquad T_i^{\rm cv}(A_i^{\rm cv})^{-1}(T_i^{\rm cv})^*
                         =\eta_i\mathcal C_{(i+1)}.
\end{equation}
These component estimates do not claim spatial decay.
\end{proposition}

\begin{proof}
The real reference inverse in the Gaussian structure has gap
$\gamma_r$, and its perturbation has norm at most
$\widehat\epsilon\le\gamma_r/2$.  Hence
\[
 \operatorname{Re}(B^*,C_v^{-1}B)\ge\tfrac12\gamma_r\|B\|^2,
 \qquad \|C_v\|\le2/\gamma_r.
\]
Conjugation covariance makes $C_v$ and its inverse real at
real $v$; they are symmetric, so the displayed inequality
makes them positive definite.  The potentials are real there.
Thus the unnormalized density in \eqref{IV-67-14} is strictly
positive on an open box containing zero.  Its integral is
positive, proving the probability assertion.

Place the insertion $B^\alpha(b)$ in the elementary domain
$\square(b)$ and set all its other members to zero.  Its norm
is at most $n_f$, since
$t/(1+M^{-3}t^2)\le M^{3/2}/2$.
Add $\sum_{a=1}^nz_aB_a$ as a single quadratic-sort origin,
on the polydisc $|z_a|<w_\star/(n n_f)$.  The resulting
system has three origins and is admitted by the Gaussian
cluster theorem.  Its connected terms and the whole integral
are holomorphic in $z$, and the nonzero-integral identity
identifies their sum with the logarithm of that integral.
After division by its value at zero, the mixed derivative
is exactly the joint cumulant.

Remove the cutoff and the two bounded-potential origins,
retaining the insertion origin.  The Gaussian-part comparison
\eqref{eq:cluster-gaussian-part} bounds the difference of connected
terms on $X$ by
$(K_\Pi n_i^V+K_P^\Pi p_{\star,i})e^{-\lambda_{\rm out}d_{i+1}(X)}$
throughout the insertion polydisc.  Cauchy's formula adds the
factor $(n n_f/w_\star)^n$.  The mixed derivative vanishes
unless $X$ contains the entry block of every insertion;
anchoring at one such block bounds the sum over $X$ by
$C_{E,0}$.  For $1\le n\le4$ this proves
\eqref{IV-67-57}, because the uncut generating function is
$\exp(\frac12\sum_{a,b}z_az_bC_v(a,b))$.

The whole covariance is the finite matrix
$C+ C\Gamma\Gamma^*C$ from the conditioning identity;
its entries are holomorphic on the chart ball.  For $v\ne0$
apply Schwarz's lemma to
$C_{\zeta v}(a,b)-C_0(a,b)$ on
$|\zeta|<r_\alpha/\|v\|$, where its modulus is at most
$2c_C$.  This gives \eqref{IV-67-59}.
Finally, the canonical provision at the trivial pair has
the actual constrained Hessian used in \cref{IV-44-21}.
That proposition and its curvature covariance identity give
\eqref{IV-67-60} with the counting pairings used here.
\end{proof}

Differentiating the normalized one-step integral also supplies
the insertions used in \cref{sec:covariance-transport}:
\begin{equation}\label{eq:step-insertion-expectation}
 \begin{split}
 \mathsf{Sum}[\dot{\mathsf{St}}_i[f]](P[v])
      &=\nu_{i,v}((V'_i[f])^{\rm tot}),\\
 \mathsf{Sum}[K^{(i+1)}](P[v])
      &=\langle (V'_i[w_1])^{\rm tot};
                     (V'_i[w_2])^{\rm tot}\rangle_{\nu_{i,v}}.
 \end{split}
\end{equation}
The normalization term is independent of the inserted array,
so its derivative vanishes.  The integral is nonzero on a
common parameter neighbourhood and the fluctuation box is
bounded; thus differentiation of its logarithm is legitimate.
Equations \eqref{IV-67-48}--\eqref{IV-67-49} identify these
insertions with the actual lower-background differences.

\section{Gaussian evaluation of the transported covariance}
\label{sec:gaussian-covariance-transport}

Take $2\le\mathsf k<T\le K$, with the preceding canonical
data and exact covariance reduction through $T$.  Suppose
$w_a(p)=f_a(x(p))$, where each $f_a$ is real Lipschitz and
supported in a ball of radius $d\ge a_K$, and the two nonempty
plaquette supports have corner separation $D_{\mathsf k}\ge3M$
in level-$\mathsf k$ units.  Put
\begin{equation}\label{IV-68-3}
 \begin{gathered}
 \widehat d_i=d/a_{K-i},\quad P_i=81\widehat d_i^3,
 \quad \upsilon_i=(a_{K-i+1}x_{i-1})^{-1},\\
 u_a^{[i]}=\upsilon_i\|w_a\|_\infty,
 \quad \widehat m_a=u_a^{[\mathsf k+1]},
 \quad Y_i=1+A_1y_i^{\mathfrak p_0}.
 \end{gathered}
\end{equation}
Assume the running bounds through $T$:
\begin{equation}\label{IV-68-4}
 |x_j-L^{-(j-i)}x_i|\le\bar B\quad(i\le j\le T),\qquad
 x_T\ge4\bar B,\quad y_T\ge\max(2\mathfrak p_0,1+\log L).
\end{equation}
Put $n_a=\#\operatorname{supp}w_a$.
Each support contains at most $81(d/a_K)^3$ plaquettes:
for each of three directions its barycentres lie in a shifted
grid with at most $(2d/a_K+1)^3\le27(d/a_K)^3$ points in
the enclosing cube.  Consequently $n_a\eta_i^3\le P_i$.
The same coordinate count for the source corners gives
\[
 N_i(w_a)\le(2\widehat d_i/M+3)^3,
 \qquad N_i(w_a)\le L^{3(i'-i)}N_{i'}(w_a)\quad(i\le i').
\]
The coupling bounds give, for $n=j-i$,
\begin{equation}\label{eq:short-coupling-comparison}
 \begin{gathered}
 \tfrac34L^nx_j\le x_i\le\tfrac54L^nx_j,
 \quad y_j\le y_i\le(n+1)y_j,\quad \delta_i\le\delta_j,
 \quad\varepsilon_i\le\varepsilon_j,\\
 \upsilon_j\le\tfrac54\upsilon_{\mathsf k+1},\qquad
 |\upsilon_j-\upsilon_{\mathsf k+1}|
     \le\upsilon_{\mathsf k+1}\bar B/x_{j-1}
       \quad(\mathsf k<j\le T).
 \end{gathered}
\end{equation}
The first line follows directly by comparison with $x_j$;
monotonicity of $x^{-1/2}(\log x)^{\mathfrak p_0}$ holds
for $\log x\ge2\mathfrak p_0$.  For the last line divide
$|x_{j-1}-L^{-(j-1-\mathsf k)}x_{\mathsf k}|\le\bar B$
by the corresponding physical inverse coupling.

Let $r_i$ be the descending radii in \eqref{IV-67-51}, and set
\[
 c_R=\max\left\{\tfrac85\sqrt{5L/4}\,
                    c_{\mathfrak h}/\rho_{\rm ch},\rho_{\rm fl}^{-1}\right\},
 \qquad e_i^{\rm co}=\epsilon_i^{\rm cu}
       +\tfrac{16}5\sqrt{5L/4}\,
                         c_Cc_{\mathfrak h}\delta_i/r_\alpha.
\]
In addition to the place conditions assume
\begin{equation}\label{IV-68-10}
 \begin{gathered}
 \ell\ge N_{\mathsf H},\qquad
 c_R\delta_i\le\tfrac12,\qquad e_i^{\rm co}\le1,\\
 [\tfrac{48}5\sqrt{5/4}L^{-3/2}c_{\mathfrak h}
                   +8C_{\rm db}]\delta_i
                     <\varepsilon_{\mathcal W}\alpha_0
                    \quad(\mathsf k\le i<T).
 \end{gathered}
\end{equation}
These imply all three hypotheses of the descent proposition
with $r=r_{i+1}$, and the Taylor radius there has inverse
at most $c_R\delta_i$.

For a colour-valued plaquette field $F$, define
\begin{equation}\label{IV-68-13}
 Q^{(j)}(F)=\upsilon_ja_K^{-1}
     \sum_{p,q}w_1(p)w_2(q)\mathcal C_{(j)}(p,q)
                                                   \langle F_p,F_q\rangle.
\end{equation}
The Gaussian contribution after transport from $j$ through $i$ is
\begin{equation}\label{IV-68-14}
 \begin{split}
 n^G_{j,i}={}&\tfrac{N^2-1}{2}\upsilon_j^2
            \sum_{p,q}w_1(p)w_2(q)\mathcal C_{(j)}(p,q)^2\\
 &+(N^2-1)\sum_{j<l\le i}\upsilon_j\upsilon_l
            \sum_{p,q}w_1(p)w_2(q)
                       \mathcal C_{(j)}(p,q)\mathcal C_{(l)}(p,q).
 \end{split}
\end{equation}
Here $N=2$; retaining $N^2-1$ displays the colour multiplicity.
There are fixed finite constants $C_F,C'_F$ for which the errors
are defined by
\begin{equation}\label{IV-68-16}
 \begin{split}
 e_{j,j}&=C_FY_{j-1}^4P_{j-1}^2
           (\delta_{j-1}+e_{j-1}^{\rm co})u_1^{[j]}u_2^{[j]},\\
 e_{j,i+1}&=e_{j,i}+C'_FY_i^2P_{j-1}P_i
           (\delta_i+e_i^{\rm co})u_1^{[j]}u_2^{[i+1]}.
 \end{split}
\end{equation}
Their values are fixed by the estimates in the following proof,
independently of $i,j,K,\ell$ and the sources.

\begin{proposition}\label{IV-68-17}
Under these hypotheses, for $\mathsf k<j\le i\le T$ and
$A\in\mathsf{Cc}_i(r_i)$,
\begin{equation}\label{IV-68-18}
 \left|S_i^{(j)}(e^{i\eta_iA})-n^G_{j,i}
                 -Q^{(j)}(\operatorname{ang}^{[i]}A)\right|
                                                   \le e_{j,i}.
\end{equation}
In particular $|S_T^{(j)}(\mathbf1)-n^G_{j,T}|\le e_{j,T}$.
\end{proposition}

\begin{proof}
At height $i$, the lowered potential lies in
$\mathsf{Cc}_i(r_i)$ for every fluctuation in the support of
$\nu=\nu_{i,v_A}$.  This is a probability measure.  Write
$F=\operatorname{ang}^{[i+1]}A$ and $T=T_i^{\rm cv}$.
The linearized lower angle is $F+\sigma_ig_iTB$.
The bounds already proved give
\[
 \sum_b|T(p,b)|\le C^{\rm row}_{\rm tp}\eta_i^2,
 \quad |F_p|\le2L^{-2}\eta_i^2r_{i+1},
 \quad g_i|TB(p)|\le2c_{\mathfrak h}\eta_i^2\delta_i.
\]
The means and third and fourth cumulants of the components of
$B$ are at most $\epsilon_i^{\rm cu}$; their second cumulants
differ from $C_0$ by at most $e=e_i^{\rm co}\le1$.
Furthermore $|C_0(a,b)|\le c_C$.

For the initial step take $i=j-1$.  The insertion identities give
$S_j^{(j)}(\widetilde U)=\langle\Sigma_1;\Sigma_2\rangle_\nu$,
where $\Sigma_a$ is \eqref{IV-67-49}.  Split it as
$\Sigma_a=L_a+Q_a+R_a$, with
\[
 L_a=\sigma_ig_ia_K^{-1}\sum_pw_a(p)\langle F_p,TB(p)\rangle,
 \qquad Q_a=\tfrac12g_i^2a_K^{-1}\sum_pw_a(p)|TB(p)|^2.
\]
If $L_a=\sum c_a(b,\alpha)B^\alpha(b)$ and
$Q_a=\frac12\sum K_a(b,b')\sum_\alpha B^\alpha(b)B^\alpha(b')$,
the row bound, $n_a\eta_i^3\le P_i$ and
$a_K^{-1}\eta_i=x_i\upsilon_{i+1}$ give
\begin{align*}
 \|R_a\|_\infty&\le2M_{\rm pl}^2c_R^3\delta_i^3x_iP_i u_a^{[j]},\\
 \|L_a\|_\infty&\le c_{\mathfrak h}^2\delta_i^2x_iP_i u_a^{[j]},
 &\|Q_a\|_\infty&\le2c_{\mathfrak h}^2\delta_i^2x_iP_i u_a^{[j]},\\
 \|c_a\|_1&\le C^{\rm row}_{\rm tp}c_{\mathfrak h}
                     A_1y_i^{\mathfrak p_0}P_i u_a^{[j]},
 &\|K_a\|_1&\le(C^{\rm row}_{\rm tp})^2P_i u_a^{[j]}.
\end{align*}
The covariance bound $2\|G\|_\infty\|H\|_\infty$ handles
every term containing $R_a$.  Using $c_R\delta_i\le1/2$,
their sum is at most
\[
 [24c_{\mathfrak h}^2M_{\rm pl}^2c_R^3
                       +4M_{\rm pl}^4c_R^5]
       \delta_i(\delta_i^2x_i)^2P_i^2u_1^{[j]}u_2^{[j]}.
\]

The remaining calculation uses only four cumulants.  Differentiating
$\exp(\log\nu(e^{\sum z_aB_a}))$ expresses each moment as the
sum over partitions of the component indices.  For a covariance
of two products, exactly the partitions with a block crossing
the two products remain.  Thus
\[
 |\langle B_1;B_2B_3\rangle_\nu|
          \le(3+2c_C)\epsilon_i^{\rm cu},
\]
and the covariance of $B_1B_2$ with $B_3B_4$ differs from
$C_0(1,3)C_0(2,4)+C_0(1,4)C_0(2,3)$ by at most
$(11+8c_C)e$.  Indeed the eleven crossing partitions give
one fourth cumulant, four third-cumulant times mean terms,
two covariance products and four covariance times two means;
the error is bounded by
$\epsilon+4\epsilon^2+4(c_C+1)\epsilon^2+2e(2c_C+e)$,
with $\epsilon\le e\le1$.

Replacing the covariance in $\langle L_1;L_2\rangle$ by $C_0$
gives $Q^{(j)}(F)$ with error at most
$e\|c_1\|_1\|c_2\|_1$.  The two mixed linear-quadratic
covariances cost at most
$\frac32(3+2c_C)\epsilon(\|c_1\|_1\|K_2\|_1+
\|c_2\|_1\|K_1\|_1)$.  The quadratic-quadratic covariance
has error at most
$\frac94(11+8c_C)e\|K_1\|_1\|K_2\|_1$.
Its two Gaussian contractions each use
$TC_0T^*=\eta_i\mathcal C_{(j)}\otimes I$;
the colour sum is $N^2-1$, and their total is $n^G_{j,j}$.
Here $g_i^2\eta_i/a_K=\upsilon_j$ and $\sigma_i^2=1$.
Since $x_i\delta_i^2\le Y_i^2$, these estimates prove the
initial case with
\begin{align*}
 C_F={}&24c_{\mathfrak h}^2M_{\rm pl}^2c_R^3
       +4M_{\rm pl}^4c_R^5
       +(C^{\rm row}_{\rm tp})^2c_{\mathfrak h}^2\\
 &+3(3+2c_C)(C^{\rm row}_{\rm tp})^3c_{\mathfrak h}
       +\tfrac94(11+8c_C)(C^{\rm row}_{\rm tp})^4.
\end{align*}

For the induction step, the first insertion identity and the
exact action difference imply
\[
 S_{i+1}^{(j)}(\widetilde U)
                =\int S_i^{(j)}(\widetilde U^{\rm lw})\,d\nu.
\]
The constant old total cancels the subtracted action in the
new carried family.  Integrating the induction hypothesis
therefore costs only $e_{j,i}$.  Let $\overline Q^{(j)}$ be
the absolute coefficient sum in \eqref{IV-68-13}.
For any two fields,
\[
 |Q^{(j)}(F')-Q^{(j)}(F'')|
 \le\overline Q^{(j)}(\|F'\|_\infty+\|F''\|_\infty)
                                      \|F'-F''\|_\infty.
\]
The angle remainder in \eqref{IV-67-44} then bounds replacement
of the true lower angle by $F+\sigma_ig_iTB$ by
$22M_{\rm gr}c_{\mathfrak h}c_R^2\delta_i^3
\overline Q^{(j)}\eta_i^4$.
Its linear expectation costs at most
$C^{\rm row}_{\rm tp}c_{\mathfrak h}g_i\delta_i e
\overline Q^{(j)}\eta_i^4$.
The second moment of two field components differs from $C_0$
by at most $e+\epsilon^2\le2e$; hence its quadratic expectation
equals $n^G_{j,i+1}-n^G_{j,i}$ with error at most
$6(C^{\rm row}_{\rm tp})^2g_i^2e\overline Q^{(j)}\eta_i^4$.
The slice bound gives
\[
 \overline Q^{(j)}\eta_i^4
       \le C_{\rm tp}P_{j-1}P_i x_i u_1^{[j]}u_2^{[i+1]}.
\]
Since $x_i\delta_i^3\le Y_i^2\delta_i$,
$x_ig_i\delta_i\le Y_i$ and $x_ig_i^2=1$, the induction
follows with
\[
 C'_F=C_{\rm tp}[22M_{\rm gr}c_{\mathfrak h}c_R^2
              +C^{\rm row}_{\rm tp}c_{\mathfrak h}
              +6(C^{\rm row}_{\rm tp})^2].
\]
At $T$, the identity potential belongs to $\mathsf{Cc}_T(0)$
and its angle vanishes, proving the final assertion.
\end{proof}

\subsection{The seven errors in the covariance}

Assume in addition $S_\omega\theta_i\le1$,
$s_i\le s^\star_\omega/2$ and $\theta_i\le\theta_{T-1}$
for $\mathsf k\le i<T$, and choose $\mathsf k<j_1\le T$.
Let $\mathrm{Ft}_T^{(j)}$ be the right side of
\eqref{IV-67-31}, $W_T=1-m_T(\mathcal V_T)/Z$,
$C'_{\rm tp}=C_{\rm tp}/(1-L^{-3})$, and
\[
 \overline n_j^G=\tfrac{25}{16}(N^2-1)
       C_{\rm tp}C'_{\rm tp}P_{j-1}^2\widehat m_1\widehat m_2.
\]
Define the nonnegative errors
\begin{equation}\label{IV-68-23}
 \begin{split}
 E_2&=2\sum_{j_1<j\le T}X^*_{j-1,1}X^*_{j-1,2},\qquad
 E_3=\sum_{\mathsf k<j\le j_1}\mathrm{Ft}_T^{(j)},\\
 E_4&=W_T\sum_{\mathsf k<j\le j_1}(\overline n_j^G+e_{j,T}),
 \qquad E_5=\sum_{\mathsf k<j\le j_1}e_{j,T},\qquad
 E_6=\sum_{j_1<j\le T}\overline n_j^G,\\
 E_7&=\tfrac{N^2-1}{2}\tfrac94(C'_{\rm tp})^2P_{\mathsf k}^2
              \frac{\bar B}{x_T}\widehat m_1\widehat m_2,\\
 E_8&=\tfrac{N^2-1}{2}C^{\rm bg}_{\rm tp}P_T
       (2C'_{\rm tp}P_{\mathsf k}+C^{\rm bg}_{\rm tp}P_T)
                                         \widehat m_1\widehat m_2.
 \end{split}
\end{equation}

\begin{theorem}\label{IV-68-24}
Under these hypotheses,
\begin{equation}\label{IV-68-25}
 \left|\langle\mathcal O(f_1);\mathcal O(f_2)\rangle
    -\tfrac{N^2-1}{2}\upsilon_{\mathsf k+1}^2
                        \mathcal K_{K,K-\mathsf k}(f_1,f_2)\right|
       \le|\mathrm{Rest}_T|+\sum_{r=2}^8E_r.
\end{equation}
\end{theorem}

\begin{proof}
The slice estimate, $u_a^{[j]}\le5\widehat m_a/4$ and
$P_{l-1}=L^{-3(l-j)}P_{j-1}$ give
$|n^G_{j,T}|\le\overline n_j^G$ by a geometric sum.
In the exact covariance reduction \eqref{IV-65-50}, the
steps $j>j_1$ cost $E_2$ by positivity and the covariance
supremum bound.  For $j\le j_1$, transport to level $T$,
replace its background by the identity using
\eqref{IV-67-31}, and apply \cref{IV-68-17}.
Since the top measure has mass $1-W_T$, these three changes
cost $E_3,E_4,E_5$.  Restoring the Gaussian terms for
$j>j_1$ costs $E_6$.

The complete Gaussian sum is exactly
\[
 \sum_{\mathsf k<j\le T}n^G_{j,T}
  =\tfrac{N^2-1}{2}\sum_{p,q}w_1(p)w_2(q)
          \left(\sum_{\mathsf k<j\le T}
                         \upsilon_j\mathcal C_{(j)}(p,q)\right)^2.
\]
Put $D=\Pi_{\mathsf k}-\Pi_T$ and
$\upsilon=\upsilon_{\mathsf k+1}$.  By
\eqref{eq:short-coupling-comparison} the weighted slice sum
differs from $\upsilon D$ by at most
$\upsilon(\bar B/x_T)C'_{\rm tp}\eta_{\mathsf k}^3$;
the sum of their absolute values is at most
$\frac94\upsilon C'_{\rm tp}\eta_{\mathsf k}^3$.
The identity $a^2-b^2=(a-b)(a+b)$ and the plaquette counts
give $E_7$.  Finally
$C^{\rm bg}_{K,K-\mathsf k}=D+C^{\rm bg}_{K,K-T}$.
The bound $|C^{\rm bg}_{K,K-T}|\le C^{\rm bg}_{\rm tp}\eta_T^3$
and the same square identity give $E_8$.
The resulting square is precisely the kernel sum in
\eqref{IV-68-25}.
\end{proof}


\subsection{Summing the errors}

Write $p=\mathfrak p_0$, $y=y_{\mathsf k}$,
$q=\widehat d_{\mathsf k}$, $D=D_{\mathsf k}$ and $n_T=T-\mathsf k$
in this subsection.  Constants denoted by $c$ or $C$ depend only on
the fixed construction constants.  They do not depend on
the sources, the levels, the coupling sequence or the torus.
We continue to assume $q\ge M$, $D\ge3M$ and all the
hypotheses of \cref{IV-68-24}.
Here $d_*(t)=d_*(t;t,2B_3\varepsilon_t)$, and the symbols
$\rho_a^T,\rho_{12}^T$ below denote the $\rho_a,\rho_{12}$
of \eqref{IV-65-51} at the final level $T$.
The regularity assumptions include
$\varepsilon_t\le c_A\alpha_0/(8MB_3)$ and
$L^{m+K-t}\ge8M+4$ for $1\le t\le T$.
They make both entries defining $d_*(t)+1$ at least $4$,
so $d_*(t)\ge3$.

\begin{lemma}\label{IV-68-28}
There are constants $C_{\rm sz},C_{\mathsf M},c_{\rm up},c_I$
with the following properties.  Put
$r_{i,a}=\sum_{X,z}|\dot E^{(i)}_a(X,\mathbf1,z)|$.
For $j\le\mathsf k$, $n=\mathsf k-j$,
\begin{equation}\label{IV-68-29}
\begin{aligned}
 X^*_{j,a}+I^*_{j,a}+\sum_{j<i\le\mathsf k}r_{i,a}
 &\le C_{\rm sz}q^3 L^{4n}(n+1)^{2p}x_{\mathsf k}y^{2p}
                         \widehat m_a,\\
 M_j&\le C_{\mathsf M}q^3L^{4n}(n+1)^{2p}x_{\mathsf k}y^{2p}
               e^{-(D/M-2)}\widehat m_1\widehat m_2.
\end{aligned}
\end{equation}
For $\mathsf k\le j\le T$,
\begin{equation}\label{IV-68-30}
\begin{gathered}
 d_*(j)\ge d_*(T)\ge3,\qquad
 X^*_{j,a}\le c_Xy^{2p}\widehat d_j^3\widehat m_a,
 \qquad c_X=\tfrac{405}{2}B_3^2A_0^2,\\
 I^*_{j,a}\le i_T\widehat m_a,\qquad
 i_T=c_Iq^3y^{3p}
       \max\{x_T^{-1/2},\sqrt{x_{\mathsf k}}e^{-d_*(T)}\}.
\end{gathered}
\end{equation}
Define
\[
 b_n=(C_{\rm sz}+c_{\rm up}n_T)q^3L^{4n}(n+1)^{2p}
                         x_{\mathsf k}y^{2p},\qquad
 m_*=C_{\mathsf M}q^3x_{\mathsf k}y^{2p}.
\]
Then, for $0\le j\le T$ and $n=[\mathsf k-j]_+$,
\[
 Y_{j,a}\le B^T_{j,a}\le b_n\widehat m_a,\qquad
 M'_j\le[m_*L^{4n}(n+1)^{2p}+2n_Tb_0^2]
                                      \widehat m_1\widehat m_2.
\]
\end{lemma}

\begin{proof}
Locality of $\dot E_a^{(i)}$ requires its polymer to meet a
plaquette in the source support.  Anchor it in a cell containing
a corner of that plaquette.  There are $N_i(w_a)$ possible cells,
and at most $8M^3(d_i(X)+1)$ choices of its marked site.
The entropy sum at the output decay rate therefore gives
\[
 r_{i,a}\le8M^3C_{{\rm E},1}N_i(w_a)d_{i,a}.
\]
The first-derivative recurrence bounds $d_{i,a}$ by a fixed
multiple of
$\delta_{i-1}|w_a|_\infty/a_{K-i+1}$.
For $i=\mathsf k-n'$ this is at most
\[
 C L^{(n'+1)/2}(n'+2)^p\sqrt{x_{\mathsf k}}y^p\widehat m_a;
\]
the number of anchors is at most $125M^{-3}q^3L^{3n'}$.
For $i>\mathsf k$, the corresponding product is at most
$c_{\rm up}q^3\sqrt{x_{\mathsf k}}y^p\widehat m_a$.
The sum over $0\le n'<n$ is bounded by
$Cq^3L^{7n/2}(n+1)^p\sqrt{x_{\mathsf k}}y^p\widehat m_a$.

For $j=\mathsf k-n\ge1$, the plaquette count and the
coupling comparison give
$X^*_{j,a}\le216B_3^2A_0^2q^3L^{3n}(n+1)^{2p}
y^{2p}\widehat m_a$.  At $j=0$ the bound is
$162q^3L^{4n}x_{\mathsf k}\widehat m_a$.
In the first-derivative remainder, its bracket is at most
$2$ because $d_*(j)\ge3$ and $2B_3\varepsilon_j\le\alpha_0$.
It is consequently bounded by
$Cq^3L^{7n/2}(n+1)^p\sqrt{x_{\mathsf k}}y^p\widehat m_a$.
These estimates prove the first line of \eqref{IV-68-29}.

For the second line, put $u=D_j/M=L^nD/M\ge3$ in the
mixed-derivative sum.  Its term at distance $h$ below $j$
has $\nu_{j-h-1}/\nu_j\le\frac43L^{h+1}(h+2)^{2p}$.
Since
$L^hu-2\ge(u-2)+3(L^h-1)$, its sum is at most
\[
 \tfrac43\nu_je^{-(u-2)}
       \sum_{h\ge0}L^{4h+1}(h+2)^{2p}e^{-3(L^h-1)}.
\]
The series converges.  Moreover
$\nu_j\le C L^n(n+1)^{2p}x_{\mathsf k}y^{2p}
\widehat m_1\widehat m_2$ and $N_j(w_1)\le125M^{-3}q^3L^{3n}$.
This proves the required estimate, including $M_0=0$.

Above $\mathsf k$, both entries defining $d_*(j)$ decrease
as $j$ increases.  The direct bound for $X^*$ gives the
first estimate in \eqref{IV-68-30}.  Substitution into the
remainder estimate gives
\[
 I^*_{j,a}\le C(2\widehat d_j/M+3)^3
 [x_j^{-1/2}y_j^{3p}+e^{-d_*(j)}\sqrt{x_j}y_j^p]
                                                    \widehat m_a.
\]
Use $(2\widehat d_j/M+3)^3\le125M^{-3}q^3$ and
$x_T\le x_j\le x_{\mathsf k}$.  This proves the displayed
bound for $i_T$.  The same formula with its original bracket
bounded by $2$ gives
$I^*_{j,a}\le Cq^3\sqrt{x_{\mathsf k}}y^p\widehat m_a$.
Finally $B^T_{j,a}=X^*_{j,a}+I^*_{j,a}+\sum_{j<i\le T}r_{i,a}$,
and the definition of $M'_j$ adds at most $2n_Tb_0^2$ above
$\mathsf k$.  Increasing the fixed constants proves the last
two assertions.
\end{proof}

Fix the parameter $0<\eth<1$ used in
\cref{lem:iv-early-history-tail}, and put
\begin{equation}\label{IV-68-32}
 c'=(1-\eth)c_\star,\qquad
 B(x)=C_{\rm maj}^{\Sigma}+2+c_{\rm low}+\tfrac92\log x.
\end{equation}
We call $\ell=2L^m$ small over the running sequence if
\begin{equation}\label{IV-68-33}
 MR(x_K)\le L^m,\qquad LM\le L^m,\qquad
 \ell\ge N_{\mathsf H},\qquad
 B(x_K)\ell^3\le\tfrac12c'A_0^2y_K^{2p}.
\end{equation}

\begin{lemma}\label{IV-68-34}
Assume the actual history hypotheses of
\cref{lem:iv-early-history-tail} and a small side.  Then
\begin{equation}\label{IV-68-35}
 w_t\le W_t=\sum_{i=1}^tw_i
 \le2M^{-3}\ell^3L^{3(K-t)}e^{-c'A_0^2y_t^{2p}/2}.
\end{equation}
\end{lemma}

\begin{proof}
The normalized history measure gives
$W_t=\widetilde\nu_K\{Z_t\ne\varnothing\}$ by
\eqref{eq:first-large-field-mass}.  Divide
\eqref{eq:iv-early-history-tail} by the partition-function
lower bound \eqref{eq:actual-partition-bounds}, and use the product bound
\eqref{II-203-25}.  This gives
\[
 W_t\le2M^{-3}\ell^3L^{3(K-t)}
       \exp[-c'A_0^2y_t^{2p}+B(x_K)\ell^3].
\]
Now $y_t\ge y_K$ and the last condition of
\eqref{IV-68-33} absorb half of the negative exponent.
\end{proof}

\begin{lemma}\label{IV-68-36}
Assume also
$y_K\ge2\log(1+\bar B)$ and $c'pA_0^2y_K^{2p-1}\ge13$.
Set
\[
 c_p=\sup_{a\ge0}(a+2)^{4p}L^{-a},\qquad
 t_T=2M^{-3}\ell^3L^{3(K-\mathsf k)}e^{-c'A_0^2y_T^{2p}/2},
 \qquad z_T=c_Xy^{2p}\widehat d_T^3+i_T.
\]
Then
\begin{equation}\label{IV-68-37}
\begin{gathered}
 W_T\le t_T,\qquad
 \rho_a^T\le t_T(2L^4c_p+n_T)b_0\widehat m_a,\\
 \rho_{12}^T\le t_T(2L^8c_p+n_T)
              [m_*+(2n_T+1)b_0^2]\widehat m_1\widehat m_2.
\end{gathered}
\end{equation}
In particular $|\mathrm{Rest}_T|\le R_T\widehat m_1\widehat m_2$,
where
\begin{equation}\label{IV-68-38}
\begin{aligned}
R_T={}&m_*e^{-(D/M-2)}
 +2n_T(2c_Xy^{2p}q^3i_T+i_T^2)+2z_T^2\\
 &+t_T(2L^8c_p+n_T)
 [m_*+(2n_T+1)b_0^2+2b_0z_T
                    +t_T(2L^8c_p+n_T)b_0^2].
\end{aligned}
\end{equation}
\end{lemma}

\begin{proof}
The running-coupling rate estimate, with exponent $10$, is
\[
 L^{3a}e^{-c'A_0^2y_{\mathsf k-a}^{2p}/2}
 \le L^{-10a}e^{-c'A_0^2y_{\mathsf k}^{2p}/2}.
\]
Here the loss in the logarithmic coupling increment must be
controlled before using the power $2p$. For $a\ge1$, put
$\Delta=y_{\mathsf k-a}-y_{\mathsf k}$ and
$v=\bar B/x_{\mathsf k}\le1/4$. The running-coupling bound gives
\[
 \Delta\ge a\log L+\log(1-v)>2a,\qquad
 a\log L-\Delta\le\tfrac43\bar B e^{-y_{\mathsf k}}.
\]
Since $y_K\ge2\log(1+\bar B)$ and
$z e^{-z/2}\le2/e$ for $z>0$, it follows that
\[
 y_{\mathsf k}(a\log L-\Delta)
       \le\frac8{3e}\frac{\bar B}{1+\bar B}<1.
\]
Set $N=2p\ge2$. The linear and quadratic terms of the
power expansion give
\begin{align*}
 y_{\mathsf k-a}^{N}-y_{\mathsf k}^{N}
 &\ge N y_{\mathsf k}^{N-1}a\log L\notag\\
 &\quad+N y_{\mathsf k}^{N-2}
     \left[\frac{N-1}{2}\Delta^2
               -y_{\mathsf k}(a\log L-\Delta)\right]
 \ge N y_{\mathsf k}^{N-1}a\log L.
\end{align*}
The bracket is positive because $\Delta>2a$ and $a\ge1$.
The case $a=0$ is immediate. Multiplication by $c'A_0^2/2$
gives a rate at least $c'pA_0^2y_K^{2p-1}\ge13$;
the volume factor $L^{3a}$ leaves the stated rate $10$.
For $t\le\mathsf k$, put $a=\mathsf k-t$.
The preceding lemma gives $w_t\le t_TL^{-10a}$, while
\cref{IV-68-28} gives
\[
 B^T_{t-1,b}\le b_0L^4c_pL^{5a}\widehat m_b,
\quad
 M'_{t-1}+B^T_{t-1,1}B^T_{t-1,2}
 \le[m_*+(2n_T+1)b_0^2]L^8c_pL^{9a}
                              \widehat m_1\widehat m_2.
\]
Sum these products using $\sum_{a\ge0}L^{-a}\le2$.
For the remaining $n_T$ levels use $w_t\le t_T$ and $n=0$
in \cref{IV-68-28}.  This proves \eqref{IV-68-37}.
Insert these estimates, \eqref{IV-68-29} at $j=\mathsf k$,
and \eqref{IV-68-30} into the exact remainder
\eqref{IV-65-51}.  The $n_T$ cross terms give the second
term of \eqref{IV-68-38}; the top covariance gives $2z_T^2$;
the linear and quadratic tail terms give its last line.
\end{proof}

We now choose the two upper levels:
\begin{equation}\label{IV-68-40}
 j_1=\mathsf k+\lceil\log_L(qy^p)\rceil,\qquad
 T=\mathsf k+\lceil\log_L(qy^{3p})\rceil,\qquad n_T=T-\mathsf k.
\end{equation}
They are used only when $T\le K$.
They satisfy
$\widehat d_{j_1}\le y^{-p}$,
$\widehat d_T\le y^{-3p}$,
$L^{-(T-j_1)}<Ly^{-2p}$ and $qy^{3p}\le L^{n_T}<Lqy^{3p}$.

\begin{lemma}\label{IV-68-42}
There are fixed constants $C_2,C_3,C'_3,C''_3,C_5,C_6,C_7,C_8$
such that, after division by $\widehat m_1\widehat m_2$,
the seven errors in \eqref{IV-68-23} satisfy
\begin{equation}\label{IV-68-43}
\begin{aligned}
 E_2&\le C_2y^{-2p},& E_6&\le C_6y^{-6p},\\
 E_5&\le C_5n_T^2(1+A_1y^p)^4q^6\bar e_T,
 &E_7&\le C_7q^6/x_T,\\
 E_8&\le C_8(q^3+1)y^{-9p},\\
 E_4&\le W_T[C_6n_Tq^6+C_5n_T^2(1+A_1y^p)^4q^6\bar e_T],\\
 E_3&\le C_3(n_Ty^{-5p}+y^{-2p})
       +C'_3n_Tqy^{8p}x_T^{-1/2}
       +C''_3n_Tq^3y^{11p}x_{\mathsf k}e^{-d_*(T)},
\end{aligned}
\end{equation}
where $\bar e_T=\max_{\mathsf k\le i<T}(\delta_i+e_i^{\rm co})$.
\end{lemma}

\begin{proof}
For $E_2$, sum $\widehat d_j^6$ over $j\ge j_1$ in
\eqref{IV-68-30}; its first term is at most $y^{-6p}$,
leaving $y^{4p}y^{-6p}$.  The identical geometric sum in
the Gaussian bound gives $E_6$.  The recurrence for $e_{j,i}$
has at most $n_T^2$ terms, each bounded by
$C(1+A_1y^p)^4q^6\bar e_T$; this proves $E_5$ and, together
with the Gaussian size, $E_4$.  Substituting
$P_{\mathsf k}=81q^3$, $P_T\le81y^{-9p}$ into
$E_7,E_8$ gives their displayed bounds.

For $E_3$, use $k_j\le\nu_{j-1}/4
\le Cx_{j-1}y^{2p}\widehat m_1\widehat m_2$ in
\eqref{IV-67-31}.  Its initial quadratic term contains
$x_{j-1}\varepsilon_T^2L^{-4(T-j)}N_j(w_1)$.
Use $x_{j-1}\le\frac54L^{T-j+1}x_T$,
$\varepsilon_T^2x_T\le A_0^2y^{2p}$,
$N_j(w_1)\le(2\widehat d_j/M+3)^3$ and
$L^{-(T-j)}\widehat d_j=\widehat d_T$.
The cell-dependent term is bounded, after summing over
$j\le j_1$, by $Cn_Ty^{-5p}$.
The constant term uses
\[
 \sum_{j\le j_1}L^{-3(T-j)}
 \le(1-L^{-3})^{-1}L^3y^{-6p}
\]
and is bounded by $Cy^{-2p}$.
The initial exceptional-polymer term is at most
$Cq^3x_{\mathsf k}y^{2p}e^{-d_*(T)}$ per level.
For the carried quadratic term use $N_T(w_1)\le125$,
$\sum_{h\ge0}L^{3h}L^{-4h}\le(1-L^{-1})^{-1}$,
$\theta_{T-1}\le Cx_T^{-1/2}y^p$ and
$x_{j-1}\varepsilon_T^2\le Cqy^{5p}$.
It is at most $Cqy^{8p}x_T^{-1/2}$ per level.
For the carried exceptional term use $S_\omega\theta_{T-1}\le1$
and $\sum_{0\le h<n_T-1}L^{3h}<q^3y^{9p}$;
this sum is empty when $n_T=1$.
It is bounded by
$Cq^3x_{\mathsf k}y^{11p}e^{-d_*(T)}$ per level.
These four estimates prove the stated bound for $E_3$.
\end{proof}

\begin{corollary}\label{IV-68-44}
Set
\begin{equation}\label{IV-68-45}
 s_T=\max\{x_T^{-1/2},x_{\mathsf k}e^{-d_*(T)},
 x_{\mathsf k}e^{-(D/M-2)},\bar e_T,x_{\mathsf k}^2t_T\}.
\end{equation}
If $s_T\le1$, there is a fixed $C_B$ such that
\begin{multline}\label{IV-68-46}
\left|\langle\mathcal O(f_1);\mathcal O(f_2)\rangle
 -\tfrac{N^2-1}{2}\upsilon_{\mathsf k+1}^2
                         \mathcal K_{K,K-\mathsf k}(f_1,f_2)\right|
\\
\le\Big[(C_2+C_3)y^{-2p}+C_3n_Ty^{-5p}+C_6y^{-6p}\\
 +C_8(q^3+1)y^{-9p}+4c_X^2y^{-14p}
             +C_Bn_T^4q^6y^{15p}s_T\Big]
                         \widehat m_1\widehat m_2.
\end{multline}
\end{corollary}

\begin{proof}
Insert \eqref{IV-68-43} and \eqref{IV-68-38} into
\eqref{IV-68-25}.  All terms in the seven-error bound
other than the first four displayed terms of
\eqref{IV-68-46} are bounded by a fixed multiple of
$n_T^4q^6y^{15p}s_T$.  Here $x_T^{-1}\le s_T^2\le s_T$,
$W_T\le t_T\le s_T$ and $1+A_1y^p\le(1+A_1)y^p$.
For the remainder use
$i_T\le c_Iq^3y^{3p}s_T$,
$z_T\le c_Xy^{-7p}+i_T$ and
$b_0\le Cn_Tq^3x_{\mathsf k}y^{2p}$.
The top term satisfies
$2z_T^2\le4c_X^2y^{-14p}+4i_T^2$.
In the last line of \eqref{IV-68-38}, its square bracket
is at most $Cn_T^3q^6y^{5p}x_{\mathsf k}^2$;
the remaining factor is at most $Cn_Tt_T$.
Thus $x_{\mathsf k}^2t_T\le s_T$ controls precisely the
volume-dependent tail.  The other terms are smaller
monomials with one of the five factors in \eqref{IV-68-45}.
Increasing $C_B$ proves the assertion.
\end{proof}


\section{The short-distance law on a small torus}
\label{sec:small-torus-law}

We now verify the preceding hypotheses for the actual
first-passage construction.  Write $p=\mathfrak p_0$ throughout
this section.  All structural constants and
canonical maps are fixed as in \cref{III-5-11} and
Theorem~\ref{II-203-9}.  In the weighted iteration we use
the admissible-radii row and the constants fixed in
\cref{sec:weighted-small-field}.  In particular $K_\Lambda$
and $K_{\mathbf M}$ have their weighted values.  Increasing
the lower bound on the inverse reference coupling does not
change these maps or their counterterms.

\subsection{The witness depth and the working levels}

Let $0<g\le\gamma$.  The first-passage endpoint is
$X_g=L(g^{-2}+B_\sharp)$, and put
\begin{equation}\label{IV-68-48}
 C_{W,0}=\max\{1,4L^2M/A_0^2\},\qquad c_W=C_{W,0}A_0^2.
\end{equation}
For integers $s\ge0$ define
\begin{equation}\label{IV-49-6}
\begin{gathered}
 \bar x_s=L^{s+1}(g^{-2}+\bar B),\qquad
 Y_s=\log\bar x_s,\qquad C_W(s)=c_WY_s^{2p},\\
 s(d)=\min\{s\in\mathbb Z_{\ge0}:L^sd\ge C_W(s)\},
 \qquad \check d=L^{s(d)}d.
\end{gathered}
\end{equation}
The minimum exists because an exponential eventually exceeds
a polynomial.  For $0<d<1$, its value is positive and
\[
 C_W(s(d))\le\check d<LC_W(s(d)).
\]
Indeed minimality gives $L^{s-1}d<C_W(s-1)\le C_W(s)$.
The witness depth is non-increasing in $d$ and tends to
infinity as $d\downarrow0$.

For a first-passage datum $(K,\beta)$ with physical label
$\mathsf b=a_K\beta$, assume $s=s(d)\le K$ and set
\begin{equation}\label{IV-68-50}
 \mathsf k=\max\{K-s,2\},\qquad
 s_{\mathsf k}=K-\mathsf k=\min\{s,K-2\},\qquad
 \mathsf L_s=L^2C_W(s)Y_s^{3p}.
\end{equation}
Thus the construction always uses a height at least two,
including when $K=s$ or $K=s+1$.

\begin{lemma}\label{IV-68-51}
For $0<d<1$ and $s(d)\le K$, one has $s\ge3$,
$2\le\mathsf k\le K-1$ and $s-2\le s_{\mathsf k}\le s$.
With $q=d/a_{s_{\mathsf k}}$,
\begin{equation}\label{eq:working-room}
 4MY_s^{2p}\le L^{-2}C_W(s)\le q<LC_W(s),\qquad
 a_K/d\le C_W(s)^{-1},\qquad d<LC_W(s)L^{-s}.
\end{equation}
On a side with $LM\le L^m$ one has $20d\le\ell$.
For a test pair in $\mathfrak T_{T_\ell}(d;\vartheta)$
with nonzero sampled weights,
\[
 D_{\mathsf k}\ge q/\sqrt3-1\ge3M,\qquad
 D_{\mathsf k}/M-2\ge\tfrac4{\sqrt3}Y_s^{2p}-3.
\]
Furthermore
\begin{equation}\label{eq:working-coupling}
 L^{s_{\mathsf k}}g^{-2}<x_{\mathsf k}\le\bar x_s,
 \quad Y_s-\log(2L^3)<y_{\mathsf k}\le Y_s,
 \quad |(\mathsf b\upsilon_{\mathsf k+1})^2-1|
                      \le3\bar Bg^2L^{2-s},
 \quad\upsilon_{\mathsf k+1}^2\le\tfrac98\mathsf b^{-2}.
\end{equation}
If $j_1,T$ are chosen by \eqref{IV-68-40}, then
$L^{n_T}<\mathsf L_s$.  If $\mathsf L_s\le L^{s-3}$ and
$g^{-2}\ge4\bar B$, then
\begin{equation}\label{eq:working-top}
 \mathsf k<j_1\le T\le K-2,\qquad n_T\le s-4,\qquad
 L^{K-T}>L^{s-2}/\mathsf L_s,\qquad
 x_T>\frac{4L^{s-2}g^{-2}}{5\mathsf L_s}.
\end{equation}
\end{lemma}

\begin{proof}
The normalization conditions give $g^{-2}\ge4$,
$\bar Bg^2\le13/24$, $L\ge13$ and $Y_0>2$.
Since $L^{-2}C_W(0)>1>d$, no index $s\le2$ can be a
witness.  The formula for $\mathsf k$ gives
$s-2\le s_{\mathsf k}\le s$; multiplying the bounds on
$\check d$ by $L^{s_{\mathsf k}-s}$ proves
\eqref{eq:working-room}.  Also $\ell\ge2LM>20d$.
The two supports have distance at least $d$.
Their plaquette corners have sup distance at least
$d/\sqrt3-a_K$.  Divide by $a_{s_{\mathsf k}}\ge a_K$
and use $q\ge4MY_s^{2p}$ to obtain the two separation bounds.

The first-passage condition gives $x_{K-1}>X_g$.
The backwards running estimate gives
\[
 x_{\mathsf k}-\bar B
 \ge L^{s_{\mathsf k}-1}(x_{K-1}-\bar B)
 >L^{s_{\mathsf k}}g^{-2},
\]
because $LB_\sharp-\bar B=(L-2)\bar B\ge0$.
The upper running estimate gives
$x_{\mathsf k}\le\bar x_{s_{\mathsf k}}\le\bar x_s$.
Since $\bar B<g^{-2}$ and $s_{\mathsf k}\ge s-2$,
$\bar x_s<2L^3x_{\mathsf k}$, proving the logarithmic bound.
Finally
\[
 \upsilon_{\mathsf k+1}^{-1}=a_{s_{\mathsf k}}x_{\mathsf k},
 \quad |\upsilon_{\mathsf k+1}^{-1}-\mathsf b|
                    \le\bar B a_{s_{\mathsf k}},
 \quad\upsilon_{\mathsf k+1}<g^2.
\]
Thus $|\mathsf b\upsilon_{\mathsf k+1}-1|
\le\bar Bg^2L^{2-s}\le1/24$; squaring gives
\eqref{eq:working-coupling}.
For the last assertion,
$L^{n_T}<Lqy_{\mathsf k}^{3p}<\mathsf L_s$.
Under its stated upper bound $n_T<s-3$, whence $T\le K-2$
and the bound on $L^{K-T}$ follows.
The running estimate
$x_T\ge\frac45L^{-n_T}x_{\mathsf k}$ proves the last inequality.
\end{proof}

\subsection{A single choice of the small-coupling threshold}

Let $\varepsilon(t)=A_0t^{-1/2}(\log t)^p$,
$\delta(t)=A_1t^{-1/2}(\log t)^p$ and
$\theta(t)=K_{\mathbf M}(1+\beta_s)\delta(t)/\alpha_0$.
Define $e^{\rm co}(t)$ by the formulas for
$e_i^{\rm co}$ and \eqref{IV-67-53}, replacing $x_i$ by $t$.
In particular its exponential term is a fixed multiple of
$\exp[-(LM)^{-3}A_1^2(\log t)^{2p}]$.
These functions are continuous, decrease on
$[e^{2p},\infty)$ and tend to zero.

Choose $x_W\ge\max\{\gamma^{-2},4\bar B,e^{2p},eL\}$
so that the following hold for every $t\ge x_W$:
\begin{equation}\label{IV-68-54}
\begin{gathered}
 2\varepsilon(t)\le\varepsilon_{\mathcal W}\alpha_0,
 \quad C_0^\flat S_A\theta(t)\le1,
 \quad\tfrac43C_1^\flat L T(3)\theta(t)\le1,\\
 (1+\beta_s)\delta(t)/\alpha_0\le\tfrac12s^\star_\omega,
 \quad\varepsilon(t)\le c_A\alpha_0/(8MB_3),
 \quad3K_P\delta(t)\le3w_\star/8,\\
 S_\omega\theta(t)\le1,
 \quad c_R\delta(t)\le\tfrac12,
 \quad e^{\rm co}(t)\le1,\\
 2[\tfrac{48}{5}\sqrt{5/4}\,L^{-3/2}c_{\mathfrak h}
                 +8C_{\rm db}]\delta(t)
                      \le\varepsilon_{\mathcal W}\alpha_0,\\
 c'pA_0^2(\log t)^{2p-1}\ge13,
 \quad2M(\log t)^r\ge N_{\mathsf H},\\
 4\delta_{\rm tp}L(L-1)M(\log t)^{2p}
                         \ge\sqrt3(3\log L+\log2).
\end{gathered}
\end{equation}
Require also $\log t\ge\max\{1+\log L,2\log(1+\bar B)\}$
and $(\log t)^r\ge C_{\rm s}$,
and the history-summation thresholds preceding
\cref{lem:iv-early-history-tail} at the fixed parameter $\eth$.
All these requirements can be imposed simultaneously:
the quantities to be made small tend to zero, the positive
powers of $\log t$ tend to infinity, and
$S_A,S_\omega,T(3)$ are finite.
Set $\gamma_W^{\rm ref}=x_W^{-1/2}$.

\begin{lemma}\label{IV-68-56}
Let $g\le\gamma_W^{\rm ref}$, let $(K,\beta)$ be a
first-passage datum with $K\ge3$, and suppose
$MR(x_K)\le L^m$.  The canonical construction satisfies
all the place, progression, density and regularity
hypotheses used above through level $K$.
Consequently, for $2\le\mathsf k<T\le K$, real Lipschitz
weights supported in balls of radius $d\ge a_K$,
$\widehat d_{\mathsf k}\ge M$, $D_{\mathsf k}\ge3M$ and
$\mathsf k<j_1\le T$, the conclusions of
\cref{IV-68-24,IV-68-28} apply.
The conclusion of \cref{IV-68-42} applies when $j_1,T$
are chosen by \eqref{IV-68-40}.
If the side is small, \cref{IV-68-36} also applies,
and \cref{IV-68-44} applies at these attached levels
when its remaining hypothesis $s_T\le1$ holds.
\end{lemma}

\begin{proof}
By \cref{III-5-11},
$x_0>\cdots>x_K>g^{-2}\ge x_W$;
the sequence and every initial segment are admissible.
The torus condition $MR(x_K)\le L^m$ makes the same
finite lattice admissible over every segment.  The canonical
one-lattice maps, radii and response therefore give the
actual places of Theorems~\ref{II-53-39} and~\ref{II-127-29} at all these
heights.  Their auxiliary datum is common to the heights,
and the octet property holds on every required smaller
lattice, as required by the canonical-place hypotheses
of those theorems.

The ordinary families satisfy $\|E^{(j)}\|\le E_0$ and
$|e_j|\le B_\sharp$.  Their weighted bare arrays consequently
have norm at most
\[
 (E_0+\tfrac32\alpha_0^2B_\sharp)
           \sum_{n\ge1}(n+1)L^{-(n-1)}<N_\star.
\]
The weighted radii and ceiling in
\cref{sec:weighted-small-field}, together with
\eqref{IV-68-54}, put the fourth potential below $w_\star/4$
and leave the required perturbation room.  The intrinsic
potential is at most $3w_\star/8$.  Thus the Lipschitz
and Cauchy derivative bounds proved there apply to these
same actual maps; no change of counterterms is involved.
Consistency of the ordinary recursion identifies
$x_{j+1}=L^{-1}x_j-e_{j+1}$ with the counterterm used in
the weighted step.  The Gaussian family is the same
localized family in the recursion and in the exact density
identity, so their flat terms cancel as in the proof of
Proposition~\ref{IV-65-25}.

The first two lines of \eqref{IV-68-54} give the regular
background and first-derivative hypotheses.  Since
$D_j\ge D_{\mathsf k}\ge3M$ for $j\le\mathsf k$,
the decreasing function $T(\cdot)$ satisfies
$T(D_j/M)\le T(3)$; hence the mixed
derivative hypotheses hold.  Also $d_*(t)\ge3$:
its analytic-radius entry is at least $4$, and
$L^{m+K-t}\ge LM\ge8M+4$ controls its geometric entry.
The next two lines give the chain-radius, covariance and
strict lowering-domain conditions in \eqref{IV-68-10};
the factor $2$ makes that domain inequality strict.
The weighted transport hypothesis $S_\omega\theta_i\le1$
is explicit.  The running bounds follow from
\cref{III-5-11}.  Finally $x_K>e$ implies $R(x_K)\ge L$,
so $L^m\ge LM$, and
$\ell\ge2MR(x_K)\ge2My_K^r\ge N_{\mathsf H}$.

The actual history construction has the same running
sequence and the finite-family tests required by
\cref{lem:iv-early-history-tail}.  Its small-field cutoffs
obey $\varepsilon_j\le\varepsilon_\flat$ by the fixed
construction calibration.  The additionally imposed
history thresholds and the last two rate conditions verify
\cref{IV-68-34,IV-68-36}.  This proves all the assertions.
\end{proof}

Put
$C_{\rm sd}=8L^6(C_{\rm maj}^{\Sigma}+2+c_{\rm low}+9/2)$.
Choose $Y_{\rm sd}$ so that for all $y\ge Y_{\rm sd}$,
\[
 C_{\rm sd}y\,[MR(e^y)]^3
                          \le\tfrac12c'A_0^2y^{2p}.
\]
Such a choice exists since $R(e^y)<Ly^r$ and
$2p>3r+2$.  Define
\begin{equation}\label{IV-68-55}
 \gamma_W^{\rm sd}=\min\{\gamma_W^{\rm ref},e^{-Y_{\rm sd}/2}\}.
\end{equation}

\begin{lemma}\label{IV-68-57}
If $g\le\gamma_W^{\rm sd}$, each side satisfying
$MR(x_K)\le L^m\le L^2MR(x_K)$ is small over the
first-passage sequence.  This includes both the least
admissible side $2L^{m_{\mathfrak L}(g)}$ and
the fixed reference side $2MR(X_g)$.
\end{lemma}

\begin{proof}
The first three conditions in \eqref{IV-68-33} follow from
\cref{IV-68-56}.  The upper side bound gives
$B(x_K)\ell^3\le C_{\rm sd}y_K[MR(x_K)]^3$,
and $y_K>\log g^{-2}\ge Y_{\rm sd}$ gives the fourth.

For the two specific sides recall
$g^{-2}<x_K\le X_g$ and the normalization inequalities
$X_g+\bar B<2Lx_K$ and
$\log(2L)\le(L^{1/r}-1)\log g^{-2}$.
Thus
\[
 R(x_K)\le R(X_g)\le R(X_g+\bar B)\le LR(x_K).
\]
The least side is determined by
\[
 L^{m_{\mathfrak L}(g)}
 =LM\max\{R(X_g+\bar B),L^{\lceil\log_LC_{\rm s}\rceil}\}.
\]
The imposed condition
$C_{\rm s}\le(\log x_W)^r$ gives
$L^{\lceil\log_LC_{\rm s}\rceil}<LC_{\rm s}\le LR(x_K)$.
This proves the required upper and lower bounds for both
sides.  Neither side depends on $K$ or $\beta$.
\end{proof}

\subsection{Errors depending only on the physical separation}

For $s\ge3$ set
\begin{equation}\label{IV-68-59}
\begin{gathered}
 Y_s^-=Y_s-\log(2L^3),\qquad
 x_s^-={4L^{s-2}g^{-2}}/{(5\mathsf L_s)},\\
 d_s^-=\min\left\{
 \frac{c_A\alpha_0\sqrt{x_s^-}}{2MB_3A_0Y_s^p},
 \frac{L^{s-1}}{2\mathsf L_s}-\frac2M\right\}-1.
\end{gathered}
\end{equation}
Where $x_s^-\ge e^{2p}$, put
\begin{equation}\label{IV-68-60}
\begin{split}
 \sigma_g(s)=\max\big\{&(x_s^-)^{-1/2},
 \bar x_se^{-d_s^-},\bar x_se^{3-4Y_s^{2p}/\sqrt3},
 \delta(x_s^-)+e^{\rm co}(x_s^-),\\
 &M^{-3}c'A_0^2Y_s^{2p}L^{3s}\bar x_s^2
                   e^{-c'A_0^2(\log x_s^-)^{2p}/2}\big\},
\end{split}
\end{equation}
and
\begin{equation}\label{IV-68-61}
\begin{split}
 \epsilon_g(s)={}&(C_2+C_3)(Y_s^-)^{-2p}
       +C_3s(Y_s^-)^{-5p}+C_6(Y_s^-)^{-6p}\\
 &+C_8[(LC_W(s))^3+1](Y_s^-)^{-9p}
       +4c_X^2(Y_s^-)^{-14p}\\
 &+C_Bs^4(LC_W(s))^6Y_s^{15p}\sigma_g(s).
\end{split}
\end{equation}
There is $s_1(g)\ge3$ such that for every $s\ge s_1(g)$,
\begin{equation}\label{IV-68-62}
 Y_s^->0,\qquad\mathsf L_s\le L^{s-3},\qquad x_s^-\ge e^{2p},\qquad
 \sigma_g(s)\le1.
\end{equation}
Indeed $Y_s$ grows linearly, $\mathsf L_s$ is a polynomial
of degree $5p$ in $Y_s$, and $x_s^-$ grows exponentially
divided by that polynomial.  Both entries defining $d_s^-$
grow faster than every polynomial.  The first and fourth
terms of $\sigma_g$ decay exponentially times a polynomial;
the remaining terms have a negative exponent whose magnitude
grows at least as $Y_s^{2p}$, with $2p>1$, overcoming
the factors exponential in $s$.  Thus
$s^a\sigma_g(s)\to0$ for every $a\ge0$, and
\begin{equation}\label{eq:depth-error-decay}
 \sup_{s\ge s_1(g)}Y_s^{2p}\epsilon_g(s)<\infty.
\end{equation}

\begin{proposition}\label{IV-68-67}
Let $g\le\gamma_W^{\rm ref}$, let $(K,\beta)$ be a
first-passage datum and let $\ell$ be small over its
canonical sequence.  If $s_1(g)\le s(d)\le K$ and
$(f,h)\in\mathfrak T_{T_\ell}(d;\vartheta)$, then
\begin{equation}\label{IV-68-68}
 \left|\langle\mathcal O(f);\mathcal O(h)\rangle_{K,\ell,\beta}
 -\tfrac{N^2-1}{2}\upsilon_{\mathsf k+1}^2
                 \mathcal K_{K,s_{\mathsf k}}(f,h)\right|
 \le\epsilon_g(s(d))\upsilon_{\mathsf k+1}^2
                                      \|f\|_\infty\|h\|_\infty.
\end{equation}
\end{proposition}

\begin{proof}
If a sampled weight vanishes, the observable and kernel
sum involving it both vanish.  Otherwise
\cref{IV-68-51,IV-68-56} verify all hypotheses of
\cref{IV-68-44} at the attached levels, with $T\le K-2$.
We check its five small factors rather than infer their
smallness from the final cutoff.
The first obeys $x_T^{-1/2}<(x_s^-)^{-1/2}$ by
\eqref{eq:working-top}.  For the second,
$\varepsilon_T\le A_0(x_s^-)^{-1/2}Y_s^p$ and
$L^{m+K-T}>ML^{s-1}/\mathsf L_s$ give
$d_*(T)\ge d_s^-$.
The separation estimate in \cref{IV-68-51} controls the
third.  Monotonicity of $\delta$ and $e^{\rm co}$ controls
the fourth since $x_i>x_T>x_s^-$ for $i<T$.
Finally smallness of the side and $B(x_K)\ge2$ imply
\[
 2M^{-3}\ell^3\le M^{-3}c'A_0^2y_K^{2p}.
\]
Using $K-\mathsf k=s_{\mathsf k}\le s$,
$y_K\le Y_s$, $y_T>\log x_s^-$ and
$x_{\mathsf k}\le\bar x_s$ bounds $x_{\mathsf k}^2t_T$
by the fifth entry of \eqref{IV-68-60}.
Thus $s_T\le\sigma_g(s)\le1$.
In \eqref{IV-68-46} use
$Y_s^-<y_{\mathsf k}\le Y_s$, $n_T\le s$ and
$q<LC_W(s)$.  Its bracket is at most $\epsilon_g(s)$,
and $\widehat m_1\widehat m_2
\le\upsilon_{\mathsf k+1}^2\|f\|_\infty\|h\|_\infty$.
\end{proof}

Put $c_2=3(N^2-1)/(16\pi^2)$ and define
\begin{equation}\label{IV-68-63}
\begin{split}
 \varrho_{g,\vartheta}(s)={}&3\pi^2\,10^6\vartheta^{-2}C_{\rm tm}
 \left[\frac1{\vartheta C_W(s)}+(LC_W(s)L^{-s})^3
       +\sup_{t\ge L^{-2}C_W(s)}t^3e^{-\delta_{\rm tp}Lt/\sqrt3}\right]\\
 &+3\bar Bg^2L^{2-s}
          +\tfrac98\,10^6\vartheta^{-2}c_2^{-1}\epsilon_g(s).
\end{split}
\end{equation}
Choose $s_0(g,\vartheta)\ge s_1(g)$ so that
$\varrho_{g,\vartheta}(s)\le1/2$ for every $s\ge s_0$,
and put
\begin{equation}\label{IV-68-64}
 d_0^{\rm ref}(g,\vartheta)=L^{-s_0(g,\vartheta)-1},\qquad
 \bar\varepsilon_{\rm ref}(d;g,\vartheta)
       =\sup_{s'\ge s(d)}\varrho_{g,\vartheta}(s').
\end{equation}
These quantities exist because \eqref{eq:depth-error-decay}
and the explicit formula give
$\sup_{s\ge s_0}Y_s^{2p}\varrho_{g,\vartheta}(s)<\infty$.
The function $\bar\varepsilon_{\rm ref}$ is non-decreasing
in $d$, takes values in $[0,1/2]$ on $(0,d_0^{\rm ref}]$,
and satisfies
\begin{equation}\label{IV-68-66}
 \bar\varepsilon_{\rm ref}(d;g,\vartheta)
                  \le C(g,\vartheta)(\log(1/d))^{-2p}.
\end{equation}
Indeed $d\le L^{-s_0-1}$ implies $s(d)>s_0$,
while $L^{s(d)}d\ge C_W(s(d))\ge1$ implies
$Y_{s(d)}\ge\log(1/d)$; apply the preceding uniform bound
at every $s'\ge s(d)$.

\begin{theorem}[Short-distance two-point law on a small torus]
\label{IV-68-69}
Let $0<g\le\gamma_W^{\rm ref}$ and $0<\vartheta<1$.
For every first-passage datum $(K,\beta)$ at $g$, with
$\mathsf b=a_K\beta$, every side small over its running
sequence, every $0<d\le d_0^{\rm ref}(g,\vartheta)$
with $s(d)\le K$, and every
$(f,h)\in\mathfrak T_{T_\ell}(d;\vartheta)$,
\begin{equation}\label{IV-68-71}
 \left|\langle\mathcal O(f);\mathcal O(h)\rangle_{K,\ell,\beta}
                -c_2\mathsf b^{-2}\mathcal K_{\rm free}(f,h)\right|
 \le\bar\varepsilon_{\rm ref}(d;g,\vartheta)
                  c_2\mathsf b^{-2}\mathcal K_{\rm free}(f,h).
\end{equation}
The constants and error function are independent of
$K,\beta,\ell$ and the test pair.
\end{theorem}

\begin{proof}
Set $s=s(d)$ and $\upsilon=\upsilon_{\mathsf k+1}$.
Subtract and add the lattice kernel with coefficient
$(N^2-1)\upsilon^2/2$, and then the free kernel with that
coefficient.  The three errors are
\[
 \epsilon_g(s)\upsilon^2\|f\|_\infty\|h\|_\infty,
 \quad\tfrac{N^2-1}{2}\upsilon^2
  |\mathcal K_{K,s_{\mathsf k}}-\tfrac3{8\pi^2}\mathcal K_{\rm free}|,
 \quad c_2\mathsf b^{-2}|(\mathsf b\upsilon)^2-1|
                                    \mathcal K_{\rm free}.
\]
The first is bounded by \cref{IV-68-67}.
To apply \cref{IV-45-38} to the second, observe that
$8a_K\le d$ by \eqref{eq:working-room} and
$20d\le\ell$ by \cref{IV-68-51}.  Its background-tail
condition holds because $q\ge4MY_s^{2p}$ and the last
inequality of \eqref{IV-68-54} holds at $x_W$.
That theorem therefore bounds the kernel error by
\[
 C_{\rm tm}\left[\frac1{\vartheta C_W(s)}
 +(LC_W(s)L^{-s})^3
 +\sup_{t\ge L^{-2}C_W(s)}t^3e^{-\delta_{\rm tp}Lt/\sqrt3}\right]
                                  \|f\|_\infty\|h\|_\infty.
\]
For the third use \eqref{eq:working-coupling}.
Finally $\mathcal K_{\rm free}\ge10^{-6}\vartheta^2
\|f\|_\infty\|h\|_\infty$ and
$\upsilon^2\le9\mathsf b^{-2}/8$.
After division by $c_2\mathsf b^{-2}\mathcal K_{\rm free}$,
the three bounds add to \eqref{IV-68-63}; here
$(9/8)(N^2-1)/(2c_2)=3\pi^2$.
Use \eqref{IV-68-64} to conclude.
\end{proof}

\begin{corollary}\label{IV-68-72}
If $0<g\le\gamma_W^{\rm sd}$, the conclusions
\eqref{IV-68-68} and \eqref{IV-68-71} hold for every
first-passage datum and every admitted test pair and depth
on the least torus at $g$, on the reference torus of side
$2MR(X_g)$, and on every side satisfying
$MR(x_K)\le L^m\le L^2MR(x_K)$.
\end{corollary}

\begin{proof}
These sides are small by \cref{IV-68-57}.
Apply \cref{IV-68-67,IV-68-69} with their respective
depth conditions.  In particular the reference torus
is the same for all first-passage data at $g$.
\end{proof}


\section{Transfer of the two-point law}
\label{sec:two-point-transfer}

The short-distance calculation on one torus can be transferred to
every admissible torus and every periodic Gibbs state. The required
comparison gains a factor equal to the volume of the support.
After accounting for the two Lipschitz norms, its error is of order
$d$ for test pairs at distance $d$.

Fix the source radius $\varrho$ and the construction constants of
\cref{III-111-27}, and let $0<g\le\gamma_\varrho$.
For $(K,\beta)\in\mathfrak d(g)$, a \emph{state at $g$} will mean
either $\widetilde\mu_{K,\ell,\beta}$ with $\ell\in\mathfrak L(g)$,
or a periodic Gibbs state at $(K,\beta)$.
Every such law is invariant under $a_K\mathbb Z^3$:
for the finite torus this follows by relabelling bonds, and
for a periodic Gibbs state it is \cref{IV-54-17}.
The covariance $S^T_2[\mathfrak s](f,h)$ is consequently invariant
under simultaneous lattice translations. Periodic lifting gives
\begin{equation}\label{IV-69-5}
 S^T_2[\widetilde\mu_{K,\ell,\beta}](f,h)
       =S^T_{2;K,\ell,\beta}(f^{\langle\ell\rangle},h^{\langle\ell\rangle})
\end{equation}
by \eqref{IV-17-14}.

Use the number $\mathsf Q_1$ of \eqref{IV-19-22}, which is
independent of $(K,\beta)$ and the torus side, and put
\begin{equation}\label{IV-51-14}
 c^{\rm tr}_2=4\log2\,(2\mathsf Q_1+2)^2\varrho^{-2},\qquad
 C^{\rm tr}_2(\vartheta)
       =2(50/3)^3c^{\rm tr}_2(40M/\vartheta)^2.
\end{equation}

\begin{proposition}[Comparison between states]\label{IV-51-20}\label{IV-51-26}
Let $\mathfrak s,\mathfrak s'$ be two states at $g$ for the same
$(K,\beta)\in\mathfrak d(g)$. If $a_K\le r\le1/20$ and
$f,h\in\mathrm{Lip}_c(\mathbb R^3)$ vanish off $\mathbb B(x_0,r)$,
then
\begin{equation}\label{IV-51-21}
 |S^T_2[\mathfrak s](f,h)-S^T_2[\mathfrak s'](f,h)|
       \le2c^{\rm tr}_2(r/0.3)^3\|f\|_\sharp\|h\|_\sharp.
\end{equation}
In particular, if $a_K\le5d\le1/20$, both functions vanish
off one ball of radius $5d$, and
$\operatorname{Lip}f\le\|f\|_\infty/(\vartheta d)$,
$\operatorname{Lip}h\le\|h\|_\infty/(\vartheta d)$ with
$0<\vartheta\le1$, then
\begin{equation}\label{IV-51-22}
 |S^T_2[\mathfrak s](f,h)-S^T_2[\mathfrak s'](f,h)|
       \le C^{\rm tr}_2(\vartheta)d\|f\|_\infty\|h\|_\infty.
\end{equation}
No separation or sign condition on $f,h$ is required.
\end{proposition}
\begin{proof}
First let $F,H$ vanish off $\mathbb B(0,1)$. The board of side
$L$ centered at zero has walls at distance $L/2\ge6.5$ from
zero. If neither entry is zero, \cref{IV-19-24} gives a single
normalizer for $(F,H)$, of radius $\varrho$ and bound
$\mathsf Q_1$, on the unbounded set $\mathfrak L(g)$.
Both entries occupy the same cell, so the distinct-cell
condition of \cref{IV-19-13} is vacuous. That proposition applies
to all the torus laws and all periodic Gibbs states. In
\eqref{IV-19-17}, the normalizer derivative is the same for
both laws and cancels on subtraction. Thus the bilinear form
$D=S^T_2[\mathfrak s]-S^T_2[\mathfrak s']$ satisfies
\begin{equation}\label{IV-51-19}
 |D(F,H)|\le2c^{\rm tr}_2\|F\|_\sharp\|H\|_\sharp.
\end{equation}
If an entry is zero the same statement follows by bilinearity.

We pack translates of $\mathbb B(x_0,r)$ into this unit ball
using lattice translations. Set $\delta=a_K\lceil2r/a_K\rceil$,
so $2r\le\delta<3r$ and $\delta\mathbb Z^3\subset a_K\mathbb Z^3$.
Let $P$ consist of $v\in\delta\mathbb Z^3$ such that every
coordinate of $x_0+v$ has modulus at most $0.575-r$.
The balls $\mathbb B(x_0+v,r)$ have disjoint interiors and lie
in the unit ball, since $3(0.575)^2<1$.
Each interval of length $1.15-2r\ge1.05$ contains at least
$\lfloor(1.15-2r)/\delta\rfloor>0.35/r-1\ge0.3/r$ points
of a prescribed translate of $\delta\mathbb Z$.
Taking the product over the coordinates yields
\[
 N:=\#P\ge(0.3/r)^3.
\]

For independent signs $\epsilon_v=\pm1$ define
$F_\epsilon=\sum_{v\in P}\epsilon_v\tau_vf$ and
$H_\epsilon=\sum_{v\in P}\epsilon_v\tau_vh$.
Their supremum norms are no larger than those of $f,h$, and
their Lipschitz constants are no larger either. To verify the
latter assertion, note that each function vanishes on the
boundary of its supporting ball. On a segment joining points
in different balls, bound the endpoint values by the distances
to the respective exit and entry points. Their sum is at most
the segment length. The same-ball and outside-ball cases give
the same Lipschitz bound. Thus
$\|F_\epsilon\|_\sharp\le\|f\|_\sharp$ and
$\|H_\epsilon\|_\sharp\le\|h\|_\sharp$.
Bilinearity and $\mathbb E(\epsilon_v\epsilon_w)=\mathbf1_{v=w}$ give
\[
 \mathbb E D(F_\epsilon,H_\epsilon)
       =\sum_{v\in P}D(\tau_vf,\tau_vh)=N D(f,h),
\]
where the last equality uses lattice translation invariance.
Apply \eqref{IV-51-19} to every choice of signs and divide by
$N$ to prove \eqref{IV-51-21}.
For the last assertion put $r=5d$. Since $d\le1/100$,
$\vartheta\le1$ and $M\ge1$, each sharp norm is bounded by
$40M/(\vartheta d)$ times the corresponding supremum norm.
Substitution in \eqref{IV-51-21} gives \eqref{IV-51-22}.
\end{proof}



\subsection{The estimate on every admissible volume}

Write $\mathfrak T_E^1(d;\vartheta)$ for the pairs in
$\mathfrak T_E(d;\vartheta)$ whose entries are $C^1$.
For $b>0$, say that a covariance $Y$ obeys the two-point estimate
at $(f,h)$ with label $b$ and tolerance $\epsilon$ if
\begin{equation}\label{IV-69-3}
 |Y-\mathsf c_2b^{-2}\mathcal K_{\rm free}(f,h)|
       \le\epsilon\mathsf c_2b^{-2}\mathcal K_{\rm free}(f,h),
 \qquad \mathsf c_2=\frac{3(N^2-1)}{16\pi^2}=\frac9{16\pi^2}.
\end{equation}
Here $N=2$ is fixed by the gauge group. The positive lower bound
in \eqref{IV-45-9} makes this equivalent to the relative-error
form of the short-distance calculation.

\begin{lemma}[Lifting and dilation of test pairs]\label{IV-69-4}\label{IV-69-9}
If $20d<\ell$, periodization takes a pair in
$\mathfrak T_{\mathbb R^3}(d;\vartheta)$ to a pair in
$\mathfrak T_{T_\ell}(d;\vartheta)$ with the same free kernel.
Conversely, every pair in the latter set has such a lift.
Both assertions preserve $C^1$ regularity.
For $f_t(x)=f(x/t)$, $h_t(x)=h(x/t)$, $t>0$,
\begin{equation}\label{IV-69-10}
 (f_t,h_t)\in\mathfrak T_{\mathbb R^3}(td;\vartheta),\qquad
 \mathcal K_{\rm free}(f_t,h_t)=\mathcal K_{\rm free}(f,h).
\end{equation}
If $n\ge0$, $\ell=L^n\ell'$ and $t=L^{-n}$, then
\[
 S^T_2[\widetilde\mu_{K+n,\ell',\beta}](f_t,h_t)
       =L^{2n}S^T_2[\widetilde\mu_{K,\ell,\beta}](f,h).
\]
Thus the estimate with label $L^{-n}b$ on the left transfers,
with the same tolerance, to the estimate with label $b$ on the right.
\end{lemma}
\begin{proof}
The two supports lie in one ball of radius $5d$. Since
$10d<\ell/2$, that ball lifts isometrically to $\mathbb R^3$;
all distances between points in the two supports are preserved.
The entries vanish before the boundary of a larger lifting chart,
so lifting and periodizing preserve supremum norms, Lipschitz
constants, integrals and $C^1$ regularity. This proves the first
assertions and the equality of the free kernels. Under dilation,
the integrals gain $t^3$, the Lipschitz constants gain $t^{-1}$,
and distances gain $t$. In the free kernel the factor $t^6$
from the two integrations cancels $t^{-6}$ from the kernel.
Finally, \eqref{IV-23-17} and \eqref{IV-23-20}, with $z=-n$,
give the displayed covariance identity. Since
$(L^{-n}b)^{-2}=L^{2n}b^{-2}$, division by $L^{2n}$ proves
the last assertion.
\end{proof}

We use the room constant of \eqref{IV-68-48} and the witness
depth of \eqref{IV-49-6}. To make their dependence on the
reference coupling explicit, write
\[
 \begin{gathered}
 \bar x_s[g]=L^{s+1}(g^{-2}+\bar B),\qquad
 C_W(s)[g]=C_{W,0}A_0^2(\log\bar x_s[g])^{2\mathfrak p_0},\qquad
 \\
 s(d)[g]=\min\{s\ge0:L^sd\ge C_W(s)[g]\},
 \end{gathered}
\]
where the minimum is over integers. The exponential $L^s$
eventually exceeds the polynomial in $s$ on the right, so this
minimum is finite. We abbreviate $s(d)[g]$ to $s(d)$ when the
reference coupling is fixed. Recall that
$C_{W,0}A_0^2\ge4L^2M$ and $X_g=L(g^{-2}+B_\sharp)$.

For $0<g\le\gamma$, $0<\vartheta<1$ and $0<d\le1/100$,
let $\mathsf e_{g,\vartheta}(d)$ be the infimum of the tolerances
for which \eqref{IV-69-3} holds on the torus
$\ell=2M R(X_g)$, for every $(K,\beta)\in\mathfrak d(g)$
with $s(d)\le K$ and every pair of type $(d;\vartheta)$,
with label $a_K\beta$. The infimum of the empty set is infinity.
If the infimum is finite, it itself is an admissible tolerance:
for every datum and pair the defining inequality is closed in
$\epsilon$ and has a positive coefficient on its right-hand side.

By \cref{IV-68-72}, the set
\[
 \mathfrak G=\{\gamma'\in(0,\gamma]:
     \mathsf e_{g,\vartheta}(d)\longrightarrow0\text{ as }d\to0
     \text{ for every }0<g\le\gamma',\ 0<\vartheta<1\}
\]
is nonempty. It is downward closed. Put
\begin{equation}\label{IV-69-13}
 \widehat\gamma_W=\tfrac12\sup\mathfrak G.
\end{equation}
An element of $\mathfrak G$ exceeds $\widehat\gamma_W$, so
$\widehat\gamma_W\in\mathfrak G$ and
$0<\widehat\gamma_W\le\gamma/2$. Define
\begin{equation}\label{IV-69-14}
 \gamma_W=
 \left[L^2\left(\bar B+
       \max\{\widehat\gamma_W^{-2},\gamma^{-4},e^{C_{\rm s}}\}
                 \right)\right]^{-1/2},\qquad
 \check g=(L^{-2}g^{-2}-\bar B)^{-1/2}\quad(0<g\le\gamma_W).
\end{equation}
These choices give
\begin{equation}\label{IV-69-15}
 \begin{gathered}
 \check g^{-2}+\bar B=L^{-2}g^{-2},\qquad
 X_{\check g}=L^{-1}g^{-2}-\bar B,\qquad
 \check g\le\widehat\gamma_W\le\gamma,\\
 \check g^{-2}\ge\gamma^{-4},\qquad
 \log\check g^{-2}\ge C_{\rm s},\qquad
 g^{-2}>L^2\bar B,\qquad g<\check g,\qquad X_{\check g}\le X_g.
 \end{gathered}
\end{equation}
Indeed, $\check g^{-2}$ is at least each term in the maximum
in \eqref{IV-69-14}; the identity for $X_{\check g}$ uses
$X_u+\bar B=L(u^{-2}+\bar B)$. In particular
$0<\gamma_W\le\gamma$.

\begin{lemma}[Changing the reference coupling]\label{IV-69-19}
Let $0<g\le\gamma_W$ and $(K,\beta)\in\mathfrak d(g)$.
Then $K'=K(\beta,\check g)$ exists and $n=K'-K\in\{2,3\}$.
The datum $(K+n,\beta)$ belongs to $\mathfrak d(\check g)$
and its label is $L^{-n}a_K\beta$. Moreover,
\begin{equation}\label{eq:carried-side-depth}
 L^n2M R(X_{\check g})\in\mathfrak L(g)\quad(n\ge2),\qquad
 s(L^{-n}d)[\check g]\le s(d)[g]+n+1
       \quad(0<d<1,\ n=2,3).
\end{equation}
\end{lemma}
\begin{proof}
Since $\beta>X_g\ge X_{\check g}$, \cref{III-5-11} gives
both first-passage indices, with identical running couplings
on their common prefix. By \eqref{III-5-12} and
\eqref{III-5-13},
\[
 \begin{gathered}
 g^{-2}<x_K\le X_g,\qquad x_{K'}\le X_{\check g}<x_j\ (j<K'),\\
 |x_j-L^{-(j-i)}x_i|
        \le\bar B(1-L^{-(j-i)})\quad(0\le i\le j\le K').
 \end{gathered}
\]
If $K'\le K$, monotonicity on the common prefix gives
$x_{K'}\ge x_K>g^{-2}>X_{\check g}$, a contradiction.
At the next step,
$x_{K+1}\ge L^{-1}x_K-\bar B(1-L^{-1})>X_{\check g}$,
so $K'\ne K+1$. If $K'\ge K+4$, then
\[
 X_{\check g}<x_{K+3}\le L^{-3}X_g+\bar B
       =L^{-2}(g^{-2}+B_\sharp)+\bar B.
\]
As $B_\sharp\le\bar B$, this implies
$(L-1)g^{-2}<(2L^2+1)\bar B\le L^2(L-1)\bar B$,
contrary to \eqref{IV-69-15}. This proves the index assertion.

For the side assertion put $X=X_{\check g}$ and
$X'=X_g+\bar B$. Then $X\ge\gamma^{-4}\ge16$ and
$\log X\ge2\log\gamma^{-2}$. Using
$\bar B=B_\sharp L/(L-1)$ and $L\ge13$ gives
\[
 X'=L^3\check g^{-2}+(L^3+L)\bar B
       \le4L^3(\check g^{-2}+B_\sharp)=4L^2X.
\]
The standing calibration
$(L^{1/r}-1)\log\gamma^{-2}\ge\log(2L)$ therefore yields
$\log X'\le L^{1/r}\log X$. Since $R$ is the smallest
power of $L$ not less than the $r$th power of the logarithm,
\[
 R(X')\le L R(X),\qquad
 R(X)\ge(\log X)^r\ge\log X\ge\log\check g^{-2}\ge C_{\rm s}.
\]
If $R(X)=L^u$, the definition of $m_{\mathfrak L}(g)$ gives
$m_{\mathfrak L}(g)\le m'+u+2$. Hence
$L^n2M R(X)=2L^{n+m'+u}\in\mathfrak L(g)$ for every $n\ge2$.

Finally, let $s=s(d)[g]$ and $Y=\log\bar x_s[g]$.
Then $L^sd\ge C_W(s)[g]$ and
$Y\ge2\log\gamma^{-2}>4\mathfrak p_0$ by the standing
calibration $\log\gamma^{-2}>2\mathfrak p_0$.
For $n=2,3$,
\[
 \bar x_{s+n+1}[\check g]=L^{s+n}g^{-2}\le L^2\bar x_s[g],
\]
and consequently
\[
 (\log\bar x_{s+n+1}[\check g])^{2\mathfrak p_0}
 \le Y^{2\mathfrak p_0}
       \exp(4\mathfrak p_0\log L/Y)
 \le L Y^{2\mathfrak p_0}.
\]
It follows that
$C_W(s+n+1)[\check g]\le L C_W(s)[g]
\le L^{s+n+1}L^{-n}d$. Thus $s+n+1$ is an admissible
integer in the definition of $s(L^{-n}d)[\check g]$.
\end{proof}

For $0<g\le\gamma_W$ and $0<\vartheta<1$, put
\begin{equation}\label{IV-69-17}
 \begin{split}
 \mathsf C_{\rm law}(g,\vartheta)
   &=\frac{10^6 C^{\rm tr}_2(\vartheta)\bar x_0[g]^2}
                       {\mathsf c_2\vartheta^2},\\
 \bar\varepsilon_{g,\vartheta}(d)
   &=\sup_{0<d'\le L^{-2}d}\mathsf e_{\check g,\vartheta}(d')
                   +\mathsf C_{\rm law}(g,\vartheta)d,
       \qquad 0<d\le1/100,
 \end{split}
\end{equation}
and
\begin{equation}\label{IV-69-18}
 d_0(g,\vartheta)=\tfrac12\sup\{d\in(0,1/100]:
                         \bar\varepsilon_{g,\vartheta}(d)\le1/3\}.
\end{equation}
Since $\check g\le\widehat\gamma_W$, the supremum in
\eqref{IV-69-17} tends to zero with $d$. The tolerance is
nondecreasing and tends to zero. The set in \eqref{IV-69-18}
is therefore nonempty, $0<d_0(g,\vartheta)\le1/200$, and
$\bar\varepsilon_{g,\vartheta}(d)\le1/3$ for $d\le d_0$:
there is an element of that set strictly exceeding its half-supremum,
and monotonicity gives the claim. All these quantities are independent
of cutoff, volume, state, pair and position.

\begin{proposition}[An admissible reference torus]\label{IV-69-20}
Let $0<g\le\gamma_W$, $(K,\beta)\in\mathfrak d(g)$,
$0<\vartheta<1$ and $0<d\le1/100$, with $s(d)<K$.
Put $n=K(\beta,\check g)-K$, $\ell'=2M R(X_{\check g})$
and $\ell=L^n\ell'$. Then $\ell\in\mathfrak L(g)$.
If $\epsilon_0=\sup_{0<d'\le L^{-2}d}
\mathsf e_{\check g,\vartheta}(d')$ is finite, the lifted
torus law at $(K,\ell,\beta)$ obeys \eqref{IV-69-3} for
every pair of type $(d;\vartheta)$, with label $a_K\beta$
and tolerance $\epsilon_0$.
\end{proposition}
\begin{proof}
The side assertion is \cref{IV-69-19}. Put $d'=L^{-n}d$.
Since $n\in\{2,3\}$, $d'\le L^{-2}d$ and
\[
 s(d')[\check g]\le s(d)[g]+n+1\le K+n.
\]
The finite number $\mathsf e_{\check g,\vartheta}(d')$
is an admissible tolerance for $(K+n,\beta)\in\mathfrak d(\check g)$
on $\ell'$. The lifting and dilation statements of
\cref{IV-69-4} transfer this estimate to $(K,\ell,\beta)$
at distance $d$, with label $a_K\beta$. Its tolerance does
not exceed $\epsilon_0$.
\end{proof}

\begin{theorem}[Uniform two-point estimate]\label{IV-69-21}
Let $0<g\le\gamma_W$, $0<\vartheta<1$,
$(K,\beta)\in\mathfrak d(g)$, and
$0<d\le d_0(g,\vartheta)$ with $s(d)<K$.
For every state $\mathfrak s$ at $g$ and every
$(f,h)\in\mathfrak T_{\mathbb R^3}(d;\vartheta)$,
\begin{equation}\label{IV-69-22}
 S^T_2[\mathfrak s](f,h)
       =\mathsf c_2b^{-2}\mathcal K_{\rm free}(f,h)(1+\eta),
 \qquad b=a_K\beta,\qquad
 |\eta|\le\bar\varepsilon_{g,\vartheta}(d)\le1/3.
\end{equation}
The same assertion holds on every torus $\ell\in\mathfrak L(g)$
for pairs in $\mathfrak T_{T_\ell}(d;\vartheta)$.
In particular every such covariance is positive.
\end{theorem}
\begin{proof}
The reference state of \cref{IV-69-20} is available, since
the first term of \eqref{IV-69-17} is finite. The witness
inequality and $C_{W,0}A_0^2\ge4L^2M$ give
$a_K<a_{s(d)}<d$, whereas $5d\le1/40<1/20$.
The pair is supported in one ball of radius $5d$, so
\eqref{IV-51-22} bounds the difference from the reference
covariance by $C^{\rm tr}_2(\vartheta)d\|f\|_\infty\|h\|_\infty$.
Dividing by the lower bound in \eqref{IV-45-9} shows that
the added relative error is at most
\[
 \frac{10^6 C^{\rm tr}_2(\vartheta)b^2d}
                    {\mathsf c_2\vartheta^2}
       \le\mathsf C_{\rm law}(g,\vartheta)d,
\]
because $b\le\bar x_0[g]$ by \eqref{IV-22-4}.
This is precisely the remaining term of \eqref{IV-69-17}.
Lifting proves the torus assertion. Finally $1+\eta\ge2/3$
and the free kernel is positive.
\end{proof}



\subsection{The label in continuum limits}

We first record that smooth test pairs exist for every allowed shape
parameter, including pairs separated by the reflection plane.

\begin{lemma}[Smooth reflection pairs]\label{lem:smooth-reflection-pair}
For every $0<\vartheta<1$ and $d>0$ there is
$f\in C_c^\infty(\mathbb R^3)$ supported in
$\mathbb B(2de_3,d)$ such that $(\theta f,f)$ belongs to
$\mathfrak T_{\mathbb R^3}^1(d;\vartheta)$, where
$\theta(x_1,x_2,x_3)=(x_1,x_2,-x_3)$.
\end{lemma}
\begin{proof}
Put $x=2de_3$ and
\[
 r_0=\max\{\vartheta,(3\vartheta/\pi)^{1/3}\},\quad
 r=(1+r_0)/2,\quad
 \epsilon=\tfrac12\min\{1-r,r-\vartheta\}.
\]
Then $r_0<r<1$, $0<\epsilon<r$,
$r+\epsilon<1$, $r-\epsilon>\vartheta$ and
$(\pi/3)r^3>\vartheta$.
Convolve $g_r(y)=\max\{0,1-|y-x|/(rd)\}$ with a nonnegative
smooth mollifier of integral one supported in $\mathbb B(0,\epsilon d)$.
The resulting function $f$ satisfies
\[
 0\le f\le1,\quad f(x)\ge1-\epsilon/r>0,\quad
 \operatorname{Lip}f\le(rd)^{-1},\quad
 \int f=(\pi/3)r^3d^3,\quad
 \operatorname{supp}f\subset\mathbb B(x,(r+\epsilon)d).
\]
The integral follows by polar integration of the cone; convolution
preserves it. Since $\|f\|_\infty\ge(r-\epsilon)/r>\vartheta/r$,
the required Lipschitz inequality holds, as does
$\int f\ge\vartheta\|f\|_\infty d^3$.
Reflection preserves these bounds. The two supports lie in balls
of radius $d$ centered at $\pm x$, so their distance is at least
$2d$; it is at most $4d$ because $f(x)>0$.
All conditions defining the test pair are satisfied.
\end{proof}

\begin{lemma}[Limits and recovery of the label]\label{IV-69-27}
For fixed $(f,h)$, inequality \eqref{IV-69-3} passes to the
limit if $Y_j\to Y$ and $b_j\to b>0$, provided it holds
along an infinite subsequence with a fixed tolerance.

More generally, fix $(f,h)\in\mathfrak T_{\mathbb R^3}(d_1;\vartheta)$,
$t_i\to0$ with $t_i>0$, and $\epsilon_i\to0$ with
$0\le\epsilon_i<1$. Suppose for each $i$ that the covariances
$Y_{j,i}$ of $(f_{t_i},h_{t_i})$ converge to a finite $y_i$,
and for all sufficiently large $j$ obey \eqref{IV-69-3}
with label $b_j>0$ and tolerance $\epsilon_i$.
Then
\begin{equation}\label{IV-69-28}
 \lim_j b_j^{-2}
       =\lim_i\frac{y_i}{\mathsf c_2\mathcal K_{\rm free}(f,h)}
       =u\in[0,\infty).
\end{equation}
The limit $u$ is positive if and only if $\sup_j b_j<\infty$;
in this case $b_j\to u^{-1/2}$.
\end{lemma}
\begin{proof}
The first assertion follows by taking limits on both sides
of \eqref{IV-69-3} along the stated subsequence.
For the second put $k=\mathsf c_2\mathcal K_{\rm free}(f,h)>0$
and $u_j=b_j^{-2}$. Dilation invariance of the free kernel gives,
for each fixed $i$ and all sufficiently large $j$,
\[
 (1-\epsilon_i)ku_j\le Y_{j,i}\le(1+\epsilon_i)ku_j.
\]
Taking an upper limit in the left inequality and a lower limit
in the right yields
\[
 (1-\epsilon_i)k\limsup_j u_j
       \le y_i\le(1+\epsilon_i)k\liminf_j u_j.
\]
The first inequality makes the upper limit finite. Letting
$i\to\infty$ shows that the upper and lower limits coincide,
at a number $u\ge0$. The same inequalities then give $y_i\to ku$.
If the $b_j$ are bounded above, $u_j$ have a positive lower
bound and $u>0$. Conversely a positive limit of the positive
$u_j$ gives a positive lower bound for the entire sequence,
so $\sup_j b_j<\infty$ and $b_j\to u^{-1/2}$.
\end{proof}

\begin{theorem}[The two-point estimate in the continuum]\label{IV-69-29}
Let $0<g\le\gamma_W$. Let $S$ be a limit at $g$ on a fixed
torus $T_\ell$ with $\ell\in\mathfrak L(g)$, a limit along
growing tori, or a limit through periodic Gibbs states.
For the fixed torus, defining sequences at $g$ are required to
consist of first-passage data in $\mathfrak d(g)$.
Write $E=T_\ell$ in the first case and $E=\mathbb R^3$ otherwise.
Then there is a single label $b(S)\in I^{\rm lab}(g)$, and
$a_{K_j}\beta_j\to b(S)$ along every defining sequence at $g$.
For every $0<\vartheta<1$, $0<d\le d_0(g,\vartheta)$ and
$(f,h)\in\mathfrak T_E^1(d;\vartheta)$,
\begin{equation}\label{IV-69-30}
 S_2(f,h)=\mathsf c_2b(S)^{-2}\mathcal K_{\rm free}(f,h)(1+\eta),
 \qquad |\eta|\le\bar\varepsilon_{g,\vartheta}(d)\le1/3.
\end{equation}
In particular,
\begin{equation}\label{eq:continuum-two-point-positive}
 S_2(f,h)\ge\tfrac23\,10^{-6}\vartheta^2\mathsf c_2b(S)^{-2}
                         \|f\|_\infty\|h\|_\infty>0.
\end{equation}
For any $(f,h)\in\mathfrak T_{\mathbb R^3}^1(d_0(g,\vartheta);\vartheta)$,
\begin{equation}\label{IV-69-31}
 b(S)^{-2}=\lim_{t\to0}
       \frac{S_2\langle f_t,h_t\rangle}
                    {\mathsf c_2\mathcal K_{\rm free}(f,h)}.
\end{equation}
Here $S_2\langle f,h\rangle=S_2(f,h)$ on $\mathbb R^3$,
and $S_2\langle f,h\rangle=
S_2(f^{\langle\ell\rangle},h^{\langle\ell\rangle})$ on the torus.
Thus the two-point function determines the label. If the same
family on the same space is also a limit at
$0<g_1\le\gamma_W$, its label at $g_1$ equals $b(S)$.
\end{theorem}
\begin{proof}
Given any defining sequence, use its lifted torus laws or its
periodic Gibbs laws, denoted by $\mathfrak s_j$. For every fixed
$C^1$ test pair, their covariances converge to
$S_2\langle f,h\rangle$. On a fixed torus this follows from
\eqref{IV-69-5} and \cref{IV-69-4}; in the other two cases
it is the definition of convergence on separated tuples.
For each fixed $d>0$, $s(d)<K_j$ eventually. Hence
\cref{IV-69-21} applies with tolerance
$\bar\varepsilon_{g,\vartheta}(d)$, while all labels
$b_j=a_{K_j}\beta_j$ lie in the compact positive interval
$I^{\rm lab}(g)$.

Choose one smooth test pair at $d_1=d_0(g,\vartheta)$ using
\cref{lem:smooth-reflection-pair}, and choose any sequence
$t_i\to0$ in $(0,1]$. The errors
$\epsilon_i=\bar\varepsilon_{g,\vartheta}(t_id_1)$ are at most
$1/3$ and tend to zero. The second assertion of
\cref{IV-69-27} shows that $b_j$ converges to a positive $b$
and recovers $b^{-2}$ from these limiting covariances.
That recovery uses only $S_2$ and the fixed pair, so it gives
the same $b$ for every defining sequence. Compactness of
$I^{\rm lab}(g)$ puts this number in that interval.

The first assertion of \cref{IV-69-27} now passes the
finite-cutoff estimate to every fixed test pair. On the torus,
lift an arbitrary test pair by \cref{IV-69-4}. This proves
\eqref{IV-69-30}; its lower bound and \eqref{IV-45-9} give
\eqref{eq:continuum-two-point-positive}.
Apply \eqref{IV-69-30} to the dilates of any pair at $d_1$.
Since their free kernels agree and their error tends to zero,
this gives the full limit \eqref{IV-69-31}.
If a second coupling $g_1$ is available, rescale a single
pair to a radius not exceeding either threshold. The two
recovery formulas then have identical right-hand sides,
so the positive labels coincide.
\end{proof}

The restriction to a small reference coupling does not restrict
the separation of positive labels on a fixed torus. Dilation
puts any two positive labels into the range just treated.

\begin{corollary}[Positive labels distinguish torus limits]
\label{cor:positive-torus-labels}
Let $\ell=2L^m$, $m\in\mathbb Z$. A limit point of the cutoff
cumulants on $T_\ell$ has at most one positive label, allowing
all defining sequences with $K_j\to\infty$ and
$a_{K_j}\beta_j\to b>0$. Consequently two such limit points
with different positive labels are different.
\end{corollary}
\begin{proof}
We give the first-passage reduction explicitly. For $b>0$ put
$g_z=(L^zb/2)^{-1/2}$. For all sufficiently large integers
$z\ge0$, $g_z\le\gamma_W$. Since
$\gamma^{-2}>\bar B$, this implies
\[
 g_z^{-2}+\bar B<L^zb<X_{g_z}-\bar B
       =(L/2)L^zb+(L-2)\bar B.
\]
Thus $L^zb$ lies in the interior of $I(g_z)$.
We also have $L^z\ell\in\mathfrak L(g_z)$ for all sufficiently
large $z$. To check the volume condition, set
$c=(L/2)b+L\bar B>0$ and
$u=L(g_z^{-2}+\bar B)\le cL^z$.
The inequality $\log u\le2r u^{1/(2r)}$ gives
\[
 (\log u)^r\le(2r)^rc^{1/2}L^{z/2}.
\]
Choose $z$ so large that
\[
 L^{z/2}\ge L^{2-m}M\max\{(2r)^rc^{1/2},C_{\rm s}\}.
\]
Then $L^{m+z}\ge L^2M\max\{(\log u)^r,C_{\rm s}\}$.
From \eqref{III-1-14}, $R(u)\le L(\log u)^r$ and
$L^{\lceil\log_L C_{\rm s}\rceil}\le LC_{\rm s}$,
the least admissible side is at most
$2L^2M\max\{(\log u)^r,C_{\rm s}\}$.
This proves the asserted membership.

Suppose now that $S$ has label $b$ along $(K_j,\beta_j)$.
For fixed $z$ as above, $L^za_{K_j}\beta_j\to L^zb$
lies eventually in $I(g_z)$, and $K_j-z\ge1$ eventually.
By \eqref{III-5-20}, $(K_j-z,\beta_j)$ is then a
first-passage datum at $g_z$. Exact finite dilation
\eqref{IV-12-5} gives a continuum limit
\[
 S^{[z]}_n(Y)=L^{-zn}S_n(Y^{[-z]})
       \quad\text{on }T_{L^z\ell},
\]
whose label at $g_z$ is $L^zb$.
If $S$ also has positive label $b'$, choose one sufficiently
large $z$ that works for both labels. The same family
$S^{[z]}$ then has labels $L^zb$ and $L^zb'$ at the two
reference couplings $g_z$ and $g'_z$, both at most $\gamma_W$.
The coupling independence in \cref{IV-69-29} gives
$L^zb=L^zb'$, hence $b=b'$. The last assertion follows
by applying this conclusion to a putatively equal pair of limits.
\end{proof}

\subsection{Periodic-state limits and their components}

The covariance estimate is unaffected by the centering law.
This observation is needed because the component approximation
uses one law for integration and another for centering.

\begin{proposition}[Transfer with an independent centering law]
\label{IV-69-36}
Let $0<g\le\gamma_W$, $(K_j,\beta_j)\in\mathfrak d(g)$,
$K_j\to\infty$ and
$\rho_j\in\mathcal G^{\rm per}_{K_j,\beta_j}$.
If the covariances of $\rho_j$ converge on every separated
$C^1$ pair, their limit has a unique label
$b=\lim_j a_{K_j}\beta_j\in I^{\rm lab}(g)$ and satisfies
\eqref{IV-69-30}, \eqref{eq:continuum-two-point-positive}
and \eqref{IV-69-31}.
In particular, this holds for the second cumulant of any
functional
\[
 \sigma(X)=\lim_j\rho_j(X_{K_j}^{\nu_j})
       \quad\text{for every separated tuple }X,
\]
where $\nu_j$ are arbitrary laws on the corresponding lattices.
The label depends only on the sequence $(K_j,\beta_j)$.
\end{proposition}
\begin{proof}
The proof of \cref{IV-69-29} used only convergence of the
covariances, the uniform estimate and the compact label interval;
it therefore proves the first assertion.
For the second, write $A_j(f)=\mathcal O_{K_j}(f)$ and
$c_j(f)=\nu_j(A_j(f))$. At each cutoff the exact identity
\[
 \rho_j((A_j(f)-c_j(f))(A_j(h)-c_j(h)))
 -\rho_j(A_j(f)-c_j(f))\rho_j(A_j(h)-c_j(h))
       =S_2^T[\rho_j](f,h)
\]
cancels all centering terms. The assumed convergence on the
one- and two-entry tuples makes its left side tend to
$\sigma(f,h)-\sigma(f)\sigma(h)=S_2^\sigma(f,h)$.
Thus the first assertion applies. In particular the recovered
label is the limit of the prescribed sequence of physical labels.
\end{proof}

\begin{lemma}[A vector orthogonal to the vacuum]\label{IV-69-33}
Let $\sigma$ be a normalized real functional with reflection
positivity at $x_3=0$. If $f\in C_c^1(\{x_3>0\})$ and
$S_2^\sigma(\theta f,f)>0$, its reflection Hilbert space
has dimension at least two.
\end{lemma}
\begin{proof}
Write $\Omega=[\varnothing]$, so $\|\Omega\|=1$.
The positive reflection form is Hermitian, and $\sigma$ is real,
so $\sigma(\theta f)=\sigma(f)=c\in\mathbb R$.
The vector $v=[(f)]-c\Omega$ therefore satisfies
\[
 \langle\Omega,v\rangle=0,\qquad
 \|v\|^2=\sigma(\theta f,f)-c^2=S_2^\sigma(\theta f,f)>0.
\]
Thus $\Omega$ and $v$ are nonzero orthogonal vectors.
\end{proof}

\begin{corollary}[Nontriviality]\label{IV-69-35}
For every limit on $\mathbb R^3$ covered by \cref{IV-69-29},
the centered functional $\sigma_S$ has a reflection Hilbert
space of dimension at least two.
\end{corollary}
\begin{proof}
The centered functional is translation invariant and reflection
positive at a coordinate mirror by \cref{IV-24-23}; for periodic
Gibbs limits its hypotheses follow from \cref{IV-54-17}.
Take the smooth reflection pair of
\cref{lem:smooth-reflection-pair} with $d=d_0(g,\vartheta)$.
Its covariance is positive by
\eqref{eq:continuum-two-point-positive}. Since
$S_2^{\sigma_S}=S_2$, \cref{IV-69-33} applies.
\end{proof}

\begin{corollary}[The same label in every component]\label{IV-69-39}
Let $0<g\le\gamma_W$, $(K_j,\beta_j)\in\mathfrak d(g)$,
$K_j\to\infty$ and
$\nu_j\in\mathcal G^{\rm per}_{K_j,\beta_j}$. Suppose the
frame $\mathfrak j=(K_j,\nu_j)_j$ converges to $\sigma$.
Then $a_{K_j}\beta_j\to b\in I^{\rm lab}(g)$, and every
limit $\sigma'$ in this frame satisfies the two-point estimate
with this same label and the same tolerance
$\bar\varepsilon_{g,\vartheta}(d)$.
For every canonical decomposition
$(\mathsf X,\lambda,(\sigma_\xi)_{\xi\in\mathsf X})$ of $\sigma$,
every $\xi\in\mathsf X$, every $0<\vartheta<1$,
$0<d\le d_0(g,\vartheta)$ and
$(f,h)\in\mathfrak T_{\mathbb R^3}^1(d;\vartheta)$,
\begin{equation}\label{IV-69-40}
 S_2^{\sigma_\xi}(f,h)
     =\mathsf c_2b^{-2}\mathcal K_{\rm free}(f,h)(1+\eta_\xi),
 \qquad |\eta_\xi|\le\bar\varepsilon_{g,\vartheta}(d)\le1/3.
\end{equation}
Every component recovers $b$ by \eqref{IV-69-31}, has the
positive lower bound \eqref{eq:continuum-two-point-positive},
and has a reflection Hilbert space of dimension at least two.
For a volume-first limit $S$, this label is $b(S)$.
\end{corollary}
\begin{proof}
By \cref{prop:gibbs-frame-classes}, the periodic Gibbs classes
form a class system with the required uniform bound system.
Every limit in the frame has approximants
$\rho_j\in\mathcal G^{\rm per}_{K_j,\beta_j}$ with
$\sigma'(X)=\lim_j\rho_j(X_{K_j}^{\nu_j})$ along the full
sequence. Apply \cref{IV-69-36}; since $\sigma$ itself is
such a limit, the physical labels converge, and their limit
is the same for every $\sigma'$.

By \cref{IV-35-36} there is a canonical decomposition whose
every component is a limit in this frame along its full sequence.
Uniqueness, \cref{IV-36-9}, identifies the component functionals
of any canonical decomposition with those of this one, at every
point of the compact component space. Thus the preceding
application of \cref{IV-69-36} proves \eqref{IV-69-40}, the
recovery formula and positivity for every component.
Each is reflection positive at the coordinate mirror by
\cref{IV-36-12}. The smooth reflection pair and
\cref{IV-69-33} give the dimension assertion.

For a volume-first limit, the centered moments of its defining
states converge to $\sigma_S$: in \eqref{IV-9-14} each moment
is a finite sum of products of convergent cumulants. Its defining
states therefore give the stated convergent frame. The label
identified above equals $b(S)$ by \cref{IV-69-29}.
\end{proof}

\appendix
\section{Relation to earlier constructions}\label{sec:literature-comparison}

The observable in this paper is the centered, rescaled scalar plaquette
energy. Its continuum correlations are defined on the diagonal-flat test
spaces of \cref{thm:scalar-os}. To compare this construction with earlier
work, one must distinguish the observable, the estimates used to remove
the cutoffs, and the symmetry and reconstruction conclusions. These
distinctions also separate the fixed-torus convergence of
\cref{thm:fixed-torus-continuum}, the subsequential infinite-volume
construction of \cref{thm:gibbs-continuum}, and the two-point asymptotic of
\cref{thm:short-distance-law}.

\subsection{Renormalization and ultraviolet stability}

Ba{\l}aban's work supplies the lattice gauge-theoretic renormalization
framework used here: block averaging, regular gauge configurations,
background fields and propagators
\cite{B-CMP98,B-CMP99a,B-CMP99b}, together with effective-action
expansions \cite{B-CMP109,B-CMP119}. The construction and use of these
objects occupy a substantial part of the present proof. Their role is to
retain locality through successive integrations and to control the
dependence of the resulting effective interactions on the background
field.

The ultraviolet stability theorem for three-dimensional pure gauge
theory \cite[Theorem~1 and equations~(5)--(6)]{B-CMP102a} gives upper and
lower bounds on the renormalized densities and, after stopping at a
fixed coarse scale, cutoff-uniform bounds on the normalized partition
function. The normalization, including the constant terms in the
action, is part of that statement. These bounds control the size of the
functional integral; they do not by themselves assert convergence of
the correlations of a specified continuum observable.

Here that further task is carried out with local complex sources for
the plaquette energy. Differentiation in several separated sources
selects terms whose regions of dependence connect their supports. The
locality estimates then make the fine-scale contributions summable.
Comparison of adjacent cutoffs gives convergence on a fixed torus.
Contact subtractions enter this comparison, while the limiting
correlation distributions are formulated on diagonal-flat test spaces.
Thus the source estimates and cutoff
comparison in
\cref{sec:local-source-energies,sec:iv-cutoff-summation,sec:iv-all-label-cutoff-limit}
are essential to \cref{thm:fixed-torus-continuum}. Ultraviolet stability
alone would not replace them.

\subsection{Abelian Higgs models and QED}

King's construction of the Abelian Higgs model in dimensions two and
three also distinguishes stability from convergence. The finite-volume
continuum limit is obtained by comparing effective actions and making
source-dependent generating functionals Cauchy
\cite[pp.~649, 655--657]{King1986I}. The infinite-volume paper proves
continuum and volume limits with its specified boundary conditions, and
establishes translation invariance, reflection positivity and entire
dependence on finitely many source parameters
\cite[Theorems~2.1--2.3]{King1986II}. The paragraph following its
Theorem~2.3 explicitly leaves ergodicity, rotation invariance and
regularity as further requirements for the full Osterwalder--Schrader
conclusion. Its comparison of rotated lattices is a separate step;
the ensuing passage to full rotations uses an expected
shape-independent periodic infinite-volume limit
\cite[Theorem~2.4 and equation~(2.24)]{King1986II}. The matter field,
photon mass and boundary conditions in these results are part of their
hypotheses; the present nonabelian pure-gauge construction is not an
application of those theorems.

The fixed-positive-photon-mass scalar-electrodynamics preprint
\cite[Theorem~1.1(A)]{ScalarQED-Main} concerns the Abelian curvature as a
two-form. Its Gaussian bound permits a distribution-valued law with
moments tested on the full Schwartz space. The rotation supplement
\cite[Theorem~A]{ScalarQED-Rotations} compares two lattice orientations
through tori with mass walls. In the present paper, rotations are
obtained instead from coordinate and diagonal reflection positivity,
a spectral restriction, and analytic continuation; see
\cref{sec:gibbs-continuum-rotation}. The scalar-electrodynamics result
keeps the photon mass positive and leaves nontriviality open
\cite[\S\S~1.4 and~23.3]{ScalarQED-Main}.

The fermionic QED3 construction in
\cite[Theorem~A]{QED3FieldStrength} likewise concerns an Abelian
field-strength tensor, under its weak-charge and positive-reference-mass
assumptions for the stated chamber-fermion lattice model. Its Gaussian
estimate controls all polynomial moments on the full Schwartz space;
its non-Gaussianity conclusion uses a separate fourth-cumulant estimate.
Neither estimate transfers to the scalar Yang--Mills observable by a
change of notation. The latter has a quadratic classical counterpart,
and the present theorem leaves contact terms on the coincidence
diagonals unspecified. The two constructions therefore require
different estimates even before their different matter and gauge
content is taken into account.

\subsection{Four-dimensional Yang--Mills comparisons}

The conditional four-dimensional $SU(2)$ construction
\cite[Theorems~1.1 and~1.3]{ScalarQED-YM4Conditional} also uses scalar
plaquette densities and separated correlations. It assumes a coupling
law and uniform separated-point bounds. An additional hypothesis
comparing partition functions on cubic and twisted tori supplies diagonal
reflection positivity for selected volume-first limits. Its
all-moment rotation conclusion additionally assumes holomorphic
continuation, growth bounds and positivity of the continued Gram
kernels. These hypotheses isolate several tasks also present here,
but do not establish the corresponding four-dimensional source
estimates or analytic continuation. The clustering conclusion is
toward the projection onto invariant vectors; the rank of that
projection is not proved to be one.

A different four-dimensional construction treats bounded
gauge-invariant block observables at a fixed physical block scale
\cite[Theorem~1.1 and \S~1.6]{ScalarQED-YM4Blocks}. Boundedness gives
joint subsequential laws for arbitrary coupling sequences. The
conclusions include hyperoctahedral symmetry, coordinate reflection
positivity for observables whose entire regions of dependence lie in
the relevant half-space, and a positive Hamiltonian. Removing the
physical block scale, obtaining full rotations, proving nontriviality
and establishing vacuum uniqueness remain separate problems there.

The fixed block size is consequential: these are bounded extended
observables, whereas the normalization of our plaquette-energy field
depends on the cutoff. Compactness of the link space does not control
that normalization. In the present three-dimensional construction,
\cref{thm:fixed-torus-continuum,thm:gibbs-continuum,thm:short-distance-law}
provide, respectively, the continuum correlations, the component
structure and the short-distance information. They make no
four-dimensional assertion.

\subsection{Reconstruction and nontriviality}

The passage from suitable Euclidean correlations to a relativistic
quantum field is the Osterwalder--Schrader reconstruction theorem
\cite{OS1973,OS1975}. Diagonal-flat test functions and the associated
freedom in contact terms already occur in that framework
\cite[\S\S~2--3]{OS1973}. Reconstruction requires the appropriate
regularity and growth bounds, covariance and positivity; reflection
positivity by itself is insufficient. Clustering supplies the
additional unique-vacuum conclusion \cite[\S~IV.1]{OS1975}.

The reconstruction sections of this paper verify these requirements
for the scalar correlations constructed here. Their substantive
inputs include polynomial energy bounds, tempered boundary values,
Euclidean covariance and clustering of the canonical components.
The unique-vacuum conclusion applies to these components; it does
not assert that an arbitrary invariant mixture has a unique vacuum
or that the infinite-volume limit is unique.

Finally, under the additional small-reference-coupling hypothesis of
\cref{thm:short-distance-law}, the nonzero two-point asymptotic proves
nontriviality and recovers the physical inverse coupling. This is
stronger than the existence of a possibly trivial reconstructed
space, but it does not prove non-Gaussianity or a mass gap. Nor does
the scalar theorem construct a tensor-valued nonabelian field strength.
These distinctions specify the scope of the present result independently
of the general reconstruction theorem.

\printbibliography
\end{document}
